% tikz_preview.tex - render the five paper figures from ../figures/ inside a minimal
% IEEE-class two-column document, at the size and font they have in the paper, to check
% that each one fits its column and comes out at the height it is supposed to.
%
%   compile:  lualatex --interaction=nonstopmode tikz_preview.tex
%
% lualatex, not pdflatex: the sweep figures run to ~100k lines with 9 to 12 axes and
% exceed pdflatex's main_memory.
%
% The .tikz are post processed (tikz_make_adjustable.m) so fonts, the legend box and
% the size follow the knobs in paperfig_preamble.tex. Nothing is overridden here: the
% figures render at the preamble defaults, which is what the paper uses, so this
% preview is a faithful check rather than a differently-tuned one.
\documentclass[journal]{IEEEtran}
\usepackage{graphicx}
\usepackage{pgfplots}
\pgfplotsset{compat=newest}
% The paper sets maths to 80 % of the body size for its page budget. Keep it, or every
% axis label and tick comes out larger here than it does in the paper.
\DeclareMathSizes{10}{8}{6}{5}

\newcommand{\figdir}{../figures}
% The sweep figures embed their colour-bar rasters by bare filename.
\graphicspath{{\figdir/}}
% adjustable font / legend / size knobs (\figscale, \figfont, \figlegendlw, ...) and
% the \fcoi notation the axis labels use:
% paperfig_preamble.tex
% Shared, adjustable styling knobs for the .tikz figures in this folder (exported by
% paper_export via matlab2tikz and post processed by tikz_make_adjustable.m).
% \input this in the document preamble, AFTER \usepackage{pgfplots}:
%
%   \usepackage{pgfplots}\pgfplotsset{compat=newest}
%   % paperfig_preamble.tex
% Shared, adjustable styling knobs for the .tikz figures in this folder (exported by
% paper_export via matlab2tikz and post processed by tikz_make_adjustable.m).
% \input this in the document preamble, AFTER \usepackage{pgfplots}:
%
%   \usepackage{pgfplots}\pgfplotsset{compat=newest}
%   % paperfig_preamble.tex
% Shared, adjustable styling knobs for the .tikz figures in this folder (exported by
% paper_export via matlab2tikz and post processed by tikz_make_adjustable.m).
% \input this in the document preamble, AFTER \usepackage{pgfplots}:
%
%   \usepackage{pgfplots}\pgfplotsset{compat=newest}
%   \input{paperfig_preamble.tex}
%
% To restyle, redefine any knob below in your preamble (after the \input), or right
% before a single \input{...tikz} to restyle just that figure. Every figure picks the
% change up at compile time; nothing has to be regenerated in MATLAB.

% ======== size: one multiplier on every baked axis dimension (width, height, and the
% multi panel offsets), so the whole figure grows or shrinks while the data stays put.
% The figures are exported at their final paper size, so this is 1 in the paper.
\providecommand{\figscale}{1}

% ======== fonts. \figfont is the base for the whole picture; the rest break out the
% individual elements and default to \figfont, so you can size, say, the legend text
% independently of the tick labels.
\providecommand{\figfont}{\footnotesize}      % base font for the whole figure
\providecommand{\figticfont}{\figfont}        % tick labels
\providecommand{\figlabelfont}{\figfont}      % axis labels (the x and y labels)
\providecommand{\figtitlefont}{\figfont}      % axis titles (the WECC / SC500 headers)
\providecommand{\figlegendfont}{\figfont}     % legend text
\providecommand{\figcornerfont}{\figfont\bfseries}  % bold corner tags

% ======== title gap: pull axis titles (the WECC / SC500 / Texas2000 column headers)
% down toward the plots. matlab2tikz anchors the title just above the axis with a
% default gap that looks too large once the figure is small; a small negative yshift
% closes it. Override per figure before an \input if a figure needs a different gap.
\providecommand{\figtitleshift}{-4pt}

% ======== zoom inset tick shift (Fig. 7). Its three insets sit on top of the settled
% COI trace of the panel behind them, and unshifted x tick numbers were cut by it.
% 1.5pt clears the trace by about 1pt; it is near the ceiling, since only about 3.1pt
% separate the digit tops from the inset's bottom border.
\providecommand{\figinsettickshift}{1.5pt}

% ======== legend box (applies to the real pgfplots legends and the hand drawn per
% node legend alike).
\providecommand{\figlegendlw}{0.6pt}          % legend box stroke width
\providecommand{\figlegenddraw}{black}        % legend box stroke colour
\providecommand{\figlegendfill}{white}        % legend box fill colour

% ======== notation used inside the axis labels. The paper defines these itself; the
% \providecommand makes the figures \input-able in a bare document as well.
\providecommand{\coi}{\mathrm{COI}}
\providecommand{\fcoi}{f_{\coi}}
\providecommand{\wcoi}{\omega_{\coi}}

\pgfplotsset{
  % plain fixed-decimal tick labels (0.05), never sci form (5e-2) or a 10^-2 common factor
  scaled ticks=false,
  every tick label/.append style={/pgf/number format/.cd, fixed, precision=4, /tikz/.cd},
  every axis label/.append style={font=\figlabelfont},
  every axis title/.append style={font=\figtitlefont, yshift=\figtitleshift},
  every axis legend/.append style={font=\figlegendfont, inner xsep=2pt, inner ysep=1pt, nodes={inner sep=1pt}},
}

%
% To restyle, redefine any knob below in your preamble (after the \input), or right
% before a single \input{...tikz} to restyle just that figure. Every figure picks the
% change up at compile time; nothing has to be regenerated in MATLAB.

% ======== size: one multiplier on every baked axis dimension (width, height, and the
% multi panel offsets), so the whole figure grows or shrinks while the data stays put.
% The figures are exported at their final paper size, so this is 1 in the paper.
\providecommand{\figscale}{1}

% ======== fonts. \figfont is the base for the whole picture; the rest break out the
% individual elements and default to \figfont, so you can size, say, the legend text
% independently of the tick labels.
\providecommand{\figfont}{\footnotesize}      % base font for the whole figure
\providecommand{\figticfont}{\figfont}        % tick labels
\providecommand{\figlabelfont}{\figfont}      % axis labels (the x and y labels)
\providecommand{\figtitlefont}{\figfont}      % axis titles (the WECC / SC500 headers)
\providecommand{\figlegendfont}{\figfont}     % legend text
\providecommand{\figcornerfont}{\figfont\bfseries}  % bold corner tags

% ======== title gap: pull axis titles (the WECC / SC500 / Texas2000 column headers)
% down toward the plots. matlab2tikz anchors the title just above the axis with a
% default gap that looks too large once the figure is small; a small negative yshift
% closes it. Override per figure before an \input if a figure needs a different gap.
\providecommand{\figtitleshift}{-4pt}

% ======== zoom inset tick shift (Fig. 7). Its three insets sit on top of the settled
% COI trace of the panel behind them, and unshifted x tick numbers were cut by it.
% 1.5pt clears the trace by about 1pt; it is near the ceiling, since only about 3.1pt
% separate the digit tops from the inset's bottom border.
\providecommand{\figinsettickshift}{1.5pt}

% ======== legend box (applies to the real pgfplots legends and the hand drawn per
% node legend alike).
\providecommand{\figlegendlw}{0.6pt}          % legend box stroke width
\providecommand{\figlegenddraw}{black}        % legend box stroke colour
\providecommand{\figlegendfill}{white}        % legend box fill colour

% ======== notation used inside the axis labels. The paper defines these itself; the
% \providecommand makes the figures \input-able in a bare document as well.
\providecommand{\coi}{\mathrm{COI}}
\providecommand{\fcoi}{f_{\coi}}
\providecommand{\wcoi}{\omega_{\coi}}

\pgfplotsset{
  % plain fixed-decimal tick labels (0.05), never sci form (5e-2) or a 10^-2 common factor
  scaled ticks=false,
  every tick label/.append style={/pgf/number format/.cd, fixed, precision=4, /tikz/.cd},
  every axis label/.append style={font=\figlabelfont},
  every axis title/.append style={font=\figtitlefont, yshift=\figtitleshift},
  every axis legend/.append style={font=\figlegendfont, inner xsep=2pt, inner ysep=1pt, nodes={inner sep=1pt}},
}

%
% To restyle, redefine any knob below in your preamble (after the \input), or right
% before a single \input{...tikz} to restyle just that figure. Every figure picks the
% change up at compile time; nothing has to be regenerated in MATLAB.

% ======== size: one multiplier on every baked axis dimension (width, height, and the
% multi panel offsets), so the whole figure grows or shrinks while the data stays put.
% The figures are exported at their final paper size, so this is 1 in the paper.
\providecommand{\figscale}{1}

% ======== fonts. \figfont is the base for the whole picture; the rest break out the
% individual elements and default to \figfont, so you can size, say, the legend text
% independently of the tick labels.
\providecommand{\figfont}{\footnotesize}      % base font for the whole figure
\providecommand{\figticfont}{\figfont}        % tick labels
\providecommand{\figlabelfont}{\figfont}      % axis labels (the x and y labels)
\providecommand{\figtitlefont}{\figfont}      % axis titles (the WECC / SC500 headers)
\providecommand{\figlegendfont}{\figfont}     % legend text
\providecommand{\figcornerfont}{\figfont\bfseries}  % bold corner tags

% ======== title gap: pull axis titles (the WECC / SC500 / Texas2000 column headers)
% down toward the plots. matlab2tikz anchors the title just above the axis with a
% default gap that looks too large once the figure is small; a small negative yshift
% closes it. Override per figure before an \input if a figure needs a different gap.
\providecommand{\figtitleshift}{-4pt}

% ======== zoom inset tick shift (Fig. 7). Its three insets sit on top of the settled
% COI trace of the panel behind them, and unshifted x tick numbers were cut by it.
% 1.5pt clears the trace by about 1pt; it is near the ceiling, since only about 3.1pt
% separate the digit tops from the inset's bottom border.
\providecommand{\figinsettickshift}{1.5pt}

% ======== legend box (applies to the real pgfplots legends and the hand drawn per
% node legend alike).
\providecommand{\figlegendlw}{0.6pt}          % legend box stroke width
\providecommand{\figlegenddraw}{black}        % legend box stroke colour
\providecommand{\figlegendfill}{white}        % legend box fill colour

% ======== notation used inside the axis labels. The paper defines these itself; the
% \providecommand makes the figures \input-able in a bare document as well.
\providecommand{\coi}{\mathrm{COI}}
\providecommand{\fcoi}{f_{\coi}}
\providecommand{\wcoi}{\omega_{\coi}}

\pgfplotsset{
  % plain fixed-decimal tick labels (0.05), never sci form (5e-2) or a 10^-2 common factor
  scaled ticks=false,
  every tick label/.append style={/pgf/number format/.cd, fixed, precision=4, /tikz/.cd},
  every axis label/.append style={font=\figlabelfont},
  every axis title/.append style={font=\figtitlefont, yshift=\figtitleshift},
  every axis legend/.append style={font=\figlegendfont, inner xsep=2pt, inner ysep=1pt, nodes={inner sep=1pt}},
}


\begin{document}
\title{Paper figures --- PGF/TikZ render check}
\author{auto-generated preview of \texttt{make\_paper\_figures} output}
\maketitle

\begin{abstract}
Each figure is the \texttt{.tikz} written by \texttt{paper\_export}, placed at its
paper size and font. Figures 4, 5(a) and 5(b) are single-column; Figure 7 spans the
full text width, as it does in the paper.
\end{abstract}

\section{Twelve-bus example ($K = 20\,I$)}

\begin{figure}[!t]\centering
  % This file was created by matlab2tikz.
%
\definecolor{mycolor1}{rgb}{0.70000,0.70000,0.75000}%
\definecolor{mycolor2}{rgb}{0.57647,0.00000,0.00000}%
\definecolor{mycolor3}{rgb}{0.00000,0.32157,0.57647}%
%
\begin{tikzpicture}[font=\figfont]

\begin{axis}[%
width={\figscale*1.171in},
height={\figscale*0.788in},
at={(\figscale*0.455in,\figscale*0.31in)},
scale only axis,
separate axis lines,
every outer x axis line/.append style={black},
every x tick label/.append style={font=\figticfont},
every x tick/.append style={black},
xmin=-150,
xmax=-6,
xtick={-150,   -6},
xlabel={$\mathrm{Re}(\lambda)$ in rad/s},
every outer y axis line/.append style={black},
every y tick label/.append style={font=\figticfont},
every y tick/.append style={black},
ymin=-40,
ymax=40,
ytick={-40,   0,  40},
ylabel={$\mathrm{Im}(\lambda)$ in rad/s},
axis background/.style={fill=white}
]
\addplot [color=white!85!black, line width=0.6pt, forget plot]
  table[row sep=crcr]{%
-140	-40\\
-140	40\\
};
\addplot [color=white!85!black, line width=0.6pt, forget plot]
  table[row sep=crcr]{%
-130	-40\\
-130	40\\
};
\addplot [color=white!85!black, line width=0.6pt, forget plot]
  table[row sep=crcr]{%
-120	-40\\
-120	40\\
};
\addplot [color=white!85!black, line width=0.6pt, forget plot]
  table[row sep=crcr]{%
-110	-40\\
-110	40\\
};
\addplot [color=white!85!black, line width=0.6pt, forget plot]
  table[row sep=crcr]{%
-100	-40\\
-100	40\\
};
\addplot [color=white!85!black, line width=0.6pt, forget plot]
  table[row sep=crcr]{%
-90	-40\\
-90	40\\
};
\addplot [color=white!85!black, line width=0.6pt, forget plot]
  table[row sep=crcr]{%
-80	-40\\
-80	40\\
};
\addplot [color=white!85!black, line width=0.6pt, forget plot]
  table[row sep=crcr]{%
-70	-40\\
-70	40\\
};
\addplot [color=white!85!black, line width=0.6pt, forget plot]
  table[row sep=crcr]{%
-60	-40\\
-60	40\\
};
\addplot [color=white!85!black, line width=0.6pt, forget plot]
  table[row sep=crcr]{%
-50	-40\\
-50	40\\
};
\addplot [color=white!85!black, line width=0.6pt, forget plot]
  table[row sep=crcr]{%
-40	-40\\
-40	40\\
};
\addplot [color=white!85!black, line width=0.6pt, forget plot]
  table[row sep=crcr]{%
-30	-40\\
-30	40\\
};
\addplot [color=white!85!black, line width=0.6pt, forget plot]
  table[row sep=crcr]{%
-20	-40\\
-20	40\\
};
\addplot [color=white!85!black, line width=0.6pt, forget plot]
  table[row sep=crcr]{%
-10	-40\\
-10	40\\
};
\addplot [color=white!85!black, line width=0.6pt, forget plot]
  table[row sep=crcr]{%
-150	-20\\
-6	-20\\
};
\addplot [color=white!85!black, line width=0.6pt, forget plot]
  table[row sep=crcr]{%
-150	0\\
-6	0\\
};
\addplot [color=white!85!black, line width=0.6pt, forget plot]
  table[row sep=crcr]{%
-150	20\\
-6	20\\
};
\addplot [color=mycolor1, dashed, line width=0.6pt, forget plot]
  table[row sep=crcr]{%
-150	840\\
-146.939	840\\
-143.878	840\\
-140.816	840\\
-137.755	840\\
-134.694	840\\
-131.633	840\\
-128.571	840\\
-125.51	840\\
-122.449	840\\
-119.388	840\\
-116.327	840\\
-113.265	840\\
-110.204	840\\
-107.143	840\\
-104.082	840\\
-101.02	840\\
-97.9592	840\\
-94.898	840\\
-91.8367	840\\
-88.7755	840\\
-85.7143	840\\
-82.6531	822.388\\
-79.5918	791.929\\
-76.5306	761.47\\
-73.4694	731.011\\
-70.4082	700.552\\
-67.3469	670.094\\
-64.2857	639.635\\
-61.2245	609.176\\
-58.1633	578.717\\
-55.102	548.258\\
-52.0408	517.8\\
-48.9796	487.341\\
-45.9184	456.882\\
-42.8571	426.423\\
-39.7959	395.964\\
-36.7347	365.506\\
-33.6735	335.047\\
-30.6122	304.588\\
-27.551	274.129\\
-24.4898	243.67\\
-21.4286	213.212\\
-18.3673	182.753\\
-15.3061	152.294\\
-12.2449	121.835\\
-9.18367	91.3764\\
-6.12245	60.9176\\
};
\addplot [color=mycolor1, dashed, line width=0.6pt, forget plot]
  table[row sep=crcr]{%
-150	-840\\
-146.939	-840\\
-143.878	-840\\
-140.816	-840\\
-137.755	-840\\
-134.694	-840\\
-131.633	-840\\
-128.571	-840\\
-125.51	-840\\
-122.449	-840\\
-119.388	-840\\
-116.327	-840\\
-113.265	-840\\
-110.204	-840\\
-107.143	-840\\
-104.082	-840\\
-101.02	-840\\
-97.9592	-840\\
-94.898	-840\\
-91.8367	-840\\
-88.7755	-840\\
-85.7143	-840\\
-82.6531	-822.388\\
-79.5918	-791.929\\
-76.5306	-761.47\\
-73.4694	-731.011\\
-70.4082	-700.552\\
-67.3469	-670.094\\
-64.2857	-639.635\\
-61.2245	-609.176\\
-58.1633	-578.717\\
-55.102	-548.258\\
-52.0408	-517.8\\
-48.9796	-487.341\\
-45.9184	-456.882\\
-42.8571	-426.423\\
-39.7959	-395.964\\
-36.7347	-365.506\\
-33.6735	-335.047\\
-30.6122	-304.588\\
-27.551	-274.129\\
-24.4898	-243.67\\
-21.4286	-213.212\\
-18.3673	-182.753\\
-15.3061	-152.294\\
-12.2449	-121.835\\
-9.18367	-91.3764\\
-6.12245	-60.9176\\
};
\addplot [color=mycolor2, line width=0.9pt, only marks, mark size=2.3pt, mark=x, mark options={solid, mycolor2}, forget plot]
  table[row sep=crcr]{%
-138.896	0\\
-111.113	0\\
-44.789	0\\
-9.18391	0\\
};
\addplot [color=mycolor3, line width=0.4pt, only marks, mark size=1.4pt, mark=*, mark options={solid, fill=mycolor3, mycolor3}, forget plot]
  table[row sep=crcr]{%
-78.5657	0\\
-64.8779	0\\
-6.72803	31.0734\\
-6.72803	-31.0734\\
-9.3304	16.538\\
-9.3304	-16.538\\
-20.2816	0\\
-6.15461	7.66302\\
-6.15461	-7.66302\\
-10.1832	0\\
};
\end{axis}

\begin{axis}[%
width={\figscale*1.171in},
height={\figscale*0.788in},
at={(\figscale*1.93in,\figscale*0.31in)},
scale only axis,
separate axis lines,
every outer x axis line/.append style={black},
every x tick label/.append style={font=\figticfont},
every x tick/.append style={black},
xmin=-6,
xmax=1,
xtick={-5,  0},
xlabel={$\mathrm{Re}(\lambda)$ in rad/s},
every outer y axis line/.append style={black},
every y tick label/.append style={font=\figticfont},
every y tick/.append style={black},
ymin=-40,
ymax=40,
ytick={-40,   0,  40},
axis background/.style={fill=white},
legend style={legend cell align=left, align=left, draw=\figlegenddraw, fill=\figlegendfill, line width=\figlegendlw}
]
\addplot [color=white!85!black, line width=0.6pt, forget plot]
  table[row sep=crcr]{%
-6	-20\\
1	-20\\
};
\addplot [color=white!85!black, line width=0.6pt, forget plot]
  table[row sep=crcr]{%
-6	0\\
1	0\\
};
\addplot [color=white!85!black, line width=0.6pt, forget plot]
  table[row sep=crcr]{%
-6	20\\
1	20\\
};
\addplot [color=mycolor1, dashed, line width=0.6pt, forget plot]
  table[row sep=crcr]{%
-6	59.6992\\
-5.87755	58.4809\\
-5.7551	57.2625\\
-5.63265	56.0442\\
-5.5102	54.8258\\
-5.38776	53.6075\\
-5.26531	52.3891\\
-5.14286	51.1708\\
-5.02041	49.9524\\
-4.89796	48.7341\\
-4.77551	47.5157\\
-4.65306	46.2974\\
-4.53061	45.079\\
-4.40816	43.8607\\
-4.28571	42.6423\\
-4.16327	41.424\\
-4.04082	40.2056\\
-3.91837	38.9873\\
-3.79592	37.7689\\
-3.67347	36.5506\\
-3.55102	35.3322\\
-3.42857	34.1139\\
-3.30612	32.8955\\
-3.18367	31.6772\\
-3.06122	30.4588\\
-2.93878	29.2404\\
-2.81633	28.0221\\
-2.69388	26.8037\\
-2.57143	25.5854\\
-2.44898	24.367\\
-2.32653	23.1487\\
-2.20408	21.9303\\
-2.08163	20.712\\
-1.95918	19.4936\\
-1.83673	18.2753\\
-1.71429	17.0569\\
-1.59184	15.8386\\
-1.46939	14.6202\\
-1.34694	13.4019\\
-1.22449	12.1835\\
-1.10204	10.9652\\
-0.979592	9.74682\\
-0.857143	8.52846\\
-0.734694	7.31011\\
-0.612245	6.09176\\
-0.489796	4.87341\\
-0.367347	3.65506\\
-0.244898	2.4367\\
-0.122449	1.21835\\
0	-0\\
};
\addplot [color=mycolor1, dashed, line width=0.6pt, forget plot]
  table[row sep=crcr]{%
-6	-59.6992\\
-5.87755	-58.4809\\
-5.7551	-57.2625\\
-5.63265	-56.0442\\
-5.5102	-54.8258\\
-5.38776	-53.6075\\
-5.26531	-52.3891\\
-5.14286	-51.1708\\
-5.02041	-49.9524\\
-4.89796	-48.7341\\
-4.77551	-47.5157\\
-4.65306	-46.2974\\
-4.53061	-45.079\\
-4.40816	-43.8607\\
-4.28571	-42.6423\\
-4.16327	-41.424\\
-4.04082	-40.2056\\
-3.91837	-38.9873\\
-3.79592	-37.7689\\
-3.67347	-36.5506\\
-3.55102	-35.3322\\
-3.42857	-34.1139\\
-3.30612	-32.8955\\
-3.18367	-31.6772\\
-3.06122	-30.4588\\
-2.93878	-29.2404\\
-2.81633	-28.0221\\
-2.69388	-26.8037\\
-2.57143	-25.5854\\
-2.44898	-24.367\\
-2.32653	-23.1487\\
-2.20408	-21.9303\\
-2.08163	-20.712\\
-1.95918	-19.4936\\
-1.83673	-18.2753\\
-1.71429	-17.0569\\
-1.59184	-15.8386\\
-1.46939	-14.6202\\
-1.34694	-13.4019\\
-1.22449	-12.1835\\
-1.10204	-10.9652\\
-0.979592	-9.74682\\
-0.857143	-8.52846\\
-0.734694	-7.31011\\
-0.612245	-6.09176\\
-0.489796	-4.87341\\
-0.367347	-3.65506\\
-0.244898	-2.4367\\
-0.122449	-1.21835\\
0	0\\
};
\addplot [color=mycolor1, dashed, line width=0.6pt, forget plot]
  table[row sep=crcr]{%
-3	-40\\
-3	40\\
};
\addplot [color=mycolor2, line width=0.9pt, only marks, mark size=2.3pt, mark=x, mark options={solid, mycolor2}]
  table[row sep=crcr]{%
-3.58618	31.5142\\
-3.58618	-31.5142\\
-3.00045	29.5084\\
-3.00045	-29.5084\\
-3.2766	18.506\\
-3.2766	-18.506\\
-5.09251	16.1516\\
-5.09251	-16.1516\\
-4.62745	8.55739\\
-4.62745	-8.55739\\
};
\addlegendentry{$K = 0$}

\addplot [color=mycolor3, line width=0.4pt, only marks, mark size=1.4pt, mark=*, mark options={solid, fill=mycolor3, mycolor3}]
  table[row sep=crcr]{%
-4.86459	29.2594\\
-4.86459	-29.2594\\
-5.07133	17.8782\\
-5.07133	-17.8782\\
};
\addlegendentry{$K = 20\,I$}

\end{axis}
\end{tikzpicture}%
  \caption{Fig.~4 --- small-signal spectrum, $K=0$ against $K=20\,I$.}
\end{figure}

\begin{figure}[!t]\centering
  % This file was created by matlab2tikz.
%
\definecolor{mycolor1}{rgb}{0.87843,0.67059,0.67059}%
\definecolor{mycolor2}{rgb}{0.67989,0.24523,0.24523}%
\definecolor{mycolor3}{rgb}{0.63656,0.15679,0.15679}%
\definecolor{mycolor4}{rgb}{0.81501,0.55362,0.55362}%
\definecolor{mycolor5}{rgb}{0.48525,0.00988,0.00988}%
\definecolor{mycolor6}{rgb}{0.51695,0.01229,0.01229}%
\definecolor{mycolor7}{rgb}{0.59104,0.07546,0.07546}%
\definecolor{mycolor8}{rgb}{0.72393,0.34173,0.34173}%
\definecolor{mycolor9}{rgb}{0.77270,0.45370,0.45370}%
\definecolor{mycolor10}{rgb}{0.55454,0.02805,0.02805}%
\definecolor{mycolor11}{rgb}{0.93333,0.98824,0.92157}%
\definecolor{mycolor12}{rgb}{0.33216,0.67989,0.24523}%
\definecolor{mycolor13}{rgb}{0.25274,0.63656,0.15679}%
\definecolor{mycolor14}{rgb}{0.60590,0.81501,0.55362}%
\definecolor{mycolor15}{rgb}{0.10495,0.48525,0.00988}%
\definecolor{mycolor16}{rgb}{0.11322,0.51695,0.01229}%
\definecolor{mycolor17}{rgb}{0.17857,0.59104,0.07546}%
\definecolor{mycolor18}{rgb}{0.41817,0.72393,0.34173}%
\definecolor{mycolor19}{rgb}{0.51750,0.77270,0.45370}%
\definecolor{mycolor20}{rgb}{0.13335,0.55454,0.02805}%
%
\begin{tikzpicture}[font=\figfont]

\begin{axis}[%
width={\figscale*1.32in},
height={\figscale*0.95in},
at={(\figscale*0.455in,\figscale*0.31in)},
scale only axis,
separate axis lines,
every outer x axis line/.append style={black},
every x tick label/.append style={font=\figticfont},
every x tick/.append style={black},
xmin=-0.125,
xmax=2,
xlabel={Time since step in s},
every outer y axis line/.append style={black},
every y tick label/.append style={font=\figticfont},
every y tick/.append style={black},
ymin=-0.2,
ymax=0.1,
ylabel={$f - f_0$ in Hz},
axis background/.style={fill=white},
xmajorgrids,
ymajorgrids,
grid style={white!55!black, opacity=0.3}
]

\addplot[area legend, draw=none, fill=mycolor1, fill opacity=0.9, forget plot]
table[row sep=crcr] {%
x	y\\
0	-0.01256\\
0.00501253	-0.0403283\\
0.0100251	-0.0463046\\
0.0150376	-0.063637\\
0.0200501	-0.0827103\\
0.0250627	-0.10035\\
0.0300752	-0.11652\\
0.0350877	-0.131007\\
0.0401003	-0.143844\\
0.0451128	-0.154867\\
0.0501253	-0.164162\\
0.0551378	-0.17162\\
0.0601504	-0.177368\\
0.0651629	-0.181357\\
0.0701754	-0.182839\\
0.075188	-0.184212\\
0.0802005	-0.185483\\
0.085213	-0.186652\\
0.0902256	-0.187726\\
0.0952381	-0.188703\\
0.100251	-0.189591\\
0.105263	-0.190388\\
0.110276	-0.1911\\
0.115288	-0.191728\\
0.120301	-0.192277\\
0.125313	-0.192747\\
0.130326	-0.193142\\
0.135338	-0.193464\\
0.140351	-0.193717\\
0.145363	-0.193901\\
0.150376	-0.194021\\
0.155388	-0.194077\\
0.160401	-0.194074\\
0.165414	-0.194014\\
0.170426	-0.193898\\
0.175439	-0.19373\\
0.180451	-0.193513\\
0.185464	-0.193248\\
0.190476	-0.192938\\
0.195489	-0.192586\\
0.200501	-0.192194\\
0.205514	-0.191765\\
0.210526	-0.191302\\
0.215539	-0.190806\\
0.220551	-0.190279\\
0.225564	-0.189726\\
0.230576	-0.189146\\
0.235589	-0.188544\\
0.240602	-0.18792\\
0.245614	-0.187277\\
0.250627	-0.186617\\
0.255639	-0.185942\\
0.260652	-0.185253\\
0.265664	-0.184537\\
0.270677	-0.18378\\
0.275689	-0.182988\\
0.280702	-0.182166\\
0.285714	-0.181319\\
0.290727	-0.18045\\
0.295739	-0.179564\\
0.300752	-0.178666\\
0.305764	-0.177759\\
0.310777	-0.176848\\
0.315789	-0.175936\\
0.320802	-0.175028\\
0.325815	-0.174127\\
0.330827	-0.173236\\
0.33584	-0.17236\\
0.340852	-0.171501\\
0.345865	-0.170663\\
0.350877	-0.169848\\
0.35589	-0.169062\\
0.360902	-0.168306\\
0.365915	-0.167584\\
0.370927	-0.166898\\
0.37594	-0.166252\\
0.380952	-0.165649\\
0.385965	-0.165092\\
0.390977	-0.164583\\
0.39599	-0.164127\\
0.401003	-0.163725\\
0.406015	-0.163382\\
0.411028	-0.163087\\
0.41604	-0.162826\\
0.421053	-0.162592\\
0.426065	-0.162379\\
0.431078	-0.162179\\
0.43609	-0.161987\\
0.441103	-0.161795\\
0.446115	-0.161598\\
0.451128	-0.161388\\
0.45614	-0.161161\\
0.461153	-0.160921\\
0.466165	-0.160681\\
0.471178	-0.160438\\
0.47619	-0.160193\\
0.481203	-0.159945\\
0.486216	-0.159693\\
0.491228	-0.159436\\
0.496241	-0.159175\\
0.501253	-0.158908\\
0.506266	-0.158636\\
0.511278	-0.158357\\
0.516291	-0.158072\\
0.521303	-0.157779\\
0.526316	-0.157479\\
0.531328	-0.157171\\
0.536341	-0.156855\\
0.541353	-0.156531\\
0.546366	-0.156198\\
0.551378	-0.155855\\
0.556391	-0.155502\\
0.561404	-0.15514\\
0.566416	-0.154767\\
0.571429	-0.154383\\
0.576441	-0.153988\\
0.581454	-0.153581\\
0.586466	-0.153161\\
0.591479	-0.152729\\
0.596491	-0.152283\\
0.601504	-0.151804\\
0.606516	-0.151276\\
0.611529	-0.150706\\
0.616541	-0.150097\\
0.621554	-0.149454\\
0.626566	-0.148783\\
0.631579	-0.148089\\
0.636591	-0.147376\\
0.641604	-0.146649\\
0.646617	-0.145913\\
0.651629	-0.145173\\
0.656642	-0.144434\\
0.661654	-0.143699\\
0.666667	-0.142975\\
0.671679	-0.142266\\
0.676692	-0.141576\\
0.681704	-0.14091\\
0.686717	-0.140272\\
0.691729	-0.139668\\
0.696742	-0.139101\\
0.701754	-0.138577\\
0.706767	-0.1381\\
0.711779	-0.137674\\
0.716792	-0.137304\\
0.721805	-0.136995\\
0.726817	-0.136751\\
0.73183	-0.136576\\
0.736842	-0.136475\\
0.741855	-0.136452\\
0.746867	-0.136543\\
0.75188	-0.136757\\
0.756892	-0.137063\\
0.761905	-0.137431\\
0.766917	-0.137828\\
0.77193	-0.138226\\
0.776942	-0.138591\\
0.781955	-0.138894\\
0.786967	-0.139102\\
0.79198	-0.139186\\
0.796992	-0.139186\\
0.802005	-0.139165\\
0.807018	-0.139123\\
0.81203	-0.139061\\
0.817043	-0.138981\\
0.822055	-0.138882\\
0.827068	-0.138765\\
0.83208	-0.138633\\
0.837093	-0.138485\\
0.842105	-0.138323\\
0.847118	-0.138147\\
0.85213	-0.137959\\
0.857143	-0.137759\\
0.862155	-0.137549\\
0.867168	-0.13733\\
0.87218	-0.137102\\
0.877193	-0.136867\\
0.882206	-0.136625\\
0.887218	-0.136378\\
0.892231	-0.136127\\
0.897243	-0.135874\\
0.902256	-0.135618\\
0.907268	-0.135361\\
0.912281	-0.135104\\
0.917293	-0.134849\\
0.922306	-0.134596\\
0.927318	-0.134347\\
0.932331	-0.134102\\
0.937343	-0.133863\\
0.942356	-0.1336\\
0.947368	-0.133294\\
0.952381	-0.132956\\
0.957393	-0.132602\\
0.962406	-0.132244\\
0.967419	-0.131896\\
0.972431	-0.131572\\
0.977444	-0.131284\\
0.982456	-0.131046\\
0.987469	-0.130872\\
0.992481	-0.130775\\
0.997494	-0.130724\\
1.00251	-0.130679\\
1.00752	-0.13064\\
1.01253	-0.130605\\
1.01754	-0.130575\\
1.02256	-0.13055\\
1.02757	-0.130529\\
1.03258	-0.130511\\
1.03759	-0.130497\\
1.04261	-0.130485\\
1.04762	-0.130476\\
1.05263	-0.130468\\
1.05764	-0.130462\\
1.06266	-0.130458\\
1.06767	-0.130454\\
1.07268	-0.13045\\
1.07769	-0.130446\\
1.08271	-0.130442\\
1.08772	-0.130436\\
1.09273	-0.13043\\
1.09774	-0.130421\\
1.10276	-0.130411\\
1.10777	-0.130398\\
1.11278	-0.130382\\
1.11779	-0.130362\\
1.12281	-0.13034\\
1.12782	-0.130313\\
1.13283	-0.130282\\
1.13784	-0.130242\\
1.14286	-0.130192\\
1.14787	-0.13013\\
1.15288	-0.130059\\
1.15789	-0.12998\\
1.16291	-0.129892\\
1.16792	-0.129797\\
1.17293	-0.129695\\
1.17794	-0.129588\\
1.18296	-0.129476\\
1.18797	-0.129361\\
1.19298	-0.129242\\
1.19799	-0.129121\\
1.20301	-0.128999\\
1.20802	-0.128877\\
1.21303	-0.128755\\
1.21805	-0.128634\\
1.22306	-0.128516\\
1.22807	-0.1284\\
1.23308	-0.128289\\
1.2381	-0.128182\\
1.24311	-0.128081\\
1.24812	-0.127987\\
1.25313	-0.1279\\
1.25815	-0.127821\\
1.26316	-0.127752\\
1.26817	-0.127692\\
1.27318	-0.12764\\
1.2782	-0.127591\\
1.28321	-0.127545\\
1.28822	-0.127501\\
1.29323	-0.127459\\
1.29825	-0.127419\\
1.30326	-0.127381\\
1.30827	-0.127344\\
1.31328	-0.127308\\
1.3183	-0.127273\\
1.32331	-0.127237\\
1.32832	-0.127202\\
1.33333	-0.127168\\
1.33835	-0.127134\\
1.34336	-0.127102\\
1.34837	-0.12707\\
1.35338	-0.12704\\
1.3584	-0.127011\\
1.36341	-0.126982\\
1.36842	-0.126955\\
1.37343	-0.126928\\
1.37845	-0.126902\\
1.38346	-0.126876\\
1.38847	-0.126851\\
1.39348	-0.126827\\
1.3985	-0.126803\\
1.40351	-0.126779\\
1.40852	-0.126756\\
1.41353	-0.126732\\
1.41855	-0.126709\\
1.42356	-0.126685\\
1.42857	-0.126661\\
1.43358	-0.126638\\
1.4386	-0.126613\\
1.44361	-0.126588\\
1.44862	-0.126563\\
1.45363	-0.126537\\
1.45865	-0.126511\\
1.46366	-0.126483\\
1.46867	-0.126455\\
1.47368	-0.126424\\
1.4787	-0.126391\\
1.48371	-0.126354\\
1.48872	-0.126316\\
1.49373	-0.126275\\
1.49875	-0.126232\\
1.50376	-0.126187\\
1.50877	-0.126141\\
1.51378	-0.126094\\
1.5188	-0.126046\\
1.52381	-0.125998\\
1.52882	-0.125949\\
1.53383	-0.125901\\
1.53885	-0.125853\\
1.54386	-0.125805\\
1.54887	-0.125759\\
1.55388	-0.125714\\
1.5589	-0.125671\\
1.56391	-0.125629\\
1.56892	-0.12559\\
1.57393	-0.125553\\
1.57895	-0.12552\\
1.58396	-0.125489\\
1.58897	-0.125462\\
1.59398	-0.125439\\
1.599	-0.12542\\
1.60401	-0.125405\\
1.60902	-0.125396\\
1.61404	-0.125391\\
1.61905	-0.125392\\
1.62406	-0.125397\\
1.62907	-0.125405\\
1.63409	-0.125414\\
1.6391	-0.125424\\
1.64411	-0.125434\\
1.64912	-0.125442\\
1.65414	-0.125448\\
1.65915	-0.125449\\
1.66416	-0.125446\\
1.66917	-0.12544\\
1.67419	-0.125431\\
1.6792	-0.125422\\
1.68421	-0.125412\\
1.68922	-0.1254\\
1.69424	-0.125387\\
1.69925	-0.125373\\
1.70426	-0.125358\\
1.70927	-0.125343\\
1.71429	-0.125326\\
1.7193	-0.125309\\
1.72431	-0.125291\\
1.72932	-0.125272\\
1.73434	-0.125253\\
1.73935	-0.125233\\
1.74436	-0.125213\\
1.74937	-0.125192\\
1.75439	-0.125171\\
1.7594	-0.12515\\
1.76441	-0.125129\\
1.76942	-0.125107\\
1.77444	-0.125085\\
1.77945	-0.125063\\
1.78446	-0.125041\\
1.78947	-0.12502\\
1.79449	-0.124998\\
1.7995	-0.124977\\
1.80451	-0.124953\\
1.80952	-0.124926\\
1.81454	-0.124896\\
1.81955	-0.124864\\
1.82456	-0.124831\\
1.82957	-0.124798\\
1.83459	-0.124765\\
1.8396	-0.124734\\
1.84461	-0.124704\\
1.84962	-0.124678\\
1.85464	-0.124656\\
1.85965	-0.124638\\
1.86466	-0.124623\\
1.86967	-0.124609\\
1.87469	-0.124594\\
1.8797	-0.12458\\
1.88471	-0.124566\\
1.88972	-0.124553\\
1.89474	-0.12454\\
1.89975	-0.124528\\
1.90476	-0.124516\\
1.90977	-0.124506\\
1.91479	-0.124496\\
1.9198	-0.124487\\
1.92481	-0.124479\\
1.92982	-0.124472\\
1.93484	-0.124466\\
1.93985	-0.124462\\
1.94486	-0.12446\\
1.94987	-0.124461\\
1.95489	-0.124464\\
1.9599	-0.124469\\
1.96491	-0.124477\\
1.96992	-0.124485\\
1.97494	-0.124496\\
1.97995	-0.124508\\
1.98496	-0.124522\\
1.98997	-0.124536\\
1.99499	-0.124552\\
2	-0.124569\\
2	-0.123803\\
1.99499	-0.123831\\
1.98997	-0.123857\\
1.98496	-0.123882\\
1.97995	-0.123907\\
1.97494	-0.123929\\
1.96992	-0.123951\\
1.96491	-0.12397\\
1.9599	-0.123988\\
1.95489	-0.124004\\
1.94987	-0.124018\\
1.94486	-0.124029\\
1.93985	-0.124038\\
1.93484	-0.124045\\
1.92982	-0.124049\\
1.92481	-0.124053\\
1.9198	-0.124055\\
1.91479	-0.124056\\
1.90977	-0.124056\\
1.90476	-0.124055\\
1.89975	-0.124054\\
1.89474	-0.124051\\
1.88972	-0.124048\\
1.88471	-0.124044\\
1.8797	-0.124039\\
1.87469	-0.124034\\
1.86967	-0.124029\\
1.86466	-0.124023\\
1.85965	-0.124016\\
1.85464	-0.124007\\
1.84962	-0.123993\\
1.84461	-0.123975\\
1.8396	-0.123954\\
1.83459	-0.12393\\
1.82957	-0.123905\\
1.82456	-0.12388\\
1.81955	-0.123854\\
1.81454	-0.12383\\
1.80952	-0.123807\\
1.80451	-0.123788\\
1.7995	-0.123771\\
1.79449	-0.123757\\
1.78947	-0.123743\\
1.78446	-0.123728\\
1.77945	-0.123714\\
1.77444	-0.123699\\
1.76942	-0.123684\\
1.76441	-0.12367\\
1.7594	-0.123656\\
1.75439	-0.123642\\
1.74937	-0.123628\\
1.74436	-0.123615\\
1.73935	-0.123603\\
1.73434	-0.123591\\
1.72932	-0.12358\\
1.72431	-0.123569\\
1.7193	-0.12356\\
1.71429	-0.123551\\
1.70927	-0.123544\\
1.70426	-0.123537\\
1.69925	-0.123532\\
1.69424	-0.123528\\
1.68922	-0.123526\\
1.68421	-0.123524\\
1.6792	-0.123525\\
1.67419	-0.123526\\
1.66917	-0.12353\\
1.66416	-0.123535\\
1.65915	-0.123544\\
1.65414	-0.123558\\
1.64912	-0.123576\\
1.64411	-0.123598\\
1.6391	-0.123621\\
1.63409	-0.123644\\
1.62907	-0.123668\\
1.62406	-0.123689\\
1.61905	-0.123708\\
1.61404	-0.123723\\
1.60902	-0.123734\\
1.60401	-0.123738\\
1.599	-0.123738\\
1.59398	-0.123734\\
1.58897	-0.123725\\
1.58396	-0.123713\\
1.57895	-0.123696\\
1.57393	-0.123677\\
1.56892	-0.123654\\
1.56391	-0.123628\\
1.5589	-0.1236\\
1.55388	-0.12357\\
1.54887	-0.123537\\
1.54386	-0.123503\\
1.53885	-0.123467\\
1.53383	-0.12343\\
1.52882	-0.123393\\
1.52381	-0.123354\\
1.5188	-0.123315\\
1.51378	-0.123277\\
1.50877	-0.123238\\
1.50376	-0.1232\\
1.49875	-0.123163\\
1.49373	-0.123127\\
1.48872	-0.123093\\
1.48371	-0.12306\\
1.4787	-0.123029\\
1.47368	-0.123001\\
1.46867	-0.122975\\
1.46366	-0.122951\\
1.45865	-0.122928\\
1.45363	-0.122906\\
1.44862	-0.122884\\
1.44361	-0.122863\\
1.4386	-0.122842\\
1.43358	-0.122822\\
1.42857	-0.122803\\
1.42356	-0.122783\\
1.41855	-0.122765\\
1.41353	-0.122747\\
1.40852	-0.122729\\
1.40351	-0.122711\\
1.3985	-0.122694\\
1.39348	-0.122677\\
1.38847	-0.122661\\
1.38346	-0.122645\\
1.37845	-0.122629\\
1.37343	-0.122613\\
1.36842	-0.122597\\
1.36341	-0.122582\\
1.3584	-0.122566\\
1.35338	-0.12255\\
1.34837	-0.122534\\
1.34336	-0.122519\\
1.33835	-0.122502\\
1.33333	-0.122486\\
1.32832	-0.122469\\
1.32331	-0.122453\\
1.3183	-0.122437\\
1.31328	-0.122422\\
1.30827	-0.122407\\
1.30326	-0.122391\\
1.29825	-0.122375\\
1.29323	-0.122357\\
1.28822	-0.122338\\
1.28321	-0.122317\\
1.2782	-0.122294\\
1.27318	-0.122269\\
1.26817	-0.12224\\
1.26316	-0.122204\\
1.25815	-0.122157\\
1.25313	-0.122102\\
1.24812	-0.122038\\
1.24311	-0.121965\\
1.2381	-0.121886\\
1.23308	-0.1218\\
1.22807	-0.121709\\
1.22306	-0.121613\\
1.21805	-0.121513\\
1.21303	-0.12141\\
1.20802	-0.121304\\
1.20301	-0.121197\\
1.19799	-0.121088\\
1.19298	-0.12098\\
1.18797	-0.120872\\
1.18296	-0.120766\\
1.17794	-0.120662\\
1.17293	-0.120561\\
1.16792	-0.120464\\
1.16291	-0.120371\\
1.15789	-0.120284\\
1.15288	-0.120203\\
1.14787	-0.12013\\
1.14286	-0.120064\\
1.13784	-0.120007\\
1.13283	-0.11996\\
1.12782	-0.11992\\
1.12281	-0.119882\\
1.11779	-0.119846\\
1.11278	-0.119813\\
1.10777	-0.119782\\
1.10276	-0.119753\\
1.09774	-0.119725\\
1.09273	-0.119699\\
1.08772	-0.119673\\
1.08271	-0.119649\\
1.07769	-0.119625\\
1.07268	-0.119601\\
1.06767	-0.119578\\
1.06266	-0.119555\\
1.05764	-0.119531\\
1.05263	-0.119507\\
1.04762	-0.119483\\
1.04261	-0.119458\\
1.03759	-0.119432\\
1.03258	-0.119405\\
1.02757	-0.119377\\
1.02256	-0.119347\\
1.01754	-0.119315\\
1.01253	-0.119282\\
1.00752	-0.119246\\
1.00251	-0.119208\\
0.997494	-0.119167\\
0.992481	-0.119123\\
0.987469	-0.119036\\
0.982456	-0.118875\\
0.977444	-0.118653\\
0.972431	-0.118384\\
0.967419	-0.118081\\
0.962406	-0.117758\\
0.957393	-0.117427\\
0.952381	-0.117103\\
0.947368	-0.116798\\
0.942356	-0.116525\\
0.937343	-0.116299\\
0.932331	-0.116098\\
0.927318	-0.115893\\
0.922306	-0.115685\\
0.917293	-0.115474\\
0.912281	-0.115261\\
0.907268	-0.115047\\
0.902256	-0.114832\\
0.897243	-0.114618\\
0.892231	-0.114406\\
0.887218	-0.114195\\
0.882206	-0.113987\\
0.877193	-0.113782\\
0.87218	-0.113581\\
0.867168	-0.113386\\
0.862155	-0.113195\\
0.857143	-0.113012\\
0.85213	-0.112836\\
0.847118	-0.112667\\
0.842105	-0.112508\\
0.837093	-0.112359\\
0.83208	-0.11222\\
0.827068	-0.112092\\
0.822055	-0.111977\\
0.817043	-0.111876\\
0.81203	-0.111788\\
0.807018	-0.111716\\
0.802005	-0.11166\\
0.796992	-0.111621\\
0.79198	-0.111601\\
0.786967	-0.111661\\
0.781955	-0.111843\\
0.776942	-0.112116\\
0.77193	-0.11245\\
0.766917	-0.112814\\
0.761905	-0.113177\\
0.756892	-0.113508\\
0.75188	-0.113778\\
0.746867	-0.113955\\
0.741855	-0.114008\\
0.736842	-0.11395\\
0.73183	-0.113813\\
0.726817	-0.113605\\
0.721805	-0.113328\\
0.716792	-0.112989\\
0.711779	-0.112592\\
0.706767	-0.112143\\
0.701754	-0.111645\\
0.696742	-0.111104\\
0.691729	-0.110525\\
0.686717	-0.109913\\
0.681704	-0.109273\\
0.676692	-0.108609\\
0.671679	-0.107926\\
0.666667	-0.107229\\
0.661654	-0.106523\\
0.656642	-0.105812\\
0.651629	-0.105102\\
0.646617	-0.104396\\
0.641604	-0.103701\\
0.636591	-0.103019\\
0.631579	-0.102357\\
0.626566	-0.101718\\
0.621554	-0.101107\\
0.616541	-0.100528\\
0.611529	-0.099987\\
0.606516	-0.0994871\\
0.601504	-0.0990331\\
0.596491	-0.0986295\\
0.591479	-0.0982599\\
0.586466	-0.0979043\\
0.581454	-0.0975612\\
0.576441	-0.0972292\\
0.571429	-0.0969068\\
0.566416	-0.0965926\\
0.561404	-0.0962849\\
0.556391	-0.0959823\\
0.551378	-0.0956832\\
0.546366	-0.095386\\
0.541353	-0.0950891\\
0.536341	-0.094791\\
0.531328	-0.0944899\\
0.526316	-0.0941844\\
0.521303	-0.0938727\\
0.516291	-0.0935533\\
0.511278	-0.0932244\\
0.506266	-0.0928847\\
0.501253	-0.0925323\\
0.496241	-0.0921661\\
0.491228	-0.091784\\
0.486216	-0.091385\\
0.481203	-0.0909672\\
0.47619	-0.0905297\\
0.471178	-0.0900705\\
0.466165	-0.0895889\\
0.461153	-0.089083\\
0.45614	-0.0885521\\
0.451128	-0.0880057\\
0.446115	-0.0874511\\
0.441103	-0.0868809\\
0.43609	-0.0862895\\
0.431078	-0.0856697\\
0.426065	-0.085016\\
0.421053	-0.0843213\\
0.41604	-0.0835802\\
0.411028	-0.0827857\\
0.406015	-0.0819324\\
0.401003	-0.0810047\\
0.39599	-0.0799954\\
0.390977	-0.0789072\\
0.385965	-0.077745\\
0.380952	-0.0765116\\
0.37594	-0.0752118\\
0.370927	-0.0738486\\
0.365915	-0.0724266\\
0.360902	-0.0709489\\
0.35589	-0.0694198\\
0.350877	-0.0678427\\
0.345865	-0.0662216\\
0.340852	-0.0645598\\
0.33584	-0.0628613\\
0.330827	-0.0611293\\
0.325815	-0.0593674\\
0.320802	-0.0575787\\
0.315789	-0.0557669\\
0.310777	-0.0539347\\
0.305764	-0.0520856\\
0.300752	-0.0502221\\
0.295739	-0.0483474\\
0.290727	-0.0464639\\
0.285714	-0.0445746\\
0.280702	-0.0426815\\
0.275689	-0.0407874\\
0.270677	-0.0388942\\
0.265664	-0.0370043\\
0.260652	-0.0351193\\
0.255639	-0.0332245\\
0.250627	-0.0313019\\
0.245614	-0.02935\\
0.240602	-0.0273667\\
0.235589	-0.0253503\\
0.230576	-0.0232985\\
0.225564	-0.0212095\\
0.220551	-0.0190806\\
0.215539	-0.0169101\\
0.210526	-0.0146953\\
0.205514	-0.012434\\
0.200501	-0.0101236\\
0.195489	-0.00776183\\
0.190476	-0.00534618\\
0.185464	-0.002874\\
0.180451	-0.000343019\\
0.175439	0.00224968\\
0.170426	0.00490612\\
0.165414	0.00762951\\
0.160401	0.0104216\\
0.155388	0.0132859\\
0.150376	0.0162238\\
0.145363	0.0192392\\
0.140351	0.022333\\
0.135338	0.0255097\\
0.130326	0.0287698\\
0.125313	0.0321179\\
0.120301	0.0355544\\
0.115288	0.0390844\\
0.110276	0.0427078\\
0.105263	0.0464302\\
0.100251	0.0502511\\
0.0952381	0.0541769\\
0.0902256	0.0582065\\
0.085213	0.0623472\\
0.0802005	0.0665971\\
0.075188	0.0709645\\
0.0701754	0.075447\\
0.0651629	0.0800535\\
0.0601504	0.082382\\
0.0551378	0.0831965\\
0.0501253	0.0825409\\
0.0451128	0.080308\\
0.0401003	0.0765951\\
0.0350877	0.0713365\\
0.0300752	0.0646779\\
0.0250627	0.056597\\
0.0200501	0.0472732\\
0.0150376	0.0367258\\
0.0100251	0.0280852\\
0.00501253	0.0310077\\
0	0.01256\\
}--cycle;
\addplot [color=mycolor2, line width=1.0pt, forget plot]
  table[row sep=crcr]{%
-0.124	0\\
-0.122	0\\
-0.12	0\\
-0.118	0\\
-0.116	0\\
-0.114	0\\
-0.112	0\\
-0.11	0\\
-0.108	0\\
-0.106	0\\
-0.104	0\\
-0.102	0\\
-0.1	0\\
-0.098	0\\
-0.096	0\\
-0.094	0\\
-0.092	0\\
-0.09	0\\
-0.088	0\\
-0.086	0\\
-0.084	0\\
-0.082	0\\
-0.08	0\\
-0.078	0\\
-0.076	0\\
-0.074	0\\
-0.072	0\\
-0.07	0\\
-0.068	0\\
-0.066	0\\
-0.064	0\\
-0.062	0\\
-0.06	0\\
-0.058	0\\
-0.056	0\\
-0.054	0\\
-0.052	0\\
-0.05	0\\
-0.048	0\\
-0.046	0\\
-0.044	0\\
-0.042	0\\
-0.04	0\\
-0.038	0\\
-0.036	0\\
-0.034	0\\
-0.032	0\\
-0.03	0\\
-0.028	0\\
-0.026	0\\
-0.024	0\\
-0.022	0\\
-0.02	0\\
-0.018	0\\
-0.016	0\\
-0.014	0\\
-0.012	0\\
-0.01	0\\
-0.008	0\\
-0.006	0\\
-0.004	0\\
-0.002	0\\
0	0\\
0.002	-0.0002405\\
0.004	-0.000500288\\
0.006	-0.000791079\\
0.008	-0.0011171\\
0.01	-0.00147799\\
0.012	-0.00187086\\
0.014	-0.00229179\\
0.016	-0.00273671\\
0.018	-0.00320202\\
0.02	-0.0036848\\
0.022	-0.00418296\\
0.024	-0.00469517\\
0.026	-0.00522081\\
0.028	-0.0057598\\
0.03	-0.00631251\\
0.032	-0.00687961\\
0.034	-0.007462\\
0.036	-0.00806067\\
0.038	-0.00867667\\
0.04	-0.00931106\\
0.042	-0.00996484\\
0.044	-0.010639\\
0.046	-0.0113343\\
0.048	-0.0120518\\
0.05	-0.0127919\\
0.052	-0.0135554\\
0.054	-0.0143428\\
0.056	-0.0151546\\
0.058	-0.0159912\\
0.06	-0.0168528\\
0.062	-0.0177397\\
0.064	-0.018652\\
0.066	-0.0195899\\
0.068	-0.0205534\\
0.07	-0.0215425\\
0.072	-0.022557\\
0.074	-0.0235969\\
0.076	-0.0246618\\
0.078	-0.0257516\\
0.08	-0.0268659\\
0.082	-0.0280043\\
0.084	-0.0291664\\
0.086	-0.0303518\\
0.088	-0.0315598\\
0.09	-0.03279\\
0.092	-0.0340416\\
0.094	-0.035314\\
0.096	-0.0366065\\
0.098	-0.0379183\\
0.1	-0.0392487\\
0.102	-0.0405967\\
0.104	-0.0419617\\
0.106	-0.0433425\\
0.108	-0.0447384\\
0.11	-0.0461483\\
0.112	-0.0475714\\
0.114	-0.0490064\\
0.116	-0.0504525\\
0.118	-0.0519086\\
0.12	-0.0533736\\
0.122	-0.0548464\\
0.124	-0.056326\\
0.126	-0.0578111\\
0.128	-0.0593008\\
0.13	-0.0607938\\
0.132	-0.062289\\
0.134	-0.0637853\\
0.136	-0.0652816\\
0.138	-0.0667766\\
0.14	-0.0682693\\
0.142	-0.0697585\\
0.144	-0.0712431\\
0.146	-0.072722\\
0.148	-0.0741941\\
0.15	-0.0756581\\
0.152	-0.0771132\\
0.154	-0.0785581\\
0.156	-0.0799918\\
0.158	-0.0814133\\
0.16	-0.0828216\\
0.162	-0.0842156\\
0.164	-0.0855944\\
0.166	-0.0869569\\
0.168	-0.0883024\\
0.17	-0.0896298\\
0.172	-0.0909382\\
0.174	-0.0922269\\
0.176	-0.093495\\
0.178	-0.0947417\\
0.18	-0.0959662\\
0.182	-0.0971678\\
0.184	-0.0983458\\
0.186	-0.0994995\\
0.188	-0.100628\\
0.19	-0.101732\\
0.192	-0.102809\\
0.194	-0.103859\\
0.196	-0.104883\\
0.198	-0.105879\\
0.2	-0.106847\\
0.202	-0.107786\\
0.204	-0.108697\\
0.206	-0.109579\\
0.208	-0.110431\\
0.21	-0.111254\\
0.212	-0.112047\\
0.214	-0.11281\\
0.216	-0.113543\\
0.218	-0.114246\\
0.22	-0.114919\\
0.222	-0.115561\\
0.224	-0.116174\\
0.226	-0.116756\\
0.228	-0.117307\\
0.23	-0.117829\\
0.232	-0.118321\\
0.234	-0.118783\\
0.236	-0.119216\\
0.238	-0.119619\\
0.24	-0.119993\\
0.242	-0.120339\\
0.244	-0.120656\\
0.246	-0.120945\\
0.248	-0.121206\\
0.25	-0.12144\\
0.252	-0.121647\\
0.254	-0.121828\\
0.256	-0.121982\\
0.258	-0.122111\\
0.26	-0.122215\\
0.262	-0.122295\\
0.264	-0.12235\\
0.266	-0.122382\\
0.268	-0.122391\\
0.27	-0.122378\\
0.272	-0.122343\\
0.274	-0.122288\\
0.276	-0.122211\\
0.278	-0.122115\\
0.28	-0.122\\
0.282	-0.121866\\
0.284	-0.121714\\
0.286	-0.121546\\
0.288	-0.12136\\
0.29	-0.121159\\
0.292	-0.120943\\
0.294	-0.120712\\
0.296	-0.120467\\
0.298	-0.120209\\
0.3	-0.119939\\
0.302	-0.119657\\
0.304	-0.119364\\
0.306	-0.119061\\
0.308	-0.118748\\
0.31	-0.118426\\
0.312	-0.118095\\
0.314	-0.117757\\
0.316	-0.117412\\
0.318	-0.11706\\
0.32	-0.116703\\
0.322	-0.116341\\
0.324	-0.115974\\
0.326	-0.115604\\
0.328	-0.11523\\
0.33	-0.114854\\
0.332	-0.114476\\
0.334	-0.114097\\
0.336	-0.113717\\
0.338	-0.113337\\
0.34	-0.112957\\
0.342	-0.112578\\
0.344	-0.1122\\
0.346	-0.111825\\
0.348	-0.111452\\
0.35	-0.111082\\
0.352	-0.110716\\
0.354	-0.110354\\
0.356	-0.109996\\
0.358	-0.109644\\
0.36	-0.109297\\
0.362	-0.108955\\
0.364	-0.108621\\
0.366	-0.108293\\
0.368	-0.107972\\
0.37	-0.107658\\
0.372	-0.107353\\
0.374	-0.107056\\
0.376	-0.106768\\
0.378	-0.106489\\
0.38	-0.106219\\
0.382	-0.105959\\
0.384	-0.105708\\
0.386	-0.105469\\
0.388	-0.105239\\
0.39	-0.105021\\
0.392	-0.104814\\
0.394	-0.104618\\
0.396	-0.104433\\
0.398	-0.104261\\
0.4	-0.1041\\
0.402	-0.103952\\
0.404	-0.103817\\
0.406	-0.103693\\
0.408	-0.103583\\
0.41	-0.103486\\
0.412	-0.103401\\
0.414	-0.10333\\
0.416	-0.103272\\
0.418	-0.103227\\
0.42	-0.103196\\
0.422	-0.103179\\
0.424	-0.103174\\
0.426	-0.103184\\
0.428	-0.103207\\
0.43	-0.103244\\
0.432	-0.103295\\
0.434	-0.103359\\
0.436	-0.103437\\
0.438	-0.103529\\
0.44	-0.103634\\
0.442	-0.103753\\
0.444	-0.103885\\
0.446	-0.104031\\
0.448	-0.10419\\
0.45	-0.104361\\
0.452	-0.104546\\
0.454	-0.104744\\
0.456	-0.104955\\
0.458	-0.105178\\
0.46	-0.105414\\
0.462	-0.105661\\
0.464	-0.105921\\
0.466	-0.106192\\
0.468	-0.106475\\
0.47	-0.10677\\
0.472	-0.107075\\
0.474	-0.107391\\
0.476	-0.107717\\
0.478	-0.108054\\
0.48	-0.108401\\
0.482	-0.108757\\
0.484	-0.109122\\
0.486	-0.109496\\
0.488	-0.109879\\
0.49	-0.11027\\
0.492	-0.110669\\
0.494	-0.111075\\
0.496	-0.111488\\
0.498	-0.111908\\
0.5	-0.112334\\
0.502	-0.112766\\
0.504	-0.113203\\
0.506	-0.113645\\
0.508	-0.114092\\
0.51	-0.114543\\
0.512	-0.114998\\
0.514	-0.115456\\
0.516	-0.115916\\
0.518	-0.116379\\
0.52	-0.116844\\
0.522	-0.11731\\
0.524	-0.117776\\
0.526	-0.118244\\
0.528	-0.118711\\
0.53	-0.119178\\
0.532	-0.119644\\
0.534	-0.120108\\
0.536	-0.120571\\
0.538	-0.121031\\
0.54	-0.121488\\
0.542	-0.121942\\
0.544	-0.122392\\
0.546	-0.122839\\
0.548	-0.12328\\
0.55	-0.123717\\
0.552	-0.124148\\
0.554	-0.124573\\
0.556	-0.124992\\
0.558	-0.125404\\
0.56	-0.12581\\
0.562	-0.126208\\
0.564	-0.126598\\
0.566	-0.12698\\
0.568	-0.127353\\
0.57	-0.127718\\
0.572	-0.128073\\
0.574	-0.128419\\
0.576	-0.128756\\
0.578	-0.129082\\
0.58	-0.129399\\
0.582	-0.129704\\
0.584	-0.129999\\
0.586	-0.130284\\
0.588	-0.130557\\
0.59	-0.130819\\
0.592	-0.131069\\
0.594	-0.131308\\
0.596	-0.131535\\
0.598	-0.13175\\
0.6	-0.131954\\
0.602	-0.132145\\
0.604	-0.132324\\
0.606	-0.132492\\
0.608	-0.132647\\
0.61	-0.13279\\
0.612	-0.132922\\
0.614	-0.133041\\
0.616	-0.133148\\
0.618	-0.133243\\
0.62	-0.133326\\
0.622	-0.133398\\
0.624	-0.133458\\
0.626	-0.133507\\
0.628	-0.133544\\
0.63	-0.13357\\
0.632	-0.133585\\
0.634	-0.133589\\
0.636	-0.133583\\
0.638	-0.133566\\
0.64	-0.133539\\
0.642	-0.133502\\
0.644	-0.133455\\
0.646	-0.133399\\
0.648	-0.133334\\
0.65	-0.13326\\
0.652	-0.133177\\
0.654	-0.133086\\
0.656	-0.132987\\
0.658	-0.13288\\
0.66	-0.132766\\
0.662	-0.132645\\
0.664	-0.132518\\
0.666	-0.132384\\
0.668	-0.132244\\
0.67	-0.132098\\
0.672	-0.131947\\
0.674	-0.131791\\
0.676	-0.131631\\
0.678	-0.131466\\
0.68	-0.131297\\
0.682	-0.131125\\
0.684	-0.130949\\
0.686	-0.130771\\
0.688	-0.13059\\
0.69	-0.130407\\
0.692	-0.130222\\
0.694	-0.130035\\
0.696	-0.129848\\
0.698	-0.129659\\
0.7	-0.12947\\
0.702	-0.12928\\
0.704	-0.129091\\
0.706	-0.128902\\
0.708	-0.128713\\
0.71	-0.128526\\
0.712	-0.128339\\
0.714	-0.128154\\
0.716	-0.12797\\
0.718	-0.127789\\
0.72	-0.12761\\
0.722	-0.127433\\
0.724	-0.127258\\
0.726	-0.127086\\
0.728	-0.126918\\
0.73	-0.126752\\
0.732	-0.12659\\
0.734	-0.126431\\
0.736	-0.126276\\
0.738	-0.126124\\
0.74	-0.125977\\
0.742	-0.125833\\
0.744	-0.125693\\
0.746	-0.125558\\
0.748	-0.125427\\
0.75	-0.1253\\
0.752	-0.125178\\
0.754	-0.12506\\
0.756	-0.124947\\
0.758	-0.124838\\
0.76	-0.124734\\
0.762	-0.124634\\
0.764	-0.124539\\
0.766	-0.124449\\
0.768	-0.124363\\
0.77	-0.124282\\
0.772	-0.124205\\
0.774	-0.124132\\
0.776	-0.124065\\
0.778	-0.124001\\
0.78	-0.123942\\
0.782	-0.123887\\
0.784	-0.123837\\
0.786	-0.12379\\
0.788	-0.123748\\
0.79	-0.12371\\
0.792	-0.123675\\
0.794	-0.123645\\
0.796	-0.123618\\
0.798	-0.123595\\
0.8	-0.123575\\
0.802	-0.123559\\
0.804	-0.123546\\
0.806	-0.123537\\
0.808	-0.12353\\
0.81	-0.123527\\
0.812	-0.123526\\
0.814	-0.123529\\
0.816	-0.123534\\
0.818	-0.123541\\
0.82	-0.123551\\
0.822	-0.123564\\
0.824	-0.123579\\
0.826	-0.123596\\
0.828	-0.123615\\
0.83	-0.123636\\
0.832	-0.123658\\
0.834	-0.123683\\
0.836	-0.123709\\
0.838	-0.123737\\
0.84	-0.123766\\
0.842	-0.123797\\
0.844	-0.123829\\
0.846	-0.123862\\
0.848	-0.123896\\
0.85	-0.123932\\
0.852	-0.123968\\
0.854	-0.124005\\
0.856	-0.124043\\
0.858	-0.124082\\
0.86	-0.124121\\
0.862	-0.124161\\
0.864	-0.124202\\
0.866	-0.124242\\
0.868	-0.124284\\
0.87	-0.124326\\
0.872	-0.124368\\
0.874	-0.12441\\
0.876	-0.124452\\
0.878	-0.124495\\
0.88	-0.124538\\
0.882	-0.124581\\
0.884	-0.124624\\
0.886	-0.124666\\
0.888	-0.124709\\
0.89	-0.124752\\
0.892	-0.124795\\
0.894	-0.124837\\
0.896	-0.124879\\
0.898	-0.124921\\
0.9	-0.124963\\
0.902	-0.125005\\
0.904	-0.125046\\
0.906	-0.125087\\
0.908	-0.125127\\
0.91	-0.125167\\
0.912	-0.125207\\
0.914	-0.125246\\
0.916	-0.125285\\
0.918	-0.125323\\
0.92	-0.125361\\
0.922	-0.125398\\
0.924	-0.125435\\
0.926	-0.125471\\
0.928	-0.125506\\
0.93	-0.125541\\
0.932	-0.125575\\
0.934	-0.125608\\
0.936	-0.125641\\
0.938	-0.125673\\
0.94	-0.125703\\
0.942	-0.125733\\
0.944	-0.125763\\
0.946	-0.125791\\
0.948	-0.125818\\
0.95	-0.125845\\
0.952	-0.12587\\
0.954	-0.125894\\
0.956	-0.125918\\
0.958	-0.12594\\
0.96	-0.125961\\
0.962	-0.125981\\
0.964	-0.126\\
0.966	-0.126017\\
0.968	-0.126033\\
0.97	-0.126048\\
0.972	-0.126062\\
0.974	-0.126074\\
0.976	-0.126085\\
0.978	-0.126095\\
0.98	-0.126103\\
0.982	-0.126109\\
0.984	-0.126115\\
0.986	-0.126118\\
0.988	-0.12612\\
0.99	-0.126121\\
0.992	-0.12612\\
0.994	-0.126118\\
0.996	-0.126114\\
0.998	-0.126108\\
1	-0.126101\\
1.002	-0.126092\\
1.004	-0.126081\\
1.006	-0.126069\\
1.008	-0.126055\\
1.01	-0.12604\\
1.012	-0.126023\\
1.014	-0.126004\\
1.016	-0.125984\\
1.018	-0.125962\\
1.02	-0.125939\\
1.022	-0.125914\\
1.024	-0.125888\\
1.026	-0.12586\\
1.028	-0.12583\\
1.03	-0.125799\\
1.032	-0.125767\\
1.034	-0.125733\\
1.036	-0.125698\\
1.038	-0.125662\\
1.04	-0.125624\\
1.042	-0.125586\\
1.044	-0.125546\\
1.046	-0.125505\\
1.048	-0.125463\\
1.05	-0.12542\\
1.052	-0.125376\\
1.054	-0.125331\\
1.056	-0.125285\\
1.058	-0.125239\\
1.06	-0.125192\\
1.062	-0.125144\\
1.064	-0.125096\\
1.066	-0.125047\\
1.068	-0.124998\\
1.07	-0.124949\\
1.072	-0.124899\\
1.074	-0.124849\\
1.076	-0.124799\\
1.078	-0.12475\\
1.08	-0.1247\\
1.082	-0.12465\\
1.084	-0.124601\\
1.086	-0.124552\\
1.088	-0.124504\\
1.09	-0.124456\\
1.092	-0.124408\\
1.094	-0.124361\\
1.096	-0.124315\\
1.098	-0.12427\\
1.1	-0.124225\\
1.102	-0.124182\\
1.104	-0.124139\\
1.106	-0.124098\\
1.108	-0.124057\\
1.11	-0.124018\\
1.112	-0.12398\\
1.114	-0.123944\\
1.116	-0.123909\\
1.118	-0.123875\\
1.12	-0.123843\\
1.122	-0.123812\\
1.124	-0.123783\\
1.126	-0.123755\\
1.128	-0.12373\\
1.13	-0.123706\\
1.132	-0.123683\\
1.134	-0.123663\\
1.136	-0.123644\\
1.138	-0.123627\\
1.14	-0.123611\\
1.142	-0.123598\\
1.144	-0.123587\\
1.146	-0.123577\\
1.148	-0.123569\\
1.15	-0.123563\\
1.152	-0.123559\\
1.154	-0.123557\\
1.156	-0.123556\\
1.158	-0.123558\\
1.16	-0.123561\\
1.162	-0.123566\\
1.164	-0.123572\\
1.166	-0.123581\\
1.168	-0.123591\\
1.17	-0.123603\\
1.172	-0.123616\\
1.174	-0.123631\\
1.176	-0.123647\\
1.178	-0.123665\\
1.18	-0.123685\\
1.182	-0.123705\\
1.184	-0.123727\\
1.186	-0.123751\\
1.188	-0.123775\\
1.19	-0.123801\\
1.192	-0.123828\\
1.194	-0.123855\\
1.196	-0.123884\\
1.198	-0.123914\\
1.2	-0.123944\\
1.202	-0.123975\\
1.204	-0.124007\\
1.206	-0.12404\\
1.208	-0.124073\\
1.21	-0.124107\\
1.212	-0.124141\\
1.214	-0.124175\\
1.216	-0.124209\\
1.218	-0.124244\\
1.22	-0.124279\\
1.222	-0.124314\\
1.224	-0.124349\\
1.226	-0.124384\\
1.228	-0.124419\\
1.23	-0.124454\\
1.232	-0.124488\\
1.234	-0.124522\\
1.236	-0.124556\\
1.238	-0.124589\\
1.24	-0.124622\\
1.242	-0.124654\\
1.244	-0.124686\\
1.246	-0.124717\\
1.248	-0.124747\\
1.25	-0.124777\\
1.252	-0.124806\\
1.254	-0.124834\\
1.256	-0.124861\\
1.258	-0.124887\\
1.26	-0.124913\\
1.262	-0.124937\\
1.264	-0.124961\\
1.266	-0.124983\\
1.268	-0.125005\\
1.27	-0.125025\\
1.272	-0.125045\\
1.274	-0.125063\\
1.276	-0.125081\\
1.278	-0.125097\\
1.28	-0.125112\\
1.282	-0.125126\\
1.284	-0.125139\\
1.286	-0.125151\\
1.288	-0.125162\\
1.29	-0.125172\\
1.292	-0.125181\\
1.294	-0.125189\\
1.296	-0.125195\\
1.298	-0.125201\\
1.3	-0.125206\\
1.302	-0.12521\\
1.304	-0.125212\\
1.306	-0.125214\\
1.308	-0.125215\\
1.31	-0.125215\\
1.312	-0.125214\\
1.314	-0.125212\\
1.316	-0.12521\\
1.318	-0.125207\\
1.32	-0.125203\\
1.322	-0.125198\\
1.324	-0.125192\\
1.326	-0.125186\\
1.328	-0.125179\\
1.33	-0.125172\\
1.332	-0.125164\\
1.334	-0.125156\\
1.336	-0.125147\\
1.338	-0.125137\\
1.34	-0.125127\\
1.342	-0.125117\\
1.344	-0.125107\\
1.346	-0.125096\\
1.348	-0.125084\\
1.35	-0.125073\\
1.352	-0.125061\\
1.354	-0.125049\\
1.356	-0.125037\\
1.358	-0.125025\\
1.36	-0.125012\\
1.362	-0.125\\
1.364	-0.124987\\
1.366	-0.124975\\
1.368	-0.124962\\
1.37	-0.124949\\
1.372	-0.124937\\
1.374	-0.124924\\
1.376	-0.124911\\
1.378	-0.124899\\
1.38	-0.124887\\
1.382	-0.124874\\
1.384	-0.124862\\
1.386	-0.12485\\
1.388	-0.124838\\
1.39	-0.124826\\
1.392	-0.124815\\
1.394	-0.124803\\
1.396	-0.124792\\
1.398	-0.124781\\
1.4	-0.12477\\
1.402	-0.12476\\
1.404	-0.124749\\
1.406	-0.124739\\
1.408	-0.124729\\
1.41	-0.124719\\
1.412	-0.124709\\
1.414	-0.1247\\
1.416	-0.124691\\
1.418	-0.124681\\
1.42	-0.124673\\
1.422	-0.124664\\
1.424	-0.124655\\
1.426	-0.124647\\
1.428	-0.124639\\
1.43	-0.124631\\
1.432	-0.124623\\
1.434	-0.124616\\
1.436	-0.124608\\
1.438	-0.124601\\
1.44	-0.124594\\
1.442	-0.124587\\
1.444	-0.12458\\
1.446	-0.124573\\
1.448	-0.124566\\
1.45	-0.12456\\
1.452	-0.124553\\
1.454	-0.124547\\
1.456	-0.12454\\
1.458	-0.124534\\
1.46	-0.124528\\
1.462	-0.124522\\
1.464	-0.124516\\
1.466	-0.12451\\
1.468	-0.124505\\
1.47	-0.124499\\
1.472	-0.124493\\
1.474	-0.124488\\
1.476	-0.124482\\
1.478	-0.124477\\
1.48	-0.124471\\
1.482	-0.124466\\
1.484	-0.124461\\
1.486	-0.124455\\
1.488	-0.12445\\
1.49	-0.124445\\
1.492	-0.12444\\
1.494	-0.124435\\
1.496	-0.12443\\
1.498	-0.124425\\
1.5	-0.124421\\
1.502	-0.124416\\
1.504	-0.124412\\
1.506	-0.124407\\
1.508	-0.124403\\
1.51	-0.124399\\
1.512	-0.124395\\
1.514	-0.124391\\
1.516	-0.124387\\
1.518	-0.124383\\
1.52	-0.12438\\
1.522	-0.124376\\
1.524	-0.124373\\
1.526	-0.12437\\
1.528	-0.124367\\
1.53	-0.124364\\
1.532	-0.124361\\
1.534	-0.124359\\
1.536	-0.124357\\
1.538	-0.124355\\
1.54	-0.124353\\
1.542	-0.124351\\
1.544	-0.12435\\
1.546	-0.124349\\
1.548	-0.124348\\
1.55	-0.124347\\
1.552	-0.124347\\
1.554	-0.124346\\
1.556	-0.124346\\
1.558	-0.124347\\
1.56	-0.124347\\
1.562	-0.124348\\
1.564	-0.124349\\
1.566	-0.12435\\
1.568	-0.124351\\
1.57	-0.124353\\
1.572	-0.124354\\
1.574	-0.124356\\
1.576	-0.124359\\
1.578	-0.124361\\
1.58	-0.124364\\
1.582	-0.124367\\
1.584	-0.12437\\
1.586	-0.124373\\
1.588	-0.124376\\
1.59	-0.12438\\
1.592	-0.124384\\
1.594	-0.124387\\
1.596	-0.124392\\
1.598	-0.124396\\
1.6	-0.1244\\
1.602	-0.124404\\
1.604	-0.124409\\
1.606	-0.124413\\
1.608	-0.124418\\
1.61	-0.124423\\
1.612	-0.124427\\
1.614	-0.124432\\
1.616	-0.124437\\
1.618	-0.124442\\
1.62	-0.124446\\
1.622	-0.124451\\
1.624	-0.124456\\
1.626	-0.124461\\
1.628	-0.124465\\
1.63	-0.12447\\
1.632	-0.124474\\
1.634	-0.124478\\
1.636	-0.124482\\
1.638	-0.124486\\
1.64	-0.12449\\
1.642	-0.124494\\
1.644	-0.124498\\
1.646	-0.124501\\
1.648	-0.124504\\
1.65	-0.124507\\
1.652	-0.12451\\
1.654	-0.124512\\
1.656	-0.124514\\
1.658	-0.124516\\
1.66	-0.124518\\
1.662	-0.124519\\
1.664	-0.12452\\
1.666	-0.124521\\
1.668	-0.124522\\
1.67	-0.124522\\
1.672	-0.124522\\
1.674	-0.124521\\
1.676	-0.124521\\
1.678	-0.12452\\
1.68	-0.124518\\
1.682	-0.124517\\
1.684	-0.124515\\
1.686	-0.124512\\
1.688	-0.12451\\
1.69	-0.124507\\
1.692	-0.124503\\
1.694	-0.1245\\
1.696	-0.124496\\
1.698	-0.124492\\
1.7	-0.124487\\
1.702	-0.124482\\
1.704	-0.124477\\
1.706	-0.124472\\
1.708	-0.124466\\
1.71	-0.12446\\
1.712	-0.124454\\
1.714	-0.124448\\
1.716	-0.124441\\
1.718	-0.124434\\
1.72	-0.124427\\
1.722	-0.12442\\
1.724	-0.124412\\
1.726	-0.124405\\
1.728	-0.124397\\
1.73	-0.124389\\
1.732	-0.124381\\
1.734	-0.124373\\
1.736	-0.124365\\
1.738	-0.124356\\
1.74	-0.124348\\
1.742	-0.124339\\
1.744	-0.124331\\
1.746	-0.124322\\
1.748	-0.124313\\
1.75	-0.124305\\
1.752	-0.124296\\
1.754	-0.124287\\
1.756	-0.124279\\
1.758	-0.12427\\
1.76	-0.124262\\
1.762	-0.124253\\
1.764	-0.124245\\
1.766	-0.124237\\
1.768	-0.124229\\
1.77	-0.124221\\
1.772	-0.124213\\
1.774	-0.124205\\
1.776	-0.124198\\
1.778	-0.12419\\
1.78	-0.124183\\
1.782	-0.124176\\
1.784	-0.124169\\
1.786	-0.124162\\
1.788	-0.124156\\
1.79	-0.12415\\
1.792	-0.124144\\
1.794	-0.124138\\
1.796	-0.124132\\
1.798	-0.124127\\
1.8	-0.124122\\
1.802	-0.124117\\
1.804	-0.124113\\
1.806	-0.124108\\
1.808	-0.124104\\
1.81	-0.1241\\
1.812	-0.124097\\
1.814	-0.124093\\
1.816	-0.12409\\
1.818	-0.124088\\
1.82	-0.124085\\
1.822	-0.124083\\
1.824	-0.12408\\
1.826	-0.124078\\
1.828	-0.124077\\
1.83	-0.124075\\
1.832	-0.124074\\
1.834	-0.124073\\
1.836	-0.124072\\
1.838	-0.124072\\
1.84	-0.124071\\
1.842	-0.124071\\
1.844	-0.124071\\
1.846	-0.124071\\
1.848	-0.124071\\
1.85	-0.124072\\
1.852	-0.124072\\
1.854	-0.124073\\
1.856	-0.124074\\
1.858	-0.124075\\
1.86	-0.124076\\
1.862	-0.124077\\
1.864	-0.124078\\
1.866	-0.124079\\
1.868	-0.124081\\
1.87	-0.124082\\
1.872	-0.124084\\
1.874	-0.124086\\
1.876	-0.124087\\
1.878	-0.124089\\
1.88	-0.12409\\
1.882	-0.124092\\
1.884	-0.124094\\
1.886	-0.124096\\
1.888	-0.124097\\
1.89	-0.124099\\
1.892	-0.124101\\
1.894	-0.124102\\
1.896	-0.124104\\
1.898	-0.124105\\
1.9	-0.124107\\
1.902	-0.124108\\
1.904	-0.12411\\
1.906	-0.124111\\
1.908	-0.124112\\
1.91	-0.124114\\
1.912	-0.124115\\
1.914	-0.124116\\
1.916	-0.124117\\
1.918	-0.124117\\
1.92	-0.124118\\
1.922	-0.124119\\
1.924	-0.124119\\
1.926	-0.12412\\
1.928	-0.12412\\
1.93	-0.124121\\
1.932	-0.124121\\
1.934	-0.124121\\
1.936	-0.124121\\
1.938	-0.124121\\
1.94	-0.124121\\
1.942	-0.12412\\
1.944	-0.12412\\
1.946	-0.124119\\
1.948	-0.124119\\
1.95	-0.124118\\
1.952	-0.124117\\
1.954	-0.124117\\
1.956	-0.124116\\
1.958	-0.124115\\
1.96	-0.124113\\
1.962	-0.124112\\
1.964	-0.124111\\
1.966	-0.12411\\
1.968	-0.124108\\
1.97	-0.124107\\
1.972	-0.124105\\
1.974	-0.124104\\
1.976	-0.124102\\
1.978	-0.124101\\
1.98	-0.124099\\
1.982	-0.124097\\
1.984	-0.124095\\
1.986	-0.124094\\
1.988	-0.124092\\
1.99	-0.12409\\
1.992	-0.124088\\
1.994	-0.124086\\
1.996	-0.124084\\
1.998	-0.124082\\
2	-0.12408\\
};
\addplot [color=mycolor3, line width=1.0pt, forget plot]
  table[row sep=crcr]{%
-0.124	0\\
-0.122	0\\
-0.12	0\\
-0.118	0\\
-0.116	0\\
-0.114	0\\
-0.112	0\\
-0.11	0\\
-0.108	0\\
-0.106	0\\
-0.104	0\\
-0.102	0\\
-0.1	0\\
-0.098	0\\
-0.096	0\\
-0.094	0\\
-0.092	0\\
-0.09	0\\
-0.088	0\\
-0.086	0\\
-0.084	0\\
-0.082	0\\
-0.08	0\\
-0.078	0\\
-0.076	0\\
-0.074	0\\
-0.072	0\\
-0.07	0\\
-0.068	0\\
-0.066	0\\
-0.064	0\\
-0.062	0\\
-0.06	0\\
-0.058	0\\
-0.056	0\\
-0.054	0\\
-0.052	0\\
-0.05	0\\
-0.048	0\\
-0.046	0\\
-0.044	0\\
-0.042	0\\
-0.04	0\\
-0.038	0\\
-0.036	0\\
-0.034	0\\
-0.032	0\\
-0.03	0\\
-0.028	0\\
-0.026	0\\
-0.024	0\\
-0.022	0\\
-0.02	0\\
-0.018	0\\
-0.016	0\\
-0.014	0\\
-0.012	0\\
-0.01	0\\
-0.008	0\\
-0.006	0\\
-0.004	0\\
-0.002	0\\
0	0\\
0.002	-0.000425411\\
0.004	-0.000880873\\
0.006	-0.00138591\\
0.008	-0.00194657\\
0.01	-0.00256065\\
0.012	-0.00322143\\
0.014	-0.00392032\\
0.016	-0.00464854\\
0.018	-0.00539815\\
0.02	-0.00616252\\
0.022	-0.00693653\\
0.024	-0.00771653\\
0.026	-0.00850016\\
0.028	-0.00928613\\
0.03	-0.010074\\
0.032	-0.010864\\
0.034	-0.0116567\\
0.036	-0.012453\\
0.038	-0.0132541\\
0.04	-0.014061\\
0.042	-0.0148748\\
0.044	-0.0156966\\
0.046	-0.0165275\\
0.048	-0.0173684\\
0.05	-0.0182201\\
0.052	-0.0190834\\
0.054	-0.0199589\\
0.056	-0.0208472\\
0.058	-0.0217487\\
0.06	-0.022664\\
0.062	-0.0235932\\
0.064	-0.0245369\\
0.066	-0.025495\\
0.068	-0.0264678\\
0.07	-0.0274555\\
0.072	-0.0284579\\
0.074	-0.0294752\\
0.076	-0.0305073\\
0.078	-0.031554\\
0.08	-0.0326152\\
0.082	-0.0336906\\
0.084	-0.0347801\\
0.086	-0.0358833\\
0.088	-0.0369999\\
0.09	-0.0381295\\
0.092	-0.0392716\\
0.094	-0.0404259\\
0.096	-0.0415918\\
0.098	-0.0427688\\
0.1	-0.0439564\\
0.102	-0.045154\\
0.104	-0.0463609\\
0.106	-0.0475766\\
0.108	-0.0488003\\
0.11	-0.0500315\\
0.112	-0.0512694\\
0.114	-0.0525133\\
0.116	-0.0537625\\
0.118	-0.0550163\\
0.12	-0.0562739\\
0.122	-0.0575345\\
0.124	-0.0587975\\
0.126	-0.0600621\\
0.128	-0.0613274\\
0.13	-0.0625928\\
0.132	-0.0638574\\
0.134	-0.0651205\\
0.136	-0.0663814\\
0.138	-0.0676392\\
0.14	-0.0688933\\
0.142	-0.0701428\\
0.144	-0.0713871\\
0.146	-0.0726255\\
0.148	-0.0738571\\
0.15	-0.0750813\\
0.152	-0.0762974\\
0.154	-0.0775048\\
0.156	-0.0787027\\
0.158	-0.0798905\\
0.16	-0.0810676\\
0.162	-0.0822333\\
0.164	-0.0833871\\
0.166	-0.0845283\\
0.168	-0.0856564\\
0.17	-0.0867709\\
0.172	-0.0878711\\
0.174	-0.0889567\\
0.176	-0.090027\\
0.178	-0.0910816\\
0.18	-0.0921201\\
0.182	-0.0931419\\
0.184	-0.0941468\\
0.186	-0.0951342\\
0.188	-0.0961038\\
0.19	-0.0970552\\
0.192	-0.0979881\\
0.194	-0.0989022\\
0.196	-0.0997971\\
0.198	-0.100673\\
0.2	-0.101528\\
0.202	-0.102364\\
0.204	-0.10318\\
0.206	-0.103975\\
0.208	-0.104749\\
0.21	-0.105503\\
0.212	-0.106236\\
0.214	-0.106948\\
0.216	-0.107638\\
0.218	-0.108307\\
0.22	-0.108955\\
0.222	-0.109581\\
0.224	-0.110186\\
0.226	-0.110769\\
0.228	-0.11133\\
0.23	-0.11187\\
0.232	-0.112388\\
0.234	-0.112884\\
0.236	-0.113359\\
0.238	-0.113812\\
0.24	-0.114244\\
0.242	-0.114654\\
0.244	-0.115043\\
0.246	-0.11541\\
0.248	-0.115757\\
0.25	-0.116082\\
0.252	-0.116387\\
0.254	-0.116671\\
0.256	-0.116934\\
0.258	-0.117177\\
0.26	-0.1174\\
0.262	-0.117603\\
0.264	-0.117786\\
0.266	-0.11795\\
0.268	-0.118094\\
0.27	-0.11822\\
0.272	-0.118327\\
0.274	-0.118415\\
0.276	-0.118486\\
0.278	-0.118538\\
0.28	-0.118573\\
0.282	-0.118591\\
0.284	-0.118593\\
0.286	-0.118577\\
0.288	-0.118546\\
0.29	-0.118499\\
0.292	-0.118437\\
0.294	-0.11836\\
0.296	-0.118268\\
0.298	-0.118163\\
0.3	-0.118044\\
0.302	-0.117912\\
0.304	-0.117767\\
0.306	-0.117609\\
0.308	-0.117441\\
0.31	-0.11726\\
0.312	-0.117069\\
0.314	-0.116868\\
0.316	-0.116657\\
0.318	-0.116437\\
0.32	-0.116208\\
0.322	-0.115971\\
0.324	-0.115726\\
0.326	-0.115474\\
0.328	-0.115215\\
0.33	-0.114951\\
0.332	-0.114681\\
0.334	-0.114406\\
0.336	-0.114126\\
0.338	-0.113843\\
0.34	-0.113557\\
0.342	-0.113268\\
0.344	-0.112977\\
0.346	-0.112684\\
0.348	-0.11239\\
0.35	-0.112096\\
0.352	-0.111802\\
0.354	-0.111509\\
0.356	-0.111216\\
0.358	-0.110926\\
0.36	-0.110638\\
0.362	-0.110353\\
0.364	-0.110071\\
0.366	-0.109793\\
0.368	-0.10952\\
0.37	-0.109251\\
0.372	-0.108988\\
0.374	-0.108731\\
0.376	-0.10848\\
0.378	-0.108236\\
0.38	-0.107999\\
0.382	-0.10777\\
0.384	-0.107549\\
0.386	-0.107337\\
0.388	-0.107134\\
0.39	-0.10694\\
0.392	-0.106755\\
0.394	-0.106581\\
0.396	-0.106417\\
0.398	-0.106264\\
0.4	-0.106122\\
0.402	-0.105991\\
0.404	-0.105871\\
0.406	-0.105763\\
0.408	-0.105668\\
0.41	-0.105584\\
0.412	-0.105513\\
0.414	-0.105454\\
0.416	-0.105408\\
0.418	-0.105375\\
0.42	-0.105355\\
0.422	-0.105347\\
0.424	-0.105353\\
0.426	-0.105372\\
0.428	-0.105404\\
0.43	-0.105449\\
0.432	-0.105507\\
0.434	-0.105578\\
0.436	-0.105662\\
0.438	-0.10576\\
0.44	-0.10587\\
0.442	-0.105992\\
0.444	-0.106128\\
0.446	-0.106275\\
0.448	-0.106435\\
0.45	-0.106607\\
0.452	-0.106791\\
0.454	-0.106987\\
0.456	-0.107194\\
0.458	-0.107412\\
0.46	-0.107641\\
0.462	-0.107881\\
0.464	-0.108131\\
0.466	-0.108391\\
0.468	-0.108661\\
0.47	-0.10894\\
0.472	-0.109228\\
0.474	-0.109525\\
0.476	-0.10983\\
0.478	-0.110143\\
0.48	-0.110464\\
0.482	-0.110792\\
0.484	-0.111127\\
0.486	-0.111468\\
0.488	-0.111816\\
0.49	-0.112168\\
0.492	-0.112527\\
0.494	-0.112889\\
0.496	-0.113257\\
0.498	-0.113628\\
0.5	-0.114003\\
0.502	-0.11438\\
0.504	-0.114761\\
0.506	-0.115144\\
0.508	-0.115529\\
0.51	-0.115915\\
0.512	-0.116302\\
0.514	-0.11669\\
0.516	-0.117079\\
0.518	-0.117467\\
0.52	-0.117855\\
0.522	-0.118241\\
0.524	-0.118627\\
0.526	-0.119011\\
0.528	-0.119393\\
0.53	-0.119772\\
0.532	-0.120149\\
0.534	-0.120523\\
0.536	-0.120894\\
0.538	-0.121261\\
0.54	-0.121624\\
0.542	-0.121983\\
0.544	-0.122337\\
0.546	-0.122687\\
0.548	-0.123032\\
0.55	-0.123371\\
0.552	-0.123705\\
0.554	-0.124033\\
0.556	-0.124355\\
0.558	-0.124671\\
0.56	-0.12498\\
0.562	-0.125283\\
0.564	-0.125579\\
0.566	-0.125869\\
0.568	-0.126151\\
0.57	-0.126426\\
0.572	-0.126694\\
0.574	-0.126955\\
0.576	-0.127208\\
0.578	-0.127454\\
0.58	-0.127692\\
0.582	-0.127922\\
0.584	-0.128144\\
0.586	-0.128359\\
0.588	-0.128566\\
0.59	-0.128765\\
0.592	-0.128956\\
0.594	-0.129139\\
0.596	-0.129315\\
0.598	-0.129482\\
0.6	-0.129642\\
0.602	-0.129793\\
0.604	-0.129937\\
0.606	-0.130074\\
0.608	-0.130202\\
0.61	-0.130323\\
0.612	-0.130436\\
0.614	-0.130542\\
0.616	-0.13064\\
0.618	-0.130731\\
0.62	-0.130815\\
0.622	-0.130891\\
0.624	-0.13096\\
0.626	-0.131023\\
0.628	-0.131078\\
0.63	-0.131127\\
0.632	-0.131169\\
0.634	-0.131204\\
0.636	-0.131233\\
0.638	-0.131256\\
0.64	-0.131272\\
0.642	-0.131283\\
0.644	-0.131287\\
0.646	-0.131286\\
0.648	-0.131279\\
0.65	-0.131266\\
0.652	-0.131248\\
0.654	-0.131225\\
0.656	-0.131197\\
0.658	-0.131163\\
0.66	-0.131125\\
0.662	-0.131082\\
0.664	-0.131035\\
0.666	-0.130983\\
0.668	-0.130926\\
0.67	-0.130866\\
0.672	-0.130802\\
0.674	-0.130733\\
0.676	-0.130662\\
0.678	-0.130586\\
0.68	-0.130507\\
0.682	-0.130425\\
0.684	-0.130339\\
0.686	-0.130251\\
0.688	-0.130159\\
0.69	-0.130065\\
0.692	-0.129968\\
0.694	-0.129869\\
0.696	-0.129768\\
0.698	-0.129664\\
0.7	-0.129558\\
0.702	-0.12945\\
0.704	-0.129341\\
0.706	-0.12923\\
0.708	-0.129117\\
0.71	-0.129003\\
0.712	-0.128888\\
0.714	-0.128771\\
0.716	-0.128654\\
0.718	-0.128535\\
0.72	-0.128416\\
0.722	-0.128297\\
0.724	-0.128177\\
0.726	-0.128056\\
0.728	-0.127935\\
0.73	-0.127814\\
0.732	-0.127694\\
0.734	-0.127573\\
0.736	-0.127453\\
0.738	-0.127333\\
0.74	-0.127213\\
0.742	-0.127094\\
0.744	-0.126976\\
0.746	-0.126859\\
0.748	-0.126742\\
0.75	-0.126627\\
0.752	-0.126513\\
0.754	-0.1264\\
0.756	-0.126289\\
0.758	-0.126179\\
0.76	-0.12607\\
0.762	-0.125964\\
0.764	-0.125859\\
0.766	-0.125756\\
0.768	-0.125655\\
0.77	-0.125557\\
0.772	-0.12546\\
0.774	-0.125366\\
0.776	-0.125273\\
0.778	-0.125184\\
0.78	-0.125097\\
0.782	-0.125012\\
0.784	-0.12493\\
0.786	-0.124851\\
0.788	-0.124775\\
0.79	-0.124701\\
0.792	-0.12463\\
0.794	-0.124563\\
0.796	-0.124498\\
0.798	-0.124436\\
0.8	-0.124377\\
0.802	-0.124321\\
0.804	-0.124269\\
0.806	-0.12422\\
0.808	-0.124173\\
0.81	-0.12413\\
0.812	-0.124091\\
0.814	-0.124054\\
0.816	-0.124021\\
0.818	-0.12399\\
0.82	-0.123964\\
0.822	-0.12394\\
0.824	-0.123919\\
0.826	-0.123902\\
0.828	-0.123888\\
0.83	-0.123876\\
0.832	-0.123868\\
0.834	-0.123863\\
0.836	-0.123861\\
0.838	-0.123862\\
0.84	-0.123866\\
0.842	-0.123873\\
0.844	-0.123882\\
0.846	-0.123894\\
0.848	-0.123909\\
0.85	-0.123927\\
0.852	-0.123947\\
0.854	-0.123969\\
0.856	-0.123994\\
0.858	-0.124021\\
0.86	-0.12405\\
0.862	-0.124081\\
0.864	-0.124115\\
0.866	-0.12415\\
0.868	-0.124187\\
0.87	-0.124225\\
0.872	-0.124266\\
0.874	-0.124307\\
0.876	-0.124351\\
0.878	-0.124395\\
0.88	-0.124441\\
0.882	-0.124487\\
0.884	-0.124535\\
0.886	-0.124584\\
0.888	-0.124633\\
0.89	-0.124683\\
0.892	-0.124733\\
0.894	-0.124784\\
0.896	-0.124835\\
0.898	-0.124886\\
0.9	-0.124937\\
0.902	-0.124989\\
0.904	-0.12504\\
0.906	-0.125091\\
0.908	-0.125141\\
0.91	-0.125191\\
0.912	-0.125241\\
0.914	-0.12529\\
0.916	-0.125338\\
0.918	-0.125385\\
0.92	-0.125431\\
0.922	-0.125477\\
0.924	-0.125521\\
0.926	-0.125564\\
0.928	-0.125606\\
0.93	-0.125646\\
0.932	-0.125685\\
0.934	-0.125723\\
0.936	-0.125759\\
0.938	-0.125794\\
0.94	-0.125827\\
0.942	-0.125858\\
0.944	-0.125887\\
0.946	-0.125915\\
0.948	-0.125941\\
0.95	-0.125965\\
0.952	-0.125987\\
0.954	-0.126007\\
0.956	-0.126026\\
0.958	-0.126042\\
0.96	-0.126057\\
0.962	-0.126069\\
0.964	-0.12608\\
0.966	-0.126089\\
0.968	-0.126095\\
0.97	-0.1261\\
0.972	-0.126103\\
0.974	-0.126104\\
0.976	-0.126103\\
0.978	-0.1261\\
0.98	-0.126095\\
0.982	-0.126088\\
0.984	-0.12608\\
0.986	-0.12607\\
0.988	-0.126058\\
0.99	-0.126044\\
0.992	-0.126029\\
0.994	-0.126012\\
0.996	-0.125994\\
0.998	-0.125974\\
1	-0.125953\\
1.002	-0.12593\\
1.004	-0.125906\\
1.006	-0.125881\\
1.008	-0.125855\\
1.01	-0.125827\\
1.012	-0.125799\\
1.014	-0.125769\\
1.016	-0.125738\\
1.018	-0.125707\\
1.02	-0.125674\\
1.022	-0.125641\\
1.024	-0.125607\\
1.026	-0.125573\\
1.028	-0.125538\\
1.03	-0.125503\\
1.032	-0.125467\\
1.034	-0.12543\\
1.036	-0.125394\\
1.038	-0.125357\\
1.04	-0.12532\\
1.042	-0.125283\\
1.044	-0.125246\\
1.046	-0.125209\\
1.048	-0.125172\\
1.05	-0.125135\\
1.052	-0.125098\\
1.054	-0.125061\\
1.056	-0.125025\\
1.058	-0.124989\\
1.06	-0.124953\\
1.062	-0.124918\\
1.064	-0.124883\\
1.066	-0.124849\\
1.068	-0.124816\\
1.07	-0.124783\\
1.072	-0.12475\\
1.074	-0.124718\\
1.076	-0.124687\\
1.078	-0.124657\\
1.08	-0.124627\\
1.082	-0.124598\\
1.084	-0.12457\\
1.086	-0.124543\\
1.088	-0.124516\\
1.09	-0.124491\\
1.092	-0.124466\\
1.094	-0.124442\\
1.096	-0.124419\\
1.098	-0.124397\\
1.1	-0.124376\\
1.102	-0.124356\\
1.104	-0.124336\\
1.106	-0.124318\\
1.108	-0.1243\\
1.11	-0.124284\\
1.112	-0.124268\\
1.114	-0.124253\\
1.116	-0.12424\\
1.118	-0.124227\\
1.12	-0.124215\\
1.122	-0.124203\\
1.124	-0.124193\\
1.126	-0.124184\\
1.128	-0.124175\\
1.13	-0.124167\\
1.132	-0.12416\\
1.134	-0.124154\\
1.136	-0.124149\\
1.138	-0.124144\\
1.14	-0.124141\\
1.142	-0.124138\\
1.144	-0.124135\\
1.146	-0.124134\\
1.148	-0.124133\\
1.15	-0.124133\\
1.152	-0.124133\\
1.154	-0.124134\\
1.156	-0.124136\\
1.158	-0.124138\\
1.16	-0.124141\\
1.162	-0.124145\\
1.164	-0.124149\\
1.166	-0.124154\\
1.168	-0.124159\\
1.17	-0.124164\\
1.172	-0.124171\\
1.174	-0.124177\\
1.176	-0.124184\\
1.178	-0.124192\\
1.18	-0.1242\\
1.182	-0.124208\\
1.184	-0.124217\\
1.186	-0.124226\\
1.188	-0.124236\\
1.19	-0.124246\\
1.192	-0.124256\\
1.194	-0.124266\\
1.196	-0.124277\\
1.198	-0.124289\\
1.2	-0.1243\\
1.202	-0.124312\\
1.204	-0.124324\\
1.206	-0.124336\\
1.208	-0.124349\\
1.21	-0.124361\\
1.212	-0.124374\\
1.214	-0.124387\\
1.216	-0.124401\\
1.218	-0.124414\\
1.22	-0.124428\\
1.222	-0.124442\\
1.224	-0.124456\\
1.226	-0.12447\\
1.228	-0.124485\\
1.23	-0.124499\\
1.232	-0.124514\\
1.234	-0.124529\\
1.236	-0.124543\\
1.238	-0.124558\\
1.24	-0.124573\\
1.242	-0.124588\\
1.244	-0.124603\\
1.246	-0.124619\\
1.248	-0.124634\\
1.25	-0.124649\\
1.252	-0.124664\\
1.254	-0.124679\\
1.256	-0.124695\\
1.258	-0.12471\\
1.26	-0.124725\\
1.262	-0.12474\\
1.264	-0.124755\\
1.266	-0.12477\\
1.268	-0.124785\\
1.27	-0.124799\\
1.272	-0.124814\\
1.274	-0.124828\\
1.276	-0.124843\\
1.278	-0.124857\\
1.28	-0.124871\\
1.282	-0.124885\\
1.284	-0.124898\\
1.286	-0.124912\\
1.288	-0.124925\\
1.29	-0.124938\\
1.292	-0.12495\\
1.294	-0.124963\\
1.296	-0.124975\\
1.298	-0.124986\\
1.3	-0.124998\\
1.302	-0.125009\\
1.304	-0.125019\\
1.306	-0.12503\\
1.308	-0.12504\\
1.31	-0.125049\\
1.312	-0.125059\\
1.314	-0.125067\\
1.316	-0.125076\\
1.318	-0.125083\\
1.32	-0.125091\\
1.322	-0.125098\\
1.324	-0.125104\\
1.326	-0.12511\\
1.328	-0.125116\\
1.33	-0.125121\\
1.332	-0.125125\\
1.334	-0.125129\\
1.336	-0.125133\\
1.338	-0.125136\\
1.34	-0.125138\\
1.342	-0.12514\\
1.344	-0.125141\\
1.346	-0.125142\\
1.348	-0.125142\\
1.35	-0.125142\\
1.352	-0.125141\\
1.354	-0.12514\\
1.356	-0.125138\\
1.358	-0.125135\\
1.36	-0.125132\\
1.362	-0.125129\\
1.364	-0.125125\\
1.366	-0.12512\\
1.368	-0.125115\\
1.37	-0.125109\\
1.372	-0.125103\\
1.374	-0.125096\\
1.376	-0.125089\\
1.378	-0.125082\\
1.38	-0.125074\\
1.382	-0.125066\\
1.384	-0.125057\\
1.386	-0.125047\\
1.388	-0.125038\\
1.39	-0.125028\\
1.392	-0.125017\\
1.394	-0.125007\\
1.396	-0.124996\\
1.398	-0.124984\\
1.4	-0.124973\\
1.402	-0.124961\\
1.404	-0.124949\\
1.406	-0.124936\\
1.408	-0.124924\\
1.41	-0.124911\\
1.412	-0.124898\\
1.414	-0.124885\\
1.416	-0.124872\\
1.418	-0.124859\\
1.42	-0.124845\\
1.422	-0.124832\\
1.424	-0.124818\\
1.426	-0.124805\\
1.428	-0.124792\\
1.43	-0.124778\\
1.432	-0.124765\\
1.434	-0.124751\\
1.436	-0.124738\\
1.438	-0.124725\\
1.44	-0.124712\\
1.442	-0.124699\\
1.444	-0.124687\\
1.446	-0.124674\\
1.448	-0.124662\\
1.45	-0.12465\\
1.452	-0.124638\\
1.454	-0.124627\\
1.456	-0.124615\\
1.458	-0.124604\\
1.46	-0.124593\\
1.462	-0.124583\\
1.464	-0.124573\\
1.466	-0.124563\\
1.468	-0.124554\\
1.47	-0.124545\\
1.472	-0.124536\\
1.474	-0.124528\\
1.476	-0.12452\\
1.478	-0.124512\\
1.48	-0.124505\\
1.482	-0.124498\\
1.484	-0.124491\\
1.486	-0.124485\\
1.488	-0.12448\\
1.49	-0.124474\\
1.492	-0.124469\\
1.494	-0.124465\\
1.496	-0.124461\\
1.498	-0.124457\\
1.5	-0.124454\\
1.502	-0.124451\\
1.504	-0.124448\\
1.506	-0.124446\\
1.508	-0.124444\\
1.51	-0.124443\\
1.512	-0.124442\\
1.514	-0.124441\\
1.516	-0.12444\\
1.518	-0.12444\\
1.52	-0.12444\\
1.522	-0.124441\\
1.524	-0.124441\\
1.526	-0.124442\\
1.528	-0.124444\\
1.53	-0.124445\\
1.532	-0.124447\\
1.534	-0.124449\\
1.536	-0.124451\\
1.538	-0.124454\\
1.54	-0.124456\\
1.542	-0.124459\\
1.544	-0.124462\\
1.546	-0.124465\\
1.548	-0.124468\\
1.55	-0.124471\\
1.552	-0.124475\\
1.554	-0.124478\\
1.556	-0.124482\\
1.558	-0.124485\\
1.56	-0.124489\\
1.562	-0.124493\\
1.564	-0.124496\\
1.566	-0.1245\\
1.568	-0.124504\\
1.57	-0.124508\\
1.572	-0.124511\\
1.574	-0.124515\\
1.576	-0.124518\\
1.578	-0.124522\\
1.58	-0.124525\\
1.582	-0.124529\\
1.584	-0.124532\\
1.586	-0.124535\\
1.588	-0.124538\\
1.59	-0.124541\\
1.592	-0.124544\\
1.594	-0.124546\\
1.596	-0.124549\\
1.598	-0.124551\\
1.6	-0.124553\\
1.602	-0.124556\\
1.604	-0.124557\\
1.606	-0.124559\\
1.608	-0.124561\\
1.61	-0.124562\\
1.612	-0.124563\\
1.614	-0.124564\\
1.616	-0.124565\\
1.618	-0.124566\\
1.62	-0.124566\\
1.622	-0.124566\\
1.624	-0.124567\\
1.626	-0.124566\\
1.628	-0.124566\\
1.63	-0.124566\\
1.632	-0.124565\\
1.634	-0.124564\\
1.636	-0.124563\\
1.638	-0.124562\\
1.64	-0.124561\\
1.642	-0.124559\\
1.644	-0.124558\\
1.646	-0.124556\\
1.648	-0.124554\\
1.65	-0.124552\\
1.652	-0.12455\\
1.654	-0.124547\\
1.656	-0.124545\\
1.658	-0.124542\\
1.66	-0.12454\\
1.662	-0.124537\\
1.664	-0.124534\\
1.666	-0.124531\\
1.668	-0.124528\\
1.67	-0.124524\\
1.672	-0.124521\\
1.674	-0.124517\\
1.676	-0.124514\\
1.678	-0.12451\\
1.68	-0.124507\\
1.682	-0.124503\\
1.684	-0.124499\\
1.686	-0.124495\\
1.688	-0.124491\\
1.69	-0.124487\\
1.692	-0.124483\\
1.694	-0.124479\\
1.696	-0.124475\\
1.698	-0.124471\\
1.7	-0.124467\\
1.702	-0.124463\\
1.704	-0.124459\\
1.706	-0.124454\\
1.708	-0.12445\\
1.71	-0.124446\\
1.712	-0.124442\\
1.714	-0.124438\\
1.716	-0.124433\\
1.718	-0.124429\\
1.72	-0.124425\\
1.722	-0.124421\\
1.724	-0.124417\\
1.726	-0.124413\\
1.728	-0.124408\\
1.73	-0.124404\\
1.732	-0.1244\\
1.734	-0.124396\\
1.736	-0.124392\\
1.738	-0.124388\\
1.74	-0.124384\\
1.742	-0.12438\\
1.744	-0.124376\\
1.746	-0.124372\\
1.748	-0.124368\\
1.75	-0.124364\\
1.752	-0.12436\\
1.754	-0.124356\\
1.756	-0.124352\\
1.758	-0.124348\\
1.76	-0.124344\\
1.762	-0.12434\\
1.764	-0.124336\\
1.766	-0.124333\\
1.768	-0.124329\\
1.77	-0.124325\\
1.772	-0.124321\\
1.774	-0.124317\\
1.776	-0.124314\\
1.778	-0.12431\\
1.78	-0.124306\\
1.782	-0.124303\\
1.784	-0.124299\\
1.786	-0.124295\\
1.788	-0.124292\\
1.79	-0.124288\\
1.792	-0.124284\\
1.794	-0.124281\\
1.796	-0.124277\\
1.798	-0.124274\\
1.8	-0.12427\\
1.802	-0.124267\\
1.804	-0.124263\\
1.806	-0.12426\\
1.808	-0.124257\\
1.81	-0.124253\\
1.812	-0.12425\\
1.814	-0.124246\\
1.816	-0.124243\\
1.818	-0.12424\\
1.82	-0.124237\\
1.822	-0.124234\\
1.824	-0.12423\\
1.826	-0.124227\\
1.828	-0.124224\\
1.83	-0.124221\\
1.832	-0.124218\\
1.834	-0.124215\\
1.836	-0.124212\\
1.838	-0.12421\\
1.84	-0.124207\\
1.842	-0.124204\\
1.844	-0.124201\\
1.846	-0.124199\\
1.848	-0.124196\\
1.85	-0.124194\\
1.852	-0.124191\\
1.854	-0.124189\\
1.856	-0.124187\\
1.858	-0.124185\\
1.86	-0.124182\\
1.862	-0.12418\\
1.864	-0.124178\\
1.866	-0.124176\\
1.868	-0.124175\\
1.87	-0.124173\\
1.872	-0.124171\\
1.874	-0.124169\\
1.876	-0.124168\\
1.878	-0.124167\\
1.88	-0.124165\\
1.882	-0.124164\\
1.884	-0.124163\\
1.886	-0.124161\\
1.888	-0.12416\\
1.89	-0.124159\\
1.892	-0.124159\\
1.894	-0.124158\\
1.896	-0.124157\\
1.898	-0.124156\\
1.9	-0.124156\\
1.902	-0.124155\\
1.904	-0.124155\\
1.906	-0.124154\\
1.908	-0.124154\\
1.91	-0.124154\\
1.912	-0.124154\\
1.914	-0.124153\\
1.916	-0.124153\\
1.918	-0.124153\\
1.92	-0.124153\\
1.922	-0.124154\\
1.924	-0.124154\\
1.926	-0.124154\\
1.928	-0.124154\\
1.93	-0.124154\\
1.932	-0.124155\\
1.934	-0.124155\\
1.936	-0.124155\\
1.938	-0.124156\\
1.94	-0.124156\\
1.942	-0.124157\\
1.944	-0.124157\\
1.946	-0.124157\\
1.948	-0.124158\\
1.95	-0.124158\\
1.952	-0.124159\\
1.954	-0.124159\\
1.956	-0.12416\\
1.958	-0.12416\\
1.96	-0.12416\\
1.962	-0.124161\\
1.964	-0.124161\\
1.966	-0.124161\\
1.968	-0.124162\\
1.97	-0.124162\\
1.972	-0.124162\\
1.974	-0.124162\\
1.976	-0.124162\\
1.978	-0.124163\\
1.98	-0.124163\\
1.982	-0.124163\\
1.984	-0.124162\\
1.986	-0.124162\\
1.988	-0.124162\\
1.99	-0.124162\\
1.992	-0.124161\\
1.994	-0.124161\\
1.996	-0.12416\\
1.998	-0.12416\\
2	-0.124159\\
};
\addplot [color=mycolor4, line width=1.0pt, forget plot]
  table[row sep=crcr]{%
-0.124	0\\
-0.122	0\\
-0.12	0\\
-0.118	0\\
-0.116	0\\
-0.114	0\\
-0.112	0\\
-0.11	0\\
-0.108	0\\
-0.106	0\\
-0.104	0\\
-0.102	0\\
-0.1	0\\
-0.098	0\\
-0.096	0\\
-0.094	0\\
-0.092	0\\
-0.09	0\\
-0.088	0\\
-0.086	0\\
-0.084	0\\
-0.082	0\\
-0.08	0\\
-0.078	0\\
-0.076	0\\
-0.074	0\\
-0.072	0\\
-0.07	0\\
-0.068	0\\
-0.066	0\\
-0.064	0\\
-0.062	0\\
-0.06	0\\
-0.058	0\\
-0.056	0\\
-0.054	0\\
-0.052	0\\
-0.05	0\\
-0.048	0\\
-0.046	0\\
-0.044	0\\
-0.042	0\\
-0.04	0\\
-0.038	0\\
-0.036	0\\
-0.034	0\\
-0.032	0\\
-0.03	0\\
-0.028	0\\
-0.026	0\\
-0.024	0\\
-0.022	0\\
-0.02	0\\
-0.018	0\\
-0.016	0\\
-0.014	0\\
-0.012	0\\
-0.01	0\\
-0.008	0\\
-0.006	0\\
-0.004	0\\
-0.002	0\\
0	0\\
0.002	-0.027024\\
0.004	-0.0404998\\
0.006	-0.0449368\\
0.008	-0.0437969\\
0.01	-0.0396165\\
0.012	-0.0341591\\
0.014	-0.0285753\\
0.016	-0.0235503\\
0.018	-0.0194345\\
0.02	-0.0163489\\
0.022	-0.0142691\\
0.024	-0.0130877\\
0.026	-0.0126581\\
0.028	-0.0128243\\
0.03	-0.013439\\
0.032	-0.0143739\\
0.034	-0.0155231\\
0.036	-0.0168041\\
0.038	-0.0181557\\
0.04	-0.0195345\\
0.042	-0.0209119\\
0.044	-0.0222705\\
0.046	-0.023601\\
0.048	-0.0249\\
0.05	-0.0261675\\
0.052	-0.0274058\\
0.054	-0.0286182\\
0.056	-0.0298084\\
0.058	-0.0309797\\
0.06	-0.0321355\\
0.062	-0.0332781\\
0.064	-0.0344098\\
0.066	-0.0355321\\
0.068	-0.0366462\\
0.07	-0.0377529\\
0.072	-0.0388528\\
0.074	-0.0399463\\
0.076	-0.0410338\\
0.078	-0.0421152\\
0.08	-0.0431909\\
0.082	-0.0442609\\
0.084	-0.0453253\\
0.086	-0.0463841\\
0.088	-0.0474376\\
0.09	-0.0484857\\
0.092	-0.0495287\\
0.094	-0.0505665\\
0.096	-0.0515993\\
0.098	-0.0526271\\
0.1	-0.0536502\\
0.102	-0.0546684\\
0.104	-0.055682\\
0.106	-0.056691\\
0.108	-0.0576953\\
0.11	-0.0586951\\
0.112	-0.0596903\\
0.114	-0.0606809\\
0.116	-0.0616669\\
0.118	-0.0626483\\
0.12	-0.063625\\
0.122	-0.0645969\\
0.124	-0.0655641\\
0.126	-0.0665263\\
0.128	-0.0674835\\
0.13	-0.0684356\\
0.132	-0.0693825\\
0.134	-0.070324\\
0.136	-0.0712599\\
0.138	-0.0721901\\
0.14	-0.0731145\\
0.142	-0.0740329\\
0.144	-0.074945\\
0.146	-0.0758508\\
0.148	-0.0767499\\
0.15	-0.0776422\\
0.152	-0.0785275\\
0.154	-0.0794056\\
0.156	-0.0802762\\
0.158	-0.0811391\\
0.16	-0.0819941\\
0.162	-0.082841\\
0.164	-0.0836795\\
0.166	-0.0845093\\
0.168	-0.0853303\\
0.17	-0.0861422\\
0.172	-0.0869448\\
0.174	-0.0877378\\
0.176	-0.0885209\\
0.178	-0.089294\\
0.18	-0.0900567\\
0.182	-0.090809\\
0.184	-0.0915504\\
0.186	-0.0922809\\
0.188	-0.0930001\\
0.19	-0.0937079\\
0.192	-0.094404\\
0.194	-0.0950882\\
0.196	-0.0957604\\
0.198	-0.0964203\\
0.2	-0.0970677\\
0.202	-0.0977024\\
0.204	-0.0983243\\
0.206	-0.0989332\\
0.208	-0.099529\\
0.21	-0.100111\\
0.212	-0.10068\\
0.214	-0.101236\\
0.216	-0.101777\\
0.218	-0.102305\\
0.22	-0.102819\\
0.222	-0.103319\\
0.224	-0.103805\\
0.226	-0.104276\\
0.228	-0.104733\\
0.23	-0.105176\\
0.232	-0.105605\\
0.234	-0.106019\\
0.236	-0.106419\\
0.238	-0.106805\\
0.24	-0.107176\\
0.242	-0.107533\\
0.244	-0.107876\\
0.246	-0.108204\\
0.248	-0.108519\\
0.25	-0.108819\\
0.252	-0.109105\\
0.254	-0.109377\\
0.256	-0.109636\\
0.258	-0.109881\\
0.26	-0.110112\\
0.262	-0.11033\\
0.264	-0.110535\\
0.266	-0.110727\\
0.268	-0.110906\\
0.27	-0.111073\\
0.272	-0.111227\\
0.274	-0.111368\\
0.276	-0.111498\\
0.278	-0.111617\\
0.28	-0.111723\\
0.282	-0.111819\\
0.284	-0.111904\\
0.286	-0.111978\\
0.288	-0.112041\\
0.29	-0.112095\\
0.292	-0.112139\\
0.294	-0.112174\\
0.296	-0.112199\\
0.298	-0.112216\\
0.3	-0.112224\\
0.302	-0.112225\\
0.304	-0.112217\\
0.306	-0.112202\\
0.308	-0.112181\\
0.31	-0.112152\\
0.312	-0.112117\\
0.314	-0.112077\\
0.316	-0.11203\\
0.318	-0.111979\\
0.32	-0.111923\\
0.322	-0.111862\\
0.324	-0.111798\\
0.326	-0.111729\\
0.328	-0.111658\\
0.33	-0.111583\\
0.332	-0.111506\\
0.334	-0.111427\\
0.336	-0.111345\\
0.338	-0.111263\\
0.34	-0.111179\\
0.342	-0.111094\\
0.344	-0.11101\\
0.346	-0.110925\\
0.348	-0.11084\\
0.35	-0.110756\\
0.352	-0.110673\\
0.354	-0.110591\\
0.356	-0.11051\\
0.358	-0.110432\\
0.36	-0.110355\\
0.362	-0.110281\\
0.364	-0.11021\\
0.366	-0.110142\\
0.368	-0.110077\\
0.37	-0.110016\\
0.372	-0.109958\\
0.374	-0.109904\\
0.376	-0.109855\\
0.378	-0.10981\\
0.38	-0.109769\\
0.382	-0.109734\\
0.384	-0.109703\\
0.386	-0.109678\\
0.388	-0.109657\\
0.39	-0.109643\\
0.392	-0.109634\\
0.394	-0.10963\\
0.396	-0.109633\\
0.398	-0.109641\\
0.4	-0.109655\\
0.402	-0.109676\\
0.404	-0.109702\\
0.406	-0.109735\\
0.408	-0.109774\\
0.41	-0.10982\\
0.412	-0.109871\\
0.414	-0.109929\\
0.416	-0.109994\\
0.418	-0.110065\\
0.42	-0.110142\\
0.422	-0.110225\\
0.424	-0.110315\\
0.426	-0.110411\\
0.428	-0.110513\\
0.43	-0.110621\\
0.432	-0.110735\\
0.434	-0.110855\\
0.436	-0.110981\\
0.438	-0.111112\\
0.44	-0.11125\\
0.442	-0.111393\\
0.444	-0.111541\\
0.446	-0.111695\\
0.448	-0.111854\\
0.45	-0.112017\\
0.452	-0.112186\\
0.454	-0.11236\\
0.456	-0.112538\\
0.458	-0.11272\\
0.46	-0.112907\\
0.462	-0.113098\\
0.464	-0.113293\\
0.466	-0.113492\\
0.468	-0.113695\\
0.47	-0.113901\\
0.472	-0.11411\\
0.474	-0.114322\\
0.476	-0.114537\\
0.478	-0.114755\\
0.48	-0.114976\\
0.482	-0.115198\\
0.484	-0.115424\\
0.486	-0.115651\\
0.488	-0.115879\\
0.49	-0.11611\\
0.492	-0.116342\\
0.494	-0.116575\\
0.496	-0.11681\\
0.498	-0.117045\\
0.5	-0.117281\\
0.502	-0.117518\\
0.504	-0.117755\\
0.506	-0.117992\\
0.508	-0.11823\\
0.51	-0.118467\\
0.512	-0.118704\\
0.514	-0.118941\\
0.516	-0.119177\\
0.518	-0.119412\\
0.52	-0.119647\\
0.522	-0.11988\\
0.524	-0.120113\\
0.526	-0.120344\\
0.528	-0.120573\\
0.53	-0.120801\\
0.532	-0.121028\\
0.534	-0.121252\\
0.536	-0.121475\\
0.538	-0.121695\\
0.54	-0.121914\\
0.542	-0.12213\\
0.544	-0.122344\\
0.546	-0.122555\\
0.548	-0.122763\\
0.55	-0.122969\\
0.552	-0.123173\\
0.554	-0.123373\\
0.556	-0.123571\\
0.558	-0.123765\\
0.56	-0.123956\\
0.562	-0.124145\\
0.564	-0.12433\\
0.566	-0.124511\\
0.568	-0.12469\\
0.57	-0.124865\\
0.572	-0.125036\\
0.574	-0.125204\\
0.576	-0.125368\\
0.578	-0.125529\\
0.58	-0.125686\\
0.582	-0.12584\\
0.584	-0.12599\\
0.586	-0.126136\\
0.588	-0.126278\\
0.59	-0.126416\\
0.592	-0.126551\\
0.594	-0.126682\\
0.596	-0.126808\\
0.598	-0.126931\\
0.6	-0.12705\\
0.602	-0.127166\\
0.604	-0.127277\\
0.606	-0.127384\\
0.608	-0.127487\\
0.61	-0.127587\\
0.612	-0.127682\\
0.614	-0.127774\\
0.616	-0.127862\\
0.618	-0.127945\\
0.62	-0.128025\\
0.622	-0.128101\\
0.624	-0.128173\\
0.626	-0.128241\\
0.628	-0.128305\\
0.63	-0.128366\\
0.632	-0.128423\\
0.634	-0.128475\\
0.636	-0.128524\\
0.638	-0.12857\\
0.64	-0.128612\\
0.642	-0.12865\\
0.644	-0.128684\\
0.646	-0.128715\\
0.648	-0.128742\\
0.65	-0.128765\\
0.652	-0.128786\\
0.654	-0.128802\\
0.656	-0.128816\\
0.658	-0.128825\\
0.66	-0.128832\\
0.662	-0.128836\\
0.664	-0.128836\\
0.666	-0.128833\\
0.668	-0.128827\\
0.67	-0.128818\\
0.672	-0.128805\\
0.674	-0.12879\\
0.676	-0.128773\\
0.678	-0.128752\\
0.68	-0.128728\\
0.682	-0.128702\\
0.684	-0.128674\\
0.686	-0.128643\\
0.688	-0.128609\\
0.69	-0.128573\\
0.692	-0.128535\\
0.694	-0.128494\\
0.696	-0.128452\\
0.698	-0.128407\\
0.7	-0.12836\\
0.702	-0.128312\\
0.704	-0.128261\\
0.706	-0.128209\\
0.708	-0.128155\\
0.71	-0.1281\\
0.712	-0.128043\\
0.714	-0.127985\\
0.716	-0.127926\\
0.718	-0.127865\\
0.72	-0.127803\\
0.722	-0.12774\\
0.724	-0.127676\\
0.726	-0.127612\\
0.728	-0.127546\\
0.73	-0.12748\\
0.732	-0.127414\\
0.734	-0.127347\\
0.736	-0.127279\\
0.738	-0.127211\\
0.74	-0.127143\\
0.742	-0.127075\\
0.744	-0.127007\\
0.746	-0.126939\\
0.748	-0.126871\\
0.75	-0.126803\\
0.752	-0.126735\\
0.754	-0.126668\\
0.756	-0.126601\\
0.758	-0.126535\\
0.76	-0.12647\\
0.762	-0.126405\\
0.764	-0.126341\\
0.766	-0.126277\\
0.768	-0.126215\\
0.77	-0.126154\\
0.772	-0.126093\\
0.774	-0.126034\\
0.776	-0.125976\\
0.778	-0.125919\\
0.78	-0.125863\\
0.782	-0.125809\\
0.784	-0.125756\\
0.786	-0.125704\\
0.788	-0.125654\\
0.79	-0.125605\\
0.792	-0.125558\\
0.794	-0.125513\\
0.796	-0.125469\\
0.798	-0.125426\\
0.8	-0.125386\\
0.802	-0.125347\\
0.804	-0.12531\\
0.806	-0.125274\\
0.808	-0.12524\\
0.81	-0.125208\\
0.812	-0.125178\\
0.814	-0.125149\\
0.816	-0.125122\\
0.818	-0.125097\\
0.82	-0.125074\\
0.822	-0.125052\\
0.824	-0.125032\\
0.826	-0.125014\\
0.828	-0.124997\\
0.83	-0.124982\\
0.832	-0.124969\\
0.834	-0.124957\\
0.836	-0.124947\\
0.838	-0.124938\\
0.84	-0.124931\\
0.842	-0.124926\\
0.844	-0.124922\\
0.846	-0.124919\\
0.848	-0.124918\\
0.85	-0.124918\\
0.852	-0.124919\\
0.854	-0.124921\\
0.856	-0.124925\\
0.858	-0.12493\\
0.86	-0.124936\\
0.862	-0.124943\\
0.864	-0.124951\\
0.866	-0.12496\\
0.868	-0.124969\\
0.87	-0.12498\\
0.872	-0.124991\\
0.874	-0.125003\\
0.876	-0.125016\\
0.878	-0.125029\\
0.88	-0.125043\\
0.882	-0.125057\\
0.884	-0.125072\\
0.886	-0.125087\\
0.888	-0.125102\\
0.89	-0.125118\\
0.892	-0.125133\\
0.894	-0.125149\\
0.896	-0.125165\\
0.898	-0.125181\\
0.9	-0.125197\\
0.902	-0.125213\\
0.904	-0.125229\\
0.906	-0.125245\\
0.908	-0.12526\\
0.91	-0.125276\\
0.912	-0.125291\\
0.914	-0.125305\\
0.916	-0.125319\\
0.918	-0.125333\\
0.92	-0.125347\\
0.922	-0.125359\\
0.924	-0.125372\\
0.926	-0.125384\\
0.928	-0.125395\\
0.93	-0.125405\\
0.932	-0.125415\\
0.934	-0.125425\\
0.936	-0.125433\\
0.938	-0.125441\\
0.94	-0.125449\\
0.942	-0.125455\\
0.944	-0.125461\\
0.946	-0.125466\\
0.948	-0.12547\\
0.95	-0.125474\\
0.952	-0.125476\\
0.954	-0.125478\\
0.956	-0.12548\\
0.958	-0.12548\\
0.96	-0.12548\\
0.962	-0.125479\\
0.964	-0.125477\\
0.966	-0.125474\\
0.968	-0.125471\\
0.97	-0.125467\\
0.972	-0.125462\\
0.974	-0.125456\\
0.976	-0.12545\\
0.978	-0.125443\\
0.98	-0.125436\\
0.982	-0.125428\\
0.984	-0.125419\\
0.986	-0.125409\\
0.988	-0.125399\\
0.99	-0.125389\\
0.992	-0.125378\\
0.994	-0.125366\\
0.996	-0.125354\\
0.998	-0.125342\\
1	-0.125329\\
1.002	-0.125315\\
1.004	-0.125301\\
1.006	-0.125287\\
1.008	-0.125273\\
1.01	-0.125258\\
1.012	-0.125243\\
1.014	-0.125228\\
1.016	-0.125212\\
1.018	-0.125196\\
1.02	-0.12518\\
1.022	-0.125164\\
1.024	-0.125148\\
1.026	-0.125132\\
1.028	-0.125115\\
1.03	-0.125099\\
1.032	-0.125082\\
1.034	-0.125066\\
1.036	-0.125049\\
1.038	-0.125033\\
1.04	-0.125016\\
1.042	-0.125\\
1.044	-0.124984\\
1.046	-0.124968\\
1.048	-0.124952\\
1.05	-0.124936\\
1.052	-0.12492\\
1.054	-0.124905\\
1.056	-0.124889\\
1.058	-0.124874\\
1.06	-0.124859\\
1.062	-0.124845\\
1.064	-0.12483\\
1.066	-0.124816\\
1.068	-0.124802\\
1.07	-0.124789\\
1.072	-0.124775\\
1.074	-0.124762\\
1.076	-0.12475\\
1.078	-0.124737\\
1.08	-0.124725\\
1.082	-0.124714\\
1.084	-0.124702\\
1.086	-0.124691\\
1.088	-0.12468\\
1.09	-0.12467\\
1.092	-0.12466\\
1.094	-0.12465\\
1.096	-0.124641\\
1.098	-0.124631\\
1.1	-0.124623\\
1.102	-0.124614\\
1.104	-0.124606\\
1.106	-0.124598\\
1.108	-0.124591\\
1.11	-0.124584\\
1.112	-0.124577\\
1.114	-0.12457\\
1.116	-0.124564\\
1.118	-0.124558\\
1.12	-0.124553\\
1.122	-0.124547\\
1.124	-0.124543\\
1.126	-0.124538\\
1.128	-0.124534\\
1.13	-0.124529\\
1.132	-0.124526\\
1.134	-0.124522\\
1.136	-0.124519\\
1.138	-0.124516\\
1.14	-0.124513\\
1.142	-0.124511\\
1.144	-0.124509\\
1.146	-0.124507\\
1.148	-0.124505\\
1.15	-0.124504\\
1.152	-0.124503\\
1.154	-0.124502\\
1.156	-0.124501\\
1.158	-0.124501\\
1.16	-0.124501\\
1.162	-0.124501\\
1.164	-0.124501\\
1.166	-0.124502\\
1.168	-0.124502\\
1.17	-0.124503\\
1.172	-0.124505\\
1.174	-0.124506\\
1.176	-0.124508\\
1.178	-0.12451\\
1.18	-0.124512\\
1.182	-0.124514\\
1.184	-0.124516\\
1.186	-0.124519\\
1.188	-0.124522\\
1.19	-0.124525\\
1.192	-0.124528\\
1.194	-0.124532\\
1.196	-0.124535\\
1.198	-0.124539\\
1.2	-0.124543\\
1.202	-0.124548\\
1.204	-0.124552\\
1.206	-0.124557\\
1.208	-0.124562\\
1.21	-0.124567\\
1.212	-0.124572\\
1.214	-0.124577\\
1.216	-0.124583\\
1.218	-0.124588\\
1.22	-0.124594\\
1.222	-0.1246\\
1.224	-0.124606\\
1.226	-0.124613\\
1.228	-0.124619\\
1.23	-0.124626\\
1.232	-0.124632\\
1.234	-0.124639\\
1.236	-0.124646\\
1.238	-0.124653\\
1.24	-0.12466\\
1.242	-0.124668\\
1.244	-0.124675\\
1.246	-0.124683\\
1.248	-0.12469\\
1.25	-0.124698\\
1.252	-0.124706\\
1.254	-0.124713\\
1.256	-0.124721\\
1.258	-0.124729\\
1.26	-0.124737\\
1.262	-0.124745\\
1.264	-0.124753\\
1.266	-0.124761\\
1.268	-0.124768\\
1.27	-0.124776\\
1.272	-0.124784\\
1.274	-0.124792\\
1.276	-0.1248\\
1.278	-0.124807\\
1.28	-0.124815\\
1.282	-0.124823\\
1.284	-0.12483\\
1.286	-0.124838\\
1.288	-0.124845\\
1.29	-0.124852\\
1.292	-0.124859\\
1.294	-0.124866\\
1.296	-0.124872\\
1.298	-0.124879\\
1.3	-0.124885\\
1.302	-0.124892\\
1.304	-0.124898\\
1.306	-0.124903\\
1.308	-0.124909\\
1.31	-0.124914\\
1.312	-0.124919\\
1.314	-0.124924\\
1.316	-0.124929\\
1.318	-0.124933\\
1.32	-0.124937\\
1.322	-0.124941\\
1.324	-0.124945\\
1.326	-0.124948\\
1.328	-0.124951\\
1.33	-0.124954\\
1.332	-0.124957\\
1.334	-0.124959\\
1.336	-0.124961\\
1.338	-0.124962\\
1.34	-0.124963\\
1.342	-0.124964\\
1.344	-0.124965\\
1.346	-0.124965\\
1.348	-0.124965\\
1.35	-0.124965\\
1.352	-0.124965\\
1.354	-0.124964\\
1.356	-0.124963\\
1.358	-0.124961\\
1.36	-0.124959\\
1.362	-0.124957\\
1.364	-0.124955\\
1.366	-0.124952\\
1.368	-0.124949\\
1.37	-0.124946\\
1.372	-0.124943\\
1.374	-0.124939\\
1.376	-0.124935\\
1.378	-0.124931\\
1.38	-0.124926\\
1.382	-0.124921\\
1.384	-0.124916\\
1.386	-0.124911\\
1.388	-0.124906\\
1.39	-0.1249\\
1.392	-0.124895\\
1.394	-0.124889\\
1.396	-0.124882\\
1.398	-0.124876\\
1.4	-0.12487\\
1.402	-0.124863\\
1.404	-0.124857\\
1.406	-0.12485\\
1.408	-0.124843\\
1.41	-0.124836\\
1.412	-0.124829\\
1.414	-0.124822\\
1.416	-0.124815\\
1.418	-0.124808\\
1.42	-0.1248\\
1.422	-0.124793\\
1.424	-0.124786\\
1.426	-0.124778\\
1.428	-0.124771\\
1.43	-0.124764\\
1.432	-0.124757\\
1.434	-0.12475\\
1.436	-0.124742\\
1.438	-0.124735\\
1.44	-0.124728\\
1.442	-0.124721\\
1.444	-0.124715\\
1.446	-0.124708\\
1.448	-0.124701\\
1.45	-0.124695\\
1.452	-0.124688\\
1.454	-0.124682\\
1.456	-0.124676\\
1.458	-0.12467\\
1.46	-0.124664\\
1.462	-0.124658\\
1.464	-0.124652\\
1.466	-0.124647\\
1.468	-0.124642\\
1.47	-0.124637\\
1.472	-0.124632\\
1.474	-0.124627\\
1.476	-0.124622\\
1.478	-0.124618\\
1.48	-0.124614\\
1.482	-0.12461\\
1.484	-0.124606\\
1.486	-0.124602\\
1.488	-0.124598\\
1.49	-0.124595\\
1.492	-0.124592\\
1.494	-0.124589\\
1.496	-0.124586\\
1.498	-0.124583\\
1.5	-0.124581\\
1.502	-0.124578\\
1.504	-0.124576\\
1.506	-0.124574\\
1.508	-0.124572\\
1.51	-0.12457\\
1.512	-0.124569\\
1.514	-0.124567\\
1.516	-0.124566\\
1.518	-0.124564\\
1.52	-0.124563\\
1.522	-0.124562\\
1.524	-0.124561\\
1.526	-0.124561\\
1.528	-0.12456\\
1.53	-0.124559\\
1.532	-0.124559\\
1.534	-0.124558\\
1.536	-0.124558\\
1.538	-0.124558\\
1.54	-0.124557\\
1.542	-0.124557\\
1.544	-0.124557\\
1.546	-0.124557\\
1.548	-0.124557\\
1.55	-0.124557\\
1.552	-0.124557\\
1.554	-0.124557\\
1.556	-0.124557\\
1.558	-0.124557\\
1.56	-0.124557\\
1.562	-0.124557\\
1.564	-0.124557\\
1.566	-0.124556\\
1.568	-0.124556\\
1.57	-0.124556\\
1.572	-0.124556\\
1.574	-0.124556\\
1.576	-0.124556\\
1.578	-0.124556\\
1.58	-0.124555\\
1.582	-0.124555\\
1.584	-0.124555\\
1.586	-0.124554\\
1.588	-0.124554\\
1.59	-0.124553\\
1.592	-0.124552\\
1.594	-0.124552\\
1.596	-0.124551\\
1.598	-0.12455\\
1.6	-0.124549\\
1.602	-0.124548\\
1.604	-0.124547\\
1.606	-0.124546\\
1.608	-0.124545\\
1.61	-0.124544\\
1.612	-0.124543\\
1.614	-0.124541\\
1.616	-0.12454\\
1.618	-0.124538\\
1.62	-0.124537\\
1.622	-0.124535\\
1.624	-0.124533\\
1.626	-0.124531\\
1.628	-0.124529\\
1.63	-0.124528\\
1.632	-0.124526\\
1.634	-0.124523\\
1.636	-0.124521\\
1.638	-0.124519\\
1.64	-0.124517\\
1.642	-0.124515\\
1.644	-0.124512\\
1.646	-0.12451\\
1.648	-0.124508\\
1.65	-0.124505\\
1.652	-0.124503\\
1.654	-0.1245\\
1.656	-0.124497\\
1.658	-0.124495\\
1.66	-0.124492\\
1.662	-0.12449\\
1.664	-0.124487\\
1.666	-0.124484\\
1.668	-0.124481\\
1.67	-0.124479\\
1.672	-0.124476\\
1.674	-0.124473\\
1.676	-0.12447\\
1.678	-0.124467\\
1.68	-0.124465\\
1.682	-0.124462\\
1.684	-0.124459\\
1.686	-0.124456\\
1.688	-0.124453\\
1.69	-0.12445\\
1.692	-0.124447\\
1.694	-0.124444\\
1.696	-0.124442\\
1.698	-0.124439\\
1.7	-0.124436\\
1.702	-0.124433\\
1.704	-0.12443\\
1.706	-0.124427\\
1.708	-0.124424\\
1.71	-0.124422\\
1.712	-0.124419\\
1.714	-0.124416\\
1.716	-0.124413\\
1.718	-0.12441\\
1.72	-0.124407\\
1.722	-0.124405\\
1.724	-0.124402\\
1.726	-0.124399\\
1.728	-0.124396\\
1.73	-0.124393\\
1.732	-0.124391\\
1.734	-0.124388\\
1.736	-0.124385\\
1.738	-0.124382\\
1.74	-0.12438\\
1.742	-0.124377\\
1.744	-0.124374\\
1.746	-0.124371\\
1.748	-0.124369\\
1.75	-0.124366\\
1.752	-0.124363\\
1.754	-0.124361\\
1.756	-0.124358\\
1.758	-0.124355\\
1.76	-0.124352\\
1.762	-0.12435\\
1.764	-0.124347\\
1.766	-0.124344\\
1.768	-0.124342\\
1.77	-0.124339\\
1.772	-0.124336\\
1.774	-0.124333\\
1.776	-0.124331\\
1.778	-0.124328\\
1.78	-0.124325\\
1.782	-0.124322\\
1.784	-0.12432\\
1.786	-0.124317\\
1.788	-0.124314\\
1.79	-0.124312\\
1.792	-0.124309\\
1.794	-0.124306\\
1.796	-0.124303\\
1.798	-0.124301\\
1.8	-0.124298\\
1.802	-0.124295\\
1.804	-0.124293\\
1.806	-0.12429\\
1.808	-0.124287\\
1.81	-0.124285\\
1.812	-0.124282\\
1.814	-0.124279\\
1.816	-0.124277\\
1.818	-0.124274\\
1.82	-0.124272\\
1.822	-0.124269\\
1.824	-0.124266\\
1.826	-0.124264\\
1.828	-0.124261\\
1.83	-0.124259\\
1.832	-0.124256\\
1.834	-0.124254\\
1.836	-0.124251\\
1.838	-0.124249\\
1.84	-0.124247\\
1.842	-0.124244\\
1.844	-0.124242\\
1.846	-0.12424\\
1.848	-0.124237\\
1.85	-0.124235\\
1.852	-0.124233\\
1.854	-0.124231\\
1.856	-0.124229\\
1.858	-0.124226\\
1.86	-0.124224\\
1.862	-0.124222\\
1.864	-0.12422\\
1.866	-0.124218\\
1.868	-0.124217\\
1.87	-0.124215\\
1.872	-0.124213\\
1.874	-0.124211\\
1.876	-0.124209\\
1.878	-0.124208\\
1.88	-0.124206\\
1.882	-0.124204\\
1.884	-0.124203\\
1.886	-0.124201\\
1.888	-0.1242\\
1.89	-0.124199\\
1.892	-0.124197\\
1.894	-0.124196\\
1.896	-0.124194\\
1.898	-0.124193\\
1.9	-0.124192\\
1.902	-0.124191\\
1.904	-0.12419\\
1.906	-0.124189\\
1.908	-0.124188\\
1.91	-0.124187\\
1.912	-0.124186\\
1.914	-0.124185\\
1.916	-0.124184\\
1.918	-0.124183\\
1.92	-0.124182\\
1.922	-0.124181\\
1.924	-0.12418\\
1.926	-0.124179\\
1.928	-0.124179\\
1.93	-0.124178\\
1.932	-0.124177\\
1.934	-0.124176\\
1.936	-0.124176\\
1.938	-0.124175\\
1.94	-0.124174\\
1.942	-0.124174\\
1.944	-0.124173\\
1.946	-0.124172\\
1.948	-0.124172\\
1.95	-0.124171\\
1.952	-0.12417\\
1.954	-0.12417\\
1.956	-0.124169\\
1.958	-0.124168\\
1.96	-0.124168\\
1.962	-0.124167\\
1.964	-0.124166\\
1.966	-0.124165\\
1.968	-0.124165\\
1.97	-0.124164\\
1.972	-0.124163\\
1.974	-0.124162\\
1.976	-0.124161\\
1.978	-0.12416\\
1.98	-0.12416\\
1.982	-0.124159\\
1.984	-0.124158\\
1.986	-0.124157\\
1.988	-0.124156\\
1.99	-0.124155\\
1.992	-0.124153\\
1.994	-0.124152\\
1.996	-0.124151\\
1.998	-0.12415\\
2	-0.124149\\
};
\addplot [color=mycolor5, line width=1.0pt, forget plot]
  table[row sep=crcr]{%
-0.124	0\\
-0.122	0\\
-0.12	0\\
-0.118	0\\
-0.116	0\\
-0.114	0\\
-0.112	0\\
-0.11	0\\
-0.108	0\\
-0.106	0\\
-0.104	0\\
-0.102	0\\
-0.1	0\\
-0.098	0\\
-0.096	0\\
-0.094	0\\
-0.092	0\\
-0.09	0\\
-0.088	0\\
-0.086	0\\
-0.084	0\\
-0.082	0\\
-0.08	0\\
-0.078	0\\
-0.076	0\\
-0.074	0\\
-0.072	0\\
-0.07	0\\
-0.068	0\\
-0.066	0\\
-0.064	0\\
-0.062	0\\
-0.06	0\\
-0.058	0\\
-0.056	0\\
-0.054	0\\
-0.052	0\\
-0.05	0\\
-0.048	0\\
-0.046	0\\
-0.044	0\\
-0.042	0\\
-0.04	0\\
-0.038	0\\
-0.036	0\\
-0.034	0\\
-0.032	0\\
-0.03	0\\
-0.028	0\\
-0.026	0\\
-0.024	0\\
-0.022	0\\
-0.02	0\\
-0.018	0\\
-0.016	0\\
-0.014	0\\
-0.012	0\\
-0.01	0\\
-0.008	0\\
-0.006	0\\
-0.004	0\\
-0.002	0\\
0	0\\
0.002	-0.000417186\\
0.004	-0.000831266\\
0.006	-0.001244\\
0.008	-0.00165701\\
0.01	-0.00207181\\
0.012	-0.00248988\\
0.014	-0.00291259\\
0.016	-0.00334131\\
0.018	-0.00377735\\
0.02	-0.00422199\\
0.022	-0.00467646\\
0.024	-0.00514198\\
0.026	-0.00561971\\
0.028	-0.00611077\\
0.03	-0.00661623\\
0.032	-0.00713713\\
0.034	-0.00767445\\
0.036	-0.00822914\\
0.038	-0.00880207\\
0.04	-0.0093941\\
0.042	-0.010006\\
0.044	-0.0106385\\
0.046	-0.0112923\\
0.048	-0.011968\\
0.05	-0.0126663\\
0.052	-0.0133875\\
0.054	-0.0141323\\
0.056	-0.0149011\\
0.058	-0.0156941\\
0.06	-0.0165118\\
0.062	-0.0173544\\
0.064	-0.0182222\\
0.066	-0.0191152\\
0.068	-0.0200337\\
0.07	-0.0209777\\
0.072	-0.0219473\\
0.074	-0.0229424\\
0.076	-0.0239629\\
0.078	-0.0250087\\
0.08	-0.0260798\\
0.082	-0.0271758\\
0.084	-0.0282965\\
0.086	-0.0294417\\
0.088	-0.030611\\
0.09	-0.0318039\\
0.092	-0.0330202\\
0.094	-0.0342593\\
0.096	-0.0355208\\
0.098	-0.0368041\\
0.1	-0.0381086\\
0.102	-0.0394338\\
0.104	-0.0407789\\
0.106	-0.0421433\\
0.108	-0.0435264\\
0.11	-0.0449273\\
0.112	-0.0463454\\
0.114	-0.0477797\\
0.116	-0.0492296\\
0.118	-0.0506941\\
0.12	-0.0521723\\
0.122	-0.0536634\\
0.124	-0.0551665\\
0.126	-0.0566806\\
0.128	-0.0582048\\
0.13	-0.059738\\
0.132	-0.0612794\\
0.134	-0.0628278\\
0.136	-0.0643824\\
0.138	-0.065942\\
0.14	-0.0675056\\
0.142	-0.0690722\\
0.144	-0.0706407\\
0.146	-0.07221\\
0.148	-0.0737792\\
0.15	-0.0753471\\
0.152	-0.0769127\\
0.154	-0.0784749\\
0.156	-0.0800326\\
0.158	-0.0815848\\
0.16	-0.0831304\\
0.162	-0.0846684\\
0.164	-0.0861978\\
0.166	-0.0877175\\
0.168	-0.0892265\\
0.17	-0.0907237\\
0.172	-0.0922082\\
0.174	-0.0936791\\
0.176	-0.0951353\\
0.178	-0.0965759\\
0.18	-0.0979999\\
0.182	-0.0994066\\
0.184	-0.100795\\
0.186	-0.102164\\
0.188	-0.103513\\
0.19	-0.104841\\
0.192	-0.106148\\
0.194	-0.107432\\
0.196	-0.108693\\
0.198	-0.10993\\
0.2	-0.111143\\
0.202	-0.11233\\
0.204	-0.113491\\
0.206	-0.114626\\
0.208	-0.115735\\
0.21	-0.116815\\
0.212	-0.117868\\
0.214	-0.118892\\
0.216	-0.119887\\
0.218	-0.120853\\
0.22	-0.121789\\
0.222	-0.122695\\
0.224	-0.123571\\
0.226	-0.124417\\
0.228	-0.125231\\
0.23	-0.126015\\
0.232	-0.126768\\
0.234	-0.127489\\
0.236	-0.128179\\
0.238	-0.128838\\
0.24	-0.129465\\
0.242	-0.130062\\
0.244	-0.130627\\
0.246	-0.131161\\
0.248	-0.131663\\
0.25	-0.132135\\
0.252	-0.132577\\
0.254	-0.132988\\
0.256	-0.133368\\
0.258	-0.133719\\
0.26	-0.134041\\
0.262	-0.134333\\
0.264	-0.134596\\
0.266	-0.134831\\
0.268	-0.135038\\
0.27	-0.135218\\
0.272	-0.13537\\
0.274	-0.135496\\
0.276	-0.135596\\
0.278	-0.135671\\
0.28	-0.135721\\
0.282	-0.135746\\
0.284	-0.135748\\
0.286	-0.135727\\
0.288	-0.135684\\
0.29	-0.135619\\
0.292	-0.135534\\
0.294	-0.135428\\
0.296	-0.135302\\
0.298	-0.135158\\
0.3	-0.134996\\
0.302	-0.134817\\
0.304	-0.134621\\
0.306	-0.134409\\
0.308	-0.134183\\
0.31	-0.133942\\
0.312	-0.133688\\
0.314	-0.133422\\
0.316	-0.133143\\
0.318	-0.132854\\
0.32	-0.132554\\
0.322	-0.132245\\
0.324	-0.131927\\
0.326	-0.131602\\
0.328	-0.131269\\
0.33	-0.13093\\
0.332	-0.130585\\
0.334	-0.130236\\
0.336	-0.129883\\
0.338	-0.129526\\
0.34	-0.129167\\
0.342	-0.128806\\
0.344	-0.128443\\
0.346	-0.12808\\
0.348	-0.127718\\
0.35	-0.127356\\
0.352	-0.126996\\
0.354	-0.126638\\
0.356	-0.126282\\
0.358	-0.12593\\
0.36	-0.125582\\
0.362	-0.125239\\
0.364	-0.1249\\
0.366	-0.124567\\
0.368	-0.12424\\
0.37	-0.12392\\
0.372	-0.123606\\
0.374	-0.1233\\
0.376	-0.123002\\
0.378	-0.122712\\
0.38	-0.122431\\
0.382	-0.122159\\
0.384	-0.121896\\
0.386	-0.121643\\
0.388	-0.1214\\
0.39	-0.121167\\
0.392	-0.120944\\
0.394	-0.120732\\
0.396	-0.120531\\
0.398	-0.120341\\
0.4	-0.120162\\
0.402	-0.119995\\
0.404	-0.119839\\
0.406	-0.119694\\
0.408	-0.119561\\
0.41	-0.11944\\
0.412	-0.119331\\
0.414	-0.119234\\
0.416	-0.119148\\
0.418	-0.119074\\
0.42	-0.119012\\
0.422	-0.118961\\
0.424	-0.118922\\
0.426	-0.118894\\
0.428	-0.118878\\
0.43	-0.118873\\
0.432	-0.118879\\
0.434	-0.118896\\
0.436	-0.118924\\
0.438	-0.118962\\
0.44	-0.119011\\
0.442	-0.11907\\
0.444	-0.119139\\
0.446	-0.119217\\
0.448	-0.119305\\
0.45	-0.119403\\
0.452	-0.119509\\
0.454	-0.119624\\
0.456	-0.119747\\
0.458	-0.119879\\
0.46	-0.120018\\
0.462	-0.120164\\
0.464	-0.120318\\
0.466	-0.120479\\
0.468	-0.120646\\
0.47	-0.120819\\
0.472	-0.120999\\
0.474	-0.121183\\
0.476	-0.121373\\
0.478	-0.121568\\
0.48	-0.121767\\
0.482	-0.121971\\
0.484	-0.122178\\
0.486	-0.122388\\
0.488	-0.122602\\
0.49	-0.122819\\
0.492	-0.123037\\
0.494	-0.123258\\
0.496	-0.123481\\
0.498	-0.123705\\
0.5	-0.12393\\
0.502	-0.124156\\
0.504	-0.124382\\
0.506	-0.124608\\
0.508	-0.124834\\
0.51	-0.125059\\
0.512	-0.125283\\
0.514	-0.125507\\
0.516	-0.125728\\
0.518	-0.125948\\
0.52	-0.126166\\
0.522	-0.126381\\
0.524	-0.126594\\
0.526	-0.126804\\
0.528	-0.12701\\
0.53	-0.127214\\
0.532	-0.127413\\
0.534	-0.127609\\
0.536	-0.1278\\
0.538	-0.127987\\
0.54	-0.12817\\
0.542	-0.128348\\
0.544	-0.128521\\
0.546	-0.128688\\
0.548	-0.128851\\
0.55	-0.129008\\
0.552	-0.129159\\
0.554	-0.129304\\
0.556	-0.129444\\
0.558	-0.129577\\
0.56	-0.129704\\
0.562	-0.129825\\
0.564	-0.12994\\
0.566	-0.130048\\
0.568	-0.130149\\
0.57	-0.130244\\
0.572	-0.130332\\
0.574	-0.130413\\
0.576	-0.130488\\
0.578	-0.130555\\
0.58	-0.130616\\
0.582	-0.13067\\
0.584	-0.130717\\
0.586	-0.130757\\
0.588	-0.13079\\
0.59	-0.130816\\
0.592	-0.130835\\
0.594	-0.130847\\
0.596	-0.130853\\
0.598	-0.130851\\
0.6	-0.130843\\
0.602	-0.130829\\
0.604	-0.130807\\
0.606	-0.130779\\
0.608	-0.130745\\
0.61	-0.130705\\
0.612	-0.130658\\
0.614	-0.130605\\
0.616	-0.130546\\
0.618	-0.130481\\
0.62	-0.13041\\
0.622	-0.130334\\
0.624	-0.130252\\
0.626	-0.130165\\
0.628	-0.130072\\
0.63	-0.129974\\
0.632	-0.129872\\
0.634	-0.129764\\
0.636	-0.129652\\
0.638	-0.129536\\
0.64	-0.129415\\
0.642	-0.12929\\
0.644	-0.129162\\
0.646	-0.129029\\
0.648	-0.128893\\
0.65	-0.128754\\
0.652	-0.128611\\
0.654	-0.128466\\
0.656	-0.128317\\
0.658	-0.128166\\
0.66	-0.128013\\
0.662	-0.127857\\
0.664	-0.1277\\
0.666	-0.127541\\
0.668	-0.12738\\
0.67	-0.127218\\
0.672	-0.127054\\
0.674	-0.12689\\
0.676	-0.126725\\
0.678	-0.126559\\
0.68	-0.126393\\
0.682	-0.126227\\
0.684	-0.126061\\
0.686	-0.125895\\
0.688	-0.125729\\
0.69	-0.125564\\
0.692	-0.1254\\
0.694	-0.125237\\
0.696	-0.125075\\
0.698	-0.124915\\
0.7	-0.124756\\
0.702	-0.124599\\
0.704	-0.124444\\
0.706	-0.124291\\
0.708	-0.12414\\
0.71	-0.123991\\
0.712	-0.123845\\
0.714	-0.123702\\
0.716	-0.123562\\
0.718	-0.123425\\
0.72	-0.123291\\
0.722	-0.12316\\
0.724	-0.123033\\
0.726	-0.122909\\
0.728	-0.122789\\
0.73	-0.122673\\
0.732	-0.122561\\
0.734	-0.122452\\
0.736	-0.122348\\
0.738	-0.122248\\
0.74	-0.122152\\
0.742	-0.122061\\
0.744	-0.121974\\
0.746	-0.121892\\
0.748	-0.121814\\
0.75	-0.121741\\
0.752	-0.121672\\
0.754	-0.121609\\
0.756	-0.12155\\
0.758	-0.121495\\
0.76	-0.121446\\
0.762	-0.121401\\
0.764	-0.121361\\
0.766	-0.121326\\
0.768	-0.121296\\
0.77	-0.121271\\
0.772	-0.12125\\
0.774	-0.121234\\
0.776	-0.121223\\
0.778	-0.121217\\
0.78	-0.121215\\
0.782	-0.121219\\
0.784	-0.121226\\
0.786	-0.121238\\
0.788	-0.121255\\
0.79	-0.121276\\
0.792	-0.121301\\
0.794	-0.121331\\
0.796	-0.121364\\
0.798	-0.121402\\
0.8	-0.121444\\
0.802	-0.121489\\
0.804	-0.121539\\
0.806	-0.121592\\
0.808	-0.121649\\
0.81	-0.121709\\
0.812	-0.121772\\
0.814	-0.121839\\
0.816	-0.121909\\
0.818	-0.121982\\
0.82	-0.122057\\
0.822	-0.122136\\
0.824	-0.122217\\
0.826	-0.1223\\
0.828	-0.122386\\
0.83	-0.122474\\
0.832	-0.122563\\
0.834	-0.122655\\
0.836	-0.122749\\
0.838	-0.122844\\
0.84	-0.12294\\
0.842	-0.123038\\
0.844	-0.123137\\
0.846	-0.123238\\
0.848	-0.123339\\
0.85	-0.12344\\
0.852	-0.123543\\
0.854	-0.123645\\
0.856	-0.123749\\
0.858	-0.123852\\
0.86	-0.123955\\
0.862	-0.124058\\
0.864	-0.124161\\
0.866	-0.124264\\
0.868	-0.124366\\
0.87	-0.124467\\
0.872	-0.124568\\
0.874	-0.124668\\
0.876	-0.124767\\
0.878	-0.124865\\
0.88	-0.124961\\
0.882	-0.125056\\
0.884	-0.12515\\
0.886	-0.125242\\
0.888	-0.125333\\
0.89	-0.125421\\
0.892	-0.125508\\
0.894	-0.125593\\
0.896	-0.125676\\
0.898	-0.125757\\
0.9	-0.125836\\
0.902	-0.125912\\
0.904	-0.125986\\
0.906	-0.126058\\
0.908	-0.126127\\
0.91	-0.126194\\
0.912	-0.126259\\
0.914	-0.12632\\
0.916	-0.126379\\
0.918	-0.126436\\
0.92	-0.126489\\
0.922	-0.12654\\
0.924	-0.126588\\
0.926	-0.126634\\
0.928	-0.126676\\
0.93	-0.126716\\
0.932	-0.126753\\
0.934	-0.126787\\
0.936	-0.126819\\
0.938	-0.126847\\
0.94	-0.126873\\
0.942	-0.126896\\
0.944	-0.126916\\
0.946	-0.126933\\
0.948	-0.126948\\
0.95	-0.12696\\
0.952	-0.126969\\
0.954	-0.126975\\
0.956	-0.126979\\
0.958	-0.126981\\
0.96	-0.126979\\
0.962	-0.126976\\
0.964	-0.12697\\
0.966	-0.126961\\
0.968	-0.12695\\
0.97	-0.126937\\
0.972	-0.126922\\
0.974	-0.126905\\
0.976	-0.126885\\
0.978	-0.126864\\
0.98	-0.12684\\
0.982	-0.126815\\
0.984	-0.126787\\
0.986	-0.126758\\
0.988	-0.126728\\
0.99	-0.126695\\
0.992	-0.126661\\
0.994	-0.126626\\
0.996	-0.126589\\
0.998	-0.126551\\
1	-0.126512\\
1.002	-0.126472\\
1.004	-0.12643\\
1.006	-0.126388\\
1.008	-0.126344\\
1.01	-0.1263\\
1.012	-0.126255\\
1.014	-0.126209\\
1.016	-0.126163\\
1.018	-0.126116\\
1.02	-0.126069\\
1.022	-0.126021\\
1.024	-0.125973\\
1.026	-0.125925\\
1.028	-0.125876\\
1.03	-0.125828\\
1.032	-0.125779\\
1.034	-0.12573\\
1.036	-0.125682\\
1.038	-0.125633\\
1.04	-0.125585\\
1.042	-0.125537\\
1.044	-0.12549\\
1.046	-0.125443\\
1.048	-0.125396\\
1.05	-0.12535\\
1.052	-0.125305\\
1.054	-0.12526\\
1.056	-0.125216\\
1.058	-0.125172\\
1.06	-0.12513\\
1.062	-0.125088\\
1.064	-0.125047\\
1.066	-0.125007\\
1.068	-0.124967\\
1.07	-0.124929\\
1.072	-0.124892\\
1.074	-0.124856\\
1.076	-0.124821\\
1.078	-0.124787\\
1.08	-0.124754\\
1.082	-0.124723\\
1.084	-0.124692\\
1.086	-0.124663\\
1.088	-0.124635\\
1.09	-0.124608\\
1.092	-0.124582\\
1.094	-0.124558\\
1.096	-0.124535\\
1.098	-0.124513\\
1.1	-0.124493\\
1.102	-0.124474\\
1.104	-0.124456\\
1.106	-0.124439\\
1.108	-0.124424\\
1.11	-0.12441\\
1.112	-0.124397\\
1.114	-0.124386\\
1.116	-0.124376\\
1.118	-0.124367\\
1.12	-0.12436\\
1.122	-0.124353\\
1.124	-0.124348\\
1.126	-0.124345\\
1.128	-0.124342\\
1.13	-0.124341\\
1.132	-0.124341\\
1.134	-0.124341\\
1.136	-0.124344\\
1.138	-0.124347\\
1.14	-0.124351\\
1.142	-0.124357\\
1.144	-0.124363\\
1.146	-0.12437\\
1.148	-0.124379\\
1.15	-0.124388\\
1.152	-0.124399\\
1.154	-0.12441\\
1.156	-0.124422\\
1.158	-0.124435\\
1.16	-0.124449\\
1.162	-0.124464\\
1.164	-0.124479\\
1.166	-0.124495\\
1.168	-0.124512\\
1.17	-0.12453\\
1.172	-0.124548\\
1.174	-0.124566\\
1.176	-0.124586\\
1.178	-0.124605\\
1.18	-0.124626\\
1.182	-0.124646\\
1.184	-0.124668\\
1.186	-0.124689\\
1.188	-0.124711\\
1.19	-0.124733\\
1.192	-0.124755\\
1.194	-0.124778\\
1.196	-0.124801\\
1.198	-0.124824\\
1.2	-0.124847\\
1.202	-0.12487\\
1.204	-0.124894\\
1.206	-0.124917\\
1.208	-0.12494\\
1.21	-0.124964\\
1.212	-0.124987\\
1.214	-0.12501\\
1.216	-0.125033\\
1.218	-0.125056\\
1.22	-0.125078\\
1.222	-0.125101\\
1.224	-0.125123\\
1.226	-0.125145\\
1.228	-0.125166\\
1.23	-0.125187\\
1.232	-0.125208\\
1.234	-0.125228\\
1.236	-0.125248\\
1.238	-0.125268\\
1.24	-0.125287\\
1.242	-0.125305\\
1.244	-0.125323\\
1.246	-0.125341\\
1.248	-0.125357\\
1.25	-0.125374\\
1.252	-0.125389\\
1.254	-0.125404\\
1.256	-0.125419\\
1.258	-0.125432\\
1.26	-0.125445\\
1.262	-0.125458\\
1.264	-0.125469\\
1.266	-0.12548\\
1.268	-0.12549\\
1.27	-0.1255\\
1.272	-0.125509\\
1.274	-0.125516\\
1.276	-0.125524\\
1.278	-0.12553\\
1.28	-0.125536\\
1.282	-0.12554\\
1.284	-0.125544\\
1.286	-0.125548\\
1.288	-0.12555\\
1.29	-0.125552\\
1.292	-0.125553\\
1.294	-0.125553\\
1.296	-0.125552\\
1.298	-0.12555\\
1.3	-0.125548\\
1.302	-0.125545\\
1.304	-0.125541\\
1.306	-0.125537\\
1.308	-0.125531\\
1.31	-0.125525\\
1.312	-0.125519\\
1.314	-0.125511\\
1.316	-0.125503\\
1.318	-0.125494\\
1.32	-0.125484\\
1.322	-0.125474\\
1.324	-0.125463\\
1.326	-0.125452\\
1.328	-0.12544\\
1.33	-0.125427\\
1.332	-0.125414\\
1.334	-0.1254\\
1.336	-0.125385\\
1.338	-0.12537\\
1.34	-0.125355\\
1.342	-0.125339\\
1.344	-0.125323\\
1.346	-0.125306\\
1.348	-0.125289\\
1.35	-0.125271\\
1.352	-0.125254\\
1.354	-0.125235\\
1.356	-0.125217\\
1.358	-0.125198\\
1.36	-0.125179\\
1.362	-0.125159\\
1.364	-0.12514\\
1.366	-0.12512\\
1.368	-0.1251\\
1.37	-0.12508\\
1.372	-0.12506\\
1.374	-0.12504\\
1.376	-0.12502\\
1.378	-0.124999\\
1.38	-0.124979\\
1.382	-0.124959\\
1.384	-0.124938\\
1.386	-0.124918\\
1.388	-0.124898\\
1.39	-0.124878\\
1.392	-0.124858\\
1.394	-0.124838\\
1.396	-0.124819\\
1.398	-0.124799\\
1.4	-0.12478\\
1.402	-0.124761\\
1.404	-0.124742\\
1.406	-0.124724\\
1.408	-0.124706\\
1.41	-0.124688\\
1.412	-0.124671\\
1.414	-0.124654\\
1.416	-0.124637\\
1.418	-0.12462\\
1.42	-0.124605\\
1.422	-0.124589\\
1.424	-0.124574\\
1.426	-0.124559\\
1.428	-0.124545\\
1.43	-0.124531\\
1.432	-0.124518\\
1.434	-0.124505\\
1.436	-0.124493\\
1.438	-0.124481\\
1.44	-0.12447\\
1.442	-0.124459\\
1.444	-0.124449\\
1.446	-0.124439\\
1.448	-0.12443\\
1.45	-0.124421\\
1.452	-0.124413\\
1.454	-0.124406\\
1.456	-0.124399\\
1.458	-0.124392\\
1.46	-0.124386\\
1.462	-0.124381\\
1.464	-0.124376\\
1.466	-0.124371\\
1.468	-0.124368\\
1.47	-0.124364\\
1.472	-0.124361\\
1.474	-0.124359\\
1.476	-0.124357\\
1.478	-0.124356\\
1.48	-0.124355\\
1.482	-0.124355\\
1.484	-0.124355\\
1.486	-0.124356\\
1.488	-0.124357\\
1.49	-0.124358\\
1.492	-0.12436\\
1.494	-0.124363\\
1.496	-0.124365\\
1.498	-0.124369\\
1.5	-0.124372\\
1.502	-0.124376\\
1.504	-0.12438\\
1.506	-0.124385\\
1.508	-0.12439\\
1.51	-0.124395\\
1.512	-0.124401\\
1.514	-0.124407\\
1.516	-0.124413\\
1.518	-0.124419\\
1.52	-0.124426\\
1.522	-0.124432\\
1.524	-0.12444\\
1.526	-0.124447\\
1.528	-0.124454\\
1.53	-0.124462\\
1.532	-0.124469\\
1.534	-0.124477\\
1.536	-0.124485\\
1.538	-0.124493\\
1.54	-0.124501\\
1.542	-0.124509\\
1.544	-0.124517\\
1.546	-0.124526\\
1.548	-0.124534\\
1.55	-0.124542\\
1.552	-0.12455\\
1.554	-0.124559\\
1.556	-0.124567\\
1.558	-0.124575\\
1.56	-0.124583\\
1.562	-0.124591\\
1.564	-0.124599\\
1.566	-0.124607\\
1.568	-0.124614\\
1.57	-0.124622\\
1.572	-0.124629\\
1.574	-0.124637\\
1.576	-0.124644\\
1.578	-0.124651\\
1.58	-0.124658\\
1.582	-0.124664\\
1.584	-0.124671\\
1.586	-0.124677\\
1.588	-0.124683\\
1.59	-0.124688\\
1.592	-0.124694\\
1.594	-0.124699\\
1.596	-0.124704\\
1.598	-0.124709\\
1.6	-0.124714\\
1.602	-0.124718\\
1.604	-0.124722\\
1.606	-0.124726\\
1.608	-0.124729\\
1.61	-0.124733\\
1.612	-0.124736\\
1.614	-0.124738\\
1.616	-0.124741\\
1.618	-0.124743\\
1.62	-0.124745\\
1.622	-0.124746\\
1.624	-0.124748\\
1.626	-0.124749\\
1.628	-0.12475\\
1.63	-0.12475\\
1.632	-0.12475\\
1.634	-0.12475\\
1.636	-0.12475\\
1.638	-0.124749\\
1.64	-0.124748\\
1.642	-0.124747\\
1.644	-0.124746\\
1.646	-0.124744\\
1.648	-0.124742\\
1.65	-0.12474\\
1.652	-0.124738\\
1.654	-0.124735\\
1.656	-0.124733\\
1.658	-0.124729\\
1.66	-0.124726\\
1.662	-0.124723\\
1.664	-0.124719\\
1.666	-0.124715\\
1.668	-0.124711\\
1.67	-0.124707\\
1.672	-0.124702\\
1.674	-0.124697\\
1.676	-0.124692\\
1.678	-0.124687\\
1.68	-0.124682\\
1.682	-0.124677\\
1.684	-0.124671\\
1.686	-0.124666\\
1.688	-0.12466\\
1.69	-0.124654\\
1.692	-0.124648\\
1.694	-0.124642\\
1.696	-0.124636\\
1.698	-0.12463\\
1.7	-0.124623\\
1.702	-0.124617\\
1.704	-0.12461\\
1.706	-0.124603\\
1.708	-0.124597\\
1.71	-0.12459\\
1.712	-0.124583\\
1.714	-0.124576\\
1.716	-0.124569\\
1.718	-0.124563\\
1.72	-0.124556\\
1.722	-0.124549\\
1.724	-0.124542\\
1.726	-0.124535\\
1.728	-0.124528\\
1.73	-0.124521\\
1.732	-0.124514\\
1.734	-0.124507\\
1.736	-0.1245\\
1.738	-0.124494\\
1.74	-0.124487\\
1.742	-0.12448\\
1.744	-0.124474\\
1.746	-0.124467\\
1.748	-0.12446\\
1.75	-0.124454\\
1.752	-0.124448\\
1.754	-0.124441\\
1.756	-0.124435\\
1.758	-0.124429\\
1.76	-0.124423\\
1.762	-0.124417\\
1.764	-0.124411\\
1.766	-0.124405\\
1.768	-0.1244\\
1.77	-0.124394\\
1.772	-0.124389\\
1.774	-0.124383\\
1.776	-0.124378\\
1.778	-0.124373\\
1.78	-0.124368\\
1.782	-0.124363\\
1.784	-0.124359\\
1.786	-0.124354\\
1.788	-0.12435\\
1.79	-0.124346\\
1.792	-0.124341\\
1.794	-0.124337\\
1.796	-0.124334\\
1.798	-0.12433\\
1.8	-0.124326\\
1.802	-0.124323\\
1.804	-0.12432\\
1.806	-0.124316\\
1.808	-0.124313\\
1.81	-0.124311\\
1.812	-0.124308\\
1.814	-0.124305\\
1.816	-0.124303\\
1.818	-0.1243\\
1.82	-0.124298\\
1.822	-0.124296\\
1.824	-0.124294\\
1.826	-0.124293\\
1.828	-0.124291\\
1.83	-0.12429\\
1.832	-0.124288\\
1.834	-0.124287\\
1.836	-0.124286\\
1.838	-0.124285\\
1.84	-0.124284\\
1.842	-0.124283\\
1.844	-0.124283\\
1.846	-0.124282\\
1.848	-0.124282\\
1.85	-0.124282\\
1.852	-0.124281\\
1.854	-0.124281\\
1.856	-0.124281\\
1.858	-0.124282\\
1.86	-0.124282\\
1.862	-0.124282\\
1.864	-0.124282\\
1.866	-0.124283\\
1.868	-0.124283\\
1.87	-0.124284\\
1.872	-0.124285\\
1.874	-0.124286\\
1.876	-0.124286\\
1.878	-0.124287\\
1.88	-0.124288\\
1.882	-0.124289\\
1.884	-0.12429\\
1.886	-0.124291\\
1.888	-0.124292\\
1.89	-0.124293\\
1.892	-0.124295\\
1.894	-0.124296\\
1.896	-0.124297\\
1.898	-0.124298\\
1.9	-0.124299\\
1.902	-0.124301\\
1.904	-0.124302\\
1.906	-0.124303\\
1.908	-0.124304\\
1.91	-0.124305\\
1.912	-0.124307\\
1.914	-0.124308\\
1.916	-0.124309\\
1.918	-0.12431\\
1.92	-0.124311\\
1.922	-0.124312\\
1.924	-0.124313\\
1.926	-0.124314\\
1.928	-0.124315\\
1.93	-0.124316\\
1.932	-0.124317\\
1.934	-0.124317\\
1.936	-0.124318\\
1.938	-0.124319\\
1.94	-0.124319\\
1.942	-0.12432\\
1.944	-0.12432\\
1.946	-0.12432\\
1.948	-0.124321\\
1.95	-0.124321\\
1.952	-0.124321\\
1.954	-0.124321\\
1.956	-0.124321\\
1.958	-0.124321\\
1.96	-0.124321\\
1.962	-0.12432\\
1.964	-0.12432\\
1.966	-0.124319\\
1.968	-0.124319\\
1.97	-0.124318\\
1.972	-0.124317\\
1.974	-0.124316\\
1.976	-0.124315\\
1.978	-0.124314\\
1.98	-0.124313\\
1.982	-0.124312\\
1.984	-0.12431\\
1.986	-0.124309\\
1.988	-0.124307\\
1.99	-0.124306\\
1.992	-0.124304\\
1.994	-0.124302\\
1.996	-0.1243\\
1.998	-0.124298\\
2	-0.124296\\
};
\addplot [color=mycolor6, line width=1.0pt, forget plot]
  table[row sep=crcr]{%
-0.124	0\\
-0.122	0\\
-0.12	0\\
-0.118	0\\
-0.116	0\\
-0.114	0\\
-0.112	0\\
-0.11	0\\
-0.108	0\\
-0.106	0\\
-0.104	0\\
-0.102	0\\
-0.1	0\\
-0.098	0\\
-0.096	0\\
-0.094	0\\
-0.092	0\\
-0.09	0\\
-0.088	0\\
-0.086	0\\
-0.084	0\\
-0.082	0\\
-0.08	0\\
-0.078	0\\
-0.076	0\\
-0.074	0\\
-0.072	0\\
-0.07	0\\
-0.068	0\\
-0.066	0\\
-0.064	0\\
-0.062	0\\
-0.06	0\\
-0.058	0\\
-0.056	0\\
-0.054	0\\
-0.052	0\\
-0.05	0\\
-0.048	0\\
-0.046	0\\
-0.044	0\\
-0.042	0\\
-0.04	0\\
-0.038	0\\
-0.036	0\\
-0.034	0\\
-0.032	0\\
-0.03	0\\
-0.028	0\\
-0.026	0\\
-0.024	0\\
-0.022	0\\
-0.02	0\\
-0.018	0\\
-0.016	0\\
-0.014	0\\
-0.012	0\\
-0.01	0\\
-0.008	0\\
-0.006	0\\
-0.004	0\\
-0.002	0\\
0	0\\
0.002	-0.00423308\\
0.004	-0.00842395\\
0.006	-0.0125756\\
0.008	-0.0166899\\
0.01	-0.0207682\\
0.012	-0.0248112\\
0.014	-0.0288197\\
0.016	-0.032794\\
0.018	-0.0367346\\
0.02	-0.0406418\\
0.022	-0.0445159\\
0.024	-0.0483571\\
0.026	-0.0521656\\
0.028	-0.0559413\\
0.03	-0.0596842\\
0.032	-0.063394\\
0.034	-0.0670703\\
0.036	-0.0707126\\
0.038	-0.0743201\\
0.04	-0.0778918\\
0.042	-0.0814267\\
0.044	-0.0849235\\
0.046	-0.0883806\\
0.048	-0.0917963\\
0.05	-0.095169\\
0.052	-0.0984964\\
0.054	-0.101776\\
0.056	-0.105007\\
0.058	-0.108185\\
0.06	-0.111308\\
0.062	-0.114373\\
0.064	-0.117378\\
0.066	-0.12032\\
0.068	-0.123194\\
0.07	-0.125999\\
0.072	-0.128731\\
0.074	-0.131386\\
0.076	-0.133962\\
0.078	-0.136455\\
0.08	-0.138861\\
0.082	-0.141177\\
0.084	-0.1434\\
0.086	-0.145527\\
0.088	-0.147553\\
0.09	-0.149477\\
0.092	-0.151294\\
0.094	-0.153002\\
0.096	-0.154598\\
0.098	-0.156078\\
0.1	-0.157441\\
0.102	-0.158684\\
0.104	-0.159804\\
0.106	-0.160799\\
0.108	-0.161668\\
0.11	-0.162407\\
0.112	-0.163017\\
0.114	-0.163495\\
0.116	-0.163839\\
0.118	-0.164051\\
0.12	-0.164127\\
0.122	-0.164069\\
0.124	-0.163877\\
0.126	-0.163549\\
0.128	-0.163087\\
0.13	-0.162491\\
0.132	-0.161762\\
0.134	-0.160902\\
0.136	-0.159912\\
0.138	-0.158793\\
0.14	-0.157548\\
0.142	-0.15618\\
0.144	-0.154691\\
0.146	-0.153083\\
0.148	-0.15136\\
0.15	-0.149526\\
0.152	-0.147584\\
0.154	-0.145538\\
0.156	-0.143392\\
0.158	-0.141152\\
0.16	-0.138821\\
0.162	-0.136404\\
0.164	-0.133907\\
0.166	-0.131335\\
0.168	-0.128693\\
0.17	-0.125987\\
0.172	-0.123223\\
0.174	-0.120407\\
0.176	-0.117545\\
0.178	-0.114644\\
0.18	-0.111708\\
0.182	-0.108746\\
0.184	-0.105763\\
0.186	-0.102766\\
0.188	-0.0997612\\
0.19	-0.0967556\\
0.192	-0.0937558\\
0.194	-0.0907681\\
0.196	-0.0877993\\
0.198	-0.0848558\\
0.2	-0.081944\\
0.202	-0.0790704\\
0.204	-0.0762413\\
0.206	-0.0734627\\
0.208	-0.0707408\\
0.21	-0.0680814\\
0.212	-0.0654904\\
0.214	-0.0629733\\
0.216	-0.0605356\\
0.218	-0.0581825\\
0.22	-0.0559189\\
0.222	-0.0537498\\
0.224	-0.0516797\\
0.226	-0.0497129\\
0.228	-0.0478535\\
0.23	-0.0461054\\
0.232	-0.044472\\
0.234	-0.0429568\\
0.236	-0.0415627\\
0.238	-0.0402924\\
0.24	-0.0391484\\
0.242	-0.0381328\\
0.244	-0.0372474\\
0.246	-0.0364937\\
0.248	-0.035873\\
0.25	-0.0353862\\
0.252	-0.0350338\\
0.254	-0.0348162\\
0.256	-0.0347334\\
0.258	-0.034785\\
0.26	-0.0349704\\
0.262	-0.0352888\\
0.264	-0.0357388\\
0.266	-0.036319\\
0.268	-0.0370276\\
0.27	-0.0378625\\
0.272	-0.0388215\\
0.274	-0.0399019\\
0.276	-0.0411009\\
0.278	-0.0424154\\
0.28	-0.043842\\
0.282	-0.0453773\\
0.284	-0.0470175\\
0.286	-0.0487586\\
0.288	-0.0505964\\
0.29	-0.0525267\\
0.292	-0.054545\\
0.294	-0.0566466\\
0.296	-0.0588267\\
0.298	-0.0610804\\
0.3	-0.0634027\\
0.302	-0.0657886\\
0.304	-0.0682326\\
0.306	-0.0707297\\
0.308	-0.0732744\\
0.31	-0.0758614\\
0.312	-0.0784853\\
0.314	-0.0811406\\
0.316	-0.083822\\
0.318	-0.0865239\\
0.32	-0.089241\\
0.322	-0.091968\\
0.324	-0.0946995\\
0.326	-0.0974303\\
0.328	-0.100155\\
0.33	-0.102869\\
0.332	-0.105567\\
0.334	-0.108244\\
0.336	-0.110896\\
0.338	-0.113517\\
0.34	-0.116104\\
0.342	-0.118652\\
0.344	-0.121156\\
0.346	-0.123613\\
0.348	-0.126019\\
0.35	-0.128371\\
0.352	-0.130664\\
0.354	-0.132895\\
0.356	-0.135062\\
0.358	-0.137161\\
0.36	-0.139191\\
0.362	-0.141147\\
0.364	-0.143028\\
0.366	-0.144832\\
0.368	-0.146558\\
0.37	-0.148202\\
0.372	-0.149765\\
0.374	-0.151245\\
0.376	-0.15264\\
0.378	-0.15395\\
0.38	-0.155175\\
0.382	-0.156314\\
0.384	-0.157367\\
0.386	-0.158334\\
0.388	-0.159214\\
0.39	-0.16001\\
0.392	-0.16072\\
0.394	-0.161347\\
0.396	-0.16189\\
0.398	-0.162352\\
0.4	-0.162732\\
0.402	-0.163034\\
0.404	-0.163257\\
0.406	-0.163405\\
0.408	-0.163479\\
0.41	-0.16348\\
0.412	-0.163412\\
0.414	-0.163275\\
0.416	-0.163073\\
0.418	-0.162808\\
0.42	-0.162482\\
0.422	-0.162098\\
0.424	-0.161658\\
0.426	-0.161166\\
0.428	-0.160623\\
0.43	-0.160033\\
0.432	-0.159397\\
0.434	-0.15872\\
0.436	-0.158004\\
0.438	-0.157252\\
0.44	-0.156465\\
0.442	-0.155648\\
0.444	-0.154804\\
0.446	-0.153933\\
0.448	-0.153041\\
0.45	-0.152128\\
0.452	-0.151198\\
0.454	-0.150253\\
0.456	-0.149297\\
0.458	-0.14833\\
0.46	-0.147356\\
0.462	-0.146377\\
0.464	-0.145395\\
0.466	-0.144413\\
0.468	-0.143432\\
0.47	-0.142454\\
0.472	-0.141481\\
0.474	-0.140515\\
0.476	-0.139558\\
0.478	-0.138612\\
0.48	-0.137677\\
0.482	-0.136755\\
0.484	-0.135847\\
0.486	-0.134956\\
0.488	-0.13408\\
0.49	-0.133223\\
0.492	-0.132384\\
0.494	-0.131564\\
0.496	-0.130764\\
0.498	-0.129985\\
0.5	-0.129227\\
0.502	-0.128491\\
0.504	-0.127776\\
0.506	-0.127083\\
0.508	-0.126413\\
0.51	-0.125764\\
0.512	-0.125138\\
0.514	-0.124534\\
0.516	-0.123952\\
0.518	-0.123392\\
0.52	-0.122853\\
0.522	-0.122335\\
0.524	-0.121838\\
0.526	-0.121361\\
0.528	-0.120904\\
0.53	-0.120466\\
0.532	-0.120047\\
0.534	-0.119646\\
0.536	-0.119262\\
0.538	-0.118895\\
0.54	-0.118545\\
0.542	-0.11821\\
0.544	-0.117889\\
0.546	-0.117583\\
0.548	-0.11729\\
0.55	-0.11701\\
0.552	-0.116742\\
0.554	-0.116486\\
0.556	-0.11624\\
0.558	-0.116004\\
0.56	-0.115778\\
0.562	-0.115561\\
0.564	-0.115352\\
0.566	-0.115151\\
0.568	-0.114957\\
0.57	-0.11477\\
0.572	-0.114589\\
0.574	-0.114415\\
0.576	-0.114246\\
0.578	-0.114082\\
0.58	-0.113924\\
0.582	-0.11377\\
0.584	-0.113621\\
0.586	-0.113477\\
0.588	-0.113337\\
0.59	-0.113201\\
0.592	-0.113069\\
0.594	-0.112942\\
0.596	-0.112819\\
0.598	-0.112701\\
0.6	-0.112587\\
0.602	-0.112478\\
0.604	-0.112374\\
0.606	-0.112275\\
0.608	-0.112182\\
0.61	-0.112095\\
0.612	-0.112013\\
0.614	-0.111939\\
0.616	-0.111871\\
0.618	-0.111811\\
0.62	-0.111758\\
0.622	-0.111714\\
0.624	-0.111679\\
0.626	-0.111652\\
0.628	-0.111635\\
0.63	-0.111629\\
0.632	-0.111633\\
0.634	-0.111648\\
0.636	-0.111675\\
0.638	-0.111713\\
0.64	-0.111765\\
0.642	-0.111829\\
0.644	-0.111907\\
0.646	-0.111998\\
0.648	-0.112104\\
0.65	-0.112224\\
0.652	-0.11236\\
0.654	-0.11251\\
0.656	-0.112676\\
0.658	-0.112858\\
0.66	-0.113056\\
0.662	-0.11327\\
0.664	-0.1135\\
0.666	-0.113747\\
0.668	-0.11401\\
0.67	-0.11429\\
0.672	-0.114587\\
0.674	-0.114899\\
0.676	-0.115228\\
0.678	-0.115573\\
0.68	-0.115934\\
0.682	-0.116311\\
0.684	-0.116703\\
0.686	-0.11711\\
0.688	-0.117531\\
0.69	-0.117967\\
0.692	-0.118416\\
0.694	-0.118878\\
0.696	-0.119353\\
0.698	-0.119839\\
0.7	-0.120337\\
0.702	-0.120845\\
0.704	-0.121363\\
0.706	-0.121889\\
0.708	-0.122424\\
0.71	-0.122966\\
0.712	-0.123514\\
0.714	-0.124067\\
0.716	-0.124625\\
0.718	-0.125186\\
0.72	-0.125749\\
0.722	-0.126314\\
0.724	-0.126879\\
0.726	-0.127443\\
0.728	-0.128006\\
0.73	-0.128566\\
0.732	-0.129122\\
0.734	-0.129672\\
0.736	-0.130217\\
0.738	-0.130754\\
0.74	-0.131283\\
0.742	-0.131803\\
0.744	-0.132313\\
0.746	-0.132811\\
0.748	-0.133296\\
0.75	-0.133769\\
0.752	-0.134227\\
0.754	-0.134669\\
0.756	-0.135096\\
0.758	-0.135505\\
0.76	-0.135896\\
0.762	-0.136269\\
0.764	-0.136622\\
0.766	-0.136954\\
0.768	-0.137266\\
0.77	-0.137557\\
0.772	-0.137825\\
0.774	-0.13807\\
0.776	-0.138293\\
0.778	-0.138491\\
0.78	-0.138666\\
0.782	-0.138816\\
0.784	-0.138942\\
0.786	-0.139043\\
0.788	-0.139119\\
0.79	-0.139169\\
0.792	-0.139195\\
0.794	-0.139195\\
0.796	-0.139171\\
0.798	-0.139121\\
0.8	-0.139047\\
0.802	-0.138948\\
0.804	-0.138824\\
0.806	-0.138677\\
0.808	-0.138506\\
0.81	-0.138313\\
0.812	-0.138096\\
0.814	-0.137858\\
0.816	-0.137598\\
0.818	-0.137317\\
0.82	-0.137016\\
0.822	-0.136696\\
0.824	-0.136357\\
0.826	-0.136001\\
0.828	-0.135627\\
0.83	-0.135237\\
0.832	-0.134832\\
0.834	-0.134413\\
0.836	-0.13398\\
0.838	-0.133535\\
0.84	-0.133079\\
0.842	-0.132613\\
0.844	-0.132137\\
0.846	-0.131653\\
0.848	-0.131162\\
0.85	-0.130665\\
0.852	-0.130163\\
0.854	-0.129657\\
0.856	-0.129148\\
0.858	-0.128637\\
0.86	-0.128126\\
0.862	-0.127615\\
0.864	-0.127105\\
0.866	-0.126598\\
0.868	-0.126095\\
0.87	-0.125596\\
0.872	-0.125102\\
0.874	-0.124615\\
0.876	-0.124135\\
0.878	-0.123663\\
0.88	-0.123201\\
0.882	-0.122748\\
0.884	-0.122306\\
0.886	-0.121876\\
0.888	-0.121458\\
0.89	-0.121053\\
0.892	-0.120662\\
0.894	-0.120284\\
0.896	-0.119922\\
0.898	-0.119575\\
0.9	-0.119244\\
0.902	-0.118929\\
0.904	-0.118631\\
0.906	-0.11835\\
0.908	-0.118086\\
0.91	-0.11784\\
0.912	-0.117612\\
0.914	-0.117402\\
0.916	-0.11721\\
0.918	-0.117037\\
0.92	-0.116881\\
0.922	-0.116744\\
0.924	-0.116626\\
0.926	-0.116525\\
0.928	-0.116443\\
0.93	-0.116379\\
0.932	-0.116332\\
0.934	-0.116303\\
0.936	-0.116291\\
0.938	-0.116295\\
0.94	-0.116317\\
0.942	-0.116354\\
0.944	-0.116408\\
0.946	-0.116476\\
0.948	-0.116559\\
0.95	-0.116657\\
0.952	-0.116769\\
0.954	-0.116894\\
0.956	-0.117031\\
0.958	-0.117181\\
0.96	-0.117342\\
0.962	-0.117514\\
0.964	-0.117696\\
0.966	-0.117889\\
0.968	-0.11809\\
0.97	-0.118299\\
0.972	-0.118516\\
0.974	-0.118741\\
0.976	-0.118971\\
0.978	-0.119208\\
0.98	-0.119449\\
0.982	-0.119695\\
0.984	-0.119945\\
0.986	-0.120197\\
0.988	-0.120453\\
0.99	-0.12071\\
0.992	-0.120968\\
0.994	-0.121227\\
0.996	-0.121486\\
0.998	-0.121745\\
1	-0.122002\\
1.002	-0.122258\\
1.004	-0.122512\\
1.006	-0.122764\\
1.008	-0.123012\\
1.01	-0.123257\\
1.012	-0.123497\\
1.014	-0.123734\\
1.016	-0.123966\\
1.018	-0.124192\\
1.02	-0.124414\\
1.022	-0.124629\\
1.024	-0.124839\\
1.026	-0.125042\\
1.028	-0.125239\\
1.03	-0.125429\\
1.032	-0.125612\\
1.034	-0.125788\\
1.036	-0.125957\\
1.038	-0.126119\\
1.04	-0.126273\\
1.042	-0.12642\\
1.044	-0.12656\\
1.046	-0.126692\\
1.048	-0.126816\\
1.05	-0.126933\\
1.052	-0.127043\\
1.054	-0.127145\\
1.056	-0.12724\\
1.058	-0.127328\\
1.06	-0.127409\\
1.062	-0.127483\\
1.064	-0.12755\\
1.066	-0.127611\\
1.068	-0.127665\\
1.07	-0.127713\\
1.072	-0.127755\\
1.074	-0.127791\\
1.076	-0.127822\\
1.078	-0.127847\\
1.08	-0.127867\\
1.082	-0.127882\\
1.084	-0.127893\\
1.086	-0.127899\\
1.088	-0.127901\\
1.09	-0.127899\\
1.092	-0.127893\\
1.094	-0.127884\\
1.096	-0.127871\\
1.098	-0.127856\\
1.1	-0.127838\\
1.102	-0.127817\\
1.104	-0.127794\\
1.106	-0.127769\\
1.108	-0.127743\\
1.11	-0.127714\\
1.112	-0.127684\\
1.114	-0.127653\\
1.116	-0.127621\\
1.118	-0.127588\\
1.12	-0.127554\\
1.122	-0.12752\\
1.124	-0.127485\\
1.126	-0.12745\\
1.128	-0.127415\\
1.13	-0.127379\\
1.132	-0.127344\\
1.134	-0.127309\\
1.136	-0.127274\\
1.138	-0.127239\\
1.14	-0.127205\\
1.142	-0.127171\\
1.144	-0.127138\\
1.146	-0.127104\\
1.148	-0.127072\\
1.15	-0.127039\\
1.152	-0.127007\\
1.154	-0.126976\\
1.156	-0.126945\\
1.158	-0.126914\\
1.16	-0.126883\\
1.162	-0.126853\\
1.164	-0.126823\\
1.166	-0.126793\\
1.168	-0.126763\\
1.17	-0.126734\\
1.172	-0.126704\\
1.174	-0.126673\\
1.176	-0.126643\\
1.178	-0.126612\\
1.18	-0.126581\\
1.182	-0.12655\\
1.184	-0.126517\\
1.186	-0.126484\\
1.188	-0.126451\\
1.19	-0.126416\\
1.192	-0.12638\\
1.194	-0.126344\\
1.196	-0.126306\\
1.198	-0.126267\\
1.2	-0.126226\\
1.202	-0.126185\\
1.204	-0.126141\\
1.206	-0.126097\\
1.208	-0.12605\\
1.21	-0.126002\\
1.212	-0.125953\\
1.214	-0.125902\\
1.216	-0.125849\\
1.218	-0.125794\\
1.22	-0.125738\\
1.222	-0.125679\\
1.224	-0.12562\\
1.226	-0.125558\\
1.228	-0.125494\\
1.23	-0.125429\\
1.232	-0.125363\\
1.234	-0.125294\\
1.236	-0.125224\\
1.238	-0.125153\\
1.24	-0.12508\\
1.242	-0.125006\\
1.244	-0.124931\\
1.246	-0.124854\\
1.248	-0.124776\\
1.25	-0.124698\\
1.252	-0.124619\\
1.254	-0.124538\\
1.256	-0.124458\\
1.258	-0.124377\\
1.26	-0.124295\\
1.262	-0.124214\\
1.264	-0.124132\\
1.266	-0.124051\\
1.268	-0.12397\\
1.27	-0.123889\\
1.272	-0.123809\\
1.274	-0.12373\\
1.276	-0.123652\\
1.278	-0.123575\\
1.28	-0.123499\\
1.282	-0.123425\\
1.284	-0.123352\\
1.286	-0.123281\\
1.288	-0.123212\\
1.29	-0.123146\\
1.292	-0.123081\\
1.294	-0.123019\\
1.296	-0.12296\\
1.298	-0.122903\\
1.3	-0.12285\\
1.302	-0.122799\\
1.304	-0.122752\\
1.306	-0.122707\\
1.308	-0.122667\\
1.31	-0.122629\\
1.312	-0.122596\\
1.314	-0.122566\\
1.316	-0.12254\\
1.318	-0.122518\\
1.32	-0.1225\\
1.322	-0.122485\\
1.324	-0.122475\\
1.326	-0.122469\\
1.328	-0.122467\\
1.33	-0.12247\\
1.332	-0.122476\\
1.334	-0.122487\\
1.336	-0.122502\\
1.338	-0.122521\\
1.34	-0.122544\\
1.342	-0.122571\\
1.344	-0.122602\\
1.346	-0.122638\\
1.348	-0.122677\\
1.35	-0.12272\\
1.352	-0.122766\\
1.354	-0.122817\\
1.356	-0.122871\\
1.358	-0.122928\\
1.36	-0.122988\\
1.362	-0.123052\\
1.364	-0.123119\\
1.366	-0.123188\\
1.368	-0.123261\\
1.37	-0.123336\\
1.372	-0.123413\\
1.374	-0.123493\\
1.376	-0.123574\\
1.378	-0.123658\\
1.38	-0.123743\\
1.382	-0.12383\\
1.384	-0.123918\\
1.386	-0.124007\\
1.388	-0.124097\\
1.39	-0.124187\\
1.392	-0.124279\\
1.394	-0.12437\\
1.396	-0.124462\\
1.398	-0.124553\\
1.4	-0.124645\\
1.402	-0.124736\\
1.404	-0.124826\\
1.406	-0.124915\\
1.408	-0.125003\\
1.41	-0.12509\\
1.412	-0.125176\\
1.414	-0.12526\\
1.416	-0.125342\\
1.418	-0.125422\\
1.42	-0.1255\\
1.422	-0.125576\\
1.424	-0.12565\\
1.426	-0.125721\\
1.428	-0.125789\\
1.43	-0.125855\\
1.432	-0.125917\\
1.434	-0.125977\\
1.436	-0.126033\\
1.438	-0.126086\\
1.44	-0.126136\\
1.442	-0.126183\\
1.444	-0.126225\\
1.446	-0.126265\\
1.448	-0.126301\\
1.45	-0.126333\\
1.452	-0.126361\\
1.454	-0.126386\\
1.456	-0.126408\\
1.458	-0.126425\\
1.46	-0.126439\\
1.462	-0.126449\\
1.464	-0.126455\\
1.466	-0.126458\\
1.468	-0.126457\\
1.47	-0.126453\\
1.472	-0.126445\\
1.474	-0.126434\\
1.476	-0.126419\\
1.478	-0.126401\\
1.48	-0.12638\\
1.482	-0.126356\\
1.484	-0.126329\\
1.486	-0.126298\\
1.488	-0.126265\\
1.49	-0.126229\\
1.492	-0.126191\\
1.494	-0.12615\\
1.496	-0.126107\\
1.498	-0.126062\\
1.5	-0.126014\\
1.502	-0.125965\\
1.504	-0.125914\\
1.506	-0.125861\\
1.508	-0.125807\\
1.51	-0.125751\\
1.512	-0.125694\\
1.514	-0.125636\\
1.516	-0.125577\\
1.518	-0.125517\\
1.52	-0.125456\\
1.522	-0.125395\\
1.524	-0.125334\\
1.526	-0.125273\\
1.528	-0.125211\\
1.53	-0.125149\\
1.532	-0.125088\\
1.534	-0.125027\\
1.536	-0.124966\\
1.538	-0.124906\\
1.54	-0.124846\\
1.542	-0.124788\\
1.544	-0.12473\\
1.546	-0.124673\\
1.548	-0.124617\\
1.55	-0.124563\\
1.552	-0.12451\\
1.554	-0.124458\\
1.556	-0.124407\\
1.558	-0.124358\\
1.56	-0.124311\\
1.562	-0.124266\\
1.564	-0.124222\\
1.566	-0.124179\\
1.568	-0.124139\\
1.57	-0.124101\\
1.572	-0.124064\\
1.574	-0.124029\\
1.576	-0.123997\\
1.578	-0.123966\\
1.58	-0.123937\\
1.582	-0.12391\\
1.584	-0.123885\\
1.586	-0.123862\\
1.588	-0.123841\\
1.59	-0.123822\\
1.592	-0.123805\\
1.594	-0.12379\\
1.596	-0.123777\\
1.598	-0.123765\\
1.6	-0.123756\\
1.602	-0.123748\\
1.604	-0.123741\\
1.606	-0.123737\\
1.608	-0.123734\\
1.61	-0.123732\\
1.612	-0.123732\\
1.614	-0.123734\\
1.616	-0.123737\\
1.618	-0.123741\\
1.62	-0.123746\\
1.622	-0.123753\\
1.624	-0.123761\\
1.626	-0.12377\\
1.628	-0.123779\\
1.63	-0.12379\\
1.632	-0.123802\\
1.634	-0.123814\\
1.636	-0.123827\\
1.638	-0.123841\\
1.64	-0.123855\\
1.642	-0.12387\\
1.644	-0.123885\\
1.646	-0.1239\\
1.648	-0.123916\\
1.65	-0.123932\\
1.652	-0.123949\\
1.654	-0.123965\\
1.656	-0.123982\\
1.658	-0.123998\\
1.66	-0.124014\\
1.662	-0.124031\\
1.664	-0.124047\\
1.666	-0.124063\\
1.668	-0.124079\\
1.67	-0.124095\\
1.672	-0.12411\\
1.674	-0.124125\\
1.676	-0.12414\\
1.678	-0.124154\\
1.68	-0.124168\\
1.682	-0.124181\\
1.684	-0.124194\\
1.686	-0.124207\\
1.688	-0.124219\\
1.69	-0.124231\\
1.692	-0.124242\\
1.694	-0.124252\\
1.696	-0.124263\\
1.698	-0.124272\\
1.7	-0.124281\\
1.702	-0.12429\\
1.704	-0.124298\\
1.706	-0.124306\\
1.708	-0.124314\\
1.71	-0.124321\\
1.712	-0.124327\\
1.714	-0.124333\\
1.716	-0.124339\\
1.718	-0.124344\\
1.72	-0.124349\\
1.722	-0.124354\\
1.724	-0.124358\\
1.726	-0.124362\\
1.728	-0.124366\\
1.73	-0.12437\\
1.732	-0.124373\\
1.734	-0.124377\\
1.736	-0.12438\\
1.738	-0.124383\\
1.74	-0.124386\\
1.742	-0.124388\\
1.744	-0.124391\\
1.746	-0.124394\\
1.748	-0.124397\\
1.75	-0.124399\\
1.752	-0.124402\\
1.754	-0.124405\\
1.756	-0.124408\\
1.758	-0.124411\\
1.76	-0.124414\\
1.762	-0.124417\\
1.764	-0.124421\\
1.766	-0.124424\\
1.768	-0.124428\\
1.77	-0.124432\\
1.772	-0.124436\\
1.774	-0.12444\\
1.776	-0.124445\\
1.778	-0.124449\\
1.78	-0.124454\\
1.782	-0.124459\\
1.784	-0.124464\\
1.786	-0.12447\\
1.788	-0.124475\\
1.79	-0.124481\\
1.792	-0.124487\\
1.794	-0.124493\\
1.796	-0.124499\\
1.798	-0.124505\\
1.8	-0.124511\\
1.802	-0.124518\\
1.804	-0.124524\\
1.806	-0.124531\\
1.808	-0.124538\\
1.81	-0.124544\\
1.812	-0.124551\\
1.814	-0.124557\\
1.816	-0.124564\\
1.818	-0.12457\\
1.82	-0.124576\\
1.822	-0.124583\\
1.824	-0.124589\\
1.826	-0.124594\\
1.828	-0.1246\\
1.83	-0.124605\\
1.832	-0.12461\\
1.834	-0.124615\\
1.836	-0.124619\\
1.838	-0.124623\\
1.84	-0.124627\\
1.842	-0.12463\\
1.844	-0.124633\\
1.846	-0.124636\\
1.848	-0.124638\\
1.85	-0.124639\\
1.852	-0.12464\\
1.854	-0.12464\\
1.856	-0.12464\\
1.858	-0.124639\\
1.86	-0.124638\\
1.862	-0.124636\\
1.864	-0.124634\\
1.866	-0.124631\\
1.868	-0.124627\\
1.87	-0.124623\\
1.872	-0.124618\\
1.874	-0.124612\\
1.876	-0.124606\\
1.878	-0.124599\\
1.88	-0.124591\\
1.882	-0.124583\\
1.884	-0.124575\\
1.886	-0.124565\\
1.888	-0.124555\\
1.89	-0.124544\\
1.892	-0.124533\\
1.894	-0.124522\\
1.896	-0.124509\\
1.898	-0.124496\\
1.9	-0.124483\\
1.902	-0.124469\\
1.904	-0.124455\\
1.906	-0.12444\\
1.908	-0.124425\\
1.91	-0.124409\\
1.912	-0.124393\\
1.914	-0.124377\\
1.916	-0.12436\\
1.918	-0.124343\\
1.92	-0.124326\\
1.922	-0.124308\\
1.924	-0.124291\\
1.926	-0.124273\\
1.928	-0.124255\\
1.93	-0.124237\\
1.932	-0.124219\\
1.934	-0.124201\\
1.936	-0.124183\\
1.938	-0.124165\\
1.94	-0.124147\\
1.942	-0.12413\\
1.944	-0.124112\\
1.946	-0.124095\\
1.948	-0.124078\\
1.95	-0.124061\\
1.952	-0.124044\\
1.954	-0.124028\\
1.956	-0.124012\\
1.958	-0.123997\\
1.96	-0.123982\\
1.962	-0.123967\\
1.964	-0.123953\\
1.966	-0.12394\\
1.968	-0.123926\\
1.97	-0.123914\\
1.972	-0.123902\\
1.974	-0.123891\\
1.976	-0.12388\\
1.978	-0.12387\\
1.98	-0.12386\\
1.982	-0.123851\\
1.984	-0.123843\\
1.986	-0.123836\\
1.988	-0.123829\\
1.99	-0.123823\\
1.992	-0.123818\\
1.994	-0.123813\\
1.996	-0.123809\\
1.998	-0.123806\\
2	-0.123803\\
};
\addplot [color=mycolor7, line width=1.0pt, forget plot]
  table[row sep=crcr]{%
-0.124	0\\
-0.122	0\\
-0.12	0\\
-0.118	0\\
-0.116	0\\
-0.114	0\\
-0.112	0\\
-0.11	0\\
-0.108	0\\
-0.106	0\\
-0.104	0\\
-0.102	0\\
-0.1	0\\
-0.098	0\\
-0.096	0\\
-0.094	0\\
-0.092	0\\
-0.09	0\\
-0.088	0\\
-0.086	0\\
-0.084	0\\
-0.082	0\\
-0.08	0\\
-0.078	0\\
-0.076	0\\
-0.074	0\\
-0.072	0\\
-0.07	0\\
-0.068	0\\
-0.066	0\\
-0.064	0\\
-0.062	0\\
-0.06	0\\
-0.058	0\\
-0.056	0\\
-0.054	0\\
-0.052	0\\
-0.05	0\\
-0.048	0\\
-0.046	0\\
-0.044	0\\
-0.042	0\\
-0.04	0\\
-0.038	0\\
-0.036	0\\
-0.034	0\\
-0.032	0\\
-0.03	0\\
-0.028	0\\
-0.026	0\\
-0.024	0\\
-0.022	0\\
-0.02	0\\
-0.018	0\\
-0.016	0\\
-0.014	0\\
-0.012	0\\
-0.01	0\\
-0.008	0\\
-0.006	0\\
-0.004	0\\
-0.002	0\\
0	0\\
0.002	-0.00892065\\
0.004	-0.0177258\\
0.006	-0.0264052\\
0.008	-0.0349469\\
0.01	-0.0433382\\
0.012	-0.051566\\
0.014	-0.0596178\\
0.016	-0.0674812\\
0.018	-0.0751447\\
0.02	-0.0825974\\
0.022	-0.0898292\\
0.024	-0.0968309\\
0.026	-0.103594\\
0.028	-0.110111\\
0.03	-0.116374\\
0.032	-0.122378\\
0.034	-0.128117\\
0.036	-0.133586\\
0.038	-0.138781\\
0.04	-0.143698\\
0.042	-0.148335\\
0.044	-0.15269\\
0.046	-0.156761\\
0.048	-0.160547\\
0.05	-0.164048\\
0.052	-0.167265\\
0.054	-0.170198\\
0.056	-0.172849\\
0.058	-0.175221\\
0.06	-0.177316\\
0.062	-0.179138\\
0.064	-0.18069\\
0.066	-0.181976\\
0.068	-0.183003\\
0.07	-0.183774\\
0.072	-0.184297\\
0.074	-0.184577\\
0.076	-0.184621\\
0.078	-0.184435\\
0.08	-0.184029\\
0.082	-0.183408\\
0.084	-0.182582\\
0.086	-0.181558\\
0.088	-0.180345\\
0.09	-0.178953\\
0.092	-0.177389\\
0.094	-0.175664\\
0.096	-0.173786\\
0.098	-0.171766\\
0.1	-0.169612\\
0.102	-0.167334\\
0.104	-0.164942\\
0.106	-0.162445\\
0.108	-0.159854\\
0.11	-0.157177\\
0.112	-0.154425\\
0.114	-0.151606\\
0.116	-0.14873\\
0.118	-0.145807\\
0.12	-0.142845\\
0.122	-0.139852\\
0.124	-0.136839\\
0.126	-0.133812\\
0.128	-0.130781\\
0.13	-0.127753\\
0.132	-0.124737\\
0.134	-0.121738\\
0.136	-0.118766\\
0.138	-0.115825\\
0.14	-0.112923\\
0.142	-0.110067\\
0.144	-0.107261\\
0.146	-0.104511\\
0.148	-0.101823\\
0.15	-0.0992014\\
0.152	-0.0966503\\
0.154	-0.094174\\
0.156	-0.0917761\\
0.158	-0.0894602\\
0.16	-0.0872291\\
0.162	-0.0850856\\
0.164	-0.0830319\\
0.166	-0.08107\\
0.168	-0.0792014\\
0.17	-0.0774276\\
0.172	-0.0757493\\
0.174	-0.0741673\\
0.176	-0.0726819\\
0.178	-0.071293\\
0.18	-0.0700003\\
0.182	-0.0688034\\
0.184	-0.0677014\\
0.186	-0.0666932\\
0.188	-0.0657774\\
0.19	-0.0649526\\
0.192	-0.0642169\\
0.194	-0.0635684\\
0.196	-0.0630048\\
0.198	-0.0625239\\
0.2	-0.0621232\\
0.202	-0.0618\\
0.204	-0.0615516\\
0.206	-0.0613751\\
0.208	-0.0612675\\
0.21	-0.0612258\\
0.212	-0.0612468\\
0.214	-0.0613275\\
0.216	-0.0614645\\
0.218	-0.0616546\\
0.22	-0.0618947\\
0.222	-0.0621815\\
0.224	-0.0625117\\
0.226	-0.0628823\\
0.228	-0.0632899\\
0.23	-0.0637317\\
0.232	-0.0642044\\
0.234	-0.0647053\\
0.236	-0.0652315\\
0.238	-0.0657801\\
0.24	-0.0663485\\
0.242	-0.0669343\\
0.244	-0.0675349\\
0.246	-0.0681481\\
0.248	-0.0687718\\
0.25	-0.0694038\\
0.252	-0.0700424\\
0.254	-0.0706857\\
0.256	-0.0713323\\
0.258	-0.0719806\\
0.26	-0.0726295\\
0.262	-0.0732777\\
0.264	-0.0739243\\
0.266	-0.0745686\\
0.268	-0.0752097\\
0.27	-0.0758473\\
0.272	-0.0764809\\
0.274	-0.0771103\\
0.276	-0.0777355\\
0.278	-0.0783565\\
0.28	-0.0789734\\
0.282	-0.0795867\\
0.284	-0.0801967\\
0.286	-0.080804\\
0.288	-0.0814092\\
0.29	-0.0820131\\
0.292	-0.0826165\\
0.294	-0.0832204\\
0.296	-0.0838258\\
0.298	-0.0844337\\
0.3	-0.0850453\\
0.302	-0.0856617\\
0.304	-0.0862843\\
0.306	-0.0869142\\
0.308	-0.0875527\\
0.31	-0.0882011\\
0.312	-0.0888608\\
0.314	-0.089533\\
0.316	-0.090219\\
0.318	-0.09092\\
0.32	-0.0916374\\
0.322	-0.0923723\\
0.324	-0.0931259\\
0.326	-0.0938992\\
0.328	-0.0946932\\
0.33	-0.0955091\\
0.332	-0.0963475\\
0.334	-0.0972095\\
0.336	-0.0980955\\
0.338	-0.0990064\\
0.34	-0.0999426\\
0.342	-0.100905\\
0.344	-0.101893\\
0.346	-0.102907\\
0.348	-0.103947\\
0.35	-0.105014\\
0.352	-0.106107\\
0.354	-0.107226\\
0.356	-0.10837\\
0.358	-0.109539\\
0.36	-0.110732\\
0.362	-0.111949\\
0.364	-0.113188\\
0.366	-0.114448\\
0.368	-0.115729\\
0.37	-0.117029\\
0.372	-0.118347\\
0.374	-0.11968\\
0.376	-0.121029\\
0.378	-0.122389\\
0.38	-0.123761\\
0.382	-0.125142\\
0.384	-0.126529\\
0.386	-0.127922\\
0.388	-0.129317\\
0.39	-0.130713\\
0.392	-0.132106\\
0.394	-0.133496\\
0.396	-0.134878\\
0.398	-0.136251\\
0.4	-0.137613\\
0.402	-0.13896\\
0.404	-0.14029\\
0.406	-0.141601\\
0.408	-0.14289\\
0.41	-0.144154\\
0.412	-0.145391\\
0.414	-0.146598\\
0.416	-0.147772\\
0.418	-0.148912\\
0.42	-0.150015\\
0.422	-0.151078\\
0.424	-0.1521\\
0.426	-0.153077\\
0.428	-0.154007\\
0.43	-0.154889\\
0.432	-0.155721\\
0.434	-0.1565\\
0.436	-0.157225\\
0.438	-0.157894\\
0.44	-0.158506\\
0.442	-0.159058\\
0.444	-0.15955\\
0.446	-0.159981\\
0.448	-0.160349\\
0.45	-0.160654\\
0.452	-0.160894\\
0.454	-0.161069\\
0.456	-0.161179\\
0.458	-0.161223\\
0.46	-0.1612\\
0.462	-0.161112\\
0.464	-0.160958\\
0.466	-0.160738\\
0.468	-0.160452\\
0.47	-0.160102\\
0.472	-0.159688\\
0.474	-0.15921\\
0.476	-0.158671\\
0.478	-0.15807\\
0.48	-0.15741\\
0.482	-0.156691\\
0.484	-0.155915\\
0.486	-0.155085\\
0.488	-0.154201\\
0.49	-0.153266\\
0.492	-0.152282\\
0.494	-0.151251\\
0.496	-0.150175\\
0.498	-0.149057\\
0.5	-0.147899\\
0.502	-0.146704\\
0.504	-0.145475\\
0.506	-0.144213\\
0.508	-0.142922\\
0.51	-0.141605\\
0.512	-0.140265\\
0.514	-0.138904\\
0.516	-0.137526\\
0.518	-0.136133\\
0.52	-0.134729\\
0.522	-0.133316\\
0.524	-0.131898\\
0.526	-0.130478\\
0.528	-0.129059\\
0.53	-0.127643\\
0.532	-0.126234\\
0.534	-0.124834\\
0.536	-0.123447\\
0.538	-0.122075\\
0.54	-0.120721\\
0.542	-0.119388\\
0.544	-0.118079\\
0.546	-0.116796\\
0.548	-0.115541\\
0.55	-0.114317\\
0.552	-0.113127\\
0.554	-0.111972\\
0.556	-0.110855\\
0.558	-0.109777\\
0.56	-0.108741\\
0.562	-0.107749\\
0.564	-0.106801\\
0.566	-0.1059\\
0.568	-0.105047\\
0.57	-0.104244\\
0.572	-0.10349\\
0.574	-0.102789\\
0.576	-0.102139\\
0.578	-0.101543\\
0.58	-0.101\\
0.582	-0.100512\\
0.584	-0.100078\\
0.586	-0.0996989\\
0.588	-0.0993748\\
0.59	-0.0991055\\
0.592	-0.098891\\
0.594	-0.0987309\\
0.596	-0.0986247\\
0.598	-0.098572\\
0.6	-0.0985721\\
0.602	-0.0986242\\
0.604	-0.0987273\\
0.606	-0.0988806\\
0.608	-0.0990829\\
0.61	-0.0993329\\
0.612	-0.0996295\\
0.614	-0.0999711\\
0.616	-0.100356\\
0.618	-0.100784\\
0.62	-0.101251\\
0.622	-0.101758\\
0.624	-0.102301\\
0.626	-0.10288\\
0.628	-0.103491\\
0.63	-0.104134\\
0.632	-0.104806\\
0.634	-0.105506\\
0.636	-0.106231\\
0.638	-0.106979\\
0.64	-0.107748\\
0.642	-0.108536\\
0.644	-0.109341\\
0.646	-0.110161\\
0.648	-0.110994\\
0.65	-0.111837\\
0.652	-0.112688\\
0.654	-0.113547\\
0.656	-0.114409\\
0.658	-0.115274\\
0.66	-0.11614\\
0.662	-0.117004\\
0.664	-0.117865\\
0.666	-0.118721\\
0.668	-0.11957\\
0.67	-0.12041\\
0.672	-0.12124\\
0.674	-0.122058\\
0.676	-0.122863\\
0.678	-0.123653\\
0.68	-0.124427\\
0.682	-0.125183\\
0.684	-0.12592\\
0.686	-0.126638\\
0.688	-0.127334\\
0.69	-0.128009\\
0.692	-0.12866\\
0.694	-0.129288\\
0.696	-0.129891\\
0.698	-0.130469\\
0.7	-0.131021\\
0.702	-0.131547\\
0.704	-0.132047\\
0.706	-0.132519\\
0.708	-0.132964\\
0.71	-0.133382\\
0.712	-0.133772\\
0.714	-0.134134\\
0.716	-0.134469\\
0.718	-0.134777\\
0.72	-0.135057\\
0.722	-0.13531\\
0.724	-0.135537\\
0.726	-0.135737\\
0.728	-0.135912\\
0.73	-0.136061\\
0.732	-0.136186\\
0.734	-0.136286\\
0.736	-0.136363\\
0.738	-0.136417\\
0.74	-0.136449\\
0.742	-0.13646\\
0.744	-0.13645\\
0.746	-0.136421\\
0.748	-0.136373\\
0.75	-0.136306\\
0.752	-0.136223\\
0.754	-0.136124\\
0.756	-0.13601\\
0.758	-0.135881\\
0.76	-0.135739\\
0.762	-0.135585\\
0.764	-0.135419\\
0.766	-0.135243\\
0.768	-0.135057\\
0.77	-0.134863\\
0.772	-0.134661\\
0.774	-0.134452\\
0.776	-0.134237\\
0.778	-0.134018\\
0.78	-0.133793\\
0.782	-0.133566\\
0.784	-0.133336\\
0.786	-0.133104\\
0.788	-0.132871\\
0.79	-0.132637\\
0.792	-0.132403\\
0.794	-0.13217\\
0.796	-0.131939\\
0.798	-0.131709\\
0.8	-0.131482\\
0.802	-0.131257\\
0.804	-0.131036\\
0.806	-0.130818\\
0.808	-0.130605\\
0.81	-0.130395\\
0.812	-0.13019\\
0.814	-0.12999\\
0.816	-0.129795\\
0.818	-0.129605\\
0.82	-0.129421\\
0.822	-0.129241\\
0.824	-0.129067\\
0.826	-0.128899\\
0.828	-0.128735\\
0.83	-0.128577\\
0.832	-0.128424\\
0.834	-0.128277\\
0.836	-0.128134\\
0.838	-0.127995\\
0.84	-0.127862\\
0.842	-0.127733\\
0.844	-0.127608\\
0.846	-0.127486\\
0.848	-0.127369\\
0.85	-0.127254\\
0.852	-0.127143\\
0.854	-0.127034\\
0.856	-0.126928\\
0.858	-0.126824\\
0.86	-0.126722\\
0.862	-0.126621\\
0.864	-0.126521\\
0.866	-0.126422\\
0.868	-0.126323\\
0.87	-0.126225\\
0.872	-0.126126\\
0.874	-0.126027\\
0.876	-0.125928\\
0.878	-0.125827\\
0.88	-0.125725\\
0.882	-0.125622\\
0.884	-0.125516\\
0.886	-0.125409\\
0.888	-0.1253\\
0.89	-0.125189\\
0.892	-0.125075\\
0.894	-0.124958\\
0.896	-0.124839\\
0.898	-0.124717\\
0.9	-0.124592\\
0.902	-0.124465\\
0.904	-0.124334\\
0.906	-0.124201\\
0.908	-0.124065\\
0.91	-0.123926\\
0.912	-0.123785\\
0.914	-0.12364\\
0.916	-0.123494\\
0.918	-0.123345\\
0.92	-0.123195\\
0.922	-0.123042\\
0.924	-0.122888\\
0.926	-0.122732\\
0.928	-0.122575\\
0.93	-0.122418\\
0.932	-0.12226\\
0.934	-0.122101\\
0.936	-0.121943\\
0.938	-0.121785\\
0.94	-0.121629\\
0.942	-0.121473\\
0.944	-0.121319\\
0.946	-0.121167\\
0.948	-0.121017\\
0.95	-0.12087\\
0.952	-0.120726\\
0.954	-0.120586\\
0.956	-0.12045\\
0.958	-0.120318\\
0.96	-0.120192\\
0.962	-0.12007\\
0.964	-0.119954\\
0.966	-0.119844\\
0.968	-0.11974\\
0.97	-0.119644\\
0.972	-0.119554\\
0.974	-0.119472\\
0.976	-0.119397\\
0.978	-0.119331\\
0.98	-0.119273\\
0.982	-0.119224\\
0.984	-0.119184\\
0.986	-0.119153\\
0.988	-0.119132\\
0.99	-0.11912\\
0.992	-0.119118\\
0.994	-0.119125\\
0.996	-0.119143\\
0.998	-0.119171\\
1	-0.11921\\
1.002	-0.119258\\
1.004	-0.119317\\
1.006	-0.119386\\
1.008	-0.119466\\
1.01	-0.119556\\
1.012	-0.119655\\
1.014	-0.119765\\
1.016	-0.119885\\
1.018	-0.120014\\
1.02	-0.120153\\
1.022	-0.120301\\
1.024	-0.120458\\
1.026	-0.120624\\
1.028	-0.120798\\
1.03	-0.120981\\
1.032	-0.121171\\
1.034	-0.121368\\
1.036	-0.121573\\
1.038	-0.121784\\
1.04	-0.122002\\
1.042	-0.122225\\
1.044	-0.122453\\
1.046	-0.122687\\
1.048	-0.122924\\
1.05	-0.123166\\
1.052	-0.123411\\
1.054	-0.123658\\
1.056	-0.123908\\
1.058	-0.12416\\
1.06	-0.124413\\
1.062	-0.124666\\
1.064	-0.12492\\
1.066	-0.125173\\
1.068	-0.125425\\
1.07	-0.125676\\
1.072	-0.125925\\
1.074	-0.126171\\
1.076	-0.126414\\
1.078	-0.126654\\
1.08	-0.126889\\
1.082	-0.127119\\
1.084	-0.127345\\
1.086	-0.127565\\
1.088	-0.127779\\
1.09	-0.127986\\
1.092	-0.128187\\
1.094	-0.12838\\
1.096	-0.128566\\
1.098	-0.128743\\
1.1	-0.128913\\
1.102	-0.129073\\
1.104	-0.129224\\
1.106	-0.129366\\
1.108	-0.129499\\
1.11	-0.129622\\
1.112	-0.129734\\
1.114	-0.129837\\
1.116	-0.129929\\
1.118	-0.130011\\
1.12	-0.130082\\
1.122	-0.130143\\
1.124	-0.130193\\
1.126	-0.130232\\
1.128	-0.130261\\
1.13	-0.130279\\
1.132	-0.130286\\
1.134	-0.130283\\
1.136	-0.13027\\
1.138	-0.130246\\
1.14	-0.130212\\
1.142	-0.130168\\
1.144	-0.130114\\
1.146	-0.130051\\
1.148	-0.129978\\
1.15	-0.129897\\
1.152	-0.129806\\
1.154	-0.129707\\
1.156	-0.1296\\
1.158	-0.129486\\
1.16	-0.129363\\
1.162	-0.129234\\
1.164	-0.129098\\
1.166	-0.128955\\
1.168	-0.128807\\
1.17	-0.128653\\
1.172	-0.128494\\
1.174	-0.12833\\
1.176	-0.128162\\
1.178	-0.12799\\
1.18	-0.127815\\
1.182	-0.127637\\
1.184	-0.127456\\
1.186	-0.127273\\
1.188	-0.127089\\
1.19	-0.126903\\
1.192	-0.126717\\
1.194	-0.12653\\
1.196	-0.126343\\
1.198	-0.126157\\
1.2	-0.125972\\
1.202	-0.125788\\
1.204	-0.125605\\
1.206	-0.125425\\
1.208	-0.125247\\
1.21	-0.125072\\
1.212	-0.1249\\
1.214	-0.124732\\
1.216	-0.124567\\
1.218	-0.124406\\
1.22	-0.12425\\
1.222	-0.124099\\
1.224	-0.123952\\
1.226	-0.12381\\
1.228	-0.123674\\
1.23	-0.123543\\
1.232	-0.123418\\
1.234	-0.123299\\
1.236	-0.123185\\
1.238	-0.123078\\
1.24	-0.122977\\
1.242	-0.122883\\
1.244	-0.122795\\
1.246	-0.122713\\
1.248	-0.122638\\
1.25	-0.122569\\
1.252	-0.122507\\
1.254	-0.122451\\
1.256	-0.122402\\
1.258	-0.122359\\
1.26	-0.122323\\
1.262	-0.122292\\
1.264	-0.122268\\
1.266	-0.12225\\
1.268	-0.122238\\
1.27	-0.122232\\
1.272	-0.122231\\
1.274	-0.122236\\
1.276	-0.122246\\
1.278	-0.122262\\
1.28	-0.122282\\
1.282	-0.122307\\
1.284	-0.122337\\
1.286	-0.122371\\
1.288	-0.122409\\
1.29	-0.122451\\
1.292	-0.122497\\
1.294	-0.122546\\
1.296	-0.122598\\
1.298	-0.122653\\
1.3	-0.122712\\
1.302	-0.122772\\
1.304	-0.122835\\
1.306	-0.1229\\
1.308	-0.122966\\
1.31	-0.123034\\
1.312	-0.123104\\
1.314	-0.123174\\
1.316	-0.123246\\
1.318	-0.123318\\
1.32	-0.12339\\
1.322	-0.123463\\
1.324	-0.123535\\
1.326	-0.123608\\
1.328	-0.12368\\
1.33	-0.123751\\
1.332	-0.123822\\
1.334	-0.123892\\
1.336	-0.123961\\
1.338	-0.124029\\
1.34	-0.124095\\
1.342	-0.12416\\
1.344	-0.124223\\
1.346	-0.124285\\
1.348	-0.124345\\
1.35	-0.124403\\
1.352	-0.124459\\
1.354	-0.124513\\
1.356	-0.124565\\
1.358	-0.124614\\
1.36	-0.124662\\
1.362	-0.124707\\
1.364	-0.124751\\
1.366	-0.124791\\
1.368	-0.12483\\
1.37	-0.124867\\
1.372	-0.124901\\
1.374	-0.124933\\
1.376	-0.124963\\
1.378	-0.124991\\
1.38	-0.125016\\
1.382	-0.12504\\
1.384	-0.125061\\
1.386	-0.125081\\
1.388	-0.125099\\
1.39	-0.125115\\
1.392	-0.125129\\
1.394	-0.125142\\
1.396	-0.125153\\
1.398	-0.125163\\
1.4	-0.125171\\
1.402	-0.125178\\
1.404	-0.125184\\
1.406	-0.125188\\
1.408	-0.125192\\
1.41	-0.125195\\
1.412	-0.125197\\
1.414	-0.125198\\
1.416	-0.125199\\
1.418	-0.125199\\
1.42	-0.125199\\
1.422	-0.125198\\
1.424	-0.125197\\
1.426	-0.125196\\
1.428	-0.125195\\
1.43	-0.125193\\
1.432	-0.125192\\
1.434	-0.125191\\
1.436	-0.12519\\
1.438	-0.125189\\
1.44	-0.125189\\
1.442	-0.125189\\
1.444	-0.125189\\
1.446	-0.12519\\
1.448	-0.125191\\
1.45	-0.125193\\
1.452	-0.125195\\
1.454	-0.125198\\
1.456	-0.125201\\
1.458	-0.125205\\
1.46	-0.125209\\
1.462	-0.125214\\
1.464	-0.12522\\
1.466	-0.125226\\
1.468	-0.125232\\
1.47	-0.125239\\
1.472	-0.125247\\
1.474	-0.125255\\
1.476	-0.125263\\
1.478	-0.125272\\
1.48	-0.125281\\
1.482	-0.12529\\
1.484	-0.125299\\
1.486	-0.125309\\
1.488	-0.125319\\
1.49	-0.125328\\
1.492	-0.125338\\
1.494	-0.125348\\
1.496	-0.125357\\
1.498	-0.125367\\
1.5	-0.125376\\
1.502	-0.125384\\
1.504	-0.125392\\
1.506	-0.1254\\
1.508	-0.125407\\
1.51	-0.125414\\
1.512	-0.125419\\
1.514	-0.125424\\
1.516	-0.125429\\
1.518	-0.125432\\
1.52	-0.125434\\
1.522	-0.125435\\
1.524	-0.125435\\
1.526	-0.125435\\
1.528	-0.125432\\
1.53	-0.125429\\
1.532	-0.125424\\
1.534	-0.125418\\
1.536	-0.125411\\
1.538	-0.125402\\
1.54	-0.125392\\
1.542	-0.12538\\
1.544	-0.125367\\
1.546	-0.125353\\
1.548	-0.125337\\
1.55	-0.125319\\
1.552	-0.1253\\
1.554	-0.12528\\
1.556	-0.125258\\
1.558	-0.125235\\
1.56	-0.12521\\
1.562	-0.125183\\
1.564	-0.125156\\
1.566	-0.125127\\
1.568	-0.125096\\
1.57	-0.125065\\
1.572	-0.125032\\
1.574	-0.124998\\
1.576	-0.124963\\
1.578	-0.124927\\
1.58	-0.12489\\
1.582	-0.124852\\
1.584	-0.124813\\
1.586	-0.124773\\
1.588	-0.124733\\
1.59	-0.124692\\
1.592	-0.12465\\
1.594	-0.124608\\
1.596	-0.124566\\
1.598	-0.124523\\
1.6	-0.124481\\
1.602	-0.124438\\
1.604	-0.124395\\
1.606	-0.124352\\
1.608	-0.12431\\
1.61	-0.124267\\
1.612	-0.124226\\
1.614	-0.124184\\
1.616	-0.124143\\
1.618	-0.124103\\
1.62	-0.124064\\
1.622	-0.124025\\
1.624	-0.123987\\
1.626	-0.123951\\
1.628	-0.123915\\
1.63	-0.123881\\
1.632	-0.123847\\
1.634	-0.123816\\
1.636	-0.123785\\
1.638	-0.123756\\
1.64	-0.123729\\
1.642	-0.123703\\
1.644	-0.123678\\
1.646	-0.123656\\
1.648	-0.123635\\
1.65	-0.123616\\
1.652	-0.123599\\
1.654	-0.123583\\
1.656	-0.12357\\
1.658	-0.123558\\
1.66	-0.123548\\
1.662	-0.123541\\
1.664	-0.123535\\
1.666	-0.123531\\
1.668	-0.123529\\
1.67	-0.123529\\
1.672	-0.123531\\
1.674	-0.123534\\
1.676	-0.12354\\
1.678	-0.123548\\
1.68	-0.123557\\
1.682	-0.123568\\
1.684	-0.123581\\
1.686	-0.123596\\
1.688	-0.123612\\
1.69	-0.12363\\
1.692	-0.12365\\
1.694	-0.123671\\
1.696	-0.123693\\
1.698	-0.123717\\
1.7	-0.123742\\
1.702	-0.123769\\
1.704	-0.123796\\
1.706	-0.123825\\
1.708	-0.123854\\
1.71	-0.123885\\
1.712	-0.123917\\
1.714	-0.123949\\
1.716	-0.123982\\
1.718	-0.124015\\
1.72	-0.124049\\
1.722	-0.124084\\
1.724	-0.124118\\
1.726	-0.124154\\
1.728	-0.124189\\
1.73	-0.124224\\
1.732	-0.124259\\
1.734	-0.124294\\
1.736	-0.124329\\
1.738	-0.124364\\
1.74	-0.124399\\
1.742	-0.124433\\
1.744	-0.124466\\
1.746	-0.124499\\
1.748	-0.124531\\
1.75	-0.124563\\
1.752	-0.124593\\
1.754	-0.124623\\
1.756	-0.124652\\
1.758	-0.12468\\
1.76	-0.124707\\
1.762	-0.124733\\
1.764	-0.124758\\
1.766	-0.124782\\
1.768	-0.124804\\
1.77	-0.124825\\
1.772	-0.124845\\
1.774	-0.124863\\
1.776	-0.124881\\
1.778	-0.124896\\
1.78	-0.124911\\
1.782	-0.124924\\
1.784	-0.124935\\
1.786	-0.124946\\
1.788	-0.124954\\
1.79	-0.124962\\
1.792	-0.124968\\
1.794	-0.124972\\
1.796	-0.124975\\
1.798	-0.124977\\
1.8	-0.124977\\
1.802	-0.124976\\
1.804	-0.124974\\
1.806	-0.12497\\
1.808	-0.124965\\
1.81	-0.124959\\
1.812	-0.124952\\
1.814	-0.124943\\
1.816	-0.124933\\
1.818	-0.124923\\
1.82	-0.124911\\
1.822	-0.124898\\
1.824	-0.124884\\
1.826	-0.12487\\
1.828	-0.124854\\
1.83	-0.124838\\
1.832	-0.124821\\
1.834	-0.124803\\
1.836	-0.124785\\
1.838	-0.124766\\
1.84	-0.124747\\
1.842	-0.124727\\
1.844	-0.124706\\
1.846	-0.124686\\
1.848	-0.124665\\
1.85	-0.124644\\
1.852	-0.124622\\
1.854	-0.124601\\
1.856	-0.124579\\
1.858	-0.124558\\
1.86	-0.124536\\
1.862	-0.124515\\
1.864	-0.124493\\
1.866	-0.124472\\
1.868	-0.124451\\
1.87	-0.12443\\
1.872	-0.12441\\
1.874	-0.12439\\
1.876	-0.12437\\
1.878	-0.124351\\
1.88	-0.124332\\
1.882	-0.124313\\
1.884	-0.124295\\
1.886	-0.124278\\
1.888	-0.124261\\
1.89	-0.124245\\
1.892	-0.124229\\
1.894	-0.124214\\
1.896	-0.124199\\
1.898	-0.124185\\
1.9	-0.124172\\
1.902	-0.12416\\
1.904	-0.124148\\
1.906	-0.124136\\
1.908	-0.124126\\
1.91	-0.124116\\
1.912	-0.124107\\
1.914	-0.124098\\
1.916	-0.12409\\
1.918	-0.124083\\
1.92	-0.124076\\
1.922	-0.12407\\
1.924	-0.124064\\
1.926	-0.12406\\
1.928	-0.124055\\
1.93	-0.124051\\
1.932	-0.124048\\
1.934	-0.124045\\
1.936	-0.124043\\
1.938	-0.124041\\
1.94	-0.12404\\
1.942	-0.124039\\
1.944	-0.124039\\
1.946	-0.124039\\
1.948	-0.124039\\
1.95	-0.124039\\
1.952	-0.12404\\
1.954	-0.124041\\
1.956	-0.124043\\
1.958	-0.124044\\
1.96	-0.124046\\
1.962	-0.124048\\
1.964	-0.12405\\
1.966	-0.124052\\
1.968	-0.124054\\
1.97	-0.124057\\
1.972	-0.124059\\
1.974	-0.124061\\
1.976	-0.124064\\
1.978	-0.124066\\
1.98	-0.124069\\
1.982	-0.124071\\
1.984	-0.124073\\
1.986	-0.124075\\
1.988	-0.124077\\
1.99	-0.124079\\
1.992	-0.124081\\
1.994	-0.124082\\
1.996	-0.124084\\
1.998	-0.124085\\
2	-0.124086\\
};
\addplot [color=mycolor8, line width=1.0pt, forget plot]
  table[row sep=crcr]{%
-0.124	0\\
-0.122	0\\
-0.12	0\\
-0.118	0\\
-0.116	0\\
-0.114	0\\
-0.112	0\\
-0.11	0\\
-0.108	0\\
-0.106	0\\
-0.104	0\\
-0.102	0\\
-0.1	0\\
-0.098	0\\
-0.096	0\\
-0.094	0\\
-0.092	0\\
-0.09	0\\
-0.088	0\\
-0.086	0\\
-0.084	0\\
-0.082	0\\
-0.08	0\\
-0.078	0\\
-0.076	0\\
-0.074	0\\
-0.072	0\\
-0.07	0\\
-0.068	0\\
-0.066	0\\
-0.064	0\\
-0.062	0\\
-0.06	0\\
-0.058	0\\
-0.056	0\\
-0.054	0\\
-0.052	0\\
-0.05	0\\
-0.048	0\\
-0.046	0\\
-0.044	0\\
-0.042	0\\
-0.04	0\\
-0.038	0\\
-0.036	0\\
-0.034	0\\
-0.032	0\\
-0.03	0\\
-0.028	0\\
-0.026	0\\
-0.024	0\\
-0.022	0\\
-0.02	0\\
-0.018	0\\
-0.016	0\\
-0.014	0\\
-0.012	0\\
-0.01	0\\
-0.008	0\\
-0.006	0\\
-0.004	0\\
-0.002	0\\
0	0\\
0.002	-0.006565\\
0.004	-0.0111084\\
0.006	-0.0143255\\
0.008	-0.0165834\\
0.01	-0.0180911\\
0.012	-0.0189899\\
0.014	-0.0193971\\
0.016	-0.0194231\\
0.018	-0.0191742\\
0.02	-0.0187508\\
0.022	-0.0182433\\
0.024	-0.0177282\\
0.026	-0.0172673\\
0.028	-0.0169061\\
0.03	-0.0166756\\
0.032	-0.0165937\\
0.034	-0.0166674\\
0.036	-0.0168955\\
0.038	-0.0172707\\
0.04	-0.0177816\\
0.042	-0.0184146\\
0.044	-0.0191553\\
0.046	-0.0199891\\
0.048	-0.0209023\\
0.05	-0.0218823\\
0.052	-0.0229181\\
0.054	-0.0239999\\
0.056	-0.0251193\\
0.058	-0.0262695\\
0.06	-0.0274444\\
0.062	-0.0286394\\
0.064	-0.0298503\\
0.066	-0.0310741\\
0.068	-0.0323079\\
0.07	-0.0335495\\
0.072	-0.0347972\\
0.074	-0.0360494\\
0.076	-0.0373048\\
0.078	-0.0385624\\
0.08	-0.0398212\\
0.082	-0.0410804\\
0.084	-0.0423393\\
0.086	-0.0435974\\
0.088	-0.044854\\
0.09	-0.0461087\\
0.092	-0.0473611\\
0.094	-0.0486107\\
0.096	-0.0498573\\
0.098	-0.0511003\\
0.1	-0.0523396\\
0.102	-0.0535747\\
0.104	-0.0548054\\
0.106	-0.0560312\\
0.108	-0.057252\\
0.11	-0.0584672\\
0.112	-0.0596767\\
0.114	-0.0608801\\
0.116	-0.062077\\
0.118	-0.0632669\\
0.12	-0.0644497\\
0.122	-0.0656248\\
0.124	-0.0667918\\
0.126	-0.0679504\\
0.128	-0.0691001\\
0.13	-0.0702404\\
0.132	-0.071371\\
0.134	-0.0724914\\
0.136	-0.073601\\
0.138	-0.0746995\\
0.14	-0.0757864\\
0.142	-0.0768612\\
0.144	-0.0779234\\
0.146	-0.0789726\\
0.148	-0.0800082\\
0.15	-0.0810298\\
0.152	-0.0820369\\
0.154	-0.0830291\\
0.156	-0.0840058\\
0.158	-0.0849667\\
0.16	-0.0859112\\
0.162	-0.086839\\
0.164	-0.0877495\\
0.166	-0.0886424\\
0.168	-0.0895173\\
0.17	-0.0903738\\
0.172	-0.0912115\\
0.174	-0.09203\\
0.176	-0.092829\\
0.178	-0.0936081\\
0.18	-0.0943671\\
0.182	-0.0951057\\
0.184	-0.0958236\\
0.186	-0.0965205\\
0.188	-0.0971962\\
0.19	-0.0978505\\
0.192	-0.0984832\\
0.194	-0.0990942\\
0.196	-0.0996833\\
0.198	-0.10025\\
0.2	-0.100795\\
0.202	-0.101318\\
0.204	-0.101819\\
0.206	-0.102297\\
0.208	-0.102753\\
0.21	-0.103187\\
0.212	-0.103599\\
0.214	-0.103989\\
0.216	-0.104357\\
0.218	-0.104703\\
0.22	-0.105027\\
0.222	-0.10533\\
0.224	-0.105612\\
0.226	-0.105872\\
0.228	-0.106112\\
0.23	-0.106331\\
0.232	-0.10653\\
0.234	-0.106709\\
0.236	-0.106868\\
0.238	-0.107008\\
0.24	-0.107129\\
0.242	-0.107231\\
0.244	-0.107316\\
0.246	-0.107382\\
0.248	-0.107432\\
0.25	-0.107464\\
0.252	-0.107481\\
0.254	-0.107481\\
0.256	-0.107466\\
0.258	-0.107437\\
0.26	-0.107393\\
0.262	-0.107336\\
0.264	-0.107265\\
0.266	-0.107182\\
0.268	-0.107087\\
0.27	-0.10698\\
0.272	-0.106863\\
0.274	-0.106735\\
0.276	-0.106598\\
0.278	-0.106451\\
0.28	-0.106297\\
0.282	-0.106135\\
0.284	-0.105965\\
0.286	-0.105789\\
0.288	-0.105607\\
0.29	-0.10542\\
0.292	-0.105228\\
0.294	-0.105033\\
0.296	-0.104834\\
0.298	-0.104632\\
0.3	-0.104428\\
0.302	-0.104222\\
0.304	-0.104016\\
0.306	-0.103809\\
0.308	-0.103603\\
0.31	-0.103397\\
0.312	-0.103193\\
0.314	-0.102991\\
0.316	-0.102792\\
0.318	-0.102596\\
0.32	-0.102403\\
0.322	-0.102215\\
0.324	-0.102031\\
0.326	-0.101852\\
0.328	-0.10168\\
0.33	-0.101513\\
0.332	-0.101353\\
0.334	-0.101201\\
0.336	-0.101056\\
0.338	-0.100919\\
0.34	-0.10079\\
0.342	-0.10067\\
0.344	-0.100559\\
0.346	-0.100458\\
0.348	-0.100366\\
0.35	-0.100285\\
0.352	-0.100214\\
0.354	-0.100154\\
0.356	-0.100105\\
0.358	-0.100068\\
0.36	-0.100041\\
0.362	-0.100027\\
0.364	-0.100024\\
0.366	-0.100034\\
0.368	-0.100056\\
0.37	-0.10009\\
0.372	-0.100137\\
0.374	-0.100196\\
0.376	-0.100268\\
0.378	-0.100353\\
0.38	-0.10045\\
0.382	-0.100561\\
0.384	-0.100684\\
0.386	-0.10082\\
0.388	-0.100969\\
0.39	-0.10113\\
0.392	-0.101304\\
0.394	-0.101491\\
0.396	-0.10169\\
0.398	-0.101902\\
0.4	-0.102126\\
0.402	-0.102362\\
0.404	-0.10261\\
0.406	-0.10287\\
0.408	-0.103141\\
0.41	-0.103423\\
0.412	-0.103717\\
0.414	-0.104021\\
0.416	-0.104337\\
0.418	-0.104662\\
0.42	-0.104997\\
0.422	-0.105343\\
0.424	-0.105697\\
0.426	-0.106061\\
0.428	-0.106434\\
0.43	-0.106815\\
0.432	-0.107204\\
0.434	-0.107601\\
0.436	-0.108005\\
0.438	-0.108416\\
0.44	-0.108834\\
0.442	-0.109258\\
0.444	-0.109688\\
0.446	-0.110123\\
0.448	-0.110563\\
0.45	-0.111008\\
0.452	-0.111457\\
0.454	-0.111909\\
0.456	-0.112365\\
0.458	-0.112824\\
0.46	-0.113285\\
0.462	-0.113748\\
0.464	-0.114213\\
0.466	-0.114679\\
0.468	-0.115145\\
0.47	-0.115612\\
0.472	-0.116078\\
0.474	-0.116544\\
0.476	-0.117009\\
0.478	-0.117472\\
0.48	-0.117933\\
0.482	-0.118392\\
0.484	-0.118848\\
0.486	-0.119301\\
0.488	-0.11975\\
0.49	-0.120195\\
0.492	-0.120636\\
0.494	-0.121072\\
0.496	-0.121503\\
0.498	-0.121929\\
0.5	-0.122348\\
0.502	-0.122761\\
0.504	-0.123168\\
0.506	-0.123567\\
0.508	-0.123959\\
0.51	-0.124344\\
0.512	-0.124721\\
0.514	-0.12509\\
0.516	-0.12545\\
0.518	-0.125801\\
0.52	-0.126143\\
0.522	-0.126477\\
0.524	-0.1268\\
0.526	-0.127114\\
0.528	-0.127418\\
0.53	-0.127712\\
0.532	-0.127996\\
0.534	-0.128269\\
0.536	-0.128531\\
0.538	-0.128783\\
0.54	-0.129024\\
0.542	-0.129254\\
0.544	-0.129473\\
0.546	-0.12968\\
0.548	-0.129877\\
0.55	-0.130062\\
0.552	-0.130235\\
0.554	-0.130398\\
0.556	-0.130549\\
0.558	-0.130688\\
0.56	-0.130817\\
0.562	-0.130933\\
0.564	-0.131039\\
0.566	-0.131133\\
0.568	-0.131217\\
0.57	-0.131289\\
0.572	-0.13135\\
0.574	-0.1314\\
0.576	-0.131439\\
0.578	-0.131468\\
0.58	-0.131486\\
0.582	-0.131494\\
0.584	-0.131492\\
0.586	-0.13148\\
0.588	-0.131458\\
0.59	-0.131426\\
0.592	-0.131385\\
0.594	-0.131335\\
0.596	-0.131276\\
0.598	-0.131208\\
0.6	-0.131132\\
0.602	-0.131048\\
0.604	-0.130955\\
0.606	-0.130855\\
0.608	-0.130748\\
0.61	-0.130633\\
0.612	-0.130512\\
0.614	-0.130384\\
0.616	-0.130249\\
0.618	-0.130109\\
0.62	-0.129963\\
0.622	-0.129812\\
0.624	-0.129656\\
0.626	-0.129495\\
0.628	-0.12933\\
0.63	-0.129161\\
0.632	-0.128988\\
0.634	-0.128811\\
0.636	-0.128632\\
0.638	-0.12845\\
0.64	-0.128265\\
0.642	-0.128078\\
0.644	-0.12789\\
0.646	-0.127699\\
0.648	-0.127508\\
0.65	-0.127316\\
0.652	-0.127124\\
0.654	-0.126931\\
0.656	-0.126738\\
0.658	-0.126546\\
0.66	-0.126354\\
0.662	-0.126163\\
0.664	-0.125973\\
0.666	-0.125785\\
0.668	-0.125599\\
0.67	-0.125415\\
0.672	-0.125233\\
0.674	-0.125054\\
0.676	-0.124878\\
0.678	-0.124704\\
0.68	-0.124534\\
0.682	-0.124367\\
0.684	-0.124205\\
0.686	-0.124046\\
0.688	-0.123891\\
0.69	-0.12374\\
0.692	-0.123594\\
0.694	-0.123453\\
0.696	-0.123316\\
0.698	-0.123185\\
0.7	-0.123058\\
0.702	-0.122937\\
0.704	-0.122821\\
0.706	-0.122711\\
0.708	-0.122606\\
0.71	-0.122507\\
0.712	-0.122414\\
0.714	-0.122327\\
0.716	-0.122245\\
0.718	-0.12217\\
0.72	-0.122101\\
0.722	-0.122037\\
0.724	-0.12198\\
0.726	-0.121929\\
0.728	-0.121884\\
0.73	-0.121845\\
0.732	-0.121812\\
0.734	-0.121785\\
0.736	-0.121765\\
0.738	-0.12175\\
0.74	-0.121741\\
0.742	-0.121739\\
0.744	-0.121742\\
0.746	-0.121751\\
0.748	-0.121766\\
0.75	-0.121786\\
0.752	-0.121812\\
0.754	-0.121843\\
0.756	-0.12188\\
0.758	-0.121922\\
0.76	-0.121969\\
0.762	-0.122022\\
0.764	-0.122079\\
0.766	-0.12214\\
0.768	-0.122207\\
0.77	-0.122278\\
0.772	-0.122353\\
0.774	-0.122432\\
0.776	-0.122515\\
0.778	-0.122603\\
0.78	-0.122693\\
0.782	-0.122788\\
0.784	-0.122885\\
0.786	-0.122986\\
0.788	-0.12309\\
0.79	-0.123197\\
0.792	-0.123306\\
0.794	-0.123418\\
0.796	-0.123532\\
0.798	-0.123648\\
0.8	-0.123766\\
0.802	-0.123885\\
0.804	-0.124006\\
0.806	-0.124129\\
0.808	-0.124252\\
0.81	-0.124377\\
0.812	-0.124502\\
0.814	-0.124628\\
0.816	-0.124755\\
0.818	-0.124881\\
0.82	-0.125008\\
0.822	-0.125134\\
0.824	-0.125261\\
0.826	-0.125386\\
0.828	-0.125511\\
0.83	-0.125636\\
0.832	-0.125759\\
0.834	-0.125881\\
0.836	-0.126002\\
0.838	-0.126121\\
0.84	-0.126239\\
0.842	-0.126355\\
0.844	-0.12647\\
0.846	-0.126582\\
0.848	-0.126692\\
0.85	-0.1268\\
0.852	-0.126905\\
0.854	-0.127008\\
0.856	-0.127108\\
0.858	-0.127205\\
0.86	-0.1273\\
0.862	-0.127391\\
0.864	-0.12748\\
0.866	-0.127565\\
0.868	-0.127647\\
0.87	-0.127726\\
0.872	-0.127802\\
0.874	-0.127873\\
0.876	-0.127942\\
0.878	-0.128007\\
0.88	-0.128068\\
0.882	-0.128125\\
0.884	-0.128179\\
0.886	-0.128229\\
0.888	-0.128275\\
0.89	-0.128317\\
0.892	-0.128356\\
0.894	-0.12839\\
0.896	-0.128421\\
0.898	-0.128448\\
0.9	-0.128471\\
0.902	-0.12849\\
0.904	-0.128505\\
0.906	-0.128516\\
0.908	-0.128523\\
0.91	-0.128527\\
0.912	-0.128527\\
0.914	-0.128523\\
0.916	-0.128515\\
0.918	-0.128503\\
0.92	-0.128488\\
0.922	-0.128469\\
0.924	-0.128447\\
0.926	-0.128421\\
0.928	-0.128392\\
0.93	-0.128359\\
0.932	-0.128323\\
0.934	-0.128283\\
0.936	-0.128241\\
0.938	-0.128195\\
0.94	-0.128147\\
0.942	-0.128095\\
0.944	-0.128041\\
0.946	-0.127984\\
0.948	-0.127924\\
0.95	-0.127862\\
0.952	-0.127797\\
0.954	-0.127729\\
0.956	-0.12766\\
0.958	-0.127588\\
0.96	-0.127514\\
0.962	-0.127439\\
0.964	-0.127361\\
0.966	-0.127282\\
0.968	-0.127201\\
0.97	-0.127118\\
0.972	-0.127034\\
0.974	-0.126948\\
0.976	-0.126862\\
0.978	-0.126774\\
0.98	-0.126685\\
0.982	-0.126596\\
0.984	-0.126505\\
0.986	-0.126414\\
0.988	-0.126323\\
0.99	-0.126231\\
0.992	-0.126139\\
0.994	-0.126046\\
0.996	-0.125953\\
0.998	-0.125861\\
1	-0.125768\\
1.002	-0.125676\\
1.004	-0.125584\\
1.006	-0.125492\\
1.008	-0.125401\\
1.01	-0.125311\\
1.012	-0.125221\\
1.014	-0.125132\\
1.016	-0.125044\\
1.018	-0.124957\\
1.02	-0.124872\\
1.022	-0.124787\\
1.024	-0.124703\\
1.026	-0.124621\\
1.028	-0.124541\\
1.03	-0.124462\\
1.032	-0.124384\\
1.034	-0.124308\\
1.036	-0.124234\\
1.038	-0.124162\\
1.04	-0.124092\\
1.042	-0.124023\\
1.044	-0.123957\\
1.046	-0.123892\\
1.048	-0.12383\\
1.05	-0.12377\\
1.052	-0.123712\\
1.054	-0.123656\\
1.056	-0.123602\\
1.058	-0.123551\\
1.06	-0.123502\\
1.062	-0.123456\\
1.064	-0.123412\\
1.066	-0.12337\\
1.068	-0.123331\\
1.07	-0.123295\\
1.072	-0.123261\\
1.074	-0.123229\\
1.076	-0.1232\\
1.078	-0.123173\\
1.08	-0.123149\\
1.082	-0.123128\\
1.084	-0.123109\\
1.086	-0.123092\\
1.088	-0.123078\\
1.09	-0.123067\\
1.092	-0.123058\\
1.094	-0.123051\\
1.096	-0.123047\\
1.098	-0.123046\\
1.1	-0.123046\\
1.102	-0.123049\\
1.104	-0.123055\\
1.106	-0.123062\\
1.108	-0.123072\\
1.11	-0.123084\\
1.112	-0.123099\\
1.114	-0.123115\\
1.116	-0.123133\\
1.118	-0.123154\\
1.12	-0.123176\\
1.122	-0.123201\\
1.124	-0.123227\\
1.126	-0.123255\\
1.128	-0.123285\\
1.13	-0.123316\\
1.132	-0.123349\\
1.134	-0.123384\\
1.136	-0.12342\\
1.138	-0.123457\\
1.14	-0.123496\\
1.142	-0.123537\\
1.144	-0.123578\\
1.146	-0.123621\\
1.148	-0.123664\\
1.15	-0.123709\\
1.152	-0.123755\\
1.154	-0.123801\\
1.156	-0.123849\\
1.158	-0.123897\\
1.16	-0.123945\\
1.162	-0.123995\\
1.164	-0.124045\\
1.166	-0.124095\\
1.168	-0.124146\\
1.17	-0.124197\\
1.172	-0.124248\\
1.174	-0.124299\\
1.176	-0.12435\\
1.178	-0.124402\\
1.18	-0.124453\\
1.182	-0.124505\\
1.184	-0.124556\\
1.186	-0.124606\\
1.188	-0.124657\\
1.19	-0.124707\\
1.192	-0.124757\\
1.194	-0.124806\\
1.196	-0.124855\\
1.198	-0.124903\\
1.2	-0.12495\\
1.202	-0.124996\\
1.204	-0.125042\\
1.206	-0.125087\\
1.208	-0.125131\\
1.21	-0.125174\\
1.212	-0.125217\\
1.214	-0.125258\\
1.216	-0.125298\\
1.218	-0.125337\\
1.22	-0.125374\\
1.222	-0.125411\\
1.224	-0.125447\\
1.226	-0.125481\\
1.228	-0.125514\\
1.23	-0.125545\\
1.232	-0.125575\\
1.234	-0.125604\\
1.236	-0.125632\\
1.238	-0.125658\\
1.24	-0.125683\\
1.242	-0.125706\\
1.244	-0.125728\\
1.246	-0.125748\\
1.248	-0.125767\\
1.25	-0.125785\\
1.252	-0.125801\\
1.254	-0.125815\\
1.256	-0.125828\\
1.258	-0.12584\\
1.26	-0.12585\\
1.262	-0.125859\\
1.264	-0.125866\\
1.266	-0.125871\\
1.268	-0.125876\\
1.27	-0.125879\\
1.272	-0.12588\\
1.274	-0.12588\\
1.276	-0.125879\\
1.278	-0.125876\\
1.28	-0.125872\\
1.282	-0.125867\\
1.284	-0.125861\\
1.286	-0.125853\\
1.288	-0.125844\\
1.29	-0.125834\\
1.292	-0.125822\\
1.294	-0.12581\\
1.296	-0.125796\\
1.298	-0.125781\\
1.3	-0.125766\\
1.302	-0.125749\\
1.304	-0.125731\\
1.306	-0.125712\\
1.308	-0.125693\\
1.31	-0.125672\\
1.312	-0.125651\\
1.314	-0.125629\\
1.316	-0.125606\\
1.318	-0.125583\\
1.32	-0.125559\\
1.322	-0.125534\\
1.324	-0.125509\\
1.326	-0.125483\\
1.328	-0.125457\\
1.33	-0.12543\\
1.332	-0.125403\\
1.334	-0.125375\\
1.336	-0.125347\\
1.338	-0.125319\\
1.34	-0.125291\\
1.342	-0.125262\\
1.344	-0.125233\\
1.346	-0.125204\\
1.348	-0.125175\\
1.35	-0.125146\\
1.352	-0.125117\\
1.354	-0.125088\\
1.356	-0.125059\\
1.358	-0.12503\\
1.36	-0.125002\\
1.362	-0.124973\\
1.364	-0.124945\\
1.366	-0.124917\\
1.368	-0.124889\\
1.37	-0.124861\\
1.372	-0.124834\\
1.374	-0.124807\\
1.376	-0.124781\\
1.378	-0.124755\\
1.38	-0.124729\\
1.382	-0.124704\\
1.384	-0.12468\\
1.386	-0.124656\\
1.388	-0.124632\\
1.39	-0.124609\\
1.392	-0.124587\\
1.394	-0.124566\\
1.396	-0.124544\\
1.398	-0.124524\\
1.4	-0.124504\\
1.402	-0.124485\\
1.404	-0.124467\\
1.406	-0.124449\\
1.408	-0.124432\\
1.41	-0.124416\\
1.412	-0.124401\\
1.414	-0.124386\\
1.416	-0.124372\\
1.418	-0.124359\\
1.42	-0.124346\\
1.422	-0.124334\\
1.424	-0.124323\\
1.426	-0.124313\\
1.428	-0.124304\\
1.43	-0.124295\\
1.432	-0.124287\\
1.434	-0.12428\\
1.436	-0.124274\\
1.438	-0.124268\\
1.44	-0.124263\\
1.442	-0.124259\\
1.444	-0.124256\\
1.446	-0.124253\\
1.448	-0.124251\\
1.45	-0.12425\\
1.452	-0.12425\\
1.454	-0.12425\\
1.456	-0.124251\\
1.458	-0.124252\\
1.46	-0.124254\\
1.462	-0.124257\\
1.464	-0.12426\\
1.466	-0.124264\\
1.468	-0.124269\\
1.47	-0.124274\\
1.472	-0.12428\\
1.474	-0.124286\\
1.476	-0.124293\\
1.478	-0.1243\\
1.48	-0.124308\\
1.482	-0.124316\\
1.484	-0.124325\\
1.486	-0.124334\\
1.488	-0.124343\\
1.49	-0.124353\\
1.492	-0.124364\\
1.494	-0.124374\\
1.496	-0.124385\\
1.498	-0.124396\\
1.5	-0.124408\\
1.502	-0.124419\\
1.504	-0.124431\\
1.506	-0.124443\\
1.508	-0.124456\\
1.51	-0.124468\\
1.512	-0.124481\\
1.514	-0.124494\\
1.516	-0.124506\\
1.518	-0.124519\\
1.52	-0.124532\\
1.522	-0.124546\\
1.524	-0.124559\\
1.526	-0.124572\\
1.528	-0.124585\\
1.53	-0.124598\\
1.532	-0.124611\\
1.534	-0.124624\\
1.536	-0.124637\\
1.538	-0.124649\\
1.54	-0.124662\\
1.542	-0.124674\\
1.544	-0.124687\\
1.546	-0.124699\\
1.548	-0.124711\\
1.55	-0.124722\\
1.552	-0.124734\\
1.554	-0.124745\\
1.556	-0.124756\\
1.558	-0.124766\\
1.56	-0.124777\\
1.562	-0.124787\\
1.564	-0.124797\\
1.566	-0.124806\\
1.568	-0.124815\\
1.57	-0.124824\\
1.572	-0.124832\\
1.574	-0.12484\\
1.576	-0.124848\\
1.578	-0.124855\\
1.58	-0.124862\\
1.582	-0.124869\\
1.584	-0.124875\\
1.586	-0.12488\\
1.588	-0.124886\\
1.59	-0.12489\\
1.592	-0.124895\\
1.594	-0.124899\\
1.596	-0.124902\\
1.598	-0.124905\\
1.6	-0.124908\\
1.602	-0.12491\\
1.604	-0.124912\\
1.606	-0.124913\\
1.608	-0.124914\\
1.61	-0.124915\\
1.612	-0.124915\\
1.614	-0.124914\\
1.616	-0.124913\\
1.618	-0.124912\\
1.62	-0.124911\\
1.622	-0.124908\\
1.624	-0.124906\\
1.626	-0.124903\\
1.628	-0.1249\\
1.63	-0.124896\\
1.632	-0.124892\\
1.634	-0.124887\\
1.636	-0.124882\\
1.638	-0.124877\\
1.64	-0.124871\\
1.642	-0.124865\\
1.644	-0.124859\\
1.646	-0.124852\\
1.648	-0.124845\\
1.65	-0.124838\\
1.652	-0.12483\\
1.654	-0.124822\\
1.656	-0.124814\\
1.658	-0.124805\\
1.66	-0.124796\\
1.662	-0.124787\\
1.664	-0.124778\\
1.666	-0.124769\\
1.668	-0.124759\\
1.67	-0.124749\\
1.672	-0.124739\\
1.674	-0.124729\\
1.676	-0.124718\\
1.678	-0.124707\\
1.68	-0.124697\\
1.682	-0.124686\\
1.684	-0.124675\\
1.686	-0.124664\\
1.688	-0.124652\\
1.69	-0.124641\\
1.692	-0.12463\\
1.694	-0.124618\\
1.696	-0.124607\\
1.698	-0.124596\\
1.7	-0.124584\\
1.702	-0.124573\\
1.704	-0.124561\\
1.706	-0.12455\\
1.708	-0.124538\\
1.71	-0.124527\\
1.712	-0.124516\\
1.714	-0.124504\\
1.716	-0.124493\\
1.718	-0.124482\\
1.72	-0.124471\\
1.722	-0.124461\\
1.724	-0.12445\\
1.726	-0.124439\\
1.728	-0.124429\\
1.73	-0.124419\\
1.732	-0.124409\\
1.734	-0.124399\\
1.736	-0.124389\\
1.738	-0.124379\\
1.74	-0.12437\\
1.742	-0.124361\\
1.744	-0.124352\\
1.746	-0.124343\\
1.748	-0.124335\\
1.75	-0.124327\\
1.752	-0.124319\\
1.754	-0.124311\\
1.756	-0.124303\\
1.758	-0.124296\\
1.76	-0.124289\\
1.762	-0.124282\\
1.764	-0.124276\\
1.766	-0.12427\\
1.768	-0.124264\\
1.77	-0.124258\\
1.772	-0.124253\\
1.774	-0.124247\\
1.776	-0.124243\\
1.778	-0.124238\\
1.78	-0.124234\\
1.782	-0.12423\\
1.784	-0.124226\\
1.786	-0.124222\\
1.788	-0.124219\\
1.79	-0.124216\\
1.792	-0.124213\\
1.794	-0.124211\\
1.796	-0.124209\\
1.798	-0.124207\\
1.8	-0.124205\\
1.802	-0.124204\\
1.804	-0.124203\\
1.806	-0.124202\\
1.808	-0.124201\\
1.81	-0.124201\\
1.812	-0.1242\\
1.814	-0.1242\\
1.816	-0.124201\\
1.818	-0.124201\\
1.82	-0.124202\\
1.822	-0.124203\\
1.824	-0.124204\\
1.826	-0.124205\\
1.828	-0.124206\\
1.83	-0.124208\\
1.832	-0.124209\\
1.834	-0.124211\\
1.836	-0.124213\\
1.838	-0.124216\\
1.84	-0.124218\\
1.842	-0.12422\\
1.844	-0.124223\\
1.846	-0.124226\\
1.848	-0.124228\\
1.85	-0.124231\\
1.852	-0.124234\\
1.854	-0.124237\\
1.856	-0.12424\\
1.858	-0.124244\\
1.86	-0.124247\\
1.862	-0.12425\\
1.864	-0.124254\\
1.866	-0.124257\\
1.868	-0.12426\\
1.87	-0.124264\\
1.872	-0.124267\\
1.874	-0.124271\\
1.876	-0.124274\\
1.878	-0.124277\\
1.88	-0.124281\\
1.882	-0.124284\\
1.884	-0.124288\\
1.886	-0.124291\\
1.888	-0.124294\\
1.89	-0.124297\\
1.892	-0.1243\\
1.894	-0.124303\\
1.896	-0.124306\\
1.898	-0.124309\\
1.9	-0.124312\\
1.902	-0.124315\\
1.904	-0.124318\\
1.906	-0.12432\\
1.908	-0.124323\\
1.91	-0.124325\\
1.912	-0.124327\\
1.914	-0.124329\\
1.916	-0.124331\\
1.918	-0.124333\\
1.92	-0.124335\\
1.922	-0.124337\\
1.924	-0.124338\\
1.926	-0.12434\\
1.928	-0.124341\\
1.93	-0.124342\\
1.932	-0.124343\\
1.934	-0.124344\\
1.936	-0.124344\\
1.938	-0.124345\\
1.94	-0.124345\\
1.942	-0.124346\\
1.944	-0.124346\\
1.946	-0.124346\\
1.948	-0.124345\\
1.95	-0.124345\\
1.952	-0.124345\\
1.954	-0.124344\\
1.956	-0.124343\\
1.958	-0.124342\\
1.96	-0.124341\\
1.962	-0.12434\\
1.964	-0.124338\\
1.966	-0.124337\\
1.968	-0.124335\\
1.97	-0.124333\\
1.972	-0.124331\\
1.974	-0.124329\\
1.976	-0.124327\\
1.978	-0.124325\\
1.98	-0.124322\\
1.982	-0.12432\\
1.984	-0.124317\\
1.986	-0.124314\\
1.988	-0.124311\\
1.99	-0.124308\\
1.992	-0.124305\\
1.994	-0.124302\\
1.996	-0.124299\\
1.998	-0.124295\\
2	-0.124292\\
};
\addplot [color=mycolor9, line width=1.0pt, forget plot]
  table[row sep=crcr]{%
-0.124	0\\
-0.122	0\\
-0.12	0\\
-0.118	0\\
-0.116	0\\
-0.114	0\\
-0.112	0\\
-0.11	0\\
-0.108	0\\
-0.106	0\\
-0.104	0\\
-0.102	0\\
-0.1	0\\
-0.098	0\\
-0.096	0\\
-0.094	0\\
-0.092	0\\
-0.09	0\\
-0.088	0\\
-0.086	0\\
-0.084	0\\
-0.082	0\\
-0.08	0\\
-0.078	0\\
-0.076	0\\
-0.074	0\\
-0.072	0\\
-0.07	0\\
-0.068	0\\
-0.066	0\\
-0.064	0\\
-0.062	0\\
-0.06	0\\
-0.058	0\\
-0.056	0\\
-0.054	0\\
-0.052	0\\
-0.05	0\\
-0.048	0\\
-0.046	0\\
-0.044	0\\
-0.042	0\\
-0.04	0\\
-0.038	0\\
-0.036	0\\
-0.034	0\\
-0.032	0\\
-0.03	0\\
-0.028	0\\
-0.026	0\\
-0.024	0\\
-0.022	0\\
-0.02	0\\
-0.018	0\\
-0.016	0\\
-0.014	0\\
-0.012	0\\
-0.01	0\\
-0.008	0\\
-0.006	0\\
-0.004	0\\
-0.002	0\\
0	0\\
0.002	-0.014631\\
0.004	-0.0213542\\
0.006	-0.0235765\\
0.008	-0.0234252\\
0.01	-0.0221853\\
0.012	-0.0205973\\
0.014	-0.0190586\\
0.016	-0.0177578\\
0.018	-0.0167624\\
0.02	-0.0160746\\
0.022	-0.0156654\\
0.024	-0.0154953\\
0.026	-0.0155244\\
0.028	-0.0157177\\
0.03	-0.016047\\
0.032	-0.0164906\\
0.034	-0.0170322\\
0.036	-0.0176596\\
0.038	-0.0183638\\
0.04	-0.0191374\\
0.042	-0.0199745\\
0.044	-0.0208695\\
0.046	-0.0218172\\
0.048	-0.0228127\\
0.05	-0.023851\\
0.052	-0.0249273\\
0.054	-0.0260368\\
0.056	-0.0271751\\
0.058	-0.0283381\\
0.06	-0.029522\\
0.062	-0.0307232\\
0.064	-0.0319386\\
0.066	-0.0331655\\
0.068	-0.0344015\\
0.07	-0.0356446\\
0.072	-0.0368929\\
0.074	-0.0381451\\
0.076	-0.0393997\\
0.078	-0.0406559\\
0.08	-0.0419125\\
0.082	-0.043169\\
0.084	-0.0444245\\
0.086	-0.0456786\\
0.088	-0.0469307\\
0.09	-0.0481804\\
0.092	-0.0494272\\
0.094	-0.0506709\\
0.096	-0.0519109\\
0.098	-0.053147\\
0.1	-0.0543788\\
0.102	-0.0556061\\
0.104	-0.0568284\\
0.106	-0.0580454\\
0.108	-0.0592567\\
0.11	-0.0604621\\
0.112	-0.0616611\\
0.114	-0.0628535\\
0.116	-0.0640387\\
0.118	-0.0652164\\
0.12	-0.0663862\\
0.122	-0.0675477\\
0.124	-0.0687004\\
0.126	-0.0698439\\
0.128	-0.0709778\\
0.13	-0.0721017\\
0.132	-0.073215\\
0.134	-0.0743173\\
0.136	-0.0754081\\
0.138	-0.0764869\\
0.14	-0.0775534\\
0.142	-0.0786069\\
0.144	-0.079647\\
0.146	-0.0806733\\
0.148	-0.0816852\\
0.15	-0.0826823\\
0.152	-0.0836642\\
0.154	-0.0846303\\
0.156	-0.0855802\\
0.158	-0.0865135\\
0.16	-0.0874297\\
0.162	-0.0883285\\
0.164	-0.0892093\\
0.166	-0.0900718\\
0.168	-0.0909157\\
0.17	-0.0917405\\
0.172	-0.0925459\\
0.174	-0.0933316\\
0.176	-0.0940973\\
0.178	-0.0948426\\
0.18	-0.0955673\\
0.182	-0.0962711\\
0.184	-0.0969538\\
0.186	-0.0976152\\
0.188	-0.098255\\
0.19	-0.0988732\\
0.192	-0.0994696\\
0.194	-0.100044\\
0.196	-0.100596\\
0.198	-0.101127\\
0.2	-0.101635\\
0.202	-0.10212\\
0.204	-0.102584\\
0.206	-0.103025\\
0.208	-0.103445\\
0.21	-0.103842\\
0.212	-0.104217\\
0.214	-0.10457\\
0.216	-0.104902\\
0.218	-0.105212\\
0.22	-0.105501\\
0.222	-0.105769\\
0.224	-0.106015\\
0.226	-0.106241\\
0.228	-0.106447\\
0.23	-0.106633\\
0.232	-0.106799\\
0.234	-0.106945\\
0.236	-0.107073\\
0.238	-0.107181\\
0.24	-0.107272\\
0.242	-0.107345\\
0.244	-0.1074\\
0.246	-0.107439\\
0.248	-0.107461\\
0.25	-0.107467\\
0.252	-0.107457\\
0.254	-0.107433\\
0.256	-0.107394\\
0.258	-0.107342\\
0.26	-0.107276\\
0.262	-0.107197\\
0.264	-0.107106\\
0.266	-0.107004\\
0.268	-0.10689\\
0.27	-0.106766\\
0.272	-0.106633\\
0.274	-0.10649\\
0.276	-0.106338\\
0.278	-0.106179\\
0.28	-0.106012\\
0.282	-0.105839\\
0.284	-0.105659\\
0.286	-0.105474\\
0.288	-0.105284\\
0.29	-0.10509\\
0.292	-0.104893\\
0.294	-0.104692\\
0.296	-0.104489\\
0.298	-0.104285\\
0.3	-0.104079\\
0.302	-0.103873\\
0.304	-0.103667\\
0.306	-0.103461\\
0.308	-0.103257\\
0.31	-0.103055\\
0.312	-0.102855\\
0.314	-0.102658\\
0.316	-0.102464\\
0.318	-0.102275\\
0.32	-0.10209\\
0.322	-0.10191\\
0.324	-0.101736\\
0.326	-0.101567\\
0.328	-0.101406\\
0.33	-0.101251\\
0.332	-0.101104\\
0.334	-0.100964\\
0.336	-0.100833\\
0.338	-0.10071\\
0.34	-0.100597\\
0.342	-0.100493\\
0.344	-0.100398\\
0.346	-0.100314\\
0.348	-0.10024\\
0.35	-0.100176\\
0.352	-0.100124\\
0.354	-0.100083\\
0.356	-0.100053\\
0.358	-0.100034\\
0.36	-0.100028\\
0.362	-0.100034\\
0.364	-0.100052\\
0.366	-0.100082\\
0.368	-0.100124\\
0.37	-0.10018\\
0.372	-0.100248\\
0.374	-0.100328\\
0.376	-0.100421\\
0.378	-0.100528\\
0.38	-0.100647\\
0.382	-0.100778\\
0.384	-0.100923\\
0.386	-0.10108\\
0.388	-0.10125\\
0.39	-0.101433\\
0.392	-0.101628\\
0.394	-0.101836\\
0.396	-0.102055\\
0.398	-0.102287\\
0.4	-0.102531\\
0.402	-0.102787\\
0.404	-0.103054\\
0.406	-0.103333\\
0.408	-0.103623\\
0.41	-0.103924\\
0.412	-0.104235\\
0.414	-0.104557\\
0.416	-0.104889\\
0.418	-0.105231\\
0.42	-0.105582\\
0.422	-0.105943\\
0.424	-0.106312\\
0.426	-0.106691\\
0.428	-0.107077\\
0.43	-0.107471\\
0.432	-0.107873\\
0.434	-0.108281\\
0.436	-0.108697\\
0.438	-0.109118\\
0.44	-0.109546\\
0.442	-0.109979\\
0.444	-0.110418\\
0.446	-0.110861\\
0.448	-0.111308\\
0.45	-0.111759\\
0.452	-0.112214\\
0.454	-0.112671\\
0.456	-0.113132\\
0.458	-0.113594\\
0.46	-0.114058\\
0.462	-0.114523\\
0.464	-0.114989\\
0.466	-0.115455\\
0.468	-0.115922\\
0.47	-0.116387\\
0.472	-0.116852\\
0.474	-0.117316\\
0.476	-0.117777\\
0.478	-0.118237\\
0.48	-0.118694\\
0.482	-0.119147\\
0.484	-0.119598\\
0.486	-0.120044\\
0.488	-0.120486\\
0.49	-0.120924\\
0.492	-0.121356\\
0.494	-0.121783\\
0.496	-0.122204\\
0.498	-0.122619\\
0.5	-0.123028\\
0.502	-0.123429\\
0.504	-0.123824\\
0.506	-0.124211\\
0.508	-0.12459\\
0.51	-0.124961\\
0.512	-0.125324\\
0.514	-0.125678\\
0.516	-0.126024\\
0.518	-0.12636\\
0.52	-0.126686\\
0.522	-0.127003\\
0.524	-0.12731\\
0.526	-0.127608\\
0.528	-0.127895\\
0.53	-0.128171\\
0.532	-0.128437\\
0.534	-0.128692\\
0.536	-0.128937\\
0.538	-0.12917\\
0.54	-0.129393\\
0.542	-0.129604\\
0.544	-0.129804\\
0.546	-0.129993\\
0.548	-0.13017\\
0.55	-0.130336\\
0.552	-0.130491\\
0.554	-0.130634\\
0.556	-0.130766\\
0.558	-0.130887\\
0.56	-0.130996\\
0.562	-0.131094\\
0.564	-0.131181\\
0.566	-0.131257\\
0.568	-0.131322\\
0.57	-0.131376\\
0.572	-0.131419\\
0.574	-0.131451\\
0.576	-0.131473\\
0.578	-0.131485\\
0.58	-0.131486\\
0.582	-0.131477\\
0.584	-0.131459\\
0.586	-0.13143\\
0.588	-0.131392\\
0.59	-0.131345\\
0.592	-0.131289\\
0.594	-0.131225\\
0.596	-0.131151\\
0.598	-0.13107\\
0.6	-0.13098\\
0.602	-0.130883\\
0.604	-0.130778\\
0.606	-0.130666\\
0.608	-0.130547\\
0.61	-0.130421\\
0.612	-0.130289\\
0.614	-0.130151\\
0.616	-0.130007\\
0.618	-0.129858\\
0.62	-0.129704\\
0.622	-0.129545\\
0.624	-0.129381\\
0.626	-0.129213\\
0.628	-0.129042\\
0.63	-0.128867\\
0.632	-0.128688\\
0.634	-0.128507\\
0.636	-0.128324\\
0.638	-0.128138\\
0.64	-0.12795\\
0.642	-0.127761\\
0.644	-0.12757\\
0.646	-0.127378\\
0.648	-0.127186\\
0.65	-0.126993\\
0.652	-0.126801\\
0.654	-0.126609\\
0.656	-0.126417\\
0.658	-0.126226\\
0.66	-0.126036\\
0.662	-0.125848\\
0.664	-0.125661\\
0.666	-0.125476\\
0.668	-0.125294\\
0.67	-0.125114\\
0.672	-0.124937\\
0.674	-0.124763\\
0.676	-0.124592\\
0.678	-0.124424\\
0.68	-0.12426\\
0.682	-0.1241\\
0.684	-0.123944\\
0.686	-0.123792\\
0.688	-0.123645\\
0.69	-0.123502\\
0.692	-0.123364\\
0.694	-0.123231\\
0.696	-0.123103\\
0.698	-0.12298\\
0.7	-0.122862\\
0.702	-0.12275\\
0.704	-0.122644\\
0.706	-0.122543\\
0.708	-0.122448\\
0.71	-0.122359\\
0.712	-0.122276\\
0.714	-0.122198\\
0.716	-0.122127\\
0.718	-0.122062\\
0.72	-0.122002\\
0.722	-0.121949\\
0.724	-0.121902\\
0.726	-0.121861\\
0.728	-0.121826\\
0.73	-0.121798\\
0.732	-0.121775\\
0.734	-0.121758\\
0.736	-0.121748\\
0.738	-0.121743\\
0.74	-0.121744\\
0.742	-0.121751\\
0.744	-0.121764\\
0.746	-0.121783\\
0.748	-0.121807\\
0.75	-0.121836\\
0.752	-0.121871\\
0.754	-0.121912\\
0.756	-0.121957\\
0.758	-0.122007\\
0.76	-0.122063\\
0.762	-0.122123\\
0.764	-0.122188\\
0.766	-0.122257\\
0.768	-0.122331\\
0.77	-0.122408\\
0.772	-0.12249\\
0.774	-0.122576\\
0.776	-0.122666\\
0.778	-0.122759\\
0.78	-0.122855\\
0.782	-0.122955\\
0.784	-0.123058\\
0.786	-0.123163\\
0.788	-0.123272\\
0.79	-0.123382\\
0.792	-0.123495\\
0.794	-0.123611\\
0.796	-0.123728\\
0.798	-0.123847\\
0.8	-0.123967\\
0.802	-0.124089\\
0.804	-0.124212\\
0.806	-0.124337\\
0.808	-0.124462\\
0.81	-0.124587\\
0.812	-0.124713\\
0.814	-0.12484\\
0.816	-0.124966\\
0.818	-0.125093\\
0.82	-0.125219\\
0.822	-0.125345\\
0.824	-0.12547\\
0.826	-0.125594\\
0.828	-0.125718\\
0.83	-0.12584\\
0.832	-0.125961\\
0.834	-0.126081\\
0.836	-0.1262\\
0.838	-0.126316\\
0.84	-0.126431\\
0.842	-0.126544\\
0.844	-0.126654\\
0.846	-0.126763\\
0.848	-0.126869\\
0.85	-0.126972\\
0.852	-0.127073\\
0.854	-0.127172\\
0.856	-0.127267\\
0.858	-0.127359\\
0.86	-0.127449\\
0.862	-0.127535\\
0.864	-0.127618\\
0.866	-0.127698\\
0.868	-0.127775\\
0.87	-0.127848\\
0.872	-0.127917\\
0.874	-0.127983\\
0.876	-0.128046\\
0.878	-0.128104\\
0.88	-0.128159\\
0.882	-0.12821\\
0.884	-0.128258\\
0.886	-0.128301\\
0.888	-0.128341\\
0.89	-0.128377\\
0.892	-0.128409\\
0.894	-0.128437\\
0.896	-0.128461\\
0.898	-0.128481\\
0.9	-0.128498\\
0.902	-0.12851\\
0.904	-0.128519\\
0.906	-0.128524\\
0.908	-0.128525\\
0.91	-0.128522\\
0.912	-0.128515\\
0.914	-0.128505\\
0.916	-0.128491\\
0.918	-0.128473\\
0.92	-0.128452\\
0.922	-0.128427\\
0.924	-0.128399\\
0.926	-0.128368\\
0.928	-0.128333\\
0.93	-0.128295\\
0.932	-0.128253\\
0.934	-0.128209\\
0.936	-0.128161\\
0.938	-0.128111\\
0.94	-0.128057\\
0.942	-0.128001\\
0.944	-0.127942\\
0.946	-0.127881\\
0.948	-0.127817\\
0.95	-0.12775\\
0.952	-0.127682\\
0.954	-0.127611\\
0.956	-0.127538\\
0.958	-0.127463\\
0.96	-0.127386\\
0.962	-0.127307\\
0.964	-0.127226\\
0.966	-0.127144\\
0.968	-0.127061\\
0.97	-0.126976\\
0.972	-0.12689\\
0.974	-0.126803\\
0.976	-0.126714\\
0.978	-0.126625\\
0.98	-0.126535\\
0.982	-0.126444\\
0.984	-0.126353\\
0.986	-0.126261\\
0.988	-0.126169\\
0.99	-0.126077\\
0.992	-0.125984\\
0.994	-0.125892\\
0.996	-0.125799\\
0.998	-0.125707\\
1	-0.125615\\
1.002	-0.125523\\
1.004	-0.125432\\
1.006	-0.125341\\
1.008	-0.125251\\
1.01	-0.125162\\
1.012	-0.125074\\
1.014	-0.124987\\
1.016	-0.124901\\
1.018	-0.124816\\
1.02	-0.124732\\
1.022	-0.12465\\
1.024	-0.124568\\
1.026	-0.124489\\
1.028	-0.124411\\
1.03	-0.124335\\
1.032	-0.12426\\
1.034	-0.124187\\
1.036	-0.124116\\
1.038	-0.124047\\
1.04	-0.12398\\
1.042	-0.123915\\
1.044	-0.123852\\
1.046	-0.123791\\
1.048	-0.123732\\
1.05	-0.123676\\
1.052	-0.123621\\
1.054	-0.12357\\
1.056	-0.12352\\
1.058	-0.123473\\
1.06	-0.123428\\
1.062	-0.123386\\
1.064	-0.123346\\
1.066	-0.123308\\
1.068	-0.123274\\
1.07	-0.123241\\
1.072	-0.123211\\
1.074	-0.123184\\
1.076	-0.123159\\
1.078	-0.123137\\
1.08	-0.123117\\
1.082	-0.123099\\
1.084	-0.123085\\
1.086	-0.123072\\
1.088	-0.123063\\
1.09	-0.123055\\
1.092	-0.12305\\
1.094	-0.123048\\
1.096	-0.123048\\
1.098	-0.12305\\
1.1	-0.123054\\
1.102	-0.123061\\
1.104	-0.12307\\
1.106	-0.123082\\
1.108	-0.123095\\
1.11	-0.123111\\
1.112	-0.123129\\
1.114	-0.123148\\
1.116	-0.12317\\
1.118	-0.123194\\
1.12	-0.123219\\
1.122	-0.123247\\
1.124	-0.123276\\
1.126	-0.123307\\
1.128	-0.123339\\
1.13	-0.123373\\
1.132	-0.123409\\
1.134	-0.123446\\
1.136	-0.123484\\
1.138	-0.123524\\
1.14	-0.123565\\
1.142	-0.123607\\
1.144	-0.123651\\
1.146	-0.123695\\
1.148	-0.12374\\
1.15	-0.123786\\
1.152	-0.123834\\
1.154	-0.123881\\
1.156	-0.12393\\
1.158	-0.123979\\
1.16	-0.124029\\
1.162	-0.124079\\
1.164	-0.124129\\
1.166	-0.12418\\
1.168	-0.124231\\
1.17	-0.124282\\
1.172	-0.124334\\
1.174	-0.124385\\
1.176	-0.124436\\
1.178	-0.124488\\
1.18	-0.124539\\
1.182	-0.12459\\
1.184	-0.12464\\
1.186	-0.12469\\
1.188	-0.12474\\
1.19	-0.124789\\
1.192	-0.124838\\
1.194	-0.124886\\
1.196	-0.124934\\
1.198	-0.124981\\
1.2	-0.125027\\
1.202	-0.125072\\
1.204	-0.125116\\
1.206	-0.12516\\
1.208	-0.125202\\
1.21	-0.125243\\
1.212	-0.125284\\
1.214	-0.125323\\
1.216	-0.125361\\
1.218	-0.125398\\
1.22	-0.125434\\
1.222	-0.125469\\
1.224	-0.125502\\
1.226	-0.125534\\
1.228	-0.125565\\
1.23	-0.125594\\
1.232	-0.125622\\
1.234	-0.125648\\
1.236	-0.125674\\
1.238	-0.125697\\
1.24	-0.12572\\
1.242	-0.125741\\
1.244	-0.12576\\
1.246	-0.125778\\
1.248	-0.125794\\
1.25	-0.125809\\
1.252	-0.125823\\
1.254	-0.125835\\
1.256	-0.125846\\
1.258	-0.125855\\
1.26	-0.125862\\
1.262	-0.125869\\
1.264	-0.125873\\
1.266	-0.125877\\
1.268	-0.125879\\
1.27	-0.125879\\
1.272	-0.125879\\
1.274	-0.125876\\
1.276	-0.125873\\
1.278	-0.125868\\
1.28	-0.125862\\
1.282	-0.125855\\
1.284	-0.125846\\
1.286	-0.125836\\
1.288	-0.125825\\
1.29	-0.125813\\
1.292	-0.1258\\
1.294	-0.125785\\
1.296	-0.12577\\
1.298	-0.125754\\
1.3	-0.125736\\
1.302	-0.125718\\
1.304	-0.125699\\
1.306	-0.125679\\
1.308	-0.125658\\
1.31	-0.125636\\
1.312	-0.125613\\
1.314	-0.12559\\
1.316	-0.125566\\
1.318	-0.125542\\
1.32	-0.125517\\
1.322	-0.125491\\
1.324	-0.125465\\
1.326	-0.125438\\
1.328	-0.125411\\
1.33	-0.125384\\
1.332	-0.125356\\
1.334	-0.125328\\
1.336	-0.1253\\
1.338	-0.125271\\
1.34	-0.125243\\
1.342	-0.125214\\
1.344	-0.125185\\
1.346	-0.125156\\
1.348	-0.125127\\
1.35	-0.125098\\
1.352	-0.125069\\
1.354	-0.12504\\
1.356	-0.125011\\
1.358	-0.124983\\
1.36	-0.124954\\
1.362	-0.124926\\
1.364	-0.124898\\
1.366	-0.124871\\
1.368	-0.124843\\
1.37	-0.124816\\
1.372	-0.12479\\
1.374	-0.124764\\
1.376	-0.124738\\
1.378	-0.124713\\
1.38	-0.124688\\
1.382	-0.124664\\
1.384	-0.12464\\
1.386	-0.124617\\
1.388	-0.124595\\
1.39	-0.124573\\
1.392	-0.124552\\
1.394	-0.124531\\
1.396	-0.124511\\
1.398	-0.124492\\
1.4	-0.124473\\
1.402	-0.124455\\
1.404	-0.124438\\
1.406	-0.124422\\
1.408	-0.124406\\
1.41	-0.124391\\
1.412	-0.124377\\
1.414	-0.124363\\
1.416	-0.124351\\
1.418	-0.124339\\
1.42	-0.124328\\
1.422	-0.124317\\
1.424	-0.124307\\
1.426	-0.124299\\
1.428	-0.12429\\
1.43	-0.124283\\
1.432	-0.124276\\
1.434	-0.124271\\
1.436	-0.124265\\
1.438	-0.124261\\
1.44	-0.124257\\
1.442	-0.124255\\
1.444	-0.124252\\
1.446	-0.124251\\
1.448	-0.12425\\
1.45	-0.12425\\
1.452	-0.124251\\
1.454	-0.124252\\
1.456	-0.124254\\
1.458	-0.124257\\
1.46	-0.12426\\
1.462	-0.124264\\
1.464	-0.124268\\
1.466	-0.124273\\
1.468	-0.124278\\
1.47	-0.124285\\
1.472	-0.124291\\
1.474	-0.124298\\
1.476	-0.124306\\
1.478	-0.124314\\
1.48	-0.124322\\
1.482	-0.124331\\
1.484	-0.124341\\
1.486	-0.12435\\
1.488	-0.12436\\
1.49	-0.124371\\
1.492	-0.124382\\
1.494	-0.124393\\
1.496	-0.124404\\
1.498	-0.124416\\
1.5	-0.124427\\
1.502	-0.12444\\
1.504	-0.124452\\
1.506	-0.124464\\
1.508	-0.124477\\
1.51	-0.12449\\
1.512	-0.124502\\
1.514	-0.124515\\
1.516	-0.124528\\
1.518	-0.124541\\
1.52	-0.124554\\
1.522	-0.124567\\
1.524	-0.124581\\
1.526	-0.124594\\
1.528	-0.124607\\
1.53	-0.12462\\
1.532	-0.124632\\
1.534	-0.124645\\
1.536	-0.124658\\
1.538	-0.12467\\
1.54	-0.124682\\
1.542	-0.124695\\
1.544	-0.124706\\
1.546	-0.124718\\
1.548	-0.12473\\
1.55	-0.124741\\
1.552	-0.124752\\
1.554	-0.124763\\
1.556	-0.124773\\
1.558	-0.124783\\
1.56	-0.124793\\
1.562	-0.124803\\
1.564	-0.124812\\
1.566	-0.124821\\
1.568	-0.124829\\
1.57	-0.124837\\
1.572	-0.124845\\
1.574	-0.124853\\
1.576	-0.12486\\
1.578	-0.124866\\
1.58	-0.124872\\
1.582	-0.124878\\
1.584	-0.124884\\
1.586	-0.124889\\
1.588	-0.124893\\
1.59	-0.124897\\
1.592	-0.124901\\
1.594	-0.124904\\
1.596	-0.124907\\
1.598	-0.124909\\
1.6	-0.124911\\
1.602	-0.124913\\
1.604	-0.124914\\
1.606	-0.124914\\
1.608	-0.124915\\
1.61	-0.124914\\
1.612	-0.124914\\
1.614	-0.124912\\
1.616	-0.124911\\
1.618	-0.124909\\
1.62	-0.124906\\
1.622	-0.124904\\
1.624	-0.1249\\
1.626	-0.124897\\
1.628	-0.124893\\
1.63	-0.124888\\
1.632	-0.124884\\
1.634	-0.124878\\
1.636	-0.124873\\
1.638	-0.124867\\
1.64	-0.124861\\
1.642	-0.124854\\
1.644	-0.124847\\
1.646	-0.12484\\
1.648	-0.124832\\
1.65	-0.124825\\
1.652	-0.124816\\
1.654	-0.124808\\
1.656	-0.124799\\
1.658	-0.12479\\
1.66	-0.124781\\
1.662	-0.124772\\
1.664	-0.124762\\
1.666	-0.124752\\
1.668	-0.124742\\
1.67	-0.124732\\
1.672	-0.124721\\
1.674	-0.124711\\
1.676	-0.1247\\
1.678	-0.124689\\
1.68	-0.124678\\
1.682	-0.124667\\
1.684	-0.124656\\
1.686	-0.124645\\
1.688	-0.124634\\
1.69	-0.124622\\
1.692	-0.124611\\
1.694	-0.124599\\
1.696	-0.124588\\
1.698	-0.124576\\
1.7	-0.124565\\
1.702	-0.124554\\
1.704	-0.124542\\
1.706	-0.124531\\
1.708	-0.124519\\
1.71	-0.124508\\
1.712	-0.124497\\
1.714	-0.124486\\
1.716	-0.124475\\
1.718	-0.124464\\
1.72	-0.124454\\
1.722	-0.124443\\
1.724	-0.124433\\
1.726	-0.124422\\
1.728	-0.124412\\
1.73	-0.124402\\
1.732	-0.124392\\
1.734	-0.124383\\
1.736	-0.124373\\
1.738	-0.124364\\
1.74	-0.124355\\
1.742	-0.124346\\
1.744	-0.124338\\
1.746	-0.12433\\
1.748	-0.124321\\
1.75	-0.124314\\
1.752	-0.124306\\
1.754	-0.124299\\
1.756	-0.124292\\
1.758	-0.124285\\
1.76	-0.124278\\
1.762	-0.124272\\
1.764	-0.124266\\
1.766	-0.12426\\
1.768	-0.124255\\
1.77	-0.124249\\
1.772	-0.124244\\
1.774	-0.12424\\
1.776	-0.124235\\
1.778	-0.124231\\
1.78	-0.124227\\
1.782	-0.124224\\
1.784	-0.12422\\
1.786	-0.124217\\
1.788	-0.124214\\
1.79	-0.124212\\
1.792	-0.12421\\
1.794	-0.124208\\
1.796	-0.124206\\
1.798	-0.124204\\
1.8	-0.124203\\
1.802	-0.124202\\
1.804	-0.124201\\
1.806	-0.124201\\
1.808	-0.124201\\
1.81	-0.124201\\
1.812	-0.124201\\
1.814	-0.124201\\
1.816	-0.124202\\
1.818	-0.124202\\
1.82	-0.124203\\
1.822	-0.124205\\
1.824	-0.124206\\
1.826	-0.124207\\
1.828	-0.124209\\
1.83	-0.124211\\
1.832	-0.124213\\
1.834	-0.124215\\
1.836	-0.124217\\
1.838	-0.12422\\
1.84	-0.124222\\
1.842	-0.124225\\
1.844	-0.124228\\
1.846	-0.12423\\
1.848	-0.124233\\
1.85	-0.124236\\
1.852	-0.12424\\
1.854	-0.124243\\
1.856	-0.124246\\
1.858	-0.124249\\
1.86	-0.124253\\
1.862	-0.124256\\
1.864	-0.124259\\
1.866	-0.124263\\
1.868	-0.124266\\
1.87	-0.124269\\
1.872	-0.124273\\
1.874	-0.124276\\
1.876	-0.12428\\
1.878	-0.124283\\
1.88	-0.124286\\
1.882	-0.12429\\
1.884	-0.124293\\
1.886	-0.124296\\
1.888	-0.124299\\
1.89	-0.124302\\
1.892	-0.124305\\
1.894	-0.124308\\
1.896	-0.124311\\
1.898	-0.124314\\
1.9	-0.124317\\
1.902	-0.124319\\
1.904	-0.124322\\
1.906	-0.124324\\
1.908	-0.124327\\
1.91	-0.124329\\
1.912	-0.124331\\
1.914	-0.124333\\
1.916	-0.124334\\
1.918	-0.124336\\
1.92	-0.124338\\
1.922	-0.124339\\
1.924	-0.12434\\
1.926	-0.124342\\
1.928	-0.124343\\
1.93	-0.124343\\
1.932	-0.124344\\
1.934	-0.124345\\
1.936	-0.124345\\
1.938	-0.124345\\
1.94	-0.124346\\
1.942	-0.124346\\
1.944	-0.124345\\
1.946	-0.124345\\
1.948	-0.124345\\
1.95	-0.124344\\
1.952	-0.124343\\
1.954	-0.124342\\
1.956	-0.124341\\
1.958	-0.12434\\
1.96	-0.124339\\
1.962	-0.124337\\
1.964	-0.124336\\
1.966	-0.124334\\
1.968	-0.124332\\
1.97	-0.12433\\
1.972	-0.124328\\
1.974	-0.124326\\
1.976	-0.124323\\
1.978	-0.124321\\
1.98	-0.124318\\
1.982	-0.124315\\
1.984	-0.124312\\
1.986	-0.124309\\
1.988	-0.124306\\
1.99	-0.124303\\
1.992	-0.1243\\
1.994	-0.124296\\
1.996	-0.124293\\
1.998	-0.124289\\
2	-0.124285\\
};
\addplot [color=mycolor10, line width=1.0pt, forget plot]
  table[row sep=crcr]{%
-0.124	0\\
-0.122	0\\
-0.12	0\\
-0.118	0\\
-0.116	0\\
-0.114	0\\
-0.112	0\\
-0.11	0\\
-0.108	0\\
-0.106	0\\
-0.104	0\\
-0.102	0\\
-0.1	0\\
-0.098	0\\
-0.096	0\\
-0.094	0\\
-0.092	0\\
-0.09	0\\
-0.088	0\\
-0.086	0\\
-0.084	0\\
-0.082	0\\
-0.08	0\\
-0.078	0\\
-0.076	0\\
-0.074	0\\
-0.072	0\\
-0.07	0\\
-0.068	0\\
-0.066	0\\
-0.064	0\\
-0.062	0\\
-0.06	0\\
-0.058	0\\
-0.056	0\\
-0.054	0\\
-0.052	0\\
-0.05	0\\
-0.048	0\\
-0.046	0\\
-0.044	0\\
-0.042	0\\
-0.04	0\\
-0.038	0\\
-0.036	0\\
-0.034	0\\
-0.032	0\\
-0.03	0\\
-0.028	0\\
-0.026	0\\
-0.024	0\\
-0.022	0\\
-0.02	0\\
-0.018	0\\
-0.016	0\\
-0.014	0\\
-0.012	0\\
-0.01	0\\
-0.008	0\\
-0.006	0\\
-0.004	0\\
-0.002	0\\
0	0\\
0.002	-0.000641175\\
0.004	-0.00129733\\
0.006	-0.00198055\\
0.008	-0.00269567\\
0.01	-0.0034434\\
0.012	-0.0042223\\
0.014	-0.00503\\
0.016	-0.00586385\\
0.018	-0.00672131\\
0.02	-0.00760015\\
0.022	-0.0084985\\
0.024	-0.00941487\\
0.026	-0.0103481\\
0.028	-0.0112973\\
0.03	-0.0122618\\
0.032	-0.0132414\\
0.034	-0.0142356\\
0.036	-0.0152445\\
0.038	-0.016268\\
0.04	-0.0173062\\
0.042	-0.0183593\\
0.044	-0.0194274\\
0.046	-0.0205108\\
0.048	-0.0216098\\
0.05	-0.0227245\\
0.052	-0.0238552\\
0.054	-0.0250021\\
0.056	-0.0261653\\
0.058	-0.027345\\
0.06	-0.0285413\\
0.062	-0.0297542\\
0.064	-0.0309837\\
0.066	-0.0322297\\
0.068	-0.033492\\
0.07	-0.0347705\\
0.072	-0.036065\\
0.074	-0.037375\\
0.076	-0.0387001\\
0.078	-0.04004\\
0.08	-0.041394\\
0.082	-0.0427617\\
0.084	-0.0441422\\
0.086	-0.045535\\
0.088	-0.0469392\\
0.09	-0.048354\\
0.092	-0.0497786\\
0.094	-0.051212\\
0.096	-0.0526533\\
0.098	-0.0541013\\
0.1	-0.0555552\\
0.102	-0.0570138\\
0.104	-0.0584759\\
0.106	-0.0599405\\
0.108	-0.0614063\\
0.11	-0.0628721\\
0.112	-0.0643369\\
0.114	-0.0657992\\
0.116	-0.0672578\\
0.118	-0.0687116\\
0.12	-0.0701592\\
0.122	-0.0715995\\
0.124	-0.073031\\
0.126	-0.0744527\\
0.128	-0.0758632\\
0.13	-0.0772613\\
0.132	-0.0786459\\
0.134	-0.0800156\\
0.136	-0.0813693\\
0.138	-0.0827059\\
0.14	-0.0840242\\
0.142	-0.0853231\\
0.144	-0.0866015\\
0.146	-0.0878585\\
0.148	-0.0890929\\
0.15	-0.0903037\\
0.152	-0.0914902\\
0.154	-0.0926512\\
0.156	-0.0937861\\
0.158	-0.0948939\\
0.16	-0.0959739\\
0.162	-0.0970254\\
0.164	-0.0980476\\
0.166	-0.09904\\
0.168	-0.100002\\
0.17	-0.100933\\
0.172	-0.101832\\
0.174	-0.102699\\
0.176	-0.103534\\
0.178	-0.104337\\
0.18	-0.105106\\
0.182	-0.105842\\
0.184	-0.106545\\
0.186	-0.107213\\
0.188	-0.107849\\
0.19	-0.10845\\
0.192	-0.109018\\
0.194	-0.109551\\
0.196	-0.110051\\
0.198	-0.110518\\
0.2	-0.110951\\
0.202	-0.111351\\
0.204	-0.111718\\
0.206	-0.112053\\
0.208	-0.112355\\
0.21	-0.112625\\
0.212	-0.112864\\
0.214	-0.113072\\
0.216	-0.113249\\
0.218	-0.113397\\
0.22	-0.113514\\
0.222	-0.113604\\
0.224	-0.113664\\
0.226	-0.113698\\
0.228	-0.113704\\
0.23	-0.113684\\
0.232	-0.113639\\
0.234	-0.113569\\
0.236	-0.113475\\
0.238	-0.113358\\
0.24	-0.113219\\
0.242	-0.113058\\
0.244	-0.112877\\
0.246	-0.112675\\
0.248	-0.112455\\
0.25	-0.112217\\
0.252	-0.111961\\
0.254	-0.111689\\
0.256	-0.111402\\
0.258	-0.111101\\
0.26	-0.110786\\
0.262	-0.110458\\
0.264	-0.110118\\
0.266	-0.109768\\
0.268	-0.109408\\
0.27	-0.109039\\
0.272	-0.108661\\
0.274	-0.108277\\
0.276	-0.107886\\
0.278	-0.10749\\
0.28	-0.10709\\
0.282	-0.106686\\
0.284	-0.106279\\
0.286	-0.10587\\
0.288	-0.105461\\
0.29	-0.105051\\
0.292	-0.104641\\
0.294	-0.104234\\
0.296	-0.103828\\
0.298	-0.103425\\
0.3	-0.103027\\
0.302	-0.102632\\
0.304	-0.102243\\
0.306	-0.10186\\
0.308	-0.101484\\
0.31	-0.101115\\
0.312	-0.100754\\
0.314	-0.100401\\
0.316	-0.100058\\
0.318	-0.099725\\
0.32	-0.0994022\\
0.322	-0.0990905\\
0.324	-0.0987903\\
0.326	-0.0985022\\
0.328	-0.0982267\\
0.33	-0.0979642\\
0.332	-0.0977153\\
0.334	-0.0974803\\
0.336	-0.0972597\\
0.338	-0.0970539\\
0.34	-0.0968632\\
0.342	-0.096688\\
0.344	-0.0965286\\
0.346	-0.0963854\\
0.348	-0.0962585\\
0.35	-0.0961483\\
0.352	-0.096055\\
0.354	-0.0959788\\
0.356	-0.09592\\
0.358	-0.0958786\\
0.36	-0.0958549\\
0.362	-0.0958489\\
0.364	-0.0958608\\
0.366	-0.0958907\\
0.368	-0.0959386\\
0.37	-0.0960046\\
0.372	-0.0960886\\
0.374	-0.0961908\\
0.376	-0.096311\\
0.378	-0.0964492\\
0.38	-0.0966054\\
0.382	-0.0967794\\
0.384	-0.0969712\\
0.386	-0.0971807\\
0.388	-0.0974076\\
0.39	-0.0976519\\
0.392	-0.0979133\\
0.394	-0.0981917\\
0.396	-0.0984867\\
0.398	-0.0987982\\
0.4	-0.0991259\\
0.402	-0.0994696\\
0.404	-0.0998288\\
0.406	-0.100203\\
0.408	-0.100593\\
0.41	-0.100997\\
0.412	-0.101415\\
0.414	-0.101848\\
0.416	-0.102293\\
0.418	-0.102752\\
0.42	-0.103223\\
0.422	-0.103707\\
0.424	-0.104202\\
0.426	-0.104709\\
0.428	-0.105226\\
0.43	-0.105753\\
0.432	-0.106291\\
0.434	-0.106837\\
0.436	-0.107392\\
0.438	-0.107955\\
0.44	-0.108526\\
0.442	-0.109104\\
0.444	-0.109688\\
0.446	-0.110278\\
0.448	-0.110873\\
0.45	-0.111473\\
0.452	-0.112077\\
0.454	-0.112685\\
0.456	-0.113295\\
0.458	-0.113908\\
0.46	-0.114522\\
0.462	-0.115138\\
0.464	-0.115753\\
0.466	-0.116369\\
0.468	-0.116984\\
0.47	-0.117597\\
0.472	-0.118208\\
0.474	-0.118816\\
0.476	-0.119421\\
0.478	-0.120023\\
0.48	-0.120619\\
0.482	-0.121211\\
0.484	-0.121797\\
0.486	-0.122376\\
0.488	-0.122949\\
0.49	-0.123514\\
0.492	-0.124072\\
0.494	-0.124621\\
0.496	-0.12516\\
0.498	-0.125691\\
0.5	-0.126211\\
0.502	-0.12672\\
0.504	-0.127219\\
0.506	-0.127706\\
0.508	-0.128181\\
0.51	-0.128643\\
0.512	-0.129093\\
0.514	-0.12953\\
0.516	-0.129953\\
0.518	-0.130362\\
0.52	-0.130756\\
0.522	-0.131136\\
0.524	-0.131502\\
0.526	-0.131851\\
0.528	-0.132186\\
0.53	-0.132504\\
0.532	-0.132807\\
0.534	-0.133094\\
0.536	-0.133364\\
0.538	-0.133618\\
0.54	-0.133855\\
0.542	-0.134075\\
0.544	-0.134279\\
0.546	-0.134466\\
0.548	-0.134636\\
0.55	-0.134789\\
0.552	-0.134925\\
0.554	-0.135044\\
0.556	-0.135146\\
0.558	-0.135231\\
0.56	-0.1353\\
0.562	-0.135353\\
0.564	-0.135388\\
0.566	-0.135408\\
0.568	-0.135412\\
0.57	-0.135399\\
0.572	-0.135371\\
0.574	-0.135328\\
0.576	-0.13527\\
0.578	-0.135196\\
0.58	-0.135108\\
0.582	-0.135006\\
0.584	-0.13489\\
0.586	-0.134761\\
0.588	-0.134618\\
0.59	-0.134463\\
0.592	-0.134295\\
0.594	-0.134116\\
0.596	-0.133925\\
0.598	-0.133722\\
0.6	-0.13351\\
0.602	-0.133287\\
0.604	-0.133054\\
0.606	-0.132812\\
0.608	-0.132562\\
0.61	-0.132303\\
0.612	-0.132037\\
0.614	-0.131764\\
0.616	-0.131483\\
0.618	-0.131197\\
0.62	-0.130905\\
0.622	-0.130609\\
0.624	-0.130307\\
0.626	-0.130002\\
0.628	-0.129693\\
0.63	-0.129382\\
0.632	-0.129068\\
0.634	-0.128752\\
0.636	-0.128434\\
0.638	-0.128116\\
0.64	-0.127798\\
0.642	-0.12748\\
0.644	-0.127163\\
0.646	-0.126846\\
0.648	-0.126532\\
0.65	-0.12622\\
0.652	-0.12591\\
0.654	-0.125604\\
0.656	-0.125301\\
0.658	-0.125002\\
0.66	-0.124708\\
0.662	-0.124418\\
0.664	-0.124134\\
0.666	-0.123855\\
0.668	-0.123582\\
0.67	-0.123316\\
0.672	-0.123056\\
0.674	-0.122804\\
0.676	-0.122559\\
0.678	-0.122321\\
0.68	-0.122092\\
0.682	-0.12187\\
0.684	-0.121657\\
0.686	-0.121453\\
0.688	-0.121257\\
0.69	-0.121071\\
0.692	-0.120894\\
0.694	-0.120726\\
0.696	-0.120568\\
0.698	-0.12042\\
0.7	-0.120281\\
0.702	-0.120152\\
0.704	-0.120034\\
0.706	-0.119925\\
0.708	-0.119826\\
0.71	-0.119738\\
0.712	-0.119659\\
0.714	-0.119591\\
0.716	-0.119533\\
0.718	-0.119484\\
0.72	-0.119446\\
0.722	-0.119417\\
0.724	-0.119399\\
0.726	-0.11939\\
0.728	-0.119391\\
0.73	-0.119401\\
0.732	-0.11942\\
0.734	-0.119449\\
0.736	-0.119487\\
0.738	-0.119533\\
0.74	-0.119589\\
0.742	-0.119652\\
0.744	-0.119724\\
0.746	-0.119804\\
0.748	-0.119892\\
0.75	-0.119987\\
0.752	-0.12009\\
0.754	-0.1202\\
0.756	-0.120316\\
0.758	-0.12044\\
0.76	-0.120569\\
0.762	-0.120704\\
0.764	-0.120846\\
0.766	-0.120992\\
0.768	-0.121144\\
0.77	-0.1213\\
0.772	-0.121462\\
0.774	-0.121627\\
0.776	-0.121796\\
0.778	-0.121969\\
0.78	-0.122146\\
0.782	-0.122325\\
0.784	-0.122507\\
0.786	-0.122692\\
0.788	-0.122879\\
0.79	-0.123068\\
0.792	-0.123258\\
0.794	-0.12345\\
0.796	-0.123642\\
0.798	-0.123835\\
0.8	-0.124029\\
0.802	-0.124223\\
0.804	-0.124417\\
0.806	-0.12461\\
0.808	-0.124803\\
0.81	-0.124994\\
0.812	-0.125185\\
0.814	-0.125374\\
0.816	-0.125562\\
0.818	-0.125747\\
0.82	-0.125931\\
0.822	-0.126112\\
0.824	-0.12629\\
0.826	-0.126466\\
0.828	-0.126639\\
0.83	-0.126808\\
0.832	-0.126974\\
0.834	-0.127137\\
0.836	-0.127296\\
0.838	-0.127451\\
0.84	-0.127602\\
0.842	-0.127749\\
0.844	-0.127892\\
0.846	-0.12803\\
0.848	-0.128164\\
0.85	-0.128293\\
0.852	-0.128417\\
0.854	-0.128536\\
0.856	-0.128651\\
0.858	-0.12876\\
0.86	-0.128864\\
0.862	-0.128963\\
0.864	-0.129056\\
0.866	-0.129145\\
0.868	-0.129228\\
0.87	-0.129305\\
0.872	-0.129377\\
0.874	-0.129444\\
0.876	-0.129505\\
0.878	-0.12956\\
0.88	-0.12961\\
0.882	-0.129655\\
0.884	-0.129694\\
0.886	-0.129728\\
0.888	-0.129756\\
0.89	-0.129778\\
0.892	-0.129796\\
0.894	-0.129808\\
0.896	-0.129814\\
0.898	-0.129816\\
0.9	-0.129812\\
0.902	-0.129803\\
0.904	-0.129789\\
0.906	-0.12977\\
0.908	-0.129746\\
0.91	-0.129717\\
0.912	-0.129683\\
0.914	-0.129645\\
0.916	-0.129602\\
0.918	-0.129554\\
0.92	-0.129503\\
0.922	-0.129446\\
0.924	-0.129386\\
0.926	-0.129321\\
0.928	-0.129253\\
0.93	-0.129181\\
0.932	-0.129105\\
0.934	-0.129025\\
0.936	-0.128942\\
0.938	-0.128855\\
0.94	-0.128765\\
0.942	-0.128672\\
0.944	-0.128577\\
0.946	-0.128478\\
0.948	-0.128376\\
0.95	-0.128272\\
0.952	-0.128165\\
0.954	-0.128057\\
0.956	-0.127945\\
0.958	-0.127832\\
0.96	-0.127717\\
0.962	-0.127601\\
0.964	-0.127482\\
0.966	-0.127362\\
0.968	-0.127241\\
0.97	-0.127119\\
0.972	-0.126995\\
0.974	-0.126871\\
0.976	-0.126746\\
0.978	-0.12662\\
0.98	-0.126494\\
0.982	-0.126368\\
0.984	-0.126241\\
0.986	-0.126114\\
0.988	-0.125987\\
0.99	-0.125861\\
0.992	-0.125735\\
0.994	-0.125609\\
0.996	-0.125484\\
0.998	-0.12536\\
1	-0.125237\\
1.002	-0.125114\\
1.004	-0.124993\\
1.006	-0.124873\\
1.008	-0.124755\\
1.01	-0.124638\\
1.012	-0.124522\\
1.014	-0.124409\\
1.016	-0.124297\\
1.018	-0.124187\\
1.02	-0.124079\\
1.022	-0.123974\\
1.024	-0.12387\\
1.026	-0.12377\\
1.028	-0.123671\\
1.03	-0.123575\\
1.032	-0.123482\\
1.034	-0.123391\\
1.036	-0.123304\\
1.038	-0.123219\\
1.04	-0.123137\\
1.042	-0.123059\\
1.044	-0.122983\\
1.046	-0.12291\\
1.048	-0.122841\\
1.05	-0.122775\\
1.052	-0.122713\\
1.054	-0.122653\\
1.056	-0.122598\\
1.058	-0.122545\\
1.06	-0.122496\\
1.062	-0.122451\\
1.064	-0.122409\\
1.066	-0.122371\\
1.068	-0.122336\\
1.07	-0.122305\\
1.072	-0.122278\\
1.074	-0.122254\\
1.076	-0.122234\\
1.078	-0.122217\\
1.08	-0.122204\\
1.082	-0.122195\\
1.084	-0.122189\\
1.086	-0.122186\\
1.088	-0.122188\\
1.09	-0.122192\\
1.092	-0.1222\\
1.094	-0.122211\\
1.096	-0.122226\\
1.098	-0.122244\\
1.1	-0.122265\\
1.102	-0.12229\\
1.104	-0.122317\\
1.106	-0.122348\\
1.108	-0.122382\\
1.11	-0.122418\\
1.112	-0.122458\\
1.114	-0.1225\\
1.116	-0.122544\\
1.118	-0.122592\\
1.12	-0.122642\\
1.122	-0.122694\\
1.124	-0.122749\\
1.126	-0.122805\\
1.128	-0.122864\\
1.13	-0.122925\\
1.132	-0.122988\\
1.134	-0.123053\\
1.136	-0.123119\\
1.138	-0.123188\\
1.14	-0.123257\\
1.142	-0.123328\\
1.144	-0.1234\\
1.146	-0.123474\\
1.148	-0.123548\\
1.15	-0.123623\\
1.152	-0.1237\\
1.154	-0.123776\\
1.156	-0.123854\\
1.158	-0.123932\\
1.16	-0.12401\\
1.162	-0.124089\\
1.164	-0.124167\\
1.166	-0.124246\\
1.168	-0.124325\\
1.17	-0.124403\\
1.172	-0.124481\\
1.174	-0.124559\\
1.176	-0.124636\\
1.178	-0.124713\\
1.18	-0.124789\\
1.182	-0.124864\\
1.184	-0.124938\\
1.186	-0.125011\\
1.188	-0.125083\\
1.19	-0.125154\\
1.192	-0.125224\\
1.194	-0.125292\\
1.196	-0.125359\\
1.198	-0.125424\\
1.2	-0.125488\\
1.202	-0.12555\\
1.204	-0.12561\\
1.206	-0.125668\\
1.208	-0.125725\\
1.21	-0.125779\\
1.212	-0.125832\\
1.214	-0.125883\\
1.216	-0.125931\\
1.218	-0.125978\\
1.22	-0.126022\\
1.222	-0.126064\\
1.224	-0.126104\\
1.226	-0.126141\\
1.228	-0.126176\\
1.23	-0.126209\\
1.232	-0.12624\\
1.234	-0.126268\\
1.236	-0.126294\\
1.238	-0.126317\\
1.24	-0.126338\\
1.242	-0.126357\\
1.244	-0.126374\\
1.246	-0.126388\\
1.248	-0.1264\\
1.25	-0.126409\\
1.252	-0.126416\\
1.254	-0.126421\\
1.256	-0.126424\\
1.258	-0.126424\\
1.26	-0.126422\\
1.262	-0.126418\\
1.264	-0.126412\\
1.266	-0.126404\\
1.268	-0.126393\\
1.27	-0.126381\\
1.272	-0.126367\\
1.274	-0.126351\\
1.276	-0.126333\\
1.278	-0.126313\\
1.28	-0.126291\\
1.282	-0.126268\\
1.284	-0.126243\\
1.286	-0.126217\\
1.288	-0.126189\\
1.29	-0.126159\\
1.292	-0.126128\\
1.294	-0.126096\\
1.296	-0.126063\\
1.298	-0.126028\\
1.3	-0.125993\\
1.302	-0.125956\\
1.304	-0.125918\\
1.306	-0.12588\\
1.308	-0.12584\\
1.31	-0.1258\\
1.312	-0.125759\\
1.314	-0.125717\\
1.316	-0.125675\\
1.318	-0.125633\\
1.32	-0.12559\\
1.322	-0.125546\\
1.324	-0.125503\\
1.326	-0.125459\\
1.328	-0.125415\\
1.33	-0.12537\\
1.332	-0.125326\\
1.334	-0.125282\\
1.336	-0.125238\\
1.338	-0.125194\\
1.34	-0.12515\\
1.342	-0.125107\\
1.344	-0.125063\\
1.346	-0.125021\\
1.348	-0.124978\\
1.35	-0.124936\\
1.352	-0.124895\\
1.354	-0.124854\\
1.356	-0.124814\\
1.358	-0.124775\\
1.36	-0.124736\\
1.362	-0.124698\\
1.364	-0.12466\\
1.366	-0.124624\\
1.368	-0.124588\\
1.37	-0.124554\\
1.372	-0.12452\\
1.374	-0.124487\\
1.376	-0.124455\\
1.378	-0.124424\\
1.38	-0.124394\\
1.382	-0.124366\\
1.384	-0.124338\\
1.386	-0.124311\\
1.388	-0.124286\\
1.39	-0.124261\\
1.392	-0.124238\\
1.394	-0.124215\\
1.396	-0.124194\\
1.398	-0.124174\\
1.4	-0.124156\\
1.402	-0.124138\\
1.404	-0.124121\\
1.406	-0.124106\\
1.408	-0.124092\\
1.41	-0.124079\\
1.412	-0.124067\\
1.414	-0.124056\\
1.416	-0.124046\\
1.418	-0.124038\\
1.42	-0.12403\\
1.422	-0.124024\\
1.424	-0.124019\\
1.426	-0.124015\\
1.428	-0.124011\\
1.43	-0.124009\\
1.432	-0.124008\\
1.434	-0.124008\\
1.436	-0.124009\\
1.438	-0.124011\\
1.44	-0.124014\\
1.442	-0.124017\\
1.444	-0.124022\\
1.446	-0.124027\\
1.448	-0.124034\\
1.45	-0.124041\\
1.452	-0.124049\\
1.454	-0.124057\\
1.456	-0.124067\\
1.458	-0.124077\\
1.46	-0.124088\\
1.462	-0.1241\\
1.464	-0.124112\\
1.466	-0.124125\\
1.468	-0.124138\\
1.47	-0.124152\\
1.472	-0.124167\\
1.474	-0.124182\\
1.476	-0.124197\\
1.478	-0.124213\\
1.48	-0.124229\\
1.482	-0.124246\\
1.484	-0.124263\\
1.486	-0.124281\\
1.488	-0.124298\\
1.49	-0.124316\\
1.492	-0.124335\\
1.494	-0.124353\\
1.496	-0.124372\\
1.498	-0.124391\\
1.5	-0.12441\\
1.502	-0.124429\\
1.504	-0.124448\\
1.506	-0.124467\\
1.508	-0.124487\\
1.51	-0.124506\\
1.512	-0.124525\\
1.514	-0.124545\\
1.516	-0.124564\\
1.518	-0.124583\\
1.52	-0.124602\\
1.522	-0.124621\\
1.524	-0.12464\\
1.526	-0.124658\\
1.528	-0.124677\\
1.53	-0.124695\\
1.532	-0.124713\\
1.534	-0.12473\\
1.536	-0.124747\\
1.538	-0.124764\\
1.54	-0.124781\\
1.542	-0.124798\\
1.544	-0.124814\\
1.546	-0.124829\\
1.548	-0.124844\\
1.55	-0.124859\\
1.552	-0.124874\\
1.554	-0.124888\\
1.556	-0.124901\\
1.558	-0.124914\\
1.56	-0.124927\\
1.562	-0.124939\\
1.564	-0.12495\\
1.566	-0.124961\\
1.568	-0.124972\\
1.57	-0.124982\\
1.572	-0.124991\\
1.574	-0.125\\
1.576	-0.125008\\
1.578	-0.125016\\
1.58	-0.125023\\
1.582	-0.12503\\
1.584	-0.125036\\
1.586	-0.125041\\
1.588	-0.125046\\
1.59	-0.12505\\
1.592	-0.125054\\
1.594	-0.125057\\
1.596	-0.125059\\
1.598	-0.125061\\
1.6	-0.125062\\
1.602	-0.125062\\
1.604	-0.125062\\
1.606	-0.125062\\
1.608	-0.12506\\
1.61	-0.125058\\
1.612	-0.125056\\
1.614	-0.125053\\
1.616	-0.125049\\
1.618	-0.125045\\
1.62	-0.12504\\
1.622	-0.125035\\
1.624	-0.125029\\
1.626	-0.125023\\
1.628	-0.125016\\
1.63	-0.125008\\
1.632	-0.125\\
1.634	-0.124992\\
1.636	-0.124983\\
1.638	-0.124973\\
1.64	-0.124963\\
1.642	-0.124953\\
1.644	-0.124942\\
1.646	-0.12493\\
1.648	-0.124919\\
1.65	-0.124907\\
1.652	-0.124894\\
1.654	-0.124881\\
1.656	-0.124868\\
1.658	-0.124855\\
1.66	-0.124841\\
1.662	-0.124827\\
1.664	-0.124812\\
1.666	-0.124798\\
1.668	-0.124783\\
1.67	-0.124768\\
1.672	-0.124753\\
1.674	-0.124737\\
1.676	-0.124721\\
1.678	-0.124706\\
1.68	-0.12469\\
1.682	-0.124674\\
1.684	-0.124658\\
1.686	-0.124642\\
1.688	-0.124625\\
1.69	-0.124609\\
1.692	-0.124593\\
1.694	-0.124577\\
1.696	-0.124561\\
1.698	-0.124544\\
1.7	-0.124528\\
1.702	-0.124512\\
1.704	-0.124497\\
1.706	-0.124481\\
1.708	-0.124465\\
1.71	-0.12445\\
1.712	-0.124435\\
1.714	-0.12442\\
1.716	-0.124405\\
1.718	-0.12439\\
1.72	-0.124376\\
1.722	-0.124362\\
1.724	-0.124348\\
1.726	-0.124334\\
1.728	-0.124321\\
1.73	-0.124308\\
1.732	-0.124296\\
1.734	-0.124284\\
1.736	-0.124272\\
1.738	-0.12426\\
1.74	-0.124249\\
1.742	-0.124238\\
1.744	-0.124228\\
1.746	-0.124218\\
1.748	-0.124208\\
1.75	-0.124199\\
1.752	-0.12419\\
1.754	-0.124182\\
1.756	-0.124174\\
1.758	-0.124167\\
1.76	-0.124159\\
1.762	-0.124153\\
1.764	-0.124147\\
1.766	-0.124141\\
1.768	-0.124136\\
1.77	-0.124131\\
1.772	-0.124126\\
1.774	-0.124123\\
1.776	-0.124119\\
1.778	-0.124116\\
1.78	-0.124113\\
1.782	-0.124111\\
1.784	-0.124109\\
1.786	-0.124108\\
1.788	-0.124107\\
1.79	-0.124106\\
1.792	-0.124106\\
1.794	-0.124106\\
1.796	-0.124107\\
1.798	-0.124108\\
1.8	-0.124109\\
1.802	-0.124111\\
1.804	-0.124113\\
1.806	-0.124115\\
1.808	-0.124118\\
1.81	-0.124121\\
1.812	-0.124124\\
1.814	-0.124128\\
1.816	-0.124131\\
1.818	-0.124136\\
1.82	-0.12414\\
1.822	-0.124144\\
1.824	-0.124149\\
1.826	-0.124154\\
1.828	-0.12416\\
1.83	-0.124165\\
1.832	-0.124171\\
1.834	-0.124176\\
1.836	-0.124182\\
1.838	-0.124188\\
1.84	-0.124194\\
1.842	-0.124201\\
1.844	-0.124207\\
1.846	-0.124213\\
1.848	-0.12422\\
1.85	-0.124226\\
1.852	-0.124233\\
1.854	-0.124239\\
1.856	-0.124246\\
1.858	-0.124252\\
1.86	-0.124259\\
1.862	-0.124265\\
1.864	-0.124272\\
1.866	-0.124278\\
1.868	-0.124285\\
1.87	-0.124291\\
1.872	-0.124297\\
1.874	-0.124303\\
1.876	-0.124309\\
1.878	-0.124315\\
1.88	-0.124321\\
1.882	-0.124326\\
1.884	-0.124332\\
1.886	-0.124337\\
1.888	-0.124342\\
1.89	-0.124347\\
1.892	-0.124352\\
1.894	-0.124357\\
1.896	-0.124361\\
1.898	-0.124365\\
1.9	-0.124369\\
1.902	-0.124373\\
1.904	-0.124377\\
1.906	-0.12438\\
1.908	-0.124383\\
1.91	-0.124386\\
1.912	-0.124389\\
1.914	-0.124391\\
1.916	-0.124393\\
1.918	-0.124395\\
1.92	-0.124397\\
1.922	-0.124399\\
1.924	-0.1244\\
1.926	-0.124401\\
1.928	-0.124402\\
1.93	-0.124402\\
1.932	-0.124403\\
1.934	-0.124403\\
1.936	-0.124403\\
1.938	-0.124402\\
1.94	-0.124402\\
1.942	-0.124401\\
1.944	-0.1244\\
1.946	-0.124399\\
1.948	-0.124397\\
1.95	-0.124396\\
1.952	-0.124394\\
1.954	-0.124391\\
1.956	-0.124389\\
1.958	-0.124387\\
1.96	-0.124384\\
1.962	-0.124381\\
1.964	-0.124378\\
1.966	-0.124375\\
1.968	-0.124371\\
1.97	-0.124368\\
1.972	-0.124364\\
1.974	-0.12436\\
1.976	-0.124356\\
1.978	-0.124352\\
1.98	-0.124347\\
1.982	-0.124343\\
1.984	-0.124338\\
1.986	-0.124334\\
1.988	-0.124329\\
1.99	-0.124324\\
1.992	-0.124319\\
1.994	-0.124314\\
1.996	-0.124308\\
1.998	-0.124303\\
2	-0.124298\\
};
\addplot [color=black, line width=1.4pt, forget plot]
  table[row sep=crcr]{%
-0.124	0\\
-0.122	0\\
-0.12	0\\
-0.118	0\\
-0.116	0\\
-0.114	0\\
-0.112	0\\
-0.11	0\\
-0.108	0\\
-0.106	0\\
-0.104	0\\
-0.102	0\\
-0.1	0\\
-0.098	0\\
-0.096	0\\
-0.094	0\\
-0.092	0\\
-0.09	0\\
-0.088	0\\
-0.086	0\\
-0.084	0\\
-0.082	0\\
-0.08	0\\
-0.078	0\\
-0.076	0\\
-0.074	0\\
-0.072	0\\
-0.07	0\\
-0.068	0\\
-0.066	0\\
-0.064	0\\
-0.062	0\\
-0.06	0\\
-0.058	0\\
-0.056	0\\
-0.054	0\\
-0.052	0\\
-0.05	0\\
-0.048	0\\
-0.046	0\\
-0.044	0\\
-0.042	0\\
-0.04	0\\
-0.038	0\\
-0.036	0\\
-0.034	0\\
-0.032	0\\
-0.03	0\\
-0.028	0\\
-0.026	0\\
-0.024	0\\
-0.022	0\\
-0.02	0\\
-0.018	0\\
-0.016	0\\
-0.014	0\\
-0.012	0\\
-0.01	0\\
-0.008	0\\
-0.006	0\\
-0.004	0\\
-0.002	0\\
0	0\\
0.002	-0.00190406\\
0.004	-0.00374712\\
0.006	-0.00555085\\
0.008	-0.00732857\\
0.01	-0.00908785\\
0.012	-0.0108325\\
0.014	-0.0125641\\
0.016	-0.0142825\\
0.018	-0.0159871\\
0.02	-0.0176766\\
0.022	-0.01935\\
0.024	-0.0210061\\
0.026	-0.0226439\\
0.028	-0.0242627\\
0.03	-0.0258617\\
0.032	-0.0274407\\
0.034	-0.0289991\\
0.036	-0.030537\\
0.038	-0.0320541\\
0.04	-0.0335505\\
0.042	-0.0350263\\
0.044	-0.0364814\\
0.046	-0.037916\\
0.048	-0.0393302\\
0.05	-0.0407242\\
0.052	-0.0420982\\
0.054	-0.0434524\\
0.056	-0.0447868\\
0.058	-0.0461017\\
0.06	-0.0473973\\
0.062	-0.0486737\\
0.064	-0.0499312\\
0.066	-0.05117\\
0.068	-0.0523902\\
0.07	-0.0535921\\
0.072	-0.0547759\\
0.074	-0.0559416\\
0.076	-0.0570897\\
0.078	-0.0582202\\
0.08	-0.0593333\\
0.082	-0.0604293\\
0.084	-0.0615083\\
0.086	-0.0625706\\
0.088	-0.0636162\\
0.09	-0.0646455\\
0.092	-0.0656586\\
0.094	-0.0666557\\
0.096	-0.0676369\\
0.098	-0.0686025\\
0.1	-0.0695526\\
0.102	-0.0704875\\
0.104	-0.0714073\\
0.106	-0.0723121\\
0.108	-0.0732021\\
0.11	-0.0740776\\
0.112	-0.0749387\\
0.114	-0.0757856\\
0.116	-0.0766184\\
0.118	-0.0774372\\
0.12	-0.0782424\\
0.122	-0.079034\\
0.124	-0.0798122\\
0.126	-0.0805772\\
0.128	-0.0813291\\
0.13	-0.0820681\\
0.132	-0.0827943\\
0.134	-0.083508\\
0.136	-0.0842093\\
0.138	-0.0848983\\
0.14	-0.0855752\\
0.142	-0.0862402\\
0.144	-0.0868935\\
0.146	-0.0875351\\
0.148	-0.0881653\\
0.15	-0.0887842\\
0.152	-0.089392\\
0.154	-0.0899888\\
0.156	-0.0905749\\
0.158	-0.0911503\\
0.16	-0.0917152\\
0.162	-0.0922698\\
0.164	-0.0928143\\
0.166	-0.0933488\\
0.168	-0.0938734\\
0.17	-0.0943884\\
0.172	-0.0948939\\
0.174	-0.0953901\\
0.176	-0.095877\\
0.178	-0.096355\\
0.18	-0.0968241\\
0.182	-0.0972845\\
0.184	-0.0977363\\
0.186	-0.0981798\\
0.188	-0.0986151\\
0.19	-0.0990423\\
0.192	-0.0994616\\
0.194	-0.0998732\\
0.196	-0.100277\\
0.198	-0.100674\\
0.2	-0.101063\\
0.202	-0.101445\\
0.204	-0.101821\\
0.206	-0.102189\\
0.208	-0.102551\\
0.21	-0.102907\\
0.212	-0.103256\\
0.214	-0.103599\\
0.216	-0.103936\\
0.218	-0.104267\\
0.22	-0.104592\\
0.222	-0.104912\\
0.224	-0.105226\\
0.226	-0.105535\\
0.228	-0.105839\\
0.23	-0.106138\\
0.232	-0.106432\\
0.234	-0.106721\\
0.236	-0.107006\\
0.238	-0.107286\\
0.24	-0.107562\\
0.242	-0.107833\\
0.244	-0.108101\\
0.246	-0.108364\\
0.248	-0.108624\\
0.25	-0.10888\\
0.252	-0.109133\\
0.254	-0.109382\\
0.256	-0.109628\\
0.258	-0.10987\\
0.26	-0.110109\\
0.262	-0.110345\\
0.264	-0.110579\\
0.266	-0.110809\\
0.268	-0.111037\\
0.27	-0.111262\\
0.272	-0.111484\\
0.274	-0.111704\\
0.276	-0.111922\\
0.278	-0.112137\\
0.28	-0.11235\\
0.282	-0.112561\\
0.284	-0.11277\\
0.286	-0.112976\\
0.288	-0.113181\\
0.29	-0.113384\\
0.292	-0.113585\\
0.294	-0.113784\\
0.296	-0.113982\\
0.298	-0.114177\\
0.3	-0.114372\\
0.302	-0.114564\\
0.304	-0.114755\\
0.306	-0.114945\\
0.308	-0.115133\\
0.31	-0.115319\\
0.312	-0.115505\\
0.314	-0.115688\\
0.316	-0.115871\\
0.318	-0.116052\\
0.32	-0.116232\\
0.322	-0.11641\\
0.324	-0.116588\\
0.326	-0.116763\\
0.328	-0.116938\\
0.33	-0.117111\\
0.332	-0.117284\\
0.334	-0.117454\\
0.336	-0.117624\\
0.338	-0.117792\\
0.34	-0.11796\\
0.342	-0.118125\\
0.344	-0.11829\\
0.346	-0.118453\\
0.348	-0.118615\\
0.35	-0.118776\\
0.352	-0.118935\\
0.354	-0.119093\\
0.356	-0.11925\\
0.358	-0.119405\\
0.36	-0.119559\\
0.362	-0.119711\\
0.364	-0.119862\\
0.366	-0.120011\\
0.368	-0.120159\\
0.37	-0.120306\\
0.372	-0.120451\\
0.374	-0.120594\\
0.376	-0.120736\\
0.378	-0.120876\\
0.38	-0.121015\\
0.382	-0.121152\\
0.384	-0.121287\\
0.386	-0.121421\\
0.388	-0.121552\\
0.39	-0.121683\\
0.392	-0.121811\\
0.394	-0.121937\\
0.396	-0.122062\\
0.398	-0.122184\\
0.4	-0.122305\\
0.402	-0.122424\\
0.404	-0.122541\\
0.406	-0.122656\\
0.408	-0.122769\\
0.41	-0.12288\\
0.412	-0.122989\\
0.414	-0.123096\\
0.416	-0.123201\\
0.418	-0.123304\\
0.42	-0.123405\\
0.422	-0.123504\\
0.424	-0.1236\\
0.426	-0.123695\\
0.428	-0.123787\\
0.43	-0.123877\\
0.432	-0.123965\\
0.434	-0.124051\\
0.436	-0.124135\\
0.438	-0.124216\\
0.44	-0.124295\\
0.442	-0.124373\\
0.444	-0.124448\\
0.446	-0.12452\\
0.448	-0.124591\\
0.45	-0.12466\\
0.452	-0.124726\\
0.454	-0.12479\\
0.456	-0.124852\\
0.458	-0.124912\\
0.46	-0.12497\\
0.462	-0.125026\\
0.464	-0.125079\\
0.466	-0.125131\\
0.468	-0.12518\\
0.47	-0.125228\\
0.472	-0.125273\\
0.474	-0.125316\\
0.476	-0.125358\\
0.478	-0.125397\\
0.48	-0.125435\\
0.482	-0.12547\\
0.484	-0.125504\\
0.486	-0.125536\\
0.488	-0.125566\\
0.49	-0.125594\\
0.492	-0.12562\\
0.494	-0.125645\\
0.496	-0.125668\\
0.498	-0.125689\\
0.5	-0.125709\\
0.502	-0.125727\\
0.504	-0.125744\\
0.506	-0.125758\\
0.508	-0.125772\\
0.51	-0.125784\\
0.512	-0.125794\\
0.514	-0.125804\\
0.516	-0.125811\\
0.518	-0.125818\\
0.52	-0.125823\\
0.522	-0.125827\\
0.524	-0.12583\\
0.526	-0.125832\\
0.528	-0.125832\\
0.53	-0.125832\\
0.532	-0.12583\\
0.534	-0.125828\\
0.536	-0.125824\\
0.538	-0.12582\\
0.54	-0.125814\\
0.542	-0.125808\\
0.544	-0.125801\\
0.546	-0.125793\\
0.548	-0.125785\\
0.55	-0.125776\\
0.552	-0.125766\\
0.554	-0.125756\\
0.556	-0.125745\\
0.558	-0.125733\\
0.56	-0.125721\\
0.562	-0.125709\\
0.564	-0.125696\\
0.566	-0.125683\\
0.568	-0.125669\\
0.57	-0.125655\\
0.572	-0.125641\\
0.574	-0.125627\\
0.576	-0.125612\\
0.578	-0.125597\\
0.58	-0.125582\\
0.582	-0.125567\\
0.584	-0.125552\\
0.586	-0.125536\\
0.588	-0.125521\\
0.59	-0.125506\\
0.592	-0.12549\\
0.594	-0.125475\\
0.596	-0.12546\\
0.598	-0.125445\\
0.6	-0.12543\\
0.602	-0.125415\\
0.604	-0.1254\\
0.606	-0.125385\\
0.608	-0.125371\\
0.61	-0.125357\\
0.612	-0.125343\\
0.614	-0.125329\\
0.616	-0.125316\\
0.618	-0.125303\\
0.62	-0.12529\\
0.622	-0.125278\\
0.624	-0.125266\\
0.626	-0.125254\\
0.628	-0.125242\\
0.63	-0.125231\\
0.632	-0.125221\\
0.634	-0.12521\\
0.636	-0.1252\\
0.638	-0.125191\\
0.64	-0.125182\\
0.642	-0.125173\\
0.644	-0.125165\\
0.646	-0.125157\\
0.648	-0.12515\\
0.65	-0.125143\\
0.652	-0.125136\\
0.654	-0.12513\\
0.656	-0.125124\\
0.658	-0.125119\\
0.66	-0.125115\\
0.662	-0.12511\\
0.664	-0.125106\\
0.666	-0.125103\\
0.668	-0.1251\\
0.67	-0.125098\\
0.672	-0.125095\\
0.674	-0.125094\\
0.676	-0.125092\\
0.678	-0.125092\\
0.68	-0.125091\\
0.682	-0.125091\\
0.684	-0.125092\\
0.686	-0.125092\\
0.688	-0.125093\\
0.69	-0.125095\\
0.692	-0.125097\\
0.694	-0.125099\\
0.696	-0.125102\\
0.698	-0.125105\\
0.7	-0.125108\\
0.702	-0.125112\\
0.704	-0.125115\\
0.706	-0.12512\\
0.708	-0.125124\\
0.71	-0.125129\\
0.712	-0.125134\\
0.714	-0.125139\\
0.716	-0.125144\\
0.718	-0.12515\\
0.72	-0.125156\\
0.722	-0.125162\\
0.724	-0.125168\\
0.726	-0.125175\\
0.728	-0.125181\\
0.73	-0.125188\\
0.732	-0.125195\\
0.734	-0.125202\\
0.736	-0.125209\\
0.738	-0.125216\\
0.74	-0.125224\\
0.742	-0.125231\\
0.744	-0.125238\\
0.746	-0.125246\\
0.748	-0.125253\\
0.75	-0.12526\\
0.752	-0.125268\\
0.754	-0.125275\\
0.756	-0.125282\\
0.758	-0.12529\\
0.76	-0.125297\\
0.762	-0.125304\\
0.764	-0.125311\\
0.766	-0.125318\\
0.768	-0.125325\\
0.77	-0.125331\\
0.772	-0.125338\\
0.774	-0.125344\\
0.776	-0.125351\\
0.778	-0.125357\\
0.78	-0.125363\\
0.782	-0.125368\\
0.784	-0.125374\\
0.786	-0.125379\\
0.788	-0.125384\\
0.79	-0.125389\\
0.792	-0.125393\\
0.794	-0.125398\\
0.796	-0.125402\\
0.798	-0.125406\\
0.8	-0.125409\\
0.802	-0.125413\\
0.804	-0.125416\\
0.806	-0.125418\\
0.808	-0.125421\\
0.81	-0.125423\\
0.812	-0.125425\\
0.814	-0.125426\\
0.816	-0.125428\\
0.818	-0.125429\\
0.82	-0.125429\\
0.822	-0.125429\\
0.824	-0.12543\\
0.826	-0.125429\\
0.828	-0.125429\\
0.83	-0.125428\\
0.832	-0.125426\\
0.834	-0.125425\\
0.836	-0.125423\\
0.838	-0.125421\\
0.84	-0.125418\\
0.842	-0.125416\\
0.844	-0.125413\\
0.846	-0.125409\\
0.848	-0.125406\\
0.85	-0.125402\\
0.852	-0.125398\\
0.854	-0.125393\\
0.856	-0.125388\\
0.858	-0.125384\\
0.86	-0.125378\\
0.862	-0.125373\\
0.864	-0.125367\\
0.866	-0.125361\\
0.868	-0.125355\\
0.87	-0.125349\\
0.872	-0.125342\\
0.874	-0.125335\\
0.876	-0.125328\\
0.878	-0.125321\\
0.88	-0.125314\\
0.882	-0.125307\\
0.884	-0.125299\\
0.886	-0.125291\\
0.888	-0.125284\\
0.89	-0.125276\\
0.892	-0.125268\\
0.894	-0.125259\\
0.896	-0.125251\\
0.898	-0.125243\\
0.9	-0.125235\\
0.902	-0.125226\\
0.904	-0.125218\\
0.906	-0.125209\\
0.908	-0.125201\\
0.91	-0.125192\\
0.912	-0.125184\\
0.914	-0.125175\\
0.916	-0.125167\\
0.918	-0.125158\\
0.92	-0.12515\\
0.922	-0.125142\\
0.924	-0.125134\\
0.926	-0.125125\\
0.928	-0.125117\\
0.93	-0.125109\\
0.932	-0.125102\\
0.934	-0.125094\\
0.936	-0.125086\\
0.938	-0.125079\\
0.94	-0.125071\\
0.942	-0.125064\\
0.944	-0.125057\\
0.946	-0.12505\\
0.948	-0.125044\\
0.95	-0.125037\\
0.952	-0.125031\\
0.954	-0.125025\\
0.956	-0.125019\\
0.958	-0.125013\\
0.96	-0.125007\\
0.962	-0.125002\\
0.964	-0.124997\\
0.966	-0.124992\\
0.968	-0.124987\\
0.97	-0.124983\\
0.972	-0.124979\\
0.974	-0.124975\\
0.976	-0.124971\\
0.978	-0.124967\\
0.98	-0.124964\\
0.982	-0.124961\\
0.984	-0.124958\\
0.986	-0.124956\\
0.988	-0.124953\\
0.99	-0.124951\\
0.992	-0.124949\\
0.994	-0.124948\\
0.996	-0.124946\\
0.998	-0.124945\\
1	-0.124944\\
1.002	-0.124943\\
1.004	-0.124943\\
1.006	-0.124943\\
1.008	-0.124943\\
1.01	-0.124943\\
1.012	-0.124943\\
1.014	-0.124944\\
1.016	-0.124945\\
1.018	-0.124946\\
1.02	-0.124947\\
1.022	-0.124948\\
1.024	-0.12495\\
1.026	-0.124951\\
1.028	-0.124953\\
1.03	-0.124955\\
1.032	-0.124958\\
1.034	-0.12496\\
1.036	-0.124962\\
1.038	-0.124965\\
1.04	-0.124968\\
1.042	-0.124971\\
1.044	-0.124974\\
1.046	-0.124977\\
1.048	-0.12498\\
1.05	-0.124983\\
1.052	-0.124987\\
1.054	-0.12499\\
1.056	-0.124994\\
1.058	-0.124997\\
1.06	-0.125001\\
1.062	-0.125005\\
1.064	-0.125009\\
1.066	-0.125012\\
1.068	-0.125016\\
1.07	-0.12502\\
1.072	-0.125024\\
1.074	-0.125028\\
1.076	-0.125032\\
1.078	-0.125036\\
1.08	-0.12504\\
1.082	-0.125044\\
1.084	-0.125048\\
1.086	-0.125051\\
1.088	-0.125055\\
1.09	-0.125059\\
1.092	-0.125063\\
1.094	-0.125066\\
1.096	-0.12507\\
1.098	-0.125074\\
1.1	-0.125077\\
1.102	-0.125081\\
1.104	-0.125084\\
1.106	-0.125087\\
1.108	-0.12509\\
1.11	-0.125093\\
1.112	-0.125096\\
1.114	-0.125099\\
1.116	-0.125102\\
1.118	-0.125105\\
1.12	-0.125107\\
1.122	-0.12511\\
1.124	-0.125112\\
1.126	-0.125114\\
1.128	-0.125116\\
1.13	-0.125118\\
1.132	-0.12512\\
1.134	-0.125122\\
1.136	-0.125124\\
1.138	-0.125125\\
1.14	-0.125126\\
1.142	-0.125127\\
1.144	-0.125129\\
1.146	-0.125129\\
1.148	-0.12513\\
1.15	-0.125131\\
1.152	-0.125131\\
1.154	-0.125132\\
1.156	-0.125132\\
1.158	-0.125132\\
1.16	-0.125132\\
1.162	-0.125132\\
1.164	-0.125131\\
1.166	-0.125131\\
1.168	-0.12513\\
1.17	-0.125129\\
1.172	-0.125128\\
1.174	-0.125127\\
1.176	-0.125126\\
1.178	-0.125125\\
1.18	-0.125123\\
1.182	-0.125122\\
1.184	-0.12512\\
1.186	-0.125118\\
1.188	-0.125116\\
1.19	-0.125114\\
1.192	-0.125112\\
1.194	-0.12511\\
1.196	-0.125107\\
1.198	-0.125105\\
1.2	-0.125102\\
1.202	-0.125099\\
1.204	-0.125097\\
1.206	-0.125094\\
1.208	-0.125091\\
1.21	-0.125087\\
1.212	-0.125084\\
1.214	-0.125081\\
1.216	-0.125077\\
1.218	-0.125074\\
1.22	-0.12507\\
1.222	-0.125066\\
1.224	-0.125063\\
1.226	-0.125059\\
1.228	-0.125055\\
1.23	-0.125051\\
1.232	-0.125047\\
1.234	-0.125043\\
1.236	-0.125039\\
1.238	-0.125034\\
1.24	-0.12503\\
1.242	-0.125026\\
1.244	-0.125021\\
1.246	-0.125017\\
1.248	-0.125012\\
1.25	-0.125008\\
1.252	-0.125003\\
1.254	-0.124999\\
1.256	-0.124994\\
1.258	-0.12499\\
1.26	-0.124985\\
1.262	-0.12498\\
1.264	-0.124976\\
1.266	-0.124971\\
1.268	-0.124966\\
1.27	-0.124962\\
1.272	-0.124957\\
1.274	-0.124952\\
1.276	-0.124948\\
1.278	-0.124943\\
1.28	-0.124938\\
1.282	-0.124934\\
1.284	-0.124929\\
1.286	-0.124925\\
1.288	-0.12492\\
1.29	-0.124915\\
1.292	-0.124911\\
1.294	-0.124906\\
1.296	-0.124902\\
1.298	-0.124897\\
1.3	-0.124893\\
1.302	-0.124889\\
1.304	-0.124884\\
1.306	-0.12488\\
1.308	-0.124876\\
1.31	-0.124872\\
1.312	-0.124867\\
1.314	-0.124863\\
1.316	-0.124859\\
1.318	-0.124855\\
1.32	-0.124851\\
1.322	-0.124848\\
1.324	-0.124844\\
1.326	-0.12484\\
1.328	-0.124836\\
1.33	-0.124833\\
1.332	-0.124829\\
1.334	-0.124826\\
1.336	-0.124822\\
1.338	-0.124819\\
1.34	-0.124815\\
1.342	-0.124812\\
1.344	-0.124809\\
1.346	-0.124806\\
1.348	-0.124803\\
1.35	-0.1248\\
1.352	-0.124797\\
1.354	-0.124794\\
1.356	-0.124792\\
1.358	-0.124789\\
1.36	-0.124786\\
1.362	-0.124784\\
1.364	-0.124781\\
1.366	-0.124779\\
1.368	-0.124776\\
1.37	-0.124774\\
1.372	-0.124772\\
1.374	-0.12477\\
1.376	-0.124768\\
1.378	-0.124766\\
1.38	-0.124764\\
1.382	-0.124762\\
1.384	-0.12476\\
1.386	-0.124758\\
1.388	-0.124757\\
1.39	-0.124755\\
1.392	-0.124753\\
1.394	-0.124752\\
1.396	-0.12475\\
1.398	-0.124749\\
1.4	-0.124748\\
1.402	-0.124746\\
1.404	-0.124745\\
1.406	-0.124744\\
1.408	-0.124742\\
1.41	-0.124741\\
1.412	-0.12474\\
1.414	-0.124739\\
1.416	-0.124738\\
1.418	-0.124737\\
1.42	-0.124736\\
1.422	-0.124735\\
1.424	-0.124734\\
1.426	-0.124733\\
1.428	-0.124732\\
1.43	-0.124731\\
1.432	-0.124731\\
1.434	-0.12473\\
1.436	-0.124729\\
1.438	-0.124728\\
1.44	-0.124727\\
1.442	-0.124726\\
1.444	-0.124726\\
1.446	-0.124725\\
1.448	-0.124724\\
1.45	-0.124723\\
1.452	-0.124722\\
1.454	-0.124721\\
1.456	-0.124721\\
1.458	-0.12472\\
1.46	-0.124719\\
1.462	-0.124718\\
1.464	-0.124717\\
1.466	-0.124716\\
1.468	-0.124715\\
1.47	-0.124714\\
1.472	-0.124713\\
1.474	-0.124712\\
1.476	-0.124711\\
1.478	-0.12471\\
1.48	-0.124709\\
1.482	-0.124708\\
1.484	-0.124707\\
1.486	-0.124706\\
1.488	-0.124705\\
1.49	-0.124703\\
1.492	-0.124702\\
1.494	-0.124701\\
1.496	-0.124699\\
1.498	-0.124698\\
1.5	-0.124696\\
1.502	-0.124695\\
1.504	-0.124693\\
1.506	-0.124692\\
1.508	-0.12469\\
1.51	-0.124689\\
1.512	-0.124687\\
1.514	-0.124685\\
1.516	-0.124683\\
1.518	-0.124682\\
1.52	-0.12468\\
1.522	-0.124678\\
1.524	-0.124676\\
1.526	-0.124674\\
1.528	-0.124672\\
1.53	-0.12467\\
1.532	-0.124668\\
1.534	-0.124665\\
1.536	-0.124663\\
1.538	-0.124661\\
1.54	-0.124659\\
1.542	-0.124656\\
1.544	-0.124654\\
1.546	-0.124652\\
1.548	-0.124649\\
1.55	-0.124647\\
1.552	-0.124644\\
1.554	-0.124642\\
1.556	-0.124639\\
1.558	-0.124637\\
1.56	-0.124634\\
1.562	-0.124631\\
1.564	-0.124629\\
1.566	-0.124626\\
1.568	-0.124623\\
1.57	-0.12462\\
1.572	-0.124618\\
1.574	-0.124615\\
1.576	-0.124612\\
1.578	-0.124609\\
1.58	-0.124606\\
1.582	-0.124604\\
1.584	-0.124601\\
1.586	-0.124598\\
1.588	-0.124595\\
1.59	-0.124592\\
1.592	-0.124589\\
1.594	-0.124586\\
1.596	-0.124583\\
1.598	-0.124581\\
1.6	-0.124578\\
1.602	-0.124575\\
1.604	-0.124572\\
1.606	-0.124569\\
1.608	-0.124566\\
1.61	-0.124563\\
1.612	-0.12456\\
1.614	-0.124557\\
1.616	-0.124555\\
1.618	-0.124552\\
1.62	-0.124549\\
1.622	-0.124546\\
1.624	-0.124543\\
1.626	-0.12454\\
1.628	-0.124538\\
1.63	-0.124535\\
1.632	-0.124532\\
1.634	-0.124529\\
1.636	-0.124527\\
1.638	-0.124524\\
1.64	-0.124521\\
1.642	-0.124519\\
1.644	-0.124516\\
1.646	-0.124513\\
1.648	-0.124511\\
1.65	-0.124508\\
1.652	-0.124506\\
1.654	-0.124503\\
1.656	-0.124501\\
1.658	-0.124498\\
1.66	-0.124496\\
1.662	-0.124493\\
1.664	-0.124491\\
1.666	-0.124488\\
1.668	-0.124486\\
1.67	-0.124484\\
1.672	-0.124481\\
1.674	-0.124479\\
1.676	-0.124477\\
1.678	-0.124475\\
1.68	-0.124473\\
1.682	-0.12447\\
1.684	-0.124468\\
1.686	-0.124466\\
1.688	-0.124464\\
1.69	-0.124462\\
1.692	-0.12446\\
1.694	-0.124458\\
1.696	-0.124456\\
1.698	-0.124454\\
1.7	-0.124452\\
1.702	-0.12445\\
1.704	-0.124448\\
1.706	-0.124446\\
1.708	-0.124444\\
1.71	-0.124443\\
1.712	-0.124441\\
1.714	-0.124439\\
1.716	-0.124437\\
1.718	-0.124436\\
1.72	-0.124434\\
1.722	-0.124432\\
1.724	-0.12443\\
1.726	-0.124429\\
1.728	-0.124427\\
1.73	-0.124425\\
1.732	-0.124424\\
1.734	-0.124422\\
1.736	-0.124421\\
1.738	-0.124419\\
1.74	-0.124417\\
1.742	-0.124416\\
1.744	-0.124414\\
1.746	-0.124413\\
1.748	-0.124411\\
1.75	-0.12441\\
1.752	-0.124408\\
1.754	-0.124407\\
1.756	-0.124405\\
1.758	-0.124404\\
1.76	-0.124402\\
1.762	-0.124401\\
1.764	-0.124399\\
1.766	-0.124398\\
1.768	-0.124397\\
1.77	-0.124395\\
1.772	-0.124394\\
1.774	-0.124392\\
1.776	-0.124391\\
1.778	-0.124389\\
1.78	-0.124388\\
1.782	-0.124387\\
1.784	-0.124385\\
1.786	-0.124384\\
1.788	-0.124382\\
1.79	-0.124381\\
1.792	-0.124379\\
1.794	-0.124378\\
1.796	-0.124377\\
1.798	-0.124375\\
1.8	-0.124374\\
1.802	-0.124372\\
1.804	-0.124371\\
1.806	-0.124369\\
1.808	-0.124368\\
1.81	-0.124366\\
1.812	-0.124365\\
1.814	-0.124363\\
1.816	-0.124362\\
1.818	-0.12436\\
1.82	-0.124359\\
1.822	-0.124357\\
1.824	-0.124356\\
1.826	-0.124354\\
1.828	-0.124353\\
1.83	-0.124351\\
1.832	-0.12435\\
1.834	-0.124348\\
1.836	-0.124347\\
1.838	-0.124345\\
1.84	-0.124343\\
1.842	-0.124342\\
1.844	-0.12434\\
1.846	-0.124339\\
1.848	-0.124337\\
1.85	-0.124335\\
1.852	-0.124334\\
1.854	-0.124332\\
1.856	-0.12433\\
1.858	-0.124329\\
1.86	-0.124327\\
1.862	-0.124325\\
1.864	-0.124324\\
1.866	-0.124322\\
1.868	-0.12432\\
1.87	-0.124318\\
1.872	-0.124317\\
1.874	-0.124315\\
1.876	-0.124313\\
1.878	-0.124311\\
1.88	-0.124309\\
1.882	-0.124307\\
1.884	-0.124306\\
1.886	-0.124304\\
1.888	-0.124302\\
1.89	-0.1243\\
1.892	-0.124298\\
1.894	-0.124296\\
1.896	-0.124294\\
1.898	-0.124292\\
1.9	-0.12429\\
1.902	-0.124289\\
1.904	-0.124287\\
1.906	-0.124285\\
1.908	-0.124283\\
1.91	-0.124281\\
1.912	-0.124279\\
1.914	-0.124277\\
1.916	-0.124275\\
1.918	-0.124273\\
1.92	-0.124271\\
1.922	-0.124269\\
1.924	-0.124266\\
1.926	-0.124264\\
1.928	-0.124262\\
1.93	-0.12426\\
1.932	-0.124258\\
1.934	-0.124256\\
1.936	-0.124254\\
1.938	-0.124252\\
1.94	-0.12425\\
1.942	-0.124248\\
1.944	-0.124246\\
1.946	-0.124243\\
1.948	-0.124241\\
1.95	-0.124239\\
1.952	-0.124237\\
1.954	-0.124235\\
1.956	-0.124233\\
1.958	-0.124231\\
1.96	-0.124229\\
1.962	-0.124226\\
1.964	-0.124224\\
1.966	-0.124222\\
1.968	-0.12422\\
1.97	-0.124218\\
1.972	-0.124216\\
1.974	-0.124214\\
1.976	-0.124212\\
1.978	-0.124209\\
1.98	-0.124207\\
1.982	-0.124205\\
1.984	-0.124203\\
1.986	-0.124201\\
1.988	-0.124199\\
1.99	-0.124197\\
1.992	-0.124195\\
1.994	-0.124193\\
1.996	-0.12419\\
1.998	-0.124188\\
2	-0.124186\\
};
\addplot [color=black, dotted, line width=0.6pt, forget plot]
  table[row sep=crcr]{%
-0.125	-0.125832\\
2	-0.125832\\
};
\draw[line width=\figlegendlw, fill=\figlegendfill, draw=\figlegenddraw] (axis cs:1.13004,-0.0234737) rectangle (axis cs:1.9575,0.094);
\addplot [color=black, line width=1.4pt, forget plot]
  table[row sep=crcr]{%
1.17141	0.0705053\\
1.4362	0.0705053\\
};
\node[right, align=left, inner sep=0]
at (axis cs:1.461,0.071) {$f_{\mathrm{COI}}$};
\addplot [color=mycolor5, line width=1.6pt, forget plot]
  table[row sep=crcr]{%
1.17141	0.0352632\\
1.20083	0.0352632\\
};
\addplot [color=mycolor6, line width=1.6pt, forget plot]
  table[row sep=crcr]{%
1.20083	0.0352632\\
1.23025	0.0352632\\
};
\addplot [color=mycolor10, line width=1.6pt, forget plot]
  table[row sep=crcr]{%
1.23025	0.0352632\\
1.25967	0.0352632\\
};
\addplot [color=mycolor7, line width=1.6pt, forget plot]
  table[row sep=crcr]{%
1.25967	0.0352632\\
1.28909	0.0352632\\
};
\addplot [color=mycolor3, line width=1.6pt, forget plot]
  table[row sep=crcr]{%
1.28909	0.0352632\\
1.31852	0.0352632\\
};
\addplot [color=mycolor2, line width=1.6pt, forget plot]
  table[row sep=crcr]{%
1.31852	0.0352632\\
1.34794	0.0352632\\
};
\addplot [color=mycolor8, line width=1.6pt, forget plot]
  table[row sep=crcr]{%
1.34794	0.0352632\\
1.37736	0.0352632\\
};
\addplot [color=mycolor9, line width=1.6pt, forget plot]
  table[row sep=crcr]{%
1.37736	0.0352632\\
1.40678	0.0352632\\
};
\addplot [color=mycolor4, line width=1.6pt, forget plot]
  table[row sep=crcr]{%
1.40678	0.0352632\\
1.4362	0.0352632\\
};
\node[right, align=left, inner sep=0]
at (axis cs:1.461,0.035) {$f_i$};
\addplot [color=black, dotted, line width=0.6pt, forget plot]
  table[row sep=crcr]{%
1.17141	2.10526e-05\\
1.4362	2.10526e-05\\
};
\node[right, align=left, inner sep=0]
at (axis cs:1.461,0) {$N_{\mathrm{COI}}^{0}$};
\end{axis}

\begin{axis}[%
width={\figscale*1.32in},
height={\figscale*0.95in},
at={(\figscale*1.985in,\figscale*0.31in)},
scale only axis,
separate axis lines,
every outer x axis line/.append style={black},
every x tick label/.append style={font=\figticfont},
every x tick/.append style={black},
xmin=-0.125,
xmax=2,
xlabel={Time since step in s},
every outer y axis line/.append style={black},
every y tick label/.append style={font=\figticfont},
every y tick/.append style={black},
ymin=-0.2,
ymax=0.1,
axis background/.style={fill=white},
xmajorgrids,
ymajorgrids,
grid style={white!55!black, opacity=0.3}
]

\addplot[area legend, draw=none, fill=mycolor1, fill opacity=0.9, forget plot]
table[row sep=crcr] {%
x	y\\
0	-0.01256\\
0.00501253	-0.0403463\\
0.0100251	-0.0463599\\
0.0150376	-0.0637061\\
0.0200501	-0.0827629\\
0.0250627	-0.100364\\
0.0300752	-0.116487\\
0.0350877	-0.130925\\
0.0401003	-0.143713\\
0.0451128	-0.15469\\
0.0501253	-0.163942\\
0.0551378	-0.17136\\
0.0601504	-0.177071\\
0.0651629	-0.181024\\
0.0701754	-0.182473\\
0.075188	-0.183814\\
0.0802005	-0.185057\\
0.085213	-0.186198\\
0.0902256	-0.187246\\
0.0952381	-0.188199\\
0.100251	-0.189064\\
0.105263	-0.18984\\
0.110276	-0.190534\\
0.115288	-0.191146\\
0.120301	-0.19168\\
0.125313	-0.192137\\
0.130326	-0.192523\\
0.135338	-0.192837\\
0.140351	-0.193085\\
0.145363	-0.193268\\
0.150376	-0.193389\\
0.155388	-0.19345\\
0.160401	-0.193456\\
0.165414	-0.193406\\
0.170426	-0.193306\\
0.175439	-0.193156\\
0.180451	-0.19296\\
0.185464	-0.19272\\
0.190476	-0.192439\\
0.195489	-0.192118\\
0.200501	-0.19176\\
0.205514	-0.191368\\
0.210526	-0.190943\\
0.215539	-0.190488\\
0.220551	-0.190005\\
0.225564	-0.189495\\
0.230576	-0.188961\\
0.235589	-0.188405\\
0.240602	-0.187828\\
0.245614	-0.187231\\
0.250627	-0.186617\\
0.255639	-0.185987\\
0.260652	-0.185342\\
0.265664	-0.184668\\
0.270677	-0.183951\\
0.275689	-0.183198\\
0.280702	-0.182411\\
0.285714	-0.181595\\
0.290727	-0.180755\\
0.295739	-0.179895\\
0.300752	-0.179018\\
0.305764	-0.178129\\
0.310777	-0.177232\\
0.315789	-0.17633\\
0.320802	-0.175428\\
0.325815	-0.174528\\
0.330827	-0.173635\\
0.33584	-0.172752\\
0.340852	-0.171883\\
0.345865	-0.171031\\
0.350877	-0.170199\\
0.35589	-0.169392\\
0.360902	-0.168612\\
0.365915	-0.167864\\
0.370927	-0.167149\\
0.37594	-0.166473\\
0.380952	-0.165838\\
0.385965	-0.165247\\
0.390977	-0.164704\\
0.39599	-0.164212\\
0.401003	-0.163775\\
0.406015	-0.163397\\
0.411028	-0.163068\\
0.41604	-0.162773\\
0.421053	-0.162507\\
0.426065	-0.162263\\
0.431078	-0.162034\\
0.43609	-0.161815\\
0.441103	-0.161598\\
0.446115	-0.161379\\
0.451128	-0.161151\\
0.45614	-0.160907\\
0.461153	-0.160654\\
0.466165	-0.160403\\
0.471178	-0.160154\\
0.47619	-0.159905\\
0.481203	-0.159656\\
0.486216	-0.159407\\
0.491228	-0.159156\\
0.496241	-0.158903\\
0.501253	-0.158648\\
0.506266	-0.15839\\
0.511278	-0.158128\\
0.516291	-0.157862\\
0.521303	-0.157591\\
0.526316	-0.157314\\
0.531328	-0.157032\\
0.536341	-0.156743\\
0.541353	-0.156446\\
0.546366	-0.156142\\
0.551378	-0.155829\\
0.556391	-0.155507\\
0.561404	-0.155176\\
0.566416	-0.154834\\
0.571429	-0.15448\\
0.576441	-0.154115\\
0.581454	-0.153738\\
0.586466	-0.153347\\
0.591479	-0.152942\\
0.596491	-0.152523\\
0.601504	-0.152069\\
0.606516	-0.151566\\
0.611529	-0.151017\\
0.616541	-0.150428\\
0.621554	-0.149804\\
0.626566	-0.14915\\
0.631579	-0.14847\\
0.636591	-0.147769\\
0.641604	-0.147052\\
0.646617	-0.146325\\
0.651629	-0.145591\\
0.656642	-0.144855\\
0.661654	-0.144122\\
0.666667	-0.143398\\
0.671679	-0.142686\\
0.676692	-0.141991\\
0.681704	-0.141318\\
0.686717	-0.140673\\
0.691729	-0.140059\\
0.696742	-0.139481\\
0.701754	-0.138943\\
0.706767	-0.138452\\
0.711779	-0.13801\\
0.716792	-0.137624\\
0.721805	-0.137297\\
0.726817	-0.137034\\
0.73183	-0.136841\\
0.736842	-0.13672\\
0.741855	-0.136678\\
0.746867	-0.136749\\
0.75188	-0.136943\\
0.756892	-0.13723\\
0.761905	-0.137578\\
0.766917	-0.137957\\
0.77193	-0.138336\\
0.776942	-0.138684\\
0.781955	-0.138971\\
0.786967	-0.139164\\
0.79198	-0.139234\\
0.796992	-0.139221\\
0.802005	-0.139188\\
0.807018	-0.139135\\
0.81203	-0.139064\\
0.817043	-0.138976\\
0.822055	-0.138871\\
0.827068	-0.13875\\
0.83208	-0.138615\\
0.837093	-0.138466\\
0.842105	-0.138303\\
0.847118	-0.138129\\
0.85213	-0.137944\\
0.857143	-0.137748\\
0.862155	-0.137543\\
0.867168	-0.13733\\
0.87218	-0.13711\\
0.877193	-0.136883\\
0.882206	-0.136651\\
0.887218	-0.136414\\
0.892231	-0.136174\\
0.897243	-0.135931\\
0.902256	-0.135686\\
0.907268	-0.135441\\
0.912281	-0.135196\\
0.917293	-0.134952\\
0.922306	-0.134711\\
0.927318	-0.134472\\
0.932331	-0.134238\\
0.937343	-0.134008\\
0.942356	-0.133755\\
0.947368	-0.133456\\
0.952381	-0.133126\\
0.957393	-0.132779\\
0.962406	-0.132426\\
0.967419	-0.132083\\
0.972431	-0.131761\\
0.977444	-0.131476\\
0.982456	-0.131239\\
0.987469	-0.131065\\
0.992481	-0.130967\\
0.997494	-0.130913\\
1.00251	-0.130864\\
1.00752	-0.13082\\
1.01253	-0.13078\\
1.01754	-0.130744\\
1.02256	-0.130711\\
1.02757	-0.130681\\
1.03258	-0.130654\\
1.03759	-0.130629\\
1.04261	-0.130607\\
1.04762	-0.130586\\
1.05263	-0.130567\\
1.05764	-0.13055\\
1.06266	-0.130533\\
1.06767	-0.130516\\
1.07268	-0.1305\\
1.07769	-0.130484\\
1.08271	-0.130467\\
1.08772	-0.13045\\
1.09273	-0.130431\\
1.09774	-0.130412\\
1.10276	-0.13039\\
1.10777	-0.130366\\
1.11278	-0.13034\\
1.11779	-0.130312\\
1.12281	-0.13028\\
1.12782	-0.130245\\
1.13283	-0.130206\\
1.13784	-0.13016\\
1.14286	-0.130103\\
1.14787	-0.130037\\
1.15288	-0.129961\\
1.15789	-0.129878\\
1.16291	-0.129787\\
1.16792	-0.129689\\
1.17293	-0.129586\\
1.17794	-0.129478\\
1.18296	-0.129367\\
1.18797	-0.129251\\
1.19298	-0.129134\\
1.19799	-0.129015\\
1.20301	-0.128895\\
1.20802	-0.128776\\
1.21303	-0.128657\\
1.21805	-0.12854\\
1.22306	-0.128426\\
1.22807	-0.128315\\
1.23308	-0.128209\\
1.2381	-0.128107\\
1.24311	-0.128012\\
1.24812	-0.127923\\
1.25313	-0.127842\\
1.25815	-0.127769\\
1.26316	-0.127706\\
1.26817	-0.127652\\
1.27318	-0.127606\\
1.2782	-0.127563\\
1.28321	-0.127522\\
1.28822	-0.127484\\
1.29323	-0.127448\\
1.29825	-0.127413\\
1.30326	-0.12738\\
1.30827	-0.127348\\
1.31328	-0.127317\\
1.3183	-0.127285\\
1.32331	-0.127254\\
1.32832	-0.127223\\
1.33333	-0.127191\\
1.33835	-0.12716\\
1.34336	-0.12713\\
1.34837	-0.127101\\
1.35338	-0.127072\\
1.3584	-0.127044\\
1.36341	-0.127017\\
1.36842	-0.12699\\
1.37343	-0.126963\\
1.37845	-0.126936\\
1.38346	-0.12691\\
1.38847	-0.126884\\
1.39348	-0.126859\\
1.3985	-0.126833\\
1.40351	-0.126807\\
1.40852	-0.126781\\
1.41353	-0.126756\\
1.41855	-0.12673\\
1.42356	-0.126703\\
1.42857	-0.126677\\
1.43358	-0.12665\\
1.4386	-0.126623\\
1.44361	-0.126595\\
1.44862	-0.126567\\
1.45363	-0.126538\\
1.45865	-0.126509\\
1.46366	-0.126478\\
1.46867	-0.126447\\
1.47368	-0.126414\\
1.4787	-0.126378\\
1.48371	-0.12634\\
1.48872	-0.126299\\
1.49373	-0.126256\\
1.49875	-0.126212\\
1.50376	-0.126166\\
1.50877	-0.126119\\
1.51378	-0.126071\\
1.5188	-0.126023\\
1.52381	-0.125974\\
1.52882	-0.125926\\
1.53383	-0.125878\\
1.53885	-0.12583\\
1.54386	-0.125783\\
1.54887	-0.125738\\
1.55388	-0.125694\\
1.5589	-0.125652\\
1.56391	-0.125612\\
1.56892	-0.125575\\
1.57393	-0.12554\\
1.57895	-0.125508\\
1.58396	-0.125479\\
1.58897	-0.125454\\
1.59398	-0.125433\\
1.599	-0.125416\\
1.60401	-0.125404\\
1.60902	-0.125396\\
1.61404	-0.125394\\
1.61905	-0.125397\\
1.62406	-0.125404\\
1.62907	-0.125413\\
1.63409	-0.125424\\
1.6391	-0.125436\\
1.64411	-0.125447\\
1.64912	-0.125457\\
1.65414	-0.125464\\
1.65915	-0.125466\\
1.66416	-0.125464\\
1.66917	-0.125458\\
1.67419	-0.125451\\
1.6792	-0.125442\\
1.68421	-0.125432\\
1.68922	-0.12542\\
1.69424	-0.125407\\
1.69925	-0.125394\\
1.70426	-0.125379\\
1.70927	-0.125362\\
1.71429	-0.125346\\
1.7193	-0.125328\\
1.72431	-0.125309\\
1.72932	-0.12529\\
1.73434	-0.125269\\
1.73935	-0.125249\\
1.74436	-0.125227\\
1.74937	-0.125206\\
1.75439	-0.125184\\
1.7594	-0.125161\\
1.76441	-0.125138\\
1.76942	-0.125116\\
1.77444	-0.125092\\
1.77945	-0.125069\\
1.78446	-0.125046\\
1.78947	-0.125023\\
1.79449	-0.125001\\
1.7995	-0.124978\\
1.80451	-0.124953\\
1.80952	-0.124925\\
1.81454	-0.124894\\
1.81955	-0.124861\\
1.82456	-0.124827\\
1.82957	-0.124793\\
1.83459	-0.12476\\
1.8396	-0.124727\\
1.84461	-0.124697\\
1.84962	-0.124671\\
1.85464	-0.124648\\
1.85965	-0.12463\\
1.86466	-0.124614\\
1.86967	-0.124599\\
1.87469	-0.124584\\
1.8797	-0.12457\\
1.88471	-0.124556\\
1.88972	-0.124543\\
1.89474	-0.12453\\
1.89975	-0.124518\\
1.90476	-0.124507\\
1.90977	-0.124496\\
1.91479	-0.124487\\
1.9198	-0.124478\\
1.92481	-0.124471\\
1.92982	-0.124464\\
1.93484	-0.124459\\
1.93985	-0.124455\\
1.94486	-0.124455\\
1.94987	-0.124456\\
1.95489	-0.12446\\
1.9599	-0.124466\\
1.96491	-0.124473\\
1.96992	-0.124483\\
1.97494	-0.124494\\
1.97995	-0.124507\\
1.98496	-0.124521\\
1.98997	-0.124537\\
1.99499	-0.124553\\
2	-0.12457\\
2	-0.124327\\
1.99499	-0.124332\\
1.98997	-0.124337\\
1.98496	-0.124342\\
1.97995	-0.124347\\
1.97494	-0.124352\\
1.96992	-0.124357\\
1.96491	-0.124362\\
1.9599	-0.124367\\
1.95489	-0.124372\\
1.94987	-0.124377\\
1.94486	-0.124382\\
1.93985	-0.124387\\
1.93484	-0.124392\\
1.92982	-0.124397\\
1.92481	-0.124402\\
1.9198	-0.124407\\
1.91479	-0.124412\\
1.90977	-0.124417\\
1.90476	-0.124422\\
1.89975	-0.124427\\
1.89474	-0.124432\\
1.88972	-0.124437\\
1.88471	-0.124442\\
1.8797	-0.124447\\
1.87469	-0.124452\\
1.86967	-0.124457\\
1.86466	-0.124462\\
1.85965	-0.124467\\
1.85464	-0.124473\\
1.84962	-0.124478\\
1.84461	-0.124483\\
1.8396	-0.124488\\
1.83459	-0.124493\\
1.82957	-0.124498\\
1.82456	-0.124503\\
1.81955	-0.124508\\
1.81454	-0.124513\\
1.80952	-0.124518\\
1.80451	-0.124523\\
1.7995	-0.124528\\
1.79449	-0.124533\\
1.78947	-0.124538\\
1.78446	-0.124543\\
1.77945	-0.124548\\
1.77444	-0.124553\\
1.76942	-0.124558\\
1.76441	-0.124563\\
1.7594	-0.124568\\
1.75439	-0.124573\\
1.74937	-0.124578\\
1.74436	-0.124584\\
1.73935	-0.124589\\
1.73434	-0.124594\\
1.72932	-0.124599\\
1.72431	-0.124604\\
1.7193	-0.124609\\
1.71429	-0.124614\\
1.70927	-0.12462\\
1.70426	-0.124625\\
1.69925	-0.12463\\
1.69424	-0.124635\\
1.68922	-0.124641\\
1.68421	-0.124646\\
1.6792	-0.124651\\
1.67419	-0.124657\\
1.66917	-0.124662\\
1.66416	-0.124667\\
1.65915	-0.124673\\
1.65414	-0.124678\\
1.64912	-0.124684\\
1.64411	-0.124689\\
1.6391	-0.124695\\
1.63409	-0.1247\\
1.62907	-0.124706\\
1.62406	-0.124711\\
1.61905	-0.124717\\
1.61404	-0.124722\\
1.60902	-0.124728\\
1.60401	-0.124733\\
1.599	-0.124739\\
1.59398	-0.124745\\
1.58897	-0.12475\\
1.58396	-0.124756\\
1.57895	-0.124761\\
1.57393	-0.124767\\
1.56892	-0.124773\\
1.56391	-0.124778\\
1.5589	-0.124784\\
1.55388	-0.12479\\
1.54887	-0.124795\\
1.54386	-0.124801\\
1.53885	-0.124807\\
1.53383	-0.124812\\
1.52882	-0.124818\\
1.52381	-0.124824\\
1.5188	-0.12483\\
1.51378	-0.124835\\
1.50877	-0.124841\\
1.50376	-0.124848\\
1.49875	-0.124855\\
1.49373	-0.124863\\
1.48872	-0.124871\\
1.48371	-0.12488\\
1.4787	-0.124889\\
1.47368	-0.124898\\
1.46867	-0.124908\\
1.46366	-0.124918\\
1.45865	-0.124929\\
1.45363	-0.124939\\
1.44862	-0.12495\\
1.44361	-0.124961\\
1.4386	-0.124972\\
1.43358	-0.124983\\
1.42857	-0.124994\\
1.42356	-0.125005\\
1.41855	-0.125016\\
1.41353	-0.125027\\
1.40852	-0.125037\\
1.40351	-0.125048\\
1.3985	-0.125058\\
1.39348	-0.125068\\
1.38847	-0.125078\\
1.38346	-0.125088\\
1.37845	-0.125097\\
1.37343	-0.125106\\
1.36842	-0.125114\\
1.36341	-0.125122\\
1.3584	-0.125129\\
1.35338	-0.125135\\
1.34837	-0.125141\\
1.34336	-0.125145\\
1.33835	-0.125144\\
1.33333	-0.125139\\
1.32832	-0.125133\\
1.32331	-0.125126\\
1.3183	-0.125118\\
1.31328	-0.125112\\
1.30827	-0.125108\\
1.30326	-0.125107\\
1.29825	-0.12511\\
1.29323	-0.125117\\
1.28822	-0.125126\\
1.28321	-0.125136\\
1.2782	-0.125148\\
1.27318	-0.125162\\
1.26817	-0.125177\\
1.26316	-0.125194\\
1.25815	-0.125211\\
1.25313	-0.12523\\
1.24812	-0.12525\\
1.24311	-0.125271\\
1.2381	-0.125293\\
1.23308	-0.125316\\
1.22807	-0.125339\\
1.22306	-0.125363\\
1.21805	-0.125388\\
1.21303	-0.125413\\
1.20802	-0.125438\\
1.20301	-0.125463\\
1.19799	-0.125489\\
1.19298	-0.125514\\
1.18797	-0.12554\\
1.18296	-0.125565\\
1.17794	-0.12559\\
1.17293	-0.125615\\
1.16792	-0.12564\\
1.16291	-0.125663\\
1.15789	-0.125687\\
1.15288	-0.125709\\
1.14787	-0.125731\\
1.14286	-0.125752\\
1.13784	-0.125773\\
1.13283	-0.125795\\
1.12782	-0.125816\\
1.12281	-0.125838\\
1.11779	-0.125859\\
1.11278	-0.125881\\
1.10777	-0.125902\\
1.10276	-0.125924\\
1.09774	-0.125946\\
1.09273	-0.125968\\
1.08772	-0.125991\\
1.08271	-0.126014\\
1.07769	-0.126037\\
1.07268	-0.12606\\
1.06767	-0.126084\\
1.06266	-0.126108\\
1.05764	-0.126133\\
1.05263	-0.126158\\
1.04762	-0.126183\\
1.04261	-0.126209\\
1.03759	-0.126236\\
1.03258	-0.126263\\
1.02757	-0.126291\\
1.02256	-0.126319\\
1.01754	-0.126348\\
1.01253	-0.126378\\
1.00752	-0.126409\\
1.00251	-0.12644\\
0.997494	-0.126476\\
0.992481	-0.126518\\
0.987469	-0.126564\\
0.982456	-0.126612\\
0.977444	-0.126661\\
0.972431	-0.126707\\
0.967419	-0.126748\\
0.962406	-0.126783\\
0.957393	-0.126809\\
0.952381	-0.126824\\
0.947368	-0.126817\\
0.942356	-0.126784\\
0.937343	-0.126732\\
0.932331	-0.126667\\
0.927318	-0.126595\\
0.922306	-0.126523\\
0.917293	-0.126458\\
0.912281	-0.126407\\
0.907268	-0.126375\\
0.902256	-0.126369\\
0.897243	-0.126391\\
0.892231	-0.126435\\
0.887218	-0.1265\\
0.882206	-0.126585\\
0.877193	-0.126686\\
0.87218	-0.126803\\
0.867168	-0.126934\\
0.862155	-0.127076\\
0.857143	-0.127229\\
0.85213	-0.127389\\
0.847118	-0.127556\\
0.842105	-0.127727\\
0.837093	-0.127901\\
0.83208	-0.128075\\
0.827068	-0.128249\\
0.822055	-0.12842\\
0.817043	-0.128586\\
0.81203	-0.128746\\
0.807018	-0.128898\\
0.802005	-0.12904\\
0.796992	-0.129175\\
0.79198	-0.129308\\
0.786967	-0.12944\\
0.781955	-0.129569\\
0.776942	-0.129697\\
0.77193	-0.129824\\
0.766917	-0.12995\\
0.761905	-0.130076\\
0.756892	-0.1302\\
0.75188	-0.130325\\
0.746867	-0.130449\\
0.741855	-0.130573\\
0.736842	-0.130698\\
0.73183	-0.130824\\
0.726817	-0.130951\\
0.721805	-0.131079\\
0.716792	-0.131208\\
0.711779	-0.131339\\
0.706767	-0.131472\\
0.701754	-0.131606\\
0.696742	-0.131744\\
0.691729	-0.131884\\
0.686717	-0.132026\\
0.681704	-0.132172\\
0.676692	-0.132321\\
0.671679	-0.132473\\
0.666667	-0.132629\\
0.661654	-0.132789\\
0.656642	-0.132953\\
0.651629	-0.133121\\
0.646617	-0.133293\\
0.641604	-0.133469\\
0.636591	-0.13365\\
0.631579	-0.133836\\
0.626566	-0.134026\\
0.621554	-0.134221\\
0.616541	-0.134421\\
0.611529	-0.134626\\
0.606516	-0.134836\\
0.601504	-0.135058\\
0.596491	-0.1353\\
0.591479	-0.135559\\
0.586466	-0.135834\\
0.581454	-0.136123\\
0.576441	-0.136425\\
0.571429	-0.136738\\
0.566416	-0.137061\\
0.561404	-0.137391\\
0.556391	-0.137727\\
0.551378	-0.138068\\
0.546366	-0.138411\\
0.541353	-0.138754\\
0.536341	-0.139097\\
0.531328	-0.139437\\
0.526316	-0.139772\\
0.521303	-0.140101\\
0.516291	-0.140422\\
0.511278	-0.140734\\
0.506266	-0.141033\\
0.501253	-0.141319\\
0.496241	-0.14159\\
0.491228	-0.141843\\
0.486216	-0.142078\\
0.481203	-0.142292\\
0.47619	-0.142484\\
0.471178	-0.142652\\
0.466165	-0.142794\\
0.461153	-0.142908\\
0.45614	-0.143014\\
0.451128	-0.143123\\
0.446115	-0.143226\\
0.441103	-0.143314\\
0.43609	-0.143378\\
0.431078	-0.143408\\
0.426065	-0.143394\\
0.421053	-0.143328\\
0.41604	-0.143199\\
0.411028	-0.142999\\
0.406015	-0.142731\\
0.401003	-0.142412\\
0.39599	-0.142051\\
0.390977	-0.141659\\
0.385965	-0.141245\\
0.380952	-0.140819\\
0.37594	-0.140392\\
0.370927	-0.139972\\
0.365915	-0.13957\\
0.360902	-0.139196\\
0.35589	-0.138975\\
0.350877	-0.139013\\
0.345865	-0.139288\\
0.340852	-0.13978\\
0.33584	-0.140467\\
0.330827	-0.141328\\
0.325815	-0.142341\\
0.320802	-0.143485\\
0.315789	-0.14474\\
0.310777	-0.146082\\
0.305764	-0.147492\\
0.300752	-0.148946\\
0.295739	-0.150424\\
0.290727	-0.151904\\
0.285714	-0.153364\\
0.280702	-0.154781\\
0.275689	-0.156135\\
0.270677	-0.157403\\
0.265664	-0.158563\\
0.260652	-0.159592\\
0.255639	-0.160542\\
0.250627	-0.161474\\
0.245614	-0.16239\\
0.240602	-0.163286\\
0.235589	-0.164162\\
0.230576	-0.165016\\
0.225564	-0.165847\\
0.220551	-0.166652\\
0.215539	-0.16743\\
0.210526	-0.168179\\
0.205514	-0.168897\\
0.200501	-0.169582\\
0.195489	-0.170231\\
0.190476	-0.170843\\
0.185464	-0.171414\\
0.180451	-0.171944\\
0.175439	-0.17243\\
0.170426	-0.172868\\
0.165414	-0.173257\\
0.160401	-0.173595\\
0.155388	-0.173878\\
0.150376	-0.174104\\
0.145363	-0.174271\\
0.140351	-0.174376\\
0.135338	-0.174416\\
0.130326	-0.17439\\
0.125313	-0.174292\\
0.120301	-0.174123\\
0.115288	-0.173877\\
0.110276	-0.173555\\
0.105263	-0.173151\\
0.100251	-0.172664\\
0.0952381	-0.172089\\
0.0902256	-0.171427\\
0.085213	-0.170671\\
0.0802005	-0.169822\\
0.075188	-0.168872\\
0.0701754	-0.167824\\
0.0651629	-0.16667\\
0.0601504	-0.165412\\
0.0551378	-0.160712\\
0.0501253	-0.154439\\
0.0451128	-0.146434\\
0.0401003	-0.136758\\
0.0350877	-0.125283\\
0.0300752	-0.112126\\
0.0250627	-0.0972\\
0.0200501	-0.0806733\\
0.0150376	-0.0625122\\
0.0100251	-0.0429347\\
0.00501253	-0.032104\\
0	-0.0109093\\
}--cycle;

\addplot[area legend, draw=none, fill=mycolor1, fill opacity=0.9, forget plot]
table[row sep=crcr] {%
x	y\\
0	0.0109093\\
0.00501253	0.0227473\\
0.0100251	0.0246048\\
0.0150376	0.0354628\\
0.0200501	0.0451311\\
0.0250627	0.053419\\
0.0300752	0.0603506\\
0.0350877	0.0657782\\
0.0401003	0.0697712\\
0.0451128	0.0722283\\
0.0501253	0.0732581\\
0.0551378	0.0728092\\
0.0601504	0.0710206\\
0.0651629	0.0660328\\
0.0701754	0.0611648\\
0.075188	0.0564196\\
0.0802005	0.0517886\\
0.085213	0.0472739\\
0.0902256	0.0428676\\
0.0952381	0.0385713\\
0.100251	0.0343776\\
0.105263	0.0302875\\
0.110276	0.0262946\\
0.115288	0.0223991\\
0.120301	0.0185953\\
0.125313	0.0148828\\
0.130326	0.0112565\\
0.135338	0.00771582\\
0.140351	0.00425584\\
0.145363	0.000875648\\
0.150376	-0.00242912\\
0.155388	-0.00565976\\
0.160401	-0.00882021\\
0.165414	-0.0119121\\
0.170426	-0.0149389\\
0.175439	-0.0179026\\
0.180451	-0.0208063\\
0.185464	-0.023652\\
0.190476	-0.0264426\\
0.195489	-0.0291803\\
0.200501	-0.0318677\\
0.205514	-0.0345069\\
0.210526	-0.0371003\\
0.215539	-0.0396501\\
0.220551	-0.0421583\\
0.225564	-0.0446272\\
0.230576	-0.0470584\\
0.235589	-0.0494541\\
0.240602	-0.0518159\\
0.245614	-0.0541458\\
0.250627	-0.056445\\
0.255639	-0.0587154\\
0.260652	-0.0609581\\
0.265664	-0.0632402\\
0.270677	-0.0656136\\
0.275689	-0.0680589\\
0.280702	-0.0705552\\
0.285714	-0.0730827\\
0.290727	-0.0756207\\
0.295739	-0.0781491\\
0.300752	-0.0806468\\
0.305764	-0.0830938\\
0.310777	-0.0854687\\
0.315789	-0.0877514\\
0.320802	-0.0899205\\
0.325815	-0.0919555\\
0.330827	-0.093835\\
0.33584	-0.0955385\\
0.340852	-0.0970444\\
0.345865	-0.0983321\\
0.350877	-0.0993799\\
0.35589	-0.100167\\
0.360902	-0.100672\\
0.365915	-0.101001\\
0.370927	-0.101278\\
0.37594	-0.101514\\
0.380952	-0.101719\\
0.385965	-0.101902\\
0.390977	-0.102073\\
0.39599	-0.102242\\
0.401003	-0.102419\\
0.406015	-0.102613\\
0.411028	-0.102835\\
0.41604	-0.103101\\
0.421053	-0.103415\\
0.426065	-0.103768\\
0.431078	-0.104151\\
0.43609	-0.104554\\
0.441103	-0.104968\\
0.446115	-0.105385\\
0.451128	-0.105796\\
0.45614	-0.106191\\
0.461153	-0.106561\\
0.466165	-0.10692\\
0.471178	-0.107288\\
0.47619	-0.107662\\
0.481203	-0.108042\\
0.486216	-0.108427\\
0.491228	-0.108816\\
0.496241	-0.109208\\
0.501253	-0.109601\\
0.506266	-0.109996\\
0.511278	-0.11039\\
0.516291	-0.110783\\
0.521303	-0.111174\\
0.526316	-0.111562\\
0.531328	-0.111946\\
0.536341	-0.112324\\
0.541353	-0.112697\\
0.546366	-0.113062\\
0.551378	-0.113419\\
0.556391	-0.113767\\
0.561404	-0.114105\\
0.566416	-0.114432\\
0.571429	-0.114746\\
0.576441	-0.115047\\
0.581454	-0.115333\\
0.586466	-0.115604\\
0.591479	-0.115857\\
0.596491	-0.116093\\
0.601504	-0.11631\\
0.606516	-0.116506\\
0.611529	-0.116689\\
0.616541	-0.116867\\
0.621554	-0.11704\\
0.626566	-0.117208\\
0.631579	-0.117372\\
0.636591	-0.117532\\
0.641604	-0.117687\\
0.646617	-0.11784\\
0.651629	-0.117989\\
0.656642	-0.118135\\
0.661654	-0.118278\\
0.666667	-0.118419\\
0.671679	-0.118558\\
0.676692	-0.118694\\
0.681704	-0.118828\\
0.686717	-0.11896\\
0.691729	-0.119091\\
0.696742	-0.11922\\
0.701754	-0.119348\\
0.706767	-0.119474\\
0.711779	-0.1196\\
0.716792	-0.119725\\
0.721805	-0.119849\\
0.726817	-0.119972\\
0.73183	-0.120094\\
0.736842	-0.120217\\
0.741855	-0.120339\\
0.746867	-0.120461\\
0.75188	-0.120583\\
0.756892	-0.120705\\
0.761905	-0.120827\\
0.766917	-0.120949\\
0.77193	-0.121073\\
0.776942	-0.121196\\
0.781955	-0.121321\\
0.786967	-0.121446\\
0.79198	-0.121573\\
0.796992	-0.121701\\
0.802005	-0.12183\\
0.807018	-0.121965\\
0.81203	-0.122109\\
0.817043	-0.122261\\
0.822055	-0.122418\\
0.827068	-0.122579\\
0.83208	-0.122741\\
0.837093	-0.122904\\
0.842105	-0.123065\\
0.847118	-0.123222\\
0.85213	-0.123375\\
0.857143	-0.12352\\
0.862155	-0.123656\\
0.867168	-0.123782\\
0.87218	-0.123895\\
0.877193	-0.123995\\
0.882206	-0.124078\\
0.887218	-0.124144\\
0.892231	-0.12419\\
0.897243	-0.124216\\
0.902256	-0.124218\\
0.907268	-0.124193\\
0.912281	-0.124142\\
0.917293	-0.12407\\
0.922306	-0.123986\\
0.927318	-0.123895\\
0.932331	-0.123805\\
0.937343	-0.123721\\
0.942356	-0.12365\\
0.947368	-0.1236\\
0.952381	-0.123576\\
0.957393	-0.123573\\
0.962406	-0.123583\\
0.967419	-0.123602\\
0.972431	-0.123629\\
0.977444	-0.12366\\
0.982456	-0.123694\\
0.987469	-0.123729\\
0.992481	-0.123763\\
0.997494	-0.123793\\
1.00251	-0.123817\\
1.00752	-0.123838\\
1.01253	-0.123858\\
1.01754	-0.123879\\
1.02256	-0.123899\\
1.02757	-0.123919\\
1.03258	-0.123939\\
1.03759	-0.123958\\
1.04261	-0.123978\\
1.04762	-0.123997\\
1.05263	-0.124016\\
1.05764	-0.124035\\
1.06266	-0.124054\\
1.06767	-0.124073\\
1.07268	-0.124091\\
1.07769	-0.12411\\
1.08271	-0.124128\\
1.08772	-0.124146\\
1.09273	-0.124163\\
1.09774	-0.124181\\
1.10276	-0.124198\\
1.10777	-0.124215\\
1.11278	-0.124232\\
1.11779	-0.124248\\
1.12281	-0.124264\\
1.12782	-0.12428\\
1.13283	-0.124296\\
1.13784	-0.124312\\
1.14286	-0.124327\\
1.14787	-0.124342\\
1.15288	-0.124358\\
1.15789	-0.124373\\
1.16291	-0.12439\\
1.16792	-0.124406\\
1.17293	-0.124423\\
1.17794	-0.12444\\
1.18296	-0.124457\\
1.18797	-0.124474\\
1.19298	-0.124491\\
1.19799	-0.124508\\
1.20301	-0.124525\\
1.20802	-0.124541\\
1.21303	-0.124556\\
1.21805	-0.124572\\
1.22306	-0.124586\\
1.22807	-0.1246\\
1.23308	-0.124613\\
1.2381	-0.124625\\
1.24311	-0.124636\\
1.24812	-0.124646\\
1.25313	-0.124655\\
1.25815	-0.124663\\
1.26316	-0.12467\\
1.26817	-0.124675\\
1.27318	-0.124678\\
1.2782	-0.124681\\
1.28321	-0.124681\\
1.28822	-0.12468\\
1.29323	-0.124677\\
1.29825	-0.124672\\
1.30326	-0.124664\\
1.30827	-0.124651\\
1.31328	-0.124635\\
1.3183	-0.124617\\
1.32331	-0.124598\\
1.32832	-0.124579\\
1.33333	-0.124561\\
1.33835	-0.124546\\
1.34336	-0.124533\\
1.34837	-0.124525\\
1.35338	-0.12452\\
1.3584	-0.124515\\
1.36341	-0.124511\\
1.36842	-0.124508\\
1.37343	-0.124505\\
1.37845	-0.124503\\
1.38346	-0.124501\\
1.38847	-0.1245\\
1.39348	-0.124499\\
1.3985	-0.124499\\
1.40351	-0.124499\\
1.40852	-0.124499\\
1.41353	-0.124499\\
1.41855	-0.1245\\
1.42356	-0.1245\\
1.42857	-0.124501\\
1.43358	-0.124502\\
1.4386	-0.124503\\
1.44361	-0.124504\\
1.44862	-0.124505\\
1.45363	-0.124506\\
1.45865	-0.124506\\
1.46366	-0.124507\\
1.46867	-0.124507\\
1.47368	-0.124507\\
1.4787	-0.124506\\
1.48371	-0.124505\\
1.48872	-0.124504\\
1.49373	-0.124502\\
1.49875	-0.1245\\
1.50376	-0.124497\\
1.50877	-0.124494\\
1.51378	-0.12449\\
1.5188	-0.124486\\
1.52381	-0.124481\\
1.52882	-0.124477\\
1.53383	-0.124473\\
1.53885	-0.124468\\
1.54386	-0.124464\\
1.54887	-0.124459\\
1.55388	-0.124455\\
1.5589	-0.12445\\
1.56391	-0.124445\\
1.56892	-0.124441\\
1.57393	-0.124436\\
1.57895	-0.124431\\
1.58396	-0.124427\\
1.58897	-0.124422\\
1.59398	-0.124417\\
1.599	-0.124412\\
1.60401	-0.124408\\
1.60902	-0.124403\\
1.61404	-0.124398\\
1.61905	-0.124393\\
1.62406	-0.124388\\
1.62907	-0.124384\\
1.63409	-0.124379\\
1.6391	-0.124374\\
1.64411	-0.124369\\
1.64912	-0.124364\\
1.65414	-0.124359\\
1.65915	-0.124354\\
1.66416	-0.12435\\
1.66917	-0.124345\\
1.67419	-0.12434\\
1.6792	-0.124335\\
1.68421	-0.12433\\
1.68922	-0.124326\\
1.69424	-0.124321\\
1.69925	-0.124316\\
1.70426	-0.124311\\
1.70927	-0.124306\\
1.71429	-0.124302\\
1.7193	-0.124297\\
1.72431	-0.124292\\
1.72932	-0.124288\\
1.73434	-0.124283\\
1.73935	-0.124278\\
1.74436	-0.124274\\
1.74937	-0.124269\\
1.75439	-0.124264\\
1.7594	-0.12426\\
1.76441	-0.124255\\
1.76942	-0.12425\\
1.77444	-0.124246\\
1.77945	-0.124241\\
1.78446	-0.124237\\
1.78947	-0.124232\\
1.79449	-0.124228\\
1.7995	-0.124223\\
1.80451	-0.124219\\
1.80952	-0.124214\\
1.81454	-0.12421\\
1.81955	-0.124205\\
1.82456	-0.124201\\
1.82957	-0.124196\\
1.83459	-0.124192\\
1.8396	-0.124187\\
1.84461	-0.124183\\
1.84962	-0.124178\\
1.85464	-0.124174\\
1.85965	-0.12417\\
1.86466	-0.124165\\
1.86967	-0.124161\\
1.87469	-0.124157\\
1.8797	-0.124152\\
1.88471	-0.124148\\
1.88972	-0.124144\\
1.89474	-0.124139\\
1.89975	-0.124135\\
1.90476	-0.124131\\
1.90977	-0.124126\\
1.91479	-0.124122\\
1.9198	-0.124118\\
1.92481	-0.124113\\
1.92982	-0.124109\\
1.93484	-0.124105\\
1.93985	-0.1241\\
1.94486	-0.124096\\
1.94987	-0.124092\\
1.95489	-0.124087\\
1.9599	-0.124083\\
1.96491	-0.124079\\
1.96992	-0.124074\\
1.97494	-0.12407\\
1.97995	-0.124066\\
1.98496	-0.124062\\
1.98997	-0.124057\\
1.99499	-0.124053\\
2	-0.124049\\
2	-0.123805\\
1.99499	-0.123831\\
1.98997	-0.123857\\
1.98496	-0.123882\\
1.97995	-0.123905\\
1.97494	-0.123927\\
1.96992	-0.123948\\
1.96491	-0.123967\\
1.9599	-0.123984\\
1.95489	-0.123999\\
1.94987	-0.124012\\
1.94486	-0.124023\\
1.93985	-0.124032\\
1.93484	-0.124038\\
1.92982	-0.124042\\
1.92481	-0.124045\\
1.9198	-0.124046\\
1.91479	-0.124047\\
1.90977	-0.124047\\
1.90476	-0.124046\\
1.89975	-0.124044\\
1.89474	-0.124041\\
1.88972	-0.124038\\
1.88471	-0.124034\\
1.8797	-0.12403\\
1.87469	-0.124025\\
1.86967	-0.124019\\
1.86466	-0.124014\\
1.85965	-0.124008\\
1.85464	-0.123999\\
1.84962	-0.123986\\
1.84461	-0.123968\\
1.8396	-0.123948\\
1.83459	-0.123925\\
1.82957	-0.123901\\
1.82456	-0.123876\\
1.81955	-0.123851\\
1.81454	-0.123828\\
1.80952	-0.123807\\
1.80451	-0.123788\\
1.7995	-0.123773\\
1.79449	-0.12376\\
1.78947	-0.123747\\
1.78446	-0.123733\\
1.77945	-0.12372\\
1.77444	-0.123706\\
1.76942	-0.123693\\
1.76441	-0.12368\\
1.7594	-0.123667\\
1.75439	-0.123654\\
1.74937	-0.123642\\
1.74436	-0.12363\\
1.73935	-0.123618\\
1.73434	-0.123607\\
1.72932	-0.123597\\
1.72431	-0.123587\\
1.7193	-0.123579\\
1.71429	-0.123571\\
1.70927	-0.123564\\
1.70426	-0.123558\\
1.69925	-0.123553\\
1.69424	-0.123549\\
1.68922	-0.123546\\
1.68421	-0.123545\\
1.6792	-0.123544\\
1.67419	-0.123546\\
1.66917	-0.123549\\
1.66416	-0.123553\\
1.65915	-0.123561\\
1.65414	-0.123574\\
1.64912	-0.123591\\
1.64411	-0.123611\\
1.6391	-0.123632\\
1.63409	-0.123655\\
1.62907	-0.123676\\
1.62406	-0.123696\\
1.61905	-0.123713\\
1.61404	-0.123726\\
1.60902	-0.123734\\
1.60401	-0.123737\\
1.599	-0.123735\\
1.59398	-0.123728\\
1.58897	-0.123718\\
1.58396	-0.123703\\
1.57895	-0.123685\\
1.57393	-0.123663\\
1.56892	-0.123639\\
1.56391	-0.123612\\
1.5589	-0.123582\\
1.55388	-0.12355\\
1.54887	-0.123516\\
1.54386	-0.123481\\
1.53885	-0.123445\\
1.53383	-0.123407\\
1.52882	-0.123369\\
1.52381	-0.123331\\
1.5188	-0.123292\\
1.51378	-0.123254\\
1.50877	-0.123216\\
1.50376	-0.123179\\
1.49875	-0.123143\\
1.49373	-0.123109\\
1.48872	-0.123076\\
1.48371	-0.123045\\
1.4787	-0.123017\\
1.47368	-0.122991\\
1.46867	-0.122968\\
1.46366	-0.122947\\
1.45865	-0.122926\\
1.45363	-0.122907\\
1.44862	-0.122888\\
1.44361	-0.12287\\
1.4386	-0.122852\\
1.43358	-0.122835\\
1.42857	-0.122818\\
1.42356	-0.122802\\
1.41855	-0.122786\\
1.41353	-0.12277\\
1.40852	-0.122755\\
1.40351	-0.122739\\
1.3985	-0.122724\\
1.39348	-0.122709\\
1.38847	-0.122694\\
1.38346	-0.122679\\
1.37845	-0.122663\\
1.37343	-0.122648\\
1.36842	-0.122632\\
1.36341	-0.122616\\
1.3584	-0.1226\\
1.35338	-0.122583\\
1.34837	-0.122565\\
1.34336	-0.122547\\
1.33835	-0.122529\\
1.33333	-0.12251\\
1.32832	-0.12249\\
1.32331	-0.12247\\
1.3183	-0.12245\\
1.31328	-0.122431\\
1.30827	-0.122411\\
1.30326	-0.12239\\
1.29825	-0.122369\\
1.29323	-0.122346\\
1.28822	-0.122322\\
1.28321	-0.122295\\
1.2782	-0.122266\\
1.27318	-0.122234\\
1.26817	-0.122199\\
1.26316	-0.122158\\
1.25815	-0.122105\\
1.25313	-0.122044\\
1.24812	-0.121974\\
1.24311	-0.121896\\
1.2381	-0.121811\\
1.23308	-0.12172\\
1.22807	-0.121624\\
1.22306	-0.121523\\
1.21805	-0.121419\\
1.21303	-0.121312\\
1.20802	-0.121203\\
1.20301	-0.121093\\
1.19799	-0.120982\\
1.19298	-0.120872\\
1.18797	-0.120763\\
1.18296	-0.120656\\
1.17794	-0.120552\\
1.17293	-0.120452\\
1.16792	-0.120356\\
1.16291	-0.120266\\
1.15789	-0.120182\\
1.15288	-0.120105\\
1.14787	-0.120036\\
1.14286	-0.119976\\
1.13784	-0.119925\\
1.13283	-0.119885\\
1.12782	-0.119852\\
1.12281	-0.119822\\
1.11779	-0.119796\\
1.11278	-0.119772\\
1.10777	-0.119751\\
1.10276	-0.119732\\
1.09774	-0.119715\\
1.09273	-0.1197\\
1.08772	-0.119687\\
1.08271	-0.119674\\
1.07769	-0.119662\\
1.07268	-0.119651\\
1.06767	-0.11964\\
1.06266	-0.11963\\
1.05764	-0.119618\\
1.05263	-0.119607\\
1.04762	-0.119594\\
1.04261	-0.11958\\
1.03759	-0.119565\\
1.03258	-0.119548\\
1.02757	-0.119529\\
1.02256	-0.119507\\
1.01754	-0.119483\\
1.01253	-0.119457\\
1.00752	-0.119426\\
1.00251	-0.119393\\
0.997494	-0.119356\\
0.992481	-0.119314\\
0.987469	-0.119228\\
0.982456	-0.119068\\
0.977444	-0.118845\\
0.972431	-0.118574\\
0.967419	-0.118268\\
0.962406	-0.11794\\
0.957393	-0.117604\\
0.952381	-0.117273\\
0.947368	-0.11696\\
0.942356	-0.11668\\
0.937343	-0.116444\\
0.932331	-0.116234\\
0.927318	-0.116018\\
0.922306	-0.115799\\
0.917293	-0.115576\\
0.912281	-0.115352\\
0.907268	-0.115127\\
0.902256	-0.114901\\
0.897243	-0.114676\\
0.892231	-0.114452\\
0.887218	-0.114231\\
0.882206	-0.114012\\
0.877193	-0.113798\\
0.87218	-0.113589\\
0.867168	-0.113386\\
0.862155	-0.113189\\
0.857143	-0.113\\
0.85213	-0.11282\\
0.847118	-0.112649\\
0.842105	-0.112488\\
0.837093	-0.112339\\
0.83208	-0.112202\\
0.827068	-0.112077\\
0.822055	-0.111967\\
0.817043	-0.111871\\
0.81203	-0.111791\\
0.807018	-0.111728\\
0.802005	-0.111683\\
0.796992	-0.111655\\
0.79198	-0.111648\\
0.786967	-0.111722\\
0.781955	-0.11192\\
0.776942	-0.112209\\
0.77193	-0.112561\\
0.766917	-0.112942\\
0.761905	-0.113324\\
0.756892	-0.113675\\
0.75188	-0.113964\\
0.746867	-0.114161\\
0.741855	-0.114234\\
0.736842	-0.114195\\
0.73183	-0.114078\\
0.726817	-0.113888\\
0.721805	-0.11363\\
0.716792	-0.113309\\
0.711779	-0.112928\\
0.706767	-0.112494\\
0.701754	-0.112011\\
0.696742	-0.111483\\
0.691729	-0.110916\\
0.686717	-0.110314\\
0.681704	-0.109682\\
0.676692	-0.109024\\
0.671679	-0.108345\\
0.666667	-0.107651\\
0.661654	-0.106945\\
0.656642	-0.106233\\
0.651629	-0.105519\\
0.646617	-0.104808\\
0.641604	-0.104104\\
0.636591	-0.103413\\
0.631579	-0.102738\\
0.626566	-0.102084\\
0.621554	-0.101457\\
0.616541	-0.10086\\
0.611529	-0.100298\\
0.606516	-0.0997764\\
0.601504	-0.0992987\\
0.596491	-0.0988698\\
0.591479	-0.0984736\\
0.586466	-0.0980902\\
0.581454	-0.0977182\\
0.576441	-0.0973566\\
0.571429	-0.0970039\\
0.566416	-0.096659\\
0.561404	-0.0963205\\
0.556391	-0.0959872\\
0.551378	-0.0956576\\
0.546366	-0.0953306\\
0.541353	-0.0950046\\
0.536341	-0.0946784\\
0.531328	-0.0943505\\
0.526316	-0.0940198\\
0.521303	-0.0936845\\
0.516291	-0.0933436\\
0.511278	-0.0929955\\
0.506266	-0.0926389\\
0.501253	-0.0922722\\
0.496241	-0.0918943\\
0.491228	-0.0915036\\
0.486216	-0.0910989\\
0.481203	-0.0906785\\
0.47619	-0.0902415\\
0.471178	-0.0897861\\
0.466165	-0.0893114\\
0.461153	-0.0888157\\
0.45614	-0.088298\\
0.451128	-0.087768\\
0.446115	-0.0872325\\
0.441103	-0.0866842\\
0.43609	-0.0861174\\
0.431078	-0.0855245\\
0.426065	-0.0848998\\
0.421053	-0.0842359\\
0.41604	-0.0835272\\
0.411028	-0.0827663\\
0.406015	-0.0819475\\
0.401003	-0.081055\\
0.39599	-0.0800809\\
0.390977	-0.0790278\\
0.385965	-0.0779001\\
0.380952	-0.0767004\\
0.37594	-0.0754329\\
0.370927	-0.0741003\\
0.365915	-0.0727069\\
0.360902	-0.0712554\\
0.35589	-0.06975\\
0.350877	-0.0681934\\
0.345865	-0.0665898\\
0.340852	-0.0649419\\
0.33584	-0.0632537\\
0.330827	-0.0615279\\
0.325815	-0.0597686\\
0.320802	-0.0579783\\
0.315789	-0.0561608\\
0.310777	-0.0543188\\
0.305764	-0.0524559\\
0.300752	-0.0505745\\
0.295739	-0.0486782\\
0.290727	-0.0467692\\
0.285714	-0.0448511\\
0.280702	-0.0429258\\
0.275689	-0.0409966\\
0.270677	-0.0390655\\
0.265664	-0.0371355\\
0.260652	-0.0352083\\
0.255639	-0.0332697\\
0.250627	-0.031302\\
0.245614	-0.0293042\\
0.240602	-0.0272744\\
0.235589	-0.0252116\\
0.230576	-0.0231136\\
0.225564	-0.0209791\\
0.220551	-0.018806\\
0.215539	-0.0165925\\
0.210526	-0.0143366\\
0.205514	-0.0120362\\
0.200501	-0.00968923\\
0.195489	-0.00729344\\
0.190476	-0.00484662\\
0.185464	-0.00234632\\
0.180451	0.000209603\\
0.175439	0.00282393\\
0.170426	0.00549865\\
0.165414	0.00823689\\
0.160401	0.0110404\\
0.155388	0.0139128\\
0.150376	0.0168555\\
0.145363	0.0198724\\
0.140351	0.0229647\\
0.135338	0.0261369\\
0.130326	0.0293896\\
0.125313	0.0327278\\
0.120301	0.0361519\\
0.115288	0.0396672\\
0.110276	0.0432738\\
0.105263	0.0469774\\
0.100251	0.0507777\\
0.0952381	0.0546811\\
0.0902256	0.0586866\\
0.085213	0.0628013\\
0.0802005	0.0670238\\
0.075188	0.0713618\\
0.0701754	0.0758132\\
0.0651629	0.0803865\\
0.0601504	0.0826796\\
0.0551378	0.0834564\\
0.0501253	0.0827607\\
0.0451128	0.0804848\\
0.0401003	0.0767263\\
0.0350877	0.0714194\\
0.0300752	0.0647112\\
0.0250627	0.0565828\\
0.0200501	0.0472206\\
0.0150376	0.0366567\\
0.0100251	0.02803\\
0.00501253	0.0309896\\
0	0.01256\\
}--cycle;

\addplot[area legend, draw=none, fill=mycolor11, fill opacity=0.9, forget plot]
table[row sep=crcr] {%
x	y\\
0	-0.0109093\\
0.00501253	-0.032104\\
0.0100251	-0.0429347\\
0.0150376	-0.0625122\\
0.0200501	-0.0806733\\
0.0250627	-0.0972\\
0.0300752	-0.112126\\
0.0350877	-0.125283\\
0.0401003	-0.136758\\
0.0451128	-0.146434\\
0.0501253	-0.154439\\
0.0551378	-0.160712\\
0.0601504	-0.165412\\
0.0651629	-0.16667\\
0.0701754	-0.167824\\
0.075188	-0.168872\\
0.0802005	-0.169822\\
0.085213	-0.170671\\
0.0902256	-0.171427\\
0.0952381	-0.172089\\
0.100251	-0.172664\\
0.105263	-0.173151\\
0.110276	-0.173555\\
0.115288	-0.173877\\
0.120301	-0.174123\\
0.125313	-0.174292\\
0.130326	-0.17439\\
0.135338	-0.174416\\
0.140351	-0.174376\\
0.145363	-0.174271\\
0.150376	-0.174104\\
0.155388	-0.173878\\
0.160401	-0.173595\\
0.165414	-0.173257\\
0.170426	-0.172868\\
0.175439	-0.17243\\
0.180451	-0.171944\\
0.185464	-0.171414\\
0.190476	-0.170843\\
0.195489	-0.170231\\
0.200501	-0.169582\\
0.205514	-0.168897\\
0.210526	-0.168179\\
0.215539	-0.16743\\
0.220551	-0.166652\\
0.225564	-0.165847\\
0.230576	-0.165016\\
0.235589	-0.164162\\
0.240602	-0.163286\\
0.245614	-0.16239\\
0.250627	-0.161474\\
0.255639	-0.160542\\
0.260652	-0.159592\\
0.265664	-0.158563\\
0.270677	-0.157403\\
0.275689	-0.156135\\
0.280702	-0.154781\\
0.285714	-0.153364\\
0.290727	-0.151904\\
0.295739	-0.150424\\
0.300752	-0.148946\\
0.305764	-0.147492\\
0.310777	-0.146082\\
0.315789	-0.14474\\
0.320802	-0.143485\\
0.325815	-0.142341\\
0.330827	-0.141328\\
0.33584	-0.140467\\
0.340852	-0.13978\\
0.345865	-0.139288\\
0.350877	-0.139013\\
0.35589	-0.138975\\
0.360902	-0.139196\\
0.365915	-0.13957\\
0.370927	-0.139972\\
0.37594	-0.140392\\
0.380952	-0.140819\\
0.385965	-0.141245\\
0.390977	-0.141659\\
0.39599	-0.142051\\
0.401003	-0.142412\\
0.406015	-0.142731\\
0.411028	-0.142999\\
0.41604	-0.143199\\
0.421053	-0.143328\\
0.426065	-0.143394\\
0.431078	-0.143408\\
0.43609	-0.143378\\
0.441103	-0.143314\\
0.446115	-0.143226\\
0.451128	-0.143123\\
0.45614	-0.143014\\
0.461153	-0.142908\\
0.466165	-0.142794\\
0.471178	-0.142652\\
0.47619	-0.142484\\
0.481203	-0.142292\\
0.486216	-0.142078\\
0.491228	-0.141843\\
0.496241	-0.14159\\
0.501253	-0.141319\\
0.506266	-0.141033\\
0.511278	-0.140734\\
0.516291	-0.140422\\
0.521303	-0.140101\\
0.526316	-0.139772\\
0.531328	-0.139437\\
0.536341	-0.139097\\
0.541353	-0.138754\\
0.546366	-0.138411\\
0.551378	-0.138068\\
0.556391	-0.137727\\
0.561404	-0.137391\\
0.566416	-0.137061\\
0.571429	-0.136738\\
0.576441	-0.136425\\
0.581454	-0.136123\\
0.586466	-0.135834\\
0.591479	-0.135559\\
0.596491	-0.1353\\
0.601504	-0.135058\\
0.606516	-0.134836\\
0.611529	-0.134626\\
0.616541	-0.134421\\
0.621554	-0.134221\\
0.626566	-0.134026\\
0.631579	-0.133836\\
0.636591	-0.13365\\
0.641604	-0.133469\\
0.646617	-0.133293\\
0.651629	-0.133121\\
0.656642	-0.132953\\
0.661654	-0.132789\\
0.666667	-0.132629\\
0.671679	-0.132473\\
0.676692	-0.132321\\
0.681704	-0.132172\\
0.686717	-0.132026\\
0.691729	-0.131884\\
0.696742	-0.131744\\
0.701754	-0.131606\\
0.706767	-0.131472\\
0.711779	-0.131339\\
0.716792	-0.131208\\
0.721805	-0.131079\\
0.726817	-0.130951\\
0.73183	-0.130824\\
0.736842	-0.130698\\
0.741855	-0.130573\\
0.746867	-0.130449\\
0.75188	-0.130325\\
0.756892	-0.1302\\
0.761905	-0.130076\\
0.766917	-0.12995\\
0.77193	-0.129824\\
0.776942	-0.129697\\
0.781955	-0.129569\\
0.786967	-0.12944\\
0.79198	-0.129308\\
0.796992	-0.129175\\
0.802005	-0.12904\\
0.807018	-0.128898\\
0.81203	-0.128746\\
0.817043	-0.128586\\
0.822055	-0.12842\\
0.827068	-0.128249\\
0.83208	-0.128075\\
0.837093	-0.127901\\
0.842105	-0.127727\\
0.847118	-0.127556\\
0.85213	-0.127389\\
0.857143	-0.127229\\
0.862155	-0.127076\\
0.867168	-0.126934\\
0.87218	-0.126803\\
0.877193	-0.126686\\
0.882206	-0.126585\\
0.887218	-0.1265\\
0.892231	-0.126435\\
0.897243	-0.126391\\
0.902256	-0.126369\\
0.907268	-0.126375\\
0.912281	-0.126407\\
0.917293	-0.126458\\
0.922306	-0.126523\\
0.927318	-0.126595\\
0.932331	-0.126667\\
0.937343	-0.126732\\
0.942356	-0.126784\\
0.947368	-0.126817\\
0.952381	-0.126824\\
0.957393	-0.126809\\
0.962406	-0.126783\\
0.967419	-0.126748\\
0.972431	-0.126707\\
0.977444	-0.126661\\
0.982456	-0.126612\\
0.987469	-0.126564\\
0.992481	-0.126518\\
0.997494	-0.126476\\
1.00251	-0.12644\\
1.00752	-0.126409\\
1.01253	-0.126378\\
1.01754	-0.126348\\
1.02256	-0.126319\\
1.02757	-0.126291\\
1.03258	-0.126263\\
1.03759	-0.126236\\
1.04261	-0.126209\\
1.04762	-0.126183\\
1.05263	-0.126158\\
1.05764	-0.126133\\
1.06266	-0.126108\\
1.06767	-0.126084\\
1.07268	-0.12606\\
1.07769	-0.126037\\
1.08271	-0.126014\\
1.08772	-0.125991\\
1.09273	-0.125968\\
1.09774	-0.125946\\
1.10276	-0.125924\\
1.10777	-0.125902\\
1.11278	-0.125881\\
1.11779	-0.125859\\
1.12281	-0.125838\\
1.12782	-0.125816\\
1.13283	-0.125795\\
1.13784	-0.125773\\
1.14286	-0.125752\\
1.14787	-0.125731\\
1.15288	-0.125709\\
1.15789	-0.125687\\
1.16291	-0.125663\\
1.16792	-0.12564\\
1.17293	-0.125615\\
1.17794	-0.12559\\
1.18296	-0.125565\\
1.18797	-0.12554\\
1.19298	-0.125514\\
1.19799	-0.125489\\
1.20301	-0.125463\\
1.20802	-0.125438\\
1.21303	-0.125413\\
1.21805	-0.125388\\
1.22306	-0.125363\\
1.22807	-0.125339\\
1.23308	-0.125316\\
1.2381	-0.125293\\
1.24311	-0.125271\\
1.24812	-0.12525\\
1.25313	-0.12523\\
1.25815	-0.125211\\
1.26316	-0.125194\\
1.26817	-0.125177\\
1.27318	-0.125162\\
1.2782	-0.125148\\
1.28321	-0.125136\\
1.28822	-0.125126\\
1.29323	-0.125117\\
1.29825	-0.12511\\
1.30326	-0.125107\\
1.30827	-0.125108\\
1.31328	-0.125112\\
1.3183	-0.125118\\
1.32331	-0.125126\\
1.32832	-0.125133\\
1.33333	-0.125139\\
1.33835	-0.125144\\
1.34336	-0.125145\\
1.34837	-0.125141\\
1.35338	-0.125135\\
1.3584	-0.125129\\
1.36341	-0.125122\\
1.36842	-0.125114\\
1.37343	-0.125106\\
1.37845	-0.125097\\
1.38346	-0.125088\\
1.38847	-0.125078\\
1.39348	-0.125068\\
1.3985	-0.125058\\
1.40351	-0.125048\\
1.40852	-0.125037\\
1.41353	-0.125027\\
1.41855	-0.125016\\
1.42356	-0.125005\\
1.42857	-0.124994\\
1.43358	-0.124983\\
1.4386	-0.124972\\
1.44361	-0.124961\\
1.44862	-0.12495\\
1.45363	-0.124939\\
1.45865	-0.124929\\
1.46366	-0.124918\\
1.46867	-0.124908\\
1.47368	-0.124898\\
1.4787	-0.124889\\
1.48371	-0.12488\\
1.48872	-0.124871\\
1.49373	-0.124863\\
1.49875	-0.124855\\
1.50376	-0.124848\\
1.50877	-0.124841\\
1.51378	-0.124835\\
1.5188	-0.12483\\
1.52381	-0.124824\\
1.52882	-0.124818\\
1.53383	-0.124812\\
1.53885	-0.124807\\
1.54386	-0.124801\\
1.54887	-0.124795\\
1.55388	-0.12479\\
1.5589	-0.124784\\
1.56391	-0.124778\\
1.56892	-0.124773\\
1.57393	-0.124767\\
1.57895	-0.124761\\
1.58396	-0.124756\\
1.58897	-0.12475\\
1.59398	-0.124745\\
1.599	-0.124739\\
1.60401	-0.124733\\
1.60902	-0.124728\\
1.61404	-0.124722\\
1.61905	-0.124717\\
1.62406	-0.124711\\
1.62907	-0.124706\\
1.63409	-0.1247\\
1.6391	-0.124695\\
1.64411	-0.124689\\
1.64912	-0.124684\\
1.65414	-0.124678\\
1.65915	-0.124673\\
1.66416	-0.124667\\
1.66917	-0.124662\\
1.67419	-0.124657\\
1.6792	-0.124651\\
1.68421	-0.124646\\
1.68922	-0.124641\\
1.69424	-0.124635\\
1.69925	-0.12463\\
1.70426	-0.124625\\
1.70927	-0.12462\\
1.71429	-0.124614\\
1.7193	-0.124609\\
1.72431	-0.124604\\
1.72932	-0.124599\\
1.73434	-0.124594\\
1.73935	-0.124589\\
1.74436	-0.124584\\
1.74937	-0.124578\\
1.75439	-0.124573\\
1.7594	-0.124568\\
1.76441	-0.124563\\
1.76942	-0.124558\\
1.77444	-0.124553\\
1.77945	-0.124548\\
1.78446	-0.124543\\
1.78947	-0.124538\\
1.79449	-0.124533\\
1.7995	-0.124528\\
1.80451	-0.124523\\
1.80952	-0.124518\\
1.81454	-0.124513\\
1.81955	-0.124508\\
1.82456	-0.124503\\
1.82957	-0.124498\\
1.83459	-0.124493\\
1.8396	-0.124488\\
1.84461	-0.124483\\
1.84962	-0.124478\\
1.85464	-0.124473\\
1.85965	-0.124467\\
1.86466	-0.124462\\
1.86967	-0.124457\\
1.87469	-0.124452\\
1.8797	-0.124447\\
1.88471	-0.124442\\
1.88972	-0.124437\\
1.89474	-0.124432\\
1.89975	-0.124427\\
1.90476	-0.124422\\
1.90977	-0.124417\\
1.91479	-0.124412\\
1.9198	-0.124407\\
1.92481	-0.124402\\
1.92982	-0.124397\\
1.93484	-0.124392\\
1.93985	-0.124387\\
1.94486	-0.124382\\
1.94987	-0.124377\\
1.95489	-0.124372\\
1.9599	-0.124367\\
1.96491	-0.124362\\
1.96992	-0.124357\\
1.97494	-0.124352\\
1.97995	-0.124347\\
1.98496	-0.124342\\
1.98997	-0.124337\\
1.99499	-0.124332\\
2	-0.124327\\
2	-0.124049\\
1.99499	-0.124053\\
1.98997	-0.124057\\
1.98496	-0.124062\\
1.97995	-0.124066\\
1.97494	-0.12407\\
1.96992	-0.124074\\
1.96491	-0.124079\\
1.9599	-0.124083\\
1.95489	-0.124087\\
1.94987	-0.124092\\
1.94486	-0.124096\\
1.93985	-0.1241\\
1.93484	-0.124105\\
1.92982	-0.124109\\
1.92481	-0.124113\\
1.9198	-0.124118\\
1.91479	-0.124122\\
1.90977	-0.124126\\
1.90476	-0.124131\\
1.89975	-0.124135\\
1.89474	-0.124139\\
1.88972	-0.124144\\
1.88471	-0.124148\\
1.8797	-0.124152\\
1.87469	-0.124157\\
1.86967	-0.124161\\
1.86466	-0.124165\\
1.85965	-0.12417\\
1.85464	-0.124174\\
1.84962	-0.124178\\
1.84461	-0.124183\\
1.8396	-0.124187\\
1.83459	-0.124192\\
1.82957	-0.124196\\
1.82456	-0.124201\\
1.81955	-0.124205\\
1.81454	-0.12421\\
1.80952	-0.124214\\
1.80451	-0.124219\\
1.7995	-0.124223\\
1.79449	-0.124228\\
1.78947	-0.124232\\
1.78446	-0.124237\\
1.77945	-0.124241\\
1.77444	-0.124246\\
1.76942	-0.12425\\
1.76441	-0.124255\\
1.7594	-0.12426\\
1.75439	-0.124264\\
1.74937	-0.124269\\
1.74436	-0.124274\\
1.73935	-0.124278\\
1.73434	-0.124283\\
1.72932	-0.124288\\
1.72431	-0.124292\\
1.7193	-0.124297\\
1.71429	-0.124302\\
1.70927	-0.124306\\
1.70426	-0.124311\\
1.69925	-0.124316\\
1.69424	-0.124321\\
1.68922	-0.124326\\
1.68421	-0.12433\\
1.6792	-0.124335\\
1.67419	-0.12434\\
1.66917	-0.124345\\
1.66416	-0.12435\\
1.65915	-0.124354\\
1.65414	-0.124359\\
1.64912	-0.124364\\
1.64411	-0.124369\\
1.6391	-0.124374\\
1.63409	-0.124379\\
1.62907	-0.124384\\
1.62406	-0.124388\\
1.61905	-0.124393\\
1.61404	-0.124398\\
1.60902	-0.124403\\
1.60401	-0.124408\\
1.599	-0.124412\\
1.59398	-0.124417\\
1.58897	-0.124422\\
1.58396	-0.124427\\
1.57895	-0.124431\\
1.57393	-0.124436\\
1.56892	-0.124441\\
1.56391	-0.124445\\
1.5589	-0.12445\\
1.55388	-0.124455\\
1.54887	-0.124459\\
1.54386	-0.124464\\
1.53885	-0.124468\\
1.53383	-0.124473\\
1.52882	-0.124477\\
1.52381	-0.124481\\
1.5188	-0.124486\\
1.51378	-0.12449\\
1.50877	-0.124494\\
1.50376	-0.124497\\
1.49875	-0.1245\\
1.49373	-0.124502\\
1.48872	-0.124504\\
1.48371	-0.124505\\
1.4787	-0.124506\\
1.47368	-0.124507\\
1.46867	-0.124507\\
1.46366	-0.124507\\
1.45865	-0.124506\\
1.45363	-0.124506\\
1.44862	-0.124505\\
1.44361	-0.124504\\
1.4386	-0.124503\\
1.43358	-0.124502\\
1.42857	-0.124501\\
1.42356	-0.1245\\
1.41855	-0.1245\\
1.41353	-0.124499\\
1.40852	-0.124499\\
1.40351	-0.124499\\
1.3985	-0.124499\\
1.39348	-0.124499\\
1.38847	-0.1245\\
1.38346	-0.124501\\
1.37845	-0.124503\\
1.37343	-0.124505\\
1.36842	-0.124508\\
1.36341	-0.124511\\
1.3584	-0.124515\\
1.35338	-0.12452\\
1.34837	-0.124525\\
1.34336	-0.124533\\
1.33835	-0.124546\\
1.33333	-0.124561\\
1.32832	-0.124579\\
1.32331	-0.124598\\
1.3183	-0.124617\\
1.31328	-0.124635\\
1.30827	-0.124651\\
1.30326	-0.124664\\
1.29825	-0.124672\\
1.29323	-0.124677\\
1.28822	-0.12468\\
1.28321	-0.124681\\
1.2782	-0.124681\\
1.27318	-0.124678\\
1.26817	-0.124675\\
1.26316	-0.12467\\
1.25815	-0.124663\\
1.25313	-0.124655\\
1.24812	-0.124646\\
1.24311	-0.124636\\
1.2381	-0.124625\\
1.23308	-0.124613\\
1.22807	-0.1246\\
1.22306	-0.124586\\
1.21805	-0.124572\\
1.21303	-0.124556\\
1.20802	-0.124541\\
1.20301	-0.124525\\
1.19799	-0.124508\\
1.19298	-0.124491\\
1.18797	-0.124474\\
1.18296	-0.124457\\
1.17794	-0.12444\\
1.17293	-0.124423\\
1.16792	-0.124406\\
1.16291	-0.12439\\
1.15789	-0.124373\\
1.15288	-0.124358\\
1.14787	-0.124342\\
1.14286	-0.124327\\
1.13784	-0.124312\\
1.13283	-0.124296\\
1.12782	-0.12428\\
1.12281	-0.124264\\
1.11779	-0.124248\\
1.11278	-0.124232\\
1.10777	-0.124215\\
1.10276	-0.124198\\
1.09774	-0.124181\\
1.09273	-0.124163\\
1.08772	-0.124146\\
1.08271	-0.124128\\
1.07769	-0.12411\\
1.07268	-0.124091\\
1.06767	-0.124073\\
1.06266	-0.124054\\
1.05764	-0.124035\\
1.05263	-0.124016\\
1.04762	-0.123997\\
1.04261	-0.123978\\
1.03759	-0.123958\\
1.03258	-0.123939\\
1.02757	-0.123919\\
1.02256	-0.123899\\
1.01754	-0.123879\\
1.01253	-0.123858\\
1.00752	-0.123838\\
1.00251	-0.123817\\
0.997494	-0.123793\\
0.992481	-0.123763\\
0.987469	-0.123729\\
0.982456	-0.123694\\
0.977444	-0.12366\\
0.972431	-0.123629\\
0.967419	-0.123602\\
0.962406	-0.123583\\
0.957393	-0.123573\\
0.952381	-0.123576\\
0.947368	-0.1236\\
0.942356	-0.12365\\
0.937343	-0.123721\\
0.932331	-0.123805\\
0.927318	-0.123895\\
0.922306	-0.123986\\
0.917293	-0.12407\\
0.912281	-0.124142\\
0.907268	-0.124193\\
0.902256	-0.124218\\
0.897243	-0.124216\\
0.892231	-0.12419\\
0.887218	-0.124144\\
0.882206	-0.124078\\
0.877193	-0.123995\\
0.87218	-0.123895\\
0.867168	-0.123782\\
0.862155	-0.123656\\
0.857143	-0.12352\\
0.85213	-0.123375\\
0.847118	-0.123222\\
0.842105	-0.123065\\
0.837093	-0.122904\\
0.83208	-0.122741\\
0.827068	-0.122579\\
0.822055	-0.122418\\
0.817043	-0.122261\\
0.81203	-0.122109\\
0.807018	-0.121965\\
0.802005	-0.12183\\
0.796992	-0.121701\\
0.79198	-0.121573\\
0.786967	-0.121446\\
0.781955	-0.121321\\
0.776942	-0.121196\\
0.77193	-0.121073\\
0.766917	-0.120949\\
0.761905	-0.120827\\
0.756892	-0.120705\\
0.75188	-0.120583\\
0.746867	-0.120461\\
0.741855	-0.120339\\
0.736842	-0.120217\\
0.73183	-0.120094\\
0.726817	-0.119972\\
0.721805	-0.119849\\
0.716792	-0.119725\\
0.711779	-0.1196\\
0.706767	-0.119474\\
0.701754	-0.119348\\
0.696742	-0.11922\\
0.691729	-0.119091\\
0.686717	-0.11896\\
0.681704	-0.118828\\
0.676692	-0.118694\\
0.671679	-0.118558\\
0.666667	-0.118419\\
0.661654	-0.118278\\
0.656642	-0.118135\\
0.651629	-0.117989\\
0.646617	-0.11784\\
0.641604	-0.117687\\
0.636591	-0.117532\\
0.631579	-0.117372\\
0.626566	-0.117208\\
0.621554	-0.11704\\
0.616541	-0.116867\\
0.611529	-0.116689\\
0.606516	-0.116506\\
0.601504	-0.11631\\
0.596491	-0.116093\\
0.591479	-0.115857\\
0.586466	-0.115604\\
0.581454	-0.115333\\
0.576441	-0.115047\\
0.571429	-0.114746\\
0.566416	-0.114432\\
0.561404	-0.114105\\
0.556391	-0.113767\\
0.551378	-0.113419\\
0.546366	-0.113062\\
0.541353	-0.112697\\
0.536341	-0.112324\\
0.531328	-0.111946\\
0.526316	-0.111562\\
0.521303	-0.111174\\
0.516291	-0.110783\\
0.511278	-0.11039\\
0.506266	-0.109996\\
0.501253	-0.109601\\
0.496241	-0.109208\\
0.491228	-0.108816\\
0.486216	-0.108427\\
0.481203	-0.108042\\
0.47619	-0.107662\\
0.471178	-0.107288\\
0.466165	-0.10692\\
0.461153	-0.106561\\
0.45614	-0.106191\\
0.451128	-0.105796\\
0.446115	-0.105385\\
0.441103	-0.104968\\
0.43609	-0.104554\\
0.431078	-0.104151\\
0.426065	-0.103768\\
0.421053	-0.103415\\
0.41604	-0.103101\\
0.411028	-0.102835\\
0.406015	-0.102613\\
0.401003	-0.102419\\
0.39599	-0.102242\\
0.390977	-0.102073\\
0.385965	-0.101902\\
0.380952	-0.101719\\
0.37594	-0.101514\\
0.370927	-0.101278\\
0.365915	-0.101001\\
0.360902	-0.100672\\
0.35589	-0.100167\\
0.350877	-0.0993799\\
0.345865	-0.0983321\\
0.340852	-0.0970444\\
0.33584	-0.0955385\\
0.330827	-0.093835\\
0.325815	-0.0919555\\
0.320802	-0.0899205\\
0.315789	-0.0877514\\
0.310777	-0.0854687\\
0.305764	-0.0830938\\
0.300752	-0.0806468\\
0.295739	-0.0781491\\
0.290727	-0.0756207\\
0.285714	-0.0730827\\
0.280702	-0.0705552\\
0.275689	-0.0680589\\
0.270677	-0.0656136\\
0.265664	-0.0632402\\
0.260652	-0.0609581\\
0.255639	-0.0587154\\
0.250627	-0.056445\\
0.245614	-0.0541458\\
0.240602	-0.0518159\\
0.235589	-0.0494541\\
0.230576	-0.0470584\\
0.225564	-0.0446272\\
0.220551	-0.0421583\\
0.215539	-0.0396501\\
0.210526	-0.0371003\\
0.205514	-0.0345069\\
0.200501	-0.0318677\\
0.195489	-0.0291803\\
0.190476	-0.0264426\\
0.185464	-0.023652\\
0.180451	-0.0208063\\
0.175439	-0.0179026\\
0.170426	-0.0149389\\
0.165414	-0.0119121\\
0.160401	-0.00882021\\
0.155388	-0.00565976\\
0.150376	-0.00242912\\
0.145363	0.000875648\\
0.140351	0.00425584\\
0.135338	0.00771582\\
0.130326	0.0112565\\
0.125313	0.0148828\\
0.120301	0.0185953\\
0.115288	0.0223991\\
0.110276	0.0262946\\
0.105263	0.0302875\\
0.100251	0.0343776\\
0.0952381	0.0385713\\
0.0902256	0.0428676\\
0.085213	0.0472739\\
0.0802005	0.0517886\\
0.075188	0.0564196\\
0.0701754	0.0611648\\
0.0651629	0.0660328\\
0.0601504	0.0710206\\
0.0551378	0.0728092\\
0.0501253	0.0732581\\
0.0451128	0.0722283\\
0.0401003	0.0697712\\
0.0350877	0.0657782\\
0.0300752	0.0603506\\
0.0250627	0.053419\\
0.0200501	0.0451311\\
0.0150376	0.0354628\\
0.0100251	0.0246048\\
0.00501253	0.0227473\\
0	0.0109093\\
}--cycle;
\addplot [color=mycolor12, line width=1.0pt, forget plot]
  table[row sep=crcr]{%
-0.124	0\\
-0.122	0\\
-0.12	0\\
-0.118	0\\
-0.116	0\\
-0.114	0\\
-0.112	0\\
-0.11	0\\
-0.108	0\\
-0.106	0\\
-0.104	0\\
-0.102	0\\
-0.1	0\\
-0.098	0\\
-0.096	0\\
-0.094	0\\
-0.092	0\\
-0.09	0\\
-0.088	0\\
-0.086	0\\
-0.084	0\\
-0.082	0\\
-0.08	0\\
-0.078	0\\
-0.076	0\\
-0.074	0\\
-0.072	0\\
-0.07	0\\
-0.068	0\\
-0.066	0\\
-0.064	0\\
-0.062	0\\
-0.06	0\\
-0.058	0\\
-0.056	0\\
-0.054	0\\
-0.052	0\\
-0.05	0\\
-0.048	0\\
-0.046	0\\
-0.044	0\\
-0.042	0\\
-0.04	0\\
-0.038	0\\
-0.036	0\\
-0.034	0\\
-0.032	0\\
-0.03	0\\
-0.028	0\\
-0.026	0\\
-0.024	0\\
-0.022	0\\
-0.02	0\\
-0.018	0\\
-0.016	0\\
-0.014	0\\
-0.012	0\\
-0.01	0\\
-0.008	0\\
-0.006	0\\
-0.004	0\\
-0.002	0\\
0	0\\
0.002	-0.000283243\\
0.004	-0.000639969\\
0.006	-0.00105215\\
0.008	-0.00150939\\
0.01	-0.00200533\\
0.012	-0.0025359\\
0.014	-0.00309834\\
0.016	-0.00369071\\
0.018	-0.0043116\\
0.02	-0.00495994\\
0.022	-0.00563493\\
0.024	-0.00633593\\
0.026	-0.00706246\\
0.028	-0.0078141\\
0.03	-0.00859052\\
0.032	-0.00939143\\
0.034	-0.0102166\\
0.036	-0.0110657\\
0.038	-0.0119385\\
0.04	-0.0128347\\
0.042	-0.0137541\\
0.044	-0.0146965\\
0.046	-0.0156613\\
0.048	-0.0166484\\
0.05	-0.0176573\\
0.052	-0.0186876\\
0.054	-0.0197389\\
0.056	-0.0208107\\
0.058	-0.0219024\\
0.06	-0.0230136\\
0.062	-0.0241436\\
0.064	-0.0252918\\
0.066	-0.0264576\\
0.068	-0.0276402\\
0.07	-0.0288391\\
0.072	-0.0300535\\
0.074	-0.0312825\\
0.076	-0.0325255\\
0.078	-0.0337817\\
0.08	-0.0350501\\
0.082	-0.03633\\
0.084	-0.0376206\\
0.086	-0.0389209\\
0.088	-0.0402301\\
0.09	-0.0415472\\
0.092	-0.0428715\\
0.094	-0.0442019\\
0.096	-0.0455377\\
0.098	-0.0468777\\
0.1	-0.0482213\\
0.102	-0.0495673\\
0.104	-0.050915\\
0.106	-0.0522634\\
0.108	-0.0536115\\
0.11	-0.0549586\\
0.112	-0.0563036\\
0.114	-0.0576458\\
0.116	-0.0589842\\
0.118	-0.0603179\\
0.12	-0.0616462\\
0.122	-0.0629681\\
0.124	-0.0642829\\
0.126	-0.0655897\\
0.128	-0.0668878\\
0.13	-0.0681763\\
0.132	-0.0694545\\
0.134	-0.0707216\\
0.136	-0.071977\\
0.138	-0.0732199\\
0.14	-0.0744496\\
0.142	-0.0756656\\
0.144	-0.0768671\\
0.146	-0.0780535\\
0.148	-0.0792242\\
0.15	-0.0803787\\
0.152	-0.0815164\\
0.154	-0.0826368\\
0.156	-0.0837395\\
0.158	-0.0848238\\
0.16	-0.0858894\\
0.162	-0.0869359\\
0.164	-0.0879628\\
0.166	-0.0889698\\
0.168	-0.0899565\\
0.17	-0.0909226\\
0.172	-0.0918678\\
0.174	-0.0927919\\
0.176	-0.0936946\\
0.178	-0.0945756\\
0.18	-0.0954348\\
0.182	-0.096272\\
0.184	-0.097087\\
0.186	-0.0978797\\
0.188	-0.0986501\\
0.19	-0.099398\\
0.192	-0.100123\\
0.194	-0.100826\\
0.196	-0.101507\\
0.198	-0.102164\\
0.2	-0.1028\\
0.202	-0.103413\\
0.204	-0.104003\\
0.206	-0.104571\\
0.208	-0.105118\\
0.21	-0.105642\\
0.212	-0.106144\\
0.214	-0.106625\\
0.216	-0.107084\\
0.218	-0.107522\\
0.22	-0.107939\\
0.222	-0.108335\\
0.224	-0.108711\\
0.226	-0.109067\\
0.228	-0.109402\\
0.23	-0.109718\\
0.232	-0.110015\\
0.234	-0.110293\\
0.236	-0.110552\\
0.238	-0.110793\\
0.24	-0.111016\\
0.242	-0.111221\\
0.244	-0.11141\\
0.246	-0.111582\\
0.248	-0.111737\\
0.25	-0.111877\\
0.252	-0.112001\\
0.254	-0.11211\\
0.256	-0.112205\\
0.258	-0.112285\\
0.26	-0.112352\\
0.262	-0.112405\\
0.264	-0.112446\\
0.266	-0.112474\\
0.268	-0.112491\\
0.27	-0.112496\\
0.272	-0.11249\\
0.274	-0.112474\\
0.276	-0.112448\\
0.278	-0.112412\\
0.28	-0.112367\\
0.282	-0.112314\\
0.284	-0.112252\\
0.286	-0.112183\\
0.288	-0.112107\\
0.29	-0.112024\\
0.292	-0.111934\\
0.294	-0.111839\\
0.296	-0.111739\\
0.298	-0.111633\\
0.3	-0.111523\\
0.302	-0.111409\\
0.304	-0.111291\\
0.306	-0.11117\\
0.308	-0.111047\\
0.31	-0.110921\\
0.312	-0.110793\\
0.314	-0.110663\\
0.316	-0.110532\\
0.318	-0.1104\\
0.32	-0.110268\\
0.322	-0.110136\\
0.324	-0.110004\\
0.326	-0.109873\\
0.328	-0.109743\\
0.33	-0.109614\\
0.332	-0.109487\\
0.334	-0.109361\\
0.336	-0.109238\\
0.338	-0.109118\\
0.34	-0.109\\
0.342	-0.108885\\
0.344	-0.108774\\
0.346	-0.108666\\
0.348	-0.108563\\
0.35	-0.108463\\
0.352	-0.108368\\
0.354	-0.108277\\
0.356	-0.108191\\
0.358	-0.10811\\
0.36	-0.108034\\
0.362	-0.107964\\
0.364	-0.107899\\
0.366	-0.10784\\
0.368	-0.107786\\
0.37	-0.107739\\
0.372	-0.107698\\
0.374	-0.107663\\
0.376	-0.107634\\
0.378	-0.107612\\
0.38	-0.107596\\
0.382	-0.107587\\
0.384	-0.107584\\
0.386	-0.107588\\
0.388	-0.1076\\
0.39	-0.107617\\
0.392	-0.107642\\
0.394	-0.107674\\
0.396	-0.107713\\
0.398	-0.107758\\
0.4	-0.107811\\
0.402	-0.107871\\
0.404	-0.107937\\
0.406	-0.108011\\
0.408	-0.108091\\
0.41	-0.108178\\
0.412	-0.108273\\
0.414	-0.108373\\
0.416	-0.108481\\
0.418	-0.108595\\
0.42	-0.108716\\
0.422	-0.108843\\
0.424	-0.108977\\
0.426	-0.109117\\
0.428	-0.109263\\
0.43	-0.109415\\
0.432	-0.109574\\
0.434	-0.109738\\
0.436	-0.109908\\
0.438	-0.110083\\
0.44	-0.110264\\
0.442	-0.110451\\
0.444	-0.110642\\
0.446	-0.110839\\
0.448	-0.111041\\
0.45	-0.111247\\
0.452	-0.111458\\
0.454	-0.111673\\
0.456	-0.111893\\
0.458	-0.112117\\
0.46	-0.112344\\
0.462	-0.112576\\
0.464	-0.11281\\
0.466	-0.113049\\
0.468	-0.11329\\
0.47	-0.113534\\
0.472	-0.113781\\
0.474	-0.114031\\
0.476	-0.114283\\
0.478	-0.114538\\
0.48	-0.114794\\
0.482	-0.115052\\
0.484	-0.115312\\
0.486	-0.115573\\
0.488	-0.115835\\
0.49	-0.116098\\
0.492	-0.116362\\
0.494	-0.116627\\
0.496	-0.116892\\
0.498	-0.117157\\
0.5	-0.117423\\
0.502	-0.117688\\
0.504	-0.117952\\
0.506	-0.118216\\
0.508	-0.118479\\
0.51	-0.118742\\
0.512	-0.119003\\
0.514	-0.119262\\
0.516	-0.11952\\
0.518	-0.119777\\
0.52	-0.120031\\
0.522	-0.120283\\
0.524	-0.120534\\
0.526	-0.120781\\
0.528	-0.121026\\
0.53	-0.121269\\
0.532	-0.121508\\
0.534	-0.121744\\
0.536	-0.121978\\
0.538	-0.122208\\
0.54	-0.122434\\
0.542	-0.122657\\
0.544	-0.122876\\
0.546	-0.123091\\
0.548	-0.123302\\
0.55	-0.123509\\
0.552	-0.123712\\
0.554	-0.12391\\
0.556	-0.124104\\
0.558	-0.124294\\
0.56	-0.124479\\
0.562	-0.124659\\
0.564	-0.124835\\
0.566	-0.125006\\
0.568	-0.125172\\
0.57	-0.125333\\
0.572	-0.125489\\
0.574	-0.12564\\
0.576	-0.125786\\
0.578	-0.125927\\
0.58	-0.126063\\
0.582	-0.126193\\
0.584	-0.126319\\
0.586	-0.126439\\
0.588	-0.126555\\
0.59	-0.126665\\
0.592	-0.12677\\
0.594	-0.12687\\
0.596	-0.126964\\
0.598	-0.127054\\
0.6	-0.127139\\
0.602	-0.127219\\
0.604	-0.127293\\
0.606	-0.127363\\
0.608	-0.127428\\
0.61	-0.127488\\
0.612	-0.127544\\
0.614	-0.127594\\
0.616	-0.12764\\
0.618	-0.127682\\
0.62	-0.127719\\
0.622	-0.127752\\
0.624	-0.127781\\
0.626	-0.127805\\
0.628	-0.127825\\
0.63	-0.127841\\
0.632	-0.127854\\
0.634	-0.127862\\
0.636	-0.127867\\
0.638	-0.127868\\
0.64	-0.127866\\
0.642	-0.12786\\
0.644	-0.127851\\
0.646	-0.127839\\
0.648	-0.127824\\
0.65	-0.127806\\
0.652	-0.127785\\
0.654	-0.127761\\
0.656	-0.127735\\
0.658	-0.127707\\
0.66	-0.127676\\
0.662	-0.127643\\
0.664	-0.127608\\
0.666	-0.127571\\
0.668	-0.127532\\
0.67	-0.127491\\
0.672	-0.127449\\
0.674	-0.127405\\
0.676	-0.12736\\
0.678	-0.127314\\
0.68	-0.127266\\
0.682	-0.127217\\
0.684	-0.127168\\
0.686	-0.127118\\
0.688	-0.127067\\
0.69	-0.127015\\
0.692	-0.126963\\
0.694	-0.12691\\
0.696	-0.126858\\
0.698	-0.126804\\
0.7	-0.126751\\
0.702	-0.126698\\
0.704	-0.126645\\
0.706	-0.126592\\
0.708	-0.126539\\
0.71	-0.126487\\
0.712	-0.126435\\
0.714	-0.126383\\
0.716	-0.126332\\
0.718	-0.126281\\
0.72	-0.126231\\
0.722	-0.126182\\
0.724	-0.126134\\
0.726	-0.126086\\
0.728	-0.126039\\
0.73	-0.125993\\
0.732	-0.125949\\
0.734	-0.125905\\
0.736	-0.125862\\
0.738	-0.12582\\
0.74	-0.125779\\
0.742	-0.12574\\
0.744	-0.125702\\
0.746	-0.125664\\
0.748	-0.125628\\
0.75	-0.125594\\
0.752	-0.12556\\
0.754	-0.125528\\
0.756	-0.125497\\
0.758	-0.125467\\
0.76	-0.125439\\
0.762	-0.125411\\
0.764	-0.125386\\
0.766	-0.125361\\
0.768	-0.125337\\
0.77	-0.125315\\
0.772	-0.125294\\
0.774	-0.125275\\
0.776	-0.125256\\
0.778	-0.125239\\
0.78	-0.125223\\
0.782	-0.125208\\
0.784	-0.125194\\
0.786	-0.125182\\
0.788	-0.12517\\
0.79	-0.12516\\
0.792	-0.125151\\
0.794	-0.125143\\
0.796	-0.125135\\
0.798	-0.125129\\
0.8	-0.125124\\
0.802	-0.125119\\
0.804	-0.125116\\
0.806	-0.125113\\
0.808	-0.125112\\
0.81	-0.125111\\
0.812	-0.125111\\
0.814	-0.125111\\
0.816	-0.125113\\
0.818	-0.125115\\
0.82	-0.125117\\
0.822	-0.125121\\
0.824	-0.125124\\
0.826	-0.125129\\
0.828	-0.125134\\
0.83	-0.125139\\
0.832	-0.125145\\
0.834	-0.125152\\
0.836	-0.125158\\
0.838	-0.125166\\
0.84	-0.125173\\
0.842	-0.125181\\
0.844	-0.125189\\
0.846	-0.125198\\
0.848	-0.125206\\
0.85	-0.125215\\
0.852	-0.125224\\
0.854	-0.125233\\
0.856	-0.125243\\
0.858	-0.125252\\
0.86	-0.125261\\
0.862	-0.125271\\
0.864	-0.125281\\
0.866	-0.12529\\
0.868	-0.1253\\
0.87	-0.125309\\
0.872	-0.125319\\
0.874	-0.125328\\
0.876	-0.125338\\
0.878	-0.125347\\
0.88	-0.125356\\
0.882	-0.125365\\
0.884	-0.125374\\
0.886	-0.125382\\
0.888	-0.125391\\
0.89	-0.125399\\
0.892	-0.125407\\
0.894	-0.125415\\
0.896	-0.125423\\
0.898	-0.12543\\
0.9	-0.125437\\
0.902	-0.125444\\
0.904	-0.125451\\
0.906	-0.125457\\
0.908	-0.125463\\
0.91	-0.125469\\
0.912	-0.125474\\
0.914	-0.125479\\
0.916	-0.125484\\
0.918	-0.125488\\
0.92	-0.125492\\
0.922	-0.125496\\
0.924	-0.1255\\
0.926	-0.125503\\
0.928	-0.125505\\
0.93	-0.125508\\
0.932	-0.12551\\
0.934	-0.125511\\
0.936	-0.125513\\
0.938	-0.125513\\
0.94	-0.125514\\
0.942	-0.125514\\
0.944	-0.125514\\
0.946	-0.125513\\
0.948	-0.125512\\
0.95	-0.125511\\
0.952	-0.125509\\
0.954	-0.125507\\
0.956	-0.125505\\
0.958	-0.125502\\
0.96	-0.125499\\
0.962	-0.125495\\
0.964	-0.125492\\
0.966	-0.125487\\
0.968	-0.125483\\
0.97	-0.125478\\
0.972	-0.125473\\
0.974	-0.125467\\
0.976	-0.125461\\
0.978	-0.125455\\
0.98	-0.125448\\
0.982	-0.125441\\
0.984	-0.125434\\
0.986	-0.125426\\
0.988	-0.125418\\
0.99	-0.12541\\
0.992	-0.125402\\
0.994	-0.125393\\
0.996	-0.125384\\
0.998	-0.125375\\
1	-0.125365\\
1.002	-0.125355\\
1.004	-0.125345\\
1.006	-0.125335\\
1.008	-0.125324\\
1.01	-0.125314\\
1.012	-0.125303\\
1.014	-0.125291\\
1.016	-0.12528\\
1.018	-0.125268\\
1.02	-0.125257\\
1.022	-0.125245\\
1.024	-0.125232\\
1.026	-0.12522\\
1.028	-0.125208\\
1.03	-0.125195\\
1.032	-0.125183\\
1.034	-0.12517\\
1.036	-0.125157\\
1.038	-0.125144\\
1.04	-0.125131\\
1.042	-0.125118\\
1.044	-0.125105\\
1.046	-0.125092\\
1.048	-0.125079\\
1.05	-0.125066\\
1.052	-0.125053\\
1.054	-0.125039\\
1.056	-0.125026\\
1.058	-0.125013\\
1.06	-0.125\\
1.062	-0.124987\\
1.064	-0.124974\\
1.066	-0.124961\\
1.068	-0.124949\\
1.07	-0.124936\\
1.072	-0.124923\\
1.074	-0.124911\\
1.076	-0.124899\\
1.078	-0.124886\\
1.08	-0.124874\\
1.082	-0.124863\\
1.084	-0.124851\\
1.086	-0.124839\\
1.088	-0.124828\\
1.09	-0.124817\\
1.092	-0.124806\\
1.094	-0.124795\\
1.096	-0.124785\\
1.098	-0.124775\\
1.1	-0.124765\\
1.102	-0.124755\\
1.104	-0.124746\\
1.106	-0.124736\\
1.108	-0.124727\\
1.11	-0.124719\\
1.112	-0.12471\\
1.114	-0.124702\\
1.116	-0.124694\\
1.118	-0.124687\\
1.12	-0.124679\\
1.122	-0.124672\\
1.124	-0.124666\\
1.126	-0.124659\\
1.128	-0.124653\\
1.13	-0.124648\\
1.132	-0.124642\\
1.134	-0.124637\\
1.136	-0.124632\\
1.138	-0.124628\\
1.14	-0.124623\\
1.142	-0.124619\\
1.144	-0.124616\\
1.146	-0.124612\\
1.148	-0.124609\\
1.15	-0.124607\\
1.152	-0.124604\\
1.154	-0.124602\\
1.156	-0.1246\\
1.158	-0.124598\\
1.16	-0.124597\\
1.162	-0.124596\\
1.164	-0.124595\\
1.166	-0.124595\\
1.168	-0.124594\\
1.17	-0.124594\\
1.172	-0.124594\\
1.174	-0.124595\\
1.176	-0.124595\\
1.178	-0.124596\\
1.18	-0.124597\\
1.182	-0.124599\\
1.184	-0.1246\\
1.186	-0.124602\\
1.188	-0.124604\\
1.19	-0.124606\\
1.192	-0.124608\\
1.194	-0.12461\\
1.196	-0.124612\\
1.198	-0.124615\\
1.2	-0.124618\\
1.202	-0.124621\\
1.204	-0.124624\\
1.206	-0.124627\\
1.208	-0.12463\\
1.21	-0.124633\\
1.212	-0.124636\\
1.214	-0.12464\\
1.216	-0.124643\\
1.218	-0.124646\\
1.22	-0.12465\\
1.222	-0.124653\\
1.224	-0.124657\\
1.226	-0.124661\\
1.228	-0.124664\\
1.23	-0.124668\\
1.232	-0.124671\\
1.234	-0.124675\\
1.236	-0.124678\\
1.238	-0.124682\\
1.24	-0.124685\\
1.242	-0.124689\\
1.244	-0.124692\\
1.246	-0.124695\\
1.248	-0.124699\\
1.25	-0.124702\\
1.252	-0.124705\\
1.254	-0.124708\\
1.256	-0.124711\\
1.258	-0.124714\\
1.26	-0.124717\\
1.262	-0.124719\\
1.264	-0.124722\\
1.266	-0.124724\\
1.268	-0.124726\\
1.27	-0.124729\\
1.272	-0.124731\\
1.274	-0.124733\\
1.276	-0.124735\\
1.278	-0.124736\\
1.28	-0.124738\\
1.282	-0.12474\\
1.284	-0.124741\\
1.286	-0.124742\\
1.288	-0.124743\\
1.29	-0.124744\\
1.292	-0.124745\\
1.294	-0.124746\\
1.296	-0.124746\\
1.298	-0.124747\\
1.3	-0.124747\\
1.302	-0.124747\\
1.304	-0.124747\\
1.306	-0.124747\\
1.308	-0.124747\\
1.31	-0.124747\\
1.312	-0.124746\\
1.314	-0.124746\\
1.316	-0.124745\\
1.318	-0.124744\\
1.32	-0.124744\\
1.322	-0.124743\\
1.324	-0.124741\\
1.326	-0.12474\\
1.328	-0.124739\\
1.33	-0.124737\\
1.332	-0.124736\\
1.334	-0.124734\\
1.336	-0.124732\\
1.338	-0.124731\\
1.34	-0.124729\\
1.342	-0.124727\\
1.344	-0.124725\\
1.346	-0.124722\\
1.348	-0.12472\\
1.35	-0.124718\\
1.352	-0.124715\\
1.354	-0.124713\\
1.356	-0.12471\\
1.358	-0.124708\\
1.36	-0.124705\\
1.362	-0.124702\\
1.364	-0.1247\\
1.366	-0.124697\\
1.368	-0.124694\\
1.37	-0.124691\\
1.372	-0.124688\\
1.374	-0.124685\\
1.376	-0.124682\\
1.378	-0.124679\\
1.38	-0.124676\\
1.382	-0.124673\\
1.384	-0.12467\\
1.386	-0.124667\\
1.388	-0.124664\\
1.39	-0.124661\\
1.392	-0.124657\\
1.394	-0.124654\\
1.396	-0.124651\\
1.398	-0.124648\\
1.4	-0.124645\\
1.402	-0.124641\\
1.404	-0.124638\\
1.406	-0.124635\\
1.408	-0.124632\\
1.41	-0.124629\\
1.412	-0.124625\\
1.414	-0.124622\\
1.416	-0.124619\\
1.418	-0.124616\\
1.42	-0.124613\\
1.422	-0.12461\\
1.424	-0.124606\\
1.426	-0.124603\\
1.428	-0.1246\\
1.43	-0.124597\\
1.432	-0.124594\\
1.434	-0.124591\\
1.436	-0.124588\\
1.438	-0.124585\\
1.44	-0.124582\\
1.442	-0.124579\\
1.444	-0.124576\\
1.446	-0.124573\\
1.448	-0.12457\\
1.45	-0.124568\\
1.452	-0.124565\\
1.454	-0.124562\\
1.456	-0.124559\\
1.458	-0.124556\\
1.46	-0.124554\\
1.462	-0.124551\\
1.464	-0.124548\\
1.466	-0.124546\\
1.468	-0.124543\\
1.47	-0.12454\\
1.472	-0.124538\\
1.474	-0.124535\\
1.476	-0.124533\\
1.478	-0.12453\\
1.48	-0.124528\\
1.482	-0.124525\\
1.484	-0.124523\\
1.486	-0.12452\\
1.488	-0.124518\\
1.49	-0.124516\\
1.492	-0.124513\\
1.494	-0.124511\\
1.496	-0.124509\\
1.498	-0.124507\\
1.5	-0.124504\\
1.502	-0.124502\\
1.504	-0.1245\\
1.506	-0.124498\\
1.508	-0.124496\\
1.51	-0.124494\\
1.512	-0.124492\\
1.514	-0.12449\\
1.516	-0.124488\\
1.518	-0.124486\\
1.52	-0.124484\\
1.522	-0.124482\\
1.524	-0.12448\\
1.526	-0.124478\\
1.528	-0.124476\\
1.53	-0.124475\\
1.532	-0.124473\\
1.534	-0.124471\\
1.536	-0.124469\\
1.538	-0.124467\\
1.54	-0.124466\\
1.542	-0.124464\\
1.544	-0.124462\\
1.546	-0.124461\\
1.548	-0.124459\\
1.55	-0.124457\\
1.552	-0.124456\\
1.554	-0.124454\\
1.556	-0.124453\\
1.558	-0.124451\\
1.56	-0.12445\\
1.562	-0.124448\\
1.564	-0.124447\\
1.566	-0.124445\\
1.568	-0.124444\\
1.57	-0.124442\\
1.572	-0.124441\\
1.574	-0.124439\\
1.576	-0.124438\\
1.578	-0.124436\\
1.58	-0.124435\\
1.582	-0.124433\\
1.584	-0.124432\\
1.586	-0.124431\\
1.588	-0.124429\\
1.59	-0.124428\\
1.592	-0.124427\\
1.594	-0.124425\\
1.596	-0.124424\\
1.598	-0.124422\\
1.6	-0.124421\\
1.602	-0.12442\\
1.604	-0.124418\\
1.606	-0.124417\\
1.608	-0.124416\\
1.61	-0.124414\\
1.612	-0.124413\\
1.614	-0.124411\\
1.616	-0.12441\\
1.618	-0.124409\\
1.62	-0.124407\\
1.622	-0.124406\\
1.624	-0.124405\\
1.626	-0.124403\\
1.628	-0.124402\\
1.63	-0.1244\\
1.632	-0.124399\\
1.634	-0.124397\\
1.636	-0.124396\\
1.638	-0.124394\\
1.64	-0.124393\\
1.642	-0.124392\\
1.644	-0.12439\\
1.646	-0.124388\\
1.648	-0.124387\\
1.65	-0.124385\\
1.652	-0.124384\\
1.654	-0.124382\\
1.656	-0.124381\\
1.658	-0.124379\\
1.66	-0.124377\\
1.662	-0.124376\\
1.664	-0.124374\\
1.666	-0.124372\\
1.668	-0.124371\\
1.67	-0.124369\\
1.672	-0.124367\\
1.674	-0.124365\\
1.676	-0.124364\\
1.678	-0.124362\\
1.68	-0.12436\\
1.682	-0.124358\\
1.684	-0.124356\\
1.686	-0.124355\\
1.688	-0.124353\\
1.69	-0.124351\\
1.692	-0.124349\\
1.694	-0.124347\\
1.696	-0.124345\\
1.698	-0.124343\\
1.7	-0.124341\\
1.702	-0.124339\\
1.704	-0.124337\\
1.706	-0.124335\\
1.708	-0.124333\\
1.71	-0.124331\\
1.712	-0.124329\\
1.714	-0.124327\\
1.716	-0.124324\\
1.718	-0.124322\\
1.72	-0.12432\\
1.722	-0.124318\\
1.724	-0.124316\\
1.726	-0.124314\\
1.728	-0.124311\\
1.73	-0.124309\\
1.732	-0.124307\\
1.734	-0.124305\\
1.736	-0.124303\\
1.738	-0.1243\\
1.74	-0.124298\\
1.742	-0.124296\\
1.744	-0.124294\\
1.746	-0.124291\\
1.748	-0.124289\\
1.75	-0.124287\\
1.752	-0.124285\\
1.754	-0.124282\\
1.756	-0.12428\\
1.758	-0.124278\\
1.76	-0.124275\\
1.762	-0.124273\\
1.764	-0.124271\\
1.766	-0.124269\\
1.768	-0.124266\\
1.77	-0.124264\\
1.772	-0.124262\\
1.774	-0.124259\\
1.776	-0.124257\\
1.778	-0.124255\\
1.78	-0.124253\\
1.782	-0.12425\\
1.784	-0.124248\\
1.786	-0.124246\\
1.788	-0.124244\\
1.79	-0.124241\\
1.792	-0.124239\\
1.794	-0.124237\\
1.796	-0.124235\\
1.798	-0.124233\\
1.8	-0.12423\\
1.802	-0.124228\\
1.804	-0.124226\\
1.806	-0.124224\\
1.808	-0.124222\\
1.81	-0.12422\\
1.812	-0.124217\\
1.814	-0.124215\\
1.816	-0.124213\\
1.818	-0.124211\\
1.82	-0.124209\\
1.822	-0.124207\\
1.824	-0.124205\\
1.826	-0.124203\\
1.828	-0.124201\\
1.83	-0.124199\\
1.832	-0.124197\\
1.834	-0.124195\\
1.836	-0.124193\\
1.838	-0.124191\\
1.84	-0.124189\\
1.842	-0.124187\\
1.844	-0.124185\\
1.846	-0.124183\\
1.848	-0.124181\\
1.85	-0.124179\\
1.852	-0.124177\\
1.854	-0.124175\\
1.856	-0.124173\\
1.858	-0.124171\\
1.86	-0.124169\\
1.862	-0.124168\\
1.864	-0.124166\\
1.866	-0.124164\\
1.868	-0.124162\\
1.87	-0.12416\\
1.872	-0.124158\\
1.874	-0.124157\\
1.876	-0.124155\\
1.878	-0.124153\\
1.88	-0.124151\\
1.882	-0.124149\\
1.884	-0.124148\\
1.886	-0.124146\\
1.888	-0.124144\\
1.89	-0.124142\\
1.892	-0.124141\\
1.894	-0.124139\\
1.896	-0.124137\\
1.898	-0.124135\\
1.9	-0.124134\\
1.902	-0.124132\\
1.904	-0.12413\\
1.906	-0.124128\\
1.908	-0.124127\\
1.91	-0.124125\\
1.912	-0.124123\\
1.914	-0.124122\\
1.916	-0.12412\\
1.918	-0.124118\\
1.92	-0.124117\\
1.922	-0.124115\\
1.924	-0.124113\\
1.926	-0.124111\\
1.928	-0.12411\\
1.93	-0.124108\\
1.932	-0.124106\\
1.934	-0.124105\\
1.936	-0.124103\\
1.938	-0.124101\\
1.94	-0.1241\\
1.942	-0.124098\\
1.944	-0.124096\\
1.946	-0.124095\\
1.948	-0.124093\\
1.95	-0.124091\\
1.952	-0.12409\\
1.954	-0.124088\\
1.956	-0.124086\\
1.958	-0.124084\\
1.96	-0.124083\\
1.962	-0.124081\\
1.964	-0.124079\\
1.966	-0.124078\\
1.968	-0.124076\\
1.97	-0.124074\\
1.972	-0.124073\\
1.974	-0.124071\\
1.976	-0.124069\\
1.978	-0.124068\\
1.98	-0.124066\\
1.982	-0.124064\\
1.984	-0.124062\\
1.986	-0.124061\\
1.988	-0.124059\\
1.99	-0.124057\\
1.992	-0.124056\\
1.994	-0.124054\\
1.996	-0.124052\\
1.998	-0.12405\\
2	-0.124049\\
};
\addplot [color=mycolor13, line width=1.0pt, forget plot]
  table[row sep=crcr]{%
-0.124	0\\
-0.122	0\\
-0.12	0\\
-0.118	0\\
-0.116	0\\
-0.114	0\\
-0.112	0\\
-0.11	0\\
-0.108	0\\
-0.106	0\\
-0.104	0\\
-0.102	0\\
-0.1	0\\
-0.098	0\\
-0.096	0\\
-0.094	0\\
-0.092	0\\
-0.09	0\\
-0.088	0\\
-0.086	0\\
-0.084	0\\
-0.082	0\\
-0.08	0\\
-0.078	0\\
-0.076	0\\
-0.074	0\\
-0.072	0\\
-0.07	0\\
-0.068	0\\
-0.066	0\\
-0.064	0\\
-0.062	0\\
-0.06	0\\
-0.058	0\\
-0.056	0\\
-0.054	0\\
-0.052	0\\
-0.05	0\\
-0.048	0\\
-0.046	0\\
-0.044	0\\
-0.042	0\\
-0.04	0\\
-0.038	0\\
-0.036	0\\
-0.034	0\\
-0.032	0\\
-0.03	0\\
-0.028	0\\
-0.026	0\\
-0.024	0\\
-0.022	0\\
-0.02	0\\
-0.018	0\\
-0.016	0\\
-0.014	0\\
-0.012	0\\
-0.01	0\\
-0.008	0\\
-0.006	0\\
-0.004	0\\
-0.002	0\\
0	0\\
0.002	-0.000454142\\
0.004	-0.000971506\\
0.006	-0.00154702\\
0.008	-0.00217446\\
0.01	-0.00284765\\
0.012	-0.00356097\\
0.014	-0.00430952\\
0.016	-0.00508916\\
0.018	-0.0058964\\
0.02	-0.00672839\\
0.022	-0.00758276\\
0.024	-0.00845759\\
0.026	-0.00935134\\
0.028	-0.0102627\\
0.03	-0.0111908\\
0.032	-0.0121347\\
0.034	-0.0130937\\
0.036	-0.0140674\\
0.038	-0.0150553\\
0.04	-0.0160571\\
0.042	-0.0170725\\
0.044	-0.0181011\\
0.046	-0.0191428\\
0.048	-0.0201974\\
0.05	-0.0212645\\
0.052	-0.0223441\\
0.054	-0.0234357\\
0.056	-0.0245393\\
0.058	-0.0256545\\
0.06	-0.026781\\
0.062	-0.0279185\\
0.064	-0.0290667\\
0.066	-0.0302251\\
0.068	-0.0313934\\
0.07	-0.0325711\\
0.072	-0.0337578\\
0.074	-0.034953\\
0.076	-0.0361562\\
0.078	-0.0373668\\
0.08	-0.0385843\\
0.082	-0.039808\\
0.084	-0.0410375\\
0.086	-0.0422721\\
0.088	-0.043511\\
0.09	-0.0447537\\
0.092	-0.0459995\\
0.094	-0.0472478\\
0.096	-0.0484977\\
0.098	-0.0497486\\
0.1	-0.0509998\\
0.102	-0.0522506\\
0.104	-0.0535003\\
0.106	-0.054748\\
0.108	-0.0559932\\
0.11	-0.0572351\\
0.112	-0.0584729\\
0.114	-0.0597059\\
0.116	-0.0609334\\
0.118	-0.0621548\\
0.12	-0.0633693\\
0.122	-0.0645763\\
0.124	-0.065775\\
0.126	-0.0669649\\
0.128	-0.0681451\\
0.13	-0.0693153\\
0.132	-0.0704746\\
0.134	-0.0716225\\
0.136	-0.0727585\\
0.138	-0.0738819\\
0.14	-0.0749923\\
0.142	-0.076089\\
0.144	-0.0771716\\
0.146	-0.0782396\\
0.148	-0.0792926\\
0.15	-0.08033\\
0.152	-0.0813516\\
0.154	-0.0823568\\
0.156	-0.0833453\\
0.158	-0.0843168\\
0.16	-0.0852709\\
0.162	-0.0862073\\
0.164	-0.0871258\\
0.166	-0.088026\\
0.168	-0.0889078\\
0.17	-0.089771\\
0.172	-0.0906153\\
0.174	-0.0914406\\
0.176	-0.0922467\\
0.178	-0.0930336\\
0.18	-0.093801\\
0.182	-0.094549\\
0.184	-0.0952775\\
0.186	-0.0959865\\
0.188	-0.0966759\\
0.19	-0.0973457\\
0.192	-0.097996\\
0.194	-0.0986267\\
0.196	-0.0992381\\
0.198	-0.09983\\
0.2	-0.100403\\
0.202	-0.100956\\
0.204	-0.101491\\
0.206	-0.102006\\
0.208	-0.102503\\
0.21	-0.102982\\
0.212	-0.103442\\
0.214	-0.103884\\
0.216	-0.104307\\
0.218	-0.104714\\
0.22	-0.105102\\
0.222	-0.105474\\
0.224	-0.105828\\
0.226	-0.106166\\
0.228	-0.106487\\
0.23	-0.106792\\
0.232	-0.107081\\
0.234	-0.107355\\
0.236	-0.107613\\
0.238	-0.107856\\
0.24	-0.108085\\
0.242	-0.1083\\
0.244	-0.1085\\
0.246	-0.108687\\
0.248	-0.108861\\
0.25	-0.109021\\
0.252	-0.109169\\
0.254	-0.109305\\
0.256	-0.109429\\
0.258	-0.109541\\
0.26	-0.109642\\
0.262	-0.109733\\
0.264	-0.109812\\
0.266	-0.109882\\
0.268	-0.109942\\
0.27	-0.109993\\
0.272	-0.110034\\
0.274	-0.110067\\
0.276	-0.110092\\
0.278	-0.110109\\
0.28	-0.110118\\
0.282	-0.11012\\
0.284	-0.110115\\
0.286	-0.110103\\
0.288	-0.110086\\
0.29	-0.110063\\
0.292	-0.110034\\
0.294	-0.11\\
0.296	-0.109962\\
0.298	-0.109919\\
0.3	-0.109872\\
0.302	-0.109821\\
0.304	-0.109767\\
0.306	-0.10971\\
0.308	-0.10965\\
0.31	-0.109588\\
0.312	-0.109524\\
0.314	-0.109458\\
0.316	-0.10939\\
0.318	-0.109322\\
0.32	-0.109252\\
0.322	-0.109182\\
0.324	-0.109112\\
0.326	-0.109042\\
0.328	-0.108972\\
0.33	-0.108903\\
0.332	-0.108835\\
0.334	-0.108768\\
0.336	-0.108702\\
0.338	-0.108638\\
0.34	-0.108576\\
0.342	-0.108516\\
0.344	-0.108459\\
0.346	-0.108404\\
0.348	-0.108352\\
0.35	-0.108303\\
0.352	-0.108258\\
0.354	-0.108216\\
0.356	-0.108177\\
0.358	-0.108143\\
0.36	-0.108113\\
0.362	-0.108087\\
0.364	-0.108065\\
0.366	-0.108048\\
0.368	-0.108035\\
0.37	-0.108028\\
0.372	-0.108025\\
0.374	-0.108028\\
0.376	-0.108036\\
0.378	-0.108049\\
0.38	-0.108068\\
0.382	-0.108092\\
0.384	-0.108122\\
0.386	-0.108157\\
0.388	-0.108199\\
0.39	-0.108246\\
0.392	-0.108299\\
0.394	-0.108358\\
0.396	-0.108423\\
0.398	-0.108494\\
0.4	-0.108571\\
0.402	-0.108654\\
0.404	-0.108743\\
0.406	-0.108838\\
0.408	-0.108939\\
0.41	-0.109046\\
0.412	-0.109159\\
0.414	-0.109277\\
0.416	-0.109402\\
0.418	-0.109532\\
0.42	-0.109668\\
0.422	-0.10981\\
0.424	-0.109957\\
0.426	-0.110109\\
0.428	-0.110267\\
0.43	-0.11043\\
0.432	-0.110599\\
0.434	-0.110772\\
0.436	-0.11095\\
0.438	-0.111133\\
0.44	-0.11132\\
0.442	-0.111512\\
0.444	-0.111709\\
0.446	-0.111909\\
0.448	-0.112113\\
0.45	-0.112322\\
0.452	-0.112534\\
0.454	-0.112749\\
0.456	-0.112968\\
0.458	-0.11319\\
0.46	-0.113415\\
0.462	-0.113642\\
0.464	-0.113873\\
0.466	-0.114105\\
0.468	-0.11434\\
0.47	-0.114577\\
0.472	-0.114816\\
0.474	-0.115056\\
0.476	-0.115298\\
0.478	-0.115541\\
0.48	-0.115785\\
0.482	-0.11603\\
0.484	-0.116276\\
0.486	-0.116522\\
0.488	-0.116768\\
0.49	-0.117014\\
0.492	-0.117261\\
0.494	-0.117506\\
0.496	-0.117752\\
0.498	-0.117996\\
0.5	-0.11824\\
0.502	-0.118483\\
0.504	-0.118724\\
0.506	-0.118964\\
0.508	-0.119203\\
0.51	-0.11944\\
0.512	-0.119675\\
0.514	-0.119907\\
0.516	-0.120138\\
0.518	-0.120366\\
0.52	-0.120592\\
0.522	-0.120815\\
0.524	-0.121035\\
0.526	-0.121252\\
0.528	-0.121467\\
0.53	-0.121678\\
0.532	-0.121886\\
0.534	-0.12209\\
0.536	-0.122291\\
0.538	-0.122489\\
0.54	-0.122682\\
0.542	-0.122872\\
0.544	-0.123059\\
0.546	-0.123241\\
0.548	-0.123419\\
0.55	-0.123593\\
0.552	-0.123763\\
0.554	-0.123929\\
0.556	-0.124091\\
0.558	-0.124249\\
0.56	-0.124402\\
0.562	-0.124551\\
0.564	-0.124696\\
0.566	-0.124836\\
0.568	-0.124972\\
0.57	-0.125104\\
0.572	-0.125231\\
0.574	-0.125354\\
0.576	-0.125472\\
0.578	-0.125587\\
0.58	-0.125696\\
0.582	-0.125802\\
0.584	-0.125903\\
0.586	-0.126\\
0.588	-0.126093\\
0.59	-0.126182\\
0.592	-0.126266\\
0.594	-0.126347\\
0.596	-0.126423\\
0.598	-0.126496\\
0.6	-0.126564\\
0.602	-0.126629\\
0.604	-0.126689\\
0.606	-0.126746\\
0.608	-0.1268\\
0.61	-0.126849\\
0.612	-0.126896\\
0.614	-0.126938\\
0.616	-0.126978\\
0.618	-0.127013\\
0.62	-0.127046\\
0.622	-0.127076\\
0.624	-0.127102\\
0.626	-0.127125\\
0.628	-0.127146\\
0.63	-0.127163\\
0.632	-0.127178\\
0.634	-0.12719\\
0.636	-0.1272\\
0.638	-0.127207\\
0.64	-0.127212\\
0.642	-0.127214\\
0.644	-0.127214\\
0.646	-0.127212\\
0.648	-0.127207\\
0.65	-0.127201\\
0.652	-0.127193\\
0.654	-0.127183\\
0.656	-0.127171\\
0.658	-0.127157\\
0.66	-0.127142\\
0.662	-0.127125\\
0.664	-0.127107\\
0.666	-0.127088\\
0.668	-0.127067\\
0.67	-0.127045\\
0.672	-0.127022\\
0.674	-0.126997\\
0.676	-0.126972\\
0.678	-0.126946\\
0.68	-0.126919\\
0.682	-0.126891\\
0.684	-0.126862\\
0.686	-0.126833\\
0.688	-0.126803\\
0.69	-0.126773\\
0.692	-0.126742\\
0.694	-0.12671\\
0.696	-0.126678\\
0.698	-0.126646\\
0.7	-0.126614\\
0.702	-0.126581\\
0.704	-0.126549\\
0.706	-0.126516\\
0.708	-0.126483\\
0.71	-0.12645\\
0.712	-0.126417\\
0.714	-0.126384\\
0.716	-0.126351\\
0.718	-0.126318\\
0.72	-0.126286\\
0.722	-0.126253\\
0.724	-0.126221\\
0.726	-0.126189\\
0.728	-0.126158\\
0.73	-0.126127\\
0.732	-0.126096\\
0.734	-0.126065\\
0.736	-0.126035\\
0.738	-0.126006\\
0.74	-0.125977\\
0.742	-0.125948\\
0.744	-0.12592\\
0.746	-0.125893\\
0.748	-0.125866\\
0.75	-0.12584\\
0.752	-0.125814\\
0.754	-0.125789\\
0.756	-0.125764\\
0.758	-0.12574\\
0.76	-0.125717\\
0.762	-0.125694\\
0.764	-0.125672\\
0.766	-0.125651\\
0.768	-0.125631\\
0.77	-0.125611\\
0.772	-0.125592\\
0.774	-0.125573\\
0.776	-0.125555\\
0.778	-0.125538\\
0.78	-0.125522\\
0.782	-0.125507\\
0.784	-0.125492\\
0.786	-0.125478\\
0.788	-0.125464\\
0.79	-0.125451\\
0.792	-0.12544\\
0.794	-0.125428\\
0.796	-0.125418\\
0.798	-0.125408\\
0.8	-0.125399\\
0.802	-0.12539\\
0.804	-0.125382\\
0.806	-0.125375\\
0.808	-0.125369\\
0.81	-0.125363\\
0.812	-0.125358\\
0.814	-0.125353\\
0.816	-0.125349\\
0.818	-0.125346\\
0.82	-0.125343\\
0.822	-0.125341\\
0.824	-0.12534\\
0.826	-0.125339\\
0.828	-0.125338\\
0.83	-0.125338\\
0.832	-0.125339\\
0.834	-0.12534\\
0.836	-0.125341\\
0.838	-0.125343\\
0.84	-0.125346\\
0.842	-0.125349\\
0.844	-0.125352\\
0.846	-0.125355\\
0.848	-0.125359\\
0.85	-0.125364\\
0.852	-0.125368\\
0.854	-0.125373\\
0.856	-0.125379\\
0.858	-0.125384\\
0.86	-0.12539\\
0.862	-0.125396\\
0.864	-0.125402\\
0.866	-0.125408\\
0.868	-0.125415\\
0.87	-0.125421\\
0.872	-0.125428\\
0.874	-0.125435\\
0.876	-0.125442\\
0.878	-0.125449\\
0.88	-0.125456\\
0.882	-0.125463\\
0.884	-0.125469\\
0.886	-0.125476\\
0.888	-0.125483\\
0.89	-0.12549\\
0.892	-0.125497\\
0.894	-0.125503\\
0.896	-0.12551\\
0.898	-0.125516\\
0.9	-0.125522\\
0.902	-0.125528\\
0.904	-0.125534\\
0.906	-0.12554\\
0.908	-0.125545\\
0.91	-0.12555\\
0.912	-0.125555\\
0.914	-0.125559\\
0.916	-0.125564\\
0.918	-0.125568\\
0.92	-0.125571\\
0.922	-0.125575\\
0.924	-0.125578\\
0.926	-0.12558\\
0.928	-0.125583\\
0.93	-0.125584\\
0.932	-0.125586\\
0.934	-0.125587\\
0.936	-0.125588\\
0.938	-0.125588\\
0.94	-0.125588\\
0.942	-0.125588\\
0.944	-0.125587\\
0.946	-0.125586\\
0.948	-0.125584\\
0.95	-0.125582\\
0.952	-0.12558\\
0.954	-0.125577\\
0.956	-0.125574\\
0.958	-0.12557\\
0.96	-0.125566\\
0.962	-0.125562\\
0.964	-0.125557\\
0.966	-0.125552\\
0.968	-0.125546\\
0.97	-0.12554\\
0.972	-0.125534\\
0.974	-0.125527\\
0.976	-0.12552\\
0.978	-0.125512\\
0.98	-0.125504\\
0.982	-0.125496\\
0.984	-0.125488\\
0.986	-0.125479\\
0.988	-0.12547\\
0.99	-0.12546\\
0.992	-0.125451\\
0.994	-0.125441\\
0.996	-0.12543\\
0.998	-0.12542\\
1	-0.125409\\
1.002	-0.125398\\
1.004	-0.125387\\
1.006	-0.125376\\
1.008	-0.125364\\
1.01	-0.125352\\
1.012	-0.12534\\
1.014	-0.125328\\
1.016	-0.125316\\
1.018	-0.125304\\
1.02	-0.125292\\
1.022	-0.125279\\
1.024	-0.125266\\
1.026	-0.125254\\
1.028	-0.125241\\
1.03	-0.125228\\
1.032	-0.125216\\
1.034	-0.125203\\
1.036	-0.12519\\
1.038	-0.125177\\
1.04	-0.125165\\
1.042	-0.125152\\
1.044	-0.12514\\
1.046	-0.125127\\
1.048	-0.125115\\
1.05	-0.125102\\
1.052	-0.12509\\
1.054	-0.125078\\
1.056	-0.125066\\
1.058	-0.125054\\
1.06	-0.125042\\
1.062	-0.125031\\
1.064	-0.125019\\
1.066	-0.125008\\
1.068	-0.124997\\
1.07	-0.124986\\
1.072	-0.124975\\
1.074	-0.124965\\
1.076	-0.124954\\
1.078	-0.124944\\
1.08	-0.124934\\
1.082	-0.124925\\
1.084	-0.124915\\
1.086	-0.124906\\
1.088	-0.124897\\
1.09	-0.124889\\
1.092	-0.12488\\
1.094	-0.124872\\
1.096	-0.124864\\
1.098	-0.124856\\
1.1	-0.124849\\
1.102	-0.124842\\
1.104	-0.124835\\
1.106	-0.124828\\
1.108	-0.124822\\
1.11	-0.124816\\
1.112	-0.12481\\
1.114	-0.124804\\
1.116	-0.124799\\
1.118	-0.124794\\
1.12	-0.124789\\
1.122	-0.124784\\
1.124	-0.12478\\
1.126	-0.124776\\
1.128	-0.124772\\
1.13	-0.124768\\
1.132	-0.124765\\
1.134	-0.124761\\
1.136	-0.124758\\
1.138	-0.124756\\
1.14	-0.124753\\
1.142	-0.124751\\
1.144	-0.124749\\
1.146	-0.124747\\
1.148	-0.124745\\
1.15	-0.124743\\
1.152	-0.124742\\
1.154	-0.124741\\
1.156	-0.12474\\
1.158	-0.124739\\
1.16	-0.124738\\
1.162	-0.124738\\
1.164	-0.124738\\
1.166	-0.124737\\
1.168	-0.124737\\
1.17	-0.124737\\
1.172	-0.124738\\
1.174	-0.124738\\
1.176	-0.124738\\
1.178	-0.124739\\
1.18	-0.12474\\
1.182	-0.12474\\
1.184	-0.124741\\
1.186	-0.124742\\
1.188	-0.124743\\
1.19	-0.124744\\
1.192	-0.124746\\
1.194	-0.124747\\
1.196	-0.124748\\
1.198	-0.12475\\
1.2	-0.124751\\
1.202	-0.124753\\
1.204	-0.124754\\
1.206	-0.124756\\
1.208	-0.124757\\
1.21	-0.124759\\
1.212	-0.124761\\
1.214	-0.124762\\
1.216	-0.124764\\
1.218	-0.124766\\
1.22	-0.124767\\
1.222	-0.124769\\
1.224	-0.124771\\
1.226	-0.124773\\
1.228	-0.124774\\
1.23	-0.124776\\
1.232	-0.124778\\
1.234	-0.124779\\
1.236	-0.124781\\
1.238	-0.124783\\
1.24	-0.124784\\
1.242	-0.124786\\
1.244	-0.124787\\
1.246	-0.124789\\
1.248	-0.12479\\
1.25	-0.124792\\
1.252	-0.124793\\
1.254	-0.124795\\
1.256	-0.124796\\
1.258	-0.124797\\
1.26	-0.124798\\
1.262	-0.1248\\
1.264	-0.124801\\
1.266	-0.124802\\
1.268	-0.124803\\
1.27	-0.124804\\
1.272	-0.124805\\
1.274	-0.124805\\
1.276	-0.124806\\
1.278	-0.124807\\
1.28	-0.124807\\
1.282	-0.124808\\
1.284	-0.124809\\
1.286	-0.124809\\
1.288	-0.124809\\
1.29	-0.12481\\
1.292	-0.12481\\
1.294	-0.12481\\
1.296	-0.12481\\
1.298	-0.12481\\
1.3	-0.12481\\
1.302	-0.12481\\
1.304	-0.12481\\
1.306	-0.12481\\
1.308	-0.124809\\
1.31	-0.124809\\
1.312	-0.124808\\
1.314	-0.124808\\
1.316	-0.124807\\
1.318	-0.124806\\
1.32	-0.124806\\
1.322	-0.124805\\
1.324	-0.124804\\
1.326	-0.124803\\
1.328	-0.124802\\
1.33	-0.124801\\
1.332	-0.1248\\
1.334	-0.124798\\
1.336	-0.124797\\
1.338	-0.124796\\
1.34	-0.124794\\
1.342	-0.124793\\
1.344	-0.124791\\
1.346	-0.12479\\
1.348	-0.124788\\
1.35	-0.124786\\
1.352	-0.124784\\
1.354	-0.124782\\
1.356	-0.12478\\
1.358	-0.124778\\
1.36	-0.124776\\
1.362	-0.124774\\
1.364	-0.124772\\
1.366	-0.12477\\
1.368	-0.124768\\
1.37	-0.124765\\
1.372	-0.124763\\
1.374	-0.12476\\
1.376	-0.124758\\
1.378	-0.124755\\
1.38	-0.124753\\
1.382	-0.12475\\
1.384	-0.124748\\
1.386	-0.124745\\
1.388	-0.124742\\
1.39	-0.12474\\
1.392	-0.124737\\
1.394	-0.124734\\
1.396	-0.124731\\
1.398	-0.124728\\
1.4	-0.124726\\
1.402	-0.124723\\
1.404	-0.12472\\
1.406	-0.124717\\
1.408	-0.124714\\
1.41	-0.124711\\
1.412	-0.124708\\
1.414	-0.124705\\
1.416	-0.124702\\
1.418	-0.124699\\
1.42	-0.124696\\
1.422	-0.124693\\
1.424	-0.12469\\
1.426	-0.124687\\
1.428	-0.124684\\
1.43	-0.124681\\
1.432	-0.124678\\
1.434	-0.124675\\
1.436	-0.124672\\
1.438	-0.124669\\
1.44	-0.124666\\
1.442	-0.124663\\
1.444	-0.12466\\
1.446	-0.124657\\
1.448	-0.124654\\
1.45	-0.124651\\
1.452	-0.124648\\
1.454	-0.124645\\
1.456	-0.124642\\
1.458	-0.12464\\
1.46	-0.124637\\
1.462	-0.124634\\
1.464	-0.124631\\
1.466	-0.124628\\
1.468	-0.124626\\
1.47	-0.124623\\
1.472	-0.12462\\
1.474	-0.124618\\
1.476	-0.124615\\
1.478	-0.124613\\
1.48	-0.12461\\
1.482	-0.124608\\
1.484	-0.124605\\
1.486	-0.124603\\
1.488	-0.1246\\
1.49	-0.124598\\
1.492	-0.124595\\
1.494	-0.124593\\
1.496	-0.124591\\
1.498	-0.124589\\
1.5	-0.124586\\
1.502	-0.124584\\
1.504	-0.124582\\
1.506	-0.12458\\
1.508	-0.124578\\
1.51	-0.124576\\
1.512	-0.124574\\
1.514	-0.124572\\
1.516	-0.12457\\
1.518	-0.124568\\
1.52	-0.124566\\
1.522	-0.124564\\
1.524	-0.124562\\
1.526	-0.124561\\
1.528	-0.124559\\
1.53	-0.124557\\
1.532	-0.124555\\
1.534	-0.124554\\
1.536	-0.124552\\
1.538	-0.12455\\
1.54	-0.124549\\
1.542	-0.124547\\
1.544	-0.124545\\
1.546	-0.124544\\
1.548	-0.124542\\
1.55	-0.124541\\
1.552	-0.124539\\
1.554	-0.124538\\
1.556	-0.124536\\
1.558	-0.124535\\
1.56	-0.124533\\
1.562	-0.124532\\
1.564	-0.12453\\
1.566	-0.124529\\
1.568	-0.124527\\
1.57	-0.124526\\
1.572	-0.124524\\
1.574	-0.124523\\
1.576	-0.124522\\
1.578	-0.12452\\
1.58	-0.124519\\
1.582	-0.124517\\
1.584	-0.124516\\
1.586	-0.124514\\
1.588	-0.124513\\
1.59	-0.124512\\
1.592	-0.12451\\
1.594	-0.124509\\
1.596	-0.124507\\
1.598	-0.124506\\
1.6	-0.124504\\
1.602	-0.124503\\
1.604	-0.124501\\
1.606	-0.1245\\
1.608	-0.124499\\
1.61	-0.124497\\
1.612	-0.124495\\
1.614	-0.124494\\
1.616	-0.124492\\
1.618	-0.124491\\
1.62	-0.124489\\
1.622	-0.124488\\
1.624	-0.124486\\
1.626	-0.124485\\
1.628	-0.124483\\
1.63	-0.124481\\
1.632	-0.12448\\
1.634	-0.124478\\
1.636	-0.124476\\
1.638	-0.124475\\
1.64	-0.124473\\
1.642	-0.124471\\
1.644	-0.12447\\
1.646	-0.124468\\
1.648	-0.124466\\
1.65	-0.124464\\
1.652	-0.124462\\
1.654	-0.124461\\
1.656	-0.124459\\
1.658	-0.124457\\
1.66	-0.124455\\
1.662	-0.124453\\
1.664	-0.124451\\
1.666	-0.124449\\
1.668	-0.124447\\
1.67	-0.124446\\
1.672	-0.124444\\
1.674	-0.124442\\
1.676	-0.12444\\
1.678	-0.124438\\
1.68	-0.124436\\
1.682	-0.124434\\
1.684	-0.124432\\
1.686	-0.12443\\
1.688	-0.124428\\
1.69	-0.124426\\
1.692	-0.124423\\
1.694	-0.124421\\
1.696	-0.124419\\
1.698	-0.124417\\
1.7	-0.124415\\
1.702	-0.124413\\
1.704	-0.124411\\
1.706	-0.124409\\
1.708	-0.124407\\
1.71	-0.124404\\
1.712	-0.124402\\
1.714	-0.1244\\
1.716	-0.124398\\
1.718	-0.124396\\
1.72	-0.124394\\
1.722	-0.124391\\
1.724	-0.124389\\
1.726	-0.124387\\
1.728	-0.124385\\
1.73	-0.124383\\
1.732	-0.124381\\
1.734	-0.124378\\
1.736	-0.124376\\
1.738	-0.124374\\
1.74	-0.124372\\
1.742	-0.12437\\
1.744	-0.124367\\
1.746	-0.124365\\
1.748	-0.124363\\
1.75	-0.124361\\
1.752	-0.124359\\
1.754	-0.124356\\
1.756	-0.124354\\
1.758	-0.124352\\
1.76	-0.12435\\
1.762	-0.124348\\
1.764	-0.124345\\
1.766	-0.124343\\
1.768	-0.124341\\
1.77	-0.124339\\
1.772	-0.124337\\
1.774	-0.124335\\
1.776	-0.124332\\
1.778	-0.12433\\
1.78	-0.124328\\
1.782	-0.124326\\
1.784	-0.124324\\
1.786	-0.124322\\
1.788	-0.12432\\
1.79	-0.124317\\
1.792	-0.124315\\
1.794	-0.124313\\
1.796	-0.124311\\
1.798	-0.124309\\
1.8	-0.124307\\
1.802	-0.124305\\
1.804	-0.124303\\
1.806	-0.124301\\
1.808	-0.124299\\
1.81	-0.124296\\
1.812	-0.124294\\
1.814	-0.124292\\
1.816	-0.12429\\
1.818	-0.124288\\
1.82	-0.124286\\
1.822	-0.124284\\
1.824	-0.124282\\
1.826	-0.12428\\
1.828	-0.124278\\
1.83	-0.124276\\
1.832	-0.124274\\
1.834	-0.124272\\
1.836	-0.12427\\
1.838	-0.124268\\
1.84	-0.124266\\
1.842	-0.124264\\
1.844	-0.124262\\
1.846	-0.12426\\
1.848	-0.124258\\
1.85	-0.124256\\
1.852	-0.124254\\
1.854	-0.124252\\
1.856	-0.12425\\
1.858	-0.124248\\
1.86	-0.124247\\
1.862	-0.124245\\
1.864	-0.124243\\
1.866	-0.124241\\
1.868	-0.124239\\
1.87	-0.124237\\
1.872	-0.124235\\
1.874	-0.124233\\
1.876	-0.124231\\
1.878	-0.124229\\
1.88	-0.124228\\
1.882	-0.124226\\
1.884	-0.124224\\
1.886	-0.124222\\
1.888	-0.12422\\
1.89	-0.124218\\
1.892	-0.124217\\
1.894	-0.124215\\
1.896	-0.124213\\
1.898	-0.124211\\
1.9	-0.124209\\
1.902	-0.124207\\
1.904	-0.124206\\
1.906	-0.124204\\
1.908	-0.124202\\
1.91	-0.1242\\
1.912	-0.124198\\
1.914	-0.124197\\
1.916	-0.124195\\
1.918	-0.124193\\
1.92	-0.124191\\
1.922	-0.12419\\
1.924	-0.124188\\
1.926	-0.124186\\
1.928	-0.124184\\
1.93	-0.124182\\
1.932	-0.124181\\
1.934	-0.124179\\
1.936	-0.124177\\
1.938	-0.124175\\
1.94	-0.124174\\
1.942	-0.124172\\
1.944	-0.12417\\
1.946	-0.124168\\
1.948	-0.124167\\
1.95	-0.124165\\
1.952	-0.124163\\
1.954	-0.124161\\
1.956	-0.12416\\
1.958	-0.124158\\
1.96	-0.124156\\
1.962	-0.124154\\
1.964	-0.124153\\
1.966	-0.124151\\
1.968	-0.124149\\
1.97	-0.124148\\
1.972	-0.124146\\
1.974	-0.124144\\
1.976	-0.124142\\
1.978	-0.124141\\
1.98	-0.124139\\
1.982	-0.124137\\
1.984	-0.124135\\
1.986	-0.124134\\
1.988	-0.124132\\
1.99	-0.12413\\
1.992	-0.124128\\
1.994	-0.124127\\
1.996	-0.124125\\
1.998	-0.124123\\
2	-0.124122\\
};
\addplot [color=mycolor14, line width=1.0pt, forget plot]
  table[row sep=crcr]{%
-0.124	0\\
-0.122	0\\
-0.12	0\\
-0.118	0\\
-0.116	0\\
-0.114	0\\
-0.112	0\\
-0.11	0\\
-0.108	0\\
-0.106	0\\
-0.104	0\\
-0.102	0\\
-0.1	0\\
-0.098	0\\
-0.096	0\\
-0.094	0\\
-0.092	0\\
-0.09	0\\
-0.088	0\\
-0.086	0\\
-0.084	0\\
-0.082	0\\
-0.08	0\\
-0.078	0\\
-0.076	0\\
-0.074	0\\
-0.072	0\\
-0.07	0\\
-0.068	0\\
-0.066	0\\
-0.064	0\\
-0.062	0\\
-0.06	0\\
-0.058	0\\
-0.056	0\\
-0.054	0\\
-0.052	0\\
-0.05	0\\
-0.048	0\\
-0.046	0\\
-0.044	0\\
-0.042	0\\
-0.04	0\\
-0.038	0\\
-0.036	0\\
-0.034	0\\
-0.032	0\\
-0.03	0\\
-0.028	0\\
-0.026	0\\
-0.024	0\\
-0.022	0\\
-0.02	0\\
-0.018	0\\
-0.016	0\\
-0.014	0\\
-0.012	0\\
-0.01	0\\
-0.008	0\\
-0.006	0\\
-0.004	0\\
-0.002	0\\
0	0\\
0.002	-0.0237246\\
0.004	-0.0325438\\
0.006	-0.0344632\\
0.008	-0.0334073\\
0.01	-0.0312697\\
0.012	-0.0289372\\
0.014	-0.0268019\\
0.016	-0.0250155\\
0.018	-0.0236159\\
0.02	-0.0225897\\
0.022	-0.0219023\\
0.024	-0.021513\\
0.026	-0.021381\\
0.028	-0.0214689\\
0.03	-0.0217436\\
0.032	-0.0221758\\
0.034	-0.0227405\\
0.036	-0.0234164\\
0.038	-0.0241852\\
0.04	-0.0250311\\
0.042	-0.0259411\\
0.044	-0.0269039\\
0.046	-0.0279099\\
0.048	-0.0289511\\
0.05	-0.0300207\\
0.052	-0.0311129\\
0.054	-0.0322228\\
0.056	-0.0333462\\
0.058	-0.0344797\\
0.06	-0.0356203\\
0.062	-0.0367654\\
0.064	-0.037913\\
0.066	-0.0390613\\
0.068	-0.0402088\\
0.07	-0.0413542\\
0.072	-0.0424965\\
0.074	-0.0436346\\
0.076	-0.0447679\\
0.078	-0.0458956\\
0.08	-0.0470173\\
0.082	-0.0481324\\
0.084	-0.0492406\\
0.086	-0.0503414\\
0.088	-0.0514346\\
0.09	-0.0525199\\
0.092	-0.053597\\
0.094	-0.0546658\\
0.096	-0.0557261\\
0.098	-0.0567776\\
0.1	-0.0578203\\
0.102	-0.058854\\
0.104	-0.0598785\\
0.106	-0.0608938\\
0.108	-0.0618996\\
0.11	-0.0628959\\
0.112	-0.0638825\\
0.114	-0.0648594\\
0.116	-0.0658264\\
0.118	-0.0667834\\
0.12	-0.0677303\\
0.122	-0.068667\\
0.124	-0.0695934\\
0.126	-0.0705094\\
0.128	-0.0714148\\
0.13	-0.0723095\\
0.132	-0.0731935\\
0.134	-0.0740667\\
0.136	-0.0749289\\
0.138	-0.07578\\
0.14	-0.0766199\\
0.142	-0.0774486\\
0.144	-0.0782659\\
0.146	-0.0790717\\
0.148	-0.079866\\
0.15	-0.0806486\\
0.152	-0.0814196\\
0.154	-0.0821786\\
0.156	-0.0829258\\
0.158	-0.0836611\\
0.16	-0.0843843\\
0.162	-0.0850953\\
0.164	-0.0857942\\
0.166	-0.0864809\\
0.168	-0.0871553\\
0.17	-0.0878173\\
0.172	-0.088467\\
0.174	-0.0891042\\
0.176	-0.089729\\
0.178	-0.0903413\\
0.18	-0.0909411\\
0.182	-0.0915284\\
0.184	-0.0921032\\
0.186	-0.0926655\\
0.188	-0.0932153\\
0.19	-0.0937526\\
0.192	-0.0942774\\
0.194	-0.0947897\\
0.196	-0.0952896\\
0.198	-0.0957772\\
0.2	-0.0962524\\
0.202	-0.0967154\\
0.204	-0.0971661\\
0.206	-0.0976047\\
0.208	-0.0980312\\
0.21	-0.0984457\\
0.212	-0.0988483\\
0.214	-0.0992392\\
0.216	-0.0996183\\
0.218	-0.0999859\\
0.22	-0.100342\\
0.222	-0.100687\\
0.224	-0.10102\\
0.226	-0.101343\\
0.228	-0.101654\\
0.23	-0.101955\\
0.232	-0.102245\\
0.234	-0.102525\\
0.236	-0.102795\\
0.238	-0.103054\\
0.24	-0.103304\\
0.242	-0.103543\\
0.244	-0.103774\\
0.246	-0.103994\\
0.248	-0.104206\\
0.25	-0.104409\\
0.252	-0.104603\\
0.254	-0.104789\\
0.256	-0.104967\\
0.258	-0.105136\\
0.26	-0.105298\\
0.262	-0.105452\\
0.264	-0.105598\\
0.266	-0.105738\\
0.268	-0.105871\\
0.27	-0.105997\\
0.272	-0.106117\\
0.274	-0.106231\\
0.276	-0.106338\\
0.278	-0.106441\\
0.28	-0.106538\\
0.282	-0.106629\\
0.284	-0.106716\\
0.286	-0.106799\\
0.288	-0.106877\\
0.29	-0.106951\\
0.292	-0.107021\\
0.294	-0.107088\\
0.296	-0.107151\\
0.298	-0.107211\\
0.3	-0.107268\\
0.302	-0.107323\\
0.304	-0.107375\\
0.306	-0.107426\\
0.308	-0.107474\\
0.31	-0.107521\\
0.312	-0.107566\\
0.314	-0.107611\\
0.316	-0.107654\\
0.318	-0.107696\\
0.32	-0.107738\\
0.322	-0.10778\\
0.324	-0.107822\\
0.326	-0.107863\\
0.328	-0.107905\\
0.33	-0.107948\\
0.332	-0.107991\\
0.334	-0.108035\\
0.336	-0.10808\\
0.338	-0.108127\\
0.34	-0.108174\\
0.342	-0.108223\\
0.344	-0.108274\\
0.346	-0.108327\\
0.348	-0.108381\\
0.35	-0.108438\\
0.352	-0.108496\\
0.354	-0.108557\\
0.356	-0.108621\\
0.358	-0.108687\\
0.36	-0.108756\\
0.362	-0.108827\\
0.364	-0.108901\\
0.366	-0.108978\\
0.368	-0.109058\\
0.37	-0.109141\\
0.372	-0.109226\\
0.374	-0.109315\\
0.376	-0.109408\\
0.378	-0.109503\\
0.38	-0.109601\\
0.382	-0.109703\\
0.384	-0.109808\\
0.386	-0.109916\\
0.388	-0.110027\\
0.39	-0.110142\\
0.392	-0.11026\\
0.394	-0.110381\\
0.396	-0.110505\\
0.398	-0.110632\\
0.4	-0.110763\\
0.402	-0.110897\\
0.404	-0.111033\\
0.406	-0.111173\\
0.408	-0.111315\\
0.41	-0.111461\\
0.412	-0.111609\\
0.414	-0.11176\\
0.416	-0.111914\\
0.418	-0.11207\\
0.42	-0.112229\\
0.422	-0.11239\\
0.424	-0.112554\\
0.426	-0.11272\\
0.428	-0.112888\\
0.43	-0.113058\\
0.432	-0.11323\\
0.434	-0.113404\\
0.436	-0.11358\\
0.438	-0.113757\\
0.44	-0.113937\\
0.442	-0.114117\\
0.444	-0.114299\\
0.446	-0.114482\\
0.448	-0.114667\\
0.45	-0.114852\\
0.452	-0.115038\\
0.454	-0.115226\\
0.456	-0.115414\\
0.458	-0.115602\\
0.46	-0.115791\\
0.462	-0.11598\\
0.464	-0.11617\\
0.466	-0.11636\\
0.468	-0.116549\\
0.47	-0.116739\\
0.472	-0.116929\\
0.474	-0.117118\\
0.476	-0.117307\\
0.478	-0.117495\\
0.48	-0.117683\\
0.482	-0.11787\\
0.484	-0.118057\\
0.486	-0.118242\\
0.488	-0.118426\\
0.49	-0.11861\\
0.492	-0.118792\\
0.494	-0.118973\\
0.496	-0.119153\\
0.498	-0.119331\\
0.5	-0.119508\\
0.502	-0.119683\\
0.504	-0.119857\\
0.506	-0.120029\\
0.508	-0.120199\\
0.51	-0.120367\\
0.512	-0.120533\\
0.514	-0.120697\\
0.516	-0.12086\\
0.518	-0.12102\\
0.52	-0.121178\\
0.522	-0.121334\\
0.524	-0.121487\\
0.526	-0.121639\\
0.528	-0.121788\\
0.53	-0.121934\\
0.532	-0.122079\\
0.534	-0.12222\\
0.536	-0.12236\\
0.538	-0.122496\\
0.54	-0.122631\\
0.542	-0.122762\\
0.544	-0.122892\\
0.546	-0.123018\\
0.548	-0.123142\\
0.55	-0.123264\\
0.552	-0.123382\\
0.554	-0.123499\\
0.556	-0.123612\\
0.558	-0.123723\\
0.56	-0.123831\\
0.562	-0.123937\\
0.564	-0.12404\\
0.566	-0.12414\\
0.568	-0.124238\\
0.57	-0.124333\\
0.572	-0.124426\\
0.574	-0.124515\\
0.576	-0.124603\\
0.578	-0.124688\\
0.58	-0.12477\\
0.582	-0.12485\\
0.584	-0.124927\\
0.586	-0.125002\\
0.588	-0.125074\\
0.59	-0.125144\\
0.592	-0.125211\\
0.594	-0.125276\\
0.596	-0.125339\\
0.598	-0.1254\\
0.6	-0.125458\\
0.602	-0.125514\\
0.604	-0.125567\\
0.606	-0.125619\\
0.608	-0.125668\\
0.61	-0.125715\\
0.612	-0.12576\\
0.614	-0.125803\\
0.616	-0.125844\\
0.618	-0.125883\\
0.62	-0.12592\\
0.622	-0.125955\\
0.624	-0.125988\\
0.626	-0.12602\\
0.628	-0.126049\\
0.63	-0.126077\\
0.632	-0.126103\\
0.634	-0.126128\\
0.636	-0.12615\\
0.638	-0.126172\\
0.64	-0.126191\\
0.642	-0.126209\\
0.644	-0.126226\\
0.646	-0.126241\\
0.648	-0.126254\\
0.65	-0.126267\\
0.652	-0.126277\\
0.654	-0.126287\\
0.656	-0.126295\\
0.658	-0.126302\\
0.66	-0.126308\\
0.662	-0.126313\\
0.664	-0.126317\\
0.666	-0.126319\\
0.668	-0.126321\\
0.67	-0.126322\\
0.672	-0.126321\\
0.674	-0.12632\\
0.676	-0.126318\\
0.678	-0.126314\\
0.68	-0.126311\\
0.682	-0.126306\\
0.684	-0.1263\\
0.686	-0.126294\\
0.688	-0.126287\\
0.69	-0.12628\\
0.692	-0.126272\\
0.694	-0.126263\\
0.696	-0.126254\\
0.698	-0.126244\\
0.7	-0.126234\\
0.702	-0.126223\\
0.704	-0.126212\\
0.706	-0.126201\\
0.708	-0.126189\\
0.71	-0.126177\\
0.712	-0.126164\\
0.714	-0.126152\\
0.716	-0.126139\\
0.718	-0.126126\\
0.72	-0.126112\\
0.722	-0.126099\\
0.724	-0.126085\\
0.726	-0.126071\\
0.728	-0.126057\\
0.73	-0.126043\\
0.732	-0.126029\\
0.734	-0.126015\\
0.736	-0.126001\\
0.738	-0.125987\\
0.74	-0.125973\\
0.742	-0.125959\\
0.744	-0.125945\\
0.746	-0.125931\\
0.748	-0.125918\\
0.75	-0.125904\\
0.752	-0.125891\\
0.754	-0.125878\\
0.756	-0.125865\\
0.758	-0.125852\\
0.76	-0.125839\\
0.762	-0.125827\\
0.764	-0.125814\\
0.766	-0.125802\\
0.768	-0.125791\\
0.77	-0.125779\\
0.772	-0.125768\\
0.774	-0.125757\\
0.776	-0.125747\\
0.778	-0.125736\\
0.78	-0.125726\\
0.782	-0.125717\\
0.784	-0.125707\\
0.786	-0.125698\\
0.788	-0.125689\\
0.79	-0.125681\\
0.792	-0.125673\\
0.794	-0.125665\\
0.796	-0.125657\\
0.798	-0.12565\\
0.8	-0.125643\\
0.802	-0.125637\\
0.804	-0.12563\\
0.806	-0.125624\\
0.808	-0.125619\\
0.81	-0.125614\\
0.812	-0.125609\\
0.814	-0.125604\\
0.816	-0.125599\\
0.818	-0.125595\\
0.82	-0.125591\\
0.822	-0.125588\\
0.824	-0.125584\\
0.826	-0.125581\\
0.828	-0.125579\\
0.83	-0.125576\\
0.832	-0.125574\\
0.834	-0.125571\\
0.836	-0.12557\\
0.838	-0.125568\\
0.84	-0.125566\\
0.842	-0.125565\\
0.844	-0.125564\\
0.846	-0.125563\\
0.848	-0.125562\\
0.85	-0.125561\\
0.852	-0.12556\\
0.854	-0.12556\\
0.856	-0.12556\\
0.858	-0.125559\\
0.86	-0.125559\\
0.862	-0.125559\\
0.864	-0.125559\\
0.866	-0.125559\\
0.868	-0.125559\\
0.87	-0.125559\\
0.872	-0.125558\\
0.874	-0.125558\\
0.876	-0.125558\\
0.878	-0.125558\\
0.88	-0.125558\\
0.882	-0.125558\\
0.884	-0.125558\\
0.886	-0.125558\\
0.888	-0.125557\\
0.89	-0.125557\\
0.892	-0.125556\\
0.894	-0.125556\\
0.896	-0.125555\\
0.898	-0.125554\\
0.9	-0.125553\\
0.902	-0.125552\\
0.904	-0.125551\\
0.906	-0.125549\\
0.908	-0.125548\\
0.91	-0.125546\\
0.912	-0.125544\\
0.914	-0.125542\\
0.916	-0.12554\\
0.918	-0.125538\\
0.92	-0.125535\\
0.922	-0.125532\\
0.924	-0.125529\\
0.926	-0.125526\\
0.928	-0.125523\\
0.93	-0.125519\\
0.932	-0.125515\\
0.934	-0.125511\\
0.936	-0.125507\\
0.938	-0.125503\\
0.94	-0.125498\\
0.942	-0.125494\\
0.944	-0.125489\\
0.946	-0.125484\\
0.948	-0.125478\\
0.95	-0.125473\\
0.952	-0.125467\\
0.954	-0.125461\\
0.956	-0.125455\\
0.958	-0.125449\\
0.96	-0.125442\\
0.962	-0.125436\\
0.964	-0.125429\\
0.966	-0.125422\\
0.968	-0.125415\\
0.97	-0.125407\\
0.972	-0.1254\\
0.974	-0.125393\\
0.976	-0.125385\\
0.978	-0.125377\\
0.98	-0.125369\\
0.982	-0.125361\\
0.984	-0.125353\\
0.986	-0.125345\\
0.988	-0.125336\\
0.99	-0.125328\\
0.992	-0.125319\\
0.994	-0.125311\\
0.996	-0.125302\\
0.998	-0.125293\\
1	-0.125284\\
1.002	-0.125276\\
1.004	-0.125267\\
1.006	-0.125258\\
1.008	-0.125249\\
1.01	-0.12524\\
1.012	-0.125231\\
1.014	-0.125222\\
1.016	-0.125213\\
1.018	-0.125204\\
1.02	-0.125195\\
1.022	-0.125186\\
1.024	-0.125177\\
1.026	-0.125169\\
1.028	-0.12516\\
1.03	-0.125151\\
1.032	-0.125142\\
1.034	-0.125134\\
1.036	-0.125125\\
1.038	-0.125117\\
1.04	-0.125108\\
1.042	-0.1251\\
1.044	-0.125092\\
1.046	-0.125083\\
1.048	-0.125075\\
1.05	-0.125067\\
1.052	-0.125059\\
1.054	-0.125052\\
1.056	-0.125044\\
1.058	-0.125037\\
1.06	-0.125029\\
1.062	-0.125022\\
1.064	-0.125015\\
1.066	-0.125008\\
1.068	-0.125001\\
1.07	-0.124994\\
1.072	-0.124987\\
1.074	-0.124981\\
1.076	-0.124974\\
1.078	-0.124968\\
1.08	-0.124962\\
1.082	-0.124956\\
1.084	-0.12495\\
1.086	-0.124945\\
1.088	-0.124939\\
1.09	-0.124934\\
1.092	-0.124929\\
1.094	-0.124923\\
1.096	-0.124918\\
1.098	-0.124914\\
1.1	-0.124909\\
1.102	-0.124904\\
1.104	-0.1249\\
1.106	-0.124896\\
1.108	-0.124892\\
1.11	-0.124888\\
1.112	-0.124884\\
1.114	-0.12488\\
1.116	-0.124877\\
1.118	-0.124873\\
1.12	-0.12487\\
1.122	-0.124867\\
1.124	-0.124864\\
1.126	-0.124861\\
1.128	-0.124858\\
1.13	-0.124855\\
1.132	-0.124852\\
1.134	-0.12485\\
1.136	-0.124848\\
1.138	-0.124845\\
1.14	-0.124843\\
1.142	-0.124841\\
1.144	-0.124839\\
1.146	-0.124837\\
1.148	-0.124836\\
1.15	-0.124834\\
1.152	-0.124832\\
1.154	-0.124831\\
1.156	-0.12483\\
1.158	-0.124828\\
1.16	-0.124827\\
1.162	-0.124826\\
1.164	-0.124825\\
1.166	-0.124824\\
1.168	-0.124823\\
1.17	-0.124822\\
1.172	-0.124821\\
1.174	-0.12482\\
1.176	-0.12482\\
1.178	-0.124819\\
1.18	-0.124818\\
1.182	-0.124818\\
1.184	-0.124817\\
1.186	-0.124817\\
1.188	-0.124817\\
1.19	-0.124816\\
1.192	-0.124816\\
1.194	-0.124816\\
1.196	-0.124816\\
1.198	-0.124815\\
1.2	-0.124815\\
1.202	-0.124815\\
1.204	-0.124815\\
1.206	-0.124815\\
1.208	-0.124815\\
1.21	-0.124815\\
1.212	-0.124815\\
1.214	-0.124815\\
1.216	-0.124815\\
1.218	-0.124815\\
1.22	-0.124815\\
1.222	-0.124815\\
1.224	-0.124815\\
1.226	-0.124815\\
1.228	-0.124815\\
1.23	-0.124815\\
1.232	-0.124815\\
1.234	-0.124815\\
1.236	-0.124815\\
1.238	-0.124815\\
1.24	-0.124815\\
1.242	-0.124815\\
1.244	-0.124815\\
1.246	-0.124815\\
1.248	-0.124815\\
1.25	-0.124815\\
1.252	-0.124815\\
1.254	-0.124815\\
1.256	-0.124815\\
1.258	-0.124815\\
1.26	-0.124815\\
1.262	-0.124814\\
1.264	-0.124814\\
1.266	-0.124814\\
1.268	-0.124814\\
1.27	-0.124814\\
1.272	-0.124813\\
1.274	-0.124813\\
1.276	-0.124813\\
1.278	-0.124812\\
1.28	-0.124812\\
1.282	-0.124811\\
1.284	-0.124811\\
1.286	-0.12481\\
1.288	-0.12481\\
1.29	-0.124809\\
1.292	-0.124809\\
1.294	-0.124808\\
1.296	-0.124807\\
1.298	-0.124807\\
1.3	-0.124806\\
1.302	-0.124805\\
1.304	-0.124804\\
1.306	-0.124804\\
1.308	-0.124803\\
1.31	-0.124802\\
1.312	-0.124801\\
1.314	-0.1248\\
1.316	-0.124799\\
1.318	-0.124798\\
1.32	-0.124796\\
1.322	-0.124795\\
1.324	-0.124794\\
1.326	-0.124793\\
1.328	-0.124791\\
1.33	-0.12479\\
1.332	-0.124789\\
1.334	-0.124787\\
1.336	-0.124786\\
1.338	-0.124784\\
1.34	-0.124783\\
1.342	-0.124781\\
1.344	-0.12478\\
1.346	-0.124778\\
1.348	-0.124776\\
1.35	-0.124775\\
1.352	-0.124773\\
1.354	-0.124771\\
1.356	-0.124769\\
1.358	-0.124767\\
1.36	-0.124765\\
1.362	-0.124764\\
1.364	-0.124762\\
1.366	-0.12476\\
1.368	-0.124758\\
1.37	-0.124756\\
1.372	-0.124753\\
1.374	-0.124751\\
1.376	-0.124749\\
1.378	-0.124747\\
1.38	-0.124745\\
1.382	-0.124743\\
1.384	-0.124741\\
1.386	-0.124738\\
1.388	-0.124736\\
1.39	-0.124734\\
1.392	-0.124731\\
1.394	-0.124729\\
1.396	-0.124727\\
1.398	-0.124724\\
1.4	-0.124722\\
1.402	-0.12472\\
1.404	-0.124717\\
1.406	-0.124715\\
1.408	-0.124713\\
1.41	-0.12471\\
1.412	-0.124708\\
1.414	-0.124705\\
1.416	-0.124703\\
1.418	-0.1247\\
1.42	-0.124698\\
1.422	-0.124696\\
1.424	-0.124693\\
1.426	-0.124691\\
1.428	-0.124688\\
1.43	-0.124686\\
1.432	-0.124683\\
1.434	-0.124681\\
1.436	-0.124679\\
1.438	-0.124676\\
1.44	-0.124674\\
1.442	-0.124671\\
1.444	-0.124669\\
1.446	-0.124667\\
1.448	-0.124664\\
1.45	-0.124662\\
1.452	-0.124659\\
1.454	-0.124657\\
1.456	-0.124655\\
1.458	-0.124652\\
1.46	-0.12465\\
1.462	-0.124648\\
1.464	-0.124646\\
1.466	-0.124643\\
1.468	-0.124641\\
1.47	-0.124639\\
1.472	-0.124637\\
1.474	-0.124634\\
1.476	-0.124632\\
1.478	-0.12463\\
1.48	-0.124628\\
1.482	-0.124626\\
1.484	-0.124624\\
1.486	-0.124621\\
1.488	-0.124619\\
1.49	-0.124617\\
1.492	-0.124615\\
1.494	-0.124613\\
1.496	-0.124611\\
1.498	-0.124609\\
1.5	-0.124607\\
1.502	-0.124605\\
1.504	-0.124603\\
1.506	-0.124601\\
1.508	-0.124599\\
1.51	-0.124597\\
1.512	-0.124595\\
1.514	-0.124594\\
1.516	-0.124592\\
1.518	-0.12459\\
1.52	-0.124588\\
1.522	-0.124586\\
1.524	-0.124584\\
1.526	-0.124582\\
1.528	-0.124581\\
1.53	-0.124579\\
1.532	-0.124577\\
1.534	-0.124575\\
1.536	-0.124573\\
1.538	-0.124572\\
1.54	-0.12457\\
1.542	-0.124568\\
1.544	-0.124566\\
1.546	-0.124565\\
1.548	-0.124563\\
1.55	-0.124561\\
1.552	-0.124559\\
1.554	-0.124558\\
1.556	-0.124556\\
1.558	-0.124554\\
1.56	-0.124552\\
1.562	-0.124551\\
1.564	-0.124549\\
1.566	-0.124547\\
1.568	-0.124546\\
1.57	-0.124544\\
1.572	-0.124542\\
1.574	-0.12454\\
1.576	-0.124539\\
1.578	-0.124537\\
1.58	-0.124535\\
1.582	-0.124534\\
1.584	-0.124532\\
1.586	-0.12453\\
1.588	-0.124528\\
1.59	-0.124527\\
1.592	-0.124525\\
1.594	-0.124523\\
1.596	-0.124521\\
1.598	-0.12452\\
1.6	-0.124518\\
1.602	-0.124516\\
1.604	-0.124514\\
1.606	-0.124513\\
1.608	-0.124511\\
1.61	-0.124509\\
1.612	-0.124507\\
1.614	-0.124505\\
1.616	-0.124504\\
1.618	-0.124502\\
1.62	-0.1245\\
1.622	-0.124498\\
1.624	-0.124496\\
1.626	-0.124494\\
1.628	-0.124492\\
1.63	-0.124491\\
1.632	-0.124489\\
1.634	-0.124487\\
1.636	-0.124485\\
1.638	-0.124483\\
1.64	-0.124481\\
1.642	-0.124479\\
1.644	-0.124477\\
1.646	-0.124475\\
1.648	-0.124473\\
1.65	-0.124472\\
1.652	-0.12447\\
1.654	-0.124468\\
1.656	-0.124466\\
1.658	-0.124464\\
1.66	-0.124462\\
1.662	-0.12446\\
1.664	-0.124458\\
1.666	-0.124456\\
1.668	-0.124454\\
1.67	-0.124452\\
1.672	-0.12445\\
1.674	-0.124448\\
1.676	-0.124446\\
1.678	-0.124444\\
1.68	-0.124442\\
1.682	-0.12444\\
1.684	-0.124438\\
1.686	-0.124436\\
1.688	-0.124434\\
1.69	-0.124432\\
1.692	-0.124429\\
1.694	-0.124427\\
1.696	-0.124425\\
1.698	-0.124423\\
1.7	-0.124421\\
1.702	-0.124419\\
1.704	-0.124417\\
1.706	-0.124415\\
1.708	-0.124413\\
1.71	-0.124411\\
1.712	-0.124409\\
1.714	-0.124407\\
1.716	-0.124405\\
1.718	-0.124403\\
1.72	-0.124401\\
1.722	-0.124398\\
1.724	-0.124396\\
1.726	-0.124394\\
1.728	-0.124392\\
1.73	-0.12439\\
1.732	-0.124388\\
1.734	-0.124386\\
1.736	-0.124384\\
1.738	-0.124382\\
1.74	-0.12438\\
1.742	-0.124378\\
1.744	-0.124376\\
1.746	-0.124373\\
1.748	-0.124371\\
1.75	-0.124369\\
1.752	-0.124367\\
1.754	-0.124365\\
1.756	-0.124363\\
1.758	-0.124361\\
1.76	-0.124359\\
1.762	-0.124357\\
1.764	-0.124355\\
1.766	-0.124353\\
1.768	-0.124351\\
1.77	-0.124349\\
1.772	-0.124347\\
1.774	-0.124345\\
1.776	-0.124342\\
1.778	-0.12434\\
1.78	-0.124338\\
1.782	-0.124336\\
1.784	-0.124334\\
1.786	-0.124332\\
1.788	-0.12433\\
1.79	-0.124328\\
1.792	-0.124326\\
1.794	-0.124324\\
1.796	-0.124322\\
1.798	-0.12432\\
1.8	-0.124318\\
1.802	-0.124316\\
1.804	-0.124314\\
1.806	-0.124312\\
1.808	-0.12431\\
1.81	-0.124308\\
1.812	-0.124306\\
1.814	-0.124304\\
1.816	-0.124302\\
1.818	-0.1243\\
1.82	-0.124298\\
1.822	-0.124296\\
1.824	-0.124294\\
1.826	-0.124292\\
1.828	-0.12429\\
1.83	-0.124288\\
1.832	-0.124286\\
1.834	-0.124284\\
1.836	-0.124282\\
1.838	-0.12428\\
1.84	-0.124278\\
1.842	-0.124277\\
1.844	-0.124275\\
1.846	-0.124273\\
1.848	-0.124271\\
1.85	-0.124269\\
1.852	-0.124267\\
1.854	-0.124265\\
1.856	-0.124263\\
1.858	-0.124261\\
1.86	-0.124259\\
1.862	-0.124257\\
1.864	-0.124255\\
1.866	-0.124253\\
1.868	-0.124252\\
1.87	-0.12425\\
1.872	-0.124248\\
1.874	-0.124246\\
1.876	-0.124244\\
1.878	-0.124242\\
1.88	-0.12424\\
1.882	-0.124238\\
1.884	-0.124237\\
1.886	-0.124235\\
1.888	-0.124233\\
1.89	-0.124231\\
1.892	-0.124229\\
1.894	-0.124227\\
1.896	-0.124225\\
1.898	-0.124224\\
1.9	-0.124222\\
1.902	-0.12422\\
1.904	-0.124218\\
1.906	-0.124216\\
1.908	-0.124214\\
1.91	-0.124213\\
1.912	-0.124211\\
1.914	-0.124209\\
1.916	-0.124207\\
1.918	-0.124205\\
1.92	-0.124203\\
1.922	-0.124202\\
1.924	-0.1242\\
1.926	-0.124198\\
1.928	-0.124196\\
1.93	-0.124194\\
1.932	-0.124193\\
1.934	-0.124191\\
1.936	-0.124189\\
1.938	-0.124187\\
1.94	-0.124185\\
1.942	-0.124184\\
1.944	-0.124182\\
1.946	-0.12418\\
1.948	-0.124178\\
1.95	-0.124176\\
1.952	-0.124175\\
1.954	-0.124173\\
1.956	-0.124171\\
1.958	-0.124169\\
1.96	-0.124167\\
1.962	-0.124166\\
1.964	-0.124164\\
1.966	-0.124162\\
1.968	-0.12416\\
1.97	-0.124159\\
1.972	-0.124157\\
1.974	-0.124155\\
1.976	-0.124153\\
1.978	-0.124151\\
1.98	-0.12415\\
1.982	-0.124148\\
1.984	-0.124146\\
1.986	-0.124144\\
1.988	-0.124142\\
1.99	-0.124141\\
1.992	-0.124139\\
1.994	-0.124137\\
1.996	-0.124135\\
1.998	-0.124134\\
2	-0.124132\\
};
\addplot [color=mycolor15, line width=1.0pt, forget plot]
  table[row sep=crcr]{%
-0.124	0\\
-0.122	0\\
-0.12	0\\
-0.118	0\\
-0.116	0\\
-0.114	0\\
-0.112	0\\
-0.11	0\\
-0.108	0\\
-0.106	0\\
-0.104	0\\
-0.102	0\\
-0.1	0\\
-0.098	0\\
-0.096	0\\
-0.094	0\\
-0.092	0\\
-0.09	0\\
-0.088	0\\
-0.086	0\\
-0.084	0\\
-0.082	0\\
-0.08	0\\
-0.078	0\\
-0.076	0\\
-0.074	0\\
-0.072	0\\
-0.07	0\\
-0.068	0\\
-0.066	0\\
-0.064	0\\
-0.062	0\\
-0.06	0\\
-0.058	0\\
-0.056	0\\
-0.054	0\\
-0.052	0\\
-0.05	0\\
-0.048	0\\
-0.046	0\\
-0.044	0\\
-0.042	0\\
-0.04	0\\
-0.038	0\\
-0.036	0\\
-0.034	0\\
-0.032	0\\
-0.03	0\\
-0.028	0\\
-0.026	0\\
-0.024	0\\
-0.022	0\\
-0.02	0\\
-0.018	0\\
-0.016	0\\
-0.014	0\\
-0.012	0\\
-0.01	0\\
-0.008	0\\
-0.006	0\\
-0.004	0\\
-0.002	0\\
0	0\\
0.002	-0.000422792\\
0.004	-0.000849858\\
0.006	-0.00127952\\
0.008	-0.00171176\\
0.01	-0.00214736\\
0.012	-0.00258748\\
0.014	-0.0030334\\
0.016	-0.00348644\\
0.018	-0.00394793\\
0.02	-0.00441913\\
0.022	-0.00490126\\
0.024	-0.00539547\\
0.026	-0.00590285\\
0.028	-0.00642444\\
0.03	-0.00696122\\
0.032	-0.00751408\\
0.034	-0.00808389\\
0.036	-0.00867143\\
0.038	-0.00927744\\
0.04	-0.0099026\\
0.042	-0.0105475\\
0.044	-0.0112128\\
0.046	-0.0118989\\
0.048	-0.0126063\\
0.05	-0.0133355\\
0.052	-0.0140867\\
0.054	-0.0148603\\
0.056	-0.0156565\\
0.058	-0.0164755\\
0.06	-0.0173176\\
0.062	-0.0181827\\
0.064	-0.019071\\
0.066	-0.0199825\\
0.068	-0.0209171\\
0.07	-0.0218748\\
0.072	-0.0228554\\
0.074	-0.0238589\\
0.076	-0.024885\\
0.078	-0.0259334\\
0.08	-0.027004\\
0.082	-0.0280963\\
0.084	-0.02921\\
0.086	-0.0303448\\
0.088	-0.0315002\\
0.09	-0.0326758\\
0.092	-0.033871\\
0.094	-0.0350854\\
0.096	-0.0363184\\
0.098	-0.0375694\\
0.1	-0.0388379\\
0.102	-0.0401231\\
0.104	-0.0414245\\
0.106	-0.0427413\\
0.108	-0.0440729\\
0.11	-0.0454185\\
0.112	-0.0467774\\
0.114	-0.0481489\\
0.116	-0.049532\\
0.118	-0.0509261\\
0.12	-0.0523304\\
0.122	-0.0537439\\
0.124	-0.0551659\\
0.126	-0.0565955\\
0.128	-0.0580318\\
0.13	-0.0594741\\
0.132	-0.0609213\\
0.134	-0.0623727\\
0.136	-0.0638273\\
0.138	-0.0652842\\
0.14	-0.0667426\\
0.142	-0.0682016\\
0.144	-0.0696602\\
0.146	-0.0711176\\
0.148	-0.0725729\\
0.15	-0.0740252\\
0.152	-0.0754735\\
0.154	-0.0769172\\
0.156	-0.0783551\\
0.158	-0.0797866\\
0.16	-0.0812107\\
0.162	-0.0826267\\
0.164	-0.0840336\\
0.166	-0.0854307\\
0.168	-0.0868172\\
0.17	-0.0881922\\
0.172	-0.089555\\
0.174	-0.0909048\\
0.176	-0.092241\\
0.178	-0.0935627\\
0.18	-0.0948693\\
0.182	-0.09616\\
0.184	-0.0974343\\
0.186	-0.0986914\\
0.188	-0.0999308\\
0.19	-0.101152\\
0.192	-0.102354\\
0.194	-0.103536\\
0.196	-0.104699\\
0.198	-0.105841\\
0.2	-0.106962\\
0.202	-0.108061\\
0.204	-0.109139\\
0.206	-0.110194\\
0.208	-0.111227\\
0.21	-0.112236\\
0.212	-0.113222\\
0.214	-0.114185\\
0.216	-0.115123\\
0.218	-0.116038\\
0.22	-0.116927\\
0.222	-0.117792\\
0.224	-0.118633\\
0.226	-0.119448\\
0.228	-0.120238\\
0.23	-0.121003\\
0.232	-0.121743\\
0.234	-0.122458\\
0.236	-0.123147\\
0.238	-0.123812\\
0.24	-0.12445\\
0.242	-0.125064\\
0.244	-0.125653\\
0.246	-0.126216\\
0.248	-0.126755\\
0.25	-0.127269\\
0.252	-0.127759\\
0.254	-0.128225\\
0.256	-0.128666\\
0.258	-0.129084\\
0.26	-0.129478\\
0.262	-0.129848\\
0.264	-0.130196\\
0.266	-0.130521\\
0.268	-0.130824\\
0.27	-0.131105\\
0.272	-0.131364\\
0.274	-0.131602\\
0.276	-0.13182\\
0.278	-0.132017\\
0.28	-0.132194\\
0.282	-0.132352\\
0.284	-0.13249\\
0.286	-0.13261\\
0.288	-0.132713\\
0.29	-0.132797\\
0.292	-0.132864\\
0.294	-0.132915\\
0.296	-0.13295\\
0.298	-0.132969\\
0.3	-0.132973\\
0.302	-0.132963\\
0.304	-0.132939\\
0.306	-0.132901\\
0.308	-0.13285\\
0.31	-0.132787\\
0.312	-0.132712\\
0.314	-0.132626\\
0.316	-0.13253\\
0.318	-0.132423\\
0.32	-0.132306\\
0.322	-0.13218\\
0.324	-0.132046\\
0.326	-0.131904\\
0.328	-0.131754\\
0.33	-0.131597\\
0.332	-0.131434\\
0.334	-0.131265\\
0.336	-0.13109\\
0.338	-0.130911\\
0.34	-0.130726\\
0.342	-0.130538\\
0.344	-0.130347\\
0.346	-0.130152\\
0.348	-0.129955\\
0.35	-0.129755\\
0.352	-0.129554\\
0.354	-0.129351\\
0.356	-0.129148\\
0.358	-0.128944\\
0.36	-0.12874\\
0.362	-0.128537\\
0.364	-0.128334\\
0.366	-0.128131\\
0.368	-0.127931\\
0.37	-0.127732\\
0.372	-0.127535\\
0.374	-0.12734\\
0.376	-0.127148\\
0.378	-0.126959\\
0.38	-0.126772\\
0.382	-0.12659\\
0.384	-0.126411\\
0.386	-0.126235\\
0.388	-0.126064\\
0.39	-0.125897\\
0.392	-0.125735\\
0.394	-0.125577\\
0.396	-0.125424\\
0.398	-0.125276\\
0.4	-0.125133\\
0.402	-0.124995\\
0.404	-0.124863\\
0.406	-0.124736\\
0.408	-0.124614\\
0.41	-0.124498\\
0.412	-0.124388\\
0.414	-0.124284\\
0.416	-0.124185\\
0.418	-0.124092\\
0.42	-0.124005\\
0.422	-0.123924\\
0.424	-0.123848\\
0.426	-0.123779\\
0.428	-0.123715\\
0.43	-0.123657\\
0.432	-0.123605\\
0.434	-0.123558\\
0.436	-0.123517\\
0.438	-0.123482\\
0.44	-0.123452\\
0.442	-0.123428\\
0.444	-0.123409\\
0.446	-0.123396\\
0.448	-0.123387\\
0.45	-0.123384\\
0.452	-0.123386\\
0.454	-0.123393\\
0.456	-0.123404\\
0.458	-0.12342\\
0.46	-0.123441\\
0.462	-0.123466\\
0.464	-0.123495\\
0.466	-0.123529\\
0.468	-0.123566\\
0.47	-0.123608\\
0.472	-0.123653\\
0.474	-0.123701\\
0.476	-0.123753\\
0.478	-0.123808\\
0.48	-0.123867\\
0.482	-0.123928\\
0.484	-0.123992\\
0.486	-0.124059\\
0.488	-0.124128\\
0.49	-0.1242\\
0.492	-0.124274\\
0.494	-0.124349\\
0.496	-0.124427\\
0.498	-0.124506\\
0.5	-0.124587\\
0.502	-0.124669\\
0.504	-0.124753\\
0.506	-0.124838\\
0.508	-0.124923\\
0.51	-0.125009\\
0.512	-0.125096\\
0.514	-0.125184\\
0.516	-0.125272\\
0.518	-0.12536\\
0.52	-0.125448\\
0.522	-0.125536\\
0.524	-0.125624\\
0.526	-0.125711\\
0.528	-0.125798\\
0.53	-0.125885\\
0.532	-0.125971\\
0.534	-0.126056\\
0.536	-0.12614\\
0.538	-0.126223\\
0.54	-0.126305\\
0.542	-0.126386\\
0.544	-0.126465\\
0.546	-0.126543\\
0.548	-0.12662\\
0.55	-0.126694\\
0.552	-0.126768\\
0.554	-0.126839\\
0.556	-0.126908\\
0.558	-0.126976\\
0.56	-0.127041\\
0.562	-0.127105\\
0.564	-0.127166\\
0.566	-0.127225\\
0.568	-0.127282\\
0.57	-0.127337\\
0.572	-0.127389\\
0.574	-0.127439\\
0.576	-0.127486\\
0.578	-0.127531\\
0.58	-0.127574\\
0.582	-0.127614\\
0.584	-0.127651\\
0.586	-0.127686\\
0.588	-0.127718\\
0.59	-0.127748\\
0.592	-0.127775\\
0.594	-0.127799\\
0.596	-0.127821\\
0.598	-0.12784\\
0.6	-0.127857\\
0.602	-0.127871\\
0.604	-0.127882\\
0.606	-0.127891\\
0.608	-0.127897\\
0.61	-0.1279\\
0.612	-0.127901\\
0.614	-0.1279\\
0.616	-0.127896\\
0.618	-0.127889\\
0.62	-0.12788\\
0.622	-0.127869\\
0.624	-0.127856\\
0.626	-0.12784\\
0.628	-0.127821\\
0.63	-0.127801\\
0.632	-0.127778\\
0.634	-0.127753\\
0.636	-0.127727\\
0.638	-0.127698\\
0.64	-0.127667\\
0.642	-0.127634\\
0.644	-0.127599\\
0.646	-0.127563\\
0.648	-0.127524\\
0.65	-0.127485\\
0.652	-0.127443\\
0.654	-0.1274\\
0.656	-0.127355\\
0.658	-0.127309\\
0.66	-0.127262\\
0.662	-0.127213\\
0.664	-0.127163\\
0.666	-0.127112\\
0.668	-0.12706\\
0.67	-0.127006\\
0.672	-0.126952\\
0.674	-0.126897\\
0.676	-0.126841\\
0.678	-0.126784\\
0.68	-0.126727\\
0.682	-0.126669\\
0.684	-0.126611\\
0.686	-0.126552\\
0.688	-0.126493\\
0.69	-0.126433\\
0.692	-0.126373\\
0.694	-0.126313\\
0.696	-0.126253\\
0.698	-0.126193\\
0.7	-0.126132\\
0.702	-0.126072\\
0.704	-0.126012\\
0.706	-0.125953\\
0.708	-0.125893\\
0.71	-0.125834\\
0.712	-0.125775\\
0.714	-0.125717\\
0.716	-0.125659\\
0.718	-0.125602\\
0.72	-0.125546\\
0.722	-0.12549\\
0.724	-0.125435\\
0.726	-0.12538\\
0.728	-0.125327\\
0.73	-0.125274\\
0.732	-0.125223\\
0.734	-0.125172\\
0.736	-0.125122\\
0.738	-0.125074\\
0.74	-0.125026\\
0.742	-0.12498\\
0.744	-0.124934\\
0.746	-0.12489\\
0.748	-0.124848\\
0.75	-0.124806\\
0.752	-0.124766\\
0.754	-0.124727\\
0.756	-0.124689\\
0.758	-0.124653\\
0.76	-0.124618\\
0.762	-0.124584\\
0.764	-0.124552\\
0.766	-0.124522\\
0.768	-0.124492\\
0.77	-0.124465\\
0.772	-0.124438\\
0.774	-0.124413\\
0.776	-0.12439\\
0.778	-0.124368\\
0.78	-0.124347\\
0.782	-0.124328\\
0.784	-0.124311\\
0.786	-0.124295\\
0.788	-0.12428\\
0.79	-0.124266\\
0.792	-0.124255\\
0.794	-0.124244\\
0.796	-0.124235\\
0.798	-0.124227\\
0.8	-0.124221\\
0.802	-0.124216\\
0.804	-0.124213\\
0.806	-0.12421\\
0.808	-0.124209\\
0.81	-0.124209\\
0.812	-0.124211\\
0.814	-0.124214\\
0.816	-0.124217\\
0.818	-0.124222\\
0.82	-0.124229\\
0.822	-0.124236\\
0.824	-0.124244\\
0.826	-0.124253\\
0.828	-0.124264\\
0.83	-0.124275\\
0.832	-0.124287\\
0.834	-0.1243\\
0.836	-0.124314\\
0.838	-0.124329\\
0.84	-0.124345\\
0.842	-0.124361\\
0.844	-0.124378\\
0.846	-0.124396\\
0.848	-0.124414\\
0.85	-0.124433\\
0.852	-0.124453\\
0.854	-0.124473\\
0.856	-0.124494\\
0.858	-0.124515\\
0.86	-0.124536\\
0.862	-0.124558\\
0.864	-0.12458\\
0.866	-0.124603\\
0.868	-0.124626\\
0.87	-0.124649\\
0.872	-0.124672\\
0.874	-0.124695\\
0.876	-0.124719\\
0.878	-0.124742\\
0.88	-0.124766\\
0.882	-0.12479\\
0.884	-0.124813\\
0.886	-0.124837\\
0.888	-0.124861\\
0.89	-0.124884\\
0.892	-0.124907\\
0.894	-0.124931\\
0.896	-0.124954\\
0.898	-0.124977\\
0.9	-0.124999\\
0.902	-0.125021\\
0.904	-0.125044\\
0.906	-0.125065\\
0.908	-0.125087\\
0.91	-0.125108\\
0.912	-0.125128\\
0.914	-0.125149\\
0.916	-0.125169\\
0.918	-0.125188\\
0.92	-0.125207\\
0.922	-0.125225\\
0.924	-0.125244\\
0.926	-0.125261\\
0.928	-0.125278\\
0.93	-0.125295\\
0.932	-0.125311\\
0.934	-0.125326\\
0.936	-0.125341\\
0.938	-0.125355\\
0.94	-0.125369\\
0.942	-0.125382\\
0.944	-0.125395\\
0.946	-0.125407\\
0.948	-0.125419\\
0.95	-0.125429\\
0.952	-0.12544\\
0.954	-0.125449\\
0.956	-0.125459\\
0.958	-0.125467\\
0.96	-0.125475\\
0.962	-0.125482\\
0.964	-0.125489\\
0.966	-0.125495\\
0.968	-0.125501\\
0.97	-0.125506\\
0.972	-0.125511\\
0.974	-0.125514\\
0.976	-0.125518\\
0.978	-0.125521\\
0.98	-0.125523\\
0.982	-0.125525\\
0.984	-0.125526\\
0.986	-0.125527\\
0.988	-0.125527\\
0.99	-0.125527\\
0.992	-0.125526\\
0.994	-0.125525\\
0.996	-0.125523\\
0.998	-0.125521\\
1	-0.125519\\
1.002	-0.125516\\
1.004	-0.125512\\
1.006	-0.125508\\
1.008	-0.125504\\
1.01	-0.1255\\
1.012	-0.125495\\
1.014	-0.12549\\
1.016	-0.125484\\
1.018	-0.125478\\
1.02	-0.125472\\
1.022	-0.125466\\
1.024	-0.125459\\
1.026	-0.125452\\
1.028	-0.125445\\
1.03	-0.125438\\
1.032	-0.12543\\
1.034	-0.125422\\
1.036	-0.125414\\
1.038	-0.125406\\
1.04	-0.125397\\
1.042	-0.125389\\
1.044	-0.12538\\
1.046	-0.125371\\
1.048	-0.125363\\
1.05	-0.125354\\
1.052	-0.125344\\
1.054	-0.125335\\
1.056	-0.125326\\
1.058	-0.125317\\
1.06	-0.125308\\
1.062	-0.125298\\
1.064	-0.125289\\
1.066	-0.12528\\
1.068	-0.12527\\
1.07	-0.125261\\
1.072	-0.125252\\
1.074	-0.125242\\
1.076	-0.125233\\
1.078	-0.125224\\
1.08	-0.125215\\
1.082	-0.125206\\
1.084	-0.125197\\
1.086	-0.125188\\
1.088	-0.12518\\
1.09	-0.125171\\
1.092	-0.125162\\
1.094	-0.125154\\
1.096	-0.125146\\
1.098	-0.125138\\
1.1	-0.12513\\
1.102	-0.125122\\
1.104	-0.125114\\
1.106	-0.125107\\
1.108	-0.125099\\
1.11	-0.125092\\
1.112	-0.125085\\
1.114	-0.125078\\
1.116	-0.125071\\
1.118	-0.125065\\
1.12	-0.125058\\
1.122	-0.125052\\
1.124	-0.125046\\
1.126	-0.12504\\
1.128	-0.125035\\
1.13	-0.125029\\
1.132	-0.125024\\
1.134	-0.125019\\
1.136	-0.125014\\
1.138	-0.125009\\
1.14	-0.125005\\
1.142	-0.125001\\
1.144	-0.124996\\
1.146	-0.124993\\
1.148	-0.124989\\
1.15	-0.124985\\
1.152	-0.124982\\
1.154	-0.124979\\
1.156	-0.124976\\
1.158	-0.124973\\
1.16	-0.12497\\
1.162	-0.124968\\
1.164	-0.124966\\
1.166	-0.124963\\
1.168	-0.124961\\
1.17	-0.12496\\
1.172	-0.124958\\
1.174	-0.124957\\
1.176	-0.124955\\
1.178	-0.124954\\
1.18	-0.124953\\
1.182	-0.124952\\
1.184	-0.124952\\
1.186	-0.124951\\
1.188	-0.124951\\
1.19	-0.12495\\
1.192	-0.12495\\
1.194	-0.12495\\
1.196	-0.12495\\
1.198	-0.12495\\
1.2	-0.12495\\
1.202	-0.124951\\
1.204	-0.124951\\
1.206	-0.124951\\
1.208	-0.124952\\
1.21	-0.124953\\
1.212	-0.124954\\
1.214	-0.124954\\
1.216	-0.124955\\
1.218	-0.124956\\
1.22	-0.124957\\
1.222	-0.124958\\
1.224	-0.124959\\
1.226	-0.12496\\
1.228	-0.124962\\
1.23	-0.124963\\
1.232	-0.124964\\
1.234	-0.124965\\
1.236	-0.124967\\
1.238	-0.124968\\
1.24	-0.124969\\
1.242	-0.12497\\
1.244	-0.124972\\
1.246	-0.124973\\
1.248	-0.124974\\
1.25	-0.124975\\
1.252	-0.124977\\
1.254	-0.124978\\
1.256	-0.124979\\
1.258	-0.12498\\
1.26	-0.124981\\
1.262	-0.124982\\
1.264	-0.124983\\
1.266	-0.124984\\
1.268	-0.124985\\
1.27	-0.124986\\
1.272	-0.124987\\
1.274	-0.124988\\
1.276	-0.124988\\
1.278	-0.124989\\
1.28	-0.124989\\
1.282	-0.12499\\
1.284	-0.12499\\
1.286	-0.124991\\
1.288	-0.124991\\
1.29	-0.124991\\
1.292	-0.124991\\
1.294	-0.124991\\
1.296	-0.124991\\
1.298	-0.124991\\
1.3	-0.124991\\
1.302	-0.12499\\
1.304	-0.12499\\
1.306	-0.124989\\
1.308	-0.124989\\
1.31	-0.124988\\
1.312	-0.124987\\
1.314	-0.124986\\
1.316	-0.124985\\
1.318	-0.124984\\
1.32	-0.124983\\
1.322	-0.124982\\
1.324	-0.12498\\
1.326	-0.124979\\
1.328	-0.124977\\
1.33	-0.124975\\
1.332	-0.124974\\
1.334	-0.124972\\
1.336	-0.12497\\
1.338	-0.124968\\
1.34	-0.124966\\
1.342	-0.124963\\
1.344	-0.124961\\
1.346	-0.124959\\
1.348	-0.124956\\
1.35	-0.124954\\
1.352	-0.124951\\
1.354	-0.124948\\
1.356	-0.124946\\
1.358	-0.124943\\
1.36	-0.12494\\
1.362	-0.124937\\
1.364	-0.124934\\
1.366	-0.124931\\
1.368	-0.124927\\
1.37	-0.124924\\
1.372	-0.124921\\
1.374	-0.124917\\
1.376	-0.124914\\
1.378	-0.12491\\
1.38	-0.124907\\
1.382	-0.124903\\
1.384	-0.1249\\
1.386	-0.124896\\
1.388	-0.124892\\
1.39	-0.124889\\
1.392	-0.124885\\
1.394	-0.124881\\
1.396	-0.124877\\
1.398	-0.124873\\
1.4	-0.124869\\
1.402	-0.124866\\
1.404	-0.124862\\
1.406	-0.124858\\
1.408	-0.124854\\
1.41	-0.12485\\
1.412	-0.124846\\
1.414	-0.124842\\
1.416	-0.124838\\
1.418	-0.124834\\
1.42	-0.12483\\
1.422	-0.124826\\
1.424	-0.124822\\
1.426	-0.124818\\
1.428	-0.124814\\
1.43	-0.12481\\
1.432	-0.124806\\
1.434	-0.124802\\
1.436	-0.124798\\
1.438	-0.124795\\
1.44	-0.124791\\
1.442	-0.124787\\
1.444	-0.124783\\
1.446	-0.124779\\
1.448	-0.124776\\
1.45	-0.124772\\
1.452	-0.124768\\
1.454	-0.124765\\
1.456	-0.124761\\
1.458	-0.124758\\
1.46	-0.124754\\
1.462	-0.124751\\
1.464	-0.124747\\
1.466	-0.124744\\
1.468	-0.12474\\
1.47	-0.124737\\
1.472	-0.124734\\
1.474	-0.124731\\
1.476	-0.124727\\
1.478	-0.124724\\
1.48	-0.124721\\
1.482	-0.124718\\
1.484	-0.124715\\
1.486	-0.124712\\
1.488	-0.124709\\
1.49	-0.124706\\
1.492	-0.124704\\
1.494	-0.124701\\
1.496	-0.124698\\
1.498	-0.124695\\
1.5	-0.124693\\
1.502	-0.12469\\
1.504	-0.124688\\
1.506	-0.124685\\
1.508	-0.124683\\
1.51	-0.12468\\
1.512	-0.124678\\
1.514	-0.124676\\
1.516	-0.124673\\
1.518	-0.124671\\
1.52	-0.124669\\
1.522	-0.124667\\
1.524	-0.124665\\
1.526	-0.124663\\
1.528	-0.124661\\
1.53	-0.124659\\
1.532	-0.124657\\
1.534	-0.124655\\
1.536	-0.124653\\
1.538	-0.124651\\
1.54	-0.124649\\
1.542	-0.124647\\
1.544	-0.124645\\
1.546	-0.124644\\
1.548	-0.124642\\
1.55	-0.12464\\
1.552	-0.124639\\
1.554	-0.124637\\
1.556	-0.124635\\
1.558	-0.124634\\
1.56	-0.124632\\
1.562	-0.124631\\
1.564	-0.124629\\
1.566	-0.124628\\
1.568	-0.124626\\
1.57	-0.124625\\
1.572	-0.124623\\
1.574	-0.124622\\
1.576	-0.12462\\
1.578	-0.124619\\
1.58	-0.124617\\
1.582	-0.124616\\
1.584	-0.124615\\
1.586	-0.124613\\
1.588	-0.124612\\
1.59	-0.12461\\
1.592	-0.124609\\
1.594	-0.124608\\
1.596	-0.124606\\
1.598	-0.124605\\
1.6	-0.124604\\
1.602	-0.124602\\
1.604	-0.124601\\
1.606	-0.124599\\
1.608	-0.124598\\
1.61	-0.124597\\
1.612	-0.124595\\
1.614	-0.124594\\
1.616	-0.124593\\
1.618	-0.124591\\
1.62	-0.12459\\
1.622	-0.124588\\
1.624	-0.124587\\
1.626	-0.124585\\
1.628	-0.124584\\
1.63	-0.124583\\
1.632	-0.124581\\
1.634	-0.12458\\
1.636	-0.124578\\
1.638	-0.124577\\
1.64	-0.124575\\
1.642	-0.124574\\
1.644	-0.124572\\
1.646	-0.124571\\
1.648	-0.124569\\
1.65	-0.124568\\
1.652	-0.124566\\
1.654	-0.124564\\
1.656	-0.124563\\
1.658	-0.124561\\
1.66	-0.12456\\
1.662	-0.124558\\
1.664	-0.124556\\
1.666	-0.124555\\
1.668	-0.124553\\
1.67	-0.124551\\
1.672	-0.124549\\
1.674	-0.124548\\
1.676	-0.124546\\
1.678	-0.124544\\
1.68	-0.124542\\
1.682	-0.124541\\
1.684	-0.124539\\
1.686	-0.124537\\
1.688	-0.124535\\
1.69	-0.124533\\
1.692	-0.124531\\
1.694	-0.12453\\
1.696	-0.124528\\
1.698	-0.124526\\
1.7	-0.124524\\
1.702	-0.124522\\
1.704	-0.12452\\
1.706	-0.124518\\
1.708	-0.124516\\
1.71	-0.124514\\
1.712	-0.124512\\
1.714	-0.12451\\
1.716	-0.124508\\
1.718	-0.124506\\
1.72	-0.124504\\
1.722	-0.124502\\
1.724	-0.1245\\
1.726	-0.124498\\
1.728	-0.124496\\
1.73	-0.124494\\
1.732	-0.124492\\
1.734	-0.124489\\
1.736	-0.124487\\
1.738	-0.124485\\
1.74	-0.124483\\
1.742	-0.124481\\
1.744	-0.124479\\
1.746	-0.124477\\
1.748	-0.124475\\
1.75	-0.124472\\
1.752	-0.12447\\
1.754	-0.124468\\
1.756	-0.124466\\
1.758	-0.124464\\
1.76	-0.124462\\
1.762	-0.124459\\
1.764	-0.124457\\
1.766	-0.124455\\
1.768	-0.124453\\
1.77	-0.124451\\
1.772	-0.124448\\
1.774	-0.124446\\
1.776	-0.124444\\
1.778	-0.124442\\
1.78	-0.12444\\
1.782	-0.124437\\
1.784	-0.124435\\
1.786	-0.124433\\
1.788	-0.124431\\
1.79	-0.124429\\
1.792	-0.124426\\
1.794	-0.124424\\
1.796	-0.124422\\
1.798	-0.12442\\
1.8	-0.124418\\
1.802	-0.124416\\
1.804	-0.124413\\
1.806	-0.124411\\
1.808	-0.124409\\
1.81	-0.124407\\
1.812	-0.124405\\
1.814	-0.124403\\
1.816	-0.124401\\
1.818	-0.124398\\
1.82	-0.124396\\
1.822	-0.124394\\
1.824	-0.124392\\
1.826	-0.12439\\
1.828	-0.124388\\
1.83	-0.124386\\
1.832	-0.124384\\
1.834	-0.124382\\
1.836	-0.124379\\
1.838	-0.124377\\
1.84	-0.124375\\
1.842	-0.124373\\
1.844	-0.124371\\
1.846	-0.124369\\
1.848	-0.124367\\
1.85	-0.124365\\
1.852	-0.124363\\
1.854	-0.124361\\
1.856	-0.124359\\
1.858	-0.124357\\
1.86	-0.124355\\
1.862	-0.124353\\
1.864	-0.124351\\
1.866	-0.124349\\
1.868	-0.124347\\
1.87	-0.124345\\
1.872	-0.124343\\
1.874	-0.124341\\
1.876	-0.124339\\
1.878	-0.124338\\
1.88	-0.124336\\
1.882	-0.124334\\
1.884	-0.124332\\
1.886	-0.12433\\
1.888	-0.124328\\
1.89	-0.124326\\
1.892	-0.124324\\
1.894	-0.124322\\
1.896	-0.124321\\
1.898	-0.124319\\
1.9	-0.124317\\
1.902	-0.124315\\
1.904	-0.124313\\
1.906	-0.124311\\
1.908	-0.124309\\
1.91	-0.124308\\
1.912	-0.124306\\
1.914	-0.124304\\
1.916	-0.124302\\
1.918	-0.1243\\
1.92	-0.124299\\
1.922	-0.124297\\
1.924	-0.124295\\
1.926	-0.124293\\
1.928	-0.124291\\
1.93	-0.12429\\
1.932	-0.124288\\
1.934	-0.124286\\
1.936	-0.124284\\
1.938	-0.124283\\
1.94	-0.124281\\
1.942	-0.124279\\
1.944	-0.124277\\
1.946	-0.124275\\
1.948	-0.124274\\
1.95	-0.124272\\
1.952	-0.12427\\
1.954	-0.124268\\
1.956	-0.124267\\
1.958	-0.124265\\
1.96	-0.124263\\
1.962	-0.124261\\
1.964	-0.12426\\
1.966	-0.124258\\
1.968	-0.124256\\
1.97	-0.124254\\
1.972	-0.124253\\
1.974	-0.124251\\
1.976	-0.124249\\
1.978	-0.124247\\
1.98	-0.124246\\
1.982	-0.124244\\
1.984	-0.124242\\
1.986	-0.12424\\
1.988	-0.124239\\
1.99	-0.124237\\
1.992	-0.124235\\
1.994	-0.124233\\
1.996	-0.124232\\
1.998	-0.12423\\
2	-0.124228\\
};
\addplot [color=mycolor16, line width=1.0pt, forget plot]
  table[row sep=crcr]{%
-0.124	0\\
-0.122	0\\
-0.12	0\\
-0.118	0\\
-0.116	0\\
-0.114	0\\
-0.112	0\\
-0.11	0\\
-0.108	0\\
-0.106	0\\
-0.104	0\\
-0.102	0\\
-0.1	0\\
-0.098	0\\
-0.096	0\\
-0.094	0\\
-0.092	0\\
-0.09	0\\
-0.088	0\\
-0.086	0\\
-0.084	0\\
-0.082	0\\
-0.08	0\\
-0.078	0\\
-0.076	0\\
-0.074	0\\
-0.072	0\\
-0.07	0\\
-0.068	0\\
-0.066	0\\
-0.064	0\\
-0.062	0\\
-0.06	0\\
-0.058	0\\
-0.056	0\\
-0.054	0\\
-0.052	0\\
-0.05	0\\
-0.048	0\\
-0.046	0\\
-0.044	0\\
-0.042	0\\
-0.04	0\\
-0.038	0\\
-0.036	0\\
-0.034	0\\
-0.032	0\\
-0.03	0\\
-0.028	0\\
-0.026	0\\
-0.024	0\\
-0.022	0\\
-0.02	0\\
-0.018	0\\
-0.016	0\\
-0.014	0\\
-0.012	0\\
-0.01	0\\
-0.008	0\\
-0.006	0\\
-0.004	0\\
-0.002	0\\
0	0\\
0.002	-0.00424271\\
0.004	-0.00844743\\
0.006	-0.0126038\\
0.008	-0.0167076\\
0.01	-0.0207579\\
0.012	-0.024755\\
0.014	-0.0287\\
0.016	-0.0325941\\
0.018	-0.0364386\\
0.02	-0.0402347\\
0.022	-0.0439833\\
0.024	-0.0476855\\
0.026	-0.0513417\\
0.028	-0.0549526\\
0.03	-0.0585184\\
0.032	-0.0620392\\
0.034	-0.0655148\\
0.036	-0.068945\\
0.038	-0.0723293\\
0.04	-0.0756671\\
0.042	-0.0789574\\
0.044	-0.0821992\\
0.046	-0.0853915\\
0.048	-0.0885328\\
0.05	-0.0916216\\
0.052	-0.0946564\\
0.054	-0.0976354\\
0.056	-0.100557\\
0.058	-0.103418\\
0.06	-0.106218\\
0.062	-0.108954\\
0.064	-0.111624\\
0.066	-0.114226\\
0.068	-0.116757\\
0.07	-0.119215\\
0.072	-0.121597\\
0.074	-0.123901\\
0.076	-0.126125\\
0.078	-0.128267\\
0.08	-0.130323\\
0.082	-0.132292\\
0.084	-0.134172\\
0.086	-0.13596\\
0.088	-0.137653\\
0.09	-0.139251\\
0.092	-0.140751\\
0.094	-0.142152\\
0.096	-0.143451\\
0.098	-0.144647\\
0.1	-0.145738\\
0.102	-0.146725\\
0.104	-0.147604\\
0.106	-0.148376\\
0.108	-0.14904\\
0.11	-0.149595\\
0.112	-0.15004\\
0.114	-0.150376\\
0.116	-0.150602\\
0.118	-0.15072\\
0.12	-0.150728\\
0.122	-0.150628\\
0.124	-0.150421\\
0.126	-0.150107\\
0.128	-0.149688\\
0.13	-0.149166\\
0.132	-0.148542\\
0.134	-0.147817\\
0.136	-0.146995\\
0.138	-0.146077\\
0.14	-0.145066\\
0.142	-0.143964\\
0.144	-0.142774\\
0.146	-0.1415\\
0.148	-0.140145\\
0.15	-0.138711\\
0.152	-0.137203\\
0.154	-0.135624\\
0.156	-0.133978\\
0.158	-0.132269\\
0.16	-0.130502\\
0.162	-0.128679\\
0.164	-0.126807\\
0.166	-0.124888\\
0.168	-0.122929\\
0.17	-0.120932\\
0.172	-0.118903\\
0.174	-0.116847\\
0.176	-0.114768\\
0.178	-0.112671\\
0.18	-0.11056\\
0.182	-0.108441\\
0.184	-0.106319\\
0.186	-0.104198\\
0.188	-0.102082\\
0.19	-0.0999768\\
0.192	-0.0978867\\
0.194	-0.0958163\\
0.196	-0.09377\\
0.198	-0.0917523\\
0.2	-0.0897675\\
0.202	-0.0878197\\
0.204	-0.085913\\
0.206	-0.0840515\\
0.208	-0.0822389\\
0.21	-0.080479\\
0.212	-0.0787753\\
0.214	-0.0771312\\
0.216	-0.0755499\\
0.218	-0.0740344\\
0.22	-0.0725878\\
0.222	-0.0712125\\
0.224	-0.0699113\\
0.226	-0.0686864\\
0.228	-0.0675399\\
0.23	-0.0664738\\
0.232	-0.0654898\\
0.234	-0.0645894\\
0.236	-0.0637741\\
0.238	-0.0630448\\
0.24	-0.0624026\\
0.242	-0.0618481\\
0.244	-0.0613819\\
0.246	-0.0610043\\
0.248	-0.0607154\\
0.25	-0.060515\\
0.252	-0.060403\\
0.254	-0.0603787\\
0.256	-0.0604417\\
0.258	-0.0605909\\
0.26	-0.0608253\\
0.262	-0.0611438\\
0.264	-0.061545\\
0.266	-0.0620273\\
0.268	-0.0625889\\
0.27	-0.0632281\\
0.272	-0.0639428\\
0.274	-0.0647308\\
0.276	-0.0655899\\
0.278	-0.0665177\\
0.28	-0.0675117\\
0.282	-0.0685692\\
0.284	-0.0696874\\
0.286	-0.0708637\\
0.288	-0.072095\\
0.29	-0.0733783\\
0.292	-0.0747107\\
0.294	-0.0760889\\
0.296	-0.0775099\\
0.298	-0.0789705\\
0.3	-0.0804673\\
0.302	-0.0819972\\
0.304	-0.0835568\\
0.306	-0.0851429\\
0.308	-0.0867522\\
0.31	-0.0883814\\
0.312	-0.0900273\\
0.314	-0.0916866\\
0.316	-0.093356\\
0.318	-0.0950325\\
0.32	-0.096713\\
0.322	-0.0983942\\
0.324	-0.100073\\
0.326	-0.101747\\
0.328	-0.103413\\
0.33	-0.105068\\
0.332	-0.10671\\
0.334	-0.108335\\
0.336	-0.109942\\
0.338	-0.111528\\
0.34	-0.11309\\
0.342	-0.114626\\
0.344	-0.116134\\
0.346	-0.117613\\
0.348	-0.119059\\
0.35	-0.120472\\
0.352	-0.12185\\
0.354	-0.12319\\
0.356	-0.124493\\
0.358	-0.125755\\
0.36	-0.126977\\
0.362	-0.128156\\
0.364	-0.129293\\
0.366	-0.130385\\
0.368	-0.131433\\
0.37	-0.132436\\
0.372	-0.133393\\
0.374	-0.134304\\
0.376	-0.135168\\
0.378	-0.135986\\
0.38	-0.136757\\
0.382	-0.137481\\
0.384	-0.138159\\
0.386	-0.138791\\
0.388	-0.139376\\
0.39	-0.139916\\
0.392	-0.140411\\
0.394	-0.140862\\
0.396	-0.141269\\
0.398	-0.141633\\
0.4	-0.141955\\
0.402	-0.142235\\
0.404	-0.142475\\
0.406	-0.142676\\
0.408	-0.142838\\
0.41	-0.142964\\
0.412	-0.143053\\
0.414	-0.143108\\
0.416	-0.143129\\
0.418	-0.143118\\
0.42	-0.143075\\
0.422	-0.143004\\
0.424	-0.142903\\
0.426	-0.142776\\
0.428	-0.142624\\
0.43	-0.142447\\
0.432	-0.142247\\
0.434	-0.142026\\
0.436	-0.141784\\
0.438	-0.141524\\
0.44	-0.141247\\
0.442	-0.140953\\
0.444	-0.140645\\
0.446	-0.140323\\
0.448	-0.139989\\
0.45	-0.139644\\
0.452	-0.139289\\
0.454	-0.138926\\
0.456	-0.138555\\
0.458	-0.138178\\
0.46	-0.137796\\
0.462	-0.13741\\
0.464	-0.13702\\
0.466	-0.136628\\
0.468	-0.136235\\
0.47	-0.135841\\
0.472	-0.135448\\
0.474	-0.135056\\
0.476	-0.134665\\
0.478	-0.134277\\
0.48	-0.133893\\
0.482	-0.133512\\
0.484	-0.133135\\
0.486	-0.132763\\
0.488	-0.132396\\
0.49	-0.132035\\
0.492	-0.13168\\
0.494	-0.131332\\
0.496	-0.13099\\
0.498	-0.130655\\
0.5	-0.130327\\
0.502	-0.130006\\
0.504	-0.129693\\
0.506	-0.129387\\
0.508	-0.129089\\
0.51	-0.128798\\
0.512	-0.128515\\
0.514	-0.12824\\
0.516	-0.127972\\
0.518	-0.127712\\
0.52	-0.127459\\
0.522	-0.127213\\
0.524	-0.126975\\
0.526	-0.126744\\
0.528	-0.126519\\
0.53	-0.126301\\
0.532	-0.12609\\
0.534	-0.125885\\
0.536	-0.125686\\
0.538	-0.125492\\
0.54	-0.125305\\
0.542	-0.125123\\
0.544	-0.124946\\
0.546	-0.124774\\
0.548	-0.124607\\
0.55	-0.124444\\
0.552	-0.124286\\
0.554	-0.124132\\
0.556	-0.123982\\
0.558	-0.123835\\
0.56	-0.123692\\
0.562	-0.123553\\
0.564	-0.123417\\
0.566	-0.123284\\
0.568	-0.123153\\
0.57	-0.123026\\
0.572	-0.122901\\
0.574	-0.122779\\
0.576	-0.12266\\
0.578	-0.122542\\
0.58	-0.122428\\
0.582	-0.122315\\
0.584	-0.122205\\
0.586	-0.122097\\
0.588	-0.121992\\
0.59	-0.121888\\
0.592	-0.121787\\
0.594	-0.121689\\
0.596	-0.121592\\
0.598	-0.121499\\
0.6	-0.121407\\
0.602	-0.121319\\
0.604	-0.121233\\
0.606	-0.12115\\
0.608	-0.12107\\
0.61	-0.120993\\
0.612	-0.120919\\
0.614	-0.120849\\
0.616	-0.120782\\
0.618	-0.120718\\
0.62	-0.120659\\
0.622	-0.120603\\
0.624	-0.120551\\
0.626	-0.120504\\
0.628	-0.120461\\
0.63	-0.120422\\
0.632	-0.120388\\
0.634	-0.120359\\
0.636	-0.120335\\
0.638	-0.120316\\
0.64	-0.120302\\
0.642	-0.120294\\
0.644	-0.120291\\
0.646	-0.120294\\
0.648	-0.120302\\
0.65	-0.120317\\
0.652	-0.120337\\
0.654	-0.120363\\
0.656	-0.120395\\
0.658	-0.120433\\
0.66	-0.120477\\
0.662	-0.120528\\
0.664	-0.120584\\
0.666	-0.120647\\
0.668	-0.120716\\
0.67	-0.120791\\
0.672	-0.120872\\
0.674	-0.120959\\
0.676	-0.121052\\
0.678	-0.121151\\
0.68	-0.121255\\
0.682	-0.121365\\
0.684	-0.121481\\
0.686	-0.121602\\
0.688	-0.121728\\
0.69	-0.121859\\
0.692	-0.121995\\
0.694	-0.122135\\
0.696	-0.12228\\
0.698	-0.12243\\
0.7	-0.122583\\
0.702	-0.12274\\
0.704	-0.1229\\
0.706	-0.123063\\
0.708	-0.12323\\
0.71	-0.123399\\
0.712	-0.12357\\
0.714	-0.123744\\
0.716	-0.123919\\
0.718	-0.124096\\
0.72	-0.124274\\
0.722	-0.124452\\
0.724	-0.124632\\
0.726	-0.124811\\
0.728	-0.12499\\
0.73	-0.125169\\
0.732	-0.125347\\
0.734	-0.125524\\
0.736	-0.125699\\
0.738	-0.125873\\
0.74	-0.126045\\
0.742	-0.126214\\
0.744	-0.126381\\
0.746	-0.126544\\
0.748	-0.126705\\
0.75	-0.126862\\
0.752	-0.127015\\
0.754	-0.127164\\
0.756	-0.127309\\
0.758	-0.127449\\
0.76	-0.127584\\
0.762	-0.127715\\
0.764	-0.12784\\
0.766	-0.12796\\
0.768	-0.128074\\
0.77	-0.128182\\
0.772	-0.128284\\
0.774	-0.128381\\
0.776	-0.128471\\
0.778	-0.128555\\
0.78	-0.128632\\
0.782	-0.128703\\
0.784	-0.128767\\
0.786	-0.128824\\
0.788	-0.128875\\
0.79	-0.128919\\
0.792	-0.128956\\
0.794	-0.128987\\
0.796	-0.129011\\
0.798	-0.129028\\
0.8	-0.129038\\
0.802	-0.129042\\
0.804	-0.129039\\
0.806	-0.12903\\
0.808	-0.129015\\
0.81	-0.128993\\
0.812	-0.128965\\
0.814	-0.128931\\
0.816	-0.128891\\
0.818	-0.128846\\
0.82	-0.128795\\
0.822	-0.128739\\
0.824	-0.128677\\
0.826	-0.128611\\
0.828	-0.128539\\
0.83	-0.128463\\
0.832	-0.128383\\
0.834	-0.128299\\
0.836	-0.128211\\
0.838	-0.128119\\
0.84	-0.128023\\
0.842	-0.127925\\
0.844	-0.127823\\
0.846	-0.127719\\
0.848	-0.127612\\
0.85	-0.127504\\
0.852	-0.127393\\
0.854	-0.12728\\
0.856	-0.127166\\
0.858	-0.127051\\
0.86	-0.126935\\
0.862	-0.126818\\
0.864	-0.126701\\
0.866	-0.126583\\
0.868	-0.126466\\
0.87	-0.126349\\
0.872	-0.126232\\
0.874	-0.126115\\
0.876	-0.126\\
0.878	-0.125886\\
0.88	-0.125773\\
0.882	-0.125662\\
0.884	-0.125552\\
0.886	-0.125444\\
0.888	-0.125338\\
0.89	-0.125234\\
0.892	-0.125133\\
0.894	-0.125034\\
0.896	-0.124938\\
0.898	-0.124844\\
0.9	-0.124754\\
0.902	-0.124666\\
0.904	-0.124581\\
0.906	-0.1245\\
0.908	-0.124422\\
0.91	-0.124347\\
0.912	-0.124275\\
0.914	-0.124207\\
0.916	-0.124142\\
0.918	-0.124081\\
0.92	-0.124024\\
0.922	-0.12397\\
0.924	-0.123919\\
0.926	-0.123873\\
0.928	-0.123829\\
0.93	-0.123789\\
0.932	-0.123753\\
0.934	-0.12372\\
0.936	-0.123691\\
0.938	-0.123665\\
0.94	-0.123643\\
0.942	-0.123624\\
0.944	-0.123608\\
0.946	-0.123595\\
0.948	-0.123585\\
0.95	-0.123579\\
0.952	-0.123575\\
0.954	-0.123574\\
0.956	-0.123576\\
0.958	-0.12358\\
0.96	-0.123587\\
0.962	-0.123597\\
0.964	-0.123609\\
0.966	-0.123623\\
0.968	-0.123639\\
0.97	-0.123657\\
0.972	-0.123678\\
0.974	-0.1237\\
0.976	-0.123723\\
0.978	-0.123749\\
0.98	-0.123775\\
0.982	-0.123803\\
0.984	-0.123833\\
0.986	-0.123863\\
0.988	-0.123895\\
0.99	-0.123927\\
0.992	-0.12396\\
0.994	-0.123994\\
0.996	-0.124029\\
0.998	-0.124064\\
1	-0.1241\\
1.002	-0.124135\\
1.004	-0.124171\\
1.006	-0.124207\\
1.008	-0.124244\\
1.01	-0.12428\\
1.012	-0.124316\\
1.014	-0.124352\\
1.016	-0.124388\\
1.018	-0.124423\\
1.02	-0.124458\\
1.022	-0.124493\\
1.024	-0.124527\\
1.026	-0.12456\\
1.028	-0.124593\\
1.03	-0.124625\\
1.032	-0.124657\\
1.034	-0.124688\\
1.036	-0.124718\\
1.038	-0.124748\\
1.04	-0.124776\\
1.042	-0.124804\\
1.044	-0.124831\\
1.046	-0.124858\\
1.048	-0.124883\\
1.05	-0.124907\\
1.052	-0.124931\\
1.054	-0.124954\\
1.056	-0.124976\\
1.058	-0.124997\\
1.06	-0.125017\\
1.062	-0.125036\\
1.064	-0.125055\\
1.066	-0.125072\\
1.068	-0.125089\\
1.07	-0.125105\\
1.072	-0.12512\\
1.074	-0.125135\\
1.076	-0.125149\\
1.078	-0.125162\\
1.08	-0.125174\\
1.082	-0.125186\\
1.084	-0.125197\\
1.086	-0.125207\\
1.088	-0.125217\\
1.09	-0.125226\\
1.092	-0.125235\\
1.094	-0.125243\\
1.096	-0.125251\\
1.098	-0.125258\\
1.1	-0.125264\\
1.102	-0.125271\\
1.104	-0.125276\\
1.106	-0.125282\\
1.108	-0.125287\\
1.11	-0.125292\\
1.112	-0.125296\\
1.114	-0.1253\\
1.116	-0.125304\\
1.118	-0.125308\\
1.12	-0.125311\\
1.122	-0.125314\\
1.124	-0.125317\\
1.126	-0.12532\\
1.128	-0.125323\\
1.13	-0.125325\\
1.132	-0.125327\\
1.134	-0.125329\\
1.136	-0.125331\\
1.138	-0.125332\\
1.14	-0.125334\\
1.142	-0.125335\\
1.144	-0.125337\\
1.146	-0.125338\\
1.148	-0.125339\\
1.15	-0.125339\\
1.152	-0.12534\\
1.154	-0.12534\\
1.156	-0.125341\\
1.158	-0.125341\\
1.16	-0.125341\\
1.162	-0.125341\\
1.164	-0.12534\\
1.166	-0.12534\\
1.168	-0.125339\\
1.17	-0.125338\\
1.172	-0.125337\\
1.174	-0.125335\\
1.176	-0.125334\\
1.178	-0.125332\\
1.18	-0.12533\\
1.182	-0.125327\\
1.184	-0.125325\\
1.186	-0.125322\\
1.188	-0.125318\\
1.19	-0.125315\\
1.192	-0.125311\\
1.194	-0.125307\\
1.196	-0.125302\\
1.198	-0.125297\\
1.2	-0.125292\\
1.202	-0.125287\\
1.204	-0.125281\\
1.206	-0.125274\\
1.208	-0.125268\\
1.21	-0.125261\\
1.212	-0.125254\\
1.214	-0.125246\\
1.216	-0.125238\\
1.218	-0.125229\\
1.22	-0.125221\\
1.222	-0.125211\\
1.224	-0.125202\\
1.226	-0.125192\\
1.228	-0.125182\\
1.23	-0.125171\\
1.232	-0.12516\\
1.234	-0.125149\\
1.236	-0.125138\\
1.238	-0.125126\\
1.24	-0.125114\\
1.242	-0.125101\\
1.244	-0.125089\\
1.246	-0.125076\\
1.248	-0.125063\\
1.25	-0.125049\\
1.252	-0.125036\\
1.254	-0.125022\\
1.256	-0.125008\\
1.258	-0.124994\\
1.26	-0.12498\\
1.262	-0.124966\\
1.264	-0.124951\\
1.266	-0.124937\\
1.268	-0.124922\\
1.27	-0.124908\\
1.272	-0.124893\\
1.274	-0.124879\\
1.276	-0.124864\\
1.278	-0.12485\\
1.28	-0.124835\\
1.282	-0.124821\\
1.284	-0.124807\\
1.286	-0.124793\\
1.288	-0.124779\\
1.29	-0.124766\\
1.292	-0.124752\\
1.294	-0.124739\\
1.296	-0.124726\\
1.298	-0.124714\\
1.3	-0.124701\\
1.302	-0.124689\\
1.304	-0.124677\\
1.306	-0.124666\\
1.308	-0.124655\\
1.31	-0.124644\\
1.312	-0.124634\\
1.314	-0.124624\\
1.316	-0.124615\\
1.318	-0.124606\\
1.32	-0.124597\\
1.322	-0.124589\\
1.324	-0.124581\\
1.326	-0.124574\\
1.328	-0.124567\\
1.33	-0.124561\\
1.332	-0.124555\\
1.334	-0.124549\\
1.336	-0.124545\\
1.338	-0.12454\\
1.34	-0.124536\\
1.342	-0.124533\\
1.344	-0.12453\\
1.346	-0.124527\\
1.348	-0.124525\\
1.35	-0.124524\\
1.352	-0.124522\\
1.354	-0.124522\\
1.356	-0.124522\\
1.358	-0.124522\\
1.36	-0.124522\\
1.362	-0.124524\\
1.364	-0.124525\\
1.366	-0.124527\\
1.368	-0.124529\\
1.37	-0.124532\\
1.372	-0.124535\\
1.374	-0.124538\\
1.376	-0.124542\\
1.378	-0.124546\\
1.38	-0.12455\\
1.382	-0.124554\\
1.384	-0.124559\\
1.386	-0.124564\\
1.388	-0.124569\\
1.39	-0.124575\\
1.392	-0.12458\\
1.394	-0.124586\\
1.396	-0.124592\\
1.398	-0.124598\\
1.4	-0.124604\\
1.402	-0.124611\\
1.404	-0.124617\\
1.406	-0.124623\\
1.408	-0.12463\\
1.41	-0.124636\\
1.412	-0.124643\\
1.414	-0.124649\\
1.416	-0.124656\\
1.418	-0.124662\\
1.42	-0.124668\\
1.422	-0.124674\\
1.424	-0.124681\\
1.426	-0.124687\\
1.428	-0.124693\\
1.43	-0.124698\\
1.432	-0.124704\\
1.434	-0.124709\\
1.436	-0.124715\\
1.438	-0.12472\\
1.44	-0.124725\\
1.442	-0.124729\\
1.444	-0.124734\\
1.446	-0.124738\\
1.448	-0.124742\\
1.45	-0.124746\\
1.452	-0.124749\\
1.454	-0.124753\\
1.456	-0.124756\\
1.458	-0.124758\\
1.46	-0.124761\\
1.462	-0.124763\\
1.464	-0.124765\\
1.466	-0.124767\\
1.468	-0.124768\\
1.47	-0.124769\\
1.472	-0.12477\\
1.474	-0.124771\\
1.476	-0.124771\\
1.478	-0.124771\\
1.48	-0.124771\\
1.482	-0.12477\\
1.484	-0.12477\\
1.486	-0.124769\\
1.488	-0.124767\\
1.49	-0.124766\\
1.492	-0.124764\\
1.494	-0.124762\\
1.496	-0.12476\\
1.498	-0.124757\\
1.5	-0.124755\\
1.502	-0.124752\\
1.504	-0.124749\\
1.506	-0.124746\\
1.508	-0.124742\\
1.51	-0.124738\\
1.512	-0.124735\\
1.514	-0.124731\\
1.516	-0.124727\\
1.518	-0.124722\\
1.52	-0.124718\\
1.522	-0.124714\\
1.524	-0.124709\\
1.526	-0.124704\\
1.528	-0.1247\\
1.53	-0.124695\\
1.532	-0.12469\\
1.534	-0.124685\\
1.536	-0.12468\\
1.538	-0.124675\\
1.54	-0.12467\\
1.542	-0.124664\\
1.544	-0.124659\\
1.546	-0.124654\\
1.548	-0.124649\\
1.55	-0.124643\\
1.552	-0.124638\\
1.554	-0.124633\\
1.556	-0.124628\\
1.558	-0.124623\\
1.56	-0.124617\\
1.562	-0.124612\\
1.564	-0.124607\\
1.566	-0.124602\\
1.568	-0.124597\\
1.57	-0.124592\\
1.572	-0.124587\\
1.574	-0.124583\\
1.576	-0.124578\\
1.578	-0.124573\\
1.58	-0.124569\\
1.582	-0.124564\\
1.584	-0.12456\\
1.586	-0.124555\\
1.588	-0.124551\\
1.59	-0.124547\\
1.592	-0.124543\\
1.594	-0.124539\\
1.596	-0.124535\\
1.598	-0.124531\\
1.6	-0.124527\\
1.602	-0.124523\\
1.604	-0.12452\\
1.606	-0.124516\\
1.608	-0.124513\\
1.61	-0.124509\\
1.612	-0.124506\\
1.614	-0.124503\\
1.616	-0.124499\\
1.618	-0.124496\\
1.62	-0.124493\\
1.622	-0.12449\\
1.624	-0.124487\\
1.626	-0.124485\\
1.628	-0.124482\\
1.63	-0.124479\\
1.632	-0.124477\\
1.634	-0.124474\\
1.636	-0.124471\\
1.638	-0.124469\\
1.64	-0.124467\\
1.642	-0.124464\\
1.644	-0.124462\\
1.646	-0.12446\\
1.648	-0.124457\\
1.65	-0.124455\\
1.652	-0.124453\\
1.654	-0.124451\\
1.656	-0.124449\\
1.658	-0.124447\\
1.66	-0.124444\\
1.662	-0.124442\\
1.664	-0.12444\\
1.666	-0.124438\\
1.668	-0.124436\\
1.67	-0.124434\\
1.672	-0.124432\\
1.674	-0.124431\\
1.676	-0.124429\\
1.678	-0.124427\\
1.68	-0.124425\\
1.682	-0.124423\\
1.684	-0.124421\\
1.686	-0.124419\\
1.688	-0.124417\\
1.69	-0.124415\\
1.692	-0.124413\\
1.694	-0.124412\\
1.696	-0.12441\\
1.698	-0.124408\\
1.7	-0.124406\\
1.702	-0.124404\\
1.704	-0.124402\\
1.706	-0.1244\\
1.708	-0.124399\\
1.71	-0.124397\\
1.712	-0.124395\\
1.714	-0.124393\\
1.716	-0.124391\\
1.718	-0.12439\\
1.72	-0.124388\\
1.722	-0.124386\\
1.724	-0.124384\\
1.726	-0.124382\\
1.728	-0.124381\\
1.73	-0.124379\\
1.732	-0.124377\\
1.734	-0.124376\\
1.736	-0.124374\\
1.738	-0.124372\\
1.74	-0.12437\\
1.742	-0.124369\\
1.744	-0.124367\\
1.746	-0.124365\\
1.748	-0.124364\\
1.75	-0.124362\\
1.752	-0.124361\\
1.754	-0.124359\\
1.756	-0.124358\\
1.758	-0.124356\\
1.76	-0.124355\\
1.762	-0.124353\\
1.764	-0.124352\\
1.766	-0.12435\\
1.768	-0.124349\\
1.77	-0.124347\\
1.772	-0.124346\\
1.774	-0.124345\\
1.776	-0.124343\\
1.778	-0.124342\\
1.78	-0.124341\\
1.782	-0.12434\\
1.784	-0.124338\\
1.786	-0.124337\\
1.788	-0.124336\\
1.79	-0.124335\\
1.792	-0.124333\\
1.794	-0.124332\\
1.796	-0.124331\\
1.798	-0.12433\\
1.8	-0.124329\\
1.802	-0.124328\\
1.804	-0.124327\\
1.806	-0.124325\\
1.808	-0.124324\\
1.81	-0.124323\\
1.812	-0.124322\\
1.814	-0.124321\\
1.816	-0.12432\\
1.818	-0.124319\\
1.82	-0.124318\\
1.822	-0.124317\\
1.824	-0.124316\\
1.826	-0.124315\\
1.828	-0.124314\\
1.83	-0.124312\\
1.832	-0.124311\\
1.834	-0.12431\\
1.836	-0.124309\\
1.838	-0.124308\\
1.84	-0.124307\\
1.842	-0.124306\\
1.844	-0.124304\\
1.846	-0.124303\\
1.848	-0.124302\\
1.85	-0.124301\\
1.852	-0.124299\\
1.854	-0.124298\\
1.856	-0.124297\\
1.858	-0.124295\\
1.86	-0.124294\\
1.862	-0.124293\\
1.864	-0.124291\\
1.866	-0.12429\\
1.868	-0.124288\\
1.87	-0.124287\\
1.872	-0.124285\\
1.874	-0.124284\\
1.876	-0.124282\\
1.878	-0.12428\\
1.88	-0.124279\\
1.882	-0.124277\\
1.884	-0.124275\\
1.886	-0.124273\\
1.888	-0.124272\\
1.89	-0.12427\\
1.892	-0.124268\\
1.894	-0.124266\\
1.896	-0.124264\\
1.898	-0.124262\\
1.9	-0.12426\\
1.902	-0.124258\\
1.904	-0.124256\\
1.906	-0.124254\\
1.908	-0.124252\\
1.91	-0.124249\\
1.912	-0.124247\\
1.914	-0.124245\\
1.916	-0.124243\\
1.918	-0.124241\\
1.92	-0.124238\\
1.922	-0.124236\\
1.924	-0.124234\\
1.926	-0.124231\\
1.928	-0.124229\\
1.93	-0.124227\\
1.932	-0.124224\\
1.934	-0.124222\\
1.936	-0.124219\\
1.938	-0.124217\\
1.94	-0.124215\\
1.942	-0.124212\\
1.944	-0.12421\\
1.946	-0.124207\\
1.948	-0.124205\\
1.95	-0.124202\\
1.952	-0.1242\\
1.954	-0.124197\\
1.956	-0.124195\\
1.958	-0.124192\\
1.96	-0.12419\\
1.962	-0.124187\\
1.964	-0.124185\\
1.966	-0.124182\\
1.968	-0.12418\\
1.97	-0.124178\\
1.972	-0.124175\\
1.974	-0.124173\\
1.976	-0.12417\\
1.978	-0.124168\\
1.98	-0.124165\\
1.982	-0.124163\\
1.984	-0.124161\\
1.986	-0.124158\\
1.988	-0.124156\\
1.99	-0.124154\\
1.992	-0.124152\\
1.994	-0.124149\\
1.996	-0.124147\\
1.998	-0.124145\\
2	-0.124143\\
};
\addplot [color=mycolor17, line width=1.0pt, forget plot]
  table[row sep=crcr]{%
-0.124	0\\
-0.122	0\\
-0.12	0\\
-0.118	0\\
-0.116	0\\
-0.114	0\\
-0.112	0\\
-0.11	0\\
-0.108	0\\
-0.106	0\\
-0.104	0\\
-0.102	0\\
-0.1	0\\
-0.098	0\\
-0.096	0\\
-0.094	0\\
-0.092	0\\
-0.09	0\\
-0.088	0\\
-0.086	0\\
-0.084	0\\
-0.082	0\\
-0.08	0\\
-0.078	0\\
-0.076	0\\
-0.074	0\\
-0.072	0\\
-0.07	0\\
-0.068	0\\
-0.066	0\\
-0.064	0\\
-0.062	0\\
-0.06	0\\
-0.058	0\\
-0.056	0\\
-0.054	0\\
-0.052	0\\
-0.05	0\\
-0.048	0\\
-0.046	0\\
-0.044	0\\
-0.042	0\\
-0.04	0\\
-0.038	0\\
-0.036	0\\
-0.034	0\\
-0.032	0\\
-0.03	0\\
-0.028	0\\
-0.026	0\\
-0.024	0\\
-0.022	0\\
-0.02	0\\
-0.018	0\\
-0.016	0\\
-0.014	0\\
-0.012	0\\
-0.01	0\\
-0.008	0\\
-0.006	0\\
-0.004	0\\
-0.002	0\\
0	0\\
0.002	-0.00891387\\
0.004	-0.017682\\
0.006	-0.0262799\\
0.008	-0.03469\\
0.01	-0.042899\\
0.012	-0.0508957\\
0.014	-0.0586704\\
0.016	-0.0662143\\
0.018	-0.0735199\\
0.02	-0.0805801\\
0.022	-0.0873886\\
0.024	-0.0939399\\
0.026	-0.100229\\
0.028	-0.106252\\
0.03	-0.112004\\
0.032	-0.117484\\
0.034	-0.122689\\
0.036	-0.127617\\
0.038	-0.132267\\
0.04	-0.136639\\
0.042	-0.140732\\
0.044	-0.144549\\
0.046	-0.148089\\
0.048	-0.151356\\
0.05	-0.154351\\
0.052	-0.157078\\
0.054	-0.15954\\
0.056	-0.161741\\
0.058	-0.163686\\
0.06	-0.165379\\
0.062	-0.166825\\
0.064	-0.168032\\
0.066	-0.169004\\
0.068	-0.169749\\
0.07	-0.170272\\
0.072	-0.170581\\
0.074	-0.170684\\
0.076	-0.170588\\
0.078	-0.170301\\
0.08	-0.169831\\
0.082	-0.169185\\
0.084	-0.168374\\
0.086	-0.167405\\
0.088	-0.166286\\
0.09	-0.165026\\
0.092	-0.163635\\
0.094	-0.16212\\
0.096	-0.160491\\
0.098	-0.158756\\
0.1	-0.156925\\
0.102	-0.155004\\
0.104	-0.153004\\
0.106	-0.150932\\
0.108	-0.148797\\
0.11	-0.146607\\
0.112	-0.144369\\
0.114	-0.142093\\
0.116	-0.139785\\
0.118	-0.137452\\
0.12	-0.135103\\
0.122	-0.132743\\
0.124	-0.130381\\
0.126	-0.128021\\
0.128	-0.125671\\
0.13	-0.123336\\
0.132	-0.121022\\
0.134	-0.118734\\
0.136	-0.116478\\
0.138	-0.114257\\
0.14	-0.112078\\
0.142	-0.109943\\
0.144	-0.107857\\
0.146	-0.105823\\
0.148	-0.103845\\
0.15	-0.101925\\
0.152	-0.100066\\
0.154	-0.0982711\\
0.156	-0.0965418\\
0.158	-0.09488\\
0.16	-0.0932872\\
0.162	-0.0917647\\
0.164	-0.0903134\\
0.166	-0.0889342\\
0.168	-0.0876275\\
0.17	-0.0863936\\
0.172	-0.0852324\\
0.174	-0.0841437\\
0.176	-0.0831272\\
0.178	-0.0821822\\
0.18	-0.0813079\\
0.182	-0.0805032\\
0.184	-0.0797671\\
0.186	-0.0790982\\
0.188	-0.0784951\\
0.19	-0.0779561\\
0.192	-0.0774795\\
0.194	-0.0770635\\
0.196	-0.0767062\\
0.198	-0.0764055\\
0.2	-0.0761594\\
0.202	-0.0759656\\
0.204	-0.075822\\
0.206	-0.0757264\\
0.208	-0.0756763\\
0.21	-0.0756696\\
0.212	-0.075704\\
0.214	-0.075777\\
0.216	-0.0758864\\
0.218	-0.0760299\\
0.22	-0.0762053\\
0.222	-0.0764102\\
0.224	-0.0766426\\
0.226	-0.0769003\\
0.228	-0.0771812\\
0.23	-0.0774833\\
0.232	-0.0778046\\
0.234	-0.0781433\\
0.236	-0.0784975\\
0.238	-0.0788657\\
0.24	-0.0792461\\
0.242	-0.0796372\\
0.244	-0.0800376\\
0.246	-0.0804459\\
0.248	-0.080861\\
0.25	-0.0812816\\
0.252	-0.0817067\\
0.254	-0.0821355\\
0.256	-0.0825669\\
0.258	-0.0830004\\
0.26	-0.0834353\\
0.262	-0.083871\\
0.264	-0.084307\\
0.266	-0.0847431\\
0.268	-0.085179\\
0.27	-0.0856145\\
0.272	-0.0860495\\
0.274	-0.086484\\
0.276	-0.0869181\\
0.278	-0.0873519\\
0.28	-0.0877857\\
0.282	-0.0882198\\
0.284	-0.0886544\\
0.286	-0.0890901\\
0.288	-0.0895271\\
0.29	-0.0899662\\
0.292	-0.0904076\\
0.294	-0.0908522\\
0.296	-0.0913003\\
0.298	-0.0917528\\
0.3	-0.0922102\\
0.302	-0.0926731\\
0.304	-0.0931423\\
0.306	-0.0936184\\
0.308	-0.0941021\\
0.31	-0.094594\\
0.312	-0.0950948\\
0.314	-0.0956051\\
0.316	-0.0961255\\
0.318	-0.0966566\\
0.32	-0.097199\\
0.322	-0.0977531\\
0.324	-0.0983195\\
0.326	-0.0988985\\
0.328	-0.0994907\\
0.33	-0.100096\\
0.332	-0.100715\\
0.334	-0.101349\\
0.336	-0.101996\\
0.338	-0.102657\\
0.34	-0.103333\\
0.342	-0.104023\\
0.344	-0.104728\\
0.346	-0.105446\\
0.348	-0.106178\\
0.35	-0.106924\\
0.352	-0.107683\\
0.354	-0.108455\\
0.356	-0.10924\\
0.358	-0.110036\\
0.36	-0.110844\\
0.362	-0.111663\\
0.364	-0.112491\\
0.366	-0.113329\\
0.368	-0.114175\\
0.37	-0.115028\\
0.372	-0.115889\\
0.374	-0.116755\\
0.376	-0.117625\\
0.378	-0.118499\\
0.38	-0.119376\\
0.382	-0.120254\\
0.384	-0.121132\\
0.386	-0.122008\\
0.388	-0.122883\\
0.39	-0.123754\\
0.392	-0.12462\\
0.394	-0.12548\\
0.396	-0.126333\\
0.398	-0.127177\\
0.4	-0.128011\\
0.402	-0.128834\\
0.404	-0.129644\\
0.406	-0.130441\\
0.408	-0.131222\\
0.41	-0.131987\\
0.412	-0.132734\\
0.414	-0.133463\\
0.416	-0.134171\\
0.418	-0.134859\\
0.42	-0.135524\\
0.422	-0.136166\\
0.424	-0.136783\\
0.426	-0.137375\\
0.428	-0.137941\\
0.43	-0.13848\\
0.432	-0.138991\\
0.434	-0.139473\\
0.436	-0.139926\\
0.438	-0.140349\\
0.44	-0.140741\\
0.442	-0.141102\\
0.444	-0.141432\\
0.446	-0.141729\\
0.448	-0.141995\\
0.45	-0.142227\\
0.452	-0.142427\\
0.454	-0.142594\\
0.456	-0.142728\\
0.458	-0.14283\\
0.46	-0.142898\\
0.462	-0.142934\\
0.464	-0.142938\\
0.466	-0.142909\\
0.468	-0.142849\\
0.47	-0.142757\\
0.472	-0.142635\\
0.474	-0.142483\\
0.476	-0.142301\\
0.478	-0.14209\\
0.48	-0.141851\\
0.482	-0.141585\\
0.484	-0.141292\\
0.486	-0.140974\\
0.488	-0.140631\\
0.49	-0.140264\\
0.492	-0.139875\\
0.494	-0.139464\\
0.496	-0.139032\\
0.498	-0.138582\\
0.5	-0.138113\\
0.502	-0.137627\\
0.504	-0.137125\\
0.506	-0.136609\\
0.508	-0.13608\\
0.51	-0.135539\\
0.512	-0.134987\\
0.514	-0.134425\\
0.516	-0.133856\\
0.518	-0.133279\\
0.52	-0.132697\\
0.522	-0.132111\\
0.524	-0.131522\\
0.526	-0.130931\\
0.528	-0.13034\\
0.53	-0.129749\\
0.532	-0.129161\\
0.534	-0.128575\\
0.536	-0.127994\\
0.538	-0.127419\\
0.54	-0.12685\\
0.542	-0.126288\\
0.544	-0.125736\\
0.546	-0.125192\\
0.548	-0.12466\\
0.55	-0.124139\\
0.552	-0.123631\\
0.554	-0.123136\\
0.556	-0.122654\\
0.558	-0.122188\\
0.56	-0.121737\\
0.562	-0.121302\\
0.564	-0.120884\\
0.566	-0.120484\\
0.568	-0.120101\\
0.57	-0.119736\\
0.572	-0.11939\\
0.574	-0.119063\\
0.576	-0.118755\\
0.578	-0.118467\\
0.58	-0.118198\\
0.582	-0.117949\\
0.584	-0.11772\\
0.586	-0.117512\\
0.588	-0.117323\\
0.59	-0.117154\\
0.592	-0.117005\\
0.594	-0.116875\\
0.596	-0.116766\\
0.598	-0.116675\\
0.6	-0.116603\\
0.602	-0.116551\\
0.604	-0.116516\\
0.606	-0.1165\\
0.608	-0.116501\\
0.61	-0.11652\\
0.612	-0.116555\\
0.614	-0.116607\\
0.616	-0.116674\\
0.618	-0.116756\\
0.62	-0.116853\\
0.622	-0.116964\\
0.624	-0.117089\\
0.626	-0.117226\\
0.628	-0.117375\\
0.63	-0.117535\\
0.632	-0.117707\\
0.634	-0.117889\\
0.636	-0.11808\\
0.638	-0.11828\\
0.64	-0.118488\\
0.642	-0.118703\\
0.644	-0.118925\\
0.646	-0.119153\\
0.648	-0.119387\\
0.65	-0.119625\\
0.652	-0.119868\\
0.654	-0.120114\\
0.656	-0.120362\\
0.658	-0.120613\\
0.66	-0.120865\\
0.662	-0.121118\\
0.664	-0.121372\\
0.666	-0.121625\\
0.668	-0.121877\\
0.67	-0.122129\\
0.672	-0.122378\\
0.674	-0.122625\\
0.676	-0.122869\\
0.678	-0.12311\\
0.68	-0.123348\\
0.682	-0.123581\\
0.684	-0.123811\\
0.686	-0.124035\\
0.688	-0.124254\\
0.69	-0.124468\\
0.692	-0.124677\\
0.694	-0.124879\\
0.696	-0.125075\\
0.698	-0.125266\\
0.7	-0.125449\\
0.702	-0.125626\\
0.704	-0.125796\\
0.706	-0.125959\\
0.708	-0.126116\\
0.71	-0.126265\\
0.712	-0.126407\\
0.714	-0.126542\\
0.716	-0.12667\\
0.718	-0.126791\\
0.72	-0.126905\\
0.722	-0.127012\\
0.724	-0.127112\\
0.726	-0.127205\\
0.728	-0.127291\\
0.73	-0.127371\\
0.732	-0.127444\\
0.734	-0.127511\\
0.736	-0.127572\\
0.738	-0.127626\\
0.74	-0.127675\\
0.742	-0.127718\\
0.744	-0.127756\\
0.746	-0.127789\\
0.748	-0.127816\\
0.75	-0.127838\\
0.752	-0.127856\\
0.754	-0.12787\\
0.756	-0.127879\\
0.758	-0.127884\\
0.76	-0.127885\\
0.762	-0.127883\\
0.764	-0.127877\\
0.766	-0.127869\\
0.768	-0.127857\\
0.77	-0.127843\\
0.772	-0.127826\\
0.774	-0.127807\\
0.776	-0.127785\\
0.778	-0.127762\\
0.78	-0.127737\\
0.782	-0.127711\\
0.784	-0.127683\\
0.786	-0.127654\\
0.788	-0.127624\\
0.79	-0.127593\\
0.792	-0.127561\\
0.794	-0.127529\\
0.796	-0.127496\\
0.798	-0.127462\\
0.8	-0.127429\\
0.802	-0.127395\\
0.804	-0.127361\\
0.806	-0.127328\\
0.808	-0.127294\\
0.81	-0.12726\\
0.812	-0.127227\\
0.814	-0.127193\\
0.816	-0.12716\\
0.818	-0.127128\\
0.82	-0.127095\\
0.822	-0.127064\\
0.824	-0.127032\\
0.826	-0.127001\\
0.828	-0.12697\\
0.83	-0.126939\\
0.832	-0.126909\\
0.834	-0.126879\\
0.836	-0.126849\\
0.838	-0.12682\\
0.84	-0.12679\\
0.842	-0.126761\\
0.844	-0.126732\\
0.846	-0.126704\\
0.848	-0.126675\\
0.85	-0.126646\\
0.852	-0.126617\\
0.854	-0.126588\\
0.856	-0.126559\\
0.858	-0.126529\\
0.86	-0.1265\\
0.862	-0.12647\\
0.864	-0.126439\\
0.866	-0.126409\\
0.868	-0.126377\\
0.87	-0.126346\\
0.872	-0.126313\\
0.874	-0.12628\\
0.876	-0.126247\\
0.878	-0.126212\\
0.88	-0.126177\\
0.882	-0.126141\\
0.884	-0.126105\\
0.886	-0.126067\\
0.888	-0.126029\\
0.89	-0.12599\\
0.892	-0.125951\\
0.894	-0.12591\\
0.896	-0.125868\\
0.898	-0.125826\\
0.9	-0.125783\\
0.902	-0.125739\\
0.904	-0.125694\\
0.906	-0.125649\\
0.908	-0.125603\\
0.91	-0.125556\\
0.912	-0.125508\\
0.914	-0.12546\\
0.916	-0.125412\\
0.918	-0.125362\\
0.92	-0.125313\\
0.922	-0.125263\\
0.924	-0.125212\\
0.926	-0.125162\\
0.928	-0.125111\\
0.93	-0.12506\\
0.932	-0.125009\\
0.934	-0.124958\\
0.936	-0.124907\\
0.938	-0.124856\\
0.94	-0.124806\\
0.942	-0.124756\\
0.944	-0.124706\\
0.946	-0.124657\\
0.948	-0.124608\\
0.95	-0.124561\\
0.952	-0.124514\\
0.954	-0.124468\\
0.956	-0.124423\\
0.958	-0.124378\\
0.96	-0.124335\\
0.962	-0.124294\\
0.964	-0.124253\\
0.966	-0.124214\\
0.968	-0.124177\\
0.97	-0.124141\\
0.972	-0.124106\\
0.974	-0.124073\\
0.976	-0.124042\\
0.978	-0.124013\\
0.98	-0.123986\\
0.982	-0.12396\\
0.984	-0.123936\\
0.986	-0.123915\\
0.988	-0.123895\\
0.99	-0.123878\\
0.992	-0.123862\\
0.994	-0.123849\\
0.996	-0.123838\\
0.998	-0.123829\\
1	-0.123822\\
1.002	-0.123817\\
1.004	-0.123815\\
1.006	-0.123815\\
1.008	-0.123817\\
1.01	-0.123821\\
1.012	-0.123827\\
1.014	-0.123835\\
1.016	-0.123846\\
1.018	-0.123858\\
1.02	-0.123873\\
1.022	-0.123889\\
1.024	-0.123908\\
1.026	-0.123928\\
1.028	-0.12395\\
1.03	-0.123974\\
1.032	-0.124\\
1.034	-0.124027\\
1.036	-0.124056\\
1.038	-0.124086\\
1.04	-0.124118\\
1.042	-0.124151\\
1.044	-0.124186\\
1.046	-0.124222\\
1.048	-0.124258\\
1.05	-0.124296\\
1.052	-0.124335\\
1.054	-0.124374\\
1.056	-0.124415\\
1.058	-0.124456\\
1.06	-0.124497\\
1.062	-0.12454\\
1.064	-0.124582\\
1.066	-0.124625\\
1.068	-0.124668\\
1.07	-0.124711\\
1.072	-0.124754\\
1.074	-0.124797\\
1.076	-0.12484\\
1.078	-0.124882\\
1.08	-0.124925\\
1.082	-0.124967\\
1.084	-0.125008\\
1.086	-0.125049\\
1.088	-0.125089\\
1.09	-0.125128\\
1.092	-0.125167\\
1.094	-0.125205\\
1.096	-0.125242\\
1.098	-0.125277\\
1.1	-0.125312\\
1.102	-0.125346\\
1.104	-0.125378\\
1.106	-0.125409\\
1.108	-0.125439\\
1.11	-0.125468\\
1.112	-0.125495\\
1.114	-0.125521\\
1.116	-0.125546\\
1.118	-0.125569\\
1.12	-0.12559\\
1.122	-0.12561\\
1.124	-0.125629\\
1.126	-0.125646\\
1.128	-0.125661\\
1.13	-0.125675\\
1.132	-0.125687\\
1.134	-0.125698\\
1.136	-0.125707\\
1.138	-0.125715\\
1.14	-0.125721\\
1.142	-0.125726\\
1.144	-0.125729\\
1.146	-0.125731\\
1.148	-0.125731\\
1.15	-0.12573\\
1.152	-0.125728\\
1.154	-0.125724\\
1.156	-0.125719\\
1.158	-0.125712\\
1.16	-0.125705\\
1.162	-0.125696\\
1.164	-0.125686\\
1.166	-0.125675\\
1.168	-0.125663\\
1.17	-0.12565\\
1.172	-0.125635\\
1.174	-0.12562\\
1.176	-0.125604\\
1.178	-0.125588\\
1.18	-0.12557\\
1.182	-0.125552\\
1.184	-0.125533\\
1.186	-0.125514\\
1.188	-0.125494\\
1.19	-0.125473\\
1.192	-0.125453\\
1.194	-0.125431\\
1.196	-0.12541\\
1.198	-0.125388\\
1.2	-0.125366\\
1.202	-0.125343\\
1.204	-0.125321\\
1.206	-0.125299\\
1.208	-0.125276\\
1.21	-0.125254\\
1.212	-0.125231\\
1.214	-0.125209\\
1.216	-0.125187\\
1.218	-0.125165\\
1.22	-0.125143\\
1.222	-0.125121\\
1.224	-0.1251\\
1.226	-0.125079\\
1.228	-0.125059\\
1.23	-0.125038\\
1.232	-0.125019\\
1.234	-0.124999\\
1.236	-0.124981\\
1.238	-0.124962\\
1.24	-0.124944\\
1.242	-0.124927\\
1.244	-0.12491\\
1.246	-0.124894\\
1.248	-0.124879\\
1.25	-0.124863\\
1.252	-0.124849\\
1.254	-0.124835\\
1.256	-0.124822\\
1.258	-0.124809\\
1.26	-0.124797\\
1.262	-0.124786\\
1.264	-0.124775\\
1.266	-0.124764\\
1.268	-0.124755\\
1.27	-0.124745\\
1.272	-0.124737\\
1.274	-0.124729\\
1.276	-0.124721\\
1.278	-0.124714\\
1.28	-0.124708\\
1.282	-0.124702\\
1.284	-0.124697\\
1.286	-0.124692\\
1.288	-0.124688\\
1.29	-0.124684\\
1.292	-0.12468\\
1.294	-0.124677\\
1.296	-0.124674\\
1.298	-0.124672\\
1.3	-0.12467\\
1.302	-0.124669\\
1.304	-0.124667\\
1.306	-0.124666\\
1.308	-0.124666\\
1.31	-0.124665\\
1.312	-0.124665\\
1.314	-0.124665\\
1.316	-0.124665\\
1.318	-0.124666\\
1.32	-0.124667\\
1.322	-0.124667\\
1.324	-0.124668\\
1.326	-0.124669\\
1.328	-0.124671\\
1.33	-0.124672\\
1.332	-0.124673\\
1.334	-0.124675\\
1.336	-0.124676\\
1.338	-0.124678\\
1.34	-0.124679\\
1.342	-0.124681\\
1.344	-0.124682\\
1.346	-0.124684\\
1.348	-0.124685\\
1.35	-0.124687\\
1.352	-0.124688\\
1.354	-0.12469\\
1.356	-0.124691\\
1.358	-0.124693\\
1.36	-0.124694\\
1.362	-0.124695\\
1.364	-0.124696\\
1.366	-0.124697\\
1.368	-0.124698\\
1.37	-0.124699\\
1.372	-0.1247\\
1.374	-0.124701\\
1.376	-0.124702\\
1.378	-0.124702\\
1.38	-0.124703\\
1.382	-0.124703\\
1.384	-0.124704\\
1.386	-0.124704\\
1.388	-0.124704\\
1.39	-0.124704\\
1.392	-0.124705\\
1.394	-0.124705\\
1.396	-0.124705\\
1.398	-0.124705\\
1.4	-0.124705\\
1.402	-0.124705\\
1.404	-0.124705\\
1.406	-0.124705\\
1.408	-0.124704\\
1.41	-0.124704\\
1.412	-0.124704\\
1.414	-0.124704\\
1.416	-0.124704\\
1.418	-0.124704\\
1.42	-0.124703\\
1.422	-0.124703\\
1.424	-0.124703\\
1.426	-0.124703\\
1.428	-0.124703\\
1.43	-0.124703\\
1.432	-0.124703\\
1.434	-0.124703\\
1.436	-0.124703\\
1.438	-0.124703\\
1.44	-0.124703\\
1.442	-0.124703\\
1.444	-0.124703\\
1.446	-0.124704\\
1.448	-0.124704\\
1.45	-0.124704\\
1.452	-0.124705\\
1.454	-0.124705\\
1.456	-0.124706\\
1.458	-0.124706\\
1.46	-0.124707\\
1.462	-0.124707\\
1.464	-0.124708\\
1.466	-0.124708\\
1.468	-0.124709\\
1.47	-0.12471\\
1.472	-0.124711\\
1.474	-0.124711\\
1.476	-0.124712\\
1.478	-0.124713\\
1.48	-0.124714\\
1.482	-0.124714\\
1.484	-0.124715\\
1.486	-0.124716\\
1.488	-0.124717\\
1.49	-0.124718\\
1.492	-0.124718\\
1.494	-0.124719\\
1.496	-0.12472\\
1.498	-0.12472\\
1.5	-0.124721\\
1.502	-0.124721\\
1.504	-0.124722\\
1.506	-0.124722\\
1.508	-0.124722\\
1.51	-0.124723\\
1.512	-0.124723\\
1.514	-0.124723\\
1.516	-0.124723\\
1.518	-0.124723\\
1.52	-0.124723\\
1.522	-0.124722\\
1.524	-0.124722\\
1.526	-0.124721\\
1.528	-0.12472\\
1.53	-0.12472\\
1.532	-0.124719\\
1.534	-0.124718\\
1.536	-0.124716\\
1.538	-0.124715\\
1.54	-0.124714\\
1.542	-0.124712\\
1.544	-0.12471\\
1.546	-0.124708\\
1.548	-0.124706\\
1.55	-0.124704\\
1.552	-0.124702\\
1.554	-0.124699\\
1.556	-0.124697\\
1.558	-0.124694\\
1.56	-0.124691\\
1.562	-0.124688\\
1.564	-0.124685\\
1.566	-0.124682\\
1.568	-0.124678\\
1.57	-0.124675\\
1.572	-0.124671\\
1.574	-0.124667\\
1.576	-0.124663\\
1.578	-0.124659\\
1.58	-0.124655\\
1.582	-0.124651\\
1.584	-0.124646\\
1.586	-0.124642\\
1.588	-0.124637\\
1.59	-0.124632\\
1.592	-0.124628\\
1.594	-0.124623\\
1.596	-0.124618\\
1.598	-0.124613\\
1.6	-0.124608\\
1.602	-0.124603\\
1.604	-0.124598\\
1.606	-0.124593\\
1.608	-0.124588\\
1.61	-0.124582\\
1.612	-0.124577\\
1.614	-0.124572\\
1.616	-0.124567\\
1.618	-0.124562\\
1.62	-0.124556\\
1.622	-0.124551\\
1.624	-0.124546\\
1.626	-0.124541\\
1.628	-0.124536\\
1.63	-0.124531\\
1.632	-0.124526\\
1.634	-0.124521\\
1.636	-0.124516\\
1.638	-0.124511\\
1.64	-0.124506\\
1.642	-0.124501\\
1.644	-0.124497\\
1.646	-0.124492\\
1.648	-0.124488\\
1.65	-0.124483\\
1.652	-0.124479\\
1.654	-0.124475\\
1.656	-0.124471\\
1.658	-0.124467\\
1.66	-0.124463\\
1.662	-0.124459\\
1.664	-0.124455\\
1.666	-0.124452\\
1.668	-0.124448\\
1.67	-0.124445\\
1.672	-0.124442\\
1.674	-0.124439\\
1.676	-0.124436\\
1.678	-0.124433\\
1.68	-0.12443\\
1.682	-0.124427\\
1.684	-0.124425\\
1.686	-0.124423\\
1.688	-0.12442\\
1.69	-0.124418\\
1.692	-0.124416\\
1.694	-0.124414\\
1.696	-0.124412\\
1.698	-0.12441\\
1.7	-0.124409\\
1.702	-0.124407\\
1.704	-0.124406\\
1.706	-0.124404\\
1.708	-0.124403\\
1.71	-0.124402\\
1.712	-0.124401\\
1.714	-0.1244\\
1.716	-0.124399\\
1.718	-0.124398\\
1.72	-0.124397\\
1.722	-0.124396\\
1.724	-0.124395\\
1.726	-0.124395\\
1.728	-0.124394\\
1.73	-0.124394\\
1.732	-0.124393\\
1.734	-0.124393\\
1.736	-0.124392\\
1.738	-0.124392\\
1.74	-0.124392\\
1.742	-0.124391\\
1.744	-0.124391\\
1.746	-0.12439\\
1.748	-0.12439\\
1.75	-0.12439\\
1.752	-0.12439\\
1.754	-0.124389\\
1.756	-0.124389\\
1.758	-0.124389\\
1.76	-0.124388\\
1.762	-0.124388\\
1.764	-0.124388\\
1.766	-0.124387\\
1.768	-0.124387\\
1.77	-0.124386\\
1.772	-0.124386\\
1.774	-0.124385\\
1.776	-0.124385\\
1.778	-0.124384\\
1.78	-0.124384\\
1.782	-0.124383\\
1.784	-0.124382\\
1.786	-0.124382\\
1.788	-0.124381\\
1.79	-0.12438\\
1.792	-0.124379\\
1.794	-0.124378\\
1.796	-0.124378\\
1.798	-0.124377\\
1.8	-0.124375\\
1.802	-0.124374\\
1.804	-0.124373\\
1.806	-0.124372\\
1.808	-0.124371\\
1.81	-0.12437\\
1.812	-0.124368\\
1.814	-0.124367\\
1.816	-0.124365\\
1.818	-0.124364\\
1.82	-0.124362\\
1.822	-0.124361\\
1.824	-0.124359\\
1.826	-0.124358\\
1.828	-0.124356\\
1.83	-0.124354\\
1.832	-0.124352\\
1.834	-0.12435\\
1.836	-0.124349\\
1.838	-0.124347\\
1.84	-0.124345\\
1.842	-0.124343\\
1.844	-0.124341\\
1.846	-0.124339\\
1.848	-0.124337\\
1.85	-0.124335\\
1.852	-0.124333\\
1.854	-0.124331\\
1.856	-0.124328\\
1.858	-0.124326\\
1.86	-0.124324\\
1.862	-0.124322\\
1.864	-0.12432\\
1.866	-0.124318\\
1.868	-0.124315\\
1.87	-0.124313\\
1.872	-0.124311\\
1.874	-0.124309\\
1.876	-0.124306\\
1.878	-0.124304\\
1.88	-0.124302\\
1.882	-0.1243\\
1.884	-0.124297\\
1.886	-0.124295\\
1.888	-0.124293\\
1.89	-0.12429\\
1.892	-0.124288\\
1.894	-0.124286\\
1.896	-0.124284\\
1.898	-0.124282\\
1.9	-0.124279\\
1.902	-0.124277\\
1.904	-0.124275\\
1.906	-0.124273\\
1.908	-0.12427\\
1.91	-0.124268\\
1.912	-0.124266\\
1.914	-0.124264\\
1.916	-0.124262\\
1.918	-0.12426\\
1.92	-0.124257\\
1.922	-0.124255\\
1.924	-0.124253\\
1.926	-0.124251\\
1.928	-0.124249\\
1.93	-0.124247\\
1.932	-0.124245\\
1.934	-0.124243\\
1.936	-0.124241\\
1.938	-0.124239\\
1.94	-0.124237\\
1.942	-0.124235\\
1.944	-0.124233\\
1.946	-0.12423\\
1.948	-0.124228\\
1.95	-0.124226\\
1.952	-0.124224\\
1.954	-0.124222\\
1.956	-0.12422\\
1.958	-0.124218\\
1.96	-0.124217\\
1.962	-0.124215\\
1.964	-0.124213\\
1.966	-0.124211\\
1.968	-0.124209\\
1.97	-0.124207\\
1.972	-0.124205\\
1.974	-0.124203\\
1.976	-0.124201\\
1.978	-0.124199\\
1.98	-0.124197\\
1.982	-0.124195\\
1.984	-0.124193\\
1.986	-0.124191\\
1.988	-0.124189\\
1.99	-0.124187\\
1.992	-0.124185\\
1.994	-0.124183\\
1.996	-0.124181\\
1.998	-0.124179\\
2	-0.124177\\
};
\addplot [color=mycolor18, line width=1.0pt, forget plot]
  table[row sep=crcr]{%
-0.124	0\\
-0.122	0\\
-0.12	0\\
-0.118	0\\
-0.116	0\\
-0.114	0\\
-0.112	0\\
-0.11	0\\
-0.108	0\\
-0.106	0\\
-0.104	0\\
-0.102	0\\
-0.1	0\\
-0.098	0\\
-0.096	0\\
-0.094	0\\
-0.092	0\\
-0.09	0\\
-0.088	0\\
-0.086	0\\
-0.084	0\\
-0.082	0\\
-0.08	0\\
-0.078	0\\
-0.076	0\\
-0.074	0\\
-0.072	0\\
-0.07	0\\
-0.068	0\\
-0.066	0\\
-0.064	0\\
-0.062	0\\
-0.06	0\\
-0.058	0\\
-0.056	0\\
-0.054	0\\
-0.052	0\\
-0.05	0\\
-0.048	0\\
-0.046	0\\
-0.044	0\\
-0.042	0\\
-0.04	0\\
-0.038	0\\
-0.036	0\\
-0.034	0\\
-0.032	0\\
-0.03	0\\
-0.028	0\\
-0.026	0\\
-0.024	0\\
-0.022	0\\
-0.02	0\\
-0.018	0\\
-0.016	0\\
-0.014	0\\
-0.012	0\\
-0.01	0\\
-0.008	0\\
-0.006	0\\
-0.004	0\\
-0.002	0\\
0	0\\
0.002	-0.00658235\\
0.004	-0.0109242\\
0.006	-0.0137217\\
0.008	-0.0155244\\
0.01	-0.0167166\\
0.012	-0.0175495\\
0.014	-0.0181815\\
0.016	-0.0187096\\
0.018	-0.0191924\\
0.02	-0.0196646\\
0.022	-0.0201464\\
0.024	-0.0206495\\
0.026	-0.0211804\\
0.028	-0.0217422\\
0.03	-0.0223364\\
0.032	-0.0229632\\
0.034	-0.0236222\\
0.036	-0.0243124\\
0.038	-0.0250327\\
0.04	-0.0257818\\
0.042	-0.0265581\\
0.044	-0.0273604\\
0.046	-0.028187\\
0.048	-0.0290366\\
0.05	-0.0299077\\
0.052	-0.030799\\
0.054	-0.0317092\\
0.056	-0.0326368\\
0.058	-0.0335808\\
0.06	-0.0345399\\
0.062	-0.035513\\
0.064	-0.036499\\
0.066	-0.0374968\\
0.068	-0.0385056\\
0.07	-0.0395242\\
0.072	-0.0405518\\
0.074	-0.0415876\\
0.076	-0.0426306\\
0.078	-0.0436801\\
0.08	-0.0447353\\
0.082	-0.0457954\\
0.084	-0.0468597\\
0.086	-0.0479275\\
0.088	-0.048998\\
0.09	-0.0500707\\
0.092	-0.0511448\\
0.094	-0.0522197\\
0.096	-0.0532948\\
0.098	-0.0543694\\
0.1	-0.055443\\
0.102	-0.056515\\
0.104	-0.0575847\\
0.106	-0.0586516\\
0.108	-0.0597151\\
0.11	-0.0607747\\
0.112	-0.0618297\\
0.114	-0.0628797\\
0.116	-0.0639242\\
0.118	-0.0649625\\
0.12	-0.0659942\\
0.122	-0.0670187\\
0.124	-0.0680355\\
0.126	-0.0690442\\
0.128	-0.0700442\\
0.13	-0.0710351\\
0.132	-0.0720164\\
0.134	-0.0729876\\
0.136	-0.0739482\\
0.138	-0.0748979\\
0.14	-0.0758362\\
0.142	-0.0767626\\
0.144	-0.0776768\\
0.146	-0.0785783\\
0.148	-0.0794668\\
0.15	-0.0803419\\
0.152	-0.0812032\\
0.154	-0.0820504\\
0.156	-0.0828832\\
0.158	-0.0837013\\
0.16	-0.0845042\\
0.162	-0.0852919\\
0.164	-0.0860639\\
0.166	-0.0868201\\
0.168	-0.0875602\\
0.17	-0.0882841\\
0.172	-0.0889914\\
0.174	-0.089682\\
0.176	-0.0903558\\
0.178	-0.0910126\\
0.18	-0.0916523\\
0.182	-0.0922747\\
0.184	-0.0928799\\
0.186	-0.0934677\\
0.188	-0.094038\\
0.19	-0.0945909\\
0.192	-0.0951263\\
0.194	-0.0956442\\
0.196	-0.0961447\\
0.198	-0.0966277\\
0.2	-0.0970934\\
0.202	-0.0975419\\
0.204	-0.0979731\\
0.206	-0.0983873\\
0.208	-0.0987846\\
0.21	-0.099165\\
0.212	-0.0995288\\
0.214	-0.0998762\\
0.216	-0.100207\\
0.218	-0.100522\\
0.22	-0.100822\\
0.222	-0.101105\\
0.224	-0.101374\\
0.226	-0.101627\\
0.228	-0.101866\\
0.23	-0.10209\\
0.232	-0.1023\\
0.234	-0.102496\\
0.236	-0.102679\\
0.238	-0.102848\\
0.24	-0.103005\\
0.242	-0.103149\\
0.244	-0.103281\\
0.246	-0.103401\\
0.248	-0.10351\\
0.25	-0.103608\\
0.252	-0.103696\\
0.254	-0.103773\\
0.256	-0.10384\\
0.258	-0.103899\\
0.26	-0.103948\\
0.262	-0.103989\\
0.264	-0.104022\\
0.266	-0.104047\\
0.268	-0.104065\\
0.27	-0.104076\\
0.272	-0.104081\\
0.274	-0.104079\\
0.276	-0.104073\\
0.278	-0.104061\\
0.28	-0.104044\\
0.282	-0.104024\\
0.284	-0.103999\\
0.286	-0.103971\\
0.288	-0.10394\\
0.29	-0.103907\\
0.292	-0.103871\\
0.294	-0.103833\\
0.296	-0.103794\\
0.298	-0.103754\\
0.3	-0.103713\\
0.302	-0.103672\\
0.304	-0.103631\\
0.306	-0.103591\\
0.308	-0.103551\\
0.31	-0.103512\\
0.312	-0.103475\\
0.314	-0.10344\\
0.316	-0.103406\\
0.318	-0.103375\\
0.32	-0.103347\\
0.322	-0.103321\\
0.324	-0.103299\\
0.326	-0.103281\\
0.328	-0.103265\\
0.33	-0.103254\\
0.332	-0.103248\\
0.334	-0.103245\\
0.336	-0.103247\\
0.338	-0.103254\\
0.34	-0.103266\\
0.342	-0.103284\\
0.344	-0.103306\\
0.346	-0.103334\\
0.348	-0.103368\\
0.35	-0.103407\\
0.352	-0.103453\\
0.354	-0.103504\\
0.356	-0.103561\\
0.358	-0.103625\\
0.36	-0.103695\\
0.362	-0.103771\\
0.364	-0.103854\\
0.366	-0.103943\\
0.368	-0.104038\\
0.37	-0.10414\\
0.372	-0.104248\\
0.374	-0.104363\\
0.376	-0.104484\\
0.378	-0.104612\\
0.38	-0.104746\\
0.382	-0.104887\\
0.384	-0.105034\\
0.386	-0.105187\\
0.388	-0.105346\\
0.39	-0.105512\\
0.392	-0.105683\\
0.394	-0.105861\\
0.396	-0.106044\\
0.398	-0.106234\\
0.4	-0.106429\\
0.402	-0.106629\\
0.404	-0.106835\\
0.406	-0.107046\\
0.408	-0.107262\\
0.41	-0.107484\\
0.412	-0.10771\\
0.414	-0.107941\\
0.416	-0.108176\\
0.418	-0.108416\\
0.42	-0.10866\\
0.422	-0.108908\\
0.424	-0.10916\\
0.426	-0.109415\\
0.428	-0.109674\\
0.43	-0.109936\\
0.432	-0.110202\\
0.434	-0.11047\\
0.436	-0.110741\\
0.438	-0.111015\\
0.44	-0.111291\\
0.442	-0.111568\\
0.444	-0.111848\\
0.446	-0.11213\\
0.448	-0.112413\\
0.45	-0.112697\\
0.452	-0.112983\\
0.454	-0.113269\\
0.456	-0.113556\\
0.458	-0.113844\\
0.46	-0.114131\\
0.462	-0.114419\\
0.464	-0.114707\\
0.466	-0.114994\\
0.468	-0.115281\\
0.47	-0.115567\\
0.472	-0.115852\\
0.474	-0.116137\\
0.476	-0.116419\\
0.478	-0.116701\\
0.48	-0.11698\\
0.482	-0.117258\\
0.484	-0.117534\\
0.486	-0.117808\\
0.488	-0.118079\\
0.49	-0.118347\\
0.492	-0.118614\\
0.494	-0.118877\\
0.496	-0.119137\\
0.498	-0.119394\\
0.5	-0.119648\\
0.502	-0.119898\\
0.504	-0.120145\\
0.506	-0.120388\\
0.508	-0.120627\\
0.51	-0.120863\\
0.512	-0.121094\\
0.514	-0.121321\\
0.516	-0.121544\\
0.518	-0.121763\\
0.52	-0.121977\\
0.522	-0.122187\\
0.524	-0.122392\\
0.526	-0.122592\\
0.528	-0.122788\\
0.53	-0.122979\\
0.532	-0.123165\\
0.534	-0.123346\\
0.536	-0.123522\\
0.538	-0.123693\\
0.54	-0.123859\\
0.542	-0.12402\\
0.544	-0.124176\\
0.546	-0.124327\\
0.548	-0.124473\\
0.55	-0.124613\\
0.552	-0.124749\\
0.554	-0.124879\\
0.556	-0.125004\\
0.558	-0.125124\\
0.56	-0.125239\\
0.562	-0.125349\\
0.564	-0.125454\\
0.566	-0.125554\\
0.568	-0.125649\\
0.57	-0.125739\\
0.572	-0.125824\\
0.574	-0.125904\\
0.576	-0.12598\\
0.578	-0.12605\\
0.58	-0.126116\\
0.582	-0.126178\\
0.584	-0.126235\\
0.586	-0.126288\\
0.588	-0.126336\\
0.59	-0.12638\\
0.592	-0.126419\\
0.594	-0.126455\\
0.596	-0.126486\\
0.598	-0.126514\\
0.6	-0.126538\\
0.602	-0.126558\\
0.604	-0.126574\\
0.606	-0.126587\\
0.608	-0.126597\\
0.61	-0.126603\\
0.612	-0.126606\\
0.614	-0.126606\\
0.616	-0.126602\\
0.618	-0.126596\\
0.62	-0.126587\\
0.622	-0.126576\\
0.624	-0.126562\\
0.626	-0.126545\\
0.628	-0.126526\\
0.63	-0.126505\\
0.632	-0.126482\\
0.634	-0.126457\\
0.636	-0.12643\\
0.638	-0.126401\\
0.64	-0.12637\\
0.642	-0.126338\\
0.644	-0.126305\\
0.646	-0.12627\\
0.648	-0.126234\\
0.65	-0.126197\\
0.652	-0.126159\\
0.654	-0.12612\\
0.656	-0.126081\\
0.658	-0.12604\\
0.66	-0.125999\\
0.662	-0.125958\\
0.664	-0.125916\\
0.666	-0.125874\\
0.668	-0.125832\\
0.67	-0.125789\\
0.672	-0.125747\\
0.674	-0.125704\\
0.676	-0.125662\\
0.678	-0.12562\\
0.68	-0.125579\\
0.682	-0.125537\\
0.684	-0.125496\\
0.686	-0.125456\\
0.688	-0.125416\\
0.69	-0.125377\\
0.692	-0.125339\\
0.694	-0.125302\\
0.696	-0.125265\\
0.698	-0.125229\\
0.7	-0.125194\\
0.702	-0.125161\\
0.704	-0.125128\\
0.706	-0.125096\\
0.708	-0.125066\\
0.71	-0.125036\\
0.712	-0.125008\\
0.714	-0.124981\\
0.716	-0.124955\\
0.718	-0.124931\\
0.72	-0.124908\\
0.722	-0.124886\\
0.724	-0.124865\\
0.726	-0.124846\\
0.728	-0.124828\\
0.73	-0.124812\\
0.732	-0.124797\\
0.734	-0.124783\\
0.736	-0.124771\\
0.738	-0.12476\\
0.74	-0.124751\\
0.742	-0.124743\\
0.744	-0.124736\\
0.746	-0.124731\\
0.748	-0.124727\\
0.75	-0.124724\\
0.752	-0.124722\\
0.754	-0.124722\\
0.756	-0.124724\\
0.758	-0.124726\\
0.76	-0.12473\\
0.762	-0.124735\\
0.764	-0.124741\\
0.766	-0.124748\\
0.768	-0.124756\\
0.77	-0.124766\\
0.772	-0.124776\\
0.774	-0.124788\\
0.776	-0.124801\\
0.778	-0.124814\\
0.78	-0.124829\\
0.782	-0.124844\\
0.784	-0.12486\\
0.786	-0.124877\\
0.788	-0.124895\\
0.79	-0.124914\\
0.792	-0.124933\\
0.794	-0.124953\\
0.796	-0.124973\\
0.798	-0.124994\\
0.8	-0.125016\\
0.802	-0.125038\\
0.804	-0.125061\\
0.806	-0.125084\\
0.808	-0.125107\\
0.81	-0.125131\\
0.812	-0.125155\\
0.814	-0.125179\\
0.816	-0.125204\\
0.818	-0.125229\\
0.82	-0.125253\\
0.822	-0.125278\\
0.824	-0.125303\\
0.826	-0.125328\\
0.828	-0.125353\\
0.83	-0.125378\\
0.832	-0.125403\\
0.834	-0.125428\\
0.836	-0.125452\\
0.838	-0.125477\\
0.84	-0.125501\\
0.842	-0.125524\\
0.844	-0.125548\\
0.846	-0.125571\\
0.848	-0.125594\\
0.85	-0.125617\\
0.852	-0.125639\\
0.854	-0.12566\\
0.856	-0.125682\\
0.858	-0.125702\\
0.86	-0.125722\\
0.862	-0.125742\\
0.864	-0.125761\\
0.866	-0.12578\\
0.868	-0.125798\\
0.87	-0.125815\\
0.872	-0.125832\\
0.874	-0.125848\\
0.876	-0.125863\\
0.878	-0.125878\\
0.88	-0.125892\\
0.882	-0.125906\\
0.884	-0.125918\\
0.886	-0.12593\\
0.888	-0.125942\\
0.89	-0.125952\\
0.892	-0.125962\\
0.894	-0.125971\\
0.896	-0.125979\\
0.898	-0.125987\\
0.9	-0.125994\\
0.902	-0.126\\
0.904	-0.126005\\
0.906	-0.126009\\
0.908	-0.126013\\
0.91	-0.126016\\
0.912	-0.126019\\
0.914	-0.12602\\
0.916	-0.126021\\
0.918	-0.126021\\
0.92	-0.126021\\
0.922	-0.126019\\
0.924	-0.126017\\
0.926	-0.126015\\
0.928	-0.126011\\
0.93	-0.126007\\
0.932	-0.126002\\
0.934	-0.125997\\
0.936	-0.125991\\
0.938	-0.125984\\
0.94	-0.125977\\
0.942	-0.125969\\
0.944	-0.125961\\
0.946	-0.125952\\
0.948	-0.125942\\
0.95	-0.125932\\
0.952	-0.125922\\
0.954	-0.125911\\
0.956	-0.125899\\
0.958	-0.125887\\
0.96	-0.125874\\
0.962	-0.125861\\
0.964	-0.125848\\
0.966	-0.125834\\
0.968	-0.12582\\
0.97	-0.125805\\
0.972	-0.125791\\
0.974	-0.125775\\
0.976	-0.12576\\
0.978	-0.125744\\
0.98	-0.125728\\
0.982	-0.125712\\
0.984	-0.125695\\
0.986	-0.125679\\
0.988	-0.125662\\
0.99	-0.125644\\
0.992	-0.125627\\
0.994	-0.12561\\
0.996	-0.125592\\
0.998	-0.125575\\
1	-0.125557\\
1.002	-0.125539\\
1.004	-0.125521\\
1.006	-0.125504\\
1.008	-0.125486\\
1.01	-0.125468\\
1.012	-0.12545\\
1.014	-0.125432\\
1.016	-0.125414\\
1.018	-0.125397\\
1.02	-0.125379\\
1.022	-0.125362\\
1.024	-0.125344\\
1.026	-0.125327\\
1.028	-0.12531\\
1.03	-0.125293\\
1.032	-0.125276\\
1.034	-0.12526\\
1.036	-0.125243\\
1.038	-0.125227\\
1.04	-0.125211\\
1.042	-0.125196\\
1.044	-0.12518\\
1.046	-0.125165\\
1.048	-0.12515\\
1.05	-0.125135\\
1.052	-0.125121\\
1.054	-0.125107\\
1.056	-0.125093\\
1.058	-0.125079\\
1.06	-0.125066\\
1.062	-0.125053\\
1.064	-0.125041\\
1.066	-0.125029\\
1.068	-0.125017\\
1.07	-0.125005\\
1.072	-0.124994\\
1.074	-0.124983\\
1.076	-0.124972\\
1.078	-0.124962\\
1.08	-0.124953\\
1.082	-0.124943\\
1.084	-0.124934\\
1.086	-0.124925\\
1.088	-0.124917\\
1.09	-0.124909\\
1.092	-0.124901\\
1.094	-0.124894\\
1.096	-0.124887\\
1.098	-0.12488\\
1.1	-0.124874\\
1.102	-0.124868\\
1.104	-0.124863\\
1.106	-0.124858\\
1.108	-0.124853\\
1.11	-0.124848\\
1.112	-0.124844\\
1.114	-0.124841\\
1.116	-0.124837\\
1.118	-0.124834\\
1.12	-0.124831\\
1.122	-0.124829\\
1.124	-0.124826\\
1.126	-0.124824\\
1.128	-0.124823\\
1.13	-0.124822\\
1.132	-0.12482\\
1.134	-0.12482\\
1.136	-0.124819\\
1.138	-0.124819\\
1.14	-0.124819\\
1.142	-0.124819\\
1.144	-0.12482\\
1.146	-0.12482\\
1.148	-0.124821\\
1.15	-0.124823\\
1.152	-0.124824\\
1.154	-0.124825\\
1.156	-0.124827\\
1.158	-0.124829\\
1.16	-0.124831\\
1.162	-0.124833\\
1.164	-0.124836\\
1.166	-0.124838\\
1.168	-0.124841\\
1.17	-0.124844\\
1.172	-0.124847\\
1.174	-0.12485\\
1.176	-0.124853\\
1.178	-0.124856\\
1.18	-0.124859\\
1.182	-0.124863\\
1.184	-0.124866\\
1.186	-0.124869\\
1.188	-0.124873\\
1.19	-0.124877\\
1.192	-0.12488\\
1.194	-0.124884\\
1.196	-0.124887\\
1.198	-0.124891\\
1.2	-0.124895\\
1.202	-0.124898\\
1.204	-0.124902\\
1.206	-0.124905\\
1.208	-0.124909\\
1.21	-0.124913\\
1.212	-0.124916\\
1.214	-0.124919\\
1.216	-0.124923\\
1.218	-0.124926\\
1.22	-0.124929\\
1.222	-0.124933\\
1.224	-0.124936\\
1.226	-0.124939\\
1.228	-0.124942\\
1.23	-0.124944\\
1.232	-0.124947\\
1.234	-0.12495\\
1.236	-0.124952\\
1.238	-0.124955\\
1.24	-0.124957\\
1.242	-0.124959\\
1.244	-0.124961\\
1.246	-0.124963\\
1.248	-0.124965\\
1.25	-0.124967\\
1.252	-0.124968\\
1.254	-0.12497\\
1.256	-0.124971\\
1.258	-0.124972\\
1.26	-0.124973\\
1.262	-0.124974\\
1.264	-0.124975\\
1.266	-0.124976\\
1.268	-0.124976\\
1.27	-0.124976\\
1.272	-0.124977\\
1.274	-0.124977\\
1.276	-0.124977\\
1.278	-0.124976\\
1.28	-0.124976\\
1.282	-0.124976\\
1.284	-0.124975\\
1.286	-0.124974\\
1.288	-0.124973\\
1.29	-0.124972\\
1.292	-0.124971\\
1.294	-0.12497\\
1.296	-0.124969\\
1.298	-0.124967\\
1.3	-0.124965\\
1.302	-0.124964\\
1.304	-0.124962\\
1.306	-0.12496\\
1.308	-0.124958\\
1.31	-0.124955\\
1.312	-0.124953\\
1.314	-0.124951\\
1.316	-0.124948\\
1.318	-0.124946\\
1.32	-0.124943\\
1.322	-0.12494\\
1.324	-0.124937\\
1.326	-0.124934\\
1.328	-0.124931\\
1.33	-0.124928\\
1.332	-0.124925\\
1.334	-0.124921\\
1.336	-0.124918\\
1.338	-0.124915\\
1.34	-0.124911\\
1.342	-0.124908\\
1.344	-0.124904\\
1.346	-0.124901\\
1.348	-0.124897\\
1.35	-0.124893\\
1.352	-0.124889\\
1.354	-0.124886\\
1.356	-0.124882\\
1.358	-0.124878\\
1.36	-0.124874\\
1.362	-0.12487\\
1.364	-0.124866\\
1.366	-0.124862\\
1.368	-0.124858\\
1.37	-0.124854\\
1.372	-0.124851\\
1.374	-0.124847\\
1.376	-0.124843\\
1.378	-0.124839\\
1.38	-0.124835\\
1.382	-0.124831\\
1.384	-0.124827\\
1.386	-0.124823\\
1.388	-0.124819\\
1.39	-0.124815\\
1.392	-0.124811\\
1.394	-0.124808\\
1.396	-0.124804\\
1.398	-0.1248\\
1.4	-0.124796\\
1.402	-0.124793\\
1.404	-0.124789\\
1.406	-0.124785\\
1.408	-0.124782\\
1.41	-0.124778\\
1.412	-0.124775\\
1.414	-0.124771\\
1.416	-0.124768\\
1.418	-0.124764\\
1.42	-0.124761\\
1.422	-0.124758\\
1.424	-0.124754\\
1.426	-0.124751\\
1.428	-0.124748\\
1.43	-0.124745\\
1.432	-0.124742\\
1.434	-0.124739\\
1.436	-0.124736\\
1.438	-0.124733\\
1.44	-0.12473\\
1.442	-0.124727\\
1.444	-0.124725\\
1.446	-0.124722\\
1.448	-0.124719\\
1.45	-0.124717\\
1.452	-0.124714\\
1.454	-0.124712\\
1.456	-0.124709\\
1.458	-0.124707\\
1.46	-0.124704\\
1.462	-0.124702\\
1.464	-0.1247\\
1.466	-0.124698\\
1.468	-0.124696\\
1.47	-0.124694\\
1.472	-0.124692\\
1.474	-0.12469\\
1.476	-0.124688\\
1.478	-0.124686\\
1.48	-0.124684\\
1.482	-0.124682\\
1.484	-0.12468\\
1.486	-0.124679\\
1.488	-0.124677\\
1.49	-0.124675\\
1.492	-0.124674\\
1.494	-0.124672\\
1.496	-0.12467\\
1.498	-0.124669\\
1.5	-0.124668\\
1.502	-0.124666\\
1.504	-0.124665\\
1.506	-0.124663\\
1.508	-0.124662\\
1.51	-0.124661\\
1.512	-0.124659\\
1.514	-0.124658\\
1.516	-0.124657\\
1.518	-0.124656\\
1.52	-0.124654\\
1.522	-0.124653\\
1.524	-0.124652\\
1.526	-0.124651\\
1.528	-0.12465\\
1.53	-0.124648\\
1.532	-0.124647\\
1.534	-0.124646\\
1.536	-0.124645\\
1.538	-0.124644\\
1.54	-0.124643\\
1.542	-0.124642\\
1.544	-0.124641\\
1.546	-0.12464\\
1.548	-0.124639\\
1.55	-0.124638\\
1.552	-0.124636\\
1.554	-0.124635\\
1.556	-0.124634\\
1.558	-0.124633\\
1.56	-0.124632\\
1.562	-0.124631\\
1.564	-0.12463\\
1.566	-0.124629\\
1.568	-0.124628\\
1.57	-0.124627\\
1.572	-0.124625\\
1.574	-0.124624\\
1.576	-0.124623\\
1.578	-0.124622\\
1.58	-0.124621\\
1.582	-0.12462\\
1.584	-0.124618\\
1.586	-0.124617\\
1.588	-0.124616\\
1.59	-0.124614\\
1.592	-0.124613\\
1.594	-0.124612\\
1.596	-0.124611\\
1.598	-0.124609\\
1.6	-0.124608\\
1.602	-0.124606\\
1.604	-0.124605\\
1.606	-0.124604\\
1.608	-0.124602\\
1.61	-0.124601\\
1.612	-0.124599\\
1.614	-0.124598\\
1.616	-0.124596\\
1.618	-0.124594\\
1.62	-0.124593\\
1.622	-0.124591\\
1.624	-0.12459\\
1.626	-0.124588\\
1.628	-0.124586\\
1.63	-0.124584\\
1.632	-0.124583\\
1.634	-0.124581\\
1.636	-0.124579\\
1.638	-0.124577\\
1.64	-0.124575\\
1.642	-0.124574\\
1.644	-0.124572\\
1.646	-0.12457\\
1.648	-0.124568\\
1.65	-0.124566\\
1.652	-0.124564\\
1.654	-0.124562\\
1.656	-0.12456\\
1.658	-0.124558\\
1.66	-0.124556\\
1.662	-0.124554\\
1.664	-0.124552\\
1.666	-0.12455\\
1.668	-0.124547\\
1.67	-0.124545\\
1.672	-0.124543\\
1.674	-0.124541\\
1.676	-0.124539\\
1.678	-0.124537\\
1.68	-0.124534\\
1.682	-0.124532\\
1.684	-0.12453\\
1.686	-0.124528\\
1.688	-0.124525\\
1.69	-0.124523\\
1.692	-0.124521\\
1.694	-0.124518\\
1.696	-0.124516\\
1.698	-0.124514\\
1.7	-0.124511\\
1.702	-0.124509\\
1.704	-0.124507\\
1.706	-0.124504\\
1.708	-0.124502\\
1.71	-0.124499\\
1.712	-0.124497\\
1.714	-0.124495\\
1.716	-0.124492\\
1.718	-0.12449\\
1.72	-0.124487\\
1.722	-0.124485\\
1.724	-0.124483\\
1.726	-0.12448\\
1.728	-0.124478\\
1.73	-0.124475\\
1.732	-0.124473\\
1.734	-0.124471\\
1.736	-0.124468\\
1.738	-0.124466\\
1.74	-0.124463\\
1.742	-0.124461\\
1.744	-0.124458\\
1.746	-0.124456\\
1.748	-0.124454\\
1.75	-0.124451\\
1.752	-0.124449\\
1.754	-0.124447\\
1.756	-0.124444\\
1.758	-0.124442\\
1.76	-0.124439\\
1.762	-0.124437\\
1.764	-0.124435\\
1.766	-0.124432\\
1.768	-0.12443\\
1.77	-0.124428\\
1.772	-0.124425\\
1.774	-0.124423\\
1.776	-0.124421\\
1.778	-0.124419\\
1.78	-0.124416\\
1.782	-0.124414\\
1.784	-0.124412\\
1.786	-0.12441\\
1.788	-0.124407\\
1.79	-0.124405\\
1.792	-0.124403\\
1.794	-0.124401\\
1.796	-0.124399\\
1.798	-0.124396\\
1.8	-0.124394\\
1.802	-0.124392\\
1.804	-0.12439\\
1.806	-0.124388\\
1.808	-0.124386\\
1.81	-0.124384\\
1.812	-0.124382\\
1.814	-0.12438\\
1.816	-0.124377\\
1.818	-0.124375\\
1.82	-0.124373\\
1.822	-0.124371\\
1.824	-0.124369\\
1.826	-0.124367\\
1.828	-0.124365\\
1.83	-0.124363\\
1.832	-0.124361\\
1.834	-0.124359\\
1.836	-0.124358\\
1.838	-0.124356\\
1.84	-0.124354\\
1.842	-0.124352\\
1.844	-0.12435\\
1.846	-0.124348\\
1.848	-0.124346\\
1.85	-0.124344\\
1.852	-0.124342\\
1.854	-0.12434\\
1.856	-0.124339\\
1.858	-0.124337\\
1.86	-0.124335\\
1.862	-0.124333\\
1.864	-0.124331\\
1.866	-0.124329\\
1.868	-0.124328\\
1.87	-0.124326\\
1.872	-0.124324\\
1.874	-0.124322\\
1.876	-0.12432\\
1.878	-0.124319\\
1.88	-0.124317\\
1.882	-0.124315\\
1.884	-0.124313\\
1.886	-0.124312\\
1.888	-0.12431\\
1.89	-0.124308\\
1.892	-0.124306\\
1.894	-0.124305\\
1.896	-0.124303\\
1.898	-0.124301\\
1.9	-0.124299\\
1.902	-0.124298\\
1.904	-0.124296\\
1.906	-0.124294\\
1.908	-0.124292\\
1.91	-0.124291\\
1.912	-0.124289\\
1.914	-0.124287\\
1.916	-0.124286\\
1.918	-0.124284\\
1.92	-0.124282\\
1.922	-0.12428\\
1.924	-0.124279\\
1.926	-0.124277\\
1.928	-0.124275\\
1.93	-0.124273\\
1.932	-0.124272\\
1.934	-0.12427\\
1.936	-0.124268\\
1.938	-0.124267\\
1.94	-0.124265\\
1.942	-0.124263\\
1.944	-0.124261\\
1.946	-0.12426\\
1.948	-0.124258\\
1.95	-0.124256\\
1.952	-0.124254\\
1.954	-0.124253\\
1.956	-0.124251\\
1.958	-0.124249\\
1.96	-0.124247\\
1.962	-0.124246\\
1.964	-0.124244\\
1.966	-0.124242\\
1.968	-0.12424\\
1.97	-0.124238\\
1.972	-0.124237\\
1.974	-0.124235\\
1.976	-0.124233\\
1.978	-0.124231\\
1.98	-0.12423\\
1.982	-0.124228\\
1.984	-0.124226\\
1.986	-0.124224\\
1.988	-0.124222\\
1.99	-0.124221\\
1.992	-0.124219\\
1.994	-0.124217\\
1.996	-0.124215\\
1.998	-0.124213\\
2	-0.124211\\
};
\addplot [color=mycolor19, line width=1.0pt, forget plot]
  table[row sep=crcr]{%
-0.124	0\\
-0.122	0\\
-0.12	0\\
-0.118	0\\
-0.116	0\\
-0.114	0\\
-0.112	0\\
-0.11	0\\
-0.108	0\\
-0.106	0\\
-0.104	0\\
-0.102	0\\
-0.1	0\\
-0.098	0\\
-0.096	0\\
-0.094	0\\
-0.092	0\\
-0.09	0\\
-0.088	0\\
-0.086	0\\
-0.084	0\\
-0.082	0\\
-0.08	0\\
-0.078	0\\
-0.076	0\\
-0.074	0\\
-0.072	0\\
-0.07	0\\
-0.068	0\\
-0.066	0\\
-0.064	0\\
-0.062	0\\
-0.06	0\\
-0.058	0\\
-0.056	0\\
-0.054	0\\
-0.052	0\\
-0.05	0\\
-0.048	0\\
-0.046	0\\
-0.044	0\\
-0.042	0\\
-0.04	0\\
-0.038	0\\
-0.036	0\\
-0.034	0\\
-0.032	0\\
-0.03	0\\
-0.028	0\\
-0.026	0\\
-0.024	0\\
-0.022	0\\
-0.02	0\\
-0.018	0\\
-0.016	0\\
-0.014	0\\
-0.012	0\\
-0.01	0\\
-0.008	0\\
-0.006	0\\
-0.004	0\\
-0.002	0\\
0	0\\
0.002	-0.0133391\\
0.004	-0.0184799\\
0.006	-0.0200856\\
0.008	-0.0202541\\
0.01	-0.0199199\\
0.012	-0.0194894\\
0.014	-0.0191307\\
0.016	-0.0189053\\
0.018	-0.0188275\\
0.02	-0.0188925\\
0.022	-0.019088\\
0.024	-0.0193999\\
0.026	-0.0198143\\
0.028	-0.0203187\\
0.03	-0.0209019\\
0.032	-0.0215543\\
0.034	-0.0222675\\
0.036	-0.023034\\
0.038	-0.0238477\\
0.04	-0.0247029\\
0.042	-0.0255951\\
0.044	-0.0265199\\
0.046	-0.0274738\\
0.048	-0.0284535\\
0.05	-0.0294563\\
0.052	-0.0304796\\
0.054	-0.0315212\\
0.056	-0.0325792\\
0.058	-0.0336517\\
0.06	-0.034737\\
0.062	-0.0358339\\
0.064	-0.0369408\\
0.066	-0.0380567\\
0.068	-0.0391803\\
0.07	-0.0403105\\
0.072	-0.0414465\\
0.074	-0.0425873\\
0.076	-0.043732\\
0.078	-0.0448798\\
0.08	-0.0460299\\
0.082	-0.0471815\\
0.084	-0.0483339\\
0.086	-0.0494865\\
0.088	-0.0506385\\
0.09	-0.0517893\\
0.092	-0.0529383\\
0.094	-0.0540847\\
0.096	-0.055228\\
0.098	-0.0563677\\
0.1	-0.057503\\
0.102	-0.0586334\\
0.104	-0.0597584\\
0.106	-0.0608774\\
0.108	-0.0619897\\
0.11	-0.0630949\\
0.112	-0.0641925\\
0.114	-0.0652818\\
0.116	-0.0663625\\
0.118	-0.0674338\\
0.12	-0.0684954\\
0.122	-0.0695468\\
0.124	-0.0705874\\
0.126	-0.0716168\\
0.128	-0.0726344\\
0.13	-0.07364\\
0.132	-0.0746329\\
0.134	-0.0756128\\
0.136	-0.0765793\\
0.138	-0.0775318\\
0.14	-0.0784701\\
0.142	-0.0793938\\
0.144	-0.0803025\\
0.146	-0.0811957\\
0.148	-0.0820733\\
0.15	-0.0829349\\
0.152	-0.0837801\\
0.154	-0.0846087\\
0.156	-0.0854204\\
0.158	-0.0862149\\
0.16	-0.0869921\\
0.162	-0.0877517\\
0.164	-0.0884934\\
0.166	-0.0892172\\
0.168	-0.0899229\\
0.17	-0.0906102\\
0.172	-0.0912792\\
0.174	-0.0919296\\
0.176	-0.0925614\\
0.178	-0.0931746\\
0.18	-0.0937691\\
0.182	-0.0943448\\
0.184	-0.0949018\\
0.186	-0.0954401\\
0.188	-0.0959597\\
0.19	-0.0964607\\
0.192	-0.0969431\\
0.194	-0.097407\\
0.196	-0.0978526\\
0.198	-0.09828\\
0.2	-0.0986893\\
0.202	-0.0990806\\
0.204	-0.0994542\\
0.206	-0.0998103\\
0.208	-0.100149\\
0.21	-0.100471\\
0.212	-0.100776\\
0.214	-0.101064\\
0.216	-0.101336\\
0.218	-0.101592\\
0.22	-0.101832\\
0.222	-0.102057\\
0.224	-0.102267\\
0.226	-0.102462\\
0.228	-0.102642\\
0.23	-0.102809\\
0.232	-0.102962\\
0.234	-0.103102\\
0.236	-0.103229\\
0.238	-0.103344\\
0.24	-0.103447\\
0.242	-0.103538\\
0.244	-0.103618\\
0.246	-0.103687\\
0.248	-0.103746\\
0.25	-0.103795\\
0.252	-0.103834\\
0.254	-0.103865\\
0.256	-0.103887\\
0.258	-0.103901\\
0.26	-0.103907\\
0.262	-0.103907\\
0.264	-0.103899\\
0.266	-0.103885\\
0.268	-0.103866\\
0.27	-0.103841\\
0.272	-0.103811\\
0.274	-0.103777\\
0.276	-0.103739\\
0.278	-0.103697\\
0.28	-0.103653\\
0.282	-0.103605\\
0.284	-0.103555\\
0.286	-0.103504\\
0.288	-0.103451\\
0.29	-0.103396\\
0.292	-0.103342\\
0.294	-0.103287\\
0.296	-0.103232\\
0.298	-0.103178\\
0.3	-0.103124\\
0.302	-0.103072\\
0.304	-0.103022\\
0.306	-0.102973\\
0.308	-0.102927\\
0.31	-0.102883\\
0.312	-0.102842\\
0.314	-0.102805\\
0.316	-0.102771\\
0.318	-0.10274\\
0.32	-0.102714\\
0.322	-0.102692\\
0.324	-0.102675\\
0.326	-0.102662\\
0.328	-0.102654\\
0.33	-0.102652\\
0.332	-0.102655\\
0.334	-0.102663\\
0.336	-0.102678\\
0.338	-0.102698\\
0.34	-0.102724\\
0.342	-0.102756\\
0.344	-0.102795\\
0.346	-0.102841\\
0.348	-0.102892\\
0.35	-0.102951\\
0.352	-0.103016\\
0.354	-0.103088\\
0.356	-0.103167\\
0.358	-0.103252\\
0.36	-0.103345\\
0.362	-0.103445\\
0.364	-0.103551\\
0.366	-0.103665\\
0.368	-0.103785\\
0.37	-0.103912\\
0.372	-0.104047\\
0.374	-0.104188\\
0.376	-0.104336\\
0.378	-0.10449\\
0.38	-0.104652\\
0.382	-0.10482\\
0.384	-0.104994\\
0.386	-0.105175\\
0.388	-0.105362\\
0.39	-0.105555\\
0.392	-0.105755\\
0.394	-0.10596\\
0.396	-0.106171\\
0.398	-0.106388\\
0.4	-0.106611\\
0.402	-0.106839\\
0.404	-0.107072\\
0.406	-0.10731\\
0.408	-0.107553\\
0.41	-0.1078\\
0.412	-0.108053\\
0.414	-0.108309\\
0.416	-0.10857\\
0.418	-0.108834\\
0.42	-0.109103\\
0.422	-0.109375\\
0.424	-0.10965\\
0.426	-0.109928\\
0.428	-0.11021\\
0.43	-0.110494\\
0.432	-0.11078\\
0.434	-0.111069\\
0.436	-0.11136\\
0.438	-0.111653\\
0.44	-0.111947\\
0.442	-0.112243\\
0.444	-0.11254\\
0.446	-0.112838\\
0.448	-0.113137\\
0.45	-0.113436\\
0.452	-0.113736\\
0.454	-0.114036\\
0.456	-0.114336\\
0.458	-0.114635\\
0.46	-0.114934\\
0.462	-0.115232\\
0.464	-0.115529\\
0.466	-0.115825\\
0.468	-0.11612\\
0.47	-0.116413\\
0.472	-0.116704\\
0.474	-0.116994\\
0.476	-0.117281\\
0.478	-0.117566\\
0.48	-0.117849\\
0.482	-0.118129\\
0.484	-0.118406\\
0.486	-0.11868\\
0.488	-0.118951\\
0.49	-0.119219\\
0.492	-0.119483\\
0.494	-0.119744\\
0.496	-0.120001\\
0.498	-0.120254\\
0.5	-0.120503\\
0.502	-0.120748\\
0.504	-0.120989\\
0.506	-0.121225\\
0.508	-0.121457\\
0.51	-0.121685\\
0.512	-0.121907\\
0.514	-0.122125\\
0.516	-0.122338\\
0.518	-0.122546\\
0.52	-0.12275\\
0.522	-0.122948\\
0.524	-0.123141\\
0.526	-0.123329\\
0.528	-0.123511\\
0.53	-0.123689\\
0.532	-0.123861\\
0.534	-0.124027\\
0.536	-0.124189\\
0.538	-0.124345\\
0.54	-0.124496\\
0.542	-0.124641\\
0.544	-0.124781\\
0.546	-0.124915\\
0.548	-0.125044\\
0.55	-0.125168\\
0.552	-0.125286\\
0.554	-0.125399\\
0.556	-0.125507\\
0.558	-0.12561\\
0.56	-0.125707\\
0.562	-0.125799\\
0.564	-0.125886\\
0.566	-0.125968\\
0.568	-0.126045\\
0.57	-0.126118\\
0.572	-0.126185\\
0.574	-0.126247\\
0.576	-0.126305\\
0.578	-0.126358\\
0.58	-0.126406\\
0.582	-0.12645\\
0.584	-0.12649\\
0.586	-0.126525\\
0.588	-0.126557\\
0.59	-0.126584\\
0.592	-0.126607\\
0.594	-0.126626\\
0.596	-0.126641\\
0.598	-0.126653\\
0.6	-0.126661\\
0.602	-0.126666\\
0.604	-0.126667\\
0.606	-0.126665\\
0.608	-0.12666\\
0.61	-0.126652\\
0.612	-0.126641\\
0.614	-0.126628\\
0.616	-0.126612\\
0.618	-0.126593\\
0.62	-0.126572\\
0.622	-0.126548\\
0.624	-0.126522\\
0.626	-0.126495\\
0.628	-0.126465\\
0.63	-0.126434\\
0.632	-0.126401\\
0.634	-0.126366\\
0.636	-0.12633\\
0.638	-0.126292\\
0.64	-0.126254\\
0.642	-0.126214\\
0.644	-0.126173\\
0.646	-0.126132\\
0.648	-0.126089\\
0.65	-0.126046\\
0.652	-0.126003\\
0.654	-0.125958\\
0.656	-0.125914\\
0.658	-0.125869\\
0.66	-0.125824\\
0.662	-0.125779\\
0.664	-0.125734\\
0.666	-0.12569\\
0.668	-0.125645\\
0.67	-0.125601\\
0.672	-0.125557\\
0.674	-0.125513\\
0.676	-0.12547\\
0.678	-0.125428\\
0.68	-0.125386\\
0.682	-0.125345\\
0.684	-0.125305\\
0.686	-0.125266\\
0.688	-0.125227\\
0.69	-0.12519\\
0.692	-0.125153\\
0.694	-0.125118\\
0.696	-0.125084\\
0.698	-0.125051\\
0.7	-0.125019\\
0.702	-0.124988\\
0.704	-0.124959\\
0.706	-0.124931\\
0.708	-0.124904\\
0.71	-0.124879\\
0.712	-0.124855\\
0.714	-0.124832\\
0.716	-0.124811\\
0.718	-0.124791\\
0.72	-0.124773\\
0.722	-0.124756\\
0.724	-0.124741\\
0.726	-0.124727\\
0.728	-0.124715\\
0.73	-0.124704\\
0.732	-0.124694\\
0.734	-0.124686\\
0.736	-0.12468\\
0.738	-0.124675\\
0.74	-0.124671\\
0.742	-0.124669\\
0.744	-0.124668\\
0.746	-0.124668\\
0.748	-0.12467\\
0.75	-0.124673\\
0.752	-0.124678\\
0.754	-0.124683\\
0.756	-0.12469\\
0.758	-0.124698\\
0.76	-0.124708\\
0.762	-0.124718\\
0.764	-0.12473\\
0.766	-0.124743\\
0.768	-0.124756\\
0.77	-0.124771\\
0.772	-0.124787\\
0.774	-0.124804\\
0.776	-0.124821\\
0.778	-0.12484\\
0.78	-0.124859\\
0.782	-0.124879\\
0.784	-0.1249\\
0.786	-0.124921\\
0.788	-0.124943\\
0.79	-0.124966\\
0.792	-0.124989\\
0.794	-0.125013\\
0.796	-0.125037\\
0.798	-0.125062\\
0.8	-0.125087\\
0.802	-0.125112\\
0.804	-0.125138\\
0.806	-0.125164\\
0.808	-0.12519\\
0.81	-0.125216\\
0.812	-0.125243\\
0.814	-0.12527\\
0.816	-0.125296\\
0.818	-0.125323\\
0.82	-0.125349\\
0.822	-0.125376\\
0.824	-0.125402\\
0.826	-0.125429\\
0.828	-0.125455\\
0.83	-0.125481\\
0.832	-0.125507\\
0.834	-0.125532\\
0.836	-0.125557\\
0.838	-0.125582\\
0.84	-0.125606\\
0.842	-0.12563\\
0.844	-0.125654\\
0.846	-0.125677\\
0.848	-0.125699\\
0.85	-0.125721\\
0.852	-0.125743\\
0.854	-0.125764\\
0.856	-0.125784\\
0.858	-0.125804\\
0.86	-0.125823\\
0.862	-0.125842\\
0.864	-0.125859\\
0.866	-0.125877\\
0.868	-0.125893\\
0.87	-0.125909\\
0.872	-0.125924\\
0.874	-0.125938\\
0.876	-0.125952\\
0.878	-0.125964\\
0.88	-0.125976\\
0.882	-0.125988\\
0.884	-0.125998\\
0.886	-0.126008\\
0.888	-0.126017\\
0.89	-0.126025\\
0.892	-0.126032\\
0.894	-0.126039\\
0.896	-0.126044\\
0.898	-0.126049\\
0.9	-0.126054\\
0.902	-0.126057\\
0.904	-0.12606\\
0.906	-0.126061\\
0.908	-0.126062\\
0.91	-0.126063\\
0.912	-0.126062\\
0.914	-0.126061\\
0.916	-0.126059\\
0.918	-0.126056\\
0.92	-0.126053\\
0.922	-0.126049\\
0.924	-0.126044\\
0.926	-0.126038\\
0.928	-0.126032\\
0.93	-0.126025\\
0.932	-0.126018\\
0.934	-0.126009\\
0.936	-0.126001\\
0.938	-0.125991\\
0.94	-0.125981\\
0.942	-0.125971\\
0.944	-0.12596\\
0.946	-0.125948\\
0.948	-0.125936\\
0.95	-0.125923\\
0.952	-0.12591\\
0.954	-0.125897\\
0.956	-0.125883\\
0.958	-0.125868\\
0.96	-0.125854\\
0.962	-0.125838\\
0.964	-0.125823\\
0.966	-0.125807\\
0.968	-0.125791\\
0.97	-0.125774\\
0.972	-0.125757\\
0.974	-0.12574\\
0.976	-0.125723\\
0.978	-0.125705\\
0.98	-0.125688\\
0.982	-0.12567\\
0.984	-0.125652\\
0.986	-0.125634\\
0.988	-0.125615\\
0.99	-0.125597\\
0.992	-0.125578\\
0.994	-0.12556\\
0.996	-0.125541\\
0.998	-0.125522\\
1	-0.125504\\
1.002	-0.125485\\
1.004	-0.125466\\
1.006	-0.125448\\
1.008	-0.125429\\
1.01	-0.125411\\
1.012	-0.125392\\
1.014	-0.125374\\
1.016	-0.125356\\
1.018	-0.125338\\
1.02	-0.12532\\
1.022	-0.125302\\
1.024	-0.125285\\
1.026	-0.125267\\
1.028	-0.12525\\
1.03	-0.125233\\
1.032	-0.125217\\
1.034	-0.1252\\
1.036	-0.125184\\
1.038	-0.125168\\
1.04	-0.125153\\
1.042	-0.125138\\
1.044	-0.125122\\
1.046	-0.125108\\
1.048	-0.125093\\
1.05	-0.125079\\
1.052	-0.125066\\
1.054	-0.125052\\
1.056	-0.125039\\
1.058	-0.125026\\
1.06	-0.125014\\
1.062	-0.125002\\
1.064	-0.12499\\
1.066	-0.124979\\
1.068	-0.124968\\
1.07	-0.124958\\
1.072	-0.124948\\
1.074	-0.124938\\
1.076	-0.124929\\
1.078	-0.12492\\
1.08	-0.124911\\
1.082	-0.124903\\
1.084	-0.124895\\
1.086	-0.124887\\
1.088	-0.12488\\
1.09	-0.124874\\
1.092	-0.124867\\
1.094	-0.124861\\
1.096	-0.124856\\
1.098	-0.12485\\
1.1	-0.124846\\
1.102	-0.124841\\
1.104	-0.124837\\
1.106	-0.124833\\
1.108	-0.12483\\
1.11	-0.124826\\
1.112	-0.124824\\
1.114	-0.124821\\
1.116	-0.124819\\
1.118	-0.124817\\
1.12	-0.124816\\
1.122	-0.124815\\
1.124	-0.124814\\
1.126	-0.124813\\
1.128	-0.124813\\
1.13	-0.124813\\
1.132	-0.124813\\
1.134	-0.124813\\
1.136	-0.124814\\
1.138	-0.124815\\
1.14	-0.124816\\
1.142	-0.124817\\
1.144	-0.124819\\
1.146	-0.124821\\
1.148	-0.124823\\
1.15	-0.124825\\
1.152	-0.124827\\
1.154	-0.12483\\
1.156	-0.124832\\
1.158	-0.124835\\
1.16	-0.124838\\
1.162	-0.124841\\
1.164	-0.124844\\
1.166	-0.124848\\
1.168	-0.124851\\
1.17	-0.124855\\
1.172	-0.124858\\
1.174	-0.124862\\
1.176	-0.124865\\
1.178	-0.124869\\
1.18	-0.124873\\
1.182	-0.124877\\
1.184	-0.124881\\
1.186	-0.124885\\
1.188	-0.124889\\
1.19	-0.124893\\
1.192	-0.124896\\
1.194	-0.1249\\
1.196	-0.124904\\
1.198	-0.124908\\
1.2	-0.124912\\
1.202	-0.124916\\
1.204	-0.124919\\
1.206	-0.124923\\
1.208	-0.124927\\
1.21	-0.12493\\
1.212	-0.124934\\
1.214	-0.124937\\
1.216	-0.124941\\
1.218	-0.124944\\
1.22	-0.124947\\
1.222	-0.12495\\
1.224	-0.124953\\
1.226	-0.124956\\
1.228	-0.124959\\
1.23	-0.124961\\
1.232	-0.124964\\
1.234	-0.124966\\
1.236	-0.124968\\
1.238	-0.12497\\
1.24	-0.124972\\
1.242	-0.124974\\
1.244	-0.124976\\
1.246	-0.124977\\
1.248	-0.124979\\
1.25	-0.12498\\
1.252	-0.124981\\
1.254	-0.124982\\
1.256	-0.124983\\
1.258	-0.124984\\
1.26	-0.124984\\
1.262	-0.124985\\
1.264	-0.124985\\
1.266	-0.124985\\
1.268	-0.124985\\
1.27	-0.124985\\
1.272	-0.124985\\
1.274	-0.124984\\
1.276	-0.124984\\
1.278	-0.124983\\
1.28	-0.124982\\
1.282	-0.124981\\
1.284	-0.12498\\
1.286	-0.124979\\
1.288	-0.124977\\
1.29	-0.124976\\
1.292	-0.124974\\
1.294	-0.124972\\
1.296	-0.12497\\
1.298	-0.124968\\
1.3	-0.124966\\
1.302	-0.124964\\
1.304	-0.124962\\
1.306	-0.124959\\
1.308	-0.124956\\
1.31	-0.124954\\
1.312	-0.124951\\
1.314	-0.124948\\
1.316	-0.124945\\
1.318	-0.124942\\
1.32	-0.124939\\
1.322	-0.124936\\
1.324	-0.124932\\
1.326	-0.124929\\
1.328	-0.124926\\
1.33	-0.124922\\
1.332	-0.124919\\
1.334	-0.124915\\
1.336	-0.124911\\
1.338	-0.124907\\
1.34	-0.124904\\
1.342	-0.1249\\
1.344	-0.124896\\
1.346	-0.124892\\
1.348	-0.124888\\
1.35	-0.124884\\
1.352	-0.12488\\
1.354	-0.124876\\
1.356	-0.124872\\
1.358	-0.124868\\
1.36	-0.124864\\
1.362	-0.12486\\
1.364	-0.124856\\
1.366	-0.124852\\
1.368	-0.124848\\
1.37	-0.124844\\
1.372	-0.12484\\
1.374	-0.124836\\
1.376	-0.124832\\
1.378	-0.124828\\
1.38	-0.124824\\
1.382	-0.12482\\
1.384	-0.124816\\
1.386	-0.124812\\
1.388	-0.124808\\
1.39	-0.124804\\
1.392	-0.1248\\
1.394	-0.124797\\
1.396	-0.124793\\
1.398	-0.124789\\
1.4	-0.124785\\
1.402	-0.124782\\
1.404	-0.124778\\
1.406	-0.124775\\
1.408	-0.124771\\
1.41	-0.124768\\
1.412	-0.124764\\
1.414	-0.124761\\
1.416	-0.124758\\
1.418	-0.124754\\
1.42	-0.124751\\
1.422	-0.124748\\
1.424	-0.124745\\
1.426	-0.124742\\
1.428	-0.124739\\
1.43	-0.124736\\
1.432	-0.124733\\
1.434	-0.12473\\
1.436	-0.124727\\
1.438	-0.124725\\
1.44	-0.124722\\
1.442	-0.124719\\
1.444	-0.124717\\
1.446	-0.124714\\
1.448	-0.124712\\
1.45	-0.124709\\
1.452	-0.124707\\
1.454	-0.124705\\
1.456	-0.124703\\
1.458	-0.1247\\
1.46	-0.124698\\
1.462	-0.124696\\
1.464	-0.124694\\
1.466	-0.124692\\
1.468	-0.12469\\
1.47	-0.124688\\
1.472	-0.124687\\
1.474	-0.124685\\
1.476	-0.124683\\
1.478	-0.124681\\
1.48	-0.12468\\
1.482	-0.124678\\
1.484	-0.124676\\
1.486	-0.124675\\
1.488	-0.124673\\
1.49	-0.124672\\
1.492	-0.12467\\
1.494	-0.124669\\
1.496	-0.124668\\
1.498	-0.124666\\
1.5	-0.124665\\
1.502	-0.124664\\
1.504	-0.124663\\
1.506	-0.124661\\
1.508	-0.12466\\
1.51	-0.124659\\
1.512	-0.124658\\
1.514	-0.124657\\
1.516	-0.124655\\
1.518	-0.124654\\
1.52	-0.124653\\
1.522	-0.124652\\
1.524	-0.124651\\
1.526	-0.12465\\
1.528	-0.124649\\
1.53	-0.124648\\
1.532	-0.124647\\
1.534	-0.124646\\
1.536	-0.124645\\
1.538	-0.124644\\
1.54	-0.124643\\
1.542	-0.124642\\
1.544	-0.124641\\
1.546	-0.12464\\
1.548	-0.124639\\
1.55	-0.124638\\
1.552	-0.124637\\
1.554	-0.124636\\
1.556	-0.124635\\
1.558	-0.124633\\
1.56	-0.124632\\
1.562	-0.124631\\
1.564	-0.12463\\
1.566	-0.124629\\
1.568	-0.124628\\
1.57	-0.124627\\
1.572	-0.124626\\
1.574	-0.124625\\
1.576	-0.124623\\
1.578	-0.124622\\
1.58	-0.124621\\
1.582	-0.12462\\
1.584	-0.124618\\
1.586	-0.124617\\
1.588	-0.124616\\
1.59	-0.124614\\
1.592	-0.124613\\
1.594	-0.124612\\
1.596	-0.12461\\
1.598	-0.124609\\
1.6	-0.124608\\
1.602	-0.124606\\
1.604	-0.124605\\
1.606	-0.124603\\
1.608	-0.124602\\
1.61	-0.1246\\
1.612	-0.124598\\
1.614	-0.124597\\
1.616	-0.124595\\
1.618	-0.124593\\
1.62	-0.124592\\
1.622	-0.12459\\
1.624	-0.124588\\
1.626	-0.124587\\
1.628	-0.124585\\
1.63	-0.124583\\
1.632	-0.124581\\
1.634	-0.124579\\
1.636	-0.124577\\
1.638	-0.124576\\
1.64	-0.124574\\
1.642	-0.124572\\
1.644	-0.12457\\
1.646	-0.124568\\
1.648	-0.124566\\
1.65	-0.124564\\
1.652	-0.124562\\
1.654	-0.124559\\
1.656	-0.124557\\
1.658	-0.124555\\
1.66	-0.124553\\
1.662	-0.124551\\
1.664	-0.124549\\
1.666	-0.124547\\
1.668	-0.124544\\
1.67	-0.124542\\
1.672	-0.12454\\
1.674	-0.124538\\
1.676	-0.124535\\
1.678	-0.124533\\
1.68	-0.124531\\
1.682	-0.124528\\
1.684	-0.124526\\
1.686	-0.124524\\
1.688	-0.124521\\
1.69	-0.124519\\
1.692	-0.124517\\
1.694	-0.124514\\
1.696	-0.124512\\
1.698	-0.124509\\
1.7	-0.124507\\
1.702	-0.124505\\
1.704	-0.124502\\
1.706	-0.1245\\
1.708	-0.124497\\
1.71	-0.124495\\
1.712	-0.124492\\
1.714	-0.12449\\
1.716	-0.124488\\
1.718	-0.124485\\
1.72	-0.124483\\
1.722	-0.12448\\
1.724	-0.124478\\
1.726	-0.124475\\
1.728	-0.124473\\
1.73	-0.12447\\
1.732	-0.124468\\
1.734	-0.124466\\
1.736	-0.124463\\
1.738	-0.124461\\
1.74	-0.124458\\
1.742	-0.124456\\
1.744	-0.124454\\
1.746	-0.124451\\
1.748	-0.124449\\
1.75	-0.124446\\
1.752	-0.124444\\
1.754	-0.124442\\
1.756	-0.124439\\
1.758	-0.124437\\
1.76	-0.124435\\
1.762	-0.124432\\
1.764	-0.12443\\
1.766	-0.124428\\
1.768	-0.124425\\
1.77	-0.124423\\
1.772	-0.124421\\
1.774	-0.124418\\
1.776	-0.124416\\
1.778	-0.124414\\
1.78	-0.124412\\
1.782	-0.124409\\
1.784	-0.124407\\
1.786	-0.124405\\
1.788	-0.124403\\
1.79	-0.124401\\
1.792	-0.124399\\
1.794	-0.124396\\
1.796	-0.124394\\
1.798	-0.124392\\
1.8	-0.12439\\
1.802	-0.124388\\
1.804	-0.124386\\
1.806	-0.124384\\
1.808	-0.124382\\
1.81	-0.12438\\
1.812	-0.124378\\
1.814	-0.124376\\
1.816	-0.124373\\
1.818	-0.124371\\
1.82	-0.124369\\
1.822	-0.124368\\
1.824	-0.124366\\
1.826	-0.124364\\
1.828	-0.124362\\
1.83	-0.12436\\
1.832	-0.124358\\
1.834	-0.124356\\
1.836	-0.124354\\
1.838	-0.124352\\
1.84	-0.12435\\
1.842	-0.124348\\
1.844	-0.124346\\
1.846	-0.124345\\
1.848	-0.124343\\
1.85	-0.124341\\
1.852	-0.124339\\
1.854	-0.124337\\
1.856	-0.124335\\
1.858	-0.124334\\
1.86	-0.124332\\
1.862	-0.12433\\
1.864	-0.124328\\
1.866	-0.124326\\
1.868	-0.124325\\
1.87	-0.124323\\
1.872	-0.124321\\
1.874	-0.124319\\
1.876	-0.124318\\
1.878	-0.124316\\
1.88	-0.124314\\
1.882	-0.124312\\
1.884	-0.124311\\
1.886	-0.124309\\
1.888	-0.124307\\
1.89	-0.124305\\
1.892	-0.124304\\
1.894	-0.124302\\
1.896	-0.1243\\
1.898	-0.124298\\
1.9	-0.124297\\
1.902	-0.124295\\
1.904	-0.124293\\
1.906	-0.124292\\
1.908	-0.12429\\
1.91	-0.124288\\
1.912	-0.124286\\
1.914	-0.124285\\
1.916	-0.124283\\
1.918	-0.124281\\
1.92	-0.12428\\
1.922	-0.124278\\
1.924	-0.124276\\
1.926	-0.124274\\
1.928	-0.124273\\
1.93	-0.124271\\
1.932	-0.124269\\
1.934	-0.124267\\
1.936	-0.124266\\
1.938	-0.124264\\
1.94	-0.124262\\
1.942	-0.12426\\
1.944	-0.124259\\
1.946	-0.124257\\
1.948	-0.124255\\
1.95	-0.124253\\
1.952	-0.124252\\
1.954	-0.12425\\
1.956	-0.124248\\
1.958	-0.124246\\
1.96	-0.124245\\
1.962	-0.124243\\
1.964	-0.124241\\
1.966	-0.124239\\
1.968	-0.124238\\
1.97	-0.124236\\
1.972	-0.124234\\
1.974	-0.124232\\
1.976	-0.12423\\
1.978	-0.124229\\
1.98	-0.124227\\
1.982	-0.124225\\
1.984	-0.124223\\
1.986	-0.124221\\
1.988	-0.12422\\
1.99	-0.124218\\
1.992	-0.124216\\
1.994	-0.124214\\
1.996	-0.124212\\
1.998	-0.12421\\
2	-0.124209\\
};
\addplot [color=mycolor20, line width=1.0pt, forget plot]
  table[row sep=crcr]{%
-0.124	0\\
-0.122	0\\
-0.12	0\\
-0.118	0\\
-0.116	0\\
-0.114	0\\
-0.112	0\\
-0.11	0\\
-0.108	0\\
-0.106	0\\
-0.104	0\\
-0.102	0\\
-0.1	0\\
-0.098	0\\
-0.096	0\\
-0.094	0\\
-0.092	0\\
-0.09	0\\
-0.088	0\\
-0.086	0\\
-0.084	0\\
-0.082	0\\
-0.08	0\\
-0.078	0\\
-0.076	0\\
-0.074	0\\
-0.072	0\\
-0.07	0\\
-0.068	0\\
-0.066	0\\
-0.064	0\\
-0.062	0\\
-0.06	0\\
-0.058	0\\
-0.056	0\\
-0.054	0\\
-0.052	0\\
-0.05	0\\
-0.048	0\\
-0.046	0\\
-0.044	0\\
-0.042	0\\
-0.04	0\\
-0.038	0\\
-0.036	0\\
-0.034	0\\
-0.032	0\\
-0.03	0\\
-0.028	0\\
-0.026	0\\
-0.024	0\\
-0.022	0\\
-0.02	0\\
-0.018	0\\
-0.016	0\\
-0.014	0\\
-0.012	0\\
-0.01	0\\
-0.008	0\\
-0.006	0\\
-0.004	0\\
-0.002	0\\
0	0\\
0.002	-0.000663291\\
0.004	-0.00136835\\
0.006	-0.00211059\\
0.008	-0.00288673\\
0.01	-0.0036944\\
0.012	-0.00453183\\
0.014	-0.0053977\\
0.016	-0.00629097\\
0.018	-0.00721081\\
0.02	-0.00815656\\
0.022	-0.00912766\\
0.024	-0.0101237\\
0.026	-0.0111441\\
0.028	-0.0121887\\
0.03	-0.013257\\
0.032	-0.0143486\\
0.034	-0.0154634\\
0.036	-0.0166008\\
0.038	-0.0177605\\
0.04	-0.0189422\\
0.042	-0.0201453\\
0.044	-0.0213694\\
0.046	-0.022614\\
0.048	-0.0238785\\
0.05	-0.0251624\\
0.052	-0.026465\\
0.054	-0.0277857\\
0.056	-0.0291236\\
0.058	-0.0304782\\
0.06	-0.0318485\\
0.062	-0.0332337\\
0.064	-0.0346329\\
0.066	-0.0360452\\
0.068	-0.0374696\\
0.07	-0.0389051\\
0.072	-0.0403507\\
0.074	-0.0418054\\
0.076	-0.0432679\\
0.078	-0.0447373\\
0.08	-0.0462123\\
0.082	-0.0476918\\
0.084	-0.0491746\\
0.086	-0.0506596\\
0.088	-0.0521455\\
0.09	-0.053631\\
0.092	-0.0551151\\
0.094	-0.0565963\\
0.096	-0.0580737\\
0.098	-0.0595457\\
0.1	-0.0610114\\
0.102	-0.0624694\\
0.104	-0.0639186\\
0.106	-0.0653578\\
0.108	-0.0667857\\
0.11	-0.0682013\\
0.112	-0.0696034\\
0.114	-0.0709909\\
0.116	-0.0723627\\
0.118	-0.0737177\\
0.12	-0.075055\\
0.122	-0.0763735\\
0.124	-0.0776722\\
0.126	-0.0789502\\
0.128	-0.0802066\\
0.13	-0.0814406\\
0.132	-0.0826513\\
0.134	-0.083838\\
0.136	-0.0849998\\
0.138	-0.0861362\\
0.14	-0.0872464\\
0.142	-0.0883299\\
0.144	-0.089386\\
0.146	-0.0904143\\
0.148	-0.0914143\\
0.15	-0.0923855\\
0.152	-0.0933275\\
0.154	-0.0942401\\
0.156	-0.0951228\\
0.158	-0.0959756\\
0.16	-0.0967981\\
0.162	-0.0975902\\
0.164	-0.0983518\\
0.166	-0.0990828\\
0.168	-0.0997832\\
0.17	-0.100453\\
0.172	-0.101092\\
0.174	-0.101701\\
0.176	-0.10228\\
0.178	-0.102828\\
0.18	-0.103347\\
0.182	-0.103836\\
0.184	-0.104295\\
0.186	-0.104726\\
0.188	-0.105128\\
0.19	-0.105502\\
0.192	-0.105848\\
0.194	-0.106166\\
0.196	-0.106458\\
0.198	-0.106723\\
0.2	-0.106963\\
0.202	-0.107177\\
0.204	-0.107367\\
0.206	-0.107533\\
0.208	-0.107675\\
0.21	-0.107794\\
0.212	-0.107892\\
0.214	-0.107968\\
0.216	-0.108023\\
0.218	-0.108058\\
0.22	-0.108074\\
0.222	-0.108071\\
0.224	-0.108051\\
0.226	-0.108014\\
0.228	-0.10796\\
0.23	-0.107891\\
0.232	-0.107807\\
0.234	-0.107709\\
0.236	-0.107598\\
0.238	-0.107474\\
0.24	-0.107339\\
0.242	-0.107192\\
0.244	-0.107036\\
0.246	-0.10687\\
0.248	-0.106696\\
0.25	-0.106514\\
0.252	-0.106324\\
0.254	-0.106128\\
0.256	-0.105927\\
0.258	-0.10572\\
0.26	-0.10551\\
0.262	-0.105295\\
0.264	-0.105078\\
0.266	-0.104858\\
0.268	-0.104637\\
0.27	-0.104415\\
0.272	-0.104193\\
0.274	-0.103971\\
0.276	-0.10375\\
0.278	-0.10353\\
0.28	-0.103313\\
0.282	-0.103097\\
0.284	-0.102886\\
0.286	-0.102677\\
0.288	-0.102473\\
0.29	-0.102273\\
0.292	-0.102079\\
0.294	-0.10189\\
0.296	-0.101707\\
0.298	-0.10153\\
0.3	-0.10136\\
0.302	-0.101197\\
0.304	-0.101042\\
0.306	-0.100894\\
0.308	-0.100754\\
0.31	-0.100623\\
0.312	-0.100501\\
0.314	-0.100387\\
0.316	-0.100283\\
0.318	-0.100188\\
0.32	-0.100103\\
0.322	-0.100027\\
0.324	-0.0999618\\
0.326	-0.0999066\\
0.328	-0.0998617\\
0.33	-0.0998272\\
0.332	-0.0998034\\
0.334	-0.0997901\\
0.336	-0.0997877\\
0.338	-0.099796\\
0.34	-0.0998152\\
0.342	-0.0998453\\
0.344	-0.0998863\\
0.346	-0.0999383\\
0.348	-0.100001\\
0.35	-0.100075\\
0.352	-0.10016\\
0.354	-0.100255\\
0.356	-0.100361\\
0.358	-0.100478\\
0.36	-0.100606\\
0.362	-0.100744\\
0.364	-0.100892\\
0.366	-0.101051\\
0.368	-0.10122\\
0.37	-0.101399\\
0.372	-0.101588\\
0.374	-0.101787\\
0.376	-0.101995\\
0.378	-0.102212\\
0.38	-0.102439\\
0.382	-0.102675\\
0.384	-0.10292\\
0.386	-0.103173\\
0.388	-0.103434\\
0.39	-0.103704\\
0.392	-0.103981\\
0.394	-0.104267\\
0.396	-0.104559\\
0.398	-0.104859\\
0.4	-0.105166\\
0.402	-0.105479\\
0.404	-0.105799\\
0.406	-0.106125\\
0.408	-0.106456\\
0.41	-0.106794\\
0.412	-0.107136\\
0.414	-0.107484\\
0.416	-0.107836\\
0.418	-0.108192\\
0.42	-0.108553\\
0.422	-0.108917\\
0.424	-0.109285\\
0.426	-0.109656\\
0.428	-0.110029\\
0.43	-0.110406\\
0.432	-0.110784\\
0.434	-0.111165\\
0.436	-0.111547\\
0.438	-0.11193\\
0.44	-0.112314\\
0.442	-0.112699\\
0.444	-0.113085\\
0.446	-0.11347\\
0.448	-0.113856\\
0.45	-0.11424\\
0.452	-0.114624\\
0.454	-0.115007\\
0.456	-0.115388\\
0.458	-0.115768\\
0.46	-0.116146\\
0.462	-0.116521\\
0.464	-0.116894\\
0.466	-0.117264\\
0.468	-0.117631\\
0.47	-0.117994\\
0.472	-0.118354\\
0.474	-0.11871\\
0.476	-0.119062\\
0.478	-0.119409\\
0.48	-0.119752\\
0.482	-0.12009\\
0.484	-0.120423\\
0.486	-0.12075\\
0.488	-0.121073\\
0.49	-0.121389\\
0.492	-0.1217\\
0.494	-0.122004\\
0.496	-0.122302\\
0.498	-0.122594\\
0.5	-0.12288\\
0.502	-0.123158\\
0.504	-0.12343\\
0.506	-0.123695\\
0.508	-0.123953\\
0.51	-0.124203\\
0.512	-0.124446\\
0.514	-0.124682\\
0.516	-0.12491\\
0.518	-0.125131\\
0.52	-0.125343\\
0.522	-0.125549\\
0.524	-0.125746\\
0.526	-0.125935\\
0.528	-0.126117\\
0.53	-0.126291\\
0.532	-0.126457\\
0.534	-0.126615\\
0.536	-0.126765\\
0.538	-0.126907\\
0.54	-0.127042\\
0.542	-0.127168\\
0.544	-0.127287\\
0.546	-0.127398\\
0.548	-0.127501\\
0.55	-0.127597\\
0.552	-0.127685\\
0.554	-0.127766\\
0.556	-0.12784\\
0.558	-0.127906\\
0.56	-0.127965\\
0.562	-0.128017\\
0.564	-0.128062\\
0.566	-0.1281\\
0.568	-0.128131\\
0.57	-0.128156\\
0.572	-0.128175\\
0.574	-0.128187\\
0.576	-0.128194\\
0.578	-0.128194\\
0.58	-0.128189\\
0.582	-0.128178\\
0.584	-0.128161\\
0.586	-0.12814\\
0.588	-0.128113\\
0.59	-0.128081\\
0.592	-0.128045\\
0.594	-0.128004\\
0.596	-0.127959\\
0.598	-0.12791\\
0.6	-0.127857\\
0.602	-0.127801\\
0.604	-0.12774\\
0.606	-0.127677\\
0.608	-0.12761\\
0.61	-0.127541\\
0.612	-0.127468\\
0.614	-0.127394\\
0.616	-0.127317\\
0.618	-0.127237\\
0.62	-0.127156\\
0.622	-0.127074\\
0.624	-0.126989\\
0.626	-0.126904\\
0.628	-0.126817\\
0.63	-0.126729\\
0.632	-0.126641\\
0.634	-0.126552\\
0.636	-0.126462\\
0.638	-0.126372\\
0.64	-0.126283\\
0.642	-0.126193\\
0.644	-0.126104\\
0.646	-0.126015\\
0.648	-0.125926\\
0.65	-0.125839\\
0.652	-0.125752\\
0.654	-0.125666\\
0.656	-0.125581\\
0.658	-0.125498\\
0.66	-0.125416\\
0.662	-0.125336\\
0.664	-0.125257\\
0.666	-0.12518\\
0.668	-0.125105\\
0.67	-0.125032\\
0.672	-0.124961\\
0.674	-0.124892\\
0.676	-0.124825\\
0.678	-0.12476\\
0.68	-0.124698\\
0.682	-0.124638\\
0.684	-0.124581\\
0.686	-0.124526\\
0.688	-0.124474\\
0.69	-0.124424\\
0.692	-0.124377\\
0.694	-0.124333\\
0.696	-0.124292\\
0.698	-0.124253\\
0.7	-0.124217\\
0.702	-0.124183\\
0.704	-0.124153\\
0.706	-0.124125\\
0.708	-0.1241\\
0.71	-0.124077\\
0.712	-0.124058\\
0.714	-0.124041\\
0.716	-0.124026\\
0.718	-0.124015\\
0.72	-0.124006\\
0.722	-0.123999\\
0.724	-0.123996\\
0.726	-0.123994\\
0.728	-0.123995\\
0.73	-0.123999\\
0.732	-0.124005\\
0.734	-0.124013\\
0.736	-0.124024\\
0.738	-0.124036\\
0.74	-0.124051\\
0.742	-0.124068\\
0.744	-0.124087\\
0.746	-0.124108\\
0.748	-0.12413\\
0.75	-0.124155\\
0.752	-0.124181\\
0.754	-0.124209\\
0.756	-0.124238\\
0.758	-0.124269\\
0.76	-0.124301\\
0.762	-0.124335\\
0.764	-0.12437\\
0.766	-0.124406\\
0.768	-0.124443\\
0.77	-0.124482\\
0.772	-0.124521\\
0.774	-0.124561\\
0.776	-0.124602\\
0.778	-0.124644\\
0.78	-0.124686\\
0.782	-0.124729\\
0.784	-0.124772\\
0.786	-0.124816\\
0.788	-0.12486\\
0.79	-0.124905\\
0.792	-0.124949\\
0.794	-0.124994\\
0.796	-0.125039\\
0.798	-0.125084\\
0.8	-0.125129\\
0.802	-0.125173\\
0.804	-0.125218\\
0.806	-0.125262\\
0.808	-0.125306\\
0.81	-0.12535\\
0.812	-0.125393\\
0.814	-0.125435\\
0.816	-0.125477\\
0.818	-0.125519\\
0.82	-0.12556\\
0.822	-0.1256\\
0.824	-0.12564\\
0.826	-0.125679\\
0.828	-0.125717\\
0.83	-0.125754\\
0.832	-0.12579\\
0.834	-0.125826\\
0.836	-0.12586\\
0.838	-0.125894\\
0.84	-0.125926\\
0.842	-0.125958\\
0.844	-0.125988\\
0.846	-0.126018\\
0.848	-0.126046\\
0.85	-0.126073\\
0.852	-0.126099\\
0.854	-0.126124\\
0.856	-0.126148\\
0.858	-0.126171\\
0.86	-0.126192\\
0.862	-0.126212\\
0.864	-0.126231\\
0.866	-0.126249\\
0.868	-0.126266\\
0.87	-0.126281\\
0.872	-0.126296\\
0.874	-0.126309\\
0.876	-0.126321\\
0.878	-0.126331\\
0.88	-0.126341\\
0.882	-0.126349\\
0.884	-0.126356\\
0.886	-0.126362\\
0.888	-0.126367\\
0.89	-0.126371\\
0.892	-0.126373\\
0.894	-0.126375\\
0.896	-0.126375\\
0.898	-0.126374\\
0.9	-0.126373\\
0.902	-0.12637\\
0.904	-0.126366\\
0.906	-0.126361\\
0.908	-0.126355\\
0.91	-0.126348\\
0.912	-0.12634\\
0.914	-0.126331\\
0.916	-0.126322\\
0.918	-0.126311\\
0.92	-0.1263\\
0.922	-0.126287\\
0.924	-0.126274\\
0.926	-0.12626\\
0.928	-0.126245\\
0.93	-0.12623\\
0.932	-0.126214\\
0.934	-0.126197\\
0.936	-0.126179\\
0.938	-0.126161\\
0.94	-0.126142\\
0.942	-0.126123\\
0.944	-0.126103\\
0.946	-0.126082\\
0.948	-0.126061\\
0.95	-0.12604\\
0.952	-0.126018\\
0.954	-0.125995\\
0.956	-0.125972\\
0.958	-0.125949\\
0.96	-0.125926\\
0.962	-0.125902\\
0.964	-0.125878\\
0.966	-0.125853\\
0.968	-0.125828\\
0.97	-0.125804\\
0.972	-0.125779\\
0.974	-0.125753\\
0.976	-0.125728\\
0.978	-0.125702\\
0.98	-0.125677\\
0.982	-0.125651\\
0.984	-0.125626\\
0.986	-0.1256\\
0.988	-0.125574\\
0.99	-0.125549\\
0.992	-0.125523\\
0.994	-0.125498\\
0.996	-0.125472\\
0.998	-0.125447\\
1	-0.125422\\
1.002	-0.125397\\
1.004	-0.125373\\
1.006	-0.125348\\
1.008	-0.125324\\
1.01	-0.1253\\
1.012	-0.125276\\
1.014	-0.125253\\
1.016	-0.12523\\
1.018	-0.125207\\
1.02	-0.125185\\
1.022	-0.125163\\
1.024	-0.125141\\
1.026	-0.12512\\
1.028	-0.125099\\
1.03	-0.125079\\
1.032	-0.125059\\
1.034	-0.125039\\
1.036	-0.12502\\
1.038	-0.125001\\
1.04	-0.124983\\
1.042	-0.124966\\
1.044	-0.124949\\
1.046	-0.124932\\
1.048	-0.124916\\
1.05	-0.1249\\
1.052	-0.124885\\
1.054	-0.124871\\
1.056	-0.124857\\
1.058	-0.124844\\
1.06	-0.124831\\
1.062	-0.124819\\
1.064	-0.124807\\
1.066	-0.124796\\
1.068	-0.124785\\
1.07	-0.124775\\
1.072	-0.124766\\
1.074	-0.124757\\
1.076	-0.124749\\
1.078	-0.124741\\
1.08	-0.124734\\
1.082	-0.124727\\
1.084	-0.124721\\
1.086	-0.124716\\
1.088	-0.124711\\
1.09	-0.124707\\
1.092	-0.124703\\
1.094	-0.124699\\
1.096	-0.124696\\
1.098	-0.124694\\
1.1	-0.124692\\
1.102	-0.124691\\
1.104	-0.12469\\
1.106	-0.12469\\
1.108	-0.12469\\
1.11	-0.12469\\
1.112	-0.124691\\
1.114	-0.124693\\
1.116	-0.124694\\
1.118	-0.124697\\
1.12	-0.124699\\
1.122	-0.124702\\
1.124	-0.124706\\
1.126	-0.124709\\
1.128	-0.124713\\
1.13	-0.124718\\
1.132	-0.124722\\
1.134	-0.124727\\
1.136	-0.124732\\
1.138	-0.124738\\
1.14	-0.124744\\
1.142	-0.12475\\
1.144	-0.124756\\
1.146	-0.124762\\
1.148	-0.124769\\
1.15	-0.124776\\
1.152	-0.124782\\
1.154	-0.124789\\
1.156	-0.124797\\
1.158	-0.124804\\
1.16	-0.124811\\
1.162	-0.124819\\
1.164	-0.124826\\
1.166	-0.124834\\
1.168	-0.124842\\
1.17	-0.124849\\
1.172	-0.124857\\
1.174	-0.124864\\
1.176	-0.124872\\
1.178	-0.12488\\
1.18	-0.124887\\
1.182	-0.124895\\
1.184	-0.124902\\
1.186	-0.12491\\
1.188	-0.124917\\
1.19	-0.124924\\
1.192	-0.124931\\
1.194	-0.124938\\
1.196	-0.124945\\
1.198	-0.124951\\
1.2	-0.124958\\
1.202	-0.124964\\
1.204	-0.12497\\
1.206	-0.124976\\
1.208	-0.124982\\
1.21	-0.124988\\
1.212	-0.124993\\
1.214	-0.124998\\
1.216	-0.125003\\
1.218	-0.125008\\
1.22	-0.125013\\
1.222	-0.125017\\
1.224	-0.125021\\
1.226	-0.125025\\
1.228	-0.125028\\
1.23	-0.125032\\
1.232	-0.125035\\
1.234	-0.125038\\
1.236	-0.12504\\
1.238	-0.125043\\
1.24	-0.125045\\
1.242	-0.125047\\
1.244	-0.125048\\
1.246	-0.12505\\
1.248	-0.125051\\
1.25	-0.125052\\
1.252	-0.125052\\
1.254	-0.125053\\
1.256	-0.125053\\
1.258	-0.125053\\
1.26	-0.125052\\
1.262	-0.125052\\
1.264	-0.125051\\
1.266	-0.12505\\
1.268	-0.125048\\
1.27	-0.125047\\
1.272	-0.125045\\
1.274	-0.125043\\
1.276	-0.125041\\
1.278	-0.125039\\
1.28	-0.125036\\
1.282	-0.125034\\
1.284	-0.125031\\
1.286	-0.125028\\
1.288	-0.125024\\
1.29	-0.125021\\
1.292	-0.125017\\
1.294	-0.125014\\
1.296	-0.12501\\
1.298	-0.125006\\
1.3	-0.125002\\
1.302	-0.124997\\
1.304	-0.124993\\
1.306	-0.124988\\
1.308	-0.124984\\
1.31	-0.124979\\
1.312	-0.124974\\
1.314	-0.124969\\
1.316	-0.124964\\
1.318	-0.124959\\
1.32	-0.124954\\
1.322	-0.124949\\
1.324	-0.124944\\
1.326	-0.124938\\
1.328	-0.124933\\
1.33	-0.124928\\
1.332	-0.124922\\
1.334	-0.124917\\
1.336	-0.124911\\
1.338	-0.124906\\
1.34	-0.1249\\
1.342	-0.124895\\
1.344	-0.124889\\
1.346	-0.124884\\
1.348	-0.124878\\
1.35	-0.124873\\
1.352	-0.124868\\
1.354	-0.124862\\
1.356	-0.124857\\
1.358	-0.124852\\
1.36	-0.124846\\
1.362	-0.124841\\
1.364	-0.124836\\
1.366	-0.124831\\
1.368	-0.124825\\
1.37	-0.12482\\
1.372	-0.124815\\
1.374	-0.124811\\
1.376	-0.124806\\
1.378	-0.124801\\
1.38	-0.124796\\
1.382	-0.124791\\
1.384	-0.124787\\
1.386	-0.124782\\
1.388	-0.124778\\
1.39	-0.124774\\
1.392	-0.124769\\
1.394	-0.124765\\
1.396	-0.124761\\
1.398	-0.124757\\
1.4	-0.124753\\
1.402	-0.124749\\
1.404	-0.124746\\
1.406	-0.124742\\
1.408	-0.124738\\
1.41	-0.124735\\
1.412	-0.124732\\
1.414	-0.124728\\
1.416	-0.124725\\
1.418	-0.124722\\
1.42	-0.124719\\
1.422	-0.124716\\
1.424	-0.124713\\
1.426	-0.12471\\
1.428	-0.124707\\
1.43	-0.124705\\
1.432	-0.124702\\
1.434	-0.1247\\
1.436	-0.124697\\
1.438	-0.124695\\
1.44	-0.124693\\
1.442	-0.124691\\
1.444	-0.124689\\
1.446	-0.124687\\
1.448	-0.124685\\
1.45	-0.124683\\
1.452	-0.124681\\
1.454	-0.124679\\
1.456	-0.124678\\
1.458	-0.124676\\
1.46	-0.124675\\
1.462	-0.124673\\
1.464	-0.124672\\
1.466	-0.124671\\
1.468	-0.124669\\
1.47	-0.124668\\
1.472	-0.124667\\
1.474	-0.124666\\
1.476	-0.124665\\
1.478	-0.124664\\
1.48	-0.124663\\
1.482	-0.124662\\
1.484	-0.124661\\
1.486	-0.12466\\
1.488	-0.124659\\
1.49	-0.124658\\
1.492	-0.124657\\
1.494	-0.124657\\
1.496	-0.124656\\
1.498	-0.124655\\
1.5	-0.124654\\
1.502	-0.124654\\
1.504	-0.124653\\
1.506	-0.124653\\
1.508	-0.124652\\
1.51	-0.124651\\
1.512	-0.124651\\
1.514	-0.12465\\
1.516	-0.12465\\
1.518	-0.124649\\
1.52	-0.124648\\
1.522	-0.124648\\
1.524	-0.124647\\
1.526	-0.124647\\
1.528	-0.124646\\
1.53	-0.124646\\
1.532	-0.124645\\
1.534	-0.124645\\
1.536	-0.124644\\
1.538	-0.124643\\
1.54	-0.124643\\
1.542	-0.124642\\
1.544	-0.124642\\
1.546	-0.124641\\
1.548	-0.12464\\
1.55	-0.12464\\
1.552	-0.124639\\
1.554	-0.124638\\
1.556	-0.124638\\
1.558	-0.124637\\
1.56	-0.124636\\
1.562	-0.124635\\
1.564	-0.124634\\
1.566	-0.124634\\
1.568	-0.124633\\
1.57	-0.124632\\
1.572	-0.124631\\
1.574	-0.12463\\
1.576	-0.124629\\
1.578	-0.124628\\
1.58	-0.124627\\
1.582	-0.124626\\
1.584	-0.124625\\
1.586	-0.124623\\
1.588	-0.124622\\
1.59	-0.124621\\
1.592	-0.12462\\
1.594	-0.124618\\
1.596	-0.124617\\
1.598	-0.124616\\
1.6	-0.124614\\
1.602	-0.124613\\
1.604	-0.124611\\
1.606	-0.12461\\
1.608	-0.124608\\
1.61	-0.124607\\
1.612	-0.124605\\
1.614	-0.124603\\
1.616	-0.124601\\
1.618	-0.1246\\
1.62	-0.124598\\
1.622	-0.124596\\
1.624	-0.124594\\
1.626	-0.124592\\
1.628	-0.12459\\
1.63	-0.124588\\
1.632	-0.124586\\
1.634	-0.124584\\
1.636	-0.124582\\
1.638	-0.12458\\
1.64	-0.124578\\
1.642	-0.124576\\
1.644	-0.124574\\
1.646	-0.124571\\
1.648	-0.124569\\
1.65	-0.124567\\
1.652	-0.124565\\
1.654	-0.124562\\
1.656	-0.12456\\
1.658	-0.124557\\
1.66	-0.124555\\
1.662	-0.124553\\
1.664	-0.12455\\
1.666	-0.124548\\
1.668	-0.124545\\
1.67	-0.124542\\
1.672	-0.12454\\
1.674	-0.124537\\
1.676	-0.124535\\
1.678	-0.124532\\
1.68	-0.12453\\
1.682	-0.124527\\
1.684	-0.124524\\
1.686	-0.124522\\
1.688	-0.124519\\
1.69	-0.124516\\
1.692	-0.124514\\
1.694	-0.124511\\
1.696	-0.124508\\
1.698	-0.124505\\
1.7	-0.124503\\
1.702	-0.1245\\
1.704	-0.124497\\
1.706	-0.124495\\
1.708	-0.124492\\
1.71	-0.124489\\
1.712	-0.124486\\
1.714	-0.124484\\
1.716	-0.124481\\
1.718	-0.124478\\
1.72	-0.124476\\
1.722	-0.124473\\
1.724	-0.12447\\
1.726	-0.124468\\
1.728	-0.124465\\
1.73	-0.124462\\
1.732	-0.12446\\
1.734	-0.124457\\
1.736	-0.124454\\
1.738	-0.124452\\
1.74	-0.124449\\
1.742	-0.124446\\
1.744	-0.124444\\
1.746	-0.124441\\
1.748	-0.124439\\
1.75	-0.124436\\
1.752	-0.124434\\
1.754	-0.124431\\
1.756	-0.124429\\
1.758	-0.124426\\
1.76	-0.124424\\
1.762	-0.124422\\
1.764	-0.124419\\
1.766	-0.124417\\
1.768	-0.124414\\
1.77	-0.124412\\
1.772	-0.12441\\
1.774	-0.124407\\
1.776	-0.124405\\
1.778	-0.124403\\
1.78	-0.124401\\
1.782	-0.124398\\
1.784	-0.124396\\
1.786	-0.124394\\
1.788	-0.124392\\
1.79	-0.12439\\
1.792	-0.124388\\
1.794	-0.124386\\
1.796	-0.124384\\
1.798	-0.124381\\
1.8	-0.124379\\
1.802	-0.124377\\
1.804	-0.124375\\
1.806	-0.124373\\
1.808	-0.124371\\
1.81	-0.12437\\
1.812	-0.124368\\
1.814	-0.124366\\
1.816	-0.124364\\
1.818	-0.124362\\
1.82	-0.12436\\
1.822	-0.124358\\
1.824	-0.124356\\
1.826	-0.124355\\
1.828	-0.124353\\
1.83	-0.124351\\
1.832	-0.124349\\
1.834	-0.124347\\
1.836	-0.124346\\
1.838	-0.124344\\
1.84	-0.124342\\
1.842	-0.12434\\
1.844	-0.124339\\
1.846	-0.124337\\
1.848	-0.124335\\
1.85	-0.124334\\
1.852	-0.124332\\
1.854	-0.12433\\
1.856	-0.124329\\
1.858	-0.124327\\
1.86	-0.124325\\
1.862	-0.124324\\
1.864	-0.124322\\
1.866	-0.124321\\
1.868	-0.124319\\
1.87	-0.124317\\
1.872	-0.124316\\
1.874	-0.124314\\
1.876	-0.124312\\
1.878	-0.124311\\
1.88	-0.124309\\
1.882	-0.124308\\
1.884	-0.124306\\
1.886	-0.124304\\
1.888	-0.124303\\
1.89	-0.124301\\
1.892	-0.124299\\
1.894	-0.124298\\
1.896	-0.124296\\
1.898	-0.124295\\
1.9	-0.124293\\
1.902	-0.124291\\
1.904	-0.12429\\
1.906	-0.124288\\
1.908	-0.124286\\
1.91	-0.124285\\
1.912	-0.124283\\
1.914	-0.124282\\
1.916	-0.12428\\
1.918	-0.124278\\
1.92	-0.124277\\
1.922	-0.124275\\
1.924	-0.124273\\
1.926	-0.124272\\
1.928	-0.12427\\
1.93	-0.124268\\
1.932	-0.124266\\
1.934	-0.124265\\
1.936	-0.124263\\
1.938	-0.124261\\
1.94	-0.12426\\
1.942	-0.124258\\
1.944	-0.124256\\
1.946	-0.124254\\
1.948	-0.124253\\
1.95	-0.124251\\
1.952	-0.124249\\
1.954	-0.124247\\
1.956	-0.124246\\
1.958	-0.124244\\
1.96	-0.124242\\
1.962	-0.12424\\
1.964	-0.124238\\
1.966	-0.124237\\
1.968	-0.124235\\
1.97	-0.124233\\
1.972	-0.124231\\
1.974	-0.124229\\
1.976	-0.124228\\
1.978	-0.124226\\
1.98	-0.124224\\
1.982	-0.124222\\
1.984	-0.12422\\
1.986	-0.124218\\
1.988	-0.124216\\
1.99	-0.124215\\
1.992	-0.124213\\
1.994	-0.124211\\
1.996	-0.124209\\
1.998	-0.124207\\
2	-0.124205\\
};
\addplot [color=black, line width=1.4pt, forget plot]
  table[row sep=crcr]{%
-0.124	0\\
-0.122	0\\
-0.12	0\\
-0.118	0\\
-0.116	0\\
-0.114	0\\
-0.112	0\\
-0.11	0\\
-0.108	0\\
-0.106	0\\
-0.104	0\\
-0.102	0\\
-0.1	0\\
-0.098	0\\
-0.096	0\\
-0.094	0\\
-0.092	0\\
-0.09	0\\
-0.088	0\\
-0.086	0\\
-0.084	0\\
-0.082	0\\
-0.08	0\\
-0.078	0\\
-0.076	0\\
-0.074	0\\
-0.072	0\\
-0.07	0\\
-0.068	0\\
-0.066	0\\
-0.064	0\\
-0.062	0\\
-0.06	0\\
-0.058	0\\
-0.056	0\\
-0.054	0\\
-0.052	0\\
-0.05	0\\
-0.048	0\\
-0.046	0\\
-0.044	0\\
-0.042	0\\
-0.04	0\\
-0.038	0\\
-0.036	0\\
-0.034	0\\
-0.032	0\\
-0.03	0\\
-0.028	0\\
-0.026	0\\
-0.024	0\\
-0.022	0\\
-0.02	0\\
-0.018	0\\
-0.016	0\\
-0.014	0\\
-0.012	0\\
-0.01	0\\
-0.008	0\\
-0.006	0\\
-0.004	0\\
-0.002	0\\
0	0\\
0.002	-0.00190601\\
0.004	-0.00375791\\
0.006	-0.005576\\
0.008	-0.00736965\\
0.01	-0.00914296\\
0.012	-0.0108975\\
0.014	-0.0126336\\
0.016	-0.0143512\\
0.018	-0.01605\\
0.02	-0.0177295\\
0.022	-0.0193895\\
0.024	-0.0210297\\
0.026	-0.0226498\\
0.028	-0.0242497\\
0.03	-0.0258292\\
0.032	-0.0273883\\
0.034	-0.0289269\\
0.036	-0.0304451\\
0.038	-0.0319429\\
0.04	-0.0334203\\
0.042	-0.0348774\\
0.044	-0.0363144\\
0.046	-0.0377313\\
0.048	-0.0391283\\
0.05	-0.0405055\\
0.052	-0.0418631\\
0.054	-0.0432013\\
0.056	-0.0445201\\
0.058	-0.0458199\\
0.06	-0.0471007\\
0.062	-0.0483627\\
0.064	-0.0496062\\
0.066	-0.0508313\\
0.068	-0.0520382\\
0.07	-0.053227\\
0.072	-0.0543981\\
0.074	-0.0555515\\
0.076	-0.0566874\\
0.078	-0.0578061\\
0.08	-0.0589078\\
0.082	-0.0599925\\
0.084	-0.0610606\\
0.086	-0.0621122\\
0.088	-0.0631474\\
0.09	-0.0641666\\
0.092	-0.0651698\\
0.094	-0.0661573\\
0.096	-0.0671292\\
0.098	-0.0680858\\
0.1	-0.0690272\\
0.102	-0.0699535\\
0.104	-0.0708651\\
0.106	-0.071762\\
0.108	-0.0726444\\
0.11	-0.0735126\\
0.112	-0.0743667\\
0.114	-0.0752068\\
0.116	-0.0760333\\
0.118	-0.0768462\\
0.12	-0.0776457\\
0.122	-0.078432\\
0.124	-0.0792053\\
0.126	-0.0799657\\
0.128	-0.0807135\\
0.13	-0.0814487\\
0.132	-0.0821717\\
0.134	-0.0828825\\
0.136	-0.0835813\\
0.138	-0.0842683\\
0.14	-0.0849437\\
0.142	-0.0856077\\
0.144	-0.0862603\\
0.146	-0.0869019\\
0.148	-0.0875324\\
0.15	-0.0881523\\
0.152	-0.0887615\\
0.154	-0.0893602\\
0.156	-0.0899487\\
0.158	-0.0905271\\
0.16	-0.0910956\\
0.162	-0.0916542\\
0.164	-0.0922033\\
0.166	-0.0927429\\
0.168	-0.0932733\\
0.17	-0.0937945\\
0.172	-0.0943067\\
0.174	-0.0948102\\
0.176	-0.095305\\
0.178	-0.0957913\\
0.18	-0.0962693\\
0.182	-0.0967392\\
0.184	-0.097201\\
0.186	-0.0976549\\
0.188	-0.0981012\\
0.19	-0.0985399\\
0.192	-0.0989712\\
0.194	-0.0993952\\
0.196	-0.0998122\\
0.198	-0.100222\\
0.2	-0.100625\\
0.202	-0.101022\\
0.204	-0.101412\\
0.206	-0.101795\\
0.208	-0.102173\\
0.21	-0.102544\\
0.212	-0.102909\\
0.214	-0.103268\\
0.216	-0.103622\\
0.218	-0.10397\\
0.22	-0.104312\\
0.222	-0.10465\\
0.224	-0.104982\\
0.226	-0.105308\\
0.228	-0.10563\\
0.23	-0.105947\\
0.232	-0.10626\\
0.234	-0.106567\\
0.236	-0.106871\\
0.238	-0.107169\\
0.24	-0.107464\\
0.242	-0.107754\\
0.244	-0.10804\\
0.246	-0.108322\\
0.248	-0.1086\\
0.25	-0.108875\\
0.252	-0.109145\\
0.254	-0.109413\\
0.256	-0.109676\\
0.258	-0.109936\\
0.26	-0.110193\\
0.262	-0.110446\\
0.264	-0.110696\\
0.266	-0.110943\\
0.268	-0.111187\\
0.27	-0.111428\\
0.272	-0.111666\\
0.274	-0.111901\\
0.276	-0.112133\\
0.278	-0.112363\\
0.28	-0.11259\\
0.282	-0.112814\\
0.284	-0.113035\\
0.286	-0.113254\\
0.288	-0.113471\\
0.29	-0.113685\\
0.292	-0.113897\\
0.294	-0.114106\\
0.296	-0.114314\\
0.298	-0.114518\\
0.3	-0.114721\\
0.302	-0.114921\\
0.304	-0.11512\\
0.306	-0.115316\\
0.308	-0.11551\\
0.31	-0.115702\\
0.312	-0.115892\\
0.314	-0.116079\\
0.316	-0.116265\\
0.318	-0.116449\\
0.32	-0.116631\\
0.322	-0.116811\\
0.324	-0.116989\\
0.326	-0.117165\\
0.328	-0.117339\\
0.33	-0.117511\\
0.332	-0.117681\\
0.334	-0.11785\\
0.336	-0.118016\\
0.338	-0.118181\\
0.34	-0.118344\\
0.342	-0.118505\\
0.344	-0.118664\\
0.346	-0.118821\\
0.348	-0.118976\\
0.35	-0.11913\\
0.352	-0.119281\\
0.354	-0.119431\\
0.356	-0.119579\\
0.358	-0.119725\\
0.36	-0.11987\\
0.362	-0.120012\\
0.364	-0.120153\\
0.366	-0.120291\\
0.368	-0.120428\\
0.37	-0.120563\\
0.372	-0.120696\\
0.374	-0.120828\\
0.376	-0.120957\\
0.378	-0.121084\\
0.38	-0.12121\\
0.382	-0.121334\\
0.384	-0.121456\\
0.386	-0.121576\\
0.388	-0.121694\\
0.39	-0.12181\\
0.392	-0.121924\\
0.394	-0.122037\\
0.396	-0.122147\\
0.398	-0.122256\\
0.4	-0.122363\\
0.402	-0.122467\\
0.404	-0.12257\\
0.406	-0.122671\\
0.408	-0.122771\\
0.41	-0.122868\\
0.412	-0.122963\\
0.414	-0.123057\\
0.416	-0.123148\\
0.418	-0.123238\\
0.42	-0.123326\\
0.422	-0.123412\\
0.424	-0.123496\\
0.426	-0.123579\\
0.428	-0.123659\\
0.43	-0.123738\\
0.432	-0.123815\\
0.434	-0.12389\\
0.436	-0.123963\\
0.438	-0.124034\\
0.44	-0.124104\\
0.442	-0.124172\\
0.444	-0.124238\\
0.446	-0.124302\\
0.448	-0.124365\\
0.45	-0.124426\\
0.452	-0.124485\\
0.454	-0.124543\\
0.456	-0.124598\\
0.458	-0.124653\\
0.46	-0.124705\\
0.462	-0.124756\\
0.464	-0.124806\\
0.466	-0.124853\\
0.468	-0.1249\\
0.47	-0.124944\\
0.472	-0.124988\\
0.474	-0.125029\\
0.476	-0.12507\\
0.478	-0.125108\\
0.48	-0.125146\\
0.482	-0.125182\\
0.484	-0.125216\\
0.486	-0.125249\\
0.488	-0.125281\\
0.49	-0.125312\\
0.492	-0.125341\\
0.494	-0.125369\\
0.496	-0.125396\\
0.498	-0.125421\\
0.5	-0.125446\\
0.502	-0.125469\\
0.504	-0.125491\\
0.506	-0.125512\\
0.508	-0.125532\\
0.51	-0.12555\\
0.512	-0.125568\\
0.514	-0.125585\\
0.516	-0.125601\\
0.518	-0.125615\\
0.52	-0.125629\\
0.522	-0.125642\\
0.524	-0.125654\\
0.526	-0.125665\\
0.528	-0.125676\\
0.53	-0.125685\\
0.532	-0.125694\\
0.534	-0.125702\\
0.536	-0.125709\\
0.538	-0.125716\\
0.54	-0.125722\\
0.542	-0.125727\\
0.544	-0.125732\\
0.546	-0.125736\\
0.548	-0.125739\\
0.55	-0.125742\\
0.552	-0.125744\\
0.554	-0.125746\\
0.556	-0.125747\\
0.558	-0.125748\\
0.56	-0.125748\\
0.562	-0.125748\\
0.564	-0.125748\\
0.566	-0.125747\\
0.568	-0.125745\\
0.57	-0.125744\\
0.572	-0.125742\\
0.574	-0.125739\\
0.576	-0.125737\\
0.578	-0.125734\\
0.58	-0.125731\\
0.582	-0.125727\\
0.584	-0.125723\\
0.586	-0.12572\\
0.588	-0.125715\\
0.59	-0.125711\\
0.592	-0.125707\\
0.594	-0.125702\\
0.596	-0.125698\\
0.598	-0.125693\\
0.6	-0.125688\\
0.602	-0.125683\\
0.604	-0.125678\\
0.606	-0.125672\\
0.608	-0.125667\\
0.61	-0.125662\\
0.612	-0.125656\\
0.614	-0.125651\\
0.616	-0.125646\\
0.618	-0.12564\\
0.62	-0.125635\\
0.622	-0.125629\\
0.624	-0.125624\\
0.626	-0.125619\\
0.628	-0.125613\\
0.63	-0.125608\\
0.632	-0.125603\\
0.634	-0.125598\\
0.636	-0.125592\\
0.638	-0.125587\\
0.64	-0.125582\\
0.642	-0.125577\\
0.644	-0.125573\\
0.646	-0.125568\\
0.648	-0.125563\\
0.65	-0.125559\\
0.652	-0.125554\\
0.654	-0.12555\\
0.656	-0.125545\\
0.658	-0.125541\\
0.66	-0.125537\\
0.662	-0.125533\\
0.664	-0.125529\\
0.666	-0.125526\\
0.668	-0.125522\\
0.67	-0.125518\\
0.672	-0.125515\\
0.674	-0.125512\\
0.676	-0.125508\\
0.678	-0.125505\\
0.68	-0.125502\\
0.682	-0.1255\\
0.684	-0.125497\\
0.686	-0.125494\\
0.688	-0.125492\\
0.69	-0.125489\\
0.692	-0.125487\\
0.694	-0.125485\\
0.696	-0.125483\\
0.698	-0.125481\\
0.7	-0.125479\\
0.702	-0.125477\\
0.704	-0.125475\\
0.706	-0.125474\\
0.708	-0.125472\\
0.71	-0.125471\\
0.712	-0.125469\\
0.714	-0.125468\\
0.716	-0.125467\\
0.718	-0.125466\\
0.72	-0.125464\\
0.722	-0.125463\\
0.724	-0.125463\\
0.726	-0.125462\\
0.728	-0.125461\\
0.73	-0.12546\\
0.732	-0.125459\\
0.734	-0.125459\\
0.736	-0.125458\\
0.738	-0.125457\\
0.74	-0.125457\\
0.742	-0.125456\\
0.744	-0.125456\\
0.746	-0.125455\\
0.748	-0.125455\\
0.75	-0.125454\\
0.752	-0.125454\\
0.754	-0.125453\\
0.756	-0.125453\\
0.758	-0.125452\\
0.76	-0.125452\\
0.762	-0.125451\\
0.764	-0.125451\\
0.766	-0.12545\\
0.768	-0.12545\\
0.77	-0.125449\\
0.772	-0.125448\\
0.774	-0.125448\\
0.776	-0.125447\\
0.778	-0.125447\\
0.78	-0.125446\\
0.782	-0.125445\\
0.784	-0.125444\\
0.786	-0.125444\\
0.788	-0.125443\\
0.79	-0.125442\\
0.792	-0.125441\\
0.794	-0.12544\\
0.796	-0.125439\\
0.798	-0.125438\\
0.8	-0.125436\\
0.802	-0.125435\\
0.804	-0.125434\\
0.806	-0.125432\\
0.808	-0.125431\\
0.81	-0.125429\\
0.812	-0.125428\\
0.814	-0.125426\\
0.816	-0.125424\\
0.818	-0.125423\\
0.82	-0.125421\\
0.822	-0.125419\\
0.824	-0.125417\\
0.826	-0.125415\\
0.828	-0.125413\\
0.83	-0.125411\\
0.832	-0.125408\\
0.834	-0.125406\\
0.836	-0.125404\\
0.838	-0.125401\\
0.84	-0.125399\\
0.842	-0.125396\\
0.844	-0.125393\\
0.846	-0.125391\\
0.848	-0.125388\\
0.85	-0.125385\\
0.852	-0.125382\\
0.854	-0.125379\\
0.856	-0.125376\\
0.858	-0.125373\\
0.86	-0.12537\\
0.862	-0.125367\\
0.864	-0.125363\\
0.866	-0.12536\\
0.868	-0.125357\\
0.87	-0.125353\\
0.872	-0.12535\\
0.874	-0.125346\\
0.876	-0.125343\\
0.878	-0.125339\\
0.88	-0.125336\\
0.882	-0.125332\\
0.884	-0.125328\\
0.886	-0.125324\\
0.888	-0.125321\\
0.89	-0.125317\\
0.892	-0.125313\\
0.894	-0.125309\\
0.896	-0.125306\\
0.898	-0.125302\\
0.9	-0.125298\\
0.902	-0.125294\\
0.904	-0.12529\\
0.906	-0.125286\\
0.908	-0.125282\\
0.91	-0.125278\\
0.912	-0.125275\\
0.914	-0.125271\\
0.916	-0.125267\\
0.918	-0.125263\\
0.92	-0.125259\\
0.922	-0.125255\\
0.924	-0.125251\\
0.926	-0.125248\\
0.928	-0.125244\\
0.93	-0.12524\\
0.932	-0.125236\\
0.934	-0.125233\\
0.936	-0.125229\\
0.938	-0.125225\\
0.94	-0.125221\\
0.942	-0.125218\\
0.944	-0.125214\\
0.946	-0.125211\\
0.948	-0.125207\\
0.95	-0.125204\\
0.952	-0.1252\\
0.954	-0.125197\\
0.956	-0.125194\\
0.958	-0.12519\\
0.96	-0.125187\\
0.962	-0.125184\\
0.964	-0.12518\\
0.966	-0.125177\\
0.968	-0.125174\\
0.97	-0.125171\\
0.972	-0.125168\\
0.974	-0.125165\\
0.976	-0.125162\\
0.978	-0.125159\\
0.98	-0.125157\\
0.982	-0.125154\\
0.984	-0.125151\\
0.986	-0.125149\\
0.988	-0.125146\\
0.99	-0.125143\\
0.992	-0.125141\\
0.994	-0.125139\\
0.996	-0.125136\\
0.998	-0.125134\\
1	-0.125131\\
1.002	-0.125129\\
1.004	-0.125127\\
1.006	-0.125125\\
1.008	-0.125123\\
1.01	-0.125121\\
1.012	-0.125119\\
1.014	-0.125117\\
1.016	-0.125115\\
1.018	-0.125113\\
1.02	-0.125111\\
1.022	-0.125109\\
1.024	-0.125108\\
1.026	-0.125106\\
1.028	-0.125104\\
1.03	-0.125103\\
1.032	-0.125101\\
1.034	-0.1251\\
1.036	-0.125098\\
1.038	-0.125097\\
1.04	-0.125095\\
1.042	-0.125094\\
1.044	-0.125093\\
1.046	-0.125091\\
1.048	-0.12509\\
1.05	-0.125089\\
1.052	-0.125087\\
1.054	-0.125086\\
1.056	-0.125085\\
1.058	-0.125084\\
1.06	-0.125083\\
1.062	-0.125082\\
1.064	-0.12508\\
1.066	-0.125079\\
1.068	-0.125078\\
1.07	-0.125077\\
1.072	-0.125076\\
1.074	-0.125075\\
1.076	-0.125074\\
1.078	-0.125073\\
1.08	-0.125072\\
1.082	-0.125071\\
1.084	-0.12507\\
1.086	-0.125069\\
1.088	-0.125068\\
1.09	-0.125067\\
1.092	-0.125066\\
1.094	-0.125065\\
1.096	-0.125064\\
1.098	-0.125063\\
1.1	-0.125062\\
1.102	-0.125061\\
1.104	-0.12506\\
1.106	-0.125059\\
1.108	-0.125058\\
1.11	-0.125058\\
1.112	-0.125057\\
1.114	-0.125056\\
1.116	-0.125054\\
1.118	-0.125053\\
1.12	-0.125052\\
1.122	-0.125051\\
1.124	-0.12505\\
1.126	-0.125049\\
1.128	-0.125048\\
1.13	-0.125047\\
1.132	-0.125046\\
1.134	-0.125045\\
1.136	-0.125044\\
1.138	-0.125043\\
1.14	-0.125041\\
1.142	-0.12504\\
1.144	-0.125039\\
1.146	-0.125038\\
1.148	-0.125036\\
1.15	-0.125035\\
1.152	-0.125034\\
1.154	-0.125033\\
1.156	-0.125031\\
1.158	-0.12503\\
1.16	-0.125029\\
1.162	-0.125027\\
1.164	-0.125026\\
1.166	-0.125024\\
1.168	-0.125023\\
1.17	-0.125021\\
1.172	-0.12502\\
1.174	-0.125018\\
1.176	-0.125017\\
1.178	-0.125015\\
1.18	-0.125014\\
1.182	-0.125012\\
1.184	-0.12501\\
1.186	-0.125009\\
1.188	-0.125007\\
1.19	-0.125005\\
1.192	-0.125004\\
1.194	-0.125002\\
1.196	-0.125\\
1.198	-0.124999\\
1.2	-0.124997\\
1.202	-0.124995\\
1.204	-0.124993\\
1.206	-0.124991\\
1.208	-0.124989\\
1.21	-0.124987\\
1.212	-0.124986\\
1.214	-0.124984\\
1.216	-0.124982\\
1.218	-0.12498\\
1.22	-0.124978\\
1.222	-0.124976\\
1.224	-0.124974\\
1.226	-0.124972\\
1.228	-0.12497\\
1.23	-0.124968\\
1.232	-0.124966\\
1.234	-0.124963\\
1.236	-0.124961\\
1.238	-0.124959\\
1.24	-0.124957\\
1.242	-0.124955\\
1.244	-0.124953\\
1.246	-0.124951\\
1.248	-0.124948\\
1.25	-0.124946\\
1.252	-0.124944\\
1.254	-0.124942\\
1.256	-0.12494\\
1.258	-0.124937\\
1.26	-0.124935\\
1.262	-0.124933\\
1.264	-0.124931\\
1.266	-0.124928\\
1.268	-0.124926\\
1.27	-0.124924\\
1.272	-0.124922\\
1.274	-0.124919\\
1.276	-0.124917\\
1.278	-0.124915\\
1.28	-0.124912\\
1.282	-0.12491\\
1.284	-0.124908\\
1.286	-0.124905\\
1.288	-0.124903\\
1.29	-0.124901\\
1.292	-0.124898\\
1.294	-0.124896\\
1.296	-0.124894\\
1.298	-0.124891\\
1.3	-0.124889\\
1.302	-0.124887\\
1.304	-0.124884\\
1.306	-0.124882\\
1.308	-0.12488\\
1.31	-0.124877\\
1.312	-0.124875\\
1.314	-0.124873\\
1.316	-0.12487\\
1.318	-0.124868\\
1.32	-0.124866\\
1.322	-0.124863\\
1.324	-0.124861\\
1.326	-0.124859\\
1.328	-0.124856\\
1.33	-0.124854\\
1.332	-0.124852\\
1.334	-0.12485\\
1.336	-0.124847\\
1.338	-0.124845\\
1.34	-0.124843\\
1.342	-0.12484\\
1.344	-0.124838\\
1.346	-0.124836\\
1.348	-0.124834\\
1.35	-0.124831\\
1.352	-0.124829\\
1.354	-0.124827\\
1.356	-0.124825\\
1.358	-0.124822\\
1.36	-0.12482\\
1.362	-0.124818\\
1.364	-0.124816\\
1.366	-0.124814\\
1.368	-0.124811\\
1.37	-0.124809\\
1.372	-0.124807\\
1.374	-0.124805\\
1.376	-0.124803\\
1.378	-0.1248\\
1.38	-0.124798\\
1.382	-0.124796\\
1.384	-0.124794\\
1.386	-0.124792\\
1.388	-0.12479\\
1.39	-0.124788\\
1.392	-0.124785\\
1.394	-0.124783\\
1.396	-0.124781\\
1.398	-0.124779\\
1.4	-0.124777\\
1.402	-0.124775\\
1.404	-0.124773\\
1.406	-0.124771\\
1.408	-0.124769\\
1.41	-0.124767\\
1.412	-0.124764\\
1.414	-0.124762\\
1.416	-0.12476\\
1.418	-0.124758\\
1.42	-0.124756\\
1.422	-0.124754\\
1.424	-0.124752\\
1.426	-0.12475\\
1.428	-0.124748\\
1.43	-0.124746\\
1.432	-0.124744\\
1.434	-0.124742\\
1.436	-0.12474\\
1.438	-0.124738\\
1.44	-0.124736\\
1.442	-0.124734\\
1.444	-0.124732\\
1.446	-0.12473\\
1.448	-0.124728\\
1.45	-0.124726\\
1.452	-0.124724\\
1.454	-0.124722\\
1.456	-0.12472\\
1.458	-0.124718\\
1.46	-0.124716\\
1.462	-0.124714\\
1.464	-0.124712\\
1.466	-0.12471\\
1.468	-0.124708\\
1.47	-0.124706\\
1.472	-0.124704\\
1.474	-0.124702\\
1.476	-0.1247\\
1.478	-0.124698\\
1.48	-0.124696\\
1.482	-0.124694\\
1.484	-0.124692\\
1.486	-0.12469\\
1.488	-0.124688\\
1.49	-0.124686\\
1.492	-0.124684\\
1.494	-0.124682\\
1.496	-0.12468\\
1.498	-0.124678\\
1.5	-0.124676\\
1.502	-0.124674\\
1.504	-0.124672\\
1.506	-0.12467\\
1.508	-0.124668\\
1.51	-0.124666\\
1.512	-0.124664\\
1.514	-0.124662\\
1.516	-0.12466\\
1.518	-0.124658\\
1.52	-0.124656\\
1.522	-0.124654\\
1.524	-0.124652\\
1.526	-0.12465\\
1.528	-0.124648\\
1.53	-0.124646\\
1.532	-0.124644\\
1.534	-0.124642\\
1.536	-0.12464\\
1.538	-0.124638\\
1.54	-0.124636\\
1.542	-0.124634\\
1.544	-0.124632\\
1.546	-0.12463\\
1.548	-0.124628\\
1.55	-0.124626\\
1.552	-0.124624\\
1.554	-0.124622\\
1.556	-0.12462\\
1.558	-0.124618\\
1.56	-0.124616\\
1.562	-0.124614\\
1.564	-0.124612\\
1.566	-0.12461\\
1.568	-0.124608\\
1.57	-0.124606\\
1.572	-0.124604\\
1.574	-0.124601\\
1.576	-0.124599\\
1.578	-0.124597\\
1.58	-0.124595\\
1.582	-0.124593\\
1.584	-0.124591\\
1.586	-0.124589\\
1.588	-0.124587\\
1.59	-0.124585\\
1.592	-0.124583\\
1.594	-0.124581\\
1.596	-0.124579\\
1.598	-0.124577\\
1.6	-0.124575\\
1.602	-0.124573\\
1.604	-0.12457\\
1.606	-0.124568\\
1.608	-0.124566\\
1.61	-0.124564\\
1.612	-0.124562\\
1.614	-0.12456\\
1.616	-0.124558\\
1.618	-0.124556\\
1.62	-0.124554\\
1.622	-0.124552\\
1.624	-0.12455\\
1.626	-0.124548\\
1.628	-0.124546\\
1.63	-0.124544\\
1.632	-0.124542\\
1.634	-0.124539\\
1.636	-0.124537\\
1.638	-0.124535\\
1.64	-0.124533\\
1.642	-0.124531\\
1.644	-0.124529\\
1.646	-0.124527\\
1.648	-0.124525\\
1.65	-0.124523\\
1.652	-0.124521\\
1.654	-0.124519\\
1.656	-0.124517\\
1.658	-0.124515\\
1.66	-0.124513\\
1.662	-0.124511\\
1.664	-0.124509\\
1.666	-0.124507\\
1.668	-0.124505\\
1.67	-0.124503\\
1.672	-0.124501\\
1.674	-0.124499\\
1.676	-0.124496\\
1.678	-0.124494\\
1.68	-0.124492\\
1.682	-0.12449\\
1.684	-0.124488\\
1.686	-0.124486\\
1.688	-0.124484\\
1.69	-0.124482\\
1.692	-0.12448\\
1.694	-0.124478\\
1.696	-0.124476\\
1.698	-0.124474\\
1.7	-0.124472\\
1.702	-0.12447\\
1.704	-0.124468\\
1.706	-0.124466\\
1.708	-0.124464\\
1.71	-0.124462\\
1.712	-0.12446\\
1.714	-0.124458\\
1.716	-0.124456\\
1.718	-0.124454\\
1.72	-0.124452\\
1.722	-0.12445\\
1.724	-0.124448\\
1.726	-0.124446\\
1.728	-0.124445\\
1.73	-0.124443\\
1.732	-0.124441\\
1.734	-0.124439\\
1.736	-0.124437\\
1.738	-0.124435\\
1.74	-0.124433\\
1.742	-0.124431\\
1.744	-0.124429\\
1.746	-0.124427\\
1.748	-0.124425\\
1.75	-0.124423\\
1.752	-0.124421\\
1.754	-0.124419\\
1.756	-0.124417\\
1.758	-0.124415\\
1.76	-0.124413\\
1.762	-0.124411\\
1.764	-0.124409\\
1.766	-0.124408\\
1.768	-0.124406\\
1.77	-0.124404\\
1.772	-0.124402\\
1.774	-0.1244\\
1.776	-0.124398\\
1.778	-0.124396\\
1.78	-0.124394\\
1.782	-0.124392\\
1.784	-0.12439\\
1.786	-0.124388\\
1.788	-0.124386\\
1.79	-0.124384\\
1.792	-0.124383\\
1.794	-0.124381\\
1.796	-0.124379\\
1.798	-0.124377\\
1.8	-0.124375\\
1.802	-0.124373\\
1.804	-0.124371\\
1.806	-0.124369\\
1.808	-0.124367\\
1.81	-0.124365\\
1.812	-0.124364\\
1.814	-0.124362\\
1.816	-0.12436\\
1.818	-0.124358\\
1.82	-0.124356\\
1.822	-0.124354\\
1.824	-0.124352\\
1.826	-0.12435\\
1.828	-0.124348\\
1.83	-0.124347\\
1.832	-0.124345\\
1.834	-0.124343\\
1.836	-0.124341\\
1.838	-0.124339\\
1.84	-0.124337\\
1.842	-0.124335\\
1.844	-0.124333\\
1.846	-0.124331\\
1.848	-0.12433\\
1.85	-0.124328\\
1.852	-0.124326\\
1.854	-0.124324\\
1.856	-0.124322\\
1.858	-0.12432\\
1.86	-0.124318\\
1.862	-0.124316\\
1.864	-0.124315\\
1.866	-0.124313\\
1.868	-0.124311\\
1.87	-0.124309\\
1.872	-0.124307\\
1.874	-0.124305\\
1.876	-0.124303\\
1.878	-0.124301\\
1.88	-0.1243\\
1.882	-0.124298\\
1.884	-0.124296\\
1.886	-0.124294\\
1.888	-0.124292\\
1.89	-0.12429\\
1.892	-0.124288\\
1.894	-0.124286\\
1.896	-0.124285\\
1.898	-0.124283\\
1.9	-0.124281\\
1.902	-0.124279\\
1.904	-0.124277\\
1.906	-0.124275\\
1.908	-0.124273\\
1.91	-0.124271\\
1.912	-0.12427\\
1.914	-0.124268\\
1.916	-0.124266\\
1.918	-0.124264\\
1.92	-0.124262\\
1.922	-0.12426\\
1.924	-0.124258\\
1.926	-0.124256\\
1.928	-0.124255\\
1.93	-0.124253\\
1.932	-0.124251\\
1.934	-0.124249\\
1.936	-0.124247\\
1.938	-0.124245\\
1.94	-0.124243\\
1.942	-0.124242\\
1.944	-0.12424\\
1.946	-0.124238\\
1.948	-0.124236\\
1.95	-0.124234\\
1.952	-0.124232\\
1.954	-0.12423\\
1.956	-0.124228\\
1.958	-0.124227\\
1.96	-0.124225\\
1.962	-0.124223\\
1.964	-0.124221\\
1.966	-0.124219\\
1.968	-0.124217\\
1.97	-0.124215\\
1.972	-0.124214\\
1.974	-0.124212\\
1.976	-0.12421\\
1.978	-0.124208\\
1.98	-0.124206\\
1.982	-0.124204\\
1.984	-0.124202\\
1.986	-0.124201\\
1.988	-0.124199\\
1.99	-0.124197\\
1.992	-0.124195\\
1.994	-0.124193\\
1.996	-0.124191\\
1.998	-0.124189\\
2	-0.124188\\
};
\addplot [color=black, dotted, line width=0.6pt, forget plot]
  table[row sep=crcr]{%
-0.125	-0.125748\\
2	-0.125748\\
};
\draw[line width=\figlegendlw, fill=\figlegendfill, draw=\figlegenddraw] (axis cs:1.13004,-0.0234737) rectangle (axis cs:1.9575,0.094);
\addplot [color=black, line width=1.4pt, forget plot]
  table[row sep=crcr]{%
1.17141	0.0705053\\
1.4362	0.0705053\\
};
\node[right, align=left, inner sep=0]
at (axis cs:1.461,0.071) {$f_{\mathrm{COI}}$};
\addplot [color=mycolor15, line width=1.6pt, forget plot]
  table[row sep=crcr]{%
1.17141	0.0352632\\
1.20083	0.0352632\\
};
\addplot [color=mycolor16, line width=1.6pt, forget plot]
  table[row sep=crcr]{%
1.20083	0.0352632\\
1.23025	0.0352632\\
};
\addplot [color=mycolor20, line width=1.6pt, forget plot]
  table[row sep=crcr]{%
1.23025	0.0352632\\
1.25967	0.0352632\\
};
\addplot [color=mycolor17, line width=1.6pt, forget plot]
  table[row sep=crcr]{%
1.25967	0.0352632\\
1.28909	0.0352632\\
};
\addplot [color=mycolor13, line width=1.6pt, forget plot]
  table[row sep=crcr]{%
1.28909	0.0352632\\
1.31852	0.0352632\\
};
\addplot [color=mycolor12, line width=1.6pt, forget plot]
  table[row sep=crcr]{%
1.31852	0.0352632\\
1.34794	0.0352632\\
};
\addplot [color=mycolor18, line width=1.6pt, forget plot]
  table[row sep=crcr]{%
1.34794	0.0352632\\
1.37736	0.0352632\\
};
\addplot [color=mycolor19, line width=1.6pt, forget plot]
  table[row sep=crcr]{%
1.37736	0.0352632\\
1.40678	0.0352632\\
};
\addplot [color=mycolor14, line width=1.6pt, forget plot]
  table[row sep=crcr]{%
1.40678	0.0352632\\
1.4362	0.0352632\\
};
\node[right, align=left, inner sep=0]
at (axis cs:1.461,0.035) {$f_i$};
\addplot [color=black, dotted, line width=0.6pt, forget plot]
  table[row sep=crcr]{%
1.17141	2.10526e-05\\
1.4362	2.10526e-05\\
};
\node[right, align=left, inner sep=0]
at (axis cs:1.461,0) {$N_{\mathrm{COI}}^{20}$};
\addplot [color=black, dotted, line width=0.6pt, forget plot]
  table[row sep=crcr]{%
-0.125	-0.125832\\
2	-0.125832\\
};
\end{axis}
\end{tikzpicture}%
  \caption{Fig.~5(a) --- per-node deviation, linearised model, $K=0$ $|$ $K=20\,I$.}
\end{figure}

\begin{figure}[!t]\centering
  % This file was created by matlab2tikz.
%
\definecolor{mycolor1}{rgb}{0.87843,0.67059,0.67059}%
\definecolor{mycolor2}{rgb}{0.63656,0.15679,0.15679}%
\definecolor{mycolor3}{rgb}{0.60075,0.09142,0.09142}%
\definecolor{mycolor4}{rgb}{0.81501,0.55362,0.55362}%
\definecolor{mycolor5}{rgb}{0.48525,0.00988,0.00988}%
\definecolor{mycolor6}{rgb}{0.53748,0.01771,0.01771}%
\definecolor{mycolor7}{rgb}{0.56781,0.04223,0.04223}%
\definecolor{mycolor8}{rgb}{0.74576,0.39131,0.39131}%
\definecolor{mycolor9}{rgb}{0.77804,0.46620,0.46620}%
\definecolor{mycolor10}{rgb}{0.51292,0.01172,0.01172}%
\definecolor{mycolor11}{rgb}{0.66892,0.22212,0.22212}%
\definecolor{mycolor12}{rgb}{0.70745,0.30499,0.30499}%
\definecolor{mycolor13}{rgb}{0.93333,0.98824,0.92157}%
\definecolor{mycolor14}{rgb}{0.25274,0.63656,0.15679}%
\definecolor{mycolor15}{rgb}{0.19329,0.60075,0.09142}%
\definecolor{mycolor16}{rgb}{0.60590,0.81501,0.55362}%
\definecolor{mycolor17}{rgb}{0.10495,0.48525,0.00988}%
\definecolor{mycolor18}{rgb}{0.12167,0.53748,0.01771}%
\definecolor{mycolor19}{rgb}{0.14735,0.56781,0.04223}%
\definecolor{mycolor20}{rgb}{0.46220,0.74576,0.39131}%
\definecolor{mycolor21}{rgb}{0.52856,0.77804,0.46620}%
\definecolor{mycolor22}{rgb}{0.11196,0.51292,0.01172}%
\definecolor{mycolor23}{rgb}{0.31148,0.66892,0.22212}%
\definecolor{mycolor24}{rgb}{0.38548,0.70745,0.30499}%
%
\begin{tikzpicture}[font=\figfont]

\begin{axis}[%
width={\figscale*1.32in},
height={\figscale*0.95in},
at={(\figscale*0.455in,\figscale*0.31in)},
scale only axis,
separate axis lines,
every outer x axis line/.append style={black},
every x tick label/.append style={font=\figticfont},
every x tick/.append style={black},
xmin=-0.125,
xmax=5,
xlabel={Time since step in s},
every outer y axis line/.append style={black},
every y tick label/.append style={font=\figticfont},
every y tick/.append style={black},
ymin=-0.2,
ymax=0.1,
ylabel={$f - f_0$ in Hz},
axis background/.style={fill=white},
xmajorgrids,
ymajorgrids,
grid style={white!55!black, opacity=0.3}
]

\addplot[area legend, draw=none, fill=mycolor1, fill opacity=0.9, forget plot]
table[row sep=crcr] {%
x	y\\
0	-0.0120484\\
0.0125313	-0.0360376\\
0.0250627	-0.0536407\\
0.037594	-0.0674355\\
0.0501253	-0.067624\\
0.0626566	-0.068093\\
0.075188	-0.0758625\\
0.0877193	-0.0838977\\
0.100251	-0.0915609\\
0.112782	-0.0982411\\
0.125313	-0.103414\\
0.137845	-0.106592\\
0.150376	-0.108507\\
0.162907	-0.11023\\
0.175439	-0.111803\\
0.18797	-0.11328\\
0.200501	-0.114678\\
0.213033	-0.116002\\
0.225564	-0.117281\\
0.238095	-0.118536\\
0.250627	-0.119755\\
0.263158	-0.12091\\
0.275689	-0.121991\\
0.288221	-0.122989\\
0.300752	-0.123872\\
0.313283	-0.124605\\
0.325815	-0.12517\\
0.338346	-0.125559\\
0.350877	-0.125749\\
0.363409	-0.125719\\
0.37594	-0.12468\\
0.388471	-0.122037\\
0.401003	-0.118059\\
0.413534	-0.113011\\
0.426065	-0.107171\\
0.438596	-0.100817\\
0.451128	-0.0942207\\
0.463659	-0.0876489\\
0.47619	-0.0813747\\
0.488722	-0.0756743\\
0.501253	-0.0708168\\
0.513784	-0.0670674\\
0.526316	-0.0646972\\
0.538847	-0.06398\\
0.551378	-0.0642311\\
0.56391	-0.0645739\\
0.576441	-0.0649953\\
0.588972	-0.0654853\\
0.601504	-0.0660302\\
0.614035	-0.0666145\\
0.626566	-0.0672264\\
0.639098	-0.0678564\\
0.651629	-0.0684921\\
0.66416	-0.0691198\\
0.676692	-0.0697278\\
0.689223	-0.0703066\\
0.701754	-0.0708445\\
0.714286	-0.071328\\
0.726817	-0.0715068\\
0.739348	-0.0711871\\
0.75188	-0.0704394\\
0.764411	-0.0693329\\
0.776942	-0.0679383\\
0.789474	-0.0663275\\
0.802005	-0.0645713\\
0.814536	-0.0627396\\
0.827068	-0.0609037\\
0.839599	-0.0591358\\
0.85213	-0.0575077\\
0.864662	-0.0560906\\
0.877193	-0.0549569\\
0.889724	-0.0541799\\
0.902256	-0.0538327\\
0.914787	-0.053988\\
0.927318	-0.0547195\\
0.93985	-0.0562797\\
0.952381	-0.0588035\\
0.964912	-0.0621889\\
0.977444	-0.0663343\\
0.989975	-0.0711386\\
1.00251	-0.0765\\
1.01504	-0.0823165\\
1.02757	-0.0884857\\
1.0401	-0.0949055\\
1.05263	-0.101473\\
1.06516	-0.108084\\
1.07769	-0.114635\\
1.09023	-0.12102\\
1.10276	-0.127135\\
1.11529	-0.132871\\
1.12782	-0.138123\\
1.14035	-0.14278\\
1.15288	-0.146734\\
1.16541	-0.149875\\
1.17794	-0.152584\\
1.19048	-0.155316\\
1.20301	-0.158063\\
1.21554	-0.16082\\
1.22807	-0.163576\\
1.2406	-0.166322\\
1.25313	-0.169049\\
1.26566	-0.171745\\
1.2782	-0.174397\\
1.29073	-0.176994\\
1.30326	-0.179523\\
1.31579	-0.18197\\
1.32832	-0.184322\\
1.34085	-0.186565\\
1.35338	-0.188687\\
1.36591	-0.190673\\
1.37845	-0.192513\\
1.39098	-0.194194\\
1.40351	-0.195707\\
1.41604	-0.197041\\
1.42857	-0.19819\\
1.4411	-0.199146\\
1.45363	-0.199905\\
1.46617	-0.199958\\
1.4787	-0.19896\\
1.49123	-0.197153\\
1.50376	-0.194779\\
1.51629	-0.192084\\
1.52882	-0.189314\\
1.54135	-0.186717\\
1.55388	-0.184542\\
1.56642	-0.183039\\
1.57895	-0.182458\\
1.59148	-0.182462\\
1.60401	-0.182541\\
1.61654	-0.182688\\
1.62907	-0.182899\\
1.6416	-0.18317\\
1.65414	-0.183496\\
1.66667	-0.183872\\
1.6792	-0.184294\\
1.69173	-0.184758\\
1.70426	-0.185259\\
1.71679	-0.185791\\
1.72932	-0.18635\\
1.74185	-0.18693\\
1.75439	-0.187524\\
1.76692	-0.188126\\
1.77945	-0.188729\\
1.79198	-0.189324\\
1.80451	-0.189904\\
1.81704	-0.190459\\
1.82957	-0.190979\\
1.84211	-0.191455\\
1.85464	-0.191876\\
1.86717	-0.192228\\
1.8797	-0.192501\\
1.89223	-0.192681\\
1.90476	-0.192754\\
1.91729	-0.192705\\
1.92982	-0.192453\\
1.94236	-0.19194\\
1.95489	-0.191189\\
1.96742	-0.190224\\
1.97995	-0.189066\\
1.99248	-0.187736\\
2.00501	-0.186255\\
2.01754	-0.184644\\
2.03008	-0.182921\\
2.04261	-0.181105\\
2.05514	-0.179215\\
2.06767	-0.177268\\
2.0802	-0.175282\\
2.09273	-0.173272\\
2.10526	-0.171255\\
2.11779	-0.169246\\
2.13033	-0.167261\\
2.14286	-0.165315\\
2.15539	-0.163422\\
2.16792	-0.161597\\
2.18045	-0.159855\\
2.19298	-0.15821\\
2.20551	-0.156677\\
2.21805	-0.15527\\
2.23058	-0.154006\\
2.24311	-0.152901\\
2.25564	-0.151971\\
2.26817	-0.151233\\
2.2807	-0.150707\\
2.29323	-0.150412\\
2.30576	-0.150224\\
2.3183	-0.149993\\
2.33083	-0.149697\\
2.34336	-0.149316\\
2.35589	-0.148831\\
2.36842	-0.148223\\
2.38095	-0.147477\\
2.39348	-0.146609\\
2.40602	-0.145649\\
2.41855	-0.144604\\
2.43108	-0.14348\\
2.44361	-0.142285\\
2.45614	-0.141028\\
2.46867	-0.139718\\
2.4812	-0.138365\\
2.49373	-0.136977\\
2.50627	-0.135564\\
2.5188	-0.134137\\
2.53133	-0.132706\\
2.54386	-0.131279\\
2.55639	-0.129866\\
2.56892	-0.128477\\
2.58145	-0.127119\\
2.59398	-0.125801\\
2.60652	-0.124529\\
2.61905	-0.123312\\
2.63158	-0.122154\\
2.64411	-0.121062\\
2.65664	-0.120039\\
2.66917	-0.119091\\
2.6817	-0.118219\\
2.69424	-0.117428\\
2.70677	-0.116717\\
2.7193	-0.116089\\
2.73183	-0.115544\\
2.74436	-0.115082\\
2.75689	-0.114615\\
2.76942	-0.114062\\
2.78195	-0.113436\\
2.79449	-0.112745\\
2.80702	-0.111999\\
2.81955	-0.111208\\
2.83208	-0.110378\\
2.84461	-0.10952\\
2.85714	-0.10864\\
2.86967	-0.107745\\
2.88221	-0.106844\\
2.89474	-0.105942\\
2.90727	-0.105046\\
2.9198	-0.104162\\
2.93233	-0.103297\\
2.94486	-0.102457\\
2.95739	-0.101648\\
2.96992	-0.100875\\
2.98246	-0.100144\\
2.99499	-0.0994615\\
3.00752	-0.0988334\\
3.02005	-0.0982659\\
3.03258	-0.0977652\\
3.04511	-0.0973379\\
3.05764	-0.096991\\
3.07018	-0.0967316\\
3.08271	-0.0965672\\
3.09524	-0.096506\\
3.10777	-0.0965563\\
3.1203	-0.0967271\\
3.13283	-0.0969448\\
3.14536	-0.0971288\\
3.15789	-0.0972783\\
3.17043	-0.0973928\\
3.18296	-0.0974726\\
3.19549	-0.0975185\\
3.20802	-0.0975317\\
3.22055	-0.0975138\\
3.23308	-0.0974672\\
3.24561	-0.0973943\\
3.25815	-0.0972982\\
3.27068	-0.0971821\\
3.28321	-0.0970495\\
3.29574	-0.0969043\\
3.30827	-0.0967504\\
3.3208	-0.0965917\\
3.33333	-0.0964323\\
3.34586	-0.0962762\\
3.3584	-0.0961275\\
3.37093	-0.09599\\
3.38346	-0.0958674\\
3.39599	-0.0957632\\
3.40852	-0.0956806\\
3.42105	-0.0956226\\
3.43358	-0.0955919\\
3.44612	-0.095591\\
3.45865	-0.0956217\\
3.47118	-0.0956858\\
3.48371	-0.0957846\\
3.49624	-0.0959391\\
3.50877	-0.096167\\
3.5213	-0.096464\\
3.53383	-0.0968259\\
3.54637	-0.0972478\\
3.5589	-0.0977247\\
3.57143	-0.0982513\\
3.58396	-0.098822\\
3.59649	-0.0994311\\
3.60902	-0.100073\\
3.62155	-0.10074\\
3.63409	-0.101428\\
3.64662	-0.10213\\
3.65915	-0.102838\\
3.67168	-0.103547\\
3.68421	-0.10425\\
3.69674	-0.104941\\
3.70927	-0.105613\\
3.7218	-0.106259\\
3.73434	-0.106873\\
3.74687	-0.107449\\
3.7594	-0.107981\\
3.77193	-0.108462\\
3.78446	-0.108888\\
3.79699	-0.109252\\
3.80952	-0.109549\\
3.82206	-0.109774\\
3.83459	-0.109922\\
3.84712	-0.109989\\
3.85965	-0.10997\\
3.87218	-0.109862\\
3.88471	-0.109661\\
3.89724	-0.109372\\
3.90977	-0.109006\\
3.92231	-0.108566\\
3.93484	-0.108056\\
3.94737	-0.107479\\
3.9599	-0.10684\\
3.97243	-0.106145\\
3.98496	-0.105397\\
3.99749	-0.104602\\
4.01003	-0.103765\\
4.02256	-0.102892\\
4.03509	-0.101988\\
4.04762	-0.101058\\
4.06015	-0.100109\\
4.07268	-0.0991468\\
4.08521	-0.0981764\\
4.09774	-0.0972039\\
4.11028	-0.096235\\
4.12281	-0.0952752\\
4.13534	-0.0943299\\
4.14787	-0.0934045\\
4.1604	-0.092504\\
4.17293	-0.0916333\\
4.18546	-0.090797\\
4.19799	-0.0899995\\
4.21053	-0.0892448\\
4.22306	-0.0885368\\
4.23559	-0.0878789\\
4.24812	-0.0872742\\
4.26065	-0.0867256\\
4.27318	-0.0862356\\
4.28571	-0.0858064\\
4.29825	-0.0854398\\
4.31078	-0.0851374\\
4.32331	-0.0848897\\
4.33584	-0.0846849\\
4.34837	-0.0845204\\
4.3609	-0.0843934\\
4.37343	-0.0843011\\
4.38596	-0.0842402\\
4.3985	-0.0842074\\
4.41103	-0.0841992\\
4.42356	-0.0842119\\
4.43609	-0.0842417\\
4.44862	-0.0842848\\
4.46115	-0.0843373\\
4.47368	-0.0843954\\
4.48622	-0.0844552\\
4.49875	-0.0845129\\
4.51128	-0.0845648\\
4.52381	-0.0846072\\
4.53634	-0.0846366\\
4.54887	-0.0846498\\
4.5614	-0.0846436\\
4.57393	-0.0846151\\
4.58647	-0.0845616\\
4.599	-0.0844806\\
4.61153	-0.0843702\\
4.62406	-0.0842282\\
4.63659	-0.0840562\\
4.64912	-0.0838561\\
4.66165	-0.0836281\\
4.67419	-0.0833723\\
4.68672	-0.0830892\\
4.69925	-0.0827797\\
4.71178	-0.0824447\\
4.72431	-0.0820857\\
4.73684	-0.0817041\\
4.74937	-0.0813017\\
4.7619	-0.0808806\\
4.77444	-0.0804428\\
4.78697	-0.0799907\\
4.7995	-0.0795266\\
4.81203	-0.0790531\\
4.82456	-0.0785727\\
4.83709	-0.0780881\\
4.84962	-0.077602\\
4.86216	-0.0771168\\
4.87469	-0.0766353\\
4.88722	-0.0761598\\
4.89975	-0.0756929\\
4.91228	-0.0752369\\
4.92481	-0.0747938\\
4.93734	-0.0743657\\
4.94987	-0.0739544\\
4.96241	-0.0735615\\
4.97494	-0.0731885\\
4.98747	-0.0728366\\
5	-0.0725067\\
5	-0.0333022\\
4.98747	-0.0329036\\
4.97494	-0.032534\\
4.96241	-0.0321926\\
4.94987	-0.0318783\\
4.93734	-0.0315898\\
4.92481	-0.0313257\\
4.91228	-0.0310841\\
4.89975	-0.0308632\\
4.88722	-0.0306609\\
4.87469	-0.0304752\\
4.86216	-0.0303036\\
4.84962	-0.0301438\\
4.83709	-0.0299934\\
4.82456	-0.0298498\\
4.81203	-0.0297105\\
4.7995	-0.0295731\\
4.78697	-0.029435\\
4.77444	-0.0292938\\
4.7619	-0.0291473\\
4.74937	-0.0289933\\
4.73684	-0.0288298\\
4.72431	-0.0286549\\
4.71178	-0.0284669\\
4.69925	-0.0282644\\
4.68672	-0.0280463\\
4.67419	-0.0278115\\
4.66165	-0.0275595\\
4.64912	-0.0272897\\
4.63659	-0.0270022\\
4.62406	-0.026697\\
4.61153	-0.0263772\\
4.599	-0.0260456\\
4.58647	-0.0257023\\
4.57393	-0.0253477\\
4.5614	-0.0249824\\
4.54887	-0.0246073\\
4.53634	-0.0242236\\
4.52381	-0.0238327\\
4.51128	-0.0234361\\
4.49875	-0.0230354\\
4.48622	-0.0226327\\
4.47368	-0.0222297\\
4.46115	-0.0218287\\
4.44862	-0.0214316\\
4.43609	-0.0210406\\
4.42356	-0.0206577\\
4.41103	-0.0202851\\
4.3985	-0.0199248\\
4.38596	-0.0195786\\
4.37343	-0.0192482\\
4.3609	-0.0189354\\
4.34837	-0.0186414\\
4.33584	-0.0183675\\
4.32331	-0.0181147\\
4.31078	-0.0178836\\
4.29825	-0.0176651\\
4.28571	-0.0174518\\
4.27318	-0.0172466\\
4.26065	-0.0170523\\
4.24812	-0.0168713\\
4.23559	-0.0167057\\
4.22306	-0.0165572\\
4.21053	-0.0164273\\
4.19799	-0.0163172\\
4.18546	-0.0162275\\
4.17293	-0.0161589\\
4.1604	-0.0161116\\
4.14787	-0.0160855\\
4.13534	-0.0160803\\
4.12281	-0.0160956\\
4.11028	-0.0161308\\
4.09774	-0.0161848\\
4.08521	-0.0162568\\
4.07268	-0.0163457\\
4.06015	-0.0164503\\
4.04762	-0.0165694\\
4.03509	-0.0167019\\
4.02256	-0.0168466\\
4.01003	-0.0170023\\
3.99749	-0.0171681\\
3.98496	-0.0173432\\
3.97243	-0.0175269\\
3.9599	-0.0177186\\
3.94737	-0.0179181\\
3.93484	-0.0181253\\
3.92231	-0.0183405\\
3.90977	-0.0185641\\
3.89724	-0.0187969\\
3.88471	-0.0190401\\
3.87218	-0.0193017\\
3.85965	-0.0195862\\
3.84712	-0.0198895\\
3.83459	-0.0202077\\
3.82206	-0.0205373\\
3.80952	-0.0208752\\
3.79699	-0.0212183\\
3.78446	-0.0215641\\
3.77193	-0.0219101\\
3.7594	-0.0222542\\
3.74687	-0.0225944\\
3.73434	-0.0229289\\
3.7218	-0.0232563\\
3.70927	-0.0235748\\
3.69674	-0.0238834\\
3.68421	-0.0241806\\
3.67168	-0.0244653\\
3.65915	-0.0247364\\
3.64662	-0.0249925\\
3.63409	-0.0252326\\
3.62155	-0.0254554\\
3.60902	-0.0256596\\
3.59649	-0.0258436\\
3.58396	-0.026006\\
3.57143	-0.0261449\\
3.5589	-0.0262585\\
3.54637	-0.0263444\\
3.53383	-0.0264004\\
3.5213	-0.0264236\\
3.50877	-0.0264112\\
3.49624	-0.0263598\\
3.48371	-0.0262659\\
3.47118	-0.0261468\\
3.45865	-0.0260227\\
3.44612	-0.0258934\\
3.43358	-0.0257581\\
3.42105	-0.0256161\\
3.40852	-0.025466\\
3.39599	-0.0253062\\
3.38346	-0.0251347\\
3.37093	-0.0249496\\
3.3584	-0.0247482\\
3.34586	-0.0245281\\
3.33333	-0.0242864\\
3.3208	-0.0240202\\
3.30827	-0.0237266\\
3.29574	-0.0234026\\
3.28321	-0.0230453\\
3.27068	-0.0226519\\
3.25815	-0.0222196\\
3.24561	-0.021746\\
3.23308	-0.021229\\
3.22055	-0.0206666\\
3.20802	-0.0200573\\
3.19549	-0.0194\\
3.18296	-0.0186942\\
3.17043	-0.0179396\\
3.15789	-0.0171365\\
3.14536	-0.0162859\\
3.13283	-0.0153892\\
3.1203	-0.0144483\\
3.10777	-0.0133864\\
3.09524	-0.0121351\\
3.08271	-0.0107093\\
3.07018	-0.00912426\\
3.05764	-0.00739599\\
3.04511	-0.00554074\\
3.03258	-0.00357518\\
3.02005	-0.00151632\\
3.00752	0.00061855\\
2.99499	0.00281191\\
2.98246	0.00504604\\
2.96992	0.00730309\\
2.95739	0.00956511\\
2.94486	0.0118141\\
2.93233	0.0140321\\
2.9198	0.0162012\\
2.90727	0.0183035\\
2.89474	0.0203213\\
2.88221	0.0222371\\
2.86967	0.0240335\\
2.85714	0.0256935\\
2.84461	0.0272003\\
2.83208	0.0285373\\
2.81955	0.0296883\\
2.80702	0.0306377\\
2.79449	0.0313701\\
2.78195	0.0318704\\
2.76942	0.0321244\\
2.75689	0.0321179\\
2.74436	0.0318377\\
2.73183	0.0313618\\
2.7193	0.0307771\\
2.70677	0.0300847\\
2.69424	0.0292861\\
2.6817	0.0283838\\
2.66917	0.0273808\\
2.65664	0.0262808\\
2.64411	0.0250885\\
2.63158	0.0238089\\
2.61905	0.0224481\\
2.60652	0.0210127\\
2.59398	0.0195102\\
2.58145	0.0179486\\
2.56892	0.0163365\\
2.55639	0.0146833\\
2.54386	0.0129986\\
2.53133	0.0112927\\
2.5188	0.00957618\\
2.50627	0.00785994\\
2.49373	0.00615499\\
2.4812	0.00447248\\
2.46867	0.0028235\\
2.45614	0.00121897\\
2.44361	-0.000330447\\
2.43108	-0.00181449\\
2.41855	-0.00322342\\
2.40602	-0.00454815\\
2.39348	-0.00578037\\
2.38095	-0.00691261\\
2.36842	-0.00796374\\
2.35589	-0.00894059\\
2.34336	-0.00981921\\
2.33083	-0.0105768\\
2.3183	-0.011192\\
2.30576	-0.0116442\\
2.29323	-0.0119146\\
2.2807	-0.0118565\\
2.26817	-0.0113543\\
2.25564	-0.0104367\\
2.24311	-0.0091335\\
2.23058	-0.00747532\\
2.21805	-0.00549363\\
2.20551	-0.00322057\\
2.19298	-0.000688832\\
2.18045	0.00206845\\
2.16792	0.00501785\\
2.15539	0.00812577\\
2.14286	0.0113585\\
2.13033	0.0146825\\
2.11779	0.0180642\\
2.10526	0.0214704\\
2.09273	0.0248681\\
2.0802	0.0282247\\
2.06767	0.0315082\\
2.05514	0.0346869\\
2.04261	0.0377296\\
2.03008	0.0406058\\
2.01754	0.0432854\\
2.00501	0.045739\\
1.99248	0.0479378\\
1.97995	0.0498537\\
1.96742	0.0514591\\
1.95489	0.052727\\
1.94236	0.0536313\\
1.92982	0.0541465\\
1.91729	0.0542475\\
1.90476	0.0539909\\
1.89223	0.0534605\\
1.8797	0.0526735\\
1.86717	0.0516474\\
1.85464	0.050401\\
1.84211	0.0489533\\
1.82957	0.0473242\\
1.81704	0.0455344\\
1.80451	0.043605\\
1.79198	0.041558\\
1.77945	0.039416\\
1.76692	0.0372022\\
1.75439	0.0349404\\
1.74185	0.032655\\
1.72932	0.0303709\\
1.71679	0.0281135\\
1.70426	0.0259087\\
1.69173	0.0237827\\
1.6792	0.0217621\\
1.66667	0.0198739\\
1.65414	0.018145\\
1.6416	0.0166027\\
1.62907	0.0152742\\
1.61654	0.0141867\\
1.60401	0.0133671\\
1.59148	0.0128422\\
1.57895	0.0126381\\
1.56642	0.0132851\\
1.55388	0.01514\\
1.54135	0.0179687\\
1.52882	0.0215357\\
1.51629	0.0256048\\
1.50376	0.0299382\\
1.49123	0.0342969\\
1.4787	0.0384404\\
1.46617	0.0421265\\
1.45363	0.0451117\\
1.4411	0.047736\\
1.42857	0.0504991\\
1.41604	0.0533939\\
1.40351	0.0564116\\
1.39098	0.0595421\\
1.37845	0.0627738\\
1.36591	0.066094\\
1.35338	0.0694892\\
1.34085	0.0729447\\
1.32832	0.0764455\\
1.31579	0.079976\\
1.30326	0.0835205\\
1.29073	0.0870632\\
1.2782	0.0905886\\
1.26566	0.0940814\\
1.25313	0.0975269\\
1.2406	0.100911\\
1.22807	0.10422\\
1.21554	0.107441\\
1.20301	0.110563\\
1.19048	0.113575\\
1.17794	0.116467\\
1.16541	0.11923\\
1.15288	0.121398\\
1.14035	0.122577\\
1.12782	0.122867\\
1.11529	0.122373\\
1.10276	0.121197\\
1.09023	0.119444\\
1.07769	0.117219\\
1.06516	0.114629\\
1.05263	0.111782\\
1.0401	0.108785\\
1.02757	0.105746\\
1.01504	0.102775\\
1.00251	0.0999815\\
0.989975	0.0974756\\
0.977444	0.0953671\\
0.964912	0.0937667\\
0.952381	0.0927852\\
0.93985	0.0925334\\
0.927318	0.0931212\\
0.914787	0.0944224\\
0.902256	0.0961939\\
0.889724	0.0983697\\
0.877193	0.100883\\
0.864662	0.103667\\
0.85213	0.106657\\
0.839599	0.109784\\
0.827068	0.112981\\
0.814536	0.116179\\
0.802005	0.119312\\
0.789474	0.122311\\
0.776942	0.125107\\
0.764411	0.127631\\
0.75188	0.129815\\
0.739348	0.13159\\
0.726817	0.132886\\
0.714286	0.133635\\
0.701754	0.134034\\
0.689223	0.134337\\
0.676692	0.134555\\
0.66416	0.1347\\
0.651629	0.13479\\
0.639098	0.134838\\
0.626566	0.134855\\
0.614035	0.134853\\
0.601504	0.134849\\
0.588972	0.134857\\
0.576441	0.134884\\
0.56391	0.134945\\
0.551378	0.135052\\
0.538847	0.135219\\
0.526316	0.136312\\
0.513784	0.139011\\
0.501253	0.143046\\
0.488722	0.148145\\
0.47619	0.15403\\
0.463659	0.160428\\
0.451128	0.167072\\
0.438596	0.173691\\
0.426065	0.180009\\
0.413534	0.185754\\
0.401003	0.190664\\
0.388471	0.19447\\
0.37594	0.196896\\
0.363409	0.197667\\
0.350877	0.197393\\
0.338346	0.196864\\
0.325815	0.196069\\
0.313283	0.195002\\
0.300752	0.193666\\
0.288221	0.192059\\
0.275689	0.190159\\
0.263158	0.187947\\
0.250627	0.185428\\
0.238095	0.182605\\
0.225564	0.179466\\
0.213033	0.17601\\
0.200501	0.172274\\
0.18797	0.168299\\
0.175439	0.164104\\
0.162907	0.159721\\
0.150376	0.155216\\
0.137845	0.150646\\
0.125313	0.144941\\
0.112782	0.137331\\
0.100251	0.128304\\
0.0877193	0.118309\\
0.075188	0.107711\\
0.0626566	0.0968459\\
0.0501253	0.0925432\\
0.037594	0.0875512\\
0.0250627	0.0675362\\
0.0125313	0.0420982\\
0	0.0120485\\
}--cycle;
\addplot [color=mycolor2, line width=1.0pt, forget plot]
  table[row sep=crcr]{%
-0.125	7.0363e-08\\
-0.1224	7.18105e-08\\
-0.1199	7.32022e-08\\
-0.1173	7.46494e-08\\
-0.1147	7.60964e-08\\
-0.1122	7.74876e-08\\
-0.1096	7.89343e-08\\
-0.1071	8.03252e-08\\
-0.1045	8.17715e-08\\
-0.1019	8.32178e-08\\
-0.0994	8.46082e-08\\
-0.0968	8.60542e-08\\
-0.0942	8.75e-08\\
-0.0917	8.889e-08\\
-0.0891	9.03355e-08\\
-0.0865	9.1781e-08\\
-0.084	9.31706e-08\\
-0.0814	9.46158e-08\\
-0.0789	9.60053e-08\\
-0.0763	9.74503e-08\\
-0.0737	9.88951e-08\\
-0.0712	1.00284e-07\\
-0.0686	1.01729e-07\\
-0.066	1.03174e-07\\
-0.0635	1.04562e-07\\
-0.0609	1.06007e-07\\
-0.0583	1.07451e-07\\
-0.0558	1.0884e-07\\
-0.0532	1.10284e-07\\
-0.0507	1.11673e-07\\
-0.0481	1.13117e-07\\
-0.0455	1.14561e-07\\
-0.043	1.15949e-07\\
-0.0404	1.17393e-07\\
-0.0378	1.18837e-07\\
-0.0353	1.20225e-07\\
-0.0327	1.21669e-07\\
-0.0301	1.23112e-07\\
-0.0276	1.245e-07\\
-0.025	1.25944e-07\\
-0.0224	1.27387e-07\\
-0.0199	1.28775e-07\\
-0.0173	1.30219e-07\\
-0.0148	1.31607e-07\\
-0.0122	1.3305e-07\\
-0.0096	1.34494e-07\\
-0.0071	1.35881e-07\\
-0.0045	1.37325e-07\\
-0.0019	1.38768e-07\\
0.0006	1.40156e-07\\
0.0032	0.000277198\\
0.0058	0.00138316\\
0.0083	0.00239605\\
0.0109	0.00339181\\
0.0134	0.0043857\\
0.016	0.00545108\\
0.0186	0.00648687\\
0.0211	0.00741343\\
0.0237	0.00831103\\
0.0263	0.00918117\\
0.0288	0.0100266\\
0.0314	0.0109246\\
0.034	0.0118268\\
0.0365	0.0126786\\
0.0391	0.0135389\\
0.0416	0.0143472\\
0.0442	0.015182\\
0.0468	0.0160207\\
0.0493	0.0168308\\
0.0519	0.0176689\\
0.0545	0.0184943\\
0.057	0.0192739\\
0.0596	0.0200754\\
0.0622	0.020875\\
0.0647	0.0216456\\
0.0673	0.022446\\
0.0698	0.0232088\\
0.0724	0.0239906\\
0.075	0.0247613\\
0.0775	0.0254971\\
0.0801	0.0262616\\
0.0827	0.027026\\
0.0852	0.0277568\\
0.0878	0.0285074\\
0.0904	0.0292462\\
0.0929	0.0299484\\
0.0955	0.030675\\
0.098	0.0313728\\
0.1006	0.0320957\\
0.1032	0.032811\\
0.1057	0.0334881\\
0.1083	0.0341816\\
0.1109	0.0348682\\
0.1134	0.035526\\
0.116	0.0362079\\
0.1186	0.0368842\\
0.1211	0.0375252\\
0.1237	0.0381806\\
0.1263	0.0388271\\
0.1288	0.0394442\\
0.1314	0.040084\\
0.1339	0.0406957\\
0.1365	0.0413246\\
0.1391	0.0419431\\
0.1416	0.0425285\\
0.1442	0.0431312\\
0.1468	0.0437312\\
0.1493	0.0443062\\
0.1519	0.0448997\\
0.1545	0.045485\\
0.157	0.046039\\
0.1596	0.0466083\\
0.1621	0.0471527\\
0.1647	0.0477178\\
0.1673	0.0482809\\
0.1698	0.0488175\\
0.1724	0.0493682\\
0.175	0.049912\\
0.1775	0.0504314\\
0.1801	0.050971\\
0.1827	0.0515106\\
0.1852	0.0520274\\
0.1878	0.0525599\\
0.1903	0.0530665\\
0.1929	0.0535895\\
0.1955	0.0541118\\
0.198	0.0546151\\
0.2006	0.0551388\\
0.2032	0.0556603\\
0.2057	0.0561577\\
0.2083	0.0566714\\
0.2109	0.057184\\
0.2134	0.0576783\\
0.216	0.0581941\\
0.2185	0.0586902\\
0.2211	0.0592038\\
0.2237	0.0597143\\
0.2262	0.0602035\\
0.2288	0.0607133\\
0.2314	0.0612254\\
0.2339	0.0617193\\
0.2365	0.0622322\\
0.2391	0.0627427\\
0.2416	0.0632314\\
0.2442	0.0637397\\
0.2467	0.0642301\\
0.2493	0.0647421\\
0.2519	0.0652546\\
0.2544	0.0657455\\
0.257	0.0662536\\
0.2596	0.0667601\\
0.2621	0.0672477\\
0.2647	0.0677562\\
0.2673	0.0682654\\
0.2698	0.0687536\\
0.2724	0.0692584\\
0.2749	0.069741\\
0.2775	0.0702415\\
0.2801	0.0707422\\
0.2826	0.0712236\\
0.2852	0.0717229\\
0.2878	0.0722187\\
0.2903	0.0726915\\
0.2929	0.0731798\\
0.2955	0.0736661\\
0.298	0.0741327\\
0.3006	0.0746159\\
0.3032	0.0750951\\
0.3057	0.0755508\\
0.3083	0.0760195\\
0.3108	0.0764662\\
0.3134	0.0769278\\
0.316	0.0773864\\
0.3185	0.0778229\\
0.3211	0.0782708\\
0.3237	0.0787118\\
0.3262	0.0791298\\
0.3288	0.0795596\\
0.3314	0.0799848\\
0.3339	0.0803887\\
0.3365	0.0808019\\
0.339	0.0811915\\
0.3416	0.0815889\\
0.3442	0.0819791\\
0.3467	0.0823485\\
0.3493	0.0827266\\
0.3519	0.0830973\\
0.3544	0.0834455\\
0.357	0.0837984\\
0.3596	0.0841425\\
0.3621	0.0844662\\
0.3647	0.0847958\\
0.3672	0.0851053\\
0.3698	0.0854183\\
0.3724	0.0857213\\
0.3749	0.0860033\\
0.3775	0.0862877\\
0.3801	0.0865639\\
0.3826	0.0868219\\
0.3852	0.0870812\\
0.3878	0.0873303\\
0.3903	0.0875599\\
0.3929	0.0877889\\
0.3954	0.0880006\\
0.398	0.0882125\\
0.4006	0.0884154\\
0.4031	0.0886011\\
0.4057	0.0887837\\
0.4083	0.088956\\
0.4108	0.0891125\\
0.4134	0.0892666\\
0.416	0.0894119\\
0.4185	0.0895427\\
0.4211	0.0896685\\
0.4236	0.0897796\\
0.4262	0.0898853\\
0.4288	0.089982\\
0.4313	0.0900667\\
0.4339	0.090146\\
0.4365	0.0902156\\
0.439	0.0902728\\
0.4416	0.0903224\\
0.4442	0.0903628\\
0.4467	0.0903935\\
0.4493	0.0904171\\
0.4519	0.0904316\\
0.4544	0.0904363\\
0.457	0.0904316\\
0.4595	0.0904183\\
0.4621	0.090396\\
0.4647	0.0903655\\
0.4672	0.0903282\\
0.4698	0.0902805\\
0.4724	0.0902233\\
0.4749	0.0901596\\
0.4775	0.0900849\\
0.4801	0.0900022\\
0.4826	0.0899151\\
0.4852	0.0898163\\
0.4877	0.0897128\\
0.4903	0.0895963\\
0.4929	0.0894711\\
0.4954	0.0893432\\
0.498	0.0892026\\
0.5006	0.0890541\\
0.5031	0.0889033\\
0.5057	0.0887379\\
0.5083	0.0885638\\
0.5108	0.0883889\\
0.5134	0.0881994\\
0.5159	0.0880101\\
0.5185	0.0878053\\
0.5211	0.0875921\\
0.5236	0.0873791\\
0.5262	0.0871494\\
0.5288	0.0869121\\
0.5313	0.0866768\\
0.5339	0.0864246\\
0.5365	0.0861643\\
0.539	0.085906\\
0.5416	0.0856292\\
0.5441	0.0853557\\
0.5467	0.0850638\\
0.5493	0.0847644\\
0.5518	0.084469\\
0.5544	0.0841538\\
0.557	0.0838302\\
0.5595	0.0835114\\
0.5621	0.0831724\\
0.5647	0.0828258\\
0.5672	0.0824852\\
0.5698	0.0821232\\
0.5723	0.0817671\\
0.5749	0.0813888\\
0.5775	0.0810025\\
0.58	0.0806239\\
0.5826	0.0802227\\
0.5852	0.0798135\\
0.5877	0.0794123\\
0.5903	0.0789868\\
0.5929	0.0785533\\
0.5954	0.0781291\\
0.598	0.0776805\\
0.6006	0.077224\\
0.6031	0.0767775\\
0.6057	0.0763049\\
0.6082	0.0758428\\
0.6108	0.0753544\\
0.6134	0.0748583\\
0.6159	0.0743741\\
0.6185	0.0738627\\
0.6211	0.0733431\\
0.6236	0.0728358\\
0.6262	0.0723003\\
0.6288	0.0717571\\
0.6313	0.0712275\\
0.6339	0.0706693\\
0.6364	0.070125\\
0.639	0.069551\\
0.6416	0.068969\\
0.6441	0.0684021\\
0.6467	0.0678051\\
0.6493	0.0672007\\
0.6518	0.0666122\\
0.6544	0.0659926\\
0.657	0.065365\\
0.6595	0.0647544\\
0.6621	0.0641121\\
0.6646	0.0634877\\
0.6672	0.0628311\\
0.6698	0.062167\\
0.6723	0.0615212\\
0.6749	0.0608423\\
0.6775	0.0601562\\
0.68	0.0594899\\
0.6826	0.0587901\\
0.6852	0.0580832\\
0.6877	0.0573967\\
0.6903	0.0566756\\
0.6928	0.0559758\\
0.6954	0.0552413\\
0.698	0.0545002\\
0.7005	0.0537814\\
0.7031	0.0530273\\
0.7057	0.0522663\\
0.7082	0.0515284\\
0.7108	0.0507546\\
0.7134	0.0499746\\
0.7159	0.0492189\\
0.7185	0.0484268\\
0.721	0.0476593\\
0.7236	0.0468549\\
0.7262	0.0460445\\
0.7287	0.0452598\\
0.7313	0.0444381\\
0.7339	0.0436108\\
0.7364	0.0428099\\
0.739	0.0419713\\
0.7416	0.0411271\\
0.7441	0.0403104\\
0.7467	0.0394558\\
0.7492	0.0386293\\
0.7518	0.0377647\\
0.7544	0.0368949\\
0.7569	0.0360538\\
0.7595	0.0351742\\
0.7621	0.0342899\\
0.7646	0.0334352\\
0.7672	0.032542\\
0.7698	0.0316441\\
0.7723	0.0307765\\
0.7749	0.0298699\\
0.7775	0.028959\\
0.78	0.0280794\\
0.7826	0.0271608\\
0.7851	0.0262737\\
0.7877	0.0253473\\
0.7903	0.024417\\
0.7928	0.023519\\
0.7954	0.0225816\\
0.798	0.021641\\
0.8005	0.0207333\\
0.8031	0.0197861\\
0.8057	0.0188357\\
0.8082	0.0179188\\
0.8108	0.0169624\\
0.8133	0.0160401\\
0.8159	0.0150784\\
0.8185	0.0141139\\
0.821	0.0131842\\
0.8236	0.0122147\\
0.8262	0.0112428\\
0.8287	0.0103063\\
0.8313	0.00933035\\
0.8339	0.00835242\\
0.8364	0.00741027\\
0.839	0.00642857\\
0.8415	0.00548298\\
0.8441	0.00449804\\
0.8467	0.00351171\\
0.8492	0.0025621\\
0.8518	0.00157331\\
0.8544	0.000583317\\
0.8569	-0.000369624\\
0.8595	-0.00136159\\
0.8621	-0.00235428\\
0.8646	-0.00330934\\
0.8672	-0.00430309\\
0.8697	-0.00525906\\
0.8723	-0.00625364\\
0.8749	-0.00724846\\
0.8774	-0.00820507\\
0.88	-0.00919983\\
0.8826	-0.0101943\\
0.8851	-0.0111503\\
0.8877	-0.0121442\\
0.8903	-0.0131376\\
0.8928	-0.0140921\\
0.8954	-0.015084\\
0.8979	-0.0160369\\
0.9005	-0.0170268\\
0.9031	-0.0180155\\
0.9056	-0.0189651\\
0.9082	-0.0199513\\
0.9108	-0.0209359\\
0.9133	-0.021881\\
0.9159	-0.0228621\\
0.9185	-0.0238413\\
0.921	-0.0247809\\
0.9236	-0.0257559\\
0.9262	-0.0267287\\
0.9287	-0.0276617\\
0.9313	-0.0286295\\
0.9338	-0.0295574\\
0.9364	-0.0305198\\
0.939	-0.0314792\\
0.9415	-0.0323988\\
0.9441	-0.033352\\
0.9467	-0.0343019\\
0.9492	-0.0352119\\
0.9518	-0.0361548\\
0.9544	-0.037094\\
0.9569	-0.0379934\\
0.9595	-0.038925\\
0.962	-0.0398168\\
0.9646	-0.0407402\\
0.9672	-0.0416591\\
0.9697	-0.0425385\\
0.9723	-0.0434486\\
0.9749	-0.044354\\
0.9774	-0.0452201\\
0.98	-0.0461158\\
0.9826	-0.0470065\\
0.9851	-0.047858\\
0.9877	-0.0487384\\
0.9902	-0.0495798\\
0.9928	-0.0504494\\
0.9954	-0.0513134\\
0.9979	-0.0521387\\
1.0005	-0.0529912\\
1.0031	-0.0538377\\
1.0056	-0.0546459\\
1.0082	-0.0554803\\
1.0108	-0.0563085\\
1.0133	-0.0570987\\
1.0159	-0.057914\\
1.0184	-0.0586917\\
1.021	-0.059494\\
1.0236	-0.0602894\\
1.0261	-0.0610477\\
1.0287	-0.0618295\\
1.0313	-0.062604\\
1.0338	-0.0633419\\
1.0364	-0.0641022\\
1.039	-0.064855\\
1.0415	-0.0655718\\
1.0441	-0.0663099\\
1.0466	-0.0670123\\
1.0492	-0.0677351\\
1.0518	-0.06845\\
1.0543	-0.06913\\
1.0569	-0.0698293\\
1.0595	-0.0705205\\
1.062	-0.0711773\\
1.0646	-0.0718523\\
1.0672	-0.072519\\
1.0697	-0.073152\\
1.0723	-0.0738021\\
1.0748	-0.0744191\\
1.0774	-0.0750523\\
1.08	-0.0756768\\
1.0825	-0.076269\\
1.0851	-0.0768762\\
1.0877	-0.0774745\\
1.0902	-0.0780413\\
1.0928	-0.078622\\
1.0954	-0.0791936\\
1.0979	-0.0797346\\
1.1005	-0.0802882\\
1.1031	-0.0808325\\
1.1056	-0.0813472\\
1.1082	-0.0818733\\
1.1107	-0.0823703\\
1.1133	-0.0828779\\
1.1159	-0.083376\\
1.1184	-0.083846\\
1.121	-0.0843254\\
1.1236	-0.0847952\\
1.1261	-0.0852379\\
1.1287	-0.0856888\\
1.1313	-0.0861299\\
1.1338	-0.0865449\\
1.1364	-0.086967\\
1.1389	-0.0873636\\
1.1415	-0.0877665\\
1.1441	-0.0881595\\
1.1466	-0.0885281\\
1.1492	-0.0889018\\
1.1518	-0.0892656\\
1.1543	-0.0896061\\
1.1569	-0.0899505\\
1.1595	-0.090285\\
1.162	-0.0905972\\
1.1646	-0.0909123\\
1.1671	-0.0912058\\
1.1697	-0.0915013\\
1.1723	-0.091787\\
1.1748	-0.0920523\\
1.1774	-0.0923185\\
1.18	-0.0925748\\
1.1825	-0.0928119\\
1.1851	-0.0930489\\
1.1877	-0.0932759\\
1.1902	-0.093485\\
1.1928	-0.0936928\\
1.1953	-0.0938834\\
1.1979	-0.0940721\\
1.2005	-0.094251\\
1.203	-0.0944139\\
1.2056	-0.0945738\\
1.2082	-0.0947241\\
1.2107	-0.0948596\\
1.2133	-0.0949911\\
1.2159	-0.0951132\\
1.2184	-0.0952216\\
1.221	-0.0953251\\
1.2235	-0.0954159\\
1.2261	-0.0955011\\
1.2287	-0.0955772\\
1.2312	-0.0956416\\
1.2338	-0.0956996\\
1.2364	-0.0957486\\
1.2389	-0.0957872\\
1.2415	-0.0958187\\
1.2441	-0.0958412\\
1.2466	-0.0958547\\
1.2492	-0.0958601\\
1.2518	-0.0958569\\
1.2543	-0.0958457\\
1.2569	-0.0958258\\
1.2594	-0.0957988\\
1.262	-0.0957626\\
1.2646	-0.0957182\\
1.2671	-0.0956678\\
1.2697	-0.0956075\\
1.2723	-0.0955393\\
1.2748	-0.0954664\\
1.2774	-0.0953829\\
1.28	-0.0952919\\
1.2825	-0.0951972\\
1.2851	-0.0950915\\
1.2876	-0.0949829\\
1.2902	-0.0948629\\
1.2928	-0.0947359\\
1.2953	-0.0946071\\
1.2979	-0.0944665\\
1.3005	-0.0943191\\
1.303	-0.0941711\\
1.3056	-0.0940108\\
1.3082	-0.0938441\\
1.3107	-0.0936779\\
1.3133	-0.093499\\
1.3158	-0.0933213\\
1.3184	-0.0931306\\
1.321	-0.0929342\\
1.3235	-0.09274\\
1.3261	-0.0925327\\
1.3287	-0.0923199\\
1.3312	-0.0921103\\
1.3338	-0.0918873\\
1.3364	-0.0916593\\
1.3389	-0.0914354\\
1.3415	-0.091198\\
1.344	-0.0909653\\
1.3466	-0.0907189\\
1.3492	-0.0904682\\
1.3517	-0.0902231\\
1.3543	-0.0899643\\
1.3569	-0.0897015\\
1.3594	-0.0894452\\
1.362	-0.0891751\\
1.3646	-0.0889014\\
1.3671	-0.088635\\
1.3697	-0.0883548\\
1.3722	-0.0880824\\
1.3748	-0.0877962\\
1.3774	-0.0875071\\
1.3799	-0.0872266\\
1.3825	-0.0869323\\
1.3851	-0.0866356\\
1.3876	-0.0863481\\
1.3902	-0.0860469\\
1.3928	-0.0857438\\
1.3953	-0.0854504\\
1.3979	-0.0851436\\
1.4005	-0.0848351\\
1.403	-0.0845371\\
1.4056	-0.0842257\\
1.4081	-0.0839252\\
1.4107	-0.0836115\\
1.4133	-0.0832969\\
1.4158	-0.0829935\\
1.4184	-0.0826772\\
1.421	-0.0823603\\
1.4235	-0.0820551\\
1.4261	-0.0817373\\
1.4287	-0.0814193\\
1.4312	-0.0811134\\
1.4338	-0.0807953\\
1.4363	-0.0804894\\
1.4389	-0.0801716\\
1.4415	-0.0798541\\
1.444	-0.0795492\\
1.4466	-0.0792327\\
1.4492	-0.0789169\\
1.4517	-0.0786139\\
1.4543	-0.0782997\\
1.4569	-0.0779865\\
1.4594	-0.0776864\\
1.462	-0.0773755\\
1.4645	-0.0770778\\
1.4671	-0.0767695\\
1.4697	-0.0764628\\
1.4722	-0.0761694\\
1.4748	-0.0758659\\
1.4774	-0.0755642\\
1.4799	-0.0752759\\
1.4825	-0.074978\\
1.4851	-0.0746821\\
1.4876	-0.0743997\\
1.4902	-0.0741081\\
1.4927	-0.07383\\
1.4953	-0.0735431\\
1.4979	-0.0732586\\
1.5004	-0.0729875\\
1.503	-0.0727082\\
1.5056	-0.0724315\\
1.5081	-0.0721681\\
1.5107	-0.0718969\\
1.5133	-0.0716286\\
1.5158	-0.0713735\\
1.5184	-0.0711111\\
1.5209	-0.0708617\\
1.5235	-0.0706054\\
1.5261	-0.0703523\\
1.5286	-0.070112\\
1.5312	-0.0698653\\
1.5338	-0.069622\\
1.5363	-0.0693912\\
1.5389	-0.0691546\\
1.5415	-0.0689214\\
1.544	-0.0687006\\
1.5466	-0.0684744\\
1.5491	-0.0682603\\
1.5517	-0.0680411\\
1.5543	-0.0678256\\
1.5568	-0.0676219\\
1.5594	-0.0674137\\
1.562	-0.0672092\\
1.5645	-0.0670161\\
1.5671	-0.066819\\
1.5697	-0.0666257\\
1.5722	-0.0664434\\
1.5748	-0.0662575\\
1.5774	-0.0660755\\
1.5799	-0.0659041\\
1.5825	-0.0657296\\
1.585	-0.0655655\\
1.5876	-0.0653985\\
1.5902	-0.0652354\\
1.5927	-0.0650822\\
1.5953	-0.0649267\\
1.5979	-0.064775\\
1.6004	-0.0646327\\
1.603	-0.0644885\\
1.6056	-0.0643481\\
1.6081	-0.0642166\\
1.6107	-0.0640836\\
1.6132	-0.0639593\\
1.6158	-0.0638336\\
1.6184	-0.0637117\\
1.6209	-0.0635979\\
1.6235	-0.0634832\\
1.6261	-0.0633721\\
1.6286	-0.0632686\\
1.6312	-0.0631646\\
1.6338	-0.0630641\\
1.6363	-0.0629707\\
1.6389	-0.062877\\
1.6414	-0.0627902\\
1.644	-0.0627032\\
1.6466	-0.0626196\\
1.6491	-0.0625423\\
1.6517	-0.0624651\\
1.6543	-0.0623912\\
1.6568	-0.062323\\
1.6594	-0.0622553\\
1.662	-0.0621906\\
1.6645	-0.0621312\\
1.6671	-0.0620724\\
1.6696	-0.0620186\\
1.6722	-0.0619655\\
1.6748	-0.0619152\\
1.6773	-0.0618695\\
1.6799	-0.0618245\\
1.6825	-0.0617822\\
1.685	-0.061744\\
1.6876	-0.0617067\\
1.6902	-0.0616719\\
1.6927	-0.0616407\\
1.6953	-0.0616105\\
1.6978	-0.0615837\\
1.7004	-0.0615579\\
1.703	-0.0615343\\
1.7055	-0.0615136\\
1.7081	-0.0614941\\
1.7107	-0.0614765\\
1.7132	-0.0614613\\
1.7158	-0.0614474\\
1.7184	-0.0614352\\
1.7209	-0.0614251\\
1.7235	-0.0614161\\
1.7261	-0.0614087\\
1.7286	-0.061403\\
1.7312	-0.0613985\\
1.7337	-0.0613954\\
1.7363	-0.0613935\\
1.7389	-0.0613927\\
1.7414	-0.0613931\\
1.744	-0.0613946\\
1.7466	-0.061397\\
1.7491	-0.0614003\\
1.7517	-0.0614045\\
1.7543	-0.0614096\\
1.7568	-0.0614151\\
1.7594	-0.0614215\\
1.7619	-0.0614283\\
1.7645	-0.0614358\\
1.7671	-0.0614438\\
1.7696	-0.0614519\\
1.7722	-0.0614606\\
1.7748	-0.0614696\\
1.7773	-0.0614785\\
1.7799	-0.0614878\\
1.7825	-0.0614973\\
1.785	-0.0615063\\
1.7876	-0.0615157\\
1.7901	-0.0615246\\
1.7927	-0.0615337\\
1.7953	-0.0615426\\
1.7978	-0.0615509\\
1.8004	-0.0615591\\
1.803	-0.0615669\\
1.8055	-0.061574\\
1.8081	-0.0615808\\
1.8107	-0.061587\\
1.8132	-0.0615924\\
1.8158	-0.0615973\\
1.8183	-0.0616014\\
1.8209	-0.0616048\\
1.8235	-0.0616073\\
1.826	-0.0616089\\
1.8286	-0.0616096\\
1.8312	-0.0616093\\
1.8337	-0.061608\\
1.8363	-0.0616056\\
1.8389	-0.061602\\
1.8414	-0.0615974\\
1.844	-0.0615914\\
1.8465	-0.0615845\\
1.8491	-0.0615759\\
1.8517	-0.061566\\
1.8542	-0.0615551\\
1.8568	-0.0615424\\
1.8594	-0.0615282\\
1.8619	-0.0615132\\
1.8645	-0.0614959\\
1.8671	-0.0614771\\
1.8696	-0.0614575\\
1.8722	-0.0614355\\
1.8747	-0.0614127\\
1.8773	-0.0613874\\
1.8799	-0.0613603\\
1.8824	-0.0613326\\
1.885	-0.061302\\
1.8876	-0.0612696\\
1.8901	-0.0612367\\
1.8927	-0.0612007\\
1.8953	-0.0611628\\
1.8978	-0.0611245\\
1.9004	-0.0610829\\
1.903	-0.0610393\\
1.9055	-0.0609956\\
1.9081	-0.0609482\\
1.9106	-0.0609007\\
1.9132	-0.0608495\\
1.9158	-0.0607962\\
1.9183	-0.0607431\\
1.9209	-0.0606859\\
1.9235	-0.0606268\\
1.926	-0.060568\\
1.9286	-0.0605049\\
1.9312	-0.0604397\\
1.9337	-0.0603752\\
1.9363	-0.0603062\\
1.9388	-0.0602379\\
1.9414	-0.0601649\\
1.944	-0.06009\\
1.9465	-0.060016\\
1.9491	-0.0599372\\
1.9517	-0.0598565\\
1.9542	-0.059777\\
1.9568	-0.0596924\\
1.9594	-0.059606\\
1.9619	-0.059521\\
1.9645	-0.0594309\\
1.967	-0.0593425\\
1.9696	-0.0592487\\
1.9722	-0.0591531\\
1.9747	-0.0590596\\
1.9773	-0.0589605\\
1.9799	-0.0588598\\
1.9824	-0.0587613\\
1.985	-0.0586572\\
1.9876	-0.0585516\\
1.9901	-0.0584484\\
1.9927	-0.0583396\\
1.9952	-0.0582336\\
1.9978	-0.0581218\\
2.0004	-0.0580086\\
2.0029	-0.0578984\\
2.0055	-0.0577825\\
2.0081	-0.0576652\\
2.0106	-0.0575512\\
2.0132	-0.0574314\\
2.0158	-0.0573105\\
2.0183	-0.0571931\\
2.0209	-0.0570699\\
2.0234	-0.0569505\\
2.026	-0.0568253\\
2.0286	-0.0566992\\
2.0311	-0.0565771\\
2.0337	-0.0564493\\
2.0363	-0.0563207\\
2.0388	-0.0561963\\
2.0414	-0.0560663\\
2.044	-0.0559357\\
2.0465	-0.0558096\\
2.0491	-0.0556779\\
2.0517	-0.0555458\\
2.0542	-0.0554185\\
2.0568	-0.0552857\\
2.0593	-0.0551578\\
2.0619	-0.0550246\\
2.0645	-0.0548913\\
2.067	-0.0547631\\
2.0696	-0.0546297\\
2.0722	-0.0544964\\
2.0747	-0.0543683\\
2.0773	-0.0542354\\
2.0799	-0.0541027\\
2.0824	-0.0539754\\
2.085	-0.0538434\\
2.0875	-0.053717\\
2.0901	-0.053586\\
2.0927	-0.0534556\\
2.0952	-0.0533309\\
2.0978	-0.053202\\
2.1004	-0.0530738\\
2.1029	-0.0529515\\
2.1055	-0.0528251\\
2.1081	-0.0526998\\
2.1106	-0.0525804\\
2.1132	-0.0524573\\
2.1157	-0.0523401\\
2.1183	-0.0522195\\
2.1209	-0.0521003\\
2.1234	-0.0519871\\
2.126	-0.0518707\\
2.1286	-0.051756\\
2.1311	-0.0516472\\
2.1337	-0.0515358\\
2.1363	-0.0514261\\
2.1388	-0.0513224\\
2.1414	-0.0512164\\
2.1439	-0.0511164\\
2.1465	-0.0510143\\
2.1491	-0.0509144\\
2.1516	-0.0508203\\
2.1542	-0.0507247\\
2.1568	-0.0506313\\
2.1593	-0.0505438\\
2.1619	-0.0504551\\
2.1645	-0.0503689\\
2.167	-0.0502884\\
2.1696	-0.0502072\\
2.1721	-0.0501316\\
2.1747	-0.0500557\\
2.1773	-0.0499824\\
2.1798	-0.0499147\\
2.1824	-0.049847\\
2.185	-0.0497822\\
2.1875	-0.0497227\\
2.1901	-0.0496638\\
2.1927	-0.0496079\\
2.1952	-0.049557\\
2.1978	-0.0495072\\
2.2004	-0.0494606\\
2.2029	-0.0494188\\
2.2055	-0.0493785\\
2.208	-0.0493429\\
2.2106	-0.0493091\\
2.2132	-0.0492787\\
2.2157	-0.0492526\\
2.2183	-0.0492289\\
2.2209	-0.0492086\\
2.2234	-0.0491924\\
2.226	-0.049179\\
2.2286	-0.0491692\\
2.2311	-0.0491631\\
2.2337	-0.0491602\\
2.2362	-0.0491609\\
2.2388	-0.0491651\\
2.2414	-0.049173\\
2.2439	-0.0491841\\
2.2465	-0.0491991\\
2.2491	-0.0492179\\
2.2516	-0.0492395\\
2.2542	-0.0492655\\
2.2568	-0.0492952\\
2.2593	-0.0493273\\
2.2619	-0.0493643\\
2.2644	-0.0494035\\
2.267	-0.0494478\\
2.2696	-0.0494958\\
2.2721	-0.0495455\\
2.2747	-0.0496009\\
2.2773	-0.0496599\\
2.2798	-0.0497202\\
2.2824	-0.0497864\\
2.285	-0.0498564\\
2.2875	-0.0499271\\
2.2901	-0.0500042\\
2.2926	-0.0500817\\
2.2952	-0.0501659\\
2.2978	-0.0502537\\
2.3003	-0.0503415\\
2.3029	-0.0504362\\
2.3055	-0.0505345\\
2.308	-0.0506323\\
2.3106	-0.0507374\\
2.3132	-0.0508459\\
2.3157	-0.0509534\\
2.3183	-0.0510685\\
2.3208	-0.0511824\\
2.3234	-0.051304\\
2.326	-0.0514289\\
2.3285	-0.051552\\
2.3311	-0.0516832\\
2.3337	-0.0518174\\
2.3362	-0.0519495\\
2.3388	-0.0520897\\
2.3414	-0.052233\\
2.3439	-0.0523736\\
2.3465	-0.0525226\\
2.349	-0.0526686\\
2.3516	-0.0528231\\
2.3542	-0.0529804\\
2.3567	-0.0531342\\
2.3593	-0.0532967\\
2.3619	-0.0534618\\
2.3644	-0.053623\\
2.367	-0.053793\\
2.3696	-0.0539654\\
2.3721	-0.0541333\\
2.3747	-0.0543103\\
2.3773	-0.0544894\\
2.3798	-0.0546637\\
2.3824	-0.054847\\
2.3849	-0.0550252\\
2.3875	-0.0552124\\
2.3901	-0.0554016\\
2.3926	-0.0555851\\
2.3952	-0.0557778\\
2.3978	-0.0559721\\
2.4003	-0.0561605\\
2.4029	-0.0563579\\
2.4055	-0.0565568\\
2.408	-0.0567494\\
2.4106	-0.056951\\
2.4131	-0.057146\\
2.4157	-0.0573501\\
2.4183	-0.0575552\\
2.4208	-0.0577534\\
2.4234	-0.0579605\\
2.426	-0.0581685\\
2.4285	-0.0583693\\
2.4311	-0.0585788\\
2.4337	-0.058789\\
2.4362	-0.0589917\\
2.4388	-0.0592029\\
2.4413	-0.0594065\\
2.4439	-0.0596186\\
2.4465	-0.059831\\
2.449	-0.0600354\\
2.4516	-0.0602481\\
2.4542	-0.0604609\\
2.4567	-0.0606655\\
2.4593	-0.0608782\\
2.4619	-0.0610907\\
2.4644	-0.0612948\\
2.467	-0.0615068\\
2.4695	-0.0617102\\
2.4721	-0.0619214\\
2.4747	-0.062132\\
2.4772	-0.0623339\\
2.4798	-0.0625432\\
2.4824	-0.0627518\\
2.4849	-0.0629516\\
2.4875	-0.0631584\\
2.4901	-0.0633643\\
2.4926	-0.0635613\\
2.4952	-0.063765\\
2.4977	-0.0639598\\
2.5003	-0.0641611\\
2.5029	-0.0643612\\
2.5054	-0.0645522\\
2.508	-0.0647494\\
2.5106	-0.064945\\
2.5131	-0.0651316\\
2.5157	-0.0653241\\
2.5183	-0.0655148\\
2.5208	-0.0656965\\
2.5234	-0.0658836\\
2.526	-0.0660688\\
2.5285	-0.066245\\
2.5311	-0.0664262\\
2.5336	-0.0665985\\
2.5362	-0.0667755\\
2.5388	-0.0669503\\
2.5413	-0.0671163\\
2.5439	-0.0672865\\
2.5465	-0.0674544\\
2.549	-0.0676135\\
2.5516	-0.0677765\\
2.5542	-0.0679369\\
2.5567	-0.0680887\\
2.5593	-0.068244\\
2.5618	-0.0683907\\
2.5644	-0.0685406\\
2.567	-0.0686877\\
2.5695	-0.0688265\\
2.5721	-0.0689679\\
2.5747	-0.0691064\\
2.5772	-0.0692368\\
2.5798	-0.0693694\\
2.5824	-0.0694989\\
2.5849	-0.0696206\\
2.5875	-0.069744\\
2.59	-0.0698597\\
2.5926	-0.0699768\\
2.5952	-0.0700907\\
2.5977	-0.0701971\\
2.6003	-0.0703046\\
2.6029	-0.0704086\\
2.6054	-0.0705056\\
2.608	-0.070603\\
2.6106	-0.070697\\
2.6131	-0.0707842\\
2.6157	-0.0708714\\
2.6182	-0.070952\\
2.6208	-0.0710324\\
2.6234	-0.0711092\\
2.6259	-0.0711797\\
2.6285	-0.0712496\\
2.6311	-0.0713158\\
2.6336	-0.0713762\\
2.6362	-0.0714353\\
2.6388	-0.0714909\\
2.6413	-0.0715409\\
2.6439	-0.0715894\\
2.6464	-0.0716325\\
2.649	-0.0716738\\
2.6516	-0.0717114\\
2.6541	-0.0717442\\
2.6567	-0.0717746\\
2.6593	-0.0718014\\
2.6618	-0.0718237\\
2.6644	-0.0718434\\
2.667	-0.0718593\\
2.6695	-0.0718713\\
2.6721	-0.0718801\\
2.6746	-0.0718852\\
2.6772	-0.0718869\\
2.6798	-0.0718851\\
2.6823	-0.0718799\\
2.6849	-0.071871\\
2.6875	-0.0718585\\
2.69	-0.0718432\\
2.6926	-0.0718238\\
2.6952	-0.0718009\\
2.6977	-0.0717756\\
2.7003	-0.0717459\\
2.7029	-0.0717127\\
2.7054	-0.0716776\\
2.708	-0.0716378\\
2.7105	-0.0715963\\
2.7131	-0.0715499\\
2.7157	-0.0715002\\
2.7182	-0.0714494\\
2.7208	-0.0713933\\
2.7234	-0.071334\\
2.7259	-0.0712741\\
2.7285	-0.0712086\\
2.7311	-0.0711401\\
2.7336	-0.0710713\\
2.7362	-0.0709969\\
2.7387	-0.0709225\\
2.7413	-0.0708422\\
2.7439	-0.0707591\\
2.7464	-0.0706765\\
2.749	-0.0705879\\
2.7516	-0.0704965\\
2.7541	-0.0704061\\
2.7567	-0.0703095\\
2.7593	-0.0702103\\
2.7618	-0.0701125\\
2.7644	-0.0700084\\
2.7669	-0.069906\\
2.7695	-0.0697971\\
2.7721	-0.069686\\
2.7746	-0.0695769\\
2.7772	-0.0694614\\
2.7798	-0.0693437\\
2.7823	-0.0692286\\
2.7849	-0.0691069\\
2.7875	-0.0689832\\
2.79	-0.0688624\\
2.7926	-0.0687351\\
2.7951	-0.068611\\
2.7977	-0.0684802\\
2.8003	-0.0683478\\
2.8028	-0.068219\\
2.8054	-0.0680835\\
2.808	-0.0679466\\
2.8105	-0.0678137\\
2.8131	-0.0676742\\
2.8157	-0.0675334\\
2.8182	-0.0673969\\
2.8208	-0.067254\\
2.8233	-0.0671155\\
2.8259	-0.0669706\\
2.8285	-0.0668248\\
2.831	-0.0666838\\
2.8336	-0.0665364\\
2.8362	-0.0663884\\
2.8387	-0.0662455\\
2.8413	-0.0660964\\
2.8439	-0.0659469\\
2.8464	-0.0658027\\
2.849	-0.0656525\\
2.8516	-0.065502\\
2.8541	-0.0653572\\
2.8567	-0.0652066\\
2.8592	-0.0650617\\
2.8618	-0.064911\\
2.8644	-0.0647606\\
2.8669	-0.0646161\\
2.8695	-0.0644661\\
2.8721	-0.0643165\\
2.8746	-0.064173\\
2.8772	-0.0640243\\
2.8798	-0.0638762\\
2.8823	-0.0637344\\
2.8849	-0.0635877\\
2.8874	-0.0634473\\
2.89	-0.0633022\\
2.8926	-0.063158\\
2.8951	-0.0630203\\
2.8977	-0.0628782\\
2.9003	-0.0627372\\
2.9028	-0.0626028\\
2.9054	-0.0624643\\
2.908	-0.0623272\\
2.9105	-0.0621966\\
2.9131	-0.0620623\\
2.9156	-0.0619346\\
2.9182	-0.0618034\\
2.9208	-0.0616739\\
2.9233	-0.061551\\
2.9259	-0.061425\\
2.9285	-0.0613009\\
2.931	-0.0611833\\
2.9336	-0.061063\\
2.9362	-0.0609447\\
2.9387	-0.0608329\\
2.9413	-0.0607188\\
2.9438	-0.0606112\\
2.9464	-0.0605015\\
2.949	-0.0603941\\
2.9515	-0.0602931\\
2.9541	-0.0601904\\
2.9567	-0.0600902\\
2.9592	-0.0599962\\
2.9618	-0.059901\\
2.9644	-0.0598084\\
2.9669	-0.0597218\\
2.9695	-0.0596345\\
2.972	-0.059553\\
2.9746	-0.0594711\\
2.9772	-0.0593919\\
2.9797	-0.0593185\\
2.9823	-0.059245\\
2.9849	-0.0591744\\
2.9874	-0.0591092\\
2.99	-0.0590445\\
2.9926	-0.0589827\\
2.9951	-0.0589261\\
2.9977	-0.0588703\\
3.0003	-0.0588176\\
3.0028	-0.0587699\\
3.0054	-0.0587233\\
3.0079	-0.0586815\\
3.0105	-0.0586411\\
3.0131	-0.0586039\\
3.0156	-0.0585712\\
3.0182	-0.0585402\\
3.0208	-0.0585126\\
3.0233	-0.058489\\
3.0259	-0.0584677\\
3.0285	-0.0584496\\
3.031	-0.0584352\\
3.0336	-0.0584235\\
3.0361	-0.0584154\\
3.0387	-0.05841\\
3.0413	-0.058408\\
3.0438	-0.058409\\
3.0464	-0.0584133\\
3.049	-0.0584209\\
3.0515	-0.0584312\\
3.0541	-0.058445\\
3.0567	-0.0584621\\
3.0592	-0.0584815\\
3.0618	-0.0585048\\
3.0643	-0.0585301\\
3.0669	-0.0585596\\
3.0695	-0.0585922\\
3.072	-0.0586264\\
3.0746	-0.0586651\\
3.0772	-0.0587067\\
3.0797	-0.0587497\\
3.0823	-0.0587972\\
3.0849	-0.0588478\\
3.0874	-0.0588991\\
3.09	-0.0589554\\
3.0925	-0.0590121\\
3.0951	-0.059074\\
3.0977	-0.0591386\\
3.1002	-0.0592033\\
3.1028	-0.0592732\\
3.1054	-0.0593459\\
3.1079	-0.0594182\\
3.1105	-0.0594959\\
3.1131	-0.0595762\\
3.1156	-0.0596557\\
3.1182	-0.0597408\\
3.1207	-0.0598249\\
3.1233	-0.0599146\\
3.1259	-0.0600067\\
3.1284	-0.0600973\\
3.131	-0.0601936\\
3.1336	-0.0602921\\
3.1361	-0.0603888\\
3.1387	-0.0604913\\
3.1413	-0.0605957\\
3.1438	-0.0606979\\
3.1464	-0.0608061\\
3.1489	-0.0609117\\
3.1515	-0.0610232\\
3.1541	-0.0611364\\
3.1566	-0.0612468\\
3.1592	-0.061363\\
3.1618	-0.0614807\\
3.1643	-0.0615952\\
3.1669	-0.0617156\\
3.1695	-0.0618372\\
3.172	-0.0619553\\
3.1746	-0.0620792\\
3.1772	-0.0622042\\
3.1797	-0.0623252\\
3.1823	-0.0624521\\
3.1848	-0.0625748\\
3.1874	-0.0627032\\
3.19	-0.0628324\\
3.1925	-0.0629571\\
3.1951	-0.0630874\\
3.1977	-0.0632182\\
3.2002	-0.0633444\\
3.2028	-0.0634759\\
3.2054	-0.0636077\\
3.2079	-0.0637346\\
3.2105	-0.0638668\\
3.213	-0.0639939\\
3.2156	-0.064126\\
3.2182	-0.0642581\\
3.2207	-0.064385\\
3.2233	-0.0645167\\
3.2259	-0.064648\\
3.2284	-0.064774\\
3.231	-0.0649046\\
3.2336	-0.0650346\\
3.2361	-0.0651591\\
3.2387	-0.0652879\\
3.2412	-0.065411\\
3.2438	-0.0655383\\
3.2464	-0.0656646\\
3.2489	-0.0657853\\
3.2515	-0.0659097\\
3.2541	-0.066033\\
3.2566	-0.0661505\\
3.2592	-0.0662715\\
3.2618	-0.0663912\\
3.2643	-0.066505\\
3.2669	-0.066622\\
3.2694	-0.066733\\
3.272	-0.0668469\\
3.2746	-0.0669592\\
3.2771	-0.0670656\\
3.2797	-0.0671745\\
3.2823	-0.0672816\\
3.2848	-0.0673828\\
3.2874	-0.0674861\\
3.29	-0.0675874\\
3.2925	-0.0676829\\
3.2951	-0.0677801\\
3.2976	-0.0678714\\
3.3002	-0.0679643\\
3.3028	-0.0680548\\
3.3053	-0.0681396\\
3.3079	-0.0682255\\
3.3105	-0.0683089\\
3.313	-0.0683867\\
3.3156	-0.0684651\\
3.3182	-0.0685408\\
3.3207	-0.0686112\\
3.3233	-0.0686816\\
3.3259	-0.0687493\\
3.3284	-0.0688118\\
3.331	-0.0688739\\
3.3335	-0.0689309\\
3.3361	-0.0689873\\
3.3387	-0.0690407\\
3.3412	-0.0690892\\
3.3438	-0.0691366\\
3.3464	-0.0691809\\
3.3489	-0.0692205\\
3.3515	-0.0692585\\
3.3541	-0.0692934\\
3.3566	-0.0693238\\
3.3592	-0.0693522\\
3.3617	-0.0693764\\
3.3643	-0.0693983\\
3.3669	-0.0694168\\
3.3694	-0.0694313\\
3.372	-0.0694431\\
3.3746	-0.0694514\\
3.3771	-0.0694561\\
3.3797	-0.0694576\\
3.3823	-0.0694555\\
3.3848	-0.0694501\\
3.3874	-0.0694411\\
3.3899	-0.0694289\\
3.3925	-0.0694128\\
3.3951	-0.0693931\\
3.3976	-0.0693706\\
3.4002	-0.0693438\\
3.4028	-0.0693132\\
3.4053	-0.0692804\\
3.4079	-0.0692426\\
3.4105	-0.0692012\\
3.413	-0.0691579\\
3.4156	-0.0691092\\
3.4181	-0.0690589\\
3.4207	-0.069003\\
3.4233	-0.0689433\\
3.4258	-0.0688825\\
3.4284	-0.0688156\\
3.431	-0.068745\\
3.4335	-0.0686737\\
3.4361	-0.0685959\\
3.4387	-0.0685144\\
3.4412	-0.0684326\\
3.4438	-0.0683439\\
3.4463	-0.0682552\\
3.4489	-0.0681594\\
3.4515	-0.0680599\\
3.454	-0.0679609\\
3.4566	-0.0678544\\
3.4592	-0.0677444\\
3.4617	-0.0676352\\
3.4643	-0.0675182\\
3.4669	-0.0673977\\
3.4694	-0.0672785\\
3.472	-0.0671512\\
3.4745	-0.0670255\\
3.4771	-0.0668914\\
3.4797	-0.066754\\
3.4822	-0.0666187\\
3.4848	-0.0664747\\
3.4874	-0.0663274\\
3.4899	-0.0661827\\
3.4925	-0.0660291\\
3.4951	-0.0658722\\
3.4976	-0.0657185\\
3.5002	-0.0655555\\
3.5028	-0.0653894\\
3.5053	-0.0652269\\
3.5079	-0.0650549\\
3.5104	-0.0648867\\
3.513	-0.064709\\
3.5156	-0.0645283\\
3.5181	-0.064352\\
3.5207	-0.0641659\\
3.5233	-0.063977\\
3.5258	-0.0637929\\
3.5284	-0.0635988\\
3.531	-0.0634022\\
3.5335	-0.0632107\\
3.5361	-0.0630091\\
3.5386	-0.062813\\
3.5412	-0.0626067\\
3.5438	-0.0623981\\
3.5463	-0.0621954\\
3.5489	-0.0619824\\
3.5515	-0.0617673\\
3.554	-0.0615585\\
3.5566	-0.0613393\\
3.5592	-0.0611181\\
3.5617	-0.0609037\\
3.5643	-0.0606788\\
3.5668	-0.0604609\\
3.5694	-0.0602326\\
3.572	-0.0600027\\
3.5745	-0.0597801\\
3.5771	-0.059547\\
3.5797	-0.0593126\\
3.5822	-0.0590858\\
3.5848	-0.0588486\\
3.5874	-0.0586102\\
3.5899	-0.0583798\\
3.5925	-0.058139\\
3.595	-0.0579066\\
3.5976	-0.0576639\\
3.6002	-0.0574202\\
3.6027	-0.0571851\\
3.6053	-0.0569399\\
3.6079	-0.0566939\\
3.6104	-0.0564568\\
3.613	-0.0562097\\
3.6156	-0.055962\\
3.6181	-0.0557235\\
3.6207	-0.0554751\\
3.6232	-0.055236\\
3.6258	-0.0549872\\
3.6284	-0.0547382\\
3.6309	-0.0544987\\
3.6335	-0.0542497\\
3.6361	-0.0540007\\
3.6386	-0.0537615\\
3.6412	-0.053513\\
3.6438	-0.0532647\\
3.6463	-0.0530264\\
3.6489	-0.052779\\
3.6515	-0.0525321\\
3.654	-0.0522953\\
3.6566	-0.0520496\\
3.6591	-0.0518142\\
3.6617	-0.0515701\\
3.6643	-0.0513269\\
3.6668	-0.051094\\
3.6694	-0.0508528\\
3.672	-0.0506126\\
3.6745	-0.0503829\\
3.6771	-0.0501452\\
3.6797	-0.0499088\\
3.6822	-0.0496828\\
3.6848	-0.0494492\\
3.6873	-0.049226\\
3.6899	-0.0489955\\
3.6925	-0.0487666\\
3.695	-0.0485482\\
3.6976	-0.0483228\\
3.7002	-0.0480993\\
3.7027	-0.0478862\\
3.7053	-0.0476666\\
3.7079	-0.047449\\
3.7104	-0.0472418\\
3.713	-0.0470285\\
3.7155	-0.0468254\\
3.7181	-0.0466166\\
3.7207	-0.0464101\\
3.7232	-0.0462139\\
3.7258	-0.0460122\\
3.7284	-0.0458131\\
3.7309	-0.0456241\\
3.7335	-0.0454302\\
3.7361	-0.045239\\
3.7386	-0.0450578\\
3.7412	-0.0448721\\
3.7437	-0.0446963\\
3.7463	-0.0445163\\
3.7489	-0.0443392\\
3.7514	-0.0441719\\
3.754	-0.0440009\\
3.7566	-0.043833\\
3.7591	-0.0436746\\
3.7617	-0.0435129\\
3.7643	-0.0433546\\
3.7668	-0.0432054\\
3.7694	-0.0430536\\
3.7719	-0.0429108\\
3.7745	-0.0427656\\
3.7771	-0.0426239\\
3.7796	-0.042491\\
3.7822	-0.0423562\\
3.7848	-0.042225\\
3.7873	-0.0421022\\
3.7899	-0.041978\\
3.7925	-0.0418575\\
3.795	-0.0417452\\
3.7976	-0.041632\\
3.8002	-0.0415225\\
3.8027	-0.0414208\\
3.8053	-0.0413187\\
3.8078	-0.0412242\\
3.8104	-0.0411296\\
3.813	-0.0410389\\
3.8155	-0.0409553\\
3.8181	-0.0408722\\
3.8207	-0.040793\\
3.8232	-0.0407206\\
3.8258	-0.0406491\\
3.8284	-0.0405815\\
3.8309	-0.0405203\\
3.8335	-0.0404605\\
3.836	-0.0404067\\
3.8386	-0.0403547\\
3.8412	-0.0403066\\
3.8437	-0.0402642\\
3.8463	-0.0402239\\
3.8489	-0.0401876\\
3.8514	-0.0401565\\
3.854	-0.040128\\
3.8566	-0.0401035\\
3.8591	-0.0400837\\
3.8617	-0.040067\\
3.8642	-0.0400546\\
3.8668	-0.0400456\\
3.8694	-0.0400405\\
3.8719	-0.0400394\\
3.8745	-0.040042\\
3.8771	-0.0400485\\
3.8796	-0.0400583\\
3.8822	-0.0400724\\
3.8848	-0.0400903\\
3.8873	-0.0401111\\
3.8899	-0.0401365\\
3.8924	-0.0401645\\
3.895	-0.0401972\\
3.8976	-0.0402337\\
3.9001	-0.0402723\\
3.9027	-0.040316\\
3.9053	-0.0403634\\
3.9078	-0.0404123\\
3.9104	-0.0404668\\
3.913	-0.0405247\\
3.9155	-0.0405838\\
3.9181	-0.0406487\\
3.9206	-0.0407143\\
3.9232	-0.0407859\\
3.9258	-0.0408608\\
3.9283	-0.0409361\\
3.9309	-0.0410176\\
3.9335	-0.0411023\\
3.936	-0.0411868\\
3.9386	-0.0412779\\
3.9412	-0.041372\\
3.9437	-0.0414655\\
3.9463	-0.0415657\\
3.9488	-0.0416649\\
3.9514	-0.0417709\\
3.954	-0.0418799\\
3.9565	-0.0419873\\
3.9591	-0.0421019\\
3.9617	-0.0422192\\
3.9642	-0.0423345\\
3.9668	-0.0424571\\
3.9694	-0.0425823\\
3.9719	-0.0427051\\
3.9745	-0.0428352\\
3.9771	-0.0429679\\
3.9796	-0.0430976\\
3.9822	-0.0432349\\
3.9847	-0.0433691\\
3.9873	-0.0435108\\
3.9899	-0.0436547\\
3.9924	-0.0437951\\
3.995	-0.0439432\\
3.9976	-0.0440932\\
4.0001	-0.0442393\\
4.0027	-0.0443932\\
4.0053	-0.0445488\\
4.0078	-0.0447002\\
4.0104	-0.0448593\\
4.0129	-0.0450139\\
4.0155	-0.0451762\\
4.0181	-0.04534\\
4.0206	-0.045499\\
4.0232	-0.0456657\\
4.0258	-0.0458338\\
4.0283	-0.0459966\\
4.0309	-0.0461672\\
4.0335	-0.046339\\
4.036	-0.0465052\\
4.0386	-0.0466791\\
4.0411	-0.0468473\\
4.0437	-0.0470231\\
4.0463	-0.0471998\\
4.0488	-0.0473705\\
4.0514	-0.0475488\\
4.054	-0.0477278\\
4.0565	-0.0479004\\
4.0591	-0.0480806\\
4.0617	-0.0482613\\
4.0642	-0.0484355\\
4.0668	-0.048617\\
4.0693	-0.0487919\\
4.0719	-0.0489741\\
4.0745	-0.0491565\\
4.077	-0.049332\\
4.0796	-0.0495147\\
4.0822	-0.0496974\\
4.0847	-0.0498731\\
4.0873	-0.0500558\\
4.0899	-0.0502384\\
4.0924	-0.0504138\\
4.095	-0.0505959\\
4.0975	-0.0507708\\
4.1001	-0.0509524\\
4.1027	-0.0511335\\
4.1052	-0.0513073\\
4.1078	-0.0514876\\
4.1104	-0.0516673\\
4.1129	-0.0518395\\
4.1155	-0.052018\\
4.1181	-0.0521958\\
4.1206	-0.052366\\
4.1232	-0.0525422\\
4.1258	-0.0527176\\
4.1283	-0.0528854\\
4.1309	-0.053059\\
4.1334	-0.0532249\\
4.136	-0.0533965\\
4.1386	-0.053567\\
4.1411	-0.0537299\\
4.1437	-0.0538981\\
4.1463	-0.0540651\\
4.1488	-0.0542245\\
4.1514	-0.054389\\
4.154	-0.0545522\\
4.1565	-0.0547078\\
4.1591	-0.0548681\\
4.1616	-0.055021\\
4.1642	-0.0551785\\
4.1668	-0.0553344\\
4.1693	-0.0554829\\
4.1719	-0.0556358\\
4.1745	-0.0557869\\
4.177	-0.0559307\\
4.1796	-0.0560786\\
4.1822	-0.0562247\\
4.1847	-0.0563636\\
4.1873	-0.0565062\\
4.1898	-0.0566416\\
4.1924	-0.0567806\\
4.195	-0.0569177\\
4.1975	-0.0570477\\
4.2001	-0.0571811\\
4.2027	-0.0573124\\
4.2052	-0.0574369\\
4.2078	-0.0575643\\
4.2104	-0.0576897\\
4.2129	-0.0578084\\
4.2155	-0.0579297\\
4.218	-0.0580445\\
4.2206	-0.0581617\\
4.2232	-0.0582768\\
4.2257	-0.0583855\\
4.2283	-0.0584964\\
4.2309	-0.0586051\\
4.2334	-0.0587076\\
4.236	-0.058812\\
4.2386	-0.0589142\\
4.2411	-0.0590104\\
4.2437	-0.0591083\\
4.2462	-0.0592002\\
4.2488	-0.0592937\\
4.2514	-0.0593849\\
4.2539	-0.0594704\\
4.2565	-0.0595572\\
4.2591	-0.0596417\\
4.2616	-0.0597208\\
4.2642	-0.0598008\\
4.2668	-0.0598786\\
4.2693	-0.0599512\\
4.2719	-0.0600245\\
4.2744	-0.0600929\\
4.277	-0.0601617\\
4.2796	-0.0602283\\
4.2821	-0.0602902\\
4.2847	-0.0603524\\
4.2873	-0.0604123\\
4.2898	-0.0604678\\
4.2924	-0.0605233\\
4.295	-0.0605767\\
4.2975	-0.0606258\\
4.3001	-0.0606748\\
4.3027	-0.0607216\\
4.3052	-0.0607645\\
4.3078	-0.060807\\
4.3103	-0.0608458\\
4.3129	-0.0608841\\
4.3155	-0.0609203\\
4.318	-0.0609531\\
4.3206	-0.0609852\\
4.3232	-0.0610152\\
4.3257	-0.0610421\\
4.3283	-0.061068\\
4.3309	-0.061092\\
4.3334	-0.0611132\\
4.336	-0.0611332\\
4.3385	-0.0611507\\
4.3411	-0.061167\\
4.3437	-0.0611813\\
4.3462	-0.0611934\\
4.3488	-0.0612041\\
4.3514	-0.0612129\\
4.3539	-0.0612198\\
4.3565	-0.0612251\\
4.3591	-0.0612287\\
4.3616	-0.0612306\\
4.3642	-0.0612308\\
4.3667	-0.0612295\\
4.3693	-0.0612265\\
4.3719	-0.0612219\\
4.3744	-0.061216\\
4.377	-0.0612084\\
4.3796	-0.0611992\\
4.3821	-0.061189\\
4.3847	-0.061177\\
4.3873	-0.0611636\\
4.3898	-0.0611493\\
4.3924	-0.0611332\\
4.3949	-0.0611165\\
4.3975	-0.0610978\\
4.4001	-0.0610779\\
4.4026	-0.0610576\\
4.4052	-0.0610353\\
4.4078	-0.0610119\\
4.4103	-0.0609884\\
4.4129	-0.0609628\\
4.4155	-0.0609363\\
4.418	-0.0609098\\
4.4206	-0.0608813\\
4.4231	-0.0608531\\
4.4257	-0.0608228\\
4.4283	-0.0607917\\
4.4308	-0.0607611\\
4.4334	-0.0607284\\
4.436	-0.060695\\
4.4385	-0.0606623\\
4.4411	-0.0606276\\
4.4437	-0.0605922\\
4.4462	-0.0605577\\
4.4488	-0.0605212\\
4.4514	-0.0604843\\
4.4539	-0.0604483\\
4.4565	-0.0604105\\
4.459	-0.0603737\\
4.4616	-0.0603351\\
4.4642	-0.0602962\\
4.4667	-0.0602586\\
4.4693	-0.0602191\\
4.4719	-0.0601795\\
4.4744	-0.0601413\\
4.477	-0.0601013\\
4.4796	-0.0600613\\
4.4821	-0.0600228\\
4.4847	-0.0599826\\
4.4872	-0.0599441\\
4.4898	-0.059904\\
4.4924	-0.0598641\\
4.4949	-0.0598257\\
4.4975	-0.059786\\
4.5001	-0.0597465\\
4.5026	-0.0597087\\
4.5052	-0.0596696\\
4.5078	-0.0596308\\
4.5103	-0.0595938\\
4.5129	-0.0595557\\
4.5154	-0.0595194\\
4.518	-0.059482\\
4.5206	-0.0594451\\
4.5231	-0.0594101\\
4.5257	-0.0593741\\
4.5283	-0.0593387\\
4.5308	-0.0593052\\
4.5334	-0.0592709\\
4.536	-0.0592373\\
4.5385	-0.0592055\\
4.5411	-0.0591731\\
4.5436	-0.0591427\\
4.5462	-0.0591117\\
4.5488	-0.0590815\\
4.5513	-0.0590532\\
4.5539	-0.0590246\\
4.5565	-0.0589967\\
4.559	-0.0589708\\
4.5616	-0.0589446\\
4.5642	-0.0589194\\
4.5667	-0.0588959\\
4.5693	-0.0588725\\
4.5718	-0.0588508\\
4.5744	-0.0588293\\
4.577	-0.0588087\\
4.5795	-0.0587899\\
4.5821	-0.0587713\\
4.5847	-0.0587537\\
4.5872	-0.0587379\\
4.5898	-0.0587224\\
4.5924	-0.058708\\
4.5949	-0.0586952\\
4.5975	-0.058683\\
4.6001	-0.0586719\\
4.6026	-0.0586623\\
4.6052	-0.0586535\\
4.6077	-0.058646\\
4.6103	-0.0586394\\
4.6129	-0.058634\\
4.6154	-0.0586299\\
4.618	-0.0586268\\
4.6206	-0.0586249\\
4.6231	-0.0586241\\
4.6257	-0.0586245\\
4.6283	-0.0586262\\
4.6308	-0.0586289\\
4.6334	-0.0586328\\
4.6359	-0.0586378\\
4.6385	-0.0586442\\
4.6411	-0.0586517\\
4.6436	-0.0586602\\
4.6462	-0.0586701\\
4.6488	-0.0586812\\
4.6513	-0.0586931\\
4.6539	-0.0587065\\
4.6565	-0.0587212\\
4.659	-0.0587365\\
4.6616	-0.0587535\\
4.6641	-0.0587709\\
4.6667	-0.0587902\\
4.6693	-0.0588107\\
4.6718	-0.0588315\\
4.6744	-0.0588542\\
4.677	-0.058878\\
4.6795	-0.058902\\
4.6821	-0.058928\\
4.6847	-0.0589552\\
4.6872	-0.0589823\\
4.6898	-0.0590115\\
4.6923	-0.0590407\\
4.6949	-0.059072\\
4.6975	-0.0591043\\
4.7	-0.0591364\\
4.7026	-0.0591707\\
4.7052	-0.059206\\
4.7077	-0.0592409\\
4.7103	-0.0592781\\
4.7129	-0.0593162\\
4.7154	-0.0593537\\
4.718	-0.0593935\\
4.7205	-0.0594327\\
4.7231	-0.0594742\\
4.7257	-0.0595166\\
4.7282	-0.0595581\\
4.7308	-0.0596021\\
4.7334	-0.0596468\\
4.7359	-0.0596905\\
4.7385	-0.0597366\\
4.7411	-0.0597834\\
4.7436	-0.0598291\\
4.7462	-0.0598772\\
4.7487	-0.059924\\
4.7513	-0.0599733\\
4.7539	-0.0600232\\
4.7564	-0.0600717\\
4.759	-0.0601226\\
4.7616	-0.060174\\
4.7641	-0.0602239\\
4.7667	-0.0602762\\
4.7693	-0.0603289\\
4.7718	-0.0603799\\
4.7744	-0.0604334\\
4.777	-0.0604871\\
4.7795	-0.0605391\\
4.7821	-0.0605934\\
4.7846	-0.0606458\\
4.7872	-0.0607006\\
4.7898	-0.0607555\\
4.7923	-0.0608085\\
4.7949	-0.0608637\\
4.7975	-0.060919\\
4.8	-0.0609723\\
4.8026	-0.0610277\\
4.8052	-0.0610831\\
4.8077	-0.0611364\\
4.8103	-0.0611917\\
4.8128	-0.0612449\\
4.8154	-0.0613\\
4.818	-0.061355\\
4.8205	-0.0614078\\
4.8231	-0.0614624\\
4.8257	-0.0615169\\
4.8282	-0.0615689\\
4.8308	-0.0616228\\
4.8334	-0.0616764\\
4.8359	-0.0617276\\
4.8385	-0.0617805\\
4.841	-0.0618309\\
4.8436	-0.0618829\\
4.8462	-0.0619345\\
4.8487	-0.0619836\\
4.8513	-0.0620342\\
4.8539	-0.0620842\\
4.8564	-0.0621318\\
4.859	-0.0621807\\
4.8616	-0.0622289\\
4.8641	-0.0622747\\
4.8667	-0.0623216\\
4.8692	-0.062366\\
4.8718	-0.0624115\\
4.8744	-0.0624563\\
4.8769	-0.0624985\\
4.8795	-0.0625417\\
4.8821	-0.0625839\\
4.8846	-0.0626238\\
4.8872	-0.0626644\\
4.8898	-0.062704\\
4.8923	-0.0627412\\
4.8949	-0.062779\\
4.8974	-0.0628144\\
4.9	-0.0628502\\
4.9026	-0.0628849\\
4.9051	-0.0629174\\
4.9077	-0.06295\\
4.9103	-0.0629816\\
4.9128	-0.0630109\\
4.9154	-0.0630402\\
4.918	-0.0630684\\
4.9205	-0.0630944\\
4.9231	-0.0631202\\
4.9257	-0.0631448\\
4.9282	-0.0631673\\
4.9308	-0.0631895\\
4.9333	-0.0632096\\
4.9359	-0.0632292\\
4.9385	-0.0632475\\
4.941	-0.0632639\\
4.9436	-0.0632796\\
4.9462	-0.063294\\
4.9487	-0.0633065\\
4.9513	-0.0633181\\
4.9539	-0.0633284\\
4.9564	-0.063337\\
4.959	-0.0633444\\
4.9615	-0.0633503\\
4.9641	-0.063355\\
4.9667	-0.0633582\\
4.9692	-0.0633599\\
4.9718	-0.0633602\\
4.9744	-0.0633591\\
4.9769	-0.0633565\\
4.9795	-0.0633525\\
4.9821	-0.0633469\\
4.9846	-0.0633401\\
4.9872	-0.0633315\\
4.9897	-0.0633219\\
4.9923	-0.0633103\\
4.9949	-0.0632972\\
4.9974	-0.0632832\\
5	-0.0632671\\
};
\addplot [color=mycolor3, line width=1.0pt, forget plot]
  table[row sep=crcr]{%
-0.125	5.46548e-08\\
-0.1224	5.58085e-08\\
-0.1199	5.69182e-08\\
-0.1173	5.80725e-08\\
-0.1147	5.92271e-08\\
-0.1122	6.03376e-08\\
-0.1096	6.14927e-08\\
-0.1071	6.26037e-08\\
-0.1045	6.37593e-08\\
-0.1019	6.49152e-08\\
-0.0994	6.60269e-08\\
-0.0968	6.71832e-08\\
-0.0942	6.83398e-08\\
-0.0917	6.94521e-08\\
-0.0891	7.06092e-08\\
-0.0865	7.17664e-08\\
-0.084	7.28793e-08\\
-0.0814	7.40369e-08\\
-0.0789	7.51501e-08\\
-0.0763	7.63081e-08\\
-0.0737	7.74662e-08\\
-0.0712	7.85799e-08\\
-0.0686	7.97383e-08\\
-0.066	8.08967e-08\\
-0.0635	8.20109e-08\\
-0.0609	8.31697e-08\\
-0.0583	8.43286e-08\\
-0.0558	8.54429e-08\\
-0.0532	8.66021e-08\\
-0.0507	8.77167e-08\\
-0.0481	8.88758e-08\\
-0.0455	9.00351e-08\\
-0.043	9.11499e-08\\
-0.0404	9.23093e-08\\
-0.0378	9.34687e-08\\
-0.0353	9.45835e-08\\
-0.0327	9.5743e-08\\
-0.0301	9.69026e-08\\
-0.0276	9.80175e-08\\
-0.025	9.91771e-08\\
-0.0224	1.00337e-07\\
-0.0199	1.01452e-07\\
-0.0173	1.02611e-07\\
-0.0148	1.03726e-07\\
-0.0122	1.04885e-07\\
-0.0096	1.06045e-07\\
-0.0071	1.07159e-07\\
-0.0045	1.08318e-07\\
-0.0019	1.09477e-07\\
0.0006	1.10592e-07\\
0.0032	0.000702492\\
0.0058	0.00314225\\
0.0083	0.00589831\\
0.0109	0.00912722\\
0.0134	0.0124159\\
0.016	0.0158081\\
0.0186	0.0190433\\
0.0211	0.0220057\\
0.0237	0.0250005\\
0.0263	0.027964\\
0.0288	0.0307865\\
0.0314	0.0336569\\
0.034	0.0364234\\
0.0365	0.038977\\
0.0391	0.0415407\\
0.0416	0.0439416\\
0.0442	0.0463839\\
0.0468	0.0487643\\
0.0493	0.0509813\\
0.0519	0.0532041\\
0.0545	0.0553471\\
0.057	0.0573446\\
0.0596	0.0593674\\
0.0622	0.0613362\\
0.0647	0.0631729\\
0.0673	0.0650182\\
0.0698	0.06673\\
0.0724	0.0684528\\
0.075	0.0701265\\
0.0775	0.0716945\\
0.0801	0.0732802\\
0.0827	0.0748145\\
0.0852	0.0762388\\
0.0878	0.0776705\\
0.0904	0.0790594\\
0.0929	0.0803607\\
0.0955	0.0816791\\
0.098	0.0829094\\
0.1006	0.0841455\\
0.1032	0.0853371\\
0.1057	0.0864459\\
0.1083	0.087567\\
0.1109	0.0886584\\
0.1134	0.0896776\\
0.116	0.090701\\
0.1186	0.0916843\\
0.1211	0.0925939\\
0.1237	0.0935081\\
0.1263	0.0943942\\
0.1288	0.0952195\\
0.1314	0.0960456\\
0.1339	0.0968047\\
0.1365	0.0975568\\
0.1391	0.0982746\\
0.1416	0.0989368\\
0.1442	0.0995983\\
0.1468	0.100229\\
0.1493	0.100802\\
0.1519	0.101361\\
0.1545	0.101883\\
0.157	0.102354\\
0.1596	0.102815\\
0.1621	0.10323\\
0.1647	0.103628\\
0.1673	0.103989\\
0.1698	0.104299\\
0.1724	0.104587\\
0.175	0.104843\\
0.1775	0.105061\\
0.1801	0.105256\\
0.1827	0.105415\\
0.1852	0.105531\\
0.1878	0.105615\\
0.1903	0.105664\\
0.1929	0.105685\\
0.1955	0.105676\\
0.198	0.105636\\
0.2006	0.105558\\
0.2032	0.105442\\
0.2057	0.1053\\
0.2083	0.105122\\
0.2109	0.104915\\
0.2134	0.104689\\
0.216	0.104421\\
0.2185	0.104131\\
0.2211	0.103798\\
0.2237	0.103436\\
0.2262	0.103064\\
0.2288	0.102652\\
0.2314	0.102212\\
0.2339	0.101762\\
0.2365	0.101265\\
0.2391	0.100742\\
0.2416	0.100218\\
0.2442	0.0996533\\
0.2467	0.0990903\\
0.2493	0.098482\\
0.2519	0.0978496\\
0.2544	0.0972207\\
0.257	0.0965487\\
0.2596	0.0958613\\
0.2621	0.0951864\\
0.2647	0.0944686\\
0.2673	0.0937329\\
0.2698	0.0930094\\
0.2724	0.0922432\\
0.2749	0.0914964\\
0.2775	0.090711\\
0.2801	0.0899162\\
0.2826	0.0891416\\
0.2852	0.0883247\\
0.2878	0.0874981\\
0.2903	0.0866973\\
0.2929	0.0858607\\
0.2955	0.0850208\\
0.298	0.0842089\\
0.3006	0.0833588\\
0.3032	0.0825035\\
0.3057	0.0816784\\
0.3083	0.0808205\\
0.3108	0.0799971\\
0.3134	0.0791417\\
0.316	0.0782858\\
0.3185	0.0774617\\
0.3211	0.0766051\\
0.3237	0.0757514\\
0.3262	0.0749352\\
0.3288	0.0740915\\
0.3314	0.0732517\\
0.3339	0.0724467\\
0.3365	0.0716123\\
0.339	0.0708146\\
0.3416	0.0699918\\
0.3442	0.069177\\
0.3467	0.0684004\\
0.3493	0.0675983\\
0.3519	0.0668012\\
0.3544	0.0660404\\
0.357	0.0652569\\
0.3596	0.0644828\\
0.3621	0.0637473\\
0.3647	0.0629901\\
0.3672	0.0622683\\
0.3698	0.0615242\\
0.3724	0.060788\\
0.3749	0.0600892\\
0.3775	0.0593725\\
0.3801	0.0586651\\
0.3826	0.0579923\\
0.3852	0.0572994\\
0.3878	0.0566142\\
0.3903	0.0559638\\
0.3929	0.0552974\\
0.3954	0.0546659\\
0.398	0.0540174\\
0.4006	0.0533762\\
0.4031	0.0527663\\
0.4057	0.0521398\\
0.4083	0.0515226\\
0.4108	0.0509382\\
0.4134	0.050339\\
0.416	0.049747\\
0.4185	0.0491839\\
0.4211	0.0486049\\
0.4236	0.0480558\\
0.4262	0.0474932\\
0.4288	0.046939\\
0.4313	0.046413\\
0.4339	0.0458717\\
0.4365	0.0453361\\
0.439	0.0448273\\
0.4416	0.0443054\\
0.4442	0.0437912\\
0.4467	0.0433033\\
0.4493	0.0428012\\
0.4519	0.0423039\\
0.4544	0.0418305\\
0.457	0.0413441\\
0.4595	0.0408825\\
0.4621	0.0404086\\
0.4647	0.0399398\\
0.4672	0.0394929\\
0.4698	0.0390319\\
0.4724	0.0385752\\
0.4749	0.0381411\\
0.4775	0.0376948\\
0.4801	0.037253\\
0.4826	0.0368316\\
0.4852	0.036396\\
0.4877	0.0359799\\
0.4903	0.0355509\\
0.4929	0.035126\\
0.4954	0.0347211\\
0.498	0.0343029\\
0.5006	0.0338866\\
0.5031	0.033488\\
0.5057	0.0330755\\
0.5083	0.0326659\\
0.5108	0.0322747\\
0.5134	0.0318702\\
0.5159	0.0314823\\
0.5185	0.0310796\\
0.5211	0.0306777\\
0.5236	0.0302924\\
0.5262	0.0298935\\
0.5288	0.0294959\\
0.5313	0.0291142\\
0.5339	0.028717\\
0.5365	0.0283192\\
0.539	0.0279367\\
0.5416	0.0275391\\
0.5441	0.0271573\\
0.5467	0.0267599\\
0.5493	0.0263615\\
0.5518	0.0259769\\
0.5544	0.0255754\\
0.557	0.0251729\\
0.5595	0.024785\\
0.5621	0.0243806\\
0.5647	0.0239743\\
0.5672	0.0235812\\
0.5698	0.0231697\\
0.5723	0.0227716\\
0.5749	0.0223556\\
0.5775	0.0219374\\
0.58	0.0215329\\
0.5826	0.0211089\\
0.5852	0.0206811\\
0.5877	0.0202661\\
0.5903	0.0198312\\
0.5929	0.0193931\\
0.5954	0.0189686\\
0.598	0.0185232\\
0.6006	0.0180731\\
0.6031	0.0176357\\
0.6057	0.0171763\\
0.6082	0.0167305\\
0.6108	0.0162628\\
0.6134	0.0157904\\
0.6159	0.0153311\\
0.6185	0.014848\\
0.6211	0.0143593\\
0.6236	0.0138844\\
0.6262	0.0133856\\
0.6288	0.0128816\\
0.6313	0.0123916\\
0.6339	0.011876\\
0.6364	0.0113743\\
0.639	0.0108467\\
0.6416	0.0103134\\
0.6441	0.00979527\\
0.6467	0.0092507\\
0.6493	0.0086998\\
0.6518	0.00816388\\
0.6544	0.00760018\\
0.657	0.00703038\\
0.6595	0.00647699\\
0.6621	0.00589568\\
0.6646	0.00533086\\
0.6672	0.00473705\\
0.6698	0.00413669\\
0.6723	0.00355353\\
0.6749	0.00294129\\
0.6775	0.00232335\\
0.68	0.00172367\\
0.6826	0.00109399\\
0.6852	0.000458084\\
0.6877	-0.000159045\\
0.6903	-0.000806371\\
0.6928	-0.00143379\\
0.6954	-0.00209142\\
0.698	-0.00275444\\
0.7005	-0.00339718\\
0.7031	-0.00407094\\
0.7057	-0.00474977\\
0.7082	-0.00540686\\
0.7108	-0.00609452\\
0.7134	-0.00678655\\
0.7159	-0.0074562\\
0.7185	-0.008157\\
0.721	-0.00883478\\
0.7236	-0.00954332\\
0.7262	-0.0102552\\
0.7287	-0.0109426\\
0.7313	-0.0116606\\
0.7339	-0.0123818\\
0.7364	-0.013078\\
0.739	-0.0138046\\
0.7416	-0.0145332\\
0.7441	-0.0152354\\
0.7467	-0.0159673\\
0.7492	-0.0166725\\
0.7518	-0.0174073\\
0.7544	-0.0181433\\
0.7569	-0.0188515\\
0.7595	-0.0195883\\
0.7621	-0.0203249\\
0.7646	-0.0210331\\
0.7672	-0.0217694\\
0.7698	-0.0225051\\
0.7723	-0.0232115\\
0.7749	-0.0239449\\
0.7775	-0.0246763\\
0.78	-0.0253777\\
0.7826	-0.0261051\\
0.7851	-0.0268024\\
0.7877	-0.0275249\\
0.7903	-0.0282442\\
0.7928	-0.0289325\\
0.7954	-0.0296444\\
0.798	-0.0303524\\
0.8005	-0.0310291\\
0.8031	-0.0317286\\
0.8057	-0.0324232\\
0.8082	-0.033086\\
0.8108	-0.0337697\\
0.8133	-0.0344216\\
0.8159	-0.0350936\\
0.8185	-0.0357594\\
0.821	-0.0363933\\
0.8236	-0.0370458\\
0.8262	-0.0376908\\
0.8287	-0.0383038\\
0.8313	-0.0389337\\
0.8339	-0.0395556\\
0.8364	-0.0401459\\
0.839	-0.0407513\\
0.8415	-0.0413251\\
0.8441	-0.0419127\\
0.8467	-0.0424909\\
0.8492	-0.0430378\\
0.8518	-0.0435972\\
0.8544	-0.0441466\\
0.8569	-0.0446651\\
0.8595	-0.045194\\
0.8621	-0.045712\\
0.8646	-0.0461999\\
0.8672	-0.0466964\\
0.8697	-0.0471634\\
0.8723	-0.0476378\\
0.8749	-0.0481004\\
0.8774	-0.0485341\\
0.88	-0.0489732\\
0.8826	-0.0494003\\
0.8851	-0.0497995\\
0.8877	-0.0502025\\
0.8903	-0.050593\\
0.8928	-0.0509564\\
0.8954	-0.0513218\\
0.8979	-0.0516609\\
0.9005	-0.052001\\
0.9031	-0.0523281\\
0.9056	-0.0526302\\
0.9082	-0.0529314\\
0.9108	-0.0532192\\
0.9133	-0.0534833\\
0.9159	-0.0537449\\
0.9185	-0.0539932\\
0.921	-0.0542193\\
0.9236	-0.0544412\\
0.9262	-0.0546496\\
0.9287	-0.0548371\\
0.9313	-0.0550189\\
0.9338	-0.0551811\\
0.9364	-0.0553367\\
0.939	-0.0554789\\
0.9415	-0.0556029\\
0.9441	-0.0557188\\
0.9467	-0.0558213\\
0.9492	-0.0559075\\
0.9518	-0.0559844\\
0.9544	-0.0560482\\
0.9569	-0.0560974\\
0.9595	-0.0561358\\
0.962	-0.0561607\\
0.9646	-0.0561741\\
0.9672	-0.0561751\\
0.9697	-0.0561644\\
0.9723	-0.0561414\\
0.9749	-0.0561061\\
0.9774	-0.0560609\\
0.98	-0.0560022\\
0.9826	-0.055932\\
0.9851	-0.0558537\\
0.9877	-0.0557612\\
0.9902	-0.0556618\\
0.9928	-0.0555477\\
0.9954	-0.0554229\\
0.9979	-0.0552929\\
1.0005	-0.0551477\\
1.0031	-0.0549925\\
1.0056	-0.054834\\
1.0082	-0.0546597\\
1.0108	-0.054476\\
1.0133	-0.0542907\\
1.0159	-0.0540892\\
1.0184	-0.0538874\\
1.021	-0.0536694\\
1.0236	-0.0534432\\
1.0261	-0.0532182\\
1.0287	-0.0529767\\
1.0313	-0.0527279\\
1.0338	-0.0524821\\
1.0364	-0.0522198\\
1.039	-0.051951\\
1.0415	-0.0516867\\
1.0441	-0.0514059\\
1.0466	-0.0511306\\
1.0492	-0.050839\\
1.0518	-0.0505424\\
1.0543	-0.0502528\\
1.0569	-0.0499472\\
1.0595	-0.0496374\\
1.062	-0.0493359\\
1.0646	-0.0490187\\
1.0672	-0.0486984\\
1.0697	-0.0483876\\
1.0723	-0.0480618\\
1.0748	-0.0477463\\
1.0774	-0.0474161\\
1.08	-0.0470841\\
1.0825	-0.0467636\\
1.0851	-0.0464292\\
1.0877	-0.0460939\\
1.0902	-0.045771\\
1.0928	-0.045435\\
1.0954	-0.045099\\
1.0979	-0.0447762\\
1.1005	-0.0444412\\
1.1031	-0.0441071\\
1.1056	-0.0437871\\
1.1082	-0.0434558\\
1.1107	-0.0431389\\
1.1133	-0.0428113\\
1.1159	-0.0424861\\
1.1184	-0.042176\\
1.121	-0.0418563\\
1.1236	-0.0415398\\
1.1261	-0.0412387\\
1.1287	-0.0409292\\
1.1313	-0.0406237\\
1.1338	-0.040334\\
1.1364	-0.0400371\\
1.1389	-0.0397561\\
1.1415	-0.0394688\\
1.1441	-0.0391866\\
1.1466	-0.0389205\\
1.1492	-0.0386492\\
1.1518	-0.0383839\\
1.1543	-0.0381347\\
1.1569	-0.0378817\\
1.1595	-0.0376352\\
1.162	-0.0374046\\
1.1646	-0.0371716\\
1.1671	-0.0369543\\
1.1697	-0.0367356\\
1.1723	-0.0365244\\
1.1748	-0.0363287\\
1.1774	-0.0361328\\
1.18	-0.035945\\
1.1825	-0.0357721\\
1.1851	-0.0356006\\
1.1877	-0.0354377\\
1.1902	-0.0352892\\
1.1928	-0.0351434\\
1.1953	-0.0350116\\
1.1979	-0.0348835\\
1.2005	-0.0347646\\
1.203	-0.0346591\\
1.2056	-0.0345586\\
1.2082	-0.0344676\\
1.2107	-0.0343892\\
1.2133	-0.0343172\\
1.2159	-0.034255\\
1.2184	-0.0342045\\
1.221	-0.0341618\\
1.2235	-0.0341302\\
1.2261	-0.0341071\\
1.2287	-0.0340943\\
1.2312	-0.0340915\\
1.2338	-0.0340986\\
1.2364	-0.0341161\\
1.2389	-0.0341426\\
1.2415	-0.0341802\\
1.2441	-0.0342282\\
1.2466	-0.0342841\\
1.2492	-0.0343524\\
1.2518	-0.034431\\
1.2543	-0.0345164\\
1.2569	-0.0346154\\
1.2594	-0.0347203\\
1.262	-0.0348394\\
1.2646	-0.0349688\\
1.2671	-0.0351029\\
1.2697	-0.0352524\\
1.2723	-0.035412\\
1.2748	-0.035575\\
1.2774	-0.0357543\\
1.28	-0.0359436\\
1.2825	-0.0361349\\
1.2851	-0.0363436\\
1.2876	-0.0365534\\
1.2902	-0.036781\\
1.2928	-0.0370183\\
1.2953	-0.0372553\\
1.2979	-0.037511\\
1.3005	-0.0377759\\
1.303	-0.0380393\\
1.3056	-0.0383221\\
1.3082	-0.0386138\\
1.3107	-0.0389026\\
1.3133	-0.0392114\\
1.3158	-0.0395163\\
1.3184	-0.0398417\\
1.321	-0.0401753\\
1.3235	-0.0405037\\
1.3261	-0.0408531\\
1.3287	-0.0412101\\
1.3312	-0.0415607\\
1.3338	-0.0419325\\
1.3364	-0.0423116\\
1.3389	-0.0426829\\
1.3415	-0.0430758\\
1.344	-0.0434599\\
1.3466	-0.0438658\\
1.3492	-0.0442781\\
1.3517	-0.0446804\\
1.3543	-0.0451046\\
1.3569	-0.0455347\\
1.3594	-0.0459534\\
1.362	-0.0463942\\
1.3646	-0.0468401\\
1.3671	-0.0472736\\
1.3697	-0.0477291\\
1.3722	-0.0481714\\
1.3748	-0.0486356\\
1.3774	-0.0491039\\
1.3799	-0.0495578\\
1.3825	-0.0500335\\
1.3851	-0.0505127\\
1.3876	-0.0509764\\
1.3902	-0.0514616\\
1.3928	-0.0519496\\
1.3953	-0.0524212\\
1.3979	-0.0529139\\
1.4005	-0.0534086\\
1.403	-0.0538861\\
1.4056	-0.0543842\\
1.4081	-0.0548644\\
1.4107	-0.0553649\\
1.4133	-0.0558663\\
1.4158	-0.056349\\
1.4184	-0.0568514\\
1.421	-0.057354\\
1.4235	-0.0578372\\
1.4261	-0.0583394\\
1.4287	-0.058841\\
1.4312	-0.0593226\\
1.4338	-0.0598225\\
1.4363	-0.0603019\\
1.4389	-0.0607991\\
1.4415	-0.0612945\\
1.444	-0.061769\\
1.4466	-0.0622603\\
1.4492	-0.0627491\\
1.4517	-0.0632167\\
1.4543	-0.0637001\\
1.4569	-0.0641804\\
1.4594	-0.064639\\
1.462	-0.0651124\\
1.4645	-0.0655641\\
1.4671	-0.0660298\\
1.4697	-0.0664913\\
1.4722	-0.0669309\\
1.4748	-0.0673834\\
1.4774	-0.067831\\
1.4799	-0.0682566\\
1.4825	-0.068694\\
1.4851	-0.0691258\\
1.4876	-0.0695356\\
1.4902	-0.069956\\
1.4927	-0.0703545\\
1.4953	-0.0707626\\
1.4979	-0.0711643\\
1.5004	-0.0715441\\
1.503	-0.0719324\\
1.5056	-0.0723136\\
1.5081	-0.0726733\\
1.5107	-0.07304\\
1.5133	-0.0733992\\
1.5158	-0.0737372\\
1.5184	-0.0740809\\
1.5209	-0.0744037\\
1.5235	-0.0747314\\
1.5261	-0.0750506\\
1.5286	-0.0753494\\
1.5312	-0.0756516\\
1.5338	-0.075945\\
1.5363	-0.0762186\\
1.5389	-0.0764942\\
1.5415	-0.0767605\\
1.544	-0.0770077\\
1.5466	-0.0772554\\
1.5491	-0.0774845\\
1.5517	-0.0777133\\
1.5543	-0.0779322\\
1.5568	-0.0781333\\
1.5594	-0.0783325\\
1.562	-0.0785217\\
1.5645	-0.0786938\\
1.5671	-0.0788628\\
1.5697	-0.0790214\\
1.5722	-0.079164\\
1.5748	-0.079302\\
1.5774	-0.0794293\\
1.5799	-0.0795417\\
1.5825	-0.0796482\\
1.585	-0.0797403\\
1.5876	-0.0798256\\
1.5902	-0.0799\\
1.5927	-0.0799613\\
1.5953	-0.0800144\\
1.5979	-0.0800565\\
1.6004	-0.0800868\\
1.603	-0.0801074\\
1.6056	-0.0801172\\
1.6081	-0.0801162\\
1.6107	-0.0801045\\
1.6132	-0.0800829\\
1.6158	-0.0800497\\
1.6184	-0.0800056\\
1.6209	-0.0799529\\
1.6235	-0.0798876\\
1.6261	-0.0798114\\
1.6286	-0.079728\\
1.6312	-0.0796307\\
1.6338	-0.0795228\\
1.6363	-0.079409\\
1.6389	-0.0792804\\
1.6414	-0.0791468\\
1.644	-0.0789977\\
1.6466	-0.0788383\\
1.6491	-0.0786753\\
1.6517	-0.0784959\\
1.6543	-0.0783064\\
1.6568	-0.0781149\\
1.6594	-0.077906\\
1.662	-0.0776874\\
1.6645	-0.0774681\\
1.6671	-0.0772307\\
1.6696	-0.0769936\\
1.6722	-0.076738\\
1.6748	-0.0764732\\
1.6773	-0.0762102\\
1.6799	-0.0759279\\
1.6825	-0.075637\\
1.685	-0.0753492\\
1.6876	-0.0750416\\
1.6902	-0.0747259\\
1.6927	-0.0744146\\
1.6953	-0.0740832\\
1.6978	-0.0737573\\
1.7004	-0.0734108\\
1.703	-0.073057\\
1.7055	-0.0727101\\
1.7081	-0.0723424\\
1.7107	-0.0719678\\
1.7132	-0.0716014\\
1.7158	-0.071214\\
1.7184	-0.0708204\\
1.7209	-0.0704363\\
1.7235	-0.070031\\
1.7261	-0.0696202\\
1.7286	-0.06922\\
1.7312	-0.0687987\\
1.7337	-0.0683889\\
1.7363	-0.067958\\
1.7389	-0.0675226\\
1.7414	-0.0670998\\
1.744	-0.0666561\\
1.7466	-0.0662085\\
1.7491	-0.0657746\\
1.7517	-0.06532\\
1.7543	-0.0648622\\
1.7568	-0.0644191\\
1.7594	-0.0639557\\
1.7619	-0.0635077\\
1.7645	-0.0630395\\
1.7671	-0.0625692\\
1.7696	-0.0621154\\
1.7722	-0.0616418\\
1.7748	-0.0611668\\
1.7773	-0.0607091\\
1.7799	-0.0602321\\
1.7825	-0.0597545\\
1.785	-0.0592949\\
1.7876	-0.0588167\\
1.7901	-0.0583569\\
1.7927	-0.0578791\\
1.7953	-0.0574017\\
1.7978	-0.0569434\\
1.8004	-0.0564677\\
1.803	-0.0559932\\
1.8055	-0.0555383\\
1.8081	-0.0550669\\
1.8107	-0.0545973\\
1.8132	-0.0541477\\
1.8158	-0.0536825\\
1.8183	-0.0532376\\
1.8209	-0.0527776\\
1.8235	-0.0523206\\
1.826	-0.0518841\\
1.8286	-0.0514336\\
1.8312	-0.0509867\\
1.8337	-0.0505606\\
1.8363	-0.0501214\\
1.8389	-0.0496865\\
1.8414	-0.0492725\\
1.844	-0.0488465\\
1.8465	-0.0484414\\
1.8491	-0.0480251\\
1.8517	-0.047614\\
1.8542	-0.0472238\\
1.8568	-0.0468235\\
1.8594	-0.0464289\\
1.8619	-0.0460552\\
1.8645	-0.0456725\\
1.8671	-0.0452961\\
1.8696	-0.0449403\\
1.8722	-0.0445768\\
1.8747	-0.0442336\\
1.8773	-0.0438835\\
1.8799	-0.0435406\\
1.8824	-0.0432176\\
1.885	-0.042889\\
1.8876	-0.042568\\
1.8901	-0.0422665\\
1.8927	-0.0419607\\
1.8953	-0.0416627\\
1.8978	-0.0413839\\
1.9004	-0.0411019\\
1.903	-0.0408283\\
1.9055	-0.0405731\\
1.9081	-0.0403161\\
1.9106	-0.0400771\\
1.9132	-0.0398372\\
1.9158	-0.039606\\
1.9183	-0.0393922\\
1.9209	-0.0391788\\
1.9235	-0.0389744\\
1.926	-0.0387866\\
1.9286	-0.0386003\\
1.9312	-0.0384234\\
1.9337	-0.0382622\\
1.9363	-0.0381038\\
1.9388	-0.0379606\\
1.9414	-0.037821\\
1.944	-0.037691\\
1.9465	-0.0375753\\
1.9491	-0.0374644\\
1.9517	-0.0373634\\
1.9542	-0.0372754\\
1.9568	-0.0371937\\
1.9594	-0.0371218\\
1.9619	-0.037062\\
1.9645	-0.0370096\\
1.967	-0.0369686\\
1.9696	-0.0369357\\
1.9722	-0.0369128\\
1.9747	-0.0369001\\
1.9773	-0.0368967\\
1.9799	-0.0369033\\
1.9824	-0.036919\\
1.985	-0.0369451\\
1.9876	-0.0369811\\
1.9901	-0.037025\\
1.9927	-0.0370804\\
1.9952	-0.037143\\
1.9978	-0.0372177\\
2.0004	-0.0373022\\
2.0029	-0.0373926\\
2.0055	-0.0374962\\
2.0081	-0.0376094\\
2.0106	-0.0377273\\
2.0132	-0.0378593\\
2.0158	-0.0380008\\
2.0183	-0.0381457\\
2.0209	-0.0383057\\
2.0234	-0.0384682\\
2.026	-0.0386463\\
2.0286	-0.0388336\\
2.0311	-0.0390222\\
2.0337	-0.0392271\\
2.0363	-0.039441\\
2.0388	-0.039655\\
2.0414	-0.039886\\
2.044	-0.0401257\\
2.0465	-0.0403642\\
2.0491	-0.0406205\\
2.0517	-0.0408852\\
2.0542	-0.0411474\\
2.0568	-0.041428\\
2.0593	-0.0417053\\
2.0619	-0.0420015\\
2.0645	-0.0423054\\
2.067	-0.0426048\\
2.0696	-0.0429235\\
2.0722	-0.0432496\\
2.0747	-0.04357\\
2.0773	-0.0439101\\
2.0799	-0.0442572\\
2.0824	-0.0445973\\
2.085	-0.0449576\\
2.0875	-0.0453102\\
2.0901	-0.0456831\\
2.0927	-0.0460623\\
2.0952	-0.0464326\\
2.0978	-0.0468235\\
2.1004	-0.0472201\\
2.1029	-0.0476067\\
2.1055	-0.0480142\\
2.1081	-0.0484268\\
2.1106	-0.0488284\\
2.1132	-0.0492509\\
2.1157	-0.0496616\\
2.1183	-0.0500931\\
2.1209	-0.0505291\\
2.1234	-0.0509523\\
2.126	-0.0513963\\
2.1286	-0.0518441\\
2.1311	-0.0522782\\
2.1337	-0.052733\\
2.1363	-0.0531911\\
2.1388	-0.0536345\\
2.1414	-0.0540984\\
2.1439	-0.054547\\
2.1465	-0.055016\\
2.1491	-0.0554873\\
2.1516	-0.0559425\\
2.1542	-0.0564177\\
2.1568	-0.0568946\\
2.1593	-0.0573546\\
2.1619	-0.0578342\\
2.1645	-0.0583149\\
2.167	-0.058778\\
2.1696	-0.0592602\\
2.1721	-0.0597243\\
2.1747	-0.0602073\\
2.1773	-0.0606902\\
2.1798	-0.0611545\\
2.1824	-0.061637\\
2.185	-0.0621189\\
2.1875	-0.0625815\\
2.1901	-0.0630616\\
2.1927	-0.0635405\\
2.1952	-0.0639997\\
2.1978	-0.0644756\\
2.2004	-0.0649497\\
2.2029	-0.0654037\\
2.2055	-0.0658735\\
2.208	-0.066323\\
2.2106	-0.0667879\\
2.2132	-0.0672499\\
2.2157	-0.0676912\\
2.2183	-0.0681469\\
2.2209	-0.0685992\\
2.2234	-0.0690306\\
2.226	-0.0694754\\
2.2286	-0.0699161\\
2.2311	-0.0703359\\
2.2337	-0.070768\\
2.2362	-0.0711791\\
2.2388	-0.0716019\\
2.2414	-0.0720196\\
2.2439	-0.0724163\\
2.2465	-0.0728236\\
2.2491	-0.0732253\\
2.2516	-0.073606\\
2.2542	-0.0739961\\
2.2568	-0.0743801\\
2.2593	-0.0747433\\
2.2619	-0.0751147\\
2.2644	-0.0754655\\
2.267	-0.0758236\\
2.2696	-0.0761748\\
2.2721	-0.0765057\\
2.2747	-0.0768427\\
2.2773	-0.0771723\\
2.2798	-0.077482\\
2.2824	-0.0777965\\
2.285	-0.0781031\\
2.2875	-0.0783904\\
2.2901	-0.0786811\\
2.2926	-0.0789528\\
2.2952	-0.0792272\\
2.2978	-0.079493\\
2.3003	-0.0797405\\
2.3029	-0.0799892\\
2.3055	-0.0802291\\
2.308	-0.0804513\\
2.3106	-0.0806735\\
2.3132	-0.0808865\\
2.3157	-0.0810826\\
2.3183	-0.0812773\\
2.3208	-0.0814557\\
2.3234	-0.0816318\\
2.326	-0.0817984\\
2.3285	-0.0819495\\
2.3311	-0.0820971\\
2.3337	-0.082235\\
2.3362	-0.0823582\\
2.3388	-0.0824768\\
2.3414	-0.0825854\\
2.3439	-0.0826805\\
2.3465	-0.0827695\\
2.349	-0.0828458\\
2.3516	-0.0829152\\
2.3542	-0.0829746\\
2.3567	-0.0830222\\
2.3593	-0.0830618\\
2.3619	-0.0830914\\
2.3644	-0.0831103\\
2.367	-0.0831202\\
2.3696	-0.08312\\
2.3721	-0.0831103\\
2.3747	-0.0830905\\
2.3773	-0.0830607\\
2.3798	-0.0830226\\
2.3824	-0.0829734\\
2.3849	-0.0829167\\
2.3875	-0.0828482\\
2.3901	-0.0827699\\
2.3926	-0.0826855\\
2.3952	-0.0825882\\
2.3978	-0.0824814\\
2.4003	-0.0823698\\
2.4029	-0.0822445\\
2.4055	-0.0821099\\
2.408	-0.0819717\\
2.4106	-0.0818191\\
2.4131	-0.0816637\\
2.4157	-0.0814935\\
2.4183	-0.0813143\\
2.4208	-0.0811339\\
2.4234	-0.0809377\\
2.426	-0.0807331\\
2.4285	-0.0805285\\
2.4311	-0.0803075\\
2.4337	-0.0800785\\
2.4362	-0.0798508\\
2.4388	-0.0796063\\
2.4413	-0.0793639\\
2.4439	-0.0791044\\
2.4465	-0.0788375\\
2.449	-0.0785741\\
2.4516	-0.0782932\\
2.4542	-0.0780053\\
2.4567	-0.0777222\\
2.4593	-0.0774212\\
2.4619	-0.0771138\\
2.4644	-0.0768124\\
2.467	-0.0764929\\
2.4695	-0.0761801\\
2.4721	-0.0758492\\
2.4747	-0.0755128\\
2.4772	-0.0751843\\
2.4798	-0.0748375\\
2.4824	-0.0744858\\
2.4849	-0.0741431\\
2.4875	-0.0737822\\
2.4901	-0.0734169\\
2.4926	-0.0730616\\
2.4952	-0.0726883\\
2.4977	-0.0723257\\
2.5003	-0.071945\\
2.5029	-0.071561\\
2.5054	-0.0711887\\
2.508	-0.0707986\\
2.5106	-0.0704057\\
2.5131	-0.0700254\\
2.5157	-0.0696277\\
2.5183	-0.0692277\\
2.5208	-0.0688413\\
2.5234	-0.0684377\\
2.526	-0.0680326\\
2.5285	-0.0676418\\
2.5311	-0.0672342\\
2.5336	-0.0668414\\
2.5362	-0.0664323\\
2.5388	-0.0660226\\
2.5413	-0.0656283\\
2.5439	-0.0652183\\
2.5465	-0.0648083\\
2.549	-0.0644143\\
2.5516	-0.0640052\\
2.5542	-0.0635967\\
2.5567	-0.0632048\\
2.5593	-0.0627983\\
2.5618	-0.0624087\\
2.5644	-0.062005\\
2.567	-0.061603\\
2.5695	-0.0612182\\
2.5721	-0.0608201\\
2.5747	-0.0604243\\
2.5772	-0.0600461\\
2.5798	-0.0596553\\
2.5824	-0.0592673\\
2.5849	-0.0588971\\
2.5875	-0.0585153\\
2.59	-0.0581513\\
2.5926	-0.0577762\\
2.5952	-0.0574049\\
2.5977	-0.0570515\\
2.6003	-0.0566879\\
2.6029	-0.0563286\\
2.6054	-0.0559872\\
2.608	-0.0556367\\
2.6106	-0.0552908\\
2.6131	-0.0549629\\
2.6157	-0.0546267\\
2.6182	-0.0543083\\
2.6208	-0.0539824\\
2.6234	-0.0536619\\
2.6259	-0.053359\\
2.6285	-0.0530495\\
2.6311	-0.0527459\\
2.6336	-0.0524596\\
2.6362	-0.0521679\\
2.6388	-0.0518823\\
2.6413	-0.0516137\\
2.6439	-0.0513407\\
2.6464	-0.0510844\\
2.649	-0.0508243\\
2.6516	-0.050571\\
2.6541	-0.0503339\\
2.6567	-0.0500942\\
2.6593	-0.0498615\\
2.6618	-0.0496445\\
2.6644	-0.049426\\
2.667	-0.0492147\\
2.6695	-0.0490185\\
2.6721	-0.0488218\\
2.6746	-0.0486397\\
2.6772	-0.0484579\\
2.6798	-0.0482837\\
2.6823	-0.0481234\\
2.6849	-0.0479644\\
2.6875	-0.0478132\\
2.69	-0.0476753\\
2.6926	-0.0475396\\
2.6952	-0.0474119\\
2.6977	-0.0472967\\
2.7003	-0.0471847\\
2.7029	-0.0470807\\
2.7054	-0.0469884\\
2.708	-0.0469004\\
2.7105	-0.0468234\\
2.7131	-0.0467514\\
2.7157	-0.0466875\\
2.7182	-0.0466337\\
2.7208	-0.0465859\\
2.7234	-0.0465462\\
2.7259	-0.0465157\\
2.7285	-0.046492\\
2.7311	-0.0464765\\
2.7336	-0.0464693\\
2.7362	-0.0464698\\
2.7387	-0.0464779\\
2.7413	-0.0464943\\
2.7439	-0.0465187\\
2.7464	-0.0465498\\
2.749	-0.0465899\\
2.7516	-0.0466381\\
2.7541	-0.0466919\\
2.7567	-0.0467556\\
2.7593	-0.0468271\\
2.7618	-0.0469032\\
2.7644	-0.04699\\
2.7669	-0.0470807\\
2.7695	-0.0471826\\
2.7721	-0.047292\\
2.7746	-0.0474043\\
2.7772	-0.0475285\\
2.7798	-0.04766\\
2.7823	-0.0477934\\
2.7849	-0.0479392\\
2.7875	-0.0480922\\
2.79	-0.048246\\
2.7926	-0.0484129\\
2.7951	-0.0485799\\
2.7977	-0.0487603\\
2.8003	-0.0489475\\
2.8028	-0.0491338\\
2.8054	-0.049334\\
2.808	-0.0495406\\
2.8105	-0.0497454\\
2.8131	-0.0499646\\
2.8157	-0.0501899\\
2.8182	-0.0504123\\
2.8208	-0.0506495\\
2.8233	-0.0508831\\
2.8259	-0.0511318\\
2.8285	-0.0513861\\
2.831	-0.0516358\\
2.8336	-0.0519009\\
2.8362	-0.0521713\\
2.8387	-0.0524362\\
2.8413	-0.0527167\\
2.8439	-0.0530021\\
2.8464	-0.0532811\\
2.849	-0.053576\\
2.8516	-0.0538754\\
2.8541	-0.0541674\\
2.8567	-0.0544755\\
2.8592	-0.0547756\\
2.8618	-0.0550917\\
2.8644	-0.0554117\\
2.8669	-0.055723\\
2.8695	-0.0560504\\
2.8721	-0.0563813\\
2.8746	-0.0567026\\
2.8772	-0.05704\\
2.8798	-0.0573804\\
2.8823	-0.0577106\\
2.8849	-0.0580567\\
2.8874	-0.058392\\
2.89	-0.0587432\\
2.8926	-0.0590967\\
2.8951	-0.0594388\\
2.8977	-0.0597966\\
2.9003	-0.0601563\\
2.9028	-0.0605038\\
2.9054	-0.0608668\\
2.908	-0.0612313\\
2.9105	-0.061583\\
2.9131	-0.06195\\
2.9156	-0.0623038\\
2.9182	-0.0626725\\
2.9208	-0.063042\\
2.9233	-0.0633977\\
2.9259	-0.0637681\\
2.9285	-0.0641387\\
2.931	-0.0644952\\
2.9336	-0.0648658\\
2.9362	-0.0652362\\
2.9387	-0.0655919\\
2.9413	-0.0659614\\
2.9438	-0.0663159\\
2.9464	-0.0666838\\
2.949	-0.0670507\\
2.9515	-0.0674023\\
2.9541	-0.0677667\\
2.9567	-0.0681297\\
2.9592	-0.068477\\
2.9618	-0.0688366\\
2.9644	-0.0691941\\
2.9669	-0.0695359\\
2.9695	-0.0698891\\
2.972	-0.0702265\\
2.9746	-0.0705748\\
2.9772	-0.0709203\\
2.9797	-0.0712499\\
2.9823	-0.0715897\\
2.9849	-0.0719262\\
2.9874	-0.0722467\\
2.99	-0.0725765\\
2.9926	-0.0729027\\
2.9951	-0.0732128\\
2.9977	-0.0735314\\
3.0003	-0.0738459\\
3.0028	-0.0741444\\
3.0054	-0.0744505\\
3.0079	-0.0747406\\
3.0105	-0.0750377\\
3.0131	-0.07533\\
3.0156	-0.0756064\\
3.0182	-0.075889\\
3.0208	-0.0761663\\
3.0233	-0.0764279\\
3.0259	-0.0766947\\
3.0285	-0.0769559\\
3.031	-0.0772016\\
3.0336	-0.0774514\\
3.0361	-0.077686\\
3.0387	-0.0779241\\
3.0413	-0.0781559\\
3.0438	-0.0783729\\
3.0464	-0.0785922\\
3.049	-0.078805\\
3.0515	-0.0790034\\
3.0541	-0.079203\\
3.0567	-0.0793958\\
3.0592	-0.0795746\\
3.0618	-0.0797537\\
3.0643	-0.0799191\\
3.0669	-0.080084\\
3.0695	-0.0802416\\
3.072	-0.0803862\\
3.0746	-0.0805291\\
3.0772	-0.0806645\\
3.0797	-0.0807874\\
3.0823	-0.0809077\\
3.0849	-0.0810201\\
3.0874	-0.0811208\\
3.09	-0.0812177\\
3.0925	-0.0813033\\
3.0951	-0.0813844\\
3.0977	-0.0814575\\
3.1002	-0.0815199\\
3.1028	-0.0815769\\
3.1054	-0.0816255\\
3.1079	-0.0816644\\
3.1105	-0.0816967\\
3.1131	-0.0817206\\
3.1156	-0.0817357\\
3.1182	-0.081743\\
3.1207	-0.0817421\\
3.1233	-0.0817329\\
3.1259	-0.0817151\\
3.1284	-0.08169\\
3.131	-0.0816555\\
3.1336	-0.0816124\\
3.1361	-0.081563\\
3.1387	-0.0815032\\
3.1413	-0.0814349\\
3.1438	-0.0813612\\
3.1464	-0.0812762\\
3.1489	-0.0811865\\
3.1515	-0.0810849\\
3.1541	-0.0809748\\
3.1566	-0.0808611\\
3.1592	-0.0807346\\
3.1618	-0.0805997\\
3.1643	-0.0804622\\
3.1669	-0.0803111\\
3.1695	-0.0801518\\
3.172	-0.079991\\
3.1746	-0.0798158\\
3.1772	-0.0796325\\
3.1797	-0.0794488\\
3.1823	-0.07925\\
3.1848	-0.0790514\\
3.1874	-0.0788373\\
3.19	-0.0786155\\
3.1925	-0.078395\\
3.1951	-0.0781584\\
3.1977	-0.0779143\\
3.2002	-0.0776728\\
3.2028	-0.0774144\\
3.2054	-0.0771489\\
3.2079	-0.076887\\
3.2105	-0.0766078\\
3.213	-0.0763329\\
3.2156	-0.0760405\\
3.2182	-0.0757414\\
3.2207	-0.0754477\\
3.2233	-0.0751361\\
3.2259	-0.0748182\\
3.2284	-0.0745067\\
3.231	-0.074177\\
3.2336	-0.0738414\\
3.2361	-0.0735133\\
3.2387	-0.0731666\\
3.2412	-0.0728282\\
3.2438	-0.072471\\
3.2464	-0.0721087\\
3.2489	-0.0717555\\
3.2515	-0.0713835\\
3.2541	-0.0710068\\
3.2566	-0.0706403\\
3.2592	-0.0702548\\
3.2618	-0.0698651\\
3.2643	-0.0694865\\
3.2669	-0.0690889\\
3.2694	-0.0687031\\
3.272	-0.0682984\\
3.2746	-0.0678901\\
3.2771	-0.0674946\\
3.2797	-0.0670801\\
3.2823	-0.0666628\\
3.2848	-0.0662588\\
3.2874	-0.0658362\\
3.29	-0.0654112\\
3.2925	-0.0650004\\
3.2951	-0.0645712\\
3.2976	-0.0641566\\
3.3002	-0.0637238\\
3.3028	-0.0632895\\
3.3053	-0.0628705\\
3.3079	-0.0624336\\
3.3105	-0.0619957\\
3.313	-0.0615739\\
3.3156	-0.0611345\\
3.3182	-0.0606946\\
3.3207	-0.0602713\\
3.3233	-0.0598309\\
3.3259	-0.0593906\\
3.3284	-0.0589674\\
3.331	-0.0585276\\
3.3335	-0.0581052\\
3.3361	-0.0576667\\
3.3387	-0.057229\\
3.3412	-0.0568092\\
3.3438	-0.0563738\\
3.3464	-0.0559398\\
3.3489	-0.0555239\\
3.3515	-0.0550931\\
3.3541	-0.0546641\\
3.3566	-0.0542536\\
3.3592	-0.0538289\\
3.3617	-0.0534227\\
3.3643	-0.0530027\\
3.3669	-0.0525855\\
3.3694	-0.0521869\\
3.372	-0.0517754\\
3.3746	-0.0513669\\
3.3771	-0.0509773\\
3.3797	-0.0505755\\
3.3823	-0.0501772\\
3.3848	-0.0497978\\
3.3874	-0.049407\\
3.3899	-0.049035\\
3.3925	-0.0486521\\
3.3951	-0.0482735\\
3.3976	-0.0479137\\
3.4002	-0.0475438\\
3.4028	-0.0471786\\
3.4053	-0.046832\\
3.4079	-0.0464763\\
3.4105	-0.0461256\\
3.413	-0.0457932\\
3.4156	-0.0454526\\
3.4181	-0.0451302\\
3.4207	-0.0448003\\
3.4233	-0.0444759\\
3.4258	-0.0441693\\
3.4284	-0.043856\\
3.431	-0.0435486\\
3.4335	-0.0432587\\
3.4361	-0.042963\\
3.4387	-0.0426735\\
3.4412	-0.042401\\
3.4438	-0.0421237\\
3.4463	-0.041863\\
3.4489	-0.0415982\\
3.4515	-0.04134\\
3.454	-0.0410978\\
3.4566	-0.0408525\\
3.4592	-0.0406138\\
3.4617	-0.0403907\\
3.4643	-0.0401653\\
3.4669	-0.0399468\\
3.4694	-0.0397432\\
3.472	-0.0395382\\
3.4745	-0.0393478\\
3.4771	-0.0391566\\
3.4797	-0.0389725\\
3.4822	-0.0388021\\
3.4848	-0.0386319\\
3.4874	-0.0384689\\
3.4899	-0.038319\\
3.4925	-0.03817\\
3.4951	-0.0380283\\
3.4976	-0.0378989\\
3.5002	-0.0377714\\
3.5028	-0.0376512\\
3.5053	-0.0375424\\
3.5079	-0.0374364\\
3.5104	-0.0373414\\
3.513	-0.0372496\\
3.5156	-0.0371651\\
3.5181	-0.0370907\\
3.5207	-0.0370205\\
3.5233	-0.0369574\\
3.5258	-0.0369036\\
3.5284	-0.0368546\\
3.531	-0.0368129\\
3.5335	-0.0367795\\
3.5361	-0.0367517\\
3.5386	-0.0367317\\
3.5412	-0.0367177\\
3.5438	-0.0367109\\
3.5463	-0.0367108\\
3.5489	-0.0367176\\
3.5515	-0.0367313\\
3.554	-0.0367509\\
3.5566	-0.036778\\
3.5592	-0.0368118\\
3.5617	-0.0368507\\
3.5643	-0.0368977\\
3.5668	-0.036949\\
3.5694	-0.0370089\\
3.572	-0.0370753\\
3.5745	-0.0371451\\
3.5771	-0.037224\\
3.5797	-0.0373092\\
3.5822	-0.037397\\
3.5848	-0.0374943\\
3.5874	-0.0375978\\
3.5899	-0.0377029\\
3.5925	-0.037818\\
3.595	-0.0379343\\
3.5976	-0.0380609\\
3.6002	-0.0381932\\
3.6027	-0.0383257\\
3.6053	-0.0384689\\
3.6079	-0.0386176\\
3.6104	-0.0387657\\
3.613	-0.0389248\\
3.6156	-0.0390892\\
3.6181	-0.039252\\
3.6207	-0.0394263\\
3.6232	-0.0395985\\
3.6258	-0.0397823\\
3.6284	-0.0399709\\
3.6309	-0.0401566\\
3.6335	-0.0403541\\
3.6361	-0.0405561\\
3.6386	-0.0407544\\
3.6412	-0.0409648\\
3.6438	-0.0411794\\
3.6463	-0.0413894\\
3.6489	-0.0416118\\
3.6515	-0.0418379\\
3.654	-0.0420589\\
3.6566	-0.0422922\\
3.6591	-0.0425199\\
3.6617	-0.04276\\
3.6643	-0.0430034\\
3.6668	-0.0432405\\
3.6694	-0.04349\\
3.672	-0.0437425\\
3.6745	-0.0439879\\
3.6771	-0.0442458\\
3.6797	-0.0445064\\
3.6822	-0.0447593\\
3.6848	-0.0450247\\
3.6873	-0.045282\\
3.6899	-0.0455518\\
3.6925	-0.0458236\\
3.695	-0.0460867\\
3.6976	-0.0463622\\
3.7002	-0.0466394\\
3.7027	-0.0469075\\
3.7053	-0.0471877\\
3.7079	-0.0474693\\
3.7104	-0.0477412\\
3.713	-0.0480251\\
3.7155	-0.0482991\\
3.7181	-0.0485849\\
3.7207	-0.0488715\\
3.7232	-0.0491477\\
3.7258	-0.0494355\\
3.7284	-0.0497237\\
3.7309	-0.0500012\\
3.7335	-0.0502899\\
3.7361	-0.0505788\\
3.7386	-0.0508566\\
3.7412	-0.0511453\\
3.7437	-0.0514227\\
3.7463	-0.0517109\\
3.7489	-0.0519985\\
3.7514	-0.0522746\\
3.754	-0.0525611\\
3.7566	-0.0528467\\
3.7591	-0.0531205\\
3.7617	-0.0534043\\
3.7643	-0.0536869\\
3.7668	-0.0539575\\
3.7694	-0.0542376\\
3.7719	-0.0545056\\
3.7745	-0.0547827\\
3.7771	-0.0550583\\
3.7796	-0.0553215\\
3.7822	-0.0555934\\
3.7848	-0.0558634\\
3.7873	-0.0561211\\
3.7899	-0.0563869\\
3.7925	-0.0566504\\
3.795	-0.0569016\\
3.7976	-0.0571604\\
3.8002	-0.0574167\\
3.8027	-0.0576606\\
3.8053	-0.0579115\\
3.8078	-0.0581501\\
3.8104	-0.0583954\\
3.813	-0.0586377\\
3.8155	-0.0588677\\
3.8181	-0.0591038\\
3.8207	-0.0593367\\
3.8232	-0.0595574\\
3.8258	-0.0597835\\
3.8284	-0.0600061\\
3.8309	-0.0602168\\
3.8335	-0.0604323\\
3.836	-0.0606359\\
3.8386	-0.060844\\
3.8412	-0.0610481\\
3.8437	-0.0612406\\
3.8463	-0.0614369\\
3.8489	-0.0616291\\
3.8514	-0.0618099\\
3.854	-0.0619938\\
3.8566	-0.0621734\\
3.8591	-0.062342\\
3.8617	-0.062513\\
3.8642	-0.0626732\\
3.8668	-0.0628353\\
3.8694	-0.0629929\\
3.8719	-0.0631401\\
3.8745	-0.0632885\\
3.8771	-0.0634323\\
3.8796	-0.063566\\
3.8822	-0.0637003\\
3.8848	-0.0638298\\
3.8873	-0.0639497\\
3.8899	-0.0640695\\
3.8924	-0.0641801\\
3.895	-0.0642902\\
3.8976	-0.0643953\\
3.9001	-0.0644916\\
3.9027	-0.0645868\\
3.9053	-0.0646769\\
3.9078	-0.0647587\\
3.9104	-0.0648388\\
3.913	-0.0649137\\
3.9155	-0.0649809\\
3.9181	-0.0650457\\
3.9206	-0.0651032\\
3.9232	-0.0651578\\
3.9258	-0.0652073\\
3.9283	-0.06525\\
3.9309	-0.0652893\\
3.9335	-0.0653235\\
3.936	-0.0653514\\
3.9386	-0.0653755\\
3.9412	-0.0653943\\
3.9437	-0.0654076\\
3.9463	-0.0654164\\
3.9488	-0.06542\\
3.9514	-0.0654188\\
3.954	-0.0654126\\
3.9565	-0.0654018\\
3.9591	-0.0653856\\
3.9617	-0.0653645\\
3.9642	-0.0653395\\
3.9668	-0.0653087\\
3.9694	-0.0652731\\
3.9719	-0.0652342\\
3.9745	-0.065189\\
3.9771	-0.0651391\\
3.9796	-0.0650866\\
3.9822	-0.0650275\\
3.9847	-0.0649663\\
3.9873	-0.0648981\\
3.9899	-0.0648254\\
3.9924	-0.0647513\\
3.995	-0.0646699\\
3.9976	-0.0645841\\
4.0001	-0.0644976\\
4.0027	-0.0644035\\
4.0053	-0.0643052\\
4.0078	-0.0642069\\
4.0104	-0.0641006\\
4.0129	-0.0639947\\
4.0155	-0.0638807\\
4.0181	-0.0637629\\
4.0206	-0.0636462\\
4.0232	-0.0635211\\
4.0258	-0.0633925\\
4.0283	-0.0632655\\
4.0309	-0.06313\\
4.0335	-0.0629912\\
4.036	-0.0628547\\
4.0386	-0.0627096\\
4.0411	-0.0625672\\
4.0437	-0.0624161\\
4.0463	-0.0622621\\
4.0488	-0.0621114\\
4.0514	-0.061952\\
4.054	-0.0617899\\
4.0565	-0.0616317\\
4.0591	-0.0614647\\
4.0617	-0.0612954\\
4.0642	-0.0611305\\
4.0668	-0.0609569\\
4.0693	-0.060788\\
4.0719	-0.0606104\\
4.0745	-0.0604309\\
4.077	-0.0602567\\
4.0796	-0.0600739\\
4.0822	-0.0598896\\
4.0847	-0.0597109\\
4.0873	-0.0595238\\
4.0899	-0.0593355\\
4.0924	-0.0591533\\
4.095	-0.0589627\\
4.0975	-0.0587787\\
4.1001	-0.0585864\\
4.1027	-0.0583934\\
4.1052	-0.0582073\\
4.1078	-0.0580132\\
4.1104	-0.0578186\\
4.1129	-0.0576313\\
4.1155	-0.0574362\\
4.1181	-0.0572411\\
4.1206	-0.0570534\\
4.1232	-0.0568584\\
4.1258	-0.0566635\\
4.1283	-0.0564765\\
4.1309	-0.0562823\\
4.1334	-0.0560961\\
4.136	-0.055903\\
4.1386	-0.0557107\\
4.1411	-0.0555265\\
4.1437	-0.0553358\\
4.1463	-0.0551462\\
4.1488	-0.0549649\\
4.1514	-0.0547775\\
4.154	-0.0545914\\
4.1565	-0.0544138\\
4.1591	-0.0542305\\
4.1616	-0.0540558\\
4.1642	-0.0538757\\
4.1668	-0.0536974\\
4.1693	-0.0535277\\
4.1719	-0.0533531\\
4.1745	-0.0531805\\
4.177	-0.0530166\\
4.1796	-0.0528482\\
4.1822	-0.0526822\\
4.1847	-0.0525248\\
4.1873	-0.0523634\\
4.1898	-0.0522107\\
4.1924	-0.0520544\\
4.195	-0.0519008\\
4.1975	-0.0517556\\
4.2001	-0.0516075\\
4.2027	-0.0514622\\
4.2052	-0.0513254\\
4.2078	-0.051186\\
4.2104	-0.0510498\\
4.2129	-0.0509218\\
4.2155	-0.0507918\\
4.218	-0.0506699\\
4.2206	-0.0505465\\
4.2232	-0.0504265\\
4.2257	-0.0503143\\
4.2283	-0.0502011\\
4.2309	-0.0500915\\
4.2334	-0.0499896\\
4.236	-0.0498871\\
4.2386	-0.0497884\\
4.2411	-0.0496971\\
4.2437	-0.0496058\\
4.2462	-0.0495217\\
4.2488	-0.049438\\
4.2514	-0.0493583\\
4.2539	-0.0492853\\
4.2565	-0.0492134\\
4.2591	-0.0491455\\
4.2616	-0.0490841\\
4.2642	-0.0490242\\
4.2668	-0.0489684\\
4.2693	-0.0489186\\
4.2719	-0.048871\\
4.2744	-0.0488291\\
4.277	-0.0487897\\
4.2796	-0.0487544\\
4.2821	-0.0487246\\
4.2847	-0.0486977\\
4.2873	-0.048675\\
4.2898	-0.0486572\\
4.2924	-0.0486429\\
4.295	-0.0486328\\
4.2975	-0.0486272\\
4.3001	-0.0486255\\
4.3027	-0.0486281\\
4.3052	-0.0486347\\
4.3078	-0.0486456\\
4.3103	-0.0486602\\
4.3129	-0.0486794\\
4.3155	-0.048703\\
4.318	-0.0487295\\
4.3206	-0.0487613\\
4.3232	-0.0487973\\
4.3257	-0.0488358\\
4.3283	-0.0488799\\
4.3309	-0.0489281\\
4.3334	-0.0489783\\
4.336	-0.0490346\\
4.3385	-0.0490925\\
4.3411	-0.0491566\\
4.3437	-0.0492248\\
4.3462	-0.049294\\
4.3488	-0.0493699\\
4.3514	-0.0494496\\
4.3539	-0.0495299\\
4.3565	-0.0496172\\
4.3591	-0.0497082\\
4.3616	-0.0497992\\
4.3642	-0.0498975\\
4.3667	-0.0499955\\
4.3693	-0.0501008\\
4.3719	-0.0502098\\
4.3744	-0.0503178\\
4.377	-0.0504335\\
4.3796	-0.0505526\\
4.3821	-0.0506703\\
4.3847	-0.0507959\\
4.3873	-0.0509248\\
4.3898	-0.0510516\\
4.3924	-0.0511866\\
4.3949	-0.0513193\\
4.3975	-0.0514602\\
4.4001	-0.0516041\\
4.4026	-0.0517451\\
4.4052	-0.0518945\\
4.4078	-0.0520467\\
4.4103	-0.0521955\\
4.4129	-0.0523529\\
4.4155	-0.0525127\\
4.418	-0.0526688\\
4.4206	-0.0528334\\
4.4231	-0.052994\\
4.4257	-0.0531631\\
4.4283	-0.0533344\\
4.4308	-0.0535012\\
4.4334	-0.0536766\\
4.436	-0.0538539\\
4.4385	-0.0540263\\
4.4411	-0.0542073\\
4.4437	-0.05439\\
4.4462	-0.0545673\\
4.4488	-0.0547532\\
4.4514	-0.0549406\\
4.4539	-0.0551222\\
4.4565	-0.0553123\\
4.459	-0.0554964\\
4.4616	-0.0556889\\
4.4642	-0.0558826\\
4.4667	-0.0560698\\
4.4693	-0.0562654\\
4.4719	-0.0564618\\
4.4744	-0.0566515\\
4.477	-0.0568494\\
4.4796	-0.057048\\
4.4821	-0.0572394\\
4.4847	-0.057439\\
4.4872	-0.0576312\\
4.4898	-0.0578314\\
4.4924	-0.0580318\\
4.4949	-0.0582246\\
4.4975	-0.0584252\\
4.5001	-0.0586258\\
4.5026	-0.0588186\\
4.5052	-0.0590188\\
4.5078	-0.0592188\\
4.5103	-0.0594108\\
4.5129	-0.05961\\
4.5154	-0.0598011\\
4.518	-0.0599992\\
4.5206	-0.0601967\\
4.5231	-0.0603859\\
4.5257	-0.0605818\\
4.5283	-0.0607768\\
4.5308	-0.0609633\\
4.5334	-0.0611563\\
4.536	-0.0613481\\
4.5385	-0.0615314\\
4.5411	-0.0617208\\
4.5436	-0.0619015\\
4.5462	-0.0620881\\
4.5488	-0.0622732\\
4.5513	-0.0624497\\
4.5539	-0.0626315\\
4.5565	-0.0628117\\
4.559	-0.0629832\\
4.5616	-0.0631597\\
4.5642	-0.0633343\\
4.5667	-0.0635002\\
4.5693	-0.0636707\\
4.5718	-0.0638327\\
4.5744	-0.0639989\\
4.577	-0.0641629\\
4.5795	-0.0643183\\
4.5821	-0.0644776\\
4.5847	-0.0646345\\
4.5872	-0.0647829\\
4.5898	-0.0649347\\
4.5924	-0.0650838\\
4.5949	-0.0652247\\
4.5975	-0.0653684\\
4.6001	-0.0655094\\
4.6026	-0.0656422\\
4.6052	-0.0657774\\
4.6077	-0.0659047\\
4.6103	-0.066034\\
4.6129	-0.0661603\\
4.6154	-0.0662788\\
4.618	-0.0663988\\
4.6206	-0.0665157\\
4.6231	-0.066625\\
4.6257	-0.0667354\\
4.6283	-0.0668425\\
4.6308	-0.0669423\\
4.6334	-0.0670427\\
4.6359	-0.0671359\\
4.6385	-0.0672295\\
4.6411	-0.0673195\\
4.6436	-0.0674027\\
4.6462	-0.0674857\\
4.6488	-0.0675651\\
4.6513	-0.0676379\\
4.6539	-0.0677101\\
4.6565	-0.0677786\\
4.659	-0.0678409\\
4.6616	-0.067902\\
4.6641	-0.0679572\\
4.6667	-0.068011\\
4.6693	-0.0680608\\
4.6718	-0.0681052\\
4.6744	-0.0681476\\
4.677	-0.0681861\\
4.6795	-0.0682195\\
4.6821	-0.0682504\\
4.6847	-0.0682775\\
4.6872	-0.0682998\\
4.6898	-0.0683192\\
4.6923	-0.0683342\\
4.6949	-0.068346\\
4.6975	-0.0683539\\
4.7	-0.0683577\\
4.7026	-0.068358\\
4.7052	-0.0683543\\
4.7077	-0.0683471\\
4.7103	-0.0683358\\
4.7129	-0.0683207\\
4.7154	-0.0683025\\
4.718	-0.0682798\\
4.7205	-0.0682544\\
4.7231	-0.0682242\\
4.7257	-0.0681902\\
4.7282	-0.068154\\
4.7308	-0.0681126\\
4.7334	-0.0680675\\
4.7359	-0.0680206\\
4.7385	-0.0679682\\
4.7411	-0.0679122\\
4.7436	-0.0678549\\
4.7462	-0.0677918\\
4.7487	-0.0677277\\
4.7513	-0.0676576\\
4.7539	-0.067584\\
4.7564	-0.0675099\\
4.759	-0.0674295\\
4.7616	-0.0673458\\
4.7641	-0.067262\\
4.7667	-0.0671717\\
4.7693	-0.067078\\
4.7718	-0.066985\\
4.7744	-0.066885\\
4.777	-0.066782\\
4.7795	-0.0666799\\
4.7821	-0.0665708\\
4.7846	-0.0664631\\
4.7872	-0.0663481\\
4.7898	-0.0662302\\
4.7923	-0.0661142\\
4.7949	-0.0659908\\
4.7975	-0.0658647\\
4.8	-0.0657408\\
4.8026	-0.0656095\\
4.8052	-0.0654755\\
4.8077	-0.0653444\\
4.8103	-0.0652056\\
4.8128	-0.0650698\\
4.8154	-0.0649264\\
4.818	-0.0647806\\
4.8205	-0.0646384\\
4.8231	-0.0644884\\
4.8257	-0.0643363\\
4.8282	-0.0641881\\
4.8308	-0.0640322\\
4.8334	-0.0638743\\
4.8359	-0.0637207\\
4.8385	-0.0635593\\
4.841	-0.0634026\\
4.8436	-0.0632379\\
4.8462	-0.0630718\\
4.8487	-0.0629106\\
4.8513	-0.0627415\\
4.8539	-0.0625712\\
4.8564	-0.0624062\\
4.859	-0.0622334\\
4.8616	-0.0620595\\
4.8641	-0.0618913\\
4.8667	-0.0617153\\
4.8692	-0.0615453\\
4.8718	-0.0613677\\
4.8744	-0.0611893\\
4.8769	-0.0610171\\
4.8795	-0.0608374\\
4.8821	-0.0606571\\
4.8846	-0.0604833\\
4.8872	-0.0603022\\
4.8898	-0.0601207\\
4.8923	-0.059946\\
4.8949	-0.0597641\\
4.8974	-0.0595891\\
4.9	-0.059407\\
4.9026	-0.0592249\\
4.9051	-0.05905\\
4.9077	-0.0588682\\
4.9103	-0.0586866\\
4.9128	-0.0585123\\
4.9154	-0.0583314\\
4.918	-0.058151\\
4.9205	-0.0579779\\
4.9231	-0.0577986\\
4.9257	-0.0576198\\
4.9282	-0.0574487\\
4.9308	-0.0572714\\
4.9333	-0.0571018\\
4.9359	-0.0569263\\
4.9385	-0.0567517\\
4.941	-0.0565849\\
4.9436	-0.0564125\\
4.9462	-0.0562412\\
4.9487	-0.0560778\\
4.9513	-0.055909\\
4.9539	-0.0557416\\
4.9564	-0.055582\\
4.959	-0.0554174\\
4.9615	-0.0552607\\
4.9641	-0.0550992\\
4.9667	-0.0549393\\
4.9692	-0.0547873\\
4.9718	-0.0546308\\
4.9744	-0.0544762\\
4.9769	-0.0543293\\
4.9795	-0.0541784\\
4.9821	-0.0540295\\
4.9846	-0.0538882\\
4.9872	-0.0537433\\
4.9897	-0.0536059\\
4.9923	-0.0534652\\
4.9949	-0.0533267\\
4.9974	-0.0531957\\
5	-0.0530616\\
};
\addplot [color=mycolor4, line width=1.0pt, forget plot]
  table[row sep=crcr]{%
-0.125	-9.09673e-08\\
-0.1224	-9.28302e-08\\
-0.1199	-9.46208e-08\\
-0.1173	-9.64823e-08\\
-0.1147	-9.83433e-08\\
-0.1122	-1.00132e-07\\
-0.1096	-1.01992e-07\\
-0.1071	-1.03779e-07\\
-0.1045	-1.05637e-07\\
-0.1019	-1.07495e-07\\
-0.0994	-1.09281e-07\\
-0.0968	-1.11137e-07\\
-0.0942	-1.12992e-07\\
-0.0917	-1.14776e-07\\
-0.0891	-1.1663e-07\\
-0.0865	-1.18484e-07\\
-0.084	-1.20266e-07\\
-0.0814	-1.22118e-07\\
-0.0789	-1.23899e-07\\
-0.0763	-1.2575e-07\\
-0.0737	-1.276e-07\\
-0.0712	-1.29379e-07\\
-0.0686	-1.31228e-07\\
-0.066	-1.33076e-07\\
-0.0635	-1.34853e-07\\
-0.0609	-1.367e-07\\
-0.0583	-1.38547e-07\\
-0.0558	-1.40321e-07\\
-0.0532	-1.42167e-07\\
-0.0507	-1.4394e-07\\
-0.0481	-1.45784e-07\\
-0.0455	-1.47627e-07\\
-0.043	-1.49399e-07\\
-0.0404	-1.51241e-07\\
-0.0378	-1.53083e-07\\
-0.0353	-1.54853e-07\\
-0.0327	-1.56693e-07\\
-0.0301	-1.58532e-07\\
-0.0276	-1.60301e-07\\
-0.025	-1.62139e-07\\
-0.0224	-1.63977e-07\\
-0.0199	-1.65743e-07\\
-0.0173	-1.6758e-07\\
-0.0148	-1.69345e-07\\
-0.0122	-1.7118e-07\\
-0.0096	-1.73015e-07\\
-0.0071	-1.74779e-07\\
-0.0045	-1.76612e-07\\
-0.0019	-1.78445e-07\\
0.0006	-1.80208e-07\\
0.0032	0.000554893\\
0.0058	0.00238121\\
0.0083	0.00419617\\
0.0109	0.00613477\\
0.0134	0.00804142\\
0.016	0.00997294\\
0.0186	0.0117672\\
0.0211	0.0133498\\
0.0237	0.0148921\\
0.0263	0.0163834\\
0.0288	0.0177923\\
0.0314	0.0192203\\
0.034	0.0205863\\
0.0365	0.0218293\\
0.0391	0.0230559\\
0.0416	0.0241887\\
0.0442	0.0253323\\
0.0468	0.0264433\\
0.0493	0.0274746\\
0.0519	0.0285023\\
0.0545	0.0294843\\
0.057	0.0303922\\
0.0596	0.0313072\\
0.0622	0.0321973\\
0.0647	0.0330288\\
0.0673	0.0338646\\
0.0698	0.0346391\\
0.0724	0.035417\\
0.075	0.0361729\\
0.0775	0.0368839\\
0.0801	0.0376081\\
0.0827	0.0383146\\
0.0852	0.0389749\\
0.0878	0.0396423\\
0.0904	0.0402938\\
0.0929	0.0409097\\
0.0955	0.0415418\\
0.098	0.0421406\\
0.1006	0.0427513\\
0.1032	0.0433481\\
0.1057	0.0439105\\
0.1083	0.0444871\\
0.1109	0.0450586\\
0.1134	0.0456034\\
0.116	0.046163\\
0.1186	0.0467127\\
0.1211	0.0472316\\
0.1237	0.0477637\\
0.1263	0.0482913\\
0.1288	0.0487957\\
0.1314	0.0493158\\
0.1339	0.049809\\
0.1365	0.050313\\
0.1391	0.0508086\\
0.1416	0.05128\\
0.1442	0.0517669\\
0.1468	0.0522503\\
0.1493	0.0527094\\
0.1519	0.0531786\\
0.1545	0.0536389\\
0.157	0.0540745\\
0.1596	0.0545229\\
0.1621	0.0549505\\
0.1647	0.05539\\
0.1673	0.0558218\\
0.1698	0.0562281\\
0.1724	0.056642\\
0.175	0.0570494\\
0.1775	0.0574366\\
0.1801	0.0578344\\
0.1827	0.0582252\\
0.1852	0.0585923\\
0.1878	0.0589648\\
0.1903	0.0593153\\
0.1929	0.0596741\\
0.1955	0.0600278\\
0.198	0.0603622\\
0.2006	0.0607019\\
0.2032	0.0610321\\
0.2057	0.0613413\\
0.2083	0.0616562\\
0.2109	0.0619659\\
0.2134	0.0622588\\
0.216	0.0625565\\
0.2185	0.0628348\\
0.2211	0.0631156\\
0.2237	0.0633891\\
0.2262	0.0636469\\
0.2288	0.0639106\\
0.2314	0.0641688\\
0.2339	0.0644104\\
0.2365	0.0646538\\
0.2391	0.0648898\\
0.2416	0.0651115\\
0.2442	0.0653379\\
0.2467	0.0655517\\
0.2493	0.0657686\\
0.2519	0.0659787\\
0.2544	0.0661741\\
0.257	0.0663716\\
0.2596	0.066565\\
0.2621	0.0667477\\
0.2647	0.0669337\\
0.2673	0.067114\\
0.2698	0.0672814\\
0.2724	0.0674497\\
0.2749	0.0676073\\
0.2775	0.0677681\\
0.2801	0.0679254\\
0.2826	0.0680724\\
0.2852	0.0682197\\
0.2878	0.068361\\
0.2903	0.068492\\
0.2929	0.0686247\\
0.2955	0.0687541\\
0.298	0.0688748\\
0.3006	0.0689951\\
0.3032	0.0691093\\
0.3057	0.0692136\\
0.3083	0.0693175\\
0.3108	0.0694139\\
0.3134	0.0695101\\
0.316	0.0696012\\
0.3185	0.0696828\\
0.3211	0.0697612\\
0.3237	0.0698337\\
0.3262	0.0698987\\
0.3288	0.0699618\\
0.3314	0.0700195\\
0.3339	0.0700688\\
0.3365	0.0701127\\
0.339	0.0701482\\
0.3416	0.0701788\\
0.3442	0.0702037\\
0.3467	0.070222\\
0.3493	0.0702341\\
0.3519	0.0702381\\
0.3544	0.0702341\\
0.357	0.0702222\\
0.3596	0.0702032\\
0.3621	0.0701785\\
0.3647	0.0701454\\
0.3672	0.0701054\\
0.3698	0.0700547\\
0.3724	0.0699948\\
0.3749	0.0699293\\
0.3775	0.0698533\\
0.3801	0.0697693\\
0.3826	0.0696799\\
0.3852	0.0695772\\
0.3878	0.0694645\\
0.3903	0.0693469\\
0.3929	0.069216\\
0.3954	0.0690818\\
0.398	0.0689331\\
0.4006	0.0687742\\
0.4031	0.0686115\\
0.4057	0.0684319\\
0.4083	0.0682425\\
0.4108	0.0680516\\
0.4134	0.0678439\\
0.416	0.067626\\
0.4185	0.0674063\\
0.4211	0.0671672\\
0.4236	0.0669275\\
0.4262	0.0666686\\
0.4288	0.0664003\\
0.4313	0.0661331\\
0.4339	0.0658448\\
0.4365	0.0655458\\
0.439	0.0652482\\
0.4416	0.064929\\
0.4442	0.0646005\\
0.4467	0.0642757\\
0.4493	0.0639282\\
0.4519	0.0635703\\
0.4544	0.0632164\\
0.457	0.0628386\\
0.4595	0.0624668\\
0.4621	0.0620713\\
0.4647	0.0616667\\
0.4672	0.0612686\\
0.4698	0.060845\\
0.4724	0.0604119\\
0.4749	0.0599871\\
0.4775	0.0595371\\
0.4801	0.0590789\\
0.4826	0.0586302\\
0.4852	0.0581547\\
0.4877	0.0576892\\
0.4903	0.0571969\\
0.4929	0.0566969\\
0.4954	0.0562091\\
0.498	0.0556945\\
0.5006	0.055172\\
0.5031	0.054662\\
0.5057	0.0541241\\
0.5083	0.0535792\\
0.5108	0.0530491\\
0.5134	0.0524914\\
0.5159	0.0519489\\
0.5185	0.0513778\\
0.5211	0.0508\\
0.5236	0.0502384\\
0.5262	0.0496487\\
0.5288	0.0490536\\
0.5313	0.0484761\\
0.5339	0.0478699\\
0.5365	0.0472578\\
0.539	0.046664\\
0.5416	0.0460417\\
0.5441	0.045439\\
0.5467	0.0448078\\
0.5493	0.0441721\\
0.5518	0.0435562\\
0.5544	0.0429112\\
0.557	0.042262\\
0.5595	0.0416343\\
0.5621	0.0409781\\
0.5647	0.0403184\\
0.5672	0.0396806\\
0.5698	0.0390137\\
0.5723	0.0383692\\
0.5749	0.0376961\\
0.5775	0.0370204\\
0.58	0.0363684\\
0.5826	0.0356878\\
0.5852	0.0350045\\
0.5877	0.0343451\\
0.5903	0.0336571\\
0.5929	0.0329675\\
0.5954	0.0323029\\
0.598	0.0316102\\
0.6006	0.0309159\\
0.6031	0.0302467\\
0.6057	0.0295493\\
0.6082	0.0288778\\
0.6108	0.0281789\\
0.6134	0.0274793\\
0.6159	0.0268061\\
0.6185	0.0261052\\
0.6211	0.0254036\\
0.6236	0.0247288\\
0.6262	0.0240271\\
0.6288	0.0233258\\
0.6313	0.0226517\\
0.6339	0.0219509\\
0.6364	0.0212772\\
0.639	0.0205771\\
0.6416	0.0198777\\
0.6441	0.0192064\\
0.6467	0.0185094\\
0.6493	0.0178136\\
0.6518	0.0171456\\
0.6544	0.0164521\\
0.657	0.0157601\\
0.6595	0.0150965\\
0.6621	0.0144084\\
0.6646	0.0137487\\
0.6672	0.0130647\\
0.6698	0.0123826\\
0.6723	0.0117289\\
0.6749	0.0110516\\
0.6775	0.0103771\\
0.68	0.00973133\\
0.6826	0.00906256\\
0.6852	0.00839665\\
0.6877	0.00775913\\
0.6903	0.00709928\\
0.6928	0.00646813\\
0.6954	0.00581537\\
0.698	0.00516633\\
0.7005	0.00454569\\
0.7031	0.00390383\\
0.7057	0.00326582\\
0.7082	0.00265626\\
0.7108	0.00202659\\
0.7134	0.00140138\\
0.7159	0.000804361\\
0.7185	0.000187744\\
0.721	-0.000400958\\
0.7236	-0.00100862\\
0.7262	-0.00161137\\
0.7287	-0.00218615\\
0.7313	-0.00277895\\
0.7339	-0.00336676\\
0.7364	-0.00392723\\
0.739	-0.00450507\\
0.7416	-0.00507755\\
0.7441	-0.00562277\\
0.7467	-0.0061843\\
0.7492	-0.00671901\\
0.7518	-0.00726971\\
0.7544	-0.00781487\\
0.7569	-0.00833367\\
0.7595	-0.00886745\\
0.7621	-0.00939524\\
0.7646	-0.0098971\\
0.7672	-0.0104133\\
0.7698	-0.0109236\\
0.7723	-0.0114086\\
0.7749	-0.0119071\\
0.7775	-0.0123994\\
0.78	-0.0128668\\
0.7826	-0.0133469\\
0.7851	-0.0138029\\
0.7877	-0.0142711\\
0.7903	-0.0147331\\
0.7928	-0.0151714\\
0.7954	-0.015621\\
0.798	-0.0160644\\
0.8005	-0.0164848\\
0.8031	-0.0169161\\
0.8057	-0.0173413\\
0.8082	-0.0177443\\
0.8108	-0.0181572\\
0.8133	-0.0185483\\
0.8159	-0.018949\\
0.8185	-0.0193438\\
0.821	-0.0197177\\
0.8236	-0.0201006\\
0.8262	-0.0204775\\
0.8287	-0.0208342\\
0.8313	-0.0211993\\
0.8339	-0.0215585\\
0.8364	-0.0218985\\
0.839	-0.0222466\\
0.8415	-0.0225758\\
0.8441	-0.0229126\\
0.8467	-0.0232437\\
0.8492	-0.0235568\\
0.8518	-0.0238773\\
0.8544	-0.0241924\\
0.8569	-0.0244903\\
0.8595	-0.0247951\\
0.8621	-0.0250945\\
0.8646	-0.0253777\\
0.8672	-0.0256673\\
0.8697	-0.0259412\\
0.8723	-0.0262214\\
0.8749	-0.0264968\\
0.8774	-0.0267572\\
0.88	-0.0270235\\
0.8826	-0.0272853\\
0.8851	-0.027533\\
0.8877	-0.0277866\\
0.8903	-0.0280359\\
0.8928	-0.0282718\\
0.8954	-0.0285132\\
0.8979	-0.0287417\\
0.9005	-0.0289756\\
0.9031	-0.029206\\
0.9056	-0.0294242\\
0.9082	-0.0296478\\
0.9108	-0.0298679\\
0.9133	-0.0300766\\
0.9159	-0.0302905\\
0.9185	-0.0305015\\
0.921	-0.0307016\\
0.9236	-0.0309071\\
0.9262	-0.0311099\\
0.9287	-0.0313024\\
0.9313	-0.0315002\\
0.9338	-0.0316882\\
0.9364	-0.0318815\\
0.939	-0.0320728\\
0.9415	-0.0322549\\
0.9441	-0.0324422\\
0.9467	-0.0326278\\
0.9492	-0.0328046\\
0.9518	-0.032987\\
0.9544	-0.0331679\\
0.9569	-0.0333406\\
0.9595	-0.033519\\
0.962	-0.0336893\\
0.9646	-0.0338654\\
0.9672	-0.0340405\\
0.9697	-0.0342081\\
0.9723	-0.0343817\\
0.9749	-0.0345547\\
0.9774	-0.0347204\\
0.98	-0.0348922\\
0.9826	-0.0350637\\
0.9851	-0.0352283\\
0.9877	-0.0353993\\
0.9902	-0.0355637\\
0.9928	-0.0357345\\
0.9954	-0.0359054\\
0.9979	-0.0360699\\
1.0005	-0.0362411\\
1.0031	-0.0364127\\
1.0056	-0.036578\\
1.0082	-0.0367504\\
1.0108	-0.0369233\\
1.0133	-0.03709\\
1.0159	-0.0372641\\
1.0184	-0.0374322\\
1.021	-0.0376078\\
1.0236	-0.0377843\\
1.0261	-0.0379548\\
1.0287	-0.0381331\\
1.0313	-0.0383123\\
1.0338	-0.0384857\\
1.0364	-0.0386672\\
1.039	-0.0388498\\
1.0415	-0.0390265\\
1.0441	-0.0392115\\
1.0466	-0.0393906\\
1.0492	-0.0395782\\
1.0518	-0.0397671\\
1.0543	-0.0399501\\
1.0569	-0.0401418\\
1.0595	-0.0403349\\
1.062	-0.0405219\\
1.0646	-0.0407179\\
1.0672	-0.0409154\\
1.0697	-0.0411067\\
1.0723	-0.0413072\\
1.0748	-0.0415014\\
1.0774	-0.0417049\\
1.08	-0.0419099\\
1.0825	-0.0421084\\
1.0851	-0.0423165\\
1.0877	-0.042526\\
1.0902	-0.042729\\
1.0928	-0.0429415\\
1.0954	-0.0431556\\
1.0979	-0.0433627\\
1.1005	-0.0435796\\
1.1031	-0.043798\\
1.1056	-0.0440093\\
1.1082	-0.0442305\\
1.1107	-0.0444444\\
1.1133	-0.0446681\\
1.1159	-0.0448932\\
1.1184	-0.0451107\\
1.121	-0.0453382\\
1.1236	-0.0455669\\
1.1261	-0.0457879\\
1.1287	-0.0460187\\
1.1313	-0.0462506\\
1.1338	-0.0464745\\
1.1364	-0.0467084\\
1.1389	-0.046934\\
1.1415	-0.0471695\\
1.1441	-0.0474058\\
1.1466	-0.0476336\\
1.1492	-0.0478712\\
1.1518	-0.0481094\\
1.1543	-0.048339\\
1.1569	-0.0485782\\
1.1595	-0.0488178\\
1.162	-0.0490485\\
1.1646	-0.0492886\\
1.1671	-0.0495197\\
1.1697	-0.0497601\\
1.1723	-0.0500007\\
1.1748	-0.0502319\\
1.1774	-0.0504723\\
1.18	-0.0507124\\
1.1825	-0.0509431\\
1.1851	-0.0511827\\
1.1877	-0.0514219\\
1.1902	-0.0516515\\
1.1928	-0.0518896\\
1.1953	-0.052118\\
1.1979	-0.0523548\\
1.2005	-0.0525908\\
1.203	-0.0528169\\
1.2056	-0.0530512\\
1.2082	-0.0532844\\
1.2107	-0.0535075\\
1.2133	-0.0537384\\
1.2159	-0.053968\\
1.2184	-0.0541876\\
1.221	-0.0544145\\
1.2235	-0.0546313\\
1.2261	-0.0548552\\
1.2287	-0.0550774\\
1.2312	-0.0552895\\
1.2338	-0.0555083\\
1.2364	-0.0557253\\
1.2389	-0.055932\\
1.2415	-0.0561451\\
1.2441	-0.0563561\\
1.2466	-0.056557\\
1.2492	-0.0567637\\
1.2518	-0.0569682\\
1.2543	-0.0571626\\
1.2569	-0.0573624\\
1.2594	-0.0575522\\
1.262	-0.0577472\\
1.2646	-0.0579396\\
1.2671	-0.0581221\\
1.2697	-0.0583093\\
1.2723	-0.0584937\\
1.2748	-0.0586685\\
1.2774	-0.0588475\\
1.28	-0.0590236\\
1.2825	-0.0591901\\
1.2851	-0.0593605\\
1.2876	-0.0595215\\
1.2902	-0.0596859\\
1.2928	-0.0598473\\
1.2953	-0.0599995\\
1.2979	-0.0601548\\
1.3005	-0.0603069\\
1.303	-0.0604501\\
1.3056	-0.060596\\
1.3082	-0.0607386\\
1.3107	-0.0608726\\
1.3133	-0.0610088\\
1.3158	-0.0611366\\
1.3184	-0.0612663\\
1.321	-0.0613927\\
1.3235	-0.0615111\\
1.3261	-0.0616309\\
1.3287	-0.0617474\\
1.3312	-0.0618563\\
1.3338	-0.0619662\\
1.3364	-0.0620728\\
1.3389	-0.0621722\\
1.3415	-0.0622722\\
1.344	-0.0623652\\
1.3466	-0.0624587\\
1.3492	-0.0625489\\
1.3517	-0.0626325\\
1.3543	-0.0627162\\
1.3569	-0.0627966\\
1.3594	-0.0628709\\
1.362	-0.062945\\
1.3646	-0.0630159\\
1.3671	-0.063081\\
1.3697	-0.0631456\\
1.3722	-0.0632048\\
1.3748	-0.0632633\\
1.3774	-0.0633188\\
1.3799	-0.0633692\\
1.3825	-0.0634187\\
1.3851	-0.0634652\\
1.3876	-0.0635072\\
1.3902	-0.063548\\
1.3928	-0.0635859\\
1.3953	-0.0636198\\
1.3979	-0.0636522\\
1.4005	-0.0636819\\
1.403	-0.063708\\
1.4056	-0.0637325\\
1.4081	-0.0637536\\
1.4107	-0.063773\\
1.4133	-0.06379\\
1.4158	-0.063804\\
1.4184	-0.0638161\\
1.421	-0.063826\\
1.4235	-0.0638333\\
1.4261	-0.0638387\\
1.4287	-0.063842\\
1.4312	-0.0638431\\
1.4338	-0.0638422\\
1.4363	-0.0638394\\
1.4389	-0.0638346\\
1.4415	-0.0638279\\
1.444	-0.0638197\\
1.4466	-0.0638095\\
1.4492	-0.0637975\\
1.4517	-0.0637844\\
1.4543	-0.0637692\\
1.4569	-0.0637524\\
1.4594	-0.0637349\\
1.462	-0.0637153\\
1.4645	-0.0636951\\
1.4671	-0.0636728\\
1.4697	-0.0636493\\
1.4722	-0.0636255\\
1.4748	-0.0635997\\
1.4774	-0.0635728\\
1.4799	-0.063546\\
1.4825	-0.0635171\\
1.4851	-0.0634873\\
1.4876	-0.0634579\\
1.4902	-0.0634266\\
1.4927	-0.0633957\\
1.4953	-0.0633629\\
1.4979	-0.0633295\\
1.5004	-0.0632969\\
1.503	-0.0632624\\
1.5056	-0.0632275\\
1.5081	-0.0631935\\
1.5107	-0.0631578\\
1.5133	-0.0631218\\
1.5158	-0.0630869\\
1.5184	-0.0630504\\
1.5209	-0.0630152\\
1.5235	-0.0629785\\
1.5261	-0.0629416\\
1.5286	-0.0629062\\
1.5312	-0.0628694\\
1.5338	-0.0628326\\
1.5363	-0.0627973\\
1.5389	-0.0627608\\
1.5415	-0.0627244\\
1.544	-0.0626897\\
1.5466	-0.0626538\\
1.5491	-0.0626196\\
1.5517	-0.0625843\\
1.5543	-0.0625494\\
1.5568	-0.0625162\\
1.5594	-0.0624821\\
1.562	-0.0624484\\
1.5645	-0.0624165\\
1.5671	-0.0623839\\
1.5697	-0.0623517\\
1.5722	-0.0623214\\
1.5748	-0.0622904\\
1.5774	-0.06226\\
1.5799	-0.0622313\\
1.5825	-0.0622022\\
1.585	-0.0621748\\
1.5876	-0.062147\\
1.5902	-0.0621199\\
1.5927	-0.0620945\\
1.5953	-0.0620688\\
1.5979	-0.0620438\\
1.6004	-0.0620205\\
1.603	-0.0619971\\
1.6056	-0.0619743\\
1.6081	-0.0619532\\
1.6107	-0.0619321\\
1.6132	-0.0619124\\
1.6158	-0.0618928\\
1.6184	-0.0618739\\
1.6209	-0.0618565\\
1.6235	-0.0618392\\
1.6261	-0.0618227\\
1.6286	-0.0618076\\
1.6312	-0.0617926\\
1.6338	-0.0617784\\
1.6363	-0.0617655\\
1.6389	-0.0617528\\
1.6414	-0.0617413\\
1.644	-0.0617301\\
1.6466	-0.0617196\\
1.6491	-0.0617102\\
1.6517	-0.0617011\\
1.6543	-0.0616927\\
1.6568	-0.0616853\\
1.6594	-0.0616783\\
1.662	-0.0616719\\
1.6645	-0.0616663\\
1.6671	-0.0616612\\
1.6696	-0.0616568\\
1.6722	-0.0616528\\
1.6748	-0.0616494\\
1.6773	-0.0616467\\
1.6799	-0.0616443\\
1.6825	-0.0616425\\
1.685	-0.0616412\\
1.6876	-0.0616403\\
1.6902	-0.0616398\\
1.6927	-0.0616397\\
1.6953	-0.0616401\\
1.6978	-0.0616407\\
1.7004	-0.0616417\\
1.703	-0.0616431\\
1.7055	-0.0616447\\
1.7081	-0.0616466\\
1.7107	-0.0616487\\
1.7132	-0.061651\\
1.7158	-0.0616536\\
1.7184	-0.0616563\\
1.7209	-0.061659\\
1.7235	-0.061662\\
1.7261	-0.061665\\
1.7286	-0.061668\\
1.7312	-0.0616712\\
1.7337	-0.0616742\\
1.7363	-0.0616773\\
1.7389	-0.0616804\\
1.7414	-0.0616832\\
1.744	-0.0616861\\
1.7466	-0.0616889\\
1.7491	-0.0616914\\
1.7517	-0.0616938\\
1.7543	-0.061696\\
1.7568	-0.0616979\\
1.7594	-0.0616996\\
1.7619	-0.061701\\
1.7645	-0.0617021\\
1.7671	-0.0617028\\
1.7696	-0.0617032\\
1.7722	-0.0617032\\
1.7748	-0.0617028\\
1.7773	-0.061702\\
1.7799	-0.0617006\\
1.7825	-0.0616988\\
1.785	-0.0616966\\
1.7876	-0.0616938\\
1.7901	-0.0616905\\
1.7927	-0.0616865\\
1.7953	-0.0616819\\
1.7978	-0.0616769\\
1.8004	-0.0616711\\
1.803	-0.0616645\\
1.8055	-0.0616576\\
1.8081	-0.0616497\\
1.8107	-0.0616411\\
1.8132	-0.0616321\\
1.8158	-0.0616219\\
1.8183	-0.0616114\\
1.8209	-0.0615998\\
1.8235	-0.0615873\\
1.826	-0.0615744\\
1.8286	-0.0615603\\
1.8312	-0.0615452\\
1.8337	-0.0615299\\
1.8363	-0.0615132\\
1.8389	-0.0614955\\
1.8414	-0.0614776\\
1.844	-0.0614581\\
1.8465	-0.0614385\\
1.8491	-0.0614171\\
1.8517	-0.0613948\\
1.8542	-0.0613724\\
1.8568	-0.0613481\\
1.8594	-0.0613229\\
1.8619	-0.0612977\\
1.8645	-0.0612705\\
1.8671	-0.0612423\\
1.8696	-0.0612143\\
1.8722	-0.0611841\\
1.8747	-0.0611541\\
1.8773	-0.0611219\\
1.8799	-0.0610887\\
1.8824	-0.0610558\\
1.885	-0.0610206\\
1.8876	-0.0609844\\
1.8901	-0.0609486\\
1.8927	-0.0609104\\
1.8953	-0.0608712\\
1.8978	-0.0608325\\
1.9004	-0.0607913\\
1.903	-0.060749\\
1.9055	-0.0607075\\
1.9081	-0.0606633\\
1.9106	-0.06062\\
1.9132	-0.0605739\\
1.9158	-0.0605269\\
1.9183	-0.0604807\\
1.9209	-0.0604319\\
1.9235	-0.060382\\
1.926	-0.0603333\\
1.9286	-0.0602817\\
1.9312	-0.0602292\\
1.9337	-0.060178\\
1.9363	-0.0601238\\
1.9388	-0.060071\\
1.9414	-0.0600152\\
1.944	-0.0599586\\
1.9465	-0.0599035\\
1.9491	-0.0598454\\
1.9517	-0.0597866\\
1.9542	-0.0597294\\
1.9568	-0.0596692\\
1.9594	-0.0596083\\
1.9619	-0.0595492\\
1.9645	-0.0594871\\
1.967	-0.0594269\\
1.9696	-0.0593637\\
1.9722	-0.0592999\\
1.9747	-0.0592382\\
1.9773	-0.0591735\\
1.9799	-0.0591083\\
1.9824	-0.0590453\\
1.985	-0.0589794\\
1.9876	-0.0589131\\
1.9901	-0.0588491\\
1.9927	-0.0587822\\
1.9952	-0.0587177\\
1.9978	-0.0586504\\
2.0004	-0.058583\\
2.0029	-0.058518\\
2.0055	-0.0584503\\
2.0081	-0.0583825\\
2.0106	-0.0583173\\
2.0132	-0.0582496\\
2.0158	-0.0581819\\
2.0183	-0.0581169\\
2.0209	-0.0580494\\
2.0234	-0.0579847\\
2.026	-0.0579177\\
2.0286	-0.0578509\\
2.0311	-0.0577869\\
2.0337	-0.0577208\\
2.0363	-0.057655\\
2.0388	-0.0575922\\
2.0414	-0.0575274\\
2.044	-0.0574631\\
2.0465	-0.0574018\\
2.0491	-0.0573387\\
2.0517	-0.0572762\\
2.0542	-0.0572168\\
2.0568	-0.0571558\\
2.0593	-0.0570979\\
2.0619	-0.0570386\\
2.0645	-0.0569801\\
2.067	-0.0569248\\
2.0696	-0.0568682\\
2.0722	-0.0568127\\
2.0747	-0.0567603\\
2.0773	-0.0567069\\
2.0799	-0.0566548\\
2.0824	-0.0566058\\
2.085	-0.0565561\\
2.0875	-0.0565095\\
2.0901	-0.0564624\\
2.0927	-0.0564168\\
2.0952	-0.0563742\\
2.0978	-0.0563315\\
2.1004	-0.0562902\\
2.1029	-0.0562521\\
2.1055	-0.0562141\\
2.1081	-0.0561777\\
2.1106	-0.0561443\\
2.1132	-0.0561114\\
2.1157	-0.0560814\\
2.1183	-0.056052\\
2.1209	-0.0560245\\
2.1234	-0.0559999\\
2.126	-0.0559763\\
2.1286	-0.0559546\\
2.1311	-0.0559357\\
2.1337	-0.0559181\\
2.1363	-0.0559025\\
2.1388	-0.0558897\\
2.1414	-0.0558784\\
2.1439	-0.0558697\\
2.1465	-0.0558628\\
2.1491	-0.0558582\\
2.1516	-0.055856\\
2.1542	-0.0558559\\
2.1568	-0.0558583\\
2.1593	-0.0558628\\
2.1619	-0.0558698\\
2.1645	-0.0558793\\
2.167	-0.0558908\\
2.1696	-0.0559052\\
2.1721	-0.0559214\\
2.1747	-0.0559407\\
2.1773	-0.0559625\\
2.1798	-0.055986\\
2.1824	-0.0560129\\
2.185	-0.0560425\\
2.1875	-0.0560733\\
2.1901	-0.056108\\
2.1927	-0.0561453\\
2.1952	-0.0561837\\
2.1978	-0.0562262\\
2.2004	-0.0562713\\
2.2029	-0.0563173\\
2.2055	-0.0563676\\
2.208	-0.0564186\\
2.2106	-0.0564742\\
2.2132	-0.0565325\\
2.2157	-0.0565911\\
2.2183	-0.0566546\\
2.2209	-0.0567208\\
2.2234	-0.0567869\\
2.226	-0.0568583\\
2.2286	-0.0569323\\
2.2311	-0.0570059\\
2.2337	-0.057085\\
2.2362	-0.0571636\\
2.2388	-0.0572477\\
2.2414	-0.0573345\\
2.2439	-0.0574203\\
2.2465	-0.057512\\
2.2491	-0.0576062\\
2.2516	-0.0576991\\
2.2542	-0.0577981\\
2.2568	-0.0578996\\
2.2593	-0.0579994\\
2.2619	-0.0581055\\
2.2644	-0.0582097\\
2.267	-0.0583204\\
2.2696	-0.0584333\\
2.2721	-0.058544\\
2.2747	-0.0586612\\
2.2773	-0.0587806\\
2.2798	-0.0588974\\
2.2824	-0.0590209\\
2.285	-0.0591465\\
2.2875	-0.0592691\\
2.2901	-0.0593985\\
2.2926	-0.0595247\\
2.2952	-0.0596578\\
2.2978	-0.0597927\\
2.3003	-0.059924\\
2.3029	-0.0600623\\
2.3055	-0.0602022\\
2.308	-0.0603383\\
2.3106	-0.0604813\\
2.3132	-0.0606258\\
2.3157	-0.0607661\\
2.3183	-0.0609133\\
2.3208	-0.0610562\\
2.3234	-0.061206\\
2.326	-0.061357\\
2.3285	-0.0615032\\
2.3311	-0.0616564\\
2.3337	-0.0618106\\
2.3362	-0.0619597\\
2.3388	-0.0621157\\
2.3414	-0.0622725\\
2.3439	-0.062424\\
2.3465	-0.0625822\\
2.349	-0.062735\\
2.3516	-0.0628944\\
2.3542	-0.0630542\\
2.3567	-0.0632084\\
2.3593	-0.063369\\
2.3619	-0.06353\\
2.3644	-0.063685\\
2.367	-0.0638463\\
2.3696	-0.0640077\\
2.3721	-0.0641629\\
2.3747	-0.0643242\\
2.3773	-0.0644854\\
2.3798	-0.0646402\\
2.3824	-0.064801\\
2.3849	-0.0649552\\
2.3875	-0.0651152\\
2.3901	-0.0652747\\
2.3926	-0.0654276\\
2.3952	-0.065586\\
2.3978	-0.0657436\\
2.4003	-0.0658945\\
2.4029	-0.0660506\\
2.4055	-0.0662058\\
2.408	-0.0663541\\
2.4106	-0.0665073\\
2.4131	-0.0666536\\
2.4157	-0.0668046\\
2.4183	-0.0669543\\
2.4208	-0.067097\\
2.4234	-0.0672441\\
2.426	-0.0673897\\
2.4285	-0.0675283\\
2.4311	-0.0676709\\
2.4337	-0.0678118\\
2.4362	-0.0679457\\
2.4388	-0.0680832\\
2.4413	-0.0682137\\
2.4439	-0.0683476\\
2.4465	-0.0684795\\
2.449	-0.0686044\\
2.4516	-0.0687322\\
2.4542	-0.068858\\
2.4567	-0.0689768\\
2.4593	-0.0690982\\
2.4619	-0.0692172\\
2.4644	-0.0693295\\
2.467	-0.0694439\\
2.4695	-0.0695515\\
2.4721	-0.069661\\
2.4747	-0.0697679\\
2.4772	-0.0698683\\
2.4798	-0.06997\\
2.4824	-0.070069\\
2.4849	-0.0701616\\
2.4875	-0.0702552\\
2.4901	-0.0703459\\
2.4926	-0.0704304\\
2.4952	-0.0705154\\
2.4977	-0.0705943\\
2.5003	-0.0706735\\
2.5029	-0.0707496\\
2.5054	-0.0708198\\
2.508	-0.0708899\\
2.5106	-0.0709567\\
2.5131	-0.071018\\
2.5157	-0.0710787\\
2.5183	-0.071136\\
2.5208	-0.0711881\\
2.5234	-0.0712391\\
2.526	-0.0712867\\
2.5285	-0.0713294\\
2.5311	-0.0713705\\
2.5336	-0.0714068\\
2.5362	-0.0714412\\
2.5388	-0.0714723\\
2.5413	-0.0714989\\
2.5439	-0.0715232\\
2.5465	-0.071544\\
2.549	-0.0715608\\
2.5516	-0.0715749\\
2.5542	-0.0715855\\
2.5567	-0.0715924\\
2.5593	-0.0715961\\
2.5618	-0.0715964\\
2.5644	-0.0715933\\
2.567	-0.0715867\\
2.5695	-0.0715771\\
2.5721	-0.0715636\\
2.5747	-0.0715467\\
2.5772	-0.0715272\\
2.5798	-0.0715034\\
2.5824	-0.0714763\\
2.5849	-0.0714469\\
2.5875	-0.0714129\\
2.59	-0.071377\\
2.5926	-0.0713364\\
2.5952	-0.0712924\\
2.5977	-0.0712468\\
2.6003	-0.0711962\\
2.6029	-0.0711422\\
2.6054	-0.0710872\\
2.608	-0.0710267\\
2.6106	-0.070963\\
2.6131	-0.0708987\\
2.6157	-0.0708287\\
2.6182	-0.0707584\\
2.6208	-0.0706821\\
2.6234	-0.0706028\\
2.6259	-0.0705235\\
2.6285	-0.0704382\\
2.6311	-0.0703498\\
2.6336	-0.070262\\
2.6362	-0.0701678\\
2.6388	-0.0700707\\
2.6413	-0.0699746\\
2.6439	-0.0698719\\
2.6464	-0.0697706\\
2.649	-0.0696625\\
2.6516	-0.0695517\\
2.6541	-0.0694427\\
2.6567	-0.0693268\\
2.6593	-0.0692084\\
2.6618	-0.0690921\\
2.6644	-0.0689688\\
2.667	-0.0688432\\
2.6695	-0.0687201\\
2.6721	-0.0685899\\
2.6746	-0.0684625\\
2.6772	-0.068328\\
2.6798	-0.0681913\\
2.6823	-0.0680579\\
2.6849	-0.0679172\\
2.6875	-0.0677746\\
2.69	-0.0676357\\
2.6926	-0.0674894\\
2.6952	-0.0673413\\
2.6977	-0.0671974\\
2.7003	-0.067046\\
2.7029	-0.0668931\\
2.7054	-0.0667447\\
2.708	-0.0665888\\
2.7105	-0.0664377\\
2.7131	-0.0662792\\
2.7157	-0.0661195\\
2.7182	-0.0659648\\
2.7208	-0.0658028\\
2.7234	-0.0656397\\
2.7259	-0.065482\\
2.7285	-0.0653171\\
2.7311	-0.0651514\\
2.7336	-0.0649913\\
2.7362	-0.0648241\\
2.7387	-0.0646627\\
2.7413	-0.0644943\\
2.7439	-0.0643254\\
2.7464	-0.0641627\\
2.749	-0.063993\\
2.7516	-0.0638231\\
2.7541	-0.0636595\\
2.7567	-0.0634893\\
2.7593	-0.063319\\
2.7618	-0.0631553\\
2.7644	-0.062985\\
2.7669	-0.0628215\\
2.7695	-0.0626517\\
2.7721	-0.0624821\\
2.7746	-0.0623195\\
2.7772	-0.0621507\\
2.7798	-0.0619825\\
2.7823	-0.0618213\\
2.7849	-0.0616543\\
2.7875	-0.0614881\\
2.79	-0.061329\\
2.7926	-0.0611644\\
2.7951	-0.061007\\
2.7977	-0.0608443\\
2.8003	-0.0606827\\
2.8028	-0.0605284\\
2.8054	-0.0603692\\
2.808	-0.0602112\\
2.8105	-0.0600606\\
2.8131	-0.0599054\\
2.8157	-0.0597517\\
2.8182	-0.0596054\\
2.8208	-0.0594548\\
2.8233	-0.0593116\\
2.8259	-0.0591644\\
2.8285	-0.0590191\\
2.831	-0.0588811\\
2.8336	-0.0587396\\
2.8362	-0.0586\\
2.8387	-0.0584678\\
2.8413	-0.0583324\\
2.8439	-0.0581992\\
2.8464	-0.0580733\\
2.849	-0.0579446\\
2.8516	-0.0578183\\
2.8541	-0.0576991\\
2.8567	-0.0575776\\
2.8592	-0.0574632\\
2.8618	-0.0573468\\
2.8644	-0.057233\\
2.8669	-0.0571261\\
2.8695	-0.0570176\\
2.8721	-0.056912\\
2.8746	-0.0568131\\
2.8772	-0.056713\\
2.8798	-0.056616\\
2.8823	-0.0565254\\
2.8849	-0.0564342\\
2.8874	-0.0563494\\
2.89	-0.0562643\\
2.8926	-0.0561823\\
2.8951	-0.0561064\\
2.8977	-0.0560307\\
2.9003	-0.0559582\\
2.9028	-0.0558916\\
2.9054	-0.0558256\\
2.908	-0.0557629\\
2.9105	-0.0557058\\
2.9131	-0.0556498\\
2.9156	-0.0555992\\
2.9182	-0.0555499\\
2.9208	-0.055504\\
2.9233	-0.0554633\\
2.9259	-0.0554243\\
2.9285	-0.0553889\\
2.931	-0.0553581\\
2.9336	-0.0553296\\
2.9362	-0.0553047\\
2.9387	-0.0552841\\
2.9413	-0.0552662\\
2.9438	-0.0552524\\
2.9464	-0.0552416\\
2.949	-0.0552343\\
2.9515	-0.0552307\\
2.9541	-0.0552306\\
2.9567	-0.055234\\
2.9592	-0.0552406\\
2.9618	-0.0552511\\
2.9644	-0.0552651\\
2.9669	-0.0552819\\
2.9695	-0.055303\\
2.972	-0.0553265\\
2.9746	-0.0553545\\
2.9772	-0.0553859\\
2.9797	-0.0554195\\
2.9823	-0.0554578\\
2.9849	-0.0554996\\
2.9874	-0.055543\\
2.99	-0.0555914\\
2.9926	-0.0556433\\
2.9951	-0.0556963\\
2.9977	-0.0557546\\
3.0003	-0.0558163\\
3.0028	-0.0558787\\
3.0054	-0.0559466\\
3.0079	-0.056015\\
3.0105	-0.0560892\\
3.0131	-0.0561665\\
3.0156	-0.0562437\\
3.0182	-0.0563269\\
3.0208	-0.0564131\\
3.0233	-0.0564987\\
3.0259	-0.0565906\\
3.0285	-0.0566853\\
3.031	-0.0567789\\
3.0336	-0.0568789\\
3.0361	-0.0569776\\
3.0387	-0.0570828\\
3.0413	-0.0571905\\
3.0438	-0.0572964\\
3.0464	-0.0574089\\
3.049	-0.0575238\\
3.0515	-0.0576363\\
3.0541	-0.0577556\\
3.0567	-0.0578771\\
3.0592	-0.0579958\\
3.0618	-0.0581213\\
3.0643	-0.0582437\\
3.0669	-0.058373\\
3.0695	-0.058504\\
3.072	-0.0586316\\
3.0746	-0.058766\\
3.0772	-0.0589019\\
3.0797	-0.059034\\
3.0823	-0.0591728\\
3.0849	-0.059313\\
3.0874	-0.0594489\\
3.09	-0.0595915\\
3.0925	-0.0597297\\
3.0951	-0.0598743\\
3.0977	-0.0600199\\
3.1002	-0.0601608\\
3.1028	-0.060308\\
3.1054	-0.0604559\\
3.1079	-0.0605988\\
3.1105	-0.0607478\\
3.1131	-0.0608973\\
3.1156	-0.0610414\\
3.1182	-0.0611915\\
3.1207	-0.061336\\
3.1233	-0.0614864\\
3.1259	-0.0616369\\
3.1284	-0.0617814\\
3.131	-0.0619316\\
3.1336	-0.0620816\\
3.1361	-0.0622254\\
3.1387	-0.0623746\\
3.1413	-0.0625234\\
3.1438	-0.0626658\\
3.1464	-0.0628133\\
3.1489	-0.0629543\\
3.1515	-0.0631002\\
3.1541	-0.0632452\\
3.1566	-0.0633836\\
3.1592	-0.0635265\\
3.1618	-0.0636683\\
3.1643	-0.0638033\\
3.1669	-0.0639425\\
3.1695	-0.0640802\\
3.172	-0.0642112\\
3.1746	-0.0643459\\
3.1772	-0.0644789\\
3.1797	-0.0646051\\
3.1823	-0.0647346\\
3.1848	-0.0648573\\
3.1874	-0.064983\\
3.19	-0.0651066\\
3.1925	-0.0652235\\
3.1951	-0.0653428\\
3.1977	-0.0654599\\
3.2002	-0.0655703\\
3.2028	-0.0656827\\
3.2054	-0.0657927\\
3.2079	-0.065896\\
3.2105	-0.0660008\\
3.213	-0.0660991\\
3.2156	-0.0661987\\
3.2182	-0.0662954\\
3.2207	-0.0663857\\
3.2233	-0.0664768\\
3.2259	-0.0665648\\
3.2284	-0.0666466\\
3.231	-0.0667286\\
3.2336	-0.0668075\\
3.2361	-0.0668803\\
3.2387	-0.0669529\\
3.2412	-0.0670195\\
3.2438	-0.0670855\\
3.2464	-0.0671481\\
3.2489	-0.0672051\\
3.2515	-0.0672609\\
3.2541	-0.0673132\\
3.2566	-0.0673602\\
3.2592	-0.0674055\\
3.2618	-0.0674471\\
3.2643	-0.0674837\\
3.2669	-0.0675182\\
3.2694	-0.0675478\\
3.272	-0.0675749\\
3.2746	-0.0675982\\
3.2771	-0.067617\\
3.2797	-0.0676329\\
3.2823	-0.0676449\\
3.2848	-0.0676528\\
3.2874	-0.0676571\\
3.29	-0.0676576\\
3.2925	-0.0676544\\
3.2951	-0.0676472\\
3.2976	-0.0676366\\
3.3002	-0.0676217\\
3.3028	-0.0676028\\
3.3053	-0.0675809\\
3.3079	-0.0675543\\
3.3105	-0.0675237\\
3.313	-0.0674906\\
3.3156	-0.0674523\\
3.3182	-0.06741\\
3.3207	-0.0673656\\
3.3233	-0.0673156\\
3.3259	-0.0672617\\
3.3284	-0.0672062\\
3.331	-0.0671446\\
3.3335	-0.0670818\\
3.3361	-0.0670126\\
3.3387	-0.0669396\\
3.3412	-0.0668658\\
3.3438	-0.0667853\\
3.3464	-0.066701\\
3.3489	-0.0666164\\
3.3515	-0.0665248\\
3.3541	-0.0664295\\
3.3566	-0.0663344\\
3.3592	-0.066232\\
3.3617	-0.06613\\
3.3643	-0.0660205\\
3.3669	-0.0659075\\
3.3694	-0.0657955\\
3.372	-0.0656757\\
3.3746	-0.0655524\\
3.3771	-0.0654307\\
3.3797	-0.0653008\\
3.3823	-0.0651677\\
3.3848	-0.0650366\\
3.3874	-0.0648971\\
3.3899	-0.0647601\\
3.3925	-0.0646145\\
3.3951	-0.0644658\\
3.3976	-0.0643201\\
3.4002	-0.0641656\\
3.4028	-0.0640083\\
3.4053	-0.0638543\\
3.4079	-0.0636914\\
3.4105	-0.0635258\\
3.413	-0.0633641\\
3.4156	-0.0631934\\
3.4181	-0.0630268\\
3.4207	-0.0628511\\
3.4233	-0.062673\\
3.4258	-0.0624995\\
3.4284	-0.0623168\\
3.431	-0.0621319\\
3.4335	-0.061952\\
3.4361	-0.0617629\\
3.4387	-0.0615718\\
3.4412	-0.0613861\\
3.4438	-0.0611912\\
3.4463	-0.061002\\
3.4489	-0.0608035\\
3.4515	-0.0606033\\
3.454	-0.0604093\\
3.4566	-0.060206\\
3.4592	-0.0600012\\
3.4617	-0.0598029\\
3.4643	-0.0595954\\
3.4669	-0.0593867\\
3.4694	-0.0591848\\
3.472	-0.0589737\\
3.4745	-0.0587698\\
3.4771	-0.0585568\\
3.4797	-0.0583428\\
3.4822	-0.0581363\\
3.4848	-0.0579207\\
3.4874	-0.0577045\\
3.4899	-0.0574961\\
3.4925	-0.0572787\\
3.4951	-0.0570609\\
3.4976	-0.0568511\\
3.5002	-0.0566326\\
3.5028	-0.0564139\\
3.5053	-0.0562034\\
3.5079	-0.0559844\\
3.5104	-0.0557738\\
3.513	-0.0555549\\
3.5156	-0.0553361\\
3.5181	-0.0551259\\
3.5207	-0.0549075\\
3.5233	-0.0546895\\
3.5258	-0.0544803\\
3.5284	-0.0542632\\
3.531	-0.0540467\\
3.5335	-0.0538391\\
3.5361	-0.053624\\
3.5386	-0.0534178\\
3.5412	-0.0532042\\
3.5438	-0.0529916\\
3.5463	-0.0527881\\
3.5489	-0.0525775\\
3.5515	-0.0523681\\
3.554	-0.0521678\\
3.5566	-0.0519608\\
3.5592	-0.0517552\\
3.5617	-0.0515588\\
3.5643	-0.0513559\\
3.5668	-0.0511624\\
3.5694	-0.0509627\\
3.572	-0.0507647\\
3.5745	-0.050576\\
3.5771	-0.0503814\\
3.5797	-0.0501888\\
3.5822	-0.0500054\\
3.5848	-0.0498166\\
3.5874	-0.0496299\\
3.5899	-0.0494524\\
3.5925	-0.0492699\\
3.595	-0.0490966\\
3.5976	-0.0489186\\
3.6002	-0.0487429\\
3.6027	-0.0485763\\
3.6053	-0.0484055\\
3.6079	-0.0482371\\
3.6104	-0.0480777\\
3.613	-0.0479145\\
3.6156	-0.047754\\
3.6181	-0.0476022\\
3.6207	-0.0474472\\
3.6232	-0.0473007\\
3.6258	-0.0471512\\
3.6284	-0.0470047\\
3.6309	-0.0468666\\
3.6335	-0.046726\\
3.6361	-0.0465884\\
3.6386	-0.0464591\\
3.6412	-0.0463277\\
3.6438	-0.0461995\\
3.6463	-0.0460793\\
3.6489	-0.0459574\\
3.6515	-0.0458389\\
3.654	-0.0457282\\
3.6566	-0.0456163\\
3.6591	-0.0455119\\
3.6617	-0.0454068\\
3.6643	-0.0453051\\
3.6668	-0.0452107\\
3.6694	-0.0451159\\
3.672	-0.0450248\\
3.6745	-0.0449405\\
3.6771	-0.0448564\\
3.6797	-0.044776\\
3.6822	-0.0447021\\
3.6848	-0.0446289\\
3.6873	-0.044562\\
3.6899	-0.0444961\\
3.6925	-0.044434\\
3.695	-0.0443777\\
3.6976	-0.044323\\
3.7002	-0.044272\\
3.7027	-0.0442266\\
3.7053	-0.0441831\\
3.7079	-0.0441435\\
3.7104	-0.0441089\\
3.713	-0.0440768\\
3.7155	-0.0440495\\
3.7181	-0.0440249\\
3.7207	-0.0440042\\
3.7232	-0.0439879\\
3.7258	-0.0439747\\
3.7284	-0.0439653\\
3.7309	-0.04396\\
3.7335	-0.0439582\\
3.7361	-0.0439602\\
3.7386	-0.0439657\\
3.7412	-0.0439753\\
3.7437	-0.043988\\
3.7463	-0.044005\\
3.7489	-0.0440258\\
3.7514	-0.0440493\\
3.754	-0.0440774\\
3.7566	-0.0441093\\
3.7591	-0.0441434\\
3.7617	-0.0441826\\
3.7643	-0.0442254\\
3.7668	-0.04427\\
3.7694	-0.0443199\\
3.7719	-0.0443713\\
3.7745	-0.0444283\\
3.7771	-0.0444888\\
3.7796	-0.0445503\\
3.7822	-0.0446177\\
3.7848	-0.0446886\\
3.7873	-0.0447599\\
3.7899	-0.0448374\\
3.7925	-0.0449183\\
3.795	-0.0449991\\
3.7976	-0.0450864\\
3.8002	-0.045177\\
3.8027	-0.045267\\
3.8053	-0.0453638\\
3.8078	-0.0454598\\
3.8104	-0.0455626\\
3.813	-0.0456684\\
3.8155	-0.0457729\\
3.8181	-0.0458845\\
3.8207	-0.045999\\
3.8232	-0.0461117\\
3.8258	-0.0462317\\
3.8284	-0.0463543\\
3.8309	-0.0464748\\
3.8335	-0.0466026\\
3.836	-0.046728\\
3.8386	-0.0468608\\
3.8412	-0.046996\\
3.8437	-0.0471283\\
3.8463	-0.0472682\\
3.8489	-0.0474103\\
3.8514	-0.0475491\\
3.854	-0.0476955\\
3.8566	-0.047844\\
3.8591	-0.0479887\\
3.8617	-0.0481411\\
3.8642	-0.0482894\\
3.8668	-0.0484455\\
3.8694	-0.0486033\\
3.8719	-0.0487567\\
3.8745	-0.0489178\\
3.8771	-0.0490804\\
3.8796	-0.0492383\\
3.8822	-0.0494038\\
3.8848	-0.0495708\\
3.8873	-0.0497325\\
3.8899	-0.0499019\\
3.8924	-0.0500659\\
3.895	-0.0502375\\
3.8976	-0.0504101\\
3.9001	-0.050577\\
3.9027	-0.0507515\\
3.9053	-0.0509267\\
3.9078	-0.0510959\\
3.9104	-0.0512726\\
3.913	-0.0514498\\
3.9155	-0.0516207\\
3.9181	-0.051799\\
3.9206	-0.0519707\\
3.9232	-0.0521496\\
3.9258	-0.0523287\\
3.9283	-0.0525012\\
3.9309	-0.0526806\\
3.9335	-0.0528601\\
3.936	-0.0530326\\
3.9386	-0.0532119\\
3.9412	-0.0533911\\
3.9437	-0.0535631\\
3.9463	-0.0537418\\
3.9488	-0.0539132\\
3.9514	-0.054091\\
3.954	-0.0542683\\
3.9565	-0.0544382\\
3.9591	-0.0546143\\
3.9617	-0.0547897\\
3.9642	-0.0549576\\
3.9668	-0.0551314\\
3.9694	-0.0553042\\
3.9719	-0.0554696\\
3.9745	-0.0556405\\
3.9771	-0.0558103\\
3.9796	-0.0559725\\
3.9822	-0.0561399\\
3.9847	-0.0562998\\
3.9873	-0.0564647\\
3.9899	-0.0566282\\
3.9924	-0.056784\\
3.995	-0.0569446\\
3.9976	-0.0571036\\
4.0001	-0.0572549\\
4.0027	-0.0574107\\
4.0053	-0.0575647\\
4.0078	-0.0577112\\
4.0104	-0.0578617\\
4.0129	-0.0580046\\
4.0155	-0.0581513\\
4.0181	-0.058296\\
4.0206	-0.0584333\\
4.0232	-0.058574\\
4.0258	-0.0587126\\
4.0283	-0.0588438\\
4.0309	-0.0589781\\
4.0335	-0.0591101\\
4.036	-0.0592349\\
4.0386	-0.0593624\\
4.0411	-0.0594828\\
4.0437	-0.0596056\\
4.0463	-0.059726\\
4.0488	-0.0598394\\
4.0514	-0.0599549\\
4.054	-0.0600679\\
4.0565	-0.060174\\
4.0591	-0.0602819\\
4.0617	-0.0603872\\
4.0642	-0.0604859\\
4.0668	-0.0605859\\
4.0693	-0.0606795\\
4.0719	-0.0607742\\
4.0745	-0.0608662\\
4.077	-0.060952\\
4.0796	-0.0610384\\
4.0822	-0.0611221\\
4.0847	-0.0612\\
4.0873	-0.0612781\\
4.0899	-0.0613534\\
4.0924	-0.0614231\\
4.095	-0.0614928\\
4.0975	-0.0615571\\
4.1001	-0.0616211\\
4.1027	-0.0616822\\
4.1052	-0.0617382\\
4.1078	-0.0617937\\
4.1104	-0.0618462\\
4.1129	-0.0618939\\
4.1155	-0.0619408\\
4.1181	-0.0619847\\
4.1206	-0.0620241\\
4.1232	-0.0620623\\
4.1258	-0.0620977\\
4.1283	-0.0621289\\
4.1309	-0.0621585\\
4.1334	-0.0621843\\
4.136	-0.0622084\\
4.1386	-0.0622295\\
4.1411	-0.0622472\\
4.1437	-0.0622628\\
4.1463	-0.0622756\\
4.1488	-0.0622853\\
4.1514	-0.0622926\\
4.154	-0.0622972\\
4.1565	-0.062299\\
4.1591	-0.0622982\\
4.1616	-0.0622949\\
4.1642	-0.0622888\\
4.1668	-0.0622801\\
4.1693	-0.0622692\\
4.1719	-0.0622554\\
4.1745	-0.062239\\
4.177	-0.0622208\\
4.1796	-0.0621994\\
4.1822	-0.0621755\\
4.1847	-0.0621502\\
4.1873	-0.0621215\\
4.1898	-0.0620917\\
4.1924	-0.0620584\\
4.195	-0.0620227\\
4.1975	-0.0619862\\
4.2001	-0.0619461\\
4.2027	-0.0619038\\
4.2052	-0.061861\\
4.2078	-0.0618144\\
4.2104	-0.0617658\\
4.2129	-0.061717\\
4.2155	-0.0616643\\
4.218	-0.0616118\\
4.2206	-0.0615553\\
4.2232	-0.0614968\\
4.2257	-0.0614389\\
4.2283	-0.0613769\\
4.2309	-0.0613131\\
4.2334	-0.0612501\\
4.236	-0.061183\\
4.2386	-0.0611143\\
4.2411	-0.0610467\\
4.2437	-0.0609749\\
4.2462	-0.0609044\\
4.2488	-0.0608297\\
4.2514	-0.0607537\\
4.2539	-0.0606793\\
4.2565	-0.0606006\\
4.2591	-0.0605207\\
4.2616	-0.0604427\\
4.2642	-0.0603605\\
4.2668	-0.0602772\\
4.2693	-0.0601961\\
4.2719	-0.0601108\\
4.2744	-0.0600279\\
4.277	-0.0599408\\
4.2796	-0.0598529\\
4.2821	-0.0597675\\
4.2847	-0.0596781\\
4.2873	-0.0595879\\
4.2898	-0.0595006\\
4.2924	-0.0594093\\
4.295	-0.0593174\\
4.2975	-0.0592286\\
4.3001	-0.0591358\\
4.3027	-0.0590426\\
4.3052	-0.0589528\\
4.3078	-0.058859\\
4.3103	-0.0587686\\
4.3129	-0.0586745\\
4.3155	-0.0585802\\
4.318	-0.0584894\\
4.3206	-0.0583951\\
4.3232	-0.0583007\\
4.3257	-0.05821\\
4.3283	-0.0581159\\
4.3309	-0.0580219\\
4.3334	-0.0579318\\
4.336	-0.0578383\\
4.3385	-0.0577487\\
4.3411	-0.0576559\\
4.3437	-0.0575635\\
4.3462	-0.0574751\\
4.3488	-0.0573836\\
4.3514	-0.0572928\\
4.3539	-0.057206\\
4.3565	-0.0571163\\
4.3591	-0.0570274\\
4.3616	-0.0569426\\
4.3642	-0.0568551\\
4.3667	-0.0567719\\
4.3693	-0.0566861\\
4.3719	-0.0566013\\
4.3744	-0.0565206\\
4.377	-0.0564377\\
4.3796	-0.0563558\\
4.3821	-0.0562781\\
4.3847	-0.0561984\\
4.3873	-0.0561199\\
4.3898	-0.0560455\\
4.3924	-0.0559694\\
4.3949	-0.0558974\\
4.3975	-0.0558238\\
4.4001	-0.0557516\\
4.4026	-0.0556834\\
4.4052	-0.055614\\
4.4078	-0.0555459\\
4.4103	-0.055482\\
4.4129	-0.0554169\\
4.4155	-0.0553534\\
4.418	-0.0552938\\
4.4206	-0.0552334\\
4.4231	-0.0551769\\
4.4257	-0.0551198\\
4.4283	-0.0550644\\
4.4308	-0.0550127\\
4.4334	-0.0549607\\
4.436	-0.0549105\\
4.4385	-0.0548639\\
4.4411	-0.0548172\\
4.4437	-0.0547724\\
4.4462	-0.0547311\\
4.4488	-0.0546899\\
4.4514	-0.0546507\\
4.4539	-0.0546148\\
4.4565	-0.0545794\\
4.459	-0.0545472\\
4.4616	-0.0545156\\
4.4642	-0.0544861\\
4.4667	-0.0544595\\
4.4693	-0.0544339\\
4.4719	-0.0544104\\
4.4744	-0.0543896\\
4.477	-0.0543701\\
4.4796	-0.0543526\\
4.4821	-0.0543378\\
4.4847	-0.0543244\\
4.4872	-0.0543135\\
4.4898	-0.0543042\\
4.4924	-0.054297\\
4.4949	-0.0542921\\
4.4975	-0.054289\\
4.5001	-0.0542881\\
4.5026	-0.0542891\\
4.5052	-0.0542923\\
4.5078	-0.0542976\\
4.5103	-0.0543047\\
4.5129	-0.0543141\\
4.5154	-0.0543251\\
4.518	-0.0543386\\
4.5206	-0.0543542\\
4.5231	-0.0543712\\
4.5257	-0.0543909\\
4.5283	-0.0544126\\
4.5308	-0.0544355\\
4.5334	-0.0544612\\
4.536	-0.054489\\
4.5385	-0.0545176\\
4.5411	-0.0545494\\
4.5436	-0.0545817\\
4.5462	-0.0546174\\
4.5488	-0.0546549\\
4.5513	-0.0546929\\
4.5539	-0.0547342\\
4.5565	-0.0547775\\
4.559	-0.0548209\\
4.5616	-0.0548678\\
4.5642	-0.0549166\\
4.5667	-0.0549652\\
4.5693	-0.0550175\\
4.5718	-0.0550694\\
4.5744	-0.0551252\\
4.577	-0.0551826\\
4.5795	-0.0552395\\
4.5821	-0.0553003\\
4.5847	-0.0553626\\
4.5872	-0.0554241\\
4.5898	-0.0554897\\
4.5924	-0.0555567\\
4.5949	-0.0556226\\
4.5975	-0.0556926\\
4.6001	-0.055764\\
4.6026	-0.0558341\\
4.6052	-0.0559082\\
4.6077	-0.0559808\\
4.6103	-0.0560576\\
4.6129	-0.0561357\\
4.6154	-0.0562119\\
4.618	-0.0562924\\
4.6206	-0.056374\\
4.6231	-0.0564535\\
4.6257	-0.0565373\\
4.6283	-0.0566221\\
4.6308	-0.0567046\\
4.6334	-0.0567914\\
4.6359	-0.0568757\\
4.6385	-0.0569642\\
4.6411	-0.0570536\\
4.6436	-0.0571403\\
4.6462	-0.0572312\\
4.6488	-0.0573229\\
4.6513	-0.0574116\\
4.6539	-0.0575045\\
4.6565	-0.057598\\
4.659	-0.0576884\\
4.6616	-0.0577829\\
4.6641	-0.0578742\\
4.6667	-0.0579696\\
4.6693	-0.0580653\\
4.6718	-0.0581576\\
4.6744	-0.0582539\\
4.677	-0.0583504\\
4.6795	-0.0584434\\
4.6821	-0.0585402\\
4.6847	-0.0586372\\
4.6872	-0.0587304\\
4.6898	-0.0588273\\
4.6923	-0.0589205\\
4.6949	-0.0590174\\
4.6975	-0.059114\\
4.7	-0.0592068\\
4.7026	-0.0593031\\
4.7052	-0.0593991\\
4.7077	-0.0594911\\
4.7103	-0.0595864\\
4.7129	-0.0596813\\
4.7154	-0.0597722\\
4.718	-0.0598661\\
4.7205	-0.059956\\
4.7231	-0.0600488\\
4.7257	-0.060141\\
4.7282	-0.0602291\\
4.7308	-0.0603199\\
4.7334	-0.06041\\
4.7359	-0.0604959\\
4.7385	-0.0605843\\
4.7411	-0.0606719\\
4.7436	-0.0607552\\
4.7462	-0.0608409\\
4.7487	-0.0609224\\
4.7513	-0.061006\\
4.7539	-0.0610886\\
4.7564	-0.0611669\\
4.759	-0.0612473\\
4.7616	-0.0613264\\
4.7641	-0.0614013\\
4.7667	-0.0614779\\
4.7693	-0.0615532\\
4.7718	-0.0616243\\
4.7744	-0.0616969\\
4.777	-0.0617681\\
4.7795	-0.0618352\\
4.7821	-0.0619035\\
4.7846	-0.0619677\\
4.7872	-0.0620329\\
4.7898	-0.0620966\\
4.7923	-0.0621563\\
4.7949	-0.0622168\\
4.7975	-0.0622756\\
4.8	-0.0623305\\
4.8026	-0.0623859\\
4.8052	-0.0624396\\
4.8077	-0.0624895\\
4.8103	-0.0625396\\
4.8128	-0.0625861\\
4.8154	-0.0626326\\
4.818	-0.0626773\\
4.8205	-0.0627184\\
4.8231	-0.0627593\\
4.8257	-0.0627982\\
4.8282	-0.0628339\\
4.8308	-0.0628689\\
4.8334	-0.062902\\
4.8359	-0.062932\\
4.8385	-0.0629611\\
4.841	-0.0629872\\
4.8436	-0.0630123\\
4.8462	-0.0630353\\
4.8487	-0.0630555\\
4.8513	-0.0630745\\
4.8539	-0.0630913\\
4.8564	-0.0631055\\
4.859	-0.0631182\\
4.8616	-0.0631288\\
4.8641	-0.063137\\
4.8667	-0.0631434\\
4.8692	-0.0631475\\
4.8718	-0.0631497\\
4.8744	-0.0631498\\
4.8769	-0.0631478\\
4.8795	-0.0631436\\
4.8821	-0.0631373\\
4.8846	-0.0631292\\
4.8872	-0.0631187\\
4.8898	-0.063106\\
4.8923	-0.0630918\\
4.8949	-0.063075\\
4.8974	-0.0630567\\
4.9	-0.0630357\\
4.9026	-0.0630126\\
4.9051	-0.0629883\\
4.9077	-0.0629611\\
4.9103	-0.0629317\\
4.9128	-0.0629015\\
4.9154	-0.0628681\\
4.918	-0.0628327\\
4.9205	-0.0627967\\
4.9231	-0.0627572\\
4.9257	-0.0627158\\
4.9282	-0.062674\\
4.9308	-0.0626287\\
4.9333	-0.0625832\\
4.9359	-0.0625341\\
4.9385	-0.062483\\
4.941	-0.0624321\\
4.9436	-0.0623773\\
4.9462	-0.0623206\\
4.9487	-0.0622644\\
4.9513	-0.0622041\\
4.9539	-0.0621421\\
4.9564	-0.0620808\\
4.959	-0.0620153\\
4.9615	-0.0619508\\
4.9641	-0.0618819\\
4.9667	-0.0618115\\
4.9692	-0.0617421\\
4.9718	-0.0616685\\
4.9744	-0.0615932\\
4.9769	-0.0615194\\
4.9795	-0.0614412\\
4.9821	-0.0613614\\
4.9846	-0.0612833\\
4.9872	-0.0612008\\
4.9897	-0.06112\\
4.9923	-0.0610347\\
4.9949	-0.0609481\\
4.9974	-0.0608636\\
5	-0.0607744\\
};
\addplot [color=mycolor5, line width=1.0pt, forget plot]
  table[row sep=crcr]{%
-0.125	4.19133e-08\\
-0.1224	4.27836e-08\\
-0.1199	4.36203e-08\\
-0.1173	4.44904e-08\\
-0.1147	4.53605e-08\\
-0.1122	4.61971e-08\\
-0.1096	4.70671e-08\\
-0.1071	4.79038e-08\\
-0.1045	4.87738e-08\\
-0.1019	4.96439e-08\\
-0.0994	5.04805e-08\\
-0.0968	5.13506e-08\\
-0.0942	5.22207e-08\\
-0.0917	5.30573e-08\\
-0.0891	5.39274e-08\\
-0.0865	5.47974e-08\\
-0.084	5.56339e-08\\
-0.0814	5.65038e-08\\
-0.0789	5.73403e-08\\
-0.0763	5.82102e-08\\
-0.0737	5.908e-08\\
-0.0712	5.99164e-08\\
-0.0686	6.07863e-08\\
-0.066	6.16561e-08\\
-0.0635	6.24925e-08\\
-0.0609	6.33623e-08\\
-0.0583	6.42322e-08\\
-0.0558	6.50685e-08\\
-0.0532	6.59384e-08\\
-0.0507	6.67748e-08\\
-0.0481	6.76446e-08\\
-0.0455	6.85145e-08\\
-0.043	6.93507e-08\\
-0.0404	7.02203e-08\\
-0.0378	7.10899e-08\\
-0.0353	7.1926e-08\\
-0.0327	7.27956e-08\\
-0.0301	7.36652e-08\\
-0.0276	7.45014e-08\\
-0.025	7.5371e-08\\
-0.0224	7.62406e-08\\
-0.0199	7.70768e-08\\
-0.0173	7.79463e-08\\
-0.0148	7.87824e-08\\
-0.0122	7.96519e-08\\
-0.0096	8.05212e-08\\
-0.0071	8.13572e-08\\
-0.0045	8.22266e-08\\
-0.0019	8.30959e-08\\
0.0006	8.39319e-08\\
0.0032	7.52131e-05\\
0.0058	0.000406261\\
0.0083	0.000788369\\
0.0109	0.00125102\\
0.0134	0.00173532\\
0.016	0.00224438\\
0.0186	0.00273667\\
0.0211	0.00319107\\
0.0237	0.00365224\\
0.0263	0.00411003\\
0.0288	0.00454807\\
0.0314	0.00499639\\
0.034	0.00543153\\
0.0365	0.00583574\\
0.0391	0.00624381\\
0.0416	0.00662802\\
0.0442	0.00702126\\
0.0468	0.00740731\\
0.0493	0.00776964\\
0.0519	0.00813583\\
0.0545	0.00849182\\
0.057	0.00882648\\
0.0596	0.00916852\\
0.0622	0.00950493\\
0.0647	0.00982223\\
0.0673	0.0101448\\
0.0698	0.0104476\\
0.0724	0.0107563\\
0.075	0.0110604\\
0.0775	0.0113493\\
0.0801	0.011646\\
0.0827	0.0119381\\
0.0852	0.012214\\
0.0878	0.0124964\\
0.0904	0.0127755\\
0.0929	0.013042\\
0.0955	0.0133175\\
0.098	0.0135803\\
0.1006	0.0138505\\
0.1032	0.0141175\\
0.1057	0.014372\\
0.1083	0.0146357\\
0.1109	0.014899\\
0.1134	0.0151514\\
0.116	0.0154125\\
0.1186	0.0156714\\
0.1211	0.0159186\\
0.1237	0.0161748\\
0.1263	0.016431\\
0.1288	0.0166774\\
0.1314	0.0169331\\
0.1339	0.0171776\\
0.1365	0.0174302\\
0.1391	0.0176816\\
0.1416	0.0179231\\
0.1442	0.0181744\\
0.1468	0.0184254\\
0.1493	0.0186658\\
0.1519	0.018914\\
0.1545	0.0191607\\
0.157	0.019397\\
0.1596	0.0196425\\
0.1621	0.0198782\\
0.1647	0.0201224\\
0.1673	0.0203648\\
0.1698	0.020596\\
0.1724	0.0208348\\
0.175	0.0210726\\
0.1775	0.0213007\\
0.1801	0.0215369\\
0.1827	0.0217713\\
0.1852	0.0219946\\
0.1878	0.0222246\\
0.1903	0.0224442\\
0.1929	0.0226715\\
0.1955	0.0228977\\
0.198	0.0231137\\
0.2006	0.0233361\\
0.2032	0.023556\\
0.2057	0.0237654\\
0.2083	0.0239817\\
0.2109	0.0241969\\
0.2134	0.0244024\\
0.216	0.0246141\\
0.2185	0.0248155\\
0.2211	0.0250226\\
0.2237	0.025228\\
0.2262	0.0254243\\
0.2288	0.0256274\\
0.2314	0.0258288\\
0.2339	0.0260206\\
0.2365	0.026218\\
0.2391	0.0264136\\
0.2416	0.0266005\\
0.2442	0.026794\\
0.2467	0.0269791\\
0.2493	0.0271702\\
0.2519	0.0273595\\
0.2544	0.02754\\
0.257	0.0277267\\
0.2596	0.0279127\\
0.2621	0.0280911\\
0.2647	0.0282758\\
0.2673	0.0284594\\
0.2698	0.0286347\\
0.2724	0.0288161\\
0.2749	0.0289901\\
0.2775	0.029171\\
0.2801	0.0293517\\
0.2826	0.0295249\\
0.2852	0.0297043\\
0.2878	0.029883\\
0.2903	0.0300545\\
0.2929	0.0302331\\
0.2955	0.030412\\
0.298	0.0305839\\
0.3006	0.0307624\\
0.3032	0.0309405\\
0.3057	0.0311116\\
0.3083	0.0312899\\
0.3108	0.0314617\\
0.3134	0.0316409\\
0.316	0.0318201\\
0.3185	0.0319923\\
0.3211	0.0321713\\
0.3237	0.0323506\\
0.3262	0.0325235\\
0.3288	0.032704\\
0.3314	0.0328848\\
0.3339	0.0330587\\
0.3365	0.0332396\\
0.339	0.0334136\\
0.3416	0.0335951\\
0.3442	0.0337772\\
0.3467	0.0339528\\
0.3493	0.0341355\\
0.3519	0.0343183\\
0.3544	0.034494\\
0.357	0.034677\\
0.3596	0.0348604\\
0.3621	0.0350373\\
0.3647	0.0352214\\
0.3672	0.0353983\\
0.3698	0.0355822\\
0.3724	0.035766\\
0.3749	0.0359429\\
0.3775	0.0361271\\
0.3801	0.0363114\\
0.3826	0.0364885\\
0.3852	0.0366724\\
0.3878	0.0368558\\
0.3903	0.0370321\\
0.3929	0.0372153\\
0.3954	0.0373914\\
0.398	0.0375743\\
0.4006	0.0377566\\
0.4031	0.0379313\\
0.4057	0.0381124\\
0.4083	0.038293\\
0.4108	0.0384664\\
0.4134	0.0386462\\
0.416	0.0388253\\
0.4185	0.0389967\\
0.4211	0.0391739\\
0.4236	0.0393436\\
0.4262	0.0395193\\
0.4288	0.0396944\\
0.4313	0.0398618\\
0.4339	0.0400347\\
0.4365	0.0402064\\
0.439	0.0403704\\
0.4416	0.0405398\\
0.4442	0.0407083\\
0.4467	0.0408691\\
0.4493	0.0410351\\
0.4519	0.0411995\\
0.4544	0.0413561\\
0.457	0.0415176\\
0.4595	0.0416716\\
0.4621	0.0418304\\
0.4647	0.0419876\\
0.4672	0.0421372\\
0.4698	0.0422909\\
0.4724	0.0424428\\
0.4749	0.0425873\\
0.4775	0.0427359\\
0.4801	0.0428828\\
0.4826	0.0430222\\
0.4852	0.0431651\\
0.4877	0.0433005\\
0.4903	0.0434394\\
0.4929	0.0435763\\
0.4954	0.0437062\\
0.498	0.0438391\\
0.5006	0.0439698\\
0.5031	0.0440932\\
0.5057	0.0442193\\
0.5083	0.0443431\\
0.5108	0.0444601\\
0.5134	0.0445796\\
0.5159	0.0446921\\
0.5185	0.0448067\\
0.5211	0.0449186\\
0.5236	0.0450239\\
0.5262	0.045131\\
0.5288	0.0452356\\
0.5313	0.0453337\\
0.5339	0.0454331\\
0.5365	0.0455297\\
0.539	0.04562\\
0.5416	0.0457112\\
0.5441	0.0457963\\
0.5467	0.0458823\\
0.5493	0.0459653\\
0.5518	0.0460423\\
0.5544	0.0461195\\
0.557	0.0461937\\
0.5595	0.0462623\\
0.5621	0.0463309\\
0.5647	0.0463964\\
0.5672	0.0464565\\
0.5698	0.0465158\\
0.5723	0.0465699\\
0.5749	0.0466232\\
0.5775	0.0466734\\
0.58	0.0467187\\
0.5826	0.0467626\\
0.5852	0.0468033\\
0.5877	0.0468392\\
0.5903	0.0468734\\
0.5929	0.0469043\\
0.5954	0.046931\\
0.598	0.0469556\\
0.6006	0.0469767\\
0.6031	0.0469937\\
0.6057	0.0470081\\
0.6082	0.0470187\\
0.6108	0.0470265\\
0.6134	0.0470309\\
0.6159	0.0470319\\
0.6185	0.0470294\\
0.6211	0.0470234\\
0.6236	0.0470144\\
0.6262	0.0470016\\
0.6288	0.0469854\\
0.6313	0.0469665\\
0.6339	0.0469434\\
0.6364	0.0469177\\
0.639	0.0468876\\
0.6416	0.0468539\\
0.6441	0.0468183\\
0.6467	0.0467778\\
0.6493	0.0467337\\
0.6518	0.046688\\
0.6544	0.0466369\\
0.657	0.0465822\\
0.6595	0.0465263\\
0.6621	0.0464649\\
0.6646	0.0464024\\
0.6672	0.046334\\
0.6698	0.046262\\
0.6723	0.0461894\\
0.6749	0.0461105\\
0.6775	0.0460282\\
0.68	0.0459458\\
0.6826	0.0458566\\
0.6852	0.045764\\
0.6877	0.0456716\\
0.6903	0.0455722\\
0.6928	0.0454734\\
0.6954	0.0453674\\
0.698	0.045258\\
0.7005	0.0451496\\
0.7031	0.0450335\\
0.7057	0.0449141\\
0.7082	0.0447962\\
0.7108	0.0446704\\
0.7134	0.0445413\\
0.7159	0.0444142\\
0.7185	0.0442787\\
0.721	0.0441455\\
0.7236	0.0440037\\
0.7262	0.0438589\\
0.7287	0.0437167\\
0.7313	0.0435659\\
0.7339	0.0434119\\
0.7364	0.043261\\
0.739	0.043101\\
0.7416	0.0429381\\
0.7441	0.0427787\\
0.7467	0.0426101\\
0.7492	0.0424453\\
0.7518	0.0422711\\
0.7544	0.0420939\\
0.7569	0.041921\\
0.7595	0.0417385\\
0.7621	0.0415533\\
0.7646	0.0413727\\
0.7672	0.0411822\\
0.7698	0.0409891\\
0.7723	0.0408009\\
0.7749	0.0406028\\
0.7775	0.0404021\\
0.78	0.0402069\\
0.7826	0.0400014\\
0.7851	0.0398016\\
0.7877	0.0395914\\
0.7903	0.0393788\\
0.7928	0.0391723\\
0.7954	0.0389553\\
0.798	0.0387362\\
0.8005	0.0385233\\
0.8031	0.0382998\\
0.8057	0.0380742\\
0.8082	0.0378553\\
0.8108	0.0376256\\
0.8133	0.0374029\\
0.8159	0.0371694\\
0.8185	0.0369339\\
0.821	0.0367056\\
0.8236	0.0364664\\
0.8262	0.0362253\\
0.8287	0.0359918\\
0.8313	0.0357473\\
0.8339	0.035501\\
0.8364	0.0352626\\
0.839	0.035013\\
0.8415	0.0347715\\
0.8441	0.0345187\\
0.8467	0.0342645\\
0.8492	0.0340185\\
0.8518	0.0337612\\
0.8544	0.0335024\\
0.8569	0.0332522\\
0.8595	0.0329907\\
0.8621	0.0327278\\
0.8646	0.0324737\\
0.8672	0.0322082\\
0.8697	0.0319516\\
0.8723	0.0316835\\
0.8749	0.0314142\\
0.8774	0.0311541\\
0.88	0.0308825\\
0.8826	0.0306097\\
0.8851	0.0303463\\
0.8877	0.0300713\\
0.8903	0.0297952\\
0.8928	0.0295287\\
0.8954	0.0292506\\
0.8979	0.0289822\\
0.9005	0.0287021\\
0.9031	0.0284209\\
0.9056	0.0281497\\
0.9082	0.0278668\\
0.9108	0.0275829\\
0.9133	0.0273091\\
0.9159	0.0270235\\
0.9185	0.026737\\
0.921	0.0264607\\
0.9236	0.0261726\\
0.9262	0.0258836\\
0.9287	0.025605\\
0.9313	0.0253145\\
0.9338	0.0250344\\
0.9364	0.0247424\\
0.939	0.0244496\\
0.9415	0.0241674\\
0.9441	0.0238731\\
0.9467	0.0235782\\
0.9492	0.0232939\\
0.9518	0.0229976\\
0.9544	0.0227005\\
0.9569	0.0224142\\
0.9595	0.0221158\\
0.962	0.0218283\\
0.9646	0.0215286\\
0.9672	0.0212282\\
0.9697	0.0209387\\
0.9723	0.020637\\
0.9749	0.0203347\\
0.9774	0.0200433\\
0.98	0.0197397\\
0.9826	0.0194354\\
0.9851	0.0191422\\
0.9877	0.0188367\\
0.9902	0.0185422\\
0.9928	0.0182354\\
0.9954	0.0179279\\
0.9979	0.0176316\\
1.0005	0.0173229\\
1.0031	0.0170134\\
1.0056	0.0167152\\
1.0082	0.0164045\\
1.0108	0.0160931\\
1.0133	0.015793\\
1.0159	0.0154803\\
1.0184	0.0151789\\
1.021	0.0148648\\
1.0236	0.01455\\
1.0261	0.0142467\\
1.0287	0.0139305\\
1.0313	0.0136136\\
1.0338	0.0133082\\
1.0364	0.0129899\\
1.039	0.0126708\\
1.0415	0.0123633\\
1.0441	0.0120428\\
1.0466	0.0117339\\
1.0492	0.0114119\\
1.0518	0.011089\\
1.0543	0.0107779\\
1.0569	0.0104535\\
1.0595	0.0101283\\
1.062	0.00981481\\
1.0646	0.009488\\
1.0672	0.00916034\\
1.0697	0.00884447\\
1.0723	0.00851511\\
1.0748	0.00819759\\
1.0774	0.00786652\\
1.08	0.00753455\\
1.0825	0.0072145\\
1.0851	0.00688073\\
1.0877	0.00654602\\
1.0902	0.00622329\\
1.0928	0.00588673\\
1.0954	0.00554921\\
1.0979	0.00522374\\
1.1005	0.00488428\\
1.1031	0.00454381\\
1.1056	0.00421548\\
1.1082	0.00387301\\
1.1107	0.00354274\\
1.1133	0.00319822\\
1.1159	0.00285265\\
1.1184	0.00251935\\
1.121	0.00217164\\
1.1236	0.00182283\\
1.1261	0.0014864\\
1.1287	0.00113541\\
1.1313	0.000783287\\
1.1338	0.000443622\\
1.1364	8.92271e-05\\
1.1389	-0.000252645\\
1.1415	-0.00060935\\
1.1441	-0.000967245\\
1.1466	-0.00131251\\
1.1492	-0.00167278\\
1.1518	-0.00203428\\
1.1543	-0.00238306\\
1.1569	-0.00274702\\
1.1595	-0.00311223\\
1.162	-0.0034646\\
1.1646	-0.00383232\\
1.1671	-0.00418713\\
1.1697	-0.00455741\\
1.1723	-0.00492902\\
1.1748	-0.00528758\\
1.1774	-0.0056618\\
1.18	-0.00603736\\
1.1825	-0.00639976\\
1.1851	-0.00677801\\
1.1877	-0.00715764\\
1.1902	-0.00752397\\
1.1928	-0.00790633\\
1.1953	-0.00827531\\
1.1979	-0.00866043\\
1.2005	-0.00904698\\
1.203	-0.00942001\\
1.2056	-0.00980939\\
1.2082	-0.0102002\\
1.2107	-0.0105774\\
1.2133	-0.0109711\\
1.2159	-0.0113662\\
1.2184	-0.0117476\\
1.221	-0.0121456\\
1.2235	-0.0125298\\
1.2261	-0.0129308\\
1.2287	-0.0133333\\
1.2312	-0.0137218\\
1.2338	-0.0141272\\
1.2364	-0.0145342\\
1.2389	-0.014927\\
1.2415	-0.015337\\
1.2441	-0.0157486\\
1.2466	-0.0161457\\
1.2492	-0.0165603\\
1.2518	-0.0169763\\
1.2543	-0.0173779\\
1.2569	-0.017797\\
1.2594	-0.0182015\\
1.262	-0.0186237\\
1.2646	-0.0190474\\
1.2671	-0.0194563\\
1.2697	-0.019883\\
1.2723	-0.0203114\\
1.2748	-0.0207247\\
1.2774	-0.021156\\
1.28	-0.021589\\
1.2825	-0.0220067\\
1.2851	-0.0224426\\
1.2876	-0.0228633\\
1.2902	-0.0233022\\
1.2928	-0.0237427\\
1.2953	-0.0241677\\
1.2979	-0.0246112\\
1.3005	-0.0250562\\
1.303	-0.0254855\\
1.3056	-0.0259335\\
1.3082	-0.026383\\
1.3107	-0.0268165\\
1.3133	-0.0272689\\
1.3158	-0.0277053\\
1.3184	-0.0281605\\
1.321	-0.0286172\\
1.3235	-0.0290578\\
1.3261	-0.0295173\\
1.3287	-0.0299782\\
1.3312	-0.0304228\\
1.3338	-0.0308865\\
1.3364	-0.0313516\\
1.3389	-0.0318\\
1.3415	-0.0322678\\
1.344	-0.0327188\\
1.3466	-0.0331892\\
1.3492	-0.0336608\\
1.3517	-0.0341156\\
1.3543	-0.0345897\\
1.3569	-0.0350652\\
1.3594	-0.0355235\\
1.362	-0.0360013\\
1.3646	-0.0364803\\
1.3671	-0.036942\\
1.3697	-0.0374233\\
1.3722	-0.0378872\\
1.3748	-0.0383707\\
1.3774	-0.0388553\\
1.3799	-0.0393223\\
1.3825	-0.039809\\
1.3851	-0.0402966\\
1.3876	-0.0407665\\
1.3902	-0.0412561\\
1.3928	-0.0417467\\
1.3953	-0.0422192\\
1.3979	-0.0427115\\
1.4005	-0.0432047\\
1.403	-0.0436797\\
1.4056	-0.0441744\\
1.4081	-0.0446509\\
1.4107	-0.0451472\\
1.4133	-0.0456441\\
1.4158	-0.0461226\\
1.4184	-0.0466209\\
1.421	-0.0471198\\
1.4235	-0.0476001\\
1.4261	-0.0481001\\
1.4287	-0.0486006\\
1.4312	-0.0490824\\
1.4338	-0.0495838\\
1.4363	-0.0500664\\
1.4389	-0.0505686\\
1.4415	-0.0510712\\
1.444	-0.0515548\\
1.4466	-0.0520579\\
1.4492	-0.0525613\\
1.4517	-0.0530455\\
1.4543	-0.0535492\\
1.4569	-0.0540531\\
1.4594	-0.0545377\\
1.462	-0.0550417\\
1.4645	-0.0555263\\
1.4671	-0.0560302\\
1.4697	-0.0565341\\
1.4722	-0.0570185\\
1.4748	-0.0575221\\
1.4774	-0.0580255\\
1.4799	-0.0585093\\
1.4825	-0.0590121\\
1.4851	-0.0595147\\
1.4876	-0.0599976\\
1.4902	-0.0604994\\
1.4927	-0.0609814\\
1.4953	-0.0614822\\
1.4979	-0.0619825\\
1.5004	-0.062463\\
1.503	-0.0629621\\
1.5056	-0.0634605\\
1.5081	-0.063939\\
1.5107	-0.0644359\\
1.5133	-0.064932\\
1.5158	-0.0654082\\
1.5184	-0.0659026\\
1.5209	-0.066377\\
1.5235	-0.0668694\\
1.5261	-0.0673607\\
1.5286	-0.0678322\\
1.5312	-0.0683213\\
1.5338	-0.0688092\\
1.5363	-0.0692772\\
1.5389	-0.0697627\\
1.5415	-0.0702467\\
1.544	-0.0707109\\
1.5466	-0.0711922\\
1.5491	-0.0716536\\
1.5517	-0.0721319\\
1.5543	-0.0726086\\
1.5568	-0.0730655\\
1.5594	-0.073539\\
1.562	-0.0740107\\
1.5645	-0.0744626\\
1.5671	-0.0749307\\
1.5697	-0.075397\\
1.5722	-0.0758435\\
1.5748	-0.0763059\\
1.5774	-0.0767663\\
1.5799	-0.077207\\
1.5825	-0.0776633\\
1.585	-0.0780999\\
1.5876	-0.0785518\\
1.5902	-0.0790014\\
1.5927	-0.0794315\\
1.5953	-0.0798765\\
1.5979	-0.0803191\\
1.6004	-0.0807423\\
1.603	-0.0811799\\
1.6056	-0.081615\\
1.6081	-0.0820309\\
1.6107	-0.0824607\\
1.6132	-0.0828715\\
1.6158	-0.083296\\
1.6184	-0.0837176\\
1.6209	-0.0841204\\
1.6235	-0.0845364\\
1.6261	-0.0849494\\
1.6286	-0.0853438\\
1.6312	-0.0857509\\
1.6338	-0.0861549\\
1.6363	-0.0865405\\
1.6389	-0.0869383\\
1.6414	-0.0873178\\
1.644	-0.0877092\\
1.6466	-0.0880974\\
1.6491	-0.0884675\\
1.6517	-0.088849\\
1.6543	-0.0892271\\
1.6568	-0.0895874\\
1.6594	-0.0899586\\
1.662	-0.0903263\\
1.6645	-0.0906764\\
1.6671	-0.091037\\
1.6696	-0.0913802\\
1.6722	-0.0917335\\
1.6748	-0.0920831\\
1.6773	-0.0924156\\
1.6799	-0.0927577\\
1.6825	-0.093096\\
1.685	-0.0934176\\
1.6876	-0.0937481\\
1.6902	-0.0940747\\
1.6927	-0.094385\\
1.6953	-0.0947038\\
1.6978	-0.0950064\\
1.7004	-0.0953172\\
1.703	-0.0956238\\
1.7055	-0.0959147\\
1.7081	-0.0962132\\
1.7107	-0.0965075\\
1.7132	-0.0967864\\
1.7158	-0.0970723\\
1.7184	-0.097354\\
1.7209	-0.0976207\\
1.7235	-0.0978939\\
1.7261	-0.0981627\\
1.7286	-0.0984171\\
1.7312	-0.0986773\\
1.7337	-0.0989234\\
1.7363	-0.0991749\\
1.7389	-0.0994219\\
1.7414	-0.0996553\\
1.744	-0.0998935\\
1.7466	-0.100127\\
1.7491	-0.100348\\
1.7517	-0.100573\\
1.7543	-0.100793\\
1.7568	-0.101\\
1.7594	-0.101212\\
1.7619	-0.101411\\
1.7645	-0.101613\\
1.7671	-0.101811\\
1.7696	-0.101996\\
1.7722	-0.102185\\
1.7748	-0.102369\\
1.7773	-0.102541\\
1.7799	-0.102716\\
1.7825	-0.102886\\
1.785	-0.103046\\
1.7876	-0.103207\\
1.7901	-0.103357\\
1.7927	-0.103509\\
1.7953	-0.103656\\
1.7978	-0.103793\\
1.8004	-0.103931\\
1.803	-0.104064\\
1.8055	-0.104188\\
1.8081	-0.104312\\
1.8107	-0.104432\\
1.8132	-0.104542\\
1.8158	-0.104653\\
1.8183	-0.104755\\
1.8209	-0.104856\\
1.8235	-0.104953\\
1.826	-0.105041\\
1.8286	-0.105129\\
1.8312	-0.105212\\
1.8337	-0.105287\\
1.8363	-0.105361\\
1.8389	-0.10543\\
1.8414	-0.105493\\
1.844	-0.105553\\
1.8465	-0.105607\\
1.8491	-0.105659\\
1.8517	-0.105706\\
1.8542	-0.105747\\
1.8568	-0.105785\\
1.8594	-0.105819\\
1.8619	-0.105848\\
1.8645	-0.105873\\
1.8671	-0.105894\\
1.8696	-0.10591\\
1.8722	-0.105922\\
1.8747	-0.10593\\
1.8773	-0.105933\\
1.8799	-0.105933\\
1.8824	-0.105928\\
1.885	-0.10592\\
1.8876	-0.105906\\
1.8901	-0.10589\\
1.8927	-0.105869\\
1.8953	-0.105843\\
1.8978	-0.105815\\
1.9004	-0.105782\\
1.903	-0.105744\\
1.9055	-0.105704\\
1.9081	-0.105659\\
1.9106	-0.105612\\
1.9132	-0.105559\\
1.9158	-0.105502\\
1.9183	-0.105443\\
1.9209	-0.105379\\
1.9235	-0.10531\\
1.926	-0.105241\\
1.9286	-0.105165\\
1.9312	-0.105086\\
1.9337	-0.105006\\
1.9363	-0.104919\\
1.9388	-0.104832\\
1.9414	-0.104739\\
1.944	-0.104641\\
1.9465	-0.104544\\
1.9491	-0.10444\\
1.9517	-0.104332\\
1.9542	-0.104226\\
1.9568	-0.104111\\
1.9594	-0.103994\\
1.9619	-0.103877\\
1.9645	-0.103753\\
1.967	-0.103631\\
1.9696	-0.103501\\
1.9722	-0.103367\\
1.9747	-0.103236\\
1.9773	-0.103096\\
1.9799	-0.102953\\
1.9824	-0.102813\\
1.985	-0.102664\\
1.9876	-0.102513\\
1.9901	-0.102364\\
1.9927	-0.102207\\
1.9952	-0.102053\\
1.9978	-0.101891\\
2.0004	-0.101725\\
2.0029	-0.101564\\
2.0055	-0.101393\\
2.0081	-0.10122\\
2.0106	-0.10105\\
2.0132	-0.100872\\
2.0158	-0.100691\\
2.0183	-0.100515\\
2.0209	-0.100329\\
2.0234	-0.100148\\
2.026	-0.0999574\\
2.0286	-0.0997646\\
2.0311	-0.0995771\\
2.0337	-0.0993798\\
2.0363	-0.0991803\\
2.0388	-0.0989864\\
2.0414	-0.0987826\\
2.044	-0.0985767\\
2.0465	-0.0983768\\
2.0491	-0.0981669\\
2.0517	-0.097955\\
2.0542	-0.0977494\\
2.0568	-0.0975336\\
2.0593	-0.0973244\\
2.0619	-0.097105\\
2.0645	-0.0968838\\
2.067	-0.0966693\\
2.0696	-0.0964446\\
2.0722	-0.0962182\\
2.0747	-0.0959989\\
2.0773	-0.0957691\\
2.0799	-0.0955378\\
2.0824	-0.0953138\\
2.085	-0.0950794\\
2.0875	-0.0948526\\
2.0901	-0.0946152\\
2.0927	-0.0943764\\
2.0952	-0.0941453\\
2.0978	-0.0939037\\
2.1004	-0.0936608\\
2.1029	-0.0934259\\
2.1055	-0.0931803\\
2.1081	-0.0929335\\
2.1106	-0.0926949\\
2.1132	-0.0924457\\
2.1157	-0.0922049\\
2.1183	-0.0919533\\
2.1209	-0.0917006\\
2.1234	-0.0914566\\
2.126	-0.0912017\\
2.1286	-0.0909458\\
2.1311	-0.0906988\\
2.1337	-0.0904409\\
2.1363	-0.090182\\
2.1388	-0.0899322\\
2.1414	-0.0896715\\
2.1439	-0.08942\\
2.1465	-0.0891575\\
2.1491	-0.0888942\\
2.1516	-0.0886402\\
2.1542	-0.0883753\\
2.1568	-0.0881096\\
2.1593	-0.0878534\\
2.1619	-0.0875862\\
2.1645	-0.0873182\\
2.167	-0.08706\\
2.1696	-0.0867907\\
2.1721	-0.0865311\\
2.1747	-0.0862605\\
2.1773	-0.0859893\\
2.1798	-0.085728\\
2.1824	-0.0854556\\
2.185	-0.0851826\\
2.1875	-0.0849197\\
2.1901	-0.0846456\\
2.1927	-0.0843711\\
2.1952	-0.0841066\\
2.1978	-0.0838311\\
2.2004	-0.0835551\\
2.2029	-0.0832893\\
2.2055	-0.0830124\\
2.208	-0.0827458\\
2.2106	-0.0824681\\
2.2132	-0.08219\\
2.2157	-0.0819222\\
2.2183	-0.0816434\\
2.2209	-0.0813642\\
2.2234	-0.0810954\\
2.226	-0.0808155\\
2.2286	-0.0805354\\
2.2311	-0.0802657\\
2.2337	-0.0799849\\
2.2362	-0.0797146\\
2.2388	-0.0794333\\
2.2414	-0.0791517\\
2.2439	-0.0788807\\
2.2465	-0.0785987\\
2.2491	-0.0783164\\
2.2516	-0.0780448\\
2.2542	-0.0777621\\
2.2568	-0.0774792\\
2.2593	-0.077207\\
2.2619	-0.0769237\\
2.2644	-0.0766512\\
2.267	-0.0763676\\
2.2696	-0.0760839\\
2.2721	-0.075811\\
2.2747	-0.075527\\
2.2773	-0.075243\\
2.2798	-0.0749697\\
2.2824	-0.0746854\\
2.285	-0.0744011\\
2.2875	-0.0741276\\
2.2901	-0.073843\\
2.2926	-0.0735694\\
2.2952	-0.0732848\\
2.2978	-0.0730001\\
2.3003	-0.0727264\\
2.3029	-0.0724417\\
2.3055	-0.0721569\\
2.308	-0.0718832\\
2.3106	-0.0715984\\
2.3132	-0.0713137\\
2.3157	-0.07104\\
2.3183	-0.0707553\\
2.3208	-0.0704817\\
2.3234	-0.0701971\\
2.326	-0.0699126\\
2.3285	-0.0696391\\
2.3311	-0.0693547\\
2.3337	-0.0690705\\
2.3362	-0.0687972\\
2.3388	-0.0685132\\
2.3414	-0.0682293\\
2.3439	-0.0679564\\
2.3465	-0.0676727\\
2.349	-0.0674001\\
2.3516	-0.0671167\\
2.3542	-0.0668336\\
2.3567	-0.0665614\\
2.3593	-0.0662786\\
2.3619	-0.065996\\
2.3644	-0.0657245\\
2.367	-0.0654423\\
2.3696	-0.0651604\\
2.3721	-0.0648895\\
2.3747	-0.0646081\\
2.3773	-0.064327\\
2.3798	-0.0640569\\
2.3824	-0.0637763\\
2.3849	-0.0635069\\
2.3875	-0.0632269\\
2.3901	-0.0629473\\
2.3926	-0.0626788\\
2.3952	-0.0623999\\
2.3978	-0.0621213\\
2.4003	-0.0618539\\
2.4029	-0.0615761\\
2.4055	-0.0612988\\
2.408	-0.0610325\\
2.4106	-0.060756\\
2.4131	-0.0604906\\
2.4157	-0.060215\\
2.4183	-0.0599399\\
2.4208	-0.0596759\\
2.4234	-0.0594018\\
2.426	-0.0591283\\
2.4285	-0.0588658\\
2.4311	-0.0585933\\
2.4337	-0.0583214\\
2.4362	-0.0580605\\
2.4388	-0.0577898\\
2.4413	-0.0575301\\
2.4439	-0.0572606\\
2.4465	-0.0569918\\
2.449	-0.0567339\\
2.4516	-0.0564664\\
2.4542	-0.0561997\\
2.4567	-0.0559438\\
2.4593	-0.0556785\\
2.4619	-0.0554139\\
2.4644	-0.0551602\\
2.467	-0.0548971\\
2.4695	-0.0546449\\
2.4721	-0.0543834\\
2.4747	-0.0541228\\
2.4772	-0.053873\\
2.4798	-0.0536141\\
2.4824	-0.053356\\
2.4849	-0.0531088\\
2.4875	-0.0528525\\
2.4901	-0.0525972\\
2.4926	-0.0523527\\
2.4952	-0.0520993\\
2.4977	-0.0518566\\
2.5003	-0.0516052\\
2.5029	-0.0513549\\
2.5054	-0.0511151\\
2.508	-0.0508668\\
2.5106	-0.0506197\\
2.5131	-0.050383\\
2.5157	-0.0501381\\
2.5183	-0.0498942\\
2.5208	-0.0496609\\
2.5234	-0.0494194\\
2.526	-0.049179\\
2.5285	-0.0489491\\
2.5311	-0.0487112\\
2.5336	-0.0484837\\
2.5362	-0.0482482\\
2.5388	-0.0480141\\
2.5413	-0.0477903\\
2.5439	-0.0475588\\
2.5465	-0.0473286\\
2.549	-0.0471086\\
2.5516	-0.0468811\\
2.5542	-0.0466551\\
2.5567	-0.0464391\\
2.5593	-0.0462159\\
2.5618	-0.0460027\\
2.5644	-0.0457823\\
2.567	-0.0455635\\
2.5695	-0.0453546\\
2.5721	-0.0451387\\
2.5747	-0.0449245\\
2.5772	-0.04472\\
2.5798	-0.0445088\\
2.5824	-0.0442993\\
2.5849	-0.0440993\\
2.5875	-0.043893\\
2.59	-0.0436962\\
2.5926	-0.0434932\\
2.5952	-0.0432918\\
2.5977	-0.0430999\\
2.6003	-0.0429019\\
2.6029	-0.0427057\\
2.6054	-0.0425188\\
2.608	-0.0423261\\
2.6106	-0.0421352\\
2.6131	-0.0419533\\
2.6157	-0.041766\\
2.6182	-0.0415877\\
2.6208	-0.041404\\
2.6234	-0.0412223\\
2.6259	-0.0410493\\
2.6285	-0.0408713\\
2.6311	-0.0406952\\
2.6336	-0.0405277\\
2.6362	-0.0403555\\
2.6388	-0.0401853\\
2.6413	-0.0400235\\
2.6439	-0.0398572\\
2.6464	-0.0396992\\
2.649	-0.0395369\\
2.6516	-0.0393766\\
2.6541	-0.0392245\\
2.6567	-0.0390683\\
2.6593	-0.0389143\\
2.6618	-0.0387681\\
2.6644	-0.0386182\\
2.667	-0.0384705\\
2.6695	-0.0383304\\
2.6721	-0.0381869\\
2.6746	-0.038051\\
2.6772	-0.0379118\\
2.6798	-0.0377748\\
2.6823	-0.0376452\\
2.6849	-0.0375126\\
2.6875	-0.0373822\\
2.69	-0.037259\\
2.6926	-0.037133\\
2.6952	-0.0370094\\
2.6977	-0.0368927\\
2.7003	-0.0367736\\
2.7029	-0.0366568\\
2.7054	-0.0365467\\
2.708	-0.0364344\\
2.7105	-0.0363288\\
2.7131	-0.0362212\\
2.7157	-0.0361159\\
2.7182	-0.036017\\
2.7208	-0.0359165\\
2.7234	-0.0358183\\
2.7259	-0.0357263\\
2.7285	-0.0356329\\
2.7311	-0.0355419\\
2.7336	-0.0354567\\
2.7362	-0.0353705\\
2.7387	-0.03529\\
2.7413	-0.0352087\\
2.7439	-0.0351298\\
2.7464	-0.0350563\\
2.749	-0.0349823\\
2.7516	-0.0349108\\
2.7541	-0.0348444\\
2.7567	-0.0347778\\
2.7593	-0.0347137\\
2.7618	-0.0346544\\
2.7644	-0.0345953\\
2.7669	-0.0345408\\
2.7695	-0.0344867\\
2.7721	-0.0344351\\
2.7746	-0.0343878\\
2.7772	-0.0343412\\
2.7798	-0.0342972\\
2.7823	-0.0342573\\
2.7849	-0.0342182\\
2.7875	-0.0341818\\
2.79	-0.0341492\\
2.7926	-0.0341178\\
2.7951	-0.0340901\\
2.7977	-0.0340638\\
2.8003	-0.03404\\
2.8028	-0.0340197\\
2.8054	-0.034001\\
2.808	-0.033985\\
2.8105	-0.033972\\
2.8131	-0.033961\\
2.8157	-0.0339527\\
2.8182	-0.0339471\\
2.8208	-0.0339438\\
2.8233	-0.0339431\\
2.8259	-0.033945\\
2.8285	-0.0339494\\
2.831	-0.0339561\\
2.8336	-0.0339657\\
2.8362	-0.0339778\\
2.8387	-0.0339919\\
2.8413	-0.0340091\\
2.8439	-0.0340289\\
2.8464	-0.0340504\\
2.849	-0.0340753\\
2.8516	-0.0341028\\
2.8541	-0.0341316\\
2.8567	-0.0341641\\
2.8592	-0.0341978\\
2.8618	-0.0342354\\
2.8644	-0.0342755\\
2.8669	-0.0343165\\
2.8695	-0.0343616\\
2.8721	-0.0344092\\
2.8746	-0.0344575\\
2.8772	-0.0345101\\
2.8798	-0.0345653\\
2.8823	-0.0346207\\
2.8849	-0.0346808\\
2.8874	-0.0347409\\
2.89	-0.0348059\\
2.8926	-0.0348734\\
2.8951	-0.0349406\\
2.8977	-0.0350129\\
2.9003	-0.0350877\\
2.9028	-0.0351619\\
2.9054	-0.0352415\\
2.908	-0.0353234\\
2.9105	-0.0354046\\
2.9131	-0.0354913\\
2.9156	-0.0355769\\
2.9182	-0.0356683\\
2.9208	-0.035762\\
2.9233	-0.0358544\\
2.9259	-0.0359527\\
2.9285	-0.0360533\\
2.931	-0.0361523\\
2.9336	-0.0362575\\
2.9362	-0.0363649\\
2.9387	-0.0364704\\
2.9413	-0.0365822\\
2.9438	-0.0366919\\
2.9464	-0.0368081\\
2.949	-0.0369265\\
2.9515	-0.0370424\\
2.9541	-0.0371651\\
2.9567	-0.0372899\\
2.9592	-0.0374119\\
2.9618	-0.0375409\\
2.9644	-0.0376719\\
2.9669	-0.0377999\\
2.9695	-0.0379349\\
2.972	-0.0380667\\
2.9746	-0.0382057\\
2.9772	-0.0383467\\
2.9797	-0.0384841\\
2.9823	-0.0386289\\
2.9849	-0.0387756\\
2.9874	-0.0389185\\
2.99	-0.0390689\\
2.9926	-0.0392211\\
2.9951	-0.0393692\\
2.9977	-0.039525\\
3.0003	-0.0396826\\
3.0028	-0.0398357\\
3.0054	-0.0399967\\
3.0079	-0.0401531\\
3.0105	-0.0403173\\
3.0131	-0.0404832\\
3.0156	-0.0406443\\
3.0182	-0.0408133\\
3.0208	-0.0409839\\
3.0233	-0.0411494\\
3.0259	-0.041323\\
3.0285	-0.0414981\\
3.031	-0.0416678\\
3.0336	-0.0418457\\
3.0361	-0.042018\\
3.0387	-0.0421986\\
3.0413	-0.0423805\\
3.0438	-0.0425567\\
3.0464	-0.0427412\\
3.049	-0.0429269\\
3.0515	-0.0431066\\
3.0541	-0.0432946\\
3.0567	-0.0434839\\
3.0592	-0.0436669\\
3.0618	-0.0438583\\
3.0643	-0.0440433\\
3.0669	-0.0442367\\
3.0695	-0.0444312\\
3.072	-0.0446191\\
3.0746	-0.0448154\\
3.0772	-0.0450126\\
3.0797	-0.0452031\\
3.0823	-0.045402\\
3.0849	-0.0456017\\
3.0874	-0.0457945\\
3.09	-0.0459957\\
3.0925	-0.0461898\\
3.0951	-0.0463924\\
3.0977	-0.0465956\\
3.1002	-0.0467916\\
3.1028	-0.046996\\
3.1054	-0.0472009\\
3.1079	-0.0473985\\
3.1105	-0.0476044\\
3.1131	-0.0478108\\
3.1156	-0.0480097\\
3.1182	-0.0482169\\
3.1207	-0.0484164\\
3.1233	-0.0486242\\
3.1259	-0.0488324\\
3.1284	-0.0490328\\
3.131	-0.0492414\\
3.1336	-0.0494502\\
3.1361	-0.0496511\\
3.1387	-0.0498602\\
3.1413	-0.0500694\\
3.1438	-0.0502706\\
3.1464	-0.0504799\\
3.1489	-0.0506812\\
3.1515	-0.0508904\\
3.1541	-0.0510996\\
3.1566	-0.0513007\\
3.1592	-0.0515097\\
3.1618	-0.0517186\\
3.1643	-0.0519192\\
3.1669	-0.0521277\\
3.1695	-0.0523359\\
3.172	-0.0525359\\
3.1746	-0.0527436\\
3.1772	-0.0529509\\
3.1797	-0.0531499\\
3.1823	-0.0533565\\
3.1848	-0.0535548\\
3.1874	-0.0537605\\
3.19	-0.0539658\\
3.1925	-0.0541627\\
3.1951	-0.054367\\
3.1977	-0.0545707\\
3.2002	-0.054766\\
3.2028	-0.0549685\\
3.2054	-0.0551704\\
3.2079	-0.0553639\\
3.2105	-0.0555645\\
3.213	-0.0557567\\
3.2156	-0.0559558\\
3.2182	-0.0561542\\
3.2207	-0.0563442\\
3.2233	-0.056541\\
3.2259	-0.0567369\\
3.2284	-0.0569246\\
3.231	-0.0571188\\
3.2336	-0.0573122\\
3.2361	-0.0574972\\
3.2387	-0.0576887\\
3.2412	-0.0578719\\
3.2438	-0.0580615\\
3.2464	-0.05825\\
3.2489	-0.0584303\\
3.2515	-0.0586168\\
3.2541	-0.0588023\\
3.2566	-0.0589795\\
3.2592	-0.0591628\\
3.2618	-0.0593449\\
3.2643	-0.0595189\\
3.2669	-0.0596988\\
3.2694	-0.0598706\\
3.272	-0.0600481\\
3.2746	-0.0602243\\
3.2771	-0.0603927\\
3.2797	-0.0605665\\
3.2823	-0.0607391\\
3.2848	-0.0609038\\
3.2874	-0.0610738\\
3.29	-0.0612425\\
3.2925	-0.0614035\\
3.2951	-0.0615696\\
3.2976	-0.0617281\\
3.3002	-0.0618915\\
3.3028	-0.0620536\\
3.3053	-0.0622081\\
3.3079	-0.0623674\\
3.3105	-0.0625253\\
3.313	-0.0626758\\
3.3156	-0.062831\\
3.3182	-0.0629846\\
3.3207	-0.063131\\
3.3233	-0.0632818\\
3.3259	-0.0634312\\
3.3284	-0.0635734\\
3.331	-0.0637198\\
3.3335	-0.0638592\\
3.3361	-0.0640027\\
3.3387	-0.0641447\\
3.3412	-0.0642799\\
3.3438	-0.0644189\\
3.3464	-0.0645564\\
3.3489	-0.0646872\\
3.3515	-0.0648217\\
3.3541	-0.0649546\\
3.3566	-0.065081\\
3.3592	-0.0652109\\
3.3617	-0.0653344\\
3.3643	-0.0654613\\
3.3669	-0.0655866\\
3.3694	-0.0657057\\
3.372	-0.065828\\
3.3746	-0.0659487\\
3.3771	-0.0660633\\
3.3797	-0.066181\\
3.3823	-0.0662971\\
3.3848	-0.0664072\\
3.3874	-0.0665202\\
3.3899	-0.0666275\\
3.3925	-0.0667374\\
3.3951	-0.0668458\\
3.3976	-0.0669486\\
3.4002	-0.067054\\
3.4028	-0.0671578\\
3.4053	-0.0672561\\
3.4079	-0.0673569\\
3.4105	-0.0674561\\
3.413	-0.06755\\
3.4156	-0.0676462\\
3.4181	-0.0677372\\
3.4207	-0.0678304\\
3.4233	-0.0679221\\
3.4258	-0.0680087\\
3.4284	-0.0680974\\
3.431	-0.0681846\\
3.4335	-0.0682669\\
3.4361	-0.0683512\\
3.4387	-0.0684339\\
3.4412	-0.068512\\
3.4438	-0.0685918\\
3.4463	-0.0686671\\
3.4489	-0.068744\\
3.4515	-0.0688195\\
3.454	-0.0688906\\
3.4566	-0.0689632\\
3.4592	-0.0690344\\
3.4617	-0.0691014\\
3.4643	-0.0691698\\
3.4669	-0.0692367\\
3.4694	-0.0692998\\
3.472	-0.069364\\
3.4745	-0.0694244\\
3.4771	-0.0694858\\
3.4797	-0.0695459\\
3.4822	-0.0696024\\
3.4848	-0.0696599\\
3.4874	-0.069716\\
3.4899	-0.0697686\\
3.4925	-0.0698221\\
3.4951	-0.0698743\\
3.4976	-0.0699232\\
3.5002	-0.0699729\\
3.5028	-0.0700212\\
3.5053	-0.0700665\\
3.5079	-0.0701123\\
3.5104	-0.0701553\\
3.513	-0.0701987\\
3.5156	-0.0702409\\
3.5181	-0.0702803\\
3.5207	-0.0703201\\
3.5233	-0.0703587\\
3.5258	-0.0703948\\
3.5284	-0.0704311\\
3.531	-0.0704662\\
3.5335	-0.0704989\\
3.5361	-0.0705319\\
3.5386	-0.0705624\\
3.5412	-0.0705931\\
3.5438	-0.0706228\\
3.5463	-0.0706502\\
3.5489	-0.0706777\\
3.5515	-0.0707041\\
3.554	-0.0707285\\
3.5566	-0.0707529\\
3.5592	-0.0707762\\
3.5617	-0.0707976\\
3.5643	-0.070819\\
3.5668	-0.0708385\\
3.5694	-0.0708579\\
3.572	-0.0708763\\
3.5745	-0.0708931\\
3.5771	-0.0709097\\
3.5797	-0.0709253\\
3.5822	-0.0709394\\
3.5848	-0.0709532\\
3.5874	-0.0709661\\
3.5899	-0.0709777\\
3.5925	-0.0709888\\
3.595	-0.0709987\\
3.5976	-0.0710082\\
3.6002	-0.0710168\\
3.6027	-0.0710242\\
3.6053	-0.0710312\\
3.6079	-0.0710374\\
3.6104	-0.0710425\\
3.613	-0.0710471\\
3.6156	-0.0710509\\
3.6181	-0.0710539\\
3.6207	-0.0710562\\
3.6232	-0.0710577\\
3.6258	-0.0710585\\
3.6284	-0.0710586\\
3.6309	-0.071058\\
3.6335	-0.0710567\\
3.6361	-0.0710547\\
3.6386	-0.0710521\\
3.6412	-0.0710487\\
3.6438	-0.0710446\\
3.6463	-0.0710401\\
3.6489	-0.0710347\\
3.6515	-0.0710287\\
3.654	-0.0710223\\
3.6566	-0.0710151\\
3.6591	-0.0710075\\
3.6617	-0.070999\\
3.6643	-0.0709899\\
3.6668	-0.0709806\\
3.6694	-0.0709703\\
3.672	-0.0709594\\
3.6745	-0.0709484\\
3.6771	-0.0709363\\
3.6797	-0.0709237\\
3.6822	-0.0709111\\
3.6848	-0.0708974\\
3.6873	-0.0708837\\
3.6899	-0.0708689\\
3.6925	-0.0708536\\
3.695	-0.0708383\\
3.6976	-0.0708219\\
3.7002	-0.070805\\
3.7027	-0.0707882\\
3.7053	-0.0707702\\
3.7079	-0.0707518\\
3.7104	-0.0707335\\
3.713	-0.070714\\
3.7155	-0.0706948\\
3.7181	-0.0706743\\
3.7207	-0.0706533\\
3.7232	-0.0706327\\
3.7258	-0.0706107\\
3.7284	-0.0705883\\
3.7309	-0.0705662\\
3.7335	-0.0705427\\
3.7361	-0.0705188\\
3.7386	-0.0704953\\
3.7412	-0.0704704\\
3.7437	-0.070446\\
3.7463	-0.0704202\\
3.7489	-0.0703938\\
3.7514	-0.070368\\
3.754	-0.0703407\\
3.7566	-0.0703129\\
3.7591	-0.0702857\\
3.7617	-0.0702569\\
3.7643	-0.0702277\\
3.7668	-0.0701991\\
3.7694	-0.0701688\\
3.7719	-0.0701393\\
3.7745	-0.0701081\\
3.7771	-0.0700764\\
3.7796	-0.0700454\\
3.7822	-0.0700127\\
3.7848	-0.0699795\\
3.7873	-0.069947\\
3.7899	-0.0699128\\
3.7925	-0.0698781\\
3.795	-0.0698442\\
3.7976	-0.0698084\\
3.8002	-0.0697721\\
3.8027	-0.0697367\\
3.8053	-0.0696993\\
3.8078	-0.0696629\\
3.8104	-0.0696245\\
3.813	-0.0695855\\
3.8155	-0.0695475\\
3.8181	-0.0695074\\
3.8207	-0.0694668\\
3.8232	-0.0694272\\
3.8258	-0.0693855\\
3.8284	-0.0693431\\
3.8309	-0.0693019\\
3.8335	-0.0692584\\
3.836	-0.0692161\\
3.8386	-0.0691714\\
3.8412	-0.0691262\\
3.8437	-0.0690821\\
3.8463	-0.0690357\\
3.8489	-0.0689886\\
3.8514	-0.0689427\\
3.854	-0.0688944\\
3.8566	-0.0688455\\
3.8591	-0.0687978\\
3.8617	-0.0687476\\
3.8642	-0.0686987\\
3.8668	-0.0686472\\
3.8694	-0.0685951\\
3.8719	-0.0685443\\
3.8745	-0.0684908\\
3.8771	-0.0684366\\
3.8796	-0.0683839\\
3.8822	-0.0683283\\
3.8848	-0.0682721\\
3.8873	-0.0682174\\
3.8899	-0.0681597\\
3.8924	-0.0681037\\
3.895	-0.0680446\\
3.8976	-0.0679849\\
3.9001	-0.0679267\\
3.9027	-0.0678655\\
3.9053	-0.0678036\\
3.9078	-0.0677433\\
3.9104	-0.0676799\\
3.913	-0.0676157\\
3.9155	-0.0675532\\
3.9181	-0.0674876\\
3.9206	-0.0674237\\
3.9232	-0.0673565\\
3.9258	-0.0672885\\
3.9283	-0.0672224\\
3.9309	-0.0671528\\
3.9335	-0.0670825\\
3.936	-0.0670141\\
3.9386	-0.0669423\\
3.9412	-0.0668696\\
3.9437	-0.0667989\\
3.9463	-0.0667246\\
3.9488	-0.0666525\\
3.9514	-0.0665766\\
3.954	-0.0664999\\
3.9565	-0.0664254\\
3.9591	-0.0663471\\
3.9617	-0.066268\\
3.9642	-0.0661911\\
3.9668	-0.0661104\\
3.9694	-0.0660288\\
3.9719	-0.0659496\\
3.9745	-0.0658664\\
3.9771	-0.0657824\\
3.9796	-0.0657008\\
3.9822	-0.0656152\\
3.9847	-0.065532\\
3.9873	-0.0654447\\
3.9899	-0.0653566\\
3.9924	-0.0652711\\
3.995	-0.0651814\\
3.9976	-0.0650908\\
4.0001	-0.0650029\\
4.0027	-0.0649108\\
4.0053	-0.0648178\\
4.0078	-0.0647276\\
4.0104	-0.064633\\
4.0129	-0.0645412\\
4.0155	-0.064445\\
4.0181	-0.064348\\
4.0206	-0.064254\\
4.0232	-0.0641554\\
4.0258	-0.064056\\
4.0283	-0.0639597\\
4.0309	-0.0638588\\
4.0335	-0.063757\\
4.036	-0.0636585\\
4.0386	-0.0635553\\
4.0411	-0.0634553\\
4.0437	-0.0633506\\
4.0463	-0.0632451\\
4.0488	-0.0631429\\
4.0514	-0.063036\\
4.054	-0.0629283\\
4.0565	-0.062824\\
4.0591	-0.0627149\\
4.0617	-0.0626051\\
4.0642	-0.0624989\\
4.0668	-0.0623877\\
4.0693	-0.0622801\\
4.0719	-0.0621675\\
4.0745	-0.0620543\\
4.077	-0.0619448\\
4.0796	-0.0618303\\
4.0822	-0.0617152\\
4.0847	-0.0616039\\
4.0873	-0.0614875\\
4.0899	-0.0613705\\
4.0924	-0.0612574\\
4.095	-0.0611393\\
4.0975	-0.0610251\\
4.1001	-0.0609058\\
4.1027	-0.0607859\\
4.1052	-0.0606702\\
4.1078	-0.0605492\\
4.1104	-0.0604278\\
4.1129	-0.0603105\\
4.1155	-0.0601881\\
4.1181	-0.0600652\\
4.1206	-0.0599465\\
4.1232	-0.0598227\\
4.1258	-0.0596984\\
4.1283	-0.0595785\\
4.1309	-0.0594534\\
4.1334	-0.0593327\\
4.136	-0.0592068\\
4.1386	-0.0590805\\
4.1411	-0.0589587\\
4.1437	-0.0588317\\
4.1463	-0.0587044\\
4.1488	-0.0585817\\
4.1514	-0.0584537\\
4.154	-0.0583255\\
4.1565	-0.0582019\\
4.1591	-0.0580732\\
4.1616	-0.0579492\\
4.1642	-0.0578199\\
4.1668	-0.0576905\\
4.1693	-0.0575659\\
4.1719	-0.057436\\
4.1745	-0.0573061\\
4.177	-0.057181\\
4.1796	-0.0570507\\
4.1822	-0.0569203\\
4.1847	-0.0567949\\
4.1873	-0.0566643\\
4.1898	-0.0565387\\
4.1924	-0.0564081\\
4.195	-0.0562774\\
4.1975	-0.0561517\\
4.2001	-0.0560209\\
4.2027	-0.0558902\\
4.2052	-0.0557646\\
4.2078	-0.0556339\\
4.2104	-0.0555034\\
4.2129	-0.0553779\\
4.2155	-0.0552475\\
4.218	-0.0551223\\
4.2206	-0.0549922\\
4.2232	-0.0548622\\
4.2257	-0.0547374\\
4.2283	-0.0546077\\
4.2309	-0.0544783\\
4.2334	-0.0543541\\
4.236	-0.0542251\\
4.2386	-0.0540964\\
4.2411	-0.0539729\\
4.2437	-0.0538447\\
4.2462	-0.0537218\\
4.2488	-0.0535943\\
4.2514	-0.0534671\\
4.2539	-0.0533451\\
4.2565	-0.0532186\\
4.2591	-0.0530925\\
4.2616	-0.0529717\\
4.2642	-0.0528464\\
4.2668	-0.0527216\\
4.2693	-0.0526021\\
4.2719	-0.0524782\\
4.2744	-0.0523595\\
4.277	-0.0522366\\
4.2796	-0.0521142\\
4.2821	-0.0519971\\
4.2847	-0.0518758\\
4.2873	-0.0517551\\
4.2898	-0.0516396\\
4.2924	-0.0515201\\
4.295	-0.0514012\\
4.2975	-0.0512875\\
4.3001	-0.0511699\\
4.3027	-0.0510529\\
4.3052	-0.0509412\\
4.3078	-0.0508256\\
4.3103	-0.0507151\\
4.3129	-0.050601\\
4.3155	-0.0504876\\
4.318	-0.0503793\\
4.3206	-0.0502675\\
4.3232	-0.0501564\\
4.3257	-0.0500504\\
4.3283	-0.049941\\
4.3309	-0.0498324\\
4.3334	-0.0497287\\
4.336	-0.0496218\\
4.3385	-0.0495199\\
4.3411	-0.0494147\\
4.3437	-0.0493105\\
4.3462	-0.0492111\\
4.3488	-0.0491087\\
4.3514	-0.0490072\\
4.3539	-0.0489106\\
4.3565	-0.048811\\
4.3591	-0.0487125\\
4.3616	-0.0486186\\
4.3642	-0.0485221\\
4.3667	-0.0484302\\
4.3693	-0.0483356\\
4.3719	-0.0482422\\
4.3744	-0.0481533\\
4.377	-0.0480619\\
4.3796	-0.0479716\\
4.3821	-0.0478858\\
4.3847	-0.0477977\\
4.3873	-0.0477107\\
4.3898	-0.0476281\\
4.3924	-0.0475433\\
4.3949	-0.0474629\\
4.3975	-0.0473804\\
4.4001	-0.047299\\
4.4026	-0.0472219\\
4.4052	-0.0471429\\
4.4078	-0.0470651\\
4.4103	-0.0469914\\
4.4129	-0.046916\\
4.4155	-0.0468418\\
4.418	-0.0467716\\
4.4206	-0.0466999\\
4.4231	-0.046632\\
4.4257	-0.0465627\\
4.4283	-0.0464947\\
4.4308	-0.0464305\\
4.4334	-0.0463649\\
4.436	-0.0463007\\
4.4385	-0.0462401\\
4.4411	-0.0461784\\
4.4437	-0.046118\\
4.4462	-0.0460611\\
4.4488	-0.0460033\\
4.4514	-0.0459468\\
4.4539	-0.0458937\\
4.4565	-0.0458397\\
4.459	-0.0457892\\
4.4616	-0.0457379\\
4.4642	-0.0456879\\
4.4667	-0.0456412\\
4.4693	-0.0455939\\
4.4719	-0.0455479\\
4.4744	-0.045505\\
4.477	-0.0454618\\
4.4796	-0.0454198\\
4.4821	-0.0453808\\
4.4847	-0.0453416\\
4.4872	-0.0453052\\
4.4898	-0.0452686\\
4.4924	-0.0452334\\
4.4949	-0.0452009\\
4.4975	-0.0451684\\
4.5001	-0.0451373\\
4.5026	-0.0451087\\
4.5052	-0.0450803\\
4.5078	-0.0450532\\
4.5103	-0.0450285\\
4.5129	-0.0450042\\
4.5154	-0.0449821\\
4.518	-0.0449604\\
4.5206	-0.0449402\\
4.5231	-0.0449219\\
4.5257	-0.0449044\\
4.5283	-0.0448881\\
4.5308	-0.0448738\\
4.5334	-0.0448603\\
4.536	-0.0448481\\
4.5385	-0.0448376\\
4.5411	-0.0448281\\
4.5436	-0.0448202\\
4.5462	-0.0448133\\
4.5488	-0.0448077\\
4.5513	-0.0448037\\
4.5539	-0.0448007\\
4.5565	-0.0447991\\
4.559	-0.0447989\\
4.5616	-0.0447998\\
4.5642	-0.0448022\\
4.5667	-0.0448056\\
4.5693	-0.0448105\\
4.5718	-0.0448164\\
4.5744	-0.0448238\\
4.577	-0.0448325\\
4.5795	-0.0448421\\
4.5821	-0.0448533\\
4.5847	-0.0448658\\
4.5872	-0.0448789\\
4.5898	-0.0448939\\
4.5924	-0.0449101\\
4.5949	-0.0449268\\
4.5975	-0.0449454\\
4.6001	-0.0449652\\
4.6026	-0.0449854\\
4.6052	-0.0450076\\
4.6077	-0.0450301\\
4.6103	-0.0450547\\
4.6129	-0.0450804\\
4.6154	-0.0451062\\
4.618	-0.0451342\\
4.6206	-0.0451634\\
4.6231	-0.0451925\\
4.6257	-0.0452239\\
4.6283	-0.0452565\\
4.6308	-0.0452888\\
4.6334	-0.0453235\\
4.6359	-0.0453579\\
4.6385	-0.0453948\\
4.6411	-0.0454327\\
4.6436	-0.0454702\\
4.6462	-0.0455102\\
4.6488	-0.0455512\\
4.6513	-0.0455917\\
4.6539	-0.0456347\\
4.6565	-0.0456788\\
4.659	-0.0457221\\
4.6616	-0.0457682\\
4.6641	-0.0458133\\
4.6667	-0.0458613\\
4.6693	-0.0459101\\
4.6718	-0.045958\\
4.6744	-0.0460088\\
4.677	-0.0460604\\
4.6795	-0.0461109\\
4.6821	-0.0461643\\
4.6847	-0.0462186\\
4.6872	-0.0462716\\
4.6898	-0.0463276\\
4.6923	-0.0463823\\
4.6949	-0.0464399\\
4.6975	-0.0464984\\
4.7	-0.0465553\\
4.7026	-0.0466154\\
4.7052	-0.0466762\\
4.7077	-0.0467354\\
4.7103	-0.0467977\\
4.7129	-0.0468608\\
4.7154	-0.0469221\\
4.718	-0.0469866\\
4.7205	-0.0470493\\
4.7231	-0.0471152\\
4.7257	-0.0471817\\
4.7282	-0.0472463\\
4.7308	-0.0473142\\
4.7334	-0.0473827\\
4.7359	-0.0474492\\
4.7385	-0.0475189\\
4.7411	-0.0475892\\
4.7436	-0.0476574\\
4.7462	-0.0477289\\
4.7487	-0.0477982\\
4.7513	-0.0478708\\
4.7539	-0.0479439\\
4.7564	-0.0480148\\
4.759	-0.0480889\\
4.7616	-0.0481636\\
4.7641	-0.0482359\\
4.7667	-0.0483115\\
4.7693	-0.0483876\\
4.7718	-0.0484612\\
4.7744	-0.0485381\\
4.777	-0.0486155\\
4.7795	-0.0486904\\
4.7821	-0.0487686\\
4.7846	-0.0488441\\
4.7872	-0.0489231\\
4.7898	-0.0490024\\
4.7923	-0.0490791\\
4.7949	-0.0491591\\
4.7975	-0.0492394\\
4.8	-0.049317\\
4.8026	-0.049398\\
4.8052	-0.0494793\\
4.8077	-0.0495577\\
4.8103	-0.0496395\\
4.8128	-0.0497184\\
4.8154	-0.0498008\\
4.818	-0.0498833\\
4.8205	-0.0499629\\
4.8231	-0.0500459\\
4.8257	-0.0501292\\
4.8282	-0.0502094\\
4.8308	-0.0502929\\
4.8334	-0.0503767\\
4.8359	-0.0504574\\
4.8385	-0.0505415\\
4.841	-0.0506225\\
4.8436	-0.0507068\\
4.8462	-0.0507913\\
4.8487	-0.0508726\\
4.8513	-0.0509573\\
4.8539	-0.0510421\\
4.8564	-0.0511236\\
4.859	-0.0512086\\
4.8616	-0.0512936\\
4.8641	-0.0513753\\
4.8667	-0.0514604\\
4.8692	-0.0515423\\
4.8718	-0.0516274\\
4.8744	-0.0517126\\
4.8769	-0.0517945\\
4.8795	-0.0518797\\
4.8821	-0.0519648\\
4.8846	-0.0520467\\
4.8872	-0.0521318\\
4.8898	-0.0522169\\
4.8923	-0.0522987\\
4.8949	-0.0523837\\
4.8974	-0.0524654\\
4.9	-0.0525503\\
4.9026	-0.0526351\\
4.9051	-0.0527166\\
4.9077	-0.0528013\\
4.9103	-0.0528859\\
4.9128	-0.0529671\\
4.9154	-0.0530514\\
4.918	-0.0531357\\
4.9205	-0.0532166\\
4.9231	-0.0533006\\
4.9257	-0.0533845\\
4.9282	-0.053465\\
4.9308	-0.0535486\\
4.9333	-0.0536288\\
4.9359	-0.0537121\\
4.9385	-0.0537952\\
4.941	-0.053875\\
4.9436	-0.0539578\\
4.9462	-0.0540404\\
4.9487	-0.0541197\\
4.9513	-0.0542019\\
4.9539	-0.054284\\
4.9564	-0.0543627\\
4.959	-0.0544443\\
4.9615	-0.0545226\\
4.9641	-0.0546039\\
4.9667	-0.0546849\\
4.9692	-0.0547626\\
4.9718	-0.0548432\\
4.9744	-0.0549235\\
4.9769	-0.0550006\\
4.9795	-0.0550805\\
4.9821	-0.0551601\\
4.9846	-0.0552365\\
4.9872	-0.0553157\\
4.9897	-0.0553915\\
4.9923	-0.0554702\\
4.9949	-0.0555486\\
4.9974	-0.0556238\\
5	-0.0557017\\
};
\addplot [color=mycolor6, line width=1.0pt, forget plot]
  table[row sep=crcr]{%
-0.125	7.04569e-08\\
-0.1224	7.19238e-08\\
-0.1199	7.33342e-08\\
-0.1173	7.48011e-08\\
-0.1147	7.62679e-08\\
-0.1122	7.76783e-08\\
-0.1096	7.91451e-08\\
-0.1071	8.05554e-08\\
-0.1045	8.2022e-08\\
-0.1019	8.34887e-08\\
-0.0994	8.4899e-08\\
-0.0968	8.63656e-08\\
-0.0942	8.78321e-08\\
-0.0917	8.92421e-08\\
-0.0891	9.07085e-08\\
-0.0865	9.21748e-08\\
-0.084	9.35847e-08\\
-0.0814	9.50509e-08\\
-0.0789	9.64606e-08\\
-0.0763	9.79266e-08\\
-0.0737	9.93925e-08\\
-0.0712	1.00802e-07\\
-0.0686	1.02268e-07\\
-0.066	1.03733e-07\\
-0.0635	1.05143e-07\\
-0.0609	1.06608e-07\\
-0.0583	1.08073e-07\\
-0.0558	1.09482e-07\\
-0.0532	1.10948e-07\\
-0.0507	1.12356e-07\\
-0.0481	1.13821e-07\\
-0.0455	1.15286e-07\\
-0.043	1.16695e-07\\
-0.0404	1.18159e-07\\
-0.0378	1.19624e-07\\
-0.0353	1.21032e-07\\
-0.0327	1.22496e-07\\
-0.0301	1.2396e-07\\
-0.0276	1.25368e-07\\
-0.025	1.26832e-07\\
-0.0224	1.28295e-07\\
-0.0199	1.29703e-07\\
-0.0173	1.31166e-07\\
-0.0148	1.32573e-07\\
-0.0122	1.34037e-07\\
-0.0096	1.355e-07\\
-0.0071	1.36906e-07\\
-0.0045	1.38369e-07\\
-0.0019	1.39832e-07\\
0.0006	1.41238e-07\\
0.0032	6.72486e-05\\
0.0058	0.000667632\\
0.0083	0.00132778\\
0.0109	0.00189105\\
0.0134	0.00243554\\
0.016	0.00306772\\
0.0186	0.00373123\\
0.0211	0.00433506\\
0.0237	0.00489505\\
0.0263	0.0054045\\
0.0288	0.0058873\\
0.0314	0.00641184\\
0.034	0.00695965\\
0.0365	0.00748905\\
0.0391	0.00802243\\
0.0416	0.00851381\\
0.0442	0.00901223\\
0.0468	0.00951208\\
0.0493	0.0100007\\
0.0519	0.0105132\\
0.0545	0.0110199\\
0.057	0.0114944\\
0.0596	0.0119753\\
0.0622	0.0124501\\
0.0647	0.0129076\\
0.0673	0.0133856\\
0.0698	0.0138427\\
0.0724	0.0143087\\
0.075	0.0147619\\
0.0775	0.0151884\\
0.0801	0.0156279\\
0.0827	0.0160672\\
0.0852	0.016488\\
0.0878	0.0169186\\
0.0904	0.0173374\\
0.0929	0.0177287\\
0.0955	0.0181281\\
0.098	0.0185095\\
0.1006	0.0189046\\
0.1032	0.0192949\\
0.1057	0.0196608\\
0.1083	0.0200294\\
0.1109	0.0203881\\
0.1134	0.0207281\\
0.116	0.0210798\\
0.1186	0.0214285\\
0.1211	0.0217573\\
0.1237	0.0220888\\
0.1263	0.0224098\\
0.1288	0.022712\\
0.1314	0.0230237\\
0.1339	0.0233222\\
0.1365	0.0236287\\
0.1391	0.0239274\\
0.1416	0.0242056\\
0.1442	0.0244876\\
0.1468	0.024766\\
0.1493	0.0250331\\
0.1519	0.0253096\\
0.1545	0.0255815\\
0.157	0.0258359\\
0.1596	0.0260932\\
0.1621	0.0263367\\
0.1647	0.0265891\\
0.1673	0.0268419\\
0.1698	0.0270834\\
0.1724	0.0273298\\
0.175	0.0275698\\
0.1775	0.0277962\\
0.1801	0.0280302\\
0.1827	0.0282652\\
0.1852	0.0284915\\
0.1878	0.0287245\\
0.1903	0.0289439\\
0.1929	0.0291673\\
0.1955	0.0293884\\
0.198	0.0296014\\
0.2006	0.0298242\\
0.2032	0.0300466\\
0.2057	0.0302572\\
0.2083	0.0304718\\
0.2109	0.0306831\\
0.2134	0.0308856\\
0.216	0.0310974\\
0.2185	0.0313015\\
0.2211	0.0315119\\
0.2237	0.0317183\\
0.2262	0.0319128\\
0.2288	0.032113\\
0.2314	0.0323132\\
0.2339	0.0325063\\
0.2365	0.0327061\\
0.2391	0.0329025\\
0.2416	0.033087\\
0.2442	0.0332752\\
0.2467	0.0334548\\
0.2493	0.0336414\\
0.2519	0.0338272\\
0.2544	0.0340032\\
0.257	0.0341816\\
0.2596	0.0343551\\
0.2621	0.0345188\\
0.2647	0.0346876\\
0.2673	0.0348552\\
0.2698	0.0350141\\
0.2724	0.0351748\\
0.2749	0.035324\\
0.2775	0.0354743\\
0.2801	0.0356216\\
0.2826	0.0357613\\
0.2852	0.0359042\\
0.2878	0.0360428\\
0.2903	0.0361705\\
0.2929	0.0362976\\
0.2955	0.0364199\\
0.298	0.0365346\\
0.3006	0.0366512\\
0.3032	0.0367638\\
0.3057	0.0368668\\
0.3083	0.0369676\\
0.3108	0.0370589\\
0.3134	0.0371498\\
0.316	0.0372373\\
0.3185	0.0373179\\
0.3211	0.0373968\\
0.3237	0.0374693\\
0.3262	0.0375329\\
0.3288	0.0375939\\
0.3314	0.037651\\
0.3339	0.0377025\\
0.3365	0.0377516\\
0.339	0.0377933\\
0.3416	0.0378304\\
0.3442	0.0378617\\
0.3467	0.0378874\\
0.3493	0.0379107\\
0.3519	0.0379301\\
0.3544	0.037944\\
0.357	0.0379526\\
0.3596	0.0379553\\
0.3621	0.0379532\\
0.3647	0.0379474\\
0.3672	0.0379387\\
0.3698	0.0379256\\
0.3724	0.0379075\\
0.3749	0.0378847\\
0.3775	0.0378562\\
0.3801	0.0378238\\
0.3826	0.0377899\\
0.3852	0.0377514\\
0.3878	0.0377089\\
0.3903	0.0376635\\
0.3929	0.0376116\\
0.3954	0.0375581\\
0.398	0.0374998\\
0.4006	0.0374389\\
0.4031	0.0373776\\
0.4057	0.0373101\\
0.4083	0.0372385\\
0.4108	0.0371662\\
0.4134	0.0370884\\
0.416	0.0370087\\
0.4185	0.0369302\\
0.4211	0.0368458\\
0.4236	0.0367615\\
0.4262	0.0366708\\
0.4288	0.0365775\\
0.4313	0.0364864\\
0.4339	0.0363904\\
0.4365	0.0362928\\
0.439	0.0361966\\
0.4416	0.0360941\\
0.4442	0.0359895\\
0.4467	0.0358876\\
0.4493	0.035781\\
0.4519	0.0356736\\
0.4544	0.035569\\
0.457	0.0354585\\
0.4595	0.0353506\\
0.4621	0.0352374\\
0.4647	0.0351238\\
0.4672	0.0350145\\
0.4698	0.0349005\\
0.4724	0.0347855\\
0.4749	0.0346738\\
0.4775	0.0345568\\
0.4801	0.0344398\\
0.4826	0.0343277\\
0.4852	0.0342115\\
0.4877	0.0340997\\
0.4903	0.0339829\\
0.4929	0.0338657\\
0.4954	0.0337531\\
0.498	0.0336368\\
0.5006	0.0335214\\
0.5031	0.0334111\\
0.5057	0.0332966\\
0.5083	0.0331821\\
0.5108	0.0330724\\
0.5134	0.0329593\\
0.5159	0.0328518\\
0.5185	0.0327413\\
0.5211	0.0326317\\
0.5236	0.0325269\\
0.5262	0.0324185\\
0.5288	0.0323112\\
0.5313	0.0322096\\
0.5339	0.0321057\\
0.5365	0.0320033\\
0.539	0.031906\\
0.5416	0.0318057\\
0.5441	0.0317105\\
0.5467	0.0316132\\
0.5493	0.0315179\\
0.5518	0.0314283\\
0.5544	0.0313367\\
0.557	0.0312465\\
0.5595	0.0311611\\
0.5621	0.0310741\\
0.5647	0.0309893\\
0.5672	0.03091\\
0.5698	0.0308296\\
0.5723	0.030754\\
0.5749	0.030677\\
0.5775	0.0306019\\
0.58	0.0305318\\
0.5826	0.0304614\\
0.5852	0.0303934\\
0.5877	0.0303301\\
0.5903	0.0302661\\
0.5929	0.030204\\
0.5954	0.0301464\\
0.598	0.0300889\\
0.6006	0.0300341\\
0.6031	0.0299837\\
0.6057	0.0299334\\
0.6082	0.0298869\\
0.6108	0.0298406\\
0.6134	0.0297967\\
0.6159	0.029757\\
0.6185	0.0297183\\
0.6211	0.0296819\\
0.6236	0.0296487\\
0.6262	0.0296163\\
0.6288	0.0295862\\
0.6313	0.0295596\\
0.6339	0.0295345\\
0.6364	0.0295125\\
0.639	0.0294918\\
0.6416	0.029473\\
0.6441	0.0294569\\
0.6467	0.0294424\\
0.6493	0.0294303\\
0.6518	0.0294209\\
0.6544	0.0294132\\
0.657	0.0294073\\
0.6595	0.0294034\\
0.6621	0.0294013\\
0.6646	0.0294013\\
0.6672	0.0294034\\
0.6698	0.0294076\\
0.6723	0.0294132\\
0.6749	0.0294206\\
0.6775	0.0294296\\
0.68	0.0294401\\
0.6826	0.0294527\\
0.6852	0.0294672\\
0.6877	0.0294826\\
0.6903	0.0294999\\
0.6928	0.0295177\\
0.6954	0.0295376\\
0.698	0.029559\\
0.7005	0.0295809\\
0.7031	0.029605\\
0.7057	0.0296302\\
0.7082	0.0296552\\
0.7108	0.0296821\\
0.7134	0.0297101\\
0.7159	0.029738\\
0.7185	0.0297679\\
0.721	0.0297973\\
0.7236	0.0298284\\
0.7262	0.0298599\\
0.7287	0.0298907\\
0.7313	0.0299232\\
0.7339	0.0299563\\
0.7364	0.0299884\\
0.739	0.0300218\\
0.7416	0.0300552\\
0.7441	0.0300872\\
0.7467	0.0301205\\
0.7492	0.0301525\\
0.7518	0.0301856\\
0.7544	0.0302184\\
0.7569	0.0302493\\
0.7595	0.0302809\\
0.7621	0.0303118\\
0.7646	0.030341\\
0.7672	0.0303707\\
0.7698	0.0303995\\
0.7723	0.0304262\\
0.7749	0.0304528\\
0.7775	0.0304781\\
0.78	0.0305013\\
0.7826	0.0305241\\
0.7851	0.0305448\\
0.7877	0.0305647\\
0.7903	0.0305828\\
0.7928	0.0305985\\
0.7954	0.0306129\\
0.798	0.0306254\\
0.8005	0.0306356\\
0.8031	0.0306441\\
0.8057	0.0306502\\
0.8082	0.0306537\\
0.8108	0.0306549\\
0.8133	0.0306536\\
0.8159	0.0306497\\
0.8185	0.0306431\\
0.821	0.0306341\\
0.8236	0.0306216\\
0.8262	0.0306059\\
0.8287	0.0305878\\
0.8313	0.0305659\\
0.8339	0.0305406\\
0.8364	0.030513\\
0.839	0.0304807\\
0.8415	0.0304461\\
0.8441	0.0304064\\
0.8467	0.0303628\\
0.8492	0.0303172\\
0.8518	0.0302658\\
0.8544	0.0302102\\
0.8569	0.0301527\\
0.8595	0.0300885\\
0.8621	0.0300199\\
0.8646	0.0299496\\
0.8672	0.0298721\\
0.8697	0.0297931\\
0.8723	0.0297062\\
0.8749	0.0296143\\
0.8774	0.0295211\\
0.88	0.0294194\\
0.8826	0.0293125\\
0.8851	0.0292048\\
0.8877	0.0290876\\
0.8903	0.0289649\\
0.8928	0.0288418\\
0.8954	0.0287082\\
0.8979	0.0285745\\
0.9005	0.02843\\
0.9031	0.0282797\\
0.9056	0.0281296\\
0.9082	0.0279677\\
0.9108	0.0277998\\
0.9133	0.0276326\\
0.9159	0.0274529\\
0.9185	0.027267\\
0.921	0.0270824\\
0.9236	0.0268842\\
0.9262	0.0266796\\
0.9287	0.0264769\\
0.9313	0.0262598\\
0.9338	0.026045\\
0.9364	0.0258152\\
0.939	0.0255789\\
0.9415	0.0253453\\
0.9441	0.0250959\\
0.9467	0.0248398\\
0.9492	0.0245873\\
0.9518	0.0243181\\
0.9544	0.0240421\\
0.9569	0.0237703\\
0.9595	0.0234808\\
0.962	0.0231961\\
0.9646	0.0228932\\
0.9672	0.0225836\\
0.9697	0.0222793\\
0.9723	0.0219561\\
0.9749	0.021626\\
0.9774	0.021302\\
0.98	0.0209582\\
0.9826	0.0206076\\
0.9851	0.020264\\
0.9877	0.0198999\\
0.9902	0.0195432\\
0.9928	0.0191655\\
0.9954	0.0187809\\
0.9979	0.0184046\\
1.0005	0.0180066\\
1.0031	0.0176018\\
1.0056	0.0172061\\
1.0082	0.016788\\
1.0108	0.016363\\
1.0133	0.0159481\\
1.0159	0.0155101\\
1.0184	0.0150828\\
1.021	0.0146318\\
1.0236	0.0141743\\
1.0261	0.0137282\\
1.0287	0.0132579\\
1.0313	0.0127813\\
1.0338	0.012317\\
1.0364	0.0118281\\
1.039	0.0113328\\
1.0415	0.0108508\\
1.0441	0.0103435\\
1.0466	0.00985004\\
1.0492	0.00933102\\
1.0518	0.00880614\\
1.0543	0.00829596\\
1.0569	0.00775974\\
1.0595	0.00721783\\
1.062	0.00669152\\
1.0646	0.00613881\\
1.0672	0.00558074\\
1.0697	0.00503914\\
1.0723	0.00447075\\
1.0748	0.00391939\\
1.0774	0.00334105\\
1.08	0.0027578\\
1.0825	0.00219248\\
1.0851	0.00159995\\
1.0877	0.00100283\\
1.0902	0.000424418\\
1.0928	-0.000181433\\
1.0954	-0.000791545\\
1.0979	-0.00138208\\
1.1005	-0.00200016\\
1.1031	-0.00262214\\
1.1056	-0.00322378\\
1.1082	-0.00385309\\
1.1107	-0.00446155\\
1.1133	-0.00509768\\
1.1159	-0.00573708\\
1.1184	-0.00635483\\
1.121	-0.00700027\\
1.1236	-0.00764861\\
1.1261	-0.00827462\\
1.1287	-0.00892825\\
1.1313	-0.00958435\\
1.1338	-0.0102174\\
1.1364	-0.0108779\\
1.1389	-0.011515\\
1.1415	-0.0121795\\
1.1441	-0.0128457\\
1.1466	-0.0134879\\
1.1492	-0.0141573\\
1.1518	-0.0148279\\
1.1543	-0.0154739\\
1.1569	-0.0161468\\
1.1595	-0.0168207\\
1.162	-0.0174693\\
1.1646	-0.0181445\\
1.1671	-0.0187942\\
1.1697	-0.0194701\\
1.1723	-0.0201463\\
1.1748	-0.0207965\\
1.1774	-0.0214726\\
1.18	-0.0221485\\
1.1825	-0.022798\\
1.1851	-0.0234729\\
1.1877	-0.0241471\\
1.1902	-0.0247946\\
1.1928	-0.025467\\
1.1953	-0.0261126\\
1.1979	-0.0267826\\
1.2005	-0.0274513\\
1.203	-0.0280928\\
1.2056	-0.0287583\\
1.2082	-0.0294219\\
1.2107	-0.0300582\\
1.2133	-0.0307178\\
1.2159	-0.0313752\\
1.2184	-0.032005\\
1.221	-0.0326576\\
1.2235	-0.0332825\\
1.2261	-0.0339298\\
1.2287	-0.0345741\\
1.2312	-0.0351907\\
1.2338	-0.0358289\\
1.2364	-0.0364638\\
1.2389	-0.0370711\\
1.2415	-0.0376992\\
1.2441	-0.0383236\\
1.2466	-0.0389204\\
1.2492	-0.0395373\\
1.2518	-0.0401501\\
1.2543	-0.0407355\\
1.2569	-0.0413401\\
1.2594	-0.0419174\\
1.262	-0.0425134\\
1.2646	-0.0431048\\
1.2671	-0.0436691\\
1.2697	-0.0442513\\
1.2723	-0.0448286\\
1.2748	-0.045379\\
1.2774	-0.0459465\\
1.28	-0.0465089\\
1.2825	-0.0470447\\
1.2851	-0.0475967\\
1.2876	-0.0481224\\
1.2902	-0.0486637\\
1.2928	-0.0491995\\
1.2953	-0.0497093\\
1.2979	-0.0502339\\
1.3005	-0.0507527\\
1.303	-0.0512461\\
1.3056	-0.0517533\\
1.3082	-0.0522546\\
1.3107	-0.0527309\\
1.3133	-0.0532203\\
1.3158	-0.053685\\
1.3184	-0.0541621\\
1.321	-0.054633\\
1.3235	-0.0550798\\
1.3261	-0.0555383\\
1.3287	-0.0559903\\
1.3312	-0.0564188\\
1.3338	-0.0568581\\
1.3364	-0.0572909\\
1.3389	-0.0577008\\
1.3415	-0.0581206\\
1.344	-0.058518\\
1.3466	-0.0589248\\
1.3492	-0.0593248\\
1.3517	-0.0597032\\
1.3543	-0.06009\\
1.3569	-0.0604702\\
1.3594	-0.0608293\\
1.362	-0.0611961\\
1.3646	-0.0615561\\
1.3671	-0.0618959\\
1.3697	-0.0622425\\
1.3722	-0.0625694\\
1.3748	-0.0629027\\
1.3774	-0.0632292\\
1.3799	-0.0635366\\
1.3825	-0.0638497\\
1.3851	-0.064156\\
1.3876	-0.0644441\\
1.3902	-0.0647371\\
1.3928	-0.0650232\\
1.3953	-0.0652921\\
1.3979	-0.065565\\
1.4005	-0.0658312\\
1.403	-0.0660809\\
1.4056	-0.0663339\\
1.4081	-0.066571\\
1.4107	-0.0668111\\
1.4133	-0.0670446\\
1.4158	-0.0672629\\
1.4184	-0.0674835\\
1.421	-0.0676976\\
1.4235	-0.0678974\\
1.4261	-0.0680988\\
1.4287	-0.0682939\\
1.4312	-0.0684755\\
1.4338	-0.0686582\\
1.4363	-0.0688279\\
1.4389	-0.0689984\\
1.4415	-0.0691627\\
1.444	-0.0693149\\
1.4466	-0.0694673\\
1.4492	-0.0696136\\
1.4517	-0.0697487\\
1.4543	-0.0698834\\
1.4569	-0.0700123\\
1.4594	-0.0701308\\
1.462	-0.0702484\\
1.4645	-0.0703562\\
1.4671	-0.0704628\\
1.4697	-0.0705638\\
1.4722	-0.0706559\\
1.4748	-0.0707463\\
1.4774	-0.0708314\\
1.4799	-0.0709082\\
1.4825	-0.0709831\\
1.4851	-0.0710528\\
1.4876	-0.0711151\\
1.4902	-0.071175\\
1.4927	-0.071228\\
1.4953	-0.0712784\\
1.4979	-0.0713241\\
1.5004	-0.0713636\\
1.503	-0.0714003\\
1.5056	-0.0714324\\
1.5081	-0.0714591\\
1.5107	-0.0714827\\
1.5133	-0.071502\\
1.5158	-0.0715166\\
1.5184	-0.0715279\\
1.5209	-0.0715349\\
1.5235	-0.0715384\\
1.5261	-0.0715381\\
1.5286	-0.0715343\\
1.5312	-0.0715267\\
1.5338	-0.0715156\\
1.5363	-0.0715017\\
1.5389	-0.0714838\\
1.5415	-0.0714628\\
1.544	-0.0714395\\
1.5466	-0.0714122\\
1.5491	-0.0713832\\
1.5517	-0.0713502\\
1.5543	-0.0713143\\
1.5568	-0.0712772\\
1.5594	-0.0712361\\
1.562	-0.0711925\\
1.5645	-0.0711482\\
1.5671	-0.0710998\\
1.5697	-0.0710492\\
1.5722	-0.0709985\\
1.5748	-0.0709437\\
1.5774	-0.070887\\
1.5799	-0.0708307\\
1.5825	-0.0707704\\
1.585	-0.0707108\\
1.5876	-0.0706473\\
1.5902	-0.0705822\\
1.5927	-0.0705184\\
1.5953	-0.0704507\\
1.5979	-0.0703817\\
1.6004	-0.0703144\\
1.603	-0.0702434\\
1.6056	-0.0701714\\
1.6081	-0.0701014\\
1.6107	-0.0700279\\
1.6132	-0.0699565\\
1.6158	-0.0698818\\
1.6184	-0.0698066\\
1.6209	-0.0697339\\
1.6235	-0.069658\\
1.6261	-0.0695818\\
1.6286	-0.0695085\\
1.6312	-0.0694323\\
1.6338	-0.069356\\
1.6363	-0.0692829\\
1.6389	-0.069207\\
1.6414	-0.0691344\\
1.644	-0.0690592\\
1.6466	-0.0689845\\
1.6491	-0.0689132\\
1.6517	-0.0688396\\
1.6543	-0.0687668\\
1.6568	-0.0686975\\
1.6594	-0.0686264\\
1.662	-0.0685562\\
1.6645	-0.0684896\\
1.6671	-0.0684215\\
1.6696	-0.0683571\\
1.6722	-0.0682914\\
1.6748	-0.068227\\
1.6773	-0.0681663\\
1.6799	-0.0681047\\
1.6825	-0.0680447\\
1.685	-0.0679884\\
1.6876	-0.0679316\\
1.6902	-0.0678765\\
1.6927	-0.0678252\\
1.6953	-0.0677736\\
1.6978	-0.0677259\\
1.7004	-0.0676782\\
1.703	-0.0676325\\
1.7055	-0.0675906\\
1.7081	-0.0675491\\
1.7107	-0.0675098\\
1.7132	-0.0674741\\
1.7158	-0.0674393\\
1.7184	-0.0674068\\
1.7209	-0.0673779\\
1.7235	-0.0673502\\
1.7261	-0.067325\\
1.7286	-0.0673031\\
1.7312	-0.0672829\\
1.7337	-0.067266\\
1.7363	-0.067251\\
1.7389	-0.0672387\\
1.7414	-0.0672295\\
1.744	-0.0672226\\
1.7466	-0.0672185\\
1.7491	-0.0672172\\
1.7517	-0.0672187\\
1.7543	-0.067223\\
1.7568	-0.06723\\
1.7594	-0.0672401\\
1.7619	-0.0672526\\
1.7645	-0.0672685\\
1.7671	-0.0672874\\
1.7696	-0.0673085\\
1.7722	-0.0673333\\
1.7748	-0.0673612\\
1.7773	-0.0673909\\
1.7799	-0.0674248\\
1.7825	-0.0674618\\
1.785	-0.0675003\\
1.7876	-0.0675433\\
1.7901	-0.0675876\\
1.7927	-0.0676367\\
1.7953	-0.0676889\\
1.7978	-0.067742\\
1.8004	-0.0678002\\
1.803	-0.0678616\\
1.8055	-0.0679234\\
1.8081	-0.0679908\\
1.8107	-0.0680612\\
1.8132	-0.0681318\\
1.8158	-0.0682082\\
1.8183	-0.0682845\\
1.8209	-0.0683667\\
1.8235	-0.068452\\
1.826	-0.0685368\\
1.8286	-0.0686278\\
1.8312	-0.0687218\\
1.8337	-0.0688149\\
1.8363	-0.0689145\\
1.8389	-0.069017\\
1.8414	-0.0691181\\
1.844	-0.0692261\\
1.8465	-0.0693324\\
1.8491	-0.0694457\\
1.8517	-0.0695617\\
1.8542	-0.0696756\\
1.8568	-0.0697967\\
1.8594	-0.0699204\\
1.8619	-0.0700416\\
1.8645	-0.0701702\\
1.8671	-0.0703011\\
1.8696	-0.0704293\\
1.8722	-0.0705649\\
1.8747	-0.0706975\\
1.8773	-0.0708375\\
1.8799	-0.0709798\\
1.8824	-0.0711185\\
1.885	-0.0712649\\
1.8876	-0.0714133\\
1.8901	-0.0715579\\
1.8927	-0.0717101\\
1.8953	-0.0718642\\
1.8978	-0.072014\\
1.9004	-0.0721716\\
1.903	-0.0723308\\
1.9055	-0.0724855\\
1.9081	-0.0726478\\
1.9106	-0.0728054\\
1.9132	-0.0729706\\
1.9158	-0.0731372\\
1.9183	-0.0732986\\
1.9209	-0.0734677\\
1.9235	-0.0736379\\
1.926	-0.0738027\\
1.9286	-0.073975\\
1.9312	-0.0741483\\
1.9337	-0.0743158\\
1.9363	-0.0744907\\
1.9388	-0.0746597\\
1.9414	-0.0748361\\
1.944	-0.075013\\
1.9465	-0.0751837\\
1.9491	-0.0753617\\
1.9517	-0.07554\\
1.9542	-0.0757118\\
1.9568	-0.0758907\\
1.9594	-0.0760698\\
1.9619	-0.0762421\\
1.9645	-0.0764213\\
1.967	-0.0765935\\
1.9696	-0.0767725\\
1.9722	-0.0769514\\
1.9747	-0.077123\\
1.9773	-0.0773013\\
1.9799	-0.0774791\\
1.9824	-0.0776496\\
1.985	-0.0778263\\
1.9876	-0.0780025\\
1.9901	-0.0781712\\
1.9927	-0.0783458\\
1.9952	-0.078513\\
1.9978	-0.0786859\\
2.0004	-0.0788578\\
2.0029	-0.0790221\\
2.0055	-0.0791919\\
2.0081	-0.0793604\\
2.0106	-0.0795213\\
2.0132	-0.0796873\\
2.0158	-0.0798518\\
2.0183	-0.0800087\\
2.0209	-0.0801703\\
2.0234	-0.0803241\\
2.026	-0.0804825\\
2.0286	-0.0806391\\
2.0311	-0.080788\\
2.0337	-0.080941\\
2.0363	-0.0810921\\
2.0388	-0.0812355\\
2.0414	-0.0813826\\
2.044	-0.0815276\\
2.0465	-0.081665\\
2.0491	-0.0818057\\
2.0517	-0.0819441\\
2.0542	-0.082075\\
2.0568	-0.0822088\\
2.0593	-0.0823351\\
2.0619	-0.082464\\
2.0645	-0.0825904\\
2.067	-0.0827095\\
2.0696	-0.0828307\\
2.0722	-0.0829492\\
2.0747	-0.0830607\\
2.0773	-0.0831738\\
2.0799	-0.0832841\\
2.0824	-0.0833874\\
2.085	-0.0834921\\
2.0875	-0.0835899\\
2.0901	-0.0836887\\
2.0927	-0.0837845\\
2.0952	-0.0838737\\
2.0978	-0.0839635\\
2.1004	-0.0840501\\
2.1029	-0.0841304\\
2.1055	-0.0842108\\
2.1081	-0.084288\\
2.1106	-0.0843591\\
2.1132	-0.0844299\\
2.1157	-0.0844948\\
2.1183	-0.0845591\\
2.1209	-0.0846201\\
2.1234	-0.0846755\\
2.126	-0.0847298\\
2.1286	-0.0847807\\
2.1311	-0.0848265\\
2.1337	-0.0848707\\
2.1363	-0.0849114\\
2.1388	-0.0849473\\
2.1414	-0.0849812\\
2.1439	-0.0850106\\
2.1465	-0.0850377\\
2.1491	-0.0850613\\
2.1516	-0.0850807\\
2.1542	-0.0850974\\
2.1568	-0.0851106\\
2.1593	-0.0851199\\
2.1619	-0.0851262\\
2.1645	-0.085129\\
2.167	-0.0851284\\
2.1696	-0.0851242\\
2.1721	-0.085117\\
2.1747	-0.085106\\
2.1773	-0.0850914\\
2.1798	-0.0850742\\
2.1824	-0.0850529\\
2.185	-0.0850281\\
2.1875	-0.0850009\\
2.1901	-0.0849693\\
2.1927	-0.0849343\\
2.1952	-0.0848974\\
2.1978	-0.0848557\\
2.2004	-0.0848105\\
2.2029	-0.084764\\
2.2055	-0.0847122\\
2.208	-0.0846594\\
2.2106	-0.0846011\\
2.2132	-0.0845396\\
2.2157	-0.0844773\\
2.2183	-0.0844094\\
2.2209	-0.0843383\\
2.2234	-0.0842668\\
2.226	-0.0841894\\
2.2286	-0.0841089\\
2.2311	-0.0840285\\
2.2337	-0.0839419\\
2.2362	-0.0838557\\
2.2388	-0.0837631\\
2.2414	-0.0836675\\
2.2439	-0.0835728\\
2.2465	-0.0834715\\
2.2491	-0.0833672\\
2.2516	-0.0832642\\
2.2542	-0.0831543\\
2.2568	-0.0830417\\
2.2593	-0.0829308\\
2.2619	-0.0828127\\
2.2644	-0.0826967\\
2.267	-0.0825734\\
2.2696	-0.0824475\\
2.2721	-0.0823241\\
2.2747	-0.0821932\\
2.2773	-0.0820598\\
2.2798	-0.0819292\\
2.2824	-0.0817911\\
2.285	-0.0816506\\
2.2875	-0.0815133\\
2.2901	-0.0813682\\
2.2926	-0.0812267\\
2.2952	-0.0810773\\
2.2978	-0.0809257\\
2.3003	-0.080778\\
2.3029	-0.0806223\\
2.3055	-0.0804647\\
2.308	-0.0803112\\
2.3106	-0.0801497\\
2.3132	-0.0799863\\
2.3157	-0.0798274\\
2.3183	-0.0796605\\
2.3208	-0.0794983\\
2.3234	-0.0793279\\
2.326	-0.0791559\\
2.3285	-0.078989\\
2.3311	-0.0788139\\
2.3337	-0.0786373\\
2.3362	-0.0784661\\
2.3388	-0.0782866\\
2.3414	-0.0781058\\
2.3439	-0.0779307\\
2.3465	-0.0777473\\
2.349	-0.0775698\\
2.3516	-0.0773841\\
2.3542	-0.0771973\\
2.3567	-0.0770166\\
2.3593	-0.0768276\\
2.3619	-0.0766377\\
2.3644	-0.0764542\\
2.367	-0.0762625\\
2.3696	-0.07607\\
2.3721	-0.0758841\\
2.3747	-0.07569\\
2.3773	-0.0754953\\
2.3798	-0.0753074\\
2.3824	-0.0751115\\
2.3849	-0.0749225\\
2.3875	-0.0747255\\
2.3901	-0.0745281\\
2.3926	-0.0743378\\
2.3952	-0.0741396\\
2.3978	-0.0739411\\
2.4003	-0.07375\\
2.4029	-0.073551\\
2.4055	-0.0733518\\
2.408	-0.0731602\\
2.4106	-0.0729608\\
2.4131	-0.0727691\\
2.4157	-0.0725697\\
2.4183	-0.0723704\\
2.4208	-0.0721788\\
2.4234	-0.0719797\\
2.426	-0.0717808\\
2.4285	-0.0715898\\
2.4311	-0.0713914\\
2.4337	-0.0711933\\
2.4362	-0.0710032\\
2.4388	-0.0708059\\
2.4413	-0.0706166\\
2.4439	-0.0704202\\
2.4465	-0.0702243\\
2.449	-0.0700366\\
2.4516	-0.0698418\\
2.4542	-0.0696478\\
2.4567	-0.0694619\\
2.4593	-0.0692692\\
2.4619	-0.0690773\\
2.4644	-0.0688935\\
2.467	-0.0687033\\
2.4695	-0.0685211\\
2.4721	-0.0683326\\
2.4747	-0.068145\\
2.4772	-0.0679656\\
2.4798	-0.0677799\\
2.4824	-0.0675954\\
2.4849	-0.0674189\\
2.4875	-0.0672365\\
2.4901	-0.0670552\\
2.4926	-0.066882\\
2.4952	-0.0667031\\
2.4977	-0.0665322\\
2.5003	-0.0663557\\
2.5029	-0.0661805\\
2.5054	-0.0660133\\
2.508	-0.0658407\\
2.5106	-0.0656694\\
2.5131	-0.0655061\\
2.5157	-0.0653376\\
2.5183	-0.0651706\\
2.5208	-0.0650113\\
2.5234	-0.0648472\\
2.526	-0.0646845\\
2.5285	-0.0645296\\
2.5311	-0.0643699\\
2.5336	-0.0642179\\
2.5362	-0.0640613\\
2.5388	-0.0639063\\
2.5413	-0.0637588\\
2.5439	-0.063607\\
2.5465	-0.0634568\\
2.549	-0.0633139\\
2.5516	-0.063167\\
2.5542	-0.0630217\\
2.5567	-0.0628836\\
2.5593	-0.0627416\\
2.5618	-0.0626067\\
2.5644	-0.062468\\
2.567	-0.0623311\\
2.5695	-0.0622011\\
2.5721	-0.0620675\\
2.5747	-0.0619357\\
2.5772	-0.0618105\\
2.5798	-0.0616821\\
2.5824	-0.0615554\\
2.5849	-0.0614351\\
2.5875	-0.0613118\\
2.59	-0.0611949\\
2.5926	-0.0610749\\
2.5952	-0.0609567\\
2.5977	-0.0608446\\
2.6003	-0.0607298\\
2.6029	-0.0606166\\
2.6054	-0.0605094\\
2.608	-0.0603996\\
2.6106	-0.0602915\\
2.6131	-0.0601891\\
2.6157	-0.0600842\\
2.6182	-0.0599849\\
2.6208	-0.0598832\\
2.6234	-0.0597832\\
2.6259	-0.0596886\\
2.6285	-0.0595917\\
2.6311	-0.0594964\\
2.6336	-0.0594063\\
2.6362	-0.059314\\
2.6388	-0.0592233\\
2.6413	-0.0591375\\
2.6439	-0.0590498\\
2.6464	-0.0589668\\
2.649	-0.058882\\
2.6516	-0.0587985\\
2.6541	-0.0587196\\
2.6567	-0.058639\\
2.6593	-0.0585596\\
2.6618	-0.0584846\\
2.6644	-0.058408\\
2.667	-0.0583326\\
2.6695	-0.0582613\\
2.6721	-0.0581883\\
2.6746	-0.0581194\\
2.6772	-0.0580488\\
2.6798	-0.0579794\\
2.6823	-0.0579137\\
2.6849	-0.0578466\\
2.6875	-0.0577804\\
2.69	-0.0577179\\
2.6926	-0.0576538\\
2.6952	-0.0575907\\
2.6977	-0.057531\\
2.7003	-0.0574698\\
2.7029	-0.0574095\\
2.7054	-0.0573523\\
2.708	-0.0572937\\
2.7105	-0.0572381\\
2.7131	-0.057181\\
2.7157	-0.0571246\\
2.7182	-0.0570712\\
2.7208	-0.0570162\\
2.7234	-0.0569619\\
2.7259	-0.0569103\\
2.7285	-0.0568571\\
2.7311	-0.0568046\\
2.7336	-0.0567545\\
2.7362	-0.0567029\\
2.7387	-0.0566537\\
2.7413	-0.056603\\
2.7439	-0.0565527\\
2.7464	-0.0565046\\
2.749	-0.0564549\\
2.7516	-0.0564055\\
2.7541	-0.0563583\\
2.7567	-0.0563094\\
2.7593	-0.0562607\\
2.7618	-0.056214\\
2.7644	-0.0561656\\
2.7669	-0.0561192\\
2.7695	-0.056071\\
2.7721	-0.0560228\\
2.7746	-0.0559765\\
2.7772	-0.0559283\\
2.7798	-0.05588\\
2.7823	-0.0558336\\
2.7849	-0.0557851\\
2.7875	-0.0557366\\
2.79	-0.0556897\\
2.7926	-0.0556408\\
2.7951	-0.0555936\\
2.7977	-0.0555443\\
2.8003	-0.0554947\\
2.8028	-0.0554467\\
2.8054	-0.0553965\\
2.808	-0.0553459\\
2.8105	-0.055297\\
2.8131	-0.0552457\\
2.8157	-0.055194\\
2.8182	-0.0551438\\
2.8208	-0.0550912\\
2.8233	-0.0550402\\
2.8259	-0.0549867\\
2.8285	-0.0549326\\
2.831	-0.0548802\\
2.8336	-0.054825\\
2.8362	-0.0547694\\
2.8387	-0.0547153\\
2.8413	-0.0546584\\
2.8439	-0.0546009\\
2.8464	-0.0545451\\
2.849	-0.0544863\\
2.8516	-0.0544269\\
2.8541	-0.0543692\\
2.8567	-0.0543085\\
2.8592	-0.0542495\\
2.8618	-0.0541874\\
2.8644	-0.0541246\\
2.8669	-0.0540635\\
2.8695	-0.0539993\\
2.8721	-0.0539343\\
2.8746	-0.0538712\\
2.8772	-0.0538047\\
2.8798	-0.0537376\\
2.8823	-0.0536723\\
2.8849	-0.0536036\\
2.8874	-0.0535369\\
2.89	-0.0534668\\
2.8926	-0.0533959\\
2.8951	-0.053327\\
2.8977	-0.0532546\\
2.9003	-0.0531814\\
2.9028	-0.0531104\\
2.9054	-0.0530357\\
2.908	-0.0529604\\
2.9105	-0.0528872\\
2.9131	-0.0528104\\
2.9156	-0.0527358\\
2.9182	-0.0526576\\
2.9208	-0.0525786\\
2.9233	-0.052502\\
2.9259	-0.0524217\\
2.9285	-0.0523407\\
2.931	-0.0522621\\
2.9336	-0.0521798\\
2.9362	-0.0520968\\
2.9387	-0.0520165\\
2.9413	-0.0519323\\
2.9438	-0.0518507\\
2.9464	-0.0517653\\
2.949	-0.0516794\\
2.9515	-0.0515962\\
2.9541	-0.0515092\\
2.9567	-0.0514216\\
2.9592	-0.0513369\\
2.9618	-0.0512484\\
2.9644	-0.0511594\\
2.9669	-0.0510734\\
2.9695	-0.0509836\\
2.972	-0.0508968\\
2.9746	-0.0508062\\
2.9772	-0.0507152\\
2.9797	-0.0506274\\
2.9823	-0.0505358\\
2.9849	-0.0504439\\
2.9874	-0.0503552\\
2.99	-0.0502628\\
2.9926	-0.0501702\\
2.9951	-0.050081\\
2.9977	-0.049988\\
3.0003	-0.049895\\
3.0028	-0.0498053\\
3.0054	-0.0497121\\
3.0079	-0.0496223\\
3.0105	-0.049529\\
3.0131	-0.0494356\\
3.0156	-0.0493459\\
3.0182	-0.0492527\\
3.0208	-0.0491595\\
3.0233	-0.04907\\
3.0259	-0.0489772\\
3.0285	-0.0488844\\
3.031	-0.0487955\\
3.0336	-0.0487032\\
3.0361	-0.0486148\\
3.0387	-0.0485231\\
3.0413	-0.0484317\\
3.0438	-0.0483442\\
3.0464	-0.0482536\\
3.049	-0.0481634\\
3.0515	-0.0480772\\
3.0541	-0.0479879\\
3.0567	-0.0478992\\
3.0592	-0.0478144\\
3.0618	-0.0477268\\
3.0643	-0.0476432\\
3.0669	-0.0475569\\
3.0695	-0.0474712\\
3.072	-0.0473896\\
3.0746	-0.0473054\\
3.0772	-0.047222\\
3.0797	-0.0471425\\
3.0823	-0.0470607\\
3.0849	-0.0469798\\
3.0874	-0.0469029\\
3.09	-0.0468238\\
3.0925	-0.0467487\\
3.0951	-0.0466716\\
3.0977	-0.0465955\\
3.1002	-0.0465234\\
3.1028	-0.0464494\\
3.1054	-0.0463766\\
3.1079	-0.0463077\\
3.1105	-0.0462372\\
3.1131	-0.046168\\
3.1156	-0.0461025\\
3.1182	-0.0460358\\
3.1207	-0.0459728\\
3.1233	-0.0459087\\
3.1259	-0.0458459\\
3.1284	-0.0457869\\
3.131	-0.0457269\\
3.1336	-0.0456684\\
3.1361	-0.0456136\\
3.1387	-0.045558\\
3.1413	-0.045504\\
3.1438	-0.0454536\\
3.1464	-0.0454027\\
3.1489	-0.0453553\\
3.1515	-0.0453077\\
3.1541	-0.0452618\\
3.1566	-0.0452192\\
3.1592	-0.0451766\\
3.1618	-0.0451358\\
3.1643	-0.0450983\\
3.1669	-0.045061\\
3.1695	-0.0450256\\
3.172	-0.0449932\\
3.1746	-0.0449615\\
3.1772	-0.0449316\\
3.1797	-0.0449047\\
3.1823	-0.0448787\\
3.1848	-0.0448555\\
3.1874	-0.0448333\\
3.19	-0.0448131\\
3.1925	-0.0447957\\
3.1951	-0.0447795\\
3.1977	-0.0447654\\
3.2002	-0.0447539\\
3.2028	-0.0447439\\
3.2054	-0.0447361\\
3.2079	-0.0447305\\
3.2105	-0.0447269\\
3.213	-0.0447255\\
3.2156	-0.0447261\\
3.2182	-0.0447289\\
3.2207	-0.0447337\\
3.2233	-0.0447409\\
3.2259	-0.0447504\\
3.2284	-0.0447616\\
3.231	-0.0447754\\
3.2336	-0.0447916\\
3.2361	-0.0448092\\
3.2387	-0.0448299\\
3.2412	-0.0448519\\
3.2438	-0.044877\\
3.2464	-0.0449045\\
3.2489	-0.0449331\\
3.2515	-0.0449651\\
3.2541	-0.0449995\\
3.2566	-0.0450348\\
3.2592	-0.0450737\\
3.2618	-0.045115\\
3.2643	-0.045157\\
3.2669	-0.0452029\\
3.2694	-0.0452493\\
3.272	-0.0452999\\
3.2746	-0.0453528\\
3.2771	-0.0454059\\
3.2797	-0.0454634\\
3.2823	-0.0455233\\
3.2848	-0.0455831\\
3.2874	-0.0456476\\
3.29	-0.0457144\\
3.2925	-0.0457808\\
3.2951	-0.0458522\\
3.2976	-0.0459231\\
3.3002	-0.045999\\
3.3028	-0.0460773\\
3.3053	-0.0461547\\
3.3079	-0.0462375\\
3.3105	-0.0463225\\
3.313	-0.0464064\\
3.3156	-0.0464959\\
3.3182	-0.0465876\\
3.3207	-0.0466779\\
3.3233	-0.046774\\
3.3259	-0.0468723\\
3.3284	-0.0469689\\
3.331	-0.0470714\\
3.3335	-0.0471721\\
3.3361	-0.0472788\\
3.3387	-0.0473877\\
3.3412	-0.0474944\\
3.3438	-0.0476074\\
3.3464	-0.0477224\\
3.3489	-0.0478349\\
3.3515	-0.0479539\\
3.3541	-0.0480749\\
3.3566	-0.0481931\\
3.3592	-0.0483179\\
3.3617	-0.0484398\\
3.3643	-0.0485683\\
3.3669	-0.0486987\\
3.3694	-0.0488259\\
3.372	-0.0489599\\
3.3746	-0.0490957\\
3.3771	-0.0492279\\
3.3797	-0.0493671\\
3.3823	-0.049508\\
3.3848	-0.049645\\
3.3874	-0.0497891\\
3.3899	-0.0499291\\
3.3925	-0.0500764\\
3.3951	-0.0502251\\
3.3976	-0.0503695\\
3.4002	-0.0505211\\
3.4028	-0.0506742\\
3.4053	-0.0508227\\
3.4079	-0.0509784\\
3.4105	-0.0511355\\
3.413	-0.0512877\\
3.4156	-0.0514472\\
3.4181	-0.0516017\\
3.4207	-0.0517635\\
3.4233	-0.0519265\\
3.4258	-0.0520842\\
3.4284	-0.0522492\\
3.431	-0.0524153\\
3.4335	-0.0525758\\
3.4361	-0.0527437\\
3.4387	-0.0529124\\
3.4412	-0.0530755\\
3.4438	-0.0532458\\
3.4463	-0.0534102\\
3.4489	-0.053582\\
3.4515	-0.0537544\\
3.454	-0.0539207\\
3.4566	-0.0540943\\
3.4592	-0.0542683\\
3.4617	-0.0544361\\
3.4643	-0.0546111\\
3.4669	-0.0547864\\
3.4694	-0.0549554\\
3.472	-0.0551313\\
3.4745	-0.0553007\\
3.4771	-0.0554771\\
3.4797	-0.0556537\\
3.4822	-0.0558235\\
3.4848	-0.0560002\\
3.4874	-0.0561769\\
3.4899	-0.0563468\\
3.4925	-0.0565234\\
3.4951	-0.0566999\\
3.4976	-0.0568694\\
3.5002	-0.0570455\\
3.5028	-0.0572213\\
3.5053	-0.05739\\
3.5079	-0.0575651\\
3.5104	-0.0577332\\
3.513	-0.0579074\\
3.5156	-0.0580812\\
3.5181	-0.0582478\\
3.5207	-0.0584204\\
3.5233	-0.0585925\\
3.5258	-0.0587572\\
3.5284	-0.0589278\\
3.531	-0.0590976\\
3.5335	-0.0592602\\
3.5361	-0.0594283\\
3.5386	-0.0595892\\
3.5412	-0.0597555\\
3.5438	-0.0599209\\
3.5463	-0.0600788\\
3.5489	-0.0602421\\
3.5515	-0.0604042\\
3.554	-0.060559\\
3.5566	-0.0607187\\
3.5592	-0.0608772\\
3.5617	-0.0610284\\
3.5643	-0.0611844\\
3.5668	-0.061333\\
3.5694	-0.0614862\\
3.572	-0.0616379\\
3.5745	-0.0617824\\
3.5771	-0.0619311\\
3.5797	-0.0620783\\
3.5822	-0.0622183\\
3.5848	-0.0623623\\
3.5874	-0.0625046\\
3.5899	-0.0626399\\
3.5925	-0.0627788\\
3.595	-0.0629107\\
3.5976	-0.0630461\\
3.6002	-0.0631796\\
3.6027	-0.0633063\\
3.6053	-0.0634361\\
3.6079	-0.063564\\
3.6104	-0.0636851\\
3.613	-0.0638091\\
3.6156	-0.0639311\\
3.6181	-0.0640464\\
3.6207	-0.0641643\\
3.6232	-0.0642757\\
3.6258	-0.0643895\\
3.6284	-0.0645011\\
3.6309	-0.0646064\\
3.6335	-0.0647137\\
3.6361	-0.0648188\\
3.6386	-0.0649177\\
3.6412	-0.0650183\\
3.6438	-0.0651167\\
3.6463	-0.0652091\\
3.6489	-0.065303\\
3.6515	-0.0653945\\
3.654	-0.0654803\\
3.6566	-0.0655671\\
3.6591	-0.0656484\\
3.6617	-0.0657306\\
3.6643	-0.0658104\\
3.6668	-0.0658848\\
3.6694	-0.0659598\\
3.672	-0.0660323\\
3.6745	-0.0660998\\
3.6771	-0.0661675\\
3.6797	-0.0662328\\
3.6822	-0.0662932\\
3.6848	-0.0663536\\
3.6873	-0.0664094\\
3.6899	-0.0664649\\
3.6925	-0.0665179\\
3.695	-0.0665665\\
3.6976	-0.0666147\\
3.7002	-0.0666603\\
3.7027	-0.0667018\\
3.7053	-0.0667425\\
3.7079	-0.0667807\\
3.7104	-0.0668151\\
3.713	-0.0668484\\
3.7155	-0.0668781\\
3.7181	-0.0669065\\
3.7207	-0.0669325\\
3.7232	-0.0669551\\
3.7258	-0.0669762\\
3.7284	-0.0669948\\
3.7309	-0.0670104\\
3.7335	-0.0670243\\
3.7361	-0.0670357\\
3.7386	-0.0670444\\
3.7412	-0.067051\\
3.7437	-0.0670552\\
3.7463	-0.0670572\\
3.7489	-0.0670568\\
3.7514	-0.0670542\\
3.754	-0.0670492\\
3.7566	-0.0670419\\
3.7591	-0.0670327\\
3.7617	-0.0670209\\
3.7643	-0.0670068\\
3.7668	-0.0669911\\
3.7694	-0.0669726\\
3.7719	-0.0669527\\
3.7745	-0.0669298\\
3.7771	-0.0669048\\
3.7796	-0.0668788\\
3.7822	-0.0668496\\
3.7848	-0.0668182\\
3.7873	-0.0667862\\
3.7899	-0.0667508\\
3.7925	-0.0667134\\
3.795	-0.0666755\\
3.7976	-0.0666342\\
3.8002	-0.0665909\\
3.8027	-0.0665475\\
3.8053	-0.0665004\\
3.8078	-0.0664534\\
3.8104	-0.0664027\\
3.813	-0.0663502\\
3.8155	-0.066298\\
3.8181	-0.066242\\
3.8207	-0.0661842\\
3.8232	-0.0661271\\
3.8258	-0.0660661\\
3.8284	-0.0660034\\
3.8309	-0.0659416\\
3.8335	-0.0658758\\
3.836	-0.065811\\
3.8386	-0.0657422\\
3.8412	-0.0656719\\
3.8437	-0.065603\\
3.8463	-0.0655299\\
3.8489	-0.0654554\\
3.8514	-0.0653826\\
3.854	-0.0653055\\
3.8566	-0.0652271\\
3.8591	-0.0651506\\
3.8617	-0.0650699\\
3.8642	-0.0649911\\
3.8668	-0.0649081\\
3.8694	-0.064824\\
3.8719	-0.0647421\\
3.8745	-0.0646559\\
3.8771	-0.0645688\\
3.8796	-0.0644841\\
3.8822	-0.0643951\\
3.8848	-0.0643052\\
3.8873	-0.0642179\\
3.8899	-0.0641264\\
3.8924	-0.0640377\\
3.895	-0.0639447\\
3.8976	-0.063851\\
3.9001	-0.0637603\\
3.9027	-0.0636653\\
3.9053	-0.0635698\\
3.9078	-0.0634774\\
3.9104	-0.0633808\\
3.913	-0.0632837\\
3.9155	-0.0631899\\
3.9181	-0.063092\\
3.9206	-0.0629975\\
3.9232	-0.0628989\\
3.9258	-0.0627999\\
3.9283	-0.0627046\\
3.9309	-0.0626052\\
3.9335	-0.0625056\\
3.936	-0.0624097\\
3.9386	-0.0623099\\
3.9412	-0.0622099\\
3.9437	-0.0621138\\
3.9463	-0.0620138\\
3.9488	-0.0619177\\
3.9514	-0.0618178\\
3.954	-0.061718\\
3.9565	-0.0616221\\
3.9591	-0.0615226\\
3.9617	-0.0614232\\
3.9642	-0.0613279\\
3.9668	-0.0612291\\
3.9694	-0.0611305\\
3.9719	-0.061036\\
3.9745	-0.0609381\\
3.9771	-0.0608406\\
3.9796	-0.0607473\\
3.9822	-0.0606506\\
3.9847	-0.0605582\\
3.9873	-0.0604625\\
3.9899	-0.0603674\\
3.9924	-0.0602765\\
3.995	-0.0601825\\
3.9976	-0.0600892\\
4.0001	-0.0600001\\
4.0027	-0.0599081\\
4.0053	-0.0598168\\
4.0078	-0.0597298\\
4.0104	-0.05964\\
4.0129	-0.0595544\\
4.0155	-0.0594662\\
4.0181	-0.0593789\\
4.0206	-0.0592957\\
4.0232	-0.0592101\\
4.0258	-0.0591255\\
4.0283	-0.059045\\
4.0309	-0.0589622\\
4.0335	-0.0588804\\
4.036	-0.0588027\\
4.0386	-0.0587229\\
4.0411	-0.0586472\\
4.0437	-0.0585695\\
4.0463	-0.0584929\\
4.0488	-0.0584203\\
4.0514	-0.058346\\
4.054	-0.0582727\\
4.0565	-0.0582035\\
4.0591	-0.0581325\\
4.0617	-0.0580628\\
4.0642	-0.057997\\
4.0668	-0.0579297\\
4.0693	-0.0578662\\
4.0719	-0.0578014\\
4.0745	-0.0577378\\
4.077	-0.0576779\\
4.0796	-0.0576169\\
4.0822	-0.0575572\\
4.0847	-0.0575011\\
4.0873	-0.057444\\
4.0899	-0.0573883\\
4.0924	-0.057336\\
4.095	-0.0572829\\
4.0975	-0.0572332\\
4.1001	-0.0571828\\
4.1027	-0.0571339\\
4.1052	-0.0570881\\
4.1078	-0.0570418\\
4.1104	-0.056997\\
4.1129	-0.0569552\\
4.1155	-0.0569131\\
4.1181	-0.0568724\\
4.1206	-0.0568346\\
4.1232	-0.0567967\\
4.1258	-0.0567602\\
4.1283	-0.0567264\\
4.1309	-0.0566927\\
4.1334	-0.0566616\\
4.136	-0.0566307\\
4.1386	-0.0566011\\
4.1411	-0.056574\\
4.1437	-0.0565473\\
4.1463	-0.0565219\\
4.1488	-0.0564988\\
4.1514	-0.0564761\\
4.154	-0.0564548\\
4.1565	-0.0564357\\
4.1591	-0.0564171\\
4.1616	-0.0564005\\
4.1642	-0.0563845\\
4.1668	-0.0563699\\
4.1693	-0.0563572\\
4.1719	-0.0563452\\
4.1745	-0.0563345\\
4.177	-0.0563255\\
4.1796	-0.0563173\\
4.1822	-0.0563105\\
4.1847	-0.0563051\\
4.1873	-0.0563007\\
4.1898	-0.0562976\\
4.1924	-0.0562956\\
4.195	-0.0562948\\
4.1975	-0.0562952\\
4.2001	-0.0562967\\
4.2027	-0.0562994\\
4.2052	-0.0563031\\
4.2078	-0.056308\\
4.2104	-0.056314\\
4.2129	-0.0563209\\
4.2155	-0.056329\\
4.218	-0.0563378\\
4.2206	-0.056348\\
4.2232	-0.0563592\\
4.2257	-0.0563708\\
4.2283	-0.0563839\\
4.2309	-0.056398\\
4.2334	-0.0564124\\
4.236	-0.0564282\\
4.2386	-0.0564449\\
4.2411	-0.0564618\\
4.2437	-0.0564802\\
4.2462	-0.0564986\\
4.2488	-0.0565186\\
4.2514	-0.0565393\\
4.2539	-0.0565599\\
4.2565	-0.056582\\
4.2591	-0.0566049\\
4.2616	-0.0566274\\
4.2642	-0.0566515\\
4.2668	-0.0566762\\
4.2693	-0.0567006\\
4.2719	-0.0567264\\
4.2744	-0.0567518\\
4.277	-0.0567786\\
4.2796	-0.056806\\
4.2821	-0.0568327\\
4.2847	-0.056861\\
4.2873	-0.0568896\\
4.2898	-0.0569175\\
4.2924	-0.0569468\\
4.295	-0.0569765\\
4.2975	-0.0570053\\
4.3001	-0.0570355\\
4.3027	-0.057066\\
4.3052	-0.0570955\\
4.3078	-0.0571263\\
4.3103	-0.0571562\\
4.3129	-0.0571873\\
4.3155	-0.0572185\\
4.318	-0.0572487\\
4.3206	-0.05728\\
4.3232	-0.0573114\\
4.3257	-0.0573416\\
4.3283	-0.057373\\
4.3309	-0.0574043\\
4.3334	-0.0574343\\
4.336	-0.0574654\\
4.3385	-0.0574952\\
4.3411	-0.0575261\\
4.3437	-0.0575567\\
4.3462	-0.0575859\\
4.3488	-0.0576161\\
4.3514	-0.0576461\\
4.3539	-0.0576745\\
4.3565	-0.0577039\\
4.3591	-0.0577329\\
4.3616	-0.0577604\\
4.3642	-0.0577886\\
4.3667	-0.0578154\\
4.3693	-0.0578428\\
4.3719	-0.0578697\\
4.3744	-0.0578952\\
4.377	-0.0579212\\
4.3796	-0.0579466\\
4.3821	-0.0579705\\
4.3847	-0.0579949\\
4.3873	-0.0580186\\
4.3898	-0.0580408\\
4.3924	-0.0580633\\
4.3949	-0.0580844\\
4.3975	-0.0581056\\
4.4001	-0.058126\\
4.4026	-0.0581451\\
4.4052	-0.0581641\\
4.4078	-0.0581824\\
4.4103	-0.0581993\\
4.4129	-0.0582161\\
4.4155	-0.058232\\
4.418	-0.0582466\\
4.4206	-0.0582609\\
4.4231	-0.0582739\\
4.4257	-0.0582865\\
4.4283	-0.0582982\\
4.4308	-0.0583087\\
4.4334	-0.0583186\\
4.436	-0.0583276\\
4.4385	-0.0583354\\
4.4411	-0.0583425\\
4.4437	-0.0583487\\
4.4462	-0.0583537\\
4.4488	-0.058358\\
4.4514	-0.0583612\\
4.4539	-0.0583634\\
4.4565	-0.0583646\\
4.459	-0.0583649\\
4.4616	-0.0583641\\
4.4642	-0.0583622\\
4.4667	-0.0583595\\
4.4693	-0.0583556\\
4.4719	-0.0583506\\
4.4744	-0.0583448\\
4.477	-0.0583377\\
4.4796	-0.0583295\\
4.4821	-0.0583206\\
4.4847	-0.0583102\\
4.4872	-0.0582993\\
4.4898	-0.0582868\\
4.4924	-0.0582732\\
4.4949	-0.0582591\\
4.4975	-0.0582433\\
4.5001	-0.0582265\\
4.5026	-0.0582092\\
4.5052	-0.0581902\\
4.5078	-0.05817\\
4.5103	-0.0581496\\
4.5129	-0.0581273\\
4.5154	-0.0581048\\
4.518	-0.0580803\\
4.5206	-0.0580548\\
4.5231	-0.0580292\\
4.5257	-0.0580015\\
4.5283	-0.0579727\\
4.5308	-0.057944\\
4.5334	-0.0579131\\
4.536	-0.0578811\\
4.5385	-0.0578494\\
4.5411	-0.0578154\\
4.5436	-0.0577817\\
4.5462	-0.0577456\\
4.5488	-0.0577085\\
4.5513	-0.0576719\\
4.5539	-0.0576328\\
4.5565	-0.0575927\\
4.559	-0.0575533\\
4.5616	-0.0575113\\
4.5642	-0.0574683\\
4.5667	-0.0574261\\
4.5693	-0.0573812\\
4.5718	-0.0573372\\
4.5744	-0.0572906\\
4.577	-0.057243\\
4.5795	-0.0571965\\
4.5821	-0.0571472\\
4.5847	-0.057097\\
4.5872	-0.057048\\
4.5898	-0.0569962\\
4.5924	-0.0569436\\
4.5949	-0.0568922\\
4.5975	-0.056838\\
4.6001	-0.0567831\\
4.6026	-0.0567296\\
4.6052	-0.0566732\\
4.6077	-0.0566182\\
4.6103	-0.0565605\\
4.6129	-0.056502\\
4.6154	-0.0564451\\
4.618	-0.0563854\\
4.6206	-0.056325\\
4.6231	-0.0562664\\
4.6257	-0.0562048\\
4.6283	-0.0561427\\
4.6308	-0.0560824\\
4.6334	-0.0560192\\
4.6359	-0.055958\\
4.6385	-0.0558939\\
4.6411	-0.0558292\\
4.6436	-0.0557667\\
4.6462	-0.0557012\\
4.6488	-0.0556354\\
4.6513	-0.0555717\\
4.6539	-0.0555051\\
4.6565	-0.0554381\\
4.659	-0.0553735\\
4.6616	-0.0553059\\
4.6641	-0.0552407\\
4.6667	-0.0551727\\
4.6693	-0.0551044\\
4.6718	-0.0550386\\
4.6744	-0.0549699\\
4.677	-0.0549011\\
4.6795	-0.0548349\\
4.6821	-0.0547658\\
4.6847	-0.0546967\\
4.6872	-0.0546301\\
4.6898	-0.0545609\\
4.6923	-0.0544943\\
4.6949	-0.0544251\\
4.6975	-0.0543558\\
4.7	-0.0542893\\
4.7026	-0.0542202\\
4.7052	-0.0541512\\
4.7077	-0.054085\\
4.7103	-0.0540162\\
4.7129	-0.0539476\\
4.7154	-0.0538819\\
4.718	-0.0538137\\
4.7205	-0.0537483\\
4.7231	-0.0536806\\
4.7257	-0.0536132\\
4.7282	-0.0535487\\
4.7308	-0.0534818\\
4.7334	-0.0534154\\
4.7359	-0.0533518\\
4.7385	-0.0532861\\
4.7411	-0.0532209\\
4.7436	-0.0531585\\
4.7462	-0.0530941\\
4.7487	-0.0530326\\
4.7513	-0.0529692\\
4.7539	-0.0529063\\
4.7564	-0.0528463\\
4.759	-0.0527845\\
4.7616	-0.0527233\\
4.7641	-0.052665\\
4.7667	-0.052605\\
4.7693	-0.0525456\\
4.7718	-0.0524892\\
4.7744	-0.0524311\\
4.777	-0.0523738\\
4.7795	-0.0523194\\
4.7821	-0.0522635\\
4.7846	-0.0522106\\
4.7872	-0.0521562\\
4.7898	-0.0521027\\
4.7923	-0.052052\\
4.7949	-0.0520001\\
4.7975	-0.0519491\\
4.8	-0.0519008\\
4.8026	-0.0518515\\
4.8052	-0.0518031\\
4.8077	-0.0517575\\
4.8103	-0.0517109\\
4.8128	-0.051667\\
4.8154	-0.0516224\\
4.818	-0.0515787\\
4.8205	-0.0515376\\
4.8231	-0.0514959\\
4.8257	-0.0514553\\
4.8282	-0.0514171\\
4.8308	-0.0513785\\
4.8334	-0.051341\\
4.8359	-0.0513059\\
4.8385	-0.0512705\\
4.841	-0.0512375\\
4.8436	-0.0512043\\
4.8462	-0.0511722\\
4.8487	-0.0511424\\
4.8513	-0.0511126\\
4.8539	-0.0510839\\
4.8564	-0.0510574\\
4.859	-0.051031\\
4.8616	-0.0510058\\
4.8641	-0.0509827\\
4.8667	-0.0509598\\
4.8692	-0.050939\\
4.8718	-0.0509185\\
4.8744	-0.0508993\\
4.8769	-0.0508819\\
4.8795	-0.0508651\\
4.8821	-0.0508495\\
4.8846	-0.0508356\\
4.8872	-0.0508225\\
4.8898	-0.0508106\\
4.8923	-0.0508003\\
4.8949	-0.0507908\\
4.8974	-0.0507829\\
4.9	-0.0507759\\
4.9026	-0.0507702\\
4.9051	-0.0507659\\
4.9077	-0.0507627\\
4.9103	-0.0507607\\
4.9128	-0.05076\\
4.9154	-0.0507605\\
4.918	-0.0507623\\
4.9205	-0.0507651\\
4.9231	-0.0507694\\
4.9257	-0.0507749\\
4.9282	-0.0507813\\
4.9308	-0.0507893\\
4.9333	-0.0507981\\
4.9359	-0.0508085\\
4.9385	-0.0508201\\
4.941	-0.0508324\\
4.9436	-0.0508465\\
4.9462	-0.0508617\\
4.9487	-0.0508775\\
4.9513	-0.0508952\\
4.9539	-0.050914\\
4.9564	-0.0509332\\
4.959	-0.0509544\\
4.9615	-0.0509758\\
4.9641	-0.0509993\\
4.9667	-0.0510239\\
4.9692	-0.0510486\\
4.9718	-0.0510755\\
4.9744	-0.0511035\\
4.9769	-0.0511314\\
4.9795	-0.0511616\\
4.9821	-0.0511928\\
4.9846	-0.0512239\\
4.9872	-0.0512573\\
4.9897	-0.0512904\\
4.9923	-0.0513258\\
4.9949	-0.0513623\\
4.9974	-0.0513983\\
5	-0.0514367\\
};
\addplot [color=mycolor7, line width=1.0pt, forget plot]
  table[row sep=crcr]{%
-0.125	7.04514e-08\\
-0.1224	7.19181e-08\\
-0.1199	7.33283e-08\\
-0.1173	7.47949e-08\\
-0.1147	7.62614e-08\\
-0.1122	7.76714e-08\\
-0.1096	7.91378e-08\\
-0.1071	8.05478e-08\\
-0.1045	8.20142e-08\\
-0.1019	8.34804e-08\\
-0.0994	8.48901e-08\\
-0.0968	8.63562e-08\\
-0.0942	8.78221e-08\\
-0.0917	8.92316e-08\\
-0.0891	9.06975e-08\\
-0.0865	9.21632e-08\\
-0.084	9.35725e-08\\
-0.0814	9.5038e-08\\
-0.0789	9.64471e-08\\
-0.0763	9.79124e-08\\
-0.0737	9.93775e-08\\
-0.0712	1.00786e-07\\
-0.0686	1.02251e-07\\
-0.066	1.03716e-07\\
-0.0635	1.05125e-07\\
-0.0609	1.06589e-07\\
-0.0583	1.08054e-07\\
-0.0558	1.09462e-07\\
-0.0532	1.10926e-07\\
-0.0507	1.12334e-07\\
-0.0481	1.13798e-07\\
-0.0455	1.15262e-07\\
-0.043	1.16669e-07\\
-0.0404	1.18133e-07\\
-0.0378	1.19596e-07\\
-0.0353	1.21003e-07\\
-0.0327	1.22466e-07\\
-0.0301	1.23929e-07\\
-0.0276	1.25336e-07\\
-0.025	1.26799e-07\\
-0.0224	1.28261e-07\\
-0.0199	1.29667e-07\\
-0.0173	1.3113e-07\\
-0.0148	1.32535e-07\\
-0.0122	1.33997e-07\\
-0.0096	1.35459e-07\\
-0.0071	1.36865e-07\\
-0.0045	1.38326e-07\\
-0.0019	1.39788e-07\\
0.0006	1.41193e-07\\
0.0032	5.6199e-05\\
0.0058	0.000601072\\
0.0083	0.00116888\\
0.0109	0.00161069\\
0.0134	0.00204193\\
0.016	0.00257726\\
0.0186	0.00316145\\
0.0211	0.00369169\\
0.0237	0.00416696\\
0.0263	0.00458208\\
0.0288	0.00497048\\
0.0314	0.00539889\\
0.034	0.00585464\\
0.0365	0.00629781\\
0.0391	0.00674132\\
0.0416	0.00714488\\
0.0442	0.00755085\\
0.0468	0.00795836\\
0.0493	0.00835944\\
0.0519	0.00878279\\
0.0545	0.00920252\\
0.057	0.00959554\\
0.0596	0.00999361\\
0.0622	0.0103876\\
0.0647	0.0107695\\
0.0673	0.0111715\\
0.0698	0.0115589\\
0.0724	0.0119565\\
0.075	0.0123456\\
0.0775	0.0127141\\
0.0801	0.013097\\
0.0827	0.0134835\\
0.0852	0.0138577\\
0.0878	0.0142448\\
0.0904	0.0146255\\
0.0929	0.0149852\\
0.0955	0.0153563\\
0.098	0.0157148\\
0.1006	0.0160906\\
0.1032	0.0164668\\
0.1057	0.0168245\\
0.1083	0.0171899\\
0.1109	0.0175505\\
0.1134	0.0178965\\
0.116	0.0182588\\
0.1186	0.0186228\\
0.1211	0.0189707\\
0.1237	0.0193271\\
0.1263	0.0196776\\
0.1288	0.020012\\
0.1314	0.0203606\\
0.1339	0.0206978\\
0.1365	0.0210483\\
0.1391	0.0213949\\
0.1416	0.0217225\\
0.1442	0.0220588\\
0.1468	0.0223942\\
0.1493	0.0227181\\
0.1519	0.0230559\\
0.1545	0.0233915\\
0.157	0.0237094\\
0.1596	0.0240347\\
0.1621	0.0243448\\
0.1647	0.0246676\\
0.1673	0.0249916\\
0.1698	0.0253025\\
0.1724	0.0256221\\
0.175	0.0259362\\
0.1775	0.0262342\\
0.1801	0.0265428\\
0.1827	0.026852\\
0.1852	0.0271495\\
0.1878	0.0274563\\
0.1903	0.0277468\\
0.1929	0.0280437\\
0.1955	0.0283375\\
0.198	0.0286195\\
0.2006	0.028913\\
0.2032	0.029205\\
0.2057	0.0294819\\
0.2083	0.0297646\\
0.2109	0.0300427\\
0.2134	0.0303081\\
0.216	0.0305836\\
0.2185	0.0308476\\
0.2211	0.0311191\\
0.2237	0.0313853\\
0.2262	0.0316361\\
0.2288	0.0318932\\
0.2314	0.0321486\\
0.2339	0.0323931\\
0.2365	0.0326449\\
0.2391	0.032892\\
0.2416	0.0331239\\
0.2442	0.03336\\
0.2467	0.0335839\\
0.2493	0.033815\\
0.2519	0.0340438\\
0.2544	0.0342598\\
0.257	0.0344785\\
0.2596	0.0346912\\
0.2621	0.0348911\\
0.2647	0.0350961\\
0.2673	0.0352984\\
0.2698	0.0354894\\
0.2724	0.0356825\\
0.2749	0.0358619\\
0.2775	0.0360425\\
0.2801	0.0362187\\
0.2826	0.036385\\
0.2852	0.0365544\\
0.2878	0.0367186\\
0.2903	0.0368701\\
0.2929	0.037021\\
0.2955	0.0371662\\
0.298	0.0373018\\
0.3006	0.037439\\
0.3032	0.0375715\\
0.3057	0.0376928\\
0.3083	0.0378119\\
0.3108	0.0379202\\
0.3134	0.0380276\\
0.316	0.0381308\\
0.3185	0.0382257\\
0.3211	0.0383187\\
0.3237	0.0384046\\
0.3262	0.0384805\\
0.3288	0.0385535\\
0.3314	0.0386216\\
0.3339	0.0386829\\
0.3365	0.0387415\\
0.339	0.0387918\\
0.3416	0.038837\\
0.3442	0.0388757\\
0.3467	0.0389078\\
0.3493	0.0389367\\
0.3519	0.0389609\\
0.3544	0.0389788\\
0.357	0.0389909\\
0.3596	0.0389963\\
0.3621	0.0389961\\
0.3647	0.0389914\\
0.3672	0.0389828\\
0.3698	0.0389691\\
0.3724	0.0389495\\
0.3749	0.0389247\\
0.3775	0.0388933\\
0.3801	0.0388571\\
0.3826	0.0388186\\
0.3852	0.0387746\\
0.3878	0.0387257\\
0.3903	0.0386735\\
0.3929	0.0386137\\
0.3954	0.0385519\\
0.398	0.0384839\\
0.4006	0.0384126\\
0.4031	0.0383404\\
0.4057	0.0382609\\
0.4083	0.0381765\\
0.4108	0.0380913\\
0.4134	0.0379992\\
0.416	0.0379043\\
0.4185	0.0378105\\
0.4211	0.0377095\\
0.4236	0.0376087\\
0.4262	0.0375\\
0.4288	0.0373882\\
0.4313	0.0372785\\
0.4339	0.0371626\\
0.4365	0.0370443\\
0.439	0.0369278\\
0.4416	0.0368037\\
0.4442	0.0366768\\
0.4467	0.0365531\\
0.4493	0.0364232\\
0.4519	0.036292\\
0.4544	0.0361642\\
0.457	0.0360291\\
0.4595	0.0358972\\
0.4621	0.0357587\\
0.4647	0.0356195\\
0.4672	0.0354852\\
0.4698	0.0353449\\
0.4724	0.0352034\\
0.4749	0.0350662\\
0.4775	0.0349225\\
0.4801	0.0347785\\
0.4826	0.0346403\\
0.4852	0.0344969\\
0.4877	0.0343588\\
0.4903	0.0342148\\
0.4929	0.0340703\\
0.4954	0.0339316\\
0.498	0.0337881\\
0.5006	0.0336455\\
0.5031	0.0335092\\
0.5057	0.0333679\\
0.5083	0.0332268\\
0.5108	0.0330917\\
0.5134	0.0329523\\
0.5159	0.0328197\\
0.5185	0.0326834\\
0.5211	0.0325483\\
0.5236	0.0324193\\
0.5262	0.0322862\\
0.5288	0.0321545\\
0.5313	0.0320297\\
0.5339	0.0319021\\
0.5365	0.0317764\\
0.539	0.0316571\\
0.5416	0.0315344\\
0.5441	0.0314181\\
0.5467	0.0312994\\
0.5493	0.0311831\\
0.5518	0.0310738\\
0.5544	0.0309622\\
0.557	0.0308526\\
0.5595	0.0307492\\
0.5621	0.0306439\\
0.5647	0.0305414\\
0.5672	0.0304456\\
0.5698	0.0303487\\
0.5723	0.0302577\\
0.5749	0.0301654\\
0.5775	0.0300755\\
0.58	0.0299919\\
0.5826	0.0299079\\
0.5852	0.029827\\
0.5877	0.0297518\\
0.5903	0.0296762\\
0.5929	0.0296031\\
0.5954	0.0295356\\
0.598	0.0294684\\
0.6006	0.0294045\\
0.6031	0.0293459\\
0.6057	0.0292877\\
0.6082	0.0292343\\
0.6108	0.0291815\\
0.6134	0.0291318\\
0.6159	0.029087\\
0.6185	0.0290436\\
0.6211	0.0290031\\
0.6236	0.0289668\\
0.6262	0.0289316\\
0.6288	0.0288994\\
0.6313	0.0288713\\
0.6339	0.0288452\\
0.6364	0.028823\\
0.639	0.0288025\\
0.6416	0.0287846\\
0.6441	0.02877\\
0.6467	0.0287575\\
0.6493	0.0287481\\
0.6518	0.0287418\\
0.6544	0.0287378\\
0.657	0.0287363\\
0.6595	0.0287371\\
0.6621	0.0287405\\
0.6646	0.0287463\\
0.6672	0.0287549\\
0.6698	0.0287661\\
0.6723	0.028779\\
0.6749	0.0287944\\
0.6775	0.0288121\\
0.68	0.0288312\\
0.6826	0.0288534\\
0.6852	0.0288778\\
0.6877	0.0289033\\
0.6903	0.0289315\\
0.6928	0.0289602\\
0.6954	0.0289919\\
0.698	0.0290255\\
0.7005	0.0290596\\
0.7031	0.0290967\\
0.7057	0.0291352\\
0.7082	0.0291735\\
0.7108	0.0292146\\
0.7134	0.029257\\
0.7159	0.0292992\\
0.7185	0.0293442\\
0.721	0.0293885\\
0.7236	0.0294354\\
0.7262	0.029483\\
0.7287	0.0295296\\
0.7313	0.0295787\\
0.7339	0.0296287\\
0.7364	0.0296772\\
0.739	0.029728\\
0.7416	0.0297789\\
0.7441	0.029828\\
0.7467	0.0298792\\
0.7492	0.0299286\\
0.7518	0.02998\\
0.7544	0.0300312\\
0.7569	0.03008\\
0.7595	0.0301302\\
0.7621	0.0301799\\
0.7646	0.0302272\\
0.7672	0.0302758\\
0.7698	0.0303235\\
0.7723	0.0303685\\
0.7749	0.030414\\
0.7775	0.0304584\\
0.78	0.0304998\\
0.7826	0.0305416\\
0.7851	0.0305805\\
0.7877	0.0306193\\
0.7903	0.0306562\\
0.7928	0.0306898\\
0.7954	0.0307229\\
0.798	0.0307539\\
0.8005	0.0307817\\
0.8031	0.0308083\\
0.8057	0.0308325\\
0.8082	0.0308532\\
0.8108	0.030872\\
0.8133	0.0308875\\
0.8159	0.0309008\\
0.8185	0.0309111\\
0.821	0.0309181\\
0.8236	0.030922\\
0.8262	0.0309224\\
0.8287	0.0309196\\
0.8313	0.0309131\\
0.8339	0.0309029\\
0.8364	0.0308896\\
0.839	0.0308717\\
0.8415	0.0308506\\
0.8441	0.0308246\\
0.8467	0.0307942\\
0.8492	0.0307609\\
0.8518	0.0307219\\
0.8544	0.0306783\\
0.8569	0.0306318\\
0.8595	0.0305786\\
0.8621	0.0305204\\
0.8646	0.0304598\\
0.8672	0.0303917\\
0.8697	0.0303214\\
0.8723	0.0302429\\
0.8749	0.0301588\\
0.8774	0.0300728\\
0.88	0.0299777\\
0.8826	0.029877\\
0.8851	0.0297747\\
0.8877	0.0296625\\
0.8903	0.0295442\\
0.8928	0.0294246\\
0.8954	0.0292941\\
0.8979	0.0291628\\
0.9005	0.0290201\\
0.9031	0.0288709\\
0.9056	0.0287213\\
0.9082	0.0285592\\
0.9108	0.0283904\\
0.9133	0.0282218\\
0.9159	0.0280397\\
0.9185	0.0278509\\
0.921	0.0276627\\
0.9236	0.0274602\\
0.9262	0.0272505\\
0.9287	0.0270421\\
0.9313	0.0268185\\
0.9338	0.0265966\\
0.9364	0.0263588\\
0.939	0.0261136\\
0.9415	0.0258709\\
0.9441	0.0256111\\
0.9467	0.025344\\
0.9492	0.02508\\
0.9518	0.0247982\\
0.9544	0.0245088\\
0.9569	0.0242234\\
0.9595	0.023919\\
0.962	0.0236192\\
0.9646	0.0232999\\
0.9672	0.022973\\
0.9697	0.0226514\\
0.9723	0.0223093\\
0.9749	0.0219596\\
0.9774	0.021616\\
0.98	0.0212511\\
0.9826	0.0208785\\
0.9851	0.020513\\
0.9877	0.0201254\\
0.9902	0.0197453\\
0.9928	0.0193425\\
0.9954	0.0189321\\
0.9979	0.0185302\\
1.0005	0.0181048\\
1.0031	0.0176719\\
1.0056	0.0172484\\
1.0082	0.0168005\\
1.0108	0.0163452\\
1.0133	0.0159003\\
1.0159	0.0154303\\
1.0184	0.0149715\\
1.021	0.0144871\\
1.0236	0.0139954\\
1.0261	0.0135158\\
1.0287	0.0130099\\
1.0313	0.0124969\\
1.0338	0.011997\\
1.0364	0.0114702\\
1.039	0.0109365\\
1.0415	0.0104169\\
1.0441	0.00986971\\
1.0466	0.00933728\\
1.0492	0.00877707\\
1.0518	0.00821033\\
1.0543	0.00765928\\
1.0569	0.0070799\\
1.0595	0.0064942\\
1.062	0.00592519\\
1.0646	0.00532744\\
1.0672	0.00472372\\
1.0697	0.00413767\\
1.0723	0.00352247\\
1.0748	0.00292554\\
1.0774	0.00229924\\
1.08	0.00166749\\
1.0825	0.00105501\\
1.0851	0.000412907\\
1.0877	-0.000234313\\
1.0902	-0.000861368\\
1.0928	-0.0015183\\
1.0954	-0.00217998\\
1.0979	-0.00282054\\
1.1005	-0.0034911\\
1.1031	-0.004166\\
1.1056	-0.00481892\\
1.1082	-0.00550198\\
1.1107	-0.00616249\\
1.1133	-0.00685313\\
1.1159	-0.00754741\\
1.1184	-0.00821827\\
1.121	-0.00891927\\
1.1236	-0.00962349\\
1.1261	-0.0103035\\
1.1287	-0.0110136\\
1.1313	-0.0117264\\
1.1338	-0.0124143\\
1.1364	-0.013132\\
1.1389	-0.0138244\\
1.1415	-0.0145465\\
1.1441	-0.0152706\\
1.1466	-0.0159685\\
1.1492	-0.016696\\
1.1518	-0.0174249\\
1.1543	-0.0181271\\
1.1569	-0.0188585\\
1.1595	-0.0195909\\
1.162	-0.0202959\\
1.1646	-0.0210298\\
1.1671	-0.0217359\\
1.1697	-0.0224705\\
1.1723	-0.0232054\\
1.1748	-0.023912\\
1.1774	-0.0246467\\
1.18	-0.0253811\\
1.1825	-0.0260868\\
1.1851	-0.02682\\
1.1877	-0.0275525\\
1.1902	-0.0282558\\
1.1928	-0.0289862\\
1.1953	-0.0296872\\
1.1979	-0.0304148\\
1.2005	-0.0311408\\
1.203	-0.0318372\\
1.2056	-0.0325594\\
1.2082	-0.0332796\\
1.2107	-0.03397\\
1.2133	-0.0346856\\
1.2159	-0.0353985\\
1.2184	-0.0360815\\
1.221	-0.036789\\
1.2235	-0.0374663\\
1.2261	-0.0381677\\
1.2287	-0.0388657\\
1.2312	-0.0395335\\
1.2338	-0.0402245\\
1.2364	-0.0409117\\
1.2389	-0.0415688\\
1.2415	-0.0422482\\
1.2441	-0.0429234\\
1.2466	-0.0435685\\
1.2492	-0.0442351\\
1.2518	-0.044897\\
1.2543	-0.0455291\\
1.2569	-0.0461816\\
1.2594	-0.0468044\\
1.262	-0.047447\\
1.2646	-0.0480845\\
1.2671	-0.0486923\\
1.2697	-0.0493192\\
1.2723	-0.0499404\\
1.2748	-0.0505325\\
1.2774	-0.0511425\\
1.28	-0.0517466\\
1.2825	-0.0523218\\
1.2851	-0.0529141\\
1.2876	-0.0534777\\
1.2902	-0.0540577\\
1.2928	-0.0546314\\
1.2953	-0.0551768\\
1.2979	-0.0557377\\
1.3005	-0.0562919\\
1.303	-0.0568185\\
1.3056	-0.0573594\\
1.3082	-0.0578935\\
1.3107	-0.0584004\\
1.3133	-0.0589208\\
1.3158	-0.0594145\\
1.3184	-0.0599209\\
1.321	-0.06042\\
1.3235	-0.0608931\\
1.3261	-0.061378\\
1.3287	-0.0618555\\
1.3312	-0.0623076\\
1.3338	-0.0627704\\
1.3364	-0.0632258\\
1.3389	-0.0636564\\
1.3415	-0.0640969\\
1.344	-0.0645132\\
1.3466	-0.0649386\\
1.3492	-0.0653563\\
1.3517	-0.0657506\\
1.3543	-0.0661531\\
1.3569	-0.0665478\\
1.3594	-0.06692\\
1.362	-0.0672994\\
1.3646	-0.067671\\
1.3671	-0.0680208\\
1.3697	-0.068377\\
1.3722	-0.068712\\
1.3748	-0.0690527\\
1.3774	-0.0693855\\
1.3799	-0.0696981\\
1.3825	-0.0700155\\
1.3851	-0.070325\\
1.3876	-0.0706152\\
1.3902	-0.0709093\\
1.3928	-0.0711956\\
1.3953	-0.0714635\\
1.3979	-0.0717345\\
1.4005	-0.0719976\\
1.403	-0.0722434\\
1.4056	-0.0724913\\
1.4081	-0.0727225\\
1.4107	-0.0729554\\
1.4133	-0.0731806\\
1.4158	-0.07339\\
1.4184	-0.0736003\\
1.421	-0.0738031\\
1.4235	-0.073991\\
1.4261	-0.0741791\\
1.4287	-0.0743598\\
1.4312	-0.0745266\\
1.4338	-0.0746928\\
1.4363	-0.0748458\\
1.4389	-0.0749979\\
1.4415	-0.0751427\\
1.444	-0.0752753\\
1.4466	-0.0754063\\
1.4492	-0.0755302\\
1.4517	-0.0756429\\
1.4543	-0.0757533\\
1.4569	-0.0758569\\
1.4594	-0.0759501\\
1.462	-0.0760406\\
1.4645	-0.0761213\\
1.4671	-0.0761988\\
1.4697	-0.0762699\\
1.4722	-0.0763322\\
1.4748	-0.0763908\\
1.4774	-0.0764431\\
1.4799	-0.0764877\\
1.4825	-0.076528\\
1.4851	-0.0765623\\
1.4876	-0.0765898\\
1.4902	-0.0766126\\
1.4927	-0.0766291\\
1.4953	-0.0766407\\
1.4979	-0.0766467\\
1.5004	-0.0766473\\
1.503	-0.0766427\\
1.5056	-0.0766327\\
1.5081	-0.0766183\\
1.5107	-0.0765982\\
1.5133	-0.0765731\\
1.5158	-0.0765442\\
1.5184	-0.0765095\\
1.5209	-0.0764717\\
1.5235	-0.0764278\\
1.5261	-0.0763794\\
1.5286	-0.0763286\\
1.5312	-0.0762715\\
1.5338	-0.0762102\\
1.5363	-0.0761474\\
1.5389	-0.076078\\
1.5415	-0.0760048\\
1.544	-0.0759307\\
1.5466	-0.0758501\\
1.5491	-0.0757691\\
1.5517	-0.0756814\\
1.5543	-0.0755904\\
1.5568	-0.0754997\\
1.5594	-0.0754022\\
1.562	-0.0753017\\
1.5645	-0.0752022\\
1.5671	-0.0750959\\
1.5697	-0.0749868\\
1.5722	-0.0748795\\
1.5748	-0.0747653\\
1.5774	-0.0746487\\
1.5799	-0.0745344\\
1.5825	-0.0744133\\
1.585	-0.0742949\\
1.5876	-0.0741698\\
1.5902	-0.0740427\\
1.5927	-0.0739189\\
1.5953	-0.0737884\\
1.5979	-0.0736564\\
1.6004	-0.0735281\\
1.603	-0.0733933\\
1.6056	-0.0732572\\
1.6081	-0.0731253\\
1.6107	-0.0729871\\
1.6132	-0.0728533\\
1.6158	-0.0727134\\
1.6184	-0.0725727\\
1.6209	-0.0724369\\
1.6235	-0.0722951\\
1.6261	-0.0721529\\
1.6286	-0.0720158\\
1.6312	-0.0718731\\
1.6338	-0.0717302\\
1.6363	-0.0715929\\
1.6389	-0.0714501\\
1.6414	-0.071313\\
1.644	-0.0711707\\
1.6466	-0.0710288\\
1.6491	-0.0708928\\
1.6517	-0.0707519\\
1.6543	-0.0706117\\
1.6568	-0.0704776\\
1.6594	-0.070339\\
1.662	-0.0702013\\
1.6645	-0.0700699\\
1.6671	-0.0699343\\
1.6696	-0.0698051\\
1.6722	-0.069672\\
1.6748	-0.0695402\\
1.6773	-0.069415\\
1.6799	-0.0692862\\
1.6825	-0.0691591\\
1.685	-0.0690385\\
1.6876	-0.0689148\\
1.6902	-0.068793\\
1.6927	-0.0686777\\
1.6953	-0.0685597\\
1.6978	-0.0684483\\
1.7004	-0.0683346\\
1.703	-0.068223\\
1.7055	-0.068118\\
1.7081	-0.0680111\\
1.7107	-0.0679066\\
1.7132	-0.0678085\\
1.7158	-0.067709\\
1.7184	-0.0676121\\
1.7209	-0.0675214\\
1.7235	-0.0674299\\
1.7261	-0.0673411\\
1.7286	-0.0672584\\
1.7312	-0.0671753\\
1.7337	-0.0670982\\
1.7363	-0.0670209\\
1.7389	-0.0669467\\
1.7414	-0.0668782\\
1.744	-0.0668101\\
1.7466	-0.0667452\\
1.7491	-0.0666858\\
1.7517	-0.0666273\\
1.7543	-0.066572\\
1.7568	-0.066522\\
1.7594	-0.0664734\\
1.7619	-0.0664298\\
1.7645	-0.0663878\\
1.7671	-0.0663492\\
1.7696	-0.0663155\\
1.7722	-0.0662838\\
1.7748	-0.0662556\\
1.7773	-0.0662319\\
1.7799	-0.0662107\\
1.7825	-0.066193\\
1.785	-0.0661795\\
1.7876	-0.0661689\\
1.7901	-0.0661622\\
1.7927	-0.0661587\\
1.7953	-0.0661588\\
1.7978	-0.0661624\\
1.8004	-0.0661697\\
1.803	-0.0661806\\
1.8055	-0.0661945\\
1.8081	-0.0662125\\
1.8107	-0.0662342\\
1.8132	-0.0662584\\
1.8158	-0.0662871\\
1.8183	-0.0663182\\
1.8209	-0.0663539\\
1.8235	-0.0663933\\
1.826	-0.0664345\\
1.8286	-0.0664808\\
1.8312	-0.0665306\\
1.8337	-0.0665818\\
1.8363	-0.0666384\\
1.8389	-0.0666985\\
1.8414	-0.0667595\\
1.844	-0.0668262\\
1.8465	-0.0668936\\
1.8491	-0.0669669\\
1.8517	-0.0670434\\
1.8542	-0.0671201\\
1.8568	-0.0672031\\
1.8594	-0.0672892\\
1.8619	-0.0673749\\
1.8645	-0.0674671\\
1.8671	-0.0675624\\
1.8696	-0.0676568\\
1.8722	-0.0677579\\
1.8747	-0.0678579\\
1.8773	-0.0679646\\
1.8799	-0.0680742\\
1.8824	-0.0681821\\
1.885	-0.068297\\
1.8876	-0.0684146\\
1.8901	-0.0685301\\
1.8927	-0.0686526\\
1.8953	-0.0687777\\
1.8978	-0.0689001\\
1.9004	-0.0690298\\
1.903	-0.0691618\\
1.9055	-0.0692908\\
1.9081	-0.0694271\\
1.9106	-0.06956\\
1.9132	-0.0697003\\
1.9158	-0.0698426\\
1.9183	-0.0699811\\
1.9209	-0.070127\\
1.9235	-0.0702746\\
1.926	-0.0704181\\
1.9286	-0.0705689\\
1.9312	-0.0707212\\
1.9337	-0.070869\\
1.9363	-0.0710241\\
1.9388	-0.0711744\\
1.9414	-0.071332\\
1.944	-0.0714907\\
1.9465	-0.0716443\\
1.9491	-0.071805\\
1.9517	-0.0719666\\
1.9542	-0.0721228\\
1.9568	-0.072286\\
1.9594	-0.0724498\\
1.9619	-0.0726079\\
1.9645	-0.0727729\\
1.967	-0.0729319\\
1.9696	-0.0730976\\
1.9722	-0.0732636\\
1.9747	-0.0734234\\
1.9773	-0.0735897\\
1.9799	-0.073756\\
1.9824	-0.0739159\\
1.985	-0.0740821\\
1.9876	-0.0742481\\
1.9901	-0.0744074\\
1.9927	-0.0745728\\
1.9952	-0.0747314\\
1.9978	-0.0748958\\
2.0004	-0.0750596\\
2.0029	-0.0752165\\
2.0055	-0.0753789\\
2.0081	-0.0755405\\
2.0106	-0.075695\\
2.0132	-0.0758548\\
2.0158	-0.0760135\\
2.0183	-0.076165\\
2.0209	-0.0763214\\
2.0234	-0.0764706\\
2.026	-0.0766245\\
2.0286	-0.0767769\\
2.0311	-0.0769221\\
2.0337	-0.0770716\\
2.0363	-0.0772194\\
2.0388	-0.0773599\\
2.0414	-0.0775044\\
2.044	-0.077647\\
2.0465	-0.0777824\\
2.0491	-0.0779213\\
2.0517	-0.0780581\\
2.0542	-0.0781877\\
2.0568	-0.0783205\\
2.0593	-0.078446\\
2.0619	-0.0785744\\
2.0645	-0.0787004\\
2.067	-0.0788194\\
2.0696	-0.0789408\\
2.0722	-0.0790597\\
2.0747	-0.0791717\\
2.0773	-0.0792856\\
2.0799	-0.0793969\\
2.0824	-0.0795014\\
2.085	-0.0796074\\
2.0875	-0.0797067\\
2.0901	-0.0798073\\
2.0927	-0.0799051\\
2.0952	-0.0799963\\
2.0978	-0.0800884\\
2.1004	-0.0801775\\
2.1029	-0.0802603\\
2.1055	-0.0803435\\
2.1081	-0.0804236\\
2.1106	-0.0804978\\
2.1132	-0.0805718\\
2.1157	-0.08064\\
2.1183	-0.0807079\\
2.1209	-0.0807725\\
2.1234	-0.0808316\\
2.126	-0.0808899\\
2.1286	-0.0809449\\
2.1311	-0.0809948\\
2.1337	-0.0810433\\
2.1363	-0.0810886\\
2.1388	-0.0811289\\
2.1414	-0.0811676\\
2.1439	-0.0812017\\
2.1465	-0.0812337\\
2.1491	-0.0812625\\
2.1516	-0.0812869\\
2.1542	-0.0813089\\
2.1568	-0.0813276\\
2.1593	-0.0813423\\
2.1619	-0.0813543\\
2.1645	-0.0813629\\
2.167	-0.081368\\
2.1696	-0.0813699\\
2.1721	-0.0813685\\
2.1747	-0.0813638\\
2.1773	-0.0813557\\
2.1798	-0.0813447\\
2.1824	-0.08133\\
2.185	-0.0813119\\
2.1875	-0.0812914\\
2.1901	-0.0812667\\
2.1927	-0.0812388\\
2.1952	-0.0812088\\
2.1978	-0.0811744\\
2.2004	-0.0811367\\
2.2029	-0.0810974\\
2.2055	-0.0810534\\
2.208	-0.081008\\
2.2106	-0.0809577\\
2.2132	-0.0809043\\
2.2157	-0.0808499\\
2.2183	-0.0807903\\
2.2209	-0.0807276\\
2.2234	-0.0806645\\
2.226	-0.0805958\\
2.2286	-0.0805241\\
2.2311	-0.0804524\\
2.2337	-0.080375\\
2.2362	-0.0802977\\
2.2388	-0.0802146\\
2.2414	-0.0801286\\
2.2439	-0.0800433\\
2.2465	-0.0799518\\
2.2491	-0.0798576\\
2.2516	-0.0797644\\
2.2542	-0.0796649\\
2.2568	-0.0795628\\
2.2593	-0.0794622\\
2.2619	-0.079355\\
2.2644	-0.0792496\\
2.267	-0.0791375\\
2.2696	-0.0790231\\
2.2721	-0.0789107\\
2.2747	-0.0787916\\
2.2773	-0.0786701\\
2.2798	-0.0785512\\
2.2824	-0.0784254\\
2.285	-0.0782974\\
2.2875	-0.0781723\\
2.2901	-0.0780401\\
2.2926	-0.0779111\\
2.2952	-0.077775\\
2.2978	-0.0776369\\
2.3003	-0.0775023\\
2.3029	-0.0773606\\
2.3055	-0.077217\\
2.308	-0.0770773\\
2.3106	-0.0769303\\
2.3132	-0.0767816\\
2.3157	-0.0766372\\
2.3183	-0.0764854\\
2.3208	-0.076338\\
2.3234	-0.0761832\\
2.326	-0.0760271\\
2.3285	-0.0758756\\
2.3311	-0.0757168\\
2.3337	-0.0755567\\
2.3362	-0.0754015\\
2.3388	-0.0752391\\
2.3414	-0.0750755\\
2.3439	-0.0749171\\
2.3465	-0.0747514\\
2.349	-0.0745912\\
2.3516	-0.0744236\\
2.3542	-0.0742552\\
2.3567	-0.0740924\\
2.3593	-0.0739223\\
2.3619	-0.0737515\\
2.3644	-0.0735867\\
2.367	-0.0734146\\
2.3696	-0.0732419\\
2.3721	-0.0730753\\
2.3747	-0.0729017\\
2.3773	-0.0727275\\
2.3798	-0.0725597\\
2.3824	-0.0723848\\
2.3849	-0.0722164\\
2.3875	-0.072041\\
2.3901	-0.0718654\\
2.3926	-0.0716964\\
2.3952	-0.0715206\\
2.3978	-0.0713446\\
2.4003	-0.0711755\\
2.4029	-0.0709996\\
2.4055	-0.0708237\\
2.408	-0.0706548\\
2.4106	-0.0704793\\
2.4131	-0.0703107\\
2.4157	-0.0701357\\
2.4183	-0.0699609\\
2.4208	-0.0697932\\
2.4234	-0.0696192\\
2.426	-0.0694456\\
2.4285	-0.0692792\\
2.4311	-0.0691066\\
2.4337	-0.0689345\\
2.4362	-0.0687697\\
2.4388	-0.0685989\\
2.4413	-0.0684353\\
2.4439	-0.0682659\\
2.4465	-0.0680973\\
2.449	-0.0679359\\
2.4516	-0.0677689\\
2.4542	-0.0676027\\
2.4567	-0.0674439\\
2.4593	-0.0672796\\
2.4619	-0.0671163\\
2.4644	-0.0669603\\
2.467	-0.066799\\
2.4695	-0.066645\\
2.4721	-0.066486\\
2.4747	-0.0663281\\
2.4772	-0.0661774\\
2.4798	-0.0660219\\
2.4824	-0.0658677\\
2.4849	-0.0657206\\
2.4875	-0.0655689\\
2.4901	-0.0654186\\
2.4926	-0.0652753\\
2.4952	-0.0651277\\
2.4977	-0.0649872\\
2.5003	-0.0648424\\
2.5029	-0.0646991\\
2.5054	-0.0645627\\
2.508	-0.0644224\\
2.5106	-0.0642837\\
2.5131	-0.0641517\\
2.5157	-0.064016\\
2.5183	-0.063882\\
2.5208	-0.0637546\\
2.5234	-0.0636238\\
2.526	-0.0634946\\
2.5285	-0.063372\\
2.5311	-0.0632462\\
2.5336	-0.0631267\\
2.5362	-0.0630043\\
2.5388	-0.0628835\\
2.5413	-0.0627691\\
2.5439	-0.0626518\\
2.5465	-0.0625362\\
2.549	-0.0624268\\
2.5516	-0.0623148\\
2.5542	-0.0622045\\
2.5567	-0.0621002\\
2.5593	-0.0619935\\
2.5618	-0.0618926\\
2.5644	-0.0617895\\
2.567	-0.0616882\\
2.5695	-0.0615924\\
2.5721	-0.0614947\\
2.5747	-0.0613987\\
2.5772	-0.0613082\\
2.5798	-0.0612158\\
2.5824	-0.0611252\\
2.5849	-0.0610398\\
2.5875	-0.0609528\\
2.59	-0.0608707\\
2.5926	-0.0607872\\
2.5952	-0.0607055\\
2.5977	-0.0606285\\
2.6003	-0.0605502\\
2.6029	-0.0604737\\
2.6054	-0.0604018\\
2.608	-0.0603286\\
2.6106	-0.0602572\\
2.6131	-0.0601901\\
2.6157	-0.060122\\
2.6182	-0.0600581\\
2.6208	-0.0599932\\
2.6234	-0.05993\\
2.6259	-0.0598708\\
2.6285	-0.0598107\\
2.6311	-0.0597522\\
2.6336	-0.0596974\\
2.6362	-0.0596419\\
2.6388	-0.059588\\
2.6413	-0.0595375\\
2.6439	-0.0594864\\
2.6464	-0.0594386\\
2.649	-0.0593904\\
2.6516	-0.0593435\\
2.6541	-0.0592997\\
2.6567	-0.0592554\\
2.6593	-0.0592124\\
2.6618	-0.0591723\\
2.6644	-0.0591319\\
2.667	-0.0590926\\
2.6695	-0.0590559\\
2.6721	-0.059019\\
2.6746	-0.0589845\\
2.6772	-0.0589496\\
2.6798	-0.0589159\\
2.6823	-0.0588844\\
2.6849	-0.0588526\\
2.6875	-0.0588218\\
2.69	-0.058793\\
2.6926	-0.0587639\\
2.6952	-0.0587357\\
2.6977	-0.0587093\\
2.7003	-0.0586827\\
2.7029	-0.0586568\\
2.7054	-0.0586325\\
2.708	-0.0586079\\
2.7105	-0.0585849\\
2.7131	-0.0585615\\
2.7157	-0.0585387\\
2.7182	-0.0585172\\
2.7208	-0.0584954\\
2.7234	-0.058474\\
2.7259	-0.0584538\\
2.7285	-0.0584331\\
2.7311	-0.0584128\\
2.7336	-0.0583935\\
2.7362	-0.0583737\\
2.7387	-0.0583549\\
2.7413	-0.0583354\\
2.7439	-0.0583161\\
2.7464	-0.0582977\\
2.749	-0.0582785\\
2.7516	-0.0582594\\
2.7541	-0.058241\\
2.7567	-0.0582218\\
2.7593	-0.0582025\\
2.7618	-0.0581838\\
2.7644	-0.0581642\\
2.7669	-0.0581452\\
2.7695	-0.0581252\\
2.7721	-0.0581049\\
2.7746	-0.0580851\\
2.7772	-0.0580641\\
2.7798	-0.0580428\\
2.7823	-0.0580219\\
2.7849	-0.0579997\\
2.7875	-0.0579771\\
2.79	-0.0579548\\
2.7926	-0.0579311\\
2.7951	-0.0579078\\
2.7977	-0.0578829\\
2.8003	-0.0578574\\
2.8028	-0.0578322\\
2.8054	-0.0578053\\
2.808	-0.0577778\\
2.8105	-0.0577505\\
2.8131	-0.0577214\\
2.8157	-0.0576915\\
2.8182	-0.0576619\\
2.8208	-0.0576303\\
2.8233	-0.0575991\\
2.8259	-0.0575657\\
2.8285	-0.0575314\\
2.831	-0.0574975\\
2.8336	-0.0574613\\
2.8362	-0.0574241\\
2.8387	-0.0573874\\
2.8413	-0.0573481\\
2.8439	-0.0573078\\
2.8464	-0.057268\\
2.849	-0.0572256\\
2.8516	-0.0571821\\
2.8541	-0.0571391\\
2.8567	-0.0570934\\
2.8592	-0.0570483\\
2.8618	-0.0570002\\
2.8644	-0.056951\\
2.8669	-0.0569025\\
2.8695	-0.056851\\
2.8721	-0.0567982\\
2.8746	-0.0567463\\
2.8772	-0.0566912\\
2.8798	-0.0566348\\
2.8823	-0.0565794\\
2.8849	-0.0565206\\
2.8874	-0.0564629\\
2.89	-0.0564017\\
2.8926	-0.0563392\\
2.8951	-0.0562779\\
2.8977	-0.056213\\
2.9003	-0.0561469\\
2.9028	-0.0560821\\
2.9054	-0.0560135\\
2.908	-0.0559437\\
2.9105	-0.0558754\\
2.9131	-0.0558032\\
2.9156	-0.0557326\\
2.9182	-0.0556581\\
2.9208	-0.0555823\\
2.9233	-0.0555083\\
2.9259	-0.0554303\\
2.9285	-0.055351\\
2.931	-0.0552738\\
2.9336	-0.0551923\\
2.9362	-0.0551097\\
2.9387	-0.0550292\\
2.9413	-0.0549445\\
2.9438	-0.0548619\\
2.9464	-0.0547751\\
2.949	-0.0546872\\
2.9515	-0.0546018\\
2.9541	-0.0545119\\
2.9567	-0.0544211\\
2.9592	-0.0543328\\
2.9618	-0.0542401\\
2.9644	-0.0541465\\
2.9669	-0.0540557\\
2.9695	-0.0539604\\
2.972	-0.0538679\\
2.9746	-0.053771\\
2.9772	-0.0536733\\
2.9797	-0.0535786\\
2.9823	-0.0534794\\
2.9849	-0.0533795\\
2.9874	-0.0532828\\
2.99	-0.0531817\\
2.9926	-0.0530799\\
2.9951	-0.0529814\\
2.9977	-0.0528785\\
3.0003	-0.0527751\\
3.0028	-0.0526752\\
3.0054	-0.0525708\\
3.0079	-0.05247\\
3.0105	-0.0523649\\
3.0131	-0.0522594\\
3.0156	-0.0521576\\
3.0182	-0.0520515\\
3.0208	-0.0519451\\
3.0233	-0.0518426\\
3.0259	-0.0517359\\
3.0285	-0.0516289\\
3.031	-0.051526\\
3.0336	-0.051419\\
3.0361	-0.051316\\
3.0387	-0.0512088\\
3.0413	-0.0511018\\
3.0438	-0.0509989\\
3.0464	-0.050892\\
3.049	-0.0507852\\
3.0515	-0.0506828\\
3.0541	-0.0505765\\
3.0567	-0.0504704\\
3.0592	-0.0503687\\
3.0618	-0.0502633\\
3.0643	-0.0501623\\
3.0669	-0.0500577\\
3.0695	-0.0499535\\
3.072	-0.0498538\\
3.0746	-0.0497507\\
3.0772	-0.0496481\\
3.0797	-0.0495501\\
3.0823	-0.0494488\\
3.0849	-0.0493482\\
3.0874	-0.0492522\\
3.09	-0.0491531\\
3.0925	-0.0490585\\
3.0951	-0.048961\\
3.0977	-0.0488644\\
3.1002	-0.0487724\\
3.1028	-0.0486777\\
3.1054	-0.048584\\
3.1079	-0.0484948\\
3.1105	-0.0484032\\
3.1131	-0.0483127\\
3.1156	-0.0482267\\
3.1182	-0.0481385\\
3.1207	-0.0480549\\
3.1233	-0.0479691\\
3.1259	-0.0478846\\
3.1284	-0.0478047\\
3.131	-0.0477229\\
3.1336	-0.0476425\\
3.1361	-0.0475666\\
3.1387	-0.0474891\\
3.1413	-0.0474131\\
3.1438	-0.0473415\\
3.1464	-0.0472685\\
3.1489	-0.0472\\
3.1515	-0.0471302\\
3.1541	-0.0470622\\
3.1566	-0.0469984\\
3.1592	-0.0469338\\
3.1618	-0.0468709\\
3.1643	-0.0468121\\
3.1669	-0.0467529\\
3.1695	-0.0466954\\
3.172	-0.046642\\
3.1746	-0.0465883\\
3.1772	-0.0465366\\
3.1797	-0.0464888\\
3.1823	-0.0464409\\
3.1848	-0.0463969\\
3.1874	-0.0463531\\
3.19	-0.0463114\\
3.1925	-0.0462733\\
3.1951	-0.0462357\\
3.1977	-0.0462004\\
3.2002	-0.0461684\\
3.2028	-0.0461373\\
3.2054	-0.0461085\\
3.2079	-0.0460829\\
3.2105	-0.0460584\\
3.213	-0.0460371\\
3.2156	-0.0460172\\
3.2182	-0.0459996\\
3.2207	-0.0459849\\
3.2233	-0.0459719\\
3.2259	-0.0459613\\
3.2284	-0.0459534\\
3.231	-0.0459475\\
3.2336	-0.045944\\
3.2361	-0.045943\\
3.2387	-0.0459443\\
3.2412	-0.0459479\\
3.2438	-0.0459541\\
3.2464	-0.0459627\\
3.2489	-0.0459734\\
3.2515	-0.0459869\\
3.2541	-0.046003\\
3.2566	-0.0460208\\
3.2592	-0.0460418\\
3.2618	-0.0460653\\
3.2643	-0.0460903\\
3.2669	-0.0461188\\
3.2694	-0.0461486\\
3.272	-0.0461821\\
3.2746	-0.0462181\\
3.2771	-0.0462552\\
3.2797	-0.0462963\\
3.2823	-0.0463398\\
3.2848	-0.0463842\\
3.2874	-0.0464328\\
3.29	-0.0464839\\
3.2925	-0.0465354\\
3.2951	-0.0465915\\
3.2976	-0.0466479\\
3.3002	-0.0467089\\
3.3028	-0.0467725\\
3.3053	-0.046836\\
3.3079	-0.0469045\\
3.3105	-0.0469755\\
3.313	-0.047046\\
3.3156	-0.0471219\\
3.3182	-0.0472001\\
3.3207	-0.0472777\\
3.3233	-0.0473607\\
3.3259	-0.0474462\\
3.3284	-0.0475306\\
3.331	-0.0476208\\
3.3335	-0.0477097\\
3.3361	-0.0478044\\
3.3387	-0.0479015\\
3.3412	-0.047997\\
3.3438	-0.0480986\\
3.3464	-0.0482024\\
3.3489	-0.0483044\\
3.3515	-0.0484126\\
3.3541	-0.048523\\
3.3566	-0.0486312\\
3.3592	-0.0487458\\
3.3617	-0.048858\\
3.3643	-0.0489768\\
3.3669	-0.0490976\\
3.3694	-0.0492157\\
3.372	-0.0493404\\
3.3746	-0.0494671\\
3.3771	-0.0495908\\
3.3797	-0.0497213\\
3.3823	-0.0498537\\
3.3848	-0.0499827\\
3.3874	-0.0501187\\
3.3899	-0.0502511\\
3.3925	-0.0503905\\
3.3951	-0.0505316\\
3.3976	-0.0506689\\
3.4002	-0.0508132\\
3.4028	-0.0509591\\
3.4053	-0.0511009\\
3.4079	-0.0512498\\
3.4105	-0.0514002\\
3.413	-0.0515461\\
3.4156	-0.0516993\\
3.4181	-0.0518478\\
3.4207	-0.0520036\\
3.4233	-0.0521606\\
3.4258	-0.0523127\\
3.4284	-0.052472\\
3.431	-0.0526325\\
3.4335	-0.0527879\\
3.4361	-0.0529504\\
3.4387	-0.053114\\
3.4412	-0.0532721\\
3.4438	-0.0534375\\
3.4463	-0.0535973\\
3.4489	-0.0537643\\
3.4515	-0.053932\\
3.454	-0.054094\\
3.4566	-0.054263\\
3.4592	-0.0544327\\
3.4617	-0.0545964\\
3.4643	-0.0547672\\
3.4669	-0.0549384\\
3.4694	-0.0551034\\
3.472	-0.0552753\\
3.4745	-0.0554409\\
3.4771	-0.0556134\\
3.4797	-0.0557862\\
3.4822	-0.0559524\\
3.4848	-0.0561254\\
3.4874	-0.0562984\\
3.4899	-0.0564648\\
3.4925	-0.0566377\\
3.4951	-0.0568106\\
3.4976	-0.0569767\\
3.5002	-0.0571493\\
3.5028	-0.0573216\\
3.5053	-0.0574869\\
3.5079	-0.0576586\\
3.5104	-0.0578232\\
3.513	-0.057994\\
3.5156	-0.0581643\\
3.5181	-0.0583275\\
3.5207	-0.0584966\\
3.5233	-0.0586651\\
3.5258	-0.0588264\\
3.5284	-0.0589934\\
3.531	-0.0591596\\
3.5335	-0.0593187\\
3.5361	-0.0594831\\
3.5386	-0.0596404\\
3.5412	-0.0598029\\
3.5438	-0.0599645\\
3.5463	-0.0601187\\
3.5489	-0.060278\\
3.5515	-0.0604361\\
3.554	-0.060587\\
3.5566	-0.0607426\\
3.5592	-0.0608969\\
3.5617	-0.0610439\\
3.5643	-0.0611955\\
3.5668	-0.0613398\\
3.5694	-0.0614884\\
3.572	-0.0616355\\
3.5745	-0.0617754\\
3.5771	-0.0619193\\
3.5797	-0.0620615\\
3.5822	-0.0621967\\
3.5848	-0.0623355\\
3.5874	-0.0624726\\
3.5899	-0.0626026\\
3.5925	-0.062736\\
3.595	-0.0628625\\
3.5976	-0.0629921\\
3.6002	-0.0631197\\
3.6027	-0.0632406\\
3.6053	-0.0633642\\
3.6079	-0.0634858\\
3.6104	-0.0636007\\
3.613	-0.0637181\\
3.6156	-0.0638334\\
3.6181	-0.0639421\\
3.6207	-0.064053\\
3.6232	-0.0641575\\
3.6258	-0.0642639\\
3.6284	-0.064368\\
3.6309	-0.064466\\
3.6335	-0.0645655\\
3.6361	-0.0646626\\
3.6386	-0.0647537\\
3.6412	-0.0648461\\
3.6438	-0.064936\\
3.6463	-0.0650201\\
3.6489	-0.0651052\\
3.6515	-0.0651877\\
3.654	-0.0652647\\
3.6566	-0.0653423\\
3.6591	-0.0654144\\
3.6617	-0.065487\\
3.6643	-0.0655569\\
3.6668	-0.0656217\\
3.6694	-0.0656866\\
3.672	-0.0657488\\
3.6745	-0.0658062\\
3.6771	-0.0658632\\
3.6797	-0.0659176\\
3.6822	-0.0659675\\
3.6848	-0.0660167\\
3.6873	-0.0660615\\
3.6899	-0.0661054\\
3.6925	-0.0661467\\
3.695	-0.0661839\\
3.6976	-0.06622\\
3.7002	-0.0662534\\
3.7027	-0.066283\\
3.7053	-0.0663112\\
3.7079	-0.0663367\\
3.7104	-0.0663587\\
3.713	-0.0663789\\
3.7155	-0.0663959\\
3.7181	-0.066411\\
3.7207	-0.0664234\\
3.7232	-0.0664329\\
3.7258	-0.0664402\\
3.7284	-0.0664449\\
3.7309	-0.0664469\\
3.7335	-0.0664465\\
3.7361	-0.0664435\\
3.7386	-0.0664382\\
3.7412	-0.0664302\\
3.7437	-0.0664201\\
3.7463	-0.0664071\\
3.7489	-0.0663917\\
3.7514	-0.0663745\\
3.754	-0.0663542\\
3.7566	-0.0663314\\
3.7591	-0.0663072\\
3.7617	-0.0662797\\
3.7643	-0.0662499\\
3.7668	-0.0662189\\
3.7694	-0.0661844\\
3.7719	-0.0661491\\
3.7745	-0.0661101\\
3.7771	-0.0660688\\
3.7796	-0.066027\\
3.7822	-0.0659814\\
3.7848	-0.0659336\\
3.7873	-0.0658856\\
3.7899	-0.0658335\\
3.7925	-0.0657794\\
3.795	-0.0657254\\
3.7976	-0.0656673\\
3.8002	-0.0656071\\
3.8027	-0.0655474\\
3.8053	-0.0654834\\
3.8078	-0.06542\\
3.8104	-0.0653523\\
3.813	-0.0652827\\
3.8155	-0.065214\\
3.8181	-0.0651409\\
3.8207	-0.0650661\\
3.8232	-0.0649925\\
3.8258	-0.0649143\\
3.8284	-0.0648345\\
3.8309	-0.0647562\\
3.8335	-0.0646733\\
3.836	-0.0645921\\
3.8386	-0.0645062\\
3.8412	-0.0644189\\
3.8437	-0.0643336\\
3.8463	-0.0642435\\
3.8489	-0.0641521\\
3.8514	-0.064063\\
3.854	-0.0639691\\
3.8566	-0.063874\\
3.8591	-0.0637814\\
3.8617	-0.063684\\
3.8642	-0.0635893\\
3.8668	-0.0634898\\
3.8694	-0.0633893\\
3.8719	-0.0632917\\
3.8745	-0.0631894\\
3.8771	-0.0630861\\
3.8796	-0.062986\\
3.8822	-0.0628811\\
3.8848	-0.0627754\\
3.8873	-0.0626731\\
3.8899	-0.062566\\
3.8924	-0.0624625\\
3.895	-0.0623542\\
3.8976	-0.0622454\\
3.9001	-0.0621402\\
3.9027	-0.0620304\\
3.9053	-0.0619202\\
3.9078	-0.0618138\\
3.9104	-0.0617029\\
3.913	-0.0615916\\
3.9155	-0.0614844\\
3.9181	-0.0613726\\
3.9206	-0.061265\\
3.9232	-0.0611529\\
3.9258	-0.0610408\\
3.9283	-0.0609328\\
3.9309	-0.0608206\\
3.9335	-0.0607084\\
3.936	-0.0606005\\
3.9386	-0.0604884\\
3.9412	-0.0603765\\
3.9437	-0.0602691\\
3.9463	-0.0601575\\
3.9488	-0.0600505\\
3.9514	-0.0599395\\
3.954	-0.0598289\\
3.9565	-0.0597229\\
3.9591	-0.059613\\
3.9617	-0.0595035\\
3.9642	-0.0593988\\
3.9668	-0.0592904\\
3.9694	-0.0591825\\
3.9719	-0.0590793\\
3.9745	-0.0589727\\
3.9771	-0.0588667\\
3.9796	-0.0587654\\
3.9822	-0.0586608\\
3.9847	-0.058561\\
3.9873	-0.058458\\
3.9899	-0.0583558\\
3.9924	-0.0582584\\
3.995	-0.0581579\\
3.9976	-0.0580584\\
4.0001	-0.0579636\\
4.0027	-0.057866\\
4.0053	-0.0577695\\
4.0078	-0.0576776\\
4.0104	-0.0575832\\
4.0129	-0.0574934\\
4.0155	-0.0574012\\
4.0181	-0.0573101\\
4.0206	-0.0572237\\
4.0232	-0.057135\\
4.0258	-0.0570475\\
4.0283	-0.0569647\\
4.0309	-0.0568798\\
4.0335	-0.0567962\\
4.036	-0.0567171\\
4.0386	-0.0566361\\
4.0411	-0.0565596\\
4.0437	-0.0564815\\
4.0463	-0.0564047\\
4.0488	-0.0563323\\
4.0514	-0.0562584\\
4.054	-0.0561861\\
4.0565	-0.0561179\\
4.0591	-0.0560485\\
4.0617	-0.0559807\\
4.0642	-0.0559169\\
4.0668	-0.0558521\\
4.0693	-0.0557913\\
4.0719	-0.0557297\\
4.0745	-0.0556697\\
4.077	-0.0556135\\
4.0796	-0.0555567\\
4.0822	-0.0555015\\
4.0847	-0.05545\\
4.0873	-0.0553982\\
4.0899	-0.055348\\
4.0924	-0.0553013\\
4.095	-0.0552544\\
4.0975	-0.0552109\\
4.1001	-0.0551674\\
4.1027	-0.0551256\\
4.1052	-0.055087\\
4.1078	-0.0550486\\
4.1104	-0.0550119\\
4.1129	-0.0549782\\
4.1155	-0.0549449\\
4.1181	-0.0549133\\
4.1206	-0.0548846\\
4.1232	-0.0548564\\
4.1258	-0.05483\\
4.1283	-0.0548062\\
4.1309	-0.0547832\\
4.1334	-0.0547626\\
4.136	-0.054743\\
4.1386	-0.054725\\
4.1411	-0.0547094\\
4.1437	-0.0546948\\
4.1463	-0.0546819\\
4.1488	-0.0546711\\
4.1514	-0.0546615\\
4.154	-0.0546536\\
4.1565	-0.0546475\\
4.1591	-0.0546428\\
4.1616	-0.0546399\\
4.1642	-0.0546384\\
4.1668	-0.0546386\\
4.1693	-0.0546403\\
4.1719	-0.0546436\\
4.1745	-0.0546484\\
4.177	-0.0546546\\
4.1796	-0.0546625\\
4.1822	-0.054672\\
4.1847	-0.0546825\\
4.1873	-0.0546949\\
4.1898	-0.0547083\\
4.1924	-0.0547236\\
4.195	-0.0547403\\
4.1975	-0.0547578\\
4.2001	-0.0547773\\
4.2027	-0.0547982\\
4.2052	-0.0548196\\
4.2078	-0.0548432\\
4.2104	-0.054868\\
4.2129	-0.0548932\\
4.2155	-0.0549206\\
4.218	-0.0549481\\
4.2206	-0.0549779\\
4.2232	-0.0550089\\
4.2257	-0.0550399\\
4.2283	-0.0550732\\
4.2309	-0.0551076\\
4.2334	-0.0551417\\
4.236	-0.0551783\\
4.2386	-0.0552159\\
4.2411	-0.055253\\
4.2437	-0.0552926\\
4.2462	-0.0553315\\
4.2488	-0.0553729\\
4.2514	-0.0554153\\
4.2539	-0.0554568\\
4.2565	-0.0555008\\
4.2591	-0.0555457\\
4.2616	-0.0555895\\
4.2642	-0.0556359\\
4.2668	-0.0556829\\
4.2693	-0.0557288\\
4.2719	-0.0557772\\
4.2744	-0.0558244\\
4.277	-0.055874\\
4.2796	-0.0559242\\
4.2821	-0.0559729\\
4.2847	-0.0560241\\
4.2873	-0.0560758\\
4.2898	-0.056126\\
4.2924	-0.0561785\\
4.295	-0.0562314\\
4.2975	-0.0562826\\
4.3001	-0.0563361\\
4.3027	-0.05639\\
4.3052	-0.056442\\
4.3078	-0.0564963\\
4.3103	-0.0565486\\
4.3129	-0.0566033\\
4.3155	-0.056658\\
4.318	-0.0567107\\
4.3206	-0.0567656\\
4.3232	-0.0568205\\
4.3257	-0.0568733\\
4.3283	-0.0569282\\
4.3309	-0.056983\\
4.3334	-0.0570356\\
4.336	-0.0570902\\
4.3385	-0.0571425\\
4.3411	-0.0571967\\
4.3437	-0.0572507\\
4.3462	-0.0573024\\
4.3488	-0.0573558\\
4.3514	-0.057409\\
4.3539	-0.0574597\\
4.3565	-0.0575122\\
4.3591	-0.0575642\\
4.3616	-0.0576138\\
4.3642	-0.0576649\\
4.3667	-0.0577136\\
4.3693	-0.0577638\\
4.3719	-0.0578133\\
4.3744	-0.0578605\\
4.377	-0.0579089\\
4.3796	-0.0579567\\
4.3821	-0.058002\\
4.3847	-0.0580485\\
4.3873	-0.0580943\\
4.3898	-0.0581376\\
4.3924	-0.0581819\\
4.3949	-0.0582237\\
4.3975	-0.0582665\\
4.4001	-0.0583084\\
4.4026	-0.0583479\\
4.4052	-0.0583881\\
4.4078	-0.0584273\\
4.4103	-0.0584643\\
4.4129	-0.0585017\\
4.4155	-0.0585382\\
4.418	-0.0585724\\
4.4206	-0.058607\\
4.4231	-0.0586392\\
4.4257	-0.0586717\\
4.4283	-0.0587032\\
4.4308	-0.0587324\\
4.4334	-0.0587618\\
4.436	-0.05879\\
4.4385	-0.058816\\
4.4411	-0.058842\\
4.4437	-0.0588669\\
4.4462	-0.0588896\\
4.4488	-0.0589121\\
4.4514	-0.0589334\\
4.4539	-0.0589528\\
4.4565	-0.0589717\\
4.459	-0.0589887\\
4.4616	-0.0590052\\
4.4642	-0.0590204\\
4.4667	-0.0590339\\
4.4693	-0.0590466\\
4.4719	-0.0590581\\
4.4744	-0.0590679\\
4.477	-0.0590768\\
4.4796	-0.0590843\\
4.4821	-0.0590904\\
4.4847	-0.0590954\\
4.4872	-0.059099\\
4.4898	-0.0591014\\
4.4924	-0.0591025\\
4.4949	-0.0591023\\
4.4975	-0.0591007\\
4.5001	-0.0590979\\
4.5026	-0.0590938\\
4.5052	-0.0590883\\
4.5078	-0.0590815\\
4.5103	-0.0590736\\
4.5129	-0.0590641\\
4.5154	-0.0590538\\
4.518	-0.0590417\\
4.5206	-0.0590282\\
4.5231	-0.059014\\
4.5257	-0.058998\\
4.5283	-0.0589806\\
4.5308	-0.0589627\\
4.5334	-0.0589428\\
4.536	-0.0589215\\
4.5385	-0.0588999\\
4.5411	-0.0588761\\
4.5436	-0.058852\\
4.5462	-0.0588257\\
4.5488	-0.0587981\\
4.5513	-0.0587704\\
4.5539	-0.0587404\\
4.5565	-0.0587091\\
4.559	-0.0586779\\
4.5616	-0.0586443\\
4.5642	-0.0586094\\
4.5667	-0.0585748\\
4.5693	-0.0585376\\
4.5718	-0.0585008\\
4.5744	-0.0584613\\
4.577	-0.0584207\\
4.5795	-0.0583806\\
4.5821	-0.0583379\\
4.5847	-0.058294\\
4.5872	-0.0582508\\
4.5898	-0.0582049\\
4.5924	-0.0581579\\
4.5949	-0.0581117\\
4.5975	-0.0580627\\
4.6001	-0.0580127\\
4.6026	-0.0579638\\
4.6052	-0.0579119\\
4.6077	-0.0578611\\
4.6103	-0.0578075\\
4.6129	-0.0577529\\
4.6154	-0.0576996\\
4.618	-0.0576433\\
4.6206	-0.0575862\\
4.6231	-0.0575306\\
4.6257	-0.0574719\\
4.6283	-0.0574125\\
4.6308	-0.0573546\\
4.6334	-0.0572937\\
4.6359	-0.0572345\\
4.6385	-0.0571723\\
4.6411	-0.0571094\\
4.6436	-0.0570483\\
4.6462	-0.0569842\\
4.6488	-0.0569195\\
4.6513	-0.0568568\\
4.6539	-0.056791\\
4.6565	-0.0567247\\
4.659	-0.0566605\\
4.6616	-0.0565933\\
4.6641	-0.0565283\\
4.6667	-0.0564602\\
4.6693	-0.0563918\\
4.6718	-0.0563256\\
4.6744	-0.0562565\\
4.677	-0.056187\\
4.6795	-0.05612\\
4.6821	-0.05605\\
4.6847	-0.0559798\\
4.6872	-0.0559121\\
4.6898	-0.0558414\\
4.6923	-0.0557734\\
4.6949	-0.0557025\\
4.6975	-0.0556315\\
4.7	-0.0555631\\
4.7026	-0.055492\\
4.7052	-0.0554208\\
4.7077	-0.0553524\\
4.7103	-0.0552812\\
4.7129	-0.0552101\\
4.7154	-0.0551418\\
4.718	-0.0550709\\
4.7205	-0.0550028\\
4.7231	-0.0549322\\
4.7257	-0.0548617\\
4.7282	-0.0547941\\
4.7308	-0.0547241\\
4.7334	-0.0546542\\
4.7359	-0.0545874\\
4.7385	-0.0545182\\
4.7411	-0.0544492\\
4.7436	-0.0543833\\
4.7462	-0.0543151\\
4.7487	-0.05425\\
4.7513	-0.0541826\\
4.7539	-0.0541157\\
4.7564	-0.0540518\\
4.759	-0.0539858\\
4.7616	-0.0539204\\
4.7641	-0.0538579\\
4.7667	-0.0537936\\
4.7693	-0.0537298\\
4.7718	-0.0536691\\
4.7744	-0.0536066\\
4.777	-0.0535447\\
4.7795	-0.0534859\\
4.7821	-0.0534254\\
4.7846	-0.0533679\\
4.7872	-0.0533088\\
4.7898	-0.0532505\\
4.7923	-0.0531952\\
4.7949	-0.0531384\\
4.7975	-0.0530825\\
4.8	-0.0530295\\
4.8026	-0.0529753\\
4.8052	-0.0529219\\
4.8077	-0.0528715\\
4.8103	-0.0528199\\
4.8128	-0.0527712\\
4.8154	-0.0527215\\
4.818	-0.0526728\\
4.8205	-0.0526268\\
4.8231	-0.0525801\\
4.8257	-0.0525343\\
4.8282	-0.0524913\\
4.8308	-0.0524476\\
4.8334	-0.0524049\\
4.8359	-0.0523649\\
4.8385	-0.0523244\\
4.841	-0.0522865\\
4.8436	-0.0522482\\
4.8462	-0.052211\\
4.8487	-0.0521763\\
4.8513	-0.0521413\\
4.8539	-0.0521076\\
4.8564	-0.0520762\\
4.859	-0.0520448\\
4.8616	-0.0520145\\
4.8641	-0.0519865\\
4.8667	-0.0519587\\
4.8692	-0.051933\\
4.8718	-0.0519076\\
4.8744	-0.0518833\\
4.8769	-0.0518612\\
4.8795	-0.0518394\\
4.8821	-0.0518189\\
4.8846	-0.0518004\\
4.8872	-0.0517824\\
4.8898	-0.0517656\\
4.8923	-0.0517507\\
4.8949	-0.0517365\\
4.8974	-0.051724\\
4.9	-0.0517123\\
4.9026	-0.0517018\\
4.9051	-0.051693\\
4.9077	-0.0516851\\
4.9103	-0.0516786\\
4.9128	-0.0516734\\
4.9154	-0.0516694\\
4.918	-0.0516666\\
4.9205	-0.0516652\\
4.9231	-0.0516649\\
4.9257	-0.051666\\
4.9282	-0.0516682\\
4.9308	-0.0516717\\
4.9333	-0.0516764\\
4.9359	-0.0516824\\
4.9385	-0.0516898\\
4.941	-0.051698\\
4.9436	-0.0517078\\
4.9462	-0.0517189\\
4.9487	-0.0517307\\
4.9513	-0.0517442\\
4.9539	-0.0517589\\
4.9564	-0.0517742\\
4.959	-0.0517913\\
4.9615	-0.0518089\\
4.9641	-0.0518284\\
4.9667	-0.051849\\
4.9692	-0.05187\\
4.9718	-0.0518929\\
4.9744	-0.051917\\
4.9769	-0.0519413\\
4.9795	-0.0519676\\
4.9821	-0.051995\\
4.9846	-0.0520225\\
4.9872	-0.0520521\\
4.9897	-0.0520816\\
4.9923	-0.0521133\\
4.9949	-0.052146\\
4.9974	-0.0521785\\
5	-0.0522133\\
};
\addplot [color=mycolor8, line width=1.0pt, forget plot]
  table[row sep=crcr]{%
-0.125	2.82238e-08\\
-0.1224	2.87521e-08\\
-0.1199	2.92588e-08\\
-0.1173	2.97846e-08\\
-0.1147	3.03092e-08\\
-0.1122	3.08126e-08\\
-0.1096	3.13347e-08\\
-0.1071	3.18356e-08\\
-0.1045	3.23556e-08\\
-0.1019	3.28742e-08\\
-0.0994	3.33719e-08\\
-0.0968	3.38885e-08\\
-0.0942	3.44037e-08\\
-0.0917	3.4898e-08\\
-0.0891	3.54112e-08\\
-0.0865	3.59231e-08\\
-0.084	3.64143e-08\\
-0.0814	3.69243e-08\\
-0.0789	3.74133e-08\\
-0.0763	3.79209e-08\\
-0.0737	3.84276e-08\\
-0.0712	3.89137e-08\\
-0.0686	3.94183e-08\\
-0.066	3.99219e-08\\
-0.0635	4.0405e-08\\
-0.0609	4.09065e-08\\
-0.0583	4.14071e-08\\
-0.0558	4.18876e-08\\
-0.0532	4.23864e-08\\
-0.0507	4.28651e-08\\
-0.0481	4.33619e-08\\
-0.0455	4.38581e-08\\
-0.043	4.43342e-08\\
-0.0404	4.48286e-08\\
-0.0378	4.53224e-08\\
-0.0353	4.57961e-08\\
-0.0327	4.62881e-08\\
-0.0301	4.67796e-08\\
-0.0276	4.72513e-08\\
-0.025	4.7741e-08\\
-0.0224	4.82305e-08\\
-0.0199	4.87002e-08\\
-0.0173	4.91879e-08\\
-0.0148	4.96567e-08\\
-0.0122	5.01435e-08\\
-0.0096	5.06295e-08\\
-0.0071	5.10966e-08\\
-0.0045	5.15818e-08\\
-0.0019	5.20662e-08\\
0.0006	5.25318e-08\\
0.0032	0.0140057\\
0.0058	0.0238141\\
0.0083	0.0314245\\
0.0109	0.0385706\\
0.0134	0.0449103\\
0.016	0.0507844\\
0.0186	0.0558252\\
0.0211	0.0600315\\
0.0237	0.0640235\\
0.0263	0.0678334\\
0.0288	0.0713625\\
0.0314	0.0748137\\
0.034	0.0779456\\
0.0365	0.0806355\\
0.0391	0.0831553\\
0.0416	0.0853785\\
0.0442	0.0875254\\
0.0468	0.0894947\\
0.0493	0.0911877\\
0.0519	0.0927183\\
0.0545	0.094023\\
0.057	0.095095\\
0.0596	0.0960526\\
0.0622	0.0968619\\
0.0647	0.097491\\
0.0673	0.0979746\\
0.0698	0.0982731\\
0.0724	0.098425\\
0.075	0.0984414\\
0.0775	0.0983471\\
0.0801	0.0981355\\
0.0827	0.0977981\\
0.0852	0.097347\\
0.0878	0.096751\\
0.0904	0.0960445\\
0.0929	0.0952798\\
0.0955	0.0944039\\
0.098	0.093479\\
0.1006	0.0924207\\
0.1032	0.0912613\\
0.1057	0.0900611\\
0.1083	0.0887417\\
0.1109	0.0873625\\
0.1134	0.0859802\\
0.116	0.0844744\\
0.1186	0.0828909\\
0.1211	0.0812975\\
0.1237	0.0795801\\
0.1263	0.0778155\\
0.1288	0.0760793\\
0.1314	0.0742264\\
0.1339	0.0723892\\
0.1365	0.0704171\\
0.1391	0.0683895\\
0.1416	0.0664004\\
0.1442	0.0642998\\
0.1468	0.062165\\
0.1493	0.0600709\\
0.1519	0.0578421\\
0.1545	0.0555626\\
0.157	0.0533326\\
0.1596	0.0509847\\
0.1621	0.0487027\\
0.1647	0.0462982\\
0.1673	0.0438521\\
0.1698	0.0414573\\
0.1724	0.0389277\\
0.175	0.0363694\\
0.1775	0.0338896\\
0.1801	0.0312882\\
0.1827	0.028656\\
0.1852	0.0260894\\
0.1878	0.0233842\\
0.1903	0.0207567\\
0.1929	0.0180062\\
0.1955	0.0152408\\
0.198	0.0125632\\
0.2006	0.00975184\\
0.2032	0.00691047\\
0.2057	0.0041552\\
0.2083	0.00127534\\
0.2109	-0.00161219\\
0.2134	-0.00439632\\
0.216	-0.00730576\\
0.2185	-0.0101214\\
0.2211	-0.0130667\\
0.2237	-0.0160213\\
0.2262	-0.0188624\\
0.2288	-0.0218134\\
0.2314	-0.0247639\\
0.2339	-0.0276054\\
0.2365	-0.0305665\\
0.2391	-0.0335282\\
0.2416	-0.036368\\
0.2442	-0.039306\\
0.2467	-0.0421159\\
0.2493	-0.0450263\\
0.2519	-0.0479279\\
0.2544	-0.0507075\\
0.257	-0.05358\\
0.2596	-0.0564251\\
0.2621	-0.0591313\\
0.2647	-0.0619167\\
0.2673	-0.0646762\\
0.2698	-0.0673056\\
0.2724	-0.0700101\\
0.2749	-0.072574\\
0.2775	-0.075196\\
0.2801	-0.0777715\\
0.2826	-0.0802068\\
0.2852	-0.082699\\
0.2878	-0.0851476\\
0.2903	-0.0874547\\
0.2929	-0.0897975\\
0.2955	-0.0920789\\
0.298	-0.0942157\\
0.3006	-0.0963819\\
0.3032	-0.0984911\\
0.3057	-0.100462\\
0.3083	-0.102444\\
0.3108	-0.104281\\
0.3134	-0.106119\\
0.316	-0.107885\\
0.3185	-0.109517\\
0.3211	-0.111145\\
0.3237	-0.112696\\
0.3262	-0.11411\\
0.3288	-0.115498\\
0.3314	-0.116802\\
0.3339	-0.11798\\
0.3365	-0.119127\\
0.339	-0.120152\\
0.3416	-0.121132\\
0.3442	-0.122021\\
0.3467	-0.122789\\
0.3493	-0.1235\\
0.3519	-0.124125\\
0.3544	-0.124645\\
0.357	-0.125097\\
0.3596	-0.125455\\
0.3621	-0.125709\\
0.3647	-0.125881\\
0.3672	-0.125962\\
0.3698	-0.125958\\
0.3724	-0.125865\\
0.3749	-0.125688\\
0.3775	-0.12541\\
0.3801	-0.125037\\
0.3826	-0.124593\\
0.3852	-0.124044\\
0.3878	-0.123407\\
0.3903	-0.122711\\
0.3929	-0.121898\\
0.3954	-0.121029\\
0.398	-0.120038\\
0.4006	-0.118962\\
0.4031	-0.117848\\
0.4057	-0.116609\\
0.4083	-0.115285\\
0.4108	-0.113932\\
0.4134	-0.112442\\
0.416	-0.110873\\
0.4185	-0.109291\\
0.4211	-0.107573\\
0.4236	-0.105849\\
0.4262	-0.103983\\
0.4288	-0.10204\\
0.4313	-0.100104\\
0.4339	-0.0980238\\
0.4365	-0.0958779\\
0.439	-0.0937538\\
0.4416	-0.0914811\\
0.4442	-0.0891439\\
0.4467	-0.0868375\\
0.4493	-0.0843813\\
0.4519	-0.0818702\\
0.4544	-0.0794061\\
0.457	-0.0767924\\
0.4595	-0.0742301\\
0.4621	-0.0715155\\
0.4647	-0.0687532\\
0.4672	-0.0660561\\
0.4698	-0.0632116\\
0.4724	-0.0603283\\
0.4749	-0.0575197\\
0.4775	-0.0545617\\
0.4801	-0.0515685\\
0.4826	-0.0486605\\
0.4852	-0.0456088\\
0.4877	-0.0426506\\
0.4903	-0.0395501\\
0.4929	-0.0364256\\
0.4954	-0.0334002\\
0.498	-0.0302348\\
0.5006	-0.0270539\\
0.5031	-0.0239835\\
0.5057	-0.0207798\\
0.5083	-0.0175661\\
0.5108	-0.0144673\\
0.5134	-0.0112378\\
0.5159	-0.00812916\\
0.5185	-0.00489602\\
0.5211	-0.00166505\\
0.5236	0.00143861\\
0.5262	0.0046626\\
0.5288	0.00788127\\
0.5313	0.0109686\\
0.5339	0.0141683\\
0.5365	0.0173541\\
0.539	0.0204026\\
0.5416	0.0235572\\
0.5441	0.0265744\\
0.5467	0.0296938\\
0.5493	0.0327915\\
0.5518	0.0357468\\
0.5544	0.0387944\\
0.557	0.0418146\\
0.5595	0.0446924\\
0.5621	0.047657\\
0.5647	0.0505904\\
0.5672	0.0533791\\
0.5698	0.0562441\\
0.5723	0.0589641\\
0.5749	0.0617561\\
0.5775	0.0645101\\
0.58	0.0671211\\
0.5826	0.0697957\\
0.5852	0.0724267\\
0.5877	0.074914\\
0.5903	0.077456\\
0.5929	0.0799523\\
0.5954	0.0823089\\
0.598	0.0847126\\
0.6006	0.0870664\\
0.6031	0.0892812\\
0.6057	0.0915336\\
0.6082	0.0936502\\
0.6108	0.0958003\\
0.6134	0.0978973\\
0.6159	0.0998622\\
0.6185	0.101851\\
0.6211	0.103783\\
0.6236	0.105586\\
0.6262	0.107407\\
0.6288	0.10917\\
0.6313	0.110812\\
0.6339	0.112461\\
0.6364	0.11399\\
0.639	0.115521\\
0.6416	0.116993\\
0.6441	0.118352\\
0.6467	0.119707\\
0.6493	0.121001\\
0.6518	0.122187\\
0.6544	0.123361\\
0.657	0.124473\\
0.6595	0.125485\\
0.6621	0.126479\\
0.6646	0.127376\\
0.6672	0.12825\\
0.6698	0.129061\\
0.6723	0.129784\\
0.6749	0.130475\\
0.6775	0.131106\\
0.68	0.131657\\
0.6826	0.13217\\
0.6852	0.132623\\
0.6877	0.133001\\
0.6903	0.133335\\
0.6928	0.133602\\
0.6954	0.133821\\
0.698	0.133983\\
0.7005	0.134083\\
0.7031	0.13413\\
0.7057	0.13412\\
0.7082	0.134056\\
0.7108	0.133936\\
0.7134	0.13376\\
0.7159	0.133539\\
0.7185	0.133255\\
0.721	0.13293\\
0.7236	0.13254\\
0.7262	0.132098\\
0.7287	0.131623\\
0.7313	0.13108\\
0.7339	0.130485\\
0.7364	0.129866\\
0.739	0.129173\\
0.7416	0.128431\\
0.7441	0.127673\\
0.7467	0.126839\\
0.7492	0.125993\\
0.7518	0.125068\\
0.7544	0.124097\\
0.7569	0.123122\\
0.7595	0.122066\\
0.7621	0.120968\\
0.7646	0.119872\\
0.7672	0.118693\\
0.7698	0.117473\\
0.7723	0.116262\\
0.7749	0.114966\\
0.7775	0.113632\\
0.78	0.112316\\
0.7826	0.110911\\
0.7851	0.109527\\
0.7877	0.108054\\
0.7903	0.106547\\
0.7928	0.105068\\
0.7954	0.103499\\
0.798	0.101899\\
0.8005	0.100333\\
0.8031	0.0986747\\
0.8057	0.0969882\\
0.8082	0.0953408\\
0.8108	0.0936019\\
0.8133	0.0919059\\
0.8159	0.0901177\\
0.8185	0.0883052\\
0.821	0.0865403\\
0.8236	0.0846828\\
0.8262	0.082804\\
0.8287	0.0809783\\
0.8313	0.0790601\\
0.8339	0.0771228\\
0.8364	0.0752424\\
0.839	0.0732694\\
0.8415	0.0713566\\
0.8441	0.069352\\
0.8467	0.0673324\\
0.8492	0.0653768\\
0.8518	0.0633294\\
0.8544	0.0612688\\
0.8569	0.059276\\
0.8595	0.0571923\\
0.8621	0.0550981\\
0.8646	0.053075\\
0.8672	0.0509617\\
0.8697	0.0489212\\
0.8723	0.0467911\\
0.8749	0.0446536\\
0.8774	0.0425922\\
0.88	0.0404425\\
0.8826	0.0382874\\
0.8851	0.0362103\\
0.8877	0.0340458\\
0.8903	0.0318776\\
0.8928	0.0297899\\
0.8954	0.0276165\\
0.8979	0.0255249\\
0.9005	0.0233481\\
0.9031	0.0211702\\
0.9056	0.0190755\\
0.9082	0.0168973\\
0.9108	0.0147198\\
0.9133	0.0126272\\
0.9159	0.0104525\\
0.9185	0.00827963\\
0.921	0.0061925\\
0.9236	0.00402469\\
0.9262	0.00186037\\
0.9287	-0.000216991\\
0.9313	-0.00237328\\
0.9338	-0.00444242\\
0.9364	-0.00658966\\
0.939	-0.0087317\\
0.9415	-0.0107859\\
0.9441	-0.0129163\\
0.9467	-0.0150401\\
0.9492	-0.0170758\\
0.9518	-0.0191863\\
0.9544	-0.0212895\\
0.9569	-0.0233046\\
0.9595	-0.0253923\\
0.962	-0.0273918\\
0.9646	-0.0294629\\
0.9672	-0.0315251\\
0.9697	-0.0334996\\
0.9723	-0.0355439\\
0.9749	-0.0375785\\
0.9774	-0.0395255\\
0.98	-0.0415402\\
0.9826	-0.0435445\\
0.9851	-0.0454618\\
0.9877	-0.0474453\\
0.9902	-0.049342\\
0.9928	-0.0513034\\
0.9954	-0.0532532\\
0.9979	-0.0551167\\
1.0005	-0.0570431\\
1.0031	-0.0589575\\
1.0056	-0.0607865\\
1.0082	-0.0626763\\
1.0108	-0.0645532\\
1.0133	-0.0663454\\
1.0159	-0.0681964\\
1.0184	-0.0699635\\
1.021	-0.071788\\
1.0236	-0.0735986\\
1.0261	-0.0753262\\
1.0287	-0.0771087\\
1.0313	-0.0788766\\
1.0338	-0.0805627\\
1.0364	-0.0823016\\
1.039	-0.0840254\\
1.0415	-0.0856683\\
1.0441	-0.0873614\\
1.0466	-0.0889744\\
1.0492	-0.0906361\\
1.0518	-0.0922815\\
1.0543	-0.093848\\
1.0569	-0.0954606\\
1.0595	-0.0970561\\
1.062	-0.0985737\\
1.0646	-0.100135\\
1.0672	-0.101678\\
1.0697	-0.103144\\
1.0723	-0.104652\\
1.0748	-0.106083\\
1.0774	-0.107553\\
1.08	-0.109004\\
1.0825	-0.11038\\
1.0851	-0.111792\\
1.0877	-0.113184\\
1.0902	-0.114502\\
1.0928	-0.115853\\
1.0954	-0.117183\\
1.0979	-0.11844\\
1.1005	-0.119727\\
1.1031	-0.120992\\
1.1056	-0.122186\\
1.1082	-0.123406\\
1.1107	-0.124556\\
1.1133	-0.12573\\
1.1159	-0.126879\\
1.1184	-0.127961\\
1.121	-0.129061\\
1.1236	-0.130137\\
1.1261	-0.131147\\
1.1287	-0.132171\\
1.1313	-0.133169\\
1.1338	-0.134103\\
1.1364	-0.135048\\
1.1389	-0.135931\\
1.1415	-0.136821\\
1.1441	-0.137683\\
1.1466	-0.138484\\
1.1492	-0.139288\\
1.1518	-0.140063\\
1.1543	-0.14078\\
1.1569	-0.141495\\
1.1595	-0.14218\\
1.162	-0.142808\\
1.1646	-0.143431\\
1.1671	-0.143999\\
1.1697	-0.144559\\
1.1723	-0.145085\\
1.1748	-0.14556\\
1.1774	-0.146021\\
1.18	-0.146448\\
1.1825	-0.146826\\
1.1851	-0.147186\\
1.1877	-0.14751\\
1.1902	-0.147789\\
1.1928	-0.148044\\
1.1953	-0.148256\\
1.1979	-0.14844\\
1.2005	-0.148588\\
1.203	-0.148696\\
1.2056	-0.148772\\
1.2082	-0.148811\\
1.2107	-0.148813\\
1.2133	-0.148779\\
1.2159	-0.148707\\
1.2184	-0.148602\\
1.221	-0.148455\\
1.2235	-0.148279\\
1.2261	-0.148057\\
1.2287	-0.147798\\
1.2312	-0.147512\\
1.2338	-0.147177\\
1.2364	-0.146804\\
1.2389	-0.146409\\
1.2415	-0.145961\\
1.2441	-0.145474\\
1.2466	-0.14497\\
1.2492	-0.144408\\
1.2518	-0.143808\\
1.2543	-0.143196\\
1.2569	-0.142522\\
1.2594	-0.141839\\
1.262	-0.141091\\
1.2646	-0.140307\\
1.2671	-0.139518\\
1.2697	-0.138661\\
1.2723	-0.137768\\
1.2748	-0.136876\\
1.2774	-0.135913\\
1.28	-0.134915\\
1.2825	-0.133923\\
1.2851	-0.132857\\
1.2876	-0.1318\\
1.2902	-0.130668\\
1.2928	-0.129504\\
1.2953	-0.128353\\
1.2979	-0.127126\\
1.3005	-0.125867\\
1.303	-0.124627\\
1.3056	-0.123309\\
1.3082	-0.121961\\
1.3107	-0.120638\\
1.3133	-0.119234\\
1.3158	-0.117859\\
1.3184	-0.116403\\
1.321	-0.11492\\
1.3235	-0.113471\\
1.3261	-0.11194\\
1.3287	-0.110386\\
1.3312	-0.10887\\
1.3338	-0.107272\\
1.3364	-0.105652\\
1.3389	-0.104076\\
1.3415	-0.102418\\
1.344	-0.100807\\
1.3466	-0.0991133\\
1.3492	-0.0974034\\
1.3517	-0.0957444\\
1.3543	-0.0940046\\
1.3569	-0.092251\\
1.3594	-0.0905527\\
1.362	-0.0887747\\
1.3646	-0.0869859\\
1.3671	-0.0852564\\
1.3697	-0.083449\\
1.3722	-0.0817036\\
1.3748	-0.0798815\\
1.3774	-0.0780534\\
1.3799	-0.0762909\\
1.3825	-0.0744539\\
1.3851	-0.072614\\
1.3876	-0.0708429\\
1.3902	-0.0690001\\
1.3928	-0.0671572\\
1.3953	-0.0653861\\
1.3979	-0.0635461\\
1.4005	-0.0617092\\
1.403	-0.0599466\\
1.4056	-0.0581184\\
1.4081	-0.0563661\\
1.4107	-0.0545505\\
1.4133	-0.0527427\\
1.4158	-0.0510128\\
1.4184	-0.0492232\\
1.421	-0.0474444\\
1.4235	-0.0457449\\
1.4261	-0.0439897\\
1.4287	-0.0422481\\
1.4312	-0.0405869\\
1.4338	-0.0388742\\
1.4363	-0.0372426\\
1.4389	-0.0355624\\
1.4415	-0.0338999\\
1.444	-0.032319\\
1.4466	-0.0306939\\
1.4492	-0.0290891\\
1.4517	-0.0275659\\
1.4543	-0.0260031\\
1.4569	-0.0244629\\
1.4594	-0.023004\\
1.462	-0.0215102\\
1.4645	-0.0200972\\
1.4671	-0.0186527\\
1.4697	-0.0172343\\
1.4722	-0.0158956\\
1.4748	-0.0145303\\
1.4774	-0.013193\\
1.4799	-0.0119341\\
1.4825	-0.0106534\\
1.4851	-0.00940251\\
1.4876	-0.00822829\\
1.4902	-0.00703733\\
1.4927	-0.0059217\\
1.4953	-0.00479262\\
1.4979	-0.00369578\\
1.5004	-0.00267197\\
1.503	-0.00163971\\
1.5056	-0.000640979\\
1.5081	0.000287355\\
1.5107	0.00121921\\
1.5133	0.00211645\\
1.5158	0.00294622\\
1.5184	0.00377461\\
1.5209	0.00453762\\
1.5235	0.00529604\\
1.5261	0.00601842\\
1.5286	0.00667882\\
1.5312	0.00732987\\
1.5338	0.00794426\\
1.5363	0.0085003\\
1.5389	0.00904231\\
1.5415	0.00954722\\
1.544	0.00999764\\
1.5466	0.0104295\\
1.5491	0.0108095\\
1.5517	0.0111681\\
1.5543	0.0114892\\
1.5568	0.0117626\\
1.5594	0.0120103\\
1.562	0.0122205\\
1.5645	0.0123874\\
1.5671	0.0125244\\
1.5697	0.0126242\\
1.5722	0.0126851\\
1.5748	0.0127121\\
1.5774	0.0127023\\
1.5799	0.0126581\\
1.5825	0.0125764\\
1.585	0.0124635\\
1.5876	0.0123105\\
1.5902	0.0121216\\
1.5927	0.0119063\\
1.5953	0.0116475\\
1.5979	0.0113535\\
1.6004	0.0110379\\
1.603	0.0106756\\
1.6056	0.010279\\
1.6081	0.00986557\\
1.6107	0.00940257\\
1.6132	0.00892595\\
1.6158	0.00839793\\
1.6184	0.00783733\\
1.6209	0.00726794\\
1.6235	0.00664459\\
1.6261	0.0059899\\
1.6286	0.00533123\\
1.6312	0.00461632\\
1.6338	0.0038714\\
1.6363	0.00312727\\
1.6389	0.00232486\\
1.6414	0.00152638\\
1.644	0.000668426\\
1.6466	-0.000217067\\
1.6491	-0.00109398\\
1.6517	-0.00203195\\
1.6543	-0.00299585\\
1.6568	-0.00394661\\
1.6594	-0.00495972\\
1.662	-0.00599706\\
1.6645	-0.00701678\\
1.6671	-0.00809989\\
1.6696	-0.00916251\\
1.6722	-0.0102891\\
1.6748	-0.0114368\\
1.6773	-0.0125598\\
1.6799	-0.0137473\\
1.6825	-0.0149541\\
1.685	-0.016132\\
1.6876	-0.0173747\\
1.6902	-0.0186347\\
1.6927	-0.019862\\
1.6953	-0.021154\\
1.6978	-0.0224107\\
1.7004	-0.023732\\
1.703	-0.0250672\\
1.7055	-0.0263636\\
1.7081	-0.0277241\\
1.7107	-0.0290963\\
1.7132	-0.0304263\\
1.7158	-0.0318196\\
1.7184	-0.0332225\\
1.7209	-0.03458\\
1.7235	-0.0359997\\
1.7261	-0.037427\\
1.7286	-0.0388057\\
1.7312	-0.0402455\\
1.7337	-0.0416349\\
1.7363	-0.0430843\\
1.7389	-0.0445376\\
1.7414	-0.0459379\\
1.744	-0.0473965\\
1.7466	-0.0488567\\
1.7491	-0.0502616\\
1.7517	-0.0517228\\
1.7543	-0.0531834\\
1.7568	-0.0545866\\
1.7594	-0.0560439\\
1.7619	-0.0574425\\
1.7645	-0.0588936\\
1.7671	-0.0603405\\
1.7696	-0.061727\\
1.7722	-0.0631634\\
1.7748	-0.0645935\\
1.7773	-0.0659618\\
1.7799	-0.0673772\\
1.7825	-0.068784\\
1.785	-0.070128\\
1.7876	-0.0715161\\
1.7901	-0.0728407\\
1.7927	-0.0742073\\
1.7953	-0.0755619\\
1.7978	-0.0768525\\
1.8004	-0.0781817\\
1.803	-0.0794968\\
1.8055	-0.0807477\\
1.8081	-0.0820336\\
1.8107	-0.0833036\\
1.8132	-0.0845092\\
1.8158	-0.0857462\\
1.8183	-0.0869189\\
1.8209	-0.0881206\\
1.8235	-0.0893033\\
1.826	-0.0904221\\
1.8286	-0.0915661\\
1.8312	-0.0926894\\
1.8337	-0.0937495\\
1.8363	-0.0948307\\
1.8389	-0.0958898\\
1.8414	-0.0968866\\
1.844	-0.0979006\\
1.8465	-0.0988531\\
1.8491	-0.0998201\\
1.8517	-0.100762\\
1.8542	-0.101645\\
1.8568	-0.102537\\
1.8594	-0.103404\\
1.8619	-0.104213\\
1.8645	-0.105027\\
1.8671	-0.105815\\
1.8696	-0.106546\\
1.8722	-0.107279\\
1.8747	-0.107958\\
1.8773	-0.108636\\
1.8799	-0.109285\\
1.8824	-0.109881\\
1.885	-0.110473\\
1.8876	-0.111035\\
1.8901	-0.111548\\
1.8927	-0.112052\\
1.8953	-0.112525\\
1.8978	-0.112952\\
1.9004	-0.113365\\
1.903	-0.113748\\
1.9055	-0.114087\\
1.9081	-0.11441\\
1.9106	-0.11469\\
1.9132	-0.114952\\
1.9158	-0.115182\\
1.9183	-0.115374\\
1.9209	-0.115543\\
1.9235	-0.115681\\
1.926	-0.115784\\
1.9286	-0.11586\\
1.9312	-0.115906\\
1.9337	-0.11592\\
1.9363	-0.115904\\
1.9388	-0.11586\\
1.9414	-0.115784\\
1.944	-0.115677\\
1.9465	-0.115546\\
1.9491	-0.11538\\
1.9517	-0.115183\\
1.9542	-0.114966\\
1.9568	-0.11471\\
1.9594	-0.114426\\
1.9619	-0.114124\\
1.9645	-0.113782\\
1.967	-0.113426\\
1.9696	-0.113028\\
1.9722	-0.112602\\
1.9747	-0.112166\\
1.9773	-0.111686\\
1.9799	-0.111179\\
1.9824	-0.110666\\
1.985	-0.110106\\
1.9876	-0.109521\\
1.9901	-0.108934\\
1.9927	-0.108299\\
1.9952	-0.107665\\
1.9978	-0.106982\\
2.0004	-0.106275\\
2.0029	-0.105574\\
2.0055	-0.104822\\
2.0081	-0.104047\\
2.0106	-0.103282\\
2.0132	-0.102466\\
2.0158	-0.101628\\
2.0183	-0.100804\\
2.0209	-0.0999276\\
2.0234	-0.0990668\\
2.026	-0.0981534\\
2.0286	-0.0972219\\
2.0311	-0.0963099\\
2.0337	-0.0953448\\
2.0363	-0.0943635\\
2.0388	-0.0934052\\
2.0414	-0.0923938\\
2.044	-0.0913679\\
2.0465	-0.0903685\\
2.0491	-0.0893161\\
2.0517	-0.0882512\\
2.0542	-0.087216\\
2.0568	-0.0861282\\
2.0593	-0.0850722\\
2.0619	-0.0839641\\
2.0645	-0.0828465\\
2.067	-0.0817637\\
2.0696	-0.0806295\\
2.0722	-0.0794879\\
2.0747	-0.0783837\\
2.0773	-0.0772292\\
2.0799	-0.0760691\\
2.0824	-0.074949\\
2.085	-0.0737799\\
2.0875	-0.0726523\\
2.0901	-0.0714766\\
2.0927	-0.0702985\\
2.0952	-0.0691641\\
2.0978	-0.0679832\\
2.1004	-0.0668017\\
2.1029	-0.0656659\\
2.1055	-0.0644852\\
2.1081	-0.063306\\
2.1106	-0.0621739\\
2.1132	-0.0609991\\
2.1157	-0.0598724\\
2.1183	-0.0587042\\
2.1209	-0.0575405\\
2.1234	-0.0564261\\
2.126	-0.0552724\\
2.1286	-0.0541249\\
2.1311	-0.0530276\\
2.1337	-0.0518935\\
2.1363	-0.0507671\\
2.1388	-0.0496918\\
2.1414	-0.048582\\
2.1439	-0.0475237\\
2.1465	-0.0464325\\
2.1491	-0.0453517\\
2.1516	-0.0443226\\
2.1542	-0.0432634\\
2.1568	-0.0422158\\
2.1593	-0.0412201\\
2.1619	-0.0401969\\
2.1645	-0.0391867\\
2.167	-0.0382282\\
2.1696	-0.0372449\\
2.1721	-0.036313\\
2.1747	-0.0353583\\
2.1773	-0.0344186\\
2.1798	-0.0335298\\
2.1824	-0.0326209\\
2.185	-0.0317282\\
2.1875	-0.0308855\\
2.1901	-0.0300256\\
2.1927	-0.029183\\
2.1952	-0.0283892\\
2.1978	-0.0275812\\
2.2004	-0.0267913\\
2.2029	-0.0260492\\
2.2055	-0.0252956\\
2.208	-0.0245888\\
2.2106	-0.0238724\\
2.2132	-0.0231755\\
2.2157	-0.0225238\\
2.2183	-0.0218654\\
2.2209	-0.021227\\
2.2234	-0.0206321\\
2.226	-0.0200333\\
2.2286	-0.0194551\\
2.2311	-0.0189185\\
2.2337	-0.0183807\\
2.2362	-0.0178834\\
2.2388	-0.0173867\\
2.2414	-0.0169111\\
2.2439	-0.0164738\\
2.2465	-0.0160398\\
2.2491	-0.0156272\\
2.2516	-0.0152505\\
2.2542	-0.0148798\\
2.2568	-0.0145306\\
2.2593	-0.0142151\\
2.2619	-0.013908\\
2.2644	-0.013633\\
2.267	-0.0133681\\
2.2696	-0.0131246\\
2.2721	-0.0129108\\
2.2747	-0.0127095\\
2.2773	-0.0125295\\
2.2798	-0.0123766\\
2.2824	-0.0122384\\
2.285	-0.0121214\\
2.2875	-0.0120288\\
2.2901	-0.0119532\\
2.2926	-0.0119003\\
2.2952	-0.0118656\\
2.2978	-0.0118518\\
2.3003	-0.0118579\\
2.3029	-0.0118843\\
2.3055	-0.0119312\\
2.308	-0.0119953\\
2.3106	-0.0120817\\
2.3132	-0.012188\\
2.3157	-0.0123089\\
2.3183	-0.0124539\\
2.3208	-0.0126116\\
2.3234	-0.0127946\\
2.326	-0.0129967\\
2.3285	-0.0132089\\
2.3311	-0.0134479\\
2.3337	-0.0137055\\
2.3362	-0.0139704\\
2.3388	-0.0142638\\
2.3414	-0.014575\\
2.3439	-0.0148911\\
2.3465	-0.0152369\\
2.349	-0.0155857\\
2.3516	-0.0159653\\
2.3542	-0.0163617\\
2.3567	-0.0167585\\
2.3593	-0.0171873\\
2.3619	-0.0176323\\
2.3644	-0.0180752\\
2.367	-0.0185511\\
2.3696	-0.0190426\\
2.3721	-0.0195294\\
2.3747	-0.0200505\\
2.3773	-0.0205862\\
2.3798	-0.021115\\
2.3824	-0.0216789\\
2.3849	-0.0222343\\
2.3875	-0.0228254\\
2.3901	-0.02343\\
2.3926	-0.0240239\\
2.3952	-0.0246542\\
2.3978	-0.0252972\\
2.4003	-0.0259272\\
2.4029	-0.0265945\\
2.4055	-0.0272737\\
2.408	-0.0279378\\
2.4106	-0.0286397\\
2.4131	-0.0293251\\
2.4157	-0.0300488\\
2.4183	-0.030783\\
2.4208	-0.0314989\\
2.4234	-0.0322534\\
2.426	-0.0330178\\
2.4285	-0.0337619\\
2.4311	-0.034545\\
2.4337	-0.0353372\\
2.4362	-0.0361073\\
2.4388	-0.0369166\\
2.4413	-0.0377028\\
2.4439	-0.0385283\\
2.4465	-0.0393617\\
2.449	-0.0401701\\
2.4516	-0.0410181\\
2.4542	-0.0418732\\
2.4567	-0.0427018\\
2.4593	-0.04357\\
2.4619	-0.0444445\\
2.4644	-0.0452911\\
2.467	-0.0461772\\
2.4695	-0.0470344\\
2.4721	-0.0479311\\
2.4747	-0.0488328\\
2.4772	-0.0497043\\
2.4798	-0.0506151\\
2.4824	-0.0515301\\
2.4849	-0.0524136\\
2.4875	-0.0533362\\
2.4901	-0.0542621\\
2.4926	-0.0551555\\
2.4952	-0.0560875\\
2.4977	-0.0569861\\
2.5003	-0.057923\\
2.5029	-0.058862\\
2.5054	-0.0597667\\
2.508	-0.0607091\\
2.5106	-0.0616529\\
2.5131	-0.0625613\\
2.5157	-0.0635068\\
2.5183	-0.0644529\\
2.5208	-0.0653628\\
2.5234	-0.0663091\\
2.526	-0.067255\\
2.5285	-0.068164\\
2.5311	-0.0691085\\
2.5336	-0.0700156\\
2.5362	-0.0709576\\
2.5388	-0.0718979\\
2.5413	-0.0728001\\
2.5439	-0.0737362\\
2.5465	-0.0746697\\
2.549	-0.0755647\\
2.5516	-0.0764923\\
2.5542	-0.0774165\\
2.5567	-0.0783017\\
2.5593	-0.0792183\\
2.5618	-0.0800955\\
2.5644	-0.0810034\\
2.567	-0.0819063\\
2.5695	-0.0827696\\
2.5721	-0.083662\\
2.5747	-0.0845486\\
2.5772	-0.0853953\\
2.5798	-0.0862696\\
2.5824	-0.0871371\\
2.5849	-0.0879647\\
2.5875	-0.0888181\\
2.59	-0.0896315\\
2.5926	-0.0904696\\
2.5952	-0.0912994\\
2.5977	-0.0920892\\
2.6003	-0.0929018\\
2.6029	-0.0937053\\
2.6054	-0.0944688\\
2.608	-0.0952533\\
2.6106	-0.0960275\\
2.6131	-0.0967621\\
2.6157	-0.0975155\\
2.6182	-0.0982295\\
2.6208	-0.0989608\\
2.6234	-0.0996804\\
2.6259	-0.100361\\
2.6285	-0.101057\\
2.6311	-0.10174\\
2.6336	-0.102384\\
2.6362	-0.103042\\
2.6388	-0.103685\\
2.6413	-0.104291\\
2.6439	-0.104907\\
2.6464	-0.105486\\
2.649	-0.106074\\
2.6516	-0.106646\\
2.6541	-0.107182\\
2.6567	-0.107724\\
2.6593	-0.10825\\
2.6618	-0.10874\\
2.6644	-0.109234\\
2.667	-0.109711\\
2.6695	-0.110154\\
2.6721	-0.110597\\
2.6746	-0.111007\\
2.6772	-0.111416\\
2.6798	-0.111806\\
2.6823	-0.112165\\
2.6849	-0.112519\\
2.6875	-0.112855\\
2.69	-0.113159\\
2.6926	-0.113457\\
2.6952	-0.113736\\
2.6977	-0.113985\\
2.7003	-0.114225\\
2.7029	-0.114445\\
2.7054	-0.114637\\
2.708	-0.114817\\
2.7105	-0.11497\\
2.7131	-0.115109\\
2.7157	-0.115227\\
2.7182	-0.115321\\
2.7208	-0.115398\\
2.7234	-0.115453\\
2.7259	-0.115486\\
2.7285	-0.115499\\
2.7311	-0.115491\\
2.7336	-0.115462\\
2.7362	-0.115411\\
2.7387	-0.115341\\
2.7413	-0.115246\\
2.7439	-0.11513\\
2.7464	-0.114997\\
2.749	-0.114838\\
2.7516	-0.114656\\
2.7541	-0.11446\\
2.7567	-0.114235\\
2.7593	-0.113988\\
2.7618	-0.113729\\
2.7644	-0.113439\\
2.7669	-0.113139\\
2.7695	-0.112805\\
2.7721	-0.11245\\
2.7746	-0.112088\\
2.7772	-0.11169\\
2.7798	-0.11127\\
2.7823	-0.110847\\
2.7849	-0.110386\\
2.7875	-0.109903\\
2.79	-0.10942\\
2.7926	-0.108896\\
2.7951	-0.108374\\
2.7977	-0.10781\\
2.8003	-0.107226\\
2.8028	-0.106646\\
2.8054	-0.106023\\
2.808	-0.105381\\
2.8105	-0.104745\\
2.8131	-0.104065\\
2.8157	-0.103367\\
2.8182	-0.102678\\
2.8208	-0.101943\\
2.8233	-0.101221\\
2.8259	-0.100452\\
2.8285	-0.0996655\\
2.831	-0.0988937\\
2.8336	-0.098075\\
2.8362	-0.0972401\\
2.8387	-0.0964226\\
2.8413	-0.0955573\\
2.8439	-0.094677\\
2.8464	-0.0938169\\
2.849	-0.0929086\\
2.8516	-0.0919866\\
2.8541	-0.0910875\\
2.8567	-0.0901398\\
2.8592	-0.089217\\
2.8618	-0.0882455\\
2.8644	-0.0872625\\
2.8669	-0.086307\\
2.8695	-0.0853029\\
2.8721	-0.0842888\\
2.8746	-0.0833046\\
2.8772	-0.0822722\\
2.8798	-0.0812313\\
2.8823	-0.0802228\\
2.8849	-0.0791665\\
2.8874	-0.0781443\\
2.89	-0.0770748\\
2.8926	-0.0759992\\
2.8951	-0.0749599\\
2.8977	-0.0738742\\
2.9003	-0.0727841\\
2.9028	-0.0717323\\
2.9054	-0.0706352\\
2.908	-0.0695352\\
2.9105	-0.0684755\\
2.9131	-0.0673717\\
2.9156	-0.0663093\\
2.9182	-0.0652038\\
2.9208	-0.0640982\\
2.9233	-0.0630355\\
2.9259	-0.0619314\\
2.9285	-0.0608288\\
2.931	-0.0597707\\
2.9336	-0.0586727\\
2.9362	-0.057578\\
2.9387	-0.0565289\\
2.9413	-0.055442\\
2.9438	-0.0544013\\
2.9464	-0.0533243\\
2.949	-0.0522531\\
2.9515	-0.051229\\
2.9541	-0.0501707\\
2.9567	-0.0491198\\
2.9592	-0.0481167\\
2.9618	-0.0470816\\
2.9644	-0.0460555\\
2.9669	-0.0450776\\
2.9695	-0.0440702\\
2.972	-0.0431113\\
2.9746	-0.0421246\\
2.9772	-0.0411491\\
2.9797	-0.0402222\\
2.9823	-0.03927\\
2.9849	-0.0383304\\
2.9874	-0.0374393\\
2.99	-0.0365256\\
2.9926	-0.0356257\\
2.9951	-0.0347739\\
2.9977	-0.0339024\\
3.0003	-0.0330458\\
3.0028	-0.0322368\\
3.0054	-0.0314108\\
3.0079	-0.0306318\\
3.0105	-0.0298379\\
3.0131	-0.0290607\\
3.0156	-0.0283295\\
3.0182	-0.0275863\\
3.0208	-0.0268608\\
3.0233	-0.0261802\\
3.0259	-0.0254905\\
3.0285	-0.0248193\\
3.031	-0.0241918\\
3.0336	-0.023558\\
3.0361	-0.0229668\\
3.0387	-0.0223712\\
3.0413	-0.0217955\\
3.0438	-0.0212608\\
3.0464	-0.0207246\\
3.049	-0.0202089\\
3.0515	-0.0197324\\
3.0541	-0.0192573\\
3.0567	-0.0188031\\
3.0592	-0.0183863\\
3.0618	-0.0179736\\
3.0643	-0.0175969\\
3.0669	-0.0172262\\
3.0695	-0.016877\\
3.072	-0.0165616\\
3.0746	-0.0162549\\
3.0772	-0.0159699\\
3.0797	-0.0157165\\
3.0823	-0.0154745\\
3.0849	-0.0152542\\
3.0874	-0.0150632\\
3.09	-0.014886\\
3.0925	-0.0147362\\
3.0951	-0.014602\\
3.0977	-0.0144896\\
3.1002	-0.0144022\\
3.1028	-0.0143327\\
3.1054	-0.0142849\\
3.1079	-0.0142593\\
3.1105	-0.014254\\
3.1131	-0.0142701\\
3.1156	-0.0143059\\
3.1182	-0.014364\\
3.1207	-0.0144398\\
3.1233	-0.0145394\\
3.1259	-0.0146599\\
3.1284	-0.0147954\\
3.131	-0.0149565\\
3.1336	-0.0151381\\
3.1361	-0.015332\\
3.1387	-0.0155533\\
3.1413	-0.0157945\\
3.1438	-0.0160452\\
3.1464	-0.016325\\
3.1489	-0.0166123\\
3.1515	-0.0169298\\
3.1541	-0.0172663\\
3.1566	-0.0176074\\
3.1592	-0.0179802\\
3.1618	-0.0183712\\
3.1643	-0.018764\\
3.1669	-0.0191899\\
3.1695	-0.0196331\\
3.172	-0.0200754\\
3.1746	-0.0205519\\
3.1772	-0.021045\\
3.1797	-0.0215344\\
3.1823	-0.022059\\
3.1848	-0.0225782\\
3.1874	-0.0231332\\
3.19	-0.0237032\\
3.1925	-0.0242651\\
3.1951	-0.0248636\\
3.1977	-0.0254761\\
3.2002	-0.026078\\
3.2028	-0.0267171\\
3.2054	-0.0273691\\
3.2079	-0.0280081\\
3.2105	-0.0286847\\
3.213	-0.0293465\\
3.2156	-0.0300462\\
3.2182	-0.0307572\\
3.2207	-0.0314511\\
3.2233	-0.032183\\
3.2259	-0.0329251\\
3.2284	-0.0336479\\
3.231	-0.0344088\\
3.2336	-0.0351787\\
3.2361	-0.0359272\\
3.2387	-0.0367137\\
3.2412	-0.0374773\\
3.2438	-0.0382789\\
3.2464	-0.0390877\\
3.2489	-0.0398716\\
3.2515	-0.0406931\\
3.2541	-0.0415206\\
3.2566	-0.0423215\\
3.2592	-0.0431595\\
3.2618	-0.0440023\\
3.2643	-0.0448168\\
3.2669	-0.0456679\\
3.2694	-0.0464896\\
3.272	-0.0473473\\
3.2746	-0.0482079\\
3.2771	-0.0490377\\
3.2797	-0.0499026\\
3.2823	-0.0507693\\
3.2848	-0.0516038\\
3.2874	-0.0524726\\
3.29	-0.0533419\\
3.2925	-0.0541779\\
3.2951	-0.055047\\
3.2976	-0.0558821\\
3.3002	-0.0567497\\
3.3028	-0.0576158\\
3.3053	-0.058447\\
3.3079	-0.0593093\\
3.3105	-0.060169\\
3.313	-0.060993\\
3.3156	-0.0618468\\
3.3182	-0.0626969\\
3.3207	-0.0635106\\
3.3233	-0.0643526\\
3.3259	-0.0651899\\
3.3284	-0.0659902\\
3.331	-0.0668172\\
3.3335	-0.067607\\
3.3361	-0.0684225\\
3.3387	-0.0692315\\
3.3412	-0.070003\\
3.3438	-0.0707984\\
3.3464	-0.0715865\\
3.3489	-0.0723368\\
3.3515	-0.0731093\\
3.3541	-0.0738734\\
3.3566	-0.0745999\\
3.3592	-0.0753467\\
3.3617	-0.0760559\\
3.3643	-0.0767841\\
3.3669	-0.0775024\\
3.3694	-0.0781834\\
3.372	-0.0788814\\
3.3746	-0.0795686\\
3.3771	-0.080219\\
3.3797	-0.0808843\\
3.3823	-0.0815381\\
3.3848	-0.0821555\\
3.3874	-0.0827858\\
3.3899	-0.0833802\\
3.3925	-0.083986\\
3.3951	-0.084579\\
3.3976	-0.0851369\\
3.4002	-0.0857041\\
3.4028	-0.0862578\\
3.4053	-0.0867773\\
3.4079	-0.0873038\\
3.4105	-0.0878163\\
3.413	-0.0882955\\
3.4156	-0.0887796\\
3.4181	-0.0892312\\
3.4207	-0.0896863\\
3.4233	-0.0901264\\
3.4258	-0.0905351\\
3.4284	-0.0909452\\
3.431	-0.0913396\\
3.4335	-0.0917042\\
3.4361	-0.0920677\\
3.4387	-0.0924154\\
3.4412	-0.0927344\\
3.4438	-0.0930504\\
3.4463	-0.0933388\\
3.4489	-0.0936226\\
3.4515	-0.0938899\\
3.454	-0.0941312\\
3.4566	-0.0943659\\
3.4592	-0.0945837\\
3.4617	-0.0947773\\
3.4643	-0.094962\\
3.4669	-0.0951297\\
3.4694	-0.095275\\
3.472	-0.0954093\\
3.4745	-0.0955223\\
3.4771	-0.095623\\
3.4797	-0.0957066\\
3.4822	-0.0957708\\
3.4848	-0.0958207\\
3.4874	-0.0958534\\
3.4899	-0.0958686\\
3.4925	-0.0958677\\
3.4951	-0.0958496\\
3.4976	-0.0958161\\
3.5002	-0.0957645\\
3.5028	-0.0956959\\
3.5053	-0.0956139\\
3.5079	-0.0955121\\
3.5104	-0.0953983\\
3.513	-0.0952634\\
3.5156	-0.0951119\\
3.5181	-0.0949505\\
3.5207	-0.0947665\\
3.5233	-0.094566\\
3.5258	-0.0943579\\
3.5284	-0.0941256\\
3.531	-0.0938772\\
3.5335	-0.0936232\\
3.5361	-0.0933436\\
3.5386	-0.09306\\
3.5412	-0.0927498\\
3.5438	-0.0924242\\
3.5463	-0.0920968\\
3.5489	-0.0917415\\
3.5515	-0.0913713\\
3.554	-0.0910015\\
3.5566	-0.0906027\\
3.5592	-0.0901895\\
3.5617	-0.0897789\\
3.5643	-0.0893382\\
3.5668	-0.0889015\\
3.5694	-0.0884341\\
3.572	-0.0879534\\
3.5745	-0.0874788\\
3.5771	-0.0869727\\
3.5797	-0.0864539\\
3.5822	-0.0859435\\
3.5848	-0.0854007\\
3.5874	-0.0848461\\
3.5899	-0.0843019\\
3.5925	-0.0837247\\
3.595	-0.0831593\\
3.5976	-0.0825606\\
3.6002	-0.0819514\\
3.6027	-0.0813559\\
3.6053	-0.0807268\\
3.6079	-0.080088\\
3.6104	-0.0794649\\
3.613	-0.0788078\\
3.6156	-0.0781419\\
3.6181	-0.0774936\\
3.6207	-0.0768113\\
3.6232	-0.0761477\\
3.6258	-0.0754501\\
3.6284	-0.0747451\\
3.6309	-0.0740606\\
3.6335	-0.0733422\\
3.6361	-0.0726174\\
3.6386	-0.0719148\\
3.6412	-0.0711784\\
3.6438	-0.0704365\\
3.6463	-0.0697184\\
3.6489	-0.0689669\\
3.6515	-0.0682109\\
3.654	-0.0674801\\
3.6566	-0.0667164\\
3.6591	-0.0659787\\
3.6617	-0.0652085\\
3.6643	-0.0644355\\
3.6668	-0.0636898\\
3.6694	-0.0629122\\
3.672	-0.0621328\\
3.6745	-0.0613819\\
3.6771	-0.0605999\\
3.6797	-0.059817\\
3.6822	-0.0590637\\
3.6848	-0.0582802\\
3.6873	-0.0575269\\
3.6899	-0.056744\\
3.6925	-0.0559618\\
3.695	-0.0552108\\
3.6976	-0.0544311\\
3.7002	-0.0536532\\
3.7027	-0.0529072\\
3.7053	-0.0521337\\
3.7079	-0.0513629\\
3.7104	-0.0506246\\
3.713	-0.04986\\
3.7155	-0.0491282\\
3.7181	-0.048371\\
3.7207	-0.0476181\\
3.7232	-0.0468984\\
3.7258	-0.0461546\\
3.7284	-0.045416\\
3.7309	-0.0447109\\
3.7335	-0.0439831\\
3.7361	-0.0432613\\
3.7386	-0.0425732\\
3.7412	-0.041864\\
3.7437	-0.0411885\\
3.7463	-0.0404928\\
3.7489	-0.0398045\\
3.7514	-0.0391498\\
3.754	-0.0384767\\
3.7566	-0.0378116\\
3.7591	-0.03718\\
3.7617	-0.0365315\\
3.7643	-0.0358918\\
3.7668	-0.0352853\\
3.7694	-0.0346636\\
3.7719	-0.0340748\\
3.7745	-0.033472\\
3.7771	-0.0328791\\
3.7796	-0.0323185\\
3.7822	-0.0317457\\
3.7848	-0.0311834\\
3.7873	-0.0306529\\
3.7899	-0.0301119\\
3.7925	-0.0295819\\
3.795	-0.029083\\
3.7976	-0.0285754\\
3.8002	-0.0280793\\
3.8027	-0.0276135\\
3.8053	-0.0271407\\
3.8078	-0.0266976\\
3.8104	-0.0262486\\
3.813	-0.0258121\\
3.8155	-0.0254041\\
3.8181	-0.0249922\\
3.8207	-0.0245931\\
3.8232	-0.0242214\\
3.8258	-0.0238476\\
3.8284	-0.0234869\\
3.8309	-0.0231525\\
3.8335	-0.0228178\\
3.836	-0.0225085\\
3.8386	-0.0222001\\
3.8412	-0.0219052\\
3.8437	-0.0216345\\
3.8463	-0.0213663\\
3.8489	-0.0211119\\
3.8514	-0.0208803\\
3.854	-0.0206529\\
3.8566	-0.0204395\\
3.8591	-0.0202474\\
3.8617	-0.0200613\\
3.8642	-0.0198955\\
3.8668	-0.0197368\\
3.8694	-0.0195922\\
3.8719	-0.0194664\\
3.8745	-0.0193493\\
3.8771	-0.0192462\\
3.8796	-0.0191603\\
3.8822	-0.0190848\\
3.8848	-0.0190233\\
3.8873	-0.0189773\\
3.8899	-0.0189433\\
3.8924	-0.0189236\\
3.895	-0.0189169\\
3.8976	-0.018924\\
3.9001	-0.0189438\\
3.9027	-0.018978\\
3.9053	-0.0190259\\
3.9078	-0.0190849\\
3.9104	-0.0191595\\
3.913	-0.0192477\\
3.9155	-0.0193451\\
3.9181	-0.0194597\\
3.9206	-0.0195823\\
3.9232	-0.0197228\\
3.9258	-0.0198765\\
3.9283	-0.0200365\\
3.9309	-0.0202156\\
3.9335	-0.0204076\\
3.936	-0.0206041\\
3.9386	-0.0208209\\
3.9412	-0.0210502\\
3.9437	-0.0212824\\
3.9463	-0.0215359\\
3.9488	-0.021791\\
3.9514	-0.0220682\\
3.954	-0.0223572\\
3.9565	-0.0226461\\
3.9591	-0.022958\\
3.9617	-0.0232813\\
3.9642	-0.0236029\\
3.9668	-0.0239482\\
3.9694	-0.0243046\\
3.9719	-0.0246574\\
3.9745	-0.0250348\\
3.9771	-0.0254227\\
3.9796	-0.0258055\\
3.9822	-0.0262135\\
3.9847	-0.0266153\\
3.9873	-0.0270427\\
3.9899	-0.0274798\\
3.9924	-0.0279091\\
3.995	-0.0283645\\
3.9976	-0.0288291\\
4.0001	-0.0292842\\
4.0027	-0.029766\\
4.0053	-0.0302563\\
4.0078	-0.0307356\\
4.0104	-0.031242\\
4.0129	-0.0317365\\
4.0155	-0.0322582\\
4.0181	-0.0327874\\
4.0206	-0.0333032\\
4.0232	-0.0338465\\
4.0258	-0.0343967\\
4.0283	-0.0349321\\
4.0309	-0.0354951\\
4.0335	-0.0360643\\
4.036	-0.0366174\\
4.0386	-0.0371982\\
4.0411	-0.0377619\\
4.0437	-0.0383534\\
4.0463	-0.03895\\
4.0488	-0.0395282\\
4.0514	-0.0401341\\
4.054	-0.0407445\\
4.0565	-0.0413353\\
4.0591	-0.0419537\\
4.0617	-0.0425757\\
4.0642	-0.0431771\\
4.0668	-0.0438058\\
4.0693	-0.0444132\\
4.0719	-0.0450475\\
4.0745	-0.0456844\\
4.077	-0.046299\\
4.0796	-0.0469402\\
4.0822	-0.0475832\\
4.0847	-0.048203\\
4.0873	-0.0488488\\
4.0899	-0.0494958\\
4.0924	-0.0501187\\
4.095	-0.0507672\\
4.0975	-0.051391\\
4.1001	-0.05204\\
4.1027	-0.0526888\\
4.1052	-0.0533124\\
4.1078	-0.0539603\\
4.1104	-0.0546074\\
4.1129	-0.0552286\\
4.1155	-0.0558734\\
4.1181	-0.0565166\\
4.1206	-0.0571334\\
4.1232	-0.0577729\\
4.1258	-0.0584101\\
4.1283	-0.0590204\\
4.1309	-0.0596525\\
4.1334	-0.0602575\\
4.136	-0.0608835\\
4.1386	-0.061506\\
4.1411	-0.0621012\\
4.1437	-0.0627163\\
4.1463	-0.0633273\\
4.1488	-0.0639107\\
4.1514	-0.0645129\\
4.154	-0.0651102\\
4.1565	-0.0656799\\
4.1591	-0.0662672\\
4.1616	-0.0668267\\
4.1642	-0.067403\\
4.1668	-0.0679734\\
4.1693	-0.0685161\\
4.1719	-0.0690743\\
4.1745	-0.0696259\\
4.177	-0.0701499\\
4.1796	-0.0706881\\
4.1822	-0.0712191\\
4.1847	-0.0717227\\
4.1873	-0.0722391\\
4.1898	-0.0727283\\
4.1924	-0.0732292\\
4.195	-0.073722\\
4.1975	-0.0741881\\
4.2001	-0.0746644\\
4.2027	-0.0751321\\
4.2052	-0.0755734\\
4.2078	-0.0760235\\
4.2104	-0.0764644\\
4.2129	-0.0768795\\
4.2155	-0.0773019\\
4.218	-0.077699\\
4.2206	-0.0781022\\
4.2232	-0.0784955\\
4.2257	-0.0788641\\
4.2283	-0.0792374\\
4.2309	-0.0796003\\
4.2334	-0.0799393\\
4.236	-0.0802813\\
4.2386	-0.0806126\\
4.2411	-0.0809209\\
4.2437	-0.0812306\\
4.2462	-0.081518\\
4.2488	-0.0818058\\
4.2514	-0.0820823\\
4.2539	-0.0823373\\
4.2565	-0.0825912\\
4.2591	-0.0828335\\
4.2616	-0.0830555\\
4.2642	-0.0832747\\
4.2668	-0.0834821\\
4.2693	-0.0836703\\
4.2719	-0.0838543\\
4.2744	-0.0840198\\
4.277	-0.0841801\\
4.2796	-0.0843283\\
4.2821	-0.0844593\\
4.2847	-0.0845836\\
4.2873	-0.0846956\\
4.2898	-0.0847917\\
4.2924	-0.0848796\\
4.295	-0.0849553\\
4.2975	-0.0850164\\
4.3001	-0.0850679\\
4.3027	-0.0851071\\
4.3052	-0.0851332\\
4.3078	-0.0851483\\
4.3103	-0.0851513\\
4.3129	-0.0851424\\
4.3155	-0.0851213\\
4.318	-0.0850895\\
4.3206	-0.0850446\\
4.3232	-0.0849876\\
4.3257	-0.0849215\\
4.3283	-0.084841\\
4.3309	-0.0847486\\
4.3334	-0.0846486\\
4.336	-0.0845331\\
4.3385	-0.084411\\
4.3411	-0.0842727\\
4.3437	-0.0841229\\
4.3462	-0.0839682\\
4.3488	-0.0837962\\
4.3514	-0.083613\\
4.3539	-0.0834264\\
4.3565	-0.0832217\\
4.3591	-0.0830061\\
4.3616	-0.0827887\\
4.3642	-0.0825523\\
4.3667	-0.0823152\\
4.3693	-0.0820585\\
4.3719	-0.0817917\\
4.3744	-0.0815258\\
4.377	-0.0812396\\
4.3796	-0.0809437\\
4.3821	-0.0806503\\
4.3847	-0.080336\\
4.3873	-0.0800126\\
4.3898	-0.0796932\\
4.3924	-0.0793524\\
4.3949	-0.0790166\\
4.3975	-0.0786592\\
4.4001	-0.0782935\\
4.4026	-0.0779344\\
4.4052	-0.0775533\\
4.4078	-0.0771646\\
4.4103	-0.0767838\\
4.4129	-0.0763809\\
4.4155	-0.0759709\\
4.418	-0.0755704\\
4.4206	-0.0751475\\
4.4231	-0.0747349\\
4.4257	-0.0742998\\
4.4283	-0.073859\\
4.4308	-0.0734298\\
4.4334	-0.0729781\\
4.436	-0.0725214\\
4.4385	-0.0720775\\
4.4411	-0.0716114\\
4.4437	-0.0711408\\
4.4462	-0.0706845\\
4.4488	-0.070206\\
4.4514	-0.0697238\\
4.4539	-0.0692569\\
4.4565	-0.0687682\\
4.459	-0.0682956\\
4.4616	-0.0678014\\
4.4642	-0.0673047\\
4.4667	-0.0668251\\
4.4693	-0.0663244\\
4.4719	-0.065822\\
4.4744	-0.0653375\\
4.477	-0.0648325\\
4.4796	-0.0643265\\
4.4821	-0.0638394\\
4.4847	-0.0633323\\
4.4872	-0.0628445\\
4.4898	-0.0623373\\
4.4924	-0.0618304\\
4.4949	-0.0613434\\
4.4975	-0.0608378\\
4.5001	-0.0603332\\
4.5026	-0.0598491\\
4.5052	-0.0593472\\
4.5078	-0.0588471\\
4.5103	-0.058368\\
4.5129	-0.0578719\\
4.5154	-0.0573973\\
4.518	-0.0569062\\
4.5206	-0.0564181\\
4.5231	-0.0559517\\
4.5257	-0.0554699\\
4.5283	-0.0549916\\
4.5308	-0.0545354\\
4.5334	-0.0540648\\
4.536	-0.0535984\\
4.5385	-0.0531541\\
4.5411	-0.0526966\\
4.5436	-0.0522613\\
4.5462	-0.0518134\\
4.5488	-0.0513708\\
4.5513	-0.0509503\\
4.5539	-0.0505185\\
4.5565	-0.0500925\\
4.559	-0.0496885\\
4.5616	-0.0492743\\
4.5642	-0.0488665\\
4.5667	-0.0484805\\
4.5693	-0.0480855\\
4.5718	-0.0477122\\
4.5744	-0.0473308\\
4.577	-0.0469565\\
4.5795	-0.0466035\\
4.5821	-0.0462436\\
4.5847	-0.0458913\\
4.5872	-0.0455598\\
4.5898	-0.0452227\\
4.5924	-0.0448936\\
4.5949	-0.0445848\\
4.5975	-0.0442716\\
4.6001	-0.0439668\\
4.6026	-0.0436816\\
4.6052	-0.0433934\\
4.6077	-0.0431244\\
4.6103	-0.0428532\\
4.6129	-0.0425908\\
4.6154	-0.0423469\\
4.618	-0.042102\\
4.6206	-0.0418662\\
4.6231	-0.0416481\\
4.6257	-0.0414303\\
4.6283	-0.0412217\\
4.6308	-0.0410299\\
4.6334	-0.0408397\\
4.6359	-0.0406656\\
4.6385	-0.0404939\\
4.6411	-0.0403317\\
4.6436	-0.0401847\\
4.6462	-0.0400413\\
4.6488	-0.0399074\\
4.6513	-0.0397877\\
4.6539	-0.0396727\\
4.6565	-0.0395673\\
4.659	-0.0394751\\
4.6616	-0.0393887\\
4.6641	-0.0393147\\
4.6667	-0.0392472\\
4.6693	-0.0391893\\
4.6718	-0.0391428\\
4.6744	-0.0391038\\
4.677	-0.0390743\\
4.6795	-0.039055\\
4.6821	-0.0390443\\
4.6847	-0.0390431\\
4.6872	-0.0390508\\
4.6898	-0.039068\\
4.6923	-0.0390934\\
4.6949	-0.0391289\\
4.6975	-0.0391737\\
4.7	-0.0392254\\
4.7026	-0.0392882\\
4.7052	-0.0393599\\
4.7077	-0.0394374\\
4.7103	-0.0395268\\
4.7129	-0.039625\\
4.7154	-0.0397276\\
4.718	-0.0398428\\
4.7205	-0.0399618\\
4.7231	-0.0400937\\
4.7257	-0.0402341\\
4.7282	-0.0403769\\
4.7308	-0.0405335\\
4.7334	-0.0406981\\
4.7359	-0.0408639\\
4.7385	-0.0410441\\
4.7411	-0.041232\\
4.7436	-0.0414199\\
4.7462	-0.0416228\\
4.7487	-0.0418247\\
4.7513	-0.0420419\\
4.7539	-0.0422663\\
4.7564	-0.0424886\\
4.759	-0.0427266\\
4.7616	-0.0429714\\
4.7641	-0.043213\\
4.7667	-0.0434706\\
4.7693	-0.0437345\\
4.7718	-0.0439942\\
4.7744	-0.0442701\\
4.777	-0.044552\\
4.7795	-0.0448285\\
4.7821	-0.0451215\\
4.7846	-0.0454084\\
4.7872	-0.045712\\
4.7898	-0.0460207\\
4.7923	-0.0463223\\
4.7949	-0.0466406\\
4.7975	-0.0469637\\
4.8	-0.0472786\\
4.8026	-0.0476103\\
4.8052	-0.0479462\\
4.8077	-0.048273\\
4.8103	-0.0486167\\
4.8128	-0.0489506\\
4.8154	-0.0493014\\
4.818	-0.0496555\\
4.8205	-0.049999\\
4.8231	-0.0503592\\
4.8257	-0.0507223\\
4.8282	-0.0510739\\
4.8308	-0.051442\\
4.8334	-0.0518125\\
4.8359	-0.0521708\\
4.8385	-0.0525453\\
4.841	-0.0529072\\
4.8436	-0.0532851\\
4.8462	-0.0536645\\
4.8487	-0.0540306\\
4.8513	-0.0544124\\
4.8539	-0.0547951\\
4.8564	-0.0551639\\
4.859	-0.055548\\
4.8616	-0.0559326\\
4.8641	-0.0563025\\
4.8667	-0.0566874\\
4.8692	-0.0570574\\
4.8718	-0.057442\\
4.8744	-0.0578261\\
4.8769	-0.0581949\\
4.8795	-0.0585777\\
4.8821	-0.0589596\\
4.8846	-0.0593257\\
4.8872	-0.0597053\\
4.8898	-0.0600834\\
4.8923	-0.0604454\\
4.8949	-0.0608202\\
4.8974	-0.0611788\\
4.9	-0.0615497\\
4.9026	-0.0619183\\
4.9051	-0.0622705\\
4.9077	-0.0626342\\
4.9103	-0.0629952\\
4.9128	-0.0633396\\
4.9154	-0.0636948\\
4.918	-0.0640469\\
4.9205	-0.0643822\\
4.9231	-0.0647275\\
4.9257	-0.0650692\\
4.9282	-0.0653941\\
4.9308	-0.0657282\\
4.9333	-0.0660457\\
4.9359	-0.0663716\\
4.9385	-0.0666933\\
4.941	-0.0669983\\
4.9436	-0.067311\\
4.9462	-0.0676189\\
4.9487	-0.0679104\\
4.9513	-0.0682086\\
4.9539	-0.0685017\\
4.9564	-0.0687785\\
4.959	-0.0690611\\
4.9615	-0.0693276\\
4.9641	-0.0695993\\
4.9667	-0.0698653\\
4.9692	-0.0701155\\
4.9718	-0.0703698\\
4.9744	-0.0706181\\
4.9769	-0.070851\\
4.9795	-0.071087\\
4.9821	-0.0713167\\
4.9846	-0.0715314\\
4.9872	-0.0717483\\
4.9897	-0.0719505\\
4.9923	-0.0721542\\
4.9949	-0.0723511\\
4.9974	-0.0725338\\
5	-0.072717\\
};
\addplot [color=mycolor9, line width=1.0pt, forget plot]
  table[row sep=crcr]{%
-0.125	5.17929e-08\\
-0.1224	5.28198e-08\\
-0.1199	5.38063e-08\\
-0.1173	5.48312e-08\\
-0.1147	5.58549e-08\\
-0.1122	5.68384e-08\\
-0.1096	5.786e-08\\
-0.1071	5.88413e-08\\
-0.1045	5.98608e-08\\
-0.1019	6.08792e-08\\
-0.0994	6.18574e-08\\
-0.0968	6.28738e-08\\
-0.0942	6.38892e-08\\
-0.0917	6.48644e-08\\
-0.0891	6.58777e-08\\
-0.0865	6.689e-08\\
-0.084	6.78624e-08\\
-0.0814	6.88727e-08\\
-0.0789	6.98432e-08\\
-0.0763	7.08516e-08\\
-0.0737	7.18591e-08\\
-0.0712	7.28269e-08\\
-0.0686	7.38327e-08\\
-0.066	7.48375e-08\\
-0.0635	7.58028e-08\\
-0.0609	7.68059e-08\\
-0.0583	7.78081e-08\\
-0.0558	7.8771e-08\\
-0.0532	7.97716e-08\\
-0.0507	8.0733e-08\\
-0.0481	8.17319e-08\\
-0.0455	8.273e-08\\
-0.043	8.36891e-08\\
-0.0404	8.46858e-08\\
-0.0378	8.56817e-08\\
-0.0353	8.66388e-08\\
-0.0327	8.76331e-08\\
-0.0301	8.86269e-08\\
-0.0276	8.95817e-08\\
-0.025	9.0574e-08\\
-0.0224	9.15657e-08\\
-0.0199	9.25187e-08\\
-0.0173	9.3509e-08\\
-0.0148	9.44607e-08\\
-0.0122	9.54499e-08\\
-0.0096	9.64384e-08\\
-0.0071	9.73884e-08\\
-0.0045	9.83759e-08\\
-0.0019	9.93628e-08\\
0.0006	1.00311e-07\\
0.0032	0.0119509\\
0.0058	0.0191745\\
0.0083	0.0241741\\
0.0109	0.0285831\\
0.0134	0.0323678\\
0.016	0.0357272\\
0.0186	0.0383793\\
0.0211	0.0403484\\
0.0237	0.0420328\\
0.0263	0.0435552\\
0.0288	0.0449326\\
0.0314	0.0462285\\
0.034	0.047298\\
0.0365	0.0480814\\
0.0391	0.0486836\\
0.0416	0.0491176\\
0.0442	0.0494656\\
0.0468	0.0497103\\
0.0493	0.0498216\\
0.0519	0.0497869\\
0.0545	0.0496023\\
0.057	0.0493083\\
0.0596	0.0489125\\
0.0622	0.0484414\\
0.0647	0.0479123\\
0.0673	0.0472695\\
0.0698	0.046559\\
0.0724	0.0457335\\
0.075	0.0448439\\
0.0775	0.0439468\\
0.0801	0.0429736\\
0.0827	0.0419524\\
0.0852	0.040918\\
0.0878	0.0397889\\
0.0904	0.038621\\
0.0929	0.0374799\\
0.0955	0.0362831\\
0.098	0.0351204\\
0.1006	0.0338899\\
0.1032	0.0326327\\
0.1057	0.0314056\\
0.1083	0.0301258\\
0.1109	0.028855\\
0.1134	0.0276432\\
0.116	0.0263862\\
0.1186	0.0251233\\
0.1211	0.0239033\\
0.1237	0.0226389\\
0.1263	0.0213916\\
0.1288	0.0202139\\
0.1314	0.0190075\\
0.1339	0.0178554\\
0.1365	0.01666\\
0.1391	0.015471\\
0.1416	0.0143437\\
0.1442	0.0131968\\
0.1468	0.0120753\\
0.1493	0.0110128\\
0.1519	0.00991538\\
0.1545	0.00882383\\
0.157	0.00778605\\
0.1596	0.00672817\\
0.1621	0.00573506\\
0.1647	0.00472208\\
0.1673	0.00371974\\
0.1698	0.00276016\\
0.1724	0.00176846\\
0.175	0.000791151\\
0.1775	-0.00012886\\
0.1801	-0.00106629\\
0.1827	-0.00199178\\
0.1852	-0.00287821\\
0.1878	-0.00379875\\
0.1903	-0.00467772\\
0.1929	-0.00557818\\
0.1955	-0.00646185\\
0.198	-0.00729989\\
0.2006	-0.00816744\\
0.2032	-0.00903623\\
0.2057	-0.00987123\\
0.2083	-0.0107327\\
0.2109	-0.0115816\\
0.2134	-0.0123867\\
0.216	-0.0132186\\
0.2185	-0.0140196\\
0.2211	-0.0148555\\
0.2237	-0.0156898\\
0.2262	-0.0164843\\
0.2288	-0.0172998\\
0.2314	-0.0181077\\
0.2339	-0.0188835\\
0.2365	-0.0196932\\
0.2391	-0.0205042\\
0.2416	-0.0212796\\
0.2442	-0.0220764\\
0.2467	-0.0228334\\
0.2493	-0.0236159\\
0.2519	-0.0243988\\
0.2544	-0.025153\\
0.257	-0.0259346\\
0.2596	-0.0267077\\
0.2621	-0.0274403\\
0.2647	-0.0281935\\
0.2673	-0.0289428\\
0.2698	-0.0296625\\
0.2724	-0.0304084\\
0.2749	-0.0311183\\
0.2775	-0.0318445\\
0.2801	-0.0325582\\
0.2826	-0.0332359\\
0.2852	-0.0339361\\
0.2878	-0.0346322\\
0.2903	-0.0352942\\
0.2929	-0.0359702\\
0.2955	-0.0366309\\
0.298	-0.0372534\\
0.3006	-0.0378916\\
0.3032	-0.0385227\\
0.3057	-0.0391216\\
0.3083	-0.0397316\\
0.3108	-0.0403024\\
0.3134	-0.0408793\\
0.316	-0.0414419\\
0.3185	-0.0419729\\
0.3211	-0.0425151\\
0.3237	-0.0430442\\
0.3262	-0.0435366\\
0.3288	-0.0440296\\
0.3314	-0.0445046\\
0.3339	-0.0449475\\
0.3365	-0.0453959\\
0.339	-0.045814\\
0.3416	-0.0462319\\
0.3442	-0.0466292\\
0.3467	-0.0469915\\
0.3493	-0.0473504\\
0.3519	-0.0476937\\
0.3544	-0.0480098\\
0.357	-0.0483215\\
0.3596	-0.0486123\\
0.3621	-0.0488707\\
0.3647	-0.0491184\\
0.3672	-0.049339\\
0.3698	-0.049552\\
0.3724	-0.0497472\\
0.3749	-0.0499154\\
0.3775	-0.0500674\\
0.3801	-0.0501957\\
0.3826	-0.0502987\\
0.3852	-0.050387\\
0.3878	-0.0504565\\
0.3903	-0.0505038\\
0.3929	-0.0505298\\
0.3954	-0.0505311\\
0.398	-0.0505087\\
0.4006	-0.0504644\\
0.4031	-0.0504025\\
0.4057	-0.0503173\\
0.4083	-0.0502086\\
0.4108	-0.0500795\\
0.4134	-0.0499196\\
0.416	-0.0497351\\
0.4185	-0.0495364\\
0.4211	-0.0493079\\
0.4236	-0.0490659\\
0.4262	-0.0487884\\
0.4288	-0.0484833\\
0.4313	-0.0481647\\
0.4339	-0.0478088\\
0.4365	-0.0474293\\
0.439	-0.0470414\\
0.4416	-0.0466123\\
0.4442	-0.0461549\\
0.4467	-0.0456882\\
0.4493	-0.0451764\\
0.4519	-0.0446391\\
0.4544	-0.0440989\\
0.457	-0.0435113\\
0.4595	-0.0429197\\
0.4621	-0.0422757\\
0.4647	-0.041603\\
0.4672	-0.0409307\\
0.4698	-0.0402062\\
0.4724	-0.0394555\\
0.4749	-0.0387075\\
0.4775	-0.0379009\\
0.4801	-0.0370649\\
0.4826	-0.0362347\\
0.4852	-0.0353452\\
0.4877	-0.0344655\\
0.4903	-0.0335242\\
0.4929	-0.0325546\\
0.4954	-0.031595\\
0.498	-0.0305693\\
0.5006	-0.0295168\\
0.5031	-0.0284808\\
0.5057	-0.0273783\\
0.5083	-0.0262493\\
0.5108	-0.0251376\\
0.5134	-0.0239546\\
0.5159	-0.0227924\\
0.5185	-0.0215598\\
0.5211	-0.0203033\\
0.5236	-0.0190723\\
0.5262	-0.0177675\\
0.5288	-0.0164373\\
0.5313	-0.0151351\\
0.5339	-0.0137584\\
0.5365	-0.0123602\\
0.539	-0.0109957\\
0.5416	-0.00955531\\
0.5441	-0.00814931\\
0.5467	-0.00666566\\
0.5493	-0.00516164\\
0.5518	-0.00369782\\
0.5544	-0.00215804\\
0.557	-0.00060052\\
0.5595	0.000914218\\
0.5621	0.00250724\\
0.5647	0.00411722\\
0.5672	0.00567964\\
0.5698	0.00731807\\
0.5723	0.00890592\\
0.5749	0.0105703\\
0.5775	0.012248\\
0.58	0.0138732\\
0.5826	0.0155742\\
0.5852	0.0172846\\
0.5877	0.018937\\
0.5903	0.0206633\\
0.5929	0.0223974\\
0.5954	0.024072\\
0.598	0.0258198\\
0.6006	0.0275721\\
0.6031	0.0292599\\
0.6057	0.0310174\\
0.6082	0.0327092\\
0.6108	0.0344706\\
0.6134	0.0362328\\
0.6159	0.0379268\\
0.6185	0.0396863\\
0.6211	0.0414422\\
0.6236	0.0431269\\
0.6262	0.0448748\\
0.6288	0.0466179\\
0.6313	0.0482882\\
0.6339	0.0500179\\
0.6364	0.0516723\\
0.639	0.0533831\\
0.6416	0.0550835\\
0.6441	0.0567084\\
0.6467	0.0583869\\
0.6493	0.0600522\\
0.6518	0.0616397\\
0.6544	0.0632751\\
0.657	0.0648942\\
0.6595	0.0664352\\
0.6621	0.0680209\\
0.6646	0.0695283\\
0.6672	0.0710764\\
0.6698	0.0726033\\
0.6723	0.0740506\\
0.6749	0.0755339\\
0.6775	0.0769944\\
0.68	0.0783764\\
0.6826	0.0797894\\
0.6852	0.0811761\\
0.6877	0.0824839\\
0.6903	0.0838167\\
0.6928	0.085072\\
0.6954	0.0863495\\
0.698	0.0875975\\
0.7005	0.0887687\\
0.7031	0.0899555\\
0.7057	0.0911098\\
0.7082	0.092189\\
0.7108	0.0932789\\
0.7134	0.0943353\\
0.7159	0.0953183\\
0.7185	0.0963056\\
0.721	0.0972206\\
0.7236	0.098136\\
0.7262	0.0990145\\
0.7287	0.099824\\
0.7313	0.100629\\
0.7339	0.101394\\
0.7364	0.102093\\
0.739	0.10278\\
0.7416	0.103427\\
0.7441	0.10401\\
0.7467	0.104577\\
0.7492	0.105083\\
0.7518	0.105567\\
0.7544	0.106009\\
0.7569	0.106393\\
0.7595	0.10675\\
0.7621	0.107064\\
0.7646	0.107325\\
0.7672	0.107553\\
0.7698	0.107736\\
0.7723	0.107869\\
0.7749	0.107964\\
0.7775	0.108014\\
0.78	0.10802\\
0.7826	0.107981\\
0.7851	0.107901\\
0.7877	0.107772\\
0.7903	0.107596\\
0.7928	0.107385\\
0.7954	0.10712\\
0.798	0.106809\\
0.8005	0.106466\\
0.8031	0.106064\\
0.8057	0.105616\\
0.8082	0.105141\\
0.8108	0.104602\\
0.8133	0.104041\\
0.8159	0.103412\\
0.8185	0.102737\\
0.821	0.102045\\
0.8236	0.101279\\
0.8262	0.100469\\
0.8287	0.0996467\\
0.8313	0.0987476\\
0.8339	0.0978034\\
0.8364	0.096853\\
0.839	0.0958208\\
0.8415	0.0947865\\
0.8441	0.0936678\\
0.8467	0.0925057\\
0.8492	0.0913474\\
0.8518	0.0901003\\
0.8544	0.0888103\\
0.8569	0.0875298\\
0.8595	0.0861571\\
0.8621	0.0847431\\
0.8646	0.0833448\\
0.8672	0.0818505\\
0.8697	0.0803756\\
0.8723	0.0788024\\
0.8749	0.07719\\
0.8774	0.0756032\\
0.88	0.0739156\\
0.8826	0.0721902\\
0.8851	0.0704961\\
0.8877	0.068698\\
0.8903	0.0668639\\
0.8928	0.065067\\
0.8954	0.0631642\\
0.8979	0.0613026\\
0.9005	0.0593336\\
0.9031	0.0573315\\
0.9056	0.0553759\\
0.9082	0.0533113\\
0.9108	0.051216\\
0.9133	0.0491732\\
0.9159	0.04702\\
0.9185	0.044838\\
0.921	0.0427137\\
0.9236	0.0404778\\
0.9262	0.0382159\\
0.9287	0.0360171\\
0.9313	0.0337065\\
0.9338	0.0314622\\
0.9364	0.0291057\\
0.939	0.0267269\\
0.9415	0.0244198\\
0.9441	0.0220006\\
0.9467	0.0195622\\
0.9492	0.0172003\\
0.9518	0.0147266\\
0.9544	0.0122361\\
0.9569	0.00982668\\
0.9595	0.00730654\\
0.962	0.00487052\\
0.9646	0.00232465\\
0.9672	-0.000233029\\
0.9697	-0.00270261\\
0.9723	-0.00528064\\
0.9749	-0.0078674\\
0.9774	-0.0103618\\
0.98	-0.0129625\\
0.9826	-0.0155688\\
0.9851	-0.0180793\\
0.9877	-0.0206939\\
0.9902	-0.0232103\\
0.9928	-0.0258287\\
0.9954	-0.0284473\\
0.9979	-0.0309646\\
1.0005	-0.0335808\\
1.0031	-0.0361943\\
1.0056	-0.0387036\\
1.0082	-0.0413084\\
1.0108	-0.043907\\
1.0133	-0.0463987\\
1.0159	-0.048982\\
1.0184	-0.0514571\\
1.021	-0.0540211\\
1.0236	-0.0565735\\
1.0261	-0.0590158\\
1.0287	-0.0615421\\
1.0313	-0.0640536\\
1.0338	-0.0664534\\
1.0364	-0.0689326\\
1.039	-0.0713938\\
1.0415	-0.0737423\\
1.0441	-0.0761647\\
1.0466	-0.0784738\\
1.0492	-0.0808534\\
1.0518	-0.0832095\\
1.0543	-0.0854521\\
1.0569	-0.0877594\\
1.0595	-0.0900401\\
1.062	-0.092207\\
1.0646	-0.0944326\\
1.0672	-0.0966287\\
1.0697	-0.0987116\\
1.0723	-0.100847\\
1.0748	-0.10287\\
1.0774	-0.10494\\
1.08	-0.106977\\
1.0825	-0.108902\\
1.0851	-0.110868\\
1.0877	-0.112798\\
1.0902	-0.114617\\
1.0928	-0.116472\\
1.0954	-0.118287\\
1.0979	-0.119994\\
1.1005	-0.121729\\
1.1031	-0.123422\\
1.1056	-0.12501\\
1.1082	-0.126619\\
1.1107	-0.128125\\
1.1133	-0.129647\\
1.1159	-0.131124\\
1.1184	-0.1325\\
1.121	-0.133886\\
1.1236	-0.135225\\
1.1261	-0.136468\\
1.1287	-0.137713\\
1.1313	-0.138909\\
1.1338	-0.140013\\
1.1364	-0.141113\\
1.1389	-0.142123\\
1.1415	-0.143125\\
1.1441	-0.144076\\
1.1466	-0.144942\\
1.1492	-0.145792\\
1.1518	-0.146591\\
1.1543	-0.147311\\
1.1569	-0.148009\\
1.1595	-0.148654\\
1.162	-0.149226\\
1.1646	-0.149769\\
1.1671	-0.150242\\
1.1697	-0.150683\\
1.1723	-0.151071\\
1.1748	-0.151395\\
1.1774	-0.151681\\
1.18	-0.151915\\
1.1825	-0.152091\\
1.1851	-0.152223\\
1.1877	-0.152304\\
1.1902	-0.152333\\
1.1928	-0.152314\\
1.1953	-0.152247\\
1.1979	-0.152128\\
1.2005	-0.151959\\
1.203	-0.15175\\
1.2056	-0.151484\\
1.2082	-0.15117\\
1.2107	-0.150822\\
1.2133	-0.150413\\
1.2159	-0.149957\\
1.2184	-0.149475\\
1.221	-0.148929\\
1.2235	-0.14836\\
1.2261	-0.147726\\
1.2287	-0.147047\\
1.2312	-0.146353\\
1.2338	-0.14559\\
1.2364	-0.144785\\
1.2389	-0.143972\\
1.2415	-0.143088\\
1.2441	-0.142164\\
1.2466	-0.141239\\
1.2492	-0.14024\\
1.2518	-0.139205\\
1.2543	-0.138176\\
1.2569	-0.137071\\
1.2594	-0.135977\\
1.262	-0.134806\\
1.2646	-0.133604\\
1.2671	-0.132418\\
1.2697	-0.131156\\
1.2723	-0.129865\\
1.2748	-0.128597\\
1.2774	-0.127252\\
1.28	-0.125881\\
1.2825	-0.12454\\
1.2851	-0.123121\\
1.2876	-0.121736\\
1.2902	-0.120274\\
1.2928	-0.118792\\
1.2953	-0.117348\\
1.2979	-0.115829\\
1.3005	-0.114292\\
1.303	-0.112799\\
1.3056	-0.111232\\
1.3082	-0.109652\\
1.3107	-0.108119\\
1.3133	-0.106514\\
1.3158	-0.104961\\
1.3184	-0.103336\\
1.321	-0.101703\\
1.3235	-0.100126\\
1.3261	-0.0984794\\
1.3287	-0.0968275\\
1.3312	-0.0952352\\
1.3338	-0.0935761\\
1.3364	-0.0919149\\
1.3389	-0.0903166\\
1.3415	-0.0886543\\
1.344	-0.0870567\\
1.3466	-0.085397\\
1.3492	-0.0837402\\
1.3517	-0.0821507\\
1.3543	-0.0805023\\
1.3569	-0.0788596\\
1.3594	-0.0772863\\
1.362	-0.0756574\\
1.3646	-0.0740368\\
1.3671	-0.0724873\\
1.3697	-0.0708857\\
1.3722	-0.0693559\\
1.3748	-0.0677763\\
1.3774	-0.0662093\\
1.3799	-0.0647149\\
1.3825	-0.0631745\\
1.3851	-0.0616488\\
1.3876	-0.0601962\\
1.3902	-0.0587014\\
1.3928	-0.0572233\\
1.3953	-0.0558184\\
1.3979	-0.0543749\\
1.4005	-0.05295\\
1.403	-0.051598\\
1.4056	-0.0502112\\
1.4081	-0.0488969\\
1.4107	-0.0475503\\
1.4133	-0.0462249\\
1.4158	-0.044971\\
1.4184	-0.0436886\\
1.421	-0.0424287\\
1.4235	-0.0412389\\
1.4261	-0.0400244\\
1.4287	-0.0388335\\
1.4312	-0.0377111\\
1.4338	-0.0365675\\
1.4363	-0.0354912\\
1.4389	-0.0343962\\
1.4415	-0.0333262\\
1.444	-0.0323213\\
1.4466	-0.0313012\\
1.4492	-0.0303068\\
1.4517	-0.029375\\
1.4543	-0.0284314\\
1.4569	-0.027514\\
1.4594	-0.0266566\\
1.462	-0.0257908\\
1.4645	-0.024983\\
1.4671	-0.0241689\\
1.4697	-0.0233813\\
1.4722	-0.0226489\\
1.4748	-0.0219132\\
1.4774	-0.0212039\\
1.4799	-0.0205469\\
1.4825	-0.0198893\\
1.4851	-0.0192581\\
1.4876	-0.0186759\\
1.4902	-0.018096\\
1.4927	-0.0175629\\
1.4953	-0.0170338\\
1.4979	-0.0165305\\
1.5004	-0.0160707\\
1.503	-0.0156175\\
1.5056	-0.0151895\\
1.5081	-0.0148016\\
1.5107	-0.0144226\\
1.5133	-0.0140682\\
1.5158	-0.0137505\\
1.5184	-0.0134439\\
1.5209	-0.0131717\\
1.5235	-0.0129119\\
1.5261	-0.0126756\\
1.5286	-0.0124702\\
1.5312	-0.0122792\\
1.5338	-0.0121109\\
1.5363	-0.0119703\\
1.5389	-0.0118457\\
1.5415	-0.0117429\\
1.544	-0.0116644\\
1.5466	-0.0116037\\
1.5491	-0.011565\\
1.5517	-0.011545\\
1.5543	-0.0115453\\
1.5568	-0.0115646\\
1.5594	-0.0116039\\
1.562	-0.0116626\\
1.5645	-0.011737\\
1.5671	-0.0118327\\
1.5697	-0.0119469\\
1.5722	-0.0120738\\
1.5748	-0.0122232\\
1.5774	-0.01239\\
1.5799	-0.0125666\\
1.5825	-0.0127667\\
1.585	-0.0129747\\
1.5876	-0.0132068\\
1.5902	-0.0134548\\
1.5927	-0.0137079\\
1.5953	-0.0139861\\
1.5979	-0.0142791\\
1.6004	-0.0145747\\
1.603	-0.014896\\
1.6056	-0.0152314\\
1.6081	-0.0155667\\
1.6107	-0.0159286\\
1.6132	-0.0162889\\
1.6158	-0.0166762\\
1.6184	-0.0170759\\
1.6209	-0.0174718\\
1.6235	-0.0178953\\
1.6261	-0.0183304\\
1.6286	-0.0187595\\
1.6312	-0.0192166\\
1.6338	-0.0196846\\
1.6363	-0.0201446\\
1.6389	-0.0206331\\
1.6414	-0.0211124\\
1.644	-0.0216205\\
1.6466	-0.0221381\\
1.6491	-0.0226447\\
1.6517	-0.0231805\\
1.6543	-0.0237252\\
1.6568	-0.0242571\\
1.6594	-0.0248185\\
1.662	-0.0253882\\
1.6645	-0.0259436\\
1.6671	-0.0265288\\
1.6696	-0.0270987\\
1.6722	-0.0276987\\
1.6748	-0.0283059\\
1.6773	-0.0288963\\
1.6799	-0.0295171\\
1.6825	-0.0301446\\
1.685	-0.030754\\
1.6876	-0.031394\\
1.6902	-0.0320401\\
1.6927	-0.032667\\
1.6953	-0.0333247\\
1.6978	-0.0339624\\
1.7004	-0.0346309\\
1.703	-0.0353048\\
1.7055	-0.0359576\\
1.7081	-0.0366414\\
1.7107	-0.0373301\\
1.7132	-0.0379968\\
1.7158	-0.0386946\\
1.7184	-0.0393968\\
1.7209	-0.040076\\
1.7235	-0.0407865\\
1.7261	-0.0415009\\
1.7286	-0.0421915\\
1.7312	-0.0429134\\
1.7337	-0.0436108\\
1.7363	-0.0443395\\
1.7389	-0.0450715\\
1.7414	-0.0457783\\
1.744	-0.0465163\\
1.7466	-0.0472572\\
1.7491	-0.0479722\\
1.7517	-0.0487183\\
1.7543	-0.0494668\\
1.7568	-0.0501887\\
1.7594	-0.0509415\\
1.7619	-0.0516673\\
1.7645	-0.052424\\
1.7671	-0.0531823\\
1.7696	-0.0539129\\
1.7722	-0.054674\\
1.7748	-0.0554364\\
1.7773	-0.0561704\\
1.7799	-0.0569347\\
1.7825	-0.0576997\\
1.785	-0.0584357\\
1.7876	-0.0592016\\
1.7901	-0.0599381\\
1.7927	-0.0607042\\
1.7953	-0.06147\\
1.7978	-0.062206\\
1.8004	-0.0629709\\
1.803	-0.063735\\
1.8055	-0.0644687\\
1.8081	-0.0652307\\
1.8107	-0.0659912\\
1.8132	-0.0667209\\
1.8158	-0.067478\\
1.8183	-0.068204\\
1.8209	-0.0689568\\
1.8235	-0.069707\\
1.826	-0.0704257\\
1.8286	-0.0711702\\
1.8312	-0.0719114\\
1.8337	-0.0726207\\
1.8363	-0.0733547\\
1.8389	-0.0740846\\
1.8414	-0.0747823\\
1.844	-0.0755034\\
1.8465	-0.0761921\\
1.8491	-0.0769033\\
1.8517	-0.077609\\
1.8542	-0.0782821\\
1.8568	-0.0789762\\
1.8594	-0.0796639\\
1.8619	-0.0803189\\
1.8645	-0.0809933\\
1.8671	-0.0816604\\
1.8696	-0.0822946\\
1.8722	-0.0829465\\
1.8747	-0.0835655\\
1.8773	-0.0842009\\
1.8799	-0.0848273\\
1.8824	-0.085421\\
1.885	-0.0860291\\
1.8876	-0.0866273\\
1.8901	-0.0871929\\
1.8927	-0.0877708\\
1.8953	-0.0883377\\
1.8978	-0.0888723\\
1.9004	-0.089417\\
1.903	-0.0899497\\
1.9055	-0.0904504\\
1.9081	-0.0909589\\
1.9106	-0.0914356\\
1.9132	-0.0919184\\
1.9158	-0.0923876\\
1.9183	-0.0928257\\
1.9209	-0.0932674\\
1.9235	-0.0936945\\
1.926	-0.0940912\\
1.9286	-0.0944888\\
1.9312	-0.0948709\\
1.9337	-0.0952234\\
1.9363	-0.0955742\\
1.9388	-0.095896\\
1.9414	-0.0962143\\
1.944	-0.0965155\\
1.9465	-0.0967888\\
1.9491	-0.0970559\\
1.9517	-0.097305\\
1.9542	-0.0975275\\
1.9568	-0.0977408\\
1.9594	-0.0979356\\
1.9619	-0.098105\\
1.9645	-0.0982624\\
1.967	-0.0983956\\
1.9696	-0.0985149\\
1.9722	-0.0986144\\
1.9747	-0.0986913\\
1.9773	-0.0987515\\
1.9799	-0.0987914\\
1.9824	-0.0988105\\
1.985	-0.09881\\
1.9876	-0.0987888\\
1.9901	-0.0987486\\
1.9927	-0.0986863\\
1.9952	-0.0986064\\
1.9978	-0.0985024\\
2.0004	-0.0983772\\
2.0029	-0.0982366\\
2.0055	-0.0980694\\
2.0081	-0.0978808\\
2.0106	-0.0976791\\
2.0132	-0.0974482\\
2.0158	-0.0971959\\
2.0183	-0.096933\\
2.0209	-0.0966385\\
2.0234	-0.0963352\\
2.026	-0.0959988\\
2.0286	-0.0956411\\
2.0311	-0.0952772\\
2.0337	-0.0948781\\
2.0363	-0.0944581\\
2.0388	-0.0940346\\
2.0414	-0.0935738\\
2.044	-0.0930926\\
2.0465	-0.0926108\\
2.0491	-0.0920899\\
2.0517	-0.0915492\\
2.0542	-0.0910107\\
2.0568	-0.0904316\\
2.0593	-0.0898567\\
2.0619	-0.0892403\\
2.0645	-0.0886052\\
2.067	-0.0879774\\
2.0696	-0.0873068\\
2.0722	-0.0866185\\
2.0747	-0.0859404\\
2.0773	-0.0852185\\
2.0799	-0.0844801\\
2.0824	-0.0837548\\
2.085	-0.0829851\\
2.0875	-0.0822304\\
2.0901	-0.0814309\\
2.0927	-0.0806168\\
2.0952	-0.0798208\\
2.0978	-0.0789797\\
2.1004	-0.0781254\\
2.1029	-0.077292\\
2.1055	-0.0764134\\
2.1081	-0.0755232\\
2.1106	-0.0746567\\
2.1132	-0.0737451\\
2.1157	-0.0728591\\
2.1183	-0.0719283\\
2.1209	-0.0709885\\
2.1234	-0.0700769\\
2.126	-0.0691212\\
2.1286	-0.0681581\\
2.1311	-0.0672258\\
2.1337	-0.0662501\\
2.1363	-0.0652689\\
2.1388	-0.0643208\\
2.1414	-0.0633305\\
2.1439	-0.0623748\\
2.1465	-0.0613777\\
2.1491	-0.0603781\\
2.1516	-0.0594151\\
2.1542	-0.0584123\\
2.1568	-0.0574088\\
2.1593	-0.0564439\\
2.1619	-0.0554408\\
2.1645	-0.054439\\
2.167	-0.0534773\\
2.1696	-0.0524794\\
2.1721	-0.0515228\\
2.1747	-0.0505314\\
2.1773	-0.0495443\\
2.1798	-0.0485996\\
2.1824	-0.0476225\\
2.185	-0.0466513\\
2.1875	-0.0457238\\
2.1901	-0.0447661\\
2.1927	-0.0438163\\
2.1952	-0.0429109\\
2.1978	-0.041978\\
2.2004	-0.0410546\\
2.2029	-0.0401762\\
2.2055	-0.039273\\
2.208	-0.0384151\\
2.2106	-0.0375342\\
2.2132	-0.0366656\\
2.2157	-0.0358424\\
2.2183	-0.0349992\\
2.2209	-0.0341697\\
2.2234	-0.0333854\\
2.226	-0.0325843\\
2.2286	-0.0317982\\
2.2311	-0.0310571\\
2.2337	-0.030302\\
2.2362	-0.0295916\\
2.2388	-0.0288692\\
2.2414	-0.0281642\\
2.2439	-0.0275029\\
2.2465	-0.026833\\
2.2491	-0.0261814\\
2.2516	-0.0255727\\
2.2542	-0.0249584\\
2.2568	-0.0243635\\
2.2593	-0.0238102\\
2.2619	-0.0232545\\
2.2644	-0.0227394\\
2.267	-0.0222239\\
2.2696	-0.0217294\\
2.2721	-0.0212739\\
2.2747	-0.0208211\\
2.2773	-0.02039\\
2.2798	-0.019996\\
2.2824	-0.0196079\\
2.285	-0.0192421\\
2.2875	-0.0189114\\
2.2901	-0.0185895\\
2.2926	-0.0183014\\
2.2952	-0.0180242\\
2.2978	-0.0177698\\
2.3003	-0.0175469\\
2.3029	-0.0173377\\
2.3055	-0.0171517\\
2.308	-0.0169946\\
2.3106	-0.0168541\\
2.3132	-0.0167368\\
2.3157	-0.016646\\
2.3183	-0.0165743\\
2.3208	-0.0165273\\
2.3234	-0.0165012\\
2.326	-0.0164983\\
2.3285	-0.0165173\\
2.3311	-0.0165597\\
2.3337	-0.016625\\
2.3362	-0.0167094\\
2.3388	-0.0168196\\
2.3414	-0.0169523\\
2.3439	-0.0171013\\
2.3465	-0.0172781\\
2.349	-0.0174691\\
2.3516	-0.0176894\\
2.3542	-0.0179316\\
2.3567	-0.018185\\
2.3593	-0.0184697\\
2.3619	-0.0187756\\
2.3644	-0.0190897\\
2.367	-0.0194368\\
2.3696	-0.0198046\\
2.3721	-0.0201775\\
2.3747	-0.020585\\
2.3773	-0.0210124\\
2.3798	-0.0214419\\
2.3824	-0.0219075\\
2.3849	-0.0223732\\
2.3875	-0.0228759\\
2.3901	-0.023397\\
2.3926	-0.0239152\\
2.3952	-0.0244716\\
2.3978	-0.0250455\\
2.4003	-0.0256135\\
2.4029	-0.0262207\\
2.4055	-0.0268444\\
2.408	-0.0274593\\
2.4106	-0.0281142\\
2.4131	-0.0287585\\
2.4157	-0.0294434\\
2.4183	-0.0301429\\
2.4208	-0.030829\\
2.4234	-0.0315561\\
2.426	-0.0322968\\
2.4285	-0.0330214\\
2.4311	-0.0337874\\
2.4337	-0.0345658\\
2.4362	-0.0353256\\
2.4388	-0.036127\\
2.4413	-0.0369081\\
2.4439	-0.037731\\
2.4465	-0.0385641\\
2.449	-0.0393746\\
2.4516	-0.0402268\\
2.4542	-0.041088\\
2.4567	-0.0419243\\
2.4593	-0.0428021\\
2.4619	-0.0436877\\
2.4644	-0.0445461\\
2.467	-0.0454457\\
2.4695	-0.0463169\\
2.4721	-0.0472289\\
2.4747	-0.0481466\\
2.4772	-0.049034\\
2.4798	-0.0499615\\
2.4824	-0.0508935\\
2.4849	-0.0517934\\
2.4875	-0.0527327\\
2.4901	-0.0536752\\
2.4926	-0.0545839\\
2.4952	-0.0555313\\
2.4977	-0.0564439\\
2.5003	-0.0573945\\
2.5029	-0.0583462\\
2.5054	-0.0592618\\
2.508	-0.0602142\\
2.5106	-0.0611664\\
2.5131	-0.0620814\\
2.5157	-0.063032\\
2.5183	-0.0639811\\
2.5208	-0.064892\\
2.5234	-0.0658372\\
2.526	-0.0667797\\
2.5285	-0.0676831\\
2.5311	-0.0686193\\
2.5336	-0.0695159\\
2.5362	-0.0704443\\
2.5388	-0.0713681\\
2.5413	-0.0722518\\
2.5439	-0.0731655\\
2.5465	-0.0740736\\
2.549	-0.074941\\
2.5516	-0.0758369\\
2.5542	-0.076726\\
2.5567	-0.0775742\\
2.5593	-0.078449\\
2.5618	-0.0792828\\
2.5644	-0.080142\\
2.567	-0.0809927\\
2.5695	-0.0818023\\
2.5721	-0.0826354\\
2.5747	-0.083459\\
2.5772	-0.0842418\\
2.5798	-0.0850459\\
2.5824	-0.0858397\\
2.5849	-0.0865929\\
2.5875	-0.0873654\\
2.59	-0.0880976\\
2.5926	-0.0888478\\
2.5952	-0.0895861\\
2.5977	-0.0902846\\
2.6003	-0.0909989\\
2.6029	-0.0917007\\
2.6054	-0.0923632\\
2.608	-0.0930394\\
2.6106	-0.0937022\\
2.6131	-0.0943267\\
2.6157	-0.0949626\\
2.6182	-0.0955607\\
2.6208	-0.0961688\\
2.6234	-0.0967623\\
2.6259	-0.0973192\\
2.6285	-0.0978837\\
2.6311	-0.0984331\\
2.6336	-0.0989469\\
2.6362	-0.0994661\\
2.6388	-0.0999697\\
2.6413	-0.100439\\
2.6439	-0.100911\\
2.6464	-0.10135\\
2.649	-0.101791\\
2.6516	-0.102215\\
2.6541	-0.102607\\
2.6567	-0.102999\\
2.6593	-0.103373\\
2.6618	-0.103718\\
2.6644	-0.104059\\
2.667	-0.104383\\
2.6695	-0.104679\\
2.6721	-0.104969\\
2.6746	-0.105231\\
2.6772	-0.105487\\
2.6798	-0.105726\\
2.6823	-0.105939\\
2.6849	-0.106142\\
2.6875	-0.106329\\
2.69	-0.106491\\
2.6926	-0.106642\\
2.6952	-0.106775\\
2.6977	-0.106887\\
2.7003	-0.106985\\
2.7029	-0.107066\\
2.7054	-0.107126\\
2.708	-0.107171\\
2.7105	-0.107198\\
2.7131	-0.107209\\
2.7157	-0.107201\\
2.7182	-0.107177\\
2.7208	-0.107135\\
2.7234	-0.107075\\
2.7259	-0.107001\\
2.7285	-0.106906\\
2.7311	-0.106794\\
2.7336	-0.10667\\
2.7362	-0.106524\\
2.7387	-0.106367\\
2.7413	-0.106187\\
2.7439	-0.10599\\
2.7464	-0.105785\\
2.749	-0.105555\\
2.7516	-0.105308\\
2.7541	-0.105055\\
2.7567	-0.104776\\
2.7593	-0.10448\\
2.7618	-0.10418\\
2.7644	-0.103853\\
2.7669	-0.103523\\
2.7695	-0.103165\\
2.7721	-0.102791\\
2.7746	-0.102417\\
2.7772	-0.102013\\
2.7798	-0.101594\\
2.7823	-0.101177\\
2.7849	-0.100729\\
2.7875	-0.100266\\
2.79	-0.0998082\\
2.7926	-0.0993181\\
2.7951	-0.0988338\\
2.7977	-0.0983168\\
2.8003	-0.0977864\\
2.8028	-0.0972641\\
2.8054	-0.0967082\\
2.808	-0.0961397\\
2.8105	-0.0955814\\
2.8131	-0.0949888\\
2.8157	-0.0943843\\
2.8182	-0.0937921\\
2.8208	-0.093165\\
2.8233	-0.0925516\\
2.8259	-0.0919031\\
2.8285	-0.091244\\
2.831	-0.0906005\\
2.8336	-0.0899216\\
2.8362	-0.0892329\\
2.8387	-0.0885618\\
2.8413	-0.087855\\
2.8439	-0.0871394\\
2.8464	-0.0864433\\
2.849	-0.0857113\\
2.8516	-0.0849714\\
2.8541	-0.0842527\\
2.8567	-0.0834983\\
2.8592	-0.0827662\\
2.8618	-0.0819984\\
2.8644	-0.0812243\\
2.8669	-0.0804743\\
2.8695	-0.0796887\\
2.8721	-0.0788977\\
2.8746	-0.0781325\\
2.8772	-0.077332\\
2.8798	-0.0765272\\
2.8823	-0.0757495\\
2.8849	-0.0749371\\
2.8874	-0.0741527\\
2.89	-0.073334\\
2.8926	-0.0725127\\
2.8951	-0.0717206\\
2.8977	-0.0708949\\
2.9003	-0.0700676\\
2.9028	-0.0692707\\
2.9054	-0.068441\\
2.908	-0.0676106\\
2.9105	-0.0668118\\
2.9131	-0.0659811\\
2.9156	-0.0651826\\
2.9182	-0.0643528\\
2.9208	-0.063524\\
2.9233	-0.0627283\\
2.9259	-0.0619024\\
2.9285	-0.0610784\\
2.931	-0.0602883\\
2.9336	-0.0594692\\
2.9362	-0.0586531\\
2.9387	-0.0578714\\
2.9413	-0.057062\\
2.9438	-0.0562874\\
2.9464	-0.0554859\\
2.949	-0.054689\\
2.9515	-0.0539272\\
2.9541	-0.0531401\\
2.9567	-0.0523584\\
2.9592	-0.0516122\\
2.9618	-0.0508421\\
2.9644	-0.0500783\\
2.9669	-0.0493501\\
2.9695	-0.0485996\\
2.972	-0.0478848\\
2.9746	-0.0471487\\
2.9772	-0.0464203\\
2.9797	-0.0457275\\
2.9823	-0.0450152\\
2.9849	-0.0443114\\
2.9874	-0.043643\\
2.99	-0.0429567\\
2.9926	-0.0422798\\
2.9951	-0.0416379\\
2.9977	-0.0409799\\
3.0003	-0.040332\\
3.0028	-0.0397187\\
3.0054	-0.0390912\\
3.0079	-0.0384979\\
3.0105	-0.0378915\\
3.0131	-0.0372963\\
3.0156	-0.0367347\\
3.0182	-0.036162\\
3.0208	-0.035601\\
3.0233	-0.0350728\\
3.0259	-0.0345354\\
3.0285	-0.0340103\\
3.031	-0.0335171\\
3.0336	-0.0330166\\
3.0361	-0.0325473\\
3.0387	-0.032072\\
3.0413	-0.0316098\\
3.0438	-0.0311779\\
3.0464	-0.0307418\\
3.049	-0.0303191\\
3.0515	-0.0299257\\
3.0541	-0.0295299\\
3.0567	-0.029148\\
3.0592	-0.028794\\
3.0618	-0.0284396\\
3.0643	-0.0281122\\
3.0669	-0.0277857\\
3.0695	-0.0274736\\
3.072	-0.0271871\\
3.0746	-0.0269033\\
3.0772	-0.0266341\\
3.0797	-0.026389\\
3.0823	-0.0261486\\
3.0849	-0.0259228\\
3.0874	-0.0257197\\
3.09	-0.025523\\
3.0925	-0.0253479\\
3.0951	-0.0251803\\
3.0977	-0.0250276\\
3.1002	-0.0248949\\
3.1028	-0.0247715\\
3.1054	-0.024663\\
3.1079	-0.0245727\\
3.1105	-0.0244935\\
3.1131	-0.0244291\\
3.1156	-0.0243812\\
3.1182	-0.024346\\
3.1207	-0.0243262\\
3.1233	-0.02432\\
3.1259	-0.0243286\\
3.1284	-0.0243507\\
3.131	-0.0243881\\
3.1336	-0.0244401\\
3.1361	-0.0245037\\
3.1387	-0.0245842\\
3.1413	-0.024679\\
3.1438	-0.0247836\\
3.1464	-0.0249065\\
3.1489	-0.0250379\\
3.1515	-0.0251884\\
3.1541	-0.0253529\\
3.1566	-0.0255241\\
3.1592	-0.0257156\\
3.1618	-0.0259209\\
3.1643	-0.026131\\
3.1669	-0.0263626\\
3.1695	-0.0266076\\
3.172	-0.0268556\\
3.1746	-0.0271263\\
3.1772	-0.0274099\\
3.1797	-0.0276946\\
3.1823	-0.0280031\\
3.1848	-0.0283115\\
3.1874	-0.0286444\\
3.19	-0.0289894\\
3.1925	-0.0293325\\
3.1951	-0.0297009\\
3.1977	-0.030081\\
3.2002	-0.0304574\\
3.2028	-0.0308599\\
3.2054	-0.0312736\\
3.2079	-0.0316818\\
3.2105	-0.0321169\\
3.213	-0.0325452\\
3.2156	-0.0330009\\
3.2182	-0.0334669\\
3.2207	-0.0339244\\
3.2233	-0.0344098\\
3.2259	-0.0349049\\
3.2284	-0.0353898\\
3.231	-0.0359032\\
3.2336	-0.0364256\\
3.2361	-0.0369362\\
3.2387	-0.0374756\\
3.2412	-0.0380021\\
3.2438	-0.0385577\\
3.2464	-0.0391212\\
3.2489	-0.0396702\\
3.2515	-0.0402485\\
3.2541	-0.0408341\\
3.2566	-0.0414037\\
3.2592	-0.0420027\\
3.2618	-0.0426082\\
3.2643	-0.0431963\\
3.2669	-0.0438138\\
3.2694	-0.044413\\
3.272	-0.0450416\\
3.2746	-0.0456755\\
3.2771	-0.0462897\\
3.2797	-0.0469332\\
3.2823	-0.0475812\\
3.2848	-0.0482082\\
3.2874	-0.0488643\\
3.29	-0.049524\\
3.2925	-0.0501617\\
3.2951	-0.050828\\
3.2976	-0.0514714\\
3.3002	-0.0521432\\
3.3028	-0.0528173\\
3.3053	-0.0534676\\
3.3079	-0.0541456\\
3.3105	-0.0548253\\
3.313	-0.05548\\
3.3156	-0.056162\\
3.3182	-0.0568446\\
3.3207	-0.0575015\\
3.3233	-0.0581848\\
3.3259	-0.0588681\\
3.3284	-0.0595247\\
3.331	-0.060207\\
3.3335	-0.0608622\\
3.3361	-0.0615423\\
3.3387	-0.0622211\\
3.3412	-0.062872\\
3.3438	-0.063547\\
3.3464	-0.0642196\\
3.3489	-0.0648639\\
3.3515	-0.0655312\\
3.3541	-0.0661953\\
3.3566	-0.0668305\\
3.3592	-0.0674876\\
3.3617	-0.0681156\\
3.3643	-0.0687645\\
3.3669	-0.0694088\\
3.3694	-0.0700237\\
3.372	-0.0706583\\
3.3746	-0.0712874\\
3.3771	-0.071887\\
3.3797	-0.0725048\\
3.3823	-0.0731163\\
3.3848	-0.0736982\\
3.3874	-0.0742967\\
3.3899	-0.0748656\\
3.3925	-0.0754501\\
3.3951	-0.076027\\
3.3976	-0.0765744\\
3.4002	-0.0771357\\
3.4028	-0.0776888\\
3.4053	-0.0782124\\
3.4079	-0.0787483\\
3.4105	-0.0792752\\
3.413	-0.0797729\\
3.4156	-0.0802812\\
3.4181	-0.0807607\\
3.4207	-0.0812495\\
3.4233	-0.081728\\
3.4258	-0.0821782\\
3.4284	-0.0826358\\
3.431	-0.0830825\\
3.4335	-0.0835015\\
3.4361	-0.083926\\
3.4387	-0.0843389\\
3.4412	-0.0847248\\
3.4438	-0.0851144\\
3.4463	-0.0854774\\
3.4489	-0.0858428\\
3.4515	-0.0861956\\
3.454	-0.0865228\\
3.4566	-0.0868503\\
3.4592	-0.0871647\\
3.4617	-0.0874545\\
3.4643	-0.0877427\\
3.4669	-0.0880173\\
3.4694	-0.0882683\\
3.472	-0.0885157\\
3.4745	-0.0887404\\
3.4771	-0.0889602\\
3.4797	-0.0891656\\
3.4822	-0.0893495\\
3.4848	-0.0895266\\
3.4874	-0.0896889\\
3.4899	-0.0898312\\
3.4925	-0.0899645\\
3.4951	-0.0900829\\
3.4976	-0.0901827\\
3.5002	-0.0902716\\
3.5028	-0.0903453\\
3.5053	-0.0904019\\
3.5079	-0.0904458\\
3.5104	-0.0904736\\
3.513	-0.0904874\\
3.5156	-0.0904858\\
3.5181	-0.0904698\\
3.5207	-0.090438\\
3.5233	-0.0903908\\
3.5258	-0.0903309\\
3.5284	-0.0902535\\
3.531	-0.0901607\\
3.5335	-0.0900569\\
3.5361	-0.089934\\
3.5386	-0.0898013\\
3.5412	-0.0896484\\
3.5438	-0.0894802\\
3.5463	-0.0893043\\
3.5489	-0.0891065\\
3.5515	-0.0888938\\
3.554	-0.0886751\\
3.5566	-0.0884332\\
3.5592	-0.0881765\\
3.5617	-0.087916\\
3.5643	-0.0876308\\
3.5668	-0.087343\\
3.5694	-0.0870297\\
3.572	-0.0867024\\
3.5745	-0.0863744\\
3.5771	-0.0860198\\
3.5797	-0.0856516\\
3.5822	-0.0852848\\
3.5848	-0.0848903\\
3.5874	-0.0844826\\
3.5899	-0.0840784\\
3.5925	-0.0836456\\
3.595	-0.0832176\\
3.5976	-0.0827604\\
3.6002	-0.082291\\
3.6027	-0.0818285\\
3.6053	-0.081336\\
3.6079	-0.0808321\\
3.6104	-0.080337\\
3.613	-0.0798114\\
3.6156	-0.0792751\\
3.6181	-0.0787495\\
3.6207	-0.078193\\
3.6232	-0.0776485\\
3.6258	-0.0770727\\
3.6284	-0.0764876\\
3.6309	-0.0759164\\
3.6335	-0.0753137\\
3.6361	-0.0747025\\
3.6386	-0.0741071\\
3.6412	-0.0734801\\
3.6438	-0.0728455\\
3.6463	-0.0722284\\
3.6489	-0.0715798\\
3.6515	-0.0709245\\
3.654	-0.0702883\\
3.6566	-0.0696208\\
3.6591	-0.0689736\\
3.6617	-0.0682952\\
3.6643	-0.0676117\\
3.6668	-0.06695\\
3.6694	-0.0662576\\
3.672	-0.065561\\
3.6745	-0.0648877\\
3.6771	-0.0641841\\
3.6797	-0.0634775\\
3.6822	-0.0627954\\
3.6848	-0.0620837\\
3.6873	-0.0613974\\
3.6899	-0.060682\\
3.6925	-0.0599652\\
3.695	-0.0592749\\
3.6976	-0.0585564\\
3.7002	-0.0578375\\
3.7027	-0.0571462\\
3.7053	-0.0564275\\
3.7079	-0.0557095\\
3.7104	-0.05502\\
3.713	-0.0543043\\
3.7155	-0.0536175\\
3.7181	-0.0529053\\
3.7207	-0.0521953\\
3.7232	-0.0515151\\
3.7258	-0.0508105\\
3.7284	-0.0501092\\
3.7309	-0.0494383\\
3.7335	-0.0487443\\
3.7361	-0.0480545\\
3.7386	-0.0473955\\
3.7412	-0.0467148\\
3.7437	-0.0460651\\
3.7463	-0.0453948\\
3.7489	-0.0447301\\
3.7514	-0.0440966\\
3.754	-0.0434439\\
3.7566	-0.0427978\\
3.7591	-0.0421829\\
3.7617	-0.0415505\\
3.7643	-0.0409253\\
3.7668	-0.0403314\\
3.7694	-0.0397215\\
3.7719	-0.0391427\\
3.7745	-0.0385489\\
3.7771	-0.0379638\\
3.7796	-0.0374095\\
3.7822	-0.036842\\
3.7848	-0.0362837\\
3.7873	-0.0357559\\
3.7899	-0.0352165\\
3.7925	-0.0346871\\
3.795	-0.0341876\\
3.7976	-0.0336783\\
3.8002	-0.0331794\\
3.8027	-0.0327099\\
3.8053	-0.0322323\\
3.8078	-0.0317835\\
3.8104	-0.0313277\\
3.813	-0.0308833\\
3.8155	-0.0304669\\
3.8181	-0.0300453\\
3.8207	-0.0296355\\
3.8232	-0.0292528\\
3.8258	-0.0288665\\
3.8284	-0.0284925\\
3.8309	-0.0281445\\
3.8335	-0.0277948\\
3.836	-0.0274703\\
3.8386	-0.0271452\\
3.8412	-0.0268328\\
3.8437	-0.0265445\\
3.8463	-0.0262573\\
3.8489	-0.025983\\
3.8514	-0.0257316\\
3.854	-0.0254828\\
3.8566	-0.0252472\\
3.8591	-0.025033\\
3.8617	-0.0248232\\
3.8642	-0.024634\\
3.8668	-0.0244501\\
3.8694	-0.0242796\\
3.8719	-0.0241281\\
3.8745	-0.0239835\\
3.8771	-0.0238523\\
3.8796	-0.0237386\\
3.8822	-0.0236334\\
3.8848	-0.0235415\\
3.8873	-0.0234655\\
3.8899	-0.0233995\\
3.8924	-0.0233484\\
3.895	-0.0233082\\
3.8976	-0.023281\\
3.9001	-0.0232672\\
3.9027	-0.0232656\\
3.9053	-0.0232769\\
3.9078	-0.0232999\\
3.9104	-0.0233363\\
3.913	-0.0233854\\
3.9155	-0.0234446\\
3.9181	-0.0235185\\
3.9206	-0.0236012\\
3.9232	-0.0236994\\
3.9258	-0.0238098\\
3.9283	-0.0239275\\
3.9309	-0.0240617\\
3.9335	-0.0242078\\
3.936	-0.0243595\\
3.9386	-0.0245287\\
3.9412	-0.0247095\\
3.9437	-0.0248941\\
3.9463	-0.0250972\\
3.9488	-0.0253031\\
3.9514	-0.025528\\
3.954	-0.0257638\\
3.9565	-0.0260007\\
3.9591	-0.0262575\\
3.9617	-0.0265247\\
3.9642	-0.0267913\\
3.9668	-0.0270786\\
3.9694	-0.0273759\\
3.9719	-0.0276709\\
3.9745	-0.0279873\\
3.9771	-0.0283131\\
3.9796	-0.0286351\\
3.9822	-0.028979\\
3.9847	-0.0293181\\
3.9873	-0.0296793\\
3.9899	-0.0300492\\
3.9924	-0.0304127\\
3.995	-0.0307989\\
3.9976	-0.0311931\\
4.0001	-0.0315795\\
4.0027	-0.0319889\\
4.0053	-0.0324057\\
4.0078	-0.0328134\\
4.0104	-0.0332443\\
4.0129	-0.0336651\\
4.0155	-0.0341093\\
4.0181	-0.03456\\
4.0206	-0.0349992\\
4.0232	-0.035462\\
4.0258	-0.0359306\\
4.0283	-0.0363866\\
4.0309	-0.0368661\\
4.0335	-0.037351\\
4.036	-0.0378219\\
4.0386	-0.0383164\\
4.0411	-0.0387962\\
4.0437	-0.0392996\\
4.0463	-0.0398073\\
4.0488	-0.0402991\\
4.0514	-0.0408144\\
4.054	-0.0413333\\
4.0565	-0.0418355\\
4.0591	-0.0423608\\
4.0617	-0.0428891\\
4.0642	-0.0433997\\
4.0668	-0.0439332\\
4.0693	-0.0444484\\
4.0719	-0.0449864\\
4.0745	-0.0455262\\
4.077	-0.046047\\
4.0796	-0.04659\\
4.0822	-0.0471344\\
4.0847	-0.0476589\\
4.0873	-0.0482052\\
4.0899	-0.0487523\\
4.0924	-0.0492788\\
4.095	-0.0498267\\
4.0975	-0.0503535\\
4.1001	-0.0509014\\
4.1027	-0.051449\\
4.1052	-0.051975\\
4.1078	-0.0525214\\
4.1104	-0.0530669\\
4.1129	-0.0535904\\
4.1155	-0.0541336\\
4.1181	-0.0546753\\
4.1206	-0.0551947\\
4.1232	-0.055733\\
4.1258	-0.0562692\\
4.1283	-0.0567828\\
4.1309	-0.0573145\\
4.1334	-0.0578233\\
4.136	-0.0583498\\
4.1386	-0.0588733\\
4.1411	-0.0593737\\
4.1437	-0.059891\\
4.1463	-0.0604047\\
4.1488	-0.0608952\\
4.1514	-0.0614016\\
4.154	-0.0619041\\
4.1565	-0.0623832\\
4.1591	-0.0628773\\
4.1616	-0.0633482\\
4.1642	-0.0638333\\
4.1668	-0.0643136\\
4.1693	-0.0647708\\
4.1719	-0.0652413\\
4.1745	-0.0657064\\
4.177	-0.0661486\\
4.1796	-0.066603\\
4.1822	-0.0670516\\
4.1847	-0.0674775\\
4.1873	-0.0679146\\
4.1898	-0.068329\\
4.1924	-0.0687539\\
4.195	-0.0691723\\
4.1975	-0.0695685\\
4.2001	-0.069974\\
4.2027	-0.0703727\\
4.2052	-0.0707496\\
4.2078	-0.0711346\\
4.2104	-0.0715125\\
4.2129	-0.0718689\\
4.2155	-0.0722324\\
4.218	-0.0725748\\
4.2206	-0.0729235\\
4.2232	-0.0732645\\
4.2257	-0.0735849\\
4.2283	-0.0739105\\
4.2309	-0.074228\\
4.2334	-0.0745257\\
4.236	-0.0748273\\
4.2386	-0.0751206\\
4.2411	-0.0753947\\
4.2437	-0.0756716\\
4.2462	-0.0759297\\
4.2488	-0.0761898\\
4.2514	-0.0764412\\
4.2539	-0.0766746\\
4.2565	-0.0769089\\
4.2591	-0.0771342\\
4.2616	-0.0773425\\
4.2642	-0.0775503\\
4.2668	-0.0777491\\
4.2693	-0.0779317\\
4.2719	-0.0781127\\
4.2744	-0.0782781\\
4.277	-0.0784411\\
4.2796	-0.0785948\\
4.2821	-0.0787339\\
4.2847	-0.0788695\\
4.2873	-0.0789958\\
4.2898	-0.0791084\\
4.2924	-0.0792163\\
4.295	-0.0793149\\
4.2975	-0.0794009\\
4.3001	-0.0794811\\
4.3027	-0.0795519\\
4.3052	-0.0796111\\
4.3078	-0.0796635\\
4.3103	-0.0797051\\
4.3129	-0.0797392\\
4.3155	-0.0797639\\
4.318	-0.0797788\\
4.3206	-0.0797853\\
4.3232	-0.0797824\\
4.3257	-0.079771\\
4.3283	-0.0797501\\
4.3309	-0.0797199\\
4.3334	-0.0796824\\
4.336	-0.0796344\\
4.3385	-0.0795797\\
4.3411	-0.079514\\
4.3437	-0.0794393\\
4.3462	-0.0793592\\
4.3488	-0.0792672\\
4.3514	-0.0791664\\
4.3539	-0.0790613\\
4.3565	-0.0789436\\
4.3591	-0.0788173\\
4.3616	-0.0786879\\
4.3642	-0.0785451\\
4.3667	-0.0784\\
4.3693	-0.078241\\
4.3719	-0.0780738\\
4.3744	-0.0779056\\
4.377	-0.0777228\\
4.3796	-0.0775323\\
4.3821	-0.0773417\\
4.3847	-0.0771361\\
4.3873	-0.076923\\
4.3898	-0.0767112\\
4.3924	-0.0764837\\
4.3949	-0.0762582\\
4.3975	-0.0760168\\
4.4001	-0.0757686\\
4.4026	-0.0755234\\
4.4052	-0.075262\\
4.4078	-0.0749941\\
4.4103	-0.0747304\\
4.4129	-0.0744501\\
4.4155	-0.0741637\\
4.418	-0.0738828\\
4.4206	-0.0735848\\
4.4231	-0.0732931\\
4.4257	-0.0729842\\
4.4283	-0.07267\\
4.4308	-0.072363\\
4.4334	-0.0720388\\
4.436	-0.0717096\\
4.4385	-0.0713887\\
4.4411	-0.0710505\\
4.4437	-0.0707079\\
4.4462	-0.0703745\\
4.4488	-0.0700238\\
4.4514	-0.0696692\\
4.4539	-0.0693247\\
4.4565	-0.0689629\\
4.459	-0.0686119\\
4.4616	-0.0682437\\
4.4642	-0.0678725\\
4.4667	-0.0675129\\
4.4693	-0.0671363\\
4.4719	-0.0667571\\
4.4744	-0.0663904\\
4.477	-0.0660069\\
4.4796	-0.0656214\\
4.4821	-0.0652491\\
4.4847	-0.0648603\\
4.4872	-0.0644852\\
4.4898	-0.0640938\\
4.4924	-0.0637014\\
4.4949	-0.0633232\\
4.4975	-0.0629292\\
4.5001	-0.0625348\\
4.5026	-0.0621551\\
4.5052	-0.0617602\\
4.5078	-0.0613652\\
4.5103	-0.0609857\\
4.5129	-0.0605914\\
4.5154	-0.0602127\\
4.518	-0.0598196\\
4.5206	-0.0594275\\
4.5231	-0.0590515\\
4.5257	-0.0586616\\
4.5283	-0.0582732\\
4.5308	-0.0579013\\
4.5334	-0.0575163\\
4.536	-0.0571332\\
4.5385	-0.0567668\\
4.5411	-0.0563881\\
4.5436	-0.0560262\\
4.5462	-0.0556524\\
4.5488	-0.0552814\\
4.5513	-0.0549274\\
4.5539	-0.0545624\\
4.5565	-0.0542006\\
4.559	-0.0538559\\
4.5616	-0.0535009\\
4.5642	-0.0531497\\
4.5667	-0.0528156\\
4.5693	-0.052472\\
4.5718	-0.0521456\\
4.5744	-0.0518104\\
4.577	-0.0514796\\
4.5795	-0.0511658\\
4.5821	-0.0508441\\
4.5847	-0.0505272\\
4.5872	-0.0502272\\
4.5898	-0.0499202\\
4.5924	-0.0496184\\
4.5949	-0.0493333\\
4.5975	-0.0490421\\
4.6001	-0.0487564\\
4.6026	-0.0484871\\
4.6052	-0.0482128\\
4.6077	-0.0479545\\
4.6103	-0.0476918\\
4.6129	-0.0474351\\
4.6154	-0.0471942\\
4.618	-0.0469498\\
4.6206	-0.0467118\\
4.6231	-0.0464891\\
4.6257	-0.0462638\\
4.6283	-0.0460452\\
4.6308	-0.0458414\\
4.6334	-0.045636\\
4.6359	-0.045445\\
4.6385	-0.0452531\\
4.6411	-0.0450683\\
4.6436	-0.0448972\\
4.6462	-0.0447262\\
4.6488	-0.0445623\\
4.6513	-0.0444116\\
4.6539	-0.044262\\
4.6565	-0.0441197\\
4.659	-0.0439897\\
4.6616	-0.0438619\\
4.6641	-0.0437459\\
4.6667	-0.0436326\\
4.6693	-0.0435268\\
4.6718	-0.0434321\\
4.6744	-0.043341\\
4.677	-0.0432574\\
4.6795	-0.0431842\\
4.6821	-0.0431155\\
4.6847	-0.0430543\\
4.6872	-0.0430026\\
4.6898	-0.0429563\\
4.6923	-0.0429189\\
4.6949	-0.0428875\\
4.6975	-0.0428636\\
4.7	-0.0428477\\
4.7026	-0.0428386\\
4.7052	-0.042837\\
4.7077	-0.0428425\\
4.7103	-0.0428555\\
4.7129	-0.042876\\
4.7154	-0.0429027\\
4.718	-0.0429376\\
4.7205	-0.0429782\\
4.7231	-0.0430275\\
4.7257	-0.043084\\
4.7282	-0.0431452\\
4.7308	-0.0432158\\
4.7334	-0.0432935\\
4.7359	-0.0433749\\
4.7385	-0.0434663\\
4.7411	-0.0435647\\
4.7436	-0.0436658\\
4.7462	-0.0437776\\
4.7487	-0.0438914\\
4.7513	-0.0440163\\
4.7539	-0.0441477\\
4.7564	-0.0442803\\
4.759	-0.0444244\\
4.7616	-0.0445749\\
4.7641	-0.0447254\\
4.7667	-0.0448881\\
4.7693	-0.0450568\\
4.7718	-0.0452247\\
4.7744	-0.045405\\
4.777	-0.0455912\\
4.7795	-0.0457755\\
4.7821	-0.0459727\\
4.7846	-0.0461675\\
4.7872	-0.0463754\\
4.7898	-0.0465886\\
4.7923	-0.0467984\\
4.7949	-0.0470216\\
4.7975	-0.0472497\\
4.8	-0.0474736\\
4.8026	-0.0477111\\
4.8052	-0.0479532\\
4.8077	-0.0481901\\
4.8103	-0.0484408\\
4.8128	-0.0486859\\
4.8154	-0.0489448\\
4.818	-0.0492076\\
4.8205	-0.049464\\
4.8231	-0.0497343\\
4.8257	-0.0500082\\
4.8282	-0.0502748\\
4.8308	-0.0505553\\
4.8334	-0.050839\\
4.8359	-0.0511146\\
4.8385	-0.0514042\\
4.841	-0.0516852\\
4.8436	-0.05198\\
4.8462	-0.0522774\\
4.8487	-0.0525655\\
4.8513	-0.0528673\\
4.8539	-0.0531712\\
4.8564	-0.0534653\\
4.859	-0.0537728\\
4.8616	-0.054082\\
4.8641	-0.0543807\\
4.8667	-0.0546926\\
4.8692	-0.0549937\\
4.8718	-0.0553079\\
4.8744	-0.055623\\
4.8769	-0.0559267\\
4.8795	-0.0562432\\
4.8821	-0.0565601\\
4.8846	-0.0568651\\
4.8872	-0.0571825\\
4.8898	-0.0574999\\
4.8923	-0.057805\\
4.8949	-0.0581221\\
4.8974	-0.0584265\\
4.9	-0.0587426\\
4.9026	-0.0590579\\
4.9051	-0.0593604\\
4.9077	-0.0596739\\
4.9103	-0.0599863\\
4.9128	-0.0602855\\
4.9154	-0.0605952\\
4.918	-0.0609033\\
4.9205	-0.061198\\
4.9231	-0.0615026\\
4.9257	-0.0618052\\
4.9282	-0.0620941\\
4.9308	-0.0623924\\
4.9333	-0.0626769\\
4.9359	-0.0629702\\
4.9385	-0.0632609\\
4.941	-0.0635377\\
4.9436	-0.0638227\\
4.9462	-0.0641045\\
4.9487	-0.0643725\\
4.9513	-0.0646478\\
4.9539	-0.0649197\\
4.9564	-0.0651777\\
4.959	-0.0654424\\
4.9615	-0.0656932\\
4.9641	-0.0659502\\
4.9667	-0.066203\\
4.9692	-0.0664421\\
4.9718	-0.0666865\\
4.9744	-0.0669265\\
4.9769	-0.0671529\\
4.9795	-0.0673838\\
4.9821	-0.0676098\\
4.9846	-0.0678226\\
4.9872	-0.0680389\\
4.9897	-0.0682421\\
4.9923	-0.0684483\\
4.9949	-0.0686493\\
4.9974	-0.0688373\\
5	-0.0690276\\
};
\addplot [color=mycolor10, line width=1.0pt, forget plot]
  table[row sep=crcr]{%
-0.125	-7.83297e-08\\
-0.1224	-7.98995e-08\\
-0.1199	-8.14078e-08\\
-0.1173	-8.29753e-08\\
-0.1147	-8.45416e-08\\
-0.1122	-8.60465e-08\\
-0.1096	-8.76105e-08\\
-0.1071	-8.91133e-08\\
-0.1045	-9.06751e-08\\
-0.1019	-9.22357e-08\\
-0.0994	-9.37353e-08\\
-0.0968	-9.52937e-08\\
-0.0942	-9.6851e-08\\
-0.0917	-9.83474e-08\\
-0.0891	-9.99025e-08\\
-0.0865	-1.01457e-07\\
-0.084	-1.0295e-07\\
-0.0814	-1.04502e-07\\
-0.0789	-1.05993e-07\\
-0.0763	-1.07543e-07\\
-0.0737	-1.09092e-07\\
-0.0712	-1.10581e-07\\
-0.0686	-1.12128e-07\\
-0.066	-1.13674e-07\\
-0.0635	-1.15159e-07\\
-0.0609	-1.16703e-07\\
-0.0583	-1.18246e-07\\
-0.0558	-1.19729e-07\\
-0.0532	-1.21271e-07\\
-0.0507	-1.22752e-07\\
-0.0481	-1.24292e-07\\
-0.0455	-1.2583e-07\\
-0.043	-1.27309e-07\\
-0.0404	-1.28846e-07\\
-0.0378	-1.30382e-07\\
-0.0353	-1.31859e-07\\
-0.0327	-1.33393e-07\\
-0.0301	-1.34927e-07\\
-0.0276	-1.36402e-07\\
-0.025	-1.37934e-07\\
-0.0224	-1.39466e-07\\
-0.0199	-1.40938e-07\\
-0.0173	-1.42468e-07\\
-0.0148	-1.43939e-07\\
-0.0122	-1.45468e-07\\
-0.0096	-1.46996e-07\\
-0.0071	-1.48465e-07\\
-0.0045	-1.49992e-07\\
-0.0019	-1.51519e-07\\
0.0006	-1.52986e-07\\
0.0032	0.000922952\\
0.0058	0.00139095\\
0.0083	0.00144641\\
0.0109	0.00136049\\
0.0134	0.00124213\\
0.016	0.00109237\\
0.0186	0.000883349\\
0.0211	0.000609611\\
0.0237	0.000265291\\
0.0263	-0.000105044\\
0.0288	-0.000460168\\
0.0314	-0.000822871\\
0.034	-0.00118906\\
0.0365	-0.00155537\\
0.0391	-0.00195339\\
0.0416	-0.00234581\\
0.0442	-0.00275443\\
0.0468	-0.00315823\\
0.0493	-0.00354306\\
0.0519	-0.00394402\\
0.0545	-0.00434849\\
0.057	-0.00473939\\
0.0596	-0.00514368\\
0.0622	-0.00554158\\
0.0647	-0.00591698\\
0.0673	-0.00630146\\
0.0698	-0.0066674\\
0.0724	-0.0070444\\
0.075	-0.00741555\\
0.0775	-0.00776389\\
0.0801	-0.0081153\\
0.0827	-0.0084559\\
0.0852	-0.00877479\\
0.0878	-0.00909866\\
0.0904	-0.00941386\\
0.0929	-0.0097066\\
0.0955	-0.00999809\\
0.098	-0.0102654\\
0.1006	-0.0105309\\
0.1032	-0.0107856\\
0.1057	-0.0110206\\
0.1083	-0.0112535\\
0.1109	-0.0114726\\
0.1134	-0.0116691\\
0.116	-0.0118594\\
0.1186	-0.0120369\\
0.1211	-0.0121969\\
0.1237	-0.0123518\\
0.1263	-0.0124937\\
0.1288	-0.0126163\\
0.1314	-0.0127294\\
0.1339	-0.0128256\\
0.1365	-0.0129141\\
0.1391	-0.0129917\\
0.1416	-0.0130554\\
0.1442	-0.0131088\\
0.1468	-0.0131484\\
0.1493	-0.0131742\\
0.1519	-0.01319\\
0.1545	-0.0131958\\
0.157	-0.0131921\\
0.1596	-0.0131775\\
0.1621	-0.0131522\\
0.1647	-0.0131144\\
0.1673	-0.0130659\\
0.1698	-0.013011\\
0.1724	-0.0129459\\
0.175	-0.0128722\\
0.1775	-0.0127923\\
0.1801	-0.0126993\\
0.1827	-0.0125969\\
0.1852	-0.0124912\\
0.1878	-0.0123747\\
0.1903	-0.0122567\\
0.1929	-0.0121269\\
0.1955	-0.0119891\\
0.198	-0.0118493\\
0.2006	-0.0116972\\
0.2032	-0.0115398\\
0.2057	-0.0113841\\
0.2083	-0.0112173\\
0.2109	-0.0110448\\
0.2134	-0.0108731\\
0.216	-0.010689\\
0.2185	-0.010508\\
0.2211	-0.0103166\\
0.2237	-0.0101221\\
0.2262	-0.00993168\\
0.2288	-0.00972934\\
0.2314	-0.00952272\\
0.2339	-0.0093209\\
0.2365	-0.00910889\\
0.2391	-0.00889532\\
0.2416	-0.00868823\\
0.2442	-0.00847037\\
0.2467	-0.00825817\\
0.2493	-0.00803515\\
0.2519	-0.0078108\\
0.2544	-0.00759468\\
0.257	-0.00736963\\
0.2596	-0.00714377\\
0.2621	-0.0069253\\
0.2647	-0.00669679\\
0.2673	-0.00646775\\
0.2698	-0.00624793\\
0.2724	-0.00602025\\
0.2749	-0.00580203\\
0.2775	-0.00557523\\
0.2801	-0.00534831\\
0.2826	-0.00513038\\
0.2852	-0.00490488\\
0.2878	-0.0046813\\
0.2903	-0.00446825\\
0.2929	-0.00424829\\
0.2955	-0.00402949\\
0.298	-0.00382026\\
0.3006	-0.00360449\\
0.3032	-0.00339142\\
0.3057	-0.0031895\\
0.3083	-0.0029824\\
0.3108	-0.00278562\\
0.3134	-0.00258318\\
0.316	-0.00238333\\
0.3185	-0.00219431\\
0.3211	-0.00200162\\
0.3237	-0.00181301\\
0.3262	-0.00163516\\
0.3288	-0.00145348\\
0.3314	-0.00127516\\
0.3339	-0.00110737\\
0.3365	-0.000937366\\
0.339	-0.00077855\\
0.3416	-0.000618071\\
0.3442	-0.000461959\\
0.3467	-0.000315761\\
0.3493	-0.000168031\\
0.3519	-2.52897e-05\\
0.3544	0.000106782\\
0.357	0.000238682\\
0.3596	0.000365352\\
0.3621	0.000482562\\
0.3647	0.000599707\\
0.3672	0.000707407\\
0.3698	0.000813766\\
0.3724	0.00091411\\
0.3749	0.00100509\\
0.3775	0.00109437\\
0.3801	0.00117846\\
0.3826	0.00125427\\
0.3852	0.00132746\\
0.3878	0.00139452\\
0.3903	0.00145323\\
0.3929	0.00150863\\
0.3954	0.00155682\\
0.398	0.00160173\\
0.4006	0.00164109\\
0.4031	0.00167338\\
0.4057	0.00170103\\
0.4083	0.00172287\\
0.4108	0.00173879\\
0.4134	0.00175035\\
0.416	0.00175678\\
0.4185	0.00175787\\
0.4211	0.0017535\\
0.4236	0.0017441\\
0.4262	0.0017293\\
0.4288	0.00170977\\
0.4313	0.00168671\\
0.4339	0.00165817\\
0.4365	0.00162477\\
0.439	0.00158805\\
0.4416	0.00154537\\
0.4442	0.00149855\\
0.4467	0.00144994\\
0.4493	0.00139573\\
0.4519	0.00133767\\
0.4544	0.00127813\\
0.457	0.00121251\\
0.4595	0.00114626\\
0.4621	0.0010745\\
0.4647	0.00100007\\
0.4672	0.00092599\\
0.4698	0.000846244\\
0.4724	0.000763823\\
0.4749	0.00068236\\
0.4775	0.000595791\\
0.4801	0.000507711\\
0.4826	0.000421731\\
0.4852	0.000330949\\
0.4877	0.000242362\\
0.4903	0.000149102\\
0.4929	5.51147e-05\\
0.4954	-3.5522e-05\\
0.498	-0.000129785\\
0.5006	-0.000223975\\
0.5031	-0.00031447\\
0.5057	-0.00040838\\
0.5083	-0.000501739\\
0.5108	-0.000590557\\
0.5134	-0.000681571\\
0.5159	-0.000767608\\
0.5185	-0.000855511\\
0.5211	-0.00094174\\
0.5236	-0.00102284\\
0.5262	-0.00110491\\
0.5288	-0.00118423\\
0.5313	-0.00125764\\
0.5339	-0.00133091\\
0.5365	-0.00140098\\
0.539	-0.00146521\\
0.5416	-0.00152842\\
0.5441	-0.00158535\\
0.5467	-0.00164021\\
0.5493	-0.00169046\\
0.5518	-0.00173435\\
0.5544	-0.00177533\\
0.557	-0.00181132\\
0.5595	-0.00184088\\
0.5621	-0.00186599\\
0.5647	-0.00188512\\
0.5672	-0.00189776\\
0.5698	-0.00190488\\
0.5723	-0.00190582\\
0.5749	-0.00190035\\
0.5775	-0.00188798\\
0.58	-0.00186929\\
0.5826	-0.00184262\\
0.5852	-0.00180853\\
0.5877	-0.0017687\\
0.5903	-0.00171976\\
0.5929	-0.00166285\\
0.5954	-0.00160034\\
0.598	-0.00152703\\
0.6006	-0.00144518\\
0.6031	-0.00135841\\
0.6057	-0.00125968\\
0.6082	-0.00115639\\
0.6108	-0.00104001\\
0.6134	-0.000914286\\
0.6159	-0.000784474\\
0.6185	-0.000640186\\
0.6211	-0.000486413\\
0.6236	-0.000329501\\
0.6262	-0.000156695\\
0.6288	2.6132e-05\\
0.6313	0.000211496\\
0.6339	0.00041424\\
0.6364	0.000618741\\
0.639	0.000841384\\
0.6416	0.00107431\\
0.6441	0.00130813\\
0.6467	0.00156168\\
0.6493	0.00182585\\
0.6518	0.00208982\\
0.6544	0.00237466\\
0.657	0.00267004\\
0.6595	0.00296411\\
0.6621	0.00328049\\
0.6646	0.0035949\\
0.6672	0.00393242\\
0.6698	0.00428058\\
0.6723	0.00462534\\
0.6749	0.00499431\\
0.6775	0.00537396\\
0.68	0.00574912\\
0.6826	0.00614973\\
0.6852	0.00656088\\
0.6877	0.00696603\\
0.6903	0.00739754\\
0.6928	0.00782222\\
0.6954	0.00827404\\
0.698	0.00873617\\
0.7005	0.00919011\\
0.7031	0.009672\\
0.7057	0.0101637\\
0.7082	0.0106458\\
0.7108	0.0111567\\
0.7134	0.0116772\\
0.7159	0.0121868\\
0.7185	0.0127258\\
0.721	0.0132527\\
0.7236	0.0138095\\
0.7262	0.0143752\\
0.7287	0.0149274\\
0.7313	0.0155101\\
0.7339	0.0161013\\
0.7364	0.0166774\\
0.739	0.0172843\\
0.7416	0.017899\\
0.7441	0.0184973\\
0.7467	0.0191269\\
0.7492	0.0197391\\
0.7518	0.0203826\\
0.7544	0.0210329\\
0.7569	0.0216642\\
0.7595	0.022327\\
0.7621	0.0229958\\
0.7646	0.0236445\\
0.7672	0.0243246\\
0.7698	0.02501\\
0.7723	0.0256737\\
0.7749	0.0263687\\
0.7775	0.0270681\\
0.78	0.0277447\\
0.7826	0.0284523\\
0.7851	0.0291361\\
0.7877	0.0298505\\
0.7903	0.030568\\
0.7928	0.0312605\\
0.7954	0.0319831\\
0.798	0.032708\\
0.8005	0.0334068\\
0.8031	0.034135\\
0.8057	0.0348644\\
0.8082	0.0355666\\
0.8108	0.0362975\\
0.8133	0.0370006\\
0.8159	0.0377317\\
0.8185	0.0384625\\
0.821	0.0391644\\
0.8236	0.0398933\\
0.8262	0.0406208\\
0.8287	0.0413187\\
0.8313	0.0420425\\
0.8339	0.0427638\\
0.8364	0.0434548\\
0.839	0.0441702\\
0.8415	0.0448548\\
0.8441	0.045563\\
0.8467	0.0462671\\
0.8492	0.0469398\\
0.8518	0.0476345\\
0.8544	0.048324\\
0.8569	0.0489815\\
0.8595	0.0496595\\
0.8621	0.0503311\\
0.8646	0.0509706\\
0.8672	0.0516287\\
0.8697	0.0522545\\
0.8723	0.0528976\\
0.8749	0.0535325\\
0.8774	0.0541349\\
0.88	0.0547527\\
0.8826	0.0553612\\
0.8851	0.0559373\\
0.8877	0.0565266\\
0.8903	0.0571054\\
0.8928	0.0576519\\
0.8954	0.0582095\\
0.8979	0.0587348\\
0.9005	0.0592696\\
0.9031	0.0597922\\
0.9056	0.0602829\\
0.9082	0.0607807\\
0.9108	0.0612652\\
0.9133	0.0617183\\
0.9159	0.062176\\
0.9185	0.0626193\\
0.921	0.0630319\\
0.9236	0.0634463\\
0.9262	0.0638454\\
0.9287	0.0642144\\
0.9313	0.0645826\\
0.9338	0.0649212\\
0.9364	0.0652572\\
0.939	0.0655761\\
0.9415	0.0658665\\
0.9441	0.0661513\\
0.9467	0.0664181\\
0.9492	0.0666576\\
0.9518	0.0668884\\
0.9544	0.0671004\\
0.9569	0.0672862\\
0.9595	0.0674603\\
0.962	0.0676091\\
0.9646	0.0677442\\
0.9672	0.0678589\\
0.9697	0.0679498\\
0.9723	0.0680238\\
0.9749	0.0680766\\
0.9774	0.0681072\\
0.98	0.0681176\\
0.9826	0.0681061\\
0.9851	0.068074\\
0.9877	0.0680185\\
0.9902	0.0679436\\
0.9928	0.0678431\\
0.9954	0.0677194\\
0.9979	0.0675783\\
1.0005	0.0674082\\
1.0031	0.0672142\\
1.0056	0.0670047\\
1.0082	0.066763\\
1.0108	0.0664966\\
1.0133	0.066217\\
1.0159	0.0659016\\
1.0184	0.0655746\\
1.021	0.0652096\\
1.0236	0.0648189\\
1.0261	0.0644188\\
1.0287	0.0639774\\
1.0313	0.0635097\\
1.0338	0.0630353\\
1.0364	0.0625159\\
1.039	0.0619699\\
1.0415	0.0614196\\
1.0441	0.060821\\
1.0466	0.0602199\\
1.0492	0.0595681\\
1.0518	0.0588891\\
1.0543	0.0582104\\
1.0569	0.0574776\\
1.0595	0.0567173\\
1.062	0.0559601\\
1.0646	0.0551455\\
1.0672	0.054303\\
1.0697	0.0534668\\
1.0723	0.0525697\\
1.0748	0.0516807\\
1.0774	0.0507287\\
1.08	0.0497488\\
1.0825	0.0487801\\
1.0851	0.0477451\\
1.0877	0.0466821\\
1.0902	0.0456335\\
1.0928	0.0445155\\
1.0954	0.0433695\\
1.0979	0.0422413\\
1.1005	0.0410405\\
1.1031	0.0398118\\
1.1056	0.0386042\\
1.1082	0.0373211\\
1.1107	0.0360613\\
1.1133	0.034724\\
1.1159	0.0333593\\
1.1184	0.0320213\\
1.121	0.0306032\\
1.1236	0.0291579\\
1.1261	0.0277428\\
1.1287	0.0262449\\
1.1313	0.0247204\\
1.1338	0.0232296\\
1.1364	0.0216534\\
1.1389	0.0201133\\
1.1415	0.0184863\\
1.1441	0.0168336\\
1.1466	0.0152207\\
1.1492	0.0135186\\
1.1518	0.0117916\\
1.1543	0.0101079\\
1.1569	0.00833299\\
1.1595	0.0065341\\
1.162	0.00478206\\
1.1646	0.00293704\\
1.1671	0.00114127\\
1.1697	-0.000748595\\
1.1723	-0.00266079\\
1.1748	-0.00452015\\
1.1774	-0.00647504\\
1.18	-0.0084511\\
1.1825	-0.0103707\\
1.1851	-0.0123871\\
1.1877	-0.0144234\\
1.1902	-0.0163997\\
1.1928	-0.0184737\\
1.1953	-0.0204854\\
1.1979	-0.0225953\\
1.2005	-0.0247227\\
1.203	-0.0267843\\
1.2056	-0.0289446\\
1.2082	-0.0311209\\
1.2107	-0.033228\\
1.2133	-0.0354339\\
1.2159	-0.0376541\\
1.2184	-0.0398019\\
1.221	-0.0420484\\
1.2235	-0.0442202\\
1.2261	-0.0464905\\
1.2287	-0.0487721\\
1.2312	-0.0509759\\
1.2338	-0.0532775\\
1.2364	-0.0555884\\
1.2389	-0.0578184\\
1.2415	-0.0601454\\
1.2441	-0.0624794\\
1.2466	-0.0647298\\
1.2492	-0.0670758\\
1.2518	-0.0694267\\
1.2543	-0.0716912\\
1.2569	-0.0740497\\
1.2594	-0.07632\\
1.262	-0.0786829\\
1.2646	-0.081047\\
1.2671	-0.0833205\\
1.2697	-0.0856845\\
1.2723	-0.0880471\\
1.2748	-0.0903169\\
1.2774	-0.0926747\\
1.28	-0.0950286\\
1.2825	-0.0972877\\
1.2851	-0.0996318\\
1.2876	-0.10188\\
1.2902	-0.104211\\
1.2928	-0.106534\\
1.2953	-0.108759\\
1.2979	-0.111063\\
1.3005	-0.113357\\
1.303	-0.115552\\
1.3056	-0.117823\\
1.3082	-0.12008\\
1.3107	-0.122237\\
1.3133	-0.124466\\
1.3158	-0.126593\\
1.3184	-0.12879\\
1.321	-0.130969\\
1.3235	-0.133046\\
1.3261	-0.135188\\
1.3287	-0.137309\\
1.3312	-0.139329\\
1.3338	-0.141407\\
1.3364	-0.143463\\
1.3389	-0.145417\\
1.3415	-0.147426\\
1.344	-0.149333\\
1.3466	-0.15129\\
1.3492	-0.15322\\
1.3517	-0.15505\\
1.3543	-0.156924\\
1.3569	-0.158769\\
1.3594	-0.160514\\
1.362	-0.162298\\
1.3646	-0.16405\\
1.3671	-0.165704\\
1.3697	-0.167391\\
1.3722	-0.16898\\
1.3748	-0.170599\\
1.3774	-0.172182\\
1.3799	-0.17367\\
1.3825	-0.175181\\
1.3851	-0.176654\\
1.3876	-0.178034\\
1.3902	-0.17943\\
1.3928	-0.180787\\
1.3953	-0.182054\\
1.3979	-0.183331\\
1.4005	-0.184567\\
1.403	-0.185715\\
1.4056	-0.186867\\
1.4081	-0.187934\\
1.4107	-0.189001\\
1.4133	-0.190024\\
1.4158	-0.190966\\
1.4184	-0.1919\\
1.421	-0.19279\\
1.4235	-0.193602\\
1.4261	-0.194401\\
1.4287	-0.195153\\
1.4312	-0.195832\\
1.4338	-0.196492\\
1.4363	-0.197082\\
1.4389	-0.197648\\
1.4415	-0.198166\\
1.444	-0.198619\\
1.4466	-0.199043\\
1.4492	-0.199418\\
1.4517	-0.199732\\
1.4543	-0.200011\\
1.4569	-0.200242\\
1.4594	-0.200417\\
1.462	-0.200551\\
1.4645	-0.200633\\
1.4671	-0.20067\\
1.4697	-0.200659\\
1.4722	-0.200601\\
1.4748	-0.200493\\
1.4774	-0.200337\\
1.4799	-0.20014\\
1.4825	-0.199888\\
1.4851	-0.199587\\
1.4876	-0.199252\\
1.4902	-0.198857\\
1.4927	-0.198431\\
1.4953	-0.197942\\
1.4979	-0.197406\\
1.5004	-0.196845\\
1.503	-0.196217\\
1.5056	-0.195542\\
1.5081	-0.19485\\
1.5107	-0.194085\\
1.5133	-0.193274\\
1.5158	-0.192453\\
1.5184	-0.191555\\
1.5209	-0.19065\\
1.5235	-0.189666\\
1.5261	-0.188639\\
1.5286	-0.187611\\
1.5312	-0.186502\\
1.5338	-0.18535\\
1.5363	-0.184204\\
1.5389	-0.182973\\
1.5415	-0.181702\\
1.544	-0.180443\\
1.5466	-0.179096\\
1.5491	-0.177764\\
1.5517	-0.176343\\
1.5543	-0.174885\\
1.5568	-0.173449\\
1.5594	-0.171922\\
1.562	-0.170359\\
1.5645	-0.168825\\
1.5671	-0.167196\\
1.5697	-0.165535\\
1.5722	-0.163908\\
1.5748	-0.162186\\
1.5774	-0.160433\\
1.5799	-0.158719\\
1.5825	-0.156909\\
1.585	-0.155142\\
1.5876	-0.153278\\
1.5902	-0.151387\\
1.5927	-0.149544\\
1.5953	-0.147604\\
1.5979	-0.145639\\
1.6004	-0.143728\\
1.603	-0.141718\\
1.6056	-0.139687\\
1.6081	-0.137714\\
1.6107	-0.135642\\
1.6132	-0.133632\\
1.6158	-0.131524\\
1.6184	-0.129398\\
1.6209	-0.127337\\
1.6235	-0.125179\\
1.6261	-0.123006\\
1.6286	-0.120902\\
1.6312	-0.118701\\
1.6338	-0.116487\\
1.6363	-0.114348\\
1.6389	-0.112112\\
1.6414	-0.109952\\
1.644	-0.107696\\
1.6466	-0.105432\\
1.6491	-0.103247\\
1.6517	-0.100968\\
1.6543	-0.0986828\\
1.6568	-0.0964803\\
1.6594	-0.0941851\\
1.662	-0.091886\\
1.6645	-0.0896724\\
1.6671	-0.0873679\\
1.6696	-0.0851505\\
1.6722	-0.0828437\\
1.6748	-0.0805368\\
1.6773	-0.0783193\\
1.6799	-0.0760145\\
1.6825	-0.0737119\\
1.685	-0.0715006\\
1.6876	-0.0692045\\
1.6902	-0.0669127\\
1.6927	-0.0647138\\
1.6953	-0.0624327\\
1.6978	-0.0602455\\
1.7004	-0.0579778\\
1.703	-0.055718\\
1.7055	-0.0535532\\
1.7081	-0.0513108\\
1.7107	-0.0490783\\
1.7132	-0.0469417\\
1.7158	-0.0447305\\
1.7184	-0.0425312\\
1.7209	-0.0404282\\
1.7235	-0.038254\\
1.7261	-0.0360934\\
1.7286	-0.0340294\\
1.7312	-0.0318975\\
1.7337	-0.0298622\\
1.7363	-0.0277613\\
1.7389	-0.025677\\
1.7414	-0.023689\\
1.744	-0.021639\\
1.7466	-0.0196073\\
1.7491	-0.0176715\\
1.7517	-0.0156772\\
1.7543	-0.0137029\\
1.7568	-0.0118237\\
1.7594	-0.00988982\\
1.7619	-0.0080505\\
1.7645	-0.00615904\\
1.7671	-0.00428991\\
1.7696	-0.00251418\\
1.7722	-0.000690223\\
1.7748	0.00111003\\
1.7773	0.00281829\\
1.7799	0.00457078\\
1.7825	0.00629828\\
1.785	0.0079354\\
1.7876	0.0096127\\
1.7901	0.0112008\\
1.7927	0.0128264\\
1.7953	0.014425\\
1.7978	0.0159364\\
1.8004	0.0174812\\
1.803	0.0189979\\
1.8055	0.0204296\\
1.8081	0.0218904\\
1.8107	0.0233222\\
1.8132	0.0246713\\
1.8158	0.0260454\\
1.8183	0.0273384\\
1.8209	0.0286535\\
1.8235	0.0299382\\
1.826	0.0311445\\
1.8286	0.0323687\\
1.8312	0.0335617\\
1.8337	0.0346791\\
1.8363	0.0358101\\
1.8389	0.0369091\\
1.8414	0.0379356\\
1.844	0.0389715\\
1.8465	0.0399368\\
1.8491	0.0409086\\
1.8517	0.0418475\\
1.8542	0.0427191\\
1.8568	0.0435929\\
1.8594	0.0444334\\
1.8619	0.0452098\\
1.8645	0.0459843\\
1.8671	0.046725\\
1.8696	0.0474052\\
1.8722	0.0480792\\
1.8747	0.0486951\\
1.8773	0.0493019\\
1.8799	0.0498745\\
1.8824	0.0503926\\
1.885	0.0508975\\
1.8876	0.0513679\\
1.8901	0.0517875\\
1.8927	0.05219\\
1.8953	0.0525577\\
1.8978	0.0528785\\
1.9004	0.0531781\\
1.903	0.0534429\\
1.9055	0.0536647\\
1.9081	0.0538613\\
1.9106	0.0540175\\
1.9132	0.0541459\\
1.9158	0.0542396\\
1.9183	0.0542969\\
1.9209	0.0543226\\
1.9235	0.0543136\\
1.926	0.0542724\\
1.9286	0.0541958\\
1.9312	0.0540847\\
1.9337	0.0539456\\
1.9363	0.0537673\\
1.9388	0.0535637\\
1.9414	0.0533186\\
1.944	0.0530395\\
1.9465	0.0527394\\
1.9491	0.0523942\\
1.9517	0.0520155\\
1.9542	0.0516199\\
1.9568	0.0511759\\
1.9594	0.0506989\\
1.9619	0.0502092\\
1.9645	0.0496679\\
1.967	0.0491167\\
1.9696	0.0485118\\
1.9722	0.0478748\\
1.9747	0.0472322\\
1.9773	0.0465329\\
1.9799	0.0458021\\
1.9824	0.04507\\
1.985	0.0442782\\
1.9876	0.0434558\\
1.9901	0.0426363\\
1.9927	0.0417545\\
1.9952	0.0408784\\
1.9978	0.0399382\\
2.0004	0.0389687\\
2.0029	0.0380091\\
2.0055	0.0369831\\
2.0081	0.0359287\\
2.0106	0.0348884\\
2.0132	0.0337793\\
2.0158	0.0326429\\
2.0183	0.0315248\\
2.0209	0.0303358\\
2.0234	0.0291678\\
2.026	0.0279276\\
2.0286	0.026662\\
2.0311	0.0254214\\
2.0337	0.0241069\\
2.0363	0.0227682\\
2.0388	0.0214584\\
2.0414	0.0200733\\
2.044	0.0186652\\
2.0465	0.0172899\\
2.0491	0.015838\\
2.0517	0.0143644\\
2.0542	0.0129276\\
2.0568	0.0114129\\
2.0593	0.00993741\\
2.0619	0.00838355\\
2.0645	0.00681045\\
2.067	0.00528018\\
2.0696	0.00367082\\
2.0722	0.00204375\\
2.0747	0.000463056\\
2.0773	-0.00119719\\
2.0799	-0.00287356\\
2.0824	-0.00450012\\
2.085	-0.00620647\\
2.0875	-0.00786083\\
2.0901	-0.00959499\\
2.0927	-0.0113425\\
2.0952	-0.0130348\\
2.0978	-0.0148067\\
2.1004	-0.0165901\\
2.1029	-0.0183154\\
2.1055	-0.0201198\\
2.1081	-0.021934\\
2.1106	-0.023687\\
2.1132	-0.0255186\\
2.1157	-0.0272872\\
2.1183	-0.0291337\\
2.1209	-0.0309869\\
2.1234	-0.0327745\\
2.126	-0.0346389\\
2.1286	-0.0365081\\
2.1311	-0.0383092\\
2.1337	-0.0401859\\
2.1363	-0.0420653\\
2.1388	-0.0438745\\
2.1414	-0.0457576\\
2.1439	-0.0475691\\
2.1465	-0.0494532\\
2.1491	-0.0513369\\
2.1516	-0.0531471\\
2.1542	-0.0550279\\
2.1568	-0.0569064\\
2.1593	-0.0587096\\
2.1619	-0.0605813\\
2.1645	-0.0624485\\
2.167	-0.0642391\\
2.1696	-0.0660957\\
2.1721	-0.0678748\\
2.1747	-0.0697182\\
2.1773	-0.0715538\\
2.1798	-0.0733109\\
2.1824	-0.0751294\\
2.185	-0.0769382\\
2.1875	-0.0786677\\
2.1901	-0.0804556\\
2.1927	-0.0822318\\
2.1952	-0.0839282\\
2.1978	-0.0856797\\
2.2004	-0.0874177\\
2.2029	-0.0890754\\
2.2055	-0.090785\\
2.208	-0.0924143\\
2.2106	-0.0940931\\
2.2132	-0.0957552\\
2.2157	-0.0973372\\
2.2183	-0.0989651\\
2.2209	-0.100575\\
2.2234	-0.102104\\
2.226	-0.103676\\
2.2286	-0.105228\\
2.2311	-0.1067\\
2.2337	-0.108211\\
2.2362	-0.109643\\
2.2388	-0.111111\\
2.2414	-0.112556\\
2.2439	-0.113923\\
2.2465	-0.115322\\
2.2491	-0.116696\\
2.2516	-0.117995\\
2.2542	-0.119321\\
2.2568	-0.120621\\
2.2593	-0.121847\\
2.2619	-0.123096\\
2.2644	-0.124271\\
2.267	-0.125467\\
2.2696	-0.126635\\
2.2721	-0.127731\\
2.2747	-0.128844\\
2.2773	-0.129928\\
2.2798	-0.130943\\
2.2824	-0.13197\\
2.285	-0.132967\\
2.2875	-0.133897\\
2.2901	-0.134835\\
2.2926	-0.135708\\
2.2952	-0.136586\\
2.2978	-0.137433\\
2.3003	-0.138218\\
2.3029	-0.139003\\
2.3055	-0.139756\\
2.308	-0.140451\\
2.3106	-0.141142\\
2.3132	-0.141801\\
2.3157	-0.142404\\
2.3183	-0.142999\\
2.3208	-0.143541\\
2.3234	-0.144073\\
2.326	-0.144572\\
2.3285	-0.145021\\
2.3311	-0.145455\\
2.3337	-0.145857\\
2.3362	-0.146213\\
2.3388	-0.14655\\
2.3414	-0.146855\\
2.3439	-0.147117\\
2.3465	-0.147357\\
2.349	-0.147558\\
2.3516	-0.147734\\
2.3542	-0.147878\\
2.3567	-0.147986\\
2.3593	-0.148066\\
2.3619	-0.148115\\
2.3644	-0.148131\\
2.367	-0.148116\\
2.3696	-0.14807\\
2.3721	-0.147996\\
2.3747	-0.147889\\
2.3773	-0.14775\\
2.3798	-0.147587\\
2.3824	-0.147388\\
2.3849	-0.147168\\
2.3875	-0.146909\\
2.3901	-0.146621\\
2.3926	-0.146315\\
2.3952	-0.145969\\
2.3978	-0.145594\\
2.4003	-0.145206\\
2.4029	-0.144775\\
2.4055	-0.144315\\
2.408	-0.143848\\
2.4106	-0.143334\\
2.4131	-0.142816\\
2.4157	-0.14225\\
2.4183	-0.141658\\
2.4208	-0.141064\\
2.4234	-0.140422\\
2.426	-0.139755\\
2.4285	-0.13909\\
2.4311	-0.138376\\
2.4337	-0.137637\\
2.4362	-0.136905\\
2.4388	-0.136121\\
2.4413	-0.135346\\
2.4439	-0.134519\\
2.4465	-0.13367\\
2.449	-0.132834\\
2.4516	-0.131945\\
2.4542	-0.131035\\
2.4567	-0.130142\\
2.4593	-0.129195\\
2.4619	-0.128229\\
2.4644	-0.127283\\
2.467	-0.126283\\
2.4695	-0.125304\\
2.4721	-0.124271\\
2.4747	-0.123221\\
2.4772	-0.122197\\
2.4798	-0.121117\\
2.4824	-0.120023\\
2.4849	-0.118958\\
2.4875	-0.117837\\
2.4901	-0.116703\\
2.4926	-0.115601\\
2.4952	-0.114443\\
2.4977	-0.11332\\
2.5003	-0.11214\\
2.5029	-0.110951\\
2.5054	-0.109798\\
2.508	-0.108589\\
2.5106	-0.107372\\
2.5131	-0.106195\\
2.5157	-0.104962\\
2.5183	-0.103723\\
2.5208	-0.102525\\
2.5234	-0.101273\\
2.526	-0.100015\\
2.5285	-0.0988011\\
2.5311	-0.0975341\\
2.5336	-0.0963121\\
2.5362	-0.0950377\\
2.5388	-0.0937603\\
2.5413	-0.0925297\\
2.5439	-0.0912478\\
2.5465	-0.0899644\\
2.549	-0.0887294\\
2.5516	-0.0874445\\
2.5542	-0.0861595\\
2.5567	-0.0849244\\
2.5593	-0.0836407\\
2.5618	-0.0824076\\
2.5644	-0.081127\\
2.567	-0.0798488\\
2.5695	-0.0786222\\
2.5721	-0.0773498\\
2.5747	-0.076081\\
2.5772	-0.0748648\\
2.5798	-0.0736045\\
2.5824	-0.0723491\\
2.5849	-0.0711471\\
2.5875	-0.0699027\\
2.59	-0.068712\\
2.5926	-0.0674802\\
2.5952	-0.0662554\\
2.5977	-0.0650848\\
2.6003	-0.063875\\
2.6029	-0.0626734\\
2.6054	-0.0615261\\
2.608	-0.0603418\\
2.6106	-0.0591668\\
2.6131	-0.0580461\\
2.6157	-0.0568904\\
2.6182	-0.055789\\
2.6208	-0.054654\\
2.6234	-0.0535302\\
2.6259	-0.0524603\\
2.6285	-0.0513591\\
2.6311	-0.0502699\\
2.6336	-0.0492343\\
2.6362	-0.0481697\\
2.6388	-0.047118\\
2.6413	-0.0461192\\
2.6439	-0.0450938\\
2.6464	-0.0441208\\
2.649	-0.0431226\\
2.6516	-0.0421387\\
2.6541	-0.0412065\\
2.6567	-0.0402514\\
2.6593	-0.0393114\\
2.6618	-0.038422\\
2.6644	-0.0375122\\
2.667	-0.0366182\\
2.6695	-0.0357735\\
2.6721	-0.0349109\\
2.6746	-0.0340969\\
2.6772	-0.0332666\\
2.6798	-0.0324529\\
2.6823	-0.0316864\\
2.6849	-0.0309061\\
2.6875	-0.0301429\\
2.69	-0.0294255\\
2.6926	-0.0286965\\
2.6952	-0.0279852\\
2.6977	-0.027318\\
2.7003	-0.0266417\\
2.7029	-0.0259835\\
2.7054	-0.0253676\\
2.708	-0.0247451\\
2.7105	-0.0241637\\
2.7131	-0.0235772\\
2.7157	-0.0230093\\
2.7182	-0.0224806\\
2.7208	-0.0219492\\
2.7234	-0.0214365\\
2.7259	-0.0209612\\
2.7285	-0.0204854\\
2.7311	-0.0200285\\
2.7336	-0.0196069\\
2.7362	-0.0191871\\
2.7387	-0.0188013\\
2.7413	-0.0184187\\
2.7439	-0.0180551\\
2.7464	-0.0177234\\
2.749	-0.017397\\
2.7516	-0.0170896\\
2.7541	-0.0168119\\
2.7567	-0.0165416\\
2.7593	-0.0162903\\
2.7618	-0.0160664\\
2.7644	-0.0158521\\
2.7669	-0.0156636\\
2.7695	-0.015486\\
2.7721	-0.0153271\\
2.7746	-0.0151918\\
2.7772	-0.0150692\\
2.7798	-0.0149651\\
2.7823	-0.0148823\\
2.7849	-0.0148141\\
2.7875	-0.0147641\\
2.79	-0.014733\\
2.7926	-0.0147182\\
2.7951	-0.0147209\\
2.7977	-0.014741\\
2.8003	-0.0147787\\
2.8028	-0.0148315\\
2.8054	-0.0149033\\
2.808	-0.0149923\\
2.8105	-0.0150939\\
2.8131	-0.0152162\\
2.8157	-0.0153551\\
2.8182	-0.0155043\\
2.8208	-0.0156756\\
2.8233	-0.0158555\\
2.8259	-0.0160584\\
2.8285	-0.0162771\\
2.831	-0.0165022\\
2.8336	-0.0167515\\
2.8362	-0.017016\\
2.8387	-0.0172845\\
2.8413	-0.0175784\\
2.8439	-0.0178869\\
2.8464	-0.0181971\\
2.849	-0.0185337\\
2.8516	-0.0188843\\
2.8541	-0.0192344\\
2.8567	-0.0196117\\
2.8592	-0.019987\\
2.8618	-0.0203902\\
2.8644	-0.0208062\\
2.8669	-0.021218\\
2.8695	-0.0216584\\
2.8721	-0.0221108\\
2.8746	-0.0225569\\
2.8772	-0.0230322\\
2.8798	-0.0235188\\
2.8823	-0.023997\\
2.8849	-0.0245049\\
2.8874	-0.0250031\\
2.89	-0.0255313\\
2.8926	-0.0260693\\
2.8951	-0.0265957\\
2.8977	-0.0271524\\
2.9003	-0.0277181\\
2.9028	-0.0282703\\
2.9054	-0.028853\\
2.908	-0.0294438\\
2.9105	-0.0300193\\
2.9131	-0.0306253\\
2.9156	-0.0312149\\
2.9182	-0.0318349\\
2.9208	-0.0324617\\
2.9233	-0.0330703\\
2.9259	-0.0337093\\
2.9285	-0.0343541\\
2.931	-0.0349793\\
2.9336	-0.0356345\\
2.9362	-0.0362947\\
2.9387	-0.0369337\\
2.9413	-0.0376025\\
2.9438	-0.0382493\\
2.9464	-0.0389255\\
2.949	-0.0396051\\
2.9515	-0.0402614\\
2.9541	-0.0409466\\
2.9567	-0.0416342\\
2.9592	-0.0422974\\
2.9618	-0.0429889\\
2.9644	-0.0436818\\
2.9669	-0.0443492\\
2.9695	-0.0450442\\
2.972	-0.045713\\
2.9746	-0.0464088\\
2.9772	-0.0471046\\
2.9797	-0.0477733\\
2.9823	-0.0484681\\
2.9849	-0.0491621\\
2.9874	-0.0498282\\
2.99	-0.0505194\\
2.9926	-0.0512089\\
2.9951	-0.0518698\\
2.9977	-0.0525549\\
3.0003	-0.0532373\\
3.0028	-0.0538907\\
3.0054	-0.0545671\\
3.0079	-0.0552141\\
3.0105	-0.0558833\\
3.0131	-0.0565485\\
3.0156	-0.057184\\
3.0182	-0.0578405\\
3.0208	-0.0584921\\
3.0233	-0.0591139\\
3.0259	-0.0597553\\
3.0285	-0.060391\\
3.031	-0.0609969\\
3.0336	-0.0616209\\
3.0361	-0.062215\\
3.0387	-0.0628264\\
3.0413	-0.0634311\\
3.0438	-0.0640058\\
3.0464	-0.0645964\\
3.049	-0.0651796\\
3.0515	-0.0657331\\
3.0541	-0.066301\\
3.0567	-0.0668608\\
3.0592	-0.0673913\\
3.0618	-0.0679346\\
3.0643	-0.0684489\\
3.0669	-0.0689751\\
3.0695	-0.0694922\\
3.072	-0.0699807\\
3.0746	-0.0704796\\
3.0772	-0.0709689\\
3.0797	-0.0714303\\
3.0823	-0.0719004\\
3.0849	-0.0723605\\
3.0874	-0.0727934\\
3.09	-0.0732335\\
3.0925	-0.0736468\\
3.0951	-0.0740663\\
3.0977	-0.0744751\\
3.1002	-0.0748581\\
3.1028	-0.0752457\\
3.1054	-0.0756223\\
3.1079	-0.0759739\\
3.1105	-0.0763287\\
3.1131	-0.0766722\\
3.1156	-0.0769917\\
3.1182	-0.0773129\\
3.1207	-0.0776109\\
3.1233	-0.0779094\\
3.1259	-0.0781964\\
3.1284	-0.0784613\\
3.131	-0.0787253\\
3.1336	-0.0789776\\
3.1361	-0.0792091\\
3.1387	-0.0794382\\
3.1413	-0.0796556\\
3.1438	-0.0798534\\
3.1464	-0.0800475\\
3.1489	-0.080223\\
3.1515	-0.0803938\\
3.1541	-0.0805529\\
3.1566	-0.0806946\\
3.1592	-0.0808304\\
3.1618	-0.0809545\\
3.1643	-0.0810627\\
3.1669	-0.0811637\\
3.1695	-0.081253\\
3.172	-0.081328\\
3.1746	-0.0813945\\
3.1772	-0.0814496\\
3.1797	-0.0814917\\
3.1823	-0.0815243\\
3.1848	-0.081545\\
3.1874	-0.0815555\\
3.19	-0.0815549\\
3.1925	-0.0815438\\
3.1951	-0.0815215\\
3.1977	-0.0814883\\
3.2002	-0.0814463\\
3.2028	-0.0813921\\
3.2054	-0.0813273\\
3.2079	-0.0812551\\
3.2105	-0.08117\\
3.213	-0.0810785\\
3.2156	-0.0809734\\
3.2182	-0.0808584\\
3.2207	-0.0807386\\
3.2233	-0.0806045\\
3.2259	-0.0804609\\
3.2284	-0.0803139\\
3.231	-0.080152\\
3.2336	-0.0799811\\
3.2361	-0.0798083\\
3.2387	-0.0796201\\
3.2412	-0.079431\\
3.2438	-0.079226\\
3.2464	-0.0790128\\
3.2489	-0.0788002\\
3.2515	-0.0785714\\
3.2541	-0.0783348\\
3.2566	-0.0781002\\
3.2592	-0.077849\\
3.2618	-0.0775907\\
3.2643	-0.0773357\\
3.2669	-0.0770638\\
3.2694	-0.0767962\\
3.272	-0.0765116\\
3.2746	-0.0762208\\
3.2771	-0.0759355\\
3.2797	-0.0756331\\
3.2823	-0.0753251\\
3.2848	-0.0750238\\
3.2874	-0.0747054\\
3.29	-0.074382\\
3.2925	-0.0740665\\
3.2951	-0.0737339\\
3.2976	-0.07341\\
3.3002	-0.073069\\
3.3028	-0.0727241\\
3.3053	-0.0723889\\
3.3079	-0.0720368\\
3.3105	-0.0716814\\
3.313	-0.0713367\\
3.3156	-0.0709753\\
3.3182	-0.0706113\\
3.3207	-0.0702589\\
3.3233	-0.0698902\\
3.3259	-0.0695194\\
3.3284	-0.0691611\\
3.331	-0.0687868\\
3.3335	-0.0684256\\
3.3361	-0.0680487\\
3.3387	-0.0676708\\
3.3412	-0.0673066\\
3.3438	-0.0669273\\
3.3464	-0.0665475\\
3.3489	-0.0661821\\
3.3515	-0.0658021\\
3.3541	-0.0654222\\
3.3566	-0.0650573\\
3.3592	-0.0646784\\
3.3617	-0.0643148\\
3.3643	-0.0639375\\
3.3669	-0.0635614\\
3.3694	-0.0632011\\
3.372	-0.0628278\\
3.3746	-0.0624562\\
3.3771	-0.0621006\\
3.3797	-0.0617329\\
3.3823	-0.0613675\\
3.3848	-0.0610183\\
3.3874	-0.0606578\\
3.3899	-0.0603137\\
3.3925	-0.0599587\\
3.3951	-0.0596068\\
3.3976	-0.0592715\\
3.4002	-0.0589262\\
3.4028	-0.0585845\\
3.4053	-0.0582594\\
3.4079	-0.0579251\\
3.4105	-0.0575948\\
3.413	-0.0572812\\
3.4156	-0.0569593\\
3.4181	-0.0566539\\
3.4207	-0.0563408\\
3.4233	-0.0560325\\
3.4258	-0.0557406\\
3.4284	-0.0554419\\
3.431	-0.0551483\\
3.4335	-0.0548709\\
3.4361	-0.0545876\\
3.4387	-0.0543098\\
3.4412	-0.0540479\\
3.4438	-0.0537811\\
3.4463	-0.05353\\
3.4489	-0.0532745\\
3.4515	-0.0530251\\
3.454	-0.0527909\\
3.4566	-0.0525533\\
3.4592	-0.052322\\
3.4617	-0.0521055\\
3.4643	-0.0518866\\
3.4669	-0.0516741\\
3.4694	-0.0514759\\
3.472	-0.0512763\\
3.4745	-0.0510905\\
3.4771	-0.0509039\\
3.4797	-0.050724\\
3.4822	-0.0505574\\
3.4848	-0.0503909\\
3.4874	-0.0502312\\
3.4899	-0.0500842\\
3.4925	-0.0499381\\
3.4951	-0.049799\\
3.4976	-0.0496718\\
3.5002	-0.0495464\\
3.5028	-0.049428\\
3.5053	-0.0493208\\
3.5079	-0.0492162\\
3.5104	-0.0491224\\
3.513	-0.0490316\\
3.5156	-0.048948\\
3.5181	-0.0488741\\
3.5207	-0.0488043\\
3.5233	-0.0487414\\
3.5258	-0.0486876\\
3.5284	-0.0486385\\
3.531	-0.0485963\\
3.5335	-0.0485623\\
3.5361	-0.0485337\\
3.5386	-0.0485127\\
3.5412	-0.0484975\\
3.5438	-0.0484891\\
3.5463	-0.0484874\\
3.5489	-0.0484921\\
3.5515	-0.0485035\\
3.554	-0.0485206\\
3.5566	-0.0485448\\
3.5592	-0.0485755\\
3.5617	-0.048611\\
3.5643	-0.0486541\\
3.5668	-0.0487014\\
3.5694	-0.0487566\\
3.572	-0.0488179\\
3.5745	-0.0488825\\
3.5771	-0.0489555\\
3.5797	-0.0490342\\
3.5822	-0.0491154\\
3.5848	-0.0492053\\
3.5874	-0.0493007\\
3.5899	-0.0493976\\
3.5925	-0.0495036\\
3.595	-0.0496105\\
3.5976	-0.0497266\\
3.6002	-0.0498477\\
3.6027	-0.0499687\\
3.6053	-0.0500993\\
3.6079	-0.0502345\\
3.6104	-0.0503688\\
3.613	-0.0505128\\
3.6156	-0.0506611\\
3.6181	-0.0508075\\
3.6207	-0.0509638\\
3.6232	-0.0511177\\
3.6258	-0.0512815\\
3.6284	-0.0514489\\
3.6309	-0.0516131\\
3.6335	-0.0517872\\
3.6361	-0.0519645\\
3.6386	-0.0521379\\
3.6412	-0.0523211\\
3.6438	-0.052507\\
3.6463	-0.0526883\\
3.6489	-0.0528793\\
3.6515	-0.0530726\\
3.654	-0.0532605\\
3.6566	-0.053458\\
3.6591	-0.0536496\\
3.6617	-0.0538506\\
3.6643	-0.0540532\\
3.6668	-0.0542494\\
3.6694	-0.0544546\\
3.672	-0.054661\\
3.6745	-0.0548603\\
3.6771	-0.0550685\\
3.6797	-0.0552773\\
3.6822	-0.0554785\\
3.6848	-0.0556882\\
3.6873	-0.0558899\\
3.6899	-0.0560998\\
3.6925	-0.0563096\\
3.695	-0.0565111\\
3.6976	-0.0567202\\
3.7002	-0.0569287\\
3.7027	-0.0571286\\
3.7053	-0.0573355\\
3.7079	-0.0575415\\
3.7104	-0.0577384\\
3.713	-0.0579419\\
3.7155	-0.0581362\\
3.7181	-0.0583367\\
3.7207	-0.0585354\\
3.7232	-0.0587246\\
3.7258	-0.0589195\\
3.7284	-0.0591121\\
3.7309	-0.0592951\\
3.7335	-0.059483\\
3.7361	-0.0596683\\
3.7386	-0.0598438\\
3.7412	-0.0600236\\
3.7437	-0.0601936\\
3.7463	-0.0603673\\
3.7489	-0.0605377\\
3.7514	-0.0606984\\
3.754	-0.060862\\
3.7566	-0.0610219\\
3.7591	-0.0611722\\
3.7617	-0.0613246\\
3.7643	-0.0614731\\
3.7668	-0.0616119\\
3.7694	-0.0617522\\
3.7719	-0.0618829\\
3.7745	-0.0620145\\
3.7771	-0.0621416\\
3.7796	-0.0622594\\
3.7822	-0.0623772\\
3.7848	-0.0624902\\
3.7873	-0.0625942\\
3.7899	-0.0626975\\
3.7925	-0.0627956\\
3.795	-0.062885\\
3.7976	-0.0629729\\
3.8002	-0.0630554\\
3.8027	-0.0631295\\
3.8053	-0.0632013\\
3.8078	-0.063265\\
3.8104	-0.0633258\\
3.813	-0.0633809\\
3.8155	-0.0634284\\
3.8181	-0.0634721\\
3.8207	-0.06351\\
3.8232	-0.0635408\\
3.8258	-0.0635671\\
3.8284	-0.0635873\\
3.8309	-0.0636011\\
3.8335	-0.0636094\\
3.836	-0.0636117\\
3.8386	-0.063608\\
3.8412	-0.0635982\\
3.8437	-0.0635829\\
3.8463	-0.063561\\
3.8489	-0.0635328\\
3.8514	-0.0634999\\
3.854	-0.0634595\\
3.8566	-0.0634129\\
3.8591	-0.0633622\\
3.8617	-0.0633034\\
3.8642	-0.063241\\
3.8668	-0.0631699\\
3.8694	-0.0630927\\
3.8719	-0.0630126\\
3.8745	-0.0629233\\
3.8771	-0.0628278\\
3.8796	-0.0627302\\
3.8822	-0.0626228\\
3.8848	-0.0625093\\
3.8873	-0.0623945\\
3.8899	-0.0622692\\
3.8924	-0.0621432\\
3.895	-0.0620064\\
3.8976	-0.0618637\\
3.9001	-0.0617211\\
3.9027	-0.0615671\\
3.9053	-0.0614075\\
3.9078	-0.0612488\\
3.9104	-0.0610782\\
3.913	-0.0609023\\
3.9155	-0.0607279\\
3.9181	-0.0605415\\
3.9206	-0.0603572\\
3.9232	-0.0601605\\
3.9258	-0.0599587\\
3.9283	-0.05976\\
3.9309	-0.0595485\\
3.9335	-0.0593322\\
3.936	-0.0591198\\
3.9386	-0.0588943\\
3.9412	-0.0586644\\
3.9437	-0.0584391\\
3.9463	-0.0582005\\
3.9488	-0.0579672\\
3.9514	-0.0577205\\
3.954	-0.0574698\\
3.9565	-0.0572251\\
3.9591	-0.0569669\\
3.9617	-0.0567051\\
3.9642	-0.05645\\
3.9668	-0.0561814\\
3.9694	-0.0559095\\
3.9719	-0.0556451\\
3.9745	-0.0553671\\
3.9771	-0.0550862\\
3.9796	-0.0548136\\
3.9822	-0.0545274\\
3.9847	-0.0542499\\
3.9873	-0.0539589\\
3.9899	-0.0536657\\
3.9924	-0.0533818\\
3.995	-0.0530846\\
3.9976	-0.0527856\\
4.0001	-0.0524965\\
4.0027	-0.0521943\\
4.0053	-0.0518908\\
4.0078	-0.0515977\\
4.0104	-0.0512919\\
4.0129	-0.0509968\\
4.0155	-0.0506892\\
4.0181	-0.0503809\\
4.0206	-0.050084\\
4.0232	-0.0497748\\
4.0258	-0.0494654\\
4.0283	-0.0491679\\
4.0309	-0.0488585\\
4.0335	-0.0485493\\
4.036	-0.0482524\\
4.0386	-0.047944\\
4.0411	-0.0476482\\
4.0437	-0.0473414\\
4.0463	-0.0470355\\
4.0488	-0.0467424\\
4.0514	-0.0464388\\
4.054	-0.0461367\\
4.0565	-0.0458476\\
4.0591	-0.0455487\\
4.0617	-0.0452516\\
4.0642	-0.0449678\\
4.0668	-0.0446748\\
4.0693	-0.0443953\\
4.0719	-0.0441069\\
4.0745	-0.0438211\\
4.077	-0.0435489\\
4.0796	-0.0432686\\
4.0822	-0.0429913\\
4.0847	-0.0427277\\
4.0873	-0.0424567\\
4.0899	-0.0421892\\
4.0924	-0.0419353\\
4.095	-0.0416749\\
4.0975	-0.0414281\\
4.1001	-0.0411754\\
4.1027	-0.0409267\\
4.1052	-0.0406916\\
4.1078	-0.0404514\\
4.1104	-0.0402156\\
4.1129	-0.0399933\\
4.1155	-0.0397667\\
4.1181	-0.039545\\
4.1206	-0.0393365\\
4.1232	-0.0391246\\
4.1258	-0.038918\\
4.1283	-0.0387243\\
4.1309	-0.0385282\\
4.1334	-0.0383449\\
4.136	-0.0381598\\
4.1386	-0.0379805\\
4.1411	-0.0378136\\
4.1437	-0.0376459\\
4.1463	-0.0374842\\
4.1488	-0.0373346\\
4.1514	-0.0371851\\
4.154	-0.037042\\
4.1565	-0.0369105\\
4.1591	-0.0367801\\
4.1616	-0.036661\\
4.1642	-0.0365437\\
4.1668	-0.0364332\\
4.1693	-0.0363333\\
4.1719	-0.0362363\\
4.1745	-0.0361463\\
4.177	-0.0360665\\
4.1796	-0.0359904\\
4.1822	-0.0359216\\
4.1847	-0.0358623\\
4.1873	-0.0358078\\
4.1898	-0.0357623\\
4.1924	-0.0357223\\
4.195	-0.0356899\\
4.1975	-0.0356657\\
4.2001	-0.0356481\\
4.2027	-0.0356381\\
4.2052	-0.0356357\\
4.2078	-0.0356407\\
4.2104	-0.0356535\\
4.2129	-0.0356731\\
4.2155	-0.0357011\\
4.218	-0.0357355\\
4.2206	-0.0357789\\
4.2232	-0.0358301\\
4.2257	-0.0358868\\
4.2283	-0.0359535\\
4.2309	-0.036028\\
4.2334	-0.0361072\\
4.236	-0.0361972\\
4.2386	-0.0362952\\
4.2411	-0.0363968\\
4.2437	-0.0365102\\
4.2462	-0.0366267\\
4.2488	-0.0367556\\
4.2514	-0.0368923\\
4.2539	-0.0370311\\
4.2565	-0.0371831\\
4.2591	-0.0373429\\
4.2616	-0.0375039\\
4.2642	-0.0376789\\
4.2668	-0.0378615\\
4.2693	-0.0380444\\
4.2719	-0.038242\\
4.2744	-0.0384391\\
4.277	-0.0386515\\
4.2796	-0.0388713\\
4.2821	-0.0390896\\
4.2847	-0.0393239\\
4.2873	-0.0395655\\
4.2898	-0.0398046\\
4.2924	-0.0400603\\
4.295	-0.0403231\\
4.2975	-0.0405823\\
4.3001	-0.0408588\\
4.3027	-0.0411421\\
4.3052	-0.0414209\\
4.3078	-0.0417174\\
4.3103	-0.0420088\\
4.3129	-0.0423182\\
4.3155	-0.042634\\
4.318	-0.0429437\\
4.3206	-0.0432719\\
4.3232	-0.0436061\\
4.3257	-0.0439333\\
4.3283	-0.0442793\\
4.3309	-0.0446311\\
4.3334	-0.0449747\\
4.336	-0.0453376\\
4.3385	-0.0456917\\
4.3411	-0.0460652\\
4.3437	-0.046444\\
4.3462	-0.0468129\\
4.3488	-0.0472016\\
4.3514	-0.047595\\
4.3539	-0.0479778\\
4.3565	-0.0483804\\
4.3591	-0.0487875\\
4.3616	-0.0491829\\
4.3642	-0.0495983\\
4.3667	-0.0500014\\
4.3693	-0.0504246\\
4.3719	-0.0508514\\
4.3744	-0.0512653\\
4.377	-0.051699\\
4.3796	-0.052136\\
4.3821	-0.0525592\\
4.3847	-0.0530023\\
4.3873	-0.0534481\\
4.3898	-0.0538794\\
4.3924	-0.0543303\\
4.3949	-0.0547661\\
4.3975	-0.0552215\\
4.4001	-0.0556788\\
4.4026	-0.0561204\\
4.4052	-0.0565812\\
4.4078	-0.0570436\\
4.4103	-0.0574895\\
4.4129	-0.0579543\\
4.4155	-0.0584202\\
4.418	-0.0588689\\
4.4206	-0.0593362\\
4.4231	-0.0597861\\
4.4257	-0.0602542\\
4.4283	-0.0607225\\
4.4308	-0.0611727\\
4.4334	-0.0616407\\
4.436	-0.0621084\\
4.4385	-0.0625575\\
4.4411	-0.0630239\\
4.4437	-0.0634894\\
4.4462	-0.063936\\
4.4488	-0.0643992\\
4.4514	-0.0648609\\
4.4539	-0.0653033\\
4.4565	-0.0657617\\
4.459	-0.0662007\\
4.4616	-0.066655\\
4.4642	-0.0671071\\
4.4667	-0.0675394\\
4.4693	-0.0679864\\
4.4719	-0.0684306\\
4.4744	-0.0688548\\
4.477	-0.0692929\\
4.4796	-0.0697277\\
4.4821	-0.0701424\\
4.4847	-0.0705701\\
4.4872	-0.0709777\\
4.4898	-0.0713977\\
4.4924	-0.0718135\\
4.4949	-0.0722092\\
4.4975	-0.0726163\\
4.5001	-0.0730187\\
4.5026	-0.0734011\\
4.5052	-0.073794\\
4.5078	-0.0741817\\
4.5103	-0.0745494\\
4.5129	-0.0749266\\
4.5154	-0.075284\\
4.518	-0.0756501\\
4.5206	-0.0760102\\
4.5231	-0.0763509\\
4.5257	-0.0766991\\
4.5283	-0.077041\\
4.5308	-0.0773637\\
4.5334	-0.0776929\\
4.536	-0.0780154\\
4.5385	-0.078319\\
4.5411	-0.0786279\\
4.5436	-0.0789183\\
4.5462	-0.0792132\\
4.5488	-0.0795008\\
4.5513	-0.0797704\\
4.5539	-0.0800433\\
4.5565	-0.0803086\\
4.559	-0.0805563\\
4.5616	-0.0808063\\
4.5642	-0.0810483\\
4.5667	-0.0812734\\
4.5693	-0.0814995\\
4.5718	-0.0817091\\
4.5744	-0.0819189\\
4.577	-0.0821204\\
4.5795	-0.082306\\
4.5821	-0.0824907\\
4.5847	-0.0826668\\
4.5872	-0.0828279\\
4.5898	-0.0829868\\
4.5924	-0.0831369\\
4.5949	-0.0832728\\
4.5975	-0.0834053\\
4.6001	-0.0835289\\
4.6026	-0.0836391\\
4.6052	-0.0837449\\
4.6077	-0.0838379\\
4.6103	-0.0839257\\
4.6129	-0.0840043\\
4.6154	-0.0840711\\
4.618	-0.0841316\\
4.6206	-0.0841827\\
4.6231	-0.0842231\\
4.6257	-0.0842559\\
4.6283	-0.0842795\\
4.6308	-0.0842933\\
4.6334	-0.0842985\\
4.6359	-0.0842947\\
4.6385	-0.0842815\\
4.6411	-0.0842591\\
4.6436	-0.0842287\\
4.6462	-0.0841879\\
4.6488	-0.0841379\\
4.6513	-0.0840811\\
4.6539	-0.0840129\\
4.6565	-0.0839355\\
4.659	-0.0838525\\
4.6616	-0.0837571\\
4.6641	-0.0836569\\
4.6667	-0.0835437\\
4.6693	-0.0834216\\
4.6718	-0.0832957\\
4.6744	-0.0831561\\
4.677	-0.0830076\\
4.6795	-0.0828565\\
4.6821	-0.0826908\\
4.6847	-0.0825165\\
4.6872	-0.0823409\\
4.6898	-0.0821498\\
4.6923	-0.0819582\\
4.6949	-0.0817507\\
4.6975	-0.081535\\
4.7	-0.0813199\\
4.7026	-0.0810883\\
4.7052	-0.0808486\\
4.7077	-0.0806108\\
4.7103	-0.0803558\\
4.7129	-0.0800932\\
4.7154	-0.0798335\\
4.718	-0.0795561\\
4.7205	-0.0792824\\
4.7231	-0.0789907\\
4.7257	-0.0786919\\
4.7282	-0.078398\\
4.7308	-0.0780855\\
4.7334	-0.0777663\\
4.7359	-0.0774532\\
4.7385	-0.0771211\\
4.7411	-0.0767827\\
4.7436	-0.0764514\\
4.7462	-0.0761009\\
4.7487	-0.0757582\\
4.7513	-0.0753962\\
4.7539	-0.0750284\\
4.7564	-0.0746696\\
4.759	-0.0742911\\
4.7616	-0.0739074\\
4.7641	-0.0735337\\
4.7667	-0.0731402\\
4.7693	-0.0727418\\
4.7718	-0.0723545\\
4.7744	-0.0719472\\
4.777	-0.0715356\\
4.7795	-0.0711359\\
4.7821	-0.0707163\\
4.7846	-0.0703093\\
4.7872	-0.0698823\\
4.7898	-0.0694518\\
4.7923	-0.0690347\\
4.7949	-0.0685977\\
4.7975	-0.0681578\\
4.8	-0.0677321\\
4.8026	-0.0672866\\
4.8052	-0.0668387\\
4.8077	-0.0664058\\
4.8103	-0.0659534\\
4.8128	-0.0655165\\
4.8154	-0.0650603\\
4.818	-0.0646024\\
4.8205	-0.0641607\\
4.8231	-0.0636999\\
4.8257	-0.0632381\\
4.8282	-0.062793\\
4.8308	-0.0623293\\
4.8334	-0.0618649\\
4.8359	-0.061418\\
4.8385	-0.0609528\\
4.841	-0.0605054\\
4.8436	-0.06004\\
4.8462	-0.0595749\\
4.8487	-0.0591279\\
4.8513	-0.0586636\\
4.8539	-0.0582\\
4.8564	-0.0577549\\
4.859	-0.0572931\\
4.8616	-0.0568324\\
4.8641	-0.0563907\\
4.8667	-0.0559328\\
4.8692	-0.0554941\\
4.8718	-0.0550396\\
4.8744	-0.0545871\\
4.8769	-0.054154\\
4.8795	-0.0537058\\
4.8821	-0.0532601\\
4.8846	-0.0528339\\
4.8872	-0.0523934\\
4.8898	-0.0519558\\
4.8923	-0.051538\\
4.8949	-0.0511065\\
4.8974	-0.0506948\\
4.9	-0.0502701\\
4.9026	-0.049849\\
4.9051	-0.0494476\\
4.9077	-0.0490341\\
4.9103	-0.0486246\\
4.9128	-0.0482348\\
4.9154	-0.0478337\\
4.918	-0.047437\\
4.9205	-0.04706\\
4.9231	-0.0466724\\
4.9257	-0.0462898\\
4.9282	-0.0459265\\
4.9308	-0.0455537\\
4.9333	-0.0452002\\
4.9359	-0.0448377\\
4.9385	-0.0444807\\
4.941	-0.0441427\\
4.9436	-0.0437968\\
4.9462	-0.0434566\\
4.9487	-0.043135\\
4.9513	-0.0428065\\
4.9539	-0.0424841\\
4.9564	-0.04218\\
4.959	-0.0418698\\
4.9615	-0.0415776\\
4.9641	-0.0412801\\
4.9667	-0.0409891\\
4.9692	-0.0407156\\
4.9718	-0.0404378\\
4.9744	-0.0401669\\
4.9769	-0.0399128\\
4.9795	-0.0396555\\
4.9821	-0.0394052\\
4.9846	-0.0391712\\
4.9872	-0.0389349\\
4.9897	-0.0387146\\
4.9923	-0.0384927\\
4.9949	-0.0382781\\
4.9974	-0.0380788\\
5	-0.0378789\\
};
\addplot [color=mycolor11, line width=1.0pt, forget plot]
  table[row sep=crcr]{%
-0.125	7.08572e-08\\
-0.1224	7.23374e-08\\
-0.1199	7.37619e-08\\
-0.1173	7.52436e-08\\
-0.1147	7.67241e-08\\
-0.1122	7.81487e-08\\
-0.1096	7.96314e-08\\
-0.1071	8.10558e-08\\
-0.1045	8.25369e-08\\
-0.1019	8.40193e-08\\
-0.0994	8.54439e-08\\
-0.0968	8.69252e-08\\
-0.0942	8.84081e-08\\
-0.0917	8.98338e-08\\
-0.0891	9.13148e-08\\
-0.0865	9.2797e-08\\
-0.084	9.42226e-08\\
-0.0814	9.5704e-08\\
-0.0789	9.71295e-08\\
-0.0763	9.8613e-08\\
-0.0737	1.00095e-07\\
-0.0712	1.01519e-07\\
-0.0686	1.03002e-07\\
-0.066	1.04485e-07\\
-0.0635	1.05909e-07\\
-0.0609	1.07393e-07\\
-0.0583	1.08876e-07\\
-0.0558	1.103e-07\\
-0.0532	1.11783e-07\\
-0.0507	1.13209e-07\\
-0.0481	1.14691e-07\\
-0.0455	1.16173e-07\\
-0.043	1.17599e-07\\
-0.0404	1.19081e-07\\
-0.0378	1.20563e-07\\
-0.0353	1.21989e-07\\
-0.0327	1.23471e-07\\
-0.0301	1.24953e-07\\
-0.0276	1.26379e-07\\
-0.025	1.27861e-07\\
-0.0224	1.29342e-07\\
-0.0199	1.30767e-07\\
-0.0173	1.32249e-07\\
-0.0148	1.33673e-07\\
-0.0122	1.35155e-07\\
-0.0096	1.36637e-07\\
-0.0071	1.38061e-07\\
-0.0045	1.39542e-07\\
-0.0019	1.41024e-07\\
0.0006	1.42448e-07\\
0.0032	-0.000302599\\
0.0058	-0.000683171\\
0.0083	0.00571042\\
0.0109	0.0201943\\
0.0134	0.0358946\\
0.016	0.0483205\\
0.0186	0.0554656\\
0.0211	0.0596321\\
0.0237	0.0640099\\
0.0263	0.0695166\\
0.0288	0.075205\\
0.0314	0.0804086\\
0.034	0.0843906\\
0.0365	0.087232\\
0.0391	0.0894432\\
0.0416	0.0908983\\
0.0442	0.0916542\\
0.0468	0.0917607\\
0.0493	0.0915505\\
0.0519	0.0912126\\
0.0545	0.0905864\\
0.057	0.0893015\\
0.0596	0.0869681\\
0.0622	0.0838\\
0.0647	0.0804958\\
0.0673	0.0772704\\
0.0698	0.0743785\\
0.0724	0.0711322\\
0.075	0.0671341\\
0.0775	0.062508\\
0.0801	0.0573335\\
0.0827	0.0524129\\
0.0852	0.0481855\\
0.0878	0.0440417\\
0.0904	0.0395966\\
0.0929	0.0347277\\
0.0955	0.0292496\\
0.098	0.0241262\\
0.1006	0.0194176\\
0.1032	0.0153299\\
0.1057	0.0115539\\
0.1083	0.00731149\\
0.1109	0.00267211\\
0.1134	-0.00178258\\
0.116	-0.00587386\\
0.1186	-0.00918459\\
0.1211	-0.0118782\\
0.1237	-0.0146681\\
0.1263	-0.0177483\\
0.1288	-0.0208313\\
0.1314	-0.0236875\\
0.1339	-0.0257313\\
0.1365	-0.0271524\\
0.1391	-0.0282672\\
0.1416	-0.0294598\\
0.1442	-0.0309095\\
0.1468	-0.0322458\\
0.1493	-0.0330027\\
0.1519	-0.0330654\\
0.1545	-0.0326295\\
0.157	-0.0321619\\
0.1596	-0.0318942\\
0.1621	-0.0317359\\
0.1647	-0.0312869\\
0.1673	-0.0302285\\
0.1698	-0.028668\\
0.1724	-0.0268342\\
0.175	-0.025164\\
0.1775	-0.0237924\\
0.1801	-0.022349\\
0.1827	-0.0205205\\
0.1852	-0.0182634\\
0.1878	-0.0156032\\
0.1903	-0.0131041\\
0.1929	-0.0108033\\
0.1955	-0.00871147\\
0.198	-0.00658672\\
0.2006	-0.00399718\\
0.2032	-0.00105492\\
0.2057	0.00181679\\
0.2083	0.00453137\\
0.2109	0.00689918\\
0.2134	0.00906113\\
0.216	0.0114965\\
0.2185	0.0141319\\
0.2211	0.0170062\\
0.2237	0.019692\\
0.2262	0.0219068\\
0.2288	0.0239113\\
0.2314	0.0258949\\
0.2339	0.0279984\\
0.2365	0.0303632\\
0.2391	0.0326498\\
0.2416	0.0345219\\
0.2442	0.0360785\\
0.2467	0.0373846\\
0.2493	0.0388098\\
0.2519	0.0404103\\
0.2544	0.0419777\\
0.257	0.0433723\\
0.2596	0.0443658\\
0.2621	0.0450186\\
0.2647	0.0456265\\
0.2673	0.0463589\\
0.2698	0.0471618\\
0.2724	0.0478881\\
0.2749	0.0482764\\
0.2775	0.0483205\\
0.2801	0.0481732\\
0.2826	0.048072\\
0.2852	0.0481007\\
0.2878	0.0481444\\
0.2903	0.0479893\\
0.2929	0.0474957\\
0.2955	0.0467407\\
0.298	0.045971\\
0.3006	0.0452917\\
0.3032	0.044719\\
0.3057	0.0440994\\
0.3083	0.0432053\\
0.3108	0.0420827\\
0.3134	0.0407917\\
0.316	0.0395754\\
0.3185	0.0385516\\
0.3211	0.0375381\\
0.3237	0.036389\\
0.3262	0.035056\\
0.3288	0.0334998\\
0.3314	0.0319547\\
0.3339	0.0306148\\
0.3365	0.029358\\
0.339	0.028139\\
0.3416	0.0267101\\
0.3442	0.0250974\\
0.3467	0.0234964\\
0.3493	0.0219448\\
0.3519	0.0205735\\
0.3544	0.019342\\
0.357	0.0179951\\
0.3596	0.0164919\\
0.3621	0.014954\\
0.3647	0.0134101\\
0.3672	0.0120907\\
0.3698	0.0108765\\
0.3724	0.00969312\\
0.3749	0.00846095\\
0.3775	0.00706455\\
0.3801	0.00566071\\
0.3826	0.00443669\\
0.3852	0.00334578\\
0.3878	0.00236248\\
0.3903	0.00138944\\
0.3929	0.00027365\\
0.3954	-0.000854704\\
0.398	-0.00195939\\
0.4006	-0.00289509\\
0.4031	-0.00364872\\
0.4057	-0.00439116\\
0.4083	-0.00520327\\
0.4108	-0.00606501\\
0.4134	-0.00695562\\
0.416	-0.00772338\\
0.4185	-0.00830606\\
0.4211	-0.00881716\\
0.4236	-0.00932711\\
0.4262	-0.00994262\\
0.4288	-0.010606\\
0.4313	-0.0111873\\
0.4339	-0.0116524\\
0.4365	-0.0119896\\
0.439	-0.0122848\\
0.4416	-0.0126556\\
0.4442	-0.0131051\\
0.4467	-0.0135393\\
0.4493	-0.0138962\\
0.4519	-0.0141229\\
0.4544	-0.0142734\\
0.457	-0.0144585\\
0.4595	-0.014717\\
0.4621	-0.0150373\\
0.4647	-0.0153189\\
0.4672	-0.0154877\\
0.4698	-0.0155698\\
0.4724	-0.0156398\\
0.4749	-0.0157725\\
0.4775	-0.01599\\
0.4801	-0.0162229\\
0.4826	-0.0163824\\
0.4852	-0.0164531\\
0.4877	-0.0164764\\
0.4903	-0.0165377\\
0.4929	-0.0166851\\
0.4954	-0.0168813\\
0.498	-0.0170654\\
0.5006	-0.0171703\\
0.5031	-0.0172067\\
0.5057	-0.0172465\\
0.5083	-0.0173569\\
0.5108	-0.0175384\\
0.5134	-0.01775\\
0.5159	-0.0179087\\
0.5185	-0.0180019\\
0.5211	-0.0180639\\
0.5236	-0.0181622\\
0.5262	-0.0183441\\
0.5288	-0.018582\\
0.5313	-0.0188003\\
0.5339	-0.0189667\\
0.5365	-0.0190804\\
0.539	-0.0191961\\
0.5416	-0.0193795\\
0.5441	-0.0196235\\
0.5467	-0.0199001\\
0.5493	-0.0201384\\
0.5518	-0.0203103\\
0.5544	-0.0204641\\
0.557	-0.0206534\\
0.5595	-0.0209005\\
0.5621	-0.0212036\\
0.5647	-0.021496\\
0.5672	-0.0217259\\
0.5698	-0.021919\\
0.5723	-0.0221069\\
0.5749	-0.022353\\
0.5775	-0.0226564\\
0.58	-0.0229631\\
0.5826	-0.0232459\\
0.5852	-0.0234728\\
0.5877	-0.0236663\\
0.5903	-0.0238922\\
0.5929	-0.0241716\\
0.5954	-0.0244727\\
0.598	-0.0247713\\
0.6006	-0.0250181\\
0.6031	-0.0252122\\
0.6057	-0.0254108\\
0.6082	-0.0256375\\
0.6108	-0.0259152\\
0.6134	-0.0262003\\
0.6159	-0.026439\\
0.6185	-0.0266345\\
0.6211	-0.0268019\\
0.6236	-0.0269775\\
0.6262	-0.027199\\
0.6288	-0.027444\\
0.6313	-0.0276624\\
0.6339	-0.0278411\\
0.6364	-0.027971\\
0.639	-0.028098\\
0.6416	-0.0282522\\
0.6441	-0.0284305\\
0.6467	-0.0286179\\
0.6493	-0.0287696\\
0.6518	-0.028869\\
0.6544	-0.0289454\\
0.657	-0.0290318\\
0.6595	-0.029144\\
0.6621	-0.0292779\\
0.6646	-0.0293905\\
0.6672	-0.0294648\\
0.6698	-0.0295\\
0.6723	-0.029526\\
0.6749	-0.0295749\\
0.6775	-0.0296501\\
0.68	-0.0297238\\
0.6826	-0.0297708\\
0.6852	-0.0297766\\
0.6877	-0.0297608\\
0.6903	-0.0297535\\
0.6928	-0.0297723\\
0.6954	-0.0298083\\
0.698	-0.0298317\\
0.7005	-0.0298213\\
0.7031	-0.0297802\\
0.7057	-0.0297345\\
0.7082	-0.0297109\\
0.7108	-0.0297111\\
0.7134	-0.0297159\\
0.7159	-0.0296995\\
0.7185	-0.0296516\\
0.721	-0.0295905\\
0.7236	-0.0295381\\
0.7262	-0.0295124\\
0.7287	-0.0295049\\
0.7313	-0.0294908\\
0.7339	-0.0294514\\
0.7364	-0.0293928\\
0.739	-0.0293321\\
0.7416	-0.0292936\\
0.7441	-0.0292807\\
0.7467	-0.0292753\\
0.7492	-0.0292566\\
0.7518	-0.0292151\\
0.7544	-0.0291642\\
0.7569	-0.0291278\\
0.7595	-0.0291161\\
0.7621	-0.0291235\\
0.7646	-0.029129\\
0.7672	-0.0291173\\
0.7698	-0.0290905\\
0.7723	-0.0290679\\
0.7749	-0.0290658\\
0.7775	-0.0290881\\
0.78	-0.0291188\\
0.7826	-0.0291413\\
0.7851	-0.0291473\\
0.7877	-0.0291479\\
0.7903	-0.0291616\\
0.7928	-0.0291975\\
0.7954	-0.0292523\\
0.798	-0.0293073\\
0.8005	-0.0293473\\
0.8031	-0.0293773\\
0.8057	-0.0294108\\
0.8082	-0.0294607\\
0.8108	-0.0295333\\
0.8133	-0.0296113\\
0.8159	-0.0296848\\
0.8185	-0.0297442\\
0.821	-0.0297964\\
0.8236	-0.0298607\\
0.8262	-0.0299446\\
0.8287	-0.0300385\\
0.8313	-0.0301354\\
0.8339	-0.0302197\\
0.8364	-0.0302901\\
0.839	-0.0303646\\
0.8415	-0.0304494\\
0.8441	-0.0305533\\
0.8467	-0.0306627\\
0.8492	-0.03076\\
0.8518	-0.0308477\\
0.8544	-0.0309287\\
0.8569	-0.0310126\\
0.8595	-0.0311137\\
0.8621	-0.0312244\\
0.8646	-0.031328\\
0.8672	-0.0314232\\
0.8697	-0.0315035\\
0.8723	-0.0315854\\
0.8749	-0.0316765\\
0.8774	-0.0317745\\
0.88	-0.0318787\\
0.8826	-0.0319738\\
0.8851	-0.0320522\\
0.8877	-0.0321259\\
0.8903	-0.0322026\\
0.8928	-0.0322854\\
0.8954	-0.0323774\\
0.8979	-0.0324619\\
0.9005	-0.0325375\\
0.9031	-0.0326011\\
0.9056	-0.0326594\\
0.9082	-0.0327257\\
0.9108	-0.0327997\\
0.9133	-0.0328714\\
0.9159	-0.0329371\\
0.9185	-0.03299\\
0.921	-0.0330333\\
0.9236	-0.0330798\\
0.9262	-0.0331331\\
0.9287	-0.0331884\\
0.9313	-0.0332416\\
0.9338	-0.0332823\\
0.9364	-0.0333143\\
0.939	-0.033343\\
0.9415	-0.0333748\\
0.9441	-0.0334139\\
0.9467	-0.0334533\\
0.9492	-0.0334841\\
0.9518	-0.0335056\\
0.9544	-0.0335204\\
0.9569	-0.0335358\\
0.9595	-0.0335577\\
0.962	-0.0335824\\
0.9646	-0.033605\\
0.9672	-0.0336188\\
0.9697	-0.0336242\\
0.9723	-0.0336275\\
0.9749	-0.033635\\
0.9774	-0.0336476\\
0.98	-0.0336616\\
0.9826	-0.03367\\
0.9851	-0.0336703\\
0.9877	-0.0336658\\
0.9902	-0.0336628\\
0.9928	-0.0336652\\
0.9954	-0.0336716\\
0.9979	-0.0336761\\
1.0005	-0.0336742\\
1.0031	-0.0336659\\
1.0056	-0.0336567\\
1.0082	-0.0336512\\
1.0108	-0.0336509\\
1.0133	-0.033652\\
1.0159	-0.0336492\\
1.0184	-0.0336405\\
1.021	-0.0336277\\
1.0236	-0.0336166\\
1.0261	-0.0336106\\
1.0287	-0.0336081\\
1.0313	-0.0336044\\
1.0338	-0.0335959\\
1.0364	-0.033582\\
1.039	-0.0335671\\
1.0415	-0.0335561\\
1.0441	-0.033549\\
1.0466	-0.0335434\\
1.0492	-0.0335344\\
1.0518	-0.03352\\
1.0543	-0.0335031\\
1.0569	-0.0334866\\
1.0595	-0.0334739\\
1.062	-0.0334645\\
1.0646	-0.0334534\\
1.0672	-0.0334378\\
1.0697	-0.0334183\\
1.0723	-0.0333967\\
1.0748	-0.0333781\\
1.0774	-0.033362\\
1.08	-0.0333466\\
1.0825	-0.0333286\\
1.0851	-0.033305\\
1.0877	-0.033278\\
1.0902	-0.0332523\\
1.0928	-0.0332282\\
1.0954	-0.0332058\\
1.0979	-0.0331827\\
1.1005	-0.0331543\\
1.1031	-0.0331214\\
1.1056	-0.033088\\
1.1082	-0.0330546\\
1.1107	-0.0330244\\
1.1133	-0.032993\\
1.1159	-0.0329583\\
1.1184	-0.0329205\\
1.121	-0.0328778\\
1.1236	-0.0328348\\
1.1261	-0.0327949\\
1.1287	-0.0327545\\
1.1313	-0.0327123\\
1.1338	-0.0326679\\
1.1364	-0.0326177\\
1.1389	-0.0325675\\
1.1415	-0.0325161\\
1.1441	-0.0324662\\
1.1466	-0.032418\\
1.1492	-0.0323651\\
1.1518	-0.0323082\\
1.1543	-0.0322507\\
1.1569	-0.0321905\\
1.1595	-0.0321318\\
1.162	-0.0320763\\
1.1646	-0.0320174\\
1.1671	-0.0319577\\
1.1697	-0.0318924\\
1.1723	-0.0318259\\
1.1748	-0.0317628\\
1.1774	-0.0316989\\
1.18	-0.0316352\\
1.1825	-0.0315724\\
1.1851	-0.0315043\\
1.1877	-0.0314344\\
1.1902	-0.0313675\\
1.1928	-0.0312998\\
1.1953	-0.0312363\\
1.1979	-0.03117\\
1.2005	-0.031102\\
1.203	-0.0310349\\
1.2056	-0.030965\\
1.2082	-0.0308968\\
1.2107	-0.0308336\\
1.2133	-0.0307691\\
1.2159	-0.0307043\\
1.2184	-0.0306408\\
1.221	-0.0305744\\
1.2235	-0.0305119\\
1.2261	-0.0304496\\
1.2287	-0.0303898\\
1.2312	-0.0303334\\
1.2338	-0.0302745\\
1.2364	-0.0302154\\
1.2389	-0.0301596\\
1.2415	-0.0301042\\
1.2441	-0.0300521\\
1.2466	-0.0300043\\
1.2492	-0.0299556\\
1.2518	-0.0299071\\
1.2543	-0.0298613\\
1.2569	-0.0298161\\
1.2594	-0.029776\\
1.262	-0.0297376\\
1.2646	-0.0297014\\
1.2671	-0.0296675\\
1.2697	-0.0296333\\
1.2723	-0.0296011\\
1.2748	-0.0295733\\
1.2774	-0.0295482\\
1.28	-0.0295263\\
1.2825	-0.029507\\
1.2851	-0.0294882\\
1.2876	-0.0294717\\
1.2902	-0.0294575\\
1.2928	-0.0294472\\
1.2953	-0.0294408\\
1.2979	-0.0294369\\
1.3005	-0.0294346\\
1.303	-0.029434\\
1.3056	-0.0294357\\
1.3082	-0.029441\\
1.3107	-0.0294499\\
1.3133	-0.0294622\\
1.3158	-0.0294762\\
1.3184	-0.0294924\\
1.321	-0.0295106\\
1.3235	-0.0295309\\
1.3261	-0.0295556\\
1.3287	-0.0295839\\
1.3312	-0.0296136\\
1.3338	-0.0296463\\
1.3364	-0.0296806\\
1.3389	-0.0297158\\
1.3415	-0.0297555\\
1.344	-0.029797\\
1.3466	-0.0298431\\
1.3492	-0.0298912\\
1.3517	-0.0299388\\
1.3543	-0.02999\\
1.3569	-0.0300437\\
1.3594	-0.0300983\\
1.362	-0.030158\\
1.3646	-0.0302199\\
1.3671	-0.0302808\\
1.3697	-0.0303454\\
1.3722	-0.0304092\\
1.3748	-0.030478\\
1.3774	-0.0305496\\
1.3799	-0.0306206\\
1.3825	-0.0306959\\
1.3851	-0.0307723\\
1.3876	-0.0308468\\
1.3902	-0.0309262\\
1.3928	-0.0310078\\
1.3953	-0.0310885\\
1.3979	-0.0311739\\
1.4005	-0.0312602\\
1.403	-0.0313439\\
1.4056	-0.0314322\\
1.4081	-0.0315187\\
1.4107	-0.0316106\\
1.4133	-0.0317042\\
1.4158	-0.031795\\
1.4184	-0.0318899\\
1.421	-0.0319856\\
1.4235	-0.0320786\\
1.4261	-0.032177\\
1.4287	-0.0322768\\
1.4312	-0.0323737\\
1.4338	-0.0324749\\
1.4363	-0.0325725\\
1.4389	-0.0326747\\
1.4415	-0.0327779\\
1.444	-0.0328784\\
1.4466	-0.0329838\\
1.4492	-0.0330897\\
1.4517	-0.0331916\\
1.4543	-0.0332978\\
1.4569	-0.0334045\\
1.4594	-0.0335081\\
1.462	-0.0336167\\
1.4645	-0.0337216\\
1.4671	-0.0338307\\
1.4697	-0.0339396\\
1.4722	-0.0340445\\
1.4748	-0.0341542\\
1.4774	-0.0342645\\
1.4799	-0.0343711\\
1.4825	-0.0344819\\
1.4851	-0.0345924\\
1.4876	-0.0346985\\
1.4902	-0.0348089\\
1.4927	-0.0349155\\
1.4953	-0.0350267\\
1.4979	-0.0351379\\
1.5004	-0.0352444\\
1.503	-0.0353548\\
1.5056	-0.0354648\\
1.5081	-0.0355707\\
1.5107	-0.035681\\
1.5133	-0.0357913\\
1.5158	-0.0358969\\
1.5184	-0.0360061\\
1.5209	-0.0361104\\
1.5235	-0.0362187\\
1.5261	-0.0363268\\
1.5286	-0.0364306\\
1.5312	-0.0365382\\
1.5338	-0.0366449\\
1.5363	-0.0367467\\
1.5389	-0.0368519\\
1.5415	-0.0369566\\
1.544	-0.037057\\
1.5466	-0.0371609\\
1.5491	-0.03726\\
1.5517	-0.0373621\\
1.5543	-0.0374631\\
1.5568	-0.0375595\\
1.5594	-0.0376591\\
1.562	-0.0377582\\
1.5645	-0.0378526\\
1.5671	-0.0379496\\
1.5697	-0.0380454\\
1.5722	-0.0381365\\
1.5748	-0.0382304\\
1.5774	-0.0383235\\
1.5799	-0.0384122\\
1.5825	-0.0385032\\
1.585	-0.0385895\\
1.5876	-0.0386779\\
1.5902	-0.0387652\\
1.5927	-0.0388483\\
1.5953	-0.0389338\\
1.5979	-0.0390181\\
1.6004	-0.0390978\\
1.603	-0.0391794\\
1.6056	-0.0392596\\
1.6081	-0.0393357\\
1.6107	-0.0394139\\
1.6132	-0.0394881\\
1.6158	-0.0395639\\
1.6184	-0.0396384\\
1.6209	-0.0397086\\
1.6235	-0.0397805\\
1.6261	-0.0398514\\
1.6286	-0.0399186\\
1.6312	-0.0399874\\
1.6338	-0.0400549\\
1.6363	-0.0401185\\
1.6389	-0.0401835\\
1.6414	-0.040245\\
1.644	-0.0403081\\
1.6466	-0.0403703\\
1.6491	-0.0404291\\
1.6517	-0.040489\\
1.6543	-0.0405478\\
1.6568	-0.0406036\\
1.6594	-0.0406607\\
1.662	-0.0407172\\
1.6645	-0.0407708\\
1.6671	-0.0408256\\
1.6696	-0.0408774\\
1.6722	-0.0409306\\
1.6748	-0.0409832\\
1.6773	-0.0410333\\
1.6799	-0.0410849\\
1.6825	-0.041136\\
1.685	-0.0411846\\
1.6876	-0.0412346\\
1.6902	-0.0412842\\
1.6927	-0.0413317\\
1.6953	-0.041381\\
1.6978	-0.0414283\\
1.7004	-0.0414772\\
1.703	-0.0415258\\
1.7055	-0.0415725\\
1.7081	-0.0416212\\
1.7107	-0.04167\\
1.7132	-0.0417172\\
1.7158	-0.0417664\\
1.7184	-0.0418158\\
1.7209	-0.0418634\\
1.7235	-0.0419133\\
1.7261	-0.0419637\\
1.7286	-0.0420128\\
1.7312	-0.0420643\\
1.7337	-0.0421144\\
1.7363	-0.042167\\
1.7389	-0.0422202\\
1.7414	-0.0422721\\
1.744	-0.0423271\\
1.7466	-0.0423829\\
1.7491	-0.0424374\\
1.7517	-0.042495\\
1.7543	-0.0425535\\
1.7568	-0.0426107\\
1.7594	-0.0426714\\
1.7619	-0.042731\\
1.7645	-0.0427941\\
1.7671	-0.0428584\\
1.7696	-0.0429213\\
1.7722	-0.0429881\\
1.7748	-0.0430562\\
1.7773	-0.0431231\\
1.7799	-0.0431942\\
1.7825	-0.0432667\\
1.785	-0.0433378\\
1.7876	-0.0434132\\
1.7901	-0.0434871\\
1.7927	-0.0435657\\
1.7953	-0.0436459\\
1.7978	-0.0437247\\
1.8004	-0.0438082\\
1.803	-0.0438933\\
1.8055	-0.0439767\\
1.8081	-0.0440652\\
1.8107	-0.0441556\\
1.8132	-0.0442442\\
1.8158	-0.0443381\\
1.8183	-0.04443\\
1.8209	-0.0445273\\
1.8235	-0.0446264\\
1.826	-0.0447235\\
1.8286	-0.0448264\\
1.8312	-0.0449311\\
1.8337	-0.0450335\\
1.8363	-0.0451417\\
1.8389	-0.0452518\\
1.8414	-0.0453594\\
1.844	-0.0454733\\
1.8465	-0.0455845\\
1.8491	-0.045702\\
1.8517	-0.0458213\\
1.8542	-0.0459377\\
1.8568	-0.0460605\\
1.8594	-0.0461852\\
1.8619	-0.0463069\\
1.8645	-0.0464353\\
1.8671	-0.0465654\\
1.8696	-0.0466921\\
1.8722	-0.0468257\\
1.8747	-0.0469558\\
1.8773	-0.0470928\\
1.8799	-0.0472317\\
1.8824	-0.0473668\\
1.885	-0.047509\\
1.8876	-0.0476528\\
1.8901	-0.0477926\\
1.8927	-0.0479397\\
1.8953	-0.0480885\\
1.8978	-0.0482332\\
1.9004	-0.0483851\\
1.903	-0.0485386\\
1.9055	-0.0486876\\
1.9081	-0.0488442\\
1.9106	-0.0489962\\
1.9132	-0.0491558\\
1.9158	-0.0493168\\
1.9183	-0.049473\\
1.9209	-0.0496369\\
1.9235	-0.0498021\\
1.926	-0.0499624\\
1.9286	-0.0501304\\
1.9312	-0.0502999\\
1.9337	-0.050464\\
1.9363	-0.050636\\
1.9388	-0.0508025\\
1.9414	-0.050977\\
1.944	-0.0511528\\
1.9465	-0.0513229\\
1.9491	-0.0515011\\
1.9517	-0.0516804\\
1.9542	-0.0518539\\
1.9568	-0.0520354\\
1.9594	-0.0522181\\
1.9619	-0.0523948\\
1.9645	-0.0525796\\
1.967	-0.0527583\\
1.9696	-0.0529452\\
1.9722	-0.053133\\
1.9747	-0.0533145\\
1.9773	-0.0535043\\
1.9799	-0.053695\\
1.9824	-0.0538792\\
1.985	-0.0540717\\
1.9876	-0.054265\\
1.9901	-0.0544517\\
1.9927	-0.0546467\\
1.9952	-0.054835\\
1.9978	-0.0550316\\
2.0004	-0.0552289\\
2.0029	-0.0554194\\
2.0055	-0.0556181\\
2.0081	-0.0558175\\
2.0106	-0.05601\\
2.0132	-0.0562108\\
2.0158	-0.0564122\\
2.0183	-0.0566064\\
2.0209	-0.056809\\
2.0234	-0.0570043\\
2.026	-0.0572079\\
2.0286	-0.0574121\\
2.0311	-0.0576089\\
2.0337	-0.0578141\\
2.0363	-0.0580197\\
2.0388	-0.0582177\\
2.0414	-0.0584241\\
2.044	-0.0586309\\
2.0465	-0.0588301\\
2.0491	-0.0590376\\
2.0517	-0.0592454\\
2.0542	-0.0594454\\
2.0568	-0.0596537\\
2.0593	-0.0598542\\
2.0619	-0.060063\\
2.0645	-0.0602719\\
2.067	-0.060473\\
2.0696	-0.0606822\\
2.0722	-0.0608915\\
2.0747	-0.0610928\\
2.0773	-0.0613023\\
2.0799	-0.0615118\\
2.0824	-0.0617133\\
2.085	-0.0619227\\
2.0875	-0.0621241\\
2.0901	-0.0623334\\
2.0927	-0.0625426\\
2.0952	-0.0627436\\
2.0978	-0.0629526\\
2.1004	-0.0631613\\
2.1029	-0.0633618\\
2.1055	-0.06357\\
2.1081	-0.0637779\\
2.1106	-0.0639776\\
2.1132	-0.0641849\\
2.1157	-0.0643838\\
2.1183	-0.0645903\\
2.1209	-0.0647964\\
2.1234	-0.0649941\\
2.126	-0.0651992\\
2.1286	-0.0654037\\
2.1311	-0.0655999\\
2.1337	-0.0658033\\
2.1363	-0.066006\\
2.1388	-0.0662003\\
2.1414	-0.0664017\\
2.1439	-0.0665946\\
2.1465	-0.0667945\\
2.1491	-0.0669936\\
2.1516	-0.0671842\\
2.1542	-0.0673815\\
2.1568	-0.067578\\
2.1593	-0.067766\\
2.1619	-0.0679605\\
2.1645	-0.068154\\
2.167	-0.0683391\\
2.1696	-0.0685305\\
2.1721	-0.0687135\\
2.1747	-0.0689026\\
2.1773	-0.0690905\\
2.1798	-0.0692701\\
2.1824	-0.0694555\\
2.185	-0.0696397\\
2.1875	-0.0698155\\
2.1901	-0.0699969\\
2.1927	-0.070177\\
2.1952	-0.0703487\\
2.1978	-0.0705259\\
2.2004	-0.0707015\\
2.2029	-0.0708689\\
2.2055	-0.0710414\\
2.208	-0.0712057\\
2.2106	-0.071375\\
2.2132	-0.0715426\\
2.2157	-0.0717021\\
2.2183	-0.0718663\\
2.2209	-0.0720286\\
2.2234	-0.0721829\\
2.226	-0.0723416\\
2.2286	-0.0724984\\
2.2311	-0.0726474\\
2.2337	-0.0728003\\
2.2362	-0.0729455\\
2.2388	-0.0730945\\
2.2414	-0.0732415\\
2.2439	-0.0733808\\
2.2465	-0.0735236\\
2.2491	-0.0736643\\
2.2516	-0.0737975\\
2.2542	-0.0739339\\
2.2568	-0.074068\\
2.2593	-0.0741949\\
2.2619	-0.0743246\\
2.2644	-0.0744471\\
2.267	-0.0745722\\
2.2696	-0.074695\\
2.2721	-0.0748108\\
2.2747	-0.0749288\\
2.2773	-0.0750445\\
2.2798	-0.0751534\\
2.2824	-0.0752642\\
2.285	-0.0753725\\
2.2875	-0.0754743\\
2.2901	-0.0755776\\
2.2926	-0.0756746\\
2.2952	-0.075773\\
2.2978	-0.0758687\\
2.3003	-0.0759584\\
2.3029	-0.076049\\
2.3055	-0.076137\\
2.308	-0.0762192\\
2.3106	-0.076302\\
2.3132	-0.0763822\\
2.3157	-0.0764567\\
2.3183	-0.0765316\\
2.3208	-0.0766011\\
2.3234	-0.0766706\\
2.326	-0.0767375\\
2.3285	-0.0767992\\
2.3311	-0.0768608\\
2.3337	-0.0769196\\
2.3362	-0.0769735\\
2.3388	-0.0770269\\
2.3414	-0.0770776\\
2.3439	-0.0771237\\
2.3465	-0.0771689\\
2.349	-0.0772099\\
2.3516	-0.0772497\\
2.3542	-0.0772868\\
2.3567	-0.0773199\\
2.3593	-0.0773516\\
2.3619	-0.0773806\\
2.3644	-0.0774059\\
2.367	-0.0774295\\
2.3696	-0.0774504\\
2.3721	-0.0774679\\
2.3747	-0.0774834\\
2.3773	-0.0774962\\
2.3798	-0.077506\\
2.3824	-0.0775136\\
2.3849	-0.0775183\\
2.3875	-0.0775206\\
2.3901	-0.0775202\\
2.3926	-0.0775174\\
2.3952	-0.0775119\\
2.3978	-0.0775037\\
2.4003	-0.0774934\\
2.4029	-0.0774802\\
2.4055	-0.0774644\\
2.408	-0.0774468\\
2.4106	-0.077426\\
2.4131	-0.0774036\\
2.4157	-0.0773779\\
2.4183	-0.0773498\\
2.4208	-0.0773204\\
2.4234	-0.0772874\\
2.426	-0.077252\\
2.4285	-0.0772157\\
2.4311	-0.0771757\\
2.4337	-0.0771333\\
2.4362	-0.0770903\\
2.4388	-0.0770434\\
2.4413	-0.0769962\\
2.4439	-0.0769449\\
2.4465	-0.0768913\\
2.449	-0.0768378\\
2.4516	-0.07678\\
2.4542	-0.0767201\\
2.4567	-0.0766605\\
2.4593	-0.0765965\\
2.4619	-0.0765304\\
2.4644	-0.076465\\
2.467	-0.0763951\\
2.4695	-0.076326\\
2.4721	-0.0762523\\
2.4747	-0.0761766\\
2.4772	-0.0761022\\
2.4798	-0.076023\\
2.4824	-0.075942\\
2.4849	-0.0758625\\
2.4875	-0.0757782\\
2.4901	-0.0756921\\
2.4926	-0.0756078\\
2.4952	-0.0755186\\
2.4977	-0.0754314\\
2.5003	-0.0753391\\
2.5029	-0.0752454\\
2.5054	-0.0751538\\
2.508	-0.0750572\\
2.5106	-0.0749593\\
2.5131	-0.0748638\\
2.5157	-0.0747633\\
2.5183	-0.0746614\\
2.5208	-0.0745623\\
2.5234	-0.0744581\\
2.526	-0.0743527\\
2.5285	-0.0742503\\
2.5311	-0.0741428\\
2.5336	-0.0740384\\
2.5362	-0.0739288\\
2.5388	-0.0738183\\
2.5413	-0.0737111\\
2.5439	-0.0735988\\
2.5465	-0.0734856\\
2.549	-0.0733761\\
2.5516	-0.0732613\\
2.5542	-0.0731458\\
2.5567	-0.0730341\\
2.5593	-0.0729173\\
2.5618	-0.0728044\\
2.5644	-0.0726864\\
2.567	-0.0725678\\
2.5695	-0.0724533\\
2.5721	-0.0723338\\
2.5747	-0.0722138\\
2.5772	-0.072098\\
2.5798	-0.0719773\\
2.5824	-0.0718562\\
2.5849	-0.0717395\\
2.5875	-0.0716178\\
2.59	-0.0715006\\
2.5926	-0.0713786\\
2.5952	-0.0712564\\
2.5977	-0.0711388\\
2.6003	-0.0710164\\
2.6029	-0.0708939\\
2.6054	-0.0707761\\
2.608	-0.0706536\\
2.6106	-0.0705312\\
2.6131	-0.0704136\\
2.6157	-0.0702913\\
2.6182	-0.0701739\\
2.6208	-0.070052\\
2.6234	-0.0699303\\
2.6259	-0.0698135\\
2.6285	-0.0696923\\
2.6311	-0.0695714\\
2.6336	-0.0694555\\
2.6362	-0.0693353\\
2.6388	-0.0692154\\
2.6413	-0.0691006\\
2.6439	-0.0689817\\
2.6464	-0.0688677\\
2.649	-0.0687498\\
2.6516	-0.0686323\\
2.6541	-0.06852\\
2.6567	-0.0684037\\
2.6593	-0.068288\\
2.6618	-0.0681774\\
2.6644	-0.0680631\\
2.667	-0.0679495\\
2.6695	-0.0678409\\
2.6721	-0.0677287\\
2.6746	-0.0676216\\
2.6772	-0.067511\\
2.6798	-0.0674013\\
2.6823	-0.0672966\\
2.6849	-0.0671886\\
2.6875	-0.0670815\\
2.69	-0.0669794\\
2.6926	-0.0668741\\
2.6952	-0.0667699\\
2.6977	-0.0666706\\
2.7003	-0.0665684\\
2.7029	-0.0664672\\
2.7054	-0.0663709\\
2.708	-0.0662719\\
2.7105	-0.0661778\\
2.7131	-0.066081\\
2.7157	-0.0659854\\
2.7182	-0.0658945\\
2.7208	-0.0658013\\
2.7234	-0.0657092\\
2.7259	-0.0656219\\
2.7285	-0.0655324\\
2.7311	-0.0654441\\
2.7336	-0.0653605\\
2.7362	-0.0652749\\
2.7387	-0.0651938\\
2.7413	-0.0651108\\
2.7439	-0.0650291\\
2.7464	-0.064952\\
2.749	-0.0648731\\
2.7516	-0.0647957\\
2.7541	-0.0647226\\
2.7567	-0.0646481\\
2.7593	-0.064575\\
2.7618	-0.0645061\\
2.7644	-0.064436\\
2.7669	-0.0643699\\
2.7695	-0.0643028\\
2.7721	-0.0642371\\
2.7746	-0.0641755\\
2.7772	-0.0641129\\
2.7798	-0.0640519\\
2.7823	-0.0639947\\
2.7849	-0.0639368\\
2.7875	-0.0638805\\
2.79	-0.0638279\\
2.7926	-0.0637747\\
2.7951	-0.0637251\\
2.7977	-0.063675\\
2.8003	-0.0636266\\
2.8028	-0.0635816\\
2.8054	-0.0635364\\
2.808	-0.0634928\\
2.8105	-0.0634524\\
2.8131	-0.0634119\\
2.8157	-0.0633731\\
2.8182	-0.0633373\\
2.8208	-0.0633016\\
2.8233	-0.0632688\\
2.8259	-0.0632363\\
2.8285	-0.0632054\\
2.831	-0.0631771\\
2.8336	-0.0631492\\
2.8362	-0.0631229\\
2.8387	-0.0630991\\
2.8413	-0.0630758\\
2.8439	-0.0630541\\
2.8464	-0.0630346\\
2.849	-0.0630158\\
2.8516	-0.0629984\\
2.8541	-0.0629831\\
2.8567	-0.0629686\\
2.8592	-0.0629559\\
2.8618	-0.0629442\\
2.8644	-0.0629338\\
2.8669	-0.0629251\\
2.8695	-0.0629173\\
2.8721	-0.0629108\\
2.8746	-0.0629057\\
2.8772	-0.0629017\\
2.8798	-0.0628989\\
2.8823	-0.0628973\\
2.8849	-0.0628967\\
2.8874	-0.0628972\\
2.89	-0.0628988\\
2.8926	-0.0629015\\
2.8951	-0.062905\\
2.8977	-0.0629096\\
2.9003	-0.0629151\\
2.9028	-0.0629213\\
2.9054	-0.0629285\\
2.908	-0.0629366\\
2.9105	-0.0629451\\
2.9131	-0.0629546\\
2.9156	-0.0629644\\
2.9182	-0.0629753\\
2.9208	-0.0629867\\
2.9233	-0.0629983\\
2.9259	-0.0630108\\
2.9285	-0.0630238\\
2.931	-0.0630366\\
2.9336	-0.0630504\\
2.9362	-0.0630644\\
2.9387	-0.0630782\\
2.9413	-0.0630927\\
2.9438	-0.0631069\\
2.9464	-0.0631217\\
2.949	-0.0631366\\
2.9515	-0.063151\\
2.9541	-0.0631659\\
2.9567	-0.0631807\\
2.9592	-0.0631948\\
2.9618	-0.0632093\\
2.9644	-0.0632236\\
2.9669	-0.063237\\
2.9695	-0.0632507\\
2.972	-0.0632634\\
2.9746	-0.0632763\\
2.9772	-0.0632886\\
2.9797	-0.0633\\
2.9823	-0.0633112\\
2.9849	-0.0633218\\
2.9874	-0.0633313\\
2.99	-0.0633404\\
2.9926	-0.0633488\\
2.9951	-0.063356\\
2.9977	-0.0633627\\
3.0003	-0.0633684\\
3.0028	-0.0633729\\
3.0054	-0.0633766\\
3.0079	-0.0633791\\
3.0105	-0.0633806\\
3.0131	-0.063381\\
3.0156	-0.0633802\\
3.0182	-0.0633781\\
3.0208	-0.0633746\\
3.0233	-0.0633701\\
3.0259	-0.063364\\
3.0285	-0.0633564\\
3.031	-0.0633478\\
3.0336	-0.0633373\\
3.0361	-0.0633258\\
3.0387	-0.0633122\\
3.0413	-0.063297\\
3.0438	-0.0632809\\
3.0464	-0.0632624\\
3.049	-0.0632422\\
3.0515	-0.0632212\\
3.0541	-0.0631975\\
3.0567	-0.063172\\
3.0592	-0.0631458\\
3.0618	-0.0631167\\
3.0643	-0.063087\\
3.0669	-0.0630542\\
3.0695	-0.0630194\\
3.072	-0.0629841\\
3.0746	-0.0629455\\
3.0772	-0.0629049\\
3.0797	-0.062864\\
3.0823	-0.0628194\\
3.0849	-0.0627728\\
3.0874	-0.062726\\
3.09	-0.0626754\\
3.0925	-0.0626247\\
3.0951	-0.06257\\
3.0977	-0.0625132\\
3.1002	-0.0624566\\
3.1028	-0.0623956\\
3.1054	-0.0623326\\
3.1079	-0.06227\\
3.1105	-0.0622028\\
3.1131	-0.0621335\\
3.1156	-0.0620649\\
3.1182	-0.0619915\\
3.1207	-0.061919\\
3.1233	-0.0618415\\
3.1259	-0.0617619\\
3.1284	-0.0616834\\
3.131	-0.0615998\\
3.1336	-0.0615141\\
3.1361	-0.0614298\\
3.1387	-0.0613401\\
3.1413	-0.0612484\\
3.1438	-0.0611584\\
3.1464	-0.0610628\\
3.1489	-0.0609691\\
3.1515	-0.0608697\\
3.1541	-0.0607684\\
3.1566	-0.0606692\\
3.1592	-0.0605642\\
3.1618	-0.0604574\\
3.1643	-0.0603529\\
3.1669	-0.0602426\\
3.1695	-0.0601305\\
3.172	-0.060021\\
3.1746	-0.0599055\\
3.1772	-0.0597884\\
3.1797	-0.0596742\\
3.1823	-0.0595538\\
3.1848	-0.0594366\\
3.1874	-0.0593132\\
3.19	-0.0591884\\
3.1925	-0.0590669\\
3.1951	-0.0589392\\
3.1977	-0.05881\\
3.2002	-0.0586846\\
3.2028	-0.0585529\\
3.2054	-0.0584199\\
3.2079	-0.0582909\\
3.2105	-0.0581555\\
3.213	-0.0580242\\
3.2156	-0.0578867\\
3.2182	-0.057748\\
3.2207	-0.0576138\\
3.2233	-0.0574732\\
3.2259	-0.0573317\\
3.2284	-0.0571948\\
3.231	-0.0570516\\
3.2336	-0.0569076\\
3.2361	-0.0567685\\
3.2387	-0.0566231\\
3.2412	-0.0564826\\
3.2438	-0.056336\\
3.2464	-0.0561888\\
3.2489	-0.0560468\\
3.2515	-0.0558987\\
3.2541	-0.0557502\\
3.2566	-0.055607\\
3.2592	-0.0554578\\
3.2618	-0.0553084\\
3.2643	-0.0551644\\
3.2669	-0.0550146\\
3.2694	-0.0548704\\
3.272	-0.0547204\\
3.2746	-0.0545703\\
3.2771	-0.0544261\\
3.2797	-0.0542762\\
3.2823	-0.0541263\\
3.2848	-0.0539824\\
3.2874	-0.053833\\
3.29	-0.0536838\\
3.2925	-0.0535407\\
3.2951	-0.0533922\\
3.2976	-0.0532498\\
3.3002	-0.0531021\\
3.3028	-0.052955\\
3.3053	-0.052814\\
3.3079	-0.052668\\
3.3105	-0.0525226\\
3.313	-0.0523834\\
3.3156	-0.0522394\\
3.3182	-0.0520962\\
3.3207	-0.0519593\\
3.3233	-0.0518177\\
3.3259	-0.051677\\
3.3284	-0.0515427\\
3.331	-0.0514039\\
3.3335	-0.0512715\\
3.3361	-0.0511348\\
3.3387	-0.0509992\\
3.3412	-0.0508699\\
3.3438	-0.0507367\\
3.3464	-0.0506046\\
3.3489	-0.0504789\\
3.3515	-0.0503493\\
3.3541	-0.0502212\\
3.3566	-0.0500993\\
3.3592	-0.0499739\\
3.3617	-0.0498547\\
3.3643	-0.0497321\\
3.3669	-0.0496112\\
3.3694	-0.0494963\\
3.372	-0.0493784\\
3.3746	-0.0492621\\
3.3771	-0.0491519\\
3.3797	-0.0490389\\
3.3823	-0.0489277\\
3.3848	-0.0488223\\
3.3874	-0.0487145\\
3.3899	-0.0486126\\
3.3925	-0.0485084\\
3.3951	-0.048406\\
3.3976	-0.0483094\\
3.4002	-0.0482107\\
3.4028	-0.0481141\\
3.4053	-0.048023\\
3.4079	-0.0479302\\
3.4105	-0.0478394\\
3.413	-0.047754\\
3.4156	-0.0476673\\
3.4181	-0.0475858\\
3.4207	-0.0475032\\
3.4233	-0.0474227\\
3.4258	-0.0473473\\
3.4284	-0.047271\\
3.431	-0.0471969\\
3.4335	-0.0471277\\
3.4361	-0.047058\\
3.4387	-0.0469904\\
3.4412	-0.0469276\\
3.4438	-0.0468646\\
3.4463	-0.046806\\
3.4489	-0.0467474\\
3.4515	-0.0466912\\
3.454	-0.0466392\\
3.4566	-0.0465875\\
3.4592	-0.0465382\\
3.4617	-0.0464929\\
3.4643	-0.0464482\\
3.4669	-0.0464059\\
3.4694	-0.0463674\\
3.472	-0.0463298\\
3.4745	-0.0462959\\
3.4771	-0.046263\\
3.4797	-0.0462325\\
3.4822	-0.0462055\\
3.4848	-0.0461797\\
3.4874	-0.0461565\\
3.4899	-0.0461364\\
3.4925	-0.046118\\
3.4951	-0.0461019\\
3.4976	-0.0460889\\
3.5002	-0.0460777\\
3.5028	-0.046069\\
3.5053	-0.0460629\\
3.5079	-0.046059\\
3.5104	-0.0460576\\
3.513	-0.0460585\\
3.5156	-0.0460619\\
3.5181	-0.0460675\\
3.5207	-0.0460757\\
3.5233	-0.0460864\\
3.5258	-0.046099\\
3.5284	-0.0461144\\
3.531	-0.0461323\\
3.5335	-0.0461519\\
3.5361	-0.0461745\\
3.5386	-0.0461986\\
3.5412	-0.0462261\\
3.5438	-0.0462559\\
3.5463	-0.0462868\\
3.5489	-0.0463214\\
3.5515	-0.0463583\\
3.554	-0.046396\\
3.5566	-0.0464375\\
3.5592	-0.0464814\\
3.5617	-0.0465258\\
3.5643	-0.0465742\\
3.5668	-0.0466229\\
3.5694	-0.0466759\\
3.572	-0.0467311\\
3.5745	-0.0467863\\
3.5771	-0.0468459\\
3.5797	-0.0469078\\
3.5822	-0.0469693\\
3.5848	-0.0470355\\
3.5874	-0.0471038\\
3.5899	-0.0471715\\
3.5925	-0.047244\\
3.595	-0.0473157\\
3.5976	-0.0473924\\
3.6002	-0.047471\\
3.6027	-0.0475486\\
3.6053	-0.0476313\\
3.6079	-0.0477159\\
3.6104	-0.0477992\\
3.613	-0.0478876\\
3.6156	-0.047978\\
3.6181	-0.0480667\\
3.6207	-0.0481608\\
3.6232	-0.0482529\\
3.6258	-0.0483505\\
3.6284	-0.0484499\\
3.6309	-0.0485471\\
3.6335	-0.0486499\\
3.6361	-0.0487544\\
3.6386	-0.0488563\\
3.6412	-0.048964\\
3.6438	-0.0490732\\
3.6463	-0.0491797\\
3.6489	-0.0492919\\
3.6515	-0.0494057\\
3.654	-0.0495164\\
3.6566	-0.0496329\\
3.6591	-0.0497462\\
3.6617	-0.0498654\\
3.6643	-0.0499859\\
3.6668	-0.0501029\\
3.6694	-0.0502258\\
3.672	-0.0503499\\
3.6745	-0.0504703\\
3.6771	-0.0505967\\
3.6797	-0.0507241\\
3.6822	-0.0508475\\
3.6848	-0.0509769\\
3.6873	-0.0511022\\
3.6899	-0.0512335\\
3.6925	-0.0513656\\
3.695	-0.0514934\\
3.6976	-0.0516271\\
3.7002	-0.0517615\\
3.7027	-0.0518915\\
3.7053	-0.0520273\\
3.7079	-0.0521637\\
3.7104	-0.0522954\\
3.713	-0.0524329\\
3.7155	-0.0525657\\
3.7181	-0.0527041\\
3.7207	-0.052843\\
3.7232	-0.052977\\
3.7258	-0.0531166\\
3.7284	-0.0532565\\
3.7309	-0.0533912\\
3.7335	-0.0535316\\
3.7361	-0.0536721\\
3.7386	-0.0538073\\
3.7412	-0.0539481\\
3.7437	-0.0540834\\
3.7463	-0.0542242\\
3.7489	-0.054365\\
3.7514	-0.0545002\\
3.754	-0.0546408\\
3.7566	-0.0547812\\
3.7591	-0.054916\\
3.7617	-0.055056\\
3.7643	-0.0551957\\
3.7668	-0.0553297\\
3.7694	-0.0554688\\
3.7719	-0.0556021\\
3.7745	-0.0557404\\
3.7771	-0.0558782\\
3.7796	-0.0560102\\
3.7822	-0.056147\\
3.7848	-0.0562832\\
3.7873	-0.0564136\\
3.7899	-0.0565485\\
3.7925	-0.0566828\\
3.795	-0.0568112\\
3.7976	-0.0569441\\
3.8002	-0.0570762\\
3.8027	-0.0572024\\
3.8053	-0.0573328\\
3.8078	-0.0574574\\
3.8104	-0.0575861\\
3.813	-0.0577138\\
3.8155	-0.0578357\\
3.8181	-0.0579615\\
3.8207	-0.0580863\\
3.8232	-0.0582053\\
3.8258	-0.0583279\\
3.8284	-0.0584495\\
3.8309	-0.0585653\\
3.8335	-0.0586846\\
3.836	-0.0587982\\
3.8386	-0.0589151\\
3.8412	-0.0590308\\
3.8437	-0.0591409\\
3.8463	-0.0592541\\
3.8489	-0.059366\\
3.8514	-0.0594723\\
3.854	-0.0595816\\
3.8566	-0.0596895\\
3.8591	-0.0597919\\
3.8617	-0.0598971\\
3.8642	-0.0599968\\
3.8668	-0.0600992\\
3.8694	-0.0602\\
3.8719	-0.0602956\\
3.8745	-0.0603936\\
3.8771	-0.0604901\\
3.8796	-0.0605814\\
3.8822	-0.0606748\\
3.8848	-0.0607668\\
3.8873	-0.0608537\\
3.8899	-0.0609425\\
3.8924	-0.0610264\\
3.895	-0.0611122\\
3.8976	-0.0611963\\
3.9001	-0.0612757\\
3.9027	-0.0613567\\
3.9053	-0.061436\\
3.9078	-0.0615108\\
3.9104	-0.0615869\\
3.913	-0.0616614\\
3.9155	-0.0617314\\
3.9181	-0.0618027\\
3.9206	-0.0618696\\
3.9232	-0.0619376\\
3.9258	-0.062004\\
3.9283	-0.0620662\\
3.9309	-0.0621293\\
3.9335	-0.0621907\\
3.936	-0.0622481\\
3.9386	-0.0623063\\
3.9412	-0.0623628\\
3.9437	-0.0624155\\
3.9463	-0.0624687\\
3.9488	-0.0625183\\
3.9514	-0.0625683\\
3.954	-0.0626166\\
3.9565	-0.0626615\\
3.9591	-0.0627066\\
3.9617	-0.0627501\\
3.9642	-0.0627904\\
3.9668	-0.0628307\\
3.9694	-0.0628693\\
3.9719	-0.062905\\
3.9745	-0.0629405\\
3.9771	-0.0629745\\
3.9796	-0.0630056\\
3.9822	-0.0630365\\
3.9847	-0.0630647\\
3.9873	-0.0630925\\
3.9899	-0.0631188\\
3.9924	-0.0631426\\
3.995	-0.0631659\\
3.9976	-0.0631877\\
4.0001	-0.0632073\\
4.0027	-0.0632262\\
4.0053	-0.0632436\\
4.0078	-0.0632591\\
4.0104	-0.0632737\\
4.0129	-0.0632864\\
4.0155	-0.0632983\\
4.0181	-0.0633087\\
4.0206	-0.0633175\\
4.0232	-0.0633253\\
4.0258	-0.0633318\\
4.0283	-0.0633367\\
4.0309	-0.0633406\\
4.0335	-0.0633432\\
4.036	-0.0633445\\
4.0386	-0.0633447\\
4.0411	-0.0633437\\
4.0437	-0.0633414\\
4.0463	-0.0633379\\
4.0488	-0.0633335\\
4.0514	-0.0633278\\
4.054	-0.0633209\\
4.0565	-0.0633132\\
4.0591	-0.0633042\\
4.0617	-0.063294\\
4.0642	-0.0632833\\
4.0668	-0.0632711\\
4.0693	-0.0632584\\
4.0719	-0.0632442\\
4.0745	-0.0632291\\
4.077	-0.0632136\\
4.0796	-0.0631966\\
4.0822	-0.0631787\\
4.0847	-0.0631606\\
4.0873	-0.063141\\
4.0899	-0.0631205\\
4.0924	-0.0631\\
4.095	-0.0630779\\
4.0975	-0.0630559\\
4.1001	-0.0630322\\
4.1027	-0.0630078\\
4.1052	-0.0629837\\
4.1078	-0.0629579\\
4.1104	-0.0629314\\
4.1129	-0.0629053\\
4.1155	-0.0628776\\
4.1181	-0.0628492\\
4.1206	-0.0628214\\
4.1232	-0.0627918\\
4.1258	-0.0627617\\
4.1283	-0.0627323\\
4.1309	-0.0627012\\
4.1334	-0.0626708\\
4.136	-0.0626387\\
4.1386	-0.0626062\\
4.1411	-0.0625745\\
4.1437	-0.0625411\\
4.1463	-0.0625073\\
4.1488	-0.0624745\\
4.1514	-0.06244\\
4.154	-0.0624052\\
4.1565	-0.0623715\\
4.1591	-0.0623361\\
4.1616	-0.0623018\\
4.1642	-0.0622658\\
4.1668	-0.0622297\\
4.1693	-0.0621947\\
4.1719	-0.0621581\\
4.1745	-0.0621214\\
4.177	-0.0620859\\
4.1796	-0.0620489\\
4.1822	-0.0620117\\
4.1847	-0.0619758\\
4.1873	-0.0619385\\
4.1898	-0.0619025\\
4.1924	-0.061865\\
4.195	-0.0618275\\
4.1975	-0.0617914\\
4.2001	-0.0617539\\
4.2027	-0.0617164\\
4.2052	-0.0616803\\
4.2078	-0.0616429\\
4.2104	-0.0616055\\
4.2129	-0.0615697\\
4.2155	-0.0615325\\
4.218	-0.0614968\\
4.2206	-0.0614599\\
4.2232	-0.0614231\\
4.2257	-0.0613879\\
4.2283	-0.0613514\\
4.2309	-0.0613151\\
4.2334	-0.0612804\\
4.236	-0.0612446\\
4.2386	-0.061209\\
4.2411	-0.0611749\\
4.2437	-0.0611398\\
4.2462	-0.0611063\\
4.2488	-0.0610718\\
4.2514	-0.0610375\\
4.2539	-0.0610049\\
4.2565	-0.0609713\\
4.2591	-0.060938\\
4.2616	-0.0609064\\
4.2642	-0.0608738\\
4.2668	-0.0608417\\
4.2693	-0.0608111\\
4.2719	-0.0607797\\
4.2744	-0.0607499\\
4.277	-0.0607194\\
4.2796	-0.0606893\\
4.2821	-0.0606607\\
4.2847	-0.0606315\\
4.2873	-0.0606027\\
4.2898	-0.0605756\\
4.2924	-0.0605478\\
4.295	-0.0605204\\
4.2975	-0.0604947\\
4.3001	-0.0604684\\
4.3027	-0.0604426\\
4.3052	-0.0604183\\
4.3078	-0.0603936\\
4.3103	-0.0603704\\
4.3129	-0.0603468\\
4.3155	-0.0603237\\
4.318	-0.0603021\\
4.3206	-0.0602802\\
4.3232	-0.0602588\\
4.3257	-0.0602389\\
4.3283	-0.0602187\\
4.3309	-0.0601991\\
4.3334	-0.0601809\\
4.336	-0.0601625\\
4.3385	-0.0601454\\
4.3411	-0.0601283\\
4.3437	-0.0601117\\
4.3462	-0.0600964\\
4.3488	-0.0600811\\
4.3514	-0.0600665\\
4.3539	-0.060053\\
4.3565	-0.0600396\\
4.3591	-0.0600268\\
4.3616	-0.0600151\\
4.3642	-0.0600036\\
4.3667	-0.0599931\\
4.3693	-0.0599829\\
4.3719	-0.0599733\\
4.3744	-0.0599646\\
4.377	-0.0599563\\
4.3796	-0.0599486\\
4.3821	-0.0599417\\
4.3847	-0.0599353\\
4.3873	-0.0599294\\
4.3898	-0.0599244\\
4.3924	-0.0599198\\
4.3949	-0.0599159\\
4.3975	-0.0599125\\
4.4001	-0.0599097\\
4.4026	-0.0599076\\
4.4052	-0.059906\\
4.4078	-0.0599049\\
4.4103	-0.0599045\\
4.4129	-0.0599046\\
4.4155	-0.0599053\\
4.418	-0.0599065\\
4.4206	-0.0599082\\
4.4231	-0.0599105\\
4.4257	-0.0599133\\
4.4283	-0.0599166\\
4.4308	-0.0599204\\
4.4334	-0.0599247\\
4.436	-0.0599296\\
4.4385	-0.0599347\\
4.4411	-0.0599405\\
4.4437	-0.0599467\\
4.4462	-0.0599532\\
4.4488	-0.0599603\\
4.4514	-0.0599678\\
4.4539	-0.0599754\\
4.4565	-0.0599837\\
4.459	-0.0599921\\
4.4616	-0.0600011\\
4.4642	-0.0600105\\
4.4667	-0.0600198\\
4.4693	-0.0600298\\
4.4719	-0.0600402\\
4.4744	-0.0600503\\
4.477	-0.0600612\\
4.4796	-0.0600723\\
4.4821	-0.0600832\\
4.4847	-0.0600947\\
4.4872	-0.060106\\
4.4898	-0.0601179\\
4.4924	-0.06013\\
4.4949	-0.0601417\\
4.4975	-0.0601541\\
4.5001	-0.0601665\\
4.5026	-0.0601785\\
4.5052	-0.0601911\\
4.5078	-0.0602037\\
4.5103	-0.0602159\\
4.5129	-0.0602286\\
4.5154	-0.0602407\\
4.518	-0.0602533\\
4.5206	-0.0602659\\
4.5231	-0.0602779\\
4.5257	-0.0602903\\
4.5283	-0.0603025\\
4.5308	-0.0603142\\
4.5334	-0.0603262\\
4.536	-0.060338\\
4.5385	-0.0603491\\
4.5411	-0.0603605\\
4.5436	-0.0603712\\
4.5462	-0.0603822\\
4.5488	-0.0603928\\
4.5513	-0.0604027\\
4.5539	-0.0604127\\
4.5565	-0.0604224\\
4.559	-0.0604313\\
4.5616	-0.0604403\\
4.5642	-0.0604488\\
4.5667	-0.0604566\\
4.5693	-0.0604643\\
4.5718	-0.0604713\\
4.5744	-0.0604781\\
4.577	-0.0604844\\
4.5795	-0.0604899\\
4.5821	-0.0604952\\
4.5847	-0.0604999\\
4.5872	-0.0605038\\
4.5898	-0.0605074\\
4.5924	-0.0605103\\
4.5949	-0.0605126\\
4.5975	-0.0605143\\
4.6001	-0.0605153\\
4.6026	-0.0605157\\
4.6052	-0.0605154\\
4.6077	-0.0605145\\
4.6103	-0.0605128\\
4.6129	-0.0605103\\
4.6154	-0.0605073\\
4.618	-0.0605034\\
4.6206	-0.0604987\\
4.6231	-0.0604935\\
4.6257	-0.0604872\\
4.6283	-0.0604801\\
4.6308	-0.0604725\\
4.6334	-0.0604638\\
4.6359	-0.0604546\\
4.6385	-0.0604442\\
4.6411	-0.060433\\
4.6436	-0.0604213\\
4.6462	-0.0604083\\
4.6488	-0.0603944\\
4.6513	-0.0603801\\
4.6539	-0.0603644\\
4.6565	-0.0603478\\
4.659	-0.0603309\\
4.6616	-0.0603124\\
4.6641	-0.0602938\\
4.6667	-0.0602734\\
4.6693	-0.0602521\\
4.6718	-0.0602307\\
4.6744	-0.0602075\\
4.677	-0.0601834\\
4.6795	-0.0601592\\
4.6821	-0.0601332\\
4.6847	-0.0601061\\
4.6872	-0.0600792\\
4.6898	-0.0600502\\
4.6923	-0.0600214\\
4.6949	-0.0599905\\
4.6975	-0.0599586\\
4.7	-0.059927\\
4.7026	-0.0598932\\
4.7052	-0.0598584\\
4.7077	-0.059824\\
4.7103	-0.0597873\\
4.7129	-0.0597496\\
4.7154	-0.0597125\\
4.718	-0.0596729\\
4.7205	-0.0596339\\
4.7231	-0.0595925\\
4.7257	-0.0595501\\
4.7282	-0.0595084\\
4.7308	-0.0594642\\
4.7334	-0.059419\\
4.7359	-0.0593747\\
4.7385	-0.0593278\\
4.7411	-0.0592799\\
4.7436	-0.059233\\
4.7462	-0.0591834\\
4.7487	-0.0591349\\
4.7513	-0.0590835\\
4.7539	-0.0590313\\
4.7564	-0.0589803\\
4.759	-0.0589265\\
4.7616	-0.0588718\\
4.7641	-0.0588185\\
4.7667	-0.0587623\\
4.7693	-0.0587053\\
4.7718	-0.0586497\\
4.7744	-0.0585912\\
4.777	-0.0585319\\
4.7795	-0.0584742\\
4.7821	-0.0584135\\
4.7846	-0.0583545\\
4.7872	-0.0582924\\
4.7898	-0.0582297\\
4.7923	-0.0581687\\
4.7949	-0.0581047\\
4.7975	-0.0580401\\
4.8	-0.0579774\\
4.8026	-0.0579116\\
4.8052	-0.0578452\\
4.8077	-0.0577808\\
4.8103	-0.0577134\\
4.8128	-0.057648\\
4.8154	-0.0575795\\
4.818	-0.0575106\\
4.8205	-0.0574439\\
4.8231	-0.057374\\
4.8257	-0.0573037\\
4.8282	-0.0572358\\
4.8308	-0.0571647\\
4.8334	-0.0570933\\
4.8359	-0.0570243\\
4.8385	-0.0569521\\
4.841	-0.0568825\\
4.8436	-0.0568098\\
4.8462	-0.0567368\\
4.8487	-0.0566664\\
4.8513	-0.0565929\\
4.8539	-0.0565193\\
4.8564	-0.0564482\\
4.859	-0.0563742\\
4.8616	-0.0563\\
4.8641	-0.0562285\\
4.8667	-0.0561541\\
4.8692	-0.0560824\\
4.8718	-0.0560078\\
4.8744	-0.0559331\\
4.8769	-0.0558613\\
4.8795	-0.0557866\\
4.8821	-0.0557119\\
4.8846	-0.0556401\\
4.8872	-0.0555655\\
4.8898	-0.0554909\\
4.8923	-0.0554193\\
4.8949	-0.055345\\
4.8974	-0.0552736\\
4.9	-0.0551995\\
4.9026	-0.0551256\\
4.9051	-0.0550547\\
4.9077	-0.0549811\\
4.9103	-0.0549078\\
4.9128	-0.0548376\\
4.9154	-0.0547647\\
4.918	-0.0546922\\
4.9205	-0.0546228\\
4.9231	-0.0545509\\
4.9257	-0.0544793\\
4.9282	-0.0544108\\
4.9308	-0.05434\\
4.9333	-0.0542722\\
4.9359	-0.0542022\\
4.9385	-0.0541325\\
4.941	-0.054066\\
4.9436	-0.0539973\\
4.9462	-0.053929\\
4.9487	-0.0538639\\
4.9513	-0.0537967\\
4.9539	-0.05373\\
4.9564	-0.0536663\\
4.959	-0.0536007\\
4.9615	-0.0535382\\
4.9641	-0.0534737\\
4.9667	-0.0534099\\
4.9692	-0.0533491\\
4.9718	-0.0532865\\
4.9744	-0.0532246\\
4.9769	-0.0531657\\
4.9795	-0.0531051\\
4.9821	-0.0530453\\
4.9846	-0.0529884\\
4.9872	-0.0529299\\
4.9897	-0.0528744\\
4.9923	-0.0528174\\
4.9949	-0.0527612\\
4.9974	-0.0527079\\
5	-0.0526532\\
};
\addplot [color=mycolor12, line width=1.0pt, forget plot]
  table[row sep=crcr]{%
-0.125	7.253e-08\\
-0.1224	7.40343e-08\\
-0.1199	7.548e-08\\
-0.1173	7.69832e-08\\
-0.1147	7.84869e-08\\
-0.1122	7.99325e-08\\
-0.1096	8.14353e-08\\
-0.1071	8.28804e-08\\
-0.1045	8.43829e-08\\
-0.1019	8.58844e-08\\
-0.0994	8.73286e-08\\
-0.0968	8.88308e-08\\
-0.0942	9.03324e-08\\
-0.0917	9.17759e-08\\
-0.0891	9.3277e-08\\
-0.0865	9.47772e-08\\
-0.084	9.62195e-08\\
-0.0814	9.77202e-08\\
-0.0789	9.91629e-08\\
-0.0763	1.00663e-07\\
-0.0737	1.02162e-07\\
-0.0712	1.03604e-07\\
-0.0686	1.05102e-07\\
-0.066	1.066e-07\\
-0.0635	1.08041e-07\\
-0.0609	1.09539e-07\\
-0.0583	1.11037e-07\\
-0.0558	1.12476e-07\\
-0.0532	1.13973e-07\\
-0.0507	1.15411e-07\\
-0.0481	1.16908e-07\\
-0.0455	1.18404e-07\\
-0.043	1.19841e-07\\
-0.0404	1.21337e-07\\
-0.0378	1.22831e-07\\
-0.0353	1.24268e-07\\
-0.0327	1.25762e-07\\
-0.0301	1.27256e-07\\
-0.0276	1.28692e-07\\
-0.025	1.30184e-07\\
-0.0224	1.31677e-07\\
-0.0199	1.33111e-07\\
-0.0173	1.34603e-07\\
-0.0148	1.36037e-07\\
-0.0122	1.37528e-07\\
-0.0096	1.39018e-07\\
-0.0071	1.40451e-07\\
-0.0045	1.41941e-07\\
-0.0019	1.4343e-07\\
0.0006	1.44861e-07\\
0.0032	-0.000379021\\
0.0058	-0.00440256\\
0.0083	-0.00684582\\
0.0109	-0.00199828\\
0.0134	0.00771015\\
0.016	0.0174284\\
0.0186	0.0231068\\
0.0211	0.0249699\\
0.0237	0.0254775\\
0.0263	0.0263705\\
0.0288	0.0278165\\
0.0314	0.0291989\\
0.034	0.0298337\\
0.0365	0.0296665\\
0.0391	0.0288459\\
0.0416	0.0274427\\
0.0442	0.0252082\\
0.0468	0.0222067\\
0.0493	0.0189166\\
0.0519	0.0154873\\
0.0545	0.0121527\\
0.057	0.00871267\\
0.0596	0.00441639\\
0.0622	-0.00076263\\
0.0647	-0.0062207\\
0.0673	-0.0117792\\
0.0698	-0.0166677\\
0.0724	-0.0214525\\
0.075	-0.0264703\\
0.0775	-0.0318882\\
0.0801	-0.0380205\\
0.0827	-0.0441002\\
0.0852	-0.0493865\\
0.0878	-0.0542167\\
0.0904	-0.0587697\\
0.0929	-0.0633674\\
0.0955	-0.0685298\\
0.098	-0.0735403\\
0.1006	-0.0782514\\
0.1032	-0.0821412\\
0.1057	-0.0852717\\
0.1083	-0.0883716\\
0.1109	-0.0916736\\
0.1134	-0.0949676\\
0.116	-0.0980786\\
0.1186	-0.100439\\
0.1211	-0.101949\\
0.1237	-0.103082\\
0.1263	-0.104216\\
0.1288	-0.10546\\
0.1314	-0.10666\\
0.1339	-0.107317\\
0.1365	-0.107248\\
0.1391	-0.106578\\
0.1416	-0.105765\\
0.1442	-0.105065\\
0.1468	-0.104472\\
0.1493	-0.103676\\
0.1519	-0.102281\\
0.1545	-0.100277\\
0.157	-0.0980679\\
0.1596	-0.0958586\\
0.1621	-0.0939601\\
0.1647	-0.0920299\\
0.1673	-0.0897928\\
0.1698	-0.0871712\\
0.1724	-0.0840958\\
0.175	-0.0810249\\
0.1775	-0.0783392\\
0.1801	-0.0757991\\
0.1827	-0.0732311\\
0.1852	-0.0704691\\
0.1878	-0.067268\\
0.1903	-0.0641145\\
0.1929	-0.0610771\\
0.1955	-0.0584131\\
0.198	-0.0560491\\
0.2006	-0.0535027\\
0.2032	-0.0507032\\
0.2057	-0.0478813\\
0.2083	-0.0451037\\
0.2109	-0.0427133\\
0.2134	-0.040754\\
0.216	-0.0388259\\
0.2185	-0.0368488\\
0.2211	-0.0346345\\
0.2237	-0.0324686\\
0.2262	-0.0306837\\
0.2288	-0.0292331\\
0.2314	-0.0280469\\
0.2339	-0.02691\\
0.2365	-0.0255886\\
0.2391	-0.0242172\\
0.2416	-0.0230805\\
0.2442	-0.0222672\\
0.2467	-0.0218095\\
0.2493	-0.0214586\\
0.2519	-0.021021\\
0.2544	-0.0204882\\
0.257	-0.0199848\\
0.2596	-0.019751\\
0.2621	-0.0198514\\
0.2647	-0.0201634\\
0.2673	-0.0204614\\
0.2698	-0.0206146\\
0.2724	-0.0207066\\
0.2749	-0.0209178\\
0.2775	-0.0214129\\
0.2801	-0.0221572\\
0.2826	-0.0229323\\
0.2852	-0.0236203\\
0.2878	-0.0241602\\
0.2903	-0.0246733\\
0.2929	-0.0253789\\
0.2955	-0.0263199\\
0.298	-0.0273397\\
0.3006	-0.0283327\\
0.3032	-0.0291504\\
0.3057	-0.0298294\\
0.3083	-0.0305904\\
0.3108	-0.0314926\\
0.3134	-0.0325829\\
0.316	-0.0336719\\
0.3185	-0.0345738\\
0.3211	-0.0353392\\
0.3237	-0.0360505\\
0.3262	-0.0368294\\
0.3288	-0.0377898\\
0.3314	-0.0388095\\
0.3339	-0.039697\\
0.3365	-0.0404376\\
0.339	-0.0410249\\
0.3416	-0.0416477\\
0.3442	-0.0423919\\
0.3467	-0.0432042\\
0.3493	-0.0440234\\
0.3519	-0.0446912\\
0.3544	-0.0451753\\
0.357	-0.0456175\\
0.3596	-0.0461293\\
0.3621	-0.0467317\\
0.3647	-0.0473976\\
0.3672	-0.0479541\\
0.3698	-0.0483746\\
0.3724	-0.0486809\\
0.3749	-0.0489841\\
0.3775	-0.0494005\\
0.3801	-0.0499051\\
0.3826	-0.0503734\\
0.3852	-0.0507405\\
0.3878	-0.0509738\\
0.3903	-0.0511561\\
0.3929	-0.0514123\\
0.3954	-0.0517609\\
0.398	-0.0521685\\
0.4006	-0.052513\\
0.4031	-0.0527275\\
0.4057	-0.0528681\\
0.4083	-0.053029\\
0.4108	-0.0532778\\
0.4134	-0.0536224\\
0.416	-0.0539648\\
0.4185	-0.0542088\\
0.4211	-0.0543641\\
0.4236	-0.0544882\\
0.4262	-0.0546841\\
0.4288	-0.0549824\\
0.4313	-0.0553158\\
0.4339	-0.0556201\\
0.4365	-0.0558297\\
0.439	-0.0559728\\
0.4416	-0.0561463\\
0.4442	-0.0564107\\
0.4467	-0.0567405\\
0.4493	-0.0570868\\
0.4519	-0.0573618\\
0.4544	-0.0575497\\
0.457	-0.0577269\\
0.4595	-0.0579532\\
0.4621	-0.0582745\\
0.4647	-0.0586374\\
0.4672	-0.0589511\\
0.4698	-0.0591969\\
0.4724	-0.0593881\\
0.4749	-0.059589\\
0.4775	-0.0598695\\
0.4801	-0.0602124\\
0.4826	-0.0605396\\
0.4852	-0.0608143\\
0.4877	-0.0610072\\
0.4903	-0.0611842\\
0.4929	-0.0614033\\
0.4954	-0.0616757\\
0.498	-0.0619847\\
0.5006	-0.0622528\\
0.5031	-0.0624362\\
0.5057	-0.0625718\\
0.5083	-0.0627123\\
0.5108	-0.0628956\\
0.5134	-0.0631279\\
0.5159	-0.0633395\\
0.5185	-0.0634941\\
0.5211	-0.0635742\\
0.5236	-0.0636219\\
0.5262	-0.063696\\
0.5288	-0.0638145\\
0.5313	-0.0639394\\
0.5339	-0.0640251\\
0.5365	-0.0640342\\
0.539	-0.0639888\\
0.5416	-0.0639364\\
0.5441	-0.0639177\\
0.5467	-0.063925\\
0.5493	-0.0639132\\
0.5518	-0.0638408\\
0.5544	-0.0636959\\
0.557	-0.0635173\\
0.5595	-0.0633589\\
0.5621	-0.0632247\\
0.5647	-0.0630941\\
0.5672	-0.0629272\\
0.5698	-0.0626849\\
0.5723	-0.0624013\\
0.5749	-0.0620971\\
0.5775	-0.0618176\\
0.58	-0.0615673\\
0.5826	-0.0612884\\
0.5852	-0.0609524\\
0.5877	-0.0605709\\
0.5903	-0.0601448\\
0.5929	-0.0597295\\
0.5954	-0.0593545\\
0.598	-0.0589673\\
0.6006	-0.0585433\\
0.6031	-0.0580802\\
0.6057	-0.0575554\\
0.6082	-0.0570446\\
0.6108	-0.0565356\\
0.6134	-0.0560451\\
0.6159	-0.0555604\\
0.6185	-0.0550114\\
0.6211	-0.054413\\
0.6236	-0.0538168\\
0.6262	-0.0532091\\
0.6288	-0.0526267\\
0.6313	-0.0520725\\
0.6339	-0.0514693\\
0.6364	-0.0508465\\
0.639	-0.0501662\\
0.6416	-0.0494849\\
0.6441	-0.0488521\\
0.6467	-0.048214\\
0.6493	-0.0475689\\
0.6518	-0.0469166\\
0.6544	-0.0462019\\
0.657	-0.0454736\\
0.6595	-0.0447882\\
0.6621	-0.0441013\\
0.6646	-0.0434506\\
0.6672	-0.0427565\\
0.6698	-0.0420291\\
0.6723	-0.0413079\\
0.6749	-0.040562\\
0.6775	-0.039841\\
0.68	-0.0391685\\
0.6826	-0.038468\\
0.6852	-0.0377436\\
0.6877	-0.0370221\\
0.6903	-0.0362651\\
0.6928	-0.035554\\
0.6954	-0.0348407\\
0.698	-0.0341408\\
0.7005	-0.0334577\\
0.7031	-0.0327238\\
0.7057	-0.0319749\\
0.7082	-0.0312621\\
0.7108	-0.0305457\\
0.7134	-0.0298519\\
0.7159	-0.0291875\\
0.7185	-0.0284807\\
0.721	-0.0277836\\
0.7236	-0.0270565\\
0.7262	-0.0263482\\
0.7287	-0.0256921\\
0.7313	-0.0250243\\
0.7339	-0.0243513\\
0.7364	-0.0236885\\
0.739	-0.0229902\\
0.7416	-0.0223022\\
0.7441	-0.0216642\\
0.7467	-0.0210228\\
0.7492	-0.0204119\\
0.7518	-0.0197667\\
0.7544	-0.0191095\\
0.7569	-0.0184792\\
0.7595	-0.0178427\\
0.7621	-0.0172319\\
0.7646	-0.0166596\\
0.7672	-0.0160635\\
0.7698	-0.0154569\\
0.7723	-0.0148698\\
0.7749	-0.0142712\\
0.7775	-0.0136971\\
0.78	-0.0131663\\
0.7826	-0.0126227\\
0.7851	-0.0120954\\
0.7877	-0.0115405\\
0.7903	-0.0109909\\
0.7928	-0.0104812\\
0.7954	-0.00997602\\
0.798	-0.00948793\\
0.8005	-0.00902123\\
0.8031	-0.00853092\\
0.8057	-0.00804098\\
0.8082	-0.00758318\\
0.8108	-0.00713101\\
0.8133	-0.0067178\\
0.8159	-0.006299\\
0.8185	-0.00587989\\
0.821	-0.00547536\\
0.8236	-0.00506269\\
0.8262	-0.00467066\\
0.8287	-0.00431741\\
0.8313	-0.00396786\\
0.8339	-0.00362445\\
0.8364	-0.00329386\\
0.839	-0.00295406\\
0.8415	-0.00264193\\
0.8441	-0.00234105\\
0.8467	-0.00206305\\
0.8492	-0.00180804\\
0.8518	-0.00154613\\
0.8544	-0.00128651\\
0.8569	-0.00104679\\
0.8595	-0.000817969\\
0.8621	-0.000613709\\
0.8646	-0.000434915\\
0.8672	-0.000257338\\
0.8697	-8.94294e-05\\
0.8723	7.8694e-05\\
0.8749	0.000230606\\
0.8774	0.000354539\\
0.88	0.000461563\\
0.8826	0.000554636\\
0.8851	0.000638676\\
0.8877	0.000721386\\
0.8903	0.000792465\\
0.8928	0.000841703\\
0.8954	0.000870091\\
0.8979	0.000880102\\
0.9005	0.00088095\\
0.9031	0.000876823\\
0.9056	0.000865402\\
0.9082	0.000837982\\
0.9108	0.000788751\\
0.9133	0.000721832\\
0.9159	0.000638803\\
0.9185	0.000548976\\
0.921	0.000457084\\
0.9236	0.000350245\\
0.9262	0.000224784\\
0.9287	8.42725e-05\\
0.9313	-7.83186e-05\\
0.9338	-0.000243714\\
0.9364	-0.000421123\\
0.939	-0.000606436\\
0.9415	-0.000798257\\
0.9441	-0.00101677\\
0.9467	-0.0012537\\
0.9492	-0.00149326\\
0.9518	-0.00174885\\
0.9544	-0.00201019\\
0.9569	-0.00227125\\
0.9595	-0.00255872\\
0.962	-0.00285235\\
0.9646	-0.00317214\\
0.9672	-0.00350047\\
0.9697	-0.00382093\\
0.9723	-0.00416102\\
0.9749	-0.00451349\\
0.9774	-0.00486792\\
0.98	-0.00525185\\
0.9826	-0.00564633\\
0.9851	-0.0060309\\
0.9877	-0.00643556\\
0.9902	-0.00683272\\
0.9928	-0.00725895\\
0.9954	-0.0077003\\
0.9979	-0.00813608\\
1.0005	-0.00859607\\
1.0031	-0.00905999\\
1.0056	-0.0095112\\
1.0082	-0.00999021\\
1.0108	-0.0104825\\
1.0133	-0.0109677\\
1.0159	-0.0114805\\
1.0184	-0.0119774\\
1.021	-0.0124975\\
1.0236	-0.013024\\
1.0261	-0.0135401\\
1.0287	-0.0140886\\
1.0313	-0.0146463\\
1.0338	-0.0151872\\
1.0364	-0.0157521\\
1.039	-0.0163205\\
1.0415	-0.0168739\\
1.0441	-0.0174592\\
1.0466	-0.018031\\
1.0492	-0.0186313\\
1.0518	-0.0192339\\
1.0543	-0.0198148\\
1.0569	-0.0204228\\
1.0595	-0.0210382\\
1.062	-0.021638\\
1.0646	-0.0222679\\
1.0672	-0.0229005\\
1.0697	-0.0235091\\
1.0723	-0.0241433\\
1.0748	-0.0247573\\
1.0774	-0.0254023\\
1.08	-0.0260534\\
1.0825	-0.0266824\\
1.0851	-0.0273365\\
1.0877	-0.02799\\
1.0902	-0.0286196\\
1.0928	-0.0292785\\
1.0954	-0.0299423\\
1.0979	-0.0305837\\
1.1005	-0.0312507\\
1.1031	-0.0319158\\
1.1056	-0.0325542\\
1.1082	-0.0332195\\
1.1107	-0.0338621\\
1.1133	-0.0345331\\
1.1159	-0.0352043\\
1.1184	-0.0358473\\
1.121	-0.0365129\\
1.1236	-0.0371771\\
1.1261	-0.0378165\\
1.1287	-0.0384829\\
1.1313	-0.0391494\\
1.1338	-0.0397876\\
1.1364	-0.0404471\\
1.1389	-0.0410776\\
1.1415	-0.0417316\\
1.1441	-0.0423853\\
1.1466	-0.0430133\\
1.1492	-0.0436637\\
1.1518	-0.0443091\\
1.1543	-0.0449244\\
1.1569	-0.0455602\\
1.1595	-0.0461935\\
1.162	-0.0468009\\
1.1646	-0.0474296\\
1.1671	-0.0480291\\
1.1697	-0.0486461\\
1.1723	-0.0492567\\
1.1748	-0.0498395\\
1.1774	-0.0504425\\
1.18	-0.0510418\\
1.1825	-0.051613\\
1.1851	-0.0521999\\
1.1877	-0.0527791\\
1.1902	-0.0533298\\
1.1928	-0.0538975\\
1.1953	-0.0544391\\
1.1979	-0.054997\\
1.2005	-0.0555476\\
1.203	-0.0560688\\
1.2056	-0.0566027\\
1.2082	-0.0571299\\
1.2107	-0.0576314\\
1.2133	-0.0581472\\
1.2159	-0.0586556\\
1.2184	-0.0591361\\
1.221	-0.0596267\\
1.2235	-0.0600904\\
1.2261	-0.0605659\\
1.2287	-0.0610349\\
1.2312	-0.061479\\
1.2338	-0.0619323\\
1.2364	-0.0623757\\
1.2389	-0.0627931\\
1.2415	-0.0632191\\
1.2441	-0.0636379\\
1.2466	-0.0640337\\
1.2492	-0.0644369\\
1.2518	-0.0648303\\
1.2543	-0.0651991\\
1.2569	-0.0655735\\
1.2594	-0.065926\\
1.262	-0.0662851\\
1.2646	-0.066636\\
1.2671	-0.0669645\\
1.2697	-0.0672962\\
1.2723	-0.067618\\
1.2748	-0.067919\\
1.2774	-0.0682244\\
1.28	-0.0685218\\
1.2825	-0.0687994\\
1.2851	-0.0690783\\
1.2876	-0.069337\\
1.2902	-0.0695969\\
1.2928	-0.0698484\\
1.2953	-0.0700828\\
1.2979	-0.0703185\\
1.3005	-0.0705449\\
1.303	-0.0707533\\
1.3056	-0.0709605\\
1.3082	-0.0711591\\
1.3107	-0.0713426\\
1.3133	-0.0715258\\
1.3158	-0.071694\\
1.3184	-0.0718596\\
1.321	-0.0720158\\
1.3235	-0.0721576\\
1.3261	-0.0722972\\
1.3287	-0.0724294\\
1.3312	-0.0725492\\
1.3338	-0.0726653\\
1.3364	-0.0727724\\
1.3389	-0.072867\\
1.3415	-0.0729576\\
1.344	-0.0730379\\
1.3466	-0.0731144\\
1.3492	-0.0731831\\
1.3517	-0.0732414\\
1.3543	-0.0732936\\
1.3569	-0.0733379\\
1.3594	-0.0733738\\
1.362	-0.0734046\\
1.3646	-0.0734286\\
1.3671	-0.0734445\\
1.3697	-0.0734533\\
1.3722	-0.0734545\\
1.3748	-0.073449\\
1.3774	-0.0734373\\
1.3799	-0.0734201\\
1.3825	-0.0733959\\
1.3851	-0.0733645\\
1.3876	-0.0733275\\
1.3902	-0.0732825\\
1.3928	-0.0732315\\
1.3953	-0.0731773\\
1.3979	-0.0731153\\
1.4005	-0.0730471\\
1.403	-0.0729753\\
1.4056	-0.0728945\\
1.4081	-0.0728115\\
1.4107	-0.0727202\\
1.4133	-0.072624\\
1.4158	-0.0725265\\
1.4184	-0.0724195\\
1.421	-0.0723068\\
1.4235	-0.0721935\\
1.4261	-0.072071\\
1.4287	-0.0719444\\
1.4312	-0.0718184\\
1.4338	-0.0716828\\
1.4363	-0.0715477\\
1.4389	-0.0714024\\
1.4415	-0.071253\\
1.444	-0.0711058\\
1.4466	-0.0709493\\
1.4492	-0.0707889\\
1.4517	-0.0706309\\
1.4543	-0.0704624\\
1.4569	-0.0702902\\
1.4594	-0.0701216\\
1.462	-0.0699435\\
1.4645	-0.0697695\\
1.4671	-0.0695854\\
1.4697	-0.069398\\
1.4722	-0.0692147\\
1.4748	-0.0690215\\
1.4774	-0.0688261\\
1.4799	-0.0686362\\
1.4825	-0.0684366\\
1.4851	-0.0682344\\
1.4876	-0.0680376\\
1.4902	-0.0678309\\
1.4927	-0.0676305\\
1.4953	-0.0674208\\
1.4979	-0.0672098\\
1.5004	-0.0670053\\
1.503	-0.066791\\
1.5056	-0.0665751\\
1.5081	-0.0663664\\
1.5107	-0.0661487\\
1.5133	-0.0659304\\
1.5158	-0.0657198\\
1.5184	-0.0654999\\
1.5209	-0.0652876\\
1.5235	-0.0650662\\
1.5261	-0.0648446\\
1.5286	-0.0646317\\
1.5312	-0.0644103\\
1.5338	-0.0641889\\
1.5363	-0.0639757\\
1.5389	-0.063754\\
1.5415	-0.0635327\\
1.544	-0.0633204\\
1.5466	-0.0631005\\
1.5491	-0.0628897\\
1.5517	-0.062671\\
1.5543	-0.0624528\\
1.5568	-0.0622437\\
1.5594	-0.0620274\\
1.562	-0.0618125\\
1.5645	-0.0616072\\
1.5671	-0.0613949\\
1.5697	-0.0611837\\
1.5722	-0.0609818\\
1.5748	-0.0607733\\
1.5774	-0.0605667\\
1.5799	-0.0603699\\
1.5825	-0.060167\\
1.585	-0.0599736\\
1.5876	-0.0597741\\
1.5902	-0.0595765\\
1.5927	-0.0593885\\
1.5953	-0.0591955\\
1.5979	-0.0590048\\
1.6004	-0.0588236\\
1.603	-0.0586372\\
1.6056	-0.0584531\\
1.6081	-0.0582784\\
1.6107	-0.0580994\\
1.6132	-0.0579299\\
1.6158	-0.0577563\\
1.6184	-0.0575853\\
1.6209	-0.0574233\\
1.6235	-0.0572575\\
1.6261	-0.0570946\\
1.6286	-0.0569409\\
1.6312	-0.0567841\\
1.6338	-0.0566302\\
1.6363	-0.0564849\\
1.6389	-0.0563367\\
1.6414	-0.0561971\\
1.644	-0.0560551\\
1.6466	-0.0559164\\
1.6491	-0.0557862\\
1.6517	-0.0556537\\
1.6543	-0.0555244\\
1.6568	-0.055403\\
1.6594	-0.0552801\\
1.662	-0.0551607\\
1.6645	-0.0550491\\
1.6671	-0.0549363\\
1.6696	-0.0548309\\
1.6722	-0.0547245\\
1.6748	-0.0546215\\
1.6773	-0.0545257\\
1.6799	-0.0544296\\
1.6825	-0.054337\\
1.685	-0.0542511\\
1.6876	-0.054165\\
1.6902	-0.0540823\\
1.6927	-0.054006\\
1.6953	-0.0539302\\
1.6978	-0.0538607\\
1.7004	-0.0537917\\
1.703	-0.053726\\
1.7055	-0.053666\\
1.7081	-0.0536069\\
1.7107	-0.0535514\\
1.7132	-0.0535012\\
1.7158	-0.0534525\\
1.7184	-0.053407\\
1.7209	-0.0533663\\
1.7235	-0.0533272\\
1.7261	-0.0532915\\
1.7286	-0.0532603\\
1.7312	-0.0532312\\
1.7337	-0.0532063\\
1.7363	-0.0531835\\
1.7389	-0.0531637\\
1.7414	-0.0531477\\
1.744	-0.0531342\\
1.7466	-0.0531239\\
1.7491	-0.053117\\
1.7517	-0.0531127\\
1.7543	-0.0531114\\
1.7568	-0.0531128\\
1.7594	-0.0531172\\
1.7619	-0.0531243\\
1.7645	-0.0531345\\
1.7671	-0.0531475\\
1.7696	-0.0531626\\
1.7722	-0.0531809\\
1.7748	-0.0532018\\
1.7773	-0.0532245\\
1.7799	-0.0532507\\
1.7825	-0.0532794\\
1.785	-0.0533094\\
1.7876	-0.0533428\\
1.7901	-0.0533772\\
1.7927	-0.0534152\\
1.7953	-0.0534556\\
1.7978	-0.0534965\\
1.8004	-0.0535412\\
1.803	-0.0535879\\
1.8055	-0.0536346\\
1.8081	-0.0536852\\
1.8107	-0.0537376\\
1.8132	-0.0537899\\
1.8158	-0.053846\\
1.8183	-0.0539014\\
1.8209	-0.0539607\\
1.8235	-0.0540214\\
1.826	-0.0540813\\
1.8286	-0.0541449\\
1.8312	-0.05421\\
1.8337	-0.0542738\\
1.8363	-0.0543412\\
1.8389	-0.0544097\\
1.8414	-0.0544765\\
1.844	-0.0545471\\
1.8465	-0.0546157\\
1.8491	-0.054688\\
1.8517	-0.054761\\
1.8542	-0.0548317\\
1.8568	-0.0549058\\
1.8594	-0.0549804\\
1.8619	-0.0550526\\
1.8645	-0.0551281\\
1.8671	-0.0552038\\
1.8696	-0.0552767\\
1.8722	-0.0553526\\
1.8747	-0.0554256\\
1.8773	-0.0555015\\
1.8799	-0.0555773\\
1.8824	-0.0556499\\
1.885	-0.0557252\\
1.8876	-0.0558\\
1.8901	-0.0558715\\
1.8927	-0.0559453\\
1.8953	-0.0560185\\
1.8978	-0.0560883\\
1.9004	-0.05616\\
1.903	-0.0562309\\
1.9055	-0.0562981\\
1.9081	-0.0563669\\
1.9106	-0.0564321\\
1.9132	-0.0564988\\
1.9158	-0.0565643\\
1.9183	-0.0566259\\
1.9209	-0.0566886\\
1.9235	-0.0567497\\
1.926	-0.056807\\
1.9286	-0.0568651\\
1.9312	-0.0569214\\
1.9337	-0.0569739\\
1.9363	-0.0570266\\
1.9388	-0.0570755\\
1.9414	-0.0571243\\
1.944	-0.0571711\\
1.9465	-0.0572141\\
1.9491	-0.0572567\\
1.9517	-0.057297\\
1.9542	-0.0573335\\
1.9568	-0.0573692\\
1.9594	-0.0574023\\
1.9619	-0.0574319\\
1.9645	-0.0574601\\
1.967	-0.0574847\\
1.9696	-0.0575077\\
1.9722	-0.0575279\\
1.9747	-0.0575447\\
1.9773	-0.0575594\\
1.9799	-0.0575712\\
1.9824	-0.0575798\\
1.985	-0.0575858\\
1.9876	-0.0575887\\
1.9901	-0.0575885\\
1.9927	-0.0575854\\
1.9952	-0.0575793\\
1.9978	-0.05757\\
2.0004	-0.0575573\\
2.0029	-0.0575421\\
2.0055	-0.057523\\
2.0081	-0.0575005\\
2.0106	-0.0574757\\
2.0132	-0.0574467\\
2.0158	-0.0574142\\
2.0183	-0.0573797\\
2.0209	-0.0573404\\
2.0234	-0.0572994\\
2.026	-0.0572532\\
2.0286	-0.0572035\\
2.0311	-0.0571525\\
2.0337	-0.0570959\\
2.0363	-0.0570357\\
2.0388	-0.0569745\\
2.0414	-0.0569073\\
2.044	-0.0568366\\
2.0465	-0.0567652\\
2.0491	-0.0566874\\
2.0517	-0.0566061\\
2.0542	-0.0565245\\
2.0568	-0.0564362\\
2.0593	-0.0563479\\
2.0619	-0.0562525\\
2.0645	-0.0561537\\
2.067	-0.0560554\\
2.0696	-0.0559496\\
2.0722	-0.0558404\\
2.0747	-0.0557321\\
2.0773	-0.0556161\\
2.0799	-0.0554966\\
2.0824	-0.0553786\\
2.085	-0.0552526\\
2.0875	-0.0551282\\
2.0901	-0.0549956\\
2.0927	-0.0548597\\
2.0952	-0.0547261\\
2.0978	-0.0545839\\
2.1004	-0.0544386\\
2.1029	-0.0542959\\
2.1055	-0.0541445\\
2.1081	-0.0539901\\
2.1106	-0.0538387\\
2.1132	-0.0536785\\
2.1157	-0.0535217\\
2.1183	-0.0533558\\
2.1209	-0.0531872\\
2.1234	-0.0530224\\
2.126	-0.0528485\\
2.1286	-0.0526719\\
2.1311	-0.0524997\\
2.1337	-0.0523182\\
2.1363	-0.0521342\\
2.1388	-0.0519551\\
2.1414	-0.0517665\\
2.1439	-0.051583\\
2.1465	-0.0513901\\
2.1491	-0.0511951\\
2.1516	-0.0510057\\
2.1542	-0.0508068\\
2.1568	-0.0506059\\
2.1593	-0.050411\\
2.1619	-0.0502067\\
2.1645	-0.0500007\\
2.167	-0.049801\\
2.1696	-0.0495919\\
2.1721	-0.0493895\\
2.1747	-0.0491776\\
2.1773	-0.0489645\\
2.1798	-0.0487584\\
2.1824	-0.0485429\\
2.185	-0.0483264\\
2.1875	-0.0481173\\
2.1901	-0.0478989\\
2.1927	-0.0476798\\
2.1952	-0.0474684\\
2.1978	-0.0472479\\
2.2004	-0.0470269\\
2.2029	-0.0468139\\
2.2055	-0.046592\\
2.208	-0.0463784\\
2.2106	-0.046156\\
2.2132	-0.0459335\\
2.2157	-0.0457196\\
2.2183	-0.0454971\\
2.2209	-0.0452747\\
2.2234	-0.045061\\
2.226	-0.0448392\\
2.2286	-0.0446177\\
2.2311	-0.0444052\\
2.2337	-0.0441847\\
2.2362	-0.0439733\\
2.2388	-0.0437542\\
2.2414	-0.0435359\\
2.2439	-0.0433269\\
2.2465	-0.0431104\\
2.2491	-0.0428951\\
2.2516	-0.0426891\\
2.2542	-0.0424761\\
2.2568	-0.0422645\\
2.2593	-0.0420623\\
2.2619	-0.0418535\\
2.2644	-0.0416543\\
2.267	-0.0414487\\
2.2696	-0.0412449\\
2.2721	-0.0410506\\
2.2747	-0.0408505\\
2.2773	-0.0406523\\
2.2798	-0.0404638\\
2.2824	-0.0402698\\
2.285	-0.0400781\\
2.2875	-0.0398959\\
2.2901	-0.0397088\\
2.2926	-0.0395313\\
2.2952	-0.0393492\\
2.2978	-0.0391697\\
2.3003	-0.0389996\\
2.3029	-0.0388255\\
2.3055	-0.0386543\\
2.308	-0.0384925\\
2.3106	-0.0383271\\
2.3132	-0.0381648\\
2.3157	-0.0380117\\
2.3183	-0.0378557\\
2.3208	-0.0377088\\
2.3234	-0.0375593\\
2.326	-0.0374132\\
2.3285	-0.0372761\\
2.3311	-0.037137\\
2.3337	-0.0370016\\
2.3362	-0.0368748\\
2.3388	-0.0367466\\
2.3414	-0.0366223\\
2.3439	-0.0365064\\
2.3465	-0.0363898\\
2.349	-0.0362814\\
2.3516	-0.0361726\\
2.3542	-0.036068\\
2.3567	-0.0359713\\
2.3593	-0.0358749\\
2.3619	-0.0357828\\
2.3644	-0.0356982\\
2.367	-0.0356146\\
2.3696	-0.0355354\\
2.3721	-0.0354635\\
2.3747	-0.035393\\
2.3773	-0.0353272\\
2.3798	-0.0352682\\
2.3824	-0.0352113\\
2.3849	-0.0351611\\
2.3875	-0.0351135\\
2.3901	-0.0350705\\
2.3926	-0.0350338\\
2.3952	-0.0350002\\
2.3978	-0.0349715\\
2.4003	-0.0349485\\
2.4029	-0.0349293\\
2.4055	-0.0349151\\
2.408	-0.034906\\
2.4106	-0.0349014\\
2.4131	-0.0349017\\
2.4157	-0.0349069\\
2.4183	-0.034917\\
2.4208	-0.0349315\\
2.4234	-0.0349515\\
2.426	-0.0349765\\
2.4285	-0.0350053\\
2.4311	-0.0350402\\
2.4337	-0.0350801\\
2.4362	-0.0351233\\
2.4388	-0.0351731\\
2.4413	-0.0352258\\
2.4439	-0.0352855\\
2.4465	-0.0353502\\
2.449	-0.0354172\\
2.4516	-0.0354918\\
2.4542	-0.0355714\\
2.4567	-0.0356526\\
2.4593	-0.035742\\
2.4619	-0.0358363\\
2.4644	-0.0359316\\
2.467	-0.0360356\\
2.4695	-0.0361401\\
2.4721	-0.0362537\\
2.4747	-0.036372\\
2.4772	-0.0364904\\
2.4798	-0.0366181\\
2.4824	-0.0367507\\
2.4849	-0.0368825\\
2.4875	-0.0370243\\
2.4901	-0.0371706\\
2.4926	-0.0373157\\
2.4952	-0.0374711\\
2.4977	-0.0376248\\
2.5003	-0.037789\\
2.5029	-0.0379576\\
2.5054	-0.0381239\\
2.508	-0.038301\\
2.5106	-0.0384825\\
2.5131	-0.038661\\
2.5157	-0.0388507\\
2.5183	-0.0390446\\
2.5208	-0.0392348\\
2.5234	-0.0394366\\
2.526	-0.0396424\\
2.5285	-0.0398439\\
2.5311	-0.0400572\\
2.5336	-0.0402659\\
2.5362	-0.0404866\\
2.5388	-0.0407109\\
2.5413	-0.0409299\\
2.5439	-0.0411612\\
2.5465	-0.0413959\\
2.549	-0.0416248\\
2.5516	-0.0418661\\
2.5542	-0.0421106\\
2.5567	-0.0423486\\
2.5593	-0.0425992\\
2.5618	-0.042843\\
2.5644	-0.0430995\\
2.567	-0.0433588\\
2.5695	-0.0436107\\
2.5721	-0.0438754\\
2.5747	-0.0441427\\
2.5772	-0.0444021\\
2.5798	-0.0446743\\
2.5824	-0.0449489\\
2.5849	-0.0452151\\
2.5875	-0.045494\\
2.59	-0.0457643\\
2.5926	-0.0460474\\
2.5952	-0.0463324\\
2.5977	-0.0466083\\
2.6003	-0.0468969\\
2.6029	-0.0471873\\
2.6054	-0.0474679\\
2.608	-0.0477614\\
2.6106	-0.0480562\\
2.6131	-0.048341\\
2.6157	-0.0486384\\
2.6182	-0.0489254\\
2.6208	-0.049225\\
2.6234	-0.0495256\\
2.6259	-0.0498155\\
2.6285	-0.0501177\\
2.6311	-0.0504207\\
2.6336	-0.0507126\\
2.6362	-0.0510167\\
2.6388	-0.0513213\\
2.6413	-0.0516144\\
2.6439	-0.0519195\\
2.6464	-0.0522131\\
2.649	-0.0525185\\
2.6516	-0.0528238\\
2.6541	-0.0531173\\
2.6567	-0.0534224\\
2.6593	-0.0537271\\
2.6618	-0.0540198\\
2.6644	-0.0543237\\
2.667	-0.054627\\
2.6695	-0.054918\\
2.6721	-0.05522\\
2.6746	-0.0555095\\
2.6772	-0.0558097\\
2.6798	-0.0561089\\
2.6823	-0.0563955\\
2.6849	-0.0566924\\
2.6875	-0.056988\\
2.69	-0.0572709\\
2.6926	-0.0575637\\
2.6952	-0.057855\\
2.6977	-0.0581334\\
2.7003	-0.0584214\\
2.7029	-0.0587075\\
2.7054	-0.0589808\\
2.708	-0.0592631\\
2.7105	-0.0595326\\
2.7131	-0.0598108\\
2.7157	-0.0600867\\
2.7182	-0.0603498\\
2.7208	-0.0606211\\
2.7234	-0.0608899\\
2.7259	-0.061146\\
2.7285	-0.0614097\\
2.7311	-0.0616707\\
2.7336	-0.0619191\\
2.7362	-0.0621745\\
2.7387	-0.0624174\\
2.7413	-0.062667\\
2.7439	-0.0629136\\
2.7464	-0.0631477\\
2.749	-0.063388\\
2.7516	-0.063625\\
2.7541	-0.0638497\\
2.7567	-0.06408\\
2.7593	-0.0643069\\
2.7618	-0.0645217\\
2.7644	-0.0647415\\
2.7669	-0.0649495\\
2.7695	-0.065162\\
2.7721	-0.0653708\\
2.7746	-0.065568\\
2.7772	-0.0657692\\
2.7798	-0.0659664\\
2.7823	-0.0661523\\
2.7849	-0.0663417\\
2.7875	-0.0665269\\
2.79	-0.0667012\\
2.7926	-0.0668782\\
2.7951	-0.0670445\\
2.7977	-0.0672133\\
2.8003	-0.0673777\\
2.8028	-0.0675318\\
2.8054	-0.0676876\\
2.808	-0.0678391\\
2.8105	-0.0679806\\
2.8131	-0.0681233\\
2.8157	-0.0682615\\
2.8182	-0.0683901\\
2.8208	-0.0685194\\
2.8233	-0.0686394\\
2.8259	-0.0687597\\
2.8285	-0.0688754\\
2.831	-0.0689822\\
2.8336	-0.0690888\\
2.8362	-0.0691908\\
2.8387	-0.0692844\\
2.8413	-0.0693771\\
2.8439	-0.0694652\\
2.8464	-0.0695454\\
2.849	-0.0696243\\
2.8516	-0.0696985\\
2.8541	-0.0697654\\
2.8567	-0.0698304\\
2.8592	-0.0698885\\
2.8618	-0.0699443\\
2.8644	-0.0699954\\
2.8669	-0.0700402\\
2.8695	-0.0700823\\
2.8721	-0.0701197\\
2.8746	-0.0701513\\
2.8772	-0.0701797\\
2.8798	-0.0702036\\
2.8823	-0.0702222\\
2.8849	-0.0702372\\
2.8874	-0.0702473\\
2.89	-0.0702535\\
2.8926	-0.0702553\\
2.8951	-0.0702529\\
2.8977	-0.070246\\
2.9003	-0.0702349\\
2.9028	-0.0702202\\
2.9054	-0.0702007\\
2.908	-0.070177\\
2.9105	-0.0701503\\
2.9131	-0.0701185\\
2.9156	-0.0700841\\
2.9182	-0.0700444\\
2.9208	-0.0700007\\
2.9233	-0.069955\\
2.9259	-0.0699037\\
2.9285	-0.0698486\\
2.931	-0.0697921\\
2.9336	-0.0697298\\
2.9362	-0.0696638\\
2.9387	-0.0695969\\
2.9413	-0.069524\\
2.9438	-0.0694506\\
2.9464	-0.069371\\
2.949	-0.0692881\\
2.9515	-0.0692053\\
2.9541	-0.0691161\\
2.9567	-0.0690239\\
2.9592	-0.0689323\\
2.9618	-0.0688341\\
2.9644	-0.0687331\\
2.9669	-0.0686333\\
2.9695	-0.0685268\\
2.972	-0.068422\\
2.9746	-0.0683104\\
2.9772	-0.0681963\\
2.9797	-0.0680844\\
2.9823	-0.0679656\\
2.9849	-0.0678447\\
2.9874	-0.0677263\\
2.99	-0.0676011\\
2.9926	-0.0674739\\
2.9951	-0.0673498\\
2.9977	-0.067219\\
3.0003	-0.0670863\\
3.0028	-0.0669572\\
3.0054	-0.0668214\\
3.0079	-0.0666894\\
3.0105	-0.0665507\\
3.0131	-0.0664107\\
3.0156	-0.066275\\
3.0182	-0.0661326\\
3.0208	-0.0659892\\
3.0233	-0.0658503\\
3.0259	-0.0657051\\
3.0285	-0.065559\\
3.031	-0.0654179\\
3.0336	-0.0652705\\
3.0361	-0.0651283\\
3.0387	-0.0649799\\
3.0413	-0.0648311\\
3.0438	-0.0646878\\
3.0464	-0.0645386\\
3.049	-0.0643893\\
3.0515	-0.0642456\\
3.0541	-0.0640963\\
3.0567	-0.0639472\\
3.0592	-0.0638039\\
3.0618	-0.0636553\\
3.0643	-0.0635127\\
3.0669	-0.063365\\
3.0695	-0.0632177\\
3.072	-0.0630768\\
3.0746	-0.0629309\\
3.0772	-0.0627858\\
3.0797	-0.0626471\\
3.0823	-0.0625038\\
3.0849	-0.0623616\\
3.0874	-0.0622259\\
3.09	-0.0620859\\
3.0925	-0.0619525\\
3.0951	-0.0618151\\
3.0977	-0.0616792\\
3.1002	-0.0615498\\
3.1028	-0.0614168\\
3.1054	-0.0612855\\
3.1079	-0.0611608\\
3.1105	-0.0610329\\
3.1131	-0.0609068\\
3.1156	-0.0607873\\
3.1182	-0.0606651\\
3.1207	-0.0605494\\
3.1233	-0.0604312\\
3.1259	-0.0603152\\
3.1284	-0.0602057\\
3.131	-0.0600941\\
3.1336	-0.0599848\\
3.1361	-0.059882\\
3.1387	-0.0597775\\
3.1413	-0.0596755\\
3.1438	-0.0595799\\
3.1464	-0.059483\\
3.1489	-0.0593923\\
3.1515	-0.0593007\\
3.1541	-0.0592119\\
3.1566	-0.0591291\\
3.1592	-0.0590458\\
3.1618	-0.0589654\\
3.1643	-0.0588909\\
3.1669	-0.0588163\\
3.1695	-0.0587448\\
3.172	-0.0586789\\
3.1746	-0.0586133\\
3.1772	-0.058551\\
3.1797	-0.058494\\
3.1823	-0.0584378\\
3.1848	-0.0583869\\
3.1874	-0.0583371\\
3.19	-0.0582906\\
3.1925	-0.0582489\\
3.1951	-0.0582089\\
3.1977	-0.0581722\\
3.2002	-0.0581402\\
3.2028	-0.0581101\\
3.2054	-0.0580835\\
3.2079	-0.0580611\\
3.2105	-0.0580411\\
3.213	-0.0580252\\
3.2156	-0.0580121\\
3.2182	-0.0580023\\
3.2207	-0.0579963\\
3.2233	-0.0579934\\
3.2259	-0.0579939\\
3.2284	-0.0579978\\
3.231	-0.0580052\\
3.2336	-0.058016\\
3.2361	-0.0580298\\
3.2387	-0.0580474\\
3.2412	-0.0580677\\
3.2438	-0.0580921\\
3.2464	-0.05812\\
3.2489	-0.0581501\\
3.2515	-0.0581847\\
3.2541	-0.0582227\\
3.2566	-0.0582624\\
3.2592	-0.058307\\
3.2618	-0.058355\\
3.2643	-0.0584043\\
3.2669	-0.0584588\\
3.2694	-0.0585142\\
3.272	-0.0585751\\
3.2746	-0.0586392\\
3.2771	-0.0587039\\
3.2797	-0.0587742\\
3.2823	-0.0588477\\
3.2848	-0.0589213\\
3.2874	-0.0590008\\
3.29	-0.0590833\\
3.2925	-0.0591655\\
3.2951	-0.0592539\\
3.2976	-0.0593416\\
3.3002	-0.0594357\\
3.3028	-0.0595325\\
3.3053	-0.0596283\\
3.3079	-0.0597305\\
3.3105	-0.0598355\\
3.313	-0.0599389\\
3.3156	-0.060049\\
3.3182	-0.0601616\\
3.3207	-0.0602722\\
3.3233	-0.0603896\\
3.3259	-0.0605094\\
3.3284	-0.0606268\\
3.331	-0.060751\\
3.3335	-0.0608726\\
3.3361	-0.0610012\\
3.3387	-0.0611318\\
3.3412	-0.0612593\\
3.3438	-0.0613938\\
3.3464	-0.0615302\\
3.3489	-0.0616631\\
3.3515	-0.0618031\\
3.3541	-0.0619447\\
3.3566	-0.0620824\\
3.3592	-0.0622272\\
3.3617	-0.0623678\\
3.3643	-0.0625155\\
3.3669	-0.0626645\\
3.3694	-0.062809\\
3.372	-0.0629604\\
3.3746	-0.063113\\
3.3771	-0.0632608\\
3.3797	-0.0634154\\
3.3823	-0.0635709\\
3.3848	-0.0637213\\
3.3874	-0.0638785\\
3.3899	-0.0640302\\
3.3925	-0.0641887\\
3.3951	-0.0643477\\
3.3976	-0.064501\\
3.4002	-0.0646609\\
3.4028	-0.0648211\\
3.4053	-0.0649753\\
3.4079	-0.0651359\\
3.4105	-0.0652966\\
3.413	-0.0654511\\
3.4156	-0.0656117\\
3.4181	-0.065766\\
3.4207	-0.0659263\\
3.4233	-0.0660863\\
3.4258	-0.0662398\\
3.4284	-0.066399\\
3.431	-0.0665576\\
3.4335	-0.0667096\\
3.4361	-0.066867\\
3.4387	-0.0670236\\
3.4412	-0.0671734\\
3.4438	-0.0673283\\
3.4463	-0.0674764\\
3.4489	-0.0676293\\
3.4515	-0.067781\\
3.454	-0.0679258\\
3.4566	-0.0680751\\
3.4592	-0.0682231\\
3.4617	-0.068364\\
3.4643	-0.0685091\\
3.4669	-0.0686526\\
3.4694	-0.068789\\
3.472	-0.0689292\\
3.4745	-0.0690623\\
3.4771	-0.0691989\\
3.4797	-0.0693336\\
3.4822	-0.0694612\\
3.4848	-0.0695919\\
3.4874	-0.0697205\\
3.4899	-0.069842\\
3.4925	-0.0699662\\
3.4951	-0.070088\\
3.4976	-0.0702029\\
3.5002	-0.07032\\
3.5028	-0.0704346\\
3.5053	-0.0705423\\
3.5079	-0.0706518\\
3.5104	-0.0707545\\
3.513	-0.0708586\\
3.5156	-0.0709599\\
3.5181	-0.0710547\\
3.5207	-0.0711503\\
3.5233	-0.071243\\
3.5258	-0.0713292\\
3.5284	-0.0714159\\
3.531	-0.0714995\\
3.5335	-0.0715768\\
3.5361	-0.0716542\\
3.5386	-0.0717254\\
3.5412	-0.0717963\\
3.5438	-0.0718638\\
3.5463	-0.0719255\\
3.5489	-0.0719863\\
3.5515	-0.0720436\\
3.554	-0.0720955\\
3.5566	-0.0721459\\
3.5592	-0.0721927\\
3.5617	-0.0722344\\
3.5643	-0.0722741\\
3.5668	-0.0723088\\
3.5694	-0.0723413\\
3.572	-0.0723701\\
3.5745	-0.0723942\\
3.5771	-0.0724155\\
3.5797	-0.0724331\\
3.5822	-0.0724463\\
3.5848	-0.0724563\\
3.5874	-0.0724625\\
3.5899	-0.0724647\\
3.5925	-0.0724632\\
3.595	-0.0724581\\
3.5976	-0.0724489\\
3.6002	-0.0724358\\
3.6027	-0.0724194\\
3.6053	-0.0723986\\
3.6079	-0.0723737\\
3.6104	-0.0723461\\
3.613	-0.0723135\\
3.6156	-0.0722769\\
3.6181	-0.0722379\\
3.6207	-0.0721935\\
3.6232	-0.0721471\\
3.6258	-0.0720949\\
3.6284	-0.0720388\\
3.6309	-0.0719811\\
3.6335	-0.0719172\\
3.6361	-0.0718493\\
3.6386	-0.0717804\\
3.6412	-0.0717049\\
3.6438	-0.0716255\\
3.6463	-0.0715455\\
3.6489	-0.0714586\\
3.6515	-0.0713677\\
3.654	-0.0712768\\
3.6566	-0.0711785\\
3.6591	-0.0710804\\
3.6617	-0.0709747\\
3.6643	-0.0708653\\
3.6668	-0.0707567\\
3.6694	-0.07064\\
3.672	-0.0705198\\
3.6745	-0.0704007\\
3.6771	-0.0702734\\
3.6797	-0.0701426\\
3.6822	-0.0700135\\
3.6848	-0.0698758\\
3.6873	-0.0697402\\
3.6899	-0.0695958\\
3.6925	-0.0694481\\
3.695	-0.069303\\
3.6976	-0.0691488\\
3.7002	-0.0689914\\
3.7027	-0.0688371\\
3.7053	-0.0686736\\
3.7079	-0.0685069\\
3.7104	-0.0683439\\
3.713	-0.0681713\\
3.7155	-0.0680027\\
3.7181	-0.0678244\\
3.7207	-0.0676434\\
3.7232	-0.0674666\\
3.7258	-0.0672802\\
3.7284	-0.0670911\\
3.7309	-0.0669068\\
3.7335	-0.0667126\\
3.7361	-0.0665159\\
3.7386	-0.0663245\\
3.7412	-0.0661231\\
3.7437	-0.0659273\\
3.7463	-0.0657215\\
3.7489	-0.0655134\\
3.7514	-0.0653113\\
3.754	-0.0650992\\
3.7566	-0.064885\\
3.7591	-0.0646772\\
3.7617	-0.0644593\\
3.7643	-0.0642396\\
3.7668	-0.0640267\\
3.7694	-0.0638036\\
3.7719	-0.0635876\\
3.7745	-0.0633615\\
3.7771	-0.0631338\\
3.7796	-0.0629137\\
3.7822	-0.0626833\\
3.7848	-0.0624518\\
3.7873	-0.062228\\
3.7899	-0.0619941\\
3.7925	-0.0617592\\
3.795	-0.0615324\\
3.7976	-0.0612957\\
3.8002	-0.061058\\
3.8027	-0.0608288\\
3.8053	-0.0605898\\
3.8078	-0.0603593\\
3.8104	-0.0601191\\
3.813	-0.0598784\\
3.8155	-0.0596465\\
3.8181	-0.059405\\
3.8207	-0.0591633\\
3.8232	-0.0589307\\
3.8258	-0.0586886\\
3.8284	-0.0584464\\
3.8309	-0.0582136\\
3.8335	-0.0579716\\
3.836	-0.057739\\
3.8386	-0.0574973\\
3.8412	-0.0572559\\
3.8437	-0.0570241\\
3.8463	-0.0567835\\
3.8489	-0.0565434\\
3.8514	-0.056313\\
3.854	-0.0560741\\
3.8566	-0.0558358\\
3.8591	-0.0556075\\
3.8617	-0.0553708\\
3.8642	-0.0551441\\
3.8668	-0.0549093\\
3.8694	-0.0546756\\
3.8719	-0.0544518\\
3.8745	-0.0542203\\
3.8771	-0.05399\\
3.8796	-0.0537697\\
3.8822	-0.053542\\
3.8848	-0.0533158\\
3.8873	-0.0530996\\
3.8899	-0.0528763\\
3.8924	-0.0526631\\
3.895	-0.052443\\
3.8976	-0.0522247\\
3.9001	-0.0520164\\
3.9027	-0.0518016\\
3.9053	-0.0515887\\
3.9078	-0.0513858\\
3.9104	-0.0511768\\
3.913	-0.0509699\\
3.9155	-0.0507729\\
3.9181	-0.0505702\\
3.9206	-0.0503773\\
3.9232	-0.050179\\
3.9258	-0.049983\\
3.9283	-0.0497968\\
3.9309	-0.0496056\\
3.9335	-0.0494168\\
3.936	-0.0492376\\
3.9386	-0.0490538\\
3.9412	-0.0488726\\
3.9437	-0.0487009\\
3.9463	-0.0485249\\
3.9488	-0.0483583\\
3.9514	-0.0481877\\
3.954	-0.04802\\
3.9565	-0.0478614\\
3.9591	-0.0476993\\
3.9617	-0.0475402\\
3.9642	-0.0473899\\
3.9668	-0.0472366\\
3.9694	-0.0470863\\
3.9719	-0.0469447\\
3.9745	-0.0468005\\
3.9771	-0.0466593\\
3.9796	-0.0465266\\
3.9822	-0.0463917\\
3.9847	-0.046265\\
3.9873	-0.0461364\\
3.9899	-0.0460111\\
3.9924	-0.0458937\\
3.995	-0.0457748\\
3.9976	-0.0456592\\
4.0001	-0.0455513\\
4.0027	-0.0454423\\
4.0053	-0.0453367\\
4.0078	-0.0452383\\
4.0104	-0.0451394\\
4.0129	-0.0450475\\
4.0155	-0.0449552\\
4.0181	-0.0448665\\
4.0206	-0.0447844\\
4.0232	-0.0447024\\
4.0258	-0.0446238\\
4.0283	-0.0445516\\
4.0309	-0.0444799\\
4.0335	-0.0444117\\
4.036	-0.0443494\\
4.0386	-0.044288\\
4.0411	-0.0442323\\
4.0437	-0.0441778\\
4.0463	-0.0441267\\
4.0488	-0.0440809\\
4.0514	-0.0440367\\
4.054	-0.043996\\
4.0565	-0.0439601\\
4.0591	-0.0439262\\
4.0617	-0.0438957\\
4.0642	-0.0438697\\
4.0668	-0.0438459\\
4.0693	-0.0438263\\
4.0719	-0.0438093\\
4.0745	-0.0437957\\
4.077	-0.0437857\\
4.0796	-0.0437787\\
4.0822	-0.043775\\
4.0847	-0.0437746\\
4.0873	-0.0437775\\
4.0899	-0.0437836\\
4.0924	-0.0437925\\
4.095	-0.043805\\
4.0975	-0.0438201\\
4.1001	-0.0438389\\
4.1027	-0.0438608\\
4.1052	-0.0438849\\
4.1078	-0.043913\\
4.1104	-0.0439442\\
4.1129	-0.043977\\
4.1155	-0.0440142\\
4.1181	-0.0440543\\
4.1206	-0.0440957\\
4.1232	-0.0441416\\
4.1258	-0.0441904\\
4.1283	-0.0442401\\
4.1309	-0.0442945\\
4.1334	-0.0443494\\
4.136	-0.0444093\\
4.1386	-0.0444718\\
4.1411	-0.0445345\\
4.1437	-0.0446023\\
4.1463	-0.0446726\\
4.1488	-0.0447427\\
4.1514	-0.044818\\
4.154	-0.0448958\\
4.1565	-0.044973\\
4.1591	-0.0450555\\
4.1616	-0.0451371\\
4.1642	-0.0452242\\
4.1668	-0.0453135\\
4.1693	-0.0454015\\
4.1719	-0.0454951\\
4.1745	-0.0455909\\
4.177	-0.0456849\\
4.1796	-0.0457847\\
4.1822	-0.0458864\\
4.1847	-0.045986\\
4.1873	-0.0460915\\
4.1898	-0.0461946\\
4.1924	-0.0463035\\
4.195	-0.0464142\\
4.1975	-0.0465223\\
4.2001	-0.0466362\\
4.2027	-0.0467517\\
4.2052	-0.0468642\\
4.2078	-0.0469826\\
4.2104	-0.0471025\\
4.2129	-0.047219\\
4.2155	-0.0473415\\
4.218	-0.0474604\\
4.2206	-0.0475853\\
4.2232	-0.0477113\\
4.2257	-0.0478336\\
4.2283	-0.0479617\\
4.2309	-0.0480908\\
4.2334	-0.0482159\\
4.236	-0.0483468\\
4.2386	-0.0484785\\
4.2411	-0.048606\\
4.2437	-0.0487392\\
4.2462	-0.0488679\\
4.2488	-0.0490024\\
4.2514	-0.0491375\\
4.2539	-0.0492678\\
4.2565	-0.0494039\\
4.2591	-0.0495403\\
4.2616	-0.0496718\\
4.2642	-0.0498089\\
4.2668	-0.0499462\\
4.2693	-0.0500784\\
4.2719	-0.0502161\\
4.2744	-0.0503486\\
4.277	-0.0504864\\
4.2796	-0.0506242\\
4.2821	-0.0507566\\
4.2847	-0.0508943\\
4.2873	-0.0510318\\
4.2898	-0.0511638\\
4.2924	-0.0513008\\
4.295	-0.0514376\\
4.2975	-0.0515688\\
4.3001	-0.0517048\\
4.3027	-0.0518404\\
4.3052	-0.0519703\\
4.3078	-0.0521048\\
4.3103	-0.0522337\\
4.3129	-0.0523671\\
4.3155	-0.0524998\\
4.318	-0.0526268\\
4.3206	-0.052758\\
4.3232	-0.0528885\\
4.3257	-0.0530132\\
4.3283	-0.053142\\
4.3309	-0.0532699\\
4.3334	-0.053392\\
4.336	-0.0535179\\
4.3385	-0.0536381\\
4.3411	-0.0537619\\
4.3437	-0.0538847\\
4.3462	-0.0540017\\
4.3488	-0.0541222\\
4.3514	-0.0542414\\
4.3539	-0.0543549\\
4.3565	-0.0544717\\
4.3591	-0.0545871\\
4.3616	-0.0546969\\
4.3642	-0.0548097\\
4.3667	-0.0549168\\
4.3693	-0.0550267\\
4.3719	-0.0551352\\
4.3744	-0.0552382\\
4.377	-0.0553437\\
4.3796	-0.0554477\\
4.3821	-0.0555462\\
4.3847	-0.0556471\\
4.3873	-0.0557463\\
4.3898	-0.0558402\\
4.3924	-0.0559362\\
4.3949	-0.056027\\
4.3975	-0.0561196\\
4.4001	-0.0562106\\
4.4026	-0.0562964\\
4.4052	-0.0563839\\
4.4078	-0.0564697\\
4.4103	-0.0565504\\
4.4129	-0.0566326\\
4.4155	-0.056713\\
4.418	-0.0567886\\
4.4206	-0.0568654\\
4.4231	-0.0569375\\
4.4257	-0.0570106\\
4.4283	-0.0570819\\
4.4308	-0.0571486\\
4.4334	-0.0572162\\
4.436	-0.0572819\\
4.4385	-0.0573432\\
4.4411	-0.0574051\\
4.4437	-0.0574652\\
4.4462	-0.0575211\\
4.4488	-0.0575774\\
4.4514	-0.0576317\\
4.4539	-0.0576822\\
4.4565	-0.0577328\\
4.459	-0.0577796\\
4.4616	-0.0578265\\
4.4642	-0.0578714\\
4.4667	-0.0579129\\
4.4693	-0.0579541\\
4.4719	-0.0579934\\
4.4744	-0.0580295\\
4.477	-0.0580651\\
4.4796	-0.0580989\\
4.4821	-0.0581296\\
4.4847	-0.0581598\\
4.4872	-0.058187\\
4.4898	-0.0582135\\
4.4924	-0.0582383\\
4.4949	-0.0582603\\
4.4975	-0.0582816\\
4.5001	-0.058301\\
4.5026	-0.058318\\
4.5052	-0.058334\\
4.5078	-0.0583483\\
4.5103	-0.0583604\\
4.5129	-0.0583714\\
4.5154	-0.0583803\\
4.518	-0.058388\\
4.5206	-0.0583941\\
4.5231	-0.0583984\\
4.5257	-0.0584013\\
4.5283	-0.0584027\\
4.5308	-0.0584026\\
4.5334	-0.0584009\\
4.536	-0.0583978\\
4.5385	-0.0583934\\
4.5411	-0.0583874\\
4.5436	-0.0583803\\
4.5462	-0.0583715\\
4.5488	-0.0583614\\
4.5513	-0.0583504\\
4.5539	-0.0583377\\
4.5565	-0.0583237\\
4.559	-0.0583091\\
4.5616	-0.0582926\\
4.5642	-0.058275\\
4.5667	-0.058257\\
4.5693	-0.0582371\\
4.5718	-0.0582169\\
4.5744	-0.0581949\\
4.577	-0.0581718\\
4.5795	-0.0581487\\
4.5821	-0.0581237\\
4.5847	-0.0580977\\
4.5872	-0.0580719\\
4.5898	-0.0580442\\
4.5924	-0.0580157\\
4.5949	-0.0579874\\
4.5975	-0.0579573\\
4.6001	-0.0579265\\
4.6026	-0.0578961\\
4.6052	-0.0578639\\
4.6077	-0.0578323\\
4.6103	-0.0577989\\
4.6129	-0.0577649\\
4.6154	-0.0577317\\
4.618	-0.0576967\\
4.6206	-0.0576612\\
4.6231	-0.0576267\\
4.6257	-0.0575904\\
4.6283	-0.0575538\\
4.6308	-0.0575183\\
4.6334	-0.057481\\
4.6359	-0.057445\\
4.6385	-0.0574073\\
4.6411	-0.0573695\\
4.6436	-0.0573329\\
4.6462	-0.0572948\\
4.6488	-0.0572566\\
4.6513	-0.0572199\\
4.6539	-0.0571817\\
4.6565	-0.0571435\\
4.659	-0.0571068\\
4.6616	-0.0570688\\
4.6641	-0.0570324\\
4.6667	-0.0569947\\
4.6693	-0.0569572\\
4.6718	-0.0569213\\
4.6744	-0.0568843\\
4.677	-0.0568477\\
4.6795	-0.0568127\\
4.6821	-0.0567767\\
4.6847	-0.0567411\\
4.6872	-0.0567073\\
4.6898	-0.0566726\\
4.6923	-0.0566397\\
4.6949	-0.056606\\
4.6975	-0.0565729\\
4.7	-0.0565416\\
4.7026	-0.0565097\\
4.7052	-0.0564784\\
4.7077	-0.0564489\\
4.7103	-0.056419\\
4.7129	-0.0563898\\
4.7154	-0.0563625\\
4.718	-0.0563349\\
4.7205	-0.056309\\
4.7231	-0.056283\\
4.7257	-0.0562579\\
4.7282	-0.0562346\\
4.7308	-0.0562112\\
4.7334	-0.0561888\\
4.7359	-0.0561681\\
4.7385	-0.0561476\\
4.7411	-0.0561282\\
4.7436	-0.0561104\\
4.7462	-0.056093\\
4.7487	-0.0560772\\
4.7513	-0.0560619\\
4.7539	-0.0560478\\
4.7564	-0.0560352\\
4.759	-0.0560233\\
4.7616	-0.0560125\\
4.7641	-0.0560033\\
4.7667	-0.0559949\\
4.7693	-0.0559877\\
4.7718	-0.0559819\\
4.7744	-0.0559771\\
4.777	-0.0559736\\
4.7795	-0.0559714\\
4.7821	-0.0559704\\
4.7846	-0.0559707\\
4.7872	-0.0559722\\
4.7898	-0.0559751\\
4.7923	-0.0559791\\
4.7949	-0.0559845\\
4.7975	-0.0559913\\
4.8	-0.0559992\\
4.8026	-0.0560086\\
4.8052	-0.0560194\\
4.8077	-0.0560311\\
4.8103	-0.0560445\\
4.8128	-0.0560588\\
4.8154	-0.056075\\
4.818	-0.0560925\\
4.8205	-0.0561106\\
4.8231	-0.0561309\\
4.8257	-0.0561525\\
4.8282	-0.0561745\\
4.8308	-0.0561988\\
4.8334	-0.0562244\\
4.8359	-0.0562504\\
4.8385	-0.0562787\\
4.841	-0.0563072\\
4.8436	-0.0563381\\
4.8462	-0.0563704\\
4.8487	-0.0564027\\
4.8513	-0.0564376\\
4.8539	-0.0564738\\
4.8564	-0.0565098\\
4.859	-0.0565485\\
4.8616	-0.0565885\\
4.8641	-0.0566282\\
4.8667	-0.0566707\\
4.8692	-0.0567128\\
4.8718	-0.0567577\\
4.8744	-0.0568038\\
4.8769	-0.0568493\\
4.8795	-0.0568978\\
4.8821	-0.0569474\\
4.8846	-0.0569963\\
4.8872	-0.0570481\\
4.8898	-0.0571011\\
4.8923	-0.0571531\\
4.8949	-0.0572083\\
4.8974	-0.0572623\\
4.9	-0.0573194\\
4.9026	-0.0573776\\
4.9051	-0.0574345\\
4.9077	-0.0574946\\
4.9103	-0.0575557\\
4.9128	-0.0576152\\
4.9154	-0.0576781\\
4.918	-0.0577418\\
4.9205	-0.0578038\\
4.9231	-0.0578692\\
4.9257	-0.0579353\\
4.9282	-0.0579996\\
4.9308	-0.0580673\\
4.9333	-0.0581329\\
4.9359	-0.0582019\\
4.9385	-0.0582716\\
4.941	-0.0583391\\
4.9436	-0.05841\\
4.9462	-0.0584814\\
4.9487	-0.0585506\\
4.9513	-0.058623\\
4.9539	-0.058696\\
4.9564	-0.0587665\\
4.959	-0.0588403\\
4.9615	-0.0589116\\
4.9641	-0.0589861\\
4.9667	-0.0590609\\
4.9692	-0.0591331\\
4.9718	-0.0592085\\
4.9744	-0.059284\\
4.9769	-0.0593569\\
4.9795	-0.0594328\\
4.9821	-0.0595088\\
4.9846	-0.059582\\
4.9872	-0.0596582\\
4.9897	-0.0597315\\
4.9923	-0.0598077\\
4.9949	-0.0598838\\
4.9974	-0.059957\\
5	-0.0600329\\
};
\addplot [color=black, line width=1.4pt, forget plot]
  table[row sep=crcr]{%
-0.125	3.72652e-08\\
-0.1224	3.80477e-08\\
-0.1199	3.88002e-08\\
-0.1173	3.95829e-08\\
-0.1147	4.03657e-08\\
-0.1122	4.11186e-08\\
-0.1096	4.19017e-08\\
-0.1071	4.26548e-08\\
-0.1045	4.34382e-08\\
-0.1019	4.42217e-08\\
-0.0994	4.49752e-08\\
-0.0968	4.5759e-08\\
-0.0942	4.6543e-08\\
-0.0917	4.72968e-08\\
-0.0891	4.8081e-08\\
-0.0865	4.88653e-08\\
-0.084	4.96195e-08\\
-0.0814	5.04038e-08\\
-0.0789	5.11582e-08\\
-0.0763	5.19429e-08\\
-0.0737	5.27276e-08\\
-0.0712	5.34822e-08\\
-0.0686	5.42671e-08\\
-0.066	5.50521e-08\\
-0.0635	5.5807e-08\\
-0.0609	5.65923e-08\\
-0.0583	5.73776e-08\\
-0.0558	5.81327e-08\\
-0.0532	5.89182e-08\\
-0.0507	5.96735e-08\\
-0.0481	6.04591e-08\\
-0.0455	6.12448e-08\\
-0.043	6.20003e-08\\
-0.0404	6.2786e-08\\
-0.0378	6.35717e-08\\
-0.0353	6.43274e-08\\
-0.0327	6.51133e-08\\
-0.0301	6.58992e-08\\
-0.0276	6.6655e-08\\
-0.025	6.74411e-08\\
-0.0224	6.82272e-08\\
-0.0199	6.89831e-08\\
-0.0173	6.97692e-08\\
-0.0148	7.05252e-08\\
-0.0122	7.13114e-08\\
-0.0096	7.20976e-08\\
-0.0071	7.28536e-08\\
-0.0045	7.36398e-08\\
-0.0019	7.44261e-08\\
0.0006	7.51821e-08\\
0.0032	0.000234051\\
0.0058	0.000819684\\
0.0083	0.00150582\\
0.0109	0.00240083\\
0.0134	0.00337091\\
0.016	0.00434857\\
0.0186	0.00520206\\
0.0211	0.00591136\\
0.0237	0.00659374\\
0.0263	0.00727125\\
0.0288	0.00792842\\
0.0314	0.00860013\\
0.034	0.00923959\\
0.0365	0.0098166\\
0.0391	0.0103797\\
0.0416	0.01089\\
0.0442	0.0113926\\
0.0468	0.0118716\\
0.0493	0.0123158\\
0.0519	0.0127642\\
0.0545	0.0131958\\
0.057	0.0135877\\
0.0596	0.0139657\\
0.0622	0.0143168\\
0.0647	0.0146394\\
0.0673	0.0149702\\
0.0698	0.0152859\\
0.0724	0.0156046\\
0.075	0.0159035\\
0.0775	0.0161695\\
0.0801	0.0164314\\
0.0827	0.0166912\\
0.0852	0.0169463\\
0.0878	0.017214\\
0.0904	0.0174743\\
0.0929	0.0177104\\
0.0955	0.0179434\\
0.098	0.0181653\\
0.1006	0.0184042\\
0.1032	0.0186533\\
0.1057	0.0188957\\
0.1083	0.0191414\\
0.1109	0.0193771\\
0.1134	0.0196003\\
0.116	0.0198392\\
0.1186	0.020091\\
0.1211	0.0203423\\
0.1237	0.020604\\
0.1263	0.0208593\\
0.1288	0.0210996\\
0.1314	0.0213525\\
0.1339	0.0216061\\
0.1365	0.0218818\\
0.1391	0.0221631\\
0.1416	0.0224305\\
0.1442	0.0227021\\
0.1468	0.0229719\\
0.1493	0.0232372\\
0.1519	0.0235239\\
0.1545	0.0238183\\
0.157	0.024101\\
0.1596	0.0243884\\
0.1621	0.0246588\\
0.1647	0.0249404\\
0.1673	0.0252288\\
0.1698	0.0255132\\
0.1724	0.0258107\\
0.175	0.0261022\\
0.1775	0.0263741\\
0.1801	0.0266516\\
0.1827	0.0269307\\
0.1852	0.0272039\\
0.1878	0.0274905\\
0.1903	0.027762\\
0.1929	0.0280347\\
0.1955	0.0282979\\
0.198	0.0285473\\
0.2006	0.0288083\\
0.2032	0.0290714\\
0.2057	0.0293219\\
0.2083	0.0295735\\
0.2109	0.0298137\\
0.2134	0.0300369\\
0.216	0.0302666\\
0.2185	0.0304884\\
0.2211	0.0307178\\
0.2237	0.0309403\\
0.2262	0.0311434\\
0.2288	0.0313441\\
0.2314	0.0315388\\
0.2339	0.0317249\\
0.2365	0.031918\\
0.2391	0.0321067\\
0.2416	0.0322791\\
0.2442	0.0324471\\
0.2467	0.0326007\\
0.2493	0.032757\\
0.2519	0.032913\\
0.2544	0.033061\\
0.257	0.0332084\\
0.2596	0.0333455\\
0.2621	0.0334684\\
0.2647	0.0335909\\
0.2673	0.0337125\\
0.2698	0.033829\\
0.2724	0.0339466\\
0.2749	0.0340523\\
0.2775	0.0341532\\
0.2801	0.0342475\\
0.2826	0.0343362\\
0.2852	0.0344289\\
0.2878	0.0345206\\
0.2903	0.0346043\\
0.2929	0.0346839\\
0.2955	0.0347565\\
0.298	0.0348231\\
0.3006	0.0348928\\
0.3032	0.0349634\\
0.3057	0.0350294\\
0.3083	0.0350929\\
0.3108	0.0351478\\
0.3134	0.0352007\\
0.316	0.0352531\\
0.3185	0.0353051\\
0.3211	0.0353596\\
0.3237	0.0354112\\
0.3262	0.035456\\
0.3288	0.0354979\\
0.3314	0.0355382\\
0.3339	0.0355783\\
0.3365	0.0356218\\
0.339	0.0356631\\
0.3416	0.0357026\\
0.3442	0.0357375\\
0.3467	0.0357686\\
0.3493	0.0358013\\
0.3519	0.0358361\\
0.3544	0.0358706\\
0.357	0.0359048\\
0.3596	0.0359351\\
0.3621	0.035961\\
0.3647	0.0359872\\
0.3672	0.0360138\\
0.3698	0.0360433\\
0.3724	0.0360727\\
0.3749	0.0360983\\
0.3775	0.0361214\\
0.3801	0.0361424\\
0.3826	0.036163\\
0.3852	0.0361863\\
0.3878	0.0362104\\
0.3903	0.0362322\\
0.3929	0.0362516\\
0.3954	0.0362674\\
0.398	0.036283\\
0.4006	0.0362997\\
0.4031	0.0363171\\
0.4057	0.0363348\\
0.4083	0.0363499\\
0.4108	0.0363613\\
0.4134	0.0363711\\
0.416	0.036381\\
0.4185	0.0363917\\
0.4211	0.0364032\\
0.4236	0.0364127\\
0.4262	0.0364195\\
0.4288	0.0364235\\
0.4313	0.0364262\\
0.4339	0.0364297\\
0.4365	0.036434\\
0.439	0.0364374\\
0.4416	0.0364385\\
0.4442	0.0364364\\
0.4467	0.0364325\\
0.4493	0.0364281\\
0.4519	0.0364244\\
0.4544	0.0364209\\
0.457	0.0364156\\
0.4595	0.0364078\\
0.4621	0.0363969\\
0.4647	0.0363849\\
0.4672	0.0363736\\
0.4698	0.0363622\\
0.4724	0.0363501\\
0.4749	0.0363361\\
0.4775	0.0363189\\
0.4801	0.0362996\\
0.4826	0.0362808\\
0.4852	0.0362617\\
0.4877	0.0362432\\
0.4903	0.0362225\\
0.4929	0.0361991\\
0.4954	0.0361744\\
0.498	0.0361477\\
0.5006	0.0361212\\
0.5031	0.0360959\\
0.5057	0.0360689\\
0.5083	0.0360399\\
0.5108	0.0360097\\
0.5134	0.0359767\\
0.5159	0.0359447\\
0.5185	0.0359117\\
0.5211	0.0358787\\
0.5236	0.0358457\\
0.5262	0.0358093\\
0.5288	0.0357708\\
0.5313	0.0357331\\
0.5339	0.035694\\
0.5365	0.0356551\\
0.539	0.0356173\\
0.5416	0.0355762\\
0.5441	0.0355349\\
0.5467	0.0354905\\
0.5493	0.0354459\\
0.5518	0.0354034\\
0.5544	0.0353591\\
0.557	0.0353138\\
0.5595	0.0352685\\
0.5621	0.0352198\\
0.5647	0.0351703\\
0.5672	0.035123\\
0.5698	0.035074\\
0.5723	0.0350264\\
0.5749	0.0349756\\
0.5775	0.034923\\
0.58	0.0348714\\
0.5826	0.0348174\\
0.5852	0.0347638\\
0.5877	0.0347122\\
0.5903	0.0346577\\
0.5929	0.0346015\\
0.5954	0.0345461\\
0.598	0.0344879\\
0.6006	0.0344298\\
0.6031	0.034374\\
0.6057	0.0343156\\
0.6082	0.0342582\\
0.6108	0.0341969\\
0.6134	0.0341346\\
0.6159	0.0340744\\
0.6185	0.034012\\
0.6211	0.0339494\\
0.6236	0.0338883\\
0.6262	0.0338234\\
0.6288	0.0337571\\
0.6313	0.0336926\\
0.6339	0.0336256\\
0.6364	0.033561\\
0.639	0.0334934\\
0.6416	0.0334245\\
0.6441	0.0333568\\
0.6467	0.0332854\\
0.6493	0.0332135\\
0.6518	0.0331444\\
0.6544	0.0330721\\
0.657	0.0329989\\
0.6595	0.0329272\\
0.6621	0.0328512\\
0.6646	0.0327774\\
0.6672	0.0327003\\
0.6698	0.032623\\
0.6723	0.0325481\\
0.6749	0.0324688\\
0.6775	0.0323881\\
0.68	0.0323095\\
0.6826	0.032227\\
0.6852	0.0321443\\
0.6877	0.0320643\\
0.6903	0.0319801\\
0.6928	0.0318977\\
0.6954	0.0318107\\
0.698	0.0317227\\
0.7005	0.0316377\\
0.7031	0.0315489\\
0.7057	0.0314592\\
0.7082	0.0313717\\
0.7108	0.0312792\\
0.7134	0.0311856\\
0.7159	0.0310948\\
0.7185	0.0309999\\
0.721	0.0309081\\
0.7236	0.0308114\\
0.7262	0.0307133\\
0.7287	0.0306176\\
0.7313	0.0305171\\
0.7339	0.0304159\\
0.7364	0.030318\\
0.739	0.0302153\\
0.7416	0.0301112\\
0.7441	0.0300097\\
0.7467	0.0299029\\
0.7492	0.0297993\\
0.7518	0.029691\\
0.7544	0.0295818\\
0.7569	0.0294757\\
0.7595	0.0293639\\
0.7621	0.0292506\\
0.7646	0.0291406\\
0.7672	0.0290254\\
0.7698	0.0289095\\
0.7723	0.028797\\
0.7749	0.0286786\\
0.7775	0.0285586\\
0.78	0.028442\\
0.7826	0.0283197\\
0.7851	0.0282014\\
0.7877	0.0280773\\
0.7903	0.027952\\
0.7928	0.02783\\
0.7954	0.0277017\\
0.798	0.0275721\\
0.8005	0.0274465\\
0.8031	0.027315\\
0.8057	0.0271823\\
0.8082	0.0270532\\
0.8108	0.0269175\\
0.8133	0.0267855\\
0.8159	0.0266471\\
0.8185	0.0265077\\
0.821	0.0263725\\
0.8236	0.0262306\\
0.8262	0.0260869\\
0.8287	0.0259473\\
0.8313	0.0258006\\
0.8339	0.0256528\\
0.8364	0.0255095\\
0.839	0.0253591\\
0.8415	0.0252129\\
0.8441	0.0250591\\
0.8467	0.0249037\\
0.8492	0.0247529\\
0.8518	0.0245948\\
0.8544	0.0244353\\
0.8569	0.0242803\\
0.8595	0.0241173\\
0.8621	0.0239525\\
0.8646	0.0237924\\
0.8672	0.0236244\\
0.8697	0.0234615\\
0.8723	0.0232904\\
0.8749	0.0231174\\
0.8774	0.0229491\\
0.88	0.0227722\\
0.8826	0.0225935\\
0.8851	0.0224201\\
0.8877	0.022238\\
0.8903	0.022054\\
0.8928	0.0218749\\
0.8954	0.0216866\\
0.8979	0.0215036\\
0.9005	0.0213114\\
0.9031	0.0211173\\
0.9056	0.0209287\\
0.9082	0.0207303\\
0.9108	0.0205295\\
0.9133	0.0203343\\
0.9159	0.0201291\\
0.9185	0.0199217\\
0.921	0.0197203\\
0.9236	0.0195084\\
0.9262	0.0192939\\
0.9287	0.0190852\\
0.9313	0.0188657\\
0.9338	0.0186525\\
0.9364	0.0184283\\
0.939	0.0182015\\
0.9415	0.0179809\\
0.9441	0.0177487\\
0.9467	0.0175137\\
0.9492	0.0172852\\
0.9518	0.0170449\\
0.9544	0.0168019\\
0.9569	0.0165655\\
0.9595	0.0163167\\
0.962	0.0160745\\
0.9646	0.0158197\\
0.9672	0.0155619\\
0.9697	0.0153113\\
0.9723	0.0150476\\
0.9749	0.0147807\\
0.9774	0.014521\\
0.98	0.0142475\\
0.9826	0.0139709\\
0.9851	0.0137019\\
0.9877	0.0134189\\
0.9902	0.0131436\\
0.9928	0.0128538\\
0.9954	0.0125604\\
0.9979	0.012275\\
1.0005	0.0119748\\
1.0031	0.0116711\\
1.0056	0.0113757\\
1.0082	0.0110648\\
1.0108	0.0107502\\
1.0133	0.0104441\\
1.0159	0.0101221\\
1.0184	0.00980896\\
1.021	0.00947965\\
1.0236	0.00914648\\
1.0261	0.00882236\\
1.0287	0.00848136\\
1.0313	0.00813641\\
1.0338	0.00780104\\
1.0364	0.00744843\\
1.039	0.00709181\\
1.0415	0.00674503\\
1.0441	0.0063803\\
1.0466	0.00602567\\
1.0492	0.00565286\\
1.0518	0.00527598\\
1.0543	0.00490972\\
1.0569	0.00452469\\
1.0595	0.00413539\\
1.062	0.00375701\\
1.0646	0.00335933\\
1.0672	0.00295747\\
1.0697	0.00256711\\
1.0723	0.00215697\\
1.0748	0.00175853\\
1.0774	0.00133985\\
1.08	0.000916841\\
1.0825	0.000506065\\
1.0851	7.46979e-05\\
1.0877	-0.000360932\\
1.0902	-0.000783888\\
1.0928	-0.00122806\\
1.0954	-0.0016766\\
1.0979	-0.00211195\\
1.1005	-0.0025689\\
1.1031	-0.00303008\\
1.1056	-0.00347755\\
1.1082	-0.00394715\\
1.1107	-0.00440278\\
1.1133	-0.00488084\\
1.1159	-0.00536312\\
1.1184	-0.00583078\\
1.121	-0.00632123\\
1.1236	-0.00681587\\
1.1261	-0.00729546\\
1.1287	-0.00779834\\
1.1313	-0.00830534\\
1.1338	-0.00879666\\
1.1364	-0.00931155\\
1.1389	-0.00981042\\
1.1415	-0.0103332\\
1.1441	-0.01086\\
1.1466	-0.0113702\\
1.1492	-0.0119047\\
1.1518	-0.0124429\\
1.1543	-0.012964\\
1.1569	-0.0135096\\
1.1595	-0.014059\\
1.162	-0.0145907\\
1.1646	-0.0151473\\
1.1671	-0.0156857\\
1.1697	-0.0162491\\
1.1723	-0.016816\\
1.1748	-0.0173642\\
1.1774	-0.0179377\\
1.18	-0.0185145\\
1.1825	-0.0190721\\
1.1851	-0.0196551\\
1.1877	-0.0202411\\
1.1902	-0.0208074\\
1.1928	-0.0213993\\
1.1953	-0.0219712\\
1.1979	-0.0225686\\
1.2005	-0.0231688\\
1.203	-0.0237483\\
1.2056	-0.0243535\\
1.2082	-0.0249612\\
1.2107	-0.0255478\\
1.2133	-0.0261601\\
1.2159	-0.0267747\\
1.2184	-0.0273676\\
1.221	-0.0279862\\
1.2235	-0.0285828\\
1.2261	-0.0292051\\
1.2287	-0.0298293\\
1.2312	-0.030431\\
1.2338	-0.0310582\\
1.2364	-0.0316869\\
1.2389	-0.0322927\\
1.2415	-0.032924\\
1.2441	-0.0335564\\
1.2466	-0.0341656\\
1.2492	-0.0348\\
1.2518	-0.0354352\\
1.2543	-0.0360466\\
1.2569	-0.0366831\\
1.2594	-0.0372956\\
1.262	-0.0379329\\
1.2646	-0.0385706\\
1.2671	-0.0391838\\
1.2697	-0.0398216\\
1.2723	-0.0404593\\
1.2748	-0.0410723\\
1.2774	-0.0417095\\
1.28	-0.0423462\\
1.2825	-0.0429579\\
1.2851	-0.0435934\\
1.2876	-0.0442036\\
1.2902	-0.0448373\\
1.2928	-0.04547\\
1.2953	-0.0460772\\
1.2979	-0.0467074\\
1.3005	-0.0473362\\
1.303	-0.0479392\\
1.3056	-0.0485647\\
1.3082	-0.0491884\\
1.3107	-0.0497862\\
1.3133	-0.0504059\\
1.3158	-0.0509997\\
1.3184	-0.0516149\\
1.321	-0.0522276\\
1.3235	-0.0528143\\
1.3261	-0.0534217\\
1.3287	-0.0540263\\
1.3312	-0.0546049\\
1.3338	-0.0552035\\
1.3364	-0.0557988\\
1.3389	-0.0563681\\
1.3415	-0.0569566\\
1.344	-0.0575192\\
1.3466	-0.0581006\\
1.3492	-0.058678\\
1.3517	-0.0592295\\
1.3543	-0.059799\\
1.3569	-0.0603643\\
1.3594	-0.0609037\\
1.362	-0.0614603\\
1.3646	-0.0620123\\
1.3671	-0.0625386\\
1.3697	-0.0630812\\
1.3722	-0.0635983\\
1.3748	-0.0641311\\
1.3774	-0.0646587\\
1.3799	-0.0651611\\
1.3825	-0.0656783\\
1.3851	-0.0661901\\
1.3876	-0.0666768\\
1.3902	-0.0671775\\
1.3928	-0.0676723\\
1.3953	-0.0681426\\
1.3979	-0.0686259\\
1.4005	-0.069103\\
1.403	-0.069556\\
1.4056	-0.070021\\
1.4081	-0.0704622\\
1.4107	-0.0709147\\
1.4133	-0.0713607\\
1.4158	-0.0717833\\
1.4184	-0.0722164\\
1.421	-0.0726427\\
1.4235	-0.0730462\\
1.4261	-0.0734591\\
1.4287	-0.0738651\\
1.4312	-0.074249\\
1.4338	-0.0746412\\
1.4363	-0.0750117\\
1.4389	-0.07539\\
1.4415	-0.0757611\\
1.444	-0.0761111\\
1.4466	-0.0764679\\
1.4492	-0.0768174\\
1.4517	-0.0771466\\
1.4543	-0.0774816\\
1.4569	-0.0778092\\
1.4594	-0.0781171\\
1.462	-0.0784301\\
1.4645	-0.0787239\\
1.4671	-0.0790221\\
1.4697	-0.0793127\\
1.4722	-0.079585\\
1.4748	-0.0798608\\
1.4774	-0.080129\\
1.4799	-0.0803797\\
1.4825	-0.080633\\
1.4851	-0.0808788\\
1.4876	-0.0811079\\
1.4902	-0.0813388\\
1.4927	-0.0815536\\
1.4953	-0.0817697\\
1.4979	-0.0819782\\
1.5004	-0.0821716\\
1.503	-0.0823653\\
1.5056	-0.0825516\\
1.5081	-0.0827238\\
1.5107	-0.0828955\\
1.5133	-0.0830598\\
1.5158	-0.0832109\\
1.5184	-0.0833609\\
1.5209	-0.0834982\\
1.5235	-0.0836339\\
1.5261	-0.0837624\\
1.5286	-0.0838793\\
1.5312	-0.0839939\\
1.5338	-0.0841014\\
1.5363	-0.0841982\\
1.5389	-0.084292\\
1.5415	-0.084379\\
1.544	-0.0844563\\
1.5466	-0.0845299\\
1.5491	-0.0845945\\
1.5517	-0.0846551\\
1.5543	-0.0847092\\
1.5568	-0.084755\\
1.5594	-0.0847965\\
1.562	-0.0848315\\
1.5645	-0.0848593\\
1.5671	-0.0848822\\
1.5697	-0.0848989\\
1.5722	-0.0849093\\
1.5748	-0.0849143\\
1.5774	-0.0849133\\
1.5799	-0.084907\\
1.5825	-0.0848948\\
1.585	-0.0848778\\
1.5876	-0.0848547\\
1.5902	-0.0848261\\
1.5927	-0.0847936\\
1.5953	-0.0847546\\
1.5979	-0.0847105\\
1.6004	-0.0846632\\
1.603	-0.0846091\\
1.6056	-0.0845502\\
1.6081	-0.084489\\
1.6107	-0.0844207\\
1.6132	-0.0843507\\
1.6158	-0.0842735\\
1.6184	-0.0841919\\
1.6209	-0.0841093\\
1.6235	-0.0840193\\
1.6261	-0.0839251\\
1.6286	-0.0838308\\
1.6312	-0.0837289\\
1.6338	-0.0836231\\
1.6363	-0.0835179\\
1.6389	-0.0834049\\
1.6414	-0.083293\\
1.644	-0.0831732\\
1.6466	-0.08305\\
1.6491	-0.0829286\\
1.6517	-0.0827993\\
1.6543	-0.0826669\\
1.6568	-0.0825368\\
1.6594	-0.0823988\\
1.662	-0.0822581\\
1.6645	-0.0821203\\
1.6671	-0.0819745\\
1.6696	-0.0818321\\
1.6722	-0.0816817\\
1.6748	-0.081529\\
1.6773	-0.0813802\\
1.6799	-0.0812235\\
1.6825	-0.0810649\\
1.685	-0.0809106\\
1.6876	-0.0807485\\
1.6902	-0.0805847\\
1.6927	-0.0804258\\
1.6953	-0.0802591\\
1.6978	-0.0800976\\
1.7004	-0.0799284\\
1.703	-0.079758\\
1.7055	-0.0795932\\
1.7081	-0.0794208\\
1.7107	-0.0792476\\
1.7132	-0.0790803\\
1.7158	-0.0789057\\
1.7184	-0.0787304\\
1.7209	-0.0785614\\
1.7235	-0.0783852\\
1.7261	-0.0782088\\
1.7286	-0.0780388\\
1.7312	-0.077862\\
1.7337	-0.0776918\\
1.7363	-0.0775149\\
1.7389	-0.0773381\\
1.7414	-0.0771683\\
1.744	-0.0769919\\
1.7466	-0.0768159\\
1.7491	-0.0766471\\
1.7517	-0.076472\\
1.7543	-0.0762976\\
1.7568	-0.0761304\\
1.7594	-0.0759574\\
1.7619	-0.0757917\\
1.7645	-0.0756204\\
1.7671	-0.07545\\
1.7696	-0.0752871\\
1.7722	-0.0751189\\
1.7748	-0.0749518\\
1.7773	-0.0747923\\
1.7799	-0.0746278\\
1.7825	-0.0744647\\
1.785	-0.0743093\\
1.7876	-0.0741491\\
1.7901	-0.0739966\\
1.7927	-0.0738397\\
1.7953	-0.0736845\\
1.7978	-0.0735369\\
1.8004	-0.0733852\\
1.803	-0.0732355\\
1.8055	-0.0730934\\
1.8081	-0.0729475\\
1.8107	-0.0728038\\
1.8132	-0.0726676\\
1.8158	-0.0725281\\
1.8183	-0.072396\\
1.8209	-0.072261\\
1.8235	-0.0721283\\
1.826	-0.0720029\\
1.8286	-0.0718749\\
1.8312	-0.0717494\\
1.8337	-0.0716311\\
1.8363	-0.0715106\\
1.8389	-0.0713927\\
1.8414	-0.0712819\\
1.844	-0.0711692\\
1.8465	-0.0710635\\
1.8491	-0.0709563\\
1.8517	-0.0708518\\
1.8542	-0.0707541\\
1.8568	-0.0706553\\
1.8594	-0.0705594\\
1.8619	-0.0704699\\
1.8645	-0.0703798\\
1.8671	-0.0702927\\
1.8696	-0.0702118\\
1.8722	-0.0701306\\
1.8747	-0.0700555\\
1.8773	-0.0699804\\
1.8799	-0.0699084\\
1.8824	-0.0698422\\
1.885	-0.0697764\\
1.8876	-0.0697138\\
1.8901	-0.0696566\\
1.8927	-0.0696003\\
1.8953	-0.0695472\\
1.8978	-0.0694992\\
1.9004	-0.0694525\\
1.903	-0.0694091\\
1.9055	-0.0693704\\
1.9081	-0.0693334\\
1.9106	-0.0693009\\
1.9132	-0.0692704\\
1.9158	-0.0692432\\
1.9183	-0.0692201\\
1.9209	-0.0691993\\
1.9235	-0.0691819\\
1.926	-0.0691682\\
1.9286	-0.0691573\\
1.9312	-0.0691496\\
1.9337	-0.0691454\\
1.9363	-0.0691442\\
1.9388	-0.0691462\\
1.9414	-0.0691515\\
1.944	-0.06916\\
1.9465	-0.0691713\\
1.9491	-0.0691862\\
1.9517	-0.0692044\\
1.9542	-0.0692249\\
1.9568	-0.0692494\\
1.9594	-0.0692771\\
1.9619	-0.0693067\\
1.9645	-0.0693406\\
1.967	-0.0693761\\
1.9696	-0.0694161\\
1.9722	-0.0694593\\
1.9747	-0.0695036\\
1.9773	-0.0695527\\
1.9799	-0.0696049\\
1.9824	-0.0696578\\
1.985	-0.0697158\\
1.9876	-0.0697767\\
1.9901	-0.0698381\\
1.9927	-0.0699047\\
1.9952	-0.0699714\\
1.9978	-0.0700435\\
2.0004	-0.0701184\\
2.0029	-0.0701931\\
2.0055	-0.0702733\\
2.0081	-0.0703562\\
2.0106	-0.0704384\\
2.0132	-0.0705264\\
2.0158	-0.070617\\
2.0183	-0.0707064\\
2.0209	-0.0708018\\
2.0234	-0.0708959\\
2.026	-0.0709959\\
2.0286	-0.0710984\\
2.0311	-0.071199\\
2.0337	-0.0713058\\
2.0363	-0.0714148\\
2.0388	-0.0715215\\
2.0414	-0.0716346\\
2.044	-0.0717497\\
2.0465	-0.0718622\\
2.0491	-0.0719811\\
2.0517	-0.0721018\\
2.0542	-0.0722195\\
2.0568	-0.0723437\\
2.0593	-0.0724646\\
2.0619	-0.072592\\
2.0645	-0.0727209\\
2.067	-0.0728462\\
2.0696	-0.072978\\
2.0722	-0.0731111\\
2.0747	-0.0732403\\
2.0773	-0.0733759\\
2.0799	-0.0735126\\
2.0824	-0.0736451\\
2.085	-0.0737838\\
2.0875	-0.0739182\\
2.0901	-0.0740587\\
2.0927	-0.0742001\\
2.0952	-0.0743367\\
2.0978	-0.0744795\\
2.1004	-0.0746228\\
2.1029	-0.0747611\\
2.1055	-0.0749053\\
2.1081	-0.0750499\\
2.1106	-0.0751892\\
2.1132	-0.0753343\\
2.1157	-0.0754739\\
2.1183	-0.0756192\\
2.1209	-0.0757644\\
2.1234	-0.075904\\
2.126	-0.0760489\\
2.1286	-0.0761935\\
2.1311	-0.0763323\\
2.1337	-0.0764762\\
2.1363	-0.0766196\\
2.1388	-0.0767568\\
2.1414	-0.076899\\
2.1439	-0.0770349\\
2.1465	-0.0771754\\
2.1491	-0.077315\\
2.1516	-0.0774483\\
2.1542	-0.0775858\\
2.1568	-0.0777221\\
2.1593	-0.077852\\
2.1619	-0.0779858\\
2.1645	-0.0781181\\
2.167	-0.0782439\\
2.1696	-0.0783731\\
2.1721	-0.0784958\\
2.1747	-0.0786217\\
2.1773	-0.0787457\\
2.1798	-0.0788631\\
2.1824	-0.0789832\\
2.185	-0.0791012\\
2.1875	-0.0792126\\
2.1901	-0.0793262\\
2.1927	-0.0794374\\
2.1952	-0.0795421\\
2.1978	-0.0796485\\
2.2004	-0.0797523\\
2.2029	-0.0798496\\
2.2055	-0.079948\\
2.208	-0.0800401\\
2.2106	-0.0801329\\
2.2132	-0.0802227\\
2.2157	-0.0803062\\
2.2183	-0.08039\\
2.2209	-0.0804705\\
2.2234	-0.0805448\\
2.226	-0.0806189\\
2.2286	-0.0806894\\
2.2311	-0.080754\\
2.2337	-0.0808177\\
2.2362	-0.0808755\\
2.2388	-0.0809319\\
2.2414	-0.0809846\\
2.2439	-0.0810317\\
2.2465	-0.0810768\\
2.2491	-0.081118\\
2.2516	-0.0811538\\
2.2542	-0.081187\\
2.2568	-0.0812162\\
2.2593	-0.0812402\\
2.2619	-0.0812611\\
2.2644	-0.0812771\\
2.267	-0.0812896\\
2.2696	-0.0812977\\
2.2721	-0.0813012\\
2.2747	-0.0813006\\
2.2773	-0.0812954\\
2.2798	-0.0812862\\
2.2824	-0.0812721\\
2.285	-0.0812533\\
2.2875	-0.0812309\\
2.2901	-0.081203\\
2.2926	-0.0811717\\
2.2952	-0.0811345\\
2.2978	-0.0810926\\
2.3003	-0.0810477\\
2.3029	-0.0809964\\
2.3055	-0.0809401\\
2.308	-0.0808815\\
2.3106	-0.0808158\\
2.3132	-0.0807451\\
2.3157	-0.0806726\\
2.3183	-0.0805924\\
2.3208	-0.0805106\\
2.3234	-0.0804207\\
2.326	-0.080326\\
2.3285	-0.0802302\\
2.3311	-0.0801258\\
2.3337	-0.0800165\\
2.3362	-0.0799068\\
2.3388	-0.0797879\\
2.3414	-0.0796641\\
2.3439	-0.0795405\\
2.3465	-0.0794072\\
2.349	-0.0792744\\
2.3516	-0.0791316\\
2.3542	-0.0789841\\
2.3567	-0.0788377\\
2.3593	-0.0786809\\
2.3619	-0.0785193\\
2.3644	-0.0783595\\
2.367	-0.0781888\\
2.3696	-0.0780135\\
2.3721	-0.0778406\\
2.3747	-0.0776564\\
2.3773	-0.0774676\\
2.3798	-0.0772819\\
2.3824	-0.0770845\\
2.3849	-0.0768906\\
2.3875	-0.0766847\\
2.3901	-0.0764745\\
2.3926	-0.0762685\\
2.3952	-0.0760501\\
2.3978	-0.0758277\\
2.4003	-0.07561\\
2.4029	-0.0753797\\
2.4055	-0.0751455\\
2.408	-0.0749167\\
2.4106	-0.074675\\
2.4131	-0.0744391\\
2.4157	-0.0741902\\
2.4183	-0.0739377\\
2.4208	-0.0736916\\
2.4234	-0.0734323\\
2.426	-0.0731696\\
2.4285	-0.0729139\\
2.4311	-0.0726448\\
2.4337	-0.0723726\\
2.4362	-0.072108\\
2.4388	-0.0718299\\
2.4413	-0.0715597\\
2.4439	-0.0712759\\
2.4465	-0.0709895\\
2.449	-0.0707115\\
2.4516	-0.0704199\\
2.4542	-0.0701258\\
2.4567	-0.0698407\\
2.4593	-0.069542\\
2.4619	-0.0692411\\
2.4644	-0.0689497\\
2.467	-0.0686447\\
2.4695	-0.0683495\\
2.4721	-0.0680407\\
2.4747	-0.0677301\\
2.4772	-0.0674299\\
2.4798	-0.067116\\
2.4824	-0.0668007\\
2.4849	-0.0664962\\
2.4875	-0.0661782\\
2.4901	-0.0658589\\
2.4926	-0.0655509\\
2.4952	-0.0652295\\
2.4977	-0.0649196\\
2.5003	-0.0645964\\
2.5029	-0.0642725\\
2.5054	-0.0639604\\
2.508	-0.0636352\\
2.5106	-0.0633095\\
2.5131	-0.062996\\
2.5157	-0.0626696\\
2.5183	-0.062343\\
2.5208	-0.0620289\\
2.5234	-0.0617022\\
2.526	-0.0613755\\
2.5285	-0.0610615\\
2.5311	-0.0607352\\
2.5336	-0.0604217\\
2.5362	-0.0600961\\
2.5388	-0.0597711\\
2.5413	-0.059459\\
2.5439	-0.0591352\\
2.5465	-0.0588121\\
2.549	-0.0585023\\
2.5516	-0.0581809\\
2.5542	-0.0578606\\
2.5567	-0.0575537\\
2.5593	-0.0572356\\
2.5618	-0.056931\\
2.5644	-0.0566155\\
2.567	-0.0563014\\
2.5695	-0.0560009\\
2.5721	-0.0556898\\
2.5747	-0.0553805\\
2.5772	-0.0550847\\
2.5798	-0.0547788\\
2.5824	-0.0544748\\
2.5849	-0.0541844\\
2.5875	-0.0538843\\
2.59	-0.0535978\\
2.5926	-0.0533019\\
2.5952	-0.0530084\\
2.5977	-0.0527282\\
2.6003	-0.0524392\\
2.6029	-0.0521527\\
2.6054	-0.0518796\\
2.608	-0.051598\\
2.6106	-0.0513191\\
2.6131	-0.0510535\\
2.6157	-0.0507799\\
2.6182	-0.0505196\\
2.6208	-0.0502516\\
2.6234	-0.0499866\\
2.6259	-0.0497346\\
2.6285	-0.0494755\\
2.6311	-0.0492195\\
2.6336	-0.0489763\\
2.6362	-0.0487265\\
2.6388	-0.0484799\\
2.6413	-0.048246\\
2.6439	-0.0480059\\
2.6464	-0.0477783\\
2.649	-0.0475449\\
2.6516	-0.047315\\
2.6541	-0.0470973\\
2.6567	-0.0468743\\
2.6593	-0.0466549\\
2.6618	-0.0464473\\
2.6644	-0.046235\\
2.667	-0.0460265\\
2.6695	-0.0458294\\
2.6721	-0.0456282\\
2.6746	-0.0454382\\
2.6772	-0.0452444\\
2.6798	-0.0450545\\
2.6823	-0.0448755\\
2.6849	-0.0446931\\
2.6875	-0.0445147\\
2.69	-0.0443468\\
2.6926	-0.0441761\\
2.6952	-0.0440094\\
2.6977	-0.0438528\\
2.7003	-0.043694\\
2.7029	-0.0435391\\
2.7054	-0.043394\\
2.708	-0.0432471\\
2.7105	-0.0431097\\
2.7131	-0.0429708\\
2.7157	-0.042836\\
2.7182	-0.0427102\\
2.7208	-0.0425834\\
2.7234	-0.0424607\\
2.7259	-0.0423466\\
2.7285	-0.042232\\
2.7311	-0.0421216\\
2.7336	-0.0420192\\
2.7362	-0.0419168\\
2.7387	-0.0418222\\
2.7413	-0.0417279\\
2.7439	-0.0416377\\
2.7464	-0.0415548\\
2.749	-0.0414726\\
2.7516	-0.0413946\\
2.7541	-0.0413234\\
2.7567	-0.0412533\\
2.7593	-0.0411873\\
2.7618	-0.0411277\\
2.7644	-0.0410697\\
2.7669	-0.0410177\\
2.7695	-0.0409676\\
2.7721	-0.0409215\\
2.7746	-0.0408809\\
2.7772	-0.0408426\\
2.7798	-0.0408082\\
2.7823	-0.0407789\\
2.7849	-0.0407522\\
2.7875	-0.0407295\\
2.79	-0.0407113\\
2.7926	-0.0406961\\
2.7951	-0.0406851\\
2.7977	-0.0406775\\
2.8003	-0.0406736\\
2.8028	-0.0406733\\
2.8054	-0.0406768\\
2.808	-0.0406839\\
2.8105	-0.0406943\\
2.8131	-0.0407086\\
2.8157	-0.0407265\\
2.8182	-0.0407472\\
2.8208	-0.0407721\\
2.8233	-0.0407994\\
2.8259	-0.0408312\\
2.8285	-0.0408665\\
2.831	-0.0409036\\
2.8336	-0.0409455\\
2.8362	-0.0409908\\
2.8387	-0.0410375\\
2.8413	-0.0410892\\
2.8439	-0.0411442\\
2.8464	-0.0412001\\
2.849	-0.0412614\\
2.8516	-0.0413257\\
2.8541	-0.0413905\\
2.8567	-0.0414608\\
2.8592	-0.0415313\\
2.8618	-0.0416075\\
2.8644	-0.0416866\\
2.8669	-0.0417654\\
2.8695	-0.0418501\\
2.8721	-0.0419376\\
2.8746	-0.0420243\\
2.8772	-0.0421172\\
2.8798	-0.0422126\\
2.8823	-0.0423069\\
2.8849	-0.0424074\\
2.8874	-0.0425065\\
2.89	-0.0426119\\
2.8926	-0.0427197\\
2.8951	-0.0428257\\
2.8977	-0.0429381\\
2.9003	-0.0430528\\
2.9028	-0.0431652\\
2.9054	-0.0432843\\
2.908	-0.0434054\\
2.9105	-0.0435238\\
2.9131	-0.043649\\
2.9156	-0.0437712\\
2.9182	-0.0439002\\
2.9208	-0.0440311\\
2.9233	-0.0441586\\
2.9259	-0.044293\\
2.9285	-0.0444291\\
2.931	-0.0445615\\
2.9336	-0.0447008\\
2.9362	-0.0448417\\
2.9387	-0.0449785\\
2.9413	-0.0451223\\
2.9438	-0.0452619\\
2.9464	-0.0454083\\
2.949	-0.0455561\\
2.9515	-0.0456993\\
2.9541	-0.0458494\\
2.9567	-0.0460007\\
2.9592	-0.0461472\\
2.9618	-0.0463005\\
2.9644	-0.0464548\\
2.9669	-0.0466041\\
2.9695	-0.0467601\\
2.972	-0.046911\\
2.9746	-0.0470686\\
2.9772	-0.0472269\\
2.9797	-0.0473797\\
2.9823	-0.0475393\\
2.9849	-0.0476994\\
2.9874	-0.0478538\\
2.99	-0.0480148\\
2.9926	-0.0481763\\
2.9951	-0.0483318\\
2.9977	-0.0484938\\
3.0003	-0.0486561\\
3.0028	-0.0488123\\
3.0054	-0.0489749\\
3.0079	-0.0491313\\
3.0105	-0.049294\\
3.0131	-0.0494567\\
3.0156	-0.049613\\
3.0182	-0.0497756\\
3.0208	-0.0499379\\
3.0233	-0.0500938\\
3.0259	-0.0502557\\
3.0285	-0.0504173\\
3.031	-0.0505723\\
3.0336	-0.0507331\\
3.0361	-0.0508874\\
3.0387	-0.0510473\\
3.0413	-0.0512067\\
3.0438	-0.0513594\\
3.0464	-0.0515176\\
3.049	-0.0516751\\
3.0515	-0.0518259\\
3.0541	-0.051982\\
3.0567	-0.0521373\\
3.0592	-0.0522858\\
3.0618	-0.0524395\\
3.0643	-0.0525863\\
3.0669	-0.0527381\\
3.0695	-0.0528889\\
3.072	-0.0530329\\
3.0746	-0.0531817\\
3.0772	-0.0533293\\
3.0797	-0.0534702\\
3.0823	-0.0536156\\
3.0849	-0.0537598\\
3.0874	-0.0538973\\
3.09	-0.054039\\
3.0925	-0.054174\\
3.0951	-0.0543132\\
3.0977	-0.054451\\
3.1002	-0.0545822\\
3.1028	-0.0547172\\
3.1054	-0.0548509\\
3.1079	-0.054978\\
3.1105	-0.0551087\\
3.1131	-0.0552379\\
3.1156	-0.0553607\\
3.1182	-0.0554869\\
3.1207	-0.0556067\\
3.1233	-0.0557298\\
3.1259	-0.0558512\\
3.1284	-0.0559665\\
3.131	-0.0560847\\
3.1336	-0.0562012\\
3.1361	-0.0563117\\
3.1387	-0.056425\\
3.1413	-0.0565365\\
3.1438	-0.0566421\\
3.1464	-0.0567503\\
3.1489	-0.0568526\\
3.1515	-0.0569573\\
3.1541	-0.0570603\\
3.1566	-0.0571576\\
3.1592	-0.0572571\\
3.1618	-0.0573548\\
3.1643	-0.0574471\\
3.1669	-0.0575413\\
3.1695	-0.0576337\\
3.172	-0.0577209\\
3.1746	-0.0578098\\
3.1772	-0.0578969\\
3.1797	-0.057979\\
3.1823	-0.0580626\\
3.1848	-0.0581413\\
3.1874	-0.0582214\\
3.19	-0.0582997\\
3.1925	-0.0583734\\
3.1951	-0.0584482\\
3.1977	-0.0585213\\
3.2002	-0.05859\\
3.2028	-0.0586597\\
3.2054	-0.0587277\\
3.2079	-0.0587914\\
3.2105	-0.0588561\\
3.213	-0.0589166\\
3.2156	-0.058978\\
3.2182	-0.0590376\\
3.2207	-0.0590935\\
3.2233	-0.0591499\\
3.2259	-0.0592048\\
3.2284	-0.059256\\
3.231	-0.0593077\\
3.2336	-0.0593579\\
3.2361	-0.0594047\\
3.2387	-0.0594519\\
3.2412	-0.0594959\\
3.2438	-0.0595401\\
3.2464	-0.0595829\\
3.2489	-0.0596228\\
3.2515	-0.0596628\\
3.2541	-0.0597014\\
3.2566	-0.0597373\\
3.2592	-0.0597733\\
3.2618	-0.059808\\
3.2643	-0.0598402\\
3.2669	-0.0598724\\
3.2694	-0.0599022\\
3.272	-0.059932\\
3.2746	-0.0599606\\
3.2771	-0.0599871\\
3.2797	-0.0600135\\
3.2823	-0.0600388\\
3.2848	-0.0600622\\
3.2874	-0.0600855\\
3.29	-0.0601077\\
3.2925	-0.0601282\\
3.2951	-0.0601486\\
3.2976	-0.0601674\\
3.3002	-0.060186\\
3.3028	-0.0602038\\
3.3053	-0.0602201\\
3.3079	-0.0602363\\
3.3105	-0.0602517\\
3.313	-0.0602658\\
3.3156	-0.0602799\\
3.3182	-0.0602932\\
3.3207	-0.0603055\\
3.3233	-0.0603176\\
3.3259	-0.0603292\\
3.3284	-0.0603398\\
3.331	-0.0603503\\
3.3335	-0.06036\\
3.3361	-0.0603696\\
3.3387	-0.0603788\\
3.3412	-0.0603872\\
3.3438	-0.0603957\\
3.3464	-0.0604038\\
3.3489	-0.0604113\\
3.3515	-0.0604189\\
3.3541	-0.0604262\\
3.3566	-0.0604331\\
3.3592	-0.06044\\
3.3617	-0.0604465\\
3.3643	-0.0604532\\
3.3669	-0.0604598\\
3.3694	-0.060466\\
3.372	-0.0604725\\
3.3746	-0.0604789\\
3.3771	-0.0604851\\
3.3797	-0.0604916\\
3.3823	-0.0604981\\
3.3848	-0.0605045\\
3.3874	-0.0605112\\
3.3899	-0.0605179\\
3.3925	-0.0605249\\
3.3951	-0.0605321\\
3.3976	-0.0605393\\
3.4002	-0.0605469\\
3.4028	-0.0605549\\
3.4053	-0.0605627\\
3.4079	-0.0605712\\
3.4105	-0.06058\\
3.413	-0.0605888\\
3.4156	-0.0605982\\
3.4181	-0.0606077\\
3.4207	-0.0606179\\
3.4233	-0.0606286\\
3.4258	-0.0606392\\
3.4284	-0.0606507\\
3.431	-0.0606627\\
3.4335	-0.0606746\\
3.4361	-0.0606875\\
3.4387	-0.060701\\
3.4412	-0.0607144\\
3.4438	-0.0607288\\
3.4463	-0.0607432\\
3.4489	-0.0607588\\
3.4515	-0.0607749\\
3.454	-0.0607909\\
3.4566	-0.0608082\\
3.4592	-0.0608261\\
3.4617	-0.0608438\\
3.4643	-0.0608629\\
3.4669	-0.0608825\\
3.4694	-0.0609021\\
3.472	-0.060923\\
3.4745	-0.0609437\\
3.4771	-0.0609659\\
3.4797	-0.0609887\\
3.4822	-0.0610113\\
3.4848	-0.0610355\\
3.4874	-0.0610602\\
3.4899	-0.0610847\\
3.4925	-0.0611108\\
3.4951	-0.0611375\\
3.4976	-0.0611639\\
3.5002	-0.0611919\\
3.5028	-0.0612206\\
3.5053	-0.0612489\\
3.5079	-0.0612789\\
3.5104	-0.0613083\\
3.513	-0.0613396\\
3.5156	-0.0613716\\
3.5181	-0.0614029\\
3.5207	-0.061436\\
3.5233	-0.0614698\\
3.5258	-0.0615029\\
3.5284	-0.061538\\
3.531	-0.0615736\\
3.5335	-0.0616084\\
3.5361	-0.0616452\\
3.5386	-0.0616811\\
3.5412	-0.0617191\\
3.5438	-0.0617576\\
3.5463	-0.0617951\\
3.5489	-0.0618347\\
3.5515	-0.0618747\\
3.554	-0.0619138\\
3.5566	-0.0619549\\
3.5592	-0.0619964\\
3.5617	-0.0620368\\
3.5643	-0.0620793\\
3.5668	-0.0621206\\
3.5694	-0.062164\\
3.572	-0.0622078\\
3.5745	-0.0622502\\
3.5771	-0.0622948\\
3.5797	-0.0623397\\
3.5822	-0.0623832\\
3.5848	-0.0624288\\
3.5874	-0.0624747\\
3.5899	-0.0625191\\
3.5925	-0.0625656\\
3.595	-0.0626105\\
3.5976	-0.0626574\\
3.6002	-0.0627046\\
3.6027	-0.0627502\\
3.6053	-0.0627977\\
3.6079	-0.0628455\\
3.6104	-0.0628915\\
3.613	-0.0629395\\
3.6156	-0.0629876\\
3.6181	-0.0630339\\
3.6207	-0.0630821\\
3.6232	-0.0631285\\
3.6258	-0.0631768\\
3.6284	-0.063225\\
3.6309	-0.0632714\\
3.6335	-0.0633196\\
3.6361	-0.0633677\\
3.6386	-0.0634139\\
3.6412	-0.0634618\\
3.6438	-0.0635096\\
3.6463	-0.0635553\\
3.6489	-0.0636027\\
3.6515	-0.0636499\\
3.654	-0.0636951\\
3.6566	-0.0637418\\
3.6591	-0.0637864\\
3.6617	-0.0638326\\
3.6643	-0.0638784\\
3.6668	-0.0639221\\
3.6694	-0.0639671\\
3.672	-0.0640118\\
3.6745	-0.0640544\\
3.6771	-0.0640982\\
3.6797	-0.0641416\\
3.6822	-0.0641828\\
3.6848	-0.0642251\\
3.6873	-0.0642653\\
3.6899	-0.0643066\\
3.6925	-0.0643473\\
3.695	-0.0643858\\
3.6976	-0.0644252\\
3.7002	-0.064464\\
3.7027	-0.0645006\\
3.7053	-0.064538\\
3.7079	-0.0645747\\
3.7104	-0.0646093\\
3.713	-0.0646445\\
3.7155	-0.0646776\\
3.7181	-0.0647111\\
3.7207	-0.0647439\\
3.7232	-0.0647746\\
3.7258	-0.0648056\\
3.7284	-0.0648357\\
3.7309	-0.0648638\\
3.7335	-0.0648921\\
3.7361	-0.0649194\\
3.7386	-0.0649448\\
3.7412	-0.0649701\\
3.7437	-0.0649935\\
3.7463	-0.0650168\\
3.7489	-0.065039\\
3.7514	-0.0650594\\
3.754	-0.0650795\\
3.7566	-0.0650984\\
3.7591	-0.0651155\\
3.7617	-0.0651322\\
3.7643	-0.0651477\\
3.7668	-0.0651615\\
3.7694	-0.0651746\\
3.7719	-0.0651861\\
3.7745	-0.0651968\\
3.7771	-0.0652062\\
3.7796	-0.0652141\\
3.7822	-0.065221\\
3.7848	-0.0652266\\
3.7873	-0.0652308\\
3.7899	-0.0652338\\
3.7925	-0.0652354\\
3.795	-0.0652357\\
3.7976	-0.0652347\\
3.8002	-0.0652322\\
3.8027	-0.0652286\\
3.8053	-0.0652234\\
3.8078	-0.0652171\\
3.8104	-0.0652091\\
3.813	-0.0651997\\
3.8155	-0.0651893\\
3.8181	-0.065177\\
3.8207	-0.0651633\\
3.8232	-0.0651487\\
3.8258	-0.0651321\\
3.8284	-0.0651139\\
3.8309	-0.0650951\\
3.8335	-0.0650741\\
3.836	-0.0650524\\
3.8386	-0.0650284\\
3.8412	-0.0650029\\
3.8437	-0.0649769\\
3.8463	-0.0649484\\
3.8489	-0.0649184\\
3.8514	-0.0648882\\
3.854	-0.0648552\\
3.8566	-0.0648206\\
3.8591	-0.064786\\
3.8617	-0.0647485\\
3.8642	-0.064711\\
3.8668	-0.0646705\\
3.8694	-0.0646285\\
3.8719	-0.0645867\\
3.8745	-0.0645417\\
3.8771	-0.0644953\\
3.8796	-0.0644491\\
3.8822	-0.0643997\\
3.8848	-0.0643488\\
3.8873	-0.0642984\\
3.8899	-0.0642446\\
3.8924	-0.0641915\\
3.895	-0.0641348\\
3.8976	-0.0640766\\
3.9001	-0.0640193\\
3.9027	-0.0639583\\
3.9053	-0.0638959\\
3.9078	-0.0638345\\
3.9104	-0.0637693\\
3.913	-0.0637028\\
3.9155	-0.0636374\\
3.9181	-0.0635682\\
3.9206	-0.0635003\\
3.9232	-0.0634284\\
3.9258	-0.0633551\\
3.9283	-0.0632834\\
3.9309	-0.0632076\\
3.9335	-0.0631305\\
3.936	-0.0630552\\
3.9386	-0.0629757\\
3.9412	-0.0628949\\
3.9437	-0.0628161\\
3.9463	-0.062733\\
3.9488	-0.0626519\\
3.9514	-0.0625665\\
3.954	-0.06248\\
3.9565	-0.0623957\\
3.9591	-0.062307\\
3.9617	-0.0622173\\
3.9642	-0.0621299\\
3.9668	-0.0620381\\
3.9694	-0.0619453\\
3.9719	-0.0618552\\
3.9745	-0.0617604\\
3.9771	-0.0616648\\
3.9796	-0.061572\\
3.9822	-0.0614746\\
3.9847	-0.0613801\\
3.9873	-0.061281\\
3.9899	-0.0611812\\
3.9924	-0.0610844\\
3.995	-0.060983\\
3.9976	-0.0608809\\
4.0001	-0.060782\\
4.0027	-0.0606786\\
4.0053	-0.0605745\\
4.0078	-0.0604738\\
4.0104	-0.0603685\\
4.0129	-0.0602667\\
4.0155	-0.0601603\\
4.0181	-0.0600534\\
4.0206	-0.0599502\\
4.0232	-0.0598424\\
4.0258	-0.0597342\\
4.0283	-0.0596298\\
4.0309	-0.0595209\\
4.0335	-0.0594116\\
4.036	-0.0593062\\
4.0386	-0.0591964\\
4.0411	-0.0590906\\
4.0437	-0.0589803\\
4.0463	-0.0588699\\
4.0488	-0.0587635\\
4.0514	-0.0586528\\
4.054	-0.058542\\
4.0565	-0.0584354\\
4.0591	-0.0583246\\
4.0617	-0.0582137\\
4.0642	-0.0581072\\
4.0668	-0.0579964\\
4.0693	-0.05789\\
4.0719	-0.0577794\\
4.0745	-0.0576691\\
4.077	-0.0575631\\
4.0796	-0.0574531\\
4.0822	-0.0573434\\
4.0847	-0.0572382\\
4.0873	-0.057129\\
4.0899	-0.0570202\\
4.0924	-0.056916\\
4.095	-0.0568079\\
4.0975	-0.0567045\\
4.1001	-0.0565973\\
4.1027	-0.0564906\\
4.1052	-0.0563885\\
4.1078	-0.0562829\\
4.1104	-0.0561779\\
4.1129	-0.0560774\\
4.1155	-0.0559736\\
4.1181	-0.0558704\\
4.1206	-0.0557719\\
4.1232	-0.0556701\\
4.1258	-0.055569\\
4.1283	-0.0554726\\
4.1309	-0.0553731\\
4.1334	-0.0552782\\
4.136	-0.0551803\\
4.1386	-0.0550833\\
4.1411	-0.0549909\\
4.1437	-0.0548957\\
4.1463	-0.0548014\\
4.1488	-0.0547117\\
4.1514	-0.0546194\\
4.154	-0.0545281\\
4.1565	-0.0544413\\
4.1591	-0.054352\\
4.1616	-0.0542673\\
4.1642	-0.0541802\\
4.1668	-0.0540942\\
4.1693	-0.0540127\\
4.1719	-0.053929\\
4.1745	-0.0538465\\
4.177	-0.0537684\\
4.1796	-0.0536883\\
4.1822	-0.0536094\\
4.1847	-0.0535348\\
4.1873	-0.0534585\\
4.1898	-0.0533863\\
4.1924	-0.0533126\\
4.195	-0.0532401\\
4.1975	-0.0531717\\
4.2001	-0.053102\\
4.2027	-0.0530335\\
4.2052	-0.0529691\\
4.2078	-0.0529034\\
4.2104	-0.0528392\\
4.2129	-0.0527787\\
4.2155	-0.0527173\\
4.218	-0.0526596\\
4.2206	-0.052601\\
4.2232	-0.0525439\\
4.2257	-0.0524904\\
4.2283	-0.0524362\\
4.2309	-0.0523835\\
4.2334	-0.0523342\\
4.236	-0.0522845\\
4.2386	-0.0522363\\
4.2411	-0.0521913\\
4.2437	-0.0521461\\
4.2462	-0.0521041\\
4.2488	-0.0520618\\
4.2514	-0.0520212\\
4.2539	-0.0519835\\
4.2565	-0.0519458\\
4.2591	-0.0519097\\
4.2616	-0.0518765\\
4.2642	-0.0518434\\
4.2668	-0.0518119\\
4.2693	-0.051783\\
4.2719	-0.0517546\\
4.2744	-0.0517286\\
4.277	-0.0517032\\
4.2796	-0.0516793\\
4.2821	-0.0516577\\
4.2847	-0.0516368\\
4.2873	-0.0516174\\
4.2898	-0.0516002\\
4.2924	-0.0515838\\
4.295	-0.0515689\\
4.2975	-0.051556\\
4.3001	-0.0515441\\
4.3027	-0.0515336\\
4.3052	-0.051525\\
4.3078	-0.0515174\\
4.3103	-0.0515115\\
4.3129	-0.0515068\\
4.3155	-0.0515035\\
4.318	-0.0515017\\
4.3206	-0.0515012\\
4.3232	-0.0515021\\
4.3257	-0.0515043\\
4.3283	-0.051508\\
4.3309	-0.051513\\
4.3334	-0.0515191\\
4.336	-0.0515267\\
4.3385	-0.0515353\\
4.3411	-0.0515455\\
4.3437	-0.051557\\
4.3462	-0.0515693\\
4.3488	-0.0515833\\
4.3514	-0.0515985\\
4.3539	-0.0516143\\
4.3565	-0.0516319\\
4.3591	-0.0516507\\
4.3616	-0.0516698\\
4.3642	-0.0516909\\
4.3667	-0.0517122\\
4.3693	-0.0517355\\
4.3719	-0.0517598\\
4.3744	-0.0517842\\
4.377	-0.0518106\\
4.3796	-0.0518381\\
4.3821	-0.0518654\\
4.3847	-0.0518948\\
4.3873	-0.0519251\\
4.3898	-0.0519551\\
4.3924	-0.0519873\\
4.3949	-0.052019\\
4.3975	-0.0520529\\
4.4001	-0.0520876\\
4.4026	-0.0521218\\
4.4052	-0.0521581\\
4.4078	-0.0521951\\
4.4103	-0.0522315\\
4.4129	-0.05227\\
4.4155	-0.0523092\\
4.418	-0.0523475\\
4.4206	-0.052388\\
4.4231	-0.0524275\\
4.4257	-0.0524692\\
4.4283	-0.0525114\\
4.4308	-0.0525525\\
4.4334	-0.0525958\\
4.436	-0.0526396\\
4.4385	-0.0526821\\
4.4411	-0.0527268\\
4.4437	-0.0527718\\
4.4462	-0.0528155\\
4.4488	-0.0528613\\
4.4514	-0.0529075\\
4.4539	-0.0529521\\
4.4565	-0.0529988\\
4.459	-0.0530439\\
4.4616	-0.0530911\\
4.4642	-0.0531385\\
4.4667	-0.0531842\\
4.4693	-0.0532319\\
4.4719	-0.0532797\\
4.4744	-0.0533258\\
4.477	-0.0533737\\
4.4796	-0.0534218\\
4.4821	-0.0534679\\
4.4847	-0.053516\\
4.4872	-0.0535621\\
4.4898	-0.05361\\
4.4924	-0.0536579\\
4.4949	-0.0537038\\
4.4975	-0.0537514\\
4.5001	-0.0537988\\
4.5026	-0.0538443\\
4.5052	-0.0538914\\
4.5078	-0.0539382\\
4.5103	-0.053983\\
4.5129	-0.0540293\\
4.5154	-0.0540735\\
4.518	-0.0541192\\
4.5206	-0.0541645\\
4.5231	-0.0542077\\
4.5257	-0.0542523\\
4.5283	-0.0542964\\
4.5308	-0.0543385\\
4.5334	-0.0543818\\
4.536	-0.0544245\\
4.5385	-0.0544652\\
4.5411	-0.054507\\
4.5436	-0.0545466\\
4.5462	-0.0545873\\
4.5488	-0.0546274\\
4.5513	-0.0546654\\
4.5539	-0.0547043\\
4.5565	-0.0547426\\
4.559	-0.0547788\\
4.5616	-0.0548157\\
4.5642	-0.054852\\
4.5667	-0.0548862\\
4.5693	-0.0549211\\
4.5718	-0.0549539\\
4.5744	-0.0549873\\
4.577	-0.0550199\\
4.5795	-0.0550505\\
4.5821	-0.0550816\\
4.5847	-0.0551119\\
4.5872	-0.0551402\\
4.5898	-0.0551688\\
4.5924	-0.0551966\\
4.5949	-0.0552225\\
4.5975	-0.0552485\\
4.6001	-0.0552737\\
4.6026	-0.0552971\\
4.6052	-0.0553205\\
4.6077	-0.0553422\\
4.6103	-0.0553638\\
4.6129	-0.0553845\\
4.6154	-0.0554035\\
4.618	-0.0554224\\
4.6206	-0.0554403\\
4.6231	-0.0554566\\
4.6257	-0.0554726\\
4.6283	-0.0554877\\
4.6308	-0.0555012\\
4.6334	-0.0555144\\
4.6359	-0.0555261\\
4.6385	-0.0555373\\
4.6411	-0.0555475\\
4.6436	-0.0555565\\
4.6462	-0.0555648\\
4.6488	-0.0555721\\
4.6513	-0.0555782\\
4.6539	-0.0555836\\
4.6565	-0.055588\\
4.659	-0.0555913\\
4.6616	-0.0555938\\
4.6641	-0.0555952\\
4.6667	-0.0555958\\
4.6693	-0.0555953\\
4.6718	-0.055594\\
4.6744	-0.0555917\\
4.677	-0.0555884\\
4.6795	-0.0555844\\
4.6821	-0.0555792\\
4.6847	-0.0555731\\
4.6872	-0.0555664\\
4.6898	-0.0555585\\
4.6923	-0.05555\\
4.6949	-0.0555403\\
4.6975	-0.0555297\\
4.7	-0.0555187\\
4.7026	-0.0555063\\
4.7052	-0.0554931\\
4.7077	-0.0554796\\
4.7103	-0.0554647\\
4.7129	-0.0554489\\
4.7154	-0.055433\\
4.718	-0.0554157\\
4.7205	-0.0553982\\
4.7231	-0.0553793\\
4.7257	-0.0553597\\
4.7282	-0.05534\\
4.7308	-0.0553189\\
4.7334	-0.055297\\
4.7359	-0.0552753\\
4.7385	-0.055252\\
4.7411	-0.0552281\\
4.7436	-0.0552044\\
4.7462	-0.0551792\\
4.7487	-0.0551543\\
4.7513	-0.0551279\\
4.7539	-0.0551008\\
4.7564	-0.0550743\\
4.759	-0.0550461\\
4.7616	-0.0550174\\
4.7641	-0.0549893\\
4.7667	-0.0549595\\
4.7693	-0.0549293\\
4.7718	-0.0548998\\
4.7744	-0.0548687\\
4.777	-0.0548372\\
4.7795	-0.0548065\\
4.7821	-0.0547742\\
4.7846	-0.0547428\\
4.7872	-0.0547098\\
4.7898	-0.0546765\\
4.7923	-0.0546442\\
4.7949	-0.0546103\\
4.7975	-0.0545762\\
4.8	-0.0545432\\
4.8026	-0.0545086\\
4.8052	-0.0544738\\
4.8077	-0.0544403\\
4.8103	-0.0544052\\
4.8128	-0.0543714\\
4.8154	-0.0543361\\
4.818	-0.0543007\\
4.8205	-0.0542666\\
4.8231	-0.0542312\\
4.8257	-0.0541957\\
4.8282	-0.0541616\\
4.8308	-0.0541262\\
4.8334	-0.0540908\\
4.8359	-0.0540569\\
4.8385	-0.0540217\\
4.841	-0.053988\\
4.8436	-0.0539531\\
4.8462	-0.0539184\\
4.8487	-0.0538851\\
4.8513	-0.0538508\\
4.8539	-0.0538167\\
4.8564	-0.0537841\\
4.859	-0.0537506\\
4.8616	-0.0537173\\
4.8641	-0.0536856\\
4.8667	-0.053653\\
4.8692	-0.053622\\
4.8718	-0.0535901\\
4.8744	-0.0535587\\
4.8769	-0.0535288\\
4.8795	-0.0534983\\
4.8821	-0.0534682\\
4.8846	-0.0534397\\
4.8872	-0.0534106\\
4.8898	-0.053382\\
4.8923	-0.0533551\\
4.8949	-0.0533276\\
4.8974	-0.0533018\\
4.9	-0.0532756\\
4.9026	-0.0532499\\
4.9051	-0.0532259\\
4.9077	-0.0532016\\
4.9103	-0.053178\\
4.9128	-0.0531559\\
4.9154	-0.0531337\\
4.918	-0.0531123\\
4.9205	-0.0530924\\
4.9231	-0.0530724\\
4.9257	-0.0530533\\
4.9282	-0.0530356\\
4.9308	-0.0530181\\
4.9333	-0.053002\\
4.9359	-0.0529862\\
4.9385	-0.0529712\\
4.941	-0.0529576\\
4.9436	-0.0529444\\
4.9462	-0.0529321\\
4.9487	-0.0529212\\
4.9513	-0.0529107\\
4.9539	-0.0529012\\
4.9564	-0.052893\\
4.959	-0.0528855\\
4.9615	-0.0528791\\
4.9641	-0.0528735\\
4.9667	-0.0528689\\
4.9692	-0.0528655\\
4.9718	-0.0528629\\
4.9744	-0.0528614\\
4.9769	-0.052861\\
4.9795	-0.0528616\\
4.9821	-0.0528632\\
4.9846	-0.0528658\\
4.9872	-0.0528696\\
4.9897	-0.0528743\\
4.9923	-0.0528803\\
4.9949	-0.0528874\\
4.9974	-0.0528952\\
5	-0.0529045\\
};
\addplot [color=black, dotted, line width=0.6pt, forget plot]
  table[row sep=crcr]{%
-0.125	-0.0849146\\
5	-0.0849146\\
};
\draw[line width=\figlegendlw, fill=\figlegendfill, draw=\figlegenddraw] (axis cs:2.90186,-0.0234737) rectangle (axis cs:4.8975,0.094);
\addplot [color=black, line width=1.4pt, forget plot]
  table[row sep=crcr]{%
3.00164	0.0705053\\
3.64024	0.0705053\\
};
\node[right, align=left, inner sep=0]
at (axis cs:3.7,0.071) {$f_{\mathrm{COI}}$};
\addplot [color=mycolor5, line width=1.6pt, forget plot]
  table[row sep=crcr]{%
3.00164	0.0352632\\
3.05969	0.0352632\\
};
\addplot [color=mycolor10, line width=1.6pt, forget plot]
  table[row sep=crcr]{%
3.05969	0.0352632\\
3.11775	0.0352632\\
};
\addplot [color=mycolor6, line width=1.6pt, forget plot]
  table[row sep=crcr]{%
3.11775	0.0352632\\
3.1758	0.0352632\\
};
\addplot [color=mycolor7, line width=1.6pt, forget plot]
  table[row sep=crcr]{%
3.1758	0.0352632\\
3.23386	0.0352632\\
};
\addplot [color=mycolor3, line width=1.6pt, forget plot]
  table[row sep=crcr]{%
3.23386	0.0352632\\
3.29191	0.0352632\\
};
\addplot [color=mycolor2, line width=1.6pt, forget plot]
  table[row sep=crcr]{%
3.29191	0.0352632\\
3.34997	0.0352632\\
};
\addplot [color=mycolor11, line width=1.6pt, forget plot]
  table[row sep=crcr]{%
3.34997	0.0352632\\
3.40802	0.0352632\\
};
\addplot [color=mycolor12, line width=1.6pt, forget plot]
  table[row sep=crcr]{%
3.40802	0.0352632\\
3.46608	0.0352632\\
};
\addplot [color=mycolor8, line width=1.6pt, forget plot]
  table[row sep=crcr]{%
3.46608	0.0352632\\
3.52413	0.0352632\\
};
\addplot [color=mycolor9, line width=1.6pt, forget plot]
  table[row sep=crcr]{%
3.52413	0.0352632\\
3.58219	0.0352632\\
};
\addplot [color=mycolor4, line width=1.6pt, forget plot]
  table[row sep=crcr]{%
3.58219	0.0352632\\
3.64024	0.0352632\\
};
\node[right, align=left, inner sep=0]
at (axis cs:3.7,0.035) {$f_i$};
\addplot [color=black, dotted, line width=0.6pt, forget plot]
  table[row sep=crcr]{%
3.00164	2.10526e-05\\
3.64024	2.10526e-05\\
};
\node[right, align=left, inner sep=0]
at (axis cs:3.7,0) {$N_{\mathrm{COI}}^{0}$};
\end{axis}

\begin{axis}[%
width={\figscale*1.32in},
height={\figscale*0.95in},
at={(\figscale*1.985in,\figscale*0.31in)},
scale only axis,
separate axis lines,
every outer x axis line/.append style={black},
every x tick label/.append style={font=\figticfont},
every x tick/.append style={black},
xmin=-0.125,
xmax=5,
xlabel={Time since step in s},
every outer y axis line/.append style={black},
every y tick label/.append style={font=\figticfont},
every y tick/.append style={black},
ymin=-0.2,
ymax=0.1,
axis background/.style={fill=white},
xmajorgrids,
ymajorgrids,
grid style={white!55!black, opacity=0.3}
]

\addplot[area legend, draw=none, fill=mycolor1, fill opacity=0.9, forget plot]
table[row sep=crcr] {%
x	y\\
0	-0.0120484\\
0.0125313	-0.0359809\\
0.0250627	-0.0534805\\
0.037594	-0.0671721\\
0.0501253	-0.0672739\\
0.0626566	-0.0742048\\
0.075188	-0.0833975\\
0.0877193	-0.0923745\\
0.100251	-0.100346\\
0.112782	-0.106548\\
0.125313	-0.110299\\
0.137845	-0.110956\\
0.150376	-0.10956\\
0.162907	-0.110664\\
0.175439	-0.112484\\
0.18797	-0.114227\\
0.200501	-0.115904\\
0.213033	-0.117518\\
0.225564	-0.119094\\
0.238095	-0.120649\\
0.250627	-0.122167\\
0.263158	-0.123618\\
0.275689	-0.124989\\
0.288221	-0.126268\\
0.300752	-0.127421\\
0.313283	-0.128411\\
0.325815	-0.12922\\
0.338346	-0.129836\\
0.350877	-0.130239\\
0.363409	-0.130406\\
0.37594	-0.12955\\
0.388471	-0.127075\\
0.401003	-0.123252\\
0.413534	-0.118349\\
0.426065	-0.112642\\
0.438596	-0.106414\\
0.451128	-0.0999365\\
0.463659	-0.0934786\\
0.47619	-0.0873153\\
0.488722	-0.0817241\\
0.501253	-0.0769754\\
0.513784	-0.073336\\
0.526316	-0.0710781\\
0.538847	-0.0704766\\
0.551378	-0.0708478\\
0.56391	-0.0713161\\
0.576441	-0.0718693\\
0.588972	-0.0724981\\
0.601504	-0.0731895\\
0.614035	-0.0739289\\
0.626566	-0.074705\\
0.639098	-0.0755088\\
0.651629	-0.0763285\\
0.66416	-0.0771506\\
0.676692	-0.0779641\\
0.689223	-0.0787596\\
0.701754	-0.0795255\\
0.714286	-0.0802483\\
0.726817	-0.080678\\
0.739348	-0.0806202\\
0.75188	-0.0801453\\
0.764411	-0.0793219\\
0.776942	-0.0782201\\
0.789474	-0.0769108\\
0.802005	-0.0754637\\
0.814536	-0.0739472\\
0.827068	-0.072431\\
0.839599	-0.0709852\\
0.85213	-0.0696792\\
0.864662	-0.0685816\\
0.877193	-0.0677616\\
0.889724	-0.0672892\\
0.902256	-0.0672335\\
0.914787	-0.0676634\\
0.927318	-0.0686478\\
0.93985	-0.0704345\\
0.952381	-0.0731537\\
0.964912	-0.0766981\\
0.977444	-0.0809608\\
0.989975	-0.0858354\\
1.00251	-0.0912151\\
1.01504	-0.0969925\\
1.02757	-0.103061\\
1.0401	-0.109313\\
1.05263	-0.115641\\
1.06516	-0.121939\\
1.07769	-0.128099\\
1.09023	-0.134014\\
1.10276	-0.139577\\
1.11529	-0.144679\\
1.12782	-0.149214\\
1.14035	-0.153074\\
1.15288	-0.156151\\
1.16541	-0.158337\\
1.17794	-0.16002\\
1.19048	-0.161658\\
1.20301	-0.16325\\
1.21554	-0.164795\\
1.22807	-0.166293\\
1.2406	-0.167742\\
1.25313	-0.169142\\
1.26566	-0.170492\\
1.2782	-0.171791\\
1.29073	-0.173038\\
1.30326	-0.174231\\
1.31579	-0.17537\\
1.32832	-0.176453\\
1.34085	-0.177479\\
1.35338	-0.178447\\
1.36591	-0.179356\\
1.37845	-0.180205\\
1.39098	-0.180992\\
1.40351	-0.181717\\
1.41604	-0.182379\\
1.42857	-0.182976\\
1.4411	-0.183508\\
1.45363	-0.183974\\
1.46617	-0.183869\\
1.4787	-0.182847\\
1.49123	-0.18115\\
1.50376	-0.179019\\
1.51629	-0.176696\\
1.52882	-0.174422\\
1.54135	-0.172438\\
1.55388	-0.170987\\
1.56642	-0.170311\\
1.57895	-0.170651\\
1.59148	-0.171659\\
1.60401	-0.172816\\
1.61654	-0.174104\\
1.62907	-0.175509\\
1.6416	-0.177014\\
1.65414	-0.178604\\
1.66667	-0.180264\\
1.6792	-0.181978\\
1.69173	-0.183731\\
1.70426	-0.185508\\
1.71679	-0.187295\\
1.72932	-0.189076\\
1.74185	-0.190837\\
1.75439	-0.192562\\
1.76692	-0.194238\\
1.77945	-0.19585\\
1.79198	-0.197384\\
1.80451	-0.198824\\
1.81704	-0.200156\\
1.82957	-0.201367\\
1.84211	-0.202442\\
1.85464	-0.203366\\
1.86717	-0.204125\\
1.8797	-0.204705\\
1.89223	-0.205092\\
1.90476	-0.20527\\
1.91729	-0.205226\\
1.92982	-0.204878\\
1.94236	-0.20417\\
1.95489	-0.203129\\
1.96742	-0.201779\\
1.97995	-0.200145\\
1.99248	-0.198254\\
2.00501	-0.196131\\
2.01754	-0.1938\\
2.03008	-0.191287\\
2.04261	-0.188617\\
2.05514	-0.185815\\
2.06767	-0.182907\\
2.0802	-0.179917\\
2.09273	-0.17687\\
2.10526	-0.173791\\
2.11779	-0.170706\\
2.13033	-0.167638\\
2.14286	-0.164614\\
2.15539	-0.161657\\
2.16792	-0.158792\\
2.18045	-0.156045\\
2.19298	-0.153439\\
2.20551	-0.150999\\
2.21805	-0.14875\\
2.23058	-0.146717\\
2.24311	-0.144923\\
2.25564	-0.143394\\
2.26817	-0.142154\\
2.2807	-0.141227\\
2.29323	-0.140638\\
2.30576	-0.140268\\
2.3183	-0.139968\\
2.33083	-0.139718\\
2.34336	-0.139498\\
2.35589	-0.139289\\
2.36842	-0.139069\\
2.38095	-0.138819\\
2.39348	-0.13855\\
2.40602	-0.138288\\
2.41855	-0.138032\\
2.43108	-0.137782\\
2.44361	-0.137537\\
2.45614	-0.137297\\
2.46867	-0.13706\\
2.4812	-0.136828\\
2.49373	-0.136598\\
2.50627	-0.136371\\
2.5188	-0.136147\\
2.53133	-0.135924\\
2.54386	-0.135703\\
2.55639	-0.135482\\
2.56892	-0.135262\\
2.58145	-0.135041\\
2.59398	-0.134821\\
2.60652	-0.134599\\
2.61905	-0.134376\\
2.63158	-0.134151\\
2.64411	-0.133923\\
2.65664	-0.133693\\
2.66917	-0.13346\\
2.6817	-0.133224\\
2.69424	-0.132983\\
2.70677	-0.132739\\
2.7193	-0.132489\\
2.73183	-0.132234\\
2.74436	-0.131974\\
2.75689	-0.131621\\
2.76942	-0.131097\\
2.78195	-0.130414\\
2.79449	-0.129583\\
2.80702	-0.128618\\
2.81955	-0.12753\\
2.83208	-0.126332\\
2.84461	-0.125035\\
2.85714	-0.123651\\
2.86967	-0.122193\\
2.88221	-0.120673\\
2.89474	-0.119103\\
2.90727	-0.117494\\
2.9198	-0.11586\\
2.93233	-0.114211\\
2.94486	-0.11256\\
2.95739	-0.110919\\
2.96992	-0.1093\\
2.98246	-0.107714\\
2.99499	-0.106174\\
3.00752	-0.104692\\
3.02005	-0.103278\\
3.03258	-0.101946\\
3.04511	-0.100708\\
3.05764	-0.0995735\\
3.07018	-0.0985561\\
3.08271	-0.0976672\\
3.09524	-0.0969184\\
3.10777	-0.0963215\\
3.1203	-0.0958883\\
3.13283	-0.0955475\\
3.14536	-0.09522\\
3.15789	-0.0949063\\
3.17043	-0.0946065\\
3.18296	-0.0943209\\
3.19549	-0.0940497\\
3.20802	-0.0937933\\
3.22055	-0.0935517\\
3.23308	-0.0933252\\
3.24561	-0.093114\\
3.25815	-0.0929183\\
3.27068	-0.0927383\\
3.28321	-0.0925741\\
3.29574	-0.092426\\
3.30827	-0.0922941\\
3.3208	-0.0921785\\
3.33333	-0.0920795\\
3.34586	-0.0919972\\
3.3584	-0.0919318\\
3.37093	-0.0918834\\
3.38346	-0.0918523\\
3.39599	-0.0918385\\
3.40852	-0.0918423\\
3.42105	-0.0918639\\
3.43358	-0.0919034\\
3.44612	-0.0919611\\
3.45865	-0.0920371\\
3.47118	-0.0921317\\
3.48371	-0.092245\\
3.49624	-0.0923974\\
3.50877	-0.0926059\\
3.5213	-0.0928666\\
3.53383	-0.0931752\\
3.54637	-0.0935277\\
3.5589	-0.0939198\\
3.57143	-0.0943475\\
3.58396	-0.0948068\\
3.59649	-0.0952934\\
3.60902	-0.0958033\\
3.62155	-0.0963325\\
3.63409	-0.096877\\
3.64662	-0.0974326\\
3.65915	-0.0979955\\
3.67168	-0.0985615\\
3.68421	-0.0991267\\
3.69674	-0.0996871\\
3.70927	-0.100239\\
3.7218	-0.100778\\
3.73434	-0.1013\\
3.74687	-0.101802\\
3.7594	-0.102279\\
3.77193	-0.102728\\
3.78446	-0.103145\\
3.79699	-0.103526\\
3.80952	-0.103866\\
3.82206	-0.104163\\
3.83459	-0.104413\\
3.84712	-0.10461\\
3.85965	-0.104753\\
3.87218	-0.104837\\
3.88471	-0.104857\\
3.89724	-0.10482\\
3.90977	-0.104734\\
3.92231	-0.104602\\
3.93484	-0.104427\\
3.94737	-0.104209\\
3.9599	-0.103953\\
3.97243	-0.10366\\
3.98496	-0.103333\\
3.99749	-0.102975\\
4.01003	-0.102587\\
4.02256	-0.102172\\
4.03509	-0.101734\\
4.04762	-0.101273\\
4.06015	-0.100793\\
4.07268	-0.100296\\
4.08521	-0.0997842\\
4.09774	-0.0992608\\
4.11028	-0.0987279\\
4.12281	-0.098188\\
4.13534	-0.0976436\\
4.14787	-0.0970971\\
4.1604	-0.0965511\\
4.17293	-0.0960081\\
4.18546	-0.0954704\\
4.19799	-0.0949406\\
4.21053	-0.0944213\\
4.22306	-0.0939147\\
4.23559	-0.0934235\\
4.24812	-0.0929501\\
4.26065	-0.0924969\\
4.27318	-0.0920665\\
4.28571	-0.0916613\\
4.29825	-0.0912838\\
4.31078	-0.0909363\\
4.32331	-0.0906107\\
4.33584	-0.0902965\\
4.34837	-0.089993\\
4.3609	-0.0896994\\
4.37343	-0.089415\\
4.38596	-0.089139\\
4.3985	-0.0888707\\
4.41103	-0.0886093\\
4.42356	-0.0883541\\
4.43609	-0.0881042\\
4.44862	-0.0878589\\
4.46115	-0.0876174\\
4.47368	-0.0873789\\
4.48622	-0.0871426\\
4.49875	-0.0869077\\
4.51128	-0.0866733\\
4.52381	-0.0864388\\
4.53634	-0.0862032\\
4.54887	-0.0859657\\
4.5614	-0.0857255\\
4.57393	-0.0854817\\
4.58647	-0.0852335\\
4.599	-0.08498\\
4.61153	-0.0847204\\
4.62406	-0.0844537\\
4.63659	-0.084182\\
4.64912	-0.083908\\
4.66165	-0.0836317\\
4.67419	-0.0833532\\
4.68672	-0.0830724\\
4.69925	-0.0827894\\
4.71178	-0.0825043\\
4.72431	-0.0822171\\
4.73684	-0.0819278\\
4.74937	-0.0816364\\
4.7619	-0.081343\\
4.77444	-0.0810475\\
4.78697	-0.0807499\\
4.7995	-0.0804503\\
4.81203	-0.0801486\\
4.82456	-0.0798448\\
4.83709	-0.079539\\
4.84962	-0.0792311\\
4.86216	-0.078921\\
4.87469	-0.0786088\\
4.88722	-0.0782945\\
4.89975	-0.0779779\\
4.91228	-0.0776592\\
4.92481	-0.0773381\\
4.93734	-0.0770148\\
4.94987	-0.0766892\\
4.96241	-0.0763612\\
4.97494	-0.0760308\\
4.98747	-0.075698\\
5	-0.0753627\\
5	-0.0585992\\
4.98747	-0.0586296\\
4.97494	-0.0586596\\
4.96241	-0.0586886\\
4.94987	-0.0587162\\
4.93734	-0.058742\\
4.92481	-0.0587654\\
4.91228	-0.058786\\
4.89975	-0.0588033\\
4.88722	-0.0588168\\
4.87469	-0.0588261\\
4.86216	-0.0588307\\
4.84962	-0.05883\\
4.83709	-0.0588237\\
4.82456	-0.0588126\\
4.81203	-0.0587981\\
4.7995	-0.058781\\
4.78697	-0.0587617\\
4.77444	-0.0587409\\
4.7619	-0.0587193\\
4.74937	-0.0586972\\
4.73684	-0.0586754\\
4.72431	-0.0586543\\
4.71178	-0.0586346\\
4.69925	-0.0586168\\
4.68672	-0.0586015\\
4.67419	-0.0585892\\
4.66165	-0.0585804\\
4.64912	-0.0585757\\
4.63659	-0.0585756\\
4.62406	-0.0585807\\
4.61153	-0.0585915\\
4.599	-0.0586084\\
4.58647	-0.058632\\
4.57393	-0.0586628\\
4.5614	-0.0587013\\
4.54887	-0.058748\\
4.53634	-0.0588027\\
4.52381	-0.0588647\\
4.51128	-0.058934\\
4.49875	-0.0590103\\
4.48622	-0.0590934\\
4.47368	-0.059183\\
4.46115	-0.059279\\
4.44862	-0.059381\\
4.43609	-0.059489\\
4.42356	-0.0596025\\
4.41103	-0.0597214\\
4.3985	-0.0598455\\
4.38596	-0.0599744\\
4.37343	-0.0601079\\
4.3609	-0.0602457\\
4.34837	-0.0603875\\
4.33584	-0.0605331\\
4.32331	-0.0606821\\
4.31078	-0.0608343\\
4.29825	-0.0609894\\
4.28571	-0.061147\\
4.27318	-0.0613068\\
4.26065	-0.0614686\\
4.24812	-0.0616321\\
4.23559	-0.0617968\\
4.22306	-0.0619625\\
4.21053	-0.0621289\\
4.19799	-0.0622956\\
4.18546	-0.0624623\\
4.17293	-0.0626286\\
4.1604	-0.0627943\\
4.14787	-0.0629591\\
4.13534	-0.0631224\\
4.12281	-0.0632841\\
4.11028	-0.0634438\\
4.09774	-0.0636012\\
4.08521	-0.0637559\\
4.07268	-0.0639076\\
4.06015	-0.064056\\
4.04762	-0.0642007\\
4.03509	-0.0643415\\
4.02256	-0.064478\\
4.01003	-0.0646098\\
3.99749	-0.0647367\\
3.98496	-0.0648584\\
3.97243	-0.0649746\\
3.9599	-0.0650849\\
3.94737	-0.0651891\\
3.93484	-0.0652869\\
3.92231	-0.0653781\\
3.90977	-0.0654622\\
3.89724	-0.0655402\\
3.88471	-0.0656126\\
3.87218	-0.0656796\\
3.85965	-0.065741\\
3.84712	-0.0657969\\
3.83459	-0.0658471\\
3.82206	-0.0658918\\
3.80952	-0.0659308\\
3.79699	-0.0659641\\
3.78446	-0.0659919\\
3.77193	-0.0660141\\
3.7594	-0.0660308\\
3.74687	-0.066042\\
3.73434	-0.0660478\\
3.7218	-0.0660483\\
3.70927	-0.0660435\\
3.69674	-0.0660337\\
3.68421	-0.0660188\\
3.67168	-0.0659991\\
3.65915	-0.0659747\\
3.64662	-0.0659457\\
3.63409	-0.0659124\\
3.62155	-0.0658748\\
3.60902	-0.0658334\\
3.59649	-0.0657881\\
3.58396	-0.0657394\\
3.57143	-0.0656874\\
3.5589	-0.0656324\\
3.54637	-0.0655746\\
3.53383	-0.0655144\\
3.5213	-0.0654521\\
3.50877	-0.0653879\\
3.49624	-0.0653221\\
3.48371	-0.0652552\\
3.47118	-0.0651874\\
3.45865	-0.0651192\\
3.44612	-0.0650508\\
3.43358	-0.0649826\\
3.42105	-0.064915\\
3.40852	-0.0648484\\
3.39599	-0.0647833\\
3.38346	-0.0647199\\
3.37093	-0.0646586\\
3.3584	-0.0646\\
3.34586	-0.0645444\\
3.33333	-0.0644922\\
3.3208	-0.0644438\\
3.30827	-0.0643996\\
3.29574	-0.0643601\\
3.28321	-0.0643257\\
3.27068	-0.0642967\\
3.25815	-0.0642736\\
3.24561	-0.0642568\\
3.23308	-0.0642463\\
3.22055	-0.0642424\\
3.20802	-0.0642453\\
3.19549	-0.0642556\\
3.18296	-0.0642735\\
3.17043	-0.0642995\\
3.15789	-0.064334\\
3.14536	-0.0643773\\
3.13283	-0.0644297\\
3.1203	-0.0644916\\
3.10777	-0.0645633\\
3.09524	-0.0646452\\
3.08271	-0.0647374\\
3.07018	-0.0648403\\
3.05764	-0.0649541\\
3.04511	-0.0650791\\
3.03258	-0.0652154\\
3.02005	-0.0653632\\
3.00752	-0.0655228\\
2.99499	-0.0656943\\
2.98246	-0.0658778\\
2.96992	-0.0660734\\
2.95739	-0.0662813\\
2.94486	-0.0665015\\
2.93233	-0.066734\\
2.9198	-0.066979\\
2.90727	-0.0672364\\
2.89474	-0.0675062\\
2.88221	-0.0677884\\
2.86967	-0.0680831\\
2.85714	-0.06839\\
2.84461	-0.0687093\\
2.83208	-0.0690434\\
2.81955	-0.0693945\\
2.80702	-0.0697624\\
2.79449	-0.0701465\\
2.78195	-0.0705465\\
2.76942	-0.0709619\\
2.75689	-0.0713921\\
2.74436	-0.0718368\\
2.73183	-0.0722956\\
2.7193	-0.0727678\\
2.70677	-0.0732531\\
2.69424	-0.073751\\
2.6817	-0.0742609\\
2.66917	-0.0747825\\
2.65664	-0.0753152\\
2.64411	-0.0758587\\
2.63158	-0.0764123\\
2.61905	-0.0769756\\
2.60652	-0.0775483\\
2.59398	-0.0781298\\
2.58145	-0.0787196\\
2.56892	-0.0793174\\
2.55639	-0.0799227\\
2.54386	-0.080535\\
2.53133	-0.081154\\
2.5188	-0.0817792\\
2.50627	-0.0824101\\
2.49373	-0.0830465\\
2.4812	-0.0836878\\
2.46867	-0.0843336\\
2.45614	-0.0849836\\
2.44361	-0.0856374\\
2.43108	-0.0862945\\
2.41855	-0.0869545\\
2.40602	-0.087617\\
2.39348	-0.0882817\\
2.38095	-0.088948\\
2.36842	-0.0896156\\
2.35589	-0.0902841\\
2.34336	-0.0909529\\
2.33083	-0.0916216\\
2.3183	-0.0922899\\
2.30576	-0.092957\\
2.29323	-0.0936227\\
2.2807	-0.0942863\\
2.26817	-0.0949474\\
2.25564	-0.0956053\\
2.24311	-0.0962594\\
2.23058	-0.0969092\\
2.21805	-0.0975541\\
2.20551	-0.0981932\\
2.19298	-0.098826\\
2.18045	-0.0994517\\
2.16792	-0.10007\\
2.15539	-0.100679\\
2.14286	-0.101278\\
2.13033	-0.101868\\
2.11779	-0.102449\\
2.10526	-0.103019\\
2.09273	-0.103578\\
2.0802	-0.104127\\
2.06767	-0.104663\\
2.05514	-0.105188\\
2.04261	-0.1057\\
2.03008	-0.106198\\
2.01754	-0.106681\\
2.00501	-0.10715\\
1.99248	-0.107602\\
1.97995	-0.108038\\
1.96742	-0.108456\\
1.95489	-0.108856\\
1.94236	-0.109236\\
1.92982	-0.109595\\
1.91729	-0.109932\\
1.90476	-0.110247\\
1.89223	-0.110538\\
1.8797	-0.110803\\
1.86717	-0.111043\\
1.85464	-0.111255\\
1.84211	-0.111438\\
1.82957	-0.111592\\
1.81704	-0.111714\\
1.80451	-0.111805\\
1.79198	-0.111862\\
1.77945	-0.111885\\
1.76692	-0.111873\\
1.75439	-0.111824\\
1.74185	-0.111738\\
1.72932	-0.111613\\
1.71679	-0.111449\\
1.70426	-0.111245\\
1.69173	-0.111\\
1.6792	-0.110714\\
1.66667	-0.110395\\
1.65414	-0.110048\\
1.6416	-0.109672\\
1.62907	-0.109265\\
1.61654	-0.108825\\
1.60401	-0.108352\\
1.59148	-0.107842\\
1.57895	-0.107297\\
1.56642	-0.106713\\
1.55388	-0.106091\\
1.54135	-0.10543\\
1.52882	-0.104728\\
1.51629	-0.103986\\
1.50376	-0.103202\\
1.49123	-0.102378\\
1.4787	-0.101513\\
1.46617	-0.100606\\
1.45363	-0.0996587\\
1.4411	-0.0986707\\
1.42857	-0.0976424\\
1.41604	-0.0965745\\
1.40351	-0.0954676\\
1.39098	-0.0943222\\
1.37845	-0.0931392\\
1.36591	-0.0919194\\
1.35338	-0.0906636\\
1.34085	-0.0893726\\
1.32832	-0.0880476\\
1.31579	-0.0866893\\
1.30326	-0.0852987\\
1.29073	-0.0838766\\
1.2782	-0.0824242\\
1.26566	-0.0809421\\
1.25313	-0.0794312\\
1.2406	-0.0778923\\
1.22807	-0.0763263\\
1.21554	-0.0747337\\
1.20301	-0.0731151\\
1.19048	-0.071471\\
1.17794	-0.0698022\\
1.16541	-0.0681089\\
1.15288	-0.0663752\\
1.14035	-0.0645923\\
1.12782	-0.0627712\\
1.11529	-0.0609227\\
1.10276	-0.0590575\\
1.09023	-0.0571862\\
1.07769	-0.0553198\\
1.06516	-0.0534689\\
1.05263	-0.0516437\\
1.0401	-0.0498548\\
1.02757	-0.0481129\\
1.01504	-0.0464283\\
1.00251	-0.044811\\
0.989975	-0.0432713\\
0.977444	-0.0418197\\
0.964912	-0.040516\\
0.952381	-0.0394072\\
0.93985	-0.038485\\
0.927318	-0.0377415\\
0.914787	-0.0371683\\
0.902256	-0.0367564\\
0.889724	-0.0364973\\
0.877193	-0.0363832\\
0.864662	-0.0364056\\
0.85213	-0.0365551\\
0.839599	-0.0368232\\
0.827068	-0.0372023\\
0.814536	-0.0376835\\
0.802005	-0.0382574\\
0.789474	-0.0389152\\
0.776942	-0.039649\\
0.764411	-0.0404499\\
0.75188	-0.0413078\\
0.739348	-0.0422133\\
0.726817	-0.0431584\\
0.714286	-0.0441337\\
0.701754	-0.0451278\\
0.689223	-0.0461308\\
0.676692	-0.0471343\\
0.66416	-0.0481284\\
0.651629	-0.0491008\\
0.639098	-0.0500408\\
0.626566	-0.0509405\\
0.614035	-0.0517898\\
0.601504	-0.0525762\\
0.588972	-0.0532891\\
0.576441	-0.0539218\\
0.56391	-0.0544655\\
0.551378	-0.054908\\
0.538847	-0.0552615\\
0.526316	-0.0555602\\
0.513784	-0.0558219\\
0.501253	-0.0560596\\
0.488722	-0.0562888\\
0.47619	-0.0565307\\
0.463659	-0.0568026\\
0.451128	-0.0571149\\
0.438596	-0.0574801\\
0.426065	-0.0579177\\
0.413534	-0.0584433\\
0.401003	-0.0590629\\
0.388471	-0.0597866\\
0.37594	-0.0607964\\
0.363409	-0.0622277\\
0.350877	-0.0640221\\
0.338346	-0.0661284\\
0.325815	-0.0685119\\
0.313283	-0.0711375\\
0.300752	-0.0739583\\
0.288221	-0.0769352\\
0.275689	-0.0800487\\
0.263158	-0.0832753\\
0.250627	-0.08657\\
0.238095	-0.0898886\\
0.225564	-0.0932036\\
0.213033	-0.096475\\
0.200501	-0.0996272\\
0.18797	-0.102579\\
0.175439	-0.105271\\
0.162907	-0.107633\\
0.150376	-0.10956\\
0.137845	-0.110956\\
0.125313	-0.110299\\
0.112782	-0.106548\\
0.100251	-0.100346\\
0.0877193	-0.0923745\\
0.075188	-0.0833975\\
0.0626566	-0.0742048\\
0.0501253	-0.0655363\\
0.037594	-0.0669288\\
0.0250627	-0.053248\\
0.0125313	-0.0341609\\
0	-0.0118443\\
}--cycle;

\addplot[area legend, draw=none, fill=mycolor1, fill opacity=0.9, forget plot]
table[row sep=crcr] {%
x	y\\
0	0.0118445\\
0.0125313	0.0403347\\
0.0250627	0.0674639\\
0.037594	0.0875713\\
0.0501253	0.0911556\\
0.0626566	0.103777\\
0.075188	0.116117\\
0.0877193	0.127632\\
0.100251	0.137831\\
0.112782	0.146196\\
0.125313	0.152125\\
0.137845	0.15498\\
0.150376	0.155847\\
0.162907	0.156255\\
0.175439	0.15621\\
0.18797	0.155705\\
0.200501	0.154772\\
0.213033	0.153452\\
0.225564	0.151763\\
0.238095	0.149732\\
0.250627	0.147419\\
0.263158	0.144896\\
0.275689	0.142221\\
0.288221	0.139448\\
0.300752	0.136654\\
0.313283	0.133922\\
0.325815	0.131312\\
0.338346	0.128878\\
0.350877	0.126686\\
0.363409	0.124801\\
0.37594	0.123272\\
0.388471	0.122143\\
0.401003	0.12128\\
0.413534	0.120511\\
0.426065	0.119814\\
0.438596	0.11916\\
0.451128	0.118534\\
0.463659	0.117923\\
0.47619	0.117305\\
0.488722	0.11666\\
0.501253	0.115971\\
0.513784	0.115228\\
0.526316	0.114414\\
0.538847	0.113508\\
0.551378	0.112496\\
0.56391	0.111352\\
0.576441	0.110063\\
0.588972	0.108635\\
0.601504	0.107077\\
0.614035	0.1054\\
0.626566	0.103612\\
0.639098	0.101718\\
0.651629	0.0997264\\
0.66416	0.0976475\\
0.676692	0.0954885\\
0.689223	0.0932552\\
0.701754	0.0909557\\
0.714286	0.0885999\\
0.726817	0.0861957\\
0.739348	0.0837499\\
0.75188	0.0812713\\
0.764411	0.0787698\\
0.776942	0.0762543\\
0.789474	0.0737325\\
0.802005	0.0712136\\
0.814536	0.068708\\
0.827068	0.0662249\\
0.839599	0.0637729\\
0.85213	0.0613613\\
0.864662	0.0590005\\
0.877193	0.0566999\\
0.889724	0.0544685\\
0.902256	0.0523158\\
0.914787	0.050252\\
0.927318	0.0482867\\
0.93985	0.046429\\
0.952381	0.0446885\\
0.964912	0.0430755\\
0.977444	0.0415997\\
0.989975	0.0402147\\
1.00251	0.0388625\\
1.01504	0.0375349\\
1.02757	0.0362237\\
1.0401	0.0349202\\
1.05263	0.0336162\\
1.06516	0.0323038\\
1.07769	0.030975\\
1.09023	0.0296215\\
1.10276	0.0282352\\
1.11529	0.0268084\\
1.12782	0.0253333\\
1.14035	0.0238019\\
1.15288	0.0222064\\
1.16541	0.0205393\\
1.17794	0.0188128\\
1.19048	0.0170461\\
1.20301	0.0152422\\
1.21554	0.0134046\\
1.22807	0.0115365\\
1.2406	0.0096413\\
1.25313	0.00772233\\
1.26566	0.00578319\\
1.2782	0.00382742\\
1.29073	0.00185853\\
1.30326	-0.000119867\\
1.31579	-0.00210407\\
1.32832	-0.00409038\\
1.34085	-0.00607518\\
1.35338	-0.00805482\\
1.36591	-0.0100256\\
1.37845	-0.0119838\\
1.39098	-0.013926\\
1.40351	-0.0158485\\
1.41604	-0.0177479\\
1.42857	-0.0196206\\
1.4411	-0.0214635\\
1.45363	-0.0232732\\
1.46617	-0.0250465\\
1.4787	-0.0267804\\
1.49123	-0.028472\\
1.50376	-0.0301185\\
1.51629	-0.0317171\\
1.52882	-0.0332653\\
1.54135	-0.0347608\\
1.55388	-0.0362012\\
1.56642	-0.0375844\\
1.57895	-0.0389084\\
1.59148	-0.0401715\\
1.60401	-0.041372\\
1.61654	-0.0425083\\
1.62907	-0.0435792\\
1.6416	-0.0445832\\
1.65414	-0.0455195\\
1.66667	-0.0463869\\
1.6792	-0.0471847\\
1.69173	-0.0479208\\
1.70426	-0.0486042\\
1.71679	-0.0492362\\
1.72932	-0.0498179\\
1.74185	-0.0503507\\
1.75439	-0.0508361\\
1.76692	-0.0512756\\
1.77945	-0.0516706\\
1.79198	-0.0520228\\
1.80451	-0.052334\\
1.81704	-0.0526056\\
1.82957	-0.0528395\\
1.84211	-0.0530375\\
1.85464	-0.0532011\\
1.86717	-0.0533323\\
1.8797	-0.0534326\\
1.89223	-0.053504\\
1.90476	-0.053548\\
1.91729	-0.0535665\\
1.92982	-0.0535611\\
1.94236	-0.0535334\\
1.95489	-0.0534851\\
1.96742	-0.0534178\\
1.97995	-0.053333\\
1.99248	-0.0532323\\
2.00501	-0.0531172\\
2.01754	-0.052989\\
2.03008	-0.0528492\\
2.04261	-0.0526991\\
2.05514	-0.0525401\\
2.06767	-0.0523733\\
2.0802	-0.0522\\
2.09273	-0.0520214\\
2.10526	-0.0518387\\
2.11779	-0.0516528\\
2.13033	-0.0514648\\
2.14286	-0.0512757\\
2.15539	-0.0510869\\
2.16792	-0.0508995\\
2.18045	-0.0507138\\
2.19298	-0.05053\\
2.20551	-0.0503485\\
2.21805	-0.0501695\\
2.23058	-0.0499931\\
2.24311	-0.0498197\\
2.25564	-0.0496492\\
2.26817	-0.0494819\\
2.2807	-0.0493177\\
2.29323	-0.0491569\\
2.30576	-0.0489994\\
2.3183	-0.0488452\\
2.33083	-0.0486944\\
2.34336	-0.048547\\
2.35589	-0.0484028\\
2.36842	-0.048262\\
2.38095	-0.0481243\\
2.39348	-0.0479898\\
2.40602	-0.0478585\\
2.41855	-0.0477301\\
2.43108	-0.0476046\\
2.44361	-0.047482\\
2.45614	-0.0473622\\
2.46867	-0.047245\\
2.4812	-0.0471304\\
2.49373	-0.0470182\\
2.50627	-0.0469085\\
2.5188	-0.046801\\
2.53133	-0.0466958\\
2.54386	-0.0465927\\
2.55639	-0.0464917\\
2.56892	-0.0463927\\
2.58145	-0.0462957\\
2.59398	-0.0462005\\
2.60652	-0.0461072\\
2.61905	-0.0460156\\
2.63158	-0.0459259\\
2.64411	-0.0458378\\
2.65664	-0.0457514\\
2.66917	-0.0456666\\
2.6817	-0.0455835\\
2.69424	-0.045502\\
2.70677	-0.0454221\\
2.7193	-0.0453437\\
2.73183	-0.0452669\\
2.74436	-0.0451915\\
2.75689	-0.0451177\\
2.76942	-0.0450454\\
2.78195	-0.0449745\\
2.79449	-0.0449049\\
2.80702	-0.0448367\\
2.81955	-0.0447698\\
2.83208	-0.0447041\\
2.84461	-0.0446396\\
2.85714	-0.0445788\\
2.86967	-0.0445245\\
2.88221	-0.0444771\\
2.89474	-0.0444366\\
2.90727	-0.0444033\\
2.9198	-0.0443773\\
2.93233	-0.0443587\\
2.94486	-0.0443476\\
2.95739	-0.0443442\\
2.96992	-0.0443485\\
2.98246	-0.0443606\\
2.99499	-0.0443806\\
3.00752	-0.0444084\\
3.02005	-0.044444\\
3.03258	-0.0444875\\
3.04511	-0.0445388\\
3.05764	-0.0445979\\
3.07018	-0.0446646\\
3.08271	-0.044739\\
3.09524	-0.0448208\\
3.10777	-0.0449099\\
3.1203	-0.0450063\\
3.13283	-0.0451096\\
3.14536	-0.0452199\\
3.15789	-0.0453368\\
3.17043	-0.0454602\\
3.18296	-0.0455898\\
3.19549	-0.0457254\\
3.20802	-0.0458668\\
3.22055	-0.0460137\\
3.23308	-0.0461658\\
3.24561	-0.0463229\\
3.25815	-0.0464844\\
3.27068	-0.0466497\\
3.28321	-0.0468184\\
3.29574	-0.0469902\\
3.30827	-0.0471648\\
3.3208	-0.0473418\\
3.33333	-0.047521\\
3.34586	-0.0477019\\
3.3584	-0.0478843\\
3.37093	-0.0480678\\
3.38346	-0.0482521\\
3.39599	-0.0484368\\
3.40852	-0.0486216\\
3.42105	-0.0488063\\
3.43358	-0.0489905\\
3.44612	-0.0491738\\
3.45865	-0.0493561\\
3.47118	-0.0495369\\
3.48371	-0.0497161\\
3.49624	-0.0498933\\
3.50877	-0.0500682\\
3.5213	-0.0502408\\
3.53383	-0.0504105\\
3.54637	-0.0505774\\
3.5589	-0.0507411\\
3.57143	-0.0509014\\
3.58396	-0.0510581\\
3.59649	-0.0512111\\
3.60902	-0.0513602\\
3.62155	-0.0515052\\
3.63409	-0.051646\\
3.64662	-0.0517824\\
3.65915	-0.0519144\\
3.67168	-0.0520418\\
3.68421	-0.0521646\\
3.69674	-0.0522827\\
3.70927	-0.0523961\\
3.7218	-0.0525046\\
3.73434	-0.0526082\\
3.74687	-0.052707\\
3.7594	-0.052801\\
3.77193	-0.05289\\
3.78446	-0.0529743\\
3.79699	-0.0530537\\
3.80952	-0.0531284\\
3.82206	-0.0531984\\
3.83459	-0.0532638\\
3.84712	-0.0533246\\
3.85965	-0.053381\\
3.87218	-0.0534331\\
3.88471	-0.053481\\
3.89724	-0.0535247\\
3.90977	-0.0535645\\
3.92231	-0.0536013\\
3.93484	-0.0536361\\
3.94737	-0.0536689\\
3.9599	-0.0536996\\
3.97243	-0.0537282\\
3.98496	-0.0537547\\
3.99749	-0.0537792\\
4.01003	-0.0538015\\
4.02256	-0.0538218\\
4.03509	-0.0538401\\
4.04762	-0.0538563\\
4.06015	-0.0538704\\
4.07268	-0.0538827\\
4.08521	-0.0538929\\
4.09774	-0.0539013\\
4.11028	-0.0539077\\
4.12281	-0.0539123\\
4.13534	-0.0539151\\
4.14787	-0.0539162\\
4.1604	-0.0539155\\
4.17293	-0.0539131\\
4.18546	-0.053909\\
4.19799	-0.0539034\\
4.21053	-0.0538961\\
4.22306	-0.0538873\\
4.23559	-0.053877\\
4.24812	-0.0538652\\
4.26065	-0.0538519\\
4.27318	-0.0538372\\
4.28571	-0.0538211\\
4.29825	-0.0538035\\
4.31078	-0.0537846\\
4.32331	-0.0537643\\
4.33584	-0.0537426\\
4.34837	-0.0537196\\
4.3609	-0.0536951\\
4.37343	-0.0536693\\
4.38596	-0.053642\\
4.3985	-0.0536133\\
4.41103	-0.0535831\\
4.42356	-0.0535515\\
4.43609	-0.0535183\\
4.44862	-0.0534835\\
4.46115	-0.0534471\\
4.47368	-0.053409\\
4.48622	-0.0533692\\
4.49875	-0.0533275\\
4.51128	-0.0532839\\
4.52381	-0.0532384\\
4.53634	-0.0531908\\
4.54887	-0.053141\\
4.5614	-0.0530885\\
4.57393	-0.0530332\\
4.58647	-0.0529757\\
4.599	-0.0529166\\
4.61153	-0.0528563\\
4.62406	-0.0527954\\
4.63659	-0.0527345\\
4.64912	-0.052674\\
4.66165	-0.0526144\\
4.67419	-0.0525564\\
4.68672	-0.0525003\\
4.69925	-0.0524468\\
4.71178	-0.0523962\\
4.72431	-0.0523491\\
4.73684	-0.0523059\\
4.74937	-0.0522672\\
4.7619	-0.0522334\\
4.77444	-0.052205\\
4.78697	-0.0521824\\
4.7995	-0.0521661\\
4.81203	-0.0521565\\
4.82456	-0.0521542\\
4.83709	-0.0521596\\
4.84962	-0.052174\\
4.86216	-0.0521982\\
4.87469	-0.0522315\\
4.88722	-0.0522733\\
4.89975	-0.0523228\\
4.91228	-0.0523796\\
4.92481	-0.0524428\\
4.93734	-0.0525119\\
4.94987	-0.0525861\\
4.96241	-0.0526649\\
4.97494	-0.0527476\\
4.98747	-0.0528334\\
5	-0.0529217\\
5	-0.0361582\\
4.98747	-0.035765\\
4.97494	-0.0353763\\
4.96241	-0.0349923\\
4.94987	-0.0346132\\
4.93734	-0.034239\\
4.92481	-0.03387\\
4.91228	-0.0335064\\
4.89975	-0.0331482\\
4.88722	-0.0327956\\
4.87469	-0.0324488\\
4.86216	-0.0321078\\
4.84962	-0.031773\\
4.83709	-0.0314443\\
4.82456	-0.0311219\\
4.81203	-0.0308061\\
4.7995	-0.0304968\\
4.78697	-0.0301942\\
4.77444	-0.0298985\\
4.7619	-0.0296097\\
4.74937	-0.029328\\
4.73684	-0.0290535\\
4.72431	-0.0287863\\
4.71178	-0.0285265\\
4.69925	-0.0282742\\
4.68672	-0.0280295\\
4.67419	-0.0277924\\
4.66165	-0.0275631\\
4.64912	-0.0273416\\
4.63659	-0.027128\\
4.62406	-0.0269224\\
4.61153	-0.0267274\\
4.599	-0.026545\\
4.58647	-0.0263742\\
4.57393	-0.0262143\\
4.5614	-0.0260643\\
4.54887	-0.0259232\\
4.53634	-0.0257902\\
4.52381	-0.0256643\\
4.51128	-0.0255446\\
4.49875	-0.0254302\\
4.48622	-0.02532\\
4.47368	-0.0252132\\
4.46115	-0.0251087\\
4.44862	-0.0250057\\
4.43609	-0.0249031\\
4.42356	-0.0247999\\
4.41103	-0.0246953\\
4.3985	-0.0245881\\
4.38596	-0.0244774\\
4.37343	-0.0243621\\
4.3609	-0.0242414\\
4.34837	-0.024114\\
4.33584	-0.0239792\\
4.32331	-0.0238357\\
4.31078	-0.0236826\\
4.29825	-0.0235091\\
4.28571	-0.0233067\\
4.27318	-0.0230775\\
4.26065	-0.0228236\\
4.24812	-0.0225472\\
4.23559	-0.0222503\\
4.22306	-0.0219351\\
4.21053	-0.0216037\\
4.19799	-0.0212583\\
4.18546	-0.0209009\\
4.17293	-0.0205337\\
4.1604	-0.0201587\\
4.14787	-0.0197781\\
4.13534	-0.019394\\
4.12281	-0.0190085\\
4.11028	-0.0186237\\
4.09774	-0.0182417\\
4.08521	-0.0178646\\
4.07268	-0.0174946\\
4.06015	-0.0171337\\
4.04762	-0.0167842\\
4.03509	-0.016448\\
4.02256	-0.0161273\\
4.01003	-0.0158243\\
3.99749	-0.015541\\
3.98496	-0.0152797\\
3.97243	-0.0150424\\
3.9599	-0.0148314\\
3.94737	-0.0146486\\
3.93484	-0.0144964\\
3.92231	-0.0143769\\
3.90977	-0.0142923\\
3.89724	-0.0142447\\
3.88471	-0.0142363\\
3.87218	-0.0142761\\
3.85965	-0.0143691\\
3.84712	-0.0145111\\
3.83459	-0.0146984\\
3.82206	-0.0149269\\
3.80952	-0.0151928\\
3.79699	-0.0154922\\
3.78446	-0.0158212\\
3.77193	-0.0161759\\
3.7594	-0.0165525\\
3.74687	-0.0169472\\
3.73434	-0.017356\\
3.7218	-0.0177751\\
3.70927	-0.0182008\\
3.69674	-0.0186293\\
3.68421	-0.0190568\\
3.67168	-0.0194795\\
3.65915	-0.0198936\\
3.64662	-0.0202955\\
3.63409	-0.0206813\\
3.62155	-0.0210475\\
3.60902	-0.0213902\\
3.59649	-0.0217059\\
3.58396	-0.0219908\\
3.57143	-0.0222412\\
3.5589	-0.0224536\\
3.54637	-0.0226243\\
3.53383	-0.0227497\\
3.5213	-0.0228262\\
3.50877	-0.0228502\\
3.49624	-0.022818\\
3.48371	-0.0227262\\
3.47118	-0.0225927\\
3.45865	-0.0224381\\
3.44612	-0.0222635\\
3.43358	-0.0220696\\
3.42105	-0.0218574\\
3.40852	-0.0216278\\
3.39599	-0.0213815\\
3.38346	-0.0211196\\
3.37093	-0.020843\\
3.3584	-0.0205525\\
3.34586	-0.0202491\\
3.33333	-0.0199336\\
3.3208	-0.0196071\\
3.30827	-0.0192703\\
3.29574	-0.0189243\\
3.28321	-0.0185699\\
3.27068	-0.0182081\\
3.25815	-0.0178397\\
3.24561	-0.0174657\\
3.23308	-0.017087\\
3.22055	-0.0167044\\
3.20802	-0.0163189\\
3.19549	-0.0159313\\
3.18296	-0.0155424\\
3.17043	-0.0151532\\
3.15789	-0.0147645\\
3.14536	-0.0143771\\
3.13283	-0.0139919\\
3.1203	-0.0136095\\
3.10777	-0.0131517\\
3.09524	-0.0125476\\
3.08271	-0.0118092\\
3.07018	-0.0109488\\
3.05764	-0.00997849\\
3.04511	-0.00891034\\
3.03258	-0.00775643\\
3.02005	-0.00652883\\
3.00752	-0.00523959\\
2.99499	-0.00390073\\
2.98246	-0.00252426\\
2.96992	-0.00112215\\
2.95739	0.000293619\\
2.94486	0.00171111\\
2.93233	0.0031184\\
2.9198	0.00450357\\
2.90727	0.00585475\\
2.89474	0.00716004\\
2.88221	0.00840762\\
2.86967	0.00958562\\
2.85714	0.0106822\\
2.84461	0.0116857\\
2.83208	0.0125841\\
2.81955	0.0133658\\
2.80702	0.014019\\
2.79449	0.0145318\\
2.78195	0.0148927\\
2.76942	0.0150897\\
2.75689	0.0151113\\
2.74436	0.0149456\\
2.73183	0.0146719\\
2.7193	0.0143775\\
2.70677	0.0140634\\
2.69424	0.0137304\\
2.6817	0.0133794\\
2.66917	0.0130112\\
2.65664	0.0126268\\
2.64411	0.0122269\\
2.63158	0.0118125\\
2.61905	0.0113844\\
2.60652	0.0109434\\
2.59398	0.0104904\\
2.58145	0.0100262\\
2.56892	0.00955166\\
2.55639	0.00906763\\
2.54386	0.00857491\\
2.53133	0.00807433\\
2.5188	0.00756669\\
2.50627	0.0070528\\
2.49373	0.00653347\\
2.4812	0.00600949\\
2.46867	0.00548168\\
2.45614	0.00495081\\
2.44361	0.00441769\\
2.43108	0.00388312\\
2.41855	0.00334788\\
2.40602	0.00281278\\
2.39348	0.00227862\\
2.38095	0.00174621\\
2.36842	0.00119101\\
2.35589	0.000601608\\
2.34336	-1.52055e-06\\
2.33083	-0.000597885\\
2.3183	-0.00116698\\
2.30576	-0.00168828\\
2.29323	-0.00214125\\
2.2807	-0.00237674\\
2.26817	-0.00227516\\
2.25564	-0.0018603\\
2.24311	-0.00115591\\
2.23058	-0.000185743\\
2.21805	0.0010265\\
2.20551	0.00245714\\
2.19298	0.00408252\\
2.18045	0.00587903\\
2.16792	0.00782309\\
2.15539	0.00989116\\
2.14286	0.0120597\\
2.13033	0.0143053\\
2.11779	0.0166045\\
2.10526	0.0189338\\
2.09273	0.0212701\\
2.0802	0.0235898\\
2.06767	0.0258699\\
2.05514	0.028087\\
2.04261	0.0302181\\
2.03008	0.0322399\\
2.01754	0.0341294\\
2.00501	0.0358636\\
1.99248	0.0374194\\
1.97995	0.038774\\
1.96742	0.0399043\\
1.95489	0.0407876\\
1.94236	0.041401\\
1.92982	0.0417218\\
1.91729	0.0417271\\
1.90476	0.041475\\
1.89223	0.04105\\
1.8797	0.0404694\\
1.86717	0.0397505\\
1.85464	0.0389105\\
1.84211	0.0379667\\
1.82957	0.0369363\\
1.81704	0.0358366\\
1.80451	0.0346849\\
1.79198	0.0334985\\
1.77945	0.0322944\\
1.76692	0.0310899\\
1.75439	0.0299022\\
1.74185	0.0287483\\
1.72932	0.0276452\\
1.71679	0.0266098\\
1.70426	0.0256591\\
1.69173	0.0248098\\
1.6792	0.0240787\\
1.66667	0.0234822\\
1.65414	0.0230368\\
1.6416	0.022759\\
1.62907	0.0226649\\
1.61654	0.0227707\\
1.60401	0.0230922\\
1.59148	0.0236453\\
1.57895	0.0244456\\
1.56642	0.026013\\
1.55388	0.0286948\\
1.54135	0.0322478\\
1.52882	0.0364284\\
1.51629	0.0409931\\
1.50376	0.0456982\\
1.49123	0.0502997\\
1.4787	0.0545537\\
1.46617	0.0582159\\
1.45363	0.0610422\\
1.4411	0.0633738\\
1.42857	0.0657128\\
1.41604	0.0680562\\
1.40351	0.0704011\\
1.39098	0.072744\\
1.37845	0.0750818\\
1.36591	0.0774112\\
1.35338	0.0797288\\
1.34085	0.0820311\\
1.32832	0.0843148\\
1.31579	0.0865765\\
1.30326	0.0888126\\
1.29073	0.0910198\\
1.2782	0.0931946\\
1.26566	0.0953336\\
1.25313	0.0974335\\
1.2406	0.0994911\\
1.22807	0.101503\\
1.21554	0.103466\\
1.20301	0.105377\\
1.19048	0.107233\\
1.17794	0.109031\\
1.16541	0.110768\\
1.15288	0.111982\\
1.14035	0.112283\\
1.12782	0.111776\\
1.11529	0.110565\\
1.10276	0.108755\\
1.09023	0.10645\\
1.07769	0.103755\\
1.06516	0.100774\\
1.05263	0.0976138\\
1.0401	0.0943779\\
1.02757	0.0911713\\
1.01504	0.088099\\
1.00251	0.0852665\\
0.989975	0.0827788\\
0.977444	0.0807407\\
0.964912	0.0792575\\
0.952381	0.078435\\
0.93985	0.0783785\\
0.927318	0.0791929\\
0.914787	0.080747\\
0.902256	0.082793\\
0.889724	0.0852604\\
0.877193	0.0880783\\
0.864662	0.0911765\\
0.85213	0.0944854\\
0.839599	0.0979348\\
0.827068	0.101454\\
0.814536	0.104972\\
0.802005	0.10842\\
0.789474	0.111728\\
0.776942	0.114825\\
0.764411	0.117642\\
0.75188	0.120109\\
0.739348	0.122157\\
0.726817	0.123715\\
0.714286	0.124715\\
0.701754	0.125353\\
0.689223	0.125884\\
0.676692	0.126318\\
0.66416	0.12667\\
0.651629	0.126954\\
0.639098	0.127186\\
0.626566	0.127376\\
0.614035	0.127539\\
0.601504	0.12769\\
0.588972	0.127844\\
0.576441	0.128011\\
0.56391	0.128202\\
0.551378	0.128435\\
0.538847	0.128723\\
0.526316	0.129931\\
0.513784	0.132742\\
0.501253	0.136887\\
0.488722	0.142095\\
0.47619	0.14809\\
0.463659	0.154599\\
0.451128	0.161356\\
0.438596	0.168094\\
0.426065	0.174538\\
0.413534	0.180417\\
0.401003	0.18547\\
0.388471	0.189431\\
0.37594	0.192025\\
0.363409	0.192979\\
0.350877	0.192903\\
0.338346	0.192586\\
0.325815	0.19202\\
0.313283	0.191196\\
0.300752	0.190117\\
0.288221	0.18878\\
0.275689	0.187161\\
0.263158	0.185239\\
0.250627	0.183016\\
0.238095	0.180493\\
0.225564	0.177653\\
0.213033	0.174494\\
0.200501	0.171048\\
0.18797	0.167352\\
0.175439	0.163423\\
0.162907	0.159287\\
0.150376	0.155847\\
0.137845	0.15498\\
0.125313	0.152125\\
0.112782	0.146196\\
0.100251	0.137831\\
0.0877193	0.127632\\
0.075188	0.116117\\
0.0626566	0.103777\\
0.0501253	0.0928932\\
0.037594	0.0878146\\
0.0250627	0.0676964\\
0.0125313	0.0421548\\
0	0.0120485\\
}--cycle;

\addplot[area legend, draw=none, fill=mycolor13, fill opacity=0.9, forget plot]
table[row sep=crcr] {%
x	y\\
0	-0.0118443\\
0.0125313	-0.0341609\\
0.0250627	-0.053248\\
0.037594	-0.0669288\\
0.0501253	-0.0655363\\
0.0626566	-0.0742048\\
0.075188	-0.0833975\\
0.0877193	-0.0923745\\
0.100251	-0.100346\\
0.112782	-0.106548\\
0.125313	-0.110299\\
0.137845	-0.110956\\
0.150376	-0.10956\\
0.162907	-0.107633\\
0.175439	-0.105271\\
0.18797	-0.102579\\
0.200501	-0.0996272\\
0.213033	-0.096475\\
0.225564	-0.0932036\\
0.238095	-0.0898886\\
0.250627	-0.08657\\
0.263158	-0.0832753\\
0.275689	-0.0800487\\
0.288221	-0.0769352\\
0.300752	-0.0739583\\
0.313283	-0.0711375\\
0.325815	-0.0685119\\
0.338346	-0.0661284\\
0.350877	-0.0640221\\
0.363409	-0.0622277\\
0.37594	-0.0607964\\
0.388471	-0.0597866\\
0.401003	-0.0590629\\
0.413534	-0.0584433\\
0.426065	-0.0579177\\
0.438596	-0.0574801\\
0.451128	-0.0571149\\
0.463659	-0.0568026\\
0.47619	-0.0565307\\
0.488722	-0.0562888\\
0.501253	-0.0560596\\
0.513784	-0.0558219\\
0.526316	-0.0555602\\
0.538847	-0.0552615\\
0.551378	-0.054908\\
0.56391	-0.0544655\\
0.576441	-0.0539218\\
0.588972	-0.0532891\\
0.601504	-0.0525762\\
0.614035	-0.0517898\\
0.626566	-0.0509405\\
0.639098	-0.0500408\\
0.651629	-0.0491008\\
0.66416	-0.0481284\\
0.676692	-0.0471343\\
0.689223	-0.0461308\\
0.701754	-0.0451278\\
0.714286	-0.0441337\\
0.726817	-0.0431584\\
0.739348	-0.0422133\\
0.75188	-0.0413078\\
0.764411	-0.0404499\\
0.776942	-0.039649\\
0.789474	-0.0389152\\
0.802005	-0.0382574\\
0.814536	-0.0376835\\
0.827068	-0.0372023\\
0.839599	-0.0368232\\
0.85213	-0.0365551\\
0.864662	-0.0364056\\
0.877193	-0.0363832\\
0.889724	-0.0364973\\
0.902256	-0.0367564\\
0.914787	-0.0371683\\
0.927318	-0.0377415\\
0.93985	-0.038485\\
0.952381	-0.0394072\\
0.964912	-0.040516\\
0.977444	-0.0418197\\
0.989975	-0.0432713\\
1.00251	-0.044811\\
1.01504	-0.0464283\\
1.02757	-0.0481129\\
1.0401	-0.0498548\\
1.05263	-0.0516437\\
1.06516	-0.0534689\\
1.07769	-0.0553198\\
1.09023	-0.0571862\\
1.10276	-0.0590575\\
1.11529	-0.0609227\\
1.12782	-0.0627712\\
1.14035	-0.0645923\\
1.15288	-0.0663752\\
1.16541	-0.0681089\\
1.17794	-0.0698022\\
1.19048	-0.071471\\
1.20301	-0.0731151\\
1.21554	-0.0747337\\
1.22807	-0.0763263\\
1.2406	-0.0778923\\
1.25313	-0.0794312\\
1.26566	-0.0809421\\
1.2782	-0.0824242\\
1.29073	-0.0838766\\
1.30326	-0.0852987\\
1.31579	-0.0866893\\
1.32832	-0.0880476\\
1.34085	-0.0893726\\
1.35338	-0.0906636\\
1.36591	-0.0919194\\
1.37845	-0.0931392\\
1.39098	-0.0943222\\
1.40351	-0.0954676\\
1.41604	-0.0965745\\
1.42857	-0.0976424\\
1.4411	-0.0986707\\
1.45363	-0.0996587\\
1.46617	-0.100606\\
1.4787	-0.101513\\
1.49123	-0.102378\\
1.50376	-0.103202\\
1.51629	-0.103986\\
1.52882	-0.104728\\
1.54135	-0.10543\\
1.55388	-0.106091\\
1.56642	-0.106713\\
1.57895	-0.107297\\
1.59148	-0.107842\\
1.60401	-0.108352\\
1.61654	-0.108825\\
1.62907	-0.109265\\
1.6416	-0.109672\\
1.65414	-0.110048\\
1.66667	-0.110395\\
1.6792	-0.110714\\
1.69173	-0.111\\
1.70426	-0.111245\\
1.71679	-0.111449\\
1.72932	-0.111613\\
1.74185	-0.111738\\
1.75439	-0.111824\\
1.76692	-0.111873\\
1.77945	-0.111885\\
1.79198	-0.111862\\
1.80451	-0.111805\\
1.81704	-0.111714\\
1.82957	-0.111592\\
1.84211	-0.111438\\
1.85464	-0.111255\\
1.86717	-0.111043\\
1.8797	-0.110803\\
1.89223	-0.110538\\
1.90476	-0.110247\\
1.91729	-0.109932\\
1.92982	-0.109595\\
1.94236	-0.109236\\
1.95489	-0.108856\\
1.96742	-0.108456\\
1.97995	-0.108038\\
1.99248	-0.107602\\
2.00501	-0.10715\\
2.01754	-0.106681\\
2.03008	-0.106198\\
2.04261	-0.1057\\
2.05514	-0.105188\\
2.06767	-0.104663\\
2.0802	-0.104127\\
2.09273	-0.103578\\
2.10526	-0.103019\\
2.11779	-0.102449\\
2.13033	-0.101868\\
2.14286	-0.101278\\
2.15539	-0.100679\\
2.16792	-0.10007\\
2.18045	-0.0994517\\
2.19298	-0.098826\\
2.20551	-0.0981932\\
2.21805	-0.0975541\\
2.23058	-0.0969092\\
2.24311	-0.0962594\\
2.25564	-0.0956053\\
2.26817	-0.0949474\\
2.2807	-0.0942863\\
2.29323	-0.0936227\\
2.30576	-0.092957\\
2.3183	-0.0922899\\
2.33083	-0.0916216\\
2.34336	-0.0909529\\
2.35589	-0.0902841\\
2.36842	-0.0896156\\
2.38095	-0.088948\\
2.39348	-0.0882817\\
2.40602	-0.087617\\
2.41855	-0.0869545\\
2.43108	-0.0862945\\
2.44361	-0.0856374\\
2.45614	-0.0849836\\
2.46867	-0.0843336\\
2.4812	-0.0836878\\
2.49373	-0.0830465\\
2.50627	-0.0824101\\
2.5188	-0.0817792\\
2.53133	-0.081154\\
2.54386	-0.080535\\
2.55639	-0.0799227\\
2.56892	-0.0793174\\
2.58145	-0.0787196\\
2.59398	-0.0781298\\
2.60652	-0.0775483\\
2.61905	-0.0769756\\
2.63158	-0.0764123\\
2.64411	-0.0758587\\
2.65664	-0.0753152\\
2.66917	-0.0747825\\
2.6817	-0.0742609\\
2.69424	-0.073751\\
2.70677	-0.0732531\\
2.7193	-0.0727678\\
2.73183	-0.0722956\\
2.74436	-0.0718368\\
2.75689	-0.0713921\\
2.76942	-0.0709619\\
2.78195	-0.0705465\\
2.79449	-0.0701465\\
2.80702	-0.0697624\\
2.81955	-0.0693945\\
2.83208	-0.0690434\\
2.84461	-0.0687093\\
2.85714	-0.06839\\
2.86967	-0.0680831\\
2.88221	-0.0677884\\
2.89474	-0.0675062\\
2.90727	-0.0672364\\
2.9198	-0.066979\\
2.93233	-0.066734\\
2.94486	-0.0665015\\
2.95739	-0.0662813\\
2.96992	-0.0660734\\
2.98246	-0.0658778\\
2.99499	-0.0656943\\
3.00752	-0.0655228\\
3.02005	-0.0653632\\
3.03258	-0.0652154\\
3.04511	-0.0650791\\
3.05764	-0.0649541\\
3.07018	-0.0648403\\
3.08271	-0.0647374\\
3.09524	-0.0646452\\
3.10777	-0.0645633\\
3.1203	-0.0644916\\
3.13283	-0.0644297\\
3.14536	-0.0643773\\
3.15789	-0.064334\\
3.17043	-0.0642995\\
3.18296	-0.0642735\\
3.19549	-0.0642556\\
3.20802	-0.0642453\\
3.22055	-0.0642424\\
3.23308	-0.0642463\\
3.24561	-0.0642568\\
3.25815	-0.0642736\\
3.27068	-0.0642967\\
3.28321	-0.0643257\\
3.29574	-0.0643601\\
3.30827	-0.0643996\\
3.3208	-0.0644438\\
3.33333	-0.0644922\\
3.34586	-0.0645444\\
3.3584	-0.0646\\
3.37093	-0.0646586\\
3.38346	-0.0647199\\
3.39599	-0.0647833\\
3.40852	-0.0648484\\
3.42105	-0.064915\\
3.43358	-0.0649826\\
3.44612	-0.0650508\\
3.45865	-0.0651192\\
3.47118	-0.0651874\\
3.48371	-0.0652552\\
3.49624	-0.0653221\\
3.50877	-0.0653879\\
3.5213	-0.0654521\\
3.53383	-0.0655144\\
3.54637	-0.0655746\\
3.5589	-0.0656324\\
3.57143	-0.0656874\\
3.58396	-0.0657394\\
3.59649	-0.0657881\\
3.60902	-0.0658334\\
3.62155	-0.0658748\\
3.63409	-0.0659124\\
3.64662	-0.0659457\\
3.65915	-0.0659747\\
3.67168	-0.0659991\\
3.68421	-0.0660188\\
3.69674	-0.0660337\\
3.70927	-0.0660435\\
3.7218	-0.0660483\\
3.73434	-0.0660478\\
3.74687	-0.066042\\
3.7594	-0.0660308\\
3.77193	-0.0660141\\
3.78446	-0.0659919\\
3.79699	-0.0659641\\
3.80952	-0.0659308\\
3.82206	-0.0658918\\
3.83459	-0.0658471\\
3.84712	-0.0657969\\
3.85965	-0.065741\\
3.87218	-0.0656796\\
3.88471	-0.0656126\\
3.89724	-0.0655402\\
3.90977	-0.0654622\\
3.92231	-0.0653781\\
3.93484	-0.0652869\\
3.94737	-0.0651891\\
3.9599	-0.0650849\\
3.97243	-0.0649746\\
3.98496	-0.0648584\\
3.99749	-0.0647367\\
4.01003	-0.0646098\\
4.02256	-0.064478\\
4.03509	-0.0643415\\
4.04762	-0.0642007\\
4.06015	-0.064056\\
4.07268	-0.0639076\\
4.08521	-0.0637559\\
4.09774	-0.0636012\\
4.11028	-0.0634438\\
4.12281	-0.0632841\\
4.13534	-0.0631224\\
4.14787	-0.0629591\\
4.1604	-0.0627943\\
4.17293	-0.0626286\\
4.18546	-0.0624623\\
4.19799	-0.0622956\\
4.21053	-0.0621289\\
4.22306	-0.0619625\\
4.23559	-0.0617968\\
4.24812	-0.0616321\\
4.26065	-0.0614686\\
4.27318	-0.0613068\\
4.28571	-0.061147\\
4.29825	-0.0609894\\
4.31078	-0.0608343\\
4.32331	-0.0606821\\
4.33584	-0.0605331\\
4.34837	-0.0603875\\
4.3609	-0.0602457\\
4.37343	-0.0601079\\
4.38596	-0.0599744\\
4.3985	-0.0598455\\
4.41103	-0.0597214\\
4.42356	-0.0596025\\
4.43609	-0.059489\\
4.44862	-0.059381\\
4.46115	-0.059279\\
4.47368	-0.059183\\
4.48622	-0.0590934\\
4.49875	-0.0590103\\
4.51128	-0.058934\\
4.52381	-0.0588647\\
4.53634	-0.0588027\\
4.54887	-0.058748\\
4.5614	-0.0587013\\
4.57393	-0.0586628\\
4.58647	-0.058632\\
4.599	-0.0586084\\
4.61153	-0.0585915\\
4.62406	-0.0585807\\
4.63659	-0.0585756\\
4.64912	-0.0585757\\
4.66165	-0.0585804\\
4.67419	-0.0585892\\
4.68672	-0.0586015\\
4.69925	-0.0586168\\
4.71178	-0.0586346\\
4.72431	-0.0586543\\
4.73684	-0.0586754\\
4.74937	-0.0586972\\
4.7619	-0.0587193\\
4.77444	-0.0587409\\
4.78697	-0.0587617\\
4.7995	-0.058781\\
4.81203	-0.0587981\\
4.82456	-0.0588126\\
4.83709	-0.0588237\\
4.84962	-0.05883\\
4.86216	-0.0588307\\
4.87469	-0.0588261\\
4.88722	-0.0588168\\
4.89975	-0.0588033\\
4.91228	-0.058786\\
4.92481	-0.0587654\\
4.93734	-0.058742\\
4.94987	-0.0587162\\
4.96241	-0.0586886\\
4.97494	-0.0586596\\
4.98747	-0.0586296\\
5	-0.0585992\\
5	-0.0529217\\
4.98747	-0.0528334\\
4.97494	-0.0527476\\
4.96241	-0.0526649\\
4.94987	-0.0525861\\
4.93734	-0.0525119\\
4.92481	-0.0524428\\
4.91228	-0.0523796\\
4.89975	-0.0523228\\
4.88722	-0.0522733\\
4.87469	-0.0522315\\
4.86216	-0.0521982\\
4.84962	-0.052174\\
4.83709	-0.0521596\\
4.82456	-0.0521542\\
4.81203	-0.0521565\\
4.7995	-0.0521661\\
4.78697	-0.0521824\\
4.77444	-0.052205\\
4.7619	-0.0522334\\
4.74937	-0.0522672\\
4.73684	-0.0523059\\
4.72431	-0.0523491\\
4.71178	-0.0523962\\
4.69925	-0.0524468\\
4.68672	-0.0525003\\
4.67419	-0.0525564\\
4.66165	-0.0526144\\
4.64912	-0.052674\\
4.63659	-0.0527345\\
4.62406	-0.0527954\\
4.61153	-0.0528563\\
4.599	-0.0529166\\
4.58647	-0.0529757\\
4.57393	-0.0530332\\
4.5614	-0.0530885\\
4.54887	-0.053141\\
4.53634	-0.0531908\\
4.52381	-0.0532384\\
4.51128	-0.0532839\\
4.49875	-0.0533275\\
4.48622	-0.0533692\\
4.47368	-0.053409\\
4.46115	-0.0534471\\
4.44862	-0.0534835\\
4.43609	-0.0535183\\
4.42356	-0.0535515\\
4.41103	-0.0535831\\
4.3985	-0.0536133\\
4.38596	-0.053642\\
4.37343	-0.0536693\\
4.3609	-0.0536951\\
4.34837	-0.0537196\\
4.33584	-0.0537426\\
4.32331	-0.0537643\\
4.31078	-0.0537846\\
4.29825	-0.0538035\\
4.28571	-0.0538211\\
4.27318	-0.0538372\\
4.26065	-0.0538519\\
4.24812	-0.0538652\\
4.23559	-0.053877\\
4.22306	-0.0538873\\
4.21053	-0.0538961\\
4.19799	-0.0539034\\
4.18546	-0.053909\\
4.17293	-0.0539131\\
4.1604	-0.0539155\\
4.14787	-0.0539162\\
4.13534	-0.0539151\\
4.12281	-0.0539123\\
4.11028	-0.0539077\\
4.09774	-0.0539013\\
4.08521	-0.0538929\\
4.07268	-0.0538827\\
4.06015	-0.0538704\\
4.04762	-0.0538563\\
4.03509	-0.0538401\\
4.02256	-0.0538218\\
4.01003	-0.0538015\\
3.99749	-0.0537792\\
3.98496	-0.0537547\\
3.97243	-0.0537282\\
3.9599	-0.0536996\\
3.94737	-0.0536689\\
3.93484	-0.0536361\\
3.92231	-0.0536013\\
3.90977	-0.0535645\\
3.89724	-0.0535247\\
3.88471	-0.053481\\
3.87218	-0.0534331\\
3.85965	-0.053381\\
3.84712	-0.0533246\\
3.83459	-0.0532638\\
3.82206	-0.0531984\\
3.80952	-0.0531284\\
3.79699	-0.0530537\\
3.78446	-0.0529743\\
3.77193	-0.05289\\
3.7594	-0.052801\\
3.74687	-0.052707\\
3.73434	-0.0526082\\
3.7218	-0.0525046\\
3.70927	-0.0523961\\
3.69674	-0.0522827\\
3.68421	-0.0521646\\
3.67168	-0.0520418\\
3.65915	-0.0519144\\
3.64662	-0.0517824\\
3.63409	-0.051646\\
3.62155	-0.0515052\\
3.60902	-0.0513602\\
3.59649	-0.0512111\\
3.58396	-0.0510581\\
3.57143	-0.0509014\\
3.5589	-0.0507411\\
3.54637	-0.0505774\\
3.53383	-0.0504105\\
3.5213	-0.0502408\\
3.50877	-0.0500682\\
3.49624	-0.0498933\\
3.48371	-0.0497161\\
3.47118	-0.0495369\\
3.45865	-0.0493561\\
3.44612	-0.0491738\\
3.43358	-0.0489905\\
3.42105	-0.0488063\\
3.40852	-0.0486216\\
3.39599	-0.0484368\\
3.38346	-0.0482521\\
3.37093	-0.0480678\\
3.3584	-0.0478843\\
3.34586	-0.0477019\\
3.33333	-0.047521\\
3.3208	-0.0473418\\
3.30827	-0.0471648\\
3.29574	-0.0469902\\
3.28321	-0.0468184\\
3.27068	-0.0466497\\
3.25815	-0.0464844\\
3.24561	-0.0463229\\
3.23308	-0.0461658\\
3.22055	-0.0460137\\
3.20802	-0.0458668\\
3.19549	-0.0457254\\
3.18296	-0.0455898\\
3.17043	-0.0454602\\
3.15789	-0.0453368\\
3.14536	-0.0452199\\
3.13283	-0.0451096\\
3.1203	-0.0450063\\
3.10777	-0.0449099\\
3.09524	-0.0448208\\
3.08271	-0.044739\\
3.07018	-0.0446646\\
3.05764	-0.0445979\\
3.04511	-0.0445388\\
3.03258	-0.0444875\\
3.02005	-0.044444\\
3.00752	-0.0444084\\
2.99499	-0.0443806\\
2.98246	-0.0443606\\
2.96992	-0.0443485\\
2.95739	-0.0443442\\
2.94486	-0.0443476\\
2.93233	-0.0443587\\
2.9198	-0.0443773\\
2.90727	-0.0444033\\
2.89474	-0.0444366\\
2.88221	-0.0444771\\
2.86967	-0.0445245\\
2.85714	-0.0445788\\
2.84461	-0.0446396\\
2.83208	-0.0447041\\
2.81955	-0.0447698\\
2.80702	-0.0448367\\
2.79449	-0.0449049\\
2.78195	-0.0449745\\
2.76942	-0.0450454\\
2.75689	-0.0451177\\
2.74436	-0.0451915\\
2.73183	-0.0452669\\
2.7193	-0.0453437\\
2.70677	-0.0454221\\
2.69424	-0.045502\\
2.6817	-0.0455835\\
2.66917	-0.0456666\\
2.65664	-0.0457514\\
2.64411	-0.0458378\\
2.63158	-0.0459259\\
2.61905	-0.0460156\\
2.60652	-0.0461072\\
2.59398	-0.0462005\\
2.58145	-0.0462957\\
2.56892	-0.0463927\\
2.55639	-0.0464917\\
2.54386	-0.0465927\\
2.53133	-0.0466958\\
2.5188	-0.046801\\
2.50627	-0.0469085\\
2.49373	-0.0470182\\
2.4812	-0.0471304\\
2.46867	-0.047245\\
2.45614	-0.0473622\\
2.44361	-0.047482\\
2.43108	-0.0476046\\
2.41855	-0.0477301\\
2.40602	-0.0478585\\
2.39348	-0.0479898\\
2.38095	-0.0481243\\
2.36842	-0.048262\\
2.35589	-0.0484028\\
2.34336	-0.048547\\
2.33083	-0.0486944\\
2.3183	-0.0488452\\
2.30576	-0.0489994\\
2.29323	-0.0491569\\
2.2807	-0.0493177\\
2.26817	-0.0494819\\
2.25564	-0.0496492\\
2.24311	-0.0498197\\
2.23058	-0.0499931\\
2.21805	-0.0501695\\
2.20551	-0.0503485\\
2.19298	-0.05053\\
2.18045	-0.0507138\\
2.16792	-0.0508995\\
2.15539	-0.0510869\\
2.14286	-0.0512757\\
2.13033	-0.0514648\\
2.11779	-0.0516528\\
2.10526	-0.0518387\\
2.09273	-0.0520214\\
2.0802	-0.0522\\
2.06767	-0.0523733\\
2.05514	-0.0525401\\
2.04261	-0.0526991\\
2.03008	-0.0528492\\
2.01754	-0.052989\\
2.00501	-0.0531172\\
1.99248	-0.0532323\\
1.97995	-0.053333\\
1.96742	-0.0534178\\
1.95489	-0.0534851\\
1.94236	-0.0535334\\
1.92982	-0.0535611\\
1.91729	-0.0535665\\
1.90476	-0.053548\\
1.89223	-0.053504\\
1.8797	-0.0534326\\
1.86717	-0.0533323\\
1.85464	-0.0532011\\
1.84211	-0.0530375\\
1.82957	-0.0528395\\
1.81704	-0.0526056\\
1.80451	-0.052334\\
1.79198	-0.0520228\\
1.77945	-0.0516706\\
1.76692	-0.0512756\\
1.75439	-0.0508361\\
1.74185	-0.0503507\\
1.72932	-0.0498179\\
1.71679	-0.0492362\\
1.70426	-0.0486042\\
1.69173	-0.0479208\\
1.6792	-0.0471847\\
1.66667	-0.0463869\\
1.65414	-0.0455195\\
1.6416	-0.0445832\\
1.62907	-0.0435792\\
1.61654	-0.0425083\\
1.60401	-0.041372\\
1.59148	-0.0401715\\
1.57895	-0.0389084\\
1.56642	-0.0375844\\
1.55388	-0.0362012\\
1.54135	-0.0347608\\
1.52882	-0.0332653\\
1.51629	-0.0317171\\
1.50376	-0.0301185\\
1.49123	-0.028472\\
1.4787	-0.0267804\\
1.46617	-0.0250465\\
1.45363	-0.0232732\\
1.4411	-0.0214635\\
1.42857	-0.0196206\\
1.41604	-0.0177479\\
1.40351	-0.0158485\\
1.39098	-0.013926\\
1.37845	-0.0119838\\
1.36591	-0.0100256\\
1.35338	-0.00805482\\
1.34085	-0.00607518\\
1.32832	-0.00409038\\
1.31579	-0.00210407\\
1.30326	-0.000119867\\
1.29073	0.00185853\\
1.2782	0.00382742\\
1.26566	0.00578319\\
1.25313	0.00772233\\
1.2406	0.0096413\\
1.22807	0.0115365\\
1.21554	0.0134046\\
1.20301	0.0152422\\
1.19048	0.0170461\\
1.17794	0.0188128\\
1.16541	0.0205393\\
1.15288	0.0222064\\
1.14035	0.0238019\\
1.12782	0.0253333\\
1.11529	0.0268084\\
1.10276	0.0282352\\
1.09023	0.0296215\\
1.07769	0.030975\\
1.06516	0.0323038\\
1.05263	0.0336162\\
1.0401	0.0349202\\
1.02757	0.0362237\\
1.01504	0.0375349\\
1.00251	0.0388625\\
0.989975	0.0402147\\
0.977444	0.0415997\\
0.964912	0.0430755\\
0.952381	0.0446885\\
0.93985	0.046429\\
0.927318	0.0482867\\
0.914787	0.050252\\
0.902256	0.0523158\\
0.889724	0.0544685\\
0.877193	0.0566999\\
0.864662	0.0590005\\
0.85213	0.0613613\\
0.839599	0.0637729\\
0.827068	0.0662249\\
0.814536	0.068708\\
0.802005	0.0712136\\
0.789474	0.0737325\\
0.776942	0.0762543\\
0.764411	0.0787698\\
0.75188	0.0812713\\
0.739348	0.0837499\\
0.726817	0.0861957\\
0.714286	0.0885999\\
0.701754	0.0909557\\
0.689223	0.0932552\\
0.676692	0.0954885\\
0.66416	0.0976475\\
0.651629	0.0997264\\
0.639098	0.101718\\
0.626566	0.103612\\
0.614035	0.1054\\
0.601504	0.107077\\
0.588972	0.108635\\
0.576441	0.110063\\
0.56391	0.111352\\
0.551378	0.112496\\
0.538847	0.113508\\
0.526316	0.114414\\
0.513784	0.115228\\
0.501253	0.115971\\
0.488722	0.11666\\
0.47619	0.117305\\
0.463659	0.117923\\
0.451128	0.118534\\
0.438596	0.11916\\
0.426065	0.119814\\
0.413534	0.120511\\
0.401003	0.12128\\
0.388471	0.122143\\
0.37594	0.123272\\
0.363409	0.124801\\
0.350877	0.126686\\
0.338346	0.128878\\
0.325815	0.131312\\
0.313283	0.133922\\
0.300752	0.136654\\
0.288221	0.139448\\
0.275689	0.142221\\
0.263158	0.144896\\
0.250627	0.147419\\
0.238095	0.149732\\
0.225564	0.151763\\
0.213033	0.153452\\
0.200501	0.154772\\
0.18797	0.155705\\
0.175439	0.15621\\
0.162907	0.156255\\
0.150376	0.155847\\
0.137845	0.15498\\
0.125313	0.152125\\
0.112782	0.146196\\
0.100251	0.137831\\
0.0877193	0.127632\\
0.075188	0.116117\\
0.0626566	0.103777\\
0.0501253	0.0911556\\
0.037594	0.0875713\\
0.0250627	0.0674639\\
0.0125313	0.0403347\\
0	0.0118445\\
}--cycle;
\addplot [color=mycolor14, line width=1.0pt, forget plot]
  table[row sep=crcr]{%
-0.125	6.8378e-08\\
-0.1224	6.97937e-08\\
-0.1199	7.1155e-08\\
-0.1173	7.25706e-08\\
-0.1147	7.3986e-08\\
-0.1122	7.5347e-08\\
-0.1096	7.67623e-08\\
-0.1071	7.81231e-08\\
-0.1045	7.95382e-08\\
-0.1019	8.09532e-08\\
-0.0994	8.23136e-08\\
-0.0968	8.37284e-08\\
-0.0942	8.51431e-08\\
-0.0917	8.65032e-08\\
-0.0891	8.79176e-08\\
-0.0865	8.9332e-08\\
-0.084	9.06919e-08\\
-0.0814	9.21061e-08\\
-0.0789	9.34659e-08\\
-0.0763	9.48798e-08\\
-0.0737	9.62937e-08\\
-0.0712	9.76531e-08\\
-0.0686	9.90668e-08\\
-0.066	1.0048e-07\\
-0.0635	1.0184e-07\\
-0.0609	1.03253e-07\\
-0.0583	1.04666e-07\\
-0.0558	1.06025e-07\\
-0.0532	1.07438e-07\\
-0.0507	1.08797e-07\\
-0.0481	1.1021e-07\\
-0.0455	1.11623e-07\\
-0.043	1.12981e-07\\
-0.0404	1.14394e-07\\
-0.0378	1.15806e-07\\
-0.0353	1.17164e-07\\
-0.0327	1.18577e-07\\
-0.0301	1.19989e-07\\
-0.0276	1.21347e-07\\
-0.025	1.22759e-07\\
-0.0224	1.24171e-07\\
-0.0199	1.25528e-07\\
-0.0173	1.2694e-07\\
-0.0148	1.28297e-07\\
-0.0122	1.29709e-07\\
-0.0096	1.3112e-07\\
-0.0071	1.32477e-07\\
-0.0045	1.33889e-07\\
-0.0019	1.353e-07\\
0.0006	1.36657e-07\\
0.0032	0.000277263\\
0.0058	0.00138356\\
0.0083	0.00239693\\
0.0109	0.00339331\\
0.0134	0.00438786\\
0.016	0.00545394\\
0.0186	0.00649038\\
0.0211	0.00741747\\
0.0237	0.00831552\\
0.0263	0.00918595\\
0.0288	0.0100315\\
0.0314	0.0109293\\
0.034	0.0118311\\
0.0365	0.0126821\\
0.0391	0.0135413\\
0.0416	0.0143481\\
0.0442	0.0151808\\
0.0468	0.0160169\\
0.0493	0.016824\\
0.0519	0.0176582\\
0.0545	0.018479\\
0.057	0.0192536\\
0.0596	0.020049\\
0.0622	0.0208416\\
0.0647	0.0216045\\
0.0673	0.0223961\\
0.0698	0.0231494\\
0.0724	0.0239201\\
0.075	0.0246786\\
0.0775	0.0254015\\
0.0801	0.0261514\\
0.0827	0.0268998\\
0.0852	0.027614\\
0.0878	0.0283457\\
0.0904	0.0290643\\
0.0929	0.0297455\\
0.0955	0.0304488\\
0.098	0.0311225\\
0.1006	0.0318189\\
0.1032	0.0325058\\
0.1057	0.0331539\\
0.1083	0.0338154\\
0.1109	0.0344683\\
0.1134	0.0350917\\
0.116	0.035736\\
0.1186	0.0363727\\
0.1211	0.0369737\\
0.1237	0.0375854\\
0.1263	0.0381861\\
0.1288	0.0387573\\
0.1314	0.039347\\
0.1339	0.0399086\\
0.1365	0.0404832\\
0.1391	0.0410452\\
0.1416	0.041574\\
0.1442	0.0421157\\
0.1468	0.0426524\\
0.1493	0.0431643\\
0.1519	0.0436898\\
0.1545	0.0442048\\
0.157	0.0446891\\
0.1596	0.0451835\\
0.1621	0.0456535\\
0.1647	0.046139\\
0.1673	0.0466201\\
0.1698	0.0470756\\
0.1724	0.0475395\\
0.175	0.0479943\\
0.1775	0.0484257\\
0.1801	0.0488715\\
0.1827	0.0493149\\
0.1852	0.0497369\\
0.1878	0.0501687\\
0.1903	0.0505761\\
0.1929	0.0509937\\
0.1955	0.0514083\\
0.198	0.0518058\\
0.2006	0.0522174\\
0.2032	0.0526244\\
0.2057	0.0530097\\
0.2083	0.0534046\\
0.2109	0.0537963\\
0.2134	0.0541722\\
0.216	0.054563\\
0.2185	0.0549369\\
0.2211	0.0553213\\
0.2237	0.0557007\\
0.2262	0.056062\\
0.2288	0.0564368\\
0.2314	0.0568121\\
0.2339	0.0571728\\
0.2365	0.0575453\\
0.2391	0.0579137\\
0.2416	0.0582643\\
0.2442	0.0586271\\
0.2467	0.0589763\\
0.2493	0.0593399\\
0.2519	0.0597023\\
0.2544	0.0600478\\
0.257	0.0604032\\
0.2596	0.0607557\\
0.2621	0.0610939\\
0.2647	0.061446\\
0.2673	0.0617976\\
0.2698	0.0621332\\
0.2724	0.0624782\\
0.2749	0.0628063\\
0.2775	0.0631452\\
0.2801	0.0634833\\
0.2826	0.0638078\\
0.2852	0.064143\\
0.2878	0.0644741\\
0.2903	0.0647879\\
0.2929	0.0651103\\
0.2955	0.0654302\\
0.298	0.0657362\\
0.3006	0.0660521\\
0.3032	0.0663636\\
0.3057	0.0666579\\
0.3083	0.0669583\\
0.3108	0.0672431\\
0.3134	0.0675362\\
0.316	0.0678261\\
0.3185	0.0681005\\
0.3211	0.0683798\\
0.3237	0.0686521\\
0.3262	0.0689082\\
0.3288	0.0691696\\
0.3314	0.0694267\\
0.3339	0.0696691\\
0.3365	0.0699147\\
0.339	0.0701435\\
0.3416	0.0703739\\
0.3442	0.0705976\\
0.3467	0.0708074\\
0.3493	0.0710199\\
0.3519	0.0712255\\
0.3544	0.0714155\\
0.357	0.0716045\\
0.3596	0.0717853\\
0.3621	0.0719526\\
0.3647	0.0721202\\
0.3672	0.0722747\\
0.3698	0.0724272\\
0.3724	0.0725705\\
0.3749	0.0726996\\
0.3775	0.0728259\\
0.3801	0.0729449\\
0.3826	0.0730525\\
0.3852	0.0731563\\
0.3878	0.0732509\\
0.3903	0.0733328\\
0.3929	0.0734093\\
0.3954	0.0734752\\
0.398	0.0735364\\
0.4006	0.0735897\\
0.4031	0.0736326\\
0.4057	0.0736678\\
0.4083	0.0736937\\
0.4108	0.0737106\\
0.4134	0.0737207\\
0.416	0.0737231\\
0.4185	0.0737175\\
0.4211	0.0737029\\
0.4236	0.07368\\
0.4262	0.0736477\\
0.4288	0.0736075\\
0.4313	0.0735618\\
0.4339	0.0735069\\
0.4365	0.0734435\\
0.439	0.0733742\\
0.4416	0.0732937\\
0.4442	0.0732052\\
0.4467	0.0731133\\
0.4493	0.0730109\\
0.4519	0.0729008\\
0.4544	0.0727872\\
0.457	0.0726608\\
0.4595	0.072532\\
0.4621	0.072391\\
0.4647	0.0722434\\
0.4672	0.0720951\\
0.4698	0.0719335\\
0.4724	0.0717641\\
0.4749	0.0715941\\
0.4775	0.0714104\\
0.4801	0.0712205\\
0.4826	0.0710319\\
0.4852	0.0708293\\
0.4877	0.0706278\\
0.4903	0.0704111\\
0.4929	0.0701876\\
0.4954	0.0699669\\
0.498	0.0697317\\
0.5006	0.0694905\\
0.5031	0.0692525\\
0.5057	0.0689984\\
0.5083	0.0687377\\
0.5108	0.0684813\\
0.5134	0.0682092\\
0.5159	0.0679425\\
0.5185	0.0676594\\
0.5211	0.06737\\
0.5236	0.0670858\\
0.5262	0.0667844\\
0.5288	0.0664776\\
0.5313	0.0661778\\
0.5339	0.0658607\\
0.5365	0.0655379\\
0.539	0.065222\\
0.5416	0.0648876\\
0.5441	0.064561\\
0.5467	0.0642165\\
0.5493	0.063867\\
0.5518	0.0635261\\
0.5544	0.063166\\
0.557	0.0628002\\
0.5595	0.0624435\\
0.5621	0.0620676\\
0.5647	0.0616871\\
0.5672	0.0613166\\
0.5698	0.0609261\\
0.5723	0.0605456\\
0.5749	0.0601446\\
0.5775	0.0597387\\
0.58	0.0593441\\
0.5826	0.0589292\\
0.5852	0.0585095\\
0.5877	0.0581011\\
0.5903	0.0576713\\
0.5929	0.0572366\\
0.5954	0.0568143\\
0.598	0.0563709\\
0.6006	0.0559229\\
0.6031	0.0554877\\
0.6057	0.0550303\\
0.6082	0.0545859\\
0.6108	0.0541192\\
0.6134	0.0536484\\
0.6159	0.0531917\\
0.6185	0.0527124\\
0.6211	0.0522286\\
0.6236	0.0517589\\
0.6262	0.0512661\\
0.6288	0.0507692\\
0.6313	0.0502878\\
0.6339	0.0497831\\
0.6364	0.0492938\\
0.639	0.0487806\\
0.6416	0.0482633\\
0.6441	0.047762\\
0.6467	0.0472371\\
0.6493	0.0467085\\
0.6518	0.0461966\\
0.6544	0.0456603\\
0.657	0.04512\\
0.6595	0.0445969\\
0.6621	0.0440495\\
0.6646	0.0435199\\
0.6672	0.0429658\\
0.6698	0.0424082\\
0.6723	0.0418686\\
0.6749	0.0413039\\
0.6775	0.0407359\\
0.68	0.0401869\\
0.6826	0.0396129\\
0.6852	0.0390359\\
0.6877	0.038478\\
0.6903	0.0378946\\
0.6928	0.0373308\\
0.6954	0.0367417\\
0.698	0.03615\\
0.7005	0.0355785\\
0.7031	0.0349815\\
0.7057	0.0343816\\
0.7082	0.0338022\\
0.7108	0.0331971\\
0.7134	0.0325898\\
0.7159	0.0320037\\
0.7185	0.0313918\\
0.721	0.0308013\\
0.7236	0.0301847\\
0.7262	0.0295659\\
0.7287	0.028969\\
0.7313	0.0283464\\
0.7339	0.0277219\\
0.7364	0.0271196\\
0.739	0.0264911\\
0.7416	0.0258607\\
0.7441	0.025253\\
0.7467	0.0246194\\
0.7492	0.0240087\\
0.7518	0.0233721\\
0.7544	0.0227338\\
0.7569	0.0221186\\
0.7595	0.0214773\\
0.7621	0.0208347\\
0.7646	0.0202157\\
0.7672	0.0195708\\
0.7698	0.0189246\\
0.7723	0.018302\\
0.7749	0.0176534\\
0.7775	0.0170037\\
0.78	0.0163782\\
0.7826	0.0157268\\
0.7851	0.0150996\\
0.7877	0.0144464\\
0.7903	0.0137922\\
0.7928	0.0131624\\
0.7954	0.0125068\\
0.798	0.0118506\\
0.8005	0.0112191\\
0.8031	0.0105617\\
0.8057	0.00990361\\
0.8082	0.0092703\\
0.8108	0.00861121\\
0.8133	0.00797716\\
0.8159	0.00731747\\
0.8185	0.00665744\\
0.821	0.00602245\\
0.8236	0.0053617\\
0.8262	0.0047007\\
0.8287	0.004065\\
0.8313	0.00340384\\
0.8339	0.00274264\\
0.8364	0.00210678\\
0.839	0.00144539\\
0.8415	0.000809382\\
0.8441	0.00014801\\
0.8467	-0.000513181\\
0.8492	-0.00114873\\
0.8518	-0.00180952\\
0.8544	-0.00247015\\
0.8569	-0.0031052\\
0.8595	-0.00376538\\
0.8621	-0.00442518\\
0.8646	-0.00505917\\
0.8672	-0.00571805\\
0.8697	-0.00635119\\
0.8723	-0.00700923\\
0.8749	-0.00766679\\
0.8774	-0.00829851\\
0.88	-0.00895481\\
0.8826	-0.00961039\\
0.8851	-0.0102401\\
0.8877	-0.0108943\\
0.8903	-0.0115478\\
0.8928	-0.0121753\\
0.8954	-0.0128271\\
0.8979	-0.013453\\
0.9005	-0.0141028\\
0.9031	-0.0147517\\
0.9056	-0.0153747\\
0.9082	-0.0160217\\
0.9108	-0.0166675\\
0.9133	-0.0172873\\
0.9159	-0.0179306\\
0.9185	-0.0185727\\
0.921	-0.019189\\
0.9236	-0.0198286\\
0.9262	-0.0204669\\
0.9287	-0.0210792\\
0.9313	-0.0217146\\
0.9338	-0.022324\\
0.9364	-0.0229564\\
0.939	-0.0235872\\
0.9415	-0.0241922\\
0.9441	-0.0248198\\
0.9467	-0.0254456\\
0.9492	-0.0260456\\
0.9518	-0.0266679\\
0.9544	-0.0272883\\
0.9569	-0.0278832\\
0.9595	-0.0284999\\
0.962	-0.029091\\
0.9646	-0.0297038\\
0.9672	-0.0303144\\
0.9697	-0.0308997\\
0.9723	-0.0315062\\
0.9749	-0.0321106\\
0.9774	-0.0326897\\
0.98	-0.0332896\\
0.9826	-0.0338872\\
0.9851	-0.0344596\\
0.9877	-0.0350526\\
0.9902	-0.0356205\\
0.9928	-0.0362087\\
0.9954	-0.0367945\\
0.9979	-0.0373552\\
1.0005	-0.0379358\\
1.0031	-0.0385138\\
1.0056	-0.0390672\\
1.0082	-0.03964\\
1.0108	-0.0402101\\
1.0133	-0.0407556\\
1.0159	-0.0413202\\
1.0184	-0.0418604\\
1.021	-0.0424194\\
1.0236	-0.0429755\\
1.0261	-0.0435075\\
1.0287	-0.0440579\\
1.0313	-0.0446052\\
1.0338	-0.0451285\\
1.0364	-0.0456698\\
1.039	-0.0462081\\
1.0415	-0.0467227\\
1.0441	-0.0472548\\
1.0466	-0.0477634\\
1.0492	-0.0482891\\
1.0518	-0.0488116\\
1.0543	-0.0493109\\
1.0569	-0.049827\\
1.0595	-0.0503398\\
1.062	-0.0508297\\
1.0646	-0.0513358\\
1.0672	-0.0518384\\
1.0697	-0.0523185\\
1.0723	-0.0528145\\
1.0748	-0.0532881\\
1.0774	-0.0537772\\
1.08	-0.0542627\\
1.0825	-0.0547261\\
1.0851	-0.0552046\\
1.0877	-0.0556795\\
1.0902	-0.0561328\\
1.0928	-0.0566005\\
1.0954	-0.0570646\\
1.0979	-0.0575074\\
1.1005	-0.0579642\\
1.1031	-0.0584173\\
1.1056	-0.0588494\\
1.1082	-0.0592951\\
1.1107	-0.0597201\\
1.1133	-0.0601583\\
1.1159	-0.0605927\\
1.1184	-0.0610068\\
1.121	-0.0614336\\
1.1236	-0.0618566\\
1.1261	-0.0622596\\
1.1287	-0.0626749\\
1.1313	-0.0630863\\
1.1338	-0.0634781\\
1.1364	-0.0638817\\
1.1389	-0.0642661\\
1.1415	-0.0646619\\
1.1441	-0.0650537\\
1.1466	-0.0654267\\
1.1492	-0.0658106\\
1.1518	-0.0661905\\
1.1543	-0.0665519\\
1.1569	-0.0669239\\
1.1595	-0.0672917\\
1.162	-0.0676416\\
1.1646	-0.0680015\\
1.1671	-0.0683436\\
1.1697	-0.0686954\\
1.1723	-0.0690431\\
1.1748	-0.0693735\\
1.1774	-0.0697131\\
1.18	-0.0700486\\
1.1825	-0.0703672\\
1.1851	-0.0706945\\
1.1877	-0.0710176\\
1.1902	-0.0713245\\
1.1928	-0.0716395\\
1.1953	-0.0719384\\
1.1979	-0.0722452\\
1.2005	-0.0725479\\
1.203	-0.0728349\\
1.2056	-0.0731293\\
1.2082	-0.0734196\\
1.2107	-0.0736948\\
1.2133	-0.0739768\\
1.2159	-0.0742547\\
1.2184	-0.074518\\
1.221	-0.0747876\\
1.2235	-0.075043\\
1.2261	-0.0753045\\
1.2287	-0.0755619\\
1.2312	-0.0758053\\
1.2338	-0.0760545\\
1.2364	-0.0762995\\
1.2389	-0.0765311\\
1.2415	-0.076768\\
1.2441	-0.0770007\\
1.2466	-0.0772206\\
1.2492	-0.0774452\\
1.2518	-0.0776657\\
1.2543	-0.0778738\\
1.2569	-0.0780862\\
1.2594	-0.0782867\\
1.262	-0.0784911\\
1.2646	-0.0786914\\
1.2671	-0.0788802\\
1.2697	-0.0790726\\
1.2723	-0.079261\\
1.2748	-0.0794383\\
1.2774	-0.0796188\\
1.28	-0.0797952\\
1.2825	-0.0799612\\
1.2851	-0.0801298\\
1.2876	-0.0802883\\
1.2902	-0.0804492\\
1.2928	-0.0806062\\
1.2953	-0.0807535\\
1.2979	-0.0809029\\
1.3005	-0.0810483\\
1.303	-0.0811846\\
1.3056	-0.0813225\\
1.3082	-0.0814566\\
1.3107	-0.081582\\
1.3133	-0.0817087\\
1.3158	-0.0818269\\
1.3184	-0.0819463\\
1.321	-0.0820619\\
1.3235	-0.0821695\\
1.3261	-0.0822779\\
1.3287	-0.0823826\\
1.3312	-0.0824799\\
1.3338	-0.0825775\\
1.3364	-0.0826715\\
1.3389	-0.0827585\\
1.3415	-0.0828456\\
1.344	-0.0829259\\
1.3466	-0.0830061\\
1.3492	-0.0830828\\
1.3517	-0.0831533\\
1.3543	-0.0832232\\
1.3569	-0.0832898\\
1.3594	-0.0833506\\
1.362	-0.0834106\\
1.3646	-0.0834673\\
1.3671	-0.0835187\\
1.3697	-0.0835689\\
1.3722	-0.0836142\\
1.3748	-0.0836581\\
1.3774	-0.0836988\\
1.3799	-0.083735\\
1.3825	-0.0837696\\
1.3851	-0.0838011\\
1.3876	-0.0838285\\
1.3902	-0.083854\\
1.3928	-0.0838765\\
1.3953	-0.0838953\\
1.3979	-0.0839119\\
1.4005	-0.0839257\\
1.403	-0.0839361\\
1.4056	-0.0839442\\
1.4081	-0.0839493\\
1.4107	-0.0839519\\
1.4133	-0.0839517\\
1.4158	-0.0839489\\
1.4184	-0.0839433\\
1.421	-0.083935\\
1.4235	-0.0839246\\
1.4261	-0.0839112\\
1.4287	-0.0838952\\
1.4312	-0.0838774\\
1.4338	-0.0838564\\
1.4363	-0.0838338\\
1.4389	-0.083808\\
1.4415	-0.0837797\\
1.444	-0.0837502\\
1.4466	-0.0837172\\
1.4492	-0.0836819\\
1.4517	-0.0836458\\
1.4543	-0.0836059\\
1.4569	-0.0835639\\
1.4594	-0.0835214\\
1.462	-0.083475\\
1.4645	-0.0834284\\
1.4671	-0.0833779\\
1.4697	-0.0833252\\
1.4722	-0.0832727\\
1.4748	-0.0832161\\
1.4774	-0.0831575\\
1.4799	-0.0830993\\
1.4825	-0.0830369\\
1.4851	-0.0829726\\
1.4876	-0.082909\\
1.4902	-0.0828411\\
1.4927	-0.0827741\\
1.4953	-0.0827027\\
1.4979	-0.0826296\\
1.5004	-0.0825576\\
1.503	-0.0824812\\
1.5056	-0.0824031\\
1.5081	-0.0823265\\
1.5107	-0.0822453\\
1.5133	-0.0821626\\
1.5158	-0.0820816\\
1.5184	-0.0819959\\
1.5209	-0.0819122\\
1.5235	-0.0818237\\
1.5261	-0.0817339\\
1.5286	-0.0816462\\
1.5312	-0.0815537\\
1.5338	-0.0814599\\
1.5363	-0.0813685\\
1.5389	-0.0812722\\
1.5415	-0.0811748\\
1.544	-0.08108\\
1.5466	-0.0809803\\
1.5491	-0.0808833\\
1.5517	-0.0807815\\
1.5543	-0.0806785\\
1.5568	-0.0805786\\
1.5594	-0.0804737\\
1.562	-0.0803678\\
1.5645	-0.0802651\\
1.5671	-0.0801574\\
1.5697	-0.0800488\\
1.5722	-0.0799436\\
1.5748	-0.0798334\\
1.5774	-0.0797223\\
1.5799	-0.0796148\\
1.5825	-0.0795023\\
1.585	-0.0793934\\
1.5876	-0.0792795\\
1.5902	-0.079165\\
1.5927	-0.0790542\\
1.5953	-0.0789384\\
1.5979	-0.0788221\\
1.6004	-0.0787096\\
1.603	-0.0785922\\
1.6056	-0.0784742\\
1.6081	-0.0783603\\
1.6107	-0.0782414\\
1.6132	-0.0781267\\
1.6158	-0.078007\\
1.6184	-0.0778869\\
1.6209	-0.0777711\\
1.6235	-0.0776503\\
1.6261	-0.0775292\\
1.6286	-0.0774124\\
1.6312	-0.0772908\\
1.6338	-0.0771688\\
1.6363	-0.0770514\\
1.6389	-0.076929\\
1.6414	-0.0768112\\
1.644	-0.0766885\\
1.6466	-0.0765657\\
1.6491	-0.0764475\\
1.6517	-0.0763244\\
1.6543	-0.0762013\\
1.6568	-0.0760828\\
1.6594	-0.0759596\\
1.662	-0.0758364\\
1.6645	-0.0757178\\
1.6671	-0.0755946\\
1.6696	-0.0754761\\
1.6722	-0.075353\\
1.6748	-0.0752299\\
1.6773	-0.0751116\\
1.6799	-0.0749887\\
1.6825	-0.0748659\\
1.685	-0.074748\\
1.6876	-0.0746255\\
1.6902	-0.0745031\\
1.6927	-0.0743857\\
1.6953	-0.0742638\\
1.6978	-0.0741467\\
1.7004	-0.0740252\\
1.703	-0.0739039\\
1.7055	-0.0737876\\
1.7081	-0.0736668\\
1.7107	-0.0735464\\
1.7132	-0.0734309\\
1.7158	-0.073311\\
1.7184	-0.0731915\\
1.7209	-0.0730769\\
1.7235	-0.0729581\\
1.7261	-0.0728396\\
1.7286	-0.0727261\\
1.7312	-0.0726084\\
1.7337	-0.0724956\\
1.7363	-0.0723786\\
1.7389	-0.0722621\\
1.7414	-0.0721505\\
1.744	-0.0720349\\
1.7466	-0.0719197\\
1.7491	-0.0718094\\
1.7517	-0.0716952\\
1.7543	-0.0715814\\
1.7568	-0.0714725\\
1.7594	-0.0713597\\
1.7619	-0.0712517\\
1.7645	-0.0711399\\
1.7671	-0.0710286\\
1.7696	-0.0709221\\
1.7722	-0.0708119\\
1.7748	-0.0707023\\
1.7773	-0.0705973\\
1.7799	-0.0704888\\
1.7825	-0.0703808\\
1.785	-0.0702775\\
1.7876	-0.0701706\\
1.7901	-0.0700684\\
1.7927	-0.0699627\\
1.7953	-0.0698576\\
1.7978	-0.0697571\\
1.8004	-0.0696531\\
1.803	-0.0695498\\
1.8055	-0.0694511\\
1.8081	-0.069349\\
1.8107	-0.0692475\\
1.8132	-0.0691506\\
1.8158	-0.0690503\\
1.8183	-0.0689546\\
1.8209	-0.0688556\\
1.8235	-0.0687573\\
1.826	-0.0686633\\
1.8286	-0.0685663\\
1.8312	-0.0684699\\
1.8337	-0.0683778\\
1.8363	-0.0682827\\
1.8389	-0.0681883\\
1.8414	-0.0680981\\
1.844	-0.068005\\
1.8465	-0.067916\\
1.8491	-0.0678242\\
1.8517	-0.067733\\
1.8542	-0.067646\\
1.8568	-0.0675562\\
1.8594	-0.067467\\
1.8619	-0.0673819\\
1.8645	-0.067294\\
1.8671	-0.0672068\\
1.8696	-0.0671236\\
1.8722	-0.0670378\\
1.8747	-0.0669559\\
1.8773	-0.0668713\\
1.8799	-0.0667875\\
1.8824	-0.0667075\\
1.885	-0.0666249\\
1.8876	-0.0665431\\
1.8901	-0.066465\\
1.8927	-0.0663844\\
1.8953	-0.0663045\\
1.8978	-0.0662284\\
1.9004	-0.0661498\\
1.903	-0.0660719\\
1.9055	-0.0659976\\
1.9081	-0.065921\\
1.9106	-0.065848\\
1.9132	-0.0657727\\
1.9158	-0.0656981\\
1.9183	-0.0656269\\
1.9209	-0.0655536\\
1.9235	-0.0654809\\
1.926	-0.0654116\\
1.9286	-0.0653402\\
1.9312	-0.0652694\\
1.9337	-0.065202\\
1.9363	-0.0651325\\
1.9388	-0.0650663\\
1.9414	-0.064998\\
1.944	-0.0649304\\
1.9465	-0.064866\\
1.9491	-0.0647996\\
1.9517	-0.0647339\\
1.9542	-0.0646713\\
1.9568	-0.0646068\\
1.9594	-0.0645429\\
1.9619	-0.064482\\
1.9645	-0.0644193\\
1.967	-0.0643596\\
1.9696	-0.0642982\\
1.9722	-0.0642373\\
1.9747	-0.0641793\\
1.9773	-0.0641196\\
1.9799	-0.0640605\\
1.9824	-0.0640043\\
1.985	-0.0639463\\
1.9876	-0.063889\\
1.9901	-0.0638344\\
1.9927	-0.0637782\\
1.9952	-0.0637247\\
1.9978	-0.0636697\\
2.0004	-0.0636152\\
2.0029	-0.0635633\\
2.0055	-0.06351\\
2.0081	-0.0634572\\
2.0106	-0.0634069\\
2.0132	-0.0633552\\
2.0158	-0.0633041\\
2.0183	-0.0632554\\
2.0209	-0.0632054\\
2.0234	-0.0631577\\
2.026	-0.0631087\\
2.0286	-0.0630603\\
2.0311	-0.0630142\\
2.0337	-0.0629668\\
2.0363	-0.0629199\\
2.0388	-0.0628753\\
2.0414	-0.0628295\\
2.044	-0.0627842\\
2.0465	-0.0627411\\
2.0491	-0.0626968\\
2.0517	-0.062653\\
2.0542	-0.0626114\\
2.0568	-0.0625686\\
2.0593	-0.062528\\
2.0619	-0.0624862\\
2.0645	-0.0624449\\
2.067	-0.0624057\\
2.0696	-0.0623653\\
2.0722	-0.0623255\\
2.0747	-0.0622877\\
2.0773	-0.0622488\\
2.0799	-0.0622104\\
2.0824	-0.062174\\
2.085	-0.0621365\\
2.0875	-0.062101\\
2.0901	-0.0620645\\
2.0927	-0.0620284\\
2.0952	-0.0619942\\
2.0978	-0.0619591\\
2.1004	-0.0619245\\
2.1029	-0.0618916\\
2.1055	-0.0618579\\
2.1081	-0.0618246\\
2.1106	-0.061793\\
2.1132	-0.0617606\\
2.1157	-0.0617299\\
2.1183	-0.0616984\\
2.1209	-0.0616673\\
2.1234	-0.0616379\\
2.126	-0.0616077\\
2.1286	-0.061578\\
2.1311	-0.0615498\\
2.1337	-0.0615209\\
2.1363	-0.0614925\\
2.1388	-0.0614655\\
2.1414	-0.0614379\\
2.1439	-0.0614118\\
2.1465	-0.0613851\\
2.1491	-0.0613588\\
2.1516	-0.0613339\\
2.1542	-0.0613084\\
2.1568	-0.0612833\\
2.1593	-0.0612596\\
2.1619	-0.0612354\\
2.1645	-0.0612116\\
2.167	-0.0611891\\
2.1696	-0.0611661\\
2.1721	-0.0611444\\
2.1747	-0.0611222\\
2.1773	-0.0611005\\
2.1798	-0.0610799\\
2.1824	-0.061059\\
2.185	-0.0610384\\
2.1875	-0.0610191\\
2.1901	-0.0609993\\
2.1927	-0.06098\\
2.1952	-0.0609618\\
2.1978	-0.0609432\\
2.2004	-0.0609251\\
2.2029	-0.060908\\
2.2055	-0.0608907\\
2.208	-0.0608743\\
2.2106	-0.0608578\\
2.2132	-0.0608416\\
2.2157	-0.0608264\\
2.2183	-0.0608111\\
2.2209	-0.0607961\\
2.2234	-0.060782\\
2.226	-0.0607678\\
2.2286	-0.060754\\
2.2311	-0.0607411\\
2.2337	-0.060728\\
2.2362	-0.0607158\\
2.2388	-0.0607035\\
2.2414	-0.0606916\\
2.2439	-0.0606806\\
2.2465	-0.0606694\\
2.2491	-0.0606587\\
2.2516	-0.0606487\\
2.2542	-0.0606388\\
2.2568	-0.0606292\\
2.2593	-0.0606203\\
2.2619	-0.0606115\\
2.2644	-0.0606033\\
2.267	-0.0605953\\
2.2696	-0.0605876\\
2.2721	-0.0605805\\
2.2747	-0.0605736\\
2.2773	-0.060567\\
2.2798	-0.0605611\\
2.2824	-0.0605552\\
2.285	-0.0605498\\
2.2875	-0.0605449\\
2.2901	-0.0605402\\
2.2926	-0.0605361\\
2.2952	-0.0605321\\
2.2978	-0.0605286\\
2.3003	-0.0605255\\
2.3029	-0.0605226\\
2.3055	-0.0605202\\
2.308	-0.0605182\\
2.3106	-0.0605165\\
2.3132	-0.0605151\\
2.3157	-0.0605142\\
2.3183	-0.0605135\\
2.3208	-0.0605133\\
2.3234	-0.0605134\\
2.326	-0.0605139\\
2.3285	-0.0605147\\
2.3311	-0.0605159\\
2.3337	-0.0605175\\
2.3362	-0.0605193\\
2.3388	-0.0605216\\
2.3414	-0.0605243\\
2.3439	-0.0605272\\
2.3465	-0.0605305\\
2.349	-0.0605341\\
2.3516	-0.0605382\\
2.3542	-0.0605426\\
2.3567	-0.0605473\\
2.3593	-0.0605524\\
2.3619	-0.0605579\\
2.3644	-0.0605635\\
2.367	-0.0605697\\
2.3696	-0.0605763\\
2.3721	-0.0605829\\
2.3747	-0.0605902\\
2.3773	-0.0605978\\
2.3798	-0.0606054\\
2.3824	-0.0606137\\
2.3849	-0.060622\\
2.3875	-0.060631\\
2.3901	-0.0606403\\
2.3926	-0.0606495\\
2.3952	-0.0606595\\
2.3978	-0.0606698\\
2.4003	-0.0606801\\
2.4029	-0.060691\\
2.4055	-0.0607024\\
2.408	-0.0607135\\
2.4106	-0.0607255\\
2.4131	-0.0607373\\
2.4157	-0.0607499\\
2.4183	-0.0607629\\
2.4208	-0.0607756\\
2.4234	-0.0607892\\
2.426	-0.0608031\\
2.4285	-0.0608167\\
2.4311	-0.0608312\\
2.4337	-0.0608461\\
2.4362	-0.0608606\\
2.4388	-0.0608761\\
2.4413	-0.0608912\\
2.4439	-0.0609072\\
2.4465	-0.0609236\\
2.449	-0.0609396\\
2.4516	-0.0609565\\
2.4542	-0.0609737\\
2.4567	-0.0609906\\
2.4593	-0.0610084\\
2.4619	-0.0610265\\
2.4644	-0.0610441\\
2.467	-0.0610628\\
2.4695	-0.061081\\
2.4721	-0.0611001\\
2.4747	-0.0611196\\
2.4772	-0.0611386\\
2.4798	-0.0611586\\
2.4824	-0.0611789\\
2.4849	-0.0611986\\
2.4875	-0.0612194\\
2.4901	-0.0612404\\
2.4926	-0.0612609\\
2.4952	-0.0612824\\
2.4977	-0.0613033\\
2.5003	-0.0613253\\
2.5029	-0.0613476\\
2.5054	-0.0613692\\
2.508	-0.0613919\\
2.5106	-0.0614149\\
2.5131	-0.0614372\\
2.5157	-0.0614606\\
2.5183	-0.0614842\\
2.5208	-0.0615071\\
2.5234	-0.0615311\\
2.526	-0.0615554\\
2.5285	-0.0615789\\
2.5311	-0.0616035\\
2.5336	-0.0616274\\
2.5362	-0.0616524\\
2.5388	-0.0616776\\
2.5413	-0.0617021\\
2.5439	-0.0617277\\
2.5465	-0.0617534\\
2.549	-0.0617783\\
2.5516	-0.0618044\\
2.5542	-0.0618307\\
2.5567	-0.0618561\\
2.5593	-0.0618827\\
2.5618	-0.0619084\\
2.5644	-0.0619353\\
2.567	-0.0619623\\
2.5695	-0.0619884\\
2.5721	-0.0620157\\
2.5747	-0.0620431\\
2.5772	-0.0620696\\
2.5798	-0.0620973\\
2.5824	-0.0621251\\
2.5849	-0.0621519\\
2.5875	-0.0621799\\
2.59	-0.0622069\\
2.5926	-0.0622352\\
2.5952	-0.0622635\\
2.5977	-0.0622908\\
2.6003	-0.0623193\\
2.6029	-0.0623478\\
2.6054	-0.0623754\\
2.608	-0.0624041\\
2.6106	-0.0624328\\
2.6131	-0.0624606\\
2.6157	-0.0624895\\
2.6182	-0.0625173\\
2.6208	-0.0625463\\
2.6234	-0.0625753\\
2.6259	-0.0626033\\
2.6285	-0.0626324\\
2.6311	-0.0626615\\
2.6336	-0.0626896\\
2.6362	-0.0627187\\
2.6388	-0.0627479\\
2.6413	-0.062776\\
2.6439	-0.0628052\\
2.6464	-0.0628333\\
2.649	-0.0628625\\
2.6516	-0.0628917\\
2.6541	-0.0629197\\
2.6567	-0.0629489\\
2.6593	-0.062978\\
2.6618	-0.0630059\\
2.6644	-0.063035\\
2.667	-0.063064\\
2.6695	-0.0630919\\
2.6721	-0.0631208\\
2.6746	-0.0631485\\
2.6772	-0.0631773\\
2.6798	-0.063206\\
2.6823	-0.0632336\\
2.6849	-0.0632622\\
2.6875	-0.0632907\\
2.69	-0.063318\\
2.6926	-0.0633463\\
2.6952	-0.0633745\\
2.6977	-0.0634016\\
2.7003	-0.0634296\\
2.7029	-0.0634575\\
2.7054	-0.0634842\\
2.708	-0.0635119\\
2.7105	-0.0635383\\
2.7131	-0.0635658\\
2.7157	-0.063593\\
2.7182	-0.0636191\\
2.7208	-0.0636461\\
2.7234	-0.0636729\\
2.7259	-0.0636986\\
2.7285	-0.0637251\\
2.7311	-0.0637514\\
2.7336	-0.0637766\\
2.7362	-0.0638026\\
2.7387	-0.0638275\\
2.7413	-0.0638531\\
2.7439	-0.0638786\\
2.7464	-0.0639029\\
2.749	-0.0639279\\
2.7516	-0.0639528\\
2.7541	-0.0639765\\
2.7567	-0.0640009\\
2.7593	-0.0640251\\
2.7618	-0.0640482\\
2.7644	-0.064072\\
2.7669	-0.0640946\\
2.7695	-0.0641179\\
2.7721	-0.064141\\
2.7746	-0.064163\\
2.7772	-0.0641855\\
2.7798	-0.0642079\\
2.7823	-0.0642291\\
2.7849	-0.0642509\\
2.7875	-0.0642724\\
2.79	-0.0642929\\
2.7926	-0.0643139\\
2.7951	-0.0643338\\
2.7977	-0.0643543\\
2.8003	-0.0643744\\
2.8028	-0.0643935\\
2.8054	-0.0644131\\
2.808	-0.0644324\\
2.8105	-0.0644507\\
2.8131	-0.0644693\\
2.8157	-0.0644877\\
2.8182	-0.0645051\\
2.8208	-0.0645228\\
2.8233	-0.0645396\\
2.8259	-0.0645567\\
2.8285	-0.0645734\\
2.831	-0.0645892\\
2.8336	-0.0646054\\
2.8362	-0.0646211\\
2.8387	-0.064636\\
2.8413	-0.0646511\\
2.8439	-0.0646658\\
2.8464	-0.0646796\\
2.849	-0.0646937\\
2.8516	-0.0647074\\
2.8541	-0.0647202\\
2.8567	-0.0647331\\
2.8592	-0.0647452\\
2.8618	-0.0647575\\
2.8644	-0.0647693\\
2.8669	-0.0647804\\
2.8695	-0.0647915\\
2.8721	-0.0648022\\
2.8746	-0.0648122\\
2.8772	-0.0648222\\
2.8798	-0.0648318\\
2.8823	-0.0648406\\
2.8849	-0.0648494\\
2.8874	-0.0648575\\
2.89	-0.0648655\\
2.8926	-0.0648731\\
2.8951	-0.0648801\\
2.8977	-0.0648869\\
2.9003	-0.0648933\\
2.9028	-0.0648991\\
2.9054	-0.0649047\\
2.908	-0.0649099\\
2.9105	-0.0649145\\
2.9131	-0.0649189\\
2.9156	-0.0649227\\
2.9182	-0.0649262\\
2.9208	-0.0649294\\
2.9233	-0.064932\\
2.9259	-0.0649343\\
2.9285	-0.0649361\\
2.931	-0.0649375\\
2.9336	-0.0649385\\
2.9362	-0.0649391\\
2.9387	-0.0649393\\
2.9413	-0.064939\\
2.9438	-0.0649384\\
2.9464	-0.0649373\\
2.949	-0.0649357\\
2.9515	-0.0649338\\
2.9541	-0.0649314\\
2.9567	-0.0649286\\
2.9592	-0.0649254\\
2.9618	-0.0649217\\
2.9644	-0.0649176\\
2.9669	-0.0649132\\
2.9695	-0.0649081\\
2.972	-0.0649029\\
2.9746	-0.064897\\
2.9772	-0.0648907\\
2.9797	-0.0648842\\
2.9823	-0.0648771\\
2.9849	-0.0648694\\
2.9874	-0.0648617\\
2.99	-0.0648532\\
2.9926	-0.0648443\\
2.9951	-0.0648352\\
2.9977	-0.0648254\\
3.0003	-0.0648152\\
3.0028	-0.0648049\\
3.0054	-0.0647938\\
3.0079	-0.0647828\\
3.0105	-0.0647708\\
3.0131	-0.0647584\\
3.0156	-0.0647461\\
3.0182	-0.0647328\\
3.0208	-0.0647191\\
3.0233	-0.0647055\\
3.0259	-0.064691\\
3.0285	-0.064676\\
3.031	-0.0646612\\
3.0336	-0.0646454\\
3.0361	-0.0646297\\
3.0387	-0.0646131\\
3.0413	-0.064596\\
3.0438	-0.0645791\\
3.0464	-0.0645612\\
3.049	-0.0645429\\
3.0515	-0.0645248\\
3.0541	-0.0645056\\
3.0567	-0.0644861\\
3.0592	-0.0644668\\
3.0618	-0.0644464\\
3.0643	-0.0644264\\
3.0669	-0.0644052\\
3.0695	-0.0643835\\
3.072	-0.0643623\\
3.0746	-0.0643399\\
3.0772	-0.0643171\\
3.0797	-0.0642947\\
3.0823	-0.0642711\\
3.0849	-0.0642471\\
3.0874	-0.0642236\\
3.09	-0.0641988\\
3.0925	-0.0641746\\
3.0951	-0.0641491\\
3.0977	-0.0641231\\
3.1002	-0.0640978\\
3.1028	-0.0640711\\
3.1054	-0.064044\\
3.1079	-0.0640176\\
3.1105	-0.0639898\\
3.1131	-0.0639616\\
3.1156	-0.0639342\\
3.1182	-0.0639053\\
3.1207	-0.0638771\\
3.1233	-0.0638475\\
3.1259	-0.0638175\\
3.1284	-0.0637884\\
3.131	-0.0637577\\
3.1336	-0.0637267\\
3.1361	-0.0636965\\
3.1387	-0.0636648\\
3.1413	-0.0636328\\
3.1438	-0.0636016\\
3.1464	-0.0635689\\
3.1489	-0.0635372\\
3.1515	-0.0635038\\
3.1541	-0.0634701\\
3.1566	-0.0634375\\
3.1592	-0.0634032\\
3.1618	-0.0633685\\
3.1643	-0.0633349\\
3.1669	-0.0632997\\
3.1695	-0.0632642\\
3.172	-0.0632297\\
3.1746	-0.0631936\\
3.1772	-0.0631572\\
3.1797	-0.0631219\\
3.1823	-0.0630849\\
3.1848	-0.063049\\
3.1874	-0.0630115\\
3.19	-0.0629736\\
3.1925	-0.062937\\
3.1951	-0.0628986\\
3.1977	-0.06286\\
3.2002	-0.0628226\\
3.2028	-0.0627835\\
3.2054	-0.0627441\\
3.2079	-0.0627059\\
3.2105	-0.062666\\
3.213	-0.0626274\\
3.2156	-0.0625871\\
3.2182	-0.0625464\\
3.2207	-0.0625072\\
3.2233	-0.0624661\\
3.2259	-0.0624248\\
3.2284	-0.0623849\\
3.231	-0.0623431\\
3.2336	-0.0623012\\
3.2361	-0.0622606\\
3.2387	-0.0622183\\
3.2412	-0.0621773\\
3.2438	-0.0621346\\
3.2464	-0.0620916\\
3.2489	-0.0620501\\
3.2515	-0.0620068\\
3.2541	-0.0619633\\
3.2566	-0.0619213\\
3.2592	-0.0618774\\
3.2618	-0.0618334\\
3.2643	-0.0617909\\
3.2669	-0.0617465\\
3.2694	-0.0617037\\
3.272	-0.061659\\
3.2746	-0.0616141\\
3.2771	-0.0615709\\
3.2797	-0.0615258\\
3.2823	-0.0614805\\
3.2848	-0.0614368\\
3.2874	-0.0613913\\
3.29	-0.0613456\\
3.2925	-0.0613016\\
3.2951	-0.0612557\\
3.2976	-0.0612114\\
3.3002	-0.0611653\\
3.3028	-0.061119\\
3.3053	-0.0610745\\
3.3079	-0.061028\\
3.3105	-0.0609815\\
3.313	-0.0609366\\
3.3156	-0.0608899\\
3.3182	-0.060843\\
3.3207	-0.0607979\\
3.3233	-0.060751\\
3.3259	-0.0607039\\
3.3284	-0.0606586\\
3.331	-0.0606114\\
3.3335	-0.0605659\\
3.3361	-0.0605186\\
3.3387	-0.0604712\\
3.3412	-0.0604256\\
3.3438	-0.0603782\\
3.3464	-0.0603306\\
3.3489	-0.0602849\\
3.3515	-0.0602373\\
3.3541	-0.0601897\\
3.3566	-0.0601438\\
3.3592	-0.0600962\\
3.3617	-0.0600503\\
3.3643	-0.0600026\\
3.3669	-0.0599548\\
3.3694	-0.0599089\\
3.372	-0.0598611\\
3.3746	-0.0598134\\
3.3771	-0.0597674\\
3.3797	-0.0597197\\
3.3823	-0.0596719\\
3.3848	-0.059626\\
3.3874	-0.0595783\\
3.3899	-0.0595324\\
3.3925	-0.0594847\\
3.3951	-0.059437\\
3.3976	-0.0593912\\
3.4002	-0.0593436\\
3.4028	-0.059296\\
3.4053	-0.0592503\\
3.4079	-0.0592028\\
3.4105	-0.0591553\\
3.413	-0.0591097\\
3.4156	-0.0590624\\
3.4181	-0.0590169\\
3.4207	-0.0589696\\
3.4233	-0.0589225\\
3.4258	-0.0588771\\
3.4284	-0.0588301\\
3.431	-0.0587831\\
3.4335	-0.058738\\
3.4361	-0.0586912\\
3.4387	-0.0586444\\
3.4412	-0.0585995\\
3.4438	-0.0585529\\
3.4463	-0.0585082\\
3.4489	-0.0584618\\
3.4515	-0.0584155\\
3.454	-0.0583711\\
3.4566	-0.058325\\
3.4592	-0.058279\\
3.4617	-0.0582348\\
3.4643	-0.058189\\
3.4669	-0.0581433\\
3.4694	-0.0580995\\
3.472	-0.0580541\\
3.4745	-0.0580105\\
3.4771	-0.0579653\\
3.4797	-0.0579202\\
3.4822	-0.057877\\
3.4848	-0.0578322\\
3.4874	-0.0577875\\
3.4899	-0.0577447\\
3.4925	-0.0577003\\
3.4951	-0.057656\\
3.4976	-0.0576136\\
3.5002	-0.0575696\\
3.5028	-0.0575258\\
3.5053	-0.0574838\\
3.5079	-0.0574403\\
3.5104	-0.0573987\\
3.513	-0.0573555\\
3.5156	-0.0573124\\
3.5181	-0.0572712\\
3.5207	-0.0572285\\
3.5233	-0.057186\\
3.5258	-0.0571453\\
3.5284	-0.0571031\\
3.531	-0.0570611\\
3.5335	-0.0570209\\
3.5361	-0.0569792\\
3.5386	-0.0569393\\
3.5412	-0.056898\\
3.5438	-0.056857\\
3.5463	-0.0568176\\
3.5489	-0.0567769\\
3.5515	-0.0567364\\
3.554	-0.0566976\\
3.5566	-0.0566575\\
3.5592	-0.0566176\\
3.5617	-0.0565794\\
3.5643	-0.0565398\\
3.5668	-0.056502\\
3.5694	-0.0564629\\
3.572	-0.056424\\
3.5745	-0.0563868\\
3.5771	-0.0563483\\
3.5797	-0.05631\\
3.5822	-0.0562734\\
3.5848	-0.0562356\\
3.5874	-0.056198\\
3.5899	-0.056162\\
3.5925	-0.0561248\\
3.595	-0.0560893\\
3.5976	-0.0560526\\
3.6002	-0.0560161\\
3.6027	-0.0559812\\
3.6053	-0.0559451\\
3.6079	-0.0559093\\
3.6104	-0.0558751\\
3.613	-0.0558398\\
3.6156	-0.0558047\\
3.6181	-0.0557711\\
3.6207	-0.0557365\\
3.6232	-0.0557034\\
3.6258	-0.0556693\\
3.6284	-0.0556353\\
3.6309	-0.055603\\
3.6335	-0.0555695\\
3.6361	-0.0555364\\
3.6386	-0.0555047\\
3.6412	-0.055472\\
3.6438	-0.0554396\\
3.6463	-0.0554086\\
3.6489	-0.0553767\\
3.6515	-0.055345\\
3.654	-0.0553148\\
3.6566	-0.0552836\\
3.6591	-0.0552539\\
3.6617	-0.0552232\\
3.6643	-0.0551928\\
3.6668	-0.0551638\\
3.6694	-0.0551339\\
3.672	-0.0551043\\
3.6745	-0.0550761\\
3.6771	-0.055047\\
3.6797	-0.0550181\\
3.6822	-0.0549906\\
3.6848	-0.0549623\\
3.6873	-0.0549353\\
3.6899	-0.0549076\\
3.6925	-0.05488\\
3.695	-0.0548538\\
3.6976	-0.0548269\\
3.7002	-0.0548001\\
3.7027	-0.0547747\\
3.7053	-0.0547485\\
3.7079	-0.0547226\\
3.7104	-0.054698\\
3.713	-0.0546726\\
3.7155	-0.0546485\\
3.7181	-0.0546236\\
3.7207	-0.0545991\\
3.7232	-0.0545757\\
3.7258	-0.0545517\\
3.7284	-0.054528\\
3.7309	-0.0545055\\
3.7335	-0.0544823\\
3.7361	-0.0544594\\
3.7386	-0.0544376\\
3.7412	-0.0544153\\
3.7437	-0.054394\\
3.7463	-0.0543722\\
3.7489	-0.0543507\\
3.7514	-0.0543303\\
3.754	-0.0543093\\
3.7566	-0.0542886\\
3.7591	-0.054269\\
3.7617	-0.0542488\\
3.7643	-0.054229\\
3.7668	-0.0542101\\
3.7694	-0.0541908\\
3.7719	-0.0541725\\
3.7745	-0.0541538\\
3.7771	-0.0541353\\
3.7796	-0.0541178\\
3.7822	-0.0540999\\
3.7848	-0.0540822\\
3.7873	-0.0540656\\
3.7899	-0.0540485\\
3.7925	-0.0540317\\
3.795	-0.0540158\\
3.7976	-0.0539995\\
3.8002	-0.0539836\\
3.8027	-0.0539685\\
3.8053	-0.0539531\\
3.8078	-0.0539385\\
3.8104	-0.0539237\\
3.813	-0.0539091\\
3.8155	-0.0538953\\
3.8181	-0.0538813\\
3.8207	-0.0538676\\
3.8232	-0.0538546\\
3.8258	-0.0538414\\
3.8284	-0.0538285\\
3.8309	-0.0538164\\
3.8335	-0.053804\\
3.836	-0.0537924\\
3.8386	-0.0537806\\
3.8412	-0.053769\\
3.8437	-0.0537582\\
3.8463	-0.0537472\\
3.8489	-0.0537365\\
3.8514	-0.0537265\\
3.854	-0.0537163\\
3.8566	-0.0537064\\
3.8591	-0.0536972\\
3.8617	-0.0536878\\
3.8642	-0.0536791\\
3.8668	-0.0536703\\
3.8694	-0.0536617\\
3.8719	-0.0536538\\
3.8745	-0.0536458\\
3.8771	-0.0536381\\
3.8796	-0.0536309\\
3.8822	-0.0536237\\
3.8848	-0.0536168\\
3.8873	-0.0536104\\
3.8899	-0.053604\\
3.8924	-0.0535981\\
3.895	-0.0535922\\
3.8976	-0.0535867\\
3.9001	-0.0535815\\
3.9027	-0.0535765\\
3.9053	-0.0535717\\
3.9078	-0.0535673\\
3.9104	-0.053563\\
3.913	-0.053559\\
3.9155	-0.0535554\\
3.9181	-0.0535519\\
3.9206	-0.0535488\\
3.9232	-0.0535458\\
3.9258	-0.0535431\\
3.9283	-0.0535407\\
3.9309	-0.0535385\\
3.9335	-0.0535366\\
3.936	-0.053535\\
3.9386	-0.0535335\\
3.9412	-0.0535323\\
3.9437	-0.0535314\\
3.9463	-0.0535307\\
3.9488	-0.0535303\\
3.9514	-0.0535301\\
3.954	-0.0535302\\
3.9565	-0.0535305\\
3.9591	-0.0535311\\
3.9617	-0.0535319\\
3.9642	-0.0535329\\
3.9668	-0.0535342\\
3.9694	-0.0535357\\
3.9719	-0.0535375\\
3.9745	-0.0535395\\
3.9771	-0.0535418\\
3.9796	-0.0535442\\
3.9822	-0.0535469\\
3.9847	-0.0535498\\
3.9873	-0.053553\\
3.9899	-0.0535565\\
3.9924	-0.0535601\\
3.995	-0.053564\\
3.9976	-0.0535682\\
4.0001	-0.0535724\\
4.0027	-0.053577\\
4.0053	-0.0535819\\
4.0078	-0.0535868\\
4.0104	-0.0535921\\
4.0129	-0.0535975\\
4.0155	-0.0536032\\
4.0181	-0.0536093\\
4.0206	-0.0536152\\
4.0232	-0.0536217\\
4.0258	-0.0536284\\
4.0283	-0.053635\\
4.0309	-0.0536421\\
4.0335	-0.0536494\\
4.036	-0.0536567\\
4.0386	-0.0536645\\
4.0411	-0.0536722\\
4.0437	-0.0536804\\
4.0463	-0.0536888\\
4.0488	-0.0536971\\
4.0514	-0.0537059\\
4.054	-0.0537149\\
4.0565	-0.0537238\\
4.0591	-0.0537333\\
4.0617	-0.053743\\
4.0642	-0.0537525\\
4.0668	-0.0537625\\
4.0693	-0.0537724\\
4.0719	-0.0537829\\
4.0745	-0.0537936\\
4.077	-0.053804\\
4.0796	-0.0538151\\
4.0822	-0.0538264\\
4.0847	-0.0538374\\
4.0873	-0.0538491\\
4.0899	-0.0538609\\
4.0924	-0.0538725\\
4.095	-0.0538847\\
4.0975	-0.0538967\\
4.1001	-0.0539093\\
4.1027	-0.0539221\\
4.1052	-0.0539346\\
4.1078	-0.0539477\\
4.1104	-0.0539611\\
4.1129	-0.0539741\\
4.1155	-0.0539878\\
4.1181	-0.0540017\\
4.1206	-0.0540152\\
4.1232	-0.0540294\\
4.1258	-0.0540438\\
4.1283	-0.0540579\\
4.1309	-0.0540726\\
4.1334	-0.0540869\\
4.136	-0.054102\\
4.1386	-0.0541173\\
4.1411	-0.0541321\\
4.1437	-0.0541477\\
4.1463	-0.0541634\\
4.1488	-0.0541787\\
4.1514	-0.0541948\\
4.154	-0.054211\\
4.1565	-0.0542268\\
4.1591	-0.0542433\\
4.1616	-0.0542593\\
4.1642	-0.0542762\\
4.1668	-0.0542932\\
4.1693	-0.0543096\\
4.1719	-0.0543269\\
4.1745	-0.0543444\\
4.177	-0.0543612\\
4.1796	-0.054379\\
4.1822	-0.0543968\\
4.1847	-0.0544141\\
4.1873	-0.0544323\\
4.1898	-0.0544498\\
4.1924	-0.0544682\\
4.195	-0.0544868\\
4.1975	-0.0545047\\
4.2001	-0.0545235\\
4.2027	-0.0545424\\
4.2052	-0.0545608\\
4.2078	-0.0545799\\
4.2104	-0.0545992\\
4.2129	-0.0546179\\
4.2155	-0.0546375\\
4.218	-0.0546564\\
4.2206	-0.0546761\\
4.2232	-0.054696\\
4.2257	-0.0547153\\
4.2283	-0.0547354\\
4.2309	-0.0547556\\
4.2334	-0.0547752\\
4.236	-0.0547956\\
4.2386	-0.0548162\\
4.2411	-0.054836\\
4.2437	-0.0548568\\
4.2462	-0.0548769\\
4.2488	-0.0548978\\
4.2514	-0.0549189\\
4.2539	-0.0549392\\
4.2565	-0.0549604\\
4.2591	-0.0549818\\
4.2616	-0.0550024\\
4.2642	-0.0550239\\
4.2668	-0.0550455\\
4.2693	-0.0550663\\
4.2719	-0.0550881\\
4.2744	-0.0551091\\
4.277	-0.055131\\
4.2796	-0.055153\\
4.2821	-0.0551743\\
4.2847	-0.0551964\\
4.2873	-0.0552186\\
4.2898	-0.0552401\\
4.2924	-0.0552624\\
4.295	-0.0552849\\
4.2975	-0.0553065\\
4.3001	-0.0553291\\
4.3027	-0.0553517\\
4.3052	-0.0553735\\
4.3078	-0.0553963\\
4.3103	-0.0554182\\
4.3129	-0.055441\\
4.3155	-0.055464\\
4.318	-0.055486\\
4.3206	-0.055509\\
4.3232	-0.0555321\\
4.3257	-0.0555543\\
4.3283	-0.0555775\\
4.3309	-0.0556006\\
4.3334	-0.055623\\
4.336	-0.0556463\\
4.3385	-0.0556687\\
4.3411	-0.055692\\
4.3437	-0.0557154\\
4.3462	-0.0557379\\
4.3488	-0.0557613\\
4.3514	-0.0557848\\
4.3539	-0.0558074\\
4.3565	-0.0558309\\
4.3591	-0.0558544\\
4.3616	-0.0558771\\
4.3642	-0.0559006\\
4.3667	-0.0559233\\
4.3693	-0.0559469\\
4.3719	-0.0559705\\
4.3744	-0.0559933\\
4.377	-0.0560169\\
4.3796	-0.0560406\\
4.3821	-0.0560633\\
4.3847	-0.056087\\
4.3873	-0.0561106\\
4.3898	-0.0561334\\
4.3924	-0.056157\\
4.3949	-0.0561798\\
4.3975	-0.0562034\\
4.4001	-0.0562271\\
4.4026	-0.0562498\\
4.4052	-0.0562734\\
4.4078	-0.056297\\
4.4103	-0.0563197\\
4.4129	-0.0563433\\
4.4155	-0.0563669\\
4.418	-0.0563895\\
4.4206	-0.0564131\\
4.4231	-0.0564357\\
4.4257	-0.0564591\\
4.4283	-0.0564826\\
4.4308	-0.0565051\\
4.4334	-0.0565285\\
4.436	-0.0565519\\
4.4385	-0.0565743\\
4.4411	-0.0565976\\
4.4437	-0.0566209\\
4.4462	-0.0566432\\
4.4488	-0.0566664\\
4.4514	-0.0566895\\
4.4539	-0.0567117\\
4.4565	-0.0567348\\
4.459	-0.0567569\\
4.4616	-0.0567798\\
4.4642	-0.0568027\\
4.4667	-0.0568247\\
4.4693	-0.0568475\\
4.4719	-0.0568702\\
4.4744	-0.056892\\
4.477	-0.0569146\\
4.4796	-0.0569372\\
4.4821	-0.0569588\\
4.4847	-0.0569812\\
4.4872	-0.0570027\\
4.4898	-0.057025\\
4.4924	-0.0570473\\
4.4949	-0.0570686\\
4.4975	-0.0570907\\
4.5001	-0.0571127\\
4.5026	-0.0571337\\
4.5052	-0.0571556\\
4.5078	-0.0571774\\
4.5103	-0.0571982\\
4.5129	-0.0572198\\
4.5154	-0.0572405\\
4.518	-0.057262\\
4.5206	-0.0572833\\
4.5231	-0.0573038\\
4.5257	-0.0573249\\
4.5283	-0.057346\\
4.5308	-0.0573662\\
4.5334	-0.0573871\\
4.536	-0.0574079\\
4.5385	-0.0574278\\
4.5411	-0.0574484\\
4.5436	-0.0574681\\
4.5462	-0.0574884\\
4.5488	-0.0575087\\
4.5513	-0.0575281\\
4.5539	-0.0575482\\
4.5565	-0.0575682\\
4.559	-0.0575873\\
4.5616	-0.057607\\
4.5642	-0.0576266\\
4.5667	-0.0576454\\
4.5693	-0.0576648\\
4.5718	-0.0576833\\
4.5744	-0.0577025\\
4.577	-0.0577215\\
4.5795	-0.0577397\\
4.5821	-0.0577585\\
4.5847	-0.0577772\\
4.5872	-0.057795\\
4.5898	-0.0578134\\
4.5924	-0.0578317\\
4.5949	-0.0578492\\
4.5975	-0.0578672\\
4.6001	-0.0578851\\
4.6026	-0.0579022\\
4.6052	-0.0579198\\
4.6077	-0.0579366\\
4.6103	-0.057954\\
4.6129	-0.0579712\\
4.6154	-0.0579876\\
4.618	-0.0580045\\
4.6206	-0.0580213\\
4.6231	-0.0580373\\
4.6257	-0.0580539\\
4.6283	-0.0580702\\
4.6308	-0.0580858\\
4.6334	-0.0581019\\
4.6359	-0.0581172\\
4.6385	-0.0581329\\
4.6411	-0.0581486\\
4.6436	-0.0581634\\
4.6462	-0.0581787\\
4.6488	-0.0581939\\
4.6513	-0.0582083\\
4.6539	-0.0582232\\
4.6565	-0.0582379\\
4.659	-0.0582518\\
4.6616	-0.0582662\\
4.6641	-0.0582799\\
4.6667	-0.0582939\\
4.6693	-0.0583078\\
4.6718	-0.058321\\
4.6744	-0.0583346\\
4.677	-0.0583479\\
4.6795	-0.0583607\\
4.6821	-0.0583737\\
4.6847	-0.0583866\\
4.6872	-0.0583989\\
4.6898	-0.0584114\\
4.6923	-0.0584233\\
4.6949	-0.0584356\\
4.6975	-0.0584476\\
4.7	-0.058459\\
4.7026	-0.0584707\\
4.7052	-0.0584822\\
4.7077	-0.0584932\\
4.7103	-0.0585043\\
4.7129	-0.0585153\\
4.7154	-0.0585258\\
4.718	-0.0585364\\
4.7205	-0.0585465\\
4.7231	-0.0585568\\
4.7257	-0.0585669\\
4.7282	-0.0585764\\
4.7308	-0.0585862\\
4.7334	-0.0585958\\
4.7359	-0.0586048\\
4.7385	-0.058614\\
4.7411	-0.0586231\\
4.7436	-0.0586316\\
4.7462	-0.0586402\\
4.7487	-0.0586484\\
4.7513	-0.0586567\\
4.7539	-0.0586648\\
4.7564	-0.0586724\\
4.759	-0.0586802\\
4.7616	-0.0586877\\
4.7641	-0.0586948\\
4.7667	-0.058702\\
4.7693	-0.058709\\
4.7718	-0.0587156\\
4.7744	-0.0587222\\
4.777	-0.0587286\\
4.7795	-0.0587346\\
4.7821	-0.0587407\\
4.7846	-0.0587463\\
4.7872	-0.058752\\
4.7898	-0.0587575\\
4.7923	-0.0587626\\
4.7949	-0.0587678\\
4.7975	-0.0587727\\
4.8	-0.0587772\\
4.8026	-0.0587818\\
4.8052	-0.0587862\\
4.8077	-0.0587902\\
4.8103	-0.0587942\\
4.8128	-0.0587978\\
4.8154	-0.0588014\\
4.818	-0.0588048\\
4.8205	-0.0588079\\
4.8231	-0.058811\\
4.8257	-0.0588138\\
4.8282	-0.0588164\\
4.8308	-0.0588188\\
4.8334	-0.0588211\\
4.8359	-0.0588231\\
4.8385	-0.058825\\
4.841	-0.0588266\\
4.8436	-0.0588282\\
4.8462	-0.0588295\\
4.8487	-0.0588306\\
4.8513	-0.0588315\\
4.8539	-0.0588323\\
4.8564	-0.0588329\\
4.859	-0.0588332\\
4.8616	-0.0588334\\
4.8641	-0.0588334\\
4.8667	-0.0588333\\
4.8692	-0.0588329\\
4.8718	-0.0588324\\
4.8744	-0.0588316\\
4.8769	-0.0588307\\
4.8795	-0.0588296\\
4.8821	-0.0588283\\
4.8846	-0.0588269\\
4.8872	-0.0588253\\
4.8898	-0.0588234\\
4.8923	-0.0588215\\
4.8949	-0.0588192\\
4.8974	-0.0588169\\
4.9	-0.0588144\\
4.9026	-0.0588116\\
4.9051	-0.0588088\\
4.9077	-0.0588056\\
4.9103	-0.0588023\\
4.9128	-0.058799\\
4.9154	-0.0587953\\
4.918	-0.0587915\\
4.9205	-0.0587876\\
4.9231	-0.0587834\\
4.9257	-0.058779\\
4.9282	-0.0587746\\
4.9308	-0.0587699\\
4.9333	-0.0587652\\
4.9359	-0.0587601\\
4.9385	-0.0587549\\
4.941	-0.0587496\\
4.9436	-0.0587441\\
4.9462	-0.0587383\\
4.9487	-0.0587326\\
4.9513	-0.0587265\\
4.9539	-0.0587202\\
4.9564	-0.058714\\
4.959	-0.0587073\\
4.9615	-0.0587008\\
4.9641	-0.0586939\\
4.9667	-0.0586867\\
4.9692	-0.0586797\\
4.9718	-0.0586723\\
4.9744	-0.0586647\\
4.9769	-0.0586572\\
4.9795	-0.0586493\\
4.9821	-0.0586412\\
4.9846	-0.0586332\\
4.9872	-0.0586248\\
4.9897	-0.0586166\\
4.9923	-0.0586078\\
4.9949	-0.0585989\\
4.9974	-0.0585903\\
5	-0.0585811\\
};
\addplot [color=mycolor15, line width=1.0pt, forget plot]
  table[row sep=crcr]{%
-0.125	5.32632e-08\\
-0.1224	5.43654e-08\\
-0.1199	5.54251e-08\\
-0.1173	5.65271e-08\\
-0.1147	5.76289e-08\\
-0.1122	5.86882e-08\\
-0.1096	5.97899e-08\\
-0.1071	6.0849e-08\\
-0.1045	6.19506e-08\\
-0.1019	6.30519e-08\\
-0.0994	6.41108e-08\\
-0.0968	6.5212e-08\\
-0.0942	6.63131e-08\\
-0.0917	6.73718e-08\\
-0.0891	6.84727e-08\\
-0.0865	6.95736e-08\\
-0.084	7.0632e-08\\
-0.0814	7.17326e-08\\
-0.0789	7.27908e-08\\
-0.0763	7.38912e-08\\
-0.0737	7.49916e-08\\
-0.0712	7.60494e-08\\
-0.0686	7.71496e-08\\
-0.066	7.82496e-08\\
-0.0635	7.93072e-08\\
-0.0609	8.04071e-08\\
-0.0583	8.15068e-08\\
-0.0558	8.25642e-08\\
-0.0532	8.36638e-08\\
-0.0507	8.4721e-08\\
-0.0481	8.58204e-08\\
-0.0455	8.69197e-08\\
-0.043	8.79766e-08\\
-0.0404	8.90757e-08\\
-0.0378	9.01747e-08\\
-0.0353	9.12313e-08\\
-0.0327	9.233e-08\\
-0.0301	9.34287e-08\\
-0.0276	9.4485e-08\\
-0.025	9.55835e-08\\
-0.0224	9.66819e-08\\
-0.0199	9.77379e-08\\
-0.0173	9.88361e-08\\
-0.0148	9.98921e-08\\
-0.0122	1.0099e-07\\
-0.0096	1.02088e-07\\
-0.0071	1.03143e-07\\
-0.0045	1.04241e-07\\
-0.0019	1.05339e-07\\
0.0006	1.06394e-07\\
0.0032	0.000715124\\
0.0058	0.00318294\\
0.0083	0.00596259\\
0.0109	0.00920475\\
0.0134	0.0124912\\
0.016	0.0158637\\
0.0186	0.0190614\\
0.0211	0.021971\\
0.0237	0.0248933\\
0.0263	0.0277667\\
0.0288	0.0304862\\
0.0314	0.0332331\\
0.034	0.0358603\\
0.0365	0.0382659\\
0.0391	0.0406618\\
0.0416	0.0428884\\
0.0442	0.0451367\\
0.0468	0.0473109\\
0.0493	0.0493186\\
0.0519	0.0513132\\
0.0545	0.053218\\
0.057	0.0549776\\
0.0596	0.0567442\\
0.0622	0.0584487\\
0.0647	0.060024\\
0.0673	0.0615905\\
0.0698	0.0630284\\
0.0724	0.0644606\\
0.075	0.0658386\\
0.0775	0.0671176\\
0.0801	0.0683985\\
0.0827	0.0696245\\
0.0852	0.0707492\\
0.0878	0.0718667\\
0.0904	0.0729391\\
0.0929	0.0739342\\
0.0955	0.0749329\\
0.098	0.0758546\\
0.1006	0.0767692\\
0.1032	0.0776391\\
0.1057	0.0784387\\
0.1083	0.0792387\\
0.1109	0.0800097\\
0.1134	0.0807219\\
0.116	0.0814275\\
0.1186	0.0820947\\
0.1211	0.0827024\\
0.1237	0.083305\\
0.1263	0.0838821\\
0.1288	0.084413\\
0.1314	0.084936\\
0.1339	0.085407\\
0.1365	0.0858631\\
0.1391	0.0862888\\
0.1416	0.0866741\\
0.1442	0.0870518\\
0.1468	0.0874036\\
0.1493	0.0877129\\
0.1519	0.088002\\
0.1545	0.0882595\\
0.157	0.0884814\\
0.1596	0.0886889\\
0.1621	0.0888658\\
0.1647	0.0890223\\
0.1673	0.089147\\
0.1698	0.0892363\\
0.1724	0.0893006\\
0.175	0.0893398\\
0.1775	0.0893555\\
0.1801	0.0893468\\
0.1827	0.0893085\\
0.1852	0.0892418\\
0.1878	0.0891429\\
0.1903	0.0890231\\
0.1929	0.0888759\\
0.1955	0.0887056\\
0.198	0.0885171\\
0.2006	0.0882924\\
0.2032	0.0880384\\
0.2057	0.0877695\\
0.2083	0.0874679\\
0.2109	0.0871457\\
0.2134	0.086815\\
0.216	0.0864466\\
0.2185	0.0860673\\
0.2211	0.0856486\\
0.2237	0.085209\\
0.2262	0.0847695\\
0.2288	0.0842953\\
0.2314	0.0838014\\
0.2339	0.0833059\\
0.2365	0.0827695\\
0.2391	0.0822147\\
0.2416	0.0816674\\
0.2442	0.0810856\\
0.2467	0.0805131\\
0.2493	0.0799021\\
0.2519	0.0792743\\
0.2544	0.0786565\\
0.257	0.0780028\\
0.2596	0.0773406\\
0.2621	0.0766963\\
0.2647	0.0760167\\
0.2673	0.0753259\\
0.2698	0.0746515\\
0.2724	0.0739424\\
0.2749	0.0732561\\
0.2775	0.0725394\\
0.2801	0.071819\\
0.2826	0.0711212\\
0.2852	0.0703894\\
0.2878	0.0696531\\
0.2903	0.0689438\\
0.2929	0.068207\\
0.2955	0.0674714\\
0.298	0.066764\\
0.3006	0.0660268\\
0.3032	0.0652883\\
0.3057	0.0645792\\
0.3083	0.0638454\\
0.3108	0.0631445\\
0.3134	0.0624195\\
0.316	0.0616971\\
0.3185	0.061004\\
0.3211	0.0602861\\
0.3237	0.0595735\\
0.3262	0.058895\\
0.3288	0.0581963\\
0.3314	0.0575033\\
0.3339	0.0568409\\
0.3365	0.0561562\\
0.339	0.0555034\\
0.3416	0.0548323\\
0.3442	0.05417\\
0.3467	0.0535404\\
0.3493	0.0528917\\
0.3519	0.0522483\\
0.3544	0.0516353\\
0.357	0.0510053\\
0.3596	0.0503844\\
0.3621	0.0497957\\
0.3647	0.0491907\\
0.3672	0.0486145\\
0.3698	0.0480209\\
0.3724	0.0474342\\
0.3749	0.046878\\
0.3775	0.0463083\\
0.3801	0.0457464\\
0.3826	0.0452121\\
0.3852	0.0446617\\
0.3878	0.0441172\\
0.3903	0.0436003\\
0.3929	0.0430707\\
0.3954	0.0425689\\
0.398	0.0420533\\
0.4006	0.0415428\\
0.4031	0.0410564\\
0.4057	0.0405561\\
0.4083	0.0400628\\
0.4108	0.0395951\\
0.4134	0.0391149\\
0.416	0.0386395\\
0.4185	0.0381862\\
0.4211	0.037719\\
0.4236	0.0372748\\
0.4262	0.036819\\
0.4288	0.036369\\
0.4313	0.0359407\\
0.4339	0.0354987\\
0.4365	0.03506\\
0.439	0.0346419\\
0.4416	0.034212\\
0.4442	0.0337873\\
0.4467	0.0333832\\
0.4493	0.0329662\\
0.4519	0.0325517\\
0.4544	0.0321558\\
0.457	0.0317479\\
0.4595	0.0313598\\
0.4621	0.0309604\\
0.4647	0.0305641\\
0.4672	0.0301853\\
0.4698	0.0297933\\
0.4724	0.0294041\\
0.4749	0.0290332\\
0.4775	0.0286512\\
0.4801	0.0282725\\
0.4826	0.0279104\\
0.4852	0.0275354\\
0.4877	0.0271766\\
0.4903	0.0268062\\
0.4929	0.026439\\
0.4954	0.026089\\
0.498	0.0257272\\
0.5006	0.0253668\\
0.5031	0.0250217\\
0.5057	0.0246646\\
0.5083	0.0243102\\
0.5108	0.0239721\\
0.5134	0.0236229\\
0.5159	0.0232885\\
0.5185	0.0229418\\
0.5211	0.0225964\\
0.5236	0.0222661\\
0.5262	0.021925\\
0.5288	0.0215862\\
0.5313	0.021262\\
0.5339	0.0209257\\
0.5365	0.0205902\\
0.539	0.0202688\\
0.5416	0.0199364\\
0.5441	0.0196187\\
0.5467	0.0192898\\
0.5493	0.0189619\\
0.5518	0.018647\\
0.5544	0.0183201\\
0.557	0.0179944\\
0.5595	0.0176827\\
0.5621	0.0173599\\
0.5647	0.017038\\
0.5672	0.0167286\\
0.5698	0.0164069\\
0.5723	0.0160981\\
0.5749	0.0157777\\
0.5775	0.0154585\\
0.58	0.0151522\\
0.5826	0.0148336\\
0.5852	0.0145147\\
0.5877	0.0142078\\
0.5903	0.0138889\\
0.5929	0.0135705\\
0.5954	0.0132647\\
0.598	0.0129465\\
0.6006	0.0126276\\
0.6031	0.0123202\\
0.6057	0.012\\
0.6082	0.011692\\
0.6108	0.0113717\\
0.6134	0.0110509\\
0.6159	0.0107417\\
0.6185	0.0104188\\
0.6211	0.0100947\\
0.6236	0.00978232\\
0.6262	0.00945691\\
0.6288	0.00913082\\
0.6313	0.0088162\\
0.6339	0.00848745\\
0.6364	0.00816975\\
0.639	0.00783793\\
0.6416	0.007505\\
0.6441	0.00718391\\
0.6467	0.00684869\\
0.6493	0.00651172\\
0.6518	0.00618579\\
0.6544	0.00584493\\
0.657	0.00550244\\
0.6595	0.00517182\\
0.6621	0.00482652\\
0.6646	0.00449278\\
0.6672	0.00414354\\
0.6698	0.00379205\\
0.6723	0.00345215\\
0.6749	0.00309696\\
0.6775	0.00274014\\
0.68	0.0023953\\
0.6826	0.00203453\\
0.6852	0.00167137\\
0.6877	0.00132003\\
0.6903	0.000952674\\
0.6928	0.00059779\\
0.6954	0.000226951\\
0.698	-0.000145949\\
0.7005	-0.0005067\\
0.7031	-0.000884192\\
0.7057	-0.00126382\\
0.7082	-0.00163057\\
0.7108	-0.00201368\\
0.7134	-0.00239864\\
0.7159	-0.00277078\\
0.7185	-0.00316001\\
0.721	-0.00353627\\
0.7236	-0.00392938\\
0.7262	-0.0043241\\
0.7287	-0.00470513\\
0.7313	-0.00510315\\
0.7339	-0.00550313\\
0.7364	-0.00588957\\
0.739	-0.00629317\\
0.7416	-0.00669821\\
0.7441	-0.0070889\\
0.7467	-0.00749659\\
0.7492	-0.00789007\\
0.7518	-0.00830089\\
0.7544	-0.00871323\\
0.7569	-0.00911088\\
0.7595	-0.00952543\\
0.7621	-0.00994094\\
0.7646	-0.0103415\\
0.7672	-0.0107593\\
0.7698	-0.0111783\\
0.7723	-0.011582\\
0.7749	-0.0120026\\
0.7775	-0.0124238\\
0.78	-0.0128293\\
0.7826	-0.0132518\\
0.7851	-0.0136587\\
0.7877	-0.0140825\\
0.7903	-0.0145067\\
0.7928	-0.0149148\\
0.7954	-0.0153392\\
0.798	-0.0157639\\
0.8005	-0.0161725\\
0.8031	-0.0165977\\
0.8057	-0.0170229\\
0.8082	-0.0174316\\
0.8108	-0.0178562\\
0.8133	-0.0182643\\
0.8159	-0.0186884\\
0.8185	-0.0191123\\
0.821	-0.0195197\\
0.8236	-0.0199427\\
0.8262	-0.020365\\
0.8287	-0.0207704\\
0.8313	-0.0211912\\
0.8339	-0.0216113\\
0.8364	-0.0220146\\
0.839	-0.0224332\\
0.8415	-0.0228346\\
0.8441	-0.0232509\\
0.8467	-0.023666\\
0.8492	-0.024064\\
0.8518	-0.0244768\\
0.8544	-0.0248884\\
0.8569	-0.0252829\\
0.8595	-0.0256915\\
0.8621	-0.0260986\\
0.8646	-0.0264884\\
0.8672	-0.0268924\\
0.8697	-0.0272793\\
0.8723	-0.0276801\\
0.8749	-0.0280789\\
0.8774	-0.0284605\\
0.88	-0.0288554\\
0.8826	-0.0292484\\
0.8851	-0.0296245\\
0.8877	-0.0300136\\
0.8903	-0.0304006\\
0.8928	-0.0307705\\
0.8954	-0.0311529\\
0.8979	-0.0315184\\
0.9005	-0.0318964\\
0.9031	-0.032272\\
0.9056	-0.032631\\
0.9082	-0.0330018\\
0.9108	-0.03337\\
0.9133	-0.0337215\\
0.9159	-0.0340846\\
0.9185	-0.0344451\\
0.921	-0.0347892\\
0.9236	-0.0351445\\
0.9262	-0.0354968\\
0.9287	-0.0358329\\
0.9313	-0.0361796\\
0.9338	-0.0365103\\
0.9364	-0.0368515\\
0.939	-0.0371897\\
0.9415	-0.0375121\\
0.9441	-0.0378443\\
0.9467	-0.0381734\\
0.9492	-0.0384869\\
0.9518	-0.03881\\
0.9544	-0.0391301\\
0.9569	-0.0394348\\
0.9595	-0.0397485\\
0.962	-0.040047\\
0.9646	-0.0403543\\
0.9672	-0.0406583\\
0.9697	-0.0409477\\
0.9723	-0.0412453\\
0.9749	-0.0415396\\
0.9774	-0.0418194\\
0.98	-0.042107\\
0.9826	-0.0423912\\
0.9851	-0.0426614\\
0.9877	-0.0429391\\
0.9902	-0.0432028\\
0.9928	-0.0434737\\
0.9954	-0.043741\\
0.9979	-0.0439949\\
1.0005	-0.0442555\\
1.0031	-0.0445127\\
1.0056	-0.0447567\\
1.0082	-0.0450071\\
1.0108	-0.0452539\\
1.0133	-0.0454879\\
1.0159	-0.0457279\\
1.0184	-0.0459554\\
1.021	-0.0461886\\
1.0236	-0.0464184\\
1.0261	-0.046636\\
1.0287	-0.0468588\\
1.0313	-0.0470782\\
1.0338	-0.0472859\\
1.0364	-0.0474986\\
1.039	-0.0477079\\
1.0415	-0.0479058\\
1.0441	-0.0481083\\
1.0466	-0.0482998\\
1.0492	-0.0484957\\
1.0518	-0.0486882\\
1.0543	-0.0488702\\
1.0569	-0.0490563\\
1.0595	-0.0492389\\
1.062	-0.0494115\\
1.0646	-0.0495877\\
1.0672	-0.0497607\\
1.0697	-0.0499241\\
1.0723	-0.0500909\\
1.0748	-0.0502483\\
1.0774	-0.0504089\\
1.08	-0.0505663\\
1.0825	-0.0507148\\
1.0851	-0.0508663\\
1.0877	-0.0510149\\
1.0902	-0.0511549\\
1.0928	-0.0512976\\
1.0954	-0.0514373\\
1.0979	-0.051569\\
1.1005	-0.0517032\\
1.1031	-0.0518346\\
1.1056	-0.0519583\\
1.1082	-0.0520843\\
1.1107	-0.0522029\\
1.1133	-0.0523237\\
1.1159	-0.0524418\\
1.1184	-0.052553\\
1.121	-0.0526662\\
1.1236	-0.052777\\
1.1261	-0.0528811\\
1.1287	-0.052987\\
1.1313	-0.0530906\\
1.1338	-0.053188\\
1.1364	-0.0532871\\
1.1389	-0.0533803\\
1.1415	-0.0534751\\
1.1441	-0.0535678\\
1.1466	-0.0536548\\
1.1492	-0.0537434\\
1.1518	-0.05383\\
1.1543	-0.0539114\\
1.1569	-0.0539943\\
1.1595	-0.0540752\\
1.162	-0.0541513\\
1.1646	-0.0542287\\
1.1671	-0.0543015\\
1.1697	-0.0543755\\
1.1723	-0.054448\\
1.1748	-0.0545162\\
1.1774	-0.0545855\\
1.18	-0.0546533\\
1.1825	-0.0547172\\
1.1851	-0.0547822\\
1.1877	-0.0548459\\
1.1902	-0.0549059\\
1.1928	-0.0549671\\
1.1953	-0.0550247\\
1.1979	-0.0550834\\
1.2005	-0.055141\\
1.203	-0.0551954\\
1.2056	-0.0552509\\
1.2082	-0.0553054\\
1.2107	-0.0553568\\
1.2133	-0.0554094\\
1.2159	-0.0554611\\
1.2184	-0.05551\\
1.221	-0.0555601\\
1.2235	-0.0556075\\
1.2261	-0.0556561\\
1.2287	-0.055704\\
1.2312	-0.0557494\\
1.2338	-0.0557961\\
1.2364	-0.0558422\\
1.2389	-0.055886\\
1.2415	-0.055931\\
1.2441	-0.0559756\\
1.2466	-0.056018\\
1.2492	-0.0560618\\
1.2518	-0.0561051\\
1.2543	-0.0561465\\
1.2569	-0.0561893\\
1.2594	-0.0562302\\
1.262	-0.0562724\\
1.2646	-0.0563145\\
1.2671	-0.0563547\\
1.2697	-0.0563965\\
1.2723	-0.0564381\\
1.2748	-0.0564781\\
1.2774	-0.0565195\\
1.28	-0.056561\\
1.2825	-0.0566008\\
1.2851	-0.0566423\\
1.2876	-0.0566822\\
1.2902	-0.0567238\\
1.2928	-0.0567654\\
1.2953	-0.0568056\\
1.2979	-0.0568475\\
1.3005	-0.0568895\\
1.303	-0.0569302\\
1.3056	-0.0569726\\
1.3082	-0.0570153\\
1.3107	-0.0570566\\
1.3133	-0.0570997\\
1.3158	-0.0571415\\
1.3184	-0.0571852\\
1.321	-0.0572292\\
1.3235	-0.0572719\\
1.3261	-0.0573166\\
1.3287	-0.0573616\\
1.3312	-0.0574053\\
1.3338	-0.0574511\\
1.3364	-0.0574973\\
1.3389	-0.0575421\\
1.3415	-0.0575891\\
1.344	-0.0576347\\
1.3466	-0.0576825\\
1.3492	-0.0577308\\
1.3517	-0.0577777\\
1.3543	-0.0578269\\
1.3569	-0.0578765\\
1.3594	-0.0579248\\
1.362	-0.0579754\\
1.3646	-0.0580265\\
1.3671	-0.0580761\\
1.3697	-0.0581281\\
1.3722	-0.0581787\\
1.3748	-0.0582318\\
1.3774	-0.0582853\\
1.3799	-0.0583373\\
1.3825	-0.0583919\\
1.3851	-0.0584469\\
1.3876	-0.0585004\\
1.3902	-0.0585564\\
1.3928	-0.058613\\
1.3953	-0.0586679\\
1.3979	-0.0587254\\
1.4005	-0.0587835\\
1.403	-0.0588398\\
1.4056	-0.0588988\\
1.4081	-0.058956\\
1.4107	-0.059016\\
1.4133	-0.0590764\\
1.4158	-0.059135\\
1.4184	-0.0591963\\
1.421	-0.0592581\\
1.4235	-0.059318\\
1.4261	-0.0593807\\
1.4287	-0.0594438\\
1.4312	-0.0595048\\
1.4338	-0.0595687\\
1.4363	-0.0596306\\
1.4389	-0.0596953\\
1.4415	-0.0597604\\
1.444	-0.0598233\\
1.4466	-0.0598891\\
1.4492	-0.0599552\\
1.4517	-0.0600191\\
1.4543	-0.0600859\\
1.4569	-0.060153\\
1.4594	-0.0602178\\
1.462	-0.0602854\\
1.4645	-0.0603507\\
1.4671	-0.0604189\\
1.4697	-0.0604873\\
1.4722	-0.0605533\\
1.4748	-0.0606221\\
1.4774	-0.0606911\\
1.4799	-0.0607577\\
1.4825	-0.060827\\
1.4851	-0.0608965\\
1.4876	-0.0609635\\
1.4902	-0.0610333\\
1.4927	-0.0611004\\
1.4953	-0.0611704\\
1.4979	-0.0612404\\
1.5004	-0.0613077\\
1.503	-0.0613778\\
1.5056	-0.061448\\
1.5081	-0.0615154\\
1.5107	-0.0615855\\
1.5133	-0.0616555\\
1.5158	-0.0617228\\
1.5184	-0.0617928\\
1.5209	-0.06186\\
1.5235	-0.0619297\\
1.5261	-0.0619994\\
1.5286	-0.0620662\\
1.5312	-0.0621355\\
1.5338	-0.0622047\\
1.5363	-0.062271\\
1.5389	-0.0623398\\
1.5415	-0.0624084\\
1.544	-0.0624741\\
1.5466	-0.0625422\\
1.5491	-0.0626074\\
1.5517	-0.0626749\\
1.5543	-0.0627421\\
1.5568	-0.0628064\\
1.5594	-0.062873\\
1.562	-0.0629392\\
1.5645	-0.0630024\\
1.5671	-0.0630679\\
1.5697	-0.0631329\\
1.5722	-0.063195\\
1.5748	-0.0632592\\
1.5774	-0.0633229\\
1.5799	-0.0633838\\
1.5825	-0.0634465\\
1.585	-0.0635064\\
1.5876	-0.0635682\\
1.5902	-0.0636294\\
1.5927	-0.0636878\\
1.5953	-0.063748\\
1.5979	-0.0638075\\
1.6004	-0.0638643\\
1.603	-0.0639227\\
1.6056	-0.0639805\\
1.6081	-0.0640354\\
1.6107	-0.064092\\
1.6132	-0.0641457\\
1.6158	-0.0642009\\
1.6184	-0.0642555\\
1.6209	-0.0643073\\
1.6235	-0.0643605\\
1.6261	-0.0644129\\
1.6286	-0.0644627\\
1.6312	-0.0645137\\
1.6338	-0.0645639\\
1.6363	-0.0646115\\
1.6389	-0.0646603\\
1.6414	-0.0647065\\
1.644	-0.0647537\\
1.6466	-0.0648001\\
1.6491	-0.0648439\\
1.6517	-0.0648888\\
1.6543	-0.0649328\\
1.6568	-0.0649743\\
1.6594	-0.0650166\\
1.662	-0.0650581\\
1.6645	-0.0650972\\
1.6671	-0.065137\\
1.6696	-0.0651744\\
1.6722	-0.0652125\\
1.6748	-0.0652497\\
1.6773	-0.0652847\\
1.6799	-0.0653201\\
1.6825	-0.0653547\\
1.685	-0.065387\\
1.6876	-0.0654198\\
1.6902	-0.0654517\\
1.6927	-0.0654814\\
1.6953	-0.0655115\\
1.6978	-0.0655395\\
1.7004	-0.0655678\\
1.703	-0.0655951\\
1.7055	-0.0656205\\
1.7081	-0.0656459\\
1.7107	-0.0656705\\
1.7132	-0.0656932\\
1.7158	-0.0657159\\
1.7184	-0.0657377\\
1.7209	-0.0657577\\
1.7235	-0.0657777\\
1.7261	-0.0657967\\
1.7286	-0.0658141\\
1.7312	-0.0658312\\
1.7337	-0.0658468\\
1.7363	-0.0658622\\
1.7389	-0.0658766\\
1.7414	-0.0658895\\
1.744	-0.0659021\\
1.7466	-0.0659137\\
1.7491	-0.0659241\\
1.7517	-0.0659339\\
1.7543	-0.0659428\\
1.7568	-0.0659506\\
1.7594	-0.0659577\\
1.7619	-0.0659637\\
1.7645	-0.0659691\\
1.7671	-0.0659736\\
1.7696	-0.0659771\\
1.7722	-0.0659798\\
1.7748	-0.0659817\\
1.7773	-0.0659826\\
1.7799	-0.0659828\\
1.7825	-0.0659821\\
1.785	-0.0659806\\
1.7876	-0.0659783\\
1.7901	-0.0659752\\
1.7927	-0.0659712\\
1.7953	-0.0659664\\
1.7978	-0.065961\\
1.8004	-0.0659546\\
1.803	-0.0659474\\
1.8055	-0.0659397\\
1.8081	-0.0659309\\
1.8107	-0.0659214\\
1.8132	-0.0659115\\
1.8158	-0.0659005\\
1.8183	-0.0658892\\
1.8209	-0.0658767\\
1.8235	-0.0658635\\
1.826	-0.0658501\\
1.8286	-0.0658355\\
1.8312	-0.0658202\\
1.8337	-0.0658048\\
1.8363	-0.0657882\\
1.8389	-0.0657709\\
1.8414	-0.0657537\\
1.844	-0.0657351\\
1.8465	-0.0657166\\
1.8491	-0.0656968\\
1.8517	-0.0656764\\
1.8542	-0.0656562\\
1.8568	-0.0656347\\
1.8594	-0.0656125\\
1.8619	-0.0655907\\
1.8645	-0.0655675\\
1.8671	-0.0655437\\
1.8696	-0.0655204\\
1.8722	-0.0654956\\
1.8747	-0.0654713\\
1.8773	-0.0654455\\
1.8799	-0.0654193\\
1.8824	-0.0653936\\
1.885	-0.0653665\\
1.8876	-0.0653389\\
1.8901	-0.065312\\
1.8927	-0.0652836\\
1.8953	-0.0652548\\
1.8978	-0.0652268\\
1.9004	-0.0651973\\
1.903	-0.0651674\\
1.9055	-0.0651383\\
1.9081	-0.0651078\\
1.9106	-0.0650781\\
1.9132	-0.0650469\\
1.9158	-0.0650154\\
1.9183	-0.0649849\\
1.9209	-0.0649529\\
1.9235	-0.0649206\\
1.926	-0.0648894\\
1.9286	-0.0648567\\
1.9312	-0.0648237\\
1.9337	-0.0647919\\
1.9363	-0.0647585\\
1.9388	-0.0647263\\
1.9414	-0.0646927\\
1.944	-0.0646589\\
1.9465	-0.0646262\\
1.9491	-0.0645922\\
1.9517	-0.064558\\
1.9542	-0.0645251\\
1.9568	-0.0644907\\
1.9594	-0.0644563\\
1.9619	-0.0644231\\
1.9645	-0.0643886\\
1.967	-0.0643554\\
1.9696	-0.0643208\\
1.9722	-0.0642862\\
1.9747	-0.0642529\\
1.9773	-0.0642183\\
1.9799	-0.0641838\\
1.9824	-0.0641506\\
1.985	-0.0641161\\
1.9876	-0.0640817\\
1.9901	-0.0640487\\
1.9927	-0.0640145\\
1.9952	-0.0639816\\
1.9978	-0.0639476\\
2.0004	-0.0639137\\
2.0029	-0.0638812\\
2.0055	-0.0638475\\
2.0081	-0.063814\\
2.0106	-0.063782\\
2.0132	-0.0637488\\
2.0158	-0.0637158\\
2.0183	-0.0636843\\
2.0209	-0.0636517\\
2.0234	-0.0636206\\
2.026	-0.0635884\\
2.0286	-0.0635565\\
2.0311	-0.063526\\
2.0337	-0.0634946\\
2.0363	-0.0634635\\
2.0388	-0.0634338\\
2.0414	-0.0634032\\
2.044	-0.0633729\\
2.0465	-0.063344\\
2.0491	-0.0633143\\
2.0517	-0.0632849\\
2.0542	-0.063257\\
2.0568	-0.0632283\\
2.0593	-0.063201\\
2.0619	-0.0631729\\
2.0645	-0.0631452\\
2.067	-0.063119\\
2.0696	-0.063092\\
2.0722	-0.0630655\\
2.0747	-0.0630403\\
2.0773	-0.0630145\\
2.0799	-0.0629891\\
2.0824	-0.0629651\\
2.085	-0.0629405\\
2.0875	-0.0629173\\
2.0901	-0.0628936\\
2.0927	-0.0628703\\
2.0952	-0.0628483\\
2.0978	-0.0628259\\
2.1004	-0.062804\\
2.1029	-0.0627833\\
2.1055	-0.0627622\\
2.1081	-0.0627417\\
2.1106	-0.0627223\\
2.1132	-0.0627027\\
2.1157	-0.0626842\\
2.1183	-0.0626655\\
2.1209	-0.0626473\\
2.1234	-0.0626303\\
2.126	-0.0626131\\
2.1286	-0.0625964\\
2.1311	-0.0625808\\
2.1337	-0.0625651\\
2.1363	-0.0625499\\
2.1388	-0.0625357\\
2.1414	-0.0625215\\
2.1439	-0.0625084\\
2.1465	-0.0624952\\
2.1491	-0.0624826\\
2.1516	-0.0624709\\
2.1542	-0.0624593\\
2.1568	-0.0624483\\
2.1593	-0.0624381\\
2.1619	-0.0624281\\
2.1645	-0.0624186\\
2.167	-0.06241\\
2.1696	-0.0624016\\
2.1721	-0.062394\\
2.1747	-0.0623866\\
2.1773	-0.0623798\\
2.1798	-0.0623737\\
2.1824	-0.0623679\\
2.185	-0.0623627\\
2.1875	-0.0623581\\
2.1901	-0.0623539\\
2.1927	-0.0623503\\
2.1952	-0.0623473\\
2.1978	-0.0623447\\
2.2004	-0.0623426\\
2.2029	-0.0623411\\
2.2055	-0.0623401\\
2.208	-0.0623396\\
2.2106	-0.0623396\\
2.2132	-0.0623402\\
2.2157	-0.0623412\\
2.2183	-0.0623427\\
2.2209	-0.0623448\\
2.2234	-0.0623473\\
2.226	-0.0623504\\
2.2286	-0.062354\\
2.2311	-0.0623579\\
2.2337	-0.0623625\\
2.2362	-0.0623674\\
2.2388	-0.062373\\
2.2414	-0.0623791\\
2.2439	-0.0623854\\
2.2465	-0.0623924\\
2.2491	-0.0623999\\
2.2516	-0.0624076\\
2.2542	-0.0624161\\
2.2568	-0.062425\\
2.2593	-0.062434\\
2.2619	-0.0624439\\
2.2644	-0.0624538\\
2.267	-0.0624645\\
2.2696	-0.0624757\\
2.2721	-0.0624869\\
2.2747	-0.0624989\\
2.2773	-0.0625114\\
2.2798	-0.0625238\\
2.2824	-0.0625371\\
2.285	-0.0625509\\
2.2875	-0.0625645\\
2.2901	-0.062579\\
2.2926	-0.0625934\\
2.2952	-0.0626087\\
2.2978	-0.0626244\\
2.3003	-0.0626399\\
2.3029	-0.0626563\\
2.3055	-0.0626732\\
2.308	-0.0626897\\
2.3106	-0.0627073\\
2.3132	-0.0627251\\
2.3157	-0.0627427\\
2.3183	-0.0627612\\
2.3208	-0.0627794\\
2.3234	-0.0627986\\
2.326	-0.0628182\\
2.3285	-0.0628372\\
2.3311	-0.0628574\\
2.3337	-0.0628778\\
2.3362	-0.0628978\\
2.3388	-0.0629188\\
2.3414	-0.0629401\\
2.3439	-0.0629608\\
2.3465	-0.0629826\\
2.349	-0.0630038\\
2.3516	-0.0630261\\
2.3542	-0.0630487\\
2.3567	-0.0630706\\
2.3593	-0.0630936\\
2.3619	-0.0631169\\
2.3644	-0.0631394\\
2.367	-0.063163\\
2.3696	-0.0631869\\
2.3721	-0.06321\\
2.3747	-0.0632342\\
2.3773	-0.0632586\\
2.3798	-0.0632822\\
2.3824	-0.0633069\\
2.3849	-0.0633308\\
2.3875	-0.0633558\\
2.3901	-0.063381\\
2.3926	-0.0634053\\
2.3952	-0.0634307\\
2.3978	-0.0634562\\
2.4003	-0.0634808\\
2.4029	-0.0635066\\
2.4055	-0.0635324\\
2.408	-0.0635573\\
2.4106	-0.0635833\\
2.4131	-0.0636083\\
2.4157	-0.0636344\\
2.4183	-0.0636606\\
2.4208	-0.0636858\\
2.4234	-0.063712\\
2.426	-0.0637383\\
2.4285	-0.0637636\\
2.4311	-0.06379\\
2.4337	-0.0638163\\
2.4362	-0.0638416\\
2.4388	-0.063868\\
2.4413	-0.0638933\\
2.4439	-0.0639196\\
2.4465	-0.0639459\\
2.449	-0.0639712\\
2.4516	-0.0639974\\
2.4542	-0.0640236\\
2.4567	-0.0640487\\
2.4593	-0.0640747\\
2.4619	-0.0641007\\
2.4644	-0.0641257\\
2.467	-0.0641515\\
2.4695	-0.0641763\\
2.4721	-0.064202\\
2.4747	-0.0642275\\
2.4772	-0.064252\\
2.4798	-0.0642773\\
2.4824	-0.0643026\\
2.4849	-0.0643267\\
2.4875	-0.0643517\\
2.4901	-0.0643765\\
2.4926	-0.0644002\\
2.4952	-0.0644248\\
2.4977	-0.0644482\\
2.5003	-0.0644724\\
2.5029	-0.0644964\\
2.5054	-0.0645194\\
2.508	-0.0645431\\
2.5106	-0.0645666\\
2.5131	-0.064589\\
2.5157	-0.0646121\\
2.5183	-0.064635\\
2.5208	-0.0646568\\
2.5234	-0.0646793\\
2.526	-0.0647015\\
2.5285	-0.0647227\\
2.5311	-0.0647445\\
2.5336	-0.0647652\\
2.5362	-0.0647865\\
2.5388	-0.0648076\\
2.5413	-0.0648276\\
2.5439	-0.0648482\\
2.5465	-0.0648685\\
2.549	-0.0648877\\
2.5516	-0.0649075\\
2.5542	-0.064927\\
2.5567	-0.0649455\\
2.5593	-0.0649644\\
2.5618	-0.0649823\\
2.5644	-0.0650007\\
2.567	-0.0650187\\
2.5695	-0.0650358\\
2.5721	-0.0650532\\
2.5747	-0.0650703\\
2.5772	-0.0650865\\
2.5798	-0.065103\\
2.5824	-0.0651192\\
2.5849	-0.0651344\\
2.5875	-0.06515\\
2.59	-0.0651646\\
2.5926	-0.0651795\\
2.5952	-0.065194\\
2.5977	-0.0652076\\
2.6003	-0.0652215\\
2.6029	-0.065235\\
2.6054	-0.0652476\\
2.608	-0.0652604\\
2.6106	-0.0652729\\
2.6131	-0.0652845\\
2.6157	-0.0652963\\
2.6182	-0.0653072\\
2.6208	-0.0653182\\
2.6234	-0.0653288\\
2.6259	-0.0653387\\
2.6285	-0.0653486\\
2.6311	-0.0653582\\
2.6336	-0.065367\\
2.6362	-0.0653757\\
2.6388	-0.0653841\\
2.6413	-0.0653919\\
2.6439	-0.0653995\\
2.6464	-0.0654065\\
2.649	-0.0654134\\
2.6516	-0.0654199\\
2.6541	-0.0654257\\
2.6567	-0.0654314\\
2.6593	-0.0654368\\
2.6618	-0.0654415\\
2.6644	-0.0654461\\
2.667	-0.0654502\\
2.6695	-0.0654538\\
2.6721	-0.0654572\\
2.6746	-0.0654601\\
2.6772	-0.0654627\\
2.6798	-0.0654649\\
2.6823	-0.0654666\\
2.6849	-0.065468\\
2.6875	-0.065469\\
2.69	-0.0654696\\
2.6926	-0.0654698\\
2.6952	-0.0654696\\
2.6977	-0.065469\\
2.7003	-0.0654681\\
2.7029	-0.0654667\\
2.7054	-0.065465\\
2.708	-0.0654628\\
2.7105	-0.0654604\\
2.7131	-0.0654575\\
2.7157	-0.0654541\\
2.7182	-0.0654505\\
2.7208	-0.0654464\\
2.7234	-0.0654419\\
2.7259	-0.0654372\\
2.7285	-0.0654319\\
2.7311	-0.0654262\\
2.7336	-0.0654204\\
2.7362	-0.065414\\
2.7387	-0.0654074\\
2.7413	-0.0654002\\
2.7439	-0.0653926\\
2.7464	-0.0653849\\
2.749	-0.0653766\\
2.7516	-0.0653678\\
2.7541	-0.0653591\\
2.7567	-0.0653496\\
2.7593	-0.0653397\\
2.7618	-0.0653299\\
2.7644	-0.0653193\\
2.7669	-0.0653087\\
2.7695	-0.0652974\\
2.7721	-0.0652857\\
2.7746	-0.0652741\\
2.7772	-0.0652617\\
2.7798	-0.0652489\\
2.7823	-0.0652363\\
2.7849	-0.0652228\\
2.7875	-0.0652089\\
2.79	-0.0651953\\
2.7926	-0.0651808\\
2.7951	-0.0651664\\
2.7977	-0.0651512\\
2.8003	-0.0651356\\
2.8028	-0.0651203\\
2.8054	-0.0651041\\
2.808	-0.0650875\\
2.8105	-0.0650713\\
2.8131	-0.065054\\
2.8157	-0.0650365\\
2.8182	-0.0650192\\
2.8208	-0.065001\\
2.8233	-0.0649832\\
2.8259	-0.0649644\\
2.8285	-0.0649452\\
2.831	-0.0649265\\
2.8336	-0.0649067\\
2.8362	-0.0648866\\
2.8387	-0.064867\\
2.8413	-0.0648464\\
2.8439	-0.0648254\\
2.8464	-0.0648049\\
2.849	-0.0647833\\
2.8516	-0.0647615\\
2.8541	-0.0647402\\
2.8567	-0.0647178\\
2.8592	-0.0646959\\
2.8618	-0.064673\\
2.8644	-0.0646497\\
2.8669	-0.0646271\\
2.8695	-0.0646033\\
2.8721	-0.0645792\\
2.8746	-0.0645558\\
2.8772	-0.0645313\\
2.8798	-0.0645064\\
2.8823	-0.0644823\\
2.8849	-0.064457\\
2.8874	-0.0644324\\
2.89	-0.0644065\\
2.8926	-0.0643805\\
2.8951	-0.0643552\\
2.8977	-0.0643286\\
2.9003	-0.0643018\\
2.9028	-0.0642759\\
2.9054	-0.0642487\\
2.908	-0.0642212\\
2.9105	-0.0641946\\
2.9131	-0.0641667\\
2.9156	-0.0641397\\
2.9182	-0.0641114\\
2.9208	-0.0640829\\
2.9233	-0.0640552\\
2.9259	-0.0640263\\
2.9285	-0.0639972\\
2.931	-0.063969\\
2.9336	-0.0639395\\
2.9362	-0.0639098\\
2.9387	-0.0638811\\
2.9413	-0.063851\\
2.9438	-0.0638219\\
2.9464	-0.0637915\\
2.949	-0.0637609\\
2.9515	-0.0637313\\
2.9541	-0.0637004\\
2.9567	-0.0636693\\
2.9592	-0.0636392\\
2.9618	-0.0636078\\
2.9644	-0.0635762\\
2.9669	-0.0635457\\
2.9695	-0.0635138\\
2.972	-0.063483\\
2.9746	-0.0634508\\
2.9772	-0.0634185\\
2.9797	-0.0633873\\
2.9823	-0.0633547\\
2.9849	-0.0633219\\
2.9874	-0.0632903\\
2.99	-0.0632573\\
2.9926	-0.0632242\\
2.9951	-0.0631923\\
2.9977	-0.0631589\\
3.0003	-0.0631254\\
3.0028	-0.0630931\\
3.0054	-0.0630594\\
3.0079	-0.0630269\\
3.0105	-0.062993\\
3.0131	-0.062959\\
3.0156	-0.0629261\\
3.0182	-0.0628919\\
3.0208	-0.0628576\\
3.0233	-0.0628245\\
3.0259	-0.06279\\
3.0285	-0.0627554\\
3.031	-0.062722\\
3.0336	-0.0626873\\
3.0361	-0.0626538\\
3.0387	-0.0626188\\
3.0413	-0.0625838\\
3.0438	-0.0625501\\
3.0464	-0.062515\\
3.049	-0.0624797\\
3.0515	-0.0624458\\
3.0541	-0.0624105\\
3.0567	-0.062375\\
3.0592	-0.0623409\\
3.0618	-0.0623054\\
3.0643	-0.0622712\\
3.0669	-0.0622355\\
3.0695	-0.0621998\\
3.072	-0.0621655\\
3.0746	-0.0621297\\
3.0772	-0.0620938\\
3.0797	-0.0620593\\
3.0823	-0.0620234\\
3.0849	-0.0619875\\
3.0874	-0.0619528\\
3.09	-0.0619168\\
3.0925	-0.0618821\\
3.0951	-0.061846\\
3.0977	-0.0618099\\
3.1002	-0.0617751\\
3.1028	-0.061739\\
3.1054	-0.0617028\\
3.1079	-0.0616679\\
3.1105	-0.0616317\\
3.1131	-0.0615954\\
3.1156	-0.0615605\\
3.1182	-0.0615242\\
3.1207	-0.0614893\\
3.1233	-0.061453\\
3.1259	-0.0614167\\
3.1284	-0.0613817\\
3.131	-0.0613454\\
3.1336	-0.0613091\\
3.1361	-0.0612741\\
3.1387	-0.0612378\\
3.1413	-0.0612014\\
3.1438	-0.0611665\\
3.1464	-0.0611301\\
3.1489	-0.0610952\\
3.1515	-0.0610589\\
3.1541	-0.0610225\\
3.1566	-0.0609876\\
3.1592	-0.0609513\\
3.1618	-0.060915\\
3.1643	-0.0608802\\
3.1669	-0.0608439\\
3.1695	-0.0608077\\
3.172	-0.0607728\\
3.1746	-0.0607366\\
3.1772	-0.0607005\\
3.1797	-0.0606657\\
3.1823	-0.0606296\\
3.1848	-0.0605948\\
3.1874	-0.0605588\\
3.19	-0.0605227\\
3.1925	-0.0604881\\
3.1951	-0.0604521\\
3.1977	-0.0604161\\
3.2002	-0.0603816\\
3.2028	-0.0603457\\
3.2054	-0.0603098\\
3.2079	-0.0602754\\
3.2105	-0.0602396\\
3.213	-0.0602052\\
3.2156	-0.0601695\\
3.2182	-0.0601339\\
3.2207	-0.0600996\\
3.2233	-0.060064\\
3.2259	-0.0600285\\
3.2284	-0.0599944\\
3.231	-0.0599589\\
3.2336	-0.0599235\\
3.2361	-0.0598896\\
3.2387	-0.0598543\\
3.2412	-0.0598204\\
3.2438	-0.0597852\\
3.2464	-0.05975\\
3.2489	-0.0597163\\
3.2515	-0.0596813\\
3.2541	-0.0596463\\
3.2566	-0.0596127\\
3.2592	-0.0595779\\
3.2618	-0.0595431\\
3.2643	-0.0595097\\
3.2669	-0.059475\\
3.2694	-0.0594417\\
3.272	-0.0594072\\
3.2746	-0.0593727\\
3.2771	-0.0593396\\
3.2797	-0.0593052\\
3.2823	-0.0592709\\
3.2848	-0.059238\\
3.2874	-0.0592039\\
3.29	-0.0591698\\
3.2925	-0.0591371\\
3.2951	-0.0591032\\
3.2976	-0.0590706\\
3.3002	-0.0590368\\
3.3028	-0.0590031\\
3.3053	-0.0589707\\
3.3079	-0.0589371\\
3.3105	-0.0589036\\
3.313	-0.0588715\\
3.3156	-0.0588381\\
3.3182	-0.0588049\\
3.3207	-0.058773\\
3.3233	-0.0587399\\
3.3259	-0.0587068\\
3.3284	-0.0586752\\
3.331	-0.0586423\\
3.3335	-0.0586108\\
3.3361	-0.0585781\\
3.3387	-0.0585455\\
3.3412	-0.0585142\\
3.3438	-0.0584818\\
3.3464	-0.0584495\\
3.3489	-0.0584184\\
3.3515	-0.0583863\\
3.3541	-0.0583542\\
3.3566	-0.0583234\\
3.3592	-0.0582915\\
3.3617	-0.058261\\
3.3643	-0.0582293\\
3.3669	-0.0581976\\
3.3694	-0.0581673\\
3.372	-0.0581359\\
3.3746	-0.0581046\\
3.3771	-0.0580746\\
3.3797	-0.0580434\\
3.3823	-0.0580124\\
3.3848	-0.0579826\\
3.3874	-0.0579518\\
3.3899	-0.0579223\\
3.3925	-0.0578916\\
3.3951	-0.0578611\\
3.3976	-0.0578319\\
3.4002	-0.0578015\\
3.4028	-0.0577713\\
3.4053	-0.0577424\\
3.4079	-0.0577124\\
3.4105	-0.0576825\\
3.413	-0.0576538\\
3.4156	-0.0576241\\
3.4181	-0.0575957\\
3.4207	-0.0575663\\
3.4233	-0.0575369\\
3.4258	-0.0575088\\
3.4284	-0.0574797\\
3.431	-0.0574507\\
3.4335	-0.0574229\\
3.4361	-0.0573941\\
3.4387	-0.0573654\\
3.4412	-0.057338\\
3.4438	-0.0573096\\
3.4463	-0.0572823\\
3.4489	-0.0572541\\
3.4515	-0.0572261\\
3.454	-0.0571992\\
3.4566	-0.0571713\\
3.4592	-0.0571436\\
3.4617	-0.0571171\\
3.4643	-0.0570897\\
3.4669	-0.0570623\\
3.4694	-0.0570361\\
3.472	-0.0570091\\
3.4745	-0.0569831\\
3.4771	-0.0569563\\
3.4797	-0.0569296\\
3.4822	-0.056904\\
3.4848	-0.0568776\\
3.4874	-0.0568513\\
3.4899	-0.0568261\\
3.4925	-0.0568001\\
3.4951	-0.0567741\\
3.4976	-0.0567493\\
3.5002	-0.0567237\\
3.5028	-0.0566982\\
3.5053	-0.0566738\\
3.5079	-0.0566485\\
3.5104	-0.0566244\\
3.513	-0.0565994\\
3.5156	-0.0565746\\
3.5181	-0.0565509\\
3.5207	-0.0565263\\
3.5233	-0.0565019\\
3.5258	-0.0564786\\
3.5284	-0.0564545\\
3.531	-0.0564305\\
3.5335	-0.0564076\\
3.5361	-0.0563839\\
3.5386	-0.0563613\\
3.5412	-0.0563379\\
3.5438	-0.0563147\\
3.5463	-0.0562925\\
3.5489	-0.0562695\\
3.5515	-0.0562467\\
3.554	-0.056225\\
3.5566	-0.0562025\\
3.5592	-0.0561801\\
3.5617	-0.0561588\\
3.5643	-0.0561368\\
3.5668	-0.0561157\\
3.5694	-0.056094\\
3.572	-0.0560725\\
3.5745	-0.0560519\\
3.5771	-0.0560306\\
3.5797	-0.0560095\\
3.5822	-0.0559894\\
3.5848	-0.0559687\\
3.5874	-0.055948\\
3.5899	-0.0559284\\
3.5925	-0.0559081\\
3.595	-0.0558888\\
3.5976	-0.0558688\\
3.6002	-0.055849\\
3.6027	-0.0558301\\
3.6053	-0.0558107\\
3.6079	-0.0557914\\
3.6104	-0.055773\\
3.613	-0.055754\\
3.6156	-0.0557352\\
3.6181	-0.0557173\\
3.6207	-0.0556989\\
3.6232	-0.0556813\\
3.6258	-0.0556632\\
3.6284	-0.0556453\\
3.6309	-0.0556282\\
3.6335	-0.0556106\\
3.6361	-0.0555932\\
3.6386	-0.0555766\\
3.6412	-0.0555595\\
3.6438	-0.0555426\\
3.6463	-0.0555266\\
3.6489	-0.05551\\
3.6515	-0.0554937\\
3.654	-0.0554781\\
3.6566	-0.0554621\\
3.6591	-0.0554469\\
3.6617	-0.0554313\\
3.6643	-0.0554158\\
3.6668	-0.0554011\\
3.6694	-0.055386\\
3.672	-0.0553711\\
3.6745	-0.055357\\
3.6771	-0.0553424\\
3.6797	-0.0553281\\
3.6822	-0.0553144\\
3.6848	-0.0553005\\
3.6873	-0.0552872\\
3.6899	-0.0552736\\
3.6925	-0.0552601\\
3.695	-0.0552474\\
3.6976	-0.0552343\\
3.7002	-0.0552215\\
3.7027	-0.0552093\\
3.7053	-0.0551968\\
3.7079	-0.0551845\\
3.7104	-0.0551729\\
3.713	-0.0551609\\
3.7155	-0.0551497\\
3.7181	-0.0551381\\
3.7207	-0.0551268\\
3.7232	-0.0551161\\
3.7258	-0.0551051\\
3.7284	-0.0550943\\
3.7309	-0.0550842\\
3.7335	-0.0550738\\
3.7361	-0.0550636\\
3.7386	-0.055054\\
3.7412	-0.0550442\\
3.7437	-0.055035\\
3.7463	-0.0550256\\
3.7489	-0.0550164\\
3.7514	-0.0550077\\
3.754	-0.0549989\\
3.7566	-0.0549903\\
3.7591	-0.0549822\\
3.7617	-0.0549739\\
3.7643	-0.0549659\\
3.7668	-0.0549584\\
3.7694	-0.0549507\\
3.7719	-0.0549436\\
3.7745	-0.0549363\\
3.7771	-0.0549293\\
3.7796	-0.0549227\\
3.7822	-0.054916\\
3.7848	-0.0549096\\
3.7873	-0.0549036\\
3.7899	-0.0548975\\
3.7925	-0.0548917\\
3.795	-0.0548862\\
3.7976	-0.0548807\\
3.8002	-0.0548755\\
3.8027	-0.0548706\\
3.8053	-0.0548657\\
3.8078	-0.0548613\\
3.8104	-0.0548568\\
3.813	-0.0548525\\
3.8155	-0.0548486\\
3.8181	-0.0548447\\
3.8207	-0.054841\\
3.8232	-0.0548377\\
3.8258	-0.0548344\\
3.8284	-0.0548313\\
3.8309	-0.0548285\\
3.8335	-0.0548258\\
3.836	-0.0548234\\
3.8386	-0.0548211\\
3.8412	-0.054819\\
3.8437	-0.0548172\\
3.8463	-0.0548155\\
3.8489	-0.054814\\
3.8514	-0.0548127\\
3.854	-0.0548116\\
3.8566	-0.0548107\\
3.8591	-0.05481\\
3.8617	-0.0548094\\
3.8642	-0.0548091\\
3.8668	-0.0548089\\
3.8694	-0.054809\\
3.8719	-0.0548092\\
3.8745	-0.0548096\\
3.8771	-0.0548102\\
3.8796	-0.054811\\
3.8822	-0.054812\\
3.8848	-0.0548132\\
3.8873	-0.0548145\\
3.8899	-0.0548161\\
3.8924	-0.0548178\\
3.895	-0.0548197\\
3.8976	-0.0548218\\
3.9001	-0.0548241\\
3.9027	-0.0548266\\
3.9053	-0.0548293\\
3.9078	-0.054832\\
3.9104	-0.0548351\\
3.913	-0.0548383\\
3.9155	-0.0548416\\
3.9181	-0.0548453\\
3.9206	-0.0548489\\
3.9232	-0.0548529\\
3.9258	-0.054857\\
3.9283	-0.0548612\\
3.9309	-0.0548657\\
3.9335	-0.0548704\\
3.936	-0.0548751\\
3.9386	-0.0548801\\
3.9412	-0.0548854\\
3.9437	-0.0548905\\
3.9463	-0.0548961\\
3.9488	-0.0549016\\
3.9514	-0.0549076\\
3.954	-0.0549137\\
3.9565	-0.0549197\\
3.9591	-0.0549261\\
3.9617	-0.0549327\\
3.9642	-0.0549392\\
3.9668	-0.0549461\\
3.9694	-0.0549532\\
3.9719	-0.0549602\\
3.9745	-0.0549677\\
3.9771	-0.0549752\\
3.9796	-0.0549827\\
3.9822	-0.0549906\\
3.9847	-0.0549984\\
3.9873	-0.0550066\\
3.9899	-0.055015\\
3.9924	-0.0550232\\
3.995	-0.0550319\\
3.9976	-0.0550408\\
4.0001	-0.0550494\\
4.0027	-0.0550586\\
4.0053	-0.0550679\\
4.0078	-0.055077\\
4.0104	-0.0550866\\
4.0129	-0.055096\\
4.0155	-0.0551059\\
4.0181	-0.055116\\
4.0206	-0.0551258\\
4.0232	-0.0551361\\
4.0258	-0.0551466\\
4.0283	-0.0551568\\
4.0309	-0.0551675\\
4.0335	-0.0551784\\
4.036	-0.055189\\
4.0386	-0.0552002\\
4.0411	-0.055211\\
4.0437	-0.0552225\\
4.0463	-0.055234\\
4.0488	-0.0552452\\
4.0514	-0.0552571\\
4.054	-0.055269\\
4.0565	-0.0552806\\
4.0591	-0.0552928\\
4.0617	-0.0553051\\
4.0642	-0.055317\\
4.0668	-0.0553296\\
4.0693	-0.0553418\\
4.0719	-0.0553545\\
4.0745	-0.0553674\\
4.077	-0.05538\\
4.0796	-0.0553931\\
4.0822	-0.0554063\\
4.0847	-0.0554191\\
4.0873	-0.0554326\\
4.0899	-0.0554461\\
4.0924	-0.0554593\\
4.095	-0.055473\\
4.0975	-0.0554864\\
4.1001	-0.0555003\\
4.1027	-0.0555144\\
4.1052	-0.055528\\
4.1078	-0.0555422\\
4.1104	-0.0555566\\
4.1129	-0.0555705\\
4.1155	-0.055585\\
4.1181	-0.0555996\\
4.1206	-0.0556137\\
4.1232	-0.0556285\\
4.1258	-0.0556434\\
4.1283	-0.0556578\\
4.1309	-0.0556728\\
4.1334	-0.0556873\\
4.136	-0.0557025\\
4.1386	-0.0557178\\
4.1411	-0.0557325\\
4.1437	-0.0557479\\
4.1463	-0.0557634\\
4.1488	-0.0557784\\
4.1514	-0.055794\\
4.154	-0.0558097\\
4.1565	-0.0558248\\
4.1591	-0.0558406\\
4.1616	-0.0558559\\
4.1642	-0.0558718\\
4.1668	-0.0558878\\
4.1693	-0.0559032\\
4.1719	-0.0559193\\
4.1745	-0.0559355\\
4.177	-0.055951\\
4.1796	-0.0559673\\
4.1822	-0.0559836\\
4.1847	-0.0559993\\
4.1873	-0.0560157\\
4.1898	-0.0560315\\
4.1924	-0.0560479\\
4.195	-0.0560644\\
4.1975	-0.0560803\\
4.2001	-0.0560969\\
4.2027	-0.0561135\\
4.2052	-0.0561295\\
4.2078	-0.0561462\\
4.2104	-0.0561629\\
4.2129	-0.056179\\
4.2155	-0.0561957\\
4.218	-0.0562118\\
4.2206	-0.0562286\\
4.2232	-0.0562455\\
4.2257	-0.0562616\\
4.2283	-0.0562785\\
4.2309	-0.0562954\\
4.2334	-0.0563116\\
4.236	-0.0563285\\
4.2386	-0.0563454\\
4.2411	-0.0563616\\
4.2437	-0.0563786\\
4.2462	-0.0563948\\
4.2488	-0.0564118\\
4.2514	-0.0564287\\
4.2539	-0.056445\\
4.2565	-0.0564619\\
4.2591	-0.0564788\\
4.2616	-0.0564951\\
4.2642	-0.056512\\
4.2668	-0.0565289\\
4.2693	-0.0565451\\
4.2719	-0.056562\\
4.2744	-0.0565782\\
4.277	-0.0565951\\
4.2796	-0.0566119\\
4.2821	-0.0566281\\
4.2847	-0.0566449\\
4.2873	-0.0566617\\
4.2898	-0.0566778\\
4.2924	-0.0566945\\
4.295	-0.0567112\\
4.2975	-0.0567272\\
4.3001	-0.0567439\\
4.3027	-0.0567605\\
4.3052	-0.0567764\\
4.3078	-0.056793\\
4.3103	-0.0568089\\
4.3129	-0.0568254\\
4.3155	-0.0568418\\
4.318	-0.0568576\\
4.3206	-0.0568739\\
4.3232	-0.0568903\\
4.3257	-0.0569059\\
4.3283	-0.0569221\\
4.3309	-0.0569383\\
4.3334	-0.0569538\\
4.336	-0.0569699\\
4.3385	-0.0569853\\
4.3411	-0.0570013\\
4.3437	-0.0570173\\
4.3462	-0.0570325\\
4.3488	-0.0570483\\
4.3514	-0.0570641\\
4.3539	-0.0570792\\
4.3565	-0.0570948\\
4.3591	-0.0571104\\
4.3616	-0.0571253\\
4.3642	-0.0571407\\
4.3667	-0.0571555\\
4.3693	-0.0571708\\
4.3719	-0.0571861\\
4.3744	-0.0572007\\
4.377	-0.0572158\\
4.3796	-0.0572308\\
4.3821	-0.0572452\\
4.3847	-0.0572601\\
4.3873	-0.0572749\\
4.3898	-0.057289\\
4.3924	-0.0573037\\
4.3949	-0.0573177\\
4.3975	-0.0573321\\
4.4001	-0.0573465\\
4.4026	-0.0573603\\
4.4052	-0.0573745\\
4.4078	-0.0573887\\
4.4103	-0.0574022\\
4.4129	-0.0574162\\
4.4155	-0.0574301\\
4.418	-0.0574433\\
4.4206	-0.057457\\
4.4231	-0.0574701\\
4.4257	-0.0574836\\
4.4283	-0.057497\\
4.4308	-0.0575098\\
4.4334	-0.0575231\\
4.436	-0.0575362\\
4.4385	-0.0575487\\
4.4411	-0.0575617\\
4.4437	-0.0575745\\
4.4462	-0.0575867\\
4.4488	-0.0575994\\
4.4514	-0.0576119\\
4.4539	-0.0576238\\
4.4565	-0.0576361\\
4.459	-0.0576479\\
4.4616	-0.05766\\
4.4642	-0.0576719\\
4.4667	-0.0576834\\
4.4693	-0.0576951\\
4.4719	-0.0577068\\
4.4744	-0.0577179\\
4.477	-0.0577293\\
4.4796	-0.0577407\\
4.4821	-0.0577514\\
4.4847	-0.0577625\\
4.4872	-0.0577731\\
4.4898	-0.057784\\
4.4924	-0.0577947\\
4.4949	-0.0578049\\
4.4975	-0.0578155\\
4.5001	-0.0578258\\
4.5026	-0.0578357\\
4.5052	-0.0578459\\
4.5078	-0.0578559\\
4.5103	-0.0578654\\
4.5129	-0.0578752\\
4.5154	-0.0578845\\
4.518	-0.0578941\\
4.5206	-0.0579035\\
4.5231	-0.0579124\\
4.5257	-0.0579216\\
4.5283	-0.0579307\\
4.5308	-0.0579392\\
4.5334	-0.057948\\
4.536	-0.0579567\\
4.5385	-0.0579649\\
4.5411	-0.0579733\\
4.5436	-0.0579813\\
4.5462	-0.0579895\\
4.5488	-0.0579975\\
4.5513	-0.0580051\\
4.5539	-0.0580129\\
4.5565	-0.0580205\\
4.559	-0.0580277\\
4.5616	-0.0580351\\
4.5642	-0.0580424\\
4.5667	-0.0580492\\
4.5693	-0.0580562\\
4.5718	-0.0580628\\
4.5744	-0.0580695\\
4.577	-0.0580761\\
4.5795	-0.0580823\\
4.5821	-0.0580886\\
4.5847	-0.0580948\\
4.5872	-0.0581006\\
4.5898	-0.0581066\\
4.5924	-0.0581123\\
4.5949	-0.0581178\\
4.5975	-0.0581233\\
4.6001	-0.0581287\\
4.6026	-0.0581337\\
4.6052	-0.0581388\\
4.6077	-0.0581436\\
4.6103	-0.0581485\\
4.6129	-0.0581532\\
4.6154	-0.0581576\\
4.618	-0.058162\\
4.6206	-0.0581663\\
4.6231	-0.0581703\\
4.6257	-0.0581744\\
4.6283	-0.0581783\\
4.6308	-0.0581819\\
4.6334	-0.0581855\\
4.6359	-0.0581889\\
4.6385	-0.0581922\\
4.6411	-0.0581955\\
4.6436	-0.0581984\\
4.6462	-0.0582014\\
4.6488	-0.0582042\\
4.6513	-0.0582068\\
4.6539	-0.0582093\\
4.6565	-0.0582117\\
4.659	-0.0582139\\
4.6616	-0.0582161\\
4.6641	-0.058218\\
4.6667	-0.0582199\\
4.6693	-0.0582216\\
4.6718	-0.0582231\\
4.6744	-0.0582246\\
4.677	-0.0582259\\
4.6795	-0.0582271\\
4.6821	-0.0582282\\
4.6847	-0.0582291\\
4.6872	-0.0582299\\
4.6898	-0.0582306\\
4.6923	-0.0582311\\
4.6949	-0.0582315\\
4.6975	-0.0582318\\
4.7	-0.0582319\\
4.7026	-0.0582319\\
4.7052	-0.0582318\\
4.7077	-0.0582316\\
4.7103	-0.0582312\\
4.7129	-0.0582307\\
4.7154	-0.0582301\\
4.718	-0.0582293\\
4.7205	-0.0582285\\
4.7231	-0.0582274\\
4.7257	-0.0582263\\
4.7282	-0.0582251\\
4.7308	-0.0582237\\
4.7334	-0.0582221\\
4.7359	-0.0582205\\
4.7385	-0.0582187\\
4.7411	-0.0582168\\
4.7436	-0.0582149\\
4.7462	-0.0582127\\
4.7487	-0.0582105\\
4.7513	-0.0582081\\
4.7539	-0.0582056\\
4.7564	-0.058203\\
4.759	-0.0582003\\
4.7616	-0.0581974\\
4.7641	-0.0581945\\
4.7667	-0.0581913\\
4.7693	-0.0581881\\
4.7718	-0.0581848\\
4.7744	-0.0581813\\
4.777	-0.0581777\\
4.7795	-0.0581741\\
4.7821	-0.0581703\\
4.7846	-0.0581664\\
4.7872	-0.0581624\\
4.7898	-0.0581582\\
4.7923	-0.058154\\
4.7949	-0.0581496\\
4.7975	-0.058145\\
4.8	-0.0581406\\
4.8026	-0.0581358\\
4.8052	-0.0581309\\
4.8077	-0.0581261\\
4.8103	-0.058121\\
4.8128	-0.058116\\
4.8154	-0.0581107\\
4.818	-0.0581052\\
4.8205	-0.0580999\\
4.8231	-0.0580942\\
4.8257	-0.0580885\\
4.8282	-0.0580828\\
4.8308	-0.0580769\\
4.8334	-0.0580708\\
4.8359	-0.0580649\\
4.8385	-0.0580586\\
4.841	-0.0580524\\
4.8436	-0.058046\\
4.8462	-0.0580394\\
4.8487	-0.0580329\\
4.8513	-0.0580261\\
4.8539	-0.0580193\\
4.8564	-0.0580125\\
4.859	-0.0580054\\
4.8616	-0.0579983\\
4.8641	-0.0579913\\
4.8667	-0.0579839\\
4.8692	-0.0579767\\
4.8718	-0.0579691\\
4.8744	-0.0579615\\
4.8769	-0.057954\\
4.8795	-0.0579462\\
4.8821	-0.0579383\\
4.8846	-0.0579305\\
4.8872	-0.0579224\\
4.8898	-0.0579142\\
4.8923	-0.0579062\\
4.8949	-0.0578979\\
4.8974	-0.0578897\\
4.9	-0.0578812\\
4.9026	-0.0578725\\
4.9051	-0.0578642\\
4.9077	-0.0578553\\
4.9103	-0.0578465\\
4.9128	-0.0578378\\
4.9154	-0.0578288\\
4.918	-0.0578197\\
4.9205	-0.0578108\\
4.9231	-0.0578015\\
4.9257	-0.0577921\\
4.9282	-0.0577831\\
4.9308	-0.0577735\\
4.9333	-0.0577643\\
4.9359	-0.0577547\\
4.9385	-0.0577449\\
4.941	-0.0577355\\
4.9436	-0.0577256\\
4.9462	-0.0577156\\
4.9487	-0.057706\\
4.9513	-0.0576959\\
4.9539	-0.0576857\\
4.9564	-0.0576759\\
4.959	-0.0576656\\
4.9615	-0.0576556\\
4.9641	-0.0576451\\
4.9667	-0.0576346\\
4.9692	-0.0576245\\
4.9718	-0.0576139\\
4.9744	-0.0576032\\
4.9769	-0.0575928\\
4.9795	-0.057582\\
4.9821	-0.0575712\\
4.9846	-0.0575606\\
4.9872	-0.0575497\\
4.9897	-0.057539\\
4.9923	-0.0575279\\
4.9949	-0.0575168\\
4.9974	-0.057506\\
5	-0.0574947\\
};
\addplot [color=mycolor16, line width=1.0pt, forget plot]
  table[row sep=crcr]{%
-0.125	-9.11394e-08\\
-0.1224	-9.30158e-08\\
-0.1199	-9.48197e-08\\
-0.1173	-9.66954e-08\\
-0.1147	-9.85709e-08\\
-0.1122	-1.00374e-07\\
-0.1096	-1.02249e-07\\
-0.1071	-1.04051e-07\\
-0.1045	-1.05925e-07\\
-0.1019	-1.07799e-07\\
-0.0994	-1.096e-07\\
-0.0968	-1.11473e-07\\
-0.0942	-1.13346e-07\\
-0.0917	-1.15147e-07\\
-0.0891	-1.17019e-07\\
-0.0865	-1.1889e-07\\
-0.084	-1.2069e-07\\
-0.0814	-1.22561e-07\\
-0.0789	-1.2436e-07\\
-0.0763	-1.2623e-07\\
-0.0737	-1.28101e-07\\
-0.0712	-1.29898e-07\\
-0.0686	-1.31768e-07\\
-0.066	-1.33637e-07\\
-0.0635	-1.35434e-07\\
-0.0609	-1.37302e-07\\
-0.0583	-1.3917e-07\\
-0.0558	-1.40966e-07\\
-0.0532	-1.42834e-07\\
-0.0507	-1.44629e-07\\
-0.0481	-1.46496e-07\\
-0.0455	-1.48363e-07\\
-0.043	-1.50157e-07\\
-0.0404	-1.52023e-07\\
-0.0378	-1.53889e-07\\
-0.0353	-1.55682e-07\\
-0.0327	-1.57547e-07\\
-0.0301	-1.59412e-07\\
-0.0276	-1.61204e-07\\
-0.025	-1.63068e-07\\
-0.0224	-1.64932e-07\\
-0.0199	-1.66723e-07\\
-0.0173	-1.68586e-07\\
-0.0148	-1.70377e-07\\
-0.0122	-1.7224e-07\\
-0.0096	-1.74102e-07\\
-0.0071	-1.75892e-07\\
-0.0045	-1.77753e-07\\
-0.0019	-1.79614e-07\\
0.0006	-1.81403e-07\\
0.0032	0.000582595\\
0.0058	0.00247862\\
0.0083	0.00436924\\
0.0109	0.00638686\\
0.0134	0.00836554\\
0.016	0.0103653\\
0.0186	0.0122193\\
0.0211	0.0138502\\
0.0237	0.0154326\\
0.0263	0.0169535\\
0.0288	0.018381\\
0.0314	0.0198186\\
0.034	0.0211843\\
0.0365	0.0224178\\
0.0391	0.0236251\\
0.0416	0.0247303\\
0.0442	0.025836\\
0.0468	0.0269\\
0.0493	0.0278778\\
0.0519	0.0288414\\
0.0545	0.0297509\\
0.057	0.0305811\\
0.0596	0.0314074\\
0.0622	0.0322008\\
0.0647	0.0329321\\
0.0673	0.0336565\\
0.0698	0.0343168\\
0.0724	0.0349691\\
0.075	0.0355924\\
0.0775	0.0361695\\
0.0801	0.036748\\
0.0827	0.0373026\\
0.0852	0.037811\\
0.0878	0.0383146\\
0.0904	0.0387965\\
0.0929	0.0392441\\
0.0955	0.0396958\\
0.098	0.0401163\\
0.1006	0.0405367\\
0.1032	0.0409384\\
0.1057	0.0413089\\
0.1083	0.0416816\\
0.1109	0.0420449\\
0.1134	0.0423858\\
0.116	0.0427295\\
0.1186	0.0430596\\
0.1211	0.0433642\\
0.1237	0.0436702\\
0.1263	0.0439685\\
0.1288	0.0442496\\
0.1314	0.0445348\\
0.1339	0.0447998\\
0.1365	0.045064\\
0.1391	0.0453177\\
0.1416	0.0455544\\
0.1442	0.0457954\\
0.1468	0.046031\\
0.1493	0.0462505\\
0.1519	0.0464691\\
0.1545	0.0466774\\
0.157	0.0468698\\
0.1596	0.0470642\\
0.1621	0.0472467\\
0.1647	0.0474306\\
0.1673	0.0476063\\
0.1698	0.0477659\\
0.1724	0.0479229\\
0.175	0.0480733\\
0.1775	0.0482132\\
0.1801	0.0483539\\
0.1827	0.0484878\\
0.1852	0.0486083\\
0.1878	0.0487245\\
0.1903	0.0488292\\
0.1929	0.0489329\\
0.1955	0.0490323\\
0.198	0.0491229\\
0.2006	0.04921\\
0.2032	0.0492887\\
0.2057	0.0493572\\
0.2083	0.0494229\\
0.2109	0.0494847\\
0.2134	0.0495406\\
0.216	0.0495935\\
0.2185	0.0496378\\
0.2211	0.049677\\
0.2237	0.0497107\\
0.2262	0.0497397\\
0.2288	0.0497674\\
0.2314	0.0497917\\
0.2339	0.0498106\\
0.2365	0.0498245\\
0.2391	0.0498334\\
0.2416	0.049839\\
0.2442	0.0498433\\
0.2467	0.0498459\\
0.2493	0.049846\\
0.2519	0.0498421\\
0.2544	0.0498344\\
0.257	0.0498238\\
0.2596	0.0498121\\
0.2621	0.0498007\\
0.2647	0.049788\\
0.2673	0.0497732\\
0.2698	0.0497562\\
0.2724	0.0497363\\
0.2749	0.0497164\\
0.2775	0.0496963\\
0.2801	0.0496767\\
0.2826	0.0496573\\
0.2852	0.0496356\\
0.2878	0.049612\\
0.2903	0.0495887\\
0.2929	0.0495651\\
0.2955	0.0495427\\
0.298	0.0495218\\
0.3006	0.0494994\\
0.3032	0.0494758\\
0.3057	0.0494522\\
0.3083	0.0494279\\
0.3108	0.0494057\\
0.3134	0.0493838\\
0.316	0.049362\\
0.3185	0.0493401\\
0.3211	0.0493163\\
0.3237	0.049292\\
0.3262	0.0492693\\
0.3288	0.0492467\\
0.3314	0.0492245\\
0.3339	0.0492025\\
0.3365	0.0491781\\
0.339	0.0491536\\
0.3416	0.0491278\\
0.3442	0.0491024\\
0.3467	0.0490782\\
0.3493	0.0490523\\
0.3519	0.0490247\\
0.3544	0.0489963\\
0.357	0.0489654\\
0.3596	0.0489339\\
0.3621	0.0489034\\
0.3647	0.0488706\\
0.3672	0.0488374\\
0.3698	0.0488002\\
0.3724	0.0487606\\
0.3749	0.0487209\\
0.3775	0.0486785\\
0.3801	0.0486347\\
0.3826	0.0485906\\
0.3852	0.0485417\\
0.3878	0.0484896\\
0.3903	0.0484369\\
0.3929	0.04838\\
0.3954	0.0483235\\
0.398	0.0482624\\
0.4006	0.048198\\
0.4031	0.0481325\\
0.4057	0.0480608\\
0.4083	0.0479861\\
0.4108	0.0479119\\
0.4134	0.0478321\\
0.416	0.047749\\
0.4185	0.0476652\\
0.4211	0.0475739\\
0.4236	0.0474826\\
0.4262	0.0473846\\
0.4288	0.0472836\\
0.4313	0.0471833\\
0.4339	0.047075\\
0.4365	0.0469622\\
0.439	0.0468498\\
0.4416	0.0467292\\
0.4442	0.0466054\\
0.4467	0.0464831\\
0.4493	0.0463522\\
0.4519	0.0462169\\
0.4544	0.0460824\\
0.457	0.0459386\\
0.4595	0.045797\\
0.4621	0.0456464\\
0.4647	0.0454921\\
0.4672	0.0453399\\
0.4698	0.0451771\\
0.4724	0.04501\\
0.4749	0.0448457\\
0.4775	0.0446716\\
0.4801	0.044494\\
0.4826	0.0443196\\
0.4852	0.044134\\
0.4877	0.0439514\\
0.4903	0.0437576\\
0.4929	0.0435604\\
0.4954	0.0433676\\
0.498	0.0431637\\
0.5006	0.0429558\\
0.5031	0.0427519\\
0.5057	0.0425359\\
0.5083	0.0423163\\
0.5108	0.0421022\\
0.5134	0.0418762\\
0.5159	0.0416556\\
0.5185	0.0414222\\
0.5211	0.0411849\\
0.5236	0.0409532\\
0.5262	0.040709\\
0.5288	0.0404618\\
0.5313	0.0402209\\
0.5339	0.0399668\\
0.5365	0.0397089\\
0.539	0.0394573\\
0.5416	0.0391924\\
0.5441	0.0389348\\
0.5467	0.0386639\\
0.5493	0.0383897\\
0.5518	0.0381225\\
0.5544	0.037841\\
0.557	0.037556\\
0.5595	0.037279\\
0.5621	0.0369881\\
0.5647	0.0366941\\
0.5672	0.036408\\
0.5698	0.036107\\
0.5723	0.0358142\\
0.5749	0.0355066\\
0.5775	0.035196\\
0.58	0.0348945\\
0.5826	0.0345778\\
0.5852	0.0342576\\
0.5877	0.0339463\\
0.5903	0.0336194\\
0.5929	0.0332894\\
0.5954	0.0329693\\
0.598	0.0326335\\
0.6006	0.0322943\\
0.6031	0.0319648\\
0.6057	0.0316189\\
0.6082	0.0312833\\
0.6108	0.0309314\\
0.6134	0.0305766\\
0.6159	0.0302325\\
0.6185	0.0298713\\
0.6211	0.0295067\\
0.6236	0.0291531\\
0.6262	0.0287825\\
0.6288	0.0284091\\
0.6313	0.0280472\\
0.6339	0.0276677\\
0.6364	0.0272997\\
0.639	0.0269138\\
0.6416	0.026525\\
0.6441	0.0261485\\
0.6467	0.0257542\\
0.6493	0.0253569\\
0.6518	0.0249719\\
0.6544	0.0245685\\
0.657	0.0241621\\
0.6595	0.0237688\\
0.6621	0.0233572\\
0.6646	0.0229588\\
0.6672	0.0225417\\
0.6698	0.0221216\\
0.6723	0.0217149\\
0.6749	0.0212894\\
0.6775	0.0208614\\
0.68	0.0204476\\
0.6826	0.0200146\\
0.6852	0.0195789\\
0.6877	0.0191574\\
0.6903	0.0187166\\
0.6928	0.0182905\\
0.6954	0.0178453\\
0.698	0.0173977\\
0.7005	0.0169651\\
0.7031	0.0165127\\
0.7057	0.016058\\
0.7082	0.0156188\\
0.7108	0.01516\\
0.7134	0.0146993\\
0.7159	0.0142543\\
0.7185	0.0137895\\
0.721	0.0133405\\
0.7236	0.0128716\\
0.7262	0.012401\\
0.7287	0.011947\\
0.7313	0.0114731\\
0.7339	0.0109974\\
0.7364	0.0105383\\
0.739	0.0100592\\
0.7416	0.00957852\\
0.7441	0.00911505\\
0.7467	0.00863173\\
0.7492	0.00816568\\
0.7518	0.00767955\\
0.7544	0.00719198\\
0.7569	0.00672195\\
0.7595	0.00623204\\
0.7621	0.00574112\\
0.7646	0.00526812\\
0.7672	0.00477511\\
0.7698	0.00428099\\
0.7723	0.00380492\\
0.7749	0.00330896\\
0.7775	0.0028123\\
0.78	0.0023341\\
0.7826	0.00183606\\
0.7851	0.00135645\\
0.7877	0.000856928\\
0.7903	0.000356808\\
0.7928	-0.000124498\\
0.7954	-0.000625407\\
0.798	-0.00112669\\
0.8005	-0.00160908\\
0.8031	-0.0021112\\
0.8057	-0.00261366\\
0.8082	-0.00309698\\
0.8108	-0.00359972\\
0.8133	-0.00408316\\
0.8159	-0.00458601\\
0.8185	-0.00508898\\
0.821	-0.00557269\\
0.8236	-0.0060757\\
0.8262	-0.00657855\\
0.8287	-0.00706183\\
0.8313	-0.0075642\\
0.8339	-0.00806637\\
0.8364	-0.00854905\\
0.839	-0.00905076\\
0.8415	-0.00953281\\
0.8441	-0.0100337\\
0.8467	-0.010534\\
0.8492	-0.0110146\\
0.8518	-0.0115139\\
0.8544	-0.0120127\\
0.8569	-0.0124918\\
0.8595	-0.0129893\\
0.8621	-0.0134861\\
0.8646	-0.013963\\
0.8672	-0.0144582\\
0.8697	-0.0149336\\
0.8723	-0.0154273\\
0.8749	-0.01592\\
0.8774	-0.0163929\\
0.88	-0.0168835\\
0.8826	-0.0173732\\
0.8851	-0.0178431\\
0.8877	-0.0183307\\
0.8903	-0.0188173\\
0.8928	-0.0192839\\
0.8954	-0.019768\\
0.8979	-0.0202322\\
0.9005	-0.0207137\\
0.9031	-0.021194\\
0.9056	-0.0216545\\
0.9082	-0.0221321\\
0.9108	-0.0226082\\
0.9133	-0.0230645\\
0.9159	-0.0235376\\
0.9185	-0.0240092\\
0.921	-0.0244612\\
0.9236	-0.0249297\\
0.9262	-0.0253965\\
0.9287	-0.0258438\\
0.9313	-0.0263072\\
0.9338	-0.0267511\\
0.9364	-0.0272111\\
0.939	-0.0276693\\
0.9415	-0.0281081\\
0.9441	-0.0285626\\
0.9467	-0.0290152\\
0.9492	-0.0294484\\
0.9518	-0.0298971\\
0.9544	-0.0303439\\
0.9569	-0.0307715\\
0.9595	-0.0312142\\
0.962	-0.0316378\\
0.9646	-0.0320762\\
0.9672	-0.0325125\\
0.9697	-0.03293\\
0.9723	-0.0333621\\
0.9749	-0.0337919\\
0.9774	-0.0342029\\
0.98	-0.0346282\\
0.9826	-0.0350511\\
0.9851	-0.0354555\\
0.9877	-0.0358738\\
0.9902	-0.0362738\\
0.9928	-0.0366873\\
0.9954	-0.0370983\\
0.9979	-0.0374912\\
1.0005	-0.0378974\\
1.0031	-0.038301\\
1.0056	-0.0386868\\
1.0082	-0.0390854\\
1.0108	-0.0394815\\
1.0133	-0.0398598\\
1.0159	-0.0402506\\
1.0184	-0.040624\\
1.021	-0.0410096\\
1.0236	-0.0413926\\
1.0261	-0.0417583\\
1.0287	-0.0421359\\
1.0313	-0.0425107\\
1.0338	-0.0428686\\
1.0364	-0.043238\\
1.039	-0.0436047\\
1.0415	-0.0439546\\
1.0441	-0.0443157\\
1.0466	-0.0446602\\
1.0492	-0.0450157\\
1.0518	-0.0453684\\
1.0543	-0.0457048\\
1.0569	-0.0460519\\
1.0595	-0.046396\\
1.062	-0.0467241\\
1.0646	-0.0470625\\
1.0672	-0.0473979\\
1.0697	-0.0477177\\
1.0723	-0.0480474\\
1.0748	-0.0483616\\
1.0774	-0.0486855\\
1.08	-0.0490064\\
1.0825	-0.0493121\\
1.0851	-0.0496272\\
1.0877	-0.0499393\\
1.0902	-0.0502365\\
1.0928	-0.0505427\\
1.0954	-0.0508459\\
1.0979	-0.0511346\\
1.1005	-0.0514318\\
1.1031	-0.0517261\\
1.1056	-0.0520062\\
1.1082	-0.0522946\\
1.1107	-0.052569\\
1.1133	-0.0528514\\
1.1159	-0.0531308\\
1.1184	-0.0533966\\
1.121	-0.05367\\
1.1236	-0.0539405\\
1.1261	-0.0541977\\
1.1287	-0.0544622\\
1.1313	-0.0547237\\
1.1338	-0.0549723\\
1.1364	-0.0552278\\
1.1389	-0.0554708\\
1.1415	-0.0557204\\
1.1441	-0.0559671\\
1.1466	-0.0562014\\
1.1492	-0.0564422\\
1.1518	-0.0566799\\
1.1543	-0.0569058\\
1.1569	-0.0571377\\
1.1595	-0.0573666\\
1.162	-0.057584\\
1.1646	-0.0578071\\
1.1671	-0.0580188\\
1.1697	-0.0582361\\
1.1723	-0.0584504\\
1.1748	-0.0586538\\
1.1774	-0.0588624\\
1.18	-0.059068\\
1.1825	-0.059263\\
1.1851	-0.0594629\\
1.1877	-0.05966\\
1.1902	-0.0598467\\
1.1928	-0.0600381\\
1.1953	-0.0602194\\
1.1979	-0.0604051\\
1.2005	-0.060588\\
1.203	-0.0607612\\
1.2056	-0.0609385\\
1.2082	-0.061113\\
1.2107	-0.0612781\\
1.2133	-0.0614471\\
1.2159	-0.0616133\\
1.2184	-0.0617705\\
1.221	-0.0619312\\
1.2235	-0.0620832\\
1.2261	-0.0622386\\
1.2287	-0.0623912\\
1.2312	-0.0625354\\
1.2338	-0.0626827\\
1.2364	-0.0628273\\
1.2389	-0.0629639\\
1.2415	-0.0631033\\
1.2441	-0.06324\\
1.2466	-0.063369\\
1.2492	-0.0635006\\
1.2518	-0.0636296\\
1.2543	-0.0637512\\
1.2569	-0.0638752\\
1.2594	-0.0639919\\
1.262	-0.0641109\\
1.2646	-0.0642273\\
1.2671	-0.0643368\\
1.2697	-0.0644483\\
1.2723	-0.0645574\\
1.2748	-0.0646599\\
1.2774	-0.0647641\\
1.28	-0.0648659\\
1.2825	-0.0649615\\
1.2851	-0.0650586\\
1.2876	-0.0651497\\
1.2902	-0.0652422\\
1.2928	-0.0653323\\
1.2953	-0.0654168\\
1.2979	-0.0655024\\
1.3005	-0.0655857\\
1.303	-0.0656637\\
1.3056	-0.0657427\\
1.3082	-0.0658194\\
1.3107	-0.065891\\
1.3133	-0.0659634\\
1.3158	-0.066031\\
1.3184	-0.0660992\\
1.321	-0.0661652\\
1.3235	-0.0662268\\
1.3261	-0.0662887\\
1.3287	-0.0663486\\
1.3312	-0.0664043\\
1.3338	-0.0664602\\
1.3364	-0.0665141\\
1.3389	-0.0665641\\
1.3415	-0.0666141\\
1.344	-0.0666604\\
1.3466	-0.0667067\\
1.3492	-0.0667511\\
1.3517	-0.066792\\
1.3543	-0.0668328\\
1.3569	-0.0668717\\
1.3594	-0.0669074\\
1.362	-0.0669428\\
1.3646	-0.0669764\\
1.3671	-0.0670071\\
1.3697	-0.0670373\\
1.3722	-0.0670648\\
1.3748	-0.0670918\\
1.3774	-0.0671171\\
1.3799	-0.0671398\\
1.3825	-0.0671619\\
1.3851	-0.0671825\\
1.3876	-0.0672007\\
1.3902	-0.0672182\\
1.3928	-0.0672342\\
1.3953	-0.0672481\\
1.3979	-0.0672612\\
1.4005	-0.0672728\\
1.403	-0.0672826\\
1.4056	-0.0672914\\
1.4081	-0.0672986\\
1.4107	-0.0673047\\
1.4133	-0.0673095\\
1.4158	-0.0673129\\
1.4184	-0.0673151\\
1.421	-0.0673161\\
1.4235	-0.0673159\\
1.4261	-0.0673144\\
1.4287	-0.0673117\\
1.4312	-0.0673081\\
1.4338	-0.0673031\\
1.4363	-0.0672973\\
1.4389	-0.0672901\\
1.4415	-0.0672819\\
1.444	-0.0672729\\
1.4466	-0.0672626\\
1.4492	-0.0672512\\
1.4517	-0.0672393\\
1.4543	-0.067226\\
1.4569	-0.0672118\\
1.4594	-0.0671972\\
1.462	-0.0671811\\
1.4645	-0.0671647\\
1.4671	-0.0671469\\
1.4697	-0.0671282\\
1.4722	-0.0671095\\
1.4748	-0.0670892\\
1.4774	-0.0670681\\
1.4799	-0.0670471\\
1.4825	-0.0670246\\
1.4851	-0.0670013\\
1.4876	-0.0669782\\
1.4902	-0.0669536\\
1.4927	-0.0669292\\
1.4953	-0.0669033\\
1.4979	-0.0668767\\
1.5004	-0.0668506\\
1.503	-0.0668229\\
1.5056	-0.0667946\\
1.5081	-0.0667669\\
1.5107	-0.0667375\\
1.5133	-0.0667077\\
1.5158	-0.0666785\\
1.5184	-0.0666478\\
1.5209	-0.0666178\\
1.5235	-0.0665861\\
1.5261	-0.0665541\\
1.5286	-0.0665229\\
1.5312	-0.0664901\\
1.5338	-0.066457\\
1.5363	-0.0664248\\
1.5389	-0.066391\\
1.5415	-0.0663569\\
1.544	-0.0663239\\
1.5466	-0.0662893\\
1.5491	-0.0662557\\
1.5517	-0.0662207\\
1.5543	-0.0661854\\
1.5568	-0.0661512\\
1.5594	-0.0661156\\
1.562	-0.0660797\\
1.5645	-0.0660452\\
1.5671	-0.0660091\\
1.5697	-0.0659728\\
1.5722	-0.0659379\\
1.5748	-0.0659015\\
1.5774	-0.0658651\\
1.5799	-0.06583\\
1.5825	-0.0657934\\
1.585	-0.0657583\\
1.5876	-0.0657217\\
1.5902	-0.0656851\\
1.5927	-0.06565\\
1.5953	-0.0656134\\
1.5979	-0.065577\\
1.6004	-0.0655419\\
1.603	-0.0655056\\
1.6056	-0.0654694\\
1.6081	-0.0654346\\
1.6107	-0.0653985\\
1.6132	-0.065364\\
1.6158	-0.0653282\\
1.6184	-0.0652926\\
1.6209	-0.0652585\\
1.6235	-0.0652232\\
1.6261	-0.065188\\
1.6286	-0.0651545\\
1.6312	-0.0651197\\
1.6338	-0.0650852\\
1.6363	-0.0650522\\
1.6389	-0.0650182\\
1.6414	-0.0649857\\
1.644	-0.0649521\\
1.6466	-0.0649188\\
1.6491	-0.064887\\
1.6517	-0.0648543\\
1.6543	-0.0648218\\
1.6568	-0.0647909\\
1.6594	-0.0647591\\
1.662	-0.0647275\\
1.6645	-0.0646975\\
1.6671	-0.0646666\\
1.6696	-0.0646373\\
1.6722	-0.064607\\
1.6748	-0.0645772\\
1.6773	-0.0645488\\
1.6799	-0.0645197\\
1.6825	-0.0644909\\
1.685	-0.0644637\\
1.6876	-0.0644357\\
1.6902	-0.064408\\
1.6927	-0.0643819\\
1.6953	-0.0643551\\
1.6978	-0.0643297\\
1.7004	-0.0643036\\
1.703	-0.064278\\
1.7055	-0.0642538\\
1.7081	-0.0642291\\
1.7107	-0.0642048\\
1.7132	-0.0641818\\
1.7158	-0.0641583\\
1.7184	-0.0641353\\
1.7209	-0.0641136\\
1.7235	-0.0640914\\
1.7261	-0.0640697\\
1.7286	-0.0640493\\
1.7312	-0.0640285\\
1.7337	-0.0640089\\
1.7363	-0.063989\\
1.7389	-0.0639695\\
1.7414	-0.0639513\\
1.744	-0.0639327\\
1.7466	-0.0639147\\
1.7491	-0.0638977\\
1.7517	-0.0638806\\
1.7543	-0.0638639\\
1.7568	-0.0638483\\
1.7594	-0.0638325\\
1.7619	-0.0638178\\
1.7645	-0.063803\\
1.7671	-0.0637886\\
1.7696	-0.0637752\\
1.7722	-0.0637618\\
1.7748	-0.0637488\\
1.7773	-0.0637367\\
1.7799	-0.0637246\\
1.7825	-0.063713\\
1.785	-0.0637023\\
1.7876	-0.0636916\\
1.7901	-0.0636817\\
1.7927	-0.0636719\\
1.7953	-0.0636625\\
1.7978	-0.0636539\\
1.8004	-0.0636455\\
1.803	-0.0636374\\
1.8055	-0.0636301\\
1.8081	-0.063623\\
1.8107	-0.0636163\\
1.8132	-0.0636102\\
1.8158	-0.0636043\\
1.8183	-0.0635991\\
1.8209	-0.0635941\\
1.8235	-0.0635895\\
1.826	-0.0635854\\
1.8286	-0.0635817\\
1.8312	-0.0635783\\
1.8337	-0.0635754\\
1.8363	-0.0635728\\
1.8389	-0.0635707\\
1.8414	-0.0635689\\
1.844	-0.0635675\\
1.8465	-0.0635666\\
1.8491	-0.0635659\\
1.8517	-0.0635656\\
1.8542	-0.0635657\\
1.8568	-0.0635662\\
1.8594	-0.063567\\
1.8619	-0.0635681\\
1.8645	-0.0635696\\
1.8671	-0.0635715\\
1.8696	-0.0635736\\
1.8722	-0.0635761\\
1.8747	-0.0635789\\
1.8773	-0.063582\\
1.8799	-0.0635855\\
1.8824	-0.0635892\\
1.885	-0.0635933\\
1.8876	-0.0635977\\
1.8901	-0.0636022\\
1.8927	-0.0636072\\
1.8953	-0.0636125\\
1.8978	-0.0636178\\
1.9004	-0.0636237\\
1.903	-0.0636298\\
1.9055	-0.0636359\\
1.9081	-0.0636425\\
1.9106	-0.0636491\\
1.9132	-0.0636562\\
1.9158	-0.0636635\\
1.9183	-0.0636708\\
1.9209	-0.0636786\\
1.9235	-0.0636866\\
1.926	-0.0636945\\
1.9286	-0.0637029\\
1.9312	-0.0637115\\
1.9337	-0.06372\\
1.9363	-0.063729\\
1.9388	-0.0637379\\
1.9414	-0.0637472\\
1.944	-0.0637568\\
1.9465	-0.0637661\\
1.9491	-0.063776\\
1.9517	-0.063786\\
1.9542	-0.0637958\\
1.9568	-0.0638061\\
1.9594	-0.0638165\\
1.9619	-0.0638267\\
1.9645	-0.0638374\\
1.967	-0.0638477\\
1.9696	-0.0638587\\
1.9722	-0.0638697\\
1.9747	-0.0638804\\
1.9773	-0.0638916\\
1.9799	-0.0639029\\
1.9824	-0.0639139\\
1.985	-0.0639253\\
1.9876	-0.0639369\\
1.9901	-0.063948\\
1.9927	-0.0639597\\
1.9952	-0.0639709\\
1.9978	-0.0639827\\
2.0004	-0.0639945\\
2.0029	-0.0640059\\
2.0055	-0.0640178\\
2.0081	-0.0640297\\
2.0106	-0.0640412\\
2.0132	-0.0640531\\
2.0158	-0.0640651\\
2.0183	-0.0640766\\
2.0209	-0.0640886\\
2.0234	-0.0641001\\
2.026	-0.064112\\
2.0286	-0.064124\\
2.0311	-0.0641354\\
2.0337	-0.0641474\\
2.0363	-0.0641592\\
2.0388	-0.0641706\\
2.0414	-0.0641825\\
2.044	-0.0641942\\
2.0465	-0.0642055\\
2.0491	-0.0642172\\
2.0517	-0.0642288\\
2.0542	-0.06424\\
2.0568	-0.0642515\\
2.0593	-0.0642625\\
2.0619	-0.0642739\\
2.0645	-0.0642852\\
2.067	-0.064296\\
2.0696	-0.0643072\\
2.0722	-0.0643183\\
2.0747	-0.0643288\\
2.0773	-0.0643398\\
2.0799	-0.0643506\\
2.0824	-0.0643609\\
2.085	-0.0643715\\
2.0875	-0.0643816\\
2.0901	-0.064392\\
2.0927	-0.0644023\\
2.0952	-0.0644121\\
2.0978	-0.0644222\\
2.1004	-0.0644322\\
2.1029	-0.0644416\\
2.1055	-0.0644513\\
2.1081	-0.0644609\\
2.1106	-0.06447\\
2.1132	-0.0644794\\
2.1157	-0.0644883\\
2.1183	-0.0644973\\
2.1209	-0.0645063\\
2.1234	-0.0645148\\
2.126	-0.0645234\\
2.1286	-0.064532\\
2.1311	-0.06454\\
2.1337	-0.0645483\\
2.1363	-0.0645564\\
2.1388	-0.064564\\
2.1414	-0.0645719\\
2.1439	-0.0645792\\
2.1465	-0.0645867\\
2.1491	-0.0645941\\
2.1516	-0.064601\\
2.1542	-0.0646081\\
2.1568	-0.064615\\
2.1593	-0.0646215\\
2.1619	-0.0646281\\
2.1645	-0.0646345\\
2.167	-0.0646406\\
2.1696	-0.0646467\\
2.1721	-0.0646524\\
2.1747	-0.0646583\\
2.1773	-0.0646639\\
2.1798	-0.0646692\\
2.1824	-0.0646745\\
2.185	-0.0646797\\
2.1875	-0.0646845\\
2.1901	-0.0646894\\
2.1927	-0.0646941\\
2.1952	-0.0646985\\
2.1978	-0.0647029\\
2.2004	-0.0647071\\
2.2029	-0.064711\\
2.2055	-0.0647149\\
2.208	-0.0647185\\
2.2106	-0.0647221\\
2.2132	-0.0647256\\
2.2157	-0.0647288\\
2.2183	-0.0647319\\
2.2209	-0.0647349\\
2.2234	-0.0647376\\
2.226	-0.0647403\\
2.2286	-0.0647429\\
2.2311	-0.0647452\\
2.2337	-0.0647474\\
2.2362	-0.0647494\\
2.2388	-0.0647514\\
2.2414	-0.0647532\\
2.2439	-0.0647548\\
2.2465	-0.0647563\\
2.2491	-0.0647577\\
2.2516	-0.0647589\\
2.2542	-0.06476\\
2.2568	-0.064761\\
2.2593	-0.0647618\\
2.2619	-0.0647625\\
2.2644	-0.064763\\
2.267	-0.0647635\\
2.2696	-0.0647638\\
2.2721	-0.064764\\
2.2747	-0.064764\\
2.2773	-0.064764\\
2.2798	-0.0647638\\
2.2824	-0.0647635\\
2.285	-0.0647631\\
2.2875	-0.0647626\\
2.2901	-0.0647619\\
2.2926	-0.0647612\\
2.2952	-0.0647603\\
2.2978	-0.0647594\\
2.3003	-0.0647583\\
2.3029	-0.0647572\\
2.3055	-0.0647559\\
2.308	-0.0647545\\
2.3106	-0.0647531\\
2.3132	-0.0647515\\
2.3157	-0.0647499\\
2.3183	-0.0647481\\
2.3208	-0.0647463\\
2.3234	-0.0647444\\
2.326	-0.0647423\\
2.3285	-0.0647403\\
2.3311	-0.0647381\\
2.3337	-0.0647359\\
2.3362	-0.0647336\\
2.3388	-0.0647312\\
2.3414	-0.0647287\\
2.3439	-0.0647262\\
2.3465	-0.0647236\\
2.349	-0.064721\\
2.3516	-0.0647182\\
2.3542	-0.0647154\\
2.3567	-0.0647126\\
2.3593	-0.0647096\\
2.3619	-0.0647066\\
2.3644	-0.0647037\\
2.367	-0.0647006\\
2.3696	-0.0646974\\
2.3721	-0.0646943\\
2.3747	-0.064691\\
2.3773	-0.0646877\\
2.3798	-0.0646845\\
2.3824	-0.0646811\\
2.3849	-0.0646778\\
2.3875	-0.0646744\\
2.3901	-0.0646709\\
2.3926	-0.0646675\\
2.3952	-0.0646639\\
2.3978	-0.0646603\\
2.4003	-0.0646568\\
2.4029	-0.0646532\\
2.4055	-0.0646495\\
2.408	-0.064646\\
2.4106	-0.0646422\\
2.4131	-0.0646386\\
2.4157	-0.0646349\\
2.4183	-0.0646311\\
2.4208	-0.0646275\\
2.4234	-0.0646237\\
2.426	-0.0646199\\
2.4285	-0.0646162\\
2.4311	-0.0646124\\
2.4337	-0.0646085\\
2.4362	-0.0646048\\
2.4388	-0.064601\\
2.4413	-0.0645973\\
2.4439	-0.0645934\\
2.4465	-0.0645896\\
2.449	-0.0645858\\
2.4516	-0.064582\\
2.4542	-0.0645781\\
2.4567	-0.0645744\\
2.4593	-0.0645706\\
2.4619	-0.0645667\\
2.4644	-0.0645631\\
2.467	-0.0645592\\
2.4695	-0.0645556\\
2.4721	-0.0645517\\
2.4747	-0.0645479\\
2.4772	-0.0645443\\
2.4798	-0.0645405\\
2.4824	-0.0645368\\
2.4849	-0.0645331\\
2.4875	-0.0645294\\
2.4901	-0.0645257\\
2.4926	-0.0645221\\
2.4952	-0.0645184\\
2.4977	-0.0645149\\
2.5003	-0.0645112\\
2.5029	-0.0645075\\
2.5054	-0.064504\\
2.508	-0.0645004\\
2.5106	-0.0644968\\
2.5131	-0.0644934\\
2.5157	-0.0644898\\
2.5183	-0.0644862\\
2.5208	-0.0644828\\
2.5234	-0.0644793\\
2.526	-0.0644758\\
2.5285	-0.0644724\\
2.5311	-0.064469\\
2.5336	-0.0644656\\
2.5362	-0.0644622\\
2.5388	-0.0644588\\
2.5413	-0.0644555\\
2.5439	-0.0644521\\
2.5465	-0.0644487\\
2.549	-0.0644455\\
2.5516	-0.0644421\\
2.5542	-0.0644388\\
2.5567	-0.0644356\\
2.5593	-0.0644323\\
2.5618	-0.0644291\\
2.5644	-0.0644258\\
2.567	-0.0644225\\
2.5695	-0.0644194\\
2.5721	-0.0644161\\
2.5747	-0.0644129\\
2.5772	-0.0644097\\
2.5798	-0.0644065\\
2.5824	-0.0644033\\
2.5849	-0.0644001\\
2.5875	-0.0643969\\
2.59	-0.0643938\\
2.5926	-0.0643906\\
2.5952	-0.0643874\\
2.5977	-0.0643843\\
2.6003	-0.064381\\
2.6029	-0.0643778\\
2.6054	-0.0643747\\
2.608	-0.0643714\\
2.6106	-0.0643682\\
2.6131	-0.0643651\\
2.6157	-0.0643618\\
2.6182	-0.0643586\\
2.6208	-0.0643554\\
2.6234	-0.064352\\
2.6259	-0.0643489\\
2.6285	-0.0643455\\
2.6311	-0.0643422\\
2.6336	-0.0643389\\
2.6362	-0.0643355\\
2.6388	-0.0643321\\
2.6413	-0.0643288\\
2.6439	-0.0643253\\
2.6464	-0.0643219\\
2.649	-0.0643184\\
2.6516	-0.0643149\\
2.6541	-0.0643114\\
2.6567	-0.0643078\\
2.6593	-0.0643041\\
2.6618	-0.0643005\\
2.6644	-0.0642968\\
2.667	-0.064293\\
2.6695	-0.0642893\\
2.6721	-0.0642855\\
2.6746	-0.0642817\\
2.6772	-0.0642777\\
2.6798	-0.0642737\\
2.6823	-0.0642698\\
2.6849	-0.0642656\\
2.6875	-0.0642615\\
2.69	-0.0642574\\
2.6926	-0.0642531\\
2.6952	-0.0642487\\
2.6977	-0.0642444\\
2.7003	-0.0642399\\
2.7029	-0.0642354\\
2.7054	-0.0642309\\
2.708	-0.0642262\\
2.7105	-0.0642216\\
2.7131	-0.0642167\\
2.7157	-0.0642118\\
2.7182	-0.064207\\
2.7208	-0.0642019\\
2.7234	-0.0641967\\
2.7259	-0.0641916\\
2.7285	-0.0641862\\
2.7311	-0.0641808\\
2.7336	-0.0641754\\
2.7362	-0.0641698\\
2.7387	-0.0641643\\
2.7413	-0.0641585\\
2.7439	-0.0641525\\
2.7464	-0.0641467\\
2.749	-0.0641406\\
2.7516	-0.0641343\\
2.7541	-0.0641282\\
2.7567	-0.0641218\\
2.7593	-0.0641152\\
2.7618	-0.0641087\\
2.7644	-0.0641019\\
2.7669	-0.0640953\\
2.7695	-0.0640882\\
2.7721	-0.0640811\\
2.7746	-0.0640741\\
2.7772	-0.0640667\\
2.7798	-0.0640591\\
2.7823	-0.0640517\\
2.7849	-0.0640439\\
2.7875	-0.064036\\
2.79	-0.0640282\\
2.7926	-0.06402\\
2.7951	-0.064012\\
2.7977	-0.0640035\\
2.8003	-0.0639949\\
2.8028	-0.0639865\\
2.8054	-0.0639776\\
2.808	-0.0639685\\
2.8105	-0.0639597\\
2.8131	-0.0639503\\
2.8157	-0.0639408\\
2.8182	-0.0639315\\
2.8208	-0.0639217\\
2.8233	-0.0639121\\
2.8259	-0.063902\\
2.8285	-0.0638917\\
2.831	-0.0638817\\
2.8336	-0.0638711\\
2.8362	-0.0638603\\
2.8387	-0.0638498\\
2.8413	-0.0638387\\
2.8439	-0.0638274\\
2.8464	-0.0638164\\
2.849	-0.0638048\\
2.8516	-0.063793\\
2.8541	-0.0637815\\
2.8567	-0.0637694\\
2.8592	-0.0637576\\
2.8618	-0.0637451\\
2.8644	-0.0637325\\
2.8669	-0.0637201\\
2.8695	-0.0637071\\
2.8721	-0.0636939\\
2.8746	-0.0636811\\
2.8772	-0.0636675\\
2.8798	-0.0636538\\
2.8823	-0.0636404\\
2.8849	-0.0636263\\
2.8874	-0.0636126\\
2.89	-0.0635981\\
2.8926	-0.0635834\\
2.8951	-0.0635692\\
2.8977	-0.0635541\\
2.9003	-0.0635389\\
2.9028	-0.0635241\\
2.9054	-0.0635085\\
2.908	-0.0634927\\
2.9105	-0.0634773\\
2.9131	-0.0634611\\
2.9156	-0.0634454\\
2.9182	-0.0634288\\
2.9208	-0.063412\\
2.9233	-0.0633957\\
2.9259	-0.0633786\\
2.9285	-0.0633612\\
2.931	-0.0633444\\
2.9336	-0.0633266\\
2.9362	-0.0633087\\
2.9387	-0.0632913\\
2.9413	-0.0632729\\
2.9438	-0.0632551\\
2.9464	-0.0632364\\
2.949	-0.0632175\\
2.9515	-0.0631991\\
2.9541	-0.0631798\\
2.9567	-0.0631603\\
2.9592	-0.0631414\\
2.9618	-0.0631215\\
2.9644	-0.0631014\\
2.9669	-0.0630819\\
2.9695	-0.0630614\\
2.972	-0.0630416\\
2.9746	-0.0630207\\
2.9772	-0.0629996\\
2.9797	-0.0629792\\
2.9823	-0.0629578\\
2.9849	-0.0629361\\
2.9874	-0.0629151\\
2.99	-0.0628931\\
2.9926	-0.0628709\\
2.9951	-0.0628494\\
2.9977	-0.0628268\\
3.0003	-0.062804\\
3.0028	-0.0627819\\
3.0054	-0.0627587\\
3.0079	-0.0627363\\
3.0105	-0.0627127\\
3.0131	-0.062689\\
3.0156	-0.062666\\
3.0182	-0.0626419\\
3.0208	-0.0626177\\
3.0233	-0.0625942\\
3.0259	-0.0625695\\
3.0285	-0.0625447\\
3.031	-0.0625207\\
3.0336	-0.0624955\\
3.0361	-0.0624711\\
3.0387	-0.0624456\\
3.0413	-0.0624199\\
3.0438	-0.062395\\
3.0464	-0.0623689\\
3.049	-0.0623427\\
3.0515	-0.0623173\\
3.0541	-0.0622908\\
3.0567	-0.062264\\
3.0592	-0.0622382\\
3.0618	-0.0622111\\
3.0643	-0.0621849\\
3.0669	-0.0621576\\
3.0695	-0.06213\\
3.072	-0.0621034\\
3.0746	-0.0620755\\
3.0772	-0.0620475\\
3.0797	-0.0620204\\
3.0823	-0.0619921\\
3.0849	-0.0619636\\
3.0874	-0.0619361\\
3.09	-0.0619073\\
3.0925	-0.0618795\\
3.0951	-0.0618504\\
3.0977	-0.0618212\\
3.1002	-0.061793\\
3.1028	-0.0617635\\
3.1054	-0.0617338\\
3.1079	-0.0617052\\
3.1105	-0.0616753\\
3.1131	-0.0616453\\
3.1156	-0.0616163\\
3.1182	-0.061586\\
3.1207	-0.0615567\\
3.1233	-0.0615262\\
3.1259	-0.0614955\\
3.1284	-0.0614659\\
3.131	-0.061435\\
3.1336	-0.061404\\
3.1361	-0.061374\\
3.1387	-0.0613428\\
3.1413	-0.0613114\\
3.1438	-0.0612811\\
3.1464	-0.0612495\\
3.1489	-0.0612191\\
3.1515	-0.0611873\\
3.1541	-0.0611554\\
3.1566	-0.0611246\\
3.1592	-0.0610925\\
3.1618	-0.0610604\\
3.1643	-0.0610293\\
3.1669	-0.0609969\\
3.1695	-0.0609645\\
3.172	-0.0609332\\
3.1746	-0.0609006\\
3.1772	-0.0608679\\
3.1797	-0.0608363\\
3.1823	-0.0608035\\
3.1848	-0.0607718\\
3.1874	-0.0607388\\
3.19	-0.0607057\\
3.1925	-0.0606738\\
3.1951	-0.0606406\\
3.1977	-0.0606073\\
3.2002	-0.0605753\\
3.2028	-0.0605419\\
3.2054	-0.0605084\\
3.2079	-0.0604762\\
3.2105	-0.0604426\\
3.213	-0.0604103\\
3.2156	-0.0603766\\
3.2182	-0.0603429\\
3.2207	-0.0603105\\
3.2233	-0.0602767\\
3.2259	-0.0602428\\
3.2284	-0.0602103\\
3.231	-0.0601764\\
3.2336	-0.0601424\\
3.2361	-0.0601098\\
3.2387	-0.0600758\\
3.2412	-0.0600431\\
3.2438	-0.060009\\
3.2464	-0.059975\\
3.2489	-0.0599422\\
3.2515	-0.0599081\\
3.2541	-0.059874\\
3.2566	-0.0598412\\
3.2592	-0.0598071\\
3.2618	-0.0597729\\
3.2643	-0.0597401\\
3.2669	-0.059706\\
3.2694	-0.0596731\\
3.272	-0.059639\\
3.2746	-0.0596049\\
3.2771	-0.0595721\\
3.2797	-0.0595379\\
3.2823	-0.0595038\\
3.2848	-0.059471\\
3.2874	-0.059437\\
3.29	-0.0594029\\
3.2925	-0.0593702\\
3.2951	-0.0593362\\
3.2976	-0.0593035\\
3.3002	-0.0592695\\
3.3028	-0.0592356\\
3.3053	-0.059203\\
3.3079	-0.0591691\\
3.3105	-0.0591353\\
3.313	-0.0591028\\
3.3156	-0.059069\\
3.3182	-0.0590353\\
3.3207	-0.0590029\\
3.3233	-0.0589693\\
3.3259	-0.0589357\\
3.3284	-0.0589035\\
3.331	-0.05887\\
3.3335	-0.0588379\\
3.3361	-0.0588045\\
3.3387	-0.0587712\\
3.3412	-0.0587393\\
3.3438	-0.0587061\\
3.3464	-0.058673\\
3.3489	-0.0586412\\
3.3515	-0.0586082\\
3.3541	-0.0585753\\
3.3566	-0.0585437\\
3.3592	-0.058511\\
3.3617	-0.0584795\\
3.3643	-0.0584469\\
3.3669	-0.0584144\\
3.3694	-0.0583832\\
3.372	-0.0583508\\
3.3746	-0.0583185\\
3.3771	-0.0582876\\
3.3797	-0.0582555\\
3.3823	-0.0582235\\
3.3848	-0.0581928\\
3.3874	-0.058161\\
3.3899	-0.0581305\\
3.3925	-0.0580988\\
3.3951	-0.0580673\\
3.3976	-0.0580371\\
3.4002	-0.0580058\\
3.4028	-0.0579745\\
3.4053	-0.0579446\\
3.4079	-0.0579136\\
3.4105	-0.0578827\\
3.413	-0.0578531\\
3.4156	-0.0578224\\
3.4181	-0.0577931\\
3.4207	-0.0577626\\
3.4233	-0.0577323\\
3.4258	-0.0577032\\
3.4284	-0.0576731\\
3.431	-0.0576432\\
3.4335	-0.0576145\\
3.4361	-0.0575848\\
3.4387	-0.0575552\\
3.4412	-0.0575268\\
3.4438	-0.0574975\\
3.4463	-0.0574694\\
3.4489	-0.0574404\\
3.4515	-0.0574114\\
3.454	-0.0573837\\
3.4566	-0.057355\\
3.4592	-0.0573265\\
3.4617	-0.0572992\\
3.4643	-0.0572709\\
3.4669	-0.0572428\\
3.4694	-0.0572159\\
3.472	-0.0571881\\
3.4745	-0.0571615\\
3.4771	-0.0571339\\
3.4797	-0.0571065\\
3.4822	-0.0570803\\
3.4848	-0.0570532\\
3.4874	-0.0570263\\
3.4899	-0.0570005\\
3.4925	-0.0569738\\
3.4951	-0.0569473\\
3.4976	-0.056922\\
3.5002	-0.0568958\\
3.5028	-0.0568698\\
3.5053	-0.0568449\\
3.5079	-0.0568192\\
3.5104	-0.0567946\\
3.513	-0.0567692\\
3.5156	-0.0567439\\
3.5181	-0.0567198\\
3.5207	-0.0566949\\
3.5233	-0.0566701\\
3.5258	-0.0566464\\
3.5284	-0.056622\\
3.531	-0.0565977\\
3.5335	-0.0565746\\
3.5361	-0.0565506\\
3.5386	-0.0565278\\
3.5412	-0.0565041\\
3.5438	-0.0564807\\
3.5463	-0.0564583\\
3.5489	-0.0564352\\
3.5515	-0.0564123\\
3.554	-0.0563904\\
3.5566	-0.0563678\\
3.5592	-0.0563454\\
3.5617	-0.056324\\
3.5643	-0.0563019\\
3.5668	-0.0562809\\
3.5694	-0.0562592\\
3.572	-0.0562376\\
3.5745	-0.0562171\\
3.5771	-0.0561959\\
3.5797	-0.0561748\\
3.5822	-0.0561548\\
3.5848	-0.0561341\\
3.5874	-0.0561136\\
3.5899	-0.0560941\\
3.5925	-0.0560739\\
3.595	-0.0560547\\
3.5976	-0.0560349\\
3.6002	-0.0560153\\
3.6027	-0.0559967\\
3.6053	-0.0559774\\
3.6079	-0.0559583\\
3.6104	-0.0559402\\
3.613	-0.0559215\\
3.6156	-0.0559029\\
3.6181	-0.0558853\\
3.6207	-0.0558671\\
3.6232	-0.0558498\\
3.6258	-0.055832\\
3.6284	-0.0558144\\
3.6309	-0.0557976\\
3.6335	-0.0557804\\
3.6361	-0.0557633\\
3.6386	-0.055747\\
3.6412	-0.0557303\\
3.6438	-0.0557138\\
3.6463	-0.0556981\\
3.6489	-0.0556819\\
3.6515	-0.0556659\\
3.654	-0.0556507\\
3.6566	-0.0556351\\
3.6591	-0.0556202\\
3.6617	-0.055605\\
3.6643	-0.0555899\\
3.6668	-0.0555756\\
3.6694	-0.0555609\\
3.672	-0.0555463\\
3.6745	-0.0555325\\
3.6771	-0.0555184\\
3.6797	-0.0555044\\
3.6822	-0.0554911\\
3.6848	-0.0554775\\
3.6873	-0.0554646\\
3.6899	-0.0554513\\
3.6925	-0.0554382\\
3.695	-0.0554258\\
3.6976	-0.0554131\\
3.7002	-0.0554006\\
3.7027	-0.0553887\\
3.7053	-0.0553765\\
3.7079	-0.0553646\\
3.7104	-0.0553532\\
3.713	-0.0553416\\
3.7155	-0.0553306\\
3.7181	-0.0553193\\
3.7207	-0.0553082\\
3.7232	-0.0552977\\
3.7258	-0.055287\\
3.7284	-0.0552764\\
3.7309	-0.0552665\\
3.7335	-0.0552563\\
3.7361	-0.0552462\\
3.7386	-0.0552368\\
3.7412	-0.0552271\\
3.7437	-0.055218\\
3.7463	-0.0552087\\
3.7489	-0.0551996\\
3.7514	-0.055191\\
3.754	-0.0551822\\
3.7566	-0.0551736\\
3.7591	-0.0551655\\
3.7617	-0.0551572\\
3.7643	-0.0551491\\
3.7668	-0.0551415\\
3.7694	-0.0551338\\
3.7719	-0.0551265\\
3.7745	-0.0551192\\
3.7771	-0.055112\\
3.7796	-0.0551052\\
3.7822	-0.0550983\\
3.7848	-0.0550916\\
3.7873	-0.0550854\\
3.7899	-0.055079\\
3.7925	-0.0550728\\
3.795	-0.055067\\
3.7976	-0.0550612\\
3.8002	-0.0550555\\
3.8027	-0.0550502\\
3.8053	-0.0550449\\
3.8078	-0.0550399\\
3.8104	-0.0550349\\
3.813	-0.0550301\\
3.8155	-0.0550256\\
3.8181	-0.0550211\\
3.8207	-0.0550168\\
3.8232	-0.0550128\\
3.8258	-0.0550088\\
3.8284	-0.055005\\
3.8309	-0.0550014\\
3.8335	-0.0549979\\
3.836	-0.0549947\\
3.8386	-0.0549915\\
3.8412	-0.0549885\\
3.8437	-0.0549857\\
3.8463	-0.054983\\
3.8489	-0.0549805\\
3.8514	-0.0549782\\
3.854	-0.054976\\
3.8566	-0.0549739\\
3.8591	-0.0549721\\
3.8617	-0.0549703\\
3.8642	-0.0549688\\
3.8668	-0.0549673\\
3.8694	-0.0549661\\
3.8719	-0.054965\\
3.8745	-0.054964\\
3.8771	-0.0549632\\
3.8796	-0.0549626\\
3.8822	-0.054962\\
3.8848	-0.0549617\\
3.8873	-0.0549615\\
3.8899	-0.0549614\\
3.8924	-0.0549615\\
3.895	-0.0549618\\
3.8976	-0.0549622\\
3.9001	-0.0549627\\
3.9027	-0.0549634\\
3.9053	-0.0549642\\
3.9078	-0.0549652\\
3.9104	-0.0549663\\
3.913	-0.0549676\\
3.9155	-0.054969\\
3.9181	-0.0549706\\
3.9206	-0.0549722\\
3.9232	-0.0549741\\
3.9258	-0.0549761\\
3.9283	-0.0549781\\
3.9309	-0.0549804\\
3.9335	-0.0549829\\
3.936	-0.0549853\\
3.9386	-0.054988\\
3.9412	-0.0549909\\
3.9437	-0.0549938\\
3.9463	-0.0549969\\
3.9488	-0.055\\
3.9514	-0.0550034\\
3.954	-0.055007\\
3.9565	-0.0550105\\
3.9591	-0.0550143\\
3.9617	-0.0550182\\
3.9642	-0.0550221\\
3.9668	-0.0550263\\
3.9694	-0.0550307\\
3.9719	-0.055035\\
3.9745	-0.0550396\\
3.9771	-0.0550443\\
3.9796	-0.055049\\
3.9822	-0.0550539\\
3.9847	-0.0550589\\
3.9873	-0.0550641\\
3.9899	-0.0550695\\
3.9924	-0.0550747\\
3.995	-0.0550804\\
3.9976	-0.0550861\\
4.0001	-0.0550917\\
4.0027	-0.0550977\\
4.0053	-0.0551038\\
4.0078	-0.0551098\\
4.0104	-0.0551162\\
4.0129	-0.0551224\\
4.0155	-0.055129\\
4.0181	-0.0551357\\
4.0206	-0.0551422\\
4.0232	-0.0551492\\
4.0258	-0.0551562\\
4.0283	-0.0551631\\
4.0309	-0.0551704\\
4.0335	-0.0551778\\
4.036	-0.055185\\
4.0386	-0.0551926\\
4.0411	-0.0552001\\
4.0437	-0.0552079\\
4.0463	-0.0552159\\
4.0488	-0.0552236\\
4.0514	-0.0552318\\
4.054	-0.0552401\\
4.0565	-0.0552481\\
4.0591	-0.0552566\\
4.0617	-0.0552652\\
4.0642	-0.0552736\\
4.0668	-0.0552824\\
4.0693	-0.055291\\
4.0719	-0.0553\\
4.0745	-0.0553091\\
4.077	-0.0553179\\
4.0796	-0.0553272\\
4.0822	-0.0553366\\
4.0847	-0.0553458\\
4.0873	-0.0553554\\
4.0899	-0.055365\\
4.0924	-0.0553745\\
4.095	-0.0553843\\
4.0975	-0.0553939\\
4.1001	-0.055404\\
4.1027	-0.0554141\\
4.1052	-0.055424\\
4.1078	-0.0554343\\
4.1104	-0.0554448\\
4.1129	-0.0554549\\
4.1155	-0.0554655\\
4.1181	-0.0554761\\
4.1206	-0.0554865\\
4.1232	-0.0554974\\
4.1258	-0.0555083\\
4.1283	-0.0555189\\
4.1309	-0.05553\\
4.1334	-0.0555407\\
4.136	-0.055552\\
4.1386	-0.0555633\\
4.1411	-0.0555743\\
4.1437	-0.0555858\\
4.1463	-0.0555974\\
4.1488	-0.0556086\\
4.1514	-0.0556203\\
4.154	-0.0556321\\
4.1565	-0.0556435\\
4.1591	-0.0556554\\
4.1616	-0.055667\\
4.1642	-0.055679\\
4.1668	-0.0556912\\
4.1693	-0.0557029\\
4.1719	-0.0557152\\
4.1745	-0.0557275\\
4.177	-0.0557394\\
4.1796	-0.0557519\\
4.1822	-0.0557644\\
4.1847	-0.0557765\\
4.1873	-0.0557892\\
4.1898	-0.0558014\\
4.1924	-0.0558141\\
4.195	-0.055827\\
4.1975	-0.0558393\\
4.2001	-0.0558523\\
4.2027	-0.0558652\\
4.2052	-0.0558778\\
4.2078	-0.0558909\\
4.2104	-0.055904\\
4.2129	-0.0559167\\
4.2155	-0.0559299\\
4.218	-0.0559427\\
4.2206	-0.055956\\
4.2232	-0.0559694\\
4.2257	-0.0559823\\
4.2283	-0.0559958\\
4.2309	-0.0560093\\
4.2334	-0.0560223\\
4.236	-0.0560359\\
4.2386	-0.0560495\\
4.2411	-0.0560626\\
4.2437	-0.0560763\\
4.2462	-0.0560895\\
4.2488	-0.0561033\\
4.2514	-0.0561171\\
4.2539	-0.0561304\\
4.2565	-0.0561443\\
4.2591	-0.0561582\\
4.2616	-0.0561716\\
4.2642	-0.0561855\\
4.2668	-0.0561995\\
4.2693	-0.056213\\
4.2719	-0.056227\\
4.2744	-0.0562405\\
4.277	-0.0562546\\
4.2796	-0.0562687\\
4.2821	-0.0562822\\
4.2847	-0.0562964\\
4.2873	-0.0563105\\
4.2898	-0.0563242\\
4.2924	-0.0563383\\
4.295	-0.0563525\\
4.2975	-0.0563662\\
4.3001	-0.0563804\\
4.3027	-0.0563947\\
4.3052	-0.0564084\\
4.3078	-0.0564226\\
4.3103	-0.0564363\\
4.3129	-0.0564506\\
4.3155	-0.0564649\\
4.318	-0.0564786\\
4.3206	-0.0564929\\
4.3232	-0.0565072\\
4.3257	-0.0565209\\
4.3283	-0.0565352\\
4.3309	-0.0565495\\
4.3334	-0.0565632\\
4.336	-0.0565775\\
4.3385	-0.0565912\\
4.3411	-0.0566054\\
4.3437	-0.0566197\\
4.3462	-0.0566334\\
4.3488	-0.0566476\\
4.3514	-0.0566619\\
4.3539	-0.0566755\\
4.3565	-0.0566897\\
4.3591	-0.0567039\\
4.3616	-0.0567175\\
4.3642	-0.0567317\\
4.3667	-0.0567453\\
4.3693	-0.0567594\\
4.3719	-0.0567735\\
4.3744	-0.0567871\\
4.377	-0.0568011\\
4.3796	-0.0568151\\
4.3821	-0.0568286\\
4.3847	-0.0568426\\
4.3873	-0.0568565\\
4.3898	-0.0568699\\
4.3924	-0.0568838\\
4.3949	-0.0568972\\
4.3975	-0.056911\\
4.4001	-0.0569248\\
4.4026	-0.056938\\
4.4052	-0.0569518\\
4.4078	-0.0569655\\
4.4103	-0.0569786\\
4.4129	-0.0569922\\
4.4155	-0.0570058\\
4.418	-0.0570188\\
4.4206	-0.0570323\\
4.4231	-0.0570453\\
4.4257	-0.0570587\\
4.4283	-0.057072\\
4.4308	-0.0570849\\
4.4334	-0.0570981\\
4.436	-0.0571114\\
4.4385	-0.057124\\
4.4411	-0.0571372\\
4.4437	-0.0571502\\
4.4462	-0.0571628\\
4.4488	-0.0571757\\
4.4514	-0.0571886\\
4.4539	-0.057201\\
4.4565	-0.0572138\\
4.459	-0.0572261\\
4.4616	-0.0572388\\
4.4642	-0.0572514\\
4.4667	-0.0572635\\
4.4693	-0.057276\\
4.4719	-0.0572884\\
4.4744	-0.0573003\\
4.477	-0.0573126\\
4.4796	-0.0573249\\
4.4821	-0.0573366\\
4.4847	-0.0573487\\
4.4872	-0.0573602\\
4.4898	-0.0573722\\
4.4924	-0.0573841\\
4.4949	-0.0573955\\
4.4975	-0.0574072\\
4.5001	-0.0574189\\
4.5026	-0.0574301\\
4.5052	-0.0574416\\
4.5078	-0.0574531\\
4.5103	-0.057464\\
4.5129	-0.0574753\\
4.5154	-0.0574861\\
4.518	-0.0574972\\
4.5206	-0.0575082\\
4.5231	-0.0575188\\
4.5257	-0.0575297\\
4.5283	-0.0575405\\
4.5308	-0.0575508\\
4.5334	-0.0575614\\
4.536	-0.057572\\
4.5385	-0.057582\\
4.5411	-0.0575924\\
4.5436	-0.0576023\\
4.5462	-0.0576124\\
4.5488	-0.0576225\\
4.5513	-0.0576322\\
4.5539	-0.0576421\\
4.5565	-0.0576519\\
4.559	-0.0576612\\
4.5616	-0.0576709\\
4.5642	-0.0576804\\
4.5667	-0.0576894\\
4.5693	-0.0576988\\
4.5718	-0.0577077\\
4.5744	-0.0577168\\
4.577	-0.0577258\\
4.5795	-0.0577344\\
4.5821	-0.0577433\\
4.5847	-0.057752\\
4.5872	-0.0577603\\
4.5898	-0.0577688\\
4.5924	-0.0577772\\
4.5949	-0.0577852\\
4.5975	-0.0577935\\
4.6001	-0.0578016\\
4.6026	-0.0578093\\
4.6052	-0.0578171\\
4.6077	-0.0578246\\
4.6103	-0.0578323\\
4.6129	-0.0578399\\
4.6154	-0.057847\\
4.618	-0.0578544\\
4.6206	-0.0578616\\
4.6231	-0.0578685\\
4.6257	-0.0578755\\
4.6283	-0.0578824\\
4.6308	-0.0578889\\
4.6334	-0.0578956\\
4.6359	-0.0579019\\
4.6385	-0.0579083\\
4.6411	-0.0579146\\
4.6436	-0.0579206\\
4.6462	-0.0579267\\
4.6488	-0.0579327\\
4.6513	-0.0579383\\
4.6539	-0.057944\\
4.6565	-0.0579497\\
4.659	-0.057955\\
4.6616	-0.0579603\\
4.6641	-0.0579654\\
4.6667	-0.0579706\\
4.6693	-0.0579756\\
4.6718	-0.0579803\\
4.6744	-0.0579851\\
4.677	-0.0579897\\
4.6795	-0.0579941\\
4.6821	-0.0579985\\
4.6847	-0.0580028\\
4.6872	-0.0580069\\
4.6898	-0.0580109\\
4.6923	-0.0580147\\
4.6949	-0.0580185\\
4.6975	-0.0580222\\
4.7	-0.0580256\\
4.7026	-0.058029\\
4.7052	-0.0580324\\
4.7077	-0.0580354\\
4.7103	-0.0580385\\
4.7129	-0.0580414\\
4.7154	-0.0580441\\
4.718	-0.0580468\\
4.7205	-0.0580493\\
4.7231	-0.0580517\\
4.7257	-0.058054\\
4.7282	-0.0580561\\
4.7308	-0.0580582\\
4.7334	-0.0580601\\
4.7359	-0.0580619\\
4.7385	-0.0580636\\
4.7411	-0.0580651\\
4.7436	-0.0580665\\
4.7462	-0.0580678\\
4.7487	-0.0580689\\
4.7513	-0.05807\\
4.7539	-0.0580709\\
4.7564	-0.0580717\\
4.759	-0.0580724\\
4.7616	-0.0580729\\
4.7641	-0.0580733\\
4.7667	-0.0580736\\
4.7693	-0.0580738\\
4.7718	-0.0580739\\
4.7744	-0.0580738\\
4.777	-0.0580736\\
4.7795	-0.0580733\\
4.7821	-0.0580728\\
4.7846	-0.0580723\\
4.7872	-0.0580716\\
4.7898	-0.0580707\\
4.7923	-0.0580698\\
4.7949	-0.0580687\\
4.7975	-0.0580675\\
4.8	-0.0580663\\
4.8026	-0.0580648\\
4.8052	-0.0580632\\
4.8077	-0.0580616\\
4.8103	-0.0580598\\
4.8128	-0.0580579\\
4.8154	-0.0580559\\
4.818	-0.0580537\\
4.8205	-0.0580514\\
4.8231	-0.058049\\
4.8257	-0.0580465\\
4.8282	-0.0580439\\
4.8308	-0.0580411\\
4.8334	-0.0580382\\
4.8359	-0.0580353\\
4.8385	-0.0580321\\
4.841	-0.058029\\
4.8436	-0.0580256\\
4.8462	-0.0580221\\
4.8487	-0.0580186\\
4.8513	-0.0580148\\
4.8539	-0.058011\\
4.8564	-0.0580071\\
4.859	-0.058003\\
4.8616	-0.0579988\\
4.8641	-0.0579947\\
4.8667	-0.0579902\\
4.8692	-0.0579858\\
4.8718	-0.0579812\\
4.8744	-0.0579764\\
4.8769	-0.0579717\\
4.8795	-0.0579667\\
4.8821	-0.0579616\\
4.8846	-0.0579565\\
4.8872	-0.0579512\\
4.8898	-0.0579457\\
4.8923	-0.0579404\\
4.8949	-0.0579347\\
4.8974	-0.0579291\\
4.9	-0.0579232\\
4.9026	-0.0579172\\
4.9051	-0.0579114\\
4.9077	-0.0579052\\
4.9103	-0.0578989\\
4.9128	-0.0578927\\
4.9154	-0.0578862\\
4.918	-0.0578795\\
4.9205	-0.0578731\\
4.9231	-0.0578662\\
4.9257	-0.0578593\\
4.9282	-0.0578525\\
4.9308	-0.0578454\\
4.9333	-0.0578384\\
4.9359	-0.0578311\\
4.9385	-0.0578236\\
4.941	-0.0578164\\
4.9436	-0.0578088\\
4.9462	-0.0578011\\
4.9487	-0.0577935\\
4.9513	-0.0577856\\
4.9539	-0.0577776\\
4.9564	-0.0577698\\
4.959	-0.0577616\\
4.9615	-0.0577537\\
4.9641	-0.0577453\\
4.9667	-0.0577368\\
4.9692	-0.0577286\\
4.9718	-0.0577199\\
4.9744	-0.0577112\\
4.9769	-0.0577027\\
4.9795	-0.0576938\\
4.9821	-0.0576848\\
4.9846	-0.0576761\\
4.9872	-0.057667\\
4.9897	-0.0576581\\
4.9923	-0.0576487\\
4.9949	-0.0576393\\
4.9974	-0.0576302\\
5	-0.0576207\\
};
\addplot [color=mycolor17, line width=1.0pt, forget plot]
  table[row sep=crcr]{%
-0.125	3.95559e-08\\
-0.1224	4.0374e-08\\
-0.1199	4.11606e-08\\
-0.1173	4.19785e-08\\
-0.1147	4.27963e-08\\
-0.1122	4.35827e-08\\
-0.1096	4.44004e-08\\
-0.1071	4.51866e-08\\
-0.1045	4.60042e-08\\
-0.1019	4.68217e-08\\
-0.0994	4.76076e-08\\
-0.0968	4.8425e-08\\
-0.0942	4.92423e-08\\
-0.0917	5.0028e-08\\
-0.0891	5.08451e-08\\
-0.0865	5.16622e-08\\
-0.084	5.24477e-08\\
-0.0814	5.32646e-08\\
-0.0789	5.40501e-08\\
-0.0763	5.48668e-08\\
-0.0737	5.56835e-08\\
-0.0712	5.64687e-08\\
-0.0686	5.72852e-08\\
-0.066	5.81016e-08\\
-0.0635	5.88867e-08\\
-0.0609	5.97029e-08\\
-0.0583	6.05191e-08\\
-0.0558	6.13039e-08\\
-0.0532	6.21199e-08\\
-0.0507	6.29045e-08\\
-0.0481	6.37205e-08\\
-0.0455	6.45363e-08\\
-0.043	6.53207e-08\\
-0.0404	6.61365e-08\\
-0.0378	6.69521e-08\\
-0.0353	6.77363e-08\\
-0.0327	6.85518e-08\\
-0.0301	6.93671e-08\\
-0.0276	7.01511e-08\\
-0.025	7.09663e-08\\
-0.0224	7.17815e-08\\
-0.0199	7.25652e-08\\
-0.0173	7.33802e-08\\
-0.0148	7.41639e-08\\
-0.0122	7.49787e-08\\
-0.0096	7.57935e-08\\
-0.0071	7.65769e-08\\
-0.0045	7.73914e-08\\
-0.0019	7.8206e-08\\
0.0006	7.89892e-08\\
0.0032	7.78536e-05\\
0.0058	0.000417522\\
0.0083	0.000811805\\
0.0109	0.00129035\\
0.0134	0.00179261\\
0.016	0.00232283\\
0.0186	0.00283838\\
0.0211	0.00331675\\
0.0237	0.00380412\\
0.0263	0.0042892\\
0.0288	0.00475438\\
0.0314	0.00523177\\
0.034	0.00569674\\
0.0365	0.00613019\\
0.0391	0.0065691\\
0.0416	0.00698324\\
0.0442	0.0074078\\
0.0468	0.00782528\\
0.0493	0.00821788\\
0.0519	0.00861551\\
0.0545	0.0090028\\
0.057	0.00936734\\
0.0596	0.00974016\\
0.0622	0.010107\\
0.0647	0.0104531\\
0.0673	0.0108052\\
0.0698	0.011136\\
0.0724	0.0114731\\
0.075	0.011805\\
0.0775	0.01212\\
0.0801	0.0124433\\
0.0827	0.0127611\\
0.0852	0.013061\\
0.0878	0.0133675\\
0.0904	0.01367\\
0.0929	0.0139581\\
0.0955	0.0142553\\
0.098	0.014538\\
0.1006	0.0148281\\
0.1032	0.015114\\
0.1057	0.0153858\\
0.1083	0.0156665\\
0.1109	0.0159458\\
0.1134	0.0162127\\
0.116	0.0164878\\
0.1186	0.0167598\\
0.1211	0.0170185\\
0.1237	0.0172858\\
0.1263	0.017552\\
0.1288	0.0178071\\
0.1314	0.0180707\\
0.1339	0.0183218\\
0.1365	0.0185803\\
0.1391	0.0188365\\
0.1416	0.0190816\\
0.1442	0.0193355\\
0.1468	0.0195882\\
0.1493	0.0198291\\
0.1519	0.0200768\\
0.1545	0.0203218\\
0.157	0.0205555\\
0.1596	0.0207972\\
0.1621	0.0210282\\
0.1647	0.0212664\\
0.1673	0.0215018\\
0.1698	0.0217251\\
0.1724	0.0219546\\
0.175	0.022182\\
0.1775	0.022399\\
0.1801	0.0226225\\
0.1827	0.0228432\\
0.1852	0.0230521\\
0.1878	0.0232661\\
0.1903	0.0234692\\
0.1929	0.0236782\\
0.1955	0.023885\\
0.198	0.0240812\\
0.2006	0.0242819\\
0.2032	0.0244789\\
0.2057	0.0246652\\
0.2083	0.0248564\\
0.2109	0.0250452\\
0.2134	0.0252243\\
0.216	0.0254075\\
0.2185	0.0255802\\
0.2211	0.0257565\\
0.2237	0.0259298\\
0.2262	0.0260942\\
0.2288	0.0262628\\
0.2314	0.0264288\\
0.2339	0.0265853\\
0.2365	0.0267449\\
0.2391	0.0269015\\
0.2416	0.0270499\\
0.2442	0.0272022\\
0.2467	0.0273467\\
0.2493	0.0274943\\
0.2519	0.0276391\\
0.2544	0.0277758\\
0.257	0.0279157\\
0.2596	0.0280539\\
0.2621	0.0281853\\
0.2647	0.02832\\
0.2673	0.0284525\\
0.2698	0.0285778\\
0.2724	0.0287061\\
0.2749	0.0288281\\
0.2775	0.0289539\\
0.2801	0.0290784\\
0.2826	0.0291967\\
0.2852	0.0293179\\
0.2878	0.0294375\\
0.2903	0.0295514\\
0.2929	0.029669\\
0.2955	0.029786\\
0.298	0.0298975\\
0.3006	0.0300124\\
0.3032	0.0301259\\
0.3057	0.0302342\\
0.3083	0.0303462\\
0.3108	0.0304536\\
0.3134	0.0305649\\
0.316	0.0306755\\
0.3185	0.0307809\\
0.3211	0.0308898\\
0.3237	0.0309982\\
0.3262	0.0311022\\
0.3288	0.0312104\\
0.3314	0.0313182\\
0.3339	0.0314213\\
0.3365	0.0315278\\
0.339	0.0316298\\
0.3416	0.0317358\\
0.3442	0.0318417\\
0.3467	0.0319435\\
0.3493	0.0320491\\
0.3519	0.032154\\
0.3544	0.0322544\\
0.357	0.0323586\\
0.3596	0.0324628\\
0.3621	0.0325629\\
0.3647	0.0326668\\
0.3672	0.0327663\\
0.3698	0.0328691\\
0.3724	0.0329715\\
0.3749	0.0330697\\
0.3775	0.0331718\\
0.3801	0.0332737\\
0.3826	0.0333712\\
0.3852	0.033472\\
0.3878	0.0335721\\
0.3903	0.033668\\
0.3929	0.0337673\\
0.3954	0.0338626\\
0.398	0.0339612\\
0.4006	0.0340591\\
0.4031	0.0341524\\
0.4057	0.0342487\\
0.4083	0.0343445\\
0.4108	0.0344361\\
0.4134	0.0345309\\
0.416	0.0346248\\
0.4185	0.0347142\\
0.4211	0.0348062\\
0.4236	0.0348939\\
0.4262	0.0349843\\
0.4288	0.0350741\\
0.4313	0.0351596\\
0.4339	0.0352474\\
0.4365	0.035334\\
0.439	0.0354163\\
0.4416	0.0355009\\
0.4442	0.0355846\\
0.4467	0.0356641\\
0.4493	0.0357457\\
0.4519	0.035826\\
0.4544	0.0359018\\
0.457	0.0359796\\
0.4595	0.0360533\\
0.4621	0.0361288\\
0.4647	0.0362032\\
0.4672	0.0362733\\
0.4698	0.0363447\\
0.4724	0.0364147\\
0.4749	0.0364807\\
0.4775	0.0365481\\
0.4801	0.0366143\\
0.4826	0.0366764\\
0.4852	0.0367395\\
0.4877	0.0367985\\
0.4903	0.0368585\\
0.4929	0.036917\\
0.4954	0.0369719\\
0.498	0.0370275\\
0.5006	0.0370814\\
0.5031	0.0371315\\
0.5057	0.037182\\
0.5083	0.0372308\\
0.5108	0.0372763\\
0.5134	0.0373221\\
0.5159	0.0373644\\
0.5185	0.0374066\\
0.5211	0.0374469\\
0.5236	0.0374839\\
0.5262	0.0375208\\
0.5288	0.037556\\
0.5313	0.0375882\\
0.5339	0.0376197\\
0.5365	0.0376492\\
0.539	0.0376758\\
0.5416	0.0377016\\
0.5441	0.0377247\\
0.5467	0.0377469\\
0.5493	0.0377671\\
0.5518	0.0377846\\
0.5544	0.0378007\\
0.557	0.0378148\\
0.5595	0.0378266\\
0.5621	0.0378369\\
0.5647	0.0378452\\
0.5672	0.0378511\\
0.5698	0.0378551\\
0.5723	0.037857\\
0.5749	0.0378568\\
0.5775	0.0378547\\
0.58	0.0378506\\
0.5826	0.0378442\\
0.5852	0.0378356\\
0.5877	0.0378251\\
0.5903	0.037812\\
0.5929	0.0377968\\
0.5954	0.0377801\\
0.598	0.0377605\\
0.6006	0.0377387\\
0.6031	0.0377154\\
0.6057	0.0376888\\
0.6082	0.0376612\\
0.6108	0.0376302\\
0.6134	0.0375969\\
0.6159	0.0375627\\
0.6185	0.0375247\\
0.6211	0.0374843\\
0.6236	0.0374431\\
0.6262	0.037398\\
0.6288	0.0373506\\
0.6313	0.0373027\\
0.6339	0.0372504\\
0.6364	0.0371978\\
0.639	0.0371406\\
0.6416	0.037081\\
0.6441	0.0370213\\
0.6467	0.0369569\\
0.6493	0.03689\\
0.6518	0.0368233\\
0.6544	0.0367513\\
0.657	0.0366768\\
0.6595	0.0366029\\
0.6621	0.0365235\\
0.6646	0.0364449\\
0.6672	0.0363605\\
0.6698	0.0362735\\
0.6723	0.0361874\\
0.6749	0.0360954\\
0.6775	0.0360009\\
0.68	0.0359076\\
0.6826	0.0358081\\
0.6852	0.0357059\\
0.6877	0.0356052\\
0.6903	0.0354979\\
0.6928	0.0353924\\
0.6954	0.0352801\\
0.698	0.0351653\\
0.7005	0.0350524\\
0.7031	0.0349324\\
0.7057	0.0348099\\
0.7082	0.0346896\\
0.7108	0.034562\\
0.7134	0.0344319\\
0.7159	0.0343043\\
0.7185	0.0341691\\
0.721	0.0340366\\
0.7236	0.0338963\\
0.7262	0.0337535\\
0.7287	0.0336138\\
0.7313	0.033466\\
0.7339	0.0333157\\
0.7364	0.0331687\\
0.739	0.0330133\\
0.7416	0.0328554\\
0.7441	0.0327012\\
0.7467	0.0325385\\
0.7492	0.0323796\\
0.7518	0.0322119\\
0.7544	0.0320417\\
0.7569	0.0318757\\
0.7595	0.0317006\\
0.7621	0.0315231\\
0.7646	0.0313502\\
0.7672	0.0311679\\
0.7698	0.0309832\\
0.7723	0.0308032\\
0.7749	0.0306137\\
0.7775	0.0304219\\
0.78	0.0302352\\
0.7826	0.0300387\\
0.7851	0.0298476\\
0.7877	0.0296464\\
0.7903	0.0294429\\
0.7928	0.0292451\\
0.7954	0.0290371\\
0.798	0.0288268\\
0.8005	0.0286225\\
0.8031	0.0284077\\
0.8057	0.0281907\\
0.8082	0.0279799\\
0.8108	0.0277585\\
0.8133	0.0275436\\
0.8159	0.0273179\\
0.8185	0.0270901\\
0.821	0.0268689\\
0.8236	0.0266368\\
0.8262	0.0264025\\
0.8287	0.0261753\\
0.8313	0.025937\\
0.8339	0.0256966\\
0.8364	0.0254634\\
0.839	0.0252189\\
0.8415	0.0249819\\
0.8441	0.0247335\\
0.8467	0.024483\\
0.8492	0.0242404\\
0.8518	0.023986\\
0.8544	0.0237297\\
0.8569	0.0234815\\
0.8595	0.0232214\\
0.8621	0.0229595\\
0.8646	0.0227059\\
0.8672	0.0224403\\
0.8697	0.0221831\\
0.8723	0.0219139\\
0.8749	0.0216428\\
0.8774	0.0213805\\
0.88	0.021106\\
0.8826	0.0208298\\
0.8851	0.0205625\\
0.8877	0.0202828\\
0.8903	0.0200014\\
0.8928	0.0197292\\
0.8954	0.0194446\\
0.8979	0.0191693\\
0.9005	0.0188815\\
0.9031	0.0185919\\
0.9056	0.0183121\\
0.9082	0.0180194\\
0.9108	0.0177253\\
0.9133	0.017441\\
0.9159	0.0171438\\
0.9185	0.0168451\\
0.921	0.0165565\\
0.9236	0.0162549\\
0.9262	0.0159519\\
0.9287	0.0156591\\
0.9313	0.0153533\\
0.9338	0.0150579\\
0.9364	0.0147493\\
0.939	0.0144394\\
0.9415	0.0141401\\
0.9441	0.0138275\\
0.9467	0.0135137\\
0.9492	0.0132107\\
0.9518	0.0128943\\
0.9544	0.0125766\\
0.9569	0.01227\\
0.9595	0.0119499\\
0.962	0.011641\\
0.9646	0.0113185\\
0.9672	0.0109949\\
0.9697	0.0106827\\
0.9723	0.0103568\\
0.9749	0.0100298\\
0.9774	0.00971433\\
0.98	0.00938519\\
0.9826	0.00905498\\
0.9851	0.00873645\\
0.9877	0.00840415\\
0.9902	0.00808367\\
0.9928	0.00774938\\
0.9954	0.00741411\\
0.9979	0.00709082\\
1.0005	0.00675365\\
1.0031	0.00641552\\
1.0056	0.00608951\\
1.0082	0.00574956\\
1.0108	0.00540874\\
1.0133	0.00508019\\
1.0159	0.00473765\\
1.0184	0.00440747\\
1.021	0.00406325\\
1.0236	0.00371821\\
1.0261	0.00338569\\
1.0287	0.0030391\\
1.0313	0.00269174\\
1.0338	0.00235701\\
1.0364	0.00200815\\
1.039	0.00165855\\
1.0415	0.00132172\\
1.0441	0.000970738\\
1.0466	0.000632611\\
1.0492	0.000280295\\
1.0518	-7.26913e-05\\
1.0543	-0.000412719\\
1.0569	-0.000766971\\
1.0595	-0.00112184\\
1.062	-0.00146362\\
1.0646	-0.00181965\\
1.0672	-0.00217626\\
1.0697	-0.0025197\\
1.0723	-0.00287743\\
1.0748	-0.0032219\\
1.0774	-0.00358066\\
1.08	-0.00393993\\
1.0825	-0.00428585\\
1.0851	-0.00464609\\
1.0877	-0.00500682\\
1.0902	-0.00535411\\
1.0928	-0.00571573\\
1.0954	-0.00607778\\
1.0979	-0.0064263\\
1.1005	-0.00678917\\
1.1031	-0.00715245\\
1.1056	-0.00750213\\
1.1082	-0.00786617\\
1.1107	-0.00821654\\
1.1133	-0.00858127\\
1.1159	-0.00894635\\
1.1184	-0.0092977\\
1.121	-0.00966342\\
1.1236	-0.0100295\\
1.1261	-0.0103817\\
1.1287	-0.0107483\\
1.1313	-0.0111151\\
1.1338	-0.0114682\\
1.1364	-0.0118355\\
1.1389	-0.012189\\
1.1415	-0.0125569\\
1.1441	-0.0129249\\
1.1466	-0.013279\\
1.1492	-0.0136475\\
1.1518	-0.0140162\\
1.1543	-0.0143708\\
1.1569	-0.0147398\\
1.1595	-0.015109\\
1.162	-0.0154641\\
1.1646	-0.0158335\\
1.1671	-0.0161889\\
1.1697	-0.0165586\\
1.1723	-0.0169284\\
1.1748	-0.0172841\\
1.1774	-0.017654\\
1.18	-0.0180241\\
1.1825	-0.01838\\
1.1851	-0.0187502\\
1.1877	-0.0191204\\
1.1902	-0.0194765\\
1.1928	-0.0198468\\
1.1953	-0.0202029\\
1.1979	-0.0205733\\
1.2005	-0.0209437\\
1.203	-0.0212998\\
1.2056	-0.0216701\\
1.2082	-0.0220405\\
1.2107	-0.0223965\\
1.2133	-0.0227668\\
1.2159	-0.023137\\
1.2184	-0.0234929\\
1.221	-0.023863\\
1.2235	-0.0242188\\
1.2261	-0.0245888\\
1.2287	-0.0249586\\
1.2312	-0.0253141\\
1.2338	-0.0256838\\
1.2364	-0.0260533\\
1.2389	-0.0264085\\
1.2415	-0.0267778\\
1.2441	-0.0271469\\
1.2466	-0.0275017\\
1.2492	-0.0278705\\
1.2518	-0.0282392\\
1.2543	-0.0285935\\
1.2569	-0.0289618\\
1.2594	-0.0293157\\
1.262	-0.0296836\\
1.2646	-0.0300513\\
1.2671	-0.0304046\\
1.2697	-0.0307719\\
1.2723	-0.0311389\\
1.2748	-0.0314916\\
1.2774	-0.0318581\\
1.28	-0.0322243\\
1.2825	-0.0325762\\
1.2851	-0.032942\\
1.2876	-0.0332934\\
1.2902	-0.0336585\\
1.2928	-0.0340233\\
1.2953	-0.0343739\\
1.2979	-0.0347381\\
1.3005	-0.035102\\
1.303	-0.0354515\\
1.3056	-0.0358148\\
1.3082	-0.0361776\\
1.3107	-0.0365262\\
1.3133	-0.0368883\\
1.3158	-0.0372362\\
1.3184	-0.0375976\\
1.321	-0.0379586\\
1.3235	-0.0383053\\
1.3261	-0.0386655\\
1.3287	-0.0390253\\
1.3312	-0.0393708\\
1.3338	-0.0397297\\
1.3364	-0.0400882\\
1.3389	-0.0404324\\
1.3415	-0.04079\\
1.344	-0.0411333\\
1.3466	-0.0414899\\
1.3492	-0.0418461\\
1.3517	-0.042188\\
1.3543	-0.0425432\\
1.3569	-0.0428978\\
1.3594	-0.0432383\\
1.362	-0.0435919\\
1.3646	-0.0439449\\
1.3671	-0.0442839\\
1.3697	-0.0446358\\
1.3722	-0.0449737\\
1.3748	-0.0453245\\
1.3774	-0.0456747\\
1.3799	-0.0460109\\
1.3825	-0.04636\\
1.3851	-0.0467085\\
1.3876	-0.0470429\\
1.3902	-0.0473901\\
1.3928	-0.0477367\\
1.3953	-0.0480693\\
1.3979	-0.0484146\\
1.4005	-0.0487593\\
1.403	-0.04909\\
1.4056	-0.0494333\\
1.4081	-0.0497627\\
1.4107	-0.0501046\\
1.4133	-0.0504459\\
1.4158	-0.0507733\\
1.4184	-0.051113\\
1.421	-0.0514521\\
1.4235	-0.0517774\\
1.4261	-0.0521149\\
1.4287	-0.0524517\\
1.4312	-0.0527748\\
1.4338	-0.0531101\\
1.4363	-0.0534317\\
1.4389	-0.0537654\\
1.4415	-0.0540983\\
1.444	-0.0544176\\
1.4466	-0.0547488\\
1.4492	-0.0550792\\
1.4517	-0.0553961\\
1.4543	-0.0557248\\
1.4569	-0.0560527\\
1.4594	-0.0563671\\
1.462	-0.0566932\\
1.4645	-0.0570059\\
1.4671	-0.0573303\\
1.4697	-0.0576537\\
1.4722	-0.0579638\\
1.4748	-0.0582854\\
1.4774	-0.058606\\
1.4799	-0.0589134\\
1.4825	-0.0592322\\
1.4851	-0.0595499\\
1.4876	-0.0598546\\
1.4902	-0.0601704\\
1.4927	-0.0604731\\
1.4953	-0.0607869\\
1.4979	-0.0610997\\
1.5004	-0.0613995\\
1.503	-0.0617103\\
1.5056	-0.06202\\
1.5081	-0.0623168\\
1.5107	-0.0626244\\
1.5133	-0.0629309\\
1.5158	-0.0632246\\
1.5184	-0.063529\\
1.5209	-0.0638206\\
1.5235	-0.0641228\\
1.5261	-0.0644238\\
1.5286	-0.0647121\\
1.5312	-0.0650109\\
1.5338	-0.0653085\\
1.5363	-0.0655935\\
1.5389	-0.0658888\\
1.5415	-0.0661829\\
1.544	-0.0664645\\
1.5466	-0.0667562\\
1.5491	-0.0670356\\
1.5517	-0.0673249\\
1.5543	-0.0676129\\
1.5568	-0.0678887\\
1.5594	-0.0681743\\
1.562	-0.0684586\\
1.5645	-0.0687307\\
1.5671	-0.0690125\\
1.5697	-0.0692929\\
1.5722	-0.0695613\\
1.5748	-0.0698392\\
1.5774	-0.0701157\\
1.5799	-0.0703803\\
1.5825	-0.0706542\\
1.585	-0.0709162\\
1.5876	-0.0711874\\
1.5902	-0.0714572\\
1.5927	-0.0717153\\
1.5953	-0.0719823\\
1.5979	-0.072248\\
1.6004	-0.072502\\
1.603	-0.0727648\\
1.6056	-0.0730262\\
1.6081	-0.0732762\\
1.6107	-0.0735347\\
1.6132	-0.0737818\\
1.6158	-0.0740374\\
1.6184	-0.0742916\\
1.6209	-0.0745345\\
1.6235	-0.0747856\\
1.6261	-0.0750352\\
1.6286	-0.0752738\\
1.6312	-0.0755204\\
1.6338	-0.0757655\\
1.6363	-0.0759997\\
1.6389	-0.0762416\\
1.6414	-0.0764728\\
1.644	-0.0767117\\
1.6466	-0.076949\\
1.6491	-0.0771756\\
1.6517	-0.0774097\\
1.6543	-0.0776422\\
1.6568	-0.0778642\\
1.6594	-0.0780935\\
1.662	-0.0783211\\
1.6645	-0.0785384\\
1.6671	-0.0787627\\
1.6696	-0.0789769\\
1.6722	-0.0791979\\
1.6748	-0.0794173\\
1.6773	-0.0796266\\
1.6799	-0.0798426\\
1.6825	-0.0800568\\
1.685	-0.0802612\\
1.6876	-0.0804721\\
1.6902	-0.0806812\\
1.6927	-0.0808806\\
1.6953	-0.0810863\\
1.6978	-0.0812823\\
1.7004	-0.0814845\\
1.703	-0.0816849\\
1.7055	-0.0818759\\
1.7081	-0.0820727\\
1.7107	-0.0822678\\
1.7132	-0.0824536\\
1.7158	-0.082645\\
1.7184	-0.0828346\\
1.7209	-0.0830152\\
1.7235	-0.0832012\\
1.7261	-0.0833854\\
1.7286	-0.0835607\\
1.7312	-0.0837412\\
1.7337	-0.0839129\\
1.7363	-0.0840897\\
1.7389	-0.0842646\\
1.7414	-0.084431\\
1.744	-0.0846021\\
1.7466	-0.0847714\\
1.7491	-0.0849323\\
1.7517	-0.0850978\\
1.7543	-0.0852613\\
1.7568	-0.0854167\\
1.7594	-0.0855765\\
1.7619	-0.0857282\\
1.7645	-0.0858841\\
1.7671	-0.086038\\
1.7696	-0.0861842\\
1.7722	-0.0863343\\
1.7748	-0.0864824\\
1.7773	-0.0866229\\
1.7799	-0.0867671\\
1.7825	-0.0869093\\
1.785	-0.0870442\\
1.7876	-0.0871825\\
1.7901	-0.0873135\\
1.7927	-0.0874479\\
1.7953	-0.0875802\\
1.7978	-0.0877055\\
1.8004	-0.0878339\\
1.803	-0.0879602\\
1.8055	-0.0880797\\
1.8081	-0.0882021\\
1.8107	-0.0883224\\
1.8132	-0.0884361\\
1.8158	-0.0885524\\
1.8183	-0.0886622\\
1.8209	-0.0887745\\
1.8235	-0.0888846\\
1.826	-0.0889886\\
1.8286	-0.0890948\\
1.8312	-0.0891988\\
1.8337	-0.089297\\
1.8363	-0.089397\\
1.8389	-0.0894949\\
1.8414	-0.0895871\\
1.844	-0.0896809\\
1.8465	-0.0897692\\
1.8491	-0.089859\\
1.8517	-0.0899467\\
1.8542	-0.090029\\
1.8568	-0.0901126\\
1.8594	-0.0901941\\
1.8619	-0.0902705\\
1.8645	-0.090348\\
1.8671	-0.0904233\\
1.8696	-0.0904937\\
1.8722	-0.0905649\\
1.8747	-0.0906314\\
1.8773	-0.0906985\\
1.8799	-0.0907636\\
1.8824	-0.0908241\\
1.885	-0.090885\\
1.8876	-0.0909438\\
1.8901	-0.0909984\\
1.8927	-0.0910531\\
1.8953	-0.0911058\\
1.8978	-0.0911544\\
1.9004	-0.0912029\\
1.903	-0.0912494\\
1.9055	-0.0912921\\
1.9081	-0.0913344\\
1.9106	-0.0913732\\
1.9132	-0.0914115\\
1.9158	-0.0914477\\
1.9183	-0.0914805\\
1.9209	-0.0915127\\
1.9235	-0.0915427\\
1.926	-0.0915697\\
1.9286	-0.0915957\\
1.9312	-0.0916197\\
1.9337	-0.0916407\\
1.9363	-0.0916607\\
1.9388	-0.0916779\\
1.9414	-0.0916938\\
1.944	-0.0917076\\
1.9465	-0.091719\\
1.9491	-0.0917288\\
1.9517	-0.0917366\\
1.9542	-0.0917422\\
1.9568	-0.0917461\\
1.9594	-0.0917479\\
1.9619	-0.0917478\\
1.9645	-0.0917456\\
1.967	-0.0917417\\
1.9696	-0.0917356\\
1.9722	-0.0917276\\
1.9747	-0.091718\\
1.9773	-0.091706\\
1.9799	-0.0916921\\
1.9824	-0.0916769\\
1.985	-0.0916591\\
1.9876	-0.0916393\\
1.9901	-0.0916185\\
1.9927	-0.091595\\
1.9952	-0.0915705\\
1.9978	-0.0915431\\
2.0004	-0.0915138\\
2.0029	-0.0914839\\
2.0055	-0.0914508\\
2.0081	-0.0914159\\
2.0106	-0.0913805\\
2.0132	-0.0913419\\
2.0158	-0.0913013\\
2.0183	-0.0912606\\
2.0209	-0.0912164\\
2.0234	-0.0911721\\
2.026	-0.0911243\\
2.0286	-0.0910746\\
2.0311	-0.0910251\\
2.0337	-0.0909719\\
2.0363	-0.0909168\\
2.0388	-0.0908622\\
2.0414	-0.0908036\\
2.044	-0.0907432\\
2.0465	-0.0906835\\
2.0491	-0.0906196\\
2.0517	-0.090554\\
2.0542	-0.0904893\\
2.0568	-0.0904203\\
2.0593	-0.0903523\\
2.0619	-0.0902799\\
2.0645	-0.0902058\\
2.067	-0.0901329\\
2.0696	-0.0900555\\
2.0722	-0.0899764\\
2.0747	-0.0898987\\
2.0773	-0.0898164\\
2.0799	-0.0897324\\
2.0824	-0.0896501\\
2.085	-0.0895629\\
2.0875	-0.0894775\\
2.0901	-0.0893872\\
2.0927	-0.0892953\\
2.0952	-0.0892054\\
2.0978	-0.0891104\\
2.1004	-0.0890139\\
2.1029	-0.0889196\\
2.1055	-0.0888201\\
2.1081	-0.088719\\
2.1106	-0.0886204\\
2.1132	-0.0885164\\
2.1157	-0.088415\\
2.1183	-0.0883081\\
2.1209	-0.0881998\\
2.1234	-0.0880942\\
2.126	-0.0879831\\
2.1286	-0.0878705\\
2.1311	-0.087761\\
2.1337	-0.0876456\\
2.1363	-0.087529\\
2.1388	-0.0874155\\
2.1414	-0.0872961\\
2.1439	-0.0871801\\
2.1465	-0.0870582\\
2.1491	-0.0869349\\
2.1516	-0.0868152\\
2.1542	-0.0866894\\
2.1568	-0.0865623\\
2.1593	-0.086439\\
2.1619	-0.0863095\\
2.1645	-0.0861788\\
2.167	-0.0860519\\
2.1696	-0.0859188\\
2.1721	-0.0857897\\
2.1747	-0.0856543\\
2.1773	-0.0855178\\
2.1798	-0.0853854\\
2.1824	-0.0852466\\
2.185	-0.0851067\\
2.1875	-0.0849711\\
2.1901	-0.0848291\\
2.1927	-0.0846859\\
2.1952	-0.0845473\\
2.1978	-0.0844021\\
2.2004	-0.0842559\\
2.2029	-0.0841144\\
2.2055	-0.0839662\\
2.208	-0.0838228\\
2.2106	-0.0836727\\
2.2132	-0.0835217\\
2.2157	-0.0833755\\
2.2183	-0.0832226\\
2.2209	-0.0830689\\
2.2234	-0.0829201\\
2.226	-0.0827646\\
2.2286	-0.0826082\\
2.2311	-0.082457\\
2.2337	-0.082299\\
2.2362	-0.0821462\\
2.2388	-0.0819866\\
2.2414	-0.0818261\\
2.2439	-0.0816711\\
2.2465	-0.0815092\\
2.2491	-0.0813465\\
2.2516	-0.0811894\\
2.2542	-0.0810253\\
2.2568	-0.0808604\\
2.2593	-0.0807013\\
2.2619	-0.0805352\\
2.2644	-0.0803748\\
2.267	-0.0802074\\
2.2696	-0.0800393\\
2.2721	-0.0798772\\
2.2747	-0.079708\\
2.2773	-0.0795381\\
2.2798	-0.0793743\\
2.2824	-0.0792034\\
2.285	-0.079032\\
2.2875	-0.0788666\\
2.2901	-0.0786942\\
2.2926	-0.0785279\\
2.2952	-0.0783545\\
2.2978	-0.0781807\\
2.3003	-0.0780131\\
2.3029	-0.0778383\\
2.3055	-0.0776632\\
2.308	-0.0774944\\
2.3106	-0.0773185\\
2.3132	-0.0771422\\
2.3157	-0.0769723\\
2.3183	-0.0767954\\
2.3208	-0.0766249\\
2.3234	-0.0764473\\
2.326	-0.0762694\\
2.3285	-0.076098\\
2.3311	-0.0759196\\
2.3337	-0.0757408\\
2.3362	-0.0755688\\
2.3388	-0.0753896\\
2.3414	-0.0752102\\
2.3439	-0.0750375\\
2.3465	-0.0748577\\
2.349	-0.0746847\\
2.3516	-0.0745046\\
2.3542	-0.0743244\\
2.3567	-0.074151\\
2.3593	-0.0739705\\
2.3619	-0.0737899\\
2.3644	-0.0736162\\
2.367	-0.0734355\\
2.3696	-0.0732547\\
2.3721	-0.0730808\\
2.3747	-0.0729\\
2.3773	-0.0727191\\
2.3798	-0.0725452\\
2.3824	-0.0723643\\
2.3849	-0.0721904\\
2.3875	-0.0720095\\
2.3901	-0.0718288\\
2.3926	-0.071655\\
2.3952	-0.0714744\\
2.3978	-0.0712938\\
2.4003	-0.0711203\\
2.4029	-0.07094\\
2.4055	-0.0707598\\
2.408	-0.0705866\\
2.4106	-0.0704067\\
2.4131	-0.0702338\\
2.4157	-0.0700542\\
2.4183	-0.0698748\\
2.4208	-0.0697025\\
2.4234	-0.0695236\\
2.426	-0.0693448\\
2.4285	-0.0691731\\
2.4311	-0.0689949\\
2.4337	-0.0688169\\
2.4362	-0.068646\\
2.4388	-0.0684685\\
2.4413	-0.0682982\\
2.4439	-0.0681213\\
2.4465	-0.0679448\\
2.449	-0.0677753\\
2.4516	-0.0675995\\
2.4542	-0.067424\\
2.4567	-0.0672556\\
2.4593	-0.0670808\\
2.4619	-0.0669065\\
2.4644	-0.0667392\\
2.467	-0.0665656\\
2.4695	-0.0663991\\
2.4721	-0.0662264\\
2.4747	-0.0660541\\
2.4772	-0.0658889\\
2.4798	-0.0657175\\
2.4824	-0.0655467\\
2.4849	-0.0653828\\
2.4875	-0.0652129\\
2.4901	-0.0650435\\
2.4926	-0.0648811\\
2.4952	-0.0647127\\
2.4977	-0.0645513\\
2.5003	-0.0643841\\
2.5029	-0.0642173\\
2.5054	-0.0640575\\
2.508	-0.0638919\\
2.5106	-0.0637269\\
2.5131	-0.0635688\\
2.5157	-0.063405\\
2.5183	-0.0632417\\
2.5208	-0.0630854\\
2.5234	-0.0629234\\
2.526	-0.0627621\\
2.5285	-0.0626076\\
2.5311	-0.0624475\\
2.5336	-0.0622943\\
2.5362	-0.0621356\\
2.5388	-0.0619776\\
2.5413	-0.0618263\\
2.5439	-0.0616696\\
2.5465	-0.0615137\\
2.549	-0.0613644\\
2.5516	-0.0612099\\
2.5542	-0.0610562\\
2.5567	-0.060909\\
2.5593	-0.0607567\\
2.5618	-0.060611\\
2.5644	-0.0604602\\
2.567	-0.0603102\\
2.5695	-0.0601666\\
2.5721	-0.0600182\\
2.5747	-0.0598705\\
2.5772	-0.0597292\\
2.5798	-0.0595832\\
2.5824	-0.0594379\\
2.5849	-0.059299\\
2.5875	-0.0591553\\
2.59	-0.059018\\
2.5926	-0.058876\\
2.5952	-0.0587349\\
2.5977	-0.0586\\
2.6003	-0.0584605\\
2.6029	-0.058322\\
2.6054	-0.0581895\\
2.608	-0.0580527\\
2.6106	-0.0579167\\
2.6131	-0.0577869\\
2.6157	-0.0576527\\
2.6182	-0.0575245\\
2.6208	-0.0573921\\
2.6234	-0.0572606\\
2.6259	-0.057135\\
2.6285	-0.0570053\\
2.6311	-0.0568766\\
2.6336	-0.0567536\\
2.6362	-0.0566267\\
2.6388	-0.0565008\\
2.6413	-0.0563806\\
2.6439	-0.0562565\\
2.6464	-0.0561381\\
2.649	-0.0560159\\
2.6516	-0.0558946\\
2.6541	-0.055779\\
2.6567	-0.0556597\\
2.6593	-0.0555414\\
2.6618	-0.0554285\\
2.6644	-0.0553121\\
2.667	-0.0551967\\
2.6695	-0.0550867\\
2.6721	-0.0549733\\
2.6746	-0.0548651\\
2.6772	-0.0547537\\
2.6798	-0.0546432\\
2.6823	-0.054538\\
2.6849	-0.0544295\\
2.6875	-0.0543221\\
2.69	-0.0542197\\
2.6926	-0.0541143\\
2.6952	-0.0540099\\
2.6977	-0.0539105\\
2.7003	-0.0538081\\
2.7029	-0.0537068\\
2.7054	-0.0536103\\
2.708	-0.053511\\
2.7105	-0.0534165\\
2.7131	-0.0533193\\
2.7157	-0.0532231\\
2.7182	-0.0531316\\
2.7208	-0.0530374\\
2.7234	-0.0529443\\
2.7259	-0.0528558\\
2.7285	-0.0527648\\
2.7311	-0.0526749\\
2.7336	-0.0525894\\
2.7362	-0.0525015\\
2.7387	-0.052418\\
2.7413	-0.0523322\\
2.7439	-0.0522475\\
2.7464	-0.052167\\
2.749	-0.0520844\\
2.7516	-0.0520028\\
2.7541	-0.0519254\\
2.7567	-0.051846\\
2.7593	-0.0517676\\
2.7618	-0.0516932\\
2.7644	-0.0516169\\
2.7669	-0.0515445\\
2.7695	-0.0514703\\
2.7721	-0.0513972\\
2.7746	-0.0513279\\
2.7772	-0.0512569\\
2.7798	-0.0511869\\
2.7823	-0.0511207\\
2.7849	-0.0510528\\
2.7875	-0.050986\\
2.79	-0.0509228\\
2.7926	-0.0508581\\
2.7951	-0.050797\\
2.7977	-0.0507344\\
2.8003	-0.0506728\\
2.8028	-0.0506147\\
2.8054	-0.0505553\\
2.808	-0.0504969\\
2.8105	-0.0504418\\
2.8131	-0.0503855\\
2.8157	-0.0503303\\
2.8182	-0.0502782\\
2.8208	-0.0502251\\
2.8233	-0.050175\\
2.8259	-0.050124\\
2.8285	-0.050074\\
2.831	-0.0500269\\
2.8336	-0.049979\\
2.8362	-0.0499321\\
2.8387	-0.049888\\
2.8413	-0.0498432\\
2.8439	-0.0497994\\
2.8464	-0.0497583\\
2.849	-0.0497165\\
2.8516	-0.0496758\\
2.8541	-0.0496377\\
2.8567	-0.049599\\
2.8592	-0.0495628\\
2.8618	-0.0495262\\
2.8644	-0.0494906\\
2.8669	-0.0494573\\
2.8695	-0.0494237\\
2.8721	-0.0493911\\
2.8746	-0.0493608\\
2.8772	-0.0493302\\
2.8798	-0.0493006\\
2.8823	-0.0492731\\
2.8849	-0.0492455\\
2.8874	-0.0492199\\
2.89	-0.0491943\\
2.8926	-0.0491697\\
2.8951	-0.0491469\\
2.8977	-0.0491242\\
2.9003	-0.0491025\\
2.9028	-0.0490826\\
2.9054	-0.0490628\\
2.908	-0.049044\\
2.9105	-0.0490269\\
2.9131	-0.04901\\
2.9156	-0.0489947\\
2.9182	-0.0489797\\
2.9208	-0.0489656\\
2.9233	-0.048953\\
2.9259	-0.0489409\\
2.9285	-0.0489297\\
2.931	-0.0489198\\
2.9336	-0.0489104\\
2.9362	-0.048902\\
2.9387	-0.0488947\\
2.9413	-0.0488881\\
2.9438	-0.0488826\\
2.9464	-0.0488778\\
2.949	-0.048874\\
2.9515	-0.0488711\\
2.9541	-0.048869\\
2.9567	-0.0488678\\
2.9592	-0.0488675\\
2.9618	-0.048868\\
2.9644	-0.0488695\\
2.9669	-0.0488717\\
2.9695	-0.0488749\\
2.972	-0.0488788\\
2.9746	-0.0488836\\
2.9772	-0.0488894\\
2.9797	-0.0488957\\
2.9823	-0.0489032\\
2.9849	-0.0489115\\
2.9874	-0.0489202\\
2.99	-0.0489302\\
2.9926	-0.0489409\\
2.9951	-0.0489521\\
2.9977	-0.0489645\\
3.0003	-0.0489777\\
3.0028	-0.0489912\\
3.0054	-0.049006\\
3.0079	-0.049021\\
3.0105	-0.0490374\\
3.0131	-0.0490546\\
3.0156	-0.0490719\\
3.0182	-0.0490906\\
3.0208	-0.0491101\\
3.0233	-0.0491296\\
3.0259	-0.0491506\\
3.0285	-0.0491724\\
3.031	-0.049194\\
3.0336	-0.0492173\\
3.0361	-0.0492403\\
3.0387	-0.0492651\\
3.0413	-0.0492905\\
3.0438	-0.0493157\\
3.0464	-0.0493425\\
3.049	-0.0493701\\
3.0515	-0.0493973\\
3.0541	-0.0494263\\
3.0567	-0.0494559\\
3.0592	-0.0494851\\
3.0618	-0.0495162\\
3.0643	-0.0495466\\
3.0669	-0.049579\\
3.0695	-0.049612\\
3.072	-0.0496444\\
3.0746	-0.0496787\\
3.0772	-0.0497137\\
3.0797	-0.0497479\\
3.0823	-0.0497842\\
3.0849	-0.049821\\
3.0874	-0.0498571\\
3.09	-0.0498952\\
3.0925	-0.0499324\\
3.0951	-0.0499717\\
3.0977	-0.0500116\\
3.1002	-0.0500505\\
3.1028	-0.0500916\\
3.1054	-0.0501332\\
3.1079	-0.0501738\\
3.1105	-0.0502166\\
3.1131	-0.0502599\\
3.1156	-0.0503021\\
3.1182	-0.0503465\\
3.1207	-0.0503898\\
3.1233	-0.0504353\\
3.1259	-0.0504813\\
3.1284	-0.0505261\\
3.131	-0.0505731\\
3.1336	-0.0506207\\
3.1361	-0.0506669\\
3.1387	-0.0507155\\
3.1413	-0.0507645\\
3.1438	-0.0508121\\
3.1464	-0.0508621\\
3.1489	-0.0509107\\
3.1515	-0.0509616\\
3.1541	-0.051013\\
3.1566	-0.0510628\\
3.1592	-0.0511151\\
3.1618	-0.0511679\\
3.1643	-0.051219\\
3.1669	-0.0512725\\
3.1695	-0.0513265\\
3.172	-0.0513789\\
3.1746	-0.0514337\\
3.1772	-0.0514889\\
3.1797	-0.0515423\\
3.1823	-0.0515983\\
3.1848	-0.0516525\\
3.1874	-0.0517093\\
3.19	-0.0517664\\
3.1925	-0.0518216\\
3.1951	-0.0518794\\
3.1977	-0.0519376\\
3.2002	-0.0519939\\
3.2028	-0.0520527\\
3.2054	-0.0521119\\
3.2079	-0.0521691\\
3.2105	-0.0522289\\
3.213	-0.0522866\\
3.2156	-0.0523471\\
3.2182	-0.0524078\\
3.2207	-0.0524664\\
3.2233	-0.0525277\\
3.2259	-0.0525892\\
3.2284	-0.0526487\\
3.231	-0.0527107\\
3.2336	-0.0527731\\
3.2361	-0.0528332\\
3.2387	-0.0528961\\
3.2412	-0.0529567\\
3.2438	-0.05302\\
3.2464	-0.0530835\\
3.2489	-0.0531448\\
3.2515	-0.0532087\\
3.2541	-0.0532729\\
3.2566	-0.0533348\\
3.2592	-0.0533993\\
3.2618	-0.053464\\
3.2643	-0.0535264\\
3.2669	-0.0535915\\
3.2694	-0.0536542\\
3.272	-0.0537196\\
3.2746	-0.0537852\\
3.2771	-0.0538484\\
3.2797	-0.0539142\\
3.2823	-0.0539802\\
3.2848	-0.0540437\\
3.2874	-0.05411\\
3.29	-0.0541763\\
3.2925	-0.0542402\\
3.2951	-0.0543068\\
3.2976	-0.0543709\\
3.3002	-0.0544376\\
3.3028	-0.0545044\\
3.3053	-0.0545688\\
3.3079	-0.0546358\\
3.3105	-0.0547028\\
3.313	-0.0547673\\
3.3156	-0.0548345\\
3.3182	-0.0549017\\
3.3207	-0.0549664\\
3.3233	-0.0550337\\
3.3259	-0.055101\\
3.3284	-0.0551658\\
3.331	-0.0552331\\
3.3335	-0.0552979\\
3.3361	-0.0553653\\
3.3387	-0.0554327\\
3.3412	-0.0554975\\
3.3438	-0.0555648\\
3.3464	-0.0556322\\
3.3489	-0.0556969\\
3.3515	-0.0557642\\
3.3541	-0.0558315\\
3.3566	-0.0558961\\
3.3592	-0.0559633\\
3.3617	-0.0560278\\
3.3643	-0.0560949\\
3.3669	-0.0561619\\
3.3694	-0.0562263\\
3.372	-0.0562931\\
3.3746	-0.0563599\\
3.3771	-0.056424\\
3.3797	-0.0564906\\
3.3823	-0.0565571\\
3.3848	-0.056621\\
3.3874	-0.0566873\\
3.3899	-0.0567509\\
3.3925	-0.056817\\
3.3951	-0.0568829\\
3.3976	-0.0569462\\
3.4002	-0.0570119\\
3.4028	-0.0570774\\
3.4053	-0.0571403\\
3.4079	-0.0572056\\
3.4105	-0.0572706\\
3.413	-0.0573331\\
3.4156	-0.0573978\\
3.4181	-0.0574599\\
3.4207	-0.0575244\\
3.4233	-0.0575886\\
3.4258	-0.0576502\\
3.4284	-0.057714\\
3.431	-0.0577777\\
3.4335	-0.0578387\\
3.4361	-0.057902\\
3.4387	-0.057965\\
3.4412	-0.0580254\\
3.4438	-0.058088\\
3.4463	-0.058148\\
3.4489	-0.0582102\\
3.4515	-0.0582721\\
3.454	-0.0583314\\
3.4566	-0.0583928\\
3.4592	-0.058454\\
3.4617	-0.0585126\\
3.4643	-0.0585733\\
3.4669	-0.0586337\\
3.4694	-0.0586915\\
3.472	-0.0587514\\
3.4745	-0.0588087\\
3.4771	-0.058868\\
3.4797	-0.058927\\
3.4822	-0.0589835\\
3.4848	-0.059042\\
3.4874	-0.0591002\\
3.4899	-0.0591558\\
3.4925	-0.0592133\\
3.4951	-0.0592706\\
3.4976	-0.0593253\\
3.5002	-0.0593819\\
3.5028	-0.0594382\\
3.5053	-0.059492\\
3.5079	-0.0595477\\
3.5104	-0.0596009\\
3.513	-0.0596558\\
3.5156	-0.0597105\\
3.5181	-0.0597627\\
3.5207	-0.0598166\\
3.5233	-0.0598702\\
3.5258	-0.0599214\\
3.5284	-0.0599742\\
3.531	-0.0600267\\
3.5335	-0.0600769\\
3.5361	-0.0601287\\
3.5386	-0.0601781\\
3.5412	-0.0602291\\
3.5438	-0.0602798\\
3.5463	-0.0603281\\
3.5489	-0.060378\\
3.5515	-0.0604275\\
3.554	-0.0604747\\
3.5566	-0.0605234\\
3.5592	-0.0605717\\
3.5617	-0.0606178\\
3.5643	-0.0606653\\
3.5668	-0.0607105\\
3.5694	-0.0607572\\
3.572	-0.0608035\\
3.5745	-0.0608476\\
3.5771	-0.060893\\
3.5797	-0.060938\\
3.5822	-0.0609809\\
3.5848	-0.0610251\\
3.5874	-0.0610688\\
3.5899	-0.0611105\\
3.5925	-0.0611533\\
3.595	-0.0611941\\
3.5976	-0.0612361\\
3.6002	-0.0612777\\
3.6027	-0.0613172\\
3.6053	-0.0613579\\
3.6079	-0.0613981\\
3.6104	-0.0614364\\
3.613	-0.0614757\\
3.6156	-0.0615145\\
3.6181	-0.0615515\\
3.6207	-0.0615894\\
3.6232	-0.0616254\\
3.6258	-0.0616625\\
3.6284	-0.061699\\
3.6309	-0.0617337\\
3.6335	-0.0617694\\
3.6361	-0.0618045\\
3.6386	-0.0618379\\
3.6412	-0.0618721\\
3.6438	-0.0619058\\
3.6463	-0.0619378\\
3.6489	-0.0619706\\
3.6515	-0.0620029\\
3.654	-0.0620335\\
3.6566	-0.0620648\\
3.6591	-0.0620945\\
3.6617	-0.0621248\\
3.6643	-0.0621547\\
3.6668	-0.062183\\
3.6694	-0.0622119\\
3.672	-0.0622403\\
3.6745	-0.0622672\\
3.6771	-0.0622946\\
3.6797	-0.0623215\\
3.6822	-0.062347\\
3.6848	-0.0623729\\
3.6873	-0.0623974\\
3.6899	-0.0624224\\
3.6925	-0.0624468\\
3.695	-0.0624699\\
3.6976	-0.0624934\\
3.7002	-0.0625163\\
3.7027	-0.062538\\
3.7053	-0.0625599\\
3.7079	-0.0625814\\
3.7104	-0.0626016\\
3.713	-0.062622\\
3.7155	-0.0626413\\
3.7181	-0.0626607\\
3.7207	-0.0626797\\
3.7232	-0.0626975\\
3.7258	-0.0627154\\
3.7284	-0.0627329\\
3.7309	-0.0627492\\
3.7335	-0.0627656\\
3.7361	-0.0627816\\
3.7386	-0.0627965\\
3.7412	-0.0628114\\
3.7437	-0.0628253\\
3.7463	-0.0628392\\
3.7489	-0.0628527\\
3.7514	-0.0628651\\
3.754	-0.0628776\\
3.7566	-0.0628895\\
3.7591	-0.0629005\\
3.7617	-0.0629115\\
3.7643	-0.0629219\\
3.7668	-0.0629314\\
3.7694	-0.0629409\\
3.7719	-0.0629495\\
3.7745	-0.0629579\\
3.7771	-0.0629659\\
3.7796	-0.062973\\
3.7822	-0.06298\\
3.7848	-0.0629865\\
3.7873	-0.0629922\\
3.7899	-0.0629977\\
3.7925	-0.0630027\\
3.795	-0.063007\\
3.7976	-0.063011\\
3.8002	-0.0630145\\
3.8027	-0.0630174\\
3.8053	-0.0630199\\
3.8078	-0.0630219\\
3.8104	-0.0630235\\
3.813	-0.0630246\\
3.8155	-0.0630252\\
3.8181	-0.0630253\\
3.8207	-0.0630249\\
3.8232	-0.0630241\\
3.8258	-0.0630228\\
3.8284	-0.063021\\
3.8309	-0.0630189\\
3.8335	-0.0630161\\
3.836	-0.0630131\\
3.8386	-0.0630094\\
3.8412	-0.0630053\\
3.8437	-0.0630008\\
3.8463	-0.0629958\\
3.8489	-0.0629903\\
3.8514	-0.0629845\\
3.854	-0.0629781\\
3.8566	-0.0629712\\
3.8591	-0.0629641\\
3.8617	-0.0629563\\
3.8642	-0.0629483\\
3.8668	-0.0629396\\
3.8694	-0.0629305\\
3.8719	-0.0629212\\
3.8745	-0.0629112\\
3.8771	-0.0629007\\
3.8796	-0.0628902\\
3.8822	-0.0628789\\
3.8848	-0.0628671\\
3.8873	-0.0628553\\
3.8899	-0.0628427\\
3.8924	-0.0628301\\
3.895	-0.0628166\\
3.8976	-0.0628027\\
3.9001	-0.0627889\\
3.9027	-0.0627741\\
3.9053	-0.0627589\\
3.9078	-0.0627439\\
3.9104	-0.0627279\\
3.913	-0.0627115\\
3.9155	-0.0626953\\
3.9181	-0.0626781\\
3.9206	-0.0626612\\
3.9232	-0.0626432\\
3.9258	-0.0626248\\
3.9283	-0.0626067\\
3.9309	-0.0625875\\
3.9335	-0.0625679\\
3.936	-0.0625487\\
3.9386	-0.0625284\\
3.9412	-0.0625076\\
3.9437	-0.0624873\\
3.9463	-0.0624659\\
3.9488	-0.0624449\\
3.9514	-0.0624227\\
3.954	-0.0624001\\
3.9565	-0.062378\\
3.9591	-0.0623548\\
3.9617	-0.0623311\\
3.9642	-0.062308\\
3.9668	-0.0622837\\
3.9694	-0.062259\\
3.9719	-0.0622349\\
3.9745	-0.0622095\\
3.9771	-0.0621838\\
3.9796	-0.0621587\\
3.9822	-0.0621324\\
3.9847	-0.0621067\\
3.9873	-0.0620797\\
3.9899	-0.0620523\\
3.9924	-0.0620257\\
3.995	-0.0619977\\
3.9976	-0.0619694\\
4.0001	-0.0619419\\
4.0027	-0.061913\\
4.0053	-0.0618838\\
4.0078	-0.0618554\\
4.0104	-0.0618256\\
4.0129	-0.0617966\\
4.0155	-0.0617663\\
4.0181	-0.0617356\\
4.0206	-0.0617058\\
4.0232	-0.0616746\\
4.0258	-0.0616431\\
4.0283	-0.0616125\\
4.0309	-0.0615805\\
4.0335	-0.0615482\\
4.036	-0.0615168\\
4.0386	-0.061484\\
4.0411	-0.0614522\\
4.0437	-0.0614189\\
4.0463	-0.0613853\\
4.0488	-0.0613527\\
4.0514	-0.0613187\\
4.054	-0.0612844\\
4.0565	-0.0612512\\
4.0591	-0.0612164\\
4.0617	-0.0611814\\
4.0642	-0.0611475\\
4.0668	-0.0611121\\
4.0693	-0.0610778\\
4.0719	-0.0610419\\
4.0745	-0.0610059\\
4.077	-0.060971\\
4.0796	-0.0609345\\
4.0822	-0.0608978\\
4.0847	-0.0608624\\
4.0873	-0.0608253\\
4.0899	-0.0607881\\
4.0924	-0.0607521\\
4.095	-0.0607144\\
4.0975	-0.0606781\\
4.1001	-0.0606402\\
4.1027	-0.060602\\
4.1052	-0.0605652\\
4.1078	-0.0605267\\
4.1104	-0.0604881\\
4.1129	-0.0604508\\
4.1155	-0.0604119\\
4.1181	-0.0603728\\
4.1206	-0.0603351\\
4.1232	-0.0602958\\
4.1258	-0.0602563\\
4.1283	-0.0602182\\
4.1309	-0.0601784\\
4.1334	-0.0601401\\
4.136	-0.0601001\\
4.1386	-0.06006\\
4.1411	-0.0600213\\
4.1437	-0.0599809\\
4.1463	-0.0599405\\
4.1488	-0.0599015\\
4.1514	-0.0598608\\
4.154	-0.0598201\\
4.1565	-0.0597808\\
4.1591	-0.0597398\\
4.1616	-0.0597004\\
4.1642	-0.0596593\\
4.1668	-0.0596181\\
4.1693	-0.0595784\\
4.1719	-0.0595371\\
4.1745	-0.0594957\\
4.177	-0.0594558\\
4.1796	-0.0594143\\
4.1822	-0.0593727\\
4.1847	-0.0593327\\
4.1873	-0.059291\\
4.1898	-0.0592508\\
4.1924	-0.059209\\
4.195	-0.0591672\\
4.1975	-0.059127\\
4.2001	-0.0590851\\
4.2027	-0.0590431\\
4.2052	-0.0590028\\
4.2078	-0.0589608\\
4.2104	-0.0589188\\
4.2129	-0.0588784\\
4.2155	-0.0588364\\
4.218	-0.058796\\
4.2206	-0.0587539\\
4.2232	-0.0587119\\
4.2257	-0.0586715\\
4.2283	-0.0586294\\
4.2309	-0.0585874\\
4.2334	-0.058547\\
4.236	-0.058505\\
4.2386	-0.058463\\
4.2411	-0.0584226\\
4.2437	-0.0583807\\
4.2462	-0.0583404\\
4.2488	-0.0582985\\
4.2514	-0.0582566\\
4.2539	-0.0582164\\
4.2565	-0.0581747\\
4.2591	-0.0581329\\
4.2616	-0.0580928\\
4.2642	-0.0580512\\
4.2668	-0.0580096\\
4.2693	-0.0579697\\
4.2719	-0.0579282\\
4.2744	-0.0578884\\
4.277	-0.057847\\
4.2796	-0.0578057\\
4.2821	-0.0577661\\
4.2847	-0.057725\\
4.2873	-0.0576839\\
4.2898	-0.0576445\\
4.2924	-0.0576036\\
4.295	-0.0575628\\
4.2975	-0.0575236\\
4.3001	-0.057483\\
4.3027	-0.0574425\\
4.3052	-0.0574036\\
4.3078	-0.0573632\\
4.3103	-0.0573246\\
4.3129	-0.0572844\\
4.3155	-0.0572444\\
4.318	-0.057206\\
4.3206	-0.0571662\\
4.3232	-0.0571265\\
4.3257	-0.0570885\\
4.3283	-0.0570491\\
4.3309	-0.0570098\\
4.3334	-0.0569721\\
4.336	-0.056933\\
4.3385	-0.0568956\\
4.3411	-0.0568568\\
4.3437	-0.0568182\\
4.3462	-0.0567811\\
4.3488	-0.0567428\\
4.3514	-0.0567046\\
4.3539	-0.056668\\
4.3565	-0.05663\\
4.3591	-0.0565923\\
4.3616	-0.0565561\\
4.3642	-0.0565187\\
4.3667	-0.0564828\\
4.3693	-0.0564457\\
4.3719	-0.0564087\\
4.3744	-0.0563733\\
4.377	-0.0563367\\
4.3796	-0.0563003\\
4.3821	-0.0562654\\
4.3847	-0.0562293\\
4.3873	-0.0561933\\
4.3898	-0.056159\\
4.3924	-0.0561234\\
4.3949	-0.0560894\\
4.3975	-0.0560542\\
4.4001	-0.0560192\\
4.4026	-0.0559858\\
4.4052	-0.0559511\\
4.4078	-0.0559167\\
4.4103	-0.0558838\\
4.4129	-0.0558498\\
4.4155	-0.055816\\
4.418	-0.0557837\\
4.4206	-0.0557503\\
4.4231	-0.0557183\\
4.4257	-0.0556854\\
4.4283	-0.0556526\\
4.4308	-0.0556213\\
4.4334	-0.0555889\\
4.436	-0.0555568\\
4.4385	-0.0555261\\
4.4411	-0.0554944\\
4.4437	-0.0554629\\
4.4462	-0.0554329\\
4.4488	-0.0554018\\
4.4514	-0.055371\\
4.4539	-0.0553417\\
4.4565	-0.0553113\\
4.459	-0.0552824\\
4.4616	-0.0552525\\
4.4642	-0.0552228\\
4.4667	-0.0551946\\
4.4693	-0.0551654\\
4.4719	-0.0551365\\
4.4744	-0.0551089\\
4.477	-0.0550804\\
4.4796	-0.0550522\\
4.4821	-0.0550253\\
4.4847	-0.0549976\\
4.4872	-0.0549712\\
4.4898	-0.054944\\
4.4924	-0.054917\\
4.4949	-0.0548913\\
4.4975	-0.0548649\\
4.5001	-0.0548386\\
4.5026	-0.0548137\\
4.5052	-0.054788\\
4.5078	-0.0547625\\
4.5103	-0.0547383\\
4.5129	-0.0547134\\
4.5154	-0.0546896\\
4.518	-0.0546652\\
4.5206	-0.0546411\\
4.5231	-0.0546181\\
4.5257	-0.0545945\\
4.5283	-0.0545711\\
4.5308	-0.0545489\\
4.5334	-0.0545261\\
4.536	-0.0545035\\
4.5385	-0.0544821\\
4.5411	-0.05446\\
4.5436	-0.0544391\\
4.5462	-0.0544176\\
4.5488	-0.0543964\\
4.5513	-0.0543762\\
4.5539	-0.0543556\\
4.5565	-0.0543352\\
4.559	-0.0543158\\
4.5616	-0.0542959\\
4.5642	-0.0542764\\
4.5667	-0.0542578\\
4.5693	-0.0542387\\
4.5718	-0.0542207\\
4.5744	-0.0542022\\
4.577	-0.054184\\
4.5795	-0.0541668\\
4.5821	-0.0541491\\
4.5847	-0.0541317\\
4.5872	-0.0541153\\
4.5898	-0.0540985\\
4.5924	-0.0540819\\
4.5949	-0.0540663\\
4.5975	-0.0540503\\
4.6001	-0.0540346\\
4.6026	-0.0540198\\
4.6052	-0.0540046\\
4.6077	-0.0539903\\
4.6103	-0.0539757\\
4.6129	-0.0539614\\
4.6154	-0.0539479\\
4.618	-0.0539342\\
4.6206	-0.0539207\\
4.6231	-0.053908\\
4.6257	-0.0538951\\
4.6283	-0.0538825\\
4.6308	-0.0538706\\
4.6334	-0.0538586\\
4.6359	-0.0538472\\
4.6385	-0.0538357\\
4.6411	-0.0538245\\
4.6436	-0.053814\\
4.6462	-0.0538033\\
4.6488	-0.0537929\\
4.6513	-0.0537832\\
4.6539	-0.0537733\\
4.6565	-0.0537638\\
4.659	-0.0537549\\
4.6616	-0.0537459\\
4.6641	-0.0537375\\
4.6667	-0.053729\\
4.6693	-0.0537209\\
4.6718	-0.0537133\\
4.6744	-0.0537057\\
4.677	-0.0536983\\
4.6795	-0.0536915\\
4.6821	-0.0536847\\
4.6847	-0.0536782\\
4.6872	-0.0536722\\
4.6898	-0.0536663\\
4.6923	-0.0536608\\
4.6949	-0.0536553\\
4.6975	-0.0536502\\
4.7	-0.0536455\\
4.7026	-0.0536409\\
4.7052	-0.0536365\\
4.7077	-0.0536326\\
4.7103	-0.0536288\\
4.7129	-0.0536253\\
4.7154	-0.0536221\\
4.718	-0.0536191\\
4.7205	-0.0536165\\
4.7231	-0.053614\\
4.7257	-0.0536118\\
4.7282	-0.0536099\\
4.7308	-0.0536082\\
4.7334	-0.0536067\\
4.7359	-0.0536056\\
4.7385	-0.0536047\\
4.7411	-0.053604\\
4.7436	-0.0536037\\
4.7462	-0.0536035\\
4.7487	-0.0536036\\
4.7513	-0.053604\\
4.7539	-0.0536046\\
4.7564	-0.0536054\\
4.759	-0.0536066\\
4.7616	-0.0536079\\
4.7641	-0.0536095\\
4.7667	-0.0536114\\
4.7693	-0.0536135\\
4.7718	-0.0536158\\
4.7744	-0.0536184\\
4.777	-0.0536212\\
4.7795	-0.0536242\\
4.7821	-0.0536275\\
4.7846	-0.053631\\
4.7872	-0.0536348\\
4.7898	-0.0536388\\
4.7923	-0.0536429\\
4.7949	-0.0536475\\
4.7975	-0.0536522\\
4.8	-0.053657\\
4.8026	-0.0536622\\
4.8052	-0.0536677\\
4.8077	-0.0536731\\
4.8103	-0.053679\\
4.8128	-0.0536849\\
4.8154	-0.0536913\\
4.818	-0.0536979\\
4.8205	-0.0537044\\
4.8231	-0.0537114\\
4.8257	-0.0537187\\
4.8282	-0.0537258\\
4.8308	-0.0537335\\
4.8334	-0.0537414\\
4.8359	-0.0537492\\
4.8385	-0.0537575\\
4.841	-0.0537657\\
4.8436	-0.0537744\\
4.8462	-0.0537833\\
4.8487	-0.0537921\\
4.8513	-0.0538015\\
4.8539	-0.053811\\
4.8564	-0.0538204\\
4.859	-0.0538304\\
4.8616	-0.0538405\\
4.8641	-0.0538505\\
4.8667	-0.053861\\
4.8692	-0.0538713\\
4.8718	-0.0538822\\
4.8744	-0.0538933\\
4.8769	-0.0539042\\
4.8795	-0.0539157\\
4.8821	-0.0539274\\
4.8846	-0.0539388\\
4.8872	-0.0539508\\
4.8898	-0.053963\\
4.8923	-0.0539749\\
4.8949	-0.0539875\\
4.8974	-0.0539997\\
4.9	-0.0540127\\
4.9026	-0.0540257\\
4.9051	-0.0540385\\
4.9077	-0.0540519\\
4.9103	-0.0540655\\
4.9128	-0.0540787\\
4.9154	-0.0540926\\
4.918	-0.0541067\\
4.9205	-0.0541204\\
4.9231	-0.0541348\\
4.9257	-0.0541494\\
4.9282	-0.0541635\\
4.9308	-0.0541783\\
4.9333	-0.0541928\\
4.9359	-0.0542079\\
4.9385	-0.0542232\\
4.941	-0.0542381\\
4.9436	-0.0542536\\
4.9462	-0.0542694\\
4.9487	-0.0542846\\
4.9513	-0.0543006\\
4.9539	-0.0543168\\
4.9564	-0.0543324\\
4.959	-0.0543488\\
4.9615	-0.0543647\\
4.9641	-0.0543813\\
4.9667	-0.0543981\\
4.9692	-0.0544144\\
4.9718	-0.0544314\\
4.9744	-0.0544485\\
4.9769	-0.0544651\\
4.9795	-0.0544825\\
4.9821	-0.0545\\
4.9846	-0.0545169\\
4.9872	-0.0545347\\
4.9897	-0.0545518\\
4.9923	-0.0545697\\
4.9949	-0.0545878\\
4.9974	-0.0546052\\
5	-0.0546234\\
};
\addplot [color=mycolor18, line width=1.0pt, forget plot]
  table[row sep=crcr]{%
-0.125	6.84181e-08\\
-0.1224	6.98348e-08\\
-0.1199	7.1197e-08\\
-0.1173	7.26135e-08\\
-0.1147	7.40299e-08\\
-0.1122	7.53917e-08\\
-0.1096	7.6808e-08\\
-0.1071	7.81695e-08\\
-0.1045	7.95855e-08\\
-0.1019	8.10015e-08\\
-0.0994	8.23628e-08\\
-0.0968	8.37785e-08\\
-0.0942	8.51942e-08\\
-0.0917	8.65552e-08\\
-0.0891	8.79706e-08\\
-0.0865	8.93859e-08\\
-0.084	9.07467e-08\\
-0.0814	9.21617e-08\\
-0.0789	9.35223e-08\\
-0.0763	9.49372e-08\\
-0.0737	9.63519e-08\\
-0.0712	9.77122e-08\\
-0.0686	9.91268e-08\\
-0.066	1.00541e-07\\
-0.0635	1.01901e-07\\
-0.0609	1.03316e-07\\
-0.0583	1.0473e-07\\
-0.0558	1.06089e-07\\
-0.0532	1.07503e-07\\
-0.0507	1.08863e-07\\
-0.0481	1.10276e-07\\
-0.0455	1.1169e-07\\
-0.043	1.13049e-07\\
-0.0404	1.14463e-07\\
-0.0378	1.15876e-07\\
-0.0353	1.17235e-07\\
-0.0327	1.18648e-07\\
-0.0301	1.20061e-07\\
-0.0276	1.21419e-07\\
-0.025	1.22832e-07\\
-0.0224	1.24245e-07\\
-0.0199	1.25603e-07\\
-0.0173	1.27016e-07\\
-0.0148	1.28374e-07\\
-0.0122	1.29786e-07\\
-0.0096	1.31198e-07\\
-0.0071	1.32556e-07\\
-0.0045	1.33968e-07\\
-0.0019	1.3538e-07\\
0.0006	1.36737e-07\\
0.0032	6.72664e-05\\
0.0058	0.000667848\\
0.0083	0.00132863\\
0.0109	0.00189298\\
0.0134	0.00243878\\
0.016	0.00307252\\
0.0186	0.00373769\\
0.0211	0.00434312\\
0.0237	0.00490471\\
0.0263	0.00541561\\
0.0288	0.00589961\\
0.0314	0.00642516\\
0.034	0.00697364\\
0.0365	0.00750333\\
0.0391	0.00803658\\
0.0416	0.00852735\\
0.0442	0.00902457\\
0.0468	0.0095226\\
0.0493	0.0100089\\
0.0519	0.0105182\\
0.0545	0.0110208\\
0.057	0.0114908\\
0.0596	0.0119659\\
0.0622	0.0124342\\
0.0647	0.0128844\\
0.0673	0.013354\\
0.0698	0.0138019\\
0.0724	0.0142573\\
0.075	0.0146989\\
0.0775	0.015113\\
0.0801	0.0155386\\
0.0827	0.0159628\\
0.0852	0.0163679\\
0.0878	0.0167809\\
0.0904	0.0171809\\
0.0929	0.017553\\
0.0955	0.0179311\\
0.098	0.0182909\\
0.1006	0.0186622\\
0.1032	0.0190274\\
0.1057	0.019368\\
0.1083	0.019709\\
0.1109	0.0200388\\
0.1134	0.0203499\\
0.116	0.0206702\\
0.1186	0.0209863\\
0.1211	0.0212825\\
0.1237	0.021579\\
0.1263	0.0218637\\
0.1288	0.0221299\\
0.1314	0.0224031\\
0.1339	0.0226633\\
0.1365	0.022929\\
0.1391	0.0231858\\
0.1416	0.0234227\\
0.1442	0.0236606\\
0.1468	0.023894\\
0.1493	0.0241168\\
0.1519	0.0243464\\
0.1545	0.0245704\\
0.157	0.0247779\\
0.1596	0.0249856\\
0.1621	0.0251806\\
0.1647	0.0253819\\
0.1673	0.0255828\\
0.1698	0.0257738\\
0.1724	0.0259669\\
0.175	0.0261531\\
0.1775	0.0263271\\
0.1801	0.0265062\\
0.1827	0.0266858\\
0.1852	0.0268583\\
0.1878	0.027035\\
0.1903	0.0272\\
0.1929	0.0273664\\
0.1955	0.0275303\\
0.198	0.0276881\\
0.2006	0.0278533\\
0.2032	0.0280179\\
0.2057	0.028173\\
0.2083	0.0283297\\
0.2109	0.0284831\\
0.2134	0.0286301\\
0.216	0.0287841\\
0.2185	0.0289329\\
0.2211	0.0290859\\
0.2237	0.0292352\\
0.2262	0.0293751\\
0.2288	0.0295187\\
0.2314	0.0296628\\
0.2339	0.0298023\\
0.2365	0.0299469\\
0.2391	0.0300886\\
0.2416	0.0302211\\
0.2442	0.0303558\\
0.2467	0.0304845\\
0.2493	0.0306189\\
0.2519	0.0307533\\
0.2544	0.0308806\\
0.257	0.0310091\\
0.2596	0.0311337\\
0.2621	0.0312512\\
0.2647	0.0313728\\
0.2673	0.0314943\\
0.2698	0.0316098\\
0.2724	0.0317264\\
0.2749	0.0318342\\
0.2775	0.0319427\\
0.2801	0.0320492\\
0.2826	0.0321509\\
0.2852	0.0322554\\
0.2878	0.0323569\\
0.2903	0.0324503\\
0.2929	0.0325427\\
0.2955	0.0326319\\
0.298	0.0327159\\
0.3006	0.0328019\\
0.3032	0.0328854\\
0.3057	0.0329616\\
0.3083	0.033036\\
0.3108	0.0331033\\
0.3134	0.0331706\\
0.316	0.0332359\\
0.3185	0.0332966\\
0.3211	0.0333562\\
0.3237	0.0334109\\
0.3262	0.0334587\\
0.3288	0.0335048\\
0.3314	0.0335483\\
0.3339	0.0335882\\
0.3365	0.0336266\\
0.339	0.0336594\\
0.3416	0.0336886\\
0.3442	0.0337134\\
0.3467	0.0337342\\
0.3493	0.0337537\\
0.3519	0.0337706\\
0.3544	0.0337834\\
0.357	0.0337922\\
0.3596	0.0337963\\
0.3621	0.0337968\\
0.3647	0.0337949\\
0.3672	0.0337909\\
0.3698	0.033784\\
0.3724	0.0337731\\
0.3749	0.0337584\\
0.3775	0.0337393\\
0.3801	0.0337173\\
0.3826	0.0336942\\
0.3852	0.033668\\
0.3878	0.0336387\\
0.3903	0.0336067\\
0.3929	0.0335696\\
0.3954	0.033531\\
0.398	0.0334889\\
0.4006	0.0334449\\
0.4031	0.0334004\\
0.4057	0.0333509\\
0.4083	0.0332977\\
0.4108	0.0332436\\
0.4134	0.0331852\\
0.416	0.0331252\\
0.4185	0.0330659\\
0.4211	0.0330017\\
0.4236	0.0329371\\
0.4262	0.0328669\\
0.4288	0.0327943\\
0.4313	0.032723\\
0.4339	0.0326476\\
0.4365	0.0325705\\
0.439	0.0324939\\
0.4416	0.0324116\\
0.4442	0.0323269\\
0.4467	0.0322439\\
0.4493	0.0321565\\
0.4519	0.0320679\\
0.4544	0.031981\\
0.457	0.0318884\\
0.4595	0.0317971\\
0.4621	0.0317004\\
0.4647	0.0316028\\
0.4672	0.0315081\\
0.4698	0.0314085\\
0.4724	0.031307\\
0.4749	0.0312075\\
0.4775	0.0311023\\
0.4801	0.030996\\
0.4826	0.0308932\\
0.4852	0.0307856\\
0.4877	0.030681\\
0.4903	0.0305705\\
0.4929	0.0304584\\
0.4954	0.0303494\\
0.498	0.0302356\\
0.5006	0.0301213\\
0.5031	0.0300108\\
0.5057	0.0298946\\
0.5083	0.0297769\\
0.5108	0.0296627\\
0.5134	0.0295433\\
0.5159	0.0294283\\
0.5185	0.0293083\\
0.5211	0.0291876\\
0.5236	0.0290705\\
0.5262	0.0289476\\
0.5288	0.0288241\\
0.5313	0.0287051\\
0.5339	0.0285813\\
0.5365	0.0284572\\
0.539	0.0283371\\
0.5416	0.0282114\\
0.5441	0.0280899\\
0.5467	0.0279633\\
0.5493	0.0278367\\
0.5518	0.0277151\\
0.5544	0.0275883\\
0.557	0.0274608\\
0.5595	0.0273377\\
0.5621	0.0272094\\
0.5647	0.0270812\\
0.5672	0.0269582\\
0.5698	0.0268303\\
0.5723	0.026707\\
0.5749	0.0265783\\
0.5775	0.0264493\\
0.58	0.0263254\\
0.5826	0.0261969\\
0.5852	0.0260687\\
0.5877	0.0259455\\
0.5903	0.025817\\
0.5929	0.0256883\\
0.5954	0.0255646\\
0.598	0.0254364\\
0.6006	0.0253086\\
0.6031	0.0251859\\
0.6057	0.0250584\\
0.6082	0.0249356\\
0.6108	0.0248079\\
0.6134	0.0246805\\
0.6159	0.0245585\\
0.6185	0.0244321\\
0.6211	0.0243059\\
0.6236	0.0241845\\
0.6262	0.0240583\\
0.6288	0.0239323\\
0.6313	0.0238116\\
0.6339	0.0236866\\
0.6364	0.0235668\\
0.639	0.0234423\\
0.6416	0.0233179\\
0.6441	0.0231983\\
0.6467	0.0230744\\
0.6493	0.022951\\
0.6518	0.0228328\\
0.6544	0.0227102\\
0.657	0.0225876\\
0.6595	0.0224698\\
0.6621	0.0223475\\
0.6646	0.0222304\\
0.6672	0.022109\\
0.6698	0.0219881\\
0.6723	0.0218719\\
0.6749	0.0217511\\
0.6775	0.0216304\\
0.68	0.0215146\\
0.6826	0.0213946\\
0.6852	0.021275\\
0.6877	0.0211602\\
0.6903	0.0210408\\
0.6928	0.020926\\
0.6954	0.0208068\\
0.698	0.0206878\\
0.7005	0.0205737\\
0.7031	0.0204553\\
0.7057	0.0203368\\
0.7082	0.0202228\\
0.7108	0.0201042\\
0.7134	0.0199858\\
0.7159	0.0198721\\
0.7185	0.019754\\
0.721	0.0196404\\
0.7236	0.0195222\\
0.7262	0.0194037\\
0.7287	0.0192897\\
0.7313	0.0191712\\
0.7339	0.0190528\\
0.7364	0.0189389\\
0.739	0.0188202\\
0.7416	0.0187012\\
0.7441	0.0185864\\
0.7467	0.018467\\
0.7492	0.0183521\\
0.7518	0.0182324\\
0.7544	0.0181124\\
0.7569	0.0179967\\
0.7595	0.0178759\\
0.7621	0.0177547\\
0.7646	0.0176379\\
0.7672	0.0175162\\
0.7698	0.0173941\\
0.7723	0.0172762\\
0.7749	0.017153\\
0.7775	0.0170292\\
0.78	0.0169098\\
0.7826	0.0167851\\
0.7851	0.0166648\\
0.7877	0.0165391\\
0.7903	0.0164126\\
0.7928	0.0162903\\
0.7954	0.0161624\\
0.798	0.0160338\\
0.8005	0.0159097\\
0.8031	0.0157798\\
0.8057	0.0156491\\
0.8082	0.0155226\\
0.8108	0.0153901\\
0.8133	0.0152619\\
0.8159	0.0151278\\
0.8185	0.0149928\\
0.821	0.0148622\\
0.8236	0.0147252\\
0.8262	0.0145872\\
0.8287	0.0144535\\
0.8313	0.0143135\\
0.8339	0.0141724\\
0.8364	0.0140358\\
0.839	0.0138926\\
0.8415	0.0137536\\
0.8441	0.0136079\\
0.8467	0.013461\\
0.8492	0.0133187\\
0.8518	0.0131694\\
0.8544	0.0130189\\
0.8569	0.0128728\\
0.8595	0.0127195\\
0.8621	0.0125648\\
0.8646	0.0124148\\
0.8672	0.0122574\\
0.8697	0.0121048\\
0.8723	0.0119445\\
0.8749	0.0117827\\
0.8774	0.0116256\\
0.88	0.0114607\\
0.8826	0.0112943\\
0.8851	0.0111329\\
0.8877	0.0109634\\
0.8903	0.0107921\\
0.8928	0.0106258\\
0.8954	0.0104512\\
0.8979	0.0102818\\
0.9005	0.0101039\\
0.9031	0.00992427\\
0.9056	0.00974984\\
0.9082	0.00956659\\
0.9108	0.00938149\\
0.9133	0.00920178\\
0.9159	0.00901311\\
0.9185	0.0088226\\
0.921	0.00863761\\
0.9236	0.00844329\\
0.9262	0.00824697\\
0.9287	0.00805635\\
0.9313	0.0078562\\
0.9338	0.00766193\\
0.9364	0.00745793\\
0.939	0.00725188\\
0.9415	0.00705178\\
0.9441	0.00684163\\
0.9467	0.00662943\\
0.9492	0.00642347\\
0.9518	0.00620726\\
0.9544	0.00598893\\
0.9569	0.00577696\\
0.9595	0.00555438\\
0.962	0.00533833\\
0.9646	0.00511156\\
0.9672	0.00488268\\
0.9697	0.00466057\\
0.9723	0.00442741\\
0.9749	0.00419202\\
0.9774	0.00396359\\
0.98	0.00372388\\
0.9826	0.003482\\
0.9851	0.00324737\\
0.9877	0.00300117\\
0.9902	0.00276231\\
0.9928	0.00251168\\
0.9954	0.00225882\\
0.9979	0.00201361\\
1.0005	0.00175643\\
1.0031	0.00149704\\
1.0056	0.0012455\\
1.0082	0.000981668\\
1.0108	0.000715583\\
1.0133	0.00045764\\
1.0159	0.000187227\\
1.0184	-7.48597e-05\\
1.021	-0.000349613\\
1.0236	-0.000626615\\
1.0261	-0.000895081\\
1.0287	-0.00117646\\
1.0313	-0.00146001\\
1.0338	-0.0017347\\
1.0364	-0.00202252\\
1.039	-0.00231253\\
1.0415	-0.00259347\\
1.0441	-0.00288779\\
1.0466	-0.00317281\\
1.0492	-0.00347131\\
1.0518	-0.00377192\\
1.0543	-0.00406298\\
1.0569	-0.00436777\\
1.0595	-0.00467469\\
1.062	-0.00497178\\
1.0646	-0.00528276\\
1.0672	-0.00559578\\
1.0697	-0.00589867\\
1.0723	-0.00621569\\
1.0748	-0.00652243\\
1.0774	-0.00684342\\
1.08	-0.00716639\\
1.0825	-0.00747876\\
1.0851	-0.00780551\\
1.0877	-0.0081342\\
1.0902	-0.00845206\\
1.0928	-0.00878451\\
1.0954	-0.00911883\\
1.0979	-0.00944202\\
1.1005	-0.00977991\\
1.1031	-0.0101196\\
1.1056	-0.0104479\\
1.1082	-0.0107911\\
1.1107	-0.0111228\\
1.1133	-0.0114694\\
1.1159	-0.0118176\\
1.1184	-0.0121541\\
1.121	-0.0125056\\
1.1236	-0.0128587\\
1.1261	-0.0131998\\
1.1287	-0.0135561\\
1.1313	-0.0139139\\
1.1338	-0.0142593\\
1.1364	-0.01462\\
1.1389	-0.0149682\\
1.1415	-0.0153318\\
1.1441	-0.0156968\\
1.1466	-0.0160491\\
1.1492	-0.0164167\\
1.1518	-0.0167857\\
1.1543	-0.0171417\\
1.1569	-0.0175132\\
1.1595	-0.0178859\\
1.162	-0.0182454\\
1.1646	-0.0186204\\
1.1671	-0.0189821\\
1.1697	-0.0193593\\
1.1723	-0.0197376\\
1.1748	-0.0201024\\
1.1774	-0.0204827\\
1.18	-0.0208641\\
1.1825	-0.0212316\\
1.1851	-0.0216147\\
1.1877	-0.0219987\\
1.1902	-0.0223688\\
1.1928	-0.0227544\\
1.1953	-0.023126\\
1.1979	-0.0235132\\
1.2005	-0.023901\\
1.203	-0.0242746\\
1.2056	-0.0246638\\
1.2082	-0.0250535\\
1.2107	-0.0254289\\
1.2133	-0.0258198\\
1.2159	-0.0262111\\
1.2184	-0.0265879\\
1.221	-0.0269802\\
1.2235	-0.0273577\\
1.2261	-0.0277508\\
1.2287	-0.0281441\\
1.2312	-0.0285226\\
1.2338	-0.0289165\\
1.2364	-0.0293107\\
1.2389	-0.0296898\\
1.2415	-0.0300843\\
1.2441	-0.0304789\\
1.2466	-0.0308584\\
1.2492	-0.0312531\\
1.2518	-0.0316478\\
1.2543	-0.0320273\\
1.2569	-0.0324219\\
1.2594	-0.0328013\\
1.262	-0.0331957\\
1.2646	-0.03359\\
1.2671	-0.0339689\\
1.2697	-0.0343627\\
1.2723	-0.0347562\\
1.2748	-0.0351343\\
1.2774	-0.0355273\\
1.28	-0.0359198\\
1.2825	-0.0362968\\
1.2851	-0.0366885\\
1.2876	-0.0370647\\
1.2902	-0.0374554\\
1.2928	-0.0378456\\
1.2953	-0.0382202\\
1.2979	-0.0386092\\
1.3005	-0.0389975\\
1.303	-0.0393702\\
1.3056	-0.0397572\\
1.3082	-0.0401434\\
1.3107	-0.040514\\
1.3133	-0.0408986\\
1.3158	-0.0412677\\
1.3184	-0.0416506\\
1.321	-0.0420326\\
1.3235	-0.042399\\
1.3261	-0.0427792\\
1.3287	-0.0431583\\
1.3312	-0.0435218\\
1.3338	-0.0438989\\
1.3364	-0.0442748\\
1.3389	-0.0446352\\
1.3415	-0.0450089\\
1.344	-0.0453671\\
1.3466	-0.0457384\\
1.3492	-0.0461085\\
1.3517	-0.0464631\\
1.3543	-0.0468306\\
1.3569	-0.0471968\\
1.3594	-0.0475476\\
1.362	-0.0479111\\
1.3646	-0.0482732\\
1.3671	-0.04862\\
1.3697	-0.0489791\\
1.3722	-0.0493231\\
1.3748	-0.0496793\\
1.3774	-0.050034\\
1.3799	-0.0503736\\
1.3825	-0.0507252\\
1.3851	-0.0510751\\
1.3876	-0.05141\\
1.3902	-0.0517567\\
1.3928	-0.0521017\\
1.3953	-0.0524318\\
1.3979	-0.0527734\\
1.4005	-0.0531132\\
1.403	-0.0534383\\
1.4056	-0.0537746\\
1.4081	-0.0540962\\
1.4107	-0.0544289\\
1.4133	-0.0547597\\
1.4158	-0.055076\\
1.4184	-0.0554031\\
1.421	-0.0557283\\
1.4235	-0.0560391\\
1.4261	-0.0563605\\
1.4287	-0.0566798\\
1.4312	-0.056985\\
1.4338	-0.0573004\\
1.4363	-0.0576018\\
1.4389	-0.0579132\\
1.4415	-0.0582225\\
1.444	-0.058518\\
1.4466	-0.0588233\\
1.4492	-0.0591264\\
1.4517	-0.0594159\\
1.4543	-0.0597148\\
1.4569	-0.0600116\\
1.4594	-0.060295\\
1.462	-0.0605875\\
1.4645	-0.0608668\\
1.4671	-0.061155\\
1.4697	-0.061441\\
1.4722	-0.0617139\\
1.4748	-0.0619956\\
1.4774	-0.062275\\
1.4799	-0.0625416\\
1.4825	-0.0628166\\
1.4851	-0.0630893\\
1.4876	-0.0633494\\
1.4902	-0.0636176\\
1.4927	-0.0638734\\
1.4953	-0.0641371\\
1.4979	-0.0643985\\
1.5004	-0.0646477\\
1.503	-0.0649046\\
1.5056	-0.0651591\\
1.5081	-0.0654017\\
1.5107	-0.0656517\\
1.5133	-0.0658993\\
1.5158	-0.0661352\\
1.5184	-0.0663782\\
1.5209	-0.0666097\\
1.5235	-0.0668481\\
1.5261	-0.0670841\\
1.5286	-0.0673089\\
1.5312	-0.0675403\\
1.5338	-0.0677694\\
1.5363	-0.0679874\\
1.5389	-0.0682118\\
1.5415	-0.0684339\\
1.544	-0.0686452\\
1.5466	-0.0688626\\
1.5491	-0.0690694\\
1.5517	-0.0692822\\
1.5543	-0.0694927\\
1.5568	-0.0696928\\
1.5594	-0.0698986\\
1.562	-0.070102\\
1.5645	-0.0702954\\
1.5671	-0.0704943\\
1.5697	-0.0706908\\
1.5722	-0.0708775\\
1.5748	-0.0710694\\
1.5774	-0.071259\\
1.5799	-0.0714391\\
1.5825	-0.0716241\\
1.585	-0.0717998\\
1.5876	-0.0719803\\
1.5902	-0.0721585\\
1.5927	-0.0723276\\
1.5953	-0.0725013\\
1.5979	-0.0726727\\
1.6004	-0.0728353\\
1.603	-0.0730022\\
1.6056	-0.0731669\\
1.6081	-0.0733231\\
1.6107	-0.0734833\\
1.6132	-0.0736353\\
1.6158	-0.0737911\\
1.6184	-0.0739448\\
1.6209	-0.0740904\\
1.6235	-0.0742397\\
1.6261	-0.0743868\\
1.6286	-0.0745262\\
1.6312	-0.074669\\
1.6338	-0.0748097\\
1.6363	-0.0749429\\
1.6389	-0.0750794\\
1.6414	-0.0752086\\
1.644	-0.0753408\\
1.6466	-0.075471\\
1.6491	-0.0755942\\
1.6517	-0.0757203\\
1.6543	-0.0758443\\
1.6568	-0.0759615\\
1.6594	-0.0760815\\
1.662	-0.0761994\\
1.6645	-0.0763109\\
1.6671	-0.0764249\\
1.6696	-0.0765326\\
1.6722	-0.0766426\\
1.6748	-0.0767507\\
1.6773	-0.0768528\\
1.6799	-0.076957\\
1.6825	-0.0770594\\
1.685	-0.0771559\\
1.6876	-0.0772545\\
1.6902	-0.0773512\\
1.6927	-0.0774424\\
1.6953	-0.0775354\\
1.6978	-0.0776231\\
1.7004	-0.0777125\\
1.703	-0.0778\\
1.7055	-0.0778825\\
1.7081	-0.0779666\\
1.7107	-0.0780488\\
1.7132	-0.0781263\\
1.7158	-0.0782051\\
1.7184	-0.0782822\\
1.7209	-0.0783547\\
1.7235	-0.0784285\\
1.7261	-0.0785005\\
1.7286	-0.0785683\\
1.7312	-0.078637\\
1.7337	-0.0787017\\
1.7363	-0.0787672\\
1.7389	-0.0788312\\
1.7414	-0.0788912\\
1.744	-0.0789521\\
1.7466	-0.0790114\\
1.7491	-0.079067\\
1.7517	-0.0791233\\
1.7543	-0.079178\\
1.7568	-0.0792293\\
1.7594	-0.0792811\\
1.7619	-0.0793296\\
1.7645	-0.0793785\\
1.7671	-0.0794261\\
1.7696	-0.0794704\\
1.7722	-0.0795152\\
1.7748	-0.0795585\\
1.7773	-0.0795989\\
1.7799	-0.0796396\\
1.7825	-0.0796789\\
1.785	-0.0797155\\
1.7876	-0.0797522\\
1.7901	-0.0797863\\
1.7927	-0.0798204\\
1.7953	-0.0798534\\
1.7978	-0.0798838\\
1.8004	-0.0799143\\
1.803	-0.0799435\\
1.8055	-0.0799705\\
1.8081	-0.0799974\\
1.8107	-0.0800231\\
1.8132	-0.0800467\\
1.8158	-0.0800701\\
1.8183	-0.0800916\\
1.8209	-0.0801128\\
1.8235	-0.0801329\\
1.826	-0.0801512\\
1.8286	-0.0801691\\
1.8312	-0.080186\\
1.8337	-0.0802013\\
1.8363	-0.0802162\\
1.8389	-0.08023\\
1.8414	-0.0802424\\
1.844	-0.0802542\\
1.8465	-0.0802647\\
1.8491	-0.0802747\\
1.8517	-0.0802837\\
1.8542	-0.0802915\\
1.8568	-0.0802987\\
1.8594	-0.0803049\\
1.8619	-0.0803101\\
1.8645	-0.0803146\\
1.8671	-0.0803183\\
1.8696	-0.080321\\
1.8722	-0.0803229\\
1.8747	-0.080324\\
1.8773	-0.0803244\\
1.8799	-0.0803239\\
1.8824	-0.0803227\\
1.885	-0.0803206\\
1.8876	-0.0803178\\
1.8901	-0.0803144\\
1.8927	-0.0803101\\
1.8953	-0.0803051\\
1.8978	-0.0802995\\
1.9004	-0.0802931\\
1.903	-0.0802859\\
1.9055	-0.0802784\\
1.9081	-0.0802699\\
1.9106	-0.080261\\
1.9132	-0.0802512\\
1.9158	-0.0802407\\
1.9183	-0.08023\\
1.9209	-0.0802183\\
1.9235	-0.080206\\
1.926	-0.0801935\\
1.9286	-0.08018\\
1.9312	-0.0801658\\
1.9337	-0.0801517\\
1.9363	-0.0801364\\
1.9388	-0.0801212\\
1.9414	-0.0801048\\
1.944	-0.0800879\\
1.9465	-0.0800711\\
1.9491	-0.0800532\\
1.9517	-0.0800347\\
1.9542	-0.0800164\\
1.9568	-0.079997\\
1.9594	-0.079977\\
1.9619	-0.0799573\\
1.9645	-0.0799364\\
1.967	-0.0799158\\
1.9696	-0.079894\\
1.9722	-0.0798717\\
1.9747	-0.0798498\\
1.9773	-0.0798266\\
1.9799	-0.079803\\
1.9824	-0.0797799\\
1.985	-0.0797555\\
1.9876	-0.0797306\\
1.9901	-0.0797063\\
1.9927	-0.0796807\\
1.9952	-0.0796556\\
1.9978	-0.0796292\\
2.0004	-0.0796024\\
2.0029	-0.0795763\\
2.0055	-0.0795487\\
2.0081	-0.0795208\\
2.0106	-0.0794936\\
2.0132	-0.079465\\
2.0158	-0.079436\\
2.0183	-0.0794078\\
2.0209	-0.0793782\\
2.0234	-0.0793494\\
2.026	-0.0793191\\
2.0286	-0.0792884\\
2.0311	-0.0792587\\
2.0337	-0.0792274\\
2.0363	-0.0791958\\
2.0388	-0.0791651\\
2.0414	-0.079133\\
2.044	-0.0791005\\
2.0465	-0.0790689\\
2.0491	-0.0790359\\
2.0517	-0.0790025\\
2.0542	-0.0789701\\
2.0568	-0.0789362\\
2.0593	-0.0789033\\
2.0619	-0.0788688\\
2.0645	-0.078834\\
2.067	-0.0788003\\
2.0696	-0.078765\\
2.0722	-0.0787294\\
2.0747	-0.078695\\
2.0773	-0.0786589\\
2.0799	-0.0786225\\
2.0824	-0.0785873\\
2.085	-0.0785504\\
2.0875	-0.0785147\\
2.0901	-0.0784773\\
2.0927	-0.0784397\\
2.0952	-0.0784033\\
2.0978	-0.0783652\\
2.1004	-0.0783268\\
2.1029	-0.0782897\\
2.1055	-0.0782508\\
2.1081	-0.0782117\\
2.1106	-0.0781739\\
2.1132	-0.0781343\\
2.1157	-0.078096\\
2.1183	-0.0780559\\
2.1209	-0.0780157\\
2.1234	-0.0779767\\
2.126	-0.0779359\\
2.1286	-0.0778949\\
2.1311	-0.0778553\\
2.1337	-0.0778138\\
2.1363	-0.0777721\\
2.1388	-0.0777318\\
2.1414	-0.0776896\\
2.1439	-0.0776489\\
2.1465	-0.0776063\\
2.1491	-0.0775634\\
2.1516	-0.077522\\
2.1542	-0.0774786\\
2.1568	-0.0774351\\
2.1593	-0.077393\\
2.1619	-0.077349\\
2.1645	-0.0773047\\
2.167	-0.077262\\
2.1696	-0.0772172\\
2.1721	-0.077174\\
2.1747	-0.0771289\\
2.1773	-0.0770835\\
2.1798	-0.0770396\\
2.1824	-0.0769937\\
2.185	-0.0769476\\
2.1875	-0.076903\\
2.1901	-0.0768564\\
2.1927	-0.0768096\\
2.1952	-0.0767644\\
2.1978	-0.0767171\\
2.2004	-0.0766695\\
2.2029	-0.0766236\\
2.2055	-0.0765756\\
2.208	-0.0765292\\
2.2106	-0.0764807\\
2.2132	-0.076432\\
2.2157	-0.0763849\\
2.2183	-0.0763357\\
2.2209	-0.0762863\\
2.2234	-0.0762385\\
2.226	-0.0761885\\
2.2286	-0.0761383\\
2.2311	-0.0760899\\
2.2337	-0.0760392\\
2.2362	-0.0759902\\
2.2388	-0.075939\\
2.2414	-0.0758876\\
2.2439	-0.0758379\\
2.2465	-0.075786\\
2.2491	-0.0757338\\
2.2516	-0.0756834\\
2.2542	-0.0756307\\
2.2568	-0.0755778\\
2.2593	-0.0755267\\
2.2619	-0.0754732\\
2.2644	-0.0754216\\
2.267	-0.0753676\\
2.2696	-0.0753134\\
2.2721	-0.075261\\
2.2747	-0.0752063\\
2.2773	-0.0751513\\
2.2798	-0.0750982\\
2.2824	-0.0750427\\
2.285	-0.0749869\\
2.2875	-0.074933\\
2.2901	-0.0748768\\
2.2926	-0.0748224\\
2.2952	-0.0747656\\
2.2978	-0.0747085\\
2.3003	-0.0746534\\
2.3029	-0.0745958\\
2.3055	-0.0745379\\
2.308	-0.074482\\
2.3106	-0.0744236\\
2.3132	-0.0743649\\
2.3157	-0.0743083\\
2.3183	-0.0742491\\
2.3208	-0.0741919\\
2.3234	-0.0741322\\
2.326	-0.0740722\\
2.3285	-0.0740142\\
2.3311	-0.0739537\\
2.3337	-0.0738929\\
2.3362	-0.0738342\\
2.3388	-0.0737729\\
2.3414	-0.0737113\\
2.3439	-0.0736518\\
2.3465	-0.0735897\\
2.349	-0.0735297\\
2.3516	-0.0734671\\
2.3542	-0.0734041\\
2.3567	-0.0733434\\
2.3593	-0.0732799\\
2.3619	-0.0732162\\
2.3644	-0.0731547\\
2.367	-0.0730905\\
2.3696	-0.073026\\
2.3721	-0.0729637\\
2.3747	-0.0728987\\
2.3773	-0.0728334\\
2.3798	-0.0727704\\
2.3824	-0.0727046\\
2.3849	-0.0726411\\
2.3875	-0.0725748\\
2.3901	-0.0725082\\
2.3926	-0.072444\\
2.3952	-0.0723769\\
2.3978	-0.0723095\\
2.4003	-0.0722446\\
2.4029	-0.0721767\\
2.4055	-0.0721086\\
2.408	-0.0720429\\
2.4106	-0.0719744\\
2.4131	-0.0719082\\
2.4157	-0.0718391\\
2.4183	-0.0717698\\
2.4208	-0.0717029\\
2.4234	-0.0716331\\
2.426	-0.0715631\\
2.4285	-0.0714955\\
2.4311	-0.071425\\
2.4337	-0.0713543\\
2.4362	-0.071286\\
2.4388	-0.0712148\\
2.4413	-0.0711461\\
2.4439	-0.0710745\\
2.4465	-0.0710026\\
2.449	-0.0709333\\
2.4516	-0.0708609\\
2.4542	-0.0707884\\
2.4567	-0.0707184\\
2.4593	-0.0706455\\
2.4619	-0.0705723\\
2.4644	-0.0705017\\
2.467	-0.0704281\\
2.4695	-0.0703571\\
2.4721	-0.0702831\\
2.4747	-0.0702089\\
2.4772	-0.0701373\\
2.4798	-0.0700627\\
2.4824	-0.0699879\\
2.4849	-0.0699158\\
2.4875	-0.0698406\\
2.4901	-0.0697652\\
2.4926	-0.0696926\\
2.4952	-0.0696168\\
2.4977	-0.0695438\\
2.5003	-0.0694678\\
2.5029	-0.0693915\\
2.5054	-0.069318\\
2.508	-0.0692414\\
2.5106	-0.0691646\\
2.5131	-0.0690907\\
2.5157	-0.0690136\\
2.5183	-0.0689364\\
2.5208	-0.068862\\
2.5234	-0.0687844\\
2.526	-0.0687067\\
2.5285	-0.0686319\\
2.5311	-0.0685539\\
2.5336	-0.0684788\\
2.5362	-0.0684006\\
2.5388	-0.0683223\\
2.5413	-0.0682468\\
2.5439	-0.0681682\\
2.5465	-0.0680894\\
2.549	-0.0680136\\
2.5516	-0.0679347\\
2.5542	-0.0678556\\
2.5567	-0.0677795\\
2.5593	-0.0677002\\
2.5618	-0.0676239\\
2.5644	-0.0675444\\
2.567	-0.0674648\\
2.5695	-0.0673882\\
2.5721	-0.0673085\\
2.5747	-0.0672287\\
2.5772	-0.0671519\\
2.5798	-0.0670719\\
2.5824	-0.0669918\\
2.5849	-0.0669148\\
2.5875	-0.0668346\\
2.59	-0.0667575\\
2.5926	-0.0666772\\
2.5952	-0.0665969\\
2.5977	-0.0665196\\
2.6003	-0.0664392\\
2.6029	-0.0663587\\
2.6054	-0.0662813\\
2.608	-0.0662008\\
2.6106	-0.0661202\\
2.6131	-0.0660427\\
2.6157	-0.0659621\\
2.6182	-0.0658846\\
2.6208	-0.0658039\\
2.6234	-0.0657233\\
2.6259	-0.0656457\\
2.6285	-0.0655651\\
2.6311	-0.0654844\\
2.6336	-0.0654069\\
2.6362	-0.0653262\\
2.6388	-0.0652456\\
2.6413	-0.0651681\\
2.6439	-0.0650875\\
2.6464	-0.0650101\\
2.649	-0.0649295\\
2.6516	-0.064849\\
2.6541	-0.0647717\\
2.6567	-0.0646913\\
2.6593	-0.0646109\\
2.6618	-0.0645337\\
2.6644	-0.0644534\\
2.667	-0.0643732\\
2.6695	-0.0642962\\
2.6721	-0.0642161\\
2.6746	-0.0641392\\
2.6772	-0.0640592\\
2.6798	-0.0639794\\
2.6823	-0.0639027\\
2.6849	-0.063823\\
2.6875	-0.0637434\\
2.69	-0.0636669\\
2.6926	-0.0635875\\
2.6952	-0.0635082\\
2.6977	-0.0634321\\
2.7003	-0.063353\\
2.7029	-0.063274\\
2.7054	-0.0631982\\
2.708	-0.0631195\\
2.7105	-0.0630439\\
2.7131	-0.0629654\\
2.7157	-0.062887\\
2.7182	-0.0628118\\
2.7208	-0.0627338\\
2.7234	-0.0626558\\
2.7259	-0.0625811\\
2.7285	-0.0625034\\
2.7311	-0.062426\\
2.7336	-0.0623516\\
2.7362	-0.0622745\\
2.7387	-0.0622004\\
2.7413	-0.0621236\\
2.7439	-0.062047\\
2.7464	-0.0619735\\
2.749	-0.0618972\\
2.7516	-0.0618211\\
2.7541	-0.0617481\\
2.7567	-0.0616724\\
2.7593	-0.0615969\\
2.7618	-0.0615244\\
2.7644	-0.0614493\\
2.7669	-0.0613773\\
2.7695	-0.0613026\\
2.7721	-0.0612281\\
2.7746	-0.0611567\\
2.7772	-0.0610827\\
2.7798	-0.0610088\\
2.7823	-0.0609381\\
2.7849	-0.0608647\\
2.7875	-0.0607916\\
2.79	-0.0607215\\
2.7926	-0.0606488\\
2.7951	-0.0605792\\
2.7977	-0.060507\\
2.8003	-0.0604351\\
2.8028	-0.0603662\\
2.8054	-0.0602948\\
2.808	-0.0602237\\
2.8105	-0.0601556\\
2.8131	-0.0600849\\
2.8157	-0.0600146\\
2.8182	-0.0599472\\
2.8208	-0.0598774\\
2.8233	-0.0598106\\
2.8259	-0.0597414\\
2.8285	-0.0596724\\
2.831	-0.0596064\\
2.8336	-0.059538\\
2.8362	-0.05947\\
2.8387	-0.0594048\\
2.8413	-0.0593373\\
2.8439	-0.0592701\\
2.8464	-0.0592058\\
2.849	-0.0591393\\
2.8516	-0.059073\\
2.8541	-0.0590096\\
2.8567	-0.058944\\
2.8592	-0.0588812\\
2.8618	-0.0588162\\
2.8644	-0.0587515\\
2.8669	-0.0586897\\
2.8695	-0.0586256\\
2.8721	-0.058562\\
2.8746	-0.0585011\\
2.8772	-0.058438\\
2.8798	-0.0583754\\
2.8823	-0.0583154\\
2.8849	-0.0582534\\
2.8874	-0.0581942\\
2.89	-0.0581329\\
2.8926	-0.0580719\\
2.8951	-0.0580136\\
2.8977	-0.0579534\\
2.9003	-0.0578935\\
2.9028	-0.0578362\\
2.9054	-0.0577771\\
2.908	-0.0577183\\
2.9105	-0.0576621\\
2.9131	-0.057604\\
2.9156	-0.0575485\\
2.9182	-0.0574911\\
2.9208	-0.0574341\\
2.9233	-0.0573797\\
2.9259	-0.0573235\\
2.9285	-0.0572676\\
2.931	-0.0572143\\
2.9336	-0.0571592\\
2.9362	-0.0571045\\
2.9387	-0.0570523\\
2.9413	-0.0569983\\
2.9438	-0.0569468\\
2.9464	-0.0568937\\
2.949	-0.0568409\\
2.9515	-0.0567905\\
2.9541	-0.0567385\\
2.9567	-0.0566869\\
2.9592	-0.0566377\\
2.9618	-0.0565869\\
2.9644	-0.0565365\\
2.9669	-0.0564884\\
2.9695	-0.0564388\\
2.972	-0.0563915\\
2.9746	-0.0563427\\
2.9772	-0.0562943\\
2.9797	-0.0562481\\
2.9823	-0.0562005\\
2.9849	-0.0561533\\
2.9874	-0.0561084\\
2.99	-0.056062\\
2.9926	-0.0560161\\
2.9951	-0.0559723\\
2.9977	-0.0559271\\
3.0003	-0.0558824\\
3.0028	-0.0558398\\
3.0054	-0.055796\\
3.0079	-0.0557542\\
3.0105	-0.0557111\\
3.0131	-0.0556685\\
3.0156	-0.0556279\\
3.0182	-0.0555861\\
3.0208	-0.0555447\\
3.0233	-0.0555053\\
3.0259	-0.0554648\\
3.0285	-0.0554247\\
3.031	-0.0553865\\
3.0336	-0.0553472\\
3.0361	-0.0553099\\
3.0387	-0.0552714\\
3.0413	-0.0552334\\
3.0438	-0.0551973\\
3.0464	-0.0551601\\
3.049	-0.0551234\\
3.0515	-0.0550885\\
3.0541	-0.0550526\\
3.0567	-0.0550172\\
3.0592	-0.0549835\\
3.0618	-0.0549489\\
3.0643	-0.054916\\
3.0669	-0.0548822\\
3.0695	-0.0548489\\
3.072	-0.0548173\\
3.0746	-0.0547848\\
3.0772	-0.0547527\\
3.0797	-0.0547223\\
3.0823	-0.0546911\\
3.0849	-0.0546603\\
3.0874	-0.0546311\\
3.09	-0.0546012\\
3.0925	-0.0545728\\
3.0951	-0.0545438\\
3.0977	-0.0545151\\
3.1002	-0.054488\\
3.1028	-0.0544602\\
3.1054	-0.0544328\\
3.1079	-0.0544069\\
3.1105	-0.0543803\\
3.1131	-0.0543542\\
3.1156	-0.0543295\\
3.1182	-0.0543043\\
3.1207	-0.0542804\\
3.1233	-0.0542559\\
3.1259	-0.0542319\\
3.1284	-0.0542093\\
3.131	-0.0541861\\
3.1336	-0.0541634\\
3.1361	-0.0541419\\
3.1387	-0.05412\\
3.1413	-0.0540985\\
3.1438	-0.0540782\\
3.1464	-0.0540576\\
3.1489	-0.0540381\\
3.1515	-0.0540183\\
3.1541	-0.0539989\\
3.1566	-0.0539806\\
3.1592	-0.053962\\
3.1618	-0.0539438\\
3.1643	-0.0539267\\
3.1669	-0.0539094\\
3.1695	-0.0538924\\
3.172	-0.0538765\\
3.1746	-0.0538604\\
3.1772	-0.0538447\\
3.1797	-0.0538299\\
3.1823	-0.053815\\
3.1848	-0.0538011\\
3.1874	-0.0537869\\
3.19	-0.0537732\\
3.1925	-0.0537604\\
3.1951	-0.0537475\\
3.1977	-0.053735\\
3.2002	-0.0537233\\
3.2028	-0.0537116\\
3.2054	-0.0537003\\
3.2079	-0.0536898\\
3.2105	-0.0536792\\
3.213	-0.0536695\\
3.2156	-0.0536597\\
3.2182	-0.0536503\\
3.2207	-0.0536417\\
3.2233	-0.053633\\
3.2259	-0.0536248\\
3.2284	-0.0536173\\
3.231	-0.0536098\\
3.2336	-0.0536028\\
3.2361	-0.0535963\\
3.2387	-0.05359\\
3.2412	-0.0535843\\
3.2438	-0.0535787\\
3.2464	-0.0535735\\
3.2489	-0.0535688\\
3.2515	-0.0535644\\
3.2541	-0.0535603\\
3.2566	-0.0535567\\
3.2592	-0.0535534\\
3.2618	-0.0535504\\
3.2643	-0.0535478\\
3.2669	-0.0535456\\
3.2694	-0.0535437\\
3.272	-0.0535422\\
3.2746	-0.053541\\
3.2771	-0.0535401\\
3.2797	-0.0535396\\
3.2823	-0.0535395\\
3.2848	-0.0535397\\
3.2874	-0.0535403\\
3.29	-0.0535412\\
3.2925	-0.0535424\\
3.2951	-0.0535439\\
3.2976	-0.0535458\\
3.3002	-0.053548\\
3.3028	-0.0535506\\
3.3053	-0.0535535\\
3.3079	-0.0535567\\
3.3105	-0.0535603\\
3.313	-0.0535641\\
3.3156	-0.0535684\\
3.3182	-0.0535729\\
3.3207	-0.0535777\\
3.3233	-0.0535829\\
3.3259	-0.0535884\\
3.3284	-0.0535941\\
3.331	-0.0536002\\
3.3335	-0.0536065\\
3.3361	-0.0536133\\
3.3387	-0.0536204\\
3.3412	-0.0536276\\
3.3438	-0.0536353\\
3.3464	-0.0536433\\
3.3489	-0.0536513\\
3.3515	-0.05366\\
3.3541	-0.0536689\\
3.3566	-0.0536778\\
3.3592	-0.0536873\\
3.3617	-0.0536967\\
3.3643	-0.0537068\\
3.3669	-0.0537172\\
3.3694	-0.0537275\\
3.372	-0.0537385\\
3.3746	-0.0537497\\
3.3771	-0.0537608\\
3.3797	-0.0537726\\
3.3823	-0.0537846\\
3.3848	-0.0537965\\
3.3874	-0.0538091\\
3.3899	-0.0538215\\
3.3925	-0.0538347\\
3.3951	-0.0538481\\
3.3976	-0.0538612\\
3.4002	-0.0538751\\
3.4028	-0.0538893\\
3.4053	-0.0539032\\
3.4079	-0.0539179\\
3.4105	-0.0539329\\
3.413	-0.0539475\\
3.4156	-0.0539629\\
3.4181	-0.053978\\
3.4207	-0.0539939\\
3.4233	-0.0540101\\
3.4258	-0.0540258\\
3.4284	-0.0540425\\
3.431	-0.0540593\\
3.4335	-0.0540758\\
3.4361	-0.0540931\\
3.4387	-0.0541106\\
3.4412	-0.0541277\\
3.4438	-0.0541457\\
3.4463	-0.0541632\\
3.4489	-0.0541816\\
3.4515	-0.0542002\\
3.454	-0.0542183\\
3.4566	-0.0542374\\
3.4592	-0.0542566\\
3.4617	-0.0542753\\
3.4643	-0.054295\\
3.4669	-0.0543149\\
3.4694	-0.0543341\\
3.472	-0.0543544\\
3.4745	-0.054374\\
3.4771	-0.0543946\\
3.4797	-0.0544154\\
3.4822	-0.0544356\\
3.4848	-0.0544568\\
3.4874	-0.0544781\\
3.4899	-0.0544988\\
3.4925	-0.0545205\\
3.4951	-0.0545424\\
3.4976	-0.0545636\\
3.5002	-0.0545858\\
3.5028	-0.0546081\\
3.5053	-0.0546298\\
3.5079	-0.0546525\\
3.5104	-0.0546745\\
3.513	-0.0546974\\
3.5156	-0.0547206\\
3.5181	-0.054743\\
3.5207	-0.0547664\\
3.5233	-0.05479\\
3.5258	-0.0548128\\
3.5284	-0.0548367\\
3.531	-0.0548606\\
3.5335	-0.0548838\\
3.5361	-0.0549081\\
3.5386	-0.0549315\\
3.5412	-0.0549561\\
3.5438	-0.0549807\\
3.5463	-0.0550045\\
3.5489	-0.0550293\\
3.5515	-0.0550543\\
3.554	-0.0550785\\
3.5566	-0.0551037\\
3.5592	-0.055129\\
3.5617	-0.0551534\\
3.5643	-0.0551789\\
3.5668	-0.0552036\\
3.5694	-0.0552293\\
3.572	-0.0552551\\
3.5745	-0.05528\\
3.5771	-0.055306\\
3.5797	-0.055332\\
3.5822	-0.0553572\\
3.5848	-0.0553834\\
3.5874	-0.0554097\\
3.5899	-0.0554351\\
3.5925	-0.0554616\\
3.595	-0.0554871\\
3.5976	-0.0555137\\
3.6002	-0.0555404\\
3.6027	-0.0555661\\
3.6053	-0.0555929\\
3.6079	-0.0556198\\
3.6104	-0.0556456\\
3.613	-0.0556726\\
3.6156	-0.0556996\\
3.6181	-0.0557256\\
3.6207	-0.0557527\\
3.6232	-0.0557788\\
3.6258	-0.055806\\
3.6284	-0.0558333\\
3.6309	-0.0558595\\
3.6335	-0.0558868\\
3.6361	-0.0559141\\
3.6386	-0.0559404\\
3.6412	-0.0559678\\
3.6438	-0.0559952\\
3.6463	-0.0560215\\
3.6489	-0.056049\\
3.6515	-0.0560764\\
3.654	-0.0561028\\
3.6566	-0.0561303\\
3.6591	-0.0561567\\
3.6617	-0.0561842\\
3.6643	-0.0562117\\
3.6668	-0.0562381\\
3.6694	-0.0562656\\
3.672	-0.0562931\\
3.6745	-0.0563195\\
3.6771	-0.056347\\
3.6797	-0.0563744\\
3.6822	-0.0564008\\
3.6848	-0.0564282\\
3.6873	-0.0564546\\
3.6899	-0.0564819\\
3.6925	-0.0565093\\
3.695	-0.0565356\\
3.6976	-0.0565629\\
3.7002	-0.0565901\\
3.7027	-0.0566163\\
3.7053	-0.0566435\\
3.7079	-0.0566706\\
3.7104	-0.0566967\\
3.713	-0.0567238\\
3.7155	-0.0567498\\
3.7181	-0.0567767\\
3.7207	-0.0568037\\
3.7232	-0.0568295\\
3.7258	-0.0568563\\
3.7284	-0.0568831\\
3.7309	-0.0569088\\
3.7335	-0.0569354\\
3.7361	-0.056962\\
3.7386	-0.0569874\\
3.7412	-0.0570139\\
3.7437	-0.0570392\\
3.7463	-0.0570655\\
3.7489	-0.0570917\\
3.7514	-0.0571169\\
3.754	-0.0571429\\
3.7566	-0.0571689\\
3.7591	-0.0571938\\
3.7617	-0.0572196\\
3.7643	-0.0572453\\
3.7668	-0.05727\\
3.7694	-0.0572955\\
3.7719	-0.05732\\
3.7745	-0.0573453\\
3.7771	-0.0573706\\
3.7796	-0.0573947\\
3.7822	-0.0574198\\
3.7848	-0.0574447\\
3.7873	-0.0574686\\
3.7899	-0.0574933\\
3.7925	-0.0575179\\
3.795	-0.0575415\\
3.7976	-0.0575658\\
3.8002	-0.0575901\\
3.8027	-0.0576133\\
3.8053	-0.0576373\\
3.8078	-0.0576603\\
3.8104	-0.0576841\\
3.813	-0.0577078\\
3.8155	-0.0577304\\
3.8181	-0.0577538\\
3.8207	-0.057777\\
3.8232	-0.0577993\\
3.8258	-0.0578222\\
3.8284	-0.0578451\\
3.8309	-0.0578669\\
3.8335	-0.0578895\\
3.836	-0.0579111\\
3.8386	-0.0579334\\
3.8412	-0.0579555\\
3.8437	-0.0579767\\
3.8463	-0.0579985\\
3.8489	-0.0580202\\
3.8514	-0.0580409\\
3.854	-0.0580623\\
3.8566	-0.0580835\\
3.8591	-0.0581038\\
3.8617	-0.0581247\\
3.8642	-0.0581447\\
3.8668	-0.0581652\\
3.8694	-0.0581857\\
3.8719	-0.0582052\\
3.8745	-0.0582253\\
3.8771	-0.0582452\\
3.8796	-0.0582642\\
3.8822	-0.0582838\\
3.8848	-0.0583032\\
3.8873	-0.0583217\\
3.8899	-0.0583408\\
3.8924	-0.058359\\
3.895	-0.0583777\\
3.8976	-0.0583962\\
3.9001	-0.0584139\\
3.9027	-0.0584321\\
3.9053	-0.0584501\\
3.9078	-0.0584672\\
3.9104	-0.0584848\\
3.913	-0.0585023\\
3.9155	-0.0585189\\
3.9181	-0.0585359\\
3.9206	-0.0585522\\
3.9232	-0.0585689\\
3.9258	-0.0585854\\
3.9283	-0.058601\\
3.9309	-0.0586172\\
3.9335	-0.0586331\\
3.936	-0.0586482\\
3.9386	-0.0586637\\
3.9412	-0.0586791\\
3.9437	-0.0586936\\
3.9463	-0.0587086\\
3.9488	-0.0587227\\
3.9514	-0.0587373\\
3.954	-0.0587516\\
3.9565	-0.0587652\\
3.9591	-0.0587791\\
3.9617	-0.0587929\\
3.9642	-0.0588059\\
3.9668	-0.0588192\\
3.9694	-0.0588323\\
3.9719	-0.0588447\\
3.9745	-0.0588574\\
3.9771	-0.0588699\\
3.9796	-0.0588817\\
3.9822	-0.0588938\\
3.9847	-0.0589052\\
3.9873	-0.0589169\\
3.9899	-0.0589283\\
3.9924	-0.0589391\\
3.995	-0.0589501\\
3.9976	-0.058961\\
4.0001	-0.0589711\\
4.0027	-0.0589815\\
4.0053	-0.0589917\\
4.0078	-0.0590013\\
4.0104	-0.059011\\
4.0129	-0.0590202\\
4.0155	-0.0590295\\
4.0181	-0.0590386\\
4.0206	-0.0590472\\
4.0232	-0.0590559\\
4.0258	-0.0590643\\
4.0283	-0.0590722\\
4.0309	-0.0590803\\
4.0335	-0.0590881\\
4.036	-0.0590954\\
4.0386	-0.0591028\\
4.0411	-0.0591097\\
4.0437	-0.0591166\\
4.0463	-0.0591233\\
4.0488	-0.0591296\\
4.0514	-0.0591359\\
4.054	-0.059142\\
4.0565	-0.0591476\\
4.0591	-0.0591533\\
4.0617	-0.0591587\\
4.0642	-0.0591637\\
4.0668	-0.0591687\\
4.0693	-0.0591734\\
4.0719	-0.0591779\\
4.0745	-0.0591823\\
4.077	-0.0591863\\
4.0796	-0.0591902\\
4.0822	-0.0591939\\
4.0847	-0.0591973\\
4.0873	-0.0592006\\
4.0899	-0.0592036\\
4.0924	-0.0592064\\
4.095	-0.059209\\
4.0975	-0.0592114\\
4.1001	-0.0592136\\
4.1027	-0.0592156\\
4.1052	-0.0592174\\
4.1078	-0.0592189\\
4.1104	-0.0592203\\
4.1129	-0.0592215\\
4.1155	-0.0592224\\
4.1181	-0.0592232\\
4.1206	-0.0592237\\
4.1232	-0.059224\\
4.1258	-0.0592241\\
4.1283	-0.0592241\\
4.1309	-0.0592238\\
4.1334	-0.0592233\\
4.136	-0.0592226\\
4.1386	-0.0592217\\
4.1411	-0.0592207\\
4.1437	-0.0592194\\
4.1463	-0.0592179\\
4.1488	-0.0592162\\
4.1514	-0.0592143\\
4.154	-0.0592122\\
4.1565	-0.05921\\
4.1591	-0.0592075\\
4.1616	-0.0592049\\
4.1642	-0.059202\\
4.1668	-0.059199\\
4.1693	-0.0591958\\
4.1719	-0.0591923\\
4.1745	-0.0591887\\
4.177	-0.059185\\
4.1796	-0.0591809\\
4.1822	-0.0591767\\
4.1847	-0.0591724\\
4.1873	-0.0591678\\
4.1898	-0.0591632\\
4.1924	-0.0591582\\
4.195	-0.059153\\
4.1975	-0.0591479\\
4.2001	-0.0591424\\
4.2027	-0.0591367\\
4.2052	-0.059131\\
4.2078	-0.0591249\\
4.2104	-0.0591187\\
4.2129	-0.0591125\\
4.2155	-0.0591059\\
4.218	-0.0590994\\
4.2206	-0.0590924\\
4.2232	-0.0590853\\
4.2257	-0.0590783\\
4.2283	-0.0590708\\
4.2309	-0.0590632\\
4.2334	-0.0590557\\
4.236	-0.0590478\\
4.2386	-0.0590396\\
4.2411	-0.0590317\\
4.2437	-0.0590232\\
4.2462	-0.0590149\\
4.2488	-0.0590062\\
4.2514	-0.0589973\\
4.2539	-0.0589885\\
4.2565	-0.0589793\\
4.2591	-0.0589699\\
4.2616	-0.0589607\\
4.2642	-0.058951\\
4.2668	-0.0589412\\
4.2693	-0.0589316\\
4.2719	-0.0589215\\
4.2744	-0.0589116\\
4.277	-0.0589011\\
4.2796	-0.0588906\\
4.2821	-0.0588803\\
4.2847	-0.0588694\\
4.2873	-0.0588585\\
4.2898	-0.0588478\\
4.2924	-0.0588365\\
4.295	-0.0588251\\
4.2975	-0.058814\\
4.3001	-0.0588024\\
4.3027	-0.0587906\\
4.3052	-0.0587792\\
4.3078	-0.0587671\\
4.3103	-0.0587554\\
4.3129	-0.0587432\\
4.3155	-0.0587308\\
4.318	-0.0587187\\
4.3206	-0.0587061\\
4.3232	-0.0586933\\
4.3257	-0.058681\\
4.3283	-0.058668\\
4.3309	-0.0586549\\
4.3334	-0.0586422\\
4.336	-0.0586289\\
4.3385	-0.058616\\
4.3411	-0.0586025\\
4.3437	-0.0585888\\
4.3462	-0.0585756\\
4.3488	-0.0585618\\
4.3514	-0.0585479\\
4.3539	-0.0585344\\
4.3565	-0.0585203\\
4.3591	-0.058506\\
4.3616	-0.0584923\\
4.3642	-0.0584779\\
4.3667	-0.0584639\\
4.3693	-0.0584494\\
4.3719	-0.0584347\\
4.3744	-0.0584205\\
4.377	-0.0584057\\
4.3796	-0.0583908\\
4.3821	-0.0583764\\
4.3847	-0.0583613\\
4.3873	-0.0583462\\
4.3898	-0.0583315\\
4.3924	-0.0583162\\
4.3949	-0.0583015\\
4.3975	-0.058286\\
4.4001	-0.0582705\\
4.4026	-0.0582556\\
4.4052	-0.05824\\
4.4078	-0.0582243\\
4.4103	-0.0582091\\
4.4129	-0.0581933\\
4.4155	-0.0581774\\
4.418	-0.0581621\\
4.4206	-0.0581461\\
4.4231	-0.0581307\\
4.4257	-0.0581146\\
4.4283	-0.0580985\\
4.4308	-0.0580829\\
4.4334	-0.0580667\\
4.436	-0.0580504\\
4.4385	-0.0580347\\
4.4411	-0.0580183\\
4.4437	-0.0580019\\
4.4462	-0.0579861\\
4.4488	-0.0579697\\
4.4514	-0.0579531\\
4.4539	-0.0579372\\
4.4565	-0.0579206\\
4.459	-0.0579047\\
4.4616	-0.057888\\
4.4642	-0.0578714\\
4.4667	-0.0578553\\
4.4693	-0.0578386\\
4.4719	-0.0578218\\
4.4744	-0.0578057\\
4.477	-0.0577889\\
4.4796	-0.0577721\\
4.4821	-0.057756\\
4.4847	-0.0577391\\
4.4872	-0.0577229\\
4.4898	-0.0577061\\
4.4924	-0.0576892\\
4.4949	-0.057673\\
4.4975	-0.0576561\\
4.5001	-0.0576392\\
4.5026	-0.0576229\\
4.5052	-0.057606\\
4.5078	-0.0575891\\
4.5103	-0.0575729\\
4.5129	-0.057556\\
4.5154	-0.0575397\\
4.518	-0.0575228\\
4.5206	-0.0575059\\
4.5231	-0.0574897\\
4.5257	-0.0574728\\
4.5283	-0.057456\\
4.5308	-0.0574397\\
4.5334	-0.0574229\\
4.536	-0.0574061\\
4.5385	-0.0573899\\
4.5411	-0.0573731\\
4.5436	-0.057357\\
4.5462	-0.0573402\\
4.5488	-0.0573235\\
4.5513	-0.0573074\\
4.5539	-0.0572907\\
4.5565	-0.057274\\
4.559	-0.057258\\
4.5616	-0.0572414\\
4.5642	-0.0572248\\
4.5667	-0.0572089\\
4.5693	-0.0571924\\
4.5718	-0.0571765\\
4.5744	-0.05716\\
4.577	-0.0571436\\
4.5795	-0.0571278\\
4.5821	-0.0571115\\
4.5847	-0.0570952\\
4.5872	-0.0570795\\
4.5898	-0.0570633\\
4.5924	-0.0570471\\
4.5949	-0.0570315\\
4.5975	-0.0570154\\
4.6001	-0.0569994\\
4.6026	-0.056984\\
4.6052	-0.056968\\
4.6077	-0.0569527\\
4.6103	-0.0569368\\
4.6129	-0.056921\\
4.6154	-0.0569058\\
4.618	-0.0568901\\
4.6206	-0.0568744\\
4.6231	-0.0568594\\
4.6257	-0.0568439\\
4.6283	-0.0568284\\
4.6308	-0.0568135\\
4.6334	-0.0567982\\
4.6359	-0.0567834\\
4.6385	-0.0567682\\
4.6411	-0.056753\\
4.6436	-0.0567384\\
4.6462	-0.0567234\\
4.6488	-0.0567083\\
4.6513	-0.056694\\
4.6539	-0.0566791\\
4.6565	-0.0566643\\
4.659	-0.0566501\\
4.6616	-0.0566355\\
4.6641	-0.0566214\\
4.6667	-0.0566069\\
4.6693	-0.0565924\\
4.6718	-0.0565786\\
4.6744	-0.0565643\\
4.677	-0.0565501\\
4.6795	-0.0565365\\
4.6821	-0.0565224\\
4.6847	-0.0565084\\
4.6872	-0.056495\\
4.6898	-0.0564811\\
4.6923	-0.0564679\\
4.6949	-0.0564542\\
4.6975	-0.0564406\\
4.7	-0.0564275\\
4.7026	-0.0564141\\
4.7052	-0.0564007\\
4.7077	-0.0563879\\
4.7103	-0.0563747\\
4.7129	-0.0563616\\
4.7154	-0.0563491\\
4.718	-0.0563361\\
4.7205	-0.0563238\\
4.7231	-0.056311\\
4.7257	-0.0562983\\
4.7282	-0.0562862\\
4.7308	-0.0562736\\
4.7334	-0.0562612\\
4.7359	-0.0562494\\
4.7385	-0.0562371\\
4.7411	-0.056225\\
4.7436	-0.0562133\\
4.7462	-0.0562014\\
4.7487	-0.0561899\\
4.7513	-0.0561781\\
4.7539	-0.0561664\\
4.7564	-0.0561553\\
4.759	-0.0561438\\
4.7616	-0.0561323\\
4.7641	-0.0561215\\
4.7667	-0.0561102\\
4.7693	-0.0560991\\
4.7718	-0.0560885\\
4.7744	-0.0560775\\
4.777	-0.0560667\\
4.7795	-0.0560563\\
4.7821	-0.0560457\\
4.7846	-0.0560355\\
4.7872	-0.0560251\\
4.7898	-0.0560147\\
4.7923	-0.0560048\\
4.7949	-0.0559947\\
4.7975	-0.0559846\\
4.8	-0.055975\\
4.8026	-0.0559651\\
4.8052	-0.0559553\\
4.8077	-0.055946\\
4.8103	-0.0559365\\
4.8128	-0.0559273\\
4.8154	-0.055918\\
4.818	-0.0559087\\
4.8205	-0.0558999\\
4.8231	-0.0558908\\
4.8257	-0.0558818\\
4.8282	-0.0558733\\
4.8308	-0.0558645\\
4.8334	-0.0558558\\
4.8359	-0.0558476\\
4.8385	-0.0558391\\
4.841	-0.0558311\\
4.8436	-0.0558228\\
4.8462	-0.0558146\\
4.8487	-0.0558069\\
4.8513	-0.0557989\\
4.8539	-0.055791\\
4.8564	-0.0557836\\
4.859	-0.0557759\\
4.8616	-0.0557683\\
4.8641	-0.0557612\\
4.8667	-0.0557538\\
4.8692	-0.0557468\\
4.8718	-0.0557397\\
4.8744	-0.0557326\\
4.8769	-0.0557259\\
4.8795	-0.0557191\\
4.8821	-0.0557123\\
4.8846	-0.0557059\\
4.8872	-0.0556994\\
4.8898	-0.0556929\\
4.8923	-0.0556868\\
4.8949	-0.0556805\\
4.8974	-0.0556746\\
4.9	-0.0556686\\
4.9026	-0.0556626\\
4.9051	-0.055657\\
4.9077	-0.0556512\\
4.9103	-0.0556456\\
4.9128	-0.0556403\\
4.9154	-0.0556348\\
4.918	-0.0556294\\
4.9205	-0.0556244\\
4.9231	-0.0556192\\
4.9257	-0.0556142\\
4.9282	-0.0556094\\
4.9308	-0.0556045\\
4.9333	-0.0556\\
4.9359	-0.0555953\\
4.9385	-0.0555907\\
4.941	-0.0555864\\
4.9436	-0.055582\\
4.9462	-0.0555777\\
4.9487	-0.0555737\\
4.9513	-0.0555696\\
4.9539	-0.0555656\\
4.9564	-0.0555619\\
4.959	-0.0555581\\
4.9615	-0.0555545\\
4.9641	-0.0555509\\
4.9667	-0.0555474\\
4.9692	-0.0555441\\
4.9718	-0.0555407\\
4.9744	-0.0555375\\
4.9769	-0.0555345\\
4.9795	-0.0555314\\
4.9821	-0.0555284\\
4.9846	-0.0555257\\
4.9872	-0.0555229\\
4.9897	-0.0555203\\
4.9923	-0.0555177\\
4.9949	-0.0555152\\
4.9974	-0.0555128\\
5	-0.0555105\\
};
\addplot [color=mycolor19, line width=1.0pt, forget plot]
  table[row sep=crcr]{%
-0.125	6.8406e-08\\
-0.1224	6.98224e-08\\
-0.1199	7.11843e-08\\
-0.1173	7.26006e-08\\
-0.1147	7.40168e-08\\
-0.1122	7.53784e-08\\
-0.1096	7.67944e-08\\
-0.1071	7.81558e-08\\
-0.1045	7.95716e-08\\
-0.1019	8.09873e-08\\
-0.0994	8.23484e-08\\
-0.0968	8.37638e-08\\
-0.0942	8.51792e-08\\
-0.0917	8.654e-08\\
-0.0891	8.79551e-08\\
-0.0865	8.93702e-08\\
-0.084	9.07307e-08\\
-0.0814	9.21456e-08\\
-0.0789	9.3506e-08\\
-0.0763	9.49207e-08\\
-0.0737	9.63353e-08\\
-0.0712	9.76954e-08\\
-0.0686	9.91097e-08\\
-0.066	1.00524e-07\\
-0.0635	1.01884e-07\\
-0.0609	1.03298e-07\\
-0.0583	1.04712e-07\\
-0.0558	1.06071e-07\\
-0.0532	1.07485e-07\\
-0.0507	1.08844e-07\\
-0.0481	1.10258e-07\\
-0.0455	1.11671e-07\\
-0.043	1.1303e-07\\
-0.0404	1.14444e-07\\
-0.0378	1.15857e-07\\
-0.0353	1.17215e-07\\
-0.0327	1.18628e-07\\
-0.0301	1.20041e-07\\
-0.0276	1.214e-07\\
-0.025	1.22812e-07\\
-0.0224	1.24225e-07\\
-0.0199	1.25583e-07\\
-0.0173	1.26995e-07\\
-0.0148	1.28353e-07\\
-0.0122	1.29765e-07\\
-0.0096	1.31177e-07\\
-0.0071	1.32535e-07\\
-0.0045	1.33947e-07\\
-0.0019	1.35359e-07\\
0.0006	1.36716e-07\\
0.0032	5.62123e-05\\
0.0058	0.000601299\\
0.0083	0.0011698\\
0.0109	0.0016128\\
0.0134	0.00204549\\
0.016	0.00258253\\
0.0186	0.00316853\\
0.0211	0.00370052\\
0.0237	0.00417751\\
0.0263	0.00459416\\
0.0288	0.00498379\\
0.0314	0.00541319\\
0.034	0.00586953\\
0.0365	0.00631284\\
0.0391	0.00675597\\
0.0416	0.00715858\\
0.0442	0.00756291\\
0.0468	0.00796805\\
0.0493	0.00836608\\
0.0519	0.00878542\\
0.0545	0.00920023\\
0.057	0.0095876\\
0.0596	0.00997877\\
0.0622	0.0103648\\
0.0647	0.0107379\\
0.0673	0.0111298\\
0.0698	0.0115063\\
0.0724	0.0118912\\
0.075	0.0122664\\
0.0775	0.0126203\\
0.0801	0.0129868\\
0.0827	0.0133554\\
0.0852	0.0137112\\
0.0878	0.0140777\\
0.0904	0.0144365\\
0.0929	0.0147737\\
0.0955	0.0151201\\
0.098	0.0154534\\
0.1006	0.0158017\\
0.1032	0.0161488\\
0.1057	0.0164773\\
0.1083	0.0168109\\
0.1109	0.0171383\\
0.1134	0.0174511\\
0.116	0.0177774\\
0.1186	0.018104\\
0.1211	0.0184149\\
0.1237	0.0187313\\
0.1263	0.0190406\\
0.1288	0.0193341\\
0.1314	0.019639\\
0.1339	0.0199329\\
0.1365	0.0202373\\
0.1391	0.0205366\\
0.1416	0.0208177\\
0.1442	0.0211046\\
0.1468	0.0213894\\
0.1493	0.0216638\\
0.1519	0.0219491\\
0.1545	0.0222313\\
0.157	0.022497\\
0.1596	0.0227672\\
0.1621	0.0230236\\
0.1647	0.0232898\\
0.1673	0.0235564\\
0.1698	0.0238115\\
0.1724	0.0240724\\
0.175	0.0243273\\
0.1775	0.0245679\\
0.1801	0.0248162\\
0.1827	0.0250648\\
0.1852	0.0253035\\
0.1878	0.025549\\
0.1903	0.0257801\\
0.1929	0.0260152\\
0.1955	0.026247\\
0.198	0.0264692\\
0.2006	0.0267003\\
0.2032	0.02693\\
0.2057	0.0271471\\
0.2083	0.0273675\\
0.2109	0.0275835\\
0.2134	0.0277894\\
0.216	0.0280032\\
0.2185	0.0282082\\
0.2211	0.0284186\\
0.2237	0.0286241\\
0.2262	0.0288169\\
0.2288	0.0290143\\
0.2314	0.0292104\\
0.2339	0.0293985\\
0.2365	0.0295922\\
0.2391	0.0297819\\
0.2416	0.0299593\\
0.2442	0.0301395\\
0.2467	0.0303106\\
0.2493	0.0304875\\
0.2519	0.030663\\
0.2544	0.0308287\\
0.257	0.0309962\\
0.2596	0.0311586\\
0.2621	0.0313113\\
0.2647	0.0314683\\
0.2673	0.0316239\\
0.2698	0.0317711\\
0.2724	0.0319198\\
0.2749	0.0320578\\
0.2775	0.0321967\\
0.2801	0.0323326\\
0.2826	0.0324615\\
0.2852	0.0325934\\
0.2878	0.0327216\\
0.2903	0.0328399\\
0.2929	0.0329577\\
0.2955	0.0330715\\
0.298	0.0331783\\
0.3006	0.0332873\\
0.3032	0.033393\\
0.3057	0.0334902\\
0.3083	0.0335859\\
0.3108	0.0336732\\
0.3134	0.0337605\\
0.316	0.0338452\\
0.3185	0.033924\\
0.3211	0.0340018\\
0.3237	0.0340744\\
0.3262	0.034139\\
0.3288	0.034202\\
0.3314	0.0342618\\
0.3339	0.0343167\\
0.3365	0.0343702\\
0.339	0.0344172\\
0.3416	0.0344607\\
0.3442	0.0344993\\
0.3467	0.0345328\\
0.3493	0.0345648\\
0.3519	0.0345937\\
0.3544	0.0346176\\
0.357	0.0346374\\
0.3596	0.034652\\
0.3621	0.034662\\
0.3647	0.0346693\\
0.3672	0.0346737\\
0.3698	0.0346747\\
0.3724	0.0346713\\
0.3749	0.0346632\\
0.3775	0.0346504\\
0.3801	0.034634\\
0.3826	0.0346157\\
0.3852	0.0345937\\
0.3878	0.0345679\\
0.3903	0.0345388\\
0.3929	0.034504\\
0.3954	0.0344671\\
0.398	0.0344258\\
0.4006	0.034382\\
0.4031	0.034337\\
0.4057	0.0342864\\
0.4083	0.0342315\\
0.4108	0.0341751\\
0.4134	0.0341136\\
0.416	0.0340497\\
0.4185	0.0339859\\
0.4211	0.0339166\\
0.4236	0.0338463\\
0.4262	0.0337697\\
0.4288	0.03369\\
0.4313	0.0336112\\
0.4339	0.0335274\\
0.4365	0.0334412\\
0.439	0.0333554\\
0.4416	0.0332629\\
0.4442	0.0331674\\
0.4467	0.0330735\\
0.4493	0.0329742\\
0.4519	0.0328731\\
0.4544	0.0327737\\
0.457	0.0326676\\
0.4595	0.032563\\
0.4621	0.032452\\
0.4647	0.0323395\\
0.4672	0.0322301\\
0.4698	0.0321147\\
0.4724	0.0319972\\
0.4749	0.0318819\\
0.4775	0.03176\\
0.4801	0.0316366\\
0.4826	0.0315171\\
0.4852	0.0313917\\
0.4877	0.0312698\\
0.4903	0.031141\\
0.4929	0.0310105\\
0.4954	0.0308836\\
0.498	0.0307509\\
0.5006	0.0306175\\
0.5031	0.0304885\\
0.5057	0.0303529\\
0.5083	0.0302158\\
0.5108	0.0300828\\
0.5134	0.0299437\\
0.5159	0.0298097\\
0.5185	0.0296699\\
0.5211	0.0295293\\
0.5236	0.0293931\\
0.5262	0.0292504\\
0.5288	0.029107\\
0.5313	0.0289689\\
0.5339	0.0288252\\
0.5365	0.0286812\\
0.539	0.0285422\\
0.5416	0.0283968\\
0.5441	0.0282564\\
0.5467	0.0281102\\
0.5493	0.0279643\\
0.5518	0.027824\\
0.5544	0.0276779\\
0.557	0.0275313\\
0.5595	0.0273899\\
0.5621	0.0272427\\
0.5647	0.0270958\\
0.5672	0.0269549\\
0.5698	0.0268085\\
0.5723	0.0266676\\
0.5749	0.0265208\\
0.5775	0.0263739\\
0.58	0.0262329\\
0.5826	0.0260867\\
0.5852	0.025941\\
0.5877	0.0258011\\
0.5903	0.0256555\\
0.5929	0.0255099\\
0.5954	0.0253701\\
0.598	0.0252253\\
0.6006	0.0250811\\
0.6031	0.0249429\\
0.6057	0.0247993\\
0.6082	0.0246613\\
0.6108	0.024518\\
0.6134	0.0243752\\
0.6159	0.0242386\\
0.6185	0.0240971\\
0.6211	0.023956\\
0.6236	0.0238206\\
0.6262	0.0236799\\
0.6288	0.0235397\\
0.6313	0.0234055\\
0.6339	0.0232666\\
0.6364	0.0231337\\
0.639	0.0229957\\
0.6416	0.022858\\
0.6441	0.0227259\\
0.6467	0.0225891\\
0.6493	0.022453\\
0.6518	0.0223228\\
0.6544	0.0221879\\
0.657	0.0220532\\
0.6595	0.021924\\
0.6621	0.0217901\\
0.6646	0.0216619\\
0.6672	0.0215293\\
0.6698	0.0213972\\
0.6723	0.0212705\\
0.6749	0.021139\\
0.6775	0.0210078\\
0.68	0.0208821\\
0.6826	0.020752\\
0.6852	0.0206225\\
0.6877	0.0204983\\
0.6903	0.0203694\\
0.6928	0.0202456\\
0.6954	0.0201172\\
0.698	0.0199893\\
0.7005	0.0198668\\
0.7031	0.0197398\\
0.7057	0.0196129\\
0.7082	0.0194911\\
0.7108	0.0193645\\
0.7134	0.0192383\\
0.7159	0.0191172\\
0.7185	0.0189917\\
0.721	0.0188711\\
0.7236	0.0187458\\
0.7262	0.0186204\\
0.7287	0.0185\\
0.7313	0.018375\\
0.7339	0.0182502\\
0.7364	0.0181303\\
0.739	0.0180055\\
0.7416	0.0178807\\
0.7441	0.0177605\\
0.7467	0.0176355\\
0.7492	0.0175154\\
0.7518	0.0173905\\
0.7544	0.0172655\\
0.7569	0.017145\\
0.7595	0.0170194\\
0.7621	0.0168936\\
0.7646	0.0167725\\
0.7672	0.0166464\\
0.7698	0.0165201\\
0.7723	0.0163983\\
0.7749	0.0162711\\
0.7775	0.0161435\\
0.78	0.0160205\\
0.7826	0.0158922\\
0.7851	0.0157685\\
0.7877	0.0156393\\
0.7903	0.0155095\\
0.7928	0.0153841\\
0.7954	0.015253\\
0.798	0.0151215\\
0.8005	0.0149944\\
0.8031	0.0148617\\
0.8057	0.0147281\\
0.8082	0.0145989\\
0.8108	0.0144637\\
0.8133	0.014333\\
0.8159	0.0141963\\
0.8185	0.0140588\\
0.821	0.0139257\\
0.8236	0.0137862\\
0.8262	0.0136457\\
0.8287	0.0135097\\
0.8313	0.0133672\\
0.8339	0.0132238\\
0.8364	0.0130849\\
0.839	0.0129392\\
0.8415	0.012798\\
0.8441	0.0126499\\
0.8467	0.0125006\\
0.8492	0.0123559\\
0.8518	0.0122042\\
0.8544	0.0120512\\
0.8569	0.0119027\\
0.8595	0.0117468\\
0.8621	0.0115896\\
0.8646	0.011437\\
0.8672	0.011277\\
0.8697	0.0111217\\
0.8723	0.0109587\\
0.8749	0.010794\\
0.8774	0.0106342\\
0.88	0.0104663\\
0.8826	0.0102969\\
0.8851	0.0101325\\
0.8877	0.00995977\\
0.8903	0.00978527\\
0.8928	0.00961578\\
0.8954	0.00943774\\
0.8979	0.00926489\\
0.9005	0.00908338\\
0.9031	0.00890004\\
0.9056	0.00872194\\
0.9082	0.00853477\\
0.9108	0.00834565\\
0.9133	0.00816197\\
0.9159	0.00796905\\
0.9185	0.00777418\\
0.921	0.0075849\\
0.9236	0.00738599\\
0.9262	0.00718497\\
0.9287	0.00698971\\
0.9313	0.00678462\\
0.9338	0.00658546\\
0.9364	0.00637626\\
0.939	0.00616488\\
0.9415	0.00595954\\
0.9441	0.0057438\\
0.9467	0.00552588\\
0.9492	0.00531429\\
0.9518	0.00509209\\
0.9544	0.00486763\\
0.9569	0.00464964\\
0.9595	0.00442065\\
0.962	0.00419832\\
0.9646	0.00396487\\
0.9672	0.00372916\\
0.9697	0.00350035\\
0.9723	0.00326009\\
0.9749	0.00301745\\
0.9774	0.00278191\\
0.98	0.00253466\\
0.9826	0.00228509\\
0.9851	0.00204293\\
0.9877	0.00178876\\
0.9902	0.00154209\\
0.9928	0.0012832\\
0.9954	0.00102193\\
0.9979	0.000768488\\
1.0005	0.00050261\\
1.0031	0.000234364\\
1.0056	-2.5829e-05\\
1.0082	-0.0002988\\
1.0108	-0.000574172\\
1.0133	-0.000841185\\
1.0159	-0.00112118\\
1.0184	-0.00139261\\
1.021	-0.00167723\\
1.0236	-0.00196424\\
1.0261	-0.00224247\\
1.0287	-0.00253414\\
1.0313	-0.00282813\\
1.0338	-0.00311299\\
1.0364	-0.00341153\\
1.039	-0.0037124\\
1.0415	-0.00400391\\
1.0441	-0.00430936\\
1.0466	-0.00460522\\
1.0492	-0.00491512\\
1.0518	-0.00522727\\
1.0543	-0.00552954\\
1.0569	-0.00584614\\
1.0595	-0.00616499\\
1.062	-0.00647367\\
1.0646	-0.00679684\\
1.0672	-0.00712216\\
1.0697	-0.00743702\\
1.0723	-0.00776659\\
1.0748	-0.00808553\\
1.0774	-0.00841931\\
1.08	-0.00875519\\
1.0825	-0.00908008\\
1.0851	-0.00941997\\
1.0877	-0.00976191\\
1.0902	-0.0100926\\
1.0928	-0.0104385\\
1.0954	-0.0107864\\
1.0979	-0.0111228\\
1.1005	-0.0114744\\
1.1031	-0.011828\\
1.1056	-0.0121697\\
1.1082	-0.0125269\\
1.1107	-0.0128722\\
1.1133	-0.013233\\
1.1159	-0.0135956\\
1.1184	-0.0139459\\
1.121	-0.0143118\\
1.1236	-0.0146795\\
1.1261	-0.0150346\\
1.1287	-0.0154056\\
1.1313	-0.0157781\\
1.1338	-0.0161378\\
1.1364	-0.0165134\\
1.1389	-0.0168759\\
1.1415	-0.0172545\\
1.1441	-0.0176345\\
1.1466	-0.0180012\\
1.1492	-0.018384\\
1.1518	-0.0187681\\
1.1543	-0.0191387\\
1.1569	-0.0195254\\
1.1595	-0.0199133\\
1.162	-0.0202875\\
1.1646	-0.0206778\\
1.1671	-0.0210542\\
1.1697	-0.0214468\\
1.1723	-0.0218404\\
1.1748	-0.0222199\\
1.1774	-0.0226156\\
1.18	-0.0230122\\
1.1825	-0.0233945\\
1.1851	-0.0237929\\
1.1877	-0.0241922\\
1.1902	-0.0245769\\
1.1928	-0.0249779\\
1.1953	-0.0253641\\
1.1979	-0.0257664\\
1.2005	-0.0261694\\
1.203	-0.0265576\\
1.2056	-0.0269618\\
1.2082	-0.0273666\\
1.2107	-0.0277564\\
1.2133	-0.0281622\\
1.2159	-0.0285685\\
1.2184	-0.0289595\\
1.221	-0.0293665\\
1.2235	-0.0297582\\
1.2261	-0.0301658\\
1.2287	-0.0305737\\
1.2312	-0.0309662\\
1.2338	-0.0313744\\
1.2364	-0.0317828\\
1.2389	-0.0321756\\
1.2415	-0.0325842\\
1.2441	-0.0329927\\
1.2466	-0.0333856\\
1.2492	-0.0337941\\
1.2518	-0.0342024\\
1.2543	-0.034595\\
1.2569	-0.035003\\
1.2594	-0.0353952\\
1.262	-0.0358028\\
1.2646	-0.0362101\\
1.2671	-0.0366014\\
1.2697	-0.037008\\
1.2723	-0.0374142\\
1.2748	-0.0378043\\
1.2774	-0.0382095\\
1.28	-0.0386142\\
1.2825	-0.0390028\\
1.2851	-0.0394064\\
1.2876	-0.0397938\\
1.2902	-0.040196\\
1.2928	-0.0405976\\
1.2953	-0.0409829\\
1.2979	-0.0413829\\
1.3005	-0.0417821\\
1.303	-0.0421651\\
1.3056	-0.0425625\\
1.3082	-0.0429589\\
1.3107	-0.0433392\\
1.3133	-0.0437337\\
1.3158	-0.0441121\\
1.3184	-0.0445045\\
1.321	-0.0448957\\
1.3235	-0.0452709\\
1.3261	-0.0456598\\
1.3287	-0.0460476\\
1.3312	-0.0464192\\
1.3338	-0.0468044\\
1.3364	-0.0471883\\
1.3389	-0.0475562\\
1.3415	-0.0479373\\
1.344	-0.0483025\\
1.3466	-0.0486809\\
1.3492	-0.0490578\\
1.3517	-0.0494187\\
1.3543	-0.0497925\\
1.3569	-0.0501648\\
1.3594	-0.0505212\\
1.362	-0.0508903\\
1.3646	-0.0512577\\
1.3671	-0.0516093\\
1.3697	-0.0519733\\
1.3722	-0.0523217\\
1.3748	-0.0526822\\
1.3774	-0.0530409\\
1.3799	-0.0533841\\
1.3825	-0.0537392\\
1.3851	-0.0540923\\
1.3876	-0.0544301\\
1.3902	-0.0547795\\
1.3928	-0.0551269\\
1.3953	-0.055459\\
1.3979	-0.0558025\\
1.4005	-0.0561439\\
1.403	-0.0564702\\
1.4056	-0.0568075\\
1.4081	-0.0571299\\
1.4107	-0.057463\\
1.4133	-0.057794\\
1.4158	-0.0581102\\
1.4184	-0.0584369\\
1.421	-0.0587614\\
1.4235	-0.0590713\\
1.4261	-0.0593913\\
1.4287	-0.0597091\\
1.4312	-0.0600125\\
1.4338	-0.0603258\\
1.4363	-0.0606248\\
1.4389	-0.0609334\\
1.4415	-0.0612397\\
1.444	-0.061532\\
1.4466	-0.0618336\\
1.4492	-0.0621328\\
1.4517	-0.0624182\\
1.4543	-0.0627126\\
1.4569	-0.0630046\\
1.4594	-0.063283\\
1.462	-0.0635701\\
1.4645	-0.0638438\\
1.4671	-0.064126\\
1.4697	-0.0644057\\
1.4722	-0.0646722\\
1.4748	-0.064947\\
1.4774	-0.0652192\\
1.4799	-0.0654785\\
1.4825	-0.0657456\\
1.4851	-0.0660102\\
1.4876	-0.0662622\\
1.4902	-0.0665217\\
1.4927	-0.0667688\\
1.4953	-0.0670232\\
1.4979	-0.067275\\
1.5004	-0.0675146\\
1.503	-0.0677613\\
1.5056	-0.0680053\\
1.5081	-0.0682375\\
1.5107	-0.0684764\\
1.5133	-0.0687127\\
1.5158	-0.0689373\\
1.5184	-0.0691684\\
1.5209	-0.0693881\\
1.5235	-0.069614\\
1.5261	-0.0698372\\
1.5286	-0.0700493\\
1.5312	-0.0702674\\
1.5338	-0.0704828\\
1.5363	-0.0706874\\
1.5389	-0.0708976\\
1.5415	-0.0711051\\
1.544	-0.0713022\\
1.5466	-0.0715046\\
1.5491	-0.0716967\\
1.5517	-0.0718939\\
1.5543	-0.0720884\\
1.5568	-0.072273\\
1.5594	-0.0724624\\
1.562	-0.0726492\\
1.5645	-0.0728264\\
1.5671	-0.073008\\
1.5697	-0.0731871\\
1.5722	-0.0733568\\
1.5748	-0.0735307\\
1.5774	-0.0737021\\
1.5799	-0.0738645\\
1.5825	-0.0740308\\
1.585	-0.0741883\\
1.5876	-0.0743496\\
1.5902	-0.0745083\\
1.5927	-0.0746586\\
1.5953	-0.0748123\\
1.5979	-0.0749636\\
1.6004	-0.0751067\\
1.603	-0.075253\\
1.6056	-0.0753968\\
1.6081	-0.0755328\\
1.6107	-0.0756718\\
1.6132	-0.0758031\\
1.6158	-0.0759373\\
1.6184	-0.076069\\
1.6209	-0.0761934\\
1.6235	-0.0763204\\
1.6261	-0.076445\\
1.6286	-0.0765626\\
1.6312	-0.0766825\\
1.6338	-0.0768001\\
1.6363	-0.0769109\\
1.6389	-0.0770239\\
1.6414	-0.0771303\\
1.644	-0.0772387\\
1.6466	-0.0773449\\
1.6491	-0.0774448\\
1.6517	-0.0775465\\
1.6543	-0.077646\\
1.6568	-0.0777395\\
1.6594	-0.0778346\\
1.662	-0.0779274\\
1.6645	-0.0780147\\
1.6671	-0.0781033\\
1.6696	-0.0781865\\
1.6722	-0.0782709\\
1.6748	-0.0783532\\
1.6773	-0.0784303\\
1.6799	-0.0785085\\
1.6825	-0.0785846\\
1.685	-0.0786559\\
1.6876	-0.078728\\
1.6902	-0.0787981\\
1.6927	-0.0788636\\
1.6953	-0.0789298\\
1.6978	-0.0789916\\
1.7004	-0.0790539\\
1.703	-0.0791143\\
1.7055	-0.0791706\\
1.7081	-0.0792273\\
1.7107	-0.0792821\\
1.7132	-0.0793331\\
1.7158	-0.0793842\\
1.7184	-0.0794336\\
1.7209	-0.0794794\\
1.7235	-0.0795252\\
1.7261	-0.0795693\\
1.7286	-0.0796101\\
1.7312	-0.0796507\\
1.7337	-0.0796882\\
1.7363	-0.0797256\\
1.7389	-0.0797613\\
1.7414	-0.079794\\
1.744	-0.0798265\\
1.7466	-0.0798573\\
1.7491	-0.0798855\\
1.7517	-0.0799132\\
1.7543	-0.0799394\\
1.7568	-0.0799631\\
1.7594	-0.0799863\\
1.7619	-0.0800072\\
1.7645	-0.0800275\\
1.7671	-0.0800463\\
1.7696	-0.080063\\
1.7722	-0.080079\\
1.7748	-0.0800936\\
1.7773	-0.0801063\\
1.7799	-0.0801182\\
1.7825	-0.0801287\\
1.785	-0.0801376\\
1.7876	-0.0801456\\
1.7901	-0.080152\\
1.7927	-0.0801574\\
1.7953	-0.0801616\\
1.7978	-0.0801645\\
1.8004	-0.0801663\\
1.803	-0.0801668\\
1.8055	-0.0801663\\
1.8081	-0.0801646\\
1.8107	-0.0801617\\
1.8132	-0.0801579\\
1.8158	-0.0801528\\
1.8183	-0.0801469\\
1.8209	-0.0801398\\
1.8235	-0.0801315\\
1.826	-0.0801226\\
1.8286	-0.0801124\\
1.8312	-0.0801012\\
1.8337	-0.0800894\\
1.8363	-0.0800763\\
1.8389	-0.0800621\\
1.8414	-0.0800477\\
1.844	-0.0800318\\
1.8465	-0.0800156\\
1.8491	-0.0799979\\
1.8517	-0.0799794\\
1.8542	-0.0799608\\
1.8568	-0.0799406\\
1.8594	-0.0799196\\
1.8619	-0.0798986\\
1.8645	-0.0798761\\
1.8671	-0.0798527\\
1.8696	-0.0798296\\
1.8722	-0.0798048\\
1.8747	-0.0797803\\
1.8773	-0.0797541\\
1.8799	-0.0797272\\
1.8824	-0.0797007\\
1.885	-0.0796725\\
1.8876	-0.0796436\\
1.8901	-0.0796153\\
1.8927	-0.0795852\\
1.8953	-0.0795545\\
1.8978	-0.0795244\\
1.9004	-0.0794926\\
1.903	-0.0794602\\
1.9055	-0.0794285\\
1.9081	-0.079395\\
1.9106	-0.0793623\\
1.9132	-0.0793279\\
1.9158	-0.0792929\\
1.9183	-0.0792588\\
1.9209	-0.0792228\\
1.9235	-0.0791864\\
1.926	-0.079151\\
1.9286	-0.0791138\\
1.9312	-0.0790761\\
1.9337	-0.0790395\\
1.9363	-0.079001\\
1.9388	-0.0789636\\
1.9414	-0.0789244\\
1.944	-0.0788848\\
1.9465	-0.0788464\\
1.9491	-0.0788061\\
1.9517	-0.0787654\\
1.9542	-0.078726\\
1.9568	-0.0786848\\
1.9594	-0.0786432\\
1.9619	-0.078603\\
1.9645	-0.0785608\\
1.967	-0.07852\\
1.9696	-0.0784773\\
1.9722	-0.0784344\\
1.9747	-0.0783929\\
1.9773	-0.0783495\\
1.9799	-0.0783058\\
1.9824	-0.0782636\\
1.985	-0.0782195\\
1.9876	-0.0781752\\
1.9901	-0.0781325\\
1.9927	-0.0780878\\
1.9952	-0.0780446\\
1.9978	-0.0779996\\
2.0004	-0.0779544\\
2.0029	-0.0779107\\
2.0055	-0.0778652\\
2.0081	-0.0778195\\
2.0106	-0.0777755\\
2.0132	-0.0777295\\
2.0158	-0.0776834\\
2.0183	-0.077639\\
2.0209	-0.0775927\\
2.0234	-0.077548\\
2.026	-0.0775015\\
2.0286	-0.0774548\\
2.0311	-0.0774099\\
2.0337	-0.077363\\
2.0363	-0.0773161\\
2.0388	-0.0772709\\
2.0414	-0.0772238\\
2.044	-0.0771767\\
2.0465	-0.0771313\\
2.0491	-0.077084\\
2.0517	-0.0770366\\
2.0542	-0.076991\\
2.0568	-0.0769435\\
2.0593	-0.0768979\\
2.0619	-0.0768503\\
2.0645	-0.0768027\\
2.067	-0.0767569\\
2.0696	-0.0767092\\
2.0722	-0.0766615\\
2.0747	-0.0766156\\
2.0773	-0.0765678\\
2.0799	-0.07652\\
2.0824	-0.076474\\
2.085	-0.0764261\\
2.0875	-0.0763801\\
2.0901	-0.0763322\\
2.0927	-0.0762843\\
2.0952	-0.0762382\\
2.0978	-0.0761903\\
2.1004	-0.0761423\\
2.1029	-0.0760962\\
2.1055	-0.0760482\\
2.1081	-0.0760002\\
2.1106	-0.0759541\\
2.1132	-0.0759061\\
2.1157	-0.07586\\
2.1183	-0.075812\\
2.1209	-0.075764\\
2.1234	-0.0757178\\
2.126	-0.0756698\\
2.1286	-0.0756218\\
2.1311	-0.0755756\\
2.1337	-0.0755276\\
2.1363	-0.0754796\\
2.1388	-0.0754334\\
2.1414	-0.0753854\\
2.1439	-0.0753392\\
2.1465	-0.0752912\\
2.1491	-0.0752431\\
2.1516	-0.0751969\\
2.1542	-0.0751489\\
2.1568	-0.0751008\\
2.1593	-0.0750546\\
2.1619	-0.0750066\\
2.1645	-0.0749585\\
2.167	-0.0749123\\
2.1696	-0.0748642\\
2.1721	-0.074818\\
2.1747	-0.0747699\\
2.1773	-0.0747218\\
2.1798	-0.0746755\\
2.1824	-0.0746274\\
2.185	-0.0745792\\
2.1875	-0.0745329\\
2.1901	-0.0744847\\
2.1927	-0.0744365\\
2.1952	-0.0743901\\
2.1978	-0.0743419\\
2.2004	-0.0742936\\
2.2029	-0.0742472\\
2.2055	-0.0741988\\
2.208	-0.0741524\\
2.2106	-0.074104\\
2.2132	-0.0740556\\
2.2157	-0.074009\\
2.2183	-0.0739605\\
2.2209	-0.073912\\
2.2234	-0.0738653\\
2.226	-0.0738167\\
2.2286	-0.073768\\
2.2311	-0.0737212\\
2.2337	-0.0736725\\
2.2362	-0.0736256\\
2.2388	-0.0735768\\
2.2414	-0.0735279\\
2.2439	-0.0734808\\
2.2465	-0.0734318\\
2.2491	-0.0733828\\
2.2516	-0.0733356\\
2.2542	-0.0732864\\
2.2568	-0.0732372\\
2.2593	-0.0731897\\
2.2619	-0.0731404\\
2.2644	-0.0730928\\
2.267	-0.0730433\\
2.2696	-0.0729937\\
2.2721	-0.072946\\
2.2747	-0.0728963\\
2.2773	-0.0728464\\
2.2798	-0.0727985\\
2.2824	-0.0727485\\
2.285	-0.0726984\\
2.2875	-0.0726502\\
2.2901	-0.0725999\\
2.2926	-0.0725515\\
2.2952	-0.0725011\\
2.2978	-0.0724506\\
2.3003	-0.0724019\\
2.3029	-0.0723512\\
2.3055	-0.0723004\\
2.308	-0.0722514\\
2.3106	-0.0722004\\
2.3132	-0.0721493\\
2.3157	-0.0721\\
2.3183	-0.0720486\\
2.3208	-0.0719992\\
2.3234	-0.0719476\\
2.326	-0.0718959\\
2.3285	-0.0718461\\
2.3311	-0.0717942\\
2.3337	-0.0717422\\
2.3362	-0.071692\\
2.3388	-0.0716397\\
2.3414	-0.0715873\\
2.3439	-0.0715368\\
2.3465	-0.0714842\\
2.349	-0.0714334\\
2.3516	-0.0713805\\
2.3542	-0.0713275\\
2.3567	-0.0712764\\
2.3593	-0.0712231\\
2.3619	-0.0711697\\
2.3644	-0.0711182\\
2.367	-0.0710645\\
2.3696	-0.0710106\\
2.3721	-0.0709587\\
2.3747	-0.0709046\\
2.3773	-0.0708504\\
2.3798	-0.0707981\\
2.3824	-0.0707436\\
2.3849	-0.070691\\
2.3875	-0.0706362\\
2.3901	-0.0705813\\
2.3926	-0.0705283\\
2.3952	-0.0704731\\
2.3978	-0.0704177\\
2.4003	-0.0703643\\
2.4029	-0.0703086\\
2.4055	-0.0702528\\
2.408	-0.070199\\
2.4106	-0.0701429\\
2.4131	-0.0700888\\
2.4157	-0.0700324\\
2.4183	-0.0699759\\
2.4208	-0.0699214\\
2.4234	-0.0698646\\
2.426	-0.0698076\\
2.4285	-0.0697527\\
2.4311	-0.0696954\\
2.4337	-0.069638\\
2.4362	-0.0695826\\
2.4388	-0.0695249\\
2.4413	-0.0694693\\
2.4439	-0.0694113\\
2.4465	-0.0693531\\
2.449	-0.0692971\\
2.4516	-0.0692387\\
2.4542	-0.0691801\\
2.4567	-0.0691236\\
2.4593	-0.0690648\\
2.4619	-0.0690058\\
2.4644	-0.0689489\\
2.467	-0.0688897\\
2.4695	-0.0688325\\
2.4721	-0.068773\\
2.4747	-0.0687133\\
2.4772	-0.0686558\\
2.4798	-0.0685958\\
2.4824	-0.0685357\\
2.4849	-0.0684778\\
2.4875	-0.0684174\\
2.4901	-0.0683569\\
2.4926	-0.0682986\\
2.4952	-0.0682379\\
2.4977	-0.0681794\\
2.5003	-0.0681184\\
2.5029	-0.0680572\\
2.5054	-0.0679983\\
2.508	-0.067937\\
2.5106	-0.0678755\\
2.5131	-0.0678162\\
2.5157	-0.0677545\\
2.5183	-0.0676927\\
2.5208	-0.0676331\\
2.5234	-0.067571\\
2.526	-0.0675089\\
2.5285	-0.067449\\
2.5311	-0.0673866\\
2.5336	-0.0673265\\
2.5362	-0.0672639\\
2.5388	-0.0672012\\
2.5413	-0.0671408\\
2.5439	-0.0670779\\
2.5465	-0.0670149\\
2.549	-0.0669543\\
2.5516	-0.0668911\\
2.5542	-0.0668278\\
2.5567	-0.0667669\\
2.5593	-0.0667035\\
2.5618	-0.0666424\\
2.5644	-0.0665788\\
2.567	-0.0665152\\
2.5695	-0.0664539\\
2.5721	-0.0663901\\
2.5747	-0.0663262\\
2.5772	-0.0662647\\
2.5798	-0.0662006\\
2.5824	-0.0661366\\
2.5849	-0.0660749\\
2.5875	-0.0660107\\
2.59	-0.0659489\\
2.5926	-0.0658846\\
2.5952	-0.0658202\\
2.5977	-0.0657583\\
2.6003	-0.0656938\\
2.6029	-0.0656293\\
2.6054	-0.0655673\\
2.608	-0.0655027\\
2.6106	-0.0654381\\
2.6131	-0.0653759\\
2.6157	-0.0653113\\
2.6182	-0.0652491\\
2.6208	-0.0651843\\
2.6234	-0.0651196\\
2.6259	-0.0650573\\
2.6285	-0.0649925\\
2.6311	-0.0649277\\
2.6336	-0.0648654\\
2.6362	-0.0648006\\
2.6388	-0.0647358\\
2.6413	-0.0646735\\
2.6439	-0.0646087\\
2.6464	-0.0645464\\
2.649	-0.0644816\\
2.6516	-0.0644168\\
2.6541	-0.0643545\\
2.6567	-0.0642898\\
2.6593	-0.064225\\
2.6618	-0.0641628\\
2.6644	-0.0640981\\
2.667	-0.0640334\\
2.6695	-0.0639713\\
2.6721	-0.0639067\\
2.6746	-0.0638446\\
2.6772	-0.0637801\\
2.6798	-0.0637156\\
2.6823	-0.0636537\\
2.6849	-0.0635893\\
2.6875	-0.063525\\
2.69	-0.0634632\\
2.6926	-0.063399\\
2.6952	-0.0633348\\
2.6977	-0.0632732\\
2.7003	-0.0632091\\
2.7029	-0.0631452\\
2.7054	-0.0630838\\
2.708	-0.06302\\
2.7105	-0.0629587\\
2.7131	-0.062895\\
2.7157	-0.0628315\\
2.7182	-0.0627705\\
2.7208	-0.0627071\\
2.7234	-0.0626438\\
2.7259	-0.062583\\
2.7285	-0.0625199\\
2.7311	-0.0624569\\
2.7336	-0.0623965\\
2.7362	-0.0623337\\
2.7387	-0.0622735\\
2.7413	-0.0622109\\
2.7439	-0.0621485\\
2.7464	-0.0620886\\
2.749	-0.0620264\\
2.7516	-0.0619643\\
2.7541	-0.0619048\\
2.7567	-0.061843\\
2.7593	-0.0617813\\
2.7618	-0.0617221\\
2.7644	-0.0616607\\
2.7669	-0.0616018\\
2.7695	-0.0615407\\
2.7721	-0.0614798\\
2.7746	-0.0614213\\
2.7772	-0.0613607\\
2.7798	-0.0613002\\
2.7823	-0.0612422\\
2.7849	-0.061182\\
2.7875	-0.061122\\
2.79	-0.0610644\\
2.7926	-0.0610048\\
2.7951	-0.0609476\\
2.7977	-0.0608882\\
2.8003	-0.0608291\\
2.8028	-0.0607724\\
2.8054	-0.0607136\\
2.808	-0.060655\\
2.8105	-0.0605988\\
2.8131	-0.0605406\\
2.8157	-0.0604826\\
2.8182	-0.060427\\
2.8208	-0.0603694\\
2.8233	-0.0603141\\
2.8259	-0.0602569\\
2.8285	-0.0601999\\
2.831	-0.0601452\\
2.8336	-0.0600886\\
2.8362	-0.0600322\\
2.8387	-0.0599782\\
2.8413	-0.0599222\\
2.8439	-0.0598665\\
2.8464	-0.0598131\\
2.849	-0.0597578\\
2.8516	-0.0597027\\
2.8541	-0.05965\\
2.8567	-0.0595954\\
2.8592	-0.0595431\\
2.8618	-0.0594889\\
2.8644	-0.059435\\
2.8669	-0.0593834\\
2.8695	-0.0593299\\
2.8721	-0.0592767\\
2.8746	-0.0592258\\
2.8772	-0.0591731\\
2.8798	-0.0591206\\
2.8823	-0.0590704\\
2.8849	-0.0590185\\
2.8874	-0.0589687\\
2.89	-0.0589173\\
2.8926	-0.0588661\\
2.8951	-0.0588171\\
2.8977	-0.0587664\\
2.9003	-0.058716\\
2.9028	-0.0586678\\
2.9054	-0.0586179\\
2.908	-0.0585683\\
2.9105	-0.0585208\\
2.9131	-0.0584718\\
2.9156	-0.0584248\\
2.9182	-0.0583763\\
2.9208	-0.058328\\
2.9233	-0.0582819\\
2.9259	-0.0582342\\
2.9285	-0.0581868\\
2.931	-0.0581414\\
2.9336	-0.0580946\\
2.9362	-0.058048\\
2.9387	-0.0580035\\
2.9413	-0.0579575\\
2.9438	-0.0579135\\
2.9464	-0.0578681\\
2.949	-0.0578229\\
2.9515	-0.0577798\\
2.9541	-0.0577353\\
2.9567	-0.057691\\
2.9592	-0.0576487\\
2.9618	-0.0576051\\
2.9644	-0.0575617\\
2.9669	-0.0575203\\
2.9695	-0.0574775\\
2.972	-0.0574367\\
2.9746	-0.0573945\\
2.9772	-0.0573526\\
2.9797	-0.0573126\\
2.9823	-0.0572714\\
2.9849	-0.0572304\\
2.9874	-0.0571913\\
2.99	-0.057151\\
2.9926	-0.0571109\\
2.9951	-0.0570727\\
2.9977	-0.0570333\\
3.0003	-0.0569942\\
3.0028	-0.0569569\\
3.0054	-0.0569184\\
3.0079	-0.0568817\\
3.0105	-0.0568438\\
3.0131	-0.0568063\\
3.0156	-0.0567705\\
3.0182	-0.0567336\\
3.0208	-0.056697\\
3.0233	-0.0566621\\
3.0259	-0.0566261\\
3.0285	-0.0565904\\
3.031	-0.0565564\\
3.0336	-0.0565214\\
3.0361	-0.056488\\
3.0387	-0.0564537\\
3.0413	-0.0564196\\
3.0438	-0.0563871\\
3.0464	-0.0563537\\
3.049	-0.0563206\\
3.0515	-0.0562891\\
3.0541	-0.0562566\\
3.0567	-0.0562244\\
3.0592	-0.0561938\\
3.0618	-0.0561623\\
3.0643	-0.0561323\\
3.0669	-0.0561015\\
3.0695	-0.0560709\\
3.072	-0.0560419\\
3.0746	-0.056012\\
3.0772	-0.0559824\\
3.0797	-0.0559543\\
3.0823	-0.0559253\\
3.0849	-0.0558967\\
3.0874	-0.0558695\\
3.09	-0.0558415\\
3.0925	-0.0558149\\
3.0951	-0.0557876\\
3.0977	-0.0557606\\
3.1002	-0.0557349\\
3.1028	-0.0557086\\
3.1054	-0.0556825\\
3.1079	-0.0556578\\
3.1105	-0.0556324\\
3.1131	-0.0556073\\
3.1156	-0.0555835\\
3.1182	-0.0555591\\
3.1207	-0.0555359\\
3.1233	-0.0555121\\
3.1259	-0.0554886\\
3.1284	-0.0554663\\
3.131	-0.0554435\\
3.1336	-0.055421\\
3.1361	-0.0553996\\
3.1387	-0.0553777\\
3.1413	-0.0553562\\
3.1438	-0.0553357\\
3.1464	-0.0553148\\
3.1489	-0.055295\\
3.1515	-0.0552747\\
3.1541	-0.0552547\\
3.1566	-0.0552358\\
3.1592	-0.0552164\\
3.1618	-0.0551974\\
3.1643	-0.0551794\\
3.1669	-0.0551609\\
3.1695	-0.0551428\\
3.172	-0.0551257\\
3.1746	-0.0551082\\
3.1772	-0.0550911\\
3.1797	-0.0550749\\
3.1823	-0.0550583\\
3.1848	-0.0550427\\
3.1874	-0.0550267\\
3.19	-0.0550111\\
3.1925	-0.0549964\\
3.1951	-0.0549813\\
3.1977	-0.0549666\\
3.2002	-0.0549528\\
3.2028	-0.0549387\\
3.2054	-0.0549249\\
3.2079	-0.0549119\\
3.2105	-0.0548987\\
3.213	-0.0548863\\
3.2156	-0.0548737\\
3.2182	-0.0548614\\
3.2207	-0.0548498\\
3.2233	-0.0548381\\
3.2259	-0.0548267\\
3.2284	-0.054816\\
3.231	-0.0548052\\
3.2336	-0.0547947\\
3.2361	-0.0547849\\
3.2387	-0.0547749\\
3.2412	-0.0547656\\
3.2438	-0.0547563\\
3.2464	-0.0547472\\
3.2489	-0.0547388\\
3.2515	-0.0547303\\
3.2541	-0.0547221\\
3.2566	-0.0547144\\
3.2592	-0.0547068\\
3.2618	-0.0546995\\
3.2643	-0.0546927\\
3.2669	-0.0546859\\
3.2694	-0.0546796\\
3.272	-0.0546734\\
3.2746	-0.0546674\\
3.2771	-0.0546619\\
3.2797	-0.0546565\\
3.2823	-0.0546514\\
3.2848	-0.0546467\\
3.2874	-0.0546422\\
3.29	-0.0546378\\
3.2925	-0.054634\\
3.2951	-0.0546302\\
3.2976	-0.0546268\\
3.3002	-0.0546236\\
3.3028	-0.0546206\\
3.3053	-0.054618\\
3.3079	-0.0546156\\
3.3105	-0.0546134\\
3.313	-0.0546115\\
3.3156	-0.0546099\\
3.3182	-0.0546085\\
3.3207	-0.0546074\\
3.3233	-0.0546065\\
3.3259	-0.0546058\\
3.3284	-0.0546055\\
3.331	-0.0546054\\
3.3335	-0.0546055\\
3.3361	-0.0546058\\
3.3387	-0.0546065\\
3.3412	-0.0546073\\
3.3438	-0.0546084\\
3.3464	-0.0546098\\
3.3489	-0.0546113\\
3.3515	-0.0546131\\
3.3541	-0.0546152\\
3.3566	-0.0546175\\
3.3592	-0.05462\\
3.3617	-0.0546227\\
3.3643	-0.0546257\\
3.3669	-0.054629\\
3.3694	-0.0546324\\
3.372	-0.0546361\\
3.3746	-0.05464\\
3.3771	-0.0546441\\
3.3797	-0.0546485\\
3.3823	-0.0546531\\
3.3848	-0.0546578\\
3.3874	-0.0546629\\
3.3899	-0.054668\\
3.3925	-0.0546735\\
3.3951	-0.0546792\\
3.3976	-0.054685\\
3.4002	-0.0546911\\
3.4028	-0.0546975\\
3.4053	-0.0547039\\
3.4079	-0.0547107\\
3.4105	-0.0547177\\
3.413	-0.0547247\\
3.4156	-0.0547321\\
3.4181	-0.0547394\\
3.4207	-0.0547473\\
3.4233	-0.0547553\\
3.4258	-0.0547632\\
3.4284	-0.0547717\\
3.431	-0.0547803\\
3.4335	-0.0547888\\
3.4361	-0.0547978\\
3.4387	-0.054807\\
3.4412	-0.0548161\\
3.4438	-0.0548257\\
3.4463	-0.0548351\\
3.4489	-0.0548451\\
3.4515	-0.0548552\\
3.454	-0.0548652\\
3.4566	-0.0548757\\
3.4592	-0.0548864\\
3.4617	-0.0548968\\
3.4643	-0.0549079\\
3.4669	-0.0549191\\
3.4694	-0.0549301\\
3.472	-0.0549416\\
3.4745	-0.0549529\\
3.4771	-0.0549648\\
3.4797	-0.0549769\\
3.4822	-0.0549887\\
3.4848	-0.0550011\\
3.4874	-0.0550136\\
3.4899	-0.0550259\\
3.4925	-0.0550387\\
3.4951	-0.0550518\\
3.4976	-0.0550645\\
3.5002	-0.0550778\\
3.5028	-0.0550913\\
3.5053	-0.0551044\\
3.5079	-0.0551182\\
3.5104	-0.0551316\\
3.513	-0.0551457\\
3.5156	-0.0551599\\
3.5181	-0.0551737\\
3.5207	-0.0551882\\
3.5233	-0.0552029\\
3.5258	-0.0552171\\
3.5284	-0.055232\\
3.531	-0.055247\\
3.5335	-0.0552616\\
3.5361	-0.0552769\\
3.5386	-0.0552917\\
3.5412	-0.0553073\\
3.5438	-0.055323\\
3.5463	-0.0553381\\
3.5489	-0.055354\\
3.5515	-0.0553701\\
3.554	-0.0553856\\
3.5566	-0.0554019\\
3.5592	-0.0554182\\
3.5617	-0.0554341\\
3.5643	-0.0554507\\
3.5668	-0.0554667\\
3.5694	-0.0554835\\
3.572	-0.0555004\\
3.5745	-0.0555168\\
3.5771	-0.0555339\\
3.5797	-0.0555511\\
3.5822	-0.0555677\\
3.5848	-0.0555851\\
3.5874	-0.0556026\\
3.5899	-0.0556195\\
3.5925	-0.0556371\\
3.595	-0.0556542\\
3.5976	-0.055672\\
3.6002	-0.0556899\\
3.6027	-0.0557072\\
3.6053	-0.0557253\\
3.6079	-0.0557435\\
3.6104	-0.055761\\
3.613	-0.0557793\\
3.6156	-0.0557977\\
3.6181	-0.0558154\\
3.6207	-0.0558339\\
3.6232	-0.0558517\\
3.6258	-0.0558704\\
3.6284	-0.0558891\\
3.6309	-0.0559071\\
3.6335	-0.0559259\\
3.6361	-0.0559448\\
3.6386	-0.055963\\
3.6412	-0.055982\\
3.6438	-0.056001\\
3.6463	-0.0560193\\
3.6489	-0.0560384\\
3.6515	-0.0560576\\
3.654	-0.0560761\\
3.6566	-0.0560953\\
3.6591	-0.0561139\\
3.6617	-0.0561332\\
3.6643	-0.0561526\\
3.6668	-0.0561712\\
3.6694	-0.0561906\\
3.672	-0.0562101\\
3.6745	-0.0562288\\
3.6771	-0.0562484\\
3.6797	-0.0562679\\
3.6822	-0.0562867\\
3.6848	-0.0563063\\
3.6873	-0.0563251\\
3.6899	-0.0563447\\
3.6925	-0.0563643\\
3.695	-0.0563832\\
3.6976	-0.0564029\\
3.7002	-0.0564225\\
3.7027	-0.0564414\\
3.7053	-0.0564611\\
3.7079	-0.0564808\\
3.7104	-0.0564997\\
3.713	-0.0565194\\
3.7155	-0.0565383\\
3.7181	-0.056558\\
3.7207	-0.0565776\\
3.7232	-0.0565965\\
3.7258	-0.0566162\\
3.7284	-0.0566358\\
3.7309	-0.0566547\\
3.7335	-0.0566743\\
3.7361	-0.0566938\\
3.7386	-0.0567127\\
3.7412	-0.0567322\\
3.7437	-0.056751\\
3.7463	-0.0567705\\
3.7489	-0.05679\\
3.7514	-0.0568087\\
3.754	-0.0568281\\
3.7566	-0.0568475\\
3.7591	-0.0568661\\
3.7617	-0.0568854\\
3.7643	-0.0569047\\
3.7668	-0.0569232\\
3.7694	-0.0569424\\
3.7719	-0.0569608\\
3.7745	-0.0569799\\
3.7771	-0.056999\\
3.7796	-0.0570173\\
3.7822	-0.0570362\\
3.7848	-0.0570552\\
3.7873	-0.0570733\\
3.7899	-0.0570921\\
3.7925	-0.0571109\\
3.795	-0.0571289\\
3.7976	-0.0571475\\
3.8002	-0.0571661\\
3.8027	-0.0571839\\
3.8053	-0.0572024\\
3.8078	-0.05722\\
3.8104	-0.0572384\\
3.813	-0.0572566\\
3.8155	-0.0572741\\
3.8181	-0.0572922\\
3.8207	-0.0573103\\
3.8232	-0.0573275\\
3.8258	-0.0573454\\
3.8284	-0.0573632\\
3.8309	-0.0573803\\
3.8335	-0.0573979\\
3.836	-0.0574148\\
3.8386	-0.0574323\\
3.8412	-0.0574497\\
3.8437	-0.0574663\\
3.8463	-0.0574835\\
3.8489	-0.0575007\\
3.8514	-0.057517\\
3.854	-0.057534\\
3.8566	-0.0575508\\
3.8591	-0.0575669\\
3.8617	-0.0575836\\
3.8642	-0.0575995\\
3.8668	-0.057616\\
3.8694	-0.0576323\\
3.8719	-0.0576479\\
3.8745	-0.0576641\\
3.8771	-0.0576801\\
3.8796	-0.0576954\\
3.8822	-0.0577112\\
3.8848	-0.057727\\
3.8873	-0.057742\\
3.8899	-0.0577574\\
3.8924	-0.0577722\\
3.895	-0.0577875\\
3.8976	-0.0578026\\
3.9001	-0.0578171\\
3.9027	-0.057832\\
3.9053	-0.0578468\\
3.9078	-0.0578609\\
3.9104	-0.0578754\\
3.913	-0.0578898\\
3.9155	-0.0579036\\
3.9181	-0.0579177\\
3.9206	-0.0579312\\
3.9232	-0.0579452\\
3.9258	-0.0579589\\
3.9283	-0.0579721\\
3.9309	-0.0579856\\
3.9335	-0.057999\\
3.936	-0.0580118\\
3.9386	-0.0580249\\
3.9412	-0.0580379\\
3.9437	-0.0580502\\
3.9463	-0.058063\\
3.9488	-0.0580751\\
3.9514	-0.0580875\\
3.954	-0.0580998\\
3.9565	-0.0581115\\
3.9591	-0.0581235\\
3.9617	-0.0581354\\
3.9642	-0.0581467\\
3.9668	-0.0581583\\
3.9694	-0.0581697\\
3.9719	-0.0581806\\
3.9745	-0.0581918\\
3.9771	-0.0582028\\
3.9796	-0.0582132\\
3.9822	-0.0582239\\
3.9847	-0.0582341\\
3.9873	-0.0582445\\
3.9899	-0.0582548\\
3.9924	-0.0582645\\
3.995	-0.0582745\\
3.9976	-0.0582843\\
4.0001	-0.0582936\\
4.0027	-0.0583031\\
4.0053	-0.0583124\\
4.0078	-0.0583213\\
4.0104	-0.0583303\\
4.0129	-0.0583389\\
4.0155	-0.0583476\\
4.0181	-0.0583562\\
4.0206	-0.0583643\\
4.0232	-0.0583725\\
4.0258	-0.0583806\\
4.0283	-0.0583883\\
4.0309	-0.0583961\\
4.0335	-0.0584037\\
4.036	-0.0584109\\
4.0386	-0.0584182\\
4.0411	-0.0584251\\
4.0437	-0.0584321\\
4.0463	-0.0584389\\
4.0488	-0.0584453\\
4.0514	-0.0584518\\
4.054	-0.0584582\\
4.0565	-0.0584641\\
4.0591	-0.0584702\\
4.0617	-0.058476\\
4.0642	-0.0584815\\
4.0668	-0.0584871\\
4.0693	-0.0584922\\
4.0719	-0.0584974\\
4.0745	-0.0585025\\
4.077	-0.0585072\\
4.0796	-0.0585119\\
4.0822	-0.0585165\\
4.0847	-0.0585207\\
4.0873	-0.058525\\
4.0899	-0.0585291\\
4.0924	-0.0585328\\
4.095	-0.0585366\\
4.0975	-0.0585401\\
4.1001	-0.0585435\\
4.1027	-0.0585468\\
4.1052	-0.0585498\\
4.1078	-0.0585527\\
4.1104	-0.0585555\\
4.1129	-0.058558\\
4.1155	-0.0585605\\
4.1181	-0.0585628\\
4.1206	-0.0585649\\
4.1232	-0.0585669\\
4.1258	-0.0585687\\
4.1283	-0.0585703\\
4.1309	-0.0585718\\
4.1334	-0.0585731\\
4.136	-0.0585743\\
4.1386	-0.0585753\\
4.1411	-0.0585762\\
4.1437	-0.0585769\\
4.1463	-0.0585775\\
4.1488	-0.0585778\\
4.1514	-0.0585781\\
4.154	-0.0585782\\
4.1565	-0.0585781\\
4.1591	-0.0585779\\
4.1616	-0.0585775\\
4.1642	-0.058577\\
4.1668	-0.0585763\\
4.1693	-0.0585755\\
4.1719	-0.0585746\\
4.1745	-0.0585734\\
4.177	-0.0585722\\
4.1796	-0.0585707\\
4.1822	-0.0585691\\
4.1847	-0.0585675\\
4.1873	-0.0585656\\
4.1898	-0.0585636\\
4.1924	-0.0585615\\
4.195	-0.0585591\\
4.1975	-0.0585567\\
4.2001	-0.0585541\\
4.2027	-0.0585514\\
4.2052	-0.0585486\\
4.2078	-0.0585455\\
4.2104	-0.0585423\\
4.2129	-0.0585391\\
4.2155	-0.0585356\\
4.218	-0.0585322\\
4.2206	-0.0585284\\
4.2232	-0.0585245\\
4.2257	-0.0585207\\
4.2283	-0.0585165\\
4.2309	-0.0585122\\
4.2334	-0.0585079\\
4.236	-0.0585034\\
4.2386	-0.0584987\\
4.2411	-0.058494\\
4.2437	-0.058489\\
4.2462	-0.0584841\\
4.2488	-0.0584789\\
4.2514	-0.0584736\\
4.2539	-0.0584683\\
4.2565	-0.0584627\\
4.2591	-0.0584569\\
4.2616	-0.0584513\\
4.2642	-0.0584453\\
4.2668	-0.0584392\\
4.2693	-0.0584332\\
4.2719	-0.0584269\\
4.2744	-0.0584207\\
4.277	-0.0584141\\
4.2796	-0.0584074\\
4.2821	-0.0584008\\
4.2847	-0.0583939\\
4.2873	-0.0583869\\
4.2898	-0.05838\\
4.2924	-0.0583727\\
4.295	-0.0583653\\
4.2975	-0.0583581\\
4.3001	-0.0583505\\
4.3027	-0.0583428\\
4.3052	-0.0583353\\
4.3078	-0.0583274\\
4.3103	-0.0583197\\
4.3129	-0.0583115\\
4.3155	-0.0583033\\
4.318	-0.0582953\\
4.3206	-0.0582869\\
4.3232	-0.0582784\\
4.3257	-0.0582701\\
4.3283	-0.0582613\\
4.3309	-0.0582525\\
4.3334	-0.058244\\
4.336	-0.058235\\
4.3385	-0.0582262\\
4.3411	-0.058217\\
4.3437	-0.0582078\\
4.3462	-0.0581988\\
4.3488	-0.0581893\\
4.3514	-0.0581798\\
4.3539	-0.0581706\\
4.3565	-0.0581609\\
4.3591	-0.0581511\\
4.3616	-0.0581416\\
4.3642	-0.0581317\\
4.3667	-0.0581221\\
4.3693	-0.058112\\
4.3719	-0.0581018\\
4.3744	-0.058092\\
4.377	-0.0580817\\
4.3796	-0.0580713\\
4.3821	-0.0580612\\
4.3847	-0.0580507\\
4.3873	-0.0580401\\
4.3898	-0.0580299\\
4.3924	-0.0580192\\
4.3949	-0.0580088\\
4.3975	-0.057998\\
4.4001	-0.0579871\\
4.4026	-0.0579765\\
4.4052	-0.0579655\\
4.4078	-0.0579544\\
4.4103	-0.0579437\\
4.4129	-0.0579325\\
4.4155	-0.0579213\\
4.418	-0.0579104\\
4.4206	-0.0578991\\
4.4231	-0.0578881\\
4.4257	-0.0578767\\
4.4283	-0.0578652\\
4.4308	-0.0578541\\
4.4334	-0.0578425\\
4.436	-0.0578309\\
4.4385	-0.0578197\\
4.4411	-0.057808\\
4.4437	-0.0577962\\
4.4462	-0.0577849\\
4.4488	-0.0577731\\
4.4514	-0.0577612\\
4.4539	-0.0577498\\
4.4565	-0.0577378\\
4.459	-0.0577263\\
4.4616	-0.0577143\\
4.4642	-0.0577023\\
4.4667	-0.0576907\\
4.4693	-0.0576787\\
4.4719	-0.0576666\\
4.4744	-0.0576549\\
4.477	-0.0576428\\
4.4796	-0.0576306\\
4.4821	-0.0576189\\
4.4847	-0.0576066\\
4.4872	-0.0575949\\
4.4898	-0.0575826\\
4.4924	-0.0575704\\
4.4949	-0.0575586\\
4.4975	-0.0575463\\
4.5001	-0.057534\\
4.5026	-0.0575221\\
4.5052	-0.0575098\\
4.5078	-0.0574974\\
4.5103	-0.0574856\\
4.5129	-0.0574732\\
4.5154	-0.0574613\\
4.518	-0.057449\\
4.5206	-0.0574366\\
4.5231	-0.0574247\\
4.5257	-0.0574123\\
4.5283	-0.0573999\\
4.5308	-0.057388\\
4.5334	-0.0573757\\
4.536	-0.0573633\\
4.5385	-0.0573514\\
4.5411	-0.0573391\\
4.5436	-0.0573272\\
4.5462	-0.0573148\\
4.5488	-0.0573025\\
4.5513	-0.0572906\\
4.5539	-0.0572783\\
4.5565	-0.057266\\
4.559	-0.0572542\\
4.5616	-0.0572419\\
4.5642	-0.0572296\\
4.5667	-0.0572178\\
4.5693	-0.0572056\\
4.5718	-0.0571939\\
4.5744	-0.0571817\\
4.577	-0.0571695\\
4.5795	-0.0571578\\
4.5821	-0.0571456\\
4.5847	-0.0571335\\
4.5872	-0.0571219\\
4.5898	-0.0571098\\
4.5924	-0.0570978\\
4.5949	-0.0570862\\
4.5975	-0.0570742\\
4.6001	-0.0570623\\
4.6026	-0.0570508\\
4.6052	-0.0570389\\
4.6077	-0.0570275\\
4.6103	-0.0570156\\
4.6129	-0.0570038\\
4.6154	-0.0569925\\
4.618	-0.0569807\\
4.6206	-0.056969\\
4.6231	-0.0569578\\
4.6257	-0.0569461\\
4.6283	-0.0569345\\
4.6308	-0.0569234\\
4.6334	-0.0569119\\
4.6359	-0.0569008\\
4.6385	-0.0568894\\
4.6411	-0.056878\\
4.6436	-0.056867\\
4.6462	-0.0568557\\
4.6488	-0.0568444\\
4.6513	-0.0568336\\
4.6539	-0.0568224\\
4.6565	-0.0568112\\
4.659	-0.0568006\\
4.6616	-0.0567895\\
4.6641	-0.0567789\\
4.6667	-0.056768\\
4.6693	-0.056757\\
4.6718	-0.0567466\\
4.6744	-0.0567358\\
4.677	-0.056725\\
4.6795	-0.0567147\\
4.6821	-0.056704\\
4.6847	-0.0566934\\
4.6872	-0.0566832\\
4.6898	-0.0566727\\
4.6923	-0.0566627\\
4.6949	-0.0566523\\
4.6975	-0.0566419\\
4.7	-0.056632\\
4.7026	-0.0566218\\
4.7052	-0.0566116\\
4.7077	-0.0566019\\
4.7103	-0.0565918\\
4.7129	-0.0565818\\
4.7154	-0.0565722\\
4.718	-0.0565623\\
4.7205	-0.0565529\\
4.7231	-0.0565431\\
4.7257	-0.0565334\\
4.7282	-0.0565241\\
4.7308	-0.0565145\\
4.7334	-0.0565049\\
4.7359	-0.0564958\\
4.7385	-0.0564864\\
4.7411	-0.0564771\\
4.7436	-0.0564681\\
4.7462	-0.0564589\\
4.7487	-0.0564501\\
4.7513	-0.056441\\
4.7539	-0.056432\\
4.7564	-0.0564233\\
4.759	-0.0564144\\
4.7616	-0.0564056\\
4.7641	-0.0563972\\
4.7667	-0.0563885\\
4.7693	-0.0563798\\
4.7718	-0.0563716\\
4.7744	-0.056363\\
4.777	-0.0563546\\
4.7795	-0.0563465\\
4.7821	-0.0563382\\
4.7846	-0.0563303\\
4.7872	-0.0563221\\
4.7898	-0.056314\\
4.7923	-0.0563063\\
4.7949	-0.0562983\\
4.7975	-0.0562904\\
4.8	-0.0562829\\
4.8026	-0.0562751\\
4.8052	-0.0562674\\
4.8077	-0.05626\\
4.8103	-0.0562525\\
4.8128	-0.0562453\\
4.8154	-0.0562378\\
4.818	-0.0562305\\
4.8205	-0.0562235\\
4.8231	-0.0562162\\
4.8257	-0.0562091\\
4.8282	-0.0562023\\
4.8308	-0.0561953\\
4.8334	-0.0561883\\
4.8359	-0.0561817\\
4.8385	-0.0561749\\
4.841	-0.0561684\\
4.8436	-0.0561618\\
4.8462	-0.0561552\\
4.8487	-0.0561489\\
4.8513	-0.0561424\\
4.8539	-0.0561361\\
4.8564	-0.05613\\
4.859	-0.0561237\\
4.8616	-0.0561176\\
4.8641	-0.0561117\\
4.8667	-0.0561056\\
4.8692	-0.0560999\\
4.8718	-0.056094\\
4.8744	-0.0560882\\
4.8769	-0.0560826\\
4.8795	-0.0560769\\
4.8821	-0.0560713\\
4.8846	-0.056066\\
4.8872	-0.0560605\\
4.8898	-0.0560551\\
4.8923	-0.0560499\\
4.8949	-0.0560447\\
4.8974	-0.0560397\\
4.9	-0.0560345\\
4.9026	-0.0560294\\
4.9051	-0.0560246\\
4.9077	-0.0560197\\
4.9103	-0.0560148\\
4.9128	-0.0560102\\
4.9154	-0.0560055\\
4.918	-0.0560008\\
4.9205	-0.0559964\\
4.9231	-0.0559918\\
4.9257	-0.0559874\\
4.9282	-0.0559832\\
4.9308	-0.0559788\\
4.9333	-0.0559747\\
4.9359	-0.0559705\\
4.9385	-0.0559664\\
4.941	-0.0559625\\
4.9436	-0.0559585\\
4.9462	-0.0559545\\
4.9487	-0.0559508\\
4.9513	-0.055947\\
4.9539	-0.0559432\\
4.9564	-0.0559397\\
4.959	-0.0559361\\
4.9615	-0.0559327\\
4.9641	-0.0559292\\
4.9667	-0.0559258\\
4.9692	-0.0559225\\
4.9718	-0.0559192\\
4.9744	-0.055916\\
4.9769	-0.0559129\\
4.9795	-0.0559098\\
4.9821	-0.0559067\\
4.9846	-0.0559039\\
4.9872	-0.0559009\\
4.9897	-0.0558982\\
4.9923	-0.0558954\\
4.9949	-0.0558926\\
4.9974	-0.05589\\
5	-0.0558874\\
};
\addplot [color=mycolor20, line width=1.0pt, forget plot]
  table[row sep=crcr]{%
-0.125	2.42501e-08\\
-0.1224	2.47509e-08\\
-0.1199	2.52324e-08\\
-0.1173	2.57329e-08\\
-0.1147	2.62337e-08\\
-0.1122	2.67151e-08\\
-0.1096	2.72153e-08\\
-0.1071	2.76965e-08\\
-0.1045	2.8197e-08\\
-0.1019	2.86971e-08\\
-0.0994	2.91782e-08\\
-0.0968	2.96785e-08\\
-0.0942	3.01785e-08\\
-0.0917	3.06594e-08\\
-0.0891	3.11596e-08\\
-0.0865	3.16595e-08\\
-0.084	3.21403e-08\\
-0.0814	3.26404e-08\\
-0.0789	3.31208e-08\\
-0.0763	3.36206e-08\\
-0.0737	3.41204e-08\\
-0.0712	3.46009e-08\\
-0.0686	3.51004e-08\\
-0.066	3.56001e-08\\
-0.0635	3.60803e-08\\
-0.0609	3.65798e-08\\
-0.0583	3.70792e-08\\
-0.0558	3.75594e-08\\
-0.0532	3.80588e-08\\
-0.0507	3.85388e-08\\
-0.0481	3.90379e-08\\
-0.0455	3.95372e-08\\
-0.043	4.00171e-08\\
-0.0404	4.05162e-08\\
-0.0378	4.10153e-08\\
-0.0353	4.14949e-08\\
-0.0327	4.19937e-08\\
-0.0301	4.24926e-08\\
-0.0276	4.29721e-08\\
-0.025	4.34708e-08\\
-0.0224	4.39696e-08\\
-0.0199	4.44489e-08\\
-0.0173	4.49473e-08\\
-0.0148	4.54268e-08\\
-0.0122	4.59252e-08\\
-0.0096	4.64236e-08\\
-0.0071	4.69029e-08\\
-0.0045	4.74012e-08\\
-0.0019	4.78992e-08\\
0.0006	4.83783e-08\\
0.0032	0.0138662\\
0.0058	0.0232431\\
0.0083	0.030277\\
0.0109	0.0366949\\
0.0134	0.042232\\
0.016	0.0471869\\
0.0186	0.0512462\\
0.0211	0.0544699\\
0.0237	0.0574133\\
0.0263	0.0601546\\
0.0288	0.0626415\\
0.0314	0.064999\\
0.034	0.0670343\\
0.0365	0.068675\\
0.0391	0.0701156\\
0.0416	0.0713175\\
0.0442	0.0724226\\
0.0468	0.0733741\\
0.0493	0.0741146\\
0.0519	0.0746857\\
0.0545	0.0750667\\
0.057	0.0752872\\
0.0596	0.0753994\\
0.0622	0.0754062\\
0.0647	0.0753056\\
0.0673	0.0750761\\
0.0698	0.0747346\\
0.0724	0.0742701\\
0.075	0.0737211\\
0.0775	0.073132\\
0.0801	0.0724571\\
0.0827	0.071709\\
0.0852	0.0709134\\
0.0878	0.0700121\\
0.0904	0.0690545\\
0.0929	0.0680991\\
0.0955	0.0670774\\
0.098	0.0660624\\
0.1006	0.0649625\\
0.1032	0.0638144\\
0.1057	0.0626745\\
0.1083	0.0614687\\
0.1109	0.0602541\\
0.1134	0.0590771\\
0.116	0.0578335\\
0.1186	0.0565609\\
0.1211	0.055312\\
0.1237	0.0539991\\
0.1263	0.0526853\\
0.1288	0.0514251\\
0.1314	0.0501103\\
0.1339	0.0488316\\
0.1365	0.047482\\
0.1391	0.0461186\\
0.1416	0.0448063\\
0.1442	0.0434481\\
0.1468	0.0420939\\
0.1493	0.0407856\\
0.1519	0.0394099\\
0.1545	0.0380191\\
0.157	0.036676\\
0.1596	0.0352833\\
0.1621	0.0339502\\
0.1647	0.0325631\\
0.1673	0.0311649\\
0.1698	0.0298055\\
0.1724	0.0283808\\
0.175	0.0269547\\
0.1775	0.0255883\\
0.1801	0.0241698\\
0.1827	0.0227447\\
0.1852	0.0213609\\
0.1878	0.019908\\
0.1903	0.0185048\\
0.1929	0.0170475\\
0.1955	0.0155945\\
0.198	0.014196\\
0.2006	0.0127318\\
0.2032	0.011254\\
0.2057	0.00982428\\
0.2083	0.00833678\\
0.2109	0.00685447\\
0.2134	0.00543259\\
0.216	0.00395041\\
0.2185	0.00251638\\
0.2211	0.00101634\\
0.2237	-0.00048551\\
0.2262	-0.00192393\\
0.2288	-0.00341155\\
0.2314	-0.0048949\\
0.2339	-0.00632329\\
0.2365	-0.00781347\\
0.2391	-0.00930418\\
0.2416	-0.0107306\\
0.2442	-0.0122015\\
0.2467	-0.0136045\\
0.2493	-0.0150569\\
0.2519	-0.0165071\\
0.2544	-0.0178986\\
0.257	-0.0193364\\
0.2596	-0.0207578\\
0.2621	-0.0221065\\
0.2647	-0.0234937\\
0.2673	-0.02487\\
0.2698	-0.0261848\\
0.2724	-0.0275398\\
0.2749	-0.0288242\\
0.2775	-0.0301361\\
0.2801	-0.0314239\\
0.2826	-0.0326436\\
0.2852	-0.0338963\\
0.2878	-0.035132\\
0.2903	-0.0362993\\
0.2929	-0.0374858\\
0.2955	-0.038642\\
0.298	-0.0397277\\
0.3006	-0.0408339\\
0.3032	-0.0419185\\
0.3057	-0.0429382\\
0.3083	-0.0439691\\
0.3108	-0.0449283\\
0.3134	-0.045893\\
0.316	-0.0468278\\
0.3185	-0.0477019\\
0.3211	-0.0485847\\
0.3237	-0.0494367\\
0.3262	-0.0502224\\
0.3288	-0.0510033\\
0.3314	-0.0517498\\
0.3339	-0.0524387\\
0.3365	-0.0531271\\
0.339	-0.05376\\
0.3416	-0.054384\\
0.3442	-0.0549701\\
0.3467	-0.0554987\\
0.3493	-0.0560154\\
0.3519	-0.0565019\\
0.3544	-0.0569411\\
0.357	-0.0573653\\
0.3596	-0.057753\\
0.3621	-0.0580903\\
0.3647	-0.0584067\\
0.3672	-0.0586812\\
0.3698	-0.0589378\\
0.3724	-0.0591636\\
0.3749	-0.0593488\\
0.3775	-0.0595061\\
0.3801	-0.0596284\\
0.3826	-0.059716\\
0.3852	-0.0597789\\
0.3878	-0.0598136\\
0.3903	-0.0598184\\
0.3929	-0.059791\\
0.3954	-0.0597331\\
0.398	-0.0596423\\
0.4006	-0.0595236\\
0.4031	-0.0593851\\
0.4057	-0.0592149\\
0.4083	-0.0590157\\
0.4108	-0.0587954\\
0.4134	-0.0585372\\
0.416	-0.0582523\\
0.4185	-0.0579559\\
0.4211	-0.0576248\\
0.4236	-0.0572831\\
0.4262	-0.0569014\\
0.4288	-0.0564927\\
0.4313	-0.0560761\\
0.4339	-0.0556211\\
0.4365	-0.0551458\\
0.439	-0.0546691\\
0.4416	-0.0541512\\
0.4442	-0.0536095\\
0.4467	-0.0530671\\
0.4493	-0.052483\\
0.4519	-0.051881\\
0.4544	-0.0512858\\
0.457	-0.0506488\\
0.4595	-0.0500174\\
0.4621	-0.049341\\
0.4647	-0.0486461\\
0.4672	-0.047963\\
0.4698	-0.0472385\\
0.4724	-0.0464994\\
0.4749	-0.0457737\\
0.4775	-0.0450023\\
0.4801	-0.0442148\\
0.4826	-0.0434444\\
0.4852	-0.0426316\\
0.4877	-0.0418394\\
0.4903	-0.0410037\\
0.4929	-0.0401547\\
0.4954	-0.0393256\\
0.498	-0.0384517\\
0.5006	-0.0375679\\
0.5031	-0.0367103\\
0.5057	-0.0358099\\
0.5083	-0.0349\\
0.5108	-0.0340153\\
0.5134	-0.0330857\\
0.5159	-0.0321846\\
0.5185	-0.0312415\\
0.5211	-0.030293\\
0.5236	-0.0293753\\
0.5262	-0.0284141\\
0.5288	-0.0274459\\
0.5313	-0.0265095\\
0.5339	-0.0255318\\
0.5365	-0.0245515\\
0.539	-0.0236064\\
0.5416	-0.02262\\
0.5441	-0.0216677\\
0.5467	-0.0206739\\
0.5493	-0.0196782\\
0.5518	-0.0187205\\
0.5544	-0.0177248\\
0.557	-0.0167289\\
0.5595	-0.0157704\\
0.5621	-0.0147726\\
0.5647	-0.013775\\
0.5672	-0.0128174\\
0.5698	-0.0118243\\
0.5723	-0.0108721\\
0.5749	-0.00988394\\
0.5775	-0.00889757\\
0.58	-0.00795132\\
0.5826	-0.00697081\\
0.5852	-0.00599528\\
0.5877	-0.0050625\\
0.5903	-0.00409747\\
0.5929	-0.00313701\\
0.5954	-0.00221783\\
0.598	-0.0012673\\
0.6006	-0.000323634\\
0.6031	0.000576393\\
0.6057	0.00150475\\
0.6082	0.00239053\\
0.6108	0.00330492\\
0.6134	0.00421194\\
0.6159	0.00507604\\
0.6185	0.00596527\\
0.6211	0.00684442\\
0.6236	0.00768056\\
0.6262	0.00854106\\
0.6288	0.00939229\\
0.6313	0.0102013\\
0.6339	0.0110318\\
0.6364	0.0118192\\
0.639	0.0126264\\
0.6416	0.0134221\\
0.6441	0.0141768\\
0.6467	0.0149505\\
0.6493	0.0157119\\
0.6518	0.0164314\\
0.6544	0.0171664\\
0.657	0.0178882\\
0.6595	0.0185701\\
0.6621	0.0192667\\
0.6646	0.019924\\
0.6672	0.0205936\\
0.6698	0.0212483\\
0.6723	0.021864\\
0.6749	0.0224902\\
0.6775	0.0231024\\
0.68	0.0236776\\
0.6826	0.0242612\\
0.6852	0.0248291\\
0.6877	0.0253603\\
0.6903	0.0258974\\
0.6928	0.0263997\\
0.6954	0.0269073\\
0.698	0.0273996\\
0.7005	0.0278577\\
0.7031	0.0283179\\
0.7057	0.0287618\\
0.7082	0.0291735\\
0.7108	0.0295864\\
0.7134	0.0299835\\
0.7159	0.03035\\
0.7185	0.0307147\\
0.721	0.0310495\\
0.7236	0.0313814\\
0.7262	0.0316972\\
0.7287	0.0319858\\
0.7313	0.03227\\
0.7339	0.0325375\\
0.7364	0.0327787\\
0.739	0.033013\\
0.7416	0.0332309\\
0.7441	0.0334253\\
0.7467	0.0336117\\
0.7492	0.0337754\\
0.7518	0.0339293\\
0.7544	0.0340663\\
0.7569	0.0341825\\
0.7595	0.0342876\\
0.7621	0.034377\\
0.7646	0.0344478\\
0.7672	0.0345054\\
0.7698	0.0345466\\
0.7723	0.0345708\\
0.7749	0.0345804\\
0.7775	0.0345746\\
0.78	0.0345544\\
0.7826	0.034518\\
0.7851	0.0344681\\
0.7877	0.0344006\\
0.7903	0.0343175\\
0.7928	0.0342235\\
0.7954	0.0341112\\
0.798	0.0339839\\
0.8005	0.0338473\\
0.8031	0.0336902\\
0.8057	0.0335183\\
0.8082	0.0333392\\
0.8108	0.033139\\
0.8133	0.0329332\\
0.8159	0.0327053\\
0.8185	0.032463\\
0.821	0.0322165\\
0.8236	0.0319466\\
0.8262	0.0316631\\
0.8287	0.0313781\\
0.8313	0.0310686\\
0.8339	0.0307458\\
0.8364	0.0304228\\
0.839	0.0300739\\
0.8415	0.0297265\\
0.8441	0.029353\\
0.8467	0.0289672\\
0.8492	0.0285847\\
0.8518	0.0281747\\
0.8544	0.0277524\\
0.8569	0.0273352\\
0.8595	0.0268899\\
0.8621	0.0264335\\
0.8646	0.025984\\
0.8672	0.0255054\\
0.8697	0.0250346\\
0.8723	0.0245343\\
0.8749	0.0240233\\
0.8774	0.0235223\\
0.88	0.0229913\\
0.8826	0.0224502\\
0.8851	0.0219204\\
0.8877	0.0213595\\
0.8903	0.020789\\
0.8928	0.0202317\\
0.8954	0.0196432\\
0.8979	0.019069\\
0.9005	0.0184629\\
0.9031	0.017848\\
0.9056	0.0172486\\
0.9082	0.016617\\
0.9108	0.0159775\\
0.9133	0.0153552\\
0.9159	0.0147004\\
0.9185	0.0140379\\
0.921	0.0133937\\
0.9236	0.0127166\\
0.9262	0.0120326\\
0.9287	0.0113686\\
0.9313	0.0106716\\
0.9338	0.00999522\\
0.9364	0.00928545\\
0.939	0.00856947\\
0.9415	0.00787548\\
0.9441	0.00714826\\
0.9467	0.00641563\\
0.9492	0.00570614\\
0.9518	0.00496308\\
0.9544	0.00421489\\
0.9569	0.00349091\\
0.9595	0.00273356\\
0.962	0.0020013\\
0.9646	0.00123567\\
0.9672	0.00046592\\
0.9697	-0.000278003\\
0.9723	-0.00105539\\
0.9749	-0.00183626\\
0.9774	-0.00259009\\
0.98	-0.00337703\\
0.9826	-0.00416692\\
0.9851	-0.00492913\\
0.9877	-0.0057245\\
0.9902	-0.00649158\\
0.9928	-0.00729145\\
0.9954	-0.00809324\\
0.9979	-0.00886589\\
1.0005	-0.00967116\\
1.0031	-0.0104781\\
1.0056	-0.0112553\\
1.0082	-0.0120647\\
1.0108	-0.012875\\
1.0133	-0.0136548\\
1.0159	-0.0144664\\
1.0184	-0.0152474\\
1.021	-0.01606\\
1.0236	-0.0168727\\
1.0261	-0.0176542\\
1.0287	-0.0184665\\
1.0313	-0.0192785\\
1.0338	-0.0200588\\
1.0364	-0.0208697\\
1.039	-0.0216798\\
1.0415	-0.0224579\\
1.0441	-0.0232658\\
1.0466	-0.0240413\\
1.0492	-0.0248464\\
1.0518	-0.0256499\\
1.0543	-0.026421\\
1.0569	-0.027221\\
1.0595	-0.0280189\\
1.062	-0.0287839\\
1.0646	-0.0295771\\
1.0672	-0.0303679\\
1.0697	-0.0311259\\
1.0723	-0.0319114\\
1.0748	-0.032664\\
1.0774	-0.0334436\\
1.08	-0.03422\\
1.0825	-0.0349633\\
1.0851	-0.0357331\\
1.0877	-0.0364993\\
1.0902	-0.0372325\\
1.0928	-0.0379913\\
1.0954	-0.038746\\
1.0979	-0.0394678\\
1.1005	-0.0402144\\
1.1031	-0.0409568\\
1.1056	-0.0416664\\
1.1082	-0.0423999\\
1.1107	-0.0431008\\
1.1133	-0.043825\\
1.1159	-0.0445442\\
1.1184	-0.0452312\\
1.121	-0.0459407\\
1.1236	-0.046645\\
1.1261	-0.0473173\\
1.1287	-0.048011\\
1.1313	-0.0486993\\
1.1338	-0.0493558\\
1.1364	-0.0500331\\
1.1389	-0.050679\\
1.1415	-0.051345\\
1.1441	-0.0520049\\
1.1466	-0.0526338\\
1.1492	-0.0532819\\
1.1518	-0.0539238\\
1.1543	-0.0545353\\
1.1569	-0.0551649\\
1.1595	-0.0557882\\
1.162	-0.0563814\\
1.1646	-0.0569918\\
1.1671	-0.0575726\\
1.1697	-0.0581701\\
1.1723	-0.058761\\
1.1748	-0.0593227\\
1.1774	-0.0599001\\
1.18	-0.0604706\\
1.1825	-0.0610126\\
1.1851	-0.0615695\\
1.1877	-0.0621193\\
1.1902	-0.0626414\\
1.1928	-0.0631772\\
1.1953	-0.0636857\\
1.1979	-0.0642074\\
1.2005	-0.0647218\\
1.203	-0.0652096\\
1.2056	-0.0657098\\
1.2082	-0.0662025\\
1.2107	-0.0666693\\
1.2133	-0.0671475\\
1.2159	-0.0676182\\
1.2184	-0.0680637\\
1.221	-0.0685197\\
1.2235	-0.0689512\\
1.2261	-0.0693924\\
1.2287	-0.069826\\
1.2312	-0.0702358\\
1.2338	-0.0706546\\
1.2364	-0.0710657\\
1.2389	-0.0714539\\
1.2415	-0.0718501\\
1.2441	-0.0722386\\
1.2466	-0.072605\\
1.2492	-0.0729785\\
1.2518	-0.0733444\\
1.2543	-0.073689\\
1.2569	-0.07404\\
1.2594	-0.0743702\\
1.262	-0.0747061\\
1.2646	-0.0750343\\
1.2671	-0.0753428\\
1.2697	-0.0756561\\
1.2723	-0.0759619\\
1.2748	-0.0762488\\
1.2774	-0.0765397\\
1.28	-0.076823\\
1.2825	-0.0770883\\
1.2851	-0.0773569\\
1.2876	-0.0776081\\
1.2902	-0.077862\\
1.2928	-0.0781085\\
1.2953	-0.0783385\\
1.2979	-0.0785705\\
1.3005	-0.0787951\\
1.303	-0.0790042\\
1.3056	-0.0792145\\
1.3082	-0.0794175\\
1.3107	-0.079606\\
1.3133	-0.0797949\\
1.3158	-0.0799698\\
1.3184	-0.0801447\\
1.321	-0.0803126\\
1.3235	-0.0804674\\
1.3261	-0.0806216\\
1.3287	-0.0807688\\
1.3312	-0.0809039\\
1.3338	-0.0810377\\
1.3364	-0.0811647\\
1.3389	-0.0812804\\
1.3415	-0.0813943\\
1.344	-0.0814975\\
1.3466	-0.0815984\\
1.3492	-0.0816927\\
1.3517	-0.0817774\\
1.3543	-0.0818591\\
1.3569	-0.0819345\\
1.3594	-0.0820011\\
1.362	-0.0820642\\
1.3646	-0.0821211\\
1.3671	-0.08217\\
1.3697	-0.0822151\\
1.3722	-0.0822527\\
1.3748	-0.0822861\\
1.3774	-0.0823135\\
1.3799	-0.0823345\\
1.3825	-0.0823507\\
1.3851	-0.0823612\\
1.3876	-0.0823661\\
1.3902	-0.0823658\\
1.3928	-0.08236\\
1.3953	-0.0823494\\
1.3979	-0.0823331\\
1.4005	-0.0823117\\
1.403	-0.0822861\\
1.4056	-0.0822546\\
1.4081	-0.0822196\\
1.4107	-0.0821783\\
1.4133	-0.0821322\\
1.4158	-0.0820833\\
1.4184	-0.0820278\\
1.421	-0.0819677\\
1.4235	-0.0819056\\
1.4261	-0.0818367\\
1.4287	-0.0817633\\
1.4312	-0.0816886\\
1.4338	-0.0816068\\
1.4363	-0.0815242\\
1.4389	-0.0814343\\
1.4415	-0.0813404\\
1.444	-0.0812463\\
1.4466	-0.0811447\\
1.4492	-0.0810393\\
1.4517	-0.0809345\\
1.4543	-0.0808219\\
1.4569	-0.0807058\\
1.4594	-0.0805909\\
1.462	-0.080468\\
1.4645	-0.0803468\\
1.4671	-0.0802175\\
1.4697	-0.0800851\\
1.4722	-0.0799548\\
1.4748	-0.0798164\\
1.4774	-0.0796751\\
1.4799	-0.0795365\\
1.4825	-0.0793897\\
1.4851	-0.0792402\\
1.4876	-0.079094\\
1.4902	-0.0789394\\
1.4927	-0.0787885\\
1.4953	-0.0786292\\
1.4979	-0.0784676\\
1.5004	-0.0783102\\
1.503	-0.0781443\\
1.5056	-0.0779763\\
1.5081	-0.077813\\
1.5107	-0.0776412\\
1.5133	-0.0774675\\
1.5158	-0.0772989\\
1.5184	-0.0771219\\
1.5209	-0.0769501\\
1.5235	-0.0767699\\
1.5261	-0.0765883\\
1.5286	-0.0764123\\
1.5312	-0.076228\\
1.5338	-0.0760424\\
1.5363	-0.0758628\\
1.5389	-0.0756749\\
1.5415	-0.075486\\
1.544	-0.0753034\\
1.5466	-0.0751127\\
1.5491	-0.0749284\\
1.5517	-0.074736\\
1.5543	-0.0745429\\
1.5568	-0.0743566\\
1.5594	-0.0741623\\
1.562	-0.0739675\\
1.5645	-0.0737797\\
1.5671	-0.0735841\\
1.5697	-0.0733881\\
1.5722	-0.0731994\\
1.5748	-0.0730029\\
1.5774	-0.0728063\\
1.5799	-0.0726172\\
1.5825	-0.0724205\\
1.585	-0.0722314\\
1.5876	-0.0720349\\
1.5902	-0.0718384\\
1.5927	-0.0716498\\
1.5953	-0.0714538\\
1.5979	-0.0712581\\
1.6004	-0.0710703\\
1.603	-0.0708755\\
1.6056	-0.0706811\\
1.6081	-0.0704946\\
1.6107	-0.0703013\\
1.6132	-0.070116\\
1.6158	-0.069924\\
1.6184	-0.0697327\\
1.6209	-0.0695495\\
1.6235	-0.0693597\\
1.6261	-0.0691708\\
1.6286	-0.0689901\\
1.6312	-0.068803\\
1.6338	-0.068617\\
1.6363	-0.0684391\\
1.6389	-0.0682551\\
1.6414	-0.0680793\\
1.644	-0.0678976\\
1.6466	-0.0677171\\
1.6491	-0.0675447\\
1.6517	-0.0673666\\
1.6543	-0.0671899\\
1.6568	-0.0670213\\
1.6594	-0.0668473\\
1.662	-0.0666747\\
1.6645	-0.0665101\\
1.6671	-0.0663404\\
1.6696	-0.0661786\\
1.6722	-0.0660119\\
1.6748	-0.0658468\\
1.6773	-0.0656896\\
1.6799	-0.0655277\\
1.6825	-0.0653674\\
1.685	-0.0652149\\
1.6876	-0.065058\\
1.6902	-0.0649029\\
1.6927	-0.0647553\\
1.6953	-0.0646037\\
1.6978	-0.0644596\\
1.7004	-0.0643115\\
1.703	-0.0641653\\
1.7055	-0.0640264\\
1.7081	-0.0638839\\
1.7107	-0.0637433\\
1.7132	-0.06361\\
1.7158	-0.0634732\\
1.7184	-0.0633383\\
1.7209	-0.0632106\\
1.7235	-0.0630797\\
1.7261	-0.0629507\\
1.7286	-0.0628287\\
1.7312	-0.0627038\\
1.7337	-0.0625856\\
1.7363	-0.0624647\\
1.7389	-0.0623459\\
1.7414	-0.0622336\\
1.744	-0.0621188\\
1.7466	-0.0620061\\
1.7491	-0.0618998\\
1.7517	-0.0617913\\
1.7543	-0.0616849\\
1.7568	-0.0615845\\
1.7594	-0.0614823\\
1.7619	-0.0613859\\
1.7645	-0.0612878\\
1.7671	-0.0611919\\
1.7696	-0.0611016\\
1.7722	-0.0610098\\
1.7748	-0.0609202\\
1.7773	-0.060836\\
1.7799	-0.0607505\\
1.7825	-0.0606672\\
1.785	-0.060589\\
1.7876	-0.0605099\\
1.7901	-0.0604357\\
1.7927	-0.0603607\\
1.7953	-0.0602878\\
1.7978	-0.0602196\\
1.8004	-0.0601508\\
1.803	-0.0600841\\
1.8055	-0.0600219\\
1.8081	-0.0599593\\
1.8107	-0.0598987\\
1.8132	-0.0598424\\
1.8158	-0.0597859\\
1.8183	-0.0597334\\
1.8209	-0.0596809\\
1.8235	-0.0596303\\
1.826	-0.0595836\\
1.8286	-0.059537\\
1.8312	-0.0594924\\
1.8337	-0.0594513\\
1.8363	-0.0594105\\
1.8389	-0.0593716\\
1.8414	-0.0593361\\
1.844	-0.059301\\
1.8465	-0.059269\\
1.8491	-0.0592377\\
1.8517	-0.0592081\\
1.8542	-0.0591815\\
1.8568	-0.0591556\\
1.8594	-0.0591314\\
1.8619	-0.05911\\
1.8645	-0.0590894\\
1.8671	-0.0590705\\
1.8696	-0.059054\\
1.8722	-0.0590386\\
1.8747	-0.0590253\\
1.8773	-0.0590132\\
1.8799	-0.0590028\\
1.8824	-0.0589942\\
1.885	-0.058987\\
1.8876	-0.0589813\\
1.8901	-0.0589774\\
1.8927	-0.0589748\\
1.8953	-0.0589738\\
1.8978	-0.0589742\\
1.9004	-0.0589761\\
1.903	-0.0589795\\
1.9055	-0.0589841\\
1.9081	-0.0589903\\
1.9106	-0.0589976\\
1.9132	-0.0590065\\
1.9158	-0.0590168\\
1.9183	-0.059028\\
1.9209	-0.0590409\\
1.9235	-0.059055\\
1.926	-0.0590699\\
1.9286	-0.0590865\\
1.9312	-0.0591043\\
1.9337	-0.0591226\\
1.9363	-0.0591428\\
1.9388	-0.0591632\\
1.9414	-0.0591856\\
1.944	-0.0592091\\
1.9465	-0.0592326\\
1.9491	-0.0592582\\
1.9517	-0.0592847\\
1.9542	-0.0593112\\
1.9568	-0.0593397\\
1.9594	-0.0593691\\
1.9619	-0.0593983\\
1.9645	-0.0594295\\
1.967	-0.0594603\\
1.9696	-0.0594932\\
1.9722	-0.0595269\\
1.9747	-0.05956\\
1.9773	-0.0595953\\
1.9799	-0.0596313\\
1.9824	-0.0596666\\
1.985	-0.0597039\\
1.9876	-0.059742\\
1.9901	-0.0597792\\
1.9927	-0.0598185\\
1.9952	-0.0598568\\
1.9978	-0.0598973\\
2.0004	-0.0599383\\
2.0029	-0.0599782\\
2.0055	-0.0600202\\
2.0081	-0.0600627\\
2.0106	-0.060104\\
2.0132	-0.0601473\\
2.0158	-0.0601911\\
2.0183	-0.0602335\\
2.0209	-0.060278\\
2.0234	-0.0603211\\
2.026	-0.0603663\\
2.0286	-0.0604117\\
2.0311	-0.0604556\\
2.0337	-0.0605015\\
2.0363	-0.0605477\\
2.0388	-0.0605923\\
2.0414	-0.0606388\\
2.044	-0.0606855\\
2.0465	-0.0607305\\
2.0491	-0.0607775\\
2.0517	-0.0608245\\
2.0542	-0.0608698\\
2.0568	-0.060917\\
2.0593	-0.0609624\\
2.0619	-0.0610096\\
2.0645	-0.0610568\\
2.067	-0.0611022\\
2.0696	-0.0611493\\
2.0722	-0.0611964\\
2.0747	-0.0612416\\
2.0773	-0.0612886\\
2.0799	-0.0613354\\
2.0824	-0.0613802\\
2.085	-0.0614267\\
2.0875	-0.0614713\\
2.0901	-0.0615175\\
2.0927	-0.0615634\\
2.0952	-0.0616074\\
2.0978	-0.0616529\\
2.1004	-0.0616982\\
2.1029	-0.0617415\\
2.1055	-0.0617862\\
2.1081	-0.0618306\\
2.1106	-0.0618731\\
2.1132	-0.0619169\\
2.1157	-0.0619587\\
2.1183	-0.0620018\\
2.1209	-0.0620445\\
2.1234	-0.0620853\\
2.126	-0.0621273\\
2.1286	-0.0621689\\
2.1311	-0.0622085\\
2.1337	-0.0622492\\
2.1363	-0.0622895\\
2.1388	-0.0623279\\
2.1414	-0.0623673\\
2.1439	-0.0624048\\
2.1465	-0.0624432\\
2.1491	-0.0624812\\
2.1516	-0.0625173\\
2.1542	-0.0625543\\
2.1568	-0.0625907\\
2.1593	-0.0626253\\
2.1619	-0.0626607\\
2.1645	-0.0626956\\
2.167	-0.0627286\\
2.1696	-0.0627624\\
2.1721	-0.0627944\\
2.1747	-0.062827\\
2.1773	-0.0628591\\
2.1798	-0.0628894\\
2.1824	-0.0629204\\
2.185	-0.0629507\\
2.1875	-0.0629793\\
2.1901	-0.0630085\\
2.1927	-0.063037\\
2.1952	-0.0630639\\
2.1978	-0.0630912\\
2.2004	-0.0631179\\
2.2029	-0.063143\\
2.2055	-0.0631685\\
2.208	-0.0631924\\
2.2106	-0.0632167\\
2.2132	-0.0632403\\
2.2157	-0.0632624\\
2.2183	-0.0632848\\
2.2209	-0.0633065\\
2.2234	-0.0633268\\
2.226	-0.0633473\\
2.2286	-0.0633671\\
2.2311	-0.0633856\\
2.2337	-0.0634042\\
2.2362	-0.0634215\\
2.2388	-0.0634389\\
2.2414	-0.0634556\\
2.2439	-0.0634711\\
2.2465	-0.0634866\\
2.2491	-0.0635014\\
2.2516	-0.0635151\\
2.2542	-0.0635287\\
2.2568	-0.0635417\\
2.2593	-0.0635537\\
2.2619	-0.0635655\\
2.2644	-0.0635763\\
2.267	-0.0635869\\
2.2696	-0.0635969\\
2.2721	-0.063606\\
2.2747	-0.0636148\\
2.2773	-0.0636231\\
2.2798	-0.0636305\\
2.2824	-0.0636376\\
2.285	-0.0636442\\
2.2875	-0.0636499\\
2.2901	-0.0636554\\
2.2926	-0.0636601\\
2.2952	-0.0636645\\
2.2978	-0.0636683\\
2.3003	-0.0636715\\
2.3029	-0.0636743\\
2.3055	-0.0636765\\
2.308	-0.0636782\\
2.3106	-0.0636795\\
2.3132	-0.0636803\\
2.3157	-0.0636805\\
2.3183	-0.0636804\\
2.3208	-0.0636797\\
2.3234	-0.0636786\\
2.326	-0.0636771\\
2.3285	-0.0636751\\
2.3311	-0.0636727\\
2.3337	-0.0636698\\
2.3362	-0.0636666\\
2.3388	-0.0636629\\
2.3414	-0.0636588\\
2.3439	-0.0636545\\
2.3465	-0.0636496\\
2.349	-0.0636446\\
2.3516	-0.063639\\
2.3542	-0.063633\\
2.3567	-0.0636269\\
2.3593	-0.0636203\\
2.3619	-0.0636133\\
2.3644	-0.0636062\\
2.367	-0.0635986\\
2.3696	-0.0635907\\
2.3721	-0.0635828\\
2.3747	-0.0635743\\
2.3773	-0.0635656\\
2.3798	-0.0635569\\
2.3824	-0.0635477\\
2.3849	-0.0635385\\
2.3875	-0.0635288\\
2.3901	-0.0635189\\
2.3926	-0.0635091\\
2.3952	-0.0634987\\
2.3978	-0.0634882\\
2.4003	-0.0634779\\
2.4029	-0.0634669\\
2.4055	-0.0634559\\
2.408	-0.0634451\\
2.4106	-0.0634337\\
2.4131	-0.0634227\\
2.4157	-0.0634111\\
2.4183	-0.0633994\\
2.4208	-0.063388\\
2.4234	-0.0633761\\
2.426	-0.0633641\\
2.4285	-0.0633525\\
2.4311	-0.0633404\\
2.4337	-0.0633282\\
2.4362	-0.0633164\\
2.4388	-0.0633042\\
2.4413	-0.0632924\\
2.4439	-0.0632801\\
2.4465	-0.0632678\\
2.449	-0.0632559\\
2.4516	-0.0632436\\
2.4542	-0.0632314\\
2.4567	-0.0632196\\
2.4593	-0.0632074\\
2.4619	-0.0631952\\
2.4644	-0.0631836\\
2.467	-0.0631716\\
2.4695	-0.0631601\\
2.4721	-0.0631482\\
2.4747	-0.0631364\\
2.4772	-0.0631252\\
2.4798	-0.0631136\\
2.4824	-0.0631021\\
2.4849	-0.0630912\\
2.4875	-0.06308\\
2.4901	-0.0630689\\
2.4926	-0.0630584\\
2.4952	-0.0630476\\
2.4977	-0.0630374\\
2.5003	-0.0630269\\
2.5029	-0.0630166\\
2.5054	-0.0630069\\
2.508	-0.062997\\
2.5106	-0.0629872\\
2.5131	-0.0629781\\
2.5157	-0.0629687\\
2.5183	-0.0629596\\
2.5208	-0.062951\\
2.5234	-0.0629424\\
2.526	-0.0629339\\
2.5285	-0.062926\\
2.5311	-0.062918\\
2.5336	-0.0629106\\
2.5362	-0.0629031\\
2.5388	-0.0628959\\
2.5413	-0.0628892\\
2.5439	-0.0628825\\
2.5465	-0.062876\\
2.549	-0.0628701\\
2.5516	-0.0628642\\
2.5542	-0.0628586\\
2.5567	-0.0628535\\
2.5593	-0.0628484\\
2.5618	-0.0628439\\
2.5644	-0.0628394\\
2.567	-0.0628353\\
2.5695	-0.0628315\\
2.5721	-0.062828\\
2.5747	-0.0628247\\
2.5772	-0.0628219\\
2.5798	-0.0628193\\
2.5824	-0.0628169\\
2.5849	-0.062815\\
2.5875	-0.0628133\\
2.59	-0.062812\\
2.5926	-0.0628109\\
2.5952	-0.0628102\\
2.5977	-0.0628098\\
2.6003	-0.0628097\\
2.6029	-0.0628099\\
2.6054	-0.0628104\\
2.608	-0.0628113\\
2.6106	-0.0628125\\
2.6131	-0.0628139\\
2.6157	-0.0628158\\
2.6182	-0.0628178\\
2.6208	-0.0628203\\
2.6234	-0.0628231\\
2.6259	-0.0628261\\
2.6285	-0.0628296\\
2.6311	-0.0628334\\
2.6336	-0.0628373\\
2.6362	-0.0628417\\
2.6388	-0.0628464\\
2.6413	-0.0628512\\
2.6439	-0.0628566\\
2.6464	-0.062862\\
2.649	-0.062868\\
2.6516	-0.0628743\\
2.6541	-0.0628806\\
2.6567	-0.0628874\\
2.6593	-0.0628946\\
2.6618	-0.0629018\\
2.6644	-0.0629095\\
2.667	-0.0629176\\
2.6695	-0.0629256\\
2.6721	-0.0629342\\
2.6746	-0.0629428\\
2.6772	-0.0629519\\
2.6798	-0.0629614\\
2.6823	-0.0629707\\
2.6849	-0.0629807\\
2.6875	-0.0629909\\
2.69	-0.063001\\
2.6926	-0.0630118\\
2.6952	-0.0630227\\
2.6977	-0.0630335\\
2.7003	-0.063045\\
2.7029	-0.0630567\\
2.7054	-0.0630682\\
2.708	-0.0630803\\
2.7105	-0.0630922\\
2.7131	-0.0631048\\
2.7157	-0.0631176\\
2.7182	-0.06313\\
2.7208	-0.0631432\\
2.7234	-0.0631566\\
2.7259	-0.0631696\\
2.7285	-0.0631833\\
2.7311	-0.0631972\\
2.7336	-0.0632108\\
2.7362	-0.063225\\
2.7387	-0.0632388\\
2.7413	-0.0632534\\
2.7439	-0.063268\\
2.7464	-0.0632823\\
2.749	-0.0632972\\
2.7516	-0.0633123\\
2.7541	-0.0633268\\
2.7567	-0.0633421\\
2.7593	-0.0633575\\
2.7618	-0.0633724\\
2.7644	-0.063388\\
2.7669	-0.063403\\
2.7695	-0.0634187\\
2.7721	-0.0634345\\
2.7746	-0.0634498\\
2.7772	-0.0634657\\
2.7798	-0.0634816\\
2.7823	-0.063497\\
2.7849	-0.063513\\
2.7875	-0.063529\\
2.79	-0.0635445\\
2.7926	-0.0635606\\
2.7951	-0.063576\\
2.7977	-0.0635921\\
2.8003	-0.0636081\\
2.8028	-0.0636236\\
2.8054	-0.0636396\\
2.808	-0.0636555\\
2.8105	-0.0636709\\
2.8131	-0.0636868\\
2.8157	-0.0637026\\
2.8182	-0.0637177\\
2.8208	-0.0637334\\
2.8233	-0.0637485\\
2.8259	-0.063764\\
2.8285	-0.0637794\\
2.831	-0.0637942\\
2.8336	-0.0638094\\
2.8362	-0.0638245\\
2.8387	-0.0638389\\
2.8413	-0.0638538\\
2.8439	-0.0638685\\
2.8464	-0.0638825\\
2.849	-0.0638969\\
2.8516	-0.0639112\\
2.8541	-0.0639247\\
2.8567	-0.0639387\\
2.8592	-0.0639519\\
2.8618	-0.0639654\\
2.8644	-0.0639788\\
2.8669	-0.0639914\\
2.8695	-0.0640044\\
2.8721	-0.0640171\\
2.8746	-0.0640291\\
2.8772	-0.0640414\\
2.8798	-0.0640534\\
2.8823	-0.0640647\\
2.8849	-0.0640763\\
2.8874	-0.0640871\\
2.89	-0.0640981\\
2.8926	-0.0641089\\
2.8951	-0.064119\\
2.8977	-0.0641291\\
2.9003	-0.064139\\
2.9028	-0.0641483\\
2.9054	-0.0641576\\
2.908	-0.0641665\\
2.9105	-0.0641749\\
2.9131	-0.0641832\\
2.9156	-0.064191\\
2.9182	-0.0641987\\
2.9208	-0.064206\\
2.9233	-0.0642128\\
2.9259	-0.0642194\\
2.9285	-0.0642257\\
2.931	-0.0642314\\
2.9336	-0.064237\\
2.9362	-0.0642422\\
2.9387	-0.0642469\\
2.9413	-0.0642513\\
2.9438	-0.0642553\\
2.9464	-0.0642589\\
2.949	-0.0642622\\
2.9515	-0.064265\\
2.9541	-0.0642675\\
2.9567	-0.0642695\\
2.9592	-0.0642711\\
2.9618	-0.0642724\\
2.9644	-0.0642732\\
2.9669	-0.0642736\\
2.9695	-0.0642735\\
2.972	-0.0642731\\
2.9746	-0.0642722\\
2.9772	-0.0642708\\
2.9797	-0.0642691\\
2.9823	-0.0642668\\
2.9849	-0.0642641\\
2.9874	-0.0642611\\
2.99	-0.0642575\\
2.9926	-0.0642534\\
2.9951	-0.064249\\
2.9977	-0.064244\\
3.0003	-0.0642385\\
3.0028	-0.0642328\\
3.0054	-0.0642264\\
3.0079	-0.0642198\\
3.0105	-0.0642124\\
3.0131	-0.0642045\\
3.0156	-0.0641965\\
3.0182	-0.0641877\\
3.0208	-0.0641784\\
3.0233	-0.0641689\\
3.0259	-0.0641586\\
3.0285	-0.0641478\\
3.031	-0.064137\\
3.0336	-0.0641252\\
3.0361	-0.0641134\\
3.0387	-0.0641007\\
3.0413	-0.0640874\\
3.0438	-0.0640742\\
3.0464	-0.06406\\
3.049	-0.0640452\\
3.0515	-0.0640305\\
3.0541	-0.0640148\\
3.0567	-0.0639986\\
3.0592	-0.0639825\\
3.0618	-0.0639652\\
3.0643	-0.0639482\\
3.0669	-0.0639299\\
3.0695	-0.0639112\\
3.072	-0.0638927\\
3.0746	-0.063873\\
3.0772	-0.0638528\\
3.0797	-0.0638329\\
3.0823	-0.0638117\\
3.0849	-0.0637899\\
3.0874	-0.0637686\\
3.09	-0.0637459\\
3.0925	-0.0637237\\
3.0951	-0.0637\\
3.0977	-0.0636759\\
3.1002	-0.0636522\\
3.1028	-0.0636271\\
3.1054	-0.0636015\\
3.1079	-0.0635765\\
3.1105	-0.06355\\
3.1131	-0.063523\\
3.1156	-0.0634965\\
3.1182	-0.0634686\\
3.1207	-0.0634413\\
3.1233	-0.0634125\\
3.1259	-0.0633831\\
3.1284	-0.0633545\\
3.131	-0.0633243\\
3.1336	-0.0632936\\
3.1361	-0.0632637\\
3.1387	-0.0632321\\
3.1413	-0.0632001\\
3.1438	-0.063169\\
3.1464	-0.0631361\\
3.1489	-0.0631041\\
3.1515	-0.0630704\\
3.1541	-0.0630362\\
3.1566	-0.063003\\
3.1592	-0.0629681\\
3.1618	-0.0629327\\
3.1643	-0.0628983\\
3.1669	-0.0628621\\
3.1695	-0.0628255\\
3.172	-0.06279\\
3.1746	-0.0627527\\
3.1772	-0.0627149\\
3.1797	-0.0626783\\
3.1823	-0.0626398\\
3.1848	-0.0626025\\
3.1874	-0.0625633\\
3.19	-0.0625237\\
3.1925	-0.0624854\\
3.1951	-0.0624451\\
3.1977	-0.0624045\\
3.2002	-0.0623651\\
3.2028	-0.0623239\\
3.2054	-0.0622823\\
3.2079	-0.062242\\
3.2105	-0.0621998\\
3.213	-0.0621589\\
3.2156	-0.062116\\
3.2182	-0.0620729\\
3.2207	-0.0620312\\
3.2233	-0.0619875\\
3.2259	-0.0619435\\
3.2284	-0.0619009\\
3.231	-0.0618564\\
3.2336	-0.0618116\\
3.2361	-0.0617683\\
3.2387	-0.061723\\
3.2412	-0.0616792\\
3.2438	-0.0616334\\
3.2464	-0.0615874\\
3.2489	-0.0615429\\
3.2515	-0.0614964\\
3.2541	-0.0614498\\
3.2566	-0.0614047\\
3.2592	-0.0613576\\
3.2618	-0.0613103\\
3.2643	-0.0612647\\
3.2669	-0.061217\\
3.2694	-0.061171\\
3.272	-0.061123\\
3.2746	-0.0610749\\
3.2771	-0.0610284\\
3.2797	-0.06098\\
3.2823	-0.0609314\\
3.2848	-0.0608845\\
3.2874	-0.0608356\\
3.29	-0.0607866\\
3.2925	-0.0607394\\
3.2951	-0.0606901\\
3.2976	-0.0606427\\
3.3002	-0.0605933\\
3.3028	-0.0605437\\
3.3053	-0.060496\\
3.3079	-0.0604463\\
3.3105	-0.0603966\\
3.313	-0.0603487\\
3.3156	-0.0602988\\
3.3182	-0.0602488\\
3.3207	-0.0602007\\
3.3233	-0.0601507\\
3.3259	-0.0601006\\
3.3284	-0.0600525\\
3.331	-0.0600023\\
3.3335	-0.0599541\\
3.3361	-0.059904\\
3.3387	-0.0598538\\
3.3412	-0.0598056\\
3.3438	-0.0597555\\
3.3464	-0.0597054\\
3.3489	-0.0596572\\
3.3515	-0.0596071\\
3.3541	-0.0595571\\
3.3566	-0.059509\\
3.3592	-0.0594591\\
3.3617	-0.0594111\\
3.3643	-0.0593612\\
3.3669	-0.0593115\\
3.3694	-0.0592637\\
3.372	-0.0592141\\
3.3746	-0.0591645\\
3.3771	-0.0591169\\
3.3797	-0.0590676\\
3.3823	-0.0590183\\
3.3848	-0.058971\\
3.3874	-0.058922\\
3.3899	-0.0588749\\
3.3925	-0.0588261\\
3.3951	-0.0587774\\
3.3976	-0.0587307\\
3.4002	-0.0586823\\
3.4028	-0.058634\\
3.4053	-0.0585878\\
3.4079	-0.0585398\\
3.4105	-0.058492\\
3.413	-0.0584462\\
3.4156	-0.0583987\\
3.4181	-0.0583532\\
3.4207	-0.0583061\\
3.4233	-0.0582591\\
3.4258	-0.0582142\\
3.4284	-0.0581676\\
3.431	-0.0581212\\
3.4335	-0.0580768\\
3.4361	-0.0580309\\
3.4387	-0.0579851\\
3.4412	-0.0579414\\
3.4438	-0.057896\\
3.4463	-0.0578527\\
3.4489	-0.0578078\\
3.4515	-0.0577632\\
3.454	-0.0577205\\
3.4566	-0.0576764\\
3.4592	-0.0576325\\
3.4617	-0.0575905\\
3.4643	-0.0575471\\
3.4669	-0.057504\\
3.4694	-0.0574628\\
3.472	-0.0574202\\
3.4745	-0.0573794\\
3.4771	-0.0573374\\
3.4797	-0.0572956\\
3.4822	-0.0572557\\
3.4848	-0.0572144\\
3.4874	-0.0571735\\
3.4899	-0.0571344\\
3.4925	-0.057094\\
3.4951	-0.0570539\\
3.4976	-0.0570157\\
3.5002	-0.0569762\\
3.5028	-0.056937\\
3.5053	-0.0568997\\
3.5079	-0.0568611\\
3.5104	-0.0568243\\
3.513	-0.0567864\\
3.5156	-0.0567488\\
3.5181	-0.0567129\\
3.5207	-0.0566759\\
3.5233	-0.0566393\\
3.5258	-0.0566043\\
3.5284	-0.0565683\\
3.531	-0.0565326\\
3.5335	-0.0564987\\
3.5361	-0.0564637\\
3.5386	-0.0564303\\
3.5412	-0.056396\\
3.5438	-0.056362\\
3.5463	-0.0563296\\
3.5489	-0.0562963\\
3.5515	-0.0562633\\
3.554	-0.056232\\
3.5566	-0.0561997\\
3.5592	-0.0561677\\
3.5617	-0.0561374\\
3.5643	-0.0561061\\
3.5668	-0.0560764\\
3.5694	-0.0560459\\
3.572	-0.0560157\\
3.5745	-0.0559871\\
3.5771	-0.0559576\\
3.5797	-0.0559285\\
3.5822	-0.0559008\\
3.5848	-0.0558724\\
3.5874	-0.0558444\\
3.5899	-0.0558177\\
3.5925	-0.0557904\\
3.595	-0.0557645\\
3.5976	-0.0557379\\
3.6002	-0.0557116\\
3.6027	-0.0556867\\
3.6053	-0.0556612\\
3.6079	-0.055636\\
3.6104	-0.0556121\\
3.613	-0.0555877\\
3.6156	-0.0555636\\
3.6181	-0.0555407\\
3.6207	-0.0555174\\
3.6232	-0.0554952\\
3.6258	-0.0554726\\
3.6284	-0.0554502\\
3.6309	-0.0554291\\
3.6335	-0.0554075\\
3.6361	-0.0553863\\
3.6386	-0.0553662\\
3.6412	-0.0553457\\
3.6438	-0.0553255\\
3.6463	-0.0553065\\
3.6489	-0.055287\\
3.6515	-0.0552679\\
3.654	-0.0552498\\
3.6566	-0.0552314\\
3.6591	-0.055214\\
3.6617	-0.0551963\\
3.6643	-0.0551789\\
3.6668	-0.0551626\\
3.6694	-0.0551459\\
3.672	-0.0551295\\
3.6745	-0.0551142\\
3.6771	-0.0550985\\
3.6797	-0.0550832\\
3.6822	-0.0550688\\
3.6848	-0.0550541\\
3.6873	-0.0550403\\
3.6899	-0.0550264\\
3.6925	-0.0550127\\
3.695	-0.0549999\\
3.6976	-0.0549869\\
3.7002	-0.0549742\\
3.7027	-0.0549624\\
3.7053	-0.0549503\\
3.7079	-0.0549386\\
3.7104	-0.0549277\\
3.713	-0.0549166\\
3.7155	-0.0549063\\
3.7181	-0.0548958\\
3.7207	-0.0548857\\
3.7232	-0.0548763\\
3.7258	-0.0548668\\
3.7284	-0.0548576\\
3.7309	-0.054849\\
3.7335	-0.0548404\\
3.7361	-0.0548321\\
3.7386	-0.0548244\\
3.7412	-0.0548167\\
3.7437	-0.0548096\\
3.7463	-0.0548024\\
3.7489	-0.0547956\\
3.7514	-0.0547893\\
3.754	-0.054783\\
3.7566	-0.0547771\\
3.7591	-0.0547716\\
3.7617	-0.0547661\\
3.7643	-0.054761\\
3.7668	-0.0547563\\
3.7694	-0.0547517\\
3.7719	-0.0547475\\
3.7745	-0.0547434\\
3.7771	-0.0547396\\
3.7796	-0.0547362\\
3.7822	-0.0547329\\
3.7848	-0.0547298\\
3.7873	-0.0547272\\
3.7899	-0.0547246\\
3.7925	-0.0547223\\
3.795	-0.0547204\\
3.7976	-0.0547186\\
3.8002	-0.054717\\
3.8027	-0.0547157\\
3.8053	-0.0547147\\
3.8078	-0.0547138\\
3.8104	-0.0547132\\
3.813	-0.0547128\\
3.8155	-0.0547127\\
3.8181	-0.0547127\\
3.8207	-0.054713\\
3.8232	-0.0547135\\
3.8258	-0.0547142\\
3.8284	-0.0547151\\
3.8309	-0.0547162\\
3.8335	-0.0547176\\
3.836	-0.054719\\
3.8386	-0.0547208\\
3.8412	-0.0547228\\
3.8437	-0.0547248\\
3.8463	-0.0547272\\
3.8489	-0.0547297\\
3.8514	-0.0547323\\
3.854	-0.0547353\\
3.8566	-0.0547384\\
3.8591	-0.0547415\\
3.8617	-0.054745\\
3.8642	-0.0547485\\
3.8668	-0.0547524\\
3.8694	-0.0547564\\
3.8719	-0.0547604\\
3.8745	-0.0547647\\
3.8771	-0.0547692\\
3.8796	-0.0547737\\
3.8822	-0.0547786\\
3.8848	-0.0547836\\
3.8873	-0.0547885\\
3.8899	-0.0547938\\
3.8924	-0.0547991\\
3.895	-0.0548047\\
3.8976	-0.0548104\\
3.9001	-0.0548161\\
3.9027	-0.0548222\\
3.9053	-0.0548284\\
3.9078	-0.0548344\\
3.9104	-0.0548409\\
3.913	-0.0548475\\
3.9155	-0.054854\\
3.9181	-0.0548609\\
3.9206	-0.0548676\\
3.9232	-0.0548747\\
3.9258	-0.0548819\\
3.9283	-0.054889\\
3.9309	-0.0548965\\
3.9335	-0.0549041\\
3.936	-0.0549115\\
3.9386	-0.0549193\\
3.9412	-0.0549273\\
3.9437	-0.054935\\
3.9463	-0.0549432\\
3.9488	-0.0549511\\
3.9514	-0.0549595\\
3.954	-0.054968\\
3.9565	-0.0549762\\
3.9591	-0.0549849\\
3.9617	-0.0549937\\
3.9642	-0.0550022\\
3.9668	-0.0550111\\
3.9694	-0.0550202\\
3.9719	-0.0550289\\
3.9745	-0.0550382\\
3.9771	-0.0550475\\
3.9796	-0.0550565\\
3.9822	-0.055066\\
3.9847	-0.0550752\\
3.9873	-0.0550848\\
3.9899	-0.0550945\\
3.9924	-0.0551039\\
3.995	-0.0551137\\
3.9976	-0.0551237\\
4.0001	-0.0551333\\
4.0027	-0.0551434\\
4.0053	-0.0551535\\
4.0078	-0.0551633\\
4.0104	-0.0551736\\
4.0129	-0.0551835\\
4.0155	-0.0551939\\
4.0181	-0.0552043\\
4.0206	-0.0552145\\
4.0232	-0.055225\\
4.0258	-0.0552357\\
4.0283	-0.0552459\\
4.0309	-0.0552567\\
4.0335	-0.0552674\\
4.036	-0.0552779\\
4.0386	-0.0552888\\
4.0411	-0.0552993\\
4.0437	-0.0553103\\
4.0463	-0.0553213\\
4.0488	-0.0553319\\
4.0514	-0.0553431\\
4.054	-0.0553542\\
4.0565	-0.055365\\
4.0591	-0.0553762\\
4.0617	-0.0553875\\
4.0642	-0.0553984\\
4.0668	-0.0554098\\
4.0693	-0.0554207\\
4.0719	-0.0554322\\
4.0745	-0.0554436\\
4.077	-0.0554547\\
4.0796	-0.0554662\\
4.0822	-0.0554778\\
4.0847	-0.055489\\
4.0873	-0.0555006\\
4.0899	-0.0555122\\
4.0924	-0.0555235\\
4.095	-0.0555352\\
4.0975	-0.0555465\\
4.1001	-0.0555582\\
4.1027	-0.05557\\
4.1052	-0.0555814\\
4.1078	-0.0555932\\
4.1104	-0.0556051\\
4.1129	-0.0556165\\
4.1155	-0.0556284\\
4.1181	-0.0556404\\
4.1206	-0.0556519\\
4.1232	-0.0556638\\
4.1258	-0.0556758\\
4.1283	-0.0556873\\
4.1309	-0.0556994\\
4.1334	-0.0557109\\
4.136	-0.055723\\
4.1386	-0.0557351\\
4.1411	-0.0557467\\
4.1437	-0.0557588\\
4.1463	-0.0557709\\
4.1488	-0.0557826\\
4.1514	-0.0557947\\
4.154	-0.0558069\\
4.1565	-0.0558185\\
4.1591	-0.0558307\\
4.1616	-0.0558424\\
4.1642	-0.0558546\\
4.1668	-0.0558669\\
4.1693	-0.0558786\\
4.1719	-0.0558908\\
4.1745	-0.0559031\\
4.177	-0.0559148\\
4.1796	-0.0559271\\
4.1822	-0.0559394\\
4.1847	-0.0559511\\
4.1873	-0.0559634\\
4.1898	-0.0559752\\
4.1924	-0.0559875\\
4.195	-0.0559998\\
4.1975	-0.0560116\\
4.2001	-0.0560239\\
4.2027	-0.0560363\\
4.2052	-0.0560481\\
4.2078	-0.0560604\\
4.2104	-0.0560727\\
4.2129	-0.0560846\\
4.2155	-0.0560969\\
4.218	-0.0561088\\
4.2206	-0.0561211\\
4.2232	-0.0561334\\
4.2257	-0.0561453\\
4.2283	-0.0561576\\
4.2309	-0.05617\\
4.2334	-0.0561818\\
4.236	-0.0561942\\
4.2386	-0.0562065\\
4.2411	-0.0562184\\
4.2437	-0.0562307\\
4.2462	-0.0562426\\
4.2488	-0.056255\\
4.2514	-0.0562673\\
4.2539	-0.0562792\\
4.2565	-0.0562915\\
4.2591	-0.0563038\\
4.2616	-0.0563157\\
4.2642	-0.056328\\
4.2668	-0.0563404\\
4.2693	-0.0563522\\
4.2719	-0.0563646\\
4.2744	-0.0563764\\
4.277	-0.0563887\\
4.2796	-0.056401\\
4.2821	-0.0564129\\
4.2847	-0.0564252\\
4.2873	-0.0564375\\
4.2898	-0.0564493\\
4.2924	-0.0564616\\
4.295	-0.0564739\\
4.2975	-0.0564857\\
4.3001	-0.056498\\
4.3027	-0.0565103\\
4.3052	-0.0565221\\
4.3078	-0.0565343\\
4.3103	-0.0565461\\
4.3129	-0.0565583\\
4.3155	-0.0565706\\
4.318	-0.0565823\\
4.3206	-0.0565946\\
4.3232	-0.0566068\\
4.3257	-0.0566185\\
4.3283	-0.0566307\\
4.3309	-0.0566429\\
4.3334	-0.0566546\\
4.336	-0.0566668\\
4.3385	-0.0566784\\
4.3411	-0.0566906\\
4.3437	-0.0567027\\
4.3462	-0.0567144\\
4.3488	-0.0567265\\
4.3514	-0.0567386\\
4.3539	-0.0567502\\
4.3565	-0.0567623\\
4.3591	-0.0567744\\
4.3616	-0.0567859\\
4.3642	-0.056798\\
4.3667	-0.0568095\\
4.3693	-0.0568215\\
4.3719	-0.0568335\\
4.3744	-0.056845\\
4.377	-0.056857\\
4.3796	-0.0568689\\
4.3821	-0.0568804\\
4.3847	-0.0568923\\
4.3873	-0.0569042\\
4.3898	-0.0569156\\
4.3924	-0.0569275\\
4.3949	-0.0569389\\
4.3975	-0.0569507\\
4.4001	-0.0569625\\
4.4026	-0.0569738\\
4.4052	-0.0569856\\
4.4078	-0.0569973\\
4.4103	-0.0570085\\
4.4129	-0.0570202\\
4.4155	-0.0570319\\
4.418	-0.0570431\\
4.4206	-0.0570547\\
4.4231	-0.0570659\\
4.4257	-0.0570774\\
4.4283	-0.057089\\
4.4308	-0.0571001\\
4.4334	-0.0571116\\
4.436	-0.057123\\
4.4385	-0.057134\\
4.4411	-0.0571454\\
4.4437	-0.0571568\\
4.4462	-0.0571677\\
4.4488	-0.057179\\
4.4514	-0.0571903\\
4.4539	-0.0572012\\
4.4565	-0.0572124\\
4.459	-0.0572232\\
4.4616	-0.0572343\\
4.4642	-0.0572454\\
4.4667	-0.0572561\\
4.4693	-0.0572672\\
4.4719	-0.0572782\\
4.4744	-0.0572888\\
4.477	-0.0572997\\
4.4796	-0.0573106\\
4.4821	-0.0573211\\
4.4847	-0.0573319\\
4.4872	-0.0573423\\
4.4898	-0.0573531\\
4.4924	-0.0573638\\
4.4949	-0.057374\\
4.4975	-0.0573847\\
4.5001	-0.0573952\\
4.5026	-0.0574054\\
4.5052	-0.0574159\\
4.5078	-0.0574263\\
4.5103	-0.0574363\\
4.5129	-0.0574467\\
4.5154	-0.0574566\\
4.518	-0.0574669\\
4.5206	-0.0574771\\
4.5231	-0.0574869\\
4.5257	-0.057497\\
4.5283	-0.0575071\\
4.5308	-0.0575167\\
4.5334	-0.0575266\\
4.536	-0.0575365\\
4.5385	-0.057546\\
4.5411	-0.0575558\\
4.5436	-0.0575652\\
4.5462	-0.0575749\\
4.5488	-0.0575845\\
4.5513	-0.0575937\\
4.5539	-0.0576032\\
4.5565	-0.0576127\\
4.559	-0.0576217\\
4.5616	-0.057631\\
4.5642	-0.0576403\\
4.5667	-0.0576491\\
4.5693	-0.0576583\\
4.5718	-0.057667\\
4.5744	-0.057676\\
4.577	-0.0576849\\
4.5795	-0.0576935\\
4.5821	-0.0577023\\
4.5847	-0.057711\\
4.5872	-0.0577194\\
4.5898	-0.057728\\
4.5924	-0.0577365\\
4.5949	-0.0577446\\
4.5975	-0.057753\\
4.6001	-0.0577613\\
4.6026	-0.0577692\\
4.6052	-0.0577774\\
4.6077	-0.0577852\\
4.6103	-0.0577932\\
4.6129	-0.0578011\\
4.6154	-0.0578087\\
4.618	-0.0578164\\
4.6206	-0.0578241\\
4.6231	-0.0578315\\
4.6257	-0.057839\\
4.6283	-0.0578464\\
4.6308	-0.0578535\\
4.6334	-0.0578608\\
4.6359	-0.0578677\\
4.6385	-0.0578749\\
4.6411	-0.0578819\\
4.6436	-0.0578886\\
4.6462	-0.0578954\\
4.6488	-0.0579022\\
4.6513	-0.0579086\\
4.6539	-0.0579152\\
4.6565	-0.0579217\\
4.659	-0.0579279\\
4.6616	-0.0579342\\
4.6641	-0.0579402\\
4.6667	-0.0579463\\
4.6693	-0.0579523\\
4.6718	-0.057958\\
4.6744	-0.0579639\\
4.677	-0.0579696\\
4.6795	-0.0579751\\
4.6821	-0.0579806\\
4.6847	-0.0579861\\
4.6872	-0.0579912\\
4.6898	-0.0579965\\
4.6923	-0.0580015\\
4.6949	-0.0580065\\
4.6975	-0.0580115\\
4.7	-0.0580162\\
4.7026	-0.0580209\\
4.7052	-0.0580256\\
4.7077	-0.0580299\\
4.7103	-0.0580344\\
4.7129	-0.0580387\\
4.7154	-0.0580428\\
4.718	-0.0580469\\
4.7205	-0.0580508\\
4.7231	-0.0580547\\
4.7257	-0.0580585\\
4.7282	-0.058062\\
4.7308	-0.0580656\\
4.7334	-0.0580691\\
4.7359	-0.0580723\\
4.7385	-0.0580756\\
4.7411	-0.0580787\\
4.7436	-0.0580816\\
4.7462	-0.0580845\\
4.7487	-0.0580872\\
4.7513	-0.0580899\\
4.7539	-0.0580925\\
4.7564	-0.0580949\\
4.759	-0.0580972\\
4.7616	-0.0580995\\
4.7641	-0.0581015\\
4.7667	-0.0581035\\
4.7693	-0.0581054\\
4.7718	-0.0581071\\
4.7744	-0.0581088\\
4.777	-0.0581103\\
4.7795	-0.0581117\\
4.7821	-0.058113\\
4.7846	-0.0581141\\
4.7872	-0.0581152\\
4.7898	-0.0581161\\
4.7923	-0.0581169\\
4.7949	-0.0581176\\
4.7975	-0.0581182\\
4.8	-0.0581187\\
4.8026	-0.058119\\
4.8052	-0.0581193\\
4.8077	-0.0581194\\
4.8103	-0.0581194\\
4.8128	-0.0581192\\
4.8154	-0.058119\\
4.818	-0.0581186\\
4.8205	-0.0581182\\
4.8231	-0.0581175\\
4.8257	-0.0581168\\
4.8282	-0.058116\\
4.8308	-0.058115\\
4.8334	-0.0581139\\
4.8359	-0.0581127\\
4.8385	-0.0581114\\
4.841	-0.05811\\
4.8436	-0.0581084\\
4.8462	-0.0581067\\
4.8487	-0.0581049\\
4.8513	-0.058103\\
4.8539	-0.0581009\\
4.8564	-0.0580988\\
4.859	-0.0580964\\
4.8616	-0.058094\\
4.8641	-0.0580915\\
4.8667	-0.0580888\\
4.8692	-0.0580861\\
4.8718	-0.0580832\\
4.8744	-0.0580802\\
4.8769	-0.0580771\\
4.8795	-0.0580738\\
4.8821	-0.0580704\\
4.8846	-0.058067\\
4.8872	-0.0580634\\
4.8898	-0.0580596\\
4.8923	-0.0580558\\
4.8949	-0.0580518\\
4.8974	-0.0580478\\
4.9	-0.0580436\\
4.9026	-0.0580392\\
4.9051	-0.0580349\\
4.9077	-0.0580303\\
4.9103	-0.0580255\\
4.9128	-0.0580209\\
4.9154	-0.0580159\\
4.918	-0.0580108\\
4.9205	-0.0580058\\
4.9231	-0.0580005\\
4.9257	-0.0579951\\
4.9282	-0.0579897\\
4.9308	-0.0579841\\
4.9333	-0.0579785\\
4.9359	-0.0579726\\
4.9385	-0.0579666\\
4.941	-0.0579607\\
4.9436	-0.0579545\\
4.9462	-0.0579481\\
4.9487	-0.0579419\\
4.9513	-0.0579354\\
4.9539	-0.0579287\\
4.9564	-0.0579221\\
4.959	-0.0579152\\
4.9615	-0.0579085\\
4.9641	-0.0579014\\
4.9667	-0.0578941\\
4.9692	-0.0578871\\
4.9718	-0.0578796\\
4.9744	-0.0578721\\
4.9769	-0.0578647\\
4.9795	-0.0578569\\
4.9821	-0.0578491\\
4.9846	-0.0578414\\
4.9872	-0.0578333\\
4.9897	-0.0578254\\
4.9923	-0.0578172\\
4.9949	-0.0578088\\
4.9974	-0.0578006\\
5	-0.057792\\
};
\addplot [color=mycolor21, line width=1.0pt, forget plot]
  table[row sep=crcr]{%
-0.125	4.80238e-08\\
-0.1224	4.90173e-08\\
-0.1199	4.99727e-08\\
-0.1173	5.09661e-08\\
-0.1147	5.19594e-08\\
-0.1122	5.29144e-08\\
-0.1096	5.39075e-08\\
-0.1071	5.48621e-08\\
-0.1045	5.5855e-08\\
-0.1019	5.68478e-08\\
-0.0994	5.78024e-08\\
-0.0968	5.8795e-08\\
-0.0942	5.97876e-08\\
-0.0917	6.07418e-08\\
-0.0891	6.17342e-08\\
-0.0865	6.27265e-08\\
-0.084	6.36805e-08\\
-0.0814	6.46726e-08\\
-0.0789	6.56264e-08\\
-0.0763	6.66183e-08\\
-0.0737	6.761e-08\\
-0.0712	6.85634e-08\\
-0.0686	6.95551e-08\\
-0.066	7.05466e-08\\
-0.0635	7.14998e-08\\
-0.0609	7.2491e-08\\
-0.0583	7.34822e-08\\
-0.0558	7.44351e-08\\
-0.0532	7.54261e-08\\
-0.0507	7.6379e-08\\
-0.0481	7.73697e-08\\
-0.0455	7.83605e-08\\
-0.043	7.93131e-08\\
-0.0404	8.03036e-08\\
-0.0378	8.12941e-08\\
-0.0353	8.22465e-08\\
-0.0327	8.32367e-08\\
-0.0301	8.4227e-08\\
-0.0276	8.51791e-08\\
-0.025	8.6169e-08\\
-0.0224	8.71591e-08\\
-0.0199	8.81109e-08\\
-0.0173	8.91006e-08\\
-0.0148	9.00522e-08\\
-0.0122	9.10418e-08\\
-0.0096	9.20313e-08\\
-0.0071	9.29827e-08\\
-0.0045	9.3972e-08\\
-0.0019	9.49612e-08\\
0.0006	9.59124e-08\\
0.0032	0.0118757\\
0.0058	0.0188848\\
0.0083	0.023618\\
0.0109	0.0277109\\
0.0134	0.031164\\
0.016	0.0341591\\
0.0186	0.0364368\\
0.0211	0.0380453\\
0.0237	0.0393606\\
0.0263	0.0405221\\
0.0288	0.0415613\\
0.0314	0.0425149\\
0.034	0.0432535\\
0.0365	0.0437318\\
0.0391	0.0440322\\
0.0416	0.0441919\\
0.0442	0.0442718\\
0.0468	0.0442658\\
0.0493	0.0441527\\
0.0519	0.0439026\\
0.0545	0.0435216\\
0.057	0.043057\\
0.0596	0.0425029\\
0.0622	0.0418924\\
0.0647	0.0412471\\
0.0673	0.0405017\\
0.0698	0.0397099\\
0.0724	0.038818\\
0.075	0.03788\\
0.0775	0.0369529\\
0.0801	0.0359649\\
0.0827	0.0349451\\
0.0852	0.0339267\\
0.0878	0.0328291\\
0.0904	0.0317076\\
0.0929	0.030624\\
0.0955	0.0294999\\
0.098	0.0284187\\
0.1006	0.0272841\\
0.1032	0.0261339\\
0.1057	0.0250193\\
0.1083	0.0238661\\
0.1109	0.0227305\\
0.1134	0.0216564\\
0.116	0.0205495\\
0.1186	0.0194434\\
0.1211	0.0183798\\
0.1237	0.0172835\\
0.1263	0.016209\\
0.1288	0.0152014\\
0.1314	0.0141753\\
0.1339	0.0131993\\
0.1365	0.0121894\\
0.1391	0.0111881\\
0.1416	0.0102429\\
0.1442	0.00928648\\
0.1468	0.00835621\\
0.1493	0.00747785\\
0.1519	0.00657199\\
0.1545	0.00567186\\
0.157	0.00481786\\
0.1596	0.00395054\\
0.1621	0.00313985\\
0.1647	0.00231551\\
0.1673	0.00150062\\
0.1698	0.000720151\\
0.1724	-8.64348e-05\\
0.175	-0.000879822\\
0.1775	-0.0016242\\
0.1801	-0.0023803\\
0.1827	-0.00312579\\
0.1852	-0.0038403\\
0.1878	-0.00458305\\
0.1903	-0.0052919\\
0.1929	-0.0060162\\
0.1955	-0.00672443\\
0.198	-0.00739434\\
0.2006	-0.00808743\\
0.2032	-0.00878193\\
0.2057	-0.00944922\\
0.2083	-0.0101359\\
0.2109	-0.0108096\\
0.2134	-0.0114454\\
0.216	-0.0121004\\
0.2185	-0.0127303\\
0.2211	-0.0133866\\
0.2237	-0.0140395\\
0.2262	-0.0146577\\
0.2288	-0.0152877\\
0.2314	-0.0159076\\
0.2339	-0.0165001\\
0.2365	-0.0171163\\
0.2391	-0.0177305\\
0.2416	-0.0183133\\
0.2442	-0.0189061\\
0.2467	-0.0194634\\
0.2493	-0.0200341\\
0.2519	-0.020601\\
0.2544	-0.0211431\\
0.257	-0.0216996\\
0.2596	-0.0222427\\
0.2621	-0.0227495\\
0.2647	-0.0232628\\
0.2673	-0.023767\\
0.2698	-0.024246\\
0.2724	-0.0247362\\
0.2749	-0.0251951\\
0.2775	-0.0256547\\
0.2801	-0.0260961\\
0.2826	-0.026507\\
0.2852	-0.0269241\\
0.2878	-0.0273316\\
0.2903	-0.0277108\\
0.2929	-0.0280873\\
0.2955	-0.028443\\
0.298	-0.0287674\\
0.3006	-0.0290905\\
0.3032	-0.0294015\\
0.3057	-0.029688\\
0.3083	-0.0299685\\
0.3108	-0.0302181\\
0.3134	-0.0304568\\
0.316	-0.0306772\\
0.3185	-0.0308756\\
0.3211	-0.0310685\\
0.3237	-0.031245\\
0.3262	-0.0313956\\
0.3288	-0.0315305\\
0.3314	-0.0316452\\
0.3339	-0.03174\\
0.3365	-0.0318248\\
0.339	-0.0318923\\
0.3416	-0.0319449\\
0.3442	-0.0319765\\
0.3467	-0.0319874\\
0.3493	-0.0319813\\
0.3519	-0.0319609\\
0.3544	-0.0319287\\
0.357	-0.0318802\\
0.3596	-0.0318133\\
0.3621	-0.0317306\\
0.3647	-0.0316271\\
0.3672	-0.031514\\
0.3698	-0.0313845\\
0.3724	-0.0312424\\
0.3749	-0.0310917\\
0.3775	-0.0309181\\
0.3801	-0.0307278\\
0.3826	-0.0305315\\
0.3852	-0.0303165\\
0.3878	-0.0300913\\
0.3903	-0.029864\\
0.3929	-0.0296139\\
0.3954	-0.0293594\\
0.398	-0.0290818\\
0.4006	-0.0287939\\
0.4031	-0.0285093\\
0.4057	-0.028205\\
0.4083	-0.0278902\\
0.4108	-0.0275762\\
0.4134	-0.027238\\
0.416	-0.0268903\\
0.4185	-0.0265493\\
0.4211	-0.0261888\\
0.4236	-0.0258353\\
0.4262	-0.0254588\\
0.4288	-0.0250723\\
0.4313	-0.0246926\\
0.4339	-0.0242916\\
0.4365	-0.0238861\\
0.439	-0.0234919\\
0.4416	-0.0230758\\
0.4442	-0.022652\\
0.4467	-0.0222374\\
0.4493	-0.0218007\\
0.4519	-0.0213606\\
0.4544	-0.0209347\\
0.457	-0.0204881\\
0.4595	-0.0200536\\
0.4621	-0.0195958\\
0.4647	-0.0191329\\
0.4672	-0.0186849\\
0.4698	-0.0182173\\
0.4724	-0.0177479\\
0.4749	-0.0172936\\
0.4775	-0.0168167\\
0.4801	-0.0163354\\
0.4826	-0.01587\\
0.4852	-0.0153849\\
0.4877	-0.0149179\\
0.4903	-0.0144309\\
0.4929	-0.0139412\\
0.4954	-0.0134671\\
0.498	-0.0129716\\
0.5006	-0.0124751\\
0.5031	-0.0119979\\
0.5057	-0.0115017\\
0.5083	-0.0110045\\
0.5108	-0.0105243\\
0.5134	-0.0100228\\
0.5159	-0.00953986\\
0.5185	-0.00903814\\
0.5211	-0.0085375\\
0.5236	-0.0080565\\
0.5262	-0.00755551\\
0.5288	-0.00705322\\
0.5313	-0.00656967\\
0.5339	-0.0060675\\
0.5365	-0.00556705\\
0.539	-0.00508737\\
0.5416	-0.00458916\\
0.5441	-0.00410991\\
0.5467	-0.00361135\\
0.5493	-0.00311368\\
0.5518	-0.00263714\\
0.5544	-0.00214406\\
0.557	-0.001653\\
0.5595	-0.00118184\\
0.5621	-0.000692468\\
0.5647	-0.000204332\\
0.5672	0.000262772\\
0.5698	0.000745369\\
0.5723	0.00120638\\
0.5749	0.00168344\\
0.5775	0.00215877\\
0.58	0.00261411\\
0.5826	0.00308501\\
0.5852	0.00355217\\
0.5877	0.00399744\\
0.5903	0.00445685\\
0.5929	0.00491328\\
0.5954	0.00534965\\
0.598	0.00580038\\
0.6006	0.00624695\\
0.6031	0.00667173\\
0.6057	0.00710869\\
0.6082	0.00752482\\
0.6108	0.00795396\\
0.6134	0.00837931\\
0.6159	0.00878399\\
0.6185	0.00919952\\
0.6211	0.00960923\\
0.6236	0.009998\\
0.6262	0.0103976\\
0.6288	0.0107925\\
0.6313	0.0111675\\
0.6339	0.0115517\\
0.6364	0.011915\\
0.639	0.0122864\\
0.6416	0.0126519\\
0.6441	0.012998\\
0.6467	0.0133524\\
0.6493	0.0137006\\
0.6518	0.0140288\\
0.6544	0.0143629\\
0.657	0.01469\\
0.6595	0.0149984\\
0.6621	0.0153129\\
0.6646	0.015609\\
0.6672	0.0159098\\
0.6698	0.0162027\\
0.6723	0.016477\\
0.6749	0.0167549\\
0.6775	0.0170257\\
0.68	0.0172795\\
0.6826	0.0175359\\
0.6852	0.0177842\\
0.6877	0.0180149\\
0.6903	0.0182468\\
0.6928	0.0184625\\
0.6954	0.0186794\\
0.698	0.0188885\\
0.7005	0.0190817\\
0.7031	0.0192739\\
0.7057	0.0194574\\
0.7082	0.0196259\\
0.7108	0.0197933\\
0.7134	0.0199526\\
0.7159	0.0200979\\
0.7185	0.0202402\\
0.721	0.0203685\\
0.7236	0.0204931\\
0.7262	0.0206093\\
0.7287	0.0207133\\
0.7313	0.0208131\\
0.7339	0.020904\\
0.7364	0.0209827\\
0.739	0.0210555\\
0.7416	0.0211196\\
0.7441	0.0211731\\
0.7467	0.0212205\\
0.7492	0.0212578\\
0.7518	0.0212876\\
0.7544	0.0213082\\
0.7569	0.0213193\\
0.7595	0.0213225\\
0.7621	0.0213171\\
0.7646	0.0213038\\
0.7672	0.0212812\\
0.7698	0.0212494\\
0.7723	0.0212102\\
0.7749	0.0211608\\
0.7775	0.0211029\\
0.78	0.0210392\\
0.7826	0.0209645\\
0.7851	0.0208842\\
0.7877	0.0207918\\
0.7903	0.0206905\\
0.7928	0.0205851\\
0.7954	0.0204674\\
0.798	0.0203413\\
0.8005	0.0202119\\
0.8031	0.0200686\\
0.8057	0.0199166\\
0.8082	0.0197625\\
0.8108	0.0195942\\
0.8133	0.0194248\\
0.8159	0.0192405\\
0.8185	0.0190477\\
0.821	0.0188542\\
0.8236	0.0186449\\
0.8262	0.0184276\\
0.8287	0.0182112\\
0.8313	0.0179785\\
0.8339	0.0177377\\
0.8364	0.0174985\\
0.839	0.0172417\\
0.8415	0.0169875\\
0.8441	0.0167156\\
0.8467	0.0164363\\
0.8492	0.0161606\\
0.8518	0.0158662\\
0.8544	0.015564\\
0.8569	0.0152663\\
0.8595	0.0149495\\
0.8621	0.0146256\\
0.8646	0.0143075\\
0.8672	0.0139695\\
0.8697	0.0136375\\
0.8723	0.0132851\\
0.8749	0.0129257\\
0.8774	0.0125738\\
0.88	0.0122012\\
0.8826	0.0118219\\
0.8851	0.0114507\\
0.8877	0.0110579\\
0.8903	0.0106584\\
0.8928	0.0102682\\
0.8954	0.00985631\\
0.8979	0.00945436\\
0.9005	0.00903013\\
0.9031	0.00859952\\
0.9056	0.00817954\\
0.9082	0.00773683\\
0.9108	0.0072883\\
0.9133	0.00685163\\
0.9159	0.00639183\\
0.9185	0.00592619\\
0.921	0.00547299\\
0.9236	0.00499613\\
0.9262	0.00451388\\
0.9287	0.00404525\\
0.9313	0.00355278\\
0.9338	0.0030743\\
0.9364	0.00257152\\
0.939	0.00206357\\
0.9415	0.0015705\\
0.9441	0.00105306\\
0.9467	0.000531001\\
0.9492	2.46455e-05\\
0.9518	-0.00050655\\
0.9544	-0.00104237\\
0.9569	-0.00156179\\
0.9595	-0.00210613\\
0.962	-0.00263337\\
0.9646	-0.00318564\\
0.9672	-0.00374196\\
0.9697	-0.00428068\\
0.9723	-0.0048448\\
0.9749	-0.00541264\\
0.9774	-0.00596195\\
0.98	-0.00653657\\
0.9826	-0.00711457\\
0.9851	-0.00767356\\
0.9877	-0.00825817\\
0.9902	-0.00882329\\
0.9928	-0.00941392\\
0.9954	-0.0100073\\
0.9979	-0.0105805\\
1.0005	-0.0111793\\
1.0031	-0.0117808\\
1.0056	-0.0123616\\
1.0082	-0.012968\\
1.0108	-0.0135765\\
1.0133	-0.0141637\\
1.0159	-0.0147763\\
1.0184	-0.0153673\\
1.021	-0.0159839\\
1.0236	-0.0166022\\
1.0261	-0.0171983\\
1.0287	-0.0178197\\
1.0313	-0.0184425\\
1.0338	-0.0190425\\
1.0364	-0.0196679\\
1.039	-0.0202945\\
1.0415	-0.0208979\\
1.0441	-0.0215263\\
1.0466	-0.0221312\\
1.0492	-0.0227609\\
1.0518	-0.0233913\\
1.0543	-0.023998\\
1.0569	-0.0246293\\
1.0595	-0.0252609\\
1.062	-0.0258682\\
1.0646	-0.0264999\\
1.0672	-0.0271315\\
1.0697	-0.0277388\\
1.0723	-0.0283701\\
1.0748	-0.0289768\\
1.0774	-0.0296073\\
1.08	-0.0302371\\
1.0825	-0.030842\\
1.0851	-0.0314704\\
1.0877	-0.0320979\\
1.0902	-0.0327005\\
1.0928	-0.033326\\
1.0954	-0.0339503\\
1.0979	-0.0345492\\
1.1005	-0.0351708\\
1.1031	-0.035791\\
1.1056	-0.0363858\\
1.1082	-0.0370027\\
1.1107	-0.0375942\\
1.1133	-0.0382074\\
1.1159	-0.0388186\\
1.1184	-0.0394043\\
1.121	-0.0400114\\
1.1236	-0.0406162\\
1.1261	-0.0411955\\
1.1287	-0.0417955\\
1.1313	-0.0423928\\
1.1338	-0.0429647\\
1.1364	-0.0435568\\
1.1389	-0.0441235\\
1.1415	-0.0447099\\
1.1441	-0.0452933\\
1.1466	-0.0458513\\
1.1492	-0.0464284\\
1.1518	-0.0470023\\
1.1543	-0.0475509\\
1.1569	-0.0481181\\
1.1595	-0.0486818\\
1.162	-0.0492203\\
1.1646	-0.0497766\\
1.1671	-0.050308\\
1.1697	-0.0508569\\
1.1723	-0.0514019\\
1.1748	-0.0519221\\
1.1774	-0.052459\\
1.18	-0.0529917\\
1.1825	-0.0534999\\
1.1851	-0.0540242\\
1.1877	-0.0545441\\
1.1902	-0.0550398\\
1.1928	-0.0555509\\
1.1953	-0.0560379\\
1.1979	-0.0565398\\
1.2005	-0.0570369\\
1.203	-0.0575104\\
1.2056	-0.057998\\
1.2082	-0.0584808\\
1.2107	-0.0589402\\
1.2133	-0.059413\\
1.2159	-0.0598806\\
1.2184	-0.0603254\\
1.221	-0.0607829\\
1.2235	-0.0612178\\
1.2261	-0.0616649\\
1.2287	-0.0621065\\
1.2312	-0.062526\\
1.2338	-0.0629569\\
1.2364	-0.0633823\\
1.2389	-0.0637861\\
1.2415	-0.0642005\\
1.2441	-0.0646092\\
1.2466	-0.0649968\\
1.2492	-0.0653942\\
1.2518	-0.0657858\\
1.2543	-0.066157\\
1.2569	-0.0665372\\
1.2594	-0.0668972\\
1.262	-0.0672658\\
1.2646	-0.0676285\\
1.2671	-0.0679715\\
1.2697	-0.0683224\\
1.2723	-0.0686673\\
1.2748	-0.0689932\\
1.2774	-0.0693262\\
1.28	-0.069653\\
1.2825	-0.0699615\\
1.2851	-0.0702763\\
1.2876	-0.0705733\\
1.2902	-0.0708761\\
1.2928	-0.0711727\\
1.2953	-0.0714521\\
1.2979	-0.0717365\\
1.3005	-0.0720148\\
1.303	-0.0722765\\
1.3056	-0.0725427\\
1.3082	-0.0728026\\
1.3107	-0.0730467\\
1.3133	-0.0732944\\
1.3158	-0.0735268\\
1.3184	-0.0737624\\
1.321	-0.0739918\\
1.3235	-0.0742066\\
1.3261	-0.0744239\\
1.3287	-0.0746351\\
1.3312	-0.0748323\\
1.3338	-0.0750314\\
1.3364	-0.0752244\\
1.3389	-0.0754043\\
1.3415	-0.0755854\\
1.344	-0.0757539\\
1.3466	-0.0759231\\
1.3492	-0.0760864\\
1.3517	-0.0762377\\
1.3543	-0.0763893\\
1.3569	-0.076535\\
1.3594	-0.0766695\\
1.362	-0.0768037\\
1.3646	-0.076932\\
1.3671	-0.0770499\\
1.3697	-0.077167\\
1.3722	-0.0772741\\
1.3748	-0.07738\\
1.3774	-0.0774802\\
1.3799	-0.0775713\\
1.3825	-0.0776607\\
1.3851	-0.0777445\\
1.3876	-0.0778199\\
1.3902	-0.0778931\\
1.3928	-0.0779609\\
1.3953	-0.0780211\\
1.3979	-0.0780786\\
1.4005	-0.0781308\\
1.403	-0.0781762\\
1.4056	-0.0782183\\
1.4081	-0.0782542\\
1.4107	-0.0782865\\
1.4133	-0.0783139\\
1.4158	-0.0783356\\
1.4184	-0.0783535\\
1.421	-0.0783666\\
1.4235	-0.0783748\\
1.4261	-0.0783788\\
1.4287	-0.0783782\\
1.4312	-0.0783733\\
1.4338	-0.0783639\\
1.4363	-0.0783507\\
1.4389	-0.0783328\\
1.4415	-0.0783106\\
1.444	-0.0782853\\
1.4466	-0.0782549\\
1.4492	-0.0782205\\
1.4517	-0.0781837\\
1.4543	-0.0781415\\
1.4569	-0.0780955\\
1.4594	-0.0780476\\
1.462	-0.0779943\\
1.4645	-0.0779395\\
1.4671	-0.077879\\
1.4697	-0.077815\\
1.4722	-0.0777503\\
1.4748	-0.0776797\\
1.4774	-0.0776057\\
1.4799	-0.0775316\\
1.4825	-0.0774514\\
1.4851	-0.0773682\\
1.4876	-0.0772853\\
1.4902	-0.0771962\\
1.4927	-0.0771079\\
1.4953	-0.0770133\\
1.4979	-0.0769159\\
1.5004	-0.0768198\\
1.503	-0.0767174\\
1.5056	-0.0766124\\
1.5081	-0.0765092\\
1.5107	-0.0763996\\
1.5133	-0.0762876\\
1.5158	-0.0761779\\
1.5184	-0.0760617\\
1.5209	-0.075948\\
1.5235	-0.0758278\\
1.5261	-0.0757056\\
1.5286	-0.0755864\\
1.5312	-0.0754607\\
1.5338	-0.0753332\\
1.5363	-0.0752091\\
1.5389	-0.0750784\\
1.5415	-0.0749463\\
1.544	-0.0748178\\
1.5466	-0.0746829\\
1.5491	-0.0745519\\
1.5517	-0.0744144\\
1.5543	-0.0742758\\
1.5568	-0.0741414\\
1.5594	-0.0740007\\
1.562	-0.0738589\\
1.5645	-0.0737217\\
1.5671	-0.0735781\\
1.5697	-0.0734338\\
1.5722	-0.0732943\\
1.5748	-0.0731486\\
1.5774	-0.0730023\\
1.5799	-0.072861\\
1.5825	-0.0727136\\
1.585	-0.0725715\\
1.5876	-0.0724232\\
1.5902	-0.0722747\\
1.5927	-0.0721316\\
1.5953	-0.0719826\\
1.5979	-0.0718334\\
1.6004	-0.0716898\\
1.603	-0.0715405\\
1.6056	-0.0713911\\
1.6081	-0.0712476\\
1.6107	-0.0710983\\
1.6132	-0.070955\\
1.6158	-0.0708061\\
1.6184	-0.0706575\\
1.6209	-0.0705148\\
1.6235	-0.0703668\\
1.6261	-0.0702191\\
1.6286	-0.0700775\\
1.6312	-0.0699308\\
1.6338	-0.0697845\\
1.6363	-0.0696444\\
1.6389	-0.0694993\\
1.6414	-0.0693603\\
1.644	-0.0692165\\
1.6466	-0.0690733\\
1.6491	-0.0689364\\
1.6517	-0.0687948\\
1.6543	-0.0686539\\
1.6568	-0.0685194\\
1.6594	-0.0683803\\
1.662	-0.0682421\\
1.6645	-0.0681102\\
1.6671	-0.0679739\\
1.6696	-0.0678439\\
1.6722	-0.0677096\\
1.6748	-0.0675765\\
1.6773	-0.0674495\\
1.6799	-0.0673186\\
1.6825	-0.0671888\\
1.685	-0.0670651\\
1.6876	-0.0669377\\
1.6902	-0.0668115\\
1.6927	-0.0666913\\
1.6953	-0.0665676\\
1.6978	-0.0664498\\
1.7004	-0.0663287\\
1.703	-0.0662089\\
1.7055	-0.0660949\\
1.7081	-0.0659778\\
1.7107	-0.0658621\\
1.7132	-0.0657521\\
1.7158	-0.0656391\\
1.7184	-0.0655276\\
1.7209	-0.0654217\\
1.7235	-0.0653131\\
1.7261	-0.0652059\\
1.7286	-0.0651042\\
1.7312	-0.065\\
1.7337	-0.0649012\\
1.7363	-0.0647999\\
1.7389	-0.0647001\\
1.7414	-0.0646057\\
1.744	-0.0645089\\
1.7466	-0.0644137\\
1.7491	-0.0643237\\
1.7517	-0.0642315\\
1.7543	-0.064141\\
1.7568	-0.0640554\\
1.7594	-0.0639679\\
1.7619	-0.0638853\\
1.7645	-0.0638009\\
1.7671	-0.0637181\\
1.7696	-0.0636399\\
1.7722	-0.0635602\\
1.7748	-0.0634821\\
1.7773	-0.0634084\\
1.7799	-0.0633334\\
1.7825	-0.0632599\\
1.785	-0.0631908\\
1.7876	-0.0631204\\
1.7901	-0.0630542\\
1.7927	-0.0629868\\
1.7953	-0.0629211\\
1.7978	-0.0628593\\
1.8004	-0.0627965\\
1.803	-0.0627353\\
1.8055	-0.0626779\\
1.8081	-0.0626197\\
1.8107	-0.0625631\\
1.8132	-0.06251\\
1.8158	-0.0624562\\
1.8183	-0.062406\\
1.8209	-0.0623552\\
1.8235	-0.0623058\\
1.826	-0.0622597\\
1.8286	-0.0622132\\
1.8312	-0.0621682\\
1.8337	-0.0621262\\
1.8363	-0.062084\\
1.8389	-0.0620431\\
1.8414	-0.0620052\\
1.844	-0.0619671\\
1.8465	-0.0619317\\
1.8491	-0.0618962\\
1.8517	-0.0618621\\
1.8542	-0.0618306\\
1.8568	-0.0617991\\
1.8594	-0.0617689\\
1.8619	-0.0617411\\
1.8645	-0.0617134\\
1.8671	-0.061687\\
1.8696	-0.0616627\\
1.8722	-0.0616387\\
1.8747	-0.0616168\\
1.8773	-0.0615952\\
1.8799	-0.0615747\\
1.8824	-0.0615562\\
1.885	-0.061538\\
1.8876	-0.061521\\
1.8901	-0.0615056\\
1.8927	-0.0614908\\
1.8953	-0.061477\\
1.8978	-0.0614648\\
1.9004	-0.0614531\\
1.903	-0.0614425\\
1.9055	-0.0614332\\
1.9081	-0.0614245\\
1.9106	-0.0614171\\
1.9132	-0.0614104\\
1.9158	-0.0614046\\
1.9183	-0.0613999\\
1.9209	-0.0613959\\
1.9235	-0.0613929\\
1.926	-0.0613908\\
1.9286	-0.0613894\\
1.9312	-0.0613889\\
1.9337	-0.0613892\\
1.9363	-0.0613903\\
1.9388	-0.0613921\\
1.9414	-0.0613948\\
1.944	-0.0613982\\
1.9465	-0.0614021\\
1.9491	-0.061407\\
1.9517	-0.0614126\\
1.9542	-0.0614186\\
1.9568	-0.0614255\\
1.9594	-0.061433\\
1.9619	-0.0614409\\
1.9645	-0.0614497\\
1.967	-0.0614588\\
1.9696	-0.0614688\\
1.9722	-0.0614793\\
1.9747	-0.06149\\
1.9773	-0.0615016\\
1.9799	-0.0615138\\
1.9824	-0.0615259\\
1.985	-0.0615391\\
1.9876	-0.0615527\\
1.9901	-0.0615662\\
1.9927	-0.0615807\\
1.9952	-0.061595\\
1.9978	-0.0616104\\
2.0004	-0.0616261\\
2.0029	-0.0616415\\
2.0055	-0.061658\\
2.0081	-0.0616748\\
2.0106	-0.0616912\\
2.0132	-0.0617086\\
2.0158	-0.0617264\\
2.0183	-0.0617437\\
2.0209	-0.061762\\
2.0234	-0.0617798\\
2.026	-0.0617986\\
2.0286	-0.0618176\\
2.0311	-0.0618361\\
2.0337	-0.0618555\\
2.0363	-0.0618751\\
2.0388	-0.0618941\\
2.0414	-0.0619141\\
2.044	-0.0619342\\
2.0465	-0.0619536\\
2.0491	-0.061974\\
2.0517	-0.0619944\\
2.0542	-0.0620142\\
2.0568	-0.0620348\\
2.0593	-0.0620548\\
2.0619	-0.0620755\\
2.0645	-0.0620964\\
2.067	-0.0621164\\
2.0696	-0.0621373\\
2.0722	-0.0621583\\
2.0747	-0.0621784\\
2.0773	-0.0621993\\
2.0799	-0.0622202\\
2.0824	-0.0622403\\
2.085	-0.0622611\\
2.0875	-0.0622811\\
2.0901	-0.0623019\\
2.0927	-0.0623225\\
2.0952	-0.0623424\\
2.0978	-0.0623629\\
2.1004	-0.0623833\\
2.1029	-0.0624029\\
2.1055	-0.0624231\\
2.1081	-0.0624432\\
2.1106	-0.0624624\\
2.1132	-0.0624823\\
2.1157	-0.0625012\\
2.1183	-0.0625208\\
2.1209	-0.0625402\\
2.1234	-0.0625587\\
2.126	-0.0625778\\
2.1286	-0.0625967\\
2.1311	-0.0626147\\
2.1337	-0.0626332\\
2.1363	-0.0626515\\
2.1388	-0.062669\\
2.1414	-0.0626869\\
2.1439	-0.062704\\
2.1465	-0.0627215\\
2.1491	-0.0627387\\
2.1516	-0.0627551\\
2.1542	-0.062772\\
2.1568	-0.0627885\\
2.1593	-0.0628042\\
2.1619	-0.0628203\\
2.1645	-0.0628362\\
2.167	-0.0628512\\
2.1696	-0.0628665\\
2.1721	-0.062881\\
2.1747	-0.0628959\\
2.1773	-0.0629104\\
2.1798	-0.0629241\\
2.1824	-0.0629381\\
2.185	-0.0629519\\
2.1875	-0.0629648\\
2.1901	-0.062978\\
2.1927	-0.0629909\\
2.1952	-0.063003\\
2.1978	-0.0630154\\
2.2004	-0.0630274\\
2.2029	-0.0630387\\
2.2055	-0.0630502\\
2.208	-0.0630609\\
2.2106	-0.0630718\\
2.2132	-0.0630824\\
2.2157	-0.0630923\\
2.2183	-0.0631023\\
2.2209	-0.0631121\\
2.2234	-0.0631211\\
2.226	-0.0631303\\
2.2286	-0.0631391\\
2.2311	-0.0631473\\
2.2337	-0.0631556\\
2.2362	-0.0631633\\
2.2388	-0.063171\\
2.2414	-0.0631784\\
2.2439	-0.0631852\\
2.2465	-0.063192\\
2.2491	-0.0631986\\
2.2516	-0.0632046\\
2.2542	-0.0632106\\
2.2568	-0.0632163\\
2.2593	-0.0632215\\
2.2619	-0.0632266\\
2.2644	-0.0632313\\
2.267	-0.0632359\\
2.2696	-0.0632403\\
2.2721	-0.0632442\\
2.2747	-0.063248\\
2.2773	-0.0632515\\
2.2798	-0.0632547\\
2.2824	-0.0632577\\
2.285	-0.0632605\\
2.2875	-0.0632629\\
2.2901	-0.0632652\\
2.2926	-0.0632672\\
2.2952	-0.063269\\
2.2978	-0.0632706\\
2.3003	-0.0632718\\
2.3029	-0.063273\\
2.3055	-0.0632738\\
2.308	-0.0632745\\
2.3106	-0.0632749\\
2.3132	-0.0632752\\
2.3157	-0.0632752\\
2.3183	-0.063275\\
2.3208	-0.0632747\\
2.3234	-0.0632741\\
2.326	-0.0632733\\
2.3285	-0.0632724\\
2.3311	-0.0632713\\
2.3337	-0.06327\\
2.3362	-0.0632686\\
2.3388	-0.0632669\\
2.3414	-0.0632651\\
2.3439	-0.0632632\\
2.3465	-0.0632611\\
2.349	-0.0632589\\
2.3516	-0.0632565\\
2.3542	-0.063254\\
2.3567	-0.0632514\\
2.3593	-0.0632486\\
2.3619	-0.0632456\\
2.3644	-0.0632427\\
2.367	-0.0632395\\
2.3696	-0.0632362\\
2.3721	-0.0632329\\
2.3747	-0.0632294\\
2.3773	-0.0632258\\
2.3798	-0.0632223\\
2.3824	-0.0632186\\
2.3849	-0.0632149\\
2.3875	-0.0632109\\
2.3901	-0.063207\\
2.3926	-0.0632031\\
2.3952	-0.063199\\
2.3978	-0.0631948\\
2.4003	-0.0631908\\
2.4029	-0.0631865\\
2.4055	-0.0631822\\
2.408	-0.0631781\\
2.4106	-0.0631737\\
2.4131	-0.0631695\\
2.4157	-0.0631651\\
2.4183	-0.0631607\\
2.4208	-0.0631565\\
2.4234	-0.0631521\\
2.426	-0.0631477\\
2.4285	-0.0631435\\
2.4311	-0.0631391\\
2.4337	-0.0631347\\
2.4362	-0.0631305\\
2.4388	-0.0631262\\
2.4413	-0.063122\\
2.4439	-0.0631178\\
2.4465	-0.0631136\\
2.449	-0.0631095\\
2.4516	-0.0631054\\
2.4542	-0.0631013\\
2.4567	-0.0630975\\
2.4593	-0.0630935\\
2.4619	-0.0630896\\
2.4644	-0.0630859\\
2.467	-0.0630822\\
2.4695	-0.0630787\\
2.4721	-0.0630751\\
2.4747	-0.0630716\\
2.4772	-0.0630683\\
2.4798	-0.063065\\
2.4824	-0.0630618\\
2.4849	-0.0630588\\
2.4875	-0.0630557\\
2.4901	-0.0630528\\
2.4926	-0.0630502\\
2.4952	-0.0630475\\
2.4977	-0.063045\\
2.5003	-0.0630426\\
2.5029	-0.0630403\\
2.5054	-0.0630382\\
2.508	-0.0630361\\
2.5106	-0.0630342\\
2.5131	-0.0630325\\
2.5157	-0.0630308\\
2.5183	-0.0630293\\
2.5208	-0.063028\\
2.5234	-0.0630268\\
2.526	-0.0630258\\
2.5285	-0.0630249\\
2.5311	-0.0630241\\
2.5336	-0.0630235\\
2.5362	-0.0630231\\
2.5388	-0.0630228\\
2.5413	-0.0630227\\
2.5439	-0.0630227\\
2.5465	-0.0630229\\
2.549	-0.0630233\\
2.5516	-0.0630238\\
2.5542	-0.0630245\\
2.5567	-0.0630254\\
2.5593	-0.0630264\\
2.5618	-0.0630276\\
2.5644	-0.063029\\
2.567	-0.0630305\\
2.5695	-0.0630322\\
2.5721	-0.0630341\\
2.5747	-0.0630362\\
2.5772	-0.0630384\\
2.5798	-0.0630408\\
2.5824	-0.0630434\\
2.5849	-0.0630461\\
2.5875	-0.0630491\\
2.59	-0.0630521\\
2.5926	-0.0630554\\
2.5952	-0.0630589\\
2.5977	-0.0630625\\
2.6003	-0.0630663\\
2.6029	-0.0630703\\
2.6054	-0.0630744\\
2.608	-0.0630787\\
2.6106	-0.0630833\\
2.6131	-0.0630878\\
2.6157	-0.0630927\\
2.6182	-0.0630975\\
2.6208	-0.0631028\\
2.6234	-0.0631081\\
2.6259	-0.0631135\\
2.6285	-0.0631192\\
2.6311	-0.0631251\\
2.6336	-0.0631309\\
2.6362	-0.0631371\\
2.6388	-0.0631434\\
2.6413	-0.0631497\\
2.6439	-0.0631563\\
2.6464	-0.0631629\\
2.649	-0.0631698\\
2.6516	-0.0631769\\
2.6541	-0.0631839\\
2.6567	-0.0631913\\
2.6593	-0.0631988\\
2.6618	-0.0632061\\
2.6644	-0.0632139\\
2.667	-0.0632218\\
2.6695	-0.0632296\\
2.6721	-0.0632377\\
2.6746	-0.0632457\\
2.6772	-0.0632541\\
2.6798	-0.0632626\\
2.6823	-0.0632709\\
2.6849	-0.0632797\\
2.6875	-0.0632885\\
2.69	-0.0632971\\
2.6926	-0.0633062\\
2.6952	-0.0633153\\
2.6977	-0.0633242\\
2.7003	-0.0633335\\
2.7029	-0.0633429\\
2.7054	-0.063352\\
2.708	-0.0633616\\
2.7105	-0.0633708\\
2.7131	-0.0633805\\
2.7157	-0.0633903\\
2.7182	-0.0633997\\
2.7208	-0.0634096\\
2.7234	-0.0634195\\
2.7259	-0.0634291\\
2.7285	-0.0634391\\
2.7311	-0.0634491\\
2.7336	-0.0634588\\
2.7362	-0.0634689\\
2.7387	-0.0634786\\
2.7413	-0.0634887\\
2.7439	-0.0634989\\
2.7464	-0.0635087\\
2.749	-0.0635189\\
2.7516	-0.063529\\
2.7541	-0.0635388\\
2.7567	-0.063549\\
2.7593	-0.0635591\\
2.7618	-0.0635688\\
2.7644	-0.0635789\\
2.7669	-0.0635886\\
2.7695	-0.0635987\\
2.7721	-0.0636087\\
2.7746	-0.0636182\\
2.7772	-0.0636281\\
2.7798	-0.063638\\
2.7823	-0.0636474\\
2.7849	-0.0636572\\
2.7875	-0.0636668\\
2.79	-0.063676\\
2.7926	-0.0636856\\
2.7951	-0.0636946\\
2.7977	-0.063704\\
2.8003	-0.0637132\\
2.8028	-0.063722\\
2.8054	-0.0637311\\
2.808	-0.06374\\
2.8105	-0.0637485\\
2.8131	-0.0637572\\
2.8157	-0.0637658\\
2.8182	-0.0637739\\
2.8208	-0.0637822\\
2.8233	-0.0637901\\
2.8259	-0.0637981\\
2.8285	-0.063806\\
2.831	-0.0638134\\
2.8336	-0.063821\\
2.8362	-0.0638284\\
2.8387	-0.0638354\\
2.8413	-0.0638425\\
2.8439	-0.0638494\\
2.8464	-0.0638558\\
2.849	-0.0638623\\
2.8516	-0.0638687\\
2.8541	-0.0638746\\
2.8567	-0.0638805\\
2.8592	-0.063886\\
2.8618	-0.0638915\\
2.8644	-0.0638968\\
2.8669	-0.0639017\\
2.8695	-0.0639065\\
2.8721	-0.0639111\\
2.8746	-0.0639153\\
2.8772	-0.0639195\\
2.8798	-0.0639234\\
2.8823	-0.0639269\\
2.8849	-0.0639303\\
2.8874	-0.0639333\\
2.89	-0.0639362\\
2.8926	-0.0639388\\
2.8951	-0.0639411\\
2.8977	-0.0639432\\
2.9003	-0.063945\\
2.9028	-0.0639465\\
2.9054	-0.0639477\\
2.908	-0.0639487\\
2.9105	-0.0639493\\
2.9131	-0.0639497\\
2.9156	-0.0639497\\
2.9182	-0.0639495\\
2.9208	-0.063949\\
2.9233	-0.0639481\\
2.9259	-0.063947\\
2.9285	-0.0639455\\
2.931	-0.0639438\\
2.9336	-0.0639416\\
2.9362	-0.0639392\\
2.9387	-0.0639365\\
2.9413	-0.0639334\\
2.9438	-0.0639301\\
2.9464	-0.0639263\\
2.949	-0.0639222\\
2.9515	-0.0639179\\
2.9541	-0.0639131\\
2.9567	-0.0639079\\
2.9592	-0.0639026\\
2.9618	-0.0638967\\
2.9644	-0.0638905\\
2.9669	-0.0638842\\
2.9695	-0.0638772\\
2.972	-0.0638702\\
2.9746	-0.0638625\\
2.9772	-0.0638545\\
2.9797	-0.0638464\\
2.9823	-0.0638377\\
2.9849	-0.0638285\\
2.9874	-0.0638193\\
2.99	-0.0638095\\
2.9926	-0.0637992\\
2.9951	-0.0637889\\
2.9977	-0.0637779\\
3.0003	-0.0637665\\
3.0028	-0.0637551\\
3.0054	-0.063743\\
3.0079	-0.0637309\\
3.0105	-0.0637179\\
3.0131	-0.0637046\\
3.0156	-0.0636914\\
3.0182	-0.0636773\\
3.0208	-0.0636628\\
3.0233	-0.0636485\\
3.0259	-0.0636332\\
3.0285	-0.0636175\\
3.031	-0.0636021\\
3.0336	-0.0635857\\
3.0361	-0.0635695\\
3.0387	-0.0635523\\
3.0413	-0.0635346\\
3.0438	-0.0635173\\
3.0464	-0.063499\\
3.049	-0.0634802\\
3.0515	-0.0634617\\
3.0541	-0.0634422\\
3.0567	-0.0634222\\
3.0592	-0.0634027\\
3.0618	-0.063382\\
3.0643	-0.0633617\\
3.0669	-0.0633402\\
3.0695	-0.0633183\\
3.072	-0.063297\\
3.0746	-0.0632743\\
3.0772	-0.0632513\\
3.0797	-0.0632288\\
3.0823	-0.0632051\\
3.0849	-0.0631809\\
3.0874	-0.0631573\\
3.09	-0.0631325\\
3.0925	-0.0631082\\
3.0951	-0.0630825\\
3.0977	-0.0630565\\
3.1002	-0.0630312\\
3.1028	-0.0630045\\
3.1054	-0.0629774\\
3.1079	-0.062951\\
3.1105	-0.0629232\\
3.1131	-0.0628951\\
3.1156	-0.0628676\\
3.1182	-0.0628388\\
3.1207	-0.0628107\\
3.1233	-0.0627812\\
3.1259	-0.0627513\\
3.1284	-0.0627222\\
3.131	-0.0626917\\
3.1336	-0.0626608\\
3.1361	-0.0626307\\
3.1387	-0.0625992\\
3.1413	-0.0625673\\
3.1438	-0.0625363\\
3.1464	-0.0625038\\
3.1489	-0.0624723\\
3.1515	-0.0624391\\
3.1541	-0.0624057\\
3.1566	-0.0623732\\
3.1592	-0.0623392\\
3.1618	-0.0623048\\
3.1643	-0.0622715\\
3.1669	-0.0622366\\
3.1695	-0.0622013\\
3.172	-0.0621672\\
3.1746	-0.0621314\\
3.1772	-0.0620953\\
3.1797	-0.0620604\\
3.1823	-0.0620238\\
3.1848	-0.0619884\\
3.1874	-0.0619513\\
3.19	-0.0619139\\
3.1925	-0.0618777\\
3.1951	-0.0618398\\
3.1977	-0.0618017\\
3.2002	-0.0617648\\
3.2028	-0.0617263\\
3.2054	-0.0616874\\
3.2079	-0.0616499\\
3.2105	-0.0616106\\
3.213	-0.0615727\\
3.2156	-0.061533\\
3.2182	-0.0614931\\
3.2207	-0.0614546\\
3.2233	-0.0614143\\
3.2259	-0.0613738\\
3.2284	-0.0613347\\
3.231	-0.0612938\\
3.2336	-0.0612528\\
3.2361	-0.0612132\\
3.2387	-0.0611718\\
3.2412	-0.0611319\\
3.2438	-0.0610902\\
3.2464	-0.0610483\\
3.2489	-0.0610079\\
3.2515	-0.0609658\\
3.2541	-0.0609235\\
3.2566	-0.0608827\\
3.2592	-0.0608401\\
3.2618	-0.0607975\\
3.2643	-0.0607563\\
3.2669	-0.0607134\\
3.2694	-0.060672\\
3.272	-0.0606289\\
3.2746	-0.0605857\\
3.2771	-0.060544\\
3.2797	-0.0605006\\
3.2823	-0.0604571\\
3.2848	-0.0604153\\
3.2874	-0.0603716\\
3.29	-0.0603279\\
3.2925	-0.0602858\\
3.2951	-0.060242\\
3.2976	-0.0601998\\
3.3002	-0.0601558\\
3.3028	-0.0601118\\
3.3053	-0.0600695\\
3.3079	-0.0600254\\
3.3105	-0.0599814\\
3.313	-0.0599389\\
3.3156	-0.0598948\\
3.3182	-0.0598507\\
3.3207	-0.0598082\\
3.3233	-0.0597641\\
3.3259	-0.0597199\\
3.3284	-0.0596775\\
3.331	-0.0596333\\
3.3335	-0.0595909\\
3.3361	-0.0595468\\
3.3387	-0.0595027\\
3.3412	-0.0594604\\
3.3438	-0.0594164\\
3.3464	-0.0593724\\
3.3489	-0.0593301\\
3.3515	-0.0592863\\
3.3541	-0.0592425\\
3.3566	-0.0592004\\
3.3592	-0.0591567\\
3.3617	-0.0591147\\
3.3643	-0.0590712\\
3.3669	-0.0590277\\
3.3694	-0.058986\\
3.372	-0.0589427\\
3.3746	-0.0588995\\
3.3771	-0.058858\\
3.3797	-0.058815\\
3.3823	-0.0587721\\
3.3848	-0.0587309\\
3.3874	-0.0586882\\
3.3899	-0.0586473\\
3.3925	-0.0586048\\
3.3951	-0.0585625\\
3.3976	-0.0585219\\
3.4002	-0.0584799\\
3.4028	-0.0584379\\
3.4053	-0.0583978\\
3.4079	-0.0583561\\
3.4105	-0.0583146\\
3.413	-0.0582749\\
3.4156	-0.0582337\\
3.4181	-0.0581943\\
3.4207	-0.0581535\\
3.4233	-0.0581128\\
3.4258	-0.0580739\\
3.4284	-0.0580335\\
3.431	-0.0579934\\
3.4335	-0.057955\\
3.4361	-0.0579152\\
3.4387	-0.0578756\\
3.4412	-0.0578378\\
3.4438	-0.0577986\\
3.4463	-0.0577611\\
3.4489	-0.0577223\\
3.4515	-0.0576838\\
3.454	-0.0576469\\
3.4566	-0.0576087\\
3.4592	-0.0575708\\
3.4617	-0.0575346\\
3.4643	-0.0574971\\
3.4669	-0.0574598\\
3.4694	-0.0574242\\
3.472	-0.0573875\\
3.4745	-0.0573523\\
3.4771	-0.057316\\
3.4797	-0.0572799\\
3.4822	-0.0572454\\
3.4848	-0.0572098\\
3.4874	-0.0571745\\
3.4899	-0.0571408\\
3.4925	-0.0571059\\
3.4951	-0.0570713\\
3.4976	-0.0570383\\
3.5002	-0.0570043\\
3.5028	-0.0569704\\
3.5053	-0.0569382\\
3.5079	-0.0569049\\
3.5104	-0.0568732\\
3.513	-0.0568404\\
3.5156	-0.056808\\
3.5181	-0.056777\\
3.5207	-0.0567451\\
3.5233	-0.0567134\\
3.5258	-0.0566833\\
3.5284	-0.0566522\\
3.531	-0.0566214\\
3.5335	-0.056592\\
3.5361	-0.0565618\\
3.5386	-0.056533\\
3.5412	-0.0565033\\
3.5438	-0.0564739\\
3.5463	-0.056446\\
3.5489	-0.0564172\\
3.5515	-0.0563886\\
3.554	-0.0563615\\
3.5566	-0.0563336\\
3.5592	-0.056306\\
3.5617	-0.0562797\\
3.5643	-0.0562527\\
3.5668	-0.0562269\\
3.5694	-0.0562005\\
3.572	-0.0561744\\
3.5745	-0.0561495\\
3.5771	-0.056124\\
3.5797	-0.0560987\\
3.5822	-0.0560747\\
3.5848	-0.0560501\\
3.5874	-0.0560258\\
3.5899	-0.0560027\\
3.5925	-0.0559789\\
3.595	-0.0559564\\
3.5976	-0.0559332\\
3.6002	-0.0559104\\
3.6027	-0.0558888\\
3.6053	-0.0558665\\
3.6079	-0.0558446\\
3.6104	-0.0558238\\
3.613	-0.0558025\\
3.6156	-0.0557815\\
3.6181	-0.0557616\\
3.6207	-0.0557412\\
3.6232	-0.0557219\\
3.6258	-0.055702\\
3.6284	-0.0556825\\
3.6309	-0.0556641\\
3.6335	-0.0556452\\
3.6361	-0.0556266\\
3.6386	-0.055609\\
3.6412	-0.055591\\
3.6438	-0.0555732\\
3.6463	-0.0555565\\
3.6489	-0.0555394\\
3.6515	-0.0555226\\
3.654	-0.0555067\\
3.6566	-0.0554905\\
3.6591	-0.0554751\\
3.6617	-0.0554595\\
3.6643	-0.0554441\\
3.6668	-0.0554297\\
3.6694	-0.0554149\\
3.672	-0.0554004\\
3.6745	-0.0553867\\
3.6771	-0.0553728\\
3.6797	-0.0553592\\
3.6822	-0.0553464\\
3.6848	-0.0553333\\
3.6873	-0.055321\\
3.6899	-0.0553085\\
3.6925	-0.0552963\\
3.695	-0.0552848\\
3.6976	-0.0552732\\
3.7002	-0.0552618\\
3.7027	-0.0552511\\
3.7053	-0.0552403\\
3.7079	-0.0552297\\
3.7104	-0.0552198\\
3.713	-0.0552098\\
3.7155	-0.0552004\\
3.7181	-0.0551909\\
3.7207	-0.0551816\\
3.7232	-0.055173\\
3.7258	-0.0551643\\
3.7284	-0.0551558\\
3.7309	-0.0551479\\
3.7335	-0.05514\\
3.7361	-0.0551323\\
3.7386	-0.0551251\\
3.7412	-0.0551179\\
3.7437	-0.0551113\\
3.7463	-0.0551046\\
3.7489	-0.0550981\\
3.7514	-0.0550921\\
3.754	-0.0550862\\
3.7566	-0.0550805\\
3.7591	-0.0550752\\
3.7617	-0.0550699\\
3.7643	-0.0550649\\
3.7668	-0.0550603\\
3.7694	-0.0550557\\
3.7719	-0.0550516\\
3.7745	-0.0550475\\
3.7771	-0.0550436\\
3.7796	-0.0550401\\
3.7822	-0.0550366\\
3.7848	-0.0550334\\
3.7873	-0.0550306\\
3.7899	-0.0550278\\
3.7925	-0.0550252\\
3.795	-0.0550229\\
3.7976	-0.0550208\\
3.8002	-0.0550188\\
3.8027	-0.0550172\\
3.8053	-0.0550156\\
3.8078	-0.0550143\\
3.8104	-0.0550132\\
3.813	-0.0550123\\
3.8155	-0.0550116\\
3.8181	-0.055011\\
3.8207	-0.0550107\\
3.8232	-0.0550105\\
3.8258	-0.0550105\\
3.8284	-0.0550107\\
3.8309	-0.0550111\\
3.8335	-0.0550117\\
3.836	-0.0550124\\
3.8386	-0.0550133\\
3.8412	-0.0550144\\
3.8437	-0.0550156\\
3.8463	-0.0550171\\
3.8489	-0.0550187\\
3.8514	-0.0550204\\
3.854	-0.0550223\\
3.8566	-0.0550244\\
3.8591	-0.0550266\\
3.8617	-0.0550291\\
3.8642	-0.0550315\\
3.8668	-0.0550343\\
3.8694	-0.0550372\\
3.8719	-0.0550401\\
3.8745	-0.0550433\\
3.8771	-0.0550467\\
3.8796	-0.05505\\
3.8822	-0.0550536\\
3.8848	-0.0550574\\
3.8873	-0.0550612\\
3.8899	-0.0550653\\
3.8924	-0.0550693\\
3.895	-0.0550736\\
3.8976	-0.0550781\\
3.9001	-0.0550825\\
3.9027	-0.0550872\\
3.9053	-0.0550921\\
3.9078	-0.0550968\\
3.9104	-0.0551019\\
3.913	-0.0551072\\
3.9155	-0.0551123\\
3.9181	-0.0551178\\
3.9206	-0.0551231\\
3.9232	-0.0551288\\
3.9258	-0.0551346\\
3.9283	-0.0551403\\
3.9309	-0.0551463\\
3.9335	-0.0551525\\
3.936	-0.0551585\\
3.9386	-0.0551648\\
3.9412	-0.0551713\\
3.9437	-0.0551776\\
3.9463	-0.0551842\\
3.9488	-0.0551907\\
3.9514	-0.0551976\\
3.954	-0.0552045\\
3.9565	-0.0552113\\
3.9591	-0.0552184\\
3.9617	-0.0552256\\
3.9642	-0.0552326\\
3.9668	-0.05524\\
3.9694	-0.0552475\\
3.9719	-0.0552548\\
3.9745	-0.0552625\\
3.9771	-0.0552702\\
3.9796	-0.0552777\\
3.9822	-0.0552856\\
3.9847	-0.0552933\\
3.9873	-0.0553013\\
3.9899	-0.0553094\\
3.9924	-0.0553173\\
3.995	-0.0553256\\
3.9976	-0.0553339\\
4.0001	-0.055342\\
4.0027	-0.0553505\\
4.0053	-0.055359\\
4.0078	-0.0553673\\
4.0104	-0.055376\\
4.0129	-0.0553844\\
4.0155	-0.0553932\\
4.0181	-0.055402\\
4.0206	-0.0554106\\
4.0232	-0.0554196\\
4.0258	-0.0554286\\
4.0283	-0.0554374\\
4.0309	-0.0554465\\
4.0335	-0.0554557\\
4.036	-0.0554646\\
4.0386	-0.0554739\\
4.0411	-0.0554829\\
4.0437	-0.0554924\\
4.0463	-0.0555018\\
4.0488	-0.055511\\
4.0514	-0.0555205\\
4.054	-0.0555301\\
4.0565	-0.0555394\\
4.0591	-0.0555491\\
4.0617	-0.0555588\\
4.0642	-0.0555682\\
4.0668	-0.055578\\
4.0693	-0.0555875\\
4.0719	-0.0555974\\
4.0745	-0.0556073\\
4.077	-0.0556169\\
4.0796	-0.0556269\\
4.0822	-0.055637\\
4.0847	-0.0556467\\
4.0873	-0.0556568\\
4.0899	-0.0556669\\
4.0924	-0.0556767\\
4.095	-0.0556869\\
4.0975	-0.0556968\\
4.1001	-0.055707\\
4.1027	-0.0557173\\
4.1052	-0.0557273\\
4.1078	-0.0557376\\
4.1104	-0.055748\\
4.1129	-0.055758\\
4.1155	-0.0557685\\
4.1181	-0.055779\\
4.1206	-0.055789\\
4.1232	-0.0557996\\
4.1258	-0.0558101\\
4.1283	-0.0558203\\
4.1309	-0.0558309\\
4.1334	-0.0558411\\
4.136	-0.0558517\\
4.1386	-0.0558623\\
4.1411	-0.0558726\\
4.1437	-0.0558833\\
4.1463	-0.055894\\
4.1488	-0.0559043\\
4.1514	-0.0559151\\
4.154	-0.0559258\\
4.1565	-0.0559362\\
4.1591	-0.055947\\
4.1616	-0.0559574\\
4.1642	-0.0559682\\
4.1668	-0.0559791\\
4.1693	-0.0559895\\
4.1719	-0.0560004\\
4.1745	-0.0560113\\
4.177	-0.0560218\\
4.1796	-0.0560327\\
4.1822	-0.0560436\\
4.1847	-0.0560541\\
4.1873	-0.0560651\\
4.1898	-0.0560756\\
4.1924	-0.0560866\\
4.195	-0.0560975\\
4.1975	-0.0561081\\
4.2001	-0.0561191\\
4.2027	-0.0561301\\
4.2052	-0.0561407\\
4.2078	-0.0561517\\
4.2104	-0.0561628\\
4.2129	-0.0561734\\
4.2155	-0.0561844\\
4.218	-0.056195\\
4.2206	-0.0562061\\
4.2232	-0.0562172\\
4.2257	-0.0562278\\
4.2283	-0.0562389\\
4.2309	-0.05625\\
4.2334	-0.0562606\\
4.236	-0.0562717\\
4.2386	-0.0562828\\
4.2411	-0.0562935\\
4.2437	-0.0563046\\
4.2462	-0.0563152\\
4.2488	-0.0563263\\
4.2514	-0.0563374\\
4.2539	-0.0563481\\
4.2565	-0.0563592\\
4.2591	-0.0563703\\
4.2616	-0.056381\\
4.2642	-0.0563921\\
4.2668	-0.0564032\\
4.2693	-0.0564139\\
4.2719	-0.056425\\
4.2744	-0.0564357\\
4.277	-0.0564468\\
4.2796	-0.0564579\\
4.2821	-0.0564686\\
4.2847	-0.0564797\\
4.2873	-0.0564908\\
4.2898	-0.0565014\\
4.2924	-0.0565125\\
4.295	-0.0565236\\
4.2975	-0.0565343\\
4.3001	-0.0565453\\
4.3027	-0.0565564\\
4.3052	-0.0565671\\
4.3078	-0.0565781\\
4.3103	-0.0565887\\
4.3129	-0.0565998\\
4.3155	-0.0566108\\
4.318	-0.0566214\\
4.3206	-0.0566325\\
4.3232	-0.0566435\\
4.3257	-0.0566541\\
4.3283	-0.0566651\\
4.3309	-0.0566761\\
4.3334	-0.0566866\\
4.336	-0.0566976\\
4.3385	-0.0567082\\
4.3411	-0.0567191\\
4.3437	-0.0567301\\
4.3462	-0.0567406\\
4.3488	-0.0567515\\
4.3514	-0.0567624\\
4.3539	-0.0567729\\
4.3565	-0.0567838\\
4.3591	-0.0567946\\
4.3616	-0.0568051\\
4.3642	-0.0568159\\
4.3667	-0.0568263\\
4.3693	-0.0568371\\
4.3719	-0.0568479\\
4.3744	-0.0568583\\
4.377	-0.0568691\\
4.3796	-0.0568798\\
4.3821	-0.0568901\\
4.3847	-0.0569009\\
4.3873	-0.0569116\\
4.3898	-0.0569218\\
4.3924	-0.0569325\\
4.3949	-0.0569427\\
4.3975	-0.0569533\\
4.4001	-0.0569639\\
4.4026	-0.0569741\\
4.4052	-0.0569847\\
4.4078	-0.0569952\\
4.4103	-0.0570053\\
4.4129	-0.0570158\\
4.4155	-0.0570263\\
4.418	-0.0570363\\
4.4206	-0.0570467\\
4.4231	-0.0570567\\
4.4257	-0.0570671\\
4.4283	-0.0570775\\
4.4308	-0.0570874\\
4.4334	-0.0570977\\
4.436	-0.0571079\\
4.4385	-0.0571178\\
4.4411	-0.057128\\
4.4437	-0.0571382\\
4.4462	-0.0571479\\
4.4488	-0.057158\\
4.4514	-0.0571681\\
4.4539	-0.0571778\\
4.4565	-0.0571878\\
4.459	-0.0571974\\
4.4616	-0.0572074\\
4.4642	-0.0572173\\
4.4667	-0.0572268\\
4.4693	-0.0572367\\
4.4719	-0.0572465\\
4.4744	-0.0572559\\
4.477	-0.0572657\\
4.4796	-0.0572754\\
4.4821	-0.0572847\\
4.4847	-0.0572943\\
4.4872	-0.0573036\\
4.4898	-0.0573131\\
4.4924	-0.0573226\\
4.4949	-0.0573317\\
4.4975	-0.0573412\\
4.5001	-0.0573506\\
4.5026	-0.0573596\\
4.5052	-0.0573689\\
4.5078	-0.0573781\\
4.5103	-0.057387\\
4.5129	-0.0573962\\
4.5154	-0.0574049\\
4.518	-0.057414\\
4.5206	-0.0574231\\
4.5231	-0.0574317\\
4.5257	-0.0574406\\
4.5283	-0.0574495\\
4.5308	-0.057458\\
4.5334	-0.0574668\\
4.536	-0.0574755\\
4.5385	-0.0574839\\
4.5411	-0.0574925\\
4.5436	-0.0575008\\
4.5462	-0.0575093\\
4.5488	-0.0575178\\
4.5513	-0.0575258\\
4.5539	-0.0575342\\
4.5565	-0.0575425\\
4.559	-0.0575504\\
4.5616	-0.0575586\\
4.5642	-0.0575667\\
4.5667	-0.0575745\\
4.5693	-0.0575825\\
4.5718	-0.0575901\\
4.5744	-0.057598\\
4.577	-0.0576058\\
4.5795	-0.0576133\\
4.5821	-0.057621\\
4.5847	-0.0576286\\
4.5872	-0.0576359\\
4.5898	-0.0576434\\
4.5924	-0.0576508\\
4.5949	-0.0576579\\
4.5975	-0.0576652\\
4.6001	-0.0576725\\
4.6026	-0.0576794\\
4.6052	-0.0576864\\
4.6077	-0.0576932\\
4.6103	-0.0577002\\
4.6129	-0.057707\\
4.6154	-0.0577136\\
4.618	-0.0577203\\
4.6206	-0.057727\\
4.6231	-0.0577333\\
4.6257	-0.0577398\\
4.6283	-0.0577463\\
4.6308	-0.0577524\\
4.6334	-0.0577587\\
4.6359	-0.0577646\\
4.6385	-0.0577708\\
4.6411	-0.0577768\\
4.6436	-0.0577826\\
4.6462	-0.0577885\\
4.6488	-0.0577943\\
4.6513	-0.0577998\\
4.6539	-0.0578054\\
4.6565	-0.057811\\
4.659	-0.0578163\\
4.6616	-0.0578217\\
4.6641	-0.0578268\\
4.6667	-0.057832\\
4.6693	-0.0578371\\
4.6718	-0.057842\\
4.6744	-0.057847\\
4.677	-0.0578519\\
4.6795	-0.0578565\\
4.6821	-0.0578612\\
4.6847	-0.0578658\\
4.6872	-0.0578702\\
4.6898	-0.0578746\\
4.6923	-0.0578788\\
4.6949	-0.0578831\\
4.6975	-0.0578872\\
4.7	-0.0578911\\
4.7026	-0.0578951\\
4.7052	-0.057899\\
4.7077	-0.0579027\\
4.7103	-0.0579064\\
4.7129	-0.05791\\
4.7154	-0.0579134\\
4.718	-0.0579168\\
4.7205	-0.05792\\
4.7231	-0.0579232\\
4.7257	-0.0579263\\
4.7282	-0.0579292\\
4.7308	-0.0579322\\
4.7334	-0.057935\\
4.7359	-0.0579376\\
4.7385	-0.0579403\\
4.7411	-0.0579428\\
4.7436	-0.0579452\\
4.7462	-0.0579475\\
4.7487	-0.0579497\\
4.7513	-0.0579518\\
4.7539	-0.0579538\\
4.7564	-0.0579557\\
4.759	-0.0579575\\
4.7616	-0.0579593\\
4.7641	-0.0579608\\
4.7667	-0.0579624\\
4.7693	-0.0579638\\
4.7718	-0.057965\\
4.7744	-0.0579663\\
4.777	-0.0579674\\
4.7795	-0.0579683\\
4.7821	-0.0579692\\
4.7846	-0.05797\\
4.7872	-0.0579707\\
4.7898	-0.0579713\\
4.7923	-0.0579718\\
4.7949	-0.0579721\\
4.7975	-0.0579724\\
4.8	-0.0579726\\
4.8026	-0.0579726\\
4.8052	-0.0579726\\
4.8077	-0.0579724\\
4.8103	-0.0579722\\
4.8128	-0.0579718\\
4.8154	-0.0579713\\
4.818	-0.0579707\\
4.8205	-0.0579701\\
4.8231	-0.0579693\\
4.8257	-0.0579684\\
4.8282	-0.0579674\\
4.8308	-0.0579663\\
4.8334	-0.0579651\\
4.8359	-0.0579638\\
4.8385	-0.0579623\\
4.841	-0.0579608\\
4.8436	-0.0579592\\
4.8462	-0.0579574\\
4.8487	-0.0579556\\
4.8513	-0.0579536\\
4.8539	-0.0579516\\
4.8564	-0.0579494\\
4.859	-0.0579471\\
4.8616	-0.0579447\\
4.8641	-0.0579423\\
4.8667	-0.0579397\\
4.8692	-0.0579371\\
4.8718	-0.0579342\\
4.8744	-0.0579313\\
4.8769	-0.0579283\\
4.8795	-0.0579252\\
4.8821	-0.0579219\\
4.8846	-0.0579187\\
4.8872	-0.0579152\\
4.8898	-0.0579116\\
4.8923	-0.057908\\
4.8949	-0.0579043\\
4.8974	-0.0579005\\
4.9	-0.0578965\\
4.9026	-0.0578924\\
4.9051	-0.0578883\\
4.9077	-0.057884\\
4.9103	-0.0578796\\
4.9128	-0.0578752\\
4.9154	-0.0578706\\
4.918	-0.0578659\\
4.9205	-0.0578612\\
4.9231	-0.0578563\\
4.9257	-0.0578512\\
4.9282	-0.0578463\\
4.9308	-0.057841\\
4.9333	-0.0578359\\
4.9359	-0.0578304\\
4.9385	-0.0578249\\
4.941	-0.0578194\\
4.9436	-0.0578137\\
4.9462	-0.0578078\\
4.9487	-0.0578021\\
4.9513	-0.0577961\\
4.9539	-0.0577899\\
4.9564	-0.0577839\\
4.959	-0.0577776\\
4.9615	-0.0577714\\
4.9641	-0.0577649\\
4.9667	-0.0577582\\
4.9692	-0.0577518\\
4.9718	-0.0577449\\
4.9744	-0.057738\\
4.9769	-0.0577313\\
4.9795	-0.0577242\\
4.9821	-0.057717\\
4.9846	-0.05771\\
4.9872	-0.0577026\\
4.9897	-0.0576954\\
4.9923	-0.0576879\\
4.9949	-0.0576802\\
4.9974	-0.0576728\\
5	-0.057665\\
};
\addplot [color=mycolor22, line width=1.0pt, forget plot]
  table[row sep=crcr]{%
-0.125	-7.93892e-08\\
-0.1224	-8.10232e-08\\
-0.1199	-8.2594e-08\\
-0.1173	-8.42274e-08\\
-0.1147	-8.58605e-08\\
-0.1122	-8.74304e-08\\
-0.1096	-8.90629e-08\\
-0.1071	-9.06324e-08\\
-0.1045	-9.22642e-08\\
-0.1019	-9.38958e-08\\
-0.0994	-9.54644e-08\\
-0.0968	-9.70954e-08\\
-0.0942	-9.8726e-08\\
-0.0917	-1.00294e-07\\
-0.0891	-1.01924e-07\\
-0.0865	-1.03554e-07\\
-0.084	-1.0512e-07\\
-0.0814	-1.0675e-07\\
-0.0789	-1.08316e-07\\
-0.0763	-1.09944e-07\\
-0.0737	-1.11573e-07\\
-0.0712	-1.13138e-07\\
-0.0686	-1.14766e-07\\
-0.066	-1.16393e-07\\
-0.0635	-1.17958e-07\\
-0.0609	-1.19584e-07\\
-0.0583	-1.21211e-07\\
-0.0558	-1.22774e-07\\
-0.0532	-1.244e-07\\
-0.0507	-1.25963e-07\\
-0.0481	-1.27589e-07\\
-0.0455	-1.29214e-07\\
-0.043	-1.30776e-07\\
-0.0404	-1.324e-07\\
-0.0378	-1.34024e-07\\
-0.0353	-1.35586e-07\\
-0.0327	-1.37209e-07\\
-0.0301	-1.38832e-07\\
-0.0276	-1.40393e-07\\
-0.025	-1.42016e-07\\
-0.0224	-1.43638e-07\\
-0.0199	-1.45198e-07\\
-0.0173	-1.4682e-07\\
-0.0148	-1.48379e-07\\
-0.0122	-1.5e-07\\
-0.0096	-1.51621e-07\\
-0.0071	-1.53179e-07\\
-0.0045	-1.54799e-07\\
-0.0019	-1.5642e-07\\
0.0006	-1.57977e-07\\
0.0032	0.000931275\\
0.0058	0.00142723\\
0.0083	0.00152691\\
0.0109	0.00150419\\
0.0134	0.00146232\\
0.016	0.00140724\\
0.0186	0.00130664\\
0.0211	0.00114847\\
0.0237	0.000934789\\
0.0263	0.000704839\\
0.0288	0.000493228\\
0.0314	0.000287919\\
0.034	8.65314e-05\\
0.0365	-0.000115263\\
0.0391	-0.000336813\\
0.0416	-0.000555124\\
0.0442	-0.000778727\\
0.0468	-0.000994151\\
0.0493	-0.00119528\\
0.0519	-0.00140319\\
0.0545	-0.00161327\\
0.057	-0.00181659\\
0.0596	-0.00202563\\
0.0622	-0.00222862\\
0.0647	-0.0024174\\
0.0673	-0.00260906\\
0.0698	-0.0027913\\
0.0724	-0.00297946\\
0.075	-0.00316445\\
0.0775	-0.00333666\\
0.0801	-0.00350819\\
0.0827	-0.00367261\\
0.0852	-0.0038259\\
0.0878	-0.0039818\\
0.0904	-0.00413368\\
0.0929	-0.00427397\\
0.0955	-0.00441191\\
0.098	-0.00453653\\
0.1006	-0.00465914\\
0.1032	-0.00477649\\
0.1057	-0.00488501\\
0.1083	-0.00499227\\
0.1109	-0.0050919\\
0.1134	-0.00517946\\
0.116	-0.00526262\\
0.1186	-0.0053395\\
0.1211	-0.00540885\\
0.1237	-0.0054761\\
0.1263	-0.00553691\\
0.1288	-0.00558798\\
0.1314	-0.00563328\\
0.1339	-0.00567062\\
0.1365	-0.0057047\\
0.1391	-0.00573471\\
0.1416	-0.00575904\\
0.1442	-0.00577827\\
0.1468	-0.00579056\\
0.1493	-0.00579651\\
0.1519	-0.00579825\\
0.1545	-0.00579679\\
0.157	-0.00579244\\
0.1596	-0.00578378\\
0.1621	-0.00577048\\
0.1647	-0.00575148\\
0.1673	-0.00572845\\
0.1698	-0.00570401\\
0.1724	-0.00567692\\
0.175	-0.00564759\\
0.1775	-0.00561622\\
0.1801	-0.00557975\\
0.1827	-0.00554001\\
0.1852	-0.00550005\\
0.1878	-0.00545783\\
0.1903	-0.00541667\\
0.1929	-0.00537241\\
0.1955	-0.00532575\\
0.198	-0.00527861\\
0.2006	-0.00522821\\
0.2032	-0.00517774\\
0.2057	-0.00512969\\
0.2083	-0.00507976\\
0.2109	-0.00502889\\
0.2134	-0.00497859\\
0.216	-0.00492531\\
0.2185	-0.00487422\\
0.2211	-0.00482215\\
0.2237	-0.00477117\\
0.2262	-0.00472242\\
0.2288	-0.0046712\\
0.2314	-0.00461935\\
0.2339	-0.00456965\\
0.2365	-0.00451912\\
0.2391	-0.00447021\\
0.2416	-0.00442427\\
0.2442	-0.00437676\\
0.2467	-0.00433085\\
0.2493	-0.00428316\\
0.2519	-0.00423643\\
0.2544	-0.00419309\\
0.257	-0.00414962\\
0.2596	-0.00410705\\
0.2621	-0.00406625\\
0.2647	-0.00402382\\
0.2673	-0.003982\\
0.2698	-0.00394312\\
0.2724	-0.00390438\\
0.2749	-0.00386837\\
0.2775	-0.00383143\\
0.2801	-0.00379453\\
0.2826	-0.00375929\\
0.2852	-0.00372358\\
0.2878	-0.00368941\\
0.2903	-0.00365794\\
0.2929	-0.00362603\\
0.2955	-0.0035943\\
0.298	-0.00356382\\
0.3006	-0.00353263\\
0.3032	-0.0035026\\
0.3057	-0.00347508\\
0.3083	-0.0034475\\
0.3108	-0.00342135\\
0.3134	-0.00339414\\
0.316	-0.00336709\\
0.3185	-0.00334179\\
0.3211	-0.00331667\\
0.3237	-0.00329273\\
0.3262	-0.00327032\\
0.3288	-0.0032471\\
0.3314	-0.00322385\\
0.3339	-0.00320186\\
0.3365	-0.0031799\\
0.339	-0.00315988\\
0.3416	-0.00313991\\
0.3442	-0.00312025\\
0.3467	-0.00310131\\
0.3493	-0.00308173\\
0.3519	-0.00306277\\
0.3544	-0.0030455\\
0.357	-0.00302851\\
0.3596	-0.00301208\\
0.3621	-0.00299637\\
0.3647	-0.00298004\\
0.3672	-0.00296465\\
0.3698	-0.00294943\\
0.3724	-0.00293522\\
0.3749	-0.00292227\\
0.3775	-0.0029091\\
0.3801	-0.00289595\\
0.3826	-0.00288343\\
0.3852	-0.00287095\\
0.3878	-0.00285936\\
0.3903	-0.00284903\\
0.3929	-0.00283878\\
0.3954	-0.00282903\\
0.398	-0.0028189\\
0.4006	-0.00280907\\
0.4031	-0.00280025\\
0.4057	-0.0027919\\
0.4083	-0.00278417\\
0.4108	-0.00277696\\
0.4134	-0.00276945\\
0.416	-0.00276199\\
0.4185	-0.00275521\\
0.4211	-0.00274881\\
0.4236	-0.00274326\\
0.4262	-0.00273781\\
0.4288	-0.00273234\\
0.4313	-0.00272697\\
0.4339	-0.00272148\\
0.4365	-0.00271639\\
0.439	-0.002712\\
0.4416	-0.00270775\\
0.4442	-0.00270348\\
0.4467	-0.00269915\\
0.4493	-0.00269446\\
0.4519	-0.00268983\\
0.4544	-0.00268565\\
0.457	-0.00268154\\
0.4595	-0.00267755\\
0.4621	-0.00267306\\
0.4647	-0.00266813\\
0.4672	-0.00266312\\
0.4698	-0.00265787\\
0.4724	-0.00265266\\
0.4749	-0.00264752\\
0.4775	-0.00264176\\
0.4801	-0.00263535\\
0.4826	-0.00262862\\
0.4852	-0.00262121\\
0.4877	-0.00261386\\
0.4903	-0.00260594\\
0.4929	-0.00259749\\
0.4954	-0.00258861\\
0.498	-0.0025785\\
0.5006	-0.00256762\\
0.5031	-0.00255662\\
0.5057	-0.00254469\\
0.5083	-0.00253209\\
0.5108	-0.00251911\\
0.5134	-0.00250454\\
0.5159	-0.0024895\\
0.5185	-0.00247296\\
0.5211	-0.00245565\\
0.5236	-0.00243824\\
0.5262	-0.00241913\\
0.5288	-0.00239877\\
0.5313	-0.00237795\\
0.5339	-0.0023551\\
0.5365	-0.00233119\\
0.539	-0.00230727\\
0.5416	-0.00228128\\
0.5441	-0.00225504\\
0.5467	-0.0022263\\
0.5493	-0.0021961\\
0.5518	-0.00216581\\
0.5544	-0.00213314\\
0.557	-0.00209926\\
0.5595	-0.00206538\\
0.5621	-0.0020286\\
0.5647	-0.00199021\\
0.5672	-0.00195186\\
0.5698	-0.00191062\\
0.5723	-0.00186973\\
0.5749	-0.00182582\\
0.5775	-0.00178034\\
0.58	-0.00173502\\
0.5826	-0.00168625\\
0.5852	-0.00163596\\
0.5877	-0.00158627\\
0.5903	-0.00153321\\
0.5929	-0.0014786\\
0.5954	-0.00142452\\
0.598	-0.00136659\\
0.6006	-0.00130705\\
0.6031	-0.0012484\\
0.6057	-0.00118601\\
0.6082	-0.00112464\\
0.6108	-0.00105925\\
0.6134	-0.000992183\\
0.6159	-0.000926145\\
0.6185	-0.000855984\\
0.6211	-0.000784425\\
0.6236	-0.000714312\\
0.6262	-0.000639949\\
0.6288	-0.000564015\\
0.6313	-0.000489506\\
0.6339	-0.000410556\\
0.6364	-0.00033336\\
0.639	-0.000251816\\
0.6416	-0.000168955\\
0.6441	-8.79548e-05\\
0.6467	-2.28472e-06\\
0.6493	8.47943e-05\\
0.6518	0.000169736\\
0.6544	0.000259224\\
0.657	0.000349857\\
0.6595	0.000438137\\
0.6621	0.000531191\\
0.6646	0.000621858\\
0.6672	0.000717297\\
0.6698	0.000813777\\
0.6723	0.000907452\\
0.6749	0.00100583\\
0.6775	0.00110523\\
0.68	0.00120181\\
0.6826	0.00130323\\
0.6852	0.00140552\\
0.6877	0.00150461\\
0.6903	0.00160838\\
0.6928	0.00170887\\
0.6954	0.00181416\\
0.698	0.00192023\\
0.7005	0.00202287\\
0.7031	0.00213017\\
0.7057	0.00223794\\
0.7082	0.00234201\\
0.7108	0.00245074\\
0.7134	0.00255998\\
0.7159	0.00266544\\
0.7185	0.00277545\\
0.721	0.00288143\\
0.7236	0.00299182\\
0.7262	0.0031024\\
0.7287	0.00320892\\
0.7313	0.00331986\\
0.7339	0.00343088\\
0.7364	0.0035376\\
0.739	0.00364845\\
0.7416	0.00375916\\
0.7441	0.00386548\\
0.7467	0.00397592\\
0.7492	0.0040819\\
0.7518	0.00419182\\
0.7544	0.0043013\\
0.7569	0.00440612\\
0.7595	0.00451465\\
0.7621	0.00462269\\
0.7646	0.00472608\\
0.7672	0.00483299\\
0.7698	0.00493916\\
0.7723	0.00504048\\
0.7749	0.00514502\\
0.7775	0.0052487\\
0.78	0.00534757\\
0.7826	0.00544947\\
0.7851	0.00554645\\
0.7877	0.00564619\\
0.7903	0.00574471\\
0.7928	0.00583828\\
0.7954	0.00593438\\
0.798	0.00602919\\
0.8005	0.00611905\\
0.8031	0.00621105\\
0.8057	0.00630148\\
0.8082	0.00638692\\
0.8108	0.00647419\\
0.8133	0.00655654\\
0.8159	0.0066405\\
0.8185	0.00672265\\
0.821	0.00679982\\
0.8236	0.00687815\\
0.8262	0.00695447\\
0.8287	0.00702595\\
0.8313	0.00709826\\
0.8339	0.00716842\\
0.8364	0.00723374\\
0.839	0.00729939\\
0.8415	0.00736028\\
0.8441	0.00742125\\
0.8467	0.00747978\\
0.8492	0.00753372\\
0.8518	0.00758727\\
0.8544	0.00763814\\
0.8569	0.00768447\\
0.8595	0.00772995\\
0.8621	0.00777263\\
0.8646	0.007811\\
0.8672	0.00784804\\
0.8697	0.00788081\\
0.8723	0.00791189\\
0.8749	0.00793984\\
0.8774	0.00796374\\
0.88	0.00798549\\
0.8826	0.00800401\\
0.8851	0.00801868\\
0.8877	0.0080306\\
0.8903	0.00803906\\
0.8928	0.00804392\\
0.8954	0.00804553\\
0.8979	0.00804373\\
0.9005	0.0080383\\
0.9031	0.00802918\\
0.9056	0.00801686\\
0.9082	0.00800033\\
0.9108	0.00797997\\
0.9133	0.00795676\\
0.9159	0.00792879\\
0.9185	0.00789685\\
0.921	0.00786232\\
0.9236	0.0078224\\
0.9262	0.00777837\\
0.9287	0.00773213\\
0.9313	0.00767995\\
0.9338	0.00762579\\
0.9364	0.00756525\\
0.939	0.00750036\\
0.9415	0.00743384\\
0.9441	0.00736035\\
0.9467	0.00728244\\
0.9492	0.00720332\\
0.9518	0.00711661\\
0.9544	0.00702532\\
0.9569	0.00693319\\
0.9595	0.00683285\\
0.962	0.00673198\\
0.9646	0.00662251\\
0.9672	0.00650832\\
0.9697	0.00639403\\
0.9723	0.00627046\\
0.9749	0.00614207\\
0.9774	0.00601405\\
0.98	0.00587616\\
0.9826	0.00573339\\
0.9851	0.00559148\\
0.9877	0.00543902\\
0.9902	0.00528772\\
0.9928	0.00512547\\
0.9954	0.00495822\\
0.9979	0.00479268\\
1.0005	0.00461557\\
1.0031	0.00443338\\
1.0056	0.00425338\\
1.0082	0.00406116\\
1.0108	0.00386382\\
1.0133	0.00366925\\
1.0159	0.00346186\\
1.0184	0.00325757\\
1.021	0.00304003\\
1.0236	0.00281729\\
1.0261	0.00259821\\
1.0287	0.00236528\\
1.0313	0.00212716\\
1.0338	0.00189328\\
1.0364	0.00164492\\
1.039	0.00139132\\
1.0415	0.00114253\\
1.0441	0.000878669\\
1.0466	0.000620036\\
1.0492	0.000345937\\
1.0518	6.66094e-05\\
1.0543	-0.000206915\\
1.0569	-0.000496516\\
1.0595	-0.000791339\\
1.062	-0.00107973\\
1.0646	-0.00138475\\
1.0672	-0.00169497\\
1.0697	-0.00199815\\
1.0723	-0.00231855\\
1.0748	-0.00263151\\
1.0774	-0.00296205\\
1.08	-0.00329771\\
1.0825	-0.0036253\\
1.0851	-0.003971\\
1.0877	-0.0043218\\
1.0902	-0.00466392\\
1.0928	-0.00502468\\
1.0954	-0.00539047\\
1.0979	-0.00574692\\
1.1005	-0.00612252\\
1.1031	-0.00650309\\
1.1056	-0.00687371\\
1.1082	-0.00726398\\
1.1107	-0.00764385\\
1.1133	-0.00804369\\
1.1159	-0.00844835\\
1.1184	-0.00884199\\
1.121	-0.00925606\\
1.1236	-0.00967489\\
1.1261	-0.0100821\\
1.1287	-0.0105101\\
1.1313	-0.0109427\\
1.1338	-0.0113631\\
1.1364	-0.0118047\\
1.1389	-0.0122337\\
1.1415	-0.0126842\\
1.1441	-0.0131391\\
1.1466	-0.0135807\\
1.1492	-0.0140443\\
1.1518	-0.0145121\\
1.1543	-0.014966\\
1.1569	-0.0154422\\
1.1595	-0.0159226\\
1.162	-0.0163883\\
1.1646	-0.0168767\\
1.1671	-0.0173501\\
1.1697	-0.0178464\\
1.1723	-0.0183465\\
1.1748	-0.0188311\\
1.1774	-0.0193388\\
1.18	-0.0198503\\
1.1825	-0.0203456\\
1.1851	-0.0208642\\
1.1877	-0.0213864\\
1.1902	-0.0218919\\
1.1928	-0.022421\\
1.1953	-0.022933\\
1.1979	-0.0234686\\
1.2005	-0.0240076\\
1.203	-0.0245289\\
1.2056	-0.025074\\
1.2082	-0.0256223\\
1.2107	-0.0261523\\
1.2133	-0.0267064\\
1.2159	-0.0272634\\
1.2184	-0.0278015\\
1.221	-0.0283639\\
1.2235	-0.0289072\\
1.2261	-0.0294747\\
1.2287	-0.0300447\\
1.2312	-0.0305951\\
1.2338	-0.0311698\\
1.2364	-0.0317468\\
1.2389	-0.0323037\\
1.2415	-0.032885\\
1.2441	-0.0334683\\
1.2466	-0.034031\\
1.2492	-0.0346181\\
1.2518	-0.0352069\\
1.2543	-0.0357748\\
1.2569	-0.036367\\
1.2594	-0.0369379\\
1.262	-0.0375331\\
1.2646	-0.0381296\\
1.2671	-0.0387045\\
1.2697	-0.0393035\\
1.2723	-0.0399037\\
1.2748	-0.0404818\\
1.2774	-0.041084\\
1.28	-0.041687\\
1.2825	-0.0422676\\
1.2851	-0.0428721\\
1.2876	-0.0434539\\
1.2902	-0.0440595\\
1.2928	-0.0446656\\
1.2953	-0.0452486\\
1.2979	-0.0458552\\
1.3005	-0.046462\\
1.303	-0.0470455\\
1.3056	-0.0476523\\
1.3082	-0.048259\\
1.3107	-0.0488422\\
1.3133	-0.0494484\\
1.3158	-0.0500309\\
1.3184	-0.0506362\\
1.321	-0.051241\\
1.3235	-0.0518219\\
1.3261	-0.0524252\\
1.3287	-0.0530277\\
1.3312	-0.0536061\\
1.3338	-0.0542066\\
1.3364	-0.054806\\
1.3389	-0.0553811\\
1.3415	-0.0559779\\
1.344	-0.0565505\\
1.3466	-0.0571444\\
1.3492	-0.0577368\\
1.3517	-0.0583047\\
1.3543	-0.0588936\\
1.3569	-0.0594806\\
1.3594	-0.0600432\\
1.362	-0.0606262\\
1.3646	-0.061207\\
1.3671	-0.0617633\\
1.3697	-0.0623396\\
1.3722	-0.0628914\\
1.3748	-0.0634628\\
1.3774	-0.0640315\\
1.3799	-0.0645758\\
1.3825	-0.0651391\\
1.3851	-0.0656995\\
1.3876	-0.0662355\\
1.3902	-0.0667899\\
1.3928	-0.0673412\\
1.3953	-0.0678681\\
1.3979	-0.0684128\\
1.4005	-0.0689541\\
1.403	-0.0694712\\
1.4056	-0.0700055\\
1.4081	-0.0705157\\
1.4107	-0.0710427\\
1.4133	-0.0715657\\
1.4158	-0.0720649\\
1.4184	-0.0725801\\
1.421	-0.0730912\\
1.4235	-0.0735786\\
1.4261	-0.0740814\\
1.4287	-0.0745798\\
1.4312	-0.0750548\\
1.4338	-0.0755444\\
1.4363	-0.0760108\\
1.4389	-0.0764913\\
1.4415	-0.076967\\
1.444	-0.0774199\\
1.4466	-0.0778861\\
1.4492	-0.0783473\\
1.4517	-0.0787861\\
1.4543	-0.0792373\\
1.4569	-0.0796834\\
1.4594	-0.0801074\\
1.462	-0.0805431\\
1.4645	-0.080957\\
1.4671	-0.0813821\\
1.4697	-0.0818017\\
1.4722	-0.0821999\\
1.4748	-0.0826085\\
1.4774	-0.0830114\\
1.4799	-0.0833934\\
1.4825	-0.083785\\
1.4851	-0.0841707\\
1.4876	-0.084536\\
1.4902	-0.08491\\
1.4927	-0.085264\\
1.4953	-0.0856262\\
1.4979	-0.0859822\\
1.5004	-0.0863188\\
1.503	-0.0866627\\
1.5056	-0.0870003\\
1.5081	-0.087319\\
1.5107	-0.0876442\\
1.5133	-0.0879631\\
1.5158	-0.0882636\\
1.5184	-0.0885698\\
1.5209	-0.088858\\
1.5235	-0.0891514\\
1.5261	-0.0894383\\
1.5286	-0.0897078\\
1.5312	-0.0899817\\
1.5338	-0.0902489\\
1.5363	-0.0904995\\
1.5389	-0.0907535\\
1.5415	-0.0910009\\
1.544	-0.0912323\\
1.5466	-0.0914663\\
1.5491	-0.0916849\\
1.5517	-0.0919056\\
1.5543	-0.0921195\\
1.5568	-0.0923186\\
1.5594	-0.0925191\\
1.562	-0.0927126\\
1.5645	-0.0928922\\
1.5671	-0.0930722\\
1.5697	-0.0932454\\
1.5722	-0.0934054\\
1.5748	-0.093565\\
1.5774	-0.0937177\\
1.5799	-0.093858\\
1.5825	-0.0939971\\
1.585	-0.0941244\\
1.5876	-0.0942501\\
1.5902	-0.0943688\\
1.5927	-0.0944765\\
1.5953	-0.0945818\\
1.5979	-0.0946802\\
1.6004	-0.0947683\\
1.603	-0.0948533\\
1.6056	-0.0949315\\
1.6081	-0.0950002\\
1.6107	-0.095065\\
1.6132	-0.095121\\
1.6158	-0.0951726\\
1.6184	-0.0952175\\
1.6209	-0.0952543\\
1.6235	-0.0952861\\
1.6261	-0.0953112\\
1.6286	-0.0953291\\
1.6312	-0.0953413\\
1.6338	-0.095347\\
1.6363	-0.0953463\\
1.6389	-0.0953392\\
1.6414	-0.0953263\\
1.644	-0.0953066\\
1.6466	-0.0952805\\
1.6491	-0.0952495\\
1.6517	-0.095211\\
1.6543	-0.0951663\\
1.6568	-0.0951175\\
1.6594	-0.0950607\\
1.662	-0.0949978\\
1.6645	-0.0949315\\
1.6671	-0.0948568\\
1.6696	-0.0947793\\
1.6722	-0.0946929\\
1.6748	-0.0946006\\
1.6773	-0.0945064\\
1.6799	-0.0944028\\
1.6825	-0.0942935\\
1.685	-0.094183\\
1.6876	-0.0940627\\
1.6902	-0.0939369\\
1.6927	-0.0938107\\
1.6953	-0.0936742\\
1.6978	-0.0935379\\
1.7004	-0.093391\\
1.703	-0.0932388\\
1.7055	-0.0930877\\
1.7081	-0.0929255\\
1.7107	-0.0927583\\
1.7132	-0.0925928\\
1.7158	-0.092416\\
1.7184	-0.0922343\\
1.7209	-0.0920551\\
1.7235	-0.0918642\\
1.7261	-0.0916687\\
1.7286	-0.0914764\\
1.7312	-0.091272\\
1.7337	-0.0910714\\
1.7363	-0.0908584\\
1.7389	-0.0906413\\
1.7414	-0.0904285\\
1.744	-0.0902032\\
1.7466	-0.0899739\\
1.7491	-0.0897497\\
1.7517	-0.0895128\\
1.7543	-0.089272\\
1.7568	-0.089037\\
1.7594	-0.088789\\
1.7619	-0.0885472\\
1.7645	-0.0882923\\
1.7671	-0.088034\\
1.7696	-0.0877825\\
1.7722	-0.0875178\\
1.7748	-0.0872499\\
1.7773	-0.0869894\\
1.7799	-0.0867155\\
1.7825	-0.0864387\\
1.785	-0.0861698\\
1.7876	-0.0858875\\
1.7901	-0.0856135\\
1.7927	-0.0853259\\
1.7953	-0.0850359\\
1.7978	-0.0847547\\
1.8004	-0.0844599\\
1.803	-0.0841629\\
1.8055	-0.0838752\\
1.8081	-0.0835738\\
1.8107	-0.0832705\\
1.8132	-0.082977\\
1.8158	-0.0826699\\
1.8183	-0.0823728\\
1.8209	-0.0820623\\
1.8235	-0.08175\\
1.826	-0.0814484\\
1.8286	-0.0811332\\
1.8312	-0.0808165\\
1.8337	-0.0805109\\
1.8363	-0.0801917\\
1.8389	-0.0798714\\
1.8414	-0.0795624\\
1.844	-0.07924\\
1.8465	-0.0789292\\
1.8491	-0.0786051\\
1.8517	-0.0782802\\
1.8542	-0.0779671\\
1.8568	-0.0776409\\
1.8594	-0.0773142\\
1.8619	-0.0769995\\
1.8645	-0.0766719\\
1.8671	-0.076344\\
1.8696	-0.0760284\\
1.8722	-0.0757\\
1.8747	-0.0753842\\
1.8773	-0.0750557\\
1.8799	-0.0747272\\
1.8824	-0.0744115\\
1.885	-0.0740833\\
1.8876	-0.0737554\\
1.8901	-0.0734404\\
1.8927	-0.0731132\\
1.8953	-0.0727864\\
1.8978	-0.0724727\\
1.9004	-0.072147\\
1.903	-0.071822\\
1.9055	-0.0715102\\
1.9081	-0.0711867\\
1.9106	-0.0708764\\
1.9132	-0.0705546\\
1.9158	-0.0702338\\
1.9183	-0.0699263\\
1.9209	-0.0696077\\
1.9235	-0.0692901\\
1.926	-0.068986\\
1.9286	-0.068671\\
1.9312	-0.0683573\\
1.9337	-0.068057\\
1.9363	-0.0677461\\
1.9388	-0.0674486\\
1.9414	-0.0671408\\
1.944	-0.0668346\\
1.9465	-0.0665417\\
1.9491	-0.0662389\\
1.9517	-0.0659378\\
1.9542	-0.0656501\\
1.9568	-0.0653527\\
1.9594	-0.0650572\\
1.9619	-0.064775\\
1.9645	-0.0644836\\
1.967	-0.0642053\\
1.9696	-0.0639179\\
1.9722	-0.0636328\\
1.9747	-0.0633607\\
1.9773	-0.06308\\
1.9799	-0.0628016\\
1.9824	-0.0625361\\
1.985	-0.0622624\\
1.9876	-0.0619911\\
1.9901	-0.0617326\\
1.9927	-0.0614663\\
1.9952	-0.0612125\\
1.9978	-0.0609512\\
2.0004	-0.0606926\\
2.0029	-0.0604464\\
2.0055	-0.060193\\
2.0081	-0.0599424\\
2.0106	-0.059704\\
2.0132	-0.0594588\\
2.0158	-0.0592165\\
2.0183	-0.0589862\\
2.0209	-0.0587496\\
2.0234	-0.0585248\\
2.026	-0.0582939\\
2.0286	-0.058066\\
2.0311	-0.0578498\\
2.0337	-0.0576278\\
2.0363	-0.0574089\\
2.0388	-0.0572014\\
2.0414	-0.0569886\\
2.044	-0.0567789\\
2.0465	-0.0565803\\
2.0491	-0.0563769\\
2.0517	-0.0561766\\
2.0542	-0.0559871\\
2.0568	-0.0557932\\
2.0593	-0.0556098\\
2.0619	-0.0554222\\
2.0645	-0.055238\\
2.067	-0.0550639\\
2.0696	-0.0548862\\
2.0722	-0.0547117\\
2.0747	-0.0545472\\
2.0773	-0.0543793\\
2.0799	-0.0542148\\
2.0824	-0.0540598\\
2.085	-0.0539019\\
2.0875	-0.0537532\\
2.0901	-0.053602\\
2.0927	-0.0534542\\
2.0952	-0.0533152\\
2.0978	-0.0531741\\
2.1004	-0.0530364\\
2.1029	-0.0529072\\
2.1055	-0.0527762\\
2.1081	-0.0526486\\
2.1106	-0.0525292\\
2.1132	-0.0524083\\
2.1157	-0.0522954\\
2.1183	-0.0521813\\
2.1209	-0.0520706\\
2.1234	-0.0519674\\
2.126	-0.0518634\\
2.1286	-0.0517629\\
2.1311	-0.0516695\\
2.1337	-0.0515756\\
2.1363	-0.0514852\\
2.1388	-0.0514015\\
2.1414	-0.0513177\\
2.1439	-0.0512404\\
2.1465	-0.0511632\\
2.1491	-0.0510895\\
2.1516	-0.0510217\\
2.1542	-0.0509546\\
2.1568	-0.0508907\\
2.1593	-0.0508325\\
2.1619	-0.0507752\\
2.1645	-0.0507212\\
2.167	-0.0506724\\
2.1696	-0.0506248\\
2.1721	-0.0505822\\
2.1747	-0.050541\\
2.1773	-0.050503\\
2.1798	-0.0504696\\
2.1824	-0.050438\\
2.185	-0.0504095\\
2.1875	-0.0503851\\
2.1901	-0.0503629\\
2.1927	-0.0503437\\
2.1952	-0.0503283\\
2.1978	-0.0503152\\
2.2004	-0.0503052\\
2.2029	-0.0502985\\
2.2055	-0.0502945\\
2.208	-0.0502935\\
2.2106	-0.0502953\\
2.2132	-0.0503002\\
2.2157	-0.0503076\\
2.2183	-0.0503181\\
2.2209	-0.0503316\\
2.2234	-0.0503472\\
2.226	-0.0503663\\
2.2286	-0.0503882\\
2.2311	-0.0504118\\
2.2337	-0.0504392\\
2.2362	-0.050468\\
2.2388	-0.0505007\\
2.2414	-0.0505361\\
2.2439	-0.0505726\\
2.2465	-0.0506131\\
2.2491	-0.0506563\\
2.2516	-0.0507002\\
2.2542	-0.0507484\\
2.2568	-0.0507991\\
2.2593	-0.0508502\\
2.2619	-0.0509057\\
2.2644	-0.0509614\\
2.267	-0.0510217\\
2.2696	-0.0510843\\
2.2721	-0.0511467\\
2.2747	-0.0512139\\
2.2773	-0.0512834\\
2.2798	-0.0513523\\
2.2824	-0.0514261\\
2.285	-0.0515021\\
2.2875	-0.0515771\\
2.2901	-0.0516573\\
2.2926	-0.0517363\\
2.2952	-0.0518205\\
2.2978	-0.0519067\\
2.3003	-0.0519915\\
2.3029	-0.0520816\\
2.3055	-0.0521735\\
2.308	-0.0522637\\
2.3106	-0.0523593\\
2.3132	-0.0524567\\
2.3157	-0.0525521\\
2.3183	-0.0526529\\
2.3208	-0.0527515\\
2.3234	-0.0528557\\
2.326	-0.0529615\\
2.3285	-0.0530647\\
2.3311	-0.0531736\\
2.3337	-0.053284\\
2.3362	-0.0533916\\
2.3388	-0.0535049\\
2.3414	-0.0536196\\
2.3439	-0.0537312\\
2.3465	-0.0538487\\
2.349	-0.0539628\\
2.3516	-0.0540827\\
2.3542	-0.0542039\\
2.3567	-0.0543216\\
2.3593	-0.0544451\\
2.3619	-0.0545698\\
2.3644	-0.0546907\\
2.367	-0.0548175\\
2.3696	-0.0549453\\
2.3721	-0.0550691\\
2.3747	-0.0551988\\
2.3773	-0.0553295\\
2.3798	-0.0554559\\
2.3824	-0.0555883\\
2.3849	-0.0557163\\
2.3875	-0.0558502\\
2.3901	-0.0559849\\
2.3926	-0.056115\\
2.3952	-0.056251\\
2.3978	-0.0563876\\
2.4003	-0.0565196\\
2.4029	-0.0566574\\
2.4055	-0.0567957\\
2.408	-0.0569292\\
2.4106	-0.0570684\\
2.4131	-0.0572028\\
2.4157	-0.0573428\\
2.4183	-0.0574833\\
2.4208	-0.0576186\\
2.4234	-0.0577597\\
2.426	-0.057901\\
2.4285	-0.0580371\\
2.4311	-0.0581788\\
2.4337	-0.0583207\\
2.4362	-0.0584572\\
2.4388	-0.0585993\\
2.4413	-0.058736\\
2.4439	-0.0588782\\
2.4465	-0.0590204\\
2.449	-0.0591571\\
2.4516	-0.0592992\\
2.4542	-0.0594412\\
2.4567	-0.0595776\\
2.4593	-0.0597193\\
2.4619	-0.0598609\\
2.4644	-0.0599968\\
2.467	-0.0601378\\
2.4695	-0.0602732\\
2.4721	-0.0604137\\
2.4747	-0.0605539\\
2.4772	-0.0606883\\
2.4798	-0.0608277\\
2.4824	-0.0609667\\
2.4849	-0.0610999\\
2.4875	-0.061238\\
2.4901	-0.0613755\\
2.4926	-0.0615073\\
2.4952	-0.0616438\\
2.4977	-0.0617744\\
2.5003	-0.0619097\\
2.5029	-0.0620444\\
2.5054	-0.0621733\\
2.508	-0.0623066\\
2.5106	-0.0624392\\
2.5131	-0.062566\\
2.5157	-0.0626971\\
2.5183	-0.0628274\\
2.5208	-0.062952\\
2.5234	-0.0630807\\
2.526	-0.0632086\\
2.5285	-0.0633306\\
2.5311	-0.0634567\\
2.5336	-0.0635771\\
2.5362	-0.0637013\\
2.5388	-0.0638246\\
2.5413	-0.0639422\\
2.5439	-0.0640635\\
2.5465	-0.0641838\\
2.549	-0.0642984\\
2.5516	-0.0644166\\
2.5542	-0.0645337\\
2.5567	-0.0646452\\
2.5593	-0.0647601\\
2.5618	-0.0648695\\
2.5644	-0.0649821\\
2.567	-0.0650936\\
2.5695	-0.0651996\\
2.5721	-0.0653087\\
2.5747	-0.0654165\\
2.5772	-0.065519\\
2.5798	-0.0656244\\
2.5824	-0.0657285\\
2.5849	-0.0658273\\
2.5875	-0.0659289\\
2.59	-0.0660253\\
2.5926	-0.0661242\\
2.5952	-0.0662218\\
2.5977	-0.0663144\\
2.6003	-0.0664093\\
2.6029	-0.0665028\\
2.6054	-0.0665914\\
2.608	-0.0666822\\
2.6106	-0.0667715\\
2.6131	-0.0668561\\
2.6157	-0.0669427\\
2.6182	-0.0670245\\
2.6208	-0.0671082\\
2.6234	-0.0671904\\
2.6259	-0.0672681\\
2.6285	-0.0673474\\
2.6311	-0.0674253\\
2.6336	-0.0674987\\
2.6362	-0.0675736\\
2.6388	-0.067647\\
2.6413	-0.0677161\\
2.6439	-0.0677865\\
2.6464	-0.0678528\\
2.649	-0.0679202\\
2.6516	-0.0679861\\
2.6541	-0.0680479\\
2.6567	-0.0681108\\
2.6593	-0.0681721\\
2.6618	-0.0682295\\
2.6644	-0.0682878\\
2.667	-0.0683444\\
2.6695	-0.0683975\\
2.6721	-0.0684511\\
2.6746	-0.0685012\\
2.6772	-0.0685517\\
2.6798	-0.0686007\\
2.6823	-0.0686463\\
2.6849	-0.0686923\\
2.6875	-0.0687366\\
2.69	-0.0687778\\
2.6926	-0.0688191\\
2.6952	-0.0688588\\
2.6977	-0.0688955\\
2.7003	-0.0689322\\
2.7029	-0.0689673\\
2.7054	-0.0689996\\
2.708	-0.0690316\\
2.7105	-0.069061\\
2.7131	-0.0690901\\
2.7157	-0.0691175\\
2.7182	-0.0691425\\
2.7208	-0.069167\\
2.7234	-0.06919\\
2.7259	-0.0692106\\
2.7285	-0.0692306\\
2.7311	-0.0692491\\
2.7336	-0.0692654\\
2.7362	-0.0692809\\
2.7387	-0.0692944\\
2.7413	-0.069307\\
2.7439	-0.0693182\\
2.7464	-0.0693275\\
2.749	-0.0693357\\
2.7516	-0.0693425\\
2.7541	-0.0693477\\
2.7567	-0.0693516\\
2.7593	-0.0693542\\
2.7618	-0.0693552\\
2.7644	-0.069355\\
2.7669	-0.0693534\\
2.7695	-0.0693504\\
2.7721	-0.069346\\
2.7746	-0.0693405\\
2.7772	-0.0693334\\
2.7798	-0.069325\\
2.7823	-0.0693156\\
2.7849	-0.0693046\\
2.7875	-0.0692922\\
2.79	-0.069279\\
2.7926	-0.0692641\\
2.7951	-0.0692485\\
2.7977	-0.069231\\
2.8003	-0.0692123\\
2.8028	-0.0691931\\
2.8054	-0.0691719\\
2.808	-0.0691495\\
2.8105	-0.0691268\\
2.8131	-0.069102\\
2.8157	-0.069076\\
2.8182	-0.0690499\\
2.8208	-0.0690217\\
2.8233	-0.0689934\\
2.8259	-0.0689628\\
2.8285	-0.0689312\\
2.831	-0.0688997\\
2.8336	-0.0688658\\
2.8362	-0.0688309\\
2.8387	-0.0687963\\
2.8413	-0.0687593\\
2.8439	-0.0687213\\
2.8464	-0.0686837\\
2.849	-0.0686436\\
2.8516	-0.0686026\\
2.8541	-0.0685621\\
2.8567	-0.0685191\\
2.8592	-0.0684769\\
2.8618	-0.068432\\
2.8644	-0.0683863\\
2.8669	-0.0683414\\
2.8695	-0.0682938\\
2.8721	-0.0682453\\
2.8746	-0.0681979\\
2.8772	-0.0681478\\
2.8798	-0.0680968\\
2.8823	-0.068047\\
2.8849	-0.0679944\\
2.8874	-0.067943\\
2.89	-0.0678889\\
2.8926	-0.067834\\
2.8951	-0.0677805\\
2.8977	-0.0677241\\
2.9003	-0.067667\\
2.9028	-0.0676115\\
2.9054	-0.067553\\
2.908	-0.0674939\\
2.9105	-0.0674364\\
2.9131	-0.067376\\
2.9156	-0.0673173\\
2.9182	-0.0672557\\
2.9208	-0.0671934\\
2.9233	-0.067133\\
2.9259	-0.0670697\\
2.9285	-0.0670057\\
2.931	-0.0669438\\
2.9336	-0.0668788\\
2.9362	-0.0668133\\
2.9387	-0.0667499\\
2.9413	-0.0666834\\
2.9438	-0.0666191\\
2.9464	-0.0665518\\
2.949	-0.066484\\
2.9515	-0.0664185\\
2.9541	-0.0663499\\
2.9567	-0.0662809\\
2.9592	-0.0662142\\
2.9618	-0.0661445\\
2.9644	-0.0660744\\
2.9669	-0.0660067\\
2.9695	-0.065936\\
2.972	-0.0658678\\
2.9746	-0.0657965\\
2.9772	-0.0657249\\
2.9797	-0.0656558\\
2.9823	-0.0655837\\
2.9849	-0.0655114\\
2.9874	-0.0654416\\
2.99	-0.0653689\\
2.9926	-0.0652959\\
2.9951	-0.0652256\\
2.9977	-0.0651523\\
3.0003	-0.0650788\\
3.0028	-0.065008\\
3.0054	-0.0649343\\
3.0079	-0.0648632\\
3.0105	-0.0647893\\
3.0131	-0.0647152\\
3.0156	-0.0646439\\
3.0182	-0.0645696\\
3.0208	-0.0644954\\
3.0233	-0.0644239\\
3.0259	-0.0643495\\
3.0285	-0.0642751\\
3.031	-0.0642036\\
3.0336	-0.0641292\\
3.0361	-0.0640576\\
3.0387	-0.0639832\\
3.0413	-0.0639089\\
3.0438	-0.0638374\\
3.0464	-0.0637632\\
3.049	-0.063689\\
3.0515	-0.0636177\\
3.0541	-0.0635437\\
3.0567	-0.0634697\\
3.0592	-0.0633987\\
3.0618	-0.063325\\
3.0643	-0.0632542\\
3.0669	-0.0631808\\
3.0695	-0.0631074\\
3.072	-0.0630371\\
3.0746	-0.0629641\\
3.0772	-0.0628912\\
3.0797	-0.0628214\\
3.0823	-0.0627489\\
3.0849	-0.0626766\\
3.0874	-0.0626073\\
3.09	-0.0625354\\
3.0925	-0.0624665\\
3.0951	-0.0623951\\
3.0977	-0.062324\\
3.1002	-0.0622557\\
3.1028	-0.0621851\\
3.1054	-0.0621147\\
3.1079	-0.0620472\\
3.1105	-0.0619773\\
3.1131	-0.0619077\\
3.1156	-0.0618411\\
3.1182	-0.0617721\\
3.1207	-0.061706\\
3.1233	-0.0616376\\
3.1259	-0.0615695\\
3.1284	-0.0615043\\
3.131	-0.0614369\\
3.1336	-0.0613698\\
3.1361	-0.0613056\\
3.1387	-0.0612392\\
3.1413	-0.0611731\\
3.1438	-0.0611099\\
3.1464	-0.0610445\\
3.1489	-0.060982\\
3.1515	-0.0609174\\
3.1541	-0.0608531\\
3.1566	-0.0607917\\
3.1592	-0.0607282\\
3.1618	-0.0606651\\
3.1643	-0.0606048\\
3.1669	-0.0605425\\
3.1695	-0.0604806\\
3.172	-0.0604215\\
3.1746	-0.0603604\\
3.1772	-0.0602997\\
3.1797	-0.0602417\\
3.1823	-0.0601819\\
3.1848	-0.0601247\\
3.1874	-0.0600657\\
3.19	-0.0600071\\
3.1925	-0.0599512\\
3.1951	-0.0598935\\
3.1977	-0.0598362\\
3.2002	-0.0597815\\
3.2028	-0.0597251\\
3.2054	-0.0596692\\
3.2079	-0.0596158\\
3.2105	-0.0595607\\
3.213	-0.0595082\\
3.2156	-0.059454\\
3.2182	-0.0594003\\
3.2207	-0.059349\\
3.2233	-0.0592962\\
3.2259	-0.0592439\\
3.2284	-0.0591939\\
3.231	-0.0591425\\
3.2336	-0.0590915\\
3.2361	-0.0590429\\
3.2387	-0.0589928\\
3.2412	-0.0589451\\
3.2438	-0.058896\\
3.2464	-0.0588473\\
3.2489	-0.0588009\\
3.2515	-0.0587531\\
3.2541	-0.0587058\\
3.2566	-0.0586608\\
3.2592	-0.0586144\\
3.2618	-0.0585685\\
3.2643	-0.0585248\\
3.2669	-0.0584798\\
3.2694	-0.0584369\\
3.272	-0.0583929\\
3.2746	-0.0583492\\
3.2771	-0.0583077\\
3.2797	-0.058265\\
3.2823	-0.0582228\\
3.2848	-0.0581826\\
3.2874	-0.0581412\\
3.29	-0.0581003\\
3.2925	-0.0580615\\
3.2951	-0.0580215\\
3.2976	-0.0579835\\
3.3002	-0.0579444\\
3.3028	-0.0579057\\
3.3053	-0.057869\\
3.3079	-0.0578312\\
3.3105	-0.0577939\\
3.313	-0.0577585\\
3.3156	-0.057722\\
3.3182	-0.057686\\
3.3207	-0.0576518\\
3.3233	-0.0576167\\
3.3259	-0.057582\\
3.3284	-0.057549\\
3.331	-0.0575151\\
3.3335	-0.057483\\
3.3361	-0.05745\\
3.3387	-0.0574174\\
3.3412	-0.0573864\\
3.3438	-0.0573547\\
3.3464	-0.0573233\\
3.3489	-0.0572935\\
3.3515	-0.057263\\
3.3541	-0.0572328\\
3.3566	-0.0572042\\
3.3592	-0.0571749\\
3.3617	-0.0571471\\
3.3643	-0.0571185\\
3.3669	-0.0570903\\
3.3694	-0.0570636\\
3.372	-0.0570362\\
3.3746	-0.0570091\\
3.3771	-0.0569835\\
3.3797	-0.0569572\\
3.3823	-0.0569313\\
3.3848	-0.0569068\\
3.3874	-0.0568816\\
3.3899	-0.0568577\\
3.3925	-0.0568332\\
3.3951	-0.0568091\\
3.3976	-0.0567863\\
3.4002	-0.0567629\\
3.4028	-0.0567398\\
3.4053	-0.056718\\
3.4079	-0.0566956\\
3.4105	-0.0566735\\
3.413	-0.0566526\\
3.4156	-0.0566312\\
3.4181	-0.056611\\
3.4207	-0.0565902\\
3.4233	-0.0565698\\
3.4258	-0.0565504\\
3.4284	-0.0565306\\
3.431	-0.0565111\\
3.4335	-0.0564927\\
3.4361	-0.0564738\\
3.4387	-0.0564552\\
3.4412	-0.0564376\\
3.4438	-0.0564196\\
3.4463	-0.0564025\\
3.4489	-0.0563851\\
3.4515	-0.0563679\\
3.454	-0.0563517\\
3.4566	-0.056335\\
3.4592	-0.0563187\\
3.4617	-0.0563032\\
3.4643	-0.0562874\\
3.4669	-0.0562719\\
3.4694	-0.0562571\\
3.472	-0.0562421\\
3.4745	-0.0562279\\
3.4771	-0.0562133\\
3.4797	-0.056199\\
3.4822	-0.0561854\\
3.4848	-0.0561716\\
3.4874	-0.056158\\
3.4899	-0.0561451\\
3.4925	-0.056132\\
3.4951	-0.0561191\\
3.4976	-0.0561068\\
3.5002	-0.0560944\\
3.5028	-0.0560821\\
3.5053	-0.0560705\\
3.5079	-0.0560586\\
3.5104	-0.0560474\\
3.513	-0.0560359\\
3.5156	-0.0560247\\
3.5181	-0.056014\\
3.5207	-0.0560032\\
3.5233	-0.0559925\\
3.5258	-0.0559824\\
3.5284	-0.055972\\
3.531	-0.0559619\\
3.5335	-0.0559523\\
3.5361	-0.0559425\\
3.5386	-0.0559333\\
3.5412	-0.0559238\\
3.5438	-0.0559145\\
3.5463	-0.0559057\\
3.5489	-0.0558968\\
3.5515	-0.055888\\
3.554	-0.0558796\\
3.5566	-0.0558711\\
3.5592	-0.0558627\\
3.5617	-0.0558548\\
3.5643	-0.0558468\\
3.5668	-0.0558391\\
3.5694	-0.0558313\\
3.572	-0.0558237\\
3.5745	-0.0558164\\
3.5771	-0.055809\\
3.5797	-0.0558018\\
3.5822	-0.0557949\\
3.5848	-0.0557879\\
3.5874	-0.055781\\
3.5899	-0.0557745\\
3.5925	-0.0557678\\
3.595	-0.0557615\\
3.5976	-0.055755\\
3.6002	-0.0557487\\
3.6027	-0.0557427\\
3.6053	-0.0557366\\
3.6079	-0.0557306\\
3.6104	-0.0557249\\
3.613	-0.0557191\\
3.6156	-0.0557134\\
3.6181	-0.055708\\
3.6207	-0.0557024\\
3.6232	-0.0556972\\
3.6258	-0.0556919\\
3.6284	-0.0556866\\
3.6309	-0.0556816\\
3.6335	-0.0556765\\
3.6361	-0.0556715\\
3.6386	-0.0556668\\
3.6412	-0.055662\\
3.6438	-0.0556572\\
3.6463	-0.0556527\\
3.6489	-0.0556481\\
3.6515	-0.0556436\\
3.654	-0.0556393\\
3.6566	-0.0556349\\
3.6591	-0.0556307\\
3.6617	-0.0556264\\
3.6643	-0.0556223\\
3.6668	-0.0556183\\
3.6694	-0.0556142\\
3.672	-0.0556102\\
3.6745	-0.0556065\\
3.6771	-0.0556026\\
3.6797	-0.0555988\\
3.6822	-0.0555952\\
3.6848	-0.0555915\\
3.6873	-0.055588\\
3.6899	-0.0555844\\
3.6925	-0.0555808\\
3.695	-0.0555775\\
3.6976	-0.0555741\\
3.7002	-0.0555707\\
3.7027	-0.0555675\\
3.7053	-0.0555643\\
3.7079	-0.055561\\
3.7104	-0.055558\\
3.713	-0.0555549\\
3.7155	-0.0555519\\
3.7181	-0.0555489\\
3.7207	-0.0555459\\
3.7232	-0.0555431\\
3.7258	-0.0555402\\
3.7284	-0.0555374\\
3.7309	-0.0555347\\
3.7335	-0.0555319\\
3.7361	-0.0555292\\
3.7386	-0.0555267\\
3.7412	-0.055524\\
3.7437	-0.0555215\\
3.7463	-0.055519\\
3.7489	-0.0555165\\
3.7514	-0.0555142\\
3.754	-0.0555117\\
3.7566	-0.0555094\\
3.7591	-0.0555071\\
3.7617	-0.0555048\\
3.7643	-0.0555026\\
3.7668	-0.0555004\\
3.7694	-0.0554983\\
3.7719	-0.0554962\\
3.7745	-0.0554941\\
3.7771	-0.0554921\\
3.7796	-0.0554901\\
3.7822	-0.0554882\\
3.7848	-0.0554862\\
3.7873	-0.0554844\\
3.7899	-0.0554826\\
3.7925	-0.0554807\\
3.795	-0.055479\\
3.7976	-0.0554773\\
3.8002	-0.0554756\\
3.8027	-0.055474\\
3.8053	-0.0554724\\
3.8078	-0.0554709\\
3.8104	-0.0554694\\
3.813	-0.0554679\\
3.8155	-0.0554665\\
3.8181	-0.0554651\\
3.8207	-0.0554637\\
3.8232	-0.0554624\\
3.8258	-0.0554611\\
3.8284	-0.0554599\\
3.8309	-0.0554587\\
3.8335	-0.0554576\\
3.836	-0.0554565\\
3.8386	-0.0554554\\
3.8412	-0.0554544\\
3.8437	-0.0554535\\
3.8463	-0.0554525\\
3.8489	-0.0554516\\
3.8514	-0.0554508\\
3.854	-0.05545\\
3.8566	-0.0554493\\
3.8591	-0.0554486\\
3.8617	-0.0554479\\
3.8642	-0.0554473\\
3.8668	-0.0554468\\
3.8694	-0.0554462\\
3.8719	-0.0554458\\
3.8745	-0.0554454\\
3.8771	-0.055445\\
3.8796	-0.0554447\\
3.8822	-0.0554444\\
3.8848	-0.0554442\\
3.8873	-0.055444\\
3.8899	-0.0554439\\
3.8924	-0.0554439\\
3.895	-0.0554438\\
3.8976	-0.0554439\\
3.9001	-0.055444\\
3.9027	-0.0554441\\
3.9053	-0.0554443\\
3.9078	-0.0554446\\
3.9104	-0.0554449\\
3.913	-0.0554452\\
3.9155	-0.0554457\\
3.9181	-0.0554461\\
3.9206	-0.0554467\\
3.9232	-0.0554473\\
3.9258	-0.0554479\\
3.9283	-0.0554486\\
3.9309	-0.0554494\\
3.9335	-0.0554502\\
3.936	-0.0554511\\
3.9386	-0.055452\\
3.9412	-0.055453\\
3.9437	-0.0554541\\
3.9463	-0.0554552\\
3.9488	-0.0554564\\
3.9514	-0.0554576\\
3.954	-0.055459\\
3.9565	-0.0554603\\
3.9591	-0.0554618\\
3.9617	-0.0554633\\
3.9642	-0.0554648\\
3.9668	-0.0554665\\
3.9694	-0.0554682\\
3.9719	-0.0554699\\
3.9745	-0.0554718\\
3.9771	-0.0554737\\
3.9796	-0.0554756\\
3.9822	-0.0554776\\
3.9847	-0.0554797\\
3.9873	-0.0554819\\
3.9899	-0.0554842\\
3.9924	-0.0554864\\
3.995	-0.0554888\\
3.9976	-0.0554913\\
4.0001	-0.0554937\\
4.0027	-0.0554963\\
4.0053	-0.055499\\
4.0078	-0.0555017\\
4.0104	-0.0555045\\
4.0129	-0.0555073\\
4.0155	-0.0555103\\
4.0181	-0.0555133\\
4.0206	-0.0555163\\
4.0232	-0.0555195\\
4.0258	-0.0555228\\
4.0283	-0.055526\\
4.0309	-0.0555294\\
4.0335	-0.0555329\\
4.036	-0.0555364\\
4.0386	-0.05554\\
4.0411	-0.0555436\\
4.0437	-0.0555474\\
4.0463	-0.0555512\\
4.0488	-0.055555\\
4.0514	-0.055559\\
4.054	-0.0555631\\
4.0565	-0.0555672\\
4.0591	-0.0555714\\
4.0617	-0.0555757\\
4.0642	-0.05558\\
4.0668	-0.0555844\\
4.0693	-0.0555888\\
4.0719	-0.0555934\\
4.0745	-0.0555981\\
4.077	-0.0556027\\
4.0796	-0.0556076\\
4.0822	-0.0556125\\
4.0847	-0.0556173\\
4.0873	-0.0556224\\
4.0899	-0.0556276\\
4.0924	-0.0556326\\
4.095	-0.055638\\
4.0975	-0.0556431\\
4.1001	-0.0556486\\
4.1027	-0.0556542\\
4.1052	-0.0556596\\
4.1078	-0.0556653\\
4.1104	-0.055671\\
4.1129	-0.0556767\\
4.1155	-0.0556826\\
4.1181	-0.0556886\\
4.1206	-0.0556944\\
4.1232	-0.0557006\\
4.1258	-0.0557068\\
4.1283	-0.0557129\\
4.1309	-0.0557193\\
4.1334	-0.0557255\\
4.136	-0.055732\\
4.1386	-0.0557386\\
4.1411	-0.055745\\
4.1437	-0.0557517\\
4.1463	-0.0557586\\
4.1488	-0.0557652\\
4.1514	-0.0557722\\
4.154	-0.0557792\\
4.1565	-0.055786\\
4.1591	-0.0557932\\
4.1616	-0.0558002\\
4.1642	-0.0558075\\
4.1668	-0.0558149\\
4.1693	-0.0558221\\
4.1719	-0.0558296\\
4.1745	-0.0558372\\
4.177	-0.0558446\\
4.1796	-0.0558523\\
4.1822	-0.0558601\\
4.1847	-0.0558677\\
4.1873	-0.0558757\\
4.1898	-0.0558834\\
4.1924	-0.0558914\\
4.195	-0.0558996\\
4.1975	-0.0559075\\
4.2001	-0.0559157\\
4.2027	-0.0559241\\
4.2052	-0.0559322\\
4.2078	-0.0559406\\
4.2104	-0.0559491\\
4.2129	-0.0559574\\
4.2155	-0.0559661\\
4.218	-0.0559744\\
4.2206	-0.0559832\\
4.2232	-0.055992\\
4.2257	-0.0560006\\
4.2283	-0.0560095\\
4.2309	-0.0560185\\
4.2334	-0.0560272\\
4.236	-0.0560364\\
4.2386	-0.0560455\\
4.2411	-0.0560544\\
4.2437	-0.0560637\\
4.2462	-0.0560727\\
4.2488	-0.0560821\\
4.2514	-0.0560915\\
4.2539	-0.0561007\\
4.2565	-0.0561102\\
4.2591	-0.0561198\\
4.2616	-0.0561291\\
4.2642	-0.0561388\\
4.2668	-0.0561485\\
4.2693	-0.0561579\\
4.2719	-0.0561678\\
4.2744	-0.0561773\\
4.277	-0.0561872\\
4.2796	-0.0561972\\
4.2821	-0.0562068\\
4.2847	-0.0562169\\
4.2873	-0.056227\\
4.2898	-0.0562368\\
4.2924	-0.056247\\
4.295	-0.0562572\\
4.2975	-0.0562671\\
4.3001	-0.0562774\\
4.3027	-0.0562877\\
4.3052	-0.0562977\\
4.3078	-0.0563081\\
4.3103	-0.0563182\\
4.3129	-0.0563287\\
4.3155	-0.0563392\\
4.318	-0.0563493\\
4.3206	-0.0563599\\
4.3232	-0.0563705\\
4.3257	-0.0563808\\
4.3283	-0.0563914\\
4.3309	-0.0564021\\
4.3334	-0.0564124\\
4.336	-0.0564232\\
4.3385	-0.0564336\\
4.3411	-0.0564444\\
4.3437	-0.0564552\\
4.3462	-0.0564656\\
4.3488	-0.0564765\\
4.3514	-0.0564874\\
4.3539	-0.0564979\\
4.3565	-0.0565088\\
4.3591	-0.0565198\\
4.3616	-0.0565303\\
4.3642	-0.0565413\\
4.3667	-0.0565519\\
4.3693	-0.0565629\\
4.3719	-0.056574\\
4.3744	-0.0565846\\
4.377	-0.0565956\\
4.3796	-0.0566067\\
4.3821	-0.0566173\\
4.3847	-0.0566284\\
4.3873	-0.0566395\\
4.3898	-0.0566502\\
4.3924	-0.0566613\\
4.3949	-0.056672\\
4.3975	-0.0566831\\
4.4001	-0.0566942\\
4.4026	-0.0567049\\
4.4052	-0.0567161\\
4.4078	-0.0567272\\
4.4103	-0.0567379\\
4.4129	-0.056749\\
4.4155	-0.0567602\\
4.418	-0.0567709\\
4.4206	-0.056782\\
4.4231	-0.0567927\\
4.4257	-0.0568038\\
4.4283	-0.0568149\\
4.4308	-0.0568255\\
4.4334	-0.0568366\\
4.436	-0.0568477\\
4.4385	-0.0568584\\
4.4411	-0.0568694\\
4.4437	-0.0568805\\
4.4462	-0.0568911\\
4.4488	-0.0569021\\
4.4514	-0.0569131\\
4.4539	-0.0569236\\
4.4565	-0.0569346\\
4.459	-0.0569451\\
4.4616	-0.0569561\\
4.4642	-0.056967\\
4.4667	-0.0569775\\
4.4693	-0.0569883\\
4.4719	-0.0569992\\
4.4744	-0.0570096\\
4.477	-0.0570204\\
4.4796	-0.0570312\\
4.4821	-0.0570415\\
4.4847	-0.0570522\\
4.4872	-0.0570625\\
4.4898	-0.0570732\\
4.4924	-0.0570838\\
4.4949	-0.057094\\
4.4975	-0.0571046\\
4.5001	-0.0571152\\
4.5026	-0.0571253\\
4.5052	-0.0571358\\
4.5078	-0.0571462\\
4.5103	-0.0571562\\
4.5129	-0.0571666\\
4.5154	-0.0571765\\
4.518	-0.0571868\\
4.5206	-0.0571971\\
4.5231	-0.0572069\\
4.5257	-0.0572171\\
4.5283	-0.0572272\\
4.5308	-0.0572369\\
4.5334	-0.057247\\
4.536	-0.057257\\
4.5385	-0.0572666\\
4.5411	-0.0572765\\
4.5436	-0.057286\\
4.5462	-0.0572958\\
4.5488	-0.0573056\\
4.5513	-0.0573149\\
4.5539	-0.0573246\\
4.5565	-0.0573342\\
4.559	-0.0573435\\
4.5616	-0.057353\\
4.5642	-0.0573624\\
4.5667	-0.0573715\\
4.5693	-0.0573809\\
4.5718	-0.0573898\\
4.5744	-0.0573991\\
4.577	-0.0574083\\
4.5795	-0.057417\\
4.5821	-0.0574261\\
4.5847	-0.0574351\\
4.5872	-0.0574438\\
4.5898	-0.0574527\\
4.5924	-0.0574615\\
4.5949	-0.0574699\\
4.5975	-0.0574787\\
4.6001	-0.0574873\\
4.6026	-0.0574956\\
4.6052	-0.0575041\\
4.6077	-0.0575122\\
4.6103	-0.0575206\\
4.6129	-0.0575289\\
4.6154	-0.0575369\\
4.618	-0.0575451\\
4.6206	-0.0575532\\
4.6231	-0.0575609\\
4.6257	-0.0575689\\
4.6283	-0.0575768\\
4.6308	-0.0575843\\
4.6334	-0.0575921\\
4.6359	-0.0575995\\
4.6385	-0.0576071\\
4.6411	-0.0576147\\
4.6436	-0.0576218\\
4.6462	-0.0576292\\
4.6488	-0.0576366\\
4.6513	-0.0576435\\
4.6539	-0.0576507\\
4.6565	-0.0576578\\
4.659	-0.0576645\\
4.6616	-0.0576714\\
4.6641	-0.057678\\
4.6667	-0.0576848\\
4.6693	-0.0576915\\
4.6718	-0.0576978\\
4.6744	-0.0577044\\
4.677	-0.0577108\\
4.6795	-0.0577169\\
4.6821	-0.0577232\\
4.6847	-0.0577294\\
4.6872	-0.0577353\\
4.6898	-0.0577413\\
4.6923	-0.057747\\
4.6949	-0.0577529\\
4.6975	-0.0577586\\
4.7	-0.0577641\\
4.7026	-0.0577697\\
4.7052	-0.0577752\\
4.7077	-0.0577804\\
4.7103	-0.0577857\\
4.7129	-0.057791\\
4.7154	-0.0577959\\
4.718	-0.057801\\
4.7205	-0.0578057\\
4.7231	-0.0578106\\
4.7257	-0.0578154\\
4.7282	-0.0578199\\
4.7308	-0.0578245\\
4.7334	-0.057829\\
4.7359	-0.0578333\\
4.7385	-0.0578376\\
4.7411	-0.0578418\\
4.7436	-0.0578458\\
4.7462	-0.0578498\\
4.7487	-0.0578536\\
4.7513	-0.0578575\\
4.7539	-0.0578612\\
4.7564	-0.0578648\\
4.759	-0.0578683\\
4.7616	-0.0578718\\
4.7641	-0.057875\\
4.7667	-0.0578783\\
4.7693	-0.0578815\\
4.7718	-0.0578844\\
4.7744	-0.0578874\\
4.777	-0.0578903\\
4.7795	-0.0578929\\
4.7821	-0.0578956\\
4.7846	-0.0578981\\
4.7872	-0.0579006\\
4.7898	-0.057903\\
4.7923	-0.0579052\\
4.7949	-0.0579073\\
4.7975	-0.0579094\\
4.8	-0.0579113\\
4.8026	-0.0579132\\
4.8052	-0.057915\\
4.8077	-0.0579166\\
4.8103	-0.0579182\\
4.8128	-0.0579196\\
4.8154	-0.057921\\
4.818	-0.0579222\\
4.8205	-0.0579233\\
4.8231	-0.0579244\\
4.8257	-0.0579254\\
4.8282	-0.0579262\\
4.8308	-0.057927\\
4.8334	-0.0579276\\
4.8359	-0.0579281\\
4.8385	-0.0579286\\
4.841	-0.0579289\\
4.8436	-0.0579291\\
4.8462	-0.0579293\\
4.8487	-0.0579293\\
4.8513	-0.0579292\\
4.8539	-0.0579291\\
4.8564	-0.0579288\\
4.859	-0.0579284\\
4.8616	-0.0579279\\
4.8641	-0.0579274\\
4.8667	-0.0579267\\
4.8692	-0.0579259\\
4.8718	-0.057925\\
4.8744	-0.057924\\
4.8769	-0.0579229\\
4.8795	-0.0579217\\
4.8821	-0.0579204\\
4.8846	-0.057919\\
4.8872	-0.0579175\\
4.8898	-0.0579159\\
4.8923	-0.0579142\\
4.8949	-0.0579124\\
4.8974	-0.0579105\\
4.9	-0.0579084\\
4.9026	-0.0579063\\
4.9051	-0.0579041\\
4.9077	-0.0579018\\
4.9103	-0.0578993\\
4.9128	-0.0578968\\
4.9154	-0.0578942\\
4.918	-0.0578914\\
4.9205	-0.0578886\\
4.9231	-0.0578857\\
4.9257	-0.0578826\\
4.9282	-0.0578795\\
4.9308	-0.0578762\\
4.9333	-0.057873\\
4.9359	-0.0578695\\
4.9385	-0.0578659\\
4.941	-0.0578623\\
4.9436	-0.0578586\\
4.9462	-0.0578547\\
4.9487	-0.0578508\\
4.9513	-0.0578467\\
4.9539	-0.0578425\\
4.9564	-0.0578384\\
4.959	-0.057834\\
4.9615	-0.0578297\\
4.9641	-0.0578251\\
4.9667	-0.0578204\\
4.9692	-0.0578157\\
4.9718	-0.0578109\\
4.9744	-0.0578059\\
4.9769	-0.057801\\
4.9795	-0.0577958\\
4.9821	-0.0577905\\
4.9846	-0.0577853\\
4.9872	-0.0577798\\
4.9897	-0.0577744\\
4.9923	-0.0577688\\
4.9949	-0.057763\\
4.9974	-0.0577573\\
5	-0.0577514\\
};
\addplot [color=mycolor23, line width=1.0pt, forget plot]
  table[row sep=crcr]{%
-0.125	6.91544e-08\\
-0.1224	7.05853e-08\\
-0.1199	7.1962e-08\\
-0.1173	7.33938e-08\\
-0.1147	7.48242e-08\\
-0.1122	7.62005e-08\\
-0.1096	7.76324e-08\\
-0.1071	7.90076e-08\\
-0.1045	8.04374e-08\\
-0.1019	8.18684e-08\\
-0.0994	8.32436e-08\\
-0.0968	8.46736e-08\\
-0.0942	8.6105e-08\\
-0.0917	8.74805e-08\\
-0.0891	8.89095e-08\\
-0.0865	9.03396e-08\\
-0.084	9.17149e-08\\
-0.0814	9.3144e-08\\
-0.0789	9.4519e-08\\
-0.0763	9.59496e-08\\
-0.0737	9.73784e-08\\
-0.0712	9.87522e-08\\
-0.0686	1.00182e-07\\
-0.066	1.01611e-07\\
-0.0635	1.02985e-07\\
-0.0609	1.04415e-07\\
-0.0583	1.05844e-07\\
-0.0558	1.07217e-07\\
-0.0532	1.08646e-07\\
-0.0507	1.1002e-07\\
-0.0481	1.11448e-07\\
-0.0455	1.12877e-07\\
-0.043	1.14251e-07\\
-0.0404	1.15679e-07\\
-0.0378	1.17107e-07\\
-0.0353	1.18481e-07\\
-0.0327	1.19909e-07\\
-0.0301	1.21336e-07\\
-0.0276	1.2271e-07\\
-0.025	1.24138e-07\\
-0.0224	1.25565e-07\\
-0.0199	1.26938e-07\\
-0.0173	1.28365e-07\\
-0.0148	1.29737e-07\\
-0.0122	1.31165e-07\\
-0.0096	1.32593e-07\\
-0.0071	1.33965e-07\\
-0.0045	1.35391e-07\\
-0.0019	1.36819e-07\\
0.0006	1.3819e-07\\
0.0032	-0.000301668\\
0.0058	-0.000675854\\
0.0083	0.00572706\\
0.0109	0.0202205\\
0.0134	0.0359287\\
0.016	0.0483612\\
0.0186	0.0555119\\
0.0211	0.0596826\\
0.0237	0.0640634\\
0.0263	0.069571\\
0.0288	0.0752579\\
0.0314	0.0804571\\
0.034	0.0844313\\
0.0365	0.0872617\\
0.0391	0.0894573\\
0.0416	0.0908929\\
0.0442	0.0916238\\
0.0468	0.0917002\\
0.0493	0.0914559\\
0.0519	0.0910764\\
0.0545	0.090402\\
0.057	0.0890642\\
0.0596	0.0866686\\
0.0622	0.083431\\
0.0647	0.0800528\\
0.0673	0.0767426\\
0.0698	0.0737612\\
0.0724	0.0704133\\
0.075	0.0663048\\
0.0775	0.0615643\\
0.0801	0.0562619\\
0.0827	0.0512043\\
0.0852	0.0468365\\
0.0878	0.0425376\\
0.0904	0.0379283\\
0.0929	0.0328927\\
0.0955	0.0272323\\
0.098	0.021925\\
0.1006	0.0170163\\
0.1032	0.0127198\\
0.1057	0.00873522\\
0.1083	0.00426795\\
0.1109	-0.000604417\\
0.1134	-0.00529072\\
0.116	-0.00963044\\
0.1186	-0.0131967\\
0.1211	-0.016142\\
0.1237	-0.0191994\\
0.1263	-0.0225529\\
0.1288	-0.0259043\\
0.1314	-0.0290449\\
0.1339	-0.031367\\
0.1365	-0.0330811\\
0.1391	-0.0344917\\
0.1416	-0.0359712\\
0.1442	-0.0377218\\
0.1468	-0.039362\\
0.1493	-0.0404133\\
0.1519	-0.0407831\\
0.1545	-0.040654\\
0.157	-0.0404802\\
0.1596	-0.0405173\\
0.1621	-0.0406515\\
0.1647	-0.0405065\\
0.1673	-0.0397506\\
0.1698	-0.0384779\\
0.1724	-0.0369388\\
0.175	-0.0355586\\
0.1775	-0.0344622\\
0.1801	-0.033302\\
0.1827	-0.0317531\\
0.1852	-0.0297596\\
0.1878	-0.0273664\\
0.1903	-0.0251164\\
0.1929	-0.0230674\\
0.1955	-0.0212217\\
0.198	-0.0193284\\
0.2006	-0.016973\\
0.2032	-0.0142558\\
0.2057	-0.0115906\\
0.2083	-0.00908081\\
0.2109	-0.00690932\\
0.2134	-0.00492954\\
0.216	-0.0026764\\
0.2185	-0.000207367\\
0.2211	0.00250547\\
0.2237	0.00504197\\
0.2262	0.00712339\\
0.2288	0.00899741\\
0.2314	0.0108581\\
0.2339	0.0128518\\
0.2365	0.0151137\\
0.2391	0.0173104\\
0.2416	0.0191077\\
0.2442	0.0205961\\
0.2467	0.0218436\\
0.2493	0.0232156\\
0.2519	0.0247729\\
0.2544	0.0263106\\
0.257	0.0276869\\
0.2596	0.0286728\\
0.2621	0.0293256\\
0.2647	0.0299398\\
0.2673	0.0306864\\
0.2698	0.031513\\
0.2724	0.0322761\\
0.2749	0.0327102\\
0.2775	0.0328099\\
0.2801	0.0327235\\
0.2826	0.032686\\
0.2852	0.0327888\\
0.2878	0.0329172\\
0.2903	0.0328536\\
0.2929	0.0324633\\
0.2955	0.0318163\\
0.298	0.0311538\\
0.3006	0.0305907\\
0.3032	0.0301421\\
0.3057	0.0296507\\
0.3083	0.0288981\\
0.3108	0.0279159\\
0.3134	0.026773\\
0.316	0.0257068\\
0.3185	0.0248318\\
0.3211	0.0239801\\
0.3237	0.0230005\\
0.3262	0.0218352\\
0.3288	0.0204548\\
0.3314	0.0190854\\
0.3339	0.0179159\\
0.3365	0.0168408\\
0.339	0.0158024\\
0.3416	0.0145662\\
0.3442	0.013148\\
0.3467	0.0117328\\
0.3493	0.0103734\\
0.3519	0.00919579\\
0.3544	0.00815442\\
0.357	0.00700982\\
0.3596	0.00571112\\
0.3621	0.00436847\\
0.3647	0.00302492\\
0.3672	0.00189671\\
0.3698	0.000883102\\
0.3724	-9.62597e-05\\
0.3749	-0.00112988\\
0.3775	-0.00232055\\
0.3801	-0.00352224\\
0.3826	-0.00455531\\
0.3852	-0.00544868\\
0.3878	-0.0062327\\
0.3903	-0.00701194\\
0.3929	-0.00792618\\
0.3954	-0.00886386\\
0.398	-0.00977507\\
0.4006	-0.0105209\\
0.4031	-0.0110924\\
0.4057	-0.0116441\\
0.4083	-0.0122648\\
0.4108	-0.0129449\\
0.4134	-0.0136516\\
0.416	-0.0142408\\
0.4185	-0.0146546\\
0.4211	-0.0149901\\
0.4236	-0.0153307\\
0.4262	-0.0157713\\
0.4288	-0.0162644\\
0.4313	-0.0166878\\
0.4339	-0.0169935\\
0.4365	-0.0171734\\
0.439	-0.0173174\\
0.4416	-0.0175315\\
0.4442	-0.0178276\\
0.4467	-0.01812\\
0.4493	-0.0183356\\
0.4519	-0.018425\\
0.4544	-0.0184446\\
0.457	-0.0184939\\
0.4595	-0.0186236\\
0.4621	-0.0188149\\
0.4647	-0.0189739\\
0.4672	-0.0190302\\
0.4698	-0.0189981\\
0.4724	-0.0189543\\
0.4749	-0.0189786\\
0.4775	-0.0190866\\
0.4801	-0.0192158\\
0.4826	-0.0192816\\
0.4852	-0.0192591\\
0.4877	-0.019194\\
0.4903	-0.0191638\\
0.4929	-0.0192212\\
0.4954	-0.0193351\\
0.498	-0.0194396\\
0.5006	-0.0194702\\
0.5031	-0.0194375\\
0.5057	-0.0194057\\
0.5083	-0.0194449\\
0.5108	-0.0195604\\
0.5134	-0.0197082\\
0.5159	-0.0198106\\
0.5185	-0.0198489\\
0.5211	-0.019857\\
0.5236	-0.019903\\
0.5262	-0.0200312\\
0.5288	-0.0202186\\
0.5313	-0.0203929\\
0.5339	-0.0205178\\
0.5365	-0.0205917\\
0.539	-0.0206685\\
0.5416	-0.020811\\
0.5441	-0.0210169\\
0.5467	-0.0212572\\
0.5493	-0.0214634\\
0.5518	-0.0216069\\
0.5544	-0.0217313\\
0.557	-0.0218896\\
0.5595	-0.0221067\\
0.5621	-0.0223803\\
0.5647	-0.0226468\\
0.5672	-0.0228547\\
0.5698	-0.0230256\\
0.5723	-0.023191\\
0.5749	-0.0234121\\
0.5775	-0.0236907\\
0.58	-0.0239757\\
0.5826	-0.0242389\\
0.5852	-0.0244478\\
0.5877	-0.0246233\\
0.5903	-0.0248287\\
0.5929	-0.0250862\\
0.5954	-0.0253672\\
0.598	-0.0256473\\
0.6006	-0.0258778\\
0.6031	-0.0260563\\
0.6057	-0.0262372\\
0.6082	-0.0264448\\
0.6108	-0.0267022\\
0.6134	-0.0269685\\
0.6159	-0.0271913\\
0.6185	-0.0273715\\
0.6211	-0.0275224\\
0.6236	-0.0276802\\
0.6262	-0.0278818\\
0.6288	-0.0281073\\
0.6313	-0.0283088\\
0.6339	-0.0284715\\
0.6364	-0.028586\\
0.639	-0.0286954\\
0.6416	-0.02883\\
0.6441	-0.0289889\\
0.6467	-0.0291572\\
0.6493	-0.0292916\\
0.6518	-0.0293753\\
0.6544	-0.0294346\\
0.657	-0.0295018\\
0.6595	-0.0295941\\
0.6621	-0.0297076\\
0.6646	-0.0298022\\
0.6672	-0.0298592\\
0.6698	-0.029877\\
0.6723	-0.0298849\\
0.6749	-0.0299131\\
0.6775	-0.029967\\
0.68	-0.0300211\\
0.6826	-0.0300492\\
0.6852	-0.0300369\\
0.6877	-0.0300028\\
0.6903	-0.0299749\\
0.6928	-0.0299724\\
0.6954	-0.0299864\\
0.698	-0.029989\\
0.7005	-0.0299598\\
0.7031	-0.029899\\
0.7057	-0.0298321\\
0.7082	-0.0297865\\
0.7108	-0.0297631\\
0.7134	-0.0297449\\
0.7159	-0.0297075\\
0.7185	-0.0296382\\
0.721	-0.0295558\\
0.7236	-0.0294794\\
0.7262	-0.0294284\\
0.7287	-0.0293964\\
0.7313	-0.0293576\\
0.7339	-0.0292944\\
0.7364	-0.0292125\\
0.739	-0.0291262\\
0.7416	-0.0290604\\
0.7441	-0.0290204\\
0.7467	-0.0289871\\
0.7492	-0.0289424\\
0.7518	-0.0288741\\
0.7544	-0.0287953\\
0.7569	-0.0287306\\
0.7595	-0.0286881\\
0.7621	-0.0286644\\
0.7646	-0.0286406\\
0.7672	-0.0285989\\
0.7698	-0.0285417\\
0.7723	-0.0284886\\
0.7749	-0.0284532\\
0.7775	-0.0284416\\
0.78	-0.0284398\\
0.7826	-0.0284291\\
0.7851	-0.0284034\\
0.7877	-0.0283701\\
0.7903	-0.0283487\\
0.7928	-0.0283499\\
0.7954	-0.0283685\\
0.798	-0.0283877\\
0.8005	-0.0283939\\
0.8031	-0.0283885\\
0.8057	-0.0283857\\
0.8082	-0.0283999\\
0.8108	-0.0284349\\
0.8133	-0.0284771\\
0.8159	-0.0285142\\
0.8185	-0.0285377\\
0.821	-0.028555\\
0.8236	-0.0285824\\
0.8262	-0.0286289\\
0.8287	-0.0286872\\
0.8313	-0.0287479\\
0.8339	-0.028797\\
0.8364	-0.0288337\\
0.839	-0.028873\\
0.8415	-0.0289238\\
0.8441	-0.0289925\\
0.8467	-0.0290677\\
0.8492	-0.0291332\\
0.8518	-0.0291886\\
0.8544	-0.0292377\\
0.8569	-0.029291\\
0.8595	-0.0293606\\
0.8621	-0.0294406\\
0.8646	-0.0295161\\
0.8672	-0.0295832\\
0.8697	-0.0296372\\
0.8723	-0.0296924\\
0.8749	-0.0297572\\
0.8774	-0.0298307\\
0.88	-0.0299108\\
0.8826	-0.0299833\\
0.8851	-0.0300411\\
0.8877	-0.0300943\\
0.8903	-0.0301512\\
0.8928	-0.0302157\\
0.8954	-0.03029\\
0.8979	-0.030359\\
0.9005	-0.0304199\\
0.9031	-0.0304701\\
0.9056	-0.0305162\\
0.9082	-0.0305709\\
0.9108	-0.0306343\\
0.9133	-0.0306973\\
0.9159	-0.0307557\\
0.9185	-0.0308026\\
0.921	-0.0308412\\
0.9236	-0.0308837\\
0.9262	-0.0309341\\
0.9287	-0.030988\\
0.9313	-0.0310413\\
0.9338	-0.0310835\\
0.9364	-0.0311182\\
0.939	-0.0311507\\
0.9415	-0.0311872\\
0.9441	-0.0312323\\
0.9467	-0.0312792\\
0.9492	-0.0313187\\
0.9518	-0.0313504\\
0.9544	-0.0313767\\
0.9569	-0.0314039\\
0.9595	-0.0314391\\
0.962	-0.0314779\\
0.9646	-0.0315166\\
0.9672	-0.0315479\\
0.9697	-0.031571\\
0.9723	-0.0315938\\
0.9749	-0.0316217\\
0.9774	-0.031655\\
0.98	-0.0316919\\
0.9826	-0.0317245\\
0.9851	-0.0317491\\
0.9877	-0.0317709\\
0.9902	-0.031794\\
0.9928	-0.0318246\\
0.9954	-0.0318603\\
0.9979	-0.0318942\\
1.0005	-0.0319241\\
1.0031	-0.0319486\\
1.0056	-0.0319718\\
1.0082	-0.0320008\\
1.0108	-0.0320362\\
1.0133	-0.0320727\\
1.0159	-0.032108\\
1.0184	-0.0321368\\
1.021	-0.0321641\\
1.0236	-0.0321939\\
1.0261	-0.0322282\\
1.0287	-0.0322687\\
1.0313	-0.0323093\\
1.0338	-0.0323444\\
1.0364	-0.032377\\
1.039	-0.0324094\\
1.0415	-0.0324448\\
1.0441	-0.0324871\\
1.0466	-0.0325302\\
1.0492	-0.032573\\
1.0518	-0.0326114\\
1.0543	-0.0326463\\
1.0569	-0.0326845\\
1.0595	-0.0327278\\
1.062	-0.0327732\\
1.0646	-0.0328205\\
1.0672	-0.0328643\\
1.0697	-0.032903\\
1.0723	-0.0329429\\
1.0748	-0.0329845\\
1.0774	-0.0330321\\
1.08	-0.0330816\\
1.0825	-0.0331273\\
1.0851	-0.0331709\\
1.0877	-0.0332122\\
1.0902	-0.0332531\\
1.0928	-0.0332993\\
1.0954	-0.0333485\\
1.0979	-0.0333954\\
1.1005	-0.0334409\\
1.1031	-0.0334829\\
1.1056	-0.0335225\\
1.1082	-0.033566\\
1.1107	-0.0336109\\
1.1133	-0.0336586\\
1.1159	-0.0337041\\
1.1184	-0.0337443\\
1.121	-0.0337837\\
1.1236	-0.0338237\\
1.1261	-0.0338646\\
1.1287	-0.0339091\\
1.1313	-0.0339528\\
1.1338	-0.0339918\\
1.1364	-0.0340291\\
1.1389	-0.0340638\\
1.1415	-0.0341014\\
1.1441	-0.0341413\\
1.1466	-0.0341802\\
1.1492	-0.0342186\\
1.1518	-0.0342536\\
1.1543	-0.0342849\\
1.1569	-0.0343177\\
1.1595	-0.0343524\\
1.162	-0.0343872\\
1.1646	-0.0344226\\
1.1671	-0.034454\\
1.1697	-0.0344838\\
1.1723	-0.0345125\\
1.1748	-0.0345411\\
1.1774	-0.0345728\\
1.18	-0.0346049\\
1.1825	-0.0346341\\
1.1851	-0.0346618\\
1.1877	-0.0346875\\
1.1902	-0.0347124\\
1.1928	-0.03474\\
1.1953	-0.0347678\\
1.1979	-0.0347964\\
1.2005	-0.0348227\\
1.203	-0.034846\\
1.2056	-0.0348694\\
1.2082	-0.034894\\
1.2107	-0.0349193\\
1.2133	-0.0349464\\
1.2159	-0.0349723\\
1.2184	-0.0349953\\
1.221	-0.0350179\\
1.2235	-0.03504\\
1.2261	-0.0350647\\
1.2287	-0.0350909\\
1.2312	-0.0351161\\
1.2338	-0.0351409\\
1.2364	-0.0351642\\
1.2389	-0.0351863\\
1.2415	-0.0352106\\
1.2441	-0.0352368\\
1.2466	-0.0352628\\
1.2492	-0.0352894\\
1.2518	-0.0353146\\
1.2543	-0.0353381\\
1.2569	-0.0353632\\
1.2594	-0.0353891\\
1.262	-0.0354176\\
1.2646	-0.0354465\\
1.2671	-0.0354734\\
1.2697	-0.0355005\\
1.2723	-0.0355277\\
1.2748	-0.0355552\\
1.2774	-0.0355855\\
1.28	-0.035617\\
1.2825	-0.0356472\\
1.2851	-0.0356777\\
1.2876	-0.0357066\\
1.2902	-0.0357376\\
1.2928	-0.0357702\\
1.2953	-0.0358031\\
1.2979	-0.0358379\\
1.3005	-0.0358723\\
1.303	-0.0359047\\
1.3056	-0.0359387\\
1.3082	-0.035974\\
1.3107	-0.0360096\\
1.3133	-0.0360477\\
1.3158	-0.0360844\\
1.3184	-0.0361221\\
1.321	-0.0361596\\
1.3235	-0.0361964\\
1.3261	-0.0362362\\
1.3287	-0.0362775\\
1.3312	-0.0363177\\
1.3338	-0.0363593\\
1.3364	-0.0364005\\
1.3389	-0.0364405\\
1.3415	-0.0364831\\
1.344	-0.0365256\\
1.3466	-0.0365709\\
1.3492	-0.0366162\\
1.3517	-0.0366595\\
1.3543	-0.0367045\\
1.3569	-0.0367501\\
1.3594	-0.0367954\\
1.362	-0.0368437\\
1.3646	-0.0368926\\
1.3671	-0.0369394\\
1.3697	-0.0369879\\
1.3722	-0.0370349\\
1.3748	-0.0370847\\
1.3774	-0.0371359\\
1.3799	-0.037186\\
1.3825	-0.0372384\\
1.3851	-0.0372905\\
1.3876	-0.0373406\\
1.3902	-0.0373934\\
1.3928	-0.0374475\\
1.3953	-0.0375006\\
1.3979	-0.0375564\\
1.4005	-0.0376121\\
1.403	-0.0376657\\
1.4056	-0.0377217\\
1.4081	-0.0377765\\
1.4107	-0.0378347\\
1.4133	-0.0378939\\
1.4158	-0.037951\\
1.4184	-0.0380104\\
1.421	-0.03807\\
1.4235	-0.0381278\\
1.4261	-0.0381892\\
1.4287	-0.0382517\\
1.4312	-0.0383124\\
1.4338	-0.0383756\\
1.4363	-0.0384365\\
1.4389	-0.0385003\\
1.4415	-0.038565\\
1.444	-0.0386283\\
1.4466	-0.0386951\\
1.4492	-0.0387624\\
1.4517	-0.0388272\\
1.4543	-0.0388948\\
1.4569	-0.0389632\\
1.4594	-0.03903\\
1.462	-0.0391005\\
1.4645	-0.0391691\\
1.4671	-0.0392407\\
1.4697	-0.0393124\\
1.4722	-0.039382\\
1.4748	-0.0394552\\
1.4774	-0.0395297\\
1.4799	-0.0396021\\
1.4825	-0.039678\\
1.4851	-0.0397542\\
1.4876	-0.0398277\\
1.4902	-0.039905\\
1.4927	-0.0399803\\
1.4953	-0.0400598\\
1.4979	-0.0401399\\
1.5004	-0.0402174\\
1.503	-0.0402983\\
1.5056	-0.0403798\\
1.5081	-0.040459\\
1.5107	-0.0405425\\
1.5133	-0.0406269\\
1.5158	-0.0407086\\
1.5184	-0.0407939\\
1.5209	-0.0408763\\
1.5235	-0.0409627\\
1.5261	-0.0410502\\
1.5286	-0.0411353\\
1.5312	-0.0412244\\
1.5338	-0.0413141\\
1.5363	-0.0414006\\
1.5389	-0.0414911\\
1.5415	-0.0415824\\
1.544	-0.0416712\\
1.5466	-0.0417644\\
1.5491	-0.0418545\\
1.5517	-0.0419486\\
1.5543	-0.0420431\\
1.5568	-0.0421345\\
1.5594	-0.0422305\\
1.562	-0.0423274\\
1.5645	-0.0424212\\
1.5671	-0.0425191\\
1.5697	-0.0426173\\
1.5722	-0.0427123\\
1.5748	-0.0428117\\
1.5774	-0.042912\\
1.5799	-0.0430091\\
1.5825	-0.0431105\\
1.585	-0.0432083\\
1.5876	-0.0433104\\
1.5902	-0.043413\\
1.5927	-0.0435124\\
1.5953	-0.0436164\\
1.5979	-0.043721\\
1.6004	-0.0438219\\
1.603	-0.043927\\
1.6056	-0.0440326\\
1.6081	-0.0441346\\
1.6107	-0.0442413\\
1.6132	-0.0443446\\
1.6158	-0.0444522\\
1.6184	-0.0445602\\
1.6209	-0.0446641\\
1.6235	-0.0447727\\
1.6261	-0.0448818\\
1.6286	-0.0449873\\
1.6312	-0.0450974\\
1.6338	-0.0452078\\
1.6363	-0.045314\\
1.6389	-0.0454247\\
1.6414	-0.0455316\\
1.644	-0.0456433\\
1.6466	-0.0457555\\
1.6491	-0.0458635\\
1.6517	-0.045976\\
1.6543	-0.0460887\\
1.6568	-0.0461973\\
1.6594	-0.0463107\\
1.662	-0.0464245\\
1.6645	-0.0465342\\
1.6671	-0.0466484\\
1.6696	-0.0467583\\
1.6722	-0.0468728\\
1.6748	-0.0469876\\
1.6773	-0.0470983\\
1.6799	-0.0472137\\
1.6825	-0.0473294\\
1.685	-0.0474406\\
1.6876	-0.0475563\\
1.6902	-0.0476722\\
1.6927	-0.047784\\
1.6953	-0.0479005\\
1.6978	-0.0480127\\
1.7004	-0.0481295\\
1.703	-0.0482463\\
1.7055	-0.0483586\\
1.7081	-0.0484757\\
1.7107	-0.0485929\\
1.7132	-0.0487059\\
1.7158	-0.0488235\\
1.7184	-0.0489411\\
1.7209	-0.0490541\\
1.7235	-0.0491717\\
1.7261	-0.0492896\\
1.7286	-0.049403\\
1.7312	-0.0495211\\
1.7337	-0.0496347\\
1.7363	-0.0497527\\
1.7389	-0.0498708\\
1.7414	-0.0499844\\
1.744	-0.0501026\\
1.7466	-0.050221\\
1.7491	-0.0503349\\
1.7517	-0.0504532\\
1.7543	-0.0505714\\
1.7568	-0.0506851\\
1.7594	-0.0508034\\
1.7619	-0.0509172\\
1.7645	-0.0510356\\
1.7671	-0.0511539\\
1.7696	-0.0512676\\
1.7722	-0.0513857\\
1.7748	-0.0515039\\
1.7773	-0.0516175\\
1.7799	-0.0517357\\
1.7825	-0.0518538\\
1.785	-0.0519672\\
1.7876	-0.052085\\
1.7901	-0.0521983\\
1.7927	-0.052316\\
1.7953	-0.0524337\\
1.7978	-0.0525469\\
1.8004	-0.0526644\\
1.803	-0.0527817\\
1.8055	-0.0528944\\
1.8081	-0.0530115\\
1.8107	-0.0531285\\
1.8132	-0.053241\\
1.8158	-0.0533578\\
1.8183	-0.05347\\
1.8209	-0.0535864\\
1.8235	-0.0537027\\
1.826	-0.0538144\\
1.8286	-0.0539305\\
1.8312	-0.0540464\\
1.8337	-0.0541577\\
1.8363	-0.0542732\\
1.8389	-0.0543884\\
1.8414	-0.0544991\\
1.844	-0.0546141\\
1.8465	-0.0547245\\
1.8491	-0.0548392\\
1.8517	-0.0549535\\
1.8542	-0.0550633\\
1.8568	-0.0551772\\
1.8594	-0.0552909\\
1.8619	-0.0554001\\
1.8645	-0.0555134\\
1.8671	-0.0556265\\
1.8696	-0.055735\\
1.8722	-0.0558475\\
1.8747	-0.0559555\\
1.8773	-0.0560677\\
1.8799	-0.0561796\\
1.8824	-0.0562869\\
1.885	-0.0563983\\
1.8876	-0.0565093\\
1.8901	-0.0566159\\
1.8927	-0.0567264\\
1.8953	-0.0568368\\
1.8978	-0.0569426\\
1.9004	-0.0570524\\
1.903	-0.0571618\\
1.9055	-0.0572668\\
1.9081	-0.0573756\\
1.9106	-0.0574801\\
1.9132	-0.0575885\\
1.9158	-0.0576965\\
1.9183	-0.0578001\\
1.9209	-0.0579075\\
1.9235	-0.0580146\\
1.926	-0.0581173\\
1.9286	-0.0582239\\
1.9312	-0.0583302\\
1.9337	-0.058432\\
1.9363	-0.0585376\\
1.9388	-0.0586388\\
1.9414	-0.0587437\\
1.944	-0.0588484\\
1.9465	-0.0589488\\
1.9491	-0.0590528\\
1.9517	-0.0591565\\
1.9542	-0.0592558\\
1.9568	-0.0593588\\
1.9594	-0.0594615\\
1.9619	-0.0595599\\
1.9645	-0.059662\\
1.967	-0.0597598\\
1.9696	-0.0598612\\
1.9722	-0.0599621\\
1.9747	-0.0600589\\
1.9773	-0.0601593\\
1.9799	-0.0602593\\
1.9824	-0.0603551\\
1.985	-0.0604544\\
1.9876	-0.0605533\\
1.9901	-0.0606481\\
1.9927	-0.0607463\\
1.9952	-0.0608405\\
1.9978	-0.060938\\
2.0004	-0.0610352\\
2.0029	-0.0611283\\
2.0055	-0.0612247\\
2.0081	-0.0613208\\
2.0106	-0.0614129\\
2.0132	-0.0615083\\
2.0158	-0.0616033\\
2.0183	-0.0616943\\
2.0209	-0.0617886\\
2.0234	-0.0618788\\
2.026	-0.0619723\\
2.0286	-0.0620655\\
2.0311	-0.0621548\\
2.0337	-0.0622472\\
2.0363	-0.0623392\\
2.0388	-0.0624273\\
2.0414	-0.0625185\\
2.044	-0.0626094\\
2.0465	-0.0626965\\
2.0491	-0.0627866\\
2.0517	-0.0628763\\
2.0542	-0.0629622\\
2.0568	-0.0630512\\
2.0593	-0.0631363\\
2.0619	-0.0632245\\
2.0645	-0.0633123\\
2.067	-0.0633963\\
2.0696	-0.0634833\\
2.0722	-0.0635699\\
2.0747	-0.0636528\\
2.0773	-0.0637385\\
2.0799	-0.0638239\\
2.0824	-0.0639057\\
2.085	-0.0639902\\
2.0875	-0.0640712\\
2.0901	-0.0641549\\
2.0927	-0.0642383\\
2.0952	-0.064318\\
2.0978	-0.0644006\\
2.1004	-0.0644827\\
2.1029	-0.0645612\\
2.1055	-0.0646425\\
2.1081	-0.0647233\\
2.1106	-0.0648006\\
2.1132	-0.0648807\\
2.1157	-0.0649572\\
2.1183	-0.0650364\\
2.1209	-0.0651151\\
2.1234	-0.0651904\\
2.126	-0.0652682\\
2.1286	-0.0653457\\
2.1311	-0.0654197\\
2.1337	-0.0654963\\
2.1363	-0.0655724\\
2.1388	-0.0656452\\
2.1414	-0.0657204\\
2.1439	-0.0657924\\
2.1465	-0.0658668\\
2.1491	-0.0659407\\
2.1516	-0.0660114\\
2.1542	-0.0660844\\
2.1568	-0.0661569\\
2.1593	-0.0662263\\
2.1619	-0.066298\\
2.1645	-0.0663692\\
2.167	-0.0664373\\
2.1696	-0.0665076\\
2.1721	-0.0665747\\
2.1747	-0.0666441\\
2.1773	-0.066713\\
2.1798	-0.0667789\\
2.1824	-0.0668469\\
2.185	-0.0669144\\
2.1875	-0.0669789\\
2.1901	-0.0670455\\
2.1927	-0.0671116\\
2.1952	-0.0671747\\
2.1978	-0.0672399\\
2.2004	-0.0673046\\
2.2029	-0.0673663\\
2.2055	-0.0674301\\
2.208	-0.0674909\\
2.2106	-0.0675537\\
2.2132	-0.067616\\
2.2157	-0.0676754\\
2.2183	-0.0677368\\
2.2209	-0.0677976\\
2.2234	-0.0678556\\
2.226	-0.0679155\\
2.2286	-0.0679748\\
2.2311	-0.0680314\\
2.2337	-0.0680898\\
2.2362	-0.0681455\\
2.2388	-0.0682028\\
2.2414	-0.0682597\\
2.2439	-0.0683139\\
2.2465	-0.0683698\\
2.2491	-0.0684252\\
2.2516	-0.0684779\\
2.2542	-0.0685323\\
2.2568	-0.0685861\\
2.2593	-0.0686374\\
2.2619	-0.0686902\\
2.2644	-0.0687406\\
2.267	-0.0687924\\
2.2696	-0.0688436\\
2.2721	-0.0688925\\
2.2747	-0.0689427\\
2.2773	-0.0689924\\
2.2798	-0.0690397\\
2.2824	-0.0690884\\
2.285	-0.0691366\\
2.2875	-0.0691824\\
2.2901	-0.0692295\\
2.2926	-0.0692743\\
2.2952	-0.0693203\\
2.2978	-0.0693658\\
2.3003	-0.0694091\\
2.3029	-0.0694535\\
2.3055	-0.0694975\\
2.308	-0.0695392\\
2.3106	-0.069582\\
2.3132	-0.0696243\\
2.3157	-0.0696645\\
2.3183	-0.0697057\\
2.3208	-0.0697448\\
2.3234	-0.0697849\\
2.326	-0.0698245\\
2.3285	-0.0698621\\
2.3311	-0.0699006\\
2.3337	-0.0699385\\
2.3362	-0.0699745\\
2.3388	-0.0700113\\
2.3414	-0.0700476\\
2.3439	-0.070082\\
2.3465	-0.0701172\\
2.349	-0.0701505\\
2.3516	-0.0701846\\
2.3542	-0.0702181\\
2.3567	-0.0702498\\
2.3593	-0.0702822\\
2.3619	-0.070314\\
2.3644	-0.0703441\\
2.367	-0.0703749\\
2.3696	-0.070405\\
2.3721	-0.0704335\\
2.3747	-0.0704625\\
2.3773	-0.0704909\\
2.3798	-0.0705178\\
2.3824	-0.0705451\\
2.3849	-0.0705709\\
2.3875	-0.0705971\\
2.3901	-0.0706227\\
2.3926	-0.0706468\\
2.3952	-0.0706713\\
2.3978	-0.0706952\\
2.4003	-0.0707176\\
2.4029	-0.0707404\\
2.4055	-0.0707626\\
2.408	-0.0707834\\
2.4106	-0.0708045\\
2.4131	-0.0708241\\
2.4157	-0.0708441\\
2.4183	-0.0708634\\
2.4208	-0.0708814\\
2.4234	-0.0708996\\
2.426	-0.0709172\\
2.4285	-0.0709336\\
2.4311	-0.07095\\
2.4337	-0.0709659\\
2.4362	-0.0709806\\
2.4388	-0.0709953\\
2.4413	-0.0710089\\
2.4439	-0.0710225\\
2.4465	-0.0710355\\
2.449	-0.0710474\\
2.4516	-0.0710592\\
2.4542	-0.0710705\\
2.4567	-0.0710807\\
2.4593	-0.0710908\\
2.4619	-0.0711003\\
2.4644	-0.0711089\\
2.467	-0.0711173\\
2.4695	-0.0711248\\
2.4721	-0.071132\\
2.4747	-0.0711386\\
2.4772	-0.0711444\\
2.4798	-0.0711499\\
2.4824	-0.0711548\\
2.4849	-0.0711589\\
2.4875	-0.0711627\\
2.4901	-0.0711658\\
2.4926	-0.0711683\\
2.4952	-0.0711703\\
2.4977	-0.0711717\\
2.5003	-0.0711726\\
2.5029	-0.0711729\\
2.5054	-0.0711726\\
2.508	-0.0711718\\
2.5106	-0.0711703\\
2.5131	-0.0711684\\
2.5157	-0.0711659\\
2.5183	-0.0711627\\
2.5208	-0.0711592\\
2.5234	-0.0711549\\
2.526	-0.07115\\
2.5285	-0.0711448\\
2.5311	-0.0711389\\
2.5336	-0.0711326\\
2.5362	-0.0711255\\
2.5388	-0.0711178\\
2.5413	-0.0711099\\
2.5439	-0.0711011\\
2.5465	-0.0710918\\
2.549	-0.0710823\\
2.5516	-0.0710718\\
2.5542	-0.0710608\\
2.5567	-0.0710497\\
2.5593	-0.0710376\\
2.5618	-0.0710254\\
2.5644	-0.0710122\\
2.567	-0.0709984\\
2.5695	-0.0709846\\
2.5721	-0.0709698\\
2.5747	-0.0709544\\
2.5772	-0.070939\\
2.5798	-0.0709225\\
2.5824	-0.0709055\\
2.5849	-0.0708886\\
2.5875	-0.0708705\\
2.59	-0.0708526\\
2.5926	-0.0708334\\
2.5952	-0.0708137\\
2.5977	-0.0707942\\
2.6003	-0.0707735\\
2.6029	-0.0707522\\
2.6054	-0.0707312\\
2.608	-0.0707089\\
2.6106	-0.070686\\
2.6131	-0.0706636\\
2.6157	-0.0706397\\
2.6182	-0.0706162\\
2.6208	-0.0705913\\
2.6234	-0.0705659\\
2.6259	-0.0705409\\
2.6285	-0.0705145\\
2.6311	-0.0704876\\
2.6336	-0.0704612\\
2.6362	-0.0704332\\
2.6388	-0.0704048\\
2.6413	-0.070377\\
2.6439	-0.0703476\\
2.6464	-0.0703188\\
2.649	-0.0702884\\
2.6516	-0.0702576\\
2.6541	-0.0702274\\
2.6567	-0.0701956\\
2.6593	-0.0701633\\
2.6618	-0.0701317\\
2.6644	-0.0700985\\
2.667	-0.0700648\\
2.6695	-0.0700319\\
2.6721	-0.0699972\\
2.6746	-0.0699635\\
2.6772	-0.0699279\\
2.6798	-0.0698919\\
2.6823	-0.0698568\\
2.6849	-0.0698199\\
2.6875	-0.0697826\\
2.69	-0.0697462\\
2.6926	-0.069708\\
2.6952	-0.0696693\\
2.6977	-0.0696317\\
2.7003	-0.0695922\\
2.7029	-0.0695522\\
2.7054	-0.0695134\\
2.708	-0.0694726\\
2.7105	-0.0694329\\
2.7131	-0.0693913\\
2.7157	-0.0693493\\
2.7182	-0.0693084\\
2.7208	-0.0692656\\
2.7234	-0.0692223\\
2.7259	-0.0691804\\
2.7285	-0.0691363\\
2.7311	-0.0690919\\
2.7336	-0.0690488\\
2.7362	-0.0690036\\
2.7387	-0.0689598\\
2.7413	-0.0689139\\
2.7439	-0.0688675\\
2.7464	-0.0688226\\
2.749	-0.0687756\\
2.7516	-0.0687282\\
2.7541	-0.0686822\\
2.7567	-0.0686341\\
2.7593	-0.0685856\\
2.7618	-0.0685387\\
2.7644	-0.0684896\\
2.7669	-0.068442\\
2.7695	-0.0683921\\
2.7721	-0.068342\\
2.7746	-0.0682934\\
2.7772	-0.0682426\\
2.7798	-0.0681915\\
2.7823	-0.068142\\
2.7849	-0.0680903\\
2.7875	-0.0680382\\
2.79	-0.0679879\\
2.7926	-0.0679352\\
2.7951	-0.0678843\\
2.7977	-0.0678311\\
2.8003	-0.0677776\\
2.8028	-0.0677258\\
2.8054	-0.0676717\\
2.808	-0.0676174\\
2.8105	-0.0675648\\
2.8131	-0.0675099\\
2.8157	-0.0674548\\
2.8182	-0.0674015\\
2.8208	-0.0673458\\
2.8233	-0.0672921\\
2.8259	-0.0672359\\
2.8285	-0.0671795\\
2.831	-0.067125\\
2.8336	-0.0670682\\
2.8362	-0.0670111\\
2.8387	-0.066956\\
2.8413	-0.0668984\\
2.8439	-0.0668407\\
2.8464	-0.0667849\\
2.849	-0.0667268\\
2.8516	-0.0666684\\
2.8541	-0.0666121\\
2.8567	-0.0665533\\
2.8592	-0.0664966\\
2.8618	-0.0664375\\
2.8644	-0.0663781\\
2.8669	-0.0663209\\
2.8695	-0.0662613\\
2.8721	-0.0662014\\
2.8746	-0.0661437\\
2.8772	-0.0660836\\
2.8798	-0.0660233\\
2.8823	-0.0659651\\
2.8849	-0.0659045\\
2.8874	-0.0658461\\
2.89	-0.0657852\\
2.8926	-0.0657242\\
2.8951	-0.0656654\\
2.8977	-0.0656041\\
2.9003	-0.0655428\\
2.9028	-0.0654836\\
2.9054	-0.065422\\
2.908	-0.0653603\\
2.9105	-0.0653008\\
2.9131	-0.0652389\\
2.9156	-0.0651793\\
2.9182	-0.0651172\\
2.9208	-0.065055\\
2.9233	-0.0649951\\
2.9259	-0.0649328\\
2.9285	-0.0648704\\
2.931	-0.0648103\\
2.9336	-0.0647477\\
2.9362	-0.0646851\\
2.9387	-0.0646249\\
2.9413	-0.0645622\\
2.9438	-0.0645018\\
2.9464	-0.064439\\
2.949	-0.0643762\\
2.9515	-0.0643158\\
2.9541	-0.0642529\\
2.9567	-0.06419\\
2.9592	-0.0641294\\
2.9618	-0.0640665\\
2.9644	-0.0640035\\
2.9669	-0.063943\\
2.9695	-0.06388\\
2.972	-0.0638194\\
2.9746	-0.0637565\\
2.9772	-0.0636935\\
2.9797	-0.063633\\
2.9823	-0.0635701\\
2.9849	-0.0635071\\
2.9874	-0.0634467\\
2.99	-0.0633838\\
2.9926	-0.063321\\
2.9951	-0.0632606\\
2.9977	-0.0631979\\
3.0003	-0.0631352\\
3.0028	-0.063075\\
3.0054	-0.0630124\\
3.0079	-0.0629523\\
3.0105	-0.0628898\\
3.0131	-0.0628274\\
3.0156	-0.0627674\\
3.0182	-0.0627052\\
3.0208	-0.062643\\
3.0233	-0.0625833\\
3.0259	-0.0625213\\
3.0285	-0.0624593\\
3.031	-0.0623999\\
3.0336	-0.0623382\\
3.0361	-0.0622789\\
3.0387	-0.0622174\\
3.0413	-0.0621559\\
3.0438	-0.062097\\
3.0464	-0.0620358\\
3.049	-0.0619747\\
3.0515	-0.0619161\\
3.0541	-0.0618552\\
3.0567	-0.0617945\\
3.0592	-0.0617363\\
3.0618	-0.0616759\\
3.0643	-0.0616179\\
3.0669	-0.0615577\\
3.0695	-0.0614977\\
3.072	-0.0614402\\
3.0746	-0.0613805\\
3.0772	-0.061321\\
3.0797	-0.0612639\\
3.0823	-0.0612046\\
3.0849	-0.0611456\\
3.0874	-0.061089\\
3.09	-0.0610303\\
3.0925	-0.060974\\
3.0951	-0.0609157\\
3.0977	-0.0608575\\
3.1002	-0.0608018\\
3.1028	-0.060744\\
3.1054	-0.0606864\\
3.1079	-0.0606312\\
3.1105	-0.060574\\
3.1131	-0.060517\\
3.1156	-0.0604624\\
3.1182	-0.0604058\\
3.1207	-0.0603516\\
3.1233	-0.0602955\\
3.1259	-0.0602395\\
3.1284	-0.060186\\
3.131	-0.0601305\\
3.1336	-0.0600752\\
3.1361	-0.0600223\\
3.1387	-0.0599675\\
3.1413	-0.0599129\\
3.1438	-0.0598607\\
3.1464	-0.0598066\\
3.1489	-0.0597548\\
3.1515	-0.0597011\\
3.1541	-0.0596478\\
3.1566	-0.0595967\\
3.1592	-0.0595438\\
3.1618	-0.0594912\\
3.1643	-0.0594409\\
3.1669	-0.0593888\\
3.1695	-0.059337\\
3.172	-0.0592874\\
3.1746	-0.0592361\\
3.1772	-0.0591851\\
3.1797	-0.0591363\\
3.1823	-0.0590858\\
3.1848	-0.0590375\\
3.1874	-0.0589876\\
3.19	-0.0589379\\
3.1925	-0.0588904\\
3.1951	-0.0588413\\
3.1977	-0.0587926\\
3.2002	-0.0587459\\
3.2028	-0.0586977\\
3.2054	-0.0586498\\
3.2079	-0.058604\\
3.2105	-0.0585566\\
3.213	-0.0585114\\
3.2156	-0.0584647\\
3.2182	-0.0584182\\
3.2207	-0.0583739\\
3.2233	-0.058328\\
3.2259	-0.0582825\\
3.2284	-0.058239\\
3.231	-0.0581941\\
3.2336	-0.0581496\\
3.2361	-0.058107\\
3.2387	-0.058063\\
3.2412	-0.0580211\\
3.2438	-0.0579777\\
3.2464	-0.0579347\\
3.2489	-0.0578937\\
3.2515	-0.0578513\\
3.2541	-0.0578093\\
3.2566	-0.0577692\\
3.2592	-0.0577278\\
3.2618	-0.0576868\\
3.2643	-0.0576476\\
3.2669	-0.0576072\\
3.2694	-0.0575687\\
3.272	-0.057529\\
3.2746	-0.0574896\\
3.2771	-0.0574521\\
3.2797	-0.0574133\\
3.2823	-0.057375\\
3.2848	-0.0573384\\
3.2874	-0.0573007\\
3.29	-0.0572633\\
3.2925	-0.0572277\\
3.2951	-0.0571911\\
3.2976	-0.0571561\\
3.3002	-0.0571202\\
3.3028	-0.0570845\\
3.3053	-0.0570506\\
3.3079	-0.0570156\\
3.3105	-0.056981\\
3.313	-0.0569481\\
3.3156	-0.0569142\\
3.3182	-0.0568807\\
3.3207	-0.0568488\\
3.3233	-0.0568159\\
3.3259	-0.0567834\\
3.3284	-0.0567525\\
3.331	-0.0567207\\
3.3335	-0.0566905\\
3.3361	-0.0566594\\
3.3387	-0.0566287\\
3.3412	-0.0565995\\
3.3438	-0.0565695\\
3.3464	-0.0565398\\
3.3489	-0.0565117\\
3.3515	-0.0564827\\
3.3541	-0.0564541\\
3.3566	-0.056427\\
3.3592	-0.0563991\\
3.3617	-0.0563727\\
3.3643	-0.0563455\\
3.3669	-0.0563187\\
3.3694	-0.0562933\\
3.372	-0.0562672\\
3.3746	-0.0562415\\
3.3771	-0.0562171\\
3.3797	-0.0561921\\
3.3823	-0.0561674\\
3.3848	-0.056144\\
3.3874	-0.0561201\\
3.3899	-0.0560974\\
3.3925	-0.0560742\\
3.3951	-0.0560513\\
3.3976	-0.0560297\\
3.4002	-0.0560075\\
3.4028	-0.0559857\\
3.4053	-0.0559651\\
3.4079	-0.0559441\\
3.4105	-0.0559233\\
3.413	-0.0559037\\
3.4156	-0.0558837\\
3.4181	-0.0558648\\
3.4207	-0.0558455\\
3.4233	-0.0558266\\
3.4258	-0.0558087\\
3.4284	-0.0557904\\
3.431	-0.0557725\\
3.4335	-0.0557557\\
3.4361	-0.0557385\\
3.4387	-0.0557217\\
3.4412	-0.0557058\\
3.4438	-0.0556897\\
3.4463	-0.0556745\\
3.4489	-0.055659\\
3.4515	-0.055644\\
3.454	-0.0556298\\
3.4566	-0.0556154\\
3.4592	-0.0556013\\
3.4617	-0.0555881\\
3.4643	-0.0555748\\
3.4669	-0.0555618\\
3.4694	-0.0555496\\
3.472	-0.0555372\\
3.4745	-0.0555257\\
3.4771	-0.055514\\
3.4797	-0.0555027\\
3.4822	-0.0554922\\
3.4848	-0.0554815\\
3.4874	-0.0554712\\
3.4899	-0.0554616\\
3.4925	-0.055452\\
3.4951	-0.0554427\\
3.4976	-0.055434\\
3.5002	-0.0554254\\
3.5028	-0.0554171\\
3.5053	-0.0554094\\
3.5079	-0.0554017\\
3.5104	-0.0553947\\
3.513	-0.0553877\\
3.5156	-0.055381\\
3.5181	-0.0553748\\
3.5207	-0.0553688\\
3.5233	-0.0553631\\
3.5258	-0.0553579\\
3.5284	-0.0553528\\
3.531	-0.055348\\
3.5335	-0.0553437\\
3.5361	-0.0553396\\
3.5386	-0.0553359\\
3.5412	-0.0553323\\
3.5438	-0.0553291\\
3.5463	-0.0553263\\
3.5489	-0.0553237\\
3.5515	-0.0553214\\
3.554	-0.0553195\\
3.5566	-0.0553177\\
3.5592	-0.0553163\\
3.5617	-0.0553153\\
3.5643	-0.0553145\\
3.5668	-0.055314\\
3.5694	-0.0553137\\
3.572	-0.0553138\\
3.5745	-0.0553142\\
3.5771	-0.0553148\\
3.5797	-0.0553157\\
3.5822	-0.0553169\\
3.5848	-0.0553184\\
3.5874	-0.0553202\\
3.5899	-0.0553221\\
3.5925	-0.0553245\\
3.595	-0.055327\\
3.5976	-0.0553299\\
3.6002	-0.055333\\
3.6027	-0.0553363\\
3.6053	-0.05534\\
3.6079	-0.055344\\
3.6104	-0.055348\\
3.613	-0.0553525\\
3.6156	-0.0553573\\
3.6181	-0.0553621\\
3.6207	-0.0553673\\
3.6232	-0.0553726\\
3.6258	-0.0553784\\
3.6284	-0.0553844\\
3.6309	-0.0553905\\
3.6335	-0.055397\\
3.6361	-0.0554038\\
3.6386	-0.0554105\\
3.6412	-0.0554177\\
3.6438	-0.0554252\\
3.6463	-0.0554327\\
3.6489	-0.0554406\\
3.6515	-0.0554488\\
3.654	-0.0554569\\
3.6566	-0.0554656\\
3.6591	-0.0554741\\
3.6617	-0.0554832\\
3.6643	-0.0554925\\
3.6668	-0.0555017\\
3.6694	-0.0555115\\
3.672	-0.0555215\\
3.6745	-0.0555313\\
3.6771	-0.0555417\\
3.6797	-0.0555523\\
3.6822	-0.0555627\\
3.6848	-0.0555737\\
3.6873	-0.0555845\\
3.6899	-0.055596\\
3.6925	-0.0556076\\
3.695	-0.055619\\
3.6976	-0.055631\\
3.7002	-0.0556432\\
3.7027	-0.0556551\\
3.7053	-0.0556677\\
3.7079	-0.0556805\\
3.7104	-0.055693\\
3.713	-0.0557061\\
3.7155	-0.0557189\\
3.7181	-0.0557324\\
3.7207	-0.0557461\\
3.7232	-0.0557594\\
3.7258	-0.0557734\\
3.7284	-0.0557876\\
3.7309	-0.0558014\\
3.7335	-0.0558159\\
3.7361	-0.0558306\\
3.7386	-0.0558448\\
3.7412	-0.0558598\\
3.7437	-0.0558743\\
3.7463	-0.0558896\\
3.7489	-0.0559051\\
3.7514	-0.05592\\
3.754	-0.0559358\\
3.7566	-0.0559516\\
3.7591	-0.055967\\
3.7617	-0.0559831\\
3.7643	-0.0559994\\
3.7668	-0.0560152\\
3.7694	-0.0560317\\
3.7719	-0.0560478\\
3.7745	-0.0560645\\
3.7771	-0.0560815\\
3.7796	-0.0560978\\
3.7822	-0.056115\\
3.7848	-0.0561322\\
3.7873	-0.056149\\
3.7899	-0.0561664\\
3.7925	-0.056184\\
3.795	-0.0562011\\
3.7976	-0.0562189\\
3.8002	-0.0562368\\
3.8027	-0.0562541\\
3.8053	-0.0562722\\
3.8078	-0.0562897\\
3.8104	-0.056308\\
3.813	-0.0563264\\
3.8155	-0.0563442\\
3.8181	-0.0563627\\
3.8207	-0.0563814\\
3.8232	-0.0563994\\
3.8258	-0.0564182\\
3.8284	-0.056437\\
3.8309	-0.0564553\\
3.8335	-0.0564743\\
3.836	-0.0564926\\
3.8386	-0.0565118\\
3.8412	-0.056531\\
3.8437	-0.0565496\\
3.8463	-0.0565689\\
3.8489	-0.0565883\\
3.8514	-0.056607\\
3.854	-0.0566265\\
3.8566	-0.0566461\\
3.8591	-0.0566649\\
3.8617	-0.0566846\\
3.8642	-0.0567035\\
3.8668	-0.0567233\\
3.8694	-0.056743\\
3.8719	-0.0567621\\
3.8745	-0.0567819\\
3.8771	-0.0568018\\
3.8796	-0.056821\\
3.8822	-0.0568409\\
3.8848	-0.0568609\\
3.8873	-0.0568801\\
3.8899	-0.0569001\\
3.8924	-0.0569194\\
3.895	-0.0569395\\
3.8976	-0.0569595\\
3.9001	-0.0569788\\
3.9027	-0.0569989\\
3.9053	-0.057019\\
3.9078	-0.0570384\\
3.9104	-0.0570585\\
3.913	-0.0570786\\
3.9155	-0.0570979\\
3.9181	-0.057118\\
3.9206	-0.0571374\\
3.9232	-0.0571575\\
3.9258	-0.0571776\\
3.9283	-0.0571969\\
3.9309	-0.0572169\\
3.9335	-0.057237\\
3.936	-0.0572563\\
3.9386	-0.0572763\\
3.9412	-0.0572963\\
3.9437	-0.0573155\\
3.9463	-0.0573355\\
3.9488	-0.0573546\\
3.9514	-0.0573745\\
3.954	-0.0573944\\
3.9565	-0.0574135\\
3.9591	-0.0574333\\
3.9617	-0.0574531\\
3.9642	-0.057472\\
3.9668	-0.0574917\\
3.9694	-0.0575114\\
3.9719	-0.0575302\\
3.9745	-0.0575498\\
3.9771	-0.0575693\\
3.9796	-0.057588\\
3.9822	-0.0576074\\
3.9847	-0.057626\\
3.9873	-0.0576453\\
3.9899	-0.0576646\\
3.9924	-0.057683\\
3.995	-0.0577022\\
3.9976	-0.0577212\\
4.0001	-0.0577395\\
4.0027	-0.0577585\\
4.0053	-0.0577773\\
4.0078	-0.0577954\\
4.0104	-0.0578142\\
4.0129	-0.0578321\\
4.0155	-0.0578507\\
4.0181	-0.0578692\\
4.0206	-0.0578869\\
4.0232	-0.0579053\\
4.0258	-0.0579236\\
4.0283	-0.0579411\\
4.0309	-0.0579592\\
4.0335	-0.0579772\\
4.036	-0.0579945\\
4.0386	-0.0580123\\
4.0411	-0.0580294\\
4.0437	-0.0580471\\
4.0463	-0.0580647\\
4.0488	-0.0580815\\
4.0514	-0.0580989\\
4.054	-0.0581162\\
4.0565	-0.0581327\\
4.0591	-0.0581498\\
4.0617	-0.0581668\\
4.0642	-0.058183\\
4.0668	-0.0581998\\
4.0693	-0.0582159\\
4.0719	-0.0582325\\
4.0745	-0.0582489\\
4.077	-0.0582647\\
4.0796	-0.0582809\\
4.0822	-0.0582971\\
4.0847	-0.0583125\\
4.0873	-0.0583284\\
4.0899	-0.0583442\\
4.0924	-0.0583593\\
4.095	-0.0583748\\
4.0975	-0.0583897\\
4.1001	-0.058405\\
4.1027	-0.0584202\\
4.1052	-0.0584347\\
4.1078	-0.0584496\\
4.1104	-0.0584645\\
4.1129	-0.0584786\\
4.1155	-0.0584932\\
4.1181	-0.0585076\\
4.1206	-0.0585214\\
4.1232	-0.0585356\\
4.1258	-0.0585496\\
4.1283	-0.058563\\
4.1309	-0.0585768\\
4.1334	-0.0585899\\
4.136	-0.0586035\\
4.1386	-0.0586168\\
4.1411	-0.0586296\\
4.1437	-0.0586427\\
4.1463	-0.0586557\\
4.1488	-0.058668\\
4.1514	-0.0586807\\
4.154	-0.0586932\\
4.1565	-0.0587052\\
4.1591	-0.0587174\\
4.1616	-0.0587291\\
4.1642	-0.058741\\
4.1668	-0.0587529\\
4.1693	-0.0587641\\
4.1719	-0.0587756\\
4.1745	-0.058787\\
4.177	-0.0587978\\
4.1796	-0.0588089\\
4.1822	-0.0588198\\
4.1847	-0.0588302\\
4.1873	-0.0588408\\
4.1898	-0.0588509\\
4.1924	-0.0588612\\
4.195	-0.0588714\\
4.1975	-0.058881\\
4.2001	-0.0588909\\
4.2027	-0.0589006\\
4.2052	-0.0589098\\
4.2078	-0.0589192\\
4.2104	-0.0589284\\
4.2129	-0.0589372\\
4.2155	-0.0589461\\
4.218	-0.0589545\\
4.2206	-0.0589631\\
4.2232	-0.0589716\\
4.2257	-0.0589796\\
4.2283	-0.0589877\\
4.2309	-0.0589957\\
4.2334	-0.0590032\\
4.236	-0.0590108\\
4.2386	-0.0590183\\
4.2411	-0.0590253\\
4.2437	-0.0590325\\
4.2462	-0.0590392\\
4.2488	-0.0590461\\
4.2514	-0.0590527\\
4.2539	-0.059059\\
4.2565	-0.0590653\\
4.2591	-0.0590715\\
4.2616	-0.0590773\\
4.2642	-0.0590832\\
4.2668	-0.0590889\\
4.2693	-0.0590942\\
4.2719	-0.0590995\\
4.2744	-0.0591045\\
4.277	-0.0591095\\
4.2796	-0.0591144\\
4.2821	-0.0591189\\
4.2847	-0.0591234\\
4.2873	-0.0591278\\
4.2898	-0.0591318\\
4.2924	-0.0591359\\
4.295	-0.0591397\\
4.2975	-0.0591433\\
4.3001	-0.0591468\\
4.3027	-0.0591502\\
4.3052	-0.0591533\\
4.3078	-0.0591563\\
4.3103	-0.0591591\\
4.3129	-0.0591618\\
4.3155	-0.0591643\\
4.318	-0.0591666\\
4.3206	-0.0591688\\
4.3232	-0.0591708\\
4.3257	-0.0591726\\
4.3283	-0.0591743\\
4.3309	-0.0591759\\
4.3334	-0.0591772\\
4.336	-0.0591784\\
4.3385	-0.0591794\\
4.3411	-0.0591803\\
4.3437	-0.059181\\
4.3462	-0.0591816\\
4.3488	-0.059182\\
4.3514	-0.0591822\\
4.3539	-0.0591822\\
4.3565	-0.0591821\\
4.3591	-0.0591819\\
4.3616	-0.0591815\\
4.3642	-0.0591809\\
4.3667	-0.0591801\\
4.3693	-0.0591792\\
4.3719	-0.0591781\\
4.3744	-0.059177\\
4.377	-0.0591756\\
4.3796	-0.059174\\
4.3821	-0.0591723\\
4.3847	-0.0591705\\
4.3873	-0.0591684\\
4.3898	-0.0591663\\
4.3924	-0.0591639\\
4.3949	-0.0591615\\
4.3975	-0.0591589\\
4.4001	-0.059156\\
4.4026	-0.0591531\\
4.4052	-0.05915\\
4.4078	-0.0591467\\
4.4103	-0.0591434\\
4.4129	-0.0591398\\
4.4155	-0.059136\\
4.418	-0.0591322\\
4.4206	-0.0591282\\
4.4231	-0.0591241\\
4.4257	-0.0591197\\
4.4283	-0.0591152\\
4.4308	-0.0591107\\
4.4334	-0.0591058\\
4.436	-0.0591009\\
4.4385	-0.0590959\\
4.4411	-0.0590906\\
4.4437	-0.0590852\\
4.4462	-0.0590799\\
4.4488	-0.0590741\\
4.4514	-0.0590683\\
4.4539	-0.0590625\\
4.4565	-0.0590563\\
4.459	-0.0590502\\
4.4616	-0.0590438\\
4.4642	-0.0590372\\
4.4667	-0.0590307\\
4.4693	-0.0590238\\
4.4719	-0.0590168\\
4.4744	-0.0590099\\
4.477	-0.0590027\\
4.4796	-0.0589952\\
4.4821	-0.0589879\\
4.4847	-0.0589802\\
4.4872	-0.0589727\\
4.4898	-0.0589647\\
4.4924	-0.0589566\\
4.4949	-0.0589486\\
4.4975	-0.0589403\\
4.5001	-0.0589317\\
4.5026	-0.0589234\\
4.5052	-0.0589146\\
4.5078	-0.0589057\\
4.5103	-0.058897\\
4.5129	-0.0588879\\
4.5154	-0.0588789\\
4.518	-0.0588695\\
4.5206	-0.05886\\
4.5231	-0.0588507\\
4.5257	-0.0588409\\
4.5283	-0.058831\\
4.5308	-0.0588213\\
4.5334	-0.0588112\\
4.536	-0.0588009\\
4.5385	-0.0587909\\
4.5411	-0.0587804\\
4.5436	-0.0587701\\
4.5462	-0.0587594\\
4.5488	-0.0587485\\
4.5513	-0.0587379\\
4.5539	-0.0587269\\
4.5565	-0.0587156\\
4.559	-0.0587048\\
4.5616	-0.0586933\\
4.5642	-0.0586818\\
4.5667	-0.0586706\\
4.5693	-0.0586588\\
4.5718	-0.0586474\\
4.5744	-0.0586354\\
4.577	-0.0586234\\
4.5795	-0.0586117\\
4.5821	-0.0585994\\
4.5847	-0.058587\\
4.5872	-0.058575\\
4.5898	-0.0585624\\
4.5924	-0.0585497\\
4.5949	-0.0585374\\
4.5975	-0.0585246\\
4.6001	-0.0585116\\
4.6026	-0.058499\\
4.6052	-0.0584859\\
4.6077	-0.0584731\\
4.6103	-0.0584598\\
4.6129	-0.0584463\\
4.6154	-0.0584333\\
4.618	-0.0584197\\
4.6206	-0.058406\\
4.6231	-0.0583928\\
4.6257	-0.0583789\\
4.6283	-0.0583649\\
4.6308	-0.0583514\\
4.6334	-0.0583373\\
4.6359	-0.0583237\\
4.6385	-0.0583094\\
4.6411	-0.058295\\
4.6436	-0.0582812\\
4.6462	-0.0582667\\
4.6488	-0.0582521\\
4.6513	-0.058238\\
4.6539	-0.0582232\\
4.6565	-0.0582084\\
4.659	-0.0581941\\
4.6616	-0.0581792\\
4.6641	-0.0581648\\
4.6667	-0.0581497\\
4.6693	-0.0581346\\
4.6718	-0.0581199\\
4.6744	-0.0581047\\
4.677	-0.0580893\\
4.6795	-0.0580745\\
4.6821	-0.0580591\\
4.6847	-0.0580436\\
4.6872	-0.0580286\\
4.6898	-0.058013\\
4.6923	-0.0579979\\
4.6949	-0.0579822\\
4.6975	-0.0579664\\
4.7	-0.0579511\\
4.7026	-0.0579352\\
4.7052	-0.0579193\\
4.7077	-0.0579039\\
4.7103	-0.0578879\\
4.7129	-0.0578718\\
4.7154	-0.0578563\\
4.718	-0.0578401\\
4.7205	-0.0578245\\
4.7231	-0.0578083\\
4.7257	-0.057792\\
4.7282	-0.0577763\\
4.7308	-0.0577599\\
4.7334	-0.0577435\\
4.7359	-0.0577277\\
4.7385	-0.0577112\\
4.7411	-0.0576947\\
4.7436	-0.0576788\\
4.7462	-0.0576622\\
4.7487	-0.0576462\\
4.7513	-0.0576296\\
4.7539	-0.0576129\\
4.7564	-0.0575969\\
4.759	-0.0575801\\
4.7616	-0.0575634\\
4.7641	-0.0575473\\
4.7667	-0.0575305\\
4.7693	-0.0575137\\
4.7718	-0.0574975\\
4.7744	-0.0574807\\
4.777	-0.0574638\\
4.7795	-0.0574476\\
4.7821	-0.0574307\\
4.7846	-0.0574144\\
4.7872	-0.0573975\\
4.7898	-0.0573806\\
4.7923	-0.0573643\\
4.7949	-0.0573474\\
4.7975	-0.0573304\\
4.8	-0.0573141\\
4.8026	-0.0572971\\
4.8052	-0.0572802\\
4.8077	-0.0572638\\
4.8103	-0.0572468\\
4.8128	-0.0572305\\
4.8154	-0.0572135\\
4.818	-0.0571966\\
4.8205	-0.0571802\\
4.8231	-0.0571633\\
4.8257	-0.0571463\\
4.8282	-0.05713\\
4.8308	-0.057113\\
4.8334	-0.0570961\\
4.8359	-0.0570798\\
4.8385	-0.0570628\\
4.841	-0.0570466\\
4.8436	-0.0570297\\
4.8462	-0.0570128\\
4.8487	-0.0569965\\
4.8513	-0.0569797\\
4.8539	-0.0569628\\
4.8564	-0.0569466\\
4.859	-0.0569298\\
4.8616	-0.056913\\
4.8641	-0.0568969\\
4.8667	-0.0568801\\
4.8692	-0.056864\\
4.8718	-0.0568473\\
4.8744	-0.0568306\\
4.8769	-0.0568146\\
4.8795	-0.056798\\
4.8821	-0.0567814\\
4.8846	-0.0567654\\
4.8872	-0.0567489\\
4.8898	-0.0567324\\
4.8923	-0.0567165\\
4.8949	-0.0567001\\
4.8974	-0.0566843\\
4.9	-0.0566679\\
4.9026	-0.0566516\\
4.9051	-0.0566359\\
4.9077	-0.0566196\\
4.9103	-0.0566034\\
4.9128	-0.0565878\\
4.9154	-0.0565716\\
4.918	-0.0565555\\
4.9205	-0.0565401\\
4.9231	-0.056524\\
4.9257	-0.056508\\
4.9282	-0.0564927\\
4.9308	-0.0564768\\
4.9333	-0.0564616\\
4.9359	-0.0564458\\
4.9385	-0.05643\\
4.941	-0.0564149\\
4.9436	-0.0563993\\
4.9462	-0.0563837\\
4.9487	-0.0563687\\
4.9513	-0.0563532\\
4.9539	-0.0563378\\
4.9564	-0.056323\\
4.959	-0.0563077\\
4.9615	-0.056293\\
4.9641	-0.0562778\\
4.9667	-0.0562626\\
4.9692	-0.0562481\\
4.9718	-0.0562331\\
4.9744	-0.0562181\\
4.9769	-0.0562037\\
4.9795	-0.0561889\\
4.9821	-0.0561741\\
4.9846	-0.0561599\\
4.9872	-0.0561453\\
4.9897	-0.0561312\\
4.9923	-0.0561167\\
4.9949	-0.0561022\\
4.9974	-0.0560884\\
5	-0.056074\\
};
\addplot [color=mycolor24, line width=1.0pt, forget plot]
  table[row sep=crcr]{%
-0.125	6.70079e-08\\
-0.1224	6.8396e-08\\
-0.1199	6.97298e-08\\
-0.1173	7.11168e-08\\
-0.1147	7.25045e-08\\
-0.1122	7.38384e-08\\
-0.1096	7.52253e-08\\
-0.1071	7.65592e-08\\
-0.1045	7.79459e-08\\
-0.1019	7.93319e-08\\
-0.0994	8.06652e-08\\
-0.0968	8.20523e-08\\
-0.0942	8.34389e-08\\
-0.0917	8.47722e-08\\
-0.0891	8.61588e-08\\
-0.0865	8.75446e-08\\
-0.084	8.88772e-08\\
-0.0814	9.02638e-08\\
-0.0789	9.15969e-08\\
-0.0763	9.2983e-08\\
-0.0737	9.43693e-08\\
-0.0712	9.57019e-08\\
-0.0686	9.70871e-08\\
-0.066	9.8473e-08\\
-0.0635	9.98061e-08\\
-0.0609	1.01192e-07\\
-0.0583	1.02578e-07\\
-0.0558	1.0391e-07\\
-0.0532	1.05295e-07\\
-0.0507	1.06627e-07\\
-0.0481	1.08013e-07\\
-0.0455	1.09398e-07\\
-0.043	1.1073e-07\\
-0.0404	1.12115e-07\\
-0.0378	1.135e-07\\
-0.0353	1.14831e-07\\
-0.0327	1.16216e-07\\
-0.0301	1.17602e-07\\
-0.0276	1.18933e-07\\
-0.025	1.20317e-07\\
-0.0224	1.21702e-07\\
-0.0199	1.23032e-07\\
-0.0173	1.24416e-07\\
-0.0148	1.25748e-07\\
-0.0122	1.27132e-07\\
-0.0096	1.28515e-07\\
-0.0071	1.29846e-07\\
-0.0045	1.31229e-07\\
-0.0019	1.32612e-07\\
0.0006	1.33942e-07\\
0.0032	-0.000378486\\
0.0058	-0.00439638\\
0.0083	-0.00682812\\
0.0109	-0.001965\\
0.0134	0.00775973\\
0.016	0.0174953\\
0.0186	0.0231909\\
0.0211	0.0250701\\
0.0237	0.0255934\\
0.0263	0.0265004\\
0.0288	0.0279579\\
0.0314	0.0293499\\
0.034	0.0299911\\
0.0365	0.0298266\\
0.0391	0.0290043\\
0.0416	0.0275948\\
0.0442	0.0253488\\
0.0468	0.0223304\\
0.0493	0.0190187\\
0.0519	0.0155611\\
0.0545	0.0121916\\
0.057	0.00871143\\
0.0596	0.00436641\\
0.0622	-0.000868429\\
0.0647	-0.00638686\\
0.0673	-0.0120152\\
0.0698	-0.0169779\\
0.0724	-0.0218473\\
0.075	-0.0269572\\
0.0775	-0.0324708\\
0.0801	-0.0387096\\
0.0827	-0.0449029\\
0.0852	-0.0503051\\
0.0878	-0.0552625\\
0.0904	-0.0599494\\
0.0929	-0.0646817\\
0.0955	-0.0699901\\
0.098	-0.0751464\\
0.1006	-0.0800146\\
0.1032	-0.0840666\\
0.1057	-0.0873575\\
0.1083	-0.0906282\\
0.1109	-0.0941049\\
0.1134	-0.0975701\\
0.116	-0.100863\\
0.1186	-0.103407\\
0.1211	-0.105096\\
0.1237	-0.106416\\
0.1263	-0.10774\\
0.1288	-0.109165\\
0.1314	-0.110556\\
0.1339	-0.111396\\
0.1365	-0.111517\\
0.1391	-0.111035\\
0.1416	-0.110402\\
0.1442	-0.109887\\
0.1468	-0.109477\\
0.1493	-0.108855\\
0.1519	-0.107638\\
0.1545	-0.105808\\
0.157	-0.103763\\
0.1596	-0.101719\\
0.1621	-0.0999754\\
0.1647	-0.0982023\\
0.1673	-0.0961178\\
0.1698	-0.0936381\\
0.1724	-0.0907041\\
0.175	-0.0877678\\
0.1775	-0.0852053\\
0.1801	-0.0827876\\
0.1827	-0.0803362\\
0.1852	-0.0776806\\
0.1878	-0.0745831\\
0.1903	-0.0715216\\
0.1929	-0.068572\\
0.1955	-0.0659885\\
0.198	-0.0636956\\
0.2006	-0.0612168\\
0.2032	-0.0584776\\
0.2057	-0.0557058\\
0.2083	-0.0529716\\
0.2109	-0.0506164\\
0.2134	-0.0486842\\
0.216	-0.0467779\\
0.2185	-0.0448152\\
0.2211	-0.0426081\\
0.2237	-0.0404407\\
0.2262	-0.0386463\\
0.2288	-0.0371785\\
0.2314	-0.0359688\\
0.2339	-0.0348034\\
0.2365	-0.0334456\\
0.2391	-0.0320296\\
0.2416	-0.0308425\\
0.2442	-0.0299693\\
0.2467	-0.0294483\\
0.2493	-0.0290263\\
0.2519	-0.0285118\\
0.2544	-0.0278986\\
0.257	-0.0273042\\
0.2596	-0.0269724\\
0.2621	-0.0269729\\
0.2647	-0.0271764\\
0.2673	-0.0273614\\
0.2698	-0.0274006\\
0.2724	-0.0273679\\
0.2749	-0.027453\\
0.2775	-0.0278113\\
0.2801	-0.0284145\\
0.2826	-0.0290505\\
0.2852	-0.0295897\\
0.2878	-0.0299756\\
0.2903	-0.0303352\\
0.2929	-0.030876\\
0.2955	-0.031648\\
0.298	-0.0325023\\
0.3006	-0.0333201\\
0.3032	-0.0339588\\
0.3057	-0.0344611\\
0.3083	-0.0350334\\
0.3108	-0.0357505\\
0.3134	-0.0366453\\
0.316	-0.0375366\\
0.3185	-0.0382456\\
0.3211	-0.0388068\\
0.3237	-0.0393097\\
0.3262	-0.0398844\\
0.3288	-0.0406299\\
0.3314	-0.0414328\\
0.3339	-0.04211\\
0.3365	-0.0426294\\
0.339	-0.0430007\\
0.3416	-0.0433957\\
0.3442	-0.0439094\\
0.3467	-0.0444986\\
0.3493	-0.0450847\\
0.3519	-0.045518\\
0.3544	-0.0457746\\
0.357	-0.0459775\\
0.3596	-0.0462479\\
0.3621	-0.0466169\\
0.3647	-0.0470398\\
0.3672	-0.0473624\\
0.3698	-0.0475388\\
0.3724	-0.0475995\\
0.3749	-0.0476651\\
0.3775	-0.0478337\\
0.3801	-0.0480909\\
0.3826	-0.048322\\
0.3852	-0.048443\\
0.3878	-0.0484302\\
0.3903	-0.0483754\\
0.3929	-0.0483853\\
0.3954	-0.0484979\\
0.398	-0.048662\\
0.4006	-0.0487649\\
0.4031	-0.0487485\\
0.4057	-0.0486501\\
0.4083	-0.0485733\\
0.4108	-0.0485952\\
0.4134	-0.0487066\\
0.416	-0.048819\\
0.4185	-0.048845\\
0.4211	-0.0487762\\
0.4236	-0.0486872\\
0.4262	-0.0486644\\
0.4288	-0.0487476\\
0.4313	-0.0488784\\
0.4339	-0.0489766\\
0.4365	-0.0489845\\
0.439	-0.0489374\\
0.4416	-0.0489171\\
0.4442	-0.0489924\\
0.4467	-0.0491454\\
0.4493	-0.0493137\\
0.4519	-0.0494165\\
0.4544	-0.0494442\\
0.457	-0.0494599\\
0.4595	-0.0495362\\
0.4621	-0.0497076\\
0.4647	-0.0499275\\
0.4672	-0.0501103\\
0.4698	-0.0502271\\
0.4724	-0.0502957\\
0.4749	-0.0503848\\
0.4775	-0.0505558\\
0.4801	-0.0507968\\
0.4826	-0.0510337\\
0.4852	-0.0512225\\
0.4877	-0.05134\\
0.4903	-0.051446\\
0.4929	-0.0516016\\
0.4954	-0.0518205\\
0.498	-0.0520824\\
0.5006	-0.0523123\\
0.5031	-0.0524673\\
0.5057	-0.0525815\\
0.5083	-0.0527088\\
0.5108	-0.0528873\\
0.5134	-0.0531233\\
0.5159	-0.0533475\\
0.5185	-0.0535244\\
0.5211	-0.053636\\
0.5236	-0.0537221\\
0.5262	-0.0538446\\
0.5288	-0.0540206\\
0.5313	-0.0542097\\
0.5339	-0.054372\\
0.5365	-0.0544673\\
0.539	-0.0545134\\
0.5416	-0.0545648\\
0.5441	-0.0546544\\
0.5467	-0.0547834\\
0.5493	-0.0549033\\
0.5518	-0.0549668\\
0.5544	-0.0549725\\
0.557	-0.0549536\\
0.5595	-0.0549569\\
0.5621	-0.0549998\\
0.5647	-0.0550557\\
0.5672	-0.0550774\\
0.5698	-0.0550405\\
0.5723	-0.0549629\\
0.5749	-0.0548813\\
0.5775	-0.0548328\\
0.58	-0.0548129\\
0.5826	-0.0547827\\
0.5852	-0.0547046\\
0.5877	-0.0545794\\
0.5903	-0.0544279\\
0.5929	-0.0542948\\
0.5954	-0.0541987\\
0.598	-0.0541096\\
0.6006	-0.0539923\\
0.6031	-0.0538319\\
0.6057	-0.0536294\\
0.6082	-0.0534351\\
0.6108	-0.0532619\\
0.6134	-0.0531143\\
0.6159	-0.0529664\\
0.6185	-0.052775\\
0.6211	-0.052541\\
0.6236	-0.052301\\
0.6262	-0.0520692\\
0.6288	-0.0518684\\
0.6313	-0.0516869\\
0.6339	-0.0514775\\
0.6364	-0.0512388\\
0.639	-0.050963\\
0.6416	-0.0506906\\
0.6441	-0.0504549\\
0.6467	-0.0502341\\
0.6493	-0.050011\\
0.6518	-0.0497688\\
0.6544	-0.0494846\\
0.657	-0.0491899\\
0.6595	-0.0489238\\
0.6621	-0.0486756\\
0.6646	-0.0484496\\
0.6672	-0.0482002\\
0.6698	-0.0479201\\
0.6723	-0.0476308\\
0.6749	-0.0473351\\
0.6775	-0.0470652\\
0.68	-0.0468274\\
0.6826	-0.0465804\\
0.6852	-0.0463107\\
0.6877	-0.0460269\\
0.6903	-0.0457248\\
0.6928	-0.0454502\\
0.6954	-0.0451903\\
0.698	-0.0449432\\
0.7005	-0.0446951\\
0.7031	-0.0444129\\
0.7057	-0.0441139\\
0.7082	-0.0438317\\
0.7108	-0.0435605\\
0.7134	-0.0433095\\
0.7159	-0.0430691\\
0.7185	-0.042801\\
0.721	-0.042523\\
0.7236	-0.0422283\\
0.7262	-0.0419482\\
0.7287	-0.0417\\
0.7313	-0.0414529\\
0.7339	-0.0411969\\
0.7364	-0.0409313\\
0.739	-0.0406414\\
0.7416	-0.0403564\\
0.7441	-0.0401006\\
0.7467	-0.0398516\\
0.7492	-0.0396134\\
0.7518	-0.0393507\\
0.7544	-0.0390702\\
0.7569	-0.0387961\\
0.7595	-0.0385234\\
0.7621	-0.0382698\\
0.7646	-0.0380351\\
0.7672	-0.0377839\\
0.7698	-0.0375157\\
0.7723	-0.0372472\\
0.7749	-0.0369724\\
0.7775	-0.0367142\\
0.78	-0.0364803\\
0.7826	-0.0362384\\
0.7851	-0.0359943\\
0.7877	-0.0357264\\
0.7903	-0.0354552\\
0.7928	-0.0352052\\
0.7954	-0.0349619\\
0.798	-0.0347277\\
0.8005	-0.0344976\\
0.8031	-0.0342453\\
0.8057	-0.0339846\\
0.8082	-0.0337385\\
0.8108	-0.0334976\\
0.8133	-0.0332793\\
0.8159	-0.0330549\\
0.8185	-0.0328216\\
0.821	-0.0325873\\
0.8236	-0.0323423\\
0.8262	-0.0321084\\
0.8287	-0.0318984\\
0.8313	-0.0316891\\
0.8339	-0.0314772\\
0.8364	-0.0312644\\
0.839	-0.0310378\\
0.8415	-0.0308252\\
0.8441	-0.0306185\\
0.8467	-0.0304253\\
0.8492	-0.0302434\\
0.8518	-0.0300487\\
0.8544	-0.0298467\\
0.8569	-0.0296533\\
0.8595	-0.0294629\\
0.8621	-0.0292877\\
0.8646	-0.0291283\\
0.8672	-0.028962\\
0.8697	-0.0287963\\
0.8723	-0.028621\\
0.8749	-0.0284522\\
0.8774	-0.0283031\\
0.88	-0.0281609\\
0.8826	-0.0280238\\
0.8851	-0.027889\\
0.8877	-0.0277445\\
0.8903	-0.0276021\\
0.8928	-0.0274754\\
0.8954	-0.0273575\\
0.8979	-0.0272533\\
0.9005	-0.027146\\
0.9031	-0.0270349\\
0.9056	-0.026927\\
0.9082	-0.0268213\\
0.9108	-0.0267285\\
0.9133	-0.0266508\\
0.9159	-0.0265754\\
0.9185	-0.0264984\\
0.921	-0.0264219\\
0.9236	-0.026345\\
0.9262	-0.0262781\\
0.9287	-0.0262259\\
0.9313	-0.0261803\\
0.9338	-0.0261382\\
0.9364	-0.0260919\\
0.939	-0.0260454\\
0.9415	-0.0260066\\
0.9441	-0.0259775\\
0.9467	-0.0259594\\
0.9492	-0.0259467\\
0.9518	-0.0259328\\
0.9544	-0.0259172\\
0.9569	-0.0259048\\
0.9595	-0.0259005\\
0.962	-0.0259069\\
0.9646	-0.0259213\\
0.9672	-0.0259376\\
0.9697	-0.0259517\\
0.9723	-0.0259663\\
0.9749	-0.0259865\\
0.9774	-0.0260151\\
0.98	-0.0260542\\
0.9826	-0.0260978\\
0.9851	-0.0261393\\
0.9877	-0.0261812\\
0.9902	-0.0262238\\
0.9928	-0.0262754\\
0.9954	-0.0263365\\
0.9979	-0.0264016\\
1.0005	-0.0264711\\
1.0031	-0.0265392\\
1.0056	-0.0266049\\
1.0082	-0.0266778\\
1.0108	-0.0267591\\
1.0133	-0.0268448\\
1.0159	-0.0269378\\
1.0184	-0.0270271\\
1.021	-0.0271189\\
1.0236	-0.0272129\\
1.0261	-0.0273093\\
1.0287	-0.0274174\\
1.0313	-0.0275314\\
1.0338	-0.0276426\\
1.0364	-0.0277572\\
1.039	-0.027872\\
1.0415	-0.0279863\\
1.0441	-0.0281121\\
1.0466	-0.0282397\\
1.0492	-0.0283758\\
1.0518	-0.0285119\\
1.0543	-0.0286421\\
1.0569	-0.0287794\\
1.0595	-0.0289221\\
1.062	-0.0290659\\
1.0646	-0.0292204\\
1.0672	-0.0293765\\
1.0697	-0.0295259\\
1.0723	-0.0296815\\
1.0748	-0.0298345\\
1.0774	-0.0299996\\
1.08	-0.0301706\\
1.0825	-0.030338\\
1.0851	-0.0305122\\
1.0877	-0.030686\\
1.0902	-0.0308546\\
1.0928	-0.0310346\\
1.0954	-0.0312205\\
1.0979	-0.0314033\\
1.1005	-0.0315948\\
1.1031	-0.0317858\\
1.1056	-0.0319697\\
1.1082	-0.0321639\\
1.1107	-0.0323555\\
1.1133	-0.0325596\\
1.1159	-0.0327664\\
1.1184	-0.0329654\\
1.121	-0.0331719\\
1.1236	-0.0333797\\
1.1261	-0.0335831\\
1.1287	-0.0337995\\
1.1313	-0.0340195\\
1.1338	-0.034232\\
1.1364	-0.0344524\\
1.1389	-0.0346645\\
1.1415	-0.0348873\\
1.1441	-0.0351142\\
1.1466	-0.0353362\\
1.1492	-0.035569\\
1.1518	-0.0358017\\
1.1543	-0.0360249\\
1.1569	-0.0362579\\
1.1595	-0.0364938\\
1.162	-0.0367242\\
1.1646	-0.0369664\\
1.1671	-0.0371999\\
1.1697	-0.037442\\
1.1723	-0.0376838\\
1.1748	-0.0379177\\
1.1774	-0.038164\\
1.18	-0.0384133\\
1.1825	-0.0386542\\
1.1851	-0.0389042\\
1.1877	-0.0391533\\
1.1902	-0.039393\\
1.1928	-0.0396443\\
1.1953	-0.0398884\\
1.1979	-0.0401441\\
1.2005	-0.0403998\\
1.203	-0.0406446\\
1.2056	-0.0408984\\
1.2082	-0.041153\\
1.2107	-0.0413996\\
1.2133	-0.0416581\\
1.2159	-0.0419171\\
1.2184	-0.0421652\\
1.221	-0.0424219\\
1.2235	-0.0426683\\
1.2261	-0.0429256\\
1.2287	-0.0431847\\
1.2312	-0.0434346\\
1.2338	-0.043694\\
1.2364	-0.0439519\\
1.2389	-0.0441986\\
1.2415	-0.0444553\\
1.2441	-0.044713\\
1.2466	-0.0449618\\
1.2492	-0.0452204\\
1.2518	-0.0454777\\
1.2543	-0.0457233\\
1.2569	-0.0459779\\
1.2594	-0.046223\\
1.262	-0.0464787\\
1.2646	-0.0467346\\
1.2671	-0.0469796\\
1.2697	-0.0472326\\
1.2723	-0.047484\\
1.2748	-0.0477251\\
1.2774	-0.0479762\\
1.28	-0.0482276\\
1.2825	-0.0484687\\
1.2851	-0.0487177\\
1.2876	-0.0489552\\
1.2902	-0.0492008\\
1.2928	-0.0494461\\
1.2953	-0.0496821\\
1.2979	-0.0499271\\
1.3005	-0.0501708\\
1.303	-0.050403\\
1.3056	-0.0506426\\
1.3082	-0.0508811\\
1.3107	-0.0511102\\
1.3133	-0.0513482\\
1.3158	-0.051576\\
1.3184	-0.051811\\
1.321	-0.0520437\\
1.3235	-0.0522658\\
1.3261	-0.052496\\
1.3287	-0.0527258\\
1.3312	-0.0529459\\
1.3338	-0.0531732\\
1.3364	-0.0533981\\
1.3389	-0.0536124\\
1.3415	-0.0538338\\
1.344	-0.054046\\
1.3466	-0.0542659\\
1.3492	-0.0544844\\
1.3517	-0.0546924\\
1.3543	-0.0549064\\
1.3569	-0.0551185\\
1.3594	-0.0553212\\
1.362	-0.0555312\\
1.3646	-0.05574\\
1.3671	-0.0559389\\
1.3697	-0.0561434\\
1.3722	-0.0563378\\
1.3748	-0.0565383\\
1.3774	-0.0567377\\
1.3799	-0.0569283\\
1.3825	-0.0571249\\
1.3851	-0.0573191\\
1.3876	-0.0575036\\
1.3902	-0.0576933\\
1.3928	-0.0578814\\
1.3953	-0.0580612\\
1.3979	-0.0582467\\
1.4005	-0.05843\\
1.403	-0.058604\\
1.4056	-0.0587825\\
1.4081	-0.0589523\\
1.4107	-0.0591274\\
1.4133	-0.0593012\\
1.4158	-0.0594664\\
1.4184	-0.059636\\
1.421	-0.0598029\\
1.4235	-0.0599612\\
1.4261	-0.0601242\\
1.4287	-0.0602857\\
1.4312	-0.0604393\\
1.4338	-0.0605969\\
1.4363	-0.060746\\
1.4389	-0.0608986\\
1.4415	-0.0610491\\
1.444	-0.0611922\\
1.4466	-0.0613395\\
1.4492	-0.0614847\\
1.4517	-0.061622\\
1.4543	-0.0617623\\
1.4569	-0.0619002\\
1.4594	-0.062031\\
1.462	-0.0621654\\
1.4645	-0.0622929\\
1.4671	-0.0624233\\
1.4697	-0.0625511\\
1.4722	-0.0626715\\
1.4748	-0.0627947\\
1.4774	-0.0629161\\
1.4799	-0.0630312\\
1.4825	-0.0631488\\
1.4851	-0.063264\\
1.4876	-0.0633724\\
1.4902	-0.0634827\\
1.4927	-0.063587\\
1.4953	-0.0636937\\
1.4979	-0.0637985\\
1.5004	-0.0638972\\
1.503	-0.0639973\\
1.5056	-0.064095\\
1.5081	-0.0641869\\
1.5107	-0.0642807\\
1.5133	-0.0643727\\
1.5158	-0.0644592\\
1.5184	-0.0645469\\
1.5209	-0.0646289\\
1.5235	-0.0647119\\
1.5261	-0.064793\\
1.5286	-0.0648693\\
1.5312	-0.0649468\\
1.5338	-0.0650222\\
1.5363	-0.0650925\\
1.5389	-0.0651632\\
1.5415	-0.0652319\\
1.544	-0.0652963\\
1.5466	-0.0653615\\
1.5491	-0.0654224\\
1.5517	-0.0654835\\
1.5543	-0.0655424\\
1.5568	-0.065597\\
1.5594	-0.0656519\\
1.562	-0.0657051\\
1.5645	-0.0657546\\
1.5671	-0.0658042\\
1.5697	-0.0658515\\
1.5722	-0.065895\\
1.5748	-0.0659383\\
1.5774	-0.06598\\
1.5799	-0.0660185\\
1.5825	-0.0660567\\
1.585	-0.0660916\\
1.5876	-0.0661258\\
1.5902	-0.066158\\
1.5927	-0.0661874\\
1.5953	-0.0662164\\
1.5979	-0.0662438\\
1.6004	-0.0662684\\
1.603	-0.066292\\
1.6056	-0.0663137\\
1.6081	-0.0663329\\
1.6107	-0.0663513\\
1.6132	-0.0663676\\
1.6158	-0.0663829\\
1.6184	-0.0663964\\
1.6209	-0.0664077\\
1.6235	-0.0664176\\
1.6261	-0.066426\\
1.6286	-0.0664327\\
1.6312	-0.0664382\\
1.6338	-0.0664421\\
1.6363	-0.0664442\\
1.6389	-0.0664447\\
1.6414	-0.0664436\\
1.644	-0.0664412\\
1.6466	-0.0664373\\
1.6491	-0.0664323\\
1.6517	-0.0664255\\
1.6543	-0.066417\\
1.6568	-0.0664074\\
1.6594	-0.0663961\\
1.662	-0.0663834\\
1.6645	-0.06637\\
1.6671	-0.0663547\\
1.6696	-0.0663385\\
1.6722	-0.0663202\\
1.6748	-0.0663005\\
1.6773	-0.0662804\\
1.6799	-0.0662584\\
1.6825	-0.066235\\
1.685	-0.0662113\\
1.6876	-0.0661852\\
1.6902	-0.0661578\\
1.6927	-0.0661302\\
1.6953	-0.0661005\\
1.6978	-0.0660709\\
1.7004	-0.0660389\\
1.703	-0.0660056\\
1.7055	-0.0659723\\
1.7081	-0.0659365\\
1.7107	-0.0658997\\
1.7132	-0.0658634\\
1.7158	-0.0658245\\
1.7184	-0.0657845\\
1.7209	-0.0657449\\
1.7235	-0.0657026\\
1.7261	-0.0656592\\
1.7286	-0.0656167\\
1.7312	-0.0655715\\
1.7337	-0.0655271\\
1.7363	-0.0654799\\
1.7389	-0.0654315\\
1.7414	-0.0653841\\
1.744	-0.065334\\
1.7466	-0.065283\\
1.7491	-0.0652332\\
1.7517	-0.0651804\\
1.7543	-0.0651267\\
1.7568	-0.0650741\\
1.7594	-0.0650186\\
1.7619	-0.0649646\\
1.7645	-0.0649076\\
1.7671	-0.0648498\\
1.7696	-0.0647934\\
1.7722	-0.064734\\
1.7748	-0.0646737\\
1.7773	-0.0646151\\
1.7799	-0.0645535\\
1.7825	-0.0644912\\
1.785	-0.0644307\\
1.7876	-0.0643669\\
1.7901	-0.0643049\\
1.7927	-0.0642398\\
1.7953	-0.0641741\\
1.7978	-0.0641104\\
1.8004	-0.0640435\\
1.803	-0.063976\\
1.8055	-0.0639105\\
1.8081	-0.0638417\\
1.8107	-0.0637724\\
1.8132	-0.0637054\\
1.8158	-0.0636352\\
1.8183	-0.0635671\\
1.8209	-0.0634958\\
1.8235	-0.0634239\\
1.826	-0.0633544\\
1.8286	-0.0632816\\
1.8312	-0.0632085\\
1.8337	-0.0631378\\
1.8363	-0.0630638\\
1.8389	-0.0629893\\
1.8414	-0.0629172\\
1.844	-0.062842\\
1.8465	-0.0627694\\
1.8491	-0.0626935\\
1.8517	-0.0626173\\
1.8542	-0.0625436\\
1.8568	-0.0624666\\
1.8594	-0.0623893\\
1.8619	-0.0623148\\
1.8645	-0.062237\\
1.8671	-0.062159\\
1.8696	-0.0620837\\
1.8722	-0.0620051\\
1.8747	-0.0619292\\
1.8773	-0.0618502\\
1.8799	-0.061771\\
1.8824	-0.0616946\\
1.885	-0.0616151\\
1.8876	-0.0615352\\
1.8901	-0.0614583\\
1.8927	-0.0613781\\
1.8953	-0.0612979\\
1.8978	-0.0612206\\
1.9004	-0.0611401\\
1.903	-0.0610595\\
1.9055	-0.0609818\\
1.9081	-0.060901\\
1.9106	-0.0608232\\
1.9132	-0.0607422\\
1.9158	-0.0606612\\
1.9183	-0.0605833\\
1.9209	-0.0605021\\
1.9235	-0.0604209\\
1.926	-0.0603429\\
1.9286	-0.0602617\\
1.9312	-0.0601805\\
1.9337	-0.0601025\\
1.9363	-0.0600213\\
1.9388	-0.0599433\\
1.9414	-0.0598622\\
1.944	-0.0597811\\
1.9465	-0.0597032\\
1.9491	-0.0596223\\
1.9517	-0.0595415\\
1.9542	-0.0594638\\
1.9568	-0.0593831\\
1.9594	-0.0593024\\
1.9619	-0.059225\\
1.9645	-0.0591447\\
1.967	-0.0590675\\
1.9696	-0.0589874\\
1.9722	-0.0589073\\
1.9747	-0.0588305\\
1.9773	-0.0587508\\
1.9799	-0.0586713\\
1.9824	-0.058595\\
1.985	-0.0585158\\
1.9876	-0.0584368\\
1.9901	-0.058361\\
1.9927	-0.0582823\\
1.9952	-0.0582069\\
1.9978	-0.0581287\\
2.0004	-0.0580507\\
2.0029	-0.057976\\
2.0055	-0.0578984\\
2.0081	-0.0578211\\
2.0106	-0.057747\\
2.0132	-0.0576702\\
2.0158	-0.0575936\\
2.0183	-0.0575203\\
2.0209	-0.0574443\\
2.0234	-0.0573714\\
2.026	-0.0572959\\
2.0286	-0.0572208\\
2.0311	-0.0571488\\
2.0337	-0.0570742\\
2.0363	-0.057\\
2.0388	-0.0569289\\
2.0414	-0.0568552\\
2.044	-0.0567819\\
2.0465	-0.0567117\\
2.0491	-0.0566391\\
2.0517	-0.0565668\\
2.0542	-0.0564975\\
2.0568	-0.0564259\\
2.0593	-0.0563573\\
2.0619	-0.0562864\\
2.0645	-0.0562159\\
2.067	-0.0561484\\
2.0696	-0.0560786\\
2.0722	-0.0560092\\
2.0747	-0.0559427\\
2.0773	-0.055874\\
2.0799	-0.0558058\\
2.0824	-0.0557405\\
2.085	-0.055673\\
2.0875	-0.0556085\\
2.0901	-0.0555418\\
2.0927	-0.0554755\\
2.0952	-0.0554122\\
2.0978	-0.0553468\\
2.1004	-0.0552818\\
2.1029	-0.0552197\\
2.1055	-0.0551555\\
2.1081	-0.0550918\\
2.1106	-0.0550309\\
2.1132	-0.054968\\
2.1157	-0.054908\\
2.1183	-0.0548461\\
2.1209	-0.0547845\\
2.1234	-0.0547258\\
2.126	-0.0546652\\
2.1286	-0.054605\\
2.1311	-0.0545476\\
2.1337	-0.0544884\\
2.1363	-0.0544296\\
2.1388	-0.0543735\\
2.1414	-0.0543157\\
2.1439	-0.0542605\\
2.1465	-0.0542036\\
2.1491	-0.0541472\\
2.1516	-0.0540935\\
2.1542	-0.054038\\
2.1568	-0.0539831\\
2.1593	-0.0539307\\
2.1619	-0.0538767\\
2.1645	-0.0538232\\
2.167	-0.0537723\\
2.1696	-0.0537198\\
2.1721	-0.0536697\\
2.1747	-0.0536182\\
2.1773	-0.0535672\\
2.1798	-0.0535187\\
2.1824	-0.0534687\\
2.185	-0.0534192\\
2.1875	-0.0533721\\
2.1901	-0.0533236\\
2.1927	-0.0532757\\
2.1952	-0.05323\\
2.1978	-0.0531831\\
2.2004	-0.0531367\\
2.2029	-0.0530926\\
2.2055	-0.0530472\\
2.208	-0.0530041\\
2.2106	-0.0529597\\
2.2132	-0.0529159\\
2.2157	-0.0528743\\
2.2183	-0.0528315\\
2.2209	-0.0527893\\
2.2234	-0.0527491\\
2.226	-0.0527079\\
2.2286	-0.0526673\\
2.2311	-0.0526287\\
2.2337	-0.0525891\\
2.2362	-0.0525515\\
2.2388	-0.052513\\
2.2414	-0.052475\\
2.2439	-0.0524389\\
2.2465	-0.052402\\
2.2491	-0.0523656\\
2.2516	-0.0523311\\
2.2542	-0.0522958\\
2.2568	-0.052261\\
2.2593	-0.0522281\\
2.2619	-0.0521944\\
2.2644	-0.0521625\\
2.267	-0.0521298\\
2.2696	-0.0520977\\
2.2721	-0.0520674\\
2.2747	-0.0520364\\
2.2773	-0.0520059\\
2.2798	-0.0519771\\
2.2824	-0.0519477\\
2.285	-0.0519189\\
2.2875	-0.0518917\\
2.2901	-0.0518639\\
2.2926	-0.0518377\\
2.2952	-0.0518111\\
2.2978	-0.0517849\\
2.3003	-0.0517603\\
2.3029	-0.0517353\\
2.3055	-0.0517108\\
2.308	-0.0516878\\
2.3106	-0.0516644\\
2.3132	-0.0516415\\
2.3157	-0.0516201\\
2.3183	-0.0515983\\
2.3208	-0.0515779\\
2.3234	-0.0515572\\
2.326	-0.0515371\\
2.3285	-0.0515182\\
2.3311	-0.0514992\\
2.3337	-0.0514807\\
2.3362	-0.0514634\\
2.3388	-0.051446\\
2.3414	-0.0514292\\
2.3439	-0.0514135\\
2.3465	-0.0513977\\
2.349	-0.0513831\\
2.3516	-0.0513684\\
2.3542	-0.0513542\\
2.3567	-0.0513412\\
2.3593	-0.0513281\\
2.3619	-0.0513156\\
2.3644	-0.0513041\\
2.367	-0.0512927\\
2.3696	-0.0512818\\
2.3721	-0.0512719\\
2.3747	-0.0512621\\
2.3773	-0.0512529\\
2.3798	-0.0512445\\
2.3824	-0.0512363\\
2.3849	-0.051229\\
2.3875	-0.0512219\\
2.3901	-0.0512154\\
2.3926	-0.0512096\\
2.3952	-0.0512041\\
2.3978	-0.0511992\\
2.4003	-0.051195\\
2.4029	-0.0511912\\
2.4055	-0.0511878\\
2.408	-0.0511852\\
2.4106	-0.0511829\\
2.4131	-0.0511813\\
2.4157	-0.0511801\\
2.4183	-0.0511795\\
2.4208	-0.0511794\\
2.4234	-0.0511798\\
2.426	-0.0511808\\
2.4285	-0.0511823\\
2.4311	-0.0511843\\
2.4337	-0.0511869\\
2.4362	-0.0511898\\
2.4388	-0.0511935\\
2.4413	-0.0511974\\
2.4439	-0.0512021\\
2.4465	-0.0512073\\
2.449	-0.0512128\\
2.4516	-0.0512191\\
2.4542	-0.0512259\\
2.4567	-0.0512329\\
2.4593	-0.0512407\\
2.4619	-0.051249\\
2.4644	-0.0512576\\
2.467	-0.0512669\\
2.4695	-0.0512765\\
2.4721	-0.0512869\\
2.4747	-0.0512978\\
2.4772	-0.0513088\\
2.4798	-0.0513208\\
2.4824	-0.0513333\\
2.4849	-0.0513458\\
2.4875	-0.0513593\\
2.4901	-0.0513733\\
2.4926	-0.0513872\\
2.4952	-0.0514023\\
2.4977	-0.0514172\\
2.5003	-0.0514332\\
2.5029	-0.0514498\\
2.5054	-0.0514661\\
2.508	-0.0514837\\
2.5106	-0.0515017\\
2.5131	-0.0515196\\
2.5157	-0.0515386\\
2.5183	-0.0515581\\
2.5208	-0.0515774\\
2.5234	-0.0515979\\
2.526	-0.0516189\\
2.5285	-0.0516396\\
2.5311	-0.0516616\\
2.5336	-0.0516832\\
2.5362	-0.0517061\\
2.5388	-0.0517295\\
2.5413	-0.0517525\\
2.5439	-0.0517769\\
2.5465	-0.0518018\\
2.549	-0.0518262\\
2.5516	-0.051852\\
2.5542	-0.0518783\\
2.5567	-0.051904\\
2.5593	-0.0519313\\
2.5618	-0.0519579\\
2.5644	-0.051986\\
2.567	-0.0520147\\
2.5695	-0.0520426\\
2.5721	-0.0520722\\
2.5747	-0.0521022\\
2.5772	-0.0521314\\
2.5798	-0.0521623\\
2.5824	-0.0521937\\
2.5849	-0.0522243\\
2.5875	-0.0522565\\
2.59	-0.0522879\\
2.5926	-0.0523211\\
2.5952	-0.0523546\\
2.5977	-0.0523873\\
2.6003	-0.0524218\\
2.6029	-0.0524566\\
2.6054	-0.0524906\\
2.608	-0.0525263\\
2.6106	-0.0525624\\
2.6131	-0.0525976\\
2.6157	-0.0526345\\
2.6182	-0.0526705\\
2.6208	-0.0527083\\
2.6234	-0.0527465\\
2.6259	-0.0527836\\
2.6285	-0.0528227\\
2.6311	-0.0528621\\
2.6336	-0.0529004\\
2.6362	-0.0529406\\
2.6388	-0.0529812\\
2.6413	-0.0530207\\
2.6439	-0.0530621\\
2.6464	-0.0531022\\
2.649	-0.0531444\\
2.6516	-0.0531869\\
2.6541	-0.0532282\\
2.6567	-0.0532715\\
2.6593	-0.0533151\\
2.6618	-0.0533574\\
2.6644	-0.0534018\\
2.667	-0.0534466\\
2.6695	-0.0534899\\
2.6721	-0.0535354\\
2.6746	-0.0535794\\
2.6772	-0.0536255\\
2.6798	-0.053672\\
2.6823	-0.053717\\
2.6849	-0.0537642\\
2.6875	-0.0538117\\
2.69	-0.0538577\\
2.6926	-0.0539058\\
2.6952	-0.0539543\\
2.6977	-0.0540012\\
2.7003	-0.0540503\\
2.7029	-0.0540997\\
2.7054	-0.0541475\\
2.708	-0.0541975\\
2.7105	-0.0542458\\
2.7131	-0.0542964\\
2.7157	-0.0543473\\
2.7182	-0.0543965\\
2.7208	-0.054448\\
2.7234	-0.0544997\\
2.7259	-0.0545498\\
2.7285	-0.054602\\
2.7311	-0.0546546\\
2.7336	-0.0547054\\
2.7362	-0.0547585\\
2.7387	-0.0548097\\
2.7413	-0.0548633\\
2.7439	-0.0549171\\
2.7464	-0.0549691\\
2.749	-0.0550234\\
2.7516	-0.055078\\
2.7541	-0.0551306\\
2.7567	-0.0551856\\
2.7593	-0.0552408\\
2.7618	-0.0552941\\
2.7644	-0.0553498\\
2.7669	-0.0554035\\
2.7695	-0.0554595\\
2.7721	-0.0555158\\
2.7746	-0.0555701\\
2.7772	-0.0556267\\
2.7798	-0.0556835\\
2.7823	-0.0557383\\
2.7849	-0.0557955\\
2.7875	-0.0558528\\
2.79	-0.0559081\\
2.7926	-0.0559658\\
2.7951	-0.0560214\\
2.7977	-0.0560794\\
2.8003	-0.0561375\\
2.8028	-0.0561936\\
2.8054	-0.056252\\
2.808	-0.0563106\\
2.8105	-0.056367\\
2.8131	-0.0564258\\
2.8157	-0.0564847\\
2.8182	-0.0565415\\
2.8208	-0.0566006\\
2.8233	-0.0566576\\
2.8259	-0.056717\\
2.8285	-0.0567764\\
2.831	-0.0568337\\
2.8336	-0.0568933\\
2.8362	-0.056953\\
2.8387	-0.0570105\\
2.8413	-0.0570704\\
2.8439	-0.0571303\\
2.8464	-0.057188\\
2.849	-0.057248\\
2.8516	-0.0573081\\
2.8541	-0.057366\\
2.8567	-0.0574262\\
2.8592	-0.0574841\\
2.8618	-0.0575443\\
2.8644	-0.0576046\\
2.8669	-0.0576626\\
2.8695	-0.0577229\\
2.8721	-0.0577833\\
2.8746	-0.0578413\\
2.8772	-0.0579016\\
2.8798	-0.057962\\
2.8823	-0.05802\\
2.8849	-0.0580803\\
2.8874	-0.0581383\\
2.89	-0.0581986\\
2.8926	-0.0582588\\
2.8951	-0.0583167\\
2.8977	-0.0583769\\
2.9003	-0.0584371\\
2.9028	-0.0584949\\
2.9054	-0.0585549\\
2.908	-0.0586149\\
2.9105	-0.0586725\\
2.9131	-0.0587324\\
2.9156	-0.0587898\\
2.9182	-0.0588496\\
2.9208	-0.0589092\\
2.9233	-0.0589664\\
2.9259	-0.0590259\\
2.9285	-0.0590852\\
2.931	-0.0591422\\
2.9336	-0.0592013\\
2.9362	-0.0592604\\
2.9387	-0.059317\\
2.9413	-0.0593758\\
2.9438	-0.0594322\\
2.9464	-0.0594908\\
2.949	-0.0595492\\
2.9515	-0.0596052\\
2.9541	-0.0596633\\
2.9567	-0.0597213\\
2.9592	-0.0597768\\
2.9618	-0.0598345\\
2.9644	-0.059892\\
2.9669	-0.0599471\\
2.9695	-0.0600042\\
2.972	-0.060059\\
2.9746	-0.0601158\\
2.9772	-0.0601724\\
2.9797	-0.0602266\\
2.9823	-0.0602828\\
2.9849	-0.0603389\\
2.9874	-0.0603925\\
2.99	-0.0604481\\
2.9926	-0.0605035\\
2.9951	-0.0605565\\
2.9977	-0.0606114\\
3.0003	-0.0606661\\
3.0028	-0.0607185\\
3.0054	-0.0607727\\
3.0079	-0.0608247\\
3.0105	-0.0608784\\
3.0131	-0.0609319\\
3.0156	-0.0609831\\
3.0182	-0.0610361\\
3.0208	-0.0610888\\
3.0233	-0.0611393\\
3.0259	-0.0611914\\
3.0285	-0.0612434\\
3.031	-0.061293\\
3.0336	-0.0613443\\
3.0361	-0.0613934\\
3.0387	-0.0614442\\
3.0413	-0.0614947\\
3.0438	-0.0615429\\
3.0464	-0.0615928\\
3.049	-0.0616424\\
3.0515	-0.0616897\\
3.0541	-0.0617387\\
3.0567	-0.0617873\\
3.0592	-0.0618338\\
3.0618	-0.0618817\\
3.0643	-0.0619276\\
3.0669	-0.0619749\\
3.0695	-0.0620219\\
3.072	-0.0620667\\
3.0746	-0.062113\\
3.0772	-0.062159\\
3.0797	-0.0622028\\
3.0823	-0.0622481\\
3.0849	-0.062293\\
3.0874	-0.0623358\\
3.09	-0.06238\\
3.0925	-0.0624221\\
3.0951	-0.0624655\\
3.0977	-0.0625086\\
3.1002	-0.0625496\\
3.1028	-0.0625919\\
3.1054	-0.0626339\\
3.1079	-0.0626738\\
3.1105	-0.062715\\
3.1131	-0.0627557\\
3.1156	-0.0627946\\
3.1182	-0.0628346\\
3.1207	-0.0628726\\
3.1233	-0.0629118\\
3.1259	-0.0629506\\
3.1284	-0.0629875\\
3.131	-0.0630254\\
3.1336	-0.063063\\
3.1361	-0.0630987\\
3.1387	-0.0631354\\
3.1413	-0.0631717\\
3.1438	-0.0632062\\
3.1464	-0.0632417\\
3.1489	-0.0632754\\
3.1515	-0.06331\\
3.1541	-0.0633442\\
3.1566	-0.0633767\\
3.1592	-0.06341\\
3.1618	-0.0634429\\
3.1643	-0.0634741\\
3.1669	-0.0635061\\
3.1695	-0.0635377\\
3.172	-0.0635676\\
3.1746	-0.0635983\\
3.1772	-0.0636285\\
3.1797	-0.0636572\\
3.1823	-0.0636865\\
3.1848	-0.0637143\\
3.1874	-0.0637427\\
3.19	-0.0637707\\
3.1925	-0.0637972\\
3.1951	-0.0638243\\
3.1977	-0.0638509\\
3.2002	-0.0638761\\
3.2028	-0.0639018\\
3.2054	-0.063927\\
3.2079	-0.0639509\\
3.2105	-0.0639752\\
3.213	-0.0639981\\
3.2156	-0.0640215\\
3.2182	-0.0640445\\
3.2207	-0.0640661\\
3.2233	-0.0640881\\
3.2259	-0.0641096\\
3.2284	-0.0641299\\
3.231	-0.0641505\\
3.2336	-0.0641706\\
3.2361	-0.0641895\\
3.2387	-0.0642087\\
3.2412	-0.0642267\\
3.2438	-0.064245\\
3.2464	-0.0642628\\
3.2489	-0.0642794\\
3.2515	-0.0642963\\
3.2541	-0.0643127\\
3.2566	-0.064328\\
3.2592	-0.0643434\\
3.2618	-0.0643584\\
3.2643	-0.0643723\\
3.2669	-0.0643863\\
3.2694	-0.0643994\\
3.272	-0.0644125\\
3.2746	-0.0644251\\
3.2771	-0.0644367\\
3.2797	-0.0644484\\
3.2823	-0.0644596\\
3.2848	-0.06447\\
3.2874	-0.0644802\\
3.29	-0.06449\\
3.2925	-0.064499\\
3.2951	-0.0645079\\
3.2976	-0.064516\\
3.3002	-0.0645239\\
3.3028	-0.0645314\\
3.3053	-0.0645381\\
3.3079	-0.0645447\\
3.3105	-0.0645508\\
3.313	-0.0645562\\
3.3156	-0.0645614\\
3.3182	-0.0645661\\
3.3207	-0.0645702\\
3.3233	-0.064574\\
3.3259	-0.0645773\\
3.3284	-0.0645801\\
3.331	-0.0645826\\
3.3335	-0.0645845\\
3.3361	-0.064586\\
3.3387	-0.0645871\\
3.3412	-0.0645877\\
3.3438	-0.0645879\\
3.3464	-0.0645877\\
3.3489	-0.064587\\
3.3515	-0.0645859\\
3.3541	-0.0645843\\
3.3566	-0.0645824\\
3.3592	-0.06458\\
3.3617	-0.0645772\\
3.3643	-0.0645739\\
3.3669	-0.0645702\\
3.3694	-0.0645661\\
3.372	-0.0645616\\
3.3746	-0.0645565\\
3.3771	-0.0645513\\
3.3797	-0.0645454\\
3.3823	-0.0645391\\
3.3848	-0.0645327\\
3.3874	-0.0645256\\
3.3899	-0.0645183\\
3.3925	-0.0645104\\
3.3951	-0.0645021\\
3.3976	-0.0644936\\
3.4002	-0.0644845\\
3.4028	-0.0644749\\
3.4053	-0.0644653\\
3.4079	-0.064455\\
3.4105	-0.0644442\\
3.413	-0.0644335\\
3.4156	-0.0644219\\
3.4181	-0.0644105\\
3.4207	-0.0643981\\
3.4233	-0.0643854\\
3.4258	-0.0643729\\
3.4284	-0.0643594\\
3.431	-0.0643455\\
3.4335	-0.0643319\\
3.4361	-0.0643173\\
3.4387	-0.0643023\\
3.4412	-0.0642876\\
3.4438	-0.0642719\\
3.4463	-0.0642565\\
3.4489	-0.0642401\\
3.4515	-0.0642233\\
3.454	-0.0642069\\
3.4566	-0.0641894\\
3.4592	-0.0641716\\
3.4617	-0.0641541\\
3.4643	-0.0641357\\
3.4669	-0.0641168\\
3.4694	-0.0640984\\
3.472	-0.0640789\\
3.4745	-0.0640598\\
3.4771	-0.0640396\\
3.4797	-0.0640191\\
3.4822	-0.0639991\\
3.4848	-0.063978\\
3.4874	-0.0639565\\
3.4899	-0.0639356\\
3.4925	-0.0639135\\
3.4951	-0.0638912\\
3.4976	-0.0638693\\
3.5002	-0.0638463\\
3.5028	-0.0638231\\
3.5053	-0.0638004\\
3.5079	-0.0637765\\
3.5104	-0.0637533\\
3.513	-0.0637288\\
3.5156	-0.0637041\\
3.5181	-0.06368\\
3.5207	-0.0636547\\
3.5233	-0.0636291\\
3.5258	-0.0636043\\
3.5284	-0.0635782\\
3.531	-0.0635518\\
3.5335	-0.0635261\\
3.5361	-0.0634992\\
3.5386	-0.0634731\\
3.5412	-0.0634457\\
3.5438	-0.063418\\
3.5463	-0.0633912\\
3.5489	-0.063363\\
3.5515	-0.0633346\\
3.554	-0.0633071\\
3.5566	-0.0632782\\
3.5592	-0.0632491\\
3.5617	-0.0632209\\
3.5643	-0.0631913\\
3.5668	-0.0631627\\
3.5694	-0.0631327\\
3.572	-0.0631024\\
3.5745	-0.0630732\\
3.5771	-0.0630425\\
3.5797	-0.0630117\\
3.5822	-0.0629818\\
3.5848	-0.0629505\\
3.5874	-0.0629191\\
3.5899	-0.0628886\\
3.5925	-0.0628568\\
3.595	-0.0628259\\
3.5976	-0.0627937\\
3.6002	-0.0627613\\
3.6027	-0.06273\\
3.6053	-0.0626972\\
3.6079	-0.0626642\\
3.6104	-0.0626324\\
3.613	-0.0625991\\
3.6156	-0.0625656\\
3.6181	-0.0625333\\
3.6207	-0.0624995\\
3.6232	-0.0624669\\
3.6258	-0.0624328\\
3.6284	-0.0623986\\
3.6309	-0.0623655\\
3.6335	-0.062331\\
3.6361	-0.0622963\\
3.6386	-0.0622628\\
3.6412	-0.0622279\\
3.6438	-0.0621928\\
3.6463	-0.0621589\\
3.6489	-0.0621236\\
3.6515	-0.0620881\\
3.654	-0.0620539\\
3.6566	-0.0620182\\
3.6591	-0.0619837\\
3.6617	-0.0619478\\
3.6643	-0.0619117\\
3.6668	-0.061877\\
3.6694	-0.0618407\\
3.672	-0.0618043\\
3.6745	-0.0617693\\
3.6771	-0.0617327\\
3.6797	-0.0616961\\
3.6822	-0.0616607\\
3.6848	-0.0616239\\
3.6873	-0.0615884\\
3.6899	-0.0615514\\
3.6925	-0.0615143\\
3.695	-0.0614786\\
3.6976	-0.0614414\\
3.7002	-0.061404\\
3.7027	-0.0613681\\
3.7053	-0.0613307\\
3.7079	-0.0612932\\
3.7104	-0.061257\\
3.713	-0.0612194\\
3.7155	-0.0611832\\
3.7181	-0.0611454\\
3.7207	-0.0611077\\
3.7232	-0.0610713\\
3.7258	-0.0610334\\
3.7284	-0.0609955\\
3.7309	-0.060959\\
3.7335	-0.0609209\\
3.7361	-0.0608829\\
3.7386	-0.0608463\\
3.7412	-0.0608082\\
3.7437	-0.0607715\\
3.7463	-0.0607334\\
3.7489	-0.0606952\\
3.7514	-0.0606584\\
3.754	-0.0606202\\
3.7566	-0.060582\\
3.7591	-0.0605452\\
3.7617	-0.0605069\\
3.7643	-0.0604686\\
3.7668	-0.0604318\\
3.7694	-0.0603935\\
3.7719	-0.0603567\\
3.7745	-0.0603184\\
3.7771	-0.0602801\\
3.7796	-0.0602433\\
3.7822	-0.060205\\
3.7848	-0.0601667\\
3.7873	-0.0601299\\
3.7899	-0.0600917\\
3.7925	-0.0600534\\
3.795	-0.0600167\\
3.7976	-0.0599785\\
3.8002	-0.0599403\\
3.8027	-0.0599036\\
3.8053	-0.0598655\\
3.8078	-0.0598289\\
3.8104	-0.0597908\\
3.813	-0.0597527\\
3.8155	-0.0597162\\
3.8181	-0.0596782\\
3.8207	-0.0596403\\
3.8232	-0.0596039\\
3.8258	-0.059566\\
3.8284	-0.0595282\\
3.8309	-0.0594919\\
3.8335	-0.0594542\\
3.836	-0.059418\\
3.8386	-0.0593803\\
3.8412	-0.0593428\\
3.8437	-0.0593067\\
3.8463	-0.0592693\\
3.8489	-0.0592319\\
3.8514	-0.059196\\
3.854	-0.0591587\\
3.8566	-0.0591215\\
3.8591	-0.0590858\\
3.8617	-0.0590488\\
3.8642	-0.0590132\\
3.8668	-0.0589763\\
3.8694	-0.0589394\\
3.8719	-0.0589041\\
3.8745	-0.0588674\\
3.8771	-0.0588308\\
3.8796	-0.0587957\\
3.8822	-0.0587592\\
3.8848	-0.0587228\\
3.8873	-0.0586879\\
3.8899	-0.0586517\\
3.8924	-0.058617\\
3.895	-0.058581\\
3.8976	-0.0585451\\
3.9001	-0.0585106\\
3.9027	-0.0584749\\
3.9053	-0.0584393\\
3.9078	-0.0584051\\
3.9104	-0.0583696\\
3.913	-0.0583343\\
3.9155	-0.0583004\\
3.9181	-0.0582653\\
3.9206	-0.0582316\\
3.9232	-0.0581967\\
3.9258	-0.0581619\\
3.9283	-0.0581285\\
3.9309	-0.058094\\
3.9335	-0.0580595\\
3.936	-0.0580264\\
3.9386	-0.0579922\\
3.9412	-0.0579581\\
3.9437	-0.0579254\\
3.9463	-0.0578915\\
3.9488	-0.057859\\
3.9514	-0.0578254\\
3.954	-0.0577919\\
3.9565	-0.0577598\\
3.9591	-0.0577265\\
3.9617	-0.0576934\\
3.9642	-0.0576617\\
3.9668	-0.0576288\\
3.9694	-0.057596\\
3.9719	-0.0575647\\
3.9745	-0.0575322\\
3.9771	-0.0574999\\
3.9796	-0.0574689\\
3.9822	-0.0574369\\
3.9847	-0.0574062\\
3.9873	-0.0573744\\
3.9899	-0.0573428\\
3.9924	-0.0573125\\
3.995	-0.0572811\\
3.9976	-0.0572499\\
4.0001	-0.0572201\\
4.0027	-0.0571892\\
4.0053	-0.0571584\\
4.0078	-0.057129\\
4.0104	-0.0570985\\
4.0129	-0.0570694\\
4.0155	-0.0570392\\
4.0181	-0.0570092\\
4.0206	-0.0569805\\
4.0232	-0.0569508\\
4.0258	-0.0569213\\
4.0283	-0.0568931\\
4.0309	-0.0568638\\
4.0335	-0.0568348\\
4.036	-0.056807\\
4.0386	-0.0567783\\
4.0411	-0.0567508\\
4.0437	-0.0567224\\
4.0463	-0.0566942\\
4.0488	-0.0566672\\
4.0514	-0.0566393\\
4.054	-0.0566116\\
4.0565	-0.0565851\\
4.0591	-0.0565577\\
4.0617	-0.0565305\\
4.0642	-0.0565045\\
4.0668	-0.0564776\\
4.0693	-0.0564519\\
4.0719	-0.0564254\\
4.0745	-0.056399\\
4.077	-0.0563738\\
4.0796	-0.0563478\\
4.0822	-0.056322\\
4.0847	-0.0562974\\
4.0873	-0.0562719\\
4.0899	-0.0562466\\
4.0924	-0.0562225\\
4.095	-0.0561975\\
4.0975	-0.0561737\\
4.1001	-0.0561492\\
4.1027	-0.0561248\\
4.1052	-0.0561015\\
4.1078	-0.0560775\\
4.1104	-0.0560537\\
4.1129	-0.056031\\
4.1155	-0.0560075\\
4.1181	-0.0559843\\
4.1206	-0.0559621\\
4.1232	-0.0559392\\
4.1258	-0.0559165\\
4.1283	-0.0558949\\
4.1309	-0.0558726\\
4.1334	-0.0558513\\
4.136	-0.0558294\\
4.1386	-0.0558076\\
4.1411	-0.0557869\\
4.1437	-0.0557655\\
4.1463	-0.0557444\\
4.1488	-0.0557242\\
4.1514	-0.0557035\\
4.154	-0.0556829\\
4.1565	-0.0556633\\
4.1591	-0.0556431\\
4.1616	-0.0556239\\
4.1642	-0.0556041\\
4.1668	-0.0555845\\
4.1693	-0.0555659\\
4.1719	-0.0555467\\
4.1745	-0.0555277\\
4.177	-0.0555097\\
4.1796	-0.0554911\\
4.1822	-0.0554727\\
4.1847	-0.0554552\\
4.1873	-0.0554372\\
4.1898	-0.0554201\\
4.1924	-0.0554025\\
4.195	-0.0553851\\
4.1975	-0.0553686\\
4.2001	-0.0553516\\
4.2027	-0.0553349\\
4.2052	-0.055319\\
4.2078	-0.0553026\\
4.2104	-0.0552864\\
4.2129	-0.0552711\\
4.2155	-0.0552553\\
4.218	-0.0552404\\
4.2206	-0.055225\\
4.2232	-0.0552099\\
4.2257	-0.0551956\\
4.2283	-0.0551808\\
4.2309	-0.0551663\\
4.2334	-0.0551526\\
4.236	-0.0551385\\
4.2386	-0.0551246\\
4.2411	-0.0551114\\
4.2437	-0.0550979\\
4.2462	-0.0550851\\
4.2488	-0.0550721\\
4.2514	-0.0550592\\
4.2539	-0.055047\\
4.2565	-0.0550346\\
4.2591	-0.0550223\\
4.2616	-0.0550108\\
4.2642	-0.0549989\\
4.2668	-0.0549873\\
4.2693	-0.0549763\\
4.2719	-0.0549651\\
4.2744	-0.0549546\\
4.277	-0.0549438\\
4.2796	-0.0549332\\
4.2821	-0.0549232\\
4.2847	-0.054913\\
4.2873	-0.0549031\\
4.2898	-0.0548937\\
4.2924	-0.0548841\\
4.295	-0.0548748\\
4.2975	-0.054866\\
4.3001	-0.0548571\\
4.3027	-0.0548484\\
4.3052	-0.0548402\\
4.3078	-0.0548319\\
4.3103	-0.0548241\\
4.3129	-0.0548162\\
4.3155	-0.0548085\\
4.318	-0.0548013\\
4.3206	-0.054794\\
4.3232	-0.0547869\\
4.3257	-0.0547803\\
4.3283	-0.0547737\\
4.3309	-0.0547672\\
4.3334	-0.0547612\\
4.336	-0.0547551\\
4.3385	-0.0547495\\
4.3411	-0.0547439\\
4.3437	-0.0547384\\
4.3462	-0.0547334\\
4.3488	-0.0547283\\
4.3514	-0.0547235\\
4.3539	-0.054719\\
4.3565	-0.0547146\\
4.3591	-0.0547103\\
4.3616	-0.0547065\\
4.3642	-0.0547026\\
4.3667	-0.0546991\\
4.3693	-0.0546957\\
4.3719	-0.0546924\\
4.3744	-0.0546895\\
4.377	-0.0546866\\
4.3796	-0.054684\\
4.3821	-0.0546816\\
4.3847	-0.0546794\\
4.3873	-0.0546773\\
4.3898	-0.0546755\\
4.3924	-0.0546738\\
4.3949	-0.0546724\\
4.3975	-0.0546711\\
4.4001	-0.05467\\
4.4026	-0.0546691\\
4.4052	-0.0546684\\
4.4078	-0.0546679\\
4.4103	-0.0546675\\
4.4129	-0.0546674\\
4.4155	-0.0546674\\
4.418	-0.0546676\\
4.4206	-0.0546681\\
4.4231	-0.0546686\\
4.4257	-0.0546694\\
4.4283	-0.0546704\\
4.4308	-0.0546715\\
4.4334	-0.0546728\\
4.436	-0.0546743\\
4.4385	-0.054676\\
4.4411	-0.0546779\\
4.4437	-0.0546799\\
4.4462	-0.0546821\\
4.4488	-0.0546845\\
4.4514	-0.0546871\\
4.4539	-0.0546897\\
4.4565	-0.0546927\\
4.459	-0.0546957\\
4.4616	-0.054699\\
4.4642	-0.0547025\\
4.4667	-0.054706\\
4.4693	-0.0547098\\
4.4719	-0.0547138\\
4.4744	-0.0547178\\
4.477	-0.0547221\\
4.4796	-0.0547266\\
4.4821	-0.054731\\
4.4847	-0.0547359\\
4.4872	-0.0547407\\
4.4898	-0.0547458\\
4.4924	-0.0547511\\
4.4949	-0.0547564\\
4.4975	-0.054762\\
4.5001	-0.0547678\\
4.5026	-0.0547735\\
4.5052	-0.0547796\\
4.5078	-0.0547859\\
4.5103	-0.0547921\\
4.5129	-0.0547986\\
4.5154	-0.0548051\\
4.518	-0.054812\\
4.5206	-0.054819\\
4.5231	-0.0548259\\
4.5257	-0.0548332\\
4.5283	-0.0548407\\
4.5308	-0.054848\\
4.5334	-0.0548558\\
4.536	-0.0548637\\
4.5385	-0.0548714\\
4.5411	-0.0548796\\
4.5436	-0.0548876\\
4.5462	-0.0548961\\
4.5488	-0.0549047\\
4.5513	-0.0549131\\
4.5539	-0.054922\\
4.5565	-0.054931\\
4.559	-0.0549398\\
4.5616	-0.0549491\\
4.5642	-0.0549585\\
4.5667	-0.0549676\\
4.5693	-0.0549773\\
4.5718	-0.0549867\\
4.5744	-0.0549966\\
4.577	-0.0550066\\
4.5795	-0.0550164\\
4.5821	-0.0550267\\
4.5847	-0.0550371\\
4.5872	-0.0550472\\
4.5898	-0.0550578\\
4.5924	-0.0550685\\
4.5949	-0.055079\\
4.5975	-0.0550899\\
4.6001	-0.055101\\
4.6026	-0.0551118\\
4.6052	-0.0551231\\
4.6077	-0.055134\\
4.6103	-0.0551455\\
4.6129	-0.0551571\\
4.6154	-0.0551684\\
4.618	-0.0551802\\
4.6206	-0.0551921\\
4.6231	-0.0552037\\
4.6257	-0.0552158\\
4.6283	-0.055228\\
4.6308	-0.0552398\\
4.6334	-0.0552522\\
4.6359	-0.0552642\\
4.6385	-0.0552767\\
4.6411	-0.0552894\\
4.6436	-0.0553016\\
4.6462	-0.0553145\\
4.6488	-0.0553274\\
4.6513	-0.0553398\\
4.6539	-0.0553529\\
4.6565	-0.0553661\\
4.659	-0.0553788\\
4.6616	-0.0553921\\
4.6641	-0.0554049\\
4.6667	-0.0554184\\
4.6693	-0.0554319\\
4.6718	-0.0554449\\
4.6744	-0.0554586\\
4.677	-0.0554723\\
4.6795	-0.0554855\\
4.6821	-0.0554994\\
4.6847	-0.0555133\\
4.6872	-0.0555267\\
4.6898	-0.0555407\\
4.6923	-0.0555543\\
4.6949	-0.0555684\\
4.6975	-0.0555826\\
4.7	-0.0555963\\
4.7026	-0.0556105\\
4.7052	-0.0556249\\
4.7077	-0.0556387\\
4.7103	-0.0556531\\
4.7129	-0.0556676\\
4.7154	-0.0556816\\
4.718	-0.0556961\\
4.7205	-0.0557101\\
4.7231	-0.0557248\\
4.7257	-0.0557394\\
4.7282	-0.0557536\\
4.7308	-0.0557683\\
4.7334	-0.0557831\\
4.7359	-0.0557973\\
4.7385	-0.0558121\\
4.7411	-0.0558269\\
4.7436	-0.0558412\\
4.7462	-0.0558561\\
4.7487	-0.0558705\\
4.7513	-0.0558854\\
4.7539	-0.0559004\\
4.7564	-0.0559147\\
4.759	-0.0559297\\
4.7616	-0.0559447\\
4.7641	-0.0559592\\
4.7667	-0.0559742\\
4.7693	-0.0559892\\
4.7718	-0.0560037\\
4.7744	-0.0560187\\
4.777	-0.0560338\\
4.7795	-0.0560482\\
4.7821	-0.0560633\\
4.7846	-0.0560777\\
4.7872	-0.0560928\\
4.7898	-0.0561078\\
4.7923	-0.0561223\\
4.7949	-0.0561373\\
4.7975	-0.0561523\\
4.8	-0.0561668\\
4.8026	-0.0561818\\
4.8052	-0.0561968\\
4.8077	-0.0562111\\
4.8103	-0.0562261\\
4.8128	-0.0562405\\
4.8154	-0.0562554\\
4.818	-0.0562703\\
4.8205	-0.0562846\\
4.8231	-0.0562994\\
4.8257	-0.0563143\\
4.8282	-0.0563285\\
4.8308	-0.0563433\\
4.8334	-0.056358\\
4.8359	-0.0563722\\
4.8385	-0.0563869\\
4.841	-0.0564009\\
4.8436	-0.0564155\\
4.8462	-0.0564301\\
4.8487	-0.0564441\\
4.8513	-0.0564586\\
4.8539	-0.056473\\
4.8564	-0.0564869\\
4.859	-0.0565013\\
4.8616	-0.0565156\\
4.8641	-0.0565293\\
4.8667	-0.0565436\\
4.8692	-0.0565572\\
4.8718	-0.0565713\\
4.8744	-0.0565854\\
4.8769	-0.0565989\\
4.8795	-0.0566128\\
4.8821	-0.0566268\\
4.8846	-0.0566401\\
4.8872	-0.0566539\\
4.8898	-0.0566676\\
4.8923	-0.0566808\\
4.8949	-0.0566944\\
4.8974	-0.0567074\\
4.9	-0.0567209\\
4.9026	-0.0567343\\
4.9051	-0.0567471\\
4.9077	-0.0567604\\
4.9103	-0.0567736\\
4.9128	-0.0567862\\
4.9154	-0.0567993\\
4.918	-0.0568123\\
4.9205	-0.0568247\\
4.9231	-0.0568375\\
4.9257	-0.0568503\\
4.9282	-0.0568625\\
4.9308	-0.0568751\\
4.9333	-0.0568871\\
4.9359	-0.0568996\\
4.9385	-0.0569119\\
4.941	-0.0569237\\
4.9436	-0.0569359\\
4.9462	-0.056948\\
4.9487	-0.0569595\\
4.9513	-0.0569715\\
4.9539	-0.0569833\\
4.9564	-0.0569946\\
4.959	-0.0570062\\
4.9615	-0.0570173\\
4.9641	-0.0570288\\
4.9667	-0.0570402\\
4.9692	-0.057051\\
4.9718	-0.0570622\\
4.9744	-0.0570732\\
4.9769	-0.0570838\\
4.9795	-0.0570947\\
4.9821	-0.0571054\\
4.9846	-0.0571157\\
4.9872	-0.0571263\\
4.9897	-0.0571363\\
4.9923	-0.0571467\\
4.9949	-0.0571569\\
4.9974	-0.0571667\\
5	-0.0571768\\
};
\addplot [color=black, line width=1.4pt, forget plot]
  table[row sep=crcr]{%
-0.125	3.52154e-08\\
-0.1224	3.59452e-08\\
-0.1199	3.66469e-08\\
-0.1173	3.73766e-08\\
-0.1147	3.81063e-08\\
-0.1122	3.88079e-08\\
-0.1096	3.95375e-08\\
-0.1071	4.0239e-08\\
-0.1045	4.09686e-08\\
-0.1019	4.16981e-08\\
-0.0994	4.23994e-08\\
-0.0968	4.31289e-08\\
-0.0942	4.38583e-08\\
-0.0917	4.45595e-08\\
-0.0891	4.52888e-08\\
-0.0865	4.60181e-08\\
-0.084	4.67192e-08\\
-0.0814	4.74484e-08\\
-0.0789	4.81496e-08\\
-0.0763	4.88786e-08\\
-0.0737	4.96077e-08\\
-0.0712	5.03087e-08\\
-0.0686	5.10377e-08\\
-0.066	5.17666e-08\\
-0.0635	5.24675e-08\\
-0.0609	5.31963e-08\\
-0.0583	5.39252e-08\\
-0.0558	5.46259e-08\\
-0.0532	5.53547e-08\\
-0.0507	5.60553e-08\\
-0.0481	5.6784e-08\\
-0.0455	5.75126e-08\\
-0.043	5.82132e-08\\
-0.0404	5.89417e-08\\
-0.0378	5.96702e-08\\
-0.0353	6.03707e-08\\
-0.0327	6.10991e-08\\
-0.0301	6.18275e-08\\
-0.0276	6.25278e-08\\
-0.025	6.32561e-08\\
-0.0224	6.39843e-08\\
-0.0199	6.46845e-08\\
-0.0173	6.54127e-08\\
-0.0148	6.61128e-08\\
-0.0122	6.68409e-08\\
-0.0096	6.75689e-08\\
-0.0071	6.82689e-08\\
-0.0045	6.89968e-08\\
-0.0019	6.97247e-08\\
0.0006	7.04246e-08\\
0.0032	0.000237545\\
0.0058	0.000833566\\
0.0083	0.00153356\\
0.0109	0.00244569\\
0.0134	0.00343409\\
0.016	0.00443222\\
0.0186	0.0053072\\
0.0211	0.00603767\\
0.0237	0.00674228\\
0.0263	0.00744195\\
0.0288	0.00812027\\
0.0314	0.00881368\\
0.034	0.00947445\\
0.0365	0.0100714\\
0.0391	0.0106546\\
0.0416	0.0111835\\
0.0442	0.0117045\\
0.0468	0.0122009\\
0.0493	0.0126608\\
0.0519	0.0131245\\
0.0545	0.0135702\\
0.057	0.0139744\\
0.0596	0.0143638\\
0.0622	0.0147249\\
0.0647	0.0150557\\
0.0673	0.0153936\\
0.0698	0.0157147\\
0.0724	0.0160375\\
0.075	0.0163388\\
0.0775	0.0166055\\
0.0801	0.0168666\\
0.0827	0.0171239\\
0.0852	0.0173749\\
0.0878	0.0176368\\
0.0904	0.0178895\\
0.0929	0.0181167\\
0.0955	0.0183387\\
0.098	0.0185484\\
0.1006	0.0187729\\
0.1032	0.019006\\
0.1057	0.0192314\\
0.1083	0.0194577\\
0.1109	0.0196724\\
0.1134	0.0198738\\
0.116	0.0200884\\
0.1186	0.0203143\\
0.1211	0.0205393\\
0.1237	0.0207721\\
0.1263	0.0209969\\
0.1288	0.0212066\\
0.1314	0.0214261\\
0.1339	0.0216462\\
0.1365	0.0218858\\
0.1391	0.0221295\\
0.1416	0.0223596\\
0.1442	0.0225912\\
0.1468	0.0228195\\
0.1493	0.0230439\\
0.1519	0.0232868\\
0.1545	0.0235363\\
0.157	0.0237748\\
0.1596	0.0240151\\
0.1621	0.0242393\\
0.1647	0.0244718\\
0.1673	0.0247102\\
0.1698	0.0249456\\
0.1724	0.0251912\\
0.175	0.0254301\\
0.1775	0.0256506\\
0.1801	0.0258739\\
0.1827	0.026098\\
0.1852	0.0263177\\
0.1878	0.026548\\
0.1903	0.0267647\\
0.1929	0.0269799\\
0.1955	0.0271851\\
0.198	0.0273782\\
0.2006	0.0275801\\
0.2032	0.0277837\\
0.2057	0.0279766\\
0.2083	0.028168\\
0.2109	0.0283476\\
0.2134	0.0285122\\
0.216	0.0286807\\
0.2185	0.0288434\\
0.2211	0.0290111\\
0.2237	0.0291718\\
0.2262	0.0293154\\
0.2288	0.0294539\\
0.2314	0.0295864\\
0.2339	0.0297126\\
0.2365	0.0298435\\
0.2391	0.0299698\\
0.2416	0.0300824\\
0.2442	0.0301882\\
0.2467	0.0302821\\
0.2493	0.0303764\\
0.2519	0.0304705\\
0.2544	0.0305591\\
0.257	0.030645\\
0.2596	0.0307208\\
0.2621	0.030785\\
0.2647	0.0308467\\
0.2673	0.0309076\\
0.2698	0.0309662\\
0.2724	0.0310238\\
0.2749	0.0310723\\
0.2775	0.0311141\\
0.2801	0.0311496\\
0.2826	0.0311821\\
0.2852	0.0312168\\
0.2878	0.031251\\
0.2903	0.0312798\\
0.2929	0.0313028\\
0.2955	0.0313193\\
0.298	0.0313326\\
0.3006	0.0313471\\
0.3032	0.0313632\\
0.3057	0.0313774\\
0.3083	0.0313875\\
0.3108	0.0313916\\
0.3134	0.0313923\\
0.316	0.0313931\\
0.3185	0.0313959\\
0.3211	0.0314\\
0.3237	0.0314019\\
0.3262	0.0313995\\
0.3288	0.0313929\\
0.3314	0.0313854\\
0.3339	0.03138\\
0.3365	0.0313769\\
0.339	0.031374\\
0.3416	0.0313683\\
0.3442	0.0313587\\
0.3467	0.0313475\\
0.3493	0.0313369\\
0.3519	0.0313291\\
0.3544	0.0313232\\
0.357	0.0313161\\
0.3596	0.0313058\\
0.3621	0.0312933\\
0.3647	0.03128\\
0.3672	0.0312693\\
0.3698	0.0312607\\
0.3724	0.0312526\\
0.3749	0.0312428\\
0.3775	0.0312296\\
0.3801	0.0312149\\
0.3826	0.0312018\\
0.3852	0.0311904\\
0.3878	0.0311806\\
0.3903	0.0311703\\
0.3929	0.0311569\\
0.3954	0.0311416\\
0.398	0.0311254\\
0.4006	0.0311107\\
0.4031	0.0310984\\
0.4057	0.0310858\\
0.4083	0.0310711\\
0.4108	0.0310542\\
0.4134	0.0310351\\
0.416	0.0310164\\
0.4185	0.031\\
0.4211	0.0309837\\
0.4236	0.030967\\
0.4262	0.0309468\\
0.4288	0.0309242\\
0.4313	0.0309016\\
0.4339	0.0308791\\
0.4365	0.0308576\\
0.439	0.0308366\\
0.4416	0.0308126\\
0.4442	0.0307858\\
0.4467	0.0307582\\
0.4493	0.0307294\\
0.4519	0.0307015\\
0.4544	0.030675\\
0.457	0.030646\\
0.4595	0.0306154\\
0.4621	0.0305812\\
0.4647	0.0305458\\
0.4672	0.0305122\\
0.4698	0.0304778\\
0.4724	0.0304427\\
0.4749	0.0304069\\
0.4775	0.0303668\\
0.4801	0.0303249\\
0.4826	0.0302843\\
0.4852	0.0302425\\
0.4877	0.0302023\\
0.4903	0.030159\\
0.4929	0.0301132\\
0.4954	0.0300668\\
0.498	0.0300175\\
0.5006	0.0299682\\
0.5031	0.0299212\\
0.5057	0.0298715\\
0.5083	0.0298197\\
0.5108	0.0297675\\
0.5134	0.0297116\\
0.5159	0.0296573\\
0.5185	0.0296012\\
0.5211	0.0295448\\
0.5236	0.0294894\\
0.5262	0.0294294\\
0.5288	0.0293674\\
0.5313	0.0293067\\
0.5339	0.0292435\\
0.5365	0.0291805\\
0.539	0.0291192\\
0.5416	0.0290537\\
0.5441	0.0289885\\
0.5467	0.0289193\\
0.5493	0.0288495\\
0.5518	0.0287825\\
0.5544	0.0287126\\
0.557	0.0286415\\
0.5595	0.0285711\\
0.5621	0.0284961\\
0.5647	0.02842\\
0.5672	0.0283468\\
0.5698	0.0282706\\
0.5723	0.0281967\\
0.5749	0.0281181\\
0.5775	0.0280376\\
0.58	0.0279587\\
0.5826	0.0278762\\
0.5852	0.0277936\\
0.5877	0.0277139\\
0.5903	0.0276297\\
0.5929	0.0275436\\
0.5954	0.0274592\\
0.598	0.0273704\\
0.6006	0.0272813\\
0.6031	0.0271954\\
0.6057	0.0271052\\
0.6082	0.027017\\
0.6108	0.0269234\\
0.6134	0.0268283\\
0.6159	0.0267362\\
0.6185	0.0266402\\
0.6211	0.0265436\\
0.6236	0.0264495\\
0.6262	0.0263498\\
0.6288	0.0262483\\
0.6313	0.0261496\\
0.6339	0.0260465\\
0.6364	0.0259469\\
0.639	0.0258424\\
0.6416	0.0257362\\
0.6441	0.0256323\\
0.6467	0.0255227\\
0.6493	0.0254122\\
0.6518	0.0253055\\
0.6544	0.0251937\\
0.657	0.0250805\\
0.6595	0.0249699\\
0.6621	0.0248531\\
0.6646	0.0247395\\
0.6672	0.0246207\\
0.6698	0.0245011\\
0.6723	0.024385\\
0.6749	0.0242626\\
0.6775	0.0241382\\
0.68	0.0240171\\
0.6826	0.02389\\
0.6852	0.0237622\\
0.6877	0.0236383\\
0.6903	0.023508\\
0.6928	0.0233809\\
0.6954	0.0232468\\
0.698	0.0231113\\
0.7005	0.0229801\\
0.7031	0.0228427\\
0.7057	0.0227041\\
0.7082	0.022569\\
0.7108	0.0224266\\
0.7134	0.0222825\\
0.7159	0.0221428\\
0.7185	0.0219965\\
0.721	0.0218547\\
0.7236	0.0217057\\
0.7262	0.0215548\\
0.7287	0.0214078\\
0.7313	0.0212535\\
0.7339	0.021098\\
0.7364	0.0209475\\
0.739	0.0207896\\
0.7416	0.0206299\\
0.7441	0.0204744\\
0.7467	0.020311\\
0.7492	0.0201527\\
0.7518	0.0199868\\
0.7544	0.0198198\\
0.7569	0.0196576\\
0.7595	0.019487\\
0.7621	0.0193145\\
0.7646	0.0191472\\
0.7672	0.018972\\
0.7698	0.0187956\\
0.7723	0.0186247\\
0.7749	0.0184451\\
0.7775	0.0182636\\
0.78	0.0180874\\
0.7826	0.0179028\\
0.7851	0.0177242\\
0.7877	0.0175372\\
0.7903	0.0173485\\
0.7928	0.0171653\\
0.7954	0.016973\\
0.798	0.0167792\\
0.8005	0.0165916\\
0.8031	0.0163953\\
0.8057	0.0161975\\
0.8082	0.0160057\\
0.8108	0.0158044\\
0.8133	0.0156093\\
0.8159	0.0154049\\
0.8185	0.0151994\\
0.821	0.0150005\\
0.8236	0.0147921\\
0.8262	0.0145819\\
0.8287	0.0143781\\
0.8313	0.0141647\\
0.8339	0.01395\\
0.8364	0.0137424\\
0.839	0.0135251\\
0.8415	0.0133146\\
0.8441	0.0130939\\
0.8467	0.0128716\\
0.8492	0.0126566\\
0.8518	0.0124318\\
0.8544	0.0122056\\
0.8569	0.0119867\\
0.8595	0.0117574\\
0.8621	0.0115264\\
0.8646	0.011303\\
0.8672	0.0110693\\
0.8697	0.0108435\\
0.8723	0.0106073\\
0.8749	0.0103695\\
0.8774	0.0101392\\
0.88	0.00989828\\
0.8826	0.009656\\
0.8851	0.00942188\\
0.8877	0.00917715\\
0.8903	0.00893096\\
0.8928	0.00869274\\
0.8954	0.00844344\\
0.8979	0.0082024\\
0.9005	0.00795049\\
0.9031	0.00769738\\
0.9056	0.00745274\\
0.9082	0.00719685\\
0.9108	0.0069394\\
0.9133	0.00669048\\
0.9159	0.00643032\\
0.9185	0.00616898\\
0.921	0.00591652\\
0.9236	0.00565261\\
0.9262	0.0053872\\
0.9287	0.0051306\\
0.9313	0.00486241\\
0.9338	0.0046034\\
0.9364	0.00433289\\
0.939	0.0040611\\
0.9415	0.00379844\\
0.9441	0.00352384\\
0.9467	0.00324786\\
0.9492	0.00298132\\
0.9518	0.002703\\
0.9544	0.00242351\\
0.9569	0.00215353\\
0.9595	0.00187138\\
0.962	0.00159877\\
0.9646	0.00131401\\
0.9672	0.00102811\\
0.9697	0.000752143\\
0.9723	0.000463967\\
0.9749	0.000174478\\
0.9774	-0.000105154\\
0.98	-0.000397234\\
0.9826	-0.000690476\\
0.9851	-0.000973461\\
0.9877	-0.00126885\\
0.9902	-0.00155403\\
0.9928	-0.00185187\\
0.9954	-0.002151\\
0.9979	-0.00243972\\
1.0005	-0.00274105\\
1.0031	-0.00304342\\
1.0056	-0.00333521\\
1.0082	-0.00363985\\
1.0108	-0.00394572\\
1.0133	-0.00424094\\
1.0159	-0.00454902\\
1.0184	-0.00484619\\
1.021	-0.00515624\\
1.0236	-0.0054674\\
1.0261	-0.00576769\\
1.0287	-0.00608112\\
1.0313	-0.00639561\\
1.0338	-0.00669893\\
1.0364	-0.00701531\\
1.039	-0.00733269\\
1.0415	-0.00763888\\
1.0441	-0.00795839\\
1.0466	-0.0082666\\
1.0492	-0.00858808\\
1.0518	-0.00891046\\
1.0543	-0.00922129\\
1.0569	-0.00954551\\
1.0595	-0.00987075\\
1.062	-0.0101844\\
1.0646	-0.0105116\\
1.0672	-0.0108396\\
1.0697	-0.0111558\\
1.0723	-0.0114854\\
1.0748	-0.0118033\\
1.0774	-0.0121348\\
1.08	-0.0124673\\
1.0825	-0.0127877\\
1.0851	-0.0131217\\
1.0877	-0.0134564\\
1.0902	-0.0137791\\
1.0928	-0.0141155\\
1.0954	-0.0144528\\
1.0979	-0.0147778\\
1.1005	-0.0151165\\
1.1031	-0.015456\\
1.1056	-0.015783\\
1.1082	-0.0161239\\
1.1107	-0.0164525\\
1.1133	-0.0167949\\
1.1159	-0.0171379\\
1.1184	-0.0174684\\
1.121	-0.0178127\\
1.1236	-0.0181576\\
1.1261	-0.01849\\
1.1287	-0.0188363\\
1.1313	-0.0191832\\
1.1338	-0.0195173\\
1.1364	-0.0198653\\
1.1389	-0.0202004\\
1.1415	-0.0205495\\
1.1441	-0.0208993\\
1.1466	-0.0212361\\
1.1492	-0.0215868\\
1.1518	-0.0219381\\
1.1543	-0.0222762\\
1.1569	-0.0226283\\
1.1595	-0.022981\\
1.162	-0.0233205\\
1.1646	-0.0236741\\
1.1671	-0.0240144\\
1.1697	-0.0243687\\
1.1723	-0.0247233\\
1.1748	-0.0250647\\
1.1774	-0.0254201\\
1.18	-0.0257759\\
1.1825	-0.0261184\\
1.1851	-0.0264747\\
1.1877	-0.0268314\\
1.1902	-0.0271746\\
1.1928	-0.0275318\\
1.1953	-0.0278755\\
1.1979	-0.0282332\\
1.2005	-0.0285911\\
1.203	-0.0289354\\
1.2056	-0.0292936\\
1.2082	-0.029652\\
1.2107	-0.0299968\\
1.2133	-0.0303556\\
1.2159	-0.0307144\\
1.2184	-0.0310595\\
1.221	-0.0314185\\
1.2235	-0.0317637\\
1.2261	-0.0321227\\
1.2287	-0.0324819\\
1.2312	-0.0328272\\
1.2338	-0.0331863\\
1.2364	-0.0335454\\
1.2389	-0.0338906\\
1.2415	-0.0342496\\
1.2441	-0.0346085\\
1.2466	-0.0349535\\
1.2492	-0.0353122\\
1.2518	-0.0356707\\
1.2543	-0.0360153\\
1.2569	-0.0363735\\
1.2594	-0.0367177\\
1.262	-0.0370756\\
1.2646	-0.0374332\\
1.2671	-0.0377768\\
1.2697	-0.0381338\\
1.2723	-0.0384906\\
1.2748	-0.0388333\\
1.2774	-0.0391895\\
1.28	-0.0395453\\
1.2825	-0.0398871\\
1.2851	-0.0402421\\
1.2876	-0.0405831\\
1.2902	-0.0409373\\
1.2928	-0.0412911\\
1.2953	-0.0416309\\
1.2979	-0.0419838\\
1.3005	-0.0423361\\
1.303	-0.0426744\\
1.3056	-0.0430256\\
1.3082	-0.0433763\\
1.3107	-0.043713\\
1.3133	-0.0440626\\
1.3158	-0.0443981\\
1.3184	-0.0447464\\
1.321	-0.045094\\
1.3235	-0.0454276\\
1.3261	-0.0457738\\
1.3287	-0.0461193\\
1.3312	-0.0464509\\
1.3338	-0.0467949\\
1.3364	-0.0471381\\
1.3389	-0.0474674\\
1.3415	-0.047809\\
1.344	-0.0481366\\
1.3466	-0.0484766\\
1.3492	-0.0488156\\
1.3517	-0.0491407\\
1.3543	-0.049478\\
1.3569	-0.0498142\\
1.3594	-0.0501366\\
1.362	-0.050471\\
1.3646	-0.0508043\\
1.3671	-0.0511238\\
1.3697	-0.0514551\\
1.3722	-0.0517726\\
1.3748	-0.0521018\\
1.3774	-0.0524298\\
1.3799	-0.0527442\\
1.3825	-0.05307\\
1.3851	-0.0533947\\
1.3876	-0.0537057\\
1.3902	-0.0540279\\
1.3928	-0.054349\\
1.3953	-0.0546565\\
1.3979	-0.0549752\\
1.4005	-0.0552925\\
1.403	-0.0555963\\
1.4056	-0.0559111\\
1.4081	-0.0562124\\
1.4107	-0.0565245\\
1.4133	-0.0568352\\
1.4158	-0.0571327\\
1.4184	-0.0574407\\
1.421	-0.0577472\\
1.4235	-0.0580406\\
1.4261	-0.0583442\\
1.4287	-0.0586464\\
1.4312	-0.0589356\\
1.4338	-0.0592349\\
1.4363	-0.0595211\\
1.4389	-0.0598173\\
1.4415	-0.060112\\
1.444	-0.0603938\\
1.4466	-0.0606854\\
1.4492	-0.0609753\\
1.4517	-0.0612525\\
1.4543	-0.0615392\\
1.4569	-0.0618243\\
1.4594	-0.0620968\\
1.462	-0.0623785\\
1.4645	-0.0626478\\
1.4671	-0.0629262\\
1.4697	-0.0632029\\
1.4722	-0.0634672\\
1.4748	-0.0637404\\
1.4774	-0.0640119\\
1.4799	-0.0642712\\
1.4825	-0.0645391\\
1.4851	-0.0648052\\
1.4876	-0.0650593\\
1.4902	-0.0653218\\
1.4927	-0.0655725\\
1.4953	-0.0658313\\
1.4979	-0.0660883\\
1.5004	-0.0663336\\
1.503	-0.0665869\\
1.5056	-0.0668382\\
1.5081	-0.0670781\\
1.5107	-0.0673256\\
1.5133	-0.0675713\\
1.5158	-0.0678056\\
1.5184	-0.0680474\\
1.5209	-0.068278\\
1.5235	-0.0685159\\
1.5261	-0.0687518\\
1.5286	-0.0689767\\
1.5312	-0.0692087\\
1.5338	-0.0694387\\
1.5363	-0.0696579\\
1.5389	-0.0698838\\
1.5415	-0.0701078\\
1.544	-0.0703211\\
1.5466	-0.0705411\\
1.5491	-0.0707506\\
1.5517	-0.0709664\\
1.5543	-0.0711802\\
1.5568	-0.0713838\\
1.5594	-0.0715935\\
1.562	-0.0718011\\
1.5645	-0.0719987\\
1.5671	-0.0722022\\
1.5697	-0.0724036\\
1.5722	-0.0725952\\
1.5748	-0.0727924\\
1.5774	-0.0729875\\
1.5799	-0.0731731\\
1.5825	-0.073364\\
1.585	-0.0735456\\
1.5876	-0.0737324\\
1.5902	-0.073917\\
1.5927	-0.0740925\\
1.5953	-0.0742729\\
1.5979	-0.0744512\\
1.6004	-0.0746206\\
1.603	-0.0747947\\
1.6056	-0.0749667\\
1.6081	-0.07513\\
1.6107	-0.0752977\\
1.6132	-0.075457\\
1.6158	-0.0756205\\
1.6184	-0.0757819\\
1.6209	-0.075935\\
1.6235	-0.0760922\\
1.6261	-0.0762472\\
1.6286	-0.0763943\\
1.6312	-0.0765451\\
1.6338	-0.0766938\\
1.6363	-0.0768348\\
1.6389	-0.0769793\\
1.6414	-0.0771163\\
1.644	-0.0772566\\
1.6466	-0.0773949\\
1.6491	-0.0775258\\
1.6517	-0.0776599\\
1.6543	-0.0777918\\
1.6568	-0.0779167\\
1.6594	-0.0780445\\
1.662	-0.0781703\\
1.6645	-0.0782892\\
1.6671	-0.0784108\\
1.6696	-0.0785258\\
1.6722	-0.0786434\\
1.6748	-0.0787589\\
1.6773	-0.0788679\\
1.6799	-0.0789794\\
1.6825	-0.0790888\\
1.685	-0.079192\\
1.6876	-0.0792974\\
1.6902	-0.0794007\\
1.6927	-0.0794981\\
1.6953	-0.0795975\\
1.6978	-0.0796911\\
1.7004	-0.0797866\\
1.703	-0.07988\\
1.7055	-0.079968\\
1.7081	-0.0800575\\
1.7107	-0.080145\\
1.7132	-0.0802274\\
1.7158	-0.0803111\\
1.7184	-0.0803929\\
1.7209	-0.0804697\\
1.7235	-0.0805477\\
1.7261	-0.0806237\\
1.7286	-0.0806951\\
1.7312	-0.0807674\\
1.7337	-0.0808352\\
1.7363	-0.0809038\\
1.7389	-0.0809706\\
1.7414	-0.081033\\
1.744	-0.0810961\\
1.7466	-0.0811574\\
1.7491	-0.0812146\\
1.7517	-0.0812723\\
1.7543	-0.0813282\\
1.7568	-0.0813803\\
1.7594	-0.0814327\\
1.7619	-0.0814814\\
1.7645	-0.0815303\\
1.7671	-0.0815775\\
1.7696	-0.0816212\\
1.7722	-0.0816649\\
1.7748	-0.081707\\
1.7773	-0.0817458\\
1.7799	-0.0817845\\
1.7825	-0.0818216\\
1.785	-0.0818556\\
1.7876	-0.0818894\\
1.7901	-0.0819203\\
1.7927	-0.0819508\\
1.7953	-0.0819798\\
1.7978	-0.0820061\\
1.8004	-0.0820319\\
1.803	-0.0820561\\
1.8055	-0.0820779\\
1.8081	-0.082099\\
1.8107	-0.0821186\\
1.8132	-0.082136\\
1.8158	-0.0821525\\
1.8183	-0.0821671\\
1.8209	-0.0821807\\
1.8235	-0.0821929\\
1.826	-0.0822032\\
1.8286	-0.0822124\\
1.8312	-0.0822203\\
1.8337	-0.0822265\\
1.8363	-0.0822315\\
1.8389	-0.0822351\\
1.8414	-0.0822373\\
1.844	-0.0822382\\
1.8465	-0.0822378\\
1.8491	-0.082236\\
1.8517	-0.0822329\\
1.8542	-0.0822287\\
1.8568	-0.082223\\
1.8594	-0.082216\\
1.8619	-0.082208\\
1.8645	-0.0821985\\
1.8671	-0.0821877\\
1.8696	-0.0821761\\
1.8722	-0.0821629\\
1.8747	-0.082149\\
1.8773	-0.0821334\\
1.8799	-0.0821166\\
1.8824	-0.0820993\\
1.885	-0.0820801\\
1.8876	-0.0820598\\
1.8901	-0.0820392\\
1.8927	-0.0820166\\
1.8953	-0.0819929\\
1.8978	-0.0819691\\
1.9004	-0.0819433\\
1.903	-0.0819163\\
1.9055	-0.0818894\\
1.9081	-0.0818604\\
1.9106	-0.0818315\\
1.9132	-0.0818004\\
1.9158	-0.0817683\\
1.9183	-0.0817365\\
1.9209	-0.0817024\\
1.9235	-0.0816673\\
1.926	-0.0816327\\
1.9286	-0.0815957\\
1.9312	-0.0815578\\
1.9337	-0.0815204\\
1.9363	-0.0814807\\
1.9388	-0.0814416\\
1.9414	-0.0814\\
1.944	-0.0813576\\
1.9465	-0.081316\\
1.9491	-0.0812718\\
1.9517	-0.0812268\\
1.9542	-0.0811827\\
1.9568	-0.0811361\\
1.9594	-0.0810886\\
1.9619	-0.0810422\\
1.9645	-0.0809931\\
1.967	-0.0809452\\
1.9696	-0.0808946\\
1.9722	-0.0808432\\
1.9747	-0.0807931\\
1.9773	-0.0807403\\
1.9799	-0.0806867\\
1.9824	-0.0806345\\
1.985	-0.0805795\\
1.9876	-0.0805238\\
1.9901	-0.0804696\\
1.9927	-0.0804125\\
1.9952	-0.0803571\\
1.9978	-0.0802987\\
2.0004	-0.0802397\\
2.0029	-0.0801824\\
2.0055	-0.0801221\\
2.0081	-0.0800613\\
2.0106	-0.0800022\\
2.0132	-0.0799401\\
2.0158	-0.0798774\\
2.0183	-0.0798167\\
2.0209	-0.0797529\\
2.0234	-0.079691\\
2.026	-0.0796261\\
2.0286	-0.0795607\\
2.0311	-0.0794973\\
2.0337	-0.0794308\\
2.0363	-0.0793639\\
2.0388	-0.079299\\
2.0414	-0.079231\\
2.044	-0.0791626\\
2.0465	-0.0790963\\
2.0491	-0.0790269\\
2.0517	-0.078957\\
2.0542	-0.0788895\\
2.0568	-0.0788187\\
2.0593	-0.0787503\\
2.0619	-0.0786787\\
2.0645	-0.0786067\\
2.067	-0.0785371\\
2.0696	-0.0784642\\
2.0722	-0.078391\\
2.0747	-0.0783203\\
2.0773	-0.0782463\\
2.0799	-0.078172\\
2.0824	-0.0781001\\
2.085	-0.0780251\\
2.0875	-0.0779526\\
2.0901	-0.0778769\\
2.0927	-0.0778008\\
2.0952	-0.0777273\\
2.0978	-0.0776506\\
2.1004	-0.0775736\\
2.1029	-0.0774992\\
2.1055	-0.0774216\\
2.1081	-0.0773437\\
2.1106	-0.0772685\\
2.1132	-0.07719\\
2.1157	-0.0771143\\
2.1183	-0.0770353\\
2.1209	-0.076956\\
2.1234	-0.0768796\\
2.126	-0.0767998\\
2.1286	-0.0767198\\
2.1311	-0.0766427\\
2.1337	-0.0765622\\
2.1363	-0.0764815\\
2.1388	-0.0764037\\
2.1414	-0.0763226\\
2.1439	-0.0762444\\
2.1465	-0.0761629\\
2.1491	-0.0760812\\
2.1516	-0.0760025\\
2.1542	-0.0759204\\
2.1568	-0.0758382\\
2.1593	-0.0757589\\
2.1619	-0.0756764\\
2.1645	-0.0755936\\
2.167	-0.0755139\\
2.1696	-0.0754309\\
2.1721	-0.0753509\\
2.1747	-0.0752676\\
2.1773	-0.0751841\\
2.1798	-0.0751037\\
2.1824	-0.07502\\
2.185	-0.0749362\\
2.1875	-0.0748554\\
2.1901	-0.0747714\\
2.1927	-0.0746872\\
2.1952	-0.0746061\\
2.1978	-0.0745218\\
2.2004	-0.0744373\\
2.2029	-0.074356\\
2.2055	-0.0742713\\
2.208	-0.0741898\\
2.2106	-0.074105\\
2.2132	-0.0740201\\
2.2157	-0.0739384\\
2.2183	-0.0738534\\
2.2209	-0.0737683\\
2.2234	-0.0736865\\
2.226	-0.0736013\\
2.2286	-0.073516\\
2.2311	-0.073434\\
2.2337	-0.0733487\\
2.2362	-0.0732666\\
2.2388	-0.0731811\\
2.2414	-0.0730957\\
2.2439	-0.0730135\\
2.2465	-0.072928\\
2.2491	-0.0728424\\
2.2516	-0.0727602\\
2.2542	-0.0726746\\
2.2568	-0.072589\\
2.2593	-0.0725067\\
2.2619	-0.0724211\\
2.2644	-0.0723388\\
2.267	-0.0722532\\
2.2696	-0.0721675\\
2.2721	-0.0720852\\
2.2747	-0.0719996\\
2.2773	-0.071914\\
2.2798	-0.0718317\\
2.2824	-0.0717462\\
2.285	-0.0716606\\
2.2875	-0.0715783\\
2.2901	-0.0714928\\
2.2926	-0.0714106\\
2.2952	-0.0713252\\
2.2978	-0.0712397\\
2.3003	-0.0711576\\
2.3029	-0.0710722\\
2.3055	-0.0709869\\
2.308	-0.0709049\\
2.3106	-0.0708196\\
2.3132	-0.0707344\\
2.3157	-0.0706525\\
2.3183	-0.0705674\\
2.3208	-0.0704856\\
2.3234	-0.0704006\\
2.326	-0.0703156\\
2.3285	-0.070234\\
2.3311	-0.0701491\\
2.3337	-0.0700643\\
2.3362	-0.0699829\\
2.3388	-0.0698982\\
2.3414	-0.0698136\\
2.3439	-0.0697323\\
2.3465	-0.0696479\\
2.349	-0.0695667\\
2.3516	-0.0694824\\
2.3542	-0.0693982\\
2.3567	-0.0693172\\
2.3593	-0.0692331\\
2.3619	-0.0691491\\
2.3644	-0.0690684\\
2.367	-0.0689846\\
2.3696	-0.0689008\\
2.3721	-0.0688204\\
2.3747	-0.0687368\\
2.3773	-0.0686533\\
2.3798	-0.0685731\\
2.3824	-0.0684898\\
2.3849	-0.0684098\\
2.3875	-0.0683267\\
2.3901	-0.0682437\\
2.3926	-0.0681639\\
2.3952	-0.0680811\\
2.3978	-0.0679984\\
2.4003	-0.067919\\
2.4029	-0.0678364\\
2.4055	-0.0677541\\
2.408	-0.0676749\\
2.4106	-0.0675927\\
2.4131	-0.0675138\\
2.4157	-0.0674319\\
2.4183	-0.06735\\
2.4208	-0.0672714\\
2.4234	-0.0671898\\
2.426	-0.0671083\\
2.4285	-0.0670301\\
2.4311	-0.0669489\\
2.4337	-0.0668677\\
2.4362	-0.0667898\\
2.4388	-0.066709\\
2.4413	-0.0666313\\
2.4439	-0.0665507\\
2.4465	-0.0664702\\
2.449	-0.0663929\\
2.4516	-0.0663127\\
2.4542	-0.0662326\\
2.4567	-0.0661557\\
2.4593	-0.0660759\\
2.4619	-0.0659962\\
2.4644	-0.0659197\\
2.467	-0.0658403\\
2.4695	-0.0657641\\
2.4721	-0.0656849\\
2.4747	-0.065606\\
2.4772	-0.0655302\\
2.4798	-0.0654515\\
2.4824	-0.0653729\\
2.4849	-0.0652975\\
2.4875	-0.0652193\\
2.4901	-0.0651412\\
2.4926	-0.0650663\\
2.4952	-0.0649885\\
2.4977	-0.0649139\\
2.5003	-0.0648364\\
2.5029	-0.0647591\\
2.5054	-0.0646849\\
2.508	-0.064608\\
2.5106	-0.0645312\\
2.5131	-0.0644575\\
2.5157	-0.064381\\
2.5183	-0.0643047\\
2.5208	-0.0642314\\
2.5234	-0.0641555\\
2.526	-0.0640797\\
2.5285	-0.064007\\
2.5311	-0.0639315\\
2.5336	-0.0638591\\
2.5362	-0.063784\\
2.5388	-0.0637091\\
2.5413	-0.0636373\\
2.5439	-0.0635627\\
2.5465	-0.0634884\\
2.549	-0.063417\\
2.5516	-0.063343\\
2.5542	-0.0632692\\
2.5567	-0.0631985\\
2.5593	-0.063125\\
2.5618	-0.0630546\\
2.5644	-0.0629816\\
2.567	-0.0629088\\
2.5695	-0.0628389\\
2.5721	-0.0627665\\
2.5747	-0.0626943\\
2.5772	-0.062625\\
2.5798	-0.0625532\\
2.5824	-0.0624816\\
2.5849	-0.062413\\
2.5875	-0.0623418\\
2.59	-0.0622735\\
2.5926	-0.0622027\\
2.5952	-0.0621322\\
2.5977	-0.0620646\\
2.6003	-0.0619945\\
2.6029	-0.0619246\\
2.6054	-0.0618576\\
2.608	-0.0617881\\
2.6106	-0.0617189\\
2.6131	-0.0616526\\
2.6157	-0.0615838\\
2.6182	-0.0615179\\
2.6208	-0.0614496\\
2.6234	-0.0613816\\
2.6259	-0.0613164\\
2.6285	-0.0612488\\
2.6311	-0.0611814\\
2.6336	-0.0611169\\
2.6362	-0.0610501\\
2.6388	-0.0609835\\
2.6413	-0.0609197\\
2.6439	-0.0608535\\
2.6464	-0.0607902\\
2.649	-0.0607246\\
2.6516	-0.0606593\\
2.6541	-0.0605967\\
2.6567	-0.0605319\\
2.6593	-0.0604673\\
2.6618	-0.0604055\\
2.6644	-0.0603414\\
2.667	-0.0602777\\
2.6695	-0.0602166\\
2.6721	-0.0601534\\
2.6746	-0.0600928\\
2.6772	-0.0600301\\
2.6798	-0.0599677\\
2.6823	-0.059908\\
2.6849	-0.0598462\\
2.6875	-0.0597846\\
2.69	-0.0597257\\
2.6926	-0.0596647\\
2.6952	-0.059604\\
2.6977	-0.0595459\\
2.7003	-0.0594858\\
2.7029	-0.059426\\
2.7054	-0.0593687\\
2.708	-0.0593095\\
2.7105	-0.0592529\\
2.7131	-0.0591942\\
2.7157	-0.0591359\\
2.7182	-0.0590802\\
2.7208	-0.0590225\\
2.7234	-0.0589651\\
2.7259	-0.0589102\\
2.7285	-0.0588534\\
2.7311	-0.058797\\
2.7336	-0.058743\\
2.7362	-0.0586872\\
2.7387	-0.0586339\\
2.7413	-0.0585787\\
2.7439	-0.0585239\\
2.7464	-0.0584715\\
2.749	-0.0584173\\
2.7516	-0.0583635\\
2.7541	-0.058312\\
2.7567	-0.0582588\\
2.7593	-0.058206\\
2.7618	-0.0581555\\
2.7644	-0.0581034\\
2.7669	-0.0580536\\
2.7695	-0.0580021\\
2.7721	-0.057951\\
2.7746	-0.0579022\\
2.7772	-0.0578518\\
2.7798	-0.0578017\\
2.7823	-0.0577539\\
2.7849	-0.0577046\\
2.7875	-0.0576556\\
2.79	-0.0576088\\
2.7926	-0.0575605\\
2.7951	-0.0575145\\
2.7977	-0.0574669\\
2.8003	-0.0574197\\
2.8028	-0.0573747\\
2.8054	-0.0573283\\
2.808	-0.0572822\\
2.8105	-0.0572383\\
2.8131	-0.0571929\\
2.8157	-0.057148\\
2.8182	-0.0571051\\
2.8208	-0.057061\\
2.8233	-0.0570188\\
2.8259	-0.0569754\\
2.8285	-0.0569324\\
2.831	-0.0568913\\
2.8336	-0.0568491\\
2.8362	-0.0568072\\
2.8387	-0.0567673\\
2.8413	-0.0567262\\
2.8439	-0.0566855\\
2.8464	-0.0566467\\
2.849	-0.0566068\\
2.8516	-0.0565673\\
2.8541	-0.0565297\\
2.8567	-0.0564909\\
2.8592	-0.0564541\\
2.8618	-0.0564162\\
2.8644	-0.0563786\\
2.8669	-0.056343\\
2.8695	-0.0563062\\
2.8721	-0.0562699\\
2.8746	-0.0562354\\
2.8772	-0.0561999\\
2.8798	-0.0561648\\
2.8823	-0.0561315\\
2.8849	-0.0560972\\
2.8874	-0.0560647\\
2.89	-0.0560313\\
2.8926	-0.0559982\\
2.8951	-0.0559669\\
2.8977	-0.0559347\\
2.9003	-0.0559029\\
2.9028	-0.0558727\\
2.9054	-0.0558418\\
2.908	-0.0558113\\
2.9105	-0.0557823\\
2.9131	-0.0557526\\
2.9156	-0.0557245\\
2.9182	-0.0556957\\
2.9208	-0.0556672\\
2.9233	-0.0556403\\
2.9259	-0.0556127\\
2.9285	-0.0555856\\
2.931	-0.0555599\\
2.9336	-0.0555335\\
2.9362	-0.0555077\\
2.9387	-0.0554832\\
2.9413	-0.0554581\\
2.9438	-0.0554345\\
2.9464	-0.0554103\\
2.949	-0.0553865\\
2.9515	-0.0553641\\
2.9541	-0.0553411\\
2.9567	-0.0553187\\
2.9592	-0.0552974\\
2.9618	-0.0552758\\
2.9644	-0.0552546\\
2.9669	-0.0552346\\
2.9695	-0.0552142\\
2.972	-0.0551951\\
2.9746	-0.0551756\\
2.9772	-0.0551565\\
2.9797	-0.0551385\\
2.9823	-0.0551203\\
2.9849	-0.0551025\\
2.9874	-0.0550857\\
2.99	-0.0550688\\
2.9926	-0.0550522\\
2.9951	-0.0550367\\
2.9977	-0.055021\\
3.0003	-0.0550058\\
3.0028	-0.0549915\\
3.0054	-0.054977\\
3.0079	-0.0549636\\
3.0105	-0.0549499\\
3.0131	-0.0549368\\
3.0156	-0.0549245\\
3.0182	-0.0549121\\
3.0208	-0.0549002\\
3.0233	-0.0548891\\
3.0259	-0.054878\\
3.0285	-0.0548674\\
3.031	-0.0548575\\
3.0336	-0.0548476\\
3.0361	-0.0548385\\
3.0387	-0.0548295\\
3.0413	-0.0548209\\
3.0438	-0.0548129\\
3.0464	-0.0548051\\
3.049	-0.0547977\\
3.0515	-0.054791\\
3.0541	-0.0547844\\
3.0567	-0.0547781\\
3.0592	-0.0547726\\
3.0618	-0.0547672\\
3.0643	-0.0547623\\
3.0669	-0.0547577\\
3.0695	-0.0547535\\
3.072	-0.0547498\\
3.0746	-0.0547464\\
3.0772	-0.0547433\\
3.0797	-0.0547408\\
3.0823	-0.0547385\\
3.0849	-0.0547366\\
3.0874	-0.0547352\\
3.09	-0.0547341\\
3.0925	-0.0547333\\
3.0951	-0.054733\\
3.0977	-0.054733\\
3.1002	-0.0547334\\
3.1028	-0.0547341\\
3.1054	-0.0547353\\
3.1079	-0.0547367\\
3.1105	-0.0547386\\
3.1131	-0.0547408\\
3.1156	-0.0547433\\
3.1182	-0.0547463\\
3.1207	-0.0547495\\
3.1233	-0.0547531\\
3.1259	-0.0547572\\
3.1284	-0.0547614\\
3.131	-0.0547661\\
3.1336	-0.0547712\\
3.1361	-0.0547764\\
3.1387	-0.0547822\\
3.1413	-0.0547883\\
3.1438	-0.0547945\\
3.1464	-0.0548013\\
3.1489	-0.0548082\\
3.1515	-0.0548156\\
3.1541	-0.0548234\\
3.1566	-0.0548312\\
3.1592	-0.0548397\\
3.1618	-0.0548484\\
3.1643	-0.0548572\\
3.1669	-0.0548666\\
3.1695	-0.0548763\\
3.172	-0.054886\\
3.1746	-0.0548963\\
3.1772	-0.054907\\
3.1797	-0.0549175\\
3.1823	-0.0549288\\
3.1848	-0.0549399\\
3.1874	-0.0549517\\
3.19	-0.0549639\\
3.1925	-0.0549758\\
3.1951	-0.0549886\\
3.1977	-0.0550016\\
3.2002	-0.0550144\\
3.2028	-0.055028\\
3.2054	-0.0550418\\
3.2079	-0.0550554\\
3.2105	-0.0550698\\
3.213	-0.0550839\\
3.2156	-0.0550989\\
3.2182	-0.0551141\\
3.2207	-0.0551289\\
3.2233	-0.0551446\\
3.2259	-0.0551606\\
3.2284	-0.0551762\\
3.231	-0.0551927\\
3.2336	-0.0552094\\
3.2361	-0.0552257\\
3.2387	-0.0552429\\
3.2412	-0.0552597\\
3.2438	-0.0552774\\
3.2464	-0.0552953\\
3.2489	-0.0553127\\
3.2515	-0.0553311\\
3.2541	-0.0553497\\
3.2566	-0.0553677\\
3.2592	-0.0553867\\
3.2618	-0.055406\\
3.2643	-0.0554247\\
3.2669	-0.0554443\\
3.2694	-0.0554634\\
3.272	-0.0554834\\
3.2746	-0.0555036\\
3.2771	-0.0555233\\
3.2797	-0.0555439\\
3.2823	-0.0555647\\
3.2848	-0.0555849\\
3.2874	-0.0556061\\
3.29	-0.0556274\\
3.2925	-0.0556481\\
3.2951	-0.0556698\\
3.2976	-0.0556908\\
3.3002	-0.0557128\\
3.3028	-0.055735\\
3.3053	-0.0557565\\
3.3079	-0.055779\\
3.3105	-0.0558016\\
3.313	-0.0558235\\
3.3156	-0.0558465\\
3.3182	-0.0558696\\
3.3207	-0.0558919\\
3.3233	-0.0559152\\
3.3259	-0.0559387\\
3.3284	-0.0559614\\
3.331	-0.0559852\\
3.3335	-0.0560081\\
3.3361	-0.0560321\\
3.3387	-0.0560562\\
3.3412	-0.0560794\\
3.3438	-0.0561038\\
3.3464	-0.0561282\\
3.3489	-0.0561518\\
3.3515	-0.0561764\\
3.3541	-0.0562011\\
3.3566	-0.056225\\
3.3592	-0.0562499\\
3.3617	-0.0562739\\
3.3643	-0.056299\\
3.3669	-0.0563241\\
3.3694	-0.0563484\\
3.372	-0.0563737\\
3.3746	-0.056399\\
3.3771	-0.0564235\\
3.3797	-0.056449\\
3.3823	-0.0564745\\
3.3848	-0.0564992\\
3.3874	-0.0565249\\
3.3899	-0.0565496\\
3.3925	-0.0565754\\
3.3951	-0.0566012\\
3.3976	-0.056626\\
3.4002	-0.0566519\\
3.4028	-0.0566779\\
3.4053	-0.0567028\\
3.4079	-0.0567288\\
3.4105	-0.0567549\\
3.413	-0.0567799\\
3.4156	-0.056806\\
3.4181	-0.056831\\
3.4207	-0.0568571\\
3.4233	-0.0568832\\
3.4258	-0.0569083\\
3.4284	-0.0569345\\
3.431	-0.0569606\\
3.4335	-0.0569857\\
3.4361	-0.0570118\\
3.4387	-0.0570379\\
3.4412	-0.057063\\
3.4438	-0.0570891\\
3.4463	-0.0571141\\
3.4489	-0.0571402\\
3.4515	-0.0571662\\
3.454	-0.0571912\\
3.4566	-0.0572172\\
3.4592	-0.0572431\\
3.4617	-0.0572681\\
3.4643	-0.0572939\\
3.4669	-0.0573198\\
3.4694	-0.0573446\\
3.472	-0.0573703\\
3.4745	-0.057395\\
3.4771	-0.0574207\\
3.4797	-0.0574463\\
3.4822	-0.0574708\\
3.4848	-0.0574963\\
3.4874	-0.0575218\\
3.4899	-0.0575461\\
3.4925	-0.0575714\\
3.4951	-0.0575967\\
3.4976	-0.0576208\\
3.5002	-0.0576459\\
3.5028	-0.0576709\\
3.5053	-0.0576949\\
3.5079	-0.0577197\\
3.5104	-0.0577436\\
3.513	-0.0577682\\
3.5156	-0.0577928\\
3.5181	-0.0578164\\
3.5207	-0.0578408\\
3.5233	-0.0578651\\
3.5258	-0.0578883\\
3.5284	-0.0579124\\
3.531	-0.0579364\\
3.5335	-0.0579594\\
3.5361	-0.0579832\\
3.5386	-0.058006\\
3.5412	-0.0580295\\
3.5438	-0.058053\\
3.5463	-0.0580754\\
3.5489	-0.0580986\\
3.5515	-0.0581217\\
3.554	-0.0581438\\
3.5566	-0.0581667\\
3.5592	-0.0581894\\
3.5617	-0.0582111\\
3.5643	-0.0582335\\
3.5668	-0.058255\\
3.5694	-0.0582772\\
3.572	-0.0582992\\
3.5745	-0.0583203\\
3.5771	-0.058342\\
3.5797	-0.0583637\\
3.5822	-0.0583843\\
3.5848	-0.0584056\\
3.5874	-0.0584268\\
3.5899	-0.058447\\
3.5925	-0.0584679\\
3.595	-0.0584878\\
3.5976	-0.0585084\\
3.6002	-0.0585288\\
3.6027	-0.0585482\\
3.6053	-0.0585683\\
3.6079	-0.0585882\\
3.6104	-0.0586072\\
3.613	-0.0586268\\
3.6156	-0.0586462\\
3.6181	-0.0586647\\
3.6207	-0.0586838\\
3.6232	-0.058702\\
3.6258	-0.0587207\\
3.6284	-0.0587392\\
3.6309	-0.0587569\\
3.6335	-0.0587751\\
3.6361	-0.0587931\\
3.6386	-0.0588102\\
3.6412	-0.0588279\\
3.6438	-0.0588454\\
3.6463	-0.058862\\
3.6489	-0.058879\\
3.6515	-0.0588959\\
3.654	-0.058912\\
3.6566	-0.0589285\\
3.6591	-0.0589442\\
3.6617	-0.0589604\\
3.6643	-0.0589763\\
3.6668	-0.0589915\\
3.6694	-0.059007\\
3.672	-0.0590223\\
3.6745	-0.0590369\\
3.6771	-0.0590519\\
3.6797	-0.0590666\\
3.6822	-0.0590806\\
3.6848	-0.059095\\
3.6873	-0.0591086\\
3.6899	-0.0591225\\
3.6925	-0.0591362\\
3.695	-0.0591492\\
3.6976	-0.0591626\\
3.7002	-0.0591757\\
3.7027	-0.0591881\\
3.7053	-0.0592008\\
3.7079	-0.0592133\\
3.7104	-0.0592251\\
3.713	-0.0592372\\
3.7155	-0.0592486\\
3.7181	-0.0592602\\
3.7207	-0.0592716\\
3.7232	-0.0592824\\
3.7258	-0.0592934\\
3.7284	-0.0593042\\
3.7309	-0.0593144\\
3.7335	-0.0593247\\
3.7361	-0.0593349\\
3.7386	-0.0593444\\
3.7412	-0.0593541\\
3.7437	-0.0593632\\
3.7463	-0.0593725\\
3.7489	-0.0593816\\
3.7514	-0.0593901\\
3.754	-0.0593987\\
3.7566	-0.0594071\\
3.7591	-0.0594149\\
3.7617	-0.0594229\\
3.7643	-0.0594306\\
3.7668	-0.0594379\\
3.7694	-0.0594452\\
3.7719	-0.059452\\
3.7745	-0.0594589\\
3.7771	-0.0594655\\
3.7796	-0.0594717\\
3.7822	-0.0594779\\
3.7848	-0.0594839\\
3.7873	-0.0594894\\
3.7899	-0.0594949\\
3.7925	-0.0595003\\
3.795	-0.0595052\\
3.7976	-0.05951\\
3.8002	-0.0595147\\
3.8027	-0.059519\\
3.8053	-0.0595232\\
3.8078	-0.0595271\\
3.8104	-0.0595308\\
3.813	-0.0595344\\
3.8155	-0.0595376\\
3.8181	-0.0595408\\
3.8207	-0.0595437\\
3.8232	-0.0595463\\
3.8258	-0.0595487\\
3.8284	-0.059551\\
3.8309	-0.0595529\\
3.8335	-0.0595548\\
3.836	-0.0595563\\
3.8386	-0.0595577\\
3.8412	-0.0595589\\
3.8437	-0.0595598\\
3.8463	-0.0595606\\
3.8489	-0.0595611\\
3.8514	-0.0595614\\
3.854	-0.0595615\\
3.8566	-0.0595614\\
3.8591	-0.0595611\\
3.8617	-0.0595606\\
3.8642	-0.0595599\\
3.8668	-0.059559\\
3.8694	-0.0595578\\
3.8719	-0.0595565\\
3.8745	-0.059555\\
3.8771	-0.0595532\\
3.8796	-0.0595513\\
3.8822	-0.0595491\\
3.8848	-0.0595467\\
3.8873	-0.0595442\\
3.8899	-0.0595414\\
3.8924	-0.0595385\\
3.895	-0.0595353\\
3.8976	-0.0595319\\
3.9001	-0.0595285\\
3.9027	-0.0595247\\
3.9053	-0.0595207\\
3.9078	-0.0595167\\
3.9104	-0.0595123\\
3.913	-0.0595077\\
3.9155	-0.0595031\\
3.9181	-0.0594981\\
3.9206	-0.0594932\\
3.9232	-0.0594878\\
3.9258	-0.0594823\\
3.9283	-0.0594768\\
3.9309	-0.0594709\\
3.9335	-0.0594648\\
3.936	-0.0594587\\
3.9386	-0.0594522\\
3.9412	-0.0594456\\
3.9437	-0.059439\\
3.9463	-0.059432\\
3.9488	-0.059425\\
3.9514	-0.0594176\\
3.954	-0.0594101\\
3.9565	-0.0594026\\
3.9591	-0.0593947\\
3.9617	-0.0593866\\
3.9642	-0.0593787\\
3.9668	-0.0593703\\
3.9694	-0.0593617\\
3.9719	-0.0593532\\
3.9745	-0.0593443\\
3.9771	-0.0593352\\
3.9796	-0.0593262\\
3.9822	-0.0593168\\
3.9847	-0.0593076\\
3.9873	-0.0592978\\
3.9899	-0.0592879\\
3.9924	-0.0592782\\
3.995	-0.0592679\\
3.9976	-0.0592575\\
4.0001	-0.0592474\\
4.0027	-0.0592367\\
4.0053	-0.0592258\\
4.0078	-0.0592152\\
4.0104	-0.0592041\\
4.0129	-0.0591932\\
4.0155	-0.0591817\\
4.0181	-0.0591701\\
4.0206	-0.0591588\\
4.0232	-0.0591469\\
4.0258	-0.0591349\\
4.0283	-0.0591232\\
4.0309	-0.0591109\\
4.0335	-0.0590984\\
4.036	-0.0590864\\
4.0386	-0.0590736\\
4.0411	-0.0590613\\
4.0437	-0.0590483\\
4.0463	-0.0590352\\
4.0488	-0.0590225\\
4.0514	-0.0590091\\
4.054	-0.0589956\\
4.0565	-0.0589825\\
4.0591	-0.0589688\\
4.0617	-0.058955\\
4.0642	-0.0589415\\
4.0668	-0.0589274\\
4.0693	-0.0589138\\
4.0719	-0.0588995\\
4.0745	-0.058885\\
4.077	-0.0588711\\
4.0796	-0.0588564\\
4.0822	-0.0588417\\
4.0847	-0.0588274\\
4.0873	-0.0588124\\
4.0899	-0.0587973\\
4.0924	-0.0587827\\
4.095	-0.0587675\\
4.0975	-0.0587527\\
4.1001	-0.0587372\\
4.1027	-0.0587217\\
4.1052	-0.0587066\\
4.1078	-0.0586909\\
4.1104	-0.058675\\
4.1129	-0.0586597\\
4.1155	-0.0586437\\
4.1181	-0.0586276\\
4.1206	-0.058612\\
4.1232	-0.0585958\\
4.1258	-0.0585794\\
4.1283	-0.0585636\\
4.1309	-0.0585471\\
4.1334	-0.0585312\\
4.136	-0.0585145\\
4.1386	-0.0584978\\
4.1411	-0.0584817\\
4.1437	-0.0584648\\
4.1463	-0.0584479\\
4.1488	-0.0584315\\
4.1514	-0.0584145\\
4.154	-0.0583973\\
4.1565	-0.0583808\\
4.1591	-0.0583636\\
4.1616	-0.0583469\\
4.1642	-0.0583296\\
4.1668	-0.0583121\\
4.1693	-0.0582954\\
4.1719	-0.0582778\\
4.1745	-0.0582603\\
4.177	-0.0582433\\
4.1796	-0.0582257\\
4.1822	-0.0582079\\
4.1847	-0.0581909\\
4.1873	-0.0581731\\
4.1898	-0.0581559\\
4.1924	-0.0581381\\
4.195	-0.0581201\\
4.1975	-0.0581029\\
4.2001	-0.0580849\\
4.2027	-0.0580669\\
4.2052	-0.0580496\\
4.2078	-0.0580315\\
4.2104	-0.0580134\\
4.2129	-0.0579959\\
4.2155	-0.0579778\\
4.218	-0.0579603\\
4.2206	-0.0579421\\
4.2232	-0.0579239\\
4.2257	-0.0579064\\
4.2283	-0.0578881\\
4.2309	-0.0578699\\
4.2334	-0.0578523\\
4.236	-0.057834\\
4.2386	-0.0578157\\
4.2411	-0.0577981\\
4.2437	-0.0577798\\
4.2462	-0.0577622\\
4.2488	-0.0577438\\
4.2514	-0.0577255\\
4.2539	-0.0577079\\
4.2565	-0.0576895\\
4.2591	-0.0576712\\
4.2616	-0.0576536\\
4.2642	-0.0576353\\
4.2668	-0.057617\\
4.2693	-0.0575994\\
4.2719	-0.057581\\
4.2744	-0.0575635\\
4.277	-0.0575452\\
4.2796	-0.0575269\\
4.2821	-0.0575094\\
4.2847	-0.0574911\\
4.2873	-0.0574729\\
4.2898	-0.0574554\\
4.2924	-0.0574372\\
4.295	-0.0574191\\
4.2975	-0.0574016\\
4.3001	-0.0573835\\
4.3027	-0.0573655\\
4.3052	-0.0573481\\
4.3078	-0.0573301\\
4.3103	-0.0573128\\
4.3129	-0.0572948\\
4.3155	-0.0572769\\
4.318	-0.0572596\\
4.3206	-0.0572418\\
4.3232	-0.057224\\
4.3257	-0.0572068\\
4.3283	-0.0571891\\
4.3309	-0.0571714\\
4.3334	-0.0571544\\
4.336	-0.0571368\\
4.3385	-0.0571199\\
4.3411	-0.0571023\\
4.3437	-0.0570848\\
4.3462	-0.0570681\\
4.3488	-0.0570507\\
4.3514	-0.0570333\\
4.3539	-0.0570167\\
4.3565	-0.0569995\\
4.3591	-0.0569823\\
4.3616	-0.0569658\\
4.3642	-0.0569487\\
4.3667	-0.0569324\\
4.3693	-0.0569154\\
4.3719	-0.0568985\\
4.3744	-0.0568823\\
4.377	-0.0568655\\
4.3796	-0.0568488\\
4.3821	-0.0568328\\
4.3847	-0.0568162\\
4.3873	-0.0567997\\
4.3898	-0.0567839\\
4.3924	-0.0567675\\
4.3949	-0.0567518\\
4.3975	-0.0567356\\
4.4001	-0.0567194\\
4.4026	-0.056704\\
4.4052	-0.0566879\\
4.4078	-0.056672\\
4.4103	-0.0566567\\
4.4129	-0.0566409\\
4.4155	-0.0566252\\
4.418	-0.0566102\\
4.4206	-0.0565946\\
4.4231	-0.0565797\\
4.4257	-0.0565643\\
4.4283	-0.056549\\
4.4308	-0.0565344\\
4.4334	-0.0565192\\
4.436	-0.0565041\\
4.4385	-0.0564897\\
4.4411	-0.0564749\\
4.4437	-0.05646\\
4.4462	-0.0564459\\
4.4488	-0.0564313\\
4.4514	-0.0564167\\
4.4539	-0.0564028\\
4.4565	-0.0563885\\
4.459	-0.0563748\\
4.4616	-0.0563606\\
4.4642	-0.0563465\\
4.4667	-0.0563331\\
4.4693	-0.0563192\\
4.4719	-0.0563054\\
4.4744	-0.0562922\\
4.477	-0.0562787\\
4.4796	-0.0562652\\
4.4821	-0.0562523\\
4.4847	-0.056239\\
4.4872	-0.0562263\\
4.4898	-0.0562132\\
4.4924	-0.0562002\\
4.4949	-0.0561878\\
4.4975	-0.056175\\
4.5001	-0.0561623\\
4.5026	-0.0561502\\
4.5052	-0.0561377\\
4.5078	-0.0561254\\
4.5103	-0.0561136\\
4.5129	-0.0561014\\
4.5154	-0.0560898\\
4.518	-0.0560779\\
4.5206	-0.056066\\
4.5231	-0.0560547\\
4.5257	-0.0560431\\
4.5283	-0.0560316\\
4.5308	-0.0560206\\
4.5334	-0.0560094\\
4.536	-0.0559982\\
4.5385	-0.0559875\\
4.5411	-0.0559766\\
4.5436	-0.0559661\\
4.5462	-0.0559554\\
4.5488	-0.0559448\\
4.5513	-0.0559347\\
4.5539	-0.0559243\\
4.5565	-0.055914\\
4.559	-0.0559042\\
4.5616	-0.0558941\\
4.5642	-0.0558842\\
4.5667	-0.0558747\\
4.5693	-0.055865\\
4.5718	-0.0558558\\
4.5744	-0.0558463\\
4.577	-0.055837\\
4.5795	-0.0558281\\
4.5821	-0.0558189\\
4.5847	-0.0558099\\
4.5872	-0.0558014\\
4.5898	-0.0557926\\
4.5924	-0.0557839\\
4.5949	-0.0557757\\
4.5975	-0.0557673\\
4.6001	-0.055759\\
4.6026	-0.0557511\\
4.6052	-0.055743\\
4.6077	-0.0557354\\
4.6103	-0.0557275\\
4.6129	-0.0557198\\
4.6154	-0.0557125\\
4.618	-0.055705\\
4.6206	-0.0556977\\
4.6231	-0.0556907\\
4.6257	-0.0556836\\
4.6283	-0.0556766\\
4.6308	-0.05567\\
4.6334	-0.0556632\\
4.6359	-0.0556568\\
4.6385	-0.0556503\\
4.6411	-0.0556438\\
4.6436	-0.0556378\\
4.6462	-0.0556316\\
4.6488	-0.0556256\\
4.6513	-0.0556199\\
4.6539	-0.055614\\
4.6565	-0.0556083\\
4.659	-0.055603\\
4.6616	-0.0555975\\
4.6641	-0.0555924\\
4.6667	-0.0555872\\
4.6693	-0.055582\\
4.6718	-0.0555773\\
4.6744	-0.0555724\\
4.677	-0.0555676\\
4.6795	-0.0555632\\
4.6821	-0.0555586\\
4.6847	-0.0555542\\
4.6872	-0.0555501\\
4.6898	-0.055546\\
4.6923	-0.0555421\\
4.6949	-0.0555381\\
4.6975	-0.0555343\\
4.7	-0.0555307\\
4.7026	-0.0555271\\
4.7052	-0.0555237\\
4.7077	-0.0555204\\
4.7103	-0.0555172\\
4.7129	-0.0555141\\
4.7154	-0.0555112\\
4.718	-0.0555083\\
4.7205	-0.0555056\\
4.7231	-0.0555029\\
4.7257	-0.0555003\\
4.7282	-0.055498\\
4.7308	-0.0554957\\
4.7334	-0.0554934\\
4.7359	-0.0554914\\
4.7385	-0.0554894\\
4.7411	-0.0554875\\
4.7436	-0.0554858\\
4.7462	-0.0554841\\
4.7487	-0.0554826\\
4.7513	-0.0554811\\
4.7539	-0.0554798\\
4.7564	-0.0554786\\
4.759	-0.0554775\\
4.7616	-0.0554764\\
4.7641	-0.0554756\\
4.7667	-0.0554747\\
4.7693	-0.055474\\
4.7718	-0.0554735\\
4.7744	-0.055473\\
4.777	-0.0554726\\
4.7795	-0.0554723\\
4.7821	-0.0554721\\
4.7846	-0.055472\\
4.7872	-0.055472\\
4.7898	-0.0554722\\
4.7923	-0.0554724\\
4.7949	-0.0554727\\
4.7975	-0.0554731\\
4.8	-0.0554736\\
4.8026	-0.0554742\\
4.8052	-0.055475\\
4.8077	-0.0554757\\
4.8103	-0.0554767\\
4.8128	-0.0554776\\
4.8154	-0.0554787\\
4.818	-0.0554799\\
4.8205	-0.0554812\\
4.8231	-0.0554826\\
4.8257	-0.0554841\\
4.8282	-0.0554856\\
4.8308	-0.0554872\\
4.8334	-0.055489\\
4.8359	-0.0554908\\
4.8385	-0.0554927\\
4.841	-0.0554947\\
4.8436	-0.0554968\\
4.8462	-0.055499\\
4.8487	-0.0555012\\
4.8513	-0.0555036\\
4.8539	-0.055506\\
4.8564	-0.0555085\\
4.859	-0.0555111\\
4.8616	-0.0555138\\
4.8641	-0.0555165\\
4.8667	-0.0555194\\
4.8692	-0.0555223\\
4.8718	-0.0555253\\
4.8744	-0.0555284\\
4.8769	-0.0555315\\
4.8795	-0.0555348\\
4.8821	-0.0555382\\
4.8846	-0.0555415\\
4.8872	-0.055545\\
4.8898	-0.0555486\\
4.8923	-0.0555521\\
4.8949	-0.0555559\\
4.8974	-0.0555595\\
4.9	-0.0555634\\
4.9026	-0.0555674\\
4.9051	-0.0555713\\
4.9077	-0.0555754\\
4.9103	-0.0555795\\
4.9128	-0.0555836\\
4.9154	-0.0555879\\
4.918	-0.0555923\\
4.9205	-0.0555966\\
4.9231	-0.0556011\\
4.9257	-0.0556057\\
4.9282	-0.0556101\\
4.9308	-0.0556148\\
4.9333	-0.0556194\\
4.9359	-0.0556242\\
4.9385	-0.0556291\\
4.941	-0.0556339\\
4.9436	-0.0556389\\
4.9462	-0.0556439\\
4.9487	-0.0556489\\
4.9513	-0.055654\\
4.9539	-0.0556593\\
4.9564	-0.0556643\\
4.959	-0.0556697\\
4.9615	-0.0556749\\
4.9641	-0.0556803\\
4.9667	-0.0556858\\
4.9692	-0.0556912\\
4.9718	-0.0556967\\
4.9744	-0.0557024\\
4.9769	-0.0557079\\
4.9795	-0.0557136\\
4.9821	-0.0557194\\
4.9846	-0.055725\\
4.9872	-0.0557309\\
4.9897	-0.0557366\\
4.9923	-0.0557426\\
4.9949	-0.0557486\\
4.9974	-0.0557544\\
5	-0.0557605\\
};
\addplot [color=black, dotted, line width=0.6pt, forget plot]
  table[row sep=crcr]{%
-0.125	-0.0822382\\
5	-0.0822382\\
};
\draw[line width=\figlegendlw, fill=\figlegendfill, draw=\figlegenddraw] (axis cs:2.90186,-0.0234737) rectangle (axis cs:4.8975,0.094);
\addplot [color=black, line width=1.4pt, forget plot]
  table[row sep=crcr]{%
3.00164	0.0705053\\
3.64024	0.0705053\\
};
\node[right, align=left, inner sep=0]
at (axis cs:3.7,0.071) {$f_{\mathrm{COI}}$};
\addplot [color=mycolor17, line width=1.6pt, forget plot]
  table[row sep=crcr]{%
3.00164	0.0352632\\
3.05969	0.0352632\\
};
\addplot [color=mycolor22, line width=1.6pt, forget plot]
  table[row sep=crcr]{%
3.05969	0.0352632\\
3.11775	0.0352632\\
};
\addplot [color=mycolor18, line width=1.6pt, forget plot]
  table[row sep=crcr]{%
3.11775	0.0352632\\
3.1758	0.0352632\\
};
\addplot [color=mycolor19, line width=1.6pt, forget plot]
  table[row sep=crcr]{%
3.1758	0.0352632\\
3.23386	0.0352632\\
};
\addplot [color=mycolor15, line width=1.6pt, forget plot]
  table[row sep=crcr]{%
3.23386	0.0352632\\
3.29191	0.0352632\\
};
\addplot [color=mycolor14, line width=1.6pt, forget plot]
  table[row sep=crcr]{%
3.29191	0.0352632\\
3.34997	0.0352632\\
};
\addplot [color=mycolor23, line width=1.6pt, forget plot]
  table[row sep=crcr]{%
3.34997	0.0352632\\
3.40802	0.0352632\\
};
\addplot [color=mycolor24, line width=1.6pt, forget plot]
  table[row sep=crcr]{%
3.40802	0.0352632\\
3.46608	0.0352632\\
};
\addplot [color=mycolor20, line width=1.6pt, forget plot]
  table[row sep=crcr]{%
3.46608	0.0352632\\
3.52413	0.0352632\\
};
\addplot [color=mycolor21, line width=1.6pt, forget plot]
  table[row sep=crcr]{%
3.52413	0.0352632\\
3.58219	0.0352632\\
};
\addplot [color=mycolor16, line width=1.6pt, forget plot]
  table[row sep=crcr]{%
3.58219	0.0352632\\
3.64024	0.0352632\\
};
\node[right, align=left, inner sep=0]
at (axis cs:3.7,0.035) {$f_i$};
\addplot [color=black, dotted, line width=0.6pt, forget plot]
  table[row sep=crcr]{%
3.00164	2.10526e-05\\
3.64024	2.10526e-05\\
};
\node[right, align=left, inner sep=0]
at (axis cs:3.7,0) {$N_{\mathrm{COI}}^{20}$};
\addplot [color=black, dotted, line width=0.6pt, forget plot]
  table[row sep=crcr]{%
-0.125	-0.0849146\\
5	-0.0849146\\
};
\end{axis}
\end{tikzpicture}%
  \caption{Fig.~5(b) --- per-node deviation, nonlinear simulation, $K=0$ $|$ $K=20\,I$.}
\end{figure}

\section{Large-system activation sweep (WECC, SC500, Texas2000)}

\begin{figure}[!t]\centering
  % This file was created by matlab2tikz.
%
\definecolor{mycolor1}{rgb}{0.70000,0.70000,0.75000}%
\definecolor{mycolor2}{rgb}{0.57647,0.00000,0.00000}%
\definecolor{mycolor3}{rgb}{0.46620,0.67454,0.77804}%
\definecolor{mycolor4}{rgb}{0.39131,0.62716,0.74576}%
\definecolor{mycolor5}{rgb}{0.30499,0.57091,0.70745}%
\definecolor{mycolor6}{rgb}{0.22212,0.51426,0.66892}%
\definecolor{mycolor7}{rgb}{0.14661,0.45864,0.63129}%
\definecolor{mycolor8}{rgb}{0.09143,0.41309,0.60075}%
\definecolor{mycolor9}{rgb}{0.04223,0.36317,0.56781}%
\definecolor{mycolor10}{rgb}{0.01771,0.31647,0.53748}%
\definecolor{mycolor11}{rgb}{0.01172,0.27830,0.51292}%
\definecolor{mycolor12}{rgb}{0.00988,0.23521,0.48525}%
\definecolor{mycolor13}{rgb}{0.07059,0.07059,0.07059}%
\definecolor{mycolor14}{rgb}{0.00000,0.32157,0.57647}%
%
\begin{tikzpicture}[font=\figfont]

\begin{axis}[%
width={\figscale*0.747in},
height={\figscale*0.745in},
at={(\figscale*0.455in,\figscale*0.596in)},
scale only axis,
separate axis lines,
every outer x axis line/.append style={black},
every x tick label/.append style={font=\figticfont},
every x tick/.append style={black},
xmin=-15,
xmax=1,
xlabel={$\mathrm{Re}(\lambda)$ in rad/s},
every outer y axis line/.append style={black},
every y tick label/.append style={font=\figticfont},
every y tick/.append style={black},
ymin=-69.8453,
ymax=69.8453,
ylabel={$\mathrm{Im}(\lambda)$ in rad/s},
axis background/.style={fill=white},
title style={font=\figcornerfont},
title={WECC},
xmajorgrids,
ymajorgrids,
grid style={white!55!black, opacity=0.3}
]
\addplot [color=mycolor1, dashed, line width=0.6pt, forget plot]
  table[row sep=crcr]{%
-15	149.248\\
-14.6939	146.202\\
-14.3878	143.156\\
-14.0816	140.11\\
-13.7755	137.065\\
-13.4694	134.019\\
-13.1633	130.973\\
-12.8571	127.927\\
-12.551	124.881\\
-12.2449	121.835\\
-11.9388	118.789\\
-11.6327	115.743\\
-11.3265	112.698\\
-11.0204	109.652\\
-10.7143	106.606\\
-10.4082	103.56\\
-10.102	100.514\\
-9.79592	97.4682\\
-9.4898	94.4223\\
-9.18367	91.3764\\
-8.87755	88.3305\\
-8.57143	85.2846\\
-8.26531	82.2388\\
-7.95918	79.1929\\
-7.65306	76.147\\
-7.34694	73.1011\\
-7.04082	70.0552\\
-6.73469	67.0094\\
-6.42857	63.9635\\
-6.12245	60.9176\\
-5.81633	57.8717\\
-5.5102	54.8258\\
-5.20408	51.78\\
-4.89796	48.7341\\
-4.59184	45.6882\\
-4.28571	42.6423\\
-3.97959	39.5964\\
-3.67347	36.5506\\
-3.36735	33.5047\\
-3.06122	30.4588\\
-2.7551	27.4129\\
-2.44898	24.367\\
-2.14286	21.3212\\
-1.83673	18.2753\\
-1.53061	15.2294\\
-1.22449	12.1835\\
-0.918367	9.13764\\
-0.612245	6.09176\\
-0.306122	3.04588\\
0	-0\\
};
\addplot [color=mycolor1, dashed, line width=0.6pt, forget plot]
  table[row sep=crcr]{%
-15	-149.248\\
-14.6939	-146.202\\
-14.3878	-143.156\\
-14.0816	-140.11\\
-13.7755	-137.065\\
-13.4694	-134.019\\
-13.1633	-130.973\\
-12.8571	-127.927\\
-12.551	-124.881\\
-12.2449	-121.835\\
-11.9388	-118.789\\
-11.6327	-115.743\\
-11.3265	-112.698\\
-11.0204	-109.652\\
-10.7143	-106.606\\
-10.4082	-103.56\\
-10.102	-100.514\\
-9.79592	-97.4682\\
-9.4898	-94.4223\\
-9.18367	-91.3764\\
-8.87755	-88.3305\\
-8.57143	-85.2846\\
-8.26531	-82.2388\\
-7.95918	-79.1929\\
-7.65306	-76.147\\
-7.34694	-73.1011\\
-7.04082	-70.0552\\
-6.73469	-67.0094\\
-6.42857	-63.9635\\
-6.12245	-60.9176\\
-5.81633	-57.8717\\
-5.5102	-54.8258\\
-5.20408	-51.78\\
-4.89796	-48.7341\\
-4.59184	-45.6882\\
-4.28571	-42.6423\\
-3.97959	-39.5964\\
-3.67347	-36.5506\\
-3.36735	-33.5047\\
-3.06122	-30.4588\\
-2.7551	-27.4129\\
-2.44898	-24.367\\
-2.14286	-21.3212\\
-1.83673	-18.2753\\
-1.53061	-15.2294\\
-1.22449	-12.1835\\
-0.918367	-9.13764\\
-0.612245	-6.09176\\
-0.306122	-3.04588\\
0	0\\
};
\addplot [color=mycolor1, dashed, line width=0.6pt, forget plot]
  table[row sep=crcr]{%
-3	-69.8453\\
-3	69.8453\\
};
\addplot [color=white!88!black, line width=0.6pt, forget plot]
  table[row sep=crcr]{%
-15	0\\
1	0\\
};
\addplot [color=mycolor2, line width=0.9pt, only marks, mark size=2.2pt, mark=x, mark options={solid, mycolor2}, forget plot]
  table[row sep=crcr]{%
-14.1762	60.7334\\
-14.1762	-60.7334\\
-10.6557	55.4763\\
-10.6557	-55.4763\\
-10.0841	53.1442\\
-10.0841	-53.1442\\
-9.27284	51.2572\\
-9.27284	-51.2572\\
-9.64256	50.6892\\
-9.64256	-50.6892\\
-11.1007	47.1389\\
-11.1007	-47.1389\\
-10.0735	45.4314\\
-10.0735	-45.4314\\
-13.3007	5.24621\\
-13.3007	-5.24621\\
-10.1166	12.4893\\
-10.1166	-12.4893\\
-5.81257	19.1688\\
-5.81257	-19.1688\\
-8.63358	20.3575\\
-8.63358	-20.3575\\
-9.17485	22.9374\\
-9.17485	-22.9374\\
-4.09754	23.0607\\
-4.09754	-23.0607\\
-11.143	34.475\\
-11.143	-34.475\\
-10.3918	30.9184\\
-10.3918	-30.9184\\
-3.81405	29.3949\\
-3.81405	-29.3949\\
-4.61303	28.9402\\
-4.61303	-28.9402\\
-6.20325	32.4411\\
-6.20325	-32.4411\\
-6.3415	34.8909\\
-6.3415	-34.8909\\
-7.57453	38.9104\\
-7.57453	-38.9104\\
-6.0694	36.7708\\
-6.0694	-36.7708\\
-6.06509	38.2933\\
-6.06509	-38.2933\\
};
\addplot [color=mycolor3, only marks, mark size=1.2pt, mark=*, mark options={solid, fill=mycolor3, mycolor3}, forget plot]
  table[row sep=crcr]{%
-14.1762	60.7334\\
-14.1762	-60.7334\\
-10.6557	55.4763\\
-10.6557	-55.4763\\
-10.0841	53.1442\\
-10.0841	-53.1442\\
-9.27671	51.2563\\
-9.27671	-51.2563\\
-9.65385	50.6896\\
-9.65385	-50.6896\\
-11.1008	47.1387\\
-11.1008	-47.1387\\
-10.0774	45.4305\\
-10.0774	-45.4305\\
-13.3938	4.39539\\
-13.3938	-4.39539\\
-10.437	11.7682\\
-10.437	-11.7682\\
-5.27296	19.2137\\
-5.27296	-19.2137\\
-9.50705	18.7434\\
-9.50705	-18.7434\\
-10.1845	21.8949\\
-10.1845	-21.8949\\
-6.58472	26.142\\
-6.58472	-26.142\\
-6.45576	27.2918\\
-6.45576	-27.2918\\
-10.1105	28.0921\\
-10.1105	-28.0921\\
-10.6516	30.8166\\
-10.6516	-30.8166\\
-11.1146	34.4365\\
-11.1146	-34.4365\\
-6.02806	32.9062\\
-6.02806	-32.9062\\
-6.56715	34.7278\\
-6.56715	-34.7278\\
-7.57704	38.915\\
-7.57704	-38.915\\
-6.06433	36.8021\\
-6.06433	-36.8021\\
-6.10326	38.2887\\
-6.10326	-38.2887\\
};
\addplot [color=mycolor4, only marks, mark size=1.2pt, mark=*, mark options={solid, fill=mycolor4, mycolor4}, forget plot]
  table[row sep=crcr]{%
-14.1761	60.7335\\
-14.1761	-60.7335\\
-10.6557	55.4763\\
-10.6557	-55.4763\\
-10.0841	53.1442\\
-10.0841	-53.1442\\
-9.27785	51.2554\\
-9.27785	-51.2554\\
-9.6706	50.6931\\
-9.6706	-50.6931\\
-11.101	47.1384\\
-11.101	-47.1384\\
-10.0842	45.4338\\
-10.0842	-45.4338\\
-13.7043	2.04065\\
-13.7043	-2.04065\\
-10.435	11.2575\\
-10.435	-11.2575\\
-11.3245	20.2757\\
-11.3245	-20.2757\\
-9.31086	18.0687\\
-9.31086	-18.0687\\
-8.54344	19.2153\\
-8.54344	-19.2153\\
-7.66357	24.6321\\
-7.66357	-24.6321\\
-12.9294	34.9136\\
-12.9294	-34.9136\\
-10.1098	27.9556\\
-10.1098	-27.9556\\
-8.06915	28.3479\\
-8.06915	-28.3479\\
-10.5619	31.2463\\
-10.5619	-31.2463\\
-8.44481	31.8355\\
-8.44481	-31.8355\\
-7.57585	38.9171\\
-7.57585	-38.9171\\
-7.16867	35.7004\\
-7.16867	-35.7004\\
-6.11779	38.2904\\
-6.11779	-38.2904\\
-6.00731	36.8626\\
-6.00731	-36.8626\\
};
\addplot [color=mycolor5, only marks, mark size=1.2pt, mark=*, mark options={solid, fill=mycolor5, mycolor5}, forget plot]
  table[row sep=crcr]{%
-14.1761	60.7335\\
-14.1761	-60.7335\\
-10.6557	55.4763\\
-10.6557	-55.4763\\
-10.0841	53.1442\\
-10.0841	-53.1442\\
-9.28433	51.2538\\
-9.28433	-51.2538\\
-9.76713	50.7136\\
-9.76713	-50.7136\\
-11.1028	47.1381\\
-11.1028	-47.1381\\
-10.1205	45.4443\\
-10.1205	-45.4443\\
-13.7839	1.5307\\
-13.7839	-1.5307\\
-10.6127	10.9634\\
-10.6127	-10.9634\\
-9.32094	17.8446\\
-9.32094	-17.8446\\
-11.3773	20.1571\\
-11.3773	-20.1571\\
-8.91654	19.2573\\
-8.91654	-19.2573\\
-8.33667	23.3081\\
-8.33667	-23.3081\\
-10.1525	28.0124\\
-10.1525	-28.0124\\
-8.76098	29.0086\\
-8.76098	-29.0086\\
-11.6869	30.4353\\
-11.6869	-30.4353\\
-13.0507	35.0553\\
-13.0507	-35.0553\\
-8.17022	32.1515\\
-8.17022	-32.1515\\
-7.24193	37.2999\\
-7.24193	-37.2999\\
-9.10254	38.986\\
-9.10254	-38.986\\
-9.32994	37.7922\\
-9.32994	-37.7922\\
-8.87307	34.663\\
-8.87307	-34.663\\
};
\addplot [color=mycolor6, only marks, mark size=1.2pt, mark=*, mark options={solid, fill=mycolor6, mycolor6}, forget plot]
  table[row sep=crcr]{%
-14.176	60.735\\
-14.176	-60.735\\
-10.6558	55.4763\\
-10.6558	-55.4763\\
-10.0844	53.1444\\
-10.0844	-53.1444\\
-10.6142	50.9574\\
-10.6142	-50.9574\\
-9.82332	50.7252\\
-9.82332	-50.7252\\
-11.0943	47.1484\\
-11.0943	-47.1484\\
-10.4506	45.4538\\
-10.4506	-45.4538\\
-12.2799	0\\
-11.1013	10.5493\\
-11.1013	-10.5493\\
-9.8142	17.7709\\
-9.8142	-17.7709\\
-11.4222	20.0962\\
-11.4222	-20.0962\\
-9.06271	19.1348\\
-9.06271	-19.1348\\
-8.60464	22.9535\\
-8.60464	-22.9535\\
-9.75924	28.3302\\
-9.75924	-28.3302\\
-10.6019	28.1132\\
-10.6019	-28.1132\\
-12.2849	30.6248\\
-12.2849	-30.6248\\
-12.894	35.1903\\
-12.894	-35.1903\\
-8.32103	32.0628\\
-8.32103	-32.0628\\
-10.9607	38.1758\\
-10.9607	-38.1758\\
-9.5509	37.7201\\
-9.5509	-37.7201\\
-9.87485	36.3225\\
-9.87485	-36.3225\\
-9.33572	34.9987\\
-9.33572	-34.9987\\
};
\addplot [color=mycolor7, only marks, mark size=1.2pt, mark=*, mark options={solid, fill=mycolor7, mycolor7}, forget plot]
  table[row sep=crcr]{%
-14.1753	60.7226\\
-14.1753	-60.7226\\
-10.6558	55.4763\\
-10.6558	-55.4763\\
-11.2733	52.793\\
-11.2733	-52.793\\
-12.8057	50.5021\\
-12.8057	-50.5021\\
-12.686	49.8956\\
-12.686	-49.8956\\
-12.494	47.031\\
-12.494	-47.031\\
-10.7117	45.49\\
-10.7117	-45.49\\
-11.383	0\\
-11.7465	9.98804\\
-11.7465	-9.98804\\
-10.1956	17.6602\\
-10.1956	-17.6602\\
-11.4936	20.0577\\
-11.4936	-20.0577\\
-9.00614	19.0793\\
-9.00614	-19.0793\\
-8.62237	22.9372\\
-8.62237	-22.9372\\
-10.076	27.919\\
-10.076	-27.919\\
-11.0968	27.9363\\
-11.0968	-27.9363\\
-12.3906	30.819\\
-12.3906	-30.819\\
-8.6323	32.049\\
-8.6323	-32.049\\
-13.1157	35.2716\\
-13.1157	-35.2716\\
-10.9413	38.2023\\
-10.9413	-38.2023\\
-9.63255	34.7922\\
-9.63255	-34.7922\\
-9.59781	37.6677\\
-9.59781	-37.6677\\
-9.69219	36.1717\\
-9.69219	-36.1717\\
};
\addplot [color=mycolor8, only marks, mark size=1.2pt, mark=*, mark options={solid, fill=mycolor8, mycolor8}, forget plot]
  table[row sep=crcr]{%
-14.3687	60.7331\\
-14.3687	-60.7331\\
-10.6558	55.4763\\
-10.6558	-55.4763\\
-13.4757	52.3978\\
-13.4757	-52.3978\\
-12.8056	50.502\\
-12.8056	-50.502\\
-13.1704	49.8844\\
-13.1704	-49.8844\\
-12.8667	46.971\\
-12.8667	-46.971\\
-13.268	44.5594\\
-13.268	-44.5594\\
-11.3307	0\\
-11.9726	9.73381\\
-11.9726	-9.73381\\
-10.2026	17.6376\\
-10.2026	-17.6376\\
-11.5054	20.0378\\
-11.5054	-20.0378\\
-9.02945	19.0823\\
-9.02945	-19.0823\\
-8.65796	22.8905\\
-8.65796	-22.8905\\
-10.327	27.8066\\
-10.327	-27.8066\\
-11.0919	27.9381\\
-11.0919	-27.9381\\
-12.4458	30.8917\\
-12.4458	-30.8917\\
-8.92126	32.0153\\
-8.92126	-32.0153\\
-13.3228	35.3066\\
-13.3228	-35.3066\\
-9.70842	34.5537\\
-9.70842	-34.5537\\
-10.9889	38.1462\\
-10.9889	-38.1462\\
-9.59868	37.6673\\
-9.59868	-37.6673\\
-9.63373	36.1352\\
-9.63373	-36.1352\\
};
\addplot [color=mycolor9, only marks, mark size=1.2pt, mark=*, mark options={solid, fill=mycolor9, mycolor9}, forget plot]
  table[row sep=crcr]{%
-14.194	54.6783\\
-14.194	-54.6783\\
-13.5956	52.3229\\
-13.5956	-52.3229\\
-12.8056	50.502\\
-12.8056	-50.502\\
-13.1698	49.8848\\
-13.1698	-49.8848\\
-14.5054	46.3881\\
-14.5054	-46.3881\\
-13.4419	44.5161\\
-13.4419	-44.5161\\
-8.75076	0\\
-13.1374	8.91877\\
-13.1374	-8.91877\\
-12.1018	17.438\\
-12.1018	-17.438\\
-9.13599	18.8889\\
-9.13599	-18.8889\\
-11.402	19.2104\\
-11.402	-19.2104\\
-8.70677	22.8364\\
-8.70677	-22.8364\\
-10.5083	27.7539\\
-10.5083	-27.7539\\
-11.0864	27.9109\\
-11.0864	-27.9109\\
-12.5524	30.9797\\
-12.5524	-30.9797\\
-13.5119	35.4163\\
-13.5119	-35.4163\\
-9.21056	32.0509\\
-9.21056	-32.0509\\
-9.77952	34.2214\\
-9.77952	-34.2214\\
-10.9713	38.1495\\
-10.9713	-38.1495\\
-9.59966	37.6667\\
-9.59966	-37.6667\\
-9.53963	36.1237\\
-9.53963	-36.1237\\
};
\addplot [color=mycolor10, only marks, mark size=1.2pt, mark=*, mark options={solid, fill=mycolor10, mycolor10}, forget plot]
  table[row sep=crcr]{%
-14.194	54.6783\\
-14.194	-54.6783\\
-13.6132	52.3475\\
-13.6132	-52.3475\\
-12.8057	50.5018\\
-12.8057	-50.5018\\
-13.1712	49.8846\\
-13.1712	-49.8846\\
-14.6117	46.1965\\
-14.6117	-46.1965\\
-13.5516	44.5118\\
-13.5516	-44.5118\\
-8.10068	0\\
-13.5109	8.204\\
-13.5109	-8.204\\
-12.1953	17.1729\\
-12.1953	-17.1729\\
-9.13268	18.8093\\
-9.13268	-18.8093\\
-11.7301	19.5582\\
-11.7301	-19.5582\\
-8.67141	22.5915\\
-8.67141	-22.5915\\
-10.6371	27.7029\\
-10.6371	-27.7029\\
-11.0899	27.9073\\
-11.0899	-27.9073\\
-13.0955	30.2632\\
-13.0955	-30.2632\\
-13.6203	35.5139\\
-13.6203	-35.5139\\
-9.46777	32.0722\\
-9.46777	-32.0722\\
-10.9487	38.1797\\
-10.9487	-38.1797\\
-9.62832	37.6233\\
-9.62832	-37.6233\\
-9.489	36.1122\\
-9.489	-36.1122\\
-9.77181	33.9504\\
-9.77181	-33.9504\\
};
\addplot [color=mycolor11, only marks, mark size=1.2pt, mark=*, mark options={solid, fill=mycolor11, mycolor11}, forget plot]
  table[row sep=crcr]{%
-14.194	54.6782\\
-14.194	-54.6782\\
-13.6211	52.3526\\
-13.6211	-52.3526\\
-12.81	50.494\\
-12.81	-50.494\\
-13.1775	49.8863\\
-13.1775	-49.8863\\
-14.6034	46.1801\\
-14.6034	-46.1801\\
-13.6097	44.5012\\
-13.6097	-44.5012\\
-7.84396	0\\
-13.6827	7.9401\\
-13.6827	-7.9401\\
-12.0961	17.1421\\
-12.0961	-17.1421\\
-12.8873	19.5499\\
-12.8873	-19.5499\\
-9.17375	18.7866\\
-9.17375	-18.7866\\
-8.72613	22.4653\\
-8.72613	-22.4653\\
-10.6764	27.6999\\
-10.6764	-27.6999\\
-11.0808	27.9053\\
-11.0808	-27.9053\\
-13.3065	30.192\\
-13.3065	-30.192\\
-14.1057	34.3118\\
-14.1057	-34.3118\\
-11.0221	38.0928\\
-11.0221	-38.0928\\
-9.63055	37.6225\\
-9.63055	-37.6225\\
-9.52591	36.1104\\
-9.52591	-36.1104\\
-9.52642	32.0614\\
-9.52642	-32.0614\\
-9.86716	33.9888\\
-9.86716	-33.9888\\
};
\addplot [color=mycolor12, only marks, mark size=1.2pt, mark=*, mark options={solid, fill=mycolor12, mycolor12}, forget plot]
  table[row sep=crcr]{%
-14.1941	54.6781\\
-14.1941	-54.6781\\
-13.6213	52.3525\\
-13.6213	-52.3525\\
-12.8091	50.4911\\
-12.8091	-50.4911\\
-13.1776	49.8864\\
-13.1776	-49.8864\\
-14.5971	46.1752\\
-14.5971	-46.1752\\
-13.6103	44.4967\\
-13.6103	-44.4967\\
-7.31083	0\\
-13.7736	7.60154\\
-13.7736	-7.60154\\
-14.1336	19.6167\\
-14.1336	-19.6167\\
-12.2029	17.0151\\
-12.2029	-17.0151\\
-9.14852	18.5269\\
-9.14852	-18.5269\\
-8.61927	22.3511\\
-8.61927	-22.3511\\
-10.6619	27.7192\\
-10.6619	-27.7192\\
-11.0807	27.8961\\
-11.0807	-27.8961\\
-13.4633	29.6108\\
-13.4633	-29.6108\\
-14.0929	33.4288\\
-14.0929	-33.4288\\
-11.024	38.0989\\
-11.024	-38.0989\\
-9.62121	37.6128\\
-9.62121	-37.6128\\
-9.53381	36.1061\\
-9.53381	-36.1061\\
-9.56839	32.0883\\
-9.56839	-32.0883\\
-9.94893	33.9842\\
-9.94893	-33.9842\\
};
\end{axis}

\begin{axis}[%
width={\figscale*0.747in},
height={\figscale*0.745in},
at={(\figscale*1.438in,\figscale*0.596in)},
scale only axis,
separate axis lines,
every outer x axis line/.append style={black},
every x tick label/.append style={font=\figticfont},
every x tick/.append style={black},
xmin=-15,
xmax=1,
xlabel={$\mathrm{Re}(\lambda)$ in rad/s},
every outer y axis line/.append style={black},
every y tick label/.append style={font=\figticfont},
every y tick/.append style={black},
ymin=-74.2188,
ymax=74.2188,
axis background/.style={fill=white},
title style={font=\figcornerfont},
title={SC500},
xmajorgrids,
ymajorgrids,
grid style={white!55!black, opacity=0.3}
]
\addplot [color=mycolor1, dashed, line width=0.6pt, forget plot]
  table[row sep=crcr]{%
-15	149.248\\
-14.6939	146.202\\
-14.3878	143.156\\
-14.0816	140.11\\
-13.7755	137.065\\
-13.4694	134.019\\
-13.1633	130.973\\
-12.8571	127.927\\
-12.551	124.881\\
-12.2449	121.835\\
-11.9388	118.789\\
-11.6327	115.743\\
-11.3265	112.698\\
-11.0204	109.652\\
-10.7143	106.606\\
-10.4082	103.56\\
-10.102	100.514\\
-9.79592	97.4682\\
-9.4898	94.4223\\
-9.18367	91.3764\\
-8.87755	88.3305\\
-8.57143	85.2846\\
-8.26531	82.2388\\
-7.95918	79.1929\\
-7.65306	76.147\\
-7.34694	73.1011\\
-7.04082	70.0552\\
-6.73469	67.0094\\
-6.42857	63.9635\\
-6.12245	60.9176\\
-5.81633	57.8717\\
-5.5102	54.8258\\
-5.20408	51.78\\
-4.89796	48.7341\\
-4.59184	45.6882\\
-4.28571	42.6423\\
-3.97959	39.5964\\
-3.67347	36.5506\\
-3.36735	33.5047\\
-3.06122	30.4588\\
-2.7551	27.4129\\
-2.44898	24.367\\
-2.14286	21.3212\\
-1.83673	18.2753\\
-1.53061	15.2294\\
-1.22449	12.1835\\
-0.918367	9.13764\\
-0.612245	6.09176\\
-0.306122	3.04588\\
0	-0\\
};
\addplot [color=mycolor1, dashed, line width=0.6pt, forget plot]
  table[row sep=crcr]{%
-15	-149.248\\
-14.6939	-146.202\\
-14.3878	-143.156\\
-14.0816	-140.11\\
-13.7755	-137.065\\
-13.4694	-134.019\\
-13.1633	-130.973\\
-12.8571	-127.927\\
-12.551	-124.881\\
-12.2449	-121.835\\
-11.9388	-118.789\\
-11.6327	-115.743\\
-11.3265	-112.698\\
-11.0204	-109.652\\
-10.7143	-106.606\\
-10.4082	-103.56\\
-10.102	-100.514\\
-9.79592	-97.4682\\
-9.4898	-94.4223\\
-9.18367	-91.3764\\
-8.87755	-88.3305\\
-8.57143	-85.2846\\
-8.26531	-82.2388\\
-7.95918	-79.1929\\
-7.65306	-76.147\\
-7.34694	-73.1011\\
-7.04082	-70.0552\\
-6.73469	-67.0094\\
-6.42857	-63.9635\\
-6.12245	-60.9176\\
-5.81633	-57.8717\\
-5.5102	-54.8258\\
-5.20408	-51.78\\
-4.89796	-48.7341\\
-4.59184	-45.6882\\
-4.28571	-42.6423\\
-3.97959	-39.5964\\
-3.67347	-36.5506\\
-3.36735	-33.5047\\
-3.06122	-30.4588\\
-2.7551	-27.4129\\
-2.44898	-24.367\\
-2.14286	-21.3212\\
-1.83673	-18.2753\\
-1.53061	-15.2294\\
-1.22449	-12.1835\\
-0.918367	-9.13764\\
-0.612245	-6.09176\\
-0.306122	-3.04588\\
0	0\\
};
\addplot [color=mycolor1, dashed, line width=0.6pt, forget plot]
  table[row sep=crcr]{%
-3	-74.2188\\
-3	74.2188\\
};
\addplot [color=white!88!black, line width=0.6pt, forget plot]
  table[row sep=crcr]{%
-15	0\\
1	0\\
};
\addplot [color=mycolor2, line width=0.9pt, only marks, mark size=2.2pt, mark=x, mark options={solid, mycolor2}, forget plot]
  table[row sep=crcr]{%
-14.5939	64.5381\\
-14.5939	-64.5381\\
-13.3799	61.6619\\
-13.3799	-61.6619\\
-11.5339	21.0111\\
-11.5339	-21.0111\\
-3.72631	24.3778\\
-3.72631	-24.3778\\
-4.91907	24.6344\\
-4.91907	-24.6344\\
-9.12798	39.5453\\
-9.12798	-39.5453\\
-6.39436	33.9619\\
-6.39436	-33.9619\\
-4.02876	33.423\\
-4.02876	-33.423\\
-3.23242	27.6612\\
-3.23242	-27.6612\\
-7.76806	43.9947\\
-7.76806	-43.9947\\
-3.59223	29.5078\\
-3.59223	-29.5078\\
-12.6373	46.6786\\
-12.6373	-46.6786\\
-9.44342	50.2707\\
-9.44342	-50.2707\\
-3.07363	29.7669\\
-3.07363	-29.7669\\
-10.611	50.8406\\
-10.611	-50.8406\\
-12.0958	54.7327\\
-12.0958	-54.7327\\
-10.0842	53.9999\\
-10.0842	-53.9999\\
-9.48062	52.4582\\
-9.48062	-52.4582\\
-10.9224	56.1903\\
-10.9224	-56.1903\\
-10.8397	55.985\\
-10.8397	-55.985\\
-3.0004	29.8516\\
-3.0004	-29.8516\\
};
\addplot [color=mycolor3, only marks, mark size=1.2pt, mark=*, mark options={solid, fill=mycolor3, mycolor3}, forget plot]
  table[row sep=crcr]{%
-14.5939	64.5381\\
-14.5939	-64.5381\\
-13.3799	61.6619\\
-13.3799	-61.6619\\
-4.52665	32.8154\\
-4.52665	-32.8154\\
-5.78649	25.4562\\
-5.78649	-25.4562\\
-7.18486	33.6569\\
-7.18486	-33.6569\\
-9.12799	39.5453\\
-9.12799	-39.5453\\
-11.8162	20.7682\\
-11.8162	-20.7682\\
-7.76807	43.9947\\
-7.76807	-43.9947\\
-12.8099	21.9871\\
-12.8099	-21.9871\\
-12.9852	24.2485\\
-12.9852	-24.2485\\
-12.6373	46.6786\\
-12.6373	-46.6786\\
-13.3703	26.1802\\
-13.3703	-26.1802\\
-9.44342	50.2707\\
-9.44342	-50.2707\\
-14.0094	26.3156\\
-14.0094	-26.3156\\
-10.6109	50.8404\\
-10.6109	-50.8404\\
-12.1963	54.7266\\
-12.1963	-54.7266\\
-10.0842	53.9999\\
-10.0842	-53.9999\\
-9.48062	52.4582\\
-9.48062	-52.4582\\
-10.9224	56.1903\\
-10.9224	-56.1903\\
-10.8397	55.985\\
-10.8397	-55.985\\
};
\addplot [color=mycolor4, only marks, mark size=1.2pt, mark=*, mark options={solid, fill=mycolor4, mycolor4}, forget plot]
  table[row sep=crcr]{%
-14.5939	64.5381\\
-14.5939	-64.5381\\
-13.3925	61.641\\
-13.3925	-61.641\\
-11.1554	38.7891\\
-11.1554	-38.7891\\
-10.6135	50.8395\\
-10.6135	-50.8395\\
-9.07499	25.3291\\
-9.07499	-25.3291\\
-12.1436	21.1672\\
-12.1436	-21.1672\\
-13.6539	20.8552\\
-13.6539	-20.8552\\
-9.4805	52.458\\
-9.4805	-52.458\\
-14.113	30.599\\
-14.113	-30.599\\
-12.9038	24.0802\\
-12.9038	-24.0802\\
-10.0843	54.0001\\
-10.0843	-54.0001\\
-13.0228	54.6244\\
-13.0228	-54.6244\\
-13.2671	26.2167\\
-13.2671	-26.2167\\
-14.048	26.3117\\
-14.048	-26.3117\\
-10.9224	56.1903\\
-10.9224	-56.1903\\
-10.8397	55.985\\
-10.8397	-55.985\\
};
\addplot [color=mycolor5, only marks, mark size=1.2pt, mark=*, mark options={solid, fill=mycolor5, mycolor5}, forget plot]
  table[row sep=crcr]{%
-14.5939	64.5381\\
-14.5939	-64.5381\\
-13.3926	61.6403\\
-13.3926	-61.6403\\
-12.4371	54.9115\\
-12.4371	-54.9115\\
-13.1195	54.6934\\
-13.1195	-54.6934\\
-13.8516	20.9337\\
-13.8516	-20.9337\\
-12.5382	22.7048\\
-12.5382	-22.7048\\
-12.5484	25.3966\\
-12.5484	-25.3966\\
-14.1456	30.6261\\
-14.1456	-30.6261\\
-13.0194	26.4667\\
-13.0194	-26.4667\\
-14.0766	26.3312\\
-14.0766	-26.3312\\
};
\addplot [color=mycolor6, only marks, mark size=1.2pt, mark=*, mark options={solid, fill=mycolor6, mycolor6}, forget plot]
  table[row sep=crcr]{%
-14.5939	64.5379\\
-14.5939	-64.5379\\
-14.7962	0\\
-13.0955	54.6848\\
-13.0955	-54.6848\\
-13.8112	20.9246\\
-13.8112	-20.9246\\
-12.8543	21.9834\\
-12.8543	-21.9834\\
-14.1663	30.6243\\
-14.1663	-30.6243\\
-13.0619	25.3026\\
-13.0619	-25.3026\\
-13.2394	26.7479\\
-13.2394	-26.7479\\
-14.1124	26.3477\\
-14.1124	-26.3477\\
};
\addplot [color=mycolor7, only marks, mark size=1.2pt, mark=*, mark options={solid, fill=mycolor7, mycolor7}, forget plot]
  table[row sep=crcr]{%
-10.8558	0\\
-12.9396	21.2384\\
-12.9396	-21.2384\\
-14.1471	20.9435\\
-14.1471	-20.9435\\
-14.1644	30.5976\\
-14.1644	-30.5976\\
-13.2865	25.1345\\
-13.2865	-25.1345\\
-13.4635	26.6327\\
-13.4635	-26.6327\\
-14.19	26.3252\\
-14.19	-26.3252\\
};
\addplot [color=mycolor8, only marks, mark size=1.2pt, mark=*, mark options={solid, fill=mycolor8, mycolor8}, forget plot]
  table[row sep=crcr]{%
-10.7271	0\\
-13.1258	20.5996\\
-13.1258	-20.5996\\
-14.3041	21.0048\\
-14.3041	-21.0048\\
-14.161	30.5881\\
-14.161	-30.5881\\
-13.4483	25.13\\
-13.4483	-25.13\\
-13.5701	26.5947\\
-13.5701	-26.5947\\
-14.2497	26.2976\\
-14.2497	-26.2976\\
};
\addplot [color=mycolor9, only marks, mark size=1.2pt, mark=*, mark options={solid, fill=mycolor9, mycolor9}, forget plot]
  table[row sep=crcr]{%
-9.63685	0\\
-13.2772	20.3436\\
-13.2772	-20.3436\\
-14.4059	20.9909\\
-14.4059	-20.9909\\
-14.1609	30.5838\\
-14.1609	-30.5838\\
-13.5158	25.1415\\
-13.5158	-25.1415\\
-14.2839	26.2775\\
-14.2839	-26.2775\\
-13.6006	26.5779\\
-13.6006	-26.5779\\
};
\addplot [color=mycolor10, only marks, mark size=1.2pt, mark=*, mark options={solid, fill=mycolor10, mycolor10}, forget plot]
  table[row sep=crcr]{%
-9.534	0\\
-13.5908	0\\
-13.4713	20.1743\\
-13.4713	-20.1743\\
-14.5013	20.9494\\
-14.5013	-20.9494\\
-14.161	30.5808\\
-14.161	-30.5808\\
-13.5736	25.1473\\
-13.5736	-25.1473\\
-14.3087	26.2484\\
-14.3087	-26.2484\\
-13.625	26.5593\\
-13.625	-26.5593\\
};
\addplot [color=mycolor11, only marks, mark size=1.2pt, mark=*, mark options={solid, fill=mycolor11, mycolor11}, forget plot]
  table[row sep=crcr]{%
-9.42594	0\\
-13.5262	0\\
-13.6469	20.1247\\
-13.6469	-20.1247\\
-14.5777	20.9145\\
-14.5777	-20.9145\\
-14.1628	30.5802\\
-14.1628	-30.5802\\
-13.6081	25.1591\\
-13.6081	-25.1591\\
-14.3276	26.2291\\
-14.3276	-26.2291\\
-13.6362	26.551\\
-13.6362	-26.551\\
};
\addplot [color=mycolor12, only marks, mark size=1.2pt, mark=*, mark options={solid, fill=mycolor12, mycolor12}, forget plot]
  table[row sep=crcr]{%
-9.30309	0\\
-13.4082	0\\
-13.7779	20.0973\\
-13.7779	-20.0973\\
-14.6342	20.8602\\
-14.6342	-20.8602\\
-14.1638	30.5795\\
-14.1638	-30.5795\\
-13.6392	25.1588\\
-13.6392	-25.1588\\
-14.3376	26.2035\\
-14.3376	-26.2035\\
-13.6372	26.5464\\
-13.6372	-26.5464\\
};
\end{axis}

\begin{axis}[%
width={\figscale*0.747in},
height={\figscale*0.745in},
at={(\figscale*2.421in,\figscale*0.596in)},
scale only axis,
separate axis lines,
every outer x axis line/.append style={black},
every x tick label/.append style={font=\figticfont},
every x tick/.append style={black},
xmin=-15,
xmax=1,
xlabel={$\mathrm{Re}(\lambda)$ in rad/s},
every outer y axis line/.append style={black},
every y tick label/.append style={font=\figticfont},
every y tick/.append style={black},
ymin=-69.2131,
ymax=69.2131,
axis background/.style={fill=white},
title style={font=\figcornerfont},
title={Texas2000},
xmajorgrids,
ymajorgrids,
grid style={white!55!black, opacity=0.3}
]
\addplot [color=mycolor1, dashed, line width=0.6pt, forget plot]
  table[row sep=crcr]{%
-15	149.248\\
-14.6939	146.202\\
-14.3878	143.156\\
-14.0816	140.11\\
-13.7755	137.065\\
-13.4694	134.019\\
-13.1633	130.973\\
-12.8571	127.927\\
-12.551	124.881\\
-12.2449	121.835\\
-11.9388	118.789\\
-11.6327	115.743\\
-11.3265	112.698\\
-11.0204	109.652\\
-10.7143	106.606\\
-10.4082	103.56\\
-10.102	100.514\\
-9.79592	97.4682\\
-9.4898	94.4223\\
-9.18367	91.3764\\
-8.87755	88.3305\\
-8.57143	85.2846\\
-8.26531	82.2388\\
-7.95918	79.1929\\
-7.65306	76.147\\
-7.34694	73.1011\\
-7.04082	70.0552\\
-6.73469	67.0094\\
-6.42857	63.9635\\
-6.12245	60.9176\\
-5.81633	57.8717\\
-5.5102	54.8258\\
-5.20408	51.78\\
-4.89796	48.7341\\
-4.59184	45.6882\\
-4.28571	42.6423\\
-3.97959	39.5964\\
-3.67347	36.5506\\
-3.36735	33.5047\\
-3.06122	30.4588\\
-2.7551	27.4129\\
-2.44898	24.367\\
-2.14286	21.3212\\
-1.83673	18.2753\\
-1.53061	15.2294\\
-1.22449	12.1835\\
-0.918367	9.13764\\
-0.612245	6.09176\\
-0.306122	3.04588\\
0	-0\\
};
\addplot [color=mycolor1, dashed, line width=0.6pt, forget plot]
  table[row sep=crcr]{%
-15	-149.248\\
-14.6939	-146.202\\
-14.3878	-143.156\\
-14.0816	-140.11\\
-13.7755	-137.065\\
-13.4694	-134.019\\
-13.1633	-130.973\\
-12.8571	-127.927\\
-12.551	-124.881\\
-12.2449	-121.835\\
-11.9388	-118.789\\
-11.6327	-115.743\\
-11.3265	-112.698\\
-11.0204	-109.652\\
-10.7143	-106.606\\
-10.4082	-103.56\\
-10.102	-100.514\\
-9.79592	-97.4682\\
-9.4898	-94.4223\\
-9.18367	-91.3764\\
-8.87755	-88.3305\\
-8.57143	-85.2846\\
-8.26531	-82.2388\\
-7.95918	-79.1929\\
-7.65306	-76.147\\
-7.34694	-73.1011\\
-7.04082	-70.0552\\
-6.73469	-67.0094\\
-6.42857	-63.9635\\
-6.12245	-60.9176\\
-5.81633	-57.8717\\
-5.5102	-54.8258\\
-5.20408	-51.78\\
-4.89796	-48.7341\\
-4.59184	-45.6882\\
-4.28571	-42.6423\\
-3.97959	-39.5964\\
-3.67347	-36.5506\\
-3.36735	-33.5047\\
-3.06122	-30.4588\\
-2.7551	-27.4129\\
-2.44898	-24.367\\
-2.14286	-21.3212\\
-1.83673	-18.2753\\
-1.53061	-15.2294\\
-1.22449	-12.1835\\
-0.918367	-9.13764\\
-0.612245	-6.09176\\
-0.306122	-3.04588\\
0	0\\
};
\addplot [color=mycolor1, dashed, line width=0.6pt, forget plot]
  table[row sep=crcr]{%
-3	-69.2131\\
-3	69.2131\\
};
\addplot [color=white!88!black, line width=0.6pt, forget plot]
  table[row sep=crcr]{%
-15	0\\
1	0\\
};
\addplot [color=mycolor2, line width=0.9pt, only marks, mark size=2.2pt, mark=x, mark options={solid, mycolor2}, forget plot]
  table[row sep=crcr]{%
-13.572	60.1307\\
-13.572	-60.1307\\
-13.2503	58.2639\\
-13.2503	-58.2639\\
-11.6928	57.9025\\
-11.6928	-57.9025\\
-13.7438	59.8599\\
-13.7438	-59.8599\\
-12.8527	56.7651\\
-12.8527	-56.7651\\
-13.4575	57.0687\\
-13.4575	-57.0687\\
-10.5698	54.1653\\
-10.5698	-54.1653\\
-11.985	53.9002\\
-11.985	-53.9002\\
-14.2205	55.3257\\
-14.2205	-55.3257\\
-10.9904	55.906\\
-10.9904	-55.906\\
-11.2039	56.877\\
-11.2039	-56.877\\
-11.2105	55.8382\\
-11.2105	-55.8382\\
-11.3549	57.2467\\
-11.3549	-57.2467\\
-11.4528	57.4839\\
-11.4528	-57.4839\\
-11.4863	57.5666\\
-11.4863	-57.5666\\
-10.5729	53.6702\\
-10.5729	-53.6702\\
-10.8956	54.0155\\
-10.8956	-54.0155\\
-11.2295	0\\
-3.6615	0\\
-7.68909	0\\
-10.099	51.0599\\
-10.099	-51.0599\\
-8.32977	48.8965\\
-8.32977	-48.8965\\
-8.49483	48.4212\\
-8.49483	-48.4212\\
-8.38704	48.2259\\
-8.38704	-48.2259\\
-8.82866	47.8287\\
-8.82866	-47.8287\\
-8.13756	48.7263\\
-8.13756	-48.7263\\
-11.4049	57.3706\\
-11.4049	-57.3706\\
-3.37549	10.555\\
-3.37549	-10.555\\
-13.8045	41.3298\\
-13.8045	-41.3298\\
-13.7416	45.6381\\
-13.7416	-45.6381\\
-11.2854	45.0977\\
-11.2854	-45.0977\\
-9.78541	45.8327\\
-9.78541	-45.8327\\
-7.69763	46.4532\\
-7.69763	-46.4532\\
-7.57789	46.03\\
-7.57789	-46.03\\
-9.59693	45.9742\\
-9.59693	-45.9742\\
-7.7484	46.0678\\
-7.7484	-46.0678\\
-7.31151	45.9243\\
-7.31151	-45.9243\\
-7.45439	45.1151\\
-7.45439	-45.1151\\
-14.2814	19.6336\\
-14.2814	-19.6336\\
-3.37955	17.7223\\
-3.37955	-17.7223\\
-8.00933	19.1156\\
-8.00933	-19.1156\\
-14.1653	38.9652\\
-14.1653	-38.9652\\
-5.7957	40.2655\\
-5.7957	-40.2655\\
-7.84923	43.3657\\
-7.84923	-43.3657\\
-6.54733	41.5879\\
-6.54733	-41.5879\\
-7.3836	43.9962\\
-7.3836	-43.9962\\
-6.97578	44.2337\\
-6.97578	-44.2337\\
-6.7416	44.1014\\
-6.7416	-44.1014\\
-6.89637	43.9724\\
-6.89637	-43.9724\\
-7.27443	40.4406\\
-7.27443	-40.4406\\
-9.6942	40.2175\\
-9.6942	-40.2175\\
-6.99482	44.1366\\
-6.99482	-44.1366\\
-6.59464	42.8063\\
-6.59464	-42.8063\\
-8.24952	42.4786\\
-8.24952	-42.4786\\
-4.00161	18.2008\\
-4.00161	-18.2008\\
-6.74812	42.9955\\
-6.74812	-42.9955\\
-7.06041	41.8574\\
-7.06041	-41.8574\\
-7.78491	40.0346\\
-7.78491	-40.0346\\
-6.23892	38.9767\\
-6.23892	-38.9767\\
-6.45407	39.3491\\
-6.45407	-39.3491\\
-5.13189	37.6765\\
-5.13189	-37.6765\\
-7.81802	34.9412\\
-7.81802	-34.9412\\
-5.84094	20.9132\\
-5.84094	-20.9132\\
-3.20584	20.171\\
-3.20584	-20.171\\
-7.13432	23.4371\\
-7.13432	-23.4371\\
-3.34561	20.346\\
-3.34561	-20.346\\
-6.00414	35.6862\\
-6.00414	-35.6862\\
-3.00417	20.6907\\
-3.00417	-20.6907\\
-4.39277	35.9613\\
-4.39277	-35.9613\\
-5.3343	36.1721\\
-5.3343	-36.1721\\
-4.07041	21.3228\\
-4.07041	-21.3228\\
-3.03244	21.6699\\
-3.03244	-21.6699\\
-3.68057	21.575\\
-3.68057	-21.575\\
-6.82586	33.7363\\
-6.82586	-33.7363\\
-4.75038	34.5626\\
-4.75038	-34.5626\\
-8.00746	30.6151\\
-8.00746	-30.6151\\
-4.33406	33.7933\\
-4.33406	-33.7933\\
-3.46354	22.5822\\
-3.46354	-22.5822\\
-3.11221	22.5752\\
-3.11221	-22.5752\\
-4.28287	34.0219\\
-4.28287	-34.0219\\
-3.00313	22.7527\\
-3.00313	-22.7527\\
-4.69222	33.1277\\
-4.69222	-33.1277\\
-4.4649	23.6053\\
-4.4649	-23.6053\\
-6.95989	29.939\\
-6.95989	-29.939\\
-4.67092	32.5818\\
-4.67092	-32.5818\\
-3.0277	23.1458\\
-3.0277	-23.1458\\
-3.31166	23.4116\\
-3.31166	-23.4116\\
-3.7574	32.4326\\
-3.7574	-32.4326\\
-5.46641	31.7579\\
-5.46641	-31.7579\\
-6.35079	26.8278\\
-6.35079	-26.8278\\
-3.57592	32.2761\\
-3.57592	-32.2761\\
-4.50249	31.7758\\
-4.50249	-31.7758\\
-4.49513	24.6403\\
-4.49513	-24.6403\\
-6.03543	28.1359\\
-6.03543	-28.1359\\
-4.86234	31.1946\\
-4.86234	-31.1946\\
-3.06014	24.5226\\
-3.06014	-24.5226\\
-4.4933	30.8321\\
-4.4933	-30.8321\\
-3.07428	24.7312\\
-3.07428	-24.7312\\
-3.26037	30.8831\\
-3.26037	-30.8831\\
-3.65868	30.7263\\
-3.65868	-30.7263\\
-3.26435	30.2954\\
-3.26435	-30.2954\\
-3.06653	25.4052\\
-3.06653	-25.4052\\
-3.21646	25.792\\
-3.21646	-25.792\\
-3.05704	25.6955\\
-3.05704	-25.6955\\
-3.0444	26.0227\\
-3.0444	-26.0227\\
-3.00055	25.4501\\
-3.00055	-25.4501\\
-3.00037	25.5072\\
-3.00037	-25.5072\\
-3.00038	25.472\\
-3.00038	-25.472\\
-3.75577	26.4457\\
-3.75577	-26.4457\\
-3.19615	26.2922\\
-3.19615	-26.2922\\
-3.03023	26.4721\\
-3.03023	-26.4721\\
-3.00009	25.5645\\
-3.00009	-25.5645\\
-4.5132	29.3053\\
-4.5132	-29.3053\\
-4.23976	29.7542\\
-4.23976	-29.7542\\
-3.28608	26.7698\\
-3.28608	-26.7698\\
-3.00003	25.4267\\
-3.00003	-25.4267\\
-3.04626	26.9503\\
-3.04626	-26.9503\\
-4.16797	28.1362\\
-4.16797	-28.1362\\
-4.12622	28.9878\\
-4.12622	-28.9878\\
-3.5438	29.7215\\
-3.5438	-29.7215\\
-3.384	29.9199\\
-3.384	-29.9199\\
-3.52365	29.644\\
-3.52365	-29.644\\
-3.03177	29.9455\\
-3.03177	-29.9455\\
-3.12756	29.5691\\
-3.12756	-29.5691\\
-3.02684	29.614\\
-3.02684	-29.614\\
-3.04089	29.9179\\
-3.04089	-29.9179\\
-3.0191	29.6884\\
-3.0191	-29.6884\\
-3.02779	29.8597\\
-3.02779	-29.8597\\
-3.02186	29.8917\\
-3.02186	-29.8917\\
-3.01302	29.771\\
-3.01302	-29.771\\
-3.08609	29.6218\\
-3.08609	-29.6218\\
-3.02167	29.9062\\
-3.02167	-29.9062\\
-3.02646	29.7884\\
-3.02646	-29.7884\\
-3.00432	29.8051\\
-3.00432	-29.8051\\
-3.01012	29.7947\\
-3.01012	-29.7947\\
-3.02065	26.9446\\
-3.02065	-26.9446\\
-3.03425	27.0937\\
-3.03425	-27.0937\\
-3.49411	27.5448\\
-3.49411	-27.5448\\
-3.12767	27.1416\\
-3.12767	-27.1416\\
-3.01701	27.3853\\
-3.01701	-27.3853\\
-3.06334	27.3867\\
-3.06334	-27.3867\\
-3.80223	28.8418\\
-3.80223	-28.8418\\
-3.00504	29.52\\
-3.00504	-29.52\\
-3.00643	29.3135\\
-3.00643	-29.3135\\
-3.01311	29.4339\\
-3.01311	-29.4339\\
-3	29.8495\\
-3	-29.8495\\
-3.0217	29.9146\\
-3.0217	-29.9146\\
-3.49037	28.8975\\
-3.49037	-28.8975\\
-3.58452	28.4454\\
-3.58452	-28.4454\\
-3.53272	28.1571\\
-3.53272	-28.1571\\
-3.08255	27.5516\\
-3.08255	-27.5516\\
-3.18644	29.1767\\
-3.18644	-29.1767\\
-3.13625	29.2316\\
-3.13625	-29.2316\\
-3.18877	27.6126\\
-3.18877	-27.6126\\
-3.00019	27.2178\\
-3.00019	-27.2178\\
-3.27276	27.6637\\
-3.27276	-27.6637\\
-3.01401	29.1904\\
-3.01401	-29.1904\\
-3.24198	27.7866\\
-3.24198	-27.7866\\
-3.02267	29.1722\\
-3.02267	-29.1722\\
-3.02169	29.8012\\
-3.02169	-29.8012\\
-3.27395	28.1145\\
-3.27395	-28.1145\\
-3.07914	27.9679\\
-3.07914	-27.9679\\
-3.02501	29.094\\
-3.02501	-29.094\\
-3.00005	27.8981\\
-3.00005	-27.8981\\
-3.01699	28.0715\\
-3.01699	-28.0715\\
-3.01969	28.1274\\
-3.01969	-28.1274\\
-3.18214	28.878\\
-3.18214	-28.878\\
-3.02177	28.7297\\
-3.02177	-28.7297\\
-3.03904	28.8825\\
-3.03904	-28.8825\\
-3.03944	28.3484\\
-3.03944	-28.3484\\
-3.01853	29.07\\
-3.01853	-29.07\\
-3.00438	28.045\\
-3.00438	-28.045\\
-3.24763	28.8124\\
-3.24763	-28.8124\\
-3.00761	28.4471\\
-3.00761	-28.4471\\
-3.05478	28.5152\\
-3.05478	-28.5152\\
-3.16933	28.7549\\
-3.16933	-28.7549\\
-3.00669	28.5861\\
-3.00669	-28.5861\\
-3.06231	28.6249\\
-3.06231	-28.6249\\
-3.04937	28.5125\\
-3.04937	-28.5125\\
-3.01277	28.5201\\
-3.01277	-28.5201\\
-3.00003	25.4264\\
-3.00003	-25.4264\\
};
\addplot [color=mycolor3, only marks, mark size=1.2pt, mark=*, mark options={solid, fill=mycolor3, mycolor3}, forget plot]
  table[row sep=crcr]{%
-13.572	60.1307\\
-13.572	-60.1307\\
-13.7438	59.8599\\
-13.7438	-59.8599\\
-13.2506	58.2639\\
-13.2506	-58.2639\\
-11.6928	57.9025\\
-11.6928	-57.9025\\
-12.8528	56.7652\\
-12.8528	-56.7652\\
-13.4576	57.0678\\
-13.4576	-57.0678\\
-14.2205	55.3257\\
-14.2205	-55.3257\\
-10.9904	55.906\\
-10.9904	-55.906\\
-11.9853	53.9004\\
-11.9853	-53.9004\\
-10.5698	54.1653\\
-10.5698	-54.1653\\
-10.5729	53.6702\\
-10.5729	-53.6702\\
-11.2039	56.877\\
-11.2039	-56.877\\
-11.2105	55.8383\\
-11.2105	-55.8383\\
-11.3549	57.2467\\
-11.3549	-57.2467\\
-11.4528	57.4839\\
-11.4528	-57.4839\\
-10.8957	54.0155\\
-10.8957	-54.0155\\
-11.4863	57.5666\\
-11.4863	-57.5666\\
-10.0993	51.0598\\
-10.0993	-51.0598\\
-11.4049	57.3706\\
-11.4049	-57.3706\\
-13.7418	45.6381\\
-13.7418	-45.6381\\
-8.32979	48.8965\\
-8.32979	-48.8965\\
-8.49484	48.4214\\
-8.49484	-48.4214\\
-8.38704	48.2259\\
-8.38704	-48.2259\\
-8.13756	48.7263\\
-8.13756	-48.7263\\
-8.82867	47.8286\\
-8.82867	-47.8286\\
-11.2819	45.0982\\
-11.2819	-45.0982\\
-13.8044	41.3297\\
-13.8044	-41.3297\\
-9.78626	45.8333\\
-9.78626	-45.8333\\
-7.69764	46.4531\\
-7.69764	-46.4531\\
-7.57798	46.03\\
-7.57798	-46.03\\
-9.59759	45.9758\\
-9.59759	-45.9758\\
-7.74865	46.0675\\
-7.74865	-46.0675\\
-7.31151	45.9243\\
-7.31151	-45.9243\\
-14.1666	38.965\\
-14.1666	-38.965\\
-14.3472	19.7094\\
-14.3472	-19.7094\\
-3.37346	10.555\\
-3.37346	-10.555\\
-7.45439	45.1151\\
-7.45439	-45.1151\\
-9.69397	40.2171\\
-9.69397	-40.2171\\
-7.84725	43.366\\
-7.84725	-43.366\\
-7.38357	43.9962\\
-7.38357	-43.9962\\
-6.97582	44.2337\\
-6.97582	-44.2337\\
-6.7416	44.1014\\
-6.7416	-44.1014\\
-5.79568	40.2658\\
-5.79568	-40.2658\\
-6.89638	43.9723\\
-6.89638	-43.9723\\
-6.54722	41.5878\\
-6.54722	-41.5878\\
-8.25016	42.4789\\
-8.25016	-42.4789\\
-6.99481	44.1366\\
-6.99481	-44.1366\\
-6.74855	42.9953\\
-6.74855	-42.9953\\
-6.59471	42.8062\\
-6.59471	-42.8062\\
-7.27629	40.4393\\
-7.27629	-40.4393\\
-7.06008	41.8563\\
-7.06008	-41.8563\\
-14.8833	0\\
-3.53019	0\\
-3.74788	0\\
-13.7768	0\\
-12.8104	0\\
-4.14078	0\\
-11.9626	0\\
-8.77125	0\\
-6.29575	0\\
-5.27792	0\\
-4.27143	0\\
-4.79728	0\\
-4.56559	0\\
-11.6712	0\\
-11.0526	0\\
-10.5304	0\\
-10.5569	0\\
-5.13213	37.6766\\
-5.13213	-37.6766\\
-6.23834	38.9761\\
-6.23834	-38.9761\\
-6.47318	39.332\\
-6.47318	-39.332\\
-7.78483	40.035\\
-7.78483	-40.035\\
-8.19305	0\\
-6.9891	0\\
-7.66365	0\\
-7.57426	0\\
-3.38238	17.7189\\
-3.38238	-17.7189\\
-8.00809	19.155\\
-8.00809	-19.155\\
-4.00525	18.2377\\
-4.00525	-18.2377\\
-7.81752	34.9394\\
-7.81752	-34.9394\\
-6.00289	35.6893\\
-6.00289	-35.6893\\
-5.33455	36.1737\\
-5.33455	-36.1737\\
-4.39277	35.9613\\
-4.39277	-35.9613\\
-6.82343	33.7343\\
-6.82343	-33.7343\\
-4.75836	34.5434\\
-4.75836	-34.5434\\
-8.00864	30.6148\\
-8.00864	-30.6148\\
-5.84047	20.9062\\
-5.84047	-20.9062\\
-3.20611	20.1693\\
-3.20611	-20.1693\\
-3.34657	20.3451\\
-3.34657	-20.3451\\
-3.00416	20.6907\\
-3.00416	-20.6907\\
-4.33405	33.7933\\
-4.33405	-33.7933\\
-4.28323	34.0217\\
-4.28323	-34.0217\\
-4.06519	21.3271\\
-4.06519	-21.3271\\
-7.14499	23.4316\\
-7.14499	-23.4316\\
-3.68127	21.5773\\
-3.68127	-21.5773\\
-3.03274	21.67\\
-3.03274	-21.67\\
-4.66441	32.5852\\
-4.66441	-32.5852\\
-4.69137	33.1246\\
-4.69137	-33.1246\\
-7.06106	0\\
-6.93422	29.9361\\
-6.93422	-29.9361\\
-3.75761	32.4329\\
-3.75761	-32.4329\\
-3.4617	22.5882\\
-3.4617	-22.5882\\
-3.11352	22.5774\\
-3.11352	-22.5774\\
-3.57592	32.2761\\
-3.57592	-32.2761\\
-5.47545	31.7686\\
-5.47545	-31.7686\\
-3.00318	22.7527\\
-3.00318	-22.7527\\
-4.51917	31.7451\\
-4.51917	-31.7451\\
-7.59116	0\\
-4.40943	23.5934\\
-4.40943	-23.5934\\
-3.32604	23.5037\\
-3.32604	-23.5037\\
-6.35339	26.8291\\
-6.35339	-26.8291\\
-4.87182	31.1978\\
-4.87182	-31.1978\\
-6.02841	28.1152\\
-6.02841	-28.1152\\
-3.26044	30.8831\\
-3.26044	-30.8831\\
-3.65869	30.7263\\
-3.65869	-30.7263\\
-4.49336	30.8322\\
-4.49336	-30.8322\\
-4.49537	24.6396\\
-4.49537	-24.6396\\
-3.26795	30.2923\\
-3.26795	-30.2923\\
-3.06025	24.5224\\
-3.06025	-24.5224\\
-3.07467	24.7316\\
-3.07467	-24.7316\\
-4.23517	29.7562\\
-4.23517	-29.7562\\
-4.50348	29.2894\\
-4.50348	-29.2894\\
-3.04098	25.4293\\
-3.04098	-25.4293\\
-3.7563	26.4457\\
-3.7563	-26.4457\\
-3.05683	25.6946\\
-3.05683	-25.6946\\
-3.38426	29.9202\\
-3.38426	-29.9202\\
-3.21292	25.7904\\
-3.21292	-25.7904\\
-3.54726	29.7302\\
-3.54726	-29.7302\\
-4.11672	28.9972\\
-4.11672	-28.9972\\
-3.52495	29.6441\\
-3.52495	-29.6441\\
-3.00433	29.8051\\
-3.00433	-29.8051\\
-3.01916	29.6886\\
-3.01916	-29.6886\\
-3.08599	29.6224\\
-3.08599	-29.6224\\
-3.19741	26.2925\\
-3.19741	-26.2925\\
-3.11981	29.5666\\
-3.11981	-29.5666\\
-4.18061	28.1508\\
-4.18061	-28.1508\\
-3.00511	29.52\\
-3.00511	-29.52\\
-3	29.8495\\
-3	-29.8495\\
-3.13505	29.2377\\
-3.13505	-29.2377\\
-3.184	29.1722\\
-3.184	-29.1722\\
-3.044	26.0225\\
-3.044	-26.0225\\
-3.80364	28.8414\\
-3.80364	-28.8414\\
-3.48539	28.9035\\
-3.48539	-28.9035\\
-3.01851	29.07\\
-3.01851	-29.07\\
-3.1596	28.8646\\
-3.1596	-28.8646\\
-3.00041	25.5094\\
-3.00041	-25.5094\\
-3.00037	25.4733\\
-3.00037	-25.4733\\
-3.00003	25.4267\\
-3.00003	-25.4267\\
-3.03904	26.9543\\
-3.03904	-26.9543\\
-3.28624	26.7692\\
-3.28624	-26.7692\\
-3.5844	28.4458\\
-3.5844	-28.4458\\
-3.24757	28.8134\\
-3.24757	-28.8134\\
-3.03299	27.0939\\
-3.03299	-27.0939\\
-3.4694	27.5431\\
-3.4694	-27.5431\\
-3.12723	27.1388\\
-3.12723	-27.1388\\
-3.06357	27.3855\\
-3.06357	-27.3855\\
-3.08144	27.5519\\
-3.08144	-27.5519\\
-3.17459	28.7659\\
-3.17459	-28.7659\\
-3.00019	27.2178\\
-3.00019	-27.2178\\
-3.20997	27.6017\\
-3.20997	-27.6017\\
-3.5357	28.1589\\
-3.5357	-28.1589\\
-3.25168	27.6777\\
-3.25168	-27.6777\\
-3.0904	28.3193\\
-3.0904	-28.3193\\
-3.01902	28.36\\
-3.01902	-28.36\\
-3.24887	27.8211\\
-3.24887	-27.8211\\
-3.07931	28.1291\\
-3.07931	-28.1291\\
-3.04001	28.5114\\
-3.04001	-28.5114\\
-3.06917	27.9921\\
-3.06917	-27.9921\\
-3.28037	28.1206\\
-3.28037	-28.1206\\
-3.01318	28.5214\\
-3.01318	-28.5214\\
-3.00497	28.0461\\
-3.00497	-28.0461\\
-3.04943	28.5127\\
-3.04943	-28.5127\\
-3.00005	27.8981\\
-3.00005	-27.8981\\
-3.00003	25.4264\\
-3.00003	-25.4264\\
};
\addplot [color=mycolor4, only marks, mark size=1.2pt, mark=*, mark options={solid, fill=mycolor4, mycolor4}, forget plot]
  table[row sep=crcr]{%
-13.5721	60.1307\\
-13.5721	-60.1307\\
-13.7438	59.8599\\
-13.7438	-59.8599\\
-13.2509	58.2641\\
-13.2509	-58.2641\\
-14.2205	55.3257\\
-14.2205	-55.3257\\
-13.4579	57.0676\\
-13.4579	-57.0676\\
-12.8524	56.7652\\
-12.8524	-56.7652\\
-11.6928	57.9025\\
-11.6928	-57.9025\\
-10.5729	53.6702\\
-10.5729	-53.6702\\
-11.9851	53.9007\\
-11.9851	-53.9007\\
-10.5698	54.1655\\
-10.5698	-54.1655\\
-10.9904	55.906\\
-10.9904	-55.906\\
-11.2039	56.877\\
-11.2039	-56.877\\
-10.8957	54.0155\\
-10.8957	-54.0155\\
-11.2109	55.8378\\
-11.2109	-55.8378\\
-11.4528	57.4839\\
-11.4528	-57.4839\\
-11.3549	57.2467\\
-11.3549	-57.2467\\
-11.4863	57.5666\\
-11.4863	-57.5666\\
-10.0997	51.0599\\
-10.0997	-51.0599\\
-8.49468	48.4216\\
-8.49468	-48.4216\\
-8.32981	48.8965\\
-8.32981	-48.8965\\
-8.38706	48.2259\\
-8.38706	-48.2259\\
-11.4049	57.3706\\
-11.4049	-57.3706\\
-13.7371	45.6315\\
-13.7371	-45.6315\\
-8.13757	48.7263\\
-8.13757	-48.7263\\
-8.82865	47.8286\\
-8.82865	-47.8286\\
-11.2747	45.0989\\
-11.2747	-45.0989\\
-9.78491	45.8354\\
-9.78491	-45.8354\\
-7.69764	46.4531\\
-7.69764	-46.4531\\
-7.57811	46.03\\
-7.57811	-46.03\\
-7.45438	45.1151\\
-7.45438	-45.1151\\
-7.31151	45.9243\\
-7.31151	-45.9243\\
-7.74878	46.0679\\
-7.74878	-46.0679\\
-9.59641	45.9809\\
-9.59641	-45.9809\\
-13.8041	41.3306\\
-13.8041	-41.3306\\
-7.84582	43.3632\\
-7.84582	-43.3632\\
-7.38334	43.9962\\
-7.38334	-43.9962\\
-6.55305	41.5857\\
-6.55305	-41.5857\\
-6.97589	44.2336\\
-6.97589	-44.2336\\
-6.7416	44.1014\\
-6.7416	-44.1014\\
-6.89637	43.9723\\
-6.89637	-43.9723\\
-6.76915	42.9781\\
-6.76915	-42.9781\\
-8.35842	42.4049\\
-8.35842	-42.4049\\
-7.05516	41.8484\\
-7.05516	-41.8484\\
-6.59687	42.8052\\
-6.59687	-42.8052\\
-6.99478	44.1365\\
-6.99478	-44.1365\\
-14.1682	38.9713\\
-14.1682	-38.9713\\
-9.69373	40.2158\\
-9.69373	-40.2158\\
-5.79521	40.2661\\
-5.79521	-40.2661\\
-7.28294	40.4361\\
-7.28294	-40.4361\\
-7.79601	40.036\\
-7.79601	-40.036\\
-6.47379	39.3328\\
-6.47379	-39.3328\\
-6.23861	38.9764\\
-6.23861	-38.9764\\
-5.13253	37.6769\\
-5.13253	-37.6769\\
-4.39278	35.9612\\
-4.39278	-35.9612\\
-6.00045	35.6909\\
-6.00045	-35.6909\\
-5.33358	36.1783\\
-5.33358	-36.1783\\
-7.80665	34.9438\\
-7.80665	-34.9438\\
-4.75885	34.5442\\
-4.75885	-34.5442\\
-3.96823	10.683\\
-3.96823	-10.683\\
-4.00155	18.242\\
-4.00155	-18.242\\
-12.806	26.5792\\
-12.806	-26.5792\\
-3.48332	0\\
-3.73272	0\\
-4.13246	0\\
-5.1964	0\\
-4.69637	0\\
-4.56314	0\\
-4.2476	0\\
-6.82599	33.7402\\
-6.82599	-33.7402\\
-4.28345	34.0216\\
-4.28345	-34.0216\\
-3.25281	20.4961\\
-3.25281	-20.4961\\
-4.33516	33.7946\\
-4.33516	-33.7946\\
-12.3808	25.2597\\
-12.3808	-25.2597\\
-11.7458	25.8847\\
-11.7458	-25.8847\\
-4.69111	33.1243\\
-4.69111	-33.1243\\
-8.01116	30.6134\\
-8.01116	-30.6134\\
-3.52839	20.6363\\
-3.52839	-20.6363\\
-8.66808	20.1082\\
-8.66808	-20.1082\\
-12.4493	21.519\\
-12.4493	-21.519\\
-6.2647	0\\
-11.7708	22.644\\
-11.7708	-22.644\\
-3.68273	21.5772\\
-3.68273	-21.5772\\
-4.65238	32.588\\
-4.65238	-32.588\\
-3.76451	32.4159\\
-3.76451	-32.4159\\
-3.57573	32.2706\\
-3.57573	-32.2706\\
-5.81561	20.8794\\
-5.81561	-20.8794\\
-3.08913	22.5975\\
-3.08913	-22.5975\\
-5.47416	31.7744\\
-5.47416	-31.7744\\
-3.46746	22.6344\\
-3.46746	-22.6344\\
-6.90568	29.9239\\
-6.90568	-29.9239\\
-4.51995	31.7447\\
-4.51995	-31.7447\\
-4.31128	22.0824\\
-4.31128	-22.0824\\
-7.1951	23.4842\\
-7.1951	-23.4842\\
-14.8137	20.882\\
-14.8137	-20.882\\
-14.6619	0\\
-6.95881	0\\
-4.26828	23.7537\\
-4.26828	-23.7537\\
-10.0092	0\\
-12.6283	0\\
-7.452	0\\
-8.18082	0\\
-10.54	0\\
-13.5436	0\\
-7.53205	0\\
-11.6	0\\
-8.48994	0\\
-13.7768	0\\
-11.2616	0\\
-11.041	0\\
-3.36996	24.6313\\
-3.36996	-24.6313\\
-4.87494	31.1984\\
-4.87494	-31.1984\\
-4.49141	24.629\\
-4.49141	-24.629\\
-6.50992	27.0736\\
-6.50992	-27.0736\\
-3.26059	30.8834\\
-3.26059	-30.8834\\
-6.02407	28.0913\\
-6.02407	-28.0913\\
-3.6587	30.7263\\
-3.6587	-30.7263\\
-4.49368	30.8325\\
-4.49368	-30.8325\\
-3.19318	25.8323\\
-3.19318	-25.8323\\
-4.8359	24.9542\\
-4.8359	-24.9542\\
-3.26858	30.2921\\
-3.26858	-30.2921\\
-3.19549	26.2956\\
-3.19549	-26.2956\\
-3.75721	26.4457\\
-3.75721	-26.4457\\
-4.22107	29.7658\\
-4.22107	-29.7658\\
-3.28412	26.7733\\
-3.28412	-26.7733\\
-3.38133	29.9104\\
-3.38133	-29.9104\\
-4.49156	29.2777\\
-4.49156	-29.2777\\
-3.55941	29.7334\\
-3.55941	-29.7334\\
-3.68539	27.3296\\
-3.68539	-27.3296\\
-4.2468	28.208\\
-4.2468	-28.208\\
-3.2308	27.4279\\
-3.2308	-27.4279\\
-3.38345	29.6792\\
-3.38345	-29.6792\\
-3.07996	27.6234\\
-3.07996	-27.6234\\
-4.12284	28.9533\\
-4.12284	-28.9533\\
-3.12166	27.6573\\
-3.12166	-27.6573\\
-3.13664	29.2345\\
-3.13664	-29.2345\\
-3.34805	27.699\\
-3.34805	-27.699\\
-3.18639	29.162\\
-3.18639	-29.162\\
-3.29818	27.7684\\
-3.29818	-27.7684\\
-3.56014	28.9387\\
-3.56014	-28.9387\\
-3.27933	28.1204\\
-3.27933	-28.1204\\
-3.01909	28.4695\\
-3.01909	-28.4695\\
-3.28398	28.2924\\
-3.28398	-28.2924\\
-3.83913	28.8315\\
-3.83913	-28.8315\\
-3.58518	28.4421\\
-3.58518	-28.4421\\
-3.15294	28.8732\\
-3.15294	-28.8732\\
-3.04708	29.0768\\
-3.04708	-29.0768\\
-3.1859	28.6935\\
-3.1859	-28.6935\\
-3.01852	29.0701\\
-3.01852	-29.0701\\
-3.23747	28.8137\\
-3.23747	-28.8137\\
-7.06106	0\\
-7.59116	0\\
-3	29.8495\\
-3	-29.8495\\
};
\addplot [color=mycolor5, only marks, mark size=1.2pt, mark=*, mark options={solid, fill=mycolor5, mycolor5}, forget plot]
  table[row sep=crcr]{%
-13.573	60.1312\\
-13.573	-60.1312\\
-13.7438	59.8599\\
-13.7438	-59.8599\\
-11.6928	57.9025\\
-11.6928	-57.9025\\
-13.2532	58.2638\\
-13.2532	-58.2638\\
-12.8521	56.7652\\
-12.8521	-56.7652\\
-10.9904	55.9061\\
-10.9904	-55.9061\\
-13.4592	57.0671\\
-13.4592	-57.0671\\
-11.2039	56.877\\
-11.2039	-56.877\\
-10.1	51.0598\\
-10.1	-51.0598\\
-14.2208	55.3256\\
-14.2208	-55.3256\\
-10.5729	53.6701\\
-10.5729	-53.6701\\
-10.5699	54.1655\\
-10.5699	-54.1655\\
-11.9856	53.9017\\
-11.9856	-53.9017\\
-11.2114	55.8381\\
-11.2114	-55.8381\\
-11.3549	57.2467\\
-11.3549	-57.2467\\
-11.4528	57.4839\\
-11.4528	-57.4839\\
-10.8962	54.0156\\
-10.8962	-54.0156\\
-11.4863	57.5666\\
-11.4863	-57.5666\\
-8.33092	48.8961\\
-8.33092	-48.8961\\
-8.49464	48.4216\\
-8.49464	-48.4216\\
-8.38705	48.2259\\
-8.38705	-48.2259\\
-9.78298	45.8365\\
-9.78298	-45.8365\\
-8.82871	47.8285\\
-8.82871	-47.8285\\
-8.13756	48.7263\\
-8.13756	-48.7263\\
-7.70039	46.4524\\
-7.70039	-46.4524\\
-7.57857	46.0301\\
-7.57857	-46.0301\\
-11.4049	57.3706\\
-11.4049	-57.3706\\
-7.31149	45.9243\\
-7.31149	-45.9243\\
-7.74863	46.0678\\
-7.74863	-46.0678\\
-13.7322	45.6294\\
-13.7322	-45.6294\\
-11.2759	45.0971\\
-11.2759	-45.0971\\
-9.59581	45.9813\\
-9.59581	-45.9813\\
-7.46148	45.1139\\
-7.46148	-45.1139\\
-7.8462	43.3616\\
-7.8462	-43.3616\\
-7.38307	43.9963\\
-7.38307	-43.9963\\
-6.9768	44.2334\\
-6.9768	-44.2334\\
-6.74161	44.1014\\
-6.74161	-44.1014\\
-6.77049	42.9782\\
-6.77049	-42.9782\\
-6.89944	43.9731\\
-6.89944	-43.9731\\
-6.99743	44.1364\\
-6.99743	-44.1364\\
-13.8269	41.3358\\
-13.8269	-41.3358\\
-14.2125	39.0112\\
-14.2125	-39.0112\\
-6.6012	42.8046\\
-6.6012	-42.8046\\
-8.55083	42.4538\\
-8.55083	-42.4538\\
-6.55873	41.5887\\
-6.55873	-41.5887\\
-5.80013	40.2634\\
-5.80013	-40.2634\\
-7.28684	40.4367\\
-7.28684	-40.4367\\
-7.05638	41.8498\\
-7.05638	-41.8498\\
-9.67825	40.2144\\
-9.67825	-40.2144\\
-6.25758	38.9477\\
-6.25758	-38.9477\\
-7.78446	40.044\\
-7.78446	-40.044\\
-6.47211	39.3313\\
-6.47211	-39.3313\\
-5.14027	37.6663\\
-5.14027	-37.6663\\
-4.3928	35.9612\\
-4.3928	-35.9612\\
-5.63423	36.1265\\
-5.63423	-36.1265\\
-5.9994	35.6888\\
-5.9994	-35.6888\\
-3.42267	0\\
-7.92475	35.1341\\
-7.92475	-35.1341\\
-5.68798	10.1763\\
-5.68798	-10.1763\\
-4.12617	0\\
-3.72425	0\\
-7.05019	33.9778\\
-7.05019	-33.9778\\
-5.47261	34.4868\\
-5.47261	-34.4868\\
-4.28085	34.0229\\
-4.28085	-34.0229\\
-4.34659	33.8139\\
-4.34659	-33.8139\\
-6.02517	17.7884\\
-6.02517	-17.7884\\
-4.10378	21.403\\
-4.10378	-21.403\\
-3.76432	32.4167\\
-3.76432	-32.4167\\
-3.85665	31.9675\\
-3.85665	-31.9675\\
-4.68489	33.1308\\
-4.68489	-33.1308\\
-4.6535	0\\
-4.23557	0\\
-4.63819	32.5858\\
-4.63819	-32.5858\\
-5.17173	0\\
-4.56065	0\\
-6.25001	0\\
-14.9733	19.8009\\
-14.9733	-19.8009\\
-14.5396	0\\
-9.75944	0\\
-12.5452	0\\
-6.93439	0\\
-8.16959	0\\
-7.31483	0\\
-7.44223	0\\
-3.26801	24.7029\\
-3.26801	-24.7029\\
-12.6423	26.7325\\
-12.6423	-26.7325\\
-7.85046	19.0783\\
-7.85046	-19.0783\\
-8.91574	18.876\\
-8.91574	-18.876\\
-14.1881	22.6871\\
-14.1881	-22.6871\\
-10.5292	0\\
-11.5405	0\\
-3.17873	26.307\\
-3.17873	-26.307\\
-7.8321	30.961\\
-7.8321	-30.961\\
-5.47508	31.7852\\
-5.47508	-31.7852\\
-5.43488	31.6316\\
-5.43488	-31.6316\\
-6.18595	20.8695\\
-6.18595	-20.8695\\
-14.8931	20.9317\\
-14.8931	-20.9317\\
-13.5241	0\\
-11.4569	19.8422\\
-11.4569	-19.8422\\
-11.7322	25.8374\\
-11.7322	-25.8374\\
-11.4435	26.0016\\
-11.4435	-26.0016\\
-11.0302	0\\
-4.80326	31.654\\
-4.80326	-31.654\\
-8.67987	20.5203\\
-8.67987	-20.5203\\
-3.24777	27.5148\\
-3.24777	-27.5148\\
-5.99181	22.1788\\
-5.99181	-22.1788\\
-9.60298	28.3074\\
-9.60298	-28.3074\\
-9.83519	27.4503\\
-9.83519	-27.4503\\
-3.32633	27.8119\\
-3.32633	-27.8119\\
-6.83829	29.7498\\
-6.83829	-29.7498\\
-6.53959	21.6387\\
-6.53959	-21.6387\\
-9.19649	28.021\\
-9.19649	-28.021\\
-7.64989	29.2076\\
-7.64989	-29.2076\\
-13.7252	0\\
-11.178	0\\
-8.40029	0\\
-12.1916	21.8313\\
-12.1916	-21.8313\\
-11.3976	22.7956\\
-11.3976	-22.7956\\
-3.49722	29.7472\\
-3.49722	-29.7472\\
-5.93672	23.6609\\
-5.93672	-23.6609\\
-3.51885	28.2944\\
-3.51885	-28.2944\\
-14.8137	20.882\\
-14.8137	-20.882\\
-3.70159	29.7628\\
-3.70159	-29.7628\\
-3.99141	30.2646\\
-3.99141	-30.2646\\
-6.26103	24.3666\\
-6.26103	-24.3666\\
-7.57461	23.2211\\
-7.57461	-23.2211\\
-5.35362	30.3159\\
-5.35362	-30.3159\\
-4.34225	30.3652\\
-4.34225	-30.3652\\
-8.74046	27.3009\\
-8.74046	-27.3009\\
-9.65344	24.0703\\
-9.65344	-24.0703\\
-3.99591	28.4599\\
-3.99591	-28.4599\\
-3.77594	28.7718\\
-3.77594	-28.7718\\
-4.55013	30.7266\\
-4.55013	-30.7266\\
-4.30481	28.9603\\
-4.30481	-28.9603\\
-7.85156	27.9833\\
-7.85156	-27.9833\\
-8.32194	25.3601\\
-8.32194	-25.3601\\
-7.48442	28.2024\\
-7.48442	-28.2024\\
-7.59116	0\\
-4.6908	29.5476\\
-4.6908	-29.5476\\
-4.43755	29.2459\\
-4.43755	-29.2459\\
-8.04289	24.8718\\
-8.04289	-24.8718\\
-5.60487	27.9361\\
-5.60487	-27.9361\\
-7.52267	26.0369\\
-7.52267	-26.0369\\
-5.93482	28.02\\
-5.93482	-28.02\\
-7.24885	28.3305\\
-7.24885	-28.3305\\
-6.43519	28.1586\\
-6.43519	-28.1586\\
-6.40207	27.0085\\
-6.40207	-27.0085\\
-7.53274	26.8266\\
-7.53274	-26.8266\\
-7.69246	27.1081\\
-7.69246	-27.1081\\
-7.55685	27.3492\\
-7.55685	-27.3492\\
-6.89644	27.394\\
-6.89644	-27.394\\
-7.22773	27.7622\\
-7.22773	-27.7622\\
-7.12743	28.295\\
-7.12743	-28.295\\
-7.06106	0\\
};
\addplot [color=mycolor6, only marks, mark size=1.2pt, mark=*, mark options={solid, fill=mycolor6, mycolor6}, forget plot]
  table[row sep=crcr]{%
-13.5758	60.1332\\
-13.5758	-60.1332\\
-13.749	59.8613\\
-13.749	-59.8613\\
-13.2608	58.2628\\
-13.2608	-58.2628\\
-10.5699	54.1652\\
-10.5699	-54.1652\\
-13.4594	57.067\\
-13.4594	-57.067\\
-12.8588	56.7699\\
-12.8588	-56.7699\\
-11.693	57.9026\\
-11.693	-57.9026\\
-10.9904	55.9061\\
-10.9904	-55.9061\\
-11.2039	56.877\\
-11.2039	-56.877\\
-14.2212	55.3254\\
-14.2212	-55.3254\\
-11.4528	57.4839\\
-11.4528	-57.4839\\
-11.3549	57.2467\\
-11.3549	-57.2467\\
-11.2114	55.8381\\
-11.2114	-55.8381\\
-11.4863	57.5666\\
-11.4863	-57.5666\\
-11.9861	53.9034\\
-11.9861	-53.9034\\
-10.5751	53.6702\\
-10.5751	-53.6702\\
-10.8966	54.016\\
-10.8966	-54.016\\
-10.1006	51.0593\\
-10.1006	-51.0593\\
-11.4049	57.3706\\
-11.4049	-57.3706\\
-8.33638	48.8943\\
-8.33638	-48.8943\\
-8.4952	48.4218\\
-8.4952	-48.4218\\
-8.38751	48.2255\\
-8.38751	-48.2255\\
-8.13785	48.7262\\
-8.13785	-48.7262\\
-13.7309	45.6337\\
-13.7309	-45.6337\\
-8.8293	47.8292\\
-8.8293	-47.8292\\
-11.2389	45.1058\\
-11.2389	-45.1058\\
-14.2499	39.0887\\
-14.2499	-39.0887\\
-13.8263	41.3357\\
-13.8263	-41.3357\\
-9.81361	45.8855\\
-9.81361	-45.8855\\
-7.70412	46.4523\\
-7.70412	-46.4523\\
-9.59512	45.9815\\
-9.59512	-45.9815\\
-7.58004	46.0303\\
-7.58004	-46.0303\\
-7.74861	46.0681\\
-7.74861	-46.0681\\
-7.31352	45.9238\\
-7.31352	-45.9238\\
-7.50844	45.1071\\
-7.50844	-45.1071\\
-5.88476	40.2086\\
-5.88476	-40.2086\\
-6.34232	38.912\\
-6.34232	-38.912\\
-9.68043	40.2149\\
-9.68043	-40.2149\\
-6.55811	41.5868\\
-6.55811	-41.5868\\
-7.82144	43.355\\
-7.82144	-43.355\\
-6.4696	39.3297\\
-6.4696	-39.3297\\
-7.38205	43.9996\\
-7.38205	-43.9996\\
-6.7417	44.1014\\
-6.7417	-44.1014\\
-6.98333	44.2324\\
-6.98333	-44.2324\\
-7.26665	40.3515\\
-7.26665	-40.3515\\
-6.89933	43.9733\\
-6.89933	-43.9733\\
-7.97492	42.0652\\
-7.97492	-42.0652\\
-6.60514	42.8056\\
-6.60514	-42.8056\\
-6.7098	42.8853\\
-6.7098	-42.8853\\
-7.76823	40.0558\\
-7.76823	-40.0558\\
-8.7284	42.5357\\
-8.7284	-42.5357\\
-7.01625	44.134\\
-7.01625	-44.134\\
-6.93462	37.9268\\
-6.93462	-37.9268\\
-4.12174	0\\
-3.37648	0\\
-3.71963	0\\
-5.96548	35.7184\\
-5.96548	-35.7184\\
-5.63007	36.1443\\
-5.63007	-36.1443\\
-7.13154	9.03302\\
-7.13154	-9.03302\\
-5.15719	0\\
-4.63185	0\\
-4.22829	0\\
-8.10193	35.777\\
-8.10193	-35.777\\
-5.50935	34.4691\\
-5.50935	-34.4691\\
-4.55819	0\\
-7.20674	35.5091\\
-7.20674	-35.5091\\
-6.24055	0\\
-8.04359	34.2825\\
-8.04359	-34.2825\\
-6.2279	33.2213\\
-6.2279	-33.2213\\
-6.44764	17.6435\\
-6.44764	-17.6435\\
-7.18299	33.31\\
-7.18299	-33.31\\
-7.13242	33.537\\
-7.13242	-33.537\\
-14.4494	0\\
-6.91924	0\\
-9.51751	0\\
-8.1552	0\\
-7.23682	0\\
-7.39534	0\\
-8.35425	0\\
-7.70947	32.0971\\
-7.70947	-32.0971\\
-6.59223	31.9791\\
-6.59223	-31.9791\\
-12.6783	26.8096\\
-12.6783	-26.8096\\
-6.4672	31.8115\\
-6.4672	-31.8115\\
-8.05729	19.2675\\
-8.05729	-19.2675\\
-8.85154	18.8253\\
-8.85154	-18.8253\\
-7.80211	30.972\\
-7.80211	-30.972\\
-7.40086	31.2284\\
-7.40086	-31.2284\\
-6.54639	20.989\\
-6.54639	-20.989\\
-14.7631	21.9479\\
-14.7631	-21.9479\\
-7.18384	30.9109\\
-7.18384	-30.9109\\
-13.5061	0\\
-7.11685	20.2914\\
-7.11685	-20.2914\\
-12.5033	0\\
-10.5241	0\\
-11.5084	0\\
-14.7986	20.7841\\
-14.7986	-20.7841\\
-13.689	0\\
-9.48061	20.122\\
-9.48061	-20.122\\
-11.4361	19.8214\\
-11.4361	-19.8214\\
-11.1224	0\\
-6.71011	31.2615\\
-6.71011	-31.2615\\
-6.33614	22.0877\\
-6.33614	-22.0877\\
-11.023	0\\
-7.14364	21.3096\\
-7.14364	-21.3096\\
-6.15634	30.8077\\
-6.15634	-30.8077\\
-11.6625	26.0142\\
-11.6625	-26.0142\\
-11.6373	25.8053\\
-11.6373	-25.8053\\
-6.25918	30.4267\\
-6.25918	-30.4267\\
-12.1025	21.6618\\
-12.1025	-21.6618\\
-6.7352	29.9438\\
-6.7352	-29.9438\\
-6.24385	24.1675\\
-6.24385	-24.1675\\
-6.9177	30.2353\\
-6.9177	-30.2353\\
-11.4651	22.6216\\
-11.4651	-22.6216\\
-7.58566	23.1262\\
-7.58566	-23.1262\\
-7.06687	23.3045\\
-7.06687	-23.3045\\
-9.60156	28.3842\\
-9.60156	-28.3842\\
-5.97876	25.8653\\
-5.97876	-25.8653\\
-9.86976	27.4468\\
-9.86976	-27.4468\\
-7.72415	24.0821\\
-7.72415	-24.0821\\
-6.64758	29.8566\\
-6.64758	-29.8566\\
-7.06106	0\\
-9.63799	23.9868\\
-9.63799	-23.9868\\
-9.27218	27.9991\\
-9.27218	-27.9991\\
-6.04166	29.3862\\
-6.04166	-29.3862\\
-6.49992	29.3704\\
-6.49992	-29.3704\\
-8.78161	27.1624\\
-8.78161	-27.1624\\
-6.27086	29.328\\
-6.27086	-29.328\\
-8.2301	25.0554\\
-8.2301	-25.0554\\
-7.64782	28.6355\\
-7.64782	-28.6355\\
-7.98067	24.8729\\
-7.98067	-24.8729\\
-7.62357	25.8808\\
-7.62357	-25.8808\\
-6.00721	27.9376\\
-6.00721	-27.9376\\
-7.90837	27.9742\\
-7.90837	-27.9742\\
-6.15044	27.0933\\
-6.15044	-27.0933\\
-7.7986	26.5666\\
-7.7986	-26.5666\\
-6.14885	28.2159\\
-6.14885	-28.2159\\
-6.17772	27.5663\\
-6.17772	-27.5663\\
-6.47694	28.3854\\
-6.47694	-28.3854\\
-6.25439	27.7581\\
-6.25439	-27.7581\\
-6.45488	27.067\\
-6.45488	-27.067\\
-7.56431	26.8505\\
-7.56431	-26.8505\\
-6.44813	27.2649\\
-6.44813	-27.2649\\
-7.18226	28.4685\\
-7.18226	-28.4685\\
-7.03994	28.5541\\
-7.03994	-28.5541\\
-7.59841	27.2337\\
-7.59841	-27.2337\\
-7.32083	27.5354\\
-7.32083	-27.5354\\
-7.32558	28.3182\\
-7.32558	-28.3182\\
-7.4182	28.1145\\
-7.4182	-28.1145\\
-7.31364	28.2478\\
-7.31364	-28.2478\\
-7.17731	27.7868\\
-7.17731	-27.7868\\
-7.11819	28.3068\\
-7.11819	-28.3068\\
-7.59116	0\\
-14.8137	20.882\\
-14.8137	-20.882\\
};
\addplot [color=mycolor7, only marks, mark size=1.2pt, mark=*, mark options={solid, fill=mycolor7, mycolor7}, forget plot]
  table[row sep=crcr]{%
-13.5945	60.142\\
-13.5945	-60.142\\
-13.8606	59.8929\\
-13.8606	-59.8929\\
-13.3419	58.2856\\
-13.3419	-58.2856\\
-12.8608	56.7747\\
-12.8608	-56.7747\\
-13.4594	57.0651\\
-13.4594	-57.0651\\
-11.7017	57.9031\\
-11.7017	-57.9031\\
-10.9904	55.9061\\
-10.9904	-55.9061\\
-11.2039	56.877\\
-11.2039	-56.877\\
-11.3549	57.2467\\
-11.3549	-57.2467\\
-11.4528	57.4839\\
-11.4528	-57.4839\\
-11.4863	57.5666\\
-11.4863	-57.5666\\
-14.2225	55.3236\\
-14.2225	-55.3236\\
-11.2114	55.8381\\
-11.2114	-55.8381\\
-10.5593	54.1446\\
-10.5593	-54.1446\\
-10.7924	53.6644\\
-10.7924	-53.6644\\
-11.9861	53.9037\\
-11.9861	-53.9037\\
-11.0169	54.0613\\
-11.0169	-54.0613\\
-10.3371	51.167\\
-10.3371	-51.167\\
-8.39092	48.8853\\
-8.39092	-48.8853\\
-8.5963	48.4471\\
-8.5963	-48.4471\\
-11.4049	57.3706\\
-11.4049	-57.3706\\
-8.13788	48.7262\\
-8.13788	-48.7262\\
-9.19808	48.2959\\
-9.19808	-48.2959\\
-13.7312	45.6343\\
-13.7312	-45.6343\\
-9.00423	47.8368\\
-9.00423	-47.8368\\
-7.86432	46.0985\\
-7.86432	-46.0985\\
-8.32572	46.5511\\
-8.32572	-46.5511\\
-9.94875	46.3409\\
-9.94875	-46.3409\\
-14.417	39.4018\\
-14.417	-39.4018\\
-13.8333	41.3274\\
-13.8333	-41.3274\\
-11.0843	44.9721\\
-11.0843	-44.9721\\
-9.58871	45.9785\\
-9.58871	-45.9785\\
-8.2732	45.1123\\
-8.2732	-45.1123\\
-10.2982	45.4697\\
-10.2982	-45.4697\\
-13.4359	39.4867\\
-13.4359	-39.4867\\
-10.1255	45.3936\\
-10.1255	-45.3936\\
-7.89271	43.8607\\
-7.89271	-43.8607\\
-7.99972	44.041\\
-7.99972	-44.041\\
-10.7969	42.705\\
-10.7969	-42.705\\
-8.54154	39.8042\\
-8.54154	-39.8042\\
-10.7116	39.4839\\
-10.7116	-39.4839\\
-9.99263	39.7848\\
-9.99263	-39.7848\\
-9.95342	43.2805\\
-9.95342	-43.2805\\
-8.50641	41.9254\\
-8.50641	-41.9254\\
-9.76402	43.7827\\
-9.76402	-43.7827\\
-9.34052	41.05\\
-9.34052	-41.05\\
-8.72917	42.5355\\
-8.72917	-42.5355\\
-9.54339	43.5781\\
-9.54339	-43.5781\\
-9.03808	42.5729\\
-9.03808	-42.5729\\
-9.30859	42.3285\\
-9.30859	-42.3285\\
-8.85893	38.6619\\
-8.85893	-38.6619\\
-9.18356	38.8324\\
-9.18356	-38.8324\\
-7.88699	37.2425\\
-7.88699	-37.2425\\
-3.32252	0\\
-3.71739	0\\
-5.1472	0\\
-4.11671	0\\
-4.62281	0\\
-4.22437	0\\
-4.55666	0\\
-7.14935	8.97419\\
-7.14935	-8.97419\\
-7.20548	35.5062\\
-7.20548	-35.5062\\
-8.39121	35.7084\\
-8.39121	-35.7084\\
-8.83966	35.1141\\
-8.83966	-35.1141\\
-6.23039	0\\
-9.79021	34.1303\\
-9.79021	-34.1303\\
-7.66998	34.0583\\
-7.66998	-34.0583\\
-6.90571	0\\
-9.23826	33.6059\\
-9.23826	-33.6059\\
-7.08122	33.5443\\
-7.08122	-33.5443\\
-7.18817	33.3\\
-7.18817	-33.3\\
-8.32641	0\\
-7.36775	0\\
-7.1806	0\\
-7.44682	32.6449\\
-7.44682	-32.6449\\
-7.01886	17.4169\\
-7.01886	-17.4169\\
-9.32102	0\\
-8.1396	0\\
-7.4877	32.1333\\
-7.4877	-32.1333\\
-6.58744	31.9784\\
-6.58744	-31.9784\\
-6.46743	31.8117\\
-6.46743	-31.8117\\
-8.93072	18.7761\\
-8.93072	-18.7761\\
-7.93797	19.4738\\
-7.93797	-19.4738\\
-14.3794	0\\
-12.4596	0\\
-10.516	0\\
-8.29195	31.2215\\
-8.29195	-31.2215\\
-7.3729	31.2374\\
-7.3729	-31.2374\\
-10.872	29.5751\\
-10.872	-29.5751\\
-7.65768	30.8668\\
-7.65768	-30.8668\\
-12.7052	26.8393\\
-12.7052	-26.8393\\
-9.73636	29.3232\\
-9.73636	-29.3232\\
-6.59127	20.9315\\
-6.59127	-20.9315\\
-14.8137	21.8121\\
-14.8137	-21.8121\\
-6.41963	22.0399\\
-6.41963	-22.0399\\
-10.6817	19.0919\\
-10.6817	-19.0919\\
-8.03896	20.0213\\
-8.03896	-20.0213\\
-14.7658	20.7162\\
-14.7658	-20.7162\\
-7.13915	21.3028\\
-7.13915	-21.3028\\
-13.4858	0\\
-11.4698	0\\
-11.0913	0\\
-13.664	0\\
-11.0148	0\\
-11.4314	19.8201\\
-11.4314	-19.8201\\
-6.26034	30.4278\\
-6.26034	-30.4278\\
-7.06795	30.4851\\
-7.06795	-30.4851\\
-11.7217	26.0044\\
-11.7217	-26.0044\\
-11.6726	25.7685\\
-11.6726	-25.7685\\
-6.25009	24.1679\\
-6.25009	-24.1679\\
-6.91427	30.2363\\
-6.91427	-30.2363\\
-7.12047	23.3234\\
-7.12047	-23.3234\\
-12.1091	21.5944\\
-12.1091	-21.5944\\
-9.18234	22.6625\\
-9.18234	-22.6625\\
-6.62456	29.9537\\
-6.62456	-29.9537\\
-5.99299	25.8402\\
-5.99299	-25.8402\\
-11.4996	22.6072\\
-11.4996	-22.6072\\
-7.64968	23.9346\\
-7.64968	-23.9346\\
-6.0427	29.3867\\
-6.0427	-29.3867\\
-9.61757	28.4561\\
-9.61757	-28.4561\\
-9.88301	27.4505\\
-9.88301	-27.4505\\
-9.79902	23.9209\\
-9.79902	-23.9209\\
-9.31268	27.9558\\
-9.31268	-27.9558\\
-6.30875	29.3109\\
-6.30875	-29.3109\\
-8.70233	27.4636\\
-8.70233	-27.4636\\
-8.88239	27.1514\\
-8.88239	-27.1514\\
-8.00147	28.5774\\
-8.00147	-28.5774\\
-6.5298	29.3703\\
-6.5298	-29.3703\\
-7.59116	0\\
-9.02245	25.9785\\
-9.02245	-25.9785\\
-8.3343	25.5495\\
-8.3343	-25.5495\\
-8.51191	25.5422\\
-8.51191	-25.5422\\
-6.14247	27.0978\\
-6.14247	-27.0978\\
-7.70463	25.8868\\
-7.70463	-25.8868\\
-6.14585	28.2192\\
-6.14585	-28.2192\\
-6.18371	27.321\\
-6.18371	-27.321\\
-6.27484	27.7785\\
-6.27484	-27.7785\\
-6.477	27.1043\\
-6.477	-27.1043\\
-6.43767	27.195\\
-6.43767	-27.195\\
-6.56592	28.3584\\
-6.56592	-28.3584\\
-7.98058	27.9589\\
-7.98058	-27.9589\\
-7.61619	26.736\\
-7.61619	-26.736\\
-7.08586	28.6665\\
-7.08586	-28.6665\\
-8.03154	27.2376\\
-8.03154	-27.2376\\
-7.24171	28.534\\
-7.24171	-28.534\\
-7.45386	28.2686\\
-7.45386	-28.2686\\
-7.77081	27.9229\\
-7.77081	-27.9229\\
-7.19162	27.6877\\
-7.19162	-27.6877\\
-7.41899	28.053\\
-7.41899	-28.053\\
-7.15814	28.3036\\
-7.15814	-28.3036\\
-7.31741	28.2638\\
-7.31741	-28.2638\\
-7.06106	0\\
-14.8137	20.882\\
-14.8137	-20.882\\
};
\addplot [color=mycolor8, only marks, mark size=1.2pt, mark=*, mark options={solid, fill=mycolor8, mycolor8}, forget plot]
  table[row sep=crcr]{%
-13.6791	60.1853\\
-13.6791	-60.1853\\
-13.9116	59.8493\\
-13.9116	-59.8493\\
-13.7173	58.424\\
-13.7173	-58.424\\
-13.0434	56.9518\\
-13.0434	-56.9518\\
-12.6701	57.5641\\
-12.6701	-57.5641\\
-12.2702	57.2294\\
-12.2702	-57.2294\\
-11.4863	57.5666\\
-11.4863	-57.5666\\
-12.2225	55.7321\\
-12.2225	-55.7321\\
-14.1819	56.6335\\
-14.1819	-56.6335\\
-12.8833	53.3917\\
-12.8833	-53.3917\\
-13.306	53.6296\\
-13.306	-53.6296\\
-14.2301	55.3199\\
-14.2301	-55.3199\\
-12.7645	50.4072\\
-12.7645	-50.4072\\
-14.0248	56.2474\\
-14.0248	-56.2474\\
-13.9945	55.1865\\
-13.9945	-55.1865\\
-14.1767	53.0205\\
-14.1767	-53.0205\\
-13.7284	53.38\\
-13.7284	-53.38\\
-14.2264	56.7368\\
-14.2264	-56.7368\\
-11.1219	48.3373\\
-11.1219	-48.3373\\
-11.3184	47.8519\\
-11.3184	-47.8519\\
-10.9588	48.1701\\
-10.9588	-48.1701\\
-11.1186	47.6813\\
-11.1186	-47.6813\\
-11.6333	47.1978\\
-11.6333	-47.1978\\
-14.1115	44.4549\\
-14.1115	-44.4549\\
-12.533	45.1753\\
-12.533	-45.1753\\
-13.8656	45.2282\\
-13.8656	-45.2282\\
-10.5198	45.8892\\
-10.5198	-45.8892\\
-10.3997	45.4769\\
-10.3997	-45.4769\\
-10.1251	45.3935\\
-10.1251	-45.3935\\
-10.5584	45.5367\\
-10.5584	-45.5367\\
-12.4452	45.4993\\
-12.4452	-45.4993\\
-14.8359	41.3168\\
-14.8359	-41.3168\\
-10.2152	44.5665\\
-10.2152	-44.5665\\
-9.36081	41.0581\\
-9.36081	-41.0581\\
-10.6354	42.8819\\
-10.6354	-42.8819\\
-10.172	43.4371\\
-10.172	-43.4371\\
-9.78158	43.7029\\
-9.78158	-43.7029\\
-9.56214	43.5781\\
-9.56214	-43.5781\\
-9.41828	42.2716\\
-9.41828	-42.2716\\
-9.56033	42.4728\\
-9.56033	-42.4728\\
-9.7293	43.4545\\
-9.7293	-43.4545\\
-11.1717	41.7834\\
-11.1717	-41.7834\\
-9.81045	43.5966\\
-9.81045	-43.5966\\
-14.4914	39.5507\\
-14.4914	-39.5507\\
-13.4637	39.5074\\
-13.4637	-39.5074\\
-9.78066	41.2292\\
-9.78066	-41.2292\\
-8.61288	39.774\\
-8.61288	-39.774\\
-10.0429	39.7981\\
-10.0429	-39.7981\\
-10.7296	39.4641\\
-10.7296	-39.4641\\
-9.29394	38.85\\
-9.29394	-38.85\\
-8.994	38.4256\\
-8.994	-38.4256\\
-3.24586	0\\
-7.93322	37.1964\\
-7.93322	-37.1964\\
-7.21336	35.5032\\
-7.21336	-35.5032\\
-8.38922	35.7021\\
-8.38922	-35.7021\\
-8.83798	35.1463\\
-8.83798	-35.1463\\
-7.66571	34.0511\\
-7.66571	-34.0511\\
-10.5208	33.9846\\
-10.5208	-33.9846\\
-7.10155	33.5445\\
-7.10155	-33.5445\\
-9.40994	33.4634\\
-9.40994	-33.4634\\
-7.17137	8.938\\
-7.17137	-8.938\\
-3.71393	0\\
-4.1135	0\\
-4.22042	0\\
-4.61253	0\\
-4.55416	0\\
-5.14123	0\\
-6.22247	0\\
-7.34465	0\\
-6.8879	0\\
-7.01221	0\\
-7.18291	33.3017\\
-7.18291	-33.3017\\
-7.01791	17.3657\\
-7.01791	-17.3657\\
-9.11005	0\\
-8.1163	0\\
-8.30786	0\\
-7.44576	32.6517\\
-7.44576	-32.6517\\
-7.48439	32.1752\\
-7.48439	-32.1752\\
-6.58868	31.9757\\
-6.58868	-31.9757\\
-6.46759	31.8117\\
-6.46759	-31.8117\\
-6.58619	20.9072\\
-6.58619	-20.9072\\
-7.91484	19.4437\\
-7.91484	-19.4437\\
-8.94426	18.7909\\
-8.94426	-18.7909\\
-14.33	0\\
-12.4146	0\\
-10.5081	0\\
-7.36123	31.2538\\
-7.36123	-31.2538\\
-8.31988	31.2023\\
-8.31988	-31.2023\\
-10.9052	29.5809\\
-10.9052	-29.5809\\
-7.84628	30.8485\\
-7.84628	-30.8485\\
-12.7258	26.8494\\
-12.7258	-26.8494\\
-6.42573	22.012\\
-6.42573	-22.012\\
-11.2996	18.8871\\
-11.2996	-18.8871\\
-9.69206	29.3013\\
-9.69206	-29.3013\\
-14.8409	21.7726\\
-14.8409	-21.7726\\
-11.4287	0\\
-8.78268	19.8524\\
-8.78268	-19.8524\\
-13.4675	0\\
-7.52885	21.3323\\
-7.52885	-21.3323\\
-11.0717	0\\
-13.6437	0\\
-11.0035	0\\
-14.7722	20.6779\\
-14.7722	-20.6779\\
-11.4257	19.8131\\
-11.4257	-19.8131\\
-6.26025	30.4278\\
-6.26025	-30.4278\\
-6.25167	24.1729\\
-6.25167	-24.1729\\
-7.15101	23.3351\\
-7.15101	-23.3351\\
-11.7344	25.9882\\
-11.7344	-25.9882\\
-6.00003	25.8324\\
-6.00003	-25.8324\\
-6.61563	29.9889\\
-6.61563	-29.9889\\
-6.91416	30.2369\\
-6.91416	-30.2369\\
-11.6886	25.7543\\
-11.6886	-25.7543\\
-12.1174	21.5657\\
-12.1174	-21.5657\\
-7.3317	30.2851\\
-7.3317	-30.2851\\
-9.64987	28.4651\\
-9.64987	-28.4651\\
-10.1421	22.4436\\
-10.1421	-22.4436\\
-7.66369	23.9363\\
-7.66369	-23.9363\\
-11.5198	22.6074\\
-11.5198	-22.6074\\
-6.04299	29.3867\\
-6.04299	-29.3867\\
-9.86256	27.4559\\
-9.86256	-27.4559\\
-9.80017	23.889\\
-9.80017	-23.889\\
-6.31437	29.3042\\
-6.31437	-29.3042\\
-9.31568	27.9326\\
-9.31568	-27.9326\\
-6.14564	27.095\\
-6.14564	-27.095\\
-8.09815	28.6634\\
-8.09815	-28.6634\\
-9.36371	25.9863\\
-9.36371	-25.9863\\
-8.72885	27.427\\
-8.72885	-27.427\\
-8.90865	27.1364\\
-8.90865	-27.1364\\
-7.59116	0\\
-14.8137	20.882\\
-14.8137	-20.882\\
-6.53683	29.3643\\
-6.53683	-29.3643\\
-8.36915	25.5125\\
-8.36915	-25.5125\\
-8.38483	25.561\\
-8.38483	-25.561\\
-7.75445	25.8867\\
-7.75445	-25.8867\\
-6.20101	27.337\\
-6.20101	-27.337\\
-6.27649	27.7796\\
-6.27649	-27.7796\\
-6.16785	28.2254\\
-6.16785	-28.2254\\
-8.04433	27.9708\\
-8.04433	-27.9708\\
-7.91958	27.9993\\
-7.91958	-27.9993\\
-6.57368	28.3283\\
-6.57368	-28.3283\\
-6.43647	27.1353\\
-6.43647	-27.1353\\
-7.6366	26.7268\\
-7.6366	-26.7268\\
-6.49879	27.1412\\
-6.49879	-27.1412\\
-7.11133	28.6665\\
-7.11133	-28.6665\\
-8.13571	27.2068\\
-8.13571	-27.2068\\
-7.25931	28.5387\\
-7.25931	-28.5387\\
-7.48459	28.286\\
-7.48459	-28.286\\
-7.194	27.6851\\
-7.194	-27.6851\\
-7.4713	27.9383\\
-7.4713	-27.9383\\
-7.15874	28.3033\\
-7.15874	-28.3033\\
-7.30772	28.2655\\
-7.30772	-28.2655\\
-7.06106	0\\
};
\addplot [color=mycolor9, only marks, mark size=1.2pt, mark=*, mark options={solid, fill=mycolor9, mycolor9}, forget plot]
  table[row sep=crcr]{%
-14.513	57.261\\
-14.513	-57.261\\
-12.8949	50.4016\\
-12.8949	-50.4016\\
-13.3743	53.5456\\
-13.3743	-53.5456\\
-13.3815	53.0352\\
-13.3815	-53.0352\\
-13.8076	55.2817\\
-13.8076	-55.2817\\
-14.0253	56.2478\\
-14.0253	-56.2478\\
-14.004	55.1788\\
-14.004	-55.1788\\
-14.1763	56.6141\\
-14.1763	-56.6141\\
-14.2742	56.849\\
-14.2742	-56.849\\
-14.8646	53.1911\\
-14.8646	-53.1911\\
-13.7168	53.3835\\
-13.7168	-53.3835\\
-14.3077	56.931\\
-14.3077	-56.931\\
-11.1472	48.3312\\
-11.1472	-48.3312\\
-11.3231	47.8543\\
-11.3231	-47.8543\\
-11.1965	47.6522\\
-11.1965	-47.6522\\
-10.959	48.1702\\
-10.959	-48.1702\\
-11.6343	47.2005\\
-11.6343	-47.2005\\
-12.8457	45.2932\\
-12.8457	-45.2932\\
-10.5188	45.8907\\
-10.5188	-45.8907\\
-14.2264	56.7368\\
-14.2264	-56.7368\\
-10.4035	45.476\\
-10.4035	-45.476\\
-12.4128	45.4538\\
-12.4128	-45.4538\\
-10.1301	45.3867\\
-10.1301	-45.3867\\
-10.5641	45.5258\\
-10.5641	-45.5258\\
-14.0057	44.6279\\
-14.0057	-44.6279\\
-10.2696	44.5578\\
-10.2696	-44.5578\\
-10.5846	42.8706\\
-10.5846	-42.8706\\
-10.2218	43.442\\
-10.2218	-43.442\\
-9.7898	43.6986\\
-9.7898	-43.6986\\
-9.56205	43.5779\\
-9.56205	-43.5779\\
-9.41833	42.2716\\
-9.41833	-42.2716\\
-9.564	42.4586\\
-9.564	-42.4586\\
-9.71292	43.4434\\
-9.71292	-43.4434\\
-9.80857	43.5985\\
-9.80857	-43.5985\\
-13.1197	39.8555\\
-13.1197	-39.8555\\
-11.1757	41.7828\\
-11.1757	-41.7828\\
-9.37121	41.0555\\
-9.37121	-41.0555\\
-9.85229	41.2388\\
-9.85229	-41.2388\\
-8.62008	39.7764\\
-8.62008	-39.7764\\
-10.0872	39.8003\\
-10.0872	-39.8003\\
-10.7062	39.4882\\
-10.7062	-39.4882\\
-8.99721	38.4264\\
-8.99721	-38.4264\\
-9.28987	38.8328\\
-9.28987	-38.8328\\
-7.94877	37.1915\\
-7.94877	-37.1915\\
-7.21354	35.5031\\
-7.21354	-35.5031\\
-8.38318	35.7049\\
-8.38318	-35.7049\\
-8.83016	35.1572\\
-8.83016	-35.1572\\
-10.744	34.1372\\
-10.744	-34.1372\\
-7.66332	34.0512\\
-7.66332	-34.0512\\
-3.15918	0\\
-7.1687	33.3306\\
-7.1687	-33.3306\\
-9.66196	33.2091\\
-9.66196	-33.2091\\
-7.1094	33.5497\\
-7.1094	-33.5497\\
-3.71095	0\\
-7.19553	8.92099\\
-7.19553	-8.92099\\
-4.21716	0\\
-4.1105	0\\
-4.60446	0\\
-4.5512	0\\
-5.1375	0\\
-6.21286	0\\
-6.7799	0\\
-6.58882	31.9743\\
-6.58882	-31.9743\\
-7.01845	17.3356\\
-7.01845	-17.3356\\
-8.92243	0\\
-8.08926	0\\
-7.32239	0\\
-6.88552	0\\
-8.29585	0\\
-7.44413	32.6554\\
-7.44413	-32.6554\\
-7.48716	32.1981\\
-7.48716	-32.1981\\
-6.46769	31.8117\\
-6.46769	-31.8117\\
-12.742	26.8538\\
-12.742	-26.8538\\
-8.94765	18.7932\\
-8.94765	-18.7932\\
-7.91475	19.4363\\
-7.91475	-19.4363\\
-6.58443	20.9104\\
-6.58443	-20.9104\\
-10.9282	29.5891\\
-10.9282	-29.5891\\
-7.36921	31.264\\
-7.36921	-31.264\\
-6.42794	22.0099\\
-6.42794	-22.0099\\
-8.34209	31.1841\\
-8.34209	-31.1841\\
-12.365	0\\
-10.501	0\\
-7.90475	30.9387\\
-7.90475	-30.9387\\
-9.67322	29.2882\\
-9.67322	-29.2882\\
-11.7415	25.9684\\
-11.7415	-25.9684\\
-11.6624	18.9232\\
-11.6624	-18.9232\\
-8.80028	19.8446\\
-8.80028	-19.8446\\
-14.3072	0\\
-14.8937	21.732\\
-14.8937	-21.732\\
-7.54624	21.3238\\
-7.54624	-21.3238\\
-14.8084	20.6075\\
-14.8084	-20.6075\\
-13.4559	0\\
-11.3856	0\\
-9.68166	28.4532\\
-9.68166	-28.4532\\
-13.6275	0\\
-11.0624	0\\
-11.4242	19.8011\\
-11.4242	-19.8011\\
-6.26013	30.4279\\
-6.26013	-30.4279\\
-10.9892	0\\
-6.25077	24.1744\\
-6.25077	-24.1744\\
-6.91419	30.2371\\
-6.91419	-30.2371\\
-6.00893	25.8256\\
-6.00893	-25.8256\\
-7.16076	23.3342\\
-7.16076	-23.3342\\
-7.3324	30.2859\\
-7.3324	-30.2859\\
-11.7013	25.7424\\
-11.7013	-25.7424\\
-6.61781	29.991\\
-6.61781	-29.991\\
-10.1727	22.3896\\
-10.1727	-22.3896\\
-6.04312	29.3867\\
-6.04312	-29.3867\\
-7.67258	23.9426\\
-7.67258	-23.9426\\
-12.1325	21.5509\\
-12.1325	-21.5509\\
-11.5393	22.5969\\
-11.5393	-22.5969\\
-6.31958	29.3006\\
-6.31958	-29.3006\\
-9.84884	27.4375\\
-9.84884	-27.4375\\
-9.79804	23.8779\\
-9.79804	-23.8779\\
-6.54858	29.3706\\
-6.54858	-29.3706\\
-9.31879	27.9156\\
-9.31879	-27.9156\\
-8.171	28.8067\\
-8.171	-28.8067\\
-6.14617	27.0931\\
-6.14617	-27.0931\\
-6.20332	27.3397\\
-6.20332	-27.3397\\
-9.47086	25.9837\\
-9.47086	-25.9837\\
-6.27578	27.7796\\
-6.27578	-27.7796\\
-8.74483	27.413\\
-8.74483	-27.413\\
-8.95932	27.1486\\
-8.95932	-27.1486\\
-7.59116	0\\
-8.37289	25.57\\
-8.37289	-25.57\\
-8.36775	25.5031\\
-8.36775	-25.5031\\
-7.7613	25.8766\\
-7.7613	-25.8766\\
-6.16792	28.2252\\
-6.16792	-28.2252\\
-7.64708	26.7235\\
-7.64708	-26.7235\\
-8.046	27.9841\\
-8.046	-27.9841\\
-7.94947	27.9974\\
-7.94947	-27.9974\\
-6.57642	28.306\\
-6.57642	-28.306\\
-6.53457	27.136\\
-6.53457	-27.136\\
-6.40941	27.1338\\
-6.40941	-27.1338\\
-7.11679	28.6839\\
-7.11679	-28.6839\\
-7.27883	28.5686\\
-7.27883	-28.5686\\
-7.50609	28.2927\\
-7.50609	-28.2927\\
-8.12744	27.2158\\
-8.12744	-27.2158\\
-7.4821	27.8219\\
-7.4821	-27.8219\\
-7.19542	27.6836\\
-7.19542	-27.6836\\
-7.15915	28.3033\\
-7.15915	-28.3033\\
-7.30762	28.2657\\
-7.30762	-28.2657\\
-7.06106	0\\
-14.8137	20.882\\
-14.8137	-20.882\\
};
\addplot [color=mycolor10, only marks, mark size=1.2pt, mark=*, mark options={solid, fill=mycolor10, mycolor10}, forget plot]
  table[row sep=crcr]{%
-14.513	57.261\\
-14.513	-57.261\\
-12.8939	50.4014\\
-12.8939	-50.4014\\
-13.3807	53.0353\\
-13.3807	-53.0353\\
-13.3738	53.5448\\
-13.3738	-53.5448\\
-14.8822	53.2253\\
-14.8822	-53.2253\\
-13.8071	55.2816\\
-13.8071	-55.2816\\
-14.0254	56.2478\\
-14.0254	-56.2478\\
-14.1763	56.6141\\
-14.1763	-56.6141\\
-14.0172	55.1899\\
-14.0172	-55.1899\\
-14.2742	56.849\\
-14.2742	-56.849\\
-14.3077	56.931\\
-14.3077	-56.931\\
-13.7075	53.3808\\
-13.7075	-53.3808\\
-11.1472	48.3312\\
-11.1472	-48.3312\\
-11.3185	47.8537\\
-11.3185	-47.8537\\
-11.1965	47.6522\\
-11.1965	-47.6522\\
-11.6335	47.203\\
-11.6335	-47.203\\
-10.959	48.1703\\
-10.959	-48.1703\\
-12.6874	45.2412\\
-12.6874	-45.2412\\
-10.5193	45.8992\\
-10.5193	-45.8992\\
-10.4023	45.4762\\
-10.4023	-45.4762\\
-10.5654	45.5209\\
-10.5654	-45.5209\\
-14.2264	56.7368\\
-14.2264	-56.7368\\
-14.0815	44.4823\\
-14.0815	-44.4823\\
-12.412	45.4195\\
-12.412	-45.4195\\
-10.1301	45.3867\\
-10.1301	-45.3867\\
-10.2695	44.5578\\
-10.2695	-44.5578\\
-10.6179	42.8475\\
-10.6179	-42.8475\\
-10.2216	43.4412\\
-10.2216	-43.4412\\
-9.7897	43.6986\\
-9.7897	-43.6986\\
-9.56206	43.5778\\
-9.56206	-43.5778\\
-9.7129	43.4434\\
-9.7129	-43.4434\\
-9.80809	43.6009\\
-9.80809	-43.6009\\
-12.8869	39.9732\\
-12.8869	-39.9732\\
-9.56813	42.4546\\
-9.56813	-42.4546\\
-9.42166	42.2759\\
-9.42166	-42.2759\\
-11.1766	41.7811\\
-11.1766	-41.7811\\
-9.37192	41.055\\
-9.37192	-41.055\\
-8.61996	39.7767\\
-8.61996	-39.7767\\
-9.87176	41.2687\\
-9.87176	-41.2687\\
-10.086	39.7998\\
-10.086	-39.7998\\
-10.6798	39.4984\\
-10.6798	-39.4984\\
-9.0205	38.4356\\
-9.0205	-38.4356\\
-9.28673	38.8251\\
-9.28673	-38.8251\\
-7.94908	37.1913\\
-7.94908	-37.1913\\
-7.21361	35.5032\\
-7.21361	-35.5032\\
-8.37717	35.7158\\
-8.37717	-35.7158\\
-8.8261	35.153\\
-8.8261	-35.153\\
-10.7926	34.1651\\
-10.7926	-34.1651\\
-7.66367	34.0509\\
-7.66367	-34.0509\\
-7.16903	33.3301\\
-7.16903	-33.3301\\
-7.10901	33.5537\\
-7.10901	-33.5537\\
-3.09804	0\\
-7.2122	8.91017\\
-7.2122	-8.91017\\
-3.7087	0\\
-4.1081	0\\
-4.21487	0\\
-4.54843	0\\
-4.59916	0\\
-5.13532	0\\
-6.85816	0\\
-6.20007	0\\
-6.64657	0\\
-9.6377	33.2027\\
-9.6377	-33.2027\\
-7.03088	17.3304\\
-7.03088	-17.3304\\
-8.04463	0\\
-7.30248	0\\
-7.44182	32.6587\\
-7.44182	-32.6587\\
-7.47341	32.1815\\
-7.47341	-32.1815\\
-8.74881	0\\
-6.58851	31.9734\\
-6.58851	-31.9734\\
-6.46775	31.8117\\
-6.46775	-31.8117\\
-6.58493	20.9113\\
-6.58493	-20.9113\\
-7.92215	19.4263\\
-7.92215	-19.4263\\
-8.94641	18.7986\\
-8.94641	-18.7986\\
-7.36775	31.2626\\
-7.36775	-31.2626\\
-8.28749	0\\
-11.0438	29.8013\\
-11.0438	-29.8013\\
-8.35264	31.1673\\
-8.35264	-31.1673\\
-7.90555	30.9422\\
-7.90555	-30.9422\\
-12.7582	26.8501\\
-12.7582	-26.8501\\
-9.71545	29.2855\\
-9.71545	-29.2855\\
-6.42535	22.0114\\
-6.42535	-22.0114\\
-11.6896	18.9635\\
-11.6896	-18.9635\\
-8.81843	19.8615\\
-8.81843	-19.8615\\
-14.2811	0\\
-7.55395	21.3207\\
-7.55395	-21.3207\\
-12.3416	0\\
-10.4931	0\\
-11.3579	0\\
-13.4398	0\\
-13.6102	0\\
-14.84	20.5299\\
-14.84	-20.5299\\
-14.9322	21.6838\\
-14.9322	-21.6838\\
-11.0526	0\\
-10.9804	0\\
-9.68519	28.4473\\
-9.68519	-28.4473\\
-11.4247	19.7946\\
-11.4247	-19.7946\\
-6.25999	30.428\\
-6.25999	-30.428\\
-6.25094	24.1745\\
-6.25094	-24.1745\\
-7.16049	23.3336\\
-7.16049	-23.3336\\
-7.3344	30.2864\\
-7.3344	-30.2864\\
-6.91486	30.2374\\
-6.91486	-30.2374\\
-6.01025	25.8227\\
-6.01025	-25.8227\\
-6.61719	29.9916\\
-6.61719	-29.9916\\
-6.04319	29.3867\\
-6.04319	-29.3867\\
-11.7457	25.9504\\
-11.7457	-25.9504\\
-11.7108	25.7336\\
-11.7108	-25.7336\\
-10.1464	22.387\\
-10.1464	-22.387\\
-7.67094	23.9387\\
-7.67094	-23.9387\\
-12.1531	21.5339\\
-12.1531	-21.5339\\
-11.557	22.5858\\
-11.557	-22.5858\\
-6.32129	29.2978\\
-6.32129	-29.2978\\
-9.85359	27.4449\\
-9.85359	-27.4449\\
-9.31745	27.9088\\
-9.31745	-27.9088\\
-9.79737	23.8746\\
-9.79737	-23.8746\\
-6.55045	29.3724\\
-6.55045	-29.3724\\
-8.1651	28.8078\\
-8.1651	-28.8078\\
-9.57276	26.0204\\
-9.57276	-26.0204\\
-6.14343	27.0945\\
-6.14343	-27.0945\\
-8.76195	27.4107\\
-8.76195	-27.4107\\
-8.97882	27.155\\
-8.97882	-27.155\\
-6.2045	27.3417\\
-6.2045	-27.3417\\
-6.27586	27.78\\
-6.27586	-27.78\\
-6.16785	28.2251\\
-6.16785	-28.2251\\
-7.59116	0\\
-14.8137	20.882\\
-14.8137	-20.882\\
-8.29931	25.61\\
-8.29931	-25.61\\
-8.35836	25.536\\
-8.35836	-25.536\\
-7.77702	25.8805\\
-7.77702	-25.8805\\
-7.28414	28.5936\\
-7.28414	-28.5936\\
-8.04342	27.9898\\
-8.04342	-27.9898\\
-7.65012	26.7232\\
-7.65012	-26.7232\\
-6.5727	28.307\\
-6.5727	-28.307\\
-7.96781	27.9909\\
-7.96781	-27.9909\\
-6.39865	27.1344\\
-6.39865	-27.1344\\
-7.119	28.6779\\
-7.119	-28.6779\\
-6.55029	27.1326\\
-6.55029	-27.1326\\
-7.5162	28.2964\\
-7.5162	-28.2964\\
-8.12926	27.2145\\
-8.12926	-27.2145\\
-7.48012	27.8074\\
-7.48012	-27.8074\\
-7.19672	27.6826\\
-7.19672	-27.6826\\
-7.15937	28.3033\\
-7.15937	-28.3033\\
-7.30754	28.2656\\
-7.30754	-28.2656\\
-7.06106	0\\
};
\addplot [color=mycolor11, only marks, mark size=1.2pt, mark=*, mark options={solid, fill=mycolor11, mycolor11}, forget plot]
  table[row sep=crcr]{%
-13.8071	55.2815\\
-13.8071	-55.2815\\
-14.513	57.2612\\
-14.513	-57.2612\\
-14.0254	56.2478\\
-14.0254	-56.2478\\
-14.1763	56.6141\\
-14.1763	-56.6141\\
-14.2742	56.849\\
-14.2742	-56.849\\
-14.3077	56.931\\
-14.3077	-56.931\\
-12.9186	50.4142\\
-12.9186	-50.4142\\
-13.3738	53.5466\\
-13.3738	-53.5466\\
-13.3813	53.0355\\
-13.3813	-53.0355\\
-14.0264	55.1999\\
-14.0264	-55.1999\\
-14.8532	53.2492\\
-14.8532	-53.2492\\
-13.7085	53.3812\\
-13.7085	-53.3812\\
-11.1481	48.3317\\
-11.1481	-48.3317\\
-11.3176	47.8539\\
-11.3176	-47.8539\\
-11.1966	47.6522\\
-11.1966	-47.6522\\
-10.959	48.1703\\
-10.959	-48.1703\\
-11.6549	47.2159\\
-11.6549	-47.2159\\
-12.6426	45.2673\\
-12.6426	-45.2673\\
-10.5209	45.902\\
-10.5209	-45.902\\
-10.4031	45.4801\\
-10.4031	-45.4801\\
-14.2264	56.7368\\
-14.2264	-56.7368\\
-10.1298	45.3869\\
-10.1298	-45.3869\\
-12.4175	45.4114\\
-12.4175	-45.4114\\
-10.5637	45.5195\\
-10.5637	-45.5195\\
-14.0579	44.455\\
-14.0579	-44.455\\
-10.2695	44.5575\\
-10.2695	-44.5575\\
-10.6271	42.8314\\
-10.6271	-42.8314\\
-10.2202	43.4405\\
-10.2202	-43.4405\\
-9.79164	43.6985\\
-9.79164	-43.6985\\
-9.56194	43.5783\\
-9.56194	-43.5783\\
-9.56809	42.4546\\
-9.56809	-42.4546\\
-9.71175	43.4434\\
-9.71175	-43.4434\\
-9.80819	43.6008\\
-9.80819	-43.6008\\
-8.61671	39.778\\
-8.61671	-39.778\\
-12.6958	40.0411\\
-12.6958	-40.0411\\
-9.42195	42.2757\\
-9.42195	-42.2757\\
-9.37545	41.0555\\
-9.37545	-41.0555\\
-11.1773	41.7866\\
-11.1773	-41.7866\\
-10.1228	39.8199\\
-10.1228	-39.8199\\
-9.87146	41.2685\\
-9.87146	-41.2685\\
-9.28766	38.8211\\
-9.28766	-38.8211\\
-10.6298	39.5093\\
-10.6298	-39.5093\\
-9.03165	38.4414\\
-9.03165	-38.4414\\
-7.95048	37.1947\\
-7.95048	-37.1947\\
-3.05796	0\\
-3.70714	0\\
-7.2136	35.5032\\
-7.2136	-35.5032\\
-8.36796	35.7129\\
-8.36796	-35.7129\\
-8.82398	35.1507\\
-8.82398	-35.1507\\
-10.7936	34.243\\
-10.7936	-34.243\\
-7.66234	34.0511\\
-7.66234	-34.0511\\
-9.65827	33.204\\
-9.65827	-33.204\\
-7.10923	33.5534\\
-7.10923	-33.5534\\
-7.16728	33.335\\
-7.16728	-33.335\\
-7.22607	8.8978\\
-7.22607	-8.8978\\
-7.44442	32.6546\\
-7.44442	-32.6546\\
-6.58803	31.9737\\
-6.58803	-31.9737\\
-6.46779	31.8117\\
-6.46779	-31.8117\\
-7.46344	32.1822\\
-7.46344	-32.1822\\
-4.10665	0\\
-7.03551	17.3312\\
-7.03551	-17.3312\\
-4.59596	0\\
-4.21339	0\\
-4.54623	0\\
-5.13324	0\\
-6.18765	0\\
-6.54432	0\\
-7.99987	0\\
-6.84097	0\\
-7.2879	0\\
-8.64768	0\\
-8.28076	0\\
-11.0138	29.9612\\
-11.0138	-29.9612\\
-6.5845	20.9136\\
-6.5845	-20.9136\\
-7.92658	19.4283\\
-7.92658	-19.4283\\
-8.35871	31.1824\\
-8.35871	-31.1824\\
-8.95124	18.8065\\
-8.95124	-18.8065\\
-7.36453	31.2641\\
-7.36453	-31.2641\\
-7.91047	30.9626\\
-7.91047	-30.9626\\
-12.7652	26.8549\\
-12.7652	-26.8549\\
-9.71417	29.2666\\
-9.71417	-29.2666\\
-6.42426	22.0107\\
-6.42426	-22.0107\\
-11.7467	18.9663\\
-11.7467	-18.9663\\
-14.8645	20.4731\\
-14.8645	-20.4731\\
-14.2562	0\\
-10.4898	0\\
-8.82803	19.854\\
-8.82803	-19.854\\
-14.9534	21.6482\\
-14.9534	-21.6482\\
-7.55831	21.324\\
-7.55831	-21.324\\
-12.3305	0\\
-13.4244	0\\
-11.3451	0\\
-11.0446	0\\
-13.5934	0\\
-6.25985	30.4281\\
-6.25985	-30.4281\\
-10.975	0\\
-6.91486	30.2375\\
-6.91486	-30.2375\\
-7.33407	30.2864\\
-7.33407	-30.2864\\
-6.25158	24.1741\\
-6.25158	-24.1741\\
-6.61515	29.9893\\
-6.61515	-29.9893\\
-11.422	19.7891\\
-11.422	-19.7891\\
-7.15707	23.33\\
-7.15707	-23.33\\
-6.00988	25.8237\\
-6.00988	-25.8237\\
-11.7484	25.9344\\
-11.7484	-25.9344\\
-11.7177	25.7264\\
-11.7177	-25.7264\\
-10.1684	22.3856\\
-10.1684	-22.3856\\
-7.66445	23.9323\\
-7.66445	-23.9323\\
-12.1628	21.5235\\
-12.1628	-21.5235\\
-11.5709	22.5766\\
-11.5709	-22.5766\\
-6.04324	29.3867\\
-6.04324	-29.3867\\
-9.68641	28.4498\\
-9.68641	-28.4498\\
-9.85759	27.45\\
-9.85759	-27.45\\
-9.79016	23.8711\\
-9.79016	-23.8711\\
-6.32091	29.2965\\
-6.32091	-29.2965\\
-9.30451	27.9141\\
-9.30451	-27.9141\\
-8.16351	28.8266\\
-8.16351	-28.8266\\
-6.13941	27.0949\\
-6.13941	-27.0949\\
-9.55091	26.0334\\
-9.55091	-26.0334\\
-6.20383	27.343\\
-6.20383	-27.343\\
-8.77761	27.4083\\
-8.77761	-27.4083\\
-6.54788	29.37\\
-6.54788	-29.37\\
-8.97972	27.1451\\
-8.97972	-27.1451\\
-7.06106	0\\
-8.23772	25.6411\\
-8.23772	-25.6411\\
-8.36404	25.5419\\
-8.36404	-25.5419\\
-7.78091	25.8781\\
-7.78091	-25.8781\\
-6.16768	28.2258\\
-6.16768	-28.2258\\
-6.27653	27.7814\\
-6.27653	-27.7814\\
-7.65241	26.7229\\
-7.65241	-26.7229\\
-8.04432	28.0033\\
-8.04432	-28.0033\\
-7.96819	27.9847\\
-7.96819	-27.9847\\
-7.11669	28.6651\\
-7.11669	-28.6651\\
-6.57384	28.3044\\
-6.57384	-28.3044\\
-6.55492	27.1311\\
-6.55492	-27.1311\\
-6.39561	27.1377\\
-6.39561	-27.1377\\
-7.29639	28.599\\
-7.29639	-28.599\\
-7.52027	28.2941\\
-7.52027	-28.2941\\
-8.1286	27.2158\\
-8.1286	-27.2158\\
-7.48239	27.808\\
-7.48239	-27.808\\
-7.19756	27.6816\\
-7.19756	-27.6816\\
-7.15952	28.3033\\
-7.15952	-28.3033\\
-7.30753	28.2656\\
-7.30753	-28.2656\\
-7.59116	0\\
-14.8137	20.882\\
-14.8137	-20.882\\
};
\addplot [color=mycolor12, only marks, mark size=1.2pt, mark=*, mark options={solid, fill=mycolor12, mycolor12}, forget plot]
  table[row sep=crcr]{%
-14.513	57.2613\\
-14.513	-57.2613\\
-13.3815	53.0356\\
-13.3815	-53.0356\\
-13.3732	53.5479\\
-13.3732	-53.5479\\
-14.8439	53.2525\\
-14.8439	-53.2525\\
-13.8071	55.2815\\
-13.8071	-55.2815\\
-14.0254	56.2478\\
-14.0254	-56.2478\\
-14.1763	56.6141\\
-14.1763	-56.6141\\
-14.0316	55.2058\\
-14.0316	-55.2058\\
-14.2742	56.849\\
-14.2742	-56.849\\
-13.7082	53.383\\
-13.7082	-53.383\\
-14.3077	56.931\\
-14.3077	-56.931\\
-12.9246	50.4192\\
-12.9246	-50.4192\\
-11.3147	47.8533\\
-11.3147	-47.8533\\
-11.1483	48.3319\\
-11.1483	-48.3319\\
-11.1967	47.6522\\
-11.1967	-47.6522\\
-11.6636	47.2232\\
-11.6636	-47.2232\\
-10.959	48.1703\\
-10.959	-48.1703\\
-14.2264	56.7368\\
-14.2264	-56.7368\\
-14.0334	44.4328\\
-14.0334	-44.4328\\
-12.6358	45.2672\\
-12.6358	-45.2672\\
-12.3961	45.3962\\
-12.3961	-45.3962\\
-10.5215	45.9036\\
-10.5215	-45.9036\\
-10.4035	45.4821\\
-10.4035	-45.4821\\
-10.5634	45.5192\\
-10.5634	-45.5192\\
-10.1298	45.3869\\
-10.1298	-45.3869\\
-10.2696	44.5574\\
-10.2696	-44.5574\\
-10.2194	43.4399\\
-10.2194	-43.4399\\
-9.79218	43.6986\\
-9.79218	-43.6986\\
-9.5619	43.5788\\
-9.5619	-43.5788\\
-9.71173	43.4434\\
-9.71173	-43.4434\\
-9.37515	41.0545\\
-9.37515	-41.0545\\
-10.6343	42.8151\\
-10.6343	-42.8151\\
-9.42203	42.2756\\
-9.42203	-42.2756\\
-9.56807	42.4547\\
-9.56807	-42.4547\\
-9.87121	41.2685\\
-9.87121	-41.2685\\
-9.80821	43.6006\\
-9.80821	-43.6006\\
-11.1772	41.7867\\
-11.1772	-41.7867\\
-12.7034	40.0503\\
-12.7034	-40.0503\\
-8.61555	39.7781\\
-8.61555	-39.7781\\
-10.1344	39.8269\\
-10.1344	-39.8269\\
-10.6349	39.506\\
-10.6349	-39.506\\
-9.28348	38.8133\\
-9.28348	-38.8133\\
-9.03886	38.4457\\
-9.03886	-38.4457\\
-7.95057	37.1956\\
-7.95057	-37.1956\\
-3.00616	0\\
-8.36359	35.7112\\
-8.36359	-35.7112\\
-7.2136	35.5032\\
-7.2136	-35.5032\\
-10.7748	34.2765\\
-10.7748	-34.2765\\
-8.82052	35.1434\\
-8.82052	-35.1434\\
-7.66221	34.0512\\
-7.66221	-34.0512\\
-9.63879	33.2145\\
-9.63879	-33.2145\\
-7.16727	33.3353\\
-7.16727	-33.3353\\
-7.10942	33.5531\\
-7.10942	-33.5531\\
-7.23434	8.89111\\
-7.23434	-8.89111\\
-7.04314	17.3309\\
-7.04314	-17.3309\\
-3.70567	0\\
-4.10509	0\\
-4.2116	0\\
-4.5437	0\\
-4.593	0\\
-5.13148	0\\
-7.92696	0\\
-6.17145	0\\
-6.4049	0\\
-6.82506	0\\
-7.27906	0\\
-7.44551	32.6556\\
-7.44551	-32.6556\\
-8.54524	0\\
-8.27599	0\\
-7.4679	32.1743\\
-7.4679	-32.1743\\
-6.58757	31.9739\\
-6.58757	-31.9739\\
-6.46782	31.8117\\
-6.46782	-31.8117\\
-6.58524	20.9146\\
-6.58524	-20.9146\\
-7.92993	19.4297\\
-7.92993	-19.4297\\
-8.95302	18.8005\\
-8.95302	-18.8005\\
-11.0218	29.966\\
-11.0218	-29.966\\
-7.36013	31.2664\\
-7.36013	-31.2664\\
-8.35995	31.1756\\
-8.35995	-31.1756\\
-7.9141	30.9654\\
-7.9141	-30.9654\\
-12.7723	26.8542\\
-12.7723	-26.8542\\
-6.42325	22.0103\\
-6.42325	-22.0103\\
-9.70224	29.242\\
-9.70224	-29.242\\
-11.7323	18.9486\\
-11.7323	-18.9486\\
-8.83063	19.8479\\
-8.83063	-19.8479\\
-14.2333	0\\
-14.8617	20.4364\\
-14.8617	-20.4364\\
-10.4851	0\\
-7.56316	21.3242\\
-7.56316	-21.3242\\
-11.3312	0\\
-13.4123	0\\
-12.3203	0\\
-14.9533	21.6274\\
-14.9533	-21.6274\\
-11.0381	0\\
-13.581	0\\
-11.4246	19.7796\\
-11.4246	-19.7796\\
-10.9702	0\\
-6.24816	24.1777\\
-6.24816	-24.1777\\
-6.25978	30.4281\\
-6.25978	-30.4281\\
-7.15997	23.3263\\
-7.15997	-23.3263\\
-6.01094	25.8238\\
-6.01094	-25.8238\\
-11.72	25.7171\\
-11.72	-25.7171\\
-11.7441	25.9269\\
-11.7441	-25.9269\\
-6.91497	30.2376\\
-6.91497	-30.2376\\
-7.33377	30.2866\\
-7.33377	-30.2866\\
-6.04327	29.3867\\
-6.04327	-29.3867\\
-6.61312	29.9872\\
-6.61312	-29.9872\\
-10.1714	22.3862\\
-10.1714	-22.3862\\
-7.66659	23.9392\\
-7.66659	-23.9392\\
-9.69302	28.4544\\
-9.69302	-28.4544\\
-12.1723	21.518\\
-12.1723	-21.518\\
-11.5738	22.5705\\
-11.5738	-22.5705\\
-9.84517	27.4469\\
-9.84517	-27.4469\\
-9.78818	23.8624\\
-9.78818	-23.8624\\
-9.30487	27.9001\\
-9.30487	-27.9001\\
-6.32059	29.2967\\
-6.32059	-29.2967\\
-6.54826	29.3693\\
-6.54826	-29.3693\\
-8.16096	28.8282\\
-8.16096	-28.8282\\
-6.13887	27.0933\\
-6.13887	-27.0933\\
-8.78803	27.3945\\
-8.78803	-27.3945\\
-9.52578	26.0415\\
-9.52578	-26.0415\\
-6.20172	27.3438\\
-6.20172	-27.3438\\
-8.99709	27.1449\\
-8.99709	-27.1449\\
-6.16752	28.2258\\
-6.16752	-28.2258\\
-6.27674	27.7829\\
-6.27674	-27.7829\\
-7.06106	0\\
-14.8137	20.882\\
-14.8137	-20.882\\
-8.2416	25.6414\\
-8.2416	-25.6414\\
-8.36648	25.5443\\
-8.36648	-25.5443\\
-7.78307	25.8759\\
-7.78307	-25.8759\\
-8.04359	28.0055\\
-8.04359	-28.0055\\
-7.65325	26.7229\\
-7.65325	-26.7229\\
-7.96376	27.9818\\
-7.96376	-27.9818\\
-6.56117	27.132\\
-6.56117	-27.132\\
-7.11915	28.6599\\
-7.11915	-28.6599\\
-6.39339	27.1361\\
-6.39339	-27.1361\\
-6.57296	28.2996\\
-6.57296	-28.2996\\
-7.29878	28.6185\\
-7.29878	-28.6185\\
-7.52212	28.2934\\
-7.52212	-28.2934\\
-8.12945	27.2153\\
-8.12945	-27.2153\\
-7.48409	27.8084\\
-7.48409	-27.8084\\
-7.19756	27.6812\\
-7.19756	-27.6812\\
-7.15963	28.3034\\
-7.15963	-28.3034\\
-7.30748	28.2657\\
-7.30748	-28.2657\\
-7.59116	0\\
};
\end{axis}

\begin{axis}[%
width={\figscale*1.4in},
height={\figscale*0.067in},
at={(\figscale*0.697in,\figscale*0.156in)},
scale only axis,
axis on top,
separate axis lines,
every outer x axis line/.append style={black},
every x tick label/.append style={font=\figticfont},
every x tick/.append style={black},
xmin=0,
xmax=100,
xtick={0,20,40,60,80,100},
xlabel={Share of units damped in \%},
every outer y axis line/.append style={black},
every y tick label/.append style={font=\figticfont},
every y tick/.append style={black},
y dir=reverse,
ymin=0,
ymax=1,
ytick={\empty},
axis background/.style={fill=mycolor13}
]
\addplot [forget plot] graphics [xmin=-0.196078, xmax=100.196, ymin=-0.5, ymax=1.5] {fig6_eig_zoom_3cases-1.png};

\addplot[area legend, draw=none, fill=mycolor2, forget plot]
table[row sep=crcr] {%
x	y\\
0	0\\
7.14286	0\\
7.14286	1\\
0	1\\
}--cycle;
\end{axis}

\begin{axis}[%
width={\figscale*0.619in},
height={\figscale*0.2in},
at={(\figscale*2.307in,\figscale*0.025in)},
scale only axis,
separate axis lines,
every outer x axis line/.append style={black},
every x tick label/.append style={font=\figticfont},
every x tick/.append style={black},
xmin=0,
xmax=1,
xtick={\empty},
every outer y axis line/.append style={black},
every y tick label/.append style={font=\figticfont},
every y tick/.append style={black},
ymin=0,
ymax=1,
ytick={\empty},
axis background/.style={fill=white}
]
\addplot [color=mycolor2, line width=0.9pt, only marks, mark size=2.2pt, mark=x, mark options={solid, mycolor2}, forget plot]
  table[row sep=crcr]{%
0.13	0.78\\
};
\node[right, align=left, inner sep=0]
at (axis cs:0.27,0.78) {$K = 0$};
\addplot [color=mycolor14, only marks, mark size=1.4pt, mark=*, mark options={solid, fill=mycolor14, mycolor14}, forget plot]
  table[row sep=crcr]{%
0.13	0.24\\
};
\node[right, align=left, inner sep=0]
at (axis cs:0.27,0.24) {$K \neq 0$};
\end{axis}
\end{tikzpicture}%
  \caption{Fig.~6 --- small-signal spectrum over the activation sweep.}
\end{figure}

% Fig. 7 is the one DOUBLE-COLUMN figure: the 3x3 grid spans the full text width
% (figure*, not figure), exactly as in the paper.
\begin{figure*}[!t]\centering
  % This file was created by matlab2tikz.
%
\definecolor{mycolor1}{rgb}{0.57647,0.00000,0.00000}%
\definecolor{mycolor2}{rgb}{0.46620,0.67454,0.77804}%
\definecolor{mycolor3}{rgb}{0.39131,0.62716,0.74576}%
\definecolor{mycolor4}{rgb}{0.30499,0.57091,0.70745}%
\definecolor{mycolor5}{rgb}{0.22212,0.51426,0.66892}%
\definecolor{mycolor6}{rgb}{0.14661,0.45864,0.63129}%
\definecolor{mycolor7}{rgb}{0.09143,0.41309,0.60075}%
\definecolor{mycolor8}{rgb}{0.04223,0.36317,0.56781}%
\definecolor{mycolor9}{rgb}{0.01771,0.31647,0.53748}%
\definecolor{mycolor10}{rgb}{0.01172,0.27830,0.51292}%
\definecolor{mycolor11}{rgb}{0.00988,0.23521,0.48525}%
\definecolor{mycolor12}{rgb}{0.07059,0.07059,0.07059}%
%
\begin{tikzpicture}[font=\figfont]

\begin{axis}[%
width={\figscale*1.528in},
height={\figscale*0.46in},
at={(\figscale*0.931in,\figscale*1.69in)},
scale only axis,
separate axis lines,
every outer x axis line/.append style={black},
every x tick label/.append style={font=\figticfont},
every x tick/.append style={black},
xmin=-0.025,
xmax=2,
xtick={0,0.5,1,1.5,2},
xticklabels={\empty},
every outer y axis line/.append style={black},
every y tick label/.append style={font=\figticfont},
every y tick/.append style={black},
ymin=-0.05,
ymax=0,
ytick={-0.04,0},
ylabel={$\Delta\fcoi(t)$},
axis background/.style={fill=white},
title style={font=\figcornerfont},
title={WECC},
xmajorgrids,
ymajorgrids,
grid style={white!55!black, opacity=0.3},
ylabel style={rotate=-90, at={(axis description cs:-0.145,0.5)}, anchor=east}
]
\addplot [color=mycolor1, line width=1.4pt, forget plot]
  table[row sep=crcr]{%
-0.024	0\\
-0.022	0\\
-0.02	0\\
-0.018	0\\
-0.016	0\\
-0.014	0\\
-0.012	0\\
-0.01	0\\
-0.008	0\\
-0.006	0\\
-0.004	0\\
-0.002	0\\
0	0\\
0.002	-0.0019122\\
0.004	-0.00365962\\
0.006	-0.0052577\\
0.008	-0.00672084\\
0.01	-0.00806242\\
0.012	-0.00929492\\
0.014	-0.0104299\\
0.016	-0.011478\\
0.018	-0.0124492\\
0.02	-0.0133525\\
0.022	-0.0141963\\
0.024	-0.0149883\\
0.026	-0.0157352\\
0.028	-0.0164435\\
0.03	-0.0171187\\
0.032	-0.0177659\\
0.034	-0.0183895\\
0.036	-0.0189935\\
0.038	-0.0195812\\
0.04	-0.0201555\\
0.042	-0.0207188\\
0.044	-0.0212731\\
0.046	-0.0218202\\
0.048	-0.0223611\\
0.05	-0.022897\\
0.052	-0.0234283\\
0.054	-0.0239556\\
0.056	-0.024479\\
0.058	-0.0249985\\
0.06	-0.0255139\\
0.062	-0.0260251\\
0.064	-0.0265317\\
0.066	-0.0270332\\
0.068	-0.0275293\\
0.07	-0.0280194\\
0.072	-0.0285032\\
0.074	-0.0289802\\
0.076	-0.0294501\\
0.078	-0.0299125\\
0.08	-0.0303672\\
0.082	-0.030814\\
0.084	-0.0312528\\
0.086	-0.0316834\\
0.088	-0.0321058\\
0.09	-0.0325201\\
0.092	-0.0329264\\
0.094	-0.0333248\\
0.096	-0.0337155\\
0.098	-0.0340986\\
0.1	-0.0344744\\
0.102	-0.0348431\\
0.104	-0.0352048\\
0.106	-0.0355598\\
0.108	-0.0359082\\
0.11	-0.0362501\\
0.112	-0.0365857\\
0.114	-0.036915\\
0.116	-0.037238\\
0.118	-0.0375547\\
0.12	-0.037865\\
0.122	-0.0381689\\
0.124	-0.038466\\
0.126	-0.0387562\\
0.128	-0.0390394\\
0.13	-0.0393151\\
0.132	-0.0395831\\
0.134	-0.0398431\\
0.136	-0.0400948\\
0.138	-0.0403377\\
0.14	-0.0405716\\
0.142	-0.0407961\\
0.144	-0.041011\\
0.146	-0.0412158\\
0.148	-0.0414105\\
0.15	-0.0415946\\
0.152	-0.0417681\\
0.154	-0.0419308\\
0.156	-0.0420826\\
0.158	-0.0422233\\
0.16	-0.0423532\\
0.162	-0.042472\\
0.164	-0.0425801\\
0.166	-0.0426775\\
0.168	-0.0427643\\
0.17	-0.042841\\
0.172	-0.0429077\\
0.174	-0.0429648\\
0.176	-0.0430126\\
0.178	-0.0430516\\
0.18	-0.0430821\\
0.182	-0.0431047\\
0.184	-0.0431198\\
0.186	-0.043128\\
0.188	-0.0431296\\
0.19	-0.0431252\\
0.192	-0.0431154\\
0.194	-0.0431006\\
0.196	-0.0430813\\
0.198	-0.0430581\\
0.2	-0.0430313\\
0.202	-0.0430016\\
0.204	-0.0429692\\
0.206	-0.0429348\\
0.208	-0.0428985\\
0.21	-0.042861\\
0.212	-0.0428224\\
0.214	-0.0427831\\
0.216	-0.0427435\\
0.218	-0.0427039\\
0.22	-0.0426644\\
0.222	-0.0426254\\
0.224	-0.042587\\
0.226	-0.0425495\\
0.228	-0.042513\\
0.23	-0.0424777\\
0.232	-0.0424437\\
0.234	-0.0424111\\
0.236	-0.04238\\
0.238	-0.0423505\\
0.24	-0.0423227\\
0.242	-0.0422966\\
0.244	-0.0422722\\
0.246	-0.0422497\\
0.248	-0.042229\\
0.25	-0.0422101\\
0.252	-0.0421931\\
0.254	-0.0421779\\
0.256	-0.0421646\\
0.258	-0.042153\\
0.26	-0.0421433\\
0.262	-0.0421354\\
0.264	-0.0421292\\
0.266	-0.0421248\\
0.268	-0.042122\\
0.27	-0.0421209\\
0.272	-0.0421214\\
0.274	-0.0421235\\
0.276	-0.0421271\\
0.278	-0.0421323\\
0.28	-0.0421388\\
0.282	-0.0421467\\
0.284	-0.042156\\
0.286	-0.0421666\\
0.288	-0.0421784\\
0.29	-0.0421913\\
0.292	-0.0422054\\
0.294	-0.0422206\\
0.296	-0.0422367\\
0.298	-0.0422538\\
0.3	-0.0422718\\
0.302	-0.0422906\\
0.304	-0.0423102\\
0.306	-0.0423305\\
0.308	-0.0423514\\
0.31	-0.042373\\
0.312	-0.0423951\\
0.314	-0.0424177\\
0.316	-0.0424407\\
0.318	-0.0424642\\
0.32	-0.0424879\\
0.322	-0.042512\\
0.324	-0.0425363\\
0.326	-0.0425608\\
0.328	-0.0425855\\
0.33	-0.0426102\\
0.332	-0.0426351\\
0.334	-0.04266\\
0.336	-0.042685\\
0.338	-0.0427099\\
0.34	-0.0427348\\
0.342	-0.0427596\\
0.344	-0.0427843\\
0.346	-0.0428089\\
0.348	-0.0428333\\
0.35	-0.0428576\\
0.352	-0.0428817\\
0.354	-0.0429057\\
0.356	-0.0429294\\
0.358	-0.0429529\\
0.36	-0.0429762\\
0.362	-0.0429992\\
0.364	-0.043022\\
0.366	-0.0430446\\
0.368	-0.0430668\\
0.37	-0.0430888\\
0.372	-0.0431106\\
0.374	-0.043132\\
0.376	-0.0431532\\
0.378	-0.043174\\
0.38	-0.0431946\\
0.382	-0.0432149\\
0.384	-0.0432348\\
0.386	-0.0432544\\
0.388	-0.0432737\\
0.39	-0.0432927\\
0.392	-0.0433114\\
0.394	-0.0433297\\
0.396	-0.0433476\\
0.398	-0.0433652\\
0.4	-0.0433825\\
0.402	-0.0433994\\
0.404	-0.0434159\\
0.406	-0.0434321\\
0.408	-0.0434478\\
0.41	-0.0434632\\
0.412	-0.0434782\\
0.414	-0.0434927\\
0.416	-0.0435069\\
0.418	-0.0435207\\
0.42	-0.043534\\
0.422	-0.0435469\\
0.424	-0.0435594\\
0.426	-0.0435715\\
0.428	-0.0435831\\
0.43	-0.0435942\\
0.432	-0.043605\\
0.434	-0.0436153\\
0.436	-0.0436251\\
0.438	-0.0436345\\
0.44	-0.0436434\\
0.442	-0.0436519\\
0.444	-0.04366\\
0.446	-0.0436676\\
0.448	-0.0436748\\
0.45	-0.0436815\\
0.452	-0.0436878\\
0.454	-0.0436937\\
0.456	-0.0436992\\
0.458	-0.0437043\\
0.46	-0.0437089\\
0.462	-0.0437132\\
0.464	-0.0437171\\
0.466	-0.0437206\\
0.468	-0.0437238\\
0.47	-0.0437266\\
0.472	-0.043729\\
0.474	-0.0437312\\
0.476	-0.043733\\
0.478	-0.0437346\\
0.48	-0.0437358\\
0.482	-0.0437368\\
0.484	-0.0437375\\
0.486	-0.043738\\
0.488	-0.0437383\\
0.49	-0.0437384\\
0.492	-0.0437382\\
0.494	-0.0437379\\
0.496	-0.0437375\\
0.498	-0.0437369\\
0.5	-0.0437361\\
0.502	-0.0437352\\
0.504	-0.0437343\\
0.506	-0.0437332\\
0.508	-0.0437321\\
0.51	-0.0437309\\
0.512	-0.0437297\\
0.514	-0.0437285\\
0.516	-0.0437272\\
0.518	-0.0437259\\
0.52	-0.0437247\\
0.522	-0.0437234\\
0.524	-0.0437222\\
0.526	-0.043721\\
0.528	-0.0437198\\
0.53	-0.0437187\\
0.532	-0.0437177\\
0.534	-0.0437167\\
0.536	-0.0437158\\
0.538	-0.043715\\
0.54	-0.0437143\\
0.542	-0.0437137\\
0.544	-0.0437131\\
0.546	-0.0437127\\
0.548	-0.0437124\\
0.55	-0.0437121\\
0.552	-0.043712\\
0.554	-0.043712\\
0.556	-0.043712\\
0.558	-0.0437122\\
0.56	-0.0437125\\
0.562	-0.0437128\\
0.564	-0.0437133\\
0.566	-0.0437139\\
0.568	-0.0437145\\
0.57	-0.0437152\\
0.572	-0.043716\\
0.574	-0.0437169\\
0.576	-0.0437179\\
0.578	-0.0437189\\
0.58	-0.04372\\
0.582	-0.0437211\\
0.584	-0.0437223\\
0.586	-0.0437235\\
0.588	-0.0437247\\
0.59	-0.043726\\
0.592	-0.0437273\\
0.594	-0.0437286\\
0.596	-0.0437299\\
0.598	-0.0437312\\
0.6	-0.0437325\\
0.602	-0.0437337\\
0.604	-0.043735\\
0.606	-0.0437362\\
0.608	-0.0437373\\
0.61	-0.0437384\\
0.612	-0.0437395\\
0.614	-0.0437405\\
0.616	-0.0437414\\
0.618	-0.0437423\\
0.62	-0.043743\\
0.622	-0.0437437\\
0.624	-0.0437443\\
0.626	-0.0437448\\
0.628	-0.0437452\\
0.63	-0.0437454\\
0.632	-0.0437456\\
0.634	-0.0437456\\
0.636	-0.0437456\\
0.638	-0.0437454\\
0.64	-0.0437451\\
0.642	-0.0437446\\
0.644	-0.043744\\
0.646	-0.0437433\\
0.648	-0.0437425\\
0.65	-0.0437415\\
0.652	-0.0437404\\
0.654	-0.0437392\\
0.656	-0.0437378\\
0.658	-0.0437363\\
0.66	-0.0437347\\
0.662	-0.0437329\\
0.664	-0.0437311\\
0.666	-0.0437291\\
0.668	-0.0437269\\
0.67	-0.0437247\\
0.672	-0.0437224\\
0.674	-0.0437199\\
0.676	-0.0437173\\
0.678	-0.0437147\\
0.68	-0.0437119\\
0.682	-0.043709\\
0.684	-0.0437061\\
0.686	-0.0437031\\
0.688	-0.0437\\
0.69	-0.0436968\\
0.692	-0.0436936\\
0.694	-0.0436904\\
0.696	-0.0436871\\
0.698	-0.0436837\\
0.7	-0.0436803\\
0.702	-0.0436769\\
0.704	-0.0436735\\
0.706	-0.0436701\\
0.708	-0.0436667\\
0.71	-0.0436633\\
0.712	-0.0436599\\
0.714	-0.0436565\\
0.716	-0.0436531\\
0.718	-0.0436498\\
0.72	-0.0436466\\
0.722	-0.0436434\\
0.724	-0.0436402\\
0.726	-0.0436372\\
0.728	-0.0436342\\
0.73	-0.0436313\\
0.732	-0.0436284\\
0.734	-0.0436257\\
0.736	-0.0436231\\
0.738	-0.0436206\\
0.74	-0.0436183\\
0.742	-0.043616\\
0.744	-0.0436139\\
0.746	-0.043612\\
0.748	-0.0436101\\
0.75	-0.0436085\\
0.752	-0.0436069\\
0.754	-0.0436056\\
0.756	-0.0436044\\
0.758	-0.0436033\\
0.76	-0.0436025\\
0.762	-0.0436018\\
0.764	-0.0436013\\
0.766	-0.0436009\\
0.768	-0.0436008\\
0.77	-0.0436008\\
0.772	-0.043601\\
0.774	-0.0436014\\
0.776	-0.0436019\\
0.778	-0.0436027\\
0.78	-0.0436036\\
0.782	-0.0436047\\
0.784	-0.043606\\
0.786	-0.0436075\\
0.788	-0.0436092\\
0.79	-0.043611\\
0.792	-0.043613\\
0.794	-0.0436151\\
0.796	-0.0436174\\
0.798	-0.0436199\\
0.8	-0.0436225\\
0.802	-0.0436253\\
0.804	-0.0436282\\
0.806	-0.0436313\\
0.808	-0.0436345\\
0.81	-0.0436378\\
0.812	-0.0436412\\
0.814	-0.0436447\\
0.816	-0.0436483\\
0.818	-0.043652\\
0.82	-0.0436559\\
0.822	-0.0436597\\
0.824	-0.0436637\\
0.826	-0.0436677\\
0.828	-0.0436718\\
0.83	-0.0436759\\
0.832	-0.0436801\\
0.834	-0.0436842\\
0.836	-0.0436884\\
0.838	-0.0436926\\
0.84	-0.0436968\\
0.842	-0.043701\\
0.844	-0.0437052\\
0.846	-0.0437093\\
0.848	-0.0437135\\
0.85	-0.0437175\\
0.852	-0.0437215\\
0.854	-0.0437255\\
0.856	-0.0437294\\
0.858	-0.0437332\\
0.86	-0.0437369\\
0.862	-0.0437405\\
0.864	-0.043744\\
0.866	-0.0437474\\
0.868	-0.0437507\\
0.87	-0.0437539\\
0.872	-0.0437569\\
0.874	-0.0437598\\
0.876	-0.0437626\\
0.878	-0.0437652\\
0.88	-0.0437677\\
0.882	-0.04377\\
0.884	-0.0437722\\
0.886	-0.0437742\\
0.888	-0.043776\\
0.89	-0.0437776\\
0.892	-0.0437791\\
0.894	-0.0437804\\
0.896	-0.0437815\\
0.898	-0.0437825\\
0.9	-0.0437833\\
0.902	-0.0437838\\
0.904	-0.0437843\\
0.906	-0.0437845\\
0.908	-0.0437845\\
0.91	-0.0437844\\
0.912	-0.0437841\\
0.914	-0.0437836\\
0.916	-0.043783\\
0.918	-0.0437822\\
0.92	-0.0437812\\
0.922	-0.04378\\
0.924	-0.0437787\\
0.926	-0.0437773\\
0.928	-0.0437756\\
0.93	-0.0437739\\
0.932	-0.043772\\
0.934	-0.04377\\
0.936	-0.0437678\\
0.938	-0.0437655\\
0.94	-0.0437631\\
0.942	-0.0437606\\
0.944	-0.043758\\
0.946	-0.0437552\\
0.948	-0.0437524\\
0.95	-0.0437495\\
0.952	-0.0437466\\
0.954	-0.0437435\\
0.956	-0.0437404\\
0.958	-0.0437373\\
0.96	-0.0437341\\
0.962	-0.0437308\\
0.964	-0.0437275\\
0.966	-0.0437242\\
0.968	-0.0437209\\
0.97	-0.0437175\\
0.972	-0.0437142\\
0.974	-0.0437109\\
0.976	-0.0437075\\
0.978	-0.0437042\\
0.98	-0.0437009\\
0.982	-0.0436977\\
0.984	-0.0436945\\
0.986	-0.0436913\\
0.988	-0.0436882\\
0.99	-0.0436851\\
0.992	-0.0436821\\
0.994	-0.0436792\\
0.996	-0.0436763\\
0.998	-0.0436736\\
1	-0.0436709\\
1.002	-0.0436683\\
1.004	-0.0436658\\
1.006	-0.0436634\\
1.008	-0.043661\\
1.01	-0.0436588\\
1.012	-0.0436567\\
1.014	-0.0436548\\
1.016	-0.0436529\\
1.018	-0.0436511\\
1.02	-0.0436495\\
1.022	-0.043648\\
1.024	-0.0436466\\
1.026	-0.0436453\\
1.028	-0.0436442\\
1.03	-0.0436432\\
1.032	-0.0436423\\
1.034	-0.0436416\\
1.036	-0.0436409\\
1.038	-0.0436404\\
1.04	-0.04364\\
1.042	-0.0436398\\
1.044	-0.0436397\\
1.046	-0.0436397\\
1.048	-0.0436398\\
1.05	-0.04364\\
1.052	-0.0436403\\
1.054	-0.0436408\\
1.056	-0.0436414\\
1.058	-0.0436421\\
1.06	-0.0436429\\
1.062	-0.0436437\\
1.064	-0.0436447\\
1.066	-0.0436458\\
1.068	-0.043647\\
1.07	-0.0436483\\
1.072	-0.0436496\\
1.074	-0.0436511\\
1.076	-0.0436526\\
1.078	-0.0436542\\
1.08	-0.0436558\\
1.082	-0.0436575\\
1.084	-0.0436593\\
1.086	-0.0436611\\
1.088	-0.043663\\
1.09	-0.0436649\\
1.092	-0.0436669\\
1.094	-0.0436689\\
1.096	-0.0436709\\
1.098	-0.043673\\
1.1	-0.0436751\\
1.102	-0.0436772\\
1.104	-0.0436793\\
1.106	-0.0436814\\
1.108	-0.0436835\\
1.11	-0.0436856\\
1.112	-0.0436878\\
1.114	-0.0436899\\
1.116	-0.043692\\
1.118	-0.043694\\
1.12	-0.0436961\\
1.122	-0.0436981\\
1.124	-0.0437001\\
1.126	-0.0437021\\
1.128	-0.043704\\
1.13	-0.0437059\\
1.132	-0.0437077\\
1.134	-0.0437095\\
1.136	-0.0437112\\
1.138	-0.0437129\\
1.14	-0.0437145\\
1.142	-0.0437161\\
1.144	-0.0437176\\
1.146	-0.0437191\\
1.148	-0.0437205\\
1.15	-0.0437218\\
1.152	-0.043723\\
1.154	-0.0437242\\
1.156	-0.0437253\\
1.158	-0.0437263\\
1.16	-0.0437273\\
1.162	-0.0437282\\
1.164	-0.043729\\
1.166	-0.0437297\\
1.168	-0.0437303\\
1.17	-0.0437309\\
1.172	-0.0437314\\
1.174	-0.0437318\\
1.176	-0.0437322\\
1.178	-0.0437324\\
1.18	-0.0437326\\
1.182	-0.0437327\\
1.184	-0.0437328\\
1.186	-0.0437328\\
1.188	-0.0437327\\
1.19	-0.0437325\\
1.192	-0.0437322\\
1.194	-0.0437319\\
1.196	-0.0437315\\
1.198	-0.0437311\\
1.2	-0.0437306\\
1.202	-0.04373\\
1.204	-0.0437294\\
1.206	-0.0437287\\
1.208	-0.043728\\
1.21	-0.0437272\\
1.212	-0.0437264\\
1.214	-0.0437255\\
1.216	-0.0437246\\
1.218	-0.0437236\\
1.22	-0.0437226\\
1.222	-0.0437216\\
1.224	-0.0437205\\
1.226	-0.0437194\\
1.228	-0.0437182\\
1.23	-0.0437171\\
1.232	-0.0437159\\
1.234	-0.0437147\\
1.236	-0.0437135\\
1.238	-0.0437123\\
1.24	-0.043711\\
1.242	-0.0437098\\
1.244	-0.0437086\\
1.246	-0.0437073\\
1.248	-0.0437061\\
1.25	-0.0437048\\
1.252	-0.0437036\\
1.254	-0.0437023\\
1.256	-0.0437011\\
1.258	-0.0436999\\
1.26	-0.0436987\\
1.262	-0.0436976\\
1.264	-0.0436964\\
1.266	-0.0436953\\
1.268	-0.0436942\\
1.27	-0.0436931\\
1.272	-0.043692\\
1.274	-0.043691\\
1.276	-0.04369\\
1.278	-0.0436891\\
1.28	-0.0436882\\
1.282	-0.0436873\\
1.284	-0.0436864\\
1.286	-0.0436856\\
1.288	-0.0436849\\
1.29	-0.0436841\\
1.292	-0.0436835\\
1.294	-0.0436828\\
1.296	-0.0436822\\
1.298	-0.0436817\\
1.3	-0.0436812\\
1.302	-0.0436807\\
1.304	-0.0436803\\
1.306	-0.0436799\\
1.308	-0.0436796\\
1.31	-0.0436793\\
1.312	-0.0436791\\
1.314	-0.0436789\\
1.316	-0.0436787\\
1.318	-0.0436786\\
1.32	-0.0436785\\
1.322	-0.0436785\\
1.324	-0.0436785\\
1.326	-0.0436786\\
1.328	-0.0436787\\
1.33	-0.0436788\\
1.332	-0.043679\\
1.334	-0.0436792\\
1.336	-0.0436795\\
1.338	-0.0436798\\
1.34	-0.0436801\\
1.342	-0.0436804\\
1.344	-0.0436808\\
1.346	-0.0436812\\
1.348	-0.0436817\\
1.35	-0.0436821\\
1.352	-0.0436826\\
1.354	-0.0436832\\
1.356	-0.0436837\\
1.358	-0.0436842\\
1.36	-0.0436848\\
1.362	-0.0436854\\
1.364	-0.043686\\
1.366	-0.0436866\\
1.368	-0.0436873\\
1.37	-0.0436879\\
1.372	-0.0436886\\
1.374	-0.0436892\\
1.376	-0.0436899\\
1.378	-0.0436906\\
1.38	-0.0436912\\
1.382	-0.0436919\\
1.384	-0.0436926\\
1.386	-0.0436933\\
1.388	-0.0436939\\
1.39	-0.0436946\\
1.392	-0.0436953\\
1.394	-0.0436959\\
1.396	-0.0436966\\
1.398	-0.0436972\\
1.4	-0.0436978\\
1.402	-0.0436984\\
1.404	-0.043699\\
1.406	-0.0436996\\
1.408	-0.0437002\\
1.41	-0.0437007\\
1.412	-0.0437013\\
1.414	-0.0437018\\
1.416	-0.0437023\\
1.418	-0.0437028\\
1.42	-0.0437032\\
1.422	-0.0437037\\
1.424	-0.0437041\\
1.426	-0.0437045\\
1.428	-0.0437049\\
1.43	-0.0437052\\
1.432	-0.0437056\\
1.434	-0.0437059\\
1.436	-0.0437062\\
1.438	-0.0437064\\
1.44	-0.0437067\\
1.442	-0.0437069\\
1.444	-0.0437071\\
1.446	-0.0437072\\
1.448	-0.0437074\\
1.45	-0.0437075\\
1.452	-0.0437076\\
1.454	-0.0437077\\
1.456	-0.0437078\\
1.458	-0.0437078\\
1.46	-0.0437078\\
1.462	-0.0437078\\
1.464	-0.0437078\\
1.466	-0.0437078\\
1.468	-0.0437077\\
1.47	-0.0437076\\
1.472	-0.0437075\\
1.474	-0.0437074\\
1.476	-0.0437073\\
1.478	-0.0437071\\
1.48	-0.0437069\\
1.482	-0.0437068\\
1.484	-0.0437066\\
1.486	-0.0437063\\
1.488	-0.0437061\\
1.49	-0.0437059\\
1.492	-0.0437056\\
1.494	-0.0437054\\
1.496	-0.0437051\\
1.498	-0.0437048\\
1.5	-0.0437045\\
1.502	-0.0437042\\
1.504	-0.0437039\\
1.506	-0.0437036\\
1.508	-0.0437033\\
1.51	-0.043703\\
1.512	-0.0437027\\
1.514	-0.0437024\\
1.516	-0.043702\\
1.518	-0.0437017\\
1.52	-0.0437014\\
1.522	-0.043701\\
1.524	-0.0437007\\
1.526	-0.0437004\\
1.528	-0.0437\\
1.53	-0.0436997\\
1.532	-0.0436994\\
1.534	-0.0436991\\
1.536	-0.0436987\\
1.538	-0.0436984\\
1.54	-0.0436981\\
1.542	-0.0436978\\
1.544	-0.0436975\\
1.546	-0.0436972\\
1.548	-0.0436969\\
1.55	-0.0436967\\
1.552	-0.0436964\\
1.554	-0.0436961\\
1.556	-0.0436959\\
1.558	-0.0436956\\
1.56	-0.0436954\\
1.562	-0.0436952\\
1.564	-0.0436949\\
1.566	-0.0436947\\
1.568	-0.0436945\\
1.57	-0.0436944\\
1.572	-0.0436942\\
1.574	-0.043694\\
1.576	-0.0436939\\
1.578	-0.0436937\\
1.58	-0.0436936\\
1.582	-0.0436935\\
1.584	-0.0436934\\
1.586	-0.0436933\\
1.588	-0.0436932\\
1.59	-0.0436931\\
1.592	-0.043693\\
1.594	-0.043693\\
1.596	-0.0436929\\
1.598	-0.0436929\\
1.6	-0.0436929\\
1.602	-0.0436929\\
1.604	-0.0436929\\
1.606	-0.0436929\\
1.608	-0.0436929\\
1.61	-0.0436929\\
1.612	-0.0436929\\
1.614	-0.043693\\
1.616	-0.043693\\
1.618	-0.0436931\\
1.62	-0.0436932\\
1.622	-0.0436933\\
1.624	-0.0436933\\
1.626	-0.0436934\\
1.628	-0.0436935\\
1.63	-0.0436937\\
1.632	-0.0436938\\
1.634	-0.0436939\\
1.636	-0.043694\\
1.638	-0.0436942\\
1.64	-0.0436943\\
1.642	-0.0436944\\
1.644	-0.0436946\\
1.646	-0.0436947\\
1.648	-0.0436949\\
1.65	-0.0436951\\
1.652	-0.0436952\\
1.654	-0.0436954\\
1.656	-0.0436956\\
1.658	-0.0436957\\
1.66	-0.0436959\\
1.662	-0.0436961\\
1.664	-0.0436963\\
1.666	-0.0436964\\
1.668	-0.0436966\\
1.67	-0.0436968\\
1.672	-0.043697\\
1.674	-0.0436971\\
1.676	-0.0436973\\
1.678	-0.0436975\\
1.68	-0.0436977\\
1.682	-0.0436978\\
1.684	-0.043698\\
1.686	-0.0436982\\
1.688	-0.0436983\\
1.69	-0.0436985\\
1.692	-0.0436986\\
1.694	-0.0436988\\
1.696	-0.0436989\\
1.698	-0.0436991\\
1.7	-0.0436992\\
1.702	-0.0436993\\
1.704	-0.0436995\\
1.706	-0.0436996\\
1.708	-0.0436997\\
1.71	-0.0436998\\
1.712	-0.0436999\\
1.714	-0.0437\\
1.716	-0.0437001\\
1.718	-0.0437002\\
1.72	-0.0437003\\
1.722	-0.0437004\\
1.724	-0.0437004\\
1.726	-0.0437005\\
1.728	-0.0437005\\
1.73	-0.0437006\\
1.732	-0.0437006\\
1.734	-0.0437007\\
1.736	-0.0437007\\
1.738	-0.0437007\\
1.74	-0.0437008\\
1.742	-0.0437008\\
1.744	-0.0437008\\
1.746	-0.0437008\\
1.748	-0.0437008\\
1.75	-0.0437008\\
1.752	-0.0437007\\
1.754	-0.0437007\\
1.756	-0.0437007\\
1.758	-0.0437007\\
1.76	-0.0437006\\
1.762	-0.0437006\\
1.764	-0.0437005\\
1.766	-0.0437005\\
1.768	-0.0437004\\
1.77	-0.0437004\\
1.772	-0.0437003\\
1.774	-0.0437002\\
1.776	-0.0437001\\
1.778	-0.0437001\\
1.78	-0.0437\\
1.782	-0.0436999\\
1.784	-0.0436998\\
1.786	-0.0436997\\
1.788	-0.0436996\\
1.79	-0.0436995\\
1.792	-0.0436994\\
1.794	-0.0436993\\
1.796	-0.0436992\\
1.798	-0.0436991\\
1.8	-0.043699\\
1.802	-0.0436989\\
1.804	-0.0436988\\
1.806	-0.0436987\\
1.808	-0.0436986\\
1.81	-0.0436985\\
1.812	-0.0436984\\
1.814	-0.0436983\\
1.816	-0.0436982\\
1.818	-0.0436981\\
1.82	-0.043698\\
1.822	-0.0436979\\
1.824	-0.0436978\\
1.826	-0.0436977\\
1.828	-0.0436976\\
1.83	-0.0436975\\
1.832	-0.0436974\\
1.834	-0.0436973\\
1.836	-0.0436972\\
1.838	-0.0436971\\
1.84	-0.043697\\
1.842	-0.043697\\
1.844	-0.0436969\\
1.846	-0.0436968\\
1.848	-0.0436967\\
1.85	-0.0436967\\
1.852	-0.0436966\\
1.854	-0.0436966\\
1.856	-0.0436965\\
1.858	-0.0436965\\
1.86	-0.0436964\\
1.862	-0.0436964\\
1.864	-0.0436963\\
1.866	-0.0436963\\
1.868	-0.0436963\\
1.87	-0.0436962\\
1.872	-0.0436962\\
1.874	-0.0436962\\
1.876	-0.0436962\\
1.878	-0.0436962\\
1.88	-0.0436962\\
1.882	-0.0436962\\
1.884	-0.0436962\\
1.886	-0.0436962\\
1.888	-0.0436962\\
1.89	-0.0436962\\
1.892	-0.0436962\\
1.894	-0.0436963\\
1.896	-0.0436963\\
1.898	-0.0436963\\
1.9	-0.0436963\\
1.902	-0.0436964\\
1.904	-0.0436964\\
1.906	-0.0436965\\
1.908	-0.0436965\\
1.91	-0.0436965\\
1.912	-0.0436966\\
1.914	-0.0436966\\
1.916	-0.0436967\\
1.918	-0.0436967\\
1.92	-0.0436968\\
1.922	-0.0436969\\
1.924	-0.0436969\\
1.926	-0.043697\\
1.928	-0.043697\\
1.93	-0.0436971\\
1.932	-0.0436972\\
1.934	-0.0436972\\
1.936	-0.0436973\\
1.938	-0.0436974\\
1.94	-0.0436974\\
1.942	-0.0436975\\
1.944	-0.0436976\\
1.946	-0.0436976\\
1.948	-0.0436977\\
1.95	-0.0436978\\
1.952	-0.0436978\\
1.954	-0.0436979\\
1.956	-0.043698\\
1.958	-0.043698\\
1.96	-0.0436981\\
1.962	-0.0436981\\
1.964	-0.0436982\\
1.966	-0.0436983\\
1.968	-0.0436983\\
1.97	-0.0436984\\
1.972	-0.0436984\\
1.974	-0.0436985\\
1.976	-0.0436985\\
1.978	-0.0436986\\
1.98	-0.0436986\\
1.982	-0.0436986\\
1.984	-0.0436987\\
1.986	-0.0436987\\
1.988	-0.0436988\\
1.99	-0.0436988\\
1.992	-0.0436988\\
1.994	-0.0436988\\
1.996	-0.0436989\\
1.998	-0.0436989\\
2	-0.0436989\\
};
\addplot [color=mycolor2, line width=1.0pt, forget plot]
  table[row sep=crcr]{%
-0.024	0\\
-0.022	0\\
-0.02	0\\
-0.018	0\\
-0.016	0\\
-0.014	0\\
-0.012	0\\
-0.01	0\\
-0.008	0\\
-0.006	0\\
-0.004	0\\
-0.002	0\\
0	0\\
0.002	-0.0019122\\
0.004	-0.00365962\\
0.006	-0.0052577\\
0.008	-0.00672084\\
0.01	-0.00806242\\
0.012	-0.00929492\\
0.014	-0.0104299\\
0.016	-0.011478\\
0.018	-0.0124492\\
0.02	-0.0133525\\
0.022	-0.0141964\\
0.024	-0.0149883\\
0.026	-0.0157353\\
0.028	-0.0164436\\
0.03	-0.0171189\\
0.032	-0.0177661\\
0.034	-0.0183898\\
0.036	-0.0189939\\
0.038	-0.0195817\\
0.04	-0.0201561\\
0.042	-0.0207196\\
0.044	-0.0212741\\
0.046	-0.0218214\\
0.048	-0.0223626\\
0.05	-0.0228988\\
0.052	-0.0234305\\
0.054	-0.0239582\\
0.056	-0.024482\\
0.058	-0.0250021\\
0.06	-0.0255182\\
0.062	-0.02603\\
0.064	-0.0265374\\
0.066	-0.0270397\\
0.068	-0.0275367\\
0.07	-0.0280279\\
0.072	-0.0285128\\
0.074	-0.0289911\\
0.076	-0.0294623\\
0.078	-0.0299262\\
0.08	-0.0303824\\
0.082	-0.0308308\\
0.084	-0.0312713\\
0.086	-0.0317038\\
0.088	-0.0321282\\
0.09	-0.0325446\\
0.092	-0.0329531\\
0.094	-0.0333538\\
0.096	-0.0337469\\
0.098	-0.0341325\\
0.1	-0.0345109\\
0.102	-0.0348823\\
0.104	-0.0352468\\
0.106	-0.0356047\\
0.108	-0.035956\\
0.11	-0.036301\\
0.112	-0.0366396\\
0.114	-0.036972\\
0.116	-0.0372982\\
0.118	-0.0376182\\
0.12	-0.0379317\\
0.122	-0.0382388\\
0.124	-0.0385392\\
0.126	-0.0388328\\
0.128	-0.0391192\\
0.13	-0.0393981\\
0.132	-0.0396694\\
0.134	-0.0399326\\
0.136	-0.0401874\\
0.138	-0.0404334\\
0.14	-0.0406703\\
0.142	-0.0408977\\
0.144	-0.0411154\\
0.146	-0.041323\\
0.148	-0.0415203\\
0.15	-0.0417069\\
0.152	-0.0418828\\
0.154	-0.0420477\\
0.156	-0.0422016\\
0.158	-0.0423443\\
0.16	-0.0424758\\
0.162	-0.0425963\\
0.164	-0.0427058\\
0.166	-0.0428043\\
0.168	-0.0428923\\
0.17	-0.0429697\\
0.172	-0.043037\\
0.174	-0.0430945\\
0.176	-0.0431426\\
0.178	-0.0431815\\
0.18	-0.0432118\\
0.182	-0.043234\\
0.184	-0.0432484\\
0.186	-0.0432556\\
0.188	-0.0432561\\
0.19	-0.0432504\\
0.192	-0.043239\\
0.194	-0.0432224\\
0.196	-0.0432011\\
0.198	-0.0431757\\
0.2	-0.0431465\\
0.202	-0.0431141\\
0.204	-0.043079\\
0.206	-0.0430415\\
0.208	-0.043002\\
0.21	-0.0429611\\
0.212	-0.0429189\\
0.214	-0.042876\\
0.216	-0.0428325\\
0.218	-0.0427888\\
0.22	-0.0427452\\
0.222	-0.0427019\\
0.224	-0.0426592\\
0.226	-0.0426172\\
0.228	-0.0425761\\
0.23	-0.0425361\\
0.232	-0.0424973\\
0.234	-0.0424599\\
0.236	-0.042424\\
0.238	-0.0423896\\
0.24	-0.0423569\\
0.242	-0.0423258\\
0.244	-0.0422966\\
0.246	-0.0422691\\
0.248	-0.0422435\\
0.25	-0.0422198\\
0.252	-0.0421979\\
0.254	-0.0421779\\
0.256	-0.0421598\\
0.258	-0.0421436\\
0.26	-0.0421293\\
0.262	-0.0421169\\
0.264	-0.0421063\\
0.266	-0.0420975\\
0.268	-0.0420906\\
0.27	-0.0420854\\
0.272	-0.0420819\\
0.274	-0.0420802\\
0.276	-0.0420801\\
0.278	-0.0420817\\
0.28	-0.0420848\\
0.282	-0.0420894\\
0.284	-0.0420956\\
0.286	-0.0421032\\
0.288	-0.0421121\\
0.29	-0.0421224\\
0.292	-0.0421339\\
0.294	-0.0421467\\
0.296	-0.0421606\\
0.298	-0.0421757\\
0.3	-0.0421917\\
0.302	-0.0422088\\
0.304	-0.0422267\\
0.306	-0.0422456\\
0.308	-0.0422652\\
0.31	-0.0422855\\
0.312	-0.0423066\\
0.314	-0.0423283\\
0.316	-0.0423505\\
0.318	-0.0423732\\
0.32	-0.0423964\\
0.322	-0.0424201\\
0.324	-0.042444\\
0.326	-0.0424683\\
0.328	-0.0424928\\
0.33	-0.0425176\\
0.332	-0.0425425\\
0.334	-0.0425675\\
0.336	-0.0425927\\
0.338	-0.0426178\\
0.34	-0.0426431\\
0.342	-0.0426683\\
0.344	-0.0426934\\
0.346	-0.0427185\\
0.348	-0.0427435\\
0.35	-0.0427684\\
0.352	-0.0427932\\
0.354	-0.0428178\\
0.356	-0.0428422\\
0.358	-0.0428664\\
0.36	-0.0428905\\
0.362	-0.0429143\\
0.364	-0.0429379\\
0.366	-0.0429612\\
0.368	-0.0429843\\
0.37	-0.0430071\\
0.372	-0.0430297\\
0.374	-0.0430519\\
0.376	-0.0430739\\
0.378	-0.0430956\\
0.38	-0.0431171\\
0.382	-0.0431382\\
0.384	-0.043159\\
0.386	-0.0431795\\
0.388	-0.0431996\\
0.39	-0.0432195\\
0.392	-0.043239\\
0.394	-0.0432582\\
0.396	-0.0432771\\
0.398	-0.0432956\\
0.4	-0.0433138\\
0.402	-0.0433316\\
0.404	-0.0433491\\
0.406	-0.0433662\\
0.408	-0.043383\\
0.41	-0.0433994\\
0.412	-0.0434154\\
0.414	-0.043431\\
0.416	-0.0434463\\
0.418	-0.0434612\\
0.42	-0.0434757\\
0.422	-0.0434898\\
0.424	-0.0435035\\
0.426	-0.0435169\\
0.428	-0.0435298\\
0.43	-0.0435424\\
0.432	-0.0435545\\
0.434	-0.0435663\\
0.436	-0.0435776\\
0.438	-0.0435886\\
0.44	-0.0435991\\
0.442	-0.0436093\\
0.444	-0.0436191\\
0.446	-0.0436284\\
0.448	-0.0436374\\
0.45	-0.043646\\
0.452	-0.0436543\\
0.454	-0.0436621\\
0.456	-0.0436696\\
0.458	-0.0436767\\
0.46	-0.0436835\\
0.462	-0.0436899\\
0.464	-0.0436959\\
0.466	-0.0437017\\
0.468	-0.0437071\\
0.47	-0.0437121\\
0.472	-0.0437169\\
0.474	-0.0437214\\
0.476	-0.0437255\\
0.478	-0.0437294\\
0.48	-0.0437331\\
0.482	-0.0437364\\
0.484	-0.0437395\\
0.486	-0.0437424\\
0.488	-0.043745\\
0.49	-0.0437474\\
0.492	-0.0437496\\
0.494	-0.0437515\\
0.496	-0.0437533\\
0.498	-0.0437549\\
0.5	-0.0437563\\
0.502	-0.0437576\\
0.504	-0.0437587\\
0.506	-0.0437597\\
0.508	-0.0437605\\
0.51	-0.0437612\\
0.512	-0.0437617\\
0.514	-0.0437622\\
0.516	-0.0437625\\
0.518	-0.0437628\\
0.52	-0.0437629\\
0.522	-0.043763\\
0.524	-0.043763\\
0.526	-0.0437629\\
0.528	-0.0437627\\
0.53	-0.0437625\\
0.532	-0.0437622\\
0.534	-0.0437618\\
0.536	-0.0437614\\
0.538	-0.0437609\\
0.54	-0.0437604\\
0.542	-0.0437599\\
0.544	-0.0437593\\
0.546	-0.0437586\\
0.548	-0.0437579\\
0.55	-0.0437572\\
0.552	-0.0437564\\
0.554	-0.0437556\\
0.556	-0.0437548\\
0.558	-0.0437539\\
0.56	-0.043753\\
0.562	-0.0437521\\
0.564	-0.0437511\\
0.566	-0.04375\\
0.568	-0.043749\\
0.57	-0.0437479\\
0.572	-0.0437467\\
0.574	-0.0437455\\
0.576	-0.0437443\\
0.578	-0.043743\\
0.58	-0.0437417\\
0.582	-0.0437403\\
0.584	-0.0437389\\
0.586	-0.0437375\\
0.588	-0.043736\\
0.59	-0.0437344\\
0.592	-0.0437328\\
0.594	-0.0437312\\
0.596	-0.0437295\\
0.598	-0.0437278\\
0.6	-0.043726\\
0.602	-0.0437242\\
0.604	-0.0437223\\
0.606	-0.0437204\\
0.608	-0.0437184\\
0.61	-0.0437164\\
0.612	-0.0437144\\
0.614	-0.0437123\\
0.616	-0.0437102\\
0.618	-0.043708\\
0.62	-0.0437058\\
0.622	-0.0437036\\
0.624	-0.0437013\\
0.626	-0.043699\\
0.628	-0.0436967\\
0.63	-0.0436944\\
0.632	-0.043692\\
0.634	-0.0436896\\
0.636	-0.0436872\\
0.638	-0.0436848\\
0.64	-0.0436824\\
0.642	-0.04368\\
0.644	-0.0436776\\
0.646	-0.0436752\\
0.648	-0.0436728\\
0.65	-0.0436705\\
0.652	-0.0436681\\
0.654	-0.0436658\\
0.656	-0.0436635\\
0.658	-0.0436612\\
0.66	-0.0436589\\
0.662	-0.0436567\\
0.664	-0.0436546\\
0.666	-0.0436524\\
0.668	-0.0436504\\
0.67	-0.0436484\\
0.672	-0.0436464\\
0.674	-0.0436445\\
0.676	-0.0436426\\
0.678	-0.0436409\\
0.68	-0.0436392\\
0.682	-0.0436376\\
0.684	-0.043636\\
0.686	-0.0436345\\
0.688	-0.0436332\\
0.69	-0.0436319\\
0.692	-0.0436307\\
0.694	-0.0436295\\
0.696	-0.0436285\\
0.698	-0.0436276\\
0.7	-0.0436267\\
0.702	-0.043626\\
0.704	-0.0436254\\
0.706	-0.0436248\\
0.708	-0.0436244\\
0.71	-0.043624\\
0.712	-0.0436238\\
0.714	-0.0436237\\
0.716	-0.0436237\\
0.718	-0.0436237\\
0.72	-0.0436239\\
0.722	-0.0436242\\
0.724	-0.0436246\\
0.726	-0.0436251\\
0.728	-0.0436256\\
0.73	-0.0436263\\
0.732	-0.0436271\\
0.734	-0.043628\\
0.736	-0.043629\\
0.738	-0.04363\\
0.74	-0.0436312\\
0.742	-0.0436324\\
0.744	-0.0436338\\
0.746	-0.0436352\\
0.748	-0.0436367\\
0.75	-0.0436383\\
0.752	-0.0436399\\
0.754	-0.0436416\\
0.756	-0.0436434\\
0.758	-0.0436453\\
0.76	-0.0436472\\
0.762	-0.0436492\\
0.764	-0.0436512\\
0.766	-0.0436533\\
0.768	-0.0436555\\
0.77	-0.0436577\\
0.772	-0.0436599\\
0.774	-0.0436622\\
0.776	-0.0436645\\
0.778	-0.0436668\\
0.78	-0.0436691\\
0.782	-0.0436715\\
0.784	-0.0436739\\
0.786	-0.0436763\\
0.788	-0.0436787\\
0.79	-0.0436812\\
0.792	-0.0436836\\
0.794	-0.043686\\
0.796	-0.0436884\\
0.798	-0.0436908\\
0.8	-0.0436932\\
0.802	-0.0436956\\
0.804	-0.0436979\\
0.806	-0.0437003\\
0.808	-0.0437026\\
0.81	-0.0437048\\
0.812	-0.043707\\
0.814	-0.0437092\\
0.816	-0.0437113\\
0.818	-0.0437134\\
0.82	-0.0437155\\
0.822	-0.0437175\\
0.824	-0.0437194\\
0.826	-0.0437213\\
0.828	-0.0437231\\
0.83	-0.0437248\\
0.832	-0.0437265\\
0.834	-0.0437281\\
0.836	-0.0437296\\
0.838	-0.0437311\\
0.84	-0.0437324\\
0.842	-0.0437337\\
0.844	-0.043735\\
0.846	-0.0437361\\
0.848	-0.0437372\\
0.85	-0.0437381\\
0.852	-0.043739\\
0.854	-0.0437398\\
0.856	-0.0437405\\
0.858	-0.0437412\\
0.86	-0.0437417\\
0.862	-0.0437422\\
0.864	-0.0437425\\
0.866	-0.0437428\\
0.868	-0.043743\\
0.87	-0.0437431\\
0.872	-0.0437432\\
0.874	-0.0437431\\
0.876	-0.043743\\
0.878	-0.0437427\\
0.88	-0.0437424\\
0.882	-0.0437421\\
0.884	-0.0437416\\
0.886	-0.0437411\\
0.888	-0.0437405\\
0.89	-0.0437398\\
0.892	-0.043739\\
0.894	-0.0437382\\
0.896	-0.0437374\\
0.898	-0.0437364\\
0.9	-0.0437354\\
0.902	-0.0437344\\
0.904	-0.0437333\\
0.906	-0.0437321\\
0.908	-0.0437309\\
0.91	-0.0437297\\
0.912	-0.0437284\\
0.914	-0.0437271\\
0.916	-0.0437257\\
0.918	-0.0437243\\
0.92	-0.0437229\\
0.922	-0.0437215\\
0.924	-0.04372\\
0.926	-0.0437186\\
0.928	-0.0437171\\
0.93	-0.0437156\\
0.932	-0.0437141\\
0.934	-0.0437125\\
0.936	-0.043711\\
0.938	-0.0437095\\
0.94	-0.043708\\
0.942	-0.0437065\\
0.944	-0.043705\\
0.946	-0.0437036\\
0.948	-0.0437021\\
0.95	-0.0437007\\
0.952	-0.0436992\\
0.954	-0.0436979\\
0.956	-0.0436965\\
0.958	-0.0436952\\
0.96	-0.0436939\\
0.962	-0.0436926\\
0.964	-0.0436914\\
0.966	-0.0436902\\
0.968	-0.043689\\
0.97	-0.0436879\\
0.972	-0.0436869\\
0.974	-0.0436858\\
0.976	-0.0436849\\
0.978	-0.0436839\\
0.98	-0.0436831\\
0.982	-0.0436822\\
0.984	-0.0436814\\
0.986	-0.0436807\\
0.988	-0.04368\\
0.99	-0.0436794\\
0.992	-0.0436788\\
0.994	-0.0436783\\
0.996	-0.0436778\\
0.998	-0.0436774\\
1	-0.043677\\
1.002	-0.0436767\\
1.004	-0.0436764\\
1.006	-0.0436761\\
1.008	-0.043676\\
1.01	-0.0436758\\
1.012	-0.0436757\\
1.014	-0.0436757\\
1.016	-0.0436757\\
1.018	-0.0436757\\
1.02	-0.0436758\\
1.022	-0.043676\\
1.024	-0.0436761\\
1.026	-0.0436763\\
1.028	-0.0436766\\
1.03	-0.0436769\\
1.032	-0.0436772\\
1.034	-0.0436775\\
1.036	-0.0436779\\
1.038	-0.0436783\\
1.04	-0.0436787\\
1.042	-0.0436792\\
1.044	-0.0436797\\
1.046	-0.0436802\\
1.048	-0.0436807\\
1.05	-0.0436813\\
1.052	-0.0436818\\
1.054	-0.0436824\\
1.056	-0.043683\\
1.058	-0.0436836\\
1.06	-0.0436842\\
1.062	-0.0436848\\
1.064	-0.0436855\\
1.066	-0.0436861\\
1.068	-0.0436867\\
1.07	-0.0436874\\
1.072	-0.043688\\
1.074	-0.0436887\\
1.076	-0.0436893\\
1.078	-0.04369\\
1.08	-0.0436906\\
1.082	-0.0436913\\
1.084	-0.0436919\\
1.086	-0.0436925\\
1.088	-0.0436932\\
1.09	-0.0436938\\
1.092	-0.0436944\\
1.094	-0.043695\\
1.096	-0.0436956\\
1.098	-0.0436961\\
1.1	-0.0436967\\
1.102	-0.0436973\\
1.104	-0.0436978\\
1.106	-0.0436983\\
1.108	-0.0436988\\
1.11	-0.0436993\\
1.112	-0.0436998\\
1.114	-0.0437002\\
1.116	-0.0437007\\
1.118	-0.0437011\\
1.12	-0.0437015\\
1.122	-0.0437019\\
1.124	-0.0437023\\
1.126	-0.0437027\\
1.128	-0.043703\\
1.13	-0.0437033\\
1.132	-0.0437036\\
1.134	-0.0437039\\
1.136	-0.0437042\\
1.138	-0.0437045\\
1.14	-0.0437047\\
1.142	-0.0437049\\
1.144	-0.0437051\\
1.146	-0.0437053\\
1.148	-0.0437055\\
1.15	-0.0437057\\
1.152	-0.0437058\\
1.154	-0.0437059\\
1.156	-0.043706\\
1.158	-0.0437061\\
1.16	-0.0437062\\
1.162	-0.0437063\\
1.164	-0.0437063\\
1.166	-0.0437064\\
1.168	-0.0437064\\
1.17	-0.0437064\\
1.172	-0.0437064\\
1.174	-0.0437064\\
1.176	-0.0437063\\
1.178	-0.0437063\\
1.18	-0.0437062\\
1.182	-0.0437062\\
1.184	-0.0437061\\
1.186	-0.043706\\
1.188	-0.0437059\\
1.19	-0.0437058\\
1.192	-0.0437057\\
1.194	-0.0437055\\
1.196	-0.0437054\\
1.198	-0.0437053\\
1.2	-0.0437051\\
1.202	-0.0437049\\
1.204	-0.0437048\\
1.206	-0.0437046\\
1.208	-0.0437044\\
1.21	-0.0437042\\
1.212	-0.043704\\
1.214	-0.0437038\\
1.216	-0.0437036\\
1.218	-0.0437034\\
1.22	-0.0437032\\
1.222	-0.043703\\
1.224	-0.0437028\\
1.226	-0.0437025\\
1.228	-0.0437023\\
1.23	-0.0437021\\
1.232	-0.0437018\\
1.234	-0.0437016\\
1.236	-0.0437014\\
1.238	-0.0437011\\
1.24	-0.0437009\\
1.242	-0.0437007\\
1.244	-0.0437004\\
1.246	-0.0437002\\
1.248	-0.0436999\\
1.25	-0.0436997\\
1.252	-0.0436995\\
1.254	-0.0436992\\
1.256	-0.043699\\
1.258	-0.0436988\\
1.26	-0.0436986\\
1.262	-0.0436984\\
1.264	-0.0436981\\
1.266	-0.0436979\\
1.268	-0.0436977\\
1.27	-0.0436975\\
1.272	-0.0436973\\
1.274	-0.0436971\\
1.276	-0.0436969\\
1.278	-0.0436968\\
1.28	-0.0436966\\
1.282	-0.0436964\\
1.284	-0.0436962\\
1.286	-0.0436961\\
1.288	-0.0436959\\
1.29	-0.0436958\\
1.292	-0.0436957\\
1.294	-0.0436955\\
1.296	-0.0436954\\
1.298	-0.0436953\\
1.3	-0.0436952\\
1.302	-0.0436951\\
1.304	-0.043695\\
1.306	-0.0436949\\
1.308	-0.0436948\\
1.31	-0.0436947\\
1.312	-0.0436947\\
1.314	-0.0436946\\
1.316	-0.0436946\\
1.318	-0.0436945\\
1.32	-0.0436945\\
1.322	-0.0436945\\
1.324	-0.0436944\\
1.326	-0.0436944\\
1.328	-0.0436944\\
1.33	-0.0436944\\
1.332	-0.0436944\\
1.334	-0.0436944\\
1.336	-0.0436944\\
1.338	-0.0436945\\
1.34	-0.0436945\\
1.342	-0.0436945\\
1.344	-0.0436946\\
1.346	-0.0436946\\
1.348	-0.0436947\\
1.35	-0.0436947\\
1.352	-0.0436948\\
1.354	-0.0436948\\
1.356	-0.0436949\\
1.358	-0.043695\\
1.36	-0.0436951\\
1.362	-0.0436951\\
1.364	-0.0436952\\
1.366	-0.0436953\\
1.368	-0.0436954\\
1.37	-0.0436955\\
1.372	-0.0436956\\
1.374	-0.0436957\\
1.376	-0.0436957\\
1.378	-0.0436958\\
1.38	-0.0436959\\
1.382	-0.043696\\
1.384	-0.0436961\\
1.386	-0.0436962\\
1.388	-0.0436963\\
1.39	-0.0436964\\
1.392	-0.0436965\\
1.394	-0.0436966\\
1.396	-0.0436967\\
1.398	-0.0436968\\
1.4	-0.0436969\\
1.402	-0.043697\\
1.404	-0.0436971\\
1.406	-0.0436972\\
1.408	-0.0436973\\
1.41	-0.0436974\\
1.412	-0.0436975\\
1.414	-0.0436975\\
1.416	-0.0436976\\
1.418	-0.0436977\\
1.42	-0.0436978\\
1.422	-0.0436979\\
1.424	-0.0436979\\
1.426	-0.043698\\
1.428	-0.0436981\\
1.43	-0.0436981\\
1.432	-0.0436982\\
1.434	-0.0436983\\
1.436	-0.0436983\\
1.438	-0.0436984\\
1.44	-0.0436984\\
1.442	-0.0436985\\
1.444	-0.0436985\\
1.446	-0.0436986\\
1.448	-0.0436986\\
1.45	-0.0436986\\
1.452	-0.0436987\\
1.454	-0.0436987\\
1.456	-0.0436988\\
1.458	-0.0436988\\
1.46	-0.0436988\\
1.462	-0.0436988\\
1.464	-0.0436989\\
1.466	-0.0436989\\
1.468	-0.0436989\\
1.47	-0.0436989\\
1.472	-0.0436989\\
1.474	-0.0436989\\
1.476	-0.043699\\
1.478	-0.043699\\
1.48	-0.043699\\
1.482	-0.043699\\
1.484	-0.043699\\
1.486	-0.043699\\
1.488	-0.043699\\
1.49	-0.043699\\
1.492	-0.043699\\
1.494	-0.043699\\
1.496	-0.043699\\
1.498	-0.043699\\
1.5	-0.043699\\
1.502	-0.043699\\
1.504	-0.043699\\
1.506	-0.0436989\\
1.508	-0.0436989\\
1.51	-0.0436989\\
1.512	-0.0436989\\
1.514	-0.0436989\\
1.516	-0.0436989\\
1.518	-0.0436989\\
1.52	-0.0436989\\
1.522	-0.0436988\\
1.524	-0.0436988\\
1.526	-0.0436988\\
1.528	-0.0436988\\
1.53	-0.0436988\\
1.532	-0.0436988\\
1.534	-0.0436987\\
1.536	-0.0436987\\
1.538	-0.0436987\\
1.54	-0.0436987\\
1.542	-0.0436987\\
1.544	-0.0436986\\
1.546	-0.0436986\\
1.548	-0.0436986\\
1.55	-0.0436986\\
1.552	-0.0436986\\
1.554	-0.0436985\\
1.556	-0.0436985\\
1.558	-0.0436985\\
1.56	-0.0436985\\
1.562	-0.0436984\\
1.564	-0.0436984\\
1.566	-0.0436984\\
1.568	-0.0436984\\
1.57	-0.0436984\\
1.572	-0.0436983\\
1.574	-0.0436983\\
1.576	-0.0436983\\
1.578	-0.0436983\\
1.58	-0.0436982\\
1.582	-0.0436982\\
1.584	-0.0436982\\
1.586	-0.0436982\\
1.588	-0.0436981\\
1.59	-0.0436981\\
1.592	-0.0436981\\
1.594	-0.0436981\\
1.596	-0.043698\\
1.598	-0.043698\\
1.6	-0.043698\\
1.602	-0.043698\\
1.604	-0.043698\\
1.606	-0.0436979\\
1.608	-0.0436979\\
1.61	-0.0436979\\
1.612	-0.0436979\\
1.614	-0.0436978\\
1.616	-0.0436978\\
1.618	-0.0436978\\
1.62	-0.0436978\\
1.622	-0.0436978\\
1.624	-0.0436977\\
1.626	-0.0436977\\
1.628	-0.0436977\\
1.63	-0.0436977\\
1.632	-0.0436977\\
1.634	-0.0436976\\
1.636	-0.0436976\\
1.638	-0.0436976\\
1.64	-0.0436976\\
1.642	-0.0436976\\
1.644	-0.0436976\\
1.646	-0.0436975\\
1.648	-0.0436975\\
1.65	-0.0436975\\
1.652	-0.0436975\\
1.654	-0.0436975\\
1.656	-0.0436975\\
1.658	-0.0436975\\
1.66	-0.0436975\\
1.662	-0.0436975\\
1.664	-0.0436974\\
1.666	-0.0436974\\
1.668	-0.0436974\\
1.67	-0.0436974\\
1.672	-0.0436974\\
1.674	-0.0436974\\
1.676	-0.0436974\\
1.678	-0.0436974\\
1.68	-0.0436974\\
1.682	-0.0436974\\
1.684	-0.0436974\\
1.686	-0.0436974\\
1.688	-0.0436974\\
1.69	-0.0436974\\
1.692	-0.0436974\\
1.694	-0.0436974\\
1.696	-0.0436974\\
1.698	-0.0436974\\
1.7	-0.0436974\\
1.702	-0.0436974\\
1.704	-0.0436975\\
1.706	-0.0436975\\
1.708	-0.0436975\\
1.71	-0.0436975\\
1.712	-0.0436975\\
1.714	-0.0436975\\
1.716	-0.0436975\\
1.718	-0.0436975\\
1.72	-0.0436975\\
1.722	-0.0436975\\
1.724	-0.0436976\\
1.726	-0.0436976\\
1.728	-0.0436976\\
1.73	-0.0436976\\
1.732	-0.0436976\\
1.734	-0.0436976\\
1.736	-0.0436976\\
1.738	-0.0436977\\
1.74	-0.0436977\\
1.742	-0.0436977\\
1.744	-0.0436977\\
1.746	-0.0436977\\
1.748	-0.0436977\\
1.75	-0.0436978\\
1.752	-0.0436978\\
1.754	-0.0436978\\
1.756	-0.0436978\\
1.758	-0.0436978\\
1.76	-0.0436978\\
1.762	-0.0436979\\
1.764	-0.0436979\\
1.766	-0.0436979\\
1.768	-0.0436979\\
1.77	-0.0436979\\
1.772	-0.0436979\\
1.774	-0.0436979\\
1.776	-0.043698\\
1.778	-0.043698\\
1.78	-0.043698\\
1.782	-0.043698\\
1.784	-0.043698\\
1.786	-0.043698\\
1.788	-0.043698\\
1.79	-0.043698\\
1.792	-0.0436981\\
1.794	-0.0436981\\
1.796	-0.0436981\\
1.798	-0.0436981\\
1.8	-0.0436981\\
1.802	-0.0436981\\
1.804	-0.0436981\\
1.806	-0.0436981\\
1.808	-0.0436981\\
1.81	-0.0436981\\
1.812	-0.0436981\\
1.814	-0.0436981\\
1.816	-0.0436982\\
1.818	-0.0436982\\
1.82	-0.0436982\\
1.822	-0.0436982\\
1.824	-0.0436982\\
1.826	-0.0436982\\
1.828	-0.0436982\\
1.83	-0.0436982\\
1.832	-0.0436982\\
1.834	-0.0436982\\
1.836	-0.0436982\\
1.838	-0.0436982\\
1.84	-0.0436982\\
1.842	-0.0436982\\
1.844	-0.0436982\\
1.846	-0.0436982\\
1.848	-0.0436982\\
1.85	-0.0436982\\
1.852	-0.0436982\\
1.854	-0.0436982\\
1.856	-0.0436981\\
1.858	-0.0436981\\
1.86	-0.0436981\\
1.862	-0.0436981\\
1.864	-0.0436981\\
1.866	-0.0436981\\
1.868	-0.0436981\\
1.87	-0.0436981\\
1.872	-0.0436981\\
1.874	-0.0436981\\
1.876	-0.0436981\\
1.878	-0.0436981\\
1.88	-0.0436981\\
1.882	-0.0436981\\
1.884	-0.0436981\\
1.886	-0.0436981\\
1.888	-0.0436981\\
1.89	-0.043698\\
1.892	-0.043698\\
1.894	-0.043698\\
1.896	-0.043698\\
1.898	-0.043698\\
1.9	-0.043698\\
1.902	-0.043698\\
1.904	-0.043698\\
1.906	-0.043698\\
1.908	-0.043698\\
1.91	-0.043698\\
1.912	-0.043698\\
1.914	-0.043698\\
1.916	-0.0436979\\
1.918	-0.0436979\\
1.92	-0.0436979\\
1.922	-0.0436979\\
1.924	-0.0436979\\
1.926	-0.0436979\\
1.928	-0.0436979\\
1.93	-0.0436979\\
1.932	-0.0436979\\
1.934	-0.0436979\\
1.936	-0.0436979\\
1.938	-0.0436979\\
1.94	-0.0436979\\
1.942	-0.0436979\\
1.944	-0.0436979\\
1.946	-0.0436979\\
1.948	-0.0436979\\
1.95	-0.0436978\\
1.952	-0.0436978\\
1.954	-0.0436978\\
1.956	-0.0436978\\
1.958	-0.0436978\\
1.96	-0.0436978\\
1.962	-0.0436978\\
1.964	-0.0436978\\
1.966	-0.0436978\\
1.968	-0.0436978\\
1.97	-0.0436978\\
1.972	-0.0436978\\
1.974	-0.0436978\\
1.976	-0.0436978\\
1.978	-0.0436978\\
1.98	-0.0436978\\
1.982	-0.0436978\\
1.984	-0.0436978\\
1.986	-0.0436978\\
1.988	-0.0436978\\
1.99	-0.0436978\\
1.992	-0.0436978\\
1.994	-0.0436978\\
1.996	-0.0436978\\
1.998	-0.0436978\\
2	-0.0436978\\
};
\addplot [color=mycolor3, line width=1.0pt, forget plot]
  table[row sep=crcr]{%
-0.024	0\\
-0.022	0\\
-0.02	0\\
-0.018	0\\
-0.016	0\\
-0.014	0\\
-0.012	0\\
-0.01	0\\
-0.008	0\\
-0.006	0\\
-0.004	0\\
-0.002	0\\
0	0\\
0.002	-0.0019122\\
0.004	-0.00365962\\
0.006	-0.0052577\\
0.008	-0.00672083\\
0.01	-0.00806241\\
0.012	-0.00929489\\
0.014	-0.0104298\\
0.016	-0.011478\\
0.018	-0.0124491\\
0.02	-0.0133524\\
0.022	-0.0141962\\
0.024	-0.014988\\
0.026	-0.0157349\\
0.028	-0.0164431\\
0.03	-0.0171183\\
0.032	-0.0177654\\
0.034	-0.0183889\\
0.036	-0.0189928\\
0.038	-0.0195803\\
0.04	-0.0201545\\
0.042	-0.0207176\\
0.044	-0.0212719\\
0.046	-0.0218187\\
0.048	-0.0223595\\
0.05	-0.0228952\\
0.052	-0.0234263\\
0.054	-0.0239534\\
0.056	-0.0244765\\
0.058	-0.0249958\\
0.06	-0.025511\\
0.062	-0.026022\\
0.064	-0.0265283\\
0.066	-0.0270296\\
0.068	-0.0275254\\
0.07	-0.0280153\\
0.072	-0.0284988\\
0.074	-0.0289756\\
0.076	-0.0294452\\
0.078	-0.0299074\\
0.08	-0.0303618\\
0.082	-0.0308083\\
0.084	-0.0312467\\
0.086	-0.031677\\
0.088	-0.0320991\\
0.09	-0.0325131\\
0.092	-0.0329191\\
0.094	-0.0333171\\
0.096	-0.0337074\\
0.098	-0.0340902\\
0.1	-0.0344655\\
0.102	-0.0348338\\
0.104	-0.035195\\
0.106	-0.0355495\\
0.108	-0.0358973\\
0.11	-0.0362387\\
0.112	-0.0365737\\
0.114	-0.0369024\\
0.116	-0.0372247\\
0.118	-0.0375407\\
0.12	-0.0378503\\
0.122	-0.0381533\\
0.124	-0.0384496\\
0.126	-0.038739\\
0.128	-0.0390212\\
0.13	-0.039296\\
0.132	-0.0395631\\
0.134	-0.0398221\\
0.136	-0.0400727\\
0.138	-0.0403146\\
0.14	-0.0405475\\
0.142	-0.0407709\\
0.144	-0.0409847\\
0.146	-0.0411885\\
0.148	-0.041382\\
0.15	-0.0415651\\
0.152	-0.0417375\\
0.154	-0.0418991\\
0.156	-0.0420499\\
0.158	-0.0421896\\
0.16	-0.0423185\\
0.162	-0.0424364\\
0.164	-0.0425436\\
0.166	-0.0426402\\
0.168	-0.0427263\\
0.17	-0.0428022\\
0.172	-0.0428683\\
0.174	-0.0429249\\
0.176	-0.0429723\\
0.178	-0.0430109\\
0.18	-0.0430412\\
0.182	-0.0430637\\
0.184	-0.0430788\\
0.186	-0.0430871\\
0.188	-0.0430889\\
0.19	-0.043085\\
0.192	-0.0430757\\
0.194	-0.0430615\\
0.196	-0.0430431\\
0.198	-0.0430209\\
0.2	-0.0429953\\
0.202	-0.0429669\\
0.204	-0.0429361\\
0.206	-0.0429033\\
0.208	-0.0428689\\
0.21	-0.0428333\\
0.212	-0.042797\\
0.214	-0.0427601\\
0.216	-0.0427231\\
0.218	-0.0426863\\
0.22	-0.0426498\\
0.222	-0.0426139\\
0.224	-0.0425789\\
0.226	-0.0425449\\
0.228	-0.0425121\\
0.23	-0.0424806\\
0.232	-0.0424507\\
0.234	-0.0424223\\
0.236	-0.0423956\\
0.238	-0.0423707\\
0.24	-0.0423476\\
0.242	-0.0423263\\
0.244	-0.042307\\
0.246	-0.0422896\\
0.248	-0.0422742\\
0.25	-0.0422607\\
0.252	-0.0422492\\
0.254	-0.0422396\\
0.256	-0.0422319\\
0.258	-0.0422261\\
0.26	-0.0422222\\
0.262	-0.0422201\\
0.264	-0.0422199\\
0.266	-0.0422213\\
0.268	-0.0422245\\
0.27	-0.0422293\\
0.272	-0.0422358\\
0.274	-0.0422438\\
0.276	-0.0422532\\
0.278	-0.0422641\\
0.28	-0.0422764\\
0.282	-0.0422899\\
0.284	-0.0423047\\
0.286	-0.0423206\\
0.288	-0.0423377\\
0.29	-0.0423557\\
0.292	-0.0423747\\
0.294	-0.0423946\\
0.296	-0.0424153\\
0.298	-0.0424367\\
0.3	-0.0424588\\
0.302	-0.0424815\\
0.304	-0.0425047\\
0.306	-0.0425284\\
0.308	-0.0425524\\
0.31	-0.0425767\\
0.312	-0.0426013\\
0.314	-0.0426261\\
0.316	-0.042651\\
0.318	-0.042676\\
0.32	-0.0427009\\
0.322	-0.0427258\\
0.324	-0.0427507\\
0.326	-0.0427753\\
0.328	-0.0427998\\
0.33	-0.042824\\
0.332	-0.042848\\
0.334	-0.0428716\\
0.336	-0.0428949\\
0.338	-0.0429178\\
0.34	-0.0429403\\
0.342	-0.0429624\\
0.344	-0.042984\\
0.346	-0.0430052\\
0.348	-0.0430259\\
0.35	-0.0430461\\
0.352	-0.0430658\\
0.354	-0.0430849\\
0.356	-0.0431036\\
0.358	-0.0431217\\
0.36	-0.0431393\\
0.362	-0.0431565\\
0.364	-0.043173\\
0.366	-0.0431891\\
0.368	-0.0432047\\
0.37	-0.0432198\\
0.372	-0.0432344\\
0.374	-0.0432486\\
0.376	-0.0432622\\
0.378	-0.0432755\\
0.38	-0.0432883\\
0.382	-0.0433006\\
0.384	-0.0433126\\
0.386	-0.0433241\\
0.388	-0.0433353\\
0.39	-0.043346\\
0.392	-0.0433565\\
0.394	-0.0433665\\
0.396	-0.0433763\\
0.398	-0.0433857\\
0.4	-0.0433948\\
0.402	-0.0434036\\
0.404	-0.0434121\\
0.406	-0.0434203\\
0.408	-0.0434282\\
0.41	-0.0434359\\
0.412	-0.0434434\\
0.414	-0.0434506\\
0.416	-0.0434576\\
0.418	-0.0434643\\
0.42	-0.0434709\\
0.422	-0.0434772\\
0.424	-0.0434834\\
0.426	-0.0434893\\
0.428	-0.0434951\\
0.43	-0.0435006\\
0.432	-0.043506\\
0.434	-0.0435113\\
0.436	-0.0435163\\
0.438	-0.0435212\\
0.44	-0.043526\\
0.442	-0.0435306\\
0.444	-0.0435351\\
0.446	-0.0435394\\
0.448	-0.0435436\\
0.45	-0.0435476\\
0.452	-0.0435515\\
0.454	-0.0435553\\
0.456	-0.043559\\
0.458	-0.0435626\\
0.46	-0.043566\\
0.462	-0.0435694\\
0.464	-0.0435726\\
0.466	-0.0435757\\
0.468	-0.0435788\\
0.47	-0.0435817\\
0.472	-0.0435846\\
0.474	-0.0435873\\
0.476	-0.04359\\
0.478	-0.0435926\\
0.48	-0.0435952\\
0.482	-0.0435976\\
0.484	-0.0436\\
0.486	-0.0436023\\
0.488	-0.0436046\\
0.49	-0.0436068\\
0.492	-0.0436089\\
0.494	-0.043611\\
0.496	-0.043613\\
0.498	-0.0436149\\
0.5	-0.0436169\\
0.502	-0.0436187\\
0.504	-0.0436206\\
0.506	-0.0436223\\
0.508	-0.0436241\\
0.51	-0.0436258\\
0.512	-0.0436274\\
0.514	-0.0436291\\
0.516	-0.0436306\\
0.518	-0.0436322\\
0.52	-0.0436337\\
0.522	-0.0436352\\
0.524	-0.0436367\\
0.526	-0.0436381\\
0.528	-0.0436395\\
0.53	-0.0436409\\
0.532	-0.0436422\\
0.534	-0.0436435\\
0.536	-0.0436448\\
0.538	-0.043646\\
0.54	-0.0436473\\
0.542	-0.0436485\\
0.544	-0.0436496\\
0.546	-0.0436508\\
0.548	-0.0436519\\
0.55	-0.043653\\
0.552	-0.0436541\\
0.554	-0.0436551\\
0.556	-0.0436561\\
0.558	-0.0436571\\
0.56	-0.0436581\\
0.562	-0.043659\\
0.564	-0.0436599\\
0.566	-0.0436608\\
0.568	-0.0436617\\
0.57	-0.0436625\\
0.572	-0.0436633\\
0.574	-0.0436641\\
0.576	-0.0436649\\
0.578	-0.0436656\\
0.58	-0.0436664\\
0.582	-0.0436671\\
0.584	-0.0436677\\
0.586	-0.0436684\\
0.588	-0.043669\\
0.59	-0.0436696\\
0.592	-0.0436702\\
0.594	-0.0436708\\
0.596	-0.0436714\\
0.598	-0.0436719\\
0.6	-0.0436725\\
0.602	-0.043673\\
0.604	-0.0436735\\
0.606	-0.043674\\
0.608	-0.0436745\\
0.61	-0.0436749\\
0.612	-0.0436754\\
0.614	-0.0436758\\
0.616	-0.0436763\\
0.618	-0.0436767\\
0.62	-0.0436771\\
0.622	-0.0436776\\
0.624	-0.043678\\
0.626	-0.0436784\\
0.628	-0.0436788\\
0.63	-0.0436792\\
0.632	-0.0436796\\
0.634	-0.04368\\
0.636	-0.0436804\\
0.638	-0.0436809\\
0.64	-0.0436813\\
0.642	-0.0436817\\
0.644	-0.0436821\\
0.646	-0.0436825\\
0.648	-0.0436829\\
0.65	-0.0436833\\
0.652	-0.0436837\\
0.654	-0.0436841\\
0.656	-0.0436846\\
0.658	-0.043685\\
0.66	-0.0436854\\
0.662	-0.0436858\\
0.664	-0.0436862\\
0.666	-0.0436866\\
0.668	-0.0436871\\
0.67	-0.0436875\\
0.672	-0.0436879\\
0.674	-0.0436883\\
0.676	-0.0436887\\
0.678	-0.043689\\
0.68	-0.0436894\\
0.682	-0.0436898\\
0.684	-0.0436902\\
0.686	-0.0436905\\
0.688	-0.0436908\\
0.69	-0.0436912\\
0.692	-0.0436915\\
0.694	-0.0436918\\
0.696	-0.0436921\\
0.698	-0.0436923\\
0.7	-0.0436926\\
0.702	-0.0436928\\
0.704	-0.043693\\
0.706	-0.0436932\\
0.708	-0.0436934\\
0.71	-0.0436936\\
0.712	-0.0436937\\
0.714	-0.0436938\\
0.716	-0.0436939\\
0.718	-0.043694\\
0.72	-0.043694\\
0.722	-0.043694\\
0.724	-0.043694\\
0.726	-0.043694\\
0.728	-0.0436939\\
0.73	-0.0436939\\
0.732	-0.0436938\\
0.734	-0.0436937\\
0.736	-0.0436935\\
0.738	-0.0436934\\
0.74	-0.0436932\\
0.742	-0.043693\\
0.744	-0.0436928\\
0.746	-0.0436926\\
0.748	-0.0436923\\
0.75	-0.0436921\\
0.752	-0.0436918\\
0.754	-0.0436915\\
0.756	-0.0436913\\
0.758	-0.043691\\
0.76	-0.0436906\\
0.762	-0.0436903\\
0.764	-0.04369\\
0.766	-0.0436897\\
0.768	-0.0436894\\
0.77	-0.0436891\\
0.772	-0.0436887\\
0.774	-0.0436884\\
0.776	-0.0436881\\
0.778	-0.0436878\\
0.78	-0.0436875\\
0.782	-0.0436873\\
0.784	-0.043687\\
0.786	-0.0436867\\
0.788	-0.0436865\\
0.79	-0.0436863\\
0.792	-0.0436861\\
0.794	-0.0436859\\
0.796	-0.0436857\\
0.798	-0.0436856\\
0.8	-0.0436854\\
0.802	-0.0436853\\
0.804	-0.0436853\\
0.806	-0.0436852\\
0.808	-0.0436852\\
0.81	-0.0436852\\
0.812	-0.0436852\\
0.814	-0.0436852\\
0.816	-0.0436853\\
0.818	-0.0436854\\
0.82	-0.0436855\\
0.822	-0.0436857\\
0.824	-0.0436859\\
0.826	-0.0436861\\
0.828	-0.0436863\\
0.83	-0.0436866\\
0.832	-0.0436869\\
0.834	-0.0436872\\
0.836	-0.0436875\\
0.838	-0.0436879\\
0.84	-0.0436882\\
0.842	-0.0436886\\
0.844	-0.043689\\
0.846	-0.0436895\\
0.848	-0.0436899\\
0.85	-0.0436904\\
0.852	-0.0436908\\
0.854	-0.0436913\\
0.856	-0.0436918\\
0.858	-0.0436923\\
0.86	-0.0436928\\
0.862	-0.0436933\\
0.864	-0.0436938\\
0.866	-0.0436944\\
0.868	-0.0436949\\
0.87	-0.0436954\\
0.872	-0.0436959\\
0.874	-0.0436965\\
0.876	-0.043697\\
0.878	-0.0436975\\
0.88	-0.043698\\
0.882	-0.0436985\\
0.884	-0.0436989\\
0.886	-0.0436994\\
0.888	-0.0436998\\
0.89	-0.0437003\\
0.892	-0.0437007\\
0.894	-0.0437011\\
0.896	-0.0437015\\
0.898	-0.0437019\\
0.9	-0.0437022\\
0.902	-0.0437025\\
0.904	-0.0437028\\
0.906	-0.0437031\\
0.908	-0.0437034\\
0.91	-0.0437036\\
0.912	-0.0437039\\
0.914	-0.0437041\\
0.916	-0.0437042\\
0.918	-0.0437044\\
0.92	-0.0437045\\
0.922	-0.0437046\\
0.924	-0.0437047\\
0.926	-0.0437048\\
0.928	-0.0437048\\
0.93	-0.0437048\\
0.932	-0.0437048\\
0.934	-0.0437048\\
0.936	-0.0437048\\
0.938	-0.0437047\\
0.94	-0.0437046\\
0.942	-0.0437045\\
0.944	-0.0437044\\
0.946	-0.0437043\\
0.948	-0.0437041\\
0.95	-0.0437039\\
0.952	-0.0437038\\
0.954	-0.0437036\\
0.956	-0.0437034\\
0.958	-0.0437031\\
0.96	-0.0437029\\
0.962	-0.0437027\\
0.964	-0.0437024\\
0.966	-0.0437022\\
0.968	-0.0437019\\
0.97	-0.0437017\\
0.972	-0.0437014\\
0.974	-0.0437012\\
0.976	-0.0437009\\
0.978	-0.0437006\\
0.98	-0.0437004\\
0.982	-0.0437001\\
0.984	-0.0436998\\
0.986	-0.0436996\\
0.988	-0.0436993\\
0.99	-0.043699\\
0.992	-0.0436988\\
0.994	-0.0436986\\
0.996	-0.0436983\\
0.998	-0.0436981\\
1	-0.0436979\\
1.002	-0.0436976\\
1.004	-0.0436974\\
1.006	-0.0436972\\
1.008	-0.043697\\
1.01	-0.0436968\\
1.012	-0.0436967\\
1.014	-0.0436965\\
1.016	-0.0436964\\
1.018	-0.0436962\\
1.02	-0.0436961\\
1.022	-0.043696\\
1.024	-0.0436958\\
1.026	-0.0436957\\
1.028	-0.0436956\\
1.03	-0.0436956\\
1.032	-0.0436955\\
1.034	-0.0436954\\
1.036	-0.0436954\\
1.038	-0.0436953\\
1.04	-0.0436953\\
1.042	-0.0436952\\
1.044	-0.0436952\\
1.046	-0.0436952\\
1.048	-0.0436952\\
1.05	-0.0436952\\
1.052	-0.0436952\\
1.054	-0.0436952\\
1.056	-0.0436953\\
1.058	-0.0436953\\
1.06	-0.0436953\\
1.062	-0.0436954\\
1.064	-0.0436954\\
1.066	-0.0436955\\
1.068	-0.0436955\\
1.07	-0.0436956\\
1.072	-0.0436956\\
1.074	-0.0436957\\
1.076	-0.0436958\\
1.078	-0.0436958\\
1.08	-0.0436959\\
1.082	-0.043696\\
1.084	-0.043696\\
1.086	-0.0436961\\
1.088	-0.0436962\\
1.09	-0.0436963\\
1.092	-0.0436963\\
1.094	-0.0436964\\
1.096	-0.0436965\\
1.098	-0.0436966\\
1.1	-0.0436966\\
1.102	-0.0436967\\
1.104	-0.0436968\\
1.106	-0.0436969\\
1.108	-0.0436969\\
1.11	-0.043697\\
1.112	-0.0436971\\
1.114	-0.0436972\\
1.116	-0.0436972\\
1.118	-0.0436973\\
1.12	-0.0436973\\
1.122	-0.0436974\\
1.124	-0.0436975\\
1.126	-0.0436975\\
1.128	-0.0436976\\
1.13	-0.0436976\\
1.132	-0.0436977\\
1.134	-0.0436977\\
1.136	-0.0436978\\
1.138	-0.0436978\\
1.14	-0.0436979\\
1.142	-0.0436979\\
1.144	-0.043698\\
1.146	-0.043698\\
1.148	-0.0436981\\
1.15	-0.0436981\\
1.152	-0.0436981\\
1.154	-0.0436982\\
1.156	-0.0436982\\
1.158	-0.0436982\\
1.16	-0.0436983\\
1.162	-0.0436983\\
1.164	-0.0436983\\
1.166	-0.0436983\\
1.168	-0.0436984\\
1.17	-0.0436984\\
1.172	-0.0436984\\
1.174	-0.0436984\\
1.176	-0.0436985\\
1.178	-0.0436985\\
1.18	-0.0436985\\
1.182	-0.0436985\\
1.184	-0.0436985\\
1.186	-0.0436985\\
1.188	-0.0436985\\
1.19	-0.0436986\\
1.192	-0.0436986\\
1.194	-0.0436986\\
1.196	-0.0436986\\
1.198	-0.0436986\\
1.2	-0.0436986\\
1.202	-0.0436986\\
1.204	-0.0436986\\
1.206	-0.0436986\\
1.208	-0.0436986\\
1.21	-0.0436986\\
1.212	-0.0436986\\
1.214	-0.0436986\\
1.216	-0.0436986\\
1.218	-0.0436986\\
1.22	-0.0436986\\
1.222	-0.0436986\\
1.224	-0.0436986\\
1.226	-0.0436986\\
1.228	-0.0436985\\
1.23	-0.0436985\\
1.232	-0.0436985\\
1.234	-0.0436985\\
1.236	-0.0436985\\
1.238	-0.0436985\\
1.24	-0.0436985\\
1.242	-0.0436984\\
1.244	-0.0436984\\
1.246	-0.0436984\\
1.248	-0.0436984\\
1.25	-0.0436984\\
1.252	-0.0436983\\
1.254	-0.0436983\\
1.256	-0.0436983\\
1.258	-0.0436983\\
1.26	-0.0436982\\
1.262	-0.0436982\\
1.264	-0.0436982\\
1.266	-0.0436982\\
1.268	-0.0436981\\
1.27	-0.0436981\\
1.272	-0.0436981\\
1.274	-0.0436981\\
1.276	-0.043698\\
1.278	-0.043698\\
1.28	-0.043698\\
1.282	-0.043698\\
1.284	-0.043698\\
1.286	-0.0436979\\
1.288	-0.0436979\\
1.29	-0.0436979\\
1.292	-0.0436979\\
1.294	-0.0436978\\
1.296	-0.0436978\\
1.298	-0.0436978\\
1.3	-0.0436978\\
1.302	-0.0436978\\
1.304	-0.0436977\\
1.306	-0.0436977\\
1.308	-0.0436977\\
1.31	-0.0436977\\
1.312	-0.0436977\\
1.314	-0.0436977\\
1.316	-0.0436977\\
1.318	-0.0436976\\
1.32	-0.0436976\\
1.322	-0.0436976\\
1.324	-0.0436976\\
1.326	-0.0436976\\
1.328	-0.0436976\\
1.33	-0.0436976\\
1.332	-0.0436976\\
1.334	-0.0436976\\
1.336	-0.0436976\\
1.338	-0.0436976\\
1.34	-0.0436976\\
1.342	-0.0436976\\
1.344	-0.0436976\\
1.346	-0.0436976\\
1.348	-0.0436976\\
1.35	-0.0436976\\
1.352	-0.0436976\\
1.354	-0.0436976\\
1.356	-0.0436976\\
1.358	-0.0436976\\
1.36	-0.0436976\\
1.362	-0.0436977\\
1.364	-0.0436977\\
1.366	-0.0436977\\
1.368	-0.0436977\\
1.37	-0.0436977\\
1.372	-0.0436977\\
1.374	-0.0436977\\
1.376	-0.0436977\\
1.378	-0.0436977\\
1.38	-0.0436978\\
1.382	-0.0436978\\
1.384	-0.0436978\\
1.386	-0.0436978\\
1.388	-0.0436978\\
1.39	-0.0436978\\
1.392	-0.0436978\\
1.394	-0.0436978\\
1.396	-0.0436979\\
1.398	-0.0436979\\
1.4	-0.0436979\\
1.402	-0.0436979\\
1.404	-0.0436979\\
1.406	-0.0436979\\
1.408	-0.0436979\\
1.41	-0.0436979\\
1.412	-0.0436979\\
1.414	-0.0436979\\
1.416	-0.0436979\\
1.418	-0.043698\\
1.42	-0.043698\\
1.422	-0.043698\\
1.424	-0.043698\\
1.426	-0.043698\\
1.428	-0.043698\\
1.43	-0.043698\\
1.432	-0.043698\\
1.434	-0.043698\\
1.436	-0.043698\\
1.438	-0.043698\\
1.44	-0.043698\\
1.442	-0.043698\\
1.444	-0.043698\\
1.446	-0.043698\\
1.448	-0.043698\\
1.45	-0.043698\\
1.452	-0.043698\\
1.454	-0.043698\\
1.456	-0.043698\\
1.458	-0.043698\\
1.46	-0.043698\\
1.462	-0.043698\\
1.464	-0.043698\\
1.466	-0.043698\\
1.468	-0.043698\\
1.47	-0.043698\\
1.472	-0.043698\\
1.474	-0.043698\\
1.476	-0.043698\\
1.478	-0.043698\\
1.48	-0.043698\\
1.482	-0.043698\\
1.484	-0.043698\\
1.486	-0.043698\\
1.488	-0.043698\\
1.49	-0.043698\\
1.492	-0.043698\\
1.494	-0.043698\\
1.496	-0.043698\\
1.498	-0.0436979\\
1.5	-0.0436979\\
1.502	-0.0436979\\
1.504	-0.0436979\\
1.506	-0.0436979\\
1.508	-0.0436979\\
1.51	-0.0436979\\
1.512	-0.0436979\\
1.514	-0.0436979\\
1.516	-0.0436979\\
1.518	-0.0436979\\
1.52	-0.0436979\\
1.522	-0.0436979\\
1.524	-0.0436979\\
1.526	-0.0436979\\
1.528	-0.0436979\\
1.53	-0.0436979\\
1.532	-0.0436979\\
1.534	-0.0436979\\
1.536	-0.0436979\\
1.538	-0.0436979\\
1.54	-0.0436979\\
1.542	-0.0436979\\
1.544	-0.0436979\\
1.546	-0.0436979\\
1.548	-0.0436979\\
1.55	-0.0436979\\
1.552	-0.0436979\\
1.554	-0.0436979\\
1.556	-0.0436979\\
1.558	-0.0436979\\
1.56	-0.0436979\\
1.562	-0.0436979\\
1.564	-0.0436979\\
1.566	-0.0436979\\
1.568	-0.0436979\\
1.57	-0.0436979\\
1.572	-0.0436979\\
1.574	-0.0436979\\
1.576	-0.0436979\\
1.578	-0.0436979\\
1.58	-0.0436979\\
1.582	-0.0436979\\
1.584	-0.0436979\\
1.586	-0.0436979\\
1.588	-0.0436979\\
1.59	-0.0436979\\
1.592	-0.0436979\\
1.594	-0.0436979\\
1.596	-0.0436979\\
1.598	-0.0436979\\
1.6	-0.0436979\\
1.602	-0.0436979\\
1.604	-0.0436979\\
1.606	-0.0436979\\
1.608	-0.0436979\\
1.61	-0.0436979\\
1.612	-0.0436979\\
1.614	-0.0436979\\
1.616	-0.0436979\\
1.618	-0.0436979\\
1.62	-0.0436979\\
1.622	-0.0436979\\
1.624	-0.0436979\\
1.626	-0.0436979\\
1.628	-0.0436979\\
1.63	-0.0436979\\
1.632	-0.0436979\\
1.634	-0.0436979\\
1.636	-0.0436979\\
1.638	-0.0436979\\
1.64	-0.0436979\\
1.642	-0.0436979\\
1.644	-0.0436979\\
1.646	-0.0436979\\
1.648	-0.0436979\\
1.65	-0.0436979\\
1.652	-0.0436979\\
1.654	-0.0436979\\
1.656	-0.0436979\\
1.658	-0.0436979\\
1.66	-0.0436979\\
1.662	-0.0436979\\
1.664	-0.0436979\\
1.666	-0.0436979\\
1.668	-0.0436979\\
1.67	-0.0436979\\
1.672	-0.0436979\\
1.674	-0.0436979\\
1.676	-0.0436979\\
1.678	-0.0436979\\
1.68	-0.0436979\\
1.682	-0.0436979\\
1.684	-0.0436979\\
1.686	-0.0436979\\
1.688	-0.0436979\\
1.69	-0.0436979\\
1.692	-0.0436979\\
1.694	-0.0436979\\
1.696	-0.0436979\\
1.698	-0.0436979\\
1.7	-0.0436979\\
1.702	-0.0436979\\
1.704	-0.0436979\\
1.706	-0.0436979\\
1.708	-0.0436979\\
1.71	-0.0436979\\
1.712	-0.0436979\\
1.714	-0.0436979\\
1.716	-0.0436979\\
1.718	-0.0436979\\
1.72	-0.0436979\\
1.722	-0.0436979\\
1.724	-0.0436979\\
1.726	-0.0436979\\
1.728	-0.0436979\\
1.73	-0.0436979\\
1.732	-0.0436979\\
1.734	-0.0436979\\
1.736	-0.0436979\\
1.738	-0.0436979\\
1.74	-0.0436979\\
1.742	-0.0436979\\
1.744	-0.0436979\\
1.746	-0.0436979\\
1.748	-0.0436979\\
1.75	-0.0436979\\
1.752	-0.0436979\\
1.754	-0.0436979\\
1.756	-0.0436979\\
1.758	-0.0436979\\
1.76	-0.0436979\\
1.762	-0.0436979\\
1.764	-0.0436979\\
1.766	-0.0436979\\
1.768	-0.0436979\\
1.77	-0.0436979\\
1.772	-0.0436979\\
1.774	-0.0436979\\
1.776	-0.0436979\\
1.778	-0.0436979\\
1.78	-0.0436979\\
1.782	-0.0436979\\
1.784	-0.0436979\\
1.786	-0.0436979\\
1.788	-0.0436979\\
1.79	-0.0436979\\
1.792	-0.0436979\\
1.794	-0.0436979\\
1.796	-0.0436979\\
1.798	-0.0436979\\
1.8	-0.0436979\\
1.802	-0.0436979\\
1.804	-0.0436979\\
1.806	-0.0436979\\
1.808	-0.0436979\\
1.81	-0.0436979\\
1.812	-0.0436979\\
1.814	-0.0436979\\
1.816	-0.0436979\\
1.818	-0.0436979\\
1.82	-0.0436979\\
1.822	-0.0436979\\
1.824	-0.0436979\\
1.826	-0.0436979\\
1.828	-0.0436979\\
1.83	-0.0436979\\
1.832	-0.0436979\\
1.834	-0.0436979\\
1.836	-0.0436979\\
1.838	-0.0436979\\
1.84	-0.0436979\\
1.842	-0.0436979\\
1.844	-0.0436979\\
1.846	-0.0436979\\
1.848	-0.0436979\\
1.85	-0.0436979\\
1.852	-0.0436979\\
1.854	-0.0436979\\
1.856	-0.0436979\\
1.858	-0.0436979\\
1.86	-0.0436979\\
1.862	-0.0436979\\
1.864	-0.0436979\\
1.866	-0.0436979\\
1.868	-0.0436979\\
1.87	-0.0436979\\
1.872	-0.0436979\\
1.874	-0.0436979\\
1.876	-0.0436979\\
1.878	-0.0436979\\
1.88	-0.0436979\\
1.882	-0.0436979\\
1.884	-0.0436979\\
1.886	-0.0436979\\
1.888	-0.0436979\\
1.89	-0.0436979\\
1.892	-0.0436979\\
1.894	-0.0436979\\
1.896	-0.0436979\\
1.898	-0.0436979\\
1.9	-0.0436979\\
1.902	-0.0436979\\
1.904	-0.0436979\\
1.906	-0.0436979\\
1.908	-0.0436979\\
1.91	-0.0436979\\
1.912	-0.0436979\\
1.914	-0.0436979\\
1.916	-0.0436979\\
1.918	-0.0436979\\
1.92	-0.0436979\\
1.922	-0.0436979\\
1.924	-0.0436979\\
1.926	-0.0436979\\
1.928	-0.0436979\\
1.93	-0.0436979\\
1.932	-0.0436979\\
1.934	-0.0436979\\
1.936	-0.0436979\\
1.938	-0.0436979\\
1.94	-0.0436979\\
1.942	-0.0436979\\
1.944	-0.0436979\\
1.946	-0.0436979\\
1.948	-0.0436979\\
1.95	-0.0436979\\
1.952	-0.0436979\\
1.954	-0.0436979\\
1.956	-0.0436979\\
1.958	-0.0436979\\
1.96	-0.0436979\\
1.962	-0.0436979\\
1.964	-0.0436979\\
1.966	-0.0436979\\
1.968	-0.0436979\\
1.97	-0.0436979\\
1.972	-0.0436979\\
1.974	-0.0436979\\
1.976	-0.0436979\\
1.978	-0.0436979\\
1.98	-0.0436979\\
1.982	-0.0436979\\
1.984	-0.0436979\\
1.986	-0.0436979\\
1.988	-0.0436979\\
1.99	-0.0436979\\
1.992	-0.0436979\\
1.994	-0.0436979\\
1.996	-0.0436979\\
1.998	-0.0436979\\
2	-0.0436979\\
};
\addplot [color=mycolor4, line width=1.0pt, forget plot]
  table[row sep=crcr]{%
-0.024	0\\
-0.022	0\\
-0.02	0\\
-0.018	0\\
-0.016	0\\
-0.014	0\\
-0.012	0\\
-0.01	0\\
-0.008	0\\
-0.006	0\\
-0.004	0\\
-0.002	0\\
0	0\\
0.002	-0.0019122\\
0.004	-0.00365962\\
0.006	-0.0052577\\
0.008	-0.00672083\\
0.01	-0.0080624\\
0.012	-0.00929489\\
0.014	-0.0104298\\
0.016	-0.0114779\\
0.018	-0.0124491\\
0.02	-0.0133524\\
0.022	-0.0141961\\
0.024	-0.014988\\
0.026	-0.0157349\\
0.028	-0.0164431\\
0.03	-0.0171182\\
0.032	-0.0177653\\
0.034	-0.0183889\\
0.036	-0.0189927\\
0.038	-0.0195803\\
0.04	-0.0201545\\
0.042	-0.0207177\\
0.044	-0.0212719\\
0.046	-0.0218189\\
0.048	-0.0223597\\
0.05	-0.0228955\\
0.052	-0.0234268\\
0.054	-0.023954\\
0.056	-0.0244774\\
0.058	-0.0249968\\
0.06	-0.0255123\\
0.062	-0.0260236\\
0.064	-0.0265302\\
0.066	-0.0270319\\
0.068	-0.0275281\\
0.07	-0.0280185\\
0.072	-0.0285025\\
0.074	-0.0289799\\
0.076	-0.0294501\\
0.078	-0.0299129\\
0.08	-0.0303681\\
0.082	-0.0308154\\
0.084	-0.0312547\\
0.086	-0.0316858\\
0.088	-0.0321089\\
0.09	-0.0325239\\
0.092	-0.0329309\\
0.094	-0.03333\\
0.096	-0.0337214\\
0.098	-0.0341053\\
0.1	-0.0344818\\
0.102	-0.0348513\\
0.104	-0.0352137\\
0.106	-0.0355694\\
0.108	-0.0359185\\
0.11	-0.0362611\\
0.112	-0.0365973\\
0.114	-0.0369271\\
0.116	-0.0372507\\
0.118	-0.0375678\\
0.12	-0.0378785\\
0.122	-0.0381825\\
0.124	-0.0384798\\
0.126	-0.0387701\\
0.128	-0.0390532\\
0.13	-0.0393288\\
0.132	-0.0395965\\
0.134	-0.0398561\\
0.136	-0.0401072\\
0.138	-0.0403495\\
0.14	-0.0405826\\
0.142	-0.0408062\\
0.144	-0.04102\\
0.146	-0.0412237\\
0.148	-0.041417\\
0.15	-0.0415996\\
0.152	-0.0417715\\
0.154	-0.0419325\\
0.156	-0.0420824\\
0.158	-0.0422213\\
0.16	-0.042349\\
0.162	-0.0424657\\
0.164	-0.0425716\\
0.166	-0.0426666\\
0.168	-0.0427511\\
0.17	-0.0428253\\
0.172	-0.0428895\\
0.174	-0.0429441\\
0.176	-0.0429894\\
0.178	-0.0430259\\
0.18	-0.0430539\\
0.182	-0.0430741\\
0.184	-0.0430868\\
0.186	-0.0430925\\
0.188	-0.0430919\\
0.19	-0.0430853\\
0.192	-0.0430734\\
0.194	-0.0430567\\
0.196	-0.0430356\\
0.198	-0.0430107\\
0.2	-0.0429826\\
0.202	-0.0429516\\
0.204	-0.0429182\\
0.206	-0.0428829\\
0.208	-0.0428462\\
0.21	-0.0428083\\
0.212	-0.0427697\\
0.214	-0.0427307\\
0.216	-0.0426917\\
0.218	-0.042653\\
0.22	-0.0426147\\
0.222	-0.0425773\\
0.224	-0.0425408\\
0.226	-0.0425055\\
0.228	-0.0424716\\
0.23	-0.0424393\\
0.232	-0.0424086\\
0.234	-0.0423797\\
0.236	-0.0423526\\
0.238	-0.0423276\\
0.24	-0.0423045\\
0.242	-0.0422836\\
0.244	-0.0422648\\
0.246	-0.0422481\\
0.248	-0.0422336\\
0.25	-0.0422212\\
0.252	-0.042211\\
0.254	-0.0422029\\
0.256	-0.042197\\
0.258	-0.0421931\\
0.26	-0.0421913\\
0.262	-0.0421915\\
0.264	-0.0421937\\
0.266	-0.0421978\\
0.268	-0.0422037\\
0.27	-0.0422115\\
0.272	-0.0422209\\
0.274	-0.0422321\\
0.276	-0.0422448\\
0.278	-0.042259\\
0.28	-0.0422747\\
0.282	-0.0422917\\
0.284	-0.04231\\
0.286	-0.0423295\\
0.288	-0.0423501\\
0.29	-0.0423717\\
0.292	-0.0423943\\
0.294	-0.0424177\\
0.296	-0.0424419\\
0.298	-0.0424667\\
0.3	-0.0424922\\
0.302	-0.0425181\\
0.304	-0.0425444\\
0.306	-0.0425711\\
0.308	-0.042598\\
0.31	-0.0426251\\
0.312	-0.0426523\\
0.314	-0.0426795\\
0.316	-0.0427066\\
0.318	-0.0427336\\
0.32	-0.0427604\\
0.322	-0.042787\\
0.324	-0.0428132\\
0.326	-0.0428391\\
0.328	-0.0428645\\
0.33	-0.0428894\\
0.332	-0.0429139\\
0.334	-0.0429378\\
0.336	-0.042961\\
0.338	-0.0429837\\
0.34	-0.0430057\\
0.342	-0.0430271\\
0.344	-0.0430477\\
0.346	-0.0430676\\
0.348	-0.0430869\\
0.35	-0.0431054\\
0.352	-0.0431232\\
0.354	-0.0431402\\
0.356	-0.0431566\\
0.358	-0.0431722\\
0.36	-0.0431871\\
0.362	-0.0432013\\
0.364	-0.0432148\\
0.366	-0.0432277\\
0.368	-0.04324\\
0.37	-0.0432516\\
0.372	-0.0432626\\
0.374	-0.043273\\
0.376	-0.0432829\\
0.378	-0.0432922\\
0.38	-0.0433011\\
0.382	-0.0433095\\
0.384	-0.0433174\\
0.386	-0.0433249\\
0.388	-0.0433321\\
0.39	-0.0433388\\
0.392	-0.0433453\\
0.394	-0.0433514\\
0.396	-0.0433572\\
0.398	-0.0433628\\
0.4	-0.0433681\\
0.402	-0.0433733\\
0.404	-0.0433782\\
0.406	-0.043383\\
0.408	-0.0433877\\
0.41	-0.0433922\\
0.412	-0.0433966\\
0.414	-0.043401\\
0.416	-0.0434052\\
0.418	-0.0434095\\
0.42	-0.0434136\\
0.422	-0.0434178\\
0.424	-0.043422\\
0.426	-0.0434261\\
0.428	-0.0434303\\
0.43	-0.0434344\\
0.432	-0.0434386\\
0.434	-0.0434429\\
0.436	-0.0434471\\
0.438	-0.0434515\\
0.44	-0.0434558\\
0.442	-0.0434602\\
0.444	-0.0434647\\
0.446	-0.0434692\\
0.448	-0.0434738\\
0.45	-0.0434784\\
0.452	-0.043483\\
0.454	-0.0434878\\
0.456	-0.0434925\\
0.458	-0.0434973\\
0.46	-0.0435022\\
0.462	-0.043507\\
0.464	-0.043512\\
0.466	-0.0435169\\
0.468	-0.0435218\\
0.47	-0.0435268\\
0.472	-0.0435318\\
0.474	-0.0435368\\
0.476	-0.0435418\\
0.478	-0.0435467\\
0.48	-0.0435517\\
0.482	-0.0435566\\
0.484	-0.0435615\\
0.486	-0.0435664\\
0.488	-0.0435712\\
0.49	-0.043576\\
0.492	-0.0435807\\
0.494	-0.0435854\\
0.496	-0.04359\\
0.498	-0.0435946\\
0.5	-0.043599\\
0.502	-0.0436034\\
0.504	-0.0436077\\
0.506	-0.0436119\\
0.508	-0.043616\\
0.51	-0.04362\\
0.512	-0.0436239\\
0.514	-0.0436277\\
0.516	-0.0436313\\
0.518	-0.0436349\\
0.52	-0.0436383\\
0.522	-0.0436416\\
0.524	-0.0436448\\
0.526	-0.0436479\\
0.528	-0.0436509\\
0.53	-0.0436537\\
0.532	-0.0436564\\
0.534	-0.0436589\\
0.536	-0.0436614\\
0.538	-0.0436637\\
0.54	-0.0436659\\
0.542	-0.0436679\\
0.544	-0.0436699\\
0.546	-0.0436717\\
0.548	-0.0436734\\
0.55	-0.043675\\
0.552	-0.0436764\\
0.554	-0.0436777\\
0.556	-0.043679\\
0.558	-0.0436801\\
0.56	-0.0436811\\
0.562	-0.043682\\
0.564	-0.0436828\\
0.566	-0.0436836\\
0.568	-0.0436842\\
0.57	-0.0436847\\
0.572	-0.0436852\\
0.574	-0.0436856\\
0.576	-0.0436859\\
0.578	-0.0436861\\
0.58	-0.0436863\\
0.582	-0.0436864\\
0.584	-0.0436864\\
0.586	-0.0436864\\
0.588	-0.0436864\\
0.59	-0.0436863\\
0.592	-0.0436861\\
0.594	-0.0436859\\
0.596	-0.0436857\\
0.598	-0.0436855\\
0.6	-0.0436852\\
0.602	-0.0436849\\
0.604	-0.0436846\\
0.606	-0.0436843\\
0.608	-0.0436839\\
0.61	-0.0436836\\
0.612	-0.0436832\\
0.614	-0.0436829\\
0.616	-0.0436825\\
0.618	-0.0436822\\
0.62	-0.0436819\\
0.622	-0.0436815\\
0.624	-0.0436812\\
0.626	-0.0436809\\
0.628	-0.0436806\\
0.63	-0.0436803\\
0.632	-0.0436801\\
0.634	-0.0436798\\
0.636	-0.0436796\\
0.638	-0.0436794\\
0.64	-0.0436792\\
0.642	-0.043679\\
0.644	-0.0436789\\
0.646	-0.0436788\\
0.648	-0.0436787\\
0.65	-0.0436786\\
0.652	-0.0436786\\
0.654	-0.0436786\\
0.656	-0.0436786\\
0.658	-0.0436786\\
0.66	-0.0436787\\
0.662	-0.0436788\\
0.664	-0.0436788\\
0.666	-0.043679\\
0.668	-0.0436791\\
0.67	-0.0436793\\
0.672	-0.0436794\\
0.674	-0.0436796\\
0.676	-0.0436798\\
0.678	-0.0436801\\
0.68	-0.0436803\\
0.682	-0.0436805\\
0.684	-0.0436808\\
0.686	-0.0436811\\
0.688	-0.0436814\\
0.69	-0.0436816\\
0.692	-0.0436819\\
0.694	-0.0436822\\
0.696	-0.0436825\\
0.698	-0.0436829\\
0.7	-0.0436832\\
0.702	-0.0436835\\
0.704	-0.0436838\\
0.706	-0.0436841\\
0.708	-0.0436844\\
0.71	-0.0436847\\
0.712	-0.043685\\
0.714	-0.0436853\\
0.716	-0.0436856\\
0.718	-0.0436859\\
0.72	-0.0436862\\
0.722	-0.0436864\\
0.724	-0.0436867\\
0.726	-0.043687\\
0.728	-0.0436872\\
0.73	-0.0436874\\
0.732	-0.0436877\\
0.734	-0.0436879\\
0.736	-0.0436881\\
0.738	-0.0436883\\
0.74	-0.0436885\\
0.742	-0.0436886\\
0.744	-0.0436888\\
0.746	-0.0436889\\
0.748	-0.0436891\\
0.75	-0.0436892\\
0.752	-0.0436893\\
0.754	-0.0436895\\
0.756	-0.0436896\\
0.758	-0.0436897\\
0.76	-0.0436897\\
0.762	-0.0436898\\
0.764	-0.0436899\\
0.766	-0.04369\\
0.768	-0.04369\\
0.77	-0.0436901\\
0.772	-0.0436901\\
0.774	-0.0436902\\
0.776	-0.0436902\\
0.778	-0.0436903\\
0.78	-0.0436903\\
0.782	-0.0436903\\
0.784	-0.0436904\\
0.786	-0.0436904\\
0.788	-0.0436905\\
0.79	-0.0436905\\
0.792	-0.0436905\\
0.794	-0.0436906\\
0.796	-0.0436906\\
0.798	-0.0436906\\
0.8	-0.0436907\\
0.802	-0.0436907\\
0.804	-0.0436908\\
0.806	-0.0436909\\
0.808	-0.0436909\\
0.81	-0.043691\\
0.812	-0.0436911\\
0.814	-0.0436911\\
0.816	-0.0436912\\
0.818	-0.0436913\\
0.82	-0.0436914\\
0.822	-0.0436915\\
0.824	-0.0436917\\
0.826	-0.0436918\\
0.828	-0.0436919\\
0.83	-0.043692\\
0.832	-0.0436922\\
0.834	-0.0436923\\
0.836	-0.0436925\\
0.838	-0.0436926\\
0.84	-0.0436928\\
0.842	-0.043693\\
0.844	-0.0436932\\
0.846	-0.0436934\\
0.848	-0.0436935\\
0.85	-0.0436937\\
0.852	-0.0436939\\
0.854	-0.0436942\\
0.856	-0.0436944\\
0.858	-0.0436946\\
0.86	-0.0436948\\
0.862	-0.043695\\
0.864	-0.0436952\\
0.866	-0.0436955\\
0.868	-0.0436957\\
0.87	-0.0436959\\
0.872	-0.0436961\\
0.874	-0.0436964\\
0.876	-0.0436966\\
0.878	-0.0436968\\
0.88	-0.043697\\
0.882	-0.0436972\\
0.884	-0.0436975\\
0.886	-0.0436977\\
0.888	-0.0436979\\
0.89	-0.0436981\\
0.892	-0.0436983\\
0.894	-0.0436985\\
0.896	-0.0436987\\
0.898	-0.0436989\\
0.9	-0.043699\\
0.902	-0.0436992\\
0.904	-0.0436994\\
0.906	-0.0436995\\
0.908	-0.0436997\\
0.91	-0.0436998\\
0.912	-0.0437\\
0.914	-0.0437001\\
0.916	-0.0437002\\
0.918	-0.0437003\\
0.92	-0.0437004\\
0.922	-0.0437005\\
0.924	-0.0437006\\
0.926	-0.0437007\\
0.928	-0.0437008\\
0.93	-0.0437008\\
0.932	-0.0437009\\
0.934	-0.0437009\\
0.936	-0.043701\\
0.938	-0.043701\\
0.94	-0.043701\\
0.942	-0.043701\\
0.944	-0.043701\\
0.946	-0.043701\\
0.948	-0.043701\\
0.95	-0.043701\\
0.952	-0.0437009\\
0.954	-0.0437009\\
0.956	-0.0437009\\
0.958	-0.0437008\\
0.96	-0.0437007\\
0.962	-0.0437007\\
0.964	-0.0437006\\
0.966	-0.0437005\\
0.968	-0.0437005\\
0.97	-0.0437004\\
0.972	-0.0437003\\
0.974	-0.0437002\\
0.976	-0.0437001\\
0.978	-0.0437\\
0.98	-0.0436999\\
0.982	-0.0436998\\
0.984	-0.0436997\\
0.986	-0.0436996\\
0.988	-0.0436995\\
0.99	-0.0436993\\
0.992	-0.0436992\\
0.994	-0.0436991\\
0.996	-0.043699\\
0.998	-0.0436989\\
1	-0.0436988\\
1.002	-0.0436986\\
1.004	-0.0436985\\
1.006	-0.0436984\\
1.008	-0.0436983\\
1.01	-0.0436982\\
1.012	-0.0436981\\
1.014	-0.043698\\
1.016	-0.0436979\\
1.018	-0.0436978\\
1.02	-0.0436977\\
1.022	-0.0436976\\
1.024	-0.0436975\\
1.026	-0.0436974\\
1.028	-0.0436973\\
1.03	-0.0436972\\
1.032	-0.0436972\\
1.034	-0.0436971\\
1.036	-0.043697\\
1.038	-0.043697\\
1.04	-0.0436969\\
1.042	-0.0436968\\
1.044	-0.0436968\\
1.046	-0.0436967\\
1.048	-0.0436967\\
1.05	-0.0436966\\
1.052	-0.0436966\\
1.054	-0.0436966\\
1.056	-0.0436965\\
1.058	-0.0436965\\
1.06	-0.0436965\\
1.062	-0.0436965\\
1.064	-0.0436965\\
1.066	-0.0436965\\
1.068	-0.0436965\\
1.07	-0.0436964\\
1.072	-0.0436964\\
1.074	-0.0436965\\
1.076	-0.0436965\\
1.078	-0.0436965\\
1.08	-0.0436965\\
1.082	-0.0436965\\
1.084	-0.0436965\\
1.086	-0.0436965\\
1.088	-0.0436966\\
1.09	-0.0436966\\
1.092	-0.0436966\\
1.094	-0.0436966\\
1.096	-0.0436967\\
1.098	-0.0436967\\
1.1	-0.0436967\\
1.102	-0.0436968\\
1.104	-0.0436968\\
1.106	-0.0436968\\
1.108	-0.0436969\\
1.11	-0.0436969\\
1.112	-0.043697\\
1.114	-0.043697\\
1.116	-0.0436971\\
1.118	-0.0436971\\
1.12	-0.0436971\\
1.122	-0.0436972\\
1.124	-0.0436972\\
1.126	-0.0436973\\
1.128	-0.0436973\\
1.13	-0.0436974\\
1.132	-0.0436974\\
1.134	-0.0436974\\
1.136	-0.0436975\\
1.138	-0.0436975\\
1.14	-0.0436976\\
1.142	-0.0436976\\
1.144	-0.0436976\\
1.146	-0.0436977\\
1.148	-0.0436977\\
1.15	-0.0436978\\
1.152	-0.0436978\\
1.154	-0.0436978\\
1.156	-0.0436979\\
1.158	-0.0436979\\
1.16	-0.0436979\\
1.162	-0.0436979\\
1.164	-0.043698\\
1.166	-0.043698\\
1.168	-0.043698\\
1.17	-0.0436981\\
1.172	-0.0436981\\
1.174	-0.0436981\\
1.176	-0.0436981\\
1.178	-0.0436981\\
1.18	-0.0436982\\
1.182	-0.0436982\\
1.184	-0.0436982\\
1.186	-0.0436982\\
1.188	-0.0436982\\
1.19	-0.0436982\\
1.192	-0.0436982\\
1.194	-0.0436983\\
1.196	-0.0436983\\
1.198	-0.0436983\\
1.2	-0.0436983\\
1.202	-0.0436983\\
1.204	-0.0436983\\
1.206	-0.0436983\\
1.208	-0.0436983\\
1.21	-0.0436983\\
1.212	-0.0436983\\
1.214	-0.0436983\\
1.216	-0.0436983\\
1.218	-0.0436983\\
1.22	-0.0436983\\
1.222	-0.0436983\\
1.224	-0.0436983\\
1.226	-0.0436983\\
1.228	-0.0436983\\
1.23	-0.0436983\\
1.232	-0.0436983\\
1.234	-0.0436983\\
1.236	-0.0436983\\
1.238	-0.0436983\\
1.24	-0.0436982\\
1.242	-0.0436982\\
1.244	-0.0436982\\
1.246	-0.0436982\\
1.248	-0.0436982\\
1.25	-0.0436982\\
1.252	-0.0436982\\
1.254	-0.0436982\\
1.256	-0.0436982\\
1.258	-0.0436982\\
1.26	-0.0436982\\
1.262	-0.0436981\\
1.264	-0.0436981\\
1.266	-0.0436981\\
1.268	-0.0436981\\
1.27	-0.0436981\\
1.272	-0.0436981\\
1.274	-0.0436981\\
1.276	-0.0436981\\
1.278	-0.0436981\\
1.28	-0.043698\\
1.282	-0.043698\\
1.284	-0.043698\\
1.286	-0.043698\\
1.288	-0.043698\\
1.29	-0.043698\\
1.292	-0.043698\\
1.294	-0.043698\\
1.296	-0.043698\\
1.298	-0.043698\\
1.3	-0.043698\\
1.302	-0.0436979\\
1.304	-0.0436979\\
1.306	-0.0436979\\
1.308	-0.0436979\\
1.31	-0.0436979\\
1.312	-0.0436979\\
1.314	-0.0436979\\
1.316	-0.0436979\\
1.318	-0.0436979\\
1.32	-0.0436979\\
1.322	-0.0436979\\
1.324	-0.0436979\\
1.326	-0.0436979\\
1.328	-0.0436979\\
1.33	-0.0436979\\
1.332	-0.0436978\\
1.334	-0.0436978\\
1.336	-0.0436978\\
1.338	-0.0436978\\
1.34	-0.0436978\\
1.342	-0.0436978\\
1.344	-0.0436978\\
1.346	-0.0436978\\
1.348	-0.0436978\\
1.35	-0.0436978\\
1.352	-0.0436978\\
1.354	-0.0436978\\
1.356	-0.0436978\\
1.358	-0.0436978\\
1.36	-0.0436978\\
1.362	-0.0436978\\
1.364	-0.0436978\\
1.366	-0.0436978\\
1.368	-0.0436978\\
1.37	-0.0436978\\
1.372	-0.0436978\\
1.374	-0.0436978\\
1.376	-0.0436978\\
1.378	-0.0436978\\
1.38	-0.0436978\\
1.382	-0.0436978\\
1.384	-0.0436978\\
1.386	-0.0436978\\
1.388	-0.0436978\\
1.39	-0.0436978\\
1.392	-0.0436978\\
1.394	-0.0436978\\
1.396	-0.0436978\\
1.398	-0.0436978\\
1.4	-0.0436978\\
1.402	-0.0436978\\
1.404	-0.0436978\\
1.406	-0.0436978\\
1.408	-0.0436978\\
1.41	-0.0436978\\
1.412	-0.0436978\\
1.414	-0.0436978\\
1.416	-0.0436978\\
1.418	-0.0436979\\
1.42	-0.0436979\\
1.422	-0.0436979\\
1.424	-0.0436979\\
1.426	-0.0436979\\
1.428	-0.0436979\\
1.43	-0.0436979\\
1.432	-0.0436979\\
1.434	-0.0436979\\
1.436	-0.0436979\\
1.438	-0.0436979\\
1.44	-0.0436979\\
1.442	-0.0436979\\
1.444	-0.0436979\\
1.446	-0.0436979\\
1.448	-0.0436979\\
1.45	-0.0436979\\
1.452	-0.0436979\\
1.454	-0.0436979\\
1.456	-0.0436979\\
1.458	-0.0436979\\
1.46	-0.0436979\\
1.462	-0.0436979\\
1.464	-0.0436979\\
1.466	-0.0436979\\
1.468	-0.0436979\\
1.47	-0.0436979\\
1.472	-0.0436979\\
1.474	-0.0436979\\
1.476	-0.0436979\\
1.478	-0.0436979\\
1.48	-0.0436979\\
1.482	-0.0436979\\
1.484	-0.0436979\\
1.486	-0.0436979\\
1.488	-0.0436979\\
1.49	-0.0436979\\
1.492	-0.0436979\\
1.494	-0.0436979\\
1.496	-0.0436979\\
1.498	-0.0436979\\
1.5	-0.0436979\\
1.502	-0.0436979\\
1.504	-0.0436979\\
1.506	-0.0436979\\
1.508	-0.0436979\\
1.51	-0.0436979\\
1.512	-0.0436979\\
1.514	-0.0436979\\
1.516	-0.0436979\\
1.518	-0.0436979\\
1.52	-0.0436979\\
1.522	-0.0436979\\
1.524	-0.0436979\\
1.526	-0.0436979\\
1.528	-0.0436979\\
1.53	-0.0436979\\
1.532	-0.0436979\\
1.534	-0.0436979\\
1.536	-0.0436979\\
1.538	-0.0436979\\
1.54	-0.0436979\\
1.542	-0.0436979\\
1.544	-0.0436979\\
1.546	-0.0436979\\
1.548	-0.0436979\\
1.55	-0.0436979\\
1.552	-0.0436979\\
1.554	-0.0436979\\
1.556	-0.0436979\\
1.558	-0.0436979\\
1.56	-0.0436979\\
1.562	-0.0436979\\
1.564	-0.0436979\\
1.566	-0.0436979\\
1.568	-0.0436979\\
1.57	-0.0436979\\
1.572	-0.0436979\\
1.574	-0.0436979\\
1.576	-0.0436979\\
1.578	-0.0436979\\
1.58	-0.0436979\\
1.582	-0.0436979\\
1.584	-0.0436979\\
1.586	-0.0436979\\
1.588	-0.0436979\\
1.59	-0.0436979\\
1.592	-0.0436979\\
1.594	-0.0436979\\
1.596	-0.0436979\\
1.598	-0.0436979\\
1.6	-0.0436979\\
1.602	-0.0436979\\
1.604	-0.0436979\\
1.606	-0.0436979\\
1.608	-0.0436979\\
1.61	-0.0436979\\
1.612	-0.0436979\\
1.614	-0.0436979\\
1.616	-0.0436979\\
1.618	-0.0436979\\
1.62	-0.0436979\\
1.622	-0.0436979\\
1.624	-0.0436979\\
1.626	-0.0436979\\
1.628	-0.0436979\\
1.63	-0.0436979\\
1.632	-0.0436979\\
1.634	-0.0436979\\
1.636	-0.0436979\\
1.638	-0.0436979\\
1.64	-0.0436979\\
1.642	-0.0436979\\
1.644	-0.0436979\\
1.646	-0.0436979\\
1.648	-0.0436979\\
1.65	-0.0436979\\
1.652	-0.0436979\\
1.654	-0.0436979\\
1.656	-0.0436979\\
1.658	-0.0436979\\
1.66	-0.0436979\\
1.662	-0.0436979\\
1.664	-0.0436979\\
1.666	-0.0436979\\
1.668	-0.0436979\\
1.67	-0.0436979\\
1.672	-0.0436979\\
1.674	-0.0436979\\
1.676	-0.0436979\\
1.678	-0.0436979\\
1.68	-0.0436979\\
1.682	-0.0436979\\
1.684	-0.0436979\\
1.686	-0.0436979\\
1.688	-0.0436979\\
1.69	-0.0436979\\
1.692	-0.0436979\\
1.694	-0.0436979\\
1.696	-0.0436979\\
1.698	-0.0436979\\
1.7	-0.0436979\\
1.702	-0.0436979\\
1.704	-0.0436979\\
1.706	-0.0436979\\
1.708	-0.0436979\\
1.71	-0.0436979\\
1.712	-0.0436979\\
1.714	-0.0436979\\
1.716	-0.0436979\\
1.718	-0.0436979\\
1.72	-0.0436979\\
1.722	-0.0436979\\
1.724	-0.0436979\\
1.726	-0.0436979\\
1.728	-0.0436979\\
1.73	-0.0436979\\
1.732	-0.0436979\\
1.734	-0.0436979\\
1.736	-0.0436979\\
1.738	-0.0436979\\
1.74	-0.0436979\\
1.742	-0.0436979\\
1.744	-0.0436979\\
1.746	-0.0436979\\
1.748	-0.0436979\\
1.75	-0.0436979\\
1.752	-0.0436979\\
1.754	-0.0436979\\
1.756	-0.0436979\\
1.758	-0.0436979\\
1.76	-0.0436979\\
1.762	-0.0436979\\
1.764	-0.0436979\\
1.766	-0.0436979\\
1.768	-0.0436979\\
1.77	-0.0436979\\
1.772	-0.0436979\\
1.774	-0.0436979\\
1.776	-0.0436979\\
1.778	-0.0436979\\
1.78	-0.0436979\\
1.782	-0.0436979\\
1.784	-0.0436979\\
1.786	-0.0436979\\
1.788	-0.0436979\\
1.79	-0.0436979\\
1.792	-0.0436979\\
1.794	-0.0436979\\
1.796	-0.0436979\\
1.798	-0.0436979\\
1.8	-0.0436979\\
1.802	-0.0436979\\
1.804	-0.0436979\\
1.806	-0.0436979\\
1.808	-0.0436979\\
1.81	-0.0436979\\
1.812	-0.0436979\\
1.814	-0.0436979\\
1.816	-0.0436979\\
1.818	-0.0436979\\
1.82	-0.0436979\\
1.822	-0.0436979\\
1.824	-0.0436979\\
1.826	-0.0436979\\
1.828	-0.0436979\\
1.83	-0.0436979\\
1.832	-0.0436979\\
1.834	-0.0436979\\
1.836	-0.0436979\\
1.838	-0.0436979\\
1.84	-0.0436979\\
1.842	-0.0436979\\
1.844	-0.0436979\\
1.846	-0.0436979\\
1.848	-0.0436979\\
1.85	-0.0436979\\
1.852	-0.0436979\\
1.854	-0.0436979\\
1.856	-0.0436979\\
1.858	-0.0436979\\
1.86	-0.0436979\\
1.862	-0.0436979\\
1.864	-0.0436979\\
1.866	-0.0436979\\
1.868	-0.0436979\\
1.87	-0.0436979\\
1.872	-0.0436979\\
1.874	-0.0436979\\
1.876	-0.0436979\\
1.878	-0.0436979\\
1.88	-0.0436979\\
1.882	-0.0436979\\
1.884	-0.0436979\\
1.886	-0.0436979\\
1.888	-0.0436979\\
1.89	-0.0436979\\
1.892	-0.0436979\\
1.894	-0.0436979\\
1.896	-0.0436979\\
1.898	-0.0436979\\
1.9	-0.0436979\\
1.902	-0.0436979\\
1.904	-0.0436979\\
1.906	-0.0436979\\
1.908	-0.0436979\\
1.91	-0.0436979\\
1.912	-0.0436979\\
1.914	-0.0436979\\
1.916	-0.0436979\\
1.918	-0.0436979\\
1.92	-0.0436979\\
1.922	-0.0436979\\
1.924	-0.0436979\\
1.926	-0.0436979\\
1.928	-0.0436979\\
1.93	-0.0436979\\
1.932	-0.0436979\\
1.934	-0.0436979\\
1.936	-0.0436979\\
1.938	-0.0436979\\
1.94	-0.0436979\\
1.942	-0.0436979\\
1.944	-0.0436979\\
1.946	-0.0436979\\
1.948	-0.0436979\\
1.95	-0.0436979\\
1.952	-0.0436979\\
1.954	-0.0436979\\
1.956	-0.0436979\\
1.958	-0.0436979\\
1.96	-0.0436979\\
1.962	-0.0436979\\
1.964	-0.0436979\\
1.966	-0.0436979\\
1.968	-0.0436979\\
1.97	-0.0436979\\
1.972	-0.0436979\\
1.974	-0.0436979\\
1.976	-0.0436979\\
1.978	-0.0436979\\
1.98	-0.0436979\\
1.982	-0.0436979\\
1.984	-0.0436979\\
1.986	-0.0436979\\
1.988	-0.0436979\\
1.99	-0.0436979\\
1.992	-0.0436979\\
1.994	-0.0436979\\
1.996	-0.0436979\\
1.998	-0.0436979\\
2	-0.0436979\\
};
\addplot [color=mycolor5, line width=1.0pt, forget plot]
  table[row sep=crcr]{%
-0.024	0\\
-0.022	0\\
-0.02	0\\
-0.018	0\\
-0.016	0\\
-0.014	0\\
-0.012	0\\
-0.01	0\\
-0.008	0\\
-0.006	0\\
-0.004	0\\
-0.002	0\\
0	0\\
0.002	-0.0019122\\
0.004	-0.00365962\\
0.006	-0.00525769\\
0.008	-0.00672081\\
0.01	-0.00806236\\
0.012	-0.00929482\\
0.014	-0.0104297\\
0.016	-0.0114778\\
0.018	-0.0124488\\
0.02	-0.013352\\
0.022	-0.0141957\\
0.024	-0.0149874\\
0.026	-0.0157341\\
0.028	-0.0164421\\
0.03	-0.017117\\
0.032	-0.0177638\\
0.034	-0.0183869\\
0.036	-0.0189904\\
0.038	-0.0195775\\
0.04	-0.0201511\\
0.042	-0.0207138\\
0.044	-0.0212674\\
0.046	-0.0218136\\
0.048	-0.0223537\\
0.05	-0.0228886\\
0.052	-0.0234189\\
0.054	-0.0239452\\
0.056	-0.0244675\\
0.058	-0.0249858\\
0.06	-0.0255001\\
0.062	-0.0260101\\
0.064	-0.0265154\\
0.066	-0.0270157\\
0.068	-0.0275105\\
0.07	-0.0279994\\
0.072	-0.0284819\\
0.074	-0.0289577\\
0.076	-0.0294264\\
0.078	-0.0298877\\
0.08	-0.0303412\\
0.082	-0.0307869\\
0.084	-0.0312246\\
0.086	-0.0316542\\
0.088	-0.0320757\\
0.09	-0.0324892\\
0.092	-0.0328947\\
0.094	-0.0332924\\
0.096	-0.0336824\\
0.098	-0.034065\\
0.1	-0.0344403\\
0.102	-0.0348086\\
0.104	-0.03517\\
0.106	-0.0355248\\
0.108	-0.0358731\\
0.11	-0.0362149\\
0.112	-0.0365506\\
0.114	-0.03688\\
0.116	-0.0372032\\
0.118	-0.0375201\\
0.12	-0.0378308\\
0.122	-0.038135\\
0.124	-0.0384325\\
0.126	-0.0387233\\
0.128	-0.039007\\
0.13	-0.0392833\\
0.132	-0.039552\\
0.134	-0.0398127\\
0.136	-0.0400651\\
0.138	-0.0403088\\
0.14	-0.0405435\\
0.142	-0.0407689\\
0.144	-0.0409846\\
0.146	-0.0411903\\
0.148	-0.0413859\\
0.15	-0.041571\\
0.152	-0.0417454\\
0.154	-0.0419091\\
0.156	-0.0420618\\
0.158	-0.0422036\\
0.16	-0.0423345\\
0.162	-0.0424545\\
0.164	-0.0425636\\
0.166	-0.0426621\\
0.168	-0.0427501\\
0.17	-0.0428279\\
0.172	-0.0428959\\
0.174	-0.0429542\\
0.176	-0.0430033\\
0.178	-0.0430436\\
0.18	-0.0430756\\
0.182	-0.0430996\\
0.184	-0.0431162\\
0.186	-0.0431259\\
0.188	-0.0431292\\
0.19	-0.0431265\\
0.192	-0.0431184\\
0.194	-0.0431055\\
0.196	-0.0430882\\
0.198	-0.0430669\\
0.2	-0.0430423\\
0.202	-0.0430148\\
0.204	-0.0429848\\
0.206	-0.0429527\\
0.208	-0.042919\\
0.21	-0.0428841\\
0.212	-0.0428483\\
0.214	-0.042812\\
0.216	-0.0427754\\
0.218	-0.0427389\\
0.22	-0.0427027\\
0.222	-0.0426671\\
0.224	-0.0426322\\
0.226	-0.0425984\\
0.228	-0.0425657\\
0.23	-0.0425343\\
0.232	-0.0425044\\
0.234	-0.042476\\
0.236	-0.0424492\\
0.238	-0.0424242\\
0.24	-0.042401\\
0.242	-0.0423796\\
0.244	-0.0423601\\
0.246	-0.0423425\\
0.248	-0.0423269\\
0.25	-0.0423132\\
0.252	-0.0423014\\
0.254	-0.0422916\\
0.256	-0.0422837\\
0.258	-0.0422776\\
0.26	-0.0422735\\
0.262	-0.0422712\\
0.264	-0.0422706\\
0.266	-0.0422718\\
0.268	-0.0422748\\
0.27	-0.0422793\\
0.272	-0.0422855\\
0.274	-0.0422932\\
0.276	-0.0423024\\
0.278	-0.0423131\\
0.28	-0.042325\\
0.282	-0.0423383\\
0.284	-0.0423528\\
0.286	-0.0423685\\
0.288	-0.0423852\\
0.29	-0.0424029\\
0.292	-0.0424216\\
0.294	-0.0424411\\
0.296	-0.0424614\\
0.298	-0.0424824\\
0.3	-0.0425041\\
0.302	-0.0425263\\
0.304	-0.0425489\\
0.306	-0.042572\\
0.308	-0.0425954\\
0.31	-0.0426191\\
0.312	-0.042643\\
0.314	-0.042667\\
0.316	-0.0426911\\
0.318	-0.0427151\\
0.32	-0.0427391\\
0.322	-0.042763\\
0.324	-0.0427867\\
0.326	-0.0428102\\
0.328	-0.0428334\\
0.33	-0.0428563\\
0.332	-0.0428788\\
0.334	-0.042901\\
0.336	-0.0429227\\
0.338	-0.042944\\
0.34	-0.0429647\\
0.342	-0.042985\\
0.344	-0.0430047\\
0.346	-0.0430239\\
0.348	-0.0430426\\
0.35	-0.0430607\\
0.352	-0.0430782\\
0.354	-0.0430951\\
0.356	-0.0431115\\
0.358	-0.0431273\\
0.36	-0.0431425\\
0.362	-0.0431572\\
0.364	-0.0431713\\
0.366	-0.0431848\\
0.368	-0.0431979\\
0.37	-0.0432104\\
0.372	-0.0432224\\
0.374	-0.0432339\\
0.376	-0.043245\\
0.378	-0.0432556\\
0.38	-0.0432658\\
0.382	-0.0432756\\
0.384	-0.043285\\
0.386	-0.0432941\\
0.388	-0.0433027\\
0.39	-0.0433111\\
0.392	-0.0433191\\
0.394	-0.0433269\\
0.396	-0.0433344\\
0.398	-0.0433416\\
0.4	-0.0433486\\
0.402	-0.0433554\\
0.404	-0.0433619\\
0.406	-0.0433683\\
0.408	-0.0433746\\
0.41	-0.0433806\\
0.412	-0.0433866\\
0.414	-0.0433924\\
0.416	-0.0433981\\
0.418	-0.0434036\\
0.42	-0.0434091\\
0.422	-0.0434146\\
0.424	-0.0434199\\
0.426	-0.0434252\\
0.428	-0.0434304\\
0.43	-0.0434355\\
0.432	-0.0434406\\
0.434	-0.0434457\\
0.436	-0.0434507\\
0.438	-0.0434557\\
0.44	-0.0434607\\
0.442	-0.0434656\\
0.444	-0.0434705\\
0.446	-0.0434754\\
0.448	-0.0434803\\
0.45	-0.0434851\\
0.452	-0.0434899\\
0.454	-0.0434947\\
0.456	-0.0434995\\
0.458	-0.0435043\\
0.46	-0.043509\\
0.462	-0.0435137\\
0.464	-0.0435184\\
0.466	-0.043523\\
0.468	-0.0435276\\
0.47	-0.0435322\\
0.472	-0.0435367\\
0.474	-0.0435412\\
0.476	-0.0435457\\
0.478	-0.0435501\\
0.48	-0.0435544\\
0.482	-0.0435587\\
0.484	-0.043563\\
0.486	-0.0435672\\
0.488	-0.0435713\\
0.49	-0.0435754\\
0.492	-0.0435795\\
0.494	-0.0435834\\
0.496	-0.0435873\\
0.498	-0.0435911\\
0.5	-0.0435949\\
0.502	-0.0435985\\
0.504	-0.0436021\\
0.506	-0.0436056\\
0.508	-0.0436091\\
0.51	-0.0436124\\
0.512	-0.0436157\\
0.514	-0.0436188\\
0.516	-0.0436219\\
0.518	-0.0436249\\
0.52	-0.0436278\\
0.522	-0.0436307\\
0.524	-0.0436334\\
0.526	-0.043636\\
0.528	-0.0436386\\
0.53	-0.0436411\\
0.532	-0.0436434\\
0.534	-0.0436457\\
0.536	-0.0436479\\
0.538	-0.04365\\
0.54	-0.043652\\
0.542	-0.0436539\\
0.544	-0.0436557\\
0.546	-0.0436575\\
0.548	-0.0436592\\
0.55	-0.0436607\\
0.552	-0.0436622\\
0.554	-0.0436637\\
0.556	-0.043665\\
0.558	-0.0436663\\
0.56	-0.0436675\\
0.562	-0.0436686\\
0.564	-0.0436696\\
0.566	-0.0436706\\
0.568	-0.0436715\\
0.57	-0.0436724\\
0.572	-0.0436732\\
0.574	-0.0436739\\
0.576	-0.0436746\\
0.578	-0.0436752\\
0.58	-0.0436758\\
0.582	-0.0436763\\
0.584	-0.0436767\\
0.586	-0.0436772\\
0.588	-0.0436776\\
0.59	-0.0436779\\
0.592	-0.0436782\\
0.594	-0.0436785\\
0.596	-0.0436787\\
0.598	-0.043679\\
0.6	-0.0436791\\
0.602	-0.0436793\\
0.604	-0.0436795\\
0.606	-0.0436796\\
0.608	-0.0436797\\
0.61	-0.0436797\\
0.612	-0.0436798\\
0.614	-0.0436799\\
0.616	-0.0436799\\
0.618	-0.0436799\\
0.62	-0.04368\\
0.622	-0.04368\\
0.624	-0.04368\\
0.626	-0.04368\\
0.628	-0.04368\\
0.63	-0.04368\\
0.632	-0.04368\\
0.634	-0.04368\\
0.636	-0.04368\\
0.638	-0.04368\\
0.64	-0.04368\\
0.642	-0.04368\\
0.644	-0.04368\\
0.646	-0.04368\\
0.648	-0.0436801\\
0.65	-0.0436801\\
0.652	-0.0436801\\
0.654	-0.0436802\\
0.656	-0.0436802\\
0.658	-0.0436803\\
0.66	-0.0436803\\
0.662	-0.0436804\\
0.664	-0.0436805\\
0.666	-0.0436806\\
0.668	-0.0436807\\
0.67	-0.0436808\\
0.672	-0.0436809\\
0.674	-0.043681\\
0.676	-0.0436811\\
0.678	-0.0436812\\
0.68	-0.0436813\\
0.682	-0.0436815\\
0.684	-0.0436816\\
0.686	-0.0436817\\
0.688	-0.0436819\\
0.69	-0.043682\\
0.692	-0.0436822\\
0.694	-0.0436823\\
0.696	-0.0436825\\
0.698	-0.0436826\\
0.7	-0.0436828\\
0.702	-0.0436829\\
0.704	-0.0436831\\
0.706	-0.0436833\\
0.708	-0.0436834\\
0.71	-0.0436836\\
0.712	-0.0436837\\
0.714	-0.0436839\\
0.716	-0.0436841\\
0.718	-0.0436842\\
0.72	-0.0436844\\
0.722	-0.0436845\\
0.724	-0.0436847\\
0.726	-0.0436848\\
0.728	-0.043685\\
0.73	-0.0436851\\
0.732	-0.0436852\\
0.734	-0.0436854\\
0.736	-0.0436855\\
0.738	-0.0436856\\
0.74	-0.0436858\\
0.742	-0.0436859\\
0.744	-0.043686\\
0.746	-0.0436861\\
0.748	-0.0436863\\
0.75	-0.0436864\\
0.752	-0.0436865\\
0.754	-0.0436866\\
0.756	-0.0436867\\
0.758	-0.0436868\\
0.76	-0.0436869\\
0.762	-0.043687\\
0.764	-0.0436871\\
0.766	-0.0436872\\
0.768	-0.0436873\\
0.77	-0.0436875\\
0.772	-0.0436876\\
0.774	-0.0436877\\
0.776	-0.0436878\\
0.778	-0.0436879\\
0.78	-0.043688\\
0.782	-0.0436881\\
0.784	-0.0436882\\
0.786	-0.0436883\\
0.788	-0.0436884\\
0.79	-0.0436886\\
0.792	-0.0436887\\
0.794	-0.0436888\\
0.796	-0.0436889\\
0.798	-0.0436891\\
0.8	-0.0436892\\
0.802	-0.0436893\\
0.804	-0.0436895\\
0.806	-0.0436896\\
0.808	-0.0436898\\
0.81	-0.0436899\\
0.812	-0.0436901\\
0.814	-0.0436902\\
0.816	-0.0436904\\
0.818	-0.0436906\\
0.82	-0.0436908\\
0.822	-0.0436909\\
0.824	-0.0436911\\
0.826	-0.0436913\\
0.828	-0.0436915\\
0.83	-0.0436917\\
0.832	-0.0436919\\
0.834	-0.0436921\\
0.836	-0.0436923\\
0.838	-0.0436925\\
0.84	-0.0436927\\
0.842	-0.0436929\\
0.844	-0.0436931\\
0.846	-0.0436933\\
0.848	-0.0436935\\
0.85	-0.0436938\\
0.852	-0.043694\\
0.854	-0.0436942\\
0.856	-0.0436944\\
0.858	-0.0436946\\
0.86	-0.0436949\\
0.862	-0.0436951\\
0.864	-0.0436953\\
0.866	-0.0436955\\
0.868	-0.0436957\\
0.87	-0.0436959\\
0.872	-0.0436961\\
0.874	-0.0436963\\
0.876	-0.0436965\\
0.878	-0.0436968\\
0.88	-0.0436969\\
0.882	-0.0436971\\
0.884	-0.0436973\\
0.886	-0.0436975\\
0.888	-0.0436977\\
0.89	-0.0436979\\
0.892	-0.043698\\
0.894	-0.0436982\\
0.896	-0.0436984\\
0.898	-0.0436985\\
0.9	-0.0436987\\
0.902	-0.0436988\\
0.904	-0.0436989\\
0.906	-0.0436991\\
0.908	-0.0436992\\
0.91	-0.0436993\\
0.912	-0.0436994\\
0.914	-0.0436995\\
0.916	-0.0436996\\
0.918	-0.0436997\\
0.92	-0.0436998\\
0.922	-0.0436998\\
0.924	-0.0436999\\
0.926	-0.0437\\
0.928	-0.0437\\
0.93	-0.0437001\\
0.932	-0.0437001\\
0.934	-0.0437001\\
0.936	-0.0437001\\
0.938	-0.0437002\\
0.94	-0.0437002\\
0.942	-0.0437002\\
0.944	-0.0437002\\
0.946	-0.0437002\\
0.948	-0.0437001\\
0.95	-0.0437001\\
0.952	-0.0437001\\
0.954	-0.0437001\\
0.956	-0.0437\\
0.958	-0.0437\\
0.96	-0.0436999\\
0.962	-0.0436999\\
0.964	-0.0436998\\
0.966	-0.0436998\\
0.968	-0.0436997\\
0.97	-0.0436997\\
0.972	-0.0436996\\
0.974	-0.0436995\\
0.976	-0.0436994\\
0.978	-0.0436994\\
0.98	-0.0436993\\
0.982	-0.0436992\\
0.984	-0.0436991\\
0.986	-0.043699\\
0.988	-0.043699\\
0.99	-0.0436989\\
0.992	-0.0436988\\
0.994	-0.0436987\\
0.996	-0.0436986\\
0.998	-0.0436985\\
1	-0.0436984\\
1.002	-0.0436984\\
1.004	-0.0436983\\
1.006	-0.0436982\\
1.008	-0.0436981\\
1.01	-0.043698\\
1.012	-0.0436979\\
1.014	-0.0436979\\
1.016	-0.0436978\\
1.018	-0.0436977\\
1.02	-0.0436976\\
1.022	-0.0436976\\
1.024	-0.0436975\\
1.026	-0.0436974\\
1.028	-0.0436974\\
1.03	-0.0436973\\
1.032	-0.0436973\\
1.034	-0.0436972\\
1.036	-0.0436972\\
1.038	-0.0436971\\
1.04	-0.0436971\\
1.042	-0.043697\\
1.044	-0.043697\\
1.046	-0.0436969\\
1.048	-0.0436969\\
1.05	-0.0436969\\
1.052	-0.0436968\\
1.054	-0.0436968\\
1.056	-0.0436968\\
1.058	-0.0436968\\
1.06	-0.0436968\\
1.062	-0.0436967\\
1.064	-0.0436967\\
1.066	-0.0436967\\
1.068	-0.0436967\\
1.07	-0.0436967\\
1.072	-0.0436967\\
1.074	-0.0436967\\
1.076	-0.0436967\\
1.078	-0.0436967\\
1.08	-0.0436967\\
1.082	-0.0436967\\
1.084	-0.0436967\\
1.086	-0.0436968\\
1.088	-0.0436968\\
1.09	-0.0436968\\
1.092	-0.0436968\\
1.094	-0.0436968\\
1.096	-0.0436968\\
1.098	-0.0436969\\
1.1	-0.0436969\\
1.102	-0.0436969\\
1.104	-0.0436969\\
1.106	-0.043697\\
1.108	-0.043697\\
1.11	-0.043697\\
1.112	-0.0436971\\
1.114	-0.0436971\\
1.116	-0.0436971\\
1.118	-0.0436972\\
1.12	-0.0436972\\
1.122	-0.0436972\\
1.124	-0.0436972\\
1.126	-0.0436973\\
1.128	-0.0436973\\
1.13	-0.0436973\\
1.132	-0.0436974\\
1.134	-0.0436974\\
1.136	-0.0436974\\
1.138	-0.0436975\\
1.14	-0.0436975\\
1.142	-0.0436975\\
1.144	-0.0436976\\
1.146	-0.0436976\\
1.148	-0.0436976\\
1.15	-0.0436977\\
1.152	-0.0436977\\
1.154	-0.0436977\\
1.156	-0.0436977\\
1.158	-0.0436978\\
1.16	-0.0436978\\
1.162	-0.0436978\\
1.164	-0.0436978\\
1.166	-0.0436979\\
1.168	-0.0436979\\
1.17	-0.0436979\\
1.172	-0.0436979\\
1.174	-0.043698\\
1.176	-0.043698\\
1.178	-0.043698\\
1.18	-0.043698\\
1.182	-0.043698\\
1.184	-0.043698\\
1.186	-0.0436981\\
1.188	-0.0436981\\
1.19	-0.0436981\\
1.192	-0.0436981\\
1.194	-0.0436981\\
1.196	-0.0436981\\
1.198	-0.0436981\\
1.2	-0.0436981\\
1.202	-0.0436982\\
1.204	-0.0436982\\
1.206	-0.0436982\\
1.208	-0.0436982\\
1.21	-0.0436982\\
1.212	-0.0436982\\
1.214	-0.0436982\\
1.216	-0.0436982\\
1.218	-0.0436982\\
1.22	-0.0436982\\
1.222	-0.0436982\\
1.224	-0.0436982\\
1.226	-0.0436982\\
1.228	-0.0436982\\
1.23	-0.0436982\\
1.232	-0.0436982\\
1.234	-0.0436982\\
1.236	-0.0436982\\
1.238	-0.0436982\\
1.24	-0.0436982\\
1.242	-0.0436982\\
1.244	-0.0436982\\
1.246	-0.0436982\\
1.248	-0.0436982\\
1.25	-0.0436982\\
1.252	-0.0436982\\
1.254	-0.0436982\\
1.256	-0.0436981\\
1.258	-0.0436981\\
1.26	-0.0436981\\
1.262	-0.0436981\\
1.264	-0.0436981\\
1.266	-0.0436981\\
1.268	-0.0436981\\
1.27	-0.0436981\\
1.272	-0.0436981\\
1.274	-0.0436981\\
1.276	-0.0436981\\
1.278	-0.0436981\\
1.28	-0.0436981\\
1.282	-0.043698\\
1.284	-0.043698\\
1.286	-0.043698\\
1.288	-0.043698\\
1.29	-0.043698\\
1.292	-0.043698\\
1.294	-0.043698\\
1.296	-0.043698\\
1.298	-0.043698\\
1.3	-0.043698\\
1.302	-0.043698\\
1.304	-0.043698\\
1.306	-0.043698\\
1.308	-0.0436979\\
1.31	-0.0436979\\
1.312	-0.0436979\\
1.314	-0.0436979\\
1.316	-0.0436979\\
1.318	-0.0436979\\
1.32	-0.0436979\\
1.322	-0.0436979\\
1.324	-0.0436979\\
1.326	-0.0436979\\
1.328	-0.0436979\\
1.33	-0.0436979\\
1.332	-0.0436979\\
1.334	-0.0436979\\
1.336	-0.0436979\\
1.338	-0.0436979\\
1.34	-0.0436979\\
1.342	-0.0436979\\
1.344	-0.0436979\\
1.346	-0.0436978\\
1.348	-0.0436978\\
1.35	-0.0436978\\
1.352	-0.0436978\\
1.354	-0.0436978\\
1.356	-0.0436978\\
1.358	-0.0436978\\
1.36	-0.0436978\\
1.362	-0.0436978\\
1.364	-0.0436978\\
1.366	-0.0436978\\
1.368	-0.0436978\\
1.37	-0.0436978\\
1.372	-0.0436978\\
1.374	-0.0436978\\
1.376	-0.0436978\\
1.378	-0.0436978\\
1.38	-0.0436978\\
1.382	-0.0436978\\
1.384	-0.0436978\\
1.386	-0.0436978\\
1.388	-0.0436978\\
1.39	-0.0436978\\
1.392	-0.0436978\\
1.394	-0.0436978\\
1.396	-0.0436978\\
1.398	-0.0436978\\
1.4	-0.0436978\\
1.402	-0.0436978\\
1.404	-0.0436978\\
1.406	-0.0436978\\
1.408	-0.0436978\\
1.41	-0.0436978\\
1.412	-0.0436978\\
1.414	-0.0436978\\
1.416	-0.0436978\\
1.418	-0.0436979\\
1.42	-0.0436979\\
1.422	-0.0436979\\
1.424	-0.0436979\\
1.426	-0.0436979\\
1.428	-0.0436979\\
1.43	-0.0436979\\
1.432	-0.0436979\\
1.434	-0.0436979\\
1.436	-0.0436979\\
1.438	-0.0436979\\
1.44	-0.0436979\\
1.442	-0.0436979\\
1.444	-0.0436979\\
1.446	-0.0436979\\
1.448	-0.0436979\\
1.45	-0.0436979\\
1.452	-0.0436979\\
1.454	-0.0436979\\
1.456	-0.0436979\\
1.458	-0.0436979\\
1.46	-0.0436979\\
1.462	-0.0436979\\
1.464	-0.0436979\\
1.466	-0.0436979\\
1.468	-0.0436979\\
1.47	-0.0436979\\
1.472	-0.0436979\\
1.474	-0.0436979\\
1.476	-0.0436979\\
1.478	-0.0436979\\
1.48	-0.0436979\\
1.482	-0.0436979\\
1.484	-0.0436979\\
1.486	-0.0436979\\
1.488	-0.0436979\\
1.49	-0.0436979\\
1.492	-0.0436979\\
1.494	-0.0436979\\
1.496	-0.0436979\\
1.498	-0.0436979\\
1.5	-0.0436979\\
1.502	-0.0436979\\
1.504	-0.0436979\\
1.506	-0.0436979\\
1.508	-0.0436979\\
1.51	-0.0436979\\
1.512	-0.0436979\\
1.514	-0.0436979\\
1.516	-0.0436979\\
1.518	-0.0436979\\
1.52	-0.0436979\\
1.522	-0.0436979\\
1.524	-0.0436979\\
1.526	-0.0436979\\
1.528	-0.0436979\\
1.53	-0.0436979\\
1.532	-0.0436979\\
1.534	-0.0436979\\
1.536	-0.0436979\\
1.538	-0.0436979\\
1.54	-0.0436979\\
1.542	-0.0436979\\
1.544	-0.0436979\\
1.546	-0.0436979\\
1.548	-0.0436979\\
1.55	-0.0436979\\
1.552	-0.0436979\\
1.554	-0.0436979\\
1.556	-0.0436979\\
1.558	-0.0436979\\
1.56	-0.0436979\\
1.562	-0.0436979\\
1.564	-0.0436979\\
1.566	-0.0436979\\
1.568	-0.0436979\\
1.57	-0.0436979\\
1.572	-0.0436979\\
1.574	-0.0436979\\
1.576	-0.0436979\\
1.578	-0.0436979\\
1.58	-0.0436979\\
1.582	-0.0436979\\
1.584	-0.0436979\\
1.586	-0.0436979\\
1.588	-0.0436979\\
1.59	-0.0436979\\
1.592	-0.0436979\\
1.594	-0.0436979\\
1.596	-0.0436979\\
1.598	-0.0436979\\
1.6	-0.0436979\\
1.602	-0.0436979\\
1.604	-0.0436979\\
1.606	-0.0436979\\
1.608	-0.0436979\\
1.61	-0.0436979\\
1.612	-0.0436979\\
1.614	-0.0436979\\
1.616	-0.0436979\\
1.618	-0.0436979\\
1.62	-0.0436979\\
1.622	-0.0436979\\
1.624	-0.0436979\\
1.626	-0.0436979\\
1.628	-0.0436979\\
1.63	-0.0436979\\
1.632	-0.0436979\\
1.634	-0.0436979\\
1.636	-0.0436979\\
1.638	-0.0436979\\
1.64	-0.0436979\\
1.642	-0.0436979\\
1.644	-0.0436979\\
1.646	-0.0436979\\
1.648	-0.0436979\\
1.65	-0.0436979\\
1.652	-0.0436979\\
1.654	-0.0436979\\
1.656	-0.0436979\\
1.658	-0.0436979\\
1.66	-0.0436979\\
1.662	-0.0436979\\
1.664	-0.0436979\\
1.666	-0.0436979\\
1.668	-0.0436979\\
1.67	-0.0436979\\
1.672	-0.0436979\\
1.674	-0.0436979\\
1.676	-0.0436979\\
1.678	-0.0436979\\
1.68	-0.0436979\\
1.682	-0.0436979\\
1.684	-0.0436979\\
1.686	-0.0436979\\
1.688	-0.0436979\\
1.69	-0.0436979\\
1.692	-0.0436979\\
1.694	-0.0436979\\
1.696	-0.0436979\\
1.698	-0.0436979\\
1.7	-0.0436979\\
1.702	-0.0436979\\
1.704	-0.0436979\\
1.706	-0.0436979\\
1.708	-0.0436979\\
1.71	-0.0436979\\
1.712	-0.0436979\\
1.714	-0.0436979\\
1.716	-0.0436979\\
1.718	-0.0436979\\
1.72	-0.0436979\\
1.722	-0.0436979\\
1.724	-0.0436979\\
1.726	-0.0436979\\
1.728	-0.0436979\\
1.73	-0.0436979\\
1.732	-0.0436979\\
1.734	-0.0436979\\
1.736	-0.0436979\\
1.738	-0.0436979\\
1.74	-0.0436979\\
1.742	-0.0436979\\
1.744	-0.0436979\\
1.746	-0.0436979\\
1.748	-0.0436979\\
1.75	-0.0436979\\
1.752	-0.0436979\\
1.754	-0.0436979\\
1.756	-0.0436979\\
1.758	-0.0436979\\
1.76	-0.0436979\\
1.762	-0.0436979\\
1.764	-0.0436979\\
1.766	-0.0436979\\
1.768	-0.0436979\\
1.77	-0.0436979\\
1.772	-0.0436979\\
1.774	-0.0436979\\
1.776	-0.0436979\\
1.778	-0.0436979\\
1.78	-0.0436979\\
1.782	-0.0436979\\
1.784	-0.0436979\\
1.786	-0.0436979\\
1.788	-0.0436979\\
1.79	-0.0436979\\
1.792	-0.0436979\\
1.794	-0.0436979\\
1.796	-0.0436979\\
1.798	-0.0436979\\
1.8	-0.0436979\\
1.802	-0.0436979\\
1.804	-0.0436979\\
1.806	-0.0436979\\
1.808	-0.0436979\\
1.81	-0.0436979\\
1.812	-0.0436979\\
1.814	-0.0436979\\
1.816	-0.0436979\\
1.818	-0.0436979\\
1.82	-0.0436979\\
1.822	-0.0436979\\
1.824	-0.0436979\\
1.826	-0.0436979\\
1.828	-0.0436979\\
1.83	-0.0436979\\
1.832	-0.0436979\\
1.834	-0.0436979\\
1.836	-0.0436979\\
1.838	-0.0436979\\
1.84	-0.0436979\\
1.842	-0.0436979\\
1.844	-0.0436979\\
1.846	-0.0436979\\
1.848	-0.0436979\\
1.85	-0.0436979\\
1.852	-0.0436979\\
1.854	-0.0436979\\
1.856	-0.0436979\\
1.858	-0.0436979\\
1.86	-0.0436979\\
1.862	-0.0436979\\
1.864	-0.0436979\\
1.866	-0.0436979\\
1.868	-0.0436979\\
1.87	-0.0436979\\
1.872	-0.0436979\\
1.874	-0.0436979\\
1.876	-0.0436979\\
1.878	-0.0436979\\
1.88	-0.0436979\\
1.882	-0.0436979\\
1.884	-0.0436979\\
1.886	-0.0436979\\
1.888	-0.0436979\\
1.89	-0.0436979\\
1.892	-0.0436979\\
1.894	-0.0436979\\
1.896	-0.0436979\\
1.898	-0.0436979\\
1.9	-0.0436979\\
1.902	-0.0436979\\
1.904	-0.0436979\\
1.906	-0.0436979\\
1.908	-0.0436979\\
1.91	-0.0436979\\
1.912	-0.0436979\\
1.914	-0.0436979\\
1.916	-0.0436979\\
1.918	-0.0436979\\
1.92	-0.0436979\\
1.922	-0.0436979\\
1.924	-0.0436979\\
1.926	-0.0436979\\
1.928	-0.0436979\\
1.93	-0.0436979\\
1.932	-0.0436979\\
1.934	-0.0436979\\
1.936	-0.0436979\\
1.938	-0.0436979\\
1.94	-0.0436979\\
1.942	-0.0436979\\
1.944	-0.0436979\\
1.946	-0.0436979\\
1.948	-0.0436979\\
1.95	-0.0436979\\
1.952	-0.0436979\\
1.954	-0.0436979\\
1.956	-0.0436979\\
1.958	-0.0436979\\
1.96	-0.0436979\\
1.962	-0.0436979\\
1.964	-0.0436979\\
1.966	-0.0436979\\
1.968	-0.0436979\\
1.97	-0.0436979\\
1.972	-0.0436979\\
1.974	-0.0436979\\
1.976	-0.0436979\\
1.978	-0.0436979\\
1.98	-0.0436979\\
1.982	-0.0436979\\
1.984	-0.0436979\\
1.986	-0.0436979\\
1.988	-0.0436979\\
1.99	-0.0436979\\
1.992	-0.0436979\\
1.994	-0.0436979\\
1.996	-0.0436979\\
1.998	-0.0436979\\
2	-0.0436979\\
};
\addplot [color=mycolor6, line width=1.0pt, forget plot]
  table[row sep=crcr]{%
-0.024	0\\
-0.022	0\\
-0.02	0\\
-0.018	0\\
-0.016	0\\
-0.014	0\\
-0.012	0\\
-0.01	0\\
-0.008	0\\
-0.006	0\\
-0.004	0\\
-0.002	0\\
0	0\\
0.002	-0.0019122\\
0.004	-0.00365961\\
0.006	-0.00525768\\
0.008	-0.00672078\\
0.01	-0.00806231\\
0.012	-0.00929472\\
0.014	-0.0104296\\
0.016	-0.0114775\\
0.018	-0.0124485\\
0.02	-0.0133516\\
0.022	-0.014195\\
0.024	-0.0149865\\
0.026	-0.015733\\
0.028	-0.0164406\\
0.03	-0.0171152\\
0.032	-0.0177615\\
0.034	-0.0183842\\
0.036	-0.018987\\
0.038	-0.0195735\\
0.04	-0.0201463\\
0.042	-0.0207081\\
0.044	-0.0212607\\
0.046	-0.0218059\\
0.048	-0.0223447\\
0.05	-0.0228783\\
0.052	-0.0234072\\
0.054	-0.0239319\\
0.056	-0.0244525\\
0.058	-0.024969\\
0.06	-0.0254814\\
0.062	-0.0259893\\
0.064	-0.0264924\\
0.066	-0.0269904\\
0.068	-0.0274828\\
0.07	-0.0279691\\
0.072	-0.028449\\
0.074	-0.0289221\\
0.076	-0.029388\\
0.078	-0.0298464\\
0.08	-0.030297\\
0.082	-0.0307397\\
0.084	-0.0311744\\
0.086	-0.0316009\\
0.088	-0.0320194\\
0.09	-0.0324298\\
0.092	-0.0328323\\
0.094	-0.033227\\
0.096	-0.0336141\\
0.098	-0.0339938\\
0.1	-0.0343664\\
0.102	-0.034732\\
0.104	-0.0350909\\
0.106	-0.0354433\\
0.108	-0.0357893\\
0.11	-0.0361291\\
0.112	-0.0364628\\
0.114	-0.0367905\\
0.116	-0.0371122\\
0.118	-0.0374279\\
0.12	-0.0377375\\
0.122	-0.038041\\
0.124	-0.038338\\
0.126	-0.0386285\\
0.128	-0.0389123\\
0.13	-0.0391889\\
0.132	-0.0394582\\
0.134	-0.0397197\\
0.136	-0.0399733\\
0.138	-0.0402185\\
0.14	-0.0404549\\
0.142	-0.0406824\\
0.144	-0.0409004\\
0.146	-0.0411088\\
0.148	-0.0413073\\
0.15	-0.0414956\\
0.152	-0.0416735\\
0.154	-0.0418409\\
0.156	-0.0419976\\
0.158	-0.0421436\\
0.16	-0.0422789\\
0.162	-0.0424035\\
0.164	-0.0425175\\
0.166	-0.042621\\
0.168	-0.0427143\\
0.17	-0.0427974\\
0.172	-0.0428709\\
0.174	-0.0429348\\
0.176	-0.0429896\\
0.178	-0.0430358\\
0.18	-0.0430736\\
0.182	-0.0431035\\
0.184	-0.0431261\\
0.186	-0.0431417\\
0.188	-0.0431509\\
0.19	-0.0431542\\
0.192	-0.043152\\
0.194	-0.0431448\\
0.196	-0.0431332\\
0.198	-0.0431176\\
0.2	-0.0430985\\
0.202	-0.0430763\\
0.204	-0.0430515\\
0.206	-0.0430244\\
0.208	-0.0429956\\
0.21	-0.0429653\\
0.212	-0.042934\\
0.214	-0.0429018\\
0.216	-0.0428693\\
0.218	-0.0428365\\
0.22	-0.0428038\\
0.222	-0.0427714\\
0.224	-0.0427396\\
0.226	-0.0427084\\
0.228	-0.0426781\\
0.23	-0.0426489\\
0.232	-0.0426207\\
0.234	-0.0425939\\
0.236	-0.0425684\\
0.238	-0.0425443\\
0.24	-0.0425218\\
0.242	-0.0425008\\
0.244	-0.0424814\\
0.246	-0.0424636\\
0.248	-0.0424475\\
0.25	-0.042433\\
0.252	-0.0424202\\
0.254	-0.0424091\\
0.256	-0.0423997\\
0.258	-0.0423919\\
0.26	-0.0423857\\
0.262	-0.0423811\\
0.264	-0.0423781\\
0.266	-0.0423767\\
0.268	-0.0423768\\
0.27	-0.0423783\\
0.272	-0.0423813\\
0.274	-0.0423857\\
0.276	-0.0423914\\
0.278	-0.0423984\\
0.28	-0.0424066\\
0.282	-0.0424161\\
0.284	-0.0424267\\
0.286	-0.0424383\\
0.288	-0.042451\\
0.29	-0.0424646\\
0.292	-0.0424791\\
0.294	-0.0424945\\
0.296	-0.0425107\\
0.298	-0.0425275\\
0.3	-0.042545\\
0.302	-0.0425631\\
0.304	-0.0425817\\
0.306	-0.0426008\\
0.308	-0.0426203\\
0.31	-0.0426401\\
0.312	-0.0426602\\
0.314	-0.0426805\\
0.316	-0.0427009\\
0.318	-0.0427215\\
0.32	-0.0427421\\
0.322	-0.0427628\\
0.324	-0.0427834\\
0.326	-0.0428039\\
0.328	-0.0428243\\
0.33	-0.0428445\\
0.332	-0.0428645\\
0.334	-0.0428843\\
0.336	-0.0429038\\
0.338	-0.042923\\
0.34	-0.0429418\\
0.342	-0.0429604\\
0.344	-0.0429785\\
0.346	-0.0429963\\
0.348	-0.0430137\\
0.35	-0.0430307\\
0.352	-0.0430472\\
0.354	-0.0430634\\
0.356	-0.0430791\\
0.358	-0.0430944\\
0.36	-0.0431093\\
0.362	-0.0431237\\
0.364	-0.0431377\\
0.366	-0.0431514\\
0.368	-0.0431646\\
0.37	-0.0431774\\
0.372	-0.0431898\\
0.374	-0.0432019\\
0.376	-0.0432136\\
0.378	-0.0432249\\
0.38	-0.0432359\\
0.382	-0.0432465\\
0.384	-0.0432569\\
0.386	-0.043267\\
0.388	-0.0432767\\
0.39	-0.0432862\\
0.392	-0.0432954\\
0.394	-0.0433044\\
0.396	-0.0433132\\
0.398	-0.0433217\\
0.4	-0.04333\\
0.402	-0.0433381\\
0.404	-0.043346\\
0.406	-0.0433537\\
0.408	-0.0433613\\
0.41	-0.0433687\\
0.412	-0.0433759\\
0.414	-0.043383\\
0.416	-0.04339\\
0.418	-0.0433968\\
0.42	-0.0434035\\
0.422	-0.0434101\\
0.424	-0.0434166\\
0.426	-0.043423\\
0.428	-0.0434292\\
0.43	-0.0434354\\
0.432	-0.0434415\\
0.434	-0.0434474\\
0.436	-0.0434533\\
0.438	-0.0434591\\
0.44	-0.0434648\\
0.442	-0.0434704\\
0.444	-0.043476\\
0.446	-0.0434814\\
0.448	-0.0434868\\
0.45	-0.043492\\
0.452	-0.0434972\\
0.454	-0.0435023\\
0.456	-0.0435074\\
0.458	-0.0435123\\
0.46	-0.0435172\\
0.462	-0.043522\\
0.464	-0.0435267\\
0.466	-0.0435313\\
0.468	-0.0435359\\
0.47	-0.0435403\\
0.472	-0.0435447\\
0.474	-0.043549\\
0.476	-0.0435532\\
0.478	-0.0435573\\
0.48	-0.0435614\\
0.482	-0.0435654\\
0.484	-0.0435692\\
0.486	-0.043573\\
0.488	-0.0435768\\
0.49	-0.0435804\\
0.492	-0.0435839\\
0.494	-0.0435874\\
0.496	-0.0435908\\
0.498	-0.0435941\\
0.5	-0.0435973\\
0.502	-0.0436004\\
0.504	-0.0436035\\
0.506	-0.0436064\\
0.508	-0.0436093\\
0.51	-0.0436121\\
0.512	-0.0436148\\
0.514	-0.0436175\\
0.516	-0.0436201\\
0.518	-0.0436225\\
0.52	-0.043625\\
0.522	-0.0436273\\
0.524	-0.0436295\\
0.526	-0.0436317\\
0.528	-0.0436338\\
0.53	-0.0436359\\
0.532	-0.0436379\\
0.534	-0.0436398\\
0.536	-0.0436416\\
0.538	-0.0436434\\
0.54	-0.0436451\\
0.542	-0.0436467\\
0.544	-0.0436483\\
0.546	-0.0436498\\
0.548	-0.0436512\\
0.55	-0.0436526\\
0.552	-0.043654\\
0.554	-0.0436552\\
0.556	-0.0436565\\
0.558	-0.0436576\\
0.56	-0.0436588\\
0.562	-0.0436598\\
0.564	-0.0436609\\
0.566	-0.0436619\\
0.568	-0.0436628\\
0.57	-0.0436637\\
0.572	-0.0436645\\
0.574	-0.0436654\\
0.576	-0.0436661\\
0.578	-0.0436669\\
0.58	-0.0436676\\
0.582	-0.0436682\\
0.584	-0.0436689\\
0.586	-0.0436695\\
0.588	-0.0436701\\
0.59	-0.0436706\\
0.592	-0.0436711\\
0.594	-0.0436716\\
0.596	-0.0436721\\
0.598	-0.0436726\\
0.6	-0.043673\\
0.602	-0.0436734\\
0.604	-0.0436738\\
0.606	-0.0436742\\
0.608	-0.0436745\\
0.61	-0.0436749\\
0.612	-0.0436752\\
0.614	-0.0436755\\
0.616	-0.0436758\\
0.618	-0.0436761\\
0.62	-0.0436763\\
0.622	-0.0436766\\
0.624	-0.0436768\\
0.626	-0.0436771\\
0.628	-0.0436773\\
0.63	-0.0436775\\
0.632	-0.0436778\\
0.634	-0.043678\\
0.636	-0.0436782\\
0.638	-0.0436784\\
0.64	-0.0436786\\
0.642	-0.0436787\\
0.644	-0.0436789\\
0.646	-0.0436791\\
0.648	-0.0436793\\
0.65	-0.0436794\\
0.652	-0.0436796\\
0.654	-0.0436798\\
0.656	-0.0436799\\
0.658	-0.0436801\\
0.66	-0.0436802\\
0.662	-0.0436804\\
0.664	-0.0436805\\
0.666	-0.0436807\\
0.668	-0.0436808\\
0.67	-0.0436809\\
0.672	-0.0436811\\
0.674	-0.0436812\\
0.676	-0.0436814\\
0.678	-0.0436815\\
0.68	-0.0436816\\
0.682	-0.0436818\\
0.684	-0.0436819\\
0.686	-0.043682\\
0.688	-0.0436821\\
0.69	-0.0436823\\
0.692	-0.0436824\\
0.694	-0.0436825\\
0.696	-0.0436826\\
0.698	-0.0436827\\
0.7	-0.0436829\\
0.702	-0.043683\\
0.704	-0.0436831\\
0.706	-0.0436832\\
0.708	-0.0436833\\
0.71	-0.0436834\\
0.712	-0.0436835\\
0.714	-0.0436836\\
0.716	-0.0436837\\
0.718	-0.0436838\\
0.72	-0.0436839\\
0.722	-0.043684\\
0.724	-0.0436841\\
0.726	-0.0436842\\
0.728	-0.0436843\\
0.73	-0.0436844\\
0.732	-0.0436845\\
0.734	-0.0436846\\
0.736	-0.0436847\\
0.738	-0.0436848\\
0.74	-0.0436849\\
0.742	-0.043685\\
0.744	-0.0436851\\
0.746	-0.0436852\\
0.748	-0.0436853\\
0.75	-0.0436854\\
0.752	-0.0436855\\
0.754	-0.0436856\\
0.756	-0.0436856\\
0.758	-0.0436857\\
0.76	-0.0436858\\
0.762	-0.043686\\
0.764	-0.0436861\\
0.766	-0.0436862\\
0.768	-0.0436863\\
0.77	-0.0436864\\
0.772	-0.0436865\\
0.774	-0.0436866\\
0.776	-0.0436867\\
0.778	-0.0436869\\
0.78	-0.043687\\
0.782	-0.0436871\\
0.784	-0.0436872\\
0.786	-0.0436874\\
0.788	-0.0436875\\
0.79	-0.0436877\\
0.792	-0.0436878\\
0.794	-0.043688\\
0.796	-0.0436881\\
0.798	-0.0436883\\
0.8	-0.0436884\\
0.802	-0.0436886\\
0.804	-0.0436888\\
0.806	-0.043689\\
0.808	-0.0436891\\
0.81	-0.0436893\\
0.812	-0.0436895\\
0.814	-0.0436897\\
0.816	-0.0436899\\
0.818	-0.0436901\\
0.82	-0.0436903\\
0.822	-0.0436905\\
0.824	-0.0436907\\
0.826	-0.0436909\\
0.828	-0.0436911\\
0.83	-0.0436913\\
0.832	-0.0436915\\
0.834	-0.0436917\\
0.836	-0.043692\\
0.838	-0.0436922\\
0.84	-0.0436924\\
0.842	-0.0436926\\
0.844	-0.0436929\\
0.846	-0.0436931\\
0.848	-0.0436933\\
0.85	-0.0436935\\
0.852	-0.0436938\\
0.854	-0.043694\\
0.856	-0.0436942\\
0.858	-0.0436944\\
0.86	-0.0436946\\
0.862	-0.0436949\\
0.864	-0.0436951\\
0.866	-0.0436953\\
0.868	-0.0436955\\
0.87	-0.0436957\\
0.872	-0.0436959\\
0.874	-0.0436961\\
0.876	-0.0436963\\
0.878	-0.0436965\\
0.88	-0.0436967\\
0.882	-0.0436969\\
0.884	-0.043697\\
0.886	-0.0436972\\
0.888	-0.0436974\\
0.89	-0.0436976\\
0.892	-0.0436977\\
0.894	-0.0436979\\
0.896	-0.043698\\
0.898	-0.0436982\\
0.9	-0.0436983\\
0.902	-0.0436984\\
0.904	-0.0436985\\
0.906	-0.0436987\\
0.908	-0.0436988\\
0.91	-0.0436989\\
0.912	-0.043699\\
0.914	-0.0436991\\
0.916	-0.0436991\\
0.918	-0.0436992\\
0.92	-0.0436993\\
0.922	-0.0436994\\
0.924	-0.0436994\\
0.926	-0.0436995\\
0.928	-0.0436995\\
0.93	-0.0436996\\
0.932	-0.0436996\\
0.934	-0.0436996\\
0.936	-0.0436996\\
0.938	-0.0436997\\
0.94	-0.0436997\\
0.942	-0.0436997\\
0.944	-0.0436997\\
0.946	-0.0436997\\
0.948	-0.0436997\\
0.95	-0.0436996\\
0.952	-0.0436996\\
0.954	-0.0436996\\
0.956	-0.0436996\\
0.958	-0.0436995\\
0.96	-0.0436995\\
0.962	-0.0436994\\
0.964	-0.0436994\\
0.966	-0.0436994\\
0.968	-0.0436993\\
0.97	-0.0436992\\
0.972	-0.0436992\\
0.974	-0.0436991\\
0.976	-0.0436991\\
0.978	-0.043699\\
0.98	-0.0436989\\
0.982	-0.0436989\\
0.984	-0.0436988\\
0.986	-0.0436987\\
0.988	-0.0436987\\
0.99	-0.0436986\\
0.992	-0.0436985\\
0.994	-0.0436984\\
0.996	-0.0436984\\
0.998	-0.0436983\\
1	-0.0436982\\
1.002	-0.0436982\\
1.004	-0.0436981\\
1.006	-0.043698\\
1.008	-0.0436979\\
1.01	-0.0436979\\
1.012	-0.0436978\\
1.014	-0.0436977\\
1.016	-0.0436977\\
1.018	-0.0436976\\
1.02	-0.0436976\\
1.022	-0.0436975\\
1.024	-0.0436974\\
1.026	-0.0436974\\
1.028	-0.0436973\\
1.03	-0.0436973\\
1.032	-0.0436972\\
1.034	-0.0436972\\
1.036	-0.0436971\\
1.038	-0.0436971\\
1.04	-0.0436971\\
1.042	-0.043697\\
1.044	-0.043697\\
1.046	-0.043697\\
1.048	-0.0436969\\
1.05	-0.0436969\\
1.052	-0.0436969\\
1.054	-0.0436968\\
1.056	-0.0436968\\
1.058	-0.0436968\\
1.06	-0.0436968\\
1.062	-0.0436968\\
1.064	-0.0436968\\
1.066	-0.0436968\\
1.068	-0.0436968\\
1.07	-0.0436967\\
1.072	-0.0436967\\
1.074	-0.0436967\\
1.076	-0.0436967\\
1.078	-0.0436968\\
1.08	-0.0436968\\
1.082	-0.0436968\\
1.084	-0.0436968\\
1.086	-0.0436968\\
1.088	-0.0436968\\
1.09	-0.0436968\\
1.092	-0.0436968\\
1.094	-0.0436968\\
1.096	-0.0436969\\
1.098	-0.0436969\\
1.1	-0.0436969\\
1.102	-0.0436969\\
1.104	-0.0436969\\
1.106	-0.043697\\
1.108	-0.043697\\
1.11	-0.043697\\
1.112	-0.043697\\
1.114	-0.0436971\\
1.116	-0.0436971\\
1.118	-0.0436971\\
1.12	-0.0436972\\
1.122	-0.0436972\\
1.124	-0.0436972\\
1.126	-0.0436972\\
1.128	-0.0436973\\
1.13	-0.0436973\\
1.132	-0.0436973\\
1.134	-0.0436974\\
1.136	-0.0436974\\
1.138	-0.0436974\\
1.14	-0.0436974\\
1.142	-0.0436975\\
1.144	-0.0436975\\
1.146	-0.0436975\\
1.148	-0.0436976\\
1.15	-0.0436976\\
1.152	-0.0436976\\
1.154	-0.0436976\\
1.156	-0.0436977\\
1.158	-0.0436977\\
1.16	-0.0436977\\
1.162	-0.0436978\\
1.164	-0.0436978\\
1.166	-0.0436978\\
1.168	-0.0436978\\
1.17	-0.0436978\\
1.172	-0.0436979\\
1.174	-0.0436979\\
1.176	-0.0436979\\
1.178	-0.0436979\\
1.18	-0.0436979\\
1.182	-0.043698\\
1.184	-0.043698\\
1.186	-0.043698\\
1.188	-0.043698\\
1.19	-0.043698\\
1.192	-0.043698\\
1.194	-0.0436981\\
1.196	-0.0436981\\
1.198	-0.0436981\\
1.2	-0.0436981\\
1.202	-0.0436981\\
1.204	-0.0436981\\
1.206	-0.0436981\\
1.208	-0.0436981\\
1.21	-0.0436981\\
1.212	-0.0436981\\
1.214	-0.0436982\\
1.216	-0.0436982\\
1.218	-0.0436982\\
1.22	-0.0436982\\
1.222	-0.0436982\\
1.224	-0.0436982\\
1.226	-0.0436982\\
1.228	-0.0436982\\
1.23	-0.0436982\\
1.232	-0.0436982\\
1.234	-0.0436982\\
1.236	-0.0436982\\
1.238	-0.0436982\\
1.24	-0.0436982\\
1.242	-0.0436982\\
1.244	-0.0436982\\
1.246	-0.0436982\\
1.248	-0.0436982\\
1.25	-0.0436982\\
1.252	-0.0436981\\
1.254	-0.0436981\\
1.256	-0.0436981\\
1.258	-0.0436981\\
1.26	-0.0436981\\
1.262	-0.0436981\\
1.264	-0.0436981\\
1.266	-0.0436981\\
1.268	-0.0436981\\
1.27	-0.0436981\\
1.272	-0.0436981\\
1.274	-0.0436981\\
1.276	-0.0436981\\
1.278	-0.0436981\\
1.28	-0.0436981\\
1.282	-0.043698\\
1.284	-0.043698\\
1.286	-0.043698\\
1.288	-0.043698\\
1.29	-0.043698\\
1.292	-0.043698\\
1.294	-0.043698\\
1.296	-0.043698\\
1.298	-0.043698\\
1.3	-0.043698\\
1.302	-0.043698\\
1.304	-0.043698\\
1.306	-0.043698\\
1.308	-0.043698\\
1.31	-0.0436979\\
1.312	-0.0436979\\
1.314	-0.0436979\\
1.316	-0.0436979\\
1.318	-0.0436979\\
1.32	-0.0436979\\
1.322	-0.0436979\\
1.324	-0.0436979\\
1.326	-0.0436979\\
1.328	-0.0436979\\
1.33	-0.0436979\\
1.332	-0.0436979\\
1.334	-0.0436979\\
1.336	-0.0436979\\
1.338	-0.0436979\\
1.34	-0.0436979\\
1.342	-0.0436979\\
1.344	-0.0436979\\
1.346	-0.0436978\\
1.348	-0.0436978\\
1.35	-0.0436978\\
1.352	-0.0436978\\
1.354	-0.0436978\\
1.356	-0.0436978\\
1.358	-0.0436978\\
1.36	-0.0436978\\
1.362	-0.0436978\\
1.364	-0.0436978\\
1.366	-0.0436978\\
1.368	-0.0436978\\
1.37	-0.0436978\\
1.372	-0.0436978\\
1.374	-0.0436978\\
1.376	-0.0436978\\
1.378	-0.0436978\\
1.38	-0.0436978\\
1.382	-0.0436978\\
1.384	-0.0436978\\
1.386	-0.0436978\\
1.388	-0.0436978\\
1.39	-0.0436978\\
1.392	-0.0436978\\
1.394	-0.0436978\\
1.396	-0.0436978\\
1.398	-0.0436978\\
1.4	-0.0436978\\
1.402	-0.0436978\\
1.404	-0.0436978\\
1.406	-0.0436978\\
1.408	-0.0436978\\
1.41	-0.0436978\\
1.412	-0.0436978\\
1.414	-0.0436978\\
1.416	-0.0436978\\
1.418	-0.0436978\\
1.42	-0.0436979\\
1.422	-0.0436979\\
1.424	-0.0436979\\
1.426	-0.0436979\\
1.428	-0.0436979\\
1.43	-0.0436979\\
1.432	-0.0436979\\
1.434	-0.0436979\\
1.436	-0.0436979\\
1.438	-0.0436979\\
1.44	-0.0436979\\
1.442	-0.0436979\\
1.444	-0.0436979\\
1.446	-0.0436979\\
1.448	-0.0436979\\
1.45	-0.0436979\\
1.452	-0.0436979\\
1.454	-0.0436979\\
1.456	-0.0436979\\
1.458	-0.0436979\\
1.46	-0.0436979\\
1.462	-0.0436979\\
1.464	-0.0436979\\
1.466	-0.0436979\\
1.468	-0.0436979\\
1.47	-0.0436979\\
1.472	-0.0436979\\
1.474	-0.0436979\\
1.476	-0.0436979\\
1.478	-0.0436979\\
1.48	-0.0436979\\
1.482	-0.0436979\\
1.484	-0.0436979\\
1.486	-0.0436979\\
1.488	-0.0436979\\
1.49	-0.0436979\\
1.492	-0.0436979\\
1.494	-0.0436979\\
1.496	-0.0436979\\
1.498	-0.0436979\\
1.5	-0.0436979\\
1.502	-0.0436979\\
1.504	-0.0436979\\
1.506	-0.0436979\\
1.508	-0.0436979\\
1.51	-0.0436979\\
1.512	-0.0436979\\
1.514	-0.0436979\\
1.516	-0.0436979\\
1.518	-0.0436979\\
1.52	-0.0436979\\
1.522	-0.0436979\\
1.524	-0.0436979\\
1.526	-0.0436979\\
1.528	-0.0436979\\
1.53	-0.0436979\\
1.532	-0.0436979\\
1.534	-0.0436979\\
1.536	-0.0436979\\
1.538	-0.0436979\\
1.54	-0.0436979\\
1.542	-0.0436979\\
1.544	-0.0436979\\
1.546	-0.0436979\\
1.548	-0.0436979\\
1.55	-0.0436979\\
1.552	-0.0436979\\
1.554	-0.0436979\\
1.556	-0.0436979\\
1.558	-0.0436979\\
1.56	-0.0436979\\
1.562	-0.0436979\\
1.564	-0.0436979\\
1.566	-0.0436979\\
1.568	-0.0436979\\
1.57	-0.0436979\\
1.572	-0.0436979\\
1.574	-0.0436979\\
1.576	-0.0436979\\
1.578	-0.0436979\\
1.58	-0.0436979\\
1.582	-0.0436979\\
1.584	-0.0436979\\
1.586	-0.0436979\\
1.588	-0.0436979\\
1.59	-0.0436979\\
1.592	-0.0436979\\
1.594	-0.0436979\\
1.596	-0.0436979\\
1.598	-0.0436979\\
1.6	-0.0436979\\
1.602	-0.0436979\\
1.604	-0.0436979\\
1.606	-0.0436979\\
1.608	-0.0436979\\
1.61	-0.0436979\\
1.612	-0.0436979\\
1.614	-0.0436979\\
1.616	-0.0436979\\
1.618	-0.0436979\\
1.62	-0.0436979\\
1.622	-0.0436979\\
1.624	-0.0436979\\
1.626	-0.0436979\\
1.628	-0.0436979\\
1.63	-0.0436979\\
1.632	-0.0436979\\
1.634	-0.0436979\\
1.636	-0.0436979\\
1.638	-0.0436979\\
1.64	-0.0436979\\
1.642	-0.0436979\\
1.644	-0.0436979\\
1.646	-0.0436979\\
1.648	-0.0436979\\
1.65	-0.0436979\\
1.652	-0.0436979\\
1.654	-0.0436979\\
1.656	-0.0436979\\
1.658	-0.0436979\\
1.66	-0.0436979\\
1.662	-0.0436979\\
1.664	-0.0436979\\
1.666	-0.0436979\\
1.668	-0.0436979\\
1.67	-0.0436979\\
1.672	-0.0436979\\
1.674	-0.0436979\\
1.676	-0.0436979\\
1.678	-0.0436979\\
1.68	-0.0436979\\
1.682	-0.0436979\\
1.684	-0.0436979\\
1.686	-0.0436979\\
1.688	-0.0436979\\
1.69	-0.0436979\\
1.692	-0.0436979\\
1.694	-0.0436979\\
1.696	-0.0436979\\
1.698	-0.0436979\\
1.7	-0.0436979\\
1.702	-0.0436979\\
1.704	-0.0436979\\
1.706	-0.0436979\\
1.708	-0.0436979\\
1.71	-0.0436979\\
1.712	-0.0436979\\
1.714	-0.0436979\\
1.716	-0.0436979\\
1.718	-0.0436979\\
1.72	-0.0436979\\
1.722	-0.0436979\\
1.724	-0.0436979\\
1.726	-0.0436979\\
1.728	-0.0436979\\
1.73	-0.0436979\\
1.732	-0.0436979\\
1.734	-0.0436979\\
1.736	-0.0436979\\
1.738	-0.0436979\\
1.74	-0.0436979\\
1.742	-0.0436979\\
1.744	-0.0436979\\
1.746	-0.0436979\\
1.748	-0.0436979\\
1.75	-0.0436979\\
1.752	-0.0436979\\
1.754	-0.0436979\\
1.756	-0.0436979\\
1.758	-0.0436979\\
1.76	-0.0436979\\
1.762	-0.0436979\\
1.764	-0.0436979\\
1.766	-0.0436979\\
1.768	-0.0436979\\
1.77	-0.0436979\\
1.772	-0.0436979\\
1.774	-0.0436979\\
1.776	-0.0436979\\
1.778	-0.0436979\\
1.78	-0.0436979\\
1.782	-0.0436979\\
1.784	-0.0436979\\
1.786	-0.0436979\\
1.788	-0.0436979\\
1.79	-0.0436979\\
1.792	-0.0436979\\
1.794	-0.0436979\\
1.796	-0.0436979\\
1.798	-0.0436979\\
1.8	-0.0436979\\
1.802	-0.0436979\\
1.804	-0.0436979\\
1.806	-0.0436979\\
1.808	-0.0436979\\
1.81	-0.0436979\\
1.812	-0.0436979\\
1.814	-0.0436979\\
1.816	-0.0436979\\
1.818	-0.0436979\\
1.82	-0.0436979\\
1.822	-0.0436979\\
1.824	-0.0436979\\
1.826	-0.0436979\\
1.828	-0.0436979\\
1.83	-0.0436979\\
1.832	-0.0436979\\
1.834	-0.0436979\\
1.836	-0.0436979\\
1.838	-0.0436979\\
1.84	-0.0436979\\
1.842	-0.0436979\\
1.844	-0.0436979\\
1.846	-0.0436979\\
1.848	-0.0436979\\
1.85	-0.0436979\\
1.852	-0.0436979\\
1.854	-0.0436979\\
1.856	-0.0436979\\
1.858	-0.0436979\\
1.86	-0.0436979\\
1.862	-0.0436979\\
1.864	-0.0436979\\
1.866	-0.0436979\\
1.868	-0.0436979\\
1.87	-0.0436979\\
1.872	-0.0436979\\
1.874	-0.0436979\\
1.876	-0.0436979\\
1.878	-0.0436979\\
1.88	-0.0436979\\
1.882	-0.0436979\\
1.884	-0.0436979\\
1.886	-0.0436979\\
1.888	-0.0436979\\
1.89	-0.0436979\\
1.892	-0.0436979\\
1.894	-0.0436979\\
1.896	-0.0436979\\
1.898	-0.0436979\\
1.9	-0.0436979\\
1.902	-0.0436979\\
1.904	-0.0436979\\
1.906	-0.0436979\\
1.908	-0.0436979\\
1.91	-0.0436979\\
1.912	-0.0436979\\
1.914	-0.0436979\\
1.916	-0.0436979\\
1.918	-0.0436979\\
1.92	-0.0436979\\
1.922	-0.0436979\\
1.924	-0.0436979\\
1.926	-0.0436979\\
1.928	-0.0436979\\
1.93	-0.0436979\\
1.932	-0.0436979\\
1.934	-0.0436979\\
1.936	-0.0436979\\
1.938	-0.0436979\\
1.94	-0.0436979\\
1.942	-0.0436979\\
1.944	-0.0436979\\
1.946	-0.0436979\\
1.948	-0.0436979\\
1.95	-0.0436979\\
1.952	-0.0436979\\
1.954	-0.0436979\\
1.956	-0.0436979\\
1.958	-0.0436979\\
1.96	-0.0436979\\
1.962	-0.0436979\\
1.964	-0.0436979\\
1.966	-0.0436979\\
1.968	-0.0436979\\
1.97	-0.0436979\\
1.972	-0.0436979\\
1.974	-0.0436979\\
1.976	-0.0436979\\
1.978	-0.0436979\\
1.98	-0.0436979\\
1.982	-0.0436979\\
1.984	-0.0436979\\
1.986	-0.0436979\\
1.988	-0.0436979\\
1.99	-0.0436979\\
1.992	-0.0436979\\
1.994	-0.0436979\\
1.996	-0.0436979\\
1.998	-0.0436979\\
2	-0.0436979\\
};
\addplot [color=mycolor7, line width=1.0pt, forget plot]
  table[row sep=crcr]{%
-0.024	0\\
-0.022	0\\
-0.02	0\\
-0.018	0\\
-0.016	0\\
-0.014	0\\
-0.012	0\\
-0.01	0\\
-0.008	0\\
-0.006	0\\
-0.004	0\\
-0.002	0\\
0	0\\
0.002	-0.0019122\\
0.004	-0.00365961\\
0.006	-0.00525768\\
0.008	-0.00672077\\
0.01	-0.00806229\\
0.012	-0.00929468\\
0.014	-0.0104295\\
0.016	-0.0114774\\
0.018	-0.0124484\\
0.02	-0.0133513\\
0.022	-0.0141947\\
0.024	-0.0149861\\
0.026	-0.0157325\\
0.028	-0.01644\\
0.03	-0.0171143\\
0.032	-0.0177605\\
0.034	-0.0183829\\
0.036	-0.0189855\\
0.038	-0.0195717\\
0.04	-0.0201442\\
0.042	-0.0207055\\
0.044	-0.0212577\\
0.046	-0.0218024\\
0.048	-0.0223407\\
0.05	-0.0228737\\
0.052	-0.0234019\\
0.054	-0.0239258\\
0.056	-0.0244457\\
0.058	-0.0249614\\
0.06	-0.0254728\\
0.062	-0.0259798\\
0.064	-0.0264819\\
0.066	-0.0269787\\
0.068	-0.02747\\
0.07	-0.0279551\\
0.072	-0.0284338\\
0.074	-0.0289055\\
0.076	-0.0293701\\
0.078	-0.029827\\
0.08	-0.0302762\\
0.082	-0.0307174\\
0.084	-0.0311506\\
0.086	-0.0315756\\
0.088	-0.0319925\\
0.09	-0.0324013\\
0.092	-0.0328022\\
0.094	-0.0331953\\
0.096	-0.0335808\\
0.098	-0.033959\\
0.1	-0.03433\\
0.102	-0.034694\\
0.104	-0.0350514\\
0.106	-0.0354023\\
0.108	-0.0357469\\
0.11	-0.0360853\\
0.112	-0.0364177\\
0.114	-0.0367441\\
0.116	-0.0370647\\
0.118	-0.0373793\\
0.12	-0.0376879\\
0.122	-0.0379904\\
0.124	-0.0382866\\
0.126	-0.0385764\\
0.128	-0.0388596\\
0.13	-0.0391357\\
0.132	-0.0394046\\
0.134	-0.039666\\
0.136	-0.0399195\\
0.138	-0.0401647\\
0.14	-0.0404013\\
0.142	-0.0406291\\
0.144	-0.0408476\\
0.146	-0.0410566\\
0.148	-0.0412558\\
0.15	-0.041445\\
0.152	-0.0416239\\
0.154	-0.0417925\\
0.156	-0.0419506\\
0.158	-0.0420981\\
0.16	-0.042235\\
0.162	-0.0423614\\
0.164	-0.0424772\\
0.166	-0.0425828\\
0.168	-0.0426782\\
0.17	-0.0427636\\
0.172	-0.0428394\\
0.174	-0.0429059\\
0.176	-0.0429633\\
0.178	-0.0430122\\
0.18	-0.0430528\\
0.182	-0.0430856\\
0.184	-0.0431111\\
0.186	-0.0431297\\
0.188	-0.043142\\
0.19	-0.0431483\\
0.192	-0.0431492\\
0.194	-0.0431453\\
0.196	-0.0431368\\
0.198	-0.0431244\\
0.2	-0.0431085\\
0.202	-0.0430895\\
0.204	-0.0430678\\
0.206	-0.043044\\
0.208	-0.0430182\\
0.21	-0.042991\\
0.212	-0.0429626\\
0.214	-0.0429334\\
0.216	-0.0429036\\
0.218	-0.0428736\\
0.22	-0.0428436\\
0.222	-0.0428137\\
0.224	-0.0427843\\
0.226	-0.0427555\\
0.228	-0.0427274\\
0.23	-0.0427002\\
0.232	-0.0426741\\
0.234	-0.042649\\
0.236	-0.0426252\\
0.238	-0.0426026\\
0.24	-0.0425815\\
0.242	-0.0425617\\
0.244	-0.0425433\\
0.246	-0.0425265\\
0.248	-0.0425111\\
0.25	-0.0424973\\
0.252	-0.0424849\\
0.254	-0.0424741\\
0.256	-0.0424648\\
0.258	-0.042457\\
0.26	-0.0424507\\
0.262	-0.0424458\\
0.264	-0.0424424\\
0.266	-0.0424404\\
0.268	-0.0424398\\
0.27	-0.0424405\\
0.272	-0.0424425\\
0.274	-0.0424458\\
0.276	-0.0424503\\
0.278	-0.0424561\\
0.28	-0.0424629\\
0.282	-0.0424709\\
0.284	-0.0424799\\
0.286	-0.0424899\\
0.288	-0.0425009\\
0.29	-0.0425127\\
0.292	-0.0425254\\
0.294	-0.0425389\\
0.296	-0.0425531\\
0.298	-0.042568\\
0.3	-0.0425835\\
0.302	-0.0425995\\
0.304	-0.0426161\\
0.306	-0.0426331\\
0.308	-0.0426505\\
0.31	-0.0426683\\
0.312	-0.0426863\\
0.314	-0.0427046\\
0.316	-0.0427231\\
0.318	-0.0427417\\
0.32	-0.0427604\\
0.322	-0.0427791\\
0.324	-0.0427978\\
0.326	-0.0428165\\
0.328	-0.0428351\\
0.33	-0.0428536\\
0.332	-0.0428719\\
0.334	-0.04289\\
0.336	-0.042908\\
0.338	-0.0429257\\
0.34	-0.0429431\\
0.342	-0.0429603\\
0.344	-0.0429771\\
0.346	-0.0429937\\
0.348	-0.0430099\\
0.35	-0.0430258\\
0.352	-0.0430414\\
0.354	-0.0430566\\
0.356	-0.0430714\\
0.358	-0.0430859\\
0.36	-0.0431001\\
0.362	-0.0431139\\
0.364	-0.0431273\\
0.366	-0.0431404\\
0.368	-0.0431532\\
0.37	-0.0431656\\
0.372	-0.0431777\\
0.374	-0.0431895\\
0.376	-0.043201\\
0.378	-0.0432122\\
0.38	-0.0432231\\
0.382	-0.0432337\\
0.384	-0.0432441\\
0.386	-0.0432542\\
0.388	-0.043264\\
0.39	-0.0432736\\
0.392	-0.0432831\\
0.394	-0.0432922\\
0.396	-0.0433012\\
0.398	-0.04331\\
0.4	-0.0433186\\
0.402	-0.0433271\\
0.404	-0.0433353\\
0.406	-0.0433434\\
0.408	-0.0433514\\
0.41	-0.0433592\\
0.412	-0.0433668\\
0.414	-0.0433743\\
0.416	-0.0433817\\
0.418	-0.043389\\
0.42	-0.0433961\\
0.422	-0.0434032\\
0.424	-0.0434101\\
0.426	-0.0434169\\
0.428	-0.0434236\\
0.43	-0.0434301\\
0.432	-0.0434366\\
0.434	-0.043443\\
0.436	-0.0434492\\
0.438	-0.0434554\\
0.44	-0.0434615\\
0.442	-0.0434674\\
0.444	-0.0434733\\
0.446	-0.0434791\\
0.448	-0.0434847\\
0.45	-0.0434903\\
0.452	-0.0434957\\
0.454	-0.0435011\\
0.456	-0.0435064\\
0.458	-0.0435115\\
0.46	-0.0435166\\
0.462	-0.0435216\\
0.464	-0.0435264\\
0.466	-0.0435312\\
0.468	-0.0435359\\
0.47	-0.0435405\\
0.472	-0.0435449\\
0.474	-0.0435493\\
0.476	-0.0435536\\
0.478	-0.0435578\\
0.48	-0.0435618\\
0.482	-0.0435658\\
0.484	-0.0435697\\
0.486	-0.0435735\\
0.488	-0.0435772\\
0.49	-0.0435808\\
0.492	-0.0435843\\
0.494	-0.0435877\\
0.496	-0.043591\\
0.498	-0.0435943\\
0.5	-0.0435974\\
0.502	-0.0436005\\
0.504	-0.0436034\\
0.506	-0.0436063\\
0.508	-0.0436091\\
0.51	-0.0436118\\
0.512	-0.0436144\\
0.514	-0.043617\\
0.516	-0.0436195\\
0.518	-0.0436218\\
0.52	-0.0436242\\
0.522	-0.0436264\\
0.524	-0.0436286\\
0.526	-0.0436306\\
0.528	-0.0436327\\
0.53	-0.0436346\\
0.532	-0.0436365\\
0.534	-0.0436383\\
0.536	-0.0436401\\
0.538	-0.0436417\\
0.54	-0.0436434\\
0.542	-0.0436449\\
0.544	-0.0436465\\
0.546	-0.0436479\\
0.548	-0.0436493\\
0.55	-0.0436507\\
0.552	-0.0436519\\
0.554	-0.0436532\\
0.556	-0.0436544\\
0.558	-0.0436555\\
0.56	-0.0436567\\
0.562	-0.0436577\\
0.564	-0.0436587\\
0.566	-0.0436597\\
0.568	-0.0436607\\
0.57	-0.0436616\\
0.572	-0.0436624\\
0.574	-0.0436633\\
0.576	-0.0436641\\
0.578	-0.0436648\\
0.58	-0.0436656\\
0.582	-0.0436663\\
0.584	-0.043667\\
0.586	-0.0436676\\
0.588	-0.0436683\\
0.59	-0.0436689\\
0.592	-0.0436694\\
0.594	-0.04367\\
0.596	-0.0436705\\
0.598	-0.043671\\
0.6	-0.0436715\\
0.602	-0.043672\\
0.604	-0.0436725\\
0.606	-0.0436729\\
0.608	-0.0436733\\
0.61	-0.0436737\\
0.612	-0.0436741\\
0.614	-0.0436745\\
0.616	-0.0436749\\
0.618	-0.0436752\\
0.62	-0.0436755\\
0.622	-0.0436759\\
0.624	-0.0436762\\
0.626	-0.0436765\\
0.628	-0.0436768\\
0.63	-0.0436771\\
0.632	-0.0436773\\
0.634	-0.0436776\\
0.636	-0.0436778\\
0.638	-0.0436781\\
0.64	-0.0436783\\
0.642	-0.0436786\\
0.644	-0.0436788\\
0.646	-0.043679\\
0.648	-0.0436792\\
0.65	-0.0436794\\
0.652	-0.0436796\\
0.654	-0.0436798\\
0.656	-0.04368\\
0.658	-0.0436802\\
0.66	-0.0436803\\
0.662	-0.0436805\\
0.664	-0.0436807\\
0.666	-0.0436808\\
0.668	-0.043681\\
0.67	-0.0436811\\
0.672	-0.0436813\\
0.674	-0.0436814\\
0.676	-0.0436816\\
0.678	-0.0436817\\
0.68	-0.0436818\\
0.682	-0.043682\\
0.684	-0.0436821\\
0.686	-0.0436822\\
0.688	-0.0436823\\
0.69	-0.0436825\\
0.692	-0.0436826\\
0.694	-0.0436827\\
0.696	-0.0436828\\
0.698	-0.0436829\\
0.7	-0.043683\\
0.702	-0.0436831\\
0.704	-0.0436832\\
0.706	-0.0436833\\
0.708	-0.0436834\\
0.71	-0.0436835\\
0.712	-0.0436836\\
0.714	-0.0436837\\
0.716	-0.0436838\\
0.718	-0.0436839\\
0.72	-0.043684\\
0.722	-0.0436841\\
0.724	-0.0436842\\
0.726	-0.0436843\\
0.728	-0.0436843\\
0.73	-0.0436844\\
0.732	-0.0436845\\
0.734	-0.0436846\\
0.736	-0.0436847\\
0.738	-0.0436848\\
0.74	-0.0436849\\
0.742	-0.043685\\
0.744	-0.0436851\\
0.746	-0.0436851\\
0.748	-0.0436852\\
0.75	-0.0436853\\
0.752	-0.0436854\\
0.754	-0.0436855\\
0.756	-0.0436856\\
0.758	-0.0436857\\
0.76	-0.0436858\\
0.762	-0.0436859\\
0.764	-0.043686\\
0.766	-0.0436862\\
0.768	-0.0436863\\
0.77	-0.0436864\\
0.772	-0.0436865\\
0.774	-0.0436866\\
0.776	-0.0436868\\
0.778	-0.0436869\\
0.78	-0.043687\\
0.782	-0.0436872\\
0.784	-0.0436873\\
0.786	-0.0436874\\
0.788	-0.0436876\\
0.79	-0.0436877\\
0.792	-0.0436879\\
0.794	-0.0436881\\
0.796	-0.0436882\\
0.798	-0.0436884\\
0.8	-0.0436886\\
0.802	-0.0436887\\
0.804	-0.0436889\\
0.806	-0.0436891\\
0.808	-0.0436893\\
0.81	-0.0436895\\
0.812	-0.0436897\\
0.814	-0.0436898\\
0.816	-0.04369\\
0.818	-0.0436902\\
0.82	-0.0436904\\
0.822	-0.0436907\\
0.824	-0.0436909\\
0.826	-0.0436911\\
0.828	-0.0436913\\
0.83	-0.0436915\\
0.832	-0.0436917\\
0.834	-0.0436919\\
0.836	-0.0436921\\
0.838	-0.0436924\\
0.84	-0.0436926\\
0.842	-0.0436928\\
0.844	-0.043693\\
0.846	-0.0436932\\
0.848	-0.0436935\\
0.85	-0.0436937\\
0.852	-0.0436939\\
0.854	-0.0436941\\
0.856	-0.0436943\\
0.858	-0.0436945\\
0.86	-0.0436948\\
0.862	-0.043695\\
0.864	-0.0436952\\
0.866	-0.0436954\\
0.868	-0.0436956\\
0.87	-0.0436958\\
0.872	-0.043696\\
0.874	-0.0436962\\
0.876	-0.0436964\\
0.878	-0.0436965\\
0.88	-0.0436967\\
0.882	-0.0436969\\
0.884	-0.0436971\\
0.886	-0.0436972\\
0.888	-0.0436974\\
0.89	-0.0436975\\
0.892	-0.0436977\\
0.894	-0.0436978\\
0.896	-0.043698\\
0.898	-0.0436981\\
0.9	-0.0436982\\
0.902	-0.0436984\\
0.904	-0.0436985\\
0.906	-0.0436986\\
0.908	-0.0436987\\
0.91	-0.0436988\\
0.912	-0.0436989\\
0.914	-0.043699\\
0.916	-0.043699\\
0.918	-0.0436991\\
0.92	-0.0436992\\
0.922	-0.0436992\\
0.924	-0.0436993\\
0.926	-0.0436993\\
0.928	-0.0436994\\
0.93	-0.0436994\\
0.932	-0.0436995\\
0.934	-0.0436995\\
0.936	-0.0436995\\
0.938	-0.0436995\\
0.94	-0.0436995\\
0.942	-0.0436995\\
0.944	-0.0436995\\
0.946	-0.0436995\\
0.948	-0.0436995\\
0.95	-0.0436995\\
0.952	-0.0436995\\
0.954	-0.0436994\\
0.956	-0.0436994\\
0.958	-0.0436994\\
0.96	-0.0436994\\
0.962	-0.0436993\\
0.964	-0.0436993\\
0.966	-0.0436992\\
0.968	-0.0436992\\
0.97	-0.0436991\\
0.972	-0.0436991\\
0.974	-0.043699\\
0.976	-0.043699\\
0.978	-0.0436989\\
0.98	-0.0436988\\
0.982	-0.0436988\\
0.984	-0.0436987\\
0.986	-0.0436987\\
0.988	-0.0436986\\
0.99	-0.0436985\\
0.992	-0.0436985\\
0.994	-0.0436984\\
0.996	-0.0436983\\
0.998	-0.0436983\\
1	-0.0436982\\
1.002	-0.0436981\\
1.004	-0.0436981\\
1.006	-0.043698\\
1.008	-0.0436979\\
1.01	-0.0436979\\
1.012	-0.0436978\\
1.014	-0.0436977\\
1.016	-0.0436977\\
1.018	-0.0436976\\
1.02	-0.0436976\\
1.022	-0.0436975\\
1.024	-0.0436975\\
1.026	-0.0436974\\
1.028	-0.0436974\\
1.03	-0.0436973\\
1.032	-0.0436973\\
1.034	-0.0436972\\
1.036	-0.0436972\\
1.038	-0.0436971\\
1.04	-0.0436971\\
1.042	-0.0436971\\
1.044	-0.043697\\
1.046	-0.043697\\
1.048	-0.043697\\
1.05	-0.0436969\\
1.052	-0.0436969\\
1.054	-0.0436969\\
1.056	-0.0436969\\
1.058	-0.0436969\\
1.06	-0.0436968\\
1.062	-0.0436968\\
1.064	-0.0436968\\
1.066	-0.0436968\\
1.068	-0.0436968\\
1.07	-0.0436968\\
1.072	-0.0436968\\
1.074	-0.0436968\\
1.076	-0.0436968\\
1.078	-0.0436968\\
1.08	-0.0436968\\
1.082	-0.0436968\\
1.084	-0.0436968\\
1.086	-0.0436968\\
1.088	-0.0436968\\
1.09	-0.0436968\\
1.092	-0.0436969\\
1.094	-0.0436969\\
1.096	-0.0436969\\
1.098	-0.0436969\\
1.1	-0.0436969\\
1.102	-0.043697\\
1.104	-0.043697\\
1.106	-0.043697\\
1.108	-0.043697\\
1.11	-0.043697\\
1.112	-0.0436971\\
1.114	-0.0436971\\
1.116	-0.0436971\\
1.118	-0.0436971\\
1.12	-0.0436972\\
1.122	-0.0436972\\
1.124	-0.0436972\\
1.126	-0.0436972\\
1.128	-0.0436973\\
1.13	-0.0436973\\
1.132	-0.0436973\\
1.134	-0.0436974\\
1.136	-0.0436974\\
1.138	-0.0436974\\
1.14	-0.0436974\\
1.142	-0.0436975\\
1.144	-0.0436975\\
1.146	-0.0436975\\
1.148	-0.0436976\\
1.15	-0.0436976\\
1.152	-0.0436976\\
1.154	-0.0436976\\
1.156	-0.0436977\\
1.158	-0.0436977\\
1.16	-0.0436977\\
1.162	-0.0436977\\
1.164	-0.0436978\\
1.166	-0.0436978\\
1.168	-0.0436978\\
1.17	-0.0436978\\
1.172	-0.0436978\\
1.174	-0.0436979\\
1.176	-0.0436979\\
1.178	-0.0436979\\
1.18	-0.0436979\\
1.182	-0.0436979\\
1.184	-0.043698\\
1.186	-0.043698\\
1.188	-0.043698\\
1.19	-0.043698\\
1.192	-0.043698\\
1.194	-0.043698\\
1.196	-0.043698\\
1.198	-0.0436981\\
1.2	-0.0436981\\
1.202	-0.0436981\\
1.204	-0.0436981\\
1.206	-0.0436981\\
1.208	-0.0436981\\
1.21	-0.0436981\\
1.212	-0.0436981\\
1.214	-0.0436981\\
1.216	-0.0436981\\
1.218	-0.0436981\\
1.22	-0.0436981\\
1.222	-0.0436982\\
1.224	-0.0436982\\
1.226	-0.0436982\\
1.228	-0.0436982\\
1.23	-0.0436982\\
1.232	-0.0436982\\
1.234	-0.0436982\\
1.236	-0.0436982\\
1.238	-0.0436982\\
1.24	-0.0436982\\
1.242	-0.0436982\\
1.244	-0.0436982\\
1.246	-0.0436982\\
1.248	-0.0436982\\
1.25	-0.0436981\\
1.252	-0.0436981\\
1.254	-0.0436981\\
1.256	-0.0436981\\
1.258	-0.0436981\\
1.26	-0.0436981\\
1.262	-0.0436981\\
1.264	-0.0436981\\
1.266	-0.0436981\\
1.268	-0.0436981\\
1.27	-0.0436981\\
1.272	-0.0436981\\
1.274	-0.0436981\\
1.276	-0.0436981\\
1.278	-0.0436981\\
1.28	-0.0436981\\
1.282	-0.0436981\\
1.284	-0.043698\\
1.286	-0.043698\\
1.288	-0.043698\\
1.29	-0.043698\\
1.292	-0.043698\\
1.294	-0.043698\\
1.296	-0.043698\\
1.298	-0.043698\\
1.3	-0.043698\\
1.302	-0.043698\\
1.304	-0.043698\\
1.306	-0.043698\\
1.308	-0.043698\\
1.31	-0.0436979\\
1.312	-0.0436979\\
1.314	-0.0436979\\
1.316	-0.0436979\\
1.318	-0.0436979\\
1.32	-0.0436979\\
1.322	-0.0436979\\
1.324	-0.0436979\\
1.326	-0.0436979\\
1.328	-0.0436979\\
1.33	-0.0436979\\
1.332	-0.0436979\\
1.334	-0.0436979\\
1.336	-0.0436979\\
1.338	-0.0436979\\
1.34	-0.0436979\\
1.342	-0.0436979\\
1.344	-0.0436979\\
1.346	-0.0436979\\
1.348	-0.0436978\\
1.35	-0.0436978\\
1.352	-0.0436978\\
1.354	-0.0436978\\
1.356	-0.0436978\\
1.358	-0.0436978\\
1.36	-0.0436978\\
1.362	-0.0436978\\
1.364	-0.0436978\\
1.366	-0.0436978\\
1.368	-0.0436978\\
1.37	-0.0436978\\
1.372	-0.0436978\\
1.374	-0.0436978\\
1.376	-0.0436978\\
1.378	-0.0436978\\
1.38	-0.0436978\\
1.382	-0.0436978\\
1.384	-0.0436978\\
1.386	-0.0436978\\
1.388	-0.0436978\\
1.39	-0.0436978\\
1.392	-0.0436978\\
1.394	-0.0436978\\
1.396	-0.0436978\\
1.398	-0.0436978\\
1.4	-0.0436978\\
1.402	-0.0436978\\
1.404	-0.0436978\\
1.406	-0.0436978\\
1.408	-0.0436978\\
1.41	-0.0436978\\
1.412	-0.0436978\\
1.414	-0.0436978\\
1.416	-0.0436978\\
1.418	-0.0436978\\
1.42	-0.0436979\\
1.422	-0.0436979\\
1.424	-0.0436979\\
1.426	-0.0436979\\
1.428	-0.0436979\\
1.43	-0.0436979\\
1.432	-0.0436979\\
1.434	-0.0436979\\
1.436	-0.0436979\\
1.438	-0.0436979\\
1.44	-0.0436979\\
1.442	-0.0436979\\
1.444	-0.0436979\\
1.446	-0.0436979\\
1.448	-0.0436979\\
1.45	-0.0436979\\
1.452	-0.0436979\\
1.454	-0.0436979\\
1.456	-0.0436979\\
1.458	-0.0436979\\
1.46	-0.0436979\\
1.462	-0.0436979\\
1.464	-0.0436979\\
1.466	-0.0436979\\
1.468	-0.0436979\\
1.47	-0.0436979\\
1.472	-0.0436979\\
1.474	-0.0436979\\
1.476	-0.0436979\\
1.478	-0.0436979\\
1.48	-0.0436979\\
1.482	-0.0436979\\
1.484	-0.0436979\\
1.486	-0.0436979\\
1.488	-0.0436979\\
1.49	-0.0436979\\
1.492	-0.0436979\\
1.494	-0.0436979\\
1.496	-0.0436979\\
1.498	-0.0436979\\
1.5	-0.0436979\\
1.502	-0.0436979\\
1.504	-0.0436979\\
1.506	-0.0436979\\
1.508	-0.0436979\\
1.51	-0.0436979\\
1.512	-0.0436979\\
1.514	-0.0436979\\
1.516	-0.0436979\\
1.518	-0.0436979\\
1.52	-0.0436979\\
1.522	-0.0436979\\
1.524	-0.0436979\\
1.526	-0.0436979\\
1.528	-0.0436979\\
1.53	-0.0436979\\
1.532	-0.0436979\\
1.534	-0.0436979\\
1.536	-0.0436979\\
1.538	-0.0436979\\
1.54	-0.0436979\\
1.542	-0.0436979\\
1.544	-0.0436979\\
1.546	-0.0436979\\
1.548	-0.0436979\\
1.55	-0.0436979\\
1.552	-0.0436979\\
1.554	-0.0436979\\
1.556	-0.0436979\\
1.558	-0.0436979\\
1.56	-0.0436979\\
1.562	-0.0436979\\
1.564	-0.0436979\\
1.566	-0.0436979\\
1.568	-0.0436979\\
1.57	-0.0436979\\
1.572	-0.0436979\\
1.574	-0.0436979\\
1.576	-0.0436979\\
1.578	-0.0436979\\
1.58	-0.0436979\\
1.582	-0.0436979\\
1.584	-0.0436979\\
1.586	-0.0436979\\
1.588	-0.0436979\\
1.59	-0.0436979\\
1.592	-0.0436979\\
1.594	-0.0436979\\
1.596	-0.0436979\\
1.598	-0.0436979\\
1.6	-0.0436979\\
1.602	-0.0436979\\
1.604	-0.0436979\\
1.606	-0.0436979\\
1.608	-0.0436979\\
1.61	-0.0436979\\
1.612	-0.0436979\\
1.614	-0.0436979\\
1.616	-0.0436979\\
1.618	-0.0436979\\
1.62	-0.0436979\\
1.622	-0.0436979\\
1.624	-0.0436979\\
1.626	-0.0436979\\
1.628	-0.0436979\\
1.63	-0.0436979\\
1.632	-0.0436979\\
1.634	-0.0436979\\
1.636	-0.0436979\\
1.638	-0.0436979\\
1.64	-0.0436979\\
1.642	-0.0436979\\
1.644	-0.0436979\\
1.646	-0.0436979\\
1.648	-0.0436979\\
1.65	-0.0436979\\
1.652	-0.0436979\\
1.654	-0.0436979\\
1.656	-0.0436979\\
1.658	-0.0436979\\
1.66	-0.0436979\\
1.662	-0.0436979\\
1.664	-0.0436979\\
1.666	-0.0436979\\
1.668	-0.0436979\\
1.67	-0.0436979\\
1.672	-0.0436979\\
1.674	-0.0436979\\
1.676	-0.0436979\\
1.678	-0.0436979\\
1.68	-0.0436979\\
1.682	-0.0436979\\
1.684	-0.0436979\\
1.686	-0.0436979\\
1.688	-0.0436979\\
1.69	-0.0436979\\
1.692	-0.0436979\\
1.694	-0.0436979\\
1.696	-0.0436979\\
1.698	-0.0436979\\
1.7	-0.0436979\\
1.702	-0.0436979\\
1.704	-0.0436979\\
1.706	-0.0436979\\
1.708	-0.0436979\\
1.71	-0.0436979\\
1.712	-0.0436979\\
1.714	-0.0436979\\
1.716	-0.0436979\\
1.718	-0.0436979\\
1.72	-0.0436979\\
1.722	-0.0436979\\
1.724	-0.0436979\\
1.726	-0.0436979\\
1.728	-0.0436979\\
1.73	-0.0436979\\
1.732	-0.0436979\\
1.734	-0.0436979\\
1.736	-0.0436979\\
1.738	-0.0436979\\
1.74	-0.0436979\\
1.742	-0.0436979\\
1.744	-0.0436979\\
1.746	-0.0436979\\
1.748	-0.0436979\\
1.75	-0.0436979\\
1.752	-0.0436979\\
1.754	-0.0436979\\
1.756	-0.0436979\\
1.758	-0.0436979\\
1.76	-0.0436979\\
1.762	-0.0436979\\
1.764	-0.0436979\\
1.766	-0.0436979\\
1.768	-0.0436979\\
1.77	-0.0436979\\
1.772	-0.0436979\\
1.774	-0.0436979\\
1.776	-0.0436979\\
1.778	-0.0436979\\
1.78	-0.0436979\\
1.782	-0.0436979\\
1.784	-0.0436979\\
1.786	-0.0436979\\
1.788	-0.0436979\\
1.79	-0.0436979\\
1.792	-0.0436979\\
1.794	-0.0436979\\
1.796	-0.0436979\\
1.798	-0.0436979\\
1.8	-0.0436979\\
1.802	-0.0436979\\
1.804	-0.0436979\\
1.806	-0.0436979\\
1.808	-0.0436979\\
1.81	-0.0436979\\
1.812	-0.0436979\\
1.814	-0.0436979\\
1.816	-0.0436979\\
1.818	-0.0436979\\
1.82	-0.0436979\\
1.822	-0.0436979\\
1.824	-0.0436979\\
1.826	-0.0436979\\
1.828	-0.0436979\\
1.83	-0.0436979\\
1.832	-0.0436979\\
1.834	-0.0436979\\
1.836	-0.0436979\\
1.838	-0.0436979\\
1.84	-0.0436979\\
1.842	-0.0436979\\
1.844	-0.0436979\\
1.846	-0.0436979\\
1.848	-0.0436979\\
1.85	-0.0436979\\
1.852	-0.0436979\\
1.854	-0.0436979\\
1.856	-0.0436979\\
1.858	-0.0436979\\
1.86	-0.0436979\\
1.862	-0.0436979\\
1.864	-0.0436979\\
1.866	-0.0436979\\
1.868	-0.0436979\\
1.87	-0.0436979\\
1.872	-0.0436979\\
1.874	-0.0436979\\
1.876	-0.0436979\\
1.878	-0.0436979\\
1.88	-0.0436979\\
1.882	-0.0436979\\
1.884	-0.0436979\\
1.886	-0.0436979\\
1.888	-0.0436979\\
1.89	-0.0436979\\
1.892	-0.0436979\\
1.894	-0.0436979\\
1.896	-0.0436979\\
1.898	-0.0436979\\
1.9	-0.0436979\\
1.902	-0.0436979\\
1.904	-0.0436979\\
1.906	-0.0436979\\
1.908	-0.0436979\\
1.91	-0.0436979\\
1.912	-0.0436979\\
1.914	-0.0436979\\
1.916	-0.0436979\\
1.918	-0.0436979\\
1.92	-0.0436979\\
1.922	-0.0436979\\
1.924	-0.0436979\\
1.926	-0.0436979\\
1.928	-0.0436979\\
1.93	-0.0436979\\
1.932	-0.0436979\\
1.934	-0.0436979\\
1.936	-0.0436979\\
1.938	-0.0436979\\
1.94	-0.0436979\\
1.942	-0.0436979\\
1.944	-0.0436979\\
1.946	-0.0436979\\
1.948	-0.0436979\\
1.95	-0.0436979\\
1.952	-0.0436979\\
1.954	-0.0436979\\
1.956	-0.0436979\\
1.958	-0.0436979\\
1.96	-0.0436979\\
1.962	-0.0436979\\
1.964	-0.0436979\\
1.966	-0.0436979\\
1.968	-0.0436979\\
1.97	-0.0436979\\
1.972	-0.0436979\\
1.974	-0.0436979\\
1.976	-0.0436979\\
1.978	-0.0436979\\
1.98	-0.0436979\\
1.982	-0.0436979\\
1.984	-0.0436979\\
1.986	-0.0436979\\
1.988	-0.0436979\\
1.99	-0.0436979\\
1.992	-0.0436979\\
1.994	-0.0436979\\
1.996	-0.0436979\\
1.998	-0.0436979\\
2	-0.0436979\\
};
\addplot [color=mycolor8, line width=1.0pt, forget plot]
  table[row sep=crcr]{%
-0.024	0\\
-0.022	0\\
-0.02	0\\
-0.018	0\\
-0.016	0\\
-0.014	0\\
-0.012	0\\
-0.01	0\\
-0.008	0\\
-0.006	0\\
-0.004	0\\
-0.002	0\\
0	0\\
0.002	-0.0019122\\
0.004	-0.00365961\\
0.006	-0.00525767\\
0.008	-0.00672075\\
0.01	-0.00806224\\
0.012	-0.0092946\\
0.014	-0.0104294\\
0.016	-0.0114772\\
0.018	-0.0124481\\
0.02	-0.013351\\
0.022	-0.0141942\\
0.024	-0.0149855\\
0.026	-0.0157316\\
0.028	-0.0164389\\
0.03	-0.0171129\\
0.032	-0.0177587\\
0.034	-0.0183807\\
0.036	-0.0189828\\
0.038	-0.0195684\\
0.04	-0.0201402\\
0.042	-0.0207009\\
0.044	-0.0212522\\
0.046	-0.021796\\
0.048	-0.0223333\\
0.05	-0.0228651\\
0.052	-0.0233922\\
0.054	-0.0239148\\
0.056	-0.0244331\\
0.058	-0.0249473\\
0.06	-0.025457\\
0.062	-0.0259621\\
0.064	-0.0264623\\
0.066	-0.0269572\\
0.068	-0.0274462\\
0.07	-0.0279291\\
0.072	-0.0284054\\
0.074	-0.0288746\\
0.076	-0.0293366\\
0.078	-0.0297909\\
0.08	-0.0302373\\
0.082	-0.0306756\\
0.084	-0.0311058\\
0.086	-0.0315278\\
0.088	-0.0319417\\
0.09	-0.0323474\\
0.092	-0.0327452\\
0.094	-0.0331352\\
0.096	-0.0335176\\
0.098	-0.0338925\\
0.1	-0.0342604\\
0.102	-0.0346214\\
0.104	-0.0349757\\
0.106	-0.0353236\\
0.108	-0.0356652\\
0.11	-0.0360008\\
0.112	-0.0363304\\
0.114	-0.0366543\\
0.116	-0.0369723\\
0.118	-0.0372845\\
0.12	-0.0375908\\
0.122	-0.0378913\\
0.124	-0.0381856\\
0.126	-0.0384736\\
0.128	-0.0387552\\
0.13	-0.03903\\
0.132	-0.0392978\\
0.134	-0.0395582\\
0.136	-0.0398109\\
0.138	-0.0400556\\
0.14	-0.040292\\
0.142	-0.0405197\\
0.144	-0.0407385\\
0.146	-0.0409479\\
0.148	-0.0411478\\
0.15	-0.0413379\\
0.152	-0.0415181\\
0.154	-0.0416881\\
0.156	-0.0418478\\
0.158	-0.0419972\\
0.16	-0.0421363\\
0.162	-0.042265\\
0.164	-0.0423835\\
0.166	-0.0424919\\
0.168	-0.0425903\\
0.17	-0.042679\\
0.172	-0.0427582\\
0.174	-0.0428283\\
0.176	-0.0428895\\
0.178	-0.0429423\\
0.18	-0.042987\\
0.182	-0.0430241\\
0.184	-0.043054\\
0.186	-0.0430772\\
0.188	-0.0430941\\
0.19	-0.0431052\\
0.192	-0.043111\\
0.194	-0.0431119\\
0.196	-0.0431085\\
0.198	-0.0431011\\
0.2	-0.0430902\\
0.202	-0.0430763\\
0.204	-0.0430598\\
0.206	-0.043041\\
0.208	-0.0430203\\
0.21	-0.0429981\\
0.212	-0.0429748\\
0.214	-0.0429505\\
0.216	-0.0429257\\
0.218	-0.0429005\\
0.22	-0.0428752\\
0.222	-0.0428501\\
0.224	-0.0428252\\
0.226	-0.0428008\\
0.228	-0.042777\\
0.23	-0.042754\\
0.232	-0.0427319\\
0.234	-0.0427108\\
0.236	-0.0426907\\
0.238	-0.0426717\\
0.24	-0.042654\\
0.242	-0.0426374\\
0.244	-0.0426222\\
0.246	-0.0426082\\
0.248	-0.0425956\\
0.25	-0.0425843\\
0.252	-0.0425743\\
0.254	-0.0425656\\
0.256	-0.0425583\\
0.258	-0.0425523\\
0.26	-0.0425476\\
0.262	-0.0425441\\
0.264	-0.0425419\\
0.266	-0.042541\\
0.268	-0.0425412\\
0.27	-0.0425426\\
0.272	-0.0425451\\
0.274	-0.0425488\\
0.276	-0.0425535\\
0.278	-0.0425592\\
0.28	-0.0425659\\
0.282	-0.0425736\\
0.284	-0.0425822\\
0.286	-0.0425916\\
0.288	-0.0426018\\
0.29	-0.0426128\\
0.292	-0.0426246\\
0.294	-0.0426369\\
0.296	-0.04265\\
0.298	-0.0426635\\
0.3	-0.0426776\\
0.302	-0.0426922\\
0.304	-0.0427072\\
0.306	-0.0427226\\
0.308	-0.0427383\\
0.31	-0.0427543\\
0.312	-0.0427705\\
0.314	-0.0427869\\
0.316	-0.0428034\\
0.318	-0.04282\\
0.32	-0.0428367\\
0.322	-0.0428534\\
0.324	-0.0428701\\
0.326	-0.0428867\\
0.328	-0.0429033\\
0.33	-0.0429197\\
0.332	-0.042936\\
0.334	-0.0429521\\
0.336	-0.042968\\
0.338	-0.0429836\\
0.34	-0.0429991\\
0.342	-0.0430143\\
0.344	-0.0430292\\
0.346	-0.0430438\\
0.348	-0.0430582\\
0.35	-0.0430722\\
0.352	-0.043086\\
0.354	-0.0430994\\
0.356	-0.0431125\\
0.358	-0.0431253\\
0.36	-0.0431379\\
0.362	-0.0431501\\
0.364	-0.043162\\
0.366	-0.0431736\\
0.368	-0.0431849\\
0.37	-0.043196\\
0.372	-0.0432067\\
0.374	-0.0432172\\
0.376	-0.0432275\\
0.378	-0.0432375\\
0.38	-0.0432473\\
0.382	-0.0432568\\
0.384	-0.0432662\\
0.386	-0.0432753\\
0.388	-0.0432842\\
0.39	-0.0432929\\
0.392	-0.0433015\\
0.394	-0.0433099\\
0.396	-0.0433181\\
0.398	-0.0433262\\
0.4	-0.0433341\\
0.402	-0.0433419\\
0.404	-0.0433496\\
0.406	-0.0433571\\
0.408	-0.0433645\\
0.41	-0.0433718\\
0.412	-0.043379\\
0.414	-0.043386\\
0.416	-0.043393\\
0.418	-0.0433999\\
0.42	-0.0434066\\
0.422	-0.0434133\\
0.424	-0.0434198\\
0.426	-0.0434263\\
0.428	-0.0434327\\
0.43	-0.043439\\
0.432	-0.0434451\\
0.434	-0.0434512\\
0.436	-0.0434572\\
0.438	-0.0434631\\
0.44	-0.0434689\\
0.442	-0.0434747\\
0.444	-0.0434803\\
0.446	-0.0434858\\
0.448	-0.0434912\\
0.45	-0.0434965\\
0.452	-0.0435018\\
0.454	-0.0435069\\
0.456	-0.0435119\\
0.458	-0.0435169\\
0.46	-0.0435217\\
0.462	-0.0435264\\
0.464	-0.0435311\\
0.466	-0.0435356\\
0.468	-0.04354\\
0.47	-0.0435443\\
0.472	-0.0435486\\
0.474	-0.0435527\\
0.476	-0.0435567\\
0.478	-0.0435606\\
0.48	-0.0435645\\
0.482	-0.0435682\\
0.484	-0.0435718\\
0.486	-0.0435753\\
0.488	-0.0435788\\
0.49	-0.0435821\\
0.492	-0.0435853\\
0.494	-0.0435885\\
0.496	-0.0435915\\
0.498	-0.0435945\\
0.5	-0.0435973\\
0.502	-0.0436001\\
0.504	-0.0436028\\
0.506	-0.0436054\\
0.508	-0.043608\\
0.51	-0.0436104\\
0.512	-0.0436128\\
0.514	-0.0436151\\
0.516	-0.0436173\\
0.518	-0.0436194\\
0.52	-0.0436215\\
0.522	-0.0436235\\
0.524	-0.0436254\\
0.526	-0.0436272\\
0.528	-0.043629\\
0.53	-0.0436308\\
0.532	-0.0436324\\
0.534	-0.043634\\
0.536	-0.0436356\\
0.538	-0.0436371\\
0.54	-0.0436385\\
0.542	-0.0436399\\
0.544	-0.0436412\\
0.546	-0.0436425\\
0.548	-0.0436437\\
0.55	-0.0436449\\
0.552	-0.0436461\\
0.554	-0.0436472\\
0.556	-0.0436483\\
0.558	-0.0436493\\
0.56	-0.0436503\\
0.562	-0.0436512\\
0.564	-0.0436521\\
0.566	-0.043653\\
0.568	-0.0436539\\
0.57	-0.0436547\\
0.572	-0.0436555\\
0.574	-0.0436563\\
0.576	-0.043657\\
0.578	-0.0436577\\
0.58	-0.0436584\\
0.582	-0.0436591\\
0.584	-0.0436597\\
0.586	-0.0436603\\
0.588	-0.0436609\\
0.59	-0.0436615\\
0.592	-0.0436621\\
0.594	-0.0436626\\
0.596	-0.0436631\\
0.598	-0.0436636\\
0.6	-0.0436641\\
0.602	-0.0436646\\
0.604	-0.0436651\\
0.606	-0.0436655\\
0.608	-0.0436659\\
0.61	-0.0436664\\
0.612	-0.0436668\\
0.614	-0.0436672\\
0.616	-0.0436675\\
0.618	-0.0436679\\
0.62	-0.0436683\\
0.622	-0.0436686\\
0.624	-0.0436689\\
0.626	-0.0436693\\
0.628	-0.0436696\\
0.63	-0.0436699\\
0.632	-0.0436702\\
0.634	-0.0436705\\
0.636	-0.0436708\\
0.638	-0.043671\\
0.64	-0.0436713\\
0.642	-0.0436716\\
0.644	-0.0436718\\
0.646	-0.043672\\
0.648	-0.0436723\\
0.65	-0.0436725\\
0.652	-0.0436727\\
0.654	-0.043673\\
0.656	-0.0436732\\
0.658	-0.0436734\\
0.66	-0.0436736\\
0.662	-0.0436738\\
0.664	-0.043674\\
0.666	-0.0436741\\
0.668	-0.0436743\\
0.67	-0.0436745\\
0.672	-0.0436747\\
0.674	-0.0436748\\
0.676	-0.043675\\
0.678	-0.0436752\\
0.68	-0.0436753\\
0.682	-0.0436755\\
0.684	-0.0436756\\
0.686	-0.0436758\\
0.688	-0.0436759\\
0.69	-0.0436761\\
0.692	-0.0436762\\
0.694	-0.0436763\\
0.696	-0.0436765\\
0.698	-0.0436766\\
0.7	-0.0436768\\
0.702	-0.0436769\\
0.704	-0.043677\\
0.706	-0.0436771\\
0.708	-0.0436773\\
0.71	-0.0436774\\
0.712	-0.0436775\\
0.714	-0.0436776\\
0.716	-0.0436778\\
0.718	-0.0436779\\
0.72	-0.043678\\
0.722	-0.0436781\\
0.724	-0.0436783\\
0.726	-0.0436784\\
0.728	-0.0436785\\
0.73	-0.0436786\\
0.732	-0.0436788\\
0.734	-0.0436789\\
0.736	-0.043679\\
0.738	-0.0436791\\
0.74	-0.0436793\\
0.742	-0.0436794\\
0.744	-0.0436795\\
0.746	-0.0436797\\
0.748	-0.0436798\\
0.75	-0.0436799\\
0.752	-0.0436801\\
0.754	-0.0436802\\
0.756	-0.0436804\\
0.758	-0.0436805\\
0.76	-0.0436807\\
0.762	-0.0436808\\
0.764	-0.043681\\
0.766	-0.0436812\\
0.768	-0.0436813\\
0.77	-0.0436815\\
0.772	-0.0436817\\
0.774	-0.0436819\\
0.776	-0.043682\\
0.778	-0.0436822\\
0.78	-0.0436824\\
0.782	-0.0436826\\
0.784	-0.0436828\\
0.786	-0.043683\\
0.788	-0.0436832\\
0.79	-0.0436834\\
0.792	-0.0436836\\
0.794	-0.0436838\\
0.796	-0.043684\\
0.798	-0.0436842\\
0.8	-0.0436845\\
0.802	-0.0436847\\
0.804	-0.0436849\\
0.806	-0.0436851\\
0.808	-0.0436854\\
0.81	-0.0436856\\
0.812	-0.0436859\\
0.814	-0.0436861\\
0.816	-0.0436863\\
0.818	-0.0436866\\
0.82	-0.0436868\\
0.822	-0.0436871\\
0.824	-0.0436873\\
0.826	-0.0436876\\
0.828	-0.0436879\\
0.83	-0.0436881\\
0.832	-0.0436884\\
0.834	-0.0436886\\
0.836	-0.0436889\\
0.838	-0.0436891\\
0.84	-0.0436894\\
0.842	-0.0436897\\
0.844	-0.0436899\\
0.846	-0.0436902\\
0.848	-0.0436904\\
0.85	-0.0436907\\
0.852	-0.0436909\\
0.854	-0.0436912\\
0.856	-0.0436914\\
0.858	-0.0436917\\
0.86	-0.0436919\\
0.862	-0.0436922\\
0.864	-0.0436924\\
0.866	-0.0436927\\
0.868	-0.0436929\\
0.87	-0.0436931\\
0.872	-0.0436933\\
0.874	-0.0436936\\
0.876	-0.0436938\\
0.878	-0.043694\\
0.88	-0.0436942\\
0.882	-0.0436944\\
0.884	-0.0436946\\
0.886	-0.0436948\\
0.888	-0.043695\\
0.89	-0.0436952\\
0.892	-0.0436953\\
0.894	-0.0436955\\
0.896	-0.0436957\\
0.898	-0.0436958\\
0.9	-0.043696\\
0.902	-0.0436961\\
0.904	-0.0436963\\
0.906	-0.0436964\\
0.908	-0.0436965\\
0.91	-0.0436967\\
0.912	-0.0436968\\
0.914	-0.0436969\\
0.916	-0.043697\\
0.918	-0.0436971\\
0.92	-0.0436972\\
0.922	-0.0436973\\
0.924	-0.0436973\\
0.926	-0.0436974\\
0.928	-0.0436975\\
0.93	-0.0436975\\
0.932	-0.0436976\\
0.934	-0.0436976\\
0.936	-0.0436977\\
0.938	-0.0436977\\
0.94	-0.0436977\\
0.942	-0.0436978\\
0.944	-0.0436978\\
0.946	-0.0436978\\
0.948	-0.0436978\\
0.95	-0.0436978\\
0.952	-0.0436978\\
0.954	-0.0436978\\
0.956	-0.0436978\\
0.958	-0.0436978\\
0.96	-0.0436978\\
0.962	-0.0436978\\
0.964	-0.0436978\\
0.966	-0.0436978\\
0.968	-0.0436977\\
0.97	-0.0436977\\
0.972	-0.0436977\\
0.974	-0.0436976\\
0.976	-0.0436976\\
0.978	-0.0436976\\
0.98	-0.0436975\\
0.982	-0.0436975\\
0.984	-0.0436975\\
0.986	-0.0436974\\
0.988	-0.0436974\\
0.99	-0.0436973\\
0.992	-0.0436973\\
0.994	-0.0436972\\
0.996	-0.0436972\\
0.998	-0.0436971\\
1	-0.0436971\\
1.002	-0.043697\\
1.004	-0.043697\\
1.006	-0.043697\\
1.008	-0.0436969\\
1.01	-0.0436969\\
1.012	-0.0436968\\
1.014	-0.0436968\\
1.016	-0.0436967\\
1.018	-0.0436967\\
1.02	-0.0436967\\
1.022	-0.0436966\\
1.024	-0.0436966\\
1.026	-0.0436966\\
1.028	-0.0436965\\
1.03	-0.0436965\\
1.032	-0.0436965\\
1.034	-0.0436964\\
1.036	-0.0436964\\
1.038	-0.0436964\\
1.04	-0.0436964\\
1.042	-0.0436963\\
1.044	-0.0436963\\
1.046	-0.0436963\\
1.048	-0.0436963\\
1.05	-0.0436963\\
1.052	-0.0436963\\
1.054	-0.0436962\\
1.056	-0.0436962\\
1.058	-0.0436962\\
1.06	-0.0436962\\
1.062	-0.0436962\\
1.064	-0.0436962\\
1.066	-0.0436962\\
1.068	-0.0436962\\
1.07	-0.0436962\\
1.072	-0.0436962\\
1.074	-0.0436962\\
1.076	-0.0436962\\
1.078	-0.0436963\\
1.08	-0.0436963\\
1.082	-0.0436963\\
1.084	-0.0436963\\
1.086	-0.0436963\\
1.088	-0.0436963\\
1.09	-0.0436964\\
1.092	-0.0436964\\
1.094	-0.0436964\\
1.096	-0.0436964\\
1.098	-0.0436964\\
1.1	-0.0436965\\
1.102	-0.0436965\\
1.104	-0.0436965\\
1.106	-0.0436966\\
1.108	-0.0436966\\
1.11	-0.0436966\\
1.112	-0.0436966\\
1.114	-0.0436967\\
1.116	-0.0436967\\
1.118	-0.0436967\\
1.12	-0.0436968\\
1.122	-0.0436968\\
1.124	-0.0436968\\
1.126	-0.0436969\\
1.128	-0.0436969\\
1.13	-0.0436969\\
1.132	-0.043697\\
1.134	-0.043697\\
1.136	-0.043697\\
1.138	-0.043697\\
1.14	-0.0436971\\
1.142	-0.0436971\\
1.144	-0.0436971\\
1.146	-0.0436972\\
1.148	-0.0436972\\
1.15	-0.0436972\\
1.152	-0.0436973\\
1.154	-0.0436973\\
1.156	-0.0436973\\
1.158	-0.0436974\\
1.16	-0.0436974\\
1.162	-0.0436974\\
1.164	-0.0436974\\
1.166	-0.0436975\\
1.168	-0.0436975\\
1.17	-0.0436975\\
1.172	-0.0436976\\
1.174	-0.0436976\\
1.176	-0.0436976\\
1.178	-0.0436976\\
1.18	-0.0436976\\
1.182	-0.0436977\\
1.184	-0.0436977\\
1.186	-0.0436977\\
1.188	-0.0436977\\
1.19	-0.0436977\\
1.192	-0.0436978\\
1.194	-0.0436978\\
1.196	-0.0436978\\
1.198	-0.0436978\\
1.2	-0.0436978\\
1.202	-0.0436978\\
1.204	-0.0436979\\
1.206	-0.0436979\\
1.208	-0.0436979\\
1.21	-0.0436979\\
1.212	-0.0436979\\
1.214	-0.0436979\\
1.216	-0.0436979\\
1.218	-0.0436979\\
1.22	-0.0436979\\
1.222	-0.043698\\
1.224	-0.043698\\
1.226	-0.043698\\
1.228	-0.043698\\
1.23	-0.043698\\
1.232	-0.043698\\
1.234	-0.043698\\
1.236	-0.043698\\
1.238	-0.043698\\
1.24	-0.043698\\
1.242	-0.043698\\
1.244	-0.043698\\
1.246	-0.043698\\
1.248	-0.043698\\
1.25	-0.043698\\
1.252	-0.043698\\
1.254	-0.043698\\
1.256	-0.043698\\
1.258	-0.043698\\
1.26	-0.043698\\
1.262	-0.043698\\
1.264	-0.043698\\
1.266	-0.043698\\
1.268	-0.043698\\
1.27	-0.043698\\
1.272	-0.043698\\
1.274	-0.043698\\
1.276	-0.043698\\
1.278	-0.043698\\
1.28	-0.0436979\\
1.282	-0.0436979\\
1.284	-0.0436979\\
1.286	-0.0436979\\
1.288	-0.0436979\\
1.29	-0.0436979\\
1.292	-0.0436979\\
1.294	-0.0436979\\
1.296	-0.0436979\\
1.298	-0.0436979\\
1.3	-0.0436979\\
1.302	-0.0436979\\
1.304	-0.0436979\\
1.306	-0.0436979\\
1.308	-0.0436979\\
1.31	-0.0436979\\
1.312	-0.0436979\\
1.314	-0.0436979\\
1.316	-0.0436979\\
1.318	-0.0436978\\
1.32	-0.0436978\\
1.322	-0.0436978\\
1.324	-0.0436978\\
1.326	-0.0436978\\
1.328	-0.0436978\\
1.33	-0.0436978\\
1.332	-0.0436978\\
1.334	-0.0436978\\
1.336	-0.0436978\\
1.338	-0.0436978\\
1.34	-0.0436978\\
1.342	-0.0436978\\
1.344	-0.0436978\\
1.346	-0.0436978\\
1.348	-0.0436978\\
1.35	-0.0436978\\
1.352	-0.0436978\\
1.354	-0.0436978\\
1.356	-0.0436978\\
1.358	-0.0436978\\
1.36	-0.0436978\\
1.362	-0.0436978\\
1.364	-0.0436978\\
1.366	-0.0436978\\
1.368	-0.0436978\\
1.37	-0.0436978\\
1.372	-0.0436978\\
1.374	-0.0436978\\
1.376	-0.0436978\\
1.378	-0.0436978\\
1.38	-0.0436978\\
1.382	-0.0436978\\
1.384	-0.0436978\\
1.386	-0.0436978\\
1.388	-0.0436978\\
1.39	-0.0436978\\
1.392	-0.0436978\\
1.394	-0.0436978\\
1.396	-0.0436978\\
1.398	-0.0436978\\
1.4	-0.0436978\\
1.402	-0.0436978\\
1.404	-0.0436978\\
1.406	-0.0436978\\
1.408	-0.0436978\\
1.41	-0.0436978\\
1.412	-0.0436978\\
1.414	-0.0436978\\
1.416	-0.0436978\\
1.418	-0.0436978\\
1.42	-0.0436978\\
1.422	-0.0436978\\
1.424	-0.0436978\\
1.426	-0.0436978\\
1.428	-0.0436978\\
1.43	-0.0436978\\
1.432	-0.0436978\\
1.434	-0.0436978\\
1.436	-0.0436978\\
1.438	-0.0436978\\
1.44	-0.0436978\\
1.442	-0.0436978\\
1.444	-0.0436979\\
1.446	-0.0436979\\
1.448	-0.0436979\\
1.45	-0.0436979\\
1.452	-0.0436979\\
1.454	-0.0436979\\
1.456	-0.0436979\\
1.458	-0.0436979\\
1.46	-0.0436979\\
1.462	-0.0436979\\
1.464	-0.0436979\\
1.466	-0.0436979\\
1.468	-0.0436979\\
1.47	-0.0436979\\
1.472	-0.0436979\\
1.474	-0.0436979\\
1.476	-0.0436979\\
1.478	-0.0436979\\
1.48	-0.0436979\\
1.482	-0.0436979\\
1.484	-0.0436979\\
1.486	-0.0436979\\
1.488	-0.0436979\\
1.49	-0.0436979\\
1.492	-0.0436979\\
1.494	-0.0436979\\
1.496	-0.0436979\\
1.498	-0.0436979\\
1.5	-0.0436979\\
1.502	-0.0436979\\
1.504	-0.0436979\\
1.506	-0.0436979\\
1.508	-0.0436979\\
1.51	-0.0436979\\
1.512	-0.0436979\\
1.514	-0.0436979\\
1.516	-0.0436979\\
1.518	-0.0436979\\
1.52	-0.0436979\\
1.522	-0.0436979\\
1.524	-0.0436979\\
1.526	-0.0436979\\
1.528	-0.0436979\\
1.53	-0.0436979\\
1.532	-0.0436979\\
1.534	-0.0436979\\
1.536	-0.0436979\\
1.538	-0.0436979\\
1.54	-0.0436979\\
1.542	-0.0436979\\
1.544	-0.0436979\\
1.546	-0.0436979\\
1.548	-0.0436979\\
1.55	-0.0436979\\
1.552	-0.0436979\\
1.554	-0.0436979\\
1.556	-0.0436979\\
1.558	-0.0436979\\
1.56	-0.0436979\\
1.562	-0.0436979\\
1.564	-0.0436979\\
1.566	-0.0436979\\
1.568	-0.0436979\\
1.57	-0.0436979\\
1.572	-0.0436979\\
1.574	-0.0436979\\
1.576	-0.0436979\\
1.578	-0.0436979\\
1.58	-0.0436979\\
1.582	-0.0436979\\
1.584	-0.0436979\\
1.586	-0.0436979\\
1.588	-0.0436979\\
1.59	-0.0436979\\
1.592	-0.0436979\\
1.594	-0.0436979\\
1.596	-0.0436979\\
1.598	-0.0436979\\
1.6	-0.0436979\\
1.602	-0.0436979\\
1.604	-0.0436979\\
1.606	-0.0436979\\
1.608	-0.0436979\\
1.61	-0.0436979\\
1.612	-0.0436979\\
1.614	-0.0436979\\
1.616	-0.0436979\\
1.618	-0.0436979\\
1.62	-0.0436979\\
1.622	-0.0436979\\
1.624	-0.0436979\\
1.626	-0.0436979\\
1.628	-0.0436979\\
1.63	-0.0436979\\
1.632	-0.0436979\\
1.634	-0.0436979\\
1.636	-0.0436979\\
1.638	-0.0436979\\
1.64	-0.0436979\\
1.642	-0.0436979\\
1.644	-0.0436979\\
1.646	-0.0436979\\
1.648	-0.0436979\\
1.65	-0.0436979\\
1.652	-0.0436979\\
1.654	-0.0436979\\
1.656	-0.0436979\\
1.658	-0.0436979\\
1.66	-0.0436979\\
1.662	-0.0436979\\
1.664	-0.0436979\\
1.666	-0.0436979\\
1.668	-0.0436979\\
1.67	-0.0436979\\
1.672	-0.0436979\\
1.674	-0.0436979\\
1.676	-0.0436979\\
1.678	-0.0436979\\
1.68	-0.0436979\\
1.682	-0.0436979\\
1.684	-0.0436979\\
1.686	-0.0436979\\
1.688	-0.0436979\\
1.69	-0.0436979\\
1.692	-0.0436979\\
1.694	-0.0436979\\
1.696	-0.0436979\\
1.698	-0.0436979\\
1.7	-0.0436979\\
1.702	-0.0436979\\
1.704	-0.0436979\\
1.706	-0.0436979\\
1.708	-0.0436979\\
1.71	-0.0436979\\
1.712	-0.0436979\\
1.714	-0.0436979\\
1.716	-0.0436979\\
1.718	-0.0436979\\
1.72	-0.0436979\\
1.722	-0.0436979\\
1.724	-0.0436979\\
1.726	-0.0436979\\
1.728	-0.0436979\\
1.73	-0.0436979\\
1.732	-0.0436979\\
1.734	-0.0436979\\
1.736	-0.0436979\\
1.738	-0.0436979\\
1.74	-0.0436979\\
1.742	-0.0436979\\
1.744	-0.0436979\\
1.746	-0.0436979\\
1.748	-0.0436979\\
1.75	-0.0436979\\
1.752	-0.0436979\\
1.754	-0.0436979\\
1.756	-0.0436979\\
1.758	-0.0436979\\
1.76	-0.0436979\\
1.762	-0.0436979\\
1.764	-0.0436979\\
1.766	-0.0436979\\
1.768	-0.0436979\\
1.77	-0.0436979\\
1.772	-0.0436979\\
1.774	-0.0436979\\
1.776	-0.0436979\\
1.778	-0.0436979\\
1.78	-0.0436979\\
1.782	-0.0436979\\
1.784	-0.0436979\\
1.786	-0.0436979\\
1.788	-0.0436979\\
1.79	-0.0436979\\
1.792	-0.0436979\\
1.794	-0.0436979\\
1.796	-0.0436979\\
1.798	-0.0436979\\
1.8	-0.0436979\\
1.802	-0.0436979\\
1.804	-0.0436979\\
1.806	-0.0436979\\
1.808	-0.0436979\\
1.81	-0.0436979\\
1.812	-0.0436979\\
1.814	-0.0436979\\
1.816	-0.0436979\\
1.818	-0.0436979\\
1.82	-0.0436979\\
1.822	-0.0436979\\
1.824	-0.0436979\\
1.826	-0.0436979\\
1.828	-0.0436979\\
1.83	-0.0436979\\
1.832	-0.0436979\\
1.834	-0.0436979\\
1.836	-0.0436979\\
1.838	-0.0436979\\
1.84	-0.0436979\\
1.842	-0.0436979\\
1.844	-0.0436979\\
1.846	-0.0436979\\
1.848	-0.0436979\\
1.85	-0.0436979\\
1.852	-0.0436979\\
1.854	-0.0436979\\
1.856	-0.0436979\\
1.858	-0.0436979\\
1.86	-0.0436979\\
1.862	-0.0436979\\
1.864	-0.0436979\\
1.866	-0.0436979\\
1.868	-0.0436979\\
1.87	-0.0436979\\
1.872	-0.0436979\\
1.874	-0.0436979\\
1.876	-0.0436979\\
1.878	-0.0436979\\
1.88	-0.0436979\\
1.882	-0.0436979\\
1.884	-0.0436979\\
1.886	-0.0436979\\
1.888	-0.0436979\\
1.89	-0.0436979\\
1.892	-0.0436979\\
1.894	-0.0436979\\
1.896	-0.0436979\\
1.898	-0.0436979\\
1.9	-0.0436979\\
1.902	-0.0436979\\
1.904	-0.0436979\\
1.906	-0.0436979\\
1.908	-0.0436979\\
1.91	-0.0436979\\
1.912	-0.0436979\\
1.914	-0.0436979\\
1.916	-0.0436979\\
1.918	-0.0436979\\
1.92	-0.0436979\\
1.922	-0.0436979\\
1.924	-0.0436979\\
1.926	-0.0436979\\
1.928	-0.0436979\\
1.93	-0.0436979\\
1.932	-0.0436979\\
1.934	-0.0436979\\
1.936	-0.0436979\\
1.938	-0.0436979\\
1.94	-0.0436979\\
1.942	-0.0436979\\
1.944	-0.0436979\\
1.946	-0.0436979\\
1.948	-0.0436979\\
1.95	-0.0436979\\
1.952	-0.0436979\\
1.954	-0.0436979\\
1.956	-0.0436979\\
1.958	-0.0436979\\
1.96	-0.0436979\\
1.962	-0.0436979\\
1.964	-0.0436979\\
1.966	-0.0436979\\
1.968	-0.0436979\\
1.97	-0.0436979\\
1.972	-0.0436979\\
1.974	-0.0436979\\
1.976	-0.0436979\\
1.978	-0.0436979\\
1.98	-0.0436979\\
1.982	-0.0436979\\
1.984	-0.0436979\\
1.986	-0.0436979\\
1.988	-0.0436979\\
1.99	-0.0436979\\
1.992	-0.0436979\\
1.994	-0.0436979\\
1.996	-0.0436979\\
1.998	-0.0436979\\
2	-0.0436979\\
};
\addplot [color=mycolor9, line width=1.0pt, forget plot]
  table[row sep=crcr]{%
-0.024	0\\
-0.022	0\\
-0.02	0\\
-0.018	0\\
-0.016	0\\
-0.014	0\\
-0.012	0\\
-0.01	0\\
-0.008	0\\
-0.006	0\\
-0.004	0\\
-0.002	0\\
0	0\\
0.002	-0.00191221\\
0.004	-0.0036597\\
0.006	-0.00525797\\
0.008	-0.00672144\\
0.01	-0.00806359\\
0.012	-0.00929691\\
0.014	-0.010433\\
0.016	-0.0114827\\
0.018	-0.0124558\\
0.02	-0.0133616\\
0.022	-0.0142085\\
0.024	-0.015004\\
0.026	-0.0157554\\
0.028	-0.0164687\\
0.03	-0.0171498\\
0.032	-0.0178038\\
0.034	-0.018435\\
0.036	-0.0190475\\
0.038	-0.0196446\\
0.04	-0.0202293\\
0.042	-0.0208039\\
0.044	-0.0213704\\
0.046	-0.0219303\\
0.048	-0.022485\\
0.05	-0.0230353\\
0.052	-0.0235816\\
0.054	-0.0241243\\
0.056	-0.0246635\\
0.058	-0.025199\\
0.06	-0.0257306\\
0.062	-0.0262577\\
0.064	-0.02678\\
0.066	-0.0272967\\
0.068	-0.0278074\\
0.07	-0.0283113\\
0.072	-0.0288078\\
0.074	-0.0292964\\
0.076	-0.0297765\\
0.078	-0.0302475\\
0.08	-0.0307092\\
0.082	-0.0311611\\
0.084	-0.0316029\\
0.086	-0.0320345\\
0.088	-0.0324558\\
0.09	-0.0328667\\
0.092	-0.0332672\\
0.094	-0.0336574\\
0.096	-0.0340376\\
0.098	-0.0344079\\
0.1	-0.0347685\\
0.102	-0.0351197\\
0.104	-0.0354618\\
0.106	-0.035795\\
0.108	-0.0361198\\
0.11	-0.0364362\\
0.112	-0.0367447\\
0.114	-0.0370455\\
0.116	-0.0373387\\
0.118	-0.0376245\\
0.12	-0.0379031\\
0.122	-0.0381746\\
0.124	-0.038439\\
0.126	-0.0386964\\
0.128	-0.0389468\\
0.13	-0.0391901\\
0.132	-0.0394263\\
0.134	-0.0396553\\
0.136	-0.0398768\\
0.138	-0.0400909\\
0.14	-0.0402974\\
0.142	-0.040496\\
0.144	-0.0406867\\
0.146	-0.0408694\\
0.148	-0.0410438\\
0.15	-0.0412099\\
0.152	-0.0413676\\
0.154	-0.0415168\\
0.156	-0.0416574\\
0.158	-0.0417895\\
0.16	-0.041913\\
0.162	-0.0420281\\
0.164	-0.0421347\\
0.166	-0.042233\\
0.168	-0.0423232\\
0.17	-0.0424055\\
0.172	-0.04248\\
0.174	-0.042547\\
0.176	-0.0426068\\
0.178	-0.0426596\\
0.18	-0.0427059\\
0.182	-0.0427458\\
0.184	-0.0427798\\
0.186	-0.0428083\\
0.188	-0.0428315\\
0.19	-0.0428499\\
0.192	-0.0428638\\
0.194	-0.0428736\\
0.196	-0.0428797\\
0.198	-0.0428825\\
0.2	-0.0428822\\
0.202	-0.0428793\\
0.204	-0.0428739\\
0.206	-0.0428666\\
0.208	-0.0428575\\
0.21	-0.0428469\\
0.212	-0.0428352\\
0.214	-0.0428224\\
0.216	-0.0428089\\
0.218	-0.0427949\\
0.22	-0.0427805\\
0.222	-0.042766\\
0.224	-0.0427514\\
0.226	-0.042737\\
0.228	-0.0427227\\
0.23	-0.0427088\\
0.232	-0.0426954\\
0.234	-0.0426824\\
0.236	-0.0426701\\
0.238	-0.0426584\\
0.24	-0.0426473\\
0.242	-0.042637\\
0.244	-0.0426275\\
0.246	-0.0426187\\
0.248	-0.0426108\\
0.25	-0.0426036\\
0.252	-0.0425973\\
0.254	-0.0425918\\
0.256	-0.0425872\\
0.258	-0.0425833\\
0.26	-0.0425803\\
0.262	-0.0425781\\
0.264	-0.0425767\\
0.266	-0.0425761\\
0.268	-0.0425762\\
0.27	-0.0425771\\
0.272	-0.0425788\\
0.274	-0.0425812\\
0.276	-0.0425843\\
0.278	-0.042588\\
0.28	-0.0425925\\
0.282	-0.0425976\\
0.284	-0.0426033\\
0.286	-0.0426096\\
0.288	-0.0426165\\
0.29	-0.0426239\\
0.292	-0.0426319\\
0.294	-0.0426404\\
0.296	-0.0426494\\
0.298	-0.0426588\\
0.3	-0.0426686\\
0.302	-0.0426789\\
0.304	-0.0426895\\
0.306	-0.0427005\\
0.308	-0.0427117\\
0.31	-0.0427233\\
0.312	-0.0427352\\
0.314	-0.0427472\\
0.316	-0.0427595\\
0.318	-0.042772\\
0.32	-0.0427847\\
0.322	-0.0427975\\
0.324	-0.0428104\\
0.326	-0.0428234\\
0.328	-0.0428365\\
0.33	-0.0428496\\
0.332	-0.0428627\\
0.334	-0.0428759\\
0.336	-0.0428891\\
0.338	-0.0429023\\
0.34	-0.0429154\\
0.342	-0.0429285\\
0.344	-0.0429415\\
0.346	-0.0429544\\
0.348	-0.0429673\\
0.35	-0.0429801\\
0.352	-0.0429927\\
0.354	-0.0430053\\
0.356	-0.0430177\\
0.358	-0.0430301\\
0.36	-0.0430423\\
0.362	-0.0430543\\
0.364	-0.0430663\\
0.366	-0.0430781\\
0.368	-0.0430897\\
0.37	-0.0431012\\
0.372	-0.0431126\\
0.374	-0.0431238\\
0.376	-0.0431349\\
0.378	-0.0431458\\
0.38	-0.0431566\\
0.382	-0.0431673\\
0.384	-0.0431778\\
0.386	-0.0431882\\
0.388	-0.0431984\\
0.39	-0.0432085\\
0.392	-0.0432184\\
0.394	-0.0432282\\
0.396	-0.0432379\\
0.398	-0.0432474\\
0.4	-0.0432568\\
0.402	-0.0432661\\
0.404	-0.0432752\\
0.406	-0.0432842\\
0.408	-0.043293\\
0.41	-0.0433018\\
0.412	-0.0433103\\
0.414	-0.0433188\\
0.416	-0.0433271\\
0.418	-0.0433353\\
0.42	-0.0433434\\
0.422	-0.0433513\\
0.424	-0.0433591\\
0.426	-0.0433668\\
0.428	-0.0433743\\
0.43	-0.0433817\\
0.432	-0.0433889\\
0.434	-0.0433961\\
0.436	-0.0434031\\
0.438	-0.0434099\\
0.44	-0.0434167\\
0.442	-0.0434233\\
0.444	-0.0434297\\
0.446	-0.0434361\\
0.448	-0.0434423\\
0.45	-0.0434484\\
0.452	-0.0434543\\
0.454	-0.0434601\\
0.456	-0.0434658\\
0.458	-0.0434713\\
0.46	-0.0434767\\
0.462	-0.043482\\
0.464	-0.0434872\\
0.466	-0.0434922\\
0.468	-0.0434971\\
0.47	-0.0435019\\
0.472	-0.0435066\\
0.474	-0.0435111\\
0.476	-0.0435156\\
0.478	-0.0435199\\
0.48	-0.0435241\\
0.482	-0.0435282\\
0.484	-0.0435321\\
0.486	-0.043536\\
0.488	-0.0435398\\
0.49	-0.0435434\\
0.492	-0.043547\\
0.494	-0.0435504\\
0.496	-0.0435538\\
0.498	-0.043557\\
0.5	-0.0435602\\
0.502	-0.0435633\\
0.504	-0.0435662\\
0.506	-0.0435691\\
0.508	-0.0435719\\
0.51	-0.0435747\\
0.512	-0.0435773\\
0.514	-0.0435799\\
0.516	-0.0435824\\
0.518	-0.0435848\\
0.52	-0.0435872\\
0.522	-0.0435894\\
0.524	-0.0435917\\
0.526	-0.0435938\\
0.528	-0.0435959\\
0.53	-0.0435979\\
0.532	-0.0435999\\
0.534	-0.0436018\\
0.536	-0.0436037\\
0.538	-0.0436055\\
0.54	-0.0436073\\
0.542	-0.043609\\
0.544	-0.0436106\\
0.546	-0.0436122\\
0.548	-0.0436138\\
0.55	-0.0436154\\
0.552	-0.0436168\\
0.554	-0.0436183\\
0.556	-0.0436197\\
0.558	-0.0436211\\
0.56	-0.0436224\\
0.562	-0.0436237\\
0.564	-0.043625\\
0.566	-0.0436262\\
0.568	-0.0436274\\
0.57	-0.0436286\\
0.572	-0.0436297\\
0.574	-0.0436309\\
0.576	-0.043632\\
0.578	-0.043633\\
0.58	-0.0436341\\
0.582	-0.0436351\\
0.584	-0.0436361\\
0.586	-0.043637\\
0.588	-0.043638\\
0.59	-0.0436389\\
0.592	-0.0436398\\
0.594	-0.0436406\\
0.596	-0.0436415\\
0.598	-0.0436423\\
0.6	-0.0436431\\
0.602	-0.0436439\\
0.604	-0.0436447\\
0.606	-0.0436455\\
0.608	-0.0436462\\
0.61	-0.0436469\\
0.612	-0.0436477\\
0.614	-0.0436483\\
0.616	-0.043649\\
0.618	-0.0436497\\
0.62	-0.0436503\\
0.622	-0.0436509\\
0.624	-0.0436516\\
0.626	-0.0436522\\
0.628	-0.0436527\\
0.63	-0.0436533\\
0.632	-0.0436539\\
0.634	-0.0436544\\
0.636	-0.043655\\
0.638	-0.0436555\\
0.64	-0.043656\\
0.642	-0.0436565\\
0.644	-0.043657\\
0.646	-0.0436574\\
0.648	-0.0436579\\
0.65	-0.0436583\\
0.652	-0.0436588\\
0.654	-0.0436592\\
0.656	-0.0436596\\
0.658	-0.0436601\\
0.66	-0.0436605\\
0.662	-0.0436609\\
0.664	-0.0436612\\
0.666	-0.0436616\\
0.668	-0.043662\\
0.67	-0.0436623\\
0.672	-0.0436627\\
0.674	-0.043663\\
0.676	-0.0436634\\
0.678	-0.0436637\\
0.68	-0.043664\\
0.682	-0.0436644\\
0.684	-0.0436647\\
0.686	-0.043665\\
0.688	-0.0436653\\
0.69	-0.0436656\\
0.692	-0.0436659\\
0.694	-0.0436661\\
0.696	-0.0436664\\
0.698	-0.0436667\\
0.7	-0.043667\\
0.702	-0.0436673\\
0.704	-0.0436675\\
0.706	-0.0436678\\
0.708	-0.043668\\
0.71	-0.0436683\\
0.712	-0.0436685\\
0.714	-0.0436688\\
0.716	-0.043669\\
0.718	-0.0436693\\
0.72	-0.0436695\\
0.722	-0.0436698\\
0.724	-0.04367\\
0.726	-0.0436703\\
0.728	-0.0436705\\
0.73	-0.0436708\\
0.732	-0.043671\\
0.734	-0.0436712\\
0.736	-0.0436715\\
0.738	-0.0436717\\
0.74	-0.043672\\
0.742	-0.0436722\\
0.744	-0.0436724\\
0.746	-0.0436727\\
0.748	-0.0436729\\
0.75	-0.0436732\\
0.752	-0.0436734\\
0.754	-0.0436737\\
0.756	-0.0436739\\
0.758	-0.0436742\\
0.76	-0.0436744\\
0.762	-0.0436747\\
0.764	-0.0436749\\
0.766	-0.0436752\\
0.768	-0.0436754\\
0.77	-0.0436757\\
0.772	-0.043676\\
0.774	-0.0436762\\
0.776	-0.0436765\\
0.778	-0.0436768\\
0.78	-0.043677\\
0.782	-0.0436773\\
0.784	-0.0436776\\
0.786	-0.0436779\\
0.788	-0.0436782\\
0.79	-0.0436784\\
0.792	-0.0436787\\
0.794	-0.043679\\
0.796	-0.0436793\\
0.798	-0.0436796\\
0.8	-0.0436799\\
0.802	-0.0436802\\
0.804	-0.0436805\\
0.806	-0.0436808\\
0.808	-0.0436811\\
0.81	-0.0436814\\
0.812	-0.0436817\\
0.814	-0.043682\\
0.816	-0.0436823\\
0.818	-0.0436826\\
0.82	-0.0436829\\
0.822	-0.0436832\\
0.824	-0.0436835\\
0.826	-0.0436838\\
0.828	-0.0436842\\
0.83	-0.0436845\\
0.832	-0.0436848\\
0.834	-0.0436851\\
0.836	-0.0436854\\
0.838	-0.0436857\\
0.84	-0.043686\\
0.842	-0.0436863\\
0.844	-0.0436866\\
0.846	-0.0436869\\
0.848	-0.0436872\\
0.85	-0.0436875\\
0.852	-0.0436878\\
0.854	-0.043688\\
0.856	-0.0436883\\
0.858	-0.0436886\\
0.86	-0.0436889\\
0.862	-0.0436892\\
0.864	-0.0436894\\
0.866	-0.0436897\\
0.868	-0.04369\\
0.87	-0.0436902\\
0.872	-0.0436905\\
0.874	-0.0436907\\
0.876	-0.0436909\\
0.878	-0.0436912\\
0.88	-0.0436914\\
0.882	-0.0436916\\
0.884	-0.0436919\\
0.886	-0.0436921\\
0.888	-0.0436923\\
0.89	-0.0436925\\
0.892	-0.0436927\\
0.894	-0.0436929\\
0.896	-0.0436931\\
0.898	-0.0436933\\
0.9	-0.0436934\\
0.902	-0.0436936\\
0.904	-0.0436938\\
0.906	-0.0436939\\
0.908	-0.0436941\\
0.91	-0.0436942\\
0.912	-0.0436943\\
0.914	-0.0436945\\
0.916	-0.0436946\\
0.918	-0.0436947\\
0.92	-0.0436948\\
0.922	-0.0436949\\
0.924	-0.043695\\
0.926	-0.0436951\\
0.928	-0.0436952\\
0.93	-0.0436953\\
0.932	-0.0436954\\
0.934	-0.0436955\\
0.936	-0.0436955\\
0.938	-0.0436956\\
0.94	-0.0436957\\
0.942	-0.0436957\\
0.944	-0.0436958\\
0.946	-0.0436958\\
0.948	-0.0436958\\
0.95	-0.0436959\\
0.952	-0.0436959\\
0.954	-0.0436959\\
0.956	-0.043696\\
0.958	-0.043696\\
0.96	-0.043696\\
0.962	-0.043696\\
0.964	-0.043696\\
0.966	-0.043696\\
0.968	-0.043696\\
0.97	-0.043696\\
0.972	-0.043696\\
0.974	-0.043696\\
0.976	-0.043696\\
0.978	-0.043696\\
0.98	-0.043696\\
0.982	-0.043696\\
0.984	-0.043696\\
0.986	-0.043696\\
0.988	-0.043696\\
0.99	-0.0436959\\
0.992	-0.0436959\\
0.994	-0.0436959\\
0.996	-0.0436959\\
0.998	-0.0436959\\
1	-0.0436959\\
1.002	-0.0436958\\
1.004	-0.0436958\\
1.006	-0.0436958\\
1.008	-0.0436958\\
1.01	-0.0436958\\
1.012	-0.0436957\\
1.014	-0.0436957\\
1.016	-0.0436957\\
1.018	-0.0436957\\
1.02	-0.0436957\\
1.022	-0.0436957\\
1.024	-0.0436956\\
1.026	-0.0436956\\
1.028	-0.0436956\\
1.03	-0.0436956\\
1.032	-0.0436956\\
1.034	-0.0436956\\
1.036	-0.0436956\\
1.038	-0.0436956\\
1.04	-0.0436956\\
1.042	-0.0436956\\
1.044	-0.0436956\\
1.046	-0.0436956\\
1.048	-0.0436956\\
1.05	-0.0436956\\
1.052	-0.0436956\\
1.054	-0.0436956\\
1.056	-0.0436956\\
1.058	-0.0436956\\
1.06	-0.0436956\\
1.062	-0.0436956\\
1.064	-0.0436956\\
1.066	-0.0436956\\
1.068	-0.0436956\\
1.07	-0.0436956\\
1.072	-0.0436957\\
1.074	-0.0436957\\
1.076	-0.0436957\\
1.078	-0.0436957\\
1.08	-0.0436957\\
1.082	-0.0436957\\
1.084	-0.0436958\\
1.086	-0.0436958\\
1.088	-0.0436958\\
1.09	-0.0436958\\
1.092	-0.0436959\\
1.094	-0.0436959\\
1.096	-0.0436959\\
1.098	-0.043696\\
1.1	-0.043696\\
1.102	-0.043696\\
1.104	-0.043696\\
1.106	-0.0436961\\
1.108	-0.0436961\\
1.11	-0.0436961\\
1.112	-0.0436962\\
1.114	-0.0436962\\
1.116	-0.0436962\\
1.118	-0.0436963\\
1.12	-0.0436963\\
1.122	-0.0436963\\
1.124	-0.0436964\\
1.126	-0.0436964\\
1.128	-0.0436965\\
1.13	-0.0436965\\
1.132	-0.0436965\\
1.134	-0.0436966\\
1.136	-0.0436966\\
1.138	-0.0436966\\
1.14	-0.0436967\\
1.142	-0.0436967\\
1.144	-0.0436967\\
1.146	-0.0436968\\
1.148	-0.0436968\\
1.15	-0.0436968\\
1.152	-0.0436969\\
1.154	-0.0436969\\
1.156	-0.0436969\\
1.158	-0.043697\\
1.16	-0.043697\\
1.162	-0.043697\\
1.164	-0.0436971\\
1.166	-0.0436971\\
1.168	-0.0436971\\
1.17	-0.0436972\\
1.172	-0.0436972\\
1.174	-0.0436972\\
1.176	-0.0436972\\
1.178	-0.0436973\\
1.18	-0.0436973\\
1.182	-0.0436973\\
1.184	-0.0436974\\
1.186	-0.0436974\\
1.188	-0.0436974\\
1.19	-0.0436974\\
1.192	-0.0436974\\
1.194	-0.0436975\\
1.196	-0.0436975\\
1.198	-0.0436975\\
1.2	-0.0436975\\
1.202	-0.0436975\\
1.204	-0.0436976\\
1.206	-0.0436976\\
1.208	-0.0436976\\
1.21	-0.0436976\\
1.212	-0.0436976\\
1.214	-0.0436976\\
1.216	-0.0436977\\
1.218	-0.0436977\\
1.22	-0.0436977\\
1.222	-0.0436977\\
1.224	-0.0436977\\
1.226	-0.0436977\\
1.228	-0.0436977\\
1.23	-0.0436977\\
1.232	-0.0436977\\
1.234	-0.0436977\\
1.236	-0.0436978\\
1.238	-0.0436978\\
1.24	-0.0436978\\
1.242	-0.0436978\\
1.244	-0.0436978\\
1.246	-0.0436978\\
1.248	-0.0436978\\
1.25	-0.0436978\\
1.252	-0.0436978\\
1.254	-0.0436978\\
1.256	-0.0436978\\
1.258	-0.0436978\\
1.26	-0.0436978\\
1.262	-0.0436978\\
1.264	-0.0436978\\
1.266	-0.0436978\\
1.268	-0.0436978\\
1.27	-0.0436978\\
1.272	-0.0436978\\
1.274	-0.0436978\\
1.276	-0.0436978\\
1.278	-0.0436978\\
1.28	-0.0436978\\
1.282	-0.0436978\\
1.284	-0.0436978\\
1.286	-0.0436978\\
1.288	-0.0436978\\
1.29	-0.0436978\\
1.292	-0.0436978\\
1.294	-0.0436978\\
1.296	-0.0436978\\
1.298	-0.0436978\\
1.3	-0.0436978\\
1.302	-0.0436978\\
1.304	-0.0436978\\
1.306	-0.0436978\\
1.308	-0.0436978\\
1.31	-0.0436978\\
1.312	-0.0436977\\
1.314	-0.0436977\\
1.316	-0.0436977\\
1.318	-0.0436977\\
1.32	-0.0436977\\
1.322	-0.0436977\\
1.324	-0.0436977\\
1.326	-0.0436977\\
1.328	-0.0436977\\
1.33	-0.0436977\\
1.332	-0.0436977\\
1.334	-0.0436977\\
1.336	-0.0436977\\
1.338	-0.0436977\\
1.34	-0.0436977\\
1.342	-0.0436977\\
1.344	-0.0436977\\
1.346	-0.0436977\\
1.348	-0.0436977\\
1.35	-0.0436977\\
1.352	-0.0436977\\
1.354	-0.0436977\\
1.356	-0.0436977\\
1.358	-0.0436977\\
1.36	-0.0436977\\
1.362	-0.0436977\\
1.364	-0.0436977\\
1.366	-0.0436977\\
1.368	-0.0436977\\
1.37	-0.0436977\\
1.372	-0.0436977\\
1.374	-0.0436977\\
1.376	-0.0436977\\
1.378	-0.0436977\\
1.38	-0.0436977\\
1.382	-0.0436977\\
1.384	-0.0436977\\
1.386	-0.0436977\\
1.388	-0.0436977\\
1.39	-0.0436977\\
1.392	-0.0436977\\
1.394	-0.0436977\\
1.396	-0.0436977\\
1.398	-0.0436977\\
1.4	-0.0436977\\
1.402	-0.0436977\\
1.404	-0.0436977\\
1.406	-0.0436977\\
1.408	-0.0436978\\
1.41	-0.0436978\\
1.412	-0.0436978\\
1.414	-0.0436978\\
1.416	-0.0436978\\
1.418	-0.0436978\\
1.42	-0.0436978\\
1.422	-0.0436978\\
1.424	-0.0436978\\
1.426	-0.0436978\\
1.428	-0.0436978\\
1.43	-0.0436978\\
1.432	-0.0436978\\
1.434	-0.0436978\\
1.436	-0.0436978\\
1.438	-0.0436978\\
1.44	-0.0436978\\
1.442	-0.0436978\\
1.444	-0.0436978\\
1.446	-0.0436978\\
1.448	-0.0436978\\
1.45	-0.0436978\\
1.452	-0.0436978\\
1.454	-0.0436978\\
1.456	-0.0436978\\
1.458	-0.0436978\\
1.46	-0.0436978\\
1.462	-0.0436978\\
1.464	-0.0436978\\
1.466	-0.0436978\\
1.468	-0.0436978\\
1.47	-0.0436978\\
1.472	-0.0436978\\
1.474	-0.0436979\\
1.476	-0.0436979\\
1.478	-0.0436979\\
1.48	-0.0436979\\
1.482	-0.0436979\\
1.484	-0.0436979\\
1.486	-0.0436979\\
1.488	-0.0436979\\
1.49	-0.0436979\\
1.492	-0.0436979\\
1.494	-0.0436979\\
1.496	-0.0436979\\
1.498	-0.0436979\\
1.5	-0.0436979\\
1.502	-0.0436979\\
1.504	-0.0436979\\
1.506	-0.0436979\\
1.508	-0.0436979\\
1.51	-0.0436979\\
1.512	-0.0436979\\
1.514	-0.0436979\\
1.516	-0.0436979\\
1.518	-0.0436979\\
1.52	-0.0436979\\
1.522	-0.0436979\\
1.524	-0.0436979\\
1.526	-0.0436979\\
1.528	-0.0436979\\
1.53	-0.0436979\\
1.532	-0.0436979\\
1.534	-0.0436979\\
1.536	-0.0436979\\
1.538	-0.0436979\\
1.54	-0.0436979\\
1.542	-0.0436979\\
1.544	-0.0436979\\
1.546	-0.0436979\\
1.548	-0.0436979\\
1.55	-0.0436979\\
1.552	-0.0436979\\
1.554	-0.0436979\\
1.556	-0.0436979\\
1.558	-0.0436979\\
1.56	-0.0436979\\
1.562	-0.0436979\\
1.564	-0.0436979\\
1.566	-0.0436979\\
1.568	-0.0436979\\
1.57	-0.0436979\\
1.572	-0.0436979\\
1.574	-0.0436979\\
1.576	-0.0436979\\
1.578	-0.0436979\\
1.58	-0.0436979\\
1.582	-0.0436979\\
1.584	-0.0436979\\
1.586	-0.0436979\\
1.588	-0.0436979\\
1.59	-0.0436979\\
1.592	-0.0436979\\
1.594	-0.0436979\\
1.596	-0.0436979\\
1.598	-0.0436979\\
1.6	-0.0436979\\
1.602	-0.0436979\\
1.604	-0.0436979\\
1.606	-0.0436979\\
1.608	-0.0436979\\
1.61	-0.0436979\\
1.612	-0.0436979\\
1.614	-0.0436979\\
1.616	-0.0436979\\
1.618	-0.0436979\\
1.62	-0.0436979\\
1.622	-0.0436979\\
1.624	-0.0436979\\
1.626	-0.0436979\\
1.628	-0.0436979\\
1.63	-0.0436979\\
1.632	-0.0436979\\
1.634	-0.0436979\\
1.636	-0.0436979\\
1.638	-0.0436979\\
1.64	-0.0436979\\
1.642	-0.0436979\\
1.644	-0.0436979\\
1.646	-0.0436979\\
1.648	-0.0436979\\
1.65	-0.0436979\\
1.652	-0.0436979\\
1.654	-0.0436979\\
1.656	-0.0436979\\
1.658	-0.0436979\\
1.66	-0.0436979\\
1.662	-0.0436979\\
1.664	-0.0436979\\
1.666	-0.0436979\\
1.668	-0.0436979\\
1.67	-0.0436979\\
1.672	-0.0436979\\
1.674	-0.0436979\\
1.676	-0.0436979\\
1.678	-0.0436979\\
1.68	-0.0436979\\
1.682	-0.0436979\\
1.684	-0.0436979\\
1.686	-0.0436979\\
1.688	-0.0436979\\
1.69	-0.0436979\\
1.692	-0.0436979\\
1.694	-0.0436979\\
1.696	-0.0436979\\
1.698	-0.0436979\\
1.7	-0.0436979\\
1.702	-0.0436979\\
1.704	-0.0436979\\
1.706	-0.0436979\\
1.708	-0.0436979\\
1.71	-0.0436979\\
1.712	-0.0436979\\
1.714	-0.0436979\\
1.716	-0.0436979\\
1.718	-0.0436979\\
1.72	-0.0436979\\
1.722	-0.0436979\\
1.724	-0.0436979\\
1.726	-0.0436979\\
1.728	-0.0436979\\
1.73	-0.0436979\\
1.732	-0.0436979\\
1.734	-0.0436979\\
1.736	-0.0436979\\
1.738	-0.0436979\\
1.74	-0.0436979\\
1.742	-0.0436979\\
1.744	-0.0436979\\
1.746	-0.0436979\\
1.748	-0.0436979\\
1.75	-0.0436979\\
1.752	-0.0436979\\
1.754	-0.0436979\\
1.756	-0.0436979\\
1.758	-0.0436979\\
1.76	-0.0436979\\
1.762	-0.0436979\\
1.764	-0.0436979\\
1.766	-0.0436979\\
1.768	-0.0436979\\
1.77	-0.0436979\\
1.772	-0.0436979\\
1.774	-0.0436979\\
1.776	-0.0436979\\
1.778	-0.0436979\\
1.78	-0.0436979\\
1.782	-0.0436979\\
1.784	-0.0436979\\
1.786	-0.0436979\\
1.788	-0.0436979\\
1.79	-0.0436979\\
1.792	-0.0436979\\
1.794	-0.0436979\\
1.796	-0.0436979\\
1.798	-0.0436979\\
1.8	-0.0436979\\
1.802	-0.0436979\\
1.804	-0.0436979\\
1.806	-0.0436979\\
1.808	-0.0436979\\
1.81	-0.0436979\\
1.812	-0.0436979\\
1.814	-0.0436979\\
1.816	-0.0436979\\
1.818	-0.0436979\\
1.82	-0.0436979\\
1.822	-0.0436979\\
1.824	-0.0436979\\
1.826	-0.0436979\\
1.828	-0.0436979\\
1.83	-0.0436979\\
1.832	-0.0436979\\
1.834	-0.0436979\\
1.836	-0.0436979\\
1.838	-0.0436979\\
1.84	-0.0436979\\
1.842	-0.0436979\\
1.844	-0.0436979\\
1.846	-0.0436979\\
1.848	-0.0436979\\
1.85	-0.0436979\\
1.852	-0.0436979\\
1.854	-0.0436979\\
1.856	-0.0436979\\
1.858	-0.0436979\\
1.86	-0.0436979\\
1.862	-0.0436979\\
1.864	-0.0436979\\
1.866	-0.0436979\\
1.868	-0.0436979\\
1.87	-0.0436979\\
1.872	-0.0436979\\
1.874	-0.0436979\\
1.876	-0.0436979\\
1.878	-0.0436979\\
1.88	-0.0436979\\
1.882	-0.0436979\\
1.884	-0.0436979\\
1.886	-0.0436979\\
1.888	-0.0436979\\
1.89	-0.0436979\\
1.892	-0.0436979\\
1.894	-0.0436979\\
1.896	-0.0436979\\
1.898	-0.0436979\\
1.9	-0.0436979\\
1.902	-0.0436979\\
1.904	-0.0436979\\
1.906	-0.0436979\\
1.908	-0.0436979\\
1.91	-0.0436979\\
1.912	-0.0436979\\
1.914	-0.0436979\\
1.916	-0.0436979\\
1.918	-0.0436979\\
1.92	-0.0436979\\
1.922	-0.0436979\\
1.924	-0.0436979\\
1.926	-0.0436979\\
1.928	-0.0436979\\
1.93	-0.0436979\\
1.932	-0.0436979\\
1.934	-0.0436979\\
1.936	-0.0436979\\
1.938	-0.0436979\\
1.94	-0.0436979\\
1.942	-0.0436979\\
1.944	-0.0436979\\
1.946	-0.0436979\\
1.948	-0.0436979\\
1.95	-0.0436979\\
1.952	-0.0436979\\
1.954	-0.0436979\\
1.956	-0.0436979\\
1.958	-0.0436979\\
1.96	-0.0436979\\
1.962	-0.0436979\\
1.964	-0.0436979\\
1.966	-0.0436979\\
1.968	-0.0436979\\
1.97	-0.0436979\\
1.972	-0.0436979\\
1.974	-0.0436979\\
1.976	-0.0436979\\
1.978	-0.0436979\\
1.98	-0.0436979\\
1.982	-0.0436979\\
1.984	-0.0436979\\
1.986	-0.0436979\\
1.988	-0.0436979\\
1.99	-0.0436979\\
1.992	-0.0436979\\
1.994	-0.0436979\\
1.996	-0.0436979\\
1.998	-0.0436979\\
2	-0.0436979\\
};
\addplot [color=mycolor10, line width=1.0pt, forget plot]
  table[row sep=crcr]{%
-0.024	0\\
-0.022	0\\
-0.02	0\\
-0.018	0\\
-0.016	0\\
-0.014	0\\
-0.012	0\\
-0.01	0\\
-0.008	0\\
-0.006	0\\
-0.004	0\\
-0.002	0\\
0	0\\
0.002	-0.00191221\\
0.004	-0.00365969\\
0.006	-0.00525793\\
0.008	-0.00672136\\
0.01	-0.00806343\\
0.012	-0.00929665\\
0.014	-0.0104326\\
0.016	-0.0114821\\
0.018	-0.012455\\
0.02	-0.0133604\\
0.022	-0.0142069\\
0.024	-0.015002\\
0.026	-0.0157528\\
0.028	-0.0164656\\
0.03	-0.017146\\
0.032	-0.0177992\\
0.034	-0.0184296\\
0.036	-0.0190411\\
0.038	-0.0196371\\
0.04	-0.0202206\\
0.042	-0.020794\\
0.044	-0.0213592\\
0.046	-0.0219178\\
0.048	-0.022471\\
0.05	-0.0230196\\
0.052	-0.0235644\\
0.054	-0.0241054\\
0.056	-0.0246429\\
0.058	-0.0251767\\
0.06	-0.0257065\\
0.062	-0.0262318\\
0.064	-0.0267523\\
0.066	-0.0272672\\
0.068	-0.0277761\\
0.07	-0.0282783\\
0.072	-0.0287731\\
0.074	-0.02926\\
0.076	-0.0297385\\
0.078	-0.030208\\
0.08	-0.0306682\\
0.082	-0.0311186\\
0.084	-0.0315592\\
0.086	-0.0319896\\
0.088	-0.0324097\\
0.09	-0.0328196\\
0.092	-0.0332192\\
0.094	-0.0336086\\
0.096	-0.0339881\\
0.098	-0.0343579\\
0.1	-0.034718\\
0.102	-0.0350689\\
0.104	-0.0354108\\
0.106	-0.035744\\
0.108	-0.0360688\\
0.11	-0.0363855\\
0.112	-0.0366943\\
0.114	-0.0369955\\
0.116	-0.0372892\\
0.118	-0.0375757\\
0.12	-0.0378551\\
0.122	-0.0381274\\
0.124	-0.0383928\\
0.126	-0.0386513\\
0.128	-0.0389029\\
0.13	-0.0391474\\
0.132	-0.0393849\\
0.134	-0.0396152\\
0.136	-0.0398382\\
0.138	-0.0400538\\
0.14	-0.0402618\\
0.142	-0.040462\\
0.144	-0.0406544\\
0.146	-0.0408387\\
0.148	-0.0410148\\
0.15	-0.0411825\\
0.152	-0.0413419\\
0.154	-0.0414927\\
0.156	-0.0416351\\
0.158	-0.0417688\\
0.16	-0.041894\\
0.162	-0.0420107\\
0.164	-0.042119\\
0.166	-0.0422189\\
0.168	-0.0423107\\
0.17	-0.0423945\\
0.172	-0.0424706\\
0.174	-0.0425391\\
0.176	-0.0426004\\
0.178	-0.0426547\\
0.18	-0.0427024\\
0.182	-0.0427438\\
0.184	-0.0427792\\
0.186	-0.042809\\
0.188	-0.0428336\\
0.19	-0.0428534\\
0.192	-0.0428686\\
0.194	-0.0428798\\
0.196	-0.0428872\\
0.198	-0.0428912\\
0.2	-0.0428922\\
0.202	-0.0428905\\
0.204	-0.0428865\\
0.206	-0.0428804\\
0.208	-0.0428725\\
0.21	-0.0428631\\
0.212	-0.0428525\\
0.214	-0.042841\\
0.216	-0.0428286\\
0.218	-0.0428157\\
0.22	-0.0428024\\
0.222	-0.0427889\\
0.224	-0.0427754\\
0.226	-0.0427619\\
0.228	-0.0427487\\
0.23	-0.0427357\\
0.232	-0.0427231\\
0.234	-0.042711\\
0.236	-0.0426995\\
0.238	-0.0426885\\
0.24	-0.0426782\\
0.242	-0.0426685\\
0.244	-0.0426596\\
0.246	-0.0426514\\
0.248	-0.042644\\
0.25	-0.0426373\\
0.252	-0.0426313\\
0.254	-0.0426262\\
0.256	-0.0426219\\
0.258	-0.0426183\\
0.26	-0.0426155\\
0.262	-0.0426134\\
0.264	-0.0426121\\
0.266	-0.0426116\\
0.268	-0.0426118\\
0.27	-0.0426127\\
0.272	-0.0426143\\
0.274	-0.0426166\\
0.276	-0.0426196\\
0.278	-0.0426233\\
0.28	-0.0426276\\
0.282	-0.0426325\\
0.284	-0.042638\\
0.286	-0.0426442\\
0.288	-0.0426508\\
0.29	-0.0426581\\
0.292	-0.0426658\\
0.294	-0.042674\\
0.296	-0.0426827\\
0.298	-0.0426919\\
0.3	-0.0427014\\
0.302	-0.0427114\\
0.304	-0.0427217\\
0.306	-0.0427324\\
0.308	-0.0427434\\
0.31	-0.0427547\\
0.312	-0.0427662\\
0.314	-0.042778\\
0.316	-0.04279\\
0.318	-0.0428022\\
0.32	-0.0428145\\
0.322	-0.042827\\
0.324	-0.0428395\\
0.326	-0.0428522\\
0.328	-0.0428649\\
0.33	-0.0428777\\
0.332	-0.0428905\\
0.334	-0.0429034\\
0.336	-0.0429162\\
0.338	-0.0429289\\
0.34	-0.0429417\\
0.342	-0.0429544\\
0.344	-0.042967\\
0.346	-0.0429795\\
0.348	-0.042992\\
0.35	-0.0430043\\
0.352	-0.0430166\\
0.354	-0.0430287\\
0.356	-0.0430407\\
0.358	-0.0430526\\
0.36	-0.0430643\\
0.362	-0.0430759\\
0.364	-0.0430874\\
0.366	-0.0430987\\
0.368	-0.0431099\\
0.37	-0.0431209\\
0.372	-0.0431318\\
0.374	-0.0431425\\
0.376	-0.0431531\\
0.378	-0.0431636\\
0.38	-0.0431739\\
0.382	-0.043184\\
0.384	-0.043194\\
0.386	-0.0432039\\
0.388	-0.0432137\\
0.39	-0.0432233\\
0.392	-0.0432327\\
0.394	-0.043242\\
0.396	-0.0432512\\
0.398	-0.0432603\\
0.4	-0.0432692\\
0.402	-0.0432779\\
0.404	-0.0432866\\
0.406	-0.0432951\\
0.408	-0.0433035\\
0.41	-0.0433118\\
0.412	-0.0433199\\
0.414	-0.0433279\\
0.416	-0.0433357\\
0.418	-0.0433435\\
0.42	-0.0433511\\
0.422	-0.0433586\\
0.424	-0.0433659\\
0.426	-0.0433732\\
0.428	-0.0433803\\
0.43	-0.0433872\\
0.432	-0.0433941\\
0.434	-0.0434008\\
0.436	-0.0434074\\
0.438	-0.0434138\\
0.44	-0.0434201\\
0.442	-0.0434263\\
0.444	-0.0434324\\
0.446	-0.0434384\\
0.448	-0.0434442\\
0.45	-0.0434499\\
0.452	-0.0434555\\
0.454	-0.0434609\\
0.456	-0.0434663\\
0.458	-0.0434715\\
0.46	-0.0434766\\
0.462	-0.0434815\\
0.464	-0.0434864\\
0.466	-0.0434911\\
0.468	-0.0434958\\
0.47	-0.0435003\\
0.472	-0.0435047\\
0.474	-0.043509\\
0.476	-0.0435131\\
0.478	-0.0435172\\
0.48	-0.0435212\\
0.482	-0.0435251\\
0.484	-0.0435288\\
0.486	-0.0435325\\
0.488	-0.0435361\\
0.49	-0.0435396\\
0.492	-0.043543\\
0.494	-0.0435463\\
0.496	-0.0435495\\
0.498	-0.0435527\\
0.5	-0.0435557\\
0.502	-0.0435587\\
0.504	-0.0435616\\
0.506	-0.0435644\\
0.508	-0.0435671\\
0.51	-0.0435698\\
0.512	-0.0435724\\
0.514	-0.0435749\\
0.516	-0.0435774\\
0.518	-0.0435798\\
0.52	-0.0435821\\
0.522	-0.0435844\\
0.524	-0.0435866\\
0.526	-0.0435888\\
0.528	-0.0435909\\
0.53	-0.043593\\
0.532	-0.043595\\
0.534	-0.0435969\\
0.536	-0.0435988\\
0.538	-0.0436007\\
0.54	-0.0436025\\
0.542	-0.0436042\\
0.544	-0.0436059\\
0.546	-0.0436076\\
0.548	-0.0436093\\
0.55	-0.0436109\\
0.552	-0.0436124\\
0.554	-0.0436139\\
0.556	-0.0436154\\
0.558	-0.0436169\\
0.56	-0.0436183\\
0.562	-0.0436197\\
0.564	-0.043621\\
0.566	-0.0436223\\
0.568	-0.0436236\\
0.57	-0.0436249\\
0.572	-0.0436261\\
0.574	-0.0436273\\
0.576	-0.0436285\\
0.578	-0.0436296\\
0.58	-0.0436308\\
0.582	-0.0436318\\
0.584	-0.0436329\\
0.586	-0.043634\\
0.588	-0.043635\\
0.59	-0.043636\\
0.592	-0.0436369\\
0.594	-0.0436379\\
0.596	-0.0436388\\
0.598	-0.0436397\\
0.6	-0.0436406\\
0.602	-0.0436414\\
0.604	-0.0436423\\
0.606	-0.0436431\\
0.608	-0.0436439\\
0.61	-0.0436447\\
0.612	-0.0436454\\
0.614	-0.0436462\\
0.616	-0.0436469\\
0.618	-0.0436476\\
0.62	-0.0436483\\
0.622	-0.043649\\
0.624	-0.0436496\\
0.626	-0.0436502\\
0.628	-0.0436509\\
0.63	-0.0436515\\
0.632	-0.043652\\
0.634	-0.0436526\\
0.636	-0.0436532\\
0.638	-0.0436537\\
0.64	-0.0436542\\
0.642	-0.0436548\\
0.644	-0.0436553\\
0.646	-0.0436557\\
0.648	-0.0436562\\
0.65	-0.0436567\\
0.652	-0.0436571\\
0.654	-0.0436576\\
0.656	-0.043658\\
0.658	-0.0436584\\
0.66	-0.0436588\\
0.662	-0.0436592\\
0.664	-0.0436596\\
0.666	-0.04366\\
0.668	-0.0436603\\
0.67	-0.0436607\\
0.672	-0.043661\\
0.674	-0.0436614\\
0.676	-0.0436617\\
0.678	-0.043662\\
0.68	-0.0436624\\
0.682	-0.0436627\\
0.684	-0.043663\\
0.686	-0.0436633\\
0.688	-0.0436636\\
0.69	-0.0436639\\
0.692	-0.0436642\\
0.694	-0.0436644\\
0.696	-0.0436647\\
0.698	-0.043665\\
0.7	-0.0436653\\
0.702	-0.0436655\\
0.704	-0.0436658\\
0.706	-0.0436661\\
0.708	-0.0436663\\
0.71	-0.0436666\\
0.712	-0.0436668\\
0.714	-0.0436671\\
0.716	-0.0436673\\
0.718	-0.0436676\\
0.72	-0.0436678\\
0.722	-0.0436681\\
0.724	-0.0436683\\
0.726	-0.0436686\\
0.728	-0.0436688\\
0.73	-0.0436691\\
0.732	-0.0436693\\
0.734	-0.0436695\\
0.736	-0.0436698\\
0.738	-0.04367\\
0.74	-0.0436703\\
0.742	-0.0436705\\
0.744	-0.0436708\\
0.746	-0.043671\\
0.748	-0.0436713\\
0.75	-0.0436716\\
0.752	-0.0436718\\
0.754	-0.0436721\\
0.756	-0.0436723\\
0.758	-0.0436726\\
0.76	-0.0436729\\
0.762	-0.0436731\\
0.764	-0.0436734\\
0.766	-0.0436737\\
0.768	-0.043674\\
0.77	-0.0436742\\
0.772	-0.0436745\\
0.774	-0.0436748\\
0.776	-0.0436751\\
0.778	-0.0436754\\
0.78	-0.0436757\\
0.782	-0.0436759\\
0.784	-0.0436762\\
0.786	-0.0436765\\
0.788	-0.0436768\\
0.79	-0.0436771\\
0.792	-0.0436774\\
0.794	-0.0436777\\
0.796	-0.0436781\\
0.798	-0.0436784\\
0.8	-0.0436787\\
0.802	-0.043679\\
0.804	-0.0436793\\
0.806	-0.0436796\\
0.808	-0.0436799\\
0.81	-0.0436803\\
0.812	-0.0436806\\
0.814	-0.0436809\\
0.816	-0.0436812\\
0.818	-0.0436815\\
0.82	-0.0436819\\
0.822	-0.0436822\\
0.824	-0.0436825\\
0.826	-0.0436828\\
0.828	-0.0436831\\
0.83	-0.0436835\\
0.832	-0.0436838\\
0.834	-0.0436841\\
0.836	-0.0436844\\
0.838	-0.0436847\\
0.84	-0.0436851\\
0.842	-0.0436854\\
0.844	-0.0436857\\
0.846	-0.043686\\
0.848	-0.0436863\\
0.85	-0.0436866\\
0.852	-0.0436869\\
0.854	-0.0436872\\
0.856	-0.0436875\\
0.858	-0.0436878\\
0.86	-0.0436881\\
0.862	-0.0436883\\
0.864	-0.0436886\\
0.866	-0.0436889\\
0.868	-0.0436892\\
0.87	-0.0436894\\
0.872	-0.0436897\\
0.874	-0.0436899\\
0.876	-0.0436902\\
0.878	-0.0436904\\
0.88	-0.0436907\\
0.882	-0.0436909\\
0.884	-0.0436911\\
0.886	-0.0436913\\
0.888	-0.0436916\\
0.89	-0.0436918\\
0.892	-0.043692\\
0.894	-0.0436922\\
0.896	-0.0436924\\
0.898	-0.0436925\\
0.9	-0.0436927\\
0.902	-0.0436929\\
0.904	-0.0436931\\
0.906	-0.0436932\\
0.908	-0.0436934\\
0.91	-0.0436935\\
0.912	-0.0436937\\
0.914	-0.0436938\\
0.916	-0.0436939\\
0.918	-0.0436941\\
0.92	-0.0436942\\
0.922	-0.0436943\\
0.924	-0.0436944\\
0.926	-0.0436945\\
0.928	-0.0436946\\
0.93	-0.0436947\\
0.932	-0.0436948\\
0.934	-0.0436949\\
0.936	-0.0436949\\
0.938	-0.043695\\
0.94	-0.0436951\\
0.942	-0.0436951\\
0.944	-0.0436952\\
0.946	-0.0436952\\
0.948	-0.0436953\\
0.95	-0.0436953\\
0.952	-0.0436954\\
0.954	-0.0436954\\
0.956	-0.0436954\\
0.958	-0.0436955\\
0.96	-0.0436955\\
0.962	-0.0436955\\
0.964	-0.0436955\\
0.966	-0.0436955\\
0.968	-0.0436955\\
0.97	-0.0436956\\
0.972	-0.0436956\\
0.974	-0.0436956\\
0.976	-0.0436956\\
0.978	-0.0436956\\
0.98	-0.0436956\\
0.982	-0.0436956\\
0.984	-0.0436956\\
0.986	-0.0436955\\
0.988	-0.0436955\\
0.99	-0.0436955\\
0.992	-0.0436955\\
0.994	-0.0436955\\
0.996	-0.0436955\\
0.998	-0.0436955\\
1	-0.0436955\\
1.002	-0.0436955\\
1.004	-0.0436954\\
1.006	-0.0436954\\
1.008	-0.0436954\\
1.01	-0.0436954\\
1.012	-0.0436954\\
1.014	-0.0436954\\
1.016	-0.0436954\\
1.018	-0.0436954\\
1.02	-0.0436953\\
1.022	-0.0436953\\
1.024	-0.0436953\\
1.026	-0.0436953\\
1.028	-0.0436953\\
1.03	-0.0436953\\
1.032	-0.0436953\\
1.034	-0.0436953\\
1.036	-0.0436953\\
1.038	-0.0436953\\
1.04	-0.0436953\\
1.042	-0.0436953\\
1.044	-0.0436953\\
1.046	-0.0436953\\
1.048	-0.0436953\\
1.05	-0.0436953\\
1.052	-0.0436953\\
1.054	-0.0436953\\
1.056	-0.0436953\\
1.058	-0.0436953\\
1.06	-0.0436953\\
1.062	-0.0436954\\
1.064	-0.0436954\\
1.066	-0.0436954\\
1.068	-0.0436954\\
1.07	-0.0436954\\
1.072	-0.0436954\\
1.074	-0.0436955\\
1.076	-0.0436955\\
1.078	-0.0436955\\
1.08	-0.0436955\\
1.082	-0.0436955\\
1.084	-0.0436956\\
1.086	-0.0436956\\
1.088	-0.0436956\\
1.09	-0.0436956\\
1.092	-0.0436957\\
1.094	-0.0436957\\
1.096	-0.0436957\\
1.098	-0.0436958\\
1.1	-0.0436958\\
1.102	-0.0436958\\
1.104	-0.0436959\\
1.106	-0.0436959\\
1.108	-0.0436959\\
1.11	-0.043696\\
1.112	-0.043696\\
1.114	-0.043696\\
1.116	-0.0436961\\
1.118	-0.0436961\\
1.12	-0.0436961\\
1.122	-0.0436962\\
1.124	-0.0436962\\
1.126	-0.0436963\\
1.128	-0.0436963\\
1.13	-0.0436963\\
1.132	-0.0436964\\
1.134	-0.0436964\\
1.136	-0.0436964\\
1.138	-0.0436965\\
1.14	-0.0436965\\
1.142	-0.0436965\\
1.144	-0.0436966\\
1.146	-0.0436966\\
1.148	-0.0436967\\
1.15	-0.0436967\\
1.152	-0.0436967\\
1.154	-0.0436968\\
1.156	-0.0436968\\
1.158	-0.0436968\\
1.16	-0.0436969\\
1.162	-0.0436969\\
1.164	-0.0436969\\
1.166	-0.043697\\
1.168	-0.043697\\
1.17	-0.043697\\
1.172	-0.0436971\\
1.174	-0.0436971\\
1.176	-0.0436971\\
1.178	-0.0436971\\
1.18	-0.0436972\\
1.182	-0.0436972\\
1.184	-0.0436972\\
1.186	-0.0436972\\
1.188	-0.0436973\\
1.19	-0.0436973\\
1.192	-0.0436973\\
1.194	-0.0436973\\
1.196	-0.0436974\\
1.198	-0.0436974\\
1.2	-0.0436974\\
1.202	-0.0436974\\
1.204	-0.0436974\\
1.206	-0.0436975\\
1.208	-0.0436975\\
1.21	-0.0436975\\
1.212	-0.0436975\\
1.214	-0.0436975\\
1.216	-0.0436975\\
1.218	-0.0436976\\
1.22	-0.0436976\\
1.222	-0.0436976\\
1.224	-0.0436976\\
1.226	-0.0436976\\
1.228	-0.0436976\\
1.23	-0.0436976\\
1.232	-0.0436976\\
1.234	-0.0436976\\
1.236	-0.0436977\\
1.238	-0.0436977\\
1.24	-0.0436977\\
1.242	-0.0436977\\
1.244	-0.0436977\\
1.246	-0.0436977\\
1.248	-0.0436977\\
1.25	-0.0436977\\
1.252	-0.0436977\\
1.254	-0.0436977\\
1.256	-0.0436977\\
1.258	-0.0436977\\
1.26	-0.0436977\\
1.262	-0.0436977\\
1.264	-0.0436977\\
1.266	-0.0436977\\
1.268	-0.0436977\\
1.27	-0.0436977\\
1.272	-0.0436977\\
1.274	-0.0436977\\
1.276	-0.0436977\\
1.278	-0.0436977\\
1.28	-0.0436977\\
1.282	-0.0436977\\
1.284	-0.0436977\\
1.286	-0.0436977\\
1.288	-0.0436977\\
1.29	-0.0436977\\
1.292	-0.0436977\\
1.294	-0.0436977\\
1.296	-0.0436977\\
1.298	-0.0436977\\
1.3	-0.0436977\\
1.302	-0.0436977\\
1.304	-0.0436977\\
1.306	-0.0436977\\
1.308	-0.0436977\\
1.31	-0.0436977\\
1.312	-0.0436977\\
1.314	-0.0436977\\
1.316	-0.0436977\\
1.318	-0.0436977\\
1.32	-0.0436977\\
1.322	-0.0436977\\
1.324	-0.0436977\\
1.326	-0.0436977\\
1.328	-0.0436977\\
1.33	-0.0436977\\
1.332	-0.0436977\\
1.334	-0.0436977\\
1.336	-0.0436977\\
1.338	-0.0436977\\
1.34	-0.0436977\\
1.342	-0.0436977\\
1.344	-0.0436977\\
1.346	-0.0436977\\
1.348	-0.0436977\\
1.35	-0.0436977\\
1.352	-0.0436977\\
1.354	-0.0436977\\
1.356	-0.0436977\\
1.358	-0.0436977\\
1.36	-0.0436977\\
1.362	-0.0436977\\
1.364	-0.0436977\\
1.366	-0.0436977\\
1.368	-0.0436977\\
1.37	-0.0436977\\
1.372	-0.0436977\\
1.374	-0.0436977\\
1.376	-0.0436977\\
1.378	-0.0436977\\
1.38	-0.0436977\\
1.382	-0.0436977\\
1.384	-0.0436977\\
1.386	-0.0436977\\
1.388	-0.0436977\\
1.39	-0.0436977\\
1.392	-0.0436977\\
1.394	-0.0436977\\
1.396	-0.0436977\\
1.398	-0.0436977\\
1.4	-0.0436977\\
1.402	-0.0436977\\
1.404	-0.0436977\\
1.406	-0.0436977\\
1.408	-0.0436977\\
1.41	-0.0436977\\
1.412	-0.0436977\\
1.414	-0.0436977\\
1.416	-0.0436977\\
1.418	-0.0436977\\
1.42	-0.0436977\\
1.422	-0.0436978\\
1.424	-0.0436978\\
1.426	-0.0436978\\
1.428	-0.0436978\\
1.43	-0.0436978\\
1.432	-0.0436978\\
1.434	-0.0436978\\
1.436	-0.0436978\\
1.438	-0.0436978\\
1.44	-0.0436978\\
1.442	-0.0436978\\
1.444	-0.0436978\\
1.446	-0.0436978\\
1.448	-0.0436978\\
1.45	-0.0436978\\
1.452	-0.0436978\\
1.454	-0.0436978\\
1.456	-0.0436978\\
1.458	-0.0436978\\
1.46	-0.0436978\\
1.462	-0.0436978\\
1.464	-0.0436978\\
1.466	-0.0436978\\
1.468	-0.0436978\\
1.47	-0.0436978\\
1.472	-0.0436978\\
1.474	-0.0436978\\
1.476	-0.0436978\\
1.478	-0.0436978\\
1.48	-0.0436978\\
1.482	-0.0436978\\
1.484	-0.0436978\\
1.486	-0.0436978\\
1.488	-0.0436978\\
1.49	-0.0436979\\
1.492	-0.0436979\\
1.494	-0.0436979\\
1.496	-0.0436979\\
1.498	-0.0436979\\
1.5	-0.0436979\\
1.502	-0.0436979\\
1.504	-0.0436979\\
1.506	-0.0436979\\
1.508	-0.0436979\\
1.51	-0.0436979\\
1.512	-0.0436979\\
1.514	-0.0436979\\
1.516	-0.0436979\\
1.518	-0.0436979\\
1.52	-0.0436979\\
1.522	-0.0436979\\
1.524	-0.0436979\\
1.526	-0.0436979\\
1.528	-0.0436979\\
1.53	-0.0436979\\
1.532	-0.0436979\\
1.534	-0.0436979\\
1.536	-0.0436979\\
1.538	-0.0436979\\
1.54	-0.0436979\\
1.542	-0.0436979\\
1.544	-0.0436979\\
1.546	-0.0436979\\
1.548	-0.0436979\\
1.55	-0.0436979\\
1.552	-0.0436979\\
1.554	-0.0436979\\
1.556	-0.0436979\\
1.558	-0.0436979\\
1.56	-0.0436979\\
1.562	-0.0436979\\
1.564	-0.0436979\\
1.566	-0.0436979\\
1.568	-0.0436979\\
1.57	-0.0436979\\
1.572	-0.0436979\\
1.574	-0.0436979\\
1.576	-0.0436979\\
1.578	-0.0436979\\
1.58	-0.0436979\\
1.582	-0.0436979\\
1.584	-0.0436979\\
1.586	-0.0436979\\
1.588	-0.0436979\\
1.59	-0.0436979\\
1.592	-0.0436979\\
1.594	-0.0436979\\
1.596	-0.0436979\\
1.598	-0.0436979\\
1.6	-0.0436979\\
1.602	-0.0436979\\
1.604	-0.0436979\\
1.606	-0.0436979\\
1.608	-0.0436979\\
1.61	-0.0436979\\
1.612	-0.0436979\\
1.614	-0.0436979\\
1.616	-0.0436979\\
1.618	-0.0436979\\
1.62	-0.0436979\\
1.622	-0.0436979\\
1.624	-0.0436979\\
1.626	-0.0436979\\
1.628	-0.0436979\\
1.63	-0.0436979\\
1.632	-0.0436979\\
1.634	-0.0436979\\
1.636	-0.0436979\\
1.638	-0.0436979\\
1.64	-0.0436979\\
1.642	-0.0436979\\
1.644	-0.0436979\\
1.646	-0.0436979\\
1.648	-0.0436979\\
1.65	-0.0436979\\
1.652	-0.0436979\\
1.654	-0.0436979\\
1.656	-0.0436979\\
1.658	-0.0436979\\
1.66	-0.0436979\\
1.662	-0.0436979\\
1.664	-0.0436979\\
1.666	-0.0436979\\
1.668	-0.0436979\\
1.67	-0.0436979\\
1.672	-0.0436979\\
1.674	-0.0436979\\
1.676	-0.0436979\\
1.678	-0.0436979\\
1.68	-0.0436979\\
1.682	-0.0436979\\
1.684	-0.0436979\\
1.686	-0.0436979\\
1.688	-0.0436979\\
1.69	-0.0436979\\
1.692	-0.0436979\\
1.694	-0.0436979\\
1.696	-0.0436979\\
1.698	-0.0436979\\
1.7	-0.0436979\\
1.702	-0.0436979\\
1.704	-0.0436979\\
1.706	-0.0436979\\
1.708	-0.0436979\\
1.71	-0.0436979\\
1.712	-0.0436979\\
1.714	-0.0436979\\
1.716	-0.0436979\\
1.718	-0.0436979\\
1.72	-0.0436979\\
1.722	-0.0436979\\
1.724	-0.0436979\\
1.726	-0.0436979\\
1.728	-0.0436979\\
1.73	-0.0436979\\
1.732	-0.0436979\\
1.734	-0.0436979\\
1.736	-0.0436979\\
1.738	-0.0436979\\
1.74	-0.0436979\\
1.742	-0.0436979\\
1.744	-0.0436979\\
1.746	-0.0436979\\
1.748	-0.0436979\\
1.75	-0.0436979\\
1.752	-0.0436979\\
1.754	-0.0436979\\
1.756	-0.0436979\\
1.758	-0.0436979\\
1.76	-0.0436979\\
1.762	-0.0436979\\
1.764	-0.0436979\\
1.766	-0.0436979\\
1.768	-0.0436979\\
1.77	-0.0436979\\
1.772	-0.0436979\\
1.774	-0.0436979\\
1.776	-0.0436979\\
1.778	-0.0436979\\
1.78	-0.0436979\\
1.782	-0.0436979\\
1.784	-0.0436979\\
1.786	-0.0436979\\
1.788	-0.0436979\\
1.79	-0.0436979\\
1.792	-0.0436979\\
1.794	-0.0436979\\
1.796	-0.0436979\\
1.798	-0.0436979\\
1.8	-0.0436979\\
1.802	-0.0436979\\
1.804	-0.0436979\\
1.806	-0.0436979\\
1.808	-0.0436979\\
1.81	-0.0436979\\
1.812	-0.0436979\\
1.814	-0.0436979\\
1.816	-0.0436979\\
1.818	-0.0436979\\
1.82	-0.0436979\\
1.822	-0.0436979\\
1.824	-0.0436979\\
1.826	-0.0436979\\
1.828	-0.0436979\\
1.83	-0.0436979\\
1.832	-0.0436979\\
1.834	-0.0436979\\
1.836	-0.0436979\\
1.838	-0.0436979\\
1.84	-0.0436979\\
1.842	-0.0436979\\
1.844	-0.0436979\\
1.846	-0.0436979\\
1.848	-0.0436979\\
1.85	-0.0436979\\
1.852	-0.0436979\\
1.854	-0.0436979\\
1.856	-0.0436979\\
1.858	-0.0436979\\
1.86	-0.0436979\\
1.862	-0.0436979\\
1.864	-0.0436979\\
1.866	-0.0436979\\
1.868	-0.0436979\\
1.87	-0.0436979\\
1.872	-0.0436979\\
1.874	-0.0436979\\
1.876	-0.0436979\\
1.878	-0.0436979\\
1.88	-0.0436979\\
1.882	-0.0436979\\
1.884	-0.0436979\\
1.886	-0.0436979\\
1.888	-0.0436979\\
1.89	-0.0436979\\
1.892	-0.0436979\\
1.894	-0.0436979\\
1.896	-0.0436979\\
1.898	-0.0436979\\
1.9	-0.0436979\\
1.902	-0.0436979\\
1.904	-0.0436979\\
1.906	-0.0436979\\
1.908	-0.0436979\\
1.91	-0.0436979\\
1.912	-0.0436979\\
1.914	-0.0436979\\
1.916	-0.0436979\\
1.918	-0.0436979\\
1.92	-0.0436979\\
1.922	-0.0436979\\
1.924	-0.0436979\\
1.926	-0.0436979\\
1.928	-0.0436979\\
1.93	-0.0436979\\
1.932	-0.0436979\\
1.934	-0.0436979\\
1.936	-0.0436979\\
1.938	-0.0436979\\
1.94	-0.0436979\\
1.942	-0.0436979\\
1.944	-0.0436979\\
1.946	-0.0436979\\
1.948	-0.0436979\\
1.95	-0.0436979\\
1.952	-0.0436979\\
1.954	-0.0436979\\
1.956	-0.0436979\\
1.958	-0.0436979\\
1.96	-0.0436979\\
1.962	-0.0436979\\
1.964	-0.0436979\\
1.966	-0.0436979\\
1.968	-0.0436979\\
1.97	-0.0436979\\
1.972	-0.0436979\\
1.974	-0.0436979\\
1.976	-0.0436979\\
1.978	-0.0436979\\
1.98	-0.0436979\\
1.982	-0.0436979\\
1.984	-0.0436979\\
1.986	-0.0436979\\
1.988	-0.0436979\\
1.99	-0.0436979\\
1.992	-0.0436979\\
1.994	-0.0436979\\
1.996	-0.0436979\\
1.998	-0.0436979\\
2	-0.0436979\\
};
\addplot [color=mycolor11, line width=1.0pt, forget plot]
  table[row sep=crcr]{%
-0.024	0\\
-0.022	0\\
-0.02	0\\
-0.018	0\\
-0.016	0\\
-0.014	0\\
-0.012	0\\
-0.01	0\\
-0.008	0\\
-0.006	0\\
-0.004	0\\
-0.002	0\\
0	0\\
0.002	-0.0019124\\
0.004	-0.00366117\\
0.006	-0.00526261\\
0.008	-0.00673176\\
0.01	-0.00808245\\
0.012	-0.00932738\\
0.014	-0.0104782\\
0.016	-0.0115456\\
0.018	-0.0125392\\
0.02	-0.0134679\\
0.022	-0.0143397\\
0.024	-0.0151619\\
0.026	-0.015941\\
0.028	-0.0166827\\
0.03	-0.0173923\\
0.032	-0.0180743\\
0.034	-0.0187326\\
0.036	-0.0193709\\
0.038	-0.0199919\\
0.04	-0.0205983\\
0.042	-0.0211921\\
0.044	-0.021775\\
0.046	-0.0223484\\
0.048	-0.0229133\\
0.05	-0.0234705\\
0.052	-0.0240205\\
0.054	-0.0245636\\
0.056	-0.0251001\\
0.058	-0.0256298\\
0.06	-0.0261526\\
0.062	-0.0266685\\
0.064	-0.027177\\
0.066	-0.027678\\
0.068	-0.028171\\
0.07	-0.0286558\\
0.072	-0.0291321\\
0.074	-0.0295995\\
0.076	-0.0300579\\
0.078	-0.030507\\
0.08	-0.0309467\\
0.082	-0.0313768\\
0.084	-0.0317974\\
0.086	-0.0322083\\
0.088	-0.0326098\\
0.09	-0.0330018\\
0.092	-0.0333844\\
0.094	-0.0337579\\
0.096	-0.0341224\\
0.098	-0.0344782\\
0.1	-0.0348253\\
0.102	-0.0351642\\
0.104	-0.0354949\\
0.106	-0.0358176\\
0.108	-0.0361327\\
0.11	-0.0364402\\
0.112	-0.0367404\\
0.114	-0.0370333\\
0.116	-0.0373191\\
0.118	-0.0375978\\
0.12	-0.0378695\\
0.122	-0.0381342\\
0.124	-0.0383918\\
0.126	-0.0386425\\
0.128	-0.0388861\\
0.13	-0.0391225\\
0.132	-0.0393516\\
0.134	-0.0395735\\
0.136	-0.0397878\\
0.138	-0.0399946\\
0.14	-0.0401937\\
0.142	-0.040385\\
0.144	-0.0405685\\
0.146	-0.0407439\\
0.148	-0.0409114\\
0.15	-0.0410707\\
0.152	-0.0412219\\
0.154	-0.041365\\
0.156	-0.0415\\
0.158	-0.041627\\
0.16	-0.0417459\\
0.162	-0.0418569\\
0.164	-0.0419602\\
0.166	-0.0420559\\
0.168	-0.0421441\\
0.17	-0.0422251\\
0.172	-0.0422991\\
0.174	-0.0423663\\
0.176	-0.042427\\
0.178	-0.0424815\\
0.18	-0.04253\\
0.182	-0.0425729\\
0.184	-0.0426104\\
0.186	-0.042643\\
0.188	-0.0426708\\
0.19	-0.0426943\\
0.192	-0.0427137\\
0.194	-0.0427294\\
0.196	-0.0427417\\
0.198	-0.0427508\\
0.2	-0.0427571\\
0.202	-0.0427609\\
0.204	-0.0427624\\
0.206	-0.0427618\\
0.208	-0.0427596\\
0.21	-0.0427558\\
0.212	-0.0427506\\
0.214	-0.0427444\\
0.216	-0.0427373\\
0.218	-0.0427295\\
0.22	-0.0427211\\
0.222	-0.0427123\\
0.224	-0.0427033\\
0.226	-0.0426941\\
0.228	-0.0426849\\
0.23	-0.0426757\\
0.232	-0.0426667\\
0.234	-0.042658\\
0.236	-0.0426496\\
0.238	-0.0426416\\
0.24	-0.042634\\
0.242	-0.0426269\\
0.244	-0.0426203\\
0.246	-0.0426143\\
0.248	-0.0426089\\
0.25	-0.0426041\\
0.252	-0.0425999\\
0.254	-0.0425964\\
0.256	-0.0425935\\
0.258	-0.0425912\\
0.26	-0.0425896\\
0.262	-0.0425887\\
0.264	-0.0425885\\
0.266	-0.0425888\\
0.268	-0.0425899\\
0.27	-0.0425915\\
0.272	-0.0425938\\
0.274	-0.0425967\\
0.276	-0.0426002\\
0.278	-0.0426043\\
0.28	-0.0426089\\
0.282	-0.0426141\\
0.284	-0.0426198\\
0.286	-0.042626\\
0.288	-0.0426327\\
0.29	-0.0426399\\
0.292	-0.0426475\\
0.294	-0.0426555\\
0.296	-0.042664\\
0.298	-0.0426728\\
0.3	-0.042682\\
0.302	-0.0426915\\
0.304	-0.0427013\\
0.306	-0.0427114\\
0.308	-0.0427218\\
0.31	-0.0427324\\
0.312	-0.0427432\\
0.314	-0.0427543\\
0.316	-0.0427655\\
0.318	-0.0427769\\
0.32	-0.0427884\\
0.322	-0.0428\\
0.324	-0.0428118\\
0.326	-0.0428236\\
0.328	-0.0428355\\
0.33	-0.0428474\\
0.332	-0.0428594\\
0.334	-0.0428714\\
0.336	-0.0428834\\
0.338	-0.0428954\\
0.34	-0.0429074\\
0.342	-0.0429193\\
0.344	-0.0429313\\
0.346	-0.0429431\\
0.348	-0.0429549\\
0.35	-0.0429667\\
0.352	-0.0429783\\
0.354	-0.0429899\\
0.356	-0.0430014\\
0.358	-0.0430129\\
0.36	-0.0430242\\
0.362	-0.0430354\\
0.364	-0.0430465\\
0.366	-0.0430575\\
0.368	-0.0430685\\
0.37	-0.0430793\\
0.372	-0.04309\\
0.374	-0.0431005\\
0.376	-0.043111\\
0.378	-0.0431214\\
0.38	-0.0431316\\
0.382	-0.0431417\\
0.384	-0.0431517\\
0.386	-0.0431616\\
0.388	-0.0431713\\
0.39	-0.043181\\
0.392	-0.0431905\\
0.394	-0.0431999\\
0.396	-0.0432092\\
0.398	-0.0432184\\
0.4	-0.0432274\\
0.402	-0.0432363\\
0.404	-0.0432451\\
0.406	-0.0432538\\
0.408	-0.0432623\\
0.41	-0.0432708\\
0.412	-0.0432791\\
0.414	-0.0432873\\
0.416	-0.0432954\\
0.418	-0.0433033\\
0.42	-0.0433111\\
0.422	-0.0433188\\
0.424	-0.0433264\\
0.426	-0.0433338\\
0.428	-0.0433412\\
0.43	-0.0433484\\
0.432	-0.0433554\\
0.434	-0.0433624\\
0.436	-0.0433692\\
0.438	-0.0433759\\
0.44	-0.0433825\\
0.442	-0.0433889\\
0.444	-0.0433952\\
0.446	-0.0434014\\
0.448	-0.0434075\\
0.45	-0.0434135\\
0.452	-0.0434193\\
0.454	-0.043425\\
0.456	-0.0434306\\
0.458	-0.043436\\
0.46	-0.0434414\\
0.462	-0.0434466\\
0.464	-0.0434517\\
0.466	-0.0434567\\
0.468	-0.0434616\\
0.47	-0.0434664\\
0.472	-0.043471\\
0.474	-0.0434756\\
0.476	-0.04348\\
0.478	-0.0434844\\
0.48	-0.0434886\\
0.482	-0.0434927\\
0.484	-0.0434968\\
0.486	-0.0435007\\
0.488	-0.0435045\\
0.49	-0.0435083\\
0.492	-0.0435119\\
0.494	-0.0435155\\
0.496	-0.043519\\
0.498	-0.0435224\\
0.5	-0.0435257\\
0.502	-0.0435289\\
0.504	-0.043532\\
0.506	-0.0435351\\
0.508	-0.0435381\\
0.51	-0.043541\\
0.512	-0.0435439\\
0.514	-0.0435467\\
0.516	-0.0435494\\
0.518	-0.043552\\
0.52	-0.0435546\\
0.522	-0.0435571\\
0.524	-0.0435596\\
0.526	-0.043562\\
0.528	-0.0435644\\
0.53	-0.0435667\\
0.532	-0.0435689\\
0.534	-0.0435711\\
0.536	-0.0435733\\
0.538	-0.0435754\\
0.54	-0.0435774\\
0.542	-0.0435794\\
0.544	-0.0435814\\
0.546	-0.0435833\\
0.548	-0.0435852\\
0.55	-0.0435871\\
0.552	-0.0435889\\
0.554	-0.0435906\\
0.556	-0.0435924\\
0.558	-0.043594\\
0.56	-0.0435957\\
0.562	-0.0435973\\
0.564	-0.0435989\\
0.566	-0.0436005\\
0.568	-0.043602\\
0.57	-0.0436035\\
0.572	-0.043605\\
0.574	-0.0436064\\
0.576	-0.0436078\\
0.578	-0.0436092\\
0.58	-0.0436105\\
0.582	-0.0436119\\
0.584	-0.0436132\\
0.586	-0.0436144\\
0.588	-0.0436157\\
0.59	-0.0436169\\
0.592	-0.0436181\\
0.594	-0.0436192\\
0.596	-0.0436204\\
0.598	-0.0436215\\
0.6	-0.0436226\\
0.602	-0.0436237\\
0.604	-0.0436247\\
0.606	-0.0436258\\
0.608	-0.0436268\\
0.61	-0.0436277\\
0.612	-0.0436287\\
0.614	-0.0436297\\
0.616	-0.0436306\\
0.618	-0.0436315\\
0.62	-0.0436324\\
0.622	-0.0436332\\
0.624	-0.0436341\\
0.626	-0.0436349\\
0.628	-0.0436357\\
0.63	-0.0436365\\
0.632	-0.0436372\\
0.634	-0.043638\\
0.636	-0.0436387\\
0.638	-0.0436394\\
0.64	-0.0436401\\
0.642	-0.0436408\\
0.644	-0.0436415\\
0.646	-0.0436421\\
0.648	-0.0436428\\
0.65	-0.0436434\\
0.652	-0.043644\\
0.654	-0.0436446\\
0.656	-0.0436452\\
0.658	-0.0436457\\
0.66	-0.0436463\\
0.662	-0.0436468\\
0.664	-0.0436474\\
0.666	-0.0436479\\
0.668	-0.0436484\\
0.67	-0.0436489\\
0.672	-0.0436494\\
0.674	-0.0436499\\
0.676	-0.0436503\\
0.678	-0.0436508\\
0.68	-0.0436513\\
0.682	-0.0436517\\
0.684	-0.0436521\\
0.686	-0.0436526\\
0.688	-0.043653\\
0.69	-0.0436534\\
0.692	-0.0436538\\
0.694	-0.0436542\\
0.696	-0.0436546\\
0.698	-0.043655\\
0.7	-0.0436554\\
0.702	-0.0436557\\
0.704	-0.0436561\\
0.706	-0.0436565\\
0.708	-0.0436568\\
0.71	-0.0436572\\
0.712	-0.0436576\\
0.714	-0.0436579\\
0.716	-0.0436583\\
0.718	-0.0436586\\
0.72	-0.043659\\
0.722	-0.0436593\\
0.724	-0.0436596\\
0.726	-0.04366\\
0.728	-0.0436603\\
0.73	-0.0436607\\
0.732	-0.043661\\
0.734	-0.0436613\\
0.736	-0.0436617\\
0.738	-0.043662\\
0.74	-0.0436623\\
0.742	-0.0436627\\
0.744	-0.043663\\
0.746	-0.0436634\\
0.748	-0.0436637\\
0.75	-0.043664\\
0.752	-0.0436644\\
0.754	-0.0436647\\
0.756	-0.043665\\
0.758	-0.0436654\\
0.76	-0.0436657\\
0.762	-0.0436661\\
0.764	-0.0436664\\
0.766	-0.0436668\\
0.768	-0.0436671\\
0.77	-0.0436675\\
0.772	-0.0436678\\
0.774	-0.0436682\\
0.776	-0.0436685\\
0.778	-0.0436689\\
0.78	-0.0436693\\
0.782	-0.0436696\\
0.784	-0.04367\\
0.786	-0.0436703\\
0.788	-0.0436707\\
0.79	-0.0436711\\
0.792	-0.0436715\\
0.794	-0.0436718\\
0.796	-0.0436722\\
0.798	-0.0436726\\
0.8	-0.0436729\\
0.802	-0.0436733\\
0.804	-0.0436737\\
0.806	-0.0436741\\
0.808	-0.0436745\\
0.81	-0.0436748\\
0.812	-0.0436752\\
0.814	-0.0436756\\
0.816	-0.043676\\
0.818	-0.0436763\\
0.82	-0.0436767\\
0.822	-0.0436771\\
0.824	-0.0436775\\
0.826	-0.0436779\\
0.828	-0.0436782\\
0.83	-0.0436786\\
0.832	-0.043679\\
0.834	-0.0436793\\
0.836	-0.0436797\\
0.838	-0.0436801\\
0.84	-0.0436804\\
0.842	-0.0436808\\
0.844	-0.0436812\\
0.846	-0.0436815\\
0.848	-0.0436819\\
0.85	-0.0436822\\
0.852	-0.0436826\\
0.854	-0.0436829\\
0.856	-0.0436833\\
0.858	-0.0436836\\
0.86	-0.0436839\\
0.862	-0.0436843\\
0.864	-0.0436846\\
0.866	-0.0436849\\
0.868	-0.0436852\\
0.87	-0.0436855\\
0.872	-0.0436858\\
0.874	-0.0436861\\
0.876	-0.0436864\\
0.878	-0.0436867\\
0.88	-0.043687\\
0.882	-0.0436873\\
0.884	-0.0436875\\
0.886	-0.0436878\\
0.888	-0.043688\\
0.89	-0.0436883\\
0.892	-0.0436885\\
0.894	-0.0436888\\
0.896	-0.043689\\
0.898	-0.0436892\\
0.9	-0.0436895\\
0.902	-0.0436897\\
0.904	-0.0436899\\
0.906	-0.0436901\\
0.908	-0.0436903\\
0.91	-0.0436905\\
0.912	-0.0436906\\
0.914	-0.0436908\\
0.916	-0.043691\\
0.918	-0.0436912\\
0.92	-0.0436913\\
0.922	-0.0436915\\
0.924	-0.0436916\\
0.926	-0.0436917\\
0.928	-0.0436919\\
0.93	-0.043692\\
0.932	-0.0436921\\
0.934	-0.0436922\\
0.936	-0.0436924\\
0.938	-0.0436925\\
0.94	-0.0436926\\
0.942	-0.0436927\\
0.944	-0.0436928\\
0.946	-0.0436928\\
0.948	-0.0436929\\
0.95	-0.043693\\
0.952	-0.0436931\\
0.954	-0.0436931\\
0.956	-0.0436932\\
0.958	-0.0436933\\
0.96	-0.0436933\\
0.962	-0.0436934\\
0.964	-0.0436934\\
0.966	-0.0436935\\
0.968	-0.0436935\\
0.97	-0.0436935\\
0.972	-0.0436936\\
0.974	-0.0436936\\
0.976	-0.0436936\\
0.978	-0.0436937\\
0.98	-0.0436937\\
0.982	-0.0436937\\
0.984	-0.0436937\\
0.986	-0.0436938\\
0.988	-0.0436938\\
0.99	-0.0436938\\
0.992	-0.0436938\\
0.994	-0.0436938\\
0.996	-0.0436938\\
0.998	-0.0436939\\
1	-0.0436939\\
1.002	-0.0436939\\
1.004	-0.0436939\\
1.006	-0.0436939\\
1.008	-0.0436939\\
1.01	-0.0436939\\
1.012	-0.0436939\\
1.014	-0.0436939\\
1.016	-0.0436939\\
1.018	-0.043694\\
1.02	-0.043694\\
1.022	-0.043694\\
1.024	-0.043694\\
1.026	-0.043694\\
1.028	-0.043694\\
1.03	-0.043694\\
1.032	-0.043694\\
1.034	-0.043694\\
1.036	-0.0436941\\
1.038	-0.0436941\\
1.04	-0.0436941\\
1.042	-0.0436941\\
1.044	-0.0436941\\
1.046	-0.0436941\\
1.048	-0.0436941\\
1.05	-0.0436942\\
1.052	-0.0436942\\
1.054	-0.0436942\\
1.056	-0.0436942\\
1.058	-0.0436943\\
1.06	-0.0436943\\
1.062	-0.0436943\\
1.064	-0.0436943\\
1.066	-0.0436944\\
1.068	-0.0436944\\
1.07	-0.0436944\\
1.072	-0.0436944\\
1.074	-0.0436945\\
1.076	-0.0436945\\
1.078	-0.0436945\\
1.08	-0.0436946\\
1.082	-0.0436946\\
1.084	-0.0436946\\
1.086	-0.0436947\\
1.088	-0.0436947\\
1.09	-0.0436948\\
1.092	-0.0436948\\
1.094	-0.0436948\\
1.096	-0.0436949\\
1.098	-0.0436949\\
1.1	-0.0436949\\
1.102	-0.043695\\
1.104	-0.043695\\
1.106	-0.0436951\\
1.108	-0.0436951\\
1.11	-0.0436952\\
1.112	-0.0436952\\
1.114	-0.0436952\\
1.116	-0.0436953\\
1.118	-0.0436953\\
1.12	-0.0436954\\
1.122	-0.0436954\\
1.124	-0.0436955\\
1.126	-0.0436955\\
1.128	-0.0436955\\
1.13	-0.0436956\\
1.132	-0.0436956\\
1.134	-0.0436957\\
1.136	-0.0436957\\
1.138	-0.0436958\\
1.14	-0.0436958\\
1.142	-0.0436959\\
1.144	-0.0436959\\
1.146	-0.0436959\\
1.148	-0.043696\\
1.15	-0.043696\\
1.152	-0.0436961\\
1.154	-0.0436961\\
1.156	-0.0436962\\
1.158	-0.0436962\\
1.16	-0.0436962\\
1.162	-0.0436963\\
1.164	-0.0436963\\
1.166	-0.0436964\\
1.168	-0.0436964\\
1.17	-0.0436964\\
1.172	-0.0436965\\
1.174	-0.0436965\\
1.176	-0.0436965\\
1.178	-0.0436966\\
1.18	-0.0436966\\
1.182	-0.0436966\\
1.184	-0.0436967\\
1.186	-0.0436967\\
1.188	-0.0436967\\
1.19	-0.0436968\\
1.192	-0.0436968\\
1.194	-0.0436968\\
1.196	-0.0436969\\
1.198	-0.0436969\\
1.2	-0.0436969\\
1.202	-0.0436969\\
1.204	-0.043697\\
1.206	-0.043697\\
1.208	-0.043697\\
1.21	-0.043697\\
1.212	-0.0436971\\
1.214	-0.0436971\\
1.216	-0.0436971\\
1.218	-0.0436971\\
1.22	-0.0436971\\
1.222	-0.0436972\\
1.224	-0.0436972\\
1.226	-0.0436972\\
1.228	-0.0436972\\
1.23	-0.0436972\\
1.232	-0.0436972\\
1.234	-0.0436973\\
1.236	-0.0436973\\
1.238	-0.0436973\\
1.24	-0.0436973\\
1.242	-0.0436973\\
1.244	-0.0436973\\
1.246	-0.0436973\\
1.248	-0.0436973\\
1.25	-0.0436974\\
1.252	-0.0436974\\
1.254	-0.0436974\\
1.256	-0.0436974\\
1.258	-0.0436974\\
1.26	-0.0436974\\
1.262	-0.0436974\\
1.264	-0.0436974\\
1.266	-0.0436974\\
1.268	-0.0436974\\
1.27	-0.0436974\\
1.272	-0.0436974\\
1.274	-0.0436974\\
1.276	-0.0436974\\
1.278	-0.0436974\\
1.28	-0.0436975\\
1.282	-0.0436975\\
1.284	-0.0436975\\
1.286	-0.0436975\\
1.288	-0.0436975\\
1.29	-0.0436975\\
1.292	-0.0436975\\
1.294	-0.0436975\\
1.296	-0.0436975\\
1.298	-0.0436975\\
1.3	-0.0436975\\
1.302	-0.0436975\\
1.304	-0.0436975\\
1.306	-0.0436975\\
1.308	-0.0436975\\
1.31	-0.0436975\\
1.312	-0.0436975\\
1.314	-0.0436975\\
1.316	-0.0436975\\
1.318	-0.0436975\\
1.32	-0.0436975\\
1.322	-0.0436975\\
1.324	-0.0436975\\
1.326	-0.0436975\\
1.328	-0.0436975\\
1.33	-0.0436975\\
1.332	-0.0436975\\
1.334	-0.0436975\\
1.336	-0.0436975\\
1.338	-0.0436975\\
1.34	-0.0436975\\
1.342	-0.0436975\\
1.344	-0.0436975\\
1.346	-0.0436975\\
1.348	-0.0436975\\
1.35	-0.0436975\\
1.352	-0.0436975\\
1.354	-0.0436975\\
1.356	-0.0436975\\
1.358	-0.0436975\\
1.36	-0.0436975\\
1.362	-0.0436975\\
1.364	-0.0436975\\
1.366	-0.0436975\\
1.368	-0.0436975\\
1.37	-0.0436975\\
1.372	-0.0436975\\
1.374	-0.0436976\\
1.376	-0.0436976\\
1.378	-0.0436976\\
1.38	-0.0436976\\
1.382	-0.0436976\\
1.384	-0.0436976\\
1.386	-0.0436976\\
1.388	-0.0436976\\
1.39	-0.0436976\\
1.392	-0.0436976\\
1.394	-0.0436976\\
1.396	-0.0436976\\
1.398	-0.0436976\\
1.4	-0.0436976\\
1.402	-0.0436976\\
1.404	-0.0436976\\
1.406	-0.0436976\\
1.408	-0.0436976\\
1.41	-0.0436976\\
1.412	-0.0436976\\
1.414	-0.0436976\\
1.416	-0.0436976\\
1.418	-0.0436976\\
1.42	-0.0436976\\
1.422	-0.0436976\\
1.424	-0.0436977\\
1.426	-0.0436977\\
1.428	-0.0436977\\
1.43	-0.0436977\\
1.432	-0.0436977\\
1.434	-0.0436977\\
1.436	-0.0436977\\
1.438	-0.0436977\\
1.44	-0.0436977\\
1.442	-0.0436977\\
1.444	-0.0436977\\
1.446	-0.0436977\\
1.448	-0.0436977\\
1.45	-0.0436977\\
1.452	-0.0436977\\
1.454	-0.0436977\\
1.456	-0.0436977\\
1.458	-0.0436977\\
1.46	-0.0436977\\
1.462	-0.0436977\\
1.464	-0.0436977\\
1.466	-0.0436977\\
1.468	-0.0436977\\
1.47	-0.0436977\\
1.472	-0.0436978\\
1.474	-0.0436978\\
1.476	-0.0436978\\
1.478	-0.0436978\\
1.48	-0.0436978\\
1.482	-0.0436978\\
1.484	-0.0436978\\
1.486	-0.0436978\\
1.488	-0.0436978\\
1.49	-0.0436978\\
1.492	-0.0436978\\
1.494	-0.0436978\\
1.496	-0.0436978\\
1.498	-0.0436978\\
1.5	-0.0436978\\
1.502	-0.0436978\\
1.504	-0.0436978\\
1.506	-0.0436978\\
1.508	-0.0436978\\
1.51	-0.0436978\\
1.512	-0.0436978\\
1.514	-0.0436978\\
1.516	-0.0436978\\
1.518	-0.0436978\\
1.52	-0.0436978\\
1.522	-0.0436978\\
1.524	-0.0436978\\
1.526	-0.0436978\\
1.528	-0.0436978\\
1.53	-0.0436978\\
1.532	-0.0436978\\
1.534	-0.0436978\\
1.536	-0.0436978\\
1.538	-0.0436978\\
1.54	-0.0436978\\
1.542	-0.0436978\\
1.544	-0.0436978\\
1.546	-0.0436978\\
1.548	-0.0436978\\
1.55	-0.0436978\\
1.552	-0.0436978\\
1.554	-0.0436978\\
1.556	-0.0436978\\
1.558	-0.0436978\\
1.56	-0.0436978\\
1.562	-0.0436978\\
1.564	-0.0436978\\
1.566	-0.0436979\\
1.568	-0.0436979\\
1.57	-0.0436979\\
1.572	-0.0436979\\
1.574	-0.0436979\\
1.576	-0.0436979\\
1.578	-0.0436979\\
1.58	-0.0436979\\
1.582	-0.0436979\\
1.584	-0.0436979\\
1.586	-0.0436979\\
1.588	-0.0436979\\
1.59	-0.0436979\\
1.592	-0.0436979\\
1.594	-0.0436979\\
1.596	-0.0436979\\
1.598	-0.0436979\\
1.6	-0.0436979\\
1.602	-0.0436979\\
1.604	-0.0436979\\
1.606	-0.0436979\\
1.608	-0.0436979\\
1.61	-0.0436979\\
1.612	-0.0436979\\
1.614	-0.0436979\\
1.616	-0.0436979\\
1.618	-0.0436979\\
1.62	-0.0436979\\
1.622	-0.0436979\\
1.624	-0.0436979\\
1.626	-0.0436979\\
1.628	-0.0436979\\
1.63	-0.0436979\\
1.632	-0.0436979\\
1.634	-0.0436979\\
1.636	-0.0436979\\
1.638	-0.0436979\\
1.64	-0.0436979\\
1.642	-0.0436979\\
1.644	-0.0436979\\
1.646	-0.0436979\\
1.648	-0.0436979\\
1.65	-0.0436979\\
1.652	-0.0436979\\
1.654	-0.0436979\\
1.656	-0.0436979\\
1.658	-0.0436979\\
1.66	-0.0436979\\
1.662	-0.0436979\\
1.664	-0.0436979\\
1.666	-0.0436979\\
1.668	-0.0436979\\
1.67	-0.0436979\\
1.672	-0.0436979\\
1.674	-0.0436979\\
1.676	-0.0436979\\
1.678	-0.0436979\\
1.68	-0.0436979\\
1.682	-0.0436979\\
1.684	-0.0436979\\
1.686	-0.0436979\\
1.688	-0.0436979\\
1.69	-0.0436979\\
1.692	-0.0436979\\
1.694	-0.0436979\\
1.696	-0.0436979\\
1.698	-0.0436979\\
1.7	-0.0436979\\
1.702	-0.0436979\\
1.704	-0.0436979\\
1.706	-0.0436979\\
1.708	-0.0436979\\
1.71	-0.0436979\\
1.712	-0.0436979\\
1.714	-0.0436979\\
1.716	-0.0436979\\
1.718	-0.0436979\\
1.72	-0.0436979\\
1.722	-0.0436979\\
1.724	-0.0436979\\
1.726	-0.0436979\\
1.728	-0.0436979\\
1.73	-0.0436979\\
1.732	-0.0436979\\
1.734	-0.0436979\\
1.736	-0.0436979\\
1.738	-0.0436979\\
1.74	-0.0436979\\
1.742	-0.0436979\\
1.744	-0.0436979\\
1.746	-0.0436979\\
1.748	-0.0436979\\
1.75	-0.0436979\\
1.752	-0.0436979\\
1.754	-0.0436979\\
1.756	-0.0436979\\
1.758	-0.0436979\\
1.76	-0.0436979\\
1.762	-0.0436979\\
1.764	-0.0436979\\
1.766	-0.0436979\\
1.768	-0.0436979\\
1.77	-0.0436979\\
1.772	-0.0436979\\
1.774	-0.0436979\\
1.776	-0.0436979\\
1.778	-0.0436979\\
1.78	-0.0436979\\
1.782	-0.0436979\\
1.784	-0.0436979\\
1.786	-0.0436979\\
1.788	-0.0436979\\
1.79	-0.0436979\\
1.792	-0.0436979\\
1.794	-0.0436979\\
1.796	-0.0436979\\
1.798	-0.0436979\\
1.8	-0.0436979\\
1.802	-0.0436979\\
1.804	-0.0436979\\
1.806	-0.0436979\\
1.808	-0.0436979\\
1.81	-0.0436979\\
1.812	-0.0436979\\
1.814	-0.0436979\\
1.816	-0.0436979\\
1.818	-0.0436979\\
1.82	-0.0436979\\
1.822	-0.0436979\\
1.824	-0.0436979\\
1.826	-0.0436979\\
1.828	-0.0436979\\
1.83	-0.0436979\\
1.832	-0.0436979\\
1.834	-0.0436979\\
1.836	-0.0436979\\
1.838	-0.0436979\\
1.84	-0.0436979\\
1.842	-0.0436979\\
1.844	-0.0436979\\
1.846	-0.0436979\\
1.848	-0.0436979\\
1.85	-0.0436979\\
1.852	-0.0436979\\
1.854	-0.0436979\\
1.856	-0.0436979\\
1.858	-0.0436979\\
1.86	-0.0436979\\
1.862	-0.0436979\\
1.864	-0.0436979\\
1.866	-0.0436979\\
1.868	-0.0436979\\
1.87	-0.0436979\\
1.872	-0.0436979\\
1.874	-0.0436979\\
1.876	-0.0436979\\
1.878	-0.0436979\\
1.88	-0.0436979\\
1.882	-0.0436979\\
1.884	-0.0436979\\
1.886	-0.0436979\\
1.888	-0.0436979\\
1.89	-0.0436979\\
1.892	-0.0436979\\
1.894	-0.0436979\\
1.896	-0.0436979\\
1.898	-0.0436979\\
1.9	-0.0436979\\
1.902	-0.0436979\\
1.904	-0.0436979\\
1.906	-0.0436979\\
1.908	-0.0436979\\
1.91	-0.0436979\\
1.912	-0.0436979\\
1.914	-0.0436979\\
1.916	-0.0436979\\
1.918	-0.0436979\\
1.92	-0.0436979\\
1.922	-0.0436979\\
1.924	-0.0436979\\
1.926	-0.0436979\\
1.928	-0.0436979\\
1.93	-0.0436979\\
1.932	-0.0436979\\
1.934	-0.0436979\\
1.936	-0.0436979\\
1.938	-0.0436979\\
1.94	-0.0436979\\
1.942	-0.0436979\\
1.944	-0.0436979\\
1.946	-0.0436979\\
1.948	-0.0436979\\
1.95	-0.0436979\\
1.952	-0.0436979\\
1.954	-0.0436979\\
1.956	-0.0436979\\
1.958	-0.0436979\\
1.96	-0.0436979\\
1.962	-0.0436979\\
1.964	-0.0436979\\
1.966	-0.0436979\\
1.968	-0.0436979\\
1.97	-0.0436979\\
1.972	-0.0436979\\
1.974	-0.0436979\\
1.976	-0.0436979\\
1.978	-0.0436979\\
1.98	-0.0436979\\
1.982	-0.0436979\\
1.984	-0.0436979\\
1.986	-0.0436979\\
1.988	-0.0436979\\
1.99	-0.0436979\\
1.992	-0.0436979\\
1.994	-0.0436979\\
1.996	-0.0436979\\
1.998	-0.0436979\\
2	-0.0436979\\
};
\draw[dashed, line width=0.7pt, draw=black] (axis cs:0,-0.0437676) rectangle (axis cs:0.5,-0.0413803);
\end{axis}

\begin{axis}[%
width={\figscale*1.528in},
height={\figscale*0.46in},
at={(\figscale*2.841in,\figscale*1.69in)},
scale only axis,
separate axis lines,
every outer x axis line/.append style={black},
every x tick label/.append style={font=\figticfont},
every x tick/.append style={black},
xmin=-0.025,
xmax=2,
xtick={0,0.5,1,1.5,2},
xticklabels={\empty},
every outer y axis line/.append style={black},
every y tick label/.append style={font=\figticfont},
every y tick/.append style={black},
ymin=-0.0206901,
ymax=0,
ytick={-0.02,0},
axis background/.style={fill=white},
title style={font=\figcornerfont},
title={SC500},
xmajorgrids,
ymajorgrids,
grid style={white!55!black, opacity=0.3},
ylabel style={rotate=-90, at={(axis description cs:-0.145,0.5)}, anchor=east}
]
\addplot [color=mycolor1, line width=1.4pt, forget plot]
  table[row sep=crcr]{%
-0.024	0\\
-0.022	0\\
-0.02	0\\
-0.018	0\\
-0.016	0\\
-0.014	0\\
-0.012	0\\
-0.01	0\\
-0.008	0\\
-0.006	0\\
-0.004	0\\
-0.002	0\\
0	0\\
0.002	-0.00183136\\
0.004	-0.00337796\\
0.006	-0.00470392\\
0.008	-0.00585756\\
0.01	-0.00687514\\
0.012	-0.00778386\\
0.014	-0.00860411\\
0.016	-0.00935119\\
0.018	-0.0100367\\
0.02	-0.0106695\\
0.022	-0.0112563\\
0.024	-0.0118026\\
0.026	-0.0123126\\
0.028	-0.0127899\\
0.03	-0.0132372\\
0.032	-0.013657\\
0.034	-0.0140516\\
0.036	-0.0144226\\
0.038	-0.0147718\\
0.04	-0.0151007\\
0.042	-0.0154107\\
0.044	-0.0157029\\
0.046	-0.0159786\\
0.048	-0.0162389\\
0.05	-0.0164849\\
0.052	-0.0167174\\
0.054	-0.0169375\\
0.056	-0.017146\\
0.058	-0.0173437\\
0.06	-0.0175314\\
0.062	-0.0177098\\
0.064	-0.0178796\\
0.066	-0.0180412\\
0.068	-0.0181953\\
0.07	-0.0183424\\
0.072	-0.0184828\\
0.074	-0.0186169\\
0.076	-0.0187451\\
0.078	-0.0188675\\
0.08	-0.0189844\\
0.082	-0.0190959\\
0.084	-0.0192021\\
0.086	-0.0193032\\
0.088	-0.0193992\\
0.09	-0.0194901\\
0.092	-0.019576\\
0.094	-0.0196567\\
0.096	-0.0197323\\
0.098	-0.0198028\\
0.1	-0.0198681\\
0.102	-0.0199282\\
0.104	-0.0199831\\
0.106	-0.0200328\\
0.108	-0.0200772\\
0.11	-0.0201164\\
0.112	-0.0201504\\
0.114	-0.0201792\\
0.116	-0.0202028\\
0.118	-0.0202214\\
0.12	-0.0202349\\
0.122	-0.0202436\\
0.124	-0.0202474\\
0.126	-0.0202466\\
0.128	-0.0202412\\
0.13	-0.0202315\\
0.132	-0.0202175\\
0.134	-0.0201996\\
0.136	-0.0201777\\
0.138	-0.0201522\\
0.14	-0.0201233\\
0.142	-0.020091\\
0.144	-0.0200558\\
0.146	-0.0200177\\
0.148	-0.0199771\\
0.15	-0.019934\\
0.152	-0.0198888\\
0.154	-0.0198417\\
0.156	-0.0197929\\
0.158	-0.0197425\\
0.16	-0.0196909\\
0.162	-0.0196382\\
0.164	-0.0195846\\
0.166	-0.0195304\\
0.168	-0.0194757\\
0.17	-0.0194207\\
0.172	-0.0193656\\
0.174	-0.0193106\\
0.176	-0.0192559\\
0.178	-0.0192016\\
0.18	-0.0191479\\
0.182	-0.0190949\\
0.184	-0.0190428\\
0.186	-0.0189917\\
0.188	-0.0189418\\
0.19	-0.0188931\\
0.192	-0.0188458\\
0.194	-0.0188\\
0.196	-0.0187558\\
0.198	-0.0187132\\
0.2	-0.0186724\\
0.202	-0.0186335\\
0.204	-0.0185964\\
0.206	-0.0185613\\
0.208	-0.0185282\\
0.21	-0.0184972\\
0.212	-0.0184682\\
0.214	-0.0184413\\
0.216	-0.0184166\\
0.218	-0.018394\\
0.22	-0.0183735\\
0.222	-0.0183552\\
0.224	-0.018339\\
0.226	-0.0183249\\
0.228	-0.0183129\\
0.23	-0.018303\\
0.232	-0.0182951\\
0.234	-0.0182891\\
0.236	-0.0182851\\
0.238	-0.018283\\
0.24	-0.0182826\\
0.242	-0.0182841\\
0.244	-0.0182873\\
0.246	-0.0182921\\
0.248	-0.0182984\\
0.25	-0.0183063\\
0.252	-0.0183156\\
0.254	-0.0183262\\
0.256	-0.0183381\\
0.258	-0.0183512\\
0.26	-0.0183654\\
0.262	-0.0183806\\
0.264	-0.0183967\\
0.266	-0.0184137\\
0.268	-0.0184315\\
0.27	-0.01845\\
0.272	-0.0184691\\
0.274	-0.0184887\\
0.276	-0.0185088\\
0.278	-0.0185293\\
0.28	-0.01855\\
0.282	-0.018571\\
0.284	-0.0185922\\
0.286	-0.0186134\\
0.288	-0.0186346\\
0.29	-0.0186558\\
0.292	-0.0186769\\
0.294	-0.0186978\\
0.296	-0.0187184\\
0.298	-0.0187388\\
0.3	-0.0187588\\
0.302	-0.0187784\\
0.304	-0.0187975\\
0.306	-0.0188162\\
0.308	-0.0188344\\
0.31	-0.018852\\
0.312	-0.018869\\
0.314	-0.0188853\\
0.316	-0.018901\\
0.318	-0.0189161\\
0.32	-0.0189304\\
0.322	-0.018944\\
0.324	-0.0189569\\
0.326	-0.0189691\\
0.328	-0.0189806\\
0.33	-0.0189913\\
0.332	-0.0190012\\
0.334	-0.0190104\\
0.336	-0.0190189\\
0.338	-0.0190267\\
0.34	-0.0190338\\
0.342	-0.0190402\\
0.344	-0.0190459\\
0.346	-0.0190509\\
0.348	-0.0190553\\
0.35	-0.0190592\\
0.352	-0.0190624\\
0.354	-0.019065\\
0.356	-0.0190672\\
0.358	-0.0190688\\
0.36	-0.01907\\
0.362	-0.0190707\\
0.364	-0.019071\\
0.366	-0.0190709\\
0.368	-0.0190705\\
0.37	-0.0190697\\
0.372	-0.0190687\\
0.374	-0.0190674\\
0.376	-0.0190659\\
0.378	-0.0190641\\
0.38	-0.0190622\\
0.382	-0.0190602\\
0.384	-0.019058\\
0.386	-0.0190558\\
0.388	-0.0190534\\
0.39	-0.019051\\
0.392	-0.0190486\\
0.394	-0.0190462\\
0.396	-0.0190437\\
0.398	-0.0190413\\
0.4	-0.0190389\\
0.402	-0.0190366\\
0.404	-0.0190343\\
0.406	-0.0190321\\
0.408	-0.0190299\\
0.41	-0.0190279\\
0.412	-0.0190259\\
0.414	-0.019024\\
0.416	-0.0190222\\
0.418	-0.0190204\\
0.42	-0.0190188\\
0.422	-0.0190172\\
0.424	-0.0190157\\
0.426	-0.0190143\\
0.428	-0.019013\\
0.43	-0.0190117\\
0.432	-0.0190105\\
0.434	-0.0190093\\
0.436	-0.0190082\\
0.438	-0.0190071\\
0.44	-0.0190061\\
0.442	-0.019005\\
0.444	-0.019004\\
0.446	-0.0190029\\
0.448	-0.0190019\\
0.45	-0.0190008\\
0.452	-0.0189997\\
0.454	-0.0189985\\
0.456	-0.0189973\\
0.458	-0.018996\\
0.46	-0.0189947\\
0.462	-0.0189933\\
0.464	-0.0189918\\
0.466	-0.0189903\\
0.468	-0.0189886\\
0.47	-0.0189868\\
0.472	-0.018985\\
0.474	-0.018983\\
0.476	-0.0189809\\
0.478	-0.0189787\\
0.48	-0.0189763\\
0.482	-0.0189739\\
0.484	-0.0189713\\
0.486	-0.0189686\\
0.488	-0.0189658\\
0.49	-0.0189628\\
0.492	-0.0189598\\
0.494	-0.0189566\\
0.496	-0.0189533\\
0.498	-0.0189499\\
0.5	-0.0189464\\
0.502	-0.0189428\\
0.504	-0.0189391\\
0.506	-0.0189353\\
0.508	-0.0189314\\
0.51	-0.0189275\\
0.512	-0.0189234\\
0.514	-0.0189194\\
0.516	-0.0189152\\
0.518	-0.0189111\\
0.52	-0.0189069\\
0.522	-0.0189026\\
0.524	-0.0188984\\
0.526	-0.0188942\\
0.528	-0.0188899\\
0.53	-0.0188857\\
0.532	-0.0188816\\
0.534	-0.0188774\\
0.536	-0.0188734\\
0.538	-0.0188694\\
0.54	-0.0188654\\
0.542	-0.0188616\\
0.544	-0.0188579\\
0.546	-0.0188543\\
0.548	-0.0188508\\
0.55	-0.0188474\\
0.552	-0.0188442\\
0.554	-0.0188411\\
0.556	-0.0188383\\
0.558	-0.0188356\\
0.56	-0.018833\\
0.562	-0.0188307\\
0.564	-0.0188286\\
0.566	-0.0188267\\
0.568	-0.018825\\
0.57	-0.0188236\\
0.572	-0.0188224\\
0.574	-0.0188214\\
0.576	-0.0188207\\
0.578	-0.0188202\\
0.58	-0.01882\\
0.582	-0.01882\\
0.584	-0.0188204\\
0.586	-0.0188209\\
0.588	-0.0188218\\
0.59	-0.0188229\\
0.592	-0.0188243\\
0.594	-0.0188259\\
0.596	-0.0188278\\
0.598	-0.01883\\
0.6	-0.0188324\\
0.602	-0.0188351\\
0.604	-0.018838\\
0.606	-0.0188411\\
0.608	-0.0188445\\
0.61	-0.0188481\\
0.612	-0.0188519\\
0.614	-0.018856\\
0.616	-0.0188602\\
0.618	-0.0188646\\
0.62	-0.0188692\\
0.622	-0.0188739\\
0.624	-0.0188788\\
0.626	-0.0188838\\
0.628	-0.018889\\
0.63	-0.0188942\\
0.632	-0.0188996\\
0.634	-0.018905\\
0.636	-0.0189105\\
0.638	-0.0189161\\
0.64	-0.0189216\\
0.642	-0.0189272\\
0.644	-0.0189328\\
0.646	-0.0189384\\
0.648	-0.018944\\
0.65	-0.0189495\\
0.652	-0.0189549\\
0.654	-0.0189603\\
0.656	-0.0189656\\
0.658	-0.0189707\\
0.66	-0.0189758\\
0.662	-0.0189807\\
0.664	-0.0189855\\
0.666	-0.01899\\
0.668	-0.0189945\\
0.67	-0.0189987\\
0.672	-0.0190027\\
0.674	-0.0190066\\
0.676	-0.0190102\\
0.678	-0.0190135\\
0.68	-0.0190167\\
0.682	-0.0190195\\
0.684	-0.0190222\\
0.686	-0.0190245\\
0.688	-0.0190266\\
0.69	-0.0190284\\
0.692	-0.01903\\
0.694	-0.0190313\\
0.696	-0.0190322\\
0.698	-0.0190329\\
0.7	-0.0190334\\
0.702	-0.0190335\\
0.704	-0.0190333\\
0.706	-0.0190329\\
0.708	-0.0190322\\
0.71	-0.0190312\\
0.712	-0.0190299\\
0.714	-0.0190284\\
0.716	-0.0190266\\
0.718	-0.0190246\\
0.72	-0.0190223\\
0.722	-0.0190198\\
0.724	-0.019017\\
0.726	-0.019014\\
0.728	-0.0190109\\
0.73	-0.0190075\\
0.732	-0.0190039\\
0.734	-0.0190002\\
0.736	-0.0189963\\
0.738	-0.0189922\\
0.74	-0.018988\\
0.742	-0.0189837\\
0.744	-0.0189792\\
0.746	-0.0189747\\
0.748	-0.0189701\\
0.75	-0.0189654\\
0.752	-0.0189606\\
0.754	-0.0189558\\
0.756	-0.018951\\
0.758	-0.0189461\\
0.76	-0.0189413\\
0.762	-0.0189365\\
0.764	-0.0189316\\
0.766	-0.0189269\\
0.768	-0.0189221\\
0.77	-0.0189175\\
0.772	-0.0189129\\
0.774	-0.0189084\\
0.776	-0.018904\\
0.778	-0.0188998\\
0.78	-0.0188956\\
0.782	-0.0188916\\
0.784	-0.0188877\\
0.786	-0.018884\\
0.788	-0.0188804\\
0.79	-0.0188771\\
0.792	-0.0188738\\
0.794	-0.0188708\\
0.796	-0.018868\\
0.798	-0.0188654\\
0.8	-0.0188629\\
0.802	-0.0188607\\
0.804	-0.0188587\\
0.806	-0.0188569\\
0.808	-0.0188554\\
0.81	-0.018854\\
0.812	-0.0188529\\
0.814	-0.018852\\
0.816	-0.0188513\\
0.818	-0.0188509\\
0.82	-0.0188506\\
0.822	-0.0188506\\
0.824	-0.0188508\\
0.826	-0.0188512\\
0.828	-0.0188519\\
0.83	-0.0188527\\
0.832	-0.0188537\\
0.834	-0.018855\\
0.836	-0.0188564\\
0.838	-0.018858\\
0.84	-0.0188598\\
0.842	-0.0188617\\
0.844	-0.0188638\\
0.846	-0.0188661\\
0.848	-0.0188685\\
0.85	-0.0188711\\
0.852	-0.0188737\\
0.854	-0.0188765\\
0.856	-0.0188794\\
0.858	-0.0188824\\
0.86	-0.0188855\\
0.862	-0.0188887\\
0.864	-0.018892\\
0.866	-0.0188953\\
0.868	-0.0188986\\
0.87	-0.018902\\
0.872	-0.0189054\\
0.874	-0.0189089\\
0.876	-0.0189123\\
0.878	-0.0189157\\
0.88	-0.0189192\\
0.882	-0.0189226\\
0.884	-0.018926\\
0.886	-0.0189293\\
0.888	-0.0189326\\
0.89	-0.0189358\\
0.892	-0.018939\\
0.894	-0.0189421\\
0.896	-0.0189451\\
0.898	-0.018948\\
0.9	-0.0189509\\
0.902	-0.0189536\\
0.904	-0.0189562\\
0.906	-0.0189587\\
0.908	-0.0189611\\
0.91	-0.0189633\\
0.912	-0.0189655\\
0.914	-0.0189675\\
0.916	-0.0189693\\
0.918	-0.018971\\
0.92	-0.0189726\\
0.922	-0.018974\\
0.924	-0.0189753\\
0.926	-0.0189764\\
0.928	-0.0189774\\
0.93	-0.0189782\\
0.932	-0.0189789\\
0.934	-0.0189794\\
0.936	-0.0189798\\
0.938	-0.01898\\
0.94	-0.0189801\\
0.942	-0.01898\\
0.944	-0.0189798\\
0.946	-0.0189795\\
0.948	-0.018979\\
0.95	-0.0189784\\
0.952	-0.0189777\\
0.954	-0.0189768\\
0.956	-0.0189758\\
0.958	-0.0189747\\
0.96	-0.0189735\\
0.962	-0.0189722\\
0.964	-0.0189708\\
0.966	-0.0189693\\
0.968	-0.0189678\\
0.97	-0.0189661\\
0.972	-0.0189644\\
0.974	-0.0189626\\
0.976	-0.0189607\\
0.978	-0.0189588\\
0.98	-0.0189569\\
0.982	-0.0189548\\
0.984	-0.0189528\\
0.986	-0.0189507\\
0.988	-0.0189487\\
0.99	-0.0189466\\
0.992	-0.0189444\\
0.994	-0.0189423\\
0.996	-0.0189402\\
0.998	-0.0189381\\
1	-0.018936\\
1.002	-0.018934\\
1.004	-0.0189319\\
1.006	-0.0189299\\
1.008	-0.0189279\\
1.01	-0.018926\\
1.012	-0.0189241\\
1.014	-0.0189223\\
1.016	-0.0189205\\
1.018	-0.0189188\\
1.02	-0.0189172\\
1.022	-0.0189156\\
1.024	-0.018914\\
1.026	-0.0189126\\
1.028	-0.0189112\\
1.03	-0.0189099\\
1.032	-0.0189087\\
1.034	-0.0189076\\
1.036	-0.0189065\\
1.038	-0.0189056\\
1.04	-0.0189047\\
1.042	-0.0189039\\
1.044	-0.0189032\\
1.046	-0.0189026\\
1.048	-0.018902\\
1.05	-0.0189016\\
1.052	-0.0189012\\
1.054	-0.018901\\
1.056	-0.0189008\\
1.058	-0.0189007\\
1.06	-0.0189006\\
1.062	-0.0189007\\
1.064	-0.0189008\\
1.066	-0.018901\\
1.068	-0.0189013\\
1.07	-0.0189016\\
1.072	-0.018902\\
1.074	-0.0189025\\
1.076	-0.0189031\\
1.078	-0.0189037\\
1.08	-0.0189043\\
1.082	-0.018905\\
1.084	-0.0189058\\
1.086	-0.0189066\\
1.088	-0.0189075\\
1.09	-0.0189083\\
1.092	-0.0189093\\
1.094	-0.0189102\\
1.096	-0.0189112\\
1.098	-0.0189122\\
1.1	-0.0189133\\
1.102	-0.0189143\\
1.104	-0.0189154\\
1.106	-0.0189165\\
1.108	-0.0189176\\
1.11	-0.0189186\\
1.112	-0.0189197\\
1.114	-0.0189208\\
1.116	-0.0189219\\
1.118	-0.018923\\
1.12	-0.0189241\\
1.122	-0.0189252\\
1.124	-0.0189262\\
1.126	-0.0189272\\
1.128	-0.0189282\\
1.13	-0.0189292\\
1.132	-0.0189302\\
1.134	-0.0189311\\
1.136	-0.018932\\
1.138	-0.0189329\\
1.14	-0.0189337\\
1.142	-0.0189345\\
1.144	-0.0189353\\
1.146	-0.018936\\
1.148	-0.0189367\\
1.15	-0.0189374\\
1.152	-0.018938\\
1.154	-0.0189386\\
1.156	-0.0189391\\
1.158	-0.0189396\\
1.16	-0.0189401\\
1.162	-0.0189405\\
1.164	-0.0189409\\
1.166	-0.0189412\\
1.168	-0.0189415\\
1.17	-0.0189418\\
1.172	-0.018942\\
1.174	-0.0189422\\
1.176	-0.0189423\\
1.178	-0.0189424\\
1.18	-0.0189425\\
1.182	-0.0189425\\
1.184	-0.0189425\\
1.186	-0.0189425\\
1.188	-0.0189424\\
1.19	-0.0189423\\
1.192	-0.0189422\\
1.194	-0.018942\\
1.196	-0.0189418\\
1.198	-0.0189416\\
1.2	-0.0189414\\
1.202	-0.0189411\\
1.204	-0.0189408\\
1.206	-0.0189405\\
1.208	-0.0189402\\
1.21	-0.0189399\\
1.212	-0.0189395\\
1.214	-0.0189391\\
1.216	-0.0189387\\
1.218	-0.0189383\\
1.22	-0.0189379\\
1.222	-0.0189375\\
1.224	-0.0189371\\
1.226	-0.0189367\\
1.228	-0.0189362\\
1.23	-0.0189358\\
1.232	-0.0189354\\
1.234	-0.0189349\\
1.236	-0.0189345\\
1.238	-0.018934\\
1.24	-0.0189336\\
1.242	-0.0189332\\
1.244	-0.0189327\\
1.246	-0.0189323\\
1.248	-0.0189319\\
1.25	-0.0189315\\
1.252	-0.0189311\\
1.254	-0.0189307\\
1.256	-0.0189303\\
1.258	-0.0189299\\
1.26	-0.0189295\\
1.262	-0.0189292\\
1.264	-0.0189288\\
1.266	-0.0189285\\
1.268	-0.0189282\\
1.27	-0.0189279\\
1.272	-0.0189275\\
1.274	-0.0189273\\
1.276	-0.018927\\
1.278	-0.0189267\\
1.28	-0.0189265\\
1.282	-0.0189262\\
1.284	-0.018926\\
1.286	-0.0189258\\
1.288	-0.0189256\\
1.29	-0.0189254\\
1.292	-0.0189252\\
1.294	-0.018925\\
1.296	-0.0189249\\
1.298	-0.0189247\\
1.3	-0.0189246\\
1.302	-0.0189245\\
1.304	-0.0189244\\
1.306	-0.0189243\\
1.308	-0.0189242\\
1.31	-0.0189241\\
1.312	-0.0189241\\
1.314	-0.018924\\
1.316	-0.018924\\
1.318	-0.0189239\\
1.32	-0.0189239\\
1.322	-0.0189239\\
1.324	-0.0189239\\
1.326	-0.0189239\\
1.328	-0.0189239\\
1.33	-0.0189239\\
1.332	-0.0189239\\
1.334	-0.0189239\\
1.336	-0.0189239\\
1.338	-0.018924\\
1.34	-0.018924\\
1.342	-0.0189241\\
1.344	-0.0189241\\
1.346	-0.0189242\\
1.348	-0.0189243\\
1.35	-0.0189244\\
1.352	-0.0189244\\
1.354	-0.0189245\\
1.356	-0.0189246\\
1.358	-0.0189247\\
1.36	-0.0189248\\
1.362	-0.0189249\\
1.364	-0.018925\\
1.366	-0.0189252\\
1.368	-0.0189253\\
1.37	-0.0189254\\
1.372	-0.0189255\\
1.374	-0.0189257\\
1.376	-0.0189258\\
1.378	-0.0189259\\
1.38	-0.0189261\\
1.382	-0.0189262\\
1.384	-0.0189264\\
1.386	-0.0189265\\
1.388	-0.0189267\\
1.39	-0.0189268\\
1.392	-0.018927\\
1.394	-0.0189272\\
1.396	-0.0189273\\
1.398	-0.0189275\\
1.4	-0.0189277\\
1.402	-0.0189279\\
1.404	-0.018928\\
1.406	-0.0189282\\
1.408	-0.0189284\\
1.41	-0.0189286\\
1.412	-0.0189287\\
1.414	-0.0189289\\
1.416	-0.0189291\\
1.418	-0.0189293\\
1.42	-0.0189295\\
1.422	-0.0189297\\
1.424	-0.0189298\\
1.426	-0.01893\\
1.428	-0.0189302\\
1.43	-0.0189304\\
1.432	-0.0189306\\
1.434	-0.0189307\\
1.436	-0.0189309\\
1.438	-0.0189311\\
1.44	-0.0189312\\
1.442	-0.0189314\\
1.444	-0.0189316\\
1.446	-0.0189317\\
1.448	-0.0189319\\
1.45	-0.018932\\
1.452	-0.0189322\\
1.454	-0.0189323\\
1.456	-0.0189324\\
1.458	-0.0189325\\
1.46	-0.0189327\\
1.462	-0.0189328\\
1.464	-0.0189329\\
1.466	-0.018933\\
1.468	-0.018933\\
1.47	-0.0189331\\
1.472	-0.0189332\\
1.474	-0.0189332\\
1.476	-0.0189333\\
1.478	-0.0189333\\
1.48	-0.0189334\\
1.482	-0.0189334\\
1.484	-0.0189334\\
1.486	-0.0189334\\
1.488	-0.0189334\\
1.49	-0.0189334\\
1.492	-0.0189334\\
1.494	-0.0189334\\
1.496	-0.0189333\\
1.498	-0.0189333\\
1.5	-0.0189332\\
1.502	-0.0189331\\
1.504	-0.018933\\
1.506	-0.018933\\
1.508	-0.0189329\\
1.51	-0.0189327\\
1.512	-0.0189326\\
1.514	-0.0189325\\
1.516	-0.0189324\\
1.518	-0.0189322\\
1.52	-0.0189321\\
1.522	-0.0189319\\
1.524	-0.0189318\\
1.526	-0.0189316\\
1.528	-0.0189314\\
1.53	-0.0189312\\
1.532	-0.018931\\
1.534	-0.0189308\\
1.536	-0.0189307\\
1.538	-0.0189305\\
1.54	-0.0189303\\
1.542	-0.01893\\
1.544	-0.0189298\\
1.546	-0.0189296\\
1.548	-0.0189294\\
1.55	-0.0189292\\
1.552	-0.018929\\
1.554	-0.0189288\\
1.556	-0.0189286\\
1.558	-0.0189284\\
1.56	-0.0189282\\
1.562	-0.018928\\
1.564	-0.0189278\\
1.566	-0.0189276\\
1.568	-0.0189274\\
1.57	-0.0189272\\
1.572	-0.018927\\
1.574	-0.0189268\\
1.576	-0.0189267\\
1.578	-0.0189265\\
1.58	-0.0189264\\
1.582	-0.0189262\\
1.584	-0.0189261\\
1.586	-0.0189259\\
1.588	-0.0189258\\
1.59	-0.0189257\\
1.592	-0.0189256\\
1.594	-0.0189255\\
1.596	-0.0189254\\
1.598	-0.0189254\\
1.6	-0.0189253\\
1.602	-0.0189253\\
1.604	-0.0189252\\
1.606	-0.0189252\\
1.608	-0.0189252\\
1.61	-0.0189252\\
1.612	-0.0189252\\
1.614	-0.0189252\\
1.616	-0.0189252\\
1.618	-0.0189253\\
1.62	-0.0189253\\
1.622	-0.0189254\\
1.624	-0.0189254\\
1.626	-0.0189255\\
1.628	-0.0189256\\
1.63	-0.0189257\\
1.632	-0.0189258\\
1.634	-0.0189259\\
1.636	-0.0189261\\
1.638	-0.0189262\\
1.64	-0.0189263\\
1.642	-0.0189265\\
1.644	-0.0189266\\
1.646	-0.0189268\\
1.648	-0.018927\\
1.65	-0.0189271\\
1.652	-0.0189273\\
1.654	-0.0189275\\
1.656	-0.0189277\\
1.658	-0.0189279\\
1.66	-0.0189281\\
1.662	-0.0189283\\
1.664	-0.0189284\\
1.666	-0.0189286\\
1.668	-0.0189288\\
1.67	-0.018929\\
1.672	-0.0189292\\
1.674	-0.0189294\\
1.676	-0.0189296\\
1.678	-0.0189298\\
1.68	-0.01893\\
1.682	-0.0189302\\
1.684	-0.0189303\\
1.686	-0.0189305\\
1.688	-0.0189307\\
1.69	-0.0189308\\
1.692	-0.018931\\
1.694	-0.0189312\\
1.696	-0.0189313\\
1.698	-0.0189314\\
1.7	-0.0189316\\
1.702	-0.0189317\\
1.704	-0.0189318\\
1.706	-0.0189319\\
1.708	-0.018932\\
1.71	-0.0189321\\
1.712	-0.0189322\\
1.714	-0.0189322\\
1.716	-0.0189323\\
1.718	-0.0189323\\
1.72	-0.0189324\\
1.722	-0.0189324\\
1.724	-0.0189324\\
1.726	-0.0189324\\
1.728	-0.0189324\\
1.73	-0.0189324\\
1.732	-0.0189324\\
1.734	-0.0189324\\
1.736	-0.0189323\\
1.738	-0.0189323\\
1.74	-0.0189322\\
1.742	-0.0189321\\
1.744	-0.0189321\\
1.746	-0.018932\\
1.748	-0.0189319\\
1.75	-0.0189318\\
1.752	-0.0189317\\
1.754	-0.0189316\\
1.756	-0.0189315\\
1.758	-0.0189313\\
1.76	-0.0189312\\
1.762	-0.0189311\\
1.764	-0.0189309\\
1.766	-0.0189308\\
1.768	-0.0189307\\
1.77	-0.0189305\\
1.772	-0.0189304\\
1.774	-0.0189302\\
1.776	-0.0189301\\
1.778	-0.0189299\\
1.78	-0.0189297\\
1.782	-0.0189296\\
1.784	-0.0189294\\
1.786	-0.0189293\\
1.788	-0.0189291\\
1.79	-0.018929\\
1.792	-0.0189288\\
1.794	-0.0189287\\
1.796	-0.0189285\\
1.798	-0.0189284\\
1.8	-0.0189282\\
1.802	-0.0189281\\
1.804	-0.018928\\
1.806	-0.0189278\\
1.808	-0.0189277\\
1.81	-0.0189276\\
1.812	-0.0189275\\
1.814	-0.0189274\\
1.816	-0.0189273\\
1.818	-0.0189272\\
1.82	-0.0189271\\
1.822	-0.0189271\\
1.824	-0.018927\\
1.826	-0.0189269\\
1.828	-0.0189269\\
1.83	-0.0189268\\
1.832	-0.0189268\\
1.834	-0.0189268\\
1.836	-0.0189268\\
1.838	-0.0189267\\
1.84	-0.0189267\\
1.842	-0.0189267\\
1.844	-0.0189267\\
1.846	-0.0189268\\
1.848	-0.0189268\\
1.85	-0.0189268\\
1.852	-0.0189269\\
1.854	-0.0189269\\
1.856	-0.018927\\
1.858	-0.018927\\
1.86	-0.0189271\\
1.862	-0.0189271\\
1.864	-0.0189272\\
1.866	-0.0189273\\
1.868	-0.0189274\\
1.87	-0.0189275\\
1.872	-0.0189276\\
1.874	-0.0189277\\
1.876	-0.0189278\\
1.878	-0.0189279\\
1.88	-0.018928\\
1.882	-0.0189281\\
1.884	-0.0189282\\
1.886	-0.0189283\\
1.888	-0.0189284\\
1.89	-0.0189285\\
1.892	-0.0189287\\
1.894	-0.0189288\\
1.896	-0.0189289\\
1.898	-0.018929\\
1.9	-0.0189291\\
1.902	-0.0189292\\
1.904	-0.0189293\\
1.906	-0.0189294\\
1.908	-0.0189296\\
1.91	-0.0189297\\
1.912	-0.0189298\\
1.914	-0.0189299\\
1.916	-0.01893\\
1.918	-0.01893\\
1.92	-0.0189301\\
1.922	-0.0189302\\
1.924	-0.0189303\\
1.926	-0.0189304\\
1.928	-0.0189304\\
1.93	-0.0189305\\
1.932	-0.0189306\\
1.934	-0.0189306\\
1.936	-0.0189307\\
1.938	-0.0189307\\
1.94	-0.0189307\\
1.942	-0.0189308\\
1.944	-0.0189308\\
1.946	-0.0189308\\
1.948	-0.0189308\\
1.95	-0.0189308\\
1.952	-0.0189308\\
1.954	-0.0189308\\
1.956	-0.0189308\\
1.958	-0.0189308\\
1.96	-0.0189308\\
1.962	-0.0189308\\
1.964	-0.0189307\\
1.966	-0.0189307\\
1.968	-0.0189307\\
1.97	-0.0189306\\
1.972	-0.0189306\\
1.974	-0.0189305\\
1.976	-0.0189305\\
1.978	-0.0189304\\
1.98	-0.0189304\\
1.982	-0.0189303\\
1.984	-0.0189302\\
1.986	-0.0189302\\
1.988	-0.0189301\\
1.99	-0.01893\\
1.992	-0.0189299\\
1.994	-0.0189299\\
1.996	-0.0189298\\
1.998	-0.0189297\\
2	-0.0189296\\
};
\addplot [color=mycolor2, line width=1.0pt, forget plot]
  table[row sep=crcr]{%
-0.024	0\\
-0.022	0\\
-0.02	0\\
-0.018	0\\
-0.016	0\\
-0.014	0\\
-0.012	0\\
-0.01	0\\
-0.008	0\\
-0.006	0\\
-0.004	0\\
-0.002	0\\
0	0\\
0.002	-0.00183136\\
0.004	-0.00337796\\
0.006	-0.00470391\\
0.008	-0.00585754\\
0.01	-0.00687511\\
0.012	-0.00778382\\
0.014	-0.00860405\\
0.016	-0.00935112\\
0.018	-0.0100366\\
0.02	-0.0106694\\
0.022	-0.0112563\\
0.024	-0.0118026\\
0.026	-0.0123128\\
0.028	-0.0127901\\
0.03	-0.0132376\\
0.032	-0.0136577\\
0.034	-0.0140526\\
0.036	-0.014424\\
0.038	-0.0147737\\
0.04	-0.0151031\\
0.042	-0.0154138\\
0.044	-0.0157068\\
0.046	-0.0159835\\
0.048	-0.0162448\\
0.05	-0.016492\\
0.052	-0.0167259\\
0.054	-0.0169474\\
0.056	-0.0171576\\
0.058	-0.0173571\\
0.06	-0.0175467\\
0.062	-0.0177273\\
0.064	-0.0178993\\
0.066	-0.0180633\\
0.068	-0.0182201\\
0.07	-0.0183699\\
0.072	-0.0185132\\
0.074	-0.0186503\\
0.076	-0.0187816\\
0.078	-0.0189073\\
0.08	-0.0190276\\
0.082	-0.0191425\\
0.084	-0.0192524\\
0.086	-0.0193571\\
0.088	-0.0194568\\
0.09	-0.0195514\\
0.092	-0.019641\\
0.094	-0.0197254\\
0.096	-0.0198048\\
0.098	-0.0198789\\
0.1	-0.0199479\\
0.102	-0.0200115\\
0.104	-0.0200699\\
0.106	-0.0201228\\
0.108	-0.0201704\\
0.11	-0.0202126\\
0.112	-0.0202493\\
0.114	-0.0202807\\
0.116	-0.0203066\\
0.118	-0.0203273\\
0.12	-0.0203426\\
0.122	-0.0203527\\
0.124	-0.0203578\\
0.126	-0.0203578\\
0.128	-0.0203529\\
0.13	-0.0203433\\
0.132	-0.020329\\
0.134	-0.0203103\\
0.136	-0.0202874\\
0.138	-0.0202604\\
0.14	-0.0202295\\
0.142	-0.0201948\\
0.144	-0.0201567\\
0.146	-0.0201154\\
0.148	-0.020071\\
0.15	-0.0200237\\
0.152	-0.0199738\\
0.154	-0.0199216\\
0.156	-0.0198672\\
0.158	-0.0198108\\
0.16	-0.0197527\\
0.162	-0.0196932\\
0.164	-0.0196324\\
0.166	-0.0195705\\
0.168	-0.0195078\\
0.17	-0.0194444\\
0.172	-0.0193806\\
0.174	-0.0193166\\
0.176	-0.0192526\\
0.178	-0.0191887\\
0.18	-0.0191252\\
0.182	-0.0190623\\
0.184	-0.019\\
0.186	-0.0189386\\
0.188	-0.0188783\\
0.19	-0.0188191\\
0.192	-0.0187613\\
0.194	-0.018705\\
0.196	-0.0186503\\
0.198	-0.0185973\\
0.2	-0.0185462\\
0.202	-0.0184971\\
0.204	-0.01845\\
0.206	-0.0184051\\
0.208	-0.0183625\\
0.21	-0.0183222\\
0.212	-0.0182844\\
0.214	-0.018249\\
0.216	-0.0182161\\
0.218	-0.0181858\\
0.22	-0.0181581\\
0.222	-0.018133\\
0.224	-0.0181106\\
0.226	-0.0180908\\
0.228	-0.0180737\\
0.23	-0.0180593\\
0.232	-0.0180475\\
0.234	-0.0180383\\
0.236	-0.0180318\\
0.238	-0.0180278\\
0.24	-0.0180263\\
0.242	-0.0180273\\
0.244	-0.0180307\\
0.246	-0.0180366\\
0.248	-0.0180447\\
0.25	-0.018055\\
0.252	-0.0180675\\
0.254	-0.0180821\\
0.256	-0.0180987\\
0.258	-0.0181172\\
0.26	-0.0181375\\
0.262	-0.0181595\\
0.264	-0.0181832\\
0.266	-0.0182084\\
0.268	-0.0182351\\
0.27	-0.018263\\
0.272	-0.0182922\\
0.274	-0.0183225\\
0.276	-0.0183539\\
0.278	-0.0183861\\
0.28	-0.0184191\\
0.282	-0.0184528\\
0.284	-0.018487\\
0.286	-0.0185218\\
0.288	-0.0185568\\
0.29	-0.0185922\\
0.292	-0.0186277\\
0.294	-0.0186632\\
0.296	-0.0186987\\
0.298	-0.018734\\
0.3	-0.018769\\
0.302	-0.0188037\\
0.304	-0.018838\\
0.306	-0.0188717\\
0.308	-0.0189048\\
0.31	-0.0189372\\
0.312	-0.0189688\\
0.314	-0.0189996\\
0.316	-0.0190294\\
0.318	-0.0190583\\
0.32	-0.0190861\\
0.322	-0.0191128\\
0.324	-0.0191383\\
0.326	-0.0191626\\
0.328	-0.0191857\\
0.33	-0.0192074\\
0.332	-0.0192279\\
0.334	-0.019247\\
0.336	-0.0192647\\
0.338	-0.019281\\
0.34	-0.0192959\\
0.342	-0.0193094\\
0.344	-0.0193214\\
0.346	-0.019332\\
0.348	-0.0193412\\
0.35	-0.019349\\
0.352	-0.0193554\\
0.354	-0.0193604\\
0.356	-0.0193641\\
0.358	-0.0193664\\
0.36	-0.0193674\\
0.362	-0.0193671\\
0.364	-0.0193656\\
0.366	-0.0193628\\
0.368	-0.0193589\\
0.37	-0.0193538\\
0.372	-0.0193477\\
0.374	-0.0193405\\
0.376	-0.0193323\\
0.378	-0.0193232\\
0.38	-0.0193132\\
0.382	-0.0193023\\
0.384	-0.0192907\\
0.386	-0.0192783\\
0.388	-0.0192652\\
0.39	-0.0192515\\
0.392	-0.0192372\\
0.394	-0.0192224\\
0.396	-0.0192071\\
0.398	-0.0191915\\
0.4	-0.0191754\\
0.402	-0.0191591\\
0.404	-0.0191425\\
0.406	-0.0191257\\
0.408	-0.0191087\\
0.41	-0.0190917\\
0.412	-0.0190746\\
0.414	-0.0190575\\
0.416	-0.0190404\\
0.418	-0.0190234\\
0.42	-0.0190066\\
0.422	-0.0189899\\
0.424	-0.0189734\\
0.426	-0.0189572\\
0.428	-0.0189413\\
0.43	-0.0189257\\
0.432	-0.0189104\\
0.434	-0.0188956\\
0.436	-0.0188811\\
0.438	-0.0188671\\
0.44	-0.0188536\\
0.442	-0.0188406\\
0.444	-0.0188282\\
0.446	-0.0188162\\
0.448	-0.0188049\\
0.45	-0.0187941\\
0.452	-0.0187839\\
0.454	-0.0187743\\
0.456	-0.0187654\\
0.458	-0.0187571\\
0.46	-0.0187494\\
0.462	-0.0187424\\
0.464	-0.0187361\\
0.466	-0.0187304\\
0.468	-0.0187254\\
0.47	-0.018721\\
0.472	-0.0187173\\
0.474	-0.0187142\\
0.476	-0.0187119\\
0.478	-0.0187101\\
0.48	-0.018709\\
0.482	-0.0187085\\
0.484	-0.0187086\\
0.486	-0.0187094\\
0.488	-0.0187107\\
0.49	-0.0187126\\
0.492	-0.018715\\
0.494	-0.018718\\
0.496	-0.0187215\\
0.498	-0.0187255\\
0.5	-0.01873\\
0.502	-0.0187349\\
0.504	-0.0187403\\
0.506	-0.0187461\\
0.508	-0.0187522\\
0.51	-0.0187587\\
0.512	-0.0187656\\
0.514	-0.0187728\\
0.516	-0.0187802\\
0.518	-0.0187879\\
0.52	-0.0187958\\
0.522	-0.0188039\\
0.524	-0.0188122\\
0.526	-0.0188207\\
0.528	-0.0188293\\
0.53	-0.0188379\\
0.532	-0.0188466\\
0.534	-0.0188554\\
0.536	-0.0188642\\
0.538	-0.0188729\\
0.54	-0.0188817\\
0.542	-0.0188903\\
0.544	-0.0188989\\
0.546	-0.0189074\\
0.548	-0.0189158\\
0.55	-0.018924\\
0.552	-0.018932\\
0.554	-0.0189399\\
0.556	-0.0189475\\
0.558	-0.018955\\
0.56	-0.0189621\\
0.562	-0.0189691\\
0.564	-0.0189757\\
0.566	-0.0189821\\
0.568	-0.0189882\\
0.57	-0.018994\\
0.572	-0.0189995\\
0.574	-0.0190046\\
0.576	-0.0190095\\
0.578	-0.019014\\
0.58	-0.0190181\\
0.582	-0.019022\\
0.584	-0.0190254\\
0.586	-0.0190286\\
0.588	-0.0190314\\
0.59	-0.0190338\\
0.592	-0.0190359\\
0.594	-0.0190376\\
0.596	-0.0190391\\
0.598	-0.0190401\\
0.6	-0.0190409\\
0.602	-0.0190413\\
0.604	-0.0190414\\
0.606	-0.0190412\\
0.608	-0.0190407\\
0.61	-0.0190399\\
0.612	-0.0190389\\
0.614	-0.0190375\\
0.616	-0.0190359\\
0.618	-0.019034\\
0.62	-0.0190319\\
0.622	-0.0190296\\
0.624	-0.0190271\\
0.626	-0.0190243\\
0.628	-0.0190213\\
0.63	-0.0190182\\
0.632	-0.0190149\\
0.634	-0.0190115\\
0.636	-0.0190079\\
0.638	-0.0190042\\
0.64	-0.0190003\\
0.642	-0.0189964\\
0.644	-0.0189924\\
0.646	-0.0189883\\
0.648	-0.0189841\\
0.65	-0.0189799\\
0.652	-0.0189756\\
0.654	-0.0189713\\
0.656	-0.018967\\
0.658	-0.0189627\\
0.66	-0.0189585\\
0.662	-0.0189542\\
0.664	-0.0189499\\
0.666	-0.0189458\\
0.668	-0.0189416\\
0.67	-0.0189375\\
0.672	-0.0189335\\
0.674	-0.0189296\\
0.676	-0.0189257\\
0.678	-0.018922\\
0.68	-0.0189184\\
0.682	-0.0189148\\
0.684	-0.0189114\\
0.686	-0.0189081\\
0.688	-0.018905\\
0.69	-0.018902\\
0.692	-0.0188991\\
0.694	-0.0188964\\
0.696	-0.0188938\\
0.698	-0.0188913\\
0.7	-0.0188891\\
0.702	-0.0188869\\
0.704	-0.018885\\
0.706	-0.0188832\\
0.708	-0.0188816\\
0.71	-0.0188801\\
0.712	-0.0188788\\
0.714	-0.0188776\\
0.716	-0.0188767\\
0.718	-0.0188758\\
0.72	-0.0188752\\
0.722	-0.0188747\\
0.724	-0.0188743\\
0.726	-0.0188742\\
0.728	-0.0188741\\
0.73	-0.0188742\\
0.732	-0.0188745\\
0.734	-0.0188749\\
0.736	-0.0188754\\
0.738	-0.0188761\\
0.74	-0.0188769\\
0.742	-0.0188778\\
0.744	-0.0188788\\
0.746	-0.01888\\
0.748	-0.0188812\\
0.75	-0.0188826\\
0.752	-0.018884\\
0.754	-0.0188855\\
0.756	-0.0188871\\
0.758	-0.0188888\\
0.76	-0.0188906\\
0.762	-0.0188924\\
0.764	-0.0188943\\
0.766	-0.0188962\\
0.768	-0.0188982\\
0.77	-0.0189002\\
0.772	-0.0189023\\
0.774	-0.0189043\\
0.776	-0.0189064\\
0.778	-0.0189085\\
0.78	-0.0189106\\
0.782	-0.0189128\\
0.784	-0.0189149\\
0.786	-0.018917\\
0.788	-0.0189191\\
0.79	-0.0189211\\
0.792	-0.0189232\\
0.794	-0.0189252\\
0.796	-0.0189271\\
0.798	-0.0189291\\
0.8	-0.018931\\
0.802	-0.0189328\\
0.804	-0.0189346\\
0.806	-0.0189363\\
0.808	-0.018938\\
0.81	-0.0189396\\
0.812	-0.0189412\\
0.814	-0.0189427\\
0.816	-0.0189441\\
0.818	-0.0189454\\
0.82	-0.0189467\\
0.822	-0.0189479\\
0.824	-0.018949\\
0.826	-0.01895\\
0.828	-0.018951\\
0.83	-0.0189518\\
0.832	-0.0189526\\
0.834	-0.0189533\\
0.836	-0.018954\\
0.838	-0.0189545\\
0.84	-0.018955\\
0.842	-0.0189554\\
0.844	-0.0189557\\
0.846	-0.0189559\\
0.848	-0.0189561\\
0.85	-0.0189562\\
0.852	-0.0189562\\
0.854	-0.0189561\\
0.856	-0.018956\\
0.858	-0.0189558\\
0.86	-0.0189555\\
0.862	-0.0189552\\
0.864	-0.0189548\\
0.866	-0.0189543\\
0.868	-0.0189538\\
0.87	-0.0189532\\
0.872	-0.0189526\\
0.874	-0.0189519\\
0.876	-0.0189512\\
0.878	-0.0189505\\
0.88	-0.0189497\\
0.882	-0.0189489\\
0.884	-0.018948\\
0.886	-0.0189471\\
0.888	-0.0189462\\
0.89	-0.0189452\\
0.892	-0.0189443\\
0.894	-0.0189433\\
0.896	-0.0189423\\
0.898	-0.0189413\\
0.9	-0.0189402\\
0.902	-0.0189392\\
0.904	-0.0189382\\
0.906	-0.0189372\\
0.908	-0.0189361\\
0.91	-0.0189351\\
0.912	-0.0189341\\
0.914	-0.0189331\\
0.916	-0.0189321\\
0.918	-0.0189311\\
0.92	-0.0189301\\
0.922	-0.0189292\\
0.924	-0.0189283\\
0.926	-0.0189274\\
0.928	-0.0189265\\
0.93	-0.0189256\\
0.932	-0.0189248\\
0.934	-0.018924\\
0.936	-0.0189233\\
0.938	-0.0189225\\
0.94	-0.0189218\\
0.942	-0.0189212\\
0.944	-0.0189206\\
0.946	-0.01892\\
0.948	-0.0189194\\
0.95	-0.0189189\\
0.952	-0.0189184\\
0.954	-0.018918\\
0.956	-0.0189176\\
0.958	-0.0189173\\
0.96	-0.0189169\\
0.962	-0.0189167\\
0.964	-0.0189164\\
0.966	-0.0189162\\
0.968	-0.0189161\\
0.97	-0.018916\\
0.972	-0.0189159\\
0.974	-0.0189158\\
0.976	-0.0189158\\
0.978	-0.0189159\\
0.98	-0.0189159\\
0.982	-0.018916\\
0.984	-0.0189161\\
0.986	-0.0189163\\
0.988	-0.0189165\\
0.99	-0.0189167\\
0.992	-0.018917\\
0.994	-0.0189172\\
0.996	-0.0189176\\
0.998	-0.0189179\\
1	-0.0189182\\
1.002	-0.0189186\\
1.004	-0.018919\\
1.006	-0.0189194\\
1.008	-0.0189198\\
1.01	-0.0189203\\
1.012	-0.0189207\\
1.014	-0.0189212\\
1.016	-0.0189217\\
1.018	-0.0189222\\
1.02	-0.0189227\\
1.022	-0.0189232\\
1.024	-0.0189237\\
1.026	-0.0189242\\
1.028	-0.0189247\\
1.03	-0.0189252\\
1.032	-0.0189258\\
1.034	-0.0189263\\
1.036	-0.0189268\\
1.038	-0.0189273\\
1.04	-0.0189278\\
1.042	-0.0189283\\
1.044	-0.0189288\\
1.046	-0.0189292\\
1.048	-0.0189297\\
1.05	-0.0189302\\
1.052	-0.0189306\\
1.054	-0.018931\\
1.056	-0.0189314\\
1.058	-0.0189318\\
1.06	-0.0189322\\
1.062	-0.0189326\\
1.064	-0.0189329\\
1.066	-0.0189332\\
1.068	-0.0189335\\
1.07	-0.0189338\\
1.072	-0.0189341\\
1.074	-0.0189344\\
1.076	-0.0189346\\
1.078	-0.0189348\\
1.08	-0.018935\\
1.082	-0.0189352\\
1.084	-0.0189353\\
1.086	-0.0189354\\
1.088	-0.0189356\\
1.09	-0.0189356\\
1.092	-0.0189357\\
1.094	-0.0189358\\
1.096	-0.0189358\\
1.098	-0.0189358\\
1.1	-0.0189358\\
1.102	-0.0189358\\
1.104	-0.0189358\\
1.106	-0.0189357\\
1.108	-0.0189356\\
1.11	-0.0189355\\
1.112	-0.0189354\\
1.114	-0.0189353\\
1.116	-0.0189352\\
1.118	-0.018935\\
1.12	-0.0189349\\
1.122	-0.0189347\\
1.124	-0.0189345\\
1.126	-0.0189343\\
1.128	-0.0189341\\
1.13	-0.0189339\\
1.132	-0.0189337\\
1.134	-0.0189335\\
1.136	-0.0189332\\
1.138	-0.018933\\
1.14	-0.0189327\\
1.142	-0.0189325\\
1.144	-0.0189322\\
1.146	-0.018932\\
1.148	-0.0189317\\
1.15	-0.0189315\\
1.152	-0.0189312\\
1.154	-0.0189309\\
1.156	-0.0189307\\
1.158	-0.0189304\\
1.16	-0.0189302\\
1.162	-0.0189299\\
1.164	-0.0189297\\
1.166	-0.0189294\\
1.168	-0.0189292\\
1.17	-0.018929\\
1.172	-0.0189287\\
1.174	-0.0189285\\
1.176	-0.0189283\\
1.178	-0.0189281\\
1.18	-0.0189279\\
1.182	-0.0189277\\
1.184	-0.0189275\\
1.186	-0.0189274\\
1.188	-0.0189272\\
1.19	-0.018927\\
1.192	-0.0189269\\
1.194	-0.0189267\\
1.196	-0.0189266\\
1.198	-0.0189265\\
1.2	-0.0189264\\
1.202	-0.0189263\\
1.204	-0.0189262\\
1.206	-0.0189261\\
1.208	-0.0189261\\
1.21	-0.018926\\
1.212	-0.018926\\
1.214	-0.0189259\\
1.216	-0.0189259\\
1.218	-0.0189259\\
1.22	-0.0189259\\
1.222	-0.0189259\\
1.224	-0.0189259\\
1.226	-0.0189259\\
1.228	-0.0189259\\
1.23	-0.018926\\
1.232	-0.018926\\
1.234	-0.0189261\\
1.236	-0.0189261\\
1.238	-0.0189262\\
1.24	-0.0189263\\
1.242	-0.0189263\\
1.244	-0.0189264\\
1.246	-0.0189265\\
1.248	-0.0189266\\
1.25	-0.0189267\\
1.252	-0.0189268\\
1.254	-0.0189269\\
1.256	-0.018927\\
1.258	-0.0189271\\
1.26	-0.0189273\\
1.262	-0.0189274\\
1.264	-0.0189275\\
1.266	-0.0189276\\
1.268	-0.0189277\\
1.27	-0.0189279\\
1.272	-0.018928\\
1.274	-0.0189281\\
1.276	-0.0189282\\
1.278	-0.0189284\\
1.28	-0.0189285\\
1.282	-0.0189286\\
1.284	-0.0189287\\
1.286	-0.0189289\\
1.288	-0.018929\\
1.29	-0.0189291\\
1.292	-0.0189292\\
1.294	-0.0189293\\
1.296	-0.0189294\\
1.298	-0.0189295\\
1.3	-0.0189296\\
1.302	-0.0189297\\
1.304	-0.0189298\\
1.306	-0.0189299\\
1.308	-0.01893\\
1.31	-0.0189301\\
1.312	-0.0189302\\
1.314	-0.0189302\\
1.316	-0.0189303\\
1.318	-0.0189304\\
1.32	-0.0189304\\
1.322	-0.0189305\\
1.324	-0.0189305\\
1.326	-0.0189306\\
1.328	-0.0189306\\
1.33	-0.0189306\\
1.332	-0.0189307\\
1.334	-0.0189307\\
1.336	-0.0189307\\
1.338	-0.0189307\\
1.34	-0.0189307\\
1.342	-0.0189307\\
1.344	-0.0189307\\
1.346	-0.0189307\\
1.348	-0.0189307\\
1.35	-0.0189307\\
1.352	-0.0189307\\
1.354	-0.0189307\\
1.356	-0.0189307\\
1.358	-0.0189306\\
1.36	-0.0189306\\
1.362	-0.0189306\\
1.364	-0.0189305\\
1.366	-0.0189305\\
1.368	-0.0189304\\
1.37	-0.0189304\\
1.372	-0.0189304\\
1.374	-0.0189303\\
1.376	-0.0189303\\
1.378	-0.0189302\\
1.38	-0.0189301\\
1.382	-0.0189301\\
1.384	-0.01893\\
1.386	-0.01893\\
1.388	-0.0189299\\
1.39	-0.0189299\\
1.392	-0.0189298\\
1.394	-0.0189297\\
1.396	-0.0189297\\
1.398	-0.0189296\\
1.4	-0.0189296\\
1.402	-0.0189295\\
1.404	-0.0189294\\
1.406	-0.0189294\\
1.408	-0.0189293\\
1.41	-0.0189293\\
1.412	-0.0189292\\
1.414	-0.0189292\\
1.416	-0.0189291\\
1.418	-0.0189291\\
1.42	-0.018929\\
1.422	-0.018929\\
1.424	-0.0189289\\
1.426	-0.0189289\\
1.428	-0.0189288\\
1.43	-0.0189288\\
1.432	-0.0189287\\
1.434	-0.0189287\\
1.436	-0.0189287\\
1.438	-0.0189286\\
1.44	-0.0189286\\
1.442	-0.0189286\\
1.444	-0.0189285\\
1.446	-0.0189285\\
1.448	-0.0189285\\
1.45	-0.0189285\\
1.452	-0.0189285\\
1.454	-0.0189284\\
1.456	-0.0189284\\
1.458	-0.0189284\\
1.46	-0.0189284\\
1.462	-0.0189284\\
1.464	-0.0189284\\
1.466	-0.0189284\\
1.468	-0.0189284\\
1.47	-0.0189284\\
1.472	-0.0189284\\
1.474	-0.0189284\\
1.476	-0.0189284\\
1.478	-0.0189284\\
1.48	-0.0189284\\
1.482	-0.0189285\\
1.484	-0.0189285\\
1.486	-0.0189285\\
1.488	-0.0189285\\
1.49	-0.0189285\\
1.492	-0.0189285\\
1.494	-0.0189286\\
1.496	-0.0189286\\
1.498	-0.0189286\\
1.5	-0.0189286\\
1.502	-0.0189287\\
1.504	-0.0189287\\
1.506	-0.0189287\\
1.508	-0.0189287\\
1.51	-0.0189288\\
1.512	-0.0189288\\
1.514	-0.0189288\\
1.516	-0.0189288\\
1.518	-0.0189289\\
1.52	-0.0189289\\
1.522	-0.0189289\\
1.524	-0.018929\\
1.526	-0.018929\\
1.528	-0.018929\\
1.53	-0.018929\\
1.532	-0.0189291\\
1.534	-0.0189291\\
1.536	-0.0189291\\
1.538	-0.0189292\\
1.54	-0.0189292\\
1.542	-0.0189292\\
1.544	-0.0189292\\
1.546	-0.0189292\\
1.548	-0.0189293\\
1.55	-0.0189293\\
1.552	-0.0189293\\
1.554	-0.0189293\\
1.556	-0.0189294\\
1.558	-0.0189294\\
1.56	-0.0189294\\
1.562	-0.0189294\\
1.564	-0.0189294\\
1.566	-0.0189294\\
1.568	-0.0189294\\
1.57	-0.0189295\\
1.572	-0.0189295\\
1.574	-0.0189295\\
1.576	-0.0189295\\
1.578	-0.0189295\\
1.58	-0.0189295\\
1.582	-0.0189295\\
1.584	-0.0189295\\
1.586	-0.0189295\\
1.588	-0.0189295\\
1.59	-0.0189295\\
1.592	-0.0189295\\
1.594	-0.0189295\\
1.596	-0.0189295\\
1.598	-0.0189295\\
1.6	-0.0189295\\
1.602	-0.0189295\\
1.604	-0.0189295\\
1.606	-0.0189295\\
1.608	-0.0189295\\
1.61	-0.0189295\\
1.612	-0.0189295\\
1.614	-0.0189295\\
1.616	-0.0189295\\
1.618	-0.0189294\\
1.62	-0.0189294\\
1.622	-0.0189294\\
1.624	-0.0189294\\
1.626	-0.0189294\\
1.628	-0.0189294\\
1.63	-0.0189294\\
1.632	-0.0189294\\
1.634	-0.0189294\\
1.636	-0.0189293\\
1.638	-0.0189293\\
1.64	-0.0189293\\
1.642	-0.0189293\\
1.644	-0.0189293\\
1.646	-0.0189293\\
1.648	-0.0189293\\
1.65	-0.0189292\\
1.652	-0.0189292\\
1.654	-0.0189292\\
1.656	-0.0189292\\
1.658	-0.0189292\\
1.66	-0.0189292\\
1.662	-0.0189292\\
1.664	-0.0189292\\
1.666	-0.0189291\\
1.668	-0.0189291\\
1.67	-0.0189291\\
1.672	-0.0189291\\
1.674	-0.0189291\\
1.676	-0.0189291\\
1.678	-0.0189291\\
1.68	-0.0189291\\
1.682	-0.0189291\\
1.684	-0.018929\\
1.686	-0.018929\\
1.688	-0.018929\\
1.69	-0.018929\\
1.692	-0.018929\\
1.694	-0.018929\\
1.696	-0.018929\\
1.698	-0.018929\\
1.7	-0.018929\\
1.702	-0.018929\\
1.704	-0.018929\\
1.706	-0.018929\\
1.708	-0.018929\\
1.71	-0.018929\\
1.712	-0.018929\\
1.714	-0.018929\\
1.716	-0.018929\\
1.718	-0.018929\\
1.72	-0.018929\\
1.722	-0.018929\\
1.724	-0.018929\\
1.726	-0.018929\\
1.728	-0.018929\\
1.73	-0.018929\\
1.732	-0.018929\\
1.734	-0.018929\\
1.736	-0.018929\\
1.738	-0.018929\\
1.74	-0.018929\\
1.742	-0.018929\\
1.744	-0.018929\\
1.746	-0.018929\\
1.748	-0.018929\\
1.75	-0.018929\\
1.752	-0.018929\\
1.754	-0.018929\\
1.756	-0.018929\\
1.758	-0.0189291\\
1.76	-0.0189291\\
1.762	-0.0189291\\
1.764	-0.0189291\\
1.766	-0.0189291\\
1.768	-0.0189291\\
1.77	-0.0189291\\
1.772	-0.0189291\\
1.774	-0.0189291\\
1.776	-0.0189291\\
1.778	-0.0189291\\
1.78	-0.0189291\\
1.782	-0.0189291\\
1.784	-0.0189291\\
1.786	-0.0189292\\
1.788	-0.0189292\\
1.79	-0.0189292\\
1.792	-0.0189292\\
1.794	-0.0189292\\
1.796	-0.0189292\\
1.798	-0.0189292\\
1.8	-0.0189292\\
1.802	-0.0189292\\
1.804	-0.0189292\\
1.806	-0.0189292\\
1.808	-0.0189292\\
1.81	-0.0189292\\
1.812	-0.0189292\\
1.814	-0.0189292\\
1.816	-0.0189292\\
1.818	-0.0189292\\
1.82	-0.0189292\\
1.822	-0.0189292\\
1.824	-0.0189292\\
1.826	-0.0189292\\
1.828	-0.0189292\\
1.83	-0.0189292\\
1.832	-0.0189292\\
1.834	-0.0189292\\
1.836	-0.0189292\\
1.838	-0.0189292\\
1.84	-0.0189292\\
1.842	-0.0189292\\
1.844	-0.0189292\\
1.846	-0.0189292\\
1.848	-0.0189292\\
1.85	-0.0189292\\
1.852	-0.0189292\\
1.854	-0.0189292\\
1.856	-0.0189292\\
1.858	-0.0189292\\
1.86	-0.0189292\\
1.862	-0.0189292\\
1.864	-0.0189292\\
1.866	-0.0189292\\
1.868	-0.0189292\\
1.87	-0.0189292\\
1.872	-0.0189292\\
1.874	-0.0189292\\
1.876	-0.0189292\\
1.878	-0.0189292\\
1.88	-0.0189292\\
1.882	-0.0189292\\
1.884	-0.0189292\\
1.886	-0.0189292\\
1.888	-0.0189292\\
1.89	-0.0189292\\
1.892	-0.0189292\\
1.894	-0.0189292\\
1.896	-0.0189292\\
1.898	-0.0189292\\
1.9	-0.0189292\\
1.902	-0.0189292\\
1.904	-0.0189292\\
1.906	-0.0189292\\
1.908	-0.0189292\\
1.91	-0.0189292\\
1.912	-0.0189292\\
1.914	-0.0189292\\
1.916	-0.0189291\\
1.918	-0.0189291\\
1.92	-0.0189291\\
1.922	-0.0189291\\
1.924	-0.0189291\\
1.926	-0.0189291\\
1.928	-0.0189291\\
1.93	-0.0189291\\
1.932	-0.0189291\\
1.934	-0.0189291\\
1.936	-0.0189291\\
1.938	-0.0189291\\
1.94	-0.0189291\\
1.942	-0.0189291\\
1.944	-0.0189291\\
1.946	-0.0189291\\
1.948	-0.0189291\\
1.95	-0.0189291\\
1.952	-0.0189291\\
1.954	-0.0189291\\
1.956	-0.0189291\\
1.958	-0.0189291\\
1.96	-0.0189291\\
1.962	-0.0189291\\
1.964	-0.0189291\\
1.966	-0.0189291\\
1.968	-0.0189291\\
1.97	-0.0189291\\
1.972	-0.0189291\\
1.974	-0.0189291\\
1.976	-0.0189291\\
1.978	-0.0189291\\
1.98	-0.0189291\\
1.982	-0.0189291\\
1.984	-0.0189291\\
1.986	-0.0189291\\
1.988	-0.0189291\\
1.99	-0.0189291\\
1.992	-0.0189291\\
1.994	-0.0189291\\
1.996	-0.0189291\\
1.998	-0.0189291\\
2	-0.0189291\\
};
\addplot [color=mycolor3, line width=1.0pt, forget plot]
  table[row sep=crcr]{%
-0.024	0\\
-0.022	0\\
-0.02	0\\
-0.018	0\\
-0.016	0\\
-0.014	0\\
-0.012	0\\
-0.01	0\\
-0.008	0\\
-0.006	0\\
-0.004	0\\
-0.002	0\\
0	0\\
0.002	-0.00183136\\
0.004	-0.00337792\\
0.006	-0.00470378\\
0.008	-0.00585724\\
0.01	-0.00687453\\
0.012	-0.00778283\\
0.014	-0.00860249\\
0.016	-0.00934882\\
0.018	-0.0100334\\
0.02	-0.010665\\
0.022	-0.0112505\\
0.024	-0.0117952\\
0.026	-0.0123034\\
0.028	-0.0127786\\
0.03	-0.0132236\\
0.032	-0.0136409\\
0.034	-0.0140326\\
0.036	-0.0144005\\
0.038	-0.0147464\\
0.04	-0.0150717\\
0.042	-0.0153778\\
0.044	-0.0156659\\
0.046	-0.0159373\\
0.048	-0.0161931\\
0.05	-0.0164343\\
0.052	-0.0166619\\
0.054	-0.0168769\\
0.056	-0.0170801\\
0.058	-0.0172724\\
0.06	-0.0174546\\
0.062	-0.0176274\\
0.064	-0.0177914\\
0.066	-0.0179474\\
0.068	-0.0180958\\
0.07	-0.0182372\\
0.072	-0.0183719\\
0.074	-0.0185005\\
0.076	-0.0186232\\
0.078	-0.0187403\\
0.08	-0.018852\\
0.082	-0.0189585\\
0.084	-0.0190601\\
0.086	-0.0191567\\
0.088	-0.0192485\\
0.09	-0.0193355\\
0.092	-0.0194178\\
0.094	-0.0194953\\
0.096	-0.0195681\\
0.098	-0.0196362\\
0.1	-0.0196995\\
0.102	-0.0197582\\
0.104	-0.019812\\
0.106	-0.0198611\\
0.108	-0.0199055\\
0.11	-0.0199451\\
0.112	-0.01998\\
0.114	-0.0200103\\
0.116	-0.0200359\\
0.118	-0.020057\\
0.12	-0.0200736\\
0.122	-0.0200858\\
0.124	-0.0200938\\
0.126	-0.0200976\\
0.128	-0.0200974\\
0.13	-0.0200933\\
0.132	-0.0200855\\
0.134	-0.0200741\\
0.136	-0.0200592\\
0.138	-0.0200412\\
0.14	-0.0200201\\
0.142	-0.0199961\\
0.144	-0.0199694\\
0.146	-0.0199403\\
0.148	-0.0199088\\
0.15	-0.0198753\\
0.152	-0.0198398\\
0.154	-0.0198026\\
0.156	-0.0197638\\
0.158	-0.0197237\\
0.16	-0.0196823\\
0.162	-0.01964\\
0.164	-0.0195969\\
0.166	-0.0195531\\
0.168	-0.0195087\\
0.17	-0.0194641\\
0.172	-0.0194192\\
0.174	-0.0193743\\
0.176	-0.0193295\\
0.178	-0.0192849\\
0.18	-0.0192406\\
0.182	-0.0191969\\
0.184	-0.0191537\\
0.186	-0.0191112\\
0.188	-0.0190696\\
0.19	-0.0190288\\
0.192	-0.0189891\\
0.194	-0.0189504\\
0.196	-0.0189129\\
0.198	-0.0188767\\
0.2	-0.0188418\\
0.202	-0.0188082\\
0.204	-0.0187761\\
0.206	-0.0187454\\
0.208	-0.0187163\\
0.21	-0.0186887\\
0.212	-0.0186627\\
0.214	-0.0186384\\
0.216	-0.0186157\\
0.218	-0.0185946\\
0.22	-0.0185752\\
0.222	-0.0185575\\
0.224	-0.0185414\\
0.226	-0.018527\\
0.228	-0.0185143\\
0.23	-0.0185032\\
0.232	-0.0184937\\
0.234	-0.0184858\\
0.236	-0.0184794\\
0.238	-0.0184746\\
0.24	-0.0184713\\
0.242	-0.0184694\\
0.244	-0.0184689\\
0.246	-0.0184698\\
0.248	-0.018472\\
0.25	-0.0184755\\
0.252	-0.0184801\\
0.254	-0.018486\\
0.256	-0.0184929\\
0.258	-0.0185009\\
0.26	-0.0185098\\
0.262	-0.0185197\\
0.264	-0.0185305\\
0.266	-0.018542\\
0.268	-0.0185543\\
0.27	-0.0185673\\
0.272	-0.0185809\\
0.274	-0.018595\\
0.276	-0.0186096\\
0.278	-0.0186247\\
0.28	-0.0186402\\
0.282	-0.0186559\\
0.284	-0.0186719\\
0.286	-0.0186881\\
0.288	-0.0187044\\
0.29	-0.0187209\\
0.292	-0.0187373\\
0.294	-0.0187537\\
0.296	-0.01877\\
0.298	-0.0187862\\
0.3	-0.0188022\\
0.302	-0.018818\\
0.304	-0.0188335\\
0.306	-0.0188488\\
0.308	-0.0188636\\
0.31	-0.0188781\\
0.312	-0.0188921\\
0.314	-0.0189057\\
0.316	-0.0189188\\
0.318	-0.0189314\\
0.32	-0.0189435\\
0.322	-0.018955\\
0.324	-0.0189659\\
0.326	-0.0189762\\
0.328	-0.0189859\\
0.33	-0.018995\\
0.332	-0.0190035\\
0.334	-0.0190114\\
0.336	-0.0190186\\
0.338	-0.0190252\\
0.34	-0.0190311\\
0.342	-0.0190364\\
0.344	-0.0190411\\
0.346	-0.0190452\\
0.348	-0.0190487\\
0.35	-0.0190516\\
0.352	-0.0190539\\
0.354	-0.0190557\\
0.356	-0.019057\\
0.358	-0.0190577\\
0.36	-0.0190579\\
0.362	-0.0190576\\
0.364	-0.0190569\\
0.366	-0.0190558\\
0.368	-0.0190543\\
0.37	-0.0190524\\
0.372	-0.0190501\\
0.374	-0.0190475\\
0.376	-0.0190446\\
0.378	-0.0190414\\
0.38	-0.0190379\\
0.382	-0.0190342\\
0.384	-0.0190303\\
0.386	-0.0190263\\
0.388	-0.019022\\
0.39	-0.0190176\\
0.392	-0.0190131\\
0.394	-0.0190086\\
0.396	-0.0190039\\
0.398	-0.0189992\\
0.4	-0.0189944\\
0.402	-0.0189896\\
0.404	-0.0189848\\
0.406	-0.01898\\
0.408	-0.0189753\\
0.41	-0.0189706\\
0.412	-0.0189659\\
0.414	-0.0189613\\
0.416	-0.0189568\\
0.418	-0.0189523\\
0.42	-0.018948\\
0.422	-0.0189437\\
0.424	-0.0189396\\
0.426	-0.0189356\\
0.428	-0.0189317\\
0.43	-0.0189279\\
0.432	-0.0189243\\
0.434	-0.0189208\\
0.436	-0.0189174\\
0.438	-0.0189142\\
0.44	-0.0189111\\
0.442	-0.0189082\\
0.444	-0.0189055\\
0.446	-0.0189029\\
0.448	-0.0189004\\
0.45	-0.0188981\\
0.452	-0.018896\\
0.454	-0.018894\\
0.456	-0.0188922\\
0.458	-0.0188905\\
0.46	-0.018889\\
0.462	-0.0188876\\
0.464	-0.0188864\\
0.466	-0.0188853\\
0.468	-0.0188844\\
0.47	-0.0188836\\
0.472	-0.018883\\
0.474	-0.0188825\\
0.476	-0.0188822\\
0.478	-0.018882\\
0.48	-0.0188819\\
0.482	-0.018882\\
0.484	-0.0188821\\
0.486	-0.0188825\\
0.488	-0.0188829\\
0.49	-0.0188834\\
0.492	-0.0188841\\
0.494	-0.0188849\\
0.496	-0.0188858\\
0.498	-0.0188867\\
0.5	-0.0188878\\
0.502	-0.018889\\
0.504	-0.0188902\\
0.506	-0.0188915\\
0.508	-0.0188929\\
0.51	-0.0188944\\
0.512	-0.0188959\\
0.514	-0.0188975\\
0.516	-0.0188991\\
0.518	-0.0189008\\
0.52	-0.0189025\\
0.522	-0.0189042\\
0.524	-0.018906\\
0.526	-0.0189078\\
0.528	-0.0189096\\
0.53	-0.0189114\\
0.532	-0.0189132\\
0.534	-0.018915\\
0.536	-0.0189168\\
0.538	-0.0189186\\
0.54	-0.0189204\\
0.542	-0.0189221\\
0.544	-0.0189238\\
0.546	-0.0189255\\
0.548	-0.0189272\\
0.55	-0.0189288\\
0.552	-0.0189303\\
0.554	-0.0189318\\
0.556	-0.0189332\\
0.558	-0.0189346\\
0.56	-0.0189359\\
0.562	-0.0189372\\
0.564	-0.0189384\\
0.566	-0.0189395\\
0.568	-0.0189406\\
0.57	-0.0189415\\
0.572	-0.0189425\\
0.574	-0.0189433\\
0.576	-0.0189441\\
0.578	-0.0189448\\
0.58	-0.0189454\\
0.582	-0.0189459\\
0.584	-0.0189464\\
0.586	-0.0189468\\
0.588	-0.0189472\\
0.59	-0.0189475\\
0.592	-0.0189477\\
0.594	-0.0189478\\
0.596	-0.0189479\\
0.598	-0.0189479\\
0.6	-0.0189479\\
0.602	-0.0189478\\
0.604	-0.0189476\\
0.606	-0.0189474\\
0.608	-0.0189472\\
0.61	-0.0189469\\
0.612	-0.0189466\\
0.614	-0.0189462\\
0.616	-0.0189458\\
0.618	-0.0189453\\
0.62	-0.0189449\\
0.622	-0.0189444\\
0.624	-0.0189438\\
0.626	-0.0189433\\
0.628	-0.0189427\\
0.63	-0.0189421\\
0.632	-0.0189415\\
0.634	-0.0189409\\
0.636	-0.0189402\\
0.638	-0.0189396\\
0.64	-0.0189389\\
0.642	-0.0189383\\
0.644	-0.0189376\\
0.646	-0.018937\\
0.648	-0.0189363\\
0.65	-0.0189357\\
0.652	-0.0189351\\
0.654	-0.0189344\\
0.656	-0.0189338\\
0.658	-0.0189332\\
0.66	-0.0189326\\
0.662	-0.018932\\
0.664	-0.0189314\\
0.666	-0.0189308\\
0.668	-0.0189303\\
0.67	-0.0189298\\
0.672	-0.0189292\\
0.674	-0.0189287\\
0.676	-0.0189283\\
0.678	-0.0189278\\
0.68	-0.0189273\\
0.682	-0.0189269\\
0.684	-0.0189265\\
0.686	-0.0189261\\
0.688	-0.0189258\\
0.69	-0.0189254\\
0.692	-0.0189251\\
0.694	-0.0189248\\
0.696	-0.0189245\\
0.698	-0.0189243\\
0.7	-0.018924\\
0.702	-0.0189238\\
0.704	-0.0189236\\
0.706	-0.0189235\\
0.708	-0.0189233\\
0.71	-0.0189232\\
0.712	-0.0189231\\
0.714	-0.018923\\
0.716	-0.0189229\\
0.718	-0.0189228\\
0.72	-0.0189228\\
0.722	-0.0189228\\
0.724	-0.0189228\\
0.726	-0.0189228\\
0.728	-0.0189228\\
0.73	-0.0189229\\
0.732	-0.0189229\\
0.734	-0.018923\\
0.736	-0.0189231\\
0.738	-0.0189232\\
0.74	-0.0189233\\
0.742	-0.0189235\\
0.744	-0.0189236\\
0.746	-0.0189238\\
0.748	-0.0189239\\
0.75	-0.0189241\\
0.752	-0.0189243\\
0.754	-0.0189245\\
0.756	-0.0189247\\
0.758	-0.0189249\\
0.76	-0.0189251\\
0.762	-0.0189253\\
0.764	-0.0189255\\
0.766	-0.0189257\\
0.768	-0.018926\\
0.77	-0.0189262\\
0.772	-0.0189264\\
0.774	-0.0189266\\
0.776	-0.0189269\\
0.778	-0.0189271\\
0.78	-0.0189273\\
0.782	-0.0189275\\
0.784	-0.0189278\\
0.786	-0.018928\\
0.788	-0.0189282\\
0.79	-0.0189284\\
0.792	-0.0189286\\
0.794	-0.0189288\\
0.796	-0.018929\\
0.798	-0.0189292\\
0.8	-0.0189294\\
0.802	-0.0189296\\
0.804	-0.0189297\\
0.806	-0.0189299\\
0.808	-0.01893\\
0.81	-0.0189302\\
0.812	-0.0189303\\
0.814	-0.0189304\\
0.816	-0.0189306\\
0.818	-0.0189307\\
0.82	-0.0189308\\
0.822	-0.0189309\\
0.824	-0.018931\\
0.826	-0.018931\\
0.828	-0.0189311\\
0.83	-0.0189312\\
0.832	-0.0189312\\
0.834	-0.0189313\\
0.836	-0.0189313\\
0.838	-0.0189313\\
0.84	-0.0189314\\
0.842	-0.0189314\\
0.844	-0.0189314\\
0.846	-0.0189314\\
0.848	-0.0189314\\
0.85	-0.0189314\\
0.852	-0.0189314\\
0.854	-0.0189313\\
0.856	-0.0189313\\
0.858	-0.0189313\\
0.86	-0.0189312\\
0.862	-0.0189312\\
0.864	-0.0189311\\
0.866	-0.0189311\\
0.868	-0.018931\\
0.87	-0.018931\\
0.872	-0.0189309\\
0.874	-0.0189309\\
0.876	-0.0189308\\
0.878	-0.0189307\\
0.88	-0.0189307\\
0.882	-0.0189306\\
0.884	-0.0189305\\
0.886	-0.0189304\\
0.888	-0.0189304\\
0.89	-0.0189303\\
0.892	-0.0189302\\
0.894	-0.0189301\\
0.896	-0.0189301\\
0.898	-0.01893\\
0.9	-0.0189299\\
0.902	-0.0189298\\
0.904	-0.0189298\\
0.906	-0.0189297\\
0.908	-0.0189296\\
0.91	-0.0189295\\
0.912	-0.0189295\\
0.914	-0.0189294\\
0.916	-0.0189293\\
0.918	-0.0189293\\
0.92	-0.0189292\\
0.922	-0.0189292\\
0.924	-0.0189291\\
0.926	-0.018929\\
0.928	-0.018929\\
0.93	-0.0189289\\
0.932	-0.0189289\\
0.934	-0.0189289\\
0.936	-0.0189288\\
0.938	-0.0189288\\
0.94	-0.0189287\\
0.942	-0.0189287\\
0.944	-0.0189287\\
0.946	-0.0189286\\
0.948	-0.0189286\\
0.95	-0.0189286\\
0.952	-0.0189285\\
0.954	-0.0189285\\
0.956	-0.0189285\\
0.958	-0.0189285\\
0.96	-0.0189285\\
0.962	-0.0189285\\
0.964	-0.0189285\\
0.966	-0.0189285\\
0.968	-0.0189284\\
0.97	-0.0189284\\
0.972	-0.0189284\\
0.974	-0.0189284\\
0.976	-0.0189284\\
0.978	-0.0189284\\
0.98	-0.0189285\\
0.982	-0.0189285\\
0.984	-0.0189285\\
0.986	-0.0189285\\
0.988	-0.0189285\\
0.99	-0.0189285\\
0.992	-0.0189285\\
0.994	-0.0189285\\
0.996	-0.0189286\\
0.998	-0.0189286\\
1	-0.0189286\\
1.002	-0.0189286\\
1.004	-0.0189286\\
1.006	-0.0189287\\
1.008	-0.0189287\\
1.01	-0.0189287\\
1.012	-0.0189287\\
1.014	-0.0189288\\
1.016	-0.0189288\\
1.018	-0.0189288\\
1.02	-0.0189288\\
1.022	-0.0189289\\
1.024	-0.0189289\\
1.026	-0.0189289\\
1.028	-0.0189289\\
1.03	-0.018929\\
1.032	-0.018929\\
1.034	-0.018929\\
1.036	-0.018929\\
1.038	-0.0189291\\
1.04	-0.0189291\\
1.042	-0.0189291\\
1.044	-0.0189291\\
1.046	-0.0189291\\
1.048	-0.0189292\\
1.05	-0.0189292\\
1.052	-0.0189292\\
1.054	-0.0189292\\
1.056	-0.0189292\\
1.058	-0.0189292\\
1.06	-0.0189293\\
1.062	-0.0189293\\
1.064	-0.0189293\\
1.066	-0.0189293\\
1.068	-0.0189293\\
1.07	-0.0189293\\
1.072	-0.0189293\\
1.074	-0.0189293\\
1.076	-0.0189294\\
1.078	-0.0189294\\
1.08	-0.0189294\\
1.082	-0.0189294\\
1.084	-0.0189294\\
1.086	-0.0189294\\
1.088	-0.0189294\\
1.09	-0.0189294\\
1.092	-0.0189294\\
1.094	-0.0189294\\
1.096	-0.0189294\\
1.098	-0.0189294\\
1.1	-0.0189294\\
1.102	-0.0189294\\
1.104	-0.0189294\\
1.106	-0.0189294\\
1.108	-0.0189294\\
1.11	-0.0189294\\
1.112	-0.0189294\\
1.114	-0.0189294\\
1.116	-0.0189294\\
1.118	-0.0189294\\
1.12	-0.0189293\\
1.122	-0.0189293\\
1.124	-0.0189293\\
1.126	-0.0189293\\
1.128	-0.0189293\\
1.13	-0.0189293\\
1.132	-0.0189293\\
1.134	-0.0189293\\
1.136	-0.0189293\\
1.138	-0.0189293\\
1.14	-0.0189293\\
1.142	-0.0189293\\
1.144	-0.0189293\\
1.146	-0.0189293\\
1.148	-0.0189292\\
1.15	-0.0189292\\
1.152	-0.0189292\\
1.154	-0.0189292\\
1.156	-0.0189292\\
1.158	-0.0189292\\
1.16	-0.0189292\\
1.162	-0.0189292\\
1.164	-0.0189292\\
1.166	-0.0189292\\
1.168	-0.0189292\\
1.17	-0.0189292\\
1.172	-0.0189292\\
1.174	-0.0189292\\
1.176	-0.0189291\\
1.178	-0.0189291\\
1.18	-0.0189291\\
1.182	-0.0189291\\
1.184	-0.0189291\\
1.186	-0.0189291\\
1.188	-0.0189291\\
1.19	-0.0189291\\
1.192	-0.0189291\\
1.194	-0.0189291\\
1.196	-0.0189291\\
1.198	-0.0189291\\
1.2	-0.0189291\\
1.202	-0.0189291\\
1.204	-0.0189291\\
1.206	-0.0189291\\
1.208	-0.0189291\\
1.21	-0.0189291\\
1.212	-0.0189291\\
1.214	-0.0189291\\
1.216	-0.0189291\\
1.218	-0.0189291\\
1.22	-0.0189291\\
1.222	-0.0189291\\
1.224	-0.0189291\\
1.226	-0.0189291\\
1.228	-0.0189291\\
1.23	-0.0189291\\
1.232	-0.0189291\\
1.234	-0.0189291\\
1.236	-0.0189291\\
1.238	-0.0189291\\
1.24	-0.0189291\\
1.242	-0.0189291\\
1.244	-0.0189291\\
1.246	-0.0189291\\
1.248	-0.0189291\\
1.25	-0.0189291\\
1.252	-0.0189291\\
1.254	-0.0189291\\
1.256	-0.0189291\\
1.258	-0.0189291\\
1.26	-0.0189291\\
1.262	-0.0189291\\
1.264	-0.0189291\\
1.266	-0.0189291\\
1.268	-0.0189291\\
1.27	-0.0189291\\
1.272	-0.0189291\\
1.274	-0.0189291\\
1.276	-0.0189291\\
1.278	-0.0189291\\
1.28	-0.0189291\\
1.282	-0.0189291\\
1.284	-0.0189291\\
1.286	-0.0189291\\
1.288	-0.0189291\\
1.29	-0.0189291\\
1.292	-0.0189292\\
1.294	-0.0189292\\
1.296	-0.0189292\\
1.298	-0.0189292\\
1.3	-0.0189292\\
1.302	-0.0189292\\
1.304	-0.0189292\\
1.306	-0.0189292\\
1.308	-0.0189292\\
1.31	-0.0189292\\
1.312	-0.0189292\\
1.314	-0.0189292\\
1.316	-0.0189292\\
1.318	-0.0189292\\
1.32	-0.0189292\\
1.322	-0.0189292\\
1.324	-0.0189292\\
1.326	-0.0189292\\
1.328	-0.0189292\\
1.33	-0.0189292\\
1.332	-0.0189292\\
1.334	-0.0189292\\
1.336	-0.0189292\\
1.338	-0.0189292\\
1.34	-0.0189292\\
1.342	-0.0189292\\
1.344	-0.0189292\\
1.346	-0.0189292\\
1.348	-0.0189292\\
1.35	-0.0189292\\
1.352	-0.0189292\\
1.354	-0.0189292\\
1.356	-0.0189292\\
1.358	-0.0189292\\
1.36	-0.0189292\\
1.362	-0.0189292\\
1.364	-0.0189292\\
1.366	-0.0189292\\
1.368	-0.0189292\\
1.37	-0.0189292\\
1.372	-0.0189292\\
1.374	-0.0189292\\
1.376	-0.0189292\\
1.378	-0.0189292\\
1.38	-0.0189292\\
1.382	-0.0189292\\
1.384	-0.0189292\\
1.386	-0.0189292\\
1.388	-0.0189292\\
1.39	-0.0189292\\
1.392	-0.0189292\\
1.394	-0.0189292\\
1.396	-0.0189292\\
1.398	-0.0189292\\
1.4	-0.0189292\\
1.402	-0.0189292\\
1.404	-0.0189292\\
1.406	-0.0189292\\
1.408	-0.0189292\\
1.41	-0.0189292\\
1.412	-0.0189292\\
1.414	-0.0189292\\
1.416	-0.0189292\\
1.418	-0.0189292\\
1.42	-0.0189292\\
1.422	-0.0189292\\
1.424	-0.0189292\\
1.426	-0.0189292\\
1.428	-0.0189292\\
1.43	-0.0189292\\
1.432	-0.0189292\\
1.434	-0.0189292\\
1.436	-0.0189292\\
1.438	-0.0189292\\
1.44	-0.0189292\\
1.442	-0.0189292\\
1.444	-0.0189292\\
1.446	-0.0189292\\
1.448	-0.0189291\\
1.45	-0.0189291\\
1.452	-0.0189291\\
1.454	-0.0189291\\
1.456	-0.0189291\\
1.458	-0.0189291\\
1.46	-0.0189291\\
1.462	-0.0189291\\
1.464	-0.0189291\\
1.466	-0.0189291\\
1.468	-0.0189291\\
1.47	-0.0189291\\
1.472	-0.0189291\\
1.474	-0.0189291\\
1.476	-0.0189291\\
1.478	-0.0189291\\
1.48	-0.0189291\\
1.482	-0.0189291\\
1.484	-0.0189291\\
1.486	-0.0189291\\
1.488	-0.0189291\\
1.49	-0.0189291\\
1.492	-0.0189291\\
1.494	-0.0189291\\
1.496	-0.0189291\\
1.498	-0.0189292\\
1.5	-0.0189292\\
1.502	-0.0189292\\
1.504	-0.0189292\\
1.506	-0.0189292\\
1.508	-0.0189292\\
1.51	-0.0189292\\
1.512	-0.0189292\\
1.514	-0.0189292\\
1.516	-0.0189292\\
1.518	-0.0189292\\
1.52	-0.0189292\\
1.522	-0.0189292\\
1.524	-0.0189292\\
1.526	-0.0189292\\
1.528	-0.0189292\\
1.53	-0.0189292\\
1.532	-0.0189292\\
1.534	-0.0189292\\
1.536	-0.0189292\\
1.538	-0.0189292\\
1.54	-0.0189292\\
1.542	-0.0189292\\
1.544	-0.0189292\\
1.546	-0.0189292\\
1.548	-0.0189292\\
1.55	-0.0189292\\
1.552	-0.0189292\\
1.554	-0.0189292\\
1.556	-0.0189292\\
1.558	-0.0189292\\
1.56	-0.0189292\\
1.562	-0.0189292\\
1.564	-0.0189292\\
1.566	-0.0189292\\
1.568	-0.0189292\\
1.57	-0.0189292\\
1.572	-0.0189292\\
1.574	-0.0189292\\
1.576	-0.0189292\\
1.578	-0.0189292\\
1.58	-0.0189292\\
1.582	-0.0189292\\
1.584	-0.0189292\\
1.586	-0.0189292\\
1.588	-0.0189292\\
1.59	-0.0189292\\
1.592	-0.0189292\\
1.594	-0.0189292\\
1.596	-0.0189292\\
1.598	-0.0189292\\
1.6	-0.0189292\\
1.602	-0.0189292\\
1.604	-0.0189292\\
1.606	-0.0189292\\
1.608	-0.0189292\\
1.61	-0.0189292\\
1.612	-0.0189292\\
1.614	-0.0189292\\
1.616	-0.0189292\\
1.618	-0.0189292\\
1.62	-0.0189292\\
1.622	-0.0189292\\
1.624	-0.0189292\\
1.626	-0.0189292\\
1.628	-0.0189292\\
1.63	-0.0189292\\
1.632	-0.0189292\\
1.634	-0.0189292\\
1.636	-0.0189292\\
1.638	-0.0189292\\
1.64	-0.0189292\\
1.642	-0.0189292\\
1.644	-0.0189292\\
1.646	-0.0189292\\
1.648	-0.0189292\\
1.65	-0.0189292\\
1.652	-0.0189292\\
1.654	-0.0189292\\
1.656	-0.0189292\\
1.658	-0.0189292\\
1.66	-0.0189292\\
1.662	-0.0189292\\
1.664	-0.0189292\\
1.666	-0.0189292\\
1.668	-0.0189292\\
1.67	-0.0189292\\
1.672	-0.0189292\\
1.674	-0.0189292\\
1.676	-0.0189292\\
1.678	-0.0189292\\
1.68	-0.0189292\\
1.682	-0.0189292\\
1.684	-0.0189292\\
1.686	-0.0189292\\
1.688	-0.0189292\\
1.69	-0.0189292\\
1.692	-0.0189292\\
1.694	-0.0189292\\
1.696	-0.0189292\\
1.698	-0.0189292\\
1.7	-0.0189292\\
1.702	-0.0189292\\
1.704	-0.0189292\\
1.706	-0.0189292\\
1.708	-0.0189292\\
1.71	-0.0189292\\
1.712	-0.0189292\\
1.714	-0.0189292\\
1.716	-0.0189292\\
1.718	-0.0189292\\
1.72	-0.0189292\\
1.722	-0.0189292\\
1.724	-0.0189292\\
1.726	-0.0189292\\
1.728	-0.0189292\\
1.73	-0.0189292\\
1.732	-0.0189292\\
1.734	-0.0189292\\
1.736	-0.0189292\\
1.738	-0.0189292\\
1.74	-0.0189292\\
1.742	-0.0189292\\
1.744	-0.0189292\\
1.746	-0.0189292\\
1.748	-0.0189292\\
1.75	-0.0189292\\
1.752	-0.0189292\\
1.754	-0.0189292\\
1.756	-0.0189292\\
1.758	-0.0189292\\
1.76	-0.0189292\\
1.762	-0.0189292\\
1.764	-0.0189292\\
1.766	-0.0189292\\
1.768	-0.0189292\\
1.77	-0.0189292\\
1.772	-0.0189292\\
1.774	-0.0189292\\
1.776	-0.0189292\\
1.778	-0.0189292\\
1.78	-0.0189292\\
1.782	-0.0189292\\
1.784	-0.0189292\\
1.786	-0.0189292\\
1.788	-0.0189292\\
1.79	-0.0189292\\
1.792	-0.0189292\\
1.794	-0.0189292\\
1.796	-0.0189292\\
1.798	-0.0189292\\
1.8	-0.0189292\\
1.802	-0.0189292\\
1.804	-0.0189292\\
1.806	-0.0189292\\
1.808	-0.0189292\\
1.81	-0.0189292\\
1.812	-0.0189292\\
1.814	-0.0189292\\
1.816	-0.0189292\\
1.818	-0.0189292\\
1.82	-0.0189292\\
1.822	-0.0189292\\
1.824	-0.0189292\\
1.826	-0.0189292\\
1.828	-0.0189292\\
1.83	-0.0189292\\
1.832	-0.0189292\\
1.834	-0.0189292\\
1.836	-0.0189292\\
1.838	-0.0189292\\
1.84	-0.0189292\\
1.842	-0.0189292\\
1.844	-0.0189292\\
1.846	-0.0189292\\
1.848	-0.0189292\\
1.85	-0.0189292\\
1.852	-0.0189292\\
1.854	-0.0189292\\
1.856	-0.0189292\\
1.858	-0.0189292\\
1.86	-0.0189292\\
1.862	-0.0189292\\
1.864	-0.0189292\\
1.866	-0.0189292\\
1.868	-0.0189292\\
1.87	-0.0189292\\
1.872	-0.0189292\\
1.874	-0.0189292\\
1.876	-0.0189292\\
1.878	-0.0189292\\
1.88	-0.0189292\\
1.882	-0.0189292\\
1.884	-0.0189292\\
1.886	-0.0189292\\
1.888	-0.0189292\\
1.89	-0.0189292\\
1.892	-0.0189292\\
1.894	-0.0189292\\
1.896	-0.0189292\\
1.898	-0.0189292\\
1.9	-0.0189292\\
1.902	-0.0189292\\
1.904	-0.0189292\\
1.906	-0.0189292\\
1.908	-0.0189292\\
1.91	-0.0189292\\
1.912	-0.0189292\\
1.914	-0.0189292\\
1.916	-0.0189292\\
1.918	-0.0189292\\
1.92	-0.0189292\\
1.922	-0.0189292\\
1.924	-0.0189292\\
1.926	-0.0189292\\
1.928	-0.0189292\\
1.93	-0.0189292\\
1.932	-0.0189292\\
1.934	-0.0189292\\
1.936	-0.0189292\\
1.938	-0.0189292\\
1.94	-0.0189292\\
1.942	-0.0189292\\
1.944	-0.0189292\\
1.946	-0.0189292\\
1.948	-0.0189292\\
1.95	-0.0189292\\
1.952	-0.0189292\\
1.954	-0.0189292\\
1.956	-0.0189292\\
1.958	-0.0189292\\
1.96	-0.0189292\\
1.962	-0.0189292\\
1.964	-0.0189292\\
1.966	-0.0189292\\
1.968	-0.0189292\\
1.97	-0.0189292\\
1.972	-0.0189292\\
1.974	-0.0189292\\
1.976	-0.0189292\\
1.978	-0.0189292\\
1.98	-0.0189292\\
1.982	-0.0189292\\
1.984	-0.0189292\\
1.986	-0.0189292\\
1.988	-0.0189292\\
1.99	-0.0189292\\
1.992	-0.0189292\\
1.994	-0.0189292\\
1.996	-0.0189292\\
1.998	-0.0189292\\
2	-0.0189292\\
};
\addplot [color=mycolor4, line width=1.0pt, forget plot]
  table[row sep=crcr]{%
-0.024	0\\
-0.022	0\\
-0.02	0\\
-0.018	0\\
-0.016	0\\
-0.014	0\\
-0.012	0\\
-0.01	0\\
-0.008	0\\
-0.006	0\\
-0.004	0\\
-0.002	0\\
0	0\\
0.002	-0.00183135\\
0.004	-0.00337788\\
0.006	-0.00470365\\
0.008	-0.00585694\\
0.01	-0.00687396\\
0.012	-0.00778188\\
0.014	-0.00860104\\
0.016	-0.00934673\\
0.018	-0.0100305\\
0.02	-0.0106612\\
0.022	-0.0112456\\
0.024	-0.0117891\\
0.026	-0.012296\\
0.028	-0.0127697\\
0.03	-0.013213\\
0.032	-0.0136286\\
0.034	-0.0140185\\
0.036	-0.0143846\\
0.038	-0.0147285\\
0.04	-0.0150518\\
0.042	-0.0153558\\
0.044	-0.015642\\
0.046	-0.0159113\\
0.048	-0.0161651\\
0.05	-0.0164044\\
0.052	-0.0166301\\
0.054	-0.0168433\\
0.056	-0.0170448\\
0.058	-0.0172355\\
0.06	-0.0174162\\
0.062	-0.0175876\\
0.064	-0.0177504\\
0.066	-0.0179053\\
0.068	-0.0180528\\
0.07	-0.0181933\\
0.072	-0.0183275\\
0.074	-0.0184555\\
0.076	-0.0185779\\
0.078	-0.0186947\\
0.08	-0.0188064\\
0.082	-0.018913\\
0.084	-0.0190147\\
0.086	-0.0191115\\
0.088	-0.0192037\\
0.09	-0.0192911\\
0.092	-0.0193738\\
0.094	-0.0194519\\
0.096	-0.0195252\\
0.098	-0.0195939\\
0.1	-0.0196578\\
0.102	-0.019717\\
0.104	-0.0197714\\
0.106	-0.019821\\
0.108	-0.0198659\\
0.11	-0.019906\\
0.112	-0.0199413\\
0.114	-0.0199719\\
0.116	-0.0199978\\
0.118	-0.0200191\\
0.12	-0.0200359\\
0.122	-0.0200482\\
0.124	-0.0200562\\
0.126	-0.02006\\
0.128	-0.0200598\\
0.13	-0.0200557\\
0.132	-0.0200478\\
0.134	-0.0200363\\
0.136	-0.0200215\\
0.138	-0.0200034\\
0.14	-0.0199824\\
0.142	-0.0199585\\
0.144	-0.0199321\\
0.146	-0.0199033\\
0.148	-0.0198723\\
0.15	-0.0198393\\
0.152	-0.0198046\\
0.154	-0.0197684\\
0.156	-0.0197308\\
0.158	-0.019692\\
0.16	-0.0196523\\
0.162	-0.0196118\\
0.164	-0.0195708\\
0.166	-0.0195294\\
0.168	-0.0194877\\
0.17	-0.0194459\\
0.172	-0.0194043\\
0.174	-0.0193629\\
0.176	-0.0193218\\
0.178	-0.0192812\\
0.18	-0.0192413\\
0.182	-0.0192021\\
0.184	-0.0191637\\
0.186	-0.0191262\\
0.188	-0.0190898\\
0.19	-0.0190544\\
0.192	-0.0190203\\
0.194	-0.0189873\\
0.196	-0.0189556\\
0.198	-0.0189253\\
0.2	-0.0188963\\
0.202	-0.0188687\\
0.204	-0.0188425\\
0.206	-0.0188178\\
0.208	-0.0187945\\
0.21	-0.0187727\\
0.212	-0.0187523\\
0.214	-0.0187334\\
0.216	-0.0187159\\
0.218	-0.0186998\\
0.22	-0.0186851\\
0.222	-0.0186717\\
0.224	-0.0186598\\
0.226	-0.0186491\\
0.228	-0.0186396\\
0.23	-0.0186314\\
0.232	-0.0186244\\
0.234	-0.0186186\\
0.236	-0.0186138\\
0.238	-0.0186101\\
0.24	-0.0186074\\
0.242	-0.0186056\\
0.244	-0.0186048\\
0.246	-0.0186048\\
0.248	-0.0186056\\
0.25	-0.0186072\\
0.252	-0.0186095\\
0.254	-0.0186125\\
0.256	-0.0186162\\
0.258	-0.0186204\\
0.26	-0.0186252\\
0.262	-0.0186305\\
0.264	-0.0186363\\
0.266	-0.0186425\\
0.268	-0.0186491\\
0.27	-0.018656\\
0.272	-0.0186633\\
0.274	-0.0186709\\
0.276	-0.0186788\\
0.278	-0.0186869\\
0.28	-0.0186952\\
0.282	-0.0187037\\
0.284	-0.0187124\\
0.286	-0.0187212\\
0.288	-0.0187301\\
0.29	-0.018739\\
0.292	-0.0187481\\
0.294	-0.0187572\\
0.296	-0.0187663\\
0.298	-0.0187754\\
0.3	-0.0187845\\
0.302	-0.0187936\\
0.304	-0.0188026\\
0.306	-0.0188115\\
0.308	-0.0188203\\
0.31	-0.0188291\\
0.312	-0.0188376\\
0.314	-0.0188461\\
0.316	-0.0188544\\
0.318	-0.0188625\\
0.32	-0.0188704\\
0.322	-0.0188782\\
0.324	-0.0188857\\
0.326	-0.018893\\
0.328	-0.0189\\
0.33	-0.0189068\\
0.332	-0.0189134\\
0.334	-0.0189197\\
0.336	-0.0189257\\
0.338	-0.0189314\\
0.34	-0.0189369\\
0.342	-0.018942\\
0.344	-0.0189469\\
0.346	-0.0189515\\
0.348	-0.0189557\\
0.35	-0.0189597\\
0.352	-0.0189634\\
0.354	-0.0189668\\
0.356	-0.0189699\\
0.358	-0.0189727\\
0.36	-0.0189752\\
0.362	-0.0189775\\
0.364	-0.0189794\\
0.366	-0.0189811\\
0.368	-0.0189826\\
0.37	-0.0189837\\
0.372	-0.0189847\\
0.374	-0.0189854\\
0.376	-0.0189859\\
0.378	-0.0189861\\
0.38	-0.0189862\\
0.382	-0.0189861\\
0.384	-0.0189857\\
0.386	-0.0189852\\
0.388	-0.0189846\\
0.39	-0.0189838\\
0.392	-0.0189828\\
0.394	-0.0189818\\
0.396	-0.0189806\\
0.398	-0.0189792\\
0.4	-0.0189778\\
0.402	-0.0189764\\
0.404	-0.0189748\\
0.406	-0.0189732\\
0.408	-0.0189715\\
0.41	-0.0189697\\
0.412	-0.018968\\
0.414	-0.0189661\\
0.416	-0.0189643\\
0.418	-0.0189625\\
0.42	-0.0189606\\
0.422	-0.0189587\\
0.424	-0.0189569\\
0.426	-0.018955\\
0.428	-0.0189532\\
0.43	-0.0189514\\
0.432	-0.0189496\\
0.434	-0.0189479\\
0.436	-0.0189462\\
0.438	-0.0189445\\
0.44	-0.0189428\\
0.442	-0.0189412\\
0.444	-0.0189397\\
0.446	-0.0189382\\
0.448	-0.0189367\\
0.45	-0.0189353\\
0.452	-0.018934\\
0.454	-0.0189327\\
0.456	-0.0189314\\
0.458	-0.0189303\\
0.46	-0.0189291\\
0.462	-0.0189281\\
0.464	-0.018927\\
0.466	-0.0189261\\
0.468	-0.0189252\\
0.47	-0.0189243\\
0.472	-0.0189235\\
0.474	-0.0189228\\
0.476	-0.0189221\\
0.478	-0.0189215\\
0.48	-0.0189209\\
0.482	-0.0189204\\
0.484	-0.0189199\\
0.486	-0.0189194\\
0.488	-0.018919\\
0.49	-0.0189187\\
0.492	-0.0189184\\
0.494	-0.0189182\\
0.496	-0.018918\\
0.498	-0.0189178\\
0.5	-0.0189177\\
0.502	-0.0189176\\
0.504	-0.0189175\\
0.506	-0.0189175\\
0.508	-0.0189175\\
0.51	-0.0189176\\
0.512	-0.0189177\\
0.514	-0.0189178\\
0.516	-0.0189179\\
0.518	-0.0189181\\
0.52	-0.0189183\\
0.522	-0.0189185\\
0.524	-0.0189187\\
0.526	-0.018919\\
0.528	-0.0189192\\
0.53	-0.0189195\\
0.532	-0.0189198\\
0.534	-0.0189201\\
0.536	-0.0189205\\
0.538	-0.0189208\\
0.54	-0.0189211\\
0.542	-0.0189215\\
0.544	-0.0189218\\
0.546	-0.0189222\\
0.548	-0.0189226\\
0.55	-0.018923\\
0.552	-0.0189233\\
0.554	-0.0189237\\
0.556	-0.0189241\\
0.558	-0.0189244\\
0.56	-0.0189248\\
0.562	-0.0189251\\
0.564	-0.0189255\\
0.566	-0.0189258\\
0.568	-0.0189262\\
0.57	-0.0189265\\
0.572	-0.0189268\\
0.574	-0.0189271\\
0.576	-0.0189274\\
0.578	-0.0189277\\
0.58	-0.018928\\
0.582	-0.0189283\\
0.584	-0.0189285\\
0.586	-0.0189288\\
0.588	-0.018929\\
0.59	-0.0189292\\
0.592	-0.0189294\\
0.594	-0.0189296\\
0.596	-0.0189298\\
0.598	-0.01893\\
0.6	-0.0189302\\
0.602	-0.0189303\\
0.604	-0.0189305\\
0.606	-0.0189306\\
0.608	-0.0189307\\
0.61	-0.0189308\\
0.612	-0.0189309\\
0.614	-0.018931\\
0.616	-0.0189311\\
0.618	-0.0189311\\
0.62	-0.0189312\\
0.622	-0.0189312\\
0.624	-0.0189313\\
0.626	-0.0189313\\
0.628	-0.0189313\\
0.63	-0.0189313\\
0.632	-0.0189313\\
0.634	-0.0189313\\
0.636	-0.0189313\\
0.638	-0.0189313\\
0.64	-0.0189313\\
0.642	-0.0189313\\
0.644	-0.0189312\\
0.646	-0.0189312\\
0.648	-0.0189312\\
0.65	-0.0189311\\
0.652	-0.0189311\\
0.654	-0.018931\\
0.656	-0.018931\\
0.658	-0.0189309\\
0.66	-0.0189309\\
0.662	-0.0189308\\
0.664	-0.0189307\\
0.666	-0.0189307\\
0.668	-0.0189306\\
0.67	-0.0189305\\
0.672	-0.0189305\\
0.674	-0.0189304\\
0.676	-0.0189303\\
0.678	-0.0189303\\
0.68	-0.0189302\\
0.682	-0.0189301\\
0.684	-0.0189301\\
0.686	-0.01893\\
0.688	-0.0189299\\
0.69	-0.0189299\\
0.692	-0.0189298\\
0.694	-0.0189297\\
0.696	-0.0189297\\
0.698	-0.0189296\\
0.7	-0.0189296\\
0.702	-0.0189295\\
0.704	-0.0189295\\
0.706	-0.0189294\\
0.708	-0.0189294\\
0.71	-0.0189293\\
0.712	-0.0189293\\
0.714	-0.0189292\\
0.716	-0.0189292\\
0.718	-0.0189291\\
0.72	-0.0189291\\
0.722	-0.0189291\\
0.724	-0.018929\\
0.726	-0.018929\\
0.728	-0.018929\\
0.73	-0.0189289\\
0.732	-0.0189289\\
0.734	-0.0189289\\
0.736	-0.0189289\\
0.738	-0.0189288\\
0.74	-0.0189288\\
0.742	-0.0189288\\
0.744	-0.0189288\\
0.746	-0.0189288\\
0.748	-0.0189288\\
0.75	-0.0189288\\
0.752	-0.0189288\\
0.754	-0.0189288\\
0.756	-0.0189288\\
0.758	-0.0189288\\
0.76	-0.0189288\\
0.762	-0.0189288\\
0.764	-0.0189288\\
0.766	-0.0189288\\
0.768	-0.0189288\\
0.77	-0.0189288\\
0.772	-0.0189288\\
0.774	-0.0189288\\
0.776	-0.0189288\\
0.778	-0.0189288\\
0.78	-0.0189288\\
0.782	-0.0189288\\
0.784	-0.0189288\\
0.786	-0.0189288\\
0.788	-0.0189289\\
0.79	-0.0189289\\
0.792	-0.0189289\\
0.794	-0.0189289\\
0.796	-0.0189289\\
0.798	-0.0189289\\
0.8	-0.0189289\\
0.802	-0.0189289\\
0.804	-0.018929\\
0.806	-0.018929\\
0.808	-0.018929\\
0.81	-0.018929\\
0.812	-0.018929\\
0.814	-0.018929\\
0.816	-0.018929\\
0.818	-0.018929\\
0.82	-0.0189291\\
0.822	-0.0189291\\
0.824	-0.0189291\\
0.826	-0.0189291\\
0.828	-0.0189291\\
0.83	-0.0189291\\
0.832	-0.0189291\\
0.834	-0.0189291\\
0.836	-0.0189291\\
0.838	-0.0189291\\
0.84	-0.0189291\\
0.842	-0.0189292\\
0.844	-0.0189292\\
0.846	-0.0189292\\
0.848	-0.0189292\\
0.85	-0.0189292\\
0.852	-0.0189292\\
0.854	-0.0189292\\
0.856	-0.0189292\\
0.858	-0.0189292\\
0.86	-0.0189292\\
0.862	-0.0189292\\
0.864	-0.0189292\\
0.866	-0.0189292\\
0.868	-0.0189292\\
0.87	-0.0189292\\
0.872	-0.0189292\\
0.874	-0.0189292\\
0.876	-0.0189292\\
0.878	-0.0189292\\
0.88	-0.0189292\\
0.882	-0.0189292\\
0.884	-0.0189292\\
0.886	-0.0189292\\
0.888	-0.0189292\\
0.89	-0.0189292\\
0.892	-0.0189292\\
0.894	-0.0189292\\
0.896	-0.0189292\\
0.898	-0.0189292\\
0.9	-0.0189292\\
0.902	-0.0189292\\
0.904	-0.0189292\\
0.906	-0.0189292\\
0.908	-0.0189292\\
0.91	-0.0189292\\
0.912	-0.0189292\\
0.914	-0.0189292\\
0.916	-0.0189292\\
0.918	-0.0189292\\
0.92	-0.0189292\\
0.922	-0.0189292\\
0.924	-0.0189292\\
0.926	-0.0189292\\
0.928	-0.0189292\\
0.93	-0.0189292\\
0.932	-0.0189292\\
0.934	-0.0189292\\
0.936	-0.0189292\\
0.938	-0.0189292\\
0.94	-0.0189292\\
0.942	-0.0189292\\
0.944	-0.0189292\\
0.946	-0.0189292\\
0.948	-0.0189292\\
0.95	-0.0189292\\
0.952	-0.0189292\\
0.954	-0.0189292\\
0.956	-0.0189292\\
0.958	-0.0189292\\
0.96	-0.0189292\\
0.962	-0.0189292\\
0.964	-0.0189292\\
0.966	-0.0189292\\
0.968	-0.0189292\\
0.97	-0.0189292\\
0.972	-0.0189292\\
0.974	-0.0189292\\
0.976	-0.0189292\\
0.978	-0.0189292\\
0.98	-0.0189291\\
0.982	-0.0189291\\
0.984	-0.0189291\\
0.986	-0.0189291\\
0.988	-0.0189291\\
0.99	-0.0189291\\
0.992	-0.0189291\\
0.994	-0.0189291\\
0.996	-0.0189291\\
0.998	-0.0189291\\
1	-0.0189291\\
1.002	-0.0189291\\
1.004	-0.0189291\\
1.006	-0.0189291\\
1.008	-0.0189291\\
1.01	-0.0189291\\
1.012	-0.0189291\\
1.014	-0.0189291\\
1.016	-0.0189291\\
1.018	-0.0189291\\
1.02	-0.0189291\\
1.022	-0.0189291\\
1.024	-0.0189291\\
1.026	-0.0189291\\
1.028	-0.0189291\\
1.03	-0.0189291\\
1.032	-0.0189291\\
1.034	-0.0189291\\
1.036	-0.0189291\\
1.038	-0.0189291\\
1.04	-0.0189291\\
1.042	-0.0189291\\
1.044	-0.0189291\\
1.046	-0.0189291\\
1.048	-0.0189291\\
1.05	-0.0189291\\
1.052	-0.0189292\\
1.054	-0.0189292\\
1.056	-0.0189292\\
1.058	-0.0189292\\
1.06	-0.0189292\\
1.062	-0.0189292\\
1.064	-0.0189292\\
1.066	-0.0189292\\
1.068	-0.0189292\\
1.07	-0.0189292\\
1.072	-0.0189292\\
1.074	-0.0189292\\
1.076	-0.0189292\\
1.078	-0.0189292\\
1.08	-0.0189292\\
1.082	-0.0189292\\
1.084	-0.0189292\\
1.086	-0.0189292\\
1.088	-0.0189292\\
1.09	-0.0189292\\
1.092	-0.0189292\\
1.094	-0.0189292\\
1.096	-0.0189292\\
1.098	-0.0189292\\
1.1	-0.0189292\\
1.102	-0.0189292\\
1.104	-0.0189292\\
1.106	-0.0189292\\
1.108	-0.0189292\\
1.11	-0.0189292\\
1.112	-0.0189292\\
1.114	-0.0189292\\
1.116	-0.0189292\\
1.118	-0.0189292\\
1.12	-0.0189292\\
1.122	-0.0189292\\
1.124	-0.0189292\\
1.126	-0.0189292\\
1.128	-0.0189292\\
1.13	-0.0189292\\
1.132	-0.0189292\\
1.134	-0.0189292\\
1.136	-0.0189292\\
1.138	-0.0189292\\
1.14	-0.0189292\\
1.142	-0.0189292\\
1.144	-0.0189292\\
1.146	-0.0189292\\
1.148	-0.0189292\\
1.15	-0.0189292\\
1.152	-0.0189292\\
1.154	-0.0189292\\
1.156	-0.0189292\\
1.158	-0.0189292\\
1.16	-0.0189292\\
1.162	-0.0189292\\
1.164	-0.0189292\\
1.166	-0.0189292\\
1.168	-0.0189292\\
1.17	-0.0189292\\
1.172	-0.0189292\\
1.174	-0.0189292\\
1.176	-0.0189292\\
1.178	-0.0189292\\
1.18	-0.0189292\\
1.182	-0.0189292\\
1.184	-0.0189292\\
1.186	-0.0189292\\
1.188	-0.0189292\\
1.19	-0.0189292\\
1.192	-0.0189292\\
1.194	-0.0189292\\
1.196	-0.0189292\\
1.198	-0.0189292\\
1.2	-0.0189292\\
1.202	-0.0189292\\
1.204	-0.0189292\\
1.206	-0.0189292\\
1.208	-0.0189292\\
1.21	-0.0189292\\
1.212	-0.0189292\\
1.214	-0.0189292\\
1.216	-0.0189292\\
1.218	-0.0189292\\
1.22	-0.0189292\\
1.222	-0.0189292\\
1.224	-0.0189292\\
1.226	-0.0189292\\
1.228	-0.0189292\\
1.23	-0.0189292\\
1.232	-0.0189292\\
1.234	-0.0189292\\
1.236	-0.0189292\\
1.238	-0.0189292\\
1.24	-0.0189292\\
1.242	-0.0189292\\
1.244	-0.0189292\\
1.246	-0.0189292\\
1.248	-0.0189292\\
1.25	-0.0189292\\
1.252	-0.0189292\\
1.254	-0.0189292\\
1.256	-0.0189292\\
1.258	-0.0189292\\
1.26	-0.0189292\\
1.262	-0.0189292\\
1.264	-0.0189292\\
1.266	-0.0189292\\
1.268	-0.0189292\\
1.27	-0.0189292\\
1.272	-0.0189292\\
1.274	-0.0189292\\
1.276	-0.0189292\\
1.278	-0.0189292\\
1.28	-0.0189292\\
1.282	-0.0189292\\
1.284	-0.0189292\\
1.286	-0.0189292\\
1.288	-0.0189292\\
1.29	-0.0189292\\
1.292	-0.0189292\\
1.294	-0.0189292\\
1.296	-0.0189292\\
1.298	-0.0189292\\
1.3	-0.0189292\\
1.302	-0.0189292\\
1.304	-0.0189292\\
1.306	-0.0189292\\
1.308	-0.0189292\\
1.31	-0.0189292\\
1.312	-0.0189292\\
1.314	-0.0189292\\
1.316	-0.0189292\\
1.318	-0.0189292\\
1.32	-0.0189292\\
1.322	-0.0189292\\
1.324	-0.0189292\\
1.326	-0.0189292\\
1.328	-0.0189292\\
1.33	-0.0189292\\
1.332	-0.0189292\\
1.334	-0.0189292\\
1.336	-0.0189292\\
1.338	-0.0189292\\
1.34	-0.0189292\\
1.342	-0.0189292\\
1.344	-0.0189292\\
1.346	-0.0189292\\
1.348	-0.0189292\\
1.35	-0.0189292\\
1.352	-0.0189292\\
1.354	-0.0189292\\
1.356	-0.0189292\\
1.358	-0.0189292\\
1.36	-0.0189292\\
1.362	-0.0189292\\
1.364	-0.0189292\\
1.366	-0.0189292\\
1.368	-0.0189292\\
1.37	-0.0189292\\
1.372	-0.0189292\\
1.374	-0.0189292\\
1.376	-0.0189292\\
1.378	-0.0189292\\
1.38	-0.0189292\\
1.382	-0.0189292\\
1.384	-0.0189292\\
1.386	-0.0189292\\
1.388	-0.0189292\\
1.39	-0.0189292\\
1.392	-0.0189292\\
1.394	-0.0189292\\
1.396	-0.0189292\\
1.398	-0.0189292\\
1.4	-0.0189292\\
1.402	-0.0189292\\
1.404	-0.0189292\\
1.406	-0.0189292\\
1.408	-0.0189292\\
1.41	-0.0189292\\
1.412	-0.0189292\\
1.414	-0.0189292\\
1.416	-0.0189292\\
1.418	-0.0189292\\
1.42	-0.0189292\\
1.422	-0.0189292\\
1.424	-0.0189292\\
1.426	-0.0189292\\
1.428	-0.0189292\\
1.43	-0.0189292\\
1.432	-0.0189292\\
1.434	-0.0189292\\
1.436	-0.0189292\\
1.438	-0.0189292\\
1.44	-0.0189292\\
1.442	-0.0189292\\
1.444	-0.0189292\\
1.446	-0.0189292\\
1.448	-0.0189292\\
1.45	-0.0189292\\
1.452	-0.0189292\\
1.454	-0.0189292\\
1.456	-0.0189292\\
1.458	-0.0189292\\
1.46	-0.0189292\\
1.462	-0.0189292\\
1.464	-0.0189292\\
1.466	-0.0189292\\
1.468	-0.0189292\\
1.47	-0.0189292\\
1.472	-0.0189292\\
1.474	-0.0189292\\
1.476	-0.0189292\\
1.478	-0.0189292\\
1.48	-0.0189292\\
1.482	-0.0189292\\
1.484	-0.0189292\\
1.486	-0.0189292\\
1.488	-0.0189292\\
1.49	-0.0189292\\
1.492	-0.0189292\\
1.494	-0.0189292\\
1.496	-0.0189292\\
1.498	-0.0189292\\
1.5	-0.0189292\\
1.502	-0.0189292\\
1.504	-0.0189292\\
1.506	-0.0189292\\
1.508	-0.0189292\\
1.51	-0.0189292\\
1.512	-0.0189292\\
1.514	-0.0189292\\
1.516	-0.0189292\\
1.518	-0.0189292\\
1.52	-0.0189292\\
1.522	-0.0189292\\
1.524	-0.0189292\\
1.526	-0.0189292\\
1.528	-0.0189292\\
1.53	-0.0189292\\
1.532	-0.0189292\\
1.534	-0.0189292\\
1.536	-0.0189292\\
1.538	-0.0189292\\
1.54	-0.0189292\\
1.542	-0.0189292\\
1.544	-0.0189292\\
1.546	-0.0189292\\
1.548	-0.0189292\\
1.55	-0.0189292\\
1.552	-0.0189292\\
1.554	-0.0189292\\
1.556	-0.0189292\\
1.558	-0.0189292\\
1.56	-0.0189292\\
1.562	-0.0189292\\
1.564	-0.0189292\\
1.566	-0.0189292\\
1.568	-0.0189292\\
1.57	-0.0189292\\
1.572	-0.0189292\\
1.574	-0.0189292\\
1.576	-0.0189292\\
1.578	-0.0189292\\
1.58	-0.0189292\\
1.582	-0.0189292\\
1.584	-0.0189292\\
1.586	-0.0189292\\
1.588	-0.0189292\\
1.59	-0.0189292\\
1.592	-0.0189292\\
1.594	-0.0189292\\
1.596	-0.0189292\\
1.598	-0.0189292\\
1.6	-0.0189292\\
1.602	-0.0189292\\
1.604	-0.0189292\\
1.606	-0.0189292\\
1.608	-0.0189292\\
1.61	-0.0189292\\
1.612	-0.0189292\\
1.614	-0.0189292\\
1.616	-0.0189292\\
1.618	-0.0189292\\
1.62	-0.0189292\\
1.622	-0.0189292\\
1.624	-0.0189292\\
1.626	-0.0189292\\
1.628	-0.0189292\\
1.63	-0.0189292\\
1.632	-0.0189292\\
1.634	-0.0189292\\
1.636	-0.0189292\\
1.638	-0.0189292\\
1.64	-0.0189292\\
1.642	-0.0189292\\
1.644	-0.0189292\\
1.646	-0.0189292\\
1.648	-0.0189292\\
1.65	-0.0189292\\
1.652	-0.0189292\\
1.654	-0.0189292\\
1.656	-0.0189292\\
1.658	-0.0189292\\
1.66	-0.0189292\\
1.662	-0.0189292\\
1.664	-0.0189292\\
1.666	-0.0189292\\
1.668	-0.0189292\\
1.67	-0.0189292\\
1.672	-0.0189292\\
1.674	-0.0189292\\
1.676	-0.0189292\\
1.678	-0.0189292\\
1.68	-0.0189292\\
1.682	-0.0189292\\
1.684	-0.0189292\\
1.686	-0.0189292\\
1.688	-0.0189292\\
1.69	-0.0189292\\
1.692	-0.0189292\\
1.694	-0.0189292\\
1.696	-0.0189292\\
1.698	-0.0189292\\
1.7	-0.0189292\\
1.702	-0.0189292\\
1.704	-0.0189292\\
1.706	-0.0189292\\
1.708	-0.0189292\\
1.71	-0.0189292\\
1.712	-0.0189292\\
1.714	-0.0189292\\
1.716	-0.0189292\\
1.718	-0.0189292\\
1.72	-0.0189292\\
1.722	-0.0189292\\
1.724	-0.0189292\\
1.726	-0.0189292\\
1.728	-0.0189292\\
1.73	-0.0189292\\
1.732	-0.0189292\\
1.734	-0.0189292\\
1.736	-0.0189292\\
1.738	-0.0189292\\
1.74	-0.0189292\\
1.742	-0.0189292\\
1.744	-0.0189292\\
1.746	-0.0189292\\
1.748	-0.0189292\\
1.75	-0.0189292\\
1.752	-0.0189292\\
1.754	-0.0189292\\
1.756	-0.0189292\\
1.758	-0.0189292\\
1.76	-0.0189292\\
1.762	-0.0189292\\
1.764	-0.0189292\\
1.766	-0.0189292\\
1.768	-0.0189292\\
1.77	-0.0189292\\
1.772	-0.0189292\\
1.774	-0.0189292\\
1.776	-0.0189292\\
1.778	-0.0189292\\
1.78	-0.0189292\\
1.782	-0.0189292\\
1.784	-0.0189292\\
1.786	-0.0189292\\
1.788	-0.0189292\\
1.79	-0.0189292\\
1.792	-0.0189292\\
1.794	-0.0189292\\
1.796	-0.0189292\\
1.798	-0.0189292\\
1.8	-0.0189292\\
1.802	-0.0189292\\
1.804	-0.0189292\\
1.806	-0.0189292\\
1.808	-0.0189292\\
1.81	-0.0189292\\
1.812	-0.0189292\\
1.814	-0.0189292\\
1.816	-0.0189292\\
1.818	-0.0189292\\
1.82	-0.0189292\\
1.822	-0.0189292\\
1.824	-0.0189292\\
1.826	-0.0189292\\
1.828	-0.0189292\\
1.83	-0.0189292\\
1.832	-0.0189292\\
1.834	-0.0189292\\
1.836	-0.0189292\\
1.838	-0.0189292\\
1.84	-0.0189292\\
1.842	-0.0189292\\
1.844	-0.0189292\\
1.846	-0.0189292\\
1.848	-0.0189292\\
1.85	-0.0189292\\
1.852	-0.0189292\\
1.854	-0.0189292\\
1.856	-0.0189292\\
1.858	-0.0189292\\
1.86	-0.0189292\\
1.862	-0.0189292\\
1.864	-0.0189292\\
1.866	-0.0189292\\
1.868	-0.0189292\\
1.87	-0.0189292\\
1.872	-0.0189292\\
1.874	-0.0189292\\
1.876	-0.0189292\\
1.878	-0.0189292\\
1.88	-0.0189292\\
1.882	-0.0189292\\
1.884	-0.0189292\\
1.886	-0.0189292\\
1.888	-0.0189292\\
1.89	-0.0189292\\
1.892	-0.0189292\\
1.894	-0.0189292\\
1.896	-0.0189292\\
1.898	-0.0189292\\
1.9	-0.0189292\\
1.902	-0.0189292\\
1.904	-0.0189292\\
1.906	-0.0189292\\
1.908	-0.0189292\\
1.91	-0.0189292\\
1.912	-0.0189292\\
1.914	-0.0189292\\
1.916	-0.0189292\\
1.918	-0.0189292\\
1.92	-0.0189292\\
1.922	-0.0189292\\
1.924	-0.0189292\\
1.926	-0.0189292\\
1.928	-0.0189292\\
1.93	-0.0189292\\
1.932	-0.0189292\\
1.934	-0.0189292\\
1.936	-0.0189292\\
1.938	-0.0189292\\
1.94	-0.0189292\\
1.942	-0.0189292\\
1.944	-0.0189292\\
1.946	-0.0189292\\
1.948	-0.0189292\\
1.95	-0.0189292\\
1.952	-0.0189292\\
1.954	-0.0189292\\
1.956	-0.0189292\\
1.958	-0.0189292\\
1.96	-0.0189292\\
1.962	-0.0189292\\
1.964	-0.0189292\\
1.966	-0.0189292\\
1.968	-0.0189292\\
1.97	-0.0189292\\
1.972	-0.0189292\\
1.974	-0.0189292\\
1.976	-0.0189292\\
1.978	-0.0189292\\
1.98	-0.0189292\\
1.982	-0.0189292\\
1.984	-0.0189292\\
1.986	-0.0189292\\
1.988	-0.0189292\\
1.99	-0.0189292\\
1.992	-0.0189292\\
1.994	-0.0189292\\
1.996	-0.0189292\\
1.998	-0.0189292\\
2	-0.0189292\\
};
\addplot [color=mycolor5, line width=1.0pt, forget plot]
  table[row sep=crcr]{%
-0.024	0\\
-0.022	0\\
-0.02	0\\
-0.018	0\\
-0.016	0\\
-0.014	0\\
-0.012	0\\
-0.01	0\\
-0.008	0\\
-0.006	0\\
-0.004	0\\
-0.002	0\\
0	0\\
0.002	-0.00183135\\
0.004	-0.00337785\\
0.006	-0.00470358\\
0.008	-0.00585677\\
0.01	-0.00687365\\
0.012	-0.00778136\\
0.014	-0.00860025\\
0.016	-0.0093456\\
0.018	-0.010029\\
0.02	-0.0106592\\
0.022	-0.0112431\\
0.024	-0.011786\\
0.026	-0.0122921\\
0.028	-0.0127651\\
0.03	-0.0132077\\
0.032	-0.0136224\\
0.034	-0.0140115\\
0.036	-0.0143766\\
0.038	-0.0147196\\
0.04	-0.015042\\
0.042	-0.0153451\\
0.044	-0.0156302\\
0.046	-0.0158986\\
0.048	-0.0161513\\
0.05	-0.0163896\\
0.052	-0.0166142\\
0.054	-0.0168263\\
0.056	-0.0170267\\
0.058	-0.0172162\\
0.06	-0.0173957\\
0.062	-0.0175659\\
0.064	-0.0177274\\
0.066	-0.0178808\\
0.068	-0.0180268\\
0.07	-0.0181658\\
0.072	-0.0182983\\
0.074	-0.0184245\\
0.076	-0.018545\\
0.078	-0.0186598\\
0.08	-0.0187693\\
0.082	-0.0188736\\
0.084	-0.0189729\\
0.086	-0.0190672\\
0.088	-0.0191566\\
0.09	-0.0192413\\
0.092	-0.0193211\\
0.094	-0.0193962\\
0.096	-0.0194665\\
0.098	-0.019532\\
0.1	-0.0195928\\
0.102	-0.0196488\\
0.104	-0.0197\\
0.106	-0.0197465\\
0.108	-0.0197883\\
0.11	-0.0198253\\
0.112	-0.0198578\\
0.114	-0.0198857\\
0.116	-0.0199091\\
0.118	-0.0199281\\
0.12	-0.0199428\\
0.122	-0.0199534\\
0.124	-0.0199599\\
0.126	-0.0199627\\
0.128	-0.0199617\\
0.13	-0.0199572\\
0.132	-0.0199493\\
0.134	-0.0199384\\
0.136	-0.0199244\\
0.138	-0.0199078\\
0.14	-0.0198886\\
0.142	-0.0198671\\
0.144	-0.0198434\\
0.146	-0.0198178\\
0.148	-0.0197905\\
0.15	-0.0197617\\
0.152	-0.0197316\\
0.154	-0.0197003\\
0.156	-0.0196681\\
0.158	-0.019635\\
0.16	-0.0196014\\
0.162	-0.0195673\\
0.164	-0.0195329\\
0.166	-0.0194983\\
0.168	-0.0194637\\
0.17	-0.0194292\\
0.172	-0.0193948\\
0.174	-0.0193608\\
0.176	-0.0193272\\
0.178	-0.0192941\\
0.18	-0.0192615\\
0.182	-0.0192296\\
0.184	-0.0191984\\
0.186	-0.0191679\\
0.188	-0.0191383\\
0.19	-0.0191096\\
0.192	-0.0190818\\
0.194	-0.0190549\\
0.196	-0.0190289\\
0.198	-0.019004\\
0.2	-0.0189801\\
0.202	-0.0189572\\
0.204	-0.0189354\\
0.206	-0.0189146\\
0.208	-0.0188949\\
0.21	-0.0188762\\
0.212	-0.0188586\\
0.214	-0.018842\\
0.216	-0.0188265\\
0.218	-0.018812\\
0.22	-0.0187985\\
0.222	-0.018786\\
0.224	-0.0187744\\
0.226	-0.0187639\\
0.228	-0.0187542\\
0.23	-0.0187455\\
0.232	-0.0187377\\
0.234	-0.0187308\\
0.236	-0.0187247\\
0.238	-0.0187194\\
0.24	-0.0187149\\
0.242	-0.0187111\\
0.244	-0.0187081\\
0.246	-0.0187058\\
0.248	-0.0187042\\
0.25	-0.0187032\\
0.252	-0.0187028\\
0.254	-0.018703\\
0.256	-0.0187037\\
0.258	-0.018705\\
0.26	-0.0187068\\
0.262	-0.0187091\\
0.264	-0.0187118\\
0.266	-0.0187149\\
0.268	-0.0187184\\
0.27	-0.0187222\\
0.272	-0.0187264\\
0.274	-0.0187309\\
0.276	-0.0187356\\
0.278	-0.0187407\\
0.28	-0.0187459\\
0.282	-0.0187514\\
0.284	-0.018757\\
0.286	-0.0187628\\
0.288	-0.0187687\\
0.29	-0.0187748\\
0.292	-0.0187809\\
0.294	-0.0187871\\
0.296	-0.0187934\\
0.298	-0.0187997\\
0.3	-0.018806\\
0.302	-0.0188124\\
0.304	-0.0188187\\
0.306	-0.0188249\\
0.308	-0.0188311\\
0.31	-0.0188373\\
0.312	-0.0188434\\
0.314	-0.0188493\\
0.316	-0.0188552\\
0.318	-0.018861\\
0.32	-0.0188666\\
0.322	-0.0188721\\
0.324	-0.0188774\\
0.326	-0.0188826\\
0.328	-0.0188877\\
0.33	-0.0188925\\
0.332	-0.0188972\\
0.334	-0.0189017\\
0.336	-0.018906\\
0.338	-0.0189102\\
0.34	-0.0189141\\
0.342	-0.0189179\\
0.344	-0.0189214\\
0.346	-0.0189248\\
0.348	-0.018928\\
0.35	-0.018931\\
0.352	-0.0189338\\
0.354	-0.0189364\\
0.356	-0.0189388\\
0.358	-0.0189411\\
0.36	-0.0189431\\
0.362	-0.018945\\
0.364	-0.0189468\\
0.366	-0.0189483\\
0.368	-0.0189497\\
0.37	-0.018951\\
0.372	-0.0189521\\
0.374	-0.018953\\
0.376	-0.0189538\\
0.378	-0.0189545\\
0.38	-0.0189551\\
0.382	-0.0189555\\
0.384	-0.0189558\\
0.386	-0.018956\\
0.388	-0.0189561\\
0.39	-0.0189561\\
0.392	-0.0189561\\
0.394	-0.0189559\\
0.396	-0.0189557\\
0.398	-0.0189553\\
0.4	-0.0189549\\
0.402	-0.0189545\\
0.404	-0.018954\\
0.406	-0.0189534\\
0.408	-0.0189528\\
0.41	-0.0189522\\
0.412	-0.0189515\\
0.414	-0.0189507\\
0.416	-0.01895\\
0.418	-0.0189492\\
0.42	-0.0189484\\
0.422	-0.0189476\\
0.424	-0.0189468\\
0.426	-0.0189459\\
0.428	-0.0189451\\
0.43	-0.0189442\\
0.432	-0.0189434\\
0.434	-0.0189425\\
0.436	-0.0189417\\
0.438	-0.0189408\\
0.44	-0.01894\\
0.442	-0.0189392\\
0.444	-0.0189384\\
0.446	-0.0189376\\
0.448	-0.0189368\\
0.45	-0.018936\\
0.452	-0.0189353\\
0.454	-0.0189346\\
0.456	-0.0189339\\
0.458	-0.0189332\\
0.46	-0.0189326\\
0.462	-0.018932\\
0.464	-0.0189314\\
0.466	-0.0189308\\
0.468	-0.0189302\\
0.47	-0.0189297\\
0.472	-0.0189292\\
0.474	-0.0189288\\
0.476	-0.0189283\\
0.478	-0.0189279\\
0.48	-0.0189275\\
0.482	-0.0189271\\
0.484	-0.0189268\\
0.486	-0.0189265\\
0.488	-0.0189262\\
0.49	-0.0189259\\
0.492	-0.0189257\\
0.494	-0.0189254\\
0.496	-0.0189252\\
0.498	-0.0189251\\
0.5	-0.0189249\\
0.502	-0.0189248\\
0.504	-0.0189246\\
0.506	-0.0189245\\
0.508	-0.0189244\\
0.51	-0.0189244\\
0.512	-0.0189243\\
0.514	-0.0189243\\
0.516	-0.0189243\\
0.518	-0.0189243\\
0.52	-0.0189243\\
0.522	-0.0189243\\
0.524	-0.0189243\\
0.526	-0.0189244\\
0.528	-0.0189244\\
0.53	-0.0189245\\
0.532	-0.0189246\\
0.534	-0.0189246\\
0.536	-0.0189247\\
0.538	-0.0189248\\
0.54	-0.0189249\\
0.542	-0.018925\\
0.544	-0.0189251\\
0.546	-0.0189253\\
0.548	-0.0189254\\
0.55	-0.0189255\\
0.552	-0.0189256\\
0.554	-0.0189258\\
0.556	-0.0189259\\
0.558	-0.018926\\
0.56	-0.0189261\\
0.562	-0.0189263\\
0.564	-0.0189264\\
0.566	-0.0189265\\
0.568	-0.0189267\\
0.57	-0.0189268\\
0.572	-0.0189269\\
0.574	-0.0189271\\
0.576	-0.0189272\\
0.578	-0.0189273\\
0.58	-0.0189274\\
0.582	-0.0189275\\
0.584	-0.0189277\\
0.586	-0.0189278\\
0.588	-0.0189279\\
0.59	-0.018928\\
0.592	-0.0189281\\
0.594	-0.0189282\\
0.596	-0.0189283\\
0.598	-0.0189284\\
0.6	-0.0189285\\
0.602	-0.0189285\\
0.604	-0.0189286\\
0.606	-0.0189287\\
0.608	-0.0189288\\
0.61	-0.0189288\\
0.612	-0.0189289\\
0.614	-0.0189289\\
0.616	-0.018929\\
0.618	-0.0189291\\
0.62	-0.0189291\\
0.622	-0.0189291\\
0.624	-0.0189292\\
0.626	-0.0189292\\
0.628	-0.0189293\\
0.63	-0.0189293\\
0.632	-0.0189293\\
0.634	-0.0189293\\
0.636	-0.0189294\\
0.638	-0.0189294\\
0.64	-0.0189294\\
0.642	-0.0189294\\
0.644	-0.0189294\\
0.646	-0.0189295\\
0.648	-0.0189295\\
0.65	-0.0189295\\
0.652	-0.0189295\\
0.654	-0.0189295\\
0.656	-0.0189295\\
0.658	-0.0189295\\
0.66	-0.0189295\\
0.662	-0.0189295\\
0.664	-0.0189295\\
0.666	-0.0189295\\
0.668	-0.0189295\\
0.67	-0.0189295\\
0.672	-0.0189295\\
0.674	-0.0189295\\
0.676	-0.0189295\\
0.678	-0.0189295\\
0.68	-0.0189294\\
0.682	-0.0189294\\
0.684	-0.0189294\\
0.686	-0.0189294\\
0.688	-0.0189294\\
0.69	-0.0189294\\
0.692	-0.0189294\\
0.694	-0.0189294\\
0.696	-0.0189294\\
0.698	-0.0189294\\
0.7	-0.0189294\\
0.702	-0.0189293\\
0.704	-0.0189293\\
0.706	-0.0189293\\
0.708	-0.0189293\\
0.71	-0.0189293\\
0.712	-0.0189293\\
0.714	-0.0189293\\
0.716	-0.0189293\\
0.718	-0.0189293\\
0.72	-0.0189293\\
0.722	-0.0189293\\
0.724	-0.0189292\\
0.726	-0.0189292\\
0.728	-0.0189292\\
0.73	-0.0189292\\
0.732	-0.0189292\\
0.734	-0.0189292\\
0.736	-0.0189292\\
0.738	-0.0189292\\
0.74	-0.0189292\\
0.742	-0.0189292\\
0.744	-0.0189292\\
0.746	-0.0189292\\
0.748	-0.0189292\\
0.75	-0.0189292\\
0.752	-0.0189292\\
0.754	-0.0189292\\
0.756	-0.0189292\\
0.758	-0.0189291\\
0.76	-0.0189291\\
0.762	-0.0189291\\
0.764	-0.0189291\\
0.766	-0.0189291\\
0.768	-0.0189291\\
0.77	-0.0189291\\
0.772	-0.0189291\\
0.774	-0.0189291\\
0.776	-0.0189291\\
0.778	-0.0189291\\
0.78	-0.0189291\\
0.782	-0.0189291\\
0.784	-0.0189291\\
0.786	-0.0189291\\
0.788	-0.0189291\\
0.79	-0.0189291\\
0.792	-0.0189291\\
0.794	-0.0189291\\
0.796	-0.0189291\\
0.798	-0.0189291\\
0.8	-0.0189291\\
0.802	-0.0189291\\
0.804	-0.0189291\\
0.806	-0.0189291\\
0.808	-0.0189291\\
0.81	-0.0189291\\
0.812	-0.0189291\\
0.814	-0.0189291\\
0.816	-0.0189291\\
0.818	-0.0189291\\
0.82	-0.0189291\\
0.822	-0.0189291\\
0.824	-0.0189291\\
0.826	-0.0189291\\
0.828	-0.0189291\\
0.83	-0.0189291\\
0.832	-0.0189291\\
0.834	-0.0189291\\
0.836	-0.0189291\\
0.838	-0.0189291\\
0.84	-0.0189291\\
0.842	-0.0189291\\
0.844	-0.0189291\\
0.846	-0.0189291\\
0.848	-0.0189291\\
0.85	-0.0189291\\
0.852	-0.0189291\\
0.854	-0.0189291\\
0.856	-0.0189291\\
0.858	-0.0189291\\
0.86	-0.0189291\\
0.862	-0.0189291\\
0.864	-0.0189291\\
0.866	-0.0189291\\
0.868	-0.0189291\\
0.87	-0.0189291\\
0.872	-0.0189291\\
0.874	-0.0189291\\
0.876	-0.0189291\\
0.878	-0.0189291\\
0.88	-0.0189291\\
0.882	-0.0189291\\
0.884	-0.0189291\\
0.886	-0.0189291\\
0.888	-0.0189291\\
0.89	-0.0189291\\
0.892	-0.0189291\\
0.894	-0.0189291\\
0.896	-0.0189291\\
0.898	-0.0189291\\
0.9	-0.0189291\\
0.902	-0.0189291\\
0.904	-0.0189291\\
0.906	-0.0189291\\
0.908	-0.0189291\\
0.91	-0.0189291\\
0.912	-0.0189291\\
0.914	-0.0189291\\
0.916	-0.0189291\\
0.918	-0.0189291\\
0.92	-0.0189291\\
0.922	-0.0189291\\
0.924	-0.0189291\\
0.926	-0.0189291\\
0.928	-0.0189291\\
0.93	-0.0189291\\
0.932	-0.0189291\\
0.934	-0.0189291\\
0.936	-0.0189291\\
0.938	-0.0189291\\
0.94	-0.0189291\\
0.942	-0.0189292\\
0.944	-0.0189292\\
0.946	-0.0189292\\
0.948	-0.0189292\\
0.95	-0.0189292\\
0.952	-0.0189292\\
0.954	-0.0189292\\
0.956	-0.0189292\\
0.958	-0.0189292\\
0.96	-0.0189292\\
0.962	-0.0189292\\
0.964	-0.0189292\\
0.966	-0.0189292\\
0.968	-0.0189292\\
0.97	-0.0189292\\
0.972	-0.0189292\\
0.974	-0.0189292\\
0.976	-0.0189292\\
0.978	-0.0189292\\
0.98	-0.0189292\\
0.982	-0.0189292\\
0.984	-0.0189292\\
0.986	-0.0189292\\
0.988	-0.0189292\\
0.99	-0.0189292\\
0.992	-0.0189292\\
0.994	-0.0189292\\
0.996	-0.0189292\\
0.998	-0.0189292\\
1	-0.0189292\\
1.002	-0.0189292\\
1.004	-0.0189292\\
1.006	-0.0189292\\
1.008	-0.0189292\\
1.01	-0.0189292\\
1.012	-0.0189292\\
1.014	-0.0189292\\
1.016	-0.0189292\\
1.018	-0.0189292\\
1.02	-0.0189292\\
1.022	-0.0189292\\
1.024	-0.0189292\\
1.026	-0.0189292\\
1.028	-0.0189292\\
1.03	-0.0189292\\
1.032	-0.0189292\\
1.034	-0.0189292\\
1.036	-0.0189292\\
1.038	-0.0189292\\
1.04	-0.0189292\\
1.042	-0.0189292\\
1.044	-0.0189292\\
1.046	-0.0189292\\
1.048	-0.0189292\\
1.05	-0.0189292\\
1.052	-0.0189292\\
1.054	-0.0189292\\
1.056	-0.0189292\\
1.058	-0.0189292\\
1.06	-0.0189292\\
1.062	-0.0189292\\
1.064	-0.0189292\\
1.066	-0.0189292\\
1.068	-0.0189292\\
1.07	-0.0189292\\
1.072	-0.0189292\\
1.074	-0.0189292\\
1.076	-0.0189292\\
1.078	-0.0189292\\
1.08	-0.0189292\\
1.082	-0.0189292\\
1.084	-0.0189292\\
1.086	-0.0189292\\
1.088	-0.0189292\\
1.09	-0.0189292\\
1.092	-0.0189292\\
1.094	-0.0189292\\
1.096	-0.0189292\\
1.098	-0.0189292\\
1.1	-0.0189292\\
1.102	-0.0189292\\
1.104	-0.0189292\\
1.106	-0.0189292\\
1.108	-0.0189292\\
1.11	-0.0189292\\
1.112	-0.0189292\\
1.114	-0.0189292\\
1.116	-0.0189292\\
1.118	-0.0189292\\
1.12	-0.0189292\\
1.122	-0.0189292\\
1.124	-0.0189292\\
1.126	-0.0189292\\
1.128	-0.0189292\\
1.13	-0.0189292\\
1.132	-0.0189292\\
1.134	-0.0189292\\
1.136	-0.0189292\\
1.138	-0.0189292\\
1.14	-0.0189292\\
1.142	-0.0189292\\
1.144	-0.0189292\\
1.146	-0.0189292\\
1.148	-0.0189292\\
1.15	-0.0189292\\
1.152	-0.0189292\\
1.154	-0.0189292\\
1.156	-0.0189292\\
1.158	-0.0189292\\
1.16	-0.0189292\\
1.162	-0.0189292\\
1.164	-0.0189292\\
1.166	-0.0189292\\
1.168	-0.0189292\\
1.17	-0.0189292\\
1.172	-0.0189292\\
1.174	-0.0189292\\
1.176	-0.0189292\\
1.178	-0.0189292\\
1.18	-0.0189292\\
1.182	-0.0189292\\
1.184	-0.0189292\\
1.186	-0.0189292\\
1.188	-0.0189292\\
1.19	-0.0189292\\
1.192	-0.0189292\\
1.194	-0.0189292\\
1.196	-0.0189292\\
1.198	-0.0189292\\
1.2	-0.0189292\\
1.202	-0.0189292\\
1.204	-0.0189292\\
1.206	-0.0189292\\
1.208	-0.0189292\\
1.21	-0.0189292\\
1.212	-0.0189292\\
1.214	-0.0189292\\
1.216	-0.0189292\\
1.218	-0.0189292\\
1.22	-0.0189292\\
1.222	-0.0189292\\
1.224	-0.0189292\\
1.226	-0.0189292\\
1.228	-0.0189292\\
1.23	-0.0189292\\
1.232	-0.0189292\\
1.234	-0.0189292\\
1.236	-0.0189292\\
1.238	-0.0189292\\
1.24	-0.0189292\\
1.242	-0.0189292\\
1.244	-0.0189292\\
1.246	-0.0189292\\
1.248	-0.0189292\\
1.25	-0.0189292\\
1.252	-0.0189292\\
1.254	-0.0189292\\
1.256	-0.0189292\\
1.258	-0.0189292\\
1.26	-0.0189292\\
1.262	-0.0189292\\
1.264	-0.0189292\\
1.266	-0.0189292\\
1.268	-0.0189292\\
1.27	-0.0189292\\
1.272	-0.0189292\\
1.274	-0.0189292\\
1.276	-0.0189292\\
1.278	-0.0189292\\
1.28	-0.0189292\\
1.282	-0.0189292\\
1.284	-0.0189292\\
1.286	-0.0189292\\
1.288	-0.0189292\\
1.29	-0.0189292\\
1.292	-0.0189292\\
1.294	-0.0189292\\
1.296	-0.0189292\\
1.298	-0.0189292\\
1.3	-0.0189292\\
1.302	-0.0189292\\
1.304	-0.0189292\\
1.306	-0.0189292\\
1.308	-0.0189292\\
1.31	-0.0189292\\
1.312	-0.0189292\\
1.314	-0.0189292\\
1.316	-0.0189292\\
1.318	-0.0189292\\
1.32	-0.0189292\\
1.322	-0.0189292\\
1.324	-0.0189292\\
1.326	-0.0189292\\
1.328	-0.0189292\\
1.33	-0.0189292\\
1.332	-0.0189292\\
1.334	-0.0189292\\
1.336	-0.0189292\\
1.338	-0.0189292\\
1.34	-0.0189292\\
1.342	-0.0189292\\
1.344	-0.0189292\\
1.346	-0.0189292\\
1.348	-0.0189292\\
1.35	-0.0189292\\
1.352	-0.0189292\\
1.354	-0.0189292\\
1.356	-0.0189292\\
1.358	-0.0189292\\
1.36	-0.0189292\\
1.362	-0.0189292\\
1.364	-0.0189292\\
1.366	-0.0189292\\
1.368	-0.0189292\\
1.37	-0.0189292\\
1.372	-0.0189292\\
1.374	-0.0189292\\
1.376	-0.0189292\\
1.378	-0.0189292\\
1.38	-0.0189292\\
1.382	-0.0189292\\
1.384	-0.0189292\\
1.386	-0.0189292\\
1.388	-0.0189292\\
1.39	-0.0189292\\
1.392	-0.0189292\\
1.394	-0.0189292\\
1.396	-0.0189292\\
1.398	-0.0189292\\
1.4	-0.0189292\\
1.402	-0.0189292\\
1.404	-0.0189292\\
1.406	-0.0189292\\
1.408	-0.0189292\\
1.41	-0.0189292\\
1.412	-0.0189292\\
1.414	-0.0189292\\
1.416	-0.0189292\\
1.418	-0.0189292\\
1.42	-0.0189292\\
1.422	-0.0189292\\
1.424	-0.0189292\\
1.426	-0.0189292\\
1.428	-0.0189292\\
1.43	-0.0189292\\
1.432	-0.0189292\\
1.434	-0.0189292\\
1.436	-0.0189292\\
1.438	-0.0189292\\
1.44	-0.0189292\\
1.442	-0.0189292\\
1.444	-0.0189292\\
1.446	-0.0189292\\
1.448	-0.0189292\\
1.45	-0.0189292\\
1.452	-0.0189292\\
1.454	-0.0189292\\
1.456	-0.0189292\\
1.458	-0.0189292\\
1.46	-0.0189292\\
1.462	-0.0189292\\
1.464	-0.0189292\\
1.466	-0.0189292\\
1.468	-0.0189292\\
1.47	-0.0189292\\
1.472	-0.0189292\\
1.474	-0.0189292\\
1.476	-0.0189292\\
1.478	-0.0189292\\
1.48	-0.0189292\\
1.482	-0.0189292\\
1.484	-0.0189292\\
1.486	-0.0189292\\
1.488	-0.0189292\\
1.49	-0.0189292\\
1.492	-0.0189292\\
1.494	-0.0189292\\
1.496	-0.0189292\\
1.498	-0.0189292\\
1.5	-0.0189292\\
1.502	-0.0189292\\
1.504	-0.0189292\\
1.506	-0.0189292\\
1.508	-0.0189292\\
1.51	-0.0189292\\
1.512	-0.0189292\\
1.514	-0.0189292\\
1.516	-0.0189292\\
1.518	-0.0189292\\
1.52	-0.0189292\\
1.522	-0.0189292\\
1.524	-0.0189292\\
1.526	-0.0189292\\
1.528	-0.0189292\\
1.53	-0.0189292\\
1.532	-0.0189292\\
1.534	-0.0189292\\
1.536	-0.0189292\\
1.538	-0.0189292\\
1.54	-0.0189292\\
1.542	-0.0189292\\
1.544	-0.0189292\\
1.546	-0.0189292\\
1.548	-0.0189292\\
1.55	-0.0189292\\
1.552	-0.0189292\\
1.554	-0.0189292\\
1.556	-0.0189292\\
1.558	-0.0189292\\
1.56	-0.0189292\\
1.562	-0.0189292\\
1.564	-0.0189292\\
1.566	-0.0189292\\
1.568	-0.0189292\\
1.57	-0.0189292\\
1.572	-0.0189292\\
1.574	-0.0189292\\
1.576	-0.0189292\\
1.578	-0.0189292\\
1.58	-0.0189292\\
1.582	-0.0189292\\
1.584	-0.0189292\\
1.586	-0.0189292\\
1.588	-0.0189292\\
1.59	-0.0189292\\
1.592	-0.0189292\\
1.594	-0.0189292\\
1.596	-0.0189292\\
1.598	-0.0189292\\
1.6	-0.0189292\\
1.602	-0.0189292\\
1.604	-0.0189292\\
1.606	-0.0189292\\
1.608	-0.0189292\\
1.61	-0.0189292\\
1.612	-0.0189292\\
1.614	-0.0189292\\
1.616	-0.0189292\\
1.618	-0.0189292\\
1.62	-0.0189292\\
1.622	-0.0189292\\
1.624	-0.0189292\\
1.626	-0.0189292\\
1.628	-0.0189292\\
1.63	-0.0189292\\
1.632	-0.0189292\\
1.634	-0.0189292\\
1.636	-0.0189292\\
1.638	-0.0189292\\
1.64	-0.0189292\\
1.642	-0.0189292\\
1.644	-0.0189292\\
1.646	-0.0189292\\
1.648	-0.0189292\\
1.65	-0.0189292\\
1.652	-0.0189292\\
1.654	-0.0189292\\
1.656	-0.0189292\\
1.658	-0.0189292\\
1.66	-0.0189292\\
1.662	-0.0189292\\
1.664	-0.0189292\\
1.666	-0.0189292\\
1.668	-0.0189292\\
1.67	-0.0189292\\
1.672	-0.0189292\\
1.674	-0.0189292\\
1.676	-0.0189292\\
1.678	-0.0189292\\
1.68	-0.0189292\\
1.682	-0.0189292\\
1.684	-0.0189292\\
1.686	-0.0189292\\
1.688	-0.0189292\\
1.69	-0.0189292\\
1.692	-0.0189292\\
1.694	-0.0189292\\
1.696	-0.0189292\\
1.698	-0.0189292\\
1.7	-0.0189292\\
1.702	-0.0189292\\
1.704	-0.0189292\\
1.706	-0.0189292\\
1.708	-0.0189292\\
1.71	-0.0189292\\
1.712	-0.0189292\\
1.714	-0.0189292\\
1.716	-0.0189292\\
1.718	-0.0189292\\
1.72	-0.0189292\\
1.722	-0.0189292\\
1.724	-0.0189292\\
1.726	-0.0189292\\
1.728	-0.0189292\\
1.73	-0.0189292\\
1.732	-0.0189292\\
1.734	-0.0189292\\
1.736	-0.0189292\\
1.738	-0.0189292\\
1.74	-0.0189292\\
1.742	-0.0189292\\
1.744	-0.0189292\\
1.746	-0.0189292\\
1.748	-0.0189292\\
1.75	-0.0189292\\
1.752	-0.0189292\\
1.754	-0.0189292\\
1.756	-0.0189292\\
1.758	-0.0189292\\
1.76	-0.0189292\\
1.762	-0.0189292\\
1.764	-0.0189292\\
1.766	-0.0189292\\
1.768	-0.0189292\\
1.77	-0.0189292\\
1.772	-0.0189292\\
1.774	-0.0189292\\
1.776	-0.0189292\\
1.778	-0.0189292\\
1.78	-0.0189292\\
1.782	-0.0189292\\
1.784	-0.0189292\\
1.786	-0.0189292\\
1.788	-0.0189292\\
1.79	-0.0189292\\
1.792	-0.0189292\\
1.794	-0.0189292\\
1.796	-0.0189292\\
1.798	-0.0189292\\
1.8	-0.0189292\\
1.802	-0.0189292\\
1.804	-0.0189292\\
1.806	-0.0189292\\
1.808	-0.0189292\\
1.81	-0.0189292\\
1.812	-0.0189292\\
1.814	-0.0189292\\
1.816	-0.0189292\\
1.818	-0.0189292\\
1.82	-0.0189292\\
1.822	-0.0189292\\
1.824	-0.0189292\\
1.826	-0.0189292\\
1.828	-0.0189292\\
1.83	-0.0189292\\
1.832	-0.0189292\\
1.834	-0.0189292\\
1.836	-0.0189292\\
1.838	-0.0189292\\
1.84	-0.0189292\\
1.842	-0.0189292\\
1.844	-0.0189292\\
1.846	-0.0189292\\
1.848	-0.0189292\\
1.85	-0.0189292\\
1.852	-0.0189292\\
1.854	-0.0189292\\
1.856	-0.0189292\\
1.858	-0.0189292\\
1.86	-0.0189292\\
1.862	-0.0189292\\
1.864	-0.0189292\\
1.866	-0.0189292\\
1.868	-0.0189292\\
1.87	-0.0189292\\
1.872	-0.0189292\\
1.874	-0.0189292\\
1.876	-0.0189292\\
1.878	-0.0189292\\
1.88	-0.0189292\\
1.882	-0.0189292\\
1.884	-0.0189292\\
1.886	-0.0189292\\
1.888	-0.0189292\\
1.89	-0.0189292\\
1.892	-0.0189292\\
1.894	-0.0189292\\
1.896	-0.0189292\\
1.898	-0.0189292\\
1.9	-0.0189292\\
1.902	-0.0189292\\
1.904	-0.0189292\\
1.906	-0.0189292\\
1.908	-0.0189292\\
1.91	-0.0189292\\
1.912	-0.0189292\\
1.914	-0.0189292\\
1.916	-0.0189292\\
1.918	-0.0189292\\
1.92	-0.0189292\\
1.922	-0.0189292\\
1.924	-0.0189292\\
1.926	-0.0189292\\
1.928	-0.0189292\\
1.93	-0.0189292\\
1.932	-0.0189292\\
1.934	-0.0189292\\
1.936	-0.0189292\\
1.938	-0.0189292\\
1.94	-0.0189292\\
1.942	-0.0189292\\
1.944	-0.0189292\\
1.946	-0.0189292\\
1.948	-0.0189292\\
1.95	-0.0189292\\
1.952	-0.0189292\\
1.954	-0.0189292\\
1.956	-0.0189292\\
1.958	-0.0189292\\
1.96	-0.0189292\\
1.962	-0.0189292\\
1.964	-0.0189292\\
1.966	-0.0189292\\
1.968	-0.0189292\\
1.97	-0.0189292\\
1.972	-0.0189292\\
1.974	-0.0189292\\
1.976	-0.0189292\\
1.978	-0.0189292\\
1.98	-0.0189292\\
1.982	-0.0189292\\
1.984	-0.0189292\\
1.986	-0.0189292\\
1.988	-0.0189292\\
1.99	-0.0189292\\
1.992	-0.0189292\\
1.994	-0.0189292\\
1.996	-0.0189292\\
1.998	-0.0189292\\
2	-0.0189292\\
};
\addplot [color=mycolor6, line width=1.0pt, forget plot]
  table[row sep=crcr]{%
-0.024	0\\
-0.022	0\\
-0.02	0\\
-0.018	0\\
-0.016	0\\
-0.014	0\\
-0.012	0\\
-0.01	0\\
-0.008	0\\
-0.006	0\\
-0.004	0\\
-0.002	0\\
0	0\\
0.002	-0.00183137\\
0.004	-0.00337798\\
0.006	-0.00470399\\
0.008	-0.00585772\\
0.01	-0.00687546\\
0.012	-0.00778441\\
0.014	-0.00860498\\
0.016	-0.00935252\\
0.018	-0.0100386\\
0.02	-0.0106722\\
0.022	-0.0112601\\
0.024	-0.0118076\\
0.026	-0.0123191\\
0.028	-0.0127981\\
0.03	-0.0132474\\
0.032	-0.0136694\\
0.034	-0.0140662\\
0.036	-0.0144397\\
0.038	-0.0147914\\
0.04	-0.0151229\\
0.042	-0.0154352\\
0.044	-0.0157297\\
0.046	-0.0160074\\
0.048	-0.0162694\\
0.05	-0.0165164\\
0.052	-0.0167495\\
0.054	-0.0169695\\
0.056	-0.017177\\
0.058	-0.0173729\\
0.06	-0.0175578\\
0.062	-0.0177322\\
0.064	-0.0178969\\
0.066	-0.0180523\\
0.068	-0.0181989\\
0.07	-0.0183372\\
0.072	-0.0184675\\
0.074	-0.0185903\\
0.076	-0.0187058\\
0.078	-0.0188144\\
0.08	-0.0189163\\
0.082	-0.0190117\\
0.084	-0.0191009\\
0.086	-0.0191841\\
0.088	-0.0192613\\
0.09	-0.0193329\\
0.092	-0.0193988\\
0.094	-0.0194593\\
0.096	-0.0195146\\
0.098	-0.0195646\\
0.1	-0.0196097\\
0.102	-0.0196499\\
0.104	-0.0196854\\
0.106	-0.0197163\\
0.108	-0.0197428\\
0.11	-0.0197651\\
0.112	-0.0197833\\
0.114	-0.0197976\\
0.116	-0.0198081\\
0.118	-0.0198152\\
0.12	-0.0198189\\
0.122	-0.0198195\\
0.124	-0.0198171\\
0.126	-0.0198119\\
0.128	-0.0198042\\
0.13	-0.0197941\\
0.132	-0.0197819\\
0.134	-0.0197676\\
0.136	-0.0197515\\
0.138	-0.0197338\\
0.14	-0.0197146\\
0.142	-0.0196941\\
0.144	-0.0196725\\
0.146	-0.0196498\\
0.148	-0.0196263\\
0.15	-0.019602\\
0.152	-0.0195771\\
0.154	-0.0195516\\
0.156	-0.0195258\\
0.158	-0.0194997\\
0.16	-0.0194733\\
0.162	-0.0194468\\
0.164	-0.0194203\\
0.166	-0.0193937\\
0.168	-0.0193673\\
0.17	-0.0193409\\
0.172	-0.0193148\\
0.174	-0.0192889\\
0.176	-0.0192633\\
0.178	-0.019238\\
0.18	-0.019213\\
0.182	-0.0191885\\
0.184	-0.0191644\\
0.186	-0.0191408\\
0.188	-0.0191177\\
0.19	-0.0190951\\
0.192	-0.0190731\\
0.194	-0.0190517\\
0.196	-0.0190309\\
0.198	-0.0190107\\
0.2	-0.0189912\\
0.202	-0.0189723\\
0.204	-0.0189542\\
0.206	-0.0189368\\
0.208	-0.0189201\\
0.21	-0.0189042\\
0.212	-0.018889\\
0.214	-0.0188746\\
0.216	-0.0188609\\
0.218	-0.0188481\\
0.22	-0.018836\\
0.222	-0.0188247\\
0.224	-0.0188142\\
0.226	-0.0188044\\
0.228	-0.0187954\\
0.23	-0.0187872\\
0.232	-0.0187797\\
0.234	-0.0187729\\
0.236	-0.0187669\\
0.238	-0.0187616\\
0.24	-0.0187569\\
0.242	-0.0187529\\
0.244	-0.0187495\\
0.246	-0.0187467\\
0.248	-0.0187446\\
0.25	-0.0187429\\
0.252	-0.0187419\\
0.254	-0.0187413\\
0.256	-0.0187412\\
0.258	-0.0187416\\
0.26	-0.0187424\\
0.262	-0.0187436\\
0.264	-0.0187452\\
0.266	-0.0187472\\
0.268	-0.0187495\\
0.27	-0.0187521\\
0.272	-0.0187549\\
0.274	-0.018758\\
0.276	-0.0187614\\
0.278	-0.018765\\
0.28	-0.0187688\\
0.282	-0.0187727\\
0.284	-0.0187768\\
0.286	-0.018781\\
0.288	-0.0187854\\
0.29	-0.0187898\\
0.292	-0.0187943\\
0.294	-0.0187989\\
0.296	-0.0188036\\
0.298	-0.0188082\\
0.3	-0.0188129\\
0.302	-0.0188177\\
0.304	-0.0188224\\
0.306	-0.0188271\\
0.308	-0.0188317\\
0.31	-0.0188364\\
0.312	-0.018841\\
0.314	-0.0188455\\
0.316	-0.01885\\
0.318	-0.0188544\\
0.32	-0.0188587\\
0.322	-0.018863\\
0.324	-0.0188671\\
0.326	-0.0188712\\
0.328	-0.0188751\\
0.33	-0.018879\\
0.332	-0.0188827\\
0.334	-0.0188863\\
0.336	-0.0188898\\
0.338	-0.0188932\\
0.34	-0.0188965\\
0.342	-0.0188996\\
0.344	-0.0189026\\
0.346	-0.0189055\\
0.348	-0.0189082\\
0.35	-0.0189108\\
0.352	-0.0189133\\
0.354	-0.0189157\\
0.356	-0.0189179\\
0.358	-0.01892\\
0.36	-0.018922\\
0.362	-0.0189238\\
0.364	-0.0189256\\
0.366	-0.0189272\\
0.368	-0.0189287\\
0.37	-0.0189301\\
0.372	-0.0189314\\
0.374	-0.0189326\\
0.376	-0.0189337\\
0.378	-0.0189347\\
0.38	-0.0189355\\
0.382	-0.0189363\\
0.384	-0.0189371\\
0.386	-0.0189377\\
0.388	-0.0189382\\
0.39	-0.0189387\\
0.392	-0.0189391\\
0.394	-0.0189394\\
0.396	-0.0189397\\
0.398	-0.0189399\\
0.4	-0.01894\\
0.402	-0.0189401\\
0.404	-0.0189402\\
0.406	-0.0189401\\
0.408	-0.0189401\\
0.41	-0.01894\\
0.412	-0.0189399\\
0.414	-0.0189397\\
0.416	-0.0189395\\
0.418	-0.0189392\\
0.42	-0.018939\\
0.422	-0.0189387\\
0.424	-0.0189384\\
0.426	-0.0189381\\
0.428	-0.0189377\\
0.43	-0.0189374\\
0.432	-0.018937\\
0.434	-0.0189366\\
0.436	-0.0189362\\
0.438	-0.0189359\\
0.44	-0.0189355\\
0.442	-0.0189351\\
0.444	-0.0189347\\
0.446	-0.0189342\\
0.448	-0.0189338\\
0.45	-0.0189334\\
0.452	-0.018933\\
0.454	-0.0189326\\
0.456	-0.0189323\\
0.458	-0.0189319\\
0.46	-0.0189315\\
0.462	-0.0189311\\
0.464	-0.0189308\\
0.466	-0.0189304\\
0.468	-0.0189301\\
0.47	-0.0189298\\
0.472	-0.0189294\\
0.474	-0.0189291\\
0.476	-0.0189288\\
0.478	-0.0189285\\
0.48	-0.0189283\\
0.482	-0.018928\\
0.484	-0.0189278\\
0.486	-0.0189275\\
0.488	-0.0189273\\
0.49	-0.0189271\\
0.492	-0.0189269\\
0.494	-0.0189267\\
0.496	-0.0189265\\
0.498	-0.0189263\\
0.5	-0.0189262\\
0.502	-0.018926\\
0.504	-0.0189259\\
0.506	-0.0189258\\
0.508	-0.0189257\\
0.51	-0.0189256\\
0.512	-0.0189255\\
0.514	-0.0189254\\
0.516	-0.0189254\\
0.518	-0.0189253\\
0.52	-0.0189252\\
0.522	-0.0189252\\
0.524	-0.0189252\\
0.526	-0.0189251\\
0.528	-0.0189251\\
0.53	-0.0189251\\
0.532	-0.0189251\\
0.534	-0.0189251\\
0.536	-0.0189251\\
0.538	-0.0189251\\
0.54	-0.0189252\\
0.542	-0.0189252\\
0.544	-0.0189252\\
0.546	-0.0189253\\
0.548	-0.0189253\\
0.55	-0.0189254\\
0.552	-0.0189254\\
0.554	-0.0189255\\
0.556	-0.0189255\\
0.558	-0.0189256\\
0.56	-0.0189256\\
0.562	-0.0189257\\
0.564	-0.0189258\\
0.566	-0.0189258\\
0.568	-0.0189259\\
0.57	-0.018926\\
0.572	-0.018926\\
0.574	-0.0189261\\
0.576	-0.0189262\\
0.578	-0.0189263\\
0.58	-0.0189263\\
0.582	-0.0189264\\
0.584	-0.0189265\\
0.586	-0.0189266\\
0.588	-0.0189266\\
0.59	-0.0189267\\
0.592	-0.0189268\\
0.594	-0.0189269\\
0.596	-0.0189269\\
0.598	-0.018927\\
0.6	-0.0189271\\
0.602	-0.0189272\\
0.604	-0.0189272\\
0.606	-0.0189273\\
0.608	-0.0189274\\
0.61	-0.0189274\\
0.612	-0.0189275\\
0.614	-0.0189275\\
0.616	-0.0189276\\
0.618	-0.0189277\\
0.62	-0.0189277\\
0.622	-0.0189278\\
0.624	-0.0189278\\
0.626	-0.0189279\\
0.628	-0.0189279\\
0.63	-0.018928\\
0.632	-0.018928\\
0.634	-0.0189281\\
0.636	-0.0189281\\
0.638	-0.0189282\\
0.64	-0.0189282\\
0.642	-0.0189283\\
0.644	-0.0189283\\
0.646	-0.0189283\\
0.648	-0.0189284\\
0.65	-0.0189284\\
0.652	-0.0189285\\
0.654	-0.0189285\\
0.656	-0.0189285\\
0.658	-0.0189286\\
0.66	-0.0189286\\
0.662	-0.0189286\\
0.664	-0.0189286\\
0.666	-0.0189287\\
0.668	-0.0189287\\
0.67	-0.0189287\\
0.672	-0.0189287\\
0.674	-0.0189288\\
0.676	-0.0189288\\
0.678	-0.0189288\\
0.68	-0.0189288\\
0.682	-0.0189288\\
0.684	-0.0189289\\
0.686	-0.0189289\\
0.688	-0.0189289\\
0.69	-0.0189289\\
0.692	-0.0189289\\
0.694	-0.0189289\\
0.696	-0.0189289\\
0.698	-0.018929\\
0.7	-0.018929\\
0.702	-0.018929\\
0.704	-0.018929\\
0.706	-0.018929\\
0.708	-0.018929\\
0.71	-0.018929\\
0.712	-0.018929\\
0.714	-0.018929\\
0.716	-0.018929\\
0.718	-0.018929\\
0.72	-0.018929\\
0.722	-0.018929\\
0.724	-0.0189291\\
0.726	-0.0189291\\
0.728	-0.0189291\\
0.73	-0.0189291\\
0.732	-0.0189291\\
0.734	-0.0189291\\
0.736	-0.0189291\\
0.738	-0.0189291\\
0.74	-0.0189291\\
0.742	-0.0189291\\
0.744	-0.0189291\\
0.746	-0.0189291\\
0.748	-0.0189291\\
0.75	-0.0189291\\
0.752	-0.0189291\\
0.754	-0.0189291\\
0.756	-0.0189291\\
0.758	-0.0189291\\
0.76	-0.0189291\\
0.762	-0.0189291\\
0.764	-0.0189291\\
0.766	-0.0189291\\
0.768	-0.0189291\\
0.77	-0.0189291\\
0.772	-0.0189291\\
0.774	-0.0189291\\
0.776	-0.0189291\\
0.778	-0.0189291\\
0.78	-0.0189291\\
0.782	-0.0189291\\
0.784	-0.0189291\\
0.786	-0.0189291\\
0.788	-0.0189291\\
0.79	-0.0189291\\
0.792	-0.0189291\\
0.794	-0.0189291\\
0.796	-0.0189291\\
0.798	-0.0189291\\
0.8	-0.0189291\\
0.802	-0.0189291\\
0.804	-0.0189291\\
0.806	-0.0189291\\
0.808	-0.0189291\\
0.81	-0.0189291\\
0.812	-0.0189291\\
0.814	-0.0189291\\
0.816	-0.0189291\\
0.818	-0.0189291\\
0.82	-0.0189291\\
0.822	-0.0189291\\
0.824	-0.0189291\\
0.826	-0.0189291\\
0.828	-0.0189291\\
0.83	-0.0189291\\
0.832	-0.0189291\\
0.834	-0.0189291\\
0.836	-0.0189291\\
0.838	-0.0189291\\
0.84	-0.0189291\\
0.842	-0.0189291\\
0.844	-0.0189291\\
0.846	-0.0189291\\
0.848	-0.0189291\\
0.85	-0.0189291\\
0.852	-0.0189291\\
0.854	-0.0189291\\
0.856	-0.0189291\\
0.858	-0.0189291\\
0.86	-0.0189291\\
0.862	-0.0189291\\
0.864	-0.0189291\\
0.866	-0.0189291\\
0.868	-0.0189291\\
0.87	-0.0189291\\
0.872	-0.0189291\\
0.874	-0.0189291\\
0.876	-0.0189291\\
0.878	-0.0189291\\
0.88	-0.0189291\\
0.882	-0.0189291\\
0.884	-0.0189291\\
0.886	-0.0189291\\
0.888	-0.0189291\\
0.89	-0.0189291\\
0.892	-0.0189291\\
0.894	-0.0189291\\
0.896	-0.0189291\\
0.898	-0.0189291\\
0.9	-0.0189291\\
0.902	-0.0189291\\
0.904	-0.0189291\\
0.906	-0.0189291\\
0.908	-0.0189291\\
0.91	-0.0189291\\
0.912	-0.0189291\\
0.914	-0.0189291\\
0.916	-0.0189291\\
0.918	-0.0189291\\
0.92	-0.0189291\\
0.922	-0.0189291\\
0.924	-0.0189291\\
0.926	-0.0189291\\
0.928	-0.0189291\\
0.93	-0.0189291\\
0.932	-0.0189291\\
0.934	-0.0189291\\
0.936	-0.0189291\\
0.938	-0.0189291\\
0.94	-0.0189291\\
0.942	-0.0189291\\
0.944	-0.0189291\\
0.946	-0.0189291\\
0.948	-0.0189291\\
0.95	-0.0189291\\
0.952	-0.0189291\\
0.954	-0.0189291\\
0.956	-0.0189291\\
0.958	-0.0189291\\
0.96	-0.0189291\\
0.962	-0.0189291\\
0.964	-0.0189291\\
0.966	-0.0189291\\
0.968	-0.0189291\\
0.97	-0.0189291\\
0.972	-0.0189291\\
0.974	-0.0189291\\
0.976	-0.0189291\\
0.978	-0.0189291\\
0.98	-0.0189291\\
0.982	-0.0189291\\
0.984	-0.0189291\\
0.986	-0.0189291\\
0.988	-0.0189291\\
0.99	-0.0189291\\
0.992	-0.0189291\\
0.994	-0.0189291\\
0.996	-0.0189291\\
0.998	-0.0189291\\
1	-0.0189291\\
1.002	-0.0189291\\
1.004	-0.0189291\\
1.006	-0.0189291\\
1.008	-0.0189291\\
1.01	-0.0189291\\
1.012	-0.0189291\\
1.014	-0.0189291\\
1.016	-0.0189291\\
1.018	-0.0189291\\
1.02	-0.0189291\\
1.022	-0.0189291\\
1.024	-0.0189292\\
1.026	-0.0189292\\
1.028	-0.0189292\\
1.03	-0.0189292\\
1.032	-0.0189292\\
1.034	-0.0189292\\
1.036	-0.0189292\\
1.038	-0.0189292\\
1.04	-0.0189292\\
1.042	-0.0189292\\
1.044	-0.0189292\\
1.046	-0.0189292\\
1.048	-0.0189292\\
1.05	-0.0189292\\
1.052	-0.0189292\\
1.054	-0.0189292\\
1.056	-0.0189292\\
1.058	-0.0189292\\
1.06	-0.0189292\\
1.062	-0.0189292\\
1.064	-0.0189292\\
1.066	-0.0189292\\
1.068	-0.0189292\\
1.07	-0.0189292\\
1.072	-0.0189292\\
1.074	-0.0189292\\
1.076	-0.0189292\\
1.078	-0.0189292\\
1.08	-0.0189292\\
1.082	-0.0189292\\
1.084	-0.0189292\\
1.086	-0.0189292\\
1.088	-0.0189292\\
1.09	-0.0189292\\
1.092	-0.0189292\\
1.094	-0.0189292\\
1.096	-0.0189292\\
1.098	-0.0189292\\
1.1	-0.0189292\\
1.102	-0.0189292\\
1.104	-0.0189292\\
1.106	-0.0189292\\
1.108	-0.0189292\\
1.11	-0.0189292\\
1.112	-0.0189292\\
1.114	-0.0189292\\
1.116	-0.0189292\\
1.118	-0.0189292\\
1.12	-0.0189292\\
1.122	-0.0189292\\
1.124	-0.0189292\\
1.126	-0.0189292\\
1.128	-0.0189292\\
1.13	-0.0189292\\
1.132	-0.0189292\\
1.134	-0.0189292\\
1.136	-0.0189292\\
1.138	-0.0189292\\
1.14	-0.0189292\\
1.142	-0.0189292\\
1.144	-0.0189292\\
1.146	-0.0189292\\
1.148	-0.0189292\\
1.15	-0.0189292\\
1.152	-0.0189292\\
1.154	-0.0189292\\
1.156	-0.0189292\\
1.158	-0.0189292\\
1.16	-0.0189292\\
1.162	-0.0189292\\
1.164	-0.0189292\\
1.166	-0.0189292\\
1.168	-0.0189292\\
1.17	-0.0189292\\
1.172	-0.0189292\\
1.174	-0.0189292\\
1.176	-0.0189292\\
1.178	-0.0189292\\
1.18	-0.0189292\\
1.182	-0.0189292\\
1.184	-0.0189292\\
1.186	-0.0189292\\
1.188	-0.0189292\\
1.19	-0.0189292\\
1.192	-0.0189292\\
1.194	-0.0189292\\
1.196	-0.0189292\\
1.198	-0.0189292\\
1.2	-0.0189292\\
1.202	-0.0189292\\
1.204	-0.0189292\\
1.206	-0.0189292\\
1.208	-0.0189292\\
1.21	-0.0189292\\
1.212	-0.0189292\\
1.214	-0.0189292\\
1.216	-0.0189292\\
1.218	-0.0189292\\
1.22	-0.0189292\\
1.222	-0.0189292\\
1.224	-0.0189292\\
1.226	-0.0189292\\
1.228	-0.0189292\\
1.23	-0.0189292\\
1.232	-0.0189292\\
1.234	-0.0189292\\
1.236	-0.0189292\\
1.238	-0.0189292\\
1.24	-0.0189292\\
1.242	-0.0189292\\
1.244	-0.0189292\\
1.246	-0.0189292\\
1.248	-0.0189292\\
1.25	-0.0189292\\
1.252	-0.0189292\\
1.254	-0.0189292\\
1.256	-0.0189292\\
1.258	-0.0189292\\
1.26	-0.0189292\\
1.262	-0.0189292\\
1.264	-0.0189292\\
1.266	-0.0189292\\
1.268	-0.0189292\\
1.27	-0.0189292\\
1.272	-0.0189292\\
1.274	-0.0189292\\
1.276	-0.0189292\\
1.278	-0.0189292\\
1.28	-0.0189292\\
1.282	-0.0189292\\
1.284	-0.0189292\\
1.286	-0.0189292\\
1.288	-0.0189292\\
1.29	-0.0189292\\
1.292	-0.0189292\\
1.294	-0.0189292\\
1.296	-0.0189292\\
1.298	-0.0189292\\
1.3	-0.0189292\\
1.302	-0.0189292\\
1.304	-0.0189292\\
1.306	-0.0189292\\
1.308	-0.0189292\\
1.31	-0.0189292\\
1.312	-0.0189292\\
1.314	-0.0189292\\
1.316	-0.0189292\\
1.318	-0.0189292\\
1.32	-0.0189292\\
1.322	-0.0189292\\
1.324	-0.0189292\\
1.326	-0.0189292\\
1.328	-0.0189292\\
1.33	-0.0189292\\
1.332	-0.0189292\\
1.334	-0.0189292\\
1.336	-0.0189292\\
1.338	-0.0189292\\
1.34	-0.0189292\\
1.342	-0.0189292\\
1.344	-0.0189292\\
1.346	-0.0189292\\
1.348	-0.0189292\\
1.35	-0.0189292\\
1.352	-0.0189292\\
1.354	-0.0189292\\
1.356	-0.0189292\\
1.358	-0.0189292\\
1.36	-0.0189292\\
1.362	-0.0189292\\
1.364	-0.0189292\\
1.366	-0.0189292\\
1.368	-0.0189292\\
1.37	-0.0189292\\
1.372	-0.0189292\\
1.374	-0.0189292\\
1.376	-0.0189292\\
1.378	-0.0189292\\
1.38	-0.0189292\\
1.382	-0.0189292\\
1.384	-0.0189292\\
1.386	-0.0189292\\
1.388	-0.0189292\\
1.39	-0.0189292\\
1.392	-0.0189292\\
1.394	-0.0189292\\
1.396	-0.0189292\\
1.398	-0.0189292\\
1.4	-0.0189292\\
1.402	-0.0189292\\
1.404	-0.0189292\\
1.406	-0.0189292\\
1.408	-0.0189292\\
1.41	-0.0189292\\
1.412	-0.0189292\\
1.414	-0.0189292\\
1.416	-0.0189292\\
1.418	-0.0189292\\
1.42	-0.0189292\\
1.422	-0.0189292\\
1.424	-0.0189292\\
1.426	-0.0189292\\
1.428	-0.0189292\\
1.43	-0.0189292\\
1.432	-0.0189292\\
1.434	-0.0189292\\
1.436	-0.0189292\\
1.438	-0.0189292\\
1.44	-0.0189292\\
1.442	-0.0189292\\
1.444	-0.0189292\\
1.446	-0.0189292\\
1.448	-0.0189292\\
1.45	-0.0189292\\
1.452	-0.0189292\\
1.454	-0.0189292\\
1.456	-0.0189292\\
1.458	-0.0189292\\
1.46	-0.0189292\\
1.462	-0.0189292\\
1.464	-0.0189292\\
1.466	-0.0189292\\
1.468	-0.0189292\\
1.47	-0.0189292\\
1.472	-0.0189292\\
1.474	-0.0189292\\
1.476	-0.0189292\\
1.478	-0.0189292\\
1.48	-0.0189292\\
1.482	-0.0189292\\
1.484	-0.0189292\\
1.486	-0.0189292\\
1.488	-0.0189292\\
1.49	-0.0189292\\
1.492	-0.0189292\\
1.494	-0.0189292\\
1.496	-0.0189292\\
1.498	-0.0189292\\
1.5	-0.0189292\\
1.502	-0.0189292\\
1.504	-0.0189292\\
1.506	-0.0189292\\
1.508	-0.0189292\\
1.51	-0.0189292\\
1.512	-0.0189292\\
1.514	-0.0189292\\
1.516	-0.0189292\\
1.518	-0.0189292\\
1.52	-0.0189292\\
1.522	-0.0189292\\
1.524	-0.0189292\\
1.526	-0.0189292\\
1.528	-0.0189292\\
1.53	-0.0189292\\
1.532	-0.0189292\\
1.534	-0.0189292\\
1.536	-0.0189292\\
1.538	-0.0189292\\
1.54	-0.0189292\\
1.542	-0.0189292\\
1.544	-0.0189292\\
1.546	-0.0189292\\
1.548	-0.0189292\\
1.55	-0.0189292\\
1.552	-0.0189292\\
1.554	-0.0189292\\
1.556	-0.0189292\\
1.558	-0.0189292\\
1.56	-0.0189292\\
1.562	-0.0189292\\
1.564	-0.0189292\\
1.566	-0.0189292\\
1.568	-0.0189292\\
1.57	-0.0189292\\
1.572	-0.0189292\\
1.574	-0.0189292\\
1.576	-0.0189292\\
1.578	-0.0189292\\
1.58	-0.0189292\\
1.582	-0.0189292\\
1.584	-0.0189292\\
1.586	-0.0189292\\
1.588	-0.0189292\\
1.59	-0.0189292\\
1.592	-0.0189292\\
1.594	-0.0189292\\
1.596	-0.0189292\\
1.598	-0.0189292\\
1.6	-0.0189292\\
1.602	-0.0189292\\
1.604	-0.0189292\\
1.606	-0.0189292\\
1.608	-0.0189292\\
1.61	-0.0189292\\
1.612	-0.0189292\\
1.614	-0.0189292\\
1.616	-0.0189292\\
1.618	-0.0189292\\
1.62	-0.0189292\\
1.622	-0.0189292\\
1.624	-0.0189292\\
1.626	-0.0189292\\
1.628	-0.0189292\\
1.63	-0.0189292\\
1.632	-0.0189292\\
1.634	-0.0189292\\
1.636	-0.0189292\\
1.638	-0.0189292\\
1.64	-0.0189292\\
1.642	-0.0189292\\
1.644	-0.0189292\\
1.646	-0.0189292\\
1.648	-0.0189292\\
1.65	-0.0189292\\
1.652	-0.0189292\\
1.654	-0.0189292\\
1.656	-0.0189292\\
1.658	-0.0189292\\
1.66	-0.0189292\\
1.662	-0.0189292\\
1.664	-0.0189292\\
1.666	-0.0189292\\
1.668	-0.0189292\\
1.67	-0.0189292\\
1.672	-0.0189292\\
1.674	-0.0189292\\
1.676	-0.0189292\\
1.678	-0.0189292\\
1.68	-0.0189292\\
1.682	-0.0189292\\
1.684	-0.0189292\\
1.686	-0.0189292\\
1.688	-0.0189292\\
1.69	-0.0189292\\
1.692	-0.0189292\\
1.694	-0.0189292\\
1.696	-0.0189292\\
1.698	-0.0189292\\
1.7	-0.0189292\\
1.702	-0.0189292\\
1.704	-0.0189292\\
1.706	-0.0189292\\
1.708	-0.0189292\\
1.71	-0.0189292\\
1.712	-0.0189292\\
1.714	-0.0189292\\
1.716	-0.0189292\\
1.718	-0.0189292\\
1.72	-0.0189292\\
1.722	-0.0189292\\
1.724	-0.0189292\\
1.726	-0.0189292\\
1.728	-0.0189292\\
1.73	-0.0189292\\
1.732	-0.0189292\\
1.734	-0.0189292\\
1.736	-0.0189292\\
1.738	-0.0189292\\
1.74	-0.0189292\\
1.742	-0.0189292\\
1.744	-0.0189292\\
1.746	-0.0189292\\
1.748	-0.0189292\\
1.75	-0.0189292\\
1.752	-0.0189292\\
1.754	-0.0189292\\
1.756	-0.0189292\\
1.758	-0.0189292\\
1.76	-0.0189292\\
1.762	-0.0189292\\
1.764	-0.0189292\\
1.766	-0.0189292\\
1.768	-0.0189292\\
1.77	-0.0189292\\
1.772	-0.0189292\\
1.774	-0.0189292\\
1.776	-0.0189292\\
1.778	-0.0189292\\
1.78	-0.0189292\\
1.782	-0.0189292\\
1.784	-0.0189292\\
1.786	-0.0189292\\
1.788	-0.0189292\\
1.79	-0.0189292\\
1.792	-0.0189292\\
1.794	-0.0189292\\
1.796	-0.0189292\\
1.798	-0.0189292\\
1.8	-0.0189292\\
1.802	-0.0189292\\
1.804	-0.0189292\\
1.806	-0.0189292\\
1.808	-0.0189292\\
1.81	-0.0189292\\
1.812	-0.0189292\\
1.814	-0.0189292\\
1.816	-0.0189292\\
1.818	-0.0189292\\
1.82	-0.0189292\\
1.822	-0.0189292\\
1.824	-0.0189292\\
1.826	-0.0189292\\
1.828	-0.0189292\\
1.83	-0.0189292\\
1.832	-0.0189292\\
1.834	-0.0189292\\
1.836	-0.0189292\\
1.838	-0.0189292\\
1.84	-0.0189292\\
1.842	-0.0189292\\
1.844	-0.0189292\\
1.846	-0.0189292\\
1.848	-0.0189292\\
1.85	-0.0189292\\
1.852	-0.0189292\\
1.854	-0.0189292\\
1.856	-0.0189292\\
1.858	-0.0189292\\
1.86	-0.0189292\\
1.862	-0.0189292\\
1.864	-0.0189292\\
1.866	-0.0189292\\
1.868	-0.0189292\\
1.87	-0.0189292\\
1.872	-0.0189292\\
1.874	-0.0189292\\
1.876	-0.0189292\\
1.878	-0.0189292\\
1.88	-0.0189292\\
1.882	-0.0189292\\
1.884	-0.0189292\\
1.886	-0.0189292\\
1.888	-0.0189292\\
1.89	-0.0189292\\
1.892	-0.0189292\\
1.894	-0.0189292\\
1.896	-0.0189292\\
1.898	-0.0189292\\
1.9	-0.0189292\\
1.902	-0.0189292\\
1.904	-0.0189292\\
1.906	-0.0189292\\
1.908	-0.0189292\\
1.91	-0.0189292\\
1.912	-0.0189292\\
1.914	-0.0189292\\
1.916	-0.0189292\\
1.918	-0.0189292\\
1.92	-0.0189292\\
1.922	-0.0189292\\
1.924	-0.0189292\\
1.926	-0.0189292\\
1.928	-0.0189292\\
1.93	-0.0189292\\
1.932	-0.0189292\\
1.934	-0.0189292\\
1.936	-0.0189292\\
1.938	-0.0189292\\
1.94	-0.0189292\\
1.942	-0.0189292\\
1.944	-0.0189292\\
1.946	-0.0189292\\
1.948	-0.0189292\\
1.95	-0.0189292\\
1.952	-0.0189292\\
1.954	-0.0189292\\
1.956	-0.0189292\\
1.958	-0.0189292\\
1.96	-0.0189292\\
1.962	-0.0189292\\
1.964	-0.0189292\\
1.966	-0.0189292\\
1.968	-0.0189292\\
1.97	-0.0189292\\
1.972	-0.0189292\\
1.974	-0.0189292\\
1.976	-0.0189292\\
1.978	-0.0189292\\
1.98	-0.0189292\\
1.982	-0.0189292\\
1.984	-0.0189292\\
1.986	-0.0189292\\
1.988	-0.0189292\\
1.99	-0.0189292\\
1.992	-0.0189292\\
1.994	-0.0189292\\
1.996	-0.0189292\\
1.998	-0.0189292\\
2	-0.0189292\\
};
\addplot [color=mycolor7, line width=1.0pt, forget plot]
  table[row sep=crcr]{%
-0.024	0\\
-0.022	0\\
-0.02	0\\
-0.018	0\\
-0.016	0\\
-0.014	0\\
-0.012	0\\
-0.01	0\\
-0.008	0\\
-0.006	0\\
-0.004	0\\
-0.002	0\\
0	0\\
0.002	-0.00183141\\
0.004	-0.00337829\\
0.006	-0.00470498\\
0.008	-0.00585993\\
0.01	-0.00687955\\
0.012	-0.00779111\\
0.014	-0.00861509\\
0.016	-0.00936686\\
0.018	-0.0100581\\
0.02	-0.0106975\\
0.022	-0.0112921\\
0.024	-0.0118471\\
0.026	-0.0123667\\
0.028	-0.0128543\\
0.03	-0.0133127\\
0.032	-0.0137441\\
0.034	-0.0141506\\
0.036	-0.0145337\\
0.038	-0.0148949\\
0.04	-0.0152355\\
0.042	-0.0155565\\
0.044	-0.0158591\\
0.046	-0.0161441\\
0.048	-0.0164124\\
0.05	-0.0166647\\
0.052	-0.0169019\\
0.054	-0.0171246\\
0.056	-0.0173336\\
0.058	-0.0175296\\
0.06	-0.0177131\\
0.062	-0.0178847\\
0.064	-0.0180452\\
0.066	-0.018195\\
0.068	-0.0183348\\
0.07	-0.018465\\
0.072	-0.0185861\\
0.074	-0.0186987\\
0.076	-0.0188031\\
0.078	-0.0188999\\
0.08	-0.0189895\\
0.082	-0.0190721\\
0.084	-0.0191482\\
0.086	-0.0192182\\
0.088	-0.0192823\\
0.09	-0.0193409\\
0.092	-0.0193941\\
0.094	-0.0194424\\
0.096	-0.019486\\
0.098	-0.019525\\
0.1	-0.0195597\\
0.102	-0.0195903\\
0.104	-0.0196171\\
0.106	-0.0196402\\
0.108	-0.0196597\\
0.11	-0.019676\\
0.112	-0.019689\\
0.114	-0.019699\\
0.116	-0.0197062\\
0.118	-0.0197107\\
0.12	-0.0197125\\
0.122	-0.019712\\
0.124	-0.0197091\\
0.126	-0.0197041\\
0.128	-0.019697\\
0.13	-0.0196881\\
0.132	-0.0196773\\
0.134	-0.0196649\\
0.136	-0.0196509\\
0.138	-0.0196355\\
0.14	-0.0196188\\
0.142	-0.0196009\\
0.144	-0.0195819\\
0.146	-0.019562\\
0.148	-0.0195412\\
0.15	-0.0195196\\
0.152	-0.0194974\\
0.154	-0.0194747\\
0.156	-0.0194514\\
0.158	-0.0194279\\
0.16	-0.019404\\
0.162	-0.01938\\
0.164	-0.0193558\\
0.166	-0.0193317\\
0.168	-0.0193076\\
0.17	-0.0192836\\
0.172	-0.0192597\\
0.174	-0.0192362\\
0.176	-0.0192129\\
0.178	-0.01919\\
0.18	-0.0191675\\
0.182	-0.0191454\\
0.184	-0.0191238\\
0.186	-0.0191027\\
0.188	-0.0190822\\
0.19	-0.0190622\\
0.192	-0.0190428\\
0.194	-0.0190241\\
0.196	-0.019006\\
0.198	-0.0189886\\
0.2	-0.0189718\\
0.202	-0.0189557\\
0.204	-0.0189403\\
0.206	-0.0189256\\
0.208	-0.0189116\\
0.21	-0.0188983\\
0.212	-0.0188856\\
0.214	-0.0188737\\
0.216	-0.0188624\\
0.218	-0.0188519\\
0.22	-0.018842\\
0.222	-0.0188327\\
0.224	-0.0188241\\
0.226	-0.0188162\\
0.228	-0.0188089\\
0.23	-0.0188022\\
0.232	-0.0187961\\
0.234	-0.0187906\\
0.236	-0.0187857\\
0.238	-0.0187813\\
0.24	-0.0187774\\
0.242	-0.0187741\\
0.244	-0.0187713\\
0.246	-0.018769\\
0.248	-0.0187671\\
0.25	-0.0187657\\
0.252	-0.0187647\\
0.254	-0.0187642\\
0.256	-0.018764\\
0.258	-0.0187641\\
0.26	-0.0187647\\
0.262	-0.0187655\\
0.264	-0.0187667\\
0.266	-0.0187682\\
0.268	-0.0187699\\
0.27	-0.0187719\\
0.272	-0.0187741\\
0.274	-0.0187765\\
0.276	-0.0187791\\
0.278	-0.018782\\
0.28	-0.0187849\\
0.282	-0.0187881\\
0.284	-0.0187913\\
0.286	-0.0187947\\
0.288	-0.0187981\\
0.29	-0.0188017\\
0.292	-0.0188053\\
0.294	-0.018809\\
0.296	-0.0188128\\
0.298	-0.0188165\\
0.3	-0.0188203\\
0.302	-0.0188242\\
0.304	-0.018828\\
0.306	-0.0188318\\
0.308	-0.0188356\\
0.31	-0.0188394\\
0.312	-0.0188431\\
0.314	-0.0188468\\
0.316	-0.0188504\\
0.318	-0.018854\\
0.32	-0.0188576\\
0.322	-0.0188611\\
0.324	-0.0188645\\
0.326	-0.0188678\\
0.328	-0.0188711\\
0.33	-0.0188742\\
0.332	-0.0188773\\
0.334	-0.0188803\\
0.336	-0.0188832\\
0.338	-0.0188861\\
0.34	-0.0188888\\
0.342	-0.0188914\\
0.344	-0.018894\\
0.346	-0.0188964\\
0.348	-0.0188988\\
0.35	-0.018901\\
0.352	-0.0189032\\
0.354	-0.0189052\\
0.356	-0.0189072\\
0.358	-0.0189091\\
0.36	-0.0189109\\
0.362	-0.0189126\\
0.364	-0.0189142\\
0.366	-0.0189157\\
0.368	-0.0189171\\
0.37	-0.0189185\\
0.372	-0.0189197\\
0.374	-0.0189209\\
0.376	-0.018922\\
0.378	-0.0189231\\
0.38	-0.0189241\\
0.382	-0.0189249\\
0.384	-0.0189258\\
0.386	-0.0189265\\
0.388	-0.0189272\\
0.39	-0.0189279\\
0.392	-0.0189285\\
0.394	-0.018929\\
0.396	-0.0189295\\
0.398	-0.0189299\\
0.4	-0.0189303\\
0.402	-0.0189306\\
0.404	-0.0189309\\
0.406	-0.0189312\\
0.408	-0.0189314\\
0.41	-0.0189316\\
0.412	-0.0189318\\
0.414	-0.0189319\\
0.416	-0.018932\\
0.418	-0.018932\\
0.42	-0.0189321\\
0.422	-0.0189321\\
0.424	-0.0189321\\
0.426	-0.0189321\\
0.428	-0.018932\\
0.43	-0.018932\\
0.432	-0.0189319\\
0.434	-0.0189318\\
0.436	-0.0189317\\
0.438	-0.0189316\\
0.44	-0.0189315\\
0.442	-0.0189313\\
0.444	-0.0189312\\
0.446	-0.0189311\\
0.448	-0.0189309\\
0.45	-0.0189307\\
0.452	-0.0189306\\
0.454	-0.0189304\\
0.456	-0.0189303\\
0.458	-0.0189301\\
0.46	-0.0189299\\
0.462	-0.0189298\\
0.464	-0.0189296\\
0.466	-0.0189294\\
0.468	-0.0189293\\
0.47	-0.0189291\\
0.472	-0.018929\\
0.474	-0.0189288\\
0.476	-0.0189286\\
0.478	-0.0189285\\
0.48	-0.0189284\\
0.482	-0.0189282\\
0.484	-0.0189281\\
0.486	-0.0189279\\
0.488	-0.0189278\\
0.49	-0.0189277\\
0.492	-0.0189276\\
0.494	-0.0189275\\
0.496	-0.0189273\\
0.498	-0.0189272\\
0.5	-0.0189271\\
0.502	-0.018927\\
0.504	-0.018927\\
0.506	-0.0189269\\
0.508	-0.0189268\\
0.51	-0.0189267\\
0.512	-0.0189267\\
0.514	-0.0189266\\
0.516	-0.0189265\\
0.518	-0.0189265\\
0.52	-0.0189264\\
0.522	-0.0189264\\
0.524	-0.0189263\\
0.526	-0.0189263\\
0.528	-0.0189263\\
0.53	-0.0189262\\
0.532	-0.0189262\\
0.534	-0.0189262\\
0.536	-0.0189262\\
0.538	-0.0189262\\
0.54	-0.0189262\\
0.542	-0.0189262\\
0.544	-0.0189262\\
0.546	-0.0189262\\
0.548	-0.0189262\\
0.55	-0.0189262\\
0.552	-0.0189262\\
0.554	-0.0189262\\
0.556	-0.0189262\\
0.558	-0.0189262\\
0.56	-0.0189262\\
0.562	-0.0189263\\
0.564	-0.0189263\\
0.566	-0.0189263\\
0.568	-0.0189263\\
0.57	-0.0189264\\
0.572	-0.0189264\\
0.574	-0.0189264\\
0.576	-0.0189265\\
0.578	-0.0189265\\
0.58	-0.0189265\\
0.582	-0.0189266\\
0.584	-0.0189266\\
0.586	-0.0189267\\
0.588	-0.0189267\\
0.59	-0.0189267\\
0.592	-0.0189268\\
0.594	-0.0189268\\
0.596	-0.0189269\\
0.598	-0.0189269\\
0.6	-0.0189269\\
0.602	-0.018927\\
0.604	-0.018927\\
0.606	-0.0189271\\
0.608	-0.0189271\\
0.61	-0.0189272\\
0.612	-0.0189272\\
0.614	-0.0189272\\
0.616	-0.0189273\\
0.618	-0.0189273\\
0.62	-0.0189274\\
0.622	-0.0189274\\
0.624	-0.0189275\\
0.626	-0.0189275\\
0.628	-0.0189275\\
0.63	-0.0189276\\
0.632	-0.0189276\\
0.634	-0.0189277\\
0.636	-0.0189277\\
0.638	-0.0189277\\
0.64	-0.0189278\\
0.642	-0.0189278\\
0.644	-0.0189279\\
0.646	-0.0189279\\
0.648	-0.0189279\\
0.65	-0.018928\\
0.652	-0.018928\\
0.654	-0.018928\\
0.656	-0.0189281\\
0.658	-0.0189281\\
0.66	-0.0189282\\
0.662	-0.0189282\\
0.664	-0.0189282\\
0.666	-0.0189282\\
0.668	-0.0189283\\
0.67	-0.0189283\\
0.672	-0.0189283\\
0.674	-0.0189284\\
0.676	-0.0189284\\
0.678	-0.0189284\\
0.68	-0.0189285\\
0.682	-0.0189285\\
0.684	-0.0189285\\
0.686	-0.0189285\\
0.688	-0.0189286\\
0.69	-0.0189286\\
0.692	-0.0189286\\
0.694	-0.0189286\\
0.696	-0.0189286\\
0.698	-0.0189287\\
0.7	-0.0189287\\
0.702	-0.0189287\\
0.704	-0.0189287\\
0.706	-0.0189287\\
0.708	-0.0189288\\
0.71	-0.0189288\\
0.712	-0.0189288\\
0.714	-0.0189288\\
0.716	-0.0189288\\
0.718	-0.0189288\\
0.72	-0.0189289\\
0.722	-0.0189289\\
0.724	-0.0189289\\
0.726	-0.0189289\\
0.728	-0.0189289\\
0.73	-0.0189289\\
0.732	-0.0189289\\
0.734	-0.0189289\\
0.736	-0.0189289\\
0.738	-0.018929\\
0.74	-0.018929\\
0.742	-0.018929\\
0.744	-0.018929\\
0.746	-0.018929\\
0.748	-0.018929\\
0.75	-0.018929\\
0.752	-0.018929\\
0.754	-0.018929\\
0.756	-0.018929\\
0.758	-0.018929\\
0.76	-0.018929\\
0.762	-0.018929\\
0.764	-0.018929\\
0.766	-0.018929\\
0.768	-0.018929\\
0.77	-0.0189291\\
0.772	-0.0189291\\
0.774	-0.0189291\\
0.776	-0.0189291\\
0.778	-0.0189291\\
0.78	-0.0189291\\
0.782	-0.0189291\\
0.784	-0.0189291\\
0.786	-0.0189291\\
0.788	-0.0189291\\
0.79	-0.0189291\\
0.792	-0.0189291\\
0.794	-0.0189291\\
0.796	-0.0189291\\
0.798	-0.0189291\\
0.8	-0.0189291\\
0.802	-0.0189291\\
0.804	-0.0189291\\
0.806	-0.0189291\\
0.808	-0.0189291\\
0.81	-0.0189291\\
0.812	-0.0189291\\
0.814	-0.0189291\\
0.816	-0.0189291\\
0.818	-0.0189291\\
0.82	-0.0189291\\
0.822	-0.0189291\\
0.824	-0.0189291\\
0.826	-0.0189291\\
0.828	-0.0189291\\
0.83	-0.0189291\\
0.832	-0.0189291\\
0.834	-0.0189291\\
0.836	-0.0189291\\
0.838	-0.0189291\\
0.84	-0.0189291\\
0.842	-0.0189291\\
0.844	-0.0189291\\
0.846	-0.0189291\\
0.848	-0.0189291\\
0.85	-0.0189291\\
0.852	-0.0189291\\
0.854	-0.0189291\\
0.856	-0.0189291\\
0.858	-0.0189291\\
0.86	-0.0189291\\
0.862	-0.0189291\\
0.864	-0.0189291\\
0.866	-0.0189291\\
0.868	-0.0189291\\
0.87	-0.0189291\\
0.872	-0.0189291\\
0.874	-0.0189291\\
0.876	-0.0189291\\
0.878	-0.0189291\\
0.88	-0.0189291\\
0.882	-0.0189291\\
0.884	-0.0189291\\
0.886	-0.0189291\\
0.888	-0.0189291\\
0.89	-0.0189291\\
0.892	-0.0189291\\
0.894	-0.0189291\\
0.896	-0.0189291\\
0.898	-0.0189291\\
0.9	-0.0189291\\
0.902	-0.0189291\\
0.904	-0.0189291\\
0.906	-0.0189291\\
0.908	-0.0189291\\
0.91	-0.0189291\\
0.912	-0.0189291\\
0.914	-0.0189291\\
0.916	-0.0189291\\
0.918	-0.0189291\\
0.92	-0.0189291\\
0.922	-0.0189291\\
0.924	-0.0189291\\
0.926	-0.0189291\\
0.928	-0.0189291\\
0.93	-0.0189291\\
0.932	-0.0189291\\
0.934	-0.0189291\\
0.936	-0.0189291\\
0.938	-0.0189291\\
0.94	-0.0189291\\
0.942	-0.0189291\\
0.944	-0.0189291\\
0.946	-0.0189291\\
0.948	-0.0189291\\
0.95	-0.0189291\\
0.952	-0.0189291\\
0.954	-0.0189291\\
0.956	-0.0189291\\
0.958	-0.0189291\\
0.96	-0.0189291\\
0.962	-0.0189291\\
0.964	-0.0189291\\
0.966	-0.0189291\\
0.968	-0.0189291\\
0.97	-0.0189291\\
0.972	-0.0189291\\
0.974	-0.0189291\\
0.976	-0.0189291\\
0.978	-0.0189291\\
0.98	-0.0189291\\
0.982	-0.0189291\\
0.984	-0.0189291\\
0.986	-0.0189291\\
0.988	-0.0189291\\
0.99	-0.0189291\\
0.992	-0.0189291\\
0.994	-0.0189291\\
0.996	-0.0189291\\
0.998	-0.0189291\\
1	-0.0189291\\
1.002	-0.0189291\\
1.004	-0.0189291\\
1.006	-0.0189291\\
1.008	-0.0189291\\
1.01	-0.0189291\\
1.012	-0.0189291\\
1.014	-0.0189291\\
1.016	-0.0189291\\
1.018	-0.0189291\\
1.02	-0.0189291\\
1.022	-0.0189291\\
1.024	-0.0189291\\
1.026	-0.0189291\\
1.028	-0.0189291\\
1.03	-0.0189291\\
1.032	-0.0189291\\
1.034	-0.0189291\\
1.036	-0.0189291\\
1.038	-0.0189291\\
1.04	-0.0189291\\
1.042	-0.0189291\\
1.044	-0.0189291\\
1.046	-0.0189291\\
1.048	-0.0189291\\
1.05	-0.0189291\\
1.052	-0.0189291\\
1.054	-0.0189291\\
1.056	-0.0189291\\
1.058	-0.0189291\\
1.06	-0.0189291\\
1.062	-0.0189291\\
1.064	-0.0189292\\
1.066	-0.0189292\\
1.068	-0.0189292\\
1.07	-0.0189292\\
1.072	-0.0189292\\
1.074	-0.0189292\\
1.076	-0.0189292\\
1.078	-0.0189292\\
1.08	-0.0189292\\
1.082	-0.0189292\\
1.084	-0.0189292\\
1.086	-0.0189292\\
1.088	-0.0189292\\
1.09	-0.0189292\\
1.092	-0.0189292\\
1.094	-0.0189292\\
1.096	-0.0189292\\
1.098	-0.0189292\\
1.1	-0.0189292\\
1.102	-0.0189292\\
1.104	-0.0189292\\
1.106	-0.0189292\\
1.108	-0.0189292\\
1.11	-0.0189292\\
1.112	-0.0189292\\
1.114	-0.0189292\\
1.116	-0.0189292\\
1.118	-0.0189292\\
1.12	-0.0189292\\
1.122	-0.0189292\\
1.124	-0.0189292\\
1.126	-0.0189292\\
1.128	-0.0189292\\
1.13	-0.0189292\\
1.132	-0.0189292\\
1.134	-0.0189292\\
1.136	-0.0189292\\
1.138	-0.0189292\\
1.14	-0.0189292\\
1.142	-0.0189292\\
1.144	-0.0189292\\
1.146	-0.0189292\\
1.148	-0.0189292\\
1.15	-0.0189292\\
1.152	-0.0189292\\
1.154	-0.0189292\\
1.156	-0.0189292\\
1.158	-0.0189292\\
1.16	-0.0189292\\
1.162	-0.0189292\\
1.164	-0.0189292\\
1.166	-0.0189292\\
1.168	-0.0189292\\
1.17	-0.0189292\\
1.172	-0.0189292\\
1.174	-0.0189292\\
1.176	-0.0189292\\
1.178	-0.0189292\\
1.18	-0.0189292\\
1.182	-0.0189292\\
1.184	-0.0189292\\
1.186	-0.0189292\\
1.188	-0.0189292\\
1.19	-0.0189292\\
1.192	-0.0189292\\
1.194	-0.0189292\\
1.196	-0.0189292\\
1.198	-0.0189292\\
1.2	-0.0189292\\
1.202	-0.0189292\\
1.204	-0.0189292\\
1.206	-0.0189292\\
1.208	-0.0189292\\
1.21	-0.0189292\\
1.212	-0.0189292\\
1.214	-0.0189292\\
1.216	-0.0189292\\
1.218	-0.0189292\\
1.22	-0.0189292\\
1.222	-0.0189292\\
1.224	-0.0189292\\
1.226	-0.0189292\\
1.228	-0.0189292\\
1.23	-0.0189292\\
1.232	-0.0189292\\
1.234	-0.0189292\\
1.236	-0.0189292\\
1.238	-0.0189292\\
1.24	-0.0189292\\
1.242	-0.0189292\\
1.244	-0.0189292\\
1.246	-0.0189292\\
1.248	-0.0189292\\
1.25	-0.0189292\\
1.252	-0.0189292\\
1.254	-0.0189292\\
1.256	-0.0189292\\
1.258	-0.0189292\\
1.26	-0.0189292\\
1.262	-0.0189292\\
1.264	-0.0189292\\
1.266	-0.0189292\\
1.268	-0.0189292\\
1.27	-0.0189292\\
1.272	-0.0189292\\
1.274	-0.0189292\\
1.276	-0.0189292\\
1.278	-0.0189292\\
1.28	-0.0189292\\
1.282	-0.0189292\\
1.284	-0.0189292\\
1.286	-0.0189292\\
1.288	-0.0189292\\
1.29	-0.0189292\\
1.292	-0.0189292\\
1.294	-0.0189292\\
1.296	-0.0189292\\
1.298	-0.0189292\\
1.3	-0.0189292\\
1.302	-0.0189292\\
1.304	-0.0189292\\
1.306	-0.0189292\\
1.308	-0.0189292\\
1.31	-0.0189292\\
1.312	-0.0189292\\
1.314	-0.0189292\\
1.316	-0.0189292\\
1.318	-0.0189292\\
1.32	-0.0189292\\
1.322	-0.0189292\\
1.324	-0.0189292\\
1.326	-0.0189292\\
1.328	-0.0189292\\
1.33	-0.0189292\\
1.332	-0.0189292\\
1.334	-0.0189292\\
1.336	-0.0189292\\
1.338	-0.0189292\\
1.34	-0.0189292\\
1.342	-0.0189292\\
1.344	-0.0189292\\
1.346	-0.0189292\\
1.348	-0.0189292\\
1.35	-0.0189292\\
1.352	-0.0189292\\
1.354	-0.0189292\\
1.356	-0.0189292\\
1.358	-0.0189292\\
1.36	-0.0189292\\
1.362	-0.0189292\\
1.364	-0.0189292\\
1.366	-0.0189292\\
1.368	-0.0189292\\
1.37	-0.0189292\\
1.372	-0.0189292\\
1.374	-0.0189292\\
1.376	-0.0189292\\
1.378	-0.0189292\\
1.38	-0.0189292\\
1.382	-0.0189292\\
1.384	-0.0189292\\
1.386	-0.0189292\\
1.388	-0.0189292\\
1.39	-0.0189292\\
1.392	-0.0189292\\
1.394	-0.0189292\\
1.396	-0.0189292\\
1.398	-0.0189292\\
1.4	-0.0189292\\
1.402	-0.0189292\\
1.404	-0.0189292\\
1.406	-0.0189292\\
1.408	-0.0189292\\
1.41	-0.0189292\\
1.412	-0.0189292\\
1.414	-0.0189292\\
1.416	-0.0189292\\
1.418	-0.0189292\\
1.42	-0.0189292\\
1.422	-0.0189292\\
1.424	-0.0189292\\
1.426	-0.0189292\\
1.428	-0.0189292\\
1.43	-0.0189292\\
1.432	-0.0189292\\
1.434	-0.0189292\\
1.436	-0.0189292\\
1.438	-0.0189292\\
1.44	-0.0189292\\
1.442	-0.0189292\\
1.444	-0.0189292\\
1.446	-0.0189292\\
1.448	-0.0189292\\
1.45	-0.0189292\\
1.452	-0.0189292\\
1.454	-0.0189292\\
1.456	-0.0189292\\
1.458	-0.0189292\\
1.46	-0.0189292\\
1.462	-0.0189292\\
1.464	-0.0189292\\
1.466	-0.0189292\\
1.468	-0.0189292\\
1.47	-0.0189292\\
1.472	-0.0189292\\
1.474	-0.0189292\\
1.476	-0.0189292\\
1.478	-0.0189292\\
1.48	-0.0189292\\
1.482	-0.0189292\\
1.484	-0.0189292\\
1.486	-0.0189292\\
1.488	-0.0189292\\
1.49	-0.0189292\\
1.492	-0.0189292\\
1.494	-0.0189292\\
1.496	-0.0189292\\
1.498	-0.0189292\\
1.5	-0.0189292\\
1.502	-0.0189292\\
1.504	-0.0189292\\
1.506	-0.0189292\\
1.508	-0.0189292\\
1.51	-0.0189292\\
1.512	-0.0189292\\
1.514	-0.0189292\\
1.516	-0.0189292\\
1.518	-0.0189292\\
1.52	-0.0189292\\
1.522	-0.0189292\\
1.524	-0.0189292\\
1.526	-0.0189292\\
1.528	-0.0189292\\
1.53	-0.0189292\\
1.532	-0.0189292\\
1.534	-0.0189292\\
1.536	-0.0189292\\
1.538	-0.0189292\\
1.54	-0.0189292\\
1.542	-0.0189292\\
1.544	-0.0189292\\
1.546	-0.0189292\\
1.548	-0.0189292\\
1.55	-0.0189292\\
1.552	-0.0189292\\
1.554	-0.0189292\\
1.556	-0.0189292\\
1.558	-0.0189292\\
1.56	-0.0189292\\
1.562	-0.0189292\\
1.564	-0.0189292\\
1.566	-0.0189292\\
1.568	-0.0189292\\
1.57	-0.0189292\\
1.572	-0.0189292\\
1.574	-0.0189292\\
1.576	-0.0189292\\
1.578	-0.0189292\\
1.58	-0.0189292\\
1.582	-0.0189292\\
1.584	-0.0189292\\
1.586	-0.0189292\\
1.588	-0.0189292\\
1.59	-0.0189292\\
1.592	-0.0189292\\
1.594	-0.0189292\\
1.596	-0.0189292\\
1.598	-0.0189292\\
1.6	-0.0189292\\
1.602	-0.0189292\\
1.604	-0.0189292\\
1.606	-0.0189292\\
1.608	-0.0189292\\
1.61	-0.0189292\\
1.612	-0.0189292\\
1.614	-0.0189292\\
1.616	-0.0189292\\
1.618	-0.0189292\\
1.62	-0.0189292\\
1.622	-0.0189292\\
1.624	-0.0189292\\
1.626	-0.0189292\\
1.628	-0.0189292\\
1.63	-0.0189292\\
1.632	-0.0189292\\
1.634	-0.0189292\\
1.636	-0.0189292\\
1.638	-0.0189292\\
1.64	-0.0189292\\
1.642	-0.0189292\\
1.644	-0.0189292\\
1.646	-0.0189292\\
1.648	-0.0189292\\
1.65	-0.0189292\\
1.652	-0.0189292\\
1.654	-0.0189292\\
1.656	-0.0189292\\
1.658	-0.0189292\\
1.66	-0.0189292\\
1.662	-0.0189292\\
1.664	-0.0189292\\
1.666	-0.0189292\\
1.668	-0.0189292\\
1.67	-0.0189292\\
1.672	-0.0189292\\
1.674	-0.0189292\\
1.676	-0.0189292\\
1.678	-0.0189292\\
1.68	-0.0189292\\
1.682	-0.0189292\\
1.684	-0.0189292\\
1.686	-0.0189292\\
1.688	-0.0189292\\
1.69	-0.0189292\\
1.692	-0.0189292\\
1.694	-0.0189292\\
1.696	-0.0189292\\
1.698	-0.0189292\\
1.7	-0.0189292\\
1.702	-0.0189292\\
1.704	-0.0189292\\
1.706	-0.0189292\\
1.708	-0.0189292\\
1.71	-0.0189292\\
1.712	-0.0189292\\
1.714	-0.0189292\\
1.716	-0.0189292\\
1.718	-0.0189292\\
1.72	-0.0189292\\
1.722	-0.0189292\\
1.724	-0.0189292\\
1.726	-0.0189292\\
1.728	-0.0189292\\
1.73	-0.0189292\\
1.732	-0.0189292\\
1.734	-0.0189292\\
1.736	-0.0189292\\
1.738	-0.0189292\\
1.74	-0.0189292\\
1.742	-0.0189292\\
1.744	-0.0189292\\
1.746	-0.0189292\\
1.748	-0.0189292\\
1.75	-0.0189292\\
1.752	-0.0189292\\
1.754	-0.0189292\\
1.756	-0.0189292\\
1.758	-0.0189292\\
1.76	-0.0189292\\
1.762	-0.0189292\\
1.764	-0.0189292\\
1.766	-0.0189292\\
1.768	-0.0189292\\
1.77	-0.0189292\\
1.772	-0.0189292\\
1.774	-0.0189292\\
1.776	-0.0189292\\
1.778	-0.0189292\\
1.78	-0.0189292\\
1.782	-0.0189292\\
1.784	-0.0189292\\
1.786	-0.0189292\\
1.788	-0.0189292\\
1.79	-0.0189292\\
1.792	-0.0189292\\
1.794	-0.0189292\\
1.796	-0.0189292\\
1.798	-0.0189292\\
1.8	-0.0189292\\
1.802	-0.0189292\\
1.804	-0.0189292\\
1.806	-0.0189292\\
1.808	-0.0189292\\
1.81	-0.0189292\\
1.812	-0.0189292\\
1.814	-0.0189292\\
1.816	-0.0189292\\
1.818	-0.0189292\\
1.82	-0.0189292\\
1.822	-0.0189292\\
1.824	-0.0189292\\
1.826	-0.0189292\\
1.828	-0.0189292\\
1.83	-0.0189292\\
1.832	-0.0189292\\
1.834	-0.0189292\\
1.836	-0.0189292\\
1.838	-0.0189292\\
1.84	-0.0189292\\
1.842	-0.0189292\\
1.844	-0.0189292\\
1.846	-0.0189292\\
1.848	-0.0189292\\
1.85	-0.0189292\\
1.852	-0.0189292\\
1.854	-0.0189292\\
1.856	-0.0189292\\
1.858	-0.0189292\\
1.86	-0.0189292\\
1.862	-0.0189292\\
1.864	-0.0189292\\
1.866	-0.0189292\\
1.868	-0.0189292\\
1.87	-0.0189292\\
1.872	-0.0189292\\
1.874	-0.0189292\\
1.876	-0.0189292\\
1.878	-0.0189292\\
1.88	-0.0189292\\
1.882	-0.0189292\\
1.884	-0.0189292\\
1.886	-0.0189292\\
1.888	-0.0189292\\
1.89	-0.0189292\\
1.892	-0.0189292\\
1.894	-0.0189292\\
1.896	-0.0189292\\
1.898	-0.0189292\\
1.9	-0.0189292\\
1.902	-0.0189292\\
1.904	-0.0189292\\
1.906	-0.0189292\\
1.908	-0.0189292\\
1.91	-0.0189292\\
1.912	-0.0189292\\
1.914	-0.0189292\\
1.916	-0.0189292\\
1.918	-0.0189292\\
1.92	-0.0189292\\
1.922	-0.0189292\\
1.924	-0.0189292\\
1.926	-0.0189292\\
1.928	-0.0189292\\
1.93	-0.0189292\\
1.932	-0.0189292\\
1.934	-0.0189292\\
1.936	-0.0189292\\
1.938	-0.0189292\\
1.94	-0.0189292\\
1.942	-0.0189292\\
1.944	-0.0189292\\
1.946	-0.0189292\\
1.948	-0.0189292\\
1.95	-0.0189292\\
1.952	-0.0189292\\
1.954	-0.0189292\\
1.956	-0.0189292\\
1.958	-0.0189292\\
1.96	-0.0189292\\
1.962	-0.0189292\\
1.964	-0.0189292\\
1.966	-0.0189292\\
1.968	-0.0189292\\
1.97	-0.0189292\\
1.972	-0.0189292\\
1.974	-0.0189292\\
1.976	-0.0189292\\
1.978	-0.0189292\\
1.98	-0.0189292\\
1.982	-0.0189292\\
1.984	-0.0189292\\
1.986	-0.0189292\\
1.988	-0.0189292\\
1.99	-0.0189292\\
1.992	-0.0189292\\
1.994	-0.0189292\\
1.996	-0.0189292\\
1.998	-0.0189292\\
2	-0.0189292\\
};
\addplot [color=mycolor8, line width=1.0pt, forget plot]
  table[row sep=crcr]{%
-0.024	0\\
-0.022	0\\
-0.02	0\\
-0.018	0\\
-0.016	0\\
-0.014	0\\
-0.012	0\\
-0.01	0\\
-0.008	0\\
-0.006	0\\
-0.004	0\\
-0.002	0\\
0	0\\
0.002	-0.00183144\\
0.004	-0.00337855\\
0.006	-0.00470578\\
0.008	-0.00586167\\
0.01	-0.00688263\\
0.012	-0.00779596\\
0.014	-0.0086221\\
0.016	-0.00937638\\
0.018	-0.0100703\\
0.02	-0.0107128\\
0.022	-0.0113104\\
0.024	-0.0118685\\
0.026	-0.0123911\\
0.028	-0.0128815\\
0.03	-0.0133425\\
0.032	-0.0137761\\
0.034	-0.0141843\\
0.036	-0.0145687\\
0.038	-0.0149306\\
0.04	-0.0152713\\
0.042	-0.0155919\\
0.044	-0.0158933\\
0.046	-0.0161765\\
0.048	-0.0164423\\
0.05	-0.0166917\\
0.052	-0.0169254\\
0.054	-0.0171441\\
0.056	-0.0173487\\
0.058	-0.0175399\\
0.06	-0.0177183\\
0.062	-0.0178848\\
0.064	-0.0180399\\
0.066	-0.0181844\\
0.068	-0.0183188\\
0.07	-0.0184438\\
0.072	-0.0185599\\
0.074	-0.0186677\\
0.076	-0.0187677\\
0.078	-0.0188604\\
0.08	-0.0189461\\
0.082	-0.0190254\\
0.084	-0.0190986\\
0.086	-0.0191661\\
0.088	-0.0192282\\
0.09	-0.0192851\\
0.092	-0.0193372\\
0.094	-0.0193846\\
0.096	-0.0194277\\
0.098	-0.0194667\\
0.1	-0.0195017\\
0.102	-0.0195329\\
0.104	-0.0195604\\
0.106	-0.0195846\\
0.108	-0.0196054\\
0.11	-0.019623\\
0.112	-0.0196377\\
0.114	-0.0196494\\
0.116	-0.0196584\\
0.118	-0.0196647\\
0.12	-0.0196685\\
0.122	-0.01967\\
0.124	-0.0196692\\
0.126	-0.0196662\\
0.128	-0.0196612\\
0.13	-0.0196544\\
0.132	-0.0196458\\
0.134	-0.0196355\\
0.136	-0.0196237\\
0.138	-0.0196104\\
0.14	-0.0195959\\
0.142	-0.0195802\\
0.144	-0.0195634\\
0.146	-0.0195457\\
0.148	-0.0195271\\
0.15	-0.0195077\\
0.152	-0.0194878\\
0.154	-0.0194672\\
0.156	-0.0194462\\
0.158	-0.0194249\\
0.16	-0.0194032\\
0.162	-0.0193814\\
0.164	-0.0193594\\
0.166	-0.0193374\\
0.168	-0.0193154\\
0.17	-0.0192934\\
0.172	-0.0192716\\
0.174	-0.01925\\
0.176	-0.0192286\\
0.178	-0.0192075\\
0.18	-0.0191867\\
0.182	-0.0191663\\
0.184	-0.0191462\\
0.186	-0.0191266\\
0.188	-0.0191075\\
0.19	-0.0190888\\
0.192	-0.0190707\\
0.194	-0.0190531\\
0.196	-0.019036\\
0.198	-0.0190195\\
0.2	-0.0190035\\
0.202	-0.0189882\\
0.204	-0.0189734\\
0.206	-0.0189592\\
0.208	-0.0189457\\
0.21	-0.0189327\\
0.212	-0.0189204\\
0.214	-0.0189086\\
0.216	-0.0188975\\
0.218	-0.018887\\
0.22	-0.018877\\
0.222	-0.0188677\\
0.224	-0.018859\\
0.226	-0.0188508\\
0.228	-0.0188432\\
0.23	-0.0188362\\
0.232	-0.0188297\\
0.234	-0.0188238\\
0.236	-0.0188184\\
0.238	-0.0188135\\
0.24	-0.0188091\\
0.242	-0.0188051\\
0.244	-0.0188017\\
0.246	-0.0187986\\
0.248	-0.0187961\\
0.25	-0.0187939\\
0.252	-0.0187921\\
0.254	-0.0187907\\
0.256	-0.0187897\\
0.258	-0.018789\\
0.26	-0.0187886\\
0.262	-0.0187886\\
0.264	-0.0187888\\
0.266	-0.0187893\\
0.268	-0.0187901\\
0.27	-0.0187912\\
0.272	-0.0187924\\
0.274	-0.0187939\\
0.276	-0.0187956\\
0.278	-0.0187974\\
0.28	-0.0187994\\
0.282	-0.0188016\\
0.284	-0.0188039\\
0.286	-0.0188063\\
0.288	-0.0188089\\
0.29	-0.0188115\\
0.292	-0.0188143\\
0.294	-0.0188171\\
0.296	-0.0188199\\
0.298	-0.0188229\\
0.3	-0.0188258\\
0.302	-0.0188289\\
0.304	-0.0188319\\
0.306	-0.0188349\\
0.308	-0.018838\\
0.31	-0.0188411\\
0.312	-0.0188441\\
0.314	-0.0188471\\
0.316	-0.0188501\\
0.318	-0.0188531\\
0.32	-0.0188561\\
0.322	-0.018859\\
0.324	-0.0188619\\
0.326	-0.0188647\\
0.328	-0.0188675\\
0.33	-0.0188702\\
0.332	-0.0188728\\
0.334	-0.0188754\\
0.336	-0.018878\\
0.338	-0.0188805\\
0.34	-0.0188829\\
0.342	-0.0188852\\
0.344	-0.0188874\\
0.346	-0.0188896\\
0.348	-0.0188918\\
0.35	-0.0188938\\
0.352	-0.0188958\\
0.354	-0.0188977\\
0.356	-0.0188995\\
0.358	-0.0189012\\
0.36	-0.0189029\\
0.362	-0.0189045\\
0.364	-0.0189061\\
0.366	-0.0189075\\
0.368	-0.0189089\\
0.37	-0.0189103\\
0.372	-0.0189115\\
0.374	-0.0189127\\
0.376	-0.0189139\\
0.378	-0.018915\\
0.38	-0.018916\\
0.382	-0.0189169\\
0.384	-0.0189179\\
0.386	-0.0189187\\
0.388	-0.0189195\\
0.39	-0.0189203\\
0.392	-0.018921\\
0.394	-0.0189216\\
0.396	-0.0189222\\
0.398	-0.0189228\\
0.4	-0.0189233\\
0.402	-0.0189238\\
0.404	-0.0189243\\
0.406	-0.0189247\\
0.408	-0.0189251\\
0.41	-0.0189254\\
0.412	-0.0189258\\
0.414	-0.018926\\
0.416	-0.0189263\\
0.418	-0.0189266\\
0.42	-0.0189268\\
0.422	-0.018927\\
0.424	-0.0189271\\
0.426	-0.0189273\\
0.428	-0.0189274\\
0.43	-0.0189275\\
0.432	-0.0189276\\
0.434	-0.0189277\\
0.436	-0.0189278\\
0.438	-0.0189278\\
0.44	-0.0189279\\
0.442	-0.0189279\\
0.444	-0.0189279\\
0.446	-0.0189279\\
0.448	-0.0189279\\
0.45	-0.0189279\\
0.452	-0.0189279\\
0.454	-0.0189279\\
0.456	-0.0189278\\
0.458	-0.0189278\\
0.46	-0.0189278\\
0.462	-0.0189277\\
0.464	-0.0189277\\
0.466	-0.0189276\\
0.468	-0.0189276\\
0.47	-0.0189275\\
0.472	-0.0189274\\
0.474	-0.0189274\\
0.476	-0.0189273\\
0.478	-0.0189273\\
0.48	-0.0189272\\
0.482	-0.0189271\\
0.484	-0.0189271\\
0.486	-0.018927\\
0.488	-0.018927\\
0.49	-0.0189269\\
0.492	-0.0189268\\
0.494	-0.0189268\\
0.496	-0.0189267\\
0.498	-0.0189267\\
0.5	-0.0189266\\
0.502	-0.0189266\\
0.504	-0.0189265\\
0.506	-0.0189265\\
0.508	-0.0189264\\
0.51	-0.0189264\\
0.512	-0.0189263\\
0.514	-0.0189263\\
0.516	-0.0189263\\
0.518	-0.0189262\\
0.52	-0.0189262\\
0.522	-0.0189262\\
0.524	-0.0189261\\
0.526	-0.0189261\\
0.528	-0.0189261\\
0.53	-0.0189261\\
0.532	-0.018926\\
0.534	-0.018926\\
0.536	-0.018926\\
0.538	-0.018926\\
0.54	-0.018926\\
0.542	-0.018926\\
0.544	-0.018926\\
0.546	-0.018926\\
0.548	-0.018926\\
0.55	-0.018926\\
0.552	-0.018926\\
0.554	-0.018926\\
0.556	-0.018926\\
0.558	-0.018926\\
0.56	-0.018926\\
0.562	-0.018926\\
0.564	-0.018926\\
0.566	-0.0189261\\
0.568	-0.0189261\\
0.57	-0.0189261\\
0.572	-0.0189261\\
0.574	-0.0189261\\
0.576	-0.0189262\\
0.578	-0.0189262\\
0.58	-0.0189262\\
0.582	-0.0189262\\
0.584	-0.0189263\\
0.586	-0.0189263\\
0.588	-0.0189263\\
0.59	-0.0189263\\
0.592	-0.0189264\\
0.594	-0.0189264\\
0.596	-0.0189264\\
0.598	-0.0189265\\
0.6	-0.0189265\\
0.602	-0.0189265\\
0.604	-0.0189266\\
0.606	-0.0189266\\
0.608	-0.0189267\\
0.61	-0.0189267\\
0.612	-0.0189267\\
0.614	-0.0189268\\
0.616	-0.0189268\\
0.618	-0.0189268\\
0.62	-0.0189269\\
0.622	-0.0189269\\
0.624	-0.018927\\
0.626	-0.018927\\
0.628	-0.018927\\
0.63	-0.0189271\\
0.632	-0.0189271\\
0.634	-0.0189272\\
0.636	-0.0189272\\
0.638	-0.0189272\\
0.64	-0.0189273\\
0.642	-0.0189273\\
0.644	-0.0189274\\
0.646	-0.0189274\\
0.648	-0.0189274\\
0.65	-0.0189275\\
0.652	-0.0189275\\
0.654	-0.0189275\\
0.656	-0.0189276\\
0.658	-0.0189276\\
0.66	-0.0189277\\
0.662	-0.0189277\\
0.664	-0.0189277\\
0.666	-0.0189278\\
0.668	-0.0189278\\
0.67	-0.0189278\\
0.672	-0.0189279\\
0.674	-0.0189279\\
0.676	-0.0189279\\
0.678	-0.018928\\
0.68	-0.018928\\
0.682	-0.018928\\
0.684	-0.0189281\\
0.686	-0.0189281\\
0.688	-0.0189281\\
0.69	-0.0189282\\
0.692	-0.0189282\\
0.694	-0.0189282\\
0.696	-0.0189282\\
0.698	-0.0189283\\
0.7	-0.0189283\\
0.702	-0.0189283\\
0.704	-0.0189283\\
0.706	-0.0189284\\
0.708	-0.0189284\\
0.71	-0.0189284\\
0.712	-0.0189284\\
0.714	-0.0189285\\
0.716	-0.0189285\\
0.718	-0.0189285\\
0.72	-0.0189285\\
0.722	-0.0189286\\
0.724	-0.0189286\\
0.726	-0.0189286\\
0.728	-0.0189286\\
0.73	-0.0189286\\
0.732	-0.0189286\\
0.734	-0.0189287\\
0.736	-0.0189287\\
0.738	-0.0189287\\
0.74	-0.0189287\\
0.742	-0.0189287\\
0.744	-0.0189287\\
0.746	-0.0189287\\
0.748	-0.0189288\\
0.75	-0.0189288\\
0.752	-0.0189288\\
0.754	-0.0189288\\
0.756	-0.0189288\\
0.758	-0.0189288\\
0.76	-0.0189288\\
0.762	-0.0189288\\
0.764	-0.0189288\\
0.766	-0.0189289\\
0.768	-0.0189289\\
0.77	-0.0189289\\
0.772	-0.0189289\\
0.774	-0.0189289\\
0.776	-0.0189289\\
0.778	-0.0189289\\
0.78	-0.0189289\\
0.782	-0.0189289\\
0.784	-0.0189289\\
0.786	-0.0189289\\
0.788	-0.0189289\\
0.79	-0.0189289\\
0.792	-0.0189289\\
0.794	-0.0189289\\
0.796	-0.0189289\\
0.798	-0.0189289\\
0.8	-0.018929\\
0.802	-0.018929\\
0.804	-0.018929\\
0.806	-0.018929\\
0.808	-0.018929\\
0.81	-0.018929\\
0.812	-0.018929\\
0.814	-0.018929\\
0.816	-0.018929\\
0.818	-0.018929\\
0.82	-0.018929\\
0.822	-0.018929\\
0.824	-0.018929\\
0.826	-0.018929\\
0.828	-0.018929\\
0.83	-0.018929\\
0.832	-0.018929\\
0.834	-0.018929\\
0.836	-0.018929\\
0.838	-0.018929\\
0.84	-0.018929\\
0.842	-0.018929\\
0.844	-0.018929\\
0.846	-0.018929\\
0.848	-0.018929\\
0.85	-0.018929\\
0.852	-0.018929\\
0.854	-0.018929\\
0.856	-0.018929\\
0.858	-0.018929\\
0.86	-0.018929\\
0.862	-0.018929\\
0.864	-0.018929\\
0.866	-0.018929\\
0.868	-0.018929\\
0.87	-0.018929\\
0.872	-0.018929\\
0.874	-0.018929\\
0.876	-0.018929\\
0.878	-0.018929\\
0.88	-0.018929\\
0.882	-0.018929\\
0.884	-0.018929\\
0.886	-0.018929\\
0.888	-0.018929\\
0.89	-0.018929\\
0.892	-0.018929\\
0.894	-0.018929\\
0.896	-0.018929\\
0.898	-0.018929\\
0.9	-0.018929\\
0.902	-0.018929\\
0.904	-0.018929\\
0.906	-0.018929\\
0.908	-0.018929\\
0.91	-0.018929\\
0.912	-0.018929\\
0.914	-0.018929\\
0.916	-0.018929\\
0.918	-0.018929\\
0.92	-0.018929\\
0.922	-0.018929\\
0.924	-0.0189291\\
0.926	-0.0189291\\
0.928	-0.0189291\\
0.93	-0.0189291\\
0.932	-0.0189291\\
0.934	-0.0189291\\
0.936	-0.0189291\\
0.938	-0.0189291\\
0.94	-0.0189291\\
0.942	-0.0189291\\
0.944	-0.0189291\\
0.946	-0.0189291\\
0.948	-0.0189291\\
0.95	-0.0189291\\
0.952	-0.0189291\\
0.954	-0.0189291\\
0.956	-0.0189291\\
0.958	-0.0189291\\
0.96	-0.0189291\\
0.962	-0.0189291\\
0.964	-0.0189291\\
0.966	-0.0189291\\
0.968	-0.0189291\\
0.97	-0.0189291\\
0.972	-0.0189291\\
0.974	-0.0189291\\
0.976	-0.0189291\\
0.978	-0.0189291\\
0.98	-0.0189291\\
0.982	-0.0189291\\
0.984	-0.0189291\\
0.986	-0.0189291\\
0.988	-0.0189291\\
0.99	-0.0189291\\
0.992	-0.0189291\\
0.994	-0.0189291\\
0.996	-0.0189291\\
0.998	-0.0189291\\
1	-0.0189291\\
1.002	-0.0189291\\
1.004	-0.0189291\\
1.006	-0.0189291\\
1.008	-0.0189291\\
1.01	-0.0189291\\
1.012	-0.0189291\\
1.014	-0.0189291\\
1.016	-0.0189291\\
1.018	-0.0189291\\
1.02	-0.0189291\\
1.022	-0.0189291\\
1.024	-0.0189291\\
1.026	-0.0189291\\
1.028	-0.0189291\\
1.03	-0.0189291\\
1.032	-0.0189291\\
1.034	-0.0189291\\
1.036	-0.0189291\\
1.038	-0.0189291\\
1.04	-0.0189291\\
1.042	-0.0189291\\
1.044	-0.0189291\\
1.046	-0.0189291\\
1.048	-0.0189291\\
1.05	-0.0189291\\
1.052	-0.0189291\\
1.054	-0.0189291\\
1.056	-0.0189291\\
1.058	-0.0189291\\
1.06	-0.0189291\\
1.062	-0.0189291\\
1.064	-0.0189291\\
1.066	-0.0189291\\
1.068	-0.0189291\\
1.07	-0.0189291\\
1.072	-0.0189291\\
1.074	-0.0189291\\
1.076	-0.0189291\\
1.078	-0.0189291\\
1.08	-0.0189291\\
1.082	-0.0189291\\
1.084	-0.0189291\\
1.086	-0.0189291\\
1.088	-0.0189291\\
1.09	-0.0189291\\
1.092	-0.0189291\\
1.094	-0.0189291\\
1.096	-0.0189291\\
1.098	-0.0189291\\
1.1	-0.0189291\\
1.102	-0.0189291\\
1.104	-0.0189291\\
1.106	-0.0189291\\
1.108	-0.0189291\\
1.11	-0.0189291\\
1.112	-0.0189291\\
1.114	-0.0189291\\
1.116	-0.0189291\\
1.118	-0.0189291\\
1.12	-0.0189291\\
1.122	-0.0189291\\
1.124	-0.0189291\\
1.126	-0.0189291\\
1.128	-0.0189291\\
1.13	-0.0189291\\
1.132	-0.0189291\\
1.134	-0.0189291\\
1.136	-0.0189291\\
1.138	-0.0189291\\
1.14	-0.0189291\\
1.142	-0.0189291\\
1.144	-0.0189291\\
1.146	-0.0189291\\
1.148	-0.0189291\\
1.15	-0.0189291\\
1.152	-0.0189291\\
1.154	-0.0189291\\
1.156	-0.0189291\\
1.158	-0.0189291\\
1.16	-0.0189291\\
1.162	-0.0189291\\
1.164	-0.0189291\\
1.166	-0.0189291\\
1.168	-0.0189291\\
1.17	-0.0189291\\
1.172	-0.0189291\\
1.174	-0.0189291\\
1.176	-0.0189291\\
1.178	-0.0189291\\
1.18	-0.0189291\\
1.182	-0.0189291\\
1.184	-0.0189291\\
1.186	-0.0189291\\
1.188	-0.0189291\\
1.19	-0.0189291\\
1.192	-0.0189291\\
1.194	-0.0189291\\
1.196	-0.0189291\\
1.198	-0.0189291\\
1.2	-0.0189291\\
1.202	-0.0189291\\
1.204	-0.0189291\\
1.206	-0.0189291\\
1.208	-0.0189291\\
1.21	-0.0189292\\
1.212	-0.0189292\\
1.214	-0.0189292\\
1.216	-0.0189292\\
1.218	-0.0189292\\
1.22	-0.0189292\\
1.222	-0.0189292\\
1.224	-0.0189292\\
1.226	-0.0189292\\
1.228	-0.0189292\\
1.23	-0.0189292\\
1.232	-0.0189292\\
1.234	-0.0189292\\
1.236	-0.0189292\\
1.238	-0.0189292\\
1.24	-0.0189292\\
1.242	-0.0189292\\
1.244	-0.0189292\\
1.246	-0.0189292\\
1.248	-0.0189292\\
1.25	-0.0189292\\
1.252	-0.0189292\\
1.254	-0.0189292\\
1.256	-0.0189292\\
1.258	-0.0189292\\
1.26	-0.0189292\\
1.262	-0.0189292\\
1.264	-0.0189292\\
1.266	-0.0189292\\
1.268	-0.0189292\\
1.27	-0.0189292\\
1.272	-0.0189292\\
1.274	-0.0189292\\
1.276	-0.0189292\\
1.278	-0.0189292\\
1.28	-0.0189292\\
1.282	-0.0189292\\
1.284	-0.0189292\\
1.286	-0.0189292\\
1.288	-0.0189292\\
1.29	-0.0189292\\
1.292	-0.0189292\\
1.294	-0.0189292\\
1.296	-0.0189292\\
1.298	-0.0189292\\
1.3	-0.0189292\\
1.302	-0.0189292\\
1.304	-0.0189292\\
1.306	-0.0189292\\
1.308	-0.0189292\\
1.31	-0.0189292\\
1.312	-0.0189292\\
1.314	-0.0189292\\
1.316	-0.0189292\\
1.318	-0.0189292\\
1.32	-0.0189292\\
1.322	-0.0189292\\
1.324	-0.0189292\\
1.326	-0.0189292\\
1.328	-0.0189292\\
1.33	-0.0189292\\
1.332	-0.0189292\\
1.334	-0.0189292\\
1.336	-0.0189292\\
1.338	-0.0189292\\
1.34	-0.0189292\\
1.342	-0.0189292\\
1.344	-0.0189292\\
1.346	-0.0189292\\
1.348	-0.0189292\\
1.35	-0.0189292\\
1.352	-0.0189292\\
1.354	-0.0189292\\
1.356	-0.0189292\\
1.358	-0.0189292\\
1.36	-0.0189292\\
1.362	-0.0189292\\
1.364	-0.0189292\\
1.366	-0.0189292\\
1.368	-0.0189292\\
1.37	-0.0189292\\
1.372	-0.0189292\\
1.374	-0.0189292\\
1.376	-0.0189292\\
1.378	-0.0189292\\
1.38	-0.0189292\\
1.382	-0.0189292\\
1.384	-0.0189292\\
1.386	-0.0189292\\
1.388	-0.0189292\\
1.39	-0.0189292\\
1.392	-0.0189292\\
1.394	-0.0189292\\
1.396	-0.0189292\\
1.398	-0.0189292\\
1.4	-0.0189292\\
1.402	-0.0189292\\
1.404	-0.0189292\\
1.406	-0.0189292\\
1.408	-0.0189292\\
1.41	-0.0189292\\
1.412	-0.0189292\\
1.414	-0.0189292\\
1.416	-0.0189292\\
1.418	-0.0189292\\
1.42	-0.0189292\\
1.422	-0.0189292\\
1.424	-0.0189292\\
1.426	-0.0189292\\
1.428	-0.0189292\\
1.43	-0.0189292\\
1.432	-0.0189292\\
1.434	-0.0189292\\
1.436	-0.0189292\\
1.438	-0.0189292\\
1.44	-0.0189292\\
1.442	-0.0189292\\
1.444	-0.0189292\\
1.446	-0.0189292\\
1.448	-0.0189292\\
1.45	-0.0189292\\
1.452	-0.0189292\\
1.454	-0.0189292\\
1.456	-0.0189292\\
1.458	-0.0189292\\
1.46	-0.0189292\\
1.462	-0.0189292\\
1.464	-0.0189292\\
1.466	-0.0189292\\
1.468	-0.0189292\\
1.47	-0.0189292\\
1.472	-0.0189292\\
1.474	-0.0189292\\
1.476	-0.0189292\\
1.478	-0.0189292\\
1.48	-0.0189292\\
1.482	-0.0189292\\
1.484	-0.0189292\\
1.486	-0.0189292\\
1.488	-0.0189292\\
1.49	-0.0189292\\
1.492	-0.0189292\\
1.494	-0.0189292\\
1.496	-0.0189292\\
1.498	-0.0189292\\
1.5	-0.0189292\\
1.502	-0.0189292\\
1.504	-0.0189292\\
1.506	-0.0189292\\
1.508	-0.0189292\\
1.51	-0.0189292\\
1.512	-0.0189292\\
1.514	-0.0189292\\
1.516	-0.0189292\\
1.518	-0.0189292\\
1.52	-0.0189292\\
1.522	-0.0189292\\
1.524	-0.0189292\\
1.526	-0.0189292\\
1.528	-0.0189292\\
1.53	-0.0189292\\
1.532	-0.0189292\\
1.534	-0.0189292\\
1.536	-0.0189292\\
1.538	-0.0189292\\
1.54	-0.0189292\\
1.542	-0.0189292\\
1.544	-0.0189292\\
1.546	-0.0189292\\
1.548	-0.0189292\\
1.55	-0.0189292\\
1.552	-0.0189292\\
1.554	-0.0189292\\
1.556	-0.0189292\\
1.558	-0.0189292\\
1.56	-0.0189292\\
1.562	-0.0189292\\
1.564	-0.0189292\\
1.566	-0.0189292\\
1.568	-0.0189292\\
1.57	-0.0189292\\
1.572	-0.0189292\\
1.574	-0.0189292\\
1.576	-0.0189292\\
1.578	-0.0189292\\
1.58	-0.0189292\\
1.582	-0.0189292\\
1.584	-0.0189292\\
1.586	-0.0189292\\
1.588	-0.0189292\\
1.59	-0.0189292\\
1.592	-0.0189292\\
1.594	-0.0189292\\
1.596	-0.0189292\\
1.598	-0.0189292\\
1.6	-0.0189292\\
1.602	-0.0189292\\
1.604	-0.0189292\\
1.606	-0.0189292\\
1.608	-0.0189292\\
1.61	-0.0189292\\
1.612	-0.0189292\\
1.614	-0.0189292\\
1.616	-0.0189292\\
1.618	-0.0189292\\
1.62	-0.0189292\\
1.622	-0.0189292\\
1.624	-0.0189292\\
1.626	-0.0189292\\
1.628	-0.0189292\\
1.63	-0.0189292\\
1.632	-0.0189292\\
1.634	-0.0189292\\
1.636	-0.0189292\\
1.638	-0.0189292\\
1.64	-0.0189292\\
1.642	-0.0189292\\
1.644	-0.0189292\\
1.646	-0.0189292\\
1.648	-0.0189292\\
1.65	-0.0189292\\
1.652	-0.0189292\\
1.654	-0.0189292\\
1.656	-0.0189292\\
1.658	-0.0189292\\
1.66	-0.0189292\\
1.662	-0.0189292\\
1.664	-0.0189292\\
1.666	-0.0189292\\
1.668	-0.0189292\\
1.67	-0.0189292\\
1.672	-0.0189292\\
1.674	-0.0189292\\
1.676	-0.0189292\\
1.678	-0.0189292\\
1.68	-0.0189292\\
1.682	-0.0189292\\
1.684	-0.0189292\\
1.686	-0.0189292\\
1.688	-0.0189292\\
1.69	-0.0189292\\
1.692	-0.0189292\\
1.694	-0.0189292\\
1.696	-0.0189292\\
1.698	-0.0189292\\
1.7	-0.0189292\\
1.702	-0.0189292\\
1.704	-0.0189292\\
1.706	-0.0189292\\
1.708	-0.0189292\\
1.71	-0.0189292\\
1.712	-0.0189292\\
1.714	-0.0189292\\
1.716	-0.0189292\\
1.718	-0.0189292\\
1.72	-0.0189292\\
1.722	-0.0189292\\
1.724	-0.0189292\\
1.726	-0.0189292\\
1.728	-0.0189292\\
1.73	-0.0189292\\
1.732	-0.0189292\\
1.734	-0.0189292\\
1.736	-0.0189292\\
1.738	-0.0189292\\
1.74	-0.0189292\\
1.742	-0.0189292\\
1.744	-0.0189292\\
1.746	-0.0189292\\
1.748	-0.0189292\\
1.75	-0.0189292\\
1.752	-0.0189292\\
1.754	-0.0189292\\
1.756	-0.0189292\\
1.758	-0.0189292\\
1.76	-0.0189292\\
1.762	-0.0189292\\
1.764	-0.0189292\\
1.766	-0.0189292\\
1.768	-0.0189292\\
1.77	-0.0189292\\
1.772	-0.0189292\\
1.774	-0.0189292\\
1.776	-0.0189292\\
1.778	-0.0189292\\
1.78	-0.0189292\\
1.782	-0.0189292\\
1.784	-0.0189292\\
1.786	-0.0189292\\
1.788	-0.0189292\\
1.79	-0.0189292\\
1.792	-0.0189292\\
1.794	-0.0189292\\
1.796	-0.0189292\\
1.798	-0.0189292\\
1.8	-0.0189292\\
1.802	-0.0189292\\
1.804	-0.0189292\\
1.806	-0.0189292\\
1.808	-0.0189292\\
1.81	-0.0189292\\
1.812	-0.0189292\\
1.814	-0.0189292\\
1.816	-0.0189292\\
1.818	-0.0189292\\
1.82	-0.0189292\\
1.822	-0.0189292\\
1.824	-0.0189292\\
1.826	-0.0189292\\
1.828	-0.0189292\\
1.83	-0.0189292\\
1.832	-0.0189292\\
1.834	-0.0189292\\
1.836	-0.0189292\\
1.838	-0.0189292\\
1.84	-0.0189292\\
1.842	-0.0189292\\
1.844	-0.0189292\\
1.846	-0.0189292\\
1.848	-0.0189292\\
1.85	-0.0189292\\
1.852	-0.0189292\\
1.854	-0.0189292\\
1.856	-0.0189292\\
1.858	-0.0189292\\
1.86	-0.0189292\\
1.862	-0.0189292\\
1.864	-0.0189292\\
1.866	-0.0189292\\
1.868	-0.0189292\\
1.87	-0.0189292\\
1.872	-0.0189292\\
1.874	-0.0189292\\
1.876	-0.0189292\\
1.878	-0.0189292\\
1.88	-0.0189292\\
1.882	-0.0189292\\
1.884	-0.0189292\\
1.886	-0.0189292\\
1.888	-0.0189292\\
1.89	-0.0189292\\
1.892	-0.0189292\\
1.894	-0.0189292\\
1.896	-0.0189292\\
1.898	-0.0189292\\
1.9	-0.0189292\\
1.902	-0.0189292\\
1.904	-0.0189292\\
1.906	-0.0189292\\
1.908	-0.0189292\\
1.91	-0.0189292\\
1.912	-0.0189292\\
1.914	-0.0189292\\
1.916	-0.0189292\\
1.918	-0.0189292\\
1.92	-0.0189292\\
1.922	-0.0189292\\
1.924	-0.0189292\\
1.926	-0.0189292\\
1.928	-0.0189292\\
1.93	-0.0189292\\
1.932	-0.0189292\\
1.934	-0.0189292\\
1.936	-0.0189292\\
1.938	-0.0189292\\
1.94	-0.0189292\\
1.942	-0.0189292\\
1.944	-0.0189292\\
1.946	-0.0189292\\
1.948	-0.0189292\\
1.95	-0.0189292\\
1.952	-0.0189292\\
1.954	-0.0189292\\
1.956	-0.0189292\\
1.958	-0.0189292\\
1.96	-0.0189292\\
1.962	-0.0189292\\
1.964	-0.0189292\\
1.966	-0.0189292\\
1.968	-0.0189292\\
1.97	-0.0189292\\
1.972	-0.0189292\\
1.974	-0.0189292\\
1.976	-0.0189292\\
1.978	-0.0189292\\
1.98	-0.0189292\\
1.982	-0.0189292\\
1.984	-0.0189292\\
1.986	-0.0189292\\
1.988	-0.0189292\\
1.99	-0.0189292\\
1.992	-0.0189292\\
1.994	-0.0189292\\
1.996	-0.0189292\\
1.998	-0.0189292\\
2	-0.0189292\\
};
\addplot [color=mycolor9, line width=1.0pt, forget plot]
  table[row sep=crcr]{%
-0.024	0\\
-0.022	0\\
-0.02	0\\
-0.018	0\\
-0.016	0\\
-0.014	0\\
-0.012	0\\
-0.01	0\\
-0.008	0\\
-0.006	0\\
-0.004	0\\
-0.002	0\\
0	0\\
0.002	-0.00183154\\
0.004	-0.00337927\\
0.006	-0.00470796\\
0.008	-0.00586631\\
0.01	-0.00689082\\
0.012	-0.00780871\\
0.014	-0.0086403\\
0.016	-0.00940078\\
0.018	-0.0101015\\
0.02	-0.010751\\
0.022	-0.0113558\\
0.024	-0.0119209\\
0.026	-0.0124502\\
0.028	-0.0129468\\
0.03	-0.0134133\\
0.032	-0.0138515\\
0.034	-0.0142634\\
0.036	-0.0146505\\
0.038	-0.015014\\
0.04	-0.0153552\\
0.042	-0.0156752\\
0.044	-0.015975\\
0.046	-0.0162557\\
0.048	-0.0165182\\
0.05	-0.0167634\\
0.052	-0.0169923\\
0.054	-0.0172058\\
0.056	-0.0174046\\
0.058	-0.0175898\\
0.06	-0.0177621\\
0.062	-0.0179224\\
0.064	-0.0180713\\
0.066	-0.0182098\\
0.068	-0.0183384\\
0.07	-0.0184578\\
0.072	-0.0185687\\
0.074	-0.0186716\\
0.076	-0.0187671\\
0.078	-0.0188556\\
0.08	-0.0189376\\
0.082	-0.0190135\\
0.084	-0.0190837\\
0.086	-0.0191484\\
0.088	-0.019208\\
0.09	-0.0192628\\
0.092	-0.0193129\\
0.094	-0.0193586\\
0.096	-0.0194001\\
0.098	-0.0194376\\
0.1	-0.0194712\\
0.102	-0.0195011\\
0.104	-0.0195275\\
0.106	-0.0195505\\
0.108	-0.0195702\\
0.11	-0.0195868\\
0.112	-0.0196005\\
0.114	-0.0196113\\
0.116	-0.0196193\\
0.118	-0.0196248\\
0.12	-0.0196279\\
0.122	-0.0196286\\
0.124	-0.0196272\\
0.126	-0.0196238\\
0.128	-0.0196184\\
0.13	-0.0196112\\
0.132	-0.0196024\\
0.134	-0.0195921\\
0.136	-0.0195803\\
0.138	-0.0195673\\
0.14	-0.0195531\\
0.142	-0.0195378\\
0.144	-0.0195216\\
0.146	-0.0195045\\
0.148	-0.0194868\\
0.15	-0.0194683\\
0.152	-0.0194493\\
0.154	-0.0194299\\
0.156	-0.0194101\\
0.158	-0.01939\\
0.16	-0.0193696\\
0.162	-0.0193491\\
0.164	-0.0193286\\
0.166	-0.019308\\
0.168	-0.0192874\\
0.17	-0.019267\\
0.172	-0.0192466\\
0.174	-0.0192265\\
0.176	-0.0192066\\
0.178	-0.019187\\
0.18	-0.0191677\\
0.182	-0.0191488\\
0.184	-0.0191302\\
0.186	-0.0191121\\
0.188	-0.0190943\\
0.19	-0.0190771\\
0.192	-0.0190603\\
0.194	-0.019044\\
0.196	-0.0190283\\
0.198	-0.019013\\
0.2	-0.0189983\\
0.202	-0.0189842\\
0.204	-0.0189706\\
0.206	-0.0189575\\
0.208	-0.0189451\\
0.21	-0.0189331\\
0.212	-0.0189218\\
0.214	-0.018911\\
0.216	-0.0189008\\
0.218	-0.0188912\\
0.22	-0.0188821\\
0.222	-0.0188735\\
0.224	-0.0188655\\
0.226	-0.018858\\
0.228	-0.0188511\\
0.23	-0.0188446\\
0.232	-0.0188387\\
0.234	-0.0188332\\
0.236	-0.0188282\\
0.238	-0.0188237\\
0.24	-0.0188196\\
0.242	-0.0188159\\
0.244	-0.0188127\\
0.246	-0.0188099\\
0.248	-0.0188074\\
0.25	-0.0188053\\
0.252	-0.0188036\\
0.254	-0.0188022\\
0.256	-0.0188012\\
0.258	-0.0188004\\
0.26	-0.0188\\
0.262	-0.0187998\\
0.264	-0.0187999\\
0.266	-0.0188003\\
0.268	-0.0188008\\
0.27	-0.0188016\\
0.272	-0.0188026\\
0.274	-0.0188038\\
0.276	-0.0188052\\
0.278	-0.0188067\\
0.28	-0.0188084\\
0.282	-0.0188102\\
0.284	-0.0188122\\
0.286	-0.0188142\\
0.288	-0.0188164\\
0.29	-0.0188187\\
0.292	-0.018821\\
0.294	-0.0188234\\
0.296	-0.0188259\\
0.298	-0.0188284\\
0.3	-0.018831\\
0.302	-0.0188336\\
0.304	-0.0188362\\
0.306	-0.0188388\\
0.308	-0.0188415\\
0.31	-0.0188441\\
0.312	-0.0188468\\
0.314	-0.0188494\\
0.316	-0.018852\\
0.318	-0.0188546\\
0.32	-0.0188572\\
0.322	-0.0188598\\
0.324	-0.0188623\\
0.326	-0.0188648\\
0.328	-0.0188672\\
0.33	-0.0188696\\
0.332	-0.018872\\
0.334	-0.0188743\\
0.336	-0.0188765\\
0.338	-0.0188787\\
0.34	-0.0188808\\
0.342	-0.0188829\\
0.344	-0.0188849\\
0.346	-0.0188869\\
0.348	-0.0188888\\
0.35	-0.0188906\\
0.352	-0.0188924\\
0.354	-0.0188941\\
0.356	-0.0188958\\
0.358	-0.0188974\\
0.36	-0.0188989\\
0.362	-0.0189004\\
0.364	-0.0189018\\
0.366	-0.0189032\\
0.368	-0.0189045\\
0.37	-0.0189058\\
0.372	-0.018907\\
0.374	-0.0189081\\
0.376	-0.0189092\\
0.378	-0.0189103\\
0.38	-0.0189113\\
0.382	-0.0189122\\
0.384	-0.0189131\\
0.386	-0.018914\\
0.388	-0.0189148\\
0.39	-0.0189155\\
0.392	-0.0189163\\
0.394	-0.018917\\
0.396	-0.0189176\\
0.398	-0.0189182\\
0.4	-0.0189188\\
0.402	-0.0189193\\
0.404	-0.0189199\\
0.406	-0.0189203\\
0.408	-0.0189208\\
0.41	-0.0189212\\
0.412	-0.0189216\\
0.414	-0.018922\\
0.416	-0.0189223\\
0.418	-0.0189227\\
0.42	-0.018923\\
0.422	-0.0189233\\
0.424	-0.0189235\\
0.426	-0.0189238\\
0.428	-0.018924\\
0.43	-0.0189242\\
0.432	-0.0189244\\
0.434	-0.0189246\\
0.436	-0.0189247\\
0.438	-0.0189249\\
0.44	-0.018925\\
0.442	-0.0189251\\
0.444	-0.0189253\\
0.446	-0.0189254\\
0.448	-0.0189255\\
0.45	-0.0189256\\
0.452	-0.0189256\\
0.454	-0.0189257\\
0.456	-0.0189258\\
0.458	-0.0189258\\
0.46	-0.0189259\\
0.462	-0.0189259\\
0.464	-0.0189259\\
0.466	-0.018926\\
0.468	-0.018926\\
0.47	-0.018926\\
0.472	-0.018926\\
0.474	-0.0189261\\
0.476	-0.0189261\\
0.478	-0.0189261\\
0.48	-0.0189261\\
0.482	-0.0189261\\
0.484	-0.0189261\\
0.486	-0.0189261\\
0.488	-0.0189261\\
0.49	-0.0189261\\
0.492	-0.0189261\\
0.494	-0.0189261\\
0.496	-0.0189261\\
0.498	-0.0189261\\
0.5	-0.0189261\\
0.502	-0.0189261\\
0.504	-0.0189261\\
0.506	-0.018926\\
0.508	-0.018926\\
0.51	-0.018926\\
0.512	-0.018926\\
0.514	-0.018926\\
0.516	-0.018926\\
0.518	-0.018926\\
0.52	-0.018926\\
0.522	-0.018926\\
0.524	-0.018926\\
0.526	-0.018926\\
0.528	-0.018926\\
0.53	-0.018926\\
0.532	-0.018926\\
0.534	-0.018926\\
0.536	-0.018926\\
0.538	-0.018926\\
0.54	-0.018926\\
0.542	-0.018926\\
0.544	-0.018926\\
0.546	-0.018926\\
0.548	-0.018926\\
0.55	-0.018926\\
0.552	-0.018926\\
0.554	-0.018926\\
0.556	-0.018926\\
0.558	-0.018926\\
0.56	-0.018926\\
0.562	-0.0189261\\
0.564	-0.0189261\\
0.566	-0.0189261\\
0.568	-0.0189261\\
0.57	-0.0189261\\
0.572	-0.0189261\\
0.574	-0.0189262\\
0.576	-0.0189262\\
0.578	-0.0189262\\
0.58	-0.0189262\\
0.582	-0.0189262\\
0.584	-0.0189263\\
0.586	-0.0189263\\
0.588	-0.0189263\\
0.59	-0.0189263\\
0.592	-0.0189264\\
0.594	-0.0189264\\
0.596	-0.0189264\\
0.598	-0.0189264\\
0.6	-0.0189265\\
0.602	-0.0189265\\
0.604	-0.0189265\\
0.606	-0.0189266\\
0.608	-0.0189266\\
0.61	-0.0189266\\
0.612	-0.0189267\\
0.614	-0.0189267\\
0.616	-0.0189267\\
0.618	-0.0189268\\
0.62	-0.0189268\\
0.622	-0.0189268\\
0.624	-0.0189269\\
0.626	-0.0189269\\
0.628	-0.0189269\\
0.63	-0.018927\\
0.632	-0.018927\\
0.634	-0.018927\\
0.636	-0.0189271\\
0.638	-0.0189271\\
0.64	-0.0189271\\
0.642	-0.0189272\\
0.644	-0.0189272\\
0.646	-0.0189272\\
0.648	-0.0189273\\
0.65	-0.0189273\\
0.652	-0.0189274\\
0.654	-0.0189274\\
0.656	-0.0189274\\
0.658	-0.0189275\\
0.66	-0.0189275\\
0.662	-0.0189275\\
0.664	-0.0189276\\
0.666	-0.0189276\\
0.668	-0.0189276\\
0.67	-0.0189277\\
0.672	-0.0189277\\
0.674	-0.0189277\\
0.676	-0.0189278\\
0.678	-0.0189278\\
0.68	-0.0189278\\
0.682	-0.0189279\\
0.684	-0.0189279\\
0.686	-0.0189279\\
0.688	-0.018928\\
0.69	-0.018928\\
0.692	-0.018928\\
0.694	-0.0189281\\
0.696	-0.0189281\\
0.698	-0.0189281\\
0.7	-0.0189281\\
0.702	-0.0189282\\
0.704	-0.0189282\\
0.706	-0.0189282\\
0.708	-0.0189282\\
0.71	-0.0189283\\
0.712	-0.0189283\\
0.714	-0.0189283\\
0.716	-0.0189283\\
0.718	-0.0189284\\
0.72	-0.0189284\\
0.722	-0.0189284\\
0.724	-0.0189284\\
0.726	-0.0189285\\
0.728	-0.0189285\\
0.73	-0.0189285\\
0.732	-0.0189285\\
0.734	-0.0189285\\
0.736	-0.0189286\\
0.738	-0.0189286\\
0.74	-0.0189286\\
0.742	-0.0189286\\
0.744	-0.0189286\\
0.746	-0.0189286\\
0.748	-0.0189287\\
0.75	-0.0189287\\
0.752	-0.0189287\\
0.754	-0.0189287\\
0.756	-0.0189287\\
0.758	-0.0189287\\
0.76	-0.0189287\\
0.762	-0.0189287\\
0.764	-0.0189288\\
0.766	-0.0189288\\
0.768	-0.0189288\\
0.77	-0.0189288\\
0.772	-0.0189288\\
0.774	-0.0189288\\
0.776	-0.0189288\\
0.778	-0.0189288\\
0.78	-0.0189288\\
0.782	-0.0189288\\
0.784	-0.0189288\\
0.786	-0.0189289\\
0.788	-0.0189289\\
0.79	-0.0189289\\
0.792	-0.0189289\\
0.794	-0.0189289\\
0.796	-0.0189289\\
0.798	-0.0189289\\
0.8	-0.0189289\\
0.802	-0.0189289\\
0.804	-0.0189289\\
0.806	-0.0189289\\
0.808	-0.0189289\\
0.81	-0.0189289\\
0.812	-0.0189289\\
0.814	-0.0189289\\
0.816	-0.0189289\\
0.818	-0.0189289\\
0.82	-0.0189289\\
0.822	-0.0189289\\
0.824	-0.0189289\\
0.826	-0.0189289\\
0.828	-0.0189289\\
0.83	-0.0189289\\
0.832	-0.018929\\
0.834	-0.018929\\
0.836	-0.018929\\
0.838	-0.018929\\
0.84	-0.018929\\
0.842	-0.018929\\
0.844	-0.018929\\
0.846	-0.018929\\
0.848	-0.018929\\
0.85	-0.018929\\
0.852	-0.018929\\
0.854	-0.018929\\
0.856	-0.018929\\
0.858	-0.018929\\
0.86	-0.018929\\
0.862	-0.018929\\
0.864	-0.018929\\
0.866	-0.018929\\
0.868	-0.018929\\
0.87	-0.018929\\
0.872	-0.018929\\
0.874	-0.018929\\
0.876	-0.018929\\
0.878	-0.018929\\
0.88	-0.018929\\
0.882	-0.018929\\
0.884	-0.018929\\
0.886	-0.018929\\
0.888	-0.018929\\
0.89	-0.018929\\
0.892	-0.018929\\
0.894	-0.018929\\
0.896	-0.018929\\
0.898	-0.018929\\
0.9	-0.018929\\
0.902	-0.018929\\
0.904	-0.018929\\
0.906	-0.018929\\
0.908	-0.018929\\
0.91	-0.018929\\
0.912	-0.018929\\
0.914	-0.018929\\
0.916	-0.018929\\
0.918	-0.018929\\
0.92	-0.018929\\
0.922	-0.018929\\
0.924	-0.018929\\
0.926	-0.018929\\
0.928	-0.018929\\
0.93	-0.018929\\
0.932	-0.018929\\
0.934	-0.018929\\
0.936	-0.0189291\\
0.938	-0.0189291\\
0.94	-0.0189291\\
0.942	-0.0189291\\
0.944	-0.0189291\\
0.946	-0.0189291\\
0.948	-0.0189291\\
0.95	-0.0189291\\
0.952	-0.0189291\\
0.954	-0.0189291\\
0.956	-0.0189291\\
0.958	-0.0189291\\
0.96	-0.0189291\\
0.962	-0.0189291\\
0.964	-0.0189291\\
0.966	-0.0189291\\
0.968	-0.0189291\\
0.97	-0.0189291\\
0.972	-0.0189291\\
0.974	-0.0189291\\
0.976	-0.0189291\\
0.978	-0.0189291\\
0.98	-0.0189291\\
0.982	-0.0189291\\
0.984	-0.0189291\\
0.986	-0.0189291\\
0.988	-0.0189291\\
0.99	-0.0189291\\
0.992	-0.0189291\\
0.994	-0.0189291\\
0.996	-0.0189291\\
0.998	-0.0189291\\
1	-0.0189291\\
1.002	-0.0189291\\
1.004	-0.0189291\\
1.006	-0.0189291\\
1.008	-0.0189291\\
1.01	-0.0189291\\
1.012	-0.0189291\\
1.014	-0.0189291\\
1.016	-0.0189291\\
1.018	-0.0189291\\
1.02	-0.0189291\\
1.022	-0.0189291\\
1.024	-0.0189291\\
1.026	-0.0189291\\
1.028	-0.0189291\\
1.03	-0.0189291\\
1.032	-0.0189291\\
1.034	-0.0189291\\
1.036	-0.0189291\\
1.038	-0.0189291\\
1.04	-0.0189291\\
1.042	-0.0189291\\
1.044	-0.0189291\\
1.046	-0.0189291\\
1.048	-0.0189291\\
1.05	-0.0189291\\
1.052	-0.0189291\\
1.054	-0.0189291\\
1.056	-0.0189291\\
1.058	-0.0189291\\
1.06	-0.0189291\\
1.062	-0.0189291\\
1.064	-0.0189291\\
1.066	-0.0189291\\
1.068	-0.0189291\\
1.07	-0.0189291\\
1.072	-0.0189291\\
1.074	-0.0189291\\
1.076	-0.0189291\\
1.078	-0.0189291\\
1.08	-0.0189291\\
1.082	-0.0189291\\
1.084	-0.0189291\\
1.086	-0.0189291\\
1.088	-0.0189291\\
1.09	-0.0189291\\
1.092	-0.0189291\\
1.094	-0.0189291\\
1.096	-0.0189291\\
1.098	-0.0189291\\
1.1	-0.0189291\\
1.102	-0.0189291\\
1.104	-0.0189291\\
1.106	-0.0189291\\
1.108	-0.0189291\\
1.11	-0.0189291\\
1.112	-0.0189291\\
1.114	-0.0189291\\
1.116	-0.0189291\\
1.118	-0.0189291\\
1.12	-0.0189291\\
1.122	-0.0189291\\
1.124	-0.0189291\\
1.126	-0.0189291\\
1.128	-0.0189291\\
1.13	-0.0189291\\
1.132	-0.0189291\\
1.134	-0.0189291\\
1.136	-0.0189291\\
1.138	-0.0189291\\
1.14	-0.0189291\\
1.142	-0.0189291\\
1.144	-0.0189291\\
1.146	-0.0189291\\
1.148	-0.0189291\\
1.15	-0.0189291\\
1.152	-0.0189291\\
1.154	-0.0189291\\
1.156	-0.0189291\\
1.158	-0.0189291\\
1.16	-0.0189291\\
1.162	-0.0189291\\
1.164	-0.0189291\\
1.166	-0.0189291\\
1.168	-0.0189291\\
1.17	-0.0189291\\
1.172	-0.0189291\\
1.174	-0.0189291\\
1.176	-0.0189291\\
1.178	-0.0189291\\
1.18	-0.0189291\\
1.182	-0.0189291\\
1.184	-0.0189291\\
1.186	-0.0189291\\
1.188	-0.0189291\\
1.19	-0.0189291\\
1.192	-0.0189291\\
1.194	-0.0189291\\
1.196	-0.0189291\\
1.198	-0.0189291\\
1.2	-0.0189291\\
1.202	-0.0189291\\
1.204	-0.0189291\\
1.206	-0.0189291\\
1.208	-0.0189291\\
1.21	-0.0189291\\
1.212	-0.0189291\\
1.214	-0.0189291\\
1.216	-0.0189291\\
1.218	-0.0189291\\
1.22	-0.0189291\\
1.222	-0.0189291\\
1.224	-0.0189291\\
1.226	-0.0189291\\
1.228	-0.0189292\\
1.23	-0.0189292\\
1.232	-0.0189292\\
1.234	-0.0189292\\
1.236	-0.0189292\\
1.238	-0.0189292\\
1.24	-0.0189292\\
1.242	-0.0189292\\
1.244	-0.0189292\\
1.246	-0.0189292\\
1.248	-0.0189292\\
1.25	-0.0189292\\
1.252	-0.0189292\\
1.254	-0.0189292\\
1.256	-0.0189292\\
1.258	-0.0189292\\
1.26	-0.0189292\\
1.262	-0.0189292\\
1.264	-0.0189292\\
1.266	-0.0189292\\
1.268	-0.0189292\\
1.27	-0.0189292\\
1.272	-0.0189292\\
1.274	-0.0189292\\
1.276	-0.0189292\\
1.278	-0.0189292\\
1.28	-0.0189292\\
1.282	-0.0189292\\
1.284	-0.0189292\\
1.286	-0.0189292\\
1.288	-0.0189292\\
1.29	-0.0189292\\
1.292	-0.0189292\\
1.294	-0.0189292\\
1.296	-0.0189292\\
1.298	-0.0189292\\
1.3	-0.0189292\\
1.302	-0.0189292\\
1.304	-0.0189292\\
1.306	-0.0189292\\
1.308	-0.0189292\\
1.31	-0.0189292\\
1.312	-0.0189292\\
1.314	-0.0189292\\
1.316	-0.0189292\\
1.318	-0.0189292\\
1.32	-0.0189292\\
1.322	-0.0189292\\
1.324	-0.0189292\\
1.326	-0.0189292\\
1.328	-0.0189292\\
1.33	-0.0189292\\
1.332	-0.0189292\\
1.334	-0.0189292\\
1.336	-0.0189292\\
1.338	-0.0189292\\
1.34	-0.0189292\\
1.342	-0.0189292\\
1.344	-0.0189292\\
1.346	-0.0189292\\
1.348	-0.0189292\\
1.35	-0.0189292\\
1.352	-0.0189292\\
1.354	-0.0189292\\
1.356	-0.0189292\\
1.358	-0.0189292\\
1.36	-0.0189292\\
1.362	-0.0189292\\
1.364	-0.0189292\\
1.366	-0.0189292\\
1.368	-0.0189292\\
1.37	-0.0189292\\
1.372	-0.0189292\\
1.374	-0.0189292\\
1.376	-0.0189292\\
1.378	-0.0189292\\
1.38	-0.0189292\\
1.382	-0.0189292\\
1.384	-0.0189292\\
1.386	-0.0189292\\
1.388	-0.0189292\\
1.39	-0.0189292\\
1.392	-0.0189292\\
1.394	-0.0189292\\
1.396	-0.0189292\\
1.398	-0.0189292\\
1.4	-0.0189292\\
1.402	-0.0189292\\
1.404	-0.0189292\\
1.406	-0.0189292\\
1.408	-0.0189292\\
1.41	-0.0189292\\
1.412	-0.0189292\\
1.414	-0.0189292\\
1.416	-0.0189292\\
1.418	-0.0189292\\
1.42	-0.0189292\\
1.422	-0.0189292\\
1.424	-0.0189292\\
1.426	-0.0189292\\
1.428	-0.0189292\\
1.43	-0.0189292\\
1.432	-0.0189292\\
1.434	-0.0189292\\
1.436	-0.0189292\\
1.438	-0.0189292\\
1.44	-0.0189292\\
1.442	-0.0189292\\
1.444	-0.0189292\\
1.446	-0.0189292\\
1.448	-0.0189292\\
1.45	-0.0189292\\
1.452	-0.0189292\\
1.454	-0.0189292\\
1.456	-0.0189292\\
1.458	-0.0189292\\
1.46	-0.0189292\\
1.462	-0.0189292\\
1.464	-0.0189292\\
1.466	-0.0189292\\
1.468	-0.0189292\\
1.47	-0.0189292\\
1.472	-0.0189292\\
1.474	-0.0189292\\
1.476	-0.0189292\\
1.478	-0.0189292\\
1.48	-0.0189292\\
1.482	-0.0189292\\
1.484	-0.0189292\\
1.486	-0.0189292\\
1.488	-0.0189292\\
1.49	-0.0189292\\
1.492	-0.0189292\\
1.494	-0.0189292\\
1.496	-0.0189292\\
1.498	-0.0189292\\
1.5	-0.0189292\\
1.502	-0.0189292\\
1.504	-0.0189292\\
1.506	-0.0189292\\
1.508	-0.0189292\\
1.51	-0.0189292\\
1.512	-0.0189292\\
1.514	-0.0189292\\
1.516	-0.0189292\\
1.518	-0.0189292\\
1.52	-0.0189292\\
1.522	-0.0189292\\
1.524	-0.0189292\\
1.526	-0.0189292\\
1.528	-0.0189292\\
1.53	-0.0189292\\
1.532	-0.0189292\\
1.534	-0.0189292\\
1.536	-0.0189292\\
1.538	-0.0189292\\
1.54	-0.0189292\\
1.542	-0.0189292\\
1.544	-0.0189292\\
1.546	-0.0189292\\
1.548	-0.0189292\\
1.55	-0.0189292\\
1.552	-0.0189292\\
1.554	-0.0189292\\
1.556	-0.0189292\\
1.558	-0.0189292\\
1.56	-0.0189292\\
1.562	-0.0189292\\
1.564	-0.0189292\\
1.566	-0.0189292\\
1.568	-0.0189292\\
1.57	-0.0189292\\
1.572	-0.0189292\\
1.574	-0.0189292\\
1.576	-0.0189292\\
1.578	-0.0189292\\
1.58	-0.0189292\\
1.582	-0.0189292\\
1.584	-0.0189292\\
1.586	-0.0189292\\
1.588	-0.0189292\\
1.59	-0.0189292\\
1.592	-0.0189292\\
1.594	-0.0189292\\
1.596	-0.0189292\\
1.598	-0.0189292\\
1.6	-0.0189292\\
1.602	-0.0189292\\
1.604	-0.0189292\\
1.606	-0.0189292\\
1.608	-0.0189292\\
1.61	-0.0189292\\
1.612	-0.0189292\\
1.614	-0.0189292\\
1.616	-0.0189292\\
1.618	-0.0189292\\
1.62	-0.0189292\\
1.622	-0.0189292\\
1.624	-0.0189292\\
1.626	-0.0189292\\
1.628	-0.0189292\\
1.63	-0.0189292\\
1.632	-0.0189292\\
1.634	-0.0189292\\
1.636	-0.0189292\\
1.638	-0.0189292\\
1.64	-0.0189292\\
1.642	-0.0189292\\
1.644	-0.0189292\\
1.646	-0.0189292\\
1.648	-0.0189292\\
1.65	-0.0189292\\
1.652	-0.0189292\\
1.654	-0.0189292\\
1.656	-0.0189292\\
1.658	-0.0189292\\
1.66	-0.0189292\\
1.662	-0.0189292\\
1.664	-0.0189292\\
1.666	-0.0189292\\
1.668	-0.0189292\\
1.67	-0.0189292\\
1.672	-0.0189292\\
1.674	-0.0189292\\
1.676	-0.0189292\\
1.678	-0.0189292\\
1.68	-0.0189292\\
1.682	-0.0189292\\
1.684	-0.0189292\\
1.686	-0.0189292\\
1.688	-0.0189292\\
1.69	-0.0189292\\
1.692	-0.0189292\\
1.694	-0.0189292\\
1.696	-0.0189292\\
1.698	-0.0189292\\
1.7	-0.0189292\\
1.702	-0.0189292\\
1.704	-0.0189292\\
1.706	-0.0189292\\
1.708	-0.0189292\\
1.71	-0.0189292\\
1.712	-0.0189292\\
1.714	-0.0189292\\
1.716	-0.0189292\\
1.718	-0.0189292\\
1.72	-0.0189292\\
1.722	-0.0189292\\
1.724	-0.0189292\\
1.726	-0.0189292\\
1.728	-0.0189292\\
1.73	-0.0189292\\
1.732	-0.0189292\\
1.734	-0.0189292\\
1.736	-0.0189292\\
1.738	-0.0189292\\
1.74	-0.0189292\\
1.742	-0.0189292\\
1.744	-0.0189292\\
1.746	-0.0189292\\
1.748	-0.0189292\\
1.75	-0.0189292\\
1.752	-0.0189292\\
1.754	-0.0189292\\
1.756	-0.0189292\\
1.758	-0.0189292\\
1.76	-0.0189292\\
1.762	-0.0189292\\
1.764	-0.0189292\\
1.766	-0.0189292\\
1.768	-0.0189292\\
1.77	-0.0189292\\
1.772	-0.0189292\\
1.774	-0.0189292\\
1.776	-0.0189292\\
1.778	-0.0189292\\
1.78	-0.0189292\\
1.782	-0.0189292\\
1.784	-0.0189292\\
1.786	-0.0189292\\
1.788	-0.0189292\\
1.79	-0.0189292\\
1.792	-0.0189292\\
1.794	-0.0189292\\
1.796	-0.0189292\\
1.798	-0.0189292\\
1.8	-0.0189292\\
1.802	-0.0189292\\
1.804	-0.0189292\\
1.806	-0.0189292\\
1.808	-0.0189292\\
1.81	-0.0189292\\
1.812	-0.0189292\\
1.814	-0.0189292\\
1.816	-0.0189292\\
1.818	-0.0189292\\
1.82	-0.0189292\\
1.822	-0.0189292\\
1.824	-0.0189292\\
1.826	-0.0189292\\
1.828	-0.0189292\\
1.83	-0.0189292\\
1.832	-0.0189292\\
1.834	-0.0189292\\
1.836	-0.0189292\\
1.838	-0.0189292\\
1.84	-0.0189292\\
1.842	-0.0189292\\
1.844	-0.0189292\\
1.846	-0.0189292\\
1.848	-0.0189292\\
1.85	-0.0189292\\
1.852	-0.0189292\\
1.854	-0.0189292\\
1.856	-0.0189292\\
1.858	-0.0189292\\
1.86	-0.0189292\\
1.862	-0.0189292\\
1.864	-0.0189292\\
1.866	-0.0189292\\
1.868	-0.0189292\\
1.87	-0.0189292\\
1.872	-0.0189292\\
1.874	-0.0189292\\
1.876	-0.0189292\\
1.878	-0.0189292\\
1.88	-0.0189292\\
1.882	-0.0189292\\
1.884	-0.0189292\\
1.886	-0.0189292\\
1.888	-0.0189292\\
1.89	-0.0189292\\
1.892	-0.0189292\\
1.894	-0.0189292\\
1.896	-0.0189292\\
1.898	-0.0189292\\
1.9	-0.0189292\\
1.902	-0.0189292\\
1.904	-0.0189292\\
1.906	-0.0189292\\
1.908	-0.0189292\\
1.91	-0.0189292\\
1.912	-0.0189292\\
1.914	-0.0189292\\
1.916	-0.0189292\\
1.918	-0.0189292\\
1.92	-0.0189292\\
1.922	-0.0189292\\
1.924	-0.0189292\\
1.926	-0.0189292\\
1.928	-0.0189292\\
1.93	-0.0189292\\
1.932	-0.0189292\\
1.934	-0.0189292\\
1.936	-0.0189292\\
1.938	-0.0189292\\
1.94	-0.0189292\\
1.942	-0.0189292\\
1.944	-0.0189292\\
1.946	-0.0189292\\
1.948	-0.0189292\\
1.95	-0.0189292\\
1.952	-0.0189292\\
1.954	-0.0189292\\
1.956	-0.0189292\\
1.958	-0.0189292\\
1.96	-0.0189292\\
1.962	-0.0189292\\
1.964	-0.0189292\\
1.966	-0.0189292\\
1.968	-0.0189292\\
1.97	-0.0189292\\
1.972	-0.0189292\\
1.974	-0.0189292\\
1.976	-0.0189292\\
1.978	-0.0189292\\
1.98	-0.0189292\\
1.982	-0.0189292\\
1.984	-0.0189292\\
1.986	-0.0189292\\
1.988	-0.0189292\\
1.99	-0.0189292\\
1.992	-0.0189292\\
1.994	-0.0189292\\
1.996	-0.0189292\\
1.998	-0.0189292\\
2	-0.0189292\\
};
\addplot [color=mycolor10, line width=1.0pt, forget plot]
  table[row sep=crcr]{%
-0.024	0\\
-0.022	0\\
-0.02	0\\
-0.018	0\\
-0.016	0\\
-0.014	0\\
-0.012	0\\
-0.01	0\\
-0.008	0\\
-0.006	0\\
-0.004	0\\
-0.002	0\\
0	0\\
0.002	-0.00183165\\
0.004	-0.00338005\\
0.006	-0.00471028\\
0.008	-0.00587111\\
0.01	-0.00689895\\
0.012	-0.00782084\\
0.014	-0.00865687\\
0.016	-0.00942191\\
0.018	-0.010127\\
0.02	-0.0107805\\
0.022	-0.0113887\\
0.024	-0.0119563\\
0.026	-0.0124871\\
0.028	-0.0129842\\
0.03	-0.01345\\
0.032	-0.0138867\\
0.034	-0.014296\\
0.036	-0.0146796\\
0.038	-0.0150389\\
0.04	-0.0153752\\
0.042	-0.01569\\
0.044	-0.0159842\\
0.046	-0.0162592\\
0.048	-0.0165161\\
0.05	-0.0167557\\
0.052	-0.0169793\\
0.054	-0.0171878\\
0.056	-0.0173821\\
0.058	-0.0175631\\
0.06	-0.0177318\\
0.062	-0.0178888\\
0.064	-0.0180351\\
0.066	-0.0181712\\
0.068	-0.0182979\\
0.07	-0.0184159\\
0.072	-0.0185256\\
0.074	-0.0186277\\
0.076	-0.0187225\\
0.078	-0.0188106\\
0.08	-0.0188924\\
0.082	-0.0189682\\
0.084	-0.0190383\\
0.086	-0.0191031\\
0.088	-0.0191628\\
0.09	-0.0192177\\
0.092	-0.019268\\
0.094	-0.019314\\
0.096	-0.0193557\\
0.098	-0.0193935\\
0.1	-0.0194275\\
0.102	-0.0194578\\
0.104	-0.0194847\\
0.106	-0.0195082\\
0.108	-0.0195285\\
0.11	-0.0195458\\
0.112	-0.0195602\\
0.114	-0.0195719\\
0.116	-0.0195809\\
0.118	-0.0195875\\
0.12	-0.0195917\\
0.122	-0.0195937\\
0.124	-0.0195936\\
0.126	-0.0195915\\
0.128	-0.0195876\\
0.13	-0.019582\\
0.132	-0.0195748\\
0.134	-0.0195661\\
0.136	-0.0195561\\
0.138	-0.0195448\\
0.14	-0.0195324\\
0.142	-0.0195189\\
0.144	-0.0195045\\
0.146	-0.0194893\\
0.148	-0.0194733\\
0.15	-0.0194567\\
0.152	-0.0194395\\
0.154	-0.0194218\\
0.156	-0.0194037\\
0.158	-0.0193854\\
0.16	-0.0193667\\
0.162	-0.0193479\\
0.164	-0.019329\\
0.166	-0.01931\\
0.168	-0.019291\\
0.17	-0.0192721\\
0.172	-0.0192533\\
0.174	-0.0192346\\
0.176	-0.0192161\\
0.178	-0.0191978\\
0.18	-0.0191798\\
0.182	-0.0191621\\
0.184	-0.0191448\\
0.186	-0.0191278\\
0.188	-0.0191111\\
0.19	-0.0190949\\
0.192	-0.0190791\\
0.194	-0.0190638\\
0.196	-0.0190489\\
0.198	-0.0190345\\
0.2	-0.0190205\\
0.202	-0.0190071\\
0.204	-0.0189941\\
0.206	-0.0189817\\
0.208	-0.0189697\\
0.21	-0.0189583\\
0.212	-0.0189474\\
0.214	-0.018937\\
0.216	-0.0189271\\
0.218	-0.0189177\\
0.22	-0.0189088\\
0.222	-0.0189004\\
0.224	-0.0188925\\
0.226	-0.0188851\\
0.228	-0.0188781\\
0.23	-0.0188716\\
0.232	-0.0188656\\
0.234	-0.01886\\
0.236	-0.0188549\\
0.238	-0.0188501\\
0.24	-0.0188458\\
0.242	-0.0188419\\
0.244	-0.0188384\\
0.246	-0.0188352\\
0.248	-0.0188324\\
0.25	-0.0188299\\
0.252	-0.0188278\\
0.254	-0.018826\\
0.256	-0.0188244\\
0.258	-0.0188232\\
0.26	-0.0188223\\
0.262	-0.0188216\\
0.264	-0.0188211\\
0.266	-0.0188209\\
0.268	-0.0188209\\
0.27	-0.0188211\\
0.272	-0.0188216\\
0.274	-0.0188222\\
0.276	-0.0188229\\
0.278	-0.0188239\\
0.28	-0.0188249\\
0.282	-0.0188261\\
0.284	-0.0188275\\
0.286	-0.0188289\\
0.288	-0.0188305\\
0.29	-0.0188321\\
0.292	-0.0188339\\
0.294	-0.0188357\\
0.296	-0.0188376\\
0.298	-0.0188395\\
0.3	-0.0188415\\
0.302	-0.0188435\\
0.304	-0.0188456\\
0.306	-0.0188476\\
0.308	-0.0188497\\
0.31	-0.0188519\\
0.312	-0.018854\\
0.314	-0.0188561\\
0.316	-0.0188583\\
0.318	-0.0188604\\
0.32	-0.0188625\\
0.322	-0.0188646\\
0.324	-0.0188667\\
0.326	-0.0188687\\
0.328	-0.0188707\\
0.33	-0.0188727\\
0.332	-0.0188747\\
0.334	-0.0188766\\
0.336	-0.0188785\\
0.338	-0.0188803\\
0.34	-0.0188822\\
0.342	-0.0188839\\
0.344	-0.0188856\\
0.346	-0.0188873\\
0.348	-0.018889\\
0.35	-0.0188905\\
0.352	-0.0188921\\
0.354	-0.0188936\\
0.356	-0.018895\\
0.358	-0.0188964\\
0.36	-0.0188978\\
0.362	-0.0188991\\
0.364	-0.0189003\\
0.366	-0.0189016\\
0.368	-0.0189027\\
0.37	-0.0189039\\
0.372	-0.0189049\\
0.374	-0.018906\\
0.376	-0.018907\\
0.378	-0.0189079\\
0.38	-0.0189089\\
0.382	-0.0189097\\
0.384	-0.0189106\\
0.386	-0.0189114\\
0.388	-0.0189122\\
0.39	-0.0189129\\
0.392	-0.0189136\\
0.394	-0.0189143\\
0.396	-0.0189149\\
0.398	-0.0189155\\
0.4	-0.0189161\\
0.402	-0.0189166\\
0.404	-0.0189172\\
0.406	-0.0189177\\
0.408	-0.0189181\\
0.41	-0.0189186\\
0.412	-0.018919\\
0.414	-0.0189194\\
0.416	-0.0189198\\
0.418	-0.0189202\\
0.42	-0.0189205\\
0.422	-0.0189209\\
0.424	-0.0189212\\
0.426	-0.0189215\\
0.428	-0.0189218\\
0.43	-0.018922\\
0.432	-0.0189223\\
0.434	-0.0189225\\
0.436	-0.0189227\\
0.438	-0.018923\\
0.44	-0.0189232\\
0.442	-0.0189233\\
0.444	-0.0189235\\
0.446	-0.0189237\\
0.448	-0.0189238\\
0.45	-0.018924\\
0.452	-0.0189241\\
0.454	-0.0189243\\
0.456	-0.0189244\\
0.458	-0.0189245\\
0.46	-0.0189246\\
0.462	-0.0189247\\
0.464	-0.0189248\\
0.466	-0.0189249\\
0.468	-0.018925\\
0.47	-0.0189251\\
0.472	-0.0189251\\
0.474	-0.0189252\\
0.476	-0.0189253\\
0.478	-0.0189253\\
0.48	-0.0189254\\
0.482	-0.0189254\\
0.484	-0.0189255\\
0.486	-0.0189255\\
0.488	-0.0189256\\
0.49	-0.0189256\\
0.492	-0.0189256\\
0.494	-0.0189257\\
0.496	-0.0189257\\
0.498	-0.0189257\\
0.5	-0.0189258\\
0.502	-0.0189258\\
0.504	-0.0189258\\
0.506	-0.0189258\\
0.508	-0.0189259\\
0.51	-0.0189259\\
0.512	-0.0189259\\
0.514	-0.0189259\\
0.516	-0.0189259\\
0.518	-0.0189259\\
0.52	-0.018926\\
0.522	-0.018926\\
0.524	-0.018926\\
0.526	-0.018926\\
0.528	-0.018926\\
0.53	-0.018926\\
0.532	-0.018926\\
0.534	-0.018926\\
0.536	-0.0189261\\
0.538	-0.0189261\\
0.54	-0.0189261\\
0.542	-0.0189261\\
0.544	-0.0189261\\
0.546	-0.0189261\\
0.548	-0.0189261\\
0.55	-0.0189261\\
0.552	-0.0189261\\
0.554	-0.0189262\\
0.556	-0.0189262\\
0.558	-0.0189262\\
0.56	-0.0189262\\
0.562	-0.0189262\\
0.564	-0.0189262\\
0.566	-0.0189262\\
0.568	-0.0189263\\
0.57	-0.0189263\\
0.572	-0.0189263\\
0.574	-0.0189263\\
0.576	-0.0189263\\
0.578	-0.0189263\\
0.58	-0.0189264\\
0.582	-0.0189264\\
0.584	-0.0189264\\
0.586	-0.0189264\\
0.588	-0.0189264\\
0.59	-0.0189265\\
0.592	-0.0189265\\
0.594	-0.0189265\\
0.596	-0.0189265\\
0.598	-0.0189266\\
0.6	-0.0189266\\
0.602	-0.0189266\\
0.604	-0.0189266\\
0.606	-0.0189267\\
0.608	-0.0189267\\
0.61	-0.0189267\\
0.612	-0.0189267\\
0.614	-0.0189268\\
0.616	-0.0189268\\
0.618	-0.0189268\\
0.62	-0.0189268\\
0.622	-0.0189269\\
0.624	-0.0189269\\
0.626	-0.0189269\\
0.628	-0.018927\\
0.63	-0.018927\\
0.632	-0.018927\\
0.634	-0.0189271\\
0.636	-0.0189271\\
0.638	-0.0189271\\
0.64	-0.0189271\\
0.642	-0.0189272\\
0.644	-0.0189272\\
0.646	-0.0189272\\
0.648	-0.0189273\\
0.65	-0.0189273\\
0.652	-0.0189273\\
0.654	-0.0189274\\
0.656	-0.0189274\\
0.658	-0.0189274\\
0.66	-0.0189275\\
0.662	-0.0189275\\
0.664	-0.0189275\\
0.666	-0.0189276\\
0.668	-0.0189276\\
0.67	-0.0189276\\
0.672	-0.0189277\\
0.674	-0.0189277\\
0.676	-0.0189277\\
0.678	-0.0189278\\
0.68	-0.0189278\\
0.682	-0.0189278\\
0.684	-0.0189278\\
0.686	-0.0189279\\
0.688	-0.0189279\\
0.69	-0.0189279\\
0.692	-0.018928\\
0.694	-0.018928\\
0.696	-0.018928\\
0.698	-0.018928\\
0.7	-0.0189281\\
0.702	-0.0189281\\
0.704	-0.0189281\\
0.706	-0.0189282\\
0.708	-0.0189282\\
0.71	-0.0189282\\
0.712	-0.0189282\\
0.714	-0.0189283\\
0.716	-0.0189283\\
0.718	-0.0189283\\
0.72	-0.0189283\\
0.722	-0.0189283\\
0.724	-0.0189284\\
0.726	-0.0189284\\
0.728	-0.0189284\\
0.73	-0.0189284\\
0.732	-0.0189285\\
0.734	-0.0189285\\
0.736	-0.0189285\\
0.738	-0.0189285\\
0.74	-0.0189285\\
0.742	-0.0189285\\
0.744	-0.0189286\\
0.746	-0.0189286\\
0.748	-0.0189286\\
0.75	-0.0189286\\
0.752	-0.0189286\\
0.754	-0.0189286\\
0.756	-0.0189287\\
0.758	-0.0189287\\
0.76	-0.0189287\\
0.762	-0.0189287\\
0.764	-0.0189287\\
0.766	-0.0189287\\
0.768	-0.0189287\\
0.77	-0.0189287\\
0.772	-0.0189287\\
0.774	-0.0189288\\
0.776	-0.0189288\\
0.778	-0.0189288\\
0.78	-0.0189288\\
0.782	-0.0189288\\
0.784	-0.0189288\\
0.786	-0.0189288\\
0.788	-0.0189288\\
0.79	-0.0189288\\
0.792	-0.0189288\\
0.794	-0.0189288\\
0.796	-0.0189288\\
0.798	-0.0189289\\
0.8	-0.0189289\\
0.802	-0.0189289\\
0.804	-0.0189289\\
0.806	-0.0189289\\
0.808	-0.0189289\\
0.81	-0.0189289\\
0.812	-0.0189289\\
0.814	-0.0189289\\
0.816	-0.0189289\\
0.818	-0.0189289\\
0.82	-0.0189289\\
0.822	-0.0189289\\
0.824	-0.0189289\\
0.826	-0.0189289\\
0.828	-0.0189289\\
0.83	-0.0189289\\
0.832	-0.0189289\\
0.834	-0.0189289\\
0.836	-0.0189289\\
0.838	-0.0189289\\
0.84	-0.0189289\\
0.842	-0.0189289\\
0.844	-0.018929\\
0.846	-0.018929\\
0.848	-0.018929\\
0.85	-0.018929\\
0.852	-0.018929\\
0.854	-0.018929\\
0.856	-0.018929\\
0.858	-0.018929\\
0.86	-0.018929\\
0.862	-0.018929\\
0.864	-0.018929\\
0.866	-0.018929\\
0.868	-0.018929\\
0.87	-0.018929\\
0.872	-0.018929\\
0.874	-0.018929\\
0.876	-0.018929\\
0.878	-0.018929\\
0.88	-0.018929\\
0.882	-0.018929\\
0.884	-0.018929\\
0.886	-0.018929\\
0.888	-0.018929\\
0.89	-0.018929\\
0.892	-0.018929\\
0.894	-0.018929\\
0.896	-0.018929\\
0.898	-0.018929\\
0.9	-0.018929\\
0.902	-0.018929\\
0.904	-0.018929\\
0.906	-0.018929\\
0.908	-0.018929\\
0.91	-0.018929\\
0.912	-0.018929\\
0.914	-0.018929\\
0.916	-0.018929\\
0.918	-0.018929\\
0.92	-0.018929\\
0.922	-0.018929\\
0.924	-0.018929\\
0.926	-0.018929\\
0.928	-0.018929\\
0.93	-0.018929\\
0.932	-0.018929\\
0.934	-0.018929\\
0.936	-0.018929\\
0.938	-0.018929\\
0.94	-0.018929\\
0.942	-0.0189291\\
0.944	-0.0189291\\
0.946	-0.0189291\\
0.948	-0.0189291\\
0.95	-0.0189291\\
0.952	-0.0189291\\
0.954	-0.0189291\\
0.956	-0.0189291\\
0.958	-0.0189291\\
0.96	-0.0189291\\
0.962	-0.0189291\\
0.964	-0.0189291\\
0.966	-0.0189291\\
0.968	-0.0189291\\
0.97	-0.0189291\\
0.972	-0.0189291\\
0.974	-0.0189291\\
0.976	-0.0189291\\
0.978	-0.0189291\\
0.98	-0.0189291\\
0.982	-0.0189291\\
0.984	-0.0189291\\
0.986	-0.0189291\\
0.988	-0.0189291\\
0.99	-0.0189291\\
0.992	-0.0189291\\
0.994	-0.0189291\\
0.996	-0.0189291\\
0.998	-0.0189291\\
1	-0.0189291\\
1.002	-0.0189291\\
1.004	-0.0189291\\
1.006	-0.0189291\\
1.008	-0.0189291\\
1.01	-0.0189291\\
1.012	-0.0189291\\
1.014	-0.0189291\\
1.016	-0.0189291\\
1.018	-0.0189291\\
1.02	-0.0189291\\
1.022	-0.0189291\\
1.024	-0.0189291\\
1.026	-0.0189291\\
1.028	-0.0189291\\
1.03	-0.0189291\\
1.032	-0.0189291\\
1.034	-0.0189291\\
1.036	-0.0189291\\
1.038	-0.0189291\\
1.04	-0.0189291\\
1.042	-0.0189291\\
1.044	-0.0189291\\
1.046	-0.0189291\\
1.048	-0.0189291\\
1.05	-0.0189291\\
1.052	-0.0189291\\
1.054	-0.0189291\\
1.056	-0.0189291\\
1.058	-0.0189291\\
1.06	-0.0189291\\
1.062	-0.0189291\\
1.064	-0.0189291\\
1.066	-0.0189291\\
1.068	-0.0189291\\
1.07	-0.0189291\\
1.072	-0.0189291\\
1.074	-0.0189291\\
1.076	-0.0189291\\
1.078	-0.0189291\\
1.08	-0.0189291\\
1.082	-0.0189291\\
1.084	-0.0189291\\
1.086	-0.0189291\\
1.088	-0.0189291\\
1.09	-0.0189291\\
1.092	-0.0189291\\
1.094	-0.0189291\\
1.096	-0.0189291\\
1.098	-0.0189291\\
1.1	-0.0189291\\
1.102	-0.0189291\\
1.104	-0.0189291\\
1.106	-0.0189291\\
1.108	-0.0189291\\
1.11	-0.0189291\\
1.112	-0.0189291\\
1.114	-0.0189291\\
1.116	-0.0189291\\
1.118	-0.0189291\\
1.12	-0.0189291\\
1.122	-0.0189291\\
1.124	-0.0189291\\
1.126	-0.0189291\\
1.128	-0.0189291\\
1.13	-0.0189291\\
1.132	-0.0189291\\
1.134	-0.0189291\\
1.136	-0.0189291\\
1.138	-0.0189291\\
1.14	-0.0189291\\
1.142	-0.0189291\\
1.144	-0.0189291\\
1.146	-0.0189291\\
1.148	-0.0189291\\
1.15	-0.0189291\\
1.152	-0.0189291\\
1.154	-0.0189291\\
1.156	-0.0189291\\
1.158	-0.0189291\\
1.16	-0.0189291\\
1.162	-0.0189291\\
1.164	-0.0189291\\
1.166	-0.0189291\\
1.168	-0.0189291\\
1.17	-0.0189291\\
1.172	-0.0189291\\
1.174	-0.0189291\\
1.176	-0.0189291\\
1.178	-0.0189291\\
1.18	-0.0189291\\
1.182	-0.0189291\\
1.184	-0.0189291\\
1.186	-0.0189291\\
1.188	-0.0189291\\
1.19	-0.0189291\\
1.192	-0.0189291\\
1.194	-0.0189291\\
1.196	-0.0189291\\
1.198	-0.0189291\\
1.2	-0.0189291\\
1.202	-0.0189291\\
1.204	-0.0189291\\
1.206	-0.0189291\\
1.208	-0.0189291\\
1.21	-0.0189291\\
1.212	-0.0189291\\
1.214	-0.0189291\\
1.216	-0.0189291\\
1.218	-0.0189291\\
1.22	-0.0189291\\
1.222	-0.0189291\\
1.224	-0.0189291\\
1.226	-0.0189291\\
1.228	-0.0189291\\
1.23	-0.0189291\\
1.232	-0.0189291\\
1.234	-0.0189291\\
1.236	-0.0189291\\
1.238	-0.0189292\\
1.24	-0.0189292\\
1.242	-0.0189292\\
1.244	-0.0189292\\
1.246	-0.0189292\\
1.248	-0.0189292\\
1.25	-0.0189292\\
1.252	-0.0189292\\
1.254	-0.0189292\\
1.256	-0.0189292\\
1.258	-0.0189292\\
1.26	-0.0189292\\
1.262	-0.0189292\\
1.264	-0.0189292\\
1.266	-0.0189292\\
1.268	-0.0189292\\
1.27	-0.0189292\\
1.272	-0.0189292\\
1.274	-0.0189292\\
1.276	-0.0189292\\
1.278	-0.0189292\\
1.28	-0.0189292\\
1.282	-0.0189292\\
1.284	-0.0189292\\
1.286	-0.0189292\\
1.288	-0.0189292\\
1.29	-0.0189292\\
1.292	-0.0189292\\
1.294	-0.0189292\\
1.296	-0.0189292\\
1.298	-0.0189292\\
1.3	-0.0189292\\
1.302	-0.0189292\\
1.304	-0.0189292\\
1.306	-0.0189292\\
1.308	-0.0189292\\
1.31	-0.0189292\\
1.312	-0.0189292\\
1.314	-0.0189292\\
1.316	-0.0189292\\
1.318	-0.0189292\\
1.32	-0.0189292\\
1.322	-0.0189292\\
1.324	-0.0189292\\
1.326	-0.0189292\\
1.328	-0.0189292\\
1.33	-0.0189292\\
1.332	-0.0189292\\
1.334	-0.0189292\\
1.336	-0.0189292\\
1.338	-0.0189292\\
1.34	-0.0189292\\
1.342	-0.0189292\\
1.344	-0.0189292\\
1.346	-0.0189292\\
1.348	-0.0189292\\
1.35	-0.0189292\\
1.352	-0.0189292\\
1.354	-0.0189292\\
1.356	-0.0189292\\
1.358	-0.0189292\\
1.36	-0.0189292\\
1.362	-0.0189292\\
1.364	-0.0189292\\
1.366	-0.0189292\\
1.368	-0.0189292\\
1.37	-0.0189292\\
1.372	-0.0189292\\
1.374	-0.0189292\\
1.376	-0.0189292\\
1.378	-0.0189292\\
1.38	-0.0189292\\
1.382	-0.0189292\\
1.384	-0.0189292\\
1.386	-0.0189292\\
1.388	-0.0189292\\
1.39	-0.0189292\\
1.392	-0.0189292\\
1.394	-0.0189292\\
1.396	-0.0189292\\
1.398	-0.0189292\\
1.4	-0.0189292\\
1.402	-0.0189292\\
1.404	-0.0189292\\
1.406	-0.0189292\\
1.408	-0.0189292\\
1.41	-0.0189292\\
1.412	-0.0189292\\
1.414	-0.0189292\\
1.416	-0.0189292\\
1.418	-0.0189292\\
1.42	-0.0189292\\
1.422	-0.0189292\\
1.424	-0.0189292\\
1.426	-0.0189292\\
1.428	-0.0189292\\
1.43	-0.0189292\\
1.432	-0.0189292\\
1.434	-0.0189292\\
1.436	-0.0189292\\
1.438	-0.0189292\\
1.44	-0.0189292\\
1.442	-0.0189292\\
1.444	-0.0189292\\
1.446	-0.0189292\\
1.448	-0.0189292\\
1.45	-0.0189292\\
1.452	-0.0189292\\
1.454	-0.0189292\\
1.456	-0.0189292\\
1.458	-0.0189292\\
1.46	-0.0189292\\
1.462	-0.0189292\\
1.464	-0.0189292\\
1.466	-0.0189292\\
1.468	-0.0189292\\
1.47	-0.0189292\\
1.472	-0.0189292\\
1.474	-0.0189292\\
1.476	-0.0189292\\
1.478	-0.0189292\\
1.48	-0.0189292\\
1.482	-0.0189292\\
1.484	-0.0189292\\
1.486	-0.0189292\\
1.488	-0.0189292\\
1.49	-0.0189292\\
1.492	-0.0189292\\
1.494	-0.0189292\\
1.496	-0.0189292\\
1.498	-0.0189292\\
1.5	-0.0189292\\
1.502	-0.0189292\\
1.504	-0.0189292\\
1.506	-0.0189292\\
1.508	-0.0189292\\
1.51	-0.0189292\\
1.512	-0.0189292\\
1.514	-0.0189292\\
1.516	-0.0189292\\
1.518	-0.0189292\\
1.52	-0.0189292\\
1.522	-0.0189292\\
1.524	-0.0189292\\
1.526	-0.0189292\\
1.528	-0.0189292\\
1.53	-0.0189292\\
1.532	-0.0189292\\
1.534	-0.0189292\\
1.536	-0.0189292\\
1.538	-0.0189292\\
1.54	-0.0189292\\
1.542	-0.0189292\\
1.544	-0.0189292\\
1.546	-0.0189292\\
1.548	-0.0189292\\
1.55	-0.0189292\\
1.552	-0.0189292\\
1.554	-0.0189292\\
1.556	-0.0189292\\
1.558	-0.0189292\\
1.56	-0.0189292\\
1.562	-0.0189292\\
1.564	-0.0189292\\
1.566	-0.0189292\\
1.568	-0.0189292\\
1.57	-0.0189292\\
1.572	-0.0189292\\
1.574	-0.0189292\\
1.576	-0.0189292\\
1.578	-0.0189292\\
1.58	-0.0189292\\
1.582	-0.0189292\\
1.584	-0.0189292\\
1.586	-0.0189292\\
1.588	-0.0189292\\
1.59	-0.0189292\\
1.592	-0.0189292\\
1.594	-0.0189292\\
1.596	-0.0189292\\
1.598	-0.0189292\\
1.6	-0.0189292\\
1.602	-0.0189292\\
1.604	-0.0189292\\
1.606	-0.0189292\\
1.608	-0.0189292\\
1.61	-0.0189292\\
1.612	-0.0189292\\
1.614	-0.0189292\\
1.616	-0.0189292\\
1.618	-0.0189292\\
1.62	-0.0189292\\
1.622	-0.0189292\\
1.624	-0.0189292\\
1.626	-0.0189292\\
1.628	-0.0189292\\
1.63	-0.0189292\\
1.632	-0.0189292\\
1.634	-0.0189292\\
1.636	-0.0189292\\
1.638	-0.0189292\\
1.64	-0.0189292\\
1.642	-0.0189292\\
1.644	-0.0189292\\
1.646	-0.0189292\\
1.648	-0.0189292\\
1.65	-0.0189292\\
1.652	-0.0189292\\
1.654	-0.0189292\\
1.656	-0.0189292\\
1.658	-0.0189292\\
1.66	-0.0189292\\
1.662	-0.0189292\\
1.664	-0.0189292\\
1.666	-0.0189292\\
1.668	-0.0189292\\
1.67	-0.0189292\\
1.672	-0.0189292\\
1.674	-0.0189292\\
1.676	-0.0189292\\
1.678	-0.0189292\\
1.68	-0.0189292\\
1.682	-0.0189292\\
1.684	-0.0189292\\
1.686	-0.0189292\\
1.688	-0.0189292\\
1.69	-0.0189292\\
1.692	-0.0189292\\
1.694	-0.0189292\\
1.696	-0.0189292\\
1.698	-0.0189292\\
1.7	-0.0189292\\
1.702	-0.0189292\\
1.704	-0.0189292\\
1.706	-0.0189292\\
1.708	-0.0189292\\
1.71	-0.0189292\\
1.712	-0.0189292\\
1.714	-0.0189292\\
1.716	-0.0189292\\
1.718	-0.0189292\\
1.72	-0.0189292\\
1.722	-0.0189292\\
1.724	-0.0189292\\
1.726	-0.0189292\\
1.728	-0.0189292\\
1.73	-0.0189292\\
1.732	-0.0189292\\
1.734	-0.0189292\\
1.736	-0.0189292\\
1.738	-0.0189292\\
1.74	-0.0189292\\
1.742	-0.0189292\\
1.744	-0.0189292\\
1.746	-0.0189292\\
1.748	-0.0189292\\
1.75	-0.0189292\\
1.752	-0.0189292\\
1.754	-0.0189292\\
1.756	-0.0189292\\
1.758	-0.0189292\\
1.76	-0.0189292\\
1.762	-0.0189292\\
1.764	-0.0189292\\
1.766	-0.0189292\\
1.768	-0.0189292\\
1.77	-0.0189292\\
1.772	-0.0189292\\
1.774	-0.0189292\\
1.776	-0.0189292\\
1.778	-0.0189292\\
1.78	-0.0189292\\
1.782	-0.0189292\\
1.784	-0.0189292\\
1.786	-0.0189292\\
1.788	-0.0189292\\
1.79	-0.0189292\\
1.792	-0.0189292\\
1.794	-0.0189292\\
1.796	-0.0189292\\
1.798	-0.0189292\\
1.8	-0.0189292\\
1.802	-0.0189292\\
1.804	-0.0189292\\
1.806	-0.0189292\\
1.808	-0.0189292\\
1.81	-0.0189292\\
1.812	-0.0189292\\
1.814	-0.0189292\\
1.816	-0.0189292\\
1.818	-0.0189292\\
1.82	-0.0189292\\
1.822	-0.0189292\\
1.824	-0.0189292\\
1.826	-0.0189292\\
1.828	-0.0189292\\
1.83	-0.0189292\\
1.832	-0.0189292\\
1.834	-0.0189292\\
1.836	-0.0189292\\
1.838	-0.0189292\\
1.84	-0.0189292\\
1.842	-0.0189292\\
1.844	-0.0189292\\
1.846	-0.0189292\\
1.848	-0.0189292\\
1.85	-0.0189292\\
1.852	-0.0189292\\
1.854	-0.0189292\\
1.856	-0.0189292\\
1.858	-0.0189292\\
1.86	-0.0189292\\
1.862	-0.0189292\\
1.864	-0.0189292\\
1.866	-0.0189292\\
1.868	-0.0189292\\
1.87	-0.0189292\\
1.872	-0.0189292\\
1.874	-0.0189292\\
1.876	-0.0189292\\
1.878	-0.0189292\\
1.88	-0.0189292\\
1.882	-0.0189292\\
1.884	-0.0189292\\
1.886	-0.0189292\\
1.888	-0.0189292\\
1.89	-0.0189292\\
1.892	-0.0189292\\
1.894	-0.0189292\\
1.896	-0.0189292\\
1.898	-0.0189292\\
1.9	-0.0189292\\
1.902	-0.0189292\\
1.904	-0.0189292\\
1.906	-0.0189292\\
1.908	-0.0189292\\
1.91	-0.0189292\\
1.912	-0.0189292\\
1.914	-0.0189292\\
1.916	-0.0189292\\
1.918	-0.0189292\\
1.92	-0.0189292\\
1.922	-0.0189292\\
1.924	-0.0189292\\
1.926	-0.0189292\\
1.928	-0.0189292\\
1.93	-0.0189292\\
1.932	-0.0189292\\
1.934	-0.0189292\\
1.936	-0.0189292\\
1.938	-0.0189292\\
1.94	-0.0189292\\
1.942	-0.0189292\\
1.944	-0.0189292\\
1.946	-0.0189292\\
1.948	-0.0189292\\
1.95	-0.0189292\\
1.952	-0.0189292\\
1.954	-0.0189292\\
1.956	-0.0189292\\
1.958	-0.0189292\\
1.96	-0.0189292\\
1.962	-0.0189292\\
1.964	-0.0189292\\
1.966	-0.0189292\\
1.968	-0.0189292\\
1.97	-0.0189292\\
1.972	-0.0189292\\
1.974	-0.0189292\\
1.976	-0.0189292\\
1.978	-0.0189292\\
1.98	-0.0189292\\
1.982	-0.0189292\\
1.984	-0.0189292\\
1.986	-0.0189292\\
1.988	-0.0189292\\
1.99	-0.0189292\\
1.992	-0.0189292\\
1.994	-0.0189292\\
1.996	-0.0189292\\
1.998	-0.0189292\\
2	-0.0189292\\
};
\addplot [color=mycolor11, line width=1.0pt, forget plot]
  table[row sep=crcr]{%
-0.024	0\\
-0.022	0\\
-0.02	0\\
-0.018	0\\
-0.016	0\\
-0.014	0\\
-0.012	0\\
-0.01	0\\
-0.008	0\\
-0.006	0\\
-0.004	0\\
-0.002	0\\
0	0\\
0.002	-0.00183223\\
0.004	-0.0033837\\
0.006	-0.00472\\
0.008	-0.0058893\\
0.01	-0.00692694\\
0.012	-0.00785889\\
0.014	-0.00870434\\
0.016	-0.0094775\\
0.018	-0.010189\\
0.02	-0.0108471\\
0.022	-0.0114578\\
0.024	-0.0120262\\
0.026	-0.0125563\\
0.028	-0.0130512\\
0.03	-0.0135137\\
0.032	-0.0139463\\
0.034	-0.0143508\\
0.036	-0.0147292\\
0.038	-0.0150831\\
0.04	-0.0154141\\
0.042	-0.0157235\\
0.044	-0.0160126\\
0.046	-0.0162828\\
0.048	-0.016535\\
0.05	-0.0167705\\
0.052	-0.0169902\\
0.054	-0.0171952\\
0.056	-0.0173864\\
0.058	-0.0175646\\
0.06	-0.0177307\\
0.062	-0.0178855\\
0.064	-0.0180297\\
0.066	-0.0181639\\
0.068	-0.0182889\\
0.07	-0.0184053\\
0.072	-0.0185135\\
0.074	-0.0186141\\
0.076	-0.0187076\\
0.078	-0.0187943\\
0.08	-0.0188748\\
0.082	-0.0189493\\
0.084	-0.0190182\\
0.086	-0.0190817\\
0.088	-0.0191403\\
0.09	-0.0191941\\
0.092	-0.0192433\\
0.094	-0.0192882\\
0.096	-0.0193291\\
0.098	-0.019366\\
0.1	-0.0193991\\
0.102	-0.0194287\\
0.104	-0.0194549\\
0.106	-0.0194779\\
0.108	-0.0194978\\
0.11	-0.0195147\\
0.112	-0.0195288\\
0.114	-0.0195403\\
0.116	-0.0195493\\
0.118	-0.0195558\\
0.12	-0.0195601\\
0.122	-0.0195623\\
0.124	-0.0195624\\
0.126	-0.0195607\\
0.128	-0.0195572\\
0.13	-0.019552\\
0.132	-0.0195454\\
0.134	-0.0195373\\
0.136	-0.0195279\\
0.138	-0.0195174\\
0.14	-0.0195057\\
0.142	-0.019493\\
0.144	-0.0194795\\
0.146	-0.0194652\\
0.148	-0.0194501\\
0.15	-0.0194345\\
0.152	-0.0194183\\
0.154	-0.0194017\\
0.156	-0.0193846\\
0.158	-0.0193673\\
0.16	-0.0193498\\
0.162	-0.019332\\
0.164	-0.0193142\\
0.166	-0.0192963\\
0.168	-0.0192783\\
0.17	-0.0192605\\
0.172	-0.0192427\\
0.174	-0.0192251\\
0.176	-0.0192076\\
0.178	-0.0191904\\
0.18	-0.0191734\\
0.182	-0.0191566\\
0.184	-0.0191402\\
0.186	-0.0191242\\
0.188	-0.0191084\\
0.19	-0.0190931\\
0.192	-0.0190781\\
0.194	-0.0190636\\
0.196	-0.0190495\\
0.198	-0.0190358\\
0.2	-0.0190226\\
0.202	-0.0190099\\
0.204	-0.0189976\\
0.206	-0.0189858\\
0.208	-0.0189744\\
0.21	-0.0189635\\
0.212	-0.0189532\\
0.214	-0.0189432\\
0.216	-0.0189338\\
0.218	-0.0189248\\
0.22	-0.0189163\\
0.222	-0.0189083\\
0.224	-0.0189007\\
0.226	-0.0188936\\
0.228	-0.0188869\\
0.23	-0.0188807\\
0.232	-0.0188748\\
0.234	-0.0188694\\
0.236	-0.0188644\\
0.238	-0.0188598\\
0.24	-0.0188555\\
0.242	-0.0188517\\
0.244	-0.0188482\\
0.246	-0.018845\\
0.248	-0.0188422\\
0.25	-0.0188397\\
0.252	-0.0188375\\
0.254	-0.0188356\\
0.256	-0.018834\\
0.258	-0.0188326\\
0.26	-0.0188315\\
0.262	-0.0188307\\
0.264	-0.0188301\\
0.266	-0.0188297\\
0.268	-0.0188295\\
0.27	-0.0188295\\
0.272	-0.0188297\\
0.274	-0.0188301\\
0.276	-0.0188306\\
0.278	-0.0188313\\
0.28	-0.0188321\\
0.282	-0.0188331\\
0.284	-0.0188341\\
0.286	-0.0188353\\
0.288	-0.0188366\\
0.29	-0.018838\\
0.292	-0.0188394\\
0.294	-0.018841\\
0.296	-0.0188426\\
0.298	-0.0188442\\
0.3	-0.0188459\\
0.302	-0.0188477\\
0.304	-0.0188494\\
0.306	-0.0188513\\
0.308	-0.0188531\\
0.31	-0.0188549\\
0.312	-0.0188568\\
0.314	-0.0188587\\
0.316	-0.0188606\\
0.318	-0.0188624\\
0.32	-0.0188643\\
0.322	-0.0188661\\
0.324	-0.018868\\
0.326	-0.0188698\\
0.328	-0.0188716\\
0.33	-0.0188734\\
0.332	-0.0188752\\
0.334	-0.0188769\\
0.336	-0.0188786\\
0.338	-0.0188802\\
0.34	-0.0188819\\
0.342	-0.0188835\\
0.344	-0.018885\\
0.346	-0.0188865\\
0.348	-0.018888\\
0.35	-0.0188895\\
0.352	-0.0188909\\
0.354	-0.0188922\\
0.356	-0.0188936\\
0.358	-0.0188949\\
0.36	-0.0188961\\
0.362	-0.0188973\\
0.364	-0.0188985\\
0.366	-0.0188996\\
0.368	-0.0189007\\
0.37	-0.0189018\\
0.372	-0.0189028\\
0.374	-0.0189038\\
0.376	-0.0189047\\
0.378	-0.0189056\\
0.38	-0.0189065\\
0.382	-0.0189073\\
0.384	-0.0189081\\
0.386	-0.0189089\\
0.388	-0.0189097\\
0.39	-0.0189104\\
0.392	-0.0189111\\
0.394	-0.0189117\\
0.396	-0.0189124\\
0.398	-0.018913\\
0.4	-0.0189136\\
0.402	-0.0189141\\
0.404	-0.0189146\\
0.406	-0.0189152\\
0.408	-0.0189157\\
0.41	-0.0189161\\
0.412	-0.0189166\\
0.414	-0.018917\\
0.416	-0.0189174\\
0.418	-0.0189178\\
0.42	-0.0189182\\
0.422	-0.0189186\\
0.424	-0.0189189\\
0.426	-0.0189192\\
0.428	-0.0189195\\
0.43	-0.0189198\\
0.432	-0.0189201\\
0.434	-0.0189204\\
0.436	-0.0189207\\
0.438	-0.0189209\\
0.44	-0.0189212\\
0.442	-0.0189214\\
0.444	-0.0189216\\
0.446	-0.0189218\\
0.448	-0.018922\\
0.45	-0.0189222\\
0.452	-0.0189224\\
0.454	-0.0189226\\
0.456	-0.0189227\\
0.458	-0.0189229\\
0.46	-0.018923\\
0.462	-0.0189232\\
0.464	-0.0189233\\
0.466	-0.0189235\\
0.468	-0.0189236\\
0.47	-0.0189237\\
0.472	-0.0189238\\
0.474	-0.0189239\\
0.476	-0.018924\\
0.478	-0.0189241\\
0.48	-0.0189242\\
0.482	-0.0189243\\
0.484	-0.0189244\\
0.486	-0.0189245\\
0.488	-0.0189245\\
0.49	-0.0189246\\
0.492	-0.0189247\\
0.494	-0.0189248\\
0.496	-0.0189248\\
0.498	-0.0189249\\
0.5	-0.0189249\\
0.502	-0.018925\\
0.504	-0.018925\\
0.506	-0.0189251\\
0.508	-0.0189251\\
0.51	-0.0189252\\
0.512	-0.0189252\\
0.514	-0.0189253\\
0.516	-0.0189253\\
0.518	-0.0189253\\
0.52	-0.0189254\\
0.522	-0.0189254\\
0.524	-0.0189255\\
0.526	-0.0189255\\
0.528	-0.0189255\\
0.53	-0.0189255\\
0.532	-0.0189256\\
0.534	-0.0189256\\
0.536	-0.0189256\\
0.538	-0.0189257\\
0.54	-0.0189257\\
0.542	-0.0189257\\
0.544	-0.0189257\\
0.546	-0.0189258\\
0.548	-0.0189258\\
0.55	-0.0189258\\
0.552	-0.0189258\\
0.554	-0.0189259\\
0.556	-0.0189259\\
0.558	-0.0189259\\
0.56	-0.0189259\\
0.562	-0.0189259\\
0.564	-0.018926\\
0.566	-0.018926\\
0.568	-0.018926\\
0.57	-0.018926\\
0.572	-0.0189261\\
0.574	-0.0189261\\
0.576	-0.0189261\\
0.578	-0.0189261\\
0.58	-0.0189262\\
0.582	-0.0189262\\
0.584	-0.0189262\\
0.586	-0.0189262\\
0.588	-0.0189263\\
0.59	-0.0189263\\
0.592	-0.0189263\\
0.594	-0.0189263\\
0.596	-0.0189264\\
0.598	-0.0189264\\
0.6	-0.0189264\\
0.602	-0.0189264\\
0.604	-0.0189265\\
0.606	-0.0189265\\
0.608	-0.0189265\\
0.61	-0.0189265\\
0.612	-0.0189266\\
0.614	-0.0189266\\
0.616	-0.0189266\\
0.618	-0.0189267\\
0.62	-0.0189267\\
0.622	-0.0189267\\
0.624	-0.0189268\\
0.626	-0.0189268\\
0.628	-0.0189268\\
0.63	-0.0189268\\
0.632	-0.0189269\\
0.634	-0.0189269\\
0.636	-0.0189269\\
0.638	-0.018927\\
0.64	-0.018927\\
0.642	-0.018927\\
0.644	-0.0189271\\
0.646	-0.0189271\\
0.648	-0.0189271\\
0.65	-0.0189272\\
0.652	-0.0189272\\
0.654	-0.0189272\\
0.656	-0.0189273\\
0.658	-0.0189273\\
0.66	-0.0189273\\
0.662	-0.0189274\\
0.664	-0.0189274\\
0.666	-0.0189274\\
0.668	-0.0189275\\
0.67	-0.0189275\\
0.672	-0.0189275\\
0.674	-0.0189276\\
0.676	-0.0189276\\
0.678	-0.0189276\\
0.68	-0.0189277\\
0.682	-0.0189277\\
0.684	-0.0189277\\
0.686	-0.0189277\\
0.688	-0.0189278\\
0.69	-0.0189278\\
0.692	-0.0189278\\
0.694	-0.0189279\\
0.696	-0.0189279\\
0.698	-0.0189279\\
0.7	-0.0189279\\
0.702	-0.018928\\
0.704	-0.018928\\
0.706	-0.018928\\
0.708	-0.0189281\\
0.71	-0.0189281\\
0.712	-0.0189281\\
0.714	-0.0189281\\
0.716	-0.0189282\\
0.718	-0.0189282\\
0.72	-0.0189282\\
0.722	-0.0189282\\
0.724	-0.0189283\\
0.726	-0.0189283\\
0.728	-0.0189283\\
0.73	-0.0189283\\
0.732	-0.0189283\\
0.734	-0.0189284\\
0.736	-0.0189284\\
0.738	-0.0189284\\
0.74	-0.0189284\\
0.742	-0.0189284\\
0.744	-0.0189285\\
0.746	-0.0189285\\
0.748	-0.0189285\\
0.75	-0.0189285\\
0.752	-0.0189285\\
0.754	-0.0189285\\
0.756	-0.0189286\\
0.758	-0.0189286\\
0.76	-0.0189286\\
0.762	-0.0189286\\
0.764	-0.0189286\\
0.766	-0.0189286\\
0.768	-0.0189286\\
0.77	-0.0189287\\
0.772	-0.0189287\\
0.774	-0.0189287\\
0.776	-0.0189287\\
0.778	-0.0189287\\
0.78	-0.0189287\\
0.782	-0.0189287\\
0.784	-0.0189287\\
0.786	-0.0189287\\
0.788	-0.0189287\\
0.79	-0.0189288\\
0.792	-0.0189288\\
0.794	-0.0189288\\
0.796	-0.0189288\\
0.798	-0.0189288\\
0.8	-0.0189288\\
0.802	-0.0189288\\
0.804	-0.0189288\\
0.806	-0.0189288\\
0.808	-0.0189288\\
0.81	-0.0189288\\
0.812	-0.0189288\\
0.814	-0.0189288\\
0.816	-0.0189288\\
0.818	-0.0189289\\
0.82	-0.0189289\\
0.822	-0.0189289\\
0.824	-0.0189289\\
0.826	-0.0189289\\
0.828	-0.0189289\\
0.83	-0.0189289\\
0.832	-0.0189289\\
0.834	-0.0189289\\
0.836	-0.0189289\\
0.838	-0.0189289\\
0.84	-0.0189289\\
0.842	-0.0189289\\
0.844	-0.0189289\\
0.846	-0.0189289\\
0.848	-0.0189289\\
0.85	-0.0189289\\
0.852	-0.0189289\\
0.854	-0.0189289\\
0.856	-0.0189289\\
0.858	-0.0189289\\
0.86	-0.0189289\\
0.862	-0.0189289\\
0.864	-0.0189289\\
0.866	-0.0189289\\
0.868	-0.0189289\\
0.87	-0.018929\\
0.872	-0.018929\\
0.874	-0.018929\\
0.876	-0.018929\\
0.878	-0.018929\\
0.88	-0.018929\\
0.882	-0.018929\\
0.884	-0.018929\\
0.886	-0.018929\\
0.888	-0.018929\\
0.89	-0.018929\\
0.892	-0.018929\\
0.894	-0.018929\\
0.896	-0.018929\\
0.898	-0.018929\\
0.9	-0.018929\\
0.902	-0.018929\\
0.904	-0.018929\\
0.906	-0.018929\\
0.908	-0.018929\\
0.91	-0.018929\\
0.912	-0.018929\\
0.914	-0.018929\\
0.916	-0.018929\\
0.918	-0.018929\\
0.92	-0.018929\\
0.922	-0.018929\\
0.924	-0.018929\\
0.926	-0.018929\\
0.928	-0.018929\\
0.93	-0.018929\\
0.932	-0.018929\\
0.934	-0.018929\\
0.936	-0.018929\\
0.938	-0.018929\\
0.94	-0.018929\\
0.942	-0.018929\\
0.944	-0.018929\\
0.946	-0.018929\\
0.948	-0.018929\\
0.95	-0.018929\\
0.952	-0.018929\\
0.954	-0.018929\\
0.956	-0.0189291\\
0.958	-0.0189291\\
0.96	-0.0189291\\
0.962	-0.0189291\\
0.964	-0.0189291\\
0.966	-0.0189291\\
0.968	-0.0189291\\
0.97	-0.0189291\\
0.972	-0.0189291\\
0.974	-0.0189291\\
0.976	-0.0189291\\
0.978	-0.0189291\\
0.98	-0.0189291\\
0.982	-0.0189291\\
0.984	-0.0189291\\
0.986	-0.0189291\\
0.988	-0.0189291\\
0.99	-0.0189291\\
0.992	-0.0189291\\
0.994	-0.0189291\\
0.996	-0.0189291\\
0.998	-0.0189291\\
1	-0.0189291\\
1.002	-0.0189291\\
1.004	-0.0189291\\
1.006	-0.0189291\\
1.008	-0.0189291\\
1.01	-0.0189291\\
1.012	-0.0189291\\
1.014	-0.0189291\\
1.016	-0.0189291\\
1.018	-0.0189291\\
1.02	-0.0189291\\
1.022	-0.0189291\\
1.024	-0.0189291\\
1.026	-0.0189291\\
1.028	-0.0189291\\
1.03	-0.0189291\\
1.032	-0.0189291\\
1.034	-0.0189291\\
1.036	-0.0189291\\
1.038	-0.0189291\\
1.04	-0.0189291\\
1.042	-0.0189291\\
1.044	-0.0189291\\
1.046	-0.0189291\\
1.048	-0.0189291\\
1.05	-0.0189291\\
1.052	-0.0189291\\
1.054	-0.0189291\\
1.056	-0.0189291\\
1.058	-0.0189291\\
1.06	-0.0189291\\
1.062	-0.0189291\\
1.064	-0.0189291\\
1.066	-0.0189291\\
1.068	-0.0189291\\
1.07	-0.0189291\\
1.072	-0.0189291\\
1.074	-0.0189291\\
1.076	-0.0189291\\
1.078	-0.0189291\\
1.08	-0.0189291\\
1.082	-0.0189291\\
1.084	-0.0189291\\
1.086	-0.0189291\\
1.088	-0.0189291\\
1.09	-0.0189291\\
1.092	-0.0189291\\
1.094	-0.0189291\\
1.096	-0.0189291\\
1.098	-0.0189291\\
1.1	-0.0189291\\
1.102	-0.0189291\\
1.104	-0.0189291\\
1.106	-0.0189291\\
1.108	-0.0189291\\
1.11	-0.0189291\\
1.112	-0.0189291\\
1.114	-0.0189291\\
1.116	-0.0189291\\
1.118	-0.0189291\\
1.12	-0.0189291\\
1.122	-0.0189291\\
1.124	-0.0189291\\
1.126	-0.0189291\\
1.128	-0.0189291\\
1.13	-0.0189291\\
1.132	-0.0189291\\
1.134	-0.0189291\\
1.136	-0.0189291\\
1.138	-0.0189291\\
1.14	-0.0189291\\
1.142	-0.0189291\\
1.144	-0.0189291\\
1.146	-0.0189291\\
1.148	-0.0189291\\
1.15	-0.0189291\\
1.152	-0.0189291\\
1.154	-0.0189291\\
1.156	-0.0189291\\
1.158	-0.0189291\\
1.16	-0.0189291\\
1.162	-0.0189291\\
1.164	-0.0189291\\
1.166	-0.0189291\\
1.168	-0.0189291\\
1.17	-0.0189291\\
1.172	-0.0189291\\
1.174	-0.0189291\\
1.176	-0.0189291\\
1.178	-0.0189291\\
1.18	-0.0189291\\
1.182	-0.0189291\\
1.184	-0.0189291\\
1.186	-0.0189291\\
1.188	-0.0189291\\
1.19	-0.0189291\\
1.192	-0.0189291\\
1.194	-0.0189291\\
1.196	-0.0189291\\
1.198	-0.0189291\\
1.2	-0.0189291\\
1.202	-0.0189291\\
1.204	-0.0189291\\
1.206	-0.0189291\\
1.208	-0.0189291\\
1.21	-0.0189291\\
1.212	-0.0189291\\
1.214	-0.0189291\\
1.216	-0.0189291\\
1.218	-0.0189291\\
1.22	-0.0189291\\
1.222	-0.0189291\\
1.224	-0.0189291\\
1.226	-0.0189291\\
1.228	-0.0189291\\
1.23	-0.0189291\\
1.232	-0.0189291\\
1.234	-0.0189291\\
1.236	-0.0189291\\
1.238	-0.0189291\\
1.24	-0.0189291\\
1.242	-0.0189291\\
1.244	-0.0189291\\
1.246	-0.0189291\\
1.248	-0.0189291\\
1.25	-0.0189291\\
1.252	-0.0189291\\
1.254	-0.0189291\\
1.256	-0.0189291\\
1.258	-0.0189292\\
1.26	-0.0189292\\
1.262	-0.0189292\\
1.264	-0.0189292\\
1.266	-0.0189292\\
1.268	-0.0189292\\
1.27	-0.0189292\\
1.272	-0.0189292\\
1.274	-0.0189292\\
1.276	-0.0189292\\
1.278	-0.0189292\\
1.28	-0.0189292\\
1.282	-0.0189292\\
1.284	-0.0189292\\
1.286	-0.0189292\\
1.288	-0.0189292\\
1.29	-0.0189292\\
1.292	-0.0189292\\
1.294	-0.0189292\\
1.296	-0.0189292\\
1.298	-0.0189292\\
1.3	-0.0189292\\
1.302	-0.0189292\\
1.304	-0.0189292\\
1.306	-0.0189292\\
1.308	-0.0189292\\
1.31	-0.0189292\\
1.312	-0.0189292\\
1.314	-0.0189292\\
1.316	-0.0189292\\
1.318	-0.0189292\\
1.32	-0.0189292\\
1.322	-0.0189292\\
1.324	-0.0189292\\
1.326	-0.0189292\\
1.328	-0.0189292\\
1.33	-0.0189292\\
1.332	-0.0189292\\
1.334	-0.0189292\\
1.336	-0.0189292\\
1.338	-0.0189292\\
1.34	-0.0189292\\
1.342	-0.0189292\\
1.344	-0.0189292\\
1.346	-0.0189292\\
1.348	-0.0189292\\
1.35	-0.0189292\\
1.352	-0.0189292\\
1.354	-0.0189292\\
1.356	-0.0189292\\
1.358	-0.0189292\\
1.36	-0.0189292\\
1.362	-0.0189292\\
1.364	-0.0189292\\
1.366	-0.0189292\\
1.368	-0.0189292\\
1.37	-0.0189292\\
1.372	-0.0189292\\
1.374	-0.0189292\\
1.376	-0.0189292\\
1.378	-0.0189292\\
1.38	-0.0189292\\
1.382	-0.0189292\\
1.384	-0.0189292\\
1.386	-0.0189292\\
1.388	-0.0189292\\
1.39	-0.0189292\\
1.392	-0.0189292\\
1.394	-0.0189292\\
1.396	-0.0189292\\
1.398	-0.0189292\\
1.4	-0.0189292\\
1.402	-0.0189292\\
1.404	-0.0189292\\
1.406	-0.0189292\\
1.408	-0.0189292\\
1.41	-0.0189292\\
1.412	-0.0189292\\
1.414	-0.0189292\\
1.416	-0.0189292\\
1.418	-0.0189292\\
1.42	-0.0189292\\
1.422	-0.0189292\\
1.424	-0.0189292\\
1.426	-0.0189292\\
1.428	-0.0189292\\
1.43	-0.0189292\\
1.432	-0.0189292\\
1.434	-0.0189292\\
1.436	-0.0189292\\
1.438	-0.0189292\\
1.44	-0.0189292\\
1.442	-0.0189292\\
1.444	-0.0189292\\
1.446	-0.0189292\\
1.448	-0.0189292\\
1.45	-0.0189292\\
1.452	-0.0189292\\
1.454	-0.0189292\\
1.456	-0.0189292\\
1.458	-0.0189292\\
1.46	-0.0189292\\
1.462	-0.0189292\\
1.464	-0.0189292\\
1.466	-0.0189292\\
1.468	-0.0189292\\
1.47	-0.0189292\\
1.472	-0.0189292\\
1.474	-0.0189292\\
1.476	-0.0189292\\
1.478	-0.0189292\\
1.48	-0.0189292\\
1.482	-0.0189292\\
1.484	-0.0189292\\
1.486	-0.0189292\\
1.488	-0.0189292\\
1.49	-0.0189292\\
1.492	-0.0189292\\
1.494	-0.0189292\\
1.496	-0.0189292\\
1.498	-0.0189292\\
1.5	-0.0189292\\
1.502	-0.0189292\\
1.504	-0.0189292\\
1.506	-0.0189292\\
1.508	-0.0189292\\
1.51	-0.0189292\\
1.512	-0.0189292\\
1.514	-0.0189292\\
1.516	-0.0189292\\
1.518	-0.0189292\\
1.52	-0.0189292\\
1.522	-0.0189292\\
1.524	-0.0189292\\
1.526	-0.0189292\\
1.528	-0.0189292\\
1.53	-0.0189292\\
1.532	-0.0189292\\
1.534	-0.0189292\\
1.536	-0.0189292\\
1.538	-0.0189292\\
1.54	-0.0189292\\
1.542	-0.0189292\\
1.544	-0.0189292\\
1.546	-0.0189292\\
1.548	-0.0189292\\
1.55	-0.0189292\\
1.552	-0.0189292\\
1.554	-0.0189292\\
1.556	-0.0189292\\
1.558	-0.0189292\\
1.56	-0.0189292\\
1.562	-0.0189292\\
1.564	-0.0189292\\
1.566	-0.0189292\\
1.568	-0.0189292\\
1.57	-0.0189292\\
1.572	-0.0189292\\
1.574	-0.0189292\\
1.576	-0.0189292\\
1.578	-0.0189292\\
1.58	-0.0189292\\
1.582	-0.0189292\\
1.584	-0.0189292\\
1.586	-0.0189292\\
1.588	-0.0189292\\
1.59	-0.0189292\\
1.592	-0.0189292\\
1.594	-0.0189292\\
1.596	-0.0189292\\
1.598	-0.0189292\\
1.6	-0.0189292\\
1.602	-0.0189292\\
1.604	-0.0189292\\
1.606	-0.0189292\\
1.608	-0.0189292\\
1.61	-0.0189292\\
1.612	-0.0189292\\
1.614	-0.0189292\\
1.616	-0.0189292\\
1.618	-0.0189292\\
1.62	-0.0189292\\
1.622	-0.0189292\\
1.624	-0.0189292\\
1.626	-0.0189292\\
1.628	-0.0189292\\
1.63	-0.0189292\\
1.632	-0.0189292\\
1.634	-0.0189292\\
1.636	-0.0189292\\
1.638	-0.0189292\\
1.64	-0.0189292\\
1.642	-0.0189292\\
1.644	-0.0189292\\
1.646	-0.0189292\\
1.648	-0.0189292\\
1.65	-0.0189292\\
1.652	-0.0189292\\
1.654	-0.0189292\\
1.656	-0.0189292\\
1.658	-0.0189292\\
1.66	-0.0189292\\
1.662	-0.0189292\\
1.664	-0.0189292\\
1.666	-0.0189292\\
1.668	-0.0189292\\
1.67	-0.0189292\\
1.672	-0.0189292\\
1.674	-0.0189292\\
1.676	-0.0189292\\
1.678	-0.0189292\\
1.68	-0.0189292\\
1.682	-0.0189292\\
1.684	-0.0189292\\
1.686	-0.0189292\\
1.688	-0.0189292\\
1.69	-0.0189292\\
1.692	-0.0189292\\
1.694	-0.0189292\\
1.696	-0.0189292\\
1.698	-0.0189292\\
1.7	-0.0189292\\
1.702	-0.0189292\\
1.704	-0.0189292\\
1.706	-0.0189292\\
1.708	-0.0189292\\
1.71	-0.0189292\\
1.712	-0.0189292\\
1.714	-0.0189292\\
1.716	-0.0189292\\
1.718	-0.0189292\\
1.72	-0.0189292\\
1.722	-0.0189292\\
1.724	-0.0189292\\
1.726	-0.0189292\\
1.728	-0.0189292\\
1.73	-0.0189292\\
1.732	-0.0189292\\
1.734	-0.0189292\\
1.736	-0.0189292\\
1.738	-0.0189292\\
1.74	-0.0189292\\
1.742	-0.0189292\\
1.744	-0.0189292\\
1.746	-0.0189292\\
1.748	-0.0189292\\
1.75	-0.0189292\\
1.752	-0.0189292\\
1.754	-0.0189292\\
1.756	-0.0189292\\
1.758	-0.0189292\\
1.76	-0.0189292\\
1.762	-0.0189292\\
1.764	-0.0189292\\
1.766	-0.0189292\\
1.768	-0.0189292\\
1.77	-0.0189292\\
1.772	-0.0189292\\
1.774	-0.0189292\\
1.776	-0.0189292\\
1.778	-0.0189292\\
1.78	-0.0189292\\
1.782	-0.0189292\\
1.784	-0.0189292\\
1.786	-0.0189292\\
1.788	-0.0189292\\
1.79	-0.0189292\\
1.792	-0.0189292\\
1.794	-0.0189292\\
1.796	-0.0189292\\
1.798	-0.0189292\\
1.8	-0.0189292\\
1.802	-0.0189292\\
1.804	-0.0189292\\
1.806	-0.0189292\\
1.808	-0.0189292\\
1.81	-0.0189292\\
1.812	-0.0189292\\
1.814	-0.0189292\\
1.816	-0.0189292\\
1.818	-0.0189292\\
1.82	-0.0189292\\
1.822	-0.0189292\\
1.824	-0.0189292\\
1.826	-0.0189292\\
1.828	-0.0189292\\
1.83	-0.0189292\\
1.832	-0.0189292\\
1.834	-0.0189292\\
1.836	-0.0189292\\
1.838	-0.0189292\\
1.84	-0.0189292\\
1.842	-0.0189292\\
1.844	-0.0189292\\
1.846	-0.0189292\\
1.848	-0.0189292\\
1.85	-0.0189292\\
1.852	-0.0189292\\
1.854	-0.0189292\\
1.856	-0.0189292\\
1.858	-0.0189292\\
1.86	-0.0189292\\
1.862	-0.0189292\\
1.864	-0.0189292\\
1.866	-0.0189292\\
1.868	-0.0189292\\
1.87	-0.0189292\\
1.872	-0.0189292\\
1.874	-0.0189292\\
1.876	-0.0189292\\
1.878	-0.0189292\\
1.88	-0.0189292\\
1.882	-0.0189292\\
1.884	-0.0189292\\
1.886	-0.0189292\\
1.888	-0.0189292\\
1.89	-0.0189292\\
1.892	-0.0189292\\
1.894	-0.0189292\\
1.896	-0.0189292\\
1.898	-0.0189292\\
1.9	-0.0189292\\
1.902	-0.0189292\\
1.904	-0.0189292\\
1.906	-0.0189292\\
1.908	-0.0189292\\
1.91	-0.0189292\\
1.912	-0.0189292\\
1.914	-0.0189292\\
1.916	-0.0189292\\
1.918	-0.0189292\\
1.92	-0.0189292\\
1.922	-0.0189292\\
1.924	-0.0189292\\
1.926	-0.0189292\\
1.928	-0.0189292\\
1.93	-0.0189292\\
1.932	-0.0189292\\
1.934	-0.0189292\\
1.936	-0.0189292\\
1.938	-0.0189292\\
1.94	-0.0189292\\
1.942	-0.0189292\\
1.944	-0.0189292\\
1.946	-0.0189292\\
1.948	-0.0189292\\
1.95	-0.0189292\\
1.952	-0.0189292\\
1.954	-0.0189292\\
1.956	-0.0189292\\
1.958	-0.0189292\\
1.96	-0.0189292\\
1.962	-0.0189292\\
1.964	-0.0189292\\
1.966	-0.0189292\\
1.968	-0.0189292\\
1.97	-0.0189292\\
1.972	-0.0189292\\
1.974	-0.0189292\\
1.976	-0.0189292\\
1.978	-0.0189292\\
1.98	-0.0189292\\
1.982	-0.0189292\\
1.984	-0.0189292\\
1.986	-0.0189292\\
1.988	-0.0189292\\
1.99	-0.0189292\\
1.992	-0.0189292\\
1.994	-0.0189292\\
1.996	-0.0189292\\
1.998	-0.0189292\\
2	-0.0189292\\
};
\draw[dashed, line width=0.7pt, draw=black] (axis cs:0,-0.0206901) rectangle (axis cs:0.5,-0.017507);
\end{axis}

\begin{axis}[%
width={\figscale*1.528in},
height={\figscale*0.46in},
at={(\figscale*4.752in,\figscale*1.69in)},
scale only axis,
separate axis lines,
every outer x axis line/.append style={black},
every x tick label/.append style={font=\figticfont},
every x tick/.append style={black},
xmin=-0.025,
xmax=2,
xtick={0,0.5,1,1.5,2},
xticklabels={\empty},
every outer y axis line/.append style={black},
every y tick label/.append style={font=\figticfont},
every y tick/.append style={black},
ymin=-0.0247712,
ymax=0,
ytick={-0.02,0},
axis background/.style={fill=white},
title style={font=\figcornerfont},
title={Texas2000},
xmajorgrids,
ymajorgrids,
grid style={white!55!black, opacity=0.3},
ylabel style={rotate=-90, at={(axis description cs:-0.145,0.5)}, anchor=east}
]
\addplot [color=mycolor1, line width=1.4pt, forget plot]
  table[row sep=crcr]{%
-0.024	0\\
-0.022	0\\
-0.02	0\\
-0.018	0\\
-0.016	0\\
-0.014	0\\
-0.012	0\\
-0.01	0\\
-0.008	0\\
-0.006	0\\
-0.004	0\\
-0.002	0\\
0	0\\
0.002	-0.00184649\\
0.004	-0.00342632\\
0.006	-0.00478673\\
0.008	-0.00596703\\
0.01	-0.00699977\\
0.012	-0.00791183\\
0.014	-0.00872528\\
0.016	-0.00945819\\
0.018	-0.0101252\\
0.02	-0.0107383\\
0.022	-0.0113069\\
0.024	-0.0118387\\
0.026	-0.0123396\\
0.028	-0.0128143\\
0.03	-0.0132664\\
0.032	-0.0136986\\
0.034	-0.0141129\\
0.036	-0.0145108\\
0.038	-0.0148935\\
0.04	-0.0152615\\
0.042	-0.0156154\\
0.044	-0.0159555\\
0.046	-0.016282\\
0.048	-0.016595\\
0.05	-0.0168948\\
0.052	-0.0171813\\
0.054	-0.0174547\\
0.056	-0.0177152\\
0.058	-0.017963\\
0.06	-0.0181983\\
0.062	-0.0184213\\
0.064	-0.0186325\\
0.066	-0.0188322\\
0.068	-0.0190207\\
0.07	-0.0191986\\
0.072	-0.0193662\\
0.074	-0.019524\\
0.076	-0.0196726\\
0.078	-0.0198123\\
0.08	-0.0199438\\
0.082	-0.0200675\\
0.084	-0.0201838\\
0.086	-0.0202933\\
0.088	-0.0203963\\
0.09	-0.0204935\\
0.092	-0.0205851\\
0.094	-0.0206715\\
0.096	-0.0207533\\
0.098	-0.0208307\\
0.1	-0.020904\\
0.102	-0.0209737\\
0.104	-0.02104\\
0.106	-0.0211031\\
0.108	-0.0211633\\
0.11	-0.0212208\\
0.112	-0.0212759\\
0.114	-0.0213286\\
0.116	-0.0213791\\
0.118	-0.0214276\\
0.12	-0.0214742\\
0.122	-0.0215189\\
0.124	-0.0215618\\
0.126	-0.0216029\\
0.128	-0.0216422\\
0.13	-0.0216798\\
0.132	-0.0217157\\
0.134	-0.0217498\\
0.136	-0.0217822\\
0.138	-0.0218127\\
0.14	-0.0218414\\
0.142	-0.0218682\\
0.144	-0.021893\\
0.146	-0.0219159\\
0.148	-0.0219368\\
0.15	-0.0219556\\
0.152	-0.0219723\\
0.154	-0.021987\\
0.156	-0.0219995\\
0.158	-0.02201\\
0.16	-0.0220184\\
0.162	-0.0220247\\
0.164	-0.022029\\
0.166	-0.0220314\\
0.168	-0.0220318\\
0.17	-0.0220304\\
0.172	-0.0220273\\
0.174	-0.0220225\\
0.176	-0.0220162\\
0.178	-0.0220086\\
0.18	-0.0219997\\
0.182	-0.0219896\\
0.184	-0.0219786\\
0.186	-0.0219668\\
0.188	-0.0219544\\
0.19	-0.0219415\\
0.192	-0.0219283\\
0.194	-0.0219149\\
0.196	-0.0219017\\
0.198	-0.0218886\\
0.2	-0.0218759\\
0.202	-0.0218638\\
0.204	-0.0218524\\
0.206	-0.0218418\\
0.208	-0.0218323\\
0.21	-0.0218239\\
0.212	-0.0218168\\
0.214	-0.0218111\\
0.216	-0.0218069\\
0.218	-0.0218044\\
0.22	-0.0218035\\
0.222	-0.0218044\\
0.224	-0.0218072\\
0.226	-0.0218119\\
0.228	-0.0218185\\
0.23	-0.0218271\\
0.232	-0.0218377\\
0.234	-0.0218503\\
0.236	-0.021865\\
0.238	-0.0218816\\
0.24	-0.0219002\\
0.242	-0.0219208\\
0.244	-0.0219432\\
0.246	-0.0219676\\
0.248	-0.0219937\\
0.25	-0.0220216\\
0.252	-0.0220512\\
0.254	-0.0220823\\
0.256	-0.022115\\
0.258	-0.0221492\\
0.26	-0.0221847\\
0.262	-0.0222214\\
0.264	-0.0222593\\
0.266	-0.0222983\\
0.268	-0.0223382\\
0.27	-0.022379\\
0.272	-0.0224206\\
0.274	-0.0224628\\
0.276	-0.0225056\\
0.278	-0.0225489\\
0.28	-0.0225926\\
0.282	-0.0226365\\
0.284	-0.0226806\\
0.286	-0.0227248\\
0.288	-0.022769\\
0.29	-0.0228132\\
0.292	-0.0228572\\
0.294	-0.0229009\\
0.296	-0.0229443\\
0.298	-0.0229873\\
0.3	-0.0230299\\
0.302	-0.0230719\\
0.304	-0.0231134\\
0.306	-0.0231542\\
0.308	-0.0231943\\
0.31	-0.0232337\\
0.312	-0.0232723\\
0.314	-0.02331\\
0.316	-0.0233469\\
0.318	-0.0233828\\
0.32	-0.0234179\\
0.322	-0.0234519\\
0.324	-0.023485\\
0.326	-0.023517\\
0.328	-0.023548\\
0.33	-0.023578\\
0.332	-0.0236069\\
0.334	-0.0236347\\
0.336	-0.0236615\\
0.338	-0.0236871\\
0.34	-0.0237117\\
0.342	-0.0237352\\
0.344	-0.0237577\\
0.346	-0.023779\\
0.348	-0.0237993\\
0.35	-0.0238186\\
0.352	-0.0238368\\
0.354	-0.023854\\
0.356	-0.0238701\\
0.358	-0.0238853\\
0.36	-0.0238996\\
0.362	-0.0239129\\
0.364	-0.0239252\\
0.366	-0.0239367\\
0.368	-0.0239474\\
0.37	-0.0239572\\
0.372	-0.0239662\\
0.374	-0.0239744\\
0.376	-0.0239819\\
0.378	-0.0239887\\
0.38	-0.0239948\\
0.382	-0.0240004\\
0.384	-0.0240053\\
0.386	-0.0240097\\
0.388	-0.0240136\\
0.39	-0.024017\\
0.392	-0.02402\\
0.394	-0.0240226\\
0.396	-0.0240248\\
0.398	-0.0240268\\
0.4	-0.0240285\\
0.402	-0.02403\\
0.404	-0.0240313\\
0.406	-0.0240325\\
0.408	-0.0240336\\
0.41	-0.0240347\\
0.412	-0.0240357\\
0.414	-0.0240368\\
0.416	-0.0240379\\
0.418	-0.0240391\\
0.42	-0.0240405\\
0.422	-0.024042\\
0.424	-0.0240437\\
0.426	-0.0240457\\
0.428	-0.0240479\\
0.43	-0.0240504\\
0.432	-0.0240532\\
0.434	-0.0240563\\
0.436	-0.0240598\\
0.438	-0.0240637\\
0.44	-0.0240679\\
0.442	-0.0240725\\
0.444	-0.0240775\\
0.446	-0.024083\\
0.448	-0.0240889\\
0.45	-0.0240952\\
0.452	-0.0241019\\
0.454	-0.024109\\
0.456	-0.0241166\\
0.458	-0.0241246\\
0.46	-0.024133\\
0.462	-0.0241418\\
0.464	-0.024151\\
0.466	-0.0241606\\
0.468	-0.0241705\\
0.47	-0.0241807\\
0.472	-0.0241913\\
0.474	-0.0242022\\
0.476	-0.0242133\\
0.478	-0.0242247\\
0.48	-0.0242363\\
0.482	-0.0242481\\
0.484	-0.0242601\\
0.486	-0.0242722\\
0.488	-0.0242844\\
0.49	-0.0242966\\
0.492	-0.0243089\\
0.494	-0.0243213\\
0.496	-0.0243335\\
0.498	-0.0243458\\
0.5	-0.0243579\\
0.502	-0.0243699\\
0.504	-0.0243818\\
0.506	-0.0243935\\
0.508	-0.0244049\\
0.51	-0.0244162\\
0.512	-0.0244271\\
0.514	-0.0244377\\
0.516	-0.024448\\
0.518	-0.024458\\
0.52	-0.0244675\\
0.522	-0.0244767\\
0.524	-0.0244854\\
0.526	-0.0244937\\
0.528	-0.0245015\\
0.53	-0.0245088\\
0.532	-0.0245157\\
0.534	-0.024522\\
0.536	-0.0245278\\
0.538	-0.0245331\\
0.54	-0.0245379\\
0.542	-0.0245421\\
0.544	-0.0245457\\
0.546	-0.0245489\\
0.548	-0.0245514\\
0.55	-0.0245535\\
0.552	-0.024555\\
0.554	-0.024556\\
0.556	-0.0245565\\
0.558	-0.0245564\\
0.56	-0.0245559\\
0.562	-0.0245549\\
0.564	-0.0245534\\
0.566	-0.0245515\\
0.568	-0.0245492\\
0.57	-0.0245465\\
0.572	-0.0245433\\
0.574	-0.0245399\\
0.576	-0.024536\\
0.578	-0.0245319\\
0.58	-0.0245275\\
0.582	-0.0245228\\
0.584	-0.024518\\
0.586	-0.0245129\\
0.588	-0.0245076\\
0.59	-0.0245022\\
0.592	-0.0244967\\
0.594	-0.024491\\
0.596	-0.0244854\\
0.598	-0.0244797\\
0.6	-0.0244739\\
0.602	-0.0244683\\
0.604	-0.0244626\\
0.606	-0.0244571\\
0.608	-0.0244516\\
0.61	-0.0244463\\
0.612	-0.0244411\\
0.614	-0.0244362\\
0.616	-0.0244314\\
0.618	-0.0244268\\
0.62	-0.0244225\\
0.622	-0.0244184\\
0.624	-0.0244146\\
0.626	-0.0244111\\
0.628	-0.0244079\\
0.63	-0.0244051\\
0.632	-0.0244025\\
0.634	-0.0244003\\
0.636	-0.0243985\\
0.638	-0.0243969\\
0.64	-0.0243958\\
0.642	-0.024395\\
0.644	-0.0243946\\
0.646	-0.0243946\\
0.648	-0.0243949\\
0.65	-0.0243956\\
0.652	-0.0243967\\
0.654	-0.0243981\\
0.656	-0.0243998\\
0.658	-0.0244019\\
0.66	-0.0244044\\
0.662	-0.0244072\\
0.664	-0.0244103\\
0.666	-0.0244137\\
0.668	-0.0244174\\
0.67	-0.0244213\\
0.672	-0.0244256\\
0.674	-0.0244301\\
0.676	-0.0244348\\
0.678	-0.0244398\\
0.68	-0.0244449\\
0.682	-0.0244503\\
0.684	-0.0244558\\
0.686	-0.0244615\\
0.688	-0.0244673\\
0.69	-0.0244733\\
0.692	-0.0244793\\
0.694	-0.0244855\\
0.696	-0.0244917\\
0.698	-0.0244979\\
0.7	-0.0245042\\
0.702	-0.0245105\\
0.704	-0.0245169\\
0.706	-0.0245232\\
0.708	-0.0245295\\
0.71	-0.0245357\\
0.712	-0.0245419\\
0.714	-0.024548\\
0.716	-0.024554\\
0.718	-0.02456\\
0.72	-0.0245658\\
0.722	-0.0245715\\
0.724	-0.0245771\\
0.726	-0.0245825\\
0.728	-0.0245878\\
0.73	-0.024593\\
0.732	-0.024598\\
0.734	-0.0246028\\
0.736	-0.0246074\\
0.738	-0.0246118\\
0.74	-0.0246161\\
0.742	-0.0246202\\
0.744	-0.024624\\
0.746	-0.0246277\\
0.748	-0.0246312\\
0.75	-0.0246345\\
0.752	-0.0246375\\
0.754	-0.0246404\\
0.756	-0.0246431\\
0.758	-0.0246455\\
0.76	-0.0246478\\
0.762	-0.0246499\\
0.764	-0.0246518\\
0.766	-0.0246534\\
0.768	-0.0246549\\
0.77	-0.0246562\\
0.772	-0.0246574\\
0.774	-0.0246583\\
0.776	-0.0246591\\
0.778	-0.0246597\\
0.78	-0.0246602\\
0.782	-0.0246605\\
0.784	-0.0246606\\
0.786	-0.0246606\\
0.788	-0.0246605\\
0.79	-0.0246602\\
0.792	-0.0246599\\
0.794	-0.0246594\\
0.796	-0.0246587\\
0.798	-0.024658\\
0.8	-0.0246572\\
0.802	-0.0246563\\
0.804	-0.0246553\\
0.806	-0.0246543\\
0.808	-0.0246532\\
0.81	-0.024652\\
0.812	-0.0246507\\
0.814	-0.0246495\\
0.816	-0.0246481\\
0.818	-0.0246468\\
0.82	-0.0246454\\
0.822	-0.024644\\
0.824	-0.0246426\\
0.826	-0.0246412\\
0.828	-0.0246398\\
0.83	-0.0246384\\
0.832	-0.024637\\
0.834	-0.0246356\\
0.836	-0.0246342\\
0.838	-0.0246329\\
0.84	-0.0246316\\
0.842	-0.0246304\\
0.844	-0.0246292\\
0.846	-0.0246281\\
0.848	-0.024627\\
0.85	-0.024626\\
0.852	-0.024625\\
0.854	-0.0246241\\
0.856	-0.0246233\\
0.858	-0.0246226\\
0.86	-0.0246219\\
0.862	-0.0246213\\
0.864	-0.0246208\\
0.866	-0.0246204\\
0.868	-0.0246201\\
0.87	-0.0246199\\
0.872	-0.0246198\\
0.874	-0.0246198\\
0.876	-0.0246199\\
0.878	-0.02462\\
0.88	-0.0246203\\
0.882	-0.0246207\\
0.884	-0.0246212\\
0.886	-0.0246218\\
0.888	-0.0246225\\
0.89	-0.0246232\\
0.892	-0.0246241\\
0.894	-0.0246251\\
0.896	-0.0246262\\
0.898	-0.0246274\\
0.9	-0.0246287\\
0.902	-0.02463\\
0.904	-0.0246315\\
0.906	-0.024633\\
0.908	-0.0246347\\
0.91	-0.0246364\\
0.912	-0.0246382\\
0.914	-0.0246401\\
0.916	-0.024642\\
0.918	-0.024644\\
0.92	-0.0246461\\
0.922	-0.0246483\\
0.924	-0.0246505\\
0.926	-0.0246527\\
0.928	-0.024655\\
0.93	-0.0246574\\
0.932	-0.0246598\\
0.934	-0.0246622\\
0.936	-0.0246647\\
0.938	-0.0246672\\
0.94	-0.0246697\\
0.942	-0.0246722\\
0.944	-0.0246748\\
0.946	-0.0246773\\
0.948	-0.0246799\\
0.95	-0.0246824\\
0.952	-0.024685\\
0.954	-0.0246875\\
0.956	-0.02469\\
0.958	-0.0246925\\
0.96	-0.024695\\
0.962	-0.0246975\\
0.964	-0.0246999\\
0.966	-0.0247023\\
0.968	-0.0247046\\
0.97	-0.0247069\\
0.972	-0.0247092\\
0.974	-0.0247114\\
0.976	-0.0247135\\
0.978	-0.0247156\\
0.98	-0.0247177\\
0.982	-0.0247196\\
0.984	-0.0247215\\
0.986	-0.0247234\\
0.988	-0.0247252\\
0.99	-0.0247269\\
0.992	-0.0247285\\
0.994	-0.0247301\\
0.996	-0.0247316\\
0.998	-0.024733\\
1	-0.0247343\\
1.002	-0.0247356\\
1.004	-0.0247368\\
1.006	-0.0247379\\
1.008	-0.024739\\
1.01	-0.0247399\\
1.012	-0.0247408\\
1.014	-0.0247417\\
1.016	-0.0247424\\
1.018	-0.0247431\\
1.02	-0.0247437\\
1.022	-0.0247442\\
1.024	-0.0247447\\
1.026	-0.0247451\\
1.028	-0.0247454\\
1.03	-0.0247457\\
1.032	-0.0247459\\
1.034	-0.0247461\\
1.036	-0.0247462\\
1.038	-0.0247463\\
1.04	-0.0247463\\
1.042	-0.0247462\\
1.044	-0.0247461\\
1.046	-0.024746\\
1.048	-0.0247458\\
1.05	-0.0247456\\
1.052	-0.0247454\\
1.054	-0.0247451\\
1.056	-0.0247448\\
1.058	-0.0247445\\
1.06	-0.0247441\\
1.062	-0.0247437\\
1.064	-0.0247434\\
1.066	-0.0247429\\
1.068	-0.0247425\\
1.07	-0.0247421\\
1.072	-0.0247417\\
1.074	-0.0247412\\
1.076	-0.0247408\\
1.078	-0.0247403\\
1.08	-0.0247399\\
1.082	-0.0247395\\
1.084	-0.024739\\
1.086	-0.0247386\\
1.088	-0.0247382\\
1.09	-0.0247378\\
1.092	-0.0247374\\
1.094	-0.0247371\\
1.096	-0.0247367\\
1.098	-0.0247364\\
1.1	-0.0247361\\
1.102	-0.0247358\\
1.104	-0.0247355\\
1.106	-0.0247352\\
1.108	-0.024735\\
1.11	-0.0247348\\
1.112	-0.0247346\\
1.114	-0.0247345\\
1.116	-0.0247343\\
1.118	-0.0247342\\
1.12	-0.0247341\\
1.122	-0.0247341\\
1.124	-0.024734\\
1.126	-0.024734\\
1.128	-0.0247341\\
1.13	-0.0247341\\
1.132	-0.0247342\\
1.134	-0.0247343\\
1.136	-0.0247344\\
1.138	-0.0247345\\
1.14	-0.0247347\\
1.142	-0.0247349\\
1.144	-0.0247351\\
1.146	-0.0247353\\
1.148	-0.0247356\\
1.15	-0.0247358\\
1.152	-0.0247361\\
1.154	-0.0247364\\
1.156	-0.0247367\\
1.158	-0.0247371\\
1.16	-0.0247374\\
1.162	-0.0247378\\
1.164	-0.0247382\\
1.166	-0.0247385\\
1.168	-0.0247389\\
1.17	-0.0247393\\
1.172	-0.0247397\\
1.174	-0.0247401\\
1.176	-0.0247406\\
1.178	-0.024741\\
1.18	-0.0247414\\
1.182	-0.0247418\\
1.184	-0.0247422\\
1.186	-0.0247427\\
1.188	-0.0247431\\
1.19	-0.0247435\\
1.192	-0.0247439\\
1.194	-0.0247444\\
1.196	-0.0247448\\
1.198	-0.0247452\\
1.2	-0.0247456\\
1.202	-0.024746\\
1.204	-0.0247463\\
1.206	-0.0247467\\
1.208	-0.0247471\\
1.21	-0.0247474\\
1.212	-0.0247478\\
1.214	-0.0247481\\
1.216	-0.0247484\\
1.218	-0.0247488\\
1.22	-0.0247491\\
1.222	-0.0247493\\
1.224	-0.0247496\\
1.226	-0.0247499\\
1.228	-0.0247501\\
1.23	-0.0247504\\
1.232	-0.0247506\\
1.234	-0.0247508\\
1.236	-0.024751\\
1.238	-0.0247512\\
1.24	-0.0247513\\
1.242	-0.0247515\\
1.244	-0.0247516\\
1.246	-0.0247518\\
1.248	-0.0247519\\
1.25	-0.024752\\
1.252	-0.0247521\\
1.254	-0.0247521\\
1.256	-0.0247522\\
1.258	-0.0247523\\
1.26	-0.0247523\\
1.262	-0.0247523\\
1.264	-0.0247524\\
1.266	-0.0247524\\
1.268	-0.0247524\\
1.27	-0.0247524\\
1.272	-0.0247524\\
1.274	-0.0247523\\
1.276	-0.0247523\\
1.278	-0.0247523\\
1.28	-0.0247522\\
1.282	-0.0247522\\
1.284	-0.0247521\\
1.286	-0.0247521\\
1.288	-0.024752\\
1.29	-0.0247519\\
1.292	-0.0247519\\
1.294	-0.0247518\\
1.296	-0.0247517\\
1.298	-0.0247517\\
1.3	-0.0247516\\
1.302	-0.0247515\\
1.304	-0.0247514\\
1.306	-0.0247513\\
1.308	-0.0247513\\
1.31	-0.0247512\\
1.312	-0.0247511\\
1.314	-0.024751\\
1.316	-0.024751\\
1.318	-0.0247509\\
1.32	-0.0247509\\
1.322	-0.0247508\\
1.324	-0.0247508\\
1.326	-0.0247507\\
1.328	-0.0247507\\
1.33	-0.0247506\\
1.332	-0.0247506\\
1.334	-0.0247506\\
1.336	-0.0247506\\
1.338	-0.0247505\\
1.34	-0.0247505\\
1.342	-0.0247505\\
1.344	-0.0247505\\
1.346	-0.0247506\\
1.348	-0.0247506\\
1.35	-0.0247506\\
1.352	-0.0247507\\
1.354	-0.0247507\\
1.356	-0.0247508\\
1.358	-0.0247508\\
1.36	-0.0247509\\
1.362	-0.024751\\
1.364	-0.024751\\
1.366	-0.0247511\\
1.368	-0.0247512\\
1.37	-0.0247513\\
1.372	-0.0247514\\
1.374	-0.0247516\\
1.376	-0.0247517\\
1.378	-0.0247518\\
1.38	-0.024752\\
1.382	-0.0247521\\
1.384	-0.0247523\\
1.386	-0.0247524\\
1.388	-0.0247526\\
1.39	-0.0247528\\
1.392	-0.0247529\\
1.394	-0.0247531\\
1.396	-0.0247533\\
1.398	-0.0247535\\
1.4	-0.0247537\\
1.402	-0.0247539\\
1.404	-0.0247541\\
1.406	-0.0247543\\
1.408	-0.0247545\\
1.41	-0.0247547\\
1.412	-0.0247549\\
1.414	-0.0247551\\
1.416	-0.0247553\\
1.418	-0.0247555\\
1.42	-0.0247558\\
1.422	-0.024756\\
1.424	-0.0247562\\
1.426	-0.0247564\\
1.428	-0.0247566\\
1.43	-0.0247568\\
1.432	-0.0247571\\
1.434	-0.0247573\\
1.436	-0.0247575\\
1.438	-0.0247577\\
1.44	-0.0247579\\
1.442	-0.0247581\\
1.444	-0.0247583\\
1.446	-0.0247585\\
1.448	-0.0247587\\
1.45	-0.0247589\\
1.452	-0.0247591\\
1.454	-0.0247593\\
1.456	-0.0247595\\
1.458	-0.0247597\\
1.46	-0.0247599\\
1.462	-0.0247601\\
1.464	-0.0247603\\
1.466	-0.0247604\\
1.468	-0.0247606\\
1.47	-0.0247608\\
1.472	-0.0247609\\
1.474	-0.0247611\\
1.476	-0.0247612\\
1.478	-0.0247614\\
1.48	-0.0247615\\
1.482	-0.0247617\\
1.484	-0.0247618\\
1.486	-0.024762\\
1.488	-0.0247621\\
1.49	-0.0247622\\
1.492	-0.0247623\\
1.494	-0.0247625\\
1.496	-0.0247626\\
1.498	-0.0247627\\
1.5	-0.0247628\\
1.502	-0.0247629\\
1.504	-0.024763\\
1.506	-0.0247631\\
1.508	-0.0247632\\
1.51	-0.0247633\\
1.512	-0.0247634\\
1.514	-0.0247635\\
1.516	-0.0247636\\
1.518	-0.0247636\\
1.52	-0.0247637\\
1.522	-0.0247638\\
1.524	-0.0247639\\
1.526	-0.024764\\
1.528	-0.024764\\
1.53	-0.0247641\\
1.532	-0.0247642\\
1.534	-0.0247643\\
1.536	-0.0247643\\
1.538	-0.0247644\\
1.54	-0.0247645\\
1.542	-0.0247645\\
1.544	-0.0247646\\
1.546	-0.0247647\\
1.548	-0.0247647\\
1.55	-0.0247648\\
1.552	-0.0247649\\
1.554	-0.024765\\
1.556	-0.024765\\
1.558	-0.0247651\\
1.56	-0.0247652\\
1.562	-0.0247652\\
1.564	-0.0247653\\
1.566	-0.0247654\\
1.568	-0.0247655\\
1.57	-0.0247655\\
1.572	-0.0247656\\
1.574	-0.0247657\\
1.576	-0.0247658\\
1.578	-0.0247658\\
1.58	-0.0247659\\
1.582	-0.024766\\
1.584	-0.0247661\\
1.586	-0.0247662\\
1.588	-0.0247663\\
1.59	-0.0247663\\
1.592	-0.0247664\\
1.594	-0.0247665\\
1.596	-0.0247666\\
1.598	-0.0247667\\
1.6	-0.0247668\\
1.602	-0.0247669\\
1.604	-0.024767\\
1.606	-0.0247671\\
1.608	-0.0247672\\
1.61	-0.0247673\\
1.612	-0.0247674\\
1.614	-0.0247674\\
1.616	-0.0247675\\
1.618	-0.0247676\\
1.62	-0.0247677\\
1.622	-0.0247678\\
1.624	-0.0247679\\
1.626	-0.024768\\
1.628	-0.0247681\\
1.63	-0.0247682\\
1.632	-0.0247683\\
1.634	-0.0247684\\
1.636	-0.0247685\\
1.638	-0.0247686\\
1.64	-0.0247687\\
1.642	-0.0247687\\
1.644	-0.0247688\\
1.646	-0.0247689\\
1.648	-0.024769\\
1.65	-0.0247691\\
1.652	-0.0247692\\
1.654	-0.0247692\\
1.656	-0.0247693\\
1.658	-0.0247694\\
1.66	-0.0247695\\
1.662	-0.0247695\\
1.664	-0.0247696\\
1.666	-0.0247697\\
1.668	-0.0247697\\
1.67	-0.0247698\\
1.672	-0.0247698\\
1.674	-0.0247699\\
1.676	-0.02477\\
1.678	-0.02477\\
1.68	-0.0247701\\
1.682	-0.0247701\\
1.684	-0.0247701\\
1.686	-0.0247702\\
1.688	-0.0247702\\
1.69	-0.0247703\\
1.692	-0.0247703\\
1.694	-0.0247703\\
1.696	-0.0247703\\
1.698	-0.0247704\\
1.7	-0.0247704\\
1.702	-0.0247704\\
1.704	-0.0247704\\
1.706	-0.0247704\\
1.708	-0.0247705\\
1.71	-0.0247705\\
1.712	-0.0247705\\
1.714	-0.0247705\\
1.716	-0.0247705\\
1.718	-0.0247705\\
1.72	-0.0247705\\
1.722	-0.0247705\\
1.724	-0.0247705\\
1.726	-0.0247705\\
1.728	-0.0247705\\
1.73	-0.0247704\\
1.732	-0.0247704\\
1.734	-0.0247704\\
1.736	-0.0247704\\
1.738	-0.0247704\\
1.74	-0.0247704\\
1.742	-0.0247704\\
1.744	-0.0247703\\
1.746	-0.0247703\\
1.748	-0.0247703\\
1.75	-0.0247703\\
1.752	-0.0247703\\
1.754	-0.0247702\\
1.756	-0.0247702\\
1.758	-0.0247702\\
1.76	-0.0247702\\
1.762	-0.0247702\\
1.764	-0.0247701\\
1.766	-0.0247701\\
1.768	-0.0247701\\
1.77	-0.0247701\\
1.772	-0.02477\\
1.774	-0.02477\\
1.776	-0.02477\\
1.778	-0.02477\\
1.78	-0.02477\\
1.782	-0.0247699\\
1.784	-0.0247699\\
1.786	-0.0247699\\
1.788	-0.0247699\\
1.79	-0.0247699\\
1.792	-0.0247699\\
1.794	-0.0247699\\
1.796	-0.0247698\\
1.798	-0.0247698\\
1.8	-0.0247698\\
1.802	-0.0247698\\
1.804	-0.0247698\\
1.806	-0.0247698\\
1.808	-0.0247698\\
1.81	-0.0247698\\
1.812	-0.0247698\\
1.814	-0.0247698\\
1.816	-0.0247698\\
1.818	-0.0247698\\
1.82	-0.0247698\\
1.822	-0.0247698\\
1.824	-0.0247698\\
1.826	-0.0247698\\
1.828	-0.0247698\\
1.83	-0.0247699\\
1.832	-0.0247699\\
1.834	-0.0247699\\
1.836	-0.0247699\\
1.838	-0.0247699\\
1.84	-0.0247699\\
1.842	-0.0247699\\
1.844	-0.02477\\
1.846	-0.02477\\
1.848	-0.02477\\
1.85	-0.02477\\
1.852	-0.0247701\\
1.854	-0.0247701\\
1.856	-0.0247701\\
1.858	-0.0247701\\
1.86	-0.0247702\\
1.862	-0.0247702\\
1.864	-0.0247702\\
1.866	-0.0247702\\
1.868	-0.0247703\\
1.87	-0.0247703\\
1.872	-0.0247703\\
1.874	-0.0247704\\
1.876	-0.0247704\\
1.878	-0.0247704\\
1.88	-0.0247704\\
1.882	-0.0247705\\
1.884	-0.0247705\\
1.886	-0.0247705\\
1.888	-0.0247706\\
1.89	-0.0247706\\
1.892	-0.0247706\\
1.894	-0.0247706\\
1.896	-0.0247707\\
1.898	-0.0247707\\
1.9	-0.0247707\\
1.902	-0.0247708\\
1.904	-0.0247708\\
1.906	-0.0247708\\
1.908	-0.0247708\\
1.91	-0.0247709\\
1.912	-0.0247709\\
1.914	-0.0247709\\
1.916	-0.0247709\\
1.918	-0.024771\\
1.92	-0.024771\\
1.922	-0.024771\\
1.924	-0.024771\\
1.926	-0.024771\\
1.928	-0.024771\\
1.93	-0.0247711\\
1.932	-0.0247711\\
1.934	-0.0247711\\
1.936	-0.0247711\\
1.938	-0.0247711\\
1.94	-0.0247711\\
1.942	-0.0247711\\
1.944	-0.0247711\\
1.946	-0.0247712\\
1.948	-0.0247712\\
1.95	-0.0247712\\
1.952	-0.0247712\\
1.954	-0.0247712\\
1.956	-0.0247712\\
1.958	-0.0247712\\
1.96	-0.0247712\\
1.962	-0.0247712\\
1.964	-0.0247712\\
1.966	-0.0247712\\
1.968	-0.0247712\\
1.97	-0.0247712\\
1.972	-0.0247712\\
1.974	-0.0247712\\
1.976	-0.0247712\\
1.978	-0.0247712\\
1.98	-0.0247712\\
1.982	-0.0247712\\
1.984	-0.0247712\\
1.986	-0.0247712\\
1.988	-0.0247712\\
1.99	-0.0247712\\
1.992	-0.0247712\\
1.994	-0.0247712\\
1.996	-0.0247712\\
1.998	-0.0247712\\
2	-0.0247712\\
};
\addplot [color=mycolor2, line width=1.0pt, forget plot]
  table[row sep=crcr]{%
-0.024	0\\
-0.022	0\\
-0.02	0\\
-0.018	0\\
-0.016	0\\
-0.014	0\\
-0.012	0\\
-0.01	0\\
-0.008	0\\
-0.006	0\\
-0.004	0\\
-0.002	0\\
0	0\\
0.002	-0.00184649\\
0.004	-0.00342632\\
0.006	-0.00478673\\
0.008	-0.00596703\\
0.01	-0.00699977\\
0.012	-0.00791183\\
0.014	-0.00872528\\
0.016	-0.00945819\\
0.018	-0.0101252\\
0.02	-0.0107383\\
0.022	-0.0113069\\
0.024	-0.0118387\\
0.026	-0.0123396\\
0.028	-0.0128143\\
0.03	-0.0132664\\
0.032	-0.0136986\\
0.034	-0.0141129\\
0.036	-0.0145109\\
0.038	-0.0148936\\
0.04	-0.0152616\\
0.042	-0.0156155\\
0.044	-0.0159557\\
0.046	-0.0162822\\
0.048	-0.0165954\\
0.05	-0.0168951\\
0.052	-0.0171817\\
0.054	-0.0174553\\
0.056	-0.0177158\\
0.058	-0.0179637\\
0.06	-0.0181991\\
0.062	-0.0184223\\
0.064	-0.0186336\\
0.066	-0.0188334\\
0.068	-0.0190222\\
0.07	-0.0192002\\
0.072	-0.019368\\
0.074	-0.0195261\\
0.076	-0.0196749\\
0.078	-0.0198149\\
0.08	-0.0199467\\
0.082	-0.0200707\\
0.084	-0.0201873\\
0.086	-0.0202971\\
0.088	-0.0204006\\
0.09	-0.0204981\\
0.092	-0.0205902\\
0.094	-0.0206771\\
0.096	-0.0207593\\
0.098	-0.0208372\\
0.1	-0.0209112\\
0.102	-0.0209814\\
0.104	-0.0210482\\
0.106	-0.021112\\
0.108	-0.0211728\\
0.11	-0.021231\\
0.112	-0.0212867\\
0.114	-0.0213402\\
0.116	-0.0213915\\
0.118	-0.0214407\\
0.12	-0.0214881\\
0.122	-0.0215336\\
0.124	-0.0215773\\
0.126	-0.0216193\\
0.128	-0.0216595\\
0.13	-0.021698\\
0.132	-0.0217348\\
0.134	-0.0217698\\
0.136	-0.0218031\\
0.138	-0.0218346\\
0.14	-0.0218642\\
0.142	-0.0218919\\
0.144	-0.0219177\\
0.146	-0.0219416\\
0.148	-0.0219634\\
0.15	-0.0219832\\
0.152	-0.0220009\\
0.154	-0.0220165\\
0.156	-0.02203\\
0.158	-0.0220414\\
0.16	-0.0220507\\
0.162	-0.0220579\\
0.164	-0.0220631\\
0.166	-0.0220663\\
0.168	-0.0220676\\
0.17	-0.022067\\
0.172	-0.0220647\\
0.174	-0.0220607\\
0.176	-0.0220552\\
0.178	-0.0220482\\
0.18	-0.0220399\\
0.182	-0.0220305\\
0.184	-0.0220201\\
0.186	-0.0220089\\
0.188	-0.0219969\\
0.19	-0.0219845\\
0.192	-0.0219717\\
0.194	-0.0219587\\
0.196	-0.0219458\\
0.198	-0.021933\\
0.2	-0.0219205\\
0.202	-0.0219086\\
0.204	-0.0218973\\
0.206	-0.0218868\\
0.208	-0.0218773\\
0.21	-0.0218688\\
0.212	-0.0218617\\
0.214	-0.0218558\\
0.216	-0.0218514\\
0.218	-0.0218486\\
0.22	-0.0218474\\
0.222	-0.021848\\
0.224	-0.0218504\\
0.226	-0.0218546\\
0.228	-0.0218607\\
0.23	-0.0218688\\
0.232	-0.0218788\\
0.234	-0.0218908\\
0.236	-0.0219047\\
0.238	-0.0219206\\
0.24	-0.0219385\\
0.242	-0.0219582\\
0.244	-0.0219799\\
0.246	-0.0220034\\
0.248	-0.0220286\\
0.25	-0.0220556\\
0.252	-0.0220842\\
0.254	-0.0221144\\
0.256	-0.0221461\\
0.258	-0.0221793\\
0.26	-0.0222137\\
0.262	-0.0222494\\
0.264	-0.0222863\\
0.266	-0.0223242\\
0.268	-0.0223631\\
0.27	-0.0224029\\
0.272	-0.0224434\\
0.274	-0.0224846\\
0.276	-0.0225264\\
0.278	-0.0225686\\
0.28	-0.0226112\\
0.282	-0.0226542\\
0.284	-0.0226973\\
0.286	-0.0227405\\
0.288	-0.0227838\\
0.29	-0.022827\\
0.292	-0.0228701\\
0.294	-0.0229129\\
0.296	-0.0229555\\
0.298	-0.0229977\\
0.3	-0.0230394\\
0.302	-0.0230807\\
0.304	-0.0231214\\
0.306	-0.0231615\\
0.308	-0.023201\\
0.31	-0.0232397\\
0.312	-0.0232777\\
0.314	-0.0233149\\
0.316	-0.0233513\\
0.318	-0.0233867\\
0.32	-0.0234213\\
0.322	-0.023455\\
0.324	-0.0234877\\
0.326	-0.0235194\\
0.328	-0.0235501\\
0.33	-0.0235798\\
0.332	-0.0236084\\
0.334	-0.0236361\\
0.336	-0.0236627\\
0.338	-0.0236882\\
0.34	-0.0237127\\
0.342	-0.0237361\\
0.344	-0.0237585\\
0.346	-0.0237799\\
0.348	-0.0238002\\
0.35	-0.0238194\\
0.352	-0.0238377\\
0.354	-0.023855\\
0.356	-0.0238712\\
0.358	-0.0238865\\
0.36	-0.0239009\\
0.362	-0.0239143\\
0.364	-0.0239268\\
0.366	-0.0239385\\
0.368	-0.0239493\\
0.37	-0.0239592\\
0.372	-0.0239684\\
0.374	-0.0239768\\
0.376	-0.0239844\\
0.378	-0.0239914\\
0.38	-0.0239977\\
0.382	-0.0240034\\
0.384	-0.0240084\\
0.386	-0.0240129\\
0.388	-0.0240169\\
0.39	-0.0240204\\
0.392	-0.0240235\\
0.394	-0.0240262\\
0.396	-0.0240285\\
0.398	-0.0240305\\
0.4	-0.0240323\\
0.402	-0.0240338\\
0.404	-0.024035\\
0.406	-0.0240362\\
0.408	-0.0240372\\
0.41	-0.0240382\\
0.412	-0.0240391\\
0.414	-0.02404\\
0.416	-0.024041\\
0.418	-0.024042\\
0.42	-0.0240432\\
0.422	-0.0240444\\
0.424	-0.0240459\\
0.426	-0.0240476\\
0.428	-0.0240495\\
0.43	-0.0240516\\
0.432	-0.0240541\\
0.434	-0.0240568\\
0.436	-0.0240599\\
0.438	-0.0240633\\
0.44	-0.0240671\\
0.442	-0.0240713\\
0.444	-0.0240759\\
0.446	-0.0240809\\
0.448	-0.0240863\\
0.45	-0.0240921\\
0.452	-0.0240983\\
0.454	-0.0241049\\
0.456	-0.0241119\\
0.458	-0.0241194\\
0.46	-0.0241273\\
0.462	-0.0241355\\
0.464	-0.0241442\\
0.466	-0.0241532\\
0.468	-0.0241626\\
0.47	-0.0241723\\
0.472	-0.0241824\\
0.474	-0.0241927\\
0.476	-0.0242033\\
0.478	-0.0242142\\
0.48	-0.0242253\\
0.482	-0.0242366\\
0.484	-0.0242481\\
0.486	-0.0242598\\
0.488	-0.0242715\\
0.49	-0.0242834\\
0.492	-0.0242953\\
0.494	-0.0243072\\
0.496	-0.0243192\\
0.498	-0.0243311\\
0.5	-0.0243429\\
0.502	-0.0243546\\
0.504	-0.0243662\\
0.506	-0.0243777\\
0.508	-0.0243889\\
0.51	-0.0243999\\
0.512	-0.0244107\\
0.514	-0.0244212\\
0.516	-0.0244314\\
0.518	-0.0244412\\
0.52	-0.0244507\\
0.522	-0.0244598\\
0.524	-0.0244685\\
0.526	-0.0244768\\
0.528	-0.0244847\\
0.53	-0.0244921\\
0.532	-0.024499\\
0.534	-0.0245054\\
0.536	-0.0245114\\
0.538	-0.0245168\\
0.54	-0.0245217\\
0.542	-0.0245261\\
0.544	-0.02453\\
0.546	-0.0245333\\
0.548	-0.0245361\\
0.55	-0.0245384\\
0.552	-0.0245402\\
0.554	-0.0245415\\
0.556	-0.0245423\\
0.558	-0.0245425\\
0.56	-0.0245423\\
0.562	-0.0245416\\
0.564	-0.0245405\\
0.566	-0.0245389\\
0.568	-0.0245369\\
0.57	-0.0245345\\
0.572	-0.0245317\\
0.574	-0.0245286\\
0.576	-0.0245251\\
0.578	-0.0245214\\
0.58	-0.0245173\\
0.582	-0.024513\\
0.584	-0.0245084\\
0.586	-0.0245036\\
0.588	-0.0244987\\
0.59	-0.0244936\\
0.592	-0.0244883\\
0.594	-0.024483\\
0.596	-0.0244776\\
0.598	-0.0244722\\
0.6	-0.0244667\\
0.602	-0.0244612\\
0.604	-0.0244558\\
0.606	-0.0244505\\
0.608	-0.0244452\\
0.61	-0.0244401\\
0.612	-0.0244351\\
0.614	-0.0244303\\
0.616	-0.0244256\\
0.618	-0.0244212\\
0.62	-0.0244169\\
0.622	-0.024413\\
0.624	-0.0244092\\
0.626	-0.0244058\\
0.628	-0.0244026\\
0.63	-0.0243998\\
0.632	-0.0243972\\
0.634	-0.024395\\
0.636	-0.0243931\\
0.638	-0.0243916\\
0.64	-0.0243904\\
0.642	-0.0243895\\
0.644	-0.0243891\\
0.646	-0.0243889\\
0.648	-0.0243891\\
0.65	-0.0243897\\
0.652	-0.0243906\\
0.654	-0.0243919\\
0.656	-0.0243935\\
0.658	-0.0243955\\
0.66	-0.0243978\\
0.662	-0.0244004\\
0.664	-0.0244033\\
0.666	-0.0244065\\
0.668	-0.0244101\\
0.67	-0.0244138\\
0.672	-0.0244179\\
0.674	-0.0244222\\
0.676	-0.0244268\\
0.678	-0.0244315\\
0.68	-0.0244365\\
0.682	-0.0244417\\
0.684	-0.024447\\
0.686	-0.0244525\\
0.688	-0.0244582\\
0.69	-0.0244639\\
0.692	-0.0244698\\
0.694	-0.0244758\\
0.696	-0.0244819\\
0.698	-0.024488\\
0.7	-0.0244941\\
0.702	-0.0245003\\
0.704	-0.0245065\\
0.706	-0.0245127\\
0.708	-0.0245189\\
0.71	-0.024525\\
0.712	-0.0245311\\
0.714	-0.0245372\\
0.716	-0.0245432\\
0.718	-0.024549\\
0.72	-0.0245548\\
0.722	-0.0245605\\
0.724	-0.0245661\\
0.726	-0.0245715\\
0.728	-0.0245768\\
0.73	-0.024582\\
0.732	-0.024587\\
0.734	-0.0245918\\
0.736	-0.0245965\\
0.738	-0.024601\\
0.74	-0.0246053\\
0.742	-0.0246095\\
0.744	-0.0246134\\
0.746	-0.0246172\\
0.748	-0.0246208\\
0.75	-0.0246242\\
0.752	-0.0246274\\
0.754	-0.0246304\\
0.756	-0.0246332\\
0.758	-0.0246358\\
0.76	-0.0246382\\
0.762	-0.0246404\\
0.764	-0.0246425\\
0.766	-0.0246443\\
0.768	-0.024646\\
0.77	-0.0246475\\
0.772	-0.0246488\\
0.774	-0.0246499\\
0.776	-0.0246508\\
0.778	-0.0246516\\
0.78	-0.0246523\\
0.782	-0.0246528\\
0.784	-0.0246531\\
0.786	-0.0246533\\
0.788	-0.0246533\\
0.79	-0.0246532\\
0.792	-0.024653\\
0.794	-0.0246527\\
0.796	-0.0246523\\
0.798	-0.0246517\\
0.8	-0.024651\\
0.802	-0.0246503\\
0.804	-0.0246495\\
0.806	-0.0246486\\
0.808	-0.0246476\\
0.81	-0.0246465\\
0.812	-0.0246454\\
0.814	-0.0246442\\
0.816	-0.024643\\
0.818	-0.0246418\\
0.82	-0.0246405\\
0.822	-0.0246392\\
0.824	-0.0246379\\
0.826	-0.0246365\\
0.828	-0.0246352\\
0.83	-0.0246338\\
0.832	-0.0246325\\
0.834	-0.0246311\\
0.836	-0.0246298\\
0.838	-0.0246285\\
0.84	-0.0246273\\
0.842	-0.024626\\
0.844	-0.0246249\\
0.846	-0.0246237\\
0.848	-0.0246226\\
0.85	-0.0246216\\
0.852	-0.0246206\\
0.854	-0.0246197\\
0.856	-0.0246188\\
0.858	-0.024618\\
0.86	-0.0246173\\
0.862	-0.0246167\\
0.864	-0.0246162\\
0.866	-0.0246157\\
0.868	-0.0246153\\
0.87	-0.024615\\
0.872	-0.0246149\\
0.874	-0.0246148\\
0.876	-0.0246148\\
0.878	-0.0246149\\
0.88	-0.0246151\\
0.882	-0.0246154\\
0.884	-0.0246158\\
0.886	-0.0246163\\
0.888	-0.0246169\\
0.89	-0.0246176\\
0.892	-0.0246184\\
0.894	-0.0246193\\
0.896	-0.0246203\\
0.898	-0.0246214\\
0.9	-0.0246226\\
0.902	-0.0246239\\
0.904	-0.0246253\\
0.906	-0.0246268\\
0.908	-0.0246284\\
0.91	-0.02463\\
0.912	-0.0246318\\
0.914	-0.0246336\\
0.916	-0.0246355\\
0.918	-0.0246375\\
0.92	-0.0246396\\
0.922	-0.0246417\\
0.924	-0.0246439\\
0.926	-0.0246461\\
0.928	-0.0246484\\
0.93	-0.0246508\\
0.932	-0.0246531\\
0.934	-0.0246556\\
0.936	-0.0246581\\
0.938	-0.0246606\\
0.94	-0.0246631\\
0.942	-0.0246657\\
0.944	-0.0246683\\
0.946	-0.0246708\\
0.948	-0.0246734\\
0.95	-0.0246761\\
0.952	-0.0246787\\
0.954	-0.0246813\\
0.956	-0.0246838\\
0.958	-0.0246864\\
0.96	-0.024689\\
0.962	-0.0246915\\
0.964	-0.024694\\
0.966	-0.0246965\\
0.968	-0.0246989\\
0.97	-0.0247013\\
0.972	-0.0247037\\
0.974	-0.024706\\
0.976	-0.0247082\\
0.978	-0.0247104\\
0.98	-0.0247126\\
0.982	-0.0247147\\
0.984	-0.0247167\\
0.986	-0.0247186\\
0.988	-0.0247205\\
0.99	-0.0247224\\
0.992	-0.0247241\\
0.994	-0.0247258\\
0.996	-0.0247274\\
0.998	-0.0247289\\
1	-0.0247303\\
1.002	-0.0247317\\
1.004	-0.024733\\
1.006	-0.0247342\\
1.008	-0.0247354\\
1.01	-0.0247364\\
1.012	-0.0247374\\
1.014	-0.0247383\\
1.016	-0.0247391\\
1.018	-0.0247399\\
1.02	-0.0247406\\
1.022	-0.0247412\\
1.024	-0.0247417\\
1.026	-0.0247421\\
1.028	-0.0247425\\
1.03	-0.0247429\\
1.032	-0.0247431\\
1.034	-0.0247433\\
1.036	-0.0247434\\
1.038	-0.0247435\\
1.04	-0.0247435\\
1.042	-0.0247435\\
1.044	-0.0247434\\
1.046	-0.0247433\\
1.048	-0.0247431\\
1.05	-0.0247429\\
1.052	-0.0247427\\
1.054	-0.0247424\\
1.056	-0.024742\\
1.058	-0.0247417\\
1.06	-0.0247413\\
1.062	-0.0247409\\
1.064	-0.0247405\\
1.066	-0.02474\\
1.068	-0.0247396\\
1.07	-0.0247391\\
1.072	-0.0247386\\
1.074	-0.0247381\\
1.076	-0.0247376\\
1.078	-0.0247371\\
1.08	-0.0247366\\
1.082	-0.0247361\\
1.084	-0.0247356\\
1.086	-0.0247351\\
1.088	-0.0247346\\
1.09	-0.0247342\\
1.092	-0.0247337\\
1.094	-0.0247333\\
1.096	-0.0247328\\
1.098	-0.0247324\\
1.1	-0.024732\\
1.102	-0.0247317\\
1.104	-0.0247313\\
1.106	-0.024731\\
1.108	-0.0247307\\
1.11	-0.0247305\\
1.112	-0.0247302\\
1.114	-0.02473\\
1.116	-0.0247298\\
1.118	-0.0247297\\
1.12	-0.0247295\\
1.122	-0.0247294\\
1.124	-0.0247294\\
1.126	-0.0247293\\
1.128	-0.0247293\\
1.13	-0.0247293\\
1.132	-0.0247294\\
1.134	-0.0247295\\
1.136	-0.0247296\\
1.138	-0.0247297\\
1.14	-0.0247299\\
1.142	-0.0247301\\
1.144	-0.0247303\\
1.146	-0.0247305\\
1.148	-0.0247308\\
1.15	-0.0247311\\
1.152	-0.0247314\\
1.154	-0.0247317\\
1.156	-0.024732\\
1.158	-0.0247324\\
1.16	-0.0247328\\
1.162	-0.0247332\\
1.164	-0.0247336\\
1.166	-0.024734\\
1.168	-0.0247345\\
1.17	-0.0247349\\
1.172	-0.0247354\\
1.174	-0.0247359\\
1.176	-0.0247364\\
1.178	-0.0247369\\
1.18	-0.0247374\\
1.182	-0.0247379\\
1.184	-0.0247384\\
1.186	-0.0247389\\
1.188	-0.0247394\\
1.19	-0.0247399\\
1.192	-0.0247404\\
1.194	-0.0247409\\
1.196	-0.0247414\\
1.198	-0.0247419\\
1.2	-0.0247424\\
1.202	-0.0247429\\
1.204	-0.0247433\\
1.206	-0.0247438\\
1.208	-0.0247442\\
1.21	-0.0247447\\
1.212	-0.0247451\\
1.214	-0.0247455\\
1.216	-0.0247459\\
1.218	-0.0247463\\
1.22	-0.0247467\\
1.222	-0.0247471\\
1.224	-0.0247474\\
1.226	-0.0247477\\
1.228	-0.0247481\\
1.23	-0.0247484\\
1.232	-0.0247486\\
1.234	-0.0247489\\
1.236	-0.0247491\\
1.238	-0.0247494\\
1.24	-0.0247496\\
1.242	-0.0247498\\
1.244	-0.02475\\
1.246	-0.0247501\\
1.248	-0.0247503\\
1.25	-0.0247504\\
1.252	-0.0247505\\
1.254	-0.0247506\\
1.256	-0.0247507\\
1.258	-0.0247508\\
1.26	-0.0247508\\
1.262	-0.0247509\\
1.264	-0.0247509\\
1.266	-0.0247509\\
1.268	-0.0247509\\
1.27	-0.0247509\\
1.272	-0.0247509\\
1.274	-0.0247508\\
1.276	-0.0247508\\
1.278	-0.0247507\\
1.28	-0.0247506\\
1.282	-0.0247506\\
1.284	-0.0247505\\
1.286	-0.0247504\\
1.288	-0.0247503\\
1.29	-0.0247502\\
1.292	-0.0247501\\
1.294	-0.02475\\
1.296	-0.0247498\\
1.298	-0.0247497\\
1.3	-0.0247496\\
1.302	-0.0247495\\
1.304	-0.0247494\\
1.306	-0.0247492\\
1.308	-0.0247491\\
1.31	-0.024749\\
1.312	-0.0247489\\
1.314	-0.0247488\\
1.316	-0.0247486\\
1.318	-0.0247485\\
1.32	-0.0247484\\
1.322	-0.0247483\\
1.324	-0.0247482\\
1.326	-0.0247482\\
1.328	-0.0247481\\
1.33	-0.024748\\
1.332	-0.0247479\\
1.334	-0.0247479\\
1.336	-0.0247478\\
1.338	-0.0247478\\
1.34	-0.0247478\\
1.342	-0.0247477\\
1.344	-0.0247477\\
1.346	-0.0247477\\
1.348	-0.0247477\\
1.35	-0.0247478\\
1.352	-0.0247478\\
1.354	-0.0247478\\
1.356	-0.0247479\\
1.358	-0.0247479\\
1.36	-0.024748\\
1.362	-0.0247481\\
1.364	-0.0247482\\
1.366	-0.0247483\\
1.368	-0.0247484\\
1.37	-0.0247485\\
1.372	-0.0247486\\
1.374	-0.0247488\\
1.376	-0.0247489\\
1.378	-0.0247491\\
1.38	-0.0247493\\
1.382	-0.0247495\\
1.384	-0.0247496\\
1.386	-0.0247498\\
1.388	-0.02475\\
1.39	-0.0247503\\
1.392	-0.0247505\\
1.394	-0.0247507\\
1.396	-0.0247509\\
1.398	-0.0247512\\
1.4	-0.0247514\\
1.402	-0.0247517\\
1.404	-0.0247519\\
1.406	-0.0247522\\
1.408	-0.0247524\\
1.41	-0.0247527\\
1.412	-0.024753\\
1.414	-0.0247533\\
1.416	-0.0247535\\
1.418	-0.0247538\\
1.42	-0.0247541\\
1.422	-0.0247543\\
1.424	-0.0247546\\
1.426	-0.0247549\\
1.428	-0.0247552\\
1.43	-0.0247554\\
1.432	-0.0247557\\
1.434	-0.024756\\
1.436	-0.0247563\\
1.438	-0.0247565\\
1.44	-0.0247568\\
1.442	-0.024757\\
1.444	-0.0247573\\
1.446	-0.0247575\\
1.448	-0.0247578\\
1.45	-0.024758\\
1.452	-0.0247583\\
1.454	-0.0247585\\
1.456	-0.0247587\\
1.458	-0.0247589\\
1.46	-0.0247591\\
1.462	-0.0247594\\
1.464	-0.0247596\\
1.466	-0.0247598\\
1.468	-0.0247599\\
1.47	-0.0247601\\
1.472	-0.0247603\\
1.474	-0.0247605\\
1.476	-0.0247606\\
1.478	-0.0247608\\
1.48	-0.0247609\\
1.482	-0.0247611\\
1.484	-0.0247612\\
1.486	-0.0247613\\
1.488	-0.0247615\\
1.49	-0.0247616\\
1.492	-0.0247617\\
1.494	-0.0247618\\
1.496	-0.0247619\\
1.498	-0.024762\\
1.5	-0.0247621\\
1.502	-0.0247622\\
1.504	-0.0247623\\
1.506	-0.0247623\\
1.508	-0.0247624\\
1.51	-0.0247625\\
1.512	-0.0247625\\
1.514	-0.0247626\\
1.516	-0.0247627\\
1.518	-0.0247627\\
1.52	-0.0247628\\
1.522	-0.0247628\\
1.524	-0.0247629\\
1.526	-0.0247629\\
1.528	-0.0247629\\
1.53	-0.024763\\
1.532	-0.024763\\
1.534	-0.0247631\\
1.536	-0.0247631\\
1.538	-0.0247631\\
1.54	-0.0247632\\
1.542	-0.0247632\\
1.544	-0.0247633\\
1.546	-0.0247633\\
1.548	-0.0247633\\
1.55	-0.0247634\\
1.552	-0.0247634\\
1.554	-0.0247635\\
1.556	-0.0247635\\
1.558	-0.0247636\\
1.56	-0.0247636\\
1.562	-0.0247637\\
1.564	-0.0247637\\
1.566	-0.0247638\\
1.568	-0.0247638\\
1.57	-0.0247639\\
1.572	-0.024764\\
1.574	-0.024764\\
1.576	-0.0247641\\
1.578	-0.0247642\\
1.58	-0.0247643\\
1.582	-0.0247644\\
1.584	-0.0247644\\
1.586	-0.0247645\\
1.588	-0.0247646\\
1.59	-0.0247647\\
1.592	-0.0247648\\
1.594	-0.0247649\\
1.596	-0.024765\\
1.598	-0.0247651\\
1.6	-0.0247652\\
1.602	-0.0247653\\
1.604	-0.0247654\\
1.606	-0.0247656\\
1.608	-0.0247657\\
1.61	-0.0247658\\
1.612	-0.0247659\\
1.614	-0.024766\\
1.616	-0.0247662\\
1.618	-0.0247663\\
1.62	-0.0247664\\
1.622	-0.0247665\\
1.624	-0.0247667\\
1.626	-0.0247668\\
1.628	-0.0247669\\
1.63	-0.024767\\
1.632	-0.0247672\\
1.634	-0.0247673\\
1.636	-0.0247674\\
1.638	-0.0247676\\
1.64	-0.0247677\\
1.642	-0.0247678\\
1.644	-0.0247679\\
1.646	-0.024768\\
1.648	-0.0247682\\
1.65	-0.0247683\\
1.652	-0.0247684\\
1.654	-0.0247685\\
1.656	-0.0247686\\
1.658	-0.0247687\\
1.66	-0.0247688\\
1.662	-0.0247689\\
1.664	-0.024769\\
1.666	-0.0247691\\
1.668	-0.0247692\\
1.67	-0.0247693\\
1.672	-0.0247694\\
1.674	-0.0247694\\
1.676	-0.0247695\\
1.678	-0.0247696\\
1.68	-0.0247697\\
1.682	-0.0247697\\
1.684	-0.0247698\\
1.686	-0.0247698\\
1.688	-0.0247699\\
1.69	-0.0247699\\
1.692	-0.02477\\
1.694	-0.02477\\
1.696	-0.02477\\
1.698	-0.0247701\\
1.7	-0.0247701\\
1.702	-0.0247701\\
1.704	-0.0247701\\
1.706	-0.0247701\\
1.708	-0.0247701\\
1.71	-0.0247702\\
1.712	-0.0247702\\
1.714	-0.0247702\\
1.716	-0.0247701\\
1.718	-0.0247701\\
1.72	-0.0247701\\
1.722	-0.0247701\\
1.724	-0.0247701\\
1.726	-0.0247701\\
1.728	-0.0247701\\
1.73	-0.02477\\
1.732	-0.02477\\
1.734	-0.02477\\
1.736	-0.0247699\\
1.738	-0.0247699\\
1.74	-0.0247699\\
1.742	-0.0247698\\
1.744	-0.0247698\\
1.746	-0.0247698\\
1.748	-0.0247697\\
1.75	-0.0247697\\
1.752	-0.0247696\\
1.754	-0.0247696\\
1.756	-0.0247695\\
1.758	-0.0247695\\
1.76	-0.0247695\\
1.762	-0.0247694\\
1.764	-0.0247694\\
1.766	-0.0247693\\
1.768	-0.0247693\\
1.77	-0.0247693\\
1.772	-0.0247692\\
1.774	-0.0247692\\
1.776	-0.0247692\\
1.778	-0.0247691\\
1.78	-0.0247691\\
1.782	-0.0247691\\
1.784	-0.024769\\
1.786	-0.024769\\
1.788	-0.024769\\
1.79	-0.024769\\
1.792	-0.0247689\\
1.794	-0.0247689\\
1.796	-0.0247689\\
1.798	-0.0247689\\
1.8	-0.0247689\\
1.802	-0.0247689\\
1.804	-0.0247689\\
1.806	-0.0247689\\
1.808	-0.0247689\\
1.81	-0.0247689\\
1.812	-0.0247689\\
1.814	-0.0247689\\
1.816	-0.0247689\\
1.818	-0.0247689\\
1.82	-0.0247689\\
1.822	-0.0247689\\
1.824	-0.0247689\\
1.826	-0.024769\\
1.828	-0.024769\\
1.83	-0.024769\\
1.832	-0.024769\\
1.834	-0.0247691\\
1.836	-0.0247691\\
1.838	-0.0247691\\
1.84	-0.0247692\\
1.842	-0.0247692\\
1.844	-0.0247692\\
1.846	-0.0247693\\
1.848	-0.0247693\\
1.85	-0.0247694\\
1.852	-0.0247694\\
1.854	-0.0247694\\
1.856	-0.0247695\\
1.858	-0.0247695\\
1.86	-0.0247696\\
1.862	-0.0247696\\
1.864	-0.0247697\\
1.866	-0.0247697\\
1.868	-0.0247698\\
1.87	-0.0247698\\
1.872	-0.0247699\\
1.874	-0.0247699\\
1.876	-0.02477\\
1.878	-0.02477\\
1.88	-0.0247701\\
1.882	-0.0247701\\
1.884	-0.0247702\\
1.886	-0.0247702\\
1.888	-0.0247703\\
1.89	-0.0247703\\
1.892	-0.0247704\\
1.894	-0.0247704\\
1.896	-0.0247704\\
1.898	-0.0247705\\
1.9	-0.0247705\\
1.902	-0.0247706\\
1.904	-0.0247706\\
1.906	-0.0247706\\
1.908	-0.0247707\\
1.91	-0.0247707\\
1.912	-0.0247707\\
1.914	-0.0247708\\
1.916	-0.0247708\\
1.918	-0.0247708\\
1.92	-0.0247708\\
1.922	-0.0247709\\
1.924	-0.0247709\\
1.926	-0.0247709\\
1.928	-0.0247709\\
1.93	-0.0247709\\
1.932	-0.0247709\\
1.934	-0.0247709\\
1.936	-0.024771\\
1.938	-0.024771\\
1.94	-0.024771\\
1.942	-0.024771\\
1.944	-0.024771\\
1.946	-0.024771\\
1.948	-0.024771\\
1.95	-0.024771\\
1.952	-0.024771\\
1.954	-0.024771\\
1.956	-0.024771\\
1.958	-0.024771\\
1.96	-0.024771\\
1.962	-0.0247709\\
1.964	-0.0247709\\
1.966	-0.0247709\\
1.968	-0.0247709\\
1.97	-0.0247709\\
1.972	-0.0247709\\
1.974	-0.0247709\\
1.976	-0.0247709\\
1.978	-0.0247709\\
1.98	-0.0247708\\
1.982	-0.0247708\\
1.984	-0.0247708\\
1.986	-0.0247708\\
1.988	-0.0247708\\
1.99	-0.0247708\\
1.992	-0.0247708\\
1.994	-0.0247708\\
1.996	-0.0247707\\
1.998	-0.0247707\\
2	-0.0247707\\
};
\addplot [color=mycolor3, line width=1.0pt, forget plot]
  table[row sep=crcr]{%
-0.024	0\\
-0.022	0\\
-0.02	0\\
-0.018	0\\
-0.016	0\\
-0.014	0\\
-0.012	0\\
-0.01	0\\
-0.008	0\\
-0.006	0\\
-0.004	0\\
-0.002	0\\
0	0\\
0.002	-0.00184649\\
0.004	-0.00342632\\
0.006	-0.00478673\\
0.008	-0.00596703\\
0.01	-0.00699977\\
0.012	-0.00791183\\
0.014	-0.00872529\\
0.016	-0.0094582\\
0.018	-0.0101253\\
0.02	-0.0107383\\
0.022	-0.011307\\
0.024	-0.0118388\\
0.026	-0.0123397\\
0.028	-0.0128144\\
0.03	-0.0132666\\
0.032	-0.0136988\\
0.034	-0.0141133\\
0.036	-0.0145113\\
0.038	-0.0148941\\
0.04	-0.0152623\\
0.042	-0.0156164\\
0.044	-0.0159567\\
0.046	-0.0162835\\
0.048	-0.0165968\\
0.05	-0.0168969\\
0.052	-0.0171838\\
0.054	-0.0174577\\
0.056	-0.0177187\\
0.058	-0.017967\\
0.06	-0.0182029\\
0.062	-0.0184267\\
0.064	-0.0186386\\
0.066	-0.0188391\\
0.068	-0.0190285\\
0.07	-0.0192073\\
0.072	-0.019376\\
0.074	-0.019535\\
0.076	-0.0196847\\
0.078	-0.0198258\\
0.08	-0.0199586\\
0.082	-0.0200838\\
0.084	-0.0202016\\
0.086	-0.0203127\\
0.088	-0.0204175\\
0.09	-0.0205165\\
0.092	-0.02061\\
0.094	-0.0206985\\
0.096	-0.0207823\\
0.098	-0.0208618\\
0.1	-0.0209374\\
0.102	-0.0210094\\
0.104	-0.0210781\\
0.106	-0.0211437\\
0.108	-0.0212064\\
0.11	-0.0212665\\
0.112	-0.0213243\\
0.114	-0.0213797\\
0.116	-0.021433\\
0.118	-0.0214844\\
0.12	-0.0215338\\
0.122	-0.0215814\\
0.124	-0.0216272\\
0.126	-0.0216713\\
0.128	-0.0217136\\
0.13	-0.0217543\\
0.132	-0.0217932\\
0.134	-0.0218303\\
0.136	-0.0218656\\
0.138	-0.0218991\\
0.14	-0.0219308\\
0.142	-0.0219605\\
0.144	-0.0219882\\
0.146	-0.022014\\
0.148	-0.0220376\\
0.15	-0.0220592\\
0.152	-0.0220786\\
0.154	-0.0220959\\
0.156	-0.022111\\
0.158	-0.0221239\\
0.16	-0.0221346\\
0.162	-0.0221432\\
0.164	-0.0221496\\
0.166	-0.022154\\
0.168	-0.0221563\\
0.17	-0.0221567\\
0.172	-0.0221552\\
0.174	-0.0221519\\
0.176	-0.022147\\
0.178	-0.0221405\\
0.18	-0.0221327\\
0.182	-0.0221235\\
0.184	-0.0221133\\
0.186	-0.0221021\\
0.188	-0.02209\\
0.19	-0.0220774\\
0.192	-0.0220642\\
0.194	-0.0220508\\
0.196	-0.0220373\\
0.198	-0.0220238\\
0.2	-0.0220105\\
0.202	-0.0219976\\
0.204	-0.0219852\\
0.206	-0.0219736\\
0.208	-0.0219628\\
0.21	-0.0219529\\
0.212	-0.0219442\\
0.214	-0.0219368\\
0.216	-0.0219307\\
0.218	-0.021926\\
0.22	-0.0219229\\
0.222	-0.0219215\\
0.224	-0.0219218\\
0.226	-0.0219238\\
0.228	-0.0219277\\
0.23	-0.0219335\\
0.232	-0.0219411\\
0.234	-0.0219507\\
0.236	-0.0219621\\
0.238	-0.0219755\\
0.24	-0.0219908\\
0.242	-0.022008\\
0.244	-0.022027\\
0.246	-0.0220478\\
0.248	-0.0220704\\
0.25	-0.0220947\\
0.252	-0.0221206\\
0.254	-0.0221482\\
0.256	-0.0221772\\
0.258	-0.0222077\\
0.26	-0.0222395\\
0.262	-0.0222726\\
0.264	-0.0223069\\
0.266	-0.0223423\\
0.268	-0.0223787\\
0.27	-0.022416\\
0.272	-0.0224541\\
0.274	-0.022493\\
0.276	-0.0225325\\
0.278	-0.0225725\\
0.28	-0.022613\\
0.282	-0.0226539\\
0.284	-0.022695\\
0.286	-0.0227364\\
0.288	-0.0227778\\
0.29	-0.0228193\\
0.292	-0.0228608\\
0.294	-0.0229021\\
0.296	-0.0229432\\
0.298	-0.0229841\\
0.3	-0.0230246\\
0.302	-0.0230647\\
0.304	-0.0231044\\
0.306	-0.0231435\\
0.308	-0.0231821\\
0.31	-0.0232201\\
0.312	-0.0232574\\
0.314	-0.023294\\
0.316	-0.0233299\\
0.318	-0.023365\\
0.32	-0.0233993\\
0.322	-0.0234327\\
0.324	-0.0234653\\
0.326	-0.023497\\
0.328	-0.0235277\\
0.33	-0.0235575\\
0.332	-0.0235864\\
0.334	-0.0236143\\
0.336	-0.0236412\\
0.338	-0.0236671\\
0.34	-0.023692\\
0.342	-0.0237159\\
0.344	-0.0237388\\
0.346	-0.0237607\\
0.348	-0.0237816\\
0.35	-0.0238015\\
0.352	-0.0238204\\
0.354	-0.0238383\\
0.356	-0.0238552\\
0.358	-0.0238712\\
0.36	-0.0238863\\
0.362	-0.0239004\\
0.364	-0.0239136\\
0.366	-0.0239259\\
0.368	-0.0239373\\
0.37	-0.0239479\\
0.372	-0.0239576\\
0.374	-0.0239666\\
0.376	-0.0239748\\
0.378	-0.0239822\\
0.38	-0.023989\\
0.382	-0.023995\\
0.384	-0.0240005\\
0.386	-0.0240053\\
0.388	-0.0240095\\
0.39	-0.0240132\\
0.392	-0.0240163\\
0.394	-0.0240191\\
0.396	-0.0240213\\
0.398	-0.0240232\\
0.4	-0.0240248\\
0.402	-0.024026\\
0.404	-0.024027\\
0.406	-0.0240277\\
0.408	-0.0240282\\
0.41	-0.0240286\\
0.412	-0.0240289\\
0.414	-0.0240291\\
0.416	-0.0240292\\
0.418	-0.0240294\\
0.42	-0.0240296\\
0.422	-0.0240298\\
0.424	-0.0240302\\
0.426	-0.0240307\\
0.428	-0.0240313\\
0.43	-0.0240322\\
0.432	-0.0240333\\
0.434	-0.0240346\\
0.436	-0.0240362\\
0.438	-0.0240381\\
0.44	-0.0240403\\
0.442	-0.0240428\\
0.444	-0.0240457\\
0.446	-0.024049\\
0.448	-0.0240526\\
0.45	-0.0240566\\
0.452	-0.024061\\
0.454	-0.0240658\\
0.456	-0.0240711\\
0.458	-0.0240767\\
0.46	-0.0240827\\
0.462	-0.0240891\\
0.464	-0.0240959\\
0.466	-0.0241031\\
0.468	-0.0241107\\
0.47	-0.0241186\\
0.472	-0.0241269\\
0.474	-0.0241355\\
0.476	-0.0241444\\
0.478	-0.0241536\\
0.48	-0.0241631\\
0.482	-0.0241729\\
0.484	-0.0241828\\
0.486	-0.024193\\
0.488	-0.0242034\\
0.49	-0.0242139\\
0.492	-0.0242246\\
0.494	-0.0242354\\
0.496	-0.0242462\\
0.498	-0.0242571\\
0.5	-0.024268\\
0.502	-0.0242789\\
0.504	-0.0242897\\
0.506	-0.0243005\\
0.508	-0.0243112\\
0.51	-0.0243217\\
0.512	-0.0243321\\
0.514	-0.0243424\\
0.516	-0.0243524\\
0.518	-0.0243622\\
0.52	-0.0243717\\
0.522	-0.024381\\
0.524	-0.0243899\\
0.526	-0.0243986\\
0.528	-0.0244069\\
0.53	-0.0244148\\
0.532	-0.0244224\\
0.534	-0.0244296\\
0.536	-0.0244363\\
0.538	-0.0244427\\
0.54	-0.0244487\\
0.542	-0.0244542\\
0.544	-0.0244593\\
0.546	-0.024464\\
0.548	-0.0244682\\
0.55	-0.024472\\
0.552	-0.0244753\\
0.554	-0.0244782\\
0.556	-0.0244807\\
0.558	-0.0244827\\
0.56	-0.0244843\\
0.562	-0.0244855\\
0.564	-0.0244863\\
0.566	-0.0244868\\
0.568	-0.0244868\\
0.57	-0.0244865\\
0.572	-0.0244858\\
0.574	-0.0244849\\
0.576	-0.0244836\\
0.578	-0.024482\\
0.58	-0.0244802\\
0.582	-0.0244781\\
0.584	-0.0244758\\
0.586	-0.0244732\\
0.588	-0.0244705\\
0.59	-0.0244677\\
0.592	-0.0244647\\
0.594	-0.0244616\\
0.596	-0.0244584\\
0.598	-0.0244551\\
0.6	-0.0244518\\
0.602	-0.0244485\\
0.604	-0.0244452\\
0.606	-0.0244419\\
0.608	-0.0244387\\
0.61	-0.0244355\\
0.612	-0.0244324\\
0.614	-0.0244294\\
0.616	-0.0244266\\
0.618	-0.0244239\\
0.62	-0.0244214\\
0.622	-0.024419\\
0.624	-0.0244168\\
0.626	-0.0244149\\
0.628	-0.0244131\\
0.63	-0.0244116\\
0.632	-0.0244103\\
0.634	-0.0244093\\
0.636	-0.0244085\\
0.638	-0.024408\\
0.64	-0.0244078\\
0.642	-0.0244078\\
0.644	-0.0244081\\
0.646	-0.0244087\\
0.648	-0.0244096\\
0.65	-0.0244107\\
0.652	-0.0244121\\
0.654	-0.0244138\\
0.656	-0.0244158\\
0.658	-0.024418\\
0.66	-0.0244204\\
0.662	-0.0244232\\
0.664	-0.0244261\\
0.666	-0.0244293\\
0.668	-0.0244327\\
0.67	-0.0244363\\
0.672	-0.0244401\\
0.674	-0.0244442\\
0.676	-0.0244483\\
0.678	-0.0244527\\
0.68	-0.0244572\\
0.682	-0.0244619\\
0.684	-0.0244666\\
0.686	-0.0244715\\
0.688	-0.0244765\\
0.69	-0.0244816\\
0.692	-0.0244867\\
0.694	-0.024492\\
0.696	-0.0244972\\
0.698	-0.0245025\\
0.7	-0.0245078\\
0.702	-0.0245131\\
0.704	-0.0245185\\
0.706	-0.0245238\\
0.708	-0.024529\\
0.71	-0.0245343\\
0.712	-0.0245394\\
0.714	-0.0245446\\
0.716	-0.0245496\\
0.718	-0.0245546\\
0.72	-0.0245595\\
0.722	-0.0245642\\
0.724	-0.0245689\\
0.726	-0.0245735\\
0.728	-0.0245779\\
0.73	-0.0245822\\
0.732	-0.0245864\\
0.734	-0.0245905\\
0.736	-0.0245944\\
0.738	-0.0245982\\
0.74	-0.0246018\\
0.742	-0.0246053\\
0.744	-0.0246086\\
0.746	-0.0246118\\
0.748	-0.0246148\\
0.75	-0.0246177\\
0.752	-0.0246204\\
0.754	-0.0246229\\
0.756	-0.0246253\\
0.758	-0.0246276\\
0.76	-0.0246297\\
0.762	-0.0246316\\
0.764	-0.0246334\\
0.766	-0.024635\\
0.768	-0.0246365\\
0.77	-0.0246379\\
0.772	-0.0246391\\
0.774	-0.0246402\\
0.776	-0.0246412\\
0.778	-0.024642\\
0.78	-0.0246428\\
0.782	-0.0246434\\
0.784	-0.0246439\\
0.786	-0.0246442\\
0.788	-0.0246445\\
0.79	-0.0246447\\
0.792	-0.0246448\\
0.794	-0.0246448\\
0.796	-0.0246447\\
0.798	-0.0246445\\
0.8	-0.0246443\\
0.802	-0.0246439\\
0.804	-0.0246436\\
0.806	-0.0246431\\
0.808	-0.0246426\\
0.81	-0.0246421\\
0.812	-0.0246415\\
0.814	-0.0246409\\
0.816	-0.0246402\\
0.818	-0.0246395\\
0.82	-0.0246388\\
0.822	-0.024638\\
0.824	-0.0246373\\
0.826	-0.0246365\\
0.828	-0.0246357\\
0.83	-0.0246349\\
0.832	-0.0246341\\
0.834	-0.0246334\\
0.836	-0.0246326\\
0.838	-0.0246318\\
0.84	-0.0246311\\
0.842	-0.0246304\\
0.844	-0.0246297\\
0.846	-0.024629\\
0.848	-0.0246284\\
0.85	-0.0246278\\
0.852	-0.0246273\\
0.854	-0.0246268\\
0.856	-0.0246263\\
0.858	-0.0246259\\
0.86	-0.0246255\\
0.862	-0.0246252\\
0.864	-0.0246249\\
0.866	-0.0246248\\
0.868	-0.0246246\\
0.87	-0.0246246\\
0.872	-0.0246246\\
0.874	-0.0246246\\
0.876	-0.0246248\\
0.878	-0.024625\\
0.88	-0.0246252\\
0.882	-0.0246256\\
0.884	-0.024626\\
0.886	-0.0246265\\
0.888	-0.0246271\\
0.89	-0.0246278\\
0.892	-0.0246285\\
0.894	-0.0246293\\
0.896	-0.0246302\\
0.898	-0.0246311\\
0.9	-0.0246321\\
0.902	-0.0246332\\
0.904	-0.0246344\\
0.906	-0.0246356\\
0.908	-0.024637\\
0.91	-0.0246383\\
0.912	-0.0246398\\
0.914	-0.0246413\\
0.916	-0.0246429\\
0.918	-0.0246445\\
0.92	-0.0246462\\
0.922	-0.0246479\\
0.924	-0.0246497\\
0.926	-0.0246516\\
0.928	-0.0246535\\
0.93	-0.0246554\\
0.932	-0.0246574\\
0.934	-0.0246594\\
0.936	-0.0246614\\
0.938	-0.0246635\\
0.94	-0.0246656\\
0.942	-0.0246677\\
0.944	-0.0246699\\
0.946	-0.024672\\
0.948	-0.0246742\\
0.95	-0.0246763\\
0.952	-0.0246785\\
0.954	-0.0246807\\
0.956	-0.0246829\\
0.958	-0.024685\\
0.96	-0.0246872\\
0.962	-0.0246893\\
0.964	-0.0246914\\
0.966	-0.0246935\\
0.968	-0.0246956\\
0.97	-0.0246977\\
0.972	-0.0246997\\
0.974	-0.0247016\\
0.976	-0.0247036\\
0.978	-0.0247055\\
0.98	-0.0247073\\
0.982	-0.0247091\\
0.984	-0.0247109\\
0.986	-0.0247126\\
0.988	-0.0247143\\
0.99	-0.0247159\\
0.992	-0.0247174\\
0.994	-0.0247189\\
0.996	-0.0247203\\
0.998	-0.0247217\\
1	-0.024723\\
1.002	-0.0247242\\
1.004	-0.0247254\\
1.006	-0.0247265\\
1.008	-0.0247275\\
1.01	-0.0247285\\
1.012	-0.0247294\\
1.014	-0.0247303\\
1.016	-0.0247311\\
1.018	-0.0247318\\
1.02	-0.0247324\\
1.022	-0.024733\\
1.024	-0.0247335\\
1.026	-0.024734\\
1.028	-0.0247344\\
1.03	-0.0247348\\
1.032	-0.024735\\
1.034	-0.0247353\\
1.036	-0.0247354\\
1.038	-0.0247356\\
1.04	-0.0247356\\
1.042	-0.0247357\\
1.044	-0.0247356\\
1.046	-0.0247356\\
1.048	-0.0247355\\
1.05	-0.0247353\\
1.052	-0.0247351\\
1.054	-0.0247349\\
1.056	-0.0247346\\
1.058	-0.0247344\\
1.06	-0.024734\\
1.062	-0.0247337\\
1.064	-0.0247333\\
1.066	-0.0247329\\
1.068	-0.0247325\\
1.07	-0.0247321\\
1.072	-0.0247317\\
1.074	-0.0247312\\
1.076	-0.0247308\\
1.078	-0.0247303\\
1.08	-0.0247298\\
1.082	-0.0247294\\
1.084	-0.0247289\\
1.086	-0.0247284\\
1.088	-0.024728\\
1.09	-0.0247275\\
1.092	-0.0247271\\
1.094	-0.0247266\\
1.096	-0.0247262\\
1.098	-0.0247258\\
1.1	-0.0247254\\
1.102	-0.024725\\
1.104	-0.0247246\\
1.106	-0.0247243\\
1.108	-0.024724\\
1.11	-0.0247237\\
1.112	-0.0247234\\
1.114	-0.0247231\\
1.116	-0.0247229\\
1.118	-0.0247227\\
1.12	-0.0247225\\
1.122	-0.0247224\\
1.124	-0.0247222\\
1.126	-0.0247221\\
1.128	-0.0247221\\
1.13	-0.024722\\
1.132	-0.024722\\
1.134	-0.024722\\
1.136	-0.0247221\\
1.138	-0.0247222\\
1.14	-0.0247223\\
1.142	-0.0247224\\
1.144	-0.0247226\\
1.146	-0.0247228\\
1.148	-0.024723\\
1.15	-0.0247232\\
1.152	-0.0247235\\
1.154	-0.0247238\\
1.156	-0.0247241\\
1.158	-0.0247245\\
1.16	-0.0247249\\
1.162	-0.0247253\\
1.164	-0.0247257\\
1.166	-0.0247261\\
1.168	-0.0247266\\
1.17	-0.0247271\\
1.172	-0.0247276\\
1.174	-0.0247281\\
1.176	-0.0247286\\
1.178	-0.0247291\\
1.18	-0.0247297\\
1.182	-0.0247303\\
1.184	-0.0247309\\
1.186	-0.0247315\\
1.188	-0.0247321\\
1.19	-0.0247327\\
1.192	-0.0247333\\
1.194	-0.0247339\\
1.196	-0.0247345\\
1.198	-0.0247352\\
1.2	-0.0247358\\
1.202	-0.0247364\\
1.204	-0.0247371\\
1.206	-0.0247377\\
1.208	-0.0247383\\
1.21	-0.0247389\\
1.212	-0.0247396\\
1.214	-0.0247402\\
1.216	-0.0247408\\
1.218	-0.0247414\\
1.22	-0.024742\\
1.222	-0.0247425\\
1.224	-0.0247431\\
1.226	-0.0247437\\
1.228	-0.0247442\\
1.23	-0.0247447\\
1.232	-0.0247453\\
1.234	-0.0247458\\
1.236	-0.0247463\\
1.238	-0.0247467\\
1.24	-0.0247472\\
1.242	-0.0247476\\
1.244	-0.0247481\\
1.246	-0.0247485\\
1.248	-0.0247488\\
1.25	-0.0247492\\
1.252	-0.0247496\\
1.254	-0.0247499\\
1.256	-0.0247502\\
1.258	-0.0247505\\
1.26	-0.0247508\\
1.262	-0.0247511\\
1.264	-0.0247513\\
1.266	-0.0247515\\
1.268	-0.0247517\\
1.27	-0.0247519\\
1.272	-0.0247521\\
1.274	-0.0247522\\
1.276	-0.0247524\\
1.278	-0.0247525\\
1.28	-0.0247526\\
1.282	-0.0247526\\
1.284	-0.0247527\\
1.286	-0.0247528\\
1.288	-0.0247528\\
1.29	-0.0247528\\
1.292	-0.0247528\\
1.294	-0.0247528\\
1.296	-0.0247528\\
1.298	-0.0247527\\
1.3	-0.0247527\\
1.302	-0.0247526\\
1.304	-0.0247526\\
1.306	-0.0247525\\
1.308	-0.0247524\\
1.31	-0.0247523\\
1.312	-0.0247522\\
1.314	-0.0247521\\
1.316	-0.024752\\
1.318	-0.0247519\\
1.32	-0.0247518\\
1.322	-0.0247516\\
1.324	-0.0247515\\
1.326	-0.0247514\\
1.328	-0.0247512\\
1.33	-0.0247511\\
1.332	-0.024751\\
1.334	-0.0247509\\
1.336	-0.0247507\\
1.338	-0.0247506\\
1.34	-0.0247505\\
1.342	-0.0247504\\
1.344	-0.0247503\\
1.346	-0.0247501\\
1.348	-0.02475\\
1.35	-0.0247499\\
1.352	-0.0247499\\
1.354	-0.0247498\\
1.356	-0.0247497\\
1.358	-0.0247496\\
1.36	-0.0247496\\
1.362	-0.0247495\\
1.364	-0.0247495\\
1.366	-0.0247495\\
1.368	-0.0247495\\
1.37	-0.0247495\\
1.372	-0.0247495\\
1.374	-0.0247495\\
1.376	-0.0247495\\
1.378	-0.0247495\\
1.38	-0.0247496\\
1.382	-0.0247496\\
1.384	-0.0247497\\
1.386	-0.0247498\\
1.388	-0.0247499\\
1.39	-0.02475\\
1.392	-0.0247501\\
1.394	-0.0247503\\
1.396	-0.0247504\\
1.398	-0.0247506\\
1.4	-0.0247507\\
1.402	-0.0247509\\
1.404	-0.0247511\\
1.406	-0.0247513\\
1.408	-0.0247515\\
1.41	-0.0247517\\
1.412	-0.0247519\\
1.414	-0.0247521\\
1.416	-0.0247524\\
1.418	-0.0247526\\
1.42	-0.0247529\\
1.422	-0.0247531\\
1.424	-0.0247534\\
1.426	-0.0247537\\
1.428	-0.0247539\\
1.43	-0.0247542\\
1.432	-0.0247545\\
1.434	-0.0247548\\
1.436	-0.0247551\\
1.438	-0.0247554\\
1.44	-0.0247557\\
1.442	-0.024756\\
1.444	-0.0247563\\
1.446	-0.0247566\\
1.448	-0.0247569\\
1.45	-0.0247572\\
1.452	-0.0247575\\
1.454	-0.0247578\\
1.456	-0.0247581\\
1.458	-0.0247583\\
1.46	-0.0247586\\
1.462	-0.0247589\\
1.464	-0.0247592\\
1.466	-0.0247595\\
1.468	-0.0247598\\
1.47	-0.02476\\
1.472	-0.0247603\\
1.474	-0.0247605\\
1.476	-0.0247608\\
1.478	-0.024761\\
1.48	-0.0247613\\
1.482	-0.0247615\\
1.484	-0.0247617\\
1.486	-0.0247619\\
1.488	-0.0247621\\
1.49	-0.0247623\\
1.492	-0.0247625\\
1.494	-0.0247627\\
1.496	-0.0247629\\
1.498	-0.0247631\\
1.5	-0.0247632\\
1.502	-0.0247634\\
1.504	-0.0247635\\
1.506	-0.0247636\\
1.508	-0.0247638\\
1.51	-0.0247639\\
1.512	-0.024764\\
1.514	-0.0247641\\
1.516	-0.0247642\\
1.518	-0.0247643\\
1.52	-0.0247643\\
1.522	-0.0247644\\
1.524	-0.0247645\\
1.526	-0.0247645\\
1.528	-0.0247646\\
1.53	-0.0247646\\
1.532	-0.0247646\\
1.534	-0.0247647\\
1.536	-0.0247647\\
1.538	-0.0247647\\
1.54	-0.0247647\\
1.542	-0.0247647\\
1.544	-0.0247647\\
1.546	-0.0247647\\
1.548	-0.0247647\\
1.55	-0.0247647\\
1.552	-0.0247647\\
1.554	-0.0247647\\
1.556	-0.0247647\\
1.558	-0.0247646\\
1.56	-0.0247646\\
1.562	-0.0247646\\
1.564	-0.0247646\\
1.566	-0.0247646\\
1.568	-0.0247645\\
1.57	-0.0247645\\
1.572	-0.0247645\\
1.574	-0.0247644\\
1.576	-0.0247644\\
1.578	-0.0247644\\
1.58	-0.0247644\\
1.582	-0.0247644\\
1.584	-0.0247643\\
1.586	-0.0247643\\
1.588	-0.0247643\\
1.59	-0.0247643\\
1.592	-0.0247643\\
1.594	-0.0247643\\
1.596	-0.0247643\\
1.598	-0.0247643\\
1.6	-0.0247643\\
1.602	-0.0247643\\
1.604	-0.0247643\\
1.606	-0.0247643\\
1.608	-0.0247643\\
1.61	-0.0247643\\
1.612	-0.0247644\\
1.614	-0.0247644\\
1.616	-0.0247644\\
1.618	-0.0247645\\
1.62	-0.0247645\\
1.622	-0.0247645\\
1.624	-0.0247646\\
1.626	-0.0247646\\
1.628	-0.0247647\\
1.63	-0.0247648\\
1.632	-0.0247648\\
1.634	-0.0247649\\
1.636	-0.024765\\
1.638	-0.024765\\
1.64	-0.0247651\\
1.642	-0.0247652\\
1.644	-0.0247653\\
1.646	-0.0247653\\
1.648	-0.0247654\\
1.65	-0.0247655\\
1.652	-0.0247656\\
1.654	-0.0247657\\
1.656	-0.0247658\\
1.658	-0.0247659\\
1.66	-0.024766\\
1.662	-0.0247661\\
1.664	-0.0247662\\
1.666	-0.0247663\\
1.668	-0.0247664\\
1.67	-0.0247665\\
1.672	-0.0247666\\
1.674	-0.0247667\\
1.676	-0.0247668\\
1.678	-0.0247669\\
1.68	-0.024767\\
1.682	-0.0247671\\
1.684	-0.0247672\\
1.686	-0.0247673\\
1.688	-0.0247674\\
1.69	-0.0247675\\
1.692	-0.0247676\\
1.694	-0.0247677\\
1.696	-0.0247678\\
1.698	-0.0247679\\
1.7	-0.024768\\
1.702	-0.024768\\
1.704	-0.0247681\\
1.706	-0.0247682\\
1.708	-0.0247683\\
1.71	-0.0247683\\
1.712	-0.0247684\\
1.714	-0.0247685\\
1.716	-0.0247685\\
1.718	-0.0247686\\
1.72	-0.0247687\\
1.722	-0.0247687\\
1.724	-0.0247688\\
1.726	-0.0247688\\
1.728	-0.0247689\\
1.73	-0.0247689\\
1.732	-0.024769\\
1.734	-0.024769\\
1.736	-0.024769\\
1.738	-0.0247691\\
1.74	-0.0247691\\
1.742	-0.0247691\\
1.744	-0.0247691\\
1.746	-0.0247691\\
1.748	-0.0247692\\
1.75	-0.0247692\\
1.752	-0.0247692\\
1.754	-0.0247692\\
1.756	-0.0247692\\
1.758	-0.0247692\\
1.76	-0.0247692\\
1.762	-0.0247692\\
1.764	-0.0247692\\
1.766	-0.0247692\\
1.768	-0.0247692\\
1.77	-0.0247692\\
1.772	-0.0247692\\
1.774	-0.0247692\\
1.776	-0.0247692\\
1.778	-0.0247691\\
1.78	-0.0247691\\
1.782	-0.0247691\\
1.784	-0.0247691\\
1.786	-0.0247691\\
1.788	-0.0247691\\
1.79	-0.0247691\\
1.792	-0.024769\\
1.794	-0.024769\\
1.796	-0.024769\\
1.798	-0.024769\\
1.8	-0.024769\\
1.802	-0.024769\\
1.804	-0.024769\\
1.806	-0.024769\\
1.808	-0.0247689\\
1.81	-0.0247689\\
1.812	-0.0247689\\
1.814	-0.0247689\\
1.816	-0.0247689\\
1.818	-0.0247689\\
1.82	-0.0247689\\
1.822	-0.0247689\\
1.824	-0.0247689\\
1.826	-0.0247689\\
1.828	-0.0247689\\
1.83	-0.0247689\\
1.832	-0.0247689\\
1.834	-0.0247689\\
1.836	-0.024769\\
1.838	-0.024769\\
1.84	-0.024769\\
1.842	-0.024769\\
1.844	-0.024769\\
1.846	-0.024769\\
1.848	-0.0247691\\
1.85	-0.0247691\\
1.852	-0.0247691\\
1.854	-0.0247691\\
1.856	-0.0247692\\
1.858	-0.0247692\\
1.86	-0.0247692\\
1.862	-0.0247693\\
1.864	-0.0247693\\
1.866	-0.0247693\\
1.868	-0.0247694\\
1.87	-0.0247694\\
1.872	-0.0247695\\
1.874	-0.0247695\\
1.876	-0.0247695\\
1.878	-0.0247696\\
1.88	-0.0247696\\
1.882	-0.0247697\\
1.884	-0.0247697\\
1.886	-0.0247698\\
1.888	-0.0247698\\
1.89	-0.0247698\\
1.892	-0.0247699\\
1.894	-0.0247699\\
1.896	-0.02477\\
1.898	-0.02477\\
1.9	-0.0247701\\
1.902	-0.0247701\\
1.904	-0.0247702\\
1.906	-0.0247702\\
1.908	-0.0247703\\
1.91	-0.0247703\\
1.912	-0.0247703\\
1.914	-0.0247704\\
1.916	-0.0247704\\
1.918	-0.0247705\\
1.92	-0.0247705\\
1.922	-0.0247706\\
1.924	-0.0247706\\
1.926	-0.0247706\\
1.928	-0.0247707\\
1.93	-0.0247707\\
1.932	-0.0247707\\
1.934	-0.0247708\\
1.936	-0.0247708\\
1.938	-0.0247708\\
1.94	-0.0247708\\
1.942	-0.0247709\\
1.944	-0.0247709\\
1.946	-0.0247709\\
1.948	-0.0247709\\
1.95	-0.024771\\
1.952	-0.024771\\
1.954	-0.024771\\
1.956	-0.024771\\
1.958	-0.024771\\
1.96	-0.024771\\
1.962	-0.024771\\
1.964	-0.0247711\\
1.966	-0.0247711\\
1.968	-0.0247711\\
1.97	-0.0247711\\
1.972	-0.0247711\\
1.974	-0.0247711\\
1.976	-0.0247711\\
1.978	-0.0247711\\
1.98	-0.0247711\\
1.982	-0.0247711\\
1.984	-0.0247711\\
1.986	-0.0247711\\
1.988	-0.0247711\\
1.99	-0.0247711\\
1.992	-0.0247711\\
1.994	-0.0247711\\
1.996	-0.0247711\\
1.998	-0.0247711\\
2	-0.0247711\\
};
\addplot [color=mycolor4, line width=1.0pt, forget plot]
  table[row sep=crcr]{%
-0.024	0\\
-0.022	0\\
-0.02	0\\
-0.018	0\\
-0.016	0\\
-0.014	0\\
-0.012	0\\
-0.01	0\\
-0.008	0\\
-0.006	0\\
-0.004	0\\
-0.002	0\\
0	0\\
0.002	-0.00184649\\
0.004	-0.00342632\\
0.006	-0.00478673\\
0.008	-0.00596703\\
0.01	-0.00699977\\
0.012	-0.00791183\\
0.014	-0.00872529\\
0.016	-0.00945821\\
0.018	-0.0101253\\
0.02	-0.0107384\\
0.022	-0.011307\\
0.024	-0.0118388\\
0.026	-0.0123398\\
0.028	-0.0128146\\
0.03	-0.0132668\\
0.032	-0.0136991\\
0.034	-0.0141136\\
0.036	-0.0145118\\
0.038	-0.0148946\\
0.04	-0.0152629\\
0.042	-0.0156172\\
0.044	-0.0159577\\
0.046	-0.0162847\\
0.048	-0.0165984\\
0.05	-0.0168987\\
0.052	-0.017186\\
0.054	-0.0174603\\
0.056	-0.0177217\\
0.058	-0.0179705\\
0.06	-0.018207\\
0.062	-0.0184313\\
0.064	-0.0186439\\
0.066	-0.0188452\\
0.068	-0.0190354\\
0.07	-0.0192151\\
0.072	-0.0193848\\
0.074	-0.0195448\\
0.076	-0.0196957\\
0.078	-0.019838\\
0.08	-0.0199722\\
0.082	-0.0200987\\
0.084	-0.0202181\\
0.086	-0.0203308\\
0.088	-0.0204373\\
0.09	-0.0205381\\
0.092	-0.0206335\\
0.094	-0.020724\\
0.096	-0.02081\\
0.098	-0.0208918\\
0.1	-0.0209697\\
0.102	-0.0210441\\
0.104	-0.0211153\\
0.106	-0.0211836\\
0.108	-0.0212491\\
0.11	-0.0213121\\
0.112	-0.0213727\\
0.114	-0.0214312\\
0.116	-0.0214877\\
0.118	-0.0215423\\
0.12	-0.021595\\
0.122	-0.021646\\
0.124	-0.0216952\\
0.126	-0.0217428\\
0.128	-0.0217888\\
0.13	-0.021833\\
0.132	-0.0218756\\
0.134	-0.0219165\\
0.136	-0.0219556\\
0.138	-0.0219929\\
0.14	-0.0220284\\
0.142	-0.0220619\\
0.144	-0.0220935\\
0.146	-0.022123\\
0.148	-0.0221505\\
0.15	-0.0221758\\
0.152	-0.022199\\
0.154	-0.0222199\\
0.156	-0.0222386\\
0.158	-0.0222551\\
0.16	-0.0222693\\
0.162	-0.0222813\\
0.164	-0.022291\\
0.166	-0.0222986\\
0.168	-0.022304\\
0.17	-0.0223073\\
0.172	-0.0223086\\
0.174	-0.022308\\
0.176	-0.0223056\\
0.178	-0.0223015\\
0.18	-0.0222958\\
0.182	-0.0222886\\
0.184	-0.0222801\\
0.186	-0.0222705\\
0.188	-0.0222598\\
0.19	-0.0222483\\
0.192	-0.0222361\\
0.194	-0.0222234\\
0.196	-0.0222103\\
0.198	-0.022197\\
0.2	-0.0221837\\
0.202	-0.0221705\\
0.204	-0.0221576\\
0.206	-0.0221451\\
0.208	-0.0221332\\
0.21	-0.0221221\\
0.212	-0.0221118\\
0.214	-0.0221025\\
0.216	-0.0220943\\
0.218	-0.0220873\\
0.22	-0.0220816\\
0.222	-0.0220773\\
0.224	-0.0220744\\
0.226	-0.0220731\\
0.228	-0.0220734\\
0.23	-0.0220753\\
0.232	-0.0220788\\
0.234	-0.022084\\
0.236	-0.022091\\
0.238	-0.0220996\\
0.24	-0.02211\\
0.242	-0.022122\\
0.244	-0.0221357\\
0.246	-0.0221511\\
0.248	-0.022168\\
0.25	-0.0221866\\
0.252	-0.0222066\\
0.254	-0.0222281\\
0.256	-0.022251\\
0.258	-0.0222752\\
0.26	-0.0223007\\
0.262	-0.0223274\\
0.264	-0.0223553\\
0.266	-0.0223842\\
0.268	-0.022414\\
0.27	-0.0224448\\
0.272	-0.0224764\\
0.274	-0.0225087\\
0.276	-0.0225418\\
0.278	-0.0225754\\
0.28	-0.0226095\\
0.282	-0.022644\\
0.284	-0.0226789\\
0.286	-0.0227142\\
0.288	-0.0227496\\
0.29	-0.0227852\\
0.292	-0.0228208\\
0.294	-0.0228565\\
0.296	-0.0228922\\
0.298	-0.0229278\\
0.3	-0.0229632\\
0.302	-0.0229984\\
0.304	-0.0230333\\
0.306	-0.023068\\
0.308	-0.0231023\\
0.31	-0.0231362\\
0.312	-0.0231697\\
0.314	-0.0232027\\
0.316	-0.0232353\\
0.318	-0.0232672\\
0.32	-0.0232987\\
0.322	-0.0233295\\
0.324	-0.0233597\\
0.326	-0.0233893\\
0.328	-0.0234182\\
0.33	-0.0234465\\
0.332	-0.023474\\
0.334	-0.0235008\\
0.336	-0.0235269\\
0.338	-0.0235523\\
0.34	-0.0235769\\
0.342	-0.0236008\\
0.344	-0.0236239\\
0.346	-0.0236462\\
0.348	-0.0236678\\
0.35	-0.0236885\\
0.352	-0.0237085\\
0.354	-0.0237278\\
0.356	-0.0237462\\
0.358	-0.0237639\\
0.36	-0.0237808\\
0.362	-0.0237969\\
0.364	-0.0238123\\
0.366	-0.023827\\
0.368	-0.0238409\\
0.37	-0.0238541\\
0.372	-0.0238666\\
0.374	-0.0238783\\
0.376	-0.0238894\\
0.378	-0.0238999\\
0.38	-0.0239097\\
0.382	-0.0239189\\
0.384	-0.0239274\\
0.386	-0.0239354\\
0.388	-0.0239428\\
0.39	-0.0239497\\
0.392	-0.0239561\\
0.394	-0.023962\\
0.396	-0.0239674\\
0.398	-0.0239724\\
0.4	-0.023977\\
0.402	-0.0239812\\
0.404	-0.023985\\
0.406	-0.0239886\\
0.408	-0.0239918\\
0.41	-0.0239948\\
0.412	-0.0239976\\
0.414	-0.0240002\\
0.416	-0.0240026\\
0.418	-0.0240048\\
0.42	-0.0240069\\
0.422	-0.024009\\
0.424	-0.024011\\
0.426	-0.024013\\
0.428	-0.024015\\
0.43	-0.024017\\
0.432	-0.024019\\
0.434	-0.0240211\\
0.436	-0.0240234\\
0.438	-0.0240257\\
0.44	-0.0240282\\
0.442	-0.0240308\\
0.444	-0.0240336\\
0.446	-0.0240366\\
0.448	-0.0240398\\
0.45	-0.0240432\\
0.452	-0.0240468\\
0.454	-0.0240506\\
0.456	-0.0240547\\
0.458	-0.0240591\\
0.46	-0.0240637\\
0.462	-0.0240685\\
0.464	-0.0240737\\
0.466	-0.024079\\
0.468	-0.0240846\\
0.47	-0.0240905\\
0.472	-0.0240966\\
0.474	-0.0241029\\
0.476	-0.0241095\\
0.478	-0.0241163\\
0.48	-0.0241233\\
0.482	-0.0241305\\
0.484	-0.0241378\\
0.486	-0.0241454\\
0.488	-0.0241531\\
0.49	-0.0241609\\
0.492	-0.0241689\\
0.494	-0.0241769\\
0.496	-0.0241851\\
0.498	-0.0241933\\
0.5	-0.0242015\\
0.502	-0.0242098\\
0.504	-0.0242181\\
0.506	-0.0242264\\
0.508	-0.0242347\\
0.51	-0.024243\\
0.512	-0.0242511\\
0.514	-0.0242592\\
0.516	-0.0242672\\
0.518	-0.0242751\\
0.52	-0.0242829\\
0.522	-0.0242905\\
0.524	-0.024298\\
0.526	-0.0243052\\
0.528	-0.0243123\\
0.53	-0.0243193\\
0.532	-0.0243259\\
0.534	-0.0243324\\
0.536	-0.0243387\\
0.538	-0.0243447\\
0.54	-0.0243504\\
0.542	-0.024356\\
0.544	-0.0243612\\
0.546	-0.0243663\\
0.548	-0.024371\\
0.55	-0.0243755\\
0.552	-0.0243798\\
0.554	-0.0243837\\
0.556	-0.0243875\\
0.558	-0.024391\\
0.56	-0.0243942\\
0.562	-0.0243972\\
0.564	-0.0244\\
0.566	-0.0244025\\
0.568	-0.0244049\\
0.57	-0.024407\\
0.572	-0.0244089\\
0.574	-0.0244107\\
0.576	-0.0244122\\
0.578	-0.0244136\\
0.58	-0.0244149\\
0.582	-0.024416\\
0.584	-0.024417\\
0.586	-0.0244179\\
0.588	-0.0244187\\
0.59	-0.0244195\\
0.592	-0.0244201\\
0.594	-0.0244207\\
0.596	-0.0244213\\
0.598	-0.0244219\\
0.6	-0.0244224\\
0.602	-0.0244229\\
0.604	-0.0244235\\
0.606	-0.0244241\\
0.608	-0.0244247\\
0.61	-0.0244253\\
0.612	-0.0244261\\
0.614	-0.0244269\\
0.616	-0.0244278\\
0.618	-0.0244287\\
0.62	-0.0244298\\
0.622	-0.024431\\
0.624	-0.0244322\\
0.626	-0.0244336\\
0.628	-0.0244351\\
0.63	-0.0244368\\
0.632	-0.0244385\\
0.634	-0.0244404\\
0.636	-0.0244425\\
0.638	-0.0244446\\
0.64	-0.0244469\\
0.642	-0.0244493\\
0.644	-0.0244518\\
0.646	-0.0244544\\
0.648	-0.0244572\\
0.65	-0.0244601\\
0.652	-0.0244631\\
0.654	-0.0244662\\
0.656	-0.0244693\\
0.658	-0.0244726\\
0.66	-0.024476\\
0.662	-0.0244794\\
0.664	-0.0244829\\
0.666	-0.0244864\\
0.668	-0.02449\\
0.67	-0.0244937\\
0.672	-0.0244974\\
0.674	-0.024501\\
0.676	-0.0245048\\
0.678	-0.0245085\\
0.68	-0.0245122\\
0.682	-0.0245159\\
0.684	-0.0245195\\
0.686	-0.0245232\\
0.688	-0.0245268\\
0.69	-0.0245303\\
0.692	-0.0245338\\
0.694	-0.0245373\\
0.696	-0.0245406\\
0.698	-0.0245439\\
0.7	-0.0245471\\
0.702	-0.0245503\\
0.704	-0.0245533\\
0.706	-0.0245562\\
0.708	-0.0245591\\
0.71	-0.0245618\\
0.712	-0.0245644\\
0.714	-0.0245669\\
0.716	-0.0245694\\
0.718	-0.0245717\\
0.72	-0.0245738\\
0.722	-0.0245759\\
0.724	-0.0245779\\
0.726	-0.0245797\\
0.728	-0.0245815\\
0.73	-0.0245831\\
0.732	-0.0245847\\
0.734	-0.0245861\\
0.736	-0.0245874\\
0.738	-0.0245887\\
0.74	-0.0245898\\
0.742	-0.0245909\\
0.744	-0.0245919\\
0.746	-0.0245928\\
0.748	-0.0245936\\
0.75	-0.0245943\\
0.752	-0.024595\\
0.754	-0.0245957\\
0.756	-0.0245963\\
0.758	-0.0245968\\
0.76	-0.0245973\\
0.762	-0.0245978\\
0.764	-0.0245982\\
0.766	-0.0245986\\
0.768	-0.024599\\
0.77	-0.0245994\\
0.772	-0.0245998\\
0.774	-0.0246002\\
0.776	-0.0246005\\
0.778	-0.0246009\\
0.78	-0.0246013\\
0.782	-0.0246017\\
0.784	-0.0246021\\
0.786	-0.0246025\\
0.788	-0.024603\\
0.79	-0.0246035\\
0.792	-0.024604\\
0.794	-0.0246046\\
0.796	-0.0246052\\
0.798	-0.0246058\\
0.8	-0.0246064\\
0.802	-0.0246071\\
0.804	-0.0246079\\
0.806	-0.0246087\\
0.808	-0.0246095\\
0.81	-0.0246103\\
0.812	-0.0246112\\
0.814	-0.0246122\\
0.816	-0.0246132\\
0.818	-0.0246142\\
0.82	-0.0246153\\
0.822	-0.0246164\\
0.824	-0.0246175\\
0.826	-0.0246187\\
0.828	-0.0246199\\
0.83	-0.0246211\\
0.832	-0.0246224\\
0.834	-0.0246237\\
0.836	-0.024625\\
0.838	-0.0246263\\
0.84	-0.0246277\\
0.842	-0.0246291\\
0.844	-0.0246305\\
0.846	-0.0246319\\
0.848	-0.0246333\\
0.85	-0.0246347\\
0.852	-0.0246362\\
0.854	-0.0246376\\
0.856	-0.024639\\
0.858	-0.0246405\\
0.86	-0.0246419\\
0.862	-0.0246433\\
0.864	-0.0246448\\
0.866	-0.0246462\\
0.868	-0.0246476\\
0.87	-0.024649\\
0.872	-0.0246503\\
0.874	-0.0246517\\
0.876	-0.024653\\
0.878	-0.0246544\\
0.88	-0.0246557\\
0.882	-0.024657\\
0.884	-0.0246582\\
0.886	-0.0246594\\
0.888	-0.0246607\\
0.89	-0.0246618\\
0.892	-0.024663\\
0.894	-0.0246641\\
0.896	-0.0246652\\
0.898	-0.0246663\\
0.9	-0.0246674\\
0.902	-0.0246684\\
0.904	-0.0246694\\
0.906	-0.0246703\\
0.908	-0.0246713\\
0.91	-0.0246722\\
0.912	-0.0246731\\
0.914	-0.0246739\\
0.916	-0.0246747\\
0.918	-0.0246755\\
0.92	-0.0246763\\
0.922	-0.024677\\
0.924	-0.0246778\\
0.926	-0.0246785\\
0.928	-0.0246791\\
0.93	-0.0246798\\
0.932	-0.0246804\\
0.934	-0.024681\\
0.936	-0.0246816\\
0.938	-0.0246822\\
0.94	-0.0246827\\
0.942	-0.0246833\\
0.944	-0.0246838\\
0.946	-0.0246843\\
0.948	-0.0246848\\
0.95	-0.0246852\\
0.952	-0.0246857\\
0.954	-0.0246861\\
0.956	-0.0246866\\
0.958	-0.024687\\
0.96	-0.0246874\\
0.962	-0.0246878\\
0.964	-0.0246882\\
0.966	-0.0246886\\
0.968	-0.024689\\
0.97	-0.0246893\\
0.972	-0.0246897\\
0.974	-0.0246901\\
0.976	-0.0246904\\
0.978	-0.0246908\\
0.98	-0.0246912\\
0.982	-0.0246915\\
0.984	-0.0246919\\
0.986	-0.0246923\\
0.988	-0.0246927\\
0.99	-0.024693\\
0.992	-0.0246934\\
0.994	-0.0246938\\
0.996	-0.0246942\\
0.998	-0.0246946\\
1	-0.024695\\
1.002	-0.0246954\\
1.004	-0.0246958\\
1.006	-0.0246962\\
1.008	-0.0246966\\
1.01	-0.0246971\\
1.012	-0.0246975\\
1.014	-0.024698\\
1.016	-0.0246984\\
1.018	-0.0246989\\
1.02	-0.0246994\\
1.022	-0.0246999\\
1.024	-0.0247004\\
1.026	-0.0247009\\
1.028	-0.0247014\\
1.03	-0.0247019\\
1.032	-0.0247025\\
1.034	-0.024703\\
1.036	-0.0247036\\
1.038	-0.0247041\\
1.04	-0.0247047\\
1.042	-0.0247053\\
1.044	-0.0247059\\
1.046	-0.0247065\\
1.048	-0.0247071\\
1.05	-0.0247077\\
1.052	-0.0247083\\
1.054	-0.0247089\\
1.056	-0.0247096\\
1.058	-0.0247102\\
1.06	-0.0247108\\
1.062	-0.0247115\\
1.064	-0.0247121\\
1.066	-0.0247128\\
1.068	-0.0247134\\
1.07	-0.0247141\\
1.072	-0.0247147\\
1.074	-0.0247154\\
1.076	-0.024716\\
1.078	-0.0247166\\
1.08	-0.0247173\\
1.082	-0.0247179\\
1.084	-0.0247185\\
1.086	-0.0247192\\
1.088	-0.0247198\\
1.09	-0.0247204\\
1.092	-0.024721\\
1.094	-0.0247216\\
1.096	-0.0247221\\
1.098	-0.0247227\\
1.1	-0.0247233\\
1.102	-0.0247238\\
1.104	-0.0247244\\
1.106	-0.0247249\\
1.108	-0.0247254\\
1.11	-0.0247259\\
1.112	-0.0247264\\
1.114	-0.0247268\\
1.116	-0.0247273\\
1.118	-0.0247277\\
1.12	-0.0247282\\
1.122	-0.0247286\\
1.124	-0.024729\\
1.126	-0.0247293\\
1.128	-0.0247297\\
1.13	-0.02473\\
1.132	-0.0247304\\
1.134	-0.0247307\\
1.136	-0.024731\\
1.138	-0.0247313\\
1.14	-0.0247315\\
1.142	-0.0247318\\
1.144	-0.024732\\
1.146	-0.0247323\\
1.148	-0.0247325\\
1.15	-0.0247327\\
1.152	-0.0247329\\
1.154	-0.024733\\
1.156	-0.0247332\\
1.158	-0.0247334\\
1.16	-0.0247335\\
1.162	-0.0247336\\
1.164	-0.0247337\\
1.166	-0.0247339\\
1.168	-0.024734\\
1.17	-0.0247341\\
1.172	-0.0247342\\
1.174	-0.0247342\\
1.176	-0.0247343\\
1.178	-0.0247344\\
1.18	-0.0247345\\
1.182	-0.0247345\\
1.184	-0.0247346\\
1.186	-0.0247347\\
1.188	-0.0247348\\
1.19	-0.0247348\\
1.192	-0.0247349\\
1.194	-0.024735\\
1.196	-0.0247351\\
1.198	-0.0247351\\
1.2	-0.0247352\\
1.202	-0.0247353\\
1.204	-0.0247354\\
1.206	-0.0247355\\
1.208	-0.0247356\\
1.21	-0.0247357\\
1.212	-0.0247358\\
1.214	-0.024736\\
1.216	-0.0247361\\
1.218	-0.0247363\\
1.22	-0.0247364\\
1.222	-0.0247366\\
1.224	-0.0247367\\
1.226	-0.0247369\\
1.228	-0.0247371\\
1.23	-0.0247373\\
1.232	-0.0247375\\
1.234	-0.0247377\\
1.236	-0.024738\\
1.238	-0.0247382\\
1.24	-0.0247384\\
1.242	-0.0247387\\
1.244	-0.024739\\
1.246	-0.0247392\\
1.248	-0.0247395\\
1.25	-0.0247398\\
1.252	-0.0247401\\
1.254	-0.0247404\\
1.256	-0.0247407\\
1.258	-0.024741\\
1.26	-0.0247414\\
1.262	-0.0247417\\
1.264	-0.024742\\
1.266	-0.0247424\\
1.268	-0.0247427\\
1.27	-0.0247431\\
1.272	-0.0247434\\
1.274	-0.0247438\\
1.276	-0.0247441\\
1.278	-0.0247445\\
1.28	-0.0247448\\
1.282	-0.0247452\\
1.284	-0.0247455\\
1.286	-0.0247459\\
1.288	-0.0247462\\
1.29	-0.0247466\\
1.292	-0.0247469\\
1.294	-0.0247473\\
1.296	-0.0247476\\
1.298	-0.024748\\
1.3	-0.0247483\\
1.302	-0.0247486\\
1.304	-0.0247489\\
1.306	-0.0247492\\
1.308	-0.0247495\\
1.31	-0.0247498\\
1.312	-0.0247501\\
1.314	-0.0247504\\
1.316	-0.0247507\\
1.318	-0.0247509\\
1.32	-0.0247512\\
1.322	-0.0247514\\
1.324	-0.0247516\\
1.326	-0.0247519\\
1.328	-0.0247521\\
1.33	-0.0247523\\
1.332	-0.0247525\\
1.334	-0.0247527\\
1.336	-0.0247528\\
1.338	-0.024753\\
1.34	-0.0247532\\
1.342	-0.0247533\\
1.344	-0.0247534\\
1.346	-0.0247536\\
1.348	-0.0247537\\
1.35	-0.0247538\\
1.352	-0.0247539\\
1.354	-0.024754\\
1.356	-0.024754\\
1.358	-0.0247541\\
1.36	-0.0247542\\
1.362	-0.0247542\\
1.364	-0.0247543\\
1.366	-0.0247543\\
1.368	-0.0247544\\
1.37	-0.0247544\\
1.372	-0.0247544\\
1.374	-0.0247544\\
1.376	-0.0247544\\
1.378	-0.0247545\\
1.38	-0.0247545\\
1.382	-0.0247545\\
1.384	-0.0247545\\
1.386	-0.0247545\\
1.388	-0.0247545\\
1.39	-0.0247545\\
1.392	-0.0247545\\
1.394	-0.0247545\\
1.396	-0.0247545\\
1.398	-0.0247545\\
1.4	-0.0247545\\
1.402	-0.0247545\\
1.404	-0.0247545\\
1.406	-0.0247545\\
1.408	-0.0247545\\
1.41	-0.0247545\\
1.412	-0.0247545\\
1.414	-0.0247545\\
1.416	-0.0247546\\
1.418	-0.0247546\\
1.42	-0.0247546\\
1.422	-0.0247547\\
1.424	-0.0247547\\
1.426	-0.0247547\\
1.428	-0.0247548\\
1.43	-0.0247549\\
1.432	-0.0247549\\
1.434	-0.024755\\
1.436	-0.0247551\\
1.438	-0.0247552\\
1.44	-0.0247553\\
1.442	-0.0247553\\
1.444	-0.0247555\\
1.446	-0.0247556\\
1.448	-0.0247557\\
1.45	-0.0247558\\
1.452	-0.0247559\\
1.454	-0.0247561\\
1.456	-0.0247562\\
1.458	-0.0247563\\
1.46	-0.0247565\\
1.462	-0.0247566\\
1.464	-0.0247568\\
1.466	-0.024757\\
1.468	-0.0247571\\
1.47	-0.0247573\\
1.472	-0.0247575\\
1.474	-0.0247576\\
1.476	-0.0247578\\
1.478	-0.024758\\
1.48	-0.0247582\\
1.482	-0.0247584\\
1.484	-0.0247586\\
1.486	-0.0247588\\
1.488	-0.024759\\
1.49	-0.0247591\\
1.492	-0.0247593\\
1.494	-0.0247595\\
1.496	-0.0247597\\
1.498	-0.0247599\\
1.5	-0.0247601\\
1.502	-0.0247603\\
1.504	-0.0247605\\
1.506	-0.0247607\\
1.508	-0.0247609\\
1.51	-0.024761\\
1.512	-0.0247612\\
1.514	-0.0247614\\
1.516	-0.0247616\\
1.518	-0.0247617\\
1.52	-0.0247619\\
1.522	-0.024762\\
1.524	-0.0247622\\
1.526	-0.0247624\\
1.528	-0.0247625\\
1.53	-0.0247626\\
1.532	-0.0247628\\
1.534	-0.0247629\\
1.536	-0.024763\\
1.538	-0.0247631\\
1.54	-0.0247633\\
1.542	-0.0247634\\
1.544	-0.0247635\\
1.546	-0.0247636\\
1.548	-0.0247637\\
1.55	-0.0247638\\
1.552	-0.0247638\\
1.554	-0.0247639\\
1.556	-0.024764\\
1.558	-0.024764\\
1.56	-0.0247641\\
1.562	-0.0247642\\
1.564	-0.0247642\\
1.566	-0.0247643\\
1.568	-0.0247643\\
1.57	-0.0247643\\
1.572	-0.0247644\\
1.574	-0.0247644\\
1.576	-0.0247644\\
1.578	-0.0247644\\
1.58	-0.0247644\\
1.582	-0.0247645\\
1.584	-0.0247645\\
1.586	-0.0247645\\
1.588	-0.0247645\\
1.59	-0.0247645\\
1.592	-0.0247645\\
1.594	-0.0247645\\
1.596	-0.0247645\\
1.598	-0.0247645\\
1.6	-0.0247644\\
1.602	-0.0247644\\
1.604	-0.0247644\\
1.606	-0.0247644\\
1.608	-0.0247644\\
1.61	-0.0247644\\
1.612	-0.0247644\\
1.614	-0.0247644\\
1.616	-0.0247644\\
1.618	-0.0247643\\
1.62	-0.0247643\\
1.622	-0.0247643\\
1.624	-0.0247643\\
1.626	-0.0247643\\
1.628	-0.0247643\\
1.63	-0.0247643\\
1.632	-0.0247643\\
1.634	-0.0247643\\
1.636	-0.0247643\\
1.638	-0.0247644\\
1.64	-0.0247644\\
1.642	-0.0247644\\
1.644	-0.0247644\\
1.646	-0.0247644\\
1.648	-0.0247644\\
1.65	-0.0247645\\
1.652	-0.0247645\\
1.654	-0.0247645\\
1.656	-0.0247646\\
1.658	-0.0247646\\
1.66	-0.0247647\\
1.662	-0.0247647\\
1.664	-0.0247648\\
1.666	-0.0247648\\
1.668	-0.0247649\\
1.67	-0.0247649\\
1.672	-0.024765\\
1.674	-0.0247651\\
1.676	-0.0247651\\
1.678	-0.0247652\\
1.68	-0.0247653\\
1.682	-0.0247654\\
1.684	-0.0247654\\
1.686	-0.0247655\\
1.688	-0.0247656\\
1.69	-0.0247657\\
1.692	-0.0247658\\
1.694	-0.0247659\\
1.696	-0.0247659\\
1.698	-0.024766\\
1.7	-0.0247661\\
1.702	-0.0247662\\
1.704	-0.0247663\\
1.706	-0.0247664\\
1.708	-0.0247665\\
1.71	-0.0247666\\
1.712	-0.0247667\\
1.714	-0.0247668\\
1.716	-0.0247669\\
1.718	-0.024767\\
1.72	-0.0247671\\
1.722	-0.0247672\\
1.724	-0.0247673\\
1.726	-0.0247674\\
1.728	-0.0247675\\
1.73	-0.0247675\\
1.732	-0.0247676\\
1.734	-0.0247677\\
1.736	-0.0247678\\
1.738	-0.0247679\\
1.74	-0.024768\\
1.742	-0.024768\\
1.744	-0.0247681\\
1.746	-0.0247682\\
1.748	-0.0247683\\
1.75	-0.0247683\\
1.752	-0.0247684\\
1.754	-0.0247685\\
1.756	-0.0247685\\
1.758	-0.0247686\\
1.76	-0.0247686\\
1.762	-0.0247687\\
1.764	-0.0247687\\
1.766	-0.0247688\\
1.768	-0.0247688\\
1.77	-0.0247689\\
1.772	-0.0247689\\
1.774	-0.0247689\\
1.776	-0.024769\\
1.778	-0.024769\\
1.78	-0.024769\\
1.782	-0.024769\\
1.784	-0.0247691\\
1.786	-0.0247691\\
1.788	-0.0247691\\
1.79	-0.0247691\\
1.792	-0.0247691\\
1.794	-0.0247691\\
1.796	-0.0247691\\
1.798	-0.0247691\\
1.8	-0.0247691\\
1.802	-0.0247691\\
1.804	-0.0247691\\
1.806	-0.0247691\\
1.808	-0.0247691\\
1.81	-0.0247691\\
1.812	-0.0247691\\
1.814	-0.0247691\\
1.816	-0.0247691\\
1.818	-0.0247691\\
1.82	-0.0247691\\
1.822	-0.0247691\\
1.824	-0.0247691\\
1.826	-0.0247691\\
1.828	-0.0247691\\
1.83	-0.0247691\\
1.832	-0.0247691\\
1.834	-0.024769\\
1.836	-0.024769\\
1.838	-0.024769\\
1.84	-0.024769\\
1.842	-0.024769\\
1.844	-0.024769\\
1.846	-0.024769\\
1.848	-0.024769\\
1.85	-0.024769\\
1.852	-0.024769\\
1.854	-0.024769\\
1.856	-0.024769\\
1.858	-0.024769\\
1.86	-0.024769\\
1.862	-0.024769\\
1.864	-0.024769\\
1.866	-0.024769\\
1.868	-0.024769\\
1.87	-0.0247691\\
1.872	-0.0247691\\
1.874	-0.0247691\\
1.876	-0.0247691\\
1.878	-0.0247691\\
1.88	-0.0247691\\
1.882	-0.0247692\\
1.884	-0.0247692\\
1.886	-0.0247692\\
1.888	-0.0247692\\
1.89	-0.0247693\\
1.892	-0.0247693\\
1.894	-0.0247693\\
1.896	-0.0247694\\
1.898	-0.0247694\\
1.9	-0.0247694\\
1.902	-0.0247695\\
1.904	-0.0247695\\
1.906	-0.0247696\\
1.908	-0.0247696\\
1.91	-0.0247696\\
1.912	-0.0247697\\
1.914	-0.0247697\\
1.916	-0.0247698\\
1.918	-0.0247698\\
1.92	-0.0247699\\
1.922	-0.0247699\\
1.924	-0.02477\\
1.926	-0.02477\\
1.928	-0.0247701\\
1.93	-0.0247701\\
1.932	-0.0247701\\
1.934	-0.0247702\\
1.936	-0.0247702\\
1.938	-0.0247703\\
1.94	-0.0247703\\
1.942	-0.0247704\\
1.944	-0.0247704\\
1.946	-0.0247705\\
1.948	-0.0247705\\
1.95	-0.0247705\\
1.952	-0.0247706\\
1.954	-0.0247706\\
1.956	-0.0247707\\
1.958	-0.0247707\\
1.96	-0.0247707\\
1.962	-0.0247708\\
1.964	-0.0247708\\
1.966	-0.0247708\\
1.968	-0.0247709\\
1.97	-0.0247709\\
1.972	-0.0247709\\
1.974	-0.024771\\
1.976	-0.024771\\
1.978	-0.024771\\
1.98	-0.024771\\
1.982	-0.0247711\\
1.984	-0.0247711\\
1.986	-0.0247711\\
1.988	-0.0247711\\
1.99	-0.0247711\\
1.992	-0.0247712\\
1.994	-0.0247712\\
1.996	-0.0247712\\
1.998	-0.0247712\\
2	-0.0247712\\
};
\addplot [color=mycolor5, line width=1.0pt, forget plot]
  table[row sep=crcr]{%
-0.024	0\\
-0.022	0\\
-0.02	0\\
-0.018	0\\
-0.016	0\\
-0.014	0\\
-0.012	0\\
-0.01	0\\
-0.008	0\\
-0.006	0\\
-0.004	0\\
-0.002	0\\
0	0\\
0.002	-0.00184649\\
0.004	-0.00342631\\
0.006	-0.00478673\\
0.008	-0.00596702\\
0.01	-0.00699977\\
0.012	-0.00791183\\
0.014	-0.00872528\\
0.016	-0.0094582\\
0.018	-0.0101253\\
0.02	-0.0107383\\
0.022	-0.011307\\
0.024	-0.0118388\\
0.026	-0.0123398\\
0.028	-0.0128146\\
0.03	-0.0132668\\
0.032	-0.0136992\\
0.034	-0.0141137\\
0.036	-0.0145119\\
0.038	-0.0148949\\
0.04	-0.0152633\\
0.042	-0.0156177\\
0.044	-0.0159584\\
0.046	-0.0162856\\
0.048	-0.0165994\\
0.05	-0.0169001\\
0.052	-0.0171876\\
0.054	-0.0174622\\
0.056	-0.0177241\\
0.058	-0.0179733\\
0.06	-0.0182103\\
0.062	-0.0184352\\
0.064	-0.0186484\\
0.066	-0.0188504\\
0.068	-0.0190414\\
0.07	-0.0192219\\
0.072	-0.0193925\\
0.074	-0.0195535\\
0.076	-0.0197055\\
0.078	-0.019849\\
0.08	-0.0199844\\
0.082	-0.0201123\\
0.084	-0.0202331\\
0.086	-0.0203474\\
0.088	-0.0204555\\
0.09	-0.020558\\
0.092	-0.0206552\\
0.094	-0.0207476\\
0.096	-0.0208356\\
0.098	-0.0209195\\
0.1	-0.0209996\\
0.102	-0.0210763\\
0.104	-0.0211498\\
0.106	-0.0212205\\
0.108	-0.0212885\\
0.11	-0.0213541\\
0.112	-0.0214174\\
0.114	-0.0214787\\
0.116	-0.0215379\\
0.118	-0.0215954\\
0.12	-0.021651\\
0.122	-0.0217049\\
0.124	-0.0217572\\
0.126	-0.0218078\\
0.128	-0.0218568\\
0.13	-0.0219041\\
0.132	-0.0219497\\
0.134	-0.0219936\\
0.136	-0.0220358\\
0.138	-0.0220761\\
0.14	-0.0221146\\
0.142	-0.0221511\\
0.144	-0.0221856\\
0.146	-0.022218\\
0.148	-0.0222483\\
0.15	-0.0222764\\
0.152	-0.0223022\\
0.154	-0.0223257\\
0.156	-0.0223469\\
0.158	-0.0223657\\
0.16	-0.0223821\\
0.162	-0.0223962\\
0.164	-0.0224079\\
0.166	-0.0224173\\
0.168	-0.0224244\\
0.17	-0.0224292\\
0.172	-0.0224319\\
0.174	-0.0224325\\
0.176	-0.022431\\
0.178	-0.0224277\\
0.18	-0.0224226\\
0.182	-0.0224159\\
0.184	-0.0224077\\
0.186	-0.0223981\\
0.188	-0.0223872\\
0.19	-0.0223753\\
0.192	-0.0223625\\
0.194	-0.022349\\
0.196	-0.0223349\\
0.198	-0.0223204\\
0.2	-0.0223056\\
0.202	-0.0222908\\
0.204	-0.0222761\\
0.206	-0.0222616\\
0.208	-0.0222474\\
0.21	-0.0222338\\
0.212	-0.0222209\\
0.214	-0.0222088\\
0.216	-0.0221976\\
0.218	-0.0221874\\
0.22	-0.0221784\\
0.222	-0.0221706\\
0.224	-0.0221641\\
0.226	-0.022159\\
0.228	-0.0221554\\
0.23	-0.0221532\\
0.232	-0.0221526\\
0.234	-0.0221536\\
0.236	-0.0221563\\
0.238	-0.0221605\\
0.24	-0.0221664\\
0.242	-0.0221739\\
0.244	-0.022183\\
0.246	-0.0221938\\
0.248	-0.0222061\\
0.25	-0.02222\\
0.252	-0.0222354\\
0.254	-0.0222522\\
0.256	-0.0222705\\
0.258	-0.0222902\\
0.26	-0.0223112\\
0.262	-0.0223335\\
0.264	-0.022357\\
0.266	-0.0223816\\
0.268	-0.0224072\\
0.27	-0.0224339\\
0.272	-0.0224615\\
0.274	-0.02249\\
0.276	-0.0225193\\
0.278	-0.0225493\\
0.28	-0.02258\\
0.282	-0.0226112\\
0.284	-0.022643\\
0.286	-0.0226752\\
0.288	-0.0227079\\
0.29	-0.0227408\\
0.292	-0.022774\\
0.294	-0.0228074\\
0.296	-0.022841\\
0.298	-0.0228747\\
0.3	-0.0229084\\
0.302	-0.022942\\
0.304	-0.0229757\\
0.306	-0.0230091\\
0.308	-0.0230425\\
0.31	-0.0230756\\
0.312	-0.0231085\\
0.314	-0.0231411\\
0.316	-0.0231733\\
0.318	-0.0232053\\
0.32	-0.0232368\\
0.322	-0.0232678\\
0.324	-0.0232984\\
0.326	-0.0233286\\
0.328	-0.0233582\\
0.33	-0.0233872\\
0.332	-0.0234157\\
0.334	-0.0234436\\
0.336	-0.0234708\\
0.338	-0.0234975\\
0.34	-0.0235235\\
0.342	-0.0235488\\
0.344	-0.0235734\\
0.346	-0.0235973\\
0.348	-0.0236205\\
0.35	-0.0236429\\
0.352	-0.0236646\\
0.354	-0.0236856\\
0.356	-0.0237058\\
0.358	-0.0237253\\
0.36	-0.023744\\
0.362	-0.0237619\\
0.364	-0.023779\\
0.366	-0.0237954\\
0.368	-0.023811\\
0.37	-0.0238258\\
0.372	-0.0238399\\
0.374	-0.0238532\\
0.376	-0.0238658\\
0.378	-0.0238776\\
0.38	-0.0238887\\
0.382	-0.0238991\\
0.384	-0.0239089\\
0.386	-0.0239179\\
0.388	-0.0239263\\
0.39	-0.023934\\
0.392	-0.0239411\\
0.394	-0.0239477\\
0.396	-0.0239536\\
0.398	-0.023959\\
0.4	-0.0239639\\
0.402	-0.0239684\\
0.404	-0.0239723\\
0.406	-0.0239758\\
0.408	-0.0239789\\
0.41	-0.0239817\\
0.412	-0.0239841\\
0.414	-0.0239862\\
0.416	-0.023988\\
0.418	-0.0239896\\
0.42	-0.023991\\
0.422	-0.0239922\\
0.424	-0.0239933\\
0.426	-0.0239942\\
0.428	-0.0239951\\
0.43	-0.0239959\\
0.432	-0.0239967\\
0.434	-0.0239975\\
0.436	-0.0239984\\
0.438	-0.0239993\\
0.44	-0.0240003\\
0.442	-0.0240015\\
0.444	-0.0240028\\
0.446	-0.0240043\\
0.448	-0.0240059\\
0.45	-0.0240078\\
0.452	-0.0240098\\
0.454	-0.0240122\\
0.456	-0.0240148\\
0.458	-0.0240176\\
0.46	-0.0240208\\
0.462	-0.0240242\\
0.464	-0.0240279\\
0.466	-0.0240319\\
0.468	-0.0240363\\
0.47	-0.0240409\\
0.472	-0.0240459\\
0.474	-0.0240512\\
0.476	-0.0240567\\
0.478	-0.0240626\\
0.48	-0.0240688\\
0.482	-0.0240752\\
0.484	-0.024082\\
0.486	-0.024089\\
0.488	-0.0240962\\
0.49	-0.0241037\\
0.492	-0.0241114\\
0.494	-0.0241194\\
0.496	-0.0241275\\
0.498	-0.0241358\\
0.5	-0.0241443\\
0.502	-0.0241529\\
0.504	-0.0241616\\
0.506	-0.0241705\\
0.508	-0.0241794\\
0.51	-0.0241884\\
0.512	-0.0241974\\
0.514	-0.0242065\\
0.516	-0.0242155\\
0.518	-0.0242245\\
0.52	-0.0242335\\
0.522	-0.0242425\\
0.524	-0.0242513\\
0.526	-0.0242601\\
0.528	-0.0242688\\
0.53	-0.0242773\\
0.532	-0.0242857\\
0.534	-0.0242939\\
0.536	-0.0243019\\
0.538	-0.0243098\\
0.54	-0.0243174\\
0.542	-0.0243249\\
0.544	-0.0243321\\
0.546	-0.0243391\\
0.548	-0.0243458\\
0.55	-0.0243523\\
0.552	-0.0243585\\
0.554	-0.0243645\\
0.556	-0.0243702\\
0.558	-0.0243757\\
0.56	-0.0243808\\
0.562	-0.0243857\\
0.564	-0.0243904\\
0.566	-0.0243948\\
0.568	-0.0243989\\
0.57	-0.0244027\\
0.572	-0.0244063\\
0.574	-0.0244097\\
0.576	-0.0244128\\
0.578	-0.0244157\\
0.58	-0.0244184\\
0.582	-0.0244208\\
0.584	-0.0244231\\
0.586	-0.0244251\\
0.588	-0.024427\\
0.59	-0.0244287\\
0.592	-0.0244302\\
0.594	-0.0244316\\
0.596	-0.0244329\\
0.598	-0.024434\\
0.6	-0.0244351\\
0.602	-0.024436\\
0.604	-0.0244369\\
0.606	-0.0244376\\
0.608	-0.0244384\\
0.61	-0.024439\\
0.612	-0.0244397\\
0.614	-0.0244403\\
0.616	-0.0244409\\
0.618	-0.0244415\\
0.62	-0.0244421\\
0.622	-0.0244428\\
0.624	-0.0244434\\
0.626	-0.0244441\\
0.628	-0.0244449\\
0.63	-0.0244457\\
0.632	-0.0244466\\
0.634	-0.0244475\\
0.636	-0.0244485\\
0.638	-0.0244496\\
0.64	-0.0244507\\
0.642	-0.024452\\
0.644	-0.0244533\\
0.646	-0.0244547\\
0.648	-0.0244562\\
0.65	-0.0244578\\
0.652	-0.0244595\\
0.654	-0.0244613\\
0.656	-0.0244632\\
0.658	-0.0244651\\
0.66	-0.0244672\\
0.662	-0.0244693\\
0.664	-0.0244715\\
0.666	-0.0244738\\
0.668	-0.0244761\\
0.67	-0.0244786\\
0.672	-0.024481\\
0.674	-0.0244836\\
0.676	-0.0244862\\
0.678	-0.0244888\\
0.68	-0.0244915\\
0.682	-0.0244943\\
0.684	-0.024497\\
0.686	-0.0244998\\
0.688	-0.0245026\\
0.69	-0.0245055\\
0.692	-0.0245083\\
0.694	-0.0245111\\
0.696	-0.024514\\
0.698	-0.0245168\\
0.7	-0.0245197\\
0.702	-0.0245225\\
0.704	-0.0245253\\
0.706	-0.0245281\\
0.708	-0.0245308\\
0.71	-0.0245335\\
0.712	-0.0245362\\
0.714	-0.0245389\\
0.716	-0.0245415\\
0.718	-0.0245441\\
0.72	-0.0245466\\
0.722	-0.0245491\\
0.724	-0.0245516\\
0.726	-0.024554\\
0.728	-0.0245563\\
0.73	-0.0245586\\
0.732	-0.0245609\\
0.734	-0.0245631\\
0.736	-0.0245653\\
0.738	-0.0245674\\
0.74	-0.0245694\\
0.742	-0.0245715\\
0.744	-0.0245734\\
0.746	-0.0245754\\
0.748	-0.0245772\\
0.75	-0.0245791\\
0.752	-0.0245809\\
0.754	-0.0245826\\
0.756	-0.0245844\\
0.758	-0.0245861\\
0.76	-0.0245877\\
0.762	-0.0245893\\
0.764	-0.0245909\\
0.766	-0.0245925\\
0.768	-0.024594\\
0.77	-0.0245955\\
0.772	-0.024597\\
0.774	-0.0245984\\
0.776	-0.0245999\\
0.778	-0.0246013\\
0.78	-0.0246027\\
0.782	-0.0246041\\
0.784	-0.0246054\\
0.786	-0.0246067\\
0.788	-0.0246081\\
0.79	-0.0246094\\
0.792	-0.0246107\\
0.794	-0.0246119\\
0.796	-0.0246132\\
0.798	-0.0246145\\
0.8	-0.0246157\\
0.802	-0.0246169\\
0.804	-0.0246181\\
0.806	-0.0246193\\
0.808	-0.0246205\\
0.81	-0.0246217\\
0.812	-0.0246228\\
0.814	-0.024624\\
0.816	-0.0246251\\
0.818	-0.0246262\\
0.82	-0.0246273\\
0.822	-0.0246284\\
0.824	-0.0246294\\
0.826	-0.0246305\\
0.828	-0.0246315\\
0.83	-0.0246325\\
0.832	-0.0246335\\
0.834	-0.0246344\\
0.836	-0.0246354\\
0.838	-0.0246363\\
0.84	-0.0246372\\
0.842	-0.0246381\\
0.844	-0.024639\\
0.846	-0.0246399\\
0.848	-0.0246407\\
0.85	-0.0246415\\
0.852	-0.0246423\\
0.854	-0.0246431\\
0.856	-0.0246438\\
0.858	-0.0246446\\
0.86	-0.0246453\\
0.862	-0.024646\\
0.864	-0.0246467\\
0.866	-0.0246474\\
0.868	-0.0246481\\
0.87	-0.0246487\\
0.872	-0.0246493\\
0.874	-0.02465\\
0.876	-0.0246506\\
0.878	-0.0246512\\
0.88	-0.0246518\\
0.882	-0.0246524\\
0.884	-0.024653\\
0.886	-0.0246536\\
0.888	-0.0246541\\
0.89	-0.0246547\\
0.892	-0.0246553\\
0.894	-0.0246559\\
0.896	-0.0246564\\
0.898	-0.024657\\
0.9	-0.0246576\\
0.902	-0.0246582\\
0.904	-0.0246588\\
0.906	-0.0246593\\
0.908	-0.0246599\\
0.91	-0.0246606\\
0.912	-0.0246612\\
0.914	-0.0246618\\
0.916	-0.0246624\\
0.918	-0.0246631\\
0.92	-0.0246637\\
0.922	-0.0246644\\
0.924	-0.0246651\\
0.926	-0.0246658\\
0.928	-0.0246665\\
0.93	-0.0246672\\
0.932	-0.0246679\\
0.934	-0.0246686\\
0.936	-0.0246694\\
0.938	-0.0246702\\
0.94	-0.0246709\\
0.942	-0.0246717\\
0.944	-0.0246725\\
0.946	-0.0246733\\
0.948	-0.0246741\\
0.95	-0.0246749\\
0.952	-0.0246757\\
0.954	-0.0246766\\
0.956	-0.0246774\\
0.958	-0.0246783\\
0.96	-0.0246791\\
0.962	-0.0246799\\
0.964	-0.0246808\\
0.966	-0.0246817\\
0.968	-0.0246825\\
0.97	-0.0246834\\
0.972	-0.0246842\\
0.974	-0.0246851\\
0.976	-0.0246859\\
0.978	-0.0246868\\
0.98	-0.0246876\\
0.982	-0.0246884\\
0.984	-0.0246893\\
0.986	-0.0246901\\
0.988	-0.0246909\\
0.99	-0.0246917\\
0.992	-0.0246925\\
0.994	-0.0246933\\
0.996	-0.024694\\
0.998	-0.0246948\\
1	-0.0246955\\
1.002	-0.0246963\\
1.004	-0.024697\\
1.006	-0.0246977\\
1.008	-0.0246984\\
1.01	-0.0246991\\
1.012	-0.0246997\\
1.014	-0.0247004\\
1.016	-0.024701\\
1.018	-0.0247016\\
1.02	-0.0247022\\
1.022	-0.0247028\\
1.024	-0.0247034\\
1.026	-0.024704\\
1.028	-0.0247045\\
1.03	-0.0247051\\
1.032	-0.0247056\\
1.034	-0.0247061\\
1.036	-0.0247066\\
1.038	-0.0247071\\
1.04	-0.0247075\\
1.042	-0.024708\\
1.044	-0.0247084\\
1.046	-0.0247089\\
1.048	-0.0247093\\
1.05	-0.0247097\\
1.052	-0.0247101\\
1.054	-0.0247105\\
1.056	-0.0247109\\
1.058	-0.0247112\\
1.06	-0.0247116\\
1.062	-0.024712\\
1.064	-0.0247123\\
1.066	-0.0247127\\
1.068	-0.024713\\
1.07	-0.0247134\\
1.072	-0.0247137\\
1.074	-0.024714\\
1.076	-0.0247143\\
1.078	-0.0247147\\
1.08	-0.024715\\
1.082	-0.0247153\\
1.084	-0.0247156\\
1.086	-0.0247159\\
1.088	-0.0247162\\
1.09	-0.0247166\\
1.092	-0.0247169\\
1.094	-0.0247172\\
1.096	-0.0247175\\
1.098	-0.0247178\\
1.1	-0.0247181\\
1.102	-0.0247184\\
1.104	-0.0247187\\
1.106	-0.0247191\\
1.108	-0.0247194\\
1.11	-0.0247197\\
1.112	-0.02472\\
1.114	-0.0247203\\
1.116	-0.0247207\\
1.118	-0.024721\\
1.12	-0.0247213\\
1.122	-0.0247216\\
1.124	-0.0247219\\
1.126	-0.0247223\\
1.128	-0.0247226\\
1.13	-0.0247229\\
1.132	-0.0247233\\
1.134	-0.0247236\\
1.136	-0.0247239\\
1.138	-0.0247243\\
1.14	-0.0247246\\
1.142	-0.0247249\\
1.144	-0.0247253\\
1.146	-0.0247256\\
1.148	-0.0247259\\
1.15	-0.0247263\\
1.152	-0.0247266\\
1.154	-0.024727\\
1.156	-0.0247273\\
1.158	-0.0247276\\
1.16	-0.024728\\
1.162	-0.0247283\\
1.164	-0.0247286\\
1.166	-0.024729\\
1.168	-0.0247293\\
1.17	-0.0247296\\
1.172	-0.0247299\\
1.174	-0.0247303\\
1.176	-0.0247306\\
1.178	-0.0247309\\
1.18	-0.0247312\\
1.182	-0.0247316\\
1.184	-0.0247319\\
1.186	-0.0247322\\
1.188	-0.0247325\\
1.19	-0.0247328\\
1.192	-0.0247331\\
1.194	-0.0247334\\
1.196	-0.0247337\\
1.198	-0.024734\\
1.2	-0.0247343\\
1.202	-0.0247346\\
1.204	-0.0247349\\
1.206	-0.0247352\\
1.208	-0.0247355\\
1.21	-0.0247358\\
1.212	-0.0247361\\
1.214	-0.0247363\\
1.216	-0.0247366\\
1.218	-0.0247369\\
1.22	-0.0247372\\
1.222	-0.0247374\\
1.224	-0.0247377\\
1.226	-0.024738\\
1.228	-0.0247382\\
1.23	-0.0247385\\
1.232	-0.0247387\\
1.234	-0.024739\\
1.236	-0.0247392\\
1.238	-0.0247395\\
1.24	-0.0247397\\
1.242	-0.02474\\
1.244	-0.0247402\\
1.246	-0.0247405\\
1.248	-0.0247407\\
1.25	-0.0247409\\
1.252	-0.0247412\\
1.254	-0.0247414\\
1.256	-0.0247416\\
1.258	-0.0247418\\
1.26	-0.0247421\\
1.262	-0.0247423\\
1.264	-0.0247425\\
1.266	-0.0247427\\
1.268	-0.0247429\\
1.27	-0.0247432\\
1.272	-0.0247434\\
1.274	-0.0247436\\
1.276	-0.0247438\\
1.278	-0.024744\\
1.28	-0.0247442\\
1.282	-0.0247444\\
1.284	-0.0247446\\
1.286	-0.0247448\\
1.288	-0.024745\\
1.29	-0.0247452\\
1.292	-0.0247454\\
1.294	-0.0247456\\
1.296	-0.0247458\\
1.298	-0.0247459\\
1.3	-0.0247461\\
1.302	-0.0247463\\
1.304	-0.0247465\\
1.306	-0.0247467\\
1.308	-0.0247468\\
1.31	-0.024747\\
1.312	-0.0247472\\
1.314	-0.0247474\\
1.316	-0.0247475\\
1.318	-0.0247477\\
1.32	-0.0247479\\
1.322	-0.024748\\
1.324	-0.0247482\\
1.326	-0.0247484\\
1.328	-0.0247485\\
1.33	-0.0247487\\
1.332	-0.0247489\\
1.334	-0.024749\\
1.336	-0.0247492\\
1.338	-0.0247493\\
1.34	-0.0247495\\
1.342	-0.0247496\\
1.344	-0.0247498\\
1.346	-0.0247499\\
1.348	-0.0247501\\
1.35	-0.0247502\\
1.352	-0.0247504\\
1.354	-0.0247505\\
1.356	-0.0247507\\
1.358	-0.0247508\\
1.36	-0.0247509\\
1.362	-0.0247511\\
1.364	-0.0247512\\
1.366	-0.0247514\\
1.368	-0.0247515\\
1.37	-0.0247516\\
1.372	-0.0247518\\
1.374	-0.0247519\\
1.376	-0.0247521\\
1.378	-0.0247522\\
1.38	-0.0247523\\
1.382	-0.0247525\\
1.384	-0.0247526\\
1.386	-0.0247527\\
1.388	-0.0247529\\
1.39	-0.024753\\
1.392	-0.0247531\\
1.394	-0.0247532\\
1.396	-0.0247534\\
1.398	-0.0247535\\
1.4	-0.0247536\\
1.402	-0.0247538\\
1.404	-0.0247539\\
1.406	-0.024754\\
1.408	-0.0247542\\
1.41	-0.0247543\\
1.412	-0.0247544\\
1.414	-0.0247545\\
1.416	-0.0247547\\
1.418	-0.0247548\\
1.42	-0.0247549\\
1.422	-0.024755\\
1.424	-0.0247552\\
1.426	-0.0247553\\
1.428	-0.0247554\\
1.43	-0.0247555\\
1.432	-0.0247557\\
1.434	-0.0247558\\
1.436	-0.0247559\\
1.438	-0.024756\\
1.44	-0.0247562\\
1.442	-0.0247563\\
1.444	-0.0247564\\
1.446	-0.0247565\\
1.448	-0.0247566\\
1.45	-0.0247568\\
1.452	-0.0247569\\
1.454	-0.024757\\
1.456	-0.0247571\\
1.458	-0.0247572\\
1.46	-0.0247574\\
1.462	-0.0247575\\
1.464	-0.0247576\\
1.466	-0.0247577\\
1.468	-0.0247578\\
1.47	-0.0247579\\
1.472	-0.024758\\
1.474	-0.0247581\\
1.476	-0.0247583\\
1.478	-0.0247584\\
1.48	-0.0247585\\
1.482	-0.0247586\\
1.484	-0.0247587\\
1.486	-0.0247588\\
1.488	-0.0247589\\
1.49	-0.024759\\
1.492	-0.0247591\\
1.494	-0.0247592\\
1.496	-0.0247593\\
1.498	-0.0247594\\
1.5	-0.0247595\\
1.502	-0.0247596\\
1.504	-0.0247597\\
1.506	-0.0247598\\
1.508	-0.0247599\\
1.51	-0.02476\\
1.512	-0.0247601\\
1.514	-0.0247602\\
1.516	-0.0247602\\
1.518	-0.0247603\\
1.52	-0.0247604\\
1.522	-0.0247605\\
1.524	-0.0247606\\
1.526	-0.0247607\\
1.528	-0.0247608\\
1.53	-0.0247608\\
1.532	-0.0247609\\
1.534	-0.024761\\
1.536	-0.0247611\\
1.538	-0.0247612\\
1.54	-0.0247612\\
1.542	-0.0247613\\
1.544	-0.0247614\\
1.546	-0.0247615\\
1.548	-0.0247615\\
1.55	-0.0247616\\
1.552	-0.0247617\\
1.554	-0.0247618\\
1.556	-0.0247618\\
1.558	-0.0247619\\
1.56	-0.024762\\
1.562	-0.0247621\\
1.564	-0.0247621\\
1.566	-0.0247622\\
1.568	-0.0247623\\
1.57	-0.0247623\\
1.572	-0.0247624\\
1.574	-0.0247625\\
1.576	-0.0247626\\
1.578	-0.0247626\\
1.58	-0.0247627\\
1.582	-0.0247628\\
1.584	-0.0247628\\
1.586	-0.0247629\\
1.588	-0.024763\\
1.59	-0.024763\\
1.592	-0.0247631\\
1.594	-0.0247632\\
1.596	-0.0247632\\
1.598	-0.0247633\\
1.6	-0.0247634\\
1.602	-0.0247634\\
1.604	-0.0247635\\
1.606	-0.0247636\\
1.608	-0.0247636\\
1.61	-0.0247637\\
1.612	-0.0247637\\
1.614	-0.0247638\\
1.616	-0.0247639\\
1.618	-0.0247639\\
1.62	-0.024764\\
1.622	-0.0247641\\
1.624	-0.0247641\\
1.626	-0.0247642\\
1.628	-0.0247642\\
1.63	-0.0247643\\
1.632	-0.0247644\\
1.634	-0.0247644\\
1.636	-0.0247645\\
1.638	-0.0247645\\
1.64	-0.0247646\\
1.642	-0.0247647\\
1.644	-0.0247647\\
1.646	-0.0247648\\
1.648	-0.0247648\\
1.65	-0.0247649\\
1.652	-0.024765\\
1.654	-0.024765\\
1.656	-0.0247651\\
1.658	-0.0247651\\
1.66	-0.0247652\\
1.662	-0.0247652\\
1.664	-0.0247653\\
1.666	-0.0247653\\
1.668	-0.0247654\\
1.67	-0.0247655\\
1.672	-0.0247655\\
1.674	-0.0247656\\
1.676	-0.0247656\\
1.678	-0.0247657\\
1.68	-0.0247657\\
1.682	-0.0247658\\
1.684	-0.0247658\\
1.686	-0.0247659\\
1.688	-0.0247659\\
1.69	-0.024766\\
1.692	-0.024766\\
1.694	-0.0247661\\
1.696	-0.0247661\\
1.698	-0.0247662\\
1.7	-0.0247662\\
1.702	-0.0247663\\
1.704	-0.0247663\\
1.706	-0.0247664\\
1.708	-0.0247664\\
1.71	-0.0247664\\
1.712	-0.0247665\\
1.714	-0.0247665\\
1.716	-0.0247666\\
1.718	-0.0247666\\
1.72	-0.0247667\\
1.722	-0.0247667\\
1.724	-0.0247668\\
1.726	-0.0247668\\
1.728	-0.0247668\\
1.73	-0.0247669\\
1.732	-0.0247669\\
1.734	-0.024767\\
1.736	-0.024767\\
1.738	-0.0247671\\
1.74	-0.0247671\\
1.742	-0.0247671\\
1.744	-0.0247672\\
1.746	-0.0247672\\
1.748	-0.0247673\\
1.75	-0.0247673\\
1.752	-0.0247673\\
1.754	-0.0247674\\
1.756	-0.0247674\\
1.758	-0.0247675\\
1.76	-0.0247675\\
1.762	-0.0247675\\
1.764	-0.0247676\\
1.766	-0.0247676\\
1.768	-0.0247676\\
1.77	-0.0247677\\
1.772	-0.0247677\\
1.774	-0.0247678\\
1.776	-0.0247678\\
1.778	-0.0247678\\
1.78	-0.0247679\\
1.782	-0.0247679\\
1.784	-0.0247679\\
1.786	-0.024768\\
1.788	-0.024768\\
1.79	-0.024768\\
1.792	-0.0247681\\
1.794	-0.0247681\\
1.796	-0.0247681\\
1.798	-0.0247682\\
1.8	-0.0247682\\
1.802	-0.0247682\\
1.804	-0.0247683\\
1.806	-0.0247683\\
1.808	-0.0247683\\
1.81	-0.0247684\\
1.812	-0.0247684\\
1.814	-0.0247684\\
1.816	-0.0247685\\
1.818	-0.0247685\\
1.82	-0.0247685\\
1.822	-0.0247686\\
1.824	-0.0247686\\
1.826	-0.0247686\\
1.828	-0.0247687\\
1.83	-0.0247687\\
1.832	-0.0247687\\
1.834	-0.0247688\\
1.836	-0.0247688\\
1.838	-0.0247688\\
1.84	-0.0247688\\
1.842	-0.0247689\\
1.844	-0.0247689\\
1.846	-0.0247689\\
1.848	-0.024769\\
1.85	-0.024769\\
1.852	-0.024769\\
1.854	-0.024769\\
1.856	-0.0247691\\
1.858	-0.0247691\\
1.86	-0.0247691\\
1.862	-0.0247691\\
1.864	-0.0247692\\
1.866	-0.0247692\\
1.868	-0.0247692\\
1.87	-0.0247693\\
1.872	-0.0247693\\
1.874	-0.0247693\\
1.876	-0.0247693\\
1.878	-0.0247694\\
1.88	-0.0247694\\
1.882	-0.0247694\\
1.884	-0.0247694\\
1.886	-0.0247695\\
1.888	-0.0247695\\
1.89	-0.0247695\\
1.892	-0.0247695\\
1.894	-0.0247696\\
1.896	-0.0247696\\
1.898	-0.0247696\\
1.9	-0.0247696\\
1.902	-0.0247697\\
1.904	-0.0247697\\
1.906	-0.0247697\\
1.908	-0.0247697\\
1.91	-0.0247698\\
1.912	-0.0247698\\
1.914	-0.0247698\\
1.916	-0.0247698\\
1.918	-0.0247698\\
1.92	-0.0247699\\
1.922	-0.0247699\\
1.924	-0.0247699\\
1.926	-0.0247699\\
1.928	-0.02477\\
1.93	-0.02477\\
1.932	-0.02477\\
1.934	-0.02477\\
1.936	-0.02477\\
1.938	-0.0247701\\
1.94	-0.0247701\\
1.942	-0.0247701\\
1.944	-0.0247701\\
1.946	-0.0247701\\
1.948	-0.0247702\\
1.95	-0.0247702\\
1.952	-0.0247702\\
1.954	-0.0247702\\
1.956	-0.0247702\\
1.958	-0.0247703\\
1.96	-0.0247703\\
1.962	-0.0247703\\
1.964	-0.0247703\\
1.966	-0.0247703\\
1.968	-0.0247704\\
1.97	-0.0247704\\
1.972	-0.0247704\\
1.974	-0.0247704\\
1.976	-0.0247704\\
1.978	-0.0247705\\
1.98	-0.0247705\\
1.982	-0.0247705\\
1.984	-0.0247705\\
1.986	-0.0247705\\
1.988	-0.0247705\\
1.99	-0.0247706\\
1.992	-0.0247706\\
1.994	-0.0247706\\
1.996	-0.0247706\\
1.998	-0.0247706\\
2	-0.0247707\\
};
\addplot [color=mycolor6, line width=1.0pt, forget plot]
  table[row sep=crcr]{%
-0.024	0\\
-0.022	0\\
-0.02	0\\
-0.018	0\\
-0.016	0\\
-0.014	0\\
-0.012	0\\
-0.01	0\\
-0.008	0\\
-0.006	0\\
-0.004	0\\
-0.002	0\\
0	0\\
0.002	-0.00184649\\
0.004	-0.00342631\\
0.006	-0.00478673\\
0.008	-0.00596702\\
0.01	-0.00699976\\
0.012	-0.00791182\\
0.014	-0.00872527\\
0.016	-0.00945819\\
0.018	-0.0101253\\
0.02	-0.0107383\\
0.022	-0.011307\\
0.024	-0.0118388\\
0.026	-0.0123398\\
0.028	-0.0128146\\
0.03	-0.0132669\\
0.032	-0.0136993\\
0.034	-0.014114\\
0.036	-0.0145123\\
0.038	-0.0148953\\
0.04	-0.0152639\\
0.042	-0.0156184\\
0.044	-0.0159593\\
0.046	-0.0162867\\
0.048	-0.0166009\\
0.05	-0.0169018\\
0.052	-0.0171898\\
0.054	-0.0174648\\
0.056	-0.0177271\\
0.058	-0.0179769\\
0.06	-0.0182144\\
0.062	-0.01844\\
0.064	-0.0186539\\
0.066	-0.0188566\\
0.068	-0.0190484\\
0.07	-0.0192299\\
0.072	-0.0194014\\
0.074	-0.0195634\\
0.076	-0.0197165\\
0.078	-0.0198611\\
0.08	-0.0199977\\
0.082	-0.0201268\\
0.084	-0.020249\\
0.086	-0.0203646\\
0.088	-0.0204741\\
0.09	-0.020578\\
0.092	-0.0206767\\
0.094	-0.0207705\\
0.096	-0.02086\\
0.098	-0.0209453\\
0.1	-0.0210269\\
0.102	-0.0211051\\
0.104	-0.0211801\\
0.106	-0.0212522\\
0.108	-0.0213216\\
0.11	-0.0213885\\
0.112	-0.0214531\\
0.114	-0.0215155\\
0.116	-0.0215759\\
0.118	-0.0216343\\
0.12	-0.0216909\\
0.122	-0.0217456\\
0.124	-0.0217986\\
0.126	-0.0218498\\
0.128	-0.0218993\\
0.13	-0.0219469\\
0.132	-0.0219928\\
0.134	-0.0220367\\
0.136	-0.0220788\\
0.138	-0.0221189\\
0.14	-0.022157\\
0.142	-0.0221929\\
0.144	-0.0222267\\
0.146	-0.0222583\\
0.148	-0.0222876\\
0.15	-0.0223145\\
0.152	-0.0223391\\
0.154	-0.0223612\\
0.156	-0.0223808\\
0.158	-0.022398\\
0.16	-0.0224127\\
0.162	-0.0224249\\
0.164	-0.0224347\\
0.166	-0.0224421\\
0.168	-0.0224471\\
0.17	-0.0224498\\
0.172	-0.0224503\\
0.174	-0.0224487\\
0.176	-0.0224451\\
0.178	-0.0224396\\
0.18	-0.0224324\\
0.182	-0.0224235\\
0.184	-0.0224132\\
0.186	-0.0224016\\
0.188	-0.0223889\\
0.19	-0.0223752\\
0.192	-0.0223607\\
0.194	-0.0223457\\
0.196	-0.0223302\\
0.198	-0.0223144\\
0.2	-0.0222986\\
0.202	-0.0222829\\
0.204	-0.0222675\\
0.206	-0.0222525\\
0.208	-0.0222381\\
0.21	-0.0222245\\
0.212	-0.0222117\\
0.214	-0.0222\\
0.216	-0.0221894\\
0.218	-0.0221801\\
0.22	-0.0221722\\
0.222	-0.0221657\\
0.224	-0.0221607\\
0.226	-0.0221574\\
0.228	-0.0221557\\
0.23	-0.0221558\\
0.232	-0.0221575\\
0.234	-0.0221611\\
0.236	-0.0221664\\
0.238	-0.0221736\\
0.24	-0.0221825\\
0.242	-0.0221932\\
0.244	-0.0222056\\
0.246	-0.0222197\\
0.248	-0.0222354\\
0.25	-0.0222528\\
0.252	-0.0222717\\
0.254	-0.0222922\\
0.256	-0.022314\\
0.258	-0.0223372\\
0.26	-0.0223616\\
0.262	-0.0223873\\
0.264	-0.0224141\\
0.266	-0.0224418\\
0.268	-0.0224706\\
0.27	-0.0225001\\
0.272	-0.0225305\\
0.274	-0.0225615\\
0.276	-0.022593\\
0.278	-0.0226251\\
0.28	-0.0226576\\
0.282	-0.0226904\\
0.284	-0.0227235\\
0.286	-0.0227568\\
0.288	-0.0227901\\
0.29	-0.0228235\\
0.292	-0.0228568\\
0.294	-0.02289\\
0.296	-0.022923\\
0.298	-0.0229558\\
0.3	-0.0229883\\
0.302	-0.0230205\\
0.304	-0.0230523\\
0.306	-0.0230836\\
0.308	-0.0231145\\
0.31	-0.0231449\\
0.312	-0.0231747\\
0.314	-0.023204\\
0.316	-0.0232327\\
0.318	-0.0232607\\
0.32	-0.0232881\\
0.322	-0.0233149\\
0.324	-0.023341\\
0.326	-0.0233664\\
0.328	-0.0233911\\
0.33	-0.0234152\\
0.332	-0.0234385\\
0.334	-0.0234612\\
0.336	-0.0234831\\
0.338	-0.0235044\\
0.34	-0.023525\\
0.342	-0.0235449\\
0.344	-0.0235641\\
0.346	-0.0235827\\
0.348	-0.0236006\\
0.35	-0.0236179\\
0.352	-0.0236346\\
0.354	-0.0236506\\
0.356	-0.023666\\
0.358	-0.0236809\\
0.36	-0.0236952\\
0.362	-0.0237089\\
0.364	-0.023722\\
0.366	-0.0237347\\
0.368	-0.0237468\\
0.37	-0.0237585\\
0.372	-0.0237696\\
0.374	-0.0237803\\
0.376	-0.0237906\\
0.378	-0.0238005\\
0.38	-0.0238099\\
0.382	-0.023819\\
0.384	-0.0238276\\
0.386	-0.023836\\
0.388	-0.023844\\
0.39	-0.0238516\\
0.392	-0.023859\\
0.394	-0.0238661\\
0.396	-0.0238729\\
0.398	-0.0238795\\
0.4	-0.0238859\\
0.402	-0.023892\\
0.404	-0.0238979\\
0.406	-0.0239037\\
0.408	-0.0239093\\
0.41	-0.0239147\\
0.412	-0.02392\\
0.414	-0.0239252\\
0.416	-0.0239303\\
0.418	-0.0239353\\
0.42	-0.0239402\\
0.422	-0.0239451\\
0.424	-0.0239499\\
0.426	-0.0239547\\
0.428	-0.0239594\\
0.43	-0.0239641\\
0.432	-0.0239689\\
0.434	-0.0239736\\
0.436	-0.0239784\\
0.438	-0.0239832\\
0.44	-0.023988\\
0.442	-0.0239929\\
0.444	-0.0239978\\
0.446	-0.0240028\\
0.448	-0.0240079\\
0.45	-0.0240131\\
0.452	-0.0240183\\
0.454	-0.0240236\\
0.456	-0.024029\\
0.458	-0.0240345\\
0.46	-0.02404\\
0.462	-0.0240457\\
0.464	-0.0240515\\
0.466	-0.0240574\\
0.468	-0.0240633\\
0.47	-0.0240694\\
0.472	-0.0240755\\
0.474	-0.0240818\\
0.476	-0.0240881\\
0.478	-0.0240946\\
0.48	-0.0241011\\
0.482	-0.0241077\\
0.484	-0.0241144\\
0.486	-0.0241211\\
0.488	-0.0241279\\
0.49	-0.0241348\\
0.492	-0.0241417\\
0.494	-0.0241487\\
0.496	-0.0241557\\
0.498	-0.0241628\\
0.5	-0.0241698\\
0.502	-0.0241769\\
0.504	-0.0241841\\
0.506	-0.0241912\\
0.508	-0.0241983\\
0.51	-0.0242054\\
0.512	-0.0242125\\
0.514	-0.0242195\\
0.516	-0.0242265\\
0.518	-0.0242335\\
0.52	-0.0242404\\
0.522	-0.0242472\\
0.524	-0.024254\\
0.526	-0.0242607\\
0.528	-0.0242673\\
0.53	-0.0242738\\
0.532	-0.0242803\\
0.534	-0.0242866\\
0.536	-0.0242928\\
0.538	-0.0242989\\
0.54	-0.0243048\\
0.542	-0.0243107\\
0.544	-0.0243164\\
0.546	-0.0243219\\
0.548	-0.0243273\\
0.55	-0.0243326\\
0.552	-0.0243377\\
0.554	-0.0243427\\
0.556	-0.0243475\\
0.558	-0.0243522\\
0.56	-0.0243566\\
0.562	-0.024361\\
0.564	-0.0243652\\
0.566	-0.0243692\\
0.568	-0.024373\\
0.57	-0.0243768\\
0.572	-0.0243803\\
0.574	-0.0243837\\
0.576	-0.024387\\
0.578	-0.0243901\\
0.58	-0.0243931\\
0.582	-0.0243959\\
0.584	-0.0243986\\
0.586	-0.0244012\\
0.588	-0.0244036\\
0.59	-0.024406\\
0.592	-0.0244082\\
0.594	-0.0244103\\
0.596	-0.0244124\\
0.598	-0.0244143\\
0.6	-0.0244162\\
0.602	-0.024418\\
0.604	-0.0244197\\
0.606	-0.0244213\\
0.608	-0.0244229\\
0.61	-0.0244245\\
0.612	-0.024426\\
0.614	-0.0244274\\
0.616	-0.0244289\\
0.618	-0.0244303\\
0.62	-0.0244317\\
0.622	-0.0244331\\
0.624	-0.0244345\\
0.626	-0.0244359\\
0.628	-0.0244373\\
0.63	-0.0244388\\
0.632	-0.0244402\\
0.634	-0.0244417\\
0.636	-0.0244432\\
0.638	-0.0244447\\
0.64	-0.0244463\\
0.642	-0.0244479\\
0.644	-0.0244496\\
0.646	-0.0244513\\
0.648	-0.0244531\\
0.65	-0.0244549\\
0.652	-0.0244567\\
0.654	-0.0244586\\
0.656	-0.0244606\\
0.658	-0.0244626\\
0.66	-0.0244647\\
0.662	-0.0244668\\
0.664	-0.024469\\
0.666	-0.0244712\\
0.668	-0.0244735\\
0.67	-0.0244758\\
0.672	-0.0244781\\
0.674	-0.0244805\\
0.676	-0.024483\\
0.678	-0.0244854\\
0.68	-0.024488\\
0.682	-0.0244905\\
0.684	-0.0244931\\
0.686	-0.0244957\\
0.688	-0.0244983\\
0.69	-0.0245009\\
0.692	-0.0245036\\
0.694	-0.0245062\\
0.696	-0.0245089\\
0.698	-0.0245116\\
0.7	-0.0245143\\
0.702	-0.0245169\\
0.704	-0.0245196\\
0.706	-0.0245223\\
0.708	-0.0245249\\
0.71	-0.0245276\\
0.712	-0.0245302\\
0.714	-0.0245328\\
0.716	-0.0245353\\
0.718	-0.0245379\\
0.72	-0.0245404\\
0.722	-0.0245429\\
0.724	-0.0245453\\
0.726	-0.0245477\\
0.728	-0.0245501\\
0.73	-0.0245525\\
0.732	-0.0245548\\
0.734	-0.024557\\
0.736	-0.0245593\\
0.738	-0.0245615\\
0.74	-0.0245636\\
0.742	-0.0245657\\
0.744	-0.0245678\\
0.746	-0.0245698\\
0.748	-0.0245718\\
0.75	-0.0245737\\
0.752	-0.0245756\\
0.754	-0.0245774\\
0.756	-0.0245792\\
0.758	-0.024581\\
0.76	-0.0245827\\
0.762	-0.0245844\\
0.764	-0.024586\\
0.766	-0.0245876\\
0.768	-0.0245892\\
0.77	-0.0245907\\
0.772	-0.0245922\\
0.774	-0.0245937\\
0.776	-0.0245951\\
0.778	-0.0245965\\
0.78	-0.0245979\\
0.782	-0.0245993\\
0.784	-0.0246006\\
0.786	-0.0246018\\
0.788	-0.0246031\\
0.79	-0.0246043\\
0.792	-0.0246055\\
0.794	-0.0246067\\
0.796	-0.0246079\\
0.798	-0.024609\\
0.8	-0.0246102\\
0.802	-0.0246113\\
0.804	-0.0246123\\
0.806	-0.0246134\\
0.808	-0.0246144\\
0.81	-0.0246155\\
0.812	-0.0246165\\
0.814	-0.0246175\\
0.816	-0.0246185\\
0.818	-0.0246194\\
0.82	-0.0246204\\
0.822	-0.0246214\\
0.824	-0.0246223\\
0.826	-0.0246232\\
0.828	-0.0246241\\
0.83	-0.024625\\
0.832	-0.0246259\\
0.834	-0.0246268\\
0.836	-0.0246277\\
0.838	-0.0246285\\
0.84	-0.0246294\\
0.842	-0.0246302\\
0.844	-0.0246311\\
0.846	-0.0246319\\
0.848	-0.0246327\\
0.85	-0.0246336\\
0.852	-0.0246344\\
0.854	-0.0246352\\
0.856	-0.024636\\
0.858	-0.0246368\\
0.86	-0.0246376\\
0.862	-0.0246384\\
0.864	-0.0246392\\
0.866	-0.02464\\
0.868	-0.0246408\\
0.87	-0.0246416\\
0.872	-0.0246423\\
0.874	-0.0246431\\
0.876	-0.0246439\\
0.878	-0.0246447\\
0.88	-0.0246455\\
0.882	-0.0246462\\
0.884	-0.024647\\
0.886	-0.0246478\\
0.888	-0.0246486\\
0.89	-0.0246493\\
0.892	-0.0246501\\
0.894	-0.0246509\\
0.896	-0.0246517\\
0.898	-0.0246525\\
0.9	-0.0246533\\
0.902	-0.0246541\\
0.904	-0.0246549\\
0.906	-0.0246556\\
0.908	-0.0246564\\
0.91	-0.0246572\\
0.912	-0.024658\\
0.914	-0.0246588\\
0.916	-0.0246596\\
0.918	-0.0246605\\
0.92	-0.0246613\\
0.922	-0.0246621\\
0.924	-0.0246629\\
0.926	-0.0246637\\
0.928	-0.0246645\\
0.93	-0.0246653\\
0.932	-0.0246661\\
0.934	-0.0246669\\
0.936	-0.0246678\\
0.938	-0.0246686\\
0.94	-0.0246694\\
0.942	-0.0246702\\
0.944	-0.024671\\
0.946	-0.0246718\\
0.948	-0.0246726\\
0.95	-0.0246734\\
0.952	-0.0246742\\
0.954	-0.024675\\
0.956	-0.0246758\\
0.958	-0.0246766\\
0.96	-0.0246774\\
0.962	-0.0246782\\
0.964	-0.024679\\
0.966	-0.0246797\\
0.968	-0.0246805\\
0.97	-0.0246813\\
0.972	-0.024682\\
0.974	-0.0246828\\
0.976	-0.0246835\\
0.978	-0.0246843\\
0.98	-0.024685\\
0.982	-0.0246857\\
0.984	-0.0246864\\
0.986	-0.0246871\\
0.988	-0.0246878\\
0.99	-0.0246885\\
0.992	-0.0246892\\
0.994	-0.0246899\\
0.996	-0.0246905\\
0.998	-0.0246912\\
1	-0.0246918\\
1.002	-0.0246925\\
1.004	-0.0246931\\
1.006	-0.0246937\\
1.008	-0.0246943\\
1.01	-0.0246949\\
1.012	-0.0246955\\
1.014	-0.0246961\\
1.016	-0.0246967\\
1.018	-0.0246973\\
1.02	-0.0246978\\
1.022	-0.0246984\\
1.024	-0.0246989\\
1.026	-0.0246995\\
1.028	-0.0247\\
1.03	-0.0247005\\
1.032	-0.024701\\
1.034	-0.0247016\\
1.036	-0.0247021\\
1.038	-0.0247026\\
1.04	-0.0247031\\
1.042	-0.0247035\\
1.044	-0.024704\\
1.046	-0.0247045\\
1.048	-0.024705\\
1.05	-0.0247054\\
1.052	-0.0247059\\
1.054	-0.0247063\\
1.056	-0.0247068\\
1.058	-0.0247072\\
1.06	-0.0247077\\
1.062	-0.0247081\\
1.064	-0.0247086\\
1.066	-0.024709\\
1.068	-0.0247094\\
1.07	-0.0247098\\
1.072	-0.0247102\\
1.074	-0.0247107\\
1.076	-0.0247111\\
1.078	-0.0247115\\
1.08	-0.0247119\\
1.082	-0.0247123\\
1.084	-0.0247127\\
1.086	-0.0247131\\
1.088	-0.0247135\\
1.09	-0.0247139\\
1.092	-0.0247142\\
1.094	-0.0247146\\
1.096	-0.024715\\
1.098	-0.0247154\\
1.1	-0.0247158\\
1.102	-0.0247161\\
1.104	-0.0247165\\
1.106	-0.0247169\\
1.108	-0.0247173\\
1.11	-0.0247176\\
1.112	-0.024718\\
1.114	-0.0247183\\
1.116	-0.0247187\\
1.118	-0.0247191\\
1.12	-0.0247194\\
1.122	-0.0247198\\
1.124	-0.0247201\\
1.126	-0.0247205\\
1.128	-0.0247208\\
1.13	-0.0247212\\
1.132	-0.0247215\\
1.134	-0.0247218\\
1.136	-0.0247222\\
1.138	-0.0247225\\
1.14	-0.0247229\\
1.142	-0.0247232\\
1.144	-0.0247235\\
1.146	-0.0247239\\
1.148	-0.0247242\\
1.15	-0.0247245\\
1.152	-0.0247248\\
1.154	-0.0247252\\
1.156	-0.0247255\\
1.158	-0.0247258\\
1.16	-0.0247261\\
1.162	-0.0247264\\
1.164	-0.0247268\\
1.166	-0.0247271\\
1.168	-0.0247274\\
1.17	-0.0247277\\
1.172	-0.024728\\
1.174	-0.0247283\\
1.176	-0.0247286\\
1.178	-0.0247289\\
1.18	-0.0247292\\
1.182	-0.0247295\\
1.184	-0.0247298\\
1.186	-0.0247301\\
1.188	-0.0247304\\
1.19	-0.0247307\\
1.192	-0.024731\\
1.194	-0.0247313\\
1.196	-0.0247316\\
1.198	-0.0247319\\
1.2	-0.0247322\\
1.202	-0.0247325\\
1.204	-0.0247327\\
1.206	-0.024733\\
1.208	-0.0247333\\
1.21	-0.0247336\\
1.212	-0.0247339\\
1.214	-0.0247342\\
1.216	-0.0247344\\
1.218	-0.0247347\\
1.22	-0.024735\\
1.222	-0.0247353\\
1.224	-0.0247355\\
1.226	-0.0247358\\
1.228	-0.0247361\\
1.23	-0.0247363\\
1.232	-0.0247366\\
1.234	-0.0247368\\
1.236	-0.0247371\\
1.238	-0.0247374\\
1.24	-0.0247376\\
1.242	-0.0247379\\
1.244	-0.0247381\\
1.246	-0.0247384\\
1.248	-0.0247386\\
1.25	-0.0247389\\
1.252	-0.0247391\\
1.254	-0.0247394\\
1.256	-0.0247396\\
1.258	-0.0247398\\
1.26	-0.0247401\\
1.262	-0.0247403\\
1.264	-0.0247406\\
1.266	-0.0247408\\
1.268	-0.024741\\
1.27	-0.0247412\\
1.272	-0.0247415\\
1.274	-0.0247417\\
1.276	-0.0247419\\
1.278	-0.0247421\\
1.28	-0.0247423\\
1.282	-0.0247426\\
1.284	-0.0247428\\
1.286	-0.024743\\
1.288	-0.0247432\\
1.29	-0.0247434\\
1.292	-0.0247436\\
1.294	-0.0247438\\
1.296	-0.024744\\
1.298	-0.0247442\\
1.3	-0.0247444\\
1.302	-0.0247446\\
1.304	-0.0247448\\
1.306	-0.024745\\
1.308	-0.0247452\\
1.31	-0.0247454\\
1.312	-0.0247456\\
1.314	-0.0247458\\
1.316	-0.0247459\\
1.318	-0.0247461\\
1.32	-0.0247463\\
1.322	-0.0247465\\
1.324	-0.0247467\\
1.326	-0.0247468\\
1.328	-0.024747\\
1.33	-0.0247472\\
1.332	-0.0247474\\
1.334	-0.0247475\\
1.336	-0.0247477\\
1.338	-0.0247479\\
1.34	-0.024748\\
1.342	-0.0247482\\
1.344	-0.0247484\\
1.346	-0.0247485\\
1.348	-0.0247487\\
1.35	-0.0247489\\
1.352	-0.024749\\
1.354	-0.0247492\\
1.356	-0.0247493\\
1.358	-0.0247495\\
1.36	-0.0247496\\
1.362	-0.0247498\\
1.364	-0.0247499\\
1.366	-0.0247501\\
1.368	-0.0247502\\
1.37	-0.0247504\\
1.372	-0.0247505\\
1.374	-0.0247507\\
1.376	-0.0247508\\
1.378	-0.024751\\
1.38	-0.0247511\\
1.382	-0.0247513\\
1.384	-0.0247514\\
1.386	-0.0247516\\
1.388	-0.0247517\\
1.39	-0.0247518\\
1.392	-0.024752\\
1.394	-0.0247521\\
1.396	-0.0247522\\
1.398	-0.0247524\\
1.4	-0.0247525\\
1.402	-0.0247527\\
1.404	-0.0247528\\
1.406	-0.0247529\\
1.408	-0.0247531\\
1.41	-0.0247532\\
1.412	-0.0247533\\
1.414	-0.0247534\\
1.416	-0.0247536\\
1.418	-0.0247537\\
1.42	-0.0247538\\
1.422	-0.024754\\
1.424	-0.0247541\\
1.426	-0.0247542\\
1.428	-0.0247543\\
1.43	-0.0247545\\
1.432	-0.0247546\\
1.434	-0.0247547\\
1.436	-0.0247548\\
1.438	-0.024755\\
1.44	-0.0247551\\
1.442	-0.0247552\\
1.444	-0.0247553\\
1.446	-0.0247554\\
1.448	-0.0247556\\
1.45	-0.0247557\\
1.452	-0.0247558\\
1.454	-0.0247559\\
1.456	-0.024756\\
1.458	-0.0247561\\
1.46	-0.0247563\\
1.462	-0.0247564\\
1.464	-0.0247565\\
1.466	-0.0247566\\
1.468	-0.0247567\\
1.47	-0.0247568\\
1.472	-0.0247569\\
1.474	-0.024757\\
1.476	-0.0247572\\
1.478	-0.0247573\\
1.48	-0.0247574\\
1.482	-0.0247575\\
1.484	-0.0247576\\
1.486	-0.0247577\\
1.488	-0.0247578\\
1.49	-0.0247579\\
1.492	-0.024758\\
1.494	-0.0247581\\
1.496	-0.0247582\\
1.498	-0.0247583\\
1.5	-0.0247584\\
1.502	-0.0247585\\
1.504	-0.0247586\\
1.506	-0.0247587\\
1.508	-0.0247588\\
1.51	-0.0247589\\
1.512	-0.024759\\
1.514	-0.0247591\\
1.516	-0.0247592\\
1.518	-0.0247593\\
1.52	-0.0247594\\
1.522	-0.0247595\\
1.524	-0.0247596\\
1.526	-0.0247597\\
1.528	-0.0247598\\
1.53	-0.0247599\\
1.532	-0.02476\\
1.534	-0.0247601\\
1.536	-0.0247602\\
1.538	-0.0247602\\
1.54	-0.0247603\\
1.542	-0.0247604\\
1.544	-0.0247605\\
1.546	-0.0247606\\
1.548	-0.0247607\\
1.55	-0.0247608\\
1.552	-0.0247608\\
1.554	-0.0247609\\
1.556	-0.024761\\
1.558	-0.0247611\\
1.56	-0.0247612\\
1.562	-0.0247612\\
1.564	-0.0247613\\
1.566	-0.0247614\\
1.568	-0.0247615\\
1.57	-0.0247616\\
1.572	-0.0247616\\
1.574	-0.0247617\\
1.576	-0.0247618\\
1.578	-0.0247619\\
1.58	-0.0247619\\
1.582	-0.024762\\
1.584	-0.0247621\\
1.586	-0.0247622\\
1.588	-0.0247622\\
1.59	-0.0247623\\
1.592	-0.0247624\\
1.594	-0.0247624\\
1.596	-0.0247625\\
1.598	-0.0247626\\
1.6	-0.0247627\\
1.602	-0.0247627\\
1.604	-0.0247628\\
1.606	-0.0247629\\
1.608	-0.0247629\\
1.61	-0.024763\\
1.612	-0.0247631\\
1.614	-0.0247631\\
1.616	-0.0247632\\
1.618	-0.0247633\\
1.62	-0.0247633\\
1.622	-0.0247634\\
1.624	-0.0247635\\
1.626	-0.0247635\\
1.628	-0.0247636\\
1.63	-0.0247636\\
1.632	-0.0247637\\
1.634	-0.0247638\\
1.636	-0.0247638\\
1.638	-0.0247639\\
1.64	-0.024764\\
1.642	-0.024764\\
1.644	-0.0247641\\
1.646	-0.0247641\\
1.648	-0.0247642\\
1.65	-0.0247642\\
1.652	-0.0247643\\
1.654	-0.0247644\\
1.656	-0.0247644\\
1.658	-0.0247645\\
1.66	-0.0247645\\
1.662	-0.0247646\\
1.664	-0.0247647\\
1.666	-0.0247647\\
1.668	-0.0247648\\
1.67	-0.0247648\\
1.672	-0.0247649\\
1.674	-0.0247649\\
1.676	-0.024765\\
1.678	-0.024765\\
1.68	-0.0247651\\
1.682	-0.0247651\\
1.684	-0.0247652\\
1.686	-0.0247653\\
1.688	-0.0247653\\
1.69	-0.0247654\\
1.692	-0.0247654\\
1.694	-0.0247655\\
1.696	-0.0247655\\
1.698	-0.0247656\\
1.7	-0.0247656\\
1.702	-0.0247657\\
1.704	-0.0247657\\
1.706	-0.0247658\\
1.708	-0.0247658\\
1.71	-0.0247659\\
1.712	-0.0247659\\
1.714	-0.024766\\
1.716	-0.024766\\
1.718	-0.0247661\\
1.72	-0.0247661\\
1.722	-0.0247662\\
1.724	-0.0247662\\
1.726	-0.0247663\\
1.728	-0.0247663\\
1.73	-0.0247663\\
1.732	-0.0247664\\
1.734	-0.0247664\\
1.736	-0.0247665\\
1.738	-0.0247665\\
1.74	-0.0247666\\
1.742	-0.0247666\\
1.744	-0.0247667\\
1.746	-0.0247667\\
1.748	-0.0247668\\
1.75	-0.0247668\\
1.752	-0.0247668\\
1.754	-0.0247669\\
1.756	-0.0247669\\
1.758	-0.024767\\
1.76	-0.024767\\
1.762	-0.024767\\
1.764	-0.0247671\\
1.766	-0.0247671\\
1.768	-0.0247672\\
1.77	-0.0247672\\
1.772	-0.0247673\\
1.774	-0.0247673\\
1.776	-0.0247673\\
1.778	-0.0247674\\
1.78	-0.0247674\\
1.782	-0.0247674\\
1.784	-0.0247675\\
1.786	-0.0247675\\
1.788	-0.0247676\\
1.79	-0.0247676\\
1.792	-0.0247676\\
1.794	-0.0247677\\
1.796	-0.0247677\\
1.798	-0.0247677\\
1.8	-0.0247678\\
1.802	-0.0247678\\
1.804	-0.0247679\\
1.806	-0.0247679\\
1.808	-0.0247679\\
1.81	-0.024768\\
1.812	-0.024768\\
1.814	-0.024768\\
1.816	-0.0247681\\
1.818	-0.0247681\\
1.82	-0.0247681\\
1.822	-0.0247682\\
1.824	-0.0247682\\
1.826	-0.0247682\\
1.828	-0.0247683\\
1.83	-0.0247683\\
1.832	-0.0247683\\
1.834	-0.0247684\\
1.836	-0.0247684\\
1.838	-0.0247684\\
1.84	-0.0247685\\
1.842	-0.0247685\\
1.844	-0.0247685\\
1.846	-0.0247685\\
1.848	-0.0247686\\
1.85	-0.0247686\\
1.852	-0.0247686\\
1.854	-0.0247687\\
1.856	-0.0247687\\
1.858	-0.0247687\\
1.86	-0.0247688\\
1.862	-0.0247688\\
1.864	-0.0247688\\
1.866	-0.0247688\\
1.868	-0.0247689\\
1.87	-0.0247689\\
1.872	-0.0247689\\
1.874	-0.024769\\
1.876	-0.024769\\
1.878	-0.024769\\
1.88	-0.024769\\
1.882	-0.0247691\\
1.884	-0.0247691\\
1.886	-0.0247691\\
1.888	-0.0247691\\
1.89	-0.0247692\\
1.892	-0.0247692\\
1.894	-0.0247692\\
1.896	-0.0247692\\
1.898	-0.0247693\\
1.9	-0.0247693\\
1.902	-0.0247693\\
1.904	-0.0247694\\
1.906	-0.0247694\\
1.908	-0.0247694\\
1.91	-0.0247694\\
1.912	-0.0247695\\
1.914	-0.0247695\\
1.916	-0.0247695\\
1.918	-0.0247695\\
1.92	-0.0247696\\
1.922	-0.0247696\\
1.924	-0.0247696\\
1.926	-0.0247696\\
1.928	-0.0247696\\
1.93	-0.0247697\\
1.932	-0.0247697\\
1.934	-0.0247697\\
1.936	-0.0247697\\
1.938	-0.0247698\\
1.94	-0.0247698\\
1.942	-0.0247698\\
1.944	-0.0247698\\
1.946	-0.0247699\\
1.948	-0.0247699\\
1.95	-0.0247699\\
1.952	-0.0247699\\
1.954	-0.0247699\\
1.956	-0.02477\\
1.958	-0.02477\\
1.96	-0.02477\\
1.962	-0.02477\\
1.964	-0.02477\\
1.966	-0.0247701\\
1.968	-0.0247701\\
1.97	-0.0247701\\
1.972	-0.0247701\\
1.974	-0.0247702\\
1.976	-0.0247702\\
1.978	-0.0247702\\
1.98	-0.0247702\\
1.982	-0.0247702\\
1.984	-0.0247703\\
1.986	-0.0247703\\
1.988	-0.0247703\\
1.99	-0.0247703\\
1.992	-0.0247703\\
1.994	-0.0247703\\
1.996	-0.0247704\\
1.998	-0.0247704\\
2	-0.0247704\\
};
\addplot [color=mycolor7, line width=1.0pt, forget plot]
  table[row sep=crcr]{%
-0.024	0\\
-0.022	0\\
-0.02	0\\
-0.018	0\\
-0.016	0\\
-0.014	0\\
-0.012	0\\
-0.01	0\\
-0.008	0\\
-0.006	0\\
-0.004	0\\
-0.002	0\\
0	0\\
0.002	-0.00184649\\
0.004	-0.00342633\\
0.006	-0.00478679\\
0.008	-0.00596718\\
0.01	-0.00700007\\
0.012	-0.00791235\\
0.014	-0.00872612\\
0.016	-0.00945946\\
0.018	-0.0101271\\
0.02	-0.0107409\\
0.022	-0.0113104\\
0.024	-0.0118432\\
0.026	-0.0123455\\
0.028	-0.0128217\\
0.03	-0.0132757\\
0.032	-0.01371\\
0.034	-0.0141269\\
0.036	-0.0145276\\
0.038	-0.0149134\\
0.04	-0.015285\\
0.042	-0.0156428\\
0.044	-0.0159872\\
0.046	-0.0163185\\
0.048	-0.0166366\\
0.05	-0.0169419\\
0.052	-0.0172343\\
0.054	-0.017514\\
0.056	-0.0177812\\
0.058	-0.018036\\
0.06	-0.0182785\\
0.062	-0.0185092\\
0.064	-0.0187283\\
0.066	-0.0189361\\
0.068	-0.0191329\\
0.07	-0.0193192\\
0.072	-0.0194954\\
0.074	-0.0196619\\
0.076	-0.0198191\\
0.078	-0.0199676\\
0.08	-0.0201076\\
0.082	-0.0202398\\
0.084	-0.0203644\\
0.086	-0.020482\\
0.088	-0.020593\\
0.09	-0.0206977\\
0.092	-0.0207966\\
0.094	-0.0208901\\
0.096	-0.0209784\\
0.098	-0.021062\\
0.1	-0.0211412\\
0.102	-0.0212162\\
0.104	-0.0212873\\
0.106	-0.0213549\\
0.108	-0.0214191\\
0.11	-0.0214802\\
0.112	-0.0215383\\
0.114	-0.0215937\\
0.116	-0.0216466\\
0.118	-0.021697\\
0.12	-0.0217451\\
0.122	-0.021791\\
0.124	-0.0218348\\
0.126	-0.0218766\\
0.128	-0.0219164\\
0.13	-0.0219544\\
0.132	-0.0219904\\
0.134	-0.0220247\\
0.136	-0.0220571\\
0.138	-0.0220877\\
0.14	-0.0221165\\
0.142	-0.0221435\\
0.144	-0.0221687\\
0.146	-0.022192\\
0.148	-0.0222136\\
0.15	-0.0222334\\
0.152	-0.0222514\\
0.154	-0.0222676\\
0.156	-0.0222821\\
0.158	-0.0222947\\
0.16	-0.0223057\\
0.162	-0.0223149\\
0.164	-0.0223225\\
0.166	-0.0223285\\
0.168	-0.0223329\\
0.17	-0.0223358\\
0.172	-0.0223373\\
0.174	-0.0223375\\
0.176	-0.0223365\\
0.178	-0.0223343\\
0.18	-0.022331\\
0.182	-0.0223268\\
0.184	-0.0223218\\
0.186	-0.0223161\\
0.188	-0.0223098\\
0.19	-0.022303\\
0.192	-0.0222959\\
0.194	-0.0222885\\
0.196	-0.022281\\
0.198	-0.0222736\\
0.2	-0.0222663\\
0.202	-0.0222592\\
0.204	-0.0222525\\
0.206	-0.0222463\\
0.208	-0.0222406\\
0.21	-0.0222356\\
0.212	-0.0222313\\
0.214	-0.0222279\\
0.216	-0.0222254\\
0.218	-0.0222238\\
0.22	-0.0222233\\
0.222	-0.0222239\\
0.224	-0.0222256\\
0.226	-0.0222285\\
0.228	-0.0222326\\
0.23	-0.0222379\\
0.232	-0.0222445\\
0.234	-0.0222524\\
0.236	-0.0222614\\
0.238	-0.0222718\\
0.24	-0.0222833\\
0.242	-0.0222962\\
0.244	-0.0223102\\
0.246	-0.0223253\\
0.248	-0.0223417\\
0.25	-0.0223591\\
0.252	-0.0223776\\
0.254	-0.0223972\\
0.256	-0.0224177\\
0.258	-0.0224391\\
0.26	-0.0224614\\
0.262	-0.0224846\\
0.264	-0.0225085\\
0.266	-0.0225331\\
0.268	-0.0225584\\
0.27	-0.0225843\\
0.272	-0.0226107\\
0.274	-0.0226377\\
0.276	-0.022665\\
0.278	-0.0226927\\
0.28	-0.0227207\\
0.282	-0.022749\\
0.284	-0.0227775\\
0.286	-0.0228061\\
0.288	-0.0228348\\
0.29	-0.0228636\\
0.292	-0.0228924\\
0.294	-0.0229211\\
0.296	-0.0229498\\
0.298	-0.0229783\\
0.3	-0.0230066\\
0.302	-0.0230347\\
0.304	-0.0230626\\
0.306	-0.0230902\\
0.308	-0.0231176\\
0.31	-0.0231445\\
0.312	-0.0231711\\
0.314	-0.0231973\\
0.316	-0.0232232\\
0.318	-0.0232485\\
0.32	-0.0232735\\
0.322	-0.0232979\\
0.324	-0.0233219\\
0.326	-0.0233454\\
0.328	-0.0233683\\
0.33	-0.0233908\\
0.332	-0.0234127\\
0.334	-0.0234341\\
0.336	-0.023455\\
0.338	-0.0234753\\
0.34	-0.0234951\\
0.342	-0.0235143\\
0.344	-0.023533\\
0.346	-0.0235512\\
0.348	-0.0235688\\
0.35	-0.0235859\\
0.352	-0.0236024\\
0.354	-0.0236184\\
0.356	-0.0236339\\
0.358	-0.0236489\\
0.36	-0.0236634\\
0.362	-0.0236774\\
0.364	-0.0236909\\
0.366	-0.0237039\\
0.368	-0.0237165\\
0.37	-0.0237286\\
0.372	-0.0237402\\
0.374	-0.0237514\\
0.376	-0.0237623\\
0.378	-0.0237727\\
0.38	-0.0237827\\
0.382	-0.0237923\\
0.384	-0.0238016\\
0.386	-0.0238106\\
0.388	-0.0238192\\
0.39	-0.0238275\\
0.392	-0.0238355\\
0.394	-0.0238432\\
0.396	-0.0238506\\
0.398	-0.0238578\\
0.4	-0.0238648\\
0.402	-0.0238716\\
0.404	-0.0238781\\
0.406	-0.0238845\\
0.408	-0.0238907\\
0.41	-0.0238967\\
0.412	-0.0239026\\
0.414	-0.0239083\\
0.416	-0.023914\\
0.418	-0.0239196\\
0.42	-0.0239251\\
0.422	-0.0239305\\
0.424	-0.0239359\\
0.426	-0.0239412\\
0.428	-0.0239465\\
0.43	-0.0239518\\
0.432	-0.023957\\
0.434	-0.0239623\\
0.436	-0.0239676\\
0.438	-0.023973\\
0.44	-0.0239783\\
0.442	-0.0239837\\
0.444	-0.0239892\\
0.446	-0.0239947\\
0.448	-0.0240003\\
0.45	-0.0240059\\
0.452	-0.0240117\\
0.454	-0.0240174\\
0.456	-0.0240233\\
0.458	-0.0240293\\
0.46	-0.0240353\\
0.462	-0.0240414\\
0.464	-0.0240476\\
0.466	-0.0240539\\
0.468	-0.0240603\\
0.47	-0.0240667\\
0.472	-0.0240732\\
0.474	-0.0240798\\
0.476	-0.0240864\\
0.478	-0.0240932\\
0.48	-0.0240999\\
0.482	-0.0241068\\
0.484	-0.0241137\\
0.486	-0.0241206\\
0.488	-0.0241275\\
0.49	-0.0241345\\
0.492	-0.0241415\\
0.494	-0.0241486\\
0.496	-0.0241556\\
0.498	-0.0241626\\
0.5	-0.0241697\\
0.502	-0.0241767\\
0.504	-0.0241837\\
0.506	-0.0241906\\
0.508	-0.0241975\\
0.51	-0.0242044\\
0.512	-0.0242112\\
0.514	-0.024218\\
0.516	-0.0242246\\
0.518	-0.0242312\\
0.52	-0.0242378\\
0.522	-0.0242442\\
0.524	-0.0242506\\
0.526	-0.0242568\\
0.528	-0.0242629\\
0.53	-0.024269\\
0.532	-0.0242749\\
0.534	-0.0242807\\
0.536	-0.0242863\\
0.538	-0.0242919\\
0.54	-0.0242973\\
0.542	-0.0243026\\
0.544	-0.0243077\\
0.546	-0.0243127\\
0.548	-0.0243176\\
0.55	-0.0243223\\
0.552	-0.0243269\\
0.554	-0.0243314\\
0.556	-0.0243357\\
0.558	-0.0243399\\
0.56	-0.0243439\\
0.562	-0.0243478\\
0.564	-0.0243516\\
0.566	-0.0243553\\
0.568	-0.0243588\\
0.57	-0.0243622\\
0.572	-0.0243655\\
0.574	-0.0243686\\
0.576	-0.0243717\\
0.578	-0.0243746\\
0.58	-0.0243775\\
0.582	-0.0243802\\
0.584	-0.0243829\\
0.586	-0.0243854\\
0.588	-0.0243879\\
0.59	-0.0243903\\
0.592	-0.0243926\\
0.594	-0.0243949\\
0.596	-0.0243971\\
0.598	-0.0243992\\
0.6	-0.0244013\\
0.602	-0.0244033\\
0.604	-0.0244053\\
0.606	-0.0244072\\
0.608	-0.0244092\\
0.61	-0.0244111\\
0.612	-0.0244129\\
0.614	-0.0244148\\
0.616	-0.0244166\\
0.618	-0.0244185\\
0.62	-0.0244203\\
0.622	-0.0244221\\
0.624	-0.024424\\
0.626	-0.0244258\\
0.628	-0.0244276\\
0.63	-0.0244295\\
0.632	-0.0244314\\
0.634	-0.0244333\\
0.636	-0.0244352\\
0.638	-0.0244371\\
0.64	-0.0244391\\
0.642	-0.0244411\\
0.644	-0.0244431\\
0.646	-0.0244451\\
0.648	-0.0244472\\
0.65	-0.0244493\\
0.652	-0.0244514\\
0.654	-0.0244536\\
0.656	-0.0244558\\
0.658	-0.024458\\
0.66	-0.0244603\\
0.662	-0.0244625\\
0.664	-0.0244648\\
0.666	-0.0244672\\
0.668	-0.0244695\\
0.67	-0.0244719\\
0.672	-0.0244743\\
0.674	-0.0244767\\
0.676	-0.0244792\\
0.678	-0.0244816\\
0.68	-0.0244841\\
0.682	-0.0244866\\
0.684	-0.0244891\\
0.686	-0.0244916\\
0.688	-0.0244941\\
0.69	-0.0244967\\
0.692	-0.0244992\\
0.694	-0.0245017\\
0.696	-0.0245042\\
0.698	-0.0245068\\
0.7	-0.0245093\\
0.702	-0.0245118\\
0.704	-0.0245143\\
0.706	-0.0245168\\
0.708	-0.0245193\\
0.71	-0.0245217\\
0.712	-0.0245242\\
0.714	-0.0245266\\
0.716	-0.024529\\
0.718	-0.0245314\\
0.72	-0.0245337\\
0.722	-0.024536\\
0.724	-0.0245383\\
0.726	-0.0245406\\
0.728	-0.0245429\\
0.73	-0.0245451\\
0.732	-0.0245473\\
0.734	-0.0245494\\
0.736	-0.0245515\\
0.738	-0.0245536\\
0.74	-0.0245557\\
0.742	-0.0245577\\
0.744	-0.0245597\\
0.746	-0.0245617\\
0.748	-0.0245636\\
0.75	-0.0245655\\
0.752	-0.0245673\\
0.754	-0.0245691\\
0.756	-0.0245709\\
0.758	-0.0245726\\
0.76	-0.0245744\\
0.762	-0.024576\\
0.764	-0.0245777\\
0.766	-0.0245793\\
0.768	-0.0245809\\
0.77	-0.0245824\\
0.772	-0.024584\\
0.774	-0.0245854\\
0.776	-0.0245869\\
0.778	-0.0245883\\
0.78	-0.0245897\\
0.782	-0.0245911\\
0.784	-0.0245925\\
0.786	-0.0245938\\
0.788	-0.0245951\\
0.79	-0.0245964\\
0.792	-0.0245976\\
0.794	-0.0245989\\
0.796	-0.0246001\\
0.798	-0.0246012\\
0.8	-0.0246024\\
0.802	-0.0246036\\
0.804	-0.0246047\\
0.806	-0.0246058\\
0.808	-0.0246069\\
0.81	-0.024608\\
0.812	-0.0246091\\
0.814	-0.0246101\\
0.816	-0.0246111\\
0.818	-0.0246122\\
0.82	-0.0246132\\
0.822	-0.0246142\\
0.824	-0.0246152\\
0.826	-0.0246161\\
0.828	-0.0246171\\
0.83	-0.0246181\\
0.832	-0.024619\\
0.834	-0.02462\\
0.836	-0.0246209\\
0.838	-0.0246218\\
0.84	-0.0246228\\
0.842	-0.0246237\\
0.844	-0.0246246\\
0.846	-0.0246255\\
0.848	-0.0246264\\
0.85	-0.0246273\\
0.852	-0.0246282\\
0.854	-0.0246291\\
0.856	-0.0246299\\
0.858	-0.0246308\\
0.86	-0.0246317\\
0.862	-0.0246326\\
0.864	-0.0246334\\
0.866	-0.0246343\\
0.868	-0.0246352\\
0.87	-0.024636\\
0.872	-0.0246369\\
0.874	-0.0246378\\
0.876	-0.0246386\\
0.878	-0.0246395\\
0.88	-0.0246403\\
0.882	-0.0246412\\
0.884	-0.024642\\
0.886	-0.0246429\\
0.888	-0.0246437\\
0.89	-0.0246446\\
0.892	-0.0246454\\
0.894	-0.0246463\\
0.896	-0.0246471\\
0.898	-0.024648\\
0.9	-0.0246488\\
0.902	-0.0246497\\
0.904	-0.0246505\\
0.906	-0.0246513\\
0.908	-0.0246522\\
0.91	-0.024653\\
0.912	-0.0246538\\
0.914	-0.0246547\\
0.916	-0.0246555\\
0.918	-0.0246563\\
0.92	-0.0246572\\
0.922	-0.024658\\
0.924	-0.0246588\\
0.926	-0.0246596\\
0.928	-0.0246605\\
0.93	-0.0246613\\
0.932	-0.0246621\\
0.934	-0.0246629\\
0.936	-0.0246637\\
0.938	-0.0246645\\
0.94	-0.0246653\\
0.942	-0.0246661\\
0.944	-0.0246669\\
0.946	-0.0246677\\
0.948	-0.0246685\\
0.95	-0.0246693\\
0.952	-0.02467\\
0.954	-0.0246708\\
0.956	-0.0246716\\
0.958	-0.0246723\\
0.96	-0.0246731\\
0.962	-0.0246739\\
0.964	-0.0246746\\
0.966	-0.0246754\\
0.968	-0.0246761\\
0.97	-0.0246768\\
0.972	-0.0246776\\
0.974	-0.0246783\\
0.976	-0.024679\\
0.978	-0.0246797\\
0.98	-0.0246804\\
0.982	-0.0246811\\
0.984	-0.0246818\\
0.986	-0.0246825\\
0.988	-0.0246832\\
0.99	-0.0246839\\
0.992	-0.0246846\\
0.994	-0.0246852\\
0.996	-0.0246859\\
0.998	-0.0246865\\
1	-0.0246872\\
1.002	-0.0246878\\
1.004	-0.0246885\\
1.006	-0.0246891\\
1.008	-0.0246897\\
1.01	-0.0246903\\
1.012	-0.0246909\\
1.014	-0.0246915\\
1.016	-0.0246921\\
1.018	-0.0246927\\
1.02	-0.0246933\\
1.022	-0.0246939\\
1.024	-0.0246945\\
1.026	-0.024695\\
1.028	-0.0246956\\
1.03	-0.0246961\\
1.032	-0.0246967\\
1.034	-0.0246972\\
1.036	-0.0246978\\
1.038	-0.0246983\\
1.04	-0.0246988\\
1.042	-0.0246993\\
1.044	-0.0246999\\
1.046	-0.0247004\\
1.048	-0.0247009\\
1.05	-0.0247014\\
1.052	-0.0247019\\
1.054	-0.0247024\\
1.056	-0.0247028\\
1.058	-0.0247033\\
1.06	-0.0247038\\
1.062	-0.0247043\\
1.064	-0.0247047\\
1.066	-0.0247052\\
1.068	-0.0247056\\
1.07	-0.0247061\\
1.072	-0.0247065\\
1.074	-0.024707\\
1.076	-0.0247074\\
1.078	-0.0247079\\
1.08	-0.0247083\\
1.082	-0.0247087\\
1.084	-0.0247092\\
1.086	-0.0247096\\
1.088	-0.02471\\
1.09	-0.0247104\\
1.092	-0.0247108\\
1.094	-0.0247112\\
1.096	-0.0247116\\
1.098	-0.024712\\
1.1	-0.0247124\\
1.102	-0.0247128\\
1.104	-0.0247132\\
1.106	-0.0247136\\
1.108	-0.024714\\
1.11	-0.0247144\\
1.112	-0.0247148\\
1.114	-0.0247152\\
1.116	-0.0247155\\
1.118	-0.0247159\\
1.12	-0.0247163\\
1.122	-0.0247167\\
1.124	-0.024717\\
1.126	-0.0247174\\
1.128	-0.0247177\\
1.13	-0.0247181\\
1.132	-0.0247185\\
1.134	-0.0247188\\
1.136	-0.0247192\\
1.138	-0.0247195\\
1.14	-0.0247199\\
1.142	-0.0247202\\
1.144	-0.0247206\\
1.146	-0.0247209\\
1.148	-0.0247213\\
1.15	-0.0247216\\
1.152	-0.0247219\\
1.154	-0.0247223\\
1.156	-0.0247226\\
1.158	-0.0247229\\
1.16	-0.0247233\\
1.162	-0.0247236\\
1.164	-0.0247239\\
1.166	-0.0247243\\
1.168	-0.0247246\\
1.17	-0.0247249\\
1.172	-0.0247252\\
1.174	-0.0247255\\
1.176	-0.0247259\\
1.178	-0.0247262\\
1.18	-0.0247265\\
1.182	-0.0247268\\
1.184	-0.0247271\\
1.186	-0.0247274\\
1.188	-0.0247277\\
1.19	-0.024728\\
1.192	-0.0247283\\
1.194	-0.0247286\\
1.196	-0.0247289\\
1.198	-0.0247292\\
1.2	-0.0247295\\
1.202	-0.0247298\\
1.204	-0.0247301\\
1.206	-0.0247304\\
1.208	-0.0247307\\
1.21	-0.024731\\
1.212	-0.0247313\\
1.214	-0.0247316\\
1.216	-0.0247318\\
1.218	-0.0247321\\
1.22	-0.0247324\\
1.222	-0.0247327\\
1.224	-0.024733\\
1.226	-0.0247332\\
1.228	-0.0247335\\
1.23	-0.0247338\\
1.232	-0.0247341\\
1.234	-0.0247343\\
1.236	-0.0247346\\
1.238	-0.0247349\\
1.24	-0.0247351\\
1.242	-0.0247354\\
1.244	-0.0247356\\
1.246	-0.0247359\\
1.248	-0.0247362\\
1.25	-0.0247364\\
1.252	-0.0247367\\
1.254	-0.0247369\\
1.256	-0.0247372\\
1.258	-0.0247374\\
1.26	-0.0247377\\
1.262	-0.0247379\\
1.264	-0.0247382\\
1.266	-0.0247384\\
1.268	-0.0247386\\
1.27	-0.0247389\\
1.272	-0.0247391\\
1.274	-0.0247393\\
1.276	-0.0247396\\
1.278	-0.0247398\\
1.28	-0.02474\\
1.282	-0.0247403\\
1.284	-0.0247405\\
1.286	-0.0247407\\
1.288	-0.0247409\\
1.29	-0.0247412\\
1.292	-0.0247414\\
1.294	-0.0247416\\
1.296	-0.0247418\\
1.298	-0.024742\\
1.3	-0.0247422\\
1.302	-0.0247424\\
1.304	-0.0247426\\
1.306	-0.0247428\\
1.308	-0.0247431\\
1.31	-0.0247433\\
1.312	-0.0247435\\
1.314	-0.0247437\\
1.316	-0.0247438\\
1.318	-0.024744\\
1.32	-0.0247442\\
1.322	-0.0247444\\
1.324	-0.0247446\\
1.326	-0.0247448\\
1.328	-0.024745\\
1.33	-0.0247452\\
1.332	-0.0247454\\
1.334	-0.0247455\\
1.336	-0.0247457\\
1.338	-0.0247459\\
1.34	-0.0247461\\
1.342	-0.0247463\\
1.344	-0.0247464\\
1.346	-0.0247466\\
1.348	-0.0247468\\
1.35	-0.024747\\
1.352	-0.0247471\\
1.354	-0.0247473\\
1.356	-0.0247475\\
1.358	-0.0247476\\
1.36	-0.0247478\\
1.362	-0.024748\\
1.364	-0.0247481\\
1.366	-0.0247483\\
1.368	-0.0247484\\
1.37	-0.0247486\\
1.372	-0.0247487\\
1.374	-0.0247489\\
1.376	-0.0247491\\
1.378	-0.0247492\\
1.38	-0.0247494\\
1.382	-0.0247495\\
1.384	-0.0247497\\
1.386	-0.0247498\\
1.388	-0.02475\\
1.39	-0.0247501\\
1.392	-0.0247503\\
1.394	-0.0247504\\
1.396	-0.0247506\\
1.398	-0.0247507\\
1.4	-0.0247508\\
1.402	-0.024751\\
1.404	-0.0247511\\
1.406	-0.0247513\\
1.408	-0.0247514\\
1.41	-0.0247515\\
1.412	-0.0247517\\
1.414	-0.0247518\\
1.416	-0.024752\\
1.418	-0.0247521\\
1.42	-0.0247522\\
1.422	-0.0247524\\
1.424	-0.0247525\\
1.426	-0.0247526\\
1.428	-0.0247528\\
1.43	-0.0247529\\
1.432	-0.024753\\
1.434	-0.0247532\\
1.436	-0.0247533\\
1.438	-0.0247534\\
1.44	-0.0247535\\
1.442	-0.0247537\\
1.444	-0.0247538\\
1.446	-0.0247539\\
1.448	-0.024754\\
1.45	-0.0247542\\
1.452	-0.0247543\\
1.454	-0.0247544\\
1.456	-0.0247545\\
1.458	-0.0247547\\
1.46	-0.0247548\\
1.462	-0.0247549\\
1.464	-0.024755\\
1.466	-0.0247552\\
1.468	-0.0247553\\
1.47	-0.0247554\\
1.472	-0.0247555\\
1.474	-0.0247556\\
1.476	-0.0247557\\
1.478	-0.0247559\\
1.48	-0.024756\\
1.482	-0.0247561\\
1.484	-0.0247562\\
1.486	-0.0247563\\
1.488	-0.0247564\\
1.49	-0.0247565\\
1.492	-0.0247567\\
1.494	-0.0247568\\
1.496	-0.0247569\\
1.498	-0.024757\\
1.5	-0.0247571\\
1.502	-0.0247572\\
1.504	-0.0247573\\
1.506	-0.0247574\\
1.508	-0.0247575\\
1.51	-0.0247576\\
1.512	-0.0247577\\
1.514	-0.0247578\\
1.516	-0.0247579\\
1.518	-0.024758\\
1.52	-0.0247581\\
1.522	-0.0247582\\
1.524	-0.0247583\\
1.526	-0.0247584\\
1.528	-0.0247585\\
1.53	-0.0247586\\
1.532	-0.0247587\\
1.534	-0.0247588\\
1.536	-0.0247589\\
1.538	-0.024759\\
1.54	-0.0247591\\
1.542	-0.0247592\\
1.544	-0.0247593\\
1.546	-0.0247594\\
1.548	-0.0247595\\
1.55	-0.0247596\\
1.552	-0.0247597\\
1.554	-0.0247598\\
1.556	-0.0247598\\
1.558	-0.0247599\\
1.56	-0.02476\\
1.562	-0.0247601\\
1.564	-0.0247602\\
1.566	-0.0247603\\
1.568	-0.0247604\\
1.57	-0.0247604\\
1.572	-0.0247605\\
1.574	-0.0247606\\
1.576	-0.0247607\\
1.578	-0.0247608\\
1.58	-0.0247608\\
1.582	-0.0247609\\
1.584	-0.024761\\
1.586	-0.0247611\\
1.588	-0.0247612\\
1.59	-0.0247612\\
1.592	-0.0247613\\
1.594	-0.0247614\\
1.596	-0.0247615\\
1.598	-0.0247615\\
1.6	-0.0247616\\
1.602	-0.0247617\\
1.604	-0.0247618\\
1.606	-0.0247618\\
1.608	-0.0247619\\
1.61	-0.024762\\
1.612	-0.0247621\\
1.614	-0.0247621\\
1.616	-0.0247622\\
1.618	-0.0247623\\
1.62	-0.0247623\\
1.622	-0.0247624\\
1.624	-0.0247625\\
1.626	-0.0247625\\
1.628	-0.0247626\\
1.63	-0.0247627\\
1.632	-0.0247627\\
1.634	-0.0247628\\
1.636	-0.0247629\\
1.638	-0.0247629\\
1.64	-0.024763\\
1.642	-0.0247631\\
1.644	-0.0247631\\
1.646	-0.0247632\\
1.648	-0.0247633\\
1.65	-0.0247633\\
1.652	-0.0247634\\
1.654	-0.0247635\\
1.656	-0.0247635\\
1.658	-0.0247636\\
1.66	-0.0247636\\
1.662	-0.0247637\\
1.664	-0.0247638\\
1.666	-0.0247638\\
1.668	-0.0247639\\
1.67	-0.024764\\
1.672	-0.024764\\
1.674	-0.0247641\\
1.676	-0.0247641\\
1.678	-0.0247642\\
1.68	-0.0247642\\
1.682	-0.0247643\\
1.684	-0.0247644\\
1.686	-0.0247644\\
1.688	-0.0247645\\
1.69	-0.0247645\\
1.692	-0.0247646\\
1.694	-0.0247646\\
1.696	-0.0247647\\
1.698	-0.0247648\\
1.7	-0.0247648\\
1.702	-0.0247649\\
1.704	-0.0247649\\
1.706	-0.024765\\
1.708	-0.024765\\
1.71	-0.0247651\\
1.712	-0.0247651\\
1.714	-0.0247652\\
1.716	-0.0247652\\
1.718	-0.0247653\\
1.72	-0.0247653\\
1.722	-0.0247654\\
1.724	-0.0247654\\
1.726	-0.0247655\\
1.728	-0.0247656\\
1.73	-0.0247656\\
1.732	-0.0247657\\
1.734	-0.0247657\\
1.736	-0.0247657\\
1.738	-0.0247658\\
1.74	-0.0247658\\
1.742	-0.0247659\\
1.744	-0.0247659\\
1.746	-0.024766\\
1.748	-0.024766\\
1.75	-0.0247661\\
1.752	-0.0247661\\
1.754	-0.0247662\\
1.756	-0.0247662\\
1.758	-0.0247663\\
1.76	-0.0247663\\
1.762	-0.0247664\\
1.764	-0.0247664\\
1.766	-0.0247665\\
1.768	-0.0247665\\
1.77	-0.0247665\\
1.772	-0.0247666\\
1.774	-0.0247666\\
1.776	-0.0247667\\
1.778	-0.0247667\\
1.78	-0.0247668\\
1.782	-0.0247668\\
1.784	-0.0247668\\
1.786	-0.0247669\\
1.788	-0.0247669\\
1.79	-0.024767\\
1.792	-0.024767\\
1.794	-0.024767\\
1.796	-0.0247671\\
1.798	-0.0247671\\
1.8	-0.0247672\\
1.802	-0.0247672\\
1.804	-0.0247672\\
1.806	-0.0247673\\
1.808	-0.0247673\\
1.81	-0.0247674\\
1.812	-0.0247674\\
1.814	-0.0247674\\
1.816	-0.0247675\\
1.818	-0.0247675\\
1.82	-0.0247675\\
1.822	-0.0247676\\
1.824	-0.0247676\\
1.826	-0.0247676\\
1.828	-0.0247677\\
1.83	-0.0247677\\
1.832	-0.0247678\\
1.834	-0.0247678\\
1.836	-0.0247678\\
1.838	-0.0247679\\
1.84	-0.0247679\\
1.842	-0.0247679\\
1.844	-0.024768\\
1.846	-0.024768\\
1.848	-0.024768\\
1.85	-0.0247681\\
1.852	-0.0247681\\
1.854	-0.0247681\\
1.856	-0.0247682\\
1.858	-0.0247682\\
1.86	-0.0247682\\
1.862	-0.0247683\\
1.864	-0.0247683\\
1.866	-0.0247683\\
1.868	-0.0247683\\
1.87	-0.0247684\\
1.872	-0.0247684\\
1.874	-0.0247684\\
1.876	-0.0247685\\
1.878	-0.0247685\\
1.88	-0.0247685\\
1.882	-0.0247686\\
1.884	-0.0247686\\
1.886	-0.0247686\\
1.888	-0.0247686\\
1.89	-0.0247687\\
1.892	-0.0247687\\
1.894	-0.0247687\\
1.896	-0.0247688\\
1.898	-0.0247688\\
1.9	-0.0247688\\
1.902	-0.0247688\\
1.904	-0.0247689\\
1.906	-0.0247689\\
1.908	-0.0247689\\
1.91	-0.024769\\
1.912	-0.024769\\
1.914	-0.024769\\
1.916	-0.024769\\
1.918	-0.0247691\\
1.92	-0.0247691\\
1.922	-0.0247691\\
1.924	-0.0247691\\
1.926	-0.0247692\\
1.928	-0.0247692\\
1.93	-0.0247692\\
1.932	-0.0247693\\
1.934	-0.0247693\\
1.936	-0.0247693\\
1.938	-0.0247693\\
1.94	-0.0247694\\
1.942	-0.0247694\\
1.944	-0.0247694\\
1.946	-0.0247694\\
1.948	-0.0247695\\
1.95	-0.0247695\\
1.952	-0.0247695\\
1.954	-0.0247695\\
1.956	-0.0247695\\
1.958	-0.0247696\\
1.96	-0.0247696\\
1.962	-0.0247696\\
1.964	-0.0247696\\
1.966	-0.0247697\\
1.968	-0.0247697\\
1.97	-0.0247697\\
1.972	-0.0247697\\
1.974	-0.0247698\\
1.976	-0.0247698\\
1.978	-0.0247698\\
1.98	-0.0247698\\
1.982	-0.0247698\\
1.984	-0.0247699\\
1.986	-0.0247699\\
1.988	-0.0247699\\
1.99	-0.0247699\\
1.992	-0.02477\\
1.994	-0.02477\\
1.996	-0.02477\\
1.998	-0.02477\\
2	-0.02477\\
};
\addplot [color=mycolor8, line width=1.0pt, forget plot]
  table[row sep=crcr]{%
-0.024	0\\
-0.022	0\\
-0.02	0\\
-0.018	0\\
-0.016	0\\
-0.014	0\\
-0.012	0\\
-0.01	0\\
-0.008	0\\
-0.006	0\\
-0.004	0\\
-0.002	0\\
0	0\\
0.002	-0.0018465\\
0.004	-0.00342636\\
0.006	-0.00478689\\
0.008	-0.0059674\\
0.01	-0.00700049\\
0.012	-0.00791307\\
0.014	-0.00872724\\
0.016	-0.00946111\\
0.018	-0.0101294\\
0.02	-0.010744\\
0.022	-0.0113146\\
0.024	-0.0118486\\
0.026	-0.0123523\\
0.028	-0.0128302\\
0.03	-0.0132859\\
0.032	-0.0137224\\
0.034	-0.0141415\\
0.036	-0.0145448\\
0.038	-0.0149333\\
0.04	-0.0153078\\
0.042	-0.0156688\\
0.044	-0.0160166\\
0.046	-0.0163514\\
0.048	-0.0166732\\
0.05	-0.0169823\\
0.052	-0.0172786\\
0.054	-0.0175623\\
0.056	-0.0178334\\
0.058	-0.0180922\\
0.06	-0.0183387\\
0.062	-0.0185733\\
0.064	-0.0187962\\
0.066	-0.0190076\\
0.068	-0.0192079\\
0.07	-0.0193975\\
0.072	-0.0195767\\
0.074	-0.019746\\
0.076	-0.0199057\\
0.078	-0.0200563\\
0.08	-0.0201982\\
0.082	-0.0203318\\
0.084	-0.0204575\\
0.086	-0.0205758\\
0.088	-0.0206871\\
0.09	-0.0207918\\
0.092	-0.0208903\\
0.094	-0.0209829\\
0.096	-0.02107\\
0.098	-0.021152\\
0.1	-0.0212293\\
0.102	-0.0213021\\
0.104	-0.0213707\\
0.106	-0.0214354\\
0.108	-0.0214966\\
0.11	-0.0215544\\
0.112	-0.0216091\\
0.114	-0.0216608\\
0.116	-0.0217099\\
0.118	-0.0217564\\
0.12	-0.0218006\\
0.122	-0.0218425\\
0.124	-0.0218824\\
0.126	-0.0219202\\
0.128	-0.0219561\\
0.13	-0.0219902\\
0.132	-0.0220225\\
0.134	-0.0220531\\
0.136	-0.0220821\\
0.138	-0.0221094\\
0.14	-0.0221351\\
0.142	-0.0221591\\
0.144	-0.0221816\\
0.146	-0.0222025\\
0.148	-0.0222218\\
0.15	-0.0222396\\
0.152	-0.0222558\\
0.154	-0.0222704\\
0.156	-0.0222835\\
0.158	-0.022295\\
0.16	-0.0223051\\
0.162	-0.0223136\\
0.164	-0.0223207\\
0.166	-0.0223264\\
0.168	-0.0223307\\
0.17	-0.0223337\\
0.172	-0.0223355\\
0.174	-0.0223361\\
0.176	-0.0223356\\
0.178	-0.0223341\\
0.18	-0.0223316\\
0.182	-0.0223283\\
0.184	-0.0223243\\
0.186	-0.0223197\\
0.188	-0.0223145\\
0.19	-0.0223089\\
0.192	-0.022303\\
0.194	-0.022297\\
0.196	-0.0222908\\
0.198	-0.0222847\\
0.2	-0.0222787\\
0.202	-0.022273\\
0.204	-0.0222676\\
0.206	-0.0222626\\
0.208	-0.0222582\\
0.21	-0.0222544\\
0.212	-0.0222514\\
0.214	-0.0222491\\
0.216	-0.0222476\\
0.218	-0.0222471\\
0.22	-0.0222476\\
0.222	-0.0222491\\
0.224	-0.0222517\\
0.226	-0.0222553\\
0.228	-0.0222601\\
0.23	-0.0222661\\
0.232	-0.0222732\\
0.234	-0.0222815\\
0.236	-0.022291\\
0.238	-0.0223016\\
0.24	-0.0223134\\
0.242	-0.0223264\\
0.244	-0.0223404\\
0.246	-0.0223556\\
0.248	-0.0223718\\
0.25	-0.0223891\\
0.252	-0.0224073\\
0.254	-0.0224265\\
0.256	-0.0224466\\
0.258	-0.0224675\\
0.26	-0.0224893\\
0.262	-0.0225118\\
0.264	-0.0225349\\
0.266	-0.0225588\\
0.268	-0.0225832\\
0.27	-0.0226082\\
0.272	-0.0226336\\
0.274	-0.0226595\\
0.276	-0.0226858\\
0.278	-0.0227124\\
0.28	-0.0227392\\
0.282	-0.0227663\\
0.284	-0.0227935\\
0.286	-0.0228209\\
0.288	-0.0228483\\
0.29	-0.0228757\\
0.292	-0.0229032\\
0.294	-0.0229305\\
0.296	-0.0229578\\
0.298	-0.022985\\
0.3	-0.0230119\\
0.302	-0.0230387\\
0.304	-0.0230652\\
0.306	-0.0230915\\
0.308	-0.0231174\\
0.31	-0.0231431\\
0.312	-0.0231684\\
0.314	-0.0231933\\
0.316	-0.0232179\\
0.318	-0.0232421\\
0.32	-0.0232658\\
0.322	-0.0232891\\
0.324	-0.023312\\
0.326	-0.0233344\\
0.328	-0.0233563\\
0.33	-0.0233777\\
0.332	-0.0233987\\
0.334	-0.0234192\\
0.336	-0.0234392\\
0.338	-0.0234587\\
0.34	-0.0234777\\
0.342	-0.0234962\\
0.344	-0.0235142\\
0.346	-0.0235317\\
0.348	-0.0235488\\
0.35	-0.0235653\\
0.352	-0.0235813\\
0.354	-0.0235969\\
0.356	-0.023612\\
0.358	-0.0236266\\
0.36	-0.0236408\\
0.362	-0.0236545\\
0.364	-0.0236678\\
0.366	-0.0236806\\
0.368	-0.023693\\
0.37	-0.023705\\
0.372	-0.0237166\\
0.374	-0.0237278\\
0.376	-0.0237386\\
0.378	-0.0237491\\
0.38	-0.0237592\\
0.382	-0.0237689\\
0.384	-0.0237783\\
0.386	-0.0237874\\
0.388	-0.0237963\\
0.39	-0.0238048\\
0.392	-0.023813\\
0.394	-0.023821\\
0.396	-0.0238288\\
0.398	-0.0238363\\
0.4	-0.0238436\\
0.402	-0.0238507\\
0.404	-0.0238576\\
0.406	-0.0238643\\
0.408	-0.0238709\\
0.41	-0.0238774\\
0.412	-0.0238837\\
0.414	-0.0238899\\
0.416	-0.023896\\
0.418	-0.023902\\
0.42	-0.0239079\\
0.422	-0.0239138\\
0.424	-0.0239196\\
0.426	-0.0239254\\
0.428	-0.0239311\\
0.43	-0.0239368\\
0.432	-0.0239425\\
0.434	-0.0239483\\
0.436	-0.023954\\
0.438	-0.0239597\\
0.44	-0.0239655\\
0.442	-0.0239713\\
0.444	-0.0239771\\
0.446	-0.023983\\
0.448	-0.0239889\\
0.45	-0.0239949\\
0.452	-0.0240009\\
0.454	-0.024007\\
0.456	-0.0240132\\
0.458	-0.0240194\\
0.46	-0.0240256\\
0.462	-0.024032\\
0.464	-0.0240384\\
0.466	-0.0240448\\
0.468	-0.0240514\\
0.47	-0.0240579\\
0.472	-0.0240646\\
0.474	-0.0240713\\
0.476	-0.024078\\
0.478	-0.0240848\\
0.48	-0.0240916\\
0.482	-0.0240984\\
0.484	-0.0241053\\
0.486	-0.0241122\\
0.488	-0.0241191\\
0.49	-0.0241261\\
0.492	-0.024133\\
0.494	-0.02414\\
0.496	-0.0241469\\
0.498	-0.0241538\\
0.5	-0.0241607\\
0.502	-0.0241676\\
0.504	-0.0241745\\
0.506	-0.0241813\\
0.508	-0.024188\\
0.51	-0.0241947\\
0.512	-0.0242013\\
0.514	-0.0242079\\
0.516	-0.0242144\\
0.518	-0.0242208\\
0.52	-0.0242272\\
0.522	-0.0242334\\
0.524	-0.0242396\\
0.526	-0.0242456\\
0.528	-0.0242516\\
0.53	-0.0242574\\
0.532	-0.0242632\\
0.534	-0.0242688\\
0.536	-0.0242743\\
0.538	-0.0242797\\
0.54	-0.0242849\\
0.542	-0.02429\\
0.544	-0.024295\\
0.546	-0.0242999\\
0.548	-0.0243047\\
0.55	-0.0243093\\
0.552	-0.0243138\\
0.554	-0.0243181\\
0.556	-0.0243223\\
0.558	-0.0243265\\
0.56	-0.0243304\\
0.562	-0.0243343\\
0.564	-0.024338\\
0.566	-0.0243416\\
0.568	-0.0243452\\
0.57	-0.0243485\\
0.572	-0.0243518\\
0.574	-0.024355\\
0.576	-0.0243581\\
0.578	-0.0243611\\
0.58	-0.0243639\\
0.582	-0.0243667\\
0.584	-0.0243695\\
0.586	-0.0243721\\
0.588	-0.0243746\\
0.59	-0.0243771\\
0.592	-0.0243795\\
0.594	-0.0243819\\
0.596	-0.0243842\\
0.598	-0.0243865\\
0.6	-0.0243887\\
0.602	-0.0243908\\
0.604	-0.024393\\
0.606	-0.0243951\\
0.608	-0.0243971\\
0.61	-0.0243992\\
0.612	-0.0244012\\
0.614	-0.0244032\\
0.616	-0.0244052\\
0.618	-0.0244072\\
0.62	-0.0244092\\
0.622	-0.0244112\\
0.624	-0.0244132\\
0.626	-0.0244152\\
0.628	-0.0244172\\
0.63	-0.0244192\\
0.632	-0.0244213\\
0.634	-0.0244233\\
0.636	-0.0244254\\
0.638	-0.0244275\\
0.64	-0.0244296\\
0.642	-0.0244317\\
0.644	-0.0244339\\
0.646	-0.0244361\\
0.648	-0.0244383\\
0.65	-0.0244405\\
0.652	-0.0244428\\
0.654	-0.024445\\
0.656	-0.0244473\\
0.658	-0.0244497\\
0.66	-0.024452\\
0.662	-0.0244544\\
0.664	-0.0244568\\
0.666	-0.0244592\\
0.668	-0.0244616\\
0.67	-0.0244641\\
0.672	-0.0244666\\
0.674	-0.024469\\
0.676	-0.0244715\\
0.678	-0.024474\\
0.68	-0.0244766\\
0.682	-0.0244791\\
0.684	-0.0244816\\
0.686	-0.0244842\\
0.688	-0.0244867\\
0.69	-0.0244892\\
0.692	-0.0244918\\
0.694	-0.0244943\\
0.696	-0.0244968\\
0.698	-0.0244993\\
0.7	-0.0245018\\
0.702	-0.0245043\\
0.704	-0.0245068\\
0.706	-0.0245093\\
0.708	-0.0245117\\
0.71	-0.0245142\\
0.712	-0.0245166\\
0.714	-0.024519\\
0.716	-0.0245214\\
0.718	-0.0245237\\
0.72	-0.024526\\
0.722	-0.0245283\\
0.724	-0.0245306\\
0.726	-0.0245328\\
0.728	-0.024535\\
0.73	-0.0245372\\
0.732	-0.0245394\\
0.734	-0.0245415\\
0.736	-0.0245436\\
0.738	-0.0245456\\
0.74	-0.0245476\\
0.742	-0.0245496\\
0.744	-0.0245516\\
0.746	-0.0245535\\
0.748	-0.0245554\\
0.75	-0.0245573\\
0.752	-0.0245591\\
0.754	-0.0245609\\
0.756	-0.0245626\\
0.758	-0.0245644\\
0.76	-0.0245661\\
0.762	-0.0245677\\
0.764	-0.0245694\\
0.766	-0.024571\\
0.768	-0.0245726\\
0.77	-0.0245741\\
0.772	-0.0245756\\
0.774	-0.0245771\\
0.776	-0.0245786\\
0.778	-0.0245801\\
0.78	-0.0245815\\
0.782	-0.0245829\\
0.784	-0.0245842\\
0.786	-0.0245856\\
0.788	-0.0245869\\
0.79	-0.0245882\\
0.792	-0.0245895\\
0.794	-0.0245908\\
0.796	-0.024592\\
0.798	-0.0245932\\
0.8	-0.0245944\\
0.802	-0.0245956\\
0.804	-0.0245968\\
0.806	-0.024598\\
0.808	-0.0245991\\
0.81	-0.0246002\\
0.812	-0.0246014\\
0.814	-0.0246025\\
0.816	-0.0246036\\
0.818	-0.0246046\\
0.82	-0.0246057\\
0.822	-0.0246068\\
0.824	-0.0246078\\
0.826	-0.0246088\\
0.828	-0.0246099\\
0.83	-0.0246109\\
0.832	-0.0246119\\
0.834	-0.0246129\\
0.836	-0.0246139\\
0.838	-0.0246149\\
0.84	-0.0246159\\
0.842	-0.0246168\\
0.844	-0.0246178\\
0.846	-0.0246188\\
0.848	-0.0246197\\
0.85	-0.0246207\\
0.852	-0.0246216\\
0.854	-0.0246226\\
0.856	-0.0246235\\
0.858	-0.0246244\\
0.86	-0.0246254\\
0.862	-0.0246263\\
0.864	-0.0246272\\
0.866	-0.0246281\\
0.868	-0.024629\\
0.87	-0.0246299\\
0.872	-0.0246309\\
0.874	-0.0246318\\
0.876	-0.0246327\\
0.878	-0.0246336\\
0.88	-0.0246345\\
0.882	-0.0246353\\
0.884	-0.0246362\\
0.886	-0.0246371\\
0.888	-0.024638\\
0.89	-0.0246389\\
0.892	-0.0246398\\
0.894	-0.0246406\\
0.896	-0.0246415\\
0.898	-0.0246424\\
0.9	-0.0246433\\
0.902	-0.0246441\\
0.904	-0.024645\\
0.906	-0.0246459\\
0.908	-0.0246467\\
0.91	-0.0246476\\
0.912	-0.0246484\\
0.914	-0.0246493\\
0.916	-0.0246501\\
0.918	-0.024651\\
0.92	-0.0246518\\
0.922	-0.0246527\\
0.924	-0.0246535\\
0.926	-0.0246543\\
0.928	-0.0246552\\
0.93	-0.024656\\
0.932	-0.0246568\\
0.934	-0.0246576\\
0.936	-0.0246585\\
0.938	-0.0246593\\
0.94	-0.0246601\\
0.942	-0.0246609\\
0.944	-0.0246617\\
0.946	-0.0246625\\
0.948	-0.0246633\\
0.95	-0.0246641\\
0.952	-0.0246649\\
0.954	-0.0246657\\
0.956	-0.0246664\\
0.958	-0.0246672\\
0.96	-0.024668\\
0.962	-0.0246687\\
0.964	-0.0246695\\
0.966	-0.0246703\\
0.968	-0.024671\\
0.97	-0.0246718\\
0.972	-0.0246725\\
0.974	-0.0246732\\
0.976	-0.024674\\
0.978	-0.0246747\\
0.98	-0.0246754\\
0.982	-0.0246761\\
0.984	-0.0246768\\
0.986	-0.0246775\\
0.988	-0.0246782\\
0.99	-0.0246789\\
0.992	-0.0246796\\
0.994	-0.0246803\\
0.996	-0.0246809\\
0.998	-0.0246816\\
1	-0.0246823\\
1.002	-0.0246829\\
1.004	-0.0246836\\
1.006	-0.0246842\\
1.008	-0.0246848\\
1.01	-0.0246855\\
1.012	-0.0246861\\
1.014	-0.0246867\\
1.016	-0.0246873\\
1.018	-0.0246879\\
1.02	-0.0246885\\
1.022	-0.0246891\\
1.024	-0.0246897\\
1.026	-0.0246903\\
1.028	-0.0246909\\
1.03	-0.0246915\\
1.032	-0.024692\\
1.034	-0.0246926\\
1.036	-0.0246932\\
1.038	-0.0246937\\
1.04	-0.0246943\\
1.042	-0.0246948\\
1.044	-0.0246953\\
1.046	-0.0246959\\
1.048	-0.0246964\\
1.05	-0.0246969\\
1.052	-0.0246974\\
1.054	-0.0246979\\
1.056	-0.0246984\\
1.058	-0.0246989\\
1.06	-0.0246994\\
1.062	-0.0246999\\
1.064	-0.0247004\\
1.066	-0.0247009\\
1.068	-0.0247014\\
1.07	-0.0247019\\
1.072	-0.0247023\\
1.074	-0.0247028\\
1.076	-0.0247033\\
1.078	-0.0247037\\
1.08	-0.0247042\\
1.082	-0.0247046\\
1.084	-0.0247051\\
1.086	-0.0247055\\
1.088	-0.024706\\
1.09	-0.0247064\\
1.092	-0.0247068\\
1.094	-0.0247073\\
1.096	-0.0247077\\
1.098	-0.0247081\\
1.1	-0.0247085\\
1.102	-0.0247089\\
1.104	-0.0247094\\
1.106	-0.0247098\\
1.108	-0.0247102\\
1.11	-0.0247106\\
1.112	-0.024711\\
1.114	-0.0247114\\
1.116	-0.0247118\\
1.118	-0.0247122\\
1.12	-0.0247126\\
1.122	-0.0247129\\
1.124	-0.0247133\\
1.126	-0.0247137\\
1.128	-0.0247141\\
1.13	-0.0247145\\
1.132	-0.0247149\\
1.134	-0.0247152\\
1.136	-0.0247156\\
1.138	-0.024716\\
1.14	-0.0247163\\
1.142	-0.0247167\\
1.144	-0.0247171\\
1.146	-0.0247174\\
1.148	-0.0247178\\
1.15	-0.0247181\\
1.152	-0.0247185\\
1.154	-0.0247188\\
1.156	-0.0247192\\
1.158	-0.0247195\\
1.16	-0.0247199\\
1.162	-0.0247202\\
1.164	-0.0247206\\
1.166	-0.0247209\\
1.168	-0.0247212\\
1.17	-0.0247216\\
1.172	-0.0247219\\
1.174	-0.0247222\\
1.176	-0.0247226\\
1.178	-0.0247229\\
1.18	-0.0247232\\
1.182	-0.0247235\\
1.184	-0.0247239\\
1.186	-0.0247242\\
1.188	-0.0247245\\
1.19	-0.0247248\\
1.192	-0.0247251\\
1.194	-0.0247255\\
1.196	-0.0247258\\
1.198	-0.0247261\\
1.2	-0.0247264\\
1.202	-0.0247267\\
1.204	-0.024727\\
1.206	-0.0247273\\
1.208	-0.0247276\\
1.21	-0.0247279\\
1.212	-0.0247282\\
1.214	-0.0247285\\
1.216	-0.0247288\\
1.218	-0.0247291\\
1.22	-0.0247294\\
1.222	-0.0247297\\
1.224	-0.02473\\
1.226	-0.0247303\\
1.228	-0.0247305\\
1.23	-0.0247308\\
1.232	-0.0247311\\
1.234	-0.0247314\\
1.236	-0.0247317\\
1.238	-0.0247319\\
1.24	-0.0247322\\
1.242	-0.0247325\\
1.244	-0.0247328\\
1.246	-0.024733\\
1.248	-0.0247333\\
1.25	-0.0247336\\
1.252	-0.0247338\\
1.254	-0.0247341\\
1.256	-0.0247344\\
1.258	-0.0247346\\
1.26	-0.0247349\\
1.262	-0.0247351\\
1.264	-0.0247354\\
1.266	-0.0247356\\
1.268	-0.0247359\\
1.27	-0.0247361\\
1.272	-0.0247364\\
1.274	-0.0247366\\
1.276	-0.0247369\\
1.278	-0.0247371\\
1.28	-0.0247374\\
1.282	-0.0247376\\
1.284	-0.0247378\\
1.286	-0.0247381\\
1.288	-0.0247383\\
1.29	-0.0247385\\
1.292	-0.0247388\\
1.294	-0.024739\\
1.296	-0.0247392\\
1.298	-0.0247395\\
1.3	-0.0247397\\
1.302	-0.0247399\\
1.304	-0.0247401\\
1.306	-0.0247403\\
1.308	-0.0247406\\
1.31	-0.0247408\\
1.312	-0.024741\\
1.314	-0.0247412\\
1.316	-0.0247414\\
1.318	-0.0247416\\
1.32	-0.0247418\\
1.322	-0.024742\\
1.324	-0.0247422\\
1.326	-0.0247424\\
1.328	-0.0247426\\
1.33	-0.0247428\\
1.332	-0.024743\\
1.334	-0.0247432\\
1.336	-0.0247434\\
1.338	-0.0247436\\
1.34	-0.0247438\\
1.342	-0.024744\\
1.344	-0.0247441\\
1.346	-0.0247443\\
1.348	-0.0247445\\
1.35	-0.0247447\\
1.352	-0.0247449\\
1.354	-0.0247451\\
1.356	-0.0247452\\
1.358	-0.0247454\\
1.36	-0.0247456\\
1.362	-0.0247458\\
1.364	-0.0247459\\
1.366	-0.0247461\\
1.368	-0.0247463\\
1.37	-0.0247464\\
1.372	-0.0247466\\
1.374	-0.0247468\\
1.376	-0.0247469\\
1.378	-0.0247471\\
1.38	-0.0247473\\
1.382	-0.0247474\\
1.384	-0.0247476\\
1.386	-0.0247477\\
1.388	-0.0247479\\
1.39	-0.0247481\\
1.392	-0.0247482\\
1.394	-0.0247484\\
1.396	-0.0247485\\
1.398	-0.0247487\\
1.4	-0.0247488\\
1.402	-0.024749\\
1.404	-0.0247491\\
1.406	-0.0247493\\
1.408	-0.0247494\\
1.41	-0.0247496\\
1.412	-0.0247497\\
1.414	-0.0247499\\
1.416	-0.02475\\
1.418	-0.0247502\\
1.42	-0.0247503\\
1.422	-0.0247505\\
1.424	-0.0247506\\
1.426	-0.0247507\\
1.428	-0.0247509\\
1.43	-0.024751\\
1.432	-0.0247512\\
1.434	-0.0247513\\
1.436	-0.0247514\\
1.438	-0.0247516\\
1.44	-0.0247517\\
1.442	-0.0247518\\
1.444	-0.024752\\
1.446	-0.0247521\\
1.448	-0.0247523\\
1.45	-0.0247524\\
1.452	-0.0247525\\
1.454	-0.0247526\\
1.456	-0.0247528\\
1.458	-0.0247529\\
1.46	-0.024753\\
1.462	-0.0247532\\
1.464	-0.0247533\\
1.466	-0.0247534\\
1.468	-0.0247536\\
1.47	-0.0247537\\
1.472	-0.0247538\\
1.474	-0.0247539\\
1.476	-0.0247541\\
1.478	-0.0247542\\
1.48	-0.0247543\\
1.482	-0.0247544\\
1.484	-0.0247546\\
1.486	-0.0247547\\
1.488	-0.0247548\\
1.49	-0.0247549\\
1.492	-0.024755\\
1.494	-0.0247551\\
1.496	-0.0247553\\
1.498	-0.0247554\\
1.5	-0.0247555\\
1.502	-0.0247556\\
1.504	-0.0247557\\
1.506	-0.0247558\\
1.508	-0.024756\\
1.51	-0.0247561\\
1.512	-0.0247562\\
1.514	-0.0247563\\
1.516	-0.0247564\\
1.518	-0.0247565\\
1.52	-0.0247566\\
1.522	-0.0247567\\
1.524	-0.0247568\\
1.526	-0.0247569\\
1.528	-0.024757\\
1.53	-0.0247572\\
1.532	-0.0247573\\
1.534	-0.0247574\\
1.536	-0.0247575\\
1.538	-0.0247576\\
1.54	-0.0247577\\
1.542	-0.0247578\\
1.544	-0.0247579\\
1.546	-0.024758\\
1.548	-0.0247581\\
1.55	-0.0247582\\
1.552	-0.0247583\\
1.554	-0.0247583\\
1.556	-0.0247584\\
1.558	-0.0247585\\
1.56	-0.0247586\\
1.562	-0.0247587\\
1.564	-0.0247588\\
1.566	-0.0247589\\
1.568	-0.024759\\
1.57	-0.0247591\\
1.572	-0.0247592\\
1.574	-0.0247593\\
1.576	-0.0247594\\
1.578	-0.0247594\\
1.58	-0.0247595\\
1.582	-0.0247596\\
1.584	-0.0247597\\
1.586	-0.0247598\\
1.588	-0.0247599\\
1.59	-0.02476\\
1.592	-0.02476\\
1.594	-0.0247601\\
1.596	-0.0247602\\
1.598	-0.0247603\\
1.6	-0.0247604\\
1.602	-0.0247604\\
1.604	-0.0247605\\
1.606	-0.0247606\\
1.608	-0.0247607\\
1.61	-0.0247608\\
1.612	-0.0247608\\
1.614	-0.0247609\\
1.616	-0.024761\\
1.618	-0.0247611\\
1.62	-0.0247611\\
1.622	-0.0247612\\
1.624	-0.0247613\\
1.626	-0.0247614\\
1.628	-0.0247614\\
1.63	-0.0247615\\
1.632	-0.0247616\\
1.634	-0.0247617\\
1.636	-0.0247617\\
1.638	-0.0247618\\
1.64	-0.0247619\\
1.642	-0.0247619\\
1.644	-0.024762\\
1.646	-0.0247621\\
1.648	-0.0247622\\
1.65	-0.0247622\\
1.652	-0.0247623\\
1.654	-0.0247624\\
1.656	-0.0247624\\
1.658	-0.0247625\\
1.66	-0.0247626\\
1.662	-0.0247626\\
1.664	-0.0247627\\
1.666	-0.0247628\\
1.668	-0.0247628\\
1.67	-0.0247629\\
1.672	-0.024763\\
1.674	-0.024763\\
1.676	-0.0247631\\
1.678	-0.0247631\\
1.68	-0.0247632\\
1.682	-0.0247633\\
1.684	-0.0247633\\
1.686	-0.0247634\\
1.688	-0.0247635\\
1.69	-0.0247635\\
1.692	-0.0247636\\
1.694	-0.0247636\\
1.696	-0.0247637\\
1.698	-0.0247638\\
1.7	-0.0247638\\
1.702	-0.0247639\\
1.704	-0.0247639\\
1.706	-0.024764\\
1.708	-0.0247641\\
1.71	-0.0247641\\
1.712	-0.0247642\\
1.714	-0.0247642\\
1.716	-0.0247643\\
1.718	-0.0247644\\
1.72	-0.0247644\\
1.722	-0.0247645\\
1.724	-0.0247645\\
1.726	-0.0247646\\
1.728	-0.0247646\\
1.73	-0.0247647\\
1.732	-0.0247647\\
1.734	-0.0247648\\
1.736	-0.0247648\\
1.738	-0.0247649\\
1.74	-0.024765\\
1.742	-0.024765\\
1.744	-0.0247651\\
1.746	-0.0247651\\
1.748	-0.0247652\\
1.75	-0.0247652\\
1.752	-0.0247653\\
1.754	-0.0247653\\
1.756	-0.0247654\\
1.758	-0.0247654\\
1.76	-0.0247655\\
1.762	-0.0247655\\
1.764	-0.0247656\\
1.766	-0.0247656\\
1.768	-0.0247657\\
1.77	-0.0247657\\
1.772	-0.0247658\\
1.774	-0.0247658\\
1.776	-0.0247659\\
1.778	-0.0247659\\
1.78	-0.0247659\\
1.782	-0.024766\\
1.784	-0.024766\\
1.786	-0.0247661\\
1.788	-0.0247661\\
1.79	-0.0247662\\
1.792	-0.0247662\\
1.794	-0.0247663\\
1.796	-0.0247663\\
1.798	-0.0247663\\
1.8	-0.0247664\\
1.802	-0.0247664\\
1.804	-0.0247665\\
1.806	-0.0247665\\
1.808	-0.0247666\\
1.81	-0.0247666\\
1.812	-0.0247666\\
1.814	-0.0247667\\
1.816	-0.0247667\\
1.818	-0.0247668\\
1.82	-0.0247668\\
1.822	-0.0247668\\
1.824	-0.0247669\\
1.826	-0.0247669\\
1.828	-0.024767\\
1.83	-0.024767\\
1.832	-0.024767\\
1.834	-0.0247671\\
1.836	-0.0247671\\
1.838	-0.0247672\\
1.84	-0.0247672\\
1.842	-0.0247672\\
1.844	-0.0247673\\
1.846	-0.0247673\\
1.848	-0.0247673\\
1.85	-0.0247674\\
1.852	-0.0247674\\
1.854	-0.0247675\\
1.856	-0.0247675\\
1.858	-0.0247675\\
1.86	-0.0247676\\
1.862	-0.0247676\\
1.864	-0.0247676\\
1.866	-0.0247677\\
1.868	-0.0247677\\
1.87	-0.0247677\\
1.872	-0.0247678\\
1.874	-0.0247678\\
1.876	-0.0247678\\
1.878	-0.0247679\\
1.88	-0.0247679\\
1.882	-0.0247679\\
1.884	-0.024768\\
1.886	-0.024768\\
1.888	-0.024768\\
1.89	-0.0247681\\
1.892	-0.0247681\\
1.894	-0.0247681\\
1.896	-0.0247682\\
1.898	-0.0247682\\
1.9	-0.0247682\\
1.902	-0.0247683\\
1.904	-0.0247683\\
1.906	-0.0247683\\
1.908	-0.0247683\\
1.91	-0.0247684\\
1.912	-0.0247684\\
1.914	-0.0247684\\
1.916	-0.0247685\\
1.918	-0.0247685\\
1.92	-0.0247685\\
1.922	-0.0247686\\
1.924	-0.0247686\\
1.926	-0.0247686\\
1.928	-0.0247686\\
1.93	-0.0247687\\
1.932	-0.0247687\\
1.934	-0.0247687\\
1.936	-0.0247688\\
1.938	-0.0247688\\
1.94	-0.0247688\\
1.942	-0.0247688\\
1.944	-0.0247689\\
1.946	-0.0247689\\
1.948	-0.0247689\\
1.95	-0.024769\\
1.952	-0.024769\\
1.954	-0.024769\\
1.956	-0.024769\\
1.958	-0.0247691\\
1.96	-0.0247691\\
1.962	-0.0247691\\
1.964	-0.0247691\\
1.966	-0.0247692\\
1.968	-0.0247692\\
1.97	-0.0247692\\
1.972	-0.0247692\\
1.974	-0.0247693\\
1.976	-0.0247693\\
1.978	-0.0247693\\
1.98	-0.0247693\\
1.982	-0.0247694\\
1.984	-0.0247694\\
1.986	-0.0247694\\
1.988	-0.0247694\\
1.99	-0.0247695\\
1.992	-0.0247695\\
1.994	-0.0247695\\
1.996	-0.0247695\\
1.998	-0.0247695\\
2	-0.0247696\\
};
\addplot [color=mycolor9, line width=1.0pt, forget plot]
  table[row sep=crcr]{%
-0.024	0\\
-0.022	0\\
-0.02	0\\
-0.018	0\\
-0.016	0\\
-0.014	0\\
-0.012	0\\
-0.01	0\\
-0.008	0\\
-0.006	0\\
-0.004	0\\
-0.002	0\\
0	0\\
0.002	-0.00184652\\
0.004	-0.00342654\\
0.006	-0.00478745\\
0.008	-0.00596865\\
0.01	-0.00700279\\
0.012	-0.00791681\\
0.014	-0.00873287\\
0.016	-0.00946905\\
0.018	-0.0101401\\
0.02	-0.0107579\\
0.022	-0.0113321\\
0.024	-0.0118702\\
0.026	-0.0123782\\
0.028	-0.0128608\\
0.03	-0.0133216\\
0.032	-0.0137632\\
0.034	-0.0141876\\
0.036	-0.0145963\\
0.038	-0.0149903\\
0.04	-0.0153702\\
0.042	-0.0157364\\
0.044	-0.0160893\\
0.046	-0.0164289\\
0.048	-0.0167552\\
0.05	-0.0170685\\
0.052	-0.0173686\\
0.054	-0.0176556\\
0.056	-0.0179296\\
0.058	-0.0181908\\
0.06	-0.0184392\\
0.062	-0.0186752\\
0.064	-0.018899\\
0.066	-0.0191109\\
0.068	-0.0193112\\
0.07	-0.0195004\\
0.072	-0.0196788\\
0.074	-0.0198468\\
0.076	-0.020005\\
0.078	-0.0201538\\
0.08	-0.0202937\\
0.082	-0.020425\\
0.084	-0.0205484\\
0.086	-0.0206641\\
0.088	-0.0207728\\
0.09	-0.0208748\\
0.092	-0.0209706\\
0.094	-0.0210606\\
0.096	-0.0211451\\
0.098	-0.0212245\\
0.1	-0.0212992\\
0.102	-0.0213696\\
0.104	-0.0214359\\
0.106	-0.0214984\\
0.108	-0.0215574\\
0.11	-0.0216131\\
0.112	-0.0216658\\
0.114	-0.0217157\\
0.116	-0.021763\\
0.118	-0.0218078\\
0.12	-0.0218503\\
0.122	-0.0218906\\
0.124	-0.0219288\\
0.126	-0.0219651\\
0.128	-0.0219995\\
0.13	-0.0220321\\
0.132	-0.022063\\
0.134	-0.0220921\\
0.136	-0.0221196\\
0.138	-0.0221455\\
0.14	-0.0221697\\
0.142	-0.0221923\\
0.144	-0.0222134\\
0.146	-0.0222328\\
0.148	-0.0222508\\
0.15	-0.0222671\\
0.152	-0.0222819\\
0.154	-0.0222952\\
0.156	-0.022307\\
0.158	-0.0223173\\
0.16	-0.0223262\\
0.162	-0.0223336\\
0.164	-0.0223397\\
0.166	-0.0223444\\
0.168	-0.0223478\\
0.17	-0.02235\\
0.172	-0.022351\\
0.174	-0.022351\\
0.176	-0.0223499\\
0.178	-0.022348\\
0.18	-0.0223451\\
0.182	-0.0223416\\
0.184	-0.0223374\\
0.186	-0.0223326\\
0.188	-0.0223275\\
0.19	-0.0223219\\
0.192	-0.0223162\\
0.194	-0.0223103\\
0.196	-0.0223044\\
0.198	-0.0222986\\
0.2	-0.0222929\\
0.202	-0.0222876\\
0.204	-0.0222826\\
0.206	-0.0222781\\
0.208	-0.0222741\\
0.21	-0.0222708\\
0.212	-0.0222681\\
0.214	-0.0222663\\
0.216	-0.0222653\\
0.218	-0.0222652\\
0.22	-0.022266\\
0.222	-0.0222679\\
0.224	-0.0222707\\
0.226	-0.0222747\\
0.228	-0.0222797\\
0.23	-0.0222858\\
0.232	-0.0222931\\
0.234	-0.0223014\\
0.236	-0.0223109\\
0.238	-0.0223215\\
0.24	-0.0223332\\
0.242	-0.022346\\
0.244	-0.0223598\\
0.246	-0.0223747\\
0.248	-0.0223905\\
0.25	-0.0224074\\
0.252	-0.0224251\\
0.254	-0.0224438\\
0.256	-0.0224633\\
0.258	-0.0224836\\
0.26	-0.0225046\\
0.262	-0.0225263\\
0.264	-0.0225487\\
0.266	-0.0225718\\
0.268	-0.0225953\\
0.27	-0.0226194\\
0.272	-0.0226439\\
0.274	-0.0226687\\
0.276	-0.022694\\
0.278	-0.0227195\\
0.28	-0.0227453\\
0.282	-0.0227712\\
0.284	-0.0227974\\
0.286	-0.0228236\\
0.288	-0.0228499\\
0.29	-0.0228762\\
0.292	-0.0229024\\
0.294	-0.0229287\\
0.296	-0.0229548\\
0.298	-0.0229808\\
0.3	-0.0230066\\
0.302	-0.0230322\\
0.304	-0.0230576\\
0.306	-0.0230827\\
0.308	-0.0231076\\
0.31	-0.0231322\\
0.312	-0.0231564\\
0.314	-0.0231803\\
0.316	-0.0232039\\
0.318	-0.023227\\
0.32	-0.0232498\\
0.322	-0.0232722\\
0.324	-0.0232941\\
0.326	-0.0233156\\
0.328	-0.0233367\\
0.33	-0.0233574\\
0.332	-0.0233776\\
0.334	-0.0233973\\
0.336	-0.0234166\\
0.338	-0.0234355\\
0.34	-0.0234538\\
0.342	-0.0234718\\
0.344	-0.0234892\\
0.346	-0.0235062\\
0.348	-0.0235228\\
0.35	-0.0235389\\
0.352	-0.0235545\\
0.354	-0.0235697\\
0.356	-0.0235845\\
0.358	-0.0235989\\
0.36	-0.0236128\\
0.362	-0.0236263\\
0.364	-0.0236394\\
0.366	-0.0236521\\
0.368	-0.0236644\\
0.37	-0.0236763\\
0.372	-0.0236879\\
0.374	-0.023699\\
0.376	-0.0237099\\
0.378	-0.0237204\\
0.38	-0.0237306\\
0.382	-0.0237404\\
0.384	-0.02375\\
0.386	-0.0237592\\
0.388	-0.0237682\\
0.39	-0.0237769\\
0.392	-0.0237854\\
0.394	-0.0237936\\
0.396	-0.0238016\\
0.398	-0.0238094\\
0.4	-0.023817\\
0.402	-0.0238244\\
0.404	-0.0238316\\
0.406	-0.0238387\\
0.408	-0.0238456\\
0.41	-0.0238524\\
0.412	-0.023859\\
0.414	-0.0238656\\
0.416	-0.023872\\
0.418	-0.0238784\\
0.42	-0.0238847\\
0.422	-0.0238909\\
0.424	-0.0238971\\
0.426	-0.0239032\\
0.428	-0.0239093\\
0.43	-0.0239154\\
0.432	-0.0239215\\
0.434	-0.0239275\\
0.436	-0.0239336\\
0.438	-0.0239396\\
0.44	-0.0239457\\
0.442	-0.0239518\\
0.444	-0.023958\\
0.446	-0.0239641\\
0.448	-0.0239704\\
0.45	-0.0239766\\
0.452	-0.0239829\\
0.454	-0.0239892\\
0.456	-0.0239956\\
0.458	-0.024002\\
0.46	-0.0240085\\
0.462	-0.0240151\\
0.464	-0.0240216\\
0.466	-0.0240283\\
0.468	-0.0240349\\
0.47	-0.0240417\\
0.472	-0.0240484\\
0.474	-0.0240552\\
0.476	-0.0240621\\
0.478	-0.024069\\
0.48	-0.0240759\\
0.482	-0.0240828\\
0.484	-0.0240897\\
0.486	-0.0240967\\
0.488	-0.0241036\\
0.49	-0.0241106\\
0.492	-0.0241176\\
0.494	-0.0241245\\
0.496	-0.0241315\\
0.498	-0.0241384\\
0.5	-0.0241453\\
0.502	-0.0241521\\
0.504	-0.0241589\\
0.506	-0.0241657\\
0.508	-0.0241724\\
0.51	-0.0241791\\
0.512	-0.0241857\\
0.514	-0.0241922\\
0.516	-0.0241986\\
0.518	-0.024205\\
0.52	-0.0242113\\
0.522	-0.0242175\\
0.524	-0.0242236\\
0.526	-0.0242296\\
0.528	-0.0242355\\
0.53	-0.0242413\\
0.532	-0.0242469\\
0.534	-0.0242525\\
0.536	-0.024258\\
0.538	-0.0242633\\
0.54	-0.0242685\\
0.542	-0.0242736\\
0.544	-0.0242786\\
0.546	-0.0242834\\
0.548	-0.0242882\\
0.55	-0.0242928\\
0.552	-0.0242972\\
0.554	-0.0243016\\
0.556	-0.0243058\\
0.558	-0.02431\\
0.56	-0.024314\\
0.562	-0.0243178\\
0.564	-0.0243216\\
0.566	-0.0243253\\
0.568	-0.0243288\\
0.57	-0.0243323\\
0.572	-0.0243356\\
0.574	-0.0243389\\
0.576	-0.024342\\
0.578	-0.0243451\\
0.58	-0.024348\\
0.582	-0.0243509\\
0.584	-0.0243537\\
0.586	-0.0243565\\
0.588	-0.0243591\\
0.59	-0.0243617\\
0.592	-0.0243643\\
0.594	-0.0243668\\
0.596	-0.0243692\\
0.598	-0.0243716\\
0.6	-0.0243739\\
0.602	-0.0243762\\
0.604	-0.0243785\\
0.606	-0.0243808\\
0.608	-0.024383\\
0.61	-0.0243852\\
0.612	-0.0243874\\
0.614	-0.0243895\\
0.616	-0.0243917\\
0.618	-0.0243939\\
0.62	-0.024396\\
0.622	-0.0243982\\
0.624	-0.0244003\\
0.626	-0.0244025\\
0.628	-0.0244046\\
0.63	-0.0244068\\
0.632	-0.024409\\
0.634	-0.0244112\\
0.636	-0.0244134\\
0.638	-0.0244157\\
0.64	-0.0244179\\
0.642	-0.0244202\\
0.644	-0.0244225\\
0.646	-0.0244248\\
0.648	-0.0244272\\
0.65	-0.0244295\\
0.652	-0.0244319\\
0.654	-0.0244343\\
0.656	-0.0244367\\
0.658	-0.0244392\\
0.66	-0.0244417\\
0.662	-0.0244441\\
0.664	-0.0244466\\
0.666	-0.0244491\\
0.668	-0.0244517\\
0.67	-0.0244542\\
0.672	-0.0244568\\
0.674	-0.0244593\\
0.676	-0.0244619\\
0.678	-0.0244645\\
0.68	-0.0244671\\
0.682	-0.0244697\\
0.684	-0.0244723\\
0.686	-0.0244749\\
0.688	-0.0244775\\
0.69	-0.0244801\\
0.692	-0.0244826\\
0.694	-0.0244852\\
0.696	-0.0244878\\
0.698	-0.0244903\\
0.7	-0.0244929\\
0.702	-0.0244954\\
0.704	-0.0244979\\
0.706	-0.0245004\\
0.708	-0.0245029\\
0.71	-0.0245054\\
0.712	-0.0245078\\
0.714	-0.0245102\\
0.716	-0.0245126\\
0.718	-0.024515\\
0.72	-0.0245173\\
0.722	-0.0245196\\
0.724	-0.0245219\\
0.726	-0.0245242\\
0.728	-0.0245264\\
0.73	-0.0245286\\
0.732	-0.0245307\\
0.734	-0.0245329\\
0.736	-0.024535\\
0.738	-0.024537\\
0.74	-0.0245391\\
0.742	-0.0245411\\
0.744	-0.0245431\\
0.746	-0.024545\\
0.748	-0.0245469\\
0.75	-0.0245488\\
0.752	-0.0245506\\
0.754	-0.0245525\\
0.756	-0.0245542\\
0.758	-0.024556\\
0.76	-0.0245577\\
0.762	-0.0245594\\
0.764	-0.0245611\\
0.766	-0.0245627\\
0.768	-0.0245643\\
0.77	-0.0245659\\
0.772	-0.0245675\\
0.774	-0.024569\\
0.776	-0.0245705\\
0.778	-0.024572\\
0.78	-0.0245734\\
0.782	-0.0245749\\
0.784	-0.0245763\\
0.786	-0.0245777\\
0.788	-0.024579\\
0.79	-0.0245804\\
0.792	-0.0245817\\
0.794	-0.024583\\
0.796	-0.0245843\\
0.798	-0.0245856\\
0.8	-0.0245869\\
0.802	-0.0245881\\
0.804	-0.0245893\\
0.806	-0.0245905\\
0.808	-0.0245917\\
0.81	-0.0245929\\
0.812	-0.0245941\\
0.814	-0.0245952\\
0.816	-0.0245964\\
0.818	-0.0245975\\
0.82	-0.0245986\\
0.822	-0.0245997\\
0.824	-0.0246008\\
0.826	-0.0246019\\
0.828	-0.024603\\
0.83	-0.0246041\\
0.832	-0.0246052\\
0.834	-0.0246062\\
0.836	-0.0246073\\
0.838	-0.0246083\\
0.84	-0.0246093\\
0.842	-0.0246104\\
0.844	-0.0246114\\
0.846	-0.0246124\\
0.848	-0.0246134\\
0.85	-0.0246144\\
0.852	-0.0246154\\
0.854	-0.0246164\\
0.856	-0.0246174\\
0.858	-0.0246184\\
0.86	-0.0246193\\
0.862	-0.0246203\\
0.864	-0.0246213\\
0.866	-0.0246222\\
0.868	-0.0246232\\
0.87	-0.0246242\\
0.872	-0.0246251\\
0.874	-0.024626\\
0.876	-0.024627\\
0.878	-0.0246279\\
0.88	-0.0246289\\
0.882	-0.0246298\\
0.884	-0.0246307\\
0.886	-0.0246316\\
0.888	-0.0246326\\
0.89	-0.0246335\\
0.892	-0.0246344\\
0.894	-0.0246353\\
0.896	-0.0246362\\
0.898	-0.0246371\\
0.9	-0.024638\\
0.902	-0.0246389\\
0.904	-0.0246398\\
0.906	-0.0246407\\
0.908	-0.0246416\\
0.91	-0.0246425\\
0.912	-0.0246433\\
0.914	-0.0246442\\
0.916	-0.0246451\\
0.918	-0.024646\\
0.92	-0.0246468\\
0.922	-0.0246477\\
0.924	-0.0246486\\
0.926	-0.0246494\\
0.928	-0.0246503\\
0.93	-0.0246511\\
0.932	-0.024652\\
0.934	-0.0246528\\
0.936	-0.0246536\\
0.938	-0.0246545\\
0.94	-0.0246553\\
0.942	-0.0246561\\
0.944	-0.0246569\\
0.946	-0.0246578\\
0.948	-0.0246586\\
0.95	-0.0246594\\
0.952	-0.0246602\\
0.954	-0.024661\\
0.956	-0.0246618\\
0.958	-0.0246626\\
0.96	-0.0246634\\
0.962	-0.0246641\\
0.964	-0.0246649\\
0.966	-0.0246657\\
0.968	-0.0246664\\
0.97	-0.0246672\\
0.972	-0.024668\\
0.974	-0.0246687\\
0.976	-0.0246695\\
0.978	-0.0246702\\
0.98	-0.0246709\\
0.982	-0.0246717\\
0.984	-0.0246724\\
0.986	-0.0246731\\
0.988	-0.0246738\\
0.99	-0.0246745\\
0.992	-0.0246752\\
0.994	-0.0246759\\
0.996	-0.0246766\\
0.998	-0.0246773\\
1	-0.024678\\
1.002	-0.0246786\\
1.004	-0.0246793\\
1.006	-0.02468\\
1.008	-0.0246806\\
1.01	-0.0246813\\
1.012	-0.0246819\\
1.014	-0.0246825\\
1.016	-0.0246832\\
1.018	-0.0246838\\
1.02	-0.0246844\\
1.022	-0.024685\\
1.024	-0.0246856\\
1.026	-0.0246862\\
1.028	-0.0246868\\
1.03	-0.0246874\\
1.032	-0.024688\\
1.034	-0.0246886\\
1.036	-0.0246892\\
1.038	-0.0246897\\
1.04	-0.0246903\\
1.042	-0.0246909\\
1.044	-0.0246914\\
1.046	-0.024692\\
1.048	-0.0246925\\
1.05	-0.0246931\\
1.052	-0.0246936\\
1.054	-0.0246941\\
1.056	-0.0246946\\
1.058	-0.0246952\\
1.06	-0.0246957\\
1.062	-0.0246962\\
1.064	-0.0246967\\
1.066	-0.0246972\\
1.068	-0.0246977\\
1.07	-0.0246982\\
1.072	-0.0246987\\
1.074	-0.0246992\\
1.076	-0.0246997\\
1.078	-0.0247001\\
1.08	-0.0247006\\
1.082	-0.0247011\\
1.084	-0.0247015\\
1.086	-0.024702\\
1.088	-0.0247025\\
1.09	-0.0247029\\
1.092	-0.0247034\\
1.094	-0.0247038\\
1.096	-0.0247043\\
1.098	-0.0247047\\
1.1	-0.0247051\\
1.102	-0.0247056\\
1.104	-0.024706\\
1.106	-0.0247064\\
1.108	-0.0247069\\
1.11	-0.0247073\\
1.112	-0.0247077\\
1.114	-0.0247081\\
1.116	-0.0247085\\
1.118	-0.0247089\\
1.12	-0.0247093\\
1.122	-0.0247098\\
1.124	-0.0247102\\
1.126	-0.0247106\\
1.128	-0.0247109\\
1.13	-0.0247113\\
1.132	-0.0247117\\
1.134	-0.0247121\\
1.136	-0.0247125\\
1.138	-0.0247129\\
1.14	-0.0247133\\
1.142	-0.0247136\\
1.144	-0.024714\\
1.146	-0.0247144\\
1.148	-0.0247148\\
1.15	-0.0247151\\
1.152	-0.0247155\\
1.154	-0.0247159\\
1.156	-0.0247162\\
1.158	-0.0247166\\
1.16	-0.024717\\
1.162	-0.0247173\\
1.164	-0.0247177\\
1.166	-0.024718\\
1.168	-0.0247184\\
1.17	-0.0247187\\
1.172	-0.0247191\\
1.174	-0.0247194\\
1.176	-0.0247198\\
1.178	-0.0247201\\
1.18	-0.0247204\\
1.182	-0.0247208\\
1.184	-0.0247211\\
1.186	-0.0247214\\
1.188	-0.0247218\\
1.19	-0.0247221\\
1.192	-0.0247224\\
1.194	-0.0247228\\
1.196	-0.0247231\\
1.198	-0.0247234\\
1.2	-0.0247237\\
1.202	-0.024724\\
1.204	-0.0247244\\
1.206	-0.0247247\\
1.208	-0.024725\\
1.21	-0.0247253\\
1.212	-0.0247256\\
1.214	-0.0247259\\
1.216	-0.0247262\\
1.218	-0.0247265\\
1.22	-0.0247268\\
1.222	-0.0247271\\
1.224	-0.0247274\\
1.226	-0.0247277\\
1.228	-0.024728\\
1.23	-0.0247283\\
1.232	-0.0247286\\
1.234	-0.0247289\\
1.236	-0.0247292\\
1.238	-0.0247295\\
1.24	-0.0247298\\
1.242	-0.0247301\\
1.244	-0.0247303\\
1.246	-0.0247306\\
1.248	-0.0247309\\
1.25	-0.0247312\\
1.252	-0.0247315\\
1.254	-0.0247317\\
1.256	-0.024732\\
1.258	-0.0247323\\
1.26	-0.0247326\\
1.262	-0.0247328\\
1.264	-0.0247331\\
1.266	-0.0247333\\
1.268	-0.0247336\\
1.27	-0.0247339\\
1.272	-0.0247341\\
1.274	-0.0247344\\
1.276	-0.0247346\\
1.278	-0.0247349\\
1.28	-0.0247351\\
1.282	-0.0247354\\
1.284	-0.0247356\\
1.286	-0.0247359\\
1.288	-0.0247361\\
1.29	-0.0247364\\
1.292	-0.0247366\\
1.294	-0.0247369\\
1.296	-0.0247371\\
1.298	-0.0247373\\
1.3	-0.0247376\\
1.302	-0.0247378\\
1.304	-0.024738\\
1.306	-0.0247382\\
1.308	-0.0247385\\
1.31	-0.0247387\\
1.312	-0.0247389\\
1.314	-0.0247391\\
1.316	-0.0247394\\
1.318	-0.0247396\\
1.32	-0.0247398\\
1.322	-0.02474\\
1.324	-0.0247402\\
1.326	-0.0247404\\
1.328	-0.0247406\\
1.33	-0.0247408\\
1.332	-0.024741\\
1.334	-0.0247412\\
1.336	-0.0247415\\
1.338	-0.0247417\\
1.34	-0.0247418\\
1.342	-0.024742\\
1.344	-0.0247422\\
1.346	-0.0247424\\
1.348	-0.0247426\\
1.35	-0.0247428\\
1.352	-0.024743\\
1.354	-0.0247432\\
1.356	-0.0247434\\
1.358	-0.0247436\\
1.36	-0.0247438\\
1.362	-0.0247439\\
1.364	-0.0247441\\
1.366	-0.0247443\\
1.368	-0.0247445\\
1.37	-0.0247446\\
1.372	-0.0247448\\
1.374	-0.024745\\
1.376	-0.0247452\\
1.378	-0.0247453\\
1.38	-0.0247455\\
1.382	-0.0247457\\
1.384	-0.0247459\\
1.386	-0.024746\\
1.388	-0.0247462\\
1.39	-0.0247464\\
1.392	-0.0247465\\
1.394	-0.0247467\\
1.396	-0.0247468\\
1.398	-0.024747\\
1.4	-0.0247472\\
1.402	-0.0247473\\
1.404	-0.0247475\\
1.406	-0.0247476\\
1.408	-0.0247478\\
1.41	-0.024748\\
1.412	-0.0247481\\
1.414	-0.0247483\\
1.416	-0.0247484\\
1.418	-0.0247486\\
1.42	-0.0247487\\
1.422	-0.0247489\\
1.424	-0.024749\\
1.426	-0.0247492\\
1.428	-0.0247493\\
1.43	-0.0247495\\
1.432	-0.0247496\\
1.434	-0.0247498\\
1.436	-0.0247499\\
1.438	-0.0247501\\
1.44	-0.0247502\\
1.442	-0.0247503\\
1.444	-0.0247505\\
1.446	-0.0247506\\
1.448	-0.0247508\\
1.45	-0.0247509\\
1.452	-0.024751\\
1.454	-0.0247512\\
1.456	-0.0247513\\
1.458	-0.0247515\\
1.46	-0.0247516\\
1.462	-0.0247517\\
1.464	-0.0247519\\
1.466	-0.024752\\
1.468	-0.0247521\\
1.47	-0.0247523\\
1.472	-0.0247524\\
1.474	-0.0247525\\
1.476	-0.0247527\\
1.478	-0.0247528\\
1.48	-0.0247529\\
1.482	-0.0247531\\
1.484	-0.0247532\\
1.486	-0.0247533\\
1.488	-0.0247534\\
1.49	-0.0247536\\
1.492	-0.0247537\\
1.494	-0.0247538\\
1.496	-0.0247539\\
1.498	-0.0247541\\
1.5	-0.0247542\\
1.502	-0.0247543\\
1.504	-0.0247544\\
1.506	-0.0247545\\
1.508	-0.0247547\\
1.51	-0.0247548\\
1.512	-0.0247549\\
1.514	-0.024755\\
1.516	-0.0247551\\
1.518	-0.0247553\\
1.52	-0.0247554\\
1.522	-0.0247555\\
1.524	-0.0247556\\
1.526	-0.0247557\\
1.528	-0.0247558\\
1.53	-0.0247559\\
1.532	-0.024756\\
1.534	-0.0247561\\
1.536	-0.0247563\\
1.538	-0.0247564\\
1.54	-0.0247565\\
1.542	-0.0247566\\
1.544	-0.0247567\\
1.546	-0.0247568\\
1.548	-0.0247569\\
1.55	-0.024757\\
1.552	-0.0247571\\
1.554	-0.0247572\\
1.556	-0.0247573\\
1.558	-0.0247574\\
1.56	-0.0247575\\
1.562	-0.0247576\\
1.564	-0.0247577\\
1.566	-0.0247578\\
1.568	-0.0247579\\
1.57	-0.024758\\
1.572	-0.0247581\\
1.574	-0.0247582\\
1.576	-0.0247583\\
1.578	-0.0247584\\
1.58	-0.0247584\\
1.582	-0.0247585\\
1.584	-0.0247586\\
1.586	-0.0247587\\
1.588	-0.0247588\\
1.59	-0.0247589\\
1.592	-0.024759\\
1.594	-0.0247591\\
1.596	-0.0247592\\
1.598	-0.0247592\\
1.6	-0.0247593\\
1.602	-0.0247594\\
1.604	-0.0247595\\
1.606	-0.0247596\\
1.608	-0.0247597\\
1.61	-0.0247598\\
1.612	-0.0247598\\
1.614	-0.0247599\\
1.616	-0.02476\\
1.618	-0.0247601\\
1.62	-0.0247602\\
1.622	-0.0247602\\
1.624	-0.0247603\\
1.626	-0.0247604\\
1.628	-0.0247605\\
1.63	-0.0247606\\
1.632	-0.0247606\\
1.634	-0.0247607\\
1.636	-0.0247608\\
1.638	-0.0247609\\
1.64	-0.0247609\\
1.642	-0.024761\\
1.644	-0.0247611\\
1.646	-0.0247612\\
1.648	-0.0247612\\
1.65	-0.0247613\\
1.652	-0.0247614\\
1.654	-0.0247615\\
1.656	-0.0247615\\
1.658	-0.0247616\\
1.66	-0.0247617\\
1.662	-0.0247617\\
1.664	-0.0247618\\
1.666	-0.0247619\\
1.668	-0.024762\\
1.67	-0.024762\\
1.672	-0.0247621\\
1.674	-0.0247622\\
1.676	-0.0247622\\
1.678	-0.0247623\\
1.68	-0.0247624\\
1.682	-0.0247624\\
1.684	-0.0247625\\
1.686	-0.0247626\\
1.688	-0.0247626\\
1.69	-0.0247627\\
1.692	-0.0247628\\
1.694	-0.0247628\\
1.696	-0.0247629\\
1.698	-0.024763\\
1.7	-0.024763\\
1.702	-0.0247631\\
1.704	-0.0247631\\
1.706	-0.0247632\\
1.708	-0.0247633\\
1.71	-0.0247633\\
1.712	-0.0247634\\
1.714	-0.0247635\\
1.716	-0.0247635\\
1.718	-0.0247636\\
1.72	-0.0247636\\
1.722	-0.0247637\\
1.724	-0.0247638\\
1.726	-0.0247638\\
1.728	-0.0247639\\
1.73	-0.0247639\\
1.732	-0.024764\\
1.734	-0.024764\\
1.736	-0.0247641\\
1.738	-0.0247642\\
1.74	-0.0247642\\
1.742	-0.0247643\\
1.744	-0.0247643\\
1.746	-0.0247644\\
1.748	-0.0247644\\
1.75	-0.0247645\\
1.752	-0.0247646\\
1.754	-0.0247646\\
1.756	-0.0247647\\
1.758	-0.0247647\\
1.76	-0.0247648\\
1.762	-0.0247648\\
1.764	-0.0247649\\
1.766	-0.0247649\\
1.768	-0.024765\\
1.77	-0.024765\\
1.772	-0.0247651\\
1.774	-0.0247651\\
1.776	-0.0247652\\
1.778	-0.0247652\\
1.78	-0.0247653\\
1.782	-0.0247653\\
1.784	-0.0247654\\
1.786	-0.0247654\\
1.788	-0.0247655\\
1.79	-0.0247655\\
1.792	-0.0247656\\
1.794	-0.0247656\\
1.796	-0.0247657\\
1.798	-0.0247657\\
1.8	-0.0247658\\
1.802	-0.0247658\\
1.804	-0.0247658\\
1.806	-0.0247659\\
1.808	-0.0247659\\
1.81	-0.024766\\
1.812	-0.024766\\
1.814	-0.0247661\\
1.816	-0.0247661\\
1.818	-0.0247662\\
1.82	-0.0247662\\
1.822	-0.0247662\\
1.824	-0.0247663\\
1.826	-0.0247663\\
1.828	-0.0247664\\
1.83	-0.0247664\\
1.832	-0.0247665\\
1.834	-0.0247665\\
1.836	-0.0247665\\
1.838	-0.0247666\\
1.84	-0.0247666\\
1.842	-0.0247667\\
1.844	-0.0247667\\
1.846	-0.0247667\\
1.848	-0.0247668\\
1.85	-0.0247668\\
1.852	-0.0247669\\
1.854	-0.0247669\\
1.856	-0.0247669\\
1.858	-0.024767\\
1.86	-0.024767\\
1.862	-0.024767\\
1.864	-0.0247671\\
1.866	-0.0247671\\
1.868	-0.0247672\\
1.87	-0.0247672\\
1.872	-0.0247672\\
1.874	-0.0247673\\
1.876	-0.0247673\\
1.878	-0.0247673\\
1.88	-0.0247674\\
1.882	-0.0247674\\
1.884	-0.0247675\\
1.886	-0.0247675\\
1.888	-0.0247675\\
1.89	-0.0247676\\
1.892	-0.0247676\\
1.894	-0.0247676\\
1.896	-0.0247677\\
1.898	-0.0247677\\
1.9	-0.0247677\\
1.902	-0.0247678\\
1.904	-0.0247678\\
1.906	-0.0247678\\
1.908	-0.0247679\\
1.91	-0.0247679\\
1.912	-0.0247679\\
1.914	-0.024768\\
1.916	-0.024768\\
1.918	-0.024768\\
1.92	-0.0247681\\
1.922	-0.0247681\\
1.924	-0.0247681\\
1.926	-0.0247682\\
1.928	-0.0247682\\
1.93	-0.0247682\\
1.932	-0.0247682\\
1.934	-0.0247683\\
1.936	-0.0247683\\
1.938	-0.0247683\\
1.94	-0.0247684\\
1.942	-0.0247684\\
1.944	-0.0247684\\
1.946	-0.0247685\\
1.948	-0.0247685\\
1.95	-0.0247685\\
1.952	-0.0247685\\
1.954	-0.0247686\\
1.956	-0.0247686\\
1.958	-0.0247686\\
1.96	-0.0247687\\
1.962	-0.0247687\\
1.964	-0.0247687\\
1.966	-0.0247687\\
1.968	-0.0247688\\
1.97	-0.0247688\\
1.972	-0.0247688\\
1.974	-0.0247689\\
1.976	-0.0247689\\
1.978	-0.0247689\\
1.98	-0.0247689\\
1.982	-0.024769\\
1.984	-0.024769\\
1.986	-0.024769\\
1.988	-0.024769\\
1.99	-0.0247691\\
1.992	-0.0247691\\
1.994	-0.0247691\\
1.996	-0.0247691\\
1.998	-0.0247692\\
2	-0.0247692\\
};
\addplot [color=mycolor10, line width=1.0pt, forget plot]
  table[row sep=crcr]{%
-0.024	0\\
-0.022	0\\
-0.02	0\\
-0.018	0\\
-0.016	0\\
-0.014	0\\
-0.012	0\\
-0.01	0\\
-0.008	0\\
-0.006	0\\
-0.004	0\\
-0.002	0\\
0	0\\
0.002	-0.00184664\\
0.004	-0.00342739\\
0.006	-0.00479004\\
0.008	-0.00597422\\
0.01	-0.00701266\\
0.012	-0.00793228\\
0.014	-0.00875511\\
0.016	-0.00949912\\
0.018	-0.0101788\\
0.02	-0.010806\\
0.022	-0.0113899\\
0.024	-0.0119379\\
0.026	-0.0124559\\
0.028	-0.0129482\\
0.03	-0.0134184\\
0.032	-0.0138689\\
0.034	-0.0143016\\
0.036	-0.0147179\\
0.038	-0.0151185\\
0.04	-0.0155043\\
0.042	-0.0158755\\
0.044	-0.0162324\\
0.046	-0.0165751\\
0.048	-0.0169037\\
0.05	-0.0172183\\
0.052	-0.017519\\
0.054	-0.0178058\\
0.056	-0.0180789\\
0.058	-0.0183386\\
0.06	-0.018585\\
0.062	-0.0188184\\
0.064	-0.0190392\\
0.066	-0.0192477\\
0.068	-0.0194444\\
0.07	-0.0196297\\
0.072	-0.019804\\
0.074	-0.0199679\\
0.076	-0.0201218\\
0.078	-0.0202662\\
0.08	-0.0204017\\
0.082	-0.0205288\\
0.084	-0.0206479\\
0.086	-0.0207595\\
0.088	-0.0208641\\
0.09	-0.0209621\\
0.092	-0.0210541\\
0.094	-0.0211403\\
0.096	-0.0212212\\
0.098	-0.0212972\\
0.1	-0.0213686\\
0.102	-0.0214358\\
0.104	-0.0214991\\
0.106	-0.0215587\\
0.108	-0.0216149\\
0.11	-0.021668\\
0.112	-0.0217181\\
0.114	-0.0217656\\
0.116	-0.0218105\\
0.118	-0.021853\\
0.12	-0.0218933\\
0.122	-0.0219316\\
0.124	-0.0219678\\
0.126	-0.0220021\\
0.128	-0.0220346\\
0.13	-0.0220654\\
0.132	-0.0220945\\
0.134	-0.0221219\\
0.136	-0.0221477\\
0.138	-0.0221719\\
0.14	-0.0221945\\
0.142	-0.0222156\\
0.144	-0.0222351\\
0.146	-0.0222532\\
0.148	-0.0222697\\
0.15	-0.0222847\\
0.152	-0.0222982\\
0.154	-0.0223103\\
0.156	-0.0223209\\
0.158	-0.0223301\\
0.16	-0.022338\\
0.162	-0.0223444\\
0.164	-0.0223496\\
0.166	-0.0223534\\
0.168	-0.0223561\\
0.17	-0.0223576\\
0.172	-0.0223581\\
0.174	-0.0223575\\
0.176	-0.022356\\
0.178	-0.0223537\\
0.18	-0.0223505\\
0.182	-0.0223467\\
0.184	-0.0223424\\
0.186	-0.0223375\\
0.188	-0.0223323\\
0.19	-0.0223268\\
0.192	-0.0223211\\
0.194	-0.0223153\\
0.196	-0.0223096\\
0.198	-0.022304\\
0.2	-0.0222985\\
0.202	-0.0222934\\
0.204	-0.0222887\\
0.206	-0.0222845\\
0.208	-0.0222808\\
0.21	-0.0222778\\
0.212	-0.0222755\\
0.214	-0.0222739\\
0.216	-0.0222732\\
0.218	-0.0222734\\
0.22	-0.0222745\\
0.222	-0.0222766\\
0.224	-0.0222796\\
0.226	-0.0222837\\
0.228	-0.0222889\\
0.23	-0.0222951\\
0.232	-0.0223024\\
0.234	-0.0223108\\
0.236	-0.0223202\\
0.238	-0.0223307\\
0.24	-0.0223423\\
0.242	-0.0223549\\
0.244	-0.0223685\\
0.246	-0.0223831\\
0.248	-0.0223987\\
0.25	-0.0224152\\
0.252	-0.0224325\\
0.254	-0.0224507\\
0.256	-0.0224697\\
0.258	-0.0224895\\
0.26	-0.02251\\
0.262	-0.0225311\\
0.264	-0.0225529\\
0.266	-0.0225752\\
0.268	-0.022598\\
0.27	-0.0226213\\
0.272	-0.0226451\\
0.274	-0.0226692\\
0.276	-0.0226936\\
0.278	-0.0227183\\
0.28	-0.0227432\\
0.282	-0.0227683\\
0.284	-0.0227935\\
0.286	-0.0228188\\
0.288	-0.0228442\\
0.29	-0.0228696\\
0.292	-0.022895\\
0.294	-0.0229203\\
0.296	-0.0229455\\
0.298	-0.0229706\\
0.3	-0.0229955\\
0.302	-0.0230202\\
0.304	-0.0230447\\
0.306	-0.023069\\
0.308	-0.023093\\
0.31	-0.0231167\\
0.312	-0.0231401\\
0.314	-0.0231632\\
0.316	-0.023186\\
0.318	-0.0232083\\
0.32	-0.0232303\\
0.322	-0.023252\\
0.324	-0.0232732\\
0.326	-0.0232941\\
0.328	-0.0233145\\
0.33	-0.0233345\\
0.332	-0.0233541\\
0.334	-0.0233732\\
0.336	-0.023392\\
0.338	-0.0234103\\
0.34	-0.0234281\\
0.342	-0.0234456\\
0.344	-0.0234626\\
0.346	-0.0234792\\
0.348	-0.0234953\\
0.35	-0.023511\\
0.352	-0.0235263\\
0.354	-0.0235412\\
0.356	-0.0235557\\
0.358	-0.0235698\\
0.36	-0.0235835\\
0.362	-0.0235968\\
0.364	-0.0236097\\
0.366	-0.0236222\\
0.368	-0.0236344\\
0.37	-0.0236462\\
0.372	-0.0236576\\
0.374	-0.0236688\\
0.376	-0.0236796\\
0.378	-0.02369\\
0.38	-0.0237002\\
0.382	-0.0237101\\
0.384	-0.0237197\\
0.386	-0.023729\\
0.388	-0.0237381\\
0.39	-0.0237469\\
0.392	-0.0237555\\
0.394	-0.0237638\\
0.396	-0.023772\\
0.398	-0.0237799\\
0.4	-0.0237877\\
0.402	-0.0237953\\
0.404	-0.0238027\\
0.406	-0.02381\\
0.408	-0.0238171\\
0.41	-0.0238241\\
0.412	-0.023831\\
0.414	-0.0238378\\
0.416	-0.0238445\\
0.418	-0.0238511\\
0.42	-0.0238577\\
0.422	-0.0238641\\
0.424	-0.0238706\\
0.426	-0.023877\\
0.428	-0.0238834\\
0.43	-0.0238897\\
0.432	-0.023896\\
0.434	-0.0239024\\
0.436	-0.0239087\\
0.438	-0.023915\\
0.44	-0.0239214\\
0.442	-0.0239277\\
0.444	-0.0239341\\
0.446	-0.0239405\\
0.448	-0.023947\\
0.45	-0.0239535\\
0.452	-0.02396\\
0.454	-0.0239666\\
0.456	-0.0239732\\
0.458	-0.0239798\\
0.46	-0.0239865\\
0.462	-0.0239932\\
0.464	-0.024\\
0.466	-0.0240068\\
0.468	-0.0240137\\
0.47	-0.0240205\\
0.472	-0.0240275\\
0.474	-0.0240344\\
0.476	-0.0240414\\
0.478	-0.0240484\\
0.48	-0.0240554\\
0.482	-0.0240625\\
0.484	-0.0240695\\
0.486	-0.0240766\\
0.488	-0.0240836\\
0.49	-0.0240907\\
0.492	-0.0240977\\
0.494	-0.0241047\\
0.496	-0.0241117\\
0.498	-0.0241187\\
0.5	-0.0241257\\
0.502	-0.0241326\\
0.504	-0.0241394\\
0.506	-0.0241462\\
0.508	-0.024153\\
0.51	-0.0241597\\
0.512	-0.0241663\\
0.514	-0.0241728\\
0.516	-0.0241793\\
0.518	-0.0241857\\
0.52	-0.024192\\
0.522	-0.0241982\\
0.524	-0.0242043\\
0.526	-0.0242104\\
0.528	-0.0242163\\
0.53	-0.0242221\\
0.532	-0.0242278\\
0.534	-0.0242334\\
0.536	-0.0242388\\
0.538	-0.0242442\\
0.54	-0.0242494\\
0.542	-0.0242546\\
0.544	-0.0242596\\
0.546	-0.0242644\\
0.548	-0.0242692\\
0.55	-0.0242738\\
0.552	-0.0242784\\
0.554	-0.0242828\\
0.556	-0.0242871\\
0.558	-0.0242912\\
0.56	-0.0242953\\
0.562	-0.0242992\\
0.564	-0.0243031\\
0.566	-0.0243068\\
0.568	-0.0243104\\
0.57	-0.0243139\\
0.572	-0.0243174\\
0.574	-0.0243207\\
0.576	-0.0243239\\
0.578	-0.0243271\\
0.58	-0.0243301\\
0.582	-0.0243331\\
0.584	-0.024336\\
0.586	-0.0243389\\
0.588	-0.0243417\\
0.59	-0.0243444\\
0.592	-0.024347\\
0.594	-0.0243496\\
0.596	-0.0243522\\
0.598	-0.0243547\\
0.6	-0.0243572\\
0.602	-0.0243596\\
0.604	-0.024362\\
0.606	-0.0243644\\
0.608	-0.0243668\\
0.61	-0.0243691\\
0.612	-0.0243714\\
0.614	-0.0243737\\
0.616	-0.024376\\
0.618	-0.0243783\\
0.62	-0.0243806\\
0.622	-0.0243829\\
0.624	-0.0243852\\
0.626	-0.0243875\\
0.628	-0.0243898\\
0.63	-0.0243921\\
0.632	-0.0243944\\
0.634	-0.0243968\\
0.636	-0.0243991\\
0.638	-0.0244015\\
0.64	-0.0244039\\
0.642	-0.0244063\\
0.644	-0.0244087\\
0.646	-0.0244112\\
0.648	-0.0244136\\
0.65	-0.0244161\\
0.652	-0.0244186\\
0.654	-0.0244211\\
0.656	-0.0244237\\
0.658	-0.0244262\\
0.66	-0.0244288\\
0.662	-0.0244314\\
0.664	-0.024434\\
0.666	-0.0244366\\
0.668	-0.0244392\\
0.67	-0.0244419\\
0.672	-0.0244445\\
0.674	-0.0244472\\
0.676	-0.0244498\\
0.678	-0.0244525\\
0.68	-0.0244552\\
0.682	-0.0244578\\
0.684	-0.0244605\\
0.686	-0.0244632\\
0.688	-0.0244658\\
0.69	-0.0244685\\
0.692	-0.0244712\\
0.694	-0.0244738\\
0.696	-0.0244764\\
0.698	-0.0244791\\
0.7	-0.0244817\\
0.702	-0.0244842\\
0.704	-0.0244868\\
0.706	-0.0244894\\
0.708	-0.0244919\\
0.71	-0.0244944\\
0.712	-0.0244969\\
0.714	-0.0244994\\
0.716	-0.0245018\\
0.718	-0.0245042\\
0.72	-0.0245066\\
0.722	-0.024509\\
0.724	-0.0245113\\
0.726	-0.0245136\\
0.728	-0.0245159\\
0.73	-0.0245181\\
0.732	-0.0245203\\
0.734	-0.0245225\\
0.736	-0.0245246\\
0.738	-0.0245267\\
0.74	-0.0245288\\
0.742	-0.0245309\\
0.744	-0.0245329\\
0.746	-0.0245349\\
0.748	-0.0245368\\
0.75	-0.0245388\\
0.752	-0.0245407\\
0.754	-0.0245425\\
0.756	-0.0245443\\
0.758	-0.0245462\\
0.76	-0.0245479\\
0.762	-0.0245497\\
0.764	-0.0245514\\
0.766	-0.0245531\\
0.768	-0.0245547\\
0.77	-0.0245564\\
0.772	-0.024558\\
0.774	-0.0245595\\
0.776	-0.0245611\\
0.778	-0.0245626\\
0.78	-0.0245641\\
0.782	-0.0245656\\
0.784	-0.0245671\\
0.786	-0.0245685\\
0.788	-0.02457\\
0.79	-0.0245714\\
0.792	-0.0245727\\
0.794	-0.0245741\\
0.796	-0.0245755\\
0.798	-0.0245768\\
0.8	-0.0245781\\
0.802	-0.0245794\\
0.804	-0.0245807\\
0.806	-0.0245819\\
0.808	-0.0245832\\
0.81	-0.0245844\\
0.812	-0.0245857\\
0.814	-0.0245869\\
0.816	-0.0245881\\
0.818	-0.0245893\\
0.82	-0.0245904\\
0.822	-0.0245916\\
0.824	-0.0245928\\
0.826	-0.0245939\\
0.828	-0.0245951\\
0.83	-0.0245962\\
0.832	-0.0245973\\
0.834	-0.0245984\\
0.836	-0.0245995\\
0.838	-0.0246006\\
0.84	-0.0246017\\
0.842	-0.0246028\\
0.844	-0.0246039\\
0.846	-0.0246049\\
0.848	-0.024606\\
0.85	-0.0246071\\
0.852	-0.0246081\\
0.854	-0.0246092\\
0.856	-0.0246102\\
0.858	-0.0246112\\
0.86	-0.0246123\\
0.862	-0.0246133\\
0.864	-0.0246143\\
0.866	-0.0246153\\
0.868	-0.0246163\\
0.87	-0.0246173\\
0.872	-0.0246183\\
0.874	-0.0246193\\
0.876	-0.0246203\\
0.878	-0.0246213\\
0.88	-0.0246222\\
0.882	-0.0246232\\
0.884	-0.0246242\\
0.886	-0.0246252\\
0.888	-0.0246261\\
0.89	-0.0246271\\
0.892	-0.024628\\
0.894	-0.024629\\
0.896	-0.0246299\\
0.898	-0.0246309\\
0.9	-0.0246318\\
0.902	-0.0246327\\
0.904	-0.0246337\\
0.906	-0.0246346\\
0.908	-0.0246355\\
0.91	-0.0246364\\
0.912	-0.0246374\\
0.914	-0.0246383\\
0.916	-0.0246392\\
0.918	-0.0246401\\
0.92	-0.024641\\
0.922	-0.0246419\\
0.924	-0.0246428\\
0.926	-0.0246436\\
0.928	-0.0246445\\
0.93	-0.0246454\\
0.932	-0.0246463\\
0.934	-0.0246472\\
0.936	-0.024648\\
0.938	-0.0246489\\
0.94	-0.0246497\\
0.942	-0.0246506\\
0.944	-0.0246514\\
0.946	-0.0246523\\
0.948	-0.0246531\\
0.95	-0.024654\\
0.952	-0.0246548\\
0.954	-0.0246556\\
0.956	-0.0246564\\
0.958	-0.0246572\\
0.96	-0.0246581\\
0.962	-0.0246589\\
0.964	-0.0246597\\
0.966	-0.0246605\\
0.968	-0.0246612\\
0.97	-0.024662\\
0.972	-0.0246628\\
0.974	-0.0246636\\
0.976	-0.0246643\\
0.978	-0.0246651\\
0.98	-0.0246659\\
0.982	-0.0246666\\
0.984	-0.0246674\\
0.986	-0.0246681\\
0.988	-0.0246688\\
0.99	-0.0246696\\
0.992	-0.0246703\\
0.994	-0.024671\\
0.996	-0.0246717\\
0.998	-0.0246724\\
1	-0.0246731\\
1.002	-0.0246738\\
1.004	-0.0246745\\
1.006	-0.0246752\\
1.008	-0.0246759\\
1.01	-0.0246765\\
1.012	-0.0246772\\
1.014	-0.0246779\\
1.016	-0.0246785\\
1.018	-0.0246792\\
1.02	-0.0246798\\
1.022	-0.0246804\\
1.024	-0.0246811\\
1.026	-0.0246817\\
1.028	-0.0246823\\
1.03	-0.0246829\\
1.032	-0.0246836\\
1.034	-0.0246842\\
1.036	-0.0246848\\
1.038	-0.0246853\\
1.04	-0.0246859\\
1.042	-0.0246865\\
1.044	-0.0246871\\
1.046	-0.0246877\\
1.048	-0.0246882\\
1.05	-0.0246888\\
1.052	-0.0246894\\
1.054	-0.0246899\\
1.056	-0.0246905\\
1.058	-0.024691\\
1.06	-0.0246915\\
1.062	-0.0246921\\
1.064	-0.0246926\\
1.066	-0.0246931\\
1.068	-0.0246937\\
1.07	-0.0246942\\
1.072	-0.0246947\\
1.074	-0.0246952\\
1.076	-0.0246957\\
1.078	-0.0246962\\
1.08	-0.0246967\\
1.082	-0.0246972\\
1.084	-0.0246977\\
1.086	-0.0246982\\
1.088	-0.0246986\\
1.09	-0.0246991\\
1.092	-0.0246996\\
1.094	-0.0247001\\
1.096	-0.0247005\\
1.098	-0.024701\\
1.1	-0.0247014\\
1.102	-0.0247019\\
1.104	-0.0247024\\
1.106	-0.0247028\\
1.108	-0.0247032\\
1.11	-0.0247037\\
1.112	-0.0247041\\
1.114	-0.0247046\\
1.116	-0.024705\\
1.118	-0.0247054\\
1.12	-0.0247059\\
1.122	-0.0247063\\
1.124	-0.0247067\\
1.126	-0.0247071\\
1.128	-0.0247075\\
1.13	-0.0247079\\
1.132	-0.0247084\\
1.134	-0.0247088\\
1.136	-0.0247092\\
1.138	-0.0247096\\
1.14	-0.02471\\
1.142	-0.0247104\\
1.144	-0.0247108\\
1.146	-0.0247112\\
1.148	-0.0247115\\
1.15	-0.0247119\\
1.152	-0.0247123\\
1.154	-0.0247127\\
1.156	-0.0247131\\
1.158	-0.0247135\\
1.16	-0.0247138\\
1.162	-0.0247142\\
1.164	-0.0247146\\
1.166	-0.0247149\\
1.168	-0.0247153\\
1.17	-0.0247157\\
1.172	-0.024716\\
1.174	-0.0247164\\
1.176	-0.0247168\\
1.178	-0.0247171\\
1.18	-0.0247175\\
1.182	-0.0247178\\
1.184	-0.0247182\\
1.186	-0.0247185\\
1.188	-0.0247189\\
1.19	-0.0247192\\
1.192	-0.0247195\\
1.194	-0.0247199\\
1.196	-0.0247202\\
1.198	-0.0247206\\
1.2	-0.0247209\\
1.202	-0.0247212\\
1.204	-0.0247216\\
1.206	-0.0247219\\
1.208	-0.0247222\\
1.21	-0.0247225\\
1.212	-0.0247229\\
1.214	-0.0247232\\
1.216	-0.0247235\\
1.218	-0.0247238\\
1.22	-0.0247242\\
1.222	-0.0247245\\
1.224	-0.0247248\\
1.226	-0.0247251\\
1.228	-0.0247254\\
1.23	-0.0247257\\
1.232	-0.024726\\
1.234	-0.0247263\\
1.236	-0.0247266\\
1.238	-0.0247269\\
1.24	-0.0247272\\
1.242	-0.0247275\\
1.244	-0.0247278\\
1.246	-0.0247281\\
1.248	-0.0247284\\
1.25	-0.0247287\\
1.252	-0.024729\\
1.254	-0.0247293\\
1.256	-0.0247296\\
1.258	-0.0247298\\
1.26	-0.0247301\\
1.262	-0.0247304\\
1.264	-0.0247307\\
1.266	-0.024731\\
1.268	-0.0247312\\
1.27	-0.0247315\\
1.272	-0.0247318\\
1.274	-0.0247321\\
1.276	-0.0247323\\
1.278	-0.0247326\\
1.28	-0.0247328\\
1.282	-0.0247331\\
1.284	-0.0247334\\
1.286	-0.0247336\\
1.288	-0.0247339\\
1.29	-0.0247341\\
1.292	-0.0247344\\
1.294	-0.0247346\\
1.296	-0.0247349\\
1.298	-0.0247351\\
1.3	-0.0247354\\
1.302	-0.0247356\\
1.304	-0.0247359\\
1.306	-0.0247361\\
1.308	-0.0247363\\
1.31	-0.0247366\\
1.312	-0.0247368\\
1.314	-0.024737\\
1.316	-0.0247373\\
1.318	-0.0247375\\
1.32	-0.0247377\\
1.322	-0.0247379\\
1.324	-0.0247382\\
1.326	-0.0247384\\
1.328	-0.0247386\\
1.33	-0.0247388\\
1.332	-0.024739\\
1.334	-0.0247393\\
1.336	-0.0247395\\
1.338	-0.0247397\\
1.34	-0.0247399\\
1.342	-0.0247401\\
1.344	-0.0247403\\
1.346	-0.0247405\\
1.348	-0.0247407\\
1.35	-0.0247409\\
1.352	-0.0247411\\
1.354	-0.0247413\\
1.356	-0.0247415\\
1.358	-0.0247417\\
1.36	-0.0247419\\
1.362	-0.0247421\\
1.364	-0.0247423\\
1.366	-0.0247425\\
1.368	-0.0247427\\
1.37	-0.0247428\\
1.372	-0.024743\\
1.374	-0.0247432\\
1.376	-0.0247434\\
1.378	-0.0247436\\
1.38	-0.0247438\\
1.382	-0.0247439\\
1.384	-0.0247441\\
1.386	-0.0247443\\
1.388	-0.0247445\\
1.39	-0.0247446\\
1.392	-0.0247448\\
1.394	-0.024745\\
1.396	-0.0247452\\
1.398	-0.0247453\\
1.4	-0.0247455\\
1.402	-0.0247457\\
1.404	-0.0247458\\
1.406	-0.024746\\
1.408	-0.0247462\\
1.41	-0.0247463\\
1.412	-0.0247465\\
1.414	-0.0247467\\
1.416	-0.0247468\\
1.418	-0.024747\\
1.42	-0.0247471\\
1.422	-0.0247473\\
1.424	-0.0247475\\
1.426	-0.0247476\\
1.428	-0.0247478\\
1.43	-0.0247479\\
1.432	-0.0247481\\
1.434	-0.0247482\\
1.436	-0.0247484\\
1.438	-0.0247485\\
1.44	-0.0247487\\
1.442	-0.0247488\\
1.444	-0.024749\\
1.446	-0.0247491\\
1.448	-0.0247493\\
1.45	-0.0247494\\
1.452	-0.0247496\\
1.454	-0.0247497\\
1.456	-0.0247499\\
1.458	-0.02475\\
1.46	-0.0247502\\
1.462	-0.0247503\\
1.464	-0.0247505\\
1.466	-0.0247506\\
1.468	-0.0247507\\
1.47	-0.0247509\\
1.472	-0.024751\\
1.474	-0.0247512\\
1.476	-0.0247513\\
1.478	-0.0247514\\
1.48	-0.0247516\\
1.482	-0.0247517\\
1.484	-0.0247518\\
1.486	-0.024752\\
1.488	-0.0247521\\
1.49	-0.0247522\\
1.492	-0.0247524\\
1.494	-0.0247525\\
1.496	-0.0247526\\
1.498	-0.0247528\\
1.5	-0.0247529\\
1.502	-0.024753\\
1.504	-0.0247531\\
1.506	-0.0247533\\
1.508	-0.0247534\\
1.51	-0.0247535\\
1.512	-0.0247536\\
1.514	-0.0247538\\
1.516	-0.0247539\\
1.518	-0.024754\\
1.52	-0.0247541\\
1.522	-0.0247543\\
1.524	-0.0247544\\
1.526	-0.0247545\\
1.528	-0.0247546\\
1.53	-0.0247547\\
1.532	-0.0247548\\
1.534	-0.024755\\
1.536	-0.0247551\\
1.538	-0.0247552\\
1.54	-0.0247553\\
1.542	-0.0247554\\
1.544	-0.0247555\\
1.546	-0.0247556\\
1.548	-0.0247557\\
1.55	-0.0247559\\
1.552	-0.024756\\
1.554	-0.0247561\\
1.556	-0.0247562\\
1.558	-0.0247563\\
1.56	-0.0247564\\
1.562	-0.0247565\\
1.564	-0.0247566\\
1.566	-0.0247567\\
1.568	-0.0247568\\
1.57	-0.0247569\\
1.572	-0.024757\\
1.574	-0.0247571\\
1.576	-0.0247572\\
1.578	-0.0247573\\
1.58	-0.0247574\\
1.582	-0.0247575\\
1.584	-0.0247576\\
1.586	-0.0247577\\
1.588	-0.0247578\\
1.59	-0.0247579\\
1.592	-0.024758\\
1.594	-0.0247581\\
1.596	-0.0247581\\
1.598	-0.0247582\\
1.6	-0.0247583\\
1.602	-0.0247584\\
1.604	-0.0247585\\
1.606	-0.0247586\\
1.608	-0.0247587\\
1.61	-0.0247588\\
1.612	-0.0247589\\
1.614	-0.024759\\
1.616	-0.024759\\
1.618	-0.0247591\\
1.62	-0.0247592\\
1.622	-0.0247593\\
1.624	-0.0247594\\
1.626	-0.0247595\\
1.628	-0.0247595\\
1.63	-0.0247596\\
1.632	-0.0247597\\
1.634	-0.0247598\\
1.636	-0.0247599\\
1.638	-0.02476\\
1.64	-0.02476\\
1.642	-0.0247601\\
1.644	-0.0247602\\
1.646	-0.0247603\\
1.648	-0.0247604\\
1.65	-0.0247604\\
1.652	-0.0247605\\
1.654	-0.0247606\\
1.656	-0.0247607\\
1.658	-0.0247607\\
1.66	-0.0247608\\
1.662	-0.0247609\\
1.664	-0.024761\\
1.666	-0.024761\\
1.668	-0.0247611\\
1.67	-0.0247612\\
1.672	-0.0247613\\
1.674	-0.0247613\\
1.676	-0.0247614\\
1.678	-0.0247615\\
1.68	-0.0247616\\
1.682	-0.0247616\\
1.684	-0.0247617\\
1.686	-0.0247618\\
1.688	-0.0247618\\
1.69	-0.0247619\\
1.692	-0.024762\\
1.694	-0.024762\\
1.696	-0.0247621\\
1.698	-0.0247622\\
1.7	-0.0247622\\
1.702	-0.0247623\\
1.704	-0.0247624\\
1.706	-0.0247624\\
1.708	-0.0247625\\
1.71	-0.0247626\\
1.712	-0.0247626\\
1.714	-0.0247627\\
1.716	-0.0247628\\
1.718	-0.0247628\\
1.72	-0.0247629\\
1.722	-0.024763\\
1.724	-0.024763\\
1.726	-0.0247631\\
1.728	-0.0247632\\
1.73	-0.0247632\\
1.732	-0.0247633\\
1.734	-0.0247633\\
1.736	-0.0247634\\
1.738	-0.0247635\\
1.74	-0.0247635\\
1.742	-0.0247636\\
1.744	-0.0247636\\
1.746	-0.0247637\\
1.748	-0.0247638\\
1.75	-0.0247638\\
1.752	-0.0247639\\
1.754	-0.0247639\\
1.756	-0.024764\\
1.758	-0.024764\\
1.76	-0.0247641\\
1.762	-0.0247642\\
1.764	-0.0247642\\
1.766	-0.0247643\\
1.768	-0.0247643\\
1.77	-0.0247644\\
1.772	-0.0247644\\
1.774	-0.0247645\\
1.776	-0.0247645\\
1.778	-0.0247646\\
1.78	-0.0247646\\
1.782	-0.0247647\\
1.784	-0.0247648\\
1.786	-0.0247648\\
1.788	-0.0247649\\
1.79	-0.0247649\\
1.792	-0.024765\\
1.794	-0.024765\\
1.796	-0.0247651\\
1.798	-0.0247651\\
1.8	-0.0247652\\
1.802	-0.0247652\\
1.804	-0.0247653\\
1.806	-0.0247653\\
1.808	-0.0247654\\
1.81	-0.0247654\\
1.812	-0.0247654\\
1.814	-0.0247655\\
1.816	-0.0247655\\
1.818	-0.0247656\\
1.82	-0.0247656\\
1.822	-0.0247657\\
1.824	-0.0247657\\
1.826	-0.0247658\\
1.828	-0.0247658\\
1.83	-0.0247659\\
1.832	-0.0247659\\
1.834	-0.024766\\
1.836	-0.024766\\
1.838	-0.024766\\
1.84	-0.0247661\\
1.842	-0.0247661\\
1.844	-0.0247662\\
1.846	-0.0247662\\
1.848	-0.0247663\\
1.85	-0.0247663\\
1.852	-0.0247663\\
1.854	-0.0247664\\
1.856	-0.0247664\\
1.858	-0.0247665\\
1.86	-0.0247665\\
1.862	-0.0247665\\
1.864	-0.0247666\\
1.866	-0.0247666\\
1.868	-0.0247667\\
1.87	-0.0247667\\
1.872	-0.0247667\\
1.874	-0.0247668\\
1.876	-0.0247668\\
1.878	-0.0247669\\
1.88	-0.0247669\\
1.882	-0.0247669\\
1.884	-0.024767\\
1.886	-0.024767\\
1.888	-0.0247671\\
1.89	-0.0247671\\
1.892	-0.0247671\\
1.894	-0.0247672\\
1.896	-0.0247672\\
1.898	-0.0247672\\
1.9	-0.0247673\\
1.902	-0.0247673\\
1.904	-0.0247673\\
1.906	-0.0247674\\
1.908	-0.0247674\\
1.91	-0.0247674\\
1.912	-0.0247675\\
1.914	-0.0247675\\
1.916	-0.0247676\\
1.918	-0.0247676\\
1.92	-0.0247676\\
1.922	-0.0247677\\
1.924	-0.0247677\\
1.926	-0.0247677\\
1.928	-0.0247678\\
1.93	-0.0247678\\
1.932	-0.0247678\\
1.934	-0.0247679\\
1.936	-0.0247679\\
1.938	-0.0247679\\
1.94	-0.024768\\
1.942	-0.024768\\
1.944	-0.024768\\
1.946	-0.0247681\\
1.948	-0.0247681\\
1.95	-0.0247681\\
1.952	-0.0247681\\
1.954	-0.0247682\\
1.956	-0.0247682\\
1.958	-0.0247682\\
1.96	-0.0247683\\
1.962	-0.0247683\\
1.964	-0.0247683\\
1.966	-0.0247684\\
1.968	-0.0247684\\
1.97	-0.0247684\\
1.972	-0.0247684\\
1.974	-0.0247685\\
1.976	-0.0247685\\
1.978	-0.0247685\\
1.98	-0.0247686\\
1.982	-0.0247686\\
1.984	-0.0247686\\
1.986	-0.0247686\\
1.988	-0.0247687\\
1.99	-0.0247687\\
1.992	-0.0247687\\
1.994	-0.0247688\\
1.996	-0.0247688\\
1.998	-0.0247688\\
2	-0.0247688\\
};
\addplot [color=mycolor11, line width=1.0pt, forget plot]
  table[row sep=crcr]{%
-0.024	0\\
-0.022	0\\
-0.02	0\\
-0.018	0\\
-0.016	0\\
-0.014	0\\
-0.012	0\\
-0.01	0\\
-0.008	0\\
-0.006	0\\
-0.004	0\\
-0.002	0\\
0	0\\
0.002	-0.0018468\\
0.004	-0.00342853\\
0.006	-0.00479339\\
0.008	-0.00598113\\
0.01	-0.00702436\\
0.012	-0.0079498\\
0.014	-0.0087792\\
0.016	-0.00953022\\
0.018	-0.0102171\\
0.02	-0.0108512\\
0.022	-0.0114418\\
0.024	-0.0119959\\
0.026	-0.0125192\\
0.028	-0.0130161\\
0.03	-0.0134899\\
0.032	-0.0139432\\
0.034	-0.0143777\\
0.036	-0.014795\\
0.038	-0.0151959\\
0.04	-0.0155811\\
0.042	-0.0159512\\
0.044	-0.0163064\\
0.046	-0.016647\\
0.048	-0.0169731\\
0.05	-0.017285\\
0.052	-0.0175828\\
0.054	-0.0178667\\
0.056	-0.0181369\\
0.058	-0.0183936\\
0.06	-0.0186372\\
0.062	-0.0188679\\
0.064	-0.019086\\
0.066	-0.0192921\\
0.068	-0.0194865\\
0.07	-0.0196697\\
0.072	-0.0198421\\
0.074	-0.0200042\\
0.076	-0.0201565\\
0.078	-0.0202995\\
0.08	-0.0204336\\
0.082	-0.0205595\\
0.084	-0.0206774\\
0.086	-0.020788\\
0.088	-0.0208917\\
0.09	-0.0209888\\
0.092	-0.0210799\\
0.094	-0.0211654\\
0.096	-0.0212456\\
0.098	-0.0213208\\
0.1	-0.0213916\\
0.102	-0.0214581\\
0.104	-0.0215207\\
0.106	-0.0215797\\
0.108	-0.0216354\\
0.11	-0.0216879\\
0.112	-0.0217375\\
0.114	-0.0217844\\
0.116	-0.0218288\\
0.118	-0.0218709\\
0.12	-0.0219108\\
0.122	-0.0219485\\
0.124	-0.0219843\\
0.126	-0.0220183\\
0.128	-0.0220504\\
0.13	-0.0220808\\
0.132	-0.0221096\\
0.134	-0.0221367\\
0.136	-0.0221622\\
0.138	-0.0221861\\
0.14	-0.0222086\\
0.142	-0.0222295\\
0.144	-0.0222488\\
0.146	-0.0222667\\
0.148	-0.0222832\\
0.15	-0.0222981\\
0.152	-0.0223116\\
0.154	-0.0223237\\
0.156	-0.0223343\\
0.158	-0.0223436\\
0.16	-0.0223515\\
0.162	-0.0223581\\
0.164	-0.0223634\\
0.166	-0.0223675\\
0.168	-0.0223703\\
0.17	-0.0223721\\
0.172	-0.0223728\\
0.174	-0.0223726\\
0.176	-0.0223714\\
0.178	-0.0223694\\
0.18	-0.0223667\\
0.182	-0.0223633\\
0.184	-0.0223593\\
0.186	-0.0223549\\
0.188	-0.0223501\\
0.19	-0.022345\\
0.192	-0.0223398\\
0.194	-0.0223345\\
0.196	-0.0223291\\
0.198	-0.022324\\
0.2	-0.022319\\
0.202	-0.0223143\\
0.204	-0.0223099\\
0.206	-0.0223061\\
0.208	-0.0223028\\
0.21	-0.0223001\\
0.212	-0.022298\\
0.214	-0.0222968\\
0.216	-0.0222963\\
0.218	-0.0222966\\
0.22	-0.0222979\\
0.222	-0.0223001\\
0.224	-0.0223033\\
0.226	-0.0223074\\
0.228	-0.0223125\\
0.23	-0.0223187\\
0.232	-0.0223259\\
0.234	-0.0223342\\
0.236	-0.0223434\\
0.238	-0.0223537\\
0.24	-0.022365\\
0.242	-0.0223773\\
0.244	-0.0223906\\
0.246	-0.0224048\\
0.248	-0.0224199\\
0.25	-0.0224359\\
0.252	-0.0224527\\
0.254	-0.0224704\\
0.256	-0.0224888\\
0.258	-0.0225079\\
0.26	-0.0225278\\
0.262	-0.0225482\\
0.264	-0.0225693\\
0.266	-0.0225908\\
0.268	-0.0226129\\
0.27	-0.0226355\\
0.272	-0.0226584\\
0.274	-0.0226817\\
0.276	-0.0227053\\
0.278	-0.0227291\\
0.28	-0.0227532\\
0.282	-0.0227774\\
0.284	-0.0228018\\
0.286	-0.0228262\\
0.288	-0.0228507\\
0.29	-0.0228753\\
0.292	-0.0228998\\
0.294	-0.0229242\\
0.296	-0.0229486\\
0.298	-0.0229728\\
0.3	-0.0229969\\
0.302	-0.0230207\\
0.304	-0.0230444\\
0.306	-0.0230679\\
0.308	-0.0230911\\
0.31	-0.023114\\
0.312	-0.0231367\\
0.314	-0.023159\\
0.316	-0.023181\\
0.318	-0.0232027\\
0.32	-0.0232241\\
0.322	-0.023245\\
0.324	-0.0232656\\
0.326	-0.0232858\\
0.328	-0.0233057\\
0.33	-0.0233251\\
0.332	-0.0233441\\
0.334	-0.0233628\\
0.336	-0.023381\\
0.338	-0.0233988\\
0.34	-0.0234162\\
0.342	-0.0234332\\
0.344	-0.0234498\\
0.346	-0.023466\\
0.348	-0.0234818\\
0.35	-0.0234972\\
0.352	-0.0235122\\
0.354	-0.0235268\\
0.356	-0.023541\\
0.358	-0.0235548\\
0.36	-0.0235683\\
0.362	-0.0235814\\
0.364	-0.0235941\\
0.366	-0.0236065\\
0.368	-0.0236185\\
0.37	-0.0236302\\
0.372	-0.0236415\\
0.374	-0.0236526\\
0.376	-0.0236633\\
0.378	-0.0236737\\
0.38	-0.0236839\\
0.382	-0.0236937\\
0.384	-0.0237033\\
0.386	-0.0237127\\
0.388	-0.0237218\\
0.39	-0.0237306\\
0.392	-0.0237392\\
0.394	-0.0237477\\
0.396	-0.0237559\\
0.398	-0.0237639\\
0.4	-0.0237718\\
0.402	-0.0237794\\
0.404	-0.023787\\
0.406	-0.0237944\\
0.408	-0.0238016\\
0.41	-0.0238088\\
0.412	-0.0238158\\
0.414	-0.0238227\\
0.416	-0.0238296\\
0.418	-0.0238363\\
0.42	-0.023843\\
0.422	-0.0238496\\
0.424	-0.0238562\\
0.426	-0.0238628\\
0.428	-0.0238693\\
0.43	-0.0238758\\
0.432	-0.0238823\\
0.434	-0.0238887\\
0.436	-0.0238952\\
0.438	-0.0239017\\
0.44	-0.0239082\\
0.442	-0.0239147\\
0.444	-0.0239212\\
0.446	-0.0239277\\
0.448	-0.0239343\\
0.45	-0.0239409\\
0.452	-0.0239475\\
0.454	-0.0239542\\
0.456	-0.0239609\\
0.458	-0.0239676\\
0.46	-0.0239744\\
0.462	-0.0239812\\
0.464	-0.0239881\\
0.466	-0.023995\\
0.468	-0.0240019\\
0.47	-0.0240088\\
0.472	-0.0240158\\
0.474	-0.0240228\\
0.476	-0.0240298\\
0.478	-0.0240369\\
0.48	-0.0240439\\
0.482	-0.024051\\
0.484	-0.0240581\\
0.486	-0.0240651\\
0.488	-0.0240722\\
0.49	-0.0240793\\
0.492	-0.0240863\\
0.494	-0.0240934\\
0.496	-0.0241004\\
0.498	-0.0241073\\
0.5	-0.0241143\\
0.502	-0.0241212\\
0.504	-0.024128\\
0.506	-0.0241348\\
0.508	-0.0241416\\
0.51	-0.0241482\\
0.512	-0.0241548\\
0.514	-0.0241614\\
0.516	-0.0241678\\
0.518	-0.0241742\\
0.52	-0.0241805\\
0.522	-0.0241867\\
0.524	-0.0241928\\
0.526	-0.0241988\\
0.528	-0.0242047\\
0.53	-0.0242105\\
0.532	-0.0242162\\
0.534	-0.0242217\\
0.536	-0.0242272\\
0.538	-0.0242325\\
0.54	-0.0242378\\
0.542	-0.0242429\\
0.544	-0.0242479\\
0.546	-0.0242528\\
0.548	-0.0242575\\
0.55	-0.0242622\\
0.552	-0.0242667\\
0.554	-0.0242711\\
0.556	-0.0242754\\
0.558	-0.0242796\\
0.56	-0.0242837\\
0.562	-0.0242877\\
0.564	-0.0242915\\
0.566	-0.0242953\\
0.568	-0.024299\\
0.57	-0.0243025\\
0.572	-0.024306\\
0.574	-0.0243094\\
0.576	-0.0243126\\
0.578	-0.0243158\\
0.58	-0.024319\\
0.582	-0.024322\\
0.584	-0.024325\\
0.586	-0.0243279\\
0.588	-0.0243307\\
0.59	-0.0243335\\
0.592	-0.0243362\\
0.594	-0.0243389\\
0.596	-0.0243416\\
0.598	-0.0243441\\
0.6	-0.0243467\\
0.602	-0.0243492\\
0.604	-0.0243517\\
0.606	-0.0243541\\
0.608	-0.0243566\\
0.61	-0.024359\\
0.612	-0.0243614\\
0.614	-0.0243638\\
0.616	-0.0243662\\
0.618	-0.0243685\\
0.62	-0.0243709\\
0.622	-0.0243733\\
0.624	-0.0243757\\
0.626	-0.024378\\
0.628	-0.0243804\\
0.63	-0.0243828\\
0.632	-0.0243852\\
0.634	-0.0243877\\
0.636	-0.0243901\\
0.638	-0.0243925\\
0.64	-0.024395\\
0.642	-0.0243975\\
0.644	-0.0244\\
0.646	-0.0244025\\
0.648	-0.024405\\
0.65	-0.0244076\\
0.652	-0.0244101\\
0.654	-0.0244127\\
0.656	-0.0244153\\
0.658	-0.0244179\\
0.66	-0.0244206\\
0.662	-0.0244232\\
0.664	-0.0244259\\
0.666	-0.0244285\\
0.668	-0.0244312\\
0.67	-0.0244339\\
0.672	-0.0244366\\
0.674	-0.0244393\\
0.676	-0.024442\\
0.678	-0.0244447\\
0.68	-0.0244474\\
0.682	-0.0244501\\
0.684	-0.0244528\\
0.686	-0.0244555\\
0.688	-0.0244582\\
0.69	-0.0244609\\
0.692	-0.0244636\\
0.694	-0.0244663\\
0.696	-0.0244689\\
0.698	-0.0244716\\
0.7	-0.0244742\\
0.702	-0.0244768\\
0.704	-0.0244794\\
0.706	-0.024482\\
0.708	-0.0244846\\
0.71	-0.0244871\\
0.712	-0.0244896\\
0.714	-0.0244921\\
0.716	-0.0244945\\
0.718	-0.024497\\
0.72	-0.0244994\\
0.722	-0.0245018\\
0.724	-0.0245041\\
0.726	-0.0245064\\
0.728	-0.0245087\\
0.73	-0.024511\\
0.732	-0.0245132\\
0.734	-0.0245154\\
0.736	-0.0245176\\
0.738	-0.0245197\\
0.74	-0.0245218\\
0.742	-0.0245239\\
0.744	-0.0245259\\
0.746	-0.0245279\\
0.748	-0.0245299\\
0.75	-0.0245318\\
0.752	-0.0245338\\
0.754	-0.0245356\\
0.756	-0.0245375\\
0.758	-0.0245393\\
0.76	-0.0245411\\
0.762	-0.0245429\\
0.764	-0.0245446\\
0.766	-0.0245463\\
0.768	-0.024548\\
0.77	-0.0245497\\
0.772	-0.0245513\\
0.774	-0.0245529\\
0.776	-0.0245545\\
0.778	-0.0245561\\
0.78	-0.0245576\\
0.782	-0.0245591\\
0.784	-0.0245606\\
0.786	-0.0245621\\
0.788	-0.0245636\\
0.79	-0.024565\\
0.792	-0.0245664\\
0.794	-0.0245678\\
0.796	-0.0245692\\
0.798	-0.0245706\\
0.8	-0.0245719\\
0.802	-0.0245732\\
0.804	-0.0245746\\
0.806	-0.0245759\\
0.808	-0.0245772\\
0.81	-0.0245784\\
0.812	-0.0245797\\
0.814	-0.0245809\\
0.816	-0.0245822\\
0.818	-0.0245834\\
0.82	-0.0245846\\
0.822	-0.0245858\\
0.824	-0.024587\\
0.826	-0.0245882\\
0.828	-0.0245894\\
0.83	-0.0245906\\
0.832	-0.0245917\\
0.834	-0.0245929\\
0.836	-0.024594\\
0.838	-0.0245951\\
0.84	-0.0245963\\
0.842	-0.0245974\\
0.844	-0.0245985\\
0.846	-0.0245996\\
0.848	-0.0246007\\
0.85	-0.0246018\\
0.852	-0.0246029\\
0.854	-0.0246039\\
0.856	-0.024605\\
0.858	-0.0246061\\
0.86	-0.0246071\\
0.862	-0.0246082\\
0.864	-0.0246092\\
0.866	-0.0246103\\
0.868	-0.0246113\\
0.87	-0.0246123\\
0.872	-0.0246134\\
0.874	-0.0246144\\
0.876	-0.0246154\\
0.878	-0.0246164\\
0.88	-0.0246174\\
0.882	-0.0246184\\
0.884	-0.0246194\\
0.886	-0.0246204\\
0.888	-0.0246214\\
0.89	-0.0246224\\
0.892	-0.0246234\\
0.894	-0.0246243\\
0.896	-0.0246253\\
0.898	-0.0246263\\
0.9	-0.0246272\\
0.902	-0.0246282\\
0.904	-0.0246292\\
0.906	-0.0246301\\
0.908	-0.024631\\
0.91	-0.024632\\
0.912	-0.0246329\\
0.914	-0.0246339\\
0.916	-0.0246348\\
0.918	-0.0246357\\
0.92	-0.0246366\\
0.922	-0.0246375\\
0.924	-0.0246385\\
0.926	-0.0246394\\
0.928	-0.0246403\\
0.93	-0.0246412\\
0.932	-0.0246421\\
0.934	-0.0246429\\
0.936	-0.0246438\\
0.938	-0.0246447\\
0.94	-0.0246456\\
0.942	-0.0246465\\
0.944	-0.0246473\\
0.946	-0.0246482\\
0.948	-0.024649\\
0.95	-0.0246499\\
0.952	-0.0246507\\
0.954	-0.0246516\\
0.956	-0.0246524\\
0.958	-0.0246532\\
0.96	-0.0246541\\
0.962	-0.0246549\\
0.964	-0.0246557\\
0.966	-0.0246565\\
0.968	-0.0246573\\
0.97	-0.0246581\\
0.972	-0.0246589\\
0.974	-0.0246597\\
0.976	-0.0246605\\
0.978	-0.0246613\\
0.98	-0.024662\\
0.982	-0.0246628\\
0.984	-0.0246636\\
0.986	-0.0246643\\
0.988	-0.0246651\\
0.99	-0.0246658\\
0.992	-0.0246665\\
0.994	-0.0246673\\
0.996	-0.024668\\
0.998	-0.0246687\\
1	-0.0246694\\
1.002	-0.0246702\\
1.004	-0.0246709\\
1.006	-0.0246716\\
1.008	-0.0246722\\
1.01	-0.0246729\\
1.012	-0.0246736\\
1.014	-0.0246743\\
1.016	-0.024675\\
1.018	-0.0246756\\
1.02	-0.0246763\\
1.022	-0.0246769\\
1.024	-0.0246776\\
1.026	-0.0246782\\
1.028	-0.0246789\\
1.03	-0.0246795\\
1.032	-0.0246801\\
1.034	-0.0246807\\
1.036	-0.0246813\\
1.038	-0.024682\\
1.04	-0.0246826\\
1.042	-0.0246832\\
1.044	-0.0246837\\
1.046	-0.0246843\\
1.048	-0.0246849\\
1.05	-0.0246855\\
1.052	-0.0246861\\
1.054	-0.0246866\\
1.056	-0.0246872\\
1.058	-0.0246878\\
1.06	-0.0246883\\
1.062	-0.0246889\\
1.064	-0.0246894\\
1.066	-0.02469\\
1.068	-0.0246905\\
1.07	-0.024691\\
1.072	-0.0246915\\
1.074	-0.0246921\\
1.076	-0.0246926\\
1.078	-0.0246931\\
1.08	-0.0246936\\
1.082	-0.0246941\\
1.084	-0.0246946\\
1.086	-0.0246951\\
1.088	-0.0246956\\
1.09	-0.0246961\\
1.092	-0.0246966\\
1.094	-0.0246971\\
1.096	-0.0246976\\
1.098	-0.024698\\
1.1	-0.0246985\\
1.102	-0.024699\\
1.104	-0.0246994\\
1.106	-0.0246999\\
1.108	-0.0247004\\
1.11	-0.0247008\\
1.112	-0.0247013\\
1.114	-0.0247017\\
1.116	-0.0247022\\
1.118	-0.0247026\\
1.12	-0.0247031\\
1.122	-0.0247035\\
1.124	-0.0247039\\
1.126	-0.0247044\\
1.128	-0.0247048\\
1.13	-0.0247052\\
1.132	-0.0247056\\
1.134	-0.0247061\\
1.136	-0.0247065\\
1.138	-0.0247069\\
1.14	-0.0247073\\
1.142	-0.0247077\\
1.144	-0.0247081\\
1.146	-0.0247085\\
1.148	-0.0247089\\
1.15	-0.0247093\\
1.152	-0.0247097\\
1.154	-0.0247101\\
1.156	-0.0247105\\
1.158	-0.0247109\\
1.16	-0.0247113\\
1.162	-0.0247117\\
1.164	-0.024712\\
1.166	-0.0247124\\
1.168	-0.0247128\\
1.17	-0.0247132\\
1.172	-0.0247135\\
1.174	-0.0247139\\
1.176	-0.0247143\\
1.178	-0.0247147\\
1.18	-0.024715\\
1.182	-0.0247154\\
1.184	-0.0247157\\
1.186	-0.0247161\\
1.188	-0.0247165\\
1.19	-0.0247168\\
1.192	-0.0247172\\
1.194	-0.0247175\\
1.196	-0.0247179\\
1.198	-0.0247182\\
1.2	-0.0247186\\
1.202	-0.0247189\\
1.204	-0.0247192\\
1.206	-0.0247196\\
1.208	-0.0247199\\
1.21	-0.0247203\\
1.212	-0.0247206\\
1.214	-0.0247209\\
1.216	-0.0247213\\
1.218	-0.0247216\\
1.22	-0.0247219\\
1.222	-0.0247222\\
1.224	-0.0247226\\
1.226	-0.0247229\\
1.228	-0.0247232\\
1.23	-0.0247235\\
1.232	-0.0247238\\
1.234	-0.0247242\\
1.236	-0.0247245\\
1.238	-0.0247248\\
1.24	-0.0247251\\
1.242	-0.0247254\\
1.244	-0.0247257\\
1.246	-0.024726\\
1.248	-0.0247263\\
1.25	-0.0247266\\
1.252	-0.0247269\\
1.254	-0.0247272\\
1.256	-0.0247275\\
1.258	-0.0247278\\
1.26	-0.0247281\\
1.262	-0.0247284\\
1.264	-0.0247287\\
1.266	-0.0247289\\
1.268	-0.0247292\\
1.27	-0.0247295\\
1.272	-0.0247298\\
1.274	-0.0247301\\
1.276	-0.0247303\\
1.278	-0.0247306\\
1.28	-0.0247309\\
1.282	-0.0247312\\
1.284	-0.0247314\\
1.286	-0.0247317\\
1.288	-0.024732\\
1.29	-0.0247322\\
1.292	-0.0247325\\
1.294	-0.0247327\\
1.296	-0.024733\\
1.298	-0.0247332\\
1.3	-0.0247335\\
1.302	-0.0247337\\
1.304	-0.024734\\
1.306	-0.0247342\\
1.308	-0.0247345\\
1.31	-0.0247347\\
1.312	-0.024735\\
1.314	-0.0247352\\
1.316	-0.0247355\\
1.318	-0.0247357\\
1.32	-0.0247359\\
1.322	-0.0247362\\
1.324	-0.0247364\\
1.326	-0.0247366\\
1.328	-0.0247368\\
1.33	-0.0247371\\
1.332	-0.0247373\\
1.334	-0.0247375\\
1.336	-0.0247377\\
1.338	-0.024738\\
1.34	-0.0247382\\
1.342	-0.0247384\\
1.344	-0.0247386\\
1.346	-0.0247388\\
1.348	-0.024739\\
1.35	-0.0247392\\
1.352	-0.0247394\\
1.354	-0.0247396\\
1.356	-0.0247398\\
1.358	-0.0247401\\
1.36	-0.0247403\\
1.362	-0.0247405\\
1.364	-0.0247407\\
1.366	-0.0247408\\
1.368	-0.024741\\
1.37	-0.0247412\\
1.372	-0.0247414\\
1.374	-0.0247416\\
1.376	-0.0247418\\
1.378	-0.024742\\
1.38	-0.0247422\\
1.382	-0.0247424\\
1.384	-0.0247426\\
1.386	-0.0247427\\
1.388	-0.0247429\\
1.39	-0.0247431\\
1.392	-0.0247433\\
1.394	-0.0247435\\
1.396	-0.0247436\\
1.398	-0.0247438\\
1.4	-0.024744\\
1.402	-0.0247442\\
1.404	-0.0247443\\
1.406	-0.0247445\\
1.408	-0.0247447\\
1.41	-0.0247449\\
1.412	-0.024745\\
1.414	-0.0247452\\
1.416	-0.0247454\\
1.418	-0.0247455\\
1.42	-0.0247457\\
1.422	-0.0247459\\
1.424	-0.024746\\
1.426	-0.0247462\\
1.428	-0.0247464\\
1.43	-0.0247465\\
1.432	-0.0247467\\
1.434	-0.0247468\\
1.436	-0.024747\\
1.438	-0.0247472\\
1.44	-0.0247473\\
1.442	-0.0247475\\
1.444	-0.0247476\\
1.446	-0.0247478\\
1.448	-0.024748\\
1.45	-0.0247481\\
1.452	-0.0247483\\
1.454	-0.0247484\\
1.456	-0.0247486\\
1.458	-0.0247487\\
1.46	-0.0247489\\
1.462	-0.024749\\
1.464	-0.0247492\\
1.466	-0.0247493\\
1.468	-0.0247495\\
1.47	-0.0247496\\
1.472	-0.0247497\\
1.474	-0.0247499\\
1.476	-0.02475\\
1.478	-0.0247502\\
1.48	-0.0247503\\
1.482	-0.0247505\\
1.484	-0.0247506\\
1.486	-0.0247507\\
1.488	-0.0247509\\
1.49	-0.024751\\
1.492	-0.0247512\\
1.494	-0.0247513\\
1.496	-0.0247514\\
1.498	-0.0247516\\
1.5	-0.0247517\\
1.502	-0.0247518\\
1.504	-0.024752\\
1.506	-0.0247521\\
1.508	-0.0247522\\
1.51	-0.0247524\\
1.512	-0.0247525\\
1.514	-0.0247526\\
1.516	-0.0247527\\
1.518	-0.0247529\\
1.52	-0.024753\\
1.522	-0.0247531\\
1.524	-0.0247532\\
1.526	-0.0247534\\
1.528	-0.0247535\\
1.53	-0.0247536\\
1.532	-0.0247537\\
1.534	-0.0247539\\
1.536	-0.024754\\
1.538	-0.0247541\\
1.54	-0.0247542\\
1.542	-0.0247543\\
1.544	-0.0247544\\
1.546	-0.0247546\\
1.548	-0.0247547\\
1.55	-0.0247548\\
1.552	-0.0247549\\
1.554	-0.024755\\
1.556	-0.0247551\\
1.558	-0.0247552\\
1.56	-0.0247553\\
1.562	-0.0247555\\
1.564	-0.0247556\\
1.566	-0.0247557\\
1.568	-0.0247558\\
1.57	-0.0247559\\
1.572	-0.024756\\
1.574	-0.0247561\\
1.576	-0.0247562\\
1.578	-0.0247563\\
1.58	-0.0247564\\
1.582	-0.0247565\\
1.584	-0.0247566\\
1.586	-0.0247567\\
1.588	-0.0247568\\
1.59	-0.0247569\\
1.592	-0.024757\\
1.594	-0.0247571\\
1.596	-0.0247572\\
1.598	-0.0247573\\
1.6	-0.0247574\\
1.602	-0.0247575\\
1.604	-0.0247576\\
1.606	-0.0247577\\
1.608	-0.0247578\\
1.61	-0.0247579\\
1.612	-0.0247579\\
1.614	-0.024758\\
1.616	-0.0247581\\
1.618	-0.0247582\\
1.62	-0.0247583\\
1.622	-0.0247584\\
1.624	-0.0247585\\
1.626	-0.0247586\\
1.628	-0.0247587\\
1.63	-0.0247587\\
1.632	-0.0247588\\
1.634	-0.0247589\\
1.636	-0.024759\\
1.638	-0.0247591\\
1.64	-0.0247592\\
1.642	-0.0247593\\
1.644	-0.0247593\\
1.646	-0.0247594\\
1.648	-0.0247595\\
1.65	-0.0247596\\
1.652	-0.0247597\\
1.654	-0.0247597\\
1.656	-0.0247598\\
1.658	-0.0247599\\
1.66	-0.02476\\
1.662	-0.0247601\\
1.664	-0.0247601\\
1.666	-0.0247602\\
1.668	-0.0247603\\
1.67	-0.0247604\\
1.672	-0.0247605\\
1.674	-0.0247605\\
1.676	-0.0247606\\
1.678	-0.0247607\\
1.68	-0.0247608\\
1.682	-0.0247608\\
1.684	-0.0247609\\
1.686	-0.024761\\
1.688	-0.0247611\\
1.69	-0.0247611\\
1.692	-0.0247612\\
1.694	-0.0247613\\
1.696	-0.0247614\\
1.698	-0.0247614\\
1.7	-0.0247615\\
1.702	-0.0247616\\
1.704	-0.0247616\\
1.706	-0.0247617\\
1.708	-0.0247618\\
1.71	-0.0247618\\
1.712	-0.0247619\\
1.714	-0.024762\\
1.716	-0.0247621\\
1.718	-0.0247621\\
1.72	-0.0247622\\
1.722	-0.0247623\\
1.724	-0.0247623\\
1.726	-0.0247624\\
1.728	-0.0247625\\
1.73	-0.0247625\\
1.732	-0.0247626\\
1.734	-0.0247626\\
1.736	-0.0247627\\
1.738	-0.0247628\\
1.74	-0.0247628\\
1.742	-0.0247629\\
1.744	-0.024763\\
1.746	-0.024763\\
1.748	-0.0247631\\
1.75	-0.0247631\\
1.752	-0.0247632\\
1.754	-0.0247633\\
1.756	-0.0247633\\
1.758	-0.0247634\\
1.76	-0.0247635\\
1.762	-0.0247635\\
1.764	-0.0247636\\
1.766	-0.0247636\\
1.768	-0.0247637\\
1.77	-0.0247637\\
1.772	-0.0247638\\
1.774	-0.0247639\\
1.776	-0.0247639\\
1.778	-0.024764\\
1.78	-0.024764\\
1.782	-0.0247641\\
1.784	-0.0247641\\
1.786	-0.0247642\\
1.788	-0.0247642\\
1.79	-0.0247643\\
1.792	-0.0247644\\
1.794	-0.0247644\\
1.796	-0.0247645\\
1.798	-0.0247645\\
1.8	-0.0247646\\
1.802	-0.0247646\\
1.804	-0.0247647\\
1.806	-0.0247647\\
1.808	-0.0247648\\
1.81	-0.0247648\\
1.812	-0.0247649\\
1.814	-0.0247649\\
1.816	-0.024765\\
1.818	-0.024765\\
1.82	-0.0247651\\
1.822	-0.0247651\\
1.824	-0.0247652\\
1.826	-0.0247652\\
1.828	-0.0247653\\
1.83	-0.0247653\\
1.832	-0.0247654\\
1.834	-0.0247654\\
1.836	-0.0247654\\
1.838	-0.0247655\\
1.84	-0.0247655\\
1.842	-0.0247656\\
1.844	-0.0247656\\
1.846	-0.0247657\\
1.848	-0.0247657\\
1.85	-0.0247658\\
1.852	-0.0247658\\
1.854	-0.0247659\\
1.856	-0.0247659\\
1.858	-0.0247659\\
1.86	-0.024766\\
1.862	-0.024766\\
1.864	-0.0247661\\
1.866	-0.0247661\\
1.868	-0.0247662\\
1.87	-0.0247662\\
1.872	-0.0247662\\
1.874	-0.0247663\\
1.876	-0.0247663\\
1.878	-0.0247664\\
1.88	-0.0247664\\
1.882	-0.0247664\\
1.884	-0.0247665\\
1.886	-0.0247665\\
1.888	-0.0247666\\
1.89	-0.0247666\\
1.892	-0.0247666\\
1.894	-0.0247667\\
1.896	-0.0247667\\
1.898	-0.0247668\\
1.9	-0.0247668\\
1.902	-0.0247668\\
1.904	-0.0247669\\
1.906	-0.0247669\\
1.908	-0.024767\\
1.91	-0.024767\\
1.912	-0.024767\\
1.914	-0.0247671\\
1.916	-0.0247671\\
1.918	-0.0247671\\
1.92	-0.0247672\\
1.922	-0.0247672\\
1.924	-0.0247672\\
1.926	-0.0247673\\
1.928	-0.0247673\\
1.93	-0.0247674\\
1.932	-0.0247674\\
1.934	-0.0247674\\
1.936	-0.0247675\\
1.938	-0.0247675\\
1.94	-0.0247675\\
1.942	-0.0247676\\
1.944	-0.0247676\\
1.946	-0.0247676\\
1.948	-0.0247677\\
1.95	-0.0247677\\
1.952	-0.0247677\\
1.954	-0.0247678\\
1.956	-0.0247678\\
1.958	-0.0247678\\
1.96	-0.0247679\\
1.962	-0.0247679\\
1.964	-0.0247679\\
1.966	-0.024768\\
1.968	-0.024768\\
1.97	-0.024768\\
1.972	-0.024768\\
1.974	-0.0247681\\
1.976	-0.0247681\\
1.978	-0.0247681\\
1.98	-0.0247682\\
1.982	-0.0247682\\
1.984	-0.0247682\\
1.986	-0.0247683\\
1.988	-0.0247683\\
1.99	-0.0247683\\
1.992	-0.0247684\\
1.994	-0.0247684\\
1.996	-0.0247684\\
1.998	-0.0247684\\
2	-0.0247685\\
};
\draw[dashed, line width=0.7pt, draw=black] (axis cs:0,-0.023) rectangle (axis cs:0.5,-0.0215);
\end{axis}

\begin{axis}[%
width={\figscale*1.528in},
height={\figscale*0.46in},
at={(\figscale*0.931in,\figscale*1.125in)},
scale only axis,
separate axis lines,
every outer x axis line/.append style={black},
every x tick label/.append style={font=\figticfont},
every x tick/.append style={black},
xmin=-0.025,
xmax=2,
xtick={0,0.5,1,1.5,2},
xticklabels={\empty},
every outer y axis line/.append style={black},
every y tick label/.append style={font=\figticfont},
every y tick/.append style={black},
ymin=0,
ymax=0.201716,
ytick={0,0.2},
ylabel={$\max_i|f_i(t)-\fcoi(t)|$},
axis background/.style={fill=white},
xmajorgrids,
ymajorgrids,
grid style={white!55!black, opacity=0.3},
ylabel style={rotate=-90, at={(axis description cs:-0.145,0.5)}, anchor=east}
]
\addplot [color=mycolor1, line width=1.4pt, forget plot]
  table[row sep=crcr]{%
-0.024	0\\
-0.022	0\\
-0.02	0\\
-0.018	0\\
-0.016	0\\
-0.014	0\\
-0.012	0\\
-0.01	0\\
-0.008	0\\
-0.006	0\\
-0.004	0\\
-0.002	0\\
0	0\\
0.002	0.0391009\\
0.004	0.074092\\
0.006	0.104794\\
0.008	0.131117\\
0.01	0.153053\\
0.012	0.170674\\
0.014	0.184121\\
0.016	0.193599\\
0.018	0.199363\\
0.02	0.201716\\
0.022	0.200993\\
0.024	0.197554\\
0.026	0.191777\\
0.028	0.184044\\
0.03	0.174736\\
0.032	0.164228\\
0.034	0.152877\\
0.036	0.14102\\
0.038	0.128968\\
0.04	0.118721\\
0.042	0.124439\\
0.044	0.129422\\
0.046	0.133603\\
0.048	0.136933\\
0.05	0.139377\\
0.052	0.140916\\
0.054	0.141546\\
0.056	0.141279\\
0.058	0.14014\\
0.06	0.138166\\
0.062	0.135408\\
0.064	0.131924\\
0.066	0.127783\\
0.068	0.123062\\
0.07	0.117839\\
0.072	0.112201\\
0.074	0.106235\\
0.076	0.100027\\
0.078	0.0998835\\
0.08	0.102432\\
0.082	0.104833\\
0.084	0.107086\\
0.086	0.109189\\
0.088	0.111141\\
0.09	0.11294\\
0.092	0.114585\\
0.094	0.116076\\
0.096	0.11741\\
0.098	0.118586\\
0.1	0.119604\\
0.102	0.120461\\
0.104	0.121157\\
0.106	0.121691\\
0.108	0.122061\\
0.11	0.122265\\
0.112	0.122304\\
0.114	0.122175\\
0.116	0.121879\\
0.118	0.121416\\
0.12	0.122417\\
0.122	0.12375\\
0.124	0.124889\\
0.126	0.125836\\
0.128	0.126594\\
0.13	0.127168\\
0.132	0.127562\\
0.134	0.127781\\
0.136	0.127831\\
0.138	0.127717\\
0.14	0.127446\\
0.142	0.127023\\
0.144	0.126455\\
0.146	0.125749\\
0.148	0.124911\\
0.15	0.123948\\
0.152	0.122866\\
0.154	0.121672\\
0.156	0.120372\\
0.158	0.118973\\
0.16	0.11748\\
0.162	0.115899\\
0.164	0.114236\\
0.166	0.112496\\
0.168	0.110685\\
0.17	0.108808\\
0.172	0.10687\\
0.174	0.104874\\
0.176	0.102827\\
0.178	0.100731\\
0.18	0.0985923\\
0.182	0.0964134\\
0.184	0.0941985\\
0.186	0.0919512\\
0.188	0.0896751\\
0.19	0.0873735\\
0.192	0.0850497\\
0.194	0.0827068\\
0.196	0.0803479\\
0.198	0.0779759\\
0.2	0.0755936\\
0.202	0.0732038\\
0.204	0.0708092\\
0.206	0.0684123\\
0.208	0.0660156\\
0.21	0.0636216\\
0.212	0.0612327\\
0.214	0.058851\\
0.216	0.0564789\\
0.218	0.0541185\\
0.22	0.0517718\\
0.222	0.0494408\\
0.224	0.0471274\\
0.226	0.0448335\\
0.228	0.0425608\\
0.23	0.040311\\
0.232	0.0380857\\
0.234	0.0358865\\
0.236	0.0344272\\
0.238	0.0332559\\
0.24	0.0320742\\
0.242	0.0308843\\
0.244	0.0296879\\
0.246	0.0284874\\
0.248	0.0283227\\
0.25	0.0284892\\
0.252	0.0285506\\
0.254	0.0285095\\
0.256	0.0283685\\
0.258	0.0281304\\
0.26	0.027798\\
0.262	0.0273745\\
0.264	0.0268631\\
0.266	0.0262671\\
0.268	0.0255901\\
0.27	0.0248356\\
0.272	0.0240073\\
0.274	0.023109\\
0.276	0.0221445\\
0.278	0.021118\\
0.28	0.0200332\\
0.282	0.0188945\\
0.284	0.0177058\\
0.286	0.0164713\\
0.288	0.0151953\\
0.29	0.0146666\\
0.292	0.0143832\\
0.294	0.0141038\\
0.296	0.0142072\\
0.298	0.0151507\\
0.3	0.01605\\
0.302	0.0169051\\
0.304	0.0177159\\
0.306	0.0184825\\
0.308	0.0192051\\
0.31	0.0198836\\
0.312	0.0205184\\
0.314	0.0211095\\
0.316	0.0216571\\
0.318	0.0221614\\
0.32	0.0226228\\
0.322	0.0230415\\
0.324	0.0234178\\
0.326	0.023752\\
0.328	0.0240447\\
0.33	0.0242961\\
0.332	0.0245067\\
0.334	0.0246772\\
0.336	0.0248079\\
0.338	0.0248994\\
0.34	0.0249523\\
0.342	0.0249674\\
0.344	0.0249451\\
0.346	0.0248863\\
0.348	0.0247917\\
0.35	0.024662\\
0.352	0.024498\\
0.354	0.0248983\\
0.356	0.0254514\\
0.358	0.0259441\\
0.36	0.0263758\\
0.362	0.0267462\\
0.364	0.0270552\\
0.366	0.0273028\\
0.368	0.027489\\
0.37	0.027614\\
0.372	0.0276783\\
0.374	0.0276825\\
0.376	0.0276271\\
0.378	0.027513\\
0.38	0.027341\\
0.382	0.0271123\\
0.384	0.0268279\\
0.386	0.0264891\\
0.388	0.0260973\\
0.39	0.025654\\
0.392	0.0251608\\
0.394	0.0246192\\
0.396	0.0240311\\
0.398	0.0233983\\
0.4	0.0227228\\
0.402	0.0220064\\
0.404	0.0212513\\
0.406	0.0204596\\
0.408	0.0196334\\
0.41	0.018775\\
0.412	0.0178867\\
0.414	0.0169707\\
0.416	0.0160294\\
0.418	0.0150651\\
0.42	0.0140803\\
0.422	0.0130773\\
0.424	0.0120586\\
0.426	0.0110267\\
0.428	0.00998379\\
0.43	0.00893246\\
0.432	0.00787507\\
0.434	0.00681401\\
0.436	0.00575165\\
0.438	0.00469033\\
0.44	0.00363236\\
0.442	0.00296529\\
0.444	0.00302936\\
0.446	0.00341838\\
0.448	0.00403406\\
0.45	0.00463551\\
0.452	0.00522205\\
0.454	0.00579304\\
0.456	0.00634784\\
0.458	0.00688589\\
0.46	0.00740661\\
0.462	0.00790947\\
0.464	0.008394\\
0.466	0.00885972\\
0.468	0.00964685\\
0.47	0.010412\\
0.472	0.0111446\\
0.474	0.0118433\\
0.476	0.0125072\\
0.478	0.0131351\\
0.48	0.0137262\\
0.482	0.0142797\\
0.484	0.0147947\\
0.486	0.0152708\\
0.488	0.0157074\\
0.49	0.0161039\\
0.492	0.0164602\\
0.494	0.0167758\\
0.496	0.0170507\\
0.498	0.0172847\\
0.5	0.0174779\\
0.502	0.0176304\\
0.504	0.0177423\\
0.506	0.017814\\
0.508	0.0178457\\
0.51	0.017838\\
0.512	0.0177914\\
0.514	0.0177064\\
0.516	0.0175838\\
0.518	0.0174242\\
0.52	0.0172285\\
0.522	0.0169976\\
0.524	0.0167325\\
0.526	0.0164341\\
0.528	0.0161034\\
0.53	0.0157417\\
0.532	0.0153502\\
0.534	0.0149299\\
0.536	0.0144823\\
0.538	0.0140086\\
0.54	0.0135101\\
0.542	0.0129883\\
0.544	0.0124447\\
0.546	0.0118805\\
0.548	0.0112973\\
0.55	0.0106967\\
0.552	0.0100801\\
0.554	0.009449\\
0.556	0.00908967\\
0.558	0.00872923\\
0.56	0.00836112\\
0.562	0.007986\\
0.564	0.00760458\\
0.566	0.00721752\\
0.568	0.0068255\\
0.57	0.00642922\\
0.572	0.00602934\\
0.574	0.00562655\\
0.576	0.00522151\\
0.578	0.0048149\\
0.58	0.00440737\\
0.582	0.00399957\\
0.584	0.00359216\\
0.586	0.00318577\\
0.588	0.00278102\\
0.59	0.00263926\\
0.592	0.00325925\\
0.594	0.00386505\\
0.596	0.00445555\\
0.598	0.00502968\\
0.6	0.00558643\\
0.602	0.00612483\\
0.604	0.00664396\\
0.606	0.00714295\\
0.608	0.007621\\
0.61	0.00807735\\
0.612	0.00851131\\
0.614	0.00892224\\
0.616	0.00930954\\
0.618	0.0096727\\
0.62	0.0100113\\
0.622	0.0103248\\
0.624	0.0106129\\
0.626	0.0108754\\
0.628	0.011112\\
0.63	0.0113226\\
0.632	0.0115069\\
0.634	0.011665\\
0.636	0.0117969\\
0.638	0.0119025\\
0.64	0.0119821\\
0.642	0.0120357\\
0.644	0.0120636\\
0.646	0.0120661\\
0.648	0.0120435\\
0.65	0.0119961\\
0.652	0.0119244\\
0.654	0.0118287\\
0.656	0.0117097\\
0.658	0.0115679\\
0.66	0.0114039\\
0.662	0.0112182\\
0.664	0.0110116\\
0.666	0.0107848\\
0.668	0.0105385\\
0.67	0.0102735\\
0.672	0.00999054\\
0.674	0.0096905\\
0.676	0.00937422\\
0.678	0.00904257\\
0.68	0.00869646\\
0.682	0.0083368\\
0.684	0.00796452\\
0.686	0.00758057\\
0.688	0.00718592\\
0.69	0.00678153\\
0.692	0.00636838\\
0.694	0.00609592\\
0.696	0.00597135\\
0.698	0.00583879\\
0.7	0.00569859\\
0.702	0.00555113\\
0.704	0.00539676\\
0.706	0.00523586\\
0.708	0.00506883\\
0.71	0.00489603\\
0.712	0.00471787\\
0.714	0.00453473\\
0.716	0.00434701\\
0.718	0.0041551\\
0.72	0.00395941\\
0.722	0.00376033\\
0.724	0.00355826\\
0.726	0.00335359\\
0.728	0.00314672\\
0.73	0.00293803\\
0.732	0.00272793\\
0.734	0.00262504\\
0.736	0.00299339\\
0.738	0.00335044\\
0.74	0.00369557\\
0.742	0.00402824\\
0.744	0.00434791\\
0.746	0.00465408\\
0.748	0.00494631\\
0.75	0.00522415\\
0.752	0.00548724\\
0.754	0.00573521\\
0.756	0.00596776\\
0.758	0.0061846\\
0.76	0.00638551\\
0.762	0.00657028\\
0.764	0.00673875\\
0.766	0.00689079\\
0.768	0.0070263\\
0.77	0.00714525\\
0.772	0.0072476\\
0.774	0.00733338\\
0.776	0.00740263\\
0.778	0.00745546\\
0.78	0.00749198\\
0.782	0.00751234\\
0.784	0.00751673\\
0.786	0.00750537\\
0.788	0.0074785\\
0.79	0.00743642\\
0.792	0.00737941\\
0.794	0.00730783\\
0.796	0.00722202\\
0.798	0.00712237\\
0.8	0.0070093\\
0.802	0.00688324\\
0.804	0.00674463\\
0.806	0.00659396\\
0.808	0.00643172\\
0.81	0.00625842\\
0.812	0.00607458\\
0.814	0.00588075\\
0.816	0.00567748\\
0.818	0.00546534\\
0.82	0.00524491\\
0.822	0.00501678\\
0.824	0.00478153\\
0.826	0.00453978\\
0.828	0.00429213\\
0.83	0.00403918\\
0.832	0.00378156\\
0.834	0.00351987\\
0.836	0.0033322\\
0.838	0.00329538\\
0.84	0.00325332\\
0.842	0.00320618\\
0.844	0.00315413\\
0.846	0.00309734\\
0.848	0.003036\\
0.85	0.00297029\\
0.852	0.0029004\\
0.854	0.00282653\\
0.856	0.00274887\\
0.858	0.00266764\\
0.86	0.00258303\\
0.862	0.00249526\\
0.864	0.00240453\\
0.866	0.00231106\\
0.868	0.00221507\\
0.87	0.00211676\\
0.872	0.00201637\\
0.874	0.00191409\\
0.876	0.00188406\\
0.878	0.00209342\\
0.88	0.00229517\\
0.882	0.00248899\\
0.884	0.00267459\\
0.886	0.00285172\\
0.888	0.00302012\\
0.89	0.00317957\\
0.892	0.00332987\\
0.894	0.00347085\\
0.896	0.00360234\\
0.898	0.00372423\\
0.9	0.0038364\\
0.902	0.00393875\\
0.904	0.00403124\\
0.906	0.00411381\\
0.908	0.00418644\\
0.91	0.00424913\\
0.912	0.00430191\\
0.914	0.00434481\\
0.916	0.00437791\\
0.918	0.00440127\\
0.92	0.00441501\\
0.922	0.00441925\\
0.924	0.00441412\\
0.926	0.00439977\\
0.928	0.0043764\\
0.93	0.00434418\\
0.932	0.00430331\\
0.934	0.00425404\\
0.936	0.00419658\\
0.938	0.00413118\\
0.94	0.00405812\\
0.942	0.00397767\\
0.944	0.0038901\\
0.946	0.00379573\\
0.948	0.00369485\\
0.95	0.00358778\\
0.952	0.00347484\\
0.954	0.00335637\\
0.956	0.0032327\\
0.958	0.00310417\\
0.96	0.00297115\\
0.962	0.00283397\\
0.964	0.00269299\\
0.966	0.00254857\\
0.968	0.00240107\\
0.97	0.00225085\\
0.972	0.00209828\\
0.974	0.0019437\\
0.976	0.00178749\\
0.978	0.00171593\\
0.98	0.00170753\\
0.982	0.00169606\\
0.984	0.0016816\\
0.986	0.00166424\\
0.988	0.00164406\\
0.99	0.00162114\\
0.992	0.00159559\\
0.994	0.00156749\\
0.996	0.00153696\\
0.998	0.00150408\\
1	0.00146898\\
1.002	0.00143175\\
1.004	0.0013925\\
1.006	0.00135136\\
1.008	0.00130843\\
1.01	0.00126382\\
1.012	0.00121766\\
1.014	0.00117007\\
1.016	0.00117053\\
1.018	0.00128791\\
1.02	0.00140069\\
1.022	0.0015087\\
1.024	0.00161178\\
1.026	0.0017098\\
1.028	0.00180262\\
1.03	0.00189015\\
1.032	0.00197226\\
1.034	0.00204888\\
1.036	0.00211994\\
1.038	0.00218536\\
1.04	0.00224511\\
1.042	0.00229915\\
1.044	0.00234744\\
1.046	0.00238999\\
1.048	0.0024268\\
1.05	0.00245787\\
1.052	0.00248324\\
1.054	0.00250294\\
1.056	0.00251702\\
1.058	0.00252555\\
1.06	0.00252858\\
1.062	0.00252621\\
1.064	0.00251853\\
1.066	0.00250563\\
1.068	0.00248764\\
1.07	0.00246466\\
1.072	0.00243683\\
1.074	0.00240429\\
1.076	0.00236718\\
1.078	0.00232565\\
1.08	0.00227986\\
1.082	0.00222998\\
1.084	0.00217619\\
1.086	0.00211864\\
1.088	0.00205754\\
1.09	0.00199307\\
1.092	0.00192541\\
1.094	0.00185478\\
1.096	0.00178135\\
1.098	0.00170534\\
1.1	0.00162694\\
1.102	0.00154637\\
1.104	0.00146383\\
1.106	0.00137952\\
1.108	0.00129366\\
1.11	0.00120646\\
1.112	0.00111812\\
1.114	0.00102884\\
1.116	0.000938839\\
1.118	0.000848314\\
1.12	0.000822557\\
1.122	0.000821372\\
1.124	0.00081859\\
1.126	0.000814247\\
1.128	0.000808385\\
1.13	0.000801044\\
1.132	0.000792271\\
1.134	0.000782109\\
1.136	0.000770609\\
1.138	0.000757818\\
1.14	0.00074379\\
1.142	0.000728577\\
1.144	0.000712232\\
1.146	0.000694812\\
1.148	0.000676373\\
1.15	0.000656971\\
1.152	0.000636666\\
1.154	0.000658633\\
1.156	0.000725365\\
1.158	0.000789444\\
1.16	0.000850778\\
1.162	0.000909279\\
1.164	0.000964869\\
1.166	0.00101748\\
1.168	0.00106704\\
1.17	0.00111349\\
1.172	0.0011568\\
1.174	0.00119691\\
1.176	0.0012338\\
1.178	0.00126744\\
1.18	0.0012978\\
1.182	0.00132488\\
1.184	0.00134867\\
1.186	0.00136918\\
1.188	0.00138642\\
1.19	0.00140039\\
1.192	0.00141114\\
1.194	0.00141868\\
1.196	0.00142305\\
1.198	0.0014243\\
1.2	0.00142247\\
1.202	0.00141762\\
1.204	0.00140981\\
1.206	0.00139911\\
1.208	0.00138558\\
1.21	0.0013693\\
1.212	0.00135035\\
1.214	0.00132882\\
1.216	0.0013048\\
1.218	0.00127837\\
1.22	0.00124964\\
1.222	0.0012187\\
1.224	0.00118566\\
1.226	0.00115062\\
1.228	0.0011137\\
1.23	0.001075\\
1.232	0.00103463\\
1.234	0.000992716\\
1.236	0.000949366\\
1.238	0.000904699\\
1.24	0.000858833\\
1.242	0.000811887\\
1.244	0.00076398\\
1.246	0.000715232\\
1.248	0.000665761\\
1.25	0.000615686\\
1.252	0.000565124\\
1.254	0.000514194\\
1.256	0.000463012\\
1.258	0.00041169\\
1.26	0.000377613\\
1.262	0.00037827\\
1.264	0.000378128\\
1.266	0.000377203\\
1.268	0.000375514\\
1.27	0.00037308\\
1.272	0.000369921\\
1.274	0.000366059\\
1.276	0.000361516\\
1.278	0.000356316\\
1.28	0.000350483\\
1.282	0.000344042\\
1.284	0.00033702\\
1.286	0.000329443\\
1.288	0.000321338\\
1.29	0.00034323\\
1.292	0.000381662\\
1.294	0.000418622\\
1.296	0.000454053\\
1.298	0.000487905\\
1.3	0.000520131\\
1.302	0.000550689\\
1.304	0.00057954\\
1.306	0.000606651\\
1.308	0.000631992\\
1.31	0.000655538\\
1.312	0.000677267\\
1.314	0.000697165\\
1.316	0.000715217\\
1.318	0.000731416\\
1.32	0.000745758\\
1.322	0.000758243\\
1.324	0.000768876\\
1.326	0.000777663\\
1.328	0.000784618\\
1.33	0.000789756\\
1.332	0.000793095\\
1.334	0.00079466\\
1.336	0.000794477\\
1.338	0.000792574\\
1.34	0.000788986\\
1.342	0.000783749\\
1.344	0.0007769\\
1.346	0.000768484\\
1.348	0.000758543\\
1.35	0.000747125\\
1.352	0.000734279\\
1.354	0.000720059\\
1.356	0.000704516\\
1.358	0.000687708\\
1.36	0.000669692\\
1.362	0.000650528\\
1.364	0.000630276\\
1.366	0.000608998\\
1.368	0.000586759\\
1.37	0.000563622\\
1.372	0.000539652\\
1.374	0.000514917\\
1.376	0.000489481\\
1.378	0.000463413\\
1.38	0.00043678\\
1.382	0.000409648\\
1.384	0.000382086\\
1.386	0.000354159\\
1.388	0.000325936\\
1.39	0.000297482\\
1.392	0.000268863\\
1.394	0.000240144\\
1.396	0.000211389\\
1.398	0.000182661\\
1.4	0.000173099\\
1.402	0.000173234\\
1.404	0.000172982\\
1.406	0.00017235\\
1.408	0.000171349\\
1.41	0.000169988\\
1.412	0.000168277\\
1.414	0.000166229\\
1.416	0.000163854\\
1.418	0.000161164\\
1.42	0.000158172\\
1.422	0.000154891\\
1.424	0.000162128\\
1.426	0.000184666\\
1.428	0.000206421\\
1.43	0.000227358\\
1.432	0.000247445\\
1.434	0.000266653\\
1.436	0.000284955\\
1.438	0.000302327\\
1.44	0.000318745\\
1.442	0.000334191\\
1.444	0.000348647\\
1.446	0.000362098\\
1.448	0.000374533\\
1.45	0.00038594\\
1.452	0.000396312\\
1.454	0.000405644\\
1.456	0.000413934\\
1.458	0.000421181\\
1.46	0.000427385\\
1.462	0.000432553\\
1.464	0.000436689\\
1.466	0.000439802\\
1.468	0.000441902\\
1.47	0.000443002\\
1.472	0.000443116\\
1.474	0.000442261\\
1.476	0.000440454\\
1.478	0.000437716\\
1.48	0.000434068\\
1.482	0.000429534\\
1.484	0.000424138\\
1.486	0.000417906\\
1.488	0.000410867\\
1.49	0.000403048\\
1.492	0.000394481\\
1.494	0.000385196\\
1.496	0.000375226\\
1.498	0.000364604\\
1.5	0.000353364\\
1.502	0.000341541\\
1.504	0.000329171\\
1.506	0.000316289\\
1.508	0.000302934\\
1.51	0.000289141\\
1.512	0.000274948\\
1.514	0.000260394\\
1.516	0.000245515\\
1.518	0.000230349\\
1.52	0.000214936\\
1.522	0.000199312\\
1.524	0.000183515\\
1.526	0.000167583\\
1.528	0.000151552\\
1.53	0.00013546\\
1.532	0.000119342\\
1.534	0.000103235\\
1.536	8.71732e-05\\
1.538	8.22581e-05\\
1.54	8.20007e-05\\
1.542	8.15555e-05\\
1.544	8.0927e-05\\
1.546	8.01205e-05\\
1.548	7.91413e-05\\
1.55	7.79951e-05\\
1.552	7.6688e-05\\
1.554	7.5226e-05\\
1.556	7.36157e-05\\
1.558	7.73003e-05\\
1.56	9.03926e-05\\
1.562	0.000103067\\
1.564	0.000115302\\
1.566	0.000127078\\
1.568	0.000138379\\
1.57	0.000149186\\
1.572	0.000159484\\
1.574	0.00016926\\
1.576	0.000178501\\
1.578	0.000187195\\
1.58	0.000195333\\
1.582	0.000202906\\
1.584	0.000209907\\
1.586	0.000216331\\
1.588	0.000222172\\
1.59	0.000227428\\
1.592	0.000232098\\
1.594	0.00023618\\
1.596	0.000239676\\
1.598	0.000242588\\
1.6	0.000244919\\
1.602	0.000246674\\
1.604	0.000247859\\
1.606	0.00024848\\
1.608	0.000248547\\
1.61	0.000248067\\
1.612	0.000247051\\
1.614	0.000245511\\
1.616	0.000243458\\
1.618	0.000240905\\
1.62	0.000237868\\
1.622	0.000234359\\
1.624	0.000230396\\
1.626	0.000225993\\
1.628	0.000221169\\
1.63	0.00021594\\
1.632	0.000210326\\
1.634	0.000204344\\
1.636	0.000198014\\
1.638	0.000191356\\
1.64	0.00018439\\
1.642	0.000177136\\
1.644	0.000169615\\
1.646	0.000161848\\
1.648	0.000153856\\
1.65	0.00014566\\
1.652	0.000137281\\
1.654	0.000128742\\
1.656	0.000120063\\
1.658	0.000111265\\
1.66	0.000102371\\
1.662	9.34008e-05\\
1.664	8.43756e-05\\
1.666	7.53161e-05\\
1.668	6.62426e-05\\
1.67	5.71753e-05\\
1.672	4.8134e-05\\
1.674	4.27719e-05\\
1.676	4.23922e-05\\
1.678	4.19165e-05\\
1.68	4.13477e-05\\
1.682	4.06888e-05\\
1.684	3.99429e-05\\
1.686	3.91133e-05\\
1.688	3.82034e-05\\
1.69	3.72168e-05\\
1.692	3.68279e-05\\
1.694	4.44148e-05\\
1.696	5.1778e-05\\
1.698	5.89052e-05\\
1.7	6.57847e-05\\
1.702	7.24057e-05\\
1.704	7.87578e-05\\
1.706	8.48318e-05\\
1.708	9.06189e-05\\
1.71	9.61113e-05\\
1.712	0.000101302\\
1.714	0.000106184\\
1.716	0.000110753\\
1.718	0.000115003\\
1.72	0.000118931\\
1.722	0.000122533\\
1.724	0.000125807\\
1.726	0.000128751\\
1.728	0.000131364\\
1.73	0.000133646\\
1.732	0.000135598\\
1.734	0.00013722\\
1.736	0.000138515\\
1.738	0.000139485\\
1.74	0.000140134\\
1.742	0.000140466\\
1.744	0.000140484\\
1.746	0.000140194\\
1.748	0.000139602\\
1.75	0.000138715\\
1.752	0.000137538\\
1.754	0.00013608\\
1.756	0.000134348\\
1.758	0.000132351\\
1.76	0.000130097\\
1.762	0.000127596\\
1.764	0.000124857\\
1.766	0.000121889\\
1.768	0.000118704\\
1.77	0.000115312\\
1.772	0.000111723\\
1.774	0.00010795\\
1.776	0.000104002\\
1.778	9.98915e-05\\
1.78	9.56305e-05\\
1.782	9.12304e-05\\
1.784	8.67032e-05\\
1.786	8.20609e-05\\
1.788	7.73155e-05\\
1.79	7.24791e-05\\
1.792	6.75637e-05\\
1.794	6.25813e-05\\
1.796	5.7544e-05\\
1.798	5.24636e-05\\
1.8	4.73519e-05\\
1.802	4.22207e-05\\
1.804	3.70813e-05\\
1.806	3.19452e-05\\
1.808	2.68235e-05\\
1.81	2.44401e-05\\
1.812	2.40553e-05\\
1.814	2.36195e-05\\
1.816	2.31345e-05\\
1.818	2.26022e-05\\
1.82	2.20247e-05\\
1.822	2.14039e-05\\
1.824	2.0742e-05\\
1.826	2.00412e-05\\
1.828	2.13268e-05\\
1.83	2.5629e-05\\
1.832	2.98049e-05\\
1.834	3.38476e-05\\
1.836	3.77503e-05\\
1.838	4.15069e-05\\
1.84	4.51115e-05\\
1.842	4.85588e-05\\
1.844	5.18438e-05\\
1.846	5.49621e-05\\
1.848	5.79096e-05\\
1.85	6.06826e-05\\
1.852	6.32779e-05\\
1.854	6.56928e-05\\
1.856	6.7925e-05\\
1.858	6.99726e-05\\
1.86	7.18342e-05\\
1.862	7.35086e-05\\
1.864	7.49954e-05\\
1.866	7.62944e-05\\
1.868	7.74057e-05\\
1.87	7.833e-05\\
1.872	7.90684e-05\\
1.874	7.96222e-05\\
1.876	7.99933e-05\\
1.878	8.01838e-05\\
1.88	8.01962e-05\\
1.882	8.00335e-05\\
1.884	7.96988e-05\\
1.886	7.91956e-05\\
1.888	7.85278e-05\\
1.89	7.76994e-05\\
1.892	7.6715e-05\\
1.894	7.55791e-05\\
1.896	7.42968e-05\\
1.898	7.28731e-05\\
1.9	7.13135e-05\\
1.902	6.96236e-05\\
1.904	6.78091e-05\\
1.906	6.5876e-05\\
1.908	6.38305e-05\\
1.91	6.16788e-05\\
1.912	5.94273e-05\\
1.914	5.70826e-05\\
1.916	5.46512e-05\\
1.918	5.21399e-05\\
1.92	4.95553e-05\\
1.922	4.69044e-05\\
1.924	4.4194e-05\\
1.926	4.14308e-05\\
1.928	3.86219e-05\\
1.93	3.57739e-05\\
1.932	3.28939e-05\\
1.934	2.99885e-05\\
1.936	2.70645e-05\\
1.938	2.41285e-05\\
1.94	2.11872e-05\\
1.942	1.8247e-05\\
1.944	1.59251e-05\\
1.946	1.6436e-05\\
1.948	1.68895e-05\\
1.95	1.72847e-05\\
1.952	1.76209e-05\\
1.954	1.78977e-05\\
1.956	1.81148e-05\\
1.958	1.82721e-05\\
1.96	1.83699e-05\\
1.962	1.84083e-05\\
1.964	1.83881e-05\\
1.966	1.83098e-05\\
1.968	1.81744e-05\\
1.97	1.94876e-05\\
1.972	2.17309e-05\\
1.974	2.38908e-05\\
1.976	2.5964e-05\\
1.978	2.79474e-05\\
1.98	2.98381e-05\\
1.982	3.16334e-05\\
1.984	3.33311e-05\\
1.986	3.49289e-05\\
1.988	3.6425e-05\\
1.99	3.78177e-05\\
1.992	3.91057e-05\\
1.994	4.02879e-05\\
1.996	4.13633e-05\\
1.998	4.23314e-05\\
2	4.31918e-05\\
};
\addplot [color=mycolor2, line width=1.0pt, forget plot]
  table[row sep=crcr]{%
-0.024	0\\
-0.022	0\\
-0.02	0\\
-0.018	0\\
-0.016	0\\
-0.014	0\\
-0.012	0\\
-0.01	0\\
-0.008	0\\
-0.006	0\\
-0.004	0\\
-0.002	0\\
0	0\\
0.002	0.0391009\\
0.004	0.074092\\
0.006	0.104794\\
0.008	0.131117\\
0.01	0.153053\\
0.012	0.170674\\
0.014	0.18412\\
0.016	0.193597\\
0.018	0.19936\\
0.02	0.201712\\
0.022	0.200987\\
0.024	0.197546\\
0.026	0.191766\\
0.028	0.184029\\
0.03	0.174717\\
0.032	0.164203\\
0.034	0.152845\\
0.036	0.140979\\
0.038	0.128917\\
0.04	0.118709\\
0.042	0.124424\\
0.044	0.129403\\
0.046	0.133579\\
0.048	0.136904\\
0.05	0.139342\\
0.052	0.140874\\
0.054	0.141496\\
0.056	0.14122\\
0.058	0.14007\\
0.06	0.138084\\
0.062	0.135312\\
0.064	0.131813\\
0.066	0.127656\\
0.068	0.122915\\
0.07	0.117672\\
0.072	0.112012\\
0.074	0.106021\\
0.076	0.0997859\\
0.078	0.093394\\
0.08	0.0869295\\
0.082	0.0886518\\
0.084	0.090368\\
0.086	0.0919601\\
0.088	0.0934289\\
0.09	0.0947753\\
0.092	0.096\\
0.094	0.0971036\\
0.096	0.0980866\\
0.098	0.0989494\\
0.1	0.0996922\\
0.102	0.101544\\
0.104	0.104574\\
0.106	0.10743\\
0.108	0.110102\\
0.11	0.112587\\
0.112	0.11488\\
0.114	0.116977\\
0.116	0.118877\\
0.118	0.120578\\
0.12	0.122081\\
0.122	0.123385\\
0.124	0.124494\\
0.126	0.125409\\
0.128	0.126133\\
0.13	0.126672\\
0.132	0.127029\\
0.134	0.127209\\
0.136	0.127218\\
0.138	0.127061\\
0.14	0.126745\\
0.142	0.126276\\
0.144	0.12566\\
0.146	0.124904\\
0.148	0.124015\\
0.15	0.122999\\
0.152	0.121862\\
0.154	0.120612\\
0.156	0.119254\\
0.158	0.117796\\
0.16	0.116242\\
0.162	0.114599\\
0.164	0.112873\\
0.166	0.111069\\
0.168	0.109193\\
0.17	0.10725\\
0.172	0.105244\\
0.174	0.103182\\
0.176	0.101066\\
0.178	0.0989028\\
0.18	0.0966954\\
0.182	0.094448\\
0.184	0.0921649\\
0.186	0.0898497\\
0.188	0.0875061\\
0.19	0.0851377\\
0.192	0.0827479\\
0.194	0.08034\\
0.196	0.0779174\\
0.198	0.075483\\
0.2	0.0730399\\
0.202	0.070591\\
0.204	0.0681393\\
0.206	0.0656875\\
0.208	0.0632382\\
0.21	0.0607942\\
0.212	0.0583579\\
0.214	0.0559319\\
0.216	0.0535184\\
0.218	0.05112\\
0.22	0.0487387\\
0.222	0.0463767\\
0.224	0.0440362\\
0.226	0.0417191\\
0.228	0.0394274\\
0.23	0.0371628\\
0.232	0.0358801\\
0.234	0.0348403\\
0.236	0.033788\\
0.238	0.0327245\\
0.24	0.0316515\\
0.242	0.0305705\\
0.244	0.0294835\\
0.246	0.0283922\\
0.248	0.0272988\\
0.25	0.0262053\\
0.252	0.0251137\\
0.254	0.0240264\\
0.256	0.0229455\\
0.258	0.0218732\\
0.26	0.0208117\\
0.262	0.0197631\\
0.264	0.0187294\\
0.266	0.0177127\\
0.268	0.017246\\
0.27	0.0169437\\
0.272	0.0166468\\
0.274	0.0163551\\
0.276	0.0160687\\
0.278	0.0157874\\
0.28	0.0155113\\
0.282	0.0152403\\
0.284	0.0149743\\
0.286	0.0147132\\
0.288	0.014457\\
0.29	0.0142055\\
0.292	0.0139587\\
0.294	0.0141852\\
0.296	0.0150279\\
0.298	0.0158232\\
0.3	0.0165714\\
0.302	0.0172726\\
0.304	0.0179272\\
0.306	0.0185355\\
0.308	0.0190978\\
0.31	0.0196145\\
0.312	0.020086\\
0.314	0.0205128\\
0.316	0.0208953\\
0.318	0.021234\\
0.32	0.0215296\\
0.322	0.0217824\\
0.324	0.0219933\\
0.326	0.0221628\\
0.328	0.0222916\\
0.33	0.0223804\\
0.332	0.0224301\\
0.334	0.0224413\\
0.336	0.0224148\\
0.338	0.0223517\\
0.34	0.0222526\\
0.342	0.0221186\\
0.344	0.0219506\\
0.346	0.0217496\\
0.348	0.0215164\\
0.35	0.0212523\\
0.352	0.0209581\\
0.354	0.0206351\\
0.356	0.0202842\\
0.358	0.0199066\\
0.36	0.0195033\\
0.362	0.0190756\\
0.364	0.0186246\\
0.366	0.0181514\\
0.368	0.0176572\\
0.37	0.0171431\\
0.372	0.0166104\\
0.374	0.0160603\\
0.376	0.0154939\\
0.378	0.0149124\\
0.38	0.014317\\
0.382	0.0137088\\
0.384	0.0130892\\
0.386	0.0124591\\
0.388	0.0118199\\
0.39	0.0111725\\
0.392	0.0105183\\
0.394	0.00985829\\
0.396	0.0091936\\
0.398	0.00857876\\
0.4	0.00813755\\
0.402	0.00768684\\
0.404	0.00722778\\
0.406	0.00676148\\
0.408	0.0062891\\
0.41	0.00581176\\
0.412	0.00533061\\
0.414	0.00484679\\
0.416	0.0043614\\
0.418	0.00388018\\
0.42	0.0038779\\
0.422	0.00386622\\
0.424	0.00384519\\
0.426	0.00381488\\
0.428	0.00380734\\
0.43	0.00395127\\
0.432	0.00408572\\
0.434	0.00421036\\
0.436	0.0043249\\
0.438	0.0044291\\
0.44	0.00459094\\
0.442	0.00510116\\
0.444	0.00559583\\
0.446	0.00607445\\
0.448	0.00653657\\
0.45	0.00698176\\
0.452	0.00740962\\
0.454	0.0078198\\
0.456	0.00821196\\
0.458	0.0085858\\
0.46	0.00894106\\
0.462	0.00927752\\
0.464	0.00959496\\
0.466	0.00989323\\
0.468	0.0101722\\
0.47	0.0104317\\
0.472	0.0106718\\
0.474	0.0108923\\
0.476	0.0110933\\
0.478	0.0112747\\
0.48	0.0114367\\
0.482	0.0115793\\
0.484	0.0117026\\
0.486	0.0118067\\
0.488	0.0118919\\
0.49	0.0119583\\
0.492	0.012006\\
0.494	0.0120355\\
0.496	0.0120468\\
0.498	0.0120404\\
0.5	0.0120165\\
0.502	0.0119754\\
0.504	0.0119176\\
0.506	0.0118433\\
0.508	0.011753\\
0.51	0.0116471\\
0.512	0.0115261\\
0.514	0.0113902\\
0.516	0.0112402\\
0.518	0.0110763\\
0.52	0.0108991\\
0.522	0.0107092\\
0.524	0.0105069\\
0.526	0.0102929\\
0.528	0.0100676\\
0.53	0.00983169\\
0.532	0.00958562\\
0.534	0.00932997\\
0.536	0.00906532\\
0.538	0.00879221\\
0.54	0.00851123\\
0.542	0.00822295\\
0.544	0.00792793\\
0.546	0.00762677\\
0.548	0.00732003\\
0.55	0.00700829\\
0.552	0.00669212\\
0.554	0.00637209\\
0.556	0.00604876\\
0.558	0.0057227\\
0.56	0.00539446\\
0.562	0.0050646\\
0.564	0.00473365\\
0.566	0.00440215\\
0.568	0.00407062\\
0.57	0.00373958\\
0.572	0.00340954\\
0.574	0.003081\\
0.576	0.00275443\\
0.578	0.00243032\\
0.58	0.00210912\\
0.582	0.0020538\\
0.584	0.00224256\\
0.586	0.00242355\\
0.588	0.00259651\\
0.59	0.00276121\\
0.592	0.00291747\\
0.594	0.00306508\\
0.596	0.00320391\\
0.598	0.0033338\\
0.6	0.00345467\\
0.602	0.0035664\\
0.604	0.00366894\\
0.606	0.00376225\\
0.608	0.0038463\\
0.61	0.00392109\\
0.612	0.00398664\\
0.614	0.00404298\\
0.616	0.00409018\\
0.618	0.00412832\\
0.62	0.00415748\\
0.622	0.00417779\\
0.624	0.00418937\\
0.626	0.00419237\\
0.628	0.00418696\\
0.63	0.00417331\\
0.632	0.00415162\\
0.634	0.00418851\\
0.636	0.00430335\\
0.638	0.00440877\\
0.64	0.00450481\\
0.642	0.00459151\\
0.644	0.00466892\\
0.646	0.00473712\\
0.648	0.00479617\\
0.65	0.00484619\\
0.652	0.00488727\\
0.654	0.00491954\\
0.656	0.00494314\\
0.658	0.0049582\\
0.66	0.00496488\\
0.662	0.00496336\\
0.664	0.0049538\\
0.666	0.0049364\\
0.668	0.00491135\\
0.67	0.00487885\\
0.672	0.00483911\\
0.674	0.00479237\\
0.676	0.00473883\\
0.678	0.00467874\\
0.68	0.00461234\\
0.682	0.00453986\\
0.684	0.00446157\\
0.686	0.0043777\\
0.688	0.00428852\\
0.69	0.00419429\\
0.692	0.00409527\\
0.694	0.00399173\\
0.696	0.00388394\\
0.698	0.00377216\\
0.7	0.00365665\\
0.702	0.0035377\\
0.704	0.00341557\\
0.706	0.00329052\\
0.708	0.00316283\\
0.71	0.00303276\\
0.712	0.00290057\\
0.714	0.00276652\\
0.716	0.00263088\\
0.718	0.00249389\\
0.72	0.00235581\\
0.722	0.00221689\\
0.724	0.00207737\\
0.726	0.00193748\\
0.728	0.00179747\\
0.73	0.00165756\\
0.732	0.00151798\\
0.734	0.00137894\\
0.736	0.00127772\\
0.738	0.00134382\\
0.74	0.00140559\\
0.742	0.00146302\\
0.744	0.00151607\\
0.746	0.00156476\\
0.748	0.00160907\\
0.75	0.00164903\\
0.752	0.00168465\\
0.754	0.00171597\\
0.756	0.00174302\\
0.758	0.00176586\\
0.76	0.00178454\\
0.762	0.00179912\\
0.764	0.00180967\\
0.766	0.00181628\\
0.768	0.00181902\\
0.77	0.001818\\
0.772	0.0018133\\
0.774	0.00180503\\
0.776	0.0017933\\
0.778	0.00177821\\
0.78	0.0017599\\
0.782	0.00173848\\
0.784	0.00171408\\
0.786	0.00168682\\
0.788	0.00165683\\
0.79	0.00162427\\
0.792	0.0016763\\
0.794	0.00173128\\
0.796	0.00178238\\
0.798	0.00182961\\
0.8	0.00187297\\
0.802	0.00191247\\
0.804	0.00194814\\
0.806	0.00198\\
0.808	0.00200809\\
0.81	0.00203244\\
0.812	0.00205309\\
0.814	0.00207009\\
0.816	0.00208348\\
0.818	0.00209333\\
0.82	0.0020997\\
0.822	0.00210264\\
0.824	0.00210222\\
0.826	0.00209851\\
0.828	0.0020916\\
0.83	0.00208155\\
0.832	0.00206845\\
0.834	0.00205238\\
0.836	0.00203342\\
0.838	0.00201167\\
0.84	0.00198722\\
0.842	0.00196016\\
0.844	0.00193059\\
0.846	0.0018986\\
0.848	0.00186429\\
0.85	0.00182776\\
0.852	0.00178912\\
0.854	0.00174847\\
0.856	0.0017059\\
0.858	0.00166153\\
0.86	0.00161547\\
0.862	0.00156781\\
0.864	0.00151866\\
0.866	0.00146813\\
0.868	0.00141633\\
0.87	0.00136335\\
0.872	0.00130931\\
0.874	0.00125431\\
0.876	0.00119846\\
0.878	0.00114186\\
0.88	0.00108461\\
0.882	0.00102681\\
0.884	0.000968561\\
0.886	0.000909969\\
0.888	0.000851129\\
0.89	0.000792136\\
0.892	0.000733085\\
0.894	0.00067407\\
0.896	0.000615181\\
0.898	0.000587697\\
0.9	0.000597483\\
0.902	0.000605771\\
0.904	0.000612594\\
0.906	0.000617986\\
0.908	0.000621983\\
0.91	0.000624625\\
0.912	0.00062595\\
0.914	0.000626002\\
0.916	0.000624823\\
0.918	0.000622458\\
0.92	0.000618953\\
0.922	0.000614354\\
0.924	0.000608709\\
0.926	0.000602066\\
0.928	0.000594474\\
0.93	0.000585982\\
0.932	0.000576639\\
0.934	0.000566494\\
0.936	0.000555598\\
0.938	0.000543999\\
0.94	0.000531747\\
0.942	0.000528871\\
0.944	0.000563556\\
0.946	0.000596648\\
0.948	0.000628127\\
0.95	0.000657976\\
0.952	0.000686177\\
0.954	0.000712721\\
0.956	0.000737596\\
0.958	0.000760796\\
0.96	0.000782319\\
0.962	0.000802162\\
0.964	0.000820327\\
0.966	0.000836819\\
0.968	0.000851646\\
0.97	0.000864815\\
0.972	0.00087634\\
0.974	0.000886235\\
0.976	0.000894517\\
0.978	0.000901205\\
0.98	0.00090632\\
0.982	0.000909886\\
0.984	0.000911929\\
0.986	0.000912475\\
0.988	0.000911555\\
0.99	0.000909199\\
0.992	0.00090544\\
0.994	0.000900314\\
0.996	0.000893856\\
0.998	0.000886104\\
1	0.000877097\\
1.002	0.000866874\\
1.004	0.000855479\\
1.006	0.000842952\\
1.008	0.000829338\\
1.01	0.000814681\\
1.012	0.000799026\\
1.014	0.00078242\\
1.016	0.000764908\\
1.018	0.000746538\\
1.02	0.000727357\\
1.022	0.000707414\\
1.024	0.000686756\\
1.026	0.000665433\\
1.028	0.000643491\\
1.03	0.000620981\\
1.032	0.00059795\\
1.034	0.000574446\\
1.036	0.000550517\\
1.038	0.000526212\\
1.04	0.000501577\\
1.042	0.000476659\\
1.044	0.000451504\\
1.046	0.000426159\\
1.048	0.000400668\\
1.05	0.000375076\\
1.052	0.000349426\\
1.054	0.00032376\\
1.056	0.000298122\\
1.058	0.00027255\\
1.06	0.000247086\\
1.062	0.000221768\\
1.064	0.000196634\\
1.066	0.00017172\\
1.068	0.000158534\\
1.07	0.000159223\\
1.072	0.000159662\\
1.074	0.000159861\\
1.076	0.00015983\\
1.078	0.000159579\\
1.08	0.000159118\\
1.082	0.000158455\\
1.084	0.000157601\\
1.086	0.000156565\\
1.088	0.000155356\\
1.09	0.000153982\\
1.092	0.000152454\\
1.094	0.00015078\\
1.096	0.000154428\\
1.098	0.000171863\\
1.1	0.000188651\\
1.102	0.000204781\\
1.104	0.00022024\\
1.106	0.000235019\\
1.108	0.000249108\\
1.11	0.0002625\\
1.112	0.00027519\\
1.114	0.000287171\\
1.116	0.000298441\\
1.118	0.000308997\\
1.12	0.000318837\\
1.122	0.000327962\\
1.124	0.000336372\\
1.126	0.000344071\\
1.128	0.00035106\\
1.13	0.000357345\\
1.132	0.00036293\\
1.134	0.000367823\\
1.136	0.00037203\\
1.138	0.000375559\\
1.14	0.000378421\\
1.142	0.000380624\\
1.144	0.00038218\\
1.146	0.000383099\\
1.148	0.000383396\\
1.15	0.000383082\\
1.152	0.000382171\\
1.154	0.000380677\\
1.156	0.000378616\\
1.158	0.000376004\\
1.16	0.000372855\\
1.162	0.000369187\\
1.164	0.000365017\\
1.166	0.000360361\\
1.168	0.000355238\\
1.17	0.000349667\\
1.172	0.000343664\\
1.174	0.00033725\\
1.176	0.000330442\\
1.178	0.000323261\\
1.18	0.000315725\\
1.182	0.000307854\\
1.184	0.000299667\\
1.186	0.000291184\\
1.188	0.000282424\\
1.19	0.000273408\\
1.192	0.000264155\\
1.194	0.000254684\\
1.196	0.000245014\\
1.198	0.000235166\\
1.2	0.000225158\\
1.202	0.00021501\\
1.204	0.00020474\\
1.206	0.000194368\\
1.208	0.00018391\\
1.21	0.000173386\\
1.212	0.000162814\\
1.214	0.000152211\\
1.216	0.000141594\\
1.218	0.00013098\\
1.22	0.000120385\\
1.222	0.000109827\\
1.224	9.932e-05\\
1.226	8.888e-05\\
1.228	7.85218e-05\\
1.23	6.82599e-05\\
1.232	5.81083e-05\\
1.234	4.80805e-05\\
1.236	4.45351e-05\\
1.238	4.27619e-05\\
1.24	4.09626e-05\\
1.242	4.0237e-05\\
1.244	4.17544e-05\\
1.246	4.3218e-05\\
1.248	4.46266e-05\\
1.25	4.59792e-05\\
1.252	4.72744e-05\\
1.254	4.85113e-05\\
1.256	5.06769e-05\\
1.258	5.83451e-05\\
1.26	6.57598e-05\\
1.262	7.29147e-05\\
1.264	7.98041e-05\\
1.266	8.64228e-05\\
1.268	9.27662e-05\\
1.27	9.88301e-05\\
1.272	0.000104611\\
1.274	0.000110106\\
1.276	0.000115312\\
1.278	0.000120227\\
1.28	0.00012485\\
1.282	0.00012918\\
1.284	0.000133216\\
1.286	0.000136958\\
1.288	0.000140406\\
1.29	0.000143561\\
1.292	0.000146425\\
1.294	0.000148999\\
1.296	0.000151284\\
1.298	0.000153285\\
1.3	0.000155004\\
1.302	0.000156443\\
1.304	0.000157607\\
1.306	0.000158499\\
1.308	0.000159125\\
1.31	0.000159488\\
1.312	0.000159594\\
1.314	0.000159448\\
1.316	0.000159055\\
1.318	0.000158422\\
1.32	0.000157555\\
1.322	0.000156459\\
1.324	0.000155142\\
1.326	0.00015361\\
1.328	0.000151871\\
1.33	0.000149931\\
1.332	0.000147797\\
1.334	0.000145478\\
1.336	0.000142981\\
1.338	0.000140313\\
1.34	0.000137482\\
1.342	0.000134497\\
1.344	0.000131365\\
1.346	0.000128094\\
1.348	0.000124692\\
1.35	0.000121168\\
1.352	0.000117529\\
1.354	0.000113784\\
1.356	0.000109941\\
1.358	0.000106007\\
1.36	0.000101992\\
1.362	9.79023e-05\\
1.364	9.37467e-05\\
1.366	8.95329e-05\\
1.368	8.52687e-05\\
1.37	8.09618e-05\\
1.372	7.66198e-05\\
1.374	7.22503e-05\\
1.376	6.78606e-05\\
1.378	6.34579e-05\\
1.38	5.90495e-05\\
1.382	5.46422e-05\\
1.384	5.02428e-05\\
1.386	4.58581e-05\\
1.388	4.14945e-05\\
1.39	3.71583e-05\\
1.392	3.28556e-05\\
1.394	2.85925e-05\\
1.396	2.43746e-05\\
1.398	2.04316e-05\\
1.4	2.01638e-05\\
1.402	2.00751e-05\\
1.404	2.15502e-05\\
1.406	2.29712e-05\\
1.408	2.4336e-05\\
1.41	2.56426e-05\\
1.412	2.68891e-05\\
1.414	2.80739e-05\\
1.416	2.91953e-05\\
1.418	3.02518e-05\\
1.42	3.12422e-05\\
1.422	3.21652e-05\\
1.424	3.30199e-05\\
1.426	3.38054e-05\\
1.428	3.4521e-05\\
1.43	3.58222e-05\\
1.432	3.84777e-05\\
1.434	4.10185e-05\\
1.436	4.3443e-05\\
1.438	4.57501e-05\\
1.44	4.79385e-05\\
1.442	5.00076e-05\\
1.444	5.19564e-05\\
1.446	5.37846e-05\\
1.448	5.54919e-05\\
1.45	5.70781e-05\\
1.452	5.85434e-05\\
1.454	5.98879e-05\\
1.456	6.11122e-05\\
1.458	6.22168e-05\\
1.46	6.32025e-05\\
1.462	6.40702e-05\\
1.464	6.48211e-05\\
1.466	6.54563e-05\\
1.468	6.59773e-05\\
1.47	6.63856e-05\\
1.472	6.66829e-05\\
1.474	6.68709e-05\\
1.476	6.69517e-05\\
1.478	6.69273e-05\\
1.48	6.67999e-05\\
1.482	6.65717e-05\\
1.484	6.62451e-05\\
1.486	6.58226e-05\\
1.488	6.53068e-05\\
1.49	6.47005e-05\\
1.492	6.40062e-05\\
1.494	6.32269e-05\\
1.496	6.23655e-05\\
1.498	6.14249e-05\\
1.5	6.04082e-05\\
1.502	5.93184e-05\\
1.504	5.81587e-05\\
1.506	5.69323e-05\\
1.508	5.56423e-05\\
1.51	5.42921e-05\\
1.512	5.28849e-05\\
1.514	5.1424e-05\\
1.516	4.99127e-05\\
1.518	4.83544e-05\\
1.52	4.67524e-05\\
1.522	4.511e-05\\
1.524	4.34306e-05\\
1.526	4.17174e-05\\
1.528	3.99739e-05\\
1.53	3.82032e-05\\
1.532	3.64088e-05\\
1.534	3.45938e-05\\
1.536	3.27614e-05\\
1.538	3.09149e-05\\
1.54	2.90573e-05\\
1.542	2.71918e-05\\
1.544	2.53215e-05\\
1.546	2.34493e-05\\
1.548	2.15782e-05\\
1.55	1.9711e-05\\
1.552	1.78507e-05\\
1.554	1.6e-05\\
1.556	1.41616e-05\\
1.558	1.2338e-05\\
1.56	1.27459e-05\\
1.562	1.33918e-05\\
1.564	1.39997e-05\\
1.566	1.45692e-05\\
1.568	1.51001e-05\\
1.57	1.55921e-05\\
1.572	1.60452e-05\\
1.574	1.64593e-05\\
1.576	1.68344e-05\\
1.578	1.71707e-05\\
1.58	1.74684e-05\\
1.582	1.77278e-05\\
1.584	1.79492e-05\\
1.586	1.8133e-05\\
1.588	1.82798e-05\\
1.59	1.839e-05\\
1.592	1.84643e-05\\
1.594	1.85032e-05\\
1.596	1.85077e-05\\
1.598	1.84783e-05\\
1.6	1.87174e-05\\
1.602	1.97162e-05\\
1.604	2.06633e-05\\
1.606	2.15584e-05\\
1.608	2.24009e-05\\
1.61	2.31907e-05\\
1.612	2.39276e-05\\
1.614	2.46113e-05\\
1.616	2.5242e-05\\
1.618	2.58196e-05\\
1.62	2.63444e-05\\
1.622	2.68165e-05\\
1.624	2.72363e-05\\
1.626	2.76041e-05\\
1.628	2.79205e-05\\
1.63	2.8186e-05\\
1.632	2.84011e-05\\
1.634	2.85666e-05\\
1.636	2.86832e-05\\
1.638	2.87518e-05\\
1.64	2.87731e-05\\
1.642	2.87482e-05\\
1.644	2.86781e-05\\
1.646	2.85637e-05\\
1.648	2.84062e-05\\
1.65	2.82068e-05\\
1.652	2.79666e-05\\
1.654	2.76868e-05\\
1.656	2.73687e-05\\
1.658	2.70137e-05\\
1.66	2.66231e-05\\
1.662	2.61982e-05\\
1.664	2.57404e-05\\
1.666	2.52513e-05\\
1.668	2.47321e-05\\
1.67	2.41845e-05\\
1.672	2.36098e-05\\
1.674	2.30095e-05\\
1.676	2.23852e-05\\
1.678	2.17385e-05\\
1.68	2.10706e-05\\
1.682	2.03834e-05\\
1.684	1.96781e-05\\
1.686	1.89564e-05\\
1.688	1.82197e-05\\
1.69	1.74697e-05\\
1.692	1.67076e-05\\
1.694	1.59351e-05\\
1.696	1.51536e-05\\
1.698	1.43645e-05\\
1.7	1.35693e-05\\
1.702	1.27694e-05\\
1.704	1.19661e-05\\
1.706	1.11609e-05\\
1.708	1.0355e-05\\
1.71	9.54981e-06\\
1.712	8.74653e-06\\
1.714	7.94643e-06\\
1.716	7.15071e-06\\
1.718	6.36057e-06\\
1.72	6.02331e-06\\
1.722	6.22888e-06\\
1.724	6.41865e-06\\
1.726	6.59268e-06\\
1.728	6.75104e-06\\
1.73	6.89385e-06\\
1.732	7.02122e-06\\
1.734	7.13333e-06\\
1.736	7.23035e-06\\
1.738	7.31248e-06\\
1.74	7.37995e-06\\
1.742	7.433e-06\\
1.744	7.47189e-06\\
1.746	7.49692e-06\\
1.748	7.50838e-06\\
1.75	7.50659e-06\\
1.752	7.49188e-06\\
1.754	7.46461e-06\\
1.756	7.42513e-06\\
1.758	7.37382e-06\\
1.76	7.31107e-06\\
1.762	7.7165e-06\\
1.764	8.14346e-06\\
1.766	8.54854e-06\\
1.768	8.93159e-06\\
1.77	9.29248e-06\\
1.772	9.63112e-06\\
1.774	9.94746e-06\\
1.776	1.02415e-05\\
1.778	1.05132e-05\\
1.78	1.07627e-05\\
1.782	1.099e-05\\
1.784	1.11954e-05\\
1.786	1.13788e-05\\
1.788	1.15405e-05\\
1.79	1.16808e-05\\
1.792	1.17998e-05\\
1.794	1.18978e-05\\
1.796	1.19752e-05\\
1.798	1.20322e-05\\
1.8	1.20693e-05\\
1.802	1.20866e-05\\
1.804	1.20848e-05\\
1.806	1.20641e-05\\
1.808	1.2025e-05\\
1.81	1.19679e-05\\
1.812	1.18934e-05\\
1.814	1.18019e-05\\
1.816	1.16939e-05\\
1.818	1.157e-05\\
1.82	1.14306e-05\\
1.822	1.12762e-05\\
1.824	1.11076e-05\\
1.826	1.09251e-05\\
1.828	1.07295e-05\\
1.83	1.05212e-05\\
1.832	1.03009e-05\\
1.834	1.00691e-05\\
1.836	9.82645e-06\\
1.838	9.57356e-06\\
1.84	9.31102e-06\\
1.842	9.03945e-06\\
1.844	8.75944e-06\\
1.846	8.47161e-06\\
1.848	8.17656e-06\\
1.85	7.8749e-06\\
1.852	7.56723e-06\\
1.854	7.25415e-06\\
1.856	6.93626e-06\\
1.858	6.61414e-06\\
1.86	6.28839e-06\\
1.862	5.95958e-06\\
1.864	5.62828e-06\\
1.866	5.29505e-06\\
1.868	4.96045e-06\\
1.87	4.62503e-06\\
1.872	4.28932e-06\\
1.874	3.95384e-06\\
1.876	3.61912e-06\\
1.878	3.28564e-06\\
1.88	2.95391e-06\\
1.882	2.6244e-06\\
1.884	2.36601e-06\\
1.886	2.42458e-06\\
1.888	2.47769e-06\\
1.89	2.52541e-06\\
1.892	2.5678e-06\\
1.894	2.60494e-06\\
1.896	2.63693e-06\\
1.898	2.66384e-06\\
1.9	2.68577e-06\\
1.902	2.70283e-06\\
1.904	2.71512e-06\\
1.906	2.72276e-06\\
1.908	2.72586e-06\\
1.91	2.72455e-06\\
1.912	2.71894e-06\\
1.914	2.70918e-06\\
1.916	2.69538e-06\\
1.918	2.67769e-06\\
1.92	2.68267e-06\\
1.922	2.88981e-06\\
1.924	3.08817e-06\\
1.926	3.2776e-06\\
1.928	3.458e-06\\
1.93	3.62926e-06\\
1.932	3.7913e-06\\
1.934	3.94405e-06\\
1.936	4.08746e-06\\
1.938	4.22149e-06\\
1.94	4.34613e-06\\
1.942	4.46136e-06\\
1.944	4.5672e-06\\
1.946	4.66366e-06\\
1.948	4.7508e-06\\
1.95	4.82865e-06\\
1.952	4.89728e-06\\
1.954	4.95678e-06\\
1.956	5.00723e-06\\
1.958	5.04875e-06\\
1.96	5.08143e-06\\
1.962	5.10542e-06\\
1.964	5.12085e-06\\
1.966	5.12787e-06\\
1.968	5.12664e-06\\
1.97	5.11733e-06\\
1.972	5.10013e-06\\
1.974	5.07521e-06\\
1.976	5.04278e-06\\
1.978	5.00305e-06\\
1.98	4.95621e-06\\
1.982	4.90251e-06\\
1.984	4.84215e-06\\
1.986	4.77538e-06\\
1.988	4.70242e-06\\
1.99	4.62354e-06\\
1.992	4.53896e-06\\
1.994	4.44895e-06\\
1.996	4.35376e-06\\
1.998	4.25364e-06\\
2	4.14887e-06\\
};
\addplot [color=mycolor3, line width=1.0pt, forget plot]
  table[row sep=crcr]{%
-0.024	0\\
-0.022	0\\
-0.02	0\\
-0.018	0\\
-0.016	0\\
-0.014	0\\
-0.012	0\\
-0.01	0\\
-0.008	0\\
-0.006	0\\
-0.004	0\\
-0.002	0\\
0	0\\
0.002	0.0391009\\
0.004	0.074092\\
0.006	0.104794\\
0.008	0.131117\\
0.01	0.153053\\
0.012	0.170674\\
0.014	0.18412\\
0.016	0.193597\\
0.018	0.19936\\
0.02	0.201711\\
0.022	0.200986\\
0.024	0.197545\\
0.026	0.191764\\
0.028	0.184026\\
0.03	0.174713\\
0.032	0.164198\\
0.034	0.152838\\
0.036	0.140971\\
0.038	0.128907\\
0.04	0.118714\\
0.042	0.12443\\
0.044	0.129408\\
0.046	0.133584\\
0.048	0.136907\\
0.05	0.139343\\
0.052	0.140872\\
0.054	0.141491\\
0.056	0.141209\\
0.058	0.140052\\
0.06	0.138058\\
0.062	0.135275\\
0.064	0.131764\\
0.066	0.127591\\
0.068	0.122832\\
0.07	0.117568\\
0.072	0.111883\\
0.074	0.105863\\
0.076	0.0995958\\
0.078	0.093167\\
0.08	0.086661\\
0.082	0.0862144\\
0.084	0.0879184\\
0.086	0.0895075\\
0.088	0.0909826\\
0.09	0.0923445\\
0.092	0.0935938\\
0.094	0.0947311\\
0.096	0.0957566\\
0.098	0.0966704\\
0.1	0.0974724\\
0.102	0.0981623\\
0.104	0.0987397\\
0.106	0.0992039\\
0.108	0.0995543\\
0.11	0.0997902\\
0.112	0.100297\\
0.114	0.101823\\
0.116	0.103169\\
0.118	0.104335\\
0.12	0.105325\\
0.122	0.10614\\
0.124	0.106784\\
0.126	0.107262\\
0.128	0.10758\\
0.13	0.107741\\
0.132	0.107753\\
0.134	0.107621\\
0.136	0.107352\\
0.138	0.106953\\
0.14	0.106431\\
0.142	0.105792\\
0.144	0.105043\\
0.146	0.104192\\
0.148	0.103245\\
0.15	0.102209\\
0.152	0.101091\\
0.154	0.099897\\
0.156	0.0986328\\
0.158	0.0973047\\
0.16	0.0959182\\
0.162	0.094479\\
0.164	0.092992\\
0.166	0.0914622\\
0.168	0.0898942\\
0.17	0.0882924\\
0.172	0.0866609\\
0.174	0.0850035\\
0.176	0.083324\\
0.178	0.0816256\\
0.18	0.0799117\\
0.182	0.0781851\\
0.184	0.0764488\\
0.186	0.0747052\\
0.188	0.072957\\
0.19	0.0712065\\
0.192	0.0694558\\
0.194	0.067707\\
0.196	0.065962\\
0.198	0.0642228\\
0.2	0.062491\\
0.202	0.0607684\\
0.204	0.0590565\\
0.206	0.0573568\\
0.208	0.0556706\\
0.21	0.0539995\\
0.212	0.0523446\\
0.214	0.0507072\\
0.216	0.0490884\\
0.218	0.0474892\\
0.22	0.0459108\\
0.222	0.044354\\
0.224	0.0428198\\
0.226	0.0413089\\
0.228	0.0398221\\
0.23	0.0383602\\
0.232	0.0369237\\
0.234	0.0355132\\
0.236	0.0341292\\
0.238	0.0327722\\
0.24	0.0314425\\
0.242	0.0301404\\
0.244	0.0288663\\
0.246	0.0276202\\
0.248	0.0264025\\
0.25	0.0252131\\
0.252	0.0240522\\
0.254	0.0231202\\
0.256	0.0222677\\
0.258	0.0214234\\
0.26	0.0205894\\
0.262	0.0197672\\
0.264	0.0189587\\
0.266	0.0181654\\
0.268	0.017389\\
0.27	0.0166309\\
0.272	0.0158924\\
0.274	0.0152408\\
0.276	0.0149727\\
0.278	0.0147109\\
0.28	0.0144554\\
0.282	0.014206\\
0.284	0.0139627\\
0.286	0.0137253\\
0.288	0.0134937\\
0.29	0.0132678\\
0.292	0.0130474\\
0.294	0.0128325\\
0.296	0.0126229\\
0.298	0.0124622\\
0.3	0.0126385\\
0.302	0.0127808\\
0.304	0.012889\\
0.306	0.0129627\\
0.308	0.0130021\\
0.31	0.0130075\\
0.312	0.012979\\
0.314	0.0129172\\
0.316	0.0128227\\
0.318	0.0126963\\
0.32	0.0125388\\
0.322	0.0123513\\
0.324	0.0121349\\
0.326	0.0118907\\
0.328	0.0116202\\
0.33	0.0113246\\
0.332	0.0110055\\
0.334	0.0106645\\
0.336	0.0103031\\
0.338	0.00992309\\
0.34	0.00952612\\
0.342	0.00911396\\
0.344	0.00900385\\
0.346	0.00925356\\
0.348	0.00948671\\
0.35	0.00970294\\
0.352	0.00990199\\
0.354	0.0100836\\
0.356	0.0102476\\
0.358	0.0103939\\
0.36	0.0105224\\
0.362	0.0106331\\
0.364	0.010726\\
0.366	0.0108011\\
0.368	0.0108587\\
0.37	0.0108987\\
0.372	0.0109215\\
0.374	0.0109273\\
0.376	0.0109163\\
0.378	0.0108889\\
0.38	0.0108453\\
0.382	0.010786\\
0.384	0.0107113\\
0.386	0.0106216\\
0.388	0.0105175\\
0.39	0.0103993\\
0.392	0.0102677\\
0.394	0.010123\\
0.396	0.00996583\\
0.398	0.00979674\\
0.4	0.00961629\\
0.402	0.00942505\\
0.404	0.0092236\\
0.406	0.00901256\\
0.408	0.00879252\\
0.41	0.00856409\\
0.412	0.00832791\\
0.414	0.00808458\\
0.416	0.00783474\\
0.418	0.00757901\\
0.42	0.007318\\
0.422	0.00705236\\
0.424	0.00678268\\
0.426	0.00650958\\
0.428	0.00623368\\
0.43	0.00595556\\
0.432	0.00567582\\
0.434	0.00539503\\
0.436	0.00511377\\
0.438	0.00483258\\
0.44	0.00455202\\
0.442	0.0042726\\
0.444	0.00399484\\
0.446	0.00371923\\
0.448	0.00344626\\
0.45	0.00328111\\
0.452	0.00319787\\
0.454	0.00311556\\
0.456	0.0030342\\
0.458	0.0029538\\
0.46	0.00287439\\
0.462	0.00279597\\
0.464	0.00271858\\
0.466	0.00264221\\
0.468	0.0025669\\
0.47	0.00249265\\
0.472	0.00241949\\
0.474	0.00246398\\
0.476	0.00252005\\
0.478	0.00257238\\
0.48	0.00262094\\
0.482	0.00266574\\
0.484	0.00270678\\
0.486	0.00274406\\
0.488	0.00277761\\
0.49	0.00280744\\
0.492	0.00283357\\
0.494	0.00285605\\
0.496	0.00287491\\
0.498	0.00289019\\
0.5	0.00290194\\
0.502	0.00291021\\
0.504	0.00291507\\
0.506	0.00291657\\
0.508	0.00291477\\
0.51	0.00290976\\
0.512	0.00290161\\
0.514	0.00289038\\
0.516	0.00287617\\
0.518	0.00285906\\
0.52	0.00283913\\
0.522	0.00281648\\
0.524	0.0027912\\
0.526	0.00276337\\
0.528	0.00273311\\
0.53	0.00270051\\
0.532	0.00266566\\
0.534	0.00262868\\
0.536	0.00258966\\
0.538	0.00254871\\
0.54	0.00250593\\
0.542	0.00246142\\
0.544	0.00241531\\
0.546	0.00236768\\
0.548	0.00231865\\
0.55	0.00226832\\
0.552	0.00221679\\
0.554	0.00216418\\
0.556	0.00211059\\
0.558	0.00205612\\
0.56	0.00200087\\
0.562	0.00194495\\
0.564	0.00188845\\
0.566	0.00183147\\
0.568	0.00177411\\
0.57	0.00171647\\
0.572	0.00165863\\
0.574	0.00160069\\
0.576	0.00154273\\
0.578	0.00148485\\
0.58	0.00142712\\
0.582	0.00136962\\
0.584	0.00131244\\
0.586	0.00125565\\
0.588	0.00119932\\
0.59	0.00114353\\
0.592	0.00108833\\
0.594	0.0010338\\
0.596	0.000979991\\
0.598	0.00093306\\
0.6	0.000902319\\
0.602	0.000867365\\
0.604	0.000828482\\
0.606	0.000785972\\
0.608	0.000740151\\
0.61	0.00069135\\
0.612	0.000639913\\
0.614	0.000586188\\
0.616	0.000600269\\
0.618	0.000660265\\
0.62	0.000717818\\
0.622	0.000772857\\
0.624	0.000825319\\
0.626	0.000875147\\
0.628	0.000922292\\
0.63	0.000966709\\
0.632	0.00100836\\
0.634	0.00104723\\
0.636	0.00108327\\
0.638	0.00111649\\
0.64	0.00114686\\
0.642	0.00117439\\
0.644	0.00119907\\
0.646	0.00122092\\
0.648	0.00123994\\
0.65	0.00125617\\
0.652	0.00126962\\
0.654	0.00128032\\
0.656	0.00128832\\
0.658	0.00129364\\
0.66	0.00129634\\
0.662	0.00129648\\
0.664	0.00129409\\
0.666	0.00128925\\
0.668	0.00128201\\
0.67	0.00127245\\
0.672	0.00126063\\
0.674	0.00124663\\
0.676	0.00123053\\
0.678	0.0012124\\
0.68	0.00119233\\
0.682	0.00117041\\
0.684	0.00114672\\
0.686	0.00112136\\
0.688	0.00109441\\
0.69	0.00106597\\
0.692	0.00103613\\
0.694	0.001005\\
0.696	0.000972661\\
0.698	0.000939218\\
0.7	0.000904769\\
0.702	0.00086941\\
0.704	0.00083324\\
0.706	0.000796358\\
0.708	0.000758859\\
0.71	0.000720842\\
0.712	0.000682402\\
0.714	0.000643633\\
0.716	0.000604629\\
0.718	0.000565482\\
0.72	0.000526281\\
0.722	0.000487116\\
0.724	0.000448073\\
0.726	0.000409234\\
0.728	0.000370683\\
0.73	0.000332498\\
0.732	0.000319977\\
0.734	0.000342665\\
0.736	0.000363189\\
0.738	0.000381478\\
0.74	0.000397474\\
0.742	0.000411133\\
0.744	0.000422427\\
0.746	0.00043134\\
0.748	0.000437872\\
0.75	0.000442036\\
0.752	0.000443859\\
0.754	0.000443379\\
0.756	0.000440648\\
0.758	0.000435729\\
0.76	0.000428697\\
0.762	0.000419635\\
0.764	0.000408639\\
0.766	0.000395811\\
0.768	0.000381263\\
0.77	0.000365113\\
0.772	0.000347486\\
0.774	0.000328513\\
0.776	0.000309042\\
0.778	0.000322976\\
0.78	0.000335632\\
0.782	0.000347018\\
0.784	0.000357144\\
0.786	0.000366026\\
0.788	0.000373679\\
0.79	0.000380122\\
0.792	0.000385379\\
0.794	0.000389473\\
0.796	0.00039243\\
0.798	0.00039428\\
0.8	0.000395054\\
0.802	0.000394784\\
0.804	0.000393505\\
0.806	0.000391253\\
0.808	0.000388065\\
0.81	0.00038398\\
0.812	0.00037904\\
0.814	0.000373283\\
0.816	0.000366754\\
0.818	0.000359493\\
0.82	0.000351546\\
0.822	0.000342955\\
0.824	0.000333766\\
0.826	0.000324021\\
0.828	0.000313767\\
0.83	0.000303047\\
0.832	0.000291905\\
0.834	0.000280386\\
0.836	0.000268533\\
0.838	0.000256389\\
0.84	0.000243997\\
0.842	0.000231398\\
0.844	0.000222867\\
0.846	0.000218976\\
0.848	0.000213948\\
0.85	0.000207846\\
0.852	0.000200737\\
0.854	0.000192693\\
0.856	0.000183791\\
0.858	0.000174107\\
0.86	0.000163723\\
0.862	0.000152722\\
0.864	0.000141187\\
0.866	0.000136407\\
0.868	0.000137292\\
0.87	0.000137927\\
0.872	0.000138315\\
0.874	0.000138456\\
0.876	0.000138355\\
0.878	0.000138014\\
0.88	0.000137437\\
0.882	0.000136628\\
0.884	0.000135591\\
0.886	0.000134333\\
0.888	0.000132858\\
0.89	0.000131172\\
0.892	0.000129282\\
0.894	0.000127195\\
0.896	0.000124918\\
0.898	0.000122459\\
0.9	0.000119826\\
0.902	0.000117027\\
0.904	0.00011407\\
0.906	0.000110965\\
0.908	0.000107721\\
0.91	0.000104346\\
0.912	0.000106734\\
0.914	0.000109845\\
0.916	0.000112556\\
0.918	0.000114876\\
0.92	0.000116812\\
0.922	0.000118375\\
0.924	0.000119575\\
0.926	0.000120423\\
0.928	0.000120929\\
0.93	0.000121107\\
0.932	0.000120968\\
0.934	0.000120525\\
0.936	0.000119792\\
0.938	0.00011878\\
0.94	0.000117505\\
0.942	0.000115979\\
0.944	0.000114217\\
0.946	0.000112233\\
0.948	0.00011004\\
0.95	0.000107652\\
0.952	0.000105084\\
0.954	0.000102349\\
0.956	9.94606e-05\\
0.958	9.64329e-05\\
0.96	9.3279e-05\\
0.962	9.00123e-05\\
0.964	8.66455e-05\\
0.966	8.31913e-05\\
0.968	7.96621e-05\\
0.97	7.60699e-05\\
0.972	7.24264e-05\\
0.974	6.87431e-05\\
0.976	6.50308e-05\\
0.978	6.13002e-05\\
0.98	5.75615e-05\\
0.982	5.48306e-05\\
0.984	5.83423e-05\\
0.986	6.14066e-05\\
0.988	6.40135e-05\\
0.99	6.61565e-05\\
0.992	6.78317e-05\\
0.994	6.90384e-05\\
0.996	6.97789e-05\\
0.998	7.00582e-05\\
1	6.98841e-05\\
1.002	6.9267e-05\\
1.004	6.82198e-05\\
1.006	6.67579e-05\\
1.008	6.48987e-05\\
1.01	6.26618e-05\\
1.012	6.00687e-05\\
1.014	5.71425e-05\\
1.016	5.39079e-05\\
1.018	5.0391e-05\\
1.02	4.83219e-05\\
1.022	4.72271e-05\\
1.024	4.59643e-05\\
1.026	4.54464e-05\\
1.028	4.51062e-05\\
1.03	4.46872e-05\\
1.032	4.41935e-05\\
1.034	4.36292e-05\\
1.036	4.29984e-05\\
1.038	4.2305e-05\\
1.04	4.15532e-05\\
1.042	4.07468e-05\\
1.044	3.98899e-05\\
1.046	3.89864e-05\\
1.048	3.80401e-05\\
1.05	3.70548e-05\\
1.052	3.60341e-05\\
1.054	3.49817e-05\\
1.056	3.39011e-05\\
1.058	3.51598e-05\\
1.06	3.81668e-05\\
1.062	4.09052e-05\\
1.064	4.33643e-05\\
1.066	4.55349e-05\\
1.068	4.741e-05\\
1.07	4.89843e-05\\
1.072	5.02546e-05\\
1.074	5.12194e-05\\
1.076	5.18791e-05\\
1.078	5.2236e-05\\
1.08	5.22939e-05\\
1.082	5.20584e-05\\
1.084	5.15368e-05\\
1.086	5.07379e-05\\
1.088	4.96719e-05\\
1.09	4.83502e-05\\
1.092	4.67857e-05\\
1.094	4.49924e-05\\
1.096	4.29851e-05\\
1.098	4.07799e-05\\
1.1	3.83935e-05\\
1.102	3.58432e-05\\
1.104	3.31472e-05\\
1.106	3.03237e-05\\
1.108	2.73918e-05\\
1.11	2.43702e-05\\
1.112	2.12783e-05\\
1.114	1.81351e-05\\
1.116	1.4995e-05\\
1.118	1.54104e-05\\
1.12	1.57843e-05\\
1.122	1.61158e-05\\
1.124	1.64039e-05\\
1.126	1.66481e-05\\
1.128	1.68478e-05\\
1.13	1.70025e-05\\
1.132	1.7112e-05\\
1.134	1.71763e-05\\
1.136	1.71954e-05\\
1.138	1.71696e-05\\
1.14	1.93632e-05\\
1.142	2.14642e-05\\
1.144	2.33966e-05\\
1.146	2.51532e-05\\
1.148	2.67282e-05\\
1.15	2.81169e-05\\
1.152	2.93158e-05\\
1.154	3.03225e-05\\
1.156	3.1136e-05\\
1.158	3.17561e-05\\
1.16	3.2184e-05\\
1.162	3.24219e-05\\
1.164	3.24729e-05\\
1.166	3.23414e-05\\
1.168	3.20325e-05\\
1.17	3.15522e-05\\
1.172	3.09076e-05\\
1.174	3.01063e-05\\
1.176	2.91567e-05\\
1.178	2.80679e-05\\
1.18	2.68496e-05\\
1.182	2.5512e-05\\
1.184	2.40656e-05\\
1.186	2.25215e-05\\
1.188	2.0891e-05\\
1.19	1.91856e-05\\
1.192	1.74169e-05\\
1.194	1.55967e-05\\
1.196	1.37368e-05\\
1.198	1.21149e-05\\
1.2	1.19539e-05\\
1.202	1.17753e-05\\
1.204	1.15799e-05\\
1.206	1.13682e-05\\
1.208	1.11408e-05\\
1.21	1.08986e-05\\
1.212	1.06421e-05\\
1.214	1.03722e-05\\
1.216	1.00894e-05\\
1.218	9.79471e-06\\
1.22	9.48875e-06\\
1.222	9.17234e-06\\
1.224	1.04168e-05\\
1.226	1.16444e-05\\
1.228	1.277e-05\\
1.23	1.37897e-05\\
1.232	1.47004e-05\\
1.234	1.54997e-05\\
1.236	1.61858e-05\\
1.238	1.67577e-05\\
1.24	1.72152e-05\\
1.242	1.75585e-05\\
1.244	1.77888e-05\\
1.246	1.79076e-05\\
1.248	1.79174e-05\\
1.25	1.78208e-05\\
1.252	1.76214e-05\\
1.254	1.7323e-05\\
1.256	1.69301e-05\\
1.258	1.64475e-05\\
1.26	1.58805e-05\\
1.262	1.52347e-05\\
1.264	1.4516e-05\\
1.266	1.37307e-05\\
1.268	1.28852e-05\\
1.27	1.19861e-05\\
1.272	1.10403e-05\\
1.274	1.00546e-05\\
1.276	9.03589e-06\\
1.278	7.99123e-06\\
1.28	6.92752e-06\\
1.282	6.31283e-06\\
1.284	6.15816e-06\\
1.286	5.99042e-06\\
1.288	5.81041e-06\\
1.29	5.61892e-06\\
1.292	5.41679e-06\\
1.294	5.20486e-06\\
1.296	4.98398e-06\\
1.298	4.75501e-06\\
1.3	4.51884e-06\\
1.302	4.46369e-06\\
1.304	4.9981e-06\\
1.306	5.75643e-06\\
1.308	6.45883e-06\\
1.31	7.10252e-06\\
1.312	7.68516e-06\\
1.314	8.20483e-06\\
1.316	8.66004e-06\\
1.318	9.04973e-06\\
1.32	9.37323e-06\\
1.322	9.63033e-06\\
1.324	9.8212e-06\\
1.326	9.94639e-06\\
1.328	1.00069e-05\\
1.33	1.00039e-05\\
1.332	9.93922e-06\\
1.334	9.81475e-06\\
1.336	9.63282e-06\\
1.338	9.39599e-06\\
1.34	9.10714e-06\\
1.342	8.76934e-06\\
1.344	8.38589e-06\\
1.346	7.96031e-06\\
1.348	7.49625e-06\\
1.35	6.99752e-06\\
1.352	6.46803e-06\\
1.354	5.91179e-06\\
1.356	5.33286e-06\\
1.358	4.73534e-06\\
1.36	4.12334e-06\\
1.362	3.50095e-06\\
1.364	2.87222e-06\\
1.366	2.36184e-06\\
1.368	2.43961e-06\\
1.37	2.50721e-06\\
1.372	2.56477e-06\\
1.374	2.6124e-06\\
1.376	2.65029e-06\\
1.378	2.67861e-06\\
1.38	2.69761e-06\\
1.382	2.70751e-06\\
1.384	2.94379e-06\\
1.386	3.40722e-06\\
1.388	3.8388e-06\\
1.39	4.23668e-06\\
1.392	4.59924e-06\\
1.394	4.92514e-06\\
1.396	5.21329e-06\\
1.398	5.46284e-06\\
1.4	5.67321e-06\\
1.402	5.84405e-06\\
1.404	5.97529e-06\\
1.406	6.06707e-06\\
1.408	6.11977e-06\\
1.41	6.13402e-06\\
1.412	6.11064e-06\\
1.414	6.05066e-06\\
1.416	5.95533e-06\\
1.418	5.82605e-06\\
1.42	5.66442e-06\\
1.422	5.47218e-06\\
1.424	5.25121e-06\\
1.426	5.00353e-06\\
1.428	4.73126e-06\\
1.43	4.43663e-06\\
1.432	4.12194e-06\\
1.434	3.78956e-06\\
1.436	3.44192e-06\\
1.438	3.08146e-06\\
1.44	2.71067e-06\\
1.442	2.33201e-06\\
1.444	1.95376e-06\\
1.446	1.79658e-06\\
1.448	1.63478e-06\\
1.45	1.46938e-06\\
1.452	1.38629e-06\\
1.454	1.32277e-06\\
1.456	1.28704e-06\\
1.458	1.27984e-06\\
1.46	1.26679e-06\\
1.462	1.3684e-06\\
1.464	1.68309e-06\\
1.466	1.98054e-06\\
1.468	2.25938e-06\\
1.47	2.51837e-06\\
1.472	2.75641e-06\\
1.474	2.97257e-06\\
1.476	3.16607e-06\\
1.478	3.33626e-06\\
1.48	3.48268e-06\\
1.482	3.605e-06\\
1.484	3.70305e-06\\
1.486	3.77681e-06\\
1.488	3.8264e-06\\
1.49	3.85208e-06\\
1.492	3.85425e-06\\
1.494	3.83346e-06\\
1.496	3.79035e-06\\
1.498	3.72569e-06\\
1.5	3.64039e-06\\
1.502	3.53541e-06\\
1.504	3.41184e-06\\
1.506	3.27086e-06\\
1.508	3.11369e-06\\
1.51	2.94166e-06\\
1.512	2.75613e-06\\
1.514	2.55852e-06\\
1.516	2.35027e-06\\
1.518	2.13289e-06\\
1.52	1.90786e-06\\
1.522	1.6767e-06\\
1.524	1.44093e-06\\
1.526	1.20206e-06\\
1.528	9.72951e-07\\
1.53	8.95589e-07\\
1.532	8.15562e-07\\
1.534	7.91851e-07\\
1.536	8.08302e-07\\
1.538	8.22616e-07\\
1.54	8.34815e-07\\
1.542	8.44924e-07\\
1.544	8.53778e-07\\
1.546	1.04847e-06\\
1.548	1.23247e-06\\
1.55	1.40494e-06\\
1.552	1.56514e-06\\
1.554	1.7124e-06\\
1.556	1.84615e-06\\
1.558	1.96591e-06\\
1.56	2.07132e-06\\
1.562	2.16207e-06\\
1.564	2.23798e-06\\
1.566	2.29895e-06\\
1.568	2.34496e-06\\
1.57	2.37609e-06\\
1.572	2.39251e-06\\
1.574	2.39447e-06\\
1.576	2.3823e-06\\
1.578	2.35639e-06\\
1.58	2.31722e-06\\
1.582	2.26534e-06\\
1.584	2.20135e-06\\
1.586	2.12592e-06\\
1.588	2.03975e-06\\
1.59	1.9436e-06\\
1.592	1.83829e-06\\
1.594	1.72464e-06\\
1.596	1.60351e-06\\
1.598	1.47581e-06\\
1.6	1.34244e-06\\
1.602	1.20431e-06\\
1.604	1.06235e-06\\
1.606	9.17489e-07\\
1.608	7.70646e-07\\
1.61	6.22731e-07\\
1.612	4.74639e-07\\
1.614	3.27246e-07\\
1.616	2.82797e-07\\
1.618	3.03036e-07\\
1.62	3.21151e-07\\
1.622	3.3711e-07\\
1.624	3.70608e-07\\
1.626	4.97131e-07\\
1.628	6.17757e-07\\
1.63	7.3192e-07\\
1.632	8.39104e-07\\
1.634	9.38846e-07\\
1.636	1.03074e-06\\
1.638	1.11443e-06\\
1.64	1.18963e-06\\
1.642	1.25609e-06\\
1.644	1.31364e-06\\
1.646	1.36215e-06\\
1.648	1.40156e-06\\
1.65	1.43185e-06\\
1.652	1.45307e-06\\
1.654	1.46531e-06\\
1.656	1.46871e-06\\
1.658	1.46348e-06\\
1.66	1.44984e-06\\
1.662	1.42809e-06\\
1.664	1.39854e-06\\
1.666	1.36156e-06\\
1.668	1.31755e-06\\
1.67	1.26694e-06\\
1.672	1.21018e-06\\
1.674	1.14776e-06\\
1.676	1.08018e-06\\
1.678	1.00797e-06\\
1.68	9.3166e-07\\
1.682	8.51798e-07\\
1.684	7.68943e-07\\
1.686	6.83654e-07\\
1.688	5.96491e-07\\
1.69	5.08012e-07\\
1.692	4.18769e-07\\
1.694	4.00741e-07\\
1.696	3.94882e-07\\
1.698	3.88127e-07\\
1.7	3.80515e-07\\
1.702	3.72087e-07\\
1.704	3.62884e-07\\
1.706	3.5295e-07\\
1.708	3.4233e-07\\
1.71	3.34415e-07\\
1.712	4.0428e-07\\
1.714	4.7001e-07\\
1.716	5.31322e-07\\
1.718	5.87964e-07\\
1.72	6.3972e-07\\
1.722	6.86408e-07\\
1.724	7.27879e-07\\
1.726	7.6402e-07\\
1.728	7.94752e-07\\
1.73	8.20029e-07\\
1.732	8.39839e-07\\
1.734	8.54202e-07\\
1.736	8.63172e-07\\
1.738	8.66831e-07\\
1.74	8.65292e-07\\
1.742	8.58694e-07\\
1.744	8.47205e-07\\
1.746	8.31016e-07\\
1.748	8.10341e-07\\
1.75	7.85417e-07\\
1.752	7.56499e-07\\
1.754	7.23859e-07\\
1.756	6.87785e-07\\
1.758	6.48577e-07\\
1.76	6.06548e-07\\
1.762	5.62018e-07\\
1.764	5.15315e-07\\
1.766	4.66772e-07\\
1.768	4.16724e-07\\
1.77	3.65505e-07\\
1.772	3.13451e-07\\
1.774	2.60892e-07\\
1.776	2.08153e-07\\
1.778	1.55552e-07\\
1.78	1.17716e-07\\
1.782	1.24419e-07\\
1.784	1.30557e-07\\
1.786	1.3613e-07\\
1.788	1.41142e-07\\
1.79	1.45595e-07\\
1.792	1.84256e-07\\
1.794	2.25655e-07\\
1.796	2.64617e-07\\
1.798	3.00972e-07\\
1.8	3.34569e-07\\
1.802	3.65277e-07\\
1.804	3.92986e-07\\
1.806	4.17607e-07\\
1.808	4.3907e-07\\
1.81	4.57327e-07\\
1.812	4.7235e-07\\
1.814	4.84131e-07\\
1.816	4.92682e-07\\
1.818	4.98034e-07\\
1.82	5.00235e-07\\
1.822	4.99352e-07\\
1.824	4.95468e-07\\
1.826	4.88683e-07\\
1.828	4.79111e-07\\
1.83	4.6688e-07\\
1.832	4.5213e-07\\
1.834	4.35014e-07\\
1.836	4.15694e-07\\
1.838	3.94342e-07\\
1.84	3.71138e-07\\
1.842	3.46268e-07\\
1.844	3.19924e-07\\
1.846	2.92302e-07\\
1.848	2.63599e-07\\
1.85	2.34017e-07\\
1.852	2.03755e-07\\
1.854	1.73013e-07\\
1.856	1.41989e-07\\
1.858	1.10876e-07\\
1.86	7.98647e-08\\
1.862	5.59403e-08\\
1.864	5.08154e-08\\
1.866	4.57597e-08\\
1.868	4.14637e-08\\
1.87	6.74409e-08\\
1.872	9.42142e-08\\
1.874	1.19755e-07\\
1.876	1.43938e-07\\
1.878	1.66652e-07\\
1.88	1.87796e-07\\
1.882	2.07279e-07\\
1.884	2.25025e-07\\
1.886	2.40969e-07\\
1.888	2.5506e-07\\
1.89	2.67256e-07\\
1.892	2.7753e-07\\
1.894	2.85867e-07\\
1.896	2.92264e-07\\
1.898	2.96728e-07\\
1.9	2.99281e-07\\
1.902	2.99951e-07\\
1.904	2.98782e-07\\
1.906	2.95824e-07\\
1.908	2.91138e-07\\
1.91	2.84794e-07\\
1.912	2.7687e-07\\
1.914	2.67451e-07\\
1.916	2.56631e-07\\
1.918	2.44507e-07\\
1.92	2.31185e-07\\
1.922	2.16773e-07\\
1.924	2.01384e-07\\
1.926	1.85133e-07\\
1.928	1.68139e-07\\
1.93	1.50522e-07\\
1.932	1.32403e-07\\
1.934	1.13902e-07\\
1.936	9.51407e-08\\
1.938	7.62929e-08\\
1.94	6.79136e-08\\
1.942	5.93789e-08\\
1.944	5.07409e-08\\
1.946	4.63789e-08\\
1.948	4.60564e-08\\
1.95	4.56499e-08\\
1.952	5.06511e-08\\
1.954	6.67765e-08\\
1.956	8.21339e-08\\
1.958	9.66492e-08\\
1.96	1.10255e-07\\
1.962	1.22892e-07\\
1.964	1.34506e-07\\
1.966	1.45052e-07\\
1.968	1.54493e-07\\
1.97	1.62796e-07\\
1.972	1.69939e-07\\
1.974	1.75907e-07\\
1.976	1.80691e-07\\
1.978	1.8429e-07\\
1.98	1.86711e-07\\
1.982	1.87966e-07\\
1.984	1.88075e-07\\
1.986	1.87064e-07\\
1.988	1.84965e-07\\
1.99	1.81817e-07\\
1.992	1.77661e-07\\
1.994	1.72546e-07\\
1.996	1.66526e-07\\
1.998	1.59657e-07\\
2	1.51999e-07\\
};
\addplot [color=mycolor4, line width=1.0pt, forget plot]
  table[row sep=crcr]{%
-0.024	0\\
-0.022	0\\
-0.02	0\\
-0.018	0\\
-0.016	0\\
-0.014	0\\
-0.012	0\\
-0.01	0\\
-0.008	0\\
-0.006	0\\
-0.004	0\\
-0.002	0\\
0	0\\
0.002	0.0391009\\
0.004	0.074092\\
0.006	0.104794\\
0.008	0.131117\\
0.01	0.153053\\
0.012	0.170674\\
0.014	0.18412\\
0.016	0.193597\\
0.018	0.19936\\
0.02	0.201711\\
0.022	0.200986\\
0.024	0.197544\\
0.026	0.191763\\
0.028	0.184025\\
0.03	0.174712\\
0.032	0.164196\\
0.034	0.152837\\
0.036	0.140969\\
0.038	0.128904\\
0.04	0.118712\\
0.042	0.124427\\
0.044	0.129405\\
0.046	0.13358\\
0.048	0.136902\\
0.05	0.139337\\
0.052	0.140865\\
0.054	0.141482\\
0.056	0.141199\\
0.058	0.140041\\
0.06	0.138045\\
0.062	0.13526\\
0.064	0.131747\\
0.066	0.127572\\
0.068	0.12281\\
0.07	0.117543\\
0.072	0.111855\\
0.074	0.105832\\
0.076	0.0995607\\
0.078	0.0931279\\
0.08	0.0866174\\
0.082	0.0860941\\
0.084	0.0877876\\
0.086	0.0893647\\
0.088	0.0908264\\
0.09	0.0921732\\
0.092	0.0934057\\
0.094	0.0945245\\
0.096	0.0955297\\
0.098	0.0964214\\
0.1	0.0971995\\
0.102	0.0978637\\
0.104	0.0984136\\
0.106	0.0988487\\
0.108	0.0991683\\
0.11	0.0993717\\
0.112	0.0999624\\
0.114	0.101462\\
0.116	0.102781\\
0.118	0.103919\\
0.12	0.104879\\
0.122	0.105664\\
0.124	0.106278\\
0.126	0.106726\\
0.128	0.107012\\
0.13	0.107142\\
0.132	0.107122\\
0.134	0.106959\\
0.136	0.106659\\
0.138	0.106229\\
0.14	0.105675\\
0.142	0.105006\\
0.144	0.104229\\
0.146	0.103349\\
0.148	0.102375\\
0.15	0.101312\\
0.152	0.100169\\
0.154	0.0989508\\
0.156	0.0976642\\
0.158	0.0963152\\
0.16	0.0949097\\
0.162	0.0934531\\
0.164	0.0919507\\
0.166	0.0904076\\
0.168	0.0888283\\
0.17	0.0872174\\
0.172	0.0855791\\
0.174	0.0839172\\
0.176	0.0822356\\
0.178	0.0805376\\
0.18	0.0788266\\
0.182	0.0771055\\
0.184	0.0753772\\
0.186	0.0736444\\
0.188	0.0719095\\
0.19	0.070175\\
0.192	0.068443\\
0.194	0.0667155\\
0.196	0.0649946\\
0.198	0.0632821\\
0.2	0.0615796\\
0.202	0.0598889\\
0.204	0.0582113\\
0.206	0.0565485\\
0.208	0.0549017\\
0.21	0.0532722\\
0.212	0.0516613\\
0.214	0.05007\\
0.216	0.0484993\\
0.218	0.0469504\\
0.22	0.045424\\
0.222	0.043921\\
0.224	0.0424421\\
0.226	0.0409881\\
0.228	0.0395596\\
0.23	0.0381571\\
0.232	0.0367811\\
0.234	0.0354321\\
0.236	0.0341102\\
0.238	0.0328159\\
0.24	0.0315494\\
0.242	0.0303107\\
0.244	0.0291\\
0.246	0.0279174\\
0.248	0.0267628\\
0.25	0.0256361\\
0.252	0.0245372\\
0.254	0.023466\\
0.256	0.0224701\\
0.258	0.0216519\\
0.26	0.0208429\\
0.262	0.0200449\\
0.264	0.0192596\\
0.266	0.0184885\\
0.268	0.0177332\\
0.27	0.0169951\\
0.272	0.0162755\\
0.274	0.0155755\\
0.276	0.0149664\\
0.278	0.0147021\\
0.28	0.0144441\\
0.282	0.0141924\\
0.284	0.0139468\\
0.286	0.0137073\\
0.288	0.0134736\\
0.29	0.0132457\\
0.292	0.0130234\\
0.294	0.0128066\\
0.296	0.0125951\\
0.298	0.0123887\\
0.3	0.0121874\\
0.302	0.0119909\\
0.304	0.0117991\\
0.306	0.0116117\\
0.308	0.0114287\\
0.31	0.0112498\\
0.312	0.0110748\\
0.314	0.0109037\\
0.316	0.0107362\\
0.318	0.0105721\\
0.32	0.0104113\\
0.322	0.0102536\\
0.324	0.0100989\\
0.326	0.00994696\\
0.328	0.00979767\\
0.33	0.00965086\\
0.332	0.00950641\\
0.334	0.00936416\\
0.336	0.009224\\
0.338	0.00908579\\
0.34	0.00911526\\
0.342	0.00931028\\
0.344	0.00948941\\
0.346	0.00965245\\
0.348	0.00979926\\
0.35	0.00992972\\
0.352	0.0100438\\
0.354	0.0101415\\
0.356	0.0102229\\
0.358	0.0102881\\
0.36	0.0103372\\
0.362	0.0103704\\
0.364	0.0103879\\
0.366	0.01039\\
0.368	0.0103769\\
0.37	0.010349\\
0.372	0.0103067\\
0.374	0.0102502\\
0.376	0.0101801\\
0.378	0.0100967\\
0.38	0.0100004\\
0.382	0.00989182\\
0.384	0.0097713\\
0.386	0.00963937\\
0.388	0.00949654\\
0.39	0.00934332\\
0.392	0.00918024\\
0.394	0.00900783\\
0.396	0.00882664\\
0.398	0.00863719\\
0.4	0.00844005\\
0.402	0.00823576\\
0.404	0.00802488\\
0.406	0.00780795\\
0.408	0.00758552\\
0.41	0.00735814\\
0.412	0.00712635\\
0.414	0.00689069\\
0.416	0.00665169\\
0.418	0.00640988\\
0.42	0.00616576\\
0.422	0.00591985\\
0.424	0.00567265\\
0.426	0.00542463\\
0.428	0.00517629\\
0.43	0.00492807\\
0.432	0.00468043\\
0.434	0.00443381\\
0.436	0.00418864\\
0.438	0.00394531\\
0.44	0.00370423\\
0.442	0.00346577\\
0.444	0.00327912\\
0.446	0.00319777\\
0.448	0.00311774\\
0.45	0.00303903\\
0.452	0.00296167\\
0.454	0.00288564\\
0.456	0.00281096\\
0.458	0.00273763\\
0.46	0.00266564\\
0.462	0.00259501\\
0.464	0.00252572\\
0.466	0.00245778\\
0.468	0.00239119\\
0.47	0.00232593\\
0.472	0.002262\\
0.474	0.0021994\\
0.476	0.00219144\\
0.478	0.00225804\\
0.48	0.00232109\\
0.482	0.00238055\\
0.484	0.00243637\\
0.486	0.00248852\\
0.488	0.00253697\\
0.49	0.00258172\\
0.492	0.00262276\\
0.494	0.00266008\\
0.496	0.00269368\\
0.498	0.00272359\\
0.5	0.00274983\\
0.502	0.00277242\\
0.504	0.00279138\\
0.506	0.00280678\\
0.508	0.00281864\\
0.51	0.00282702\\
0.512	0.00283197\\
0.514	0.00283356\\
0.516	0.00283185\\
0.518	0.00282691\\
0.52	0.00281881\\
0.522	0.00280764\\
0.524	0.00279348\\
0.526	0.00277641\\
0.528	0.00275652\\
0.53	0.00273391\\
0.532	0.00270866\\
0.534	0.00268087\\
0.536	0.00265065\\
0.538	0.00261809\\
0.54	0.0025833\\
0.542	0.00254638\\
0.544	0.00250743\\
0.546	0.00246656\\
0.548	0.00242388\\
0.55	0.00237948\\
0.552	0.00233349\\
0.554	0.00228601\\
0.556	0.00223714\\
0.558	0.00218699\\
0.56	0.00213566\\
0.562	0.00208326\\
0.564	0.0020299\\
0.566	0.00197567\\
0.568	0.00192068\\
0.57	0.00186503\\
0.572	0.00180881\\
0.574	0.00175213\\
0.576	0.00169507\\
0.578	0.00163773\\
0.58	0.00158021\\
0.582	0.00152258\\
0.584	0.00146493\\
0.586	0.00140735\\
0.588	0.00134993\\
0.59	0.00129273\\
0.592	0.00123583\\
0.594	0.00117931\\
0.596	0.00112324\\
0.598	0.00106768\\
0.6	0.0010127\\
0.602	0.000958368\\
0.604	0.000904734\\
0.606	0.000851856\\
0.608	0.000799788\\
0.61	0.000748581\\
0.612	0.000698283\\
0.614	0.000648936\\
0.616	0.000600584\\
0.618	0.000554309\\
0.62	0.000593455\\
0.622	0.000630754\\
0.624	0.000666175\\
0.626	0.000699689\\
0.628	0.000731273\\
0.63	0.00076091\\
0.632	0.000788586\\
0.634	0.000814292\\
0.636	0.000838024\\
0.638	0.000859782\\
0.64	0.000879569\\
0.642	0.000897393\\
0.644	0.000913267\\
0.646	0.000927205\\
0.648	0.000939228\\
0.65	0.000949357\\
0.652	0.000957619\\
0.654	0.000964043\\
0.656	0.000968661\\
0.658	0.000971508\\
0.66	0.000972622\\
0.662	0.000972043\\
0.664	0.000969815\\
0.666	0.000965981\\
0.668	0.00096059\\
0.67	0.00095369\\
0.672	0.000945332\\
0.674	0.000935568\\
0.676	0.000924454\\
0.678	0.000912043\\
0.68	0.000898392\\
0.682	0.00088356\\
0.684	0.000867604\\
0.686	0.000850583\\
0.688	0.000832557\\
0.69	0.000813587\\
0.692	0.000793733\\
0.694	0.000773055\\
0.696	0.000751615\\
0.698	0.000729473\\
0.7	0.00070669\\
0.702	0.000683326\\
0.704	0.00065944\\
0.706	0.000635092\\
0.708	0.000610341\\
0.71	0.000585245\\
0.712	0.00055986\\
0.714	0.000534242\\
0.716	0.000508448\\
0.718	0.00048253\\
0.72	0.000456541\\
0.722	0.000430533\\
0.724	0.000404556\\
0.726	0.000378659\\
0.728	0.000352889\\
0.73	0.000327292\\
0.732	0.000301911\\
0.734	0.00027679\\
0.736	0.000264948\\
0.738	0.000256144\\
0.74	0.000247123\\
0.742	0.00023791\\
0.744	0.000228528\\
0.746	0.000219002\\
0.748	0.000209353\\
0.75	0.000199604\\
0.752	0.000189777\\
0.754	0.000179892\\
0.756	0.000169969\\
0.758	0.000160028\\
0.76	0.000150086\\
0.762	0.000140163\\
0.764	0.000130276\\
0.766	0.00012044\\
0.768	0.000110672\\
0.77	0.000100987\\
0.772	0.000107692\\
0.774	0.000121454\\
0.776	0.000134481\\
0.778	0.000146768\\
0.78	0.000158315\\
0.782	0.00016912\\
0.784	0.000179187\\
0.786	0.000188516\\
0.788	0.000197114\\
0.79	0.000204985\\
0.792	0.000212137\\
0.794	0.000218578\\
0.796	0.000224318\\
0.798	0.000229367\\
0.8	0.000233738\\
0.802	0.000237442\\
0.804	0.000240494\\
0.806	0.000242909\\
0.808	0.000244701\\
0.81	0.000245889\\
0.812	0.000246489\\
0.814	0.000246518\\
0.816	0.000245996\\
0.818	0.000244941\\
0.82	0.000243374\\
0.822	0.000241314\\
0.824	0.000238782\\
0.826	0.000235799\\
0.828	0.000232386\\
0.83	0.000228565\\
0.832	0.000224357\\
0.834	0.000219785\\
0.836	0.000214869\\
0.838	0.000209632\\
0.84	0.000204097\\
0.842	0.000198284\\
0.844	0.000192216\\
0.846	0.000185915\\
0.848	0.000179401\\
0.85	0.000172696\\
0.852	0.000165822\\
0.854	0.000158799\\
0.856	0.000151647\\
0.858	0.000144386\\
0.86	0.000137036\\
0.862	0.000129616\\
0.864	0.000124192\\
0.866	0.00012361\\
0.868	0.000122857\\
0.87	0.000121937\\
0.872	0.000120857\\
0.874	0.000119624\\
0.876	0.000118242\\
0.878	0.000116719\\
0.88	0.000115061\\
0.882	0.000113273\\
0.884	0.000111363\\
0.886	0.000109336\\
0.888	0.000107198\\
0.89	0.000104957\\
0.892	0.000102618\\
0.894	0.000100187\\
0.896	9.76711e-05\\
0.898	9.50763e-05\\
0.9	9.24087e-05\\
0.902	8.96745e-05\\
0.904	8.68797e-05\\
0.906	8.40304e-05\\
0.908	8.11326e-05\\
0.91	7.8192e-05\\
0.912	7.52144e-05\\
0.914	7.22057e-05\\
0.916	6.91713e-05\\
0.918	6.61169e-05\\
0.92	6.30477e-05\\
0.922	5.9969e-05\\
0.924	5.68861e-05\\
0.926	5.96814e-05\\
0.928	6.24672e-05\\
0.93	6.50268e-05\\
0.932	6.73617e-05\\
0.934	6.94742e-05\\
0.936	7.13669e-05\\
0.938	7.30426e-05\\
0.94	7.45048e-05\\
0.942	7.57572e-05\\
0.944	7.68039e-05\\
0.946	7.76494e-05\\
0.948	7.82984e-05\\
0.95	7.87559e-05\\
0.952	7.90272e-05\\
0.954	7.91179e-05\\
0.956	7.90338e-05\\
0.958	7.87808e-05\\
0.96	7.83651e-05\\
0.962	7.77931e-05\\
0.964	7.70712e-05\\
0.966	7.62061e-05\\
0.968	7.52045e-05\\
0.97	7.40731e-05\\
0.972	7.28189e-05\\
0.974	7.14488e-05\\
0.976	6.99697e-05\\
0.978	6.83887e-05\\
0.98	6.67127e-05\\
0.982	6.49487e-05\\
0.984	6.31035e-05\\
0.986	6.11842e-05\\
0.988	5.91975e-05\\
0.99	5.71502e-05\\
0.992	5.50489e-05\\
0.994	5.29002e-05\\
0.996	5.07105e-05\\
0.998	4.84863e-05\\
1	4.62336e-05\\
1.002	4.39585e-05\\
1.004	4.1667e-05\\
1.006	3.93647e-05\\
1.008	3.70572e-05\\
1.01	3.47499e-05\\
1.012	3.2448e-05\\
1.014	3.14434e-05\\
1.016	3.17502e-05\\
1.018	3.19891e-05\\
1.02	3.21619e-05\\
1.022	3.22702e-05\\
1.024	3.23159e-05\\
1.026	3.23008e-05\\
1.028	3.22268e-05\\
1.03	3.2096e-05\\
1.032	3.19103e-05\\
1.034	3.16718e-05\\
1.036	3.13826e-05\\
1.038	3.10448e-05\\
1.04	3.06607e-05\\
1.042	3.02323e-05\\
1.044	2.97619e-05\\
1.046	2.92517e-05\\
1.048	2.87039e-05\\
1.05	2.81208e-05\\
1.052	2.75045e-05\\
1.054	2.68573e-05\\
1.056	2.61813e-05\\
1.058	2.54789e-05\\
1.06	2.4752e-05\\
1.062	2.4003e-05\\
1.064	2.32339e-05\\
1.066	2.24468e-05\\
1.068	2.16438e-05\\
1.07	2.08269e-05\\
1.072	1.99981e-05\\
1.074	1.91594e-05\\
1.076	1.83126e-05\\
1.078	1.84859e-05\\
1.08	1.89633e-05\\
1.082	1.93787e-05\\
1.084	1.97334e-05\\
1.086	2.00287e-05\\
1.088	2.02661e-05\\
1.09	2.04473e-05\\
1.092	2.05739e-05\\
1.094	2.06476e-05\\
1.096	2.06703e-05\\
1.098	2.06437e-05\\
1.1	2.05698e-05\\
1.102	2.04505e-05\\
1.104	2.02878e-05\\
1.106	2.00836e-05\\
1.108	1.984e-05\\
1.11	1.95591e-05\\
1.112	1.92429e-05\\
1.114	1.88935e-05\\
1.116	1.85129e-05\\
1.118	1.81033e-05\\
1.12	1.76666e-05\\
1.122	1.72049e-05\\
1.124	1.67203e-05\\
1.126	1.62147e-05\\
1.128	1.56901e-05\\
1.13	1.51486e-05\\
1.132	1.45919e-05\\
1.134	1.40221e-05\\
1.136	1.34409e-05\\
1.138	1.28502e-05\\
1.14	1.22517e-05\\
1.142	1.16472e-05\\
1.144	1.10382e-05\\
1.146	1.04265e-05\\
1.148	9.81352e-06\\
1.15	9.20081e-06\\
1.152	8.58982e-06\\
1.154	7.98194e-06\\
1.156	7.42592e-06\\
1.158	7.41867e-06\\
1.16	7.56487e-06\\
1.162	7.75316e-06\\
1.164	7.92094e-06\\
1.166	8.0686e-06\\
1.168	8.19658e-06\\
1.17	8.30532e-06\\
1.172	8.39531e-06\\
1.174	8.46703e-06\\
1.176	8.52099e-06\\
1.178	8.55773e-06\\
1.18	8.57777e-06\\
1.182	8.58167e-06\\
1.184	8.57e-06\\
1.186	8.54332e-06\\
1.188	8.50222e-06\\
1.19	8.44729e-06\\
1.192	8.37912e-06\\
1.194	8.29831e-06\\
1.196	8.20546e-06\\
1.198	8.10116e-06\\
1.2	7.98602e-06\\
1.202	7.86063e-06\\
1.204	7.7256e-06\\
1.206	7.58152e-06\\
1.208	7.42898e-06\\
1.21	7.26855e-06\\
1.212	7.10083e-06\\
1.214	6.92638e-06\\
1.216	6.74575e-06\\
1.218	6.55951e-06\\
1.22	6.3682e-06\\
1.222	6.17236e-06\\
1.224	5.9725e-06\\
1.226	5.76913e-06\\
1.228	5.56276e-06\\
1.23	5.35387e-06\\
1.232	5.22812e-06\\
1.234	5.24315e-06\\
1.236	5.24516e-06\\
1.238	5.23467e-06\\
1.24	5.21219e-06\\
1.242	5.17823e-06\\
1.244	5.13335e-06\\
1.246	5.07807e-06\\
1.248	5.01293e-06\\
1.25	4.93849e-06\\
1.252	4.85529e-06\\
1.254	4.76387e-06\\
1.256	4.66479e-06\\
1.258	4.55858e-06\\
1.26	4.44577e-06\\
1.262	4.32691e-06\\
1.264	4.20253e-06\\
1.266	4.07313e-06\\
1.268	3.93923e-06\\
1.27	3.80134e-06\\
1.272	3.65995e-06\\
1.274	3.51553e-06\\
1.276	3.36856e-06\\
1.278	3.2195e-06\\
1.28	3.06879e-06\\
1.282	2.91686e-06\\
1.284	2.76413e-06\\
1.286	2.62724e-06\\
1.288	2.61915e-06\\
1.29	2.60476e-06\\
1.292	2.58433e-06\\
1.294	2.55816e-06\\
1.296	2.52654e-06\\
1.298	2.48976e-06\\
1.3	2.44814e-06\\
1.302	2.40197e-06\\
1.304	2.35155e-06\\
1.306	2.2972e-06\\
1.308	2.23921e-06\\
1.31	2.17789e-06\\
1.312	2.1226e-06\\
1.314	2.06494e-06\\
1.316	2.00377e-06\\
1.318	1.93941e-06\\
1.32	1.87216e-06\\
1.322	1.80233e-06\\
1.324	1.73021e-06\\
1.326	1.70271e-06\\
1.328	1.73543e-06\\
1.33	1.76372e-06\\
1.332	1.78769e-06\\
1.334	1.80746e-06\\
1.336	1.82314e-06\\
1.338	1.83486e-06\\
1.34	1.84274e-06\\
1.342	1.84692e-06\\
1.344	1.84751e-06\\
1.346	1.84465e-06\\
1.348	1.83849e-06\\
1.35	1.82914e-06\\
1.352	1.81674e-06\\
1.354	1.80144e-06\\
1.356	1.78336e-06\\
1.358	1.76265e-06\\
1.36	1.73944e-06\\
1.362	1.71387e-06\\
1.364	1.68607e-06\\
1.366	1.65618e-06\\
1.368	1.62433e-06\\
1.37	1.59065e-06\\
1.372	1.55527e-06\\
1.374	1.51833e-06\\
1.376	1.47995e-06\\
1.378	1.44025e-06\\
1.38	1.39936e-06\\
1.382	1.35741e-06\\
1.384	1.3145e-06\\
1.386	1.27076e-06\\
1.388	1.22629e-06\\
1.39	1.18121e-06\\
1.392	1.13563e-06\\
1.394	1.08965e-06\\
1.396	1.04337e-06\\
1.398	9.96889e-07\\
1.4	9.50303e-07\\
1.402	9.03702e-07\\
1.404	8.66174e-07\\
1.406	8.33571e-07\\
1.408	8.00053e-07\\
1.41	7.65749e-07\\
1.412	7.30779e-07\\
1.414	6.95265e-07\\
1.416	6.59323e-07\\
1.418	6.44804e-07\\
1.42	6.44087e-07\\
1.422	6.41644e-07\\
1.424	6.37561e-07\\
1.426	6.31925e-07\\
1.428	6.24824e-07\\
1.43	6.16344e-07\\
1.432	6.06576e-07\\
1.434	5.95605e-07\\
1.436	5.84516e-07\\
1.438	5.77108e-07\\
1.44	5.68869e-07\\
1.442	5.59857e-07\\
1.444	5.50129e-07\\
1.446	5.39742e-07\\
1.448	5.28751e-07\\
1.45	5.17209e-07\\
1.452	5.05168e-07\\
1.454	4.92676e-07\\
1.456	4.79782e-07\\
1.458	4.66532e-07\\
1.46	4.5297e-07\\
1.462	4.39139e-07\\
1.464	4.2508e-07\\
1.466	4.10832e-07\\
1.468	3.96431e-07\\
1.47	3.81914e-07\\
1.472	3.67313e-07\\
1.474	3.52661e-07\\
1.476	3.37988e-07\\
1.478	3.23323e-07\\
1.48	3.26854e-07\\
1.482	3.36923e-07\\
1.484	3.45963e-07\\
1.486	3.53997e-07\\
1.488	3.61048e-07\\
1.49	3.67141e-07\\
1.492	3.72302e-07\\
1.494	3.76558e-07\\
1.496	3.79936e-07\\
1.498	3.82464e-07\\
1.5	3.8417e-07\\
1.502	3.85084e-07\\
1.504	3.85235e-07\\
1.506	3.84652e-07\\
1.508	3.83366e-07\\
1.51	3.81408e-07\\
1.512	3.78807e-07\\
1.514	3.75594e-07\\
1.516	3.71799e-07\\
1.518	3.67453e-07\\
1.52	3.62586e-07\\
1.522	3.57229e-07\\
1.524	3.51411e-07\\
1.526	3.45161e-07\\
1.528	3.3851e-07\\
1.53	3.31485e-07\\
1.532	3.24116e-07\\
1.534	3.16431e-07\\
1.536	3.08456e-07\\
1.538	3.0022e-07\\
1.54	2.91748e-07\\
1.542	2.83066e-07\\
1.544	2.742e-07\\
1.546	2.65174e-07\\
1.548	2.56012e-07\\
1.55	2.46737e-07\\
1.552	2.37371e-07\\
1.554	2.27937e-07\\
1.556	2.18455e-07\\
1.558	2.08946e-07\\
1.56	1.99429e-07\\
1.562	1.89922e-07\\
1.564	1.80444e-07\\
1.566	1.71012e-07\\
1.568	1.61641e-07\\
1.57	1.52348e-07\\
1.572	1.43147e-07\\
1.574	1.34052e-07\\
1.576	1.25076e-07\\
1.578	1.16232e-07\\
1.58	1.1427e-07\\
1.582	1.15453e-07\\
1.584	1.16404e-07\\
1.586	1.17131e-07\\
1.588	1.17639e-07\\
1.59	1.17935e-07\\
1.592	1.18026e-07\\
1.594	1.17919e-07\\
1.596	1.17621e-07\\
1.598	1.17138e-07\\
1.6	1.16478e-07\\
1.602	1.15648e-07\\
1.604	1.14654e-07\\
1.606	1.13506e-07\\
1.608	1.12208e-07\\
1.61	1.10769e-07\\
1.612	1.09196e-07\\
1.614	1.07497e-07\\
1.616	1.05677e-07\\
1.618	1.03746e-07\\
1.62	1.0371e-07\\
1.622	1.03683e-07\\
1.624	1.03346e-07\\
1.626	1.02707e-07\\
1.628	1.01779e-07\\
1.63	1.00574e-07\\
1.632	9.91025e-08\\
1.634	9.73783e-08\\
1.636	9.54138e-08\\
1.638	9.32222e-08\\
1.64	9.08168e-08\\
1.642	8.82111e-08\\
1.644	8.54187e-08\\
1.646	8.24533e-08\\
1.648	7.93288e-08\\
1.65	7.6059e-08\\
1.652	7.26576e-08\\
1.654	6.91384e-08\\
1.656	6.93444e-08\\
1.658	7.00343e-08\\
1.66	7.05711e-08\\
1.662	7.09604e-08\\
1.664	7.12078e-08\\
1.666	7.13189e-08\\
1.668	7.12994e-08\\
1.67	7.1155e-08\\
1.672	7.08912e-08\\
1.674	7.05136e-08\\
1.676	7.0028e-08\\
1.678	6.94396e-08\\
1.68	6.87541e-08\\
1.682	6.79769e-08\\
1.684	6.71133e-08\\
1.686	6.61685e-08\\
1.688	6.51477e-08\\
1.69	6.40561e-08\\
1.692	6.28986e-08\\
1.694	6.16801e-08\\
1.696	6.04053e-08\\
1.698	5.90791e-08\\
1.7	5.77058e-08\\
1.702	5.629e-08\\
1.704	5.4836e-08\\
1.706	5.33479e-08\\
1.708	5.18298e-08\\
1.71	5.02856e-08\\
1.712	4.87192e-08\\
1.714	4.71342e-08\\
1.716	4.55342e-08\\
1.718	4.39225e-08\\
1.72	4.23025e-08\\
1.722	4.0908e-08\\
1.724	4.23004e-08\\
1.726	4.35484e-08\\
1.728	4.4653e-08\\
1.73	4.5615e-08\\
1.732	4.64362e-08\\
1.734	4.71182e-08\\
1.736	4.76631e-08\\
1.738	4.80735e-08\\
1.74	4.8352e-08\\
1.742	4.85016e-08\\
1.744	4.85255e-08\\
1.746	4.84271e-08\\
1.748	4.82102e-08\\
1.75	4.78786e-08\\
1.752	4.74364e-08\\
1.754	4.68877e-08\\
1.756	4.6237e-08\\
1.758	4.54888e-08\\
1.76	4.46475e-08\\
1.762	4.37181e-08\\
1.764	4.27053e-08\\
1.766	4.16139e-08\\
1.768	4.04489e-08\\
1.77	3.92152e-08\\
1.772	3.79179e-08\\
1.774	3.65619e-08\\
1.776	3.51523e-08\\
1.778	3.36941e-08\\
1.78	3.21921e-08\\
1.782	3.06515e-08\\
1.784	2.90769e-08\\
1.786	2.74732e-08\\
1.788	2.58453e-08\\
1.79	2.41977e-08\\
1.792	2.2535e-08\\
1.794	2.11396e-08\\
1.796	2.07036e-08\\
1.798	2.02488e-08\\
1.8	1.97763e-08\\
1.802	1.92872e-08\\
1.804	1.87827e-08\\
1.806	1.82638e-08\\
1.808	1.77316e-08\\
1.81	1.71872e-08\\
1.812	1.66317e-08\\
1.814	1.60661e-08\\
1.816	1.54915e-08\\
1.818	1.50009e-08\\
1.82	1.50475e-08\\
1.822	1.50059e-08\\
1.824	1.48795e-08\\
1.826	1.4672e-08\\
1.828	1.43874e-08\\
1.83	1.40299e-08\\
1.832	1.36041e-08\\
1.834	1.3115e-08\\
1.836	1.25674e-08\\
1.838	1.23718e-08\\
1.84	1.24162e-08\\
1.842	1.31727e-08\\
1.844	1.40595e-08\\
1.846	1.48911e-08\\
1.848	1.5667e-08\\
1.85	1.63871e-08\\
1.852	1.7051e-08\\
1.854	1.76589e-08\\
1.856	1.82109e-08\\
1.858	1.87072e-08\\
1.86	1.91482e-08\\
1.862	1.95345e-08\\
1.864	1.98665e-08\\
1.866	2.01451e-08\\
1.868	2.03711e-08\\
1.87	2.05454e-08\\
1.872	2.06689e-08\\
1.874	2.07429e-08\\
1.876	2.07684e-08\\
1.878	2.07468e-08\\
1.88	2.06793e-08\\
1.882	2.05674e-08\\
1.884	2.04126e-08\\
1.886	2.02163e-08\\
1.888	1.99802e-08\\
1.89	1.97058e-08\\
1.892	1.93949e-08\\
1.894	1.9049e-08\\
1.896	1.867e-08\\
1.898	1.82596e-08\\
1.9	1.78196e-08\\
1.902	1.73517e-08\\
1.904	1.68578e-08\\
1.906	1.63397e-08\\
1.908	1.57993e-08\\
1.91	1.52382e-08\\
1.912	1.46584e-08\\
1.914	1.40617e-08\\
1.916	1.34498e-08\\
1.918	1.28246e-08\\
1.92	1.21878e-08\\
1.922	1.15411e-08\\
1.924	1.08863e-08\\
1.926	1.0225e-08\\
1.928	9.55891e-09\\
1.93	8.88964e-09\\
1.932	8.21875e-09\\
1.934	7.5478e-09\\
1.936	6.87829e-09\\
1.938	6.21167e-09\\
1.94	6.10371e-09\\
1.942	6.00795e-09\\
1.944	5.90187e-09\\
1.946	5.78605e-09\\
1.948	5.66108e-09\\
1.95	5.52758e-09\\
1.952	5.38614e-09\\
1.954	5.23738e-09\\
1.956	5.08191e-09\\
1.958	4.92033e-09\\
1.96	4.75327e-09\\
1.962	4.58133e-09\\
1.964	4.40511e-09\\
1.966	4.22522e-09\\
1.968	4.04222e-09\\
1.97	3.85672e-09\\
1.972	3.66927e-09\\
1.974	3.94201e-09\\
1.976	4.34703e-09\\
1.978	4.73136e-09\\
1.98	5.09474e-09\\
1.982	5.43691e-09\\
1.984	5.75771e-09\\
1.986	6.057e-09\\
1.988	6.33471e-09\\
1.99	6.59082e-09\\
1.992	6.82536e-09\\
1.994	7.03839e-09\\
1.996	7.23005e-09\\
1.998	7.40048e-09\\
2	7.54991e-09\\
};
\addplot [color=mycolor5, line width=1.0pt, forget plot]
  table[row sep=crcr]{%
-0.024	0\\
-0.022	0\\
-0.02	0\\
-0.018	0\\
-0.016	0\\
-0.014	0\\
-0.012	0\\
-0.01	0\\
-0.008	0\\
-0.006	0\\
-0.004	0\\
-0.002	0\\
0	0\\
0.002	0.0391009\\
0.004	0.074092\\
0.006	0.104794\\
0.008	0.131117\\
0.01	0.153053\\
0.012	0.170674\\
0.014	0.18412\\
0.016	0.193597\\
0.018	0.19936\\
0.02	0.201711\\
0.022	0.200986\\
0.024	0.197544\\
0.026	0.191763\\
0.028	0.184025\\
0.03	0.174711\\
0.032	0.164196\\
0.034	0.152836\\
0.036	0.140968\\
0.038	0.128903\\
0.04	0.118709\\
0.042	0.124423\\
0.044	0.1294\\
0.046	0.133574\\
0.048	0.136895\\
0.05	0.139328\\
0.052	0.140854\\
0.054	0.141469\\
0.056	0.141184\\
0.058	0.140023\\
0.06	0.138023\\
0.062	0.135235\\
0.064	0.131717\\
0.066	0.127538\\
0.068	0.122771\\
0.07	0.117499\\
0.072	0.111804\\
0.074	0.105774\\
0.076	0.0994958\\
0.078	0.0930551\\
0.08	0.0865362\\
0.082	0.085057\\
0.084	0.0867084\\
0.086	0.0882446\\
0.088	0.0896667\\
0.09	0.0909754\\
0.092	0.0921717\\
0.094	0.0932559\\
0.096	0.0942287\\
0.098	0.0950901\\
0.1	0.0958403\\
0.102	0.0964791\\
0.104	0.0970062\\
0.106	0.0974212\\
0.108	0.0977235\\
0.11	0.0979127\\
0.112	0.0986842\\
0.114	0.10017\\
0.116	0.101477\\
0.118	0.102606\\
0.12	0.103559\\
0.122	0.10434\\
0.124	0.104953\\
0.126	0.105401\\
0.128	0.105691\\
0.13	0.105827\\
0.132	0.105816\\
0.134	0.105664\\
0.136	0.105378\\
0.138	0.104965\\
0.14	0.104431\\
0.142	0.103784\\
0.144	0.103031\\
0.146	0.102178\\
0.148	0.101233\\
0.15	0.100203\\
0.152	0.0990933\\
0.154	0.0979113\\
0.156	0.096663\\
0.158	0.0953543\\
0.16	0.0939909\\
0.162	0.0925782\\
0.164	0.0911213\\
0.166	0.089625\\
0.168	0.088094\\
0.17	0.0865326\\
0.172	0.0849448\\
0.174	0.0833344\\
0.176	0.0817049\\
0.178	0.0800597\\
0.18	0.0784018\\
0.182	0.0767342\\
0.184	0.0750596\\
0.186	0.0733804\\
0.188	0.071699\\
0.19	0.0700176\\
0.192	0.0683383\\
0.194	0.066663\\
0.196	0.0649935\\
0.198	0.0633315\\
0.2	0.0616786\\
0.202	0.0600364\\
0.204	0.0584062\\
0.206	0.0567894\\
0.208	0.0551873\\
0.21	0.053601\\
0.212	0.0520316\\
0.214	0.0504803\\
0.216	0.048948\\
0.218	0.0474357\\
0.22	0.0459441\\
0.222	0.0444741\\
0.224	0.0430264\\
0.226	0.0416017\\
0.228	0.0402005\\
0.23	0.0388235\\
0.232	0.0374711\\
0.234	0.0361437\\
0.236	0.0348417\\
0.238	0.0335653\\
0.24	0.0323149\\
0.242	0.0310905\\
0.244	0.0298924\\
0.246	0.0287205\\
0.248	0.027575\\
0.25	0.0264558\\
0.252	0.0253628\\
0.254	0.0242959\\
0.256	0.0232551\\
0.258	0.02224\\
0.26	0.0212506\\
0.262	0.0203707\\
0.264	0.0196099\\
0.266	0.0188623\\
0.268	0.0181292\\
0.27	0.0174119\\
0.272	0.0167119\\
0.274	0.0160301\\
0.276	0.0153676\\
0.278	0.0147994\\
0.28	0.0145497\\
0.282	0.0143057\\
0.284	0.0140671\\
0.286	0.013834\\
0.288	0.0136061\\
0.29	0.0133834\\
0.292	0.0131657\\
0.294	0.012953\\
0.296	0.0127449\\
0.298	0.0125415\\
0.3	0.0123426\\
0.302	0.012148\\
0.304	0.0119575\\
0.306	0.0117711\\
0.308	0.0115885\\
0.31	0.0114096\\
0.312	0.0112343\\
0.314	0.0110624\\
0.316	0.0108938\\
0.318	0.0107283\\
0.32	0.0105658\\
0.322	0.0104061\\
0.324	0.0102492\\
0.326	0.0100948\\
0.328	0.00994282\\
0.33	0.0097932\\
0.332	0.00964578\\
0.334	0.00950046\\
0.336	0.00935713\\
0.338	0.00921569\\
0.34	0.00907605\\
0.342	0.00918111\\
0.344	0.00934646\\
0.346	0.00949713\\
0.348	0.00963295\\
0.35	0.0097538\\
0.352	0.00985962\\
0.354	0.00995037\\
0.356	0.0100261\\
0.358	0.0100868\\
0.36	0.0101326\\
0.362	0.0101637\\
0.364	0.0101803\\
0.366	0.0101825\\
0.368	0.0101705\\
0.37	0.0101448\\
0.372	0.0101054\\
0.374	0.0100529\\
0.376	0.00998746\\
0.378	0.00990952\\
0.38	0.00981947\\
0.382	0.00971771\\
0.384	0.00960466\\
0.386	0.00948076\\
0.388	0.00934647\\
0.39	0.00920224\\
0.392	0.00904857\\
0.394	0.00888592\\
0.396	0.00871479\\
0.398	0.00853568\\
0.4	0.00834908\\
0.402	0.00815552\\
0.404	0.00795549\\
0.406	0.0077495\\
0.408	0.00753807\\
0.41	0.0073217\\
0.412	0.0071009\\
0.414	0.00687618\\
0.416	0.00664802\\
0.418	0.00641694\\
0.42	0.0061834\\
0.422	0.00594791\\
0.424	0.00571092\\
0.426	0.00547291\\
0.428	0.00523432\\
0.43	0.00499561\\
0.432	0.00475721\\
0.434	0.00451953\\
0.436	0.00428299\\
0.438	0.00404799\\
0.44	0.00381491\\
0.442	0.00358412\\
0.444	0.00343861\\
0.446	0.00336078\\
0.448	0.00328422\\
0.45	0.00320893\\
0.452	0.0031349\\
0.454	0.00306214\\
0.456	0.00299063\\
0.458	0.00292038\\
0.46	0.00285138\\
0.462	0.00278362\\
0.464	0.0027171\\
0.466	0.00265181\\
0.468	0.00258774\\
0.47	0.00252488\\
0.472	0.00246323\\
0.474	0.00240277\\
0.476	0.00234349\\
0.478	0.00228539\\
0.48	0.00231947\\
0.482	0.00237645\\
0.484	0.00243005\\
0.486	0.00248022\\
0.488	0.00252696\\
0.49	0.00257024\\
0.492	0.00261004\\
0.494	0.00264638\\
0.496	0.00267925\\
0.498	0.00270866\\
0.5	0.00273462\\
0.502	0.00275717\\
0.504	0.00277632\\
0.506	0.00279211\\
0.508	0.00280457\\
0.51	0.00281376\\
0.512	0.00281971\\
0.514	0.00282248\\
0.516	0.00282213\\
0.518	0.00281873\\
0.52	0.00281232\\
0.522	0.002803\\
0.524	0.00279081\\
0.526	0.00277586\\
0.528	0.0027582\\
0.53	0.00273793\\
0.532	0.00271512\\
0.534	0.00268987\\
0.536	0.00266227\\
0.538	0.00263239\\
0.54	0.00260034\\
0.542	0.00256622\\
0.544	0.0025301\\
0.546	0.0024921\\
0.548	0.00245231\\
0.55	0.00241081\\
0.552	0.00236773\\
0.554	0.00232314\\
0.556	0.00227715\\
0.558	0.00222986\\
0.56	0.00218136\\
0.562	0.00213175\\
0.564	0.00208113\\
0.566	0.00202959\\
0.568	0.00197723\\
0.57	0.00192415\\
0.572	0.00187043\\
0.574	0.00181616\\
0.576	0.00176144\\
0.578	0.00170636\\
0.58	0.00165099\\
0.582	0.00159543\\
0.584	0.00153976\\
0.586	0.00148405\\
0.588	0.00142839\\
0.59	0.00137285\\
0.592	0.00131751\\
0.594	0.00126244\\
0.596	0.00120771\\
0.598	0.00115338\\
0.6	0.00109952\\
0.602	0.00104619\\
0.604	0.000993456\\
0.606	0.000941366\\
0.608	0.000889977\\
0.61	0.00083934\\
0.612	0.000789505\\
0.614	0.000740517\\
0.616	0.00069242\\
0.618	0.000645255\\
0.62	0.000617698\\
0.622	0.00065295\\
0.624	0.000686404\\
0.626	0.000718034\\
0.628	0.000747819\\
0.63	0.000775741\\
0.632	0.000801789\\
0.634	0.000825953\\
0.636	0.000848229\\
0.638	0.000868618\\
0.64	0.000887124\\
0.642	0.000903754\\
0.644	0.000918521\\
0.646	0.00093144\\
0.648	0.00094253\\
0.65	0.000951812\\
0.652	0.000959312\\
0.654	0.000965059\\
0.656	0.000969084\\
0.658	0.000971419\\
0.66	0.000972103\\
0.662	0.000971174\\
0.664	0.000968673\\
0.666	0.000964644\\
0.668	0.000959132\\
0.67	0.000952183\\
0.672	0.000943848\\
0.674	0.000934176\\
0.676	0.000923219\\
0.678	0.000911029\\
0.68	0.000897662\\
0.682	0.000883172\\
0.684	0.000867614\\
0.686	0.000851046\\
0.688	0.000833525\\
0.69	0.000815107\\
0.692	0.00079585\\
0.694	0.000775813\\
0.696	0.000755053\\
0.698	0.000733628\\
0.7	0.000711595\\
0.702	0.000689013\\
0.704	0.000665937\\
0.706	0.000642423\\
0.708	0.000618529\\
0.71	0.000594308\\
0.712	0.000569814\\
0.714	0.000545102\\
0.716	0.000520222\\
0.718	0.000495226\\
0.72	0.000470165\\
0.722	0.000445086\\
0.724	0.000420037\\
0.726	0.000395065\\
0.728	0.000370214\\
0.73	0.000345527\\
0.732	0.000321046\\
0.734	0.000296811\\
0.736	0.000272861\\
0.738	0.000249232\\
0.74	0.00022596\\
0.742	0.000211594\\
0.744	0.000202998\\
0.746	0.00019423\\
0.748	0.000185313\\
0.75	0.00017627\\
0.752	0.000167122\\
0.754	0.000157891\\
0.756	0.000148597\\
0.758	0.00013926\\
0.76	0.0001299\\
0.762	0.000120535\\
0.764	0.000111183\\
0.766	0.00010186\\
0.768	9.25851e-05\\
0.77	8.36061e-05\\
0.772	8.23993e-05\\
0.774	8.92663e-05\\
0.776	0.000102059\\
0.778	0.000114152\\
0.78	0.000125546\\
0.782	0.000136239\\
0.784	0.000146233\\
0.786	0.000155531\\
0.788	0.000164138\\
0.79	0.000172058\\
0.792	0.000179298\\
0.794	0.000185866\\
0.796	0.000191771\\
0.798	0.000197023\\
0.8	0.000201633\\
0.802	0.000205613\\
0.804	0.000208977\\
0.806	0.000211738\\
0.808	0.00021391\\
0.81	0.000215509\\
0.812	0.000216552\\
0.814	0.000217055\\
0.816	0.000217036\\
0.818	0.000216512\\
0.82	0.000215502\\
0.822	0.000214024\\
0.824	0.000212099\\
0.826	0.000209745\\
0.828	0.000206983\\
0.83	0.000203832\\
0.832	0.000200313\\
0.834	0.000196446\\
0.836	0.000192252\\
0.838	0.000187751\\
0.84	0.000182963\\
0.842	0.00017791\\
0.844	0.000172611\\
0.846	0.000167086\\
0.848	0.000161356\\
0.85	0.000155441\\
0.852	0.000149359\\
0.854	0.00014313\\
0.856	0.000136774\\
0.858	0.000131794\\
0.86	0.00013159\\
0.862	0.000131197\\
0.864	0.000130622\\
0.866	0.00012987\\
0.868	0.000128948\\
0.87	0.00012786\\
0.872	0.000126613\\
0.874	0.000125214\\
0.876	0.000123669\\
0.878	0.000121984\\
0.88	0.000120165\\
0.882	0.000118219\\
0.884	0.000116152\\
0.886	0.000113972\\
0.888	0.000111684\\
0.89	0.000109295\\
0.892	0.000106811\\
0.894	0.00010424\\
0.896	0.000101587\\
0.898	9.88583e-05\\
0.9	9.60613e-05\\
0.902	9.32019e-05\\
0.904	9.02864e-05\\
0.906	8.7321e-05\\
0.908	8.43117e-05\\
0.91	8.12646e-05\\
0.912	7.81856e-05\\
0.914	7.50806e-05\\
0.916	7.19552e-05\\
0.918	6.88151e-05\\
0.92	6.56658e-05\\
0.922	6.25125e-05\\
0.924	5.93606e-05\\
0.926	5.6215e-05\\
0.928	5.76335e-05\\
0.93	6.00516e-05\\
0.932	6.2269e-05\\
0.934	6.42875e-05\\
0.936	6.61092e-05\\
0.938	6.77366e-05\\
0.94	6.91727e-05\\
0.942	7.04205e-05\\
0.944	7.14837e-05\\
0.946	7.2366e-05\\
0.948	7.30715e-05\\
0.95	7.36046e-05\\
0.952	7.39698e-05\\
0.954	7.41719e-05\\
0.956	7.42159e-05\\
0.958	7.41069e-05\\
0.96	7.38504e-05\\
0.962	7.34518e-05\\
0.964	7.29167e-05\\
0.966	7.22508e-05\\
0.968	7.146e-05\\
0.97	7.05502e-05\\
0.972	6.95273e-05\\
0.974	6.83973e-05\\
0.976	6.71662e-05\\
0.978	6.58402e-05\\
0.98	6.44252e-05\\
0.982	6.29274e-05\\
0.984	6.13527e-05\\
0.986	5.97071e-05\\
0.988	5.79966e-05\\
0.99	5.6227e-05\\
0.992	5.44041e-05\\
0.994	5.25337e-05\\
0.996	5.06213e-05\\
0.998	4.86725e-05\\
1	4.66927e-05\\
1.002	4.46871e-05\\
1.004	4.2661e-05\\
1.006	4.06194e-05\\
1.008	3.85672e-05\\
1.01	3.6509e-05\\
1.012	3.44496e-05\\
1.014	3.23932e-05\\
1.016	3.03443e-05\\
1.018	2.87106e-05\\
1.02	2.88651e-05\\
1.022	2.89586e-05\\
1.024	2.8993e-05\\
1.026	2.89701e-05\\
1.028	2.8892e-05\\
1.03	2.87608e-05\\
1.032	2.85785e-05\\
1.034	2.83472e-05\\
1.036	2.80691e-05\\
1.038	2.77462e-05\\
1.04	2.73808e-05\\
1.042	2.6975e-05\\
1.044	2.6531e-05\\
1.046	2.6051e-05\\
1.048	2.55371e-05\\
1.05	2.49915e-05\\
1.052	2.44164e-05\\
1.054	2.38138e-05\\
1.056	2.31859e-05\\
1.058	2.25348e-05\\
1.06	2.18625e-05\\
1.062	2.11711e-05\\
1.064	2.04625e-05\\
1.066	1.97388e-05\\
1.068	1.90018e-05\\
1.07	1.82535e-05\\
1.072	1.74956e-05\\
1.074	1.673e-05\\
1.076	1.59584e-05\\
1.078	1.51824e-05\\
1.08	1.47811e-05\\
1.082	1.53066e-05\\
1.084	1.5778e-05\\
1.086	1.61962e-05\\
1.088	1.65624e-05\\
1.09	1.68776e-05\\
1.092	1.71431e-05\\
1.094	1.73601e-05\\
1.096	1.753e-05\\
1.098	1.7654e-05\\
1.1	1.77338e-05\\
1.102	1.77706e-05\\
1.104	1.7766e-05\\
1.106	1.77217e-05\\
1.108	1.76391e-05\\
1.11	1.75198e-05\\
1.112	1.73656e-05\\
1.114	1.71779e-05\\
1.116	1.69586e-05\\
1.118	1.67092e-05\\
1.12	1.64315e-05\\
1.122	1.6127e-05\\
1.124	1.57975e-05\\
1.126	1.54447e-05\\
1.128	1.50701e-05\\
1.13	1.46755e-05\\
1.132	1.42624e-05\\
1.134	1.38325e-05\\
1.136	1.33873e-05\\
1.138	1.29284e-05\\
1.14	1.24573e-05\\
1.142	1.19756e-05\\
1.144	1.14847e-05\\
1.146	1.09861e-05\\
1.148	1.04811e-05\\
1.15	9.9711e-06\\
1.152	9.45746e-06\\
1.154	8.94144e-06\\
1.156	8.4243e-06\\
1.158	7.90722e-06\\
1.16	7.39136e-06\\
1.162	7.17222e-06\\
1.164	7.31357e-06\\
1.166	7.43722e-06\\
1.168	7.54358e-06\\
1.17	7.63306e-06\\
1.172	7.70608e-06\\
1.174	7.7631e-06\\
1.176	7.80455e-06\\
1.178	7.83093e-06\\
1.18	7.8427e-06\\
1.182	7.84036e-06\\
1.184	7.82441e-06\\
1.186	7.79534e-06\\
1.188	7.75369e-06\\
1.19	7.69995e-06\\
1.192	7.63466e-06\\
1.194	7.55833e-06\\
1.196	7.47149e-06\\
1.198	7.37467e-06\\
1.2	7.26838e-06\\
1.202	7.15316e-06\\
1.204	7.02953e-06\\
1.206	6.89799e-06\\
1.208	6.75907e-06\\
1.21	6.61327e-06\\
1.212	6.4611e-06\\
1.214	6.30305e-06\\
1.216	6.13962e-06\\
1.218	5.97128e-06\\
1.22	5.79852e-06\\
1.222	5.62178e-06\\
1.224	5.44154e-06\\
1.226	5.25824e-06\\
1.228	5.07231e-06\\
1.23	4.88418e-06\\
1.232	4.69426e-06\\
1.234	4.51474e-06\\
1.236	4.57781e-06\\
1.238	4.62849e-06\\
1.24	4.66712e-06\\
1.242	4.69407e-06\\
1.244	4.70971e-06\\
1.246	4.71444e-06\\
1.248	4.70865e-06\\
1.25	4.69275e-06\\
1.252	4.66716e-06\\
1.254	4.63229e-06\\
1.256	4.58859e-06\\
1.258	4.53647e-06\\
1.26	4.47638e-06\\
1.262	4.40876e-06\\
1.264	4.33403e-06\\
1.266	4.25264e-06\\
1.268	4.16503e-06\\
1.27	4.07162e-06\\
1.272	3.97286e-06\\
1.274	3.86917e-06\\
1.276	3.76096e-06\\
1.278	3.64867e-06\\
1.28	3.5327e-06\\
1.282	3.41345e-06\\
1.284	3.29132e-06\\
1.286	3.16669e-06\\
1.288	3.03996e-06\\
1.29	2.91148e-06\\
1.292	2.78161e-06\\
1.294	2.65071e-06\\
1.296	2.51912e-06\\
1.298	2.38715e-06\\
1.3	2.25513e-06\\
1.302	2.12335e-06\\
1.304	1.99212e-06\\
1.306	1.86171e-06\\
1.308	1.73238e-06\\
1.31	1.60439e-06\\
1.312	1.47798e-06\\
1.314	1.35338e-06\\
1.316	1.26437e-06\\
1.318	1.3152e-06\\
1.32	1.36158e-06\\
1.322	1.40357e-06\\
1.324	1.44126e-06\\
1.326	1.47474e-06\\
1.328	1.50411e-06\\
1.33	1.52946e-06\\
1.332	1.5509e-06\\
1.334	1.56852e-06\\
1.336	1.58245e-06\\
1.338	1.59278e-06\\
1.34	1.59964e-06\\
1.342	1.60315e-06\\
1.344	1.60342e-06\\
1.346	1.60057e-06\\
1.348	1.59473e-06\\
1.35	1.58602e-06\\
1.352	1.57457e-06\\
1.354	1.5605e-06\\
1.356	1.54393e-06\\
1.358	1.525e-06\\
1.36	1.50384e-06\\
1.362	1.48056e-06\\
1.364	1.45529e-06\\
1.366	1.42815e-06\\
1.368	1.39928e-06\\
1.37	1.36878e-06\\
1.372	1.33679e-06\\
1.374	1.30342e-06\\
1.376	1.26879e-06\\
1.378	1.23302e-06\\
1.38	1.19621e-06\\
1.382	1.15848e-06\\
1.384	1.11994e-06\\
1.386	1.0807e-06\\
1.388	1.04085e-06\\
1.39	1.0005e-06\\
1.392	9.94204e-07\\
1.394	9.89655e-07\\
1.396	9.82835e-07\\
1.398	9.73856e-07\\
1.4	9.62828e-07\\
1.402	9.49863e-07\\
1.404	9.35073e-07\\
1.406	9.18568e-07\\
1.408	9.00461e-07\\
1.41	8.80861e-07\\
1.412	8.5988e-07\\
1.414	8.37625e-07\\
1.416	8.14204e-07\\
1.418	7.89723e-07\\
1.42	7.64287e-07\\
1.422	7.37998e-07\\
1.424	7.10956e-07\\
1.426	6.8326e-07\\
1.428	6.55005e-07\\
1.43	6.26286e-07\\
1.432	5.97192e-07\\
1.434	5.67812e-07\\
1.436	5.38231e-07\\
1.438	5.08531e-07\\
1.44	4.78792e-07\\
1.442	4.4909e-07\\
1.444	4.19497e-07\\
1.446	3.90084e-07\\
1.448	3.74753e-07\\
1.45	3.73574e-07\\
1.452	3.71753e-07\\
1.454	3.69317e-07\\
1.456	3.66288e-07\\
1.458	3.62694e-07\\
1.46	3.58558e-07\\
1.462	3.53907e-07\\
1.464	3.48765e-07\\
1.466	3.43159e-07\\
1.468	3.37112e-07\\
1.47	3.30652e-07\\
1.472	3.23801e-07\\
1.474	3.16587e-07\\
1.476	3.09032e-07\\
1.478	3.01162e-07\\
1.48	3.08732e-07\\
1.482	3.1665e-07\\
1.484	3.23674e-07\\
1.486	3.29825e-07\\
1.488	3.35126e-07\\
1.49	3.39601e-07\\
1.492	3.43273e-07\\
1.494	3.46167e-07\\
1.496	3.48308e-07\\
1.498	3.49722e-07\\
1.5	3.50435e-07\\
1.502	3.50473e-07\\
1.504	3.49862e-07\\
1.506	3.48629e-07\\
1.508	3.46801e-07\\
1.51	3.44406e-07\\
1.512	3.41469e-07\\
1.514	3.38018e-07\\
1.516	3.3408e-07\\
1.518	3.29681e-07\\
1.52	3.24848e-07\\
1.522	3.19607e-07\\
1.524	3.13983e-07\\
1.526	3.08004e-07\\
1.528	3.01693e-07\\
1.53	2.95076e-07\\
1.532	2.88178e-07\\
1.534	2.81022e-07\\
1.536	2.73632e-07\\
1.538	2.66031e-07\\
1.54	2.58241e-07\\
1.542	2.50285e-07\\
1.544	2.42184e-07\\
1.546	2.33959e-07\\
1.548	2.25629e-07\\
1.55	2.17216e-07\\
1.552	2.08737e-07\\
1.554	2.00211e-07\\
1.556	1.91656e-07\\
1.558	1.83089e-07\\
1.56	1.74526e-07\\
1.562	1.65983e-07\\
1.564	1.57476e-07\\
1.566	1.49017e-07\\
1.568	1.40622e-07\\
1.57	1.32303e-07\\
1.572	1.24072e-07\\
1.574	1.15942e-07\\
1.576	1.07923e-07\\
1.578	1.00027e-07\\
1.58	9.22616e-08\\
1.582	8.91324e-08\\
1.584	9.12694e-08\\
1.586	9.31596e-08\\
1.588	9.48074e-08\\
1.59	9.62178e-08\\
1.592	9.73957e-08\\
1.594	9.83468e-08\\
1.596	9.90767e-08\\
1.598	9.95914e-08\\
1.6	9.98972e-08\\
1.602	1e-07\\
1.604	9.99079e-08\\
1.606	9.96264e-08\\
1.608	9.91629e-08\\
1.61	9.85246e-08\\
1.612	9.77188e-08\\
1.614	9.67528e-08\\
1.616	9.56341e-08\\
1.618	9.43702e-08\\
1.62	9.29687e-08\\
1.622	9.14372e-08\\
1.624	8.97833e-08\\
1.626	8.80147e-08\\
1.628	8.61389e-08\\
1.63	8.41634e-08\\
1.632	8.20958e-08\\
1.634	7.99436e-08\\
1.636	7.7714e-08\\
1.638	7.54143e-08\\
1.64	7.30517e-08\\
1.642	7.06333e-08\\
1.644	6.81658e-08\\
1.646	6.56562e-08\\
1.648	6.31109e-08\\
1.65	6.15365e-08\\
1.652	6.25028e-08\\
1.654	6.33137e-08\\
1.656	6.39742e-08\\
1.658	6.44891e-08\\
1.66	6.48635e-08\\
1.662	6.51024e-08\\
1.664	6.52109e-08\\
1.666	6.51941e-08\\
1.668	6.50573e-08\\
1.67	6.48056e-08\\
1.672	6.44442e-08\\
1.674	6.39784e-08\\
1.676	6.34133e-08\\
1.678	6.27541e-08\\
1.68	6.20059e-08\\
1.682	6.11739e-08\\
1.684	6.02631e-08\\
1.686	5.92784e-08\\
1.688	5.82249e-08\\
1.69	5.71074e-08\\
1.692	5.59307e-08\\
1.694	5.46994e-08\\
1.696	5.34182e-08\\
1.698	5.20916e-08\\
1.7	5.07241e-08\\
1.702	4.93198e-08\\
1.704	4.78831e-08\\
1.706	4.64179e-08\\
1.708	4.49284e-08\\
1.71	4.34182e-08\\
1.712	4.18912e-08\\
1.714	4.0351e-08\\
1.716	3.88009e-08\\
1.718	3.72445e-08\\
1.72	3.56848e-08\\
1.722	3.41249e-08\\
1.724	3.25678e-08\\
1.726	3.10164e-08\\
1.728	2.94732e-08\\
1.73	2.79409e-08\\
1.732	2.64218e-08\\
1.734	2.49182e-08\\
1.736	2.34324e-08\\
1.738	2.19662e-08\\
1.74	2.05217e-08\\
1.742	2.04339e-08\\
1.744	2.07801e-08\\
1.746	2.10717e-08\\
1.748	2.13101e-08\\
1.75	2.14971e-08\\
1.752	2.16343e-08\\
1.754	2.17235e-08\\
1.756	2.17664e-08\\
1.758	2.17648e-08\\
1.76	2.17205e-08\\
1.762	2.16352e-08\\
1.764	2.15109e-08\\
1.766	2.13493e-08\\
1.768	2.11521e-08\\
1.77	2.09214e-08\\
1.772	2.09419e-08\\
1.774	2.09328e-08\\
1.776	2.08647e-08\\
1.778	2.07396e-08\\
1.78	2.05597e-08\\
1.782	2.03273e-08\\
1.784	2.00445e-08\\
1.786	1.97137e-08\\
1.788	1.93373e-08\\
1.79	1.89178e-08\\
1.792	1.84576e-08\\
1.794	1.79592e-08\\
1.796	1.74252e-08\\
1.798	1.68582e-08\\
1.8	1.62606e-08\\
1.802	1.5635e-08\\
1.804	1.49841e-08\\
1.806	1.43103e-08\\
1.808	1.36161e-08\\
1.81	1.29041e-08\\
1.812	1.21766e-08\\
1.814	1.14362e-08\\
1.816	1.06852e-08\\
1.818	1.05839e-08\\
1.82	1.07453e-08\\
1.822	1.08826e-08\\
1.824	1.09967e-08\\
1.826	1.10884e-08\\
1.828	1.11586e-08\\
1.83	1.1208e-08\\
1.832	1.12375e-08\\
1.834	1.12478e-08\\
1.836	1.12399e-08\\
1.838	1.12146e-08\\
1.84	1.11725e-08\\
1.842	1.11146e-08\\
1.844	1.10415e-08\\
1.846	1.09541e-08\\
1.848	1.08532e-08\\
1.85	1.07394e-08\\
1.852	1.06135e-08\\
1.854	1.04762e-08\\
1.856	1.03283e-08\\
1.858	1.01703e-08\\
1.86	1.00031e-08\\
1.862	9.82716e-09\\
1.864	9.6432e-09\\
1.866	9.45183e-09\\
1.868	9.25367e-09\\
1.87	9.04928e-09\\
1.872	8.83926e-09\\
1.874	8.62414e-09\\
1.876	8.40446e-09\\
1.878	8.18074e-09\\
1.88	8.20507e-09\\
1.882	8.42943e-09\\
1.884	8.6272e-09\\
1.886	8.79865e-09\\
1.888	8.94408e-09\\
1.89	9.06387e-09\\
1.892	9.15843e-09\\
1.894	9.22824e-09\\
1.896	9.27382e-09\\
1.898	9.29573e-09\\
1.9	9.29459e-09\\
1.902	9.27104e-09\\
1.904	9.22576e-09\\
1.906	9.15948e-09\\
1.908	9.07294e-09\\
1.91	8.96692e-09\\
1.912	8.84222e-09\\
1.914	8.69966e-09\\
1.916	8.54009e-09\\
1.918	8.36437e-09\\
1.92	8.17337e-09\\
1.922	7.96798e-09\\
1.924	7.74909e-09\\
1.926	7.51761e-09\\
1.928	7.27444e-09\\
1.93	7.02048e-09\\
1.932	6.75666e-09\\
1.934	6.48387e-09\\
1.936	6.20301e-09\\
1.938	5.91497e-09\\
1.94	5.62065e-09\\
1.942	5.3209e-09\\
1.944	5.0166e-09\\
1.946	4.70859e-09\\
1.948	4.3977e-09\\
1.95	4.08475e-09\\
1.952	3.77053e-09\\
1.954	3.45581e-09\\
1.956	3.14136e-09\\
1.958	2.87082e-09\\
1.96	2.76617e-09\\
1.962	2.6605e-09\\
1.964	2.55409e-09\\
1.966	2.44717e-09\\
1.968	2.34e-09\\
1.97	2.23279e-09\\
1.972	2.12577e-09\\
1.974	2.01915e-09\\
1.976	1.91313e-09\\
1.978	1.80789e-09\\
1.98	1.70363e-09\\
1.982	1.60051e-09\\
1.984	1.49868e-09\\
1.986	1.47731e-09\\
1.988	1.54158e-09\\
1.99	1.60143e-09\\
1.992	1.76719e-09\\
1.994	1.9681e-09\\
1.996	2.15931e-09\\
1.998	2.34065e-09\\
2	2.51197e-09\\
};
\addplot [color=mycolor6, line width=1.0pt, forget plot]
  table[row sep=crcr]{%
-0.024	0\\
-0.022	0\\
-0.02	0\\
-0.018	0\\
-0.016	0\\
-0.014	0\\
-0.012	0\\
-0.01	0\\
-0.008	0\\
-0.006	0\\
-0.004	0\\
-0.002	0\\
0	0\\
0.002	0.0391009\\
0.004	0.074092\\
0.006	0.104794\\
0.008	0.131117\\
0.01	0.153053\\
0.012	0.170674\\
0.014	0.18412\\
0.016	0.193597\\
0.018	0.19936\\
0.02	0.201711\\
0.022	0.200986\\
0.024	0.197545\\
0.026	0.191763\\
0.028	0.184025\\
0.03	0.174712\\
0.032	0.164196\\
0.034	0.152836\\
0.036	0.140969\\
0.038	0.128903\\
0.04	0.118709\\
0.042	0.124423\\
0.044	0.129399\\
0.046	0.133573\\
0.048	0.136893\\
0.05	0.139326\\
0.052	0.140851\\
0.054	0.141464\\
0.056	0.141177\\
0.058	0.140014\\
0.06	0.138013\\
0.062	0.135223\\
0.064	0.131702\\
0.066	0.127519\\
0.068	0.122749\\
0.07	0.117473\\
0.072	0.111774\\
0.074	0.105739\\
0.076	0.0994555\\
0.078	0.0930088\\
0.08	0.0864834\\
0.082	0.0840929\\
0.084	0.0857004\\
0.086	0.0871937\\
0.088	0.0885738\\
0.09	0.0898419\\
0.092	0.0909989\\
0.094	0.0920455\\
0.096	0.0929825\\
0.098	0.0938102\\
0.1	0.0945287\\
0.102	0.0951382\\
0.104	0.0956385\\
0.106	0.0960294\\
0.108	0.0963103\\
0.11	0.0964811\\
0.112	0.0974443\\
0.114	0.0989132\\
0.116	0.100205\\
0.118	0.101322\\
0.12	0.102265\\
0.122	0.103038\\
0.124	0.103646\\
0.126	0.104092\\
0.128	0.104381\\
0.13	0.10452\\
0.132	0.104514\\
0.134	0.104369\\
0.136	0.104093\\
0.138	0.103692\\
0.14	0.103173\\
0.142	0.102543\\
0.144	0.101809\\
0.146	0.100978\\
0.148	0.100057\\
0.15	0.0990528\\
0.152	0.0979716\\
0.154	0.0968198\\
0.156	0.0956036\\
0.158	0.0943288\\
0.16	0.0930009\\
0.162	0.0916253\\
0.164	0.090207\\
0.166	0.0887508\\
0.168	0.087261\\
0.17	0.085742\\
0.172	0.0841976\\
0.174	0.0826315\\
0.176	0.0810471\\
0.178	0.0794477\\
0.18	0.0778362\\
0.182	0.0762153\\
0.184	0.0745876\\
0.186	0.0729556\\
0.188	0.0713215\\
0.19	0.0696874\\
0.192	0.0680551\\
0.194	0.0664266\\
0.196	0.0648035\\
0.198	0.0631874\\
0.2	0.0615799\\
0.202	0.0599824\\
0.204	0.0583961\\
0.206	0.0568225\\
0.208	0.0552625\\
0.21	0.0537174\\
0.212	0.0521882\\
0.214	0.0506759\\
0.216	0.0491814\\
0.218	0.0477056\\
0.22	0.0462494\\
0.222	0.0448134\\
0.224	0.0433983\\
0.226	0.0420048\\
0.228	0.0406335\\
0.23	0.0392849\\
0.232	0.0379594\\
0.234	0.0366575\\
0.236	0.0353796\\
0.238	0.0341258\\
0.24	0.0328965\\
0.242	0.0316918\\
0.244	0.0305119\\
0.246	0.0293569\\
0.248	0.0282268\\
0.25	0.0271216\\
0.252	0.0260413\\
0.254	0.0249859\\
0.256	0.0239552\\
0.258	0.022949\\
0.26	0.0219673\\
0.262	0.0210097\\
0.264	0.0200762\\
0.266	0.0191665\\
0.268	0.0183739\\
0.27	0.0176793\\
0.272	0.017001\\
0.274	0.01634\\
0.276	0.0156973\\
0.278	0.0150739\\
0.28	0.0146117\\
0.282	0.0143641\\
0.284	0.0141217\\
0.286	0.0138844\\
0.288	0.0136521\\
0.29	0.0134247\\
0.292	0.013202\\
0.294	0.012984\\
0.296	0.0127706\\
0.298	0.0125616\\
0.3	0.0123569\\
0.302	0.0121564\\
0.304	0.01196\\
0.306	0.0117674\\
0.308	0.0115787\\
0.31	0.0113937\\
0.312	0.0112122\\
0.314	0.0110341\\
0.316	0.0108594\\
0.318	0.0106878\\
0.32	0.0105193\\
0.322	0.0103537\\
0.324	0.0101909\\
0.326	0.0100308\\
0.328	0.00987339\\
0.33	0.00971844\\
0.332	0.00956588\\
0.334	0.00941563\\
0.336	0.0092676\\
0.338	0.00912168\\
0.34	0.00907669\\
0.342	0.00924031\\
0.344	0.00939057\\
0.346	0.00952724\\
0.348	0.00965017\\
0.35	0.00975923\\
0.352	0.00985434\\
0.354	0.00993545\\
0.356	0.0100026\\
0.358	0.0100558\\
0.36	0.0100951\\
0.362	0.0101207\\
0.364	0.0101327\\
0.366	0.0101313\\
0.368	0.0101167\\
0.37	0.0100892\\
0.372	0.0100489\\
0.374	0.00999626\\
0.376	0.00993152\\
0.378	0.00985501\\
0.38	0.00976709\\
0.382	0.00966813\\
0.384	0.00955851\\
0.386	0.00943863\\
0.388	0.00930891\\
0.39	0.00916977\\
0.392	0.00902165\\
0.394	0.00886498\\
0.396	0.00870021\\
0.398	0.00852781\\
0.4	0.00834824\\
0.402	0.00816196\\
0.404	0.00796945\\
0.406	0.00777117\\
0.408	0.0075676\\
0.41	0.00735922\\
0.412	0.00714648\\
0.414	0.00692987\\
0.416	0.00670985\\
0.418	0.00648688\\
0.42	0.00626142\\
0.422	0.00603392\\
0.424	0.00580483\\
0.426	0.00557458\\
0.428	0.00534361\\
0.43	0.00511234\\
0.432	0.00488117\\
0.434	0.00465052\\
0.436	0.00442077\\
0.438	0.0041923\\
0.44	0.00396548\\
0.442	0.00374066\\
0.444	0.0035182\\
0.446	0.00339363\\
0.448	0.00332205\\
0.45	0.00325164\\
0.452	0.0031824\\
0.454	0.00311431\\
0.456	0.00304736\\
0.458	0.00298155\\
0.46	0.00291686\\
0.462	0.00285328\\
0.464	0.00279081\\
0.466	0.00272944\\
0.468	0.00266914\\
0.47	0.00260992\\
0.472	0.00255176\\
0.474	0.00249465\\
0.476	0.00243857\\
0.478	0.00238352\\
0.48	0.00242315\\
0.482	0.00247519\\
0.484	0.00252396\\
0.486	0.00256943\\
0.488	0.00261158\\
0.49	0.0026504\\
0.492	0.00268587\\
0.494	0.00271801\\
0.496	0.00274682\\
0.498	0.0027723\\
0.5	0.00279447\\
0.502	0.00281336\\
0.504	0.002829\\
0.506	0.00284142\\
0.508	0.00285066\\
0.51	0.00285675\\
0.512	0.00285976\\
0.514	0.00285974\\
0.516	0.00285673\\
0.518	0.00285081\\
0.52	0.00284203\\
0.522	0.00283046\\
0.524	0.00281618\\
0.526	0.00279925\\
0.528	0.00277977\\
0.53	0.00275779\\
0.532	0.00273342\\
0.534	0.00270672\\
0.536	0.0026778\\
0.538	0.00264672\\
0.54	0.0026136\\
0.542	0.0025785\\
0.544	0.00254154\\
0.546	0.00250279\\
0.548	0.00246235\\
0.55	0.00242032\\
0.552	0.00237679\\
0.554	0.00233185\\
0.556	0.0022856\\
0.558	0.00223814\\
0.56	0.00218955\\
0.562	0.00213993\\
0.564	0.00208937\\
0.566	0.00203796\\
0.568	0.0019858\\
0.57	0.00193297\\
0.572	0.00187956\\
0.574	0.00182567\\
0.576	0.00177137\\
0.578	0.00171675\\
0.58	0.00166189\\
0.582	0.00160688\\
0.584	0.00155178\\
0.586	0.00149669\\
0.588	0.00144167\\
0.59	0.0013868\\
0.592	0.00133214\\
0.594	0.00127777\\
0.596	0.00122376\\
0.598	0.00117015\\
0.6	0.00111703\\
0.602	0.00106444\\
0.604	0.00101245\\
0.606	0.000961101\\
0.608	0.000910452\\
0.61	0.000860551\\
0.612	0.000811443\\
0.614	0.000763174\\
0.616	0.000715784\\
0.618	0.000669311\\
0.62	0.000623794\\
0.622	0.000648111\\
0.624	0.000681232\\
0.626	0.000712613\\
0.628	0.000742228\\
0.63	0.000770057\\
0.632	0.000796083\\
0.634	0.000820295\\
0.636	0.000842685\\
0.638	0.00086325\\
0.64	0.000881989\\
0.642	0.000898908\\
0.644	0.000914015\\
0.646	0.000927321\\
0.648	0.000938843\\
0.65	0.000948598\\
0.652	0.000956609\\
0.654	0.0009629\\
0.656	0.0009675\\
0.658	0.000970441\\
0.66	0.000971754\\
0.662	0.000971476\\
0.664	0.000969646\\
0.666	0.000966304\\
0.668	0.000961493\\
0.67	0.000955257\\
0.672	0.000947643\\
0.674	0.000938699\\
0.676	0.000928474\\
0.678	0.000917019\\
0.68	0.000904387\\
0.682	0.000890629\\
0.684	0.000875801\\
0.686	0.000859956\\
0.688	0.00084315\\
0.69	0.00082544\\
0.692	0.00080688\\
0.694	0.000787529\\
0.696	0.000767442\\
0.698	0.000746676\\
0.7	0.000725287\\
0.702	0.000703332\\
0.704	0.000680866\\
0.706	0.000657946\\
0.708	0.000634625\\
0.71	0.000610958\\
0.712	0.000587\\
0.714	0.000562801\\
0.716	0.000538414\\
0.718	0.00051389\\
0.72	0.000489278\\
0.722	0.000464627\\
0.724	0.000439984\\
0.726	0.000415394\\
0.728	0.000390903\\
0.73	0.000366554\\
0.732	0.000342387\\
0.734	0.000318444\\
0.736	0.000294762\\
0.738	0.00027138\\
0.74	0.000248332\\
0.742	0.000225651\\
0.744	0.000203371\\
0.746	0.000181521\\
0.748	0.000162268\\
0.75	0.000153586\\
0.752	0.000144796\\
0.754	0.00013592\\
0.756	0.000126977\\
0.758	0.000117987\\
0.76	0.000108969\\
0.762	9.99429e-05\\
0.764	9.20734e-05\\
0.766	8.96797e-05\\
0.768	8.72933e-05\\
0.77	8.49138e-05\\
0.772	8.25411e-05\\
0.774	8.1117e-05\\
0.776	7.96406e-05\\
0.778	9.15072e-05\\
0.78	0.00010307\\
0.782	0.000113948\\
0.784	0.000124143\\
0.786	0.000133655\\
0.788	0.00014249\\
0.79	0.000150651\\
0.792	0.000158145\\
0.794	0.00016498\\
0.796	0.000171163\\
0.798	0.000176705\\
0.8	0.000181616\\
0.802	0.000185907\\
0.804	0.00018959\\
0.806	0.00019268\\
0.808	0.00019519\\
0.81	0.000197135\\
0.812	0.000198531\\
0.814	0.000199394\\
0.816	0.000199741\\
0.818	0.000199589\\
0.82	0.000198956\\
0.822	0.00019786\\
0.824	0.000196321\\
0.826	0.000194356\\
0.828	0.000191986\\
0.83	0.00018923\\
0.832	0.000186108\\
0.834	0.000182639\\
0.836	0.000178844\\
0.838	0.000174741\\
0.84	0.000170353\\
0.842	0.000165697\\
0.844	0.000160795\\
0.846	0.000155665\\
0.848	0.000150328\\
0.85	0.000144803\\
0.852	0.000140549\\
0.854	0.000140758\\
0.856	0.000140765\\
0.858	0.000140575\\
0.86	0.000140194\\
0.862	0.000139627\\
0.864	0.00013888\\
0.866	0.000137959\\
0.868	0.00013687\\
0.87	0.000135618\\
0.872	0.00013421\\
0.874	0.000132653\\
0.876	0.000130952\\
0.878	0.000129113\\
0.88	0.000127145\\
0.882	0.000125052\\
0.884	0.000122841\\
0.886	0.000120519\\
0.888	0.000118093\\
0.89	0.000115569\\
0.892	0.000112953\\
0.894	0.000110252\\
0.896	0.000107473\\
0.898	0.000104621\\
0.9	0.000101704\\
0.902	9.8728e-05\\
0.904	9.56985e-05\\
0.906	9.26221e-05\\
0.908	8.9505e-05\\
0.91	8.6353e-05\\
0.912	8.31721e-05\\
0.914	7.99681e-05\\
0.916	7.67467e-05\\
0.918	7.35135e-05\\
0.92	7.02738e-05\\
0.922	6.7033e-05\\
0.924	6.37963e-05\\
0.926	6.05687e-05\\
0.928	5.85647e-05\\
0.93	6.09116e-05\\
0.932	6.3066e-05\\
0.934	6.50293e-05\\
0.936	6.68036e-05\\
0.938	6.83911e-05\\
0.94	6.97945e-05\\
0.942	7.10165e-05\\
0.944	7.20606e-05\\
0.946	7.29302e-05\\
0.948	7.36292e-05\\
0.95	7.41616e-05\\
0.952	7.45316e-05\\
0.954	7.47438e-05\\
0.956	7.48028e-05\\
0.958	7.47136e-05\\
0.96	7.44812e-05\\
0.962	7.41107e-05\\
0.964	7.36074e-05\\
0.966	7.29769e-05\\
0.968	7.22246e-05\\
0.97	7.13561e-05\\
0.972	7.0377e-05\\
0.974	6.92931e-05\\
0.976	6.81102e-05\\
0.978	6.68339e-05\\
0.98	6.54701e-05\\
0.982	6.40246e-05\\
0.984	6.2503e-05\\
0.986	6.09111e-05\\
0.988	5.92547e-05\\
0.99	5.75393e-05\\
0.992	5.57704e-05\\
0.994	5.39536e-05\\
0.996	5.20942e-05\\
0.998	5.01977e-05\\
1	4.82691e-05\\
1.002	4.63136e-05\\
1.004	4.43361e-05\\
1.006	4.23416e-05\\
1.008	4.03347e-05\\
1.01	3.832e-05\\
1.012	3.6302e-05\\
1.014	3.4285e-05\\
1.016	3.22732e-05\\
1.018	3.02704e-05\\
1.02	2.82807e-05\\
1.022	2.68639e-05\\
1.024	2.69367e-05\\
1.026	2.69524e-05\\
1.028	2.69132e-05\\
1.03	2.68209e-05\\
1.032	2.66777e-05\\
1.034	2.64857e-05\\
1.036	2.62469e-05\\
1.038	2.59635e-05\\
1.04	2.56376e-05\\
1.042	2.52713e-05\\
1.044	2.48669e-05\\
1.046	2.44264e-05\\
1.048	2.3952e-05\\
1.05	2.34459e-05\\
1.052	2.291e-05\\
1.054	2.23466e-05\\
1.056	2.17578e-05\\
1.058	2.11455e-05\\
1.06	2.05118e-05\\
1.062	1.98587e-05\\
1.064	1.91881e-05\\
1.066	1.85021e-05\\
1.068	1.78025e-05\\
1.07	1.7091e-05\\
1.072	1.63697e-05\\
1.074	1.56401e-05\\
1.076	1.4904e-05\\
1.078	1.41632e-05\\
1.08	1.34191e-05\\
1.082	1.36233e-05\\
1.084	1.41387e-05\\
1.086	1.4602e-05\\
1.088	1.50143e-05\\
1.09	1.53765e-05\\
1.092	1.56897e-05\\
1.094	1.59551e-05\\
1.096	1.6174e-05\\
1.098	1.63475e-05\\
1.1	1.64769e-05\\
1.102	1.65638e-05\\
1.104	1.66094e-05\\
1.106	1.66152e-05\\
1.108	1.65827e-05\\
1.11	1.65134e-05\\
1.112	1.64089e-05\\
1.114	1.62706e-05\\
1.116	1.61002e-05\\
1.118	1.58993e-05\\
1.12	1.56694e-05\\
1.122	1.54121e-05\\
1.124	1.51291e-05\\
1.126	1.48219e-05\\
1.128	1.44921e-05\\
1.13	1.41413e-05\\
1.132	1.37711e-05\\
1.134	1.3383e-05\\
1.136	1.29786e-05\\
1.138	1.25594e-05\\
1.14	1.21268e-05\\
1.142	1.16823e-05\\
1.144	1.12275e-05\\
1.146	1.07636e-05\\
1.148	1.0292e-05\\
1.15	9.8142e-06\\
1.152	9.33138e-06\\
1.154	8.84485e-06\\
1.156	8.35583e-06\\
1.158	7.86551e-06\\
1.16	7.37505e-06\\
1.162	7.34628e-06\\
1.164	7.48131e-06\\
1.166	7.59906e-06\\
1.168	7.69993e-06\\
1.17	7.78429e-06\\
1.172	7.85255e-06\\
1.174	7.90515e-06\\
1.176	7.94251e-06\\
1.178	7.9651e-06\\
1.18	7.97336e-06\\
1.182	7.96777e-06\\
1.184	7.94882e-06\\
1.186	7.917e-06\\
1.188	7.87279e-06\\
1.19	7.8167e-06\\
1.192	7.74924e-06\\
1.194	7.6709e-06\\
1.196	7.58221e-06\\
1.198	7.48368e-06\\
1.2	7.37581e-06\\
1.202	7.25911e-06\\
1.204	7.1341e-06\\
1.206	7.00128e-06\\
1.208	6.86116e-06\\
1.21	6.71422e-06\\
1.212	6.56097e-06\\
1.214	6.4019e-06\\
1.216	6.23749e-06\\
1.218	6.0682e-06\\
1.22	5.89451e-06\\
1.222	5.71688e-06\\
1.224	5.53576e-06\\
1.226	5.35158e-06\\
1.228	5.16478e-06\\
1.23	4.97579e-06\\
1.232	4.785e-06\\
1.234	4.59283e-06\\
1.236	4.4452e-06\\
1.238	4.50489e-06\\
1.24	4.55271e-06\\
1.242	4.589e-06\\
1.244	4.6141e-06\\
1.246	4.62837e-06\\
1.248	4.63218e-06\\
1.25	4.62591e-06\\
1.252	4.60995e-06\\
1.254	4.58469e-06\\
1.256	4.55054e-06\\
1.258	4.50791e-06\\
1.26	4.45721e-06\\
1.262	4.39886e-06\\
1.264	4.33327e-06\\
1.266	4.26086e-06\\
1.268	4.18205e-06\\
1.27	4.09727e-06\\
1.272	4.00691e-06\\
1.274	3.91141e-06\\
1.276	3.81116e-06\\
1.278	3.70658e-06\\
1.28	3.59806e-06\\
1.282	3.48599e-06\\
1.284	3.37076e-06\\
1.286	3.25276e-06\\
1.288	3.13235e-06\\
1.29	3.0099e-06\\
1.292	2.88576e-06\\
1.294	2.76027e-06\\
1.296	2.63378e-06\\
1.298	2.5066e-06\\
1.3	2.37904e-06\\
1.302	2.25142e-06\\
1.304	2.12403e-06\\
1.306	1.99713e-06\\
1.308	1.87101e-06\\
1.31	1.74591e-06\\
1.312	1.62208e-06\\
1.314	1.49975e-06\\
1.316	1.37915e-06\\
1.318	1.26945e-06\\
1.32	1.3167e-06\\
1.322	1.35953e-06\\
1.324	1.39804e-06\\
1.326	1.43232e-06\\
1.328	1.46246e-06\\
1.33	1.48855e-06\\
1.332	1.5107e-06\\
1.334	1.52901e-06\\
1.336	1.54359e-06\\
1.338	1.55456e-06\\
1.34	1.56203e-06\\
1.342	1.56612e-06\\
1.344	1.56694e-06\\
1.346	1.56462e-06\\
1.348	1.55928e-06\\
1.35	1.55105e-06\\
1.352	1.54005e-06\\
1.354	1.52641e-06\\
1.356	1.51025e-06\\
1.358	1.49171e-06\\
1.36	1.4709e-06\\
1.362	1.44795e-06\\
1.364	1.423e-06\\
1.366	1.39616e-06\\
1.368	1.36756e-06\\
1.37	1.33733e-06\\
1.372	1.30558e-06\\
1.374	1.27243e-06\\
1.376	1.238e-06\\
1.378	1.20242e-06\\
1.38	1.16578e-06\\
1.382	1.12821e-06\\
1.384	1.08982e-06\\
1.386	1.0507e-06\\
1.388	1.01097e-06\\
1.39	9.70731e-07\\
1.392	9.30076e-07\\
1.394	9.21144e-07\\
1.396	9.17996e-07\\
1.398	9.12659e-07\\
1.4	9.05236e-07\\
1.402	8.95835e-07\\
1.404	8.84562e-07\\
1.406	8.71521e-07\\
1.408	8.56821e-07\\
1.41	8.40566e-07\\
1.412	8.22862e-07\\
1.414	8.03814e-07\\
1.416	7.83525e-07\\
1.418	7.62097e-07\\
1.42	7.39632e-07\\
1.422	7.16229e-07\\
1.424	6.91986e-07\\
1.426	6.66998e-07\\
1.428	6.41358e-07\\
1.43	6.15159e-07\\
1.432	5.88489e-07\\
1.434	5.61435e-07\\
1.436	5.34081e-07\\
1.438	5.06508e-07\\
1.44	4.78794e-07\\
1.442	4.51015e-07\\
1.444	4.23244e-07\\
1.446	3.9555e-07\\
1.448	3.68e-07\\
1.45	3.40659e-07\\
1.452	3.19882e-07\\
1.454	3.19109e-07\\
1.456	3.17764e-07\\
1.458	3.15866e-07\\
1.46	3.13437e-07\\
1.462	3.10498e-07\\
1.464	3.0707e-07\\
1.466	3.03176e-07\\
1.468	2.98836e-07\\
1.47	2.94072e-07\\
1.472	2.9431e-07\\
1.474	3.05806e-07\\
1.476	3.16322e-07\\
1.478	3.25874e-07\\
1.48	3.34481e-07\\
1.482	3.42164e-07\\
1.484	3.48943e-07\\
1.486	3.54839e-07\\
1.488	3.59877e-07\\
1.49	3.64077e-07\\
1.492	3.67466e-07\\
1.494	3.70068e-07\\
1.496	3.71908e-07\\
1.498	3.73011e-07\\
1.5	3.73405e-07\\
1.502	3.73114e-07\\
1.504	3.72167e-07\\
1.506	3.7059e-07\\
1.508	3.68411e-07\\
1.51	3.65656e-07\\
1.512	3.62352e-07\\
1.514	3.58527e-07\\
1.516	3.54208e-07\\
1.518	3.49422e-07\\
1.52	3.44196e-07\\
1.522	3.38557e-07\\
1.524	3.32529e-07\\
1.526	3.26141e-07\\
1.528	3.19417e-07\\
1.53	3.12382e-07\\
1.532	3.05062e-07\\
1.534	2.97481e-07\\
1.536	2.89664e-07\\
1.538	2.81632e-07\\
1.54	2.7341e-07\\
1.542	2.6502e-07\\
1.544	2.56483e-07\\
1.546	2.47821e-07\\
1.548	2.39055e-07\\
1.55	2.30205e-07\\
1.552	2.2129e-07\\
1.554	2.12328e-07\\
1.556	2.03338e-07\\
1.558	1.94338e-07\\
1.56	1.85344e-07\\
1.562	1.76373e-07\\
1.564	1.67439e-07\\
1.566	1.58558e-07\\
1.568	1.49744e-07\\
1.57	1.4101e-07\\
1.572	1.32369e-07\\
1.574	1.23833e-07\\
1.576	1.15414e-07\\
1.578	1.07121e-07\\
1.58	9.89664e-08\\
1.582	9.09584e-08\\
1.584	8.3106e-08\\
1.586	8.37148e-08\\
1.588	8.53682e-08\\
1.59	8.67985e-08\\
1.592	8.80102e-08\\
1.594	8.90083e-08\\
1.596	8.9798e-08\\
1.598	9.03845e-08\\
1.6	9.07737e-08\\
1.602	9.09714e-08\\
1.604	9.09837e-08\\
1.606	9.08171e-08\\
1.608	9.0478e-08\\
1.61	8.9973e-08\\
1.612	8.93089e-08\\
1.614	8.84927e-08\\
1.616	8.75313e-08\\
1.618	8.64318e-08\\
1.62	8.52013e-08\\
1.622	8.38469e-08\\
1.624	8.23759e-08\\
1.626	8.07954e-08\\
1.628	7.91127e-08\\
1.63	7.73347e-08\\
1.632	7.54687e-08\\
1.634	7.35217e-08\\
1.636	7.15006e-08\\
1.638	6.94122e-08\\
1.64	6.72635e-08\\
1.642	6.5061e-08\\
1.644	6.28112e-08\\
1.646	6.07431e-08\\
1.648	6.21326e-08\\
1.65	6.3349e-08\\
1.652	6.43969e-08\\
1.654	6.52813e-08\\
1.656	6.6007e-08\\
1.658	6.65792e-08\\
1.66	6.7003e-08\\
1.662	6.72837e-08\\
1.664	6.74265e-08\\
1.666	6.74368e-08\\
1.668	6.73199e-08\\
1.67	6.70813e-08\\
1.672	6.67264e-08\\
1.674	6.62606e-08\\
1.676	6.56893e-08\\
1.678	6.50179e-08\\
1.68	6.42517e-08\\
1.682	6.33961e-08\\
1.684	6.24565e-08\\
1.686	6.14379e-08\\
1.688	6.03455e-08\\
1.69	5.91846e-08\\
1.692	5.796e-08\\
1.694	5.66768e-08\\
1.696	5.53397e-08\\
1.698	5.39536e-08\\
1.7	5.2523e-08\\
1.702	5.10526e-08\\
1.704	4.95467e-08\\
1.706	4.80097e-08\\
1.708	4.64457e-08\\
1.71	4.48589e-08\\
1.712	4.32532e-08\\
1.714	4.16325e-08\\
1.716	4.00003e-08\\
1.718	3.83603e-08\\
1.72	3.67159e-08\\
1.722	3.50704e-08\\
1.724	3.34269e-08\\
1.726	3.17885e-08\\
1.728	3.0158e-08\\
1.73	2.85382e-08\\
1.732	2.69317e-08\\
1.734	2.53409e-08\\
1.736	2.37682e-08\\
1.738	2.22158e-08\\
1.74	2.06857e-08\\
1.742	1.91799e-08\\
1.744	1.87189e-08\\
1.746	1.89857e-08\\
1.748	1.9204e-08\\
1.75	1.93752e-08\\
1.752	1.95009e-08\\
1.754	1.95827e-08\\
1.756	1.96222e-08\\
1.758	1.96209e-08\\
1.76	1.95806e-08\\
1.762	1.95028e-08\\
1.764	1.93893e-08\\
1.766	1.92416e-08\\
1.768	1.90614e-08\\
1.77	1.88505e-08\\
1.772	1.86103e-08\\
1.774	1.83426e-08\\
1.776	1.83218e-08\\
1.778	1.82741e-08\\
1.78	1.81757e-08\\
1.782	1.80282e-08\\
1.784	1.78338e-08\\
1.786	1.75943e-08\\
1.788	1.73119e-08\\
1.79	1.69885e-08\\
1.792	1.66264e-08\\
1.794	1.62277e-08\\
1.796	1.57946e-08\\
1.798	1.53294e-08\\
1.8	1.48344e-08\\
1.802	1.43117e-08\\
1.804	1.37637e-08\\
1.806	1.31926e-08\\
1.808	1.26006e-08\\
1.81	1.19901e-08\\
1.812	1.15111e-08\\
1.814	1.17501e-08\\
1.816	1.19602e-08\\
1.818	1.21424e-08\\
1.82	1.22974e-08\\
1.822	1.24261e-08\\
1.824	1.25295e-08\\
1.826	1.26085e-08\\
1.828	1.26639e-08\\
1.83	1.26966e-08\\
1.832	1.27076e-08\\
1.834	1.26977e-08\\
1.836	1.26679e-08\\
1.838	1.2619e-08\\
1.84	1.25519e-08\\
1.842	1.24676e-08\\
1.844	1.23667e-08\\
1.846	1.22503e-08\\
1.848	1.21192e-08\\
1.85	1.19741e-08\\
1.852	1.18159e-08\\
1.854	1.16454e-08\\
1.856	1.14634e-08\\
1.858	1.12707e-08\\
1.86	1.10679e-08\\
1.862	1.0856e-08\\
1.864	1.06354e-08\\
1.866	1.04071e-08\\
1.868	1.01716e-08\\
1.87	9.92959e-09\\
1.872	9.68175e-09\\
1.874	9.4287e-09\\
1.876	9.17102e-09\\
1.878	8.90932e-09\\
1.88	8.64417e-09\\
1.882	8.37609e-09\\
1.884	8.10564e-09\\
1.886	7.83331e-09\\
1.888	7.87513e-09\\
1.89	8.00169e-09\\
1.892	8.10515e-09\\
1.894	8.1859e-09\\
1.896	8.24434e-09\\
1.898	8.28093e-09\\
1.9	8.2962e-09\\
1.902	8.29066e-09\\
1.904	8.26491e-09\\
1.906	8.21956e-09\\
1.908	8.15524e-09\\
1.91	8.07264e-09\\
1.912	7.97246e-09\\
1.914	7.85541e-09\\
1.916	7.72223e-09\\
1.918	7.5737e-09\\
1.92	7.41059e-09\\
1.922	7.2337e-09\\
1.924	7.04382e-09\\
1.926	6.84176e-09\\
1.928	6.62836e-09\\
1.93	6.40442e-09\\
1.932	6.17078e-09\\
1.934	5.92827e-09\\
1.936	5.6777e-09\\
1.938	5.41991e-09\\
1.94	5.15569e-09\\
1.942	4.88587e-09\\
1.944	4.61123e-09\\
1.946	4.33257e-09\\
1.948	4.05066e-09\\
1.95	3.76626e-09\\
1.952	3.48011e-09\\
1.954	3.19294e-09\\
1.956	2.90547e-09\\
1.958	2.68174e-09\\
1.96	2.5729e-09\\
1.962	2.46336e-09\\
1.964	2.35337e-09\\
1.966	2.24321e-09\\
1.968	2.13311e-09\\
1.97	2.0233e-09\\
1.972	1.91402e-09\\
1.974	1.80548e-09\\
1.976	1.69788e-09\\
1.978	1.59141e-09\\
1.98	1.57765e-09\\
1.982	1.65486e-09\\
1.984	1.72687e-09\\
1.986	1.79377e-09\\
1.988	1.85564e-09\\
1.99	1.91256e-09\\
1.992	1.96464e-09\\
1.994	2.01196e-09\\
1.996	2.05463e-09\\
1.998	2.18648e-09\\
2	2.34848e-09\\
};
\addplot [color=mycolor7, line width=1.0pt, forget plot]
  table[row sep=crcr]{%
-0.024	0\\
-0.022	0\\
-0.02	0\\
-0.018	0\\
-0.016	0\\
-0.014	0\\
-0.012	0\\
-0.01	0\\
-0.008	0\\
-0.006	0\\
-0.004	0\\
-0.002	0\\
0	0\\
0.002	0.0391009\\
0.004	0.074092\\
0.006	0.104794\\
0.008	0.131117\\
0.01	0.153053\\
0.012	0.170674\\
0.014	0.18412\\
0.016	0.193597\\
0.018	0.19936\\
0.02	0.201711\\
0.022	0.200986\\
0.024	0.197545\\
0.026	0.191764\\
0.028	0.184025\\
0.03	0.174712\\
0.032	0.164197\\
0.034	0.152837\\
0.036	0.140969\\
0.038	0.128904\\
0.04	0.11871\\
0.042	0.124423\\
0.044	0.1294\\
0.046	0.133573\\
0.048	0.136893\\
0.05	0.139325\\
0.052	0.14085\\
0.054	0.141463\\
0.056	0.141176\\
0.058	0.140013\\
0.06	0.138011\\
0.062	0.13522\\
0.064	0.131698\\
0.066	0.127515\\
0.068	0.122744\\
0.07	0.117467\\
0.072	0.111766\\
0.074	0.10573\\
0.076	0.0994448\\
0.078	0.0929965\\
0.08	0.0864692\\
0.082	0.0837476\\
0.084	0.0853397\\
0.086	0.0868178\\
0.088	0.0881831\\
0.09	0.0894367\\
0.092	0.0905797\\
0.094	0.0916128\\
0.096	0.0925368\\
0.098	0.093352\\
0.1	0.0940588\\
0.102	0.0946573\\
0.104	0.0951473\\
0.106	0.0955287\\
0.108	0.0958011\\
0.11	0.0959641\\
0.112	0.0969957\\
0.114	0.0984574\\
0.116	0.0997429\\
0.118	0.100854\\
0.12	0.101792\\
0.122	0.102561\\
0.124	0.103165\\
0.126	0.103608\\
0.128	0.103896\\
0.13	0.104034\\
0.132	0.104028\\
0.134	0.103884\\
0.136	0.10361\\
0.138	0.103211\\
0.14	0.102696\\
0.142	0.102071\\
0.144	0.101342\\
0.146	0.100518\\
0.148	0.0996038\\
0.15	0.0986075\\
0.152	0.0975351\\
0.154	0.0963931\\
0.156	0.0951874\\
0.158	0.0939238\\
0.16	0.092608\\
0.162	0.0912452\\
0.164	0.0898404\\
0.166	0.0883982\\
0.168	0.0869233\\
0.17	0.0854196\\
0.172	0.083891\\
0.174	0.0823413\\
0.176	0.0807738\\
0.178	0.0791916\\
0.18	0.0775976\\
0.182	0.0759947\\
0.184	0.0743852\\
0.186	0.0727716\\
0.188	0.0711561\\
0.19	0.0695407\\
0.192	0.0679272\\
0.194	0.0663175\\
0.196	0.0647133\\
0.198	0.063116\\
0.2	0.0615272\\
0.202	0.0599482\\
0.204	0.0583804\\
0.206	0.0568248\\
0.208	0.0552828\\
0.21	0.0537553\\
0.212	0.0522434\\
0.214	0.0507481\\
0.216	0.0492702\\
0.218	0.0478106\\
0.22	0.0463701\\
0.222	0.0449494\\
0.224	0.0435491\\
0.226	0.04217\\
0.228	0.0408125\\
0.23	0.0394772\\
0.232	0.0381645\\
0.234	0.0368748\\
0.236	0.0356084\\
0.238	0.0343657\\
0.24	0.0331469\\
0.242	0.0319521\\
0.244	0.0307815\\
0.246	0.0296353\\
0.248	0.0285134\\
0.25	0.0274158\\
0.252	0.0263426\\
0.254	0.0252936\\
0.256	0.0242688\\
0.258	0.0232681\\
0.26	0.0222912\\
0.262	0.0213379\\
0.264	0.0204082\\
0.266	0.0195018\\
0.268	0.0186184\\
0.27	0.0177578\\
0.272	0.0170803\\
0.274	0.0164288\\
0.276	0.0157953\\
0.278	0.0151807\\
0.28	0.0146926\\
0.282	0.0144376\\
0.284	0.0141875\\
0.286	0.0139425\\
0.288	0.0137024\\
0.29	0.0134672\\
0.292	0.0132366\\
0.294	0.0130107\\
0.296	0.0127894\\
0.298	0.0125725\\
0.3	0.0123599\\
0.302	0.0121515\\
0.304	0.0119472\\
0.306	0.011747\\
0.308	0.0115506\\
0.31	0.011358\\
0.312	0.011169\\
0.314	0.0109836\\
0.316	0.0108016\\
0.318	0.0106229\\
0.32	0.0104474\\
0.322	0.010275\\
0.324	0.0101057\\
0.326	0.00993918\\
0.328	0.0097755\\
0.33	0.00961452\\
0.332	0.00945615\\
0.334	0.00930031\\
0.336	0.00914689\\
0.338	0.00899584\\
0.34	0.00913297\\
0.342	0.00928747\\
0.344	0.00942903\\
0.346	0.00955744\\
0.348	0.00967254\\
0.35	0.0097742\\
0.352	0.00986235\\
0.354	0.00993695\\
0.356	0.00999801\\
0.358	0.0100456\\
0.36	0.0100797\\
0.362	0.0101005\\
0.364	0.0101082\\
0.366	0.0101029\\
0.368	0.0100848\\
0.37	0.0100542\\
0.372	0.0100112\\
0.374	0.00995626\\
0.376	0.00988959\\
0.378	0.00981152\\
0.38	0.0097224\\
0.382	0.00962258\\
0.384	0.00951243\\
0.386	0.00939234\\
0.388	0.00926271\\
0.39	0.00912395\\
0.392	0.00897647\\
0.394	0.00882071\\
0.396	0.00865709\\
0.398	0.00848608\\
0.4	0.0083081\\
0.402	0.00812362\\
0.404	0.00793309\\
0.406	0.00773696\\
0.408	0.0075357\\
0.41	0.00732976\\
0.412	0.0071196\\
0.414	0.00690567\\
0.416	0.00668843\\
0.418	0.00646832\\
0.42	0.0062458\\
0.422	0.0060213\\
0.424	0.00579525\\
0.426	0.00556807\\
0.428	0.0053402\\
0.43	0.00511203\\
0.432	0.00488397\\
0.434	0.00465641\\
0.436	0.00442973\\
0.438	0.0042043\\
0.44	0.00398048\\
0.442	0.00375862\\
0.444	0.00353904\\
0.446	0.00332208\\
0.448	0.00322647\\
0.45	0.00315937\\
0.452	0.0030934\\
0.454	0.00302857\\
0.456	0.00296485\\
0.458	0.00290224\\
0.46	0.00284071\\
0.462	0.00278026\\
0.464	0.00272087\\
0.466	0.00266252\\
0.468	0.00260521\\
0.47	0.00254893\\
0.472	0.00249365\\
0.474	0.00243936\\
0.476	0.00238605\\
0.478	0.00238135\\
0.48	0.00243482\\
0.482	0.00248511\\
0.484	0.00253218\\
0.486	0.00257602\\
0.488	0.00261659\\
0.49	0.00265389\\
0.492	0.00268792\\
0.494	0.00271866\\
0.496	0.00274613\\
0.498	0.00277033\\
0.5	0.00279129\\
0.502	0.00280902\\
0.504	0.00282356\\
0.506	0.00283494\\
0.508	0.00284319\\
0.51	0.00284836\\
0.512	0.0028505\\
0.514	0.00284965\\
0.516	0.00284588\\
0.518	0.00283925\\
0.52	0.00282981\\
0.522	0.00281763\\
0.524	0.00280279\\
0.526	0.00278536\\
0.528	0.00276541\\
0.53	0.00274302\\
0.532	0.00271827\\
0.534	0.00269125\\
0.536	0.00266204\\
0.538	0.00263072\\
0.54	0.00259739\\
0.542	0.00256213\\
0.544	0.00252503\\
0.546	0.00248618\\
0.548	0.00244568\\
0.55	0.00240362\\
0.552	0.00236008\\
0.554	0.00231517\\
0.556	0.00226898\\
0.558	0.00222159\\
0.56	0.00217311\\
0.562	0.00212361\\
0.564	0.0020732\\
0.566	0.00202197\\
0.568	0.00196999\\
0.57	0.00191737\\
0.572	0.00186419\\
0.574	0.00181053\\
0.576	0.00175648\\
0.578	0.00170212\\
0.58	0.00164753\\
0.582	0.0015928\\
0.584	0.001538\\
0.586	0.00148321\\
0.588	0.0014285\\
0.59	0.00137394\\
0.592	0.00131961\\
0.594	0.00126556\\
0.596	0.00121187\\
0.598	0.0011586\\
0.6	0.0011058\\
0.602	0.00105355\\
0.604	0.00100189\\
0.606	0.000950872\\
0.608	0.000900554\\
0.61	0.000850983\\
0.612	0.000802204\\
0.614	0.00075426\\
0.616	0.000707192\\
0.618	0.00066104\\
0.62	0.000615838\\
0.622	0.000629358\\
0.624	0.000661999\\
0.626	0.000692952\\
0.628	0.00072219\\
0.63	0.000749693\\
0.632	0.000775444\\
0.634	0.000799429\\
0.636	0.00082164\\
0.638	0.000842074\\
0.64	0.000860728\\
0.642	0.000877607\\
0.644	0.000892718\\
0.646	0.00090607\\
0.648	0.000917679\\
0.65	0.000927562\\
0.652	0.000935739\\
0.654	0.000942234\\
0.656	0.000947073\\
0.658	0.000950287\\
0.66	0.000951906\\
0.662	0.000951965\\
0.664	0.000950502\\
0.666	0.000947555\\
0.668	0.000943165\\
0.67	0.000937376\\
0.672	0.000930231\\
0.674	0.000921778\\
0.676	0.000912064\\
0.678	0.000901139\\
0.68	0.000889053\\
0.682	0.000875857\\
0.684	0.000861604\\
0.686	0.000846347\\
0.688	0.00083014\\
0.69	0.000813037\\
0.692	0.000795092\\
0.694	0.000776362\\
0.696	0.000756901\\
0.698	0.000736765\\
0.7	0.000716007\\
0.702	0.000694685\\
0.704	0.000672851\\
0.706	0.00065056\\
0.708	0.000627866\\
0.71	0.000604822\\
0.712	0.00058148\\
0.714	0.000557892\\
0.716	0.000534108\\
0.718	0.000510178\\
0.72	0.000486151\\
0.722	0.000462074\\
0.724	0.000437994\\
0.726	0.000413955\\
0.728	0.000390002\\
0.73	0.000366177\\
0.732	0.00034252\\
0.734	0.000319072\\
0.736	0.00029587\\
0.738	0.00027295\\
0.74	0.000250349\\
0.742	0.000228098\\
0.744	0.00020623\\
0.746	0.000184774\\
0.748	0.00016376\\
0.75	0.00015193\\
0.752	0.00014312\\
0.754	0.000134227\\
0.756	0.000125269\\
0.758	0.000116266\\
0.76	0.000107239\\
0.762	9.8204e-05\\
0.764	9.15944e-05\\
0.766	8.92156e-05\\
0.768	8.68447e-05\\
0.77	8.44815e-05\\
0.772	8.2126e-05\\
0.774	8.00632e-05\\
0.776	7.86359e-05\\
0.778	8.4474e-05\\
0.78	9.59564e-05\\
0.782	0.000106771\\
0.784	0.000116918\\
0.786	0.0001264\\
0.788	0.00013522\\
0.79	0.000143382\\
0.792	0.000150891\\
0.794	0.000157756\\
0.796	0.000163983\\
0.798	0.000169582\\
0.8	0.000174563\\
0.802	0.000178936\\
0.804	0.000182714\\
0.806	0.000185909\\
0.808	0.000188536\\
0.81	0.000190607\\
0.812	0.000192139\\
0.814	0.000193147\\
0.816	0.000193647\\
0.818	0.000193656\\
0.82	0.000193192\\
0.822	0.000192271\\
0.824	0.000190912\\
0.826	0.000189134\\
0.828	0.000186956\\
0.83	0.000184395\\
0.832	0.000181472\\
0.834	0.000178205\\
0.836	0.000174615\\
0.838	0.00017072\\
0.84	0.00016654\\
0.842	0.000162094\\
0.844	0.000157402\\
0.846	0.000152483\\
0.848	0.000147355\\
0.85	0.000142038\\
0.852	0.000141397\\
0.854	0.000141581\\
0.856	0.000141563\\
0.858	0.000141349\\
0.86	0.000140945\\
0.862	0.000140356\\
0.864	0.000139587\\
0.866	0.000138645\\
0.868	0.000137535\\
0.87	0.000136264\\
0.872	0.000134838\\
0.874	0.000133262\\
0.876	0.000131544\\
0.878	0.00012969\\
0.88	0.000127706\\
0.882	0.000125599\\
0.884	0.000123374\\
0.886	0.00012104\\
0.888	0.000118602\\
0.89	0.000116066\\
0.892	0.00011344\\
0.894	0.00011073\\
0.896	0.000107943\\
0.898	0.000105083\\
0.9	0.00010216\\
0.902	9.9177e-05\\
0.904	9.61422e-05\\
0.906	9.30612e-05\\
0.908	8.99401e-05\\
0.91	8.67849e-05\\
0.912	8.36015e-05\\
0.914	8.03956e-05\\
0.916	7.71729e-05\\
0.918	7.39389e-05\\
0.92	7.06991e-05\\
0.922	6.74587e-05\\
0.924	6.42228e-05\\
0.926	6.09965e-05\\
0.928	5.77845e-05\\
0.93	5.91853e-05\\
0.932	6.13485e-05\\
0.934	6.33256e-05\\
0.936	6.51182e-05\\
0.938	6.67285e-05\\
0.94	6.81589e-05\\
0.942	6.94122e-05\\
0.944	7.04913e-05\\
0.946	7.13996e-05\\
0.948	7.21406e-05\\
0.95	7.27183e-05\\
0.952	7.31366e-05\\
0.954	7.33999e-05\\
0.956	7.35126e-05\\
0.958	7.34794e-05\\
0.96	7.3305e-05\\
0.962	7.29946e-05\\
0.964	7.25531e-05\\
0.966	7.19858e-05\\
0.968	7.12979e-05\\
0.97	7.0495e-05\\
0.972	6.95823e-05\\
0.974	6.85656e-05\\
0.976	6.74502e-05\\
0.978	6.62418e-05\\
0.98	6.4946e-05\\
0.982	6.35683e-05\\
0.984	6.21144e-05\\
0.986	6.05898e-05\\
0.988	5.89999e-05\\
0.99	5.73504e-05\\
0.992	5.56466e-05\\
0.994	5.38937e-05\\
0.996	5.20972e-05\\
0.998	5.02622e-05\\
1	4.83937e-05\\
1.002	4.64968e-05\\
1.004	4.45763e-05\\
1.006	4.2637e-05\\
1.008	4.06835e-05\\
1.01	3.87203e-05\\
1.012	3.67518e-05\\
1.014	3.47822e-05\\
1.016	3.28156e-05\\
1.018	3.08559e-05\\
1.02	2.8907e-05\\
1.022	2.69724e-05\\
1.024	2.64688e-05\\
1.026	2.64914e-05\\
1.028	2.64592e-05\\
1.03	2.63744e-05\\
1.032	2.62389e-05\\
1.034	2.60549e-05\\
1.036	2.58244e-05\\
1.038	2.55496e-05\\
1.04	2.52325e-05\\
1.042	2.48753e-05\\
1.044	2.44802e-05\\
1.046	2.40492e-05\\
1.048	2.35844e-05\\
1.05	2.30881e-05\\
1.052	2.25623e-05\\
1.054	2.2009e-05\\
1.056	2.14304e-05\\
1.058	2.08285e-05\\
1.06	2.02053e-05\\
1.062	1.95627e-05\\
1.064	1.89027e-05\\
1.066	1.82273e-05\\
1.068	1.75383e-05\\
1.07	1.68375e-05\\
1.072	1.61268e-05\\
1.074	1.54078e-05\\
1.076	1.46823e-05\\
1.078	1.39519e-05\\
1.08	1.32183e-05\\
1.082	1.28286e-05\\
1.084	1.33665e-05\\
1.086	1.38535e-05\\
1.088	1.42903e-05\\
1.09	1.46778e-05\\
1.092	1.50172e-05\\
1.094	1.53094e-05\\
1.096	1.55556e-05\\
1.098	1.5757e-05\\
1.1	1.59148e-05\\
1.102	1.60303e-05\\
1.104	1.61049e-05\\
1.106	1.61399e-05\\
1.108	1.61368e-05\\
1.11	1.60969e-05\\
1.112	1.60217e-05\\
1.114	1.59128e-05\\
1.116	1.57716e-05\\
1.118	1.55997e-05\\
1.12	1.53986e-05\\
1.122	1.51698e-05\\
1.124	1.49148e-05\\
1.126	1.46353e-05\\
1.128	1.43327e-05\\
1.13	1.40086e-05\\
1.132	1.36645e-05\\
1.134	1.33019e-05\\
1.136	1.29223e-05\\
1.138	1.25273e-05\\
1.14	1.21181e-05\\
1.142	1.16964e-05\\
1.144	1.12634e-05\\
1.146	1.08206e-05\\
1.148	1.03694e-05\\
1.15	9.91096e-06\\
1.152	9.4467e-06\\
1.154	8.97782e-06\\
1.156	8.50556e-06\\
1.158	8.0311e-06\\
1.16	7.55557e-06\\
1.162	7.28517e-06\\
1.164	7.41873e-06\\
1.166	7.53522e-06\\
1.168	7.63501e-06\\
1.17	7.71849e-06\\
1.172	7.78606e-06\\
1.174	7.83815e-06\\
1.176	7.87518e-06\\
1.178	7.89759e-06\\
1.18	7.90584e-06\\
1.182	7.90041e-06\\
1.184	7.88175e-06\\
1.186	7.85037e-06\\
1.188	7.80673e-06\\
1.19	7.75135e-06\\
1.192	7.68471e-06\\
1.194	7.60732e-06\\
1.196	7.51968e-06\\
1.198	7.42229e-06\\
1.2	7.31567e-06\\
1.202	7.20032e-06\\
1.204	7.07673e-06\\
1.206	6.94541e-06\\
1.208	6.80686e-06\\
1.21	6.66156e-06\\
1.212	6.51001e-06\\
1.214	6.35269e-06\\
1.216	6.19007e-06\\
1.218	6.02262e-06\\
1.22	5.85081e-06\\
1.222	5.67509e-06\\
1.224	5.4959e-06\\
1.226	5.31368e-06\\
1.228	5.12885e-06\\
1.23	4.94184e-06\\
1.232	4.75304e-06\\
1.234	4.56286e-06\\
1.236	4.37167e-06\\
1.238	4.34414e-06\\
1.24	4.40162e-06\\
1.242	4.44764e-06\\
1.244	4.48253e-06\\
1.246	4.50661e-06\\
1.248	4.52024e-06\\
1.25	4.52376e-06\\
1.252	4.51755e-06\\
1.254	4.50198e-06\\
1.256	4.47743e-06\\
1.258	4.4443e-06\\
1.26	4.40296e-06\\
1.262	4.35383e-06\\
1.264	4.29729e-06\\
1.266	4.23376e-06\\
1.268	4.16363e-06\\
1.27	4.08732e-06\\
1.272	4.00521e-06\\
1.274	3.91771e-06\\
1.276	3.82522e-06\\
1.278	3.72814e-06\\
1.28	3.62684e-06\\
1.282	3.52172e-06\\
1.284	3.41315e-06\\
1.286	3.30151e-06\\
1.288	3.18717e-06\\
1.29	3.07047e-06\\
1.292	2.95177e-06\\
1.294	2.83142e-06\\
1.296	2.70974e-06\\
1.298	2.58705e-06\\
1.3	2.46367e-06\\
1.302	2.3399e-06\\
1.304	2.21604e-06\\
1.306	2.09236e-06\\
1.308	1.96913e-06\\
1.31	1.84661e-06\\
1.312	1.72505e-06\\
1.314	1.60469e-06\\
1.316	1.48574e-06\\
1.318	1.36842e-06\\
1.32	1.29546e-06\\
1.322	1.3382e-06\\
1.324	1.37667e-06\\
1.326	1.41094e-06\\
1.328	1.4411e-06\\
1.33	1.46726e-06\\
1.332	1.48952e-06\\
1.334	1.50797e-06\\
1.336	1.52274e-06\\
1.338	1.53391e-06\\
1.34	1.54162e-06\\
1.342	1.54598e-06\\
1.344	1.5471e-06\\
1.346	1.54511e-06\\
1.348	1.54012e-06\\
1.35	1.53226e-06\\
1.352	1.52165e-06\\
1.354	1.50843e-06\\
1.356	1.4927e-06\\
1.358	1.4746e-06\\
1.36	1.45426e-06\\
1.362	1.43179e-06\\
1.364	1.40732e-06\\
1.366	1.38098e-06\\
1.368	1.35289e-06\\
1.37	1.32317e-06\\
1.372	1.29194e-06\\
1.374	1.25932e-06\\
1.376	1.22542e-06\\
1.378	1.19037e-06\\
1.38	1.15427e-06\\
1.382	1.11723e-06\\
1.384	1.07936e-06\\
1.386	1.04078e-06\\
1.388	1.00157e-06\\
1.39	9.61851e-07\\
1.392	9.21709e-07\\
1.394	8.97643e-07\\
1.396	8.9783e-07\\
1.398	8.95772e-07\\
1.4	8.91568e-07\\
1.402	8.85316e-07\\
1.404	8.77118e-07\\
1.406	8.67076e-07\\
1.408	8.5529e-07\\
1.41	8.41863e-07\\
1.412	8.26895e-07\\
1.414	8.10488e-07\\
1.416	7.92741e-07\\
1.418	7.73755e-07\\
1.42	7.53628e-07\\
1.422	7.32456e-07\\
1.424	7.10335e-07\\
1.426	6.8736e-07\\
1.428	6.63622e-07\\
1.43	6.39211e-07\\
1.432	6.14217e-07\\
1.434	5.88725e-07\\
1.436	5.62818e-07\\
1.438	5.36578e-07\\
1.44	5.10084e-07\\
1.442	4.83412e-07\\
1.444	4.56636e-07\\
1.446	4.29827e-07\\
1.448	4.03051e-07\\
1.45	3.76375e-07\\
1.452	3.49861e-07\\
1.454	3.23568e-07\\
1.456	3.03948e-07\\
1.458	3.02087e-07\\
1.46	2.99722e-07\\
1.462	2.96873e-07\\
1.464	2.93561e-07\\
1.466	2.89805e-07\\
1.468	2.85627e-07\\
1.47	2.81047e-07\\
1.472	2.92141e-07\\
1.474	3.03632e-07\\
1.476	3.14149e-07\\
1.478	3.23712e-07\\
1.48	3.32337e-07\\
1.482	3.40045e-07\\
1.484	3.46855e-07\\
1.486	3.52789e-07\\
1.488	3.5787e-07\\
1.49	3.62119e-07\\
1.492	3.65561e-07\\
1.494	3.68221e-07\\
1.496	3.70122e-07\\
1.498	3.71291e-07\\
1.5	3.71752e-07\\
1.502	3.71533e-07\\
1.504	3.7066e-07\\
1.506	3.69158e-07\\
1.508	3.67055e-07\\
1.51	3.64378e-07\\
1.512	3.61153e-07\\
1.514	3.57408e-07\\
1.516	3.53169e-07\\
1.518	3.48463e-07\\
1.52	3.43317e-07\\
1.522	3.37756e-07\\
1.524	3.31808e-07\\
1.526	3.25496e-07\\
1.528	3.18848e-07\\
1.53	3.11889e-07\\
1.532	3.04642e-07\\
1.534	2.97132e-07\\
1.536	2.89384e-07\\
1.538	2.8142e-07\\
1.54	2.73263e-07\\
1.542	2.64936e-07\\
1.544	2.5646e-07\\
1.546	2.47857e-07\\
1.548	2.39146e-07\\
1.55	2.30349e-07\\
1.552	2.21484e-07\\
1.554	2.12571e-07\\
1.556	2.03627e-07\\
1.558	1.94669e-07\\
1.56	1.85715e-07\\
1.562	1.76781e-07\\
1.564	1.67882e-07\\
1.566	1.59034e-07\\
1.568	1.50249e-07\\
1.57	1.41542e-07\\
1.572	1.32926e-07\\
1.574	1.26329e-07\\
1.576	1.20064e-07\\
1.578	1.13759e-07\\
1.58	1.07433e-07\\
1.582	1.01103e-07\\
1.584	9.47854e-08\\
1.586	8.84978e-08\\
1.588	8.22553e-08\\
1.59	8.33859e-08\\
1.592	8.4547e-08\\
1.594	8.55037e-08\\
1.596	8.62608e-08\\
1.598	8.68236e-08\\
1.6	8.71975e-08\\
1.602	8.73881e-08\\
1.604	8.74013e-08\\
1.606	8.72433e-08\\
1.608	8.69201e-08\\
1.61	8.64382e-08\\
1.612	8.58041e-08\\
1.614	8.50243e-08\\
1.616	8.41056e-08\\
1.618	8.30548e-08\\
1.62	8.18785e-08\\
1.622	8.05837e-08\\
1.624	7.91773e-08\\
1.626	7.76661e-08\\
1.628	7.6057e-08\\
1.63	7.43569e-08\\
1.632	7.25724e-08\\
1.634	7.07104e-08\\
1.636	6.87774e-08\\
1.638	6.67801e-08\\
1.64	6.4725e-08\\
1.642	6.26183e-08\\
1.644	6.04664e-08\\
1.646	6.02321e-08\\
1.648	6.16296e-08\\
1.65	6.28536e-08\\
1.652	6.39088e-08\\
1.654	6.48001e-08\\
1.656	6.55323e-08\\
1.658	6.61107e-08\\
1.66	6.65402e-08\\
1.662	6.68262e-08\\
1.664	6.69739e-08\\
1.666	6.69888e-08\\
1.668	6.68761e-08\\
1.67	6.66412e-08\\
1.672	6.62897e-08\\
1.674	6.58269e-08\\
1.676	6.52582e-08\\
1.678	6.45892e-08\\
1.68	6.3825e-08\\
1.682	6.29712e-08\\
1.684	6.20329e-08\\
1.686	6.10155e-08\\
1.688	5.99241e-08\\
1.69	5.87639e-08\\
1.692	5.75399e-08\\
1.694	5.62571e-08\\
1.696	5.49202e-08\\
1.698	5.35342e-08\\
1.7	5.21037e-08\\
1.702	5.06332e-08\\
1.704	4.91272e-08\\
1.706	4.75901e-08\\
1.708	4.6026e-08\\
1.71	4.44392e-08\\
1.712	4.28334e-08\\
1.714	4.12127e-08\\
1.716	3.95807e-08\\
1.718	3.7941e-08\\
1.72	3.6297e-08\\
1.722	3.4652e-08\\
1.724	3.30093e-08\\
1.726	3.13718e-08\\
1.728	2.97424e-08\\
1.73	2.8124e-08\\
1.732	2.6519e-08\\
1.734	2.49301e-08\\
1.736	2.33595e-08\\
1.738	2.18095e-08\\
1.74	2.0282e-08\\
1.742	1.87791e-08\\
1.744	1.7831e-08\\
1.746	1.81001e-08\\
1.748	1.83229e-08\\
1.75	1.85007e-08\\
1.752	1.86351e-08\\
1.754	1.87274e-08\\
1.756	1.87792e-08\\
1.758	1.8792e-08\\
1.76	1.87674e-08\\
1.762	1.87068e-08\\
1.764	1.86119e-08\\
1.766	1.84841e-08\\
1.768	1.83251e-08\\
1.77	1.81363e-08\\
1.772	1.79194e-08\\
1.774	1.76758e-08\\
1.776	1.74071e-08\\
1.778	1.71148e-08\\
1.78	1.68003e-08\\
1.782	1.64652e-08\\
1.784	1.6111e-08\\
1.786	1.57389e-08\\
1.788	1.53505e-08\\
1.79	1.49472e-08\\
1.792	1.47164e-08\\
1.794	1.44706e-08\\
1.796	1.419e-08\\
1.798	1.38763e-08\\
1.8	1.35315e-08\\
1.802	1.31576e-08\\
1.804	1.27565e-08\\
1.806	1.23302e-08\\
1.808	1.18808e-08\\
1.81	1.14101e-08\\
1.812	1.16399e-08\\
1.814	1.18684e-08\\
1.816	1.2068e-08\\
1.818	1.22397e-08\\
1.82	1.23842e-08\\
1.822	1.25026e-08\\
1.824	1.25957e-08\\
1.826	1.26644e-08\\
1.828	1.27097e-08\\
1.83	1.27324e-08\\
1.832	1.27336e-08\\
1.834	1.2714e-08\\
1.836	1.26746e-08\\
1.838	1.26162e-08\\
1.84	1.25399e-08\\
1.842	1.24465e-08\\
1.844	1.23368e-08\\
1.846	1.22117e-08\\
1.848	1.20721e-08\\
1.85	1.19189e-08\\
1.852	1.17527e-08\\
1.854	1.15745e-08\\
1.856	1.1385e-08\\
1.858	1.11851e-08\\
1.86	1.09754e-08\\
1.862	1.07567e-08\\
1.864	1.05298e-08\\
1.866	1.02953e-08\\
1.868	1.0054e-08\\
1.87	9.80646e-09\\
1.872	9.55338e-09\\
1.874	9.29537e-09\\
1.876	9.03305e-09\\
1.878	8.76701e-09\\
1.88	8.49781e-09\\
1.882	8.22601e-09\\
1.884	7.95213e-09\\
1.886	7.67668e-09\\
1.888	7.40015e-09\\
1.89	7.12301e-09\\
1.892	6.84572e-09\\
1.894	6.66974e-09\\
1.896	6.76209e-09\\
1.898	6.8349e-09\\
1.9	6.88854e-09\\
1.902	6.92339e-09\\
1.904	6.93987e-09\\
1.906	6.93844e-09\\
1.908	6.9196e-09\\
1.91	6.88386e-09\\
1.912	6.83178e-09\\
1.914	6.76392e-09\\
1.916	6.68089e-09\\
1.918	6.58331e-09\\
1.92	6.4718e-09\\
1.922	6.34704e-09\\
1.924	6.20967e-09\\
1.926	6.06039e-09\\
1.928	5.89989e-09\\
1.93	5.72886e-09\\
1.932	5.548e-09\\
1.934	5.35804e-09\\
1.936	5.15967e-09\\
1.938	4.95361e-09\\
1.94	4.74058e-09\\
1.942	4.52128e-09\\
1.944	4.2964e-09\\
1.946	4.06666e-09\\
1.948	3.83274e-09\\
1.95	3.59531e-09\\
1.952	3.35505e-09\\
1.954	3.11262e-09\\
1.956	2.86865e-09\\
1.958	2.70792e-09\\
1.96	2.60295e-09\\
1.962	2.49703e-09\\
1.964	2.39043e-09\\
1.966	2.28339e-09\\
1.968	2.17617e-09\\
1.97	2.06901e-09\\
1.972	1.96213e-09\\
1.974	1.85575e-09\\
1.976	1.75008e-09\\
1.978	1.64531e-09\\
1.98	1.61906e-09\\
1.982	1.6912e-09\\
1.984	1.75818e-09\\
1.986	1.82009e-09\\
1.988	1.87701e-09\\
1.99	1.92905e-09\\
1.992	1.9763e-09\\
1.994	2.01886e-09\\
1.996	2.05685e-09\\
1.998	2.09036e-09\\
2	2.11951e-09\\
};
\addplot [color=mycolor8, line width=1.0pt, forget plot]
  table[row sep=crcr]{%
-0.024	0\\
-0.022	0\\
-0.02	0\\
-0.018	0\\
-0.016	0\\
-0.014	0\\
-0.012	0\\
-0.01	0\\
-0.008	0\\
-0.006	0\\
-0.004	0\\
-0.002	0\\
0	0\\
0.002	0.0391009\\
0.004	0.074092\\
0.006	0.104794\\
0.008	0.131117\\
0.01	0.153053\\
0.012	0.170674\\
0.014	0.18412\\
0.016	0.193597\\
0.018	0.19936\\
0.02	0.201712\\
0.022	0.200986\\
0.024	0.197545\\
0.026	0.191764\\
0.028	0.184026\\
0.03	0.174713\\
0.032	0.164198\\
0.034	0.152838\\
0.036	0.14097\\
0.038	0.128905\\
0.04	0.118711\\
0.042	0.124424\\
0.044	0.129401\\
0.046	0.133574\\
0.048	0.136894\\
0.05	0.139326\\
0.052	0.14085\\
0.054	0.141463\\
0.056	0.141175\\
0.058	0.140011\\
0.06	0.138009\\
0.062	0.135217\\
0.064	0.131694\\
0.066	0.127509\\
0.068	0.122737\\
0.07	0.117457\\
0.072	0.111755\\
0.074	0.105716\\
0.076	0.0994281\\
0.078	0.0929766\\
0.08	0.0864458\\
0.082	0.083126\\
0.084	0.0846863\\
0.086	0.0861328\\
0.088	0.0874666\\
0.09	0.0886891\\
0.092	0.0898013\\
0.094	0.0908044\\
0.096	0.0916989\\
0.098	0.0924856\\
0.1	0.0931647\\
0.102	0.0937365\\
0.104	0.0942009\\
0.106	0.0945579\\
0.108	0.0948073\\
0.11	0.0949486\\
0.112	0.0961285\\
0.114	0.0975715\\
0.116	0.0988394\\
0.118	0.0999338\\
0.12	0.100857\\
0.122	0.101612\\
0.124	0.102204\\
0.126	0.102635\\
0.128	0.102913\\
0.13	0.103042\\
0.132	0.103029\\
0.134	0.10288\\
0.136	0.102601\\
0.138	0.102199\\
0.14	0.101683\\
0.142	0.101057\\
0.144	0.10033\\
0.146	0.0995079\\
0.148	0.0985981\\
0.15	0.0976072\\
0.152	0.0965415\\
0.154	0.0954075\\
0.156	0.094211\\
0.158	0.0929579\\
0.16	0.0916537\\
0.162	0.0903037\\
0.164	0.0889127\\
0.166	0.0874855\\
0.168	0.0860264\\
0.17	0.0845395\\
0.172	0.0830287\\
0.174	0.0814975\\
0.176	0.0799493\\
0.178	0.0783871\\
0.18	0.0768138\\
0.182	0.0752321\\
0.184	0.0736444\\
0.186	0.0720531\\
0.188	0.0704602\\
0.19	0.0688677\\
0.192	0.0672776\\
0.194	0.0656914\\
0.196	0.0641108\\
0.198	0.0625374\\
0.2	0.0609725\\
0.202	0.0594174\\
0.204	0.0578733\\
0.206	0.0563415\\
0.208	0.0548231\\
0.21	0.053319\\
0.212	0.0518304\\
0.214	0.0503579\\
0.216	0.0489026\\
0.218	0.0474653\\
0.22	0.0460466\\
0.222	0.0446474\\
0.224	0.0432681\\
0.226	0.0419095\\
0.228	0.040572\\
0.23	0.0392562\\
0.232	0.0379624\\
0.234	0.0366911\\
0.236	0.0354426\\
0.238	0.0342171\\
0.24	0.033015\\
0.242	0.0318362\\
0.244	0.0306811\\
0.246	0.0295497\\
0.248	0.0284421\\
0.25	0.0273581\\
0.252	0.0262979\\
0.254	0.0252614\\
0.256	0.0242485\\
0.258	0.0232589\\
0.26	0.0222927\\
0.262	0.0213497\\
0.264	0.0204295\\
0.266	0.0195321\\
0.268	0.0186573\\
0.27	0.0178047\\
0.272	0.0172614\\
0.274	0.0170467\\
0.276	0.0168341\\
0.278	0.0166235\\
0.28	0.0164149\\
0.282	0.0162083\\
0.284	0.0160038\\
0.286	0.0158013\\
0.288	0.0156008\\
0.29	0.0154022\\
0.292	0.0152057\\
0.294	0.015011\\
0.296	0.0148183\\
0.298	0.0146276\\
0.3	0.0144387\\
0.302	0.0142516\\
0.304	0.0140664\\
0.306	0.013883\\
0.308	0.0137014\\
0.31	0.0135215\\
0.312	0.0133433\\
0.314	0.0131669\\
0.316	0.0129921\\
0.318	0.0128189\\
0.32	0.0126474\\
0.322	0.0124775\\
0.324	0.0123091\\
0.326	0.0121423\\
0.328	0.011977\\
0.33	0.0118132\\
0.332	0.0116508\\
0.334	0.01149\\
0.336	0.0113306\\
0.338	0.0111726\\
0.34	0.0110161\\
0.342	0.010861\\
0.344	0.0107073\\
0.346	0.010555\\
0.348	0.0104041\\
0.35	0.0102546\\
0.352	0.0101065\\
0.354	0.0099983\\
0.356	0.0100574\\
0.358	0.0101035\\
0.36	0.0101366\\
0.362	0.0101568\\
0.364	0.0101643\\
0.366	0.0101592\\
0.368	0.0101418\\
0.37	0.0101122\\
0.372	0.0100708\\
0.374	0.0100178\\
0.376	0.00995338\\
0.378	0.00987799\\
0.38	0.00979191\\
0.382	0.00969548\\
0.384	0.00958907\\
0.386	0.00947305\\
0.388	0.00934779\\
0.39	0.00921369\\
0.392	0.00907116\\
0.394	0.00892061\\
0.396	0.00876246\\
0.398	0.00859712\\
0.4	0.00842504\\
0.402	0.00824665\\
0.404	0.00806239\\
0.406	0.00787268\\
0.408	0.00767798\\
0.41	0.00747873\\
0.412	0.00727536\\
0.414	0.00706831\\
0.416	0.00685802\\
0.418	0.00664492\\
0.42	0.00642944\\
0.422	0.006212\\
0.424	0.00599301\\
0.426	0.00577289\\
0.428	0.00555636\\
0.43	0.0054652\\
0.432	0.00537542\\
0.434	0.005287\\
0.436	0.00519993\\
0.438	0.00511418\\
0.44	0.00502976\\
0.442	0.00494663\\
0.444	0.00486479\\
0.446	0.00478422\\
0.448	0.00470489\\
0.45	0.00462681\\
0.452	0.00454993\\
0.454	0.00447426\\
0.456	0.00439978\\
0.458	0.00432645\\
0.46	0.00425428\\
0.462	0.00418325\\
0.464	0.00411332\\
0.466	0.0040445\\
0.468	0.00397676\\
0.47	0.00391009\\
0.472	0.00384447\\
0.474	0.00377988\\
0.476	0.00371631\\
0.478	0.00365374\\
0.48	0.00359216\\
0.482	0.00353155\\
0.484	0.00347189\\
0.486	0.00341318\\
0.488	0.0033554\\
0.49	0.00329853\\
0.492	0.00324255\\
0.494	0.00318746\\
0.496	0.00313325\\
0.498	0.00309312\\
0.5	0.00311428\\
0.502	0.00313217\\
0.504	0.0031468\\
0.506	0.00315821\\
0.508	0.00316644\\
0.51	0.00317153\\
0.512	0.00317354\\
0.514	0.00317251\\
0.516	0.00316849\\
0.518	0.00316156\\
0.52	0.00315177\\
0.522	0.00313919\\
0.524	0.0031239\\
0.526	0.00310596\\
0.528	0.00308545\\
0.53	0.00306246\\
0.532	0.00303706\\
0.534	0.00300933\\
0.536	0.00297936\\
0.538	0.00294724\\
0.54	0.00291306\\
0.542	0.0028769\\
0.544	0.00283886\\
0.546	0.00279903\\
0.548	0.0027575\\
0.55	0.00271436\\
0.552	0.0026697\\
0.554	0.00262363\\
0.556	0.00257622\\
0.558	0.00252759\\
0.56	0.00247781\\
0.562	0.00242698\\
0.564	0.00237519\\
0.566	0.00232254\\
0.568	0.00226911\\
0.57	0.002215\\
0.572	0.00216028\\
0.574	0.00210506\\
0.576	0.0020494\\
0.578	0.0019934\\
0.58	0.00193714\\
0.582	0.0018807\\
0.584	0.00182416\\
0.586	0.00176759\\
0.588	0.00171107\\
0.59	0.00165467\\
0.592	0.00159846\\
0.594	0.00154252\\
0.596	0.0014869\\
0.598	0.00143167\\
0.6	0.00137689\\
0.602	0.00132263\\
0.604	0.00126893\\
0.606	0.00121586\\
0.608	0.00116346\\
0.61	0.00113901\\
0.612	0.00111845\\
0.614	0.00109825\\
0.616	0.0010784\\
0.618	0.0010589\\
0.62	0.00103974\\
0.622	0.00102092\\
0.624	0.00100242\\
0.626	0.000984255\\
0.628	0.000966405\\
0.63	0.000948869\\
0.632	0.000948442\\
0.634	0.000970038\\
0.636	0.000989927\\
0.638	0.0010081\\
0.64	0.00102457\\
0.642	0.00103932\\
0.644	0.00105238\\
0.646	0.00106373\\
0.648	0.00107341\\
0.65	0.00108142\\
0.652	0.00108779\\
0.654	0.00109254\\
0.656	0.00109569\\
0.658	0.00109727\\
0.66	0.00109731\\
0.662	0.00109584\\
0.664	0.00109291\\
0.666	0.00108855\\
0.668	0.00108279\\
0.67	0.00107568\\
0.672	0.00106726\\
0.674	0.00105757\\
0.676	0.00104667\\
0.678	0.0010346\\
0.68	0.00102141\\
0.682	0.00100714\\
0.684	0.000991852\\
0.686	0.000975593\\
0.688	0.000958416\\
0.69	0.000940373\\
0.692	0.000921516\\
0.694	0.0009019\\
0.696	0.000881577\\
0.698	0.000860601\\
0.7	0.000839025\\
0.702	0.000816902\\
0.704	0.000794286\\
0.706	0.000771227\\
0.708	0.00074778\\
0.71	0.000723994\\
0.712	0.000699921\\
0.714	0.000675611\\
0.716	0.000651112\\
0.718	0.000626474\\
0.72	0.000601743\\
0.722	0.000576965\\
0.724	0.000552186\\
0.726	0.000527449\\
0.728	0.000502796\\
0.73	0.00047827\\
0.732	0.00045391\\
0.734	0.000429754\\
0.736	0.00040584\\
0.738	0.000382203\\
0.74	0.000358876\\
0.742	0.000339024\\
0.744	0.000332966\\
0.746	0.000327025\\
0.748	0.000321199\\
0.75	0.000315485\\
0.752	0.000309881\\
0.754	0.000304386\\
0.756	0.000298996\\
0.758	0.00029371\\
0.76	0.000288525\\
0.762	0.00028344\\
0.764	0.000278452\\
0.766	0.000273559\\
0.768	0.000268759\\
0.77	0.000264051\\
0.772	0.000259432\\
0.774	0.0002549\\
0.776	0.000250454\\
0.778	0.000246091\\
0.78	0.000241811\\
0.782	0.00023761\\
0.784	0.000233487\\
0.786	0.000229441\\
0.788	0.00022547\\
0.79	0.000221572\\
0.792	0.000217746\\
0.794	0.00021399\\
0.796	0.000210302\\
0.798	0.000206681\\
0.8	0.000203126\\
0.802	0.000199635\\
0.804	0.000196207\\
0.806	0.000192839\\
0.808	0.000189532\\
0.81	0.000186284\\
0.812	0.000183093\\
0.814	0.000179959\\
0.816	0.000176879\\
0.818	0.000177514\\
0.82	0.000180707\\
0.822	0.000183664\\
0.824	0.000186384\\
0.826	0.000188868\\
0.828	0.000191116\\
0.83	0.00019313\\
0.832	0.00019491\\
0.834	0.000196459\\
0.836	0.000197778\\
0.838	0.000198871\\
0.84	0.00019974\\
0.842	0.000200388\\
0.844	0.000200819\\
0.846	0.000201035\\
0.848	0.000201042\\
0.85	0.000200844\\
0.852	0.000200444\\
0.854	0.000199848\\
0.856	0.00019906\\
0.858	0.000198086\\
0.86	0.00019693\\
0.862	0.000195599\\
0.864	0.000194097\\
0.866	0.00019243\\
0.868	0.000190604\\
0.87	0.000188626\\
0.872	0.000186501\\
0.874	0.000184235\\
0.876	0.000181834\\
0.878	0.000179305\\
0.88	0.000176655\\
0.882	0.000173888\\
0.884	0.000171013\\
0.886	0.000168035\\
0.888	0.00016496\\
0.89	0.000161796\\
0.892	0.000158548\\
0.894	0.000155223\\
0.896	0.000151827\\
0.898	0.000148367\\
0.9	0.000144848\\
0.902	0.000141278\\
0.904	0.000137662\\
0.906	0.000134006\\
0.908	0.000130316\\
0.91	0.000126599\\
0.912	0.000122859\\
0.914	0.000119103\\
0.916	0.000115336\\
0.918	0.000111564\\
0.92	0.000107792\\
0.922	0.000104025\\
0.924	0.000100268\\
0.926	9.65257e-05\\
0.928	9.28037e-05\\
0.93	8.91062e-05\\
0.932	8.63295e-05\\
0.934	8.78749e-05\\
0.936	8.92487e-05\\
0.938	9.04525e-05\\
0.94	9.14884e-05\\
0.942	9.2359e-05\\
0.944	9.30669e-05\\
0.946	9.36152e-05\\
0.948	9.4007e-05\\
0.95	9.4246e-05\\
0.952	9.43359e-05\\
0.954	9.42805e-05\\
0.956	9.40842e-05\\
0.958	9.37511e-05\\
0.96	9.32858e-05\\
0.962	9.26929e-05\\
0.964	9.19771e-05\\
0.966	9.11435e-05\\
0.968	9.01968e-05\\
0.97	8.91423e-05\\
0.972	8.79851e-05\\
0.974	8.67303e-05\\
0.976	8.53831e-05\\
0.978	8.3949e-05\\
0.98	8.24331e-05\\
0.982	8.08407e-05\\
0.984	7.91772e-05\\
0.986	7.74478e-05\\
0.988	7.56577e-05\\
0.99	7.38122e-05\\
0.992	7.19164e-05\\
0.994	6.99754e-05\\
0.996	6.79941e-05\\
0.998	6.59776e-05\\
1	6.39307e-05\\
1.002	6.18582e-05\\
1.004	5.97646e-05\\
1.006	5.76547e-05\\
1.008	5.55327e-05\\
1.01	5.34031e-05\\
1.012	5.127e-05\\
1.014	4.91375e-05\\
1.016	4.70095e-05\\
1.018	4.48898e-05\\
1.02	4.2782e-05\\
1.022	4.06897e-05\\
1.024	3.86161e-05\\
1.026	3.65646e-05\\
1.028	3.45382e-05\\
1.03	3.25397e-05\\
1.032	3.05719e-05\\
1.034	2.86374e-05\\
1.036	2.67387e-05\\
1.038	2.49678e-05\\
1.04	2.45319e-05\\
1.042	2.41044e-05\\
1.044	2.36852e-05\\
1.046	2.32742e-05\\
1.048	2.28713e-05\\
1.05	2.24762e-05\\
1.052	2.20889e-05\\
1.054	2.17091e-05\\
1.056	2.13368e-05\\
1.058	2.09718e-05\\
1.06	2.0614e-05\\
1.062	2.02632e-05\\
1.064	1.99193e-05\\
1.066	1.95822e-05\\
1.068	1.92516e-05\\
1.07	1.89275e-05\\
1.072	1.86098e-05\\
1.074	1.82982e-05\\
1.076	1.79927e-05\\
1.078	1.76931e-05\\
1.08	1.73993e-05\\
1.082	1.71111e-05\\
1.084	1.68285e-05\\
1.086	1.65513e-05\\
1.088	1.62793e-05\\
1.09	1.60125e-05\\
1.092	1.57507e-05\\
1.094	1.54938e-05\\
1.096	1.52417e-05\\
1.098	1.49943e-05\\
1.1	1.47514e-05\\
1.102	1.45129e-05\\
1.104	1.42787e-05\\
1.106	1.40487e-05\\
1.108	1.38229e-05\\
1.11	1.3601e-05\\
1.112	1.3383e-05\\
1.114	1.31688e-05\\
1.116	1.29583e-05\\
1.118	1.27513e-05\\
1.12	1.25479e-05\\
1.122	1.23479e-05\\
1.124	1.21512e-05\\
1.126	1.19577e-05\\
1.128	1.17674e-05\\
1.13	1.15801e-05\\
1.132	1.13958e-05\\
1.134	1.12144e-05\\
1.136	1.10359e-05\\
1.138	1.08601e-05\\
1.14	1.0687e-05\\
1.142	1.05166e-05\\
1.144	1.0579e-05\\
1.146	1.07901e-05\\
1.148	1.09838e-05\\
1.15	1.11604e-05\\
1.152	1.13199e-05\\
1.154	1.14625e-05\\
1.156	1.15884e-05\\
1.158	1.16977e-05\\
1.16	1.17907e-05\\
1.162	1.18676e-05\\
1.164	1.19288e-05\\
1.166	1.19746e-05\\
1.168	1.20052e-05\\
1.17	1.2021e-05\\
1.172	1.20224e-05\\
1.174	1.20097e-05\\
1.176	1.19833e-05\\
1.178	1.19437e-05\\
1.18	1.18912e-05\\
1.182	1.18263e-05\\
1.184	1.17493e-05\\
1.186	1.16609e-05\\
1.188	1.15613e-05\\
1.19	1.1451e-05\\
1.192	1.13306e-05\\
1.194	1.12005e-05\\
1.196	1.10611e-05\\
1.198	1.09129e-05\\
1.2	1.07565e-05\\
1.202	1.05922e-05\\
1.204	1.04205e-05\\
1.206	1.02419e-05\\
1.208	1.00569e-05\\
1.21	9.86597e-06\\
1.212	9.6695e-06\\
1.214	9.46799e-06\\
1.216	9.26188e-06\\
1.218	9.05161e-06\\
1.22	8.83762e-06\\
1.222	8.62035e-06\\
1.224	8.40022e-06\\
1.226	8.17763e-06\\
1.228	7.953e-06\\
1.23	7.72672e-06\\
1.232	7.49918e-06\\
1.234	7.27076e-06\\
1.236	7.04182e-06\\
1.238	6.81272e-06\\
1.24	6.58381e-06\\
1.242	6.36313e-06\\
1.244	6.37039e-06\\
1.246	6.36767e-06\\
1.248	6.35528e-06\\
1.25	6.33354e-06\\
1.252	6.30277e-06\\
1.254	6.26332e-06\\
1.256	6.21553e-06\\
1.258	6.15975e-06\\
1.26	6.09636e-06\\
1.262	6.02571e-06\\
1.264	5.94817e-06\\
1.266	5.86411e-06\\
1.268	5.77391e-06\\
1.27	5.67794e-06\\
1.272	5.57657e-06\\
1.274	5.47018e-06\\
1.276	5.35914e-06\\
1.278	5.24383e-06\\
1.28	5.12459e-06\\
1.282	5.00181e-06\\
1.284	4.87582e-06\\
1.286	4.747e-06\\
1.288	4.61568e-06\\
1.29	4.4822e-06\\
1.292	4.34689e-06\\
1.294	4.21008e-06\\
1.296	4.07209e-06\\
1.298	3.93322e-06\\
1.3	3.79378e-06\\
1.302	3.65404e-06\\
1.304	3.5143e-06\\
1.306	3.37483e-06\\
1.308	3.23588e-06\\
1.31	3.0977e-06\\
1.312	2.96054e-06\\
1.314	2.82461e-06\\
1.316	2.69015e-06\\
1.318	2.55735e-06\\
1.32	2.42641e-06\\
1.322	2.29752e-06\\
1.324	2.17085e-06\\
1.326	2.04656e-06\\
1.328	1.95422e-06\\
1.33	1.91773e-06\\
1.332	1.88202e-06\\
1.334	1.84708e-06\\
1.336	1.8129e-06\\
1.338	1.77946e-06\\
1.34	1.74675e-06\\
1.342	1.71475e-06\\
1.344	1.68345e-06\\
1.346	1.65284e-06\\
1.348	1.6229e-06\\
1.35	1.59362e-06\\
1.352	1.56499e-06\\
1.354	1.53699e-06\\
1.356	1.50961e-06\\
1.358	1.48284e-06\\
1.36	1.45666e-06\\
1.362	1.43107e-06\\
1.364	1.40604e-06\\
1.366	1.38156e-06\\
1.368	1.35763e-06\\
1.37	1.33422e-06\\
1.372	1.31133e-06\\
1.374	1.28895e-06\\
1.376	1.26705e-06\\
1.378	1.24564e-06\\
1.38	1.22469e-06\\
1.382	1.20419e-06\\
1.384	1.18414e-06\\
1.386	1.16452e-06\\
1.388	1.14533e-06\\
1.39	1.12654e-06\\
1.392	1.10814e-06\\
1.394	1.09014e-06\\
1.396	1.0725e-06\\
1.398	1.05524e-06\\
1.4	1.03832e-06\\
1.402	1.02176e-06\\
1.404	1.00552e-06\\
1.406	9.89611e-07\\
1.408	9.74015e-07\\
1.41	9.58723e-07\\
1.412	9.43727e-07\\
1.414	9.29017e-07\\
1.416	9.14584e-07\\
1.418	9.00421e-07\\
1.42	8.86517e-07\\
1.422	8.72866e-07\\
1.424	8.59459e-07\\
1.426	8.46288e-07\\
1.428	8.33346e-07\\
1.43	8.20625e-07\\
1.432	8.08119e-07\\
1.434	7.95821e-07\\
1.436	7.83723e-07\\
1.438	7.71819e-07\\
1.44	7.60104e-07\\
1.442	7.48571e-07\\
1.444	7.37214e-07\\
1.446	7.26029e-07\\
1.448	7.15009e-07\\
1.45	7.0415e-07\\
1.452	6.93446e-07\\
1.454	6.82894e-07\\
1.456	6.72488e-07\\
1.458	6.62224e-07\\
1.46	6.52099e-07\\
1.462	6.42108e-07\\
1.464	6.32248e-07\\
1.466	6.22515e-07\\
1.468	6.12907e-07\\
1.47	6.08056e-07\\
1.472	6.1512e-07\\
1.474	6.21335e-07\\
1.476	6.26713e-07\\
1.478	6.31265e-07\\
1.48	6.35007e-07\\
1.482	6.37954e-07\\
1.484	6.40122e-07\\
1.486	6.41528e-07\\
1.488	6.42191e-07\\
1.49	6.42129e-07\\
1.492	6.41364e-07\\
1.494	6.39915e-07\\
1.496	6.37803e-07\\
1.498	6.35051e-07\\
1.5	6.3168e-07\\
1.502	6.27714e-07\\
1.504	6.23176e-07\\
1.506	6.18088e-07\\
1.508	6.12474e-07\\
1.51	6.06359e-07\\
1.512	5.99766e-07\\
1.514	5.92719e-07\\
1.516	5.85241e-07\\
1.518	5.77358e-07\\
1.52	5.69093e-07\\
1.522	5.6047e-07\\
1.524	5.51512e-07\\
1.526	5.42242e-07\\
1.528	5.32685e-07\\
1.53	5.22863e-07\\
1.532	5.12798e-07\\
1.534	5.02513e-07\\
1.536	4.9203e-07\\
1.538	4.8137e-07\\
1.54	4.70554e-07\\
1.542	4.59603e-07\\
1.544	4.48537e-07\\
1.546	4.37376e-07\\
1.548	4.26139e-07\\
1.55	4.14844e-07\\
1.552	4.0351e-07\\
1.554	3.92154e-07\\
1.556	3.80794e-07\\
1.558	3.69445e-07\\
1.56	3.58123e-07\\
1.562	3.46844e-07\\
1.564	3.35622e-07\\
1.566	3.24472e-07\\
1.568	3.13406e-07\\
1.57	3.02438e-07\\
1.572	2.9158e-07\\
1.574	2.80843e-07\\
1.576	2.70239e-07\\
1.578	2.59777e-07\\
1.58	2.49467e-07\\
1.582	2.3932e-07\\
1.584	2.29343e-07\\
1.586	2.19544e-07\\
1.588	2.10978e-07\\
1.59	2.06868e-07\\
1.592	2.02836e-07\\
1.594	1.98879e-07\\
1.596	1.94998e-07\\
1.598	1.91191e-07\\
1.6	1.87458e-07\\
1.602	1.83799e-07\\
1.604	1.80212e-07\\
1.606	1.76696e-07\\
1.608	1.73252e-07\\
1.61	1.69876e-07\\
1.612	1.6657e-07\\
1.614	1.63332e-07\\
1.616	1.60161e-07\\
1.618	1.57055e-07\\
1.62	1.54015e-07\\
1.622	1.5104e-07\\
1.624	1.48127e-07\\
1.626	1.45277e-07\\
1.628	1.42487e-07\\
1.63	1.39758e-07\\
1.632	1.37088e-07\\
1.634	1.34477e-07\\
1.636	1.31922e-07\\
1.638	1.29424e-07\\
1.64	1.2698e-07\\
1.642	1.24591e-07\\
1.644	1.22254e-07\\
1.646	1.1997e-07\\
1.648	1.17736e-07\\
1.65	1.15553e-07\\
1.652	1.13418e-07\\
1.654	1.11331e-07\\
1.656	1.09291e-07\\
1.658	1.07297e-07\\
1.66	1.05348e-07\\
1.662	1.03442e-07\\
1.664	1.01579e-07\\
1.666	9.97586e-08\\
1.668	9.79787e-08\\
1.67	9.62387e-08\\
1.672	9.45377e-08\\
1.674	9.28748e-08\\
1.676	9.12489e-08\\
1.678	8.96593e-08\\
1.68	8.81049e-08\\
1.682	8.6585e-08\\
1.684	8.50985e-08\\
1.686	8.36447e-08\\
1.688	8.22227e-08\\
1.69	8.08317e-08\\
1.692	7.94708e-08\\
1.694	7.81393e-08\\
1.696	7.68362e-08\\
1.698	7.5561e-08\\
1.7	7.43127e-08\\
1.702	7.30907e-08\\
1.704	7.18942e-08\\
1.706	7.07224e-08\\
1.708	6.95748e-08\\
1.71	6.84505e-08\\
1.712	6.7349e-08\\
1.714	6.62695e-08\\
1.716	6.52115e-08\\
1.718	6.41742e-08\\
1.72	6.31571e-08\\
1.722	6.21596e-08\\
1.724	6.11812e-08\\
1.726	6.02211e-08\\
1.728	5.9279e-08\\
1.73	5.83542e-08\\
1.732	5.74463e-08\\
1.734	5.65547e-08\\
1.736	5.56789e-08\\
1.738	5.48185e-08\\
1.74	5.3973e-08\\
1.742	5.3142e-08\\
1.744	5.23251e-08\\
1.746	5.15217e-08\\
1.748	5.07315e-08\\
1.75	4.99542e-08\\
1.752	4.91893e-08\\
1.754	4.84365e-08\\
1.756	4.76955e-08\\
1.758	4.69658e-08\\
1.76	4.62472e-08\\
1.762	4.55394e-08\\
1.764	4.4842e-08\\
1.766	4.41548e-08\\
1.768	4.34775e-08\\
1.77	4.28099e-08\\
1.772	4.21516e-08\\
1.774	4.15025e-08\\
1.776	4.08623e-08\\
1.778	4.02308e-08\\
1.78	3.96077e-08\\
1.782	3.8993e-08\\
1.784	3.83863e-08\\
1.786	3.77876e-08\\
1.788	3.71966e-08\\
1.79	3.66132e-08\\
1.792	3.60372e-08\\
1.794	3.54685e-08\\
1.796	3.4907e-08\\
1.798	3.43525e-08\\
1.8	3.38048e-08\\
1.802	3.3264e-08\\
1.804	3.27298e-08\\
1.806	3.22022e-08\\
1.808	3.16811e-08\\
1.81	3.11663e-08\\
1.812	3.06579e-08\\
1.814	3.01557e-08\\
1.816	2.96597e-08\\
1.818	2.91697e-08\\
1.82	2.86858e-08\\
1.822	2.82863e-08\\
1.824	2.81405e-08\\
1.826	2.7976e-08\\
1.828	2.77935e-08\\
1.83	2.75937e-08\\
1.832	2.73774e-08\\
1.834	2.71453e-08\\
1.836	2.6898e-08\\
1.838	2.66363e-08\\
1.84	2.6361e-08\\
1.842	2.60726e-08\\
1.844	2.57721e-08\\
1.846	2.54599e-08\\
1.848	2.5137e-08\\
1.85	2.48038e-08\\
1.852	2.44612e-08\\
1.854	2.41097e-08\\
1.856	2.375e-08\\
1.858	2.33829e-08\\
1.86	2.30088e-08\\
1.862	2.26285e-08\\
1.864	2.22425e-08\\
1.866	2.18514e-08\\
1.868	2.14558e-08\\
1.87	2.10563e-08\\
1.872	2.06534e-08\\
1.874	2.02476e-08\\
1.876	1.98396e-08\\
1.878	1.94296e-08\\
1.88	1.90184e-08\\
1.882	1.86062e-08\\
1.884	1.81937e-08\\
1.886	1.77812e-08\\
1.888	1.73691e-08\\
1.89	1.69578e-08\\
1.892	1.65478e-08\\
1.894	1.61394e-08\\
1.896	1.5733e-08\\
1.898	1.53289e-08\\
1.9	1.49275e-08\\
1.902	1.45291e-08\\
1.904	1.41339e-08\\
1.906	1.37423e-08\\
1.908	1.33545e-08\\
1.91	1.29708e-08\\
1.912	1.25914e-08\\
1.914	1.22166e-08\\
1.916	1.18803e-08\\
1.918	1.16525e-08\\
1.92	1.14291e-08\\
1.922	1.121e-08\\
1.924	1.09952e-08\\
1.926	1.07846e-08\\
1.928	1.05782e-08\\
1.93	1.03759e-08\\
1.932	1.01777e-08\\
1.934	9.98346e-09\\
1.936	9.79318e-09\\
1.938	9.60678e-09\\
1.94	9.42422e-09\\
1.942	9.24543e-09\\
1.944	9.07035e-09\\
1.946	8.89892e-09\\
1.948	8.73108e-09\\
1.95	8.56677e-09\\
1.952	8.40592e-09\\
1.954	8.24849e-09\\
1.956	8.0944e-09\\
1.958	7.94359e-09\\
1.96	7.796e-09\\
1.962	7.65158e-09\\
1.964	7.51025e-09\\
1.966	7.37196e-09\\
1.968	7.23665e-09\\
1.97	7.10425e-09\\
1.972	6.97471e-09\\
1.974	6.84796e-09\\
1.976	6.72394e-09\\
1.978	6.6026e-09\\
1.98	6.48388e-09\\
1.982	6.36772e-09\\
1.984	6.25405e-09\\
1.986	6.14283e-09\\
1.988	6.03399e-09\\
1.99	5.92748e-09\\
1.992	5.82325e-09\\
1.994	5.72124e-09\\
1.996	5.62139e-09\\
1.998	5.52366e-09\\
2	5.42799e-09\\
};
\addplot [color=mycolor9, line width=1.0pt, forget plot]
  table[row sep=crcr]{%
-0.024	0\\
-0.022	0\\
-0.02	0\\
-0.018	0\\
-0.016	0\\
-0.014	0\\
-0.012	0\\
-0.01	0\\
-0.008	0\\
-0.006	0\\
-0.004	0\\
-0.002	0\\
0	0\\
0.002	0.0391008\\
0.004	0.0740917\\
0.006	0.104793\\
0.008	0.131113\\
0.01	0.153044\\
0.012	0.170655\\
0.014	0.184088\\
0.016	0.193544\\
0.018	0.199279\\
0.02	0.201592\\
0.022	0.200817\\
0.024	0.197313\\
0.026	0.191454\\
0.028	0.183623\\
0.03	0.174199\\
0.032	0.163553\\
0.034	0.152043\\
0.036	0.140005\\
0.038	0.127748\\
0.04	0.115554\\
0.042	0.114998\\
0.044	0.119402\\
0.046	0.12306\\
0.048	0.125929\\
0.05	0.127986\\
0.052	0.129217\\
0.054	0.129626\\
0.056	0.12923\\
0.058	0.128057\\
0.06	0.126146\\
0.062	0.123548\\
0.064	0.120321\\
0.066	0.116531\\
0.068	0.112247\\
0.07	0.107546\\
0.072	0.102504\\
0.074	0.0972001\\
0.076	0.0917117\\
0.078	0.0861151\\
0.08	0.0819681\\
0.082	0.083472\\
0.084	0.0848464\\
0.086	0.0860938\\
0.088	0.0872167\\
0.09	0.0882177\\
0.092	0.0890993\\
0.094	0.0898642\\
0.096	0.0905147\\
0.098	0.091053\\
0.1	0.0914813\\
0.102	0.0918014\\
0.104	0.0920153\\
0.106	0.0921245\\
0.108	0.0921305\\
0.11	0.0920346\\
0.112	0.0918383\\
0.114	0.0915427\\
0.116	0.0924466\\
0.118	0.0935295\\
0.12	0.094464\\
0.122	0.0952529\\
0.124	0.0958997\\
0.126	0.096408\\
0.128	0.0967822\\
0.13	0.0970267\\
0.132	0.0971465\\
0.134	0.0971467\\
0.136	0.0970325\\
0.138	0.0968094\\
0.14	0.0964828\\
0.142	0.0960582\\
0.144	0.095541\\
0.146	0.0949368\\
0.148	0.0942508\\
0.15	0.0934883\\
0.152	0.0926543\\
0.154	0.0917537\\
0.156	0.0907913\\
0.158	0.0897717\\
0.16	0.0886993\\
0.162	0.0875782\\
0.164	0.0864125\\
0.166	0.0852059\\
0.168	0.0839621\\
0.17	0.0826846\\
0.172	0.0813765\\
0.174	0.0800411\\
0.176	0.0786811\\
0.178	0.0772994\\
0.18	0.0758987\\
0.182	0.0744814\\
0.184	0.0730499\\
0.186	0.0716065\\
0.188	0.0701534\\
0.19	0.0686926\\
0.192	0.0672262\\
0.194	0.0657559\\
0.196	0.0642836\\
0.198	0.062811\\
0.2	0.0613399\\
0.202	0.0598717\\
0.204	0.0584081\\
0.206	0.0569505\\
0.208	0.0555003\\
0.21	0.0540588\\
0.212	0.0526274\\
0.214	0.0512072\\
0.216	0.0497996\\
0.218	0.0484054\\
0.22	0.0470259\\
0.222	0.0456619\\
0.224	0.0443144\\
0.226	0.0429843\\
0.228	0.0416724\\
0.23	0.0403793\\
0.232	0.0391058\\
0.234	0.0378525\\
0.236	0.03662\\
0.238	0.0354086\\
0.24	0.034219\\
0.242	0.0330513\\
0.244	0.0319061\\
0.246	0.0307834\\
0.248	0.0296836\\
0.25	0.0286068\\
0.252	0.0275532\\
0.254	0.0265227\\
0.256	0.0255156\\
0.258	0.0245317\\
0.26	0.023571\\
0.262	0.0226336\\
0.264	0.0217192\\
0.266	0.0208279\\
0.268	0.0199594\\
0.27	0.0193029\\
0.272	0.0187028\\
0.274	0.0181149\\
0.276	0.01754\\
0.278	0.0169791\\
0.28	0.0164326\\
0.282	0.0162373\\
0.284	0.0160437\\
0.286	0.015852\\
0.288	0.015662\\
0.29	0.0154737\\
0.292	0.0152871\\
0.294	0.0151022\\
0.296	0.014919\\
0.298	0.0147375\\
0.3	0.0145575\\
0.302	0.0143792\\
0.304	0.0142025\\
0.306	0.0140274\\
0.308	0.0138538\\
0.31	0.0136817\\
0.312	0.0135112\\
0.314	0.0133422\\
0.316	0.0131746\\
0.318	0.0130086\\
0.32	0.012844\\
0.322	0.0126808\\
0.324	0.0125191\\
0.326	0.0123587\\
0.328	0.0121998\\
0.33	0.0120423\\
0.332	0.0118862\\
0.334	0.0117314\\
0.336	0.0115781\\
0.338	0.011426\\
0.34	0.0112754\\
0.342	0.0111261\\
0.344	0.0109782\\
0.346	0.0108316\\
0.348	0.0106864\\
0.35	0.0105604\\
0.352	0.0106395\\
0.354	0.0107063\\
0.356	0.0107605\\
0.358	0.0108024\\
0.36	0.0108318\\
0.362	0.0108489\\
0.364	0.0108538\\
0.366	0.0108467\\
0.368	0.0108276\\
0.37	0.0107968\\
0.372	0.0107546\\
0.374	0.0107011\\
0.376	0.0106366\\
0.378	0.0105613\\
0.38	0.0104757\\
0.382	0.01038\\
0.384	0.0102745\\
0.386	0.0101595\\
0.388	0.0100356\\
0.39	0.0099029\\
0.392	0.00976193\\
0.394	0.00961303\\
0.396	0.00945663\\
0.398	0.00929311\\
0.4	0.00912289\\
0.402	0.0089464\\
0.404	0.00876404\\
0.406	0.00857625\\
0.408	0.00838344\\
0.41	0.00818606\\
0.412	0.00798451\\
0.414	0.00777922\\
0.416	0.00757063\\
0.418	0.00735914\\
0.42	0.00714517\\
0.422	0.00692913\\
0.424	0.00671144\\
0.426	0.00649248\\
0.428	0.00627265\\
0.43	0.00605234\\
0.432	0.00583192\\
0.434	0.00573369\\
0.436	0.00564755\\
0.438	0.00556259\\
0.44	0.0054788\\
0.442	0.00539618\\
0.444	0.00531471\\
0.446	0.00523437\\
0.448	0.00515515\\
0.45	0.00507703\\
0.452	0.00500001\\
0.454	0.00492407\\
0.456	0.00484918\\
0.458	0.00477535\\
0.46	0.00470256\\
0.462	0.00463078\\
0.464	0.00456001\\
0.466	0.00449024\\
0.468	0.00442145\\
0.47	0.00435363\\
0.472	0.00428676\\
0.474	0.00422083\\
0.476	0.00415583\\
0.478	0.00409175\\
0.48	0.00402857\\
0.482	0.00396629\\
0.484	0.00390488\\
0.486	0.00384434\\
0.488	0.00378466\\
0.49	0.00372582\\
0.492	0.00368769\\
0.494	0.0037009\\
0.496	0.00371119\\
0.498	0.00371856\\
0.5	0.00372304\\
0.502	0.00372465\\
0.504	0.00372344\\
0.506	0.00371943\\
0.508	0.00371267\\
0.51	0.00370321\\
0.512	0.00369108\\
0.514	0.00367635\\
0.516	0.00365906\\
0.518	0.00363928\\
0.52	0.00361707\\
0.522	0.00359248\\
0.524	0.0035656\\
0.526	0.00353647\\
0.528	0.00350519\\
0.53	0.00347182\\
0.532	0.00343643\\
0.534	0.0033991\\
0.536	0.00335991\\
0.538	0.00331893\\
0.54	0.00327626\\
0.542	0.00323197\\
0.544	0.00318614\\
0.546	0.00313885\\
0.548	0.0030902\\
0.55	0.00304025\\
0.552	0.0029891\\
0.554	0.00293684\\
0.556	0.00288353\\
0.558	0.00282928\\
0.56	0.00277415\\
0.562	0.00271824\\
0.564	0.00266163\\
0.566	0.0026044\\
0.568	0.00254662\\
0.57	0.00248838\\
0.572	0.00242976\\
0.574	0.00237083\\
0.576	0.00231167\\
0.578	0.00225236\\
0.58	0.00219297\\
0.582	0.00213356\\
0.584	0.00207422\\
0.586	0.002015\\
0.588	0.00195598\\
0.59	0.00189722\\
0.592	0.00183878\\
0.594	0.00178072\\
0.596	0.00172311\\
0.598	0.00166599\\
0.6	0.00160942\\
0.602	0.00155347\\
0.604	0.00149816\\
0.606	0.00145286\\
0.608	0.00142887\\
0.61	0.00140526\\
0.612	0.00138202\\
0.614	0.00135916\\
0.616	0.00133666\\
0.618	0.00131452\\
0.62	0.00129274\\
0.622	0.00127131\\
0.624	0.00125021\\
0.626	0.00122946\\
0.628	0.00120904\\
0.63	0.00118895\\
0.632	0.00116918\\
0.634	0.00114974\\
0.636	0.00113292\\
0.638	0.00114741\\
0.64	0.00116034\\
0.642	0.00117169\\
0.644	0.00118147\\
0.646	0.00118969\\
0.648	0.00119636\\
0.65	0.0012015\\
0.652	0.00120512\\
0.654	0.00120725\\
0.656	0.0012079\\
0.658	0.00120711\\
0.66	0.00120489\\
0.662	0.00120129\\
0.664	0.00119633\\
0.666	0.00119005\\
0.668	0.00118248\\
0.67	0.00117367\\
0.672	0.00116366\\
0.674	0.00115248\\
0.676	0.00114017\\
0.678	0.0011268\\
0.68	0.00111239\\
0.682	0.00109699\\
0.684	0.00108066\\
0.686	0.00106344\\
0.688	0.00104538\\
0.69	0.00102653\\
0.692	0.00100693\\
0.694	0.000986645\\
0.696	0.000965717\\
0.698	0.000944197\\
0.7	0.000922136\\
0.702	0.000899582\\
0.704	0.000876587\\
0.706	0.000853198\\
0.708	0.000829464\\
0.71	0.000805435\\
0.712	0.000781156\\
0.714	0.000756676\\
0.716	0.00073204\\
0.718	0.000707294\\
0.72	0.00068248\\
0.722	0.000657644\\
0.724	0.000632826\\
0.726	0.000608068\\
0.728	0.00058341\\
0.73	0.000558889\\
0.732	0.000534544\\
0.734	0.00051041\\
0.736	0.000487088\\
0.738	0.000479043\\
0.74	0.000471138\\
0.742	0.000463371\\
0.744	0.000455741\\
0.746	0.000448243\\
0.748	0.000440877\\
0.75	0.00043364\\
0.752	0.00042653\\
0.754	0.000419543\\
0.756	0.000412679\\
0.758	0.000405934\\
0.76	0.000399307\\
0.762	0.000392795\\
0.764	0.000386396\\
0.766	0.000380108\\
0.768	0.000373929\\
0.77	0.000367858\\
0.772	0.000361891\\
0.774	0.000356026\\
0.776	0.000350263\\
0.778	0.000344599\\
0.78	0.000339032\\
0.782	0.00033356\\
0.784	0.000328181\\
0.786	0.000322894\\
0.788	0.000317697\\
0.79	0.000312588\\
0.792	0.000307565\\
0.794	0.000302627\\
0.796	0.000297771\\
0.798	0.000292998\\
0.8	0.000288304\\
0.802	0.000283688\\
0.804	0.00027915\\
0.806	0.000274686\\
0.808	0.000270297\\
0.81	0.000265979\\
0.812	0.000261733\\
0.814	0.000257557\\
0.816	0.000253449\\
0.818	0.000249409\\
0.82	0.000250756\\
0.822	0.000252247\\
0.824	0.000253524\\
0.826	0.000254588\\
0.828	0.000255439\\
0.83	0.000256079\\
0.832	0.00025651\\
0.834	0.000256733\\
0.836	0.000256752\\
0.838	0.000256569\\
0.84	0.000256187\\
0.842	0.00025561\\
0.844	0.00025484\\
0.846	0.000253883\\
0.848	0.000252741\\
0.85	0.00025142\\
0.852	0.000249923\\
0.854	0.000248256\\
0.856	0.000246424\\
0.858	0.000244431\\
0.86	0.000242283\\
0.862	0.000239985\\
0.864	0.000237543\\
0.866	0.000234962\\
0.868	0.000232248\\
0.87	0.000229407\\
0.872	0.000226445\\
0.874	0.000223367\\
0.876	0.00022018\\
0.878	0.00021689\\
0.88	0.000213502\\
0.882	0.000210024\\
0.884	0.00020646\\
0.886	0.000202817\\
0.888	0.000199102\\
0.89	0.000195319\\
0.892	0.000191476\\
0.894	0.000187578\\
0.896	0.000183632\\
0.898	0.000179642\\
0.9	0.000175615\\
0.902	0.000171556\\
0.904	0.000167472\\
0.906	0.000163367\\
0.908	0.000159248\\
0.91	0.000155119\\
0.912	0.000150986\\
0.914	0.000146853\\
0.916	0.000142727\\
0.918	0.000138611\\
0.92	0.00013451\\
0.922	0.00013043\\
0.924	0.000126374\\
0.926	0.000122347\\
0.928	0.000118352\\
0.93	0.000114395\\
0.932	0.000110479\\
0.934	0.000109076\\
0.936	0.000110034\\
0.938	0.000110838\\
0.94	0.000111488\\
0.942	0.000111988\\
0.944	0.000112338\\
0.946	0.000112543\\
0.948	0.000112605\\
0.95	0.000112526\\
0.952	0.000112311\\
0.954	0.000111962\\
0.956	0.000111484\\
0.958	0.00011088\\
0.96	0.000110154\\
0.962	0.000109311\\
0.964	0.000108355\\
0.966	0.000107291\\
0.968	0.000106122\\
0.97	0.000104854\\
0.972	0.000103491\\
0.974	0.000102038\\
0.976	0.0001005\\
0.978	9.88816e-05\\
0.98	9.7188e-05\\
0.982	9.54238e-05\\
0.984	9.35941e-05\\
0.986	9.17038e-05\\
0.988	8.97576e-05\\
0.99	8.77606e-05\\
0.992	8.57174e-05\\
0.994	8.36329e-05\\
0.996	8.15117e-05\\
0.998	7.93585e-05\\
1	7.71779e-05\\
1.002	7.49744e-05\\
1.004	7.27524e-05\\
1.006	7.05161e-05\\
1.008	6.82698e-05\\
1.01	6.60175e-05\\
1.012	6.37634e-05\\
1.014	6.15112e-05\\
1.016	5.92646e-05\\
1.018	5.70274e-05\\
1.02	5.4803e-05\\
1.022	5.25948e-05\\
1.024	5.0406e-05\\
1.026	4.82396e-05\\
1.028	4.60987e-05\\
1.03	4.41942e-05\\
1.032	4.3474e-05\\
1.034	4.27663e-05\\
1.036	4.2071e-05\\
1.038	4.13878e-05\\
1.04	4.07165e-05\\
1.042	4.00571e-05\\
1.044	3.94092e-05\\
1.046	3.87726e-05\\
1.048	3.81473e-05\\
1.05	3.75329e-05\\
1.052	3.69294e-05\\
1.054	3.63364e-05\\
1.056	3.57539e-05\\
1.058	3.51817e-05\\
1.06	3.46196e-05\\
1.062	3.40673e-05\\
1.064	3.35247e-05\\
1.066	3.29917e-05\\
1.068	3.2468e-05\\
1.07	3.19535e-05\\
1.072	3.14481e-05\\
1.074	3.09514e-05\\
1.076	3.04634e-05\\
1.078	2.99839e-05\\
1.08	2.95128e-05\\
1.082	2.90498e-05\\
1.084	2.85948e-05\\
1.086	2.81477e-05\\
1.088	2.77082e-05\\
1.09	2.72763e-05\\
1.092	2.68517e-05\\
1.094	2.64343e-05\\
1.096	2.60241e-05\\
1.098	2.56207e-05\\
1.1	2.52242e-05\\
1.102	2.48342e-05\\
1.104	2.44508e-05\\
1.106	2.40737e-05\\
1.108	2.37028e-05\\
1.11	2.33381e-05\\
1.112	2.29793e-05\\
1.114	2.26263e-05\\
1.116	2.2279e-05\\
1.118	2.19373e-05\\
1.12	2.16011e-05\\
1.122	2.12703e-05\\
1.124	2.09446e-05\\
1.126	2.06241e-05\\
1.128	2.03087e-05\\
1.13	1.99981e-05\\
1.132	1.96923e-05\\
1.134	1.93913e-05\\
1.136	1.90948e-05\\
1.138	1.88029e-05\\
1.14	1.85154e-05\\
1.142	1.82322e-05\\
1.144	1.79532e-05\\
1.146	1.76784e-05\\
1.148	1.74176e-05\\
1.15	1.74663e-05\\
1.152	1.75007e-05\\
1.154	1.75207e-05\\
1.156	1.75266e-05\\
1.158	1.75185e-05\\
1.16	1.74966e-05\\
1.162	1.74613e-05\\
1.164	1.74126e-05\\
1.166	1.7351e-05\\
1.168	1.72767e-05\\
1.17	1.71899e-05\\
1.172	1.70912e-05\\
1.174	1.69807e-05\\
1.176	1.68588e-05\\
1.178	1.6726e-05\\
1.18	1.65825e-05\\
1.182	1.64288e-05\\
1.184	1.62653e-05\\
1.186	1.60924e-05\\
1.188	1.59105e-05\\
1.19	1.572e-05\\
1.192	1.55214e-05\\
1.194	1.5315e-05\\
1.196	1.51014e-05\\
1.198	1.48808e-05\\
1.2	1.46539e-05\\
1.202	1.44209e-05\\
1.204	1.41824e-05\\
1.206	1.39388e-05\\
1.208	1.36904e-05\\
1.21	1.34378e-05\\
1.212	1.31814e-05\\
1.214	1.29214e-05\\
1.216	1.26585e-05\\
1.218	1.2393e-05\\
1.22	1.21252e-05\\
1.222	1.18556e-05\\
1.224	1.15846e-05\\
1.226	1.13125e-05\\
1.228	1.10397e-05\\
1.23	1.07665e-05\\
1.232	1.04934e-05\\
1.234	1.02206e-05\\
1.236	9.94851e-06\\
1.238	9.67742e-06\\
1.24	9.40765e-06\\
1.242	9.13948e-06\\
1.244	9.03486e-06\\
1.246	8.99402e-06\\
1.248	8.94386e-06\\
1.25	8.88471e-06\\
1.252	8.81686e-06\\
1.254	8.74067e-06\\
1.256	8.65646e-06\\
1.258	8.56459e-06\\
1.26	8.46539e-06\\
1.262	8.35924e-06\\
1.264	8.2465e-06\\
1.266	8.12753e-06\\
1.268	8.00269e-06\\
1.27	7.87237e-06\\
1.272	7.73693e-06\\
1.274	7.59674e-06\\
1.276	7.45217e-06\\
1.278	7.30359e-06\\
1.28	7.15137e-06\\
1.282	6.99585e-06\\
1.284	6.83741e-06\\
1.286	6.67639e-06\\
1.288	6.51315e-06\\
1.29	6.34801e-06\\
1.292	6.18132e-06\\
1.294	6.0134e-06\\
1.296	5.84457e-06\\
1.298	5.67514e-06\\
1.3	5.50542e-06\\
1.302	5.3357e-06\\
1.304	5.16625e-06\\
1.306	4.99737e-06\\
1.308	4.8293e-06\\
1.31	4.66231e-06\\
1.312	4.53173e-06\\
1.314	4.45512e-06\\
1.316	4.37987e-06\\
1.318	4.30596e-06\\
1.32	4.23337e-06\\
1.322	4.16207e-06\\
1.324	4.09206e-06\\
1.326	4.02331e-06\\
1.328	3.95581e-06\\
1.33	3.88953e-06\\
1.332	3.82446e-06\\
1.334	3.76058e-06\\
1.336	3.69787e-06\\
1.338	3.63631e-06\\
1.34	3.57588e-06\\
1.342	3.51656e-06\\
1.344	3.45834e-06\\
1.346	3.4012e-06\\
1.348	3.34511e-06\\
1.35	3.29007e-06\\
1.352	3.23605e-06\\
1.354	3.18303e-06\\
1.356	3.131e-06\\
1.358	3.07993e-06\\
1.36	3.02982e-06\\
1.362	2.98063e-06\\
1.364	2.93237e-06\\
1.366	2.88499e-06\\
1.368	2.8385e-06\\
1.37	2.79288e-06\\
1.372	2.74809e-06\\
1.374	2.70414e-06\\
1.376	2.66099e-06\\
1.378	2.61865e-06\\
1.38	2.57708e-06\\
1.382	2.53627e-06\\
1.384	2.49621e-06\\
1.386	2.45688e-06\\
1.388	2.41826e-06\\
1.39	2.38035e-06\\
1.392	2.34311e-06\\
1.394	2.30655e-06\\
1.396	2.27064e-06\\
1.398	2.23536e-06\\
1.4	2.20072e-06\\
1.402	2.16668e-06\\
1.404	2.13324e-06\\
1.406	2.10038e-06\\
1.408	2.06809e-06\\
1.41	2.03636e-06\\
1.412	2.00517e-06\\
1.414	1.97451e-06\\
1.416	1.94437e-06\\
1.418	1.91474e-06\\
1.42	1.88559e-06\\
1.422	1.85693e-06\\
1.424	1.82874e-06\\
1.426	1.80102e-06\\
1.428	1.77373e-06\\
1.43	1.74689e-06\\
1.432	1.72048e-06\\
1.434	1.69448e-06\\
1.436	1.66889e-06\\
1.438	1.6437e-06\\
1.44	1.6189e-06\\
1.442	1.59448e-06\\
1.444	1.57043e-06\\
1.446	1.54674e-06\\
1.448	1.52341e-06\\
1.45	1.50043e-06\\
1.452	1.47778e-06\\
1.454	1.45547e-06\\
1.456	1.43348e-06\\
1.458	1.41181e-06\\
1.46	1.39045e-06\\
1.462	1.3694e-06\\
1.464	1.34864e-06\\
1.466	1.32818e-06\\
1.468	1.308e-06\\
1.47	1.28811e-06\\
1.472	1.26849e-06\\
1.474	1.24914e-06\\
1.476	1.23005e-06\\
1.478	1.21123e-06\\
1.48	1.19267e-06\\
1.482	1.17435e-06\\
1.484	1.15628e-06\\
1.486	1.13846e-06\\
1.488	1.12717e-06\\
1.49	1.11836e-06\\
1.492	1.109e-06\\
1.494	1.09913e-06\\
1.496	1.08875e-06\\
1.498	1.07789e-06\\
1.5	1.06657e-06\\
1.502	1.05481e-06\\
1.504	1.04263e-06\\
1.506	1.03006e-06\\
1.508	1.0171e-06\\
1.51	1.0038e-06\\
1.512	9.90159e-07\\
1.514	9.7621e-07\\
1.516	9.61972e-07\\
1.518	9.47468e-07\\
1.52	9.32719e-07\\
1.522	9.17747e-07\\
1.524	9.02573e-07\\
1.526	8.87219e-07\\
1.528	8.71706e-07\\
1.53	8.56055e-07\\
1.532	8.40285e-07\\
1.534	8.24418e-07\\
1.536	8.08472e-07\\
1.538	7.92468e-07\\
1.54	7.76425e-07\\
1.542	7.60359e-07\\
1.544	7.44291e-07\\
1.546	7.28237e-07\\
1.548	7.12214e-07\\
1.55	6.96239e-07\\
1.552	6.80328e-07\\
1.554	6.64496e-07\\
1.556	6.48759e-07\\
1.558	6.35263e-07\\
1.56	6.24652e-07\\
1.562	6.14201e-07\\
1.564	6.03911e-07\\
1.566	5.93779e-07\\
1.568	5.83804e-07\\
1.57	5.73985e-07\\
1.572	5.6432e-07\\
1.574	5.54807e-07\\
1.576	5.45446e-07\\
1.578	5.36234e-07\\
1.58	5.2717e-07\\
1.582	5.18253e-07\\
1.584	5.09481e-07\\
1.586	5.00852e-07\\
1.588	4.92366e-07\\
1.59	4.84019e-07\\
1.592	4.75812e-07\\
1.594	4.67742e-07\\
1.596	4.59807e-07\\
1.598	4.52006e-07\\
1.6	4.44338e-07\\
1.602	4.368e-07\\
1.604	4.29392e-07\\
1.606	4.22111e-07\\
1.608	4.14956e-07\\
1.61	4.07925e-07\\
1.612	4.01017e-07\\
1.614	3.94229e-07\\
1.616	3.87561e-07\\
1.618	3.81011e-07\\
1.62	3.74576e-07\\
1.622	3.68255e-07\\
1.624	3.62048e-07\\
1.626	3.55951e-07\\
1.628	3.49963e-07\\
1.63	3.44083e-07\\
1.632	3.38308e-07\\
1.634	3.32638e-07\\
1.636	3.27071e-07\\
1.638	3.21604e-07\\
1.64	3.16237e-07\\
1.642	3.10967e-07\\
1.644	3.05793e-07\\
1.646	3.00713e-07\\
1.648	2.95727e-07\\
1.65	2.90831e-07\\
1.652	2.86025e-07\\
1.654	2.81307e-07\\
1.656	2.76675e-07\\
1.658	2.72128e-07\\
1.66	2.67665e-07\\
1.662	2.63283e-07\\
1.664	2.58981e-07\\
1.666	2.54758e-07\\
1.668	2.50612e-07\\
1.67	2.46541e-07\\
1.672	2.42545e-07\\
1.674	2.38622e-07\\
1.676	2.3477e-07\\
1.678	2.30989e-07\\
1.68	2.27275e-07\\
1.682	2.23629e-07\\
1.684	2.20049e-07\\
1.686	2.16533e-07\\
1.688	2.13081e-07\\
1.69	2.0969e-07\\
1.692	2.0636e-07\\
1.694	2.03089e-07\\
1.696	1.99877e-07\\
1.698	1.96721e-07\\
1.7	1.93621e-07\\
1.702	1.90575e-07\\
1.704	1.87583e-07\\
1.706	1.84643e-07\\
1.708	1.81754e-07\\
1.71	1.78915e-07\\
1.712	1.76125e-07\\
1.714	1.73383e-07\\
1.716	1.70687e-07\\
1.718	1.68038e-07\\
1.72	1.65433e-07\\
1.722	1.62872e-07\\
1.724	1.60354e-07\\
1.726	1.57878e-07\\
1.728	1.55443e-07\\
1.73	1.53048e-07\\
1.732	1.50693e-07\\
1.734	1.48376e-07\\
1.736	1.46096e-07\\
1.738	1.43853e-07\\
1.74	1.41647e-07\\
1.742	1.39475e-07\\
1.744	1.37338e-07\\
1.746	1.35235e-07\\
1.748	1.33164e-07\\
1.75	1.31127e-07\\
1.752	1.2912e-07\\
1.754	1.27145e-07\\
1.756	1.252e-07\\
1.758	1.23285e-07\\
1.76	1.21399e-07\\
1.762	1.19542e-07\\
1.764	1.17712e-07\\
1.766	1.1591e-07\\
1.768	1.14135e-07\\
1.77	1.12386e-07\\
1.772	1.10663e-07\\
1.774	1.08965e-07\\
1.776	1.07292e-07\\
1.778	1.05644e-07\\
1.78	1.04019e-07\\
1.782	1.02418e-07\\
1.784	1.0084e-07\\
1.786	9.92852e-08\\
1.788	9.77523e-08\\
1.79	9.62413e-08\\
1.792	9.47519e-08\\
1.794	9.32836e-08\\
1.796	9.18361e-08\\
1.798	9.04092e-08\\
1.8	8.90024e-08\\
1.802	8.76155e-08\\
1.804	8.62481e-08\\
1.806	8.48999e-08\\
1.808	8.35707e-08\\
1.81	8.22602e-08\\
1.812	8.09681e-08\\
1.814	7.96942e-08\\
1.816	7.84382e-08\\
1.818	7.71999e-08\\
1.82	7.59789e-08\\
1.822	7.47752e-08\\
1.824	7.35884e-08\\
1.826	7.24184e-08\\
1.828	7.12649e-08\\
1.83	7.01278e-08\\
1.832	6.90068e-08\\
1.834	6.79018e-08\\
1.836	6.68125e-08\\
1.838	6.57387e-08\\
1.84	6.46804e-08\\
1.842	6.36373e-08\\
1.844	6.26092e-08\\
1.846	6.1596e-08\\
1.848	6.05976e-08\\
1.85	5.96136e-08\\
1.852	5.86441e-08\\
1.854	5.76888e-08\\
1.856	5.67475e-08\\
1.858	5.58202e-08\\
1.86	5.49067e-08\\
1.862	5.40068e-08\\
1.864	5.31204e-08\\
1.866	5.22474e-08\\
1.868	5.13875e-08\\
1.87	5.05407e-08\\
1.872	4.97069e-08\\
1.874	4.88858e-08\\
1.876	4.80773e-08\\
1.878	4.72814e-08\\
1.88	4.64978e-08\\
1.882	4.57265e-08\\
1.884	4.49673e-08\\
1.886	4.42201e-08\\
1.888	4.34846e-08\\
1.89	4.27609e-08\\
1.892	4.20488e-08\\
1.894	4.13481e-08\\
1.896	4.06587e-08\\
1.898	3.99804e-08\\
1.9	3.93132e-08\\
1.902	3.86569e-08\\
1.904	3.80114e-08\\
1.906	3.73765e-08\\
1.908	3.67521e-08\\
1.91	3.61381e-08\\
1.912	3.55343e-08\\
1.914	3.49407e-08\\
1.916	3.4357e-08\\
1.918	3.37831e-08\\
1.92	3.3219e-08\\
1.922	3.26645e-08\\
1.924	3.21194e-08\\
1.926	3.15836e-08\\
1.928	3.1057e-08\\
1.93	3.05394e-08\\
1.932	3.00308e-08\\
1.934	2.9531e-08\\
1.936	2.90398e-08\\
1.938	2.85572e-08\\
1.94	2.8083e-08\\
1.942	2.7617e-08\\
1.944	2.71592e-08\\
1.946	2.67094e-08\\
1.948	2.62675e-08\\
1.95	2.58333e-08\\
1.952	2.54068e-08\\
1.954	2.49878e-08\\
1.956	2.45762e-08\\
1.958	2.41718e-08\\
1.96	2.37746e-08\\
1.962	2.33843e-08\\
1.964	2.3001e-08\\
1.966	2.26245e-08\\
1.968	2.22546e-08\\
1.97	2.18912e-08\\
1.972	2.15343e-08\\
1.974	2.11836e-08\\
1.976	2.08392e-08\\
1.978	2.05008e-08\\
1.98	2.01683e-08\\
1.982	1.98418e-08\\
1.984	1.9521e-08\\
1.986	1.92058e-08\\
1.988	1.88961e-08\\
1.99	1.85918e-08\\
1.992	1.82929e-08\\
1.994	1.79992e-08\\
1.996	1.77106e-08\\
1.998	1.7427e-08\\
2	1.71483e-08\\
};
\addplot [color=mycolor10, line width=1.0pt, forget plot]
  table[row sep=crcr]{%
-0.024	0\\
-0.022	0\\
-0.02	0\\
-0.018	0\\
-0.016	0\\
-0.014	0\\
-0.012	0\\
-0.01	0\\
-0.008	0\\
-0.006	0\\
-0.004	0\\
-0.002	0\\
0	0\\
0.002	0.0391008\\
0.004	0.0740917\\
0.006	0.104793\\
0.008	0.131113\\
0.01	0.153044\\
0.012	0.170656\\
0.014	0.184088\\
0.016	0.193544\\
0.018	0.199279\\
0.02	0.201592\\
0.022	0.200817\\
0.024	0.197312\\
0.026	0.191454\\
0.028	0.183622\\
0.03	0.174197\\
0.032	0.16355\\
0.034	0.15204\\
0.036	0.14\\
0.038	0.127742\\
0.04	0.115545\\
0.042	0.114948\\
0.044	0.119348\\
0.046	0.123001\\
0.048	0.125867\\
0.05	0.127918\\
0.052	0.129145\\
0.054	0.129551\\
0.056	0.12915\\
0.058	0.127973\\
0.06	0.126058\\
0.062	0.123456\\
0.064	0.120225\\
0.066	0.116431\\
0.068	0.112143\\
0.07	0.107438\\
0.072	0.102393\\
0.074	0.0970846\\
0.076	0.0915922\\
0.078	0.0859916\\
0.08	0.0816818\\
0.082	0.0831771\\
0.084	0.0845433\\
0.086	0.0857828\\
0.088	0.0868983\\
0.09	0.0878922\\
0.092	0.0887672\\
0.094	0.0895257\\
0.096	0.0901703\\
0.098	0.0907032\\
0.1	0.0911264\\
0.102	0.0914419\\
0.104	0.0916516\\
0.106	0.0917569\\
0.108	0.0917595\\
0.11	0.0916606\\
0.112	0.0914616\\
0.114	0.0911637\\
0.116	0.0921069\\
0.118	0.093187\\
0.12	0.094119\\
0.122	0.0949057\\
0.124	0.0955505\\
0.126	0.096057\\
0.128	0.0964296\\
0.13	0.0966728\\
0.132	0.0967914\\
0.134	0.0967907\\
0.136	0.0966757\\
0.138	0.096452\\
0.14	0.096125\\
0.142	0.0957001\\
0.144	0.0951829\\
0.146	0.0945788\\
0.148	0.093893\\
0.15	0.0931309\\
0.152	0.0922974\\
0.154	0.0913975\\
0.156	0.0904359\\
0.158	0.0894173\\
0.16	0.088346\\
0.162	0.0872262\\
0.164	0.086062\\
0.166	0.084857\\
0.168	0.0836151\\
0.17	0.0823395\\
0.172	0.0810336\\
0.174	0.0797004\\
0.176	0.0783429\\
0.178	0.076964\\
0.18	0.0755661\\
0.182	0.0741519\\
0.184	0.0727237\\
0.186	0.0712838\\
0.188	0.0698344\\
0.19	0.0683775\\
0.192	0.066915\\
0.194	0.065449\\
0.196	0.0639812\\
0.198	0.0625133\\
0.2	0.061047\\
0.202	0.0595839\\
0.204	0.0581255\\
0.206	0.0566733\\
0.208	0.0552286\\
0.21	0.0537929\\
0.212	0.0523674\\
0.214	0.0509533\\
0.216	0.0495518\\
0.218	0.048164\\
0.22	0.046791\\
0.222	0.0454336\\
0.224	0.0440929\\
0.226	0.0427696\\
0.228	0.0414646\\
0.23	0.0401786\\
0.232	0.0389123\\
0.234	0.0376662\\
0.236	0.0364409\\
0.238	0.035237\\
0.24	0.0340547\\
0.242	0.0328946\\
0.244	0.0317569\\
0.246	0.0306418\\
0.248	0.0295496\\
0.25	0.0284804\\
0.252	0.0274345\\
0.254	0.0264117\\
0.256	0.0254123\\
0.258	0.0244361\\
0.26	0.0234831\\
0.262	0.0225534\\
0.264	0.0216467\\
0.266	0.020763\\
0.268	0.0199021\\
0.27	0.0192468\\
0.272	0.0186532\\
0.274	0.0180717\\
0.276	0.0175032\\
0.278	0.0169485\\
0.28	0.0164081\\
0.282	0.016014\\
0.284	0.0158247\\
0.286	0.0156374\\
0.288	0.0154518\\
0.29	0.0152681\\
0.292	0.0150861\\
0.294	0.0149059\\
0.296	0.0147274\\
0.298	0.0145507\\
0.3	0.0143757\\
0.302	0.0142023\\
0.304	0.0140306\\
0.306	0.0138605\\
0.308	0.013692\\
0.31	0.013525\\
0.312	0.0133597\\
0.314	0.0131958\\
0.316	0.0130335\\
0.318	0.0128727\\
0.32	0.0127134\\
0.322	0.0125555\\
0.324	0.0123991\\
0.326	0.0122441\\
0.328	0.0120905\\
0.33	0.0119383\\
0.332	0.0117875\\
0.334	0.0116381\\
0.336	0.01149\\
0.338	0.0113433\\
0.34	0.011198\\
0.342	0.011054\\
0.344	0.0109113\\
0.346	0.01077\\
0.348	0.0106512\\
0.35	0.0107418\\
0.352	0.01082\\
0.354	0.0108856\\
0.356	0.0109386\\
0.358	0.010979\\
0.36	0.011007\\
0.362	0.0110225\\
0.364	0.0110257\\
0.366	0.0110168\\
0.368	0.0109959\\
0.37	0.0109633\\
0.372	0.0109191\\
0.374	0.0108635\\
0.376	0.010797\\
0.378	0.0107196\\
0.38	0.0106319\\
0.382	0.010534\\
0.384	0.0104263\\
0.386	0.0103092\\
0.388	0.0101831\\
0.39	0.0100482\\
0.392	0.00990509\\
0.394	0.00975405\\
0.396	0.0095955\\
0.398	0.00942988\\
0.4	0.00925759\\
0.402	0.00907905\\
0.404	0.00889469\\
0.406	0.00870495\\
0.408	0.00851024\\
0.41	0.008311\\
0.412	0.00810765\\
0.414	0.00790063\\
0.416	0.00769036\\
0.418	0.00747726\\
0.42	0.00726176\\
0.422	0.00704425\\
0.424	0.00682516\\
0.426	0.00660488\\
0.428	0.00638381\\
0.43	0.00616233\\
0.432	0.00594082\\
0.434	0.00584036\\
0.436	0.00575612\\
0.438	0.00567297\\
0.44	0.00559092\\
0.442	0.00550993\\
0.444	0.00543001\\
0.446	0.00535114\\
0.448	0.00527331\\
0.45	0.0051965\\
0.452	0.0051207\\
0.454	0.0050459\\
0.456	0.00497209\\
0.458	0.00489925\\
0.46	0.00482738\\
0.462	0.00475645\\
0.464	0.00468647\\
0.466	0.00461741\\
0.468	0.00454926\\
0.47	0.00448202\\
0.472	0.00441568\\
0.474	0.00435021\\
0.476	0.00428562\\
0.478	0.00422188\\
0.48	0.004159\\
0.482	0.00409695\\
0.484	0.00403574\\
0.486	0.00397534\\
0.488	0.00391576\\
0.49	0.00385698\\
0.492	0.00382462\\
0.494	0.00383675\\
0.496	0.00384597\\
0.498	0.0038523\\
0.5	0.00385577\\
0.502	0.00385639\\
0.504	0.00385421\\
0.506	0.00384926\\
0.508	0.00384157\\
0.51	0.0038312\\
0.512	0.00381819\\
0.514	0.00380258\\
0.516	0.00378445\\
0.518	0.00376383\\
0.52	0.0037408\\
0.522	0.00371541\\
0.524	0.00368773\\
0.526	0.00365783\\
0.528	0.00362578\\
0.53	0.00359164\\
0.532	0.0035555\\
0.534	0.00351743\\
0.536	0.00347751\\
0.538	0.00343581\\
0.54	0.00339241\\
0.542	0.0033474\\
0.544	0.00330086\\
0.546	0.00325286\\
0.548	0.00320349\\
0.55	0.00315283\\
0.552	0.00310097\\
0.554	0.00304798\\
0.556	0.00299396\\
0.558	0.00293898\\
0.56	0.00288313\\
0.562	0.00282649\\
0.564	0.00276915\\
0.566	0.00271117\\
0.568	0.00265264\\
0.57	0.00259365\\
0.572	0.00253427\\
0.574	0.00247457\\
0.576	0.00241464\\
0.578	0.00235454\\
0.58	0.00229435\\
0.582	0.00223414\\
0.584	0.00217399\\
0.586	0.00211395\\
0.588	0.0020541\\
0.59	0.00199451\\
0.592	0.00193522\\
0.594	0.00187631\\
0.596	0.00181783\\
0.598	0.00175985\\
0.6	0.0017024\\
0.602	0.00164556\\
0.604	0.00158937\\
0.606	0.00155492\\
0.608	0.00153025\\
0.61	0.00150595\\
0.612	0.00148203\\
0.614	0.00145848\\
0.616	0.00143529\\
0.618	0.00141246\\
0.62	0.00138998\\
0.622	0.00136785\\
0.624	0.00134606\\
0.626	0.0013246\\
0.628	0.00130348\\
0.63	0.00128268\\
0.632	0.00126221\\
0.634	0.00124205\\
0.636	0.0012222\\
0.638	0.00121033\\
0.64	0.00122185\\
0.642	0.00123182\\
0.644	0.00124024\\
0.646	0.00124713\\
0.648	0.00125249\\
0.65	0.00125634\\
0.652	0.0012587\\
0.654	0.00125959\\
0.656	0.00125903\\
0.658	0.00125706\\
0.66	0.0012537\\
0.662	0.00124898\\
0.664	0.00124294\\
0.666	0.0012356\\
0.668	0.00122702\\
0.67	0.00121722\\
0.672	0.00120625\\
0.674	0.00119415\\
0.676	0.00118096\\
0.678	0.00116673\\
0.68	0.0011515\\
0.682	0.00113531\\
0.684	0.00111823\\
0.686	0.00110029\\
0.688	0.00108154\\
0.69	0.00106203\\
0.692	0.00104181\\
0.694	0.00102093\\
0.696	0.000999436\\
0.698	0.000977384\\
0.7	0.000954819\\
0.702	0.000931791\\
0.704	0.000908349\\
0.706	0.000884542\\
0.708	0.000860416\\
0.71	0.00083602\\
0.712	0.0008114\\
0.714	0.000786602\\
0.716	0.000761672\\
0.718	0.000736652\\
0.72	0.000711588\\
0.722	0.00068652\\
0.724	0.00066149\\
0.726	0.000636537\\
0.728	0.000611702\\
0.73	0.00058702\\
0.732	0.000562529\\
0.734	0.00055207\\
0.736	0.000543238\\
0.738	0.000534555\\
0.74	0.000526017\\
0.742	0.000517622\\
0.744	0.000509368\\
0.746	0.000501253\\
0.748	0.000493274\\
0.75	0.00048543\\
0.752	0.000477717\\
0.754	0.000470134\\
0.756	0.000462677\\
0.758	0.000455346\\
0.76	0.000448138\\
0.762	0.000441051\\
0.764	0.000434082\\
0.766	0.00042723\\
0.768	0.000420493\\
0.77	0.000413867\\
0.772	0.000407352\\
0.774	0.000400946\\
0.776	0.000394646\\
0.778	0.00038845\\
0.78	0.000382357\\
0.782	0.000376365\\
0.784	0.000370471\\
0.786	0.000364674\\
0.788	0.000358973\\
0.79	0.000353365\\
0.792	0.000347849\\
0.794	0.000342423\\
0.796	0.000337086\\
0.798	0.000331835\\
0.8	0.000326669\\
0.802	0.000321587\\
0.804	0.000316587\\
0.806	0.000311667\\
0.808	0.000306827\\
0.81	0.000302064\\
0.812	0.000297377\\
0.814	0.000292765\\
0.816	0.000288227\\
0.818	0.000283761\\
0.82	0.000279365\\
0.822	0.000276172\\
0.824	0.000277183\\
0.826	0.000277983\\
0.828	0.000278573\\
0.83	0.000278953\\
0.832	0.000279127\\
0.834	0.000279095\\
0.836	0.000278861\\
0.838	0.000278427\\
0.84	0.000277796\\
0.842	0.000276972\\
0.844	0.000275958\\
0.846	0.000274757\\
0.848	0.000273375\\
0.85	0.000271815\\
0.852	0.000270082\\
0.854	0.00026818\\
0.856	0.000266115\\
0.858	0.000263892\\
0.86	0.000261515\\
0.862	0.000258991\\
0.864	0.000256325\\
0.866	0.000253523\\
0.868	0.000250589\\
0.87	0.000247531\\
0.872	0.000244353\\
0.874	0.000241063\\
0.876	0.000237665\\
0.878	0.000234166\\
0.88	0.000230573\\
0.882	0.00022689\\
0.884	0.000223125\\
0.886	0.000219284\\
0.888	0.000215372\\
0.89	0.000211395\\
0.892	0.00020736\\
0.894	0.000203273\\
0.896	0.000199139\\
0.898	0.000194964\\
0.9	0.000190755\\
0.902	0.000186516\\
0.904	0.000182255\\
0.906	0.000177975\\
0.908	0.000173683\\
0.91	0.000169383\\
0.912	0.000165082\\
0.914	0.000160784\\
0.916	0.000156494\\
0.918	0.000152217\\
0.92	0.000147958\\
0.922	0.000143721\\
0.924	0.000139511\\
0.926	0.000135332\\
0.928	0.000131188\\
0.93	0.000127084\\
0.932	0.000123022\\
0.934	0.000119007\\
0.936	0.00011963\\
0.938	0.000120267\\
0.94	0.000120756\\
0.942	0.000121099\\
0.944	0.000121298\\
0.946	0.000121356\\
0.948	0.000121275\\
0.95	0.00012106\\
0.952	0.000120712\\
0.954	0.000120235\\
0.956	0.000119634\\
0.958	0.000118911\\
0.96	0.000118071\\
0.962	0.000117117\\
0.964	0.000116055\\
0.966	0.000114888\\
0.968	0.000113621\\
0.97	0.000112258\\
0.972	0.000110803\\
0.974	0.000109263\\
0.976	0.00010764\\
0.978	0.00010594\\
0.98	0.000104168\\
0.982	0.000102329\\
0.984	0.000100426\\
0.986	9.8466e-05\\
0.988	9.64524e-05\\
0.99	9.43901e-05\\
0.992	9.2284e-05\\
0.994	9.01386e-05\\
0.996	8.79585e-05\\
0.998	8.57482e-05\\
1	8.35123e-05\\
1.002	8.12549e-05\\
1.004	7.89805e-05\\
1.006	7.66932e-05\\
1.008	7.4397e-05\\
1.01	7.20961e-05\\
1.012	6.97941e-05\\
1.014	6.7495e-05\\
1.016	6.52024e-05\\
1.018	6.29197e-05\\
1.02	6.06505e-05\\
1.022	5.83978e-05\\
1.024	5.6165e-05\\
1.026	5.49511e-05\\
1.028	5.40837e-05\\
1.03	5.32306e-05\\
1.032	5.23918e-05\\
1.034	5.15668e-05\\
1.036	5.07556e-05\\
1.038	4.99579e-05\\
1.04	4.91735e-05\\
1.042	4.84022e-05\\
1.044	4.76438e-05\\
1.046	4.68981e-05\\
1.048	4.61649e-05\\
1.05	4.54439e-05\\
1.052	4.47351e-05\\
1.054	4.40382e-05\\
1.056	4.3353e-05\\
1.058	4.26793e-05\\
1.06	4.20169e-05\\
1.062	4.13656e-05\\
1.064	4.07253e-05\\
1.066	4.00957e-05\\
1.068	3.94766e-05\\
1.07	3.88679e-05\\
1.072	3.82694e-05\\
1.074	3.76809e-05\\
1.076	3.71023e-05\\
1.078	3.65332e-05\\
1.08	3.59736e-05\\
1.082	3.54234e-05\\
1.084	3.48822e-05\\
1.086	3.435e-05\\
1.088	3.38265e-05\\
1.09	3.33117e-05\\
1.092	3.28053e-05\\
1.094	3.23072e-05\\
1.096	3.18172e-05\\
1.098	3.13351e-05\\
1.1	3.08609e-05\\
1.102	3.03943e-05\\
1.104	2.99353e-05\\
1.106	2.94836e-05\\
1.108	2.90391e-05\\
1.11	2.86016e-05\\
1.112	2.81712e-05\\
1.114	2.77474e-05\\
1.116	2.73304e-05\\
1.118	2.69199e-05\\
1.12	2.65157e-05\\
1.122	2.61179e-05\\
1.124	2.57262e-05\\
1.126	2.53405e-05\\
1.128	2.49607e-05\\
1.13	2.45867e-05\\
1.132	2.42183e-05\\
1.134	2.38556e-05\\
1.136	2.34982e-05\\
1.138	2.31463e-05\\
1.14	2.27995e-05\\
1.142	2.2458e-05\\
1.144	2.21214e-05\\
1.146	2.17899e-05\\
1.148	2.14631e-05\\
1.15	2.11412e-05\\
1.152	2.08239e-05\\
1.154	2.05112e-05\\
1.156	2.03416e-05\\
1.158	2.02872e-05\\
1.16	2.02199e-05\\
1.162	2.014e-05\\
1.164	2.00476e-05\\
1.166	1.99432e-05\\
1.168	1.98269e-05\\
1.17	1.9699e-05\\
1.172	1.956e-05\\
1.174	1.941e-05\\
1.176	1.92495e-05\\
1.178	1.90789e-05\\
1.18	1.88985e-05\\
1.182	1.87086e-05\\
1.184	1.85097e-05\\
1.186	1.83022e-05\\
1.188	1.80865e-05\\
1.19	1.78629e-05\\
1.192	1.76319e-05\\
1.194	1.7394e-05\\
1.196	1.71494e-05\\
1.198	1.68987e-05\\
1.2	1.66423e-05\\
1.202	1.63806e-05\\
1.204	1.6114e-05\\
1.206	1.58429e-05\\
1.208	1.55678e-05\\
1.21	1.5289e-05\\
1.212	1.5007e-05\\
1.214	1.47222e-05\\
1.216	1.44349e-05\\
1.218	1.41455e-05\\
1.22	1.38546e-05\\
1.222	1.35623e-05\\
1.224	1.32691e-05\\
1.226	1.29753e-05\\
1.228	1.26814e-05\\
1.23	1.23876e-05\\
1.232	1.20943e-05\\
1.234	1.18017e-05\\
1.236	1.15104e-05\\
1.238	1.12204e-05\\
1.24	1.09322e-05\\
1.242	1.06459e-05\\
1.244	1.0362e-05\\
1.246	1.02221e-05\\
1.248	1.01586e-05\\
1.25	1.00864e-05\\
1.252	1.00059e-05\\
1.254	9.9173e-06\\
1.256	9.82101e-06\\
1.258	9.71732e-06\\
1.26	9.60657e-06\\
1.262	9.4891e-06\\
1.264	9.36525e-06\\
1.266	9.23537e-06\\
1.268	9.09983e-06\\
1.27	8.95896e-06\\
1.272	8.81314e-06\\
1.274	8.6627e-06\\
1.276	8.50802e-06\\
1.278	8.34944e-06\\
1.28	8.18731e-06\\
1.282	8.02199e-06\\
1.284	7.85382e-06\\
1.286	7.68313e-06\\
1.288	7.51028e-06\\
1.29	7.33558e-06\\
1.292	7.15937e-06\\
1.294	6.98195e-06\\
1.296	6.80365e-06\\
1.298	6.62477e-06\\
1.3	6.46003e-06\\
1.302	6.35529e-06\\
1.304	6.25227e-06\\
1.306	6.15096e-06\\
1.308	6.05133e-06\\
1.31	5.95336e-06\\
1.312	5.85703e-06\\
1.314	5.76232e-06\\
1.316	5.6692e-06\\
1.318	5.57766e-06\\
1.32	5.48767e-06\\
1.322	5.39921e-06\\
1.324	5.31225e-06\\
1.326	5.22679e-06\\
1.328	5.14278e-06\\
1.33	5.06022e-06\\
1.332	4.97908e-06\\
1.334	4.89934e-06\\
1.336	4.82098e-06\\
1.338	4.74398e-06\\
1.34	4.66831e-06\\
1.342	4.59396e-06\\
1.344	4.5209e-06\\
1.346	4.44911e-06\\
1.348	4.37858e-06\\
1.35	4.30928e-06\\
1.352	4.24118e-06\\
1.354	4.17428e-06\\
1.356	4.10855e-06\\
1.358	4.04396e-06\\
1.36	3.98051e-06\\
1.362	3.91816e-06\\
1.364	3.8569e-06\\
1.366	3.79671e-06\\
1.368	3.73758e-06\\
1.37	3.67947e-06\\
1.372	3.62237e-06\\
1.374	3.56627e-06\\
1.376	3.51113e-06\\
1.378	3.45696e-06\\
1.38	3.40371e-06\\
1.382	3.35139e-06\\
1.384	3.29997e-06\\
1.386	3.24942e-06\\
1.388	3.19974e-06\\
1.39	3.15091e-06\\
1.392	3.10291e-06\\
1.394	3.05572e-06\\
1.396	3.00933e-06\\
1.398	2.96371e-06\\
1.4	2.91886e-06\\
1.402	2.87475e-06\\
1.404	2.83138e-06\\
1.406	2.78872e-06\\
1.408	2.74677e-06\\
1.41	2.7055e-06\\
1.412	2.6649e-06\\
1.414	2.62496e-06\\
1.416	2.58566e-06\\
1.418	2.54699e-06\\
1.42	2.50894e-06\\
1.422	2.47149e-06\\
1.424	2.43463e-06\\
1.426	2.39835e-06\\
1.428	2.36263e-06\\
1.43	2.32746e-06\\
1.432	2.29284e-06\\
1.434	2.25875e-06\\
1.436	2.22517e-06\\
1.438	2.1921e-06\\
1.44	2.15953e-06\\
1.442	2.12745e-06\\
1.444	2.09584e-06\\
1.446	2.0647e-06\\
1.448	2.03402e-06\\
1.45	2.00379e-06\\
1.452	1.97399e-06\\
1.454	1.94463e-06\\
1.456	1.91569e-06\\
1.458	1.88716e-06\\
1.46	1.85904e-06\\
1.462	1.83132e-06\\
1.464	1.80398e-06\\
1.466	1.77704e-06\\
1.468	1.75047e-06\\
1.47	1.72427e-06\\
1.472	1.69844e-06\\
1.474	1.67296e-06\\
1.476	1.64783e-06\\
1.478	1.62305e-06\\
1.48	1.59862e-06\\
1.482	1.57451e-06\\
1.484	1.55074e-06\\
1.486	1.52729e-06\\
1.488	1.50416e-06\\
1.49	1.48134e-06\\
1.492	1.45884e-06\\
1.494	1.43664e-06\\
1.496	1.41474e-06\\
1.498	1.39314e-06\\
1.5	1.37183e-06\\
1.502	1.35081e-06\\
1.504	1.33008e-06\\
1.506	1.30963e-06\\
1.508	1.28946e-06\\
1.51	1.26957e-06\\
1.512	1.24994e-06\\
1.514	1.23059e-06\\
1.516	1.2115e-06\\
1.518	1.19267e-06\\
1.52	1.17411e-06\\
1.522	1.1558e-06\\
1.524	1.13774e-06\\
1.526	1.11994e-06\\
1.528	1.10238e-06\\
1.53	1.08508e-06\\
1.532	1.06801e-06\\
1.534	1.05119e-06\\
1.536	1.0346e-06\\
1.538	1.01825e-06\\
1.54	1.00213e-06\\
1.542	9.8625e-07\\
1.544	9.70594e-07\\
1.546	9.55165e-07\\
1.548	9.39961e-07\\
1.55	9.24978e-07\\
1.552	9.10216e-07\\
1.554	8.95671e-07\\
1.556	8.81342e-07\\
1.558	8.67227e-07\\
1.56	8.53322e-07\\
1.562	8.39626e-07\\
1.564	8.26137e-07\\
1.566	8.12853e-07\\
1.568	7.99771e-07\\
1.57	7.8689e-07\\
1.572	7.74207e-07\\
1.574	7.61719e-07\\
1.576	7.49426e-07\\
1.578	7.37324e-07\\
1.58	7.25412e-07\\
1.582	7.13687e-07\\
1.584	7.02148e-07\\
1.586	6.90791e-07\\
1.588	6.79616e-07\\
1.59	6.68619e-07\\
1.592	6.57799e-07\\
1.594	6.47153e-07\\
1.596	6.3668e-07\\
1.598	6.26377e-07\\
1.6	6.16242e-07\\
1.602	6.06272e-07\\
1.604	5.96467e-07\\
1.606	5.86822e-07\\
1.608	5.77338e-07\\
1.61	5.6801e-07\\
1.612	5.58837e-07\\
1.614	5.49818e-07\\
1.616	5.40949e-07\\
1.618	5.32229e-07\\
1.62	5.23655e-07\\
1.622	5.15225e-07\\
1.624	5.06938e-07\\
1.626	4.98791e-07\\
1.628	4.90782e-07\\
1.63	4.82909e-07\\
1.632	4.7517e-07\\
1.634	4.67562e-07\\
1.636	4.60084e-07\\
1.638	4.52734e-07\\
1.64	4.45509e-07\\
1.642	4.38408e-07\\
1.644	4.31429e-07\\
1.646	4.24569e-07\\
1.648	4.17826e-07\\
1.65	4.11199e-07\\
1.652	4.04686e-07\\
1.654	3.98285e-07\\
1.656	3.91993e-07\\
1.658	3.85809e-07\\
1.66	3.79731e-07\\
1.662	3.73757e-07\\
1.664	3.67886e-07\\
1.666	3.62115e-07\\
1.668	3.56443e-07\\
1.67	3.50868e-07\\
1.672	3.45387e-07\\
1.674	3.4e-07\\
1.676	3.34705e-07\\
1.678	3.295e-07\\
1.68	3.24383e-07\\
1.682	3.19353e-07\\
1.684	3.14408e-07\\
1.686	3.09546e-07\\
1.688	3.04766e-07\\
1.69	3.00066e-07\\
1.692	2.95445e-07\\
1.694	2.90902e-07\\
1.696	2.86434e-07\\
1.698	2.8204e-07\\
1.7	2.77719e-07\\
1.702	2.7347e-07\\
1.704	2.6929e-07\\
1.706	2.65179e-07\\
1.708	2.61136e-07\\
1.71	2.57159e-07\\
1.712	2.53246e-07\\
1.714	2.49397e-07\\
1.716	2.4561e-07\\
1.718	2.41884e-07\\
1.72	2.38217e-07\\
1.722	2.3461e-07\\
1.724	2.31059e-07\\
1.726	2.27565e-07\\
1.728	2.24127e-07\\
1.73	2.20742e-07\\
1.732	2.17411e-07\\
1.734	2.14131e-07\\
1.736	2.10903e-07\\
1.738	2.07724e-07\\
1.74	2.04595e-07\\
1.742	2.01514e-07\\
1.744	1.9848e-07\\
1.746	1.95492e-07\\
1.748	1.9255e-07\\
1.75	1.89653e-07\\
1.752	1.86799e-07\\
1.754	1.83988e-07\\
1.756	1.81219e-07\\
1.758	1.78491e-07\\
1.76	1.75804e-07\\
1.762	1.73157e-07\\
1.764	1.70549e-07\\
1.766	1.67979e-07\\
1.768	1.65447e-07\\
1.77	1.62952e-07\\
1.772	1.60494e-07\\
1.774	1.58071e-07\\
1.776	1.55683e-07\\
1.778	1.5333e-07\\
1.78	1.51011e-07\\
1.782	1.48726e-07\\
1.784	1.46473e-07\\
1.786	1.44252e-07\\
1.788	1.42063e-07\\
1.79	1.39906e-07\\
1.792	1.37779e-07\\
1.794	1.35683e-07\\
1.796	1.33616e-07\\
1.798	1.31579e-07\\
1.8	1.29571e-07\\
1.802	1.27591e-07\\
1.804	1.2564e-07\\
1.806	1.23716e-07\\
1.808	1.21819e-07\\
1.81	1.19949e-07\\
1.812	1.18106e-07\\
1.814	1.16289e-07\\
1.816	1.14498e-07\\
1.818	1.12732e-07\\
1.82	1.10991e-07\\
1.822	1.09275e-07\\
1.824	1.07583e-07\\
1.826	1.05916e-07\\
1.828	1.04272e-07\\
1.83	1.02652e-07\\
1.832	1.01055e-07\\
1.834	9.94805e-08\\
1.836	9.7929e-08\\
1.838	9.63999e-08\\
1.84	9.48928e-08\\
1.842	9.34076e-08\\
1.844	9.19439e-08\\
1.846	9.05015e-08\\
1.848	8.90802e-08\\
1.85	8.76796e-08\\
1.852	8.62996e-08\\
1.854	8.49398e-08\\
1.856	8.36001e-08\\
1.858	8.22803e-08\\
1.86	8.098e-08\\
1.862	7.96991e-08\\
1.864	7.84373e-08\\
1.866	7.71944e-08\\
1.868	7.59702e-08\\
1.87	7.47645e-08\\
1.872	7.3577e-08\\
1.874	7.24076e-08\\
1.876	7.12559e-08\\
1.878	7.01219e-08\\
1.88	6.90053e-08\\
1.882	6.79058e-08\\
1.884	6.68233e-08\\
1.886	6.57576e-08\\
1.888	6.47085e-08\\
1.89	6.36757e-08\\
1.892	6.26591e-08\\
1.894	6.16585e-08\\
1.896	6.06736e-08\\
1.898	5.97042e-08\\
1.9	5.87502e-08\\
1.902	5.78114e-08\\
1.904	5.68875e-08\\
1.906	5.59784e-08\\
1.908	5.50839e-08\\
1.91	5.42037e-08\\
1.912	5.33377e-08\\
1.914	5.24857e-08\\
1.916	5.16475e-08\\
1.918	5.08229e-08\\
1.92	5.00118e-08\\
1.922	4.92138e-08\\
1.924	4.84289e-08\\
1.926	4.76568e-08\\
1.928	4.68974e-08\\
1.93	4.61505e-08\\
1.932	4.54158e-08\\
1.934	4.46933e-08\\
1.936	4.39827e-08\\
1.938	4.32839e-08\\
1.94	4.25966e-08\\
1.942	4.19207e-08\\
1.944	4.1256e-08\\
1.946	4.06023e-08\\
1.948	3.99595e-08\\
1.95	3.93274e-08\\
1.952	3.87058e-08\\
1.954	3.80945e-08\\
1.956	3.74934e-08\\
1.958	3.69023e-08\\
1.96	3.63211e-08\\
1.962	3.57495e-08\\
1.964	3.51874e-08\\
1.966	3.46347e-08\\
1.968	3.40912e-08\\
1.97	3.35568e-08\\
1.972	3.30312e-08\\
1.974	3.25143e-08\\
1.976	3.2006e-08\\
1.978	3.15061e-08\\
1.98	3.10145e-08\\
1.982	3.0531e-08\\
1.984	3.00556e-08\\
1.986	2.95879e-08\\
1.988	2.9128e-08\\
1.99	2.86756e-08\\
1.992	2.82307e-08\\
1.994	2.7793e-08\\
1.996	2.73625e-08\\
1.998	2.6939e-08\\
2	2.65225e-08\\
};
\addplot [color=mycolor11, line width=1.0pt, forget plot]
  table[row sep=crcr]{%
-0.024	0\\
-0.022	0\\
-0.02	0\\
-0.018	0\\
-0.016	0\\
-0.014	0\\
-0.012	0\\
-0.01	0\\
-0.008	0\\
-0.006	0\\
-0.004	0\\
-0.002	0\\
0	0\\
0.002	0.038829\\
0.004	0.073082\\
0.006	0.102694\\
0.008	0.127685\\
0.01	0.148152\\
0.012	0.164257\\
0.014	0.176228\\
0.016	0.184337\\
0.018	0.1889\\
0.02	0.190263\\
0.022	0.188793\\
0.024	0.184868\\
0.026	0.178873\\
0.028	0.171185\\
0.03	0.162173\\
0.032	0.152188\\
0.034	0.141558\\
0.036	0.130587\\
0.038	0.119547\\
0.04	0.108679\\
0.042	0.110169\\
0.044	0.114146\\
0.046	0.117418\\
0.048	0.119954\\
0.05	0.121735\\
0.052	0.122757\\
0.054	0.123027\\
0.056	0.122566\\
0.058	0.121402\\
0.06	0.119576\\
0.062	0.117136\\
0.064	0.114137\\
0.066	0.110639\\
0.068	0.106707\\
0.07	0.10241\\
0.072	0.0978155\\
0.074	0.0929946\\
0.076	0.0880159\\
0.078	0.0829462\\
0.08	0.0787001\\
0.082	0.0801571\\
0.084	0.0814954\\
0.086	0.0827164\\
0.088	0.0838217\\
0.09	0.0848129\\
0.092	0.0856915\\
0.094	0.0864591\\
0.096	0.0871172\\
0.098	0.0876671\\
0.1	0.0881104\\
0.102	0.0884483\\
0.104	0.0886819\\
0.106	0.0888125\\
0.108	0.0888413\\
0.11	0.0887693\\
0.112	0.0885979\\
0.114	0.0883282\\
0.116	0.0889122\\
0.118	0.0899833\\
0.12	0.0909137\\
0.122	0.0917056\\
0.124	0.0923619\\
0.126	0.0928859\\
0.128	0.0932815\\
0.13	0.0935526\\
0.132	0.0937037\\
0.134	0.0937395\\
0.136	0.0936647\\
0.138	0.0934845\\
0.14	0.0932038\\
0.142	0.0928278\\
0.144	0.0923617\\
0.146	0.0918106\\
0.148	0.0911796\\
0.15	0.0904736\\
0.152	0.0896975\\
0.154	0.0888562\\
0.156	0.0879543\\
0.158	0.0869961\\
0.16	0.0859861\\
0.162	0.0849283\\
0.164	0.0838267\\
0.166	0.0826852\\
0.168	0.0815072\\
0.17	0.0802963\\
0.172	0.0790556\\
0.174	0.0777884\\
0.176	0.0764974\\
0.178	0.0751855\\
0.18	0.0738554\\
0.182	0.0725094\\
0.184	0.07115\\
0.186	0.0697793\\
0.188	0.0683996\\
0.19	0.0670128\\
0.192	0.0656207\\
0.194	0.0642253\\
0.196	0.0628281\\
0.198	0.0614309\\
0.2	0.0600352\\
0.202	0.0586425\\
0.204	0.0572541\\
0.206	0.0558714\\
0.208	0.0544957\\
0.21	0.0531282\\
0.212	0.0517699\\
0.214	0.050422\\
0.216	0.0490856\\
0.218	0.0477616\\
0.22	0.0464509\\
0.222	0.0451543\\
0.224	0.0438728\\
0.226	0.042607\\
0.228	0.0413576\\
0.23	0.0401254\\
0.232	0.0389109\\
0.234	0.0377147\\
0.236	0.0365373\\
0.238	0.0353791\\
0.24	0.0342406\\
0.242	0.0331221\\
0.244	0.032024\\
0.246	0.0309466\\
0.248	0.0298901\\
0.25	0.0288547\\
0.252	0.0278406\\
0.254	0.0268479\\
0.256	0.0258768\\
0.258	0.0249273\\
0.26	0.0239994\\
0.262	0.0230932\\
0.264	0.0222086\\
0.266	0.0213457\\
0.268	0.0205043\\
0.27	0.0196845\\
0.272	0.0190287\\
0.274	0.0184951\\
0.276	0.0179728\\
0.278	0.0174621\\
0.28	0.0169636\\
0.282	0.0164775\\
0.284	0.0160042\\
0.286	0.0156262\\
0.288	0.0154518\\
0.29	0.0152791\\
0.292	0.015108\\
0.294	0.0149386\\
0.296	0.0147708\\
0.298	0.0146046\\
0.3	0.0144399\\
0.302	0.0142768\\
0.304	0.0141153\\
0.306	0.0139552\\
0.308	0.0137965\\
0.31	0.0136394\\
0.312	0.0134836\\
0.314	0.0133292\\
0.316	0.0131762\\
0.318	0.0130246\\
0.32	0.0128743\\
0.322	0.0127253\\
0.324	0.0125777\\
0.326	0.0124313\\
0.328	0.0122861\\
0.33	0.0121423\\
0.332	0.0119996\\
0.334	0.0118582\\
0.336	0.0117181\\
0.338	0.0115791\\
0.34	0.0114413\\
0.342	0.0113048\\
0.344	0.0111694\\
0.346	0.0110352\\
0.348	0.0110636\\
0.35	0.0111496\\
0.352	0.0112238\\
0.354	0.0112862\\
0.356	0.0113367\\
0.358	0.0113754\\
0.36	0.0114022\\
0.362	0.0114172\\
0.364	0.0114206\\
0.366	0.0114123\\
0.368	0.0113927\\
0.37	0.0113617\\
0.372	0.0113197\\
0.374	0.0112668\\
0.376	0.0112033\\
0.378	0.0111294\\
0.38	0.0110455\\
0.382	0.0109518\\
0.384	0.0108485\\
0.386	0.0107361\\
0.388	0.010615\\
0.39	0.0104853\\
0.392	0.0103475\\
0.394	0.010202\\
0.396	0.0100492\\
0.398	0.00988935\\
0.4	0.00972297\\
0.402	0.00955041\\
0.404	0.0093721\\
0.406	0.00918842\\
0.408	0.0089998\\
0.41	0.00880664\\
0.412	0.00860935\\
0.414	0.00840834\\
0.416	0.00820403\\
0.418	0.00799681\\
0.42	0.00778709\\
0.422	0.00757526\\
0.424	0.00736172\\
0.426	0.00714685\\
0.428	0.00693104\\
0.43	0.00671465\\
0.432	0.00649805\\
0.434	0.0062816\\
0.436	0.0061878\\
0.438	0.00610521\\
0.44	0.00602361\\
0.442	0.00594296\\
0.444	0.00586328\\
0.446	0.00578453\\
0.448	0.00570673\\
0.45	0.00562985\\
0.452	0.00555388\\
0.454	0.00547881\\
0.456	0.00540464\\
0.458	0.00533136\\
0.46	0.00525895\\
0.462	0.0051874\\
0.464	0.00511671\\
0.466	0.00504687\\
0.468	0.00497786\\
0.47	0.00490968\\
0.472	0.00484232\\
0.474	0.00477576\\
0.476	0.00471001\\
0.478	0.00464506\\
0.48	0.00458088\\
0.482	0.00451748\\
0.484	0.00445486\\
0.486	0.00439299\\
0.488	0.00435528\\
0.49	0.00436659\\
0.492	0.00437507\\
0.494	0.00438073\\
0.496	0.00438358\\
0.498	0.00438365\\
0.5	0.00438095\\
0.502	0.00437552\\
0.504	0.00436739\\
0.506	0.00435659\\
0.508	0.00434316\\
0.51	0.00432716\\
0.512	0.00430862\\
0.514	0.00428759\\
0.516	0.00426413\\
0.518	0.00423829\\
0.52	0.00421014\\
0.522	0.00417974\\
0.524	0.00414714\\
0.526	0.00411241\\
0.528	0.00407563\\
0.53	0.00403687\\
0.532	0.00399618\\
0.534	0.00395366\\
0.536	0.00390937\\
0.538	0.00386339\\
0.54	0.00381579\\
0.542	0.00376666\\
0.544	0.00371607\\
0.546	0.00366411\\
0.548	0.00361084\\
0.55	0.00355636\\
0.552	0.00350074\\
0.554	0.00344407\\
0.556	0.00338642\\
0.558	0.00332788\\
0.56	0.00326852\\
0.562	0.00320842\\
0.564	0.00314767\\
0.566	0.00308634\\
0.568	0.0030245\\
0.57	0.00296225\\
0.572	0.00289964\\
0.574	0.00283676\\
0.576	0.00277368\\
0.578	0.00271047\\
0.58	0.00264721\\
0.582	0.00258395\\
0.584	0.00252078\\
0.586	0.00245774\\
0.588	0.00239492\\
0.59	0.00233237\\
0.592	0.00227015\\
0.594	0.00220833\\
0.596	0.00214695\\
0.598	0.00208608\\
0.6	0.00202576\\
0.602	0.00196606\\
0.604	0.00190701\\
0.606	0.00184866\\
0.608	0.00180757\\
0.61	0.00178085\\
0.612	0.00175451\\
0.614	0.00172854\\
0.616	0.00170295\\
0.618	0.00167773\\
0.62	0.00165287\\
0.622	0.00162837\\
0.624	0.00160422\\
0.626	0.00158042\\
0.628	0.00155696\\
0.63	0.00153384\\
0.632	0.00151105\\
0.634	0.00148859\\
0.636	0.00146646\\
0.638	0.00144465\\
0.64	0.00142315\\
0.642	0.00140893\\
0.644	0.00141445\\
0.646	0.00141854\\
0.648	0.00142121\\
0.65	0.00142247\\
0.652	0.00142234\\
0.654	0.00142084\\
0.656	0.00141798\\
0.658	0.0014138\\
0.66	0.00140832\\
0.662	0.00140158\\
0.664	0.00139359\\
0.666	0.0013844\\
0.668	0.00137403\\
0.67	0.00136253\\
0.672	0.00134993\\
0.674	0.00133628\\
0.676	0.0013216\\
0.678	0.00130595\\
0.68	0.00128937\\
0.682	0.0012719\\
0.684	0.00125358\\
0.686	0.00123446\\
0.688	0.00121459\\
0.69	0.00119401\\
0.692	0.00117277\\
0.694	0.00115091\\
0.696	0.00112849\\
0.698	0.00110554\\
0.7	0.00108212\\
0.702	0.00105828\\
0.704	0.00103405\\
0.706	0.00100948\\
0.708	0.000984629\\
0.71	0.00095953\\
0.712	0.00093423\\
0.714	0.000908773\\
0.716	0.000883202\\
0.718	0.000857559\\
0.72	0.000831885\\
0.722	0.000806219\\
0.724	0.000780602\\
0.726	0.000755071\\
0.728	0.00073345\\
0.73	0.000722523\\
0.732	0.000711765\\
0.734	0.000701174\\
0.736	0.000690747\\
0.738	0.000680482\\
0.74	0.000670376\\
0.742	0.000660427\\
0.744	0.000650633\\
0.746	0.000640991\\
0.748	0.000631499\\
0.75	0.000622154\\
0.752	0.000612954\\
0.754	0.000603898\\
0.756	0.000594982\\
0.758	0.000586205\\
0.76	0.000577564\\
0.762	0.000569057\\
0.764	0.000560682\\
0.766	0.000552437\\
0.768	0.000544319\\
0.77	0.000536327\\
0.772	0.000528459\\
0.774	0.000520712\\
0.776	0.000513084\\
0.778	0.000505574\\
0.78	0.000498179\\
0.782	0.000490898\\
0.784	0.000483728\\
0.786	0.000476668\\
0.788	0.000469717\\
0.79	0.000462871\\
0.792	0.000456129\\
0.794	0.00044949\\
0.796	0.000442952\\
0.798	0.000436513\\
0.8	0.000430171\\
0.802	0.000423925\\
0.804	0.000417774\\
0.806	0.000411714\\
0.808	0.000405746\\
0.81	0.000399867\\
0.812	0.000394076\\
0.814	0.000388372\\
0.816	0.000382752\\
0.818	0.000377216\\
0.82	0.000371763\\
0.822	0.00036639\\
0.824	0.000361096\\
0.826	0.000355881\\
0.828	0.000355459\\
0.83	0.000354962\\
0.832	0.000354276\\
0.834	0.000353401\\
0.836	0.000352339\\
0.838	0.000351093\\
0.84	0.000349665\\
0.842	0.000348059\\
0.844	0.000346277\\
0.846	0.000344323\\
0.848	0.0003422\\
0.85	0.000339914\\
0.852	0.000337467\\
0.854	0.000334864\\
0.856	0.00033211\\
0.858	0.000329209\\
0.86	0.000326166\\
0.862	0.000322987\\
0.864	0.000319676\\
0.866	0.000316239\\
0.868	0.000312681\\
0.87	0.000309008\\
0.872	0.000305225\\
0.874	0.000301337\\
0.876	0.000297352\\
0.878	0.000293273\\
0.88	0.000289108\\
0.882	0.000284861\\
0.884	0.000280539\\
0.886	0.000276148\\
0.888	0.000271693\\
0.89	0.00026718\\
0.892	0.000262615\\
0.894	0.000258003\\
0.896	0.000253351\\
0.898	0.000248663\\
0.9	0.000243946\\
0.902	0.000239204\\
0.904	0.000234444\\
0.906	0.000229671\\
0.908	0.000224889\\
0.91	0.000220104\\
0.912	0.000215321\\
0.914	0.000210544\\
0.916	0.00020578\\
0.918	0.000201031\\
0.92	0.000196304\\
0.922	0.000191601\\
0.924	0.000186928\\
0.926	0.000182288\\
0.928	0.000177686\\
0.93	0.000173126\\
0.932	0.00016861\\
0.934	0.000164143\\
0.936	0.000159728\\
0.938	0.000157354\\
0.94	0.000155054\\
0.942	0.000153223\\
0.944	0.000152982\\
0.946	0.00015261\\
0.948	0.00015211\\
0.95	0.000151485\\
0.952	0.000150739\\
0.954	0.000149873\\
0.956	0.000148892\\
0.958	0.000147799\\
0.96	0.000146598\\
0.962	0.000145291\\
0.964	0.000143885\\
0.966	0.000142381\\
0.968	0.000140785\\
0.97	0.000139101\\
0.972	0.000137333\\
0.974	0.000135485\\
0.976	0.000133562\\
0.978	0.000131568\\
0.98	0.000129507\\
0.982	0.000127386\\
0.984	0.000125206\\
0.986	0.000122975\\
0.988	0.000120695\\
0.99	0.000118371\\
0.992	0.000116008\\
0.994	0.00011361\\
0.996	0.000111181\\
0.998	0.000108727\\
1	0.00010625\\
1.002	0.000103756\\
1.004	0.000101247\\
1.006	9.87297e-05\\
1.008	9.62061e-05\\
1.01	9.36805e-05\\
1.012	9.11567e-05\\
1.014	8.95883e-05\\
1.016	8.82678e-05\\
1.018	8.69673e-05\\
1.02	8.56863e-05\\
1.022	8.44247e-05\\
1.024	8.31822e-05\\
1.026	8.19586e-05\\
1.028	8.07535e-05\\
1.03	7.95668e-05\\
1.032	7.83982e-05\\
1.034	7.72475e-05\\
1.036	7.61143e-05\\
1.038	7.49985e-05\\
1.04	7.38998e-05\\
1.042	7.28179e-05\\
1.044	7.17527e-05\\
1.046	7.07038e-05\\
1.048	6.96711e-05\\
1.05	6.86543e-05\\
1.052	6.76532e-05\\
1.054	6.66675e-05\\
1.056	6.5697e-05\\
1.058	6.47415e-05\\
1.06	6.38007e-05\\
1.062	6.28745e-05\\
1.064	6.19625e-05\\
1.066	6.10646e-05\\
1.068	6.01805e-05\\
1.07	5.93101e-05\\
1.072	5.84531e-05\\
1.074	5.76092e-05\\
1.076	5.67784e-05\\
1.078	5.59603e-05\\
1.08	5.51547e-05\\
1.082	5.43615e-05\\
1.084	5.35804e-05\\
1.086	5.28112e-05\\
1.088	5.20538e-05\\
1.09	5.13079e-05\\
1.092	5.05733e-05\\
1.094	4.98499e-05\\
1.096	4.91374e-05\\
1.098	4.84356e-05\\
1.1	4.77444e-05\\
1.102	4.70636e-05\\
1.104	4.63931e-05\\
1.106	4.57325e-05\\
1.108	4.50818e-05\\
1.11	4.44407e-05\\
1.112	4.38092e-05\\
1.114	4.31869e-05\\
1.116	4.25739e-05\\
1.118	4.19698e-05\\
1.12	4.13747e-05\\
1.122	4.07881e-05\\
1.124	4.02102e-05\\
1.126	3.96406e-05\\
1.128	3.90792e-05\\
1.13	3.8526e-05\\
1.132	3.79807e-05\\
1.134	3.74432e-05\\
1.136	3.69134e-05\\
1.138	3.63911e-05\\
1.14	3.58762e-05\\
1.142	3.53686e-05\\
1.144	3.48682e-05\\
1.146	3.43748e-05\\
1.148	3.38884e-05\\
1.15	3.34087e-05\\
1.152	3.29357e-05\\
1.154	3.24693e-05\\
1.156	3.20093e-05\\
1.158	3.15558e-05\\
1.16	3.11084e-05\\
1.162	3.06672e-05\\
1.164	3.02321e-05\\
1.166	2.9803e-05\\
1.168	2.93797e-05\\
1.17	2.89621e-05\\
1.172	2.85503e-05\\
1.174	2.8144e-05\\
1.176	2.7767e-05\\
1.178	2.74989e-05\\
1.18	2.7223e-05\\
1.182	2.69396e-05\\
1.184	2.66489e-05\\
1.186	2.63513e-05\\
1.188	2.60472e-05\\
1.19	2.57368e-05\\
1.192	2.54206e-05\\
1.194	2.50989e-05\\
1.196	2.47721e-05\\
1.198	2.44404e-05\\
1.2	2.41044e-05\\
1.202	2.37643e-05\\
1.204	2.34204e-05\\
1.206	2.30733e-05\\
1.208	2.27232e-05\\
1.21	2.23705e-05\\
1.212	2.20155e-05\\
1.214	2.16587e-05\\
1.216	2.13003e-05\\
1.218	2.09407e-05\\
1.22	2.05803e-05\\
1.222	2.02194e-05\\
1.224	1.98583e-05\\
1.226	1.94973e-05\\
1.228	1.91368e-05\\
1.23	1.87771e-05\\
1.232	1.84704e-05\\
1.234	1.82001e-05\\
1.236	1.79335e-05\\
1.238	1.76706e-05\\
1.24	1.74114e-05\\
1.242	1.71556e-05\\
1.244	1.69035e-05\\
1.246	1.66548e-05\\
1.248	1.64095e-05\\
1.25	1.61677e-05\\
1.252	1.59292e-05\\
1.254	1.56941e-05\\
1.256	1.54623e-05\\
1.258	1.52337e-05\\
1.26	1.50083e-05\\
1.262	1.47862e-05\\
1.264	1.45671e-05\\
1.266	1.43512e-05\\
1.268	1.41384e-05\\
1.27	1.39286e-05\\
1.272	1.37218e-05\\
1.274	1.3518e-05\\
1.276	1.33172e-05\\
1.278	1.31192e-05\\
1.28	1.29241e-05\\
1.282	1.27319e-05\\
1.284	1.25424e-05\\
1.286	1.23558e-05\\
1.288	1.21718e-05\\
1.29	1.19906e-05\\
1.292	1.18121e-05\\
1.294	1.16362e-05\\
1.296	1.14629e-05\\
1.298	1.12922e-05\\
1.3	1.1124e-05\\
1.302	1.09584e-05\\
1.304	1.07953e-05\\
1.306	1.06346e-05\\
1.308	1.04763e-05\\
1.31	1.03204e-05\\
1.312	1.01669e-05\\
1.314	1.00158e-05\\
1.316	9.86691e-06\\
1.318	9.72033e-06\\
1.32	9.57598e-06\\
1.322	9.43386e-06\\
1.324	9.29391e-06\\
1.326	9.15613e-06\\
1.328	9.02047e-06\\
1.33	8.88691e-06\\
1.332	8.75541e-06\\
1.334	8.62596e-06\\
1.336	8.49852e-06\\
1.338	8.37306e-06\\
1.34	8.24956e-06\\
1.342	8.12798e-06\\
1.344	8.0083e-06\\
1.346	7.89049e-06\\
1.348	7.77452e-06\\
1.35	7.66036e-06\\
1.352	7.54799e-06\\
1.354	7.43738e-06\\
1.356	7.32851e-06\\
1.358	7.22134e-06\\
1.36	7.11584e-06\\
1.362	7.012e-06\\
1.364	6.90979e-06\\
1.366	6.80918e-06\\
1.368	6.71014e-06\\
1.37	6.61265e-06\\
1.372	6.51668e-06\\
1.374	6.42221e-06\\
1.376	6.32921e-06\\
1.378	6.23766e-06\\
1.38	6.14753e-06\\
1.382	6.05881e-06\\
1.384	5.97146e-06\\
1.386	5.88546e-06\\
1.388	5.80079e-06\\
1.39	5.71742e-06\\
1.392	5.63534e-06\\
1.394	5.55452e-06\\
1.396	5.47493e-06\\
1.398	5.39656e-06\\
1.4	5.31939e-06\\
1.402	5.24339e-06\\
1.404	5.16854e-06\\
1.406	5.09482e-06\\
1.408	5.02221e-06\\
1.41	4.9507e-06\\
1.412	4.88025e-06\\
1.414	4.81086e-06\\
1.416	4.7425e-06\\
1.418	4.67515e-06\\
1.42	4.6088e-06\\
1.422	4.54343e-06\\
1.424	4.47901e-06\\
1.426	4.41554e-06\\
1.428	4.353e-06\\
1.43	4.29136e-06\\
1.432	4.23061e-06\\
1.434	4.17075e-06\\
1.436	4.11174e-06\\
1.438	4.05358e-06\\
1.44	3.99625e-06\\
1.442	3.93973e-06\\
1.444	3.88402e-06\\
1.446	3.82909e-06\\
1.448	3.77494e-06\\
1.45	3.72155e-06\\
1.452	3.6689e-06\\
1.454	3.61699e-06\\
1.456	3.5658e-06\\
1.458	3.51532e-06\\
1.46	3.46554e-06\\
1.462	3.41645e-06\\
1.464	3.36803e-06\\
1.466	3.32028e-06\\
1.468	3.27317e-06\\
1.47	3.22672e-06\\
1.472	3.18089e-06\\
1.474	3.13569e-06\\
1.476	3.0911e-06\\
1.478	3.04712e-06\\
1.48	3.00373e-06\\
1.482	2.96093e-06\\
1.484	2.9187e-06\\
1.486	2.87705e-06\\
1.488	2.83595e-06\\
1.49	2.79541e-06\\
1.492	2.75541e-06\\
1.494	2.71595e-06\\
1.496	2.67702e-06\\
1.498	2.63862e-06\\
1.5	2.60073e-06\\
1.502	2.56334e-06\\
1.504	2.52646e-06\\
1.506	2.49008e-06\\
1.508	2.45418e-06\\
1.51	2.41877e-06\\
1.512	2.38383e-06\\
1.514	2.34937e-06\\
1.516	2.31537e-06\\
1.518	2.28183e-06\\
1.52	2.24874e-06\\
1.522	2.2161e-06\\
1.524	2.18391e-06\\
1.526	2.15215e-06\\
1.528	2.12082e-06\\
1.53	2.08993e-06\\
1.532	2.05945e-06\\
1.534	2.02939e-06\\
1.536	1.99974e-06\\
1.538	1.97051e-06\\
1.54	1.94167e-06\\
1.542	1.91323e-06\\
1.544	1.88519e-06\\
1.546	1.85754e-06\\
1.548	1.83027e-06\\
1.55	1.80338e-06\\
1.552	1.77687e-06\\
1.554	1.75073e-06\\
1.556	1.72496e-06\\
1.558	1.69955e-06\\
1.56	1.6745e-06\\
1.562	1.64981e-06\\
1.564	1.62546e-06\\
1.566	1.60147e-06\\
1.568	1.57782e-06\\
1.57	1.5545e-06\\
1.572	1.53152e-06\\
1.574	1.50888e-06\\
1.576	1.48656e-06\\
1.578	1.46456e-06\\
1.58	1.44289e-06\\
1.582	1.42152e-06\\
1.584	1.40048e-06\\
1.586	1.37974e-06\\
1.588	1.3593e-06\\
1.59	1.33916e-06\\
1.592	1.31932e-06\\
1.594	1.29978e-06\\
1.596	1.28052e-06\\
1.598	1.26155e-06\\
1.6	1.24287e-06\\
1.602	1.22446e-06\\
1.604	1.20632e-06\\
1.606	1.18846e-06\\
1.608	1.17086e-06\\
1.61	1.15353e-06\\
1.612	1.13646e-06\\
1.614	1.11965e-06\\
1.616	1.10309e-06\\
1.618	1.08678e-06\\
1.62	1.07072e-06\\
1.622	1.0549e-06\\
1.624	1.03933e-06\\
1.626	1.02398e-06\\
1.628	1.00888e-06\\
1.63	9.93999e-07\\
1.632	9.79348e-07\\
1.634	9.64921e-07\\
1.636	9.50714e-07\\
1.638	9.36725e-07\\
1.64	9.22949e-07\\
1.642	9.09384e-07\\
1.644	8.96026e-07\\
1.646	8.82874e-07\\
1.648	8.69922e-07\\
1.65	8.57169e-07\\
1.652	8.44611e-07\\
1.654	8.32246e-07\\
1.656	8.2007e-07\\
1.658	8.08081e-07\\
1.66	7.96276e-07\\
1.662	7.84651e-07\\
1.664	7.73204e-07\\
1.666	7.61932e-07\\
1.668	7.50833e-07\\
1.67	7.39903e-07\\
1.672	7.2914e-07\\
1.674	7.18541e-07\\
1.676	7.08104e-07\\
1.678	6.97826e-07\\
1.68	6.87705e-07\\
1.682	6.77737e-07\\
1.684	6.67921e-07\\
1.686	6.58253e-07\\
1.688	6.48732e-07\\
1.69	6.39355e-07\\
1.692	6.3012e-07\\
1.694	6.21024e-07\\
1.696	6.12065e-07\\
1.698	6.03241e-07\\
1.7	5.9455e-07\\
1.702	5.85989e-07\\
1.704	5.77556e-07\\
1.706	5.69249e-07\\
1.708	5.61066e-07\\
1.71	5.53005e-07\\
1.712	5.45064e-07\\
1.714	5.3724e-07\\
1.716	5.29533e-07\\
1.718	5.21939e-07\\
1.72	5.14458e-07\\
1.722	5.07086e-07\\
1.724	4.99823e-07\\
1.726	4.92667e-07\\
1.728	4.85615e-07\\
1.73	4.78667e-07\\
1.732	4.71819e-07\\
1.734	4.65072e-07\\
1.736	4.58422e-07\\
1.738	4.51868e-07\\
1.74	4.4541e-07\\
1.742	4.39044e-07\\
1.744	4.32771e-07\\
1.746	4.26587e-07\\
1.748	4.20493e-07\\
1.75	4.14485e-07\\
1.752	4.08564e-07\\
1.754	4.02726e-07\\
1.756	3.96973e-07\\
1.758	3.913e-07\\
1.76	3.85709e-07\\
1.762	3.80196e-07\\
1.764	3.74762e-07\\
1.766	3.69404e-07\\
1.768	3.64122e-07\\
1.77	3.58915e-07\\
1.772	3.5378e-07\\
1.774	3.48718e-07\\
1.776	3.43726e-07\\
1.778	3.38805e-07\\
1.78	3.33952e-07\\
1.782	3.29167e-07\\
1.784	3.24449e-07\\
1.786	3.19797e-07\\
1.788	3.1521e-07\\
1.79	3.10686e-07\\
1.792	3.06226e-07\\
1.794	3.01827e-07\\
1.796	2.9749e-07\\
1.798	2.93212e-07\\
1.8	2.88995e-07\\
1.802	2.84835e-07\\
1.804	2.80734e-07\\
1.806	2.76689e-07\\
1.808	2.72701e-07\\
1.81	2.68767e-07\\
1.812	2.64889e-07\\
1.814	2.61064e-07\\
1.816	2.57292e-07\\
1.818	2.53573e-07\\
1.82	2.49905e-07\\
1.822	2.46289e-07\\
1.824	2.42722e-07\\
1.826	2.39205e-07\\
1.828	2.35738e-07\\
1.83	2.32318e-07\\
1.832	2.28946e-07\\
1.834	2.25622e-07\\
1.836	2.22343e-07\\
1.838	2.19111e-07\\
1.84	2.15924e-07\\
1.842	2.12781e-07\\
1.844	2.09683e-07\\
1.846	2.06628e-07\\
1.848	2.03616e-07\\
1.85	2.00647e-07\\
1.852	1.97719e-07\\
1.854	1.94833e-07\\
1.856	1.91988e-07\\
1.858	1.89183e-07\\
1.86	1.86418e-07\\
1.862	1.83692e-07\\
1.864	1.81005e-07\\
1.866	1.78356e-07\\
1.868	1.75745e-07\\
1.87	1.73172e-07\\
1.872	1.70635e-07\\
1.874	1.68135e-07\\
1.876	1.6567e-07\\
1.878	1.63242e-07\\
1.88	1.60848e-07\\
1.882	1.58489e-07\\
1.884	1.56164e-07\\
1.886	1.53872e-07\\
1.888	1.51614e-07\\
1.89	1.49389e-07\\
1.892	1.47196e-07\\
1.894	1.45035e-07\\
1.896	1.42906e-07\\
1.898	1.40807e-07\\
1.9	1.3874e-07\\
1.902	1.36703e-07\\
1.904	1.34696e-07\\
1.906	1.32718e-07\\
1.908	1.30769e-07\\
1.91	1.2885e-07\\
1.912	1.26958e-07\\
1.914	1.25095e-07\\
1.916	1.23259e-07\\
1.918	1.2145e-07\\
1.92	1.19668e-07\\
1.922	1.17913e-07\\
1.924	1.16184e-07\\
1.926	1.1448e-07\\
1.928	1.12802e-07\\
1.93	1.11149e-07\\
1.932	1.0952e-07\\
1.934	1.07916e-07\\
1.936	1.06336e-07\\
1.938	1.04779e-07\\
1.94	1.03246e-07\\
1.942	1.01735e-07\\
1.944	1.00247e-07\\
1.946	9.87817e-08\\
1.948	9.7338e-08\\
1.95	9.59159e-08\\
1.952	9.45152e-08\\
1.954	9.31354e-08\\
1.956	9.17763e-08\\
1.958	9.04375e-08\\
1.96	8.91188e-08\\
1.962	8.78199e-08\\
1.964	8.65404e-08\\
1.966	8.528e-08\\
1.968	8.40385e-08\\
1.97	8.28156e-08\\
1.972	8.16109e-08\\
1.974	8.04243e-08\\
1.976	7.92554e-08\\
1.978	7.81039e-08\\
1.98	7.69697e-08\\
1.982	7.58523e-08\\
1.984	7.47516e-08\\
1.986	7.36673e-08\\
1.988	7.25991e-08\\
1.99	7.15468e-08\\
1.992	7.05101e-08\\
1.994	6.94888e-08\\
1.996	6.84827e-08\\
1.998	6.74914e-08\\
2	6.65149e-08\\
};
\end{axis}

\begin{axis}[%
width={\figscale*1.528in},
height={\figscale*0.46in},
at={(\figscale*2.841in,\figscale*1.125in)},
scale only axis,
separate axis lines,
every outer x axis line/.append style={black},
every x tick label/.append style={font=\figticfont},
every x tick/.append style={black},
xmin=-0.025,
xmax=2,
xtick={0,0.5,1,1.5,2},
xticklabels={\empty},
every outer y axis line/.append style={black},
every y tick label/.append style={font=\figticfont},
every y tick/.append style={black},
ymin=0,
ymax=0.0572362,
ytick={0,0.05},
axis background/.style={fill=white},
xmajorgrids,
ymajorgrids,
grid style={white!55!black, opacity=0.3},
ylabel style={rotate=-90, at={(axis description cs:-0.145,0.5)}, anchor=east}
]
\addplot [color=mycolor1, line width=1.4pt, forget plot]
  table[row sep=crcr]{%
-0.024	0\\
-0.022	0\\
-0.02	0\\
-0.018	0\\
-0.016	0\\
-0.014	0\\
-0.012	0\\
-0.01	0\\
-0.008	0\\
-0.006	0\\
-0.004	0\\
-0.002	0\\
0	0\\
0.002	0.00916184\\
0.004	0.0147068\\
0.006	0.0198658\\
0.008	0.0237563\\
0.01	0.0264987\\
0.012	0.0282393\\
0.014	0.0291285\\
0.016	0.0293161\\
0.018	0.028959\\
0.02	0.0291501\\
0.022	0.0292758\\
0.024	0.0303387\\
0.026	0.0328616\\
0.028	0.0353071\\
0.03	0.0376644\\
0.032	0.0399233\\
0.034	0.0420746\\
0.036	0.0441097\\
0.038	0.0460205\\
0.04	0.0477997\\
0.042	0.0494406\\
0.044	0.0509372\\
0.046	0.0522841\\
0.048	0.0534765\\
0.05	0.0545103\\
0.052	0.0553821\\
0.054	0.0560891\\
0.056	0.0566291\\
0.058	0.0570008\\
0.06	0.0572032\\
0.062	0.0572362\\
0.064	0.0571003\\
0.066	0.0567967\\
0.068	0.0563271\\
0.07	0.0556939\\
0.072	0.0549001\\
0.074	0.0539493\\
0.076	0.0528457\\
0.078	0.051594\\
0.08	0.0501995\\
0.082	0.048668\\
0.084	0.0470057\\
0.086	0.0452194\\
0.088	0.043316\\
0.09	0.0413031\\
0.092	0.0391886\\
0.094	0.0369806\\
0.096	0.0346876\\
0.098	0.0323181\\
0.1	0.0298811\\
0.102	0.0273856\\
0.104	0.0250678\\
0.106	0.0245539\\
0.108	0.0239887\\
0.11	0.0233743\\
0.112	0.0227128\\
0.114	0.0220064\\
0.116	0.0212577\\
0.118	0.0204691\\
0.12	0.019643\\
0.122	0.019207\\
0.124	0.0194165\\
0.126	0.0196799\\
0.128	0.0199091\\
0.13	0.0201032\\
0.132	0.0202619\\
0.134	0.0203844\\
0.136	0.0204704\\
0.138	0.0205194\\
0.14	0.0205312\\
0.142	0.0210121\\
0.144	0.0227402\\
0.146	0.0243564\\
0.148	0.0258565\\
0.15	0.0272371\\
0.152	0.028495\\
0.154	0.0296277\\
0.156	0.030633\\
0.158	0.0315093\\
0.16	0.0322555\\
0.162	0.0328707\\
0.164	0.0333549\\
0.166	0.0337082\\
0.168	0.0339314\\
0.17	0.0340256\\
0.172	0.0339923\\
0.174	0.0338337\\
0.176	0.0335521\\
0.178	0.0331504\\
0.18	0.0326318\\
0.182	0.0319999\\
0.184	0.0312587\\
0.186	0.0304123\\
0.188	0.0294654\\
0.19	0.0284229\\
0.192	0.0272899\\
0.194	0.0260718\\
0.196	0.0247743\\
0.198	0.0234031\\
0.2	0.0219643\\
0.202	0.020464\\
0.204	0.0189085\\
0.206	0.0173043\\
0.208	0.0156578\\
0.21	0.0140381\\
0.212	0.0140376\\
0.214	0.0139946\\
0.216	0.0139095\\
0.218	0.0137835\\
0.22	0.0136173\\
0.222	0.0134121\\
0.224	0.0131691\\
0.226	0.0128893\\
0.228	0.0125743\\
0.23	0.0122255\\
0.232	0.0118442\\
0.234	0.0114321\\
0.236	0.0109908\\
0.238	0.010522\\
0.24	0.0111848\\
0.242	0.012589\\
0.244	0.0139302\\
0.246	0.0152044\\
0.248	0.0164079\\
0.25	0.0175373\\
0.252	0.0185896\\
0.254	0.0195619\\
0.256	0.0204519\\
0.258	0.0212574\\
0.26	0.0219766\\
0.262	0.0226082\\
0.264	0.0231509\\
0.266	0.023604\\
0.268	0.023967\\
0.27	0.0242398\\
0.272	0.0244226\\
0.274	0.024516\\
0.276	0.0245207\\
0.278	0.0244379\\
0.28	0.024269\\
0.282	0.0240158\\
0.284	0.0236803\\
0.286	0.0232646\\
0.288	0.0227714\\
0.29	0.0222035\\
0.292	0.0215638\\
0.294	0.0208555\\
0.296	0.0200821\\
0.298	0.0192472\\
0.3	0.0183546\\
0.302	0.0174082\\
0.304	0.0164121\\
0.306	0.0153705\\
0.308	0.0142877\\
0.31	0.0131681\\
0.312	0.0120162\\
0.314	0.0108366\\
0.316	0.00963376\\
0.318	0.00948972\\
0.32	0.00958629\\
0.322	0.00965208\\
0.324	0.00968726\\
0.326	0.00969208\\
0.328	0.00966689\\
0.33	0.0096121\\
0.332	0.00952825\\
0.334	0.00941592\\
0.336	0.0092758\\
0.338	0.00910863\\
0.34	0.00891527\\
0.342	0.00869659\\
0.344	0.00845357\\
0.346	0.00818725\\
0.348	0.00912998\\
0.35	0.0100661\\
0.352	0.010954\\
0.354	0.0117911\\
0.356	0.0125752\\
0.358	0.0133041\\
0.36	0.013976\\
0.362	0.0145892\\
0.364	0.0151424\\
0.366	0.0156342\\
0.368	0.0160639\\
0.37	0.0164306\\
0.372	0.0167339\\
0.374	0.0169736\\
0.376	0.0171496\\
0.378	0.0172622\\
0.38	0.0173117\\
0.382	0.0172989\\
0.384	0.0172246\\
0.386	0.0170899\\
0.388	0.0168961\\
0.39	0.0166445\\
0.392	0.016337\\
0.394	0.0159753\\
0.396	0.0155615\\
0.398	0.0150976\\
0.4	0.014586\\
0.402	0.0140292\\
0.404	0.0134298\\
0.406	0.0127903\\
0.408	0.0121137\\
0.41	0.0114029\\
0.412	0.0106608\\
0.414	0.00989047\\
0.416	0.00966763\\
0.418	0.00945173\\
0.42	0.00921467\\
0.422	0.00895734\\
0.424	0.00868064\\
0.426	0.00838555\\
0.428	0.00807305\\
0.43	0.00774418\\
0.432	0.00739999\\
0.434	0.00704156\\
0.436	0.00690853\\
0.438	0.00692113\\
0.44	0.00691165\\
0.442	0.00688037\\
0.444	0.00682761\\
0.446	0.00675377\\
0.448	0.0066593\\
0.45	0.00654473\\
0.452	0.00641063\\
0.454	0.00659526\\
0.456	0.00724754\\
0.458	0.00786546\\
0.46	0.00844727\\
0.462	0.00899137\\
0.464	0.00949633\\
0.466	0.00996086\\
0.468	0.0103838\\
0.47	0.0107643\\
0.472	0.0111015\\
0.474	0.0113948\\
0.476	0.0116436\\
0.478	0.0118478\\
0.48	0.0120071\\
0.482	0.0121217\\
0.484	0.0121916\\
0.486	0.0122172\\
0.488	0.012199\\
0.49	0.0121377\\
0.492	0.0120341\\
0.494	0.011889\\
0.496	0.0117036\\
0.498	0.011479\\
0.5	0.0112165\\
0.502	0.0109176\\
0.504	0.0105839\\
0.506	0.0102169\\
0.508	0.00981852\\
0.51	0.00939045\\
0.512	0.00893467\\
0.514	0.00845316\\
0.516	0.00794798\\
0.518	0.00742125\\
0.52	0.00687512\\
0.522	0.006858\\
0.524	0.00687319\\
0.526	0.00687084\\
0.528	0.0068512\\
0.53	0.00681457\\
0.532	0.0067613\\
0.534	0.00669175\\
0.536	0.00660635\\
0.538	0.00650554\\
0.54	0.0063898\\
0.542	0.00625965\\
0.544	0.00611563\\
0.546	0.00595831\\
0.548	0.00578828\\
0.55	0.00560615\\
0.552	0.00541259\\
0.554	0.00520823\\
0.556	0.00507328\\
0.558	0.00501509\\
0.56	0.00494145\\
0.562	0.005161\\
0.564	0.00559027\\
0.566	0.00599388\\
0.568	0.00637072\\
0.57	0.00671981\\
0.572	0.00704027\\
0.574	0.00733136\\
0.576	0.00759242\\
0.578	0.00782293\\
0.58	0.00802247\\
0.582	0.00819074\\
0.584	0.00832758\\
0.586	0.0084329\\
0.588	0.00850677\\
0.59	0.00854933\\
0.592	0.00856086\\
0.594	0.00854174\\
0.596	0.00849244\\
0.598	0.00841355\\
0.6	0.00830576\\
0.602	0.00816982\\
0.604	0.00800661\\
0.606	0.00781709\\
0.608	0.00760227\\
0.61	0.00736328\\
0.612	0.00710128\\
0.614	0.00681753\\
0.616	0.00651334\\
0.618	0.00619006\\
0.62	0.00584911\\
0.622	0.00549195\\
0.624	0.00512008\\
0.626	0.00473501\\
0.628	0.00433831\\
0.63	0.00393155\\
0.632	0.00378117\\
0.634	0.00388037\\
0.636	0.00396828\\
0.638	0.00404487\\
0.64	0.00411011\\
0.642	0.00416402\\
0.644	0.00420662\\
0.646	0.00423799\\
0.648	0.00425822\\
0.65	0.00426744\\
0.652	0.00426578\\
0.654	0.00425343\\
0.656	0.00423059\\
0.658	0.00419746\\
0.66	0.0041543\\
0.662	0.00410137\\
0.664	0.00403895\\
0.666	0.00396736\\
0.668	0.00388691\\
0.67	0.00394824\\
0.672	0.0042249\\
0.674	0.00448263\\
0.676	0.00472076\\
0.678	0.00493871\\
0.68	0.00513598\\
0.682	0.00531213\\
0.684	0.00546683\\
0.686	0.00559981\\
0.688	0.0057109\\
0.69	0.0058\\
0.692	0.00586709\\
0.694	0.00591223\\
0.696	0.00593555\\
0.698	0.00593728\\
0.7	0.0059177\\
0.702	0.00587717\\
0.704	0.00581613\\
0.706	0.00573508\\
0.708	0.00563458\\
0.71	0.00551526\\
0.712	0.00537781\\
0.714	0.00522296\\
0.716	0.00505152\\
0.718	0.00486432\\
0.72	0.00466225\\
0.722	0.00444625\\
0.724	0.00421729\\
0.726	0.00397635\\
0.728	0.00372448\\
0.73	0.00346272\\
0.732	0.00319216\\
0.734	0.00291389\\
0.736	0.00262901\\
0.738	0.00233864\\
0.74	0.0020439\\
0.742	0.001943\\
0.744	0.00203979\\
0.746	0.00216874\\
0.748	0.00228959\\
0.75	0.00240204\\
0.752	0.0025058\\
0.754	0.00260064\\
0.756	0.00268633\\
0.758	0.00276269\\
0.76	0.00282957\\
0.762	0.00288684\\
0.764	0.00293442\\
0.766	0.00297224\\
0.768	0.00300028\\
0.77	0.00301854\\
0.772	0.00302706\\
0.774	0.00302591\\
0.776	0.00301518\\
0.778	0.00299499\\
0.78	0.00316555\\
0.782	0.00332475\\
0.784	0.00346972\\
0.786	0.00360014\\
0.788	0.00371572\\
0.79	0.00381626\\
0.792	0.0039016\\
0.794	0.00397164\\
0.796	0.00402634\\
0.798	0.00406571\\
0.8	0.00408982\\
0.802	0.00409879\\
0.804	0.0040928\\
0.806	0.00407208\\
0.808	0.0040369\\
0.81	0.00398759\\
0.812	0.00392452\\
0.814	0.00384811\\
0.816	0.00375881\\
0.818	0.00365714\\
0.82	0.00354362\\
0.822	0.00341883\\
0.824	0.00328338\\
0.826	0.00313791\\
0.828	0.00298307\\
0.83	0.00281956\\
0.832	0.00264809\\
0.834	0.00246939\\
0.836	0.00228421\\
0.838	0.00209331\\
0.84	0.00189746\\
0.842	0.00169743\\
0.844	0.00149401\\
0.846	0.00128799\\
0.848	0.00118185\\
0.85	0.00124144\\
0.852	0.0013556\\
0.854	0.00146481\\
0.856	0.00156837\\
0.858	0.00166597\\
0.86	0.00175735\\
0.862	0.00184227\\
0.864	0.00192055\\
0.866	0.00199199\\
0.868	0.00205643\\
0.87	0.00211374\\
0.872	0.00216381\\
0.874	0.00220656\\
0.876	0.00224191\\
0.878	0.00226983\\
0.88	0.0022903\\
0.882	0.00230334\\
0.884	0.00230898\\
0.886	0.00230727\\
0.888	0.00237324\\
0.89	0.0024671\\
0.892	0.00255066\\
0.894	0.00262376\\
0.896	0.00268627\\
0.898	0.0027381\\
0.9	0.00277921\\
0.902	0.00280961\\
0.904	0.00282931\\
0.906	0.0028384\\
0.908	0.00283699\\
0.91	0.00282522\\
0.912	0.00280327\\
0.914	0.00277137\\
0.916	0.00272977\\
0.918	0.00267874\\
0.92	0.0026186\\
0.922	0.00254969\\
0.924	0.00247239\\
0.926	0.00238709\\
0.928	0.00229421\\
0.93	0.00219418\\
0.932	0.00208748\\
0.934	0.00197457\\
0.936	0.00185596\\
0.938	0.00173216\\
0.94	0.00160368\\
0.942	0.00147107\\
0.944	0.00133485\\
0.946	0.00119559\\
0.948	0.00105381\\
0.95	0.000910093\\
0.952	0.000814398\\
0.954	0.000808241\\
0.956	0.000799641\\
0.958	0.000788658\\
0.96	0.00084294\\
0.962	0.0009327\\
0.964	0.00101888\\
0.966	0.00110096\\
0.968	0.00117871\\
0.97	0.00125192\\
0.972	0.00132043\\
0.974	0.00138405\\
0.976	0.00144264\\
0.978	0.00149606\\
0.98	0.0015442\\
0.982	0.00158697\\
0.984	0.00162427\\
0.986	0.00165606\\
0.988	0.00168228\\
0.99	0.00170292\\
0.992	0.00171796\\
0.994	0.00172846\\
0.996	0.00178769\\
0.998	0.00183955\\
1	0.00188395\\
1.002	0.00192084\\
1.004	0.00195016\\
1.006	0.00197192\\
1.008	0.00198614\\
1.01	0.00199285\\
1.012	0.00199213\\
1.014	0.00198409\\
1.016	0.00196885\\
1.018	0.00194655\\
1.02	0.00191738\\
1.022	0.00188152\\
1.024	0.0018392\\
1.026	0.00179065\\
1.028	0.00173614\\
1.03	0.00167593\\
1.032	0.00161034\\
1.034	0.00153966\\
1.036	0.00146422\\
1.038	0.00138436\\
1.04	0.00130044\\
1.042	0.0012128\\
1.044	0.00112183\\
1.046	0.00102789\\
1.048	0.000931373\\
1.05	0.000832664\\
1.052	0.000732153\\
1.054	0.000667354\\
1.056	0.00064751\\
1.058	0.000626528\\
1.06	0.000604487\\
1.062	0.000581462\\
1.064	0.000559497\\
1.066	0.000555937\\
1.068	0.000550704\\
1.07	0.00055034\\
1.072	0.000619477\\
1.074	0.000686012\\
1.076	0.000749639\\
1.078	0.000810191\\
1.08	0.000867508\\
1.082	0.000921443\\
1.084	0.000971864\\
1.086	0.00101865\\
1.088	0.00106169\\
1.09	0.00110089\\
1.092	0.00113617\\
1.094	0.00116746\\
1.096	0.0011958\\
1.098	0.00124288\\
1.1	0.00128478\\
1.102	0.00132143\\
1.104	0.00135274\\
1.106	0.00137867\\
1.108	0.0013992\\
1.11	0.00141432\\
1.112	0.00142403\\
1.114	0.00142837\\
1.116	0.00142739\\
1.118	0.00142115\\
1.12	0.00140976\\
1.122	0.00139331\\
1.124	0.00137192\\
1.126	0.00134575\\
1.128	0.00131493\\
1.13	0.00127964\\
1.132	0.00124007\\
1.134	0.00119642\\
1.136	0.00114889\\
1.138	0.0010977\\
1.14	0.0010431\\
1.142	0.000985319\\
1.144	0.000924612\\
1.146	0.000861238\\
1.148	0.000795461\\
1.15	0.000727553\\
1.152	0.000657788\\
1.154	0.000586446\\
1.156	0.000513807\\
1.158	0.000440155\\
1.16	0.000398259\\
1.162	0.000399038\\
1.164	0.000398933\\
1.166	0.000397963\\
1.168	0.000396152\\
1.17	0.000393521\\
1.172	0.000390095\\
1.174	0.0003859\\
1.176	0.000380962\\
1.178	0.000375309\\
1.18	0.000368967\\
1.182	0.000422653\\
1.184	0.000485026\\
1.186	0.000544873\\
1.188	0.000602004\\
1.19	0.000656243\\
1.192	0.000707426\\
1.194	0.000755402\\
1.196	0.000800033\\
1.198	0.000841195\\
1.2	0.00087878\\
1.202	0.000912692\\
1.204	0.000942851\\
1.206	0.000969191\\
1.208	0.000991661\\
1.21	0.00101022\\
1.212	0.00102486\\
1.214	0.00103556\\
1.216	0.00104234\\
1.218	0.00104521\\
1.22	0.0010442\\
1.222	0.00103938\\
1.224	0.0010308\\
1.226	0.00101854\\
1.228	0.00100269\\
1.23	0.000983343\\
1.232	0.000960615\\
1.234	0.000934628\\
1.236	0.000905517\\
1.238	0.000873425\\
1.24	0.000838503\\
1.242	0.000800913\\
1.244	0.000760824\\
1.246	0.00071841\\
1.248	0.000673855\\
1.25	0.000627347\\
1.252	0.000585548\\
1.254	0.000547775\\
1.256	0.000508716\\
1.258	0.000468512\\
1.26	0.000427305\\
1.262	0.000385239\\
1.264	0.000342459\\
1.266	0.00029911\\
1.268	0.000255339\\
1.27	0.00021129\\
1.272	0.00019798\\
1.274	0.000202749\\
1.276	0.00020678\\
1.278	0.000210076\\
1.28	0.000212638\\
1.282	0.000219029\\
1.284	0.000268192\\
1.286	0.000315816\\
1.288	0.000361744\\
1.29	0.000405829\\
1.292	0.000447932\\
1.294	0.000487922\\
1.296	0.000525677\\
1.298	0.000561087\\
1.3	0.000594048\\
1.302	0.000624468\\
1.304	0.000652267\\
1.306	0.000677373\\
1.308	0.000699726\\
1.31	0.000719276\\
1.312	0.000735983\\
1.314	0.00074982\\
1.316	0.000760769\\
1.318	0.000768823\\
1.32	0.000773986\\
1.322	0.000776273\\
1.324	0.000775708\\
1.326	0.000772325\\
1.328	0.000766169\\
1.33	0.000757294\\
1.332	0.000745763\\
1.334	0.000731649\\
1.336	0.000715032\\
1.338	0.000696001\\
1.34	0.000674653\\
1.342	0.000651091\\
1.344	0.000625427\\
1.346	0.000597777\\
1.348	0.000579039\\
1.35	0.00056239\\
1.352	0.00054415\\
1.354	0.000524397\\
1.356	0.000503214\\
1.358	0.000480687\\
1.36	0.000456906\\
1.362	0.000431963\\
1.364	0.000405952\\
1.366	0.000378971\\
1.368	0.000351119\\
1.37	0.000322496\\
1.372	0.000293205\\
1.374	0.000263347\\
1.376	0.000233026\\
1.378	0.000202345\\
1.38	0.000171408\\
1.382	0.000140318\\
1.384	0.000120429\\
1.386	0.000158137\\
1.388	0.000194852\\
1.39	0.000230454\\
1.392	0.000264826\\
1.394	0.000297858\\
1.396	0.000329444\\
1.398	0.000359488\\
1.4	0.000387897\\
1.402	0.000414588\\
1.404	0.000439483\\
1.406	0.000462512\\
1.408	0.000483614\\
1.41	0.000502732\\
1.412	0.000519822\\
1.414	0.000534843\\
1.416	0.000547765\\
1.418	0.000558565\\
1.42	0.000567227\\
1.422	0.000573744\\
1.424	0.000578117\\
1.426	0.000580354\\
1.428	0.000580471\\
1.43	0.000578491\\
1.432	0.000574445\\
1.434	0.00056837\\
1.436	0.000560312\\
1.438	0.000550322\\
1.44	0.000538457\\
1.442	0.000524782\\
1.444	0.000509365\\
1.446	0.000492283\\
1.448	0.000473615\\
1.45	0.000453446\\
1.452	0.000435031\\
1.454	0.000427766\\
1.456	0.000419218\\
1.458	0.00040943\\
1.46	0.000398451\\
1.462	0.000386332\\
1.464	0.000373126\\
1.466	0.000358891\\
1.468	0.000343686\\
1.47	0.000327572\\
1.472	0.000310615\\
1.474	0.000292878\\
1.476	0.000274431\\
1.478	0.000255342\\
1.48	0.000235681\\
1.482	0.000215518\\
1.484	0.000194927\\
1.486	0.000173979\\
1.488	0.000152747\\
1.49	0.000131303\\
1.492	0.000135569\\
1.494	0.000162433\\
1.496	0.000188417\\
1.498	0.000213436\\
1.5	0.000237412\\
1.502	0.000260269\\
1.504	0.000281938\\
1.506	0.000302355\\
1.508	0.00032146\\
1.51	0.000339199\\
1.512	0.000355524\\
1.514	0.000370392\\
1.516	0.000383766\\
1.518	0.000395615\\
1.52	0.000405912\\
1.522	0.000414639\\
1.524	0.000421782\\
1.526	0.000427333\\
1.528	0.000431289\\
1.53	0.000433655\\
1.532	0.00043444\\
1.534	0.000433658\\
1.536	0.00043133\\
1.538	0.000427481\\
1.54	0.000422143\\
1.542	0.000415351\\
1.544	0.000407147\\
1.546	0.000397574\\
1.548	0.000386683\\
1.55	0.000374527\\
1.552	0.000361165\\
1.554	0.000346658\\
1.556	0.00033107\\
1.558	0.00031447\\
1.56	0.000305483\\
1.562	0.000301784\\
1.564	0.000297157\\
1.566	0.000291629\\
1.568	0.000285231\\
1.57	0.000277995\\
1.572	0.000269956\\
1.574	0.00026115\\
1.576	0.000251617\\
1.578	0.000241397\\
1.58	0.000230532\\
1.582	0.000219067\\
1.584	0.000207046\\
1.586	0.000194517\\
1.588	0.000181526\\
1.59	0.000168123\\
1.592	0.000154355\\
1.594	0.000140274\\
1.596	0.000125928\\
1.598	0.000111368\\
1.6	0.00012983\\
1.602	0.000148865\\
1.604	0.000167155\\
1.606	0.000184645\\
1.608	0.000201279\\
1.61	0.000217009\\
1.612	0.000231788\\
1.614	0.000245574\\
1.616	0.000258328\\
1.618	0.000270018\\
1.62	0.000280613\\
1.622	0.000290088\\
1.624	0.00029842\\
1.626	0.000305594\\
1.628	0.000311597\\
1.63	0.00031642\\
1.632	0.00032006\\
1.634	0.000322517\\
1.636	0.000323794\\
1.638	0.000323902\\
1.64	0.000322853\\
1.642	0.000320663\\
1.644	0.000317353\\
1.646	0.000312948\\
1.648	0.000307476\\
1.65	0.000300968\\
1.652	0.000293459\\
1.654	0.000284988\\
1.656	0.000275594\\
1.658	0.000265323\\
1.66	0.000254221\\
1.662	0.000242337\\
1.664	0.000229721\\
1.666	0.000216428\\
1.668	0.000211445\\
1.67	0.000209785\\
1.672	0.000207466\\
1.674	0.000204505\\
1.676	0.00020092\\
1.678	0.000196731\\
1.68	0.00019196\\
1.682	0.000186631\\
1.684	0.000180768\\
1.686	0.000174398\\
1.688	0.000167549\\
1.69	0.000160249\\
1.692	0.000152528\\
1.694	0.000144418\\
1.696	0.00013595\\
1.698	0.000127156\\
1.7	0.00011807\\
1.702	0.000108725\\
1.704	9.91551e-05\\
1.706	0.000101096\\
1.708	0.000115054\\
1.71	0.000128445\\
1.712	0.000141227\\
1.714	0.000153363\\
1.716	0.000164816\\
1.718	0.000175553\\
1.72	0.000185544\\
1.722	0.000194761\\
1.724	0.000203182\\
1.726	0.000210784\\
1.728	0.000217549\\
1.73	0.000223464\\
1.732	0.000228517\\
1.734	0.0002327\\
1.736	0.000236006\\
1.738	0.000238436\\
1.74	0.000239989\\
1.742	0.000240671\\
1.744	0.000240489\\
1.746	0.000239452\\
1.748	0.000237575\\
1.75	0.000234874\\
1.752	0.000231368\\
1.754	0.000227078\\
1.756	0.000222029\\
1.758	0.000216246\\
1.76	0.00020976\\
1.762	0.0002026\\
1.764	0.000194801\\
1.766	0.000186397\\
1.768	0.000177425\\
1.77	0.000167923\\
1.772	0.000157931\\
1.774	0.00014749\\
1.776	0.000144768\\
1.778	0.000144234\\
1.78	0.000143238\\
1.782	0.00014179\\
1.784	0.000139901\\
1.786	0.000137583\\
1.788	0.000134851\\
1.79	0.000131718\\
1.792	0.000128201\\
1.794	0.000124315\\
1.796	0.00012008\\
1.798	0.000115513\\
1.8	0.000110635\\
1.802	0.000105465\\
1.804	0.000100024\\
1.806	9.43336e-05\\
1.808	8.84162e-05\\
1.81	8.22939e-05\\
1.812	7.75143e-05\\
1.814	8.77486e-05\\
1.816	9.75551e-05\\
1.818	0.000106904\\
1.82	0.000115766\\
1.822	0.000124117\\
1.824	0.000131932\\
1.826	0.00013919\\
1.828	0.000145871\\
1.83	0.000151957\\
1.832	0.000157435\\
1.834	0.000162291\\
1.836	0.000166515\\
1.838	0.0001701\\
1.84	0.00017304\\
1.842	0.000175332\\
1.844	0.000176975\\
1.846	0.000177971\\
1.848	0.000178323\\
1.85	0.000178038\\
1.852	0.000177123\\
1.854	0.00017559\\
1.856	0.000173451\\
1.858	0.00017072\\
1.86	0.000167413\\
1.862	0.00016355\\
1.864	0.00015915\\
1.866	0.000154234\\
1.868	0.000148827\\
1.87	0.000142952\\
1.872	0.000136637\\
1.874	0.000129908\\
1.876	0.000122795\\
1.878	0.000115326\\
1.88	0.000107532\\
1.882	9.94459e-05\\
1.884	9.80944e-05\\
1.886	9.81554e-05\\
1.888	9.78972e-05\\
1.89	9.73252e-05\\
1.892	9.64456e-05\\
1.894	9.52658e-05\\
1.896	9.37937e-05\\
1.898	9.20383e-05\\
1.9	9.00094e-05\\
1.902	8.77174e-05\\
1.904	8.51735e-05\\
1.906	8.23894e-05\\
1.908	7.93778e-05\\
1.91	7.61515e-05\\
1.912	7.27242e-05\\
1.914	6.91097e-05\\
1.916	6.53225e-05\\
1.918	6.13774e-05\\
1.92	6.63672e-05\\
1.922	7.35526e-05\\
1.924	8.03954e-05\\
1.926	8.68753e-05\\
1.928	9.29733e-05\\
1.93	9.86722e-05\\
1.932	0.000103956\\
1.934	0.000108812\\
1.936	0.000113226\\
1.938	0.000117188\\
1.94	0.00012069\\
1.942	0.000123724\\
1.944	0.000126285\\
1.946	0.000128369\\
1.948	0.000129974\\
1.95	0.0001311\\
1.952	0.000131748\\
1.954	0.000131922\\
1.956	0.000131626\\
1.958	0.000130867\\
1.96	0.000129653\\
1.962	0.000127993\\
1.964	0.000125898\\
1.966	0.00012338\\
1.968	0.000120454\\
1.97	0.000117135\\
1.972	0.000113437\\
1.974	0.00010938\\
1.976	0.000104982\\
1.978	0.000100262\\
1.98	9.52395e-05\\
1.982	8.99371e-05\\
1.984	8.43764e-05\\
1.986	7.858e-05\\
1.988	7.25712e-05\\
1.99	6.63736e-05\\
1.992	6.55904e-05\\
1.994	6.59447e-05\\
1.996	6.60812e-05\\
1.998	6.60026e-05\\
2	6.57123e-05\\
};
\addplot [color=mycolor2, line width=1.0pt, forget plot]
  table[row sep=crcr]{%
-0.024	0\\
-0.022	0\\
-0.02	0\\
-0.018	0\\
-0.016	0\\
-0.014	0\\
-0.012	0\\
-0.01	0\\
-0.008	0\\
-0.006	0\\
-0.004	0\\
-0.002	0\\
0	0\\
0.002	0.00916184\\
0.004	0.0147068\\
0.006	0.0198658\\
0.008	0.0237561\\
0.01	0.0264983\\
0.012	0.0282385\\
0.014	0.0291271\\
0.016	0.0293137\\
0.018	0.0289565\\
0.02	0.0291464\\
0.022	0.0292712\\
0.024	0.0293698\\
0.026	0.029564\\
0.028	0.0301127\\
0.03	0.0303591\\
0.032	0.0308408\\
0.034	0.0322232\\
0.036	0.0335129\\
0.038	0.0347088\\
0.04	0.0358101\\
0.042	0.0368161\\
0.044	0.0377265\\
0.046	0.0385414\\
0.048	0.0392607\\
0.05	0.0398848\\
0.052	0.0404144\\
0.054	0.04085\\
0.056	0.0411927\\
0.058	0.0414436\\
0.06	0.0416038\\
0.062	0.0416749\\
0.064	0.0416585\\
0.066	0.0415562\\
0.068	0.0413701\\
0.07	0.0411023\\
0.072	0.0407549\\
0.074	0.0403304\\
0.076	0.0398313\\
0.078	0.0392603\\
0.08	0.0386201\\
0.082	0.0379138\\
0.084	0.0371443\\
0.086	0.0363148\\
0.088	0.0354285\\
0.09	0.0344887\\
0.092	0.0334989\\
0.094	0.0324624\\
0.096	0.0313828\\
0.098	0.0302637\\
0.1	0.0291085\\
0.102	0.0279209\\
0.104	0.0267045\\
0.106	0.0254627\\
0.108	0.0241993\\
0.11	0.0229178\\
0.112	0.0216215\\
0.114	0.0208166\\
0.116	0.0201395\\
0.118	0.0194291\\
0.12	0.018688\\
0.122	0.0179187\\
0.124	0.0171236\\
0.126	0.0166298\\
0.128	0.0163601\\
0.13	0.0160573\\
0.132	0.0157225\\
0.134	0.015357\\
0.136	0.0149619\\
0.138	0.0145386\\
0.14	0.0140885\\
0.142	0.0136129\\
0.144	0.0131135\\
0.146	0.0125917\\
0.148	0.0120491\\
0.15	0.0114874\\
0.152	0.0109081\\
0.154	0.010313\\
0.156	0.00970374\\
0.158	0.00908206\\
0.16	0.00844969\\
0.162	0.00780839\\
0.164	0.00765377\\
0.166	0.00833996\\
0.168	0.00898192\\
0.17	0.00957929\\
0.172	0.0101319\\
0.174	0.0106395\\
0.176	0.0111023\\
0.178	0.0115203\\
0.18	0.0118938\\
0.182	0.0122231\\
0.184	0.0125085\\
0.186	0.0127508\\
0.188	0.0129504\\
0.19	0.0131081\\
0.192	0.0132247\\
0.194	0.013301\\
0.196	0.0133381\\
0.198	0.0133369\\
0.2	0.0132986\\
0.202	0.0132243\\
0.204	0.0131152\\
0.206	0.0129726\\
0.208	0.0127979\\
0.21	0.0125923\\
0.212	0.0123573\\
0.214	0.0120943\\
0.216	0.0118048\\
0.218	0.0114904\\
0.22	0.0111524\\
0.222	0.0107926\\
0.224	0.0104124\\
0.226	0.0100327\\
0.228	0.00985471\\
0.23	0.00965267\\
0.232	0.00942774\\
0.234	0.0091811\\
0.236	0.00891393\\
0.238	0.00866382\\
0.24	0.00865151\\
0.242	0.00861372\\
0.244	0.008551\\
0.246	0.00846397\\
0.248	0.00835333\\
0.25	0.00821979\\
0.252	0.00806418\\
0.254	0.00788732\\
0.256	0.00769013\\
0.258	0.00747354\\
0.26	0.00723855\\
0.262	0.00698618\\
0.264	0.00671749\\
0.266	0.00643358\\
0.268	0.00613557\\
0.27	0.0058246\\
0.272	0.00550184\\
0.274	0.00516848\\
0.276	0.0048257\\
0.278	0.00447472\\
0.28	0.00411673\\
0.282	0.00375295\\
0.284	0.00338458\\
0.286	0.00349666\\
0.288	0.00380309\\
0.29	0.00409172\\
0.292	0.00436223\\
0.294	0.00461433\\
0.296	0.00484777\\
0.298	0.00506239\\
0.3	0.00525806\\
0.302	0.00543471\\
0.304	0.00559233\\
0.306	0.00573096\\
0.308	0.00585067\\
0.31	0.00595162\\
0.312	0.00603397\\
0.314	0.00609797\\
0.316	0.00614387\\
0.318	0.00617201\\
0.32	0.00618273\\
0.322	0.00617643\\
0.324	0.00615354\\
0.326	0.00611452\\
0.328	0.00605988\\
0.33	0.00599015\\
0.332	0.00590587\\
0.334	0.00580763\\
0.336	0.00569604\\
0.338	0.00557173\\
0.34	0.00543534\\
0.342	0.00528753\\
0.344	0.00521114\\
0.346	0.00513777\\
0.348	0.0050519\\
0.35	0.0049787\\
0.352	0.00499866\\
0.354	0.00500311\\
0.356	0.00499234\\
0.358	0.0049667\\
0.36	0.00492655\\
0.362	0.00487231\\
0.364	0.00480442\\
0.366	0.00472338\\
0.368	0.00462969\\
0.37	0.0045239\\
0.372	0.00440658\\
0.374	0.00427832\\
0.376	0.00413975\\
0.378	0.00399151\\
0.38	0.00383425\\
0.382	0.00366864\\
0.384	0.00349538\\
0.386	0.00333429\\
0.388	0.00320821\\
0.39	0.00307629\\
0.392	0.002939\\
0.394	0.0027968\\
0.396	0.00265017\\
0.398	0.00249959\\
0.4	0.00234551\\
0.402	0.00218843\\
0.404	0.0020288\\
0.406	0.00186711\\
0.408	0.00170382\\
0.41	0.00168162\\
0.412	0.00182244\\
0.414	0.00195489\\
0.416	0.00207883\\
0.418	0.0021941\\
0.42	0.00230062\\
0.422	0.0023983\\
0.424	0.0024871\\
0.426	0.00256698\\
0.428	0.00263795\\
0.43	0.00270004\\
0.432	0.00275327\\
0.434	0.00279774\\
0.436	0.00283352\\
0.438	0.00286074\\
0.44	0.00287952\\
0.442	0.00289002\\
0.444	0.00289241\\
0.446	0.00288689\\
0.448	0.00287365\\
0.45	0.00285293\\
0.452	0.00282496\\
0.454	0.00278999\\
0.456	0.00274829\\
0.458	0.00270014\\
0.46	0.00264582\\
0.462	0.00258564\\
0.464	0.00257913\\
0.466	0.00259249\\
0.468	0.0025975\\
0.47	0.00259434\\
0.472	0.00258323\\
0.474	0.00256439\\
0.476	0.00253807\\
0.478	0.00250454\\
0.48	0.00246409\\
0.482	0.00241701\\
0.484	0.00236362\\
0.486	0.00230425\\
0.488	0.00223924\\
0.49	0.00219639\\
0.492	0.00217393\\
0.494	0.00214616\\
0.496	0.00211326\\
0.498	0.00207541\\
0.5	0.00203282\\
0.502	0.00198569\\
0.504	0.00193424\\
0.506	0.00187869\\
0.508	0.00181927\\
0.51	0.00175622\\
0.512	0.00168978\\
0.514	0.00162018\\
0.516	0.00154769\\
0.518	0.00147255\\
0.52	0.00139501\\
0.522	0.00131532\\
0.524	0.00123375\\
0.526	0.00115054\\
0.528	0.00106596\\
0.53	0.000980245\\
0.532	0.000893657\\
0.534	0.000813537\\
0.536	0.000878256\\
0.538	0.000939057\\
0.54	0.000995872\\
0.542	0.00104865\\
0.544	0.00109734\\
0.546	0.00114193\\
0.548	0.00118238\\
0.55	0.0012187\\
0.552	0.00125088\\
0.554	0.00127895\\
0.556	0.00130293\\
0.558	0.00132285\\
0.56	0.00133876\\
0.562	0.00135073\\
0.564	0.0013588\\
0.566	0.00136306\\
0.568	0.00136359\\
0.57	0.00136047\\
0.572	0.00135381\\
0.574	0.00134371\\
0.576	0.00133027\\
0.578	0.00131362\\
0.58	0.00129388\\
0.582	0.00127117\\
0.584	0.00124563\\
0.586	0.00123686\\
0.588	0.00122672\\
0.59	0.00121357\\
0.592	0.00119753\\
0.594	0.0011787\\
0.596	0.0011572\\
0.598	0.00114646\\
0.6	0.00115852\\
0.602	0.00116718\\
0.604	0.00117248\\
0.606	0.00117448\\
0.608	0.00117323\\
0.61	0.00116881\\
0.612	0.00116128\\
0.614	0.00115074\\
0.616	0.00113727\\
0.618	0.00112097\\
0.62	0.00110193\\
0.622	0.00108028\\
0.624	0.00105611\\
0.626	0.00102955\\
0.628	0.00100071\\
0.63	0.000969729\\
0.632	0.000936727\\
0.634	0.000901837\\
0.636	0.000865195\\
0.638	0.000826937\\
0.64	0.000787201\\
0.642	0.000746126\\
0.644	0.000703854\\
0.646	0.000660526\\
0.648	0.000616284\\
0.65	0.000571268\\
0.652	0.000525618\\
0.654	0.000479476\\
0.656	0.000432977\\
0.658	0.00038626\\
0.66	0.00041396\\
0.662	0.00044355\\
0.664	0.000471307\\
0.666	0.0004972\\
0.668	0.000521204\\
0.67	0.0005433\\
0.672	0.000563474\\
0.674	0.000581717\\
0.676	0.000598024\\
0.678	0.000612397\\
0.68	0.000624843\\
0.682	0.000635373\\
0.684	0.000644001\\
0.686	0.000650748\\
0.688	0.000655639\\
0.69	0.000658703\\
0.692	0.000659971\\
0.694	0.000659481\\
0.696	0.000657274\\
0.698	0.000653392\\
0.7	0.000647883\\
0.702	0.000640798\\
0.704	0.000632189\\
0.706	0.000622113\\
0.708	0.000613472\\
0.71	0.000610958\\
0.712	0.000606892\\
0.714	0.000601318\\
0.716	0.000594283\\
0.718	0.00058992\\
0.72	0.000597388\\
0.722	0.000603018\\
0.724	0.000606832\\
0.726	0.000608859\\
0.728	0.000609132\\
0.73	0.000607688\\
0.732	0.000604568\\
0.734	0.000599817\\
0.736	0.000593483\\
0.738	0.000585621\\
0.74	0.000576284\\
0.742	0.000565531\\
0.744	0.000553425\\
0.746	0.000540028\\
0.748	0.000525408\\
0.75	0.000509632\\
0.752	0.000492771\\
0.754	0.000474897\\
0.756	0.000456083\\
0.758	0.000436403\\
0.76	0.000415933\\
0.762	0.000394747\\
0.764	0.000372924\\
0.766	0.000350538\\
0.768	0.000327666\\
0.77	0.000304385\\
0.772	0.00028077\\
0.774	0.000256896\\
0.776	0.000232837\\
0.778	0.000208665\\
0.78	0.000184453\\
0.782	0.000186862\\
0.784	0.000203044\\
0.786	0.000218344\\
0.788	0.000232743\\
0.79	0.00024622\\
0.792	0.000258759\\
0.794	0.000270348\\
0.796	0.000280974\\
0.798	0.000290631\\
0.8	0.000299312\\
0.802	0.000307016\\
0.804	0.000313742\\
0.806	0.000319492\\
0.808	0.000324272\\
0.81	0.000328089\\
0.812	0.000330952\\
0.814	0.000332874\\
0.816	0.00033387\\
0.818	0.000333954\\
0.82	0.000333147\\
0.822	0.000331467\\
0.824	0.000328939\\
0.826	0.000325585\\
0.828	0.000321431\\
0.83	0.000316505\\
0.832	0.000313492\\
0.834	0.000312233\\
0.836	0.000310158\\
0.838	0.000307289\\
0.84	0.000303651\\
0.842	0.000299269\\
0.844	0.00030103\\
0.846	0.000302016\\
0.848	0.000302119\\
0.85	0.000301358\\
0.852	0.000299758\\
0.854	0.000297343\\
0.856	0.000294139\\
0.858	0.000290175\\
0.86	0.000285481\\
0.862	0.000280087\\
0.864	0.000274027\\
0.866	0.000267335\\
0.868	0.000260044\\
0.87	0.000252191\\
0.872	0.000243813\\
0.874	0.000234945\\
0.876	0.000225627\\
0.878	0.000215895\\
0.88	0.000205789\\
0.882	0.000195346\\
0.884	0.000184606\\
0.886	0.000173607\\
0.888	0.000162387\\
0.89	0.000150983\\
0.892	0.000139435\\
0.894	0.000127777\\
0.896	0.000116048\\
0.898	0.000104282\\
0.9	9.25149e-05\\
0.902	8.23327e-05\\
0.904	9.15408e-05\\
0.906	0.000100305\\
0.908	0.00010861\\
0.91	0.000116442\\
0.912	0.00012379\\
0.914	0.000130644\\
0.916	0.000136995\\
0.918	0.000142836\\
0.92	0.00014816\\
0.922	0.000152965\\
0.924	0.000157248\\
0.926	0.000161007\\
0.928	0.000164243\\
0.93	0.000166958\\
0.932	0.000169155\\
0.934	0.000170839\\
0.936	0.000172016\\
0.938	0.000172692\\
0.94	0.000172877\\
0.942	0.00017258\\
0.944	0.000171812\\
0.946	0.000170585\\
0.948	0.000168912\\
0.95	0.000166807\\
0.952	0.000164285\\
0.954	0.000161361\\
0.956	0.000160496\\
0.958	0.000159372\\
0.96	0.000157828\\
0.962	0.000155878\\
0.964	0.000153536\\
0.966	0.000150816\\
0.968	0.000147736\\
0.97	0.000144311\\
0.972	0.000140559\\
0.974	0.000138057\\
0.976	0.000136644\\
0.978	0.000134886\\
0.98	0.000132798\\
0.982	0.000130395\\
0.984	0.000127693\\
0.986	0.000124707\\
0.988	0.000121453\\
0.99	0.000117949\\
0.992	0.000114211\\
0.994	0.000110256\\
0.996	0.000106101\\
0.998	0.000101765\\
1	9.72627e-05\\
1.002	9.43493e-05\\
1.004	9.31151e-05\\
1.006	9.16361e-05\\
1.008	8.99169e-05\\
1.01	8.79627e-05\\
1.012	8.57796e-05\\
1.014	8.33743e-05\\
1.016	8.07541e-05\\
1.018	7.79272e-05\\
1.02	7.49023e-05\\
1.022	7.16889e-05\\
1.024	6.82968e-05\\
1.026	6.47367e-05\\
1.028	6.10196e-05\\
1.03	5.7157e-05\\
1.032	6.07762e-05\\
1.034	6.44344e-05\\
1.036	6.78276e-05\\
1.038	7.09522e-05\\
1.04	7.38052e-05\\
1.042	7.63847e-05\\
1.044	7.86893e-05\\
1.046	8.07189e-05\\
1.048	8.24737e-05\\
1.05	8.3955e-05\\
1.052	8.51648e-05\\
1.054	8.61058e-05\\
1.056	8.67814e-05\\
1.058	8.71957e-05\\
1.06	8.73536e-05\\
1.062	8.72605e-05\\
1.064	8.69224e-05\\
1.066	8.6346e-05\\
1.068	8.55383e-05\\
1.07	8.4507e-05\\
1.072	8.32604e-05\\
1.074	8.18068e-05\\
1.076	8.01554e-05\\
1.078	7.92044e-05\\
1.08	7.8473e-05\\
1.082	7.75375e-05\\
1.084	7.64055e-05\\
1.086	7.50849e-05\\
1.088	7.35842e-05\\
1.09	7.19118e-05\\
1.092	7.00769e-05\\
1.094	6.80885e-05\\
1.096	6.67659e-05\\
1.098	6.81372e-05\\
1.1	6.92826e-05\\
1.102	7.02008e-05\\
1.104	7.08912e-05\\
1.106	7.13541e-05\\
1.108	7.15905e-05\\
1.11	7.16019e-05\\
1.112	7.13908e-05\\
1.114	7.09604e-05\\
1.116	7.03146e-05\\
1.118	6.94579e-05\\
1.12	6.83954e-05\\
1.122	6.71331e-05\\
1.124	6.56773e-05\\
1.126	6.40351e-05\\
1.128	6.22141e-05\\
1.13	6.02223e-05\\
1.132	5.80683e-05\\
1.134	5.57612e-05\\
1.136	5.33103e-05\\
1.138	5.07256e-05\\
1.14	4.80172e-05\\
1.142	4.51955e-05\\
1.144	4.22713e-05\\
1.146	3.92554e-05\\
1.148	3.6159e-05\\
1.15	3.29934e-05\\
1.152	2.97698e-05\\
1.154	3.04961e-05\\
1.156	3.21467e-05\\
1.158	3.3665e-05\\
1.16	3.50499e-05\\
1.162	3.63009e-05\\
1.164	3.74178e-05\\
1.166	3.84007e-05\\
1.168	3.92502e-05\\
1.17	3.99674e-05\\
1.172	4.05535e-05\\
1.174	4.10103e-05\\
1.176	4.13396e-05\\
1.178	4.1544e-05\\
1.18	4.16259e-05\\
1.182	4.15884e-05\\
1.184	4.14345e-05\\
1.186	4.11679e-05\\
1.188	4.07922e-05\\
1.19	4.03113e-05\\
1.192	3.97292e-05\\
1.194	3.90504e-05\\
1.196	3.82794e-05\\
1.198	3.74206e-05\\
1.2	3.75817e-05\\
1.202	3.90518e-05\\
1.204	4.03676e-05\\
1.206	4.15272e-05\\
1.208	4.25294e-05\\
1.21	4.33733e-05\\
1.212	4.4059e-05\\
1.214	4.45867e-05\\
1.216	4.49575e-05\\
1.218	4.51727e-05\\
1.22	4.52345e-05\\
1.222	4.51453e-05\\
1.224	4.4908e-05\\
1.226	4.45261e-05\\
1.228	4.40036e-05\\
1.23	4.33447e-05\\
1.232	4.25541e-05\\
1.234	4.16371e-05\\
1.236	4.05989e-05\\
1.238	3.94454e-05\\
1.24	3.81828e-05\\
1.242	3.68173e-05\\
1.244	3.53556e-05\\
1.246	3.38044e-05\\
1.248	3.2171e-05\\
1.25	3.04623e-05\\
1.252	2.86858e-05\\
1.254	2.68489e-05\\
1.256	2.4959e-05\\
1.258	2.30238e-05\\
1.26	2.10507e-05\\
1.262	1.90475e-05\\
1.264	1.70215e-05\\
1.266	1.49803e-05\\
1.268	1.37913e-05\\
1.27	1.28486e-05\\
1.272	1.2024e-05\\
1.274	1.28615e-05\\
1.276	1.36414e-05\\
1.278	1.43633e-05\\
1.28	1.50266e-05\\
1.282	1.56314e-05\\
1.284	1.61774e-05\\
1.286	1.66649e-05\\
1.288	1.70942e-05\\
1.29	1.74657e-05\\
1.292	1.77801e-05\\
1.294	1.8038e-05\\
1.296	1.82404e-05\\
1.298	1.83883e-05\\
1.3	1.84828e-05\\
1.302	1.85252e-05\\
1.304	1.85168e-05\\
1.306	1.90737e-05\\
1.308	2.00544e-05\\
1.31	2.09433e-05\\
1.312	2.17393e-05\\
1.314	2.24417e-05\\
1.316	2.30502e-05\\
1.318	2.35646e-05\\
1.32	2.39854e-05\\
1.322	2.43131e-05\\
1.324	2.45487e-05\\
1.326	2.46935e-05\\
1.328	2.4749e-05\\
1.33	2.47172e-05\\
1.332	2.46002e-05\\
1.334	2.44003e-05\\
1.336	2.41204e-05\\
1.338	2.37632e-05\\
1.34	2.33318e-05\\
1.342	2.28297e-05\\
1.344	2.22602e-05\\
1.346	2.1627e-05\\
1.348	2.09341e-05\\
1.35	2.01852e-05\\
1.352	1.93845e-05\\
1.354	1.85361e-05\\
1.356	1.76442e-05\\
1.358	1.67132e-05\\
1.36	1.57474e-05\\
1.362	1.47512e-05\\
1.364	1.38665e-05\\
1.366	1.36702e-05\\
1.368	1.34422e-05\\
1.37	1.31835e-05\\
1.372	1.28952e-05\\
1.374	1.25784e-05\\
1.376	1.22343e-05\\
1.378	1.18642e-05\\
1.38	1.14693e-05\\
1.382	1.10509e-05\\
1.384	1.06106e-05\\
1.386	1.01497e-05\\
1.388	9.66961e-06\\
1.39	9.17193e-06\\
1.392	8.65814e-06\\
1.394	8.12981e-06\\
1.396	7.58849e-06\\
1.398	7.03579e-06\\
1.4	6.47329e-06\\
1.402	5.90261e-06\\
1.404	6.19608e-06\\
1.406	6.72248e-06\\
1.408	7.37696e-06\\
1.41	7.98567e-06\\
1.412	8.5476e-06\\
1.414	9.06196e-06\\
1.416	9.52816e-06\\
1.418	9.94582e-06\\
1.42	1.03147e-05\\
1.422	1.06349e-05\\
1.424	1.09066e-05\\
1.426	1.113e-05\\
1.428	1.13059e-05\\
1.43	1.14348e-05\\
1.432	1.15177e-05\\
1.434	1.15556e-05\\
1.436	1.15496e-05\\
1.438	1.15012e-05\\
1.44	1.14117e-05\\
1.442	1.12827e-05\\
1.444	1.11159e-05\\
1.446	1.0913e-05\\
1.448	1.0676e-05\\
1.45	1.04068e-05\\
1.452	1.01073e-05\\
1.454	9.77977e-06\\
1.456	9.42622e-06\\
1.458	9.04886e-06\\
1.46	8.6499e-06\\
1.462	8.23156e-06\\
1.464	8.02203e-06\\
1.466	8.16279e-06\\
1.468	8.27829e-06\\
1.47	8.36861e-06\\
1.472	8.43387e-06\\
1.474	8.47426e-06\\
1.476	8.49006e-06\\
1.478	8.4816e-06\\
1.48	8.44928e-06\\
1.482	8.39357e-06\\
1.484	8.31498e-06\\
1.486	8.21411e-06\\
1.488	8.0916e-06\\
1.49	7.94813e-06\\
1.492	7.78446e-06\\
1.494	7.60137e-06\\
1.496	7.39971e-06\\
1.498	7.18035e-06\\
1.5	6.94423e-06\\
1.502	6.6923e-06\\
1.504	6.42556e-06\\
1.506	6.14501e-06\\
1.508	5.85171e-06\\
1.51	5.54673e-06\\
1.512	5.23114e-06\\
1.514	4.90606e-06\\
1.516	4.57259e-06\\
1.518	4.23186e-06\\
1.52	3.88499e-06\\
1.522	3.65486e-06\\
1.524	3.80647e-06\\
1.526	3.9363e-06\\
1.528	4.04462e-06\\
1.53	4.13174e-06\\
1.532	4.19809e-06\\
1.534	4.24417e-06\\
1.536	4.27053e-06\\
1.538	4.27781e-06\\
1.54	4.2667e-06\\
1.542	4.23795e-06\\
1.544	4.19238e-06\\
1.546	4.13084e-06\\
1.548	4.05421e-06\\
1.55	3.96344e-06\\
1.552	3.96738e-06\\
1.554	4.01619e-06\\
1.556	4.05428e-06\\
1.558	4.08174e-06\\
1.56	4.09867e-06\\
1.562	4.10523e-06\\
1.564	4.10155e-06\\
1.566	4.08782e-06\\
1.568	4.06423e-06\\
1.57	4.03099e-06\\
1.572	3.99674e-06\\
1.574	4.10295e-06\\
1.576	4.2427e-06\\
1.578	4.36693e-06\\
1.58	4.47554e-06\\
1.582	4.56849e-06\\
1.584	4.64581e-06\\
1.586	4.70756e-06\\
1.588	4.75386e-06\\
1.59	4.78486e-06\\
1.592	4.80076e-06\\
1.594	4.80181e-06\\
1.596	4.78831e-06\\
1.598	4.76059e-06\\
1.6	4.719e-06\\
1.602	4.66396e-06\\
1.604	4.59591e-06\\
1.606	4.51532e-06\\
1.608	4.42269e-06\\
1.61	4.31855e-06\\
1.612	4.20347e-06\\
1.614	4.07802e-06\\
1.616	3.94281e-06\\
1.618	3.79846e-06\\
1.62	3.6456e-06\\
1.622	3.4849e-06\\
1.624	3.31702e-06\\
1.626	3.14263e-06\\
1.628	2.96241e-06\\
1.63	2.77706e-06\\
1.632	2.58726e-06\\
1.634	2.39369e-06\\
1.636	2.19706e-06\\
1.638	1.99804e-06\\
1.64	1.79731e-06\\
1.642	1.59553e-06\\
1.644	1.39337e-06\\
1.646	1.19147e-06\\
1.648	1.08773e-06\\
1.65	1.20962e-06\\
1.652	1.32625e-06\\
1.654	1.43738e-06\\
1.656	1.54281e-06\\
1.658	1.64234e-06\\
1.66	1.73579e-06\\
1.662	1.823e-06\\
1.664	1.90385e-06\\
1.666	1.97819e-06\\
1.668	2.04595e-06\\
1.67	2.10703e-06\\
1.672	2.16137e-06\\
1.674	2.20892e-06\\
1.676	2.24967e-06\\
1.678	2.2836e-06\\
1.68	2.31073e-06\\
1.682	2.33108e-06\\
1.684	2.34472e-06\\
1.686	2.3517e-06\\
1.688	2.35211e-06\\
1.69	2.34604e-06\\
1.692	2.33363e-06\\
1.694	2.31499e-06\\
1.696	2.31869e-06\\
1.698	2.36143e-06\\
1.7	2.39572e-06\\
1.702	2.42165e-06\\
1.704	2.43935e-06\\
1.706	2.44896e-06\\
1.708	2.45065e-06\\
1.71	2.44461e-06\\
1.712	2.43105e-06\\
1.714	2.41021e-06\\
1.716	2.38233e-06\\
1.718	2.34769e-06\\
1.72	2.30656e-06\\
1.722	2.25925e-06\\
1.724	2.20606e-06\\
1.726	2.14731e-06\\
1.728	2.08334e-06\\
1.73	2.0145e-06\\
1.732	1.94113e-06\\
1.734	1.86358e-06\\
1.736	1.78222e-06\\
1.738	1.69741e-06\\
1.74	1.60951e-06\\
1.742	1.5189e-06\\
1.744	1.42595e-06\\
1.746	1.331e-06\\
1.748	1.23444e-06\\
1.75	1.13661e-06\\
1.752	1.03787e-06\\
1.754	9.38558e-07\\
1.756	8.39024e-07\\
1.758	7.83626e-07\\
1.76	7.60382e-07\\
1.762	7.35424e-07\\
1.764	7.58649e-07\\
1.766	7.93464e-07\\
1.768	8.2506e-07\\
1.77	8.53252e-07\\
1.772	8.77878e-07\\
1.774	8.98793e-07\\
1.776	9.1587e-07\\
1.778	9.3465e-07\\
1.78	9.89634e-07\\
1.782	1.04067e-06\\
1.784	1.08768e-06\\
1.786	1.13059e-06\\
1.788	1.16936e-06\\
1.79	1.20395e-06\\
1.792	1.23433e-06\\
1.794	1.26048e-06\\
1.796	1.28242e-06\\
1.798	1.30015e-06\\
1.8	1.3137e-06\\
1.802	1.32311e-06\\
1.804	1.32842e-06\\
1.806	1.32971e-06\\
1.808	1.32704e-06\\
1.81	1.3205e-06\\
1.812	1.31018e-06\\
1.814	1.29619e-06\\
1.816	1.27863e-06\\
1.818	1.25764e-06\\
1.82	1.23334e-06\\
1.822	1.22214e-06\\
1.824	1.20967e-06\\
1.826	1.19387e-06\\
1.828	1.17486e-06\\
1.83	1.15276e-06\\
1.832	1.12771e-06\\
1.834	1.09984e-06\\
1.836	1.06931e-06\\
1.838	1.03626e-06\\
1.84	1.00084e-06\\
1.842	9.63219e-07\\
1.844	9.23555e-07\\
1.846	8.82016e-07\\
1.848	8.38767e-07\\
1.85	7.9398e-07\\
1.852	7.49866e-07\\
1.854	7.1292e-07\\
1.856	7.45255e-07\\
1.858	7.88707e-07\\
1.86	8.2907e-07\\
1.862	8.66191e-07\\
1.864	8.9993e-07\\
1.866	9.30163e-07\\
1.868	9.56782e-07\\
1.87	9.79695e-07\\
1.872	9.98825e-07\\
1.874	1.01411e-06\\
1.876	1.02551e-06\\
1.878	1.033e-06\\
1.88	1.03657e-06\\
1.882	1.03622e-06\\
1.884	1.03197e-06\\
1.886	1.02387e-06\\
1.888	1.01197e-06\\
1.89	9.9634e-07\\
1.892	9.7706e-07\\
1.894	9.54232e-07\\
1.896	9.27969e-07\\
1.898	8.98395e-07\\
1.9	8.65651e-07\\
1.902	8.29887e-07\\
1.904	7.91263e-07\\
1.906	7.49952e-07\\
1.908	7.06136e-07\\
1.91	6.60004e-07\\
1.912	6.66218e-07\\
1.914	6.75365e-07\\
1.916	6.82224e-07\\
1.918	6.86825e-07\\
1.92	6.89209e-07\\
1.922	6.89421e-07\\
1.924	6.8751e-07\\
1.926	6.83535e-07\\
1.928	6.77556e-07\\
1.93	6.69642e-07\\
1.932	6.59863e-07\\
1.934	6.48296e-07\\
1.936	6.3502e-07\\
1.938	6.20118e-07\\
1.94	6.03679e-07\\
1.942	5.92263e-07\\
1.944	5.82625e-07\\
1.946	5.71465e-07\\
1.948	5.58859e-07\\
1.95	5.44882e-07\\
1.952	5.29617e-07\\
1.954	5.13143e-07\\
1.956	5.47881e-07\\
1.958	5.86182e-07\\
1.96	6.21889e-07\\
1.962	6.54894e-07\\
1.964	6.85102e-07\\
1.966	7.12428e-07\\
1.968	7.368e-07\\
1.97	7.58159e-07\\
1.972	7.76459e-07\\
1.974	7.91665e-07\\
1.976	8.03756e-07\\
1.978	8.12724e-07\\
1.98	8.18572e-07\\
1.982	8.21316e-07\\
1.984	8.20984e-07\\
1.986	8.17616e-07\\
1.988	8.11262e-07\\
1.99	8.01986e-07\\
1.992	7.89859e-07\\
1.994	7.74966e-07\\
1.996	7.574e-07\\
1.998	7.37262e-07\\
2	7.14664e-07\\
};
\addplot [color=mycolor3, line width=1.0pt, forget plot]
  table[row sep=crcr]{%
-0.024	0\\
-0.022	0\\
-0.02	0\\
-0.018	0\\
-0.016	0\\
-0.014	0\\
-0.012	0\\
-0.01	0\\
-0.008	0\\
-0.006	0\\
-0.004	0\\
-0.002	0\\
0	0\\
0.002	0.00916185\\
0.004	0.0147068\\
0.006	0.0198659\\
0.008	0.0237562\\
0.01	0.0264985\\
0.012	0.0282387\\
0.014	0.0291273\\
0.016	0.0293138\\
0.018	0.0289573\\
0.02	0.0291472\\
0.022	0.0292724\\
0.024	0.029371\\
0.026	0.0295666\\
0.028	0.0301154\\
0.03	0.0303618\\
0.032	0.0303135\\
0.034	0.0301565\\
0.036	0.031263\\
0.038	0.0322763\\
0.04	0.0331967\\
0.042	0.0340247\\
0.044	0.0347608\\
0.046	0.0354059\\
0.048	0.0359612\\
0.05	0.0364278\\
0.052	0.036807\\
0.054	0.0371004\\
0.056	0.0373097\\
0.058	0.0374364\\
0.06	0.0374827\\
0.062	0.0374504\\
0.064	0.0373416\\
0.066	0.0371587\\
0.068	0.036904\\
0.07	0.0365798\\
0.072	0.0361889\\
0.074	0.0357337\\
0.076	0.0352172\\
0.078	0.0346421\\
0.08	0.0340113\\
0.082	0.0333279\\
0.084	0.0325949\\
0.086	0.0318155\\
0.088	0.0309927\\
0.09	0.03013\\
0.092	0.0292304\\
0.094	0.0282972\\
0.096	0.0273338\\
0.098	0.0263435\\
0.1	0.0253293\\
0.102	0.0242947\\
0.104	0.0232429\\
0.106	0.0221769\\
0.108	0.0211\\
0.11	0.0200152\\
0.112	0.0189254\\
0.114	0.0178337\\
0.116	0.0172939\\
0.118	0.016791\\
0.12	0.0162658\\
0.122	0.0157204\\
0.124	0.0151568\\
0.126	0.014577\\
0.128	0.0139831\\
0.13	0.0135489\\
0.132	0.0133412\\
0.134	0.0131111\\
0.136	0.0128595\\
0.138	0.0125877\\
0.14	0.0122968\\
0.142	0.0119882\\
0.144	0.0116631\\
0.146	0.0113228\\
0.148	0.010996\\
0.15	0.0108305\\
0.152	0.0106576\\
0.154	0.0104778\\
0.156	0.0102915\\
0.158	0.0100993\\
0.16	0.00990147\\
0.162	0.00969869\\
0.164	0.00949143\\
0.166	0.0092802\\
0.168	0.00906551\\
0.17	0.00884789\\
0.172	0.00862786\\
0.174	0.0084059\\
0.176	0.00818253\\
0.178	0.00795823\\
0.18	0.00773346\\
0.182	0.00750869\\
0.184	0.00728436\\
0.186	0.00706088\\
0.188	0.00706475\\
0.19	0.00714501\\
0.192	0.00719797\\
0.194	0.00722464\\
0.196	0.0072261\\
0.198	0.00720343\\
0.2	0.00715773\\
0.202	0.00709015\\
0.204	0.00700185\\
0.206	0.00689398\\
0.208	0.00676773\\
0.21	0.00662428\\
0.212	0.00646482\\
0.214	0.00629051\\
0.216	0.00610254\\
0.218	0.00590208\\
0.22	0.00569028\\
0.222	0.00546828\\
0.224	0.00523719\\
0.226	0.00499812\\
0.228	0.00475214\\
0.23	0.00450031\\
0.232	0.00424365\\
0.234	0.00415427\\
0.236	0.00405491\\
0.238	0.00394551\\
0.24	0.00382679\\
0.242	0.00369946\\
0.244	0.00356424\\
0.246	0.00342185\\
0.248	0.003273\\
0.25	0.00311838\\
0.252	0.00303628\\
0.254	0.00301728\\
0.256	0.0029889\\
0.258	0.00295158\\
0.26	0.00290576\\
0.262	0.0028519\\
0.264	0.00279046\\
0.266	0.00272191\\
0.268	0.00264672\\
0.27	0.00256537\\
0.272	0.00247832\\
0.274	0.00238605\\
0.276	0.00228904\\
0.278	0.00218775\\
0.28	0.00208265\\
0.282	0.0019742\\
0.284	0.00186285\\
0.286	0.00174906\\
0.288	0.00183794\\
0.29	0.00194779\\
0.292	0.00204832\\
0.294	0.00213958\\
0.296	0.00222166\\
0.298	0.00229465\\
0.3	0.00235867\\
0.302	0.00241387\\
0.304	0.00246043\\
0.306	0.00249851\\
0.308	0.00252833\\
0.31	0.0025501\\
0.312	0.00256406\\
0.314	0.00257046\\
0.316	0.00256955\\
0.318	0.00256161\\
0.32	0.00254692\\
0.322	0.00252577\\
0.324	0.00249846\\
0.326	0.00246529\\
0.328	0.00242657\\
0.33	0.00238262\\
0.332	0.00233376\\
0.334	0.0022803\\
0.336	0.00222256\\
0.338	0.00216087\\
0.34	0.00209554\\
0.342	0.00202689\\
0.344	0.00195525\\
0.346	0.00188091\\
0.348	0.00183707\\
0.35	0.00180792\\
0.352	0.00177475\\
0.354	0.00173777\\
0.356	0.00169724\\
0.358	0.00165337\\
0.36	0.0016064\\
0.362	0.00155657\\
0.364	0.00150412\\
0.366	0.00144928\\
0.368	0.00139228\\
0.37	0.00133337\\
0.372	0.00129071\\
0.374	0.00127849\\
0.376	0.00126304\\
0.378	0.00124448\\
0.38	0.00122298\\
0.382	0.00119869\\
0.384	0.00117177\\
0.386	0.00114237\\
0.388	0.00111067\\
0.39	0.00107683\\
0.392	0.001041\\
0.394	0.00100337\\
0.396	0.000964102\\
0.398	0.000923354\\
0.4	0.000881296\\
0.402	0.000838091\\
0.404	0.000793904\\
0.406	0.000748894\\
0.408	0.000703218\\
0.41	0.000657032\\
0.412	0.000610486\\
0.414	0.000563727\\
0.416	0.000548735\\
0.418	0.000579836\\
0.42	0.000608096\\
0.422	0.000633539\\
0.424	0.000656197\\
0.426	0.00067611\\
0.428	0.000693321\\
0.43	0.00070788\\
0.432	0.000719845\\
0.434	0.000729275\\
0.436	0.000736236\\
0.438	0.0007408\\
0.44	0.00074304\\
0.442	0.000743034\\
0.444	0.000740865\\
0.446	0.000736617\\
0.448	0.000730379\\
0.45	0.000722239\\
0.452	0.000712291\\
0.454	0.000700629\\
0.456	0.000687348\\
0.458	0.000672545\\
0.46	0.000656317\\
0.462	0.000638764\\
0.464	0.00062355\\
0.466	0.000619787\\
0.468	0.00061434\\
0.47	0.000607286\\
0.472	0.000598698\\
0.474	0.000588657\\
0.476	0.000577242\\
0.478	0.000564533\\
0.48	0.000550613\\
0.482	0.000535565\\
0.484	0.000519472\\
0.486	0.000502418\\
0.488	0.000484486\\
0.49	0.00046576\\
0.492	0.000454098\\
0.494	0.000450847\\
0.496	0.000446394\\
0.498	0.000440793\\
0.5	0.000434101\\
0.502	0.000426373\\
0.504	0.000417669\\
0.506	0.000408049\\
0.508	0.000397571\\
0.51	0.000386295\\
0.512	0.000374283\\
0.514	0.000361595\\
0.516	0.000348291\\
0.518	0.000334431\\
0.52	0.000320074\\
0.522	0.000305279\\
0.524	0.000290105\\
0.526	0.000274608\\
0.528	0.000258844\\
0.53	0.000242868\\
0.532	0.000226734\\
0.534	0.000210493\\
0.536	0.000194197\\
0.538	0.000177894\\
0.54	0.000185716\\
0.542	0.000195061\\
0.544	0.0002035\\
0.546	0.000211042\\
0.548	0.000217698\\
0.55	0.000223481\\
0.552	0.000228406\\
0.554	0.000232489\\
0.556	0.000235751\\
0.558	0.000238212\\
0.56	0.000239893\\
0.562	0.000240819\\
0.564	0.000241015\\
0.566	0.000240507\\
0.568	0.000239322\\
0.57	0.000237488\\
0.572	0.000235036\\
0.574	0.000231993\\
0.576	0.000228392\\
0.578	0.000224262\\
0.58	0.000219635\\
0.582	0.000214543\\
0.584	0.000213514\\
0.586	0.000212865\\
0.588	0.00021163\\
0.59	0.000209832\\
0.592	0.000207498\\
0.594	0.000204651\\
0.596	0.000201319\\
0.598	0.000197528\\
0.6	0.000193305\\
0.602	0.000188677\\
0.604	0.000183672\\
0.606	0.000178316\\
0.608	0.000172639\\
0.61	0.000166666\\
0.612	0.000160427\\
0.614	0.000159528\\
0.616	0.000158503\\
0.618	0.00015706\\
0.62	0.000155219\\
0.622	0.000152999\\
0.624	0.000150418\\
0.626	0.000147496\\
0.628	0.000144253\\
0.63	0.00014071\\
0.632	0.000136886\\
0.634	0.000132802\\
0.636	0.00012848\\
0.638	0.000123938\\
0.64	0.000119198\\
0.642	0.00011428\\
0.644	0.000109205\\
0.646	0.000103991\\
0.648	9.86599e-05\\
0.65	9.32297e-05\\
0.652	8.77196e-05\\
0.654	8.21484e-05\\
0.656	7.6534e-05\\
0.658	7.08941e-05\\
0.66	6.52457e-05\\
0.662	5.96054e-05\\
0.664	6.12184e-05\\
0.666	6.41333e-05\\
0.668	6.67601e-05\\
0.67	6.91017e-05\\
0.672	7.11614e-05\\
0.674	7.29432e-05\\
0.676	7.4452e-05\\
0.678	7.56929e-05\\
0.68	7.6672e-05\\
0.682	7.73956e-05\\
0.684	7.78708e-05\\
0.686	7.81049e-05\\
0.688	7.81058e-05\\
0.69	7.78817e-05\\
0.692	7.74413e-05\\
0.694	7.67933e-05\\
0.696	7.59472e-05\\
0.698	7.49122e-05\\
0.7	7.3698e-05\\
0.702	7.23145e-05\\
0.704	7.11287e-05\\
0.706	7.12482e-05\\
0.708	7.11589e-05\\
0.71	7.08686e-05\\
0.712	7.03852e-05\\
0.714	6.97167e-05\\
0.716	6.88717e-05\\
0.718	6.78589e-05\\
0.72	6.6687e-05\\
0.722	6.53651e-05\\
0.724	6.39024e-05\\
0.726	6.23079e-05\\
0.728	6.05909e-05\\
0.73	5.87608e-05\\
0.732	5.68269e-05\\
0.734	5.47983e-05\\
0.736	5.42592e-05\\
0.738	5.40046e-05\\
0.74	5.36061e-05\\
0.742	5.30698e-05\\
0.744	5.24022e-05\\
0.746	5.16097e-05\\
0.748	5.0699e-05\\
0.75	4.96768e-05\\
0.752	4.85499e-05\\
0.754	4.73253e-05\\
0.756	4.60097e-05\\
0.758	4.46102e-05\\
0.76	4.31336e-05\\
0.762	4.15868e-05\\
0.764	3.99766e-05\\
0.766	3.83099e-05\\
0.768	3.65933e-05\\
0.77	3.48334e-05\\
0.772	3.30366e-05\\
0.774	3.12094e-05\\
0.776	2.93579e-05\\
0.778	2.74883e-05\\
0.78	2.56063e-05\\
0.782	2.37178e-05\\
0.784	2.18282e-05\\
0.786	1.9943e-05\\
0.788	2.04933e-05\\
0.79	2.14037e-05\\
0.792	2.22201e-05\\
0.794	2.29438e-05\\
0.796	2.35758e-05\\
0.798	2.41178e-05\\
0.8	2.45715e-05\\
0.802	2.49387e-05\\
0.804	2.52215e-05\\
0.806	2.54222e-05\\
0.808	2.55431e-05\\
0.81	2.55868e-05\\
0.812	2.55559e-05\\
0.814	2.54533e-05\\
0.816	2.52818e-05\\
0.818	2.50444e-05\\
0.82	2.47441e-05\\
0.822	2.43841e-05\\
0.824	2.39675e-05\\
0.826	2.34976e-05\\
0.828	2.34543e-05\\
0.83	2.3478e-05\\
0.832	2.34343e-05\\
0.834	2.33258e-05\\
0.836	2.31551e-05\\
0.838	2.29248e-05\\
0.84	2.26377e-05\\
0.842	2.22966e-05\\
0.844	2.19044e-05\\
0.846	2.1464e-05\\
0.848	2.09783e-05\\
0.85	2.04503e-05\\
0.852	1.98829e-05\\
0.854	1.92792e-05\\
0.856	1.86421e-05\\
0.858	1.81255e-05\\
0.86	1.80985e-05\\
0.862	1.80214e-05\\
0.864	1.78962e-05\\
0.866	1.77248e-05\\
0.868	1.75095e-05\\
0.87	1.72522e-05\\
0.872	1.69552e-05\\
0.874	1.66207e-05\\
0.876	1.62509e-05\\
0.878	1.58481e-05\\
0.88	1.54145e-05\\
0.882	1.49525e-05\\
0.884	1.44644e-05\\
0.886	1.39523e-05\\
0.888	1.34187e-05\\
0.89	1.28657e-05\\
0.892	1.22956e-05\\
0.894	1.17106e-05\\
0.896	1.11129e-05\\
0.898	1.05046e-05\\
0.9	9.88775e-06\\
0.902	9.2645e-06\\
0.904	8.6368e-06\\
0.906	8.00657e-06\\
0.908	7.37571e-06\\
0.91	6.74601e-06\\
0.912	6.73546e-06\\
0.914	7.03321e-06\\
0.916	7.3006e-06\\
0.918	7.53792e-06\\
0.92	7.74556e-06\\
0.922	7.92395e-06\\
0.924	8.07361e-06\\
0.926	8.19512e-06\\
0.928	8.28909e-06\\
0.93	8.35624e-06\\
0.932	8.39729e-06\\
0.934	8.41303e-06\\
0.936	8.40431e-06\\
0.938	8.372e-06\\
0.94	8.317e-06\\
0.942	8.24027e-06\\
0.944	8.14277e-06\\
0.946	8.0255e-06\\
0.948	7.88949e-06\\
0.95	7.73578e-06\\
0.952	7.74218e-06\\
0.954	7.74919e-06\\
0.956	7.73384e-06\\
0.958	7.69696e-06\\
0.96	7.63941e-06\\
0.962	7.56209e-06\\
0.964	7.46592e-06\\
0.966	7.35183e-06\\
0.968	7.22079e-06\\
0.97	7.07376e-06\\
0.972	6.91174e-06\\
0.974	6.73571e-06\\
0.976	6.54669e-06\\
0.978	6.34566e-06\\
0.98	6.13364e-06\\
0.982	5.9858e-06\\
0.984	5.97704e-06\\
0.986	5.95168e-06\\
0.988	5.91039e-06\\
0.99	5.85384e-06\\
0.992	5.78275e-06\\
0.994	5.69782e-06\\
0.996	5.59979e-06\\
0.998	5.48941e-06\\
1	5.36743e-06\\
1.002	5.2346e-06\\
1.004	5.09169e-06\\
1.006	4.93946e-06\\
1.008	4.77868e-06\\
1.01	4.61009e-06\\
1.012	4.43446e-06\\
1.014	4.25253e-06\\
1.016	4.06505e-06\\
1.018	3.87273e-06\\
1.02	3.6763e-06\\
1.022	3.47645e-06\\
1.024	3.27387e-06\\
1.026	3.06923e-06\\
1.028	2.86319e-06\\
1.03	2.65636e-06\\
1.032	2.44936e-06\\
1.034	2.24279e-06\\
1.036	2.23489e-06\\
1.038	2.33134e-06\\
1.04	2.41772e-06\\
1.042	2.49413e-06\\
1.044	2.56071e-06\\
1.046	2.61763e-06\\
1.048	2.66506e-06\\
1.05	2.70322e-06\\
1.052	2.73232e-06\\
1.054	2.75261e-06\\
1.056	2.76434e-06\\
1.058	2.76779e-06\\
1.06	2.76324e-06\\
1.062	2.75099e-06\\
1.064	2.73136e-06\\
1.066	2.70465e-06\\
1.068	2.67122e-06\\
1.07	2.63138e-06\\
1.072	2.58548e-06\\
1.074	2.53387e-06\\
1.076	2.53202e-06\\
1.078	2.53343e-06\\
1.08	2.5276e-06\\
1.082	2.51481e-06\\
1.084	2.49533e-06\\
1.086	2.46947e-06\\
1.088	2.43752e-06\\
1.09	2.39979e-06\\
1.092	2.35658e-06\\
1.094	2.30823e-06\\
1.096	2.25504e-06\\
1.098	2.19735e-06\\
1.1	2.13546e-06\\
1.102	2.06971e-06\\
1.104	2.00042e-06\\
1.106	1.95164e-06\\
1.108	1.94995e-06\\
1.11	1.94282e-06\\
1.112	1.93045e-06\\
1.114	1.91307e-06\\
1.116	1.89089e-06\\
1.118	1.86415e-06\\
1.12	1.83308e-06\\
1.122	1.79793e-06\\
1.124	1.75894e-06\\
1.126	1.71635e-06\\
1.128	1.67042e-06\\
1.13	1.62138e-06\\
1.132	1.5695e-06\\
1.134	1.51501e-06\\
1.136	1.45816e-06\\
1.138	1.39919e-06\\
1.14	1.33836e-06\\
1.142	1.27589e-06\\
1.144	1.21202e-06\\
1.146	1.14699e-06\\
1.148	1.08101e-06\\
1.15	1.01431e-06\\
1.152	9.47104e-07\\
1.154	8.79599e-07\\
1.156	8.11998e-07\\
1.158	7.44495e-07\\
1.16	7.28106e-07\\
1.162	7.59665e-07\\
1.164	7.87955e-07\\
1.166	8.1301e-07\\
1.168	8.34873e-07\\
1.17	8.53595e-07\\
1.172	8.69234e-07\\
1.174	8.81853e-07\\
1.176	8.91523e-07\\
1.178	8.98322e-07\\
1.18	9.02329e-07\\
1.182	9.03634e-07\\
1.184	9.02327e-07\\
1.186	8.98503e-07\\
1.188	8.92263e-07\\
1.19	8.83709e-07\\
1.192	8.72947e-07\\
1.194	8.60087e-07\\
1.196	8.45237e-07\\
1.198	8.28511e-07\\
1.2	8.26302e-07\\
1.202	8.26769e-07\\
1.204	8.24865e-07\\
1.206	8.2068e-07\\
1.208	8.14308e-07\\
1.21	8.05844e-07\\
1.212	7.95387e-07\\
1.214	7.83037e-07\\
1.216	7.68898e-07\\
1.218	7.53073e-07\\
1.22	7.35668e-07\\
1.222	7.16789e-07\\
1.224	6.96542e-07\\
1.226	6.75033e-07\\
1.228	6.52369e-07\\
1.23	6.33656e-07\\
1.232	6.333e-07\\
1.234	6.31163e-07\\
1.236	6.27316e-07\\
1.238	6.2183e-07\\
1.24	6.14779e-07\\
1.242	6.0624e-07\\
1.244	5.9629e-07\\
1.246	5.85008e-07\\
1.248	5.72473e-07\\
1.25	5.58765e-07\\
1.252	5.43966e-07\\
1.254	5.28155e-07\\
1.256	5.11415e-07\\
1.258	4.93825e-07\\
1.26	4.75466e-07\\
1.262	4.56416e-07\\
1.264	4.36755e-07\\
1.266	4.1656e-07\\
1.268	3.95907e-07\\
1.27	3.74871e-07\\
1.272	3.53525e-07\\
1.274	3.31939e-07\\
1.276	3.10185e-07\\
1.278	2.88328e-07\\
1.28	2.66435e-07\\
1.282	2.44568e-07\\
1.284	2.36921e-07\\
1.286	2.47191e-07\\
1.288	2.56393e-07\\
1.29	2.64538e-07\\
1.292	2.71642e-07\\
1.294	2.77722e-07\\
1.296	2.82797e-07\\
1.298	2.86888e-07\\
1.3	2.9002e-07\\
1.302	2.92217e-07\\
1.304	2.93507e-07\\
1.306	2.9392e-07\\
1.308	2.93484e-07\\
1.31	2.92232e-07\\
1.312	2.90196e-07\\
1.314	2.87411e-07\\
1.316	2.83911e-07\\
1.318	2.79731e-07\\
1.32	2.74908e-07\\
1.322	2.69478e-07\\
1.324	2.68356e-07\\
1.326	2.68524e-07\\
1.328	2.67924e-07\\
1.33	2.66587e-07\\
1.332	2.64541e-07\\
1.334	2.61818e-07\\
1.336	2.58449e-07\\
1.338	2.54468e-07\\
1.34	2.49906e-07\\
1.342	2.44798e-07\\
1.344	2.39178e-07\\
1.346	2.33079e-07\\
1.348	2.26535e-07\\
1.35	2.19582e-07\\
1.352	2.12253e-07\\
1.354	2.05187e-07\\
1.356	2.0518e-07\\
1.358	2.04591e-07\\
1.36	2.03443e-07\\
1.362	2.01758e-07\\
1.364	1.99561e-07\\
1.366	1.96876e-07\\
1.368	1.93728e-07\\
1.37	1.90141e-07\\
1.372	1.86143e-07\\
1.374	1.81758e-07\\
1.376	1.77013e-07\\
1.378	1.71934e-07\\
1.38	1.66547e-07\\
1.382	1.6088e-07\\
1.384	1.54956e-07\\
1.386	1.48804e-07\\
1.388	1.42447e-07\\
1.39	1.35912e-07\\
1.392	1.29224e-07\\
1.394	1.22406e-07\\
1.396	1.15484e-07\\
1.398	1.08479e-07\\
1.4	1.01416e-07\\
1.402	9.43163e-08\\
1.404	8.72014e-08\\
1.406	8.00923e-08\\
1.408	7.65197e-08\\
1.41	7.9902e-08\\
1.412	8.29379e-08\\
1.414	8.56309e-08\\
1.416	8.79856e-08\\
1.418	9.00073e-08\\
1.42	9.17019e-08\\
1.422	9.30762e-08\\
1.424	9.41375e-08\\
1.426	9.48941e-08\\
1.428	9.53544e-08\\
1.43	9.55276e-08\\
1.432	9.54234e-08\\
1.434	9.50519e-08\\
1.436	9.44235e-08\\
1.438	9.35492e-08\\
1.44	9.244e-08\\
1.442	9.11076e-08\\
1.444	8.95634e-08\\
1.446	8.78196e-08\\
1.448	8.73615e-08\\
1.45	8.74234e-08\\
1.452	8.72342e-08\\
1.454	8.68034e-08\\
1.456	8.61407e-08\\
1.458	8.52564e-08\\
1.46	8.41609e-08\\
1.462	8.28647e-08\\
1.464	8.13787e-08\\
1.466	7.9714e-08\\
1.468	7.78817e-08\\
1.47	7.58928e-08\\
1.472	7.37588e-08\\
1.474	7.14909e-08\\
1.476	6.91003e-08\\
1.478	6.65981e-08\\
1.48	6.64589e-08\\
1.482	6.62709e-08\\
1.484	6.59013e-08\\
1.486	6.53576e-08\\
1.488	6.46477e-08\\
1.49	6.37794e-08\\
1.492	6.2761e-08\\
1.494	6.16006e-08\\
1.496	6.03068e-08\\
1.498	5.88879e-08\\
1.5	5.73525e-08\\
1.502	5.5709e-08\\
1.504	5.39661e-08\\
1.506	5.21322e-08\\
1.508	5.02159e-08\\
1.51	4.82254e-08\\
1.512	4.61692e-08\\
1.514	4.40553e-08\\
1.516	4.18919e-08\\
1.518	3.96869e-08\\
1.52	3.7448e-08\\
1.522	3.51828e-08\\
1.524	3.28987e-08\\
1.526	3.06027e-08\\
1.528	2.83019e-08\\
1.53	2.6003e-08\\
1.532	2.48525e-08\\
1.534	2.59514e-08\\
1.536	2.69373e-08\\
1.538	2.78113e-08\\
1.54	2.8575e-08\\
1.542	2.92301e-08\\
1.544	2.97788e-08\\
1.546	3.02231e-08\\
1.548	3.05656e-08\\
1.55	3.08091e-08\\
1.552	3.09562e-08\\
1.554	3.10101e-08\\
1.556	3.0974e-08\\
1.558	3.08511e-08\\
1.56	3.0645e-08\\
1.562	3.03592e-08\\
1.564	2.99974e-08\\
1.566	2.95632e-08\\
1.568	2.90607e-08\\
1.57	2.84935e-08\\
1.572	2.83448e-08\\
1.574	2.83643e-08\\
1.576	2.83026e-08\\
1.578	2.81627e-08\\
1.58	2.79478e-08\\
1.582	2.76612e-08\\
1.584	2.73063e-08\\
1.586	2.68864e-08\\
1.588	2.64052e-08\\
1.59	2.58661e-08\\
1.592	2.52728e-08\\
1.594	2.46288e-08\\
1.596	2.39378e-08\\
1.598	2.32035e-08\\
1.6	2.24294e-08\\
1.602	2.16191e-08\\
1.604	2.15055e-08\\
1.606	2.1449e-08\\
1.608	2.13336e-08\\
1.61	2.11617e-08\\
1.612	2.09357e-08\\
1.614	2.06583e-08\\
1.616	2.03319e-08\\
1.618	1.99595e-08\\
1.62	1.95435e-08\\
1.622	1.90868e-08\\
1.624	1.8592e-08\\
1.626	1.80621e-08\\
1.628	1.74996e-08\\
1.63	1.69075e-08\\
1.632	1.62883e-08\\
1.634	1.56449e-08\\
1.636	1.498e-08\\
1.638	1.42961e-08\\
1.64	1.3596e-08\\
1.642	1.28821e-08\\
1.644	1.21571e-08\\
1.646	1.14233e-08\\
1.648	1.06832e-08\\
1.65	9.93915e-09\\
1.652	9.19335e-09\\
1.654	8.448e-09\\
1.656	8.0308e-09\\
1.658	8.38911e-09\\
1.66	8.71084e-09\\
1.662	8.99636e-09\\
1.664	9.24617e-09\\
1.666	9.4608e-09\\
1.668	9.64092e-09\\
1.67	9.78722e-09\\
1.672	9.90051e-09\\
1.674	9.98162e-09\\
1.676	1.00315e-08\\
1.678	1.0051e-08\\
1.68	1.00413e-08\\
1.682	1.00034e-08\\
1.684	9.93843e-09\\
1.686	9.8475e-09\\
1.688	9.7318e-09\\
1.69	9.59255e-09\\
1.692	9.43099e-09\\
1.694	9.24835e-09\\
1.696	9.20277e-09\\
1.698	9.20973e-09\\
1.7	9.19025e-09\\
1.702	9.14531e-09\\
1.704	9.07596e-09\\
1.706	8.98325e-09\\
1.708	8.8683e-09\\
1.71	8.73222e-09\\
1.712	8.57614e-09\\
1.714	8.40123e-09\\
1.716	8.20866e-09\\
1.718	7.99961e-09\\
1.72	7.77526e-09\\
1.722	7.53679e-09\\
1.724	7.28538e-09\\
1.726	7.02222e-09\\
1.728	6.96716e-09\\
1.73	6.94883e-09\\
1.732	6.9114e-09\\
1.734	6.85565e-09\\
1.736	6.7824e-09\\
1.738	6.69249e-09\\
1.74	6.58676e-09\\
1.742	6.46607e-09\\
1.744	6.33132e-09\\
1.746	6.18339e-09\\
1.748	6.02315e-09\\
1.75	5.85152e-09\\
1.752	5.66939e-09\\
1.754	5.47765e-09\\
1.756	5.27718e-09\\
1.758	5.06887e-09\\
1.76	4.85359e-09\\
1.762	4.6322e-09\\
1.764	4.40555e-09\\
1.766	4.17447e-09\\
1.768	3.93977e-09\\
1.77	3.70226e-09\\
1.772	3.46269e-09\\
1.774	3.22184e-09\\
1.776	2.98042e-09\\
1.778	2.73914e-09\\
1.78	2.60237e-09\\
1.782	2.71863e-09\\
1.784	2.82301e-09\\
1.786	2.91564e-09\\
1.788	2.99667e-09\\
1.79	3.06629e-09\\
1.792	3.1247e-09\\
1.794	3.17215e-09\\
1.796	3.20889e-09\\
1.798	3.23519e-09\\
1.8	3.25137e-09\\
1.802	3.25772e-09\\
1.804	3.2546e-09\\
1.806	3.24233e-09\\
1.808	3.2213e-09\\
1.81	3.19186e-09\\
1.812	3.15441e-09\\
1.814	3.10933e-09\\
1.816	3.05704e-09\\
1.818	2.99792e-09\\
1.82	2.98432e-09\\
1.822	2.98675e-09\\
1.824	2.98062e-09\\
1.826	2.96624e-09\\
1.828	2.94395e-09\\
1.83	2.91408e-09\\
1.832	2.87701e-09\\
1.834	2.83309e-09\\
1.836	2.78268e-09\\
1.838	2.72616e-09\\
1.84	2.66391e-09\\
1.842	2.59632e-09\\
1.844	2.52375e-09\\
1.846	2.4466e-09\\
1.848	2.36525e-09\\
1.85	2.28007e-09\\
1.852	2.25935e-09\\
1.854	2.25362e-09\\
1.856	2.24169e-09\\
1.858	2.22381e-09\\
1.86	2.20024e-09\\
1.862	2.17124e-09\\
1.864	2.13711e-09\\
1.866	2.09811e-09\\
1.868	2.05453e-09\\
1.87	2.00665e-09\\
1.872	1.95477e-09\\
1.874	1.89918e-09\\
1.876	1.84016e-09\\
1.878	1.77801e-09\\
1.88	1.71301e-09\\
1.882	1.64546e-09\\
1.884	1.57563e-09\\
1.886	1.5038e-09\\
1.888	1.43025e-09\\
1.89	1.35526e-09\\
1.892	1.27908e-09\\
1.894	1.20197e-09\\
1.896	1.12419e-09\\
1.898	1.04598e-09\\
1.9	9.67588e-10\\
1.902	8.89234e-10\\
1.904	8.42326e-10\\
1.906	8.80226e-10\\
1.908	9.14274e-10\\
1.91	9.44513e-10\\
1.912	9.70989e-10\\
1.914	9.93763e-10\\
1.916	1.0129e-09\\
1.918	1.02848e-09\\
1.92	1.04058e-09\\
1.922	1.04928e-09\\
1.924	1.0547e-09\\
1.926	1.05692e-09\\
1.928	1.05606e-09\\
1.93	1.05222e-09\\
1.932	1.04553e-09\\
1.934	1.0361e-09\\
1.936	1.02406e-09\\
1.938	1.00953e-09\\
1.94	9.92655e-10\\
1.942	9.73552e-10\\
1.944	9.6916e-10\\
1.946	9.69996e-10\\
1.948	9.68041e-10\\
1.95	9.63402e-10\\
1.952	9.56187e-10\\
1.954	9.46508e-10\\
1.956	9.34481e-10\\
1.958	9.20223e-10\\
1.96	9.03856e-10\\
1.962	8.855e-10\\
1.964	8.6528e-10\\
1.966	8.43318e-10\\
1.968	8.1974e-10\\
1.97	7.9467e-10\\
1.972	7.68234e-10\\
1.974	7.40554e-10\\
1.976	7.33114e-10\\
1.978	7.31212e-10\\
1.98	7.27297e-10\\
1.982	7.21454e-10\\
1.984	7.13767e-10\\
1.986	7.04325e-10\\
1.988	6.93216e-10\\
1.99	6.80533e-10\\
1.992	6.66367e-10\\
1.994	6.50811e-10\\
1.996	6.3396e-10\\
1.998	6.15908e-10\\
2	5.96749e-10\\
};
\addplot [color=mycolor4, line width=1.0pt, forget plot]
  table[row sep=crcr]{%
-0.024	0\\
-0.022	0\\
-0.02	0\\
-0.018	0\\
-0.016	0\\
-0.014	0\\
-0.012	0\\
-0.01	0\\
-0.008	0\\
-0.006	0\\
-0.004	0\\
-0.002	0\\
0	0\\
0.002	0.00916185\\
0.004	0.0147068\\
0.006	0.019866\\
0.008	0.0237565\\
0.01	0.0264989\\
0.012	0.0282394\\
0.014	0.0291282\\
0.016	0.0293151\\
0.018	0.0289591\\
0.02	0.0291495\\
0.022	0.0292754\\
0.024	0.0293744\\
0.026	0.0295711\\
0.028	0.0301204\\
0.03	0.0303673\\
0.032	0.0303193\\
0.034	0.0299894\\
0.036	0.0304733\\
0.038	0.031424\\
0.04	0.0322821\\
0.042	0.0330485\\
0.044	0.033724\\
0.046	0.0343098\\
0.048	0.0348071\\
0.05	0.0352174\\
0.052	0.0355422\\
0.054	0.0357833\\
0.056	0.0359425\\
0.058	0.0360216\\
0.06	0.0360229\\
0.062	0.0359485\\
0.064	0.0358006\\
0.066	0.0355818\\
0.068	0.0352946\\
0.07	0.0349415\\
0.072	0.0345253\\
0.074	0.0340489\\
0.076	0.0335151\\
0.078	0.0329269\\
0.08	0.0322875\\
0.082	0.0316\\
0.084	0.0308676\\
0.086	0.0300935\\
0.088	0.029281\\
0.09	0.0284335\\
0.092	0.0275543\\
0.094	0.0266468\\
0.096	0.0257142\\
0.098	0.0247599\\
0.1	0.0237872\\
0.102	0.0227993\\
0.104	0.0217995\\
0.106	0.0207908\\
0.108	0.0197763\\
0.11	0.018759\\
0.112	0.0177418\\
0.114	0.0167275\\
0.116	0.0160648\\
0.118	0.0156387\\
0.12	0.015195\\
0.122	0.0147355\\
0.124	0.014262\\
0.126	0.0137762\\
0.128	0.01328\\
0.13	0.012775\\
0.132	0.0123938\\
0.134	0.01227\\
0.136	0.0121285\\
0.138	0.0119699\\
0.14	0.0117953\\
0.142	0.0116055\\
0.144	0.0114014\\
0.146	0.011184\\
0.148	0.0109541\\
0.15	0.0107126\\
0.152	0.0104606\\
0.154	0.0101988\\
0.156	0.00992809\\
0.158	0.00964946\\
0.16	0.00936373\\
0.162	0.00907174\\
0.164	0.00877433\\
0.166	0.00847232\\
0.168	0.00817639\\
0.17	0.00794804\\
0.172	0.00772002\\
0.174	0.00749282\\
0.176	0.00726694\\
0.178	0.00704284\\
0.18	0.00682094\\
0.182	0.00660167\\
0.184	0.0063854\\
0.186	0.00617249\\
0.188	0.00596324\\
0.19	0.00575796\\
0.192	0.00555689\\
0.194	0.00536026\\
0.196	0.00516827\\
0.198	0.00498109\\
0.2	0.00479884\\
0.202	0.00465261\\
0.204	0.00458295\\
0.206	0.00449993\\
0.208	0.00440457\\
0.21	0.00429789\\
0.212	0.00418087\\
0.214	0.00405448\\
0.216	0.00391967\\
0.218	0.00377736\\
0.22	0.00362844\\
0.222	0.00347378\\
0.224	0.00331421\\
0.226	0.00315054\\
0.228	0.00298354\\
0.23	0.00281396\\
0.232	0.00264249\\
0.234	0.00246981\\
0.236	0.00236915\\
0.238	0.00227591\\
0.24	0.00222399\\
0.242	0.00218015\\
0.244	0.00213049\\
0.246	0.00207542\\
0.248	0.00201539\\
0.25	0.00195081\\
0.252	0.00188209\\
0.254	0.00180964\\
0.256	0.00173386\\
0.258	0.00165514\\
0.26	0.00157386\\
0.262	0.00150713\\
0.264	0.00153755\\
0.266	0.00156157\\
0.268	0.00157941\\
0.27	0.00159131\\
0.272	0.00159751\\
0.274	0.00159827\\
0.276	0.00159384\\
0.278	0.00158449\\
0.28	0.00157046\\
0.282	0.00155203\\
0.284	0.00152946\\
0.286	0.00150302\\
0.288	0.00147296\\
0.29	0.00143956\\
0.292	0.00140307\\
0.294	0.00136375\\
0.296	0.00132186\\
0.298	0.00130656\\
0.3	0.00133866\\
0.302	0.0013657\\
0.304	0.0013878\\
0.306	0.00140513\\
0.308	0.00141783\\
0.31	0.00142606\\
0.312	0.00143\\
0.314	0.00142981\\
0.316	0.00142568\\
0.318	0.00141778\\
0.32	0.00140631\\
0.322	0.00139145\\
0.324	0.00137339\\
0.326	0.00135232\\
0.328	0.00132844\\
0.33	0.00130193\\
0.332	0.001273\\
0.334	0.00124183\\
0.336	0.00120862\\
0.338	0.00117354\\
0.34	0.00113679\\
0.342	0.00109855\\
0.344	0.001059\\
0.346	0.00101831\\
0.348	0.000976658\\
0.35	0.000934205\\
0.352	0.000891114\\
0.354	0.000847541\\
0.356	0.000819213\\
0.358	0.000807315\\
0.36	0.00079371\\
0.362	0.000778511\\
0.364	0.000761831\\
0.366	0.000743784\\
0.368	0.000724482\\
0.37	0.000704036\\
0.372	0.000682556\\
0.374	0.000660149\\
0.376	0.000636922\\
0.378	0.000612978\\
0.38	0.000588419\\
0.382	0.000563342\\
0.384	0.000537843\\
0.386	0.000512015\\
0.388	0.000485946\\
0.39	0.000459721\\
0.392	0.000436566\\
0.394	0.000438905\\
0.396	0.000439961\\
0.398	0.000439789\\
0.4	0.000438442\\
0.402	0.000435978\\
0.404	0.000432453\\
0.406	0.000427923\\
0.408	0.000422447\\
0.41	0.000416081\\
0.412	0.000408884\\
0.414	0.000400912\\
0.416	0.000392222\\
0.418	0.00038287\\
0.42	0.000372913\\
0.422	0.000362404\\
0.424	0.000351398\\
0.426	0.000339948\\
0.428	0.000328106\\
0.43	0.000315922\\
0.432	0.000303445\\
0.434	0.000290724\\
0.436	0.000277806\\
0.438	0.000264734\\
0.44	0.000260639\\
0.442	0.000260956\\
0.444	0.000260439\\
0.446	0.000259131\\
0.448	0.000257076\\
0.45	0.00025432\\
0.452	0.000250904\\
0.454	0.000246874\\
0.456	0.000242272\\
0.458	0.000237142\\
0.46	0.000231525\\
0.462	0.000225463\\
0.464	0.000218997\\
0.466	0.000212167\\
0.468	0.000205012\\
0.47	0.000197569\\
0.472	0.000189876\\
0.474	0.000181968\\
0.476	0.000173881\\
0.478	0.000170492\\
0.48	0.000169407\\
0.482	0.000167849\\
0.484	0.000165845\\
0.486	0.00016342\\
0.488	0.000160602\\
0.49	0.000157416\\
0.492	0.00015389\\
0.494	0.00015005\\
0.496	0.000145921\\
0.498	0.00014153\\
0.5	0.000136902\\
0.502	0.000132062\\
0.504	0.000127035\\
0.506	0.000121844\\
0.508	0.000116514\\
0.51	0.000111067\\
0.512	0.000105525\\
0.514	9.99099e-05\\
0.516	9.42424e-05\\
0.518	8.85424e-05\\
0.52	8.28293e-05\\
0.522	7.9549e-05\\
0.524	8.02316e-05\\
0.526	8.0637e-05\\
0.528	8.07781e-05\\
0.53	8.06675e-05\\
0.532	8.03184e-05\\
0.534	7.97436e-05\\
0.536	7.89563e-05\\
0.538	7.79694e-05\\
0.54	7.67962e-05\\
0.542	7.54493e-05\\
0.544	7.39417e-05\\
0.546	7.22859e-05\\
0.548	7.04945e-05\\
0.55	6.85796e-05\\
0.552	6.65533e-05\\
0.554	6.44272e-05\\
0.556	6.22128e-05\\
0.558	5.99212e-05\\
0.56	5.75631e-05\\
0.562	5.72105e-05\\
0.564	5.75222e-05\\
0.566	5.7646e-05\\
0.568	5.75906e-05\\
0.57	5.73649e-05\\
0.572	5.69777e-05\\
0.574	5.6438e-05\\
0.576	5.5755e-05\\
0.578	5.49376e-05\\
0.58	5.39949e-05\\
0.582	5.29359e-05\\
0.584	5.17695e-05\\
0.586	5.05045e-05\\
0.588	4.91496e-05\\
0.59	4.77133e-05\\
0.592	4.62041e-05\\
0.594	4.46301e-05\\
0.596	4.29992e-05\\
0.598	4.13194e-05\\
0.6	4.07225e-05\\
0.602	4.05508e-05\\
0.604	4.02666e-05\\
0.606	3.98763e-05\\
0.608	3.93858e-05\\
0.61	3.88015e-05\\
0.612	3.81296e-05\\
0.614	3.7376e-05\\
0.616	3.65471e-05\\
0.618	3.56487e-05\\
0.62	3.46868e-05\\
0.622	3.36674e-05\\
0.624	3.2596e-05\\
0.626	3.14784e-05\\
0.628	3.03199e-05\\
0.63	2.91261e-05\\
0.632	2.79021e-05\\
0.634	2.66529e-05\\
0.636	2.53834e-05\\
0.638	2.40983e-05\\
0.64	2.28021e-05\\
0.642	2.14993e-05\\
0.644	2.01939e-05\\
0.646	1.88899e-05\\
0.648	1.75912e-05\\
0.65	1.63013e-05\\
0.652	1.60929e-05\\
0.654	1.61614e-05\\
0.656	1.6179e-05\\
0.658	1.6148e-05\\
0.66	1.60709e-05\\
0.662	1.59503e-05\\
0.664	1.57887e-05\\
0.666	1.55886e-05\\
0.668	1.53524e-05\\
0.67	1.50827e-05\\
0.672	1.47819e-05\\
0.674	1.44525e-05\\
0.676	1.40968e-05\\
0.678	1.37172e-05\\
0.68	1.33161e-05\\
0.682	1.28957e-05\\
0.684	1.24581e-05\\
0.686	1.20057e-05\\
0.688	1.15404e-05\\
0.69	1.13018e-05\\
0.692	1.13436e-05\\
0.694	1.13489e-05\\
0.696	1.13195e-05\\
0.698	1.12571e-05\\
0.7	1.11637e-05\\
0.702	1.1041e-05\\
0.704	1.08908e-05\\
0.706	1.0715e-05\\
0.708	1.05153e-05\\
0.71	1.02935e-05\\
0.712	1.00514e-05\\
0.714	9.79074e-06\\
0.716	9.51316e-06\\
0.718	9.22038e-06\\
0.72	8.91401e-06\\
0.722	8.77324e-06\\
0.724	8.79081e-06\\
0.726	8.78067e-06\\
0.728	8.74414e-06\\
0.73	8.68258e-06\\
0.732	8.59736e-06\\
0.734	8.48986e-06\\
0.736	8.36147e-06\\
0.738	8.21357e-06\\
0.74	8.04755e-06\\
0.742	7.86477e-06\\
0.744	7.6666e-06\\
0.746	7.45438e-06\\
0.748	7.22943e-06\\
0.75	6.99306e-06\\
0.752	6.74655e-06\\
0.754	6.49114e-06\\
0.756	6.22806e-06\\
0.758	5.95847e-06\\
0.76	5.68354e-06\\
0.762	5.40437e-06\\
0.764	5.12202e-06\\
0.766	4.83753e-06\\
0.768	4.55188e-06\\
0.77	4.26602e-06\\
0.772	3.98084e-06\\
0.774	3.69718e-06\\
0.776	3.41587e-06\\
0.778	3.13765e-06\\
0.78	2.96789e-06\\
0.782	2.98415e-06\\
0.784	2.99063e-06\\
0.786	2.98782e-06\\
0.788	2.9762e-06\\
0.79	2.95626e-06\\
0.792	2.9285e-06\\
0.794	2.89339e-06\\
0.796	2.85144e-06\\
0.798	2.80311e-06\\
0.8	2.74888e-06\\
0.802	2.68922e-06\\
0.804	2.62459e-06\\
0.806	2.55544e-06\\
0.808	2.4822e-06\\
0.81	2.40532e-06\\
0.812	2.3252e-06\\
0.814	2.24225e-06\\
0.816	2.25572e-06\\
0.818	2.26378e-06\\
0.82	2.26448e-06\\
0.822	2.25817e-06\\
0.824	2.24523e-06\\
0.826	2.22605e-06\\
0.828	2.20099e-06\\
0.83	2.17044e-06\\
0.832	2.13476e-06\\
0.834	2.09433e-06\\
0.836	2.04952e-06\\
0.838	2.00068e-06\\
0.84	1.94817e-06\\
0.842	1.89235e-06\\
0.844	1.85974e-06\\
0.846	1.87237e-06\\
0.848	1.8787e-06\\
0.85	1.87903e-06\\
0.852	1.87365e-06\\
0.854	1.86284e-06\\
0.856	1.84691e-06\\
0.858	1.82615e-06\\
0.86	1.80085e-06\\
0.862	1.77132e-06\\
0.864	1.73784e-06\\
0.866	1.70071e-06\\
0.868	1.66022e-06\\
0.87	1.61664e-06\\
0.872	1.57026e-06\\
0.874	1.52136e-06\\
0.876	1.4702e-06\\
0.878	1.41705e-06\\
0.88	1.36215e-06\\
0.882	1.30577e-06\\
0.884	1.24814e-06\\
0.886	1.18949e-06\\
0.888	1.13006e-06\\
0.89	1.07006e-06\\
0.892	1.0097e-06\\
0.894	9.49172e-07\\
0.896	8.88677e-07\\
0.898	8.28393e-07\\
0.9	7.68494e-07\\
0.902	7.09142e-07\\
0.904	6.50492e-07\\
0.906	5.92688e-07\\
0.908	5.40168e-07\\
0.91	5.43881e-07\\
0.912	5.45759e-07\\
0.914	5.45891e-07\\
0.916	5.44366e-07\\
0.918	5.41274e-07\\
0.92	5.36705e-07\\
0.922	5.30749e-07\\
0.924	5.23496e-07\\
0.926	5.15035e-07\\
0.928	5.05455e-07\\
0.93	4.94842e-07\\
0.932	4.83282e-07\\
0.934	4.70861e-07\\
0.936	4.5766e-07\\
0.938	4.43762e-07\\
0.94	4.3776e-07\\
0.942	4.3292e-07\\
0.944	4.3353e-07\\
0.946	4.33882e-07\\
0.948	4.3287e-07\\
0.95	4.30567e-07\\
0.952	4.27044e-07\\
0.954	4.22372e-07\\
0.956	4.16625e-07\\
0.958	4.09874e-07\\
0.96	4.0219e-07\\
0.962	3.93644e-07\\
0.964	3.84305e-07\\
0.966	3.82226e-07\\
0.968	3.8712e-07\\
0.97	3.90566e-07\\
0.972	3.92624e-07\\
0.974	3.93356e-07\\
0.976	3.92825e-07\\
0.978	3.91094e-07\\
0.98	3.88228e-07\\
0.982	3.84292e-07\\
0.984	3.7935e-07\\
0.986	3.73468e-07\\
0.988	3.66712e-07\\
0.99	3.59144e-07\\
0.992	3.5083e-07\\
0.994	3.41832e-07\\
0.996	3.32212e-07\\
0.998	3.22033e-07\\
1	3.11353e-07\\
1.002	3.00231e-07\\
1.004	2.88725e-07\\
1.006	2.76888e-07\\
1.008	2.64776e-07\\
1.01	2.5244e-07\\
1.012	2.3993e-07\\
1.014	2.27293e-07\\
1.016	2.14576e-07\\
1.018	2.01823e-07\\
1.02	1.89075e-07\\
1.022	1.76373e-07\\
1.024	1.63753e-07\\
1.026	1.51251e-07\\
1.028	1.389e-07\\
1.03	1.26732e-07\\
1.032	1.14775e-07\\
1.034	1.03056e-07\\
1.036	9.28452e-08\\
1.038	9.35709e-08\\
1.04	9.39725e-08\\
1.042	9.4066e-08\\
1.044	9.38675e-08\\
1.046	9.41584e-08\\
1.048	9.53504e-08\\
1.05	9.6165e-08\\
1.052	9.66184e-08\\
1.054	9.67271e-08\\
1.056	9.65082e-08\\
1.058	9.59788e-08\\
1.06	9.51564e-08\\
1.062	9.40585e-08\\
1.064	9.27026e-08\\
1.066	9.11062e-08\\
1.068	8.92869e-08\\
1.07	8.7262e-08\\
1.072	8.50486e-08\\
1.074	8.26638e-08\\
1.076	8.16232e-08\\
1.078	8.08062e-08\\
1.08	7.97769e-08\\
1.082	7.85496e-08\\
1.084	7.71388e-08\\
1.086	7.55583e-08\\
1.088	7.67185e-08\\
1.09	7.81113e-08\\
1.092	7.9193e-08\\
1.094	7.99756e-08\\
1.096	8.04713e-08\\
1.098	8.06927e-08\\
1.1	8.06524e-08\\
1.102	8.03635e-08\\
1.104	7.9839e-08\\
1.106	7.9092e-08\\
1.108	7.81357e-08\\
1.11	7.69833e-08\\
1.112	7.56479e-08\\
1.114	7.41425e-08\\
1.116	7.24802e-08\\
1.118	7.06737e-08\\
1.12	6.87356e-08\\
1.122	6.66783e-08\\
1.124	6.45141e-08\\
1.126	6.22549e-08\\
1.128	5.99123e-08\\
1.13	5.74977e-08\\
1.132	5.50221e-08\\
1.134	5.24962e-08\\
1.136	4.99303e-08\\
1.138	4.73345e-08\\
1.14	4.47183e-08\\
1.142	4.20908e-08\\
1.144	3.9461e-08\\
1.146	3.68371e-08\\
1.148	3.42271e-08\\
1.15	3.16385e-08\\
1.152	2.90785e-08\\
1.154	2.65537e-08\\
1.156	2.40703e-08\\
1.158	2.16342e-08\\
1.16	1.92508e-08\\
1.162	1.94024e-08\\
1.164	1.98306e-08\\
1.166	2.01719e-08\\
1.168	2.04296e-08\\
1.17	2.06068e-08\\
1.172	2.07069e-08\\
1.174	2.07334e-08\\
1.176	2.06899e-08\\
1.178	2.058e-08\\
1.18	2.04071e-08\\
1.182	2.01751e-08\\
1.184	1.98876e-08\\
1.186	1.95482e-08\\
1.188	1.91607e-08\\
1.19	1.87286e-08\\
1.192	1.82555e-08\\
1.194	1.77452e-08\\
1.196	1.7201e-08\\
1.198	1.66264e-08\\
1.2	1.60249e-08\\
1.202	1.53997e-08\\
1.204	1.50238e-08\\
1.206	1.47982e-08\\
1.208	1.50064e-08\\
1.21	1.54226e-08\\
1.212	1.57686e-08\\
1.214	1.60467e-08\\
1.216	1.62597e-08\\
1.218	1.641e-08\\
1.22	1.65004e-08\\
1.222	1.65336e-08\\
1.224	1.65125e-08\\
1.226	1.64397e-08\\
1.228	1.63183e-08\\
1.23	1.6151e-08\\
1.232	1.59407e-08\\
1.234	1.56903e-08\\
1.236	1.54026e-08\\
1.238	1.50805e-08\\
1.24	1.47266e-08\\
1.242	1.43439e-08\\
1.244	1.39348e-08\\
1.246	1.35022e-08\\
1.248	1.30485e-08\\
1.25	1.25763e-08\\
1.252	1.20881e-08\\
1.254	1.15861e-08\\
1.256	1.10727e-08\\
1.258	1.05502e-08\\
1.26	1.00205e-08\\
1.262	9.48582e-09\\
1.264	8.94804e-09\\
1.266	8.40904e-09\\
1.268	7.87057e-09\\
1.27	7.33432e-09\\
1.272	6.80186e-09\\
1.274	6.27469e-09\\
1.276	5.7542e-09\\
1.278	5.2417e-09\\
1.28	4.7384e-09\\
1.282	4.26554e-09\\
1.284	4.35153e-09\\
1.286	4.41861e-09\\
1.288	4.4675e-09\\
1.29	4.49896e-09\\
1.292	4.51376e-09\\
1.294	4.51269e-09\\
1.296	4.49656e-09\\
1.298	4.46619e-09\\
1.3	4.42239e-09\\
1.302	4.36599e-09\\
1.304	4.29782e-09\\
1.306	4.21869e-09\\
1.308	4.12944e-09\\
1.31	4.03086e-09\\
1.312	3.92376e-09\\
1.314	3.80893e-09\\
1.316	3.68713e-09\\
1.318	3.55914e-09\\
1.32	3.42568e-09\\
1.322	3.28747e-09\\
1.324	3.14522e-09\\
1.326	2.9996e-09\\
1.328	2.86927e-09\\
1.33	2.97928e-09\\
1.332	3.07426e-09\\
1.334	3.15466e-09\\
1.336	3.22095e-09\\
1.338	3.27361e-09\\
1.34	3.31315e-09\\
1.342	3.34009e-09\\
1.344	3.35496e-09\\
1.346	3.35832e-09\\
1.348	3.35072e-09\\
1.35	3.33273e-09\\
1.352	3.30491e-09\\
1.354	3.26785e-09\\
1.356	3.22212e-09\\
1.358	3.16829e-09\\
1.36	3.10694e-09\\
1.362	3.03863e-09\\
1.364	2.96394e-09\\
1.366	2.88341e-09\\
1.368	2.79759e-09\\
1.37	2.70704e-09\\
1.372	2.61226e-09\\
1.374	2.51379e-09\\
1.376	2.41212e-09\\
1.378	2.30773e-09\\
1.38	2.20112e-09\\
1.382	2.09272e-09\\
1.384	1.98299e-09\\
1.386	1.87234e-09\\
1.388	1.76119e-09\\
1.39	1.64992e-09\\
1.392	1.5389e-09\\
1.394	1.42848e-09\\
1.396	1.31899e-09\\
1.398	1.21073e-09\\
1.4	1.10402e-09\\
1.402	9.99105e-10\\
1.404	9.58842e-10\\
1.406	9.72601e-10\\
1.408	9.82445e-10\\
1.41	9.88532e-10\\
1.412	9.91025e-10\\
1.414	9.90089e-10\\
1.416	9.85895e-10\\
1.418	9.78616e-10\\
1.42	9.68424e-10\\
1.422	9.55497e-10\\
1.424	9.4001e-10\\
1.426	9.22139e-10\\
1.428	9.0206e-10\\
1.43	8.79947e-10\\
1.432	8.55974e-10\\
1.434	8.3031e-10\\
1.436	8.03125e-10\\
1.438	7.74583e-10\\
1.44	7.44846e-10\\
1.442	7.14072e-10\\
1.444	6.82416e-10\\
1.446	6.50025e-10\\
1.448	6.17046e-10\\
1.45	5.88488e-10\\
1.452	6.11095e-10\\
1.454	6.30557e-10\\
1.456	6.46972e-10\\
1.458	6.60442e-10\\
1.46	6.71073e-10\\
1.462	6.78978e-10\\
1.464	6.84268e-10\\
1.466	6.87062e-10\\
1.468	6.87477e-10\\
1.47	6.85634e-10\\
1.472	6.81655e-10\\
1.474	6.7566e-10\\
1.476	6.67774e-10\\
1.478	6.58117e-10\\
1.48	6.46811e-10\\
1.482	6.33977e-10\\
1.484	6.19734e-10\\
1.486	6.04199e-10\\
1.488	5.8749e-10\\
1.49	5.69718e-10\\
1.492	5.50996e-10\\
1.494	5.31432e-10\\
1.496	5.11131e-10\\
1.498	4.90197e-10\\
1.5	4.68728e-10\\
1.502	4.4682e-10\\
1.504	4.24567e-10\\
1.506	4.02058e-10\\
1.508	3.79376e-10\\
1.51	3.56605e-10\\
1.512	3.33821e-10\\
1.514	3.11099e-10\\
1.516	2.88508e-10\\
1.518	2.66114e-10\\
1.52	2.4398e-10\\
1.522	2.22163e-10\\
1.524	2.16224e-10\\
1.526	2.18765e-10\\
1.528	2.20454e-10\\
1.53	2.21327e-10\\
1.532	2.21422e-10\\
1.534	2.20777e-10\\
1.536	2.19432e-10\\
1.538	2.17425e-10\\
1.54	2.14796e-10\\
1.542	2.11584e-10\\
1.544	2.0783e-10\\
1.546	2.03572e-10\\
1.548	1.9885e-10\\
1.55	1.93702e-10\\
1.552	1.88166e-10\\
1.554	1.82279e-10\\
1.556	1.76079e-10\\
1.558	1.69601e-10\\
1.56	1.6288e-10\\
1.562	1.5595e-10\\
1.564	1.48844e-10\\
1.566	1.41595e-10\\
1.568	1.34233e-10\\
1.57	1.26787e-10\\
1.572	1.23976e-10\\
1.574	1.28565e-10\\
1.576	1.32501e-10\\
1.578	1.35804e-10\\
1.58	1.38493e-10\\
1.582	1.40591e-10\\
1.584	1.42119e-10\\
1.586	1.43101e-10\\
1.588	1.43561e-10\\
1.59	1.43522e-10\\
1.592	1.43011e-10\\
1.594	1.42051e-10\\
1.596	1.40668e-10\\
1.598	1.38888e-10\\
1.6	1.36737e-10\\
1.602	1.34239e-10\\
1.604	1.31422e-10\\
1.606	1.28308e-10\\
1.608	1.24925e-10\\
1.61	1.21297e-10\\
1.612	1.17447e-10\\
1.614	1.134e-10\\
1.616	1.09179e-10\\
1.618	1.04806e-10\\
1.62	1.00305e-10\\
1.622	9.56964e-11\\
1.624	9.10008e-11\\
1.626	8.62384e-11\\
1.628	8.14285e-11\\
1.63	7.65895e-11\\
1.632	7.1739e-11\\
1.634	6.68938e-11\\
1.636	6.20699e-11\\
1.638	5.72821e-11\\
1.64	5.25448e-11\\
1.642	4.83561e-11\\
1.644	4.90906e-11\\
1.646	4.96275e-11\\
1.648	4.99745e-11\\
1.65	5.01397e-11\\
1.652	5.01312e-11\\
1.654	4.99574e-11\\
1.656	4.96268e-11\\
1.658	4.9148e-11\\
1.66	4.85298e-11\\
1.662	4.77809e-11\\
1.664	4.69101e-11\\
1.666	4.59261e-11\\
1.668	4.48377e-11\\
1.67	4.36535e-11\\
1.672	4.2382e-11\\
1.674	4.10318e-11\\
1.676	3.96111e-11\\
1.678	3.81281e-11\\
1.68	3.65906e-11\\
1.682	3.50065e-11\\
1.684	3.33833e-11\\
1.686	3.17283e-11\\
1.688	3.00485e-11\\
1.69	2.83507e-11\\
1.692	2.66416e-11\\
1.694	2.74288e-11\\
1.696	2.82733e-11\\
1.698	2.89824e-11\\
1.7	2.95606e-11\\
1.702	3.00127e-11\\
1.704	3.03434e-11\\
1.706	3.05577e-11\\
1.708	3.06609e-11\\
1.71	3.0658e-11\\
1.712	3.05545e-11\\
1.714	3.03556e-11\\
1.716	3.00668e-11\\
1.718	2.96935e-11\\
1.72	2.9241e-11\\
1.722	2.87149e-11\\
1.724	2.81204e-11\\
1.726	2.74629e-11\\
1.728	2.67476e-11\\
1.73	2.59798e-11\\
1.732	2.51645e-11\\
1.734	2.43066e-11\\
1.736	2.34111e-11\\
1.738	2.24828e-11\\
1.74	2.15262e-11\\
1.742	2.05458e-11\\
1.744	1.95461e-11\\
1.746	1.85311e-11\\
1.748	1.7505e-11\\
1.75	1.64716e-11\\
1.752	1.54345e-11\\
1.754	1.43975e-11\\
1.756	1.33637e-11\\
1.758	1.23365e-11\\
1.76	1.13188e-11\\
1.762	1.10196e-11\\
1.764	1.11733e-11\\
1.766	1.12832e-11\\
1.768	1.13511e-11\\
1.77	1.13787e-11\\
1.772	1.13681e-11\\
1.774	1.1321e-11\\
1.776	1.12394e-11\\
1.778	1.11253e-11\\
1.78	1.09805e-11\\
1.782	1.08072e-11\\
1.784	1.06071e-11\\
1.786	1.03824e-11\\
1.788	1.01348e-11\\
1.79	9.86634e-12\\
1.792	9.57887e-12\\
1.794	9.27422e-12\\
1.796	8.95422e-12\\
1.798	8.62062e-12\\
1.8	8.27517e-12\\
1.802	7.91955e-12\\
1.804	7.55541e-12\\
1.806	7.18433e-12\\
1.808	6.80783e-12\\
1.81	6.4274e-12\\
1.812	6.04443e-12\\
1.814	6.12189e-12\\
1.816	6.30984e-12\\
1.818	6.46771e-12\\
1.82	6.59645e-12\\
1.822	6.69704e-12\\
1.824	6.7705e-12\\
1.826	6.81793e-12\\
1.828	6.84044e-12\\
1.83	6.83912e-12\\
1.832	6.81516e-12\\
1.834	6.76972e-12\\
1.836	6.70403e-12\\
1.838	6.61926e-12\\
1.84	6.51666e-12\\
1.842	6.39741e-12\\
1.844	6.26272e-12\\
1.846	6.1138e-12\\
1.848	5.95188e-12\\
1.85	5.77808e-12\\
1.852	5.5936e-12\\
1.854	5.39954e-12\\
1.856	5.19704e-12\\
1.858	4.9872e-12\\
1.86	4.77106e-12\\
1.862	4.54967e-12\\
1.864	4.324e-12\\
1.866	4.09502e-12\\
1.868	3.86364e-12\\
1.87	3.63077e-12\\
1.872	3.39725e-12\\
1.874	3.16389e-12\\
1.876	2.93144e-12\\
1.878	2.70064e-12\\
1.88	2.47217e-12\\
1.882	2.49455e-12\\
1.884	2.53225e-12\\
1.886	2.55991e-12\\
1.888	2.57792e-12\\
1.89	2.58669e-12\\
1.892	2.58666e-12\\
1.894	2.5782e-12\\
1.896	2.56177e-12\\
1.898	2.5378e-12\\
1.9	2.50673e-12\\
1.902	2.46902e-12\\
1.904	2.4251e-12\\
1.906	2.37543e-12\\
1.908	2.32044e-12\\
1.91	2.26055e-12\\
1.912	2.19623e-12\\
1.914	2.12788e-12\\
1.916	2.0559e-12\\
1.918	1.98074e-12\\
1.92	1.9028e-12\\
1.922	1.82248e-12\\
1.924	1.74013e-12\\
1.926	1.65611e-12\\
1.928	1.57079e-12\\
1.93	1.48451e-12\\
1.932	1.3976e-12\\
1.934	1.40284e-12\\
1.936	1.4438e-12\\
1.938	1.47813e-12\\
1.94	1.50603e-12\\
1.942	1.52774e-12\\
1.944	1.54345e-12\\
1.946	1.55342e-12\\
1.948	1.55787e-12\\
1.95	1.55708e-12\\
1.952	1.5513e-12\\
1.954	1.54078e-12\\
1.956	1.52579e-12\\
1.958	1.50657e-12\\
1.96	1.4834e-12\\
1.962	1.45654e-12\\
1.964	1.42625e-12\\
1.966	1.39279e-12\\
1.968	1.35642e-12\\
1.97	1.31739e-12\\
1.972	1.27598e-12\\
1.974	1.23239e-12\\
1.976	1.18688e-12\\
1.978	1.13968e-12\\
1.98	1.09102e-12\\
1.982	1.04116e-12\\
1.984	9.90286e-13\\
1.986	9.38626e-13\\
1.988	8.86365e-13\\
1.99	8.33687e-13\\
1.992	7.80797e-13\\
1.994	7.27881e-13\\
1.996	6.75086e-13\\
1.998	6.22605e-13\\
2	5.70585e-13\\
};
\addplot [color=mycolor5, line width=1.0pt, forget plot]
  table[row sep=crcr]{%
-0.024	0\\
-0.022	0\\
-0.02	0\\
-0.018	0\\
-0.016	0\\
-0.014	0\\
-0.012	0\\
-0.01	0\\
-0.008	0\\
-0.006	0\\
-0.004	0\\
-0.002	0\\
0	0\\
0.002	0.00916185\\
0.004	0.0147068\\
0.006	0.019866\\
0.008	0.0237566\\
0.01	0.026499\\
0.012	0.0282395\\
0.014	0.0291284\\
0.016	0.0293152\\
0.018	0.0289594\\
0.02	0.0291496\\
0.022	0.0292755\\
0.024	0.0293743\\
0.026	0.0295714\\
0.028	0.0301203\\
0.03	0.0303665\\
0.032	0.0303177\\
0.034	0.0299867\\
0.036	0.0301477\\
0.038	0.0310682\\
0.04	0.0318953\\
0.042	0.03263\\
0.044	0.0332731\\
0.046	0.0338259\\
0.048	0.0342897\\
0.05	0.0346661\\
0.052	0.0349566\\
0.054	0.0351632\\
0.056	0.0352878\\
0.058	0.0353325\\
0.06	0.0352996\\
0.062	0.0351913\\
0.064	0.0350103\\
0.066	0.0347591\\
0.068	0.0344404\\
0.07	0.0340571\\
0.072	0.0336122\\
0.074	0.0331087\\
0.076	0.0325497\\
0.078	0.0319385\\
0.08	0.0312783\\
0.082	0.0305726\\
0.084	0.0298246\\
0.086	0.0290379\\
0.088	0.028216\\
0.09	0.0273623\\
0.092	0.0264802\\
0.094	0.0255734\\
0.096	0.0246451\\
0.098	0.0236989\\
0.1	0.022738\\
0.102	0.0217657\\
0.104	0.0207852\\
0.106	0.0197997\\
0.108	0.0188121\\
0.11	0.0178254\\
0.112	0.0168424\\
0.114	0.0158657\\
0.116	0.0151438\\
0.118	0.0147532\\
0.12	0.014348\\
0.122	0.0139298\\
0.124	0.0135003\\
0.126	0.013061\\
0.128	0.0126135\\
0.13	0.0121593\\
0.132	0.0116999\\
0.134	0.0115387\\
0.136	0.01145\\
0.138	0.0113459\\
0.14	0.0112272\\
0.142	0.0110946\\
0.144	0.0109487\\
0.146	0.0107903\\
0.148	0.0106201\\
0.15	0.0104386\\
0.152	0.0102468\\
0.154	0.0100452\\
0.156	0.00983456\\
0.158	0.00961558\\
0.16	0.00938894\\
0.162	0.00915532\\
0.164	0.00891541\\
0.166	0.00866987\\
0.168	0.00841935\\
0.17	0.00816451\\
0.172	0.00790598\\
0.174	0.00764437\\
0.176	0.0073803\\
0.178	0.00711437\\
0.18	0.00684715\\
0.182	0.0065792\\
0.184	0.00631106\\
0.186	0.00604328\\
0.188	0.00577634\\
0.19	0.00559178\\
0.192	0.00544997\\
0.194	0.00531112\\
0.196	0.00517525\\
0.198	0.00504235\\
0.2	0.00491244\\
0.202	0.00478549\\
0.204	0.00466147\\
0.206	0.00454035\\
0.208	0.00442209\\
0.21	0.00430664\\
0.212	0.00419393\\
0.214	0.00408392\\
0.216	0.00397654\\
0.218	0.00387172\\
0.22	0.00376939\\
0.222	0.00366949\\
0.224	0.00357194\\
0.226	0.00347668\\
0.228	0.00338364\\
0.23	0.00329275\\
0.232	0.00320395\\
0.234	0.00311717\\
0.236	0.00303235\\
0.238	0.00294944\\
0.24	0.00286838\\
0.242	0.00278913\\
0.244	0.00271163\\
0.246	0.00263583\\
0.248	0.00256171\\
0.25	0.00248921\\
0.252	0.00241831\\
0.254	0.00234898\\
0.256	0.00228118\\
0.258	0.0022149\\
0.26	0.0021501\\
0.262	0.00208677\\
0.264	0.00202488\\
0.266	0.00196443\\
0.268	0.0019054\\
0.27	0.00184777\\
0.272	0.00179154\\
0.274	0.00173668\\
0.276	0.00168319\\
0.278	0.00163107\\
0.28	0.00158029\\
0.282	0.00153086\\
0.284	0.00148275\\
0.286	0.00143596\\
0.288	0.00139048\\
0.29	0.0013463\\
0.292	0.0013034\\
0.294	0.00126177\\
0.296	0.0012214\\
0.298	0.001242\\
0.3	0.00126005\\
0.302	0.00127397\\
0.304	0.00128391\\
0.306	0.00129003\\
0.308	0.00129247\\
0.31	0.0012914\\
0.312	0.00128698\\
0.314	0.00127937\\
0.316	0.00126875\\
0.318	0.00125527\\
0.32	0.00123912\\
0.322	0.00122046\\
0.324	0.00119947\\
0.326	0.00117631\\
0.328	0.00115114\\
0.33	0.00112415\\
0.332	0.00109549\\
0.334	0.00106533\\
0.336	0.00103382\\
0.338	0.00100112\\
0.34	0.000967376\\
0.342	0.000932743\\
0.344	0.000897363\\
0.346	0.000861374\\
0.348	0.000824909\\
0.35	0.000788098\\
0.352	0.000751065\\
0.354	0.000713928\\
0.356	0.0006768\\
0.358	0.000665564\\
0.36	0.000655043\\
0.362	0.000643381\\
0.364	0.000630662\\
0.366	0.000616965\\
0.368	0.000602373\\
0.37	0.000586965\\
0.372	0.000570819\\
0.374	0.000554013\\
0.376	0.000536624\\
0.378	0.000518724\\
0.38	0.000500386\\
0.382	0.000481681\\
0.384	0.000462677\\
0.386	0.000443438\\
0.388	0.00042403\\
0.39	0.000404512\\
0.392	0.000384943\\
0.394	0.000365379\\
0.396	0.000345873\\
0.398	0.000332842\\
0.4	0.000332741\\
0.402	0.000331949\\
0.404	0.000330493\\
0.406	0.000328403\\
0.408	0.000325709\\
0.41	0.000322961\\
0.412	0.000322054\\
0.414	0.000320443\\
0.416	0.000318162\\
0.418	0.000315248\\
0.42	0.000311738\\
0.422	0.000307668\\
0.424	0.000303073\\
0.426	0.00029799\\
0.428	0.000292455\\
0.43	0.000286502\\
0.432	0.000280165\\
0.434	0.00027348\\
0.436	0.000266479\\
0.438	0.000259195\\
0.44	0.00025166\\
0.442	0.000243905\\
0.444	0.000235959\\
0.446	0.000227853\\
0.448	0.000219615\\
0.45	0.000211271\\
0.452	0.000202849\\
0.454	0.000194373\\
0.456	0.000185868\\
0.458	0.000177358\\
0.46	0.000168863\\
0.462	0.000160406\\
0.464	0.000152006\\
0.466	0.000143682\\
0.468	0.000135451\\
0.47	0.000127331\\
0.472	0.000119337\\
0.474	0.000111483\\
0.476	0.000103782\\
0.478	9.62474e-05\\
0.48	8.88898e-05\\
0.482	8.56434e-05\\
0.484	8.49514e-05\\
0.486	8.43037e-05\\
0.488	8.33749e-05\\
0.49	8.2183e-05\\
0.492	8.07457e-05\\
0.494	7.90808e-05\\
0.496	7.72056e-05\\
0.498	7.51374e-05\\
0.5	7.28933e-05\\
0.502	7.049e-05\\
0.504	6.79439e-05\\
0.506	6.52709e-05\\
0.508	6.24868e-05\\
0.51	5.96068e-05\\
0.512	5.66456e-05\\
0.514	5.36175e-05\\
0.516	5.17194e-05\\
0.518	5.0133e-05\\
0.52	4.85858e-05\\
0.522	4.70773e-05\\
0.524	4.56069e-05\\
0.526	4.41741e-05\\
0.528	4.27784e-05\\
0.53	4.14194e-05\\
0.532	4.00966e-05\\
0.534	3.88096e-05\\
0.536	3.75582e-05\\
0.538	3.63418e-05\\
0.54	3.51601e-05\\
0.542	3.54819e-05\\
0.544	3.66907e-05\\
0.546	3.77395e-05\\
0.548	3.8633e-05\\
0.55	3.93758e-05\\
0.552	3.9973e-05\\
0.554	4.043e-05\\
0.556	4.0752e-05\\
0.558	4.09449e-05\\
0.56	4.10142e-05\\
0.562	4.09658e-05\\
0.564	4.08057e-05\\
0.566	4.054e-05\\
0.568	4.01745e-05\\
0.57	3.97155e-05\\
0.572	3.91689e-05\\
0.574	3.85408e-05\\
0.576	3.78371e-05\\
0.578	3.70639e-05\\
0.58	3.6227e-05\\
0.582	3.53321e-05\\
0.584	3.4385e-05\\
0.586	3.33911e-05\\
0.588	3.23558e-05\\
0.59	3.12845e-05\\
0.592	3.01823e-05\\
0.594	2.9054e-05\\
0.596	2.79046e-05\\
0.598	2.67385e-05\\
0.6	2.60952e-05\\
0.602	2.59855e-05\\
0.604	2.58108e-05\\
0.606	2.55749e-05\\
0.608	2.52813e-05\\
0.61	2.49336e-05\\
0.612	2.45354e-05\\
0.614	2.40905e-05\\
0.616	2.36022e-05\\
0.618	2.30742e-05\\
0.62	2.25098e-05\\
0.622	2.19125e-05\\
0.624	2.12855e-05\\
0.626	2.06322e-05\\
0.628	1.99556e-05\\
0.63	1.92589e-05\\
0.632	1.85451e-05\\
0.634	1.78169e-05\\
0.636	1.70772e-05\\
0.638	1.63286e-05\\
0.64	1.55738e-05\\
0.642	1.48151e-05\\
0.644	1.4055e-05\\
0.646	1.32957e-05\\
0.648	1.25392e-05\\
0.65	1.17876e-05\\
0.652	1.10428e-05\\
0.654	1.03066e-05\\
0.656	9.58064e-06\\
0.658	8.86649e-06\\
0.66	8.16561e-06\\
0.662	7.47933e-06\\
0.664	7.48468e-06\\
0.666	7.63146e-06\\
0.668	7.75282e-06\\
0.67	7.84961e-06\\
0.672	7.9227e-06\\
0.674	7.97299e-06\\
0.676	8.00144e-06\\
0.678	8.00902e-06\\
0.68	7.99671e-06\\
0.682	7.96554e-06\\
0.684	7.91652e-06\\
0.686	7.85069e-06\\
0.688	7.76909e-06\\
0.69	7.67275e-06\\
0.692	7.56269e-06\\
0.694	7.43994e-06\\
0.696	7.30551e-06\\
0.698	7.16039e-06\\
0.7	7.00556e-06\\
0.702	6.84195e-06\\
0.704	6.6705e-06\\
0.706	6.4921e-06\\
0.708	6.30763e-06\\
0.71	6.11791e-06\\
0.712	5.92376e-06\\
0.714	5.72593e-06\\
0.716	5.52517e-06\\
0.718	5.32217e-06\\
0.72	5.11758e-06\\
0.722	4.91205e-06\\
0.724	4.70614e-06\\
0.726	4.50042e-06\\
0.728	4.2954e-06\\
0.73	4.09155e-06\\
0.732	3.88932e-06\\
0.734	3.68911e-06\\
0.736	3.51589e-06\\
0.738	3.47294e-06\\
0.74	3.41978e-06\\
0.742	3.35714e-06\\
0.744	3.28575e-06\\
0.746	3.20633e-06\\
0.748	3.11959e-06\\
0.75	3.02622e-06\\
0.752	2.9269e-06\\
0.754	2.8223e-06\\
0.756	2.71307e-06\\
0.758	2.59983e-06\\
0.76	2.48319e-06\\
0.762	2.36374e-06\\
0.764	2.24205e-06\\
0.766	2.11865e-06\\
0.768	1.99406e-06\\
0.77	1.86879e-06\\
0.772	1.7433e-06\\
0.774	1.61803e-06\\
0.776	1.49341e-06\\
0.778	1.36983e-06\\
0.78	1.24767e-06\\
0.782	1.12727e-06\\
0.784	1.08549e-06\\
0.786	1.12806e-06\\
0.788	1.16529e-06\\
0.79	1.19731e-06\\
0.792	1.22425e-06\\
0.794	1.24624e-06\\
0.796	1.26344e-06\\
0.798	1.27601e-06\\
0.8	1.28411e-06\\
0.802	1.28793e-06\\
0.804	1.28764e-06\\
0.806	1.28344e-06\\
0.808	1.27553e-06\\
0.81	1.2641e-06\\
0.812	1.24935e-06\\
0.814	1.2315e-06\\
0.816	1.21074e-06\\
0.818	1.18728e-06\\
0.82	1.16132e-06\\
0.822	1.13308e-06\\
0.824	1.10275e-06\\
0.826	1.07053e-06\\
0.828	1.03663e-06\\
0.83	1.00123e-06\\
0.832	9.64529e-07\\
0.834	9.26706e-07\\
0.836	8.87943e-07\\
0.838	8.60601e-07\\
0.84	8.4825e-07\\
0.842	8.55363e-07\\
0.844	8.63596e-07\\
0.846	8.6891e-07\\
0.848	8.71435e-07\\
0.85	8.71305e-07\\
0.852	8.68652e-07\\
0.854	8.63611e-07\\
0.856	8.56317e-07\\
0.858	8.46906e-07\\
0.86	8.35511e-07\\
0.862	8.22267e-07\\
0.864	8.07306e-07\\
0.866	7.90759e-07\\
0.868	7.72755e-07\\
0.87	7.53422e-07\\
0.872	7.32885e-07\\
0.874	7.11265e-07\\
0.876	6.88682e-07\\
0.878	6.65252e-07\\
0.88	6.41088e-07\\
0.882	6.16299e-07\\
0.884	5.90991e-07\\
0.886	5.65267e-07\\
0.888	5.39225e-07\\
0.89	5.1296e-07\\
0.892	4.86563e-07\\
0.894	4.60119e-07\\
0.896	4.33711e-07\\
0.898	4.07418e-07\\
0.9	3.81314e-07\\
0.902	3.55468e-07\\
0.904	3.29947e-07\\
0.906	3.04811e-07\\
0.908	2.80118e-07\\
0.91	2.55922e-07\\
0.912	2.32272e-07\\
0.914	2.09212e-07\\
0.916	1.89797e-07\\
0.918	1.93434e-07\\
0.92	1.96109e-07\\
0.922	1.97863e-07\\
0.924	1.98742e-07\\
0.926	1.9879e-07\\
0.928	1.98051e-07\\
0.93	1.96571e-07\\
0.932	1.94397e-07\\
0.934	1.91572e-07\\
0.936	1.88142e-07\\
0.938	1.84153e-07\\
0.94	1.79649e-07\\
0.942	1.74673e-07\\
0.944	1.69268e-07\\
0.946	1.63478e-07\\
0.948	1.64014e-07\\
0.95	1.67273e-07\\
0.952	1.70001e-07\\
0.954	1.72217e-07\\
0.956	1.73943e-07\\
0.958	1.752e-07\\
0.96	1.76008e-07\\
0.962	1.76388e-07\\
0.964	1.76361e-07\\
0.966	1.75945e-07\\
0.968	1.7516e-07\\
0.97	1.74026e-07\\
0.972	1.72561e-07\\
0.974	1.70784e-07\\
0.976	1.68712e-07\\
0.978	1.66364e-07\\
0.98	1.63757e-07\\
0.982	1.60907e-07\\
0.984	1.57832e-07\\
0.986	1.54548e-07\\
0.988	1.51071e-07\\
0.99	1.47416e-07\\
0.992	1.436e-07\\
0.994	1.39638e-07\\
0.996	1.35543e-07\\
0.998	1.31331e-07\\
1	1.27016e-07\\
1.002	1.22613e-07\\
1.004	1.18133e-07\\
1.006	1.13592e-07\\
1.008	1.09002e-07\\
1.01	1.04375e-07\\
1.012	9.97246e-08\\
1.014	9.50618e-08\\
1.016	9.03984e-08\\
1.018	8.57457e-08\\
1.02	8.11144e-08\\
1.022	7.6515e-08\\
1.024	7.19573e-08\\
1.026	6.7451e-08\\
1.028	6.30051e-08\\
1.03	5.8628e-08\\
1.032	5.74001e-08\\
1.034	5.78323e-08\\
1.036	5.80598e-08\\
1.038	5.80916e-08\\
1.04	5.79371e-08\\
1.042	5.76055e-08\\
1.044	5.71066e-08\\
1.046	5.64501e-08\\
1.048	5.56457e-08\\
1.05	5.47033e-08\\
1.052	5.36329e-08\\
1.054	5.24441e-08\\
1.056	5.11469e-08\\
1.058	4.9751e-08\\
1.06	4.82657e-08\\
1.062	4.67006e-08\\
1.064	4.50649e-08\\
1.066	4.33675e-08\\
1.068	4.16172e-08\\
1.07	3.98225e-08\\
1.072	3.79918e-08\\
1.074	3.61328e-08\\
1.076	3.42534e-08\\
1.078	3.23608e-08\\
1.08	3.0462e-08\\
1.082	2.85639e-08\\
1.084	2.79593e-08\\
1.086	2.842e-08\\
1.088	2.87689e-08\\
1.09	2.90105e-08\\
1.092	2.91495e-08\\
1.094	2.91905e-08\\
1.096	2.91383e-08\\
1.098	2.89977e-08\\
1.1	2.87736e-08\\
1.102	2.84709e-08\\
1.104	2.80945e-08\\
1.106	2.76492e-08\\
1.108	2.71399e-08\\
1.11	2.65715e-08\\
1.112	2.59486e-08\\
1.114	2.5276e-08\\
1.116	2.45584e-08\\
1.118	2.38003e-08\\
1.12	2.3006e-08\\
1.122	2.21801e-08\\
1.124	2.13267e-08\\
1.126	2.04499e-08\\
1.128	1.95538e-08\\
1.13	1.90189e-08\\
1.132	1.85416e-08\\
1.134	1.80418e-08\\
1.136	1.75221e-08\\
1.138	1.6985e-08\\
1.14	1.64328e-08\\
1.142	1.58677e-08\\
1.144	1.52917e-08\\
1.146	1.47069e-08\\
1.148	1.41151e-08\\
1.15	1.35181e-08\\
1.152	1.29174e-08\\
1.154	1.23148e-08\\
1.156	1.17116e-08\\
1.158	1.11093e-08\\
1.16	1.05091e-08\\
1.162	1.00912e-08\\
1.164	1.00015e-08\\
1.166	9.87939e-09\\
1.168	9.72683e-09\\
1.17	9.546e-09\\
1.172	9.33899e-09\\
1.174	9.10791e-09\\
1.176	8.8548e-09\\
1.178	8.58172e-09\\
1.18	8.29064e-09\\
1.182	7.98352e-09\\
1.184	7.66228e-09\\
1.186	7.32875e-09\\
1.188	6.98477e-09\\
1.19	6.63206e-09\\
1.192	6.27232e-09\\
1.194	5.90719e-09\\
1.196	5.53822e-09\\
1.198	5.1669e-09\\
1.2	4.79468e-09\\
1.202	4.42292e-09\\
1.204	4.48964e-09\\
1.206	4.60348e-09\\
1.208	4.6936e-09\\
1.21	4.76098e-09\\
1.212	4.80664e-09\\
1.214	4.83162e-09\\
1.216	4.83697e-09\\
1.218	4.82376e-09\\
1.22	4.79306e-09\\
1.222	4.74595e-09\\
1.224	4.68351e-09\\
1.226	4.60681e-09\\
1.228	4.51692e-09\\
1.23	4.4149e-09\\
1.232	4.30177e-09\\
1.234	4.17857e-09\\
1.236	4.04629e-09\\
1.238	3.90592e-09\\
1.24	3.77507e-09\\
1.242	3.82621e-09\\
1.244	3.86502e-09\\
1.246	3.89199e-09\\
1.248	3.90764e-09\\
1.25	3.91247e-09\\
1.252	3.90704e-09\\
1.254	3.89188e-09\\
1.256	3.86753e-09\\
1.258	3.83454e-09\\
1.26	3.79346e-09\\
1.262	3.74483e-09\\
1.264	3.68921e-09\\
1.266	3.62712e-09\\
1.268	3.5591e-09\\
1.27	3.48567e-09\\
1.272	3.40733e-09\\
1.274	3.3246e-09\\
1.276	3.23793e-09\\
1.278	3.14782e-09\\
1.28	3.0547e-09\\
1.282	2.95902e-09\\
1.284	2.86118e-09\\
1.286	2.76159e-09\\
1.288	2.66063e-09\\
1.29	2.55865e-09\\
1.292	2.456e-09\\
1.294	2.353e-09\\
1.296	2.24993e-09\\
1.298	2.14709e-09\\
1.3	2.05397e-09\\
1.302	1.97274e-09\\
1.304	1.88921e-09\\
1.306	1.80379e-09\\
1.308	1.71686e-09\\
1.31	1.62877e-09\\
1.312	1.54215e-09\\
1.314	1.50311e-09\\
1.316	1.46276e-09\\
1.318	1.42122e-09\\
1.32	1.37865e-09\\
1.322	1.33515e-09\\
1.324	1.29086e-09\\
1.326	1.24591e-09\\
1.328	1.20797e-09\\
1.33	1.21241e-09\\
1.332	1.21275e-09\\
1.334	1.20916e-09\\
1.336	1.20186e-09\\
1.338	1.19103e-09\\
1.34	1.17689e-09\\
1.342	1.15963e-09\\
1.344	1.13945e-09\\
1.346	1.11657e-09\\
1.348	1.09117e-09\\
1.35	1.06347e-09\\
1.352	1.03367e-09\\
1.354	1.00196e-09\\
1.356	9.68527e-10\\
1.358	9.33574e-10\\
1.36	8.97285e-10\\
1.362	8.59842e-10\\
1.364	8.21423e-10\\
1.366	7.82202e-10\\
1.368	7.42347e-10\\
1.37	7.0202e-10\\
1.372	6.61377e-10\\
1.374	6.20569e-10\\
1.376	5.79738e-10\\
1.378	5.39023e-10\\
1.38	4.98554e-10\\
1.382	4.81265e-10\\
1.384	4.91754e-10\\
1.386	4.99959e-10\\
1.388	5.05972e-10\\
1.39	5.09884e-10\\
1.392	5.11791e-10\\
1.394	5.11789e-10\\
1.396	5.09976e-10\\
1.398	5.06452e-10\\
1.4	5.01316e-10\\
1.402	4.96327e-10\\
1.404	4.9473e-10\\
1.406	4.91908e-10\\
1.408	4.87925e-10\\
1.41	4.82842e-10\\
1.412	4.76724e-10\\
1.414	4.69634e-10\\
1.416	4.61637e-10\\
1.418	4.52795e-10\\
1.42	4.43174e-10\\
1.422	4.32835e-10\\
1.424	4.21843e-10\\
1.426	4.10259e-10\\
1.428	3.98144e-10\\
1.43	3.85559e-10\\
1.432	3.72562e-10\\
1.434	3.59212e-10\\
1.436	3.45563e-10\\
1.438	3.31672e-10\\
1.44	3.17589e-10\\
1.442	3.03366e-10\\
1.444	2.95378e-10\\
1.446	2.92616e-10\\
1.448	2.89244e-10\\
1.45	2.85298e-10\\
1.452	2.80812e-10\\
1.454	2.7582e-10\\
1.456	2.70357e-10\\
1.458	2.64457e-10\\
1.46	2.58154e-10\\
1.462	2.5148e-10\\
1.464	2.44469e-10\\
1.466	2.37152e-10\\
1.468	2.29561e-10\\
1.47	2.21727e-10\\
1.472	2.1368e-10\\
1.474	2.0545e-10\\
1.476	1.97064e-10\\
1.478	1.88551e-10\\
1.48	1.79938e-10\\
1.482	1.71249e-10\\
1.484	1.62511e-10\\
1.486	1.53747e-10\\
1.488	1.44979e-10\\
1.49	1.36231e-10\\
1.492	1.27523e-10\\
1.494	1.18874e-10\\
1.496	1.14676e-10\\
1.498	1.17708e-10\\
1.5	1.20209e-10\\
1.502	1.22195e-10\\
1.504	1.23686e-10\\
1.506	1.247e-10\\
1.508	1.25259e-10\\
1.51	1.25383e-10\\
1.512	1.25092e-10\\
1.514	1.2441e-10\\
1.516	1.23356e-10\\
1.518	1.21954e-10\\
1.52	1.20225e-10\\
1.522	1.18193e-10\\
1.524	1.15877e-10\\
1.526	1.13302e-10\\
1.528	1.10488e-10\\
1.53	1.07456e-10\\
1.532	1.04228e-10\\
1.534	1.00825e-10\\
1.536	9.72667e-11\\
1.538	9.35727e-11\\
1.54	8.97622e-11\\
1.542	8.5854e-11\\
1.544	8.18762e-11\\
1.546	8.17398e-11\\
1.548	8.14315e-11\\
1.55	8.09597e-11\\
1.552	8.03324e-11\\
1.554	7.95576e-11\\
1.556	7.86431e-11\\
1.558	7.75966e-11\\
1.56	7.64257e-11\\
1.562	7.51381e-11\\
1.564	7.37411e-11\\
1.566	7.22421e-11\\
1.568	7.06484e-11\\
1.57	6.89671e-11\\
1.572	6.72053e-11\\
1.574	6.53701e-11\\
1.576	6.53276e-11\\
1.578	6.52386e-11\\
1.58	6.49925e-11\\
1.582	6.45969e-11\\
1.584	6.4059e-11\\
1.586	6.33863e-11\\
1.588	6.25864e-11\\
1.59	6.16668e-11\\
1.592	6.06348e-11\\
1.594	5.94981e-11\\
1.596	5.82639e-11\\
1.598	5.69397e-11\\
1.6	5.55327e-11\\
1.602	5.40501e-11\\
1.604	5.2499e-11\\
1.606	5.08862e-11\\
1.608	4.92186e-11\\
1.61	4.75029e-11\\
1.612	4.57455e-11\\
1.614	4.39528e-11\\
1.616	4.2131e-11\\
1.618	4.0286e-11\\
1.62	3.84235e-11\\
1.622	3.65492e-11\\
1.624	3.46684e-11\\
1.626	3.27863e-11\\
1.628	3.09078e-11\\
1.63	2.90375e-11\\
1.632	2.71801e-11\\
1.634	2.53396e-11\\
1.636	2.40102e-11\\
1.638	2.37621e-11\\
1.64	2.3445e-11\\
1.642	2.30634e-11\\
1.644	2.2622e-11\\
1.646	2.21251e-11\\
1.648	2.15774e-11\\
1.65	2.09832e-11\\
1.652	2.03469e-11\\
1.654	1.96728e-11\\
1.656	1.8965e-11\\
1.658	1.82277e-11\\
1.66	1.74648e-11\\
1.662	1.66802e-11\\
1.664	1.58777e-11\\
1.666	1.50609e-11\\
1.668	1.42332e-11\\
1.67	1.33981e-11\\
1.672	1.25588e-11\\
1.674	1.20006e-11\\
1.676	1.24349e-11\\
1.678	1.28079e-11\\
1.68	1.31213e-11\\
1.682	1.33772e-11\\
1.684	1.35776e-11\\
1.686	1.37247e-11\\
1.688	1.38205e-11\\
1.69	1.38674e-11\\
1.692	1.38677e-11\\
1.694	1.38237e-11\\
1.696	1.37377e-11\\
1.698	1.36121e-11\\
1.7	1.34493e-11\\
1.702	1.32516e-11\\
1.704	1.30213e-11\\
1.706	1.27609e-11\\
1.708	1.24727e-11\\
1.71	1.21589e-11\\
1.712	1.18664e-11\\
1.714	1.18822e-11\\
1.716	1.18652e-11\\
1.718	1.1817e-11\\
1.72	1.17392e-11\\
1.722	1.16334e-11\\
1.724	1.15014e-11\\
1.726	1.13446e-11\\
1.728	1.11648e-11\\
1.73	1.09635e-11\\
1.732	1.07423e-11\\
1.734	1.05027e-11\\
1.736	1.02463e-11\\
1.738	9.97458e-12\\
1.74	9.68898e-12\\
1.742	9.39092e-12\\
1.744	9.08184e-12\\
1.746	8.76306e-12\\
1.748	8.43592e-12\\
1.75	8.10171e-12\\
1.752	7.76164e-12\\
1.754	7.41694e-12\\
1.756	7.06871e-12\\
1.758	6.71808e-12\\
1.76	6.3661e-12\\
1.762	6.0138e-12\\
1.764	5.66213e-12\\
1.766	5.31201e-12\\
1.768	4.9643e-12\\
1.77	4.64816e-12\\
1.772	4.50361e-12\\
1.774	4.35178e-12\\
1.776	4.19357e-12\\
1.778	4.02982e-12\\
1.78	3.86137e-12\\
1.782	3.68903e-12\\
1.784	3.51355e-12\\
1.786	3.33567e-12\\
1.788	3.15612e-12\\
1.79	2.97558e-12\\
1.792	2.8238e-12\\
1.794	2.69265e-12\\
1.796	2.74065e-12\\
1.798	2.84292e-12\\
1.8	2.93177e-12\\
1.802	3.00757e-12\\
1.804	3.07069e-12\\
1.806	3.12151e-12\\
1.808	3.16044e-12\\
1.81	3.18796e-12\\
1.812	3.20448e-12\\
1.814	3.21051e-12\\
1.816	3.20649e-12\\
1.818	3.19288e-12\\
1.82	3.17023e-12\\
1.822	3.13899e-12\\
1.824	3.09969e-12\\
1.826	3.05282e-12\\
1.828	2.9989e-12\\
1.83	2.93841e-12\\
1.832	2.87189e-12\\
1.834	2.79978e-12\\
1.836	2.72262e-12\\
1.838	2.64088e-12\\
1.84	2.55505e-12\\
1.842	2.46558e-12\\
1.844	2.37291e-12\\
1.846	2.2775e-12\\
1.848	2.2515e-12\\
1.85	2.25123e-12\\
1.852	2.24491e-12\\
1.854	2.23285e-12\\
1.856	2.21536e-12\\
1.858	2.19276e-12\\
1.86	2.1654e-12\\
1.862	2.13358e-12\\
1.864	2.09758e-12\\
1.866	2.05774e-12\\
1.868	2.01434e-12\\
1.87	1.96769e-12\\
1.872	1.91808e-12\\
1.874	1.86581e-12\\
1.876	1.81112e-12\\
1.878	1.75431e-12\\
1.88	1.69563e-12\\
1.882	1.63534e-12\\
1.884	1.57369e-12\\
1.886	1.5109e-12\\
1.888	1.44722e-12\\
1.89	1.38288e-12\\
1.892	1.31807e-12\\
1.894	1.25299e-12\\
1.896	1.18784e-12\\
1.898	1.12278e-12\\
1.9	1.05802e-12\\
1.902	9.93685e-13\\
1.904	9.29952e-13\\
1.906	8.66955e-13\\
1.908	8.04826e-13\\
1.91	7.43698e-13\\
1.912	6.83696e-13\\
1.914	6.24917e-13\\
1.916	5.72467e-13\\
1.918	5.49964e-13\\
1.92	5.30952e-13\\
1.922	5.54037e-13\\
1.924	5.74265e-13\\
1.926	5.91705e-13\\
1.928	6.06463e-13\\
1.93	6.18633e-13\\
1.932	6.28308e-13\\
1.934	6.35585e-13\\
1.936	6.40555e-13\\
1.938	6.43322e-13\\
1.94	6.4397e-13\\
1.942	6.42596e-13\\
1.944	6.39342e-13\\
1.946	6.343e-13\\
1.948	6.27568e-13\\
1.95	6.19247e-13\\
1.952	6.09461e-13\\
1.954	5.98287e-13\\
1.956	5.85856e-13\\
1.958	5.7225e-13\\
1.96	5.57531e-13\\
1.962	5.41865e-13\\
1.964	5.25324e-13\\
1.966	5.07972e-13\\
1.968	4.89938e-13\\
1.97	4.71307e-13\\
1.972	4.52167e-13\\
1.974	4.32569e-13\\
1.976	4.12611e-13\\
1.978	3.92374e-13\\
1.98	3.71942e-13\\
1.982	3.66749e-13\\
1.984	3.68304e-13\\
1.986	3.68768e-13\\
1.988	3.68187e-13\\
1.99	3.66636e-13\\
1.992	3.64149e-13\\
1.994	3.60779e-13\\
1.996	3.56591e-13\\
1.998	3.51646e-13\\
2	3.45963e-13\\
};
\addplot [color=mycolor6, line width=1.0pt, forget plot]
  table[row sep=crcr]{%
-0.024	0\\
-0.022	0\\
-0.02	0\\
-0.018	0\\
-0.016	0\\
-0.014	0\\
-0.012	0\\
-0.01	0\\
-0.008	0\\
-0.006	0\\
-0.004	0\\
-0.002	0\\
0	0\\
0.002	0.00916182\\
0.004	0.0147064\\
0.006	0.0198638\\
0.008	0.0237502\\
0.01	0.0264851\\
0.012	0.0282134\\
0.014	0.0290841\\
0.016	0.0292463\\
0.018	0.0288985\\
0.02	0.0290636\\
0.022	0.0291689\\
0.024	0.0292335\\
0.026	0.0294184\\
0.028	0.0299268\\
0.03	0.0301269\\
0.032	0.0300265\\
0.034	0.0296386\\
0.036	0.0306076\\
0.038	0.0315397\\
0.04	0.0323726\\
0.042	0.0331068\\
0.044	0.0337429\\
0.046	0.0342819\\
0.048	0.0347251\\
0.05	0.035074\\
0.052	0.0353303\\
0.054	0.0354962\\
0.056	0.0355738\\
0.058	0.0355656\\
0.06	0.0354745\\
0.062	0.0353031\\
0.064	0.0350548\\
0.066	0.0347326\\
0.068	0.03434\\
0.07	0.0338805\\
0.072	0.0333579\\
0.074	0.0327759\\
0.076	0.0321384\\
0.078	0.0314493\\
0.08	0.0307126\\
0.082	0.0299324\\
0.084	0.0291127\\
0.086	0.0282575\\
0.088	0.0273708\\
0.09	0.0264566\\
0.092	0.0255189\\
0.094	0.0245615\\
0.096	0.0235881\\
0.098	0.0226024\\
0.1	0.021608\\
0.102	0.0206082\\
0.104	0.0196063\\
0.106	0.0186055\\
0.108	0.0176088\\
0.11	0.0166189\\
0.112	0.0156386\\
0.114	0.0146704\\
0.116	0.0140279\\
0.118	0.0136483\\
0.12	0.0132584\\
0.122	0.0128595\\
0.124	0.012453\\
0.126	0.0120403\\
0.128	0.0116226\\
0.13	0.0112012\\
0.132	0.0107772\\
0.134	0.0105755\\
0.136	0.0104944\\
0.138	0.0104109\\
0.14	0.010315\\
0.142	0.0102071\\
0.144	0.0100878\\
0.146	0.00995759\\
0.148	0.00981696\\
0.15	0.0096665\\
0.152	0.00950673\\
0.154	0.00933822\\
0.156	0.00916151\\
0.158	0.00897717\\
0.16	0.00878574\\
0.162	0.00858781\\
0.164	0.00838391\\
0.166	0.00817461\\
0.168	0.00796047\\
0.17	0.00774204\\
0.172	0.00751985\\
0.174	0.00729446\\
0.176	0.00706638\\
0.178	0.00683613\\
0.18	0.00660424\\
0.182	0.00637118\\
0.184	0.00613745\\
0.186	0.00596436\\
0.188	0.00585881\\
0.19	0.00575443\\
0.192	0.00565126\\
0.194	0.00554931\\
0.196	0.00544861\\
0.198	0.00534916\\
0.2	0.00525099\\
0.202	0.00515411\\
0.204	0.00505851\\
0.206	0.00496419\\
0.208	0.00487117\\
0.21	0.00477944\\
0.212	0.004689\\
0.214	0.00459984\\
0.216	0.00451194\\
0.218	0.00442532\\
0.22	0.00433995\\
0.222	0.00425584\\
0.224	0.00417295\\
0.226	0.0040913\\
0.228	0.00401086\\
0.23	0.00393162\\
0.232	0.00385358\\
0.234	0.00377671\\
0.236	0.00370102\\
0.238	0.00362648\\
0.24	0.00355309\\
0.242	0.00348083\\
0.244	0.00340969\\
0.246	0.00333967\\
0.248	0.00327076\\
0.25	0.00320294\\
0.252	0.0031362\\
0.254	0.00307054\\
0.256	0.00300596\\
0.258	0.00294243\\
0.26	0.00287996\\
0.262	0.00281854\\
0.264	0.00275817\\
0.266	0.00269883\\
0.268	0.00264053\\
0.27	0.00258325\\
0.272	0.002527\\
0.274	0.00247177\\
0.276	0.00241755\\
0.278	0.00236434\\
0.28	0.00231213\\
0.282	0.00226092\\
0.284	0.0022107\\
0.286	0.00216147\\
0.288	0.00211323\\
0.29	0.00206596\\
0.292	0.00201966\\
0.294	0.00197432\\
0.296	0.00192994\\
0.298	0.00188651\\
0.3	0.00184401\\
0.302	0.00180244\\
0.304	0.00176179\\
0.306	0.00172205\\
0.308	0.00168321\\
0.31	0.00164525\\
0.312	0.00160816\\
0.314	0.00157194\\
0.316	0.00153657\\
0.318	0.00150202\\
0.32	0.0014683\\
0.322	0.00143539\\
0.324	0.00140326\\
0.326	0.00137191\\
0.328	0.00134133\\
0.33	0.00131148\\
0.332	0.00128236\\
0.334	0.00125396\\
0.336	0.00122626\\
0.338	0.00119923\\
0.34	0.00117287\\
0.342	0.00114715\\
0.344	0.00112207\\
0.346	0.0010976\\
0.348	0.00107373\\
0.35	0.00105045\\
0.352	0.00102774\\
0.354	0.00100557\\
0.356	0.00098395\\
0.358	0.000962849\\
0.36	0.000942257\\
0.362	0.00092216\\
0.364	0.000902543\\
0.366	0.000883393\\
0.368	0.000864696\\
0.37	0.00084644\\
0.372	0.000828612\\
0.374	0.0008112\\
0.376	0.000794191\\
0.378	0.000777575\\
0.38	0.000761341\\
0.382	0.000745476\\
0.384	0.000729971\\
0.386	0.000714816\\
0.388	0.000700001\\
0.39	0.000685517\\
0.392	0.000671353\\
0.394	0.000657502\\
0.396	0.000643954\\
0.398	0.000630702\\
0.4	0.000617738\\
0.402	0.000605054\\
0.404	0.000592642\\
0.406	0.000580495\\
0.408	0.000568607\\
0.41	0.000556971\\
0.412	0.000545579\\
0.414	0.000534427\\
0.416	0.000523509\\
0.418	0.000512817\\
0.42	0.000502347\\
0.422	0.000492094\\
0.424	0.000482052\\
0.426	0.000472215\\
0.428	0.00046258\\
0.43	0.000453142\\
0.432	0.000443895\\
0.434	0.000434836\\
0.436	0.00042596\\
0.438	0.000417263\\
0.44	0.000408741\\
0.442	0.00040039\\
0.444	0.000392206\\
0.446	0.000384186\\
0.448	0.000376326\\
0.45	0.000368622\\
0.452	0.000361071\\
0.454	0.00035367\\
0.456	0.000346415\\
0.458	0.000339304\\
0.46	0.000332333\\
0.462	0.0003255\\
0.464	0.000318801\\
0.466	0.000312234\\
0.468	0.000305795\\
0.47	0.000299483\\
0.472	0.000293295\\
0.474	0.000287227\\
0.476	0.000281279\\
0.478	0.000275447\\
0.48	0.000269728\\
0.482	0.000264121\\
0.484	0.000258624\\
0.486	0.000253234\\
0.488	0.000247949\\
0.49	0.000242768\\
0.492	0.000237687\\
0.494	0.000232706\\
0.496	0.000227822\\
0.498	0.000223033\\
0.5	0.000218338\\
0.502	0.000213735\\
0.504	0.000209222\\
0.506	0.000204798\\
0.508	0.00020046\\
0.51	0.000196208\\
0.512	0.00019204\\
0.514	0.000187954\\
0.516	0.000183948\\
0.518	0.000180022\\
0.52	0.000176175\\
0.522	0.000172403\\
0.524	0.000168707\\
0.526	0.000165085\\
0.528	0.000161536\\
0.53	0.000158058\\
0.532	0.00015465\\
0.534	0.000151312\\
0.536	0.000148041\\
0.538	0.000144837\\
0.54	0.000141698\\
0.542	0.000138624\\
0.544	0.000135613\\
0.546	0.000132664\\
0.548	0.000129776\\
0.55	0.000126948\\
0.552	0.00012418\\
0.554	0.000121469\\
0.556	0.000118815\\
0.558	0.000116217\\
0.56	0.000113675\\
0.562	0.000111186\\
0.564	0.00010875\\
0.566	0.000106366\\
0.568	0.000104034\\
0.57	0.000101751\\
0.572	9.95179e-05\\
0.574	9.7333e-05\\
0.576	9.51954e-05\\
0.578	9.31044e-05\\
0.58	9.1059e-05\\
0.582	8.90583e-05\\
0.584	8.71015e-05\\
0.586	8.51877e-05\\
0.588	8.33161e-05\\
0.59	8.14857e-05\\
0.592	7.96959e-05\\
0.594	7.79457e-05\\
0.596	7.62343e-05\\
0.598	7.4561e-05\\
0.6	7.2925e-05\\
0.602	7.13254e-05\\
0.604	6.97615e-05\\
0.606	6.82326e-05\\
0.608	6.67379e-05\\
0.61	6.52767e-05\\
0.612	6.38482e-05\\
0.614	6.24517e-05\\
0.616	6.10866e-05\\
0.618	5.97521e-05\\
0.62	5.84475e-05\\
0.622	5.71723e-05\\
0.624	5.59257e-05\\
0.626	5.47071e-05\\
0.628	5.35159e-05\\
0.63	5.23514e-05\\
0.632	5.1213e-05\\
0.634	5.01002e-05\\
0.636	4.90124e-05\\
0.638	4.7949e-05\\
0.64	4.69094e-05\\
0.642	4.58931e-05\\
0.644	4.48995e-05\\
0.646	4.39282e-05\\
0.648	4.29786e-05\\
0.65	4.20502e-05\\
0.652	4.11425e-05\\
0.654	4.0255e-05\\
0.656	3.93873e-05\\
0.658	3.8539e-05\\
0.66	3.77094e-05\\
0.662	3.68983e-05\\
0.664	3.61052e-05\\
0.666	3.53297e-05\\
0.668	3.45714e-05\\
0.67	3.38298e-05\\
0.672	3.31046e-05\\
0.674	3.23954e-05\\
0.676	3.17018e-05\\
0.678	3.10235e-05\\
0.68	3.03601e-05\\
0.682	2.97113e-05\\
0.684	2.90767e-05\\
0.686	2.8456e-05\\
0.688	2.78489e-05\\
0.69	2.72551e-05\\
0.692	2.66742e-05\\
0.694	2.6106e-05\\
0.696	2.55502e-05\\
0.698	2.50065e-05\\
0.7	2.44745e-05\\
0.702	2.39542e-05\\
0.704	2.34451e-05\\
0.706	2.2947e-05\\
0.708	2.24597e-05\\
0.71	2.19829e-05\\
0.712	2.15164e-05\\
0.714	2.106e-05\\
0.716	2.06133e-05\\
0.718	2.01763e-05\\
0.72	1.97487e-05\\
0.722	1.93302e-05\\
0.724	1.89207e-05\\
0.726	1.85199e-05\\
0.728	1.81278e-05\\
0.73	1.77439e-05\\
0.732	1.73683e-05\\
0.734	1.70007e-05\\
0.736	1.66409e-05\\
0.738	1.62887e-05\\
0.74	1.5944e-05\\
0.742	1.56066e-05\\
0.744	1.52764e-05\\
0.746	1.49531e-05\\
0.748	1.46367e-05\\
0.75	1.4327e-05\\
0.752	1.40238e-05\\
0.754	1.3727e-05\\
0.756	1.34364e-05\\
0.758	1.3152e-05\\
0.76	1.28735e-05\\
0.762	1.26009e-05\\
0.764	1.2334e-05\\
0.766	1.20728e-05\\
0.768	1.1817e-05\\
0.77	1.15665e-05\\
0.772	1.13213e-05\\
0.774	1.10813e-05\\
0.776	1.08462e-05\\
0.778	1.06161e-05\\
0.78	1.03908e-05\\
0.782	1.01702e-05\\
0.784	9.9542e-06\\
0.786	9.74272e-06\\
0.788	9.53565e-06\\
0.79	9.33292e-06\\
0.792	9.13441e-06\\
0.794	8.94005e-06\\
0.796	8.74975e-06\\
0.798	8.56343e-06\\
0.8	8.38099e-06\\
0.802	8.20236e-06\\
0.804	8.02747e-06\\
0.806	7.85622e-06\\
0.808	7.68856e-06\\
0.81	7.5244e-06\\
0.812	7.36367e-06\\
0.814	7.20631e-06\\
0.816	7.05224e-06\\
0.818	6.90139e-06\\
0.82	6.7537e-06\\
0.822	6.60911e-06\\
0.824	6.46754e-06\\
0.826	6.32895e-06\\
0.828	6.19326e-06\\
0.83	6.06043e-06\\
0.832	5.93038e-06\\
0.834	5.80307e-06\\
0.836	5.67843e-06\\
0.838	5.55642e-06\\
0.84	5.43698e-06\\
0.842	5.32006e-06\\
0.844	5.2056e-06\\
0.846	5.09356e-06\\
0.848	4.98389e-06\\
0.85	4.87654e-06\\
0.852	4.77145e-06\\
0.854	4.6686e-06\\
0.856	4.56792e-06\\
0.858	4.46938e-06\\
0.86	4.37292e-06\\
0.862	4.27852e-06\\
0.864	4.18613e-06\\
0.866	4.09569e-06\\
0.868	4.00719e-06\\
0.87	3.92057e-06\\
0.872	3.8358e-06\\
0.874	3.75284e-06\\
0.876	3.67165e-06\\
0.878	3.59219e-06\\
0.88	3.51444e-06\\
0.882	3.43835e-06\\
0.884	3.36389e-06\\
0.886	3.29103e-06\\
0.888	3.21973e-06\\
0.89	3.14997e-06\\
0.892	3.0817e-06\\
0.894	3.01491e-06\\
0.896	2.94955e-06\\
0.898	2.8856e-06\\
0.9	2.82303e-06\\
0.902	2.76181e-06\\
0.904	2.70191e-06\\
0.906	2.64331e-06\\
0.908	2.58597e-06\\
0.91	2.52988e-06\\
0.912	2.475e-06\\
0.914	2.42131e-06\\
0.916	2.36878e-06\\
0.918	2.3174e-06\\
0.92	2.26713e-06\\
0.922	2.21795e-06\\
0.924	2.16984e-06\\
0.926	2.12278e-06\\
0.928	2.07675e-06\\
0.93	2.03172e-06\\
0.932	1.98767e-06\\
0.934	1.94458e-06\\
0.936	1.90243e-06\\
0.938	1.8612e-06\\
0.94	1.82087e-06\\
0.942	1.78142e-06\\
0.944	1.74284e-06\\
0.946	1.7051e-06\\
0.948	1.66818e-06\\
0.95	1.63207e-06\\
0.952	1.59676e-06\\
0.954	1.56221e-06\\
0.956	1.52843e-06\\
0.958	1.49538e-06\\
0.96	1.46306e-06\\
0.962	1.43145e-06\\
0.964	1.40053e-06\\
0.966	1.37028e-06\\
0.968	1.3407e-06\\
0.97	1.31177e-06\\
0.972	1.28348e-06\\
0.974	1.2558e-06\\
0.976	1.22873e-06\\
0.978	1.20226e-06\\
0.98	1.17636e-06\\
0.982	1.15104e-06\\
0.984	1.12627e-06\\
0.986	1.10204e-06\\
0.988	1.07834e-06\\
0.99	1.05516e-06\\
0.992	1.03249e-06\\
0.994	1.01031e-06\\
0.996	9.88618e-07\\
0.998	9.67401e-07\\
1	9.46648e-07\\
1.002	9.26349e-07\\
1.004	9.06492e-07\\
1.006	8.87069e-07\\
1.008	8.68069e-07\\
1.01	8.49483e-07\\
1.012	8.31303e-07\\
1.014	8.13518e-07\\
1.016	7.96119e-07\\
1.018	7.79099e-07\\
1.02	7.62449e-07\\
1.022	7.4616e-07\\
1.024	7.30224e-07\\
1.026	7.14634e-07\\
1.028	6.99381e-07\\
1.03	6.84458e-07\\
1.032	6.69858e-07\\
1.034	6.55573e-07\\
1.036	6.41596e-07\\
1.038	6.27921e-07\\
1.04	6.1454e-07\\
1.042	6.01447e-07\\
1.044	5.88636e-07\\
1.046	5.761e-07\\
1.048	5.63833e-07\\
1.05	5.51829e-07\\
1.052	5.40083e-07\\
1.054	5.28588e-07\\
1.056	5.17338e-07\\
1.058	5.0633e-07\\
1.06	4.95556e-07\\
1.062	4.85012e-07\\
1.064	4.74693e-07\\
1.066	4.64594e-07\\
1.068	4.54709e-07\\
1.07	4.45035e-07\\
1.072	4.35566e-07\\
1.074	4.26299e-07\\
1.076	4.17228e-07\\
1.078	4.08349e-07\\
1.08	3.99659e-07\\
1.082	3.91153e-07\\
1.084	3.82826e-07\\
1.086	3.74676e-07\\
1.088	3.66698e-07\\
1.09	3.58888e-07\\
1.092	3.51244e-07\\
1.094	3.4376e-07\\
1.096	3.36435e-07\\
1.098	3.29264e-07\\
1.1	3.22244e-07\\
1.102	3.15372e-07\\
1.104	3.08645e-07\\
1.106	3.0206e-07\\
1.108	2.95613e-07\\
1.11	2.89302e-07\\
1.112	2.83123e-07\\
1.114	2.77075e-07\\
1.116	2.71154e-07\\
1.118	2.65358e-07\\
1.12	2.59683e-07\\
1.122	2.54128e-07\\
1.124	2.4869e-07\\
1.126	2.43367e-07\\
1.128	2.38155e-07\\
1.13	2.33053e-07\\
1.132	2.28058e-07\\
1.134	2.23169e-07\\
1.136	2.18383e-07\\
1.138	2.13698e-07\\
1.14	2.09111e-07\\
1.142	2.04621e-07\\
1.144	2.00226e-07\\
1.146	1.95924e-07\\
1.148	1.91712e-07\\
1.15	1.8759e-07\\
1.152	1.83554e-07\\
1.154	1.79604e-07\\
1.156	1.75738e-07\\
1.158	1.71953e-07\\
1.16	1.68249e-07\\
1.162	1.64623e-07\\
1.164	1.61075e-07\\
1.166	1.57601e-07\\
1.168	1.54202e-07\\
1.17	1.50874e-07\\
1.172	1.47618e-07\\
1.174	1.44431e-07\\
1.176	1.41312e-07\\
1.178	1.38259e-07\\
1.18	1.35272e-07\\
1.182	1.32348e-07\\
1.184	1.29487e-07\\
1.186	1.26687e-07\\
1.188	1.23948e-07\\
1.19	1.21266e-07\\
1.192	1.18643e-07\\
1.194	1.16076e-07\\
1.196	1.13564e-07\\
1.198	1.11106e-07\\
1.2	1.08701e-07\\
1.202	1.06348e-07\\
1.204	1.04045e-07\\
1.206	1.01792e-07\\
1.208	9.95882e-08\\
1.21	9.74318e-08\\
1.212	9.53219e-08\\
1.214	9.32578e-08\\
1.216	9.12384e-08\\
1.218	8.92627e-08\\
1.22	8.73299e-08\\
1.222	8.54391e-08\\
1.224	8.35893e-08\\
1.226	8.17798e-08\\
1.228	8.00096e-08\\
1.23	7.82779e-08\\
1.232	7.65839e-08\\
1.234	7.49268e-08\\
1.236	7.33058e-08\\
1.238	7.17202e-08\\
1.24	7.01691e-08\\
1.242	6.86519e-08\\
1.244	6.71678e-08\\
1.246	6.57161e-08\\
1.248	6.42961e-08\\
1.25	6.29072e-08\\
1.252	6.15485e-08\\
1.254	6.02196e-08\\
1.256	5.89197e-08\\
1.258	5.76483e-08\\
1.26	5.64046e-08\\
1.262	5.51881e-08\\
1.264	5.39982e-08\\
1.266	5.28343e-08\\
1.268	5.16959e-08\\
1.27	5.05824e-08\\
1.272	4.94931e-08\\
1.274	4.84277e-08\\
1.276	4.73856e-08\\
1.278	4.63662e-08\\
1.28	4.53691e-08\\
1.282	4.43938e-08\\
1.284	4.34397e-08\\
1.286	4.25065e-08\\
1.288	4.15936e-08\\
1.29	4.07006e-08\\
1.292	3.98271e-08\\
1.294	3.89726e-08\\
1.296	3.81367e-08\\
1.298	3.7319e-08\\
1.3	3.6519e-08\\
1.302	3.57365e-08\\
1.304	3.49709e-08\\
1.306	3.4222e-08\\
1.308	3.34893e-08\\
1.31	3.27725e-08\\
1.312	3.20713e-08\\
1.314	3.13852e-08\\
1.316	3.07139e-08\\
1.318	3.00572e-08\\
1.32	2.94147e-08\\
1.322	2.8786e-08\\
1.324	2.81709e-08\\
1.326	2.75691e-08\\
1.328	2.69802e-08\\
1.33	2.64041e-08\\
1.332	2.58403e-08\\
1.334	2.52886e-08\\
1.336	2.47488e-08\\
1.338	2.42206e-08\\
1.34	2.37037e-08\\
1.342	2.31979e-08\\
1.344	2.27029e-08\\
1.346	2.22185e-08\\
1.348	2.17445e-08\\
1.35	2.12806e-08\\
1.352	2.08266e-08\\
1.354	2.03824e-08\\
1.356	1.99475e-08\\
1.358	1.9522e-08\\
1.36	1.91056e-08\\
1.362	1.8698e-08\\
1.364	1.8299e-08\\
1.366	1.79086e-08\\
1.368	1.75265e-08\\
1.37	1.71525e-08\\
1.372	1.67864e-08\\
1.374	1.64281e-08\\
1.376	1.60775e-08\\
1.378	1.57342e-08\\
1.38	1.53983e-08\\
1.382	1.50694e-08\\
1.384	1.47476e-08\\
1.386	1.44325e-08\\
1.388	1.41241e-08\\
1.39	1.38223e-08\\
1.392	1.35268e-08\\
1.394	1.32377e-08\\
1.396	1.29546e-08\\
1.398	1.26775e-08\\
1.4	1.24063e-08\\
1.402	1.21408e-08\\
1.404	1.1881e-08\\
1.406	1.16266e-08\\
1.408	1.13776e-08\\
1.41	1.11339e-08\\
1.412	1.08954e-08\\
1.414	1.06619e-08\\
1.416	1.04333e-08\\
1.418	1.02096e-08\\
1.42	9.99067e-09\\
1.422	9.77634e-09\\
1.424	9.56655e-09\\
1.426	9.36121e-09\\
1.428	9.16023e-09\\
1.43	8.96351e-09\\
1.432	8.77096e-09\\
1.434	8.58251e-09\\
1.436	8.39805e-09\\
1.438	8.21751e-09\\
1.44	8.04081e-09\\
1.442	7.86787e-09\\
1.444	7.69861e-09\\
1.446	7.53295e-09\\
1.448	7.37081e-09\\
1.45	7.21214e-09\\
1.452	7.05684e-09\\
1.454	6.90486e-09\\
1.456	6.75612e-09\\
1.458	6.61055e-09\\
1.46	6.4681e-09\\
1.462	6.32868e-09\\
1.464	6.19226e-09\\
1.466	6.05875e-09\\
1.468	5.92809e-09\\
1.47	5.80024e-09\\
1.472	5.67513e-09\\
1.474	5.5527e-09\\
1.476	5.4329e-09\\
1.478	5.31567e-09\\
1.48	5.20096e-09\\
1.482	5.08871e-09\\
1.484	4.97888e-09\\
1.486	4.87142e-09\\
1.488	4.76626e-09\\
1.49	4.66337e-09\\
1.492	4.5627e-09\\
1.494	4.4642e-09\\
1.496	4.36782e-09\\
1.498	4.27353e-09\\
1.5	4.18127e-09\\
1.502	4.091e-09\\
1.504	4.00268e-09\\
1.506	3.91628e-09\\
1.508	3.83174e-09\\
1.51	3.74903e-09\\
1.512	3.66811e-09\\
1.514	3.58895e-09\\
1.516	3.5115e-09\\
1.518	3.43572e-09\\
1.52	3.36159e-09\\
1.522	3.28907e-09\\
1.524	3.21812e-09\\
1.526	3.14871e-09\\
1.528	3.0808e-09\\
1.53	3.01437e-09\\
1.532	2.94938e-09\\
1.534	2.88579e-09\\
1.536	2.82359e-09\\
1.538	2.76274e-09\\
1.54	2.70322e-09\\
1.542	2.64498e-09\\
1.544	2.58801e-09\\
1.546	2.53227e-09\\
1.548	2.47775e-09\\
1.55	2.42441e-09\\
1.552	2.37222e-09\\
1.554	2.32117e-09\\
1.556	2.27123e-09\\
1.558	2.22237e-09\\
1.56	2.17458e-09\\
1.562	2.12782e-09\\
1.564	2.08207e-09\\
1.566	2.03732e-09\\
1.568	1.99354e-09\\
1.57	1.9507e-09\\
1.572	1.9088e-09\\
1.574	1.8678e-09\\
1.576	1.82769e-09\\
1.578	1.78845e-09\\
1.58	1.75006e-09\\
1.582	1.7125e-09\\
1.584	1.67576e-09\\
1.586	1.63981e-09\\
1.588	1.60463e-09\\
1.59	1.57022e-09\\
1.592	1.53655e-09\\
1.594	1.50361e-09\\
1.596	1.47138e-09\\
1.598	1.43984e-09\\
1.6	1.40898e-09\\
1.602	1.3788e-09\\
1.604	1.34926e-09\\
1.606	1.32035e-09\\
1.608	1.29207e-09\\
1.61	1.2644e-09\\
1.612	1.23733e-09\\
1.614	1.21084e-09\\
1.616	1.18491e-09\\
1.618	1.15955e-09\\
1.62	1.13473e-09\\
1.622	1.11044e-09\\
1.624	1.08667e-09\\
1.626	1.06342e-09\\
1.628	1.04066e-09\\
1.63	1.01839e-09\\
1.632	9.966e-10\\
1.634	9.75276e-10\\
1.636	9.54408e-10\\
1.638	9.33986e-10\\
1.64	9.14002e-10\\
1.642	8.94446e-10\\
1.644	8.75308e-10\\
1.646	8.56579e-10\\
1.648	8.3825e-10\\
1.65	8.20313e-10\\
1.652	8.0276e-10\\
1.654	7.85581e-10\\
1.656	7.68769e-10\\
1.658	7.52316e-10\\
1.66	7.36214e-10\\
1.662	7.20456e-10\\
1.664	7.05034e-10\\
1.666	6.8994e-10\\
1.668	6.75169e-10\\
1.67	6.60713e-10\\
1.672	6.46565e-10\\
1.674	6.32718e-10\\
1.676	6.19167e-10\\
1.678	6.05904e-10\\
1.68	5.92924e-10\\
1.682	5.80221e-10\\
1.684	5.67788e-10\\
1.686	5.55621e-10\\
1.688	5.43712e-10\\
1.69	5.32058e-10\\
1.692	5.20651e-10\\
1.694	5.09488e-10\\
1.696	4.98563e-10\\
1.698	4.8787e-10\\
1.7	4.77405e-10\\
1.702	4.67164e-10\\
1.704	4.5714e-10\\
1.706	4.47331e-10\\
1.708	4.3773e-10\\
1.71	4.28335e-10\\
1.712	4.19139e-10\\
1.714	4.1014e-10\\
1.716	4.01333e-10\\
1.718	3.92714e-10\\
1.72	3.84278e-10\\
1.722	3.76023e-10\\
1.724	3.67944e-10\\
1.726	3.60038e-10\\
1.728	3.523e-10\\
1.73	3.44728e-10\\
1.732	3.37318e-10\\
1.734	3.30066e-10\\
1.736	3.22969e-10\\
1.738	3.16024e-10\\
1.74	3.09228e-10\\
1.742	3.02577e-10\\
1.744	2.96068e-10\\
1.746	2.89699e-10\\
1.748	2.83466e-10\\
1.75	2.77367e-10\\
1.752	2.71398e-10\\
1.754	2.65557e-10\\
1.756	2.59842e-10\\
1.758	2.54249e-10\\
1.76	2.48776e-10\\
1.762	2.43421e-10\\
1.764	2.38181e-10\\
1.766	2.33053e-10\\
1.768	2.28035e-10\\
1.77	2.23125e-10\\
1.772	2.18321e-10\\
1.774	2.1362e-10\\
1.776	2.0902e-10\\
1.778	2.04519e-10\\
1.78	2.00115e-10\\
1.782	1.95805e-10\\
1.784	1.91589e-10\\
1.786	1.87463e-10\\
1.788	1.83426e-10\\
1.79	1.79476e-10\\
1.792	1.75611e-10\\
1.794	1.71829e-10\\
1.796	1.68129e-10\\
1.798	1.64509e-10\\
1.8	1.60966e-10\\
1.802	1.57501e-10\\
1.804	1.54109e-10\\
1.806	1.50791e-10\\
1.808	1.47545e-10\\
1.81	1.44369e-10\\
1.812	1.41261e-10\\
1.814	1.3822e-10\\
1.816	1.35245e-10\\
1.818	1.32334e-10\\
1.82	1.29486e-10\\
1.822	1.267e-10\\
1.824	1.23974e-10\\
1.826	1.21306e-10\\
1.828	1.18696e-10\\
1.83	1.16143e-10\\
1.832	1.13645e-10\\
1.834	1.112e-10\\
1.836	1.08809e-10\\
1.838	1.06469e-10\\
1.84	1.04179e-10\\
1.842	1.01939e-10\\
1.844	9.97478e-11\\
1.846	9.76035e-11\\
1.848	9.55056e-11\\
1.85	9.3453e-11\\
1.852	9.14446e-11\\
1.854	8.94796e-11\\
1.856	8.7557e-11\\
1.858	8.56759e-11\\
1.86	8.38354e-11\\
1.862	8.20346e-11\\
1.864	8.02727e-11\\
1.866	7.85487e-11\\
1.868	7.68619e-11\\
1.87	7.52115e-11\\
1.872	7.35967e-11\\
1.874	7.20167e-11\\
1.876	7.04707e-11\\
1.878	6.89581e-11\\
1.88	6.7478e-11\\
1.882	6.60298e-11\\
1.884	6.46128e-11\\
1.886	6.32263e-11\\
1.888	6.18697e-11\\
1.89	6.05423e-11\\
1.892	5.92434e-11\\
1.894	5.79725e-11\\
1.896	5.67289e-11\\
1.898	5.55121e-11\\
1.9	5.43215e-11\\
1.902	5.31564e-11\\
1.904	5.20164e-11\\
1.906	5.09009e-11\\
1.908	4.98094e-11\\
1.91	4.87413e-11\\
1.912	4.76962e-11\\
1.914	4.66735e-11\\
1.916	4.56727e-11\\
1.918	4.46935e-11\\
1.92	4.37353e-11\\
1.922	4.27976e-11\\
1.924	4.18801e-11\\
1.926	4.09822e-11\\
1.928	4.01036e-11\\
1.93	3.92439e-11\\
1.932	3.84026e-11\\
1.934	3.75793e-11\\
1.936	3.67736e-11\\
1.938	3.59852e-11\\
1.94	3.52138e-11\\
1.942	3.44588e-11\\
1.944	3.372e-11\\
1.946	3.29971e-11\\
1.948	3.22896e-11\\
1.95	3.15974e-11\\
1.952	3.09199e-11\\
1.954	3.02569e-11\\
1.956	2.96081e-11\\
1.958	2.89733e-11\\
1.96	2.8352e-11\\
1.962	2.77439e-11\\
1.964	2.71489e-11\\
1.966	2.65667e-11\\
1.968	2.59969e-11\\
1.97	2.54393e-11\\
1.972	2.48936e-11\\
1.974	2.43596e-11\\
1.976	2.3837e-11\\
1.978	2.33256e-11\\
1.98	2.28251e-11\\
1.982	2.23354e-11\\
1.984	2.18561e-11\\
1.986	2.13871e-11\\
1.988	2.09281e-11\\
1.99	2.04789e-11\\
1.992	2.00394e-11\\
1.994	1.96092e-11\\
1.996	1.91883e-11\\
1.998	1.87764e-11\\
2	1.83733e-11\\
};
\addplot [color=mycolor7, line width=1.0pt, forget plot]
  table[row sep=crcr]{%
-0.024	0\\
-0.022	0\\
-0.02	0\\
-0.018	0\\
-0.016	0\\
-0.014	0\\
-0.012	0\\
-0.01	0\\
-0.008	0\\
-0.006	0\\
-0.004	0\\
-0.002	0\\
0	0\\
0.002	0.00916175\\
0.004	0.0147058\\
0.006	0.0198607\\
0.008	0.0237421\\
0.01	0.0264677\\
0.012	0.0281815\\
0.014	0.0290317\\
0.016	0.0291664\\
0.018	0.0287765\\
0.02	0.0288963\\
0.022	0.028967\\
0.024	0.0289749\\
0.026	0.0286487\\
0.028	0.028026\\
0.03	0.0276602\\
0.032	0.0289044\\
0.034	0.0300446\\
0.036	0.0310798\\
0.038	0.0320094\\
0.04	0.0328336\\
0.042	0.0335527\\
0.044	0.0341676\\
0.046	0.0346795\\
0.048	0.03509\\
0.05	0.0354011\\
0.052	0.035615\\
0.054	0.0357344\\
0.056	0.0357622\\
0.058	0.0357015\\
0.06	0.0355556\\
0.062	0.0353282\\
0.064	0.0350231\\
0.066	0.0346442\\
0.068	0.0341954\\
0.07	0.033681\\
0.072	0.0331053\\
0.074	0.0324724\\
0.076	0.0317868\\
0.078	0.0310527\\
0.08	0.0302746\\
0.082	0.0294566\\
0.084	0.028603\\
0.086	0.027718\\
0.088	0.0268057\\
0.09	0.0258701\\
0.092	0.024915\\
0.094	0.0239443\\
0.096	0.0229615\\
0.098	0.0219703\\
0.1	0.0209738\\
0.102	0.0199753\\
0.104	0.018978\\
0.106	0.0179846\\
0.108	0.0169979\\
0.11	0.0160206\\
0.112	0.0150549\\
0.114	0.0141032\\
0.116	0.0135189\\
0.118	0.0131468\\
0.12	0.0127656\\
0.122	0.0123766\\
0.124	0.0119811\\
0.126	0.0115803\\
0.128	0.0111753\\
0.13	0.0107673\\
0.132	0.0103574\\
0.134	0.010124\\
0.136	0.0100563\\
0.138	0.00997781\\
0.14	0.00988756\\
0.142	0.00978612\\
0.144	0.00967399\\
0.146	0.00955169\\
0.148	0.00941974\\
0.15	0.00927869\\
0.152	0.00912905\\
0.154	0.00897137\\
0.156	0.00880617\\
0.158	0.00863399\\
0.16	0.00845537\\
0.162	0.00827082\\
0.164	0.00808086\\
0.166	0.00788603\\
0.168	0.00768681\\
0.17	0.00748372\\
0.172	0.00727725\\
0.174	0.00706788\\
0.176	0.00685609\\
0.178	0.00664233\\
0.18	0.00642706\\
0.182	0.00621071\\
0.184	0.00605181\\
0.186	0.00594156\\
0.188	0.00583268\\
0.19	0.00572522\\
0.192	0.00561921\\
0.194	0.00551468\\
0.196	0.00541165\\
0.198	0.00531014\\
0.2	0.00521017\\
0.202	0.00511173\\
0.204	0.00501482\\
0.206	0.00491944\\
0.208	0.00482559\\
0.21	0.00473325\\
0.212	0.0046424\\
0.214	0.00455303\\
0.216	0.00446512\\
0.218	0.00437865\\
0.22	0.0042936\\
0.222	0.00420994\\
0.224	0.00412765\\
0.226	0.00404671\\
0.228	0.00396709\\
0.23	0.00388877\\
0.232	0.00381173\\
0.234	0.00373595\\
0.236	0.0036614\\
0.238	0.00358807\\
0.24	0.00351593\\
0.242	0.00344498\\
0.244	0.0033752\\
0.246	0.00330656\\
0.248	0.00323906\\
0.25	0.00317268\\
0.252	0.00310741\\
0.254	0.00304324\\
0.256	0.00298017\\
0.258	0.00291817\\
0.26	0.00285725\\
0.262	0.00279739\\
0.264	0.00273859\\
0.266	0.00268084\\
0.268	0.00262414\\
0.27	0.00256846\\
0.272	0.00251382\\
0.274	0.0024602\\
0.276	0.00240759\\
0.278	0.002356\\
0.28	0.0023054\\
0.282	0.00225579\\
0.284	0.00220716\\
0.286	0.00215951\\
0.288	0.00211283\\
0.29	0.0020671\\
0.292	0.00202231\\
0.294	0.00197846\\
0.296	0.00193554\\
0.298	0.00189352\\
0.3	0.00185241\\
0.302	0.00181218\\
0.304	0.00177283\\
0.306	0.00173434\\
0.308	0.0016967\\
0.31	0.00165989\\
0.312	0.00162391\\
0.314	0.00158873\\
0.316	0.00155434\\
0.318	0.00152072\\
0.32	0.00148787\\
0.322	0.00145576\\
0.324	0.00142439\\
0.326	0.00139373\\
0.328	0.00136378\\
0.33	0.0013345\\
0.332	0.0013059\\
0.334	0.00127796\\
0.336	0.00125066\\
0.338	0.00122398\\
0.34	0.00119791\\
0.342	0.00117244\\
0.344	0.00114755\\
0.346	0.00112323\\
0.348	0.00109946\\
0.35	0.00107624\\
0.352	0.00105354\\
0.354	0.00103135\\
0.356	0.00100966\\
0.358	0.000988461\\
0.36	0.000967736\\
0.362	0.000947473\\
0.364	0.000927662\\
0.366	0.000908289\\
0.368	0.000889345\\
0.37	0.000870817\\
0.372	0.000852696\\
0.374	0.00083497\\
0.376	0.00081763\\
0.378	0.000800666\\
0.38	0.000784068\\
0.382	0.000767827\\
0.384	0.000751935\\
0.386	0.000736383\\
0.388	0.000721162\\
0.39	0.000706264\\
0.392	0.000691682\\
0.394	0.000677408\\
0.396	0.000663435\\
0.398	0.000649755\\
0.4	0.000636362\\
0.402	0.00062325\\
0.404	0.000610411\\
0.406	0.00059784\\
0.408	0.00058553\\
0.41	0.000573477\\
0.412	0.000561674\\
0.414	0.000550115\\
0.416	0.000538796\\
0.418	0.000527711\\
0.42	0.000516855\\
0.422	0.000506224\\
0.424	0.000495812\\
0.426	0.000485615\\
0.428	0.000475628\\
0.43	0.000465847\\
0.432	0.000456267\\
0.434	0.000446885\\
0.436	0.000437696\\
0.438	0.000428696\\
0.44	0.00041988\\
0.442	0.000411246\\
0.444	0.000402789\\
0.446	0.000394506\\
0.448	0.000386392\\
0.45	0.000378445\\
0.452	0.00037066\\
0.454	0.000363035\\
0.456	0.000355565\\
0.458	0.000348248\\
0.46	0.00034108\\
0.462	0.000334058\\
0.464	0.000327178\\
0.466	0.000320439\\
0.468	0.000313836\\
0.47	0.000307368\\
0.472	0.00030103\\
0.474	0.00029482\\
0.476	0.000288736\\
0.478	0.000282775\\
0.48	0.000276934\\
0.482	0.000271211\\
0.484	0.000265603\\
0.486	0.000260107\\
0.488	0.000254722\\
0.49	0.000249445\\
0.492	0.000244274\\
0.494	0.000239206\\
0.496	0.00023424\\
0.498	0.000229373\\
0.5	0.000224604\\
0.502	0.00021993\\
0.504	0.00021535\\
0.506	0.000210862\\
0.508	0.000206463\\
0.51	0.000202152\\
0.512	0.000197928\\
0.514	0.000193788\\
0.516	0.000189731\\
0.518	0.000185756\\
0.52	0.00018186\\
0.522	0.000178043\\
0.524	0.000174302\\
0.526	0.000170637\\
0.528	0.000167046\\
0.53	0.000163527\\
0.532	0.00016008\\
0.534	0.000156702\\
0.536	0.000153393\\
0.538	0.000150151\\
0.54	0.000146975\\
0.542	0.000143863\\
0.544	0.000140816\\
0.546	0.00013783\\
0.548	0.000134906\\
0.55	0.000132042\\
0.552	0.000129236\\
0.554	0.000126489\\
0.556	0.000123798\\
0.558	0.000121163\\
0.56	0.000118583\\
0.562	0.000116056\\
0.564	0.000113582\\
0.566	0.000111159\\
0.568	0.000108787\\
0.57	0.000106465\\
0.572	0.000104191\\
0.574	0.000101965\\
0.576	9.97863e-05\\
0.578	9.76531e-05\\
0.58	9.5565e-05\\
0.582	9.35211e-05\\
0.584	9.15204e-05\\
0.586	8.95622e-05\\
0.588	8.76455e-05\\
0.59	8.57696e-05\\
0.592	8.39336e-05\\
0.594	8.21368e-05\\
0.596	8.03783e-05\\
0.598	7.86574e-05\\
0.6	7.69734e-05\\
0.602	7.53253e-05\\
0.604	7.37126e-05\\
0.606	7.21345e-05\\
0.608	7.05902e-05\\
0.61	6.90791e-05\\
0.612	6.76005e-05\\
0.614	6.61538e-05\\
0.616	6.47381e-05\\
0.618	6.33529e-05\\
0.62	6.19976e-05\\
0.622	6.06715e-05\\
0.624	5.9374e-05\\
0.626	5.81045e-05\\
0.628	5.68623e-05\\
0.63	5.5647e-05\\
0.632	5.44579e-05\\
0.634	5.32945e-05\\
0.636	5.21563e-05\\
0.638	5.10426e-05\\
0.64	4.9953e-05\\
0.642	4.88869e-05\\
0.644	4.78439e-05\\
0.646	4.68234e-05\\
0.648	4.58249e-05\\
0.65	4.48481e-05\\
0.652	4.38923e-05\\
0.654	4.29572e-05\\
0.656	4.20424e-05\\
0.658	4.11473e-05\\
0.66	4.02715e-05\\
0.662	3.94147e-05\\
0.664	3.85764e-05\\
0.666	3.77562e-05\\
0.668	3.69537e-05\\
0.67	3.61686e-05\\
0.672	3.54004e-05\\
0.674	3.46488e-05\\
0.676	3.39134e-05\\
0.678	3.31939e-05\\
0.68	3.24899e-05\\
0.682	3.18012e-05\\
0.684	3.11272e-05\\
0.686	3.04679e-05\\
0.688	2.98227e-05\\
0.69	2.91914e-05\\
0.692	2.85737e-05\\
0.694	2.79694e-05\\
0.696	2.7378e-05\\
0.698	2.67994e-05\\
0.7	2.62332e-05\\
0.702	2.56792e-05\\
0.704	2.5137e-05\\
0.706	2.46066e-05\\
0.708	2.40875e-05\\
0.71	2.35795e-05\\
0.712	2.30824e-05\\
0.714	2.2596e-05\\
0.716	2.212e-05\\
0.718	2.16542e-05\\
0.72	2.11983e-05\\
0.722	2.07522e-05\\
0.724	2.03155e-05\\
0.726	1.98882e-05\\
0.728	1.947e-05\\
0.73	1.90608e-05\\
0.732	1.86602e-05\\
0.734	1.82681e-05\\
0.736	1.78844e-05\\
0.738	1.75088e-05\\
0.74	1.71411e-05\\
0.742	1.67813e-05\\
0.744	1.64291e-05\\
0.746	1.60843e-05\\
0.748	1.57468e-05\\
0.75	1.54164e-05\\
0.752	1.5093e-05\\
0.754	1.47764e-05\\
0.756	1.44665e-05\\
0.758	1.41631e-05\\
0.76	1.38661e-05\\
0.762	1.35753e-05\\
0.764	1.32906e-05\\
0.766	1.30119e-05\\
0.768	1.2739e-05\\
0.77	1.24718e-05\\
0.772	1.22102e-05\\
0.774	1.19541e-05\\
0.776	1.17034e-05\\
0.778	1.14579e-05\\
0.78	1.12175e-05\\
0.782	1.09821e-05\\
0.784	1.07516e-05\\
0.786	1.05259e-05\\
0.788	1.0305e-05\\
0.79	1.00886e-05\\
0.792	9.87672e-06\\
0.794	9.66925e-06\\
0.796	9.4661e-06\\
0.798	9.26716e-06\\
0.8	9.07237e-06\\
0.802	8.88161e-06\\
0.804	8.69482e-06\\
0.806	8.51191e-06\\
0.808	8.3328e-06\\
0.81	8.1574e-06\\
0.812	7.98565e-06\\
0.814	7.81746e-06\\
0.816	7.65276e-06\\
0.818	7.49148e-06\\
0.82	7.33354e-06\\
0.822	7.17889e-06\\
0.824	7.02745e-06\\
0.826	6.87915e-06\\
0.828	6.73394e-06\\
0.83	6.59174e-06\\
0.832	6.45249e-06\\
0.834	6.31615e-06\\
0.836	6.18263e-06\\
0.838	6.05189e-06\\
0.84	5.92388e-06\\
0.842	5.79852e-06\\
0.844	5.67578e-06\\
0.846	5.55559e-06\\
0.848	5.43791e-06\\
0.85	5.32269e-06\\
0.852	5.20986e-06\\
0.854	5.09939e-06\\
0.856	4.99123e-06\\
0.858	4.88532e-06\\
0.86	4.78163e-06\\
0.862	4.68011e-06\\
0.864	4.58071e-06\\
0.866	4.4834e-06\\
0.868	4.38812e-06\\
0.87	4.29484e-06\\
0.872	4.20351e-06\\
0.874	4.11411e-06\\
0.876	4.02658e-06\\
0.878	3.94089e-06\\
0.88	3.857e-06\\
0.882	3.77487e-06\\
0.884	3.69448e-06\\
0.886	3.61578e-06\\
0.888	3.53873e-06\\
0.89	3.46332e-06\\
0.892	3.38949e-06\\
0.894	3.31723e-06\\
0.896	3.24649e-06\\
0.898	3.17725e-06\\
0.9	3.10947e-06\\
0.902	3.04313e-06\\
0.904	2.9782e-06\\
0.906	2.91464e-06\\
0.908	2.85243e-06\\
0.91	2.79154e-06\\
0.912	2.73195e-06\\
0.914	2.67362e-06\\
0.916	2.61654e-06\\
0.918	2.56067e-06\\
0.92	2.50599e-06\\
0.922	2.45247e-06\\
0.924	2.4001e-06\\
0.926	2.34885e-06\\
0.928	2.29869e-06\\
0.93	2.2496e-06\\
0.932	2.20156e-06\\
0.934	2.15455e-06\\
0.936	2.10854e-06\\
0.938	2.06352e-06\\
0.94	2.01946e-06\\
0.942	1.97635e-06\\
0.944	1.93416e-06\\
0.946	1.89288e-06\\
0.948	1.85248e-06\\
0.95	1.81295e-06\\
0.952	1.77427e-06\\
0.954	1.73642e-06\\
0.956	1.69938e-06\\
0.958	1.66314e-06\\
0.96	1.62768e-06\\
0.962	1.59298e-06\\
0.964	1.55903e-06\\
0.966	1.5258e-06\\
0.968	1.49329e-06\\
0.97	1.46148e-06\\
0.972	1.43036e-06\\
0.974	1.3999e-06\\
0.976	1.3701e-06\\
0.978	1.34095e-06\\
0.98	1.31241e-06\\
0.982	1.2845e-06\\
0.984	1.25718e-06\\
0.986	1.23045e-06\\
0.988	1.2043e-06\\
0.99	1.17871e-06\\
0.992	1.15367e-06\\
0.994	1.12916e-06\\
0.996	1.10519e-06\\
0.998	1.08173e-06\\
1	1.05877e-06\\
1.002	1.03631e-06\\
1.004	1.01433e-06\\
1.006	9.92818e-07\\
1.008	9.71771e-07\\
1.01	9.51175e-07\\
1.012	9.3102e-07\\
1.014	9.11298e-07\\
1.016	8.91998e-07\\
1.018	8.73111e-07\\
1.02	8.54629e-07\\
1.022	8.36542e-07\\
1.024	8.18842e-07\\
1.026	8.0152e-07\\
1.028	7.84568e-07\\
1.03	7.67978e-07\\
1.032	7.51742e-07\\
1.034	7.35853e-07\\
1.036	7.20302e-07\\
1.038	7.05083e-07\\
1.04	6.90187e-07\\
1.042	6.75609e-07\\
1.044	6.6134e-07\\
1.046	6.47375e-07\\
1.048	6.33707e-07\\
1.05	6.20329e-07\\
1.052	6.07234e-07\\
1.054	5.94418e-07\\
1.056	5.81873e-07\\
1.058	5.69594e-07\\
1.06	5.57575e-07\\
1.062	5.4581e-07\\
1.064	5.34295e-07\\
1.066	5.23022e-07\\
1.068	5.11988e-07\\
1.07	5.01187e-07\\
1.072	4.90614e-07\\
1.074	4.80264e-07\\
1.076	4.70132e-07\\
1.078	4.60214e-07\\
1.08	4.50505e-07\\
1.082	4.41e-07\\
1.084	4.31695e-07\\
1.086	4.22586e-07\\
1.088	4.13669e-07\\
1.09	4.0494e-07\\
1.092	3.96393e-07\\
1.094	3.88027e-07\\
1.096	3.79836e-07\\
1.098	3.71818e-07\\
1.1	3.63967e-07\\
1.102	3.56282e-07\\
1.104	3.48757e-07\\
1.106	3.41391e-07\\
1.108	3.34179e-07\\
1.11	3.27119e-07\\
1.112	3.20206e-07\\
1.114	3.13439e-07\\
1.116	3.06813e-07\\
1.118	3.00327e-07\\
1.12	2.93976e-07\\
1.122	2.87759e-07\\
1.124	2.81672e-07\\
1.126	2.75712e-07\\
1.128	2.69878e-07\\
1.13	2.64166e-07\\
1.132	2.58573e-07\\
1.134	2.53098e-07\\
1.136	2.47738e-07\\
1.138	2.4249e-07\\
1.14	2.37352e-07\\
1.142	2.32322e-07\\
1.144	2.27398e-07\\
1.146	2.22577e-07\\
1.148	2.17857e-07\\
1.15	2.13236e-07\\
1.152	2.08713e-07\\
1.154	2.04284e-07\\
1.156	1.99949e-07\\
1.158	1.95704e-07\\
1.16	1.91549e-07\\
1.162	1.87482e-07\\
1.164	1.835e-07\\
1.166	1.79601e-07\\
1.168	1.75785e-07\\
1.17	1.7205e-07\\
1.172	1.68393e-07\\
1.174	1.64813e-07\\
1.176	1.61308e-07\\
1.178	1.57878e-07\\
1.18	1.5452e-07\\
1.182	1.51233e-07\\
1.184	1.48015e-07\\
1.186	1.44866e-07\\
1.188	1.41783e-07\\
1.19	1.38765e-07\\
1.192	1.35811e-07\\
1.194	1.32919e-07\\
1.196	1.30089e-07\\
1.198	1.27319e-07\\
1.2	1.24608e-07\\
1.202	1.21954e-07\\
1.204	1.19356e-07\\
1.206	1.16814e-07\\
1.208	1.14325e-07\\
1.21	1.1189e-07\\
1.212	1.09506e-07\\
1.214	1.07173e-07\\
1.216	1.04889e-07\\
1.218	1.02655e-07\\
1.22	1.00467e-07\\
1.222	9.83266e-08\\
1.224	9.62314e-08\\
1.226	9.4181e-08\\
1.228	9.21742e-08\\
1.23	9.02103e-08\\
1.232	8.82882e-08\\
1.234	8.64071e-08\\
1.236	8.45663e-08\\
1.238	8.27647e-08\\
1.24	8.10015e-08\\
1.242	7.92761e-08\\
1.244	7.75875e-08\\
1.246	7.5935e-08\\
1.248	7.43178e-08\\
1.25	7.27353e-08\\
1.252	7.11865e-08\\
1.254	6.96709e-08\\
1.256	6.81877e-08\\
1.258	6.67362e-08\\
1.26	6.53157e-08\\
1.262	6.39257e-08\\
1.264	6.25654e-08\\
1.266	6.12342e-08\\
1.268	5.99315e-08\\
1.27	5.86567e-08\\
1.272	5.74092e-08\\
1.274	5.61884e-08\\
1.276	5.49936e-08\\
1.278	5.38245e-08\\
1.28	5.26804e-08\\
1.282	5.15607e-08\\
1.284	5.04651e-08\\
1.286	4.93928e-08\\
1.288	4.83435e-08\\
1.29	4.73167e-08\\
1.292	4.63118e-08\\
1.294	4.53283e-08\\
1.296	4.4366e-08\\
1.298	4.34241e-08\\
1.3	4.25024e-08\\
1.302	4.16005e-08\\
1.304	4.07177e-08\\
1.306	3.98539e-08\\
1.308	3.90084e-08\\
1.31	3.81811e-08\\
1.312	3.73713e-08\\
1.314	3.65789e-08\\
1.316	3.58034e-08\\
1.318	3.50444e-08\\
1.32	3.43016e-08\\
1.322	3.35746e-08\\
1.324	3.28631e-08\\
1.326	3.21668e-08\\
1.328	3.14853e-08\\
1.33	3.08183e-08\\
1.332	3.01655e-08\\
1.334	2.95266e-08\\
1.336	2.89012e-08\\
1.338	2.82892e-08\\
1.34	2.76902e-08\\
1.342	2.71039e-08\\
1.344	2.65301e-08\\
1.346	2.59685e-08\\
1.348	2.54187e-08\\
1.35	2.48807e-08\\
1.352	2.43541e-08\\
1.354	2.38386e-08\\
1.356	2.33341e-08\\
1.358	2.28402e-08\\
1.36	2.23568e-08\\
1.362	2.18837e-08\\
1.364	2.14206e-08\\
1.366	2.09673e-08\\
1.368	2.05236e-08\\
1.37	2.00893e-08\\
1.372	1.96642e-08\\
1.374	1.9248e-08\\
1.376	1.88407e-08\\
1.378	1.8442e-08\\
1.38	1.80517e-08\\
1.382	1.76697e-08\\
1.384	1.72957e-08\\
1.386	1.69296e-08\\
1.388	1.65713e-08\\
1.39	1.62206e-08\\
1.392	1.58772e-08\\
1.394	1.55411e-08\\
1.396	1.52121e-08\\
1.398	1.489e-08\\
1.4	1.45748e-08\\
1.402	1.42661e-08\\
1.404	1.3964e-08\\
1.406	1.36683e-08\\
1.408	1.33788e-08\\
1.41	1.30955e-08\\
1.412	1.28181e-08\\
1.414	1.25465e-08\\
1.416	1.22807e-08\\
1.418	1.20205e-08\\
1.42	1.17658e-08\\
1.422	1.15164e-08\\
1.424	1.12723e-08\\
1.426	1.10334e-08\\
1.428	1.07995e-08\\
1.43	1.05705e-08\\
1.432	1.03464e-08\\
1.434	1.0127e-08\\
1.436	9.91225e-09\\
1.438	9.70201e-09\\
1.44	9.49621e-09\\
1.442	9.29476e-09\\
1.444	9.09756e-09\\
1.446	8.90452e-09\\
1.448	8.71556e-09\\
1.45	8.5306e-09\\
1.452	8.34954e-09\\
1.454	8.1723e-09\\
1.456	7.99881e-09\\
1.458	7.82899e-09\\
1.46	7.66275e-09\\
1.462	7.50003e-09\\
1.464	7.34076e-09\\
1.466	7.18485e-09\\
1.468	7.03223e-09\\
1.47	6.88285e-09\\
1.472	6.73663e-09\\
1.474	6.59351e-09\\
1.476	6.45341e-09\\
1.478	6.31628e-09\\
1.48	6.18206e-09\\
1.482	6.05067e-09\\
1.484	5.92208e-09\\
1.486	5.7962e-09\\
1.488	5.673e-09\\
1.49	5.55241e-09\\
1.492	5.43437e-09\\
1.494	5.31884e-09\\
1.496	5.20576e-09\\
1.498	5.09508e-09\\
1.5	4.98675e-09\\
1.502	4.88072e-09\\
1.504	4.77694e-09\\
1.506	4.67536e-09\\
1.508	4.57595e-09\\
1.51	4.47864e-09\\
1.512	4.3834e-09\\
1.514	4.29019e-09\\
1.516	4.19895e-09\\
1.518	4.10966e-09\\
1.52	4.02226e-09\\
1.522	3.93673e-09\\
1.524	3.85301e-09\\
1.526	3.77107e-09\\
1.528	3.69088e-09\\
1.53	3.6124e-09\\
1.532	3.53558e-09\\
1.534	3.4604e-09\\
1.536	3.38682e-09\\
1.538	3.31481e-09\\
1.54	3.24433e-09\\
1.542	3.17535e-09\\
1.544	3.10784e-09\\
1.546	3.04176e-09\\
1.548	2.9771e-09\\
1.55	2.91381e-09\\
1.552	2.85187e-09\\
1.554	2.79125e-09\\
1.556	2.73192e-09\\
1.558	2.67386e-09\\
1.56	2.61703e-09\\
1.562	2.56141e-09\\
1.564	2.50698e-09\\
1.566	2.45371e-09\\
1.568	2.40157e-09\\
1.57	2.35055e-09\\
1.572	2.30061e-09\\
1.574	2.25174e-09\\
1.576	2.2039e-09\\
1.578	2.15709e-09\\
1.58	2.11127e-09\\
1.582	2.06643e-09\\
1.584	2.02255e-09\\
1.586	1.9796e-09\\
1.588	1.93756e-09\\
1.59	1.89642e-09\\
1.592	1.85616e-09\\
1.594	1.81675e-09\\
1.596	1.77819e-09\\
1.598	1.74044e-09\\
1.6	1.7035e-09\\
1.602	1.66734e-09\\
1.604	1.63196e-09\\
1.606	1.59732e-09\\
1.608	1.56343e-09\\
1.61	1.53025e-09\\
1.612	1.49778e-09\\
1.614	1.46601e-09\\
1.616	1.43491e-09\\
1.618	1.40447e-09\\
1.62	1.37467e-09\\
1.622	1.34551e-09\\
1.624	1.31697e-09\\
1.626	1.28904e-09\\
1.628	1.2617e-09\\
1.63	1.23495e-09\\
1.632	1.20876e-09\\
1.634	1.18312e-09\\
1.636	1.15804e-09\\
1.638	1.13348e-09\\
1.64	1.10945e-09\\
1.642	1.08592e-09\\
1.644	1.0629e-09\\
1.646	1.04036e-09\\
1.648	1.01831e-09\\
1.65	9.9672e-10\\
1.652	9.7559e-10\\
1.654	9.54908e-10\\
1.656	9.34665e-10\\
1.658	9.14851e-10\\
1.66	8.95458e-10\\
1.662	8.76476e-10\\
1.664	8.57896e-10\\
1.666	8.39711e-10\\
1.668	8.21911e-10\\
1.67	8.04488e-10\\
1.672	7.87435e-10\\
1.674	7.70743e-10\\
1.676	7.54404e-10\\
1.678	7.38413e-10\\
1.68	7.22759e-10\\
1.682	7.07438e-10\\
1.684	6.92441e-10\\
1.686	6.77762e-10\\
1.688	6.63394e-10\\
1.69	6.4933e-10\\
1.692	6.35564e-10\\
1.694	6.2209e-10\\
1.696	6.089e-10\\
1.698	5.9599e-10\\
1.7	5.83354e-10\\
1.702	5.70985e-10\\
1.704	5.58878e-10\\
1.706	5.47027e-10\\
1.708	5.35427e-10\\
1.71	5.24073e-10\\
1.712	5.12959e-10\\
1.714	5.02081e-10\\
1.716	4.91432e-10\\
1.718	4.81009e-10\\
1.72	4.70807e-10\\
1.722	4.60821e-10\\
1.724	4.51046e-10\\
1.726	4.41478e-10\\
1.728	4.32113e-10\\
1.73	4.22946e-10\\
1.732	4.13973e-10\\
1.734	4.0519e-10\\
1.736	3.96593e-10\\
1.738	3.88178e-10\\
1.74	3.79941e-10\\
1.742	3.71879e-10\\
1.744	3.63987e-10\\
1.746	3.56263e-10\\
1.748	3.48702e-10\\
1.75	3.41301e-10\\
1.752	3.34057e-10\\
1.754	3.26966e-10\\
1.756	3.20026e-10\\
1.758	3.13233e-10\\
1.76	3.06583e-10\\
1.762	3.00075e-10\\
1.764	2.93704e-10\\
1.766	2.87469e-10\\
1.768	2.81365e-10\\
1.77	2.75391e-10\\
1.772	2.69544e-10\\
1.774	2.6382e-10\\
1.776	2.58218e-10\\
1.778	2.52735e-10\\
1.78	2.47368e-10\\
1.782	2.42115e-10\\
1.784	2.36973e-10\\
1.786	2.3194e-10\\
1.788	2.27014e-10\\
1.79	2.22192e-10\\
1.792	2.17473e-10\\
1.794	2.12854e-10\\
1.796	2.08333e-10\\
1.798	2.03908e-10\\
1.8	1.99576e-10\\
1.802	1.95337e-10\\
1.804	1.91188e-10\\
1.806	1.87127e-10\\
1.808	1.83151e-10\\
1.81	1.79261e-10\\
1.812	1.75453e-10\\
1.814	1.71726e-10\\
1.816	1.68078e-10\\
1.818	1.64507e-10\\
1.82	1.61012e-10\\
1.822	1.57592e-10\\
1.824	1.54244e-10\\
1.826	1.50968e-10\\
1.828	1.47761e-10\\
1.83	1.44622e-10\\
1.832	1.41549e-10\\
1.834	1.38542e-10\\
1.836	1.35599e-10\\
1.838	1.32719e-10\\
1.84	1.299e-10\\
1.842	1.2714e-10\\
1.844	1.2444e-10\\
1.846	1.21796e-10\\
1.848	1.19209e-10\\
1.85	1.16677e-10\\
1.852	1.14199e-10\\
1.854	1.11773e-10\\
1.856	1.09399e-10\\
1.858	1.07076e-10\\
1.86	1.04802e-10\\
1.862	1.02576e-10\\
1.864	1.00397e-10\\
1.866	9.82654e-11\\
1.868	9.61786e-11\\
1.87	9.41362e-11\\
1.872	9.21372e-11\\
1.874	9.01807e-11\\
1.876	8.82657e-11\\
1.878	8.63915e-11\\
1.88	8.45572e-11\\
1.882	8.27618e-11\\
1.884	8.10046e-11\\
1.886	7.92847e-11\\
1.888	7.76014e-11\\
1.89	7.59539e-11\\
1.892	7.43414e-11\\
1.894	7.27632e-11\\
1.896	7.12185e-11\\
1.898	6.97066e-11\\
1.9	6.82269e-11\\
1.902	6.67787e-11\\
1.904	6.53612e-11\\
1.906	6.39738e-11\\
1.908	6.2616e-11\\
1.91	6.1287e-11\\
1.912	5.99862e-11\\
1.914	5.87131e-11\\
1.916	5.74671e-11\\
1.918	5.62475e-11\\
1.92	5.50538e-11\\
1.922	5.38855e-11\\
1.924	5.2742e-11\\
1.926	5.16228e-11\\
1.928	5.05274e-11\\
1.93	4.94552e-11\\
1.932	4.84058e-11\\
1.934	4.73787e-11\\
1.936	4.63734e-11\\
1.938	4.53894e-11\\
1.94	4.44264e-11\\
1.942	4.34838e-11\\
1.944	4.25612e-11\\
1.946	4.16582e-11\\
1.948	4.07744e-11\\
1.95	3.99093e-11\\
1.952	3.90626e-11\\
1.954	3.82339e-11\\
1.956	3.74228e-11\\
1.958	3.66288e-11\\
1.96	3.58518e-11\\
1.962	3.50912e-11\\
1.964	3.43468e-11\\
1.966	3.36181e-11\\
1.968	3.29049e-11\\
1.97	3.22069e-11\\
1.972	3.15236e-11\\
1.974	3.08549e-11\\
1.976	3.02004e-11\\
1.978	2.95597e-11\\
1.98	2.89326e-11\\
1.982	2.83189e-11\\
1.984	2.77181e-11\\
1.986	2.71302e-11\\
1.988	2.65547e-11\\
1.99	2.59913e-11\\
1.992	2.544e-11\\
1.994	2.49003e-11\\
1.996	2.4372e-11\\
1.998	2.3855e-11\\
2	2.33489e-11\\
};
\addplot [color=mycolor8, line width=1.0pt, forget plot]
  table[row sep=crcr]{%
-0.024	0\\
-0.022	0\\
-0.02	0\\
-0.018	0\\
-0.016	0\\
-0.014	0\\
-0.012	0\\
-0.01	0\\
-0.008	0\\
-0.006	0\\
-0.004	0\\
-0.002	0\\
0	0\\
0.002	0.00916168\\
0.004	0.014705\\
0.006	0.0198559\\
0.008	0.023729\\
0.01	0.0264398\\
0.012	0.0281307\\
0.014	0.028949\\
0.016	0.0290425\\
0.018	0.0285938\\
0.02	0.0286511\\
0.022	0.0283017\\
0.024	0.0276115\\
0.026	0.0266456\\
0.028	0.0263637\\
0.03	0.0276986\\
0.032	0.028928\\
0.034	0.0300506\\
0.036	0.0310657\\
0.038	0.0319731\\
0.04	0.0327731\\
0.042	0.0334665\\
0.044	0.0340545\\
0.046	0.0345387\\
0.048	0.0349211\\
0.05	0.035204\\
0.052	0.0353901\\
0.054	0.0354823\\
0.056	0.0354837\\
0.058	0.035398\\
0.06	0.0352286\\
0.062	0.0349793\\
0.064	0.0346542\\
0.066	0.0342573\\
0.068	0.0337927\\
0.07	0.0332646\\
0.072	0.0326775\\
0.074	0.0320355\\
0.076	0.0313429\\
0.078	0.0306041\\
0.08	0.0298233\\
0.082	0.0290047\\
0.084	0.0281525\\
0.086	0.0272707\\
0.088	0.0263633\\
0.09	0.0254342\\
0.092	0.0244871\\
0.094	0.0235257\\
0.096	0.0225536\\
0.098	0.0215741\\
0.1	0.0205904\\
0.102	0.0196056\\
0.104	0.0186228\\
0.106	0.0176446\\
0.108	0.0166738\\
0.11	0.0157129\\
0.112	0.0147641\\
0.114	0.0138297\\
0.116	0.0132132\\
0.118	0.0128546\\
0.12	0.0124874\\
0.122	0.0121129\\
0.124	0.0117321\\
0.126	0.0113464\\
0.128	0.0109569\\
0.13	0.0105646\\
0.132	0.0101707\\
0.134	0.00987277\\
0.136	0.00981845\\
0.138	0.00975229\\
0.14	0.00967475\\
0.142	0.00958631\\
0.144	0.00948743\\
0.146	0.0093786\\
0.148	0.00926031\\
0.15	0.00913305\\
0.152	0.00899732\\
0.154	0.0088536\\
0.156	0.00870241\\
0.158	0.00854422\\
0.16	0.00837954\\
0.162	0.00820886\\
0.164	0.00803266\\
0.166	0.00785144\\
0.168	0.00766566\\
0.17	0.0074758\\
0.172	0.00728232\\
0.174	0.00708569\\
0.176	0.00688635\\
0.178	0.00668474\\
0.18	0.0064813\\
0.182	0.00627644\\
0.184	0.00607058\\
0.186	0.00591285\\
0.188	0.00581497\\
0.19	0.00571822\\
0.192	0.00562263\\
0.194	0.00552823\\
0.196	0.00543505\\
0.198	0.00534311\\
0.2	0.00525242\\
0.202	0.00516299\\
0.204	0.00507483\\
0.206	0.00498794\\
0.208	0.00490233\\
0.21	0.00481798\\
0.212	0.00473489\\
0.214	0.00465305\\
0.216	0.00457245\\
0.218	0.00449307\\
0.22	0.0044149\\
0.222	0.00433793\\
0.224	0.00426213\\
0.226	0.00418749\\
0.228	0.00411399\\
0.23	0.00404162\\
0.232	0.00397034\\
0.234	0.00390016\\
0.236	0.00383103\\
0.238	0.00376296\\
0.24	0.00369592\\
0.242	0.00362989\\
0.244	0.00356485\\
0.246	0.0035008\\
0.248	0.00343771\\
0.25	0.00337558\\
0.252	0.00331439\\
0.254	0.00325412\\
0.256	0.00319477\\
0.258	0.00313632\\
0.26	0.00307876\\
0.262	0.00302209\\
0.264	0.00296629\\
0.266	0.00291136\\
0.268	0.00285728\\
0.27	0.00280406\\
0.272	0.00275168\\
0.274	0.00270015\\
0.276	0.00264944\\
0.278	0.00259956\\
0.28	0.0025505\\
0.282	0.00250225\\
0.284	0.00245482\\
0.286	0.00240819\\
0.288	0.00236235\\
0.29	0.00231731\\
0.292	0.00227306\\
0.294	0.00222958\\
0.296	0.00218689\\
0.298	0.00214495\\
0.3	0.00210378\\
0.302	0.00206337\\
0.304	0.0020237\\
0.306	0.00198477\\
0.308	0.00194658\\
0.31	0.0019091\\
0.312	0.00187234\\
0.314	0.00183629\\
0.316	0.00180094\\
0.318	0.00176627\\
0.32	0.00173228\\
0.322	0.00169896\\
0.324	0.00166629\\
0.326	0.00163427\\
0.328	0.00160289\\
0.33	0.00157214\\
0.332	0.001542\\
0.334	0.00151246\\
0.336	0.00148352\\
0.338	0.00145516\\
0.34	0.00142736\\
0.342	0.00140013\\
0.344	0.00137345\\
0.346	0.0013473\\
0.348	0.00132168\\
0.35	0.00129657\\
0.352	0.00127196\\
0.354	0.00124785\\
0.356	0.00122422\\
0.358	0.00120106\\
0.36	0.00117836\\
0.362	0.00115611\\
0.364	0.0011343\\
0.366	0.00111292\\
0.368	0.00109197\\
0.37	0.00107142\\
0.372	0.00105127\\
0.374	0.00103152\\
0.376	0.00101215\\
0.378	0.000993161\\
0.38	0.000974536\\
0.382	0.000956269\\
0.384	0.000938354\\
0.386	0.000920782\\
0.388	0.000903545\\
0.39	0.000886638\\
0.392	0.000870051\\
0.394	0.00085378\\
0.396	0.000837817\\
0.398	0.000822155\\
0.4	0.000806789\\
0.402	0.000791712\\
0.404	0.000776918\\
0.406	0.000762402\\
0.408	0.000748159\\
0.41	0.000734182\\
0.412	0.000720467\\
0.414	0.000707008\\
0.416	0.000693801\\
0.418	0.00068084\\
0.42	0.000668122\\
0.422	0.00065564\\
0.424	0.000643391\\
0.426	0.000631371\\
0.428	0.000619575\\
0.43	0.000607999\\
0.432	0.000596638\\
0.434	0.000585489\\
0.436	0.000574548\\
0.438	0.000563812\\
0.44	0.000553275\\
0.442	0.000542934\\
0.444	0.000532787\\
0.446	0.000522828\\
0.448	0.000513055\\
0.45	0.000503465\\
0.452	0.000494053\\
0.454	0.000484816\\
0.456	0.000475752\\
0.458	0.000466856\\
0.46	0.000458127\\
0.462	0.000449559\\
0.464	0.000441151\\
0.466	0.0004329\\
0.468	0.000424802\\
0.47	0.000416855\\
0.472	0.000409055\\
0.474	0.0004014\\
0.476	0.000393887\\
0.478	0.000386514\\
0.48	0.000379277\\
0.482	0.000372175\\
0.484	0.000365203\\
0.486	0.000358361\\
0.488	0.000351645\\
0.49	0.000345054\\
0.492	0.000338584\\
0.494	0.000332233\\
0.496	0.000326\\
0.498	0.000319881\\
0.5	0.000313875\\
0.502	0.00030798\\
0.504	0.000302193\\
0.506	0.000296513\\
0.508	0.000290937\\
0.51	0.000285464\\
0.512	0.000280091\\
0.514	0.000274817\\
0.516	0.00026964\\
0.518	0.000264558\\
0.52	0.000259569\\
0.522	0.000254672\\
0.524	0.000249865\\
0.526	0.000245147\\
0.528	0.000240515\\
0.53	0.000235968\\
0.532	0.000231505\\
0.534	0.000227125\\
0.536	0.000222825\\
0.538	0.000218604\\
0.54	0.000214461\\
0.542	0.000210395\\
0.544	0.000206403\\
0.546	0.000202486\\
0.548	0.000198641\\
0.55	0.000194867\\
0.552	0.000191164\\
0.554	0.000187529\\
0.556	0.000183962\\
0.558	0.000180461\\
0.56	0.000177025\\
0.562	0.000173654\\
0.564	0.000170345\\
0.566	0.000167098\\
0.568	0.000163913\\
0.57	0.000160786\\
0.572	0.000157719\\
0.574	0.000154709\\
0.576	0.000151756\\
0.578	0.000148858\\
0.58	0.000146015\\
0.582	0.000143226\\
0.584	0.000140489\\
0.586	0.000137804\\
0.588	0.00013517\\
0.59	0.000132586\\
0.592	0.00013005\\
0.594	0.000127563\\
0.596	0.000125124\\
0.598	0.00012273\\
0.6	0.000120382\\
0.602	0.000118079\\
0.604	0.00011582\\
0.606	0.000113604\\
0.608	0.00011143\\
0.61	0.000109297\\
0.612	0.000107206\\
0.614	0.000105154\\
0.616	0.000103142\\
0.618	0.000101168\\
0.62	9.9232e-05\\
0.622	9.7333e-05\\
0.624	9.54704e-05\\
0.626	9.36435e-05\\
0.628	9.18516e-05\\
0.63	9.00941e-05\\
0.632	8.83703e-05\\
0.634	8.66796e-05\\
0.636	8.50213e-05\\
0.638	8.33949e-05\\
0.64	8.17997e-05\\
0.642	8.02351e-05\\
0.644	7.87006e-05\\
0.646	7.71955e-05\\
0.648	7.57194e-05\\
0.65	7.42716e-05\\
0.652	7.28517e-05\\
0.654	7.1459e-05\\
0.656	7.00931e-05\\
0.658	6.87535e-05\\
0.66	6.74396e-05\\
0.662	6.6151e-05\\
0.664	6.48871e-05\\
0.666	6.36475e-05\\
0.668	6.24318e-05\\
0.67	6.12395e-05\\
0.672	6.007e-05\\
0.674	5.89231e-05\\
0.676	5.77982e-05\\
0.678	5.66949e-05\\
0.68	5.56128e-05\\
0.682	5.45516e-05\\
0.684	5.35107e-05\\
0.686	5.24898e-05\\
0.688	5.14886e-05\\
0.69	5.05066e-05\\
0.692	4.95435e-05\\
0.694	4.85989e-05\\
0.696	4.76725e-05\\
0.698	4.67638e-05\\
0.7	4.58727e-05\\
0.702	4.49986e-05\\
0.704	4.41413e-05\\
0.706	4.33005e-05\\
0.708	4.24758e-05\\
0.71	4.1667e-05\\
0.712	4.08737e-05\\
0.714	4.00956e-05\\
0.716	3.93325e-05\\
0.718	3.85839e-05\\
0.72	3.78498e-05\\
0.722	3.71297e-05\\
0.724	3.64234e-05\\
0.726	3.57306e-05\\
0.728	3.50511e-05\\
0.73	3.43847e-05\\
0.732	3.3731e-05\\
0.734	3.30897e-05\\
0.736	3.24608e-05\\
0.738	3.18439e-05\\
0.74	3.12388e-05\\
0.742	3.06452e-05\\
0.744	3.0063e-05\\
0.746	2.94919e-05\\
0.748	2.89317e-05\\
0.75	2.83821e-05\\
0.752	2.78431e-05\\
0.754	2.73143e-05\\
0.756	2.67956e-05\\
0.758	2.62868e-05\\
0.76	2.57877e-05\\
0.762	2.52981e-05\\
0.764	2.48178e-05\\
0.766	2.43466e-05\\
0.768	2.38844e-05\\
0.77	2.3431e-05\\
0.772	2.29861e-05\\
0.774	2.25498e-05\\
0.776	2.21217e-05\\
0.778	2.17017e-05\\
0.78	2.12897e-05\\
0.782	2.08855e-05\\
0.784	2.0489e-05\\
0.786	2.00999e-05\\
0.788	1.97183e-05\\
0.79	1.93439e-05\\
0.792	1.89765e-05\\
0.794	1.86162e-05\\
0.796	1.82626e-05\\
0.798	1.79157e-05\\
0.8	1.75754e-05\\
0.802	1.72415e-05\\
0.804	1.6914e-05\\
0.806	1.65926e-05\\
0.808	1.62773e-05\\
0.81	1.5968e-05\\
0.812	1.56645e-05\\
0.814	1.53668e-05\\
0.816	1.50746e-05\\
0.818	1.4788e-05\\
0.82	1.45069e-05\\
0.822	1.4231e-05\\
0.824	1.39603e-05\\
0.826	1.36948e-05\\
0.828	1.34343e-05\\
0.83	1.31787e-05\\
0.832	1.29279e-05\\
0.834	1.26818e-05\\
0.836	1.24405e-05\\
0.838	1.22036e-05\\
0.84	1.19713e-05\\
0.842	1.17433e-05\\
0.844	1.15197e-05\\
0.846	1.13003e-05\\
0.848	1.1085e-05\\
0.85	1.08738e-05\\
0.852	1.06666e-05\\
0.854	1.04634e-05\\
0.856	1.02639e-05\\
0.858	1.00683e-05\\
0.86	9.87635e-06\\
0.862	9.68804e-06\\
0.864	9.50329e-06\\
0.866	9.32205e-06\\
0.868	9.14423e-06\\
0.87	8.96979e-06\\
0.872	8.79866e-06\\
0.874	8.63077e-06\\
0.876	8.46607e-06\\
0.878	8.30449e-06\\
0.88	8.14597e-06\\
0.882	7.99047e-06\\
0.884	7.83792e-06\\
0.886	7.68826e-06\\
0.888	7.54145e-06\\
0.89	7.39742e-06\\
0.892	7.25614e-06\\
0.894	7.11754e-06\\
0.896	6.98157e-06\\
0.898	6.84819e-06\\
0.9	6.71735e-06\\
0.902	6.589e-06\\
0.904	6.46309e-06\\
0.906	6.33958e-06\\
0.908	6.21843e-06\\
0.91	6.09958e-06\\
0.912	5.98299e-06\\
0.914	5.86863e-06\\
0.916	5.75645e-06\\
0.918	5.6464e-06\\
0.92	5.53846e-06\\
0.922	5.43257e-06\\
0.924	5.32871e-06\\
0.926	5.22683e-06\\
0.928	5.12689e-06\\
0.93	5.02887e-06\\
0.932	4.93271e-06\\
0.934	4.83839e-06\\
0.936	4.74588e-06\\
0.938	4.65513e-06\\
0.94	4.56612e-06\\
0.942	4.47881e-06\\
0.944	4.39317e-06\\
0.946	4.30917e-06\\
0.948	4.22678e-06\\
0.95	4.14596e-06\\
0.952	4.06669e-06\\
0.954	3.98894e-06\\
0.956	3.91268e-06\\
0.958	3.83787e-06\\
0.96	3.7645e-06\\
0.962	3.69254e-06\\
0.964	3.62195e-06\\
0.966	3.55272e-06\\
0.968	3.48481e-06\\
0.97	3.4182e-06\\
0.972	3.35287e-06\\
0.974	3.2888e-06\\
0.976	3.22595e-06\\
0.978	3.1643e-06\\
0.98	3.10384e-06\\
0.982	3.04454e-06\\
0.984	2.98637e-06\\
0.986	2.92932e-06\\
0.988	2.87336e-06\\
0.99	2.81847e-06\\
0.992	2.76464e-06\\
0.994	2.71184e-06\\
0.996	2.66005e-06\\
0.998	2.60925e-06\\
1	2.55943e-06\\
1.002	2.51056e-06\\
1.004	2.46263e-06\\
1.006	2.41562e-06\\
1.008	2.36951e-06\\
1.01	2.32428e-06\\
1.012	2.27992e-06\\
1.014	2.23641e-06\\
1.016	2.19373e-06\\
1.018	2.15187e-06\\
1.02	2.11081e-06\\
1.022	2.07054e-06\\
1.024	2.03104e-06\\
1.026	1.99229e-06\\
1.028	1.95429e-06\\
1.03	1.91701e-06\\
1.032	1.88045e-06\\
1.034	1.84459e-06\\
1.036	1.80941e-06\\
1.038	1.77491e-06\\
1.04	1.74106e-06\\
1.042	1.70786e-06\\
1.044	1.6753e-06\\
1.046	1.64336e-06\\
1.048	1.61203e-06\\
1.05	1.5813e-06\\
1.052	1.55116e-06\\
1.054	1.52159e-06\\
1.056	1.49258e-06\\
1.058	1.46414e-06\\
1.06	1.43623e-06\\
1.062	1.40885e-06\\
1.064	1.382e-06\\
1.066	1.35566e-06\\
1.068	1.32983e-06\\
1.07	1.30448e-06\\
1.072	1.27962e-06\\
1.074	1.25524e-06\\
1.076	1.23132e-06\\
1.078	1.20785e-06\\
1.08	1.18484e-06\\
1.082	1.16226e-06\\
1.084	1.14011e-06\\
1.086	1.11838e-06\\
1.088	1.09707e-06\\
1.09	1.07616e-06\\
1.092	1.05565e-06\\
1.094	1.03554e-06\\
1.096	1.0158e-06\\
1.098	9.96443e-07\\
1.1	9.77453e-07\\
1.102	9.58824e-07\\
1.104	9.4055e-07\\
1.106	9.22624e-07\\
1.108	9.05039e-07\\
1.11	8.87788e-07\\
1.112	8.70866e-07\\
1.114	8.54266e-07\\
1.116	8.37982e-07\\
1.118	8.22007e-07\\
1.12	8.06337e-07\\
1.122	7.90965e-07\\
1.124	7.75885e-07\\
1.126	7.61092e-07\\
1.128	7.46581e-07\\
1.13	7.32345e-07\\
1.132	7.18381e-07\\
1.134	7.04682e-07\\
1.136	6.91243e-07\\
1.138	6.78061e-07\\
1.14	6.65129e-07\\
1.142	6.52443e-07\\
1.144	6.39999e-07\\
1.146	6.27791e-07\\
1.148	6.15815e-07\\
1.15	6.04068e-07\\
1.152	5.92544e-07\\
1.154	5.81239e-07\\
1.156	5.70149e-07\\
1.158	5.59271e-07\\
1.16	5.48599e-07\\
1.162	5.38131e-07\\
1.164	5.27861e-07\\
1.166	5.17788e-07\\
1.168	5.07906e-07\\
1.17	4.98212e-07\\
1.172	4.88703e-07\\
1.174	4.79374e-07\\
1.176	4.70224e-07\\
1.178	4.61248e-07\\
1.18	4.52442e-07\\
1.182	4.43805e-07\\
1.184	4.35332e-07\\
1.186	4.2702e-07\\
1.188	4.18867e-07\\
1.19	4.10869e-07\\
1.192	4.03024e-07\\
1.194	3.95328e-07\\
1.196	3.87778e-07\\
1.198	3.80373e-07\\
1.2	3.73109e-07\\
1.202	3.65984e-07\\
1.204	3.58994e-07\\
1.206	3.52137e-07\\
1.208	3.45412e-07\\
1.21	3.38815e-07\\
1.212	3.32343e-07\\
1.214	3.25995e-07\\
1.216	3.19768e-07\\
1.218	3.1366e-07\\
1.22	3.07669e-07\\
1.222	3.01792e-07\\
1.224	2.96027e-07\\
1.226	2.90372e-07\\
1.228	2.84825e-07\\
1.23	2.79384e-07\\
1.232	2.74047e-07\\
1.234	2.68812e-07\\
1.236	2.63677e-07\\
1.238	2.5864e-07\\
1.24	2.53699e-07\\
1.242	2.48853e-07\\
1.244	2.44099e-07\\
1.246	2.39436e-07\\
1.248	2.34862e-07\\
1.25	2.30376e-07\\
1.252	2.25975e-07\\
1.254	2.21659e-07\\
1.256	2.17424e-07\\
1.258	2.13271e-07\\
1.26	2.09198e-07\\
1.262	2.05202e-07\\
1.264	2.01282e-07\\
1.266	1.97438e-07\\
1.268	1.93666e-07\\
1.27	1.89967e-07\\
1.272	1.86339e-07\\
1.274	1.8278e-07\\
1.276	1.79289e-07\\
1.278	1.75865e-07\\
1.28	1.72507e-07\\
1.282	1.69212e-07\\
1.284	1.65981e-07\\
1.286	1.62811e-07\\
1.288	1.59702e-07\\
1.29	1.56653e-07\\
1.292	1.53661e-07\\
1.294	1.50727e-07\\
1.296	1.47849e-07\\
1.298	1.45026e-07\\
1.3	1.42257e-07\\
1.302	1.39541e-07\\
1.304	1.36877e-07\\
1.306	1.34264e-07\\
1.308	1.31701e-07\\
1.31	1.29187e-07\\
1.312	1.26721e-07\\
1.314	1.24302e-07\\
1.316	1.21929e-07\\
1.318	1.19601e-07\\
1.32	1.17318e-07\\
1.322	1.15079e-07\\
1.324	1.12883e-07\\
1.326	1.10728e-07\\
1.328	1.08615e-07\\
1.33	1.06542e-07\\
1.332	1.04508e-07\\
1.334	1.02514e-07\\
1.336	1.00557e-07\\
1.338	9.86384e-08\\
1.34	9.6756e-08\\
1.342	9.49096e-08\\
1.344	9.30985e-08\\
1.346	9.13219e-08\\
1.348	8.95793e-08\\
1.35	8.78699e-08\\
1.352	8.61932e-08\\
1.354	8.45486e-08\\
1.356	8.29353e-08\\
1.358	8.13528e-08\\
1.36	7.98006e-08\\
1.362	7.82779e-08\\
1.364	7.67844e-08\\
1.366	7.53193e-08\\
1.368	7.38823e-08\\
1.37	7.24726e-08\\
1.372	7.10899e-08\\
1.374	6.97335e-08\\
1.376	6.8403e-08\\
1.378	6.70979e-08\\
1.38	6.58178e-08\\
1.382	6.4562e-08\\
1.384	6.33302e-08\\
1.386	6.21219e-08\\
1.388	6.09367e-08\\
1.39	5.97741e-08\\
1.392	5.86336e-08\\
1.394	5.75149e-08\\
1.396	5.64176e-08\\
1.398	5.53412e-08\\
1.4	5.42853e-08\\
1.402	5.32495e-08\\
1.404	5.22336e-08\\
1.406	5.12369e-08\\
1.408	5.02593e-08\\
1.41	4.93004e-08\\
1.412	4.83597e-08\\
1.414	4.7437e-08\\
1.416	4.65318e-08\\
1.418	4.5644e-08\\
1.42	4.4773e-08\\
1.422	4.39187e-08\\
1.424	4.30806e-08\\
1.426	4.22586e-08\\
1.428	4.14522e-08\\
1.43	4.06612e-08\\
1.432	3.98853e-08\\
1.434	3.91241e-08\\
1.436	3.83775e-08\\
1.438	3.76451e-08\\
1.44	3.69267e-08\\
1.442	3.6222e-08\\
1.444	3.55307e-08\\
1.446	3.48526e-08\\
1.448	3.41875e-08\\
1.45	3.3535e-08\\
1.452	3.2895e-08\\
1.454	3.22671e-08\\
1.456	3.16513e-08\\
1.458	3.10472e-08\\
1.46	3.04546e-08\\
1.462	2.98733e-08\\
1.464	2.93031e-08\\
1.466	2.87438e-08\\
1.468	2.81951e-08\\
1.47	2.76569e-08\\
1.472	2.7129e-08\\
1.474	2.66111e-08\\
1.476	2.61032e-08\\
1.478	2.56049e-08\\
1.48	2.51161e-08\\
1.482	2.46366e-08\\
1.484	2.41663e-08\\
1.486	2.3705e-08\\
1.488	2.32524e-08\\
1.49	2.28085e-08\\
1.492	2.23731e-08\\
1.494	2.1946e-08\\
1.496	2.1527e-08\\
1.498	2.1116e-08\\
1.5	2.07128e-08\\
1.502	2.03174e-08\\
1.504	1.99295e-08\\
1.506	1.9549e-08\\
1.508	1.91758e-08\\
1.51	1.88096e-08\\
1.512	1.84505e-08\\
1.514	1.80982e-08\\
1.516	1.77527e-08\\
1.518	1.74137e-08\\
1.52	1.70813e-08\\
1.522	1.67551e-08\\
1.524	1.64352e-08\\
1.526	1.61214e-08\\
1.528	1.58136e-08\\
1.53	1.55117e-08\\
1.532	1.52155e-08\\
1.534	1.4925e-08\\
1.536	1.464e-08\\
1.538	1.43605e-08\\
1.54	1.40863e-08\\
1.542	1.38173e-08\\
1.544	1.35535e-08\\
1.546	1.32947e-08\\
1.548	1.30409e-08\\
1.55	1.27919e-08\\
1.552	1.25476e-08\\
1.554	1.23081e-08\\
1.556	1.2073e-08\\
1.558	1.18425e-08\\
1.56	1.16164e-08\\
1.562	1.13946e-08\\
1.564	1.11771e-08\\
1.566	1.09637e-08\\
1.568	1.07543e-08\\
1.57	1.0549e-08\\
1.572	1.03476e-08\\
1.574	1.015e-08\\
1.576	9.95625e-09\\
1.578	9.76616e-09\\
1.58	9.5797e-09\\
1.582	9.3968e-09\\
1.584	9.2174e-09\\
1.586	9.04142e-09\\
1.588	8.86881e-09\\
1.59	8.69949e-09\\
1.592	8.5334e-09\\
1.594	8.37049e-09\\
1.596	8.21069e-09\\
1.598	8.05394e-09\\
1.6	7.90018e-09\\
1.602	7.74936e-09\\
1.604	7.60142e-09\\
1.606	7.45631e-09\\
1.608	7.31396e-09\\
1.61	7.17434e-09\\
1.612	7.03738e-09\\
1.614	6.90304e-09\\
1.616	6.77126e-09\\
1.618	6.642e-09\\
1.62	6.51521e-09\\
1.622	6.39084e-09\\
1.624	6.26885e-09\\
1.626	6.14918e-09\\
1.628	6.0318e-09\\
1.63	5.91666e-09\\
1.632	5.80372e-09\\
1.634	5.69294e-09\\
1.636	5.58427e-09\\
1.638	5.47768e-09\\
1.64	5.37312e-09\\
1.642	5.27056e-09\\
1.644	5.16996e-09\\
1.646	5.07128e-09\\
1.648	4.97448e-09\\
1.65	4.87953e-09\\
1.652	4.78639e-09\\
1.654	4.69503e-09\\
1.656	4.60542e-09\\
1.658	4.51751e-09\\
1.66	4.43129e-09\\
1.662	4.34671e-09\\
1.664	4.26374e-09\\
1.666	4.18236e-09\\
1.668	4.10253e-09\\
1.67	4.02423e-09\\
1.672	3.94742e-09\\
1.674	3.87208e-09\\
1.676	3.79818e-09\\
1.678	3.72568e-09\\
1.68	3.65457e-09\\
1.682	3.58482e-09\\
1.684	3.5164e-09\\
1.686	3.44928e-09\\
1.688	3.38345e-09\\
1.69	3.31887e-09\\
1.692	3.25553e-09\\
1.694	3.19339e-09\\
1.696	3.13244e-09\\
1.698	3.07265e-09\\
1.7	3.01401e-09\\
1.702	2.95648e-09\\
1.704	2.90005e-09\\
1.706	2.8447e-09\\
1.708	2.79041e-09\\
1.71	2.73715e-09\\
1.712	2.6849e-09\\
1.714	2.63366e-09\\
1.716	2.58339e-09\\
1.718	2.53408e-09\\
1.72	2.48571e-09\\
1.722	2.43827e-09\\
1.724	2.39173e-09\\
1.726	2.34608e-09\\
1.728	2.3013e-09\\
1.73	2.25738e-09\\
1.732	2.21429e-09\\
1.734	2.17203e-09\\
1.736	2.13057e-09\\
1.738	2.0899e-09\\
1.74	2.05001e-09\\
1.742	2.01088e-09\\
1.744	1.9725e-09\\
1.746	1.93485e-09\\
1.748	1.89792e-09\\
1.75	1.86169e-09\\
1.752	1.82615e-09\\
1.754	1.7913e-09\\
1.756	1.7571e-09\\
1.758	1.72356e-09\\
1.76	1.69066e-09\\
1.762	1.65839e-09\\
1.764	1.62674e-09\\
1.766	1.59568e-09\\
1.768	1.56523e-09\\
1.77	1.53535e-09\\
1.772	1.50604e-09\\
1.774	1.47729e-09\\
1.776	1.44909e-09\\
1.778	1.42143e-09\\
1.78	1.3943e-09\\
1.782	1.36768e-09\\
1.784	1.34157e-09\\
1.786	1.31596e-09\\
1.788	1.29084e-09\\
1.79	1.2662e-09\\
1.792	1.24203e-09\\
1.794	1.21832e-09\\
1.796	1.19506e-09\\
1.798	1.17225e-09\\
1.8	1.14987e-09\\
1.802	1.12792e-09\\
1.804	1.10639e-09\\
1.806	1.08527e-09\\
1.808	1.06455e-09\\
1.81	1.04423e-09\\
1.812	1.02429e-09\\
1.814	1.00474e-09\\
1.816	9.8556e-10\\
1.818	9.66745e-10\\
1.82	9.4829e-10\\
1.822	9.30187e-10\\
1.824	9.12429e-10\\
1.826	8.95011e-10\\
1.828	8.77925e-10\\
1.83	8.61165e-10\\
1.832	8.44725e-10\\
1.834	8.28599e-10\\
1.836	8.1278e-10\\
1.838	7.97264e-10\\
1.84	7.82044e-10\\
1.842	7.67114e-10\\
1.844	7.52469e-10\\
1.846	7.38104e-10\\
1.848	7.24014e-10\\
1.85	7.10192e-10\\
1.852	6.96634e-10\\
1.854	6.83335e-10\\
1.856	6.7029e-10\\
1.858	6.57493e-10\\
1.86	6.44942e-10\\
1.862	6.32629e-10\\
1.864	6.20552e-10\\
1.866	6.08706e-10\\
1.868	5.97085e-10\\
1.87	5.85687e-10\\
1.872	5.74506e-10\\
1.874	5.63538e-10\\
1.876	5.5278e-10\\
1.878	5.42227e-10\\
1.88	5.31876e-10\\
1.882	5.21722e-10\\
1.884	5.11763e-10\\
1.886	5.01993e-10\\
1.888	4.9241e-10\\
1.89	4.8301e-10\\
1.892	4.73789e-10\\
1.894	4.64745e-10\\
1.896	4.55873e-10\\
1.898	4.4717e-10\\
1.9	4.38634e-10\\
1.902	4.3026e-10\\
1.904	4.22047e-10\\
1.906	4.1399e-10\\
1.908	4.06087e-10\\
1.91	3.98335e-10\\
1.912	3.90731e-10\\
1.914	3.83272e-10\\
1.916	3.75955e-10\\
1.918	3.68779e-10\\
1.92	3.61739e-10\\
1.922	3.54833e-10\\
1.924	3.4806e-10\\
1.926	3.41416e-10\\
1.928	3.34898e-10\\
1.93	3.28505e-10\\
1.932	3.22235e-10\\
1.934	3.16083e-10\\
1.936	3.1005e-10\\
1.938	3.04131e-10\\
1.94	2.98325e-10\\
1.942	2.92631e-10\\
1.944	2.87045e-10\\
1.946	2.81565e-10\\
1.948	2.76191e-10\\
1.95	2.70918e-10\\
1.952	2.65747e-10\\
1.954	2.60674e-10\\
1.956	2.55698e-10\\
1.958	2.50817e-10\\
1.96	2.4603e-10\\
1.962	2.41333e-10\\
1.964	2.36726e-10\\
1.966	2.32208e-10\\
1.968	2.27775e-10\\
1.97	2.23427e-10\\
1.972	2.19162e-10\\
1.974	2.14979e-10\\
1.976	2.10875e-10\\
1.978	2.0685e-10\\
1.98	2.02901e-10\\
1.982	1.99028e-10\\
1.984	1.95229e-10\\
1.986	1.91503e-10\\
1.988	1.87847e-10\\
1.99	1.84261e-10\\
1.992	1.80744e-10\\
1.994	1.77294e-10\\
1.996	1.7391e-10\\
1.998	1.7059e-10\\
2	1.67334e-10\\
};
\addplot [color=mycolor9, line width=1.0pt, forget plot]
  table[row sep=crcr]{%
-0.024	0\\
-0.022	0\\
-0.02	0\\
-0.018	0\\
-0.016	0\\
-0.014	0\\
-0.012	0\\
-0.01	0\\
-0.008	0\\
-0.006	0\\
-0.004	0\\
-0.002	0\\
0	0\\
0.002	0.00916153\\
0.004	0.0147036\\
0.006	0.0198484\\
0.008	0.0237095\\
0.01	0.0263996\\
0.012	0.0280596\\
0.014	0.0288363\\
0.016	0.0288778\\
0.018	0.0283298\\
0.02	0.0273308\\
0.022	0.026009\\
0.024	0.0251931\\
0.026	0.0251402\\
0.028	0.0265757\\
0.03	0.0279031\\
0.032	0.0291209\\
0.034	0.0302281\\
0.036	0.0312244\\
0.038	0.03211\\
0.04	0.0328858\\
0.042	0.033553\\
0.044	0.0341133\\
0.046	0.0345688\\
0.048	0.0349219\\
0.05	0.0351754\\
0.052	0.0353324\\
0.054	0.035396\\
0.056	0.0353698\\
0.058	0.0352575\\
0.06	0.0350629\\
0.062	0.0347899\\
0.064	0.0344427\\
0.066	0.0340253\\
0.068	0.033542\\
0.07	0.032997\\
0.072	0.0323947\\
0.074	0.0317392\\
0.076	0.0310348\\
0.078	0.0302858\\
0.08	0.0294964\\
0.082	0.0286707\\
0.084	0.0278128\\
0.086	0.0269267\\
0.088	0.0260163\\
0.09	0.0250855\\
0.092	0.0241379\\
0.094	0.0231771\\
0.096	0.0222066\\
0.098	0.0212298\\
0.1	0.0202499\\
0.102	0.0192698\\
0.104	0.0182926\\
0.106	0.017321\\
0.108	0.0163576\\
0.11	0.0154049\\
0.112	0.0144652\\
0.114	0.0135406\\
0.116	0.0129318\\
0.118	0.0125785\\
0.12	0.0122173\\
0.122	0.0118495\\
0.124	0.0114761\\
0.126	0.0110983\\
0.128	0.0107173\\
0.13	0.010334\\
0.132	0.00994959\\
0.134	0.0096228\\
0.136	0.00957248\\
0.138	0.00951084\\
0.14	0.00943829\\
0.142	0.00935529\\
0.144	0.00926229\\
0.146	0.00915974\\
0.148	0.0090481\\
0.15	0.00892783\\
0.152	0.00879941\\
0.154	0.0086633\\
0.156	0.00851998\\
0.158	0.0083699\\
0.16	0.00821355\\
0.162	0.00805138\\
0.164	0.00788387\\
0.166	0.00771146\\
0.168	0.00753461\\
0.17	0.00735378\\
0.172	0.00716941\\
0.174	0.00698192\\
0.176	0.00679176\\
0.178	0.00659934\\
0.18	0.00640506\\
0.182	0.00620934\\
0.184	0.00601256\\
0.186	0.00589707\\
0.188	0.00579902\\
0.19	0.00570218\\
0.192	0.00560656\\
0.194	0.00551221\\
0.196	0.00541914\\
0.198	0.00532736\\
0.2	0.00523689\\
0.202	0.00514773\\
0.204	0.00505989\\
0.206	0.00497337\\
0.208	0.00488815\\
0.21	0.00480424\\
0.212	0.00472162\\
0.214	0.00464029\\
0.216	0.00456022\\
0.218	0.00448141\\
0.22	0.00440383\\
0.222	0.00432747\\
0.224	0.0042523\\
0.226	0.00417832\\
0.228	0.00410549\\
0.23	0.0040338\\
0.232	0.00396323\\
0.234	0.00389375\\
0.236	0.00382535\\
0.238	0.00375802\\
0.24	0.00369172\\
0.242	0.00362644\\
0.244	0.00356216\\
0.246	0.00349888\\
0.248	0.00343656\\
0.25	0.0033752\\
0.252	0.00331478\\
0.254	0.00325529\\
0.256	0.00319671\\
0.258	0.00313904\\
0.26	0.00308226\\
0.262	0.00302635\\
0.264	0.00297132\\
0.266	0.00291716\\
0.268	0.00286384\\
0.27	0.00281137\\
0.272	0.00275974\\
0.274	0.00270893\\
0.276	0.00265895\\
0.278	0.00260979\\
0.28	0.00256143\\
0.282	0.00251388\\
0.284	0.00246712\\
0.286	0.00242115\\
0.288	0.00237596\\
0.29	0.00233156\\
0.292	0.00228792\\
0.294	0.00224504\\
0.296	0.00220292\\
0.298	0.00216155\\
0.3	0.00212092\\
0.302	0.00208102\\
0.304	0.00204185\\
0.306	0.0020034\\
0.308	0.00196566\\
0.31	0.00192862\\
0.312	0.00189227\\
0.314	0.00185661\\
0.316	0.00182162\\
0.318	0.00178729\\
0.32	0.00175362\\
0.322	0.00172059\\
0.324	0.0016882\\
0.326	0.00165643\\
0.328	0.00162527\\
0.33	0.00159472\\
0.332	0.00156476\\
0.334	0.00153538\\
0.336	0.00150657\\
0.338	0.00147833\\
0.34	0.00145063\\
0.342	0.00142347\\
0.344	0.00139684\\
0.346	0.00137072\\
0.348	0.00134511\\
0.35	0.00132\\
0.352	0.00129538\\
0.354	0.00127123\\
0.356	0.00124754\\
0.358	0.00122431\\
0.36	0.00120153\\
0.362	0.00117919\\
0.364	0.00115727\\
0.366	0.00113577\\
0.368	0.00111467\\
0.37	0.00109398\\
0.372	0.00107368\\
0.374	0.00105376\\
0.376	0.00103422\\
0.378	0.00101505\\
0.38	0.000996231\\
0.382	0.000977768\\
0.384	0.00095965\\
0.386	0.000941869\\
0.388	0.000924419\\
0.39	0.000907294\\
0.392	0.000890486\\
0.394	0.00087399\\
0.396	0.000857799\\
0.398	0.000841907\\
0.4	0.000826309\\
0.402	0.000810999\\
0.404	0.000795972\\
0.406	0.000781221\\
0.408	0.000766743\\
0.41	0.000752531\\
0.412	0.000738582\\
0.414	0.000724889\\
0.416	0.000711449\\
0.418	0.000698257\\
0.42	0.000685308\\
0.422	0.000672598\\
0.424	0.000660122\\
0.426	0.000647877\\
0.428	0.000635858\\
0.43	0.000624062\\
0.432	0.000612483\\
0.434	0.000601119\\
0.436	0.000589965\\
0.438	0.000579017\\
0.44	0.000568273\\
0.442	0.000557727\\
0.444	0.000547377\\
0.446	0.000537219\\
0.448	0.000527249\\
0.45	0.000517464\\
0.452	0.000507861\\
0.454	0.000498436\\
0.456	0.000489186\\
0.458	0.000480108\\
0.46	0.000471198\\
0.462	0.000462453\\
0.464	0.000453871\\
0.466	0.000445447\\
0.468	0.00043718\\
0.47	0.000429066\\
0.472	0.000421102\\
0.474	0.000413286\\
0.476	0.000405614\\
0.478	0.000398084\\
0.48	0.000390693\\
0.482	0.000383439\\
0.484	0.000376318\\
0.486	0.000369329\\
0.488	0.000362469\\
0.49	0.000355735\\
0.492	0.000349124\\
0.494	0.000342636\\
0.496	0.000336266\\
0.498	0.000330014\\
0.5	0.000323876\\
0.502	0.000317851\\
0.504	0.000311937\\
0.506	0.000306131\\
0.508	0.000300431\\
0.51	0.000294835\\
0.512	0.000289342\\
0.514	0.00028395\\
0.516	0.000278656\\
0.518	0.000273459\\
0.52	0.000268357\\
0.522	0.000263348\\
0.524	0.000258431\\
0.526	0.000253604\\
0.528	0.000248866\\
0.53	0.000244214\\
0.532	0.000239647\\
0.534	0.000235163\\
0.536	0.000230762\\
0.538	0.000226442\\
0.54	0.000222201\\
0.542	0.000218038\\
0.544	0.000213951\\
0.546	0.000209939\\
0.548	0.000206001\\
0.55	0.000202135\\
0.552	0.000198341\\
0.554	0.000194617\\
0.556	0.000190961\\
0.558	0.000187373\\
0.56	0.000183851\\
0.562	0.000180394\\
0.564	0.000177001\\
0.566	0.000173671\\
0.568	0.000170403\\
0.57	0.000167196\\
0.572	0.000164048\\
0.574	0.000160959\\
0.576	0.000157927\\
0.578	0.000154952\\
0.58	0.000152032\\
0.582	0.000149167\\
0.584	0.000146355\\
0.586	0.000143596\\
0.588	0.000140888\\
0.59	0.000138231\\
0.592	0.000135624\\
0.594	0.000133066\\
0.596	0.000130556\\
0.598	0.000128093\\
0.6	0.000125676\\
0.602	0.000123305\\
0.604	0.000120978\\
0.606	0.000118695\\
0.608	0.000116455\\
0.61	0.000114257\\
0.612	0.000112101\\
0.614	0.000109985\\
0.616	0.00010791\\
0.618	0.000105873\\
0.62	0.000103875\\
0.622	0.000101915\\
0.624	9.99912e-05\\
0.626	9.81042e-05\\
0.628	9.62528e-05\\
0.63	9.44364e-05\\
0.632	9.26544e-05\\
0.634	9.0906e-05\\
0.636	8.91908e-05\\
0.638	8.75079e-05\\
0.64	8.58569e-05\\
0.642	8.42372e-05\\
0.644	8.26481e-05\\
0.646	8.10891e-05\\
0.648	7.95595e-05\\
0.65	7.8059e-05\\
0.652	7.65868e-05\\
0.654	7.51425e-05\\
0.656	7.37256e-05\\
0.658	7.23355e-05\\
0.66	7.09717e-05\\
0.662	6.96337e-05\\
0.664	6.83211e-05\\
0.666	6.70333e-05\\
0.668	6.57699e-05\\
0.67	6.45305e-05\\
0.672	6.33145e-05\\
0.674	6.21215e-05\\
0.676	6.09512e-05\\
0.678	5.9803e-05\\
0.68	5.86766e-05\\
0.682	5.75714e-05\\
0.684	5.64873e-05\\
0.686	5.54236e-05\\
0.688	5.43801e-05\\
0.69	5.33563e-05\\
0.692	5.2352e-05\\
0.694	5.13666e-05\\
0.696	5.03999e-05\\
0.698	4.94515e-05\\
0.7	4.8521e-05\\
0.702	4.76082e-05\\
0.704	4.67126e-05\\
0.706	4.5834e-05\\
0.708	4.4972e-05\\
0.71	4.41262e-05\\
0.712	4.32965e-05\\
0.714	4.24825e-05\\
0.716	4.16838e-05\\
0.718	4.09003e-05\\
0.72	4.01315e-05\\
0.722	3.93773e-05\\
0.724	3.86373e-05\\
0.726	3.79113e-05\\
0.728	3.7199e-05\\
0.73	3.65001e-05\\
0.732	3.58144e-05\\
0.734	3.51416e-05\\
0.736	3.44816e-05\\
0.738	3.38339e-05\\
0.74	3.31985e-05\\
0.742	3.2575e-05\\
0.744	3.19633e-05\\
0.746	3.13631e-05\\
0.748	3.07741e-05\\
0.75	3.01963e-05\\
0.752	2.96293e-05\\
0.754	2.9073e-05\\
0.756	2.85271e-05\\
0.758	2.79915e-05\\
0.76	2.7466e-05\\
0.762	2.69503e-05\\
0.764	2.64443e-05\\
0.766	2.59478e-05\\
0.768	2.54606e-05\\
0.77	2.49826e-05\\
0.772	2.45135e-05\\
0.774	2.40532e-05\\
0.776	2.36015e-05\\
0.778	2.31583e-05\\
0.78	2.27234e-05\\
0.782	2.22966e-05\\
0.784	2.18778e-05\\
0.786	2.14669e-05\\
0.788	2.10637e-05\\
0.79	2.0668e-05\\
0.792	2.02797e-05\\
0.794	1.98987e-05\\
0.796	1.95248e-05\\
0.798	1.91579e-05\\
0.8	1.87978e-05\\
0.802	1.84445e-05\\
0.804	1.80978e-05\\
0.806	1.77576e-05\\
0.808	1.74237e-05\\
0.81	1.70961e-05\\
0.812	1.67746e-05\\
0.814	1.64591e-05\\
0.816	1.61496e-05\\
0.818	1.58458e-05\\
0.82	1.55477e-05\\
0.822	1.52552e-05\\
0.824	1.49681e-05\\
0.826	1.46864e-05\\
0.828	1.441e-05\\
0.83	1.41387e-05\\
0.832	1.38726e-05\\
0.834	1.36114e-05\\
0.836	1.33551e-05\\
0.838	1.31036e-05\\
0.84	1.28568e-05\\
0.842	1.26146e-05\\
0.844	1.2377e-05\\
0.846	1.21438e-05\\
0.848	1.1915e-05\\
0.85	1.16905e-05\\
0.852	1.14701e-05\\
0.854	1.12539e-05\\
0.856	1.10418e-05\\
0.858	1.08336e-05\\
0.86	1.06294e-05\\
0.862	1.0429e-05\\
0.864	1.02323e-05\\
0.866	1.00393e-05\\
0.868	9.85e-06\\
0.87	9.6642e-06\\
0.872	9.4819e-06\\
0.874	9.30303e-06\\
0.876	9.12751e-06\\
0.878	8.95529e-06\\
0.88	8.78631e-06\\
0.882	8.62051e-06\\
0.884	8.45782e-06\\
0.886	8.2982e-06\\
0.888	8.14157e-06\\
0.89	7.9879e-06\\
0.892	7.83711e-06\\
0.894	7.68917e-06\\
0.896	7.54401e-06\\
0.898	7.40159e-06\\
0.9	7.26184e-06\\
0.902	7.12473e-06\\
0.904	6.99021e-06\\
0.906	6.85822e-06\\
0.908	6.72871e-06\\
0.91	6.60165e-06\\
0.912	6.47699e-06\\
0.914	6.35467e-06\\
0.916	6.23467e-06\\
0.918	6.11693e-06\\
0.92	6.0014e-06\\
0.922	5.88806e-06\\
0.924	5.77686e-06\\
0.926	5.66776e-06\\
0.928	5.56072e-06\\
0.93	5.4557e-06\\
0.932	5.35266e-06\\
0.934	5.25157e-06\\
0.936	5.15239e-06\\
0.938	5.05508e-06\\
0.94	4.95962e-06\\
0.942	4.86595e-06\\
0.944	4.77406e-06\\
0.946	4.6839e-06\\
0.948	4.59545e-06\\
0.95	4.50867e-06\\
0.952	4.42353e-06\\
0.954	4.34e-06\\
0.956	4.25804e-06\\
0.958	4.17764e-06\\
0.96	4.09876e-06\\
0.962	4.02137e-06\\
0.964	3.94545e-06\\
0.966	3.87096e-06\\
0.968	3.79788e-06\\
0.97	3.72618e-06\\
0.972	3.65583e-06\\
0.974	3.58682e-06\\
0.976	3.51912e-06\\
0.978	3.45269e-06\\
0.98	3.38752e-06\\
0.982	3.32358e-06\\
0.984	3.26085e-06\\
0.986	3.19931e-06\\
0.988	3.13893e-06\\
0.99	3.0797e-06\\
0.992	3.02158e-06\\
0.994	2.96457e-06\\
0.996	2.90863e-06\\
0.998	2.85375e-06\\
1	2.7999e-06\\
1.002	2.74708e-06\\
1.004	2.69525e-06\\
1.006	2.6444e-06\\
1.008	2.59452e-06\\
1.01	2.54557e-06\\
1.012	2.49755e-06\\
1.014	2.45044e-06\\
1.016	2.40422e-06\\
1.018	2.35887e-06\\
1.02	2.31438e-06\\
1.022	2.27073e-06\\
1.024	2.22791e-06\\
1.026	2.18589e-06\\
1.028	2.14466e-06\\
1.03	2.10422e-06\\
1.032	2.06454e-06\\
1.034	2.02561e-06\\
1.036	1.98741e-06\\
1.038	1.94993e-06\\
1.04	1.91316e-06\\
1.042	1.87709e-06\\
1.044	1.84169e-06\\
1.046	1.80696e-06\\
1.048	1.77289e-06\\
1.05	1.73946e-06\\
1.052	1.70667e-06\\
1.054	1.67449e-06\\
1.056	1.64291e-06\\
1.058	1.61194e-06\\
1.06	1.58154e-06\\
1.062	1.55172e-06\\
1.064	1.52246e-06\\
1.066	1.49376e-06\\
1.068	1.46559e-06\\
1.07	1.43796e-06\\
1.072	1.41084e-06\\
1.074	1.38424e-06\\
1.076	1.35814e-06\\
1.078	1.33253e-06\\
1.08	1.30741e-06\\
1.082	1.28275e-06\\
1.084	1.25856e-06\\
1.086	1.23483e-06\\
1.088	1.21154e-06\\
1.09	1.1887e-06\\
1.092	1.16628e-06\\
1.094	1.14429e-06\\
1.096	1.12271e-06\\
1.098	1.10153e-06\\
1.1	1.08076e-06\\
1.102	1.06038e-06\\
1.104	1.04038e-06\\
1.106	1.02075e-06\\
1.108	1.0015e-06\\
1.11	9.82609e-07\\
1.112	9.64075e-07\\
1.114	9.45889e-07\\
1.116	9.28046e-07\\
1.118	9.10539e-07\\
1.12	8.93362e-07\\
1.122	8.76509e-07\\
1.124	8.59972e-07\\
1.126	8.43748e-07\\
1.128	8.27829e-07\\
1.13	8.1221e-07\\
1.132	7.96885e-07\\
1.134	7.81848e-07\\
1.136	7.67095e-07\\
1.138	7.5262e-07\\
1.14	7.38418e-07\\
1.142	7.24483e-07\\
1.144	7.10811e-07\\
1.146	6.97397e-07\\
1.148	6.84235e-07\\
1.15	6.71321e-07\\
1.152	6.58651e-07\\
1.154	6.46219e-07\\
1.156	6.34022e-07\\
1.158	6.22055e-07\\
1.16	6.10313e-07\\
1.162	5.98792e-07\\
1.164	5.87489e-07\\
1.166	5.76398e-07\\
1.168	5.65517e-07\\
1.17	5.54841e-07\\
1.172	5.44366e-07\\
1.174	5.34089e-07\\
1.176	5.24006e-07\\
1.178	5.14112e-07\\
1.18	5.04406e-07\\
1.182	4.94882e-07\\
1.184	4.85538e-07\\
1.186	4.7637e-07\\
1.188	4.67375e-07\\
1.19	4.5855e-07\\
1.192	4.49891e-07\\
1.194	4.41396e-07\\
1.196	4.33061e-07\\
1.198	4.24883e-07\\
1.2	4.1686e-07\\
1.202	4.08988e-07\\
1.204	4.01264e-07\\
1.206	3.93686e-07\\
1.208	3.86252e-07\\
1.21	3.78957e-07\\
1.212	3.718e-07\\
1.214	3.64779e-07\\
1.216	3.5789e-07\\
1.218	3.51131e-07\\
1.22	3.44499e-07\\
1.222	3.37993e-07\\
1.224	3.3161e-07\\
1.226	3.25347e-07\\
1.228	3.19202e-07\\
1.23	3.13174e-07\\
1.232	3.07259e-07\\
1.234	3.01456e-07\\
1.236	2.95763e-07\\
1.238	2.90177e-07\\
1.24	2.84696e-07\\
1.242	2.7932e-07\\
1.244	2.74044e-07\\
1.246	2.68869e-07\\
1.248	2.63791e-07\\
1.25	2.58809e-07\\
1.252	2.53921e-07\\
1.254	2.49125e-07\\
1.256	2.4442e-07\\
1.258	2.39804e-07\\
1.26	2.35275e-07\\
1.262	2.30832e-07\\
1.264	2.26473e-07\\
1.266	2.22196e-07\\
1.268	2.17999e-07\\
1.27	2.13883e-07\\
1.272	2.09843e-07\\
1.274	2.05881e-07\\
1.276	2.01993e-07\\
1.278	1.98178e-07\\
1.28	1.94436e-07\\
1.282	1.90764e-07\\
1.284	1.87162e-07\\
1.286	1.83627e-07\\
1.288	1.8016e-07\\
1.29	1.76758e-07\\
1.292	1.7342e-07\\
1.294	1.70145e-07\\
1.296	1.66932e-07\\
1.298	1.6378e-07\\
1.3	1.60688e-07\\
1.302	1.57654e-07\\
1.304	1.54677e-07\\
1.306	1.51756e-07\\
1.308	1.48891e-07\\
1.31	1.46079e-07\\
1.312	1.43321e-07\\
1.314	1.40615e-07\\
1.316	1.3796e-07\\
1.318	1.35355e-07\\
1.32	1.328e-07\\
1.322	1.30292e-07\\
1.324	1.27832e-07\\
1.326	1.25419e-07\\
1.328	1.23051e-07\\
1.33	1.20728e-07\\
1.332	1.18448e-07\\
1.334	1.16212e-07\\
1.336	1.14018e-07\\
1.338	1.11865e-07\\
1.34	1.09753e-07\\
1.342	1.07681e-07\\
1.344	1.05648e-07\\
1.346	1.03654e-07\\
1.348	1.01697e-07\\
1.35	9.9777e-08\\
1.352	9.78933e-08\\
1.354	9.60452e-08\\
1.356	9.4232e-08\\
1.358	9.2453e-08\\
1.36	9.07076e-08\\
1.362	8.89951e-08\\
1.364	8.7315e-08\\
1.366	8.56666e-08\\
1.368	8.40494e-08\\
1.37	8.24626e-08\\
1.372	8.09058e-08\\
1.374	7.93784e-08\\
1.376	7.78799e-08\\
1.378	7.64096e-08\\
1.38	7.49671e-08\\
1.382	7.35518e-08\\
1.384	7.21632e-08\\
1.386	7.08009e-08\\
1.388	6.94642e-08\\
1.39	6.81528e-08\\
1.392	6.68662e-08\\
1.394	6.56038e-08\\
1.396	6.43652e-08\\
1.398	6.31501e-08\\
1.4	6.19579e-08\\
1.402	6.07881e-08\\
1.404	5.96405e-08\\
1.406	5.85145e-08\\
1.408	5.74098e-08\\
1.41	5.63259e-08\\
1.412	5.52625e-08\\
1.414	5.42191e-08\\
1.416	5.31954e-08\\
1.418	5.21911e-08\\
1.42	5.12057e-08\\
1.422	5.02389e-08\\
1.424	4.92904e-08\\
1.426	4.83597e-08\\
1.428	4.74467e-08\\
1.43	4.65508e-08\\
1.432	4.56719e-08\\
1.434	4.48095e-08\\
1.436	4.39635e-08\\
1.438	4.31334e-08\\
1.44	4.23189e-08\\
1.442	4.15199e-08\\
1.444	4.07359e-08\\
1.446	3.99667e-08\\
1.448	3.92121e-08\\
1.45	3.84716e-08\\
1.452	3.77452e-08\\
1.454	3.70325e-08\\
1.456	3.63332e-08\\
1.458	3.56471e-08\\
1.46	3.4974e-08\\
1.462	3.43136e-08\\
1.464	3.36656e-08\\
1.466	3.30299e-08\\
1.468	3.24061e-08\\
1.47	3.17942e-08\\
1.472	3.11938e-08\\
1.474	3.06047e-08\\
1.476	3.00268e-08\\
1.478	2.94597e-08\\
1.48	2.89034e-08\\
1.482	2.83576e-08\\
1.484	2.7822e-08\\
1.486	2.72966e-08\\
1.488	2.67811e-08\\
1.49	2.62754e-08\\
1.492	2.57791e-08\\
1.494	2.52923e-08\\
1.496	2.48146e-08\\
1.498	2.4346e-08\\
1.5	2.38862e-08\\
1.502	2.34351e-08\\
1.504	2.29925e-08\\
1.506	2.25583e-08\\
1.508	2.21322e-08\\
1.51	2.17142e-08\\
1.512	2.13041e-08\\
1.514	2.09018e-08\\
1.516	2.0507e-08\\
1.518	2.01197e-08\\
1.52	1.97397e-08\\
1.522	1.93669e-08\\
1.524	1.90011e-08\\
1.526	1.86423e-08\\
1.528	1.82902e-08\\
1.53	1.79447e-08\\
1.532	1.76058e-08\\
1.534	1.72733e-08\\
1.536	1.69471e-08\\
1.538	1.6627e-08\\
1.54	1.63129e-08\\
1.542	1.60048e-08\\
1.544	1.57026e-08\\
1.546	1.5406e-08\\
1.548	1.5115e-08\\
1.55	1.48295e-08\\
1.552	1.45495e-08\\
1.554	1.42747e-08\\
1.556	1.4005e-08\\
1.558	1.37405e-08\\
1.56	1.3481e-08\\
1.562	1.32264e-08\\
1.564	1.29766e-08\\
1.566	1.27315e-08\\
1.568	1.2491e-08\\
1.57	1.22551e-08\\
1.572	1.20237e-08\\
1.574	1.17966e-08\\
1.576	1.15738e-08\\
1.578	1.13552e-08\\
1.58	1.11407e-08\\
1.582	1.09303e-08\\
1.584	1.07239e-08\\
1.586	1.05213e-08\\
1.588	1.03226e-08\\
1.59	1.01277e-08\\
1.592	9.93639e-09\\
1.594	9.74873e-09\\
1.596	9.56461e-09\\
1.598	9.38397e-09\\
1.6	9.20674e-09\\
1.602	9.03286e-09\\
1.604	8.86226e-09\\
1.606	8.69489e-09\\
1.608	8.53067e-09\\
1.61	8.36956e-09\\
1.612	8.21149e-09\\
1.614	8.05641e-09\\
1.616	7.90426e-09\\
1.618	7.75498e-09\\
1.62	7.60852e-09\\
1.622	7.46483e-09\\
1.624	7.32385e-09\\
1.626	7.18553e-09\\
1.628	7.04983e-09\\
1.63	6.91669e-09\\
1.632	6.78606e-09\\
1.634	6.6579e-09\\
1.636	6.53217e-09\\
1.638	6.4088e-09\\
1.64	6.28777e-09\\
1.642	6.16902e-09\\
1.644	6.05252e-09\\
1.646	5.93822e-09\\
1.648	5.82607e-09\\
1.65	5.71604e-09\\
1.652	5.6081e-09\\
1.654	5.50219e-09\\
1.656	5.39828e-09\\
1.658	5.29633e-09\\
1.66	5.19631e-09\\
1.662	5.09817e-09\\
1.664	5.0019e-09\\
1.666	4.90743e-09\\
1.668	4.81476e-09\\
1.67	4.72383e-09\\
1.672	4.63462e-09\\
1.674	4.5471e-09\\
1.676	4.46123e-09\\
1.678	4.37698e-09\\
1.68	4.29432e-09\\
1.682	4.21322e-09\\
1.684	4.13365e-09\\
1.686	4.05559e-09\\
1.688	3.979e-09\\
1.69	3.90386e-09\\
1.692	3.83013e-09\\
1.694	3.7578e-09\\
1.696	3.68684e-09\\
1.698	3.61721e-09\\
1.7	3.5489e-09\\
1.702	3.48188e-09\\
1.704	3.41612e-09\\
1.706	3.35161e-09\\
1.708	3.28831e-09\\
1.71	3.22621e-09\\
1.712	3.16529e-09\\
1.714	3.10551e-09\\
1.716	3.04686e-09\\
1.718	2.98932e-09\\
1.72	2.93287e-09\\
1.722	2.87748e-09\\
1.724	2.82314e-09\\
1.726	2.76982e-09\\
1.728	2.71751e-09\\
1.73	2.66619e-09\\
1.732	2.61584e-09\\
1.734	2.56644e-09\\
1.736	2.51797e-09\\
1.738	2.47042e-09\\
1.74	2.42376e-09\\
1.742	2.37799e-09\\
1.744	2.33308e-09\\
1.746	2.28902e-09\\
1.748	2.24579e-09\\
1.75	2.20337e-09\\
1.752	2.16176e-09\\
1.754	2.12094e-09\\
1.756	2.08088e-09\\
1.758	2.04158e-09\\
1.76	2.00302e-09\\
1.762	1.9652e-09\\
1.764	1.92808e-09\\
1.766	1.89167e-09\\
1.768	1.85594e-09\\
1.77	1.82089e-09\\
1.772	1.7865e-09\\
1.774	1.75276e-09\\
1.776	1.71966e-09\\
1.778	1.68718e-09\\
1.78	1.65531e-09\\
1.782	1.62405e-09\\
1.784	1.59338e-09\\
1.786	1.56329e-09\\
1.788	1.53376e-09\\
1.79	1.50479e-09\\
1.792	1.47637e-09\\
1.794	1.44849e-09\\
1.796	1.42113e-09\\
1.798	1.39429e-09\\
1.8	1.36796e-09\\
1.802	1.34212e-09\\
1.804	1.31677e-09\\
1.806	1.2919e-09\\
1.808	1.2675e-09\\
1.81	1.24357e-09\\
1.812	1.22008e-09\\
1.814	1.19703e-09\\
1.816	1.17443e-09\\
1.818	1.15225e-09\\
1.82	1.13048e-09\\
1.822	1.10913e-09\\
1.824	1.08818e-09\\
1.826	1.06763e-09\\
1.828	1.04747e-09\\
1.83	1.02768e-09\\
1.832	1.00827e-09\\
1.834	9.89229e-10\\
1.836	9.70545e-10\\
1.838	9.52214e-10\\
1.84	9.3423e-10\\
1.842	9.16585e-10\\
1.844	8.99273e-10\\
1.846	8.82288e-10\\
1.848	8.65624e-10\\
1.85	8.49275e-10\\
1.852	8.33235e-10\\
1.854	8.17497e-10\\
1.856	8.02057e-10\\
1.858	7.86908e-10\\
1.86	7.72045e-10\\
1.862	7.57464e-10\\
1.864	7.43157e-10\\
1.866	7.29121e-10\\
1.868	7.1535e-10\\
1.87	7.01839e-10\\
1.872	6.88583e-10\\
1.874	6.75578e-10\\
1.876	6.62818e-10\\
1.878	6.50299e-10\\
1.88	6.38017e-10\\
1.882	6.25966e-10\\
1.884	6.14144e-10\\
1.886	6.02544e-10\\
1.888	5.91164e-10\\
1.89	5.79998e-10\\
1.892	5.69044e-10\\
1.894	5.58296e-10\\
1.896	5.47751e-10\\
1.898	5.37406e-10\\
1.9	5.27256e-10\\
1.902	5.17297e-10\\
1.904	5.07527e-10\\
1.906	4.97942e-10\\
1.908	4.88537e-10\\
1.91	4.7931e-10\\
1.912	4.70257e-10\\
1.914	4.61375e-10\\
1.916	4.52661e-10\\
1.918	4.44112e-10\\
1.92	4.35724e-10\\
1.922	4.27494e-10\\
1.924	4.1942e-10\\
1.926	4.11499e-10\\
1.928	4.03727e-10\\
1.93	3.96102e-10\\
1.932	3.8862e-10\\
1.934	3.81281e-10\\
1.936	3.74079e-10\\
1.938	3.67014e-10\\
1.94	3.60082e-10\\
1.942	3.53282e-10\\
1.944	3.46609e-10\\
1.946	3.40063e-10\\
1.948	3.3364e-10\\
1.95	3.27339e-10\\
1.952	3.21156e-10\\
1.954	3.15091e-10\\
1.956	3.0914e-10\\
1.958	3.03301e-10\\
1.96	2.97573e-10\\
1.962	2.91952e-10\\
1.964	2.86438e-10\\
1.966	2.81029e-10\\
1.968	2.75721e-10\\
1.97	2.70513e-10\\
1.972	2.65404e-10\\
1.974	2.60392e-10\\
1.976	2.55474e-10\\
1.978	2.50649e-10\\
1.98	2.45915e-10\\
1.982	2.4127e-10\\
1.984	2.36713e-10\\
1.986	2.32243e-10\\
1.988	2.27856e-10\\
1.99	2.23553e-10\\
1.992	2.19331e-10\\
1.994	2.15188e-10\\
1.996	2.11124e-10\\
1.998	2.07137e-10\\
2	2.03225e-10\\
};
\addplot [color=mycolor10, line width=1.0pt, forget plot]
  table[row sep=crcr]{%
-0.024	0\\
-0.022	0\\
-0.02	0\\
-0.018	0\\
-0.016	0\\
-0.014	0\\
-0.012	0\\
-0.01	0\\
-0.008	0\\
-0.006	0\\
-0.004	0\\
-0.002	0\\
0	0\\
0.002	0.00916134\\
0.004	0.0147017\\
0.006	0.0187732\\
0.008	0.0221002\\
0.01	0.0243074\\
0.012	0.025585\\
0.014	0.0261095\\
0.016	0.0260421\\
0.018	0.0255269\\
0.02	0.0252376\\
0.022	0.0249224\\
0.024	0.0248311\\
0.026	0.0252312\\
0.028	0.0266549\\
0.03	0.0279681\\
0.032	0.0291695\\
0.034	0.0302586\\
0.036	0.0312357\\
0.038	0.0321013\\
0.04	0.0328566\\
0.042	0.0335033\\
0.044	0.0340433\\
0.046	0.0344789\\
0.048	0.0348128\\
0.05	0.0350479\\
0.052	0.0351873\\
0.054	0.0352343\\
0.056	0.0351926\\
0.058	0.0350658\\
0.06	0.0348579\\
0.062	0.0345726\\
0.064	0.0342141\\
0.066	0.0337866\\
0.068	0.0332942\\
0.07	0.0327412\\
0.072	0.0321318\\
0.074	0.0314703\\
0.076	0.0307609\\
0.078	0.0300079\\
0.08	0.0292154\\
0.082	0.0283876\\
0.084	0.0275287\\
0.086	0.0266425\\
0.088	0.0257329\\
0.09	0.0248039\\
0.092	0.023859\\
0.094	0.0229019\\
0.096	0.021936\\
0.098	0.0209646\\
0.1	0.0199909\\
0.102	0.0190178\\
0.104	0.0180484\\
0.106	0.0170853\\
0.108	0.0161311\\
0.11	0.0151881\\
0.112	0.0142587\\
0.114	0.013345\\
0.116	0.0127211\\
0.118	0.0123757\\
0.12	0.012023\\
0.122	0.0116639\\
0.124	0.0112997\\
0.126	0.0109314\\
0.128	0.0105601\\
0.13	0.0101869\\
0.132	0.00981276\\
0.134	0.00944339\\
0.136	0.00939974\\
0.138	0.00934505\\
0.14	0.00927972\\
0.142	0.00920418\\
0.144	0.00911885\\
0.146	0.00902416\\
0.148	0.00892056\\
0.15	0.00880848\\
0.152	0.00868837\\
0.154	0.00856068\\
0.156	0.00842585\\
0.158	0.00828433\\
0.16	0.00813657\\
0.162	0.00798301\\
0.164	0.0078241\\
0.166	0.00766028\\
0.168	0.00749198\\
0.17	0.00731964\\
0.172	0.00714368\\
0.174	0.00696451\\
0.176	0.00678255\\
0.178	0.00659821\\
0.18	0.00641188\\
0.182	0.00622395\\
0.184	0.00603479\\
0.186	0.00585983\\
0.188	0.00576471\\
0.19	0.00567078\\
0.192	0.00557806\\
0.194	0.00548658\\
0.196	0.00539634\\
0.198	0.00530736\\
0.2	0.00521964\\
0.202	0.00513319\\
0.204	0.00504801\\
0.206	0.00496408\\
0.208	0.00488141\\
0.21	0.00479998\\
0.212	0.00471978\\
0.214	0.00464079\\
0.216	0.00456301\\
0.218	0.00448642\\
0.22	0.00441099\\
0.222	0.00433672\\
0.224	0.00426357\\
0.226	0.00419153\\
0.228	0.00412059\\
0.23	0.00405072\\
0.232	0.0039819\\
0.234	0.00391412\\
0.236	0.00384735\\
0.238	0.00378158\\
0.24	0.00371678\\
0.242	0.00365296\\
0.244	0.00359008\\
0.246	0.00352813\\
0.248	0.00346709\\
0.25	0.00340697\\
0.252	0.00334773\\
0.254	0.00328937\\
0.256	0.00323188\\
0.258	0.00317524\\
0.26	0.00311945\\
0.262	0.0030645\\
0.264	0.00301037\\
0.266	0.00295706\\
0.268	0.00290457\\
0.27	0.00285288\\
0.272	0.00280199\\
0.274	0.00275189\\
0.276	0.00270257\\
0.278	0.00265404\\
0.28	0.00260627\\
0.282	0.00255927\\
0.284	0.00251303\\
0.286	0.00246755\\
0.288	0.00242281\\
0.29	0.00237882\\
0.292	0.00233556\\
0.294	0.00229303\\
0.296	0.00225123\\
0.298	0.00221014\\
0.3	0.00216976\\
0.302	0.00213008\\
0.304	0.0020911\\
0.306	0.00205281\\
0.308	0.00201519\\
0.31	0.00197825\\
0.312	0.00194197\\
0.314	0.00190634\\
0.316	0.00187136\\
0.318	0.00183702\\
0.32	0.0018033\\
0.322	0.0017702\\
0.324	0.00173771\\
0.326	0.00170582\\
0.328	0.00167452\\
0.33	0.0016438\\
0.332	0.00161364\\
0.334	0.00158405\\
0.336	0.00155501\\
0.338	0.0015265\\
0.34	0.00149853\\
0.342	0.00147107\\
0.344	0.00144413\\
0.346	0.00141769\\
0.348	0.00139173\\
0.35	0.00136626\\
0.352	0.00134126\\
0.354	0.00131672\\
0.356	0.00129264\\
0.358	0.001269\\
0.36	0.00124579\\
0.362	0.00122302\\
0.364	0.00120066\\
0.366	0.00117871\\
0.368	0.00115716\\
0.37	0.001136\\
0.372	0.00111523\\
0.374	0.00109484\\
0.376	0.00107483\\
0.378	0.00105517\\
0.38	0.00103587\\
0.382	0.00101692\\
0.384	0.000998317\\
0.386	0.000980049\\
0.388	0.000962111\\
0.39	0.000944497\\
0.392	0.000927202\\
0.394	0.000910219\\
0.396	0.000893544\\
0.398	0.00087717\\
0.4	0.000861091\\
0.402	0.000845304\\
0.404	0.000829802\\
0.406	0.000814581\\
0.408	0.000799635\\
0.41	0.00078496\\
0.412	0.000770551\\
0.414	0.000756404\\
0.416	0.000742514\\
0.418	0.000728876\\
0.42	0.000715486\\
0.422	0.00070234\\
0.424	0.000689434\\
0.426	0.000676764\\
0.428	0.000664325\\
0.43	0.000652113\\
0.432	0.000640124\\
0.434	0.000628356\\
0.436	0.000616803\\
0.438	0.000605461\\
0.44	0.000594328\\
0.442	0.0005834\\
0.444	0.000572672\\
0.446	0.000562142\\
0.448	0.000551805\\
0.45	0.000541658\\
0.452	0.000531699\\
0.454	0.000521922\\
0.456	0.000512326\\
0.458	0.000502906\\
0.46	0.00049366\\
0.462	0.000484584\\
0.464	0.000475674\\
0.466	0.000466929\\
0.468	0.000458345\\
0.47	0.000449919\\
0.472	0.000441647\\
0.474	0.000433528\\
0.476	0.000425557\\
0.478	0.000417733\\
0.48	0.000410052\\
0.482	0.000402512\\
0.484	0.00039511\\
0.486	0.000387844\\
0.488	0.000380711\\
0.49	0.000373708\\
0.492	0.000366833\\
0.494	0.000360083\\
0.496	0.000353457\\
0.498	0.000346951\\
0.5	0.000340564\\
0.502	0.000334293\\
0.504	0.000328137\\
0.506	0.000322092\\
0.508	0.000316158\\
0.51	0.000310331\\
0.512	0.00030461\\
0.514	0.000298993\\
0.516	0.000293479\\
0.518	0.000288064\\
0.52	0.000282748\\
0.522	0.000277528\\
0.524	0.000272403\\
0.526	0.000267371\\
0.528	0.00026243\\
0.53	0.00025758\\
0.532	0.000252817\\
0.534	0.000248141\\
0.536	0.00024355\\
0.538	0.000239042\\
0.54	0.000234617\\
0.542	0.000230272\\
0.544	0.000226006\\
0.546	0.000221818\\
0.548	0.000217707\\
0.55	0.00021367\\
0.552	0.000209707\\
0.554	0.000205817\\
0.556	0.000201997\\
0.558	0.000198248\\
0.56	0.000194567\\
0.562	0.000190954\\
0.564	0.000187407\\
0.566	0.000183925\\
0.568	0.000180507\\
0.57	0.000177152\\
0.572	0.000173858\\
0.574	0.000170625\\
0.576	0.000167452\\
0.578	0.000164337\\
0.58	0.00016128\\
0.582	0.000158279\\
0.584	0.000155334\\
0.586	0.000152443\\
0.588	0.000149606\\
0.59	0.000146821\\
0.592	0.000144087\\
0.594	0.000141405\\
0.596	0.000138772\\
0.598	0.000136188\\
0.6	0.000133651\\
0.602	0.000131162\\
0.604	0.00012872\\
0.606	0.000126322\\
0.608	0.00012397\\
0.61	0.00012166\\
0.612	0.000119394\\
0.614	0.000117171\\
0.616	0.000114988\\
0.618	0.000112846\\
0.62	0.000110744\\
0.622	0.000108681\\
0.624	0.000106657\\
0.626	0.000104671\\
0.628	0.000102721\\
0.63	0.000100808\\
0.632	9.89303e-05\\
0.634	9.70878e-05\\
0.636	9.52797e-05\\
0.638	9.35053e-05\\
0.64	9.1764e-05\\
0.642	9.00553e-05\\
0.644	8.83784e-05\\
0.646	8.67328e-05\\
0.648	8.51179e-05\\
0.65	8.35332e-05\\
0.652	8.19781e-05\\
0.654	8.0452e-05\\
0.656	7.89544e-05\\
0.658	7.74848e-05\\
0.66	7.60426e-05\\
0.662	7.46273e-05\\
0.664	7.32385e-05\\
0.666	7.18756e-05\\
0.668	7.05381e-05\\
0.67	6.92257e-05\\
0.672	6.79377e-05\\
0.674	6.66738e-05\\
0.676	6.54335e-05\\
0.678	6.42163e-05\\
0.68	6.30219e-05\\
0.682	6.18497e-05\\
0.684	6.06995e-05\\
0.686	5.95707e-05\\
0.688	5.8463e-05\\
0.69	5.7376e-05\\
0.692	5.63093e-05\\
0.694	5.52625e-05\\
0.696	5.42352e-05\\
0.698	5.32271e-05\\
0.7	5.22378e-05\\
0.702	5.12669e-05\\
0.704	5.03142e-05\\
0.706	4.93792e-05\\
0.708	4.84616e-05\\
0.71	4.75612e-05\\
0.712	4.66776e-05\\
0.714	4.58104e-05\\
0.716	4.49594e-05\\
0.718	4.41242e-05\\
0.72	4.33046e-05\\
0.722	4.25002e-05\\
0.724	4.17109e-05\\
0.726	4.09362e-05\\
0.728	4.01759e-05\\
0.73	3.94298e-05\\
0.732	3.86976e-05\\
0.734	3.7979e-05\\
0.736	3.72737e-05\\
0.738	3.65816e-05\\
0.74	3.59023e-05\\
0.742	3.52356e-05\\
0.744	3.45813e-05\\
0.746	3.39392e-05\\
0.748	3.3309e-05\\
0.75	3.26905e-05\\
0.752	3.20834e-05\\
0.754	3.14877e-05\\
0.756	3.09029e-05\\
0.758	3.03291e-05\\
0.76	2.97658e-05\\
0.762	2.9213e-05\\
0.764	2.86705e-05\\
0.766	2.8138e-05\\
0.768	2.76153e-05\\
0.77	2.71024e-05\\
0.772	2.65989e-05\\
0.774	2.61048e-05\\
0.776	2.56198e-05\\
0.778	2.51438e-05\\
0.78	2.46766e-05\\
0.782	2.4218e-05\\
0.784	2.3768e-05\\
0.786	2.33262e-05\\
0.788	2.28927e-05\\
0.79	2.24671e-05\\
0.792	2.20495e-05\\
0.794	2.16395e-05\\
0.796	2.12372e-05\\
0.798	2.08423e-05\\
0.8	2.04547e-05\\
0.802	2.00742e-05\\
0.804	1.97008e-05\\
0.806	1.93344e-05\\
0.808	1.89747e-05\\
0.81	1.86216e-05\\
0.812	1.82751e-05\\
0.814	1.7935e-05\\
0.816	1.76012e-05\\
0.818	1.72736e-05\\
0.82	1.6952e-05\\
0.822	1.66364e-05\\
0.824	1.63266e-05\\
0.826	1.60226e-05\\
0.828	1.57242e-05\\
0.83	1.54314e-05\\
0.832	1.5144e-05\\
0.834	1.48618e-05\\
0.836	1.4585e-05\\
0.838	1.43132e-05\\
0.84	1.40465e-05\\
0.842	1.37848e-05\\
0.844	1.35279e-05\\
0.846	1.32758e-05\\
0.848	1.30283e-05\\
0.85	1.27854e-05\\
0.852	1.25471e-05\\
0.854	1.23132e-05\\
0.856	1.20836e-05\\
0.858	1.18583e-05\\
0.86	1.16372e-05\\
0.862	1.14202e-05\\
0.864	1.12072e-05\\
0.866	1.09982e-05\\
0.868	1.0793e-05\\
0.87	1.05917e-05\\
0.872	1.03941e-05\\
0.874	1.02002e-05\\
0.876	1.00099e-05\\
0.878	9.8232e-06\\
0.88	9.63994e-06\\
0.882	9.46008e-06\\
0.884	9.28357e-06\\
0.886	9.11036e-06\\
0.888	8.94037e-06\\
0.89	8.77354e-06\\
0.892	8.60983e-06\\
0.894	8.44917e-06\\
0.896	8.2915e-06\\
0.898	8.13677e-06\\
0.9	7.98493e-06\\
0.902	7.83592e-06\\
0.904	7.68968e-06\\
0.906	7.54618e-06\\
0.908	7.40535e-06\\
0.91	7.26715e-06\\
0.912	7.13153e-06\\
0.914	6.99844e-06\\
0.916	6.86784e-06\\
0.918	6.73967e-06\\
0.92	6.6139e-06\\
0.922	6.49047e-06\\
0.924	6.36935e-06\\
0.926	6.25049e-06\\
0.928	6.13384e-06\\
0.93	6.01938e-06\\
0.932	5.90705e-06\\
0.934	5.79683e-06\\
0.936	5.68866e-06\\
0.938	5.58251e-06\\
0.94	5.47834e-06\\
0.942	5.37612e-06\\
0.944	5.27581e-06\\
0.946	5.17737e-06\\
0.948	5.08077e-06\\
0.95	4.98598e-06\\
0.952	4.89296e-06\\
0.954	4.80167e-06\\
0.956	4.71209e-06\\
0.958	4.62418e-06\\
0.96	4.53792e-06\\
0.962	4.45327e-06\\
0.964	4.37019e-06\\
0.966	4.28867e-06\\
0.968	4.20868e-06\\
0.97	4.13017e-06\\
0.972	4.05314e-06\\
0.974	3.97754e-06\\
0.976	3.90335e-06\\
0.978	3.83055e-06\\
0.98	3.75911e-06\\
0.982	3.689e-06\\
0.984	3.6202e-06\\
0.986	3.55269e-06\\
0.988	3.48643e-06\\
0.99	3.42142e-06\\
0.992	3.35762e-06\\
0.994	3.295e-06\\
0.996	3.23356e-06\\
0.998	3.17326e-06\\
1	3.11409e-06\\
1.002	3.05603e-06\\
1.004	2.99904e-06\\
1.006	2.94312e-06\\
1.008	2.88825e-06\\
1.01	2.83439e-06\\
1.012	2.78154e-06\\
1.014	2.72968e-06\\
1.016	2.67879e-06\\
1.018	2.62884e-06\\
1.02	2.57983e-06\\
1.022	2.53173e-06\\
1.024	2.48452e-06\\
1.026	2.4382e-06\\
1.028	2.39274e-06\\
1.03	2.34813e-06\\
1.032	2.30435e-06\\
1.034	2.26138e-06\\
1.036	2.21922e-06\\
1.038	2.17784e-06\\
1.04	2.13724e-06\\
1.042	2.09739e-06\\
1.044	2.05828e-06\\
1.046	2.0199e-06\\
1.048	1.98224e-06\\
1.05	1.94528e-06\\
1.052	1.90901e-06\\
1.054	1.87341e-06\\
1.056	1.83848e-06\\
1.058	1.80419e-06\\
1.06	1.77055e-06\\
1.062	1.73753e-06\\
1.064	1.70513e-06\\
1.066	1.67334e-06\\
1.068	1.64213e-06\\
1.07	1.61151e-06\\
1.072	1.58145e-06\\
1.074	1.55196e-06\\
1.076	1.52301e-06\\
1.078	1.49461e-06\\
1.08	1.46673e-06\\
1.082	1.43938e-06\\
1.084	1.41253e-06\\
1.086	1.38618e-06\\
1.088	1.36033e-06\\
1.09	1.33495e-06\\
1.092	1.31005e-06\\
1.094	1.28561e-06\\
1.096	1.26163e-06\\
1.098	1.23809e-06\\
1.1	1.215e-06\\
1.102	1.19233e-06\\
1.104	1.17009e-06\\
1.106	1.14826e-06\\
1.108	1.12683e-06\\
1.11	1.10581e-06\\
1.112	1.08518e-06\\
1.114	1.06493e-06\\
1.116	1.04506e-06\\
1.118	1.02556e-06\\
1.12	1.00642e-06\\
1.122	9.87643e-07\\
1.124	9.69213e-07\\
1.126	9.51127e-07\\
1.128	9.33378e-07\\
1.13	9.1596e-07\\
1.132	8.98866e-07\\
1.134	8.82091e-07\\
1.136	8.65629e-07\\
1.138	8.49474e-07\\
1.14	8.3362e-07\\
1.142	8.18061e-07\\
1.144	8.02793e-07\\
1.146	7.8781e-07\\
1.148	7.73106e-07\\
1.15	7.58676e-07\\
1.152	7.44515e-07\\
1.154	7.30618e-07\\
1.156	7.16981e-07\\
1.158	7.03598e-07\\
1.16	6.90464e-07\\
1.162	6.77576e-07\\
1.164	6.64928e-07\\
1.166	6.52515e-07\\
1.168	6.40335e-07\\
1.17	6.28382e-07\\
1.172	6.16651e-07\\
1.174	6.0514e-07\\
1.176	5.93843e-07\\
1.178	5.82758e-07\\
1.18	5.71879e-07\\
1.182	5.61203e-07\\
1.184	5.50726e-07\\
1.186	5.40445e-07\\
1.188	5.30355e-07\\
1.19	5.20454e-07\\
1.192	5.10738e-07\\
1.194	5.01203e-07\\
1.196	4.91846e-07\\
1.198	4.82664e-07\\
1.2	4.73653e-07\\
1.202	4.64811e-07\\
1.204	4.56133e-07\\
1.206	4.47617e-07\\
1.208	4.39261e-07\\
1.21	4.3106e-07\\
1.212	4.23012e-07\\
1.214	4.15115e-07\\
1.216	4.07365e-07\\
1.218	3.9976e-07\\
1.22	3.92297e-07\\
1.222	3.84973e-07\\
1.224	3.77785e-07\\
1.226	3.70732e-07\\
1.228	3.63811e-07\\
1.23	3.57019e-07\\
1.232	3.50354e-07\\
1.234	3.43813e-07\\
1.236	3.37394e-07\\
1.238	3.31095e-07\\
1.24	3.24914e-07\\
1.242	3.18848e-07\\
1.244	3.12896e-07\\
1.246	3.07054e-07\\
1.248	3.01322e-07\\
1.25	2.95696e-07\\
1.252	2.90176e-07\\
1.254	2.84759e-07\\
1.256	2.79443e-07\\
1.258	2.74226e-07\\
1.26	2.69106e-07\\
1.262	2.64083e-07\\
1.264	2.59152e-07\\
1.266	2.54314e-07\\
1.268	2.49567e-07\\
1.27	2.44908e-07\\
1.272	2.40336e-07\\
1.274	2.35849e-07\\
1.276	2.31446e-07\\
1.278	2.27125e-07\\
1.28	2.22885e-07\\
1.282	2.18724e-07\\
1.284	2.14641e-07\\
1.286	2.10634e-07\\
1.288	2.06702e-07\\
1.29	2.02843e-07\\
1.292	1.99057e-07\\
1.294	1.95341e-07\\
1.296	1.91694e-07\\
1.298	1.88116e-07\\
1.3	1.84604e-07\\
1.302	1.81158e-07\\
1.304	1.77776e-07\\
1.306	1.74457e-07\\
1.308	1.712e-07\\
1.31	1.68004e-07\\
1.312	1.64868e-07\\
1.314	1.6179e-07\\
1.316	1.5877e-07\\
1.318	1.55806e-07\\
1.32	1.52898e-07\\
1.322	1.50043e-07\\
1.324	1.47242e-07\\
1.326	1.44494e-07\\
1.328	1.41796e-07\\
1.33	1.39149e-07\\
1.332	1.36552e-07\\
1.334	1.34002e-07\\
1.336	1.31501e-07\\
1.338	1.29046e-07\\
1.34	1.26637e-07\\
1.342	1.24273e-07\\
1.344	1.21953e-07\\
1.346	1.19676e-07\\
1.348	1.17442e-07\\
1.35	1.1525e-07\\
1.352	1.13098e-07\\
1.354	1.10987e-07\\
1.356	1.08915e-07\\
1.358	1.06882e-07\\
1.36	1.04886e-07\\
1.362	1.02928e-07\\
1.364	1.01007e-07\\
1.366	9.91211e-08\\
1.368	9.72707e-08\\
1.37	9.54547e-08\\
1.372	9.36727e-08\\
1.374	9.1924e-08\\
1.376	9.02079e-08\\
1.378	8.85238e-08\\
1.38	8.68711e-08\\
1.382	8.52493e-08\\
1.384	8.36578e-08\\
1.386	8.2096e-08\\
1.388	8.05633e-08\\
1.39	7.90593e-08\\
1.392	7.75833e-08\\
1.394	7.61348e-08\\
1.396	7.47134e-08\\
1.398	7.33186e-08\\
1.4	7.19498e-08\\
1.402	7.06065e-08\\
1.404	6.92883e-08\\
1.406	6.79947e-08\\
1.408	6.67252e-08\\
1.41	6.54795e-08\\
1.412	6.4257e-08\\
1.414	6.30573e-08\\
1.416	6.188e-08\\
1.418	6.07247e-08\\
1.42	5.95909e-08\\
1.422	5.84783e-08\\
1.424	5.73865e-08\\
1.426	5.63151e-08\\
1.428	5.52637e-08\\
1.43	5.42318e-08\\
1.432	5.32193e-08\\
1.434	5.22256e-08\\
1.436	5.12505e-08\\
1.438	5.02937e-08\\
1.44	4.93546e-08\\
1.442	4.84331e-08\\
1.444	4.75288e-08\\
1.446	4.66414e-08\\
1.448	4.57705e-08\\
1.45	4.49159e-08\\
1.452	4.40773e-08\\
1.454	4.32543e-08\\
1.456	4.24467e-08\\
1.458	4.16541e-08\\
1.46	4.08764e-08\\
1.462	4.01131e-08\\
1.464	3.93641e-08\\
1.466	3.86291e-08\\
1.468	3.79079e-08\\
1.47	3.72e-08\\
1.472	3.65054e-08\\
1.474	3.58238e-08\\
1.476	3.51549e-08\\
1.478	3.44985e-08\\
1.48	3.38543e-08\\
1.482	3.32222e-08\\
1.484	3.26018e-08\\
1.486	3.19931e-08\\
1.488	3.13957e-08\\
1.49	3.08094e-08\\
1.492	3.02341e-08\\
1.494	2.96696e-08\\
1.496	2.91156e-08\\
1.498	2.85719e-08\\
1.5	2.80384e-08\\
1.502	2.75148e-08\\
1.504	2.7001e-08\\
1.506	2.64969e-08\\
1.508	2.60021e-08\\
1.51	2.55165e-08\\
1.512	2.50401e-08\\
1.514	2.45725e-08\\
1.516	2.41136e-08\\
1.518	2.36634e-08\\
1.52	2.32215e-08\\
1.522	2.27879e-08\\
1.524	2.23623e-08\\
1.526	2.19448e-08\\
1.528	2.1535e-08\\
1.53	2.11329e-08\\
1.532	2.07382e-08\\
1.534	2.0351e-08\\
1.536	1.9971e-08\\
1.538	1.9598e-08\\
1.54	1.92321e-08\\
1.542	1.88729e-08\\
1.544	1.85205e-08\\
1.546	1.81747e-08\\
1.548	1.78353e-08\\
1.55	1.75022e-08\\
1.552	1.71754e-08\\
1.554	1.68547e-08\\
1.556	1.65399e-08\\
1.558	1.62311e-08\\
1.56	1.5928e-08\\
1.562	1.56305e-08\\
1.564	1.53387e-08\\
1.566	1.50522e-08\\
1.568	1.47712e-08\\
1.57	1.44953e-08\\
1.572	1.42246e-08\\
1.574	1.3959e-08\\
1.576	1.36984e-08\\
1.578	1.34426e-08\\
1.58	1.31915e-08\\
1.582	1.29452e-08\\
1.584	1.27035e-08\\
1.586	1.24662e-08\\
1.588	1.22335e-08\\
1.59	1.2005e-08\\
1.592	1.17808e-08\\
1.594	1.15608e-08\\
1.596	1.1345e-08\\
1.598	1.11331e-08\\
1.6	1.09252e-08\\
1.602	1.07212e-08\\
1.604	1.0521e-08\\
1.606	1.03245e-08\\
1.608	1.01317e-08\\
1.61	9.94254e-09\\
1.612	9.75688e-09\\
1.614	9.57468e-09\\
1.616	9.39589e-09\\
1.618	9.22043e-09\\
1.62	9.04825e-09\\
1.622	8.87929e-09\\
1.624	8.71348e-09\\
1.626	8.55077e-09\\
1.628	8.3911e-09\\
1.63	8.23441e-09\\
1.632	8.08064e-09\\
1.634	7.92975e-09\\
1.636	7.78167e-09\\
1.638	7.63636e-09\\
1.64	7.49377e-09\\
1.642	7.35383e-09\\
1.644	7.21651e-09\\
1.646	7.08175e-09\\
1.648	6.94951e-09\\
1.65	6.81974e-09\\
1.652	6.69239e-09\\
1.654	6.56742e-09\\
1.656	6.44479e-09\\
1.658	6.32444e-09\\
1.66	6.20634e-09\\
1.662	6.09045e-09\\
1.664	5.97672e-09\\
1.666	5.86511e-09\\
1.668	5.75559e-09\\
1.67	5.64812e-09\\
1.672	5.54265e-09\\
1.674	5.43915e-09\\
1.676	5.33758e-09\\
1.678	5.23791e-09\\
1.68	5.1401e-09\\
1.682	5.04412e-09\\
1.684	4.94992e-09\\
1.686	4.85749e-09\\
1.688	4.76679e-09\\
1.69	4.67777e-09\\
1.692	4.59042e-09\\
1.694	4.5047e-09\\
1.696	4.42059e-09\\
1.698	4.33804e-09\\
1.7	4.25703e-09\\
1.702	4.17754e-09\\
1.704	4.09953e-09\\
1.706	4.02298e-09\\
1.708	3.94785e-09\\
1.71	3.87413e-09\\
1.712	3.80179e-09\\
1.714	3.7308e-09\\
1.716	3.66113e-09\\
1.718	3.59276e-09\\
1.72	3.52567e-09\\
1.722	3.45983e-09\\
1.724	3.39523e-09\\
1.726	3.33183e-09\\
1.728	3.26961e-09\\
1.73	3.20855e-09\\
1.732	3.14864e-09\\
1.734	3.08984e-09\\
1.736	3.03214e-09\\
1.738	2.97552e-09\\
1.74	2.91995e-09\\
1.742	2.86543e-09\\
1.744	2.81192e-09\\
1.746	2.75941e-09\\
1.748	2.70788e-09\\
1.75	2.65731e-09\\
1.752	2.60769e-09\\
1.754	2.55899e-09\\
1.756	2.51121e-09\\
1.758	2.46431e-09\\
1.76	2.41829e-09\\
1.762	2.37313e-09\\
1.764	2.32882e-09\\
1.766	2.28533e-09\\
1.768	2.24265e-09\\
1.77	2.20077e-09\\
1.772	2.15968e-09\\
1.774	2.11934e-09\\
1.776	2.07977e-09\\
1.778	2.04093e-09\\
1.78	2.00282e-09\\
1.782	1.96542e-09\\
1.784	1.92871e-09\\
1.786	1.89269e-09\\
1.788	1.85735e-09\\
1.79	1.82267e-09\\
1.792	1.78863e-09\\
1.794	1.75523e-09\\
1.796	1.72245e-09\\
1.798	1.69028e-09\\
1.8	1.65872e-09\\
1.802	1.62774e-09\\
1.804	1.59734e-09\\
1.806	1.56751e-09\\
1.808	1.53824e-09\\
1.81	1.50952e-09\\
1.812	1.48133e-09\\
1.814	1.45366e-09\\
1.816	1.42652e-09\\
1.818	1.39988e-09\\
1.82	1.37373e-09\\
1.822	1.34808e-09\\
1.824	1.3229e-09\\
1.826	1.2982e-09\\
1.828	1.27396e-09\\
1.83	1.25017e-09\\
1.832	1.22682e-09\\
1.834	1.20391e-09\\
1.836	1.18143e-09\\
1.838	1.15936e-09\\
1.84	1.13771e-09\\
1.842	1.11647e-09\\
1.844	1.09562e-09\\
1.846	1.07516e-09\\
1.848	1.05508e-09\\
1.85	1.03537e-09\\
1.852	1.01604e-09\\
1.854	9.97064e-10\\
1.856	9.78444e-10\\
1.858	9.60171e-10\\
1.86	9.4224e-10\\
1.862	9.24644e-10\\
1.864	9.07376e-10\\
1.866	8.90431e-10\\
1.868	8.73803e-10\\
1.87	8.57484e-10\\
1.872	8.41471e-10\\
1.874	8.25757e-10\\
1.876	8.10336e-10\\
1.878	7.95203e-10\\
1.88	7.80352e-10\\
1.882	7.65779e-10\\
1.884	7.51479e-10\\
1.886	7.37445e-10\\
1.888	7.23673e-10\\
1.89	7.10159e-10\\
1.892	6.96896e-10\\
1.894	6.83882e-10\\
1.896	6.71111e-10\\
1.898	6.58578e-10\\
1.9	6.46279e-10\\
1.902	6.3421e-10\\
1.904	6.22366e-10\\
1.906	6.10743e-10\\
1.908	5.99338e-10\\
1.91	5.88145e-10\\
1.912	5.77162e-10\\
1.914	5.66383e-10\\
1.916	5.55806e-10\\
1.918	5.45427e-10\\
1.92	5.35241e-10\\
1.922	5.25245e-10\\
1.924	5.15436e-10\\
1.926	5.05811e-10\\
1.928	4.96365e-10\\
1.93	4.87095e-10\\
1.932	4.77999e-10\\
1.934	4.69072e-10\\
1.936	4.60313e-10\\
1.938	4.51716e-10\\
1.94	4.43281e-10\\
1.942	4.35002e-10\\
1.944	4.26879e-10\\
1.946	4.18907e-10\\
1.948	4.11084e-10\\
1.95	4.03407e-10\\
1.952	3.95874e-10\\
1.954	3.88481e-10\\
1.956	3.81226e-10\\
1.958	3.74107e-10\\
1.96	3.6712e-10\\
1.962	3.60264e-10\\
1.964	3.53537e-10\\
1.966	3.46934e-10\\
1.968	3.40455e-10\\
1.97	3.34098e-10\\
1.972	3.27858e-10\\
1.974	3.21736e-10\\
1.976	3.15727e-10\\
1.978	3.09831e-10\\
1.98	3.04045e-10\\
1.982	2.98367e-10\\
1.984	2.92795e-10\\
1.986	2.87327e-10\\
1.988	2.81962e-10\\
1.99	2.76696e-10\\
1.992	2.71529e-10\\
1.994	2.66458e-10\\
1.996	2.61482e-10\\
1.998	2.56599e-10\\
2	2.51807e-10\\
};
\addplot [color=mycolor11, line width=1.0pt, forget plot]
  table[row sep=crcr]{%
-0.024	0\\
-0.022	0\\
-0.02	0\\
-0.018	0\\
-0.016	0\\
-0.014	0\\
-0.012	0\\
-0.01	0\\
-0.008	0\\
-0.006	0\\
-0.004	0\\
-0.002	0\\
0	0\\
0.002	0.00898937\\
0.004	0.0141913\\
0.006	0.0187978\\
0.008	0.0221228\\
0.01	0.0243173\\
0.012	0.0255719\\
0.014	0.0260652\\
0.016	0.0259606\\
0.018	0.0254055\\
0.02	0.025098\\
0.022	0.0247699\\
0.024	0.0246417\\
0.026	0.0253389\\
0.028	0.0267558\\
0.03	0.0280607\\
0.032	0.0292528\\
0.034	0.0303318\\
0.036	0.0312979\\
0.038	0.032152\\
0.04	0.0328954\\
0.042	0.0335297\\
0.044	0.0340572\\
0.046	0.0344802\\
0.048	0.0348013\\
0.05	0.0350237\\
0.052	0.0351504\\
0.054	0.0351849\\
0.056	0.0351309\\
0.058	0.034992\\
0.06	0.0347722\\
0.062	0.0344755\\
0.064	0.034106\\
0.066	0.0336679\\
0.068	0.0331653\\
0.07	0.0326027\\
0.072	0.0319842\\
0.074	0.0313142\\
0.076	0.0305969\\
0.078	0.0298367\\
0.08	0.0290376\\
0.082	0.028204\\
0.084	0.0273398\\
0.086	0.026449\\
0.088	0.0255357\\
0.09	0.0246035\\
0.092	0.0236562\\
0.094	0.0226973\\
0.096	0.0217303\\
0.098	0.0207585\\
0.1	0.019785\\
0.102	0.0188129\\
0.104	0.017845\\
0.106	0.0168841\\
0.108	0.0159326\\
0.11	0.0149929\\
0.112	0.0140673\\
0.114	0.0131579\\
0.116	0.0125407\\
0.118	0.012198\\
0.12	0.0118483\\
0.122	0.0114926\\
0.124	0.0111323\\
0.126	0.0107682\\
0.128	0.0104016\\
0.13	0.0100333\\
0.132	0.00966435\\
0.134	0.00934525\\
0.136	0.00923661\\
0.138	0.00918471\\
0.14	0.00912249\\
0.142	0.00905036\\
0.144	0.00896872\\
0.146	0.00887799\\
0.148	0.0087786\\
0.15	0.00867097\\
0.152	0.00855553\\
0.154	0.00843271\\
0.156	0.00830295\\
0.158	0.00816666\\
0.16	0.00802429\\
0.162	0.00787626\\
0.164	0.00772301\\
0.166	0.00756494\\
0.168	0.00740249\\
0.17	0.00723606\\
0.172	0.00706607\\
0.174	0.00689293\\
0.176	0.00671702\\
0.178	0.00653873\\
0.18	0.00635845\\
0.182	0.00617656\\
0.184	0.0059934\\
0.186	0.00585268\\
0.188	0.00575921\\
0.19	0.00566692\\
0.192	0.00557583\\
0.194	0.00548595\\
0.196	0.0053973\\
0.198	0.00530989\\
0.2	0.00522372\\
0.202	0.00513879\\
0.204	0.0050551\\
0.206	0.00497265\\
0.208	0.00489142\\
0.21	0.0048114\\
0.212	0.00473259\\
0.214	0.00465496\\
0.216	0.00457851\\
0.218	0.0045032\\
0.22	0.00442904\\
0.222	0.00435599\\
0.224	0.00428404\\
0.226	0.00421317\\
0.228	0.00414335\\
0.23	0.00407458\\
0.232	0.00400683\\
0.234	0.00394008\\
0.236	0.00387431\\
0.238	0.00380951\\
0.24	0.00374566\\
0.242	0.00368274\\
0.244	0.00362073\\
0.246	0.00355963\\
0.248	0.00349942\\
0.25	0.00344008\\
0.252	0.0033816\\
0.254	0.00332397\\
0.256	0.00326718\\
0.258	0.00321122\\
0.26	0.00315607\\
0.262	0.00310174\\
0.264	0.00304821\\
0.266	0.00299547\\
0.268	0.00294352\\
0.27	0.00289235\\
0.272	0.00284195\\
0.274	0.00279232\\
0.276	0.00274345\\
0.278	0.00269533\\
0.28	0.00264796\\
0.282	0.00260133\\
0.284	0.00255544\\
0.286	0.00251028\\
0.288	0.00246585\\
0.29	0.00242214\\
0.292	0.00237914\\
0.294	0.00233684\\
0.296	0.00229525\\
0.298	0.00225436\\
0.3	0.00221415\\
0.302	0.00217462\\
0.304	0.00213577\\
0.306	0.00209758\\
0.308	0.00206005\\
0.31	0.00202318\\
0.312	0.00198694\\
0.314	0.00195135\\
0.316	0.00191638\\
0.318	0.00188202\\
0.32	0.00184828\\
0.322	0.00181513\\
0.324	0.00178258\\
0.326	0.00175061\\
0.328	0.00171921\\
0.33	0.00168838\\
0.332	0.0016581\\
0.334	0.00162836\\
0.336	0.00159916\\
0.338	0.00157048\\
0.34	0.00154232\\
0.342	0.00151467\\
0.344	0.00148751\\
0.346	0.00146084\\
0.348	0.00143465\\
0.35	0.00140893\\
0.352	0.00138367\\
0.354	0.00135887\\
0.356	0.0013345\\
0.358	0.00131058\\
0.36	0.00128708\\
0.362	0.00126399\\
0.364	0.00124132\\
0.366	0.00121905\\
0.368	0.00119718\\
0.37	0.00117569\\
0.372	0.00115459\\
0.374	0.00113385\\
0.376	0.00111349\\
0.378	0.00109348\\
0.38	0.00107383\\
0.382	0.00105452\\
0.384	0.00103555\\
0.386	0.00101692\\
0.388	0.000998614\\
0.39	0.000980632\\
0.392	0.000962967\\
0.394	0.000945614\\
0.396	0.000928566\\
0.398	0.00091182\\
0.4	0.000895369\\
0.402	0.000879209\\
0.404	0.000863335\\
0.406	0.000847742\\
0.408	0.000832425\\
0.41	0.00081738\\
0.412	0.000802602\\
0.414	0.000788087\\
0.416	0.00077383\\
0.418	0.000759828\\
0.42	0.000746075\\
0.422	0.000732569\\
0.424	0.000719304\\
0.426	0.000706276\\
0.428	0.000693483\\
0.43	0.000680919\\
0.432	0.000668581\\
0.434	0.000656465\\
0.436	0.000644567\\
0.438	0.000632884\\
0.44	0.000621411\\
0.442	0.000610146\\
0.444	0.000599085\\
0.446	0.000588224\\
0.448	0.000577559\\
0.45	0.000567087\\
0.452	0.000556805\\
0.454	0.000546709\\
0.456	0.000536796\\
0.458	0.000527063\\
0.46	0.000517506\\
0.462	0.000508123\\
0.464	0.000498909\\
0.466	0.000489862\\
0.468	0.000480979\\
0.47	0.000472258\\
0.472	0.000463693\\
0.474	0.000455284\\
0.476	0.000447027\\
0.478	0.000438919\\
0.48	0.000430958\\
0.482	0.00042314\\
0.484	0.000415464\\
0.486	0.000407926\\
0.488	0.000400523\\
0.49	0.000393255\\
0.492	0.000386117\\
0.494	0.000379107\\
0.496	0.000372223\\
0.498	0.000365464\\
0.5	0.000358825\\
0.502	0.000352306\\
0.504	0.000345904\\
0.506	0.000339616\\
0.508	0.000333442\\
0.51	0.000327378\\
0.512	0.000321423\\
0.514	0.000315575\\
0.516	0.000309831\\
0.518	0.00030419\\
0.52	0.000298651\\
0.522	0.00029321\\
0.524	0.000287867\\
0.526	0.00028262\\
0.528	0.000277467\\
0.53	0.000272407\\
0.532	0.000267437\\
0.534	0.000262556\\
0.536	0.000257763\\
0.538	0.000253056\\
0.54	0.000248433\\
0.542	0.000243894\\
0.544	0.000239436\\
0.546	0.000235059\\
0.548	0.00023076\\
0.55	0.000226539\\
0.552	0.000222394\\
0.554	0.000218323\\
0.556	0.000214326\\
0.558	0.000210402\\
0.56	0.000206548\\
0.562	0.000202764\\
0.564	0.000199049\\
0.566	0.000195401\\
0.568	0.000191819\\
0.57	0.000188302\\
0.572	0.000184849\\
0.574	0.000181459\\
0.576	0.000178131\\
0.578	0.000174863\\
0.58	0.000171654\\
0.582	0.000168505\\
0.584	0.000165412\\
0.586	0.000162376\\
0.588	0.000159396\\
0.59	0.000156469\\
0.592	0.000153597\\
0.594	0.000150777\\
0.596	0.000148008\\
0.598	0.000145291\\
0.6	0.000142623\\
0.602	0.000140004\\
0.604	0.000137432\\
0.606	0.000134909\\
0.608	0.000132431\\
0.61	0.000129999\\
0.612	0.000127611\\
0.614	0.000125268\\
0.616	0.000122967\\
0.618	0.000120709\\
0.62	0.000118492\\
0.622	0.000116316\\
0.624	0.00011418\\
0.626	0.000112083\\
0.628	0.000110025\\
0.63	0.000108005\\
0.632	0.000106022\\
0.634	0.000104075\\
0.636	0.000102164\\
0.638	0.000100289\\
0.64	9.84475e-05\\
0.642	9.66403e-05\\
0.644	9.48663e-05\\
0.646	9.3125e-05\\
0.648	9.14157e-05\\
0.65	8.97379e-05\\
0.652	8.8091e-05\\
0.654	8.64744e-05\\
0.656	8.48876e-05\\
0.658	8.333e-05\\
0.66	8.1801e-05\\
0.662	8.03002e-05\\
0.664	7.88271e-05\\
0.666	7.7381e-05\\
0.668	7.59616e-05\\
0.67	7.45683e-05\\
0.672	7.32006e-05\\
0.674	7.18582e-05\\
0.676	7.05404e-05\\
0.678	6.92469e-05\\
0.68	6.79772e-05\\
0.682	6.67308e-05\\
0.684	6.55074e-05\\
0.686	6.43065e-05\\
0.688	6.31277e-05\\
0.69	6.19706e-05\\
0.692	6.08347e-05\\
0.694	5.97197e-05\\
0.696	5.86253e-05\\
0.698	5.75509e-05\\
0.7	5.64963e-05\\
0.702	5.54611e-05\\
0.704	5.44449e-05\\
0.706	5.34474e-05\\
0.708	5.24682e-05\\
0.71	5.15069e-05\\
0.712	5.05634e-05\\
0.714	4.96371e-05\\
0.716	4.87278e-05\\
0.718	4.78353e-05\\
0.72	4.69591e-05\\
0.722	4.60989e-05\\
0.724	4.52546e-05\\
0.726	4.44257e-05\\
0.728	4.3612e-05\\
0.73	4.28132e-05\\
0.732	4.20291e-05\\
0.734	4.12593e-05\\
0.736	4.05036e-05\\
0.738	3.97618e-05\\
0.74	3.90335e-05\\
0.742	3.83185e-05\\
0.744	3.76167e-05\\
0.746	3.69276e-05\\
0.748	3.62512e-05\\
0.75	3.55871e-05\\
0.752	3.49352e-05\\
0.754	3.42952e-05\\
0.756	3.36668e-05\\
0.758	3.305e-05\\
0.76	3.24444e-05\\
0.762	3.18499e-05\\
0.764	3.12662e-05\\
0.766	3.06932e-05\\
0.768	3.01306e-05\\
0.77	2.95783e-05\\
0.772	2.90361e-05\\
0.774	2.85038e-05\\
0.776	2.79812e-05\\
0.778	2.74681e-05\\
0.78	2.69644e-05\\
0.782	2.64699e-05\\
0.784	2.59844e-05\\
0.786	2.55078e-05\\
0.788	2.50398e-05\\
0.79	2.45804e-05\\
0.792	2.41294e-05\\
0.794	2.36866e-05\\
0.796	2.32519e-05\\
0.798	2.28251e-05\\
0.8	2.24061e-05\\
0.802	2.19947e-05\\
0.804	2.15909e-05\\
0.806	2.11944e-05\\
0.808	2.08052e-05\\
0.81	2.0423e-05\\
0.812	2.00479e-05\\
0.814	1.96796e-05\\
0.816	1.9318e-05\\
0.818	1.8963e-05\\
0.82	1.86145e-05\\
0.822	1.82724e-05\\
0.824	1.79365e-05\\
0.826	1.76068e-05\\
0.828	1.72831e-05\\
0.83	1.69653e-05\\
0.832	1.66534e-05\\
0.834	1.63471e-05\\
0.836	1.60465e-05\\
0.838	1.57513e-05\\
0.84	1.54616e-05\\
0.842	1.51772e-05\\
0.844	1.4898e-05\\
0.846	1.46238e-05\\
0.848	1.43548e-05\\
0.85	1.40906e-05\\
0.852	1.38313e-05\\
0.854	1.35768e-05\\
0.856	1.33269e-05\\
0.858	1.30816e-05\\
0.86	1.28408e-05\\
0.862	1.26044e-05\\
0.864	1.23724e-05\\
0.866	1.21446e-05\\
0.868	1.1921e-05\\
0.87	1.17015e-05\\
0.872	1.14861e-05\\
0.874	1.12746e-05\\
0.876	1.1067e-05\\
0.878	1.08632e-05\\
0.88	1.06632e-05\\
0.882	1.04668e-05\\
0.884	1.02741e-05\\
0.886	1.00849e-05\\
0.888	9.89913e-06\\
0.89	9.71683e-06\\
0.892	9.53789e-06\\
0.894	9.36223e-06\\
0.896	9.18982e-06\\
0.898	9.02057e-06\\
0.9	8.85445e-06\\
0.902	8.69138e-06\\
0.904	8.53132e-06\\
0.906	8.3742e-06\\
0.908	8.21998e-06\\
0.91	8.0686e-06\\
0.912	7.92001e-06\\
0.914	7.77416e-06\\
0.916	7.631e-06\\
0.918	7.49048e-06\\
0.92	7.35254e-06\\
0.922	7.21715e-06\\
0.924	7.08426e-06\\
0.926	6.95381e-06\\
0.928	6.82577e-06\\
0.93	6.70009e-06\\
0.932	6.57672e-06\\
0.934	6.45563e-06\\
0.936	6.33678e-06\\
0.938	6.22011e-06\\
0.94	6.1056e-06\\
0.942	5.99319e-06\\
0.944	5.88286e-06\\
0.946	5.77457e-06\\
0.948	5.66827e-06\\
0.95	5.56392e-06\\
0.952	5.46151e-06\\
0.954	5.36098e-06\\
0.956	5.2623e-06\\
0.958	5.16544e-06\\
0.96	5.07037e-06\\
0.962	4.97705e-06\\
0.964	4.88545e-06\\
0.966	4.79554e-06\\
0.968	4.70728e-06\\
0.97	4.62066e-06\\
0.972	4.53562e-06\\
0.974	4.45216e-06\\
0.976	4.37023e-06\\
0.978	4.28981e-06\\
0.98	4.21087e-06\\
0.982	4.13339e-06\\
0.984	4.05733e-06\\
0.986	3.98268e-06\\
0.988	3.9094e-06\\
0.99	3.83746e-06\\
0.992	3.76686e-06\\
0.994	3.69755e-06\\
0.996	3.62952e-06\\
0.998	3.56274e-06\\
1	3.49719e-06\\
1.002	3.43285e-06\\
1.004	3.36969e-06\\
1.006	3.30769e-06\\
1.008	3.24683e-06\\
1.01	3.1871e-06\\
1.012	3.12846e-06\\
1.014	3.0709e-06\\
1.016	3.0144e-06\\
1.018	2.95894e-06\\
1.02	2.9045e-06\\
1.022	2.85106e-06\\
1.024	2.7986e-06\\
1.026	2.74711e-06\\
1.028	2.69657e-06\\
1.03	2.64695e-06\\
1.032	2.59825e-06\\
1.034	2.55044e-06\\
1.036	2.50351e-06\\
1.038	2.45744e-06\\
1.04	2.41222e-06\\
1.042	2.36784e-06\\
1.044	2.32426e-06\\
1.046	2.28149e-06\\
1.048	2.23951e-06\\
1.05	2.19829e-06\\
1.052	2.15784e-06\\
1.054	2.11813e-06\\
1.056	2.07914e-06\\
1.058	2.04088e-06\\
1.06	2.00331e-06\\
1.062	1.96644e-06\\
1.064	1.93025e-06\\
1.066	1.89472e-06\\
1.068	1.85984e-06\\
1.07	1.82561e-06\\
1.072	1.792e-06\\
1.074	1.75901e-06\\
1.076	1.72663e-06\\
1.078	1.69485e-06\\
1.08	1.66364e-06\\
1.082	1.63301e-06\\
1.084	1.60295e-06\\
1.086	1.57344e-06\\
1.088	1.54446e-06\\
1.09	1.51603e-06\\
1.092	1.48811e-06\\
1.094	1.46071e-06\\
1.096	1.43381e-06\\
1.098	1.40741e-06\\
1.1	1.38149e-06\\
1.102	1.35605e-06\\
1.104	1.33108e-06\\
1.106	1.30656e-06\\
1.108	1.2825e-06\\
1.11	1.25888e-06\\
1.112	1.23569e-06\\
1.114	1.21293e-06\\
1.116	1.19059e-06\\
1.118	1.16866e-06\\
1.12	1.14714e-06\\
1.122	1.12601e-06\\
1.124	1.10527e-06\\
1.126	1.08491e-06\\
1.128	1.06492e-06\\
1.13	1.0453e-06\\
1.132	1.02605e-06\\
1.134	1.00714e-06\\
1.136	9.88591e-07\\
1.138	9.70378e-07\\
1.14	9.52501e-07\\
1.142	9.34953e-07\\
1.144	9.17728e-07\\
1.146	9.0082e-07\\
1.148	8.84223e-07\\
1.15	8.67932e-07\\
1.152	8.51941e-07\\
1.154	8.36244e-07\\
1.156	8.20836e-07\\
1.158	8.05713e-07\\
1.16	7.90867e-07\\
1.162	7.76295e-07\\
1.164	7.61992e-07\\
1.166	7.47952e-07\\
1.168	7.3417e-07\\
1.17	7.20642e-07\\
1.172	7.07364e-07\\
1.174	6.9433e-07\\
1.176	6.81536e-07\\
1.178	6.68978e-07\\
1.18	6.56652e-07\\
1.182	6.44552e-07\\
1.184	6.32675e-07\\
1.186	6.21017e-07\\
1.188	6.09574e-07\\
1.19	5.98342e-07\\
1.192	5.87317e-07\\
1.194	5.76495e-07\\
1.196	5.65872e-07\\
1.198	5.55445e-07\\
1.2	5.4521e-07\\
1.202	5.35164e-07\\
1.204	5.25303e-07\\
1.206	5.15624e-07\\
1.208	5.06123e-07\\
1.21	4.96797e-07\\
1.212	4.87643e-07\\
1.214	4.78657e-07\\
1.216	4.69837e-07\\
1.218	4.6118e-07\\
1.22	4.52682e-07\\
1.222	4.44341e-07\\
1.224	4.36154e-07\\
1.226	4.28117e-07\\
1.228	4.20229e-07\\
1.23	4.12486e-07\\
1.232	4.04886e-07\\
1.234	3.97425e-07\\
1.236	3.90103e-07\\
1.238	3.82915e-07\\
1.24	3.75859e-07\\
1.242	3.68934e-07\\
1.244	3.62136e-07\\
1.246	3.55464e-07\\
1.248	3.48914e-07\\
1.25	3.42485e-07\\
1.252	3.36175e-07\\
1.254	3.29981e-07\\
1.256	3.23901e-07\\
1.258	3.17933e-07\\
1.26	3.12075e-07\\
1.262	3.06325e-07\\
1.264	3.00681e-07\\
1.266	2.95141e-07\\
1.268	2.89703e-07\\
1.27	2.84365e-07\\
1.272	2.79126e-07\\
1.274	2.73983e-07\\
1.276	2.68935e-07\\
1.278	2.6398e-07\\
1.28	2.59116e-07\\
1.282	2.54342e-07\\
1.284	2.49655e-07\\
1.286	2.45055e-07\\
1.288	2.4054e-07\\
1.29	2.36108e-07\\
1.292	2.31758e-07\\
1.294	2.27488e-07\\
1.296	2.23296e-07\\
1.298	2.19182e-07\\
1.3	2.15144e-07\\
1.302	2.1118e-07\\
1.304	2.07289e-07\\
1.306	2.03469e-07\\
1.308	1.9972e-07\\
1.31	1.9604e-07\\
1.312	1.92428e-07\\
1.314	1.88883e-07\\
1.316	1.85402e-07\\
1.318	1.81986e-07\\
1.32	1.78633e-07\\
1.322	1.75342e-07\\
1.324	1.72111e-07\\
1.326	1.6894e-07\\
1.328	1.65827e-07\\
1.33	1.62771e-07\\
1.332	1.59772e-07\\
1.334	1.56828e-07\\
1.336	1.53938e-07\\
1.338	1.51102e-07\\
1.34	1.48318e-07\\
1.342	1.45585e-07\\
1.344	1.42902e-07\\
1.346	1.40269e-07\\
1.348	1.37684e-07\\
1.35	1.35147e-07\\
1.352	1.32657e-07\\
1.354	1.30212e-07\\
1.356	1.27813e-07\\
1.358	1.25458e-07\\
1.36	1.23146e-07\\
1.362	1.20877e-07\\
1.364	1.18649e-07\\
1.366	1.16463e-07\\
1.368	1.14317e-07\\
1.37	1.1221e-07\\
1.372	1.10142e-07\\
1.374	1.08113e-07\\
1.376	1.0612e-07\\
1.378	1.04165e-07\\
1.38	1.02245e-07\\
1.382	1.00361e-07\\
1.384	9.85117e-08\\
1.386	9.66963e-08\\
1.388	9.49144e-08\\
1.39	9.31653e-08\\
1.392	9.14484e-08\\
1.394	8.97632e-08\\
1.396	8.8109e-08\\
1.398	8.64852e-08\\
1.4	8.48914e-08\\
1.402	8.3327e-08\\
1.404	8.17914e-08\\
1.406	8.0284e-08\\
1.408	7.88045e-08\\
1.41	7.73522e-08\\
1.412	7.59267e-08\\
1.414	7.45274e-08\\
1.416	7.31539e-08\\
1.418	7.18057e-08\\
1.42	7.04824e-08\\
1.422	6.91835e-08\\
1.424	6.79084e-08\\
1.426	6.66569e-08\\
1.428	6.54285e-08\\
1.43	6.42226e-08\\
1.432	6.3039e-08\\
1.434	6.18772e-08\\
1.436	6.07368e-08\\
1.438	5.96175e-08\\
1.44	5.85187e-08\\
1.442	5.74402e-08\\
1.444	5.63816e-08\\
1.446	5.53425e-08\\
1.448	5.43225e-08\\
1.45	5.33213e-08\\
1.452	5.23386e-08\\
1.454	5.1374e-08\\
1.456	5.04271e-08\\
1.458	4.94977e-08\\
1.46	4.85855e-08\\
1.462	4.769e-08\\
1.464	4.68111e-08\\
1.466	4.59483e-08\\
1.468	4.51014e-08\\
1.47	4.42702e-08\\
1.472	4.34543e-08\\
1.474	4.26534e-08\\
1.476	4.18672e-08\\
1.478	4.10956e-08\\
1.48	4.03381e-08\\
1.482	3.95947e-08\\
1.484	3.88649e-08\\
1.486	3.81486e-08\\
1.488	3.74455e-08\\
1.49	3.67553e-08\\
1.492	3.60779e-08\\
1.494	3.54129e-08\\
1.496	3.47602e-08\\
1.498	3.41195e-08\\
1.5	3.34907e-08\\
1.502	3.28734e-08\\
1.504	3.22675e-08\\
1.506	3.16728e-08\\
1.508	3.1089e-08\\
1.51	3.0516e-08\\
1.512	2.99535e-08\\
1.514	2.94014e-08\\
1.516	2.88595e-08\\
1.518	2.83276e-08\\
1.52	2.78055e-08\\
1.522	2.7293e-08\\
1.524	2.67899e-08\\
1.526	2.62962e-08\\
1.528	2.58115e-08\\
1.53	2.53357e-08\\
1.532	2.48687e-08\\
1.534	2.44104e-08\\
1.536	2.39604e-08\\
1.538	2.35188e-08\\
1.54	2.30853e-08\\
1.542	2.26598e-08\\
1.544	2.22422e-08\\
1.546	2.18322e-08\\
1.548	2.14298e-08\\
1.55	2.10348e-08\\
1.552	2.06471e-08\\
1.554	2.02665e-08\\
1.556	1.9893e-08\\
1.558	1.95263e-08\\
1.56	1.91664e-08\\
1.562	1.88131e-08\\
1.564	1.84664e-08\\
1.566	1.8126e-08\\
1.568	1.77919e-08\\
1.57	1.7464e-08\\
1.572	1.71421e-08\\
1.574	1.68261e-08\\
1.576	1.6516e-08\\
1.578	1.62116e-08\\
1.58	1.59127e-08\\
1.582	1.56194e-08\\
1.584	1.53315e-08\\
1.586	1.5049e-08\\
1.588	1.47716e-08\\
1.59	1.44993e-08\\
1.592	1.4232e-08\\
1.594	1.39697e-08\\
1.596	1.37122e-08\\
1.598	1.34595e-08\\
1.6	1.32114e-08\\
1.602	1.29679e-08\\
1.604	1.27289e-08\\
1.606	1.24942e-08\\
1.608	1.2264e-08\\
1.61	1.20379e-08\\
1.612	1.1816e-08\\
1.614	1.15982e-08\\
1.616	1.13844e-08\\
1.618	1.11746e-08\\
1.62	1.09686e-08\\
1.622	1.07665e-08\\
1.624	1.0568e-08\\
1.626	1.03732e-08\\
1.628	1.0182e-08\\
1.63	9.99435e-09\\
1.632	9.81013e-09\\
1.634	9.62931e-09\\
1.636	9.45182e-09\\
1.638	9.2776e-09\\
1.64	9.1066e-09\\
1.642	8.93875e-09\\
1.644	8.77399e-09\\
1.646	8.61226e-09\\
1.648	8.45352e-09\\
1.65	8.2977e-09\\
1.652	8.14476e-09\\
1.654	7.99463e-09\\
1.656	7.84728e-09\\
1.658	7.70263e-09\\
1.66	7.56066e-09\\
1.662	7.4213e-09\\
1.664	7.28451e-09\\
1.666	7.15024e-09\\
1.668	7.01844e-09\\
1.67	6.88908e-09\\
1.672	6.7621e-09\\
1.674	6.63746e-09\\
1.676	6.51511e-09\\
1.678	6.39502e-09\\
1.68	6.27715e-09\\
1.682	6.16145e-09\\
1.684	6.04788e-09\\
1.686	5.9364e-09\\
1.688	5.82698e-09\\
1.69	5.71958e-09\\
1.692	5.61415e-09\\
1.694	5.51067e-09\\
1.696	5.40909e-09\\
1.698	5.30939e-09\\
1.7	5.21153e-09\\
1.702	5.11547e-09\\
1.704	5.02118e-09\\
1.706	4.92862e-09\\
1.708	4.83778e-09\\
1.71	4.7486e-09\\
1.712	4.66108e-09\\
1.714	4.57516e-09\\
1.716	4.49083e-09\\
1.718	4.40805e-09\\
1.72	4.3268e-09\\
1.722	4.24705e-09\\
1.724	4.16876e-09\\
1.726	4.09192e-09\\
1.728	4.0165e-09\\
1.73	3.94246e-09\\
1.732	3.86979e-09\\
1.734	3.79846e-09\\
1.736	3.72845e-09\\
1.738	3.65972e-09\\
1.74	3.59226e-09\\
1.742	3.52605e-09\\
1.744	3.46105e-09\\
1.746	3.39725e-09\\
1.748	3.33463e-09\\
1.75	3.27317e-09\\
1.752	3.21283e-09\\
1.754	3.15361e-09\\
1.756	3.09548e-09\\
1.758	3.03842e-09\\
1.76	2.98242e-09\\
1.762	2.92744e-09\\
1.764	2.87348e-09\\
1.766	2.82051e-09\\
1.768	2.76852e-09\\
1.77	2.71749e-09\\
1.772	2.6674e-09\\
1.774	2.61823e-09\\
1.776	2.56997e-09\\
1.778	2.5226e-09\\
1.78	2.4761e-09\\
1.782	2.43045e-09\\
1.784	2.38565e-09\\
1.786	2.34168e-09\\
1.788	2.29851e-09\\
1.79	2.25614e-09\\
1.792	2.21456e-09\\
1.794	2.17374e-09\\
1.796	2.13367e-09\\
1.798	2.09434e-09\\
1.8	2.05573e-09\\
1.802	2.01784e-09\\
1.804	1.98064e-09\\
1.806	1.94413e-09\\
1.808	1.90829e-09\\
1.81	1.87312e-09\\
1.812	1.83859e-09\\
1.814	1.8047e-09\\
1.816	1.77143e-09\\
1.818	1.73878e-09\\
1.82	1.70673e-09\\
1.822	1.67527e-09\\
1.824	1.64439e-09\\
1.826	1.61407e-09\\
1.828	1.58432e-09\\
1.83	1.55512e-09\\
1.832	1.52645e-09\\
1.834	1.49831e-09\\
1.836	1.47069e-09\\
1.838	1.44358e-09\\
1.84	1.41697e-09\\
1.842	1.39085e-09\\
1.844	1.36521e-09\\
1.846	1.34005e-09\\
1.848	1.31535e-09\\
1.85	1.2911e-09\\
1.852	1.2673e-09\\
1.854	1.24394e-09\\
1.856	1.22101e-09\\
1.858	1.1985e-09\\
1.86	1.17641e-09\\
1.862	1.15472e-09\\
1.864	1.13344e-09\\
1.866	1.11255e-09\\
1.868	1.09204e-09\\
1.87	1.07191e-09\\
1.872	1.05215e-09\\
1.874	1.03275e-09\\
1.876	1.01372e-09\\
1.878	9.9503e-10\\
1.88	9.76688e-10\\
1.882	9.58684e-10\\
1.884	9.41012e-10\\
1.886	9.23665e-10\\
1.888	9.06639e-10\\
1.89	8.89926e-10\\
1.892	8.73522e-10\\
1.894	8.5742e-10\\
1.896	8.41614e-10\\
1.898	8.261e-10\\
1.9	8.10872e-10\\
1.902	7.95925e-10\\
1.904	7.81253e-10\\
1.906	7.66852e-10\\
1.908	7.52716e-10\\
1.91	7.38841e-10\\
1.912	7.25222e-10\\
1.914	7.11853e-10\\
1.916	6.98731e-10\\
1.918	6.85851e-10\\
1.92	6.73209e-10\\
1.922	6.60799e-10\\
1.924	6.48618e-10\\
1.926	6.36662e-10\\
1.928	6.24926e-10\\
1.93	6.13406e-10\\
1.932	6.02099e-10\\
1.934	5.91e-10\\
1.936	5.80106e-10\\
1.938	5.69413e-10\\
1.94	5.58916e-10\\
1.942	5.48614e-10\\
1.944	5.38501e-10\\
1.946	5.28574e-10\\
1.948	5.18831e-10\\
1.95	5.09267e-10\\
1.952	4.99879e-10\\
1.954	4.90664e-10\\
1.956	4.8162e-10\\
1.958	4.72742e-10\\
1.96	4.64027e-10\\
1.962	4.55474e-10\\
1.964	4.47078e-10\\
1.966	4.38837e-10\\
1.968	4.30747e-10\\
1.97	4.22807e-10\\
1.972	4.15013e-10\\
1.974	4.07363e-10\\
1.976	3.99854e-10\\
1.978	3.92483e-10\\
1.98	3.85248e-10\\
1.982	3.78147e-10\\
1.984	3.71176e-10\\
1.986	3.64334e-10\\
1.988	3.57618e-10\\
1.99	3.51026e-10\\
1.992	3.44555e-10\\
1.994	3.38204e-10\\
1.996	3.31969e-10\\
1.998	3.2585e-10\\
2	3.19843e-10\\
};
\end{axis}

\begin{axis}[%
width={\figscale*1.528in},
height={\figscale*0.46in},
at={(\figscale*4.752in,\figscale*1.125in)},
scale only axis,
separate axis lines,
every outer x axis line/.append style={black},
every x tick label/.append style={font=\figticfont},
every x tick/.append style={black},
xmin=-0.025,
xmax=2,
xtick={0,0.5,1,1.5,2},
xticklabels={\empty},
every outer y axis line/.append style={black},
every y tick label/.append style={font=\figticfont},
every y tick/.append style={black},
ymin=0,
ymax=0.2,
ytick={0,0.2},
axis background/.style={fill=white},
xmajorgrids,
ymajorgrids,
grid style={white!55!black, opacity=0.3},
ylabel style={rotate=-90, at={(axis description cs:-0.145,0.5)}, anchor=east}
]
\addplot [color=mycolor1, line width=1.4pt, forget plot]
  table[row sep=crcr]{%
-0.024	0\\
-0.022	0\\
-0.02	0\\
-0.018	0\\
-0.016	0\\
-0.014	0\\
-0.012	0\\
-0.01	0\\
-0.008	0\\
-0.006	0\\
-0.004	0\\
-0.002	0\\
0	0\\
0.002	0.051675\\
0.004	0.0926051\\
0.006	0.123451\\
0.008	0.14509\\
0.01	0.158555\\
0.012	0.164973\\
0.014	0.165508\\
0.016	0.161316\\
0.018	0.154887\\
0.02	0.146582\\
0.022	0.136203\\
0.024	0.126018\\
0.026	0.133202\\
0.028	0.139568\\
0.03	0.145066\\
0.032	0.149659\\
0.034	0.153315\\
0.036	0.156017\\
0.038	0.157755\\
0.04	0.158531\\
0.042	0.158356\\
0.044	0.157251\\
0.046	0.155246\\
0.048	0.152378\\
0.05	0.148696\\
0.052	0.144249\\
0.054	0.139099\\
0.056	0.133308\\
0.058	0.126944\\
0.06	0.120079\\
0.062	0.112785\\
0.064	0.110638\\
0.066	0.112011\\
0.068	0.113116\\
0.07	0.113951\\
0.072	0.114517\\
0.074	0.114815\\
0.076	0.114848\\
0.078	0.114618\\
0.08	0.114128\\
0.082	0.113382\\
0.084	0.112963\\
0.086	0.112344\\
0.088	0.111496\\
0.09	0.110425\\
0.092	0.109136\\
0.094	0.107726\\
0.096	0.106114\\
0.098	0.104458\\
0.1	0.104344\\
0.102	0.104009\\
0.104	0.103454\\
0.106	0.102678\\
0.108	0.101684\\
0.11	0.100473\\
0.112	0.0990488\\
0.114	0.0974146\\
0.116	0.0955754\\
0.118	0.0935366\\
0.12	0.0913048\\
0.122	0.0888869\\
0.124	0.086291\\
0.126	0.0840568\\
0.128	0.0829837\\
0.13	0.0817803\\
0.132	0.0809944\\
0.134	0.08008\\
0.136	0.07904\\
0.138	0.0778776\\
0.14	0.0765963\\
0.142	0.0752003\\
0.144	0.0736936\\
0.146	0.0720814\\
0.148	0.0704385\\
0.15	0.0686918\\
0.152	0.0668477\\
0.154	0.0649127\\
0.156	0.0628936\\
0.158	0.0607975\\
0.16	0.0586316\\
0.162	0.0564034\\
0.164	0.0541204\\
0.166	0.0529592\\
0.168	0.0542736\\
0.17	0.0552811\\
0.172	0.055981\\
0.174	0.0563754\\
0.176	0.0564691\\
0.178	0.0562694\\
0.18	0.0557856\\
0.182	0.0550295\\
0.184	0.0540148\\
0.186	0.0527569\\
0.188	0.0512729\\
0.19	0.0495813\\
0.192	0.0477017\\
0.194	0.0456547\\
0.196	0.0434618\\
0.198	0.041145\\
0.2	0.0387264\\
0.202	0.0362285\\
0.204	0.0353356\\
0.206	0.0369511\\
0.208	0.0384302\\
0.21	0.0397635\\
0.212	0.0409428\\
0.214	0.0419613\\
0.216	0.0428132\\
0.218	0.0434938\\
0.22	0.044\\
0.222	0.0443297\\
0.224	0.0444819\\
0.226	0.0444573\\
0.228	0.0442572\\
0.23	0.0438846\\
0.232	0.0433433\\
0.234	0.0426383\\
0.236	0.0417756\\
0.238	0.0407622\\
0.24	0.039606\\
0.242	0.0383158\\
0.244	0.036901\\
0.246	0.0353719\\
0.248	0.0337393\\
0.25	0.0320144\\
0.252	0.0302091\\
0.254	0.0298543\\
0.256	0.0317445\\
0.258	0.0335406\\
0.26	0.0352368\\
0.262	0.0368279\\
0.264	0.038309\\
0.266	0.0396758\\
0.268	0.0409244\\
0.27	0.0420515\\
0.272	0.0430543\\
0.274	0.0439303\\
0.276	0.0446778\\
0.278	0.0452954\\
0.28	0.0457823\\
0.282	0.0461381\\
0.284	0.046363\\
0.286	0.0464575\\
0.288	0.0464228\\
0.29	0.0462604\\
0.292	0.0459722\\
0.294	0.0459179\\
0.296	0.0463325\\
0.298	0.0466248\\
0.3	0.0467951\\
0.302	0.0471151\\
0.304	0.047465\\
0.306	0.0476758\\
0.308	0.047748\\
0.31	0.047683\\
0.312	0.0474823\\
0.314	0.0471482\\
0.316	0.0466834\\
0.318	0.046091\\
0.32	0.0453745\\
0.322	0.044538\\
0.324	0.0435858\\
0.326	0.0425226\\
0.328	0.0413535\\
0.33	0.040084\\
0.332	0.0387198\\
0.334	0.0372669\\
0.336	0.0357315\\
0.338	0.0343708\\
0.34	0.0340745\\
0.342	0.0337067\\
0.344	0.0332691\\
0.346	0.0327635\\
0.348	0.032192\\
0.35	0.0315567\\
0.352	0.03086\\
0.354	0.0301043\\
0.356	0.0296727\\
0.358	0.0295708\\
0.36	0.029405\\
0.362	0.0291766\\
0.364	0.0288869\\
0.366	0.0285374\\
0.368	0.0281298\\
0.37	0.027666\\
0.372	0.027148\\
0.374	0.0265778\\
0.376	0.0259578\\
0.378	0.0252903\\
0.38	0.0245778\\
0.382	0.0238229\\
0.384	0.0230284\\
0.386	0.022197\\
0.388	0.0213315\\
0.39	0.0206456\\
0.392	0.0205986\\
0.394	0.0205144\\
0.396	0.0203937\\
0.398	0.020237\\
0.4	0.0200449\\
0.402	0.0198183\\
0.404	0.0195578\\
0.406	0.0197581\\
0.408	0.0201296\\
0.41	0.0204127\\
0.412	0.0206077\\
0.414	0.0207153\\
0.416	0.0207363\\
0.418	0.020672\\
0.42	0.0205412\\
0.422	0.0204189\\
0.424	0.0202484\\
0.426	0.0199986\\
0.428	0.0196713\\
0.43	0.0192688\\
0.432	0.0187935\\
0.434	0.0187272\\
0.436	0.0187737\\
0.438	0.0187883\\
0.44	0.0187713\\
0.442	0.0187231\\
0.444	0.0186443\\
0.446	0.0185354\\
0.448	0.018397\\
0.45	0.0182298\\
0.452	0.0180344\\
0.454	0.0178118\\
0.456	0.0175628\\
0.458	0.0172882\\
0.46	0.016989\\
0.462	0.0166662\\
0.464	0.0163209\\
0.466	0.0159541\\
0.468	0.0155669\\
0.47	0.0151605\\
0.472	0.0154078\\
0.474	0.0161352\\
0.476	0.0168076\\
0.478	0.0174229\\
0.48	0.0179797\\
0.482	0.0184766\\
0.484	0.0189125\\
0.486	0.0192863\\
0.488	0.0195976\\
0.49	0.0198459\\
0.492	0.0200309\\
0.494	0.0201529\\
0.496	0.020212\\
0.498	0.0202089\\
0.5	0.0201442\\
0.502	0.0200189\\
0.504	0.0198342\\
0.506	0.0195915\\
0.508	0.0195523\\
0.51	0.0199834\\
0.512	0.0203537\\
0.514	0.0206627\\
0.516	0.0209099\\
0.518	0.0210954\\
0.52	0.0212191\\
0.522	0.0212813\\
0.524	0.0212826\\
0.526	0.0212236\\
0.528	0.0211051\\
0.53	0.0209448\\
0.532	0.020798\\
0.534	0.0206331\\
0.536	0.0204127\\
0.538	0.0201382\\
0.54	0.019811\\
0.542	0.0194328\\
0.544	0.0190054\\
0.546	0.0185307\\
0.548	0.0180108\\
0.55	0.0174479\\
0.552	0.0168442\\
0.554	0.0162022\\
0.556	0.0155243\\
0.558	0.0149878\\
0.56	0.01508\\
0.562	0.0152115\\
0.564	0.0153052\\
0.566	0.0153611\\
0.568	0.0153797\\
0.57	0.0153611\\
0.572	0.0153059\\
0.574	0.0152145\\
0.576	0.0150877\\
0.578	0.0149261\\
0.58	0.0147306\\
0.582	0.0145021\\
0.584	0.0142415\\
0.586	0.01395\\
0.588	0.0136287\\
0.59	0.0132788\\
0.592	0.0129016\\
0.594	0.0124984\\
0.596	0.0120707\\
0.598	0.0116199\\
0.6	0.0111474\\
0.602	0.0107933\\
0.604	0.0105885\\
0.606	0.0103618\\
0.608	0.0101144\\
0.61	0.00984711\\
0.612	0.00956095\\
0.614	0.00925696\\
0.616	0.00917058\\
0.618	0.00955294\\
0.62	0.00989281\\
0.622	0.0101897\\
0.624	0.0104432\\
0.626	0.0106531\\
0.628	0.0108194\\
0.63	0.0109421\\
0.632	0.0110214\\
0.634	0.0110577\\
0.636	0.0110514\\
0.638	0.0110032\\
0.64	0.010927\\
0.642	0.0108652\\
0.644	0.0108001\\
0.646	0.0106981\\
0.648	0.0105574\\
0.65	0.010379\\
0.652	0.0101641\\
0.654	0.00991387\\
0.656	0.00962952\\
0.658	0.00931249\\
0.66	0.00896425\\
0.662	0.008722\\
0.664	0.00859292\\
0.666	0.00844717\\
0.668	0.00828528\\
0.67	0.00810777\\
0.672	0.00791522\\
0.674	0.00770823\\
0.676	0.00748745\\
0.678	0.00727503\\
0.68	0.00739774\\
0.682	0.00749355\\
0.684	0.00763613\\
0.686	0.00794248\\
0.688	0.00822088\\
0.69	0.00847066\\
0.692	0.00869128\\
0.694	0.00888231\\
0.696	0.00904342\\
0.698	0.00917441\\
0.7	0.00927519\\
0.702	0.00934578\\
0.704	0.00938631\\
0.706	0.009397\\
0.708	0.00937822\\
0.71	0.00933038\\
0.712	0.00925405\\
0.714	0.00914985\\
0.716	0.00901853\\
0.718	0.00886089\\
0.72	0.00867784\\
0.722	0.00852597\\
0.724	0.00881072\\
0.726	0.00909128\\
0.728	0.00934268\\
0.73	0.0095644\\
0.732	0.00975607\\
0.734	0.00991739\\
0.736	0.0100482\\
0.738	0.0101483\\
0.74	0.0102179\\
0.742	0.0102569\\
0.744	0.0102657\\
0.746	0.0102445\\
0.748	0.0101938\\
0.75	0.0101223\\
0.752	0.0100704\\
0.754	0.0100121\\
0.756	0.00993629\\
0.758	0.00983289\\
0.76	0.00970252\\
0.762	0.00954589\\
0.764	0.00936381\\
0.766	0.00915712\\
0.768	0.00892677\\
0.77	0.00867376\\
0.772	0.00839914\\
0.774	0.00810402\\
0.776	0.00778959\\
0.778	0.00745706\\
0.78	0.0072178\\
0.782	0.00720004\\
0.784	0.00716367\\
0.786	0.00710899\\
0.788	0.00703638\\
0.79	0.00694625\\
0.792	0.00704736\\
0.794	0.00715374\\
0.796	0.0072411\\
0.798	0.00730941\\
0.8	0.00735865\\
0.802	0.00738889\\
0.804	0.00740024\\
0.806	0.00739287\\
0.808	0.00736698\\
0.81	0.00732286\\
0.812	0.00726081\\
0.814	0.00718121\\
0.816	0.00708446\\
0.818	0.00697103\\
0.82	0.00684141\\
0.822	0.00669614\\
0.824	0.00653582\\
0.826	0.00636104\\
0.828	0.00617248\\
0.83	0.0059708\\
0.832	0.00575673\\
0.834	0.005531\\
0.836	0.00529439\\
0.838	0.00504768\\
0.84	0.00518491\\
0.842	0.00533603\\
0.844	0.00546528\\
0.846	0.00557252\\
0.848	0.00565766\\
0.85	0.00572073\\
0.852	0.00576179\\
0.854	0.00578098\\
0.856	0.00577854\\
0.858	0.00575474\\
0.86	0.00571338\\
0.862	0.00569402\\
0.864	0.00566429\\
0.866	0.00563048\\
0.868	0.00557658\\
0.87	0.00550302\\
0.872	0.00541028\\
0.874	0.00529892\\
0.876	0.00516953\\
0.878	0.00502277\\
0.88	0.00485934\\
0.882	0.00467998\\
0.884	0.00448549\\
0.886	0.00427668\\
0.888	0.00405443\\
0.89	0.0039955\\
0.892	0.00405986\\
0.894	0.00411464\\
0.896	0.00415982\\
0.898	0.00419538\\
0.9	0.00422134\\
0.902	0.00423774\\
0.904	0.00430708\\
0.906	0.0043824\\
0.908	0.00444301\\
0.91	0.00448887\\
0.912	0.00451999\\
0.914	0.00453644\\
0.916	0.00453833\\
0.918	0.00452583\\
0.92	0.00449915\\
0.922	0.00445856\\
0.924	0.00440437\\
0.926	0.00433695\\
0.928	0.00425668\\
0.93	0.00416402\\
0.932	0.00411361\\
0.934	0.00424953\\
0.936	0.00437084\\
0.938	0.00447729\\
0.94	0.00456867\\
0.942	0.00464484\\
0.944	0.0047057\\
0.946	0.00475122\\
0.948	0.00478142\\
0.95	0.00479636\\
0.952	0.00489349\\
0.954	0.00497592\\
0.956	0.00504254\\
0.958	0.00509331\\
0.96	0.00512823\\
0.962	0.00514736\\
0.964	0.00515081\\
0.966	0.00513875\\
0.968	0.00511138\\
0.97	0.00506923\\
0.972	0.00504908\\
0.974	0.00501637\\
0.976	0.00498813\\
0.978	0.00494541\\
0.98	0.00488853\\
0.982	0.00481782\\
0.984	0.00473366\\
0.986	0.00463647\\
0.988	0.00452672\\
0.99	0.00440489\\
0.992	0.00427151\\
0.994	0.00412714\\
0.996	0.00397236\\
0.998	0.00380779\\
1	0.00363408\\
1.002	0.00345187\\
1.004	0.00338843\\
1.006	0.00339664\\
1.008	0.00339588\\
1.01	0.00338597\\
1.012	0.00336697\\
1.014	0.00333905\\
1.016	0.0033024\\
1.018	0.00325724\\
1.02	0.0032038\\
1.022	0.00319579\\
1.024	0.00327255\\
1.026	0.00334002\\
1.028	0.00339811\\
1.03	0.00344674\\
1.032	0.00348588\\
1.034	0.00351553\\
1.036	0.00353569\\
1.038	0.00354643\\
1.04	0.0035478\\
1.042	0.00353991\\
1.044	0.00352287\\
1.046	0.00349686\\
1.048	0.00346202\\
1.05	0.00341856\\
1.052	0.00336671\\
1.054	0.0033067\\
1.056	0.00323879\\
1.058	0.00316327\\
1.06	0.00308044\\
1.062	0.00299061\\
1.064	0.0029655\\
1.066	0.00301779\\
1.068	0.00305857\\
1.07	0.00308787\\
1.072	0.00310571\\
1.074	0.00311219\\
1.076	0.00310742\\
1.078	0.00309157\\
1.08	0.00306481\\
1.082	0.00305517\\
1.084	0.0030364\\
1.086	0.00302114\\
1.088	0.00299819\\
1.09	0.00296492\\
1.092	0.00292155\\
1.094	0.00286836\\
1.096	0.00280563\\
1.098	0.0027337\\
1.1	0.0026529\\
1.102	0.00256361\\
1.104	0.00246622\\
1.106	0.00236115\\
1.108	0.00224884\\
1.11	0.00212974\\
1.112	0.00205158\\
1.114	0.00209715\\
1.116	0.00213546\\
1.118	0.00216647\\
1.12	0.00219015\\
1.122	0.00220652\\
1.124	0.00221559\\
1.126	0.00221744\\
1.128	0.00221215\\
1.13	0.00219982\\
1.132	0.00218059\\
1.134	0.0021546\\
1.136	0.00212205\\
1.138	0.00208312\\
1.14	0.00203805\\
1.142	0.00203816\\
1.144	0.00204183\\
1.146	0.00204159\\
1.148	0.0021155\\
1.15	0.00218206\\
1.152	0.00224114\\
1.154	0.00229262\\
1.156	0.00233641\\
1.158	0.00237245\\
1.16	0.00240071\\
1.162	0.00242118\\
1.164	0.00243389\\
1.166	0.00243888\\
1.168	0.00243622\\
1.17	0.00247282\\
1.172	0.00252349\\
1.174	0.00256594\\
1.176	0.00260011\\
1.178	0.00262599\\
1.18	0.00264357\\
1.182	0.00265289\\
1.184	0.00265403\\
1.186	0.00264705\\
1.188	0.00263208\\
1.19	0.00260926\\
1.192	0.00259935\\
1.194	0.00258425\\
1.196	0.00256905\\
1.198	0.00255203\\
1.2	0.00252752\\
1.202	0.00249569\\
1.204	0.00245672\\
1.206	0.00241082\\
1.208	0.00235821\\
1.21	0.00229913\\
1.212	0.00223386\\
1.214	0.00216266\\
1.216	0.00208584\\
1.218	0.00200371\\
1.22	0.00191659\\
1.222	0.00182482\\
1.224	0.00172875\\
1.226	0.00162872\\
1.228	0.00161785\\
1.23	0.00162909\\
1.232	0.0016357\\
1.234	0.00163787\\
1.236	0.00163556\\
1.238	0.00162876\\
1.24	0.00161753\\
1.242	0.00160196\\
1.244	0.00158216\\
1.246	0.00155822\\
1.248	0.00153028\\
1.25	0.00149846\\
1.252	0.00146474\\
1.254	0.00151324\\
1.256	0.00155724\\
1.258	0.00159667\\
1.26	0.00163145\\
1.262	0.00166154\\
1.264	0.0016869\\
1.266	0.00170751\\
1.268	0.00172337\\
1.27	0.00173449\\
1.272	0.00174087\\
1.274	0.00174256\\
1.276	0.00173961\\
1.278	0.00173207\\
1.28	0.00172001\\
1.282	0.00170354\\
1.284	0.00168273\\
1.286	0.0016577\\
1.288	0.00162856\\
1.29	0.00164182\\
1.292	0.00164999\\
1.294	0.00165221\\
1.296	0.00164854\\
1.298	0.00163907\\
1.3	0.00162389\\
1.302	0.00161826\\
1.304	0.00161\\
1.306	0.00159934\\
1.308	0.00159041\\
1.31	0.00157602\\
1.312	0.00155629\\
1.314	0.00153134\\
1.316	0.00150131\\
1.318	0.00146638\\
1.32	0.0014267\\
1.322	0.00138246\\
1.324	0.00133388\\
1.326	0.00128115\\
1.328	0.0012245\\
1.33	0.00116415\\
1.332	0.00110036\\
1.334	0.00108803\\
1.336	0.00108967\\
1.338	0.00108779\\
1.34	0.00108244\\
1.342	0.00107367\\
1.344	0.00106157\\
1.346	0.00104623\\
1.348	0.00104683\\
1.35	0.00105774\\
1.352	0.00106508\\
1.354	0.00106885\\
1.356	0.00106907\\
1.358	0.0010658\\
1.36	0.00105909\\
1.362	0.00105682\\
1.364	0.0010959\\
1.366	0.00113116\\
1.368	0.0011625\\
1.37	0.00118988\\
1.372	0.00121323\\
1.374	0.00123253\\
1.376	0.00124776\\
1.378	0.00125891\\
1.38	0.00126598\\
1.382	0.001269\\
1.384	0.00126801\\
1.386	0.00126305\\
1.388	0.00126345\\
1.39	0.00129373\\
1.392	0.00131972\\
1.394	0.00134137\\
1.396	0.00135867\\
1.398	0.00137159\\
1.4	0.00138016\\
1.402	0.00138439\\
1.404	0.00138431\\
1.406	0.00137998\\
1.408	0.00137145\\
1.41	0.0013588\\
1.412	0.00135317\\
1.414	0.0013465\\
1.416	0.0013362\\
1.418	0.00132961\\
1.42	0.00131907\\
1.422	0.00130464\\
1.424	0.00128642\\
1.426	0.00126451\\
1.428	0.00123902\\
1.43	0.00121008\\
1.432	0.00117781\\
1.434	0.00114236\\
1.436	0.00110388\\
1.438	0.00106252\\
1.44	0.00101847\\
1.442	0.000971872\\
1.444	0.000922924\\
1.446	0.000871806\\
1.448	0.000818708\\
1.45	0.000773423\\
1.452	0.00078432\\
1.454	0.000792912\\
1.456	0.000799198\\
1.458	0.000803184\\
1.46	0.000804968\\
1.462	0.000804548\\
1.464	0.000801889\\
1.466	0.000797026\\
1.468	0.000789998\\
1.47	0.000780852\\
1.472	0.000769641\\
1.474	0.000756422\\
1.476	0.000741259\\
1.478	0.000724222\\
1.48	0.00073596\\
1.482	0.000750932\\
1.484	0.000762939\\
1.486	0.000771976\\
1.488	0.00077805\\
1.49	0.000781177\\
1.492	0.000792168\\
1.494	0.000809466\\
1.496	0.000824407\\
1.498	0.000836974\\
1.5	0.000847155\\
1.502	0.000854949\\
1.504	0.000860358\\
1.506	0.000863393\\
1.508	0.000864768\\
1.51	0.000871564\\
1.512	0.000875212\\
1.514	0.000875739\\
1.516	0.000873182\\
1.518	0.000867586\\
1.52	0.000859007\\
1.522	0.000855529\\
1.524	0.000851982\\
1.526	0.000845511\\
1.528	0.000841132\\
1.53	0.000835161\\
1.532	0.00082637\\
1.534	0.000814819\\
1.536	0.000800582\\
1.538	0.000783737\\
1.54	0.000764372\\
1.542	0.000742581\\
1.544	0.000718464\\
1.546	0.00069213\\
1.548	0.000663691\\
1.55	0.000633266\\
1.552	0.000600979\\
1.554	0.000566959\\
1.556	0.000534948\\
1.558	0.000543381\\
1.56	0.000549918\\
1.562	0.000554557\\
1.564	0.000557301\\
1.566	0.000558164\\
1.568	0.000557162\\
1.57	0.00055432\\
1.572	0.000549669\\
1.574	0.000543245\\
1.576	0.000535091\\
1.578	0.000546452\\
1.58	0.000566797\\
1.582	0.000585121\\
1.584	0.000601384\\
1.586	0.000615554\\
1.588	0.000627607\\
1.59	0.000637525\\
1.592	0.000645298\\
1.594	0.000650924\\
1.596	0.000654409\\
1.598	0.000655764\\
1.6	0.000655009\\
1.602	0.00065217\\
1.604	0.00064728\\
1.606	0.00064038\\
1.608	0.00065375\\
1.61	0.000668947\\
1.612	0.000681884\\
1.614	0.000692542\\
1.616	0.000700913\\
1.618	0.000706995\\
1.62	0.000710792\\
1.622	0.00071232\\
1.624	0.000711597\\
1.626	0.000708652\\
1.628	0.00070352\\
1.63	0.000696242\\
1.632	0.000692701\\
1.634	0.000689595\\
1.636	0.000684382\\
1.638	0.000679748\\
1.64	0.000675208\\
1.642	0.000668633\\
1.644	0.000660067\\
1.646	0.000649561\\
1.648	0.000637171\\
1.65	0.000622958\\
1.652	0.000606991\\
1.654	0.00058934\\
1.656	0.000570083\\
1.658	0.0005493\\
1.66	0.000527077\\
1.662	0.000503502\\
1.664	0.000478667\\
1.666	0.000452667\\
1.668	0.000425601\\
1.67	0.000397567\\
1.672	0.000390942\\
1.674	0.000390649\\
1.676	0.000389457\\
1.678	0.000387379\\
1.68	0.000384431\\
1.682	0.000384921\\
1.684	0.000387572\\
1.686	0.000389136\\
1.688	0.00038962\\
1.69	0.000388995\\
1.692	0.000387276\\
1.694	0.00038825\\
1.696	0.00039742\\
1.698	0.000405036\\
1.7	0.000411093\\
1.702	0.000415588\\
1.704	0.000418525\\
1.706	0.000419913\\
1.708	0.000419766\\
1.71	0.000418104\\
1.712	0.000414951\\
1.714	0.000410335\\
1.716	0.000409514\\
1.718	0.000422178\\
1.72	0.000433214\\
1.722	0.000442604\\
1.724	0.000450336\\
1.726	0.000456407\\
1.728	0.000460814\\
1.73	0.000463564\\
1.732	0.000464668\\
1.734	0.000464141\\
1.736	0.000462005\\
1.738	0.000458287\\
1.74	0.000453017\\
1.742	0.000450492\\
1.744	0.000448607\\
1.746	0.000445197\\
1.748	0.000441449\\
1.75	0.000438736\\
1.752	0.000434548\\
1.754	0.000428917\\
1.756	0.000421879\\
1.758	0.000413474\\
1.76	0.000403745\\
1.762	0.000392742\\
1.764	0.000380515\\
1.766	0.000367119\\
1.768	0.000352612\\
1.77	0.000337055\\
1.772	0.000320511\\
1.774	0.000303048\\
1.776	0.000297321\\
1.778	0.000297439\\
1.78	0.000296563\\
1.782	0.000294707\\
1.784	0.000291888\\
1.786	0.000288127\\
1.788	0.000283448\\
1.79	0.000277879\\
1.792	0.000274224\\
1.794	0.000286073\\
1.796	0.000296881\\
1.798	0.000306625\\
1.8	0.000315282\\
1.802	0.000322837\\
1.804	0.000329275\\
1.806	0.000334588\\
1.808	0.00033877\\
1.81	0.00034182\\
1.812	0.00034374\\
1.814	0.000344538\\
1.816	0.000344222\\
1.818	0.000342806\\
1.82	0.000340309\\
1.822	0.00033675\\
1.824	0.000332153\\
1.826	0.000336432\\
1.828	0.000345082\\
1.83	0.000352549\\
1.832	0.000358822\\
1.834	0.000363893\\
1.836	0.00036776\\
1.838	0.000370423\\
1.84	0.000371888\\
1.842	0.000372162\\
1.844	0.000371259\\
1.846	0.000369195\\
1.848	0.00036599\\
1.85	0.000361667\\
1.852	0.000359392\\
1.854	0.00035779\\
1.856	0.000355086\\
1.858	0.000351553\\
1.86	0.000349512\\
1.862	0.000346407\\
1.864	0.000342259\\
1.866	0.000337096\\
1.868	0.000330944\\
1.87	0.000323836\\
1.872	0.000315806\\
1.874	0.00030689\\
1.876	0.000297127\\
1.878	0.00028656\\
1.88	0.000275232\\
1.882	0.000263188\\
1.884	0.000250476\\
1.886	0.000237145\\
1.888	0.000223245\\
1.89	0.000208828\\
1.892	0.000193948\\
1.894	0.000182832\\
1.896	0.000178842\\
1.898	0.000174225\\
1.9	0.000173877\\
1.902	0.000178489\\
1.904	0.000186603\\
1.906	0.000193943\\
1.908	0.000200493\\
1.91	0.000206241\\
1.912	0.000211177\\
1.914	0.000215294\\
1.916	0.000218588\\
1.918	0.000221057\\
1.92	0.000222704\\
1.922	0.000223532\\
1.924	0.00022355\\
1.926	0.000222765\\
1.928	0.000221192\\
1.93	0.000218844\\
1.932	0.000215739\\
1.934	0.000212975\\
1.936	0.000220061\\
1.938	0.000226296\\
1.94	0.000231669\\
1.942	0.000236173\\
1.944	0.000239803\\
1.946	0.000242558\\
1.948	0.000244439\\
1.95	0.00024545\\
1.952	0.000245598\\
1.954	0.000244894\\
1.956	0.00024335\\
1.958	0.00024098\\
1.96	0.000237803\\
1.962	0.000236165\\
1.964	0.000235146\\
1.966	0.00023333\\
1.968	0.000230733\\
1.97	0.000229228\\
1.972	0.00022723\\
1.974	0.000224478\\
1.976	0.00022099\\
1.978	0.000217604\\
1.98	0.000216709\\
1.982	0.00021523\\
1.984	0.000213177\\
1.986	0.000210563\\
1.988	0.000207401\\
1.99	0.000203707\\
1.992	0.000199498\\
1.994	0.000194791\\
1.996	0.000189605\\
1.998	0.000183961\\
2	0.00017788\\
};
\addplot [color=mycolor2, line width=1.0pt, forget plot]
  table[row sep=crcr]{%
-0.024	0\\
-0.022	0\\
-0.02	0\\
-0.018	0\\
-0.016	0\\
-0.014	0\\
-0.012	0\\
-0.01	0\\
-0.008	0\\
-0.006	0\\
-0.004	0\\
-0.002	0\\
0	0\\
0.002	0.051675\\
0.004	0.0926051\\
0.006	0.123451\\
0.008	0.14509\\
0.01	0.158555\\
0.012	0.164971\\
0.014	0.165505\\
0.016	0.161311\\
0.018	0.154882\\
0.02	0.146575\\
0.022	0.136194\\
0.024	0.126013\\
0.026	0.133196\\
0.028	0.13956\\
0.03	0.145056\\
0.032	0.149645\\
0.034	0.153299\\
0.036	0.155997\\
0.038	0.157731\\
0.04	0.158501\\
0.042	0.158321\\
0.044	0.157209\\
0.046	0.155197\\
0.048	0.152322\\
0.05	0.14863\\
0.052	0.144175\\
0.054	0.139014\\
0.056	0.133212\\
0.058	0.126837\\
0.06	0.119959\\
0.062	0.112652\\
0.064	0.110559\\
0.066	0.111923\\
0.068	0.113018\\
0.07	0.113842\\
0.072	0.114397\\
0.074	0.114684\\
0.076	0.114704\\
0.078	0.114461\\
0.08	0.113958\\
0.082	0.113198\\
0.084	0.112398\\
0.086	0.111854\\
0.088	0.111085\\
0.09	0.110095\\
0.092	0.108887\\
0.094	0.107465\\
0.096	0.105837\\
0.098	0.104006\\
0.1	0.10198\\
0.102	0.0997664\\
0.104	0.0973718\\
0.106	0.0948045\\
0.108	0.0920732\\
0.11	0.0891868\\
0.112	0.0867723\\
0.114	0.0868395\\
0.116	0.0867371\\
0.118	0.0864655\\
0.12	0.0860255\\
0.122	0.085419\\
0.124	0.084648\\
0.126	0.0837155\\
0.128	0.0826249\\
0.13	0.0814917\\
0.132	0.0806918\\
0.134	0.0797631\\
0.136	0.0787087\\
0.138	0.0775317\\
0.14	0.0762359\\
0.142	0.0748251\\
0.144	0.0733038\\
0.146	0.0716764\\
0.148	0.0700113\\
0.15	0.0682491\\
0.152	0.0663894\\
0.154	0.0644391\\
0.156	0.0624048\\
0.158	0.0602938\\
0.16	0.0581133\\
0.162	0.0558707\\
0.164	0.0535738\\
0.166	0.052418\\
0.168	0.0537534\\
0.17	0.0547829\\
0.172	0.0555058\\
0.174	0.055924\\
0.176	0.0560423\\
0.178	0.0558678\\
0.18	0.0554098\\
0.182	0.0546799\\
0.184	0.0536919\\
0.186	0.0524609\\
0.188	0.0510041\\
0.19	0.0493397\\
0.192	0.0474874\\
0.194	0.0454677\\
0.196	0.0433018\\
0.198	0.0410117\\
0.2	0.0386196\\
0.202	0.0361478\\
0.204	0.0354175\\
0.206	0.0370783\\
0.208	0.0386017\\
0.21	0.0399781\\
0.212	0.0411991\\
0.214	0.0422578\\
0.216	0.0431482\\
0.218	0.0438655\\
0.22	0.0444065\\
0.222	0.0447687\\
0.224	0.0449514\\
0.226	0.0449549\\
0.228	0.0447805\\
0.23	0.0444311\\
0.232	0.0439104\\
0.234	0.0432233\\
0.236	0.0423758\\
0.238	0.0413748\\
0.24	0.0402283\\
0.242	0.0389449\\
0.244	0.0375342\\
0.246	0.0360063\\
0.248	0.0343721\\
0.25	0.032643\\
0.252	0.0308306\\
0.254	0.0301992\\
0.256	0.0321089\\
0.258	0.0339232\\
0.26	0.0356362\\
0.262	0.0372426\\
0.264	0.0387376\\
0.266	0.0401167\\
0.268	0.041376\\
0.27	0.0425123\\
0.272	0.0435227\\
0.274	0.0444046\\
0.276	0.0451565\\
0.278	0.0457767\\
0.28	0.0462646\\
0.282	0.0466198\\
0.284	0.0468424\\
0.286	0.046933\\
0.288	0.0468929\\
0.29	0.0467234\\
0.292	0.0464267\\
0.294	0.0460051\\
0.296	0.0458071\\
0.298	0.0462301\\
0.3	0.0465326\\
0.302	0.0467146\\
0.304	0.0467767\\
0.306	0.04672\\
0.308	0.0465456\\
0.31	0.0462553\\
0.312	0.0458512\\
0.314	0.0453354\\
0.316	0.0447108\\
0.318	0.0439804\\
0.32	0.0431475\\
0.322	0.0422156\\
0.324	0.0411887\\
0.326	0.0400709\\
0.328	0.0388665\\
0.33	0.0375802\\
0.332	0.0362167\\
0.334	0.0347811\\
0.336	0.0333116\\
0.338	0.0330931\\
0.34	0.0327956\\
0.342	0.0324211\\
0.344	0.0319716\\
0.346	0.0314491\\
0.348	0.0308562\\
0.35	0.0301956\\
0.352	0.0300729\\
0.354	0.0300882\\
0.356	0.0300369\\
0.358	0.0299198\\
0.36	0.029738\\
0.362	0.0294927\\
0.364	0.0291854\\
0.366	0.0288175\\
0.368	0.0283909\\
0.37	0.0279075\\
0.372	0.0273693\\
0.374	0.0267785\\
0.376	0.0261375\\
0.378	0.0254487\\
0.38	0.0247146\\
0.382	0.023938\\
0.384	0.0231216\\
0.386	0.0222683\\
0.388	0.0213808\\
0.39	0.0207327\\
0.392	0.0206726\\
0.394	0.0205755\\
0.396	0.0204417\\
0.398	0.0202721\\
0.4	0.0200672\\
0.402	0.0198278\\
0.404	0.0195548\\
0.406	0.019249\\
0.408	0.0196784\\
0.41	0.0200878\\
0.412	0.0204096\\
0.414	0.020644\\
0.416	0.0207916\\
0.418	0.020853\\
0.42	0.0208294\\
0.422	0.020722\\
0.424	0.0205323\\
0.426	0.0202621\\
0.428	0.0199135\\
0.43	0.0194887\\
0.432	0.0189902\\
0.434	0.0187986\\
0.436	0.0188311\\
0.438	0.0188317\\
0.44	0.0188009\\
0.442	0.018739\\
0.444	0.0186466\\
0.446	0.0185243\\
0.448	0.0183727\\
0.45	0.0181925\\
0.452	0.0179845\\
0.454	0.0177496\\
0.456	0.0174885\\
0.458	0.0172023\\
0.46	0.0168919\\
0.462	0.0165583\\
0.464	0.0162026\\
0.466	0.0158258\\
0.468	0.0154292\\
0.47	0.0150262\\
0.472	0.0158226\\
0.474	0.0165648\\
0.476	0.0172506\\
0.478	0.0178782\\
0.48	0.0184459\\
0.482	0.0189524\\
0.484	0.0193964\\
0.486	0.019777\\
0.488	0.0200937\\
0.49	0.0203459\\
0.492	0.0205335\\
0.494	0.0206566\\
0.496	0.0207155\\
0.498	0.0207106\\
0.5	0.0206428\\
0.502	0.020513\\
0.504	0.0203224\\
0.506	0.0200725\\
0.508	0.0197649\\
0.51	0.0194013\\
0.512	0.0197616\\
0.514	0.0201824\\
0.516	0.0205428\\
0.518	0.0208424\\
0.52	0.0210808\\
0.522	0.021258\\
0.524	0.021374\\
0.526	0.0214292\\
0.528	0.021424\\
0.53	0.0213591\\
0.532	0.0212355\\
0.534	0.0210541\\
0.536	0.0208163\\
0.538	0.0205233\\
0.54	0.0201768\\
0.542	0.0197786\\
0.544	0.0193303\\
0.546	0.0188342\\
0.548	0.0182922\\
0.55	0.0177067\\
0.552	0.01708\\
0.554	0.0164146\\
0.556	0.0157131\\
0.558	0.0153577\\
0.56	0.0151751\\
0.562	0.0149563\\
0.564	0.0147025\\
0.566	0.0144148\\
0.568	0.0140944\\
0.57	0.0137428\\
0.572	0.0133612\\
0.574	0.0129511\\
0.576	0.0125142\\
0.578	0.012052\\
0.58	0.0117589\\
0.582	0.0118246\\
0.584	0.011862\\
0.586	0.0118713\\
0.588	0.0118529\\
0.59	0.011807\\
0.592	0.0117341\\
0.594	0.0116348\\
0.596	0.0115097\\
0.598	0.0113594\\
0.6	0.0111847\\
0.602	0.0109863\\
0.604	0.0107651\\
0.606	0.0105221\\
0.608	0.0102581\\
0.61	0.00997423\\
0.612	0.00967147\\
0.614	0.00935094\\
0.616	0.00901376\\
0.618	0.00868465\\
0.62	0.00911837\\
0.622	0.00951078\\
0.624	0.00986121\\
0.626	0.0101691\\
0.628	0.0104341\\
0.63	0.0106558\\
0.632	0.0108343\\
0.634	0.0109694\\
0.636	0.0110615\\
0.638	0.0111106\\
0.64	0.0111174\\
0.642	0.0110824\\
0.644	0.0110062\\
0.646	0.0108896\\
0.648	0.0107336\\
0.65	0.0105393\\
0.652	0.0103077\\
0.654	0.0100401\\
0.656	0.00973796\\
0.658	0.00940262\\
0.66	0.00903565\\
0.662	0.00870202\\
0.664	0.00856692\\
0.666	0.00841544\\
0.668	0.0082481\\
0.67	0.00806547\\
0.672	0.00786811\\
0.674	0.00765664\\
0.676	0.00743171\\
0.678	0.00719399\\
0.68	0.00720373\\
0.682	0.00757623\\
0.684	0.00792163\\
0.686	0.00823903\\
0.688	0.00852761\\
0.69	0.00878668\\
0.692	0.00901568\\
0.694	0.00921413\\
0.696	0.00938172\\
0.698	0.00951821\\
0.7	0.0096235\\
0.702	0.0096976\\
0.704	0.00974064\\
0.706	0.00975284\\
0.708	0.00973455\\
0.71	0.00968623\\
0.712	0.00960841\\
0.714	0.00950175\\
0.716	0.009367\\
0.718	0.00920499\\
0.72	0.00901665\\
0.722	0.00880298\\
0.724	0.00856507\\
0.726	0.00841233\\
0.728	0.00873997\\
0.73	0.00903958\\
0.732	0.00931051\\
0.734	0.00955219\\
0.736	0.00976417\\
0.738	0.00994609\\
0.74	0.0100977\\
0.742	0.0102188\\
0.744	0.0103094\\
0.746	0.0103695\\
0.748	0.0103993\\
0.75	0.010399\\
0.752	0.0103689\\
0.754	0.0103094\\
0.756	0.0102211\\
0.758	0.0101045\\
0.76	0.00996043\\
0.762	0.00978953\\
0.764	0.00959267\\
0.766	0.00937078\\
0.768	0.00912484\\
0.77	0.0088559\\
0.772	0.00856507\\
0.774	0.00825352\\
0.776	0.00792248\\
0.778	0.00757322\\
0.78	0.00749012\\
0.782	0.0074615\\
0.784	0.0074137\\
0.786	0.0073471\\
0.788	0.00726208\\
0.79	0.00715912\\
0.792	0.00703872\\
0.794	0.00690143\\
0.796	0.00674786\\
0.798	0.00657865\\
0.8	0.00639449\\
0.802	0.00619609\\
0.804	0.00598422\\
0.806	0.00575966\\
0.808	0.00552324\\
0.81	0.0052758\\
0.812	0.00520828\\
0.814	0.00523782\\
0.816	0.00525424\\
0.818	0.00525764\\
0.82	0.00524816\\
0.822	0.00522596\\
0.824	0.00519127\\
0.826	0.00514433\\
0.828	0.00508542\\
0.83	0.00501484\\
0.832	0.00493295\\
0.834	0.00484012\\
0.836	0.00473674\\
0.838	0.00462325\\
0.84	0.00463964\\
0.842	0.00485592\\
0.844	0.00505122\\
0.846	0.00522518\\
0.848	0.00537752\\
0.85	0.005508\\
0.852	0.00561649\\
0.854	0.00570292\\
0.856	0.00576728\\
0.858	0.00580964\\
0.86	0.00583016\\
0.862	0.00582904\\
0.864	0.00580655\\
0.866	0.00576305\\
0.868	0.00569894\\
0.87	0.0056147\\
0.872	0.00551085\\
0.874	0.00538797\\
0.876	0.00524673\\
0.878	0.0050878\\
0.88	0.00491193\\
0.882	0.00471991\\
0.884	0.00451258\\
0.886	0.0042908\\
0.888	0.00405549\\
0.89	0.0040258\\
0.892	0.00408666\\
0.894	0.00413802\\
0.896	0.00417984\\
0.898	0.00421213\\
0.9	0.00429941\\
0.902	0.00441021\\
0.904	0.00450597\\
0.906	0.00458654\\
0.908	0.0046518\\
0.91	0.0047017\\
0.912	0.00473623\\
0.914	0.00475546\\
0.916	0.00475949\\
0.918	0.0047485\\
0.92	0.00472269\\
0.922	0.00468234\\
0.924	0.00462776\\
0.926	0.00455932\\
0.928	0.00447742\\
0.93	0.00438252\\
0.932	0.0042751\\
0.934	0.00415571\\
0.936	0.00402491\\
0.938	0.00388331\\
0.94	0.00373152\\
0.942	0.00368478\\
0.944	0.00389098\\
0.946	0.00408382\\
0.948	0.00426284\\
0.95	0.00442758\\
0.952	0.00457767\\
0.954	0.00471279\\
0.956	0.00483264\\
0.958	0.00493702\\
0.96	0.00502574\\
0.962	0.0050987\\
0.964	0.00515581\\
0.966	0.00519708\\
0.968	0.00522254\\
0.97	0.00523228\\
0.972	0.00522645\\
0.974	0.00520522\\
0.976	0.00516883\\
0.978	0.00511758\\
0.98	0.00505179\\
0.982	0.00497182\\
0.984	0.0048781\\
0.986	0.00477108\\
0.988	0.00465125\\
0.99	0.00451914\\
0.992	0.00437531\\
0.994	0.00422035\\
0.996	0.00405488\\
0.998	0.00387956\\
1	0.00369506\\
1.002	0.00355044\\
1.004	0.00356177\\
1.006	0.00356337\\
1.008	0.00355565\\
1.01	0.00353839\\
1.012	0.00351175\\
1.014	0.00347592\\
1.016	0.00343112\\
1.018	0.00337757\\
1.02	0.00331556\\
1.022	0.00324536\\
1.024	0.00316728\\
1.026	0.00308166\\
1.028	0.00298884\\
1.03	0.00288919\\
1.032	0.0027831\\
1.034	0.00267096\\
1.036	0.00255318\\
1.038	0.00243019\\
1.04	0.00232546\\
1.042	0.00231591\\
1.044	0.00233022\\
1.046	0.00233844\\
1.048	0.00234063\\
1.05	0.00233685\\
1.052	0.00232717\\
1.054	0.0023117\\
1.056	0.00229056\\
1.058	0.00233809\\
1.06	0.00246531\\
1.062	0.00258195\\
1.064	0.00268778\\
1.066	0.00278259\\
1.068	0.00286619\\
1.07	0.00293847\\
1.072	0.00299932\\
1.074	0.00304868\\
1.076	0.00308655\\
1.078	0.00311294\\
1.08	0.0031279\\
1.082	0.00313152\\
1.084	0.00312394\\
1.086	0.00310532\\
1.088	0.00307584\\
1.09	0.00303575\\
1.092	0.0029853\\
1.094	0.00292479\\
1.096	0.00285453\\
1.098	0.00277487\\
1.1	0.00268619\\
1.102	0.00258889\\
1.104	0.00248339\\
1.106	0.00237013\\
1.108	0.00224959\\
1.11	0.00212224\\
1.112	0.0021608\\
1.114	0.00221057\\
1.116	0.00225274\\
1.118	0.00228725\\
1.12	0.00231408\\
1.122	0.00233322\\
1.124	0.00234469\\
1.126	0.00234856\\
1.128	0.0023449\\
1.13	0.00233382\\
1.132	0.00231544\\
1.134	0.00228993\\
1.136	0.00225746\\
1.138	0.00221824\\
1.14	0.00217248\\
1.142	0.00212044\\
1.144	0.00206237\\
1.146	0.00199856\\
1.148	0.00192931\\
1.15	0.00188329\\
1.152	0.00190326\\
1.154	0.00191874\\
1.156	0.00192976\\
1.158	0.00193634\\
1.16	0.00193851\\
1.162	0.00193632\\
1.164	0.00193642\\
1.166	0.00204146\\
1.168	0.00213938\\
1.17	0.00222994\\
1.172	0.00231294\\
1.174	0.00238817\\
1.176	0.0024555\\
1.178	0.00251477\\
1.18	0.00256589\\
1.182	0.00260878\\
1.184	0.0026434\\
1.186	0.00266972\\
1.188	0.00268774\\
1.19	0.00269751\\
1.192	0.00269908\\
1.194	0.00269253\\
1.196	0.00267798\\
1.198	0.00265556\\
1.2	0.00262543\\
1.202	0.00258778\\
1.204	0.00254281\\
1.206	0.00249074\\
1.208	0.00243182\\
1.21	0.00236633\\
1.212	0.00229454\\
1.214	0.00221675\\
1.216	0.00213329\\
1.218	0.00204449\\
1.22	0.00195069\\
1.222	0.00185226\\
1.224	0.00174955\\
1.226	0.00170824\\
1.228	0.00172052\\
1.23	0.00172793\\
1.232	0.0017305\\
1.234	0.00172852\\
1.236	0.00172181\\
1.238	0.00171044\\
1.24	0.00169449\\
1.242	0.00167408\\
1.244	0.00164932\\
1.246	0.00162032\\
1.248	0.00158724\\
1.25	0.00155022\\
1.252	0.00150942\\
1.254	0.00146501\\
1.256	0.00141716\\
1.258	0.00136607\\
1.26	0.00131193\\
1.262	0.00125494\\
1.264	0.00119531\\
1.266	0.00113325\\
1.268	0.00106899\\
1.27	0.00106105\\
1.272	0.00106909\\
1.274	0.00107427\\
1.276	0.0011382\\
1.278	0.00121254\\
1.28	0.00128157\\
1.282	0.00134513\\
1.284	0.00140308\\
1.286	0.00145531\\
1.288	0.00150171\\
1.29	0.00154219\\
1.292	0.00157672\\
1.294	0.00160523\\
1.296	0.00162772\\
1.298	0.00164417\\
1.3	0.00165462\\
1.302	0.0016591\\
1.304	0.00165767\\
1.306	0.00165041\\
1.308	0.0016374\\
1.31	0.00161876\\
1.312	0.00159463\\
1.314	0.00156513\\
1.316	0.00153045\\
1.318	0.00149075\\
1.32	0.00144621\\
1.322	0.00139706\\
1.324	0.00134349\\
1.326	0.00128575\\
1.328	0.00122407\\
1.33	0.00115869\\
1.332	0.00115407\\
1.334	0.00116111\\
1.336	0.00116436\\
1.338	0.00116386\\
1.34	0.00115967\\
1.342	0.00115184\\
1.344	0.00114044\\
1.346	0.00112556\\
1.348	0.00110731\\
1.35	0.00108577\\
1.352	0.00106108\\
1.354	0.00103336\\
1.356	0.00100274\\
1.358	0.00096937\\
1.36	0.000933405\\
1.362	0.000895004\\
1.364	0.000854333\\
1.366	0.000811565\\
1.368	0.000766877\\
1.37	0.000730041\\
1.372	0.000694834\\
1.374	0.000657525\\
1.376	0.000718254\\
1.378	0.000788224\\
1.38	0.000855179\\
1.382	0.000918933\\
1.384	0.000979313\\
1.386	0.00103616\\
1.388	0.00108933\\
1.39	0.00113868\\
1.392	0.00118411\\
1.394	0.0012255\\
1.396	0.00126276\\
1.398	0.00129582\\
1.4	0.00132463\\
1.402	0.00134912\\
1.404	0.00136927\\
1.406	0.00138506\\
1.408	0.0013965\\
1.41	0.00140358\\
1.412	0.00140634\\
1.414	0.00140481\\
1.416	0.00139905\\
1.418	0.00138913\\
1.42	0.00137512\\
1.422	0.0013571\\
1.424	0.0013352\\
1.426	0.00130951\\
1.428	0.00128017\\
1.43	0.00124732\\
1.432	0.00121108\\
1.434	0.00117164\\
1.436	0.00112913\\
1.438	0.00108375\\
1.44	0.00103566\\
1.442	0.000985055\\
1.444	0.000932125\\
1.446	0.000877068\\
1.448	0.000821024\\
1.45	0.00083239\\
1.452	0.000841327\\
1.454	0.000847835\\
1.456	0.000851922\\
1.458	0.000853605\\
1.46	0.000853051\\
1.462	0.00085014\\
1.464	0.00084491\\
1.466	0.000837402\\
1.468	0.000827666\\
1.47	0.000815757\\
1.472	0.000801737\\
1.474	0.000785673\\
1.476	0.000767638\\
1.478	0.000747709\\
1.48	0.000725971\\
1.482	0.000702511\\
1.484	0.000677419\\
1.486	0.000650793\\
1.488	0.000622731\\
1.49	0.000593336\\
1.492	0.000562714\\
1.494	0.00054479\\
1.496	0.00058751\\
1.498	0.000627579\\
1.5	0.0006649\\
1.502	0.000699384\\
1.504	0.000730955\\
1.506	0.000759546\\
1.508	0.000785099\\
1.51	0.00080757\\
1.512	0.000826922\\
1.514	0.000843132\\
1.516	0.000856185\\
1.518	0.000866076\\
1.52	0.000872814\\
1.522	0.000876414\\
1.524	0.000876903\\
1.526	0.000874318\\
1.528	0.000868705\\
1.53	0.000860119\\
1.532	0.000848626\\
1.534	0.000834298\\
1.536	0.000817218\\
1.538	0.000797474\\
1.54	0.000775164\\
1.542	0.000750393\\
1.544	0.000723272\\
1.546	0.000693918\\
1.548	0.000662456\\
1.55	0.000629013\\
1.552	0.000593725\\
1.554	0.00056383\\
1.556	0.000557036\\
1.558	0.000548543\\
1.56	0.000538403\\
1.562	0.000526671\\
1.564	0.000513409\\
1.566	0.000498681\\
1.568	0.000482559\\
1.57	0.000477788\\
1.572	0.000479457\\
1.574	0.000479506\\
1.576	0.0004779\\
1.578	0.000474667\\
1.58	0.000469838\\
1.582	0.000463449\\
1.584	0.000455541\\
1.586	0.000446158\\
1.588	0.000435348\\
1.59	0.000423165\\
1.592	0.000409664\\
1.594	0.000394905\\
1.596	0.00037895\\
1.598	0.000386528\\
1.6	0.000421986\\
1.602	0.0004558\\
1.604	0.000487878\\
1.606	0.000518136\\
1.608	0.000546494\\
1.61	0.000572882\\
1.612	0.000597234\\
1.614	0.000619494\\
1.616	0.000639614\\
1.618	0.00065755\\
1.62	0.00067327\\
1.622	0.000686747\\
1.624	0.000697962\\
1.626	0.000706904\\
1.628	0.000713569\\
1.63	0.000717963\\
1.632	0.000720095\\
1.634	0.000719984\\
1.636	0.000717657\\
1.638	0.000713146\\
1.64	0.00070649\\
1.642	0.000697736\\
1.644	0.000686935\\
1.646	0.000674147\\
1.648	0.000659435\\
1.65	0.00064287\\
1.652	0.000624525\\
1.654	0.000604482\\
1.656	0.000582825\\
1.658	0.000559642\\
1.66	0.000535027\\
1.662	0.000509076\\
1.664	0.000481889\\
1.666	0.000453568\\
1.668	0.000424218\\
1.67	0.00039479\\
1.672	0.000395074\\
1.674	0.00039445\\
1.676	0.000400407\\
1.678	0.000405383\\
1.68	0.000409164\\
1.682	0.000411753\\
1.684	0.000413157\\
1.686	0.000413469\\
1.688	0.000412618\\
1.69	0.000410619\\
1.692	0.000407492\\
1.694	0.000403261\\
1.696	0.000397952\\
1.698	0.000391592\\
1.7	0.000384216\\
1.702	0.000375857\\
1.704	0.000366553\\
1.706	0.000356344\\
1.708	0.000345271\\
1.71	0.00033338\\
1.712	0.000320716\\
1.714	0.000307328\\
1.716	0.000307766\\
1.718	0.000328789\\
1.72	0.000348387\\
1.722	0.000366515\\
1.724	0.00038313\\
1.726	0.000398198\\
1.728	0.000411688\\
1.73	0.000423574\\
1.732	0.000433838\\
1.734	0.000442465\\
1.736	0.000449448\\
1.738	0.000454783\\
1.74	0.000458473\\
1.742	0.000460526\\
1.744	0.000460955\\
1.746	0.000459779\\
1.748	0.00045702\\
1.75	0.000452707\\
1.752	0.000446874\\
1.754	0.000439556\\
1.756	0.000430798\\
1.758	0.000420644\\
1.76	0.000409145\\
1.762	0.000396354\\
1.764	0.000382331\\
1.766	0.000367135\\
1.768	0.000350831\\
1.77	0.000333485\\
1.772	0.000315168\\
1.774	0.00029595\\
1.776	0.000275907\\
1.778	0.000255114\\
1.78	0.000237919\\
1.782	0.000229049\\
1.784	0.000219573\\
1.786	0.000216737\\
1.788	0.000222021\\
1.79	0.000226472\\
1.792	0.000230086\\
1.794	0.000232865\\
1.796	0.00023481\\
1.798	0.000235958\\
1.8	0.000236318\\
1.802	0.000235865\\
1.804	0.000234613\\
1.806	0.000232575\\
1.808	0.000229769\\
1.81	0.000226214\\
1.812	0.000222039\\
1.814	0.000218364\\
1.816	0.000214167\\
1.818	0.000209464\\
1.82	0.000215288\\
1.822	0.00023305\\
1.824	0.000249911\\
1.826	0.000265827\\
1.828	0.000280756\\
1.83	0.00029466\\
1.832	0.000307506\\
1.834	0.000319263\\
1.836	0.000329906\\
1.838	0.000339411\\
1.84	0.000347761\\
1.842	0.000354942\\
1.844	0.000360944\\
1.846	0.00036576\\
1.848	0.000369388\\
1.85	0.000371829\\
1.852	0.00037309\\
1.854	0.00037318\\
1.856	0.000372111\\
1.858	0.0003699\\
1.86	0.000366568\\
1.862	0.000362139\\
1.864	0.000356638\\
1.866	0.000350097\\
1.868	0.000342549\\
1.87	0.00033403\\
1.872	0.000324578\\
1.874	0.000314236\\
1.876	0.000303047\\
1.878	0.000291057\\
1.88	0.000278316\\
1.882	0.000264872\\
1.884	0.000250777\\
1.886	0.000236086\\
1.888	0.000220853\\
1.89	0.000205134\\
1.892	0.000188985\\
1.894	0.000175728\\
1.896	0.00018072\\
1.898	0.000185148\\
1.9	0.000189006\\
1.902	0.000192287\\
1.904	0.00019499\\
1.906	0.000197112\\
1.908	0.000198655\\
1.91	0.00019962\\
1.912	0.00020006\\
1.914	0.000199932\\
1.916	0.000199243\\
1.918	0.000198001\\
1.92	0.000196218\\
1.922	0.000193906\\
1.924	0.000191077\\
1.926	0.000187748\\
1.928	0.000183934\\
1.93	0.000179653\\
1.932	0.000175865\\
1.934	0.000174793\\
1.936	0.000173336\\
1.938	0.000171652\\
1.94	0.000181912\\
1.942	0.000191407\\
1.944	0.000200115\\
1.946	0.000208017\\
1.948	0.000215097\\
1.95	0.000221342\\
1.952	0.000226741\\
1.954	0.000231288\\
1.956	0.000234976\\
1.958	0.000237806\\
1.96	0.000239777\\
1.962	0.000240894\\
1.964	0.000241164\\
1.966	0.000240596\\
1.968	0.000239201\\
1.97	0.000236995\\
1.972	0.000233994\\
1.974	0.000230219\\
1.976	0.00022569\\
1.978	0.000220432\\
1.98	0.00021447\\
1.982	0.000207834\\
1.984	0.000200552\\
1.986	0.000192657\\
1.988	0.000184181\\
1.99	0.000175161\\
1.992	0.000165631\\
1.994	0.00015563\\
1.996	0.000145195\\
1.998	0.000134368\\
2	0.000123187\\
};
\addplot [color=mycolor3, line width=1.0pt, forget plot]
  table[row sep=crcr]{%
-0.024	0\\
-0.022	0\\
-0.02	0\\
-0.018	0\\
-0.016	0\\
-0.014	0\\
-0.012	0\\
-0.01	0\\
-0.008	0\\
-0.006	0\\
-0.004	0\\
-0.002	0\\
0	0\\
0.002	0.051675\\
0.004	0.0926051\\
0.006	0.12345\\
0.008	0.145089\\
0.01	0.158554\\
0.012	0.16497\\
0.014	0.165503\\
0.016	0.161308\\
0.018	0.154879\\
0.02	0.14657\\
0.022	0.136188\\
0.024	0.126011\\
0.026	0.133193\\
0.028	0.139556\\
0.03	0.145051\\
0.032	0.149639\\
0.034	0.153291\\
0.036	0.155987\\
0.038	0.157719\\
0.04	0.158487\\
0.042	0.158303\\
0.044	0.157188\\
0.046	0.155171\\
0.048	0.152292\\
0.05	0.148595\\
0.052	0.144134\\
0.054	0.138968\\
0.056	0.133159\\
0.058	0.126776\\
0.06	0.11989\\
0.062	0.112575\\
0.064	0.110494\\
0.066	0.11185\\
0.068	0.112936\\
0.07	0.113752\\
0.072	0.114297\\
0.074	0.114573\\
0.076	0.114582\\
0.078	0.114327\\
0.08	0.113812\\
0.082	0.113039\\
0.084	0.112015\\
0.086	0.110744\\
0.088	0.109234\\
0.09	0.10749\\
0.092	0.105521\\
0.094	0.103334\\
0.096	0.100939\\
0.098	0.098344\\
0.1	0.0955599\\
0.102	0.0925969\\
0.104	0.0894662\\
0.106	0.0861791\\
0.108	0.0859655\\
0.11	0.0863596\\
0.112	0.0865847\\
0.114	0.08664\\
0.116	0.0865254\\
0.118	0.0862411\\
0.12	0.0857883\\
0.122	0.0851686\\
0.124	0.0843843\\
0.126	0.0834382\\
0.128	0.0823338\\
0.13	0.081215\\
0.132	0.0804019\\
0.134	0.07946\\
0.136	0.0783922\\
0.138	0.0772018\\
0.14	0.0758925\\
0.142	0.0744683\\
0.144	0.0729335\\
0.146	0.0713151\\
0.148	0.0696435\\
0.15	0.0678683\\
0.152	0.0659959\\
0.154	0.0640329\\
0.156	0.0619862\\
0.158	0.059863\\
0.16	0.0576706\\
0.162	0.0554165\\
0.164	0.0531083\\
0.166	0.0519524\\
0.168	0.0533027\\
0.17	0.0543481\\
0.172	0.0550877\\
0.174	0.0555235\\
0.176	0.05566\\
0.178	0.0555045\\
0.18	0.0550661\\
0.182	0.0543564\\
0.184	0.053389\\
0.186	0.0521791\\
0.188	0.0507437\\
0.19	0.049101\\
0.192	0.0472705\\
0.194	0.0452728\\
0.196	0.043129\\
0.198	0.0408609\\
0.2	0.0384907\\
0.202	0.0360407\\
0.204	0.0352494\\
0.206	0.0369359\\
0.208	0.0384849\\
0.21	0.039887\\
0.212	0.0411336\\
0.214	0.0422175\\
0.216	0.0431329\\
0.218	0.0438749\\
0.22	0.04444\\
0.222	0.0448259\\
0.224	0.0450316\\
0.226	0.0450574\\
0.228	0.0449046\\
0.23	0.0445759\\
0.232	0.044075\\
0.234	0.0434068\\
0.236	0.0425771\\
0.238	0.041593\\
0.24	0.0404622\\
0.242	0.0391934\\
0.244	0.037796\\
0.246	0.0362802\\
0.248	0.0346569\\
0.25	0.0329374\\
0.252	0.0311334\\
0.254	0.0306973\\
0.256	0.0326277\\
0.258	0.0344611\\
0.26	0.0361915\\
0.262	0.0378135\\
0.264	0.0393223\\
0.266	0.0407134\\
0.268	0.0419829\\
0.27	0.0431275\\
0.272	0.0441441\\
0.274	0.0450304\\
0.276	0.0457845\\
0.278	0.0464052\\
0.28	0.0468915\\
0.282	0.0472431\\
0.284	0.0474602\\
0.286	0.0475435\\
0.288	0.0474939\\
0.29	0.0473133\\
0.292	0.0470035\\
0.294	0.0465671\\
0.296	0.0460071\\
0.298	0.0453267\\
0.3	0.0445297\\
0.302	0.0436201\\
0.304	0.0426024\\
0.306	0.0414814\\
0.308	0.0402621\\
0.31	0.0389499\\
0.312	0.0375505\\
0.314	0.0360696\\
0.316	0.035266\\
0.318	0.0353921\\
0.32	0.0354347\\
0.322	0.0353947\\
0.324	0.0352728\\
0.326	0.0350702\\
0.328	0.0347881\\
0.33	0.0344282\\
0.332	0.0339922\\
0.334	0.0338979\\
0.336	0.0337455\\
0.338	0.0335116\\
0.34	0.0331977\\
0.342	0.0328057\\
0.344	0.0323377\\
0.346	0.0317958\\
0.348	0.0311826\\
0.35	0.0305009\\
0.352	0.0304909\\
0.354	0.0304916\\
0.356	0.0304247\\
0.358	0.0302913\\
0.36	0.0300923\\
0.362	0.0298291\\
0.364	0.0295032\\
0.366	0.0291162\\
0.368	0.0286699\\
0.37	0.0281662\\
0.372	0.0276073\\
0.374	0.0269955\\
0.376	0.0263332\\
0.378	0.0256228\\
0.38	0.0248671\\
0.382	0.0240686\\
0.384	0.0232304\\
0.386	0.0223552\\
0.388	0.0214461\\
0.39	0.0205061\\
0.392	0.0195383\\
0.394	0.0192715\\
0.396	0.0192545\\
0.398	0.0191465\\
0.4	0.0189493\\
0.402	0.018665\\
0.404	0.018296\\
0.406	0.0178451\\
0.408	0.0173151\\
0.41	0.0173211\\
0.412	0.0176261\\
0.414	0.0178986\\
0.416	0.0181384\\
0.418	0.0183451\\
0.42	0.0185186\\
0.422	0.0186587\\
0.424	0.0187654\\
0.426	0.0188388\\
0.428	0.018879\\
0.43	0.0188862\\
0.432	0.0188608\\
0.434	0.0188031\\
0.436	0.0187137\\
0.438	0.018593\\
0.44	0.0184416\\
0.442	0.0182603\\
0.444	0.0180499\\
0.446	0.017811\\
0.448	0.0175447\\
0.45	0.0172518\\
0.452	0.0169334\\
0.454	0.0165905\\
0.456	0.0162241\\
0.458	0.0158355\\
0.46	0.0154258\\
0.462	0.0149961\\
0.464	0.0145479\\
0.466	0.0140823\\
0.468	0.0146919\\
0.47	0.0155614\\
0.472	0.0163776\\
0.474	0.0171379\\
0.476	0.0178403\\
0.478	0.0184828\\
0.48	0.0190638\\
0.482	0.0195817\\
0.484	0.0200355\\
0.486	0.0204241\\
0.488	0.0207469\\
0.49	0.0210034\\
0.492	0.0211936\\
0.494	0.0213173\\
0.496	0.021375\\
0.498	0.0213672\\
0.5	0.0212946\\
0.502	0.0211583\\
0.504	0.0209595\\
0.506	0.0206996\\
0.508	0.0203803\\
0.51	0.0200034\\
0.512	0.019571\\
0.514	0.0190852\\
0.516	0.0185484\\
0.518	0.0179632\\
0.52	0.0173322\\
0.522	0.0166583\\
0.524	0.0159443\\
0.526	0.0151933\\
0.528	0.0144085\\
0.53	0.0144481\\
0.532	0.0147995\\
0.534	0.0151103\\
0.536	0.0153798\\
0.538	0.0156078\\
0.54	0.0157938\\
0.542	0.0159379\\
0.544	0.01604\\
0.546	0.0161002\\
0.548	0.0161188\\
0.55	0.0160962\\
0.552	0.016033\\
0.554	0.0159296\\
0.556	0.015787\\
0.558	0.0156059\\
0.56	0.0153873\\
0.562	0.0151323\\
0.564	0.014842\\
0.566	0.0145178\\
0.568	0.014161\\
0.57	0.0137729\\
0.572	0.0133553\\
0.574	0.0129095\\
0.576	0.0124374\\
0.578	0.0120069\\
0.58	0.0120892\\
0.582	0.0121423\\
0.584	0.0121664\\
0.586	0.0121618\\
0.588	0.0121288\\
0.59	0.0120678\\
0.592	0.0119794\\
0.594	0.0118641\\
0.596	0.0117227\\
0.598	0.0115557\\
0.6	0.0113639\\
0.602	0.0111483\\
0.604	0.0109098\\
0.606	0.0106492\\
0.608	0.0103676\\
0.61	0.0100661\\
0.612	0.00977633\\
0.614	0.00956063\\
0.616	0.00930546\\
0.618	0.00901234\\
0.62	0.00868295\\
0.622	0.00831905\\
0.624	0.00792249\\
0.626	0.00749524\\
0.628	0.00703934\\
0.63	0.00661041\\
0.632	0.00652967\\
0.634	0.0064492\\
0.636	0.0063691\\
0.638	0.00628949\\
0.64	0.00621049\\
0.642	0.00613221\\
0.644	0.00613739\\
0.646	0.00641283\\
0.648	0.00665945\\
0.65	0.00687688\\
0.652	0.00706486\\
0.654	0.00722323\\
0.656	0.0073519\\
0.658	0.00745089\\
0.66	0.0075203\\
0.662	0.00756031\\
0.664	0.00757122\\
0.666	0.00755338\\
0.668	0.00750723\\
0.67	0.00743331\\
0.672	0.00733222\\
0.674	0.00720464\\
0.676	0.00705133\\
0.678	0.00702827\\
0.68	0.00744527\\
0.682	0.00783536\\
0.684	0.00819747\\
0.686	0.00853062\\
0.688	0.00883395\\
0.69	0.00910673\\
0.692	0.00934836\\
0.694	0.00955833\\
0.696	0.00973628\\
0.698	0.00988196\\
0.7	0.00999524\\
0.702	0.0100761\\
0.704	0.0101247\\
0.706	0.0101412\\
0.708	0.010126\\
0.71	0.0100795\\
0.712	0.0100023\\
0.714	0.00989493\\
0.716	0.0097583\\
0.718	0.00959319\\
0.72	0.00940057\\
0.722	0.00918145\\
0.724	0.00893695\\
0.726	0.00866826\\
0.728	0.00837664\\
0.73	0.00806341\\
0.732	0.00772996\\
0.734	0.00737772\\
0.736	0.00700818\\
0.738	0.00662287\\
0.74	0.00622336\\
0.742	0.00581124\\
0.744	0.00541524\\
0.746	0.00561864\\
0.748	0.00580633\\
0.75	0.00607366\\
0.752	0.00632293\\
0.754	0.00655356\\
0.756	0.00676503\\
0.758	0.00695689\\
0.76	0.00712875\\
0.762	0.00728028\\
0.764	0.00741124\\
0.766	0.00752144\\
0.768	0.00761075\\
0.77	0.00767913\\
0.772	0.00772658\\
0.774	0.00775317\\
0.776	0.00775905\\
0.778	0.00774442\\
0.78	0.00770954\\
0.782	0.00765473\\
0.784	0.00758037\\
0.786	0.0074869\\
0.788	0.0073748\\
0.79	0.00724463\\
0.792	0.00709696\\
0.794	0.00693245\\
0.796	0.00675177\\
0.798	0.00655565\\
0.8	0.00634485\\
0.802	0.00612018\\
0.804	0.00588247\\
0.806	0.00563259\\
0.808	0.00537143\\
0.81	0.00532844\\
0.812	0.00536207\\
0.814	0.00538213\\
0.816	0.00538874\\
0.818	0.00538202\\
0.82	0.00536214\\
0.822	0.00532931\\
0.824	0.00528378\\
0.826	0.00522583\\
0.828	0.00515577\\
0.83	0.00507396\\
0.832	0.00498077\\
0.834	0.00487661\\
0.836	0.00476191\\
0.838	0.00463714\\
0.84	0.00450279\\
0.842	0.00435937\\
0.844	0.0042074\\
0.846	0.00404744\\
0.848	0.00388005\\
0.85	0.00370582\\
0.852	0.00352534\\
0.854	0.00333922\\
0.856	0.00314958\\
0.858	0.00310885\\
0.86	0.00306825\\
0.862	0.00302785\\
0.864	0.00298769\\
0.866	0.00294785\\
0.868	0.0030078\\
0.87	0.00316634\\
0.872	0.00331073\\
0.874	0.00344069\\
0.876	0.00355602\\
0.878	0.00365657\\
0.88	0.0037422\\
0.882	0.00381286\\
0.884	0.00386854\\
0.886	0.00390925\\
0.888	0.00393507\\
0.89	0.00394612\\
0.892	0.00394257\\
0.894	0.0040015\\
0.896	0.00416767\\
0.898	0.00431913\\
0.9	0.00445552\\
0.902	0.00457652\\
0.904	0.00468187\\
0.906	0.00477139\\
0.908	0.00484494\\
0.91	0.00490243\\
0.912	0.00494386\\
0.914	0.00496926\\
0.916	0.00497873\\
0.918	0.00497243\\
0.92	0.00495056\\
0.922	0.00491339\\
0.924	0.00486123\\
0.926	0.00479443\\
0.928	0.00471342\\
0.93	0.00461865\\
0.932	0.00451062\\
0.934	0.00438987\\
0.936	0.00425699\\
0.938	0.00411258\\
0.94	0.0039573\\
0.942	0.00379184\\
0.944	0.00361689\\
0.946	0.00343319\\
0.948	0.0032415\\
0.95	0.0030426\\
0.952	0.00283727\\
0.954	0.00262632\\
0.956	0.00241055\\
0.958	0.00245132\\
0.96	0.00256582\\
0.962	0.002672\\
0.964	0.00276959\\
0.966	0.00285835\\
0.968	0.00293808\\
0.97	0.00300861\\
0.972	0.00306979\\
0.974	0.00312152\\
0.976	0.00316374\\
0.978	0.00319639\\
0.98	0.00321948\\
0.982	0.00323594\\
0.984	0.00333517\\
0.986	0.00342437\\
0.988	0.00350339\\
0.99	0.00357208\\
0.992	0.00363035\\
0.994	0.00367813\\
0.996	0.0037154\\
0.998	0.00374215\\
1	0.00375842\\
1.002	0.00376429\\
1.004	0.00375983\\
1.006	0.0037452\\
1.008	0.00372054\\
1.01	0.00368604\\
1.012	0.00364193\\
1.014	0.00358845\\
1.016	0.00352586\\
1.018	0.00345447\\
1.02	0.00337459\\
1.022	0.00328656\\
1.024	0.00319076\\
1.026	0.00308755\\
1.028	0.00297736\\
1.03	0.00286058\\
1.032	0.00273766\\
1.034	0.00260905\\
1.036	0.0024752\\
1.038	0.00238844\\
1.04	0.00240896\\
1.042	0.00242313\\
1.044	0.00243098\\
1.046	0.00243256\\
1.048	0.00242793\\
1.05	0.0024172\\
1.052	0.00240045\\
1.054	0.00237781\\
1.056	0.00234942\\
1.058	0.00231543\\
1.06	0.00227601\\
1.062	0.00223135\\
1.064	0.00218163\\
1.066	0.00212708\\
1.068	0.0020679\\
1.07	0.00200434\\
1.072	0.00193663\\
1.074	0.00186503\\
1.076	0.0017898\\
1.078	0.00171121\\
1.08	0.00162953\\
1.082	0.00154785\\
1.084	0.00158477\\
1.086	0.00161481\\
1.088	0.00163796\\
1.09	0.00165423\\
1.092	0.00166365\\
1.094	0.0016663\\
1.096	0.00166224\\
1.098	0.00169374\\
1.1	0.00175981\\
1.102	0.00186043\\
1.104	0.00195418\\
1.106	0.00204082\\
1.108	0.00212012\\
1.11	0.00219189\\
1.112	0.00225596\\
1.114	0.00231219\\
1.116	0.00236049\\
1.118	0.00240078\\
1.12	0.002433\\
1.122	0.00245714\\
1.124	0.00247322\\
1.126	0.00248127\\
1.128	0.00248137\\
1.13	0.00247361\\
1.132	0.00245812\\
1.134	0.00243504\\
1.136	0.00240455\\
1.138	0.00236685\\
1.14	0.00232215\\
1.142	0.00227072\\
1.144	0.0022128\\
1.146	0.00214868\\
1.148	0.00207867\\
1.15	0.00200309\\
1.152	0.00192227\\
1.154	0.00183656\\
1.156	0.00174633\\
1.158	0.00165195\\
1.16	0.00155381\\
1.162	0.00145229\\
1.164	0.00134779\\
1.166	0.00124074\\
1.168	0.00114365\\
1.17	0.00113471\\
1.172	0.00114303\\
1.174	0.00120888\\
1.176	0.00127069\\
1.178	0.00132829\\
1.18	0.00138155\\
1.182	0.00143033\\
1.184	0.0014745\\
1.186	0.00151398\\
1.188	0.00154868\\
1.19	0.00157853\\
1.192	0.00160349\\
1.194	0.00162353\\
1.196	0.00163862\\
1.198	0.00164879\\
1.2	0.00165403\\
1.202	0.0016544\\
1.204	0.00164994\\
1.206	0.00164071\\
1.208	0.0016268\\
1.21	0.00163294\\
1.212	0.00167747\\
1.214	0.0017169\\
1.216	0.00175117\\
1.218	0.00178024\\
1.22	0.00180407\\
1.222	0.00182265\\
1.224	0.00183598\\
1.226	0.00184408\\
1.228	0.00184698\\
1.23	0.00184472\\
1.232	0.00183739\\
1.234	0.00182505\\
1.236	0.00180781\\
1.238	0.00178577\\
1.24	0.00175905\\
1.242	0.00172779\\
1.244	0.00169214\\
1.246	0.00165227\\
1.248	0.00160833\\
1.25	0.00156052\\
1.252	0.00150903\\
1.254	0.00145406\\
1.256	0.00139581\\
1.258	0.00133451\\
1.26	0.00127038\\
1.262	0.00120366\\
1.264	0.00113457\\
1.266	0.00110064\\
1.268	0.0011106\\
1.27	0.00111757\\
1.272	0.00112157\\
1.274	0.00112262\\
1.276	0.00112076\\
1.278	0.00111602\\
1.28	0.00110846\\
1.282	0.00109812\\
1.284	0.00108508\\
1.286	0.0010694\\
1.288	0.00105118\\
1.29	0.00103049\\
1.292	0.00100742\\
1.294	0.000982079\\
1.296	0.000954568\\
1.298	0.000924998\\
1.3	0.000893482\\
1.302	0.00086014\\
1.304	0.000874651\\
1.306	0.00088574\\
1.308	0.000893243\\
1.31	0.000897179\\
1.312	0.00089758\\
1.314	0.000935431\\
1.316	0.000980838\\
1.318	0.00102267\\
1.32	0.00106081\\
1.322	0.00109519\\
1.324	0.00112571\\
1.326	0.00115232\\
1.328	0.00117499\\
1.33	0.00119366\\
1.332	0.00120833\\
1.334	0.001219\\
1.336	0.00122567\\
1.338	0.00122838\\
1.34	0.00122716\\
1.342	0.00122206\\
1.344	0.00121316\\
1.346	0.00120053\\
1.348	0.00118427\\
1.35	0.00116447\\
1.352	0.00114125\\
1.354	0.00111473\\
1.356	0.00108506\\
1.358	0.00105237\\
1.36	0.00101682\\
1.362	0.000978571\\
1.364	0.000937789\\
1.366	0.00089465\\
1.368	0.000849337\\
1.37	0.000802037\\
1.372	0.000752941\\
1.374	0.000702244\\
1.376	0.000650145\\
1.378	0.000596846\\
1.38	0.000568877\\
1.382	0.00056602\\
1.384	0.000563043\\
1.386	0.000559948\\
1.388	0.000558509\\
1.39	0.000595169\\
1.392	0.00062979\\
1.394	0.000662279\\
1.396	0.000692549\\
1.398	0.000720521\\
1.4	0.000746128\\
1.402	0.000769308\\
1.404	0.000790009\\
1.406	0.00080819\\
1.408	0.000823815\\
1.41	0.000836861\\
1.412	0.000847311\\
1.414	0.000855158\\
1.416	0.000860403\\
1.418	0.000863058\\
1.42	0.00086314\\
1.422	0.000860677\\
1.424	0.000855704\\
1.426	0.000848266\\
1.428	0.000838412\\
1.43	0.000826201\\
1.432	0.0008117\\
1.434	0.000794981\\
1.436	0.000813775\\
1.438	0.000835492\\
1.44	0.00085465\\
1.442	0.000871217\\
1.444	0.000885172\\
1.446	0.000896499\\
1.448	0.000905194\\
1.45	0.000911259\\
1.452	0.000914705\\
1.454	0.000915551\\
1.456	0.000913825\\
1.458	0.000909559\\
1.46	0.000902797\\
1.462	0.000893588\\
1.464	0.000881989\\
1.466	0.000868062\\
1.468	0.000851878\\
1.47	0.000833512\\
1.472	0.000813046\\
1.474	0.000790568\\
1.476	0.00076617\\
1.478	0.00073995\\
1.48	0.000712011\\
1.482	0.000682457\\
1.484	0.000651399\\
1.486	0.000618951\\
1.488	0.000585229\\
1.49	0.000550351\\
1.492	0.000514439\\
1.494	0.000509515\\
1.496	0.000513577\\
1.498	0.000516249\\
1.5	0.000517539\\
1.502	0.000517461\\
1.504	0.000516031\\
1.506	0.000513269\\
1.508	0.000509202\\
1.51	0.000503856\\
1.512	0.000497264\\
1.514	0.000489462\\
1.516	0.000480487\\
1.518	0.000470382\\
1.52	0.000465795\\
1.522	0.000473538\\
1.524	0.000479383\\
1.526	0.000483332\\
1.528	0.000485397\\
1.53	0.000503504\\
1.532	0.000521726\\
1.534	0.000538084\\
1.536	0.000552543\\
1.538	0.000565074\\
1.54	0.000575659\\
1.542	0.000584283\\
1.544	0.00059094\\
1.546	0.000595632\\
1.548	0.000598366\\
1.55	0.000599157\\
1.552	0.000598026\\
1.554	0.000595003\\
1.556	0.00059012\\
1.558	0.00058342\\
1.56	0.000574948\\
1.562	0.000564756\\
1.564	0.000552902\\
1.566	0.00053945\\
1.568	0.000524466\\
1.57	0.000516096\\
1.572	0.000512967\\
1.574	0.000508085\\
1.576	0.000501489\\
1.578	0.000493222\\
1.58	0.000483332\\
1.582	0.00047187\\
1.584	0.000458893\\
1.586	0.00044446\\
1.588	0.000428637\\
1.59	0.000411491\\
1.592	0.000393092\\
1.594	0.000373515\\
1.596	0.000352837\\
1.598	0.000331136\\
1.6	0.000308494\\
1.602	0.000284995\\
1.604	0.000269425\\
1.606	0.000284735\\
1.608	0.00030416\\
1.61	0.00032251\\
1.612	0.000339736\\
1.614	0.000355791\\
1.616	0.000370635\\
1.618	0.000384231\\
1.62	0.000396548\\
1.622	0.000407558\\
1.624	0.000417238\\
1.626	0.000425571\\
1.628	0.000432543\\
1.63	0.000438146\\
1.632	0.000442376\\
1.634	0.000445235\\
1.636	0.000446726\\
1.638	0.000446861\\
1.64	0.000445654\\
1.642	0.000443124\\
1.644	0.000439292\\
1.646	0.000434187\\
1.648	0.000427839\\
1.65	0.000420283\\
1.652	0.000411557\\
1.654	0.000401704\\
1.656	0.000390767\\
1.658	0.000378796\\
1.66	0.000391172\\
1.662	0.000403154\\
1.664	0.000413873\\
1.666	0.000423311\\
1.668	0.000431452\\
1.67	0.000438287\\
1.672	0.000443809\\
1.674	0.000448016\\
1.676	0.00045091\\
1.678	0.000452498\\
1.68	0.00045279\\
1.682	0.000451799\\
1.684	0.000449545\\
1.686	0.000446048\\
1.688	0.000441333\\
1.69	0.00043543\\
1.692	0.000428371\\
1.694	0.00042019\\
1.696	0.000410927\\
1.698	0.000400623\\
1.7	0.000389322\\
1.702	0.000377069\\
1.704	0.000363915\\
1.706	0.000349911\\
1.708	0.000335109\\
1.71	0.000319566\\
1.712	0.000303337\\
1.714	0.00028648\\
1.716	0.000269056\\
1.718	0.000251125\\
1.72	0.000232747\\
1.722	0.000234721\\
1.724	0.000236342\\
1.726	0.000237322\\
1.728	0.000237668\\
1.73	0.000237386\\
1.732	0.000236484\\
1.734	0.000238549\\
1.736	0.000244556\\
1.738	0.000249569\\
1.74	0.000253582\\
1.742	0.000256596\\
1.744	0.000258613\\
1.746	0.000265798\\
1.748	0.000272685\\
1.75	0.000278621\\
1.752	0.000283598\\
1.754	0.000287609\\
1.756	0.000290654\\
1.758	0.000292732\\
1.76	0.000293849\\
1.762	0.000294014\\
1.764	0.000293237\\
1.766	0.000291533\\
1.768	0.000288919\\
1.77	0.000285418\\
1.772	0.000281051\\
1.774	0.000275845\\
1.776	0.00026983\\
1.778	0.000263037\\
1.78	0.000255499\\
1.782	0.000247253\\
1.784	0.000250441\\
1.786	0.000254269\\
1.788	0.000257167\\
1.79	0.000259141\\
1.792	0.000260194\\
1.794	0.000260337\\
1.796	0.000259581\\
1.798	0.00025794\\
1.8	0.000255429\\
1.802	0.00025207\\
1.804	0.000247883\\
1.806	0.000242892\\
1.808	0.000237123\\
1.81	0.000230604\\
1.812	0.000223366\\
1.814	0.00021544\\
1.816	0.000206861\\
1.818	0.000197663\\
1.82	0.000187883\\
1.822	0.000177561\\
1.824	0.000166734\\
1.826	0.000157676\\
1.828	0.000167344\\
1.83	0.000176419\\
1.832	0.000184876\\
1.834	0.000192693\\
1.836	0.000199853\\
1.838	0.000206337\\
1.84	0.000212131\\
1.842	0.000217225\\
1.844	0.000221608\\
1.846	0.000225273\\
1.848	0.000228216\\
1.85	0.000230436\\
1.852	0.000231932\\
1.854	0.000232708\\
1.856	0.00023277\\
1.858	0.000232123\\
1.86	0.00023078\\
1.862	0.000228751\\
1.864	0.000226051\\
1.866	0.000222697\\
1.868	0.000218706\\
1.87	0.000214099\\
1.872	0.000208899\\
1.874	0.000203129\\
1.876	0.000196814\\
1.878	0.000189981\\
1.88	0.000182659\\
1.882	0.000178346\\
1.884	0.0001855\\
1.886	0.000192052\\
1.888	0.00019799\\
1.89	0.000203302\\
1.892	0.00020798\\
1.894	0.000212017\\
1.896	0.000215406\\
1.898	0.000218146\\
1.9	0.000220236\\
1.902	0.000221676\\
1.904	0.000222471\\
1.906	0.000222625\\
1.908	0.000222145\\
1.91	0.000221041\\
1.912	0.000219323\\
1.914	0.000217003\\
1.916	0.000214097\\
1.918	0.00021062\\
1.92	0.00020659\\
1.922	0.000202025\\
1.924	0.000196947\\
1.926	0.000191377\\
1.928	0.000185338\\
1.93	0.000178855\\
1.932	0.000171952\\
1.934	0.000164657\\
1.936	0.000156996\\
1.938	0.000148997\\
1.94	0.00014069\\
1.942	0.000132103\\
1.944	0.000123266\\
1.946	0.000115332\\
1.948	0.000119995\\
1.95	0.000124151\\
1.952	0.000127792\\
1.954	0.000130912\\
1.956	0.000133507\\
1.958	0.000135575\\
1.96	0.000137115\\
1.962	0.00013813\\
1.964	0.000139517\\
1.966	0.000141436\\
1.968	0.00014288\\
1.97	0.000143851\\
1.972	0.000144352\\
1.974	0.000144386\\
1.976	0.000143959\\
1.978	0.000143079\\
1.98	0.000141755\\
1.982	0.000139997\\
1.984	0.000137816\\
1.986	0.000135226\\
1.988	0.00013224\\
1.99	0.000128876\\
1.992	0.000125148\\
1.994	0.000121075\\
1.996	0.000116676\\
1.998	0.00011197\\
2	0.000107552\\
};
\addplot [color=mycolor4, line width=1.0pt, forget plot]
  table[row sep=crcr]{%
-0.024	0\\
-0.022	0\\
-0.02	0\\
-0.018	0\\
-0.016	0\\
-0.014	0\\
-0.012	0\\
-0.01	0\\
-0.008	0\\
-0.006	0\\
-0.004	0\\
-0.002	0\\
0	0\\
0.002	0.051675\\
0.004	0.092605\\
0.006	0.12345\\
0.008	0.145089\\
0.01	0.158554\\
0.012	0.16497\\
0.014	0.165503\\
0.016	0.161308\\
0.018	0.154878\\
0.02	0.146569\\
0.022	0.136185\\
0.024	0.126011\\
0.026	0.133192\\
0.028	0.139555\\
0.03	0.145049\\
0.032	0.149637\\
0.034	0.153288\\
0.036	0.155983\\
0.038	0.157714\\
0.04	0.15848\\
0.042	0.158295\\
0.044	0.157178\\
0.046	0.15516\\
0.048	0.152278\\
0.05	0.148579\\
0.052	0.144115\\
0.054	0.138945\\
0.056	0.133132\\
0.058	0.126745\\
0.06	0.119855\\
0.062	0.112534\\
0.064	0.110466\\
0.066	0.111818\\
0.068	0.1129\\
0.07	0.113711\\
0.072	0.114251\\
0.074	0.114521\\
0.076	0.114524\\
0.078	0.114262\\
0.08	0.113739\\
0.082	0.112959\\
0.084	0.111926\\
0.086	0.110647\\
0.088	0.109126\\
0.09	0.107372\\
0.092	0.105392\\
0.094	0.103194\\
0.096	0.100786\\
0.098	0.0981787\\
0.1	0.0953812\\
0.102	0.0924042\\
0.104	0.0892587\\
0.106	0.0859563\\
0.108	0.082509\\
0.11	0.078929\\
0.112	0.0793731\\
0.114	0.0800218\\
0.116	0.0805243\\
0.118	0.0808789\\
0.12	0.0810844\\
0.122	0.0811403\\
0.124	0.0810466\\
0.126	0.0808039\\
0.128	0.0804132\\
0.13	0.0798761\\
0.132	0.0791948\\
0.134	0.078372\\
0.136	0.0774109\\
0.138	0.0763151\\
0.14	0.0750887\\
0.142	0.0737364\\
0.144	0.0722631\\
0.146	0.0706743\\
0.148	0.0689757\\
0.15	0.0671735\\
0.152	0.0652741\\
0.154	0.0632842\\
0.156	0.0612108\\
0.158	0.0590613\\
0.16	0.0568429\\
0.162	0.0545633\\
0.164	0.0522302\\
0.166	0.051144\\
0.168	0.0524908\\
0.17	0.0535341\\
0.172	0.0542732\\
0.174	0.05471\\
0.176	0.0548491\\
0.178	0.0546976\\
0.18	0.054265\\
0.182	0.0535626\\
0.184	0.0526041\\
0.186	0.0514048\\
0.188	0.0499814\\
0.19	0.0483524\\
0.192	0.0465373\\
0.194	0.0445564\\
0.196	0.0424309\\
0.198	0.0401827\\
0.2	0.0378338\\
0.202	0.0354064\\
0.204	0.0344934\\
0.206	0.036204\\
0.208	0.037779\\
0.21	0.0392087\\
0.212	0.0404848\\
0.214	0.0415998\\
0.216	0.0425479\\
0.218	0.043324\\
0.22	0.0439247\\
0.222	0.0443476\\
0.224	0.0445915\\
0.226	0.0446567\\
0.228	0.0445443\\
0.23	0.044257\\
0.232	0.0437983\\
0.234	0.0431731\\
0.236	0.0423871\\
0.238	0.041447\\
0.24	0.0403606\\
0.242	0.0391364\\
0.244	0.0377838\\
0.246	0.0363128\\
0.248	0.034734\\
0.25	0.0330587\\
0.252	0.0312986\\
0.254	0.0306854\\
0.256	0.0326628\\
0.258	0.0345432\\
0.26	0.0363208\\
0.262	0.0379898\\
0.264	0.0395454\\
0.266	0.0409829\\
0.268	0.0422983\\
0.27	0.0434882\\
0.272	0.0445493\\
0.274	0.0454794\\
0.276	0.0462763\\
0.278	0.0469385\\
0.28	0.0474652\\
0.282	0.0478558\\
0.284	0.0481105\\
0.286	0.0482297\\
0.288	0.0482144\\
0.29	0.0480662\\
0.292	0.047787\\
0.294	0.0473793\\
0.296	0.0468458\\
0.298	0.0461897\\
0.3	0.0454149\\
0.302	0.0445251\\
0.304	0.0435249\\
0.306	0.0424189\\
0.308	0.0412122\\
0.31	0.0399101\\
0.312	0.0385181\\
0.314	0.0370421\\
0.316	0.0354882\\
0.318	0.0338625\\
0.32	0.0321716\\
0.322	0.030422\\
0.324	0.0286203\\
0.326	0.0267734\\
0.328	0.0254215\\
0.33	0.0257136\\
0.332	0.0259455\\
0.334	0.0261171\\
0.336	0.0262285\\
0.338	0.0262803\\
0.34	0.0262729\\
0.342	0.0262072\\
0.344	0.0260841\\
0.346	0.0259049\\
0.348	0.0256709\\
0.35	0.0253836\\
0.352	0.0250447\\
0.354	0.0246561\\
0.356	0.0242197\\
0.358	0.0237376\\
0.36	0.0232121\\
0.362	0.0226456\\
0.364	0.0220404\\
0.366	0.0213991\\
0.368	0.0207244\\
0.37	0.0200189\\
0.372	0.0192854\\
0.374	0.0185267\\
0.376	0.0182296\\
0.378	0.0184281\\
0.38	0.0185486\\
0.382	0.0185919\\
0.384	0.018559\\
0.386	0.0185838\\
0.388	0.018987\\
0.39	0.0192938\\
0.392	0.0195047\\
0.394	0.0196204\\
0.396	0.019642\\
0.398	0.019571\\
0.4	0.0194089\\
0.402	0.0191579\\
0.404	0.0188202\\
0.406	0.0183983\\
0.408	0.0178952\\
0.41	0.017314\\
0.412	0.0166578\\
0.414	0.0159304\\
0.416	0.0151354\\
0.418	0.0142767\\
0.42	0.0133586\\
0.422	0.0128422\\
0.424	0.0126981\\
0.426	0.0125551\\
0.428	0.0124133\\
0.43	0.0122731\\
0.432	0.0121345\\
0.434	0.0119979\\
0.436	0.0118632\\
0.438	0.0117308\\
0.44	0.0116007\\
0.442	0.0114731\\
0.444	0.0113482\\
0.446	0.011226\\
0.448	0.0111067\\
0.45	0.0109903\\
0.452	0.0108771\\
0.454	0.0107669\\
0.456	0.01066\\
0.458	0.0105564\\
0.46	0.0104562\\
0.462	0.011448\\
0.464	0.0125014\\
0.466	0.0135115\\
0.468	0.0144749\\
0.47	0.0153887\\
0.472	0.0162501\\
0.474	0.0170565\\
0.476	0.0178057\\
0.478	0.0184955\\
0.48	0.0191241\\
0.482	0.0196899\\
0.484	0.0201915\\
0.486	0.020628\\
0.488	0.0209983\\
0.49	0.0213019\\
0.492	0.0215385\\
0.494	0.0217081\\
0.496	0.0218106\\
0.498	0.0218467\\
0.5	0.0218168\\
0.502	0.0217218\\
0.504	0.0215629\\
0.506	0.0213413\\
0.508	0.0210585\\
0.51	0.0207164\\
0.512	0.0203167\\
0.514	0.0198617\\
0.516	0.0193535\\
0.518	0.0187946\\
0.52	0.0181876\\
0.522	0.0175353\\
0.524	0.0168404\\
0.526	0.016106\\
0.528	0.0153351\\
0.53	0.014531\\
0.532	0.0136968\\
0.534	0.0128359\\
0.536	0.0119516\\
0.538	0.0110474\\
0.54	0.0101267\\
0.542	0.00919298\\
0.544	0.00871592\\
0.546	0.00870191\\
0.548	0.00871212\\
0.55	0.00877197\\
0.552	0.00881138\\
0.554	0.00883053\\
0.556	0.00882965\\
0.558	0.00880907\\
0.56	0.00876915\\
0.562	0.00871032\\
0.564	0.00863308\\
0.566	0.00853797\\
0.568	0.008495\\
0.57	0.00846929\\
0.572	0.00844214\\
0.574	0.00841354\\
0.576	0.00838342\\
0.578	0.00835177\\
0.58	0.00831856\\
0.582	0.00828376\\
0.584	0.00824737\\
0.586	0.00826583\\
0.588	0.00868594\\
0.59	0.00906024\\
0.592	0.0093881\\
0.594	0.00966906\\
0.596	0.00990284\\
0.598	0.0100893\\
0.6	0.0102286\\
0.602	0.0103208\\
0.604	0.0103665\\
0.606	0.0103661\\
0.608	0.0103204\\
0.61	0.0102303\\
0.612	0.0100967\\
0.614	0.00992082\\
0.616	0.00970397\\
0.618	0.00944757\\
0.62	0.00915316\\
0.622	0.0088224\\
0.624	0.00845706\\
0.626	0.00805899\\
0.628	0.00763017\\
0.63	0.00717262\\
0.632	0.00694655\\
0.634	0.00687997\\
0.636	0.00681303\\
0.638	0.00674582\\
0.64	0.00667841\\
0.642	0.0066109\\
0.644	0.00654338\\
0.646	0.00647592\\
0.648	0.00640863\\
0.65	0.00634157\\
0.652	0.00627484\\
0.654	0.00620852\\
0.656	0.00614268\\
0.658	0.00607742\\
0.66	0.00601279\\
0.662	0.00594889\\
0.664	0.00588578\\
0.666	0.00582352\\
0.668	0.0057622\\
0.67	0.00570186\\
0.672	0.00564257\\
0.674	0.00558439\\
0.676	0.00606021\\
0.678	0.00653308\\
0.68	0.00698182\\
0.682	0.00740513\\
0.684	0.0078018\\
0.686	0.00817073\\
0.688	0.00851095\\
0.69	0.0088216\\
0.692	0.00910194\\
0.694	0.00935134\\
0.696	0.0095693\\
0.698	0.00975545\\
0.7	0.00990951\\
0.702	0.0100313\\
0.704	0.0101209\\
0.706	0.0101783\\
0.708	0.0102038\\
0.71	0.0101976\\
0.712	0.0101602\\
0.714	0.0100921\\
0.716	0.00999395\\
0.718	0.00986646\\
0.72	0.00971046\\
0.722	0.00952688\\
0.724	0.00931671\\
0.726	0.00908104\\
0.728	0.00882102\\
0.73	0.00853788\\
0.732	0.00823292\\
0.734	0.00790748\\
0.736	0.00756299\\
0.738	0.00720089\\
0.74	0.00682269\\
0.742	0.00642992\\
0.744	0.00602414\\
0.746	0.00560696\\
0.748	0.00517997\\
0.75	0.0047448\\
0.752	0.00432967\\
0.754	0.00431699\\
0.756	0.00430461\\
0.758	0.00429245\\
0.76	0.00428045\\
0.762	0.00426857\\
0.764	0.00425673\\
0.766	0.00424489\\
0.768	0.00423299\\
0.77	0.00422097\\
0.772	0.00420879\\
0.774	0.0041964\\
0.776	0.00418375\\
0.778	0.0041708\\
0.78	0.0041575\\
0.782	0.00414381\\
0.784	0.0041297\\
0.786	0.00411513\\
0.788	0.00410007\\
0.79	0.0040845\\
0.792	0.00406837\\
0.794	0.00405168\\
0.796	0.00404629\\
0.798	0.00425223\\
0.8	0.0044366\\
0.802	0.00459906\\
0.804	0.00473936\\
0.806	0.00485734\\
0.808	0.00495291\\
0.81	0.00502607\\
0.812	0.00507689\\
0.814	0.00510552\\
0.816	0.00511219\\
0.818	0.0050972\\
0.82	0.00506093\\
0.822	0.00500382\\
0.824	0.0049264\\
0.826	0.00482923\\
0.828	0.00471296\\
0.83	0.00457829\\
0.832	0.00442598\\
0.834	0.00425682\\
0.836	0.00407169\\
0.838	0.00387147\\
0.84	0.00365712\\
0.842	0.00346847\\
0.844	0.00343787\\
0.846	0.00340699\\
0.848	0.00337586\\
0.85	0.00334453\\
0.852	0.00331301\\
0.854	0.00328136\\
0.856	0.00324961\\
0.858	0.00321779\\
0.86	0.00318594\\
0.862	0.00315411\\
0.864	0.00312232\\
0.866	0.00309061\\
0.868	0.00305901\\
0.87	0.00302757\\
0.872	0.00299632\\
0.874	0.00296529\\
0.876	0.00293452\\
0.878	0.00290402\\
0.88	0.00287385\\
0.882	0.00284403\\
0.884	0.00281457\\
0.886	0.00278553\\
0.888	0.00295078\\
0.89	0.00317394\\
0.892	0.00338525\\
0.894	0.0035841\\
0.896	0.00376995\\
0.898	0.00394227\\
0.9	0.00410063\\
0.902	0.00424463\\
0.904	0.00437393\\
0.906	0.00448825\\
0.908	0.00458737\\
0.91	0.00467113\\
0.912	0.00473941\\
0.914	0.00479217\\
0.916	0.00482942\\
0.918	0.0048512\\
0.92	0.00485765\\
0.922	0.00484892\\
0.924	0.00482524\\
0.926	0.00478687\\
0.928	0.00473415\\
0.93	0.00466743\\
0.932	0.00458712\\
0.934	0.00449369\\
0.936	0.00438762\\
0.938	0.00426946\\
0.94	0.00413977\\
0.942	0.00399916\\
0.944	0.00384826\\
0.946	0.00368773\\
0.948	0.00351827\\
0.95	0.00334058\\
0.952	0.0031554\\
0.954	0.00296347\\
0.956	0.00276556\\
0.958	0.00256244\\
0.96	0.00235489\\
0.962	0.0021437\\
0.964	0.0020836\\
0.966	0.00207533\\
0.968	0.00206731\\
0.97	0.00205952\\
0.972	0.00205194\\
0.974	0.00204452\\
0.976	0.00203727\\
0.978	0.00203014\\
0.98	0.00202312\\
0.982	0.00201618\\
0.984	0.0020093\\
0.986	0.00200246\\
0.988	0.00199564\\
0.99	0.00198881\\
0.992	0.00198195\\
0.994	0.00197504\\
0.996	0.00196806\\
0.998	0.001961\\
1	0.00195384\\
1.002	0.00194656\\
1.004	0.00193914\\
1.006	0.00194129\\
1.008	0.002046\\
1.01	0.00214036\\
1.012	0.00222419\\
1.014	0.00229734\\
1.016	0.00235972\\
1.018	0.00241125\\
1.02	0.00245192\\
1.022	0.00248173\\
1.024	0.00250075\\
1.026	0.00250906\\
1.028	0.00250679\\
1.03	0.00249411\\
1.032	0.00247121\\
1.034	0.00243833\\
1.036	0.00239573\\
1.038	0.00234371\\
1.04	0.00228259\\
1.042	0.00221274\\
1.044	0.00213452\\
1.046	0.00204834\\
1.048	0.00195463\\
1.05	0.00185384\\
1.052	0.00174643\\
1.054	0.00168312\\
1.056	0.00166988\\
1.058	0.00165646\\
1.06	0.00164287\\
1.062	0.00162912\\
1.064	0.00161522\\
1.066	0.00160119\\
1.068	0.00158704\\
1.07	0.00157278\\
1.072	0.00155843\\
1.074	0.001544\\
1.076	0.00152951\\
1.078	0.00151497\\
1.08	0.0015004\\
1.082	0.00148581\\
1.084	0.00147122\\
1.086	0.00145664\\
1.088	0.00144208\\
1.09	0.00142757\\
1.092	0.00141312\\
1.094	0.00139874\\
1.096	0.00138445\\
1.098	0.00137025\\
1.1	0.00146084\\
1.102	0.00156624\\
1.104	0.00166575\\
1.106	0.00175909\\
1.108	0.00184599\\
1.11	0.00192623\\
1.112	0.00199962\\
1.114	0.00206596\\
1.116	0.00212511\\
1.118	0.00217696\\
1.12	0.00222142\\
1.122	0.0022584\\
1.124	0.00228789\\
1.126	0.00230987\\
1.128	0.00232436\\
1.13	0.0023314\\
1.132	0.00233107\\
1.134	0.00232345\\
1.136	0.00230867\\
1.138	0.00228687\\
1.14	0.00225822\\
1.142	0.00222291\\
1.144	0.00218115\\
1.146	0.00213317\\
1.148	0.00207922\\
1.15	0.00201957\\
1.152	0.0019545\\
1.154	0.00188431\\
1.156	0.00180932\\
1.158	0.00172985\\
1.16	0.00164624\\
1.162	0.00155885\\
1.164	0.00146802\\
1.166	0.00137412\\
1.168	0.00127753\\
1.17	0.00117862\\
1.172	0.00107776\\
1.174	0.00105848\\
1.176	0.00105042\\
1.178	0.00104226\\
1.18	0.00103402\\
1.182	0.0010257\\
1.184	0.0010173\\
1.186	0.00100885\\
1.188	0.00100034\\
1.19	0.000991797\\
1.192	0.000983217\\
1.194	0.000974614\\
1.196	0.000965999\\
1.198	0.00095738\\
1.2	0.000948768\\
1.202	0.000940172\\
1.204	0.000931601\\
1.206	0.0009248\\
1.208	0.000921301\\
1.21	0.000917822\\
1.212	0.000914353\\
1.214	0.000910884\\
1.216	0.000912793\\
1.218	0.000965656\\
1.22	0.00101359\\
1.222	0.0010565\\
1.224	0.00109431\\
1.226	0.00112695\\
1.228	0.00115439\\
1.23	0.00117661\\
1.232	0.0011936\\
1.234	0.00120539\\
1.236	0.001212\\
1.238	0.00121349\\
1.24	0.00120993\\
1.242	0.00120141\\
1.244	0.00118802\\
1.246	0.0011699\\
1.248	0.00114718\\
1.25	0.00112\\
1.252	0.00108853\\
1.254	0.00105294\\
1.256	0.00101343\\
1.258	0.000970193\\
1.26	0.000923438\\
1.262	0.000873386\\
1.264	0.000820266\\
1.266	0.000799399\\
1.268	0.000793869\\
1.27	0.000788247\\
1.272	0.000782535\\
1.274	0.000776738\\
1.276	0.000770859\\
1.278	0.000764902\\
1.28	0.000758871\\
1.282	0.000752772\\
1.284	0.000746609\\
1.286	0.000740388\\
1.288	0.000734113\\
1.29	0.000727791\\
1.292	0.000721426\\
1.294	0.000715026\\
1.296	0.000708596\\
1.298	0.000702142\\
1.3	0.000695671\\
1.302	0.000689189\\
1.304	0.000682702\\
1.306	0.000676216\\
1.308	0.000669738\\
1.31	0.000669378\\
1.312	0.000721544\\
1.314	0.000770959\\
1.316	0.00081748\\
1.318	0.000860977\\
1.32	0.000901334\\
1.322	0.000938447\\
1.324	0.000972224\\
1.326	0.00100259\\
1.328	0.00102947\\
1.33	0.00105283\\
1.332	0.00107262\\
1.334	0.00108882\\
1.336	0.00110141\\
1.338	0.0011104\\
1.34	0.00111581\\
1.342	0.00111765\\
1.344	0.00111597\\
1.346	0.00111082\\
1.348	0.00110227\\
1.35	0.00109038\\
1.352	0.00107524\\
1.354	0.00105696\\
1.356	0.00103562\\
1.358	0.00101136\\
1.36	0.000984288\\
1.362	0.000954541\\
1.364	0.000922258\\
1.366	0.000887587\\
1.368	0.000850682\\
1.37	0.000811701\\
1.372	0.00077081\\
1.374	0.000728177\\
1.376	0.000683976\\
1.378	0.000638383\\
1.38	0.000591576\\
1.382	0.000553185\\
1.384	0.000549593\\
1.386	0.000545947\\
1.388	0.000542251\\
1.39	0.000538505\\
1.392	0.000534713\\
1.394	0.000530879\\
1.396	0.000527004\\
1.398	0.000523093\\
1.4	0.000519147\\
1.402	0.000515172\\
1.404	0.00051117\\
1.406	0.000507146\\
1.408	0.000503102\\
1.41	0.000499042\\
1.412	0.000494971\\
1.414	0.000490892\\
1.416	0.000486809\\
1.418	0.000482725\\
1.42	0.000478646\\
1.422	0.000474574\\
1.424	0.000470513\\
1.426	0.000466467\\
1.428	0.00046244\\
1.43	0.000473071\\
1.432	0.000494985\\
1.434	0.000514474\\
1.436	0.000531504\\
1.438	0.00054605\\
1.44	0.000558099\\
1.442	0.000567643\\
1.444	0.000574686\\
1.446	0.000579241\\
1.448	0.000581326\\
1.45	0.000580973\\
1.452	0.000578218\\
1.454	0.000573105\\
1.456	0.00056569\\
1.458	0.000556032\\
1.46	0.000544199\\
1.462	0.000530265\\
1.464	0.000514312\\
1.466	0.000496427\\
1.468	0.000476702\\
1.47	0.000455236\\
1.472	0.000432131\\
1.474	0.000407494\\
1.476	0.000390566\\
1.478	0.000388427\\
1.48	0.000386251\\
1.482	0.000384039\\
1.484	0.000381791\\
1.486	0.000379507\\
1.488	0.000377187\\
1.49	0.000374833\\
1.492	0.000372445\\
1.494	0.000370024\\
1.496	0.000367571\\
1.498	0.000365088\\
1.5	0.000362575\\
1.502	0.000360035\\
1.504	0.000357468\\
1.506	0.000354877\\
1.508	0.000352264\\
1.51	0.000349629\\
1.512	0.000346976\\
1.514	0.000344306\\
1.516	0.000341621\\
1.518	0.000338924\\
1.52	0.000336217\\
1.522	0.000333502\\
1.524	0.000356124\\
1.526	0.000379196\\
1.528	0.000400854\\
1.53	0.000421038\\
1.532	0.000439695\\
1.534	0.000456779\\
1.536	0.000472247\\
1.538	0.000486065\\
1.54	0.000498206\\
1.542	0.000508647\\
1.544	0.000517373\\
1.546	0.000524375\\
1.548	0.00052965\\
1.55	0.000533202\\
1.552	0.00053504\\
1.554	0.00053518\\
1.556	0.000533642\\
1.558	0.000530455\\
1.56	0.00052565\\
1.562	0.000519266\\
1.564	0.000511346\\
1.566	0.000501937\\
1.568	0.000491092\\
1.57	0.000478868\\
1.572	0.000465327\\
1.574	0.000450533\\
1.576	0.000434555\\
1.578	0.000417464\\
1.58	0.000399337\\
1.582	0.000380249\\
1.584	0.000360281\\
1.586	0.000339516\\
1.588	0.000318036\\
1.59	0.000295926\\
1.592	0.000281283\\
1.594	0.000279591\\
1.596	0.000277879\\
1.598	0.000276147\\
1.6	0.000274395\\
1.602	0.000272625\\
1.604	0.000270836\\
1.606	0.00026903\\
1.608	0.000267208\\
1.61	0.00026537\\
1.612	0.000263518\\
1.614	0.000261653\\
1.616	0.000259775\\
1.618	0.000257887\\
1.62	0.000255989\\
1.622	0.000254082\\
1.624	0.000252169\\
1.626	0.00025025\\
1.628	0.000248327\\
1.63	0.000246402\\
1.632	0.000244475\\
1.634	0.000242548\\
1.636	0.000240623\\
1.638	0.000238701\\
1.64	0.000236784\\
1.642	0.000234873\\
1.644	0.000241322\\
1.646	0.000250124\\
1.648	0.000257741\\
1.65	0.000264164\\
1.652	0.000269388\\
1.654	0.000273412\\
1.656	0.000276239\\
1.658	0.000277877\\
1.66	0.000278337\\
1.662	0.000277635\\
1.664	0.000275791\\
1.666	0.000272829\\
1.668	0.000268775\\
1.67	0.000263661\\
1.672	0.000257519\\
1.674	0.000250388\\
1.676	0.000242308\\
1.678	0.000233322\\
1.68	0.000223475\\
1.682	0.000212815\\
1.684	0.000201393\\
1.686	0.000196919\\
1.688	0.000195467\\
1.69	0.000194042\\
1.692	0.000192646\\
1.694	0.000191278\\
1.696	0.000189938\\
1.698	0.000188626\\
1.7	0.000187343\\
1.702	0.000186087\\
1.704	0.000184858\\
1.706	0.000183656\\
1.708	0.000182481\\
1.71	0.000181332\\
1.712	0.000180209\\
1.714	0.00017911\\
1.716	0.000178036\\
1.718	0.000176986\\
1.72	0.000175958\\
1.722	0.000174953\\
1.724	0.000173968\\
1.726	0.000173005\\
1.728	0.00017206\\
1.73	0.000171135\\
1.732	0.000170227\\
1.734	0.000169335\\
1.736	0.000176233\\
1.738	0.000186974\\
1.74	0.000197023\\
1.742	0.000206353\\
1.744	0.00021494\\
1.746	0.000222762\\
1.748	0.000229801\\
1.75	0.000236042\\
1.752	0.000241474\\
1.754	0.000246086\\
1.756	0.000249874\\
1.758	0.000252833\\
1.76	0.000254965\\
1.762	0.000256273\\
1.764	0.000256762\\
1.766	0.00025644\\
1.768	0.000255321\\
1.77	0.000253418\\
1.772	0.000250747\\
1.774	0.000247329\\
1.776	0.000243185\\
1.778	0.000238339\\
1.78	0.000232817\\
1.782	0.000226648\\
1.784	0.000219862\\
1.786	0.000212491\\
1.788	0.000204569\\
1.79	0.00019613\\
1.792	0.000187212\\
1.794	0.000177853\\
1.796	0.00016809\\
1.798	0.000157964\\
1.8	0.000147515\\
1.802	0.000142303\\
1.804	0.000141455\\
1.806	0.000140599\\
1.808	0.000139736\\
1.81	0.000138866\\
1.812	0.000137989\\
1.814	0.000137105\\
1.816	0.000136214\\
1.818	0.000135317\\
1.82	0.000134413\\
1.822	0.000133505\\
1.824	0.00013259\\
1.826	0.000131671\\
1.828	0.000130747\\
1.83	0.00012982\\
1.832	0.000128888\\
1.834	0.000127954\\
1.836	0.000127017\\
1.838	0.000126078\\
1.84	0.000125138\\
1.842	0.000124197\\
1.844	0.000123255\\
1.846	0.000122314\\
1.848	0.000121374\\
1.85	0.000120436\\
1.852	0.000119499\\
1.854	0.000118566\\
1.856	0.000117635\\
1.858	0.000121552\\
1.86	0.000124924\\
1.862	0.000127722\\
1.864	0.000129944\\
1.866	0.000131592\\
1.868	0.000132667\\
1.87	0.000133175\\
1.872	0.000133122\\
1.874	0.000132516\\
1.876	0.000131368\\
1.878	0.000129691\\
1.88	0.000127497\\
1.882	0.000124802\\
1.884	0.000121624\\
1.886	0.000117981\\
1.888	0.000113893\\
1.89	0.000109381\\
1.892	0.000104468\\
1.894	0.000101332\\
1.896	0.000100577\\
1.898	9.98346e-05\\
1.9	9.91045e-05\\
1.902	9.83871e-05\\
1.904	9.76824e-05\\
1.906	9.69906e-05\\
1.908	9.63115e-05\\
1.91	9.56451e-05\\
1.912	9.49915e-05\\
1.914	9.43506e-05\\
1.916	9.37222e-05\\
1.918	9.31061e-05\\
1.92	9.25024e-05\\
1.922	9.19106e-05\\
1.924	9.13307e-05\\
1.926	9.07624e-05\\
1.928	9.02054e-05\\
1.93	8.96594e-05\\
1.932	8.91242e-05\\
1.934	8.85993e-05\\
1.936	8.80845e-05\\
1.938	8.75794e-05\\
1.94	8.70835e-05\\
1.942	8.65966e-05\\
1.944	8.61181e-05\\
1.946	8.56477e-05\\
1.948	8.71972e-05\\
1.95	9.21842e-05\\
1.952	9.6833e-05\\
1.954	0.000101131\\
1.956	0.000105068\\
1.958	0.000108634\\
1.96	0.000111822\\
1.962	0.000114624\\
1.964	0.000117035\\
1.966	0.000119053\\
1.968	0.000120674\\
1.97	0.000121899\\
1.972	0.000122728\\
1.974	0.000123162\\
1.976	0.000123207\\
1.978	0.000122866\\
1.98	0.000122145\\
1.982	0.000121053\\
1.984	0.000119597\\
1.986	0.000117788\\
1.988	0.000115637\\
1.99	0.000113155\\
1.992	0.000110356\\
1.994	0.000107254\\
1.996	0.000103863\\
1.998	0.0001002\\
2	9.62813e-05\\
};
\addplot [color=mycolor5, line width=1.0pt, forget plot]
  table[row sep=crcr]{%
-0.024	0\\
-0.022	0\\
-0.02	0\\
-0.018	0\\
-0.016	0\\
-0.014	0\\
-0.012	0\\
-0.01	0\\
-0.008	0\\
-0.006	0\\
-0.004	0\\
-0.002	0\\
0	0\\
0.002	0.0516749\\
0.004	0.092605\\
0.006	0.12345\\
0.008	0.145089\\
0.01	0.158553\\
0.012	0.164969\\
0.014	0.165501\\
0.016	0.161304\\
0.018	0.154875\\
0.02	0.146564\\
0.022	0.136179\\
0.024	0.126009\\
0.026	0.13319\\
0.028	0.139551\\
0.03	0.145045\\
0.032	0.149631\\
0.034	0.15328\\
0.036	0.155974\\
0.038	0.157702\\
0.04	0.158466\\
0.042	0.158277\\
0.044	0.157157\\
0.046	0.155134\\
0.048	0.152248\\
0.05	0.148543\\
0.052	0.144073\\
0.054	0.138897\\
0.056	0.133077\\
0.058	0.126682\\
0.06	0.119782\\
0.062	0.112452\\
0.064	0.104766\\
0.066	0.0988253\\
0.068	0.0991951\\
0.07	0.0997621\\
0.072	0.10011\\
0.074	0.100243\\
0.076	0.100164\\
0.078	0.0998782\\
0.08	0.0993907\\
0.082	0.0987068\\
0.084	0.097832\\
0.086	0.0967726\\
0.088	0.0955348\\
0.09	0.0941257\\
0.092	0.0925523\\
0.094	0.0908221\\
0.096	0.0889429\\
0.098	0.0869229\\
0.1	0.0847703\\
0.102	0.0824939\\
0.104	0.0801026\\
0.106	0.0776054\\
0.108	0.0750117\\
0.11	0.0723309\\
0.112	0.0723711\\
0.114	0.0727505\\
0.116	0.0730187\\
0.118	0.0731756\\
0.12	0.0732215\\
0.122	0.0731569\\
0.124	0.0729828\\
0.126	0.0727987\\
0.128	0.0725297\\
0.13	0.0721413\\
0.132	0.0716359\\
0.134	0.0710163\\
0.136	0.0702853\\
0.138	0.0694466\\
0.14	0.0685039\\
0.142	0.0674614\\
0.144	0.0663234\\
0.146	0.0650946\\
0.148	0.0637801\\
0.15	0.062385\\
0.152	0.0609148\\
0.154	0.059375\\
0.156	0.0577714\\
0.158	0.0561099\\
0.16	0.0543965\\
0.162	0.0526372\\
0.164	0.0508381\\
0.166	0.0505436\\
0.168	0.0519076\\
0.17	0.0529694\\
0.172	0.0537283\\
0.174	0.0541862\\
0.176	0.0543475\\
0.178	0.0542194\\
0.18	0.053811\\
0.182	0.0531339\\
0.184	0.0522015\\
0.186	0.051029\\
0.188	0.0496332\\
0.19	0.0480324\\
0.192	0.0462458\\
0.194	0.0442938\\
0.196	0.0421976\\
0.198	0.0399788\\
0.2	0.0376594\\
0.202	0.0352616\\
0.204	0.0340882\\
0.206	0.0358347\\
0.208	0.0374465\\
0.21	0.038914\\
0.212	0.0402285\\
0.214	0.0413828\\
0.216	0.0423706\\
0.218	0.043187\\
0.22	0.0438282\\
0.222	0.0442917\\
0.224	0.0445763\\
0.226	0.044682\\
0.228	0.04461\\
0.23	0.0443626\\
0.232	0.0439434\\
0.234	0.0433569\\
0.236	0.0426088\\
0.238	0.0417058\\
0.24	0.0406554\\
0.242	0.039466\\
0.244	0.038147\\
0.246	0.036708\\
0.248	0.0351599\\
0.25	0.0335135\\
0.252	0.0317806\\
0.254	0.029973\\
0.256	0.0281029\\
0.258	0.0261827\\
0.26	0.024225\\
0.262	0.0229193\\
0.264	0.0227865\\
0.266	0.0235222\\
0.268	0.0242012\\
0.27	0.0248222\\
0.272	0.0253839\\
0.274	0.0258853\\
0.276	0.0263258\\
0.278	0.0267047\\
0.28	0.027022\\
0.282	0.0272775\\
0.284	0.0274715\\
0.286	0.0276043\\
0.288	0.0276767\\
0.29	0.0276894\\
0.292	0.0276436\\
0.294	0.0275403\\
0.296	0.0273809\\
0.298	0.0271671\\
0.3	0.0269005\\
0.302	0.026583\\
0.304	0.0262164\\
0.306	0.0258031\\
0.308	0.025345\\
0.31	0.0248446\\
0.312	0.0243043\\
0.314	0.0237266\\
0.316	0.0231139\\
0.318	0.022469\\
0.32	0.0217945\\
0.322	0.0210931\\
0.324	0.0203675\\
0.326	0.0196205\\
0.328	0.0188548\\
0.33	0.0187065\\
0.332	0.0186172\\
0.334	0.0184994\\
0.336	0.0183538\\
0.338	0.0181812\\
0.34	0.0179826\\
0.342	0.0180543\\
0.344	0.0182205\\
0.346	0.0183539\\
0.348	0.0184545\\
0.35	0.0185226\\
0.352	0.0185585\\
0.354	0.0185625\\
0.356	0.0185352\\
0.358	0.0184772\\
0.36	0.018389\\
0.362	0.0182716\\
0.364	0.0181256\\
0.366	0.0179521\\
0.368	0.017752\\
0.37	0.0175264\\
0.372	0.0172764\\
0.374	0.0175484\\
0.376	0.0178256\\
0.378	0.0180267\\
0.38	0.0181518\\
0.382	0.0182015\\
0.384	0.0181771\\
0.386	0.01808\\
0.388	0.0179124\\
0.39	0.0176766\\
0.392	0.0173756\\
0.394	0.0170126\\
0.396	0.0165912\\
0.398	0.0161153\\
0.4	0.0155889\\
0.402	0.0150166\\
0.404	0.0144028\\
0.406	0.0137523\\
0.408	0.0132669\\
0.41	0.0131608\\
0.412	0.013056\\
0.414	0.0129526\\
0.416	0.0128508\\
0.418	0.0127505\\
0.42	0.0126519\\
0.422	0.012555\\
0.424	0.01246\\
0.426	0.0123669\\
0.428	0.0122756\\
0.43	0.0121863\\
0.432	0.0120991\\
0.434	0.0120138\\
0.436	0.0119305\\
0.438	0.0118493\\
0.44	0.0117701\\
0.442	0.011693\\
0.444	0.0116178\\
0.446	0.0115446\\
0.448	0.0114734\\
0.45	0.0114042\\
0.452	0.0113368\\
0.454	0.0112712\\
0.456	0.0112075\\
0.458	0.0111454\\
0.46	0.0110851\\
0.462	0.0110264\\
0.464	0.0109692\\
0.466	0.0109135\\
0.468	0.0108593\\
0.47	0.0108063\\
0.472	0.0107547\\
0.474	0.0107042\\
0.476	0.0106549\\
0.478	0.0106065\\
0.48	0.0105592\\
0.482	0.0105127\\
0.484	0.010467\\
0.486	0.010422\\
0.488	0.0103776\\
0.49	0.0103339\\
0.492	0.0102906\\
0.494	0.0102477\\
0.496	0.0102051\\
0.498	0.0101629\\
0.5	0.0101208\\
0.502	0.0100788\\
0.504	0.010037\\
0.506	0.00999515\\
0.508	0.00995327\\
0.51	0.00991129\\
0.512	0.00986917\\
0.514	0.00982685\\
0.516	0.0097843\\
0.518	0.00974146\\
0.52	0.00969832\\
0.522	0.00965482\\
0.524	0.00961094\\
0.526	0.00956666\\
0.528	0.00952195\\
0.53	0.00947678\\
0.532	0.00943115\\
0.534	0.00938503\\
0.536	0.00933842\\
0.538	0.00929132\\
0.54	0.0092437\\
0.542	0.00919558\\
0.544	0.00914696\\
0.546	0.00909783\\
0.548	0.00904821\\
0.55	0.0089981\\
0.552	0.00894751\\
0.554	0.00889646\\
0.556	0.00884497\\
0.558	0.00879304\\
0.56	0.0087407\\
0.562	0.00868798\\
0.564	0.00863488\\
0.566	0.00858144\\
0.568	0.00852767\\
0.57	0.00847361\\
0.572	0.00841929\\
0.574	0.00836471\\
0.576	0.00830993\\
0.578	0.00825496\\
0.58	0.00819983\\
0.582	0.00814457\\
0.584	0.00808922\\
0.586	0.00803379\\
0.588	0.00797833\\
0.59	0.00792285\\
0.592	0.0078674\\
0.594	0.00781199\\
0.596	0.00775665\\
0.598	0.00770142\\
0.6	0.00764632\\
0.602	0.00759138\\
0.604	0.00753662\\
0.606	0.00748207\\
0.608	0.00742776\\
0.61	0.00737371\\
0.612	0.00731994\\
0.614	0.00726647\\
0.616	0.00721334\\
0.618	0.00716054\\
0.62	0.00710812\\
0.622	0.00705608\\
0.624	0.00700445\\
0.626	0.00695324\\
0.628	0.00690247\\
0.63	0.00685215\\
0.632	0.00680229\\
0.634	0.00675292\\
0.636	0.00670403\\
0.638	0.00665565\\
0.64	0.00660779\\
0.642	0.00656044\\
0.644	0.00651363\\
0.646	0.00646735\\
0.648	0.00642161\\
0.65	0.00637642\\
0.652	0.00633179\\
0.654	0.00628771\\
0.656	0.00624418\\
0.658	0.00620122\\
0.66	0.00615881\\
0.662	0.00611697\\
0.664	0.00607567\\
0.666	0.00603494\\
0.668	0.00599475\\
0.67	0.00595511\\
0.672	0.005916\\
0.674	0.00587744\\
0.676	0.0058394\\
0.678	0.00580188\\
0.68	0.00576487\\
0.682	0.00572837\\
0.684	0.00569236\\
0.686	0.00565683\\
0.688	0.00562177\\
0.69	0.00558717\\
0.692	0.00555302\\
0.694	0.0055193\\
0.696	0.00548599\\
0.698	0.0054531\\
0.7	0.0054206\\
0.702	0.00538847\\
0.704	0.00535671\\
0.706	0.00532529\\
0.708	0.0052942\\
0.71	0.00526342\\
0.712	0.00523294\\
0.714	0.00520274\\
0.716	0.00517281\\
0.718	0.00514313\\
0.72	0.00511368\\
0.722	0.00508444\\
0.724	0.00505541\\
0.726	0.00502656\\
0.728	0.00499787\\
0.73	0.00496934\\
0.732	0.00494095\\
0.734	0.00491268\\
0.736	0.00488452\\
0.738	0.00485645\\
0.74	0.00482847\\
0.742	0.00480055\\
0.744	0.00477269\\
0.746	0.00474488\\
0.748	0.0047171\\
0.75	0.00468934\\
0.752	0.0046616\\
0.754	0.00463386\\
0.756	0.00460612\\
0.758	0.00457837\\
0.76	0.00455061\\
0.762	0.00452283\\
0.764	0.00449502\\
0.766	0.00446717\\
0.768	0.0044393\\
0.77	0.00441139\\
0.772	0.00438345\\
0.774	0.00435546\\
0.776	0.00432744\\
0.778	0.00429939\\
0.78	0.0042713\\
0.782	0.00424317\\
0.784	0.00421502\\
0.786	0.00418685\\
0.788	0.00415865\\
0.79	0.00413044\\
0.792	0.00410222\\
0.794	0.00407399\\
0.796	0.00404577\\
0.798	0.00401756\\
0.8	0.00398937\\
0.802	0.00396121\\
0.804	0.00393309\\
0.806	0.003905\\
0.808	0.00387698\\
0.81	0.00384901\\
0.812	0.00382112\\
0.814	0.00379331\\
0.816	0.0037656\\
0.818	0.00373798\\
0.82	0.00371048\\
0.822	0.0036831\\
0.824	0.00365585\\
0.826	0.00362874\\
0.828	0.00360179\\
0.83	0.00357499\\
0.832	0.00354836\\
0.834	0.00352191\\
0.836	0.00349564\\
0.838	0.00346957\\
0.84	0.0034437\\
0.842	0.00341804\\
0.844	0.00339259\\
0.846	0.00336736\\
0.848	0.00334237\\
0.85	0.0033176\\
0.852	0.00329308\\
0.854	0.00326879\\
0.856	0.00324476\\
0.858	0.00322097\\
0.86	0.00319744\\
0.862	0.00317416\\
0.864	0.00315114\\
0.866	0.00312838\\
0.868	0.00310587\\
0.87	0.00308363\\
0.872	0.00306165\\
0.874	0.00303993\\
0.876	0.00301846\\
0.878	0.00299726\\
0.88	0.0029763\\
0.882	0.0029556\\
0.884	0.00293515\\
0.886	0.00291495\\
0.888	0.00289499\\
0.89	0.00287527\\
0.892	0.00285578\\
0.894	0.00283652\\
0.896	0.00281749\\
0.898	0.00279868\\
0.9	0.00278009\\
0.902	0.00276171\\
0.904	0.00274353\\
0.906	0.00272555\\
0.908	0.00270776\\
0.91	0.00269016\\
0.912	0.00267274\\
0.914	0.0026555\\
0.916	0.00263842\\
0.918	0.00262151\\
0.92	0.00260476\\
0.922	0.00258815\\
0.924	0.0025717\\
0.926	0.00255538\\
0.928	0.00253919\\
0.93	0.00252313\\
0.932	0.0025072\\
0.934	0.00249138\\
0.936	0.00247567\\
0.938	0.00246007\\
0.94	0.00244458\\
0.942	0.00242917\\
0.944	0.00241387\\
0.946	0.00239865\\
0.948	0.00238351\\
0.95	0.00236845\\
0.952	0.00235347\\
0.954	0.00233857\\
0.956	0.00232373\\
0.958	0.00230896\\
0.96	0.00229426\\
0.962	0.00227961\\
0.964	0.00226503\\
0.966	0.0022505\\
0.968	0.00223603\\
0.97	0.00222161\\
0.972	0.00220725\\
0.974	0.00219293\\
0.976	0.00217867\\
0.978	0.00216445\\
0.98	0.00215029\\
0.982	0.00213617\\
0.984	0.0021221\\
0.986	0.00210807\\
0.988	0.0020941\\
0.99	0.00208017\\
0.992	0.00206629\\
0.994	0.00205245\\
0.996	0.00203867\\
0.998	0.00202493\\
1	0.00201125\\
1.002	0.00199761\\
1.004	0.00198403\\
1.006	0.0019705\\
1.008	0.00195702\\
1.01	0.00194359\\
1.012	0.00193023\\
1.014	0.00191692\\
1.016	0.00190366\\
1.018	0.00189047\\
1.02	0.00187734\\
1.022	0.00186427\\
1.024	0.00185126\\
1.026	0.00183832\\
1.028	0.00182544\\
1.03	0.00181263\\
1.032	0.00179988\\
1.034	0.00178721\\
1.036	0.00177461\\
1.038	0.00176207\\
1.04	0.00174961\\
1.042	0.00173723\\
1.044	0.00172492\\
1.046	0.00171268\\
1.048	0.00170053\\
1.05	0.00168845\\
1.052	0.00167645\\
1.054	0.00166452\\
1.056	0.00165268\\
1.058	0.00164092\\
1.06	0.00162925\\
1.062	0.00161765\\
1.064	0.00160614\\
1.066	0.00159471\\
1.068	0.00158336\\
1.07	0.0015721\\
1.072	0.00156092\\
1.074	0.00154983\\
1.076	0.00153882\\
1.078	0.0015279\\
1.08	0.00151706\\
1.082	0.00150631\\
1.084	0.00149564\\
1.086	0.00148506\\
1.088	0.00147456\\
1.09	0.00146415\\
1.092	0.00145382\\
1.094	0.00144358\\
1.096	0.00143342\\
1.098	0.00142335\\
1.1	0.00141335\\
1.102	0.00140345\\
1.104	0.00139362\\
1.106	0.00138388\\
1.108	0.00137422\\
1.11	0.00136464\\
1.112	0.00135514\\
1.114	0.00134572\\
1.116	0.00133638\\
1.118	0.00132711\\
1.12	0.00131793\\
1.122	0.00130882\\
1.124	0.00129979\\
1.126	0.00129083\\
1.128	0.00128195\\
1.13	0.00127314\\
1.132	0.00126441\\
1.134	0.00125575\\
1.136	0.00124715\\
1.138	0.00123863\\
1.14	0.00123018\\
1.142	0.00122179\\
1.144	0.00121348\\
1.146	0.00120523\\
1.148	0.00119704\\
1.15	0.00118892\\
1.152	0.00118087\\
1.154	0.00117287\\
1.156	0.00116494\\
1.158	0.00115707\\
1.16	0.00114926\\
1.162	0.0011415\\
1.164	0.00113381\\
1.166	0.00112617\\
1.168	0.00111859\\
1.17	0.00111106\\
1.172	0.00110359\\
1.174	0.00109617\\
1.176	0.0010888\\
1.178	0.00108148\\
1.18	0.00107422\\
1.182	0.001067\\
1.184	0.00105984\\
1.186	0.00105272\\
1.188	0.00104565\\
1.19	0.00103862\\
1.192	0.00103165\\
1.194	0.00102471\\
1.196	0.00101782\\
1.198	0.00101098\\
1.2	0.00100417\\
1.202	0.000997414\\
1.204	0.000990695\\
1.206	0.000984018\\
1.208	0.000977381\\
1.21	0.000970784\\
1.212	0.000964226\\
1.214	0.000957707\\
1.216	0.000951227\\
1.218	0.000944785\\
1.22	0.00093838\\
1.222	0.000932012\\
1.224	0.000925681\\
1.226	0.000919386\\
1.228	0.000913128\\
1.23	0.000906905\\
1.232	0.000900717\\
1.234	0.000894565\\
1.236	0.000888447\\
1.238	0.000882364\\
1.24	0.000876316\\
1.242	0.000870302\\
1.244	0.000864322\\
1.246	0.000858375\\
1.248	0.000852463\\
1.25	0.000846585\\
1.252	0.00084074\\
1.254	0.000834929\\
1.256	0.000829151\\
1.258	0.000823407\\
1.26	0.000817697\\
1.262	0.00081202\\
1.264	0.000806376\\
1.266	0.000800766\\
1.268	0.000795189\\
1.27	0.000789646\\
1.272	0.000784137\\
1.274	0.000778661\\
1.276	0.000773219\\
1.278	0.000767811\\
1.28	0.000762436\\
1.282	0.000757096\\
1.284	0.000751789\\
1.286	0.000746517\\
1.288	0.000741278\\
1.29	0.000736074\\
1.292	0.000730904\\
1.294	0.000725769\\
1.296	0.000720668\\
1.298	0.000715601\\
1.3	0.00071057\\
1.302	0.000705572\\
1.304	0.00070061\\
1.306	0.000695682\\
1.308	0.000690789\\
1.31	0.000685931\\
1.312	0.000681107\\
1.314	0.000676319\\
1.316	0.000671565\\
1.318	0.000666846\\
1.32	0.000662162\\
1.322	0.000657512\\
1.324	0.000652898\\
1.326	0.000648317\\
1.328	0.000643772\\
1.33	0.000639261\\
1.332	0.000634785\\
1.334	0.000630342\\
1.336	0.000625934\\
1.338	0.00062156\\
1.34	0.00061722\\
1.342	0.00061313\\
1.344	0.000609199\\
1.346	0.000605297\\
1.348	0.000601422\\
1.35	0.000597574\\
1.352	0.000593753\\
1.354	0.00058996\\
1.356	0.000586193\\
1.358	0.000582453\\
1.36	0.00057874\\
1.362	0.000575053\\
1.364	0.000571392\\
1.366	0.000567757\\
1.368	0.000564148\\
1.37	0.000560564\\
1.372	0.000557005\\
1.374	0.000553472\\
1.376	0.000549963\\
1.378	0.000546479\\
1.38	0.000543019\\
1.382	0.000539583\\
1.384	0.00053617\\
1.386	0.000532781\\
1.388	0.000529416\\
1.39	0.000526073\\
1.392	0.000522753\\
1.394	0.000519455\\
1.396	0.000516179\\
1.398	0.000512925\\
1.4	0.000509692\\
1.402	0.000506481\\
1.404	0.000503291\\
1.406	0.000500121\\
1.408	0.000496971\\
1.41	0.000493842\\
1.412	0.000490733\\
1.414	0.000487643\\
1.416	0.000484572\\
1.418	0.000481521\\
1.42	0.000478488\\
1.422	0.000475474\\
1.424	0.000472478\\
1.426	0.0004695\\
1.428	0.000466541\\
1.43	0.000463599\\
1.432	0.000460674\\
1.434	0.000457767\\
1.436	0.000454876\\
1.438	0.000452003\\
1.44	0.000449146\\
1.442	0.000446306\\
1.444	0.000443482\\
1.446	0.000440674\\
1.448	0.000437883\\
1.45	0.000435107\\
1.452	0.000432347\\
1.454	0.000429603\\
1.456	0.000426874\\
1.458	0.000424161\\
1.46	0.000421463\\
1.462	0.000418781\\
1.464	0.000416113\\
1.466	0.00041346\\
1.468	0.000410823\\
1.47	0.0004082\\
1.472	0.000405592\\
1.474	0.000402999\\
1.476	0.000400421\\
1.478	0.000397857\\
1.48	0.000395308\\
1.482	0.000392774\\
1.484	0.000390254\\
1.486	0.000387748\\
1.488	0.000385257\\
1.49	0.000382781\\
1.492	0.000380318\\
1.494	0.00037787\\
1.496	0.000375437\\
1.498	0.000373018\\
1.5	0.000370613\\
1.502	0.000368222\\
1.504	0.000365846\\
1.506	0.000363484\\
1.508	0.000361136\\
1.51	0.000358802\\
1.512	0.000356482\\
1.514	0.000354177\\
1.516	0.000351886\\
1.518	0.000349609\\
1.52	0.000347345\\
1.522	0.000345096\\
1.524	0.000342861\\
1.526	0.00034064\\
1.528	0.000338433\\
1.53	0.00033624\\
1.532	0.000334061\\
1.534	0.000331896\\
1.536	0.000329744\\
1.538	0.000327606\\
1.54	0.000325482\\
1.542	0.000323372\\
1.544	0.000321275\\
1.546	0.000319192\\
1.548	0.000317122\\
1.55	0.000315066\\
1.552	0.000313024\\
1.554	0.000310994\\
1.556	0.000308978\\
1.558	0.000306976\\
1.56	0.000304986\\
1.562	0.000303009\\
1.564	0.000301046\\
1.566	0.000299095\\
1.568	0.000297158\\
1.57	0.000295233\\
1.572	0.000293321\\
1.574	0.000291422\\
1.576	0.000289535\\
1.578	0.00028766\\
1.58	0.000285799\\
1.582	0.000283949\\
1.584	0.000282112\\
1.586	0.000280287\\
1.588	0.000278474\\
1.59	0.000276673\\
1.592	0.000274884\\
1.594	0.000273106\\
1.596	0.000271341\\
1.598	0.000269587\\
1.6	0.000267845\\
1.602	0.000266114\\
1.604	0.000264395\\
1.606	0.000262687\\
1.608	0.00026099\\
1.61	0.000259304\\
1.612	0.00025763\\
1.614	0.000255966\\
1.616	0.000254313\\
1.618	0.000252671\\
1.62	0.00025104\\
1.622	0.00024942\\
1.624	0.00024781\\
1.626	0.00024621\\
1.628	0.000244621\\
1.63	0.000243042\\
1.632	0.000241474\\
1.634	0.000239915\\
1.636	0.000238367\\
1.638	0.000236855\\
1.64	0.000235357\\
1.642	0.000233868\\
1.644	0.000232388\\
1.646	0.000230918\\
1.648	0.000229456\\
1.65	0.000228004\\
1.652	0.00022656\\
1.654	0.000225126\\
1.656	0.0002237\\
1.658	0.000222283\\
1.66	0.000220875\\
1.662	0.000219476\\
1.664	0.000218086\\
1.666	0.000216704\\
1.668	0.000215331\\
1.67	0.000213966\\
1.672	0.000212611\\
1.674	0.000211263\\
1.676	0.000209924\\
1.678	0.000208594\\
1.68	0.000207272\\
1.682	0.000205958\\
1.684	0.000204652\\
1.686	0.000203355\\
1.688	0.000202066\\
1.69	0.000200786\\
1.692	0.000199513\\
1.694	0.000198248\\
1.696	0.000196992\\
1.698	0.000195743\\
1.7	0.000194502\\
1.702	0.000193269\\
1.704	0.000192044\\
1.706	0.000190827\\
1.708	0.000189618\\
1.71	0.000188416\\
1.712	0.000187222\\
1.714	0.000186036\\
1.716	0.000184857\\
1.718	0.000183685\\
1.72	0.000182521\\
1.722	0.000181365\\
1.724	0.000180216\\
1.726	0.000179074\\
1.728	0.000177939\\
1.73	0.000176812\\
1.732	0.000175691\\
1.734	0.000174578\\
1.736	0.000173472\\
1.738	0.000172373\\
1.74	0.000171281\\
1.742	0.000170196\\
1.744	0.000169118\\
1.746	0.000168047\\
1.748	0.000166982\\
1.75	0.000165924\\
1.752	0.000164873\\
1.754	0.000163828\\
1.756	0.00016279\\
1.758	0.000161759\\
1.76	0.000160734\\
1.762	0.000159716\\
1.764	0.000158704\\
1.766	0.000157698\\
1.768	0.000156699\\
1.77	0.000155705\\
1.772	0.000154719\\
1.774	0.000153738\\
1.776	0.000152763\\
1.778	0.000151795\\
1.78	0.000150832\\
1.782	0.000149876\\
1.784	0.000148925\\
1.786	0.000147981\\
1.788	0.000147042\\
1.79	0.000146109\\
1.792	0.000145182\\
1.794	0.00014426\\
1.796	0.000143345\\
1.798	0.000142435\\
1.8	0.00014153\\
1.802	0.000140631\\
1.804	0.000139738\\
1.806	0.00013885\\
1.808	0.000137968\\
1.81	0.000137091\\
1.812	0.000136219\\
1.814	0.000135353\\
1.816	0.000134493\\
1.818	0.000133637\\
1.82	0.000132787\\
1.822	0.000131942\\
1.824	0.000131102\\
1.826	0.000130267\\
1.828	0.000129438\\
1.83	0.000128613\\
1.832	0.000127794\\
1.834	0.000126979\\
1.836	0.00012617\\
1.838	0.000125366\\
1.84	0.000124566\\
1.842	0.000123772\\
1.844	0.000122982\\
1.846	0.000122197\\
1.848	0.000121417\\
1.85	0.000120642\\
1.852	0.000119872\\
1.854	0.000119106\\
1.856	0.000118345\\
1.858	0.000117589\\
1.86	0.000116838\\
1.862	0.000116091\\
1.864	0.000115349\\
1.866	0.000114611\\
1.868	0.000113878\\
1.87	0.00011315\\
1.872	0.000112426\\
1.874	0.000111707\\
1.876	0.000110992\\
1.878	0.000110282\\
1.88	0.000109576\\
1.882	0.000108874\\
1.884	0.000108177\\
1.886	0.000107485\\
1.888	0.000106796\\
1.89	0.000106112\\
1.892	0.000105433\\
1.894	0.000104757\\
1.896	0.000104086\\
1.898	0.000103419\\
1.9	0.000102757\\
1.902	0.000102098\\
1.904	0.000101444\\
1.906	0.000100794\\
1.908	0.000100148\\
1.91	9.95058e-05\\
1.912	9.8868e-05\\
1.914	9.82342e-05\\
1.916	9.76046e-05\\
1.918	9.69789e-05\\
1.92	9.63572e-05\\
1.922	9.57395e-05\\
1.924	9.51257e-05\\
1.926	9.45159e-05\\
1.928	9.391e-05\\
1.93	9.33079e-05\\
1.932	9.27098e-05\\
1.934	9.21154e-05\\
1.936	9.15249e-05\\
1.938	9.09381e-05\\
1.94	9.03551e-05\\
1.942	8.97758e-05\\
1.944	8.92003e-05\\
1.946	8.86284e-05\\
1.948	8.80602e-05\\
1.95	8.74957e-05\\
1.952	8.69348e-05\\
1.954	8.63774e-05\\
1.956	8.58237e-05\\
1.958	8.52735e-05\\
1.96	8.47268e-05\\
1.962	8.41836e-05\\
1.964	8.36439e-05\\
1.966	8.31076e-05\\
1.968	8.25748e-05\\
1.97	8.20454e-05\\
1.972	8.15194e-05\\
1.974	8.09967e-05\\
1.976	8.04774e-05\\
1.978	7.99614e-05\\
1.98	7.94487e-05\\
1.982	7.89392e-05\\
1.984	7.8433e-05\\
1.986	7.793e-05\\
1.988	7.74303e-05\\
1.99	7.69337e-05\\
1.992	7.64402e-05\\
1.994	7.59499e-05\\
1.996	7.54628e-05\\
1.998	7.49787e-05\\
2	7.44976e-05\\
};
\addplot [color=mycolor6, line width=1.0pt, forget plot]
  table[row sep=crcr]{%
-0.024	0\\
-0.022	0\\
-0.02	0\\
-0.018	0\\
-0.016	0\\
-0.014	0\\
-0.012	0\\
-0.01	0\\
-0.008	0\\
-0.006	0\\
-0.004	0\\
-0.002	0\\
0	0\\
0.002	0.0516749\\
0.004	0.092605\\
0.006	0.12345\\
0.008	0.145089\\
0.01	0.158553\\
0.012	0.164968\\
0.014	0.1655\\
0.016	0.161302\\
0.018	0.154871\\
0.02	0.146559\\
0.022	0.136172\\
0.024	0.125991\\
0.026	0.133166\\
0.028	0.139519\\
0.03	0.145004\\
0.032	0.149579\\
0.034	0.153216\\
0.036	0.155894\\
0.038	0.157605\\
0.04	0.158351\\
0.042	0.158141\\
0.044	0.156997\\
0.046	0.154949\\
0.048	0.152034\\
0.05	0.148299\\
0.052	0.143797\\
0.054	0.138585\\
0.056	0.132729\\
0.058	0.126295\\
0.06	0.119356\\
0.062	0.111984\\
0.064	0.104256\\
0.066	0.0982845\\
0.068	0.0990596\\
0.07	0.0996115\\
0.072	0.0999433\\
0.074	0.100059\\
0.076	0.0999616\\
0.078	0.0996567\\
0.08	0.0991488\\
0.082	0.0984434\\
0.084	0.0975461\\
0.086	0.0964632\\
0.088	0.095201\\
0.09	0.0937666\\
0.092	0.092167\\
0.094	0.0904099\\
0.096	0.0885032\\
0.098	0.086455\\
0.1	0.0842738\\
0.102	0.0819684\\
0.104	0.0795478\\
0.106	0.0770212\\
0.108	0.074398\\
0.11	0.0716879\\
0.112	0.0714812\\
0.114	0.0718202\\
0.116	0.0720486\\
0.118	0.0721665\\
0.12	0.0721744\\
0.122	0.0722245\\
0.124	0.0721756\\
0.126	0.0720052\\
0.128	0.0717149\\
0.13	0.0713068\\
0.132	0.0707832\\
0.134	0.0701469\\
0.136	0.0694013\\
0.138	0.0685498\\
0.14	0.0675964\\
0.142	0.0665453\\
0.144	0.065401\\
0.146	0.0641683\\
0.148	0.0628524\\
0.15	0.0614585\\
0.152	0.059992\\
0.154	0.0584586\\
0.156	0.0568641\\
0.158	0.0552145\\
0.16	0.0535156\\
0.162	0.0517737\\
0.164	0.0499947\\
0.166	0.0515491\\
0.168	0.0529626\\
0.17	0.0540667\\
0.172	0.0548604\\
0.174	0.0553453\\
0.176	0.0555256\\
0.178	0.0554081\\
0.18	0.0550019\\
0.182	0.0543184\\
0.184	0.0533709\\
0.186	0.0521747\\
0.188	0.0507468\\
0.19	0.0491053\\
0.192	0.0472699\\
0.194	0.0452612\\
0.196	0.0431007\\
0.198	0.0408103\\
0.2	0.0384126\\
0.202	0.0359301\\
0.204	0.0333855\\
0.206	0.0308011\\
0.208	0.0287386\\
0.21	0.0297056\\
0.212	0.0305797\\
0.214	0.031356\\
0.216	0.0320302\\
0.218	0.032599\\
0.22	0.0330598\\
0.222	0.0334107\\
0.224	0.0336509\\
0.226	0.0337802\\
0.228	0.0337991\\
0.23	0.0337091\\
0.232	0.0335121\\
0.234	0.0332109\\
0.236	0.032809\\
0.238	0.0323104\\
0.24	0.0317197\\
0.242	0.0310421\\
0.244	0.0302831\\
0.246	0.0294488\\
0.248	0.0285456\\
0.25	0.0275802\\
0.252	0.0265597\\
0.254	0.0254912\\
0.256	0.0244268\\
0.258	0.0242503\\
0.26	0.0240721\\
0.262	0.0238922\\
0.264	0.0237107\\
0.266	0.0242802\\
0.268	0.0249237\\
0.27	0.0255062\\
0.272	0.0260268\\
0.274	0.0264846\\
0.276	0.0268792\\
0.278	0.0272103\\
0.28	0.0274778\\
0.282	0.0276821\\
0.284	0.0278235\\
0.286	0.0279027\\
0.288	0.0279207\\
0.29	0.0278784\\
0.292	0.0277772\\
0.294	0.0276186\\
0.296	0.0274042\\
0.298	0.0271357\\
0.3	0.0268152\\
0.302	0.0264446\\
0.304	0.0260263\\
0.306	0.0255626\\
0.308	0.0250558\\
0.31	0.0245086\\
0.312	0.0239234\\
0.314	0.0233031\\
0.316	0.0226503\\
0.318	0.0219677\\
0.32	0.0212583\\
0.322	0.0205248\\
0.324	0.0197701\\
0.326	0.0189971\\
0.328	0.0182617\\
0.33	0.0181633\\
0.332	0.0180383\\
0.334	0.0178877\\
0.336	0.0177122\\
0.338	0.0175128\\
0.34	0.0172903\\
0.342	0.0170782\\
0.344	0.0172332\\
0.346	0.0173597\\
0.348	0.0174578\\
0.35	0.0175277\\
0.352	0.0175694\\
0.354	0.0175835\\
0.356	0.0175702\\
0.358	0.0175299\\
0.36	0.0174633\\
0.362	0.0173708\\
0.364	0.0172533\\
0.366	0.0171113\\
0.368	0.0169458\\
0.37	0.0167575\\
0.372	0.0165473\\
0.374	0.0163162\\
0.376	0.0160651\\
0.378	0.0157952\\
0.38	0.0155074\\
0.382	0.0152028\\
0.384	0.0148826\\
0.386	0.0145478\\
0.388	0.0142732\\
0.39	0.0141935\\
0.392	0.0141135\\
0.394	0.0140333\\
0.396	0.0139529\\
0.398	0.0138724\\
0.4	0.0137918\\
0.402	0.0137112\\
0.404	0.0136308\\
0.406	0.0135505\\
0.408	0.0134704\\
0.41	0.0133906\\
0.412	0.0133112\\
0.414	0.0132321\\
0.416	0.0131535\\
0.418	0.0130755\\
0.42	0.012998\\
0.422	0.0129211\\
0.424	0.0128449\\
0.426	0.0127694\\
0.428	0.0126946\\
0.43	0.0126206\\
0.432	0.0125474\\
0.434	0.012475\\
0.436	0.0124035\\
0.438	0.0123329\\
0.44	0.0122631\\
0.442	0.0121943\\
0.444	0.0121264\\
0.446	0.0120593\\
0.448	0.0119932\\
0.45	0.011928\\
0.452	0.0118637\\
0.454	0.0118003\\
0.456	0.0117378\\
0.458	0.0116761\\
0.46	0.0116152\\
0.462	0.0115552\\
0.464	0.011496\\
0.466	0.0114375\\
0.468	0.0113798\\
0.47	0.0113227\\
0.472	0.0112663\\
0.474	0.0112106\\
0.476	0.0111554\\
0.478	0.0111008\\
0.48	0.0110467\\
0.482	0.0109931\\
0.484	0.0109399\\
0.486	0.0108871\\
0.488	0.0108346\\
0.49	0.0107824\\
0.492	0.0107306\\
0.494	0.0106789\\
0.496	0.0106274\\
0.498	0.0105761\\
0.5	0.0105249\\
0.502	0.0104738\\
0.504	0.0104227\\
0.506	0.0103716\\
0.508	0.0103206\\
0.51	0.0102694\\
0.512	0.0102182\\
0.514	0.010167\\
0.516	0.0101156\\
0.518	0.010064\\
0.52	0.0100124\\
0.522	0.00996051\\
0.524	0.00990849\\
0.526	0.00985629\\
0.528	0.00980389\\
0.53	0.0097513\\
0.532	0.0096985\\
0.534	0.00964551\\
0.536	0.00959231\\
0.538	0.00953892\\
0.54	0.00948534\\
0.542	0.00943158\\
0.544	0.00937764\\
0.546	0.00932354\\
0.548	0.0092693\\
0.55	0.00921492\\
0.552	0.00916042\\
0.554	0.00910583\\
0.556	0.00905114\\
0.558	0.00899639\\
0.56	0.0089416\\
0.562	0.00888678\\
0.564	0.00883196\\
0.566	0.00877715\\
0.568	0.00872239\\
0.57	0.00866768\\
0.572	0.00861305\\
0.574	0.00855853\\
0.576	0.00850413\\
0.578	0.00844987\\
0.58	0.00839578\\
0.582	0.00834187\\
0.584	0.00828817\\
0.586	0.00823469\\
0.588	0.00818146\\
0.59	0.00812848\\
0.592	0.00807577\\
0.594	0.00802335\\
0.596	0.00797124\\
0.598	0.00791944\\
0.6	0.00786797\\
0.602	0.00781684\\
0.604	0.00776605\\
0.606	0.00771563\\
0.608	0.00766558\\
0.61	0.0076159\\
0.612	0.00756659\\
0.614	0.00751768\\
0.616	0.00746915\\
0.618	0.00742101\\
0.62	0.00737327\\
0.622	0.00732592\\
0.624	0.00727897\\
0.626	0.00723241\\
0.628	0.00718624\\
0.63	0.00714047\\
0.632	0.00709509\\
0.634	0.0070501\\
0.636	0.00700549\\
0.638	0.00696126\\
0.64	0.0069174\\
0.642	0.00687392\\
0.644	0.0068308\\
0.646	0.00678804\\
0.648	0.00674563\\
0.65	0.00670357\\
0.652	0.00666186\\
0.654	0.00662048\\
0.656	0.00657943\\
0.658	0.00653871\\
0.66	0.0064983\\
0.662	0.0064582\\
0.664	0.00641841\\
0.666	0.00637892\\
0.668	0.00633972\\
0.67	0.00630081\\
0.672	0.00626217\\
0.674	0.00622382\\
0.676	0.00618573\\
0.678	0.0061479\\
0.68	0.00611033\\
0.682	0.00607302\\
0.684	0.00603595\\
0.686	0.00599913\\
0.688	0.00596254\\
0.69	0.00592619\\
0.692	0.00589006\\
0.694	0.00585416\\
0.696	0.00581847\\
0.698	0.005783\\
0.7	0.00574775\\
0.702	0.00571269\\
0.704	0.00567785\\
0.706	0.0056432\\
0.708	0.00560874\\
0.71	0.00557448\\
0.712	0.0055404\\
0.714	0.00550651\\
0.716	0.0054728\\
0.718	0.00543927\\
0.72	0.00540591\\
0.722	0.00537272\\
0.724	0.0053397\\
0.726	0.00530684\\
0.728	0.00527414\\
0.73	0.00524161\\
0.732	0.00520923\\
0.734	0.005177\\
0.736	0.00514492\\
0.738	0.00511298\\
0.74	0.0050812\\
0.742	0.00504955\\
0.744	0.00501804\\
0.746	0.00498668\\
0.748	0.00495544\\
0.75	0.00492434\\
0.752	0.00489338\\
0.754	0.00486254\\
0.756	0.00483183\\
0.758	0.00480124\\
0.76	0.00477079\\
0.762	0.00474045\\
0.764	0.00471024\\
0.766	0.00468015\\
0.768	0.00465019\\
0.77	0.00462034\\
0.772	0.00459062\\
0.774	0.00456102\\
0.776	0.00453153\\
0.778	0.00450217\\
0.78	0.00447294\\
0.782	0.00444382\\
0.784	0.00441483\\
0.786	0.00438596\\
0.788	0.00435721\\
0.79	0.00432859\\
0.792	0.0043001\\
0.794	0.00427174\\
0.796	0.00424351\\
0.798	0.00421541\\
0.8	0.00418744\\
0.802	0.00415962\\
0.804	0.00413192\\
0.806	0.00410437\\
0.808	0.00407696\\
0.81	0.0040497\\
0.812	0.00402258\\
0.814	0.0039956\\
0.816	0.00396878\\
0.818	0.00394211\\
0.82	0.0039156\\
0.822	0.00388924\\
0.824	0.00386305\\
0.826	0.00383701\\
0.828	0.00381113\\
0.83	0.00378543\\
0.832	0.00375988\\
0.834	0.00373451\\
0.836	0.0037093\\
0.838	0.00368427\\
0.84	0.0036594\\
0.842	0.00363472\\
0.844	0.0036102\\
0.846	0.00358586\\
0.848	0.0035617\\
0.85	0.00353772\\
0.852	0.00351391\\
0.854	0.00349028\\
0.856	0.00346682\\
0.858	0.00344355\\
0.86	0.00342045\\
0.862	0.00339753\\
0.864	0.00337479\\
0.866	0.00335222\\
0.868	0.00332983\\
0.87	0.00330761\\
0.872	0.00328557\\
0.874	0.00326369\\
0.876	0.00324199\\
0.878	0.00322046\\
0.88	0.0031991\\
0.882	0.00317791\\
0.884	0.00315687\\
0.886	0.00313601\\
0.888	0.0031153\\
0.89	0.00309475\\
0.892	0.00307436\\
0.894	0.00305413\\
0.896	0.00303405\\
0.898	0.00301412\\
0.9	0.00299434\\
0.902	0.00297471\\
0.904	0.00295522\\
0.906	0.00293587\\
0.908	0.00291667\\
0.91	0.0028976\\
0.912	0.00287867\\
0.914	0.00285987\\
0.916	0.0028412\\
0.918	0.00282267\\
0.92	0.00280426\\
0.922	0.00278597\\
0.924	0.00276781\\
0.926	0.00274978\\
0.928	0.00273186\\
0.93	0.00271406\\
0.932	0.00269637\\
0.934	0.00267881\\
0.936	0.00266135\\
0.938	0.002644\\
0.94	0.00262677\\
0.942	0.00260964\\
0.944	0.00259262\\
0.946	0.00257571\\
0.948	0.0025589\\
0.95	0.00254219\\
0.952	0.00252558\\
0.954	0.00250908\\
0.956	0.00249267\\
0.958	0.00247637\\
0.96	0.00246016\\
0.962	0.00244405\\
0.964	0.00242803\\
0.966	0.00241211\\
0.968	0.00239628\\
0.97	0.00238055\\
0.972	0.00236491\\
0.974	0.00234936\\
0.976	0.0023339\\
0.978	0.00231854\\
0.98	0.00230326\\
0.982	0.00228807\\
0.984	0.00227298\\
0.986	0.00225797\\
0.988	0.00224305\\
0.99	0.00222822\\
0.992	0.00221347\\
0.994	0.00219882\\
0.996	0.00218424\\
0.998	0.00216976\\
1	0.00215536\\
1.002	0.00214105\\
1.004	0.00212682\\
1.006	0.00211267\\
1.008	0.00209861\\
1.01	0.00208463\\
1.012	0.00207074\\
1.014	0.00205693\\
1.016	0.0020432\\
1.018	0.00202955\\
1.02	0.00201599\\
1.022	0.00200251\\
1.024	0.00198911\\
1.026	0.00197579\\
1.028	0.00196255\\
1.03	0.00194939\\
1.032	0.00193632\\
1.034	0.00192332\\
1.036	0.0019104\\
1.038	0.00189756\\
1.04	0.0018848\\
1.042	0.00187212\\
1.044	0.00185952\\
1.046	0.001847\\
1.048	0.00183456\\
1.05	0.00182219\\
1.052	0.0018099\\
1.054	0.00179769\\
1.056	0.00178556\\
1.058	0.0017735\\
1.06	0.00176152\\
1.062	0.00174962\\
1.064	0.00173779\\
1.066	0.00172604\\
1.068	0.00171437\\
1.07	0.00170277\\
1.072	0.00169125\\
1.074	0.00167981\\
1.076	0.00166844\\
1.078	0.00165714\\
1.08	0.00164592\\
1.082	0.00163478\\
1.084	0.00162371\\
1.086	0.00161271\\
1.088	0.00160179\\
1.09	0.00159095\\
1.092	0.00158017\\
1.094	0.00156948\\
1.096	0.00155885\\
1.098	0.0015483\\
1.1	0.00153782\\
1.102	0.00152741\\
1.104	0.00151708\\
1.106	0.00150682\\
1.108	0.00149663\\
1.11	0.00148651\\
1.112	0.00147646\\
1.114	0.00146649\\
1.116	0.00145658\\
1.118	0.00144674\\
1.12	0.00143698\\
1.122	0.00142728\\
1.124	0.00141765\\
1.126	0.00140809\\
1.128	0.0013986\\
1.13	0.00138918\\
1.132	0.00137982\\
1.134	0.00137053\\
1.136	0.0013613\\
1.138	0.00135214\\
1.14	0.00134305\\
1.142	0.00133402\\
1.144	0.00132505\\
1.146	0.00131615\\
1.148	0.00130731\\
1.15	0.00129853\\
1.152	0.00128981\\
1.154	0.00128116\\
1.156	0.00127256\\
1.158	0.00126403\\
1.16	0.00125555\\
1.162	0.00124713\\
1.164	0.00123877\\
1.166	0.00123047\\
1.168	0.00122223\\
1.17	0.00121404\\
1.172	0.00120591\\
1.174	0.00119783\\
1.176	0.00118981\\
1.178	0.00118184\\
1.18	0.00117393\\
1.182	0.00116606\\
1.184	0.00115826\\
1.186	0.0011505\\
1.188	0.00114279\\
1.19	0.00113514\\
1.192	0.00112754\\
1.194	0.00111999\\
1.196	0.00111248\\
1.198	0.00110503\\
1.2	0.00109762\\
1.202	0.00109027\\
1.204	0.00108296\\
1.206	0.0010757\\
1.208	0.00106849\\
1.21	0.00106132\\
1.212	0.0010542\\
1.214	0.00104713\\
1.216	0.0010401\\
1.218	0.00103312\\
1.22	0.00102618\\
1.222	0.00101929\\
1.224	0.00101244\\
1.226	0.00100563\\
1.228	0.000998873\\
1.23	0.000992157\\
1.232	0.000985484\\
1.234	0.000978854\\
1.236	0.000972267\\
1.238	0.000965722\\
1.24	0.00095922\\
1.242	0.00095276\\
1.244	0.000946342\\
1.246	0.000939965\\
1.248	0.00093363\\
1.25	0.000927336\\
1.252	0.000921082\\
1.254	0.00091487\\
1.256	0.000908698\\
1.258	0.000902566\\
1.26	0.000896474\\
1.262	0.000890422\\
1.264	0.00088441\\
1.266	0.000878437\\
1.268	0.000872503\\
1.27	0.000866608\\
1.272	0.000860752\\
1.274	0.000854935\\
1.276	0.000849156\\
1.278	0.000843414\\
1.28	0.000837711\\
1.282	0.000832045\\
1.284	0.000826417\\
1.286	0.000820826\\
1.288	0.000815272\\
1.29	0.000809755\\
1.292	0.000804275\\
1.294	0.000798831\\
1.296	0.000793423\\
1.298	0.000788051\\
1.3	0.000782715\\
1.302	0.000777415\\
1.304	0.00077215\\
1.306	0.00076692\\
1.308	0.000761724\\
1.31	0.000756564\\
1.312	0.000751438\\
1.314	0.000746347\\
1.316	0.000741289\\
1.318	0.000736266\\
1.32	0.000731276\\
1.322	0.00072632\\
1.324	0.000721397\\
1.326	0.000716508\\
1.328	0.000711651\\
1.33	0.000706827\\
1.332	0.000702035\\
1.334	0.000697276\\
1.336	0.000692549\\
1.338	0.000687854\\
1.34	0.000683191\\
1.342	0.00067856\\
1.344	0.00067396\\
1.346	0.000669391\\
1.348	0.000664853\\
1.35	0.000660526\\
1.352	0.00065639\\
1.354	0.000652281\\
1.356	0.000648196\\
1.358	0.000644136\\
1.36	0.000640101\\
1.362	0.000636091\\
1.364	0.000632105\\
1.366	0.000628144\\
1.368	0.000624208\\
1.37	0.000620296\\
1.372	0.000616408\\
1.374	0.000612544\\
1.376	0.000608703\\
1.378	0.000604887\\
1.38	0.000601094\\
1.382	0.000597325\\
1.384	0.00059358\\
1.386	0.000589857\\
1.388	0.000586158\\
1.39	0.000582482\\
1.392	0.000578829\\
1.394	0.000575199\\
1.396	0.000571591\\
1.398	0.000568006\\
1.4	0.000564443\\
1.402	0.000560903\\
1.404	0.000557385\\
1.406	0.000553889\\
1.408	0.000550414\\
1.41	0.000546962\\
1.412	0.00054353\\
1.414	0.000540121\\
1.416	0.000536733\\
1.418	0.000533366\\
1.42	0.000530019\\
1.422	0.000526694\\
1.424	0.00052339\\
1.426	0.000520106\\
1.428	0.000516843\\
1.43	0.0005136\\
1.432	0.000510377\\
1.434	0.000507174\\
1.436	0.000503991\\
1.438	0.000500828\\
1.44	0.000497684\\
1.442	0.00049456\\
1.444	0.000491455\\
1.446	0.00048837\\
1.448	0.000485303\\
1.45	0.000482255\\
1.452	0.000479226\\
1.454	0.000476216\\
1.456	0.000473224\\
1.458	0.000470251\\
1.46	0.000467295\\
1.462	0.000464358\\
1.464	0.000461439\\
1.466	0.000458538\\
1.468	0.000455654\\
1.47	0.000452788\\
1.472	0.000449939\\
1.474	0.000447108\\
1.476	0.000444294\\
1.478	0.000441497\\
1.48	0.000438717\\
1.482	0.000435954\\
1.484	0.000433208\\
1.486	0.000430478\\
1.488	0.000427765\\
1.49	0.000425068\\
1.492	0.000422388\\
1.494	0.000419724\\
1.496	0.000417076\\
1.498	0.000414444\\
1.5	0.000411828\\
1.502	0.000409228\\
1.504	0.000406643\\
1.506	0.000404075\\
1.508	0.000401521\\
1.51	0.000398984\\
1.512	0.000396461\\
1.514	0.000393954\\
1.516	0.000391463\\
1.518	0.000388986\\
1.52	0.000386524\\
1.522	0.000384078\\
1.524	0.000381646\\
1.526	0.000379229\\
1.528	0.000376827\\
1.53	0.000374439\\
1.532	0.000372066\\
1.534	0.000369708\\
1.536	0.000367364\\
1.538	0.000365034\\
1.54	0.000362719\\
1.542	0.000360417\\
1.544	0.00035813\\
1.546	0.000355857\\
1.548	0.000353598\\
1.55	0.000351353\\
1.552	0.000349121\\
1.554	0.000346904\\
1.556	0.0003447\\
1.558	0.00034251\\
1.56	0.000340333\\
1.562	0.000338169\\
1.564	0.00033602\\
1.566	0.000333883\\
1.568	0.00033176\\
1.57	0.00032965\\
1.572	0.000327553\\
1.574	0.000325469\\
1.576	0.000323398\\
1.578	0.00032134\\
1.58	0.000319295\\
1.582	0.000317262\\
1.584	0.000315242\\
1.586	0.000313235\\
1.588	0.000311241\\
1.59	0.000309259\\
1.592	0.000307289\\
1.594	0.000305332\\
1.596	0.000303387\\
1.598	0.000301455\\
1.6	0.000299534\\
1.602	0.000297626\\
1.604	0.000295729\\
1.606	0.000293845\\
1.608	0.000291972\\
1.61	0.000290111\\
1.612	0.000288262\\
1.614	0.000286425\\
1.616	0.000284599\\
1.618	0.000282785\\
1.62	0.000280982\\
1.622	0.000279191\\
1.624	0.000277411\\
1.626	0.000275643\\
1.628	0.000273885\\
1.63	0.000272139\\
1.632	0.000270403\\
1.634	0.000268679\\
1.636	0.000266966\\
1.638	0.000265263\\
1.64	0.000263572\\
1.642	0.000261891\\
1.644	0.000260221\\
1.646	0.000258561\\
1.648	0.000256912\\
1.65	0.000255273\\
1.652	0.000253645\\
1.654	0.000252027\\
1.656	0.000250419\\
1.658	0.000248822\\
1.66	0.000247235\\
1.662	0.000245657\\
1.664	0.00024409\\
1.666	0.000242533\\
1.668	0.000240985\\
1.67	0.000239448\\
1.672	0.00023792\\
1.674	0.000236402\\
1.676	0.000234893\\
1.678	0.000233394\\
1.68	0.000231905\\
1.682	0.000230424\\
1.684	0.000228954\\
1.686	0.000227492\\
1.688	0.00022604\\
1.69	0.000224597\\
1.692	0.000223163\\
1.694	0.000221738\\
1.696	0.000220322\\
1.698	0.000218915\\
1.7	0.000217516\\
1.702	0.000216127\\
1.704	0.000214746\\
1.706	0.000213374\\
1.708	0.000212011\\
1.71	0.000210656\\
1.712	0.00020931\\
1.714	0.000207972\\
1.716	0.000206642\\
1.718	0.000205321\\
1.72	0.000204008\\
1.722	0.000202704\\
1.724	0.000201407\\
1.726	0.000200119\\
1.728	0.000198838\\
1.73	0.000197566\\
1.732	0.000196302\\
1.734	0.000195045\\
1.736	0.000193796\\
1.738	0.000192555\\
1.74	0.000191322\\
1.742	0.000190097\\
1.744	0.000188879\\
1.746	0.000187669\\
1.748	0.000186466\\
1.75	0.000185271\\
1.752	0.000184084\\
1.754	0.000182903\\
1.756	0.000181731\\
1.758	0.000180565\\
1.76	0.000179407\\
1.762	0.000178256\\
1.764	0.000177112\\
1.766	0.000175975\\
1.768	0.000174846\\
1.77	0.000173723\\
1.772	0.000172608\\
1.774	0.0001715\\
1.776	0.000170398\\
1.778	0.000169303\\
1.78	0.000168216\\
1.782	0.000167135\\
1.784	0.000166061\\
1.786	0.000164994\\
1.788	0.000163933\\
1.79	0.000162879\\
1.792	0.000161832\\
1.794	0.000160791\\
1.796	0.000159757\\
1.798	0.000158729\\
1.8	0.000157708\\
1.802	0.000156694\\
1.804	0.000155685\\
1.806	0.000154684\\
1.808	0.000153688\\
1.81	0.000152699\\
1.812	0.000151716\\
1.814	0.000150739\\
1.816	0.000149769\\
1.818	0.000148805\\
1.82	0.000147847\\
1.822	0.000146895\\
1.824	0.000145949\\
1.826	0.000145009\\
1.828	0.000144075\\
1.83	0.000143147\\
1.832	0.000142225\\
1.834	0.000141309\\
1.836	0.000140398\\
1.838	0.000139494\\
1.84	0.000138595\\
1.842	0.000137702\\
1.844	0.000136815\\
1.846	0.000135934\\
1.848	0.000135058\\
1.85	0.000134187\\
1.852	0.000133323\\
1.854	0.000132464\\
1.856	0.00013161\\
1.858	0.000130762\\
1.86	0.000129919\\
1.862	0.000129082\\
1.864	0.00012825\\
1.866	0.000127423\\
1.868	0.000126602\\
1.87	0.000125786\\
1.872	0.000124975\\
1.874	0.000124169\\
1.876	0.000123369\\
1.878	0.000122574\\
1.88	0.000121784\\
1.882	0.000120998\\
1.884	0.000120218\\
1.886	0.000119443\\
1.888	0.000118673\\
1.89	0.000117908\\
1.892	0.000117148\\
1.894	0.000116392\\
1.896	0.000115642\\
1.898	0.000114896\\
1.9	0.000114155\\
1.902	0.000113419\\
1.904	0.000112687\\
1.906	0.000111961\\
1.908	0.000111238\\
1.91	0.000110521\\
1.912	0.000109808\\
1.914	0.000109099\\
1.916	0.000108396\\
1.918	0.000107696\\
1.92	0.000107001\\
1.922	0.000106311\\
1.924	0.000105625\\
1.926	0.000104943\\
1.928	0.000104266\\
1.93	0.000103593\\
1.932	0.000102924\\
1.934	0.00010226\\
1.936	0.000101599\\
1.938	0.000100943\\
1.94	0.000100292\\
1.942	9.96439e-05\\
1.944	9.90004e-05\\
1.946	9.83609e-05\\
1.948	9.77256e-05\\
1.95	9.70943e-05\\
1.952	9.6467e-05\\
1.954	9.58436e-05\\
1.956	9.52243e-05\\
1.958	9.46089e-05\\
1.96	9.39974e-05\\
1.962	9.33898e-05\\
1.964	9.27861e-05\\
1.966	9.21862e-05\\
1.968	9.15902e-05\\
1.97	9.09979e-05\\
1.972	9.04094e-05\\
1.974	8.98247e-05\\
1.976	8.92436e-05\\
1.978	8.86663e-05\\
1.98	8.80926e-05\\
1.982	8.75226e-05\\
1.984	8.69562e-05\\
1.986	8.63935e-05\\
1.988	8.58343e-05\\
1.99	8.52787e-05\\
1.992	8.47266e-05\\
1.994	8.4178e-05\\
1.996	8.3633e-05\\
1.998	8.30914e-05\\
2	8.25532e-05\\
};
\addplot [color=mycolor7, line width=1.0pt, forget plot]
  table[row sep=crcr]{%
-0.024	0\\
-0.022	0\\
-0.02	0\\
-0.018	0\\
-0.016	0\\
-0.014	0\\
-0.012	0\\
-0.01	0\\
-0.008	0\\
-0.006	0\\
-0.004	0\\
-0.002	0\\
0	0\\
0.002	0.0516749\\
0.004	0.0926049\\
0.006	0.12345\\
0.008	0.145087\\
0.01	0.158549\\
0.012	0.164961\\
0.014	0.165487\\
0.016	0.161282\\
0.018	0.154844\\
0.02	0.14652\\
0.022	0.136117\\
0.024	0.12436\\
0.026	0.12482\\
0.028	0.130314\\
0.03	0.134985\\
0.032	0.138808\\
0.034	0.141769\\
0.036	0.14386\\
0.038	0.145084\\
0.04	0.145452\\
0.042	0.144986\\
0.044	0.143713\\
0.046	0.141669\\
0.048	0.138895\\
0.05	0.135441\\
0.052	0.131359\\
0.054	0.126706\\
0.056	0.121543\\
0.058	0.115932\\
0.06	0.10994\\
0.062	0.10363\\
0.064	0.097126\\
0.066	0.098084\\
0.068	0.0988132\\
0.07	0.0993164\\
0.072	0.0995969\\
0.074	0.0996587\\
0.076	0.0995061\\
0.078	0.099144\\
0.08	0.0985777\\
0.082	0.097813\\
0.084	0.0968561\\
0.086	0.0957134\\
0.088	0.0943921\\
0.09	0.0928993\\
0.092	0.0912428\\
0.094	0.0894306\\
0.096	0.0874711\\
0.098	0.0853728\\
0.1	0.0831447\\
0.102	0.0807961\\
0.104	0.0783362\\
0.106	0.0757746\\
0.108	0.0731213\\
0.11	0.0703859\\
0.112	0.0703862\\
0.114	0.0706864\\
0.116	0.0708794\\
0.118	0.0709654\\
0.12	0.0709453\\
0.122	0.0708893\\
0.124	0.0708244\\
0.126	0.0706429\\
0.128	0.0703465\\
0.13	0.0699373\\
0.132	0.0694178\\
0.134	0.0687908\\
0.136	0.0680596\\
0.138	0.0672276\\
0.14	0.0662987\\
0.142	0.0652771\\
0.144	0.0641672\\
0.146	0.0629735\\
0.148	0.061701\\
0.15	0.0603548\\
0.152	0.05894\\
0.154	0.057462\\
0.156	0.0559264\\
0.158	0.0543388\\
0.16	0.0527048\\
0.162	0.0510303\\
0.164	0.0493208\\
0.166	0.0475823\\
0.168	0.0458204\\
0.17	0.0440408\\
0.172	0.0422492\\
0.174	0.0412825\\
0.176	0.0413574\\
0.178	0.0412578\\
0.18	0.0409906\\
0.182	0.0405637\\
0.184	0.0399863\\
0.186	0.0392684\\
0.188	0.0384208\\
0.19	0.037455\\
0.192	0.0363831\\
0.194	0.0352173\\
0.196	0.0339703\\
0.198	0.032655\\
0.2	0.0312842\\
0.202	0.0300735\\
0.204	0.0298421\\
0.206	0.0296115\\
0.208	0.0293816\\
0.21	0.0291528\\
0.212	0.029199\\
0.214	0.0298298\\
0.216	0.0303717\\
0.218	0.0308221\\
0.22	0.0311794\\
0.222	0.0314425\\
0.224	0.031611\\
0.226	0.0316853\\
0.228	0.0316661\\
0.23	0.0315551\\
0.232	0.0313542\\
0.234	0.0310662\\
0.236	0.0306941\\
0.238	0.0302416\\
0.24	0.0297128\\
0.242	0.029112\\
0.244	0.0284442\\
0.246	0.0277144\\
0.248	0.0269282\\
0.25	0.026091\\
0.252	0.0252088\\
0.254	0.0248851\\
0.256	0.0247031\\
0.258	0.0245202\\
0.26	0.0243364\\
0.262	0.0241517\\
0.264	0.0239658\\
0.266	0.0241791\\
0.268	0.0247989\\
0.27	0.0253597\\
0.272	0.0258604\\
0.274	0.0263002\\
0.276	0.0266787\\
0.278	0.0269954\\
0.28	0.0272504\\
0.282	0.0274439\\
0.284	0.0275762\\
0.286	0.0276479\\
0.288	0.0276599\\
0.29	0.0276132\\
0.292	0.027509\\
0.294	0.0273487\\
0.296	0.0271338\\
0.298	0.0268661\\
0.3	0.0265473\\
0.302	0.0261796\\
0.304	0.0257651\\
0.306	0.0253059\\
0.308	0.0248045\\
0.31	0.0242632\\
0.312	0.0236847\\
0.314	0.0230716\\
0.316	0.0224265\\
0.318	0.021752\\
0.32	0.0210511\\
0.322	0.0203264\\
0.324	0.0195809\\
0.326	0.0188172\\
0.328	0.0180811\\
0.33	0.0179783\\
0.332	0.0178495\\
0.334	0.0176956\\
0.336	0.0175173\\
0.338	0.0173157\\
0.34	0.0170916\\
0.342	0.0168461\\
0.344	0.0169886\\
0.346	0.0171137\\
0.348	0.0172113\\
0.35	0.0172815\\
0.352	0.0173246\\
0.354	0.0173408\\
0.356	0.0173304\\
0.358	0.017294\\
0.36	0.0172319\\
0.362	0.0171449\\
0.364	0.0170335\\
0.366	0.0168985\\
0.368	0.0167407\\
0.37	0.0165607\\
0.372	0.0163596\\
0.374	0.0161381\\
0.376	0.0158974\\
0.378	0.0156383\\
0.38	0.015362\\
0.382	0.0150693\\
0.384	0.0147615\\
0.386	0.0144397\\
0.388	0.0142363\\
0.39	0.0141664\\
0.392	0.0140963\\
0.394	0.0140263\\
0.396	0.0139562\\
0.398	0.0138861\\
0.4	0.0138161\\
0.402	0.0137461\\
0.404	0.0136763\\
0.406	0.0136065\\
0.408	0.013537\\
0.41	0.0134676\\
0.412	0.0133984\\
0.414	0.0133294\\
0.416	0.0132607\\
0.418	0.0131923\\
0.42	0.0131242\\
0.422	0.0130563\\
0.424	0.0129888\\
0.426	0.0129216\\
0.428	0.0128547\\
0.43	0.0127882\\
0.432	0.0127221\\
0.434	0.0126563\\
0.436	0.012591\\
0.438	0.0125259\\
0.44	0.0124613\\
0.442	0.0123971\\
0.444	0.0123333\\
0.446	0.0122698\\
0.448	0.0122068\\
0.45	0.0121441\\
0.452	0.0120818\\
0.454	0.0120199\\
0.456	0.0119584\\
0.458	0.0118972\\
0.46	0.0118364\\
0.462	0.0117759\\
0.464	0.0117158\\
0.466	0.011656\\
0.468	0.0115965\\
0.47	0.0115373\\
0.472	0.0114784\\
0.474	0.0114198\\
0.476	0.0113615\\
0.478	0.0113035\\
0.48	0.0112457\\
0.482	0.0111881\\
0.484	0.0111308\\
0.486	0.0110737\\
0.488	0.0110168\\
0.49	0.01096\\
0.492	0.0109035\\
0.494	0.0108471\\
0.496	0.0107909\\
0.498	0.0107348\\
0.5	0.0106789\\
0.502	0.0106231\\
0.504	0.0105674\\
0.506	0.0105118\\
0.508	0.0104564\\
0.51	0.010401\\
0.512	0.0103457\\
0.514	0.0102905\\
0.516	0.0102354\\
0.518	0.0101803\\
0.52	0.0101253\\
0.522	0.0100704\\
0.524	0.0100155\\
0.526	0.00996071\\
0.528	0.00990596\\
0.53	0.00985127\\
0.532	0.00979663\\
0.534	0.00974205\\
0.536	0.00968754\\
0.538	0.00963308\\
0.54	0.00957868\\
0.542	0.00952435\\
0.544	0.00947009\\
0.546	0.0094159\\
0.548	0.00936179\\
0.55	0.00930776\\
0.552	0.00925382\\
0.554	0.00919997\\
0.556	0.00914621\\
0.558	0.00909257\\
0.56	0.00903903\\
0.562	0.00898561\\
0.564	0.00893232\\
0.566	0.00887916\\
0.568	0.00882614\\
0.57	0.00877327\\
0.572	0.00872056\\
0.574	0.00866801\\
0.576	0.00861563\\
0.578	0.00856343\\
0.58	0.00851141\\
0.582	0.00845959\\
0.584	0.00840797\\
0.586	0.00835657\\
0.588	0.00830538\\
0.59	0.00825441\\
0.592	0.00820368\\
0.594	0.00815318\\
0.596	0.00810293\\
0.598	0.00805293\\
0.6	0.00800319\\
0.602	0.00795371\\
0.604	0.00790449\\
0.606	0.00785556\\
0.608	0.0078069\\
0.61	0.00775852\\
0.612	0.00771043\\
0.614	0.00766263\\
0.616	0.00761512\\
0.618	0.00756791\\
0.62	0.007521\\
0.622	0.00747438\\
0.624	0.00742808\\
0.626	0.00738207\\
0.628	0.00733637\\
0.63	0.00729098\\
0.632	0.00724589\\
0.634	0.00720111\\
0.636	0.00715664\\
0.638	0.00711247\\
0.64	0.0070686\\
0.642	0.00702504\\
0.644	0.00698178\\
0.646	0.00693882\\
0.648	0.00689615\\
0.65	0.00685378\\
0.652	0.0068117\\
0.654	0.00676991\\
0.656	0.0067284\\
0.658	0.00668718\\
0.66	0.00664624\\
0.662	0.00660557\\
0.664	0.00656517\\
0.666	0.00652504\\
0.668	0.00648517\\
0.67	0.00644557\\
0.672	0.00640622\\
0.674	0.00636712\\
0.676	0.00632827\\
0.678	0.00628966\\
0.68	0.00625129\\
0.682	0.00621316\\
0.684	0.00617525\\
0.686	0.00613758\\
0.688	0.00610012\\
0.69	0.00606289\\
0.692	0.00602587\\
0.694	0.00598906\\
0.696	0.00595246\\
0.698	0.00591606\\
0.7	0.00587986\\
0.702	0.00584386\\
0.704	0.00580805\\
0.706	0.00577243\\
0.708	0.005737\\
0.71	0.00570176\\
0.712	0.00566669\\
0.714	0.00563181\\
0.716	0.0055971\\
0.718	0.00556257\\
0.72	0.00552821\\
0.722	0.00549401\\
0.724	0.00545999\\
0.726	0.00542613\\
0.728	0.00539244\\
0.73	0.00535891\\
0.732	0.00532555\\
0.734	0.00529234\\
0.736	0.0052593\\
0.738	0.00522641\\
0.74	0.00519368\\
0.742	0.00516111\\
0.744	0.00512869\\
0.746	0.00509643\\
0.748	0.00506432\\
0.75	0.00503237\\
0.752	0.00500058\\
0.754	0.00496893\\
0.756	0.00493745\\
0.758	0.00490611\\
0.76	0.00487494\\
0.762	0.00484391\\
0.764	0.00481304\\
0.766	0.00478232\\
0.768	0.00475176\\
0.77	0.00472136\\
0.772	0.0046911\\
0.774	0.00466101\\
0.776	0.00463107\\
0.778	0.00460128\\
0.78	0.00457165\\
0.782	0.00454218\\
0.784	0.00451287\\
0.786	0.00448371\\
0.788	0.00445471\\
0.79	0.00442587\\
0.792	0.00439718\\
0.794	0.00436866\\
0.796	0.00434029\\
0.798	0.00431209\\
0.8	0.00428404\\
0.802	0.00425616\\
0.804	0.00422843\\
0.806	0.00420087\\
0.808	0.00417346\\
0.81	0.00414622\\
0.812	0.00411913\\
0.814	0.00409221\\
0.816	0.00406545\\
0.818	0.00403885\\
0.82	0.00401242\\
0.822	0.00398614\\
0.824	0.00396003\\
0.826	0.00393408\\
0.828	0.00390828\\
0.83	0.00388265\\
0.832	0.00385719\\
0.834	0.00383188\\
0.836	0.00380673\\
0.838	0.00378175\\
0.84	0.00375692\\
0.842	0.00373226\\
0.844	0.00370775\\
0.846	0.0036834\\
0.848	0.00365922\\
0.85	0.00363519\\
0.852	0.00361132\\
0.854	0.0035876\\
0.856	0.00356405\\
0.858	0.00354065\\
0.86	0.0035174\\
0.862	0.00349431\\
0.864	0.00347138\\
0.866	0.0034486\\
0.868	0.00342597\\
0.87	0.0034035\\
0.872	0.00338117\\
0.874	0.003359\\
0.876	0.00333698\\
0.878	0.00331511\\
0.88	0.00329338\\
0.882	0.00327181\\
0.884	0.00325038\\
0.886	0.00322909\\
0.888	0.00320795\\
0.89	0.00318696\\
0.892	0.0031661\\
0.894	0.00314539\\
0.896	0.00312482\\
0.898	0.00310439\\
0.9	0.00308409\\
0.902	0.00306394\\
0.904	0.00304392\\
0.906	0.00302403\\
0.908	0.00300428\\
0.91	0.00298467\\
0.912	0.00296518\\
0.914	0.00294583\\
0.916	0.0029266\\
0.918	0.0029075\\
0.92	0.00288853\\
0.922	0.00286969\\
0.924	0.00285097\\
0.926	0.00283237\\
0.928	0.0028139\\
0.93	0.00279555\\
0.932	0.00277732\\
0.934	0.00275921\\
0.936	0.00274121\\
0.938	0.00272334\\
0.94	0.00270557\\
0.942	0.00268793\\
0.944	0.00267039\\
0.946	0.00265297\\
0.948	0.00263566\\
0.95	0.00261846\\
0.952	0.00260137\\
0.954	0.00258439\\
0.956	0.00256752\\
0.958	0.00255075\\
0.96	0.00253409\\
0.962	0.00251753\\
0.964	0.00250108\\
0.966	0.00248473\\
0.968	0.00246848\\
0.97	0.00245233\\
0.972	0.00243628\\
0.974	0.00242033\\
0.976	0.00240448\\
0.978	0.00238872\\
0.98	0.00237307\\
0.982	0.0023575\\
0.984	0.00234204\\
0.986	0.00232667\\
0.988	0.00231139\\
0.99	0.00229621\\
0.992	0.00228112\\
0.994	0.00226612\\
0.996	0.00225121\\
0.998	0.0022364\\
1	0.00222167\\
1.002	0.00220704\\
1.004	0.00219249\\
1.006	0.00217803\\
1.008	0.00216366\\
1.01	0.00214938\\
1.012	0.00213519\\
1.014	0.00212108\\
1.016	0.00210706\\
1.018	0.00209313\\
1.02	0.00207928\\
1.022	0.00206552\\
1.024	0.00205184\\
1.026	0.00203824\\
1.028	0.00202473\\
1.03	0.00201131\\
1.032	0.00199797\\
1.034	0.00198471\\
1.036	0.00197153\\
1.038	0.00195844\\
1.04	0.00194543\\
1.042	0.0019325\\
1.044	0.00191965\\
1.046	0.00190688\\
1.048	0.00189419\\
1.05	0.00188159\\
1.052	0.00186906\\
1.054	0.00185661\\
1.056	0.00184425\\
1.058	0.00183196\\
1.06	0.00181975\\
1.062	0.00180762\\
1.064	0.00179557\\
1.066	0.0017836\\
1.068	0.0017717\\
1.07	0.00175989\\
1.072	0.00174815\\
1.074	0.00173648\\
1.076	0.00172489\\
1.078	0.00171338\\
1.08	0.00170195\\
1.082	0.00169059\\
1.084	0.0016793\\
1.086	0.00166809\\
1.088	0.00165695\\
1.09	0.00164589\\
1.092	0.0016349\\
1.094	0.00162398\\
1.096	0.00161314\\
1.098	0.00160237\\
1.1	0.00159167\\
1.102	0.00158105\\
1.104	0.00157049\\
1.106	0.00156001\\
1.108	0.00154959\\
1.11	0.00153925\\
1.112	0.00152898\\
1.114	0.00151877\\
1.116	0.00150863\\
1.118	0.00149857\\
1.12	0.00148857\\
1.122	0.00147864\\
1.124	0.00146877\\
1.126	0.00145897\\
1.128	0.00144924\\
1.13	0.00143958\\
1.132	0.00142998\\
1.134	0.00142044\\
1.136	0.00141097\\
1.138	0.00140157\\
1.14	0.00139223\\
1.142	0.00138295\\
1.144	0.00137373\\
1.146	0.00136458\\
1.148	0.00135548\\
1.15	0.00134645\\
1.152	0.00133749\\
1.154	0.00132858\\
1.156	0.00131973\\
1.158	0.00131094\\
1.16	0.00130221\\
1.162	0.00129354\\
1.164	0.00128493\\
1.166	0.00127637\\
1.168	0.00126788\\
1.17	0.00125944\\
1.172	0.00125106\\
1.174	0.00124273\\
1.176	0.00123446\\
1.178	0.00122625\\
1.18	0.00121809\\
1.182	0.00120998\\
1.184	0.00120193\\
1.186	0.00119393\\
1.188	0.00118599\\
1.19	0.0011781\\
1.192	0.00117026\\
1.194	0.00116247\\
1.196	0.00115474\\
1.198	0.00114705\\
1.2	0.00113942\\
1.202	0.00113184\\
1.204	0.00112431\\
1.206	0.00111683\\
1.208	0.00110939\\
1.21	0.00110201\\
1.212	0.00109468\\
1.214	0.00108739\\
1.216	0.00108015\\
1.218	0.00107296\\
1.22	0.00106582\\
1.222	0.00105872\\
1.224	0.00105167\\
1.226	0.00104467\\
1.228	0.00103771\\
1.23	0.0010308\\
1.232	0.00102393\\
1.234	0.00101711\\
1.236	0.00101033\\
1.238	0.0010036\\
1.24	0.000996912\\
1.242	0.000990267\\
1.244	0.000983666\\
1.246	0.000977107\\
1.248	0.000970592\\
1.25	0.000964119\\
1.252	0.000957689\\
1.254	0.0009513\\
1.256	0.000944953\\
1.258	0.000938648\\
1.26	0.000932384\\
1.262	0.000926161\\
1.264	0.000919978\\
1.266	0.000913836\\
1.268	0.000907734\\
1.27	0.000901672\\
1.272	0.000895649\\
1.274	0.000889666\\
1.276	0.000883722\\
1.278	0.000877818\\
1.28	0.000871951\\
1.282	0.000866123\\
1.284	0.000860334\\
1.286	0.000854582\\
1.288	0.000848868\\
1.29	0.000843192\\
1.292	0.000837553\\
1.294	0.000831951\\
1.296	0.000826386\\
1.298	0.000820858\\
1.3	0.000815366\\
1.302	0.00080991\\
1.304	0.000804491\\
1.306	0.000799107\\
1.308	0.000793759\\
1.31	0.000788446\\
1.312	0.000783169\\
1.314	0.000777926\\
1.316	0.000772718\\
1.318	0.000767545\\
1.32	0.000762406\\
1.322	0.000757302\\
1.324	0.000752231\\
1.326	0.000747194\\
1.328	0.000742191\\
1.33	0.000737403\\
1.332	0.000732871\\
1.334	0.000728367\\
1.336	0.000723889\\
1.338	0.000719438\\
1.34	0.000715013\\
1.342	0.000710615\\
1.344	0.000706243\\
1.346	0.000701897\\
1.348	0.000697578\\
1.35	0.000693284\\
1.352	0.000689016\\
1.354	0.000684773\\
1.356	0.000680556\\
1.358	0.000676364\\
1.36	0.000672198\\
1.362	0.000668056\\
1.364	0.00066394\\
1.366	0.000659848\\
1.368	0.000655781\\
1.37	0.000651739\\
1.372	0.000647721\\
1.374	0.000643727\\
1.376	0.000639758\\
1.378	0.000635813\\
1.38	0.000631891\\
1.382	0.000627993\\
1.384	0.00062412\\
1.386	0.000620269\\
1.388	0.000616442\\
1.39	0.000612638\\
1.392	0.000608858\\
1.394	0.0006051\\
1.396	0.000601365\\
1.398	0.000597653\\
1.4	0.000593964\\
1.402	0.000590297\\
1.404	0.000586653\\
1.406	0.00058303\\
1.408	0.00057943\\
1.41	0.000575852\\
1.412	0.000572296\\
1.414	0.000568761\\
1.416	0.000565248\\
1.418	0.000561756\\
1.42	0.000558286\\
1.422	0.000554836\\
1.424	0.000551408\\
1.426	0.000548001\\
1.428	0.000544615\\
1.43	0.000541249\\
1.432	0.000537903\\
1.434	0.000534578\\
1.436	0.000531274\\
1.438	0.000527989\\
1.44	0.000524725\\
1.442	0.00052148\\
1.444	0.000518255\\
1.446	0.00051505\\
1.448	0.000511864\\
1.45	0.000508697\\
1.452	0.00050555\\
1.454	0.000502422\\
1.456	0.000499313\\
1.458	0.000496223\\
1.46	0.000493151\\
1.462	0.000490098\\
1.464	0.000487064\\
1.466	0.000484048\\
1.468	0.00048105\\
1.47	0.000478071\\
1.472	0.000475109\\
1.474	0.000472166\\
1.476	0.00046924\\
1.478	0.000466332\\
1.48	0.000463441\\
1.482	0.000460568\\
1.484	0.000457713\\
1.486	0.000454874\\
1.488	0.000452053\\
1.49	0.000449248\\
1.492	0.000446461\\
1.494	0.000443691\\
1.496	0.000440937\\
1.498	0.0004382\\
1.5	0.000435479\\
1.502	0.000432775\\
1.504	0.000430087\\
1.506	0.000427415\\
1.508	0.000424759\\
1.51	0.000422119\\
1.512	0.000419496\\
1.514	0.000416888\\
1.516	0.000414296\\
1.518	0.000411719\\
1.52	0.000409158\\
1.522	0.000406612\\
1.524	0.000404082\\
1.526	0.000401567\\
1.528	0.000399068\\
1.53	0.000396583\\
1.532	0.000394113\\
1.534	0.000391659\\
1.536	0.000389219\\
1.538	0.000386794\\
1.54	0.000384383\\
1.542	0.000381988\\
1.544	0.000379606\\
1.546	0.00037724\\
1.548	0.000374887\\
1.55	0.000372549\\
1.552	0.000370225\\
1.554	0.000367915\\
1.556	0.000365619\\
1.558	0.000363337\\
1.56	0.000361069\\
1.562	0.000358815\\
1.564	0.000356575\\
1.566	0.000354348\\
1.568	0.000352135\\
1.57	0.000349935\\
1.572	0.000347749\\
1.574	0.000345576\\
1.576	0.000343417\\
1.578	0.000341271\\
1.58	0.000339137\\
1.582	0.000337017\\
1.584	0.00033491\\
1.586	0.000332816\\
1.588	0.000330735\\
1.59	0.000328667\\
1.592	0.000326611\\
1.594	0.000324568\\
1.596	0.000322537\\
1.598	0.00032052\\
1.6	0.000318514\\
1.602	0.000316521\\
1.604	0.00031454\\
1.606	0.000312572\\
1.608	0.000310615\\
1.61	0.000308671\\
1.612	0.000306739\\
1.614	0.000304818\\
1.616	0.00030291\\
1.618	0.000301013\\
1.62	0.000299129\\
1.622	0.000297256\\
1.624	0.000295394\\
1.626	0.000293544\\
1.628	0.000291706\\
1.63	0.000289879\\
1.632	0.000288063\\
1.634	0.000286258\\
1.636	0.000284465\\
1.638	0.000282683\\
1.64	0.000280912\\
1.642	0.000279152\\
1.644	0.000277403\\
1.646	0.000275665\\
1.648	0.000273938\\
1.65	0.000272221\\
1.652	0.000270515\\
1.654	0.00026882\\
1.656	0.000267135\\
1.658	0.00026546\\
1.66	0.000263796\\
1.662	0.000262143\\
1.664	0.000260499\\
1.666	0.000258866\\
1.668	0.000257243\\
1.67	0.00025563\\
1.672	0.000254028\\
1.674	0.000252435\\
1.676	0.000250852\\
1.678	0.000249278\\
1.68	0.000247715\\
1.682	0.000246161\\
1.684	0.000244617\\
1.686	0.000243082\\
1.688	0.000241557\\
1.69	0.000240042\\
1.692	0.000238536\\
1.694	0.000237039\\
1.696	0.000235551\\
1.698	0.000234073\\
1.7	0.000232603\\
1.702	0.000231143\\
1.704	0.000229692\\
1.706	0.00022825\\
1.708	0.000226817\\
1.71	0.000225392\\
1.712	0.000223977\\
1.714	0.00022257\\
1.716	0.000221172\\
1.718	0.000219782\\
1.72	0.000218402\\
1.722	0.000217029\\
1.724	0.000215665\\
1.726	0.00021431\\
1.728	0.000212963\\
1.73	0.000211624\\
1.732	0.000210293\\
1.734	0.000208971\\
1.736	0.000207657\\
1.738	0.000206351\\
1.74	0.000205053\\
1.742	0.000203763\\
1.744	0.000202481\\
1.746	0.000201207\\
1.748	0.000199941\\
1.75	0.000198682\\
1.752	0.000197432\\
1.754	0.000196189\\
1.756	0.000194954\\
1.758	0.000193726\\
1.76	0.000192506\\
1.762	0.000191294\\
1.764	0.000190089\\
1.766	0.000188891\\
1.768	0.000187701\\
1.77	0.000186518\\
1.772	0.000185343\\
1.774	0.000184175\\
1.776	0.000183014\\
1.778	0.00018186\\
1.78	0.000180713\\
1.782	0.000179574\\
1.784	0.000178442\\
1.786	0.000177316\\
1.788	0.000176198\\
1.79	0.000175086\\
1.792	0.000173982\\
1.794	0.000172884\\
1.796	0.000171793\\
1.798	0.000170709\\
1.8	0.000169632\\
1.802	0.000168561\\
1.804	0.000167497\\
1.806	0.00016644\\
1.808	0.000165389\\
1.81	0.000164345\\
1.812	0.000163307\\
1.814	0.000162276\\
1.816	0.000161251\\
1.818	0.000160233\\
1.82	0.000159221\\
1.822	0.000158215\\
1.824	0.000157215\\
1.826	0.000156222\\
1.828	0.000155235\\
1.83	0.000154254\\
1.832	0.00015328\\
1.834	0.000152311\\
1.836	0.000151349\\
1.838	0.000150392\\
1.84	0.000149442\\
1.842	0.000148497\\
1.844	0.000147559\\
1.846	0.000146626\\
1.848	0.000145699\\
1.85	0.000144778\\
1.852	0.000143863\\
1.854	0.000142953\\
1.856	0.000142049\\
1.858	0.000141151\\
1.86	0.000140259\\
1.862	0.000139372\\
1.864	0.00013849\\
1.866	0.000137615\\
1.868	0.000136744\\
1.87	0.00013588\\
1.872	0.00013502\\
1.874	0.000134166\\
1.876	0.000133318\\
1.878	0.000132474\\
1.88	0.000131636\\
1.882	0.000130804\\
1.884	0.000129976\\
1.886	0.000129154\\
1.888	0.000128337\\
1.89	0.000127525\\
1.892	0.000126718\\
1.894	0.000125916\\
1.896	0.000125119\\
1.898	0.000124327\\
1.9	0.000123541\\
1.902	0.000122759\\
1.904	0.000121982\\
1.906	0.00012121\\
1.908	0.000120443\\
1.91	0.00011968\\
1.912	0.000118923\\
1.914	0.00011817\\
1.916	0.000117422\\
1.918	0.000116678\\
1.92	0.000115939\\
1.922	0.000115205\\
1.924	0.000114476\\
1.926	0.000113751\\
1.928	0.00011303\\
1.93	0.000112314\\
1.932	0.000111603\\
1.934	0.000110896\\
1.936	0.000110194\\
1.938	0.000109495\\
1.94	0.000108802\\
1.942	0.000108112\\
1.944	0.000107427\\
1.946	0.000106746\\
1.948	0.00010607\\
1.95	0.000105398\\
1.952	0.00010473\\
1.954	0.000104066\\
1.956	0.000103406\\
1.958	0.00010275\\
1.96	0.000102099\\
1.962	0.000101452\\
1.964	0.000100808\\
1.966	0.000100169\\
1.968	9.95336e-05\\
1.97	9.89023e-05\\
1.972	9.82749e-05\\
1.974	9.76514e-05\\
1.976	9.70319e-05\\
1.978	9.64163e-05\\
1.98	9.58045e-05\\
1.982	9.51965e-05\\
1.984	9.45924e-05\\
1.986	9.3992e-05\\
1.988	9.33954e-05\\
1.99	9.28026e-05\\
1.992	9.22135e-05\\
1.994	9.1628e-05\\
1.996	9.10463e-05\\
1.998	9.04681e-05\\
2	8.98937e-05\\
};
\addplot [color=mycolor8, line width=1.0pt, forget plot]
  table[row sep=crcr]{%
-0.024	0\\
-0.022	0\\
-0.02	0\\
-0.018	0\\
-0.016	0\\
-0.014	0\\
-0.012	0\\
-0.01	0\\
-0.008	0\\
-0.006	0\\
-0.004	0\\
-0.002	0\\
0	0\\
0.002	0.0516749\\
0.004	0.0926048\\
0.006	0.123449\\
0.008	0.145086\\
0.01	0.158547\\
0.012	0.164957\\
0.014	0.16548\\
0.016	0.161271\\
0.018	0.154828\\
0.02	0.146498\\
0.022	0.136087\\
0.024	0.12432\\
0.026	0.124839\\
0.028	0.130332\\
0.03	0.135\\
0.032	0.13882\\
0.034	0.141775\\
0.036	0.143859\\
0.038	0.145074\\
0.04	0.145432\\
0.042	0.144953\\
0.044	0.143664\\
0.046	0.141601\\
0.048	0.138807\\
0.05	0.135329\\
0.052	0.131221\\
0.054	0.12654\\
0.056	0.121346\\
0.058	0.115703\\
0.06	0.109675\\
0.062	0.103328\\
0.064	0.0970489\\
0.066	0.0979848\\
0.068	0.0986903\\
0.07	0.0991681\\
0.072	0.0994217\\
0.074	0.099455\\
0.076	0.0992726\\
0.078	0.0988796\\
0.08	0.0982813\\
0.082	0.0974837\\
0.084	0.0964933\\
0.086	0.0953166\\
0.088	0.093961\\
0.09	0.0924338\\
0.092	0.090743\\
0.094	0.0888968\\
0.096	0.0869037\\
0.098	0.0847727\\
0.1	0.0825127\\
0.102	0.0801332\\
0.104	0.0776437\\
0.106	0.075054\\
0.108	0.072374\\
0.11	0.0696137\\
0.112	0.069287\\
0.114	0.0695634\\
0.116	0.0697355\\
0.118	0.0698037\\
0.12	0.0697686\\
0.122	0.0696315\\
0.124	0.0693936\\
0.126	0.0691961\\
0.128	0.0689004\\
0.13	0.0684964\\
0.132	0.0679867\\
0.134	0.067374\\
0.136	0.0666615\\
0.138	0.0658525\\
0.14	0.0649509\\
0.142	0.0639606\\
0.144	0.0628858\\
0.146	0.061731\\
0.148	0.0605009\\
0.15	0.0592002\\
0.152	0.0578341\\
0.154	0.0564076\\
0.156	0.054926\\
0.158	0.0533947\\
0.16	0.051819\\
0.162	0.0502044\\
0.164	0.0485564\\
0.166	0.0468805\\
0.168	0.0451821\\
0.17	0.0434666\\
0.172	0.0417394\\
0.174	0.0414429\\
0.176	0.0415313\\
0.178	0.0414422\\
0.18	0.0411826\\
0.182	0.0407606\\
0.184	0.0401853\\
0.186	0.0394671\\
0.188	0.0386167\\
0.19	0.0376459\\
0.192	0.0365669\\
0.194	0.0353921\\
0.196	0.0341344\\
0.198	0.032807\\
0.2	0.0314227\\
0.202	0.0299946\\
0.204	0.0296001\\
0.206	0.0293863\\
0.208	0.0291728\\
0.21	0.0289597\\
0.212	0.0294722\\
0.214	0.030088\\
0.216	0.0306134\\
0.218	0.0310462\\
0.22	0.0313849\\
0.222	0.0316284\\
0.224	0.0317768\\
0.226	0.0318303\\
0.228	0.03179\\
0.23	0.0316577\\
0.232	0.0314354\\
0.234	0.0311261\\
0.236	0.030733\\
0.238	0.0302598\\
0.24	0.0297109\\
0.242	0.0290906\\
0.244	0.0284041\\
0.246	0.0276564\\
0.248	0.0268533\\
0.25	0.0260003\\
0.252	0.0251034\\
0.254	0.0248617\\
0.256	0.0246788\\
0.258	0.0244945\\
0.26	0.0243089\\
0.262	0.0241217\\
0.264	0.0239331\\
0.266	0.0244626\\
0.268	0.0250727\\
0.27	0.0256231\\
0.272	0.0261128\\
0.274	0.0265411\\
0.276	0.0269075\\
0.278	0.0272118\\
0.28	0.027454\\
0.282	0.0276344\\
0.284	0.0277534\\
0.286	0.0278117\\
0.288	0.0278102\\
0.29	0.02775\\
0.292	0.0276323\\
0.294	0.0274585\\
0.296	0.0272303\\
0.298	0.0269495\\
0.3	0.0266179\\
0.302	0.0262376\\
0.304	0.0258108\\
0.306	0.0253398\\
0.308	0.0248269\\
0.31	0.0242747\\
0.312	0.0236858\\
0.314	0.0230627\\
0.316	0.0224082\\
0.318	0.021725\\
0.32	0.0210159\\
0.322	0.0202837\\
0.324	0.0195312\\
0.326	0.0187613\\
0.328	0.0181098\\
0.33	0.0179989\\
0.332	0.0178624\\
0.334	0.0177011\\
0.336	0.0175159\\
0.338	0.0173076\\
0.34	0.0170773\\
0.342	0.016826\\
0.344	0.0168991\\
0.346	0.0170166\\
0.348	0.0171075\\
0.35	0.0171719\\
0.352	0.01721\\
0.354	0.0172221\\
0.356	0.0172085\\
0.358	0.0171697\\
0.36	0.0171062\\
0.362	0.0170185\\
0.364	0.0169073\\
0.366	0.0167732\\
0.368	0.016617\\
0.37	0.0164395\\
0.372	0.0162415\\
0.374	0.0160239\\
0.376	0.0157876\\
0.378	0.0155337\\
0.38	0.015263\\
0.382	0.0149767\\
0.384	0.0146757\\
0.386	0.0143611\\
0.388	0.0141233\\
0.39	0.0140581\\
0.392	0.0139928\\
0.394	0.0139274\\
0.396	0.013862\\
0.398	0.0137965\\
0.4	0.013731\\
0.402	0.0136655\\
0.404	0.0136\\
0.406	0.0135346\\
0.408	0.0134692\\
0.41	0.0134039\\
0.412	0.0133387\\
0.414	0.0132736\\
0.416	0.0132086\\
0.418	0.0131437\\
0.42	0.013079\\
0.422	0.0130144\\
0.424	0.01295\\
0.426	0.0128858\\
0.428	0.0128218\\
0.43	0.012758\\
0.432	0.0126944\\
0.434	0.012631\\
0.436	0.0125678\\
0.438	0.0125048\\
0.44	0.0124421\\
0.442	0.0123796\\
0.444	0.0123174\\
0.446	0.0122553\\
0.448	0.0121936\\
0.45	0.012132\\
0.452	0.0120707\\
0.454	0.0120097\\
0.456	0.0119489\\
0.458	0.0118883\\
0.46	0.011828\\
0.462	0.0117679\\
0.464	0.011708\\
0.466	0.0116484\\
0.468	0.011589\\
0.47	0.0115298\\
0.472	0.0114708\\
0.474	0.0114121\\
0.476	0.0113536\\
0.478	0.0112953\\
0.48	0.0112371\\
0.482	0.0111792\\
0.484	0.0111215\\
0.486	0.0110639\\
0.488	0.0110066\\
0.49	0.0109494\\
0.492	0.0108924\\
0.494	0.0108355\\
0.496	0.0107789\\
0.498	0.0107223\\
0.5	0.010666\\
0.502	0.0106097\\
0.504	0.0105537\\
0.506	0.0104978\\
0.508	0.010442\\
0.51	0.0103863\\
0.512	0.0103308\\
0.514	0.0102754\\
0.516	0.0102201\\
0.518	0.010165\\
0.52	0.01011\\
0.522	0.0100551\\
0.524	0.0100004\\
0.526	0.00994571\\
0.528	0.00989119\\
0.53	0.0098368\\
0.532	0.00978252\\
0.534	0.00972837\\
0.536	0.00967434\\
0.538	0.00962044\\
0.54	0.00956666\\
0.542	0.00951301\\
0.544	0.00945949\\
0.546	0.00940611\\
0.548	0.00935286\\
0.55	0.00929975\\
0.552	0.00924678\\
0.554	0.00919395\\
0.556	0.00914127\\
0.558	0.00908875\\
0.56	0.00903637\\
0.562	0.00898416\\
0.564	0.00893211\\
0.566	0.00888023\\
0.568	0.00882852\\
0.57	0.00877698\\
0.572	0.00872562\\
0.574	0.00867445\\
0.576	0.00862346\\
0.578	0.00857267\\
0.58	0.00852207\\
0.582	0.00847167\\
0.584	0.00842147\\
0.586	0.00837149\\
0.588	0.00832171\\
0.59	0.00827215\\
0.592	0.00822281\\
0.594	0.00817369\\
0.596	0.0081248\\
0.598	0.00807613\\
0.6	0.0080277\\
0.602	0.00797951\\
0.604	0.00793155\\
0.606	0.00788383\\
0.608	0.00783635\\
0.61	0.00778912\\
0.612	0.00774214\\
0.614	0.0076954\\
0.616	0.00764892\\
0.618	0.00760268\\
0.62	0.0075567\\
0.622	0.00751097\\
0.624	0.0074655\\
0.626	0.00742028\\
0.628	0.00737532\\
0.63	0.00733061\\
0.632	0.00728616\\
0.634	0.00724197\\
0.636	0.00719803\\
0.638	0.00715435\\
0.64	0.00711092\\
0.642	0.00706774\\
0.644	0.00702482\\
0.646	0.00698215\\
0.648	0.00693974\\
0.65	0.00689757\\
0.652	0.00685565\\
0.654	0.00681397\\
0.656	0.00677255\\
0.658	0.00673136\\
0.66	0.00669042\\
0.662	0.00664972\\
0.664	0.00660926\\
0.666	0.00656904\\
0.668	0.00652905\\
0.67	0.00648929\\
0.672	0.00644976\\
0.674	0.00641046\\
0.676	0.00637139\\
0.678	0.00633254\\
0.68	0.00629392\\
0.682	0.00625551\\
0.684	0.00621733\\
0.686	0.00617936\\
0.688	0.0061416\\
0.69	0.00610406\\
0.692	0.00606672\\
0.694	0.0060296\\
0.696	0.00599268\\
0.698	0.00595596\\
0.7	0.00591945\\
0.702	0.00588314\\
0.704	0.00584702\\
0.706	0.0058111\\
0.708	0.00577538\\
0.71	0.00573985\\
0.712	0.00570451\\
0.714	0.00566937\\
0.716	0.00563441\\
0.718	0.00559963\\
0.72	0.00556505\\
0.722	0.00553065\\
0.724	0.00549643\\
0.726	0.00546239\\
0.728	0.00542853\\
0.73	0.00539485\\
0.732	0.00536135\\
0.734	0.00532803\\
0.736	0.00529488\\
0.738	0.00526191\\
0.74	0.00522911\\
0.742	0.00519648\\
0.744	0.00516403\\
0.746	0.00513175\\
0.748	0.00509964\\
0.75	0.0050677\\
0.752	0.00503592\\
0.754	0.00500432\\
0.756	0.00497289\\
0.758	0.00494162\\
0.76	0.00491052\\
0.762	0.00487958\\
0.764	0.00484881\\
0.766	0.0048182\\
0.768	0.00478776\\
0.77	0.00475749\\
0.772	0.00472738\\
0.774	0.00469743\\
0.776	0.00466764\\
0.778	0.00463802\\
0.78	0.00460856\\
0.782	0.00457926\\
0.784	0.00455012\\
0.786	0.00452115\\
0.788	0.00449233\\
0.79	0.00446368\\
0.792	0.00443518\\
0.794	0.00440685\\
0.796	0.00437867\\
0.798	0.00435066\\
0.8	0.0043228\\
0.802	0.00429511\\
0.804	0.00426757\\
0.806	0.00424019\\
0.808	0.00421297\\
0.81	0.0041859\\
0.812	0.004159\\
0.814	0.00413225\\
0.816	0.00410566\\
0.818	0.00407922\\
0.82	0.00405295\\
0.822	0.00402682\\
0.824	0.00400086\\
0.826	0.00397505\\
0.828	0.00394939\\
0.83	0.00392389\\
0.832	0.00389854\\
0.834	0.00387335\\
0.836	0.00384831\\
0.838	0.00382343\\
0.84	0.0037987\\
0.842	0.00377412\\
0.844	0.00374969\\
0.846	0.00372541\\
0.848	0.00370129\\
0.85	0.00367732\\
0.852	0.00365349\\
0.854	0.00362982\\
0.856	0.00360629\\
0.858	0.00358292\\
0.86	0.00355969\\
0.862	0.00353661\\
0.864	0.00351368\\
0.866	0.00349089\\
0.868	0.00346825\\
0.87	0.00344575\\
0.872	0.0034234\\
0.874	0.00340119\\
0.876	0.00337913\\
0.878	0.00335721\\
0.88	0.00333543\\
0.882	0.00331379\\
0.884	0.00329229\\
0.886	0.00327093\\
0.888	0.00324971\\
0.89	0.00322862\\
0.892	0.00320768\\
0.894	0.00318687\\
0.896	0.00316619\\
0.898	0.00314565\\
0.9	0.00312525\\
0.902	0.00310498\\
0.904	0.00308484\\
0.906	0.00306483\\
0.908	0.00304495\\
0.91	0.0030252\\
0.912	0.00300558\\
0.914	0.00298609\\
0.916	0.00296672\\
0.918	0.00294748\\
0.92	0.00292836\\
0.922	0.00290937\\
0.924	0.0028905\\
0.926	0.00287176\\
0.928	0.00285313\\
0.93	0.00283463\\
0.932	0.00281624\\
0.934	0.00279797\\
0.936	0.00277982\\
0.938	0.00276179\\
0.94	0.00274387\\
0.942	0.00272607\\
0.944	0.00270838\\
0.946	0.0026908\\
0.948	0.00267334\\
0.95	0.00265599\\
0.952	0.00263875\\
0.954	0.00262161\\
0.956	0.00260459\\
0.958	0.00258767\\
0.96	0.00257086\\
0.962	0.00255416\\
0.964	0.00253756\\
0.966	0.00252107\\
0.968	0.00250468\\
0.97	0.00248839\\
0.972	0.0024722\\
0.974	0.00245612\\
0.976	0.00244013\\
0.978	0.00242425\\
0.98	0.00240847\\
0.982	0.00239278\\
0.984	0.00237719\\
0.986	0.0023617\\
0.988	0.0023463\\
0.99	0.002331\\
0.992	0.0023158\\
0.994	0.00230069\\
0.996	0.00228567\\
0.998	0.00227075\\
1	0.00225592\\
1.002	0.00224118\\
1.004	0.00222653\\
1.006	0.00221197\\
1.008	0.00219751\\
1.01	0.00218313\\
1.012	0.00216884\\
1.014	0.00215465\\
1.016	0.00214054\\
1.018	0.00212651\\
1.02	0.00211258\\
1.022	0.00209873\\
1.024	0.00208497\\
1.026	0.00207129\\
1.028	0.0020577\\
1.03	0.00204419\\
1.032	0.00203077\\
1.034	0.00201744\\
1.036	0.00200418\\
1.038	0.00199101\\
1.04	0.00197792\\
1.042	0.00196492\\
1.044	0.001952\\
1.046	0.00193915\\
1.048	0.00192639\\
1.05	0.00191371\\
1.052	0.00190112\\
1.054	0.0018886\\
1.056	0.00187616\\
1.058	0.0018638\\
1.06	0.00185152\\
1.062	0.00183931\\
1.064	0.00182719\\
1.066	0.00181514\\
1.068	0.00180317\\
1.07	0.00179128\\
1.072	0.00177947\\
1.074	0.00176773\\
1.076	0.00175607\\
1.078	0.00174448\\
1.08	0.00173297\\
1.082	0.00172153\\
1.084	0.00171017\\
1.086	0.00169888\\
1.088	0.00168766\\
1.09	0.00167652\\
1.092	0.00166545\\
1.094	0.00165446\\
1.096	0.00164353\\
1.098	0.00163268\\
1.1	0.0016219\\
1.102	0.00161119\\
1.104	0.00160055\\
1.106	0.00158998\\
1.108	0.00157948\\
1.11	0.00156905\\
1.112	0.00155869\\
1.114	0.0015484\\
1.116	0.00153817\\
1.118	0.00152802\\
1.12	0.00151793\\
1.122	0.00150791\\
1.124	0.00149795\\
1.126	0.00148806\\
1.128	0.00147824\\
1.13	0.00146848\\
1.132	0.00145879\\
1.134	0.00144916\\
1.136	0.0014396\\
1.138	0.0014301\\
1.14	0.00142066\\
1.142	0.00141129\\
1.144	0.00140198\\
1.146	0.00139273\\
1.148	0.00138354\\
1.15	0.00137441\\
1.152	0.00136535\\
1.154	0.00135634\\
1.156	0.0013474\\
1.158	0.00133851\\
1.16	0.00132968\\
1.162	0.00132091\\
1.164	0.0013122\\
1.166	0.00130355\\
1.168	0.00129496\\
1.17	0.00128642\\
1.172	0.00127794\\
1.174	0.00126952\\
1.176	0.00126115\\
1.178	0.00125283\\
1.18	0.00124458\\
1.182	0.00123637\\
1.184	0.00122822\\
1.186	0.00122013\\
1.188	0.00121209\\
1.19	0.0012041\\
1.192	0.00119616\\
1.194	0.00118828\\
1.196	0.00118045\\
1.198	0.00117267\\
1.2	0.00116494\\
1.202	0.00115726\\
1.204	0.00114963\\
1.206	0.00114205\\
1.208	0.00113453\\
1.21	0.00112705\\
1.212	0.00111962\\
1.214	0.00111224\\
1.216	0.00110491\\
1.218	0.00109762\\
1.22	0.00109038\\
1.222	0.0010832\\
1.224	0.00107605\\
1.226	0.00106896\\
1.228	0.00106191\\
1.23	0.0010549\\
1.232	0.00104795\\
1.234	0.00104103\\
1.236	0.00103417\\
1.238	0.00102734\\
1.24	0.00102057\\
1.242	0.00101383\\
1.244	0.00100714\\
1.246	0.0010005\\
1.248	0.000993893\\
1.25	0.000987333\\
1.252	0.000980815\\
1.254	0.00097434\\
1.256	0.000967908\\
1.258	0.000961517\\
1.26	0.000955167\\
1.262	0.000948859\\
1.264	0.000942593\\
1.266	0.000936367\\
1.268	0.000930181\\
1.27	0.000924036\\
1.272	0.000917931\\
1.274	0.000911866\\
1.276	0.000905841\\
1.278	0.000899854\\
1.28	0.000893907\\
1.282	0.000887999\\
1.284	0.000882129\\
1.286	0.000876298\\
1.288	0.000870505\\
1.29	0.00086475\\
1.292	0.000859032\\
1.294	0.000853352\\
1.296	0.00084771\\
1.298	0.000842104\\
1.3	0.000836535\\
1.302	0.000831003\\
1.304	0.000825507\\
1.306	0.000820047\\
1.308	0.000814623\\
1.31	0.000809235\\
1.312	0.000804276\\
1.314	0.000799438\\
1.316	0.000794629\\
1.318	0.000789848\\
1.32	0.000785095\\
1.322	0.000780369\\
1.324	0.000775671\\
1.326	0.000771\\
1.328	0.000766356\\
1.33	0.00076174\\
1.332	0.00075715\\
1.334	0.000752588\\
1.336	0.000748052\\
1.338	0.000743542\\
1.34	0.000739059\\
1.342	0.000734603\\
1.344	0.000730172\\
1.346	0.000725768\\
1.348	0.000721389\\
1.35	0.000717036\\
1.352	0.000712709\\
1.354	0.000708407\\
1.356	0.000704131\\
1.358	0.00069988\\
1.36	0.000695654\\
1.362	0.000691453\\
1.364	0.000687276\\
1.366	0.000683125\\
1.368	0.000678998\\
1.37	0.000674895\\
1.372	0.000670817\\
1.374	0.000666763\\
1.376	0.000662732\\
1.378	0.000658726\\
1.38	0.000654744\\
1.382	0.000650785\\
1.384	0.00064685\\
1.386	0.000642938\\
1.388	0.000639049\\
1.39	0.000635184\\
1.392	0.000631341\\
1.394	0.000627522\\
1.396	0.000623725\\
1.398	0.00061995\\
1.4	0.000616199\\
1.402	0.000612469\\
1.404	0.000608762\\
1.406	0.000605077\\
1.408	0.000601413\\
1.41	0.000597772\\
1.412	0.000594152\\
1.414	0.000590554\\
1.416	0.000586978\\
1.418	0.000583422\\
1.42	0.000579888\\
1.422	0.000576375\\
1.424	0.000572883\\
1.426	0.000569412\\
1.428	0.000565961\\
1.43	0.000562531\\
1.432	0.000559122\\
1.434	0.000555733\\
1.436	0.000552364\\
1.438	0.000549015\\
1.44	0.000545686\\
1.442	0.000542377\\
1.444	0.000539087\\
1.446	0.000535817\\
1.448	0.000532567\\
1.45	0.000529336\\
1.452	0.000526124\\
1.454	0.000522932\\
1.456	0.000519758\\
1.458	0.000516604\\
1.46	0.000513468\\
1.462	0.000510351\\
1.464	0.000507252\\
1.466	0.000504172\\
1.468	0.00050111\\
1.47	0.000498066\\
1.472	0.00049504\\
1.474	0.000492033\\
1.476	0.000489043\\
1.478	0.000486071\\
1.48	0.000483117\\
1.482	0.00048018\\
1.484	0.000477261\\
1.486	0.000474359\\
1.488	0.000471474\\
1.49	0.000468606\\
1.492	0.000465756\\
1.494	0.000462922\\
1.496	0.000460105\\
1.498	0.000457305\\
1.5	0.000454522\\
1.502	0.000451755\\
1.504	0.000449005\\
1.506	0.000446271\\
1.508	0.000443553\\
1.51	0.000440851\\
1.512	0.000438166\\
1.514	0.000435496\\
1.516	0.000432842\\
1.518	0.000430205\\
1.52	0.000427582\\
1.522	0.000424976\\
1.524	0.000422384\\
1.526	0.000419809\\
1.528	0.000417248\\
1.53	0.000414703\\
1.532	0.000412173\\
1.534	0.000409658\\
1.536	0.000407159\\
1.538	0.000404674\\
1.54	0.000402204\\
1.542	0.000399748\\
1.544	0.000397307\\
1.546	0.000394881\\
1.548	0.00039247\\
1.55	0.000390072\\
1.552	0.000387689\\
1.554	0.000385321\\
1.556	0.000382966\\
1.558	0.000380626\\
1.56	0.0003783\\
1.562	0.000375987\\
1.564	0.000373689\\
1.566	0.000371404\\
1.568	0.000369133\\
1.57	0.000366876\\
1.572	0.000364632\\
1.574	0.000362401\\
1.576	0.000360184\\
1.578	0.000357981\\
1.58	0.000355791\\
1.582	0.000353613\\
1.584	0.000351449\\
1.586	0.000349298\\
1.588	0.00034716\\
1.59	0.000345035\\
1.592	0.000342923\\
1.594	0.000340823\\
1.596	0.000338736\\
1.598	0.000336662\\
1.6	0.000334601\\
1.602	0.000332551\\
1.604	0.000330514\\
1.606	0.00032849\\
1.608	0.000326478\\
1.61	0.000324477\\
1.612	0.000322489\\
1.614	0.000320513\\
1.616	0.000318549\\
1.618	0.000316597\\
1.62	0.000314657\\
1.622	0.000312729\\
1.624	0.000310812\\
1.626	0.000308907\\
1.628	0.000307013\\
1.63	0.000305131\\
1.632	0.00030326\\
1.634	0.000301401\\
1.636	0.000299553\\
1.638	0.000297716\\
1.64	0.00029589\\
1.642	0.000294076\\
1.644	0.000292272\\
1.646	0.00029048\\
1.648	0.000288698\\
1.65	0.000286927\\
1.652	0.000285167\\
1.654	0.000283418\\
1.656	0.000281679\\
1.658	0.000279951\\
1.66	0.000278233\\
1.662	0.000276525\\
1.664	0.000274829\\
1.666	0.000273142\\
1.668	0.000271466\\
1.67	0.000269799\\
1.672	0.000268143\\
1.674	0.000266497\\
1.676	0.000264861\\
1.678	0.000263235\\
1.68	0.000261619\\
1.682	0.000260012\\
1.684	0.000258416\\
1.686	0.000256829\\
1.688	0.000255251\\
1.69	0.000253684\\
1.692	0.000252125\\
1.694	0.000250577\\
1.696	0.000249037\\
1.698	0.000247507\\
1.7	0.000245986\\
1.702	0.000244475\\
1.704	0.000242972\\
1.706	0.000241479\\
1.708	0.000239995\\
1.71	0.00023852\\
1.712	0.000237053\\
1.714	0.000235596\\
1.716	0.000234147\\
1.718	0.000232708\\
1.72	0.000231277\\
1.722	0.000229854\\
1.724	0.00022844\\
1.726	0.000227035\\
1.728	0.000225639\\
1.73	0.00022425\\
1.732	0.000222871\\
1.734	0.000221499\\
1.736	0.000220136\\
1.738	0.000218781\\
1.74	0.000217435\\
1.742	0.000216096\\
1.744	0.000214766\\
1.746	0.000213444\\
1.748	0.000212129\\
1.75	0.000210823\\
1.752	0.000209525\\
1.754	0.000208234\\
1.756	0.000206952\\
1.758	0.000205677\\
1.76	0.00020441\\
1.762	0.00020315\\
1.764	0.000201898\\
1.766	0.000200654\\
1.768	0.000199418\\
1.77	0.000198188\\
1.772	0.000196967\\
1.774	0.000195753\\
1.776	0.000194546\\
1.778	0.000193346\\
1.78	0.000192154\\
1.782	0.000190969\\
1.784	0.000189791\\
1.786	0.000188621\\
1.788	0.000187457\\
1.79	0.000186301\\
1.792	0.000185152\\
1.794	0.000184009\\
1.796	0.000182874\\
1.798	0.000181746\\
1.8	0.000180624\\
1.802	0.000179509\\
1.804	0.000178402\\
1.806	0.0001773\\
1.808	0.000176206\\
1.81	0.000175118\\
1.812	0.000174037\\
1.814	0.000172963\\
1.816	0.000171895\\
1.818	0.000170833\\
1.82	0.000169779\\
1.822	0.00016873\\
1.824	0.000167688\\
1.826	0.000166652\\
1.828	0.000165623\\
1.83	0.0001646\\
1.832	0.000163583\\
1.834	0.000162573\\
1.836	0.000161568\\
1.838	0.00016057\\
1.84	0.000159578\\
1.842	0.000158592\\
1.844	0.000157612\\
1.846	0.000156638\\
1.848	0.00015567\\
1.85	0.000154708\\
1.852	0.000153752\\
1.854	0.000152802\\
1.856	0.000151858\\
1.858	0.000150919\\
1.86	0.000149986\\
1.862	0.000149059\\
1.864	0.000148138\\
1.866	0.000147222\\
1.868	0.000146312\\
1.87	0.000145407\\
1.872	0.000144508\\
1.874	0.000143614\\
1.876	0.000142726\\
1.878	0.000141844\\
1.88	0.000140967\\
1.882	0.000140095\\
1.884	0.000139229\\
1.886	0.000138367\\
1.888	0.000137512\\
1.89	0.000136661\\
1.892	0.000135816\\
1.894	0.000134976\\
1.896	0.000134141\\
1.898	0.000133311\\
1.9	0.000132486\\
1.902	0.000131666\\
1.904	0.000130852\\
1.906	0.000130042\\
1.908	0.000129237\\
1.91	0.000128438\\
1.912	0.000127643\\
1.914	0.000126853\\
1.916	0.000126068\\
1.918	0.000125287\\
1.92	0.000124512\\
1.922	0.000123741\\
1.924	0.000122975\\
1.926	0.000122214\\
1.928	0.000121457\\
1.93	0.000120705\\
1.932	0.000119958\\
1.934	0.000119215\\
1.936	0.000118477\\
1.938	0.000117743\\
1.94	0.000117014\\
1.942	0.000116289\\
1.944	0.000115569\\
1.946	0.000114853\\
1.948	0.000114142\\
1.95	0.000113435\\
1.952	0.000112732\\
1.954	0.000112034\\
1.956	0.000111339\\
1.958	0.000110649\\
1.96	0.000109964\\
1.962	0.000109282\\
1.964	0.000108605\\
1.966	0.000107932\\
1.968	0.000107263\\
1.97	0.000106598\\
1.972	0.000105937\\
1.974	0.000105281\\
1.976	0.000104628\\
1.978	0.000103979\\
1.98	0.000103335\\
1.982	0.000102694\\
1.984	0.000102057\\
1.986	0.000101424\\
1.988	0.000100795\\
1.99	0.00010017\\
1.992	9.95487e-05\\
1.994	9.89312e-05\\
1.996	9.83175e-05\\
1.998	9.77076e-05\\
2	9.71014e-05\\
};
\addplot [color=mycolor9, line width=1.0pt, forget plot]
  table[row sep=crcr]{%
-0.024	0\\
-0.022	0\\
-0.02	0\\
-0.018	0\\
-0.016	0\\
-0.014	0\\
-0.012	0\\
-0.01	0\\
-0.008	0\\
-0.006	0\\
-0.004	0\\
-0.002	0\\
0	0\\
0.002	0.0516749\\
0.004	0.0926044\\
0.006	0.123448\\
0.008	0.145082\\
0.01	0.158537\\
0.012	0.164939\\
0.014	0.16545\\
0.016	0.161225\\
0.018	0.154765\\
0.02	0.14641\\
0.022	0.135972\\
0.024	0.124172\\
0.026	0.124855\\
0.028	0.130337\\
0.03	0.134992\\
0.032	0.138795\\
0.034	0.14173\\
0.036	0.14379\\
0.038	0.144978\\
0.04	0.145305\\
0.042	0.144792\\
0.044	0.143466\\
0.046	0.141365\\
0.048	0.138529\\
0.05	0.135008\\
0.052	0.130855\\
0.054	0.126128\\
0.056	0.120888\\
0.058	0.115199\\
0.06	0.109125\\
0.062	0.102734\\
0.064	0.0968311\\
0.066	0.097737\\
0.068	0.0984121\\
0.07	0.0988592\\
0.072	0.099082\\
0.074	0.0990847\\
0.076	0.098872\\
0.078	0.0984491\\
0.08	0.0978216\\
0.082	0.0969957\\
0.084	0.0959777\\
0.086	0.0947746\\
0.088	0.0933936\\
0.09	0.0918423\\
0.092	0.0901287\\
0.094	0.0882611\\
0.096	0.086248\\
0.098	0.0840984\\
0.1	0.0818214\\
0.102	0.0794264\\
0.104	0.076923\\
0.106	0.0743209\\
0.108	0.0716301\\
0.11	0.0688607\\
0.112	0.0683213\\
0.114	0.0685846\\
0.116	0.0687453\\
0.118	0.068804\\
0.12	0.0687614\\
0.122	0.0686187\\
0.124	0.0683772\\
0.126	0.0680387\\
0.128	0.0677487\\
0.13	0.0673559\\
0.132	0.0668599\\
0.134	0.0662634\\
0.136	0.0655695\\
0.138	0.0647816\\
0.14	0.0639033\\
0.142	0.0629386\\
0.144	0.0618916\\
0.146	0.0607666\\
0.148	0.0595682\\
0.15	0.0583012\\
0.152	0.0569704\\
0.154	0.0555809\\
0.156	0.0541378\\
0.158	0.0526463\\
0.16	0.0511117\\
0.162	0.0495393\\
0.164	0.0479344\\
0.166	0.0463023\\
0.168	0.0446484\\
0.17	0.042978\\
0.172	0.0412961\\
0.174	0.0414702\\
0.176	0.04156\\
0.178	0.0414712\\
0.18	0.0412109\\
0.182	0.0407873\\
0.184	0.0402097\\
0.186	0.0394882\\
0.188	0.038634\\
0.19	0.0376588\\
0.192	0.0365747\\
0.194	0.0353945\\
0.196	0.0341311\\
0.198	0.0327975\\
0.2	0.0314069\\
0.202	0.0299724\\
0.204	0.0293106\\
0.206	0.0291061\\
0.208	0.0289014\\
0.21	0.0288989\\
0.212	0.0295907\\
0.214	0.0301942\\
0.216	0.0307069\\
0.218	0.0311264\\
0.22	0.0314514\\
0.222	0.0316812\\
0.224	0.0318156\\
0.226	0.0318552\\
0.228	0.0318011\\
0.23	0.0316551\\
0.232	0.0314196\\
0.234	0.0310975\\
0.236	0.0306921\\
0.238	0.0302072\\
0.24	0.0296472\\
0.242	0.0290167\\
0.244	0.0283206\\
0.246	0.0275644\\
0.248	0.0267536\\
0.25	0.0258939\\
0.252	0.0249914\\
0.254	0.0246542\\
0.256	0.0244682\\
0.258	0.0242807\\
0.26	0.0240916\\
0.262	0.0239008\\
0.264	0.0239796\\
0.266	0.0246405\\
0.268	0.0252425\\
0.27	0.0257844\\
0.272	0.0262653\\
0.274	0.0266845\\
0.276	0.0270417\\
0.278	0.0273366\\
0.28	0.0275694\\
0.282	0.0277403\\
0.284	0.0278498\\
0.286	0.0278986\\
0.288	0.0278876\\
0.29	0.0278181\\
0.292	0.0276912\\
0.294	0.0275085\\
0.296	0.0272715\\
0.298	0.0269822\\
0.3	0.0266425\\
0.302	0.0262544\\
0.304	0.0258201\\
0.306	0.025342\\
0.308	0.0248224\\
0.31	0.024264\\
0.312	0.0236692\\
0.314	0.0230408\\
0.316	0.0223814\\
0.318	0.0216939\\
0.32	0.0209809\\
0.322	0.0202454\\
0.324	0.0194901\\
0.326	0.0187178\\
0.328	0.0181237\\
0.33	0.0180082\\
0.332	0.0178673\\
0.334	0.0177019\\
0.336	0.0175128\\
0.338	0.017301\\
0.34	0.0170675\\
0.342	0.0168132\\
0.344	0.0168456\\
0.346	0.016957\\
0.348	0.0170424\\
0.35	0.0171021\\
0.352	0.0171363\\
0.354	0.0171451\\
0.356	0.0171289\\
0.358	0.0170882\\
0.36	0.0170235\\
0.362	0.0169353\\
0.364	0.0168241\\
0.366	0.0166908\\
0.368	0.016536\\
0.37	0.0163605\\
0.372	0.0161651\\
0.374	0.0159506\\
0.376	0.0157181\\
0.378	0.0154684\\
0.38	0.0152024\\
0.382	0.0149212\\
0.384	0.0146259\\
0.386	0.0143173\\
0.388	0.0139967\\
0.39	0.0137992\\
0.392	0.0137372\\
0.394	0.013675\\
0.396	0.0136128\\
0.398	0.0135506\\
0.4	0.0134883\\
0.402	0.013426\\
0.404	0.0133637\\
0.406	0.0133014\\
0.408	0.0132391\\
0.41	0.0131768\\
0.412	0.0131145\\
0.414	0.0130523\\
0.416	0.0129902\\
0.418	0.0129281\\
0.42	0.0128661\\
0.422	0.0128042\\
0.424	0.0127423\\
0.426	0.0126806\\
0.428	0.012619\\
0.43	0.0125575\\
0.432	0.0124961\\
0.434	0.0124349\\
0.436	0.0123737\\
0.438	0.0123128\\
0.44	0.0122519\\
0.442	0.0121913\\
0.444	0.0121307\\
0.446	0.0120704\\
0.448	0.0120102\\
0.45	0.0119501\\
0.452	0.0118902\\
0.454	0.0118305\\
0.456	0.011771\\
0.458	0.0117116\\
0.46	0.0116524\\
0.462	0.0115934\\
0.464	0.0115346\\
0.466	0.0114759\\
0.468	0.0114174\\
0.47	0.0113591\\
0.472	0.0113009\\
0.474	0.0112429\\
0.476	0.0111852\\
0.478	0.0111275\\
0.48	0.0110701\\
0.482	0.0110128\\
0.484	0.0109557\\
0.486	0.0108987\\
0.488	0.010842\\
0.49	0.0107854\\
0.492	0.0107289\\
0.494	0.0106727\\
0.496	0.0106166\\
0.498	0.0105606\\
0.5	0.0105048\\
0.502	0.0104492\\
0.504	0.0103937\\
0.506	0.0103384\\
0.508	0.0102833\\
0.51	0.0102283\\
0.512	0.0101735\\
0.514	0.0101188\\
0.516	0.0100643\\
0.518	0.0100099\\
0.52	0.00995571\\
0.522	0.00990166\\
0.524	0.00984776\\
0.526	0.00979403\\
0.528	0.00974045\\
0.53	0.00968703\\
0.532	0.00963377\\
0.534	0.00958067\\
0.536	0.00952772\\
0.538	0.00947494\\
0.54	0.00942233\\
0.542	0.00936988\\
0.544	0.00931759\\
0.546	0.00926547\\
0.548	0.00921352\\
0.55	0.00916174\\
0.552	0.00911013\\
0.554	0.0090587\\
0.556	0.00900744\\
0.558	0.00895636\\
0.56	0.00890546\\
0.562	0.00885475\\
0.564	0.00880421\\
0.566	0.00875386\\
0.568	0.0087037\\
0.57	0.00865373\\
0.572	0.00860395\\
0.574	0.00855437\\
0.576	0.00850498\\
0.578	0.00845579\\
0.58	0.00840679\\
0.582	0.008358\\
0.584	0.00830942\\
0.586	0.00826103\\
0.588	0.00821286\\
0.59	0.00816489\\
0.592	0.00811713\\
0.594	0.00806958\\
0.596	0.00802224\\
0.598	0.00797512\\
0.6	0.00792821\\
0.602	0.00788152\\
0.604	0.00783505\\
0.606	0.00778879\\
0.608	0.00774276\\
0.61	0.00769694\\
0.612	0.00765134\\
0.614	0.00760597\\
0.616	0.00756082\\
0.618	0.00751589\\
0.62	0.00747118\\
0.622	0.0074267\\
0.624	0.00738244\\
0.626	0.0073384\\
0.628	0.00729459\\
0.63	0.007251\\
0.632	0.00720763\\
0.634	0.00716449\\
0.636	0.00712158\\
0.638	0.00707888\\
0.64	0.00703641\\
0.642	0.00699417\\
0.644	0.00695214\\
0.646	0.00691034\\
0.648	0.00686876\\
0.65	0.0068274\\
0.652	0.00678626\\
0.654	0.00674535\\
0.656	0.00670465\\
0.658	0.00666417\\
0.66	0.00662391\\
0.662	0.00658387\\
0.664	0.00654404\\
0.666	0.00650443\\
0.668	0.00646504\\
0.67	0.00642586\\
0.672	0.00638689\\
0.674	0.00634814\\
0.676	0.0063096\\
0.678	0.00627126\\
0.68	0.00623314\\
0.682	0.00619523\\
0.684	0.00615753\\
0.686	0.00612003\\
0.688	0.00608274\\
0.69	0.00604566\\
0.692	0.00600878\\
0.694	0.0059721\\
0.696	0.00593562\\
0.698	0.00589935\\
0.7	0.00586327\\
0.702	0.0058274\\
0.704	0.00579172\\
0.706	0.00575624\\
0.708	0.00572096\\
0.71	0.00568587\\
0.712	0.00565097\\
0.714	0.00561626\\
0.716	0.00558175\\
0.718	0.00554743\\
0.72	0.0055133\\
0.722	0.00547935\\
0.724	0.00544559\\
0.726	0.00541202\\
0.728	0.00537864\\
0.73	0.00534544\\
0.732	0.00531242\\
0.734	0.00527958\\
0.736	0.00524692\\
0.738	0.00521445\\
0.74	0.00518215\\
0.742	0.00515003\\
0.744	0.00511809\\
0.746	0.00508632\\
0.748	0.00505473\\
0.75	0.00502331\\
0.752	0.00499206\\
0.754	0.00496099\\
0.756	0.00493009\\
0.758	0.00489936\\
0.76	0.0048688\\
0.762	0.0048384\\
0.764	0.00480818\\
0.766	0.00477812\\
0.768	0.00474822\\
0.77	0.0047185\\
0.772	0.00468893\\
0.774	0.00465953\\
0.776	0.0046303\\
0.778	0.00460122\\
0.78	0.00457231\\
0.782	0.00454356\\
0.784	0.00451497\\
0.786	0.00448654\\
0.788	0.00445826\\
0.79	0.00443015\\
0.792	0.00440219\\
0.794	0.00437439\\
0.796	0.00434675\\
0.798	0.00431926\\
0.8	0.00429193\\
0.802	0.00426475\\
0.804	0.00423773\\
0.806	0.00421086\\
0.808	0.00418415\\
0.81	0.00415759\\
0.812	0.00413118\\
0.814	0.00410492\\
0.816	0.00407882\\
0.818	0.00405286\\
0.82	0.00402706\\
0.822	0.00400141\\
0.824	0.0039759\\
0.826	0.00395055\\
0.828	0.00392535\\
0.83	0.00390029\\
0.832	0.00387538\\
0.834	0.00385062\\
0.836	0.00382601\\
0.838	0.00380154\\
0.84	0.00377722\\
0.842	0.00375305\\
0.844	0.00372902\\
0.846	0.00370514\\
0.848	0.0036814\\
0.85	0.00365781\\
0.852	0.00363436\\
0.854	0.00361105\\
0.856	0.00358789\\
0.858	0.00356486\\
0.86	0.00354198\\
0.862	0.00351924\\
0.864	0.00349665\\
0.866	0.00347419\\
0.868	0.00345187\\
0.87	0.00342969\\
0.872	0.00340765\\
0.874	0.00338575\\
0.876	0.00336398\\
0.878	0.00334235\\
0.88	0.00332086\\
0.882	0.0032995\\
0.884	0.00327828\\
0.886	0.00325719\\
0.888	0.00323624\\
0.89	0.00321541\\
0.892	0.00319472\\
0.894	0.00317416\\
0.896	0.00315374\\
0.898	0.00313344\\
0.9	0.00311327\\
0.902	0.00309323\\
0.904	0.00307331\\
0.906	0.00305353\\
0.908	0.00303387\\
0.91	0.00301434\\
0.912	0.00299493\\
0.914	0.00297564\\
0.916	0.00295648\\
0.918	0.00293744\\
0.92	0.00291852\\
0.922	0.00289972\\
0.924	0.00288104\\
0.926	0.00286248\\
0.928	0.00284404\\
0.93	0.00282572\\
0.932	0.00280751\\
0.934	0.00278942\\
0.936	0.00277145\\
0.938	0.00275359\\
0.94	0.00273584\\
0.942	0.0027182\\
0.944	0.00270068\\
0.946	0.00268327\\
0.948	0.00266597\\
0.95	0.00264877\\
0.952	0.00263169\\
0.954	0.00261471\\
0.956	0.00259784\\
0.958	0.00258108\\
0.96	0.00256443\\
0.962	0.00254787\\
0.964	0.00253142\\
0.966	0.00251508\\
0.968	0.00249884\\
0.97	0.0024827\\
0.972	0.00246666\\
0.974	0.00245072\\
0.976	0.00243488\\
0.978	0.00241914\\
0.98	0.0024035\\
0.982	0.00238796\\
0.984	0.00237251\\
0.986	0.00235716\\
0.988	0.0023419\\
0.99	0.00232674\\
0.992	0.00231168\\
0.994	0.00229671\\
0.996	0.00228183\\
0.998	0.00226704\\
1	0.00225235\\
1.002	0.00223775\\
1.004	0.00222324\\
1.006	0.00220882\\
1.008	0.00219448\\
1.01	0.00218024\\
1.012	0.00216609\\
1.014	0.00215202\\
1.016	0.00213805\\
1.018	0.00212416\\
1.02	0.00211035\\
1.022	0.00209663\\
1.024	0.002083\\
1.026	0.00206945\\
1.028	0.00205599\\
1.03	0.00204261\\
1.032	0.00202932\\
1.034	0.00201611\\
1.036	0.00200298\\
1.038	0.00198993\\
1.04	0.00197697\\
1.042	0.00196409\\
1.044	0.00195129\\
1.046	0.00193856\\
1.048	0.00192592\\
1.05	0.00191336\\
1.052	0.00190088\\
1.054	0.00188848\\
1.056	0.00187616\\
1.058	0.00186391\\
1.06	0.00185174\\
1.062	0.00183965\\
1.064	0.00182764\\
1.066	0.0018157\\
1.068	0.00180384\\
1.07	0.00179205\\
1.072	0.00178034\\
1.074	0.00176871\\
1.076	0.00175715\\
1.078	0.00174566\\
1.08	0.00173425\\
1.082	0.00172291\\
1.084	0.00171165\\
1.086	0.00170046\\
1.088	0.00168934\\
1.09	0.00167829\\
1.092	0.00166731\\
1.094	0.00165641\\
1.096	0.00164557\\
1.098	0.00163481\\
1.1	0.00162411\\
1.102	0.00161349\\
1.104	0.00160293\\
1.106	0.00159245\\
1.108	0.00158203\\
1.11	0.00157168\\
1.112	0.0015614\\
1.114	0.00155118\\
1.116	0.00154103\\
1.118	0.00153095\\
1.12	0.00152094\\
1.122	0.00151099\\
1.124	0.0015011\\
1.126	0.00149128\\
1.128	0.00148153\\
1.13	0.00147184\\
1.132	0.00146221\\
1.134	0.00145264\\
1.136	0.00144314\\
1.138	0.00143371\\
1.14	0.00142433\\
1.142	0.00141502\\
1.144	0.00140576\\
1.146	0.00139657\\
1.148	0.00138744\\
1.15	0.00137837\\
1.152	0.00136936\\
1.154	0.0013604\\
1.156	0.00135151\\
1.158	0.00134267\\
1.16	0.0013339\\
1.162	0.00132518\\
1.164	0.00131652\\
1.166	0.00130791\\
1.168	0.00129936\\
1.17	0.00129087\\
1.172	0.00128243\\
1.174	0.00127405\\
1.176	0.00126573\\
1.178	0.00125746\\
1.18	0.00124924\\
1.182	0.00124108\\
1.184	0.00123297\\
1.186	0.00122491\\
1.188	0.00121691\\
1.19	0.00120895\\
1.192	0.00120105\\
1.194	0.00119321\\
1.196	0.00118541\\
1.198	0.00117766\\
1.2	0.00116997\\
1.202	0.00116232\\
1.204	0.00115473\\
1.206	0.00114718\\
1.208	0.00113968\\
1.21	0.00113224\\
1.212	0.00112484\\
1.214	0.00111748\\
1.216	0.00111018\\
1.218	0.00110292\\
1.22	0.00109572\\
1.222	0.00108855\\
1.224	0.00108144\\
1.226	0.00107437\\
1.228	0.00106734\\
1.23	0.00106037\\
1.232	0.00105343\\
1.234	0.00104654\\
1.236	0.0010397\\
1.238	0.0010329\\
1.24	0.00102615\\
1.242	0.00101944\\
1.244	0.00101277\\
1.246	0.00100614\\
1.248	0.000999561\\
1.25	0.000993022\\
1.252	0.000986525\\
1.254	0.00098007\\
1.256	0.000973658\\
1.258	0.000967286\\
1.26	0.000960956\\
1.262	0.000954667\\
1.264	0.000948419\\
1.266	0.000942212\\
1.268	0.000936044\\
1.27	0.000929917\\
1.272	0.00092383\\
1.274	0.000917782\\
1.276	0.000911773\\
1.278	0.000905803\\
1.28	0.000899872\\
1.282	0.00089398\\
1.284	0.000888126\\
1.286	0.000882311\\
1.288	0.000876533\\
1.29	0.000870792\\
1.292	0.00086509\\
1.294	0.000859424\\
1.296	0.000853795\\
1.298	0.000848484\\
1.3	0.000843477\\
1.302	0.000838498\\
1.304	0.000833548\\
1.306	0.000828627\\
1.308	0.000823733\\
1.31	0.000818868\\
1.312	0.00081403\\
1.314	0.00080922\\
1.316	0.000804437\\
1.318	0.000799682\\
1.32	0.000794954\\
1.322	0.000790253\\
1.324	0.00078558\\
1.326	0.000780933\\
1.328	0.000776312\\
1.33	0.000771719\\
1.332	0.000767151\\
1.334	0.00076261\\
1.336	0.000758095\\
1.338	0.000753607\\
1.34	0.000749144\\
1.342	0.000744706\\
1.344	0.000740295\\
1.346	0.000735909\\
1.348	0.000731548\\
1.35	0.000727212\\
1.352	0.000722902\\
1.354	0.000718616\\
1.356	0.000714355\\
1.358	0.000710119\\
1.36	0.000705908\\
1.362	0.00070172\\
1.364	0.000697557\\
1.366	0.000693419\\
1.368	0.000689304\\
1.37	0.000685213\\
1.372	0.000681146\\
1.374	0.000677103\\
1.376	0.000673083\\
1.378	0.000669086\\
1.38	0.000665113\\
1.382	0.000661163\\
1.384	0.000657235\\
1.386	0.000653331\\
1.388	0.000649449\\
1.39	0.00064559\\
1.392	0.000641754\\
1.394	0.00063794\\
1.396	0.000634148\\
1.398	0.000630378\\
1.4	0.00062663\\
1.402	0.000622904\\
1.404	0.0006192\\
1.406	0.000615517\\
1.408	0.000611856\\
1.41	0.000608216\\
1.412	0.000604597\\
1.414	0.000601\\
1.416	0.000597424\\
1.418	0.000593868\\
1.42	0.000590333\\
1.422	0.000586819\\
1.424	0.000583325\\
1.426	0.000579852\\
1.428	0.000576399\\
1.43	0.000572966\\
1.432	0.000569553\\
1.434	0.00056616\\
1.436	0.000562787\\
1.438	0.000559434\\
1.44	0.0005561\\
1.442	0.000552786\\
1.444	0.000549491\\
1.446	0.000546215\\
1.448	0.000542958\\
1.45	0.000539721\\
1.452	0.000536502\\
1.454	0.000533302\\
1.456	0.000530121\\
1.458	0.000526958\\
1.46	0.000523814\\
1.462	0.000520688\\
1.464	0.000517581\\
1.466	0.000514492\\
1.468	0.00051142\\
1.47	0.000508367\\
1.472	0.000505331\\
1.474	0.000502313\\
1.476	0.000499313\\
1.478	0.000496331\\
1.48	0.000493365\\
1.482	0.000490417\\
1.484	0.000487487\\
1.486	0.000484573\\
1.488	0.000481677\\
1.49	0.000478797\\
1.492	0.000475934\\
1.494	0.000473088\\
1.496	0.000470259\\
1.498	0.000467446\\
1.5	0.00046465\\
1.502	0.00046187\\
1.504	0.000459106\\
1.506	0.000456359\\
1.508	0.000453627\\
1.51	0.000450912\\
1.512	0.000448212\\
1.514	0.000445528\\
1.516	0.00044286\\
1.518	0.000440208\\
1.52	0.000437571\\
1.522	0.00043495\\
1.524	0.000432344\\
1.526	0.000429753\\
1.528	0.000427178\\
1.53	0.000424617\\
1.532	0.000422072\\
1.534	0.000419541\\
1.536	0.000417026\\
1.538	0.000414525\\
1.54	0.000412039\\
1.542	0.000409568\\
1.544	0.000407111\\
1.546	0.000404669\\
1.548	0.00040224\\
1.55	0.000399827\\
1.552	0.000397427\\
1.554	0.000395042\\
1.556	0.000392671\\
1.558	0.000390313\\
1.56	0.00038797\\
1.562	0.00038564\\
1.564	0.000383324\\
1.566	0.000381022\\
1.568	0.000378734\\
1.57	0.000376459\\
1.572	0.000374197\\
1.574	0.000371949\\
1.576	0.000369714\\
1.578	0.000367492\\
1.58	0.000365284\\
1.582	0.000363088\\
1.584	0.000360906\\
1.586	0.000358736\\
1.588	0.000356579\\
1.59	0.000354435\\
1.592	0.000352304\\
1.594	0.000350186\\
1.596	0.00034808\\
1.598	0.000345986\\
1.6	0.000343905\\
1.602	0.000341836\\
1.604	0.00033978\\
1.606	0.000337736\\
1.608	0.000335704\\
1.61	0.000333684\\
1.612	0.000331676\\
1.614	0.000329679\\
1.616	0.000327695\\
1.618	0.000325723\\
1.62	0.000323762\\
1.622	0.000321813\\
1.624	0.000319876\\
1.626	0.00031795\\
1.628	0.000316035\\
1.63	0.000314132\\
1.632	0.00031224\\
1.634	0.00031036\\
1.636	0.000308491\\
1.638	0.000306633\\
1.64	0.000304785\\
1.642	0.000302949\\
1.644	0.000301124\\
1.646	0.00029931\\
1.648	0.000297506\\
1.65	0.000295714\\
1.652	0.000293932\\
1.654	0.00029216\\
1.656	0.000290399\\
1.658	0.000288649\\
1.66	0.000286909\\
1.662	0.000285179\\
1.664	0.00028346\\
1.666	0.000281751\\
1.668	0.000280052\\
1.67	0.000278363\\
1.672	0.000276684\\
1.674	0.000275016\\
1.676	0.000273357\\
1.678	0.000271708\\
1.68	0.000270069\\
1.682	0.00026844\\
1.684	0.00026682\\
1.686	0.00026521\\
1.688	0.00026361\\
1.69	0.000262019\\
1.692	0.000260438\\
1.694	0.000258866\\
1.696	0.000257303\\
1.698	0.00025575\\
1.7	0.000254206\\
1.702	0.000252671\\
1.704	0.000251146\\
1.706	0.000249629\\
1.708	0.000248121\\
1.71	0.000246623\\
1.712	0.000245133\\
1.714	0.000243653\\
1.716	0.000242181\\
1.718	0.000240718\\
1.72	0.000239263\\
1.722	0.000237817\\
1.724	0.00023638\\
1.726	0.000234952\\
1.728	0.000233532\\
1.73	0.00023212\\
1.732	0.000230717\\
1.734	0.000229322\\
1.736	0.000227936\\
1.738	0.000226558\\
1.74	0.000225188\\
1.742	0.000223826\\
1.744	0.000222472\\
1.746	0.000221126\\
1.748	0.000219789\\
1.75	0.000218459\\
1.752	0.000217138\\
1.754	0.000215824\\
1.756	0.000214518\\
1.758	0.00021322\\
1.76	0.00021193\\
1.762	0.000210647\\
1.764	0.000209372\\
1.766	0.000208105\\
1.768	0.000206845\\
1.77	0.000205593\\
1.772	0.000204348\\
1.774	0.00020311\\
1.776	0.000201881\\
1.778	0.000200658\\
1.78	0.000199443\\
1.782	0.000198235\\
1.784	0.000197034\\
1.786	0.000195841\\
1.788	0.000194654\\
1.79	0.000193475\\
1.792	0.000192303\\
1.794	0.000191137\\
1.796	0.000189979\\
1.798	0.000188828\\
1.8	0.000187684\\
1.802	0.000186546\\
1.804	0.000185415\\
1.806	0.000184291\\
1.808	0.000183174\\
1.81	0.000182064\\
1.812	0.00018096\\
1.814	0.000179863\\
1.816	0.000178772\\
1.818	0.000177688\\
1.82	0.000176611\\
1.822	0.00017554\\
1.824	0.000174475\\
1.826	0.000173417\\
1.828	0.000172365\\
1.83	0.00017132\\
1.832	0.00017028\\
1.834	0.000169247\\
1.836	0.000168221\\
1.838	0.0001672\\
1.84	0.000166186\\
1.842	0.000165177\\
1.844	0.000164175\\
1.846	0.000163179\\
1.848	0.000162188\\
1.85	0.000161204\\
1.852	0.000160226\\
1.854	0.000159253\\
1.856	0.000158287\\
1.858	0.000157326\\
1.86	0.000156371\\
1.862	0.000155421\\
1.864	0.000154478\\
1.866	0.00015354\\
1.868	0.000152608\\
1.87	0.000151681\\
1.872	0.00015076\\
1.874	0.000149845\\
1.876	0.000148935\\
1.878	0.00014803\\
1.88	0.000147131\\
1.882	0.000146238\\
1.884	0.000145349\\
1.886	0.000144466\\
1.888	0.000143589\\
1.89	0.000142717\\
1.892	0.00014185\\
1.894	0.000140988\\
1.896	0.000140131\\
1.898	0.00013928\\
1.9	0.000138433\\
1.902	0.000137592\\
1.904	0.000136756\\
1.906	0.000135925\\
1.908	0.000135099\\
1.91	0.000134278\\
1.912	0.000133462\\
1.914	0.00013265\\
1.916	0.000131844\\
1.918	0.000131043\\
1.92	0.000130246\\
1.922	0.000129454\\
1.924	0.000128667\\
1.926	0.000127885\\
1.928	0.000127107\\
1.93	0.000126334\\
1.932	0.000125566\\
1.934	0.000124802\\
1.936	0.000124043\\
1.938	0.000123289\\
1.94	0.000122539\\
1.942	0.000121794\\
1.944	0.000121053\\
1.946	0.000120316\\
1.948	0.000119584\\
1.95	0.000118857\\
1.952	0.000118134\\
1.954	0.000117415\\
1.956	0.0001167\\
1.958	0.00011599\\
1.96	0.000115284\\
1.962	0.000114583\\
1.964	0.000113885\\
1.966	0.000113192\\
1.968	0.000112503\\
1.97	0.000111818\\
1.972	0.000111138\\
1.974	0.000110461\\
1.976	0.000109789\\
1.978	0.00010912\\
1.98	0.000108456\\
1.982	0.000107795\\
1.984	0.000107139\\
1.986	0.000106486\\
1.988	0.000105838\\
1.99	0.000105193\\
1.992	0.000104553\\
1.994	0.000103916\\
1.996	0.000103283\\
1.998	0.000102654\\
2	0.000102028\\
};
\addplot [color=mycolor10, line width=1.0pt, forget plot]
  table[row sep=crcr]{%
-0.024	0\\
-0.022	0\\
-0.02	0\\
-0.018	0\\
-0.016	0\\
-0.014	0\\
-0.012	0\\
-0.01	0\\
-0.008	0\\
-0.006	0\\
-0.004	0\\
-0.002	0\\
0	0\\
0.002	0.0516746\\
0.004	0.0926004\\
0.006	0.123431\\
0.008	0.145037\\
0.01	0.158444\\
0.012	0.164773\\
0.014	0.165187\\
0.016	0.160842\\
0.018	0.152851\\
0.02	0.142244\\
0.022	0.131954\\
0.024	0.120808\\
0.026	0.124731\\
0.028	0.130165\\
0.03	0.134766\\
0.032	0.138511\\
0.034	0.141386\\
0.036	0.143384\\
0.038	0.14451\\
0.04	0.144776\\
0.042	0.144204\\
0.044	0.142821\\
0.046	0.140666\\
0.048	0.137782\\
0.05	0.134217\\
0.052	0.130025\\
0.054	0.125264\\
0.056	0.119996\\
0.058	0.114285\\
0.06	0.108195\\
0.062	0.101794\\
0.064	0.0962319\\
0.066	0.0971153\\
0.068	0.0977696\\
0.07	0.0981979\\
0.072	0.0984035\\
0.074	0.0983908\\
0.076	0.0981644\\
0.078	0.0977294\\
0.08	0.0970913\\
0.082	0.0962563\\
0.084	0.0952308\\
0.086	0.0940216\\
0.088	0.0926359\\
0.09	0.0910814\\
0.092	0.0893659\\
0.094	0.0874978\\
0.096	0.0854857\\
0.098	0.0833384\\
0.1	0.081065\\
0.102	0.0786749\\
0.104	0.0761777\\
0.106	0.0735831\\
0.108	0.0709011\\
0.11	0.0681416\\
0.112	0.0674515\\
0.114	0.0677069\\
0.116	0.0678613\\
0.118	0.0679154\\
0.12	0.06787\\
0.122	0.0677261\\
0.124	0.0674851\\
0.126	0.0671852\\
0.128	0.0668951\\
0.13	0.0665011\\
0.132	0.0660056\\
0.134	0.0654113\\
0.136	0.0647214\\
0.138	0.0639392\\
0.14	0.0630682\\
0.142	0.0621125\\
0.144	0.061076\\
0.146	0.0599631\\
0.148	0.0587784\\
0.15	0.0575265\\
0.152	0.0562122\\
0.154	0.0548405\\
0.156	0.0534165\\
0.158	0.0519452\\
0.16	0.050432\\
0.162	0.0488819\\
0.164	0.0473004\\
0.166	0.0456925\\
0.168	0.0440637\\
0.17	0.0424189\\
0.172	0.0411041\\
0.174	0.0413738\\
0.176	0.0414595\\
0.178	0.041367\\
0.18	0.0411033\\
0.182	0.0406767\\
0.184	0.0400965\\
0.186	0.039373\\
0.188	0.0385171\\
0.19	0.0375407\\
0.192	0.036456\\
0.194	0.0352756\\
0.196	0.0340125\\
0.198	0.0326797\\
0.2	0.0312904\\
0.202	0.0298575\\
0.204	0.0289261\\
0.206	0.0287287\\
0.208	0.0285308\\
0.21	0.0289186\\
0.212	0.0295997\\
0.214	0.0301926\\
0.216	0.0306945\\
0.218	0.0311033\\
0.22	0.0314178\\
0.222	0.0316372\\
0.224	0.0317616\\
0.226	0.0317916\\
0.228	0.0317284\\
0.23	0.0315738\\
0.232	0.0313303\\
0.234	0.0310007\\
0.236	0.0305886\\
0.238	0.0300977\\
0.24	0.0295325\\
0.242	0.0288975\\
0.244	0.0281979\\
0.246	0.0274389\\
0.248	0.0266262\\
0.25	0.0257654\\
0.252	0.0248626\\
0.254	0.0243596\\
0.256	0.0241731\\
0.258	0.023985\\
0.26	0.0237952\\
0.262	0.0236037\\
0.264	0.0240403\\
0.266	0.0246936\\
0.268	0.025288\\
0.27	0.0258224\\
0.272	0.026296\\
0.274	0.0267082\\
0.276	0.0270585\\
0.278	0.0273468\\
0.28	0.0275732\\
0.282	0.027738\\
0.284	0.0278417\\
0.286	0.0278851\\
0.288	0.0278691\\
0.29	0.0277948\\
0.292	0.0276636\\
0.294	0.027477\\
0.296	0.0272366\\
0.298	0.0269442\\
0.3	0.0266018\\
0.302	0.0262115\\
0.304	0.0257755\\
0.306	0.025296\\
0.308	0.0247756\\
0.31	0.0242168\\
0.312	0.023622\\
0.314	0.0229941\\
0.316	0.0223356\\
0.318	0.0216493\\
0.32	0.0209381\\
0.322	0.0202047\\
0.324	0.0194519\\
0.326	0.0186826\\
0.328	0.0180876\\
0.33	0.0179708\\
0.332	0.017829\\
0.334	0.017663\\
0.336	0.0174738\\
0.338	0.0172622\\
0.34	0.0170291\\
0.342	0.0167757\\
0.344	0.0167904\\
0.346	0.0168982\\
0.348	0.0169806\\
0.35	0.0170379\\
0.352	0.0170702\\
0.354	0.0170777\\
0.356	0.0170609\\
0.358	0.0170201\\
0.36	0.0169558\\
0.362	0.0168686\\
0.364	0.0167589\\
0.366	0.0166276\\
0.368	0.0164753\\
0.37	0.0163028\\
0.372	0.0161109\\
0.374	0.0159003\\
0.376	0.0156721\\
0.378	0.015427\\
0.38	0.0151662\\
0.382	0.0148904\\
0.384	0.0146008\\
0.386	0.0142983\\
0.388	0.013984\\
0.39	0.0136589\\
0.392	0.0134474\\
0.394	0.0133875\\
0.396	0.0133277\\
0.398	0.0132678\\
0.4	0.0132079\\
0.402	0.013148\\
0.404	0.013088\\
0.406	0.0130281\\
0.408	0.0129682\\
0.41	0.0129083\\
0.412	0.0128484\\
0.414	0.0127885\\
0.416	0.0127287\\
0.418	0.012669\\
0.42	0.0126092\\
0.422	0.0125496\\
0.424	0.01249\\
0.426	0.0124304\\
0.428	0.012371\\
0.43	0.0123116\\
0.432	0.0122523\\
0.434	0.0121931\\
0.436	0.012134\\
0.438	0.012075\\
0.44	0.012016\\
0.442	0.0119572\\
0.444	0.0118985\\
0.446	0.0118399\\
0.448	0.0117814\\
0.45	0.011723\\
0.452	0.0116648\\
0.454	0.0116067\\
0.456	0.0115487\\
0.458	0.0114908\\
0.46	0.011433\\
0.462	0.0113754\\
0.464	0.0113179\\
0.466	0.0112606\\
0.468	0.0112034\\
0.47	0.0111463\\
0.472	0.0110894\\
0.474	0.0110326\\
0.476	0.010976\\
0.478	0.0109195\\
0.48	0.0108631\\
0.482	0.0108069\\
0.484	0.0107509\\
0.486	0.010695\\
0.488	0.0106392\\
0.49	0.0105837\\
0.492	0.0105282\\
0.494	0.0104729\\
0.496	0.0104178\\
0.498	0.0103629\\
0.5	0.0103081\\
0.502	0.0102534\\
0.504	0.010199\\
0.506	0.0101447\\
0.508	0.0100905\\
0.51	0.0100365\\
0.512	0.00998272\\
0.514	0.00992908\\
0.516	0.00987561\\
0.518	0.0098223\\
0.52	0.00976917\\
0.522	0.00971621\\
0.524	0.00966343\\
0.526	0.00961082\\
0.528	0.00955839\\
0.53	0.00950613\\
0.532	0.00945405\\
0.534	0.00940216\\
0.536	0.00935044\\
0.538	0.00929891\\
0.54	0.00924756\\
0.542	0.00919639\\
0.544	0.00914541\\
0.546	0.00909462\\
0.548	0.00904402\\
0.55	0.0089936\\
0.552	0.00894337\\
0.554	0.00889334\\
0.556	0.0088435\\
0.558	0.00879385\\
0.56	0.00874439\\
0.562	0.00869513\\
0.564	0.00864606\\
0.566	0.00859719\\
0.568	0.00854852\\
0.57	0.00850004\\
0.572	0.00845177\\
0.574	0.00840369\\
0.576	0.00835581\\
0.578	0.00830813\\
0.58	0.00826065\\
0.582	0.00821338\\
0.584	0.0081663\\
0.586	0.00811943\\
0.588	0.00807276\\
0.59	0.00802629\\
0.592	0.00798002\\
0.594	0.00793396\\
0.596	0.0078881\\
0.598	0.00784244\\
0.6	0.00779699\\
0.602	0.00775174\\
0.604	0.00770669\\
0.606	0.00766185\\
0.608	0.00761721\\
0.61	0.00757278\\
0.612	0.00752855\\
0.614	0.00748452\\
0.616	0.0074407\\
0.618	0.00739708\\
0.62	0.00735366\\
0.622	0.00731045\\
0.624	0.00726744\\
0.626	0.00722464\\
0.628	0.00718204\\
0.63	0.00713964\\
0.632	0.00709745\\
0.634	0.00705546\\
0.636	0.00701367\\
0.638	0.00697209\\
0.64	0.00693071\\
0.642	0.00688953\\
0.644	0.00684856\\
0.646	0.00680779\\
0.648	0.00676722\\
0.65	0.00672686\\
0.652	0.0066867\\
0.654	0.00664674\\
0.656	0.00660698\\
0.658	0.00656743\\
0.66	0.00652808\\
0.662	0.00648893\\
0.664	0.00644998\\
0.666	0.00641124\\
0.668	0.0063727\\
0.67	0.00633436\\
0.672	0.00629622\\
0.674	0.00625828\\
0.676	0.00622054\\
0.678	0.00618301\\
0.68	0.00614567\\
0.682	0.00610853\\
0.684	0.0060716\\
0.686	0.00603486\\
0.688	0.00599832\\
0.69	0.00596198\\
0.692	0.00592584\\
0.694	0.00588989\\
0.696	0.00585414\\
0.698	0.00581859\\
0.7	0.00578323\\
0.702	0.00574807\\
0.704	0.0057131\\
0.706	0.00567833\\
0.708	0.00564375\\
0.71	0.00560936\\
0.712	0.00557516\\
0.714	0.00554115\\
0.716	0.00550733\\
0.718	0.0054737\\
0.72	0.00544026\\
0.722	0.005407\\
0.724	0.00537393\\
0.726	0.00534104\\
0.728	0.00530834\\
0.73	0.00527582\\
0.732	0.00524348\\
0.734	0.00521133\\
0.736	0.00517935\\
0.738	0.00514755\\
0.74	0.00511593\\
0.742	0.00508448\\
0.744	0.00505321\\
0.746	0.00502212\\
0.748	0.0049912\\
0.75	0.00496045\\
0.752	0.00492987\\
0.754	0.00489946\\
0.756	0.00486922\\
0.758	0.00483915\\
0.76	0.00480924\\
0.762	0.00477951\\
0.764	0.00474993\\
0.766	0.00472052\\
0.768	0.00469128\\
0.77	0.00466219\\
0.772	0.00463327\\
0.774	0.00460451\\
0.776	0.00457591\\
0.778	0.00454747\\
0.78	0.00451918\\
0.782	0.00449105\\
0.784	0.00446308\\
0.786	0.00443527\\
0.788	0.00440761\\
0.79	0.0043801\\
0.792	0.00435275\\
0.794	0.00432555\\
0.796	0.0042985\\
0.798	0.00427161\\
0.8	0.00424486\\
0.802	0.00421827\\
0.804	0.00419182\\
0.806	0.00416553\\
0.808	0.00413938\\
0.81	0.00411338\\
0.812	0.00408753\\
0.814	0.00406183\\
0.816	0.00403627\\
0.818	0.00401086\\
0.82	0.00398559\\
0.822	0.00396047\\
0.824	0.0039355\\
0.826	0.00391067\\
0.828	0.00388598\\
0.83	0.00386144\\
0.832	0.00383704\\
0.834	0.00381278\\
0.836	0.00378866\\
0.838	0.00376469\\
0.84	0.00374086\\
0.842	0.00371716\\
0.844	0.00369361\\
0.846	0.0036702\\
0.848	0.00364693\\
0.85	0.0036238\\
0.852	0.0036008\\
0.854	0.00357795\\
0.856	0.00355523\\
0.858	0.00353265\\
0.86	0.0035102\\
0.862	0.00348789\\
0.864	0.00346572\\
0.866	0.00344368\\
0.868	0.00342178\\
0.87	0.00340001\\
0.872	0.00337837\\
0.874	0.00335687\\
0.876	0.0033355\\
0.878	0.00331426\\
0.88	0.00329315\\
0.882	0.00327217\\
0.884	0.00325133\\
0.886	0.00323061\\
0.888	0.00321002\\
0.89	0.00318956\\
0.892	0.00316923\\
0.894	0.00314902\\
0.896	0.00312894\\
0.898	0.00310898\\
0.9	0.00308915\\
0.902	0.00306945\\
0.904	0.00304987\\
0.906	0.00303041\\
0.908	0.00301107\\
0.91	0.00299185\\
0.912	0.00297276\\
0.914	0.00295378\\
0.916	0.00293492\\
0.918	0.00291618\\
0.92	0.00289756\\
0.922	0.00287906\\
0.924	0.00286067\\
0.926	0.0028424\\
0.928	0.00282424\\
0.93	0.0028062\\
0.932	0.00278827\\
0.934	0.00277045\\
0.936	0.00275275\\
0.938	0.00273515\\
0.94	0.00271767\\
0.942	0.0027003\\
0.944	0.00268303\\
0.946	0.00266587\\
0.948	0.00264882\\
0.95	0.00263188\\
0.952	0.00261505\\
0.954	0.00259831\\
0.956	0.00258169\\
0.958	0.00256516\\
0.96	0.00254874\\
0.962	0.00253243\\
0.964	0.00251621\\
0.966	0.0025001\\
0.968	0.00248408\\
0.97	0.00246817\\
0.972	0.00245235\\
0.974	0.00243664\\
0.976	0.00242102\\
0.978	0.0024055\\
0.98	0.00239007\\
0.982	0.00237474\\
0.984	0.00235951\\
0.986	0.00234437\\
0.988	0.00232932\\
0.99	0.00231437\\
0.992	0.00229951\\
0.994	0.00228475\\
0.996	0.00227007\\
0.998	0.00225549\\
1	0.00224099\\
1.002	0.00222659\\
1.004	0.00221228\\
1.006	0.00219805\\
1.008	0.00218391\\
1.01	0.00216986\\
1.012	0.0021559\\
1.014	0.00214202\\
1.016	0.00212823\\
1.018	0.00211453\\
1.02	0.00210091\\
1.022	0.00208738\\
1.024	0.00207393\\
1.026	0.00206056\\
1.028	0.00204728\\
1.03	0.00203408\\
1.032	0.00202096\\
1.034	0.00200793\\
1.036	0.00199497\\
1.038	0.0019821\\
1.04	0.0019693\\
1.042	0.00195659\\
1.044	0.00194396\\
1.046	0.0019314\\
1.048	0.00191893\\
1.05	0.00190653\\
1.052	0.00189421\\
1.054	0.00188197\\
1.056	0.0018698\\
1.058	0.00185771\\
1.06	0.0018457\\
1.062	0.00183377\\
1.064	0.0018219\\
1.066	0.00181012\\
1.068	0.00179841\\
1.07	0.00178677\\
1.072	0.00177521\\
1.074	0.00176372\\
1.076	0.0017523\\
1.078	0.00174096\\
1.08	0.00172969\\
1.082	0.00171849\\
1.084	0.00170736\\
1.086	0.0016963\\
1.088	0.00168532\\
1.09	0.0016744\\
1.092	0.00166356\\
1.094	0.00165278\\
1.096	0.00164207\\
1.098	0.00163144\\
1.1	0.00162087\\
1.102	0.00161036\\
1.104	0.00159993\\
1.106	0.00158956\\
1.108	0.00157926\\
1.11	0.00156903\\
1.112	0.00155887\\
1.114	0.00154876\\
1.116	0.00153873\\
1.118	0.00152876\\
1.12	0.00151885\\
1.122	0.00150901\\
1.124	0.00149923\\
1.126	0.00148952\\
1.128	0.00147987\\
1.13	0.00147028\\
1.132	0.00146075\\
1.134	0.00145129\\
1.136	0.00144188\\
1.138	0.00143254\\
1.14	0.00142326\\
1.142	0.00141404\\
1.144	0.00140488\\
1.146	0.00139578\\
1.148	0.00138674\\
1.15	0.00137775\\
1.152	0.00136883\\
1.154	0.00135996\\
1.156	0.00135115\\
1.158	0.0013424\\
1.16	0.00133371\\
1.162	0.00132507\\
1.164	0.00131649\\
1.166	0.00130796\\
1.168	0.00129949\\
1.17	0.00129108\\
1.172	0.00128271\\
1.174	0.00127441\\
1.176	0.00126615\\
1.178	0.00125795\\
1.18	0.00124981\\
1.182	0.00124171\\
1.184	0.00123367\\
1.186	0.00122568\\
1.188	0.00121775\\
1.19	0.00120986\\
1.192	0.00120203\\
1.194	0.00119424\\
1.196	0.00118651\\
1.198	0.00117882\\
1.2	0.00117119\\
1.202	0.0011636\\
1.204	0.00115607\\
1.206	0.00114858\\
1.208	0.00114114\\
1.21	0.00113375\\
1.212	0.00112641\\
1.214	0.00111911\\
1.216	0.00111186\\
1.218	0.00110466\\
1.22	0.0010975\\
1.222	0.00109039\\
1.224	0.00108333\\
1.226	0.00107631\\
1.228	0.00106934\\
1.23	0.00106241\\
1.232	0.00105552\\
1.234	0.00104868\\
1.236	0.00104189\\
1.238	0.00103513\\
1.24	0.00102842\\
1.242	0.00102176\\
1.244	0.00101514\\
1.246	0.00100856\\
1.248	0.00100202\\
1.25	0.000995521\\
1.252	0.000989067\\
1.254	0.000982655\\
1.256	0.000976283\\
1.258	0.000969953\\
1.26	0.000963663\\
1.262	0.000957414\\
1.264	0.000951205\\
1.266	0.000945037\\
1.268	0.000938907\\
1.27	0.000932818\\
1.272	0.000926768\\
1.274	0.000920757\\
1.276	0.000914784\\
1.278	0.00090885\\
1.28	0.000902955\\
1.282	0.000897098\\
1.284	0.000891278\\
1.286	0.000885496\\
1.288	0.000879847\\
1.29	0.00087474\\
1.292	0.000869661\\
1.294	0.000864611\\
1.296	0.000859589\\
1.298	0.000854595\\
1.3	0.00084963\\
1.302	0.000844692\\
1.304	0.000839782\\
1.306	0.000834899\\
1.308	0.000830044\\
1.31	0.000825217\\
1.312	0.000820416\\
1.314	0.000815643\\
1.316	0.000810896\\
1.318	0.000806177\\
1.32	0.000801484\\
1.322	0.000796817\\
1.324	0.000792177\\
1.326	0.000787563\\
1.328	0.000782975\\
1.33	0.000778413\\
1.332	0.000773877\\
1.334	0.000769367\\
1.336	0.000764882\\
1.338	0.000760422\\
1.34	0.000755988\\
1.342	0.000751579\\
1.344	0.000747195\\
1.346	0.000742836\\
1.348	0.000738502\\
1.35	0.000734192\\
1.352	0.000729907\\
1.354	0.000725646\\
1.356	0.00072141\\
1.358	0.000717197\\
1.36	0.000713009\\
1.362	0.000708844\\
1.364	0.000704703\\
1.366	0.000700586\\
1.368	0.000696492\\
1.37	0.000692421\\
1.372	0.000688374\\
1.374	0.00068435\\
1.376	0.000680349\\
1.378	0.00067637\\
1.38	0.000672414\\
1.382	0.000668481\\
1.384	0.00066457\\
1.386	0.000660682\\
1.388	0.000656816\\
1.39	0.000652972\\
1.392	0.00064915\\
1.394	0.00064535\\
1.396	0.000641571\\
1.398	0.000637814\\
1.4	0.000634079\\
1.402	0.000630365\\
1.404	0.000626672\\
1.406	0.000623\\
1.408	0.000619349\\
1.41	0.000615719\\
1.412	0.00061211\\
1.414	0.000608522\\
1.416	0.000604954\\
1.418	0.000601407\\
1.42	0.00059788\\
1.422	0.000594373\\
1.424	0.000590886\\
1.426	0.000587419\\
1.428	0.000583972\\
1.43	0.000580545\\
1.432	0.000577137\\
1.434	0.000573749\\
1.436	0.000570381\\
1.438	0.000567031\\
1.44	0.000563701\\
1.442	0.00056039\\
1.444	0.000557098\\
1.446	0.000553825\\
1.448	0.000550571\\
1.45	0.000547335\\
1.452	0.000544118\\
1.454	0.000540919\\
1.456	0.000537739\\
1.458	0.000534577\\
1.46	0.000531433\\
1.462	0.000528307\\
1.464	0.000525199\\
1.466	0.000522109\\
1.468	0.000519037\\
1.47	0.000515982\\
1.472	0.000512945\\
1.474	0.000509925\\
1.476	0.000506923\\
1.478	0.000503938\\
1.48	0.00050097\\
1.482	0.000498019\\
1.484	0.000495086\\
1.486	0.000492169\\
1.488	0.000489269\\
1.49	0.000486385\\
1.492	0.000483519\\
1.494	0.000480668\\
1.496	0.000477834\\
1.498	0.000475017\\
1.5	0.000472216\\
1.502	0.000469431\\
1.504	0.000466662\\
1.506	0.000463909\\
1.508	0.000461171\\
1.51	0.00045845\\
1.512	0.000455744\\
1.514	0.000453054\\
1.516	0.00045038\\
1.518	0.000447721\\
1.52	0.000445077\\
1.522	0.000442449\\
1.524	0.000439836\\
1.526	0.000437238\\
1.528	0.000434655\\
1.53	0.000432087\\
1.532	0.000429534\\
1.534	0.000426995\\
1.536	0.000424472\\
1.538	0.000421963\\
1.54	0.000419468\\
1.542	0.000416988\\
1.544	0.000414523\\
1.546	0.000412071\\
1.548	0.000409634\\
1.55	0.000407212\\
1.552	0.000404803\\
1.554	0.000402408\\
1.556	0.000400027\\
1.558	0.00039766\\
1.56	0.000395307\\
1.562	0.000392967\\
1.564	0.000390641\\
1.566	0.000388329\\
1.568	0.00038603\\
1.57	0.000383744\\
1.572	0.000381472\\
1.574	0.000379213\\
1.576	0.000376967\\
1.578	0.000374734\\
1.58	0.000372514\\
1.582	0.000370307\\
1.584	0.000368113\\
1.586	0.000365932\\
1.588	0.000363764\\
1.59	0.000361608\\
1.592	0.000359465\\
1.594	0.000357334\\
1.596	0.000355216\\
1.598	0.00035311\\
1.6	0.000351016\\
1.602	0.000348935\\
1.604	0.000346866\\
1.606	0.000344808\\
1.608	0.000342763\\
1.61	0.00034073\\
1.612	0.000338709\\
1.614	0.000336699\\
1.616	0.000334702\\
1.618	0.000332716\\
1.62	0.000330741\\
1.622	0.000328779\\
1.624	0.000326827\\
1.626	0.000324887\\
1.628	0.000322959\\
1.63	0.000321041\\
1.632	0.000319135\\
1.634	0.00031724\\
1.636	0.000315357\\
1.638	0.000313484\\
1.64	0.000311622\\
1.642	0.000309771\\
1.644	0.000307931\\
1.646	0.000306102\\
1.648	0.000304283\\
1.65	0.000302475\\
1.652	0.000300678\\
1.654	0.000298891\\
1.656	0.000297115\\
1.658	0.000295349\\
1.66	0.000293593\\
1.662	0.000291848\\
1.664	0.000290113\\
1.666	0.000288388\\
1.668	0.000286673\\
1.67	0.000284968\\
1.672	0.000283274\\
1.674	0.000281589\\
1.676	0.000279914\\
1.678	0.000278249\\
1.68	0.000276593\\
1.682	0.000274948\\
1.684	0.000273312\\
1.686	0.000271685\\
1.688	0.000270068\\
1.69	0.000268461\\
1.692	0.000266863\\
1.694	0.000265275\\
1.696	0.000263695\\
1.698	0.000262125\\
1.7	0.000260565\\
1.702	0.000259013\\
1.704	0.00025747\\
1.706	0.000255937\\
1.708	0.000254413\\
1.71	0.000252897\\
1.712	0.000251391\\
1.714	0.000249893\\
1.716	0.000248404\\
1.718	0.000246924\\
1.72	0.000245452\\
1.722	0.000243989\\
1.724	0.000242535\\
1.726	0.000241089\\
1.728	0.000239652\\
1.73	0.000238223\\
1.732	0.000236803\\
1.734	0.000235391\\
1.736	0.000233987\\
1.738	0.000232592\\
1.74	0.000231205\\
1.742	0.000229826\\
1.744	0.000228455\\
1.746	0.000227092\\
1.748	0.000225737\\
1.75	0.00022439\\
1.752	0.000223051\\
1.754	0.00022172\\
1.756	0.000220397\\
1.758	0.000219082\\
1.76	0.000217774\\
1.762	0.000216474\\
1.764	0.000215182\\
1.766	0.000213897\\
1.768	0.00021262\\
1.77	0.00021135\\
1.772	0.000210088\\
1.774	0.000208834\\
1.776	0.000207587\\
1.778	0.000206347\\
1.78	0.000205114\\
1.782	0.000203889\\
1.784	0.000202671\\
1.786	0.00020146\\
1.788	0.000200256\\
1.79	0.00019906\\
1.792	0.00019787\\
1.794	0.000196688\\
1.796	0.000195512\\
1.798	0.000194344\\
1.8	0.000193182\\
1.802	0.000192027\\
1.804	0.000190879\\
1.806	0.000189738\\
1.808	0.000188604\\
1.81	0.000187476\\
1.812	0.000186355\\
1.814	0.000185241\\
1.816	0.000184133\\
1.818	0.000183032\\
1.82	0.000181937\\
1.822	0.000180849\\
1.824	0.000179767\\
1.826	0.000178692\\
1.828	0.000177623\\
1.83	0.00017656\\
1.832	0.000175504\\
1.834	0.000174454\\
1.836	0.00017341\\
1.838	0.000172372\\
1.84	0.000171341\\
1.842	0.000170315\\
1.844	0.000169296\\
1.846	0.000168282\\
1.848	0.000167275\\
1.85	0.000166274\\
1.852	0.000165278\\
1.854	0.000164289\\
1.856	0.000163305\\
1.858	0.000162327\\
1.86	0.000161355\\
1.862	0.000160389\\
1.864	0.000159429\\
1.866	0.000158474\\
1.868	0.000157525\\
1.87	0.000156581\\
1.872	0.000155643\\
1.874	0.000154711\\
1.876	0.000153784\\
1.878	0.000152863\\
1.88	0.000151947\\
1.882	0.000151037\\
1.884	0.000150132\\
1.886	0.000149232\\
1.888	0.000148338\\
1.89	0.000147449\\
1.892	0.000146565\\
1.894	0.000145687\\
1.896	0.000144813\\
1.898	0.000143945\\
1.9	0.000143082\\
1.902	0.000142225\\
1.904	0.000141372\\
1.906	0.000140524\\
1.908	0.000139682\\
1.91	0.000138844\\
1.912	0.000138012\\
1.914	0.000137184\\
1.916	0.000136361\\
1.918	0.000135544\\
1.92	0.000134731\\
1.922	0.000133922\\
1.924	0.000133119\\
1.926	0.000132321\\
1.928	0.000131527\\
1.93	0.000130738\\
1.932	0.000129953\\
1.934	0.000129173\\
1.936	0.000128398\\
1.938	0.000127628\\
1.94	0.000126862\\
1.942	0.000126101\\
1.944	0.000125344\\
1.946	0.000124591\\
1.948	0.000123844\\
1.95	0.0001231\\
1.952	0.000122361\\
1.954	0.000121627\\
1.956	0.000120896\\
1.958	0.000120171\\
1.96	0.000119449\\
1.962	0.000118732\\
1.964	0.000118019\\
1.966	0.00011731\\
1.968	0.000116606\\
1.97	0.000115905\\
1.972	0.000115209\\
1.974	0.000114517\\
1.976	0.000113829\\
1.978	0.000113145\\
1.98	0.000112466\\
1.982	0.00011179\\
1.984	0.000111118\\
1.986	0.000110451\\
1.988	0.000109787\\
1.99	0.000109127\\
1.992	0.000108472\\
1.994	0.00010782\\
1.996	0.000107172\\
1.998	0.000106528\\
2	0.000105888\\
};
\addplot [color=mycolor11, line width=1.0pt, forget plot]
  table[row sep=crcr]{%
-0.024	0\\
-0.022	0\\
-0.02	0\\
-0.018	0\\
-0.016	0\\
-0.014	0\\
-0.012	0\\
-0.01	0\\
-0.008	0\\
-0.006	0\\
-0.004	0\\
-0.002	0\\
0	0\\
0.002	0.0513918\\
0.004	0.0916266\\
0.006	0.121567\\
0.008	0.142255\\
0.01	0.154846\\
0.012	0.160551\\
0.014	0.160586\\
0.016	0.156131\\
0.018	0.149396\\
0.02	0.141153\\
0.022	0.131034\\
0.024	0.11973\\
0.026	0.124488\\
0.028	0.129861\\
0.03	0.134399\\
0.032	0.138081\\
0.034	0.140892\\
0.036	0.142829\\
0.038	0.143896\\
0.04	0.144108\\
0.042	0.143486\\
0.044	0.142061\\
0.046	0.139869\\
0.048	0.136954\\
0.05	0.133365\\
0.052	0.129156\\
0.054	0.124384\\
0.056	0.119112\\
0.058	0.113403\\
0.06	0.10732\\
0.062	0.10093\\
0.064	0.0955584\\
0.066	0.0964363\\
0.068	0.0970867\\
0.07	0.0975125\\
0.072	0.0977169\\
0.074	0.0977041\\
0.076	0.0974785\\
0.078	0.0970452\\
0.08	0.0964098\\
0.082	0.0955781\\
0.084	0.0945567\\
0.086	0.0933523\\
0.088	0.0919721\\
0.09	0.0904238\\
0.092	0.0887152\\
0.094	0.0868548\\
0.096	0.084851\\
0.098	0.0827127\\
0.1	0.080449\\
0.102	0.0780694\\
0.104	0.0755834\\
0.106	0.0730008\\
0.108	0.0703313\\
0.11	0.0675851\\
0.112	0.0666512\\
0.114	0.0669003\\
0.116	0.06705\\
0.118	0.0671008\\
0.12	0.0670535\\
0.122	0.0669092\\
0.124	0.0666694\\
0.126	0.0663356\\
0.128	0.0659099\\
0.13	0.0654887\\
0.132	0.0650151\\
0.134	0.0644448\\
0.136	0.0637809\\
0.138	0.0630265\\
0.14	0.0621853\\
0.142	0.061261\\
0.144	0.0602575\\
0.146	0.0591792\\
0.148	0.0580303\\
0.15	0.0568155\\
0.152	0.0555395\\
0.154	0.054207\\
0.156	0.0528231\\
0.158	0.0513926\\
0.16	0.0499208\\
0.162	0.0484127\\
0.164	0.0468733\\
0.166	0.0453079\\
0.168	0.0437215\\
0.17	0.0421192\\
0.172	0.0410033\\
0.174	0.0412751\\
0.176	0.0413633\\
0.178	0.0412735\\
0.18	0.0410128\\
0.182	0.0405895\\
0.184	0.0400128\\
0.186	0.0392929\\
0.188	0.038441\\
0.19	0.0374687\\
0.192	0.0363882\\
0.194	0.0352122\\
0.196	0.0339536\\
0.198	0.0326254\\
0.2	0.0312407\\
0.202	0.0298126\\
0.204	0.0287304\\
0.206	0.0285419\\
0.208	0.0283525\\
0.21	0.0289726\\
0.212	0.0296435\\
0.214	0.0302262\\
0.216	0.030718\\
0.218	0.031117\\
0.22	0.0314219\\
0.222	0.0316322\\
0.224	0.0317478\\
0.226	0.0317695\\
0.228	0.0316985\\
0.23	0.0315368\\
0.232	0.0312868\\
0.234	0.0309514\\
0.236	0.0305342\\
0.238	0.030039\\
0.24	0.0294702\\
0.242	0.0288325\\
0.244	0.0281308\\
0.246	0.0273707\\
0.248	0.0265575\\
0.25	0.0256972\\
0.252	0.0247955\\
0.254	0.024273\\
0.256	0.0240872\\
0.258	0.0238996\\
0.26	0.0237103\\
0.262	0.0235191\\
0.264	0.0241023\\
0.266	0.0247502\\
0.268	0.0253394\\
0.27	0.0258688\\
0.272	0.0263375\\
0.274	0.026745\\
0.276	0.0270909\\
0.278	0.027375\\
0.28	0.0275974\\
0.282	0.0277585\\
0.284	0.0278588\\
0.286	0.0278991\\
0.288	0.0278802\\
0.29	0.0278035\\
0.292	0.02767\\
0.294	0.0274815\\
0.296	0.0272395\\
0.298	0.0269459\\
0.3	0.0266025\\
0.302	0.0262116\\
0.304	0.0257753\\
0.306	0.0252958\\
0.308	0.0247758\\
0.31	0.0242175\\
0.312	0.0236237\\
0.314	0.022997\\
0.316	0.02234\\
0.318	0.0216556\\
0.32	0.0209464\\
0.322	0.0202153\\
0.324	0.0194651\\
0.326	0.0186986\\
0.328	0.0180864\\
0.33	0.0179706\\
0.332	0.0178301\\
0.334	0.0176655\\
0.336	0.0174779\\
0.338	0.0172681\\
0.34	0.017037\\
0.342	0.0167858\\
0.344	0.0167982\\
0.346	0.0169012\\
0.348	0.0169794\\
0.35	0.0170328\\
0.352	0.0170617\\
0.354	0.0170663\\
0.356	0.0170469\\
0.358	0.0170039\\
0.36	0.0169379\\
0.362	0.0168493\\
0.364	0.0167389\\
0.366	0.0166071\\
0.368	0.0164547\\
0.37	0.0162825\\
0.372	0.0160912\\
0.374	0.0158818\\
0.376	0.0156549\\
0.378	0.0154117\\
0.38	0.0151529\\
0.382	0.0148796\\
0.384	0.0145927\\
0.386	0.0142932\\
0.388	0.0139821\\
0.39	0.0136606\\
0.392	0.0133418\\
0.394	0.0132833\\
0.396	0.0132248\\
0.398	0.0131662\\
0.4	0.0131076\\
0.402	0.013049\\
0.404	0.0129904\\
0.406	0.0129318\\
0.408	0.0128731\\
0.41	0.0128145\\
0.412	0.0127559\\
0.414	0.0126973\\
0.416	0.0126387\\
0.418	0.0125802\\
0.42	0.0125217\\
0.422	0.0124632\\
0.424	0.0124048\\
0.426	0.0123464\\
0.428	0.012288\\
0.43	0.0122297\\
0.432	0.0121715\\
0.434	0.0121133\\
0.436	0.0120552\\
0.438	0.0119971\\
0.44	0.0119391\\
0.442	0.0118812\\
0.444	0.0118234\\
0.446	0.0117657\\
0.448	0.011708\\
0.45	0.0116504\\
0.452	0.011593\\
0.454	0.0115356\\
0.456	0.0114783\\
0.458	0.0114211\\
0.46	0.0113641\\
0.462	0.0113071\\
0.464	0.0112503\\
0.466	0.0111935\\
0.468	0.0111369\\
0.47	0.0110804\\
0.472	0.011024\\
0.474	0.0109678\\
0.476	0.0109117\\
0.478	0.0108557\\
0.48	0.0107998\\
0.482	0.0107441\\
0.484	0.0106886\\
0.486	0.0106332\\
0.488	0.0105779\\
0.49	0.0105227\\
0.492	0.0104678\\
0.494	0.010413\\
0.496	0.0103583\\
0.498	0.0103038\\
0.5	0.0102494\\
0.502	0.0101953\\
0.504	0.0101412\\
0.506	0.0100874\\
0.508	0.0100337\\
0.51	0.00998021\\
0.512	0.00992688\\
0.514	0.00987372\\
0.516	0.00982075\\
0.518	0.00976795\\
0.52	0.00971533\\
0.522	0.0096629\\
0.524	0.00961065\\
0.526	0.00955859\\
0.528	0.00950672\\
0.53	0.00945503\\
0.532	0.00940354\\
0.534	0.00935223\\
0.536	0.00930112\\
0.538	0.0092502\\
0.54	0.00919948\\
0.542	0.00914895\\
0.544	0.00909862\\
0.546	0.00904849\\
0.548	0.00899855\\
0.55	0.00894881\\
0.552	0.00889927\\
0.554	0.00884993\\
0.556	0.00880078\\
0.558	0.00875184\\
0.56	0.00870309\\
0.562	0.00865455\\
0.564	0.00860621\\
0.566	0.00855806\\
0.568	0.00851012\\
0.57	0.00846237\\
0.572	0.00841483\\
0.574	0.00836748\\
0.576	0.00832033\\
0.578	0.00827338\\
0.58	0.00822663\\
0.582	0.00818008\\
0.584	0.00813372\\
0.586	0.00808756\\
0.588	0.0080416\\
0.59	0.00799583\\
0.592	0.00795026\\
0.594	0.00790489\\
0.596	0.00785971\\
0.598	0.00781472\\
0.6	0.00776993\\
0.602	0.00772533\\
0.604	0.00768092\\
0.606	0.00763671\\
0.608	0.00759268\\
0.61	0.00754885\\
0.612	0.00750521\\
0.614	0.00746176\\
0.616	0.0074185\\
0.618	0.00737544\\
0.62	0.00733256\\
0.622	0.00728987\\
0.624	0.00724737\\
0.626	0.00720507\\
0.628	0.00716295\\
0.63	0.00712102\\
0.632	0.00707928\\
0.634	0.00703773\\
0.636	0.00699638\\
0.638	0.00695521\\
0.64	0.00691423\\
0.642	0.00687344\\
0.644	0.00683284\\
0.646	0.00679244\\
0.648	0.00675222\\
0.65	0.0067122\\
0.652	0.00667237\\
0.654	0.00663273\\
0.656	0.00659328\\
0.658	0.00655403\\
0.66	0.00651497\\
0.662	0.0064761\\
0.664	0.00643743\\
0.666	0.00639894\\
0.668	0.00636066\\
0.67	0.00632257\\
0.672	0.00628467\\
0.674	0.00624697\\
0.676	0.00620946\\
0.678	0.00617215\\
0.68	0.00613503\\
0.682	0.00609811\\
0.684	0.00606139\\
0.686	0.00602486\\
0.688	0.00598853\\
0.69	0.00595239\\
0.692	0.00591645\\
0.694	0.0058807\\
0.696	0.00584514\\
0.698	0.00580979\\
0.7	0.00577462\\
0.702	0.00573965\\
0.704	0.00570488\\
0.706	0.00567029\\
0.708	0.0056359\\
0.71	0.0056017\\
0.712	0.00556769\\
0.714	0.00553387\\
0.716	0.00550025\\
0.718	0.00546681\\
0.72	0.00543356\\
0.722	0.00540049\\
0.724	0.00536761\\
0.726	0.00533492\\
0.728	0.00530242\\
0.73	0.00527009\\
0.732	0.00523795\\
0.734	0.00520599\\
0.736	0.00517421\\
0.738	0.00514261\\
0.74	0.00511119\\
0.742	0.00507995\\
0.744	0.00504888\\
0.746	0.00501798\\
0.748	0.00498726\\
0.75	0.00495672\\
0.752	0.00492634\\
0.754	0.00489614\\
0.756	0.0048661\\
0.758	0.00483623\\
0.76	0.00480653\\
0.762	0.004777\\
0.764	0.00474763\\
0.766	0.00471842\\
0.768	0.00468938\\
0.77	0.00466049\\
0.772	0.00463177\\
0.774	0.00460321\\
0.776	0.00457481\\
0.778	0.00454657\\
0.78	0.00451848\\
0.782	0.00449055\\
0.784	0.00446277\\
0.786	0.00443515\\
0.788	0.00440768\\
0.79	0.00438036\\
0.792	0.0043532\\
0.794	0.00432619\\
0.796	0.00429932\\
0.798	0.00427261\\
0.8	0.00424605\\
0.802	0.00421964\\
0.804	0.00419337\\
0.806	0.00416725\\
0.808	0.00414128\\
0.81	0.00411545\\
0.812	0.00408977\\
0.814	0.00406423\\
0.816	0.00403884\\
0.818	0.00401359\\
0.82	0.00398849\\
0.822	0.00396353\\
0.824	0.00393871\\
0.826	0.00391403\\
0.828	0.0038895\\
0.83	0.0038651\\
0.832	0.00384085\\
0.834	0.00381674\\
0.836	0.00379276\\
0.838	0.00376893\\
0.84	0.00374523\\
0.842	0.00372168\\
0.844	0.00369826\\
0.846	0.00367498\\
0.848	0.00365183\\
0.85	0.00362882\\
0.852	0.00360595\\
0.854	0.00358322\\
0.856	0.00356061\\
0.858	0.00353815\\
0.86	0.00351582\\
0.862	0.00349362\\
0.864	0.00347155\\
0.866	0.00344962\\
0.868	0.00342782\\
0.87	0.00340615\\
0.872	0.00338462\\
0.874	0.00336321\\
0.876	0.00334193\\
0.878	0.00332079\\
0.88	0.00329977\\
0.882	0.00327888\\
0.884	0.00325812\\
0.886	0.00323748\\
0.888	0.00321697\\
0.89	0.00319659\\
0.892	0.00317633\\
0.894	0.0031562\\
0.896	0.00313619\\
0.898	0.00311631\\
0.9	0.00309655\\
0.902	0.00307691\\
0.904	0.00305739\\
0.906	0.00303799\\
0.908	0.00301871\\
0.91	0.00299956\\
0.912	0.00298052\\
0.914	0.0029616\\
0.916	0.00294279\\
0.918	0.00292411\\
0.92	0.00290554\\
0.922	0.00288708\\
0.924	0.00286874\\
0.926	0.00285052\\
0.928	0.0028324\\
0.93	0.0028144\\
0.932	0.00279652\\
0.934	0.00277874\\
0.936	0.00276107\\
0.938	0.00274352\\
0.94	0.00272607\\
0.942	0.00270873\\
0.944	0.0026915\\
0.946	0.00267438\\
0.948	0.00265736\\
0.95	0.00264045\\
0.952	0.00262364\\
0.954	0.00260694\\
0.956	0.00259034\\
0.958	0.00257385\\
0.96	0.00255745\\
0.962	0.00254116\\
0.964	0.00252497\\
0.966	0.00250888\\
0.968	0.00249289\\
0.97	0.002477\\
0.972	0.00246121\\
0.974	0.00244552\\
0.976	0.00242992\\
0.978	0.00241442\\
0.98	0.00239902\\
0.982	0.00238371\\
0.984	0.0023685\\
0.986	0.00235338\\
0.988	0.00233835\\
0.99	0.00232342\\
0.992	0.00230858\\
0.994	0.00229383\\
0.996	0.00227917\\
0.998	0.00226461\\
1	0.00225013\\
1.002	0.00223575\\
1.004	0.00222145\\
1.006	0.00220724\\
1.008	0.00219312\\
1.01	0.00217909\\
1.012	0.00216514\\
1.014	0.00215128\\
1.016	0.00213751\\
1.018	0.00212382\\
1.02	0.00211022\\
1.022	0.0020967\\
1.024	0.00208326\\
1.026	0.00206991\\
1.028	0.00205664\\
1.03	0.00204346\\
1.032	0.00203036\\
1.034	0.00201733\\
1.036	0.00200439\\
1.038	0.00199153\\
1.04	0.00197875\\
1.042	0.00196605\\
1.044	0.00195343\\
1.046	0.00194089\\
1.048	0.00192843\\
1.05	0.00191604\\
1.052	0.00190374\\
1.054	0.0018915\\
1.056	0.00187935\\
1.058	0.00186727\\
1.06	0.00185527\\
1.062	0.00184335\\
1.064	0.0018315\\
1.066	0.00181972\\
1.068	0.00180802\\
1.07	0.00179639\\
1.072	0.00178484\\
1.074	0.00177336\\
1.076	0.00176195\\
1.078	0.00175061\\
1.08	0.00173935\\
1.082	0.00172815\\
1.084	0.00171703\\
1.086	0.00170598\\
1.088	0.001695\\
1.09	0.00168409\\
1.092	0.00167325\\
1.094	0.00166248\\
1.096	0.00165178\\
1.098	0.00164114\\
1.1	0.00163058\\
1.102	0.00162008\\
1.104	0.00160965\\
1.106	0.00159928\\
1.108	0.00158898\\
1.11	0.00157875\\
1.112	0.00156859\\
1.114	0.00155849\\
1.116	0.00154845\\
1.118	0.00153848\\
1.12	0.00152857\\
1.122	0.00151873\\
1.124	0.00150895\\
1.126	0.00149923\\
1.128	0.00148958\\
1.13	0.00147999\\
1.132	0.00147046\\
1.134	0.00146099\\
1.136	0.00145158\\
1.138	0.00144224\\
1.14	0.00143295\\
1.142	0.00142372\\
1.144	0.00141456\\
1.146	0.00140545\\
1.148	0.0013964\\
1.15	0.00138741\\
1.152	0.00137848\\
1.154	0.0013696\\
1.156	0.00136079\\
1.158	0.00135203\\
1.16	0.00134332\\
1.162	0.00133468\\
1.164	0.00132608\\
1.166	0.00131755\\
1.168	0.00130907\\
1.17	0.00130064\\
1.172	0.00129227\\
1.174	0.00128395\\
1.176	0.00127569\\
1.178	0.00126747\\
1.18	0.00125932\\
1.182	0.00125121\\
1.184	0.00124316\\
1.186	0.00123515\\
1.188	0.0012272\\
1.19	0.0012193\\
1.192	0.00121146\\
1.194	0.00120366\\
1.196	0.00119591\\
1.198	0.00118821\\
1.2	0.00118056\\
1.202	0.00117296\\
1.204	0.00116541\\
1.206	0.00115791\\
1.208	0.00115045\\
1.21	0.00114305\\
1.212	0.00113569\\
1.214	0.00112838\\
1.216	0.00112111\\
1.218	0.00111389\\
1.22	0.00110672\\
1.222	0.00109959\\
1.224	0.00109251\\
1.226	0.00108548\\
1.228	0.00107848\\
1.23	0.00107154\\
1.232	0.00106464\\
1.234	0.00105778\\
1.236	0.00105097\\
1.238	0.0010442\\
1.24	0.00103747\\
1.242	0.00103079\\
1.244	0.00102415\\
1.246	0.00101755\\
1.248	0.00101099\\
1.25	0.00100448\\
1.252	0.000998007\\
1.254	0.000991577\\
1.256	0.000985188\\
1.258	0.000978839\\
1.26	0.000972532\\
1.262	0.000966265\\
1.264	0.000960038\\
1.266	0.000953851\\
1.268	0.000947704\\
1.27	0.000941596\\
1.272	0.000935528\\
1.274	0.000929499\\
1.276	0.000923508\\
1.278	0.000917556\\
1.28	0.000911643\\
1.282	0.000905773\\
1.284	0.000900586\\
1.286	0.000895427\\
1.288	0.000890297\\
1.29	0.000885195\\
1.292	0.000880122\\
1.294	0.000875077\\
1.296	0.00087006\\
1.298	0.000865071\\
1.3	0.000860109\\
1.302	0.000855175\\
1.304	0.000850269\\
1.306	0.00084539\\
1.308	0.000840538\\
1.31	0.000835713\\
1.312	0.000830915\\
1.314	0.000826144\\
1.316	0.000821399\\
1.318	0.000816681\\
1.32	0.00081199\\
1.322	0.000807324\\
1.324	0.000802685\\
1.326	0.000798072\\
1.328	0.000793484\\
1.33	0.000788922\\
1.332	0.000784386\\
1.334	0.000779875\\
1.336	0.000775389\\
1.338	0.000770928\\
1.34	0.000766493\\
1.342	0.000762082\\
1.344	0.000757696\\
1.346	0.000753335\\
1.348	0.000748998\\
1.35	0.000744686\\
1.352	0.000740397\\
1.354	0.000736133\\
1.356	0.000731893\\
1.358	0.000727677\\
1.36	0.000723484\\
1.362	0.000719315\\
1.364	0.000715169\\
1.366	0.000711047\\
1.368	0.000706948\\
1.37	0.000702871\\
1.372	0.000698818\\
1.374	0.000694788\\
1.376	0.00069078\\
1.378	0.000686795\\
1.38	0.000682833\\
1.382	0.000678892\\
1.384	0.000674974\\
1.386	0.000671078\\
1.388	0.000667204\\
1.39	0.000663352\\
1.392	0.000659521\\
1.394	0.000655712\\
1.396	0.000651924\\
1.398	0.000648158\\
1.4	0.000644413\\
1.402	0.000640689\\
1.404	0.000636987\\
1.406	0.000633305\\
1.408	0.000629643\\
1.41	0.000626003\\
1.412	0.000622383\\
1.414	0.000618783\\
1.416	0.000615204\\
1.418	0.000611645\\
1.42	0.000608106\\
1.422	0.000604587\\
1.424	0.000601088\\
1.426	0.000597608\\
1.428	0.000594149\\
1.43	0.000590708\\
1.432	0.000587288\\
1.434	0.000583886\\
1.436	0.000580504\\
1.438	0.000577141\\
1.44	0.000573797\\
1.442	0.000570472\\
1.444	0.000567165\\
1.446	0.000563878\\
1.448	0.000560609\\
1.45	0.000557358\\
1.452	0.000554126\\
1.454	0.000550912\\
1.456	0.000547716\\
1.458	0.000544539\\
1.46	0.000541379\\
1.462	0.000538238\\
1.464	0.000535114\\
1.466	0.000532008\\
1.468	0.000528919\\
1.47	0.000525848\\
1.472	0.000522794\\
1.474	0.000519758\\
1.476	0.000516739\\
1.478	0.000513737\\
1.48	0.000510752\\
1.482	0.000507785\\
1.484	0.000504834\\
1.486	0.000501899\\
1.488	0.000498982\\
1.49	0.000496081\\
1.492	0.000493196\\
1.494	0.000490328\\
1.496	0.000487477\\
1.498	0.000484641\\
1.5	0.000481822\\
1.502	0.000479019\\
1.504	0.000476231\\
1.506	0.00047346\\
1.508	0.000470705\\
1.51	0.000467965\\
1.512	0.000465241\\
1.514	0.000462532\\
1.516	0.000459839\\
1.518	0.000457161\\
1.52	0.000454499\\
1.522	0.000451851\\
1.524	0.000449219\\
1.526	0.000446602\\
1.528	0.000444\\
1.53	0.000441413\\
1.532	0.000438841\\
1.534	0.000436283\\
1.536	0.00043374\\
1.538	0.000431212\\
1.54	0.000428698\\
1.542	0.000426198\\
1.544	0.000423713\\
1.546	0.000421242\\
1.548	0.000418785\\
1.55	0.000416343\\
1.552	0.000413914\\
1.554	0.0004115\\
1.556	0.000409099\\
1.558	0.000406712\\
1.56	0.000404338\\
1.562	0.000401979\\
1.564	0.000399633\\
1.566	0.0003973\\
1.568	0.000394981\\
1.57	0.000392675\\
1.572	0.000390382\\
1.574	0.000388103\\
1.576	0.000385837\\
1.578	0.000383583\\
1.58	0.000381343\\
1.582	0.000379116\\
1.584	0.000376901\\
1.586	0.000374699\\
1.588	0.00037251\\
1.59	0.000370333\\
1.592	0.000368169\\
1.594	0.000366018\\
1.596	0.000363879\\
1.598	0.000361752\\
1.6	0.000359637\\
1.602	0.000357535\\
1.604	0.000355445\\
1.606	0.000353367\\
1.608	0.0003513\\
1.61	0.000349246\\
1.612	0.000347204\\
1.614	0.000345173\\
1.616	0.000343154\\
1.618	0.000341147\\
1.62	0.000339151\\
1.622	0.000337167\\
1.624	0.000335194\\
1.626	0.000333232\\
1.628	0.000331282\\
1.63	0.000329343\\
1.632	0.000327416\\
1.634	0.000325499\\
1.636	0.000323594\\
1.638	0.000321699\\
1.64	0.000319816\\
1.642	0.000317943\\
1.644	0.000316081\\
1.646	0.00031423\\
1.648	0.00031239\\
1.65	0.00031056\\
1.652	0.000308741\\
1.654	0.000306932\\
1.656	0.000305134\\
1.658	0.000303346\\
1.66	0.000301568\\
1.662	0.000299801\\
1.664	0.000298044\\
1.666	0.000296297\\
1.668	0.00029456\\
1.67	0.000292833\\
1.672	0.000291117\\
1.674	0.00028941\\
1.676	0.000287713\\
1.678	0.000286026\\
1.68	0.000284348\\
1.682	0.00028268\\
1.684	0.000281022\\
1.686	0.000279374\\
1.688	0.000277735\\
1.69	0.000276105\\
1.692	0.000274485\\
1.694	0.000272875\\
1.696	0.000271273\\
1.698	0.000269681\\
1.7	0.000268098\\
1.702	0.000266525\\
1.704	0.00026496\\
1.706	0.000263404\\
1.708	0.000261858\\
1.71	0.00026032\\
1.712	0.000258792\\
1.714	0.000257272\\
1.716	0.000255761\\
1.718	0.000254259\\
1.72	0.000252765\\
1.722	0.00025128\\
1.724	0.000249804\\
1.726	0.000248336\\
1.728	0.000246877\\
1.73	0.000245426\\
1.732	0.000243984\\
1.734	0.00024255\\
1.736	0.000241125\\
1.738	0.000239707\\
1.74	0.000238298\\
1.742	0.000236897\\
1.744	0.000235504\\
1.746	0.00023412\\
1.748	0.000232743\\
1.75	0.000231375\\
1.752	0.000230014\\
1.754	0.000228661\\
1.756	0.000227316\\
1.758	0.000225979\\
1.76	0.00022465\\
1.762	0.000223328\\
1.764	0.000222014\\
1.766	0.000220708\\
1.768	0.000219409\\
1.77	0.000218118\\
1.772	0.000216835\\
1.774	0.000215559\\
1.776	0.00021429\\
1.778	0.000213029\\
1.78	0.000211775\\
1.782	0.000210528\\
1.784	0.000209289\\
1.786	0.000208057\\
1.788	0.000206832\\
1.79	0.000205614\\
1.792	0.000204403\\
1.794	0.0002032\\
1.796	0.000202003\\
1.798	0.000200813\\
1.8	0.000199631\\
1.802	0.000198455\\
1.804	0.000197286\\
1.806	0.000196124\\
1.808	0.000194968\\
1.81	0.00019382\\
1.812	0.000192678\\
1.814	0.000191542\\
1.816	0.000190414\\
1.818	0.000189292\\
1.82	0.000188176\\
1.822	0.000187067\\
1.824	0.000185965\\
1.826	0.000184869\\
1.828	0.000183779\\
1.83	0.000182696\\
1.832	0.000181619\\
1.834	0.000180548\\
1.836	0.000179484\\
1.838	0.000178425\\
1.84	0.000177373\\
1.842	0.000176327\\
1.844	0.000175288\\
1.846	0.000174254\\
1.848	0.000173226\\
1.85	0.000172205\\
1.852	0.000171189\\
1.854	0.000170179\\
1.856	0.000169175\\
1.858	0.000168177\\
1.86	0.000167185\\
1.862	0.000166199\\
1.864	0.000165218\\
1.866	0.000164243\\
1.868	0.000163274\\
1.87	0.000162311\\
1.872	0.000161353\\
1.874	0.000160401\\
1.876	0.000159454\\
1.878	0.000158513\\
1.88	0.000157577\\
1.882	0.000156647\\
1.884	0.000155722\\
1.886	0.000154803\\
1.888	0.000153889\\
1.89	0.000152981\\
1.892	0.000152077\\
1.894	0.000151179\\
1.896	0.000150287\\
1.898	0.000149399\\
1.9	0.000148517\\
1.902	0.00014764\\
1.904	0.000146768\\
1.906	0.000145901\\
1.908	0.000145039\\
1.91	0.000144182\\
1.912	0.00014333\\
1.914	0.000142484\\
1.916	0.000141642\\
1.918	0.000140805\\
1.92	0.000139973\\
1.922	0.000139146\\
1.924	0.000138324\\
1.926	0.000137506\\
1.928	0.000136694\\
1.93	0.000135886\\
1.932	0.000135083\\
1.934	0.000134284\\
1.936	0.00013349\\
1.938	0.000132701\\
1.94	0.000131917\\
1.942	0.000131137\\
1.944	0.000130362\\
1.946	0.000129591\\
1.948	0.000128825\\
1.95	0.000128063\\
1.952	0.000127306\\
1.954	0.000126553\\
1.956	0.000125804\\
1.958	0.00012506\\
1.96	0.000124321\\
1.962	0.000123585\\
1.964	0.000122854\\
1.966	0.000122128\\
1.968	0.000121405\\
1.97	0.000120687\\
1.972	0.000119973\\
1.974	0.000119263\\
1.976	0.000118557\\
1.978	0.000117856\\
1.98	0.000117159\\
1.982	0.000116465\\
1.984	0.000115776\\
1.986	0.000115091\\
1.988	0.00011441\\
1.99	0.000113733\\
1.992	0.00011306\\
1.994	0.000112391\\
1.996	0.000111725\\
1.998	0.000111064\\
2	0.000110406\\
};
\end{axis}

\begin{axis}[%
width={\figscale*1.528in},
height={\figscale*0.46in},
at={(\figscale*0.931in,\figscale*0.56in)},
scale only axis,
separate axis lines,
every outer x axis line/.append style={black},
every x tick label/.append style={font=\figticfont},
every x tick/.append style={black},
xmin=-0.025,
xmax=2,
xlabel={Time since step in s},
every outer y axis line/.append style={black},
every y tick label/.append style={font=\figticfont},
every y tick/.append style={black},
ymin=-0.0359538,
ymax=0.00159155,
ytick={-0.03,0},
ylabel={$\Delta\max_i|f_i(t)-\fcoi(t)|$},
axis background/.style={fill=white},
xmajorgrids,
ymajorgrids,
grid style={white!55!black, opacity=0.3},
ylabel style={rotate=-90, at={(axis description cs:-0.145,0.5)}, anchor=east}
]
\addplot [color=mycolor1, line width=1.4pt, forget plot]
  table[row sep=crcr]{%
-0.024	0\\
-0.022	0\\
-0.02	0\\
-0.018	0\\
-0.016	0\\
-0.014	0\\
-0.012	0\\
-0.01	0\\
-0.008	0\\
-0.006	0\\
-0.004	0\\
-0.002	0\\
0	0\\
0.002	0\\
0.004	0\\
0.006	0\\
0.008	0\\
0.01	0\\
0.012	0\\
0.014	0\\
0.016	0\\
0.018	0\\
0.02	0\\
0.022	0\\
0.024	0\\
0.026	0\\
0.028	0\\
0.03	0\\
0.032	0\\
0.034	0\\
0.036	0\\
0.038	0\\
0.04	0\\
0.042	0\\
0.044	0\\
0.046	0\\
0.048	0\\
0.05	0\\
0.052	0\\
0.054	0\\
0.056	0\\
0.058	0\\
0.06	0\\
0.062	0\\
0.064	0\\
0.066	0\\
0.068	0\\
0.07	0\\
0.072	0\\
0.074	0\\
0.076	0\\
0.078	0\\
0.08	0\\
0.082	0\\
0.084	0\\
0.086	0\\
0.088	0\\
0.09	0\\
0.092	0\\
0.094	0\\
0.096	0\\
0.098	0\\
0.1	0\\
0.102	0\\
0.104	0\\
0.106	0\\
0.108	0\\
0.11	0\\
0.112	0\\
0.114	0\\
0.116	0\\
0.118	0\\
0.12	0\\
0.122	0\\
0.124	0\\
0.126	0\\
0.128	0\\
0.13	0\\
0.132	0\\
0.134	0\\
0.136	0\\
0.138	0\\
0.14	0\\
0.142	0\\
0.144	0\\
0.146	0\\
0.148	0\\
0.15	0\\
0.152	0\\
0.154	0\\
0.156	0\\
0.158	0\\
0.16	0\\
0.162	0\\
0.164	0\\
0.166	0\\
0.168	0\\
0.17	0\\
0.172	0\\
0.174	0\\
0.176	0\\
0.178	0\\
0.18	0\\
0.182	0\\
0.184	0\\
0.186	0\\
0.188	0\\
0.19	0\\
0.192	0\\
0.194	0\\
0.196	0\\
0.198	0\\
0.2	0\\
0.202	0\\
0.204	0\\
0.206	0\\
0.208	0\\
0.21	0\\
0.212	0\\
0.214	0\\
0.216	0\\
0.218	0\\
0.22	0\\
0.222	0\\
0.224	0\\
0.226	0\\
0.228	0\\
0.23	0\\
0.232	0\\
0.234	0\\
0.236	0\\
0.238	0\\
0.24	0\\
0.242	0\\
0.244	0\\
0.246	0\\
0.248	0\\
0.25	0\\
0.252	0\\
0.254	0\\
0.256	0\\
0.258	0\\
0.26	0\\
0.262	0\\
0.264	0\\
0.266	0\\
0.268	0\\
0.27	0\\
0.272	0\\
0.274	0\\
0.276	0\\
0.278	0\\
0.28	0\\
0.282	0\\
0.284	0\\
0.286	0\\
0.288	0\\
0.29	0\\
0.292	0\\
0.294	0\\
0.296	0\\
0.298	0\\
0.3	0\\
0.302	0\\
0.304	0\\
0.306	0\\
0.308	0\\
0.31	0\\
0.312	0\\
0.314	0\\
0.316	0\\
0.318	0\\
0.32	0\\
0.322	0\\
0.324	0\\
0.326	0\\
0.328	0\\
0.33	0\\
0.332	0\\
0.334	0\\
0.336	0\\
0.338	0\\
0.34	0\\
0.342	0\\
0.344	0\\
0.346	0\\
0.348	0\\
0.35	0\\
0.352	0\\
0.354	0\\
0.356	0\\
0.358	0\\
0.36	0\\
0.362	0\\
0.364	0\\
0.366	0\\
0.368	0\\
0.37	0\\
0.372	0\\
0.374	0\\
0.376	0\\
0.378	0\\
0.38	0\\
0.382	0\\
0.384	0\\
0.386	0\\
0.388	0\\
0.39	0\\
0.392	0\\
0.394	0\\
0.396	0\\
0.398	0\\
0.4	0\\
0.402	0\\
0.404	0\\
0.406	0\\
0.408	0\\
0.41	0\\
0.412	0\\
0.414	0\\
0.416	0\\
0.418	0\\
0.42	0\\
0.422	0\\
0.424	0\\
0.426	0\\
0.428	0\\
0.43	0\\
0.432	0\\
0.434	0\\
0.436	0\\
0.438	0\\
0.44	0\\
0.442	0\\
0.444	0\\
0.446	0\\
0.448	0\\
0.45	0\\
0.452	0\\
0.454	0\\
0.456	0\\
0.458	0\\
0.46	0\\
0.462	0\\
0.464	0\\
0.466	0\\
0.468	0\\
0.47	0\\
0.472	0\\
0.474	0\\
0.476	0\\
0.478	0\\
0.48	0\\
0.482	0\\
0.484	0\\
0.486	0\\
0.488	0\\
0.49	0\\
0.492	0\\
0.494	0\\
0.496	0\\
0.498	0\\
0.5	0\\
0.502	0\\
0.504	0\\
0.506	0\\
0.508	0\\
0.51	0\\
0.512	0\\
0.514	0\\
0.516	0\\
0.518	0\\
0.52	0\\
0.522	0\\
0.524	0\\
0.526	0\\
0.528	0\\
0.53	0\\
0.532	0\\
0.534	0\\
0.536	0\\
0.538	0\\
0.54	0\\
0.542	0\\
0.544	0\\
0.546	0\\
0.548	0\\
0.55	0\\
0.552	0\\
0.554	0\\
0.556	0\\
0.558	0\\
0.56	0\\
0.562	0\\
0.564	0\\
0.566	0\\
0.568	0\\
0.57	0\\
0.572	0\\
0.574	0\\
0.576	0\\
0.578	0\\
0.58	0\\
0.582	0\\
0.584	0\\
0.586	0\\
0.588	0\\
0.59	0\\
0.592	0\\
0.594	0\\
0.596	0\\
0.598	0\\
0.6	0\\
0.602	0\\
0.604	0\\
0.606	0\\
0.608	0\\
0.61	0\\
0.612	0\\
0.614	0\\
0.616	0\\
0.618	0\\
0.62	0\\
0.622	0\\
0.624	0\\
0.626	0\\
0.628	0\\
0.63	0\\
0.632	0\\
0.634	0\\
0.636	0\\
0.638	0\\
0.64	0\\
0.642	0\\
0.644	0\\
0.646	0\\
0.648	0\\
0.65	0\\
0.652	0\\
0.654	0\\
0.656	0\\
0.658	0\\
0.66	0\\
0.662	0\\
0.664	0\\
0.666	0\\
0.668	0\\
0.67	0\\
0.672	0\\
0.674	0\\
0.676	0\\
0.678	0\\
0.68	0\\
0.682	0\\
0.684	0\\
0.686	0\\
0.688	0\\
0.69	0\\
0.692	0\\
0.694	0\\
0.696	0\\
0.698	0\\
0.7	0\\
0.702	0\\
0.704	0\\
0.706	0\\
0.708	0\\
0.71	0\\
0.712	0\\
0.714	0\\
0.716	0\\
0.718	0\\
0.72	0\\
0.722	0\\
0.724	0\\
0.726	0\\
0.728	0\\
0.73	0\\
0.732	0\\
0.734	0\\
0.736	0\\
0.738	0\\
0.74	0\\
0.742	0\\
0.744	0\\
0.746	0\\
0.748	0\\
0.75	0\\
0.752	0\\
0.754	0\\
0.756	0\\
0.758	0\\
0.76	0\\
0.762	0\\
0.764	0\\
0.766	0\\
0.768	0\\
0.77	0\\
0.772	0\\
0.774	0\\
0.776	0\\
0.778	0\\
0.78	0\\
0.782	0\\
0.784	0\\
0.786	0\\
0.788	0\\
0.79	0\\
0.792	0\\
0.794	0\\
0.796	0\\
0.798	0\\
0.8	0\\
0.802	0\\
0.804	0\\
0.806	0\\
0.808	0\\
0.81	0\\
0.812	0\\
0.814	0\\
0.816	0\\
0.818	0\\
0.82	0\\
0.822	0\\
0.824	0\\
0.826	0\\
0.828	0\\
0.83	0\\
0.832	0\\
0.834	0\\
0.836	0\\
0.838	0\\
0.84	0\\
0.842	0\\
0.844	0\\
0.846	0\\
0.848	0\\
0.85	0\\
0.852	0\\
0.854	0\\
0.856	0\\
0.858	0\\
0.86	0\\
0.862	0\\
0.864	0\\
0.866	0\\
0.868	0\\
0.87	0\\
0.872	0\\
0.874	0\\
0.876	0\\
0.878	0\\
0.88	0\\
0.882	0\\
0.884	0\\
0.886	0\\
0.888	0\\
0.89	0\\
0.892	0\\
0.894	0\\
0.896	0\\
0.898	0\\
0.9	0\\
0.902	0\\
0.904	0\\
0.906	0\\
0.908	0\\
0.91	0\\
0.912	0\\
0.914	0\\
0.916	0\\
0.918	0\\
0.92	0\\
0.922	0\\
0.924	0\\
0.926	0\\
0.928	0\\
0.93	0\\
0.932	0\\
0.934	0\\
0.936	0\\
0.938	0\\
0.94	0\\
0.942	0\\
0.944	0\\
0.946	0\\
0.948	0\\
0.95	0\\
0.952	0\\
0.954	0\\
0.956	0\\
0.958	0\\
0.96	0\\
0.962	0\\
0.964	0\\
0.966	0\\
0.968	0\\
0.97	0\\
0.972	0\\
0.974	0\\
0.976	0\\
0.978	0\\
0.98	0\\
0.982	0\\
0.984	0\\
0.986	0\\
0.988	0\\
0.99	0\\
0.992	0\\
0.994	0\\
0.996	0\\
0.998	0\\
1	0\\
1.002	0\\
1.004	0\\
1.006	0\\
1.008	0\\
1.01	0\\
1.012	0\\
1.014	0\\
1.016	0\\
1.018	0\\
1.02	0\\
1.022	0\\
1.024	0\\
1.026	0\\
1.028	0\\
1.03	0\\
1.032	0\\
1.034	0\\
1.036	0\\
1.038	0\\
1.04	0\\
1.042	0\\
1.044	0\\
1.046	0\\
1.048	0\\
1.05	0\\
1.052	0\\
1.054	0\\
1.056	0\\
1.058	0\\
1.06	0\\
1.062	0\\
1.064	0\\
1.066	0\\
1.068	0\\
1.07	0\\
1.072	0\\
1.074	0\\
1.076	0\\
1.078	0\\
1.08	0\\
1.082	0\\
1.084	0\\
1.086	0\\
1.088	0\\
1.09	0\\
1.092	0\\
1.094	0\\
1.096	0\\
1.098	0\\
1.1	0\\
1.102	0\\
1.104	0\\
1.106	0\\
1.108	0\\
1.11	0\\
1.112	0\\
1.114	0\\
1.116	0\\
1.118	0\\
1.12	0\\
1.122	0\\
1.124	0\\
1.126	0\\
1.128	0\\
1.13	0\\
1.132	0\\
1.134	0\\
1.136	0\\
1.138	0\\
1.14	0\\
1.142	0\\
1.144	0\\
1.146	0\\
1.148	0\\
1.15	0\\
1.152	0\\
1.154	0\\
1.156	0\\
1.158	0\\
1.16	0\\
1.162	0\\
1.164	0\\
1.166	0\\
1.168	0\\
1.17	0\\
1.172	0\\
1.174	0\\
1.176	0\\
1.178	0\\
1.18	0\\
1.182	0\\
1.184	0\\
1.186	0\\
1.188	0\\
1.19	0\\
1.192	0\\
1.194	0\\
1.196	0\\
1.198	0\\
1.2	0\\
1.202	0\\
1.204	0\\
1.206	0\\
1.208	0\\
1.21	0\\
1.212	0\\
1.214	0\\
1.216	0\\
1.218	0\\
1.22	0\\
1.222	0\\
1.224	0\\
1.226	0\\
1.228	0\\
1.23	0\\
1.232	0\\
1.234	0\\
1.236	0\\
1.238	0\\
1.24	0\\
1.242	0\\
1.244	0\\
1.246	0\\
1.248	0\\
1.25	0\\
1.252	0\\
1.254	0\\
1.256	0\\
1.258	0\\
1.26	0\\
1.262	0\\
1.264	0\\
1.266	0\\
1.268	0\\
1.27	0\\
1.272	0\\
1.274	0\\
1.276	0\\
1.278	0\\
1.28	0\\
1.282	0\\
1.284	0\\
1.286	0\\
1.288	0\\
1.29	0\\
1.292	0\\
1.294	0\\
1.296	0\\
1.298	0\\
1.3	0\\
1.302	0\\
1.304	0\\
1.306	0\\
1.308	0\\
1.31	0\\
1.312	0\\
1.314	0\\
1.316	0\\
1.318	0\\
1.32	0\\
1.322	0\\
1.324	0\\
1.326	0\\
1.328	0\\
1.33	0\\
1.332	0\\
1.334	0\\
1.336	0\\
1.338	0\\
1.34	0\\
1.342	0\\
1.344	0\\
1.346	0\\
1.348	0\\
1.35	0\\
1.352	0\\
1.354	0\\
1.356	0\\
1.358	0\\
1.36	0\\
1.362	0\\
1.364	0\\
1.366	0\\
1.368	0\\
1.37	0\\
1.372	0\\
1.374	0\\
1.376	0\\
1.378	0\\
1.38	0\\
1.382	0\\
1.384	0\\
1.386	0\\
1.388	0\\
1.39	0\\
1.392	0\\
1.394	0\\
1.396	0\\
1.398	0\\
1.4	0\\
1.402	0\\
1.404	0\\
1.406	0\\
1.408	0\\
1.41	0\\
1.412	0\\
1.414	0\\
1.416	0\\
1.418	0\\
1.42	0\\
1.422	0\\
1.424	0\\
1.426	0\\
1.428	0\\
1.43	0\\
1.432	0\\
1.434	0\\
1.436	0\\
1.438	0\\
1.44	0\\
1.442	0\\
1.444	0\\
1.446	0\\
1.448	0\\
1.45	0\\
1.452	0\\
1.454	0\\
1.456	0\\
1.458	0\\
1.46	0\\
1.462	0\\
1.464	0\\
1.466	0\\
1.468	0\\
1.47	0\\
1.472	0\\
1.474	0\\
1.476	0\\
1.478	0\\
1.48	0\\
1.482	0\\
1.484	0\\
1.486	0\\
1.488	0\\
1.49	0\\
1.492	0\\
1.494	0\\
1.496	0\\
1.498	0\\
1.5	0\\
1.502	0\\
1.504	0\\
1.506	0\\
1.508	0\\
1.51	0\\
1.512	0\\
1.514	0\\
1.516	0\\
1.518	0\\
1.52	0\\
1.522	0\\
1.524	0\\
1.526	0\\
1.528	0\\
1.53	0\\
1.532	0\\
1.534	0\\
1.536	0\\
1.538	0\\
1.54	0\\
1.542	0\\
1.544	0\\
1.546	0\\
1.548	0\\
1.55	0\\
1.552	0\\
1.554	0\\
1.556	0\\
1.558	0\\
1.56	0\\
1.562	0\\
1.564	0\\
1.566	0\\
1.568	0\\
1.57	0\\
1.572	0\\
1.574	0\\
1.576	0\\
1.578	0\\
1.58	0\\
1.582	0\\
1.584	0\\
1.586	0\\
1.588	0\\
1.59	0\\
1.592	0\\
1.594	0\\
1.596	0\\
1.598	0\\
1.6	0\\
1.602	0\\
1.604	0\\
1.606	0\\
1.608	0\\
1.61	0\\
1.612	0\\
1.614	0\\
1.616	0\\
1.618	0\\
1.62	0\\
1.622	0\\
1.624	0\\
1.626	0\\
1.628	0\\
1.63	0\\
1.632	0\\
1.634	0\\
1.636	0\\
1.638	0\\
1.64	0\\
1.642	0\\
1.644	0\\
1.646	0\\
1.648	0\\
1.65	0\\
1.652	0\\
1.654	0\\
1.656	0\\
1.658	0\\
1.66	0\\
1.662	0\\
1.664	0\\
1.666	0\\
1.668	0\\
1.67	0\\
1.672	0\\
1.674	0\\
1.676	0\\
1.678	0\\
1.68	0\\
1.682	0\\
1.684	0\\
1.686	0\\
1.688	0\\
1.69	0\\
1.692	0\\
1.694	0\\
1.696	0\\
1.698	0\\
1.7	0\\
1.702	0\\
1.704	0\\
1.706	0\\
1.708	0\\
1.71	0\\
1.712	0\\
1.714	0\\
1.716	0\\
1.718	0\\
1.72	0\\
1.722	0\\
1.724	0\\
1.726	0\\
1.728	0\\
1.73	0\\
1.732	0\\
1.734	0\\
1.736	0\\
1.738	0\\
1.74	0\\
1.742	0\\
1.744	0\\
1.746	0\\
1.748	0\\
1.75	0\\
1.752	0\\
1.754	0\\
1.756	0\\
1.758	0\\
1.76	0\\
1.762	0\\
1.764	0\\
1.766	0\\
1.768	0\\
1.77	0\\
1.772	0\\
1.774	0\\
1.776	0\\
1.778	0\\
1.78	0\\
1.782	0\\
1.784	0\\
1.786	0\\
1.788	0\\
1.79	0\\
1.792	0\\
1.794	0\\
1.796	0\\
1.798	0\\
1.8	0\\
1.802	0\\
1.804	0\\
1.806	0\\
1.808	0\\
1.81	0\\
1.812	0\\
1.814	0\\
1.816	0\\
1.818	0\\
1.82	0\\
1.822	0\\
1.824	0\\
1.826	0\\
1.828	0\\
1.83	0\\
1.832	0\\
1.834	0\\
1.836	0\\
1.838	0\\
1.84	0\\
1.842	0\\
1.844	0\\
1.846	0\\
1.848	0\\
1.85	0\\
1.852	0\\
1.854	0\\
1.856	0\\
1.858	0\\
1.86	0\\
1.862	0\\
1.864	0\\
1.866	0\\
1.868	0\\
1.87	0\\
1.872	0\\
1.874	0\\
1.876	0\\
1.878	0\\
1.88	0\\
1.882	0\\
1.884	0\\
1.886	0\\
1.888	0\\
1.89	0\\
1.892	0\\
1.894	0\\
1.896	0\\
1.898	0\\
1.9	0\\
1.902	0\\
1.904	0\\
1.906	0\\
1.908	0\\
1.91	0\\
1.912	0\\
1.914	0\\
1.916	0\\
1.918	0\\
1.92	0\\
1.922	0\\
1.924	0\\
1.926	0\\
1.928	0\\
1.93	0\\
1.932	0\\
1.934	0\\
1.936	0\\
1.938	0\\
1.94	0\\
1.942	0\\
1.944	0\\
1.946	0\\
1.948	0\\
1.95	0\\
1.952	0\\
1.954	0\\
1.956	0\\
1.958	0\\
1.96	0\\
1.962	0\\
1.964	0\\
1.966	0\\
1.968	0\\
1.97	0\\
1.972	0\\
1.974	0\\
1.976	0\\
1.978	0\\
1.98	0\\
1.982	0\\
1.984	0\\
1.986	0\\
1.988	0\\
1.99	0\\
1.992	0\\
1.994	0\\
1.996	0\\
1.998	0\\
2	0\\
};
\addplot [color=mycolor2, line width=1.0pt, forget plot]
  table[row sep=crcr]{%
-0.024	0\\
-0.022	0\\
-0.02	0\\
-0.018	0\\
-0.016	0\\
-0.014	0\\
-0.012	0\\
-0.01	0\\
-0.008	0\\
-0.006	0\\
-0.004	0\\
-0.002	0\\
0	0\\
0.002	-4.46202e-10\\
0.004	-6.97982e-09\\
0.006	-3.46486e-08\\
0.008	-1.0768e-07\\
0.01	-2.59176e-07\\
0.012	-5.31092e-07\\
0.014	-9.74378e-07\\
0.016	-1.6492e-06\\
0.018	-2.62514e-06\\
0.02	-3.98133e-06\\
0.022	-5.80646e-06\\
0.024	-8.19855e-06\\
0.026	-1.12646e-05\\
0.028	-1.51198e-05\\
0.03	-1.98869e-05\\
0.032	-2.56949e-05\\
0.034	-3.2678e-05\\
0.036	-4.09738e-05\\
0.038	-5.07216e-05\\
0.04	-1.22914e-05\\
0.042	-1.5447e-05\\
0.044	-1.92179e-05\\
0.046	-2.36862e-05\\
0.048	-2.89394e-05\\
0.05	-3.50697e-05\\
0.052	-4.21735e-05\\
0.054	-5.03507e-05\\
0.056	-5.9704e-05\\
0.058	-7.03384e-05\\
0.06	-8.23596e-05\\
0.062	-9.58738e-05\\
0.064	-0.000110986\\
0.066	-0.0001278\\
0.068	-0.000146414\\
0.07	-0.000166926\\
0.072	-0.000189425\\
0.074	-0.000213995\\
0.076	-0.000240712\\
0.078	-0.00648947\\
0.08	-0.0155021\\
0.082	-0.0161812\\
0.084	-0.0167181\\
0.086	-0.0172292\\
0.088	-0.0177121\\
0.09	-0.0181649\\
0.092	-0.0185855\\
0.094	-0.0189721\\
0.096	-0.019323\\
0.098	-0.0196366\\
0.1	-0.0199115\\
0.102	-0.0189174\\
0.104	-0.0165829\\
0.106	-0.0142613\\
0.108	-0.0119581\\
0.11	-0.00967761\\
0.112	-0.00742335\\
0.114	-0.00519768\\
0.116	-0.00300214\\
0.118	-0.00083745\\
0.12	-0.000336405\\
0.122	-0.000364999\\
0.124	-0.000395262\\
0.126	-0.000427228\\
0.128	-0.000460928\\
0.13	-0.000496384\\
0.132	-0.000533618\\
0.134	-0.000572644\\
0.136	-0.000613469\\
0.138	-0.000656096\\
0.14	-0.000700523\\
0.142	-0.000746737\\
0.144	-0.000794723\\
0.146	-0.000844455\\
0.148	-0.000895902\\
0.15	-0.000949026\\
0.152	-0.00100378\\
0.154	-0.00106011\\
0.156	-0.00111795\\
0.158	-0.00117723\\
0.16	-0.00123788\\
0.162	-0.00129982\\
0.164	-0.00136293\\
0.166	-0.00142714\\
0.168	-0.00149232\\
0.17	-0.00155837\\
0.172	-0.00162515\\
0.174	-0.00169254\\
0.176	-0.0017604\\
0.178	-0.00182859\\
0.18	-0.00189696\\
0.182	-0.00196535\\
0.184	-0.0020336\\
0.186	-0.00210155\\
0.188	-0.00216902\\
0.19	-0.00223584\\
0.192	-0.00230183\\
0.194	-0.0023668\\
0.196	-0.00243058\\
0.198	-0.00249297\\
0.2	-0.00255378\\
0.202	-0.00261282\\
0.204	-0.0026699\\
0.206	-0.00272483\\
0.208	-0.00277741\\
0.21	-0.00282746\\
0.212	-0.00287478\\
0.214	-0.00291919\\
0.216	-0.0029605\\
0.218	-0.00299853\\
0.22	-0.0030331\\
0.222	-0.00306404\\
0.224	-0.00309118\\
0.226	-0.00311436\\
0.228	-0.00313341\\
0.23	-0.0031482\\
0.232	-0.00220568\\
0.234	-0.00104622\\
0.236	-0.000639275\\
0.238	-0.000531361\\
0.24	-0.000422772\\
0.242	-0.000313732\\
0.244	-0.000204465\\
0.246	-9.51966e-05\\
0.248	-0.00102387\\
0.25	-0.0022839\\
0.252	-0.00343689\\
0.254	-0.00448311\\
0.256	-0.00542299\\
0.258	-0.00625713\\
0.26	-0.00698628\\
0.262	-0.00761141\\
0.264	-0.00813368\\
0.266	-0.00855445\\
0.268	-0.00834408\\
0.27	-0.00789183\\
0.272	-0.00736048\\
0.274	-0.00675386\\
0.276	-0.00607587\\
0.278	-0.00533053\\
0.28	-0.00452191\\
0.282	-0.00365416\\
0.284	-0.00273148\\
0.286	-0.00175811\\
0.288	-0.000738358\\
0.29	-0.000461065\\
0.292	-0.000424501\\
0.294	8.1432e-05\\
0.296	0.000820695\\
0.298	0.000672461\\
0.3	0.000521334\\
0.302	0.00036755\\
0.304	0.000211357\\
0.306	5.30075e-05\\
0.308	-0.000107238\\
0.31	-0.000269114\\
0.312	-0.000432346\\
0.314	-0.00059666\\
0.316	-0.000761772\\
0.318	-0.000927399\\
0.32	-0.00109325\\
0.322	-0.00125904\\
0.324	-0.00142446\\
0.326	-0.00158924\\
0.328	-0.00175306\\
0.33	-0.00191565\\
0.332	-0.00207669\\
0.334	-0.00223591\\
0.336	-0.00239301\\
0.338	-0.0025477\\
0.34	-0.0026997\\
0.342	-0.00284873\\
0.344	-0.00299451\\
0.346	-0.00313677\\
0.348	-0.00327525\\
0.35	-0.00340969\\
0.352	-0.00353983\\
0.354	-0.00426323\\
0.356	-0.00516721\\
0.358	-0.00603749\\
0.36	-0.00687244\\
0.362	-0.00767061\\
0.364	-0.00843067\\
0.366	-0.00915141\\
0.368	-0.00983179\\
0.37	-0.0104709\\
0.372	-0.0110679\\
0.374	-0.0116222\\
0.376	-0.0121332\\
0.378	-0.0126006\\
0.38	-0.013024\\
0.382	-0.0134034\\
0.384	-0.0137387\\
0.386	-0.01403\\
0.388	-0.0142775\\
0.39	-0.0144815\\
0.392	-0.0146425\\
0.394	-0.0147609\\
0.396	-0.0148375\\
0.398	-0.0148196\\
0.4	-0.0145852\\
0.402	-0.0143196\\
0.404	-0.0140235\\
0.406	-0.0136981\\
0.408	-0.0133443\\
0.41	-0.0129633\\
0.412	-0.0125561\\
0.414	-0.0121239\\
0.416	-0.011668\\
0.418	-0.0111849\\
0.42	-0.0102024\\
0.422	-0.00921111\\
0.424	-0.00821346\\
0.426	-0.00721178\\
0.428	-0.00617645\\
0.43	-0.00498119\\
0.432	-0.00378935\\
0.434	-0.00260365\\
0.436	-0.00142675\\
0.438	-0.000261228\\
0.44	0.000958582\\
0.442	0.00213587\\
0.444	0.00256647\\
0.446	0.00265608\\
0.448	0.00250252\\
0.45	0.00234625\\
0.452	0.00218757\\
0.454	0.00202676\\
0.456	0.00186411\\
0.458	0.00169991\\
0.46	0.00153446\\
0.462	0.00136804\\
0.464	0.00120096\\
0.466	0.00103351\\
0.468	0.000525333\\
0.47	1.96792e-05\\
0.472	-0.000472833\\
0.474	-0.00095106\\
0.476	-0.00141392\\
0.478	-0.00186039\\
0.48	-0.0022895\\
0.482	-0.00270037\\
0.484	-0.00309216\\
0.486	-0.00346409\\
0.488	-0.00381548\\
0.49	-0.00414568\\
0.492	-0.00445413\\
0.494	-0.00474032\\
0.496	-0.00500383\\
0.498	-0.00524428\\
0.5	-0.00546139\\
0.502	-0.00565492\\
0.504	-0.0058247\\
0.506	-0.00597064\\
0.508	-0.0060927\\
0.51	-0.0061909\\
0.512	-0.00626534\\
0.514	-0.00631617\\
0.516	-0.00634361\\
0.518	-0.00634791\\
0.52	-0.00632942\\
0.522	-0.0062885\\
0.524	-0.00622559\\
0.526	-0.00614119\\
0.528	-0.00603582\\
0.53	-0.00591006\\
0.532	-0.00576455\\
0.534	-0.00559995\\
0.536	-0.00541697\\
0.538	-0.00521636\\
0.54	-0.0049989\\
0.542	-0.0047654\\
0.544	-0.00451672\\
0.546	-0.00425372\\
0.548	-0.00397731\\
0.55	-0.00368841\\
0.552	-0.00338796\\
0.554	-0.00307692\\
0.556	-0.00304092\\
0.558	-0.00300653\\
0.56	-0.00296665\\
0.562	-0.0029214\\
0.564	-0.00287093\\
0.566	-0.00281537\\
0.568	-0.00275489\\
0.57	-0.00268964\\
0.572	-0.0026198\\
0.574	-0.00254555\\
0.576	-0.00246708\\
0.578	-0.00238458\\
0.58	-0.00229825\\
0.582	-0.00194577\\
0.584	-0.0013496\\
0.586	-0.000762222\\
0.588	-0.000184515\\
0.59	0.000121955\\
0.592	-0.00034178\\
0.594	-0.000799965\\
0.596	-0.00125164\\
0.598	-0.00169588\\
0.6	-0.00213177\\
0.602	-0.00255843\\
0.604	-0.00297501\\
0.606	-0.0033807\\
0.608	-0.0037747\\
0.61	-0.00415626\\
0.612	-0.00452468\\
0.614	-0.00487926\\
0.616	-0.00521936\\
0.618	-0.00554439\\
0.62	-0.00585377\\
0.622	-0.00614699\\
0.624	-0.00642357\\
0.626	-0.00668306\\
0.628	-0.00692508\\
0.63	-0.00714926\\
0.632	-0.00735529\\
0.634	-0.0074765\\
0.636	-0.00749351\\
0.638	-0.00749374\\
0.64	-0.00747727\\
0.642	-0.0074442\\
0.644	-0.00739472\\
0.646	-0.00732902\\
0.648	-0.00724732\\
0.65	-0.00714991\\
0.652	-0.00703709\\
0.654	-0.0069092\\
0.656	-0.00676661\\
0.658	-0.00660973\\
0.66	-0.006439\\
0.662	-0.00625487\\
0.664	-0.00605784\\
0.666	-0.00584842\\
0.668	-0.00562717\\
0.67	-0.00539464\\
0.672	-0.00515143\\
0.674	-0.00489813\\
0.676	-0.00463538\\
0.678	-0.00436383\\
0.68	-0.00408412\\
0.682	-0.00379693\\
0.684	-0.00350295\\
0.686	-0.00320287\\
0.688	-0.0028974\\
0.69	-0.00258724\\
0.692	-0.00227311\\
0.694	-0.00210418\\
0.696	-0.00208741\\
0.698	-0.00206663\\
0.7	-0.00204194\\
0.702	-0.00201342\\
0.704	-0.00198119\\
0.706	-0.00194534\\
0.708	-0.00190599\\
0.71	-0.00186327\\
0.712	-0.0018173\\
0.714	-0.0017682\\
0.716	-0.00171613\\
0.718	-0.00166121\\
0.72	-0.0016036\\
0.722	-0.00154344\\
0.724	-0.00148089\\
0.726	-0.00141611\\
0.728	-0.00134925\\
0.73	-0.00128047\\
0.732	-0.00120995\\
0.734	-0.00124611\\
0.736	-0.00171567\\
0.738	-0.00200662\\
0.74	-0.00228998\\
0.742	-0.00256522\\
0.744	-0.00283183\\
0.746	-0.00308933\\
0.748	-0.00333723\\
0.75	-0.00357512\\
0.752	-0.00380259\\
0.754	-0.00401924\\
0.756	-0.00422473\\
0.758	-0.00441874\\
0.76	-0.00460098\\
0.762	-0.00477117\\
0.764	-0.00492908\\
0.766	-0.00507451\\
0.768	-0.00520728\\
0.77	-0.00532725\\
0.772	-0.0054343\\
0.774	-0.00552835\\
0.776	-0.00560934\\
0.778	-0.00567725\\
0.78	-0.00573207\\
0.782	-0.00577386\\
0.784	-0.00580265\\
0.786	-0.00581855\\
0.788	-0.00582167\\
0.79	-0.00581215\\
0.792	-0.00570311\\
0.794	-0.00557654\\
0.796	-0.00543963\\
0.798	-0.00529276\\
0.8	-0.00513634\\
0.802	-0.00497077\\
0.804	-0.00479649\\
0.806	-0.00461396\\
0.808	-0.00442363\\
0.81	-0.00422598\\
0.812	-0.00402149\\
0.814	-0.00381066\\
0.816	-0.003594\\
0.818	-0.00337201\\
0.82	-0.00314522\\
0.822	-0.00291414\\
0.824	-0.00267931\\
0.826	-0.00244127\\
0.828	-0.00220053\\
0.83	-0.00195763\\
0.832	-0.00171311\\
0.834	-0.00146749\\
0.836	-0.00129877\\
0.838	-0.00128371\\
0.84	-0.0012661\\
0.842	-0.00124602\\
0.844	-0.00122354\\
0.846	-0.00119875\\
0.848	-0.00117171\\
0.85	-0.00114253\\
0.852	-0.00111128\\
0.854	-0.00107806\\
0.856	-0.00104297\\
0.858	-0.00100611\\
0.86	-0.000967564\\
0.862	-0.00092745\\
0.864	-0.000885871\\
0.866	-0.000842932\\
0.868	-0.000798743\\
0.87	-0.000753413\\
0.872	-0.000707053\\
0.874	-0.000659775\\
0.876	-0.000685601\\
0.878	-0.000951566\\
0.88	-0.00121056\\
0.882	-0.00146218\\
0.884	-0.00170603\\
0.886	-0.00194175\\
0.888	-0.00216899\\
0.89	-0.00238743\\
0.892	-0.00259678\\
0.894	-0.00279678\\
0.896	-0.00298716\\
0.898	-0.00313653\\
0.9	-0.00323891\\
0.902	-0.00333298\\
0.904	-0.00341864\\
0.906	-0.00349582\\
0.908	-0.00356445\\
0.91	-0.00362451\\
0.912	-0.00367596\\
0.914	-0.00371881\\
0.916	-0.00375308\\
0.918	-0.00377882\\
0.92	-0.00379606\\
0.922	-0.00380489\\
0.924	-0.00380541\\
0.926	-0.00379771\\
0.928	-0.00378192\\
0.93	-0.00375819\\
0.932	-0.00372668\\
0.934	-0.00368754\\
0.936	-0.00364098\\
0.938	-0.00358718\\
0.94	-0.00352638\\
0.942	-0.0034488\\
0.944	-0.00332655\\
0.946	-0.00319908\\
0.948	-0.00306672\\
0.95	-0.0029298\\
0.952	-0.00278866\\
0.954	-0.00264365\\
0.956	-0.0024951\\
0.958	-0.00234338\\
0.96	-0.00218883\\
0.962	-0.0020318\\
0.964	-0.00187266\\
0.966	-0.00171175\\
0.968	-0.00154942\\
0.97	-0.00138604\\
0.972	-0.00122194\\
0.974	-0.00105747\\
0.976	-0.000892974\\
0.978	-0.000814722\\
0.98	-0.000801206\\
0.982	-0.000786173\\
0.984	-0.000769675\\
0.986	-0.000751767\\
0.988	-0.000732503\\
0.99	-0.000711944\\
0.992	-0.000690149\\
0.994	-0.00066718\\
0.996	-0.000643102\\
0.998	-0.00061798\\
1	-0.000591881\\
1.002	-0.000564873\\
1.004	-0.000537025\\
1.006	-0.000508406\\
1.008	-0.000479088\\
1.01	-0.000449142\\
1.012	-0.000418638\\
1.014	-0.000387649\\
1.016	-0.00040562\\
1.018	-0.000541376\\
1.02	-0.000673335\\
1.022	-0.000801284\\
1.024	-0.000925024\\
1.026	-0.00104436\\
1.028	-0.00115913\\
1.03	-0.00126916\\
1.032	-0.00137431\\
1.034	-0.00147444\\
1.036	-0.00156942\\
1.038	-0.00165915\\
1.04	-0.00174353\\
1.042	-0.00182249\\
1.044	-0.00189594\\
1.046	-0.00196383\\
1.048	-0.00202613\\
1.05	-0.0020828\\
1.052	-0.00213381\\
1.054	-0.00217918\\
1.056	-0.0022189\\
1.058	-0.00225299\\
1.06	-0.0022815\\
1.062	-0.00230444\\
1.064	-0.0023219\\
1.066	-0.00233391\\
1.068	-0.0023291\\
1.07	-0.00230544\\
1.072	-0.00227717\\
1.074	-0.00224443\\
1.076	-0.00220735\\
1.078	-0.00216607\\
1.08	-0.00212075\\
1.082	-0.00207153\\
1.084	-0.00201858\\
1.086	-0.00196208\\
1.088	-0.00190219\\
1.09	-0.00183909\\
1.092	-0.00177296\\
1.094	-0.001704\\
1.096	-0.00162692\\
1.098	-0.00153347\\
1.1	-0.00143829\\
1.102	-0.00134159\\
1.104	-0.00124359\\
1.106	-0.0011445\\
1.108	-0.00104456\\
1.11	-0.00094396\\
1.112	-0.000842927\\
1.114	-0.00074167\\
1.116	-0.000640398\\
1.118	-0.000539317\\
1.12	-0.00050372\\
1.122	-0.00049341\\
1.124	-0.000482217\\
1.126	-0.000470177\\
1.128	-0.000457325\\
1.13	-0.0004437\\
1.132	-0.00042934\\
1.134	-0.000414286\\
1.136	-0.000398579\\
1.138	-0.000382259\\
1.14	-0.000365369\\
1.142	-0.000347953\\
1.144	-0.000330053\\
1.146	-0.000311713\\
1.148	-0.000292977\\
1.15	-0.00027389\\
1.152	-0.000254495\\
1.154	-0.000277956\\
1.156	-0.000346748\\
1.158	-0.000413441\\
1.16	-0.000477923\\
1.162	-0.000540092\\
1.164	-0.000599852\\
1.166	-0.000657115\\
1.168	-0.000711798\\
1.17	-0.000763828\\
1.172	-0.000813136\\
1.174	-0.000859664\\
1.176	-0.000903359\\
1.178	-0.000944174\\
1.18	-0.000982073\\
1.182	-0.00101702\\
1.184	-0.001049\\
1.186	-0.001078\\
1.188	-0.00110399\\
1.19	-0.00112698\\
1.192	-0.00114698\\
1.194	-0.00116399\\
1.196	-0.00117804\\
1.198	-0.00118913\\
1.2	-0.00119731\\
1.202	-0.00120261\\
1.204	-0.00120507\\
1.206	-0.00120474\\
1.208	-0.00120167\\
1.21	-0.00119591\\
1.212	-0.00118754\\
1.214	-0.00117661\\
1.216	-0.0011632\\
1.218	-0.00114739\\
1.22	-0.00112925\\
1.222	-0.00110887\\
1.224	-0.00108634\\
1.226	-0.00106174\\
1.228	-0.00103518\\
1.23	-0.00100674\\
1.232	-0.000976524\\
1.234	-0.000944636\\
1.236	-0.000904831\\
1.238	-0.000861937\\
1.24	-0.00081787\\
1.242	-0.00077165\\
1.244	-0.000722226\\
1.246	-0.000672014\\
1.248	-0.000621134\\
1.25	-0.000569706\\
1.252	-0.00051785\\
1.254	-0.000465683\\
1.256	-0.000412335\\
1.258	-0.000353345\\
1.26	-0.000311853\\
1.262	-0.000305355\\
1.264	-0.000298324\\
1.266	-0.00029078\\
1.268	-0.000282748\\
1.27	-0.00027425\\
1.272	-0.00026531\\
1.274	-0.000255953\\
1.276	-0.000246204\\
1.278	-0.000236089\\
1.28	-0.000225633\\
1.282	-0.000214862\\
1.284	-0.000203804\\
1.286	-0.000192485\\
1.288	-0.000180932\\
1.29	-0.000199669\\
1.292	-0.000235237\\
1.294	-0.000269623\\
1.296	-0.000302768\\
1.298	-0.00033462\\
1.3	-0.000365128\\
1.302	-0.000394246\\
1.304	-0.000421933\\
1.306	-0.000448152\\
1.308	-0.000472867\\
1.31	-0.000496049\\
1.312	-0.000517673\\
1.314	-0.000537717\\
1.316	-0.000556161\\
1.318	-0.000572994\\
1.32	-0.000588203\\
1.322	-0.000601784\\
1.324	-0.000613734\\
1.326	-0.000624053\\
1.328	-0.000632747\\
1.33	-0.000639825\\
1.332	-0.000645298\\
1.334	-0.000649182\\
1.336	-0.000651496\\
1.338	-0.000652261\\
1.34	-0.000651504\\
1.342	-0.000649252\\
1.344	-0.000645536\\
1.346	-0.00064039\\
1.348	-0.00063385\\
1.35	-0.000625957\\
1.352	-0.00061675\\
1.354	-0.000606275\\
1.356	-0.000594576\\
1.358	-0.000581701\\
1.36	-0.000567701\\
1.362	-0.000552626\\
1.364	-0.000536529\\
1.366	-0.000519465\\
1.368	-0.00050149\\
1.37	-0.00048266\\
1.372	-0.000463032\\
1.374	-0.000442666\\
1.376	-0.000421621\\
1.378	-0.000399955\\
1.38	-0.00037773\\
1.382	-0.000355006\\
1.384	-0.000331843\\
1.386	-0.000308301\\
1.388	-0.000284442\\
1.39	-0.000260324\\
1.392	-0.000236008\\
1.394	-0.000211552\\
1.396	-0.000187015\\
1.398	-0.00016223\\
1.4	-0.000152935\\
1.402	-0.000153159\\
1.404	-0.000151432\\
1.406	-0.000149379\\
1.408	-0.000147013\\
1.41	-0.000144345\\
1.412	-0.000141388\\
1.414	-0.000138155\\
1.416	-0.000134658\\
1.418	-0.000130912\\
1.42	-0.00012693\\
1.422	-0.000122726\\
1.424	-0.000129108\\
1.426	-0.000150861\\
1.428	-0.0001719\\
1.43	-0.000191536\\
1.432	-0.000208967\\
1.434	-0.000225635\\
1.436	-0.000241512\\
1.438	-0.000256577\\
1.44	-0.000270807\\
1.442	-0.000284184\\
1.444	-0.000296691\\
1.446	-0.000308314\\
1.448	-0.000319041\\
1.45	-0.000328861\\
1.452	-0.000337769\\
1.454	-0.000345756\\
1.456	-0.000352822\\
1.458	-0.000358964\\
1.46	-0.000364183\\
1.462	-0.000368483\\
1.464	-0.000371868\\
1.466	-0.000374345\\
1.468	-0.000375925\\
1.47	-0.000376616\\
1.472	-0.000376433\\
1.474	-0.00037539\\
1.476	-0.000373502\\
1.478	-0.000370789\\
1.48	-0.000367268\\
1.482	-0.000362962\\
1.484	-0.000357893\\
1.486	-0.000352084\\
1.488	-0.00034556\\
1.49	-0.000338348\\
1.492	-0.000330474\\
1.494	-0.000321969\\
1.496	-0.00031286\\
1.498	-0.000303179\\
1.5	-0.000292955\\
1.502	-0.000282222\\
1.504	-0.000271012\\
1.506	-0.000259357\\
1.508	-0.000247292\\
1.51	-0.000234849\\
1.512	-0.000222063\\
1.514	-0.000208969\\
1.516	-0.000195602\\
1.518	-0.000181995\\
1.52	-0.000168183\\
1.522	-0.000154202\\
1.524	-0.000140084\\
1.526	-0.000125865\\
1.528	-0.000111578\\
1.53	-9.72567e-05\\
1.532	-8.29336e-05\\
1.534	-6.86414e-05\\
1.536	-5.44118e-05\\
1.538	-5.13432e-05\\
1.54	-5.29434e-05\\
1.542	-5.43637e-05\\
1.544	-5.56056e-05\\
1.546	-5.66712e-05\\
1.548	-5.75631e-05\\
1.55	-5.82841e-05\\
1.552	-5.88372e-05\\
1.554	-5.9226e-05\\
1.556	-5.94541e-05\\
1.558	-6.49623e-05\\
1.56	-7.76467e-05\\
1.562	-8.96748e-05\\
1.564	-0.000101302\\
1.566	-0.000112509\\
1.568	-0.000123279\\
1.57	-0.000133594\\
1.572	-0.000143439\\
1.574	-0.000152801\\
1.576	-0.000161667\\
1.578	-0.000170024\\
1.58	-0.000177865\\
1.582	-0.000185178\\
1.584	-0.000191958\\
1.586	-0.000198198\\
1.588	-0.000203892\\
1.59	-0.000209038\\
1.592	-0.000213633\\
1.594	-0.000217677\\
1.596	-0.000221168\\
1.598	-0.000224109\\
1.6	-0.000226202\\
1.602	-0.000226958\\
1.604	-0.000227195\\
1.606	-0.000226922\\
1.608	-0.000226146\\
1.61	-0.000224876\\
1.612	-0.000223123\\
1.614	-0.000220899\\
1.616	-0.000218216\\
1.618	-0.000215086\\
1.62	-0.000211523\\
1.622	-0.000207543\\
1.624	-0.000203159\\
1.626	-0.000198389\\
1.628	-0.000193248\\
1.63	-0.000187754\\
1.632	-0.000181925\\
1.634	-0.000175777\\
1.636	-0.000169331\\
1.638	-0.000162605\\
1.64	-0.000155617\\
1.642	-0.000148388\\
1.644	-0.000140937\\
1.646	-0.000133284\\
1.648	-0.000125449\\
1.65	-0.000117453\\
1.652	-0.000109315\\
1.654	-0.000101055\\
1.656	-9.26939e-05\\
1.658	-8.42516e-05\\
1.66	-7.5748e-05\\
1.662	-6.72026e-05\\
1.664	-5.86352e-05\\
1.666	-5.00648e-05\\
1.668	-4.15105e-05\\
1.67	-3.29909e-05\\
1.672	-2.45242e-05\\
1.674	-1.97624e-05\\
1.676	-2.00069e-05\\
1.678	-2.0178e-05\\
1.68	-2.0277e-05\\
1.682	-2.03054e-05\\
1.684	-2.02648e-05\\
1.686	-2.01569e-05\\
1.688	-1.99836e-05\\
1.69	-1.97471e-05\\
1.692	-2.01203e-05\\
1.694	-2.84797e-05\\
1.696	-3.66244e-05\\
1.698	-4.45407e-05\\
1.7	-5.22154e-05\\
1.702	-5.96363e-05\\
1.704	-6.67917e-05\\
1.706	-7.36709e-05\\
1.708	-8.02639e-05\\
1.71	-8.65615e-05\\
1.712	-9.25554e-05\\
1.714	-9.82379e-05\\
1.716	-0.000103602\\
1.718	-0.000108643\\
1.72	-0.000112908\\
1.722	-0.000116304\\
1.724	-0.000119388\\
1.726	-0.000122158\\
1.728	-0.000124613\\
1.73	-0.000126752\\
1.732	-0.000128577\\
1.734	-0.000130087\\
1.736	-0.000131285\\
1.738	-0.000132173\\
1.74	-0.000132754\\
1.742	-0.000133033\\
1.744	-0.000133012\\
1.746	-0.000132697\\
1.748	-0.000132094\\
1.75	-0.000131208\\
1.752	-0.000130046\\
1.754	-0.000128616\\
1.756	-0.000126923\\
1.758	-0.000124977\\
1.76	-0.000122786\\
1.762	-0.000119879\\
1.764	-0.000116713\\
1.766	-0.000113341\\
1.768	-0.000109773\\
1.77	-0.000106019\\
1.772	-0.000102092\\
1.774	-9.80021e-05\\
1.776	-9.37603e-05\\
1.778	-8.93783e-05\\
1.78	-8.48678e-05\\
1.782	-8.02403e-05\\
1.784	-7.55079e-05\\
1.786	-7.06821e-05\\
1.788	-6.5775e-05\\
1.79	-6.07983e-05\\
1.792	-5.57639e-05\\
1.794	-5.06835e-05\\
1.796	-4.55688e-05\\
1.798	-4.04314e-05\\
1.8	-3.52827e-05\\
1.802	-3.0134e-05\\
1.804	-2.49965e-05\\
1.806	-1.98811e-05\\
1.808	-1.47985e-05\\
1.81	-1.24722e-05\\
1.812	-1.21619e-05\\
1.814	-1.18176e-05\\
1.816	-1.14406e-05\\
1.818	-1.10323e-05\\
1.82	-1.05941e-05\\
1.822	-1.01276e-05\\
1.824	-9.6344e-06\\
1.826	-9.11606e-06\\
1.828	-1.05973e-05\\
1.83	-1.51078e-05\\
1.832	-1.95041e-05\\
1.834	-2.37785e-05\\
1.836	-2.79239e-05\\
1.838	-3.19333e-05\\
1.84	-3.58005e-05\\
1.842	-3.95193e-05\\
1.844	-4.30844e-05\\
1.846	-4.64905e-05\\
1.848	-4.9733e-05\\
1.85	-5.28077e-05\\
1.852	-5.57107e-05\\
1.854	-5.84387e-05\\
1.856	-6.09888e-05\\
1.858	-6.33585e-05\\
1.86	-6.55458e-05\\
1.862	-6.75491e-05\\
1.864	-6.93671e-05\\
1.866	-7.09993e-05\\
1.868	-7.24452e-05\\
1.87	-7.3705e-05\\
1.872	-7.4779e-05\\
1.874	-7.56684e-05\\
1.876	-7.63742e-05\\
1.878	-7.68981e-05\\
1.88	-7.72423e-05\\
1.882	-7.74091e-05\\
1.884	-7.73328e-05\\
1.886	-7.6771e-05\\
1.888	-7.60501e-05\\
1.89	-7.5174e-05\\
1.892	-7.41472e-05\\
1.894	-7.29742e-05\\
1.896	-7.16599e-05\\
1.898	-7.02093e-05\\
1.9	-6.86277e-05\\
1.902	-6.69207e-05\\
1.904	-6.5094e-05\\
1.906	-6.31533e-05\\
1.908	-6.11046e-05\\
1.91	-5.89542e-05\\
1.912	-5.67084e-05\\
1.914	-5.43734e-05\\
1.916	-5.19558e-05\\
1.918	-4.94622e-05\\
1.92	-4.68727e-05\\
1.922	-4.40146e-05\\
1.924	-4.11058e-05\\
1.926	-3.81532e-05\\
1.928	-3.51639e-05\\
1.93	-3.21447e-05\\
1.932	-2.91026e-05\\
1.934	-2.60444e-05\\
1.936	-2.2977e-05\\
1.938	-1.9907e-05\\
1.94	-1.6841e-05\\
1.942	-1.37856e-05\\
1.944	-1.13579e-05\\
1.946	-1.17724e-05\\
1.948	-1.21387e-05\\
1.95	-1.24561e-05\\
1.952	-1.27236e-05\\
1.954	-1.29409e-05\\
1.956	-1.31075e-05\\
1.958	-1.32234e-05\\
1.96	-1.32884e-05\\
1.962	-1.33029e-05\\
1.964	-1.32672e-05\\
1.966	-1.31819e-05\\
1.968	-1.30477e-05\\
1.97	-1.43703e-05\\
1.972	-1.66307e-05\\
1.974	-1.88156e-05\\
1.976	-2.09212e-05\\
1.978	-2.29443e-05\\
1.98	-2.48818e-05\\
1.982	-2.67309e-05\\
1.984	-2.84889e-05\\
1.986	-3.01535e-05\\
1.988	-3.17225e-05\\
1.99	-3.31942e-05\\
1.992	-3.45668e-05\\
1.994	-3.5839e-05\\
1.996	-3.70096e-05\\
1.998	-3.80778e-05\\
2	-3.90429e-05\\
};
\addplot [color=mycolor3, line width=1.0pt, forget plot]
  table[row sep=crcr]{%
-0.024	0\\
-0.022	0\\
-0.02	0\\
-0.018	0\\
-0.016	0\\
-0.014	0\\
-0.012	0\\
-0.01	0\\
-0.008	0\\
-0.006	0\\
-0.004	0\\
-0.002	0\\
0	0\\
0.002	-4.22851e-10\\
0.004	-7.56946e-09\\
0.006	-3.91483e-08\\
0.008	-1.24079e-07\\
0.01	-3.02059e-07\\
0.012	-6.23467e-07\\
0.014	-1.14946e-06\\
0.016	-1.95215e-06\\
0.018	-3.11482e-06\\
0.02	-4.73195e-06\\
0.022	-6.90918e-06\\
0.024	-9.76298e-06\\
0.026	-1.34202e-05\\
0.028	-1.80171e-05\\
0.03	-2.36986e-05\\
0.032	-3.06166e-05\\
0.034	-3.89285e-05\\
0.036	-4.87959e-05\\
0.038	-6.03821e-05\\
0.04	-6.53416e-06\\
0.042	-9.60591e-06\\
0.044	-1.36288e-05\\
0.046	-1.87987e-05\\
0.048	-2.53342e-05\\
0.05	-3.34769e-05\\
0.052	-4.34911e-05\\
0.054	-5.56632e-05\\
0.056	-7.03007e-05\\
0.058	-8.77303e-05\\
0.06	-0.000108297\\
0.062	-0.000132361\\
0.064	-0.000160294\\
0.066	-0.00019248\\
0.068	-0.000229307\\
0.07	-0.000271166\\
0.072	-0.000318448\\
0.074	-0.000371536\\
0.076	-0.000430806\\
0.078	-0.00671645\\
0.08	-0.0157705\\
0.082	-0.0186186\\
0.084	-0.0191678\\
0.086	-0.0196818\\
0.088	-0.0201584\\
0.09	-0.0205957\\
0.092	-0.0209916\\
0.094	-0.0213445\\
0.096	-0.021653\\
0.098	-0.0219156\\
0.1	-0.0221313\\
0.102	-0.0222989\\
0.104	-0.0224177\\
0.106	-0.022487\\
0.108	-0.0225062\\
0.11	-0.0224749\\
0.112	-0.0220061\\
0.114	-0.0203517\\
0.116	-0.0187103\\
0.118	-0.0170803\\
0.12	-0.0170924\\
0.122	-0.0176105\\
0.124	-0.0181047\\
0.126	-0.0185733\\
0.128	-0.0190145\\
0.13	-0.019427\\
0.132	-0.0198093\\
0.134	-0.0201603\\
0.136	-0.0204788\\
0.138	-0.0207639\\
0.14	-0.0210149\\
0.142	-0.021231\\
0.144	-0.0214117\\
0.146	-0.0215567\\
0.148	-0.0216656\\
0.15	-0.0217384\\
0.152	-0.0217749\\
0.154	-0.0217753\\
0.156	-0.0217396\\
0.158	-0.0216683\\
0.16	-0.0215616\\
0.162	-0.02142\\
0.164	-0.0212439\\
0.166	-0.0210341\\
0.168	-0.0207911\\
0.17	-0.0205157\\
0.172	-0.0202087\\
0.174	-0.0198707\\
0.176	-0.0195028\\
0.178	-0.0191058\\
0.18	-0.0186806\\
0.182	-0.0182283\\
0.184	-0.0177497\\
0.186	-0.017246\\
0.188	-0.016718\\
0.19	-0.016167\\
0.192	-0.0155939\\
0.194	-0.0149999\\
0.196	-0.0143859\\
0.198	-0.0137531\\
0.2	-0.0131026\\
0.202	-0.0124354\\
0.204	-0.0117527\\
0.206	-0.0110555\\
0.208	-0.010345\\
0.21	-0.00962212\\
0.212	-0.00888805\\
0.214	-0.00814384\\
0.216	-0.00739055\\
0.218	-0.00662924\\
0.22	-0.00586095\\
0.222	-0.00508673\\
0.224	-0.00430759\\
0.226	-0.00352456\\
0.228	-0.00273864\\
0.23	-0.0019508\\
0.232	-0.00116204\\
0.234	-0.000373318\\
0.236	-0.000298039\\
0.238	-0.000483703\\
0.24	-0.000631778\\
0.242	-0.00074385\\
0.244	-0.000821679\\
0.246	-0.000867191\\
0.248	-0.00192018\\
0.25	-0.00327603\\
0.252	-0.00449844\\
0.254	-0.00538936\\
0.256	-0.00610087\\
0.258	-0.00670692\\
0.26	-0.00720861\\
0.262	-0.00760727\\
0.264	-0.0079044\\
0.266	-0.0081017\\
0.268	-0.00820108\\
0.27	-0.0082047\\
0.272	-0.0081149\\
0.274	-0.00786819\\
0.276	-0.00717184\\
0.278	-0.00640701\\
0.28	-0.00557782\\
0.282	-0.00468842\\
0.284	-0.00374307\\
0.286	-0.00274604\\
0.288	-0.00170166\\
0.29	-0.0013988\\
0.292	-0.00133579\\
0.294	-0.00127125\\
0.296	-0.00158432\\
0.298	-0.00268854\\
0.3	-0.00341153\\
0.302	-0.00412422\\
0.304	-0.00482691\\
0.306	-0.0055198\\
0.308	-0.00620292\\
0.31	-0.00687619\\
0.312	-0.00753943\\
0.314	-0.0081923\\
0.316	-0.00883438\\
0.318	-0.00946514\\
0.32	-0.010084\\
0.322	-0.0106901\\
0.324	-0.0112829\\
0.326	-0.0118613\\
0.328	-0.0124245\\
0.33	-0.0129715\\
0.332	-0.0135012\\
0.334	-0.0140127\\
0.336	-0.0145047\\
0.338	-0.0149763\\
0.34	-0.0154262\\
0.342	-0.0158534\\
0.344	-0.0159413\\
0.346	-0.0156328\\
0.348	-0.015305\\
0.35	-0.014959\\
0.352	-0.014596\\
0.354	-0.0148147\\
0.356	-0.0152038\\
0.358	-0.0155501\\
0.36	-0.0158533\\
0.362	-0.0161131\\
0.364	-0.0163293\\
0.366	-0.0165016\\
0.368	-0.0166303\\
0.37	-0.0167153\\
0.372	-0.0167568\\
0.374	-0.0167552\\
0.376	-0.0167108\\
0.378	-0.0166241\\
0.38	-0.0164957\\
0.382	-0.0163263\\
0.384	-0.0161166\\
0.386	-0.0158675\\
0.388	-0.0155798\\
0.39	-0.0152547\\
0.392	-0.0148931\\
0.394	-0.0144962\\
0.396	-0.0140653\\
0.398	-0.0136016\\
0.4	-0.0131065\\
0.402	-0.0125814\\
0.404	-0.0120277\\
0.406	-0.011447\\
0.408	-0.0108409\\
0.41	-0.0102109\\
0.412	-0.00955876\\
0.414	-0.00888608\\
0.416	-0.00819461\\
0.418	-0.00748609\\
0.42	-0.00676228\\
0.422	-0.00602497\\
0.424	-0.00527597\\
0.426	-0.00451708\\
0.428	-0.00375012\\
0.43	-0.0029769\\
0.432	-0.00219925\\
0.434	-0.00141898\\
0.436	-0.00063788\\
0.438	0.000142256\\
0.44	0.00091966\\
0.442	0.00130731\\
0.444	0.00096548\\
0.446	0.000300855\\
0.448	-0.000587801\\
0.45	-0.0013544\\
0.452	-0.00202418\\
0.454	-0.00267747\\
0.456	-0.00331364\\
0.458	-0.00393208\\
0.46	-0.00453222\\
0.462	-0.0051135\\
0.464	-0.00567542\\
0.466	-0.0062175\\
0.468	-0.00707995\\
0.47	-0.00791939\\
0.472	-0.00872511\\
0.474	-0.00937936\\
0.476	-0.00998712\\
0.478	-0.0105627\\
0.48	-0.0111053\\
0.482	-0.0116139\\
0.484	-0.012088\\
0.486	-0.0125268\\
0.488	-0.0129298\\
0.49	-0.0132965\\
0.492	-0.0136266\\
0.494	-0.0139198\\
0.496	-0.0141758\\
0.498	-0.0143945\\
0.5	-0.0145759\\
0.502	-0.0147201\\
0.504	-0.0148272\\
0.506	-0.0148974\\
0.508	-0.0149309\\
0.51	-0.0149283\\
0.512	-0.0148898\\
0.514	-0.014816\\
0.516	-0.0147076\\
0.518	-0.0145652\\
0.52	-0.0143894\\
0.522	-0.0141812\\
0.524	-0.0139413\\
0.526	-0.0136707\\
0.528	-0.0133703\\
0.53	-0.0130412\\
0.532	-0.0126845\\
0.534	-0.0123012\\
0.536	-0.0118926\\
0.538	-0.0114599\\
0.54	-0.0110042\\
0.542	-0.0105269\\
0.544	-0.0100293\\
0.546	-0.00951282\\
0.548	-0.0089787\\
0.55	-0.00842839\\
0.552	-0.00786328\\
0.554	-0.00728482\\
0.556	-0.00697908\\
0.558	-0.00667311\\
0.56	-0.00636024\\
0.562	-0.00604105\\
0.564	-0.00571613\\
0.566	-0.00538604\\
0.568	-0.00505139\\
0.57	-0.00471275\\
0.572	-0.00437071\\
0.574	-0.00402586\\
0.576	-0.00367878\\
0.578	-0.00333005\\
0.58	-0.00298025\\
0.582	-0.00262995\\
0.584	-0.00227972\\
0.586	-0.00193012\\
0.588	-0.0015817\\
0.59	-0.00149573\\
0.592	-0.00217091\\
0.594	-0.00283125\\
0.596	-0.00347556\\
0.598	-0.00409662\\
0.6	-0.00468412\\
0.602	-0.00525747\\
0.604	-0.00581547\\
0.606	-0.00635698\\
0.608	-0.00688085\\
0.61	-0.007386\\
0.612	-0.0078714\\
0.614	-0.00833605\\
0.616	-0.00870927\\
0.618	-0.00901244\\
0.62	-0.00929343\\
0.622	-0.00955192\\
0.624	-0.00978762\\
0.626	-0.0100003\\
0.628	-0.0101897\\
0.63	-0.0103559\\
0.632	-0.0104985\\
0.634	-0.0106178\\
0.636	-0.0107136\\
0.638	-0.010786\\
0.64	-0.0108352\\
0.642	-0.0108613\\
0.644	-0.0108646\\
0.646	-0.0108452\\
0.648	-0.0108036\\
0.65	-0.0107399\\
0.652	-0.0106547\\
0.654	-0.0105484\\
0.656	-0.0104214\\
0.658	-0.0102743\\
0.66	-0.0101075\\
0.662	-0.00992175\\
0.664	-0.00971755\\
0.666	-0.00949557\\
0.668	-0.0092565\\
0.67	-0.00900104\\
0.672	-0.00872991\\
0.674	-0.00844387\\
0.676	-0.00814369\\
0.678	-0.00783017\\
0.68	-0.00750413\\
0.682	-0.00716639\\
0.684	-0.0068178\\
0.686	-0.00645922\\
0.688	-0.00609151\\
0.69	-0.00571556\\
0.692	-0.00533225\\
0.694	-0.00509092\\
0.696	-0.00499869\\
0.698	-0.00489957\\
0.7	-0.00479383\\
0.702	-0.00468172\\
0.704	-0.00456352\\
0.706	-0.00443951\\
0.708	-0.00430997\\
0.71	-0.00417519\\
0.712	-0.00403546\\
0.714	-0.00389109\\
0.716	-0.00374238\\
0.718	-0.00358962\\
0.72	-0.00343313\\
0.722	-0.00327322\\
0.724	-0.00311019\\
0.726	-0.00294436\\
0.728	-0.00277604\\
0.73	-0.00260554\\
0.732	-0.00240795\\
0.734	-0.00228238\\
0.736	-0.0026302\\
0.738	-0.00296896\\
0.74	-0.0032981\\
0.742	-0.00361711\\
0.744	-0.00392548\\
0.746	-0.00422274\\
0.748	-0.00450843\\
0.75	-0.00478212\\
0.752	-0.00504338\\
0.754	-0.00529183\\
0.756	-0.00552711\\
0.758	-0.00574888\\
0.76	-0.00595682\\
0.762	-0.00615065\\
0.764	-0.00633011\\
0.766	-0.00649498\\
0.768	-0.00664504\\
0.77	-0.00678013\\
0.772	-0.00690011\\
0.774	-0.00700486\\
0.776	-0.00709359\\
0.778	-0.00713248\\
0.78	-0.00715634\\
0.782	-0.00716532\\
0.784	-0.00715958\\
0.786	-0.00713934\\
0.788	-0.00710482\\
0.79	-0.00705629\\
0.792	-0.00699403\\
0.794	-0.00691835\\
0.796	-0.00682959\\
0.798	-0.00672809\\
0.8	-0.00661425\\
0.802	-0.00648845\\
0.804	-0.00635113\\
0.806	-0.00620271\\
0.808	-0.00604366\\
0.81	-0.00587444\\
0.812	-0.00569554\\
0.814	-0.00550747\\
0.816	-0.00531073\\
0.818	-0.00510585\\
0.82	-0.00489337\\
0.822	-0.00467382\\
0.824	-0.00444777\\
0.826	-0.00421576\\
0.828	-0.00397836\\
0.83	-0.00373613\\
0.832	-0.00348965\\
0.834	-0.00323948\\
0.836	-0.00306366\\
0.838	-0.00303899\\
0.84	-0.00300933\\
0.842	-0.00297479\\
0.844	-0.00293126\\
0.846	-0.00287837\\
0.848	-0.00282205\\
0.85	-0.00276244\\
0.852	-0.00269966\\
0.854	-0.00263384\\
0.856	-0.00256508\\
0.858	-0.00249353\\
0.86	-0.00241931\\
0.862	-0.00234254\\
0.864	-0.00226334\\
0.866	-0.00217466\\
0.868	-0.00207778\\
0.87	-0.00197884\\
0.872	-0.00187805\\
0.874	-0.00177563\\
0.876	-0.00174571\\
0.878	-0.00195541\\
0.88	-0.00215773\\
0.882	-0.00235236\\
0.884	-0.002539\\
0.886	-0.00271739\\
0.888	-0.00288726\\
0.89	-0.00304839\\
0.892	-0.00320059\\
0.894	-0.00334365\\
0.896	-0.00347743\\
0.898	-0.00360177\\
0.9	-0.00371657\\
0.902	-0.00382173\\
0.904	-0.00391717\\
0.906	-0.00400284\\
0.908	-0.00407872\\
0.91	-0.00414478\\
0.912	-0.00419517\\
0.914	-0.00423497\\
0.916	-0.00426535\\
0.918	-0.0042864\\
0.92	-0.0042982\\
0.922	-0.00430087\\
0.924	-0.00429454\\
0.926	-0.00427935\\
0.928	-0.00425547\\
0.93	-0.00422307\\
0.932	-0.00418235\\
0.934	-0.00413351\\
0.936	-0.00407678\\
0.938	-0.0040124\\
0.94	-0.00394062\\
0.942	-0.00386169\\
0.944	-0.00377589\\
0.946	-0.00368349\\
0.948	-0.00358481\\
0.95	-0.00348012\\
0.952	-0.00336975\\
0.954	-0.00325402\\
0.956	-0.00313324\\
0.958	-0.00300774\\
0.96	-0.00287787\\
0.962	-0.00274395\\
0.964	-0.00260634\\
0.966	-0.00246538\\
0.968	-0.00232141\\
0.97	-0.00217478\\
0.972	-0.00202585\\
0.974	-0.00187496\\
0.976	-0.00172246\\
0.978	-0.00165463\\
0.98	-0.00164996\\
0.982	-0.00164123\\
0.984	-0.00162326\\
0.986	-0.00160283\\
0.988	-0.00158004\\
0.99	-0.00155499\\
0.992	-0.00152776\\
0.994	-0.00149846\\
0.996	-0.00146718\\
0.998	-0.00143403\\
1	-0.00139909\\
1.002	-0.00136248\\
1.004	-0.00132428\\
1.006	-0.0012846\\
1.008	-0.00124353\\
1.01	-0.00120116\\
1.012	-0.0011576\\
1.014	-0.00111293\\
1.016	-0.00111662\\
1.018	-0.00123752\\
1.02	-0.00135237\\
1.022	-0.00146147\\
1.024	-0.00156582\\
1.026	-0.00166435\\
1.028	-0.00175752\\
1.03	-0.00184546\\
1.032	-0.00192807\\
1.034	-0.00200525\\
1.036	-0.00207694\\
1.038	-0.00214306\\
1.04	-0.00220356\\
1.042	-0.0022584\\
1.044	-0.00230755\\
1.046	-0.00235101\\
1.048	-0.00238876\\
1.05	-0.00242082\\
1.052	-0.00244721\\
1.054	-0.00246796\\
1.056	-0.00248312\\
1.058	-0.00249039\\
1.06	-0.00249041\\
1.062	-0.00248531\\
1.064	-0.00247517\\
1.066	-0.0024601\\
1.068	-0.00244023\\
1.07	-0.00241568\\
1.072	-0.00238658\\
1.074	-0.00235307\\
1.076	-0.0023153\\
1.078	-0.00227341\\
1.08	-0.00222757\\
1.082	-0.00217793\\
1.084	-0.00212465\\
1.086	-0.00206791\\
1.088	-0.00200787\\
1.09	-0.00194472\\
1.092	-0.00187863\\
1.094	-0.00180978\\
1.096	-0.00173836\\
1.098	-0.00166456\\
1.1	-0.00158855\\
1.102	-0.00151053\\
1.104	-0.00143068\\
1.106	-0.0013492\\
1.108	-0.00126627\\
1.11	-0.00118209\\
1.112	-0.00109684\\
1.114	-0.00101071\\
1.116	-0.000923844\\
1.118	-0.000832903\\
1.12	-0.000806773\\
1.122	-0.000805256\\
1.124	-0.000802186\\
1.126	-0.000797599\\
1.128	-0.000791537\\
1.13	-0.000784042\\
1.132	-0.000775159\\
1.134	-0.000764933\\
1.136	-0.000753413\\
1.138	-0.000740649\\
1.14	-0.000724427\\
1.142	-0.000707112\\
1.144	-0.000688836\\
1.146	-0.000669659\\
1.148	-0.000649644\\
1.15	-0.000628854\\
1.152	-0.00060735\\
1.154	-0.000628311\\
1.156	-0.000694229\\
1.158	-0.000757688\\
1.16	-0.000818594\\
1.162	-0.000876857\\
1.164	-0.000932396\\
1.166	-0.000985135\\
1.168	-0.001035\\
1.17	-0.00108194\\
1.172	-0.00112589\\
1.174	-0.00116681\\
1.176	-0.00120464\\
1.178	-0.00123937\\
1.18	-0.00127095\\
1.182	-0.00129937\\
1.184	-0.00132461\\
1.186	-0.00134666\\
1.188	-0.00136552\\
1.19	-0.00138121\\
1.192	-0.00139372\\
1.194	-0.00140308\\
1.196	-0.00140931\\
1.198	-0.00141218\\
1.2	-0.00141052\\
1.202	-0.00140585\\
1.204	-0.00139823\\
1.206	-0.00138774\\
1.208	-0.00137444\\
1.21	-0.0013584\\
1.212	-0.00133971\\
1.214	-0.00131845\\
1.216	-0.00129471\\
1.218	-0.00126858\\
1.22	-0.00124015\\
1.222	-0.00120953\\
1.224	-0.00117524\\
1.226	-0.00113898\\
1.228	-0.00110093\\
1.23	-0.00106121\\
1.232	-0.00101993\\
1.234	-0.000977216\\
1.236	-0.00093318\\
1.238	-0.000887941\\
1.24	-0.000841618\\
1.242	-0.000794329\\
1.244	-0.000746192\\
1.246	-0.000697324\\
1.248	-0.000647843\\
1.25	-0.000597865\\
1.252	-0.000547503\\
1.254	-0.000496871\\
1.256	-0.000446081\\
1.258	-0.000395243\\
1.26	-0.000361732\\
1.262	-0.000363035\\
1.264	-0.000363612\\
1.266	-0.000363472\\
1.268	-0.000362629\\
1.27	-0.000361094\\
1.272	-0.000358881\\
1.274	-0.000356004\\
1.276	-0.00035248\\
1.278	-0.000348325\\
1.28	-0.000343555\\
1.282	-0.00033773\\
1.284	-0.000330862\\
1.286	-0.000323452\\
1.288	-0.000315528\\
1.29	-0.000337611\\
1.292	-0.000376245\\
1.294	-0.000413417\\
1.296	-0.000449069\\
1.298	-0.00048315\\
1.3	-0.000515612\\
1.302	-0.000546225\\
1.304	-0.000574542\\
1.306	-0.000600894\\
1.308	-0.000625533\\
1.31	-0.000648435\\
1.312	-0.000669582\\
1.314	-0.00068896\\
1.316	-0.000706557\\
1.318	-0.000722366\\
1.32	-0.000736385\\
1.322	-0.000748613\\
1.324	-0.000759055\\
1.326	-0.000767717\\
1.328	-0.000774611\\
1.33	-0.000779752\\
1.332	-0.000783156\\
1.334	-0.000784846\\
1.336	-0.000784844\\
1.338	-0.000783178\\
1.34	-0.000779879\\
1.342	-0.000774979\\
1.344	-0.000768515\\
1.346	-0.000760523\\
1.348	-0.000751046\\
1.35	-0.000740127\\
1.352	-0.000727811\\
1.354	-0.000714147\\
1.356	-0.000699183\\
1.358	-0.000682973\\
1.36	-0.000665569\\
1.362	-0.000647027\\
1.364	-0.000627404\\
1.366	-0.000606637\\
1.368	-0.000584319\\
1.37	-0.000561115\\
1.372	-0.000537088\\
1.374	-0.000512304\\
1.376	-0.000486831\\
1.378	-0.000460735\\
1.38	-0.000434082\\
1.382	-0.000406941\\
1.384	-0.000379142\\
1.386	-0.000350752\\
1.388	-0.000322097\\
1.39	-0.000293246\\
1.392	-0.000264264\\
1.394	-0.000235219\\
1.396	-0.000206176\\
1.398	-0.000177199\\
1.4	-0.000167426\\
1.402	-0.00016739\\
1.404	-0.000167007\\
1.406	-0.000166283\\
1.408	-0.000165229\\
1.41	-0.000163854\\
1.412	-0.000162167\\
1.414	-0.000160178\\
1.416	-0.000157898\\
1.418	-0.000155338\\
1.42	-0.000152508\\
1.422	-0.000149419\\
1.424	-0.000156876\\
1.426	-0.000179663\\
1.428	-0.00020169\\
1.43	-0.000222921\\
1.432	-0.000243323\\
1.434	-0.000262864\\
1.436	-0.000281513\\
1.438	-0.000299245\\
1.44	-0.000316035\\
1.442	-0.000331859\\
1.444	-0.000346693\\
1.446	-0.000360302\\
1.448	-0.000372898\\
1.45	-0.00038447\\
1.452	-0.000394926\\
1.454	-0.000404322\\
1.456	-0.000412647\\
1.458	-0.000419901\\
1.46	-0.000426119\\
1.462	-0.000431184\\
1.464	-0.000435006\\
1.466	-0.000437821\\
1.468	-0.000439642\\
1.47	-0.000440484\\
1.472	-0.00044036\\
1.474	-0.000439288\\
1.476	-0.000437288\\
1.478	-0.00043438\\
1.48	-0.000430586\\
1.482	-0.000425929\\
1.484	-0.000420435\\
1.486	-0.000414129\\
1.488	-0.00040704\\
1.49	-0.000399196\\
1.492	-0.000390626\\
1.494	-0.000381362\\
1.496	-0.000371435\\
1.498	-0.000360878\\
1.5	-0.000349723\\
1.502	-0.000338005\\
1.504	-0.000325759\\
1.506	-0.000313019\\
1.508	-0.00029982\\
1.51	-0.000286199\\
1.512	-0.000272192\\
1.514	-0.000257835\\
1.516	-0.000243164\\
1.518	-0.000228216\\
1.52	-0.000213028\\
1.522	-0.000197635\\
1.524	-0.000182074\\
1.526	-0.000166381\\
1.528	-0.000150579\\
1.53	-0.000134564\\
1.532	-0.000118527\\
1.534	-0.000102443\\
1.536	-8.63649e-05\\
1.538	-8.14354e-05\\
1.54	-8.11659e-05\\
1.542	-8.07106e-05\\
1.544	-8.00733e-05\\
1.546	-7.9072e-05\\
1.548	-7.79088e-05\\
1.55	-7.65902e-05\\
1.552	-7.51228e-05\\
1.554	-7.35136e-05\\
1.556	-7.17695e-05\\
1.558	-7.53344e-05\\
1.56	-8.83213e-05\\
1.562	-0.000100905\\
1.564	-0.000113064\\
1.566	-0.00012478\\
1.568	-0.000136034\\
1.57	-0.00014681\\
1.572	-0.000157092\\
1.574	-0.000166866\\
1.576	-0.000176119\\
1.578	-0.000184839\\
1.58	-0.000193016\\
1.582	-0.000200641\\
1.584	-0.000207706\\
1.586	-0.000214205\\
1.588	-0.000220132\\
1.59	-0.000225485\\
1.592	-0.000230259\\
1.594	-0.000234455\\
1.596	-0.000238072\\
1.598	-0.000241112\\
1.6	-0.000243576\\
1.602	-0.00024547\\
1.604	-0.000246796\\
1.606	-0.000247563\\
1.608	-0.000247776\\
1.61	-0.000247444\\
1.612	-0.000246576\\
1.614	-0.000245183\\
1.616	-0.000243175\\
1.618	-0.000240602\\
1.62	-0.000237546\\
1.622	-0.000234022\\
1.624	-0.000230025\\
1.626	-0.000225496\\
1.628	-0.000220551\\
1.63	-0.000215208\\
1.632	-0.000209487\\
1.634	-0.000203405\\
1.636	-0.000196984\\
1.638	-0.000190242\\
1.64	-0.000183201\\
1.642	-0.00017588\\
1.644	-0.000168301\\
1.646	-0.000160486\\
1.648	-0.000152454\\
1.65	-0.000144228\\
1.652	-0.000135828\\
1.654	-0.000127276\\
1.656	-0.000118594\\
1.658	-0.000109802\\
1.66	-0.000100921\\
1.662	-9.19728e-05\\
1.664	-8.29771e-05\\
1.666	-7.39545e-05\\
1.668	-6.4925e-05\\
1.67	-5.59084e-05\\
1.672	-4.69238e-05\\
1.674	-4.16242e-05\\
1.676	-4.1312e-05\\
1.678	-4.09085e-05\\
1.68	-4.0416e-05\\
1.682	-3.9837e-05\\
1.684	-3.91739e-05\\
1.686	-3.84296e-05\\
1.688	-3.76069e-05\\
1.69	-3.67087e-05\\
1.692	-3.64092e-05\\
1.694	-4.4014e-05\\
1.696	-5.13831e-05\\
1.698	-5.85171e-05\\
1.7	-6.54042e-05\\
1.702	-7.20336e-05\\
1.704	-7.8395e-05\\
1.706	-8.44789e-05\\
1.708	-9.02766e-05\\
1.71	-9.57769e-05\\
1.712	-0.000100898\\
1.714	-0.000105714\\
1.716	-0.000110222\\
1.718	-0.000114415\\
1.72	-0.000118291\\
1.722	-0.000121847\\
1.724	-0.000125079\\
1.726	-0.000127987\\
1.728	-0.000130569\\
1.73	-0.000132826\\
1.732	-0.000134758\\
1.734	-0.000136366\\
1.736	-0.000137652\\
1.738	-0.000138619\\
1.74	-0.000139269\\
1.742	-0.000139607\\
1.744	-0.000139637\\
1.746	-0.000139363\\
1.748	-0.000138792\\
1.75	-0.000137929\\
1.752	-0.000136782\\
1.754	-0.000135356\\
1.756	-0.000133661\\
1.758	-0.000131703\\
1.76	-0.000129491\\
1.762	-0.000127034\\
1.764	-0.000124341\\
1.766	-0.000121422\\
1.768	-0.000118287\\
1.77	-0.000114946\\
1.772	-0.00011141\\
1.774	-0.000107689\\
1.776	-0.000103794\\
1.778	-9.9736e-05\\
1.78	-9.55127e-05\\
1.782	-9.1106e-05\\
1.784	-8.65727e-05\\
1.786	-8.19248e-05\\
1.788	-7.71744e-05\\
1.79	-7.23335e-05\\
1.792	-6.73794e-05\\
1.794	-6.23557e-05\\
1.796	-5.72794e-05\\
1.798	-5.21626e-05\\
1.8	-4.70174e-05\\
1.802	-4.18554e-05\\
1.804	-3.66883e-05\\
1.806	-3.15276e-05\\
1.808	-2.63844e-05\\
1.81	-2.39828e-05\\
1.812	-2.3583e-05\\
1.814	-2.31354e-05\\
1.816	-2.26418e-05\\
1.818	-2.21042e-05\\
1.82	-2.15244e-05\\
1.822	-2.09045e-05\\
1.824	-2.02465e-05\\
1.826	-1.95525e-05\\
1.828	-2.08476e-05\\
1.83	-2.51621e-05\\
1.832	-2.93528e-05\\
1.834	-3.34126e-05\\
1.836	-3.73346e-05\\
1.838	-4.11125e-05\\
1.84	-4.47404e-05\\
1.842	-4.82125e-05\\
1.844	-5.15239e-05\\
1.846	-5.46698e-05\\
1.848	-5.7646e-05\\
1.85	-6.04485e-05\\
1.852	-6.30741e-05\\
1.854	-6.55198e-05\\
1.856	-6.7783e-05\\
1.858	-6.98618e-05\\
1.86	-7.17543e-05\\
1.862	-7.34527e-05\\
1.864	-7.49446e-05\\
1.866	-7.62486e-05\\
1.868	-7.73642e-05\\
1.87	-7.82626e-05\\
1.872	-7.89742e-05\\
1.874	-7.95024e-05\\
1.876	-7.98493e-05\\
1.878	-8.00171e-05\\
1.88	-8.00084e-05\\
1.882	-7.98262e-05\\
1.884	-7.94738e-05\\
1.886	-7.89546e-05\\
1.888	-7.82727e-05\\
1.89	-7.74322e-05\\
1.892	-7.64375e-05\\
1.894	-7.52933e-05\\
1.896	-7.40045e-05\\
1.898	-7.25764e-05\\
1.9	-7.10142e-05\\
1.902	-6.93236e-05\\
1.904	-6.75103e-05\\
1.906	-6.55802e-05\\
1.908	-6.35394e-05\\
1.91	-6.1394e-05\\
1.912	-5.91505e-05\\
1.914	-5.68151e-05\\
1.916	-5.43946e-05\\
1.918	-5.18954e-05\\
1.92	-4.93242e-05\\
1.922	-4.66876e-05\\
1.924	-4.39926e-05\\
1.926	-4.12457e-05\\
1.928	-3.84537e-05\\
1.93	-3.56234e-05\\
1.932	-3.27615e-05\\
1.934	-2.98746e-05\\
1.936	-2.69693e-05\\
1.938	-2.40522e-05\\
1.94	-2.11193e-05\\
1.942	-1.81876e-05\\
1.944	-1.58743e-05\\
1.946	-1.63896e-05\\
1.948	-1.68435e-05\\
1.95	-1.72391e-05\\
1.952	-1.75703e-05\\
1.954	-1.78309e-05\\
1.956	-1.80326e-05\\
1.958	-1.81755e-05\\
1.96	-1.82596e-05\\
1.962	-1.82854e-05\\
1.964	-1.82535e-05\\
1.966	-1.81647e-05\\
1.968	-1.80199e-05\\
1.97	-1.93248e-05\\
1.972	-2.15609e-05\\
1.974	-2.37149e-05\\
1.976	-2.57833e-05\\
1.978	-2.77631e-05\\
1.98	-2.96513e-05\\
1.982	-3.14454e-05\\
1.984	-3.3143e-05\\
1.986	-3.47418e-05\\
1.988	-3.624e-05\\
1.99	-3.76359e-05\\
1.992	-3.89281e-05\\
1.994	-4.01154e-05\\
1.996	-4.11968e-05\\
1.998	-4.21718e-05\\
2	-4.30398e-05\\
};
\addplot [color=mycolor4, line width=1.0pt, forget plot]
  table[row sep=crcr]{%
-0.024	0\\
-0.022	0\\
-0.02	0\\
-0.018	0\\
-0.016	0\\
-0.014	0\\
-0.012	0\\
-0.01	0\\
-0.008	0\\
-0.006	0\\
-0.004	0\\
-0.002	0\\
0	0\\
0.002	-4.06872e-10\\
0.004	-7.6163e-09\\
0.006	-3.98869e-08\\
0.008	-1.27203e-07\\
0.01	-3.10871e-07\\
0.012	-6.43473e-07\\
0.014	-1.18904e-06\\
0.016	-2.02332e-06\\
0.018	-3.23407e-06\\
0.02	-4.9212e-06\\
0.022	-7.19682e-06\\
0.024	-1.0185e-05\\
0.026	-1.40216e-05\\
0.028	-1.88529e-05\\
0.03	-2.48355e-05\\
0.032	-3.21341e-05\\
0.034	-4.09208e-05\\
0.036	-5.13731e-05\\
0.038	-6.36716e-05\\
0.04	-8.87865e-06\\
0.042	-1.24861e-05\\
0.044	-1.71328e-05\\
0.046	-2.30241e-05\\
0.048	-3.03886e-05\\
0.05	-3.94785e-05\\
0.052	-5.05693e-05\\
0.054	-6.39591e-05\\
0.056	-7.99677e-05\\
0.058	-9.89353e-05\\
0.06	-0.00012122\\
0.062	-0.000147197\\
0.064	-0.000177255\\
0.066	-0.000211791\\
0.068	-0.000251211\\
0.07	-0.000295925\\
0.072	-0.000346341\\
0.074	-0.000402861\\
0.076	-0.00046588\\
0.078	-0.00675561\\
0.08	-0.0158141\\
0.082	-0.0187389\\
0.084	-0.0192986\\
0.086	-0.0198245\\
0.088	-0.0203147\\
0.09	-0.020767\\
0.092	-0.0211797\\
0.094	-0.0215512\\
0.096	-0.0218799\\
0.098	-0.0221647\\
0.1	-0.0224042\\
0.102	-0.0225975\\
0.104	-0.0227438\\
0.106	-0.0228422\\
0.108	-0.0228922\\
0.11	-0.0228933\\
0.112	-0.0223411\\
0.114	-0.0207129\\
0.116	-0.0190988\\
0.118	-0.0174969\\
0.12	-0.017538\\
0.122	-0.0180859\\
0.124	-0.0186104\\
0.126	-0.0191098\\
0.128	-0.0195822\\
0.13	-0.0200261\\
0.132	-0.02044\\
0.134	-0.0208224\\
0.136	-0.0211723\\
0.138	-0.0214885\\
0.14	-0.02177\\
0.142	-0.0220162\\
0.144	-0.0222263\\
0.146	-0.0223998\\
0.148	-0.0225363\\
0.15	-0.0226355\\
0.152	-0.0226972\\
0.154	-0.0227215\\
0.156	-0.0227082\\
0.158	-0.0226577\\
0.16	-0.0225701\\
0.162	-0.0224458\\
0.164	-0.0222852\\
0.166	-0.0220887\\
0.168	-0.021857\\
0.17	-0.0215907\\
0.172	-0.0212905\\
0.174	-0.020957\\
0.176	-0.0205912\\
0.178	-0.0201938\\
0.18	-0.0197657\\
0.182	-0.0193079\\
0.184	-0.0188213\\
0.186	-0.0183068\\
0.188	-0.0177656\\
0.19	-0.0171985\\
0.192	-0.0166068\\
0.194	-0.0159913\\
0.196	-0.0153533\\
0.198	-0.0146939\\
0.2	-0.014014\\
0.202	-0.013315\\
0.204	-0.0125979\\
0.206	-0.0118638\\
0.208	-0.0111139\\
0.21	-0.0103494\\
0.212	-0.0095714\\
0.214	-0.00878108\\
0.216	-0.0079796\\
0.218	-0.00716812\\
0.22	-0.0063478\\
0.222	-0.00551979\\
0.224	-0.00468525\\
0.226	-0.00384532\\
0.228	-0.00300114\\
0.23	-0.00215386\\
0.232	-0.00130459\\
0.234	-0.000454454\\
0.236	-0.000317021\\
0.238	-0.000439951\\
0.24	-0.000524878\\
0.242	-0.000573552\\
0.244	-0.000587907\\
0.246	-0.000570036\\
0.248	-0.0015599\\
0.25	-0.00285307\\
0.252	-0.0040134\\
0.254	-0.00504351\\
0.256	-0.00589838\\
0.258	-0.00647848\\
0.26	-0.00695508\\
0.262	-0.00732957\\
0.264	-0.00760348\\
0.266	-0.00777859\\
0.268	-0.00785685\\
0.27	-0.00784044\\
0.272	-0.00773176\\
0.274	-0.00753344\\
0.276	-0.00717811\\
0.278	-0.00641584\\
0.28	-0.00558909\\
0.282	-0.00470204\\
0.284	-0.00375893\\
0.286	-0.00276406\\
0.288	-0.00172175\\
0.29	-0.0014209\\
0.292	-0.00135984\\
0.294	-0.00129718\\
0.296	-0.00161212\\
0.298	-0.002762\\
0.3	-0.00386263\\
0.302	-0.00491416\\
0.304	-0.00591681\\
0.306	-0.0068708\\
0.308	-0.00777639\\
0.31	-0.00863388\\
0.312	-0.00944356\\
0.314	-0.0102058\\
0.316	-0.0109209\\
0.318	-0.0115893\\
0.32	-0.0122115\\
0.322	-0.0127878\\
0.324	-0.0133189\\
0.326	-0.0138051\\
0.328	-0.014247\\
0.33	-0.0146452\\
0.332	-0.0150003\\
0.334	-0.015313\\
0.336	-0.0155839\\
0.338	-0.0158136\\
0.34	-0.0158371\\
0.342	-0.0156571\\
0.344	-0.0154557\\
0.346	-0.0152339\\
0.348	-0.0149924\\
0.35	-0.0147322\\
0.352	-0.0144542\\
0.354	-0.0147568\\
0.356	-0.0152285\\
0.358	-0.015656\\
0.36	-0.0160386\\
0.362	-0.0163759\\
0.364	-0.0166674\\
0.366	-0.0169128\\
0.368	-0.017112\\
0.37	-0.017265\\
0.372	-0.0173717\\
0.374	-0.0174322\\
0.376	-0.017447\\
0.378	-0.0174163\\
0.38	-0.0173406\\
0.382	-0.0172204\\
0.384	-0.0170566\\
0.386	-0.0168497\\
0.388	-0.0166008\\
0.39	-0.0163107\\
0.392	-0.0159805\\
0.394	-0.0156114\\
0.396	-0.0152045\\
0.398	-0.0147612\\
0.4	-0.0142827\\
0.402	-0.0137707\\
0.404	-0.0132264\\
0.406	-0.0126516\\
0.408	-0.0120479\\
0.41	-0.0114169\\
0.412	-0.0107603\\
0.414	-0.01008\\
0.416	-0.00937766\\
0.418	-0.00865522\\
0.42	-0.00791452\\
0.422	-0.00715748\\
0.424	-0.006386\\
0.426	-0.00560203\\
0.428	-0.00480751\\
0.43	-0.00400439\\
0.432	-0.00319464\\
0.434	-0.00238019\\
0.436	-0.00156301\\
0.438	-0.000745014\\
0.44	7.18746e-05\\
0.442	0.00050048\\
0.444	0.000249764\\
0.446	-0.000220608\\
0.448	-0.000916322\\
0.45	-0.00159648\\
0.452	-0.00226038\\
0.454	-0.00290739\\
0.456	-0.00353688\\
0.458	-0.00414826\\
0.46	-0.00474096\\
0.462	-0.00531447\\
0.464	-0.00586828\\
0.466	-0.00640193\\
0.468	-0.00725566\\
0.47	-0.00808611\\
0.472	-0.0088826\\
0.474	-0.00964394\\
0.476	-0.0103157\\
0.478	-0.0108771\\
0.48	-0.0114051\\
0.482	-0.0118991\\
0.484	-0.0123584\\
0.486	-0.0127823\\
0.488	-0.0131704\\
0.49	-0.0135222\\
0.492	-0.0138374\\
0.494	-0.0141157\\
0.496	-0.014357\\
0.498	-0.0145611\\
0.5	-0.0147281\\
0.502	-0.0148579\\
0.504	-0.0149509\\
0.506	-0.0150072\\
0.508	-0.0150271\\
0.51	-0.015011\\
0.512	-0.0149594\\
0.514	-0.0148729\\
0.516	-0.0147519\\
0.518	-0.0145973\\
0.52	-0.0144097\\
0.522	-0.01419\\
0.524	-0.013939\\
0.526	-0.0136576\\
0.528	-0.0133469\\
0.53	-0.0130078\\
0.532	-0.0126415\\
0.534	-0.0122491\\
0.536	-0.0118316\\
0.538	-0.0113905\\
0.54	-0.0109268\\
0.542	-0.010442\\
0.544	-0.00993723\\
0.546	-0.00941394\\
0.548	-0.00887347\\
0.55	-0.00831722\\
0.552	-0.00774658\\
0.554	-0.00716299\\
0.556	-0.00685254\\
0.558	-0.00654224\\
0.56	-0.00622545\\
0.562	-0.00590274\\
0.564	-0.00557468\\
0.566	-0.00524184\\
0.568	-0.00490482\\
0.57	-0.00456419\\
0.572	-0.00422053\\
0.574	-0.00387442\\
0.576	-0.00352644\\
0.578	-0.00317717\\
0.58	-0.00282716\\
0.582	-0.002477\\
0.584	-0.00212723\\
0.586	-0.00177841\\
0.588	-0.00143109\\
0.59	-0.00134653\\
0.592	-0.00202342\\
0.594	-0.00268574\\
0.596	-0.00333231\\
0.598	-0.003962\\
0.6	-0.00457373\\
0.602	-0.00516646\\
0.604	-0.00573922\\
0.606	-0.00629109\\
0.608	-0.00682121\\
0.61	-0.00732877\\
0.612	-0.00781303\\
0.614	-0.0082733\\
0.616	-0.00870896\\
0.618	-0.00911839\\
0.62	-0.0094178\\
0.622	-0.00969403\\
0.624	-0.00994676\\
0.626	-0.0101757\\
0.628	-0.0103808\\
0.63	-0.0105617\\
0.632	-0.0107183\\
0.634	-0.0108507\\
0.636	-0.0109588\\
0.638	-0.0110427\\
0.64	-0.0111025\\
0.642	-0.0111383\\
0.644	-0.0111504\\
0.646	-0.0111389\\
0.648	-0.0111043\\
0.65	-0.0110467\\
0.652	-0.0109667\\
0.654	-0.0108647\\
0.656	-0.0107411\\
0.658	-0.0105964\\
0.66	-0.0104313\\
0.662	-0.0102462\\
0.664	-0.0100418\\
0.666	-0.00981884\\
0.668	-0.00957793\\
0.67	-0.0093198\\
0.672	-0.00904521\\
0.674	-0.00875493\\
0.676	-0.00844976\\
0.678	-0.00813053\\
0.68	-0.00779807\\
0.682	-0.00745324\\
0.684	-0.00709692\\
0.686	-0.00672999\\
0.688	-0.00635336\\
0.69	-0.00596794\\
0.692	-0.00557465\\
0.694	-0.00532286\\
0.696	-0.00521974\\
0.698	-0.00510932\\
0.7	-0.0049919\\
0.702	-0.0048678\\
0.704	-0.00473732\\
0.706	-0.00460077\\
0.708	-0.00445848\\
0.71	-0.00431078\\
0.712	-0.00415801\\
0.714	-0.00400048\\
0.716	-0.00383856\\
0.718	-0.00367257\\
0.72	-0.00350287\\
0.722	-0.0033298\\
0.724	-0.0031537\\
0.726	-0.00297493\\
0.728	-0.00279383\\
0.73	-0.00261074\\
0.732	-0.00242601\\
0.734	-0.00234825\\
0.736	-0.00272845\\
0.738	-0.00309429\\
0.74	-0.00344845\\
0.742	-0.00379033\\
0.744	-0.00411938\\
0.746	-0.00443508\\
0.748	-0.00473695\\
0.75	-0.00502455\\
0.752	-0.00529746\\
0.754	-0.00555532\\
0.756	-0.00579779\\
0.758	-0.00602458\\
0.76	-0.00623543\\
0.762	-0.00643012\\
0.764	-0.00660848\\
0.766	-0.00677035\\
0.768	-0.00691563\\
0.77	-0.00704426\\
0.772	-0.00713991\\
0.774	-0.00721192\\
0.776	-0.00726815\\
0.778	-0.00730869\\
0.78	-0.00733366\\
0.782	-0.00734321\\
0.784	-0.00733754\\
0.786	-0.00731685\\
0.788	-0.00728139\\
0.79	-0.00723143\\
0.792	-0.00716728\\
0.794	-0.00708925\\
0.796	-0.0069977\\
0.798	-0.00689301\\
0.8	-0.00677556\\
0.802	-0.0066458\\
0.804	-0.00650414\\
0.806	-0.00635105\\
0.808	-0.00618702\\
0.81	-0.00601253\\
0.812	-0.00582809\\
0.814	-0.00563423\\
0.816	-0.00543148\\
0.818	-0.0052204\\
0.82	-0.00500154\\
0.822	-0.00477546\\
0.824	-0.00454275\\
0.826	-0.00430398\\
0.828	-0.00405974\\
0.83	-0.00381062\\
0.832	-0.0035572\\
0.834	-0.00330008\\
0.836	-0.00311733\\
0.838	-0.00308575\\
0.84	-0.00304923\\
0.842	-0.0030079\\
0.844	-0.00296192\\
0.846	-0.00291143\\
0.848	-0.0028566\\
0.85	-0.00279759\\
0.852	-0.00273458\\
0.854	-0.00266773\\
0.856	-0.00259723\\
0.858	-0.00252325\\
0.86	-0.002446\\
0.862	-0.00236564\\
0.864	-0.00228034\\
0.866	-0.00218745\\
0.868	-0.00209221\\
0.87	-0.00199483\\
0.872	-0.00189551\\
0.874	-0.00179447\\
0.876	-0.00176582\\
0.878	-0.0019767\\
0.88	-0.00218011\\
0.882	-0.00237572\\
0.884	-0.00256323\\
0.886	-0.00274238\\
0.888	-0.00291292\\
0.89	-0.00307461\\
0.892	-0.00322725\\
0.894	-0.00337066\\
0.896	-0.00350467\\
0.898	-0.00362915\\
0.9	-0.00374399\\
0.902	-0.00384908\\
0.904	-0.00394436\\
0.906	-0.00402977\\
0.908	-0.0041053\\
0.91	-0.00417094\\
0.912	-0.00422669\\
0.914	-0.00427261\\
0.916	-0.00430874\\
0.918	-0.00433516\\
0.92	-0.00435197\\
0.922	-0.00435928\\
0.924	-0.00435723\\
0.926	-0.00434009\\
0.928	-0.00431393\\
0.93	-0.00427915\\
0.932	-0.00423595\\
0.934	-0.00418456\\
0.936	-0.00412521\\
0.938	-0.00405814\\
0.94	-0.00398362\\
0.942	-0.00390191\\
0.944	-0.0038133\\
0.946	-0.00371808\\
0.948	-0.00361655\\
0.95	-0.00350902\\
0.952	-0.00339581\\
0.954	-0.00327725\\
0.956	-0.00315366\\
0.958	-0.00302539\\
0.96	-0.00289278\\
0.962	-0.00275617\\
0.964	-0.00261592\\
0.966	-0.00247236\\
0.968	-0.00232587\\
0.97	-0.00217678\\
0.972	-0.00202546\\
0.974	-0.00187226\\
0.976	-0.00171752\\
0.978	-0.00164754\\
0.98	-0.00164081\\
0.982	-0.00163111\\
0.984	-0.0016185\\
0.986	-0.00160306\\
0.988	-0.00158486\\
0.99	-0.00156399\\
0.992	-0.00154054\\
0.994	-0.00151459\\
0.996	-0.00148625\\
0.998	-0.0014556\\
1	-0.00142274\\
1.002	-0.00138779\\
1.004	-0.00135084\\
1.006	-0.00131199\\
1.008	-0.00127137\\
1.01	-0.00122907\\
1.012	-0.00118522\\
1.014	-0.00113863\\
1.016	-0.00113878\\
1.018	-0.00125592\\
1.02	-0.00136853\\
1.022	-0.00147643\\
1.024	-0.00157946\\
1.026	-0.0016775\\
1.028	-0.0017704\\
1.03	-0.00185805\\
1.032	-0.00194035\\
1.034	-0.00201721\\
1.036	-0.00208856\\
1.038	-0.00215432\\
1.04	-0.00221445\\
1.042	-0.00226891\\
1.044	-0.00231768\\
1.046	-0.00236074\\
1.048	-0.00239809\\
1.05	-0.00242975\\
1.052	-0.00245573\\
1.054	-0.00247608\\
1.056	-0.00249084\\
1.058	-0.00250007\\
1.06	-0.00250383\\
1.062	-0.00250221\\
1.064	-0.0024953\\
1.066	-0.00248319\\
1.068	-0.002466\\
1.07	-0.00244384\\
1.072	-0.00241684\\
1.074	-0.00238513\\
1.076	-0.00234887\\
1.078	-0.00230716\\
1.08	-0.0022609\\
1.082	-0.00221061\\
1.084	-0.00215645\\
1.086	-0.00209862\\
1.088	-0.00203728\\
1.09	-0.00197262\\
1.092	-0.00190484\\
1.094	-0.00183413\\
1.096	-0.00176068\\
1.098	-0.00168469\\
1.1	-0.00160637\\
1.102	-0.00152592\\
1.104	-0.00144354\\
1.106	-0.00135944\\
1.108	-0.00127382\\
1.11	-0.0011869\\
1.112	-0.00109887\\
1.114	-0.00100995\\
1.116	-0.000920326\\
1.118	-0.000830211\\
1.12	-0.000804891\\
1.122	-0.000804167\\
1.124	-0.000801869\\
1.126	-0.000798033\\
1.128	-0.000792695\\
1.13	-0.000785896\\
1.132	-0.000777679\\
1.134	-0.000768087\\
1.136	-0.000757168\\
1.138	-0.000744968\\
1.14	-0.000731538\\
1.142	-0.00071693\\
1.144	-0.000701194\\
1.146	-0.000684385\\
1.148	-0.000666559\\
1.15	-0.00064777\\
1.152	-0.000628076\\
1.154	-0.000650651\\
1.156	-0.000717939\\
1.158	-0.000782026\\
1.16	-0.000843213\\
1.162	-0.000901526\\
1.164	-0.000956948\\
1.166	-0.00100941\\
1.168	-0.00105884\\
1.17	-0.00110519\\
1.172	-0.00114841\\
1.174	-0.00118845\\
1.176	-0.00122528\\
1.178	-0.00125888\\
1.18	-0.00128922\\
1.182	-0.0013163\\
1.184	-0.0013401\\
1.186	-0.00136064\\
1.188	-0.00137791\\
1.19	-0.00139195\\
1.192	-0.00140276\\
1.194	-0.00141038\\
1.196	-0.00141484\\
1.198	-0.0014162\\
1.2	-0.00141448\\
1.202	-0.00140976\\
1.204	-0.00140209\\
1.206	-0.00139153\\
1.208	-0.00137815\\
1.21	-0.00136203\\
1.212	-0.00134325\\
1.214	-0.0013219\\
1.216	-0.00129805\\
1.218	-0.00127181\\
1.22	-0.00124327\\
1.222	-0.00121253\\
1.224	-0.00117969\\
1.226	-0.00114486\\
1.228	-0.00110814\\
1.23	-0.00106964\\
1.232	-0.0010294\\
1.234	-0.000987473\\
1.236	-0.000944121\\
1.238	-0.000899464\\
1.24	-0.000853621\\
1.242	-0.000806709\\
1.244	-0.000758847\\
1.246	-0.000710154\\
1.248	-0.000660748\\
1.25	-0.000610747\\
1.252	-0.000560269\\
1.254	-0.000509431\\
1.256	-0.000458347\\
1.258	-0.000407132\\
1.26	-0.000373167\\
1.262	-0.000373943\\
1.264	-0.000373925\\
1.266	-0.00037313\\
1.268	-0.000371575\\
1.27	-0.000369279\\
1.272	-0.000366261\\
1.274	-0.000362543\\
1.276	-0.000358147\\
1.278	-0.000353096\\
1.28	-0.000347414\\
1.282	-0.000341126\\
1.284	-0.000334256\\
1.286	-0.000326816\\
1.288	-0.000318719\\
1.29	-0.000340625\\
1.292	-0.000379078\\
1.294	-0.000416063\\
1.296	-0.000451526\\
1.298	-0.000485415\\
1.3	-0.000517683\\
1.302	-0.000548287\\
1.304	-0.000577189\\
1.306	-0.000604354\\
1.308	-0.000629753\\
1.31	-0.00065336\\
1.312	-0.000675145\\
1.314	-0.0006951\\
1.316	-0.000713213\\
1.318	-0.000729477\\
1.32	-0.000743886\\
1.322	-0.000756441\\
1.324	-0.000767146\\
1.326	-0.000775961\\
1.328	-0.000782883\\
1.33	-0.000787992\\
1.332	-0.000791308\\
1.334	-0.000792853\\
1.336	-0.000792654\\
1.338	-0.00079074\\
1.34	-0.000787144\\
1.342	-0.000781902\\
1.344	-0.000775053\\
1.346	-0.000766639\\
1.348	-0.000756704\\
1.35	-0.000745296\\
1.352	-0.000732463\\
1.354	-0.000718257\\
1.356	-0.000702733\\
1.358	-0.000685946\\
1.36	-0.000667953\\
1.362	-0.000648814\\
1.364	-0.00062859\\
1.366	-0.000607342\\
1.368	-0.000585134\\
1.37	-0.000562031\\
1.372	-0.000538097\\
1.374	-0.000513398\\
1.376	-0.000488001\\
1.378	-0.000461973\\
1.38	-0.00043538\\
1.382	-0.000408291\\
1.384	-0.000380771\\
1.386	-0.000352889\\
1.388	-0.00032471\\
1.39	-0.000296301\\
1.392	-0.000267728\\
1.394	-0.000239055\\
1.396	-0.000210346\\
1.398	-0.000181665\\
1.4	-0.000172149\\
1.402	-0.000172331\\
1.404	-0.000172116\\
1.406	-0.000171517\\
1.408	-0.000170549\\
1.41	-0.000169222\\
1.412	-0.000167547\\
1.414	-0.000165534\\
1.416	-0.000163194\\
1.418	-0.000160519\\
1.42	-0.000157528\\
1.422	-0.00015425\\
1.424	-0.00016149\\
1.426	-0.000184034\\
1.428	-0.000205796\\
1.43	-0.000226741\\
1.432	-0.000246839\\
1.434	-0.000266058\\
1.436	-0.000284371\\
1.438	-0.00030175\\
1.44	-0.000318176\\
1.442	-0.000333631\\
1.444	-0.000348097\\
1.446	-0.000361559\\
1.448	-0.000374004\\
1.45	-0.000385422\\
1.452	-0.000395807\\
1.454	-0.000405152\\
1.456	-0.000413454\\
1.458	-0.000420714\\
1.46	-0.000426932\\
1.462	-0.000432114\\
1.464	-0.000436264\\
1.466	-0.000439391\\
1.468	-0.000441505\\
1.47	-0.00044262\\
1.472	-0.000442749\\
1.474	-0.000441908\\
1.476	-0.000440116\\
1.478	-0.000437393\\
1.48	-0.000433741\\
1.482	-0.000429197\\
1.484	-0.000423792\\
1.486	-0.000417552\\
1.488	-0.000410505\\
1.49	-0.000402681\\
1.492	-0.000394108\\
1.494	-0.000384819\\
1.496	-0.000374846\\
1.498	-0.000364221\\
1.5	-0.000352979\\
1.502	-0.000341156\\
1.504	-0.000328785\\
1.506	-0.000315905\\
1.508	-0.00030255\\
1.51	-0.00028876\\
1.512	-0.000274569\\
1.514	-0.000260018\\
1.516	-0.000245143\\
1.518	-0.000229982\\
1.52	-0.000214573\\
1.522	-0.000198954\\
1.524	-0.000183163\\
1.526	-0.000167237\\
1.528	-0.000151214\\
1.53	-0.000135128\\
1.532	-0.000119018\\
1.534	-0.000102919\\
1.536	-8.68648e-05\\
1.538	-8.19578e-05\\
1.54	-8.1709e-05\\
1.542	-8.12724e-05\\
1.544	-8.06528e-05\\
1.546	-7.98553e-05\\
1.548	-7.88853e-05\\
1.55	-7.77484e-05\\
1.552	-7.64506e-05\\
1.554	-7.49981e-05\\
1.556	-7.33972e-05\\
1.558	-7.70914e-05\\
1.56	-9.01931e-05\\
1.562	-0.000102877\\
1.564	-0.000115121\\
1.566	-0.000126907\\
1.568	-0.000138217\\
1.57	-0.000149033\\
1.572	-0.000159341\\
1.574	-0.000169126\\
1.576	-0.000178376\\
1.578	-0.000187079\\
1.58	-0.000195219\\
1.582	-0.000202791\\
1.584	-0.000209791\\
1.586	-0.000216213\\
1.588	-0.000222054\\
1.59	-0.00022731\\
1.592	-0.00023198\\
1.594	-0.000236062\\
1.596	-0.000239558\\
1.598	-0.000242471\\
1.6	-0.000244802\\
1.602	-0.000246558\\
1.604	-0.000247744\\
1.606	-0.000248367\\
1.608	-0.000248434\\
1.61	-0.000247956\\
1.612	-0.000246942\\
1.614	-0.000245403\\
1.616	-0.000243352\\
1.618	-0.000240802\\
1.62	-0.000237764\\
1.622	-0.000234255\\
1.624	-0.000230292\\
1.626	-0.00022589\\
1.628	-0.000221067\\
1.63	-0.00021584\\
1.632	-0.000210227\\
1.634	-0.000204247\\
1.636	-0.000197919\\
1.638	-0.000191263\\
1.64	-0.000184299\\
1.642	-0.000177048\\
1.644	-0.00016953\\
1.646	-0.000161765\\
1.648	-0.000153776\\
1.65	-0.000145583\\
1.652	-0.000137208\\
1.654	-0.000128673\\
1.656	-0.000119993\\
1.658	-0.000111195\\
1.66	-0.0001023\\
1.662	-9.33299e-05\\
1.664	-8.43044e-05\\
1.666	-7.52447e-05\\
1.668	-6.61713e-05\\
1.67	-5.71042e-05\\
1.672	-4.80631e-05\\
1.674	-4.27014e-05\\
1.676	-4.23221e-05\\
1.678	-4.1847e-05\\
1.68	-4.12789e-05\\
1.682	-4.06208e-05\\
1.684	-3.98758e-05\\
1.686	-3.90471e-05\\
1.688	-3.81382e-05\\
1.69	-3.71527e-05\\
1.692	-3.6765e-05\\
1.694	-4.43531e-05\\
1.696	-5.17176e-05\\
1.698	-5.88461e-05\\
1.7	-6.5727e-05\\
1.702	-7.23494e-05\\
1.704	-7.8703e-05\\
1.706	-8.47785e-05\\
1.708	-9.05671e-05\\
1.71	-9.60611e-05\\
1.712	-0.000101253\\
1.714	-0.000106137\\
1.716	-0.000110707\\
1.718	-0.000114959\\
1.72	-0.000118889\\
1.722	-0.000122492\\
1.724	-0.000125765\\
1.726	-0.000128707\\
1.728	-0.000131319\\
1.73	-0.000133601\\
1.732	-0.000135551\\
1.734	-0.000137173\\
1.736	-0.000138468\\
1.738	-0.000139437\\
1.74	-0.000140086\\
1.742	-0.000140417\\
1.744	-0.000140435\\
1.746	-0.000140146\\
1.748	-0.000139554\\
1.75	-0.000138667\\
1.752	-0.000137491\\
1.754	-0.000136033\\
1.756	-0.000134302\\
1.758	-0.000132306\\
1.76	-0.000130053\\
1.762	-0.000127552\\
1.764	-0.000124814\\
1.766	-0.000121848\\
1.768	-0.000118664\\
1.77	-0.000115273\\
1.772	-0.000111685\\
1.774	-0.000107913\\
1.776	-0.000103967\\
1.778	-9.98578e-05\\
1.78	-9.55983e-05\\
1.782	-9.11997e-05\\
1.784	-8.66741e-05\\
1.786	-8.20334e-05\\
1.788	-7.72897e-05\\
1.79	-7.24549e-05\\
1.792	-6.75412e-05\\
1.794	-6.25602e-05\\
1.796	-5.75233e-05\\
1.798	-5.24434e-05\\
1.8	-4.73322e-05\\
1.802	-4.22014e-05\\
1.804	-3.70625e-05\\
1.806	-3.19269e-05\\
1.808	-2.68057e-05\\
1.81	-2.4423e-05\\
1.812	-2.40387e-05\\
1.814	-2.36035e-05\\
1.816	-2.3119e-05\\
1.818	-2.25872e-05\\
1.82	-2.20096e-05\\
1.822	-2.13889e-05\\
1.824	-2.07271e-05\\
1.826	-2.00265e-05\\
1.828	-2.13124e-05\\
1.83	-2.5615e-05\\
1.832	-2.97913e-05\\
1.834	-3.38345e-05\\
1.836	-3.77377e-05\\
1.838	-4.14945e-05\\
1.84	-4.50991e-05\\
1.842	-4.85456e-05\\
1.844	-5.18298e-05\\
1.846	-5.49472e-05\\
1.848	-5.78939e-05\\
1.85	-6.06662e-05\\
1.852	-6.32609e-05\\
1.854	-6.56752e-05\\
1.856	-6.79068e-05\\
1.858	-6.99539e-05\\
1.86	-7.1815e-05\\
1.862	-7.34891e-05\\
1.864	-7.49756e-05\\
1.866	-7.62742e-05\\
1.868	-7.73853e-05\\
1.87	-7.83095e-05\\
1.872	-7.90477e-05\\
1.874	-7.96015e-05\\
1.876	-7.99725e-05\\
1.878	-8.0163e-05\\
1.88	-8.01756e-05\\
1.882	-8.00129e-05\\
1.884	-7.96784e-05\\
1.886	-7.91754e-05\\
1.888	-7.85078e-05\\
1.89	-7.76797e-05\\
1.892	-7.66956e-05\\
1.894	-7.55601e-05\\
1.896	-7.42781e-05\\
1.898	-7.28549e-05\\
1.9	-7.12957e-05\\
1.902	-6.96062e-05\\
1.904	-6.77922e-05\\
1.906	-6.58597e-05\\
1.908	-6.38147e-05\\
1.91	-6.16636e-05\\
1.912	-5.94127e-05\\
1.914	-5.70685e-05\\
1.916	-5.46378e-05\\
1.918	-5.2127e-05\\
1.92	-4.95432e-05\\
1.922	-4.68929e-05\\
1.924	-4.41831e-05\\
1.926	-4.14206e-05\\
1.928	-3.86123e-05\\
1.93	-3.5765e-05\\
1.932	-3.28857e-05\\
1.934	-2.99809e-05\\
1.936	-2.70576e-05\\
1.938	-2.41223e-05\\
1.94	-2.11811e-05\\
1.942	-1.8241e-05\\
1.944	-1.59192e-05\\
1.946	-1.64302e-05\\
1.948	-1.68839e-05\\
1.95	-1.72792e-05\\
1.952	-1.76155e-05\\
1.954	-1.78925e-05\\
1.956	-1.81097e-05\\
1.958	-1.82672e-05\\
1.96	-1.83651e-05\\
1.962	-1.84037e-05\\
1.964	-1.83836e-05\\
1.966	-1.83055e-05\\
1.968	-1.81703e-05\\
1.97	-1.94838e-05\\
1.972	-2.17272e-05\\
1.974	-2.38868e-05\\
1.976	-2.59596e-05\\
1.978	-2.79426e-05\\
1.98	-2.9833e-05\\
1.982	-3.1628e-05\\
1.984	-3.33253e-05\\
1.986	-3.49228e-05\\
1.988	-3.64186e-05\\
1.99	-3.78111e-05\\
1.992	-3.90989e-05\\
1.994	-4.02809e-05\\
1.996	-4.13561e-05\\
1.998	-4.2324e-05\\
2	-4.31842e-05\\
};
\addplot [color=mycolor5, line width=1.0pt, forget plot]
  table[row sep=crcr]{%
-0.024	0\\
-0.022	0\\
-0.02	0\\
-0.018	0\\
-0.016	0\\
-0.014	0\\
-0.012	0\\
-0.01	0\\
-0.008	0\\
-0.006	0\\
-0.004	0\\
-0.002	0\\
0	0\\
0.002	-1.1813e-10\\
0.004	-5.55202e-09\\
0.006	-3.36838e-08\\
0.008	-1.14187e-07\\
0.01	-2.88562e-07\\
0.012	-6.10072e-07\\
0.014	-1.14393e-06\\
0.016	-1.96763e-06\\
0.018	-3.17128e-06\\
0.02	-4.85784e-06\\
0.022	-7.14329e-06\\
0.024	-1.01564e-05\\
0.026	-1.40387e-05\\
0.028	-1.89434e-05\\
0.03	-2.50346e-05\\
0.032	-3.24865e-05\\
0.034	-4.14814e-05\\
0.036	-5.22083e-05\\
0.038	-6.48608e-05\\
0.04	-1.15881e-05\\
0.042	-1.60143e-05\\
0.044	-2.16522e-05\\
0.046	-2.87302e-05\\
0.048	-3.75014e-05\\
0.05	-4.82432e-05\\
0.052	-6.12575e-05\\
0.054	-7.68695e-05\\
0.056	-9.54267e-05\\
0.058	-0.000117297\\
0.06	-0.000142867\\
0.062	-0.000172539\\
0.064	-0.00020673\\
0.066	-0.000245863\\
0.068	-0.000290371\\
0.07	-0.000340687\\
0.072	-0.000397243\\
0.074	-0.000460464\\
0.076	-0.000530763\\
0.078	-0.00682837\\
0.08	-0.0158954\\
0.082	-0.019776\\
0.084	-0.0203778\\
0.086	-0.0209446\\
0.088	-0.0214744\\
0.09	-0.0219647\\
0.092	-0.0224138\\
0.094	-0.0228197\\
0.096	-0.0231809\\
0.098	-0.0234959\\
0.1	-0.0237634\\
0.102	-0.0239821\\
0.104	-0.0241512\\
0.106	-0.0242698\\
0.108	-0.024337\\
0.11	-0.0243523\\
0.112	-0.0236194\\
0.114	-0.0220052\\
0.116	-0.0204027\\
0.118	-0.0188101\\
0.12	-0.0188579\\
0.122	-0.01941\\
0.124	-0.0199362\\
0.126	-0.0204347\\
0.128	-0.0209036\\
0.13	-0.0213413\\
0.132	-0.0217463\\
0.134	-0.0221173\\
0.136	-0.022453\\
0.138	-0.0227523\\
0.14	-0.0230145\\
0.142	-0.0232387\\
0.144	-0.0234242\\
0.146	-0.0235707\\
0.148	-0.0236778\\
0.15	-0.0237452\\
0.152	-0.0237729\\
0.154	-0.0237609\\
0.156	-0.0237094\\
0.158	-0.0236186\\
0.16	-0.0234889\\
0.162	-0.0233207\\
0.164	-0.0231146\\
0.166	-0.0228713\\
0.168	-0.0225913\\
0.17	-0.0222755\\
0.172	-0.0219248\\
0.174	-0.0215399\\
0.176	-0.0211219\\
0.178	-0.0206717\\
0.18	-0.0201905\\
0.182	-0.0196792\\
0.184	-0.0191389\\
0.186	-0.0185708\\
0.188	-0.0179761\\
0.19	-0.0173559\\
0.192	-0.0167114\\
0.194	-0.0160439\\
0.196	-0.0153545\\
0.198	-0.0146444\\
0.2	-0.013915\\
0.202	-0.0131674\\
0.204	-0.012403\\
0.206	-0.0116229\\
0.208	-0.0108284\\
0.21	-0.0100207\\
0.212	-0.00920105\\
0.214	-0.00837072\\
0.216	-0.00753091\\
0.218	-0.00668283\\
0.22	-0.00582769\\
0.222	-0.00496669\\
0.224	-0.00410101\\
0.226	-0.00323181\\
0.228	-0.00236026\\
0.23	-0.00148749\\
0.232	-0.000614643\\
0.234	0.000257182\\
0.236	0.000414422\\
0.238	0.000309437\\
0.24	0.000240617\\
0.242	0.000206241\\
0.244	0.000204411\\
0.246	0.000233073\\
0.248	-0.000747692\\
0.25	-0.0020334\\
0.252	-0.00318786\\
0.254	-0.00421363\\
0.256	-0.00511347\\
0.258	-0.00589034\\
0.26	-0.00654738\\
0.262	-0.00700379\\
0.264	-0.00725312\\
0.266	-0.00740482\\
0.268	-0.00746091\\
0.27	-0.00742362\\
0.272	-0.00729539\\
0.274	-0.00707887\\
0.276	-0.00677692\\
0.278	-0.00631855\\
0.28	-0.00548351\\
0.282	-0.0045888\\
0.284	-0.00363865\\
0.286	-0.00263735\\
0.288	-0.00158923\\
0.29	-0.00128318\\
0.292	-0.00121749\\
0.294	-0.00115081\\
0.296	-0.00146226\\
0.298	-0.00260923\\
0.3	-0.00370746\\
0.302	-0.00475711\\
0.304	-0.00575837\\
0.306	-0.00671145\\
0.308	-0.00761658\\
0.31	-0.00847403\\
0.312	-0.00928408\\
0.314	-0.010047\\
0.316	-0.0107633\\
0.318	-0.0114331\\
0.32	-0.012057\\
0.322	-0.0126353\\
0.324	-0.0131686\\
0.326	-0.0136573\\
0.328	-0.0141018\\
0.33	-0.0145029\\
0.332	-0.014861\\
0.334	-0.0151767\\
0.336	-0.0154507\\
0.338	-0.0156837\\
0.34	-0.0158763\\
0.342	-0.0157863\\
0.344	-0.0155987\\
0.346	-0.0153892\\
0.348	-0.0151587\\
0.35	-0.0149082\\
0.352	-0.0146384\\
0.354	-0.014948\\
0.356	-0.0154253\\
0.358	-0.0158573\\
0.36	-0.0162431\\
0.362	-0.0165825\\
0.364	-0.016875\\
0.366	-0.0171203\\
0.368	-0.0173184\\
0.37	-0.0174693\\
0.372	-0.0175729\\
0.374	-0.0176296\\
0.376	-0.0176396\\
0.378	-0.0176035\\
0.38	-0.0175215\\
0.382	-0.0173946\\
0.384	-0.0172232\\
0.386	-0.0170083\\
0.388	-0.0167509\\
0.39	-0.0164518\\
0.392	-0.0161122\\
0.394	-0.0157333\\
0.396	-0.0153163\\
0.398	-0.0148627\\
0.4	-0.0143737\\
0.402	-0.0138509\\
0.404	-0.0132958\\
0.406	-0.0127101\\
0.408	-0.0120954\\
0.41	-0.0114533\\
0.412	-0.0107858\\
0.414	-0.0100945\\
0.416	-0.00938133\\
0.418	-0.00864816\\
0.42	-0.00789688\\
0.422	-0.00712942\\
0.424	-0.00634773\\
0.426	-0.00555375\\
0.428	-0.00474947\\
0.43	-0.00393685\\
0.432	-0.00311786\\
0.434	-0.00229448\\
0.436	-0.00146866\\
0.438	-0.000642336\\
0.44	0.000182552\\
0.442	0.000618824\\
0.444	0.000409254\\
0.446	-5.75932e-05\\
0.448	-0.000749836\\
0.45	-0.00142658\\
0.452	-0.00208715\\
0.454	-0.0027309\\
0.456	-0.00335721\\
0.458	-0.00396551\\
0.46	-0.00455523\\
0.462	-0.00512586\\
0.464	-0.0056769\\
0.466	-0.00620791\\
0.468	-0.00705911\\
0.47	-0.00788716\\
0.472	-0.00868137\\
0.474	-0.00944057\\
0.476	-0.0101637\\
0.478	-0.0108497\\
0.48	-0.0114067\\
0.482	-0.0119032\\
0.484	-0.0123647\\
0.486	-0.0127906\\
0.488	-0.0131804\\
0.49	-0.0135337\\
0.492	-0.0138501\\
0.494	-0.0141294\\
0.496	-0.0143714\\
0.498	-0.014576\\
0.5	-0.0147433\\
0.502	-0.0148732\\
0.504	-0.014966\\
0.506	-0.0150218\\
0.508	-0.0150412\\
0.51	-0.0150243\\
0.512	-0.0149717\\
0.514	-0.0148839\\
0.516	-0.0147616\\
0.518	-0.0146055\\
0.52	-0.0144162\\
0.522	-0.0141947\\
0.524	-0.0139417\\
0.526	-0.0136582\\
0.528	-0.0133452\\
0.53	-0.0130038\\
0.532	-0.012635\\
0.534	-0.0122401\\
0.536	-0.01182\\
0.538	-0.0113762\\
0.54	-0.0109098\\
0.542	-0.0104221\\
0.544	-0.00991455\\
0.546	-0.00938839\\
0.548	-0.00884504\\
0.55	-0.00828589\\
0.552	-0.00771235\\
0.554	-0.00712587\\
0.556	-0.00681253\\
0.558	-0.00649938\\
0.56	-0.00617976\\
0.562	-0.00585425\\
0.564	-0.00552345\\
0.566	-0.00518792\\
0.568	-0.00484827\\
0.57	-0.00450507\\
0.572	-0.00415892\\
0.574	-0.00381039\\
0.576	-0.00346007\\
0.578	-0.00310854\\
0.58	-0.00275638\\
0.582	-0.00240414\\
0.584	-0.00205241\\
0.586	-0.00170172\\
0.588	-0.00135263\\
0.59	-0.0012664\\
0.592	-0.00194173\\
0.594	-0.0026026\\
0.596	-0.00324784\\
0.598	-0.0038763\\
0.6	-0.00448691\\
0.602	-0.00507864\\
0.604	-0.0056505\\
0.606	-0.00620158\\
0.608	-0.00673102\\
0.61	-0.00723801\\
0.612	-0.00772181\\
0.614	-0.00818172\\
0.616	-0.00861712\\
0.618	-0.00902745\\
0.62	-0.00939355\\
0.622	-0.00967183\\
0.624	-0.00992654\\
0.626	-0.0101574\\
0.628	-0.0103642\\
0.63	-0.0105468\\
0.632	-0.0107051\\
0.634	-0.0108391\\
0.636	-0.0109486\\
0.638	-0.0110339\\
0.64	-0.011095\\
0.642	-0.011132\\
0.644	-0.0111451\\
0.646	-0.0111347\\
0.648	-0.011101\\
0.65	-0.0110443\\
0.652	-0.010965\\
0.654	-0.0108637\\
0.656	-0.0107407\\
0.658	-0.0105965\\
0.66	-0.0104318\\
0.662	-0.0102471\\
0.664	-0.010043\\
0.666	-0.00982018\\
0.668	-0.00957939\\
0.67	-0.00932131\\
0.672	-0.00904669\\
0.674	-0.00875632\\
0.676	-0.008451\\
0.678	-0.00813154\\
0.68	-0.0077988\\
0.682	-0.00745362\\
0.684	-0.0070969\\
0.686	-0.00672953\\
0.688	-0.0063524\\
0.69	-0.00596642\\
0.692	-0.00557253\\
0.694	-0.00532011\\
0.696	-0.0052163\\
0.698	-0.00510516\\
0.7	-0.004987\\
0.702	-0.00486211\\
0.704	-0.00473082\\
0.706	-0.00459344\\
0.708	-0.0044503\\
0.71	-0.00430172\\
0.712	-0.00414805\\
0.714	-0.00398963\\
0.716	-0.00382679\\
0.718	-0.00365988\\
0.72	-0.00348925\\
0.722	-0.00331525\\
0.724	-0.00313822\\
0.726	-0.00295853\\
0.728	-0.0027765\\
0.73	-0.00259251\\
0.732	-0.00240688\\
0.734	-0.00232823\\
0.736	-0.00272053\\
0.738	-0.0031012\\
0.74	-0.00346961\\
0.742	-0.00381665\\
0.744	-0.00414491\\
0.746	-0.00445985\\
0.748	-0.00476099\\
0.75	-0.00504788\\
0.752	-0.00532011\\
0.754	-0.00557732\\
0.756	-0.00581916\\
0.758	-0.00604534\\
0.76	-0.00625561\\
0.762	-0.00644975\\
0.764	-0.00662757\\
0.766	-0.00678893\\
0.768	-0.00693372\\
0.77	-0.00706164\\
0.772	-0.0071652\\
0.774	-0.00724411\\
0.776	-0.00730058\\
0.778	-0.00734131\\
0.78	-0.00736643\\
0.782	-0.0073761\\
0.784	-0.00737049\\
0.786	-0.00734983\\
0.788	-0.00731436\\
0.79	-0.00726436\\
0.792	-0.00720011\\
0.794	-0.00712196\\
0.796	-0.00703025\\
0.798	-0.00692535\\
0.8	-0.00680767\\
0.802	-0.00667762\\
0.804	-0.00653566\\
0.806	-0.00638223\\
0.808	-0.00621781\\
0.81	-0.00604291\\
0.812	-0.00585803\\
0.814	-0.00566369\\
0.816	-0.00546044\\
0.818	-0.00524883\\
0.82	-0.00502941\\
0.822	-0.00480275\\
0.824	-0.00456943\\
0.826	-0.00433003\\
0.828	-0.00408514\\
0.83	-0.00383535\\
0.832	-0.00358124\\
0.834	-0.00332342\\
0.836	-0.00313994\\
0.838	-0.00310763\\
0.84	-0.00307036\\
0.842	-0.00302827\\
0.844	-0.00298152\\
0.846	-0.00293026\\
0.848	-0.00287464\\
0.85	-0.00281485\\
0.852	-0.00275104\\
0.854	-0.0026834\\
0.856	-0.0026121\\
0.858	-0.00253585\\
0.86	-0.00245144\\
0.862	-0.00236406\\
0.864	-0.00227391\\
0.866	-0.00218119\\
0.868	-0.00208612\\
0.87	-0.0019889\\
0.872	-0.00188975\\
0.874	-0.00178888\\
0.876	-0.00176039\\
0.878	-0.00197144\\
0.88	-0.00217501\\
0.882	-0.00237077\\
0.884	-0.00255844\\
0.886	-0.00273775\\
0.888	-0.00290843\\
0.89	-0.00307027\\
0.892	-0.00322306\\
0.894	-0.00336661\\
0.896	-0.00350076\\
0.898	-0.00362537\\
0.9	-0.00374034\\
0.902	-0.00384555\\
0.904	-0.00394095\\
0.906	-0.00402648\\
0.908	-0.00410212\\
0.91	-0.00416787\\
0.912	-0.00422372\\
0.914	-0.00426973\\
0.916	-0.00430595\\
0.918	-0.00433246\\
0.92	-0.00434935\\
0.922	-0.00435674\\
0.924	-0.00435476\\
0.926	-0.00434356\\
0.928	-0.00431876\\
0.93	-0.00428412\\
0.932	-0.00424105\\
0.934	-0.00418975\\
0.936	-0.00413047\\
0.938	-0.00406345\\
0.94	-0.00398895\\
0.942	-0.00390725\\
0.944	-0.00381862\\
0.946	-0.00372336\\
0.948	-0.00362177\\
0.95	-0.00351417\\
0.952	-0.00340087\\
0.954	-0.00328219\\
0.956	-0.00315848\\
0.958	-0.00303007\\
0.96	-0.0028973\\
0.962	-0.00276051\\
0.964	-0.00262007\\
0.966	-0.00247632\\
0.968	-0.00232961\\
0.97	-0.0021803\\
0.972	-0.00202875\\
0.974	-0.00187531\\
0.976	-0.00172033\\
0.978	-0.00165009\\
0.98	-0.0016431\\
0.982	-0.00163313\\
0.984	-0.00162025\\
0.986	-0.00160453\\
0.988	-0.00158606\\
0.99	-0.00156492\\
0.992	-0.00154118\\
0.994	-0.00151496\\
0.996	-0.00148634\\
0.998	-0.00145541\\
1	-0.00142229\\
1.002	-0.00138706\\
1.004	-0.00134984\\
1.006	-0.00131074\\
1.008	-0.00126986\\
1.01	-0.00122731\\
1.012	-0.00118321\\
1.014	-0.00113768\\
1.016	-0.00114018\\
1.018	-0.0012592\\
1.02	-0.00137183\\
1.022	-0.00147974\\
1.024	-0.00158279\\
1.026	-0.00168083\\
1.028	-0.00177373\\
1.03	-0.00186138\\
1.032	-0.00194368\\
1.034	-0.00202054\\
1.036	-0.00209187\\
1.038	-0.00215762\\
1.04	-0.00221773\\
1.042	-0.00227217\\
1.044	-0.00232091\\
1.046	-0.00236394\\
1.048	-0.00240126\\
1.05	-0.00243288\\
1.052	-0.00245882\\
1.054	-0.00247913\\
1.056	-0.00249384\\
1.058	-0.00250301\\
1.06	-0.00250672\\
1.062	-0.00250504\\
1.064	-0.00249807\\
1.066	-0.0024859\\
1.068	-0.00246864\\
1.07	-0.00244641\\
1.072	-0.00241934\\
1.074	-0.00238756\\
1.076	-0.00235122\\
1.078	-0.00231047\\
1.08	-0.00226508\\
1.082	-0.00221468\\
1.084	-0.00216041\\
1.086	-0.00210245\\
1.088	-0.00204098\\
1.09	-0.00197619\\
1.092	-0.00190827\\
1.094	-0.00183742\\
1.096	-0.00176382\\
1.098	-0.00168768\\
1.1	-0.00160921\\
1.102	-0.0015286\\
1.104	-0.00144606\\
1.106	-0.0013618\\
1.108	-0.00127603\\
1.11	-0.00118894\\
1.112	-0.00110075\\
1.114	-0.00101166\\
1.116	-0.00092188\\
1.118	-0.000831605\\
1.12	-0.000806126\\
1.122	-0.000805245\\
1.124	-0.000802792\\
1.126	-0.000798803\\
1.128	-0.000793315\\
1.13	-0.000786369\\
1.132	-0.000778008\\
1.134	-0.000768277\\
1.136	-0.000757221\\
1.138	-0.00074489\\
1.14	-0.000731333\\
1.142	-0.000716601\\
1.144	-0.000700747\\
1.146	-0.000683826\\
1.148	-0.000665891\\
1.15	-0.000647\\
1.152	-0.000627209\\
1.154	-0.000649692\\
1.156	-0.00071694\\
1.158	-0.000781537\\
1.16	-0.000843386\\
1.162	-0.000902107\\
1.164	-0.000957555\\
1.166	-0.00101004\\
1.168	-0.00105949\\
1.17	-0.00110586\\
1.172	-0.00114909\\
1.174	-0.00118915\\
1.176	-0.001226\\
1.178	-0.0012596\\
1.18	-0.00128996\\
1.182	-0.00131704\\
1.184	-0.00134085\\
1.186	-0.00136139\\
1.188	-0.00137866\\
1.19	-0.00139269\\
1.192	-0.0014035\\
1.194	-0.00141112\\
1.196	-0.00141558\\
1.198	-0.00141692\\
1.2	-0.0014152\\
1.202	-0.00141047\\
1.204	-0.00140278\\
1.206	-0.00139221\\
1.208	-0.00137882\\
1.21	-0.00136269\\
1.212	-0.00134389\\
1.214	-0.00132252\\
1.216	-0.00129866\\
1.218	-0.0012724\\
1.22	-0.00124384\\
1.222	-0.00121308\\
1.224	-0.00118022\\
1.226	-0.00114537\\
1.228	-0.00110863\\
1.23	-0.00107011\\
1.232	-0.00102994\\
1.234	-0.000988201\\
1.236	-0.000944788\\
1.238	-0.00090007\\
1.24	-0.000854166\\
1.242	-0.000807193\\
1.244	-0.000759271\\
1.246	-0.000710517\\
1.248	-0.000661052\\
1.25	-0.000610993\\
1.252	-0.000560457\\
1.254	-0.000509562\\
1.256	-0.000458423\\
1.258	-0.000407154\\
1.26	-0.000373136\\
1.262	-0.000373861\\
1.264	-0.000373794\\
1.266	-0.00037295\\
1.268	-0.000371349\\
1.27	-0.000369008\\
1.272	-0.000365948\\
1.274	-0.00036219\\
1.276	-0.000357755\\
1.278	-0.000352667\\
1.28	-0.00034695\\
1.282	-0.000340629\\
1.284	-0.000333729\\
1.286	-0.000326276\\
1.288	-0.000318298\\
1.29	-0.000340318\\
1.292	-0.000378881\\
1.294	-0.000415971\\
1.296	-0.000451534\\
1.298	-0.000485518\\
1.3	-0.000517876\\
1.302	-0.000548566\\
1.304	-0.000577548\\
1.306	-0.000604789\\
1.308	-0.000630259\\
1.31	-0.000653933\\
1.312	-0.000675789\\
1.314	-0.000695811\\
1.316	-0.000713952\\
1.318	-0.000730101\\
1.32	-0.000744397\\
1.322	-0.00075684\\
1.324	-0.000767435\\
1.326	-0.000776189\\
1.328	-0.000783114\\
1.33	-0.000788226\\
1.332	-0.000791545\\
1.334	-0.000793092\\
1.336	-0.000792894\\
1.338	-0.000790982\\
1.34	-0.000787387\\
1.342	-0.000782146\\
1.344	-0.000775297\\
1.346	-0.000766883\\
1.348	-0.000756948\\
1.35	-0.000745539\\
1.352	-0.000732705\\
1.354	-0.000718498\\
1.356	-0.000702972\\
1.358	-0.000686183\\
1.36	-0.000668189\\
1.362	-0.000649047\\
1.364	-0.000628821\\
1.366	-0.00060757\\
1.368	-0.00058536\\
1.37	-0.000562253\\
1.372	-0.000538316\\
1.374	-0.000513613\\
1.376	-0.000488213\\
1.378	-0.00046218\\
1.38	-0.000435584\\
1.382	-0.00040849\\
1.384	-0.000380966\\
1.386	-0.000353079\\
1.388	-0.000324895\\
1.39	-0.000296482\\
1.392	-0.000267869\\
1.394	-0.000239155\\
1.396	-0.000210407\\
1.398	-0.000181688\\
1.4	-0.000172136\\
1.402	-0.000172285\\
1.404	-0.000172047\\
1.406	-0.000171432\\
1.408	-0.000170448\\
1.41	-0.000169107\\
1.412	-0.000167417\\
1.414	-0.000165391\\
1.416	-0.000163039\\
1.418	-0.000160374\\
1.42	-0.000157408\\
1.422	-0.000154153\\
1.424	-0.000161417\\
1.426	-0.000183983\\
1.428	-0.000205766\\
1.43	-0.000226732\\
1.432	-0.000246848\\
1.434	-0.000266085\\
1.436	-0.000284417\\
1.438	-0.000301818\\
1.44	-0.000318266\\
1.442	-0.000333742\\
1.444	-0.000348228\\
1.446	-0.000361708\\
1.448	-0.000374158\\
1.45	-0.000385566\\
1.452	-0.00039594\\
1.454	-0.000405275\\
1.456	-0.000413568\\
1.458	-0.000420818\\
1.46	-0.000427027\\
1.462	-0.000432199\\
1.464	-0.00043634\\
1.466	-0.000439458\\
1.468	-0.000441565\\
1.47	-0.000442671\\
1.472	-0.000442792\\
1.474	-0.000441944\\
1.476	-0.000440145\\
1.478	-0.000437415\\
1.48	-0.00043376\\
1.482	-0.000429217\\
1.484	-0.000423814\\
1.486	-0.000417576\\
1.488	-0.000410531\\
1.49	-0.000402708\\
1.492	-0.000394137\\
1.494	-0.00038485\\
1.496	-0.000374877\\
1.498	-0.000364254\\
1.5	-0.000353013\\
1.502	-0.00034119\\
1.504	-0.000328821\\
1.506	-0.000315941\\
1.508	-0.000302587\\
1.51	-0.000288797\\
1.512	-0.000274607\\
1.514	-0.000260055\\
1.516	-0.00024518\\
1.518	-0.00023002\\
1.52	-0.000214611\\
1.522	-0.000198992\\
1.524	-0.000183201\\
1.526	-0.000167275\\
1.528	-0.00015125\\
1.53	-0.000135165\\
1.532	-0.000119054\\
1.534	-0.000102954\\
1.536	-8.68996e-05\\
1.538	-8.1992e-05\\
1.54	-8.17425e-05\\
1.542	-8.13052e-05\\
1.544	-8.06848e-05\\
1.546	-7.98865e-05\\
1.548	-7.89157e-05\\
1.55	-7.77779e-05\\
1.552	-7.64792e-05\\
1.554	-7.50258e-05\\
1.556	-7.3424e-05\\
1.558	-7.71172e-05\\
1.56	-9.02181e-05\\
1.562	-0.000102901\\
1.564	-0.000115144\\
1.566	-0.000126929\\
1.568	-0.000138238\\
1.57	-0.000149054\\
1.572	-0.00015936\\
1.574	-0.000169144\\
1.576	-0.000178393\\
1.578	-0.000187095\\
1.58	-0.000195241\\
1.582	-0.000202817\\
1.584	-0.000209816\\
1.586	-0.000216237\\
1.588	-0.000222077\\
1.59	-0.000227332\\
1.592	-0.000232\\
1.594	-0.000236082\\
1.596	-0.000239577\\
1.598	-0.000242488\\
1.6	-0.000244819\\
1.602	-0.000246574\\
1.604	-0.000247759\\
1.606	-0.000248381\\
1.608	-0.000248447\\
1.61	-0.000247968\\
1.612	-0.000246953\\
1.614	-0.000245414\\
1.616	-0.000243362\\
1.618	-0.000240811\\
1.62	-0.000237775\\
1.622	-0.000234268\\
1.624	-0.000230306\\
1.626	-0.000225905\\
1.628	-0.000221083\\
1.63	-0.000215856\\
1.632	-0.000210244\\
1.634	-0.000204264\\
1.636	-0.000197937\\
1.638	-0.000191281\\
1.64	-0.000184317\\
1.642	-0.000177066\\
1.644	-0.000169547\\
1.646	-0.000161782\\
1.648	-0.000153792\\
1.65	-0.000145598\\
1.652	-0.000137219\\
1.654	-0.000128678\\
1.656	-0.000119999\\
1.658	-0.000111201\\
1.66	-0.000102306\\
1.662	-9.33357e-05\\
1.664	-8.43104e-05\\
1.666	-7.52509e-05\\
1.668	-6.61775e-05\\
1.67	-5.71105e-05\\
1.672	-4.80695e-05\\
1.674	-4.27079e-05\\
1.676	-4.23287e-05\\
1.678	-4.18537e-05\\
1.68	-4.12857e-05\\
1.682	-4.06276e-05\\
1.684	-3.98826e-05\\
1.686	-3.9054e-05\\
1.688	-3.81451e-05\\
1.69	-3.71596e-05\\
1.692	-3.6772e-05\\
1.694	-4.43601e-05\\
1.696	-5.17246e-05\\
1.698	-5.88531e-05\\
1.7	-6.5734e-05\\
1.702	-7.23564e-05\\
1.704	-7.871e-05\\
1.706	-8.47854e-05\\
1.708	-9.0574e-05\\
1.71	-9.60679e-05\\
1.712	-0.00010126\\
1.714	-0.000106144\\
1.716	-0.000110714\\
1.718	-0.000114966\\
1.72	-0.000118895\\
1.722	-0.000122499\\
1.724	-0.000125774\\
1.726	-0.00012872\\
1.728	-0.000131335\\
1.73	-0.000133618\\
1.732	-0.000135571\\
1.734	-0.000137195\\
1.736	-0.000138492\\
1.738	-0.000139464\\
1.74	-0.000140114\\
1.742	-0.000140445\\
1.744	-0.000140463\\
1.746	-0.000140173\\
1.748	-0.000139581\\
1.75	-0.000138693\\
1.752	-0.000137517\\
1.754	-0.000136059\\
1.756	-0.000134327\\
1.758	-0.000132329\\
1.76	-0.000130076\\
1.762	-0.000127574\\
1.764	-0.000124835\\
1.766	-0.000121868\\
1.768	-0.000118683\\
1.77	-0.000115291\\
1.772	-0.000111702\\
1.774	-0.000107929\\
1.776	-0.000103981\\
1.778	-9.98708e-05\\
1.78	-9.56099e-05\\
1.782	-9.12101e-05\\
1.784	-8.66832e-05\\
1.786	-8.20412e-05\\
1.788	-7.72962e-05\\
1.79	-7.24602e-05\\
1.792	-6.75452e-05\\
1.794	-6.25634e-05\\
1.796	-5.75266e-05\\
1.798	-5.24468e-05\\
1.8	-4.73357e-05\\
1.802	-4.2205e-05\\
1.804	-3.70663e-05\\
1.806	-3.19309e-05\\
1.808	-2.68098e-05\\
1.81	-2.44272e-05\\
1.812	-2.40432e-05\\
1.814	-2.36081e-05\\
1.816	-2.31238e-05\\
1.818	-2.25916e-05\\
1.82	-2.20139e-05\\
1.822	-2.1393e-05\\
1.824	-2.0731e-05\\
1.826	-2.00301e-05\\
1.828	-2.13156e-05\\
1.83	-2.56178e-05\\
1.832	-2.97937e-05\\
1.834	-3.38363e-05\\
1.836	-3.77391e-05\\
1.838	-4.14957e-05\\
1.84	-4.51003e-05\\
1.842	-4.85477e-05\\
1.844	-5.18328e-05\\
1.846	-5.49512e-05\\
1.848	-5.78987e-05\\
1.85	-6.06718e-05\\
1.852	-6.32673e-05\\
1.854	-6.56824e-05\\
1.856	-6.79147e-05\\
1.858	-6.99625e-05\\
1.86	-7.18242e-05\\
1.862	-7.34988e-05\\
1.864	-7.49858e-05\\
1.866	-7.62849e-05\\
1.868	-7.73964e-05\\
1.87	-7.83209e-05\\
1.872	-7.90595e-05\\
1.874	-7.96136e-05\\
1.876	-7.99849e-05\\
1.878	-8.01756e-05\\
1.88	-8.0188e-05\\
1.882	-8.00251e-05\\
1.884	-7.96901e-05\\
1.886	-7.91868e-05\\
1.888	-7.85188e-05\\
1.89	-7.76904e-05\\
1.892	-7.67058e-05\\
1.894	-7.55699e-05\\
1.896	-7.42875e-05\\
1.898	-7.28638e-05\\
1.9	-7.13042e-05\\
1.902	-6.96143e-05\\
1.904	-6.77999e-05\\
1.906	-6.58669e-05\\
1.908	-6.38214e-05\\
1.91	-6.16698e-05\\
1.912	-5.94185e-05\\
1.914	-5.70739e-05\\
1.916	-5.46427e-05\\
1.918	-5.21315e-05\\
1.92	-4.95472e-05\\
1.922	-4.68965e-05\\
1.924	-4.41862e-05\\
1.926	-4.14233e-05\\
1.928	-3.86146e-05\\
1.93	-3.57669e-05\\
1.932	-3.28871e-05\\
1.934	-2.9982e-05\\
1.936	-2.70583e-05\\
1.938	-2.41226e-05\\
1.94	-2.11816e-05\\
1.942	-1.82417e-05\\
1.944	-1.59201e-05\\
1.946	-1.64313e-05\\
1.948	-1.68851e-05\\
1.95	-1.72806e-05\\
1.952	-1.76171e-05\\
1.954	-1.78942e-05\\
1.956	-1.81116e-05\\
1.958	-1.82692e-05\\
1.96	-1.83671e-05\\
1.962	-1.84057e-05\\
1.964	-1.83855e-05\\
1.966	-1.83073e-05\\
1.968	-1.8172e-05\\
1.97	-1.94854e-05\\
1.972	-2.17287e-05\\
1.974	-2.38888e-05\\
1.976	-2.59621e-05\\
1.978	-2.79456e-05\\
1.98	-2.98363e-05\\
1.982	-3.16318e-05\\
1.984	-3.33296e-05\\
1.986	-3.49274e-05\\
1.988	-3.64234e-05\\
1.99	-3.78161e-05\\
1.992	-3.9104e-05\\
1.994	-4.02859e-05\\
1.996	-4.13612e-05\\
1.998	-4.23291e-05\\
2	-4.31893e-05\\
};
\addplot [color=mycolor6, line width=1.0pt, forget plot]
  table[row sep=crcr]{%
-0.024	0\\
-0.022	0\\
-0.02	0\\
-0.018	0\\
-0.016	0\\
-0.014	0\\
-0.012	0\\
-0.01	0\\
-0.008	0\\
-0.006	0\\
-0.004	0\\
-0.002	0\\
0	0\\
0.002	2.93074e-10\\
0.004	-2.49422e-09\\
0.006	-2.40547e-08\\
0.008	-9.28215e-08\\
0.01	-2.494e-07\\
0.012	-5.46454e-07\\
0.014	-1.04889e-06\\
0.016	-1.83421e-06\\
0.018	-2.9929e-06\\
0.02	-4.62877e-06\\
0.022	-6.85916e-06\\
0.024	-9.81494e-06\\
0.026	-1.36403e-05\\
0.028	-1.8492e-05\\
0.03	-2.4539e-05\\
0.032	-3.19609e-05\\
0.034	-4.09467e-05\\
0.036	-5.16934e-05\\
0.038	-6.44038e-05\\
0.04	-1.17136e-05\\
0.042	-1.64395e-05\\
0.044	-2.2479e-05\\
0.046	-3.00791e-05\\
0.048	-3.95129e-05\\
0.05	-5.10801e-05\\
0.052	-6.51056e-05\\
0.054	-8.19396e-05\\
0.056	-0.000101955\\
0.058	-0.000125548\\
0.06	-0.000153131\\
0.062	-0.000185137\\
0.064	-0.00022201\\
0.066	-0.000264205\\
0.068	-0.000312182\\
0.07	-0.000366405\\
0.072	-0.000427333\\
0.074	-0.000495419\\
0.076	-0.000571103\\
0.078	-0.00687464\\
0.08	-0.0159481\\
0.082	-0.0207402\\
0.084	-0.0213857\\
0.086	-0.0219956\\
0.088	-0.0225672\\
0.09	-0.0230983\\
0.092	-0.0235866\\
0.094	-0.0240301\\
0.096	-0.0244271\\
0.098	-0.0247759\\
0.1	-0.0250749\\
0.102	-0.025323\\
0.104	-0.0255189\\
0.106	-0.0256615\\
0.108	-0.0257502\\
0.11	-0.025784\\
0.112	-0.0248592\\
0.114	-0.023262\\
0.116	-0.0216742\\
0.118	-0.020094\\
0.12	-0.0201519\\
0.122	-0.0207116\\
0.124	-0.021243\\
0.126	-0.0217441\\
0.128	-0.0222132\\
0.13	-0.0226486\\
0.132	-0.0230488\\
0.134	-0.0234124\\
0.136	-0.0237382\\
0.138	-0.0240253\\
0.14	-0.0242727\\
0.142	-0.0244798\\
0.144	-0.0246459\\
0.146	-0.0247706\\
0.148	-0.0248537\\
0.15	-0.024895\\
0.152	-0.0248945\\
0.154	-0.0248524\\
0.156	-0.0247688\\
0.158	-0.0246442\\
0.16	-0.0244789\\
0.162	-0.0242736\\
0.164	-0.0240289\\
0.166	-0.0237455\\
0.168	-0.0234243\\
0.17	-0.0230661\\
0.172	-0.0226719\\
0.174	-0.0222428\\
0.176	-0.0217796\\
0.178	-0.0212837\\
0.18	-0.0207562\\
0.182	-0.0201981\\
0.184	-0.0196108\\
0.186	-0.0189956\\
0.188	-0.0183536\\
0.19	-0.0176861\\
0.192	-0.0169946\\
0.194	-0.0162803\\
0.196	-0.0155444\\
0.198	-0.0147885\\
0.2	-0.0140137\\
0.202	-0.0132214\\
0.204	-0.012413\\
0.206	-0.0115898\\
0.208	-0.0107531\\
0.21	-0.00990425\\
0.212	-0.00904449\\
0.214	-0.00817515\\
0.216	-0.00729751\\
0.218	-0.00641284\\
0.22	-0.00552239\\
0.222	-0.0046274\\
0.224	-0.00372908\\
0.226	-0.00282864\\
0.228	-0.00192725\\
0.23	-0.00102609\\
0.232	-0.000126283\\
0.234	0.000771042\\
0.236	0.000952321\\
0.238	0.000869913\\
0.24	0.000822217\\
0.242	0.000807517\\
0.244	0.000823931\\
0.246	0.00086942\\
0.248	-9.59121e-05\\
0.25	-0.00136756\\
0.252	-0.0025093\\
0.254	-0.00352366\\
0.256	-0.00441337\\
0.258	-0.00518135\\
0.26	-0.00583071\\
0.262	-0.00636473\\
0.264	-0.00678683\\
0.266	-0.0071006\\
0.268	-0.00721619\\
0.27	-0.00715625\\
0.272	-0.00700628\\
0.274	-0.00676896\\
0.276	-0.00644719\\
0.278	-0.0060441\\
0.28	-0.00542148\\
0.282	-0.0045303\\
0.284	-0.00358402\\
0.286	-0.00258689\\
0.288	-0.00154324\\
0.29	-0.00124192\\
0.292	-0.00118121\\
0.294	-0.00111973\\
0.296	-0.0014366\\
0.298	-0.00258915\\
0.3	-0.00369313\\
0.302	-0.00474866\\
0.304	-0.00575591\\
0.306	-0.00671507\\
0.308	-0.00762633\\
0.31	-0.00848995\\
0.312	-0.00930618\\
0.314	-0.0100753\\
0.316	-0.0107977\\
0.318	-0.0114736\\
0.32	-0.0121035\\
0.322	-0.0126878\\
0.324	-0.0132269\\
0.326	-0.0137212\\
0.328	-0.0141713\\
0.33	-0.0145776\\
0.332	-0.0149409\\
0.334	-0.0152615\\
0.336	-0.0155403\\
0.338	-0.0157777\\
0.34	-0.0158756\\
0.342	-0.0157271\\
0.344	-0.0155546\\
0.346	-0.0153591\\
0.348	-0.0151415\\
0.35	-0.0149027\\
0.352	-0.0146436\\
0.354	-0.0149629\\
0.356	-0.0154488\\
0.358	-0.0158883\\
0.36	-0.0162807\\
0.362	-0.0166255\\
0.364	-0.0169225\\
0.366	-0.0171715\\
0.368	-0.0173722\\
0.37	-0.0175248\\
0.372	-0.0176294\\
0.374	-0.0176862\\
0.376	-0.0176956\\
0.378	-0.017658\\
0.38	-0.0175739\\
0.382	-0.0174441\\
0.384	-0.0172694\\
0.386	-0.0170505\\
0.388	-0.0167884\\
0.39	-0.0164843\\
0.392	-0.0161391\\
0.394	-0.0157542\\
0.396	-0.0153309\\
0.398	-0.0148705\\
0.4	-0.0143745\\
0.402	-0.0138445\\
0.404	-0.0132819\\
0.406	-0.0126884\\
0.408	-0.0120658\\
0.41	-0.0114158\\
0.412	-0.0107402\\
0.414	-0.0100408\\
0.416	-0.0093195\\
0.418	-0.00857822\\
0.42	-0.00781887\\
0.422	-0.00704341\\
0.424	-0.00625382\\
0.426	-0.00545208\\
0.428	-0.00464018\\
0.43	-0.00382012\\
0.432	-0.0029939\\
0.434	-0.00216349\\
0.436	-0.00133088\\
0.438	-0.000498028\\
0.44	0.000333122\\
0.442	0.000775372\\
0.444	0.000488838\\
0.446	-2.47468e-05\\
0.448	-0.000712007\\
0.45	-0.00138386\\
0.452	-0.00203965\\
0.454	-0.00267873\\
0.456	-0.00330048\\
0.458	-0.00390434\\
0.46	-0.00448975\\
0.462	-0.00505619\\
0.464	-0.00560319\\
0.466	-0.00613028\\
0.468	-0.0069777\\
0.47	-0.00780212\\
0.472	-0.00859284\\
0.474	-0.0093487\\
0.476	-0.0100686\\
0.478	-0.0107516\\
0.48	-0.011303\\
0.482	-0.0118045\\
0.484	-0.0122708\\
0.486	-0.0127014\\
0.488	-0.0130958\\
0.49	-0.0134536\\
0.492	-0.0137743\\
0.494	-0.0140578\\
0.496	-0.0143038\\
0.498	-0.0145124\\
0.5	-0.0146834\\
0.502	-0.014817\\
0.504	-0.0149133\\
0.506	-0.0149725\\
0.508	-0.0149951\\
0.51	-0.0149813\\
0.512	-0.0149316\\
0.514	-0.0148467\\
0.516	-0.014727\\
0.518	-0.0145734\\
0.52	-0.0143865\\
0.522	-0.0141672\\
0.524	-0.0139163\\
0.526	-0.0136348\\
0.528	-0.0133237\\
0.53	-0.012984\\
0.532	-0.0126168\\
0.534	-0.0122232\\
0.536	-0.0118045\\
0.538	-0.0113618\\
0.54	-0.0108965\\
0.542	-0.0104098\\
0.544	-0.00990312\\
0.546	-0.00937771\\
0.548	-0.00883499\\
0.55	-0.00827638\\
0.552	-0.00770329\\
0.554	-0.00711715\\
0.556	-0.00680407\\
0.558	-0.00649109\\
0.56	-0.00617157\\
0.562	-0.00584608\\
0.564	-0.00551521\\
0.566	-0.00517956\\
0.568	-0.00483971\\
0.57	-0.00449625\\
0.572	-0.00414978\\
0.574	-0.00380088\\
0.576	-0.00345015\\
0.578	-0.00309815\\
0.58	-0.00274548\\
0.582	-0.0023927\\
0.584	-0.00204038\\
0.586	-0.00168908\\
0.588	-0.00133935\\
0.59	-0.00125246\\
0.592	-0.0019271\\
0.594	-0.00258727\\
0.596	-0.00323179\\
0.598	-0.00385953\\
0.6	-0.0044694\\
0.602	-0.00506039\\
0.604	-0.00563151\\
0.606	-0.00618185\\
0.608	-0.00671055\\
0.61	-0.0072168\\
0.612	-0.00769987\\
0.614	-0.00815906\\
0.616	-0.00859376\\
0.618	-0.00900339\\
0.62	-0.00938746\\
0.622	-0.00967667\\
0.624	-0.00993171\\
0.626	-0.0101628\\
0.628	-0.0103698\\
0.63	-0.0105525\\
0.632	-0.0107108\\
0.634	-0.0108447\\
0.636	-0.0109542\\
0.638	-0.0110393\\
0.64	-0.0111001\\
0.642	-0.0111368\\
0.644	-0.0111496\\
0.646	-0.0111388\\
0.648	-0.0111047\\
0.65	-0.0110475\\
0.652	-0.0109677\\
0.654	-0.0108658\\
0.656	-0.0107422\\
0.658	-0.0105975\\
0.66	-0.0104321\\
0.662	-0.0102467\\
0.664	-0.010042\\
0.666	-0.00981852\\
0.668	-0.00957702\\
0.67	-0.00931823\\
0.672	-0.0090429\\
0.674	-0.0087518\\
0.676	-0.00844574\\
0.678	-0.00812555\\
0.68	-0.00779207\\
0.682	-0.00744617\\
0.684	-0.00708872\\
0.686	-0.00672062\\
0.688	-0.00634277\\
0.69	-0.00595609\\
0.692	-0.0055615\\
0.694	-0.00530839\\
0.696	-0.00520391\\
0.698	-0.00509211\\
0.7	-0.00497331\\
0.702	-0.00484779\\
0.704	-0.00471589\\
0.706	-0.00457792\\
0.708	-0.0044342\\
0.71	-0.00428507\\
0.712	-0.00413087\\
0.714	-0.00397193\\
0.716	-0.00380859\\
0.718	-0.00364121\\
0.72	-0.00347013\\
0.722	-0.00329571\\
0.724	-0.00311828\\
0.726	-0.0029382\\
0.728	-0.00275582\\
0.73	-0.00257148\\
0.732	-0.00238554\\
0.734	-0.0023066\\
0.736	-0.00269863\\
0.738	-0.00307906\\
0.74	-0.00344724\\
0.742	-0.00380259\\
0.744	-0.00414454\\
0.746	-0.00447256\\
0.748	-0.00478404\\
0.75	-0.00507057\\
0.752	-0.00534244\\
0.754	-0.00559929\\
0.756	-0.00584078\\
0.758	-0.00606662\\
0.76	-0.00627654\\
0.762	-0.00647034\\
0.764	-0.00664668\\
0.766	-0.00680111\\
0.768	-0.00693901\\
0.77	-0.00706033\\
0.772	-0.00716506\\
0.774	-0.00725226\\
0.776	-0.00732299\\
0.778	-0.00736395\\
0.78	-0.00738891\\
0.782	-0.00739839\\
0.784	-0.00739258\\
0.786	-0.00737171\\
0.788	-0.00733601\\
0.79	-0.00728576\\
0.792	-0.00722127\\
0.794	-0.00714285\\
0.796	-0.00705085\\
0.798	-0.00694567\\
0.8	-0.00682769\\
0.802	-0.00669733\\
0.804	-0.00655504\\
0.806	-0.00640128\\
0.808	-0.00623653\\
0.81	-0.00606128\\
0.812	-0.00587605\\
0.814	-0.00568135\\
0.816	-0.00547774\\
0.818	-0.00526575\\
0.82	-0.00504596\\
0.822	-0.00481892\\
0.824	-0.00458521\\
0.826	-0.00434542\\
0.828	-0.00410014\\
0.83	-0.00384995\\
0.832	-0.00359545\\
0.834	-0.00333723\\
0.836	-0.00315335\\
0.838	-0.00312064\\
0.84	-0.00308297\\
0.842	-0.00304049\\
0.844	-0.00299334\\
0.846	-0.00294168\\
0.848	-0.00288567\\
0.85	-0.00282549\\
0.852	-0.00275985\\
0.854	-0.00268577\\
0.856	-0.00260811\\
0.858	-0.00252707\\
0.86	-0.00244284\\
0.862	-0.00235563\\
0.864	-0.00226565\\
0.866	-0.0021731\\
0.868	-0.0020782\\
0.87	-0.00198115\\
0.872	-0.00188216\\
0.874	-0.00178144\\
0.876	-0.00175311\\
0.878	-0.00196431\\
0.88	-0.00216803\\
0.882	-0.00236394\\
0.884	-0.00255175\\
0.886	-0.0027312\\
0.888	-0.00290202\\
0.89	-0.003064\\
0.892	-0.00321692\\
0.894	-0.00336059\\
0.896	-0.00349487\\
0.898	-0.00361961\\
0.9	-0.00373469\\
0.902	-0.00384003\\
0.904	-0.00393554\\
0.906	-0.00402118\\
0.908	-0.00409693\\
0.91	-0.00416278\\
0.912	-0.00421874\\
0.914	-0.00426484\\
0.916	-0.00430116\\
0.918	-0.00432776\\
0.92	-0.00434474\\
0.922	-0.00435221\\
0.924	-0.00435032\\
0.926	-0.00433921\\
0.928	-0.00431783\\
0.93	-0.00428326\\
0.932	-0.00424025\\
0.934	-0.00418901\\
0.936	-0.00412977\\
0.938	-0.00406279\\
0.94	-0.00398833\\
0.942	-0.00390665\\
0.944	-0.00381804\\
0.946	-0.0037228\\
0.948	-0.00362122\\
0.95	-0.00351361\\
0.952	-0.00340031\\
0.954	-0.00328162\\
0.956	-0.00315789\\
0.958	-0.00302946\\
0.96	-0.00289667\\
0.962	-0.00275985\\
0.964	-0.00261938\\
0.966	-0.00247559\\
0.968	-0.00232885\\
0.97	-0.0021795\\
0.972	-0.0020279\\
0.974	-0.00187441\\
0.976	-0.00171938\\
0.978	-0.00164909\\
0.98	-0.00164206\\
0.982	-0.00163203\\
0.984	-0.0016191\\
0.986	-0.00160333\\
0.988	-0.0015848\\
0.99	-0.0015636\\
0.992	-0.00153982\\
0.994	-0.00151354\\
0.996	-0.00148486\\
0.998	-0.00145389\\
1	-0.00142071\\
1.002	-0.00138543\\
1.004	-0.00134817\\
1.006	-0.00130902\\
1.008	-0.00126809\\
1.01	-0.0012255\\
1.012	-0.00118136\\
1.014	-0.00113578\\
1.016	-0.00113825\\
1.018	-0.00125764\\
1.02	-0.00137241\\
1.022	-0.00148183\\
1.024	-0.00158484\\
1.026	-0.00168285\\
1.028	-0.00177571\\
1.03	-0.00186332\\
1.032	-0.00194558\\
1.034	-0.0020224\\
1.036	-0.00209369\\
1.038	-0.0021594\\
1.04	-0.00221947\\
1.042	-0.00227387\\
1.044	-0.00232258\\
1.046	-0.00236557\\
1.048	-0.00240285\\
1.05	-0.00243443\\
1.052	-0.00246033\\
1.054	-0.00248059\\
1.056	-0.00249526\\
1.058	-0.0025044\\
1.06	-0.00250807\\
1.062	-0.00250635\\
1.064	-0.00249934\\
1.066	-0.00248713\\
1.068	-0.00246984\\
1.07	-0.00244757\\
1.072	-0.00242047\\
1.074	-0.00238865\\
1.076	-0.00235228\\
1.078	-0.00231149\\
1.08	-0.00226644\\
1.082	-0.00221636\\
1.084	-0.00216205\\
1.086	-0.00210404\\
1.088	-0.00204253\\
1.09	-0.00197769\\
1.092	-0.00190973\\
1.094	-0.00183882\\
1.096	-0.00176517\\
1.098	-0.00168899\\
1.1	-0.00161046\\
1.102	-0.00152981\\
1.104	-0.00144722\\
1.106	-0.00136291\\
1.108	-0.00127708\\
1.11	-0.00118995\\
1.112	-0.00110171\\
1.114	-0.00101257\\
1.116	-0.000922739\\
1.118	-0.000832415\\
1.12	-0.000806888\\
1.122	-0.00080596\\
1.124	-0.000803461\\
1.126	-0.000799425\\
1.128	-0.000793893\\
1.13	-0.000786903\\
1.132	-0.000778499\\
1.134	-0.000768726\\
1.136	-0.00075763\\
1.138	-0.000745259\\
1.14	-0.000731663\\
1.142	-0.000716894\\
1.144	-0.000701005\\
1.146	-0.000684048\\
1.148	-0.00066608\\
1.15	-0.000647157\\
1.152	-0.000627335\\
1.154	-0.000649788\\
1.156	-0.000717009\\
1.158	-0.000781579\\
1.16	-0.000843403\\
1.162	-0.000901933\\
1.164	-0.000957388\\
1.166	-0.00100988\\
1.168	-0.00105934\\
1.17	-0.00110571\\
1.172	-0.00114895\\
1.174	-0.00118901\\
1.176	-0.00122586\\
1.178	-0.00125947\\
1.18	-0.00128983\\
1.182	-0.00131691\\
1.184	-0.00134072\\
1.186	-0.00136126\\
1.188	-0.00137854\\
1.19	-0.00139258\\
1.192	-0.00140339\\
1.194	-0.00141101\\
1.196	-0.00141547\\
1.198	-0.00141681\\
1.2	-0.00141509\\
1.202	-0.00141036\\
1.204	-0.00140268\\
1.206	-0.00139211\\
1.208	-0.00137872\\
1.21	-0.00136259\\
1.212	-0.00134379\\
1.214	-0.00132242\\
1.216	-0.00129856\\
1.218	-0.0012723\\
1.22	-0.00124374\\
1.222	-0.00121298\\
1.224	-0.00118013\\
1.226	-0.00114527\\
1.228	-0.00110854\\
1.23	-0.00107002\\
1.232	-0.00102985\\
1.234	-0.000988123\\
1.236	-0.000944921\\
1.238	-0.000900194\\
1.24	-0.00085428\\
1.242	-0.000807298\\
1.244	-0.000759366\\
1.246	-0.000710604\\
1.248	-0.000661129\\
1.25	-0.00061106\\
1.252	-0.000560514\\
1.254	-0.00050961\\
1.256	-0.000458461\\
1.258	-0.000407182\\
1.26	-0.000373156\\
1.262	-0.000373871\\
1.264	-0.000373794\\
1.266	-0.000372942\\
1.268	-0.000371332\\
1.27	-0.000368983\\
1.272	-0.000365914\\
1.274	-0.000362147\\
1.276	-0.000357705\\
1.278	-0.000352609\\
1.28	-0.000346885\\
1.282	-0.000340556\\
1.284	-0.000333649\\
1.286	-0.00032619\\
1.288	-0.000318206\\
1.29	-0.00034022\\
1.292	-0.000378776\\
1.294	-0.000415861\\
1.296	-0.000451419\\
1.298	-0.000485398\\
1.3	-0.000517752\\
1.302	-0.000548438\\
1.304	-0.000577416\\
1.306	-0.000604654\\
1.308	-0.000630121\\
1.31	-0.000653792\\
1.312	-0.000675645\\
1.314	-0.000695665\\
1.316	-0.000713838\\
1.318	-0.000730147\\
1.32	-0.000744442\\
1.322	-0.000756884\\
1.324	-0.000767478\\
1.326	-0.000776231\\
1.328	-0.000783156\\
1.33	-0.000788267\\
1.332	-0.000791585\\
1.334	-0.000793131\\
1.336	-0.000792933\\
1.338	-0.00079102\\
1.34	-0.000787424\\
1.342	-0.000782183\\
1.344	-0.000775334\\
1.346	-0.000766919\\
1.348	-0.000756983\\
1.35	-0.000745574\\
1.352	-0.000732739\\
1.354	-0.000718532\\
1.356	-0.000703006\\
1.358	-0.000686217\\
1.36	-0.000668221\\
1.362	-0.00064908\\
1.364	-0.000628853\\
1.366	-0.000607602\\
1.368	-0.000585391\\
1.37	-0.000562284\\
1.372	-0.000538347\\
1.374	-0.000513644\\
1.376	-0.000488243\\
1.378	-0.000462211\\
1.38	-0.000435614\\
1.382	-0.00040852\\
1.384	-0.000380996\\
1.386	-0.000353109\\
1.388	-0.000324925\\
1.39	-0.000296511\\
1.392	-0.000267933\\
1.394	-0.000239223\\
1.396	-0.000210471\\
1.398	-0.000181749\\
1.4	-0.000172194\\
1.402	-0.000172339\\
1.404	-0.000172097\\
1.406	-0.000171479\\
1.408	-0.000170492\\
1.41	-0.000169147\\
1.412	-0.000167455\\
1.414	-0.000165425\\
1.416	-0.00016307\\
1.418	-0.000160402\\
1.42	-0.000157433\\
1.422	-0.000154175\\
1.424	-0.000161436\\
1.426	-0.000183999\\
1.428	-0.000205779\\
1.43	-0.000226743\\
1.432	-0.000246857\\
1.434	-0.000266092\\
1.436	-0.000284421\\
1.438	-0.00030182\\
1.44	-0.000318266\\
1.442	-0.00033374\\
1.444	-0.000348224\\
1.446	-0.000361703\\
1.448	-0.000374165\\
1.45	-0.000385599\\
1.452	-0.000395992\\
1.454	-0.000405325\\
1.456	-0.000413616\\
1.458	-0.000420865\\
1.46	-0.000427072\\
1.462	-0.000432242\\
1.464	-0.000436382\\
1.466	-0.000439498\\
1.468	-0.000441603\\
1.47	-0.000442708\\
1.472	-0.000442822\\
1.474	-0.000441955\\
1.476	-0.000440138\\
1.478	-0.00043739\\
1.48	-0.000433734\\
1.482	-0.000429192\\
1.484	-0.000423789\\
1.486	-0.000417551\\
1.488	-0.000410507\\
1.49	-0.000402684\\
1.492	-0.000394113\\
1.494	-0.000384826\\
1.496	-0.000374854\\
1.498	-0.000364231\\
1.5	-0.00035299\\
1.502	-0.000341168\\
1.504	-0.000328799\\
1.506	-0.000315919\\
1.508	-0.000302565\\
1.51	-0.000288775\\
1.512	-0.000274586\\
1.514	-0.000260035\\
1.516	-0.00024516\\
1.518	-0.00023\\
1.52	-0.000214592\\
1.522	-0.000198973\\
1.524	-0.000183182\\
1.526	-0.000167256\\
1.528	-0.000151233\\
1.53	-0.000135148\\
1.532	-0.000119037\\
1.534	-0.000102938\\
1.536	-8.68836e-05\\
1.538	-8.19764e-05\\
1.54	-8.17273e-05\\
1.542	-8.12905e-05\\
1.544	-8.06705e-05\\
1.546	-7.98727e-05\\
1.548	-7.89023e-05\\
1.55	-7.77649e-05\\
1.552	-7.64667e-05\\
1.554	-7.50137e-05\\
1.556	-7.34123e-05\\
1.558	-7.7106e-05\\
1.56	-9.02072e-05\\
1.562	-0.00010289\\
1.564	-0.000115134\\
1.566	-0.00012692\\
1.568	-0.000138229\\
1.57	-0.000149045\\
1.572	-0.000159352\\
1.574	-0.000169136\\
1.576	-0.000178386\\
1.578	-0.000187088\\
1.58	-0.000195234\\
1.582	-0.000202815\\
1.584	-0.000209824\\
1.586	-0.000216247\\
1.588	-0.000222087\\
1.59	-0.000227341\\
1.592	-0.00023201\\
1.594	-0.000236091\\
1.596	-0.000239586\\
1.598	-0.000242497\\
1.6	-0.000244828\\
1.602	-0.000246583\\
1.604	-0.000247768\\
1.606	-0.000248389\\
1.608	-0.000248456\\
1.61	-0.000247977\\
1.612	-0.000246962\\
1.614	-0.000245422\\
1.616	-0.00024337\\
1.618	-0.000240819\\
1.62	-0.000237782\\
1.622	-0.000234275\\
1.624	-0.000230313\\
1.626	-0.000225912\\
1.628	-0.00022109\\
1.63	-0.000215863\\
1.632	-0.00021025\\
1.634	-0.000204271\\
1.636	-0.000197943\\
1.638	-0.000191287\\
1.64	-0.000184323\\
1.642	-0.000177071\\
1.644	-0.000169552\\
1.646	-0.000161787\\
1.648	-0.000153793\\
1.65	-0.000145596\\
1.652	-0.000137217\\
1.654	-0.000128676\\
1.656	-0.000119997\\
1.658	-0.000111199\\
1.66	-0.000102304\\
1.662	-9.33336e-05\\
1.664	-8.43082e-05\\
1.666	-7.52486e-05\\
1.668	-6.61753e-05\\
1.67	-5.71082e-05\\
1.672	-4.80673e-05\\
1.674	-4.27057e-05\\
1.676	-4.23265e-05\\
1.678	-4.18514e-05\\
1.68	-4.12834e-05\\
1.682	-4.06254e-05\\
1.684	-3.98804e-05\\
1.686	-3.90518e-05\\
1.688	-3.8143e-05\\
1.69	-3.71576e-05\\
1.692	-3.677e-05\\
1.694	-4.43581e-05\\
1.696	-5.17227e-05\\
1.698	-5.88513e-05\\
1.7	-6.57322e-05\\
1.702	-7.23546e-05\\
1.704	-7.87083e-05\\
1.706	-8.47838e-05\\
1.708	-9.05725e-05\\
1.71	-9.60665e-05\\
1.712	-0.000101259\\
1.714	-0.000106143\\
1.716	-0.000110713\\
1.718	-0.000114965\\
1.72	-0.000118894\\
1.722	-0.000122498\\
1.724	-0.000125773\\
1.726	-0.000128719\\
1.728	-0.000131334\\
1.73	-0.000133618\\
1.732	-0.000135571\\
1.734	-0.000137195\\
1.736	-0.000138491\\
1.738	-0.000139463\\
1.74	-0.000140114\\
1.742	-0.000140446\\
1.744	-0.000140465\\
1.746	-0.000140175\\
1.748	-0.000139583\\
1.75	-0.000138695\\
1.752	-0.000137519\\
1.754	-0.000136061\\
1.756	-0.000134329\\
1.758	-0.000132332\\
1.76	-0.000130078\\
1.762	-0.000127576\\
1.764	-0.000124837\\
1.766	-0.00012187\\
1.768	-0.000118685\\
1.77	-0.000115293\\
1.772	-0.000111705\\
1.774	-0.000107931\\
1.776	-0.000103983\\
1.778	-9.98732e-05\\
1.78	-9.56123e-05\\
1.782	-9.12124e-05\\
1.784	-8.66854e-05\\
1.786	-8.20433e-05\\
1.788	-7.72982e-05\\
1.79	-7.24621e-05\\
1.792	-6.75471e-05\\
1.794	-6.25651e-05\\
1.796	-5.75282e-05\\
1.798	-5.24483e-05\\
1.8	-4.73371e-05\\
1.802	-4.22064e-05\\
1.804	-3.70675e-05\\
1.806	-3.1932e-05\\
1.808	-2.68109e-05\\
1.81	-2.44282e-05\\
1.812	-2.40438e-05\\
1.814	-2.36078e-05\\
1.816	-2.31226e-05\\
1.818	-2.25901e-05\\
1.82	-2.20124e-05\\
1.822	-2.13914e-05\\
1.824	-2.07295e-05\\
1.826	-2.00286e-05\\
1.828	-2.13141e-05\\
1.83	-2.56163e-05\\
1.832	-2.97922e-05\\
1.834	-3.38349e-05\\
1.836	-3.77376e-05\\
1.838	-4.14943e-05\\
1.84	-4.50989e-05\\
1.842	-4.85463e-05\\
1.844	-5.18315e-05\\
1.846	-5.49499e-05\\
1.848	-5.78975e-05\\
1.85	-6.06706e-05\\
1.852	-6.32661e-05\\
1.854	-6.56812e-05\\
1.856	-6.79136e-05\\
1.858	-6.99614e-05\\
1.86	-7.18231e-05\\
1.862	-7.34978e-05\\
1.864	-7.49848e-05\\
1.866	-7.6284e-05\\
1.868	-7.73955e-05\\
1.87	-7.83201e-05\\
1.872	-7.90587e-05\\
1.874	-7.96128e-05\\
1.876	-7.99841e-05\\
1.878	-8.01749e-05\\
1.88	-8.01876e-05\\
1.882	-8.00251e-05\\
1.884	-7.96907e-05\\
1.886	-7.91878e-05\\
1.888	-7.85199e-05\\
1.89	-7.76914e-05\\
1.892	-7.67069e-05\\
1.894	-7.5571e-05\\
1.896	-7.42886e-05\\
1.898	-7.28648e-05\\
1.9	-7.13052e-05\\
1.902	-6.96153e-05\\
1.904	-6.78008e-05\\
1.906	-6.58678e-05\\
1.908	-6.38223e-05\\
1.91	-6.16707e-05\\
1.912	-5.94193e-05\\
1.914	-5.70747e-05\\
1.916	-5.46435e-05\\
1.918	-5.21323e-05\\
1.92	-4.95479e-05\\
1.922	-4.68972e-05\\
1.924	-4.41869e-05\\
1.926	-4.1424e-05\\
1.928	-3.86152e-05\\
1.93	-3.57675e-05\\
1.932	-3.28877e-05\\
1.934	-2.99826e-05\\
1.936	-2.70588e-05\\
1.938	-2.41231e-05\\
1.94	-2.1182e-05\\
1.942	-1.82421e-05\\
1.944	-1.59205e-05\\
1.946	-1.64317e-05\\
1.948	-1.68855e-05\\
1.95	-1.72809e-05\\
1.952	-1.76174e-05\\
1.954	-1.78945e-05\\
1.956	-1.81119e-05\\
1.958	-1.82694e-05\\
1.96	-1.83673e-05\\
1.962	-1.84059e-05\\
1.964	-1.83857e-05\\
1.966	-1.83075e-05\\
1.968	-1.81722e-05\\
1.97	-1.94856e-05\\
1.972	-2.1729e-05\\
1.974	-2.3889e-05\\
1.976	-2.59623e-05\\
1.978	-2.79458e-05\\
1.98	-2.98365e-05\\
1.982	-3.16318e-05\\
1.984	-3.33293e-05\\
1.986	-3.49271e-05\\
1.988	-3.64231e-05\\
1.99	-3.78158e-05\\
1.992	-3.91038e-05\\
1.994	-4.02859e-05\\
1.996	-4.13613e-05\\
1.998	-4.23292e-05\\
2	-4.31894e-05\\
};
\addplot [color=mycolor7, line width=1.0pt, forget plot]
  table[row sep=crcr]{%
-0.024	0\\
-0.022	0\\
-0.02	0\\
-0.018	0\\
-0.016	0\\
-0.014	0\\
-0.012	0\\
-0.01	0\\
-0.008	0\\
-0.006	0\\
-0.004	0\\
-0.002	0\\
0	0\\
0.002	4.79017e-10\\
0.004	-1.09352e-09\\
0.006	-1.9581e-08\\
0.008	-8.27391e-08\\
0.01	-2.30598e-07\\
0.012	-5.15318e-07\\
0.014	-1.00136e-06\\
0.016	-1.76585e-06\\
0.018	-2.89896e-06\\
0.02	-4.50429e-06\\
0.022	-6.69908e-06\\
0.024	-9.61426e-06\\
0.026	-1.33942e-05\\
0.028	-1.81963e-05\\
0.03	-2.41901e-05\\
0.032	-3.15561e-05\\
0.034	-4.04849e-05\\
0.036	-5.11748e-05\\
0.038	-6.38309e-05\\
0.04	-1.14389e-05\\
0.042	-1.62109e-05\\
0.044	-2.23239e-05\\
0.046	-3.00306e-05\\
0.048	-3.96103e-05\\
0.05	-5.13693e-05\\
0.052	-6.56398e-05\\
0.054	-8.27792e-05\\
0.056	-0.000103169\\
0.058	-0.000127212\\
0.06	-0.000155332\\
0.062	-0.000187969\\
0.064	-0.000225575\\
0.066	-0.000268616\\
0.068	-0.00031756\\
0.07	-0.00037288\\
0.072	-0.000435044\\
0.074	-0.000504514\\
0.076	-0.000581738\\
0.078	-0.00688698\\
0.08	-0.0159624\\
0.082	-0.0210855\\
0.084	-0.0217464\\
0.086	-0.0223714\\
0.088	-0.0229579\\
0.09	-0.0235034\\
0.092	-0.0240058\\
0.094	-0.0244629\\
0.096	-0.0248728\\
0.098	-0.025234\\
0.1	-0.0255448\\
0.102	-0.0258039\\
0.104	-0.0260101\\
0.106	-0.0261622\\
0.108	-0.0262594\\
0.11	-0.026301\\
0.112	-0.0253079\\
0.114	-0.0237178\\
0.116	-0.0221365\\
0.118	-0.020562\\
0.12	-0.0206249\\
0.122	-0.0211888\\
0.124	-0.0217236\\
0.126	-0.0222273\\
0.128	-0.0226982\\
0.13	-0.0231344\\
0.132	-0.0235345\\
0.134	-0.0238972\\
0.136	-0.0242213\\
0.138	-0.0245057\\
0.14	-0.0247495\\
0.142	-0.024952\\
0.144	-0.0251127\\
0.146	-0.0252312\\
0.148	-0.0253071\\
0.15	-0.0253404\\
0.152	-0.025331\\
0.154	-0.0252792\\
0.156	-0.0251851\\
0.158	-0.0250491\\
0.16	-0.0248718\\
0.162	-0.0246537\\
0.164	-0.0243955\\
0.166	-0.024098\\
0.168	-0.0237621\\
0.17	-0.0233885\\
0.172	-0.0229785\\
0.174	-0.0225329\\
0.176	-0.022053\\
0.178	-0.0215398\\
0.18	-0.0209947\\
0.182	-0.0204187\\
0.184	-0.0198132\\
0.186	-0.0191796\\
0.188	-0.018519\\
0.19	-0.0178328\\
0.192	-0.0171225\\
0.194	-0.0163893\\
0.196	-0.0156346\\
0.198	-0.0148599\\
0.2	-0.0140664\\
0.202	-0.0132556\\
0.204	-0.0124288\\
0.206	-0.0115875\\
0.208	-0.0107328\\
0.21	-0.00986631\\
0.212	-0.00898924\\
0.214	-0.00810294\\
0.216	-0.00720871\\
0.218	-0.00630786\\
0.22	-0.00540166\\
0.222	-0.00449137\\
0.224	-0.00357823\\
0.226	-0.00266346\\
0.228	-0.00174827\\
0.23	-0.000833814\\
0.232	7.87364e-05\\
0.234	0.000988255\\
0.236	0.00118117\\
0.238	0.00110983\\
0.24	0.00107263\\
0.242	0.00106785\\
0.244	0.0010936\\
0.246	0.00114785\\
0.248	0.000190708\\
0.25	-0.00107333\\
0.252	-0.00220802\\
0.254	-0.0032159\\
0.256	-0.00409969\\
0.258	-0.0048623\\
0.26	-0.00550682\\
0.262	-0.00603653\\
0.264	-0.00645485\\
0.266	-0.00676534\\
0.268	-0.00697171\\
0.27	-0.0070778\\
0.272	-0.00692693\\
0.274	-0.00668018\\
0.276	-0.00634925\\
0.278	-0.00593726\\
0.28	-0.00534058\\
0.282	-0.00445689\\
0.284	-0.00351821\\
0.286	-0.00252878\\
0.288	-0.0014929\\
0.29	-0.00119941\\
0.292	-0.00114658\\
0.294	-0.00109301\\
0.296	-0.0014178\\
0.298	-0.00257827\\
0.3	-0.00369015\\
0.302	-0.00475355\\
0.304	-0.00576862\\
0.306	-0.00673553\\
0.308	-0.00765448\\
0.31	-0.00852568\\
0.312	-0.0093494\\
0.314	-0.0101259\\
0.316	-0.0108555\\
0.318	-0.0115386\\
0.32	-0.0121754\\
0.322	-0.0127664\\
0.324	-0.0133121\\
0.326	-0.0138128\\
0.328	-0.0142692\\
0.33	-0.0146816\\
0.332	-0.0150506\\
0.334	-0.0153769\\
0.336	-0.015661\\
0.338	-0.0159035\\
0.34	-0.0158194\\
0.342	-0.0156799\\
0.344	-0.0155161\\
0.346	-0.0153289\\
0.348	-0.0151191\\
0.35	-0.0148878\\
0.352	-0.0146356\\
0.354	-0.0149614\\
0.356	-0.0154534\\
0.358	-0.0158985\\
0.36	-0.0162961\\
0.362	-0.0166457\\
0.364	-0.016947\\
0.366	-0.0171999\\
0.368	-0.0174041\\
0.37	-0.0175598\\
0.372	-0.0176671\\
0.374	-0.0177262\\
0.376	-0.0177375\\
0.378	-0.0177015\\
0.38	-0.0176186\\
0.382	-0.0174897\\
0.384	-0.0173154\\
0.386	-0.0170968\\
0.388	-0.0168346\\
0.39	-0.0165301\\
0.392	-0.0161843\\
0.394	-0.0157985\\
0.396	-0.015374\\
0.398	-0.0149123\\
0.4	-0.0144147\\
0.402	-0.0138828\\
0.404	-0.0133182\\
0.406	-0.0127226\\
0.408	-0.0120977\\
0.41	-0.0114453\\
0.412	-0.0107671\\
0.414	-0.010065\\
0.416	-0.00934092\\
0.418	-0.00859677\\
0.42	-0.00783448\\
0.422	-0.00705603\\
0.424	-0.0062634\\
0.426	-0.00545859\\
0.428	-0.00464359\\
0.43	-0.00382043\\
0.432	-0.0029911\\
0.434	-0.0021576\\
0.436	-0.00132191\\
0.438	-0.000486023\\
0.44	0.000348126\\
0.442	0.000793326\\
0.444	0.000509685\\
0.446	-9.62983e-05\\
0.448	-0.000807591\\
0.45	-0.00147614\\
0.452	-0.00212865\\
0.454	-0.00276446\\
0.456	-0.00338299\\
0.458	-0.00398365\\
0.46	-0.00456589\\
0.462	-0.00512922\\
0.464	-0.00567313\\
0.466	-0.00619719\\
0.468	-0.00704163\\
0.47	-0.00786311\\
0.472	-0.00865095\\
0.474	-0.00940398\\
0.476	-0.0101211\\
0.478	-0.0107538\\
0.48	-0.0112914\\
0.482	-0.0117946\\
0.484	-0.0122626\\
0.486	-0.0126948\\
0.488	-0.0130908\\
0.49	-0.0134501\\
0.492	-0.0137723\\
0.494	-0.0140571\\
0.496	-0.0143045\\
0.498	-0.0145144\\
0.5	-0.0146866\\
0.502	-0.0148213\\
0.504	-0.0149187\\
0.506	-0.014979\\
0.508	-0.0150025\\
0.51	-0.0149897\\
0.512	-0.0149409\\
0.514	-0.0148568\\
0.516	-0.0147379\\
0.518	-0.014585\\
0.52	-0.0143987\\
0.522	-0.01418\\
0.524	-0.0139297\\
0.526	-0.0136487\\
0.528	-0.013338\\
0.53	-0.0129987\\
0.532	-0.0126319\\
0.534	-0.0122387\\
0.536	-0.0118202\\
0.538	-0.0113778\\
0.54	-0.0109127\\
0.542	-0.0104262\\
0.544	-0.00991963\\
0.546	-0.00939432\\
0.548	-0.00885167\\
0.55	-0.00829308\\
0.552	-0.00771999\\
0.554	-0.00713383\\
0.556	-0.0068207\\
0.558	-0.00650764\\
0.56	-0.00618801\\
0.562	-0.00586239\\
0.564	-0.00553137\\
0.566	-0.00519555\\
0.568	-0.00485551\\
0.57	-0.00451185\\
0.572	-0.00416516\\
0.574	-0.00381603\\
0.576	-0.00346504\\
0.578	-0.00311278\\
0.58	-0.00275983\\
0.582	-0.00240677\\
0.584	-0.00205416\\
0.586	-0.00170256\\
0.588	-0.00135252\\
0.59	-0.00126532\\
0.592	-0.00193964\\
0.594	-0.00259949\\
0.596	-0.00324368\\
0.598	-0.00387109\\
0.6	-0.00448063\\
0.602	-0.00507128\\
0.604	-0.00564207\\
0.606	-0.00619208\\
0.608	-0.00672045\\
0.61	-0.00722637\\
0.612	-0.00770911\\
0.614	-0.00816798\\
0.616	-0.00860235\\
0.618	-0.00901166\\
0.62	-0.00939541\\
0.622	-0.00969542\\
0.624	-0.00995094\\
0.626	-0.0101825\\
0.628	-0.0103898\\
0.63	-0.0105729\\
0.632	-0.0107315\\
0.634	-0.0108656\\
0.636	-0.0109752\\
0.638	-0.0110604\\
0.64	-0.0111214\\
0.642	-0.0111581\\
0.644	-0.0111709\\
0.646	-0.0111601\\
0.648	-0.0111258\\
0.65	-0.0110685\\
0.652	-0.0109886\\
0.654	-0.0108865\\
0.656	-0.0107627\\
0.658	-0.0106176\\
0.66	-0.010452\\
0.662	-0.0102663\\
0.664	-0.0100611\\
0.666	-0.00983727\\
0.668	-0.00959535\\
0.67	-0.00933611\\
0.672	-0.00906031\\
0.674	-0.00876872\\
0.676	-0.00846215\\
0.678	-0.00814143\\
0.68	-0.00780741\\
0.682	-0.00746094\\
0.684	-0.00710291\\
0.686	-0.00673423\\
0.688	-0.00635578\\
0.69	-0.00596849\\
0.692	-0.00557329\\
0.694	-0.00531956\\
0.696	-0.00521445\\
0.698	-0.00510203\\
0.7	-0.00498259\\
0.702	-0.00485644\\
0.704	-0.00472391\\
0.706	-0.0045853\\
0.708	-0.00444096\\
0.71	-0.00429121\\
0.712	-0.00413639\\
0.714	-0.00397684\\
0.716	-0.0038129\\
0.718	-0.00364493\\
0.72	-0.00347326\\
0.722	-0.00329826\\
0.724	-0.00312027\\
0.726	-0.00293964\\
0.728	-0.00275672\\
0.73	-0.00257186\\
0.732	-0.00238541\\
0.734	-0.00230597\\
0.736	-0.00269752\\
0.738	-0.00307749\\
0.74	-0.00344522\\
0.742	-0.00380014\\
0.744	-0.00414168\\
0.746	-0.00446931\\
0.748	-0.00478255\\
0.75	-0.00507222\\
0.752	-0.00534412\\
0.754	-0.00560098\\
0.756	-0.00584249\\
0.758	-0.00606834\\
0.76	-0.00627828\\
0.762	-0.00647208\\
0.764	-0.00664716\\
0.766	-0.00680157\\
0.768	-0.00693946\\
0.77	-0.00706076\\
0.772	-0.00716547\\
0.774	-0.00725331\\
0.776	-0.007324\\
0.778	-0.00737099\\
0.78	-0.00739602\\
0.782	-0.00740556\\
0.784	-0.00739981\\
0.786	-0.00737896\\
0.788	-0.00734328\\
0.79	-0.00729303\\
0.792	-0.00722852\\
0.794	-0.00715007\\
0.796	-0.00705803\\
0.798	-0.00695279\\
0.8	-0.00683474\\
0.802	-0.0067043\\
0.804	-0.00656192\\
0.806	-0.00640805\\
0.808	-0.00624319\\
0.81	-0.00606781\\
0.812	-0.00588244\\
0.814	-0.0056876\\
0.816	-0.00548383\\
0.818	-0.00527169\\
0.82	-0.00505172\\
0.822	-0.00482451\\
0.824	-0.00459062\\
0.826	-0.00435065\\
0.828	-0.00410517\\
0.83	-0.00385479\\
0.832	-0.00360009\\
0.834	-0.00334166\\
0.836	-0.00315758\\
0.838	-0.00312466\\
0.84	-0.00308678\\
0.842	-0.00304409\\
0.844	-0.00299673\\
0.846	-0.00294486\\
0.848	-0.00288865\\
0.85	-0.00282825\\
0.852	-0.002759\\
0.854	-0.00268495\\
0.856	-0.00260731\\
0.858	-0.00252629\\
0.86	-0.00244209\\
0.862	-0.0023549\\
0.864	-0.00226494\\
0.866	-0.00217242\\
0.868	-0.00207753\\
0.87	-0.0019805\\
0.872	-0.00188153\\
0.874	-0.00178083\\
0.876	-0.00175252\\
0.878	-0.00196373\\
0.88	-0.00216746\\
0.882	-0.00236339\\
0.884	-0.00255122\\
0.886	-0.00273068\\
0.888	-0.00290152\\
0.89	-0.0030635\\
0.892	-0.00321643\\
0.894	-0.00336012\\
0.896	-0.0034944\\
0.898	-0.00361915\\
0.9	-0.00373424\\
0.902	-0.00383958\\
0.904	-0.0039351\\
0.906	-0.00402074\\
0.908	-0.0040965\\
0.91	-0.00416234\\
0.912	-0.00421831\\
0.914	-0.00426442\\
0.916	-0.00430073\\
0.918	-0.00432733\\
0.92	-0.00434431\\
0.922	-0.00435179\\
0.924	-0.00434989\\
0.926	-0.00433878\\
0.928	-0.00431861\\
0.93	-0.00428499\\
0.932	-0.00424197\\
0.934	-0.00419071\\
0.936	-0.00413146\\
0.938	-0.00406446\\
0.94	-0.00398996\\
0.942	-0.00390826\\
0.944	-0.00381961\\
0.946	-0.00372433\\
0.948	-0.00362271\\
0.95	-0.00351506\\
0.952	-0.0034017\\
0.954	-0.00328297\\
0.956	-0.00315918\\
0.958	-0.0030307\\
0.96	-0.00289784\\
0.962	-0.00276097\\
0.964	-0.00262043\\
0.966	-0.00247658\\
0.968	-0.00232977\\
0.97	-0.00218036\\
0.972	-0.0020287\\
0.974	-0.00187514\\
0.976	-0.00172004\\
0.978	-0.00164969\\
0.98	-0.00164258\\
0.982	-0.00163249\\
0.984	-0.00161949\\
0.986	-0.00160365\\
0.988	-0.00158506\\
0.99	-0.00156379\\
0.992	-0.00153994\\
0.994	-0.0015136\\
0.996	-0.00148486\\
0.998	-0.00145382\\
1	-0.00142058\\
1.002	-0.00138525\\
1.004	-0.00134793\\
1.006	-0.00130872\\
1.008	-0.00126774\\
1.01	-0.0012251\\
1.012	-0.00118091\\
1.014	-0.00113529\\
1.016	-0.00113771\\
1.018	-0.00125706\\
1.02	-0.00137178\\
1.022	-0.00148173\\
1.024	-0.00158531\\
1.026	-0.00168331\\
1.028	-0.00177616\\
1.03	-0.00186377\\
1.032	-0.00194602\\
1.034	-0.00202283\\
1.036	-0.00209411\\
1.038	-0.00215981\\
1.04	-0.00221988\\
1.042	-0.00227427\\
1.044	-0.00232296\\
1.046	-0.00236594\\
1.048	-0.00240321\\
1.05	-0.00243478\\
1.052	-0.00246068\\
1.054	-0.00248093\\
1.056	-0.00249559\\
1.058	-0.00250472\\
1.06	-0.00250838\\
1.062	-0.00250665\\
1.064	-0.00249963\\
1.066	-0.00248741\\
1.068	-0.0024701\\
1.07	-0.00244783\\
1.072	-0.00242071\\
1.074	-0.00238888\\
1.076	-0.0023525\\
1.078	-0.0023117\\
1.08	-0.00226665\\
1.082	-0.00221716\\
1.084	-0.00216282\\
1.086	-0.00210479\\
1.088	-0.00204325\\
1.09	-0.00197839\\
1.092	-0.0019104\\
1.094	-0.00183947\\
1.096	-0.00176579\\
1.098	-0.00168958\\
1.1	-0.00161103\\
1.102	-0.00153034\\
1.104	-0.00144772\\
1.106	-0.00136338\\
1.108	-0.00127753\\
1.11	-0.00119036\\
1.112	-0.00110209\\
1.114	-0.00101293\\
1.116	-0.000923067\\
1.118	-0.000832714\\
1.12	-0.000807159\\
1.122	-0.000806202\\
1.124	-0.000803675\\
1.126	-0.000799612\\
1.128	-0.000794052\\
1.13	-0.000787036\\
1.132	-0.000778606\\
1.134	-0.000768807\\
1.136	-0.000757686\\
1.138	-0.000745291\\
1.14	-0.000731672\\
1.142	-0.00071688\\
1.144	-0.000700969\\
1.146	-0.000683991\\
1.148	-0.000666003\\
1.15	-0.00064706\\
1.152	-0.000627219\\
1.154	-0.000649655\\
1.156	-0.000716859\\
1.158	-0.000781413\\
1.16	-0.000843222\\
1.162	-0.000901994\\
1.164	-0.00095745\\
1.166	-0.00100994\\
1.168	-0.0010594\\
1.17	-0.00110578\\
1.172	-0.00114901\\
1.174	-0.00118908\\
1.176	-0.00122593\\
1.178	-0.00125954\\
1.18	-0.00128989\\
1.182	-0.00131698\\
1.184	-0.00134079\\
1.186	-0.00136133\\
1.188	-0.00137861\\
1.19	-0.00139264\\
1.192	-0.00140345\\
1.194	-0.00141107\\
1.196	-0.00141553\\
1.198	-0.00141688\\
1.2	-0.00141515\\
1.202	-0.00141042\\
1.204	-0.00140274\\
1.206	-0.00139216\\
1.208	-0.00137877\\
1.21	-0.00136264\\
1.212	-0.00134384\\
1.214	-0.00132247\\
1.216	-0.00129861\\
1.218	-0.00127235\\
1.22	-0.00124379\\
1.222	-0.00121303\\
1.224	-0.00118017\\
1.226	-0.00114531\\
1.228	-0.00110857\\
1.23	-0.00107006\\
1.232	-0.00102988\\
1.234	-0.000988153\\
1.236	-0.000944994\\
1.238	-0.000900355\\
1.24	-0.000854431\\
1.242	-0.000807439\\
1.244	-0.000759498\\
1.246	-0.000710725\\
1.248	-0.00066124\\
1.25	-0.000611162\\
1.252	-0.000560607\\
1.254	-0.000509692\\
1.256	-0.000458534\\
1.258	-0.000407246\\
1.26	-0.00037321\\
1.262	-0.000373916\\
1.264	-0.00037383\\
1.266	-0.000372969\\
1.268	-0.00037135\\
1.27	-0.000368993\\
1.272	-0.000365916\\
1.274	-0.000362141\\
1.276	-0.000357691\\
1.278	-0.000352588\\
1.28	-0.000346856\\
1.282	-0.000340521\\
1.284	-0.000333607\\
1.286	-0.000326141\\
1.288	-0.000318151\\
1.29	-0.000340159\\
1.292	-0.00037871\\
1.294	-0.00041579\\
1.296	-0.000451343\\
1.298	-0.000485318\\
1.3	-0.000517667\\
1.302	-0.000548349\\
1.304	-0.000577324\\
1.306	-0.000604559\\
1.308	-0.000630023\\
1.31	-0.000653691\\
1.312	-0.000675542\\
1.314	-0.00069556\\
1.316	-0.000713731\\
1.318	-0.000730048\\
1.32	-0.000744463\\
1.322	-0.000756905\\
1.324	-0.000767499\\
1.326	-0.000776252\\
1.328	-0.000783177\\
1.33	-0.000788288\\
1.332	-0.000791606\\
1.334	-0.000793152\\
1.336	-0.000792954\\
1.338	-0.00079104\\
1.34	-0.000787445\\
1.342	-0.000782203\\
1.344	-0.000775353\\
1.346	-0.000766938\\
1.348	-0.000757002\\
1.35	-0.000745592\\
1.352	-0.000732758\\
1.354	-0.00071855\\
1.356	-0.000703024\\
1.358	-0.000686234\\
1.36	-0.000668238\\
1.362	-0.000649096\\
1.364	-0.000628869\\
1.366	-0.000607617\\
1.368	-0.000585406\\
1.37	-0.000562299\\
1.372	-0.00053836\\
1.374	-0.000513657\\
1.376	-0.000488256\\
1.378	-0.000462223\\
1.38	-0.000435626\\
1.382	-0.000408531\\
1.384	-0.000381006\\
1.386	-0.000353119\\
1.388	-0.000324935\\
1.39	-0.00029652\\
1.392	-0.000267942\\
1.394	-0.000239247\\
1.396	-0.000210492\\
1.398	-0.000181766\\
1.4	-0.000172208\\
1.402	-0.000172349\\
1.404	-0.000172105\\
1.406	-0.000171483\\
1.408	-0.000170494\\
1.41	-0.000169146\\
1.412	-0.00016745\\
1.414	-0.000165418\\
1.416	-0.000163061\\
1.418	-0.00016039\\
1.42	-0.000157419\\
1.422	-0.000154159\\
1.424	-0.000161417\\
1.426	-0.000183979\\
1.428	-0.000205757\\
1.43	-0.000226719\\
1.432	-0.000246831\\
1.434	-0.000266065\\
1.436	-0.000284393\\
1.438	-0.00030179\\
1.44	-0.000318235\\
1.442	-0.000333708\\
1.444	-0.000348191\\
1.446	-0.000361669\\
1.448	-0.00037413\\
1.45	-0.000385563\\
1.452	-0.000395962\\
1.454	-0.000405321\\
1.456	-0.00041363\\
1.458	-0.000420878\\
1.46	-0.000427086\\
1.462	-0.000432256\\
1.464	-0.000436395\\
1.466	-0.000439512\\
1.468	-0.000441616\\
1.47	-0.000442721\\
1.472	-0.000442824\\
1.474	-0.000441957\\
1.476	-0.00044014\\
1.478	-0.000437392\\
1.48	-0.000433736\\
1.482	-0.000429194\\
1.484	-0.000423791\\
1.486	-0.000417553\\
1.488	-0.000410509\\
1.49	-0.000402686\\
1.492	-0.000394115\\
1.494	-0.000384828\\
1.496	-0.000374856\\
1.498	-0.000364232\\
1.5	-0.000352992\\
1.502	-0.000341169\\
1.504	-0.0003288\\
1.506	-0.00031592\\
1.508	-0.000302567\\
1.51	-0.000288777\\
1.512	-0.000274587\\
1.514	-0.000260036\\
1.516	-0.000245161\\
1.518	-0.000230001\\
1.52	-0.000214592\\
1.522	-0.000198974\\
1.524	-0.000183183\\
1.526	-0.000167257\\
1.528	-0.000151233\\
1.53	-0.000135148\\
1.532	-0.000119038\\
1.534	-0.000102938\\
1.536	-8.68839e-05\\
1.538	-8.19766e-05\\
1.54	-8.17275e-05\\
1.542	-8.12905e-05\\
1.544	-8.06706e-05\\
1.546	-7.98726e-05\\
1.548	-7.89022e-05\\
1.55	-7.77648e-05\\
1.552	-7.64665e-05\\
1.554	-7.50134e-05\\
1.556	-7.3412e-05\\
1.558	-7.71056e-05\\
1.56	-9.02069e-05\\
1.562	-0.00010289\\
1.564	-0.000115134\\
1.566	-0.000126919\\
1.568	-0.000138228\\
1.57	-0.000149044\\
1.572	-0.000159351\\
1.574	-0.000169134\\
1.576	-0.000178381\\
1.578	-0.000187081\\
1.58	-0.000195226\\
1.582	-0.000202805\\
1.584	-0.000209812\\
1.586	-0.000216242\\
1.588	-0.00022209\\
1.59	-0.000227345\\
1.592	-0.000232013\\
1.594	-0.000236094\\
1.596	-0.00023959\\
1.598	-0.000242501\\
1.6	-0.000244832\\
1.602	-0.000246587\\
1.604	-0.000247771\\
1.606	-0.000248393\\
1.608	-0.00024846\\
1.61	-0.00024798\\
1.612	-0.000246965\\
1.614	-0.000245426\\
1.616	-0.000243374\\
1.618	-0.000240822\\
1.62	-0.000237786\\
1.622	-0.000234279\\
1.624	-0.000230316\\
1.626	-0.000225915\\
1.628	-0.000221093\\
1.63	-0.000215866\\
1.632	-0.000210253\\
1.634	-0.000204273\\
1.636	-0.000197946\\
1.638	-0.00019129\\
1.64	-0.000184325\\
1.642	-0.000177074\\
1.644	-0.000169555\\
1.646	-0.000161788\\
1.648	-0.000153794\\
1.65	-0.000145597\\
1.652	-0.000137217\\
1.654	-0.000128677\\
1.656	-0.000119997\\
1.658	-0.000111199\\
1.66	-0.000102305\\
1.662	-9.3334e-05\\
1.664	-8.43086e-05\\
1.666	-7.52491e-05\\
1.668	-6.61757e-05\\
1.67	-5.71087e-05\\
1.672	-4.80677e-05\\
1.674	-4.27061e-05\\
1.676	-4.23269e-05\\
1.678	-4.18519e-05\\
1.68	-4.12839e-05\\
1.682	-4.06258e-05\\
1.684	-3.98808e-05\\
1.686	-3.90522e-05\\
1.688	-3.81434e-05\\
1.69	-3.7158e-05\\
1.692	-3.67704e-05\\
1.694	-4.43585e-05\\
1.696	-5.17231e-05\\
1.698	-5.88517e-05\\
1.7	-6.57326e-05\\
1.702	-7.23551e-05\\
1.704	-7.87087e-05\\
1.706	-8.47842e-05\\
1.708	-9.05729e-05\\
1.71	-9.60669e-05\\
1.712	-0.000101259\\
1.714	-0.000106143\\
1.716	-0.000110713\\
1.718	-0.000114965\\
1.72	-0.000118895\\
1.722	-0.000122498\\
1.724	-0.000125774\\
1.726	-0.000128719\\
1.728	-0.000131334\\
1.73	-0.000133618\\
1.732	-0.000135571\\
1.734	-0.000137195\\
1.736	-0.000138492\\
1.738	-0.000139464\\
1.74	-0.000140114\\
1.742	-0.000140447\\
1.744	-0.000140466\\
1.746	-0.000140176\\
1.748	-0.000139584\\
1.75	-0.000138696\\
1.752	-0.00013752\\
1.754	-0.000136062\\
1.756	-0.00013433\\
1.758	-0.000132332\\
1.76	-0.000130078\\
1.762	-0.000127577\\
1.764	-0.000124838\\
1.766	-0.000121871\\
1.768	-0.000118686\\
1.77	-0.000115294\\
1.772	-0.000111705\\
1.774	-0.000107932\\
1.776	-0.000103984\\
1.778	-9.98744e-05\\
1.78	-9.56137e-05\\
1.782	-9.12139e-05\\
1.784	-8.66871e-05\\
1.786	-8.20452e-05\\
1.788	-7.73002e-05\\
1.79	-7.24641e-05\\
1.792	-6.7549e-05\\
1.794	-6.25669e-05\\
1.796	-5.75298e-05\\
1.798	-5.24497e-05\\
1.8	-4.73384e-05\\
1.802	-4.22075e-05\\
1.804	-3.70685e-05\\
1.806	-3.19328e-05\\
1.808	-2.68116e-05\\
1.81	-2.44287e-05\\
1.812	-2.40437e-05\\
1.814	-2.36077e-05\\
1.816	-2.31224e-05\\
1.818	-2.259e-05\\
1.82	-2.20123e-05\\
1.822	-2.13914e-05\\
1.824	-2.07294e-05\\
1.826	-2.00285e-05\\
1.828	-2.13141e-05\\
1.83	-2.56163e-05\\
1.832	-2.97922e-05\\
1.834	-3.38349e-05\\
1.836	-3.77376e-05\\
1.838	-4.14943e-05\\
1.84	-4.5099e-05\\
1.842	-4.85463e-05\\
1.844	-5.18315e-05\\
1.846	-5.49499e-05\\
1.848	-5.78975e-05\\
1.85	-6.06706e-05\\
1.852	-6.32661e-05\\
1.854	-6.56813e-05\\
1.856	-6.79136e-05\\
1.858	-6.99614e-05\\
1.86	-7.18232e-05\\
1.862	-7.34979e-05\\
1.864	-7.49849e-05\\
1.866	-7.62841e-05\\
1.868	-7.73956e-05\\
1.87	-7.83202e-05\\
1.872	-7.90588e-05\\
1.874	-7.96129e-05\\
1.876	-7.99842e-05\\
1.878	-8.0175e-05\\
1.88	-8.01877e-05\\
1.882	-8.00253e-05\\
1.884	-7.96908e-05\\
1.886	-7.91879e-05\\
1.888	-7.85204e-05\\
1.89	-7.76923e-05\\
1.892	-7.67082e-05\\
1.894	-7.55725e-05\\
1.896	-7.429e-05\\
1.898	-7.28663e-05\\
1.9	-7.13066e-05\\
1.902	-6.96167e-05\\
1.904	-6.78021e-05\\
1.906	-6.58691e-05\\
1.908	-6.38236e-05\\
1.91	-6.16719e-05\\
1.912	-5.94205e-05\\
1.914	-5.70758e-05\\
1.916	-5.46445e-05\\
1.918	-5.21333e-05\\
1.92	-4.95489e-05\\
1.922	-4.68981e-05\\
1.924	-4.41878e-05\\
1.926	-4.14248e-05\\
1.928	-3.8616e-05\\
1.93	-3.57682e-05\\
1.932	-3.28883e-05\\
1.934	-2.99831e-05\\
1.936	-2.70593e-05\\
1.938	-2.41236e-05\\
1.94	-2.11824e-05\\
1.942	-1.82425e-05\\
1.944	-1.59208e-05\\
1.946	-1.6432e-05\\
1.948	-1.68857e-05\\
1.95	-1.72811e-05\\
1.952	-1.76176e-05\\
1.954	-1.78946e-05\\
1.956	-1.81119e-05\\
1.958	-1.82694e-05\\
1.96	-1.83672e-05\\
1.962	-1.84058e-05\\
1.964	-1.83857e-05\\
1.966	-1.83075e-05\\
1.968	-1.81722e-05\\
1.97	-1.94856e-05\\
1.972	-2.17289e-05\\
1.974	-2.38889e-05\\
1.976	-2.59622e-05\\
1.978	-2.79457e-05\\
1.98	-2.98364e-05\\
1.982	-3.16317e-05\\
1.984	-3.33293e-05\\
1.986	-3.4927e-05\\
1.988	-3.64231e-05\\
1.99	-3.78158e-05\\
1.992	-3.91037e-05\\
1.994	-4.02859e-05\\
1.996	-4.13613e-05\\
1.998	-4.23293e-05\\
2	-4.31896e-05\\
};
\addplot [color=mycolor8, line width=1.0pt, forget plot]
  table[row sep=crcr]{%
-0.024	0\\
-0.022	0\\
-0.02	0\\
-0.018	0\\
-0.016	0\\
-0.014	0\\
-0.012	0\\
-0.01	0\\
-0.008	0\\
-0.006	0\\
-0.004	0\\
-0.002	0\\
0	0\\
0.002	8.18602e-10\\
0.004	1.46434e-09\\
0.006	-1.14092e-08\\
0.008	-6.43129e-08\\
0.01	-1.96215e-07\\
0.012	-4.58348e-07\\
0.014	-9.14366e-07\\
0.016	-1.64072e-06\\
0.018	-2.72708e-06\\
0.02	-4.27678e-06\\
0.022	-6.40706e-06\\
0.024	-9.24925e-06\\
0.026	-1.29486e-05\\
0.028	-1.76637e-05\\
0.03	-2.35663e-05\\
0.032	-3.08396e-05\\
0.034	-3.96774e-05\\
0.036	-5.02828e-05\\
0.038	-6.28659e-05\\
0.04	-1.04498e-05\\
0.042	-1.52076e-05\\
0.044	-2.13448e-05\\
0.046	-2.91244e-05\\
0.048	-3.88377e-05\\
0.05	-5.08037e-05\\
0.052	-6.53688e-05\\
0.054	-8.2906e-05\\
0.056	-0.000103813\\
0.058	-0.000128511\\
0.06	-0.000157441\\
0.062	-0.000191064\\
0.064	-0.000229852\\
0.066	-0.000274292\\
0.068	-0.000324875\\
0.07	-0.000382094\\
0.072	-0.000446442\\
0.074	-0.000518403\\
0.076	-0.000598448\\
0.078	-0.00690686\\
0.08	-0.0159858\\
0.082	-0.0217071\\
0.084	-0.0223998\\
0.086	-0.0230565\\
0.088	-0.0236744\\
0.09	-0.0242511\\
0.092	-0.0247841\\
0.094	-0.0252713\\
0.096	-0.0257107\\
0.098	-0.0261004\\
0.1	-0.026439\\
0.102	-0.0267247\\
0.104	-0.0269565\\
0.106	-0.027133\\
0.108	-0.0272532\\
0.11	-0.0273164\\
0.112	-0.0261751\\
0.114	-0.0246037\\
0.116	-0.0230399\\
0.118	-0.0214819\\
0.12	-0.02156\\
0.122	-0.0221378\\
0.124	-0.0226853\\
0.126	-0.0232003\\
0.128	-0.0236811\\
0.13	-0.0241259\\
0.132	-0.0245333\\
0.134	-0.0249018\\
0.136	-0.0252303\\
0.138	-0.0255176\\
0.14	-0.025763\\
0.142	-0.0259657\\
0.144	-0.0261251\\
0.146	-0.0262409\\
0.148	-0.0263128\\
0.15	-0.0263407\\
0.152	-0.0263246\\
0.154	-0.0262648\\
0.156	-0.0261615\\
0.158	-0.026015\\
0.16	-0.0258261\\
0.162	-0.0255952\\
0.164	-0.0253232\\
0.166	-0.0250108\\
0.168	-0.0246589\\
0.17	-0.0242686\\
0.172	-0.0238408\\
0.174	-0.0233767\\
0.176	-0.0228775\\
0.178	-0.0223443\\
0.18	-0.0217785\\
0.182	-0.0211813\\
0.184	-0.0205541\\
0.186	-0.0198982\\
0.188	-0.0192149\\
0.19	-0.0185058\\
0.192	-0.0177721\\
0.194	-0.0170154\\
0.196	-0.0162371\\
0.198	-0.0154385\\
0.2	-0.0146212\\
0.202	-0.0137865\\
0.204	-0.0129359\\
0.206	-0.0120707\\
0.208	-0.0111925\\
0.21	-0.0103026\\
0.212	-0.00940232\\
0.214	-0.00849311\\
0.216	-0.00757629\\
0.218	-0.00665319\\
0.22	-0.00572513\\
0.222	-0.0047934\\
0.224	-0.00385927\\
0.226	-0.00292399\\
0.228	-0.00198877\\
0.23	-0.00105482\\
0.232	-0.000123308\\
0.234	0.000804613\\
0.236	0.00101536\\
0.238	0.000961264\\
0.24	0.000940717\\
0.242	0.000951991\\
0.244	0.000993197\\
0.246	0.00106229\\
0.248	0.000119385\\
0.25	-0.00113102\\
0.252	-0.00225268\\
0.254	-0.00324813\\
0.256	-0.00412007\\
0.258	-0.00487141\\
0.26	-0.00550524\\
0.262	-0.00602482\\
0.264	-0.00643354\\
0.266	-0.00673497\\
0.268	-0.00693279\\
0.27	-0.00703083\\
0.272	-0.00674586\\
0.274	-0.00606224\\
0.276	-0.00531046\\
0.278	-0.00449448\\
0.28	-0.00361833\\
0.282	-0.0026861\\
0.284	-0.00170193\\
0.286	-0.000670009\\
0.288	0.00040545\\
0.29	0.000735666\\
0.292	0.000822448\\
0.294	0.000907289\\
0.296	0.000611156\\
0.298	-0.000523188\\
0.3	-0.00161138\\
0.302	-0.00265346\\
0.304	-0.00364948\\
0.306	-0.00459952\\
0.308	-0.00550371\\
0.31	-0.00636217\\
0.312	-0.00717507\\
0.314	-0.00794261\\
0.316	-0.008665\\
0.318	-0.00934251\\
0.32	-0.00997541\\
0.322	-0.010564\\
0.324	-0.0111087\\
0.326	-0.0116098\\
0.328	-0.0120677\\
0.33	-0.0124829\\
0.332	-0.0128559\\
0.334	-0.0131872\\
0.336	-0.0134773\\
0.338	-0.0137267\\
0.34	-0.0139362\\
0.342	-0.0141064\\
0.344	-0.0142378\\
0.346	-0.0143313\\
0.348	-0.0143876\\
0.35	-0.0144073\\
0.352	-0.0143914\\
0.354	-0.0149\\
0.356	-0.015394\\
0.358	-0.0158406\\
0.36	-0.0162392\\
0.362	-0.0165894\\
0.364	-0.016891\\
0.366	-0.0171436\\
0.368	-0.0173472\\
0.37	-0.0175018\\
0.372	-0.0176075\\
0.374	-0.0176647\\
0.376	-0.0176737\\
0.378	-0.017635\\
0.38	-0.0175491\\
0.382	-0.0174168\\
0.384	-0.0172388\\
0.386	-0.0170161\\
0.388	-0.0167495\\
0.39	-0.0164403\\
0.392	-0.0160896\\
0.394	-0.0156986\\
0.396	-0.0152687\\
0.398	-0.0148012\\
0.4	-0.0142977\\
0.402	-0.0137598\\
0.404	-0.0131889\\
0.406	-0.0125869\\
0.408	-0.0119554\\
0.41	-0.0112963\\
0.412	-0.0106113\\
0.414	-0.00990235\\
0.416	-0.00917133\\
0.418	-0.00842017\\
0.42	-0.00765084\\
0.422	-0.00686533\\
0.424	-0.00606564\\
0.426	-0.00525377\\
0.428	-0.00442744\\
0.43	-0.00346726\\
0.432	-0.00249965\\
0.434	-0.00152701\\
0.436	-0.000551721\\
0.438	0.000423858\\
0.44	0.0013974\\
0.442	0.00198134\\
0.444	0.00183543\\
0.446	0.00136584\\
0.448	0.000670836\\
0.45	-8.70361e-06\\
0.452	-0.000672117\\
0.454	-0.00131877\\
0.456	-0.00194807\\
0.458	-0.00255943\\
0.46	-0.00315232\\
0.462	-0.00372623\\
0.464	-0.00428068\\
0.466	-0.00481521\\
0.468	-0.00567009\\
0.47	-0.00650195\\
0.472	-0.00730013\\
0.474	-0.00806346\\
0.476	-0.00879087\\
0.478	-0.00948136\\
0.48	-0.010134\\
0.482	-0.0107481\\
0.484	-0.0113229\\
0.486	-0.0118577\\
0.488	-0.012352\\
0.49	-0.0128054\\
0.492	-0.0132176\\
0.494	-0.0135883\\
0.496	-0.0139174\\
0.498	-0.0141916\\
0.5	-0.0143636\\
0.502	-0.0144982\\
0.504	-0.0145955\\
0.506	-0.0146557\\
0.508	-0.0146793\\
0.51	-0.0146665\\
0.512	-0.0146179\\
0.514	-0.0145339\\
0.516	-0.0144153\\
0.518	-0.0142627\\
0.52	-0.0140768\\
0.522	-0.0138585\\
0.524	-0.0136086\\
0.526	-0.0133281\\
0.528	-0.013018\\
0.53	-0.0126793\\
0.532	-0.0123131\\
0.534	-0.0119206\\
0.536	-0.0115029\\
0.538	-0.0110613\\
0.54	-0.0105971\\
0.542	-0.0101114\\
0.544	-0.00960579\\
0.546	-0.00908146\\
0.548	-0.00853985\\
0.55	-0.00798234\\
0.552	-0.00741038\\
0.554	-0.00682538\\
0.556	-0.00651345\\
0.558	-0.00620165\\
0.56	-0.00588331\\
0.562	-0.00555902\\
0.564	-0.00522938\\
0.566	-0.00489497\\
0.568	-0.00455639\\
0.57	-0.00421422\\
0.572	-0.00386906\\
0.574	-0.00352149\\
0.576	-0.00317211\\
0.578	-0.0028215\\
0.58	-0.00247023\\
0.582	-0.00211887\\
0.584	-0.001768\\
0.586	-0.00141818\\
0.588	-0.00106995\\
0.59	-0.000984589\\
0.592	-0.00166078\\
0.594	-0.00232253\\
0.596	-0.00296865\\
0.598	-0.00359802\\
0.6	-0.00420954\\
0.602	-0.0048022\\
0.604	-0.00537502\\
0.606	-0.00592709\\
0.608	-0.00645754\\
0.61	-0.00693834\\
0.612	-0.00739286\\
0.614	-0.00782399\\
0.616	-0.00823114\\
0.618	-0.0086138\\
0.62	-0.00897151\\
0.622	-0.00930386\\
0.624	-0.00961051\\
0.626	-0.00989118\\
0.628	-0.0101456\\
0.63	-0.0103737\\
0.632	-0.0105585\\
0.634	-0.010695\\
0.636	-0.0108069\\
0.638	-0.0108944\\
0.64	-0.0109575\\
0.642	-0.0109964\\
0.644	-0.0110113\\
0.646	-0.0110024\\
0.648	-0.0109701\\
0.65	-0.0109147\\
0.652	-0.0108366\\
0.654	-0.0107362\\
0.656	-0.0106141\\
0.658	-0.0104707\\
0.66	-0.0103066\\
0.662	-0.0101224\\
0.664	-0.00991873\\
0.666	-0.00969628\\
0.668	-0.00945573\\
0.67	-0.00919781\\
0.672	-0.00892328\\
0.674	-0.00863293\\
0.676	-0.00832755\\
0.678	-0.00800797\\
0.68	-0.00767505\\
0.682	-0.00732966\\
0.684	-0.00697267\\
0.686	-0.00660498\\
0.688	-0.0062275\\
0.69	-0.00584116\\
0.692	-0.00544687\\
0.694	-0.00519402\\
0.696	-0.00508977\\
0.698	-0.00497819\\
0.7	-0.00485957\\
0.702	-0.00473422\\
0.704	-0.00460247\\
0.706	-0.00446464\\
0.708	-0.00432105\\
0.71	-0.00417204\\
0.712	-0.00401794\\
0.714	-0.00385912\\
0.716	-0.00369589\\
0.718	-0.00352863\\
0.72	-0.00335767\\
0.722	-0.00318337\\
0.724	-0.00300607\\
0.726	-0.00282614\\
0.728	-0.00264392\\
0.73	-0.00245976\\
0.732	-0.00227402\\
0.734	-0.00219529\\
0.736	-0.00258755\\
0.738	-0.00296823\\
0.74	-0.0033367\\
0.742	-0.00368922\\
0.744	-0.00401494\\
0.746	-0.00432706\\
0.748	-0.00462511\\
0.75	-0.00490867\\
0.752	-0.00517735\\
0.754	-0.00543082\\
0.756	-0.00566876\\
0.758	-0.0058909\\
0.76	-0.00609699\\
0.762	-0.00628684\\
0.764	-0.0064603\\
0.766	-0.00661723\\
0.768	-0.00675754\\
0.77	-0.0068812\\
0.772	-0.00698817\\
0.774	-0.00707848\\
0.776	-0.00715218\\
0.778	-0.00720937\\
0.78	-0.00725017\\
0.782	-0.00727473\\
0.784	-0.00728324\\
0.786	-0.00727592\\
0.788	-0.00725303\\
0.79	-0.00721484\\
0.792	-0.00716167\\
0.794	-0.00709384\\
0.796	-0.00701172\\
0.798	-0.00691569\\
0.8	-0.00680618\\
0.802	-0.0066836\\
0.804	-0.00654843\\
0.806	-0.00640112\\
0.808	-0.00624219\\
0.81	-0.00607213\\
0.812	-0.00589149\\
0.814	-0.00570079\\
0.816	-0.0055006\\
0.818	-0.00528783\\
0.82	-0.0050642\\
0.822	-0.00483311\\
0.824	-0.00459515\\
0.826	-0.00435091\\
0.828	-0.00410101\\
0.83	-0.00384605\\
0.832	-0.00358665\\
0.834	-0.00332341\\
0.836	-0.00313442\\
0.838	-0.00309651\\
0.84	-0.00305358\\
0.842	-0.00300579\\
0.844	-0.00295331\\
0.846	-0.00289631\\
0.848	-0.00283496\\
0.85	-0.00276945\\
0.852	-0.00269996\\
0.854	-0.00262668\\
0.856	-0.00254981\\
0.858	-0.00246955\\
0.86	-0.0023861\\
0.862	-0.00229966\\
0.864	-0.00221043\\
0.866	-0.00211863\\
0.868	-0.00202446\\
0.87	-0.00192814\\
0.872	-0.00182987\\
0.874	-0.00172986\\
0.876	-0.00170223\\
0.878	-0.00191412\\
0.88	-0.00211852\\
0.882	-0.0023151\\
0.884	-0.00250358\\
0.886	-0.00268368\\
0.888	-0.00285516\\
0.89	-0.00301777\\
0.892	-0.00317132\\
0.894	-0.00331562\\
0.896	-0.00345052\\
0.898	-0.00357586\\
0.9	-0.00369155\\
0.902	-0.00379748\\
0.904	-0.00389358\\
0.906	-0.0039798\\
0.908	-0.00405612\\
0.91	-0.00412253\\
0.912	-0.00417905\\
0.914	-0.00422571\\
0.916	-0.00426257\\
0.918	-0.00428971\\
0.92	-0.00430722\\
0.922	-0.00431522\\
0.924	-0.00431385\\
0.926	-0.00430325\\
0.928	-0.00428359\\
0.93	-0.00425507\\
0.932	-0.00421698\\
0.934	-0.00416616\\
0.936	-0.00410733\\
0.938	-0.00404073\\
0.94	-0.00396663\\
0.942	-0.00388531\\
0.944	-0.00379704\\
0.946	-0.00370211\\
0.948	-0.00360084\\
0.95	-0.00349353\\
0.952	-0.0033805\\
0.954	-0.00326209\\
0.956	-0.00313861\\
0.958	-0.00301042\\
0.96	-0.00287786\\
0.962	-0.00274127\\
0.964	-0.00260101\\
0.966	-0.00245742\\
0.968	-0.00231087\\
0.97	-0.00216171\\
0.972	-0.00201029\\
0.974	-0.00185697\\
0.976	-0.00170211\\
0.978	-0.00163198\\
0.98	-0.00162509\\
0.982	-0.00161522\\
0.984	-0.00160243\\
0.986	-0.00158679\\
0.988	-0.0015684\\
0.99	-0.00154733\\
0.992	-0.00152367\\
0.994	-0.00149752\\
0.996	-0.00146896\\
0.998	-0.00143811\\
1	-0.00140505\\
1.002	-0.00136989\\
1.004	-0.00133274\\
1.006	-0.0012937\\
1.008	-0.00125289\\
1.01	-0.00121042\\
1.012	-0.00116639\\
1.014	-0.00112093\\
1.016	-0.00112352\\
1.018	-0.00124302\\
1.02	-0.00135791\\
1.022	-0.00146801\\
1.024	-0.00157316\\
1.026	-0.00167323\\
1.028	-0.00176809\\
1.03	-0.00185761\\
1.032	-0.00194169\\
1.034	-0.00202025\\
1.036	-0.0020932\\
1.038	-0.0021604\\
1.04	-0.00222058\\
1.042	-0.00227504\\
1.044	-0.00232376\\
1.046	-0.00236672\\
1.048	-0.00240393\\
1.05	-0.0024354\\
1.052	-0.00246115\\
1.054	-0.00248123\\
1.056	-0.00249568\\
1.058	-0.00250457\\
1.06	-0.00250797\\
1.062	-0.00250595\\
1.064	-0.00249861\\
1.066	-0.00248605\\
1.068	-0.00246839\\
1.07	-0.00244574\\
1.072	-0.00241823\\
1.074	-0.00238599\\
1.076	-0.00234919\\
1.078	-0.00230796\\
1.08	-0.00226246\\
1.082	-0.00221287\\
1.084	-0.00215936\\
1.086	-0.00210209\\
1.088	-0.00204126\\
1.09	-0.00197706\\
1.092	-0.00190966\\
1.094	-0.00183928\\
1.096	-0.00176611\\
1.098	-0.00169034\\
1.1	-0.00161219\\
1.102	-0.00153186\\
1.104	-0.00144955\\
1.106	-0.00136547\\
1.108	-0.00127984\\
1.11	-0.00119286\\
1.112	-0.00110473\\
1.114	-0.00101567\\
1.116	-0.000925881\\
1.118	-0.000835563\\
1.12	-0.000810009\\
1.122	-0.000809024\\
1.124	-0.000806438\\
1.126	-0.00080229\\
1.128	-0.000796617\\
1.13	-0.000789464\\
1.132	-0.000780875\\
1.134	-0.000770895\\
1.136	-0.000759573\\
1.138	-0.000746958\\
1.14	-0.000733103\\
1.142	-0.00071806\\
1.144	-0.000701653\\
1.146	-0.000684022\\
1.148	-0.000665389\\
1.15	-0.000645811\\
1.152	-0.000625346\\
1.154	-0.000647171\\
1.156	-0.000713776\\
1.158	-0.000777747\\
1.16	-0.000838987\\
1.162	-0.000897411\\
1.164	-0.00095294\\
1.166	-0.0010055\\
1.168	-0.00105503\\
1.17	-0.00110147\\
1.172	-0.00114478\\
1.174	-0.0011849\\
1.176	-0.00122182\\
1.178	-0.00125549\\
1.18	-0.00128591\\
1.182	-0.00131305\\
1.184	-0.00133692\\
1.186	-0.00135752\\
1.188	-0.00137485\\
1.19	-0.00138894\\
1.192	-0.00139981\\
1.194	-0.00140748\\
1.196	-0.00141199\\
1.198	-0.00141338\\
1.2	-0.00141171\\
1.202	-0.00140703\\
1.204	-0.00139939\\
1.206	-0.00138887\\
1.208	-0.00137552\\
1.21	-0.00135943\\
1.212	-0.00134068\\
1.214	-0.00131936\\
1.216	-0.00129554\\
1.218	-0.00126932\\
1.22	-0.0012408\\
1.222	-0.00121008\\
1.224	-0.00117726\\
1.226	-0.00114245\\
1.228	-0.00110575\\
1.23	-0.00106727\\
1.232	-0.00102713\\
1.234	-0.000985445\\
1.236	-0.000942324\\
1.238	-0.000897886\\
1.24	-0.000852249\\
1.242	-0.000805524\\
1.244	-0.00075761\\
1.246	-0.000708864\\
1.248	-0.000659405\\
1.25	-0.000609352\\
1.252	-0.000558822\\
1.254	-0.000507931\\
1.256	-0.000456796\\
1.258	-0.000405531\\
1.26	-0.000371516\\
1.262	-0.000372244\\
1.264	-0.00037218\\
1.266	-0.000371339\\
1.268	-0.00036974\\
1.27	-0.000367402\\
1.272	-0.000364344\\
1.274	-0.000360589\\
1.276	-0.000356157\\
1.278	-0.000351072\\
1.28	-0.000345358\\
1.282	-0.000339041\\
1.284	-0.000332144\\
1.286	-0.000324696\\
1.288	-0.000316722\\
1.29	-0.000338748\\
1.292	-0.000377315\\
1.294	-0.000414411\\
1.296	-0.000449981\\
1.298	-0.000483972\\
1.3	-0.000516337\\
1.302	-0.000547035\\
1.304	-0.000576026\\
1.306	-0.000603276\\
1.308	-0.000628756\\
1.31	-0.00065244\\
1.312	-0.000674307\\
1.314	-0.00069434\\
1.316	-0.000712527\\
1.318	-0.000728859\\
1.32	-0.000743332\\
1.322	-0.000755946\\
1.324	-0.000766705\\
1.326	-0.000775617\\
1.328	-0.000782664\\
1.33	-0.000787838\\
1.332	-0.000791213\\
1.334	-0.000792813\\
1.336	-0.000792664\\
1.338	-0.000790795\\
1.34	-0.00078724\\
1.342	-0.000782034\\
1.344	-0.000775217\\
1.346	-0.000766831\\
1.348	-0.00075692\\
1.35	-0.000745531\\
1.352	-0.000732714\\
1.354	-0.000718522\\
1.356	-0.000703007\\
1.358	-0.000686225\\
1.36	-0.000668236\\
1.362	-0.000649097\\
1.364	-0.00062887\\
1.366	-0.000607617\\
1.368	-0.000585401\\
1.37	-0.000562287\\
1.372	-0.000538341\\
1.374	-0.000513628\\
1.376	-0.000488214\\
1.378	-0.000462168\\
1.38	-0.000435555\\
1.382	-0.000408444\\
1.384	-0.000380901\\
1.386	-0.000352995\\
1.388	-0.000324791\\
1.39	-0.000296356\\
1.392	-0.000267755\\
1.394	-0.000239054\\
1.396	-0.000210317\\
1.398	-0.000181606\\
1.4	-0.000172061\\
1.402	-0.000172213\\
1.404	-0.000171976\\
1.406	-0.000171361\\
1.408	-0.000170375\\
1.41	-0.000169029\\
1.412	-0.000167334\\
1.414	-0.0001653\\
1.416	-0.000162939\\
1.418	-0.000160264\\
1.42	-0.000157286\\
1.422	-0.000154019\\
1.424	-0.000161268\\
1.426	-0.00018382\\
1.428	-0.000205587\\
1.43	-0.000226537\\
1.432	-0.000246637\\
1.434	-0.000265857\\
1.436	-0.000284172\\
1.438	-0.000301555\\
1.44	-0.000317985\\
1.442	-0.000333443\\
1.444	-0.00034791\\
1.446	-0.000361372\\
1.448	-0.000373818\\
1.45	-0.000385235\\
1.452	-0.000395619\\
1.454	-0.000404962\\
1.456	-0.000413262\\
1.458	-0.000420518\\
1.46	-0.000426733\\
1.462	-0.000431911\\
1.464	-0.000436056\\
1.466	-0.000439179\\
1.468	-0.000441289\\
1.47	-0.000442394\\
1.472	-0.000442501\\
1.474	-0.00044164\\
1.476	-0.000439828\\
1.478	-0.000437085\\
1.48	-0.000433433\\
1.482	-0.000428896\\
1.484	-0.000423498\\
1.486	-0.000417265\\
1.488	-0.000410224\\
1.49	-0.000402406\\
1.492	-0.000393839\\
1.494	-0.000384556\\
1.496	-0.000374588\\
1.498	-0.000363968\\
1.5	-0.000352732\\
1.502	-0.000340913\\
1.504	-0.000328548\\
1.506	-0.000315671\\
1.508	-0.000302321\\
1.51	-0.000288535\\
1.512	-0.000274349\\
1.514	-0.000259801\\
1.516	-0.000244929\\
1.518	-0.000229772\\
1.52	-0.000214367\\
1.522	-0.000198751\\
1.524	-0.000182963\\
1.526	-0.00016704\\
1.528	-0.000151019\\
1.53	-0.000134937\\
1.532	-0.00011883\\
1.534	-0.000102733\\
1.536	-8.66812e-05\\
1.538	-8.17767e-05\\
1.54	-8.15302e-05\\
1.542	-8.10959e-05\\
1.544	-8.04785e-05\\
1.546	-7.96831e-05\\
1.548	-7.87152e-05\\
1.55	-7.75803e-05\\
1.552	-7.62845e-05\\
1.554	-7.48338e-05\\
1.556	-7.32349e-05\\
1.558	-7.69309e-05\\
1.56	-9.00345e-05\\
1.562	-0.00010272\\
1.564	-0.000114966\\
1.566	-0.000126754\\
1.568	-0.000138065\\
1.57	-0.000148883\\
1.572	-0.000159193\\
1.574	-0.000168979\\
1.576	-0.000178231\\
1.578	-0.000186935\\
1.58	-0.000195084\\
1.582	-0.000202667\\
1.584	-0.000209678\\
1.586	-0.000216111\\
1.588	-0.000221961\\
1.59	-0.000227221\\
1.592	-0.000231895\\
1.594	-0.000235981\\
1.596	-0.000239481\\
1.598	-0.000242397\\
1.6	-0.000244731\\
1.602	-0.00024649\\
1.604	-0.000247679\\
1.606	-0.000248304\\
1.608	-0.000248373\\
1.61	-0.000247897\\
1.612	-0.000246884\\
1.614	-0.000245347\\
1.616	-0.000243298\\
1.618	-0.000240748\\
1.62	-0.000237714\\
1.622	-0.000234208\\
1.624	-0.000230247\\
1.626	-0.000225848\\
1.628	-0.000221026\\
1.63	-0.000215801\\
1.632	-0.000210189\\
1.634	-0.00020421\\
1.636	-0.000197882\\
1.638	-0.000191227\\
1.64	-0.000184263\\
1.642	-0.000177012\\
1.644	-0.000169493\\
1.646	-0.000161728\\
1.648	-0.000153738\\
1.65	-0.000145544\\
1.652	-0.000137168\\
1.654	-0.00012863\\
1.656	-0.000119953\\
1.658	-0.000111158\\
1.66	-0.000102266\\
1.662	-9.32974e-05\\
1.664	-8.4274e-05\\
1.666	-7.52163e-05\\
1.668	-6.61446e-05\\
1.67	-5.70791e-05\\
1.672	-4.80394e-05\\
1.674	-4.26791e-05\\
1.676	-4.23009e-05\\
1.678	-4.18268e-05\\
1.68	-4.12596e-05\\
1.682	-4.06022e-05\\
1.684	-3.98578e-05\\
1.686	-3.90296e-05\\
1.688	-3.81211e-05\\
1.69	-3.71359e-05\\
1.692	-3.67485e-05\\
1.694	-4.43366e-05\\
1.696	-5.17012e-05\\
1.698	-5.88297e-05\\
1.7	-6.57104e-05\\
1.702	-7.23326e-05\\
1.704	-7.8686e-05\\
1.706	-8.47611e-05\\
1.708	-9.05494e-05\\
1.71	-9.60429e-05\\
1.712	-0.000101235\\
1.714	-0.000106118\\
1.716	-0.000110688\\
1.718	-0.000114939\\
1.72	-0.000118868\\
1.722	-0.000122471\\
1.724	-0.000125746\\
1.726	-0.000128691\\
1.728	-0.000131305\\
1.73	-0.000133588\\
1.732	-0.00013554\\
1.734	-0.000137164\\
1.736	-0.00013846\\
1.738	-0.000139431\\
1.74	-0.00014008\\
1.742	-0.000140412\\
1.744	-0.000140431\\
1.746	-0.000140143\\
1.748	-0.000139552\\
1.75	-0.000138665\\
1.752	-0.000137489\\
1.754	-0.000136032\\
1.756	-0.000134301\\
1.758	-0.000132304\\
1.76	-0.000130051\\
1.762	-0.00012755\\
1.764	-0.000124812\\
1.766	-0.000121845\\
1.768	-0.000118661\\
1.77	-0.000115269\\
1.772	-0.000111681\\
1.774	-0.000107908\\
1.776	-0.000103961\\
1.778	-9.98513e-05\\
1.78	-9.55909e-05\\
1.782	-9.11914e-05\\
1.784	-8.66648e-05\\
1.786	-8.20231e-05\\
1.788	-7.72783e-05\\
1.79	-7.24425e-05\\
1.792	-6.75277e-05\\
1.794	-6.25459e-05\\
1.796	-5.75091e-05\\
1.798	-5.24293e-05\\
1.8	-4.73181e-05\\
1.802	-4.21874e-05\\
1.804	-3.70486e-05\\
1.806	-3.1913e-05\\
1.808	-2.67918e-05\\
1.81	-2.4409e-05\\
1.812	-2.40247e-05\\
1.814	-2.35894e-05\\
1.816	-2.31049e-05\\
1.818	-2.25731e-05\\
1.82	-2.1996e-05\\
1.822	-2.13756e-05\\
1.824	-2.07138e-05\\
1.826	-2.00132e-05\\
1.828	-2.1299e-05\\
1.83	-2.56014e-05\\
1.832	-2.97776e-05\\
1.834	-3.38204e-05\\
1.836	-3.77234e-05\\
1.838	-4.14802e-05\\
1.84	-4.50851e-05\\
1.842	-4.85327e-05\\
1.844	-5.18181e-05\\
1.846	-5.49367e-05\\
1.848	-5.78844e-05\\
1.85	-6.06578e-05\\
1.852	-6.32534e-05\\
1.854	-6.56687e-05\\
1.856	-6.79013e-05\\
1.858	-6.99492e-05\\
1.86	-7.18112e-05\\
1.862	-7.3486e-05\\
1.864	-7.49732e-05\\
1.866	-7.62725e-05\\
1.868	-7.73842e-05\\
1.87	-7.83089e-05\\
1.872	-7.90477e-05\\
1.874	-7.96019e-05\\
1.876	-7.99734e-05\\
1.878	-8.01644e-05\\
1.88	-8.01772e-05\\
1.882	-8.00149e-05\\
1.884	-7.96806e-05\\
1.886	-7.91778e-05\\
1.888	-7.85104e-05\\
1.89	-7.76825e-05\\
1.892	-7.66984e-05\\
1.894	-7.5563e-05\\
1.896	-7.42811e-05\\
1.898	-7.28578e-05\\
1.9	-7.12986e-05\\
1.902	-6.9609e-05\\
1.904	-6.7795e-05\\
1.906	-6.58623e-05\\
1.908	-6.38171e-05\\
1.91	-6.16658e-05\\
1.912	-5.94147e-05\\
1.914	-5.70704e-05\\
1.916	-5.46393e-05\\
1.918	-5.21282e-05\\
1.92	-4.95439e-05\\
1.922	-4.68932e-05\\
1.924	-4.4183e-05\\
1.926	-4.142e-05\\
1.928	-3.86113e-05\\
1.93	-3.57636e-05\\
1.932	-3.28837e-05\\
1.934	-2.99785e-05\\
1.936	-2.70547e-05\\
1.938	-2.41189e-05\\
1.94	-2.11778e-05\\
1.942	-1.82377e-05\\
1.944	-1.5916e-05\\
1.946	-1.64271e-05\\
1.948	-1.68808e-05\\
1.95	-1.72761e-05\\
1.952	-1.76125e-05\\
1.954	-1.78894e-05\\
1.956	-1.81067e-05\\
1.958	-1.82642e-05\\
1.96	-1.83621e-05\\
1.962	-1.84007e-05\\
1.964	-1.83805e-05\\
1.966	-1.83024e-05\\
1.968	-1.81671e-05\\
1.97	-1.94805e-05\\
1.972	-2.17239e-05\\
1.974	-2.38839e-05\\
1.976	-2.59573e-05\\
1.978	-2.79408e-05\\
1.98	-2.98316e-05\\
1.982	-3.1627e-05\\
1.984	-3.33248e-05\\
1.986	-3.49227e-05\\
1.988	-3.64189e-05\\
1.99	-3.78118e-05\\
1.992	-3.90999e-05\\
1.994	-4.02822e-05\\
1.996	-4.13577e-05\\
1.998	-4.23259e-05\\
2	-4.31863e-05\\
};
\addplot [color=mycolor9, line width=1.0pt, forget plot]
  table[row sep=crcr]{%
-0.024	0\\
-0.022	0\\
-0.02	0\\
-0.018	0\\
-0.016	0\\
-0.014	0\\
-0.012	0\\
-0.01	0\\
-0.008	0\\
-0.006	0\\
-0.004	0\\
-0.002	0\\
0	0\\
0.002	-2.58527e-08\\
0.004	-3.18905e-07\\
0.006	-1.42751e-06\\
0.008	-4.16035e-06\\
0.01	-9.54704e-06\\
0.012	-1.88074e-05\\
0.014	-3.33257e-05\\
0.016	-5.46262e-05\\
0.018	-8.43482e-05\\
0.02	-0.000124218\\
0.022	-0.000176018\\
0.024	-0.000241548\\
0.026	-0.000322586\\
0.028	-0.000420842\\
0.03	-0.000537906\\
0.032	-0.000675204\\
0.034	-0.000833939\\
0.036	-0.00101504\\
0.038	-0.00121913\\
0.04	-0.00316702\\
0.042	-0.0094419\\
0.044	-0.0100194\\
0.046	-0.0105431\\
0.048	-0.0110033\\
0.05	-0.0113911\\
0.052	-0.0116988\\
0.054	-0.0119198\\
0.056	-0.0120491\\
0.058	-0.0120831\\
0.06	-0.0120199\\
0.062	-0.0118593\\
0.064	-0.0116025\\
0.066	-0.0112527\\
0.068	-0.0108143\\
0.07	-0.0102934\\
0.072	-0.00969731\\
0.074	-0.00903461\\
0.076	-0.00831489\\
0.078	-0.0137684\\
0.08	-0.0204634\\
0.082	-0.0213611\\
0.084	-0.0222398\\
0.086	-0.0230955\\
0.088	-0.0239243\\
0.09	-0.0247225\\
0.092	-0.0254861\\
0.094	-0.0262114\\
0.096	-0.0268949\\
0.098	-0.0275331\\
0.1	-0.0281224\\
0.102	-0.0286598\\
0.104	-0.0291421\\
0.106	-0.0295664\\
0.108	-0.02993\\
0.11	-0.0302304\\
0.112	-0.0304653\\
0.114	-0.0306325\\
0.116	-0.0294327\\
0.118	-0.0278862\\
0.12	-0.027953\\
0.122	-0.0284971\\
0.124	-0.0289891\\
0.126	-0.0294277\\
0.128	-0.0298121\\
0.13	-0.0301415\\
0.132	-0.0304158\\
0.134	-0.0306348\\
0.136	-0.0307986\\
0.138	-0.0309077\\
0.14	-0.0309628\\
0.142	-0.0309645\\
0.144	-0.0309138\\
0.146	-0.0308119\\
0.148	-0.0306601\\
0.15	-0.0304596\\
0.152	-0.0302119\\
0.154	-0.0299186\\
0.156	-0.0295811\\
0.158	-0.0292012\\
0.16	-0.0287805\\
0.162	-0.0283207\\
0.164	-0.0278234\\
0.166	-0.0272904\\
0.168	-0.0267232\\
0.17	-0.0261235\\
0.172	-0.025493\\
0.174	-0.0248332\\
0.176	-0.0241457\\
0.178	-0.023432\\
0.18	-0.0226936\\
0.182	-0.021932\\
0.184	-0.0211486\\
0.186	-0.0203447\\
0.188	-0.0195217\\
0.19	-0.0186809\\
0.192	-0.0178236\\
0.194	-0.016951\\
0.196	-0.0160644\\
0.198	-0.0151649\\
0.2	-0.0142538\\
0.202	-0.0133321\\
0.204	-0.0124011\\
0.206	-0.0114618\\
0.208	-0.0105154\\
0.21	-0.00956284\\
0.212	-0.00860529\\
0.214	-0.00764379\\
0.216	-0.00667937\\
0.218	-0.00571307\\
0.22	-0.00474591\\
0.222	-0.00377887\\
0.224	-0.00281297\\
0.226	-0.00184916\\
0.228	-0.000888423\\
0.23	6.83066e-05\\
0.232	0.00102009\\
0.234	0.00196602\\
0.236	0.00219273\\
0.238	0.00215277\\
0.24	0.00214473\\
0.242	0.00216708\\
0.244	0.00221811\\
0.246	0.00229597\\
0.248	0.00136093\\
0.25	0.000117655\\
0.252	-0.000997467\\
0.254	-0.00198681\\
0.256	-0.00285295\\
0.258	-0.00359867\\
0.26	-0.00422693\\
0.262	-0.00474088\\
0.264	-0.00514383\\
0.266	-0.00543924\\
0.268	-0.00563072\\
0.27	-0.00553266\\
0.272	-0.00530448\\
0.274	-0.00499409\\
0.276	-0.0046045\\
0.278	-0.0041389\\
0.28	-0.00360067\\
0.282	-0.00265719\\
0.284	-0.00166201\\
0.286	-0.000619324\\
0.288	0.000466645\\
0.29	0.000807115\\
0.292	0.000903892\\
0.294	0.000998472\\
0.296	0.000711821\\
0.298	-0.000413293\\
0.3	-0.00149251\\
0.302	-0.00252585\\
0.304	-0.00351337\\
0.306	-0.00445515\\
0.308	-0.00535129\\
0.31	-0.00620192\\
0.312	-0.0070072\\
0.314	-0.0077673\\
0.316	-0.00848244\\
0.318	-0.00915287\\
0.32	-0.00977885\\
0.322	-0.0103607\\
0.324	-0.0108987\\
0.326	-0.0113933\\
0.328	-0.0118448\\
0.33	-0.0122538\\
0.332	-0.0126206\\
0.334	-0.0129457\\
0.336	-0.0132298\\
0.338	-0.0134733\\
0.34	-0.0136769\\
0.342	-0.0138412\\
0.344	-0.0139669\\
0.346	-0.0140547\\
0.348	-0.0141053\\
0.35	-0.0141016\\
0.352	-0.0138585\\
0.354	-0.0141921\\
0.356	-0.0146909\\
0.358	-0.0151417\\
0.36	-0.015544\\
0.362	-0.0158973\\
0.364	-0.0162014\\
0.366	-0.0164561\\
0.368	-0.0166613\\
0.37	-0.0168172\\
0.372	-0.0169238\\
0.374	-0.0169814\\
0.376	-0.0169905\\
0.378	-0.0169516\\
0.38	-0.0168653\\
0.382	-0.0167323\\
0.384	-0.0165534\\
0.386	-0.0163296\\
0.388	-0.0160618\\
0.39	-0.0157511\\
0.392	-0.0153988\\
0.394	-0.0150062\\
0.396	-0.0145745\\
0.398	-0.0141052\\
0.4	-0.0135999\\
0.402	-0.01306\\
0.404	-0.0124873\\
0.406	-0.0118833\\
0.408	-0.01125\\
0.41	-0.010589\\
0.412	-0.00990217\\
0.414	-0.00919145\\
0.416	-0.00845873\\
0.418	-0.00770596\\
0.42	-0.00693512\\
0.422	-0.0061482\\
0.424	-0.00534721\\
0.426	-0.00453418\\
0.428	-0.00371114\\
0.43	-0.00288012\\
0.432	-0.00204315\\
0.434	-0.00108032\\
0.436	-0.000104102\\
0.438	0.00087226\\
0.44	0.00184645\\
0.442	0.00243089\\
0.444	0.00228535\\
0.446	0.00181599\\
0.448	0.00112109\\
0.45	0.000441524\\
0.452	-0.00022204\\
0.454	-0.000868969\\
0.456	-0.00149866\\
0.458	-0.00211053\\
0.46	-0.00270405\\
0.462	-0.00327869\\
0.464	-0.00383398\\
0.466	-0.00436947\\
0.468	-0.0052254\\
0.47	-0.00605841\\
0.472	-0.00685784\\
0.474	-0.00762251\\
0.476	-0.00835135\\
0.478	-0.00904335\\
0.48	-0.00969763\\
0.482	-0.0103134\\
0.484	-0.0108899\\
0.486	-0.0114265\\
0.488	-0.0119227\\
0.49	-0.0123781\\
0.492	-0.0127725\\
0.494	-0.0130749\\
0.496	-0.0133395\\
0.498	-0.0135661\\
0.5	-0.0137548\\
0.502	-0.0139057\\
0.504	-0.0140188\\
0.506	-0.0140945\\
0.508	-0.014133\\
0.51	-0.0141348\\
0.512	-0.0141003\\
0.514	-0.0140301\\
0.516	-0.0139247\\
0.518	-0.0137849\\
0.52	-0.0136115\\
0.522	-0.0134052\\
0.524	-0.0131669\\
0.526	-0.0128976\\
0.528	-0.0125982\\
0.53	-0.0122699\\
0.532	-0.0119137\\
0.534	-0.0115308\\
0.536	-0.0111224\\
0.538	-0.0106896\\
0.54	-0.0102339\\
0.542	-0.00975638\\
0.544	-0.00925852\\
0.546	-0.00874164\\
0.548	-0.00820715\\
0.55	-0.00765645\\
0.552	-0.00709097\\
0.554	-0.00651217\\
0.556	-0.00620614\\
0.558	-0.00589995\\
0.56	-0.00558696\\
0.562	-0.00526776\\
0.564	-0.00494294\\
0.566	-0.00461312\\
0.568	-0.00427888\\
0.57	-0.00394084\\
0.572	-0.00359958\\
0.574	-0.00325572\\
0.576	-0.00290984\\
0.578	-0.00256254\\
0.58	-0.0022144\\
0.582	-0.00186601\\
0.584	-0.00151794\\
0.586	-0.00117077\\
0.588	-0.000825041\\
0.59	-0.00074204\\
0.592	-0.00142047\\
0.594	-0.00208433\\
0.596	-0.00273244\\
0.598	-0.0033637\\
0.6	-0.00397701\\
0.602	-0.00457136\\
0.604	-0.00514579\\
0.606	-0.00569009\\
0.608	-0.00619213\\
0.61	-0.0066721\\
0.612	-0.00712929\\
0.614	-0.00756308\\
0.616	-0.00797288\\
0.618	-0.00835818\\
0.62	-0.00871851\\
0.622	-0.00905348\\
0.624	-0.00936272\\
0.626	-0.00964597\\
0.628	-0.00990299\\
0.63	-0.0101336\\
0.632	-0.0103377\\
0.634	-0.0105153\\
0.636	-0.0106639\\
0.638	-0.0107551\\
0.64	-0.0108217\\
0.642	-0.010864\\
0.644	-0.0108822\\
0.646	-0.0108764\\
0.648	-0.0108471\\
0.65	-0.0107946\\
0.652	-0.0107192\\
0.654	-0.0106215\\
0.656	-0.0105018\\
0.658	-0.0103608\\
0.66	-0.010199\\
0.662	-0.0100169\\
0.664	-0.00981531\\
0.666	-0.00959478\\
0.668	-0.00935603\\
0.67	-0.00909982\\
0.672	-0.00882689\\
0.674	-0.00853802\\
0.676	-0.00823404\\
0.678	-0.00791577\\
0.68	-0.00758407\\
0.682	-0.0072398\\
0.684	-0.00688386\\
0.686	-0.00651713\\
0.688	-0.00614054\\
0.69	-0.00575501\\
0.692	-0.00536145\\
0.694	-0.00510927\\
0.696	-0.00500563\\
0.698	-0.00489459\\
0.7	-0.00477646\\
0.702	-0.00465154\\
0.704	-0.00452017\\
0.706	-0.00438267\\
0.708	-0.00423936\\
0.71	-0.00409059\\
0.712	-0.00393671\\
0.714	-0.00377805\\
0.716	-0.00361497\\
0.718	-0.00344781\\
0.72	-0.00327693\\
0.722	-0.00310269\\
0.724	-0.00292543\\
0.726	-0.00274552\\
0.728	-0.00256331\\
0.73	-0.00237914\\
0.732	-0.00219338\\
0.734	-0.00211463\\
0.736	-0.00250631\\
0.738	-0.00287139\\
0.74	-0.00322444\\
0.742	-0.00356487\\
0.744	-0.00389217\\
0.746	-0.00420584\\
0.748	-0.00450543\\
0.75	-0.00479051\\
0.752	-0.00506071\\
0.754	-0.00531567\\
0.756	-0.00555508\\
0.758	-0.00577867\\
0.76	-0.00598621\\
0.762	-0.00617749\\
0.764	-0.00635235\\
0.766	-0.00651068\\
0.768	-0.00665237\\
0.77	-0.00677739\\
0.772	-0.00688571\\
0.774	-0.00697735\\
0.776	-0.00705237\\
0.778	-0.00711086\\
0.78	-0.00715294\\
0.782	-0.00717878\\
0.784	-0.00718854\\
0.786	-0.00718247\\
0.788	-0.00716081\\
0.79	-0.00712383\\
0.792	-0.00707185\\
0.794	-0.0070052\\
0.796	-0.00692425\\
0.798	-0.00682938\\
0.8	-0.006721\\
0.802	-0.00659955\\
0.804	-0.00646548\\
0.806	-0.00631928\\
0.808	-0.00616142\\
0.81	-0.00599244\\
0.812	-0.00581285\\
0.814	-0.00562319\\
0.816	-0.00542403\\
0.818	-0.00521593\\
0.82	-0.00499416\\
0.822	-0.00476453\\
0.824	-0.00452801\\
0.826	-0.00428519\\
0.828	-0.00403669\\
0.83	-0.0037831\\
0.832	-0.00352505\\
0.834	-0.00326314\\
0.836	-0.00307544\\
0.838	-0.00303881\\
0.84	-0.00299714\\
0.842	-0.00295057\\
0.844	-0.00289929\\
0.846	-0.00284346\\
0.848	-0.00278326\\
0.85	-0.00271887\\
0.852	-0.00265048\\
0.854	-0.00257827\\
0.856	-0.00250245\\
0.858	-0.00242321\\
0.86	-0.00234075\\
0.862	-0.00225527\\
0.864	-0.00216699\\
0.866	-0.0020761\\
0.868	-0.00198282\\
0.87	-0.00188736\\
0.872	-0.00178992\\
0.874	-0.00169072\\
0.876	-0.00166388\\
0.878	-0.00187653\\
0.88	-0.00208167\\
0.882	-0.00227897\\
0.884	-0.00246814\\
0.886	-0.0026489\\
0.888	-0.00282102\\
0.89	-0.00298425\\
0.892	-0.00313839\\
0.894	-0.00328327\\
0.896	-0.00341871\\
0.898	-0.00354459\\
0.9	-0.00366078\\
0.902	-0.0037672\\
0.904	-0.00386377\\
0.906	-0.00395044\\
0.908	-0.00402719\\
0.91	-0.00409401\\
0.912	-0.00415092\\
0.914	-0.00419796\\
0.916	-0.00423518\\
0.918	-0.00426266\\
0.92	-0.0042805\\
0.922	-0.00428882\\
0.924	-0.00428774\\
0.926	-0.00427743\\
0.928	-0.00425804\\
0.93	-0.00422978\\
0.932	-0.00419284\\
0.934	-0.00414496\\
0.936	-0.00408654\\
0.938	-0.00402035\\
0.94	-0.00394663\\
0.942	-0.00386568\\
0.944	-0.00377776\\
0.946	-0.00368318\\
0.948	-0.00358224\\
0.95	-0.00347525\\
0.952	-0.00336253\\
0.954	-0.0032444\\
0.956	-0.00312121\\
0.958	-0.00299329\\
0.96	-0.00286099\\
0.962	-0.00272465\\
0.964	-0.00258463\\
0.966	-0.00244128\\
0.968	-0.00229495\\
0.97	-0.002146\\
0.972	-0.00199479\\
0.974	-0.00184167\\
0.976	-0.00168699\\
0.978	-0.00161705\\
0.98	-0.00161034\\
0.982	-0.00160064\\
0.984	-0.00158801\\
0.986	-0.00157254\\
0.988	-0.0015543\\
0.99	-0.00153338\\
0.992	-0.00150987\\
0.994	-0.00148386\\
0.996	-0.00145545\\
0.998	-0.00142473\\
1	-0.0013918\\
1.002	-0.00135677\\
1.004	-0.00131975\\
1.006	-0.00128084\\
1.008	-0.00124016\\
1.01	-0.00119781\\
1.012	-0.0011539\\
1.014	-0.00110856\\
1.016	-0.00111126\\
1.018	-0.00123089\\
1.02	-0.00134589\\
1.022	-0.0014561\\
1.024	-0.00156137\\
1.026	-0.00166156\\
1.028	-0.00175653\\
1.03	-0.00184595\\
1.032	-0.00192879\\
1.034	-0.00200612\\
1.036	-0.00207787\\
1.038	-0.00214398\\
1.04	-0.00220439\\
1.042	-0.00225909\\
1.044	-0.00230803\\
1.046	-0.00235122\\
1.048	-0.00238865\\
1.05	-0.00242034\\
1.052	-0.00244631\\
1.054	-0.0024666\\
1.056	-0.00248127\\
1.058	-0.00249036\\
1.06	-0.00249396\\
1.062	-0.00249215\\
1.064	-0.002485\\
1.066	-0.00247264\\
1.068	-0.00245517\\
1.07	-0.00243271\\
1.072	-0.00240539\\
1.074	-0.00237334\\
1.076	-0.00233672\\
1.078	-0.00229567\\
1.08	-0.00225035\\
1.082	-0.00220093\\
1.084	-0.00214759\\
1.086	-0.0020905\\
1.088	-0.00202983\\
1.09	-0.00196579\\
1.092	-0.00189856\\
1.094	-0.00182834\\
1.096	-0.00175532\\
1.098	-0.00167971\\
1.1	-0.00160172\\
1.102	-0.00152153\\
1.104	-0.00143938\\
1.106	-0.00135545\\
1.108	-0.00126996\\
1.11	-0.00118312\\
1.112	-0.00109514\\
1.114	-0.00100621\\
1.116	-0.00091656\\
1.118	-0.000826377\\
1.12	-0.000800956\\
1.122	-0.000800102\\
1.124	-0.000797645\\
1.126	-0.000793623\\
1.128	-0.000788076\\
1.13	-0.000781046\\
1.132	-0.000772578\\
1.134	-0.000762718\\
1.136	-0.000751514\\
1.138	-0.000739016\\
1.14	-0.000725275\\
1.142	-0.000710345\\
1.144	-0.000694279\\
1.146	-0.000677134\\
1.148	-0.000658955\\
1.15	-0.000639505\\
1.152	-0.000619165\\
1.154	-0.000641112\\
1.156	-0.000707838\\
1.158	-0.000771926\\
1.16	-0.000833281\\
1.162	-0.000891818\\
1.164	-0.000947456\\
1.166	-0.00100013\\
1.168	-0.00104976\\
1.17	-0.0010963\\
1.172	-0.00113971\\
1.174	-0.00117993\\
1.176	-0.00121694\\
1.178	-0.00125071\\
1.18	-0.00128122\\
1.182	-0.00130845\\
1.184	-0.00133241\\
1.186	-0.00135309\\
1.188	-0.0013705\\
1.19	-0.00138467\\
1.192	-0.00139562\\
1.194	-0.00140336\\
1.196	-0.00140795\\
1.198	-0.00140942\\
1.2	-0.00140782\\
1.202	-0.0014032\\
1.204	-0.00139563\\
1.206	-0.00138517\\
1.208	-0.00137189\\
1.21	-0.00135586\\
1.212	-0.00133717\\
1.214	-0.0013159\\
1.216	-0.00129214\\
1.218	-0.00126598\\
1.22	-0.00123751\\
1.222	-0.00120685\\
1.224	-0.00117408\\
1.226	-0.00113931\\
1.228	-0.00110266\\
1.23	-0.00106423\\
1.232	-0.00102414\\
1.234	-0.000982495\\
1.236	-0.000939417\\
1.238	-0.000895021\\
1.24	-0.000849425\\
1.242	-0.000802748\\
1.244	-0.000754946\\
1.246	-0.000706238\\
1.248	-0.000656817\\
1.25	-0.000606801\\
1.252	-0.000556308\\
1.254	-0.000505454\\
1.256	-0.000454355\\
1.258	-0.000403126\\
1.26	-0.000369147\\
1.262	-0.000369911\\
1.264	-0.000369881\\
1.266	-0.000369076\\
1.268	-0.000367511\\
1.27	-0.000365208\\
1.272	-0.000362184\\
1.274	-0.000358462\\
1.276	-0.000354064\\
1.278	-0.000349012\\
1.28	-0.000343332\\
1.282	-0.000337047\\
1.284	-0.000330183\\
1.286	-0.000322767\\
1.288	-0.000314825\\
1.29	-0.000336882\\
1.292	-0.000375481\\
1.294	-0.000412608\\
1.296	-0.000448208\\
1.298	-0.00048223\\
1.3	-0.000514626\\
1.302	-0.000545353\\
1.304	-0.000574374\\
1.306	-0.000601654\\
1.308	-0.000627162\\
1.31	-0.000650875\\
1.312	-0.000672736\\
1.314	-0.000692709\\
1.316	-0.000710837\\
1.318	-0.00072711\\
1.32	-0.000741525\\
1.322	-0.000754081\\
1.324	-0.000764784\\
1.326	-0.00077364\\
1.328	-0.000780662\\
1.33	-0.000785866\\
1.332	-0.000789271\\
1.334	-0.0007909\\
1.336	-0.000790779\\
1.338	-0.000788938\\
1.34	-0.00078541\\
1.342	-0.000780232\\
1.344	-0.000773442\\
1.346	-0.000765082\\
1.348	-0.000755197\\
1.35	-0.000743835\\
1.352	-0.000731043\\
1.354	-0.000716876\\
1.356	-0.000701385\\
1.358	-0.000684628\\
1.36	-0.000666663\\
1.362	-0.000647547\\
1.364	-0.000627343\\
1.366	-0.000606113\\
1.368	-0.00058392\\
1.37	-0.000560829\\
1.372	-0.000536904\\
1.374	-0.000512213\\
1.376	-0.00048682\\
1.378	-0.000460795\\
1.38	-0.000434203\\
1.382	-0.000407112\\
1.384	-0.000379589\\
1.386	-0.000351702\\
1.388	-0.000323518\\
1.39	-0.000295102\\
1.392	-0.00026652\\
1.394	-0.000237838\\
1.396	-0.000209119\\
1.398	-0.000180426\\
1.4	-0.000170899\\
1.402	-0.000171068\\
1.404	-0.000170849\\
1.406	-0.00017025\\
1.408	-0.000169281\\
1.41	-0.000167951\\
1.412	-0.000166272\\
1.414	-0.000164254\\
1.416	-0.000161909\\
1.418	-0.000159249\\
1.42	-0.000156287\\
1.422	-0.000153035\\
1.424	-0.000160299\\
1.426	-0.000182865\\
1.428	-0.000204647\\
1.43	-0.000225611\\
1.432	-0.000245725\\
1.434	-0.000264959\\
1.436	-0.000283287\\
1.438	-0.000300683\\
1.44	-0.000317126\\
1.442	-0.000332597\\
1.444	-0.000347077\\
1.446	-0.000360552\\
1.448	-0.000373009\\
1.45	-0.000384439\\
1.452	-0.000394834\\
1.454	-0.000404189\\
1.456	-0.000412501\\
1.458	-0.000419769\\
1.46	-0.000425995\\
1.462	-0.000431183\\
1.464	-0.00043534\\
1.466	-0.000438473\\
1.468	-0.000440594\\
1.47	-0.000441714\\
1.472	-0.000441848\\
1.474	-0.000441012\\
1.476	-0.000439224\\
1.478	-0.000436505\\
1.48	-0.000432876\\
1.482	-0.00042836\\
1.484	-0.000422982\\
1.486	-0.000416768\\
1.488	-0.000409739\\
1.49	-0.00040193\\
1.492	-0.000393372\\
1.494	-0.000384097\\
1.496	-0.000374137\\
1.498	-0.000363526\\
1.5	-0.000352297\\
1.502	-0.000340486\\
1.504	-0.000328128\\
1.506	-0.000315259\\
1.508	-0.000301917\\
1.51	-0.000288137\\
1.512	-0.000273958\\
1.514	-0.000259417\\
1.516	-0.000244553\\
1.518	-0.000229402\\
1.52	-0.000214003\\
1.522	-0.000198394\\
1.524	-0.000182612\\
1.526	-0.000166695\\
1.528	-0.00015068\\
1.53	-0.000134604\\
1.532	-0.000118502\\
1.534	-0.000102411\\
1.536	-8.63648e-05\\
1.538	-8.14656e-05\\
1.54	-8.12243e-05\\
1.542	-8.07951e-05\\
1.544	-8.01827e-05\\
1.546	-7.93923e-05\\
1.548	-7.84291e-05\\
1.55	-7.72989e-05\\
1.552	-7.60076e-05\\
1.554	-7.45615e-05\\
1.556	-7.29669e-05\\
1.558	-7.6665e-05\\
1.56	-8.97679e-05\\
1.562	-0.000102452\\
1.564	-0.000114698\\
1.566	-0.000126485\\
1.568	-0.000137795\\
1.57	-0.000148612\\
1.572	-0.00015892\\
1.574	-0.000168705\\
1.576	-0.000177956\\
1.578	-0.000186659\\
1.58	-0.000194806\\
1.582	-0.000202388\\
1.584	-0.000209398\\
1.586	-0.00021583\\
1.588	-0.00022168\\
1.59	-0.000226944\\
1.592	-0.000231622\\
1.594	-0.000235712\\
1.596	-0.000239216\\
1.598	-0.000242136\\
1.6	-0.000244475\\
1.602	-0.000246237\\
1.604	-0.000247429\\
1.606	-0.000248058\\
1.608	-0.000248132\\
1.61	-0.000247659\\
1.612	-0.00024665\\
1.614	-0.000245116\\
1.616	-0.00024307\\
1.618	-0.000240524\\
1.62	-0.000237493\\
1.622	-0.000233991\\
1.624	-0.000230033\\
1.626	-0.000225637\\
1.628	-0.000220819\\
1.63	-0.000215596\\
1.632	-0.000209987\\
1.634	-0.000204011\\
1.636	-0.000197687\\
1.638	-0.000191035\\
1.64	-0.000184074\\
1.642	-0.000176825\\
1.644	-0.000169309\\
1.646	-0.000161547\\
1.648	-0.00015356\\
1.65	-0.000145369\\
1.652	-0.000136995\\
1.654	-0.00012846\\
1.656	-0.000119786\\
1.658	-0.000110993\\
1.66	-0.000102103\\
1.662	-9.31376e-05\\
1.664	-8.41166e-05\\
1.666	-7.50613e-05\\
1.668	-6.5992e-05\\
1.67	-5.69288e-05\\
1.672	-4.78914e-05\\
1.674	-4.25333e-05\\
1.676	-4.21574e-05\\
1.678	-4.16855e-05\\
1.68	-4.11204e-05\\
1.682	-4.04651e-05\\
1.684	-3.97228e-05\\
1.686	-3.88967e-05\\
1.688	-3.79903e-05\\
1.69	-3.70071e-05\\
1.692	-3.66216e-05\\
1.694	-4.42117e-05\\
1.696	-5.15781e-05\\
1.698	-5.87085e-05\\
1.7	-6.55911e-05\\
1.702	-7.22151e-05\\
1.704	-7.85703e-05\\
1.706	-8.46472e-05\\
1.708	-9.04372e-05\\
1.71	-9.59324e-05\\
1.712	-0.000101126\\
1.714	-0.000106011\\
1.716	-0.000110582\\
1.718	-0.000114835\\
1.72	-0.000118765\\
1.722	-0.00012237\\
1.724	-0.000125646\\
1.726	-0.000128593\\
1.728	-0.000131209\\
1.73	-0.000133493\\
1.732	-0.000135447\\
1.734	-0.000137072\\
1.736	-0.000138369\\
1.738	-0.000139342\\
1.74	-0.000139993\\
1.742	-0.000140326\\
1.744	-0.000140346\\
1.746	-0.000140059\\
1.748	-0.000139469\\
1.75	-0.000138584\\
1.752	-0.000137409\\
1.754	-0.000135953\\
1.756	-0.000134223\\
1.758	-0.000132228\\
1.76	-0.000129976\\
1.762	-0.000127476\\
1.764	-0.000124739\\
1.766	-0.000121773\\
1.768	-0.00011859\\
1.77	-0.0001152\\
1.772	-0.000111613\\
1.774	-0.000107841\\
1.776	-0.000103894\\
1.778	-9.97859e-05\\
1.78	-9.55264e-05\\
1.782	-9.1128e-05\\
1.784	-8.66024e-05\\
1.786	-8.19616e-05\\
1.788	-7.72178e-05\\
1.79	-7.23829e-05\\
1.792	-6.74689e-05\\
1.794	-6.24881e-05\\
1.796	-5.74522e-05\\
1.798	-5.23732e-05\\
1.8	-4.72629e-05\\
1.802	-4.2133e-05\\
1.804	-3.6995e-05\\
1.806	-3.18603e-05\\
1.808	-2.67399e-05\\
1.81	-2.43579e-05\\
1.812	-2.39744e-05\\
1.814	-2.35398e-05\\
1.816	-2.30561e-05\\
1.818	-2.2525e-05\\
1.82	-2.19487e-05\\
1.822	-2.13291e-05\\
1.824	-2.06684e-05\\
1.826	-1.99688e-05\\
1.828	-2.12555e-05\\
1.83	-2.55589e-05\\
1.832	-2.97359e-05\\
1.834	-3.37797e-05\\
1.836	-3.76835e-05\\
1.838	-4.14411e-05\\
1.84	-4.50468e-05\\
1.842	-4.84952e-05\\
1.844	-5.17812e-05\\
1.846	-5.49005e-05\\
1.848	-5.7849e-05\\
1.85	-6.0623e-05\\
1.852	-6.32193e-05\\
1.854	-6.56351e-05\\
1.856	-6.78683e-05\\
1.858	-6.99168e-05\\
1.86	-7.17793e-05\\
1.862	-7.34546e-05\\
1.864	-7.49423e-05\\
1.866	-7.62421e-05\\
1.868	-7.73543e-05\\
1.87	-7.82795e-05\\
1.872	-7.90187e-05\\
1.874	-7.95733e-05\\
1.876	-7.99452e-05\\
1.878	-8.01365e-05\\
1.88	-8.01497e-05\\
1.882	-7.99878e-05\\
1.884	-7.96538e-05\\
1.886	-7.91514e-05\\
1.888	-7.84843e-05\\
1.89	-7.76567e-05\\
1.892	-7.66729e-05\\
1.894	-7.55378e-05\\
1.896	-7.42561e-05\\
1.898	-7.28331e-05\\
1.9	-7.12742e-05\\
1.902	-6.95849e-05\\
1.904	-6.77711e-05\\
1.906	-6.58386e-05\\
1.908	-6.37937e-05\\
1.91	-6.16427e-05\\
1.912	-5.93918e-05\\
1.914	-5.70476e-05\\
1.916	-5.46168e-05\\
1.918	-5.21061e-05\\
1.92	-4.95221e-05\\
1.922	-4.68718e-05\\
1.924	-4.41618e-05\\
1.926	-4.13992e-05\\
1.928	-3.85908e-05\\
1.93	-3.57434e-05\\
1.932	-3.28639e-05\\
1.934	-2.9959e-05\\
1.936	-2.70354e-05\\
1.938	-2.40999e-05\\
1.94	-2.11591e-05\\
1.942	-1.82194e-05\\
1.944	-1.58979e-05\\
1.946	-1.64093e-05\\
1.948	-1.68633e-05\\
1.95	-1.72589e-05\\
1.952	-1.75955e-05\\
1.954	-1.78727e-05\\
1.956	-1.80902e-05\\
1.958	-1.82479e-05\\
1.96	-1.83461e-05\\
1.962	-1.83849e-05\\
1.964	-1.83651e-05\\
1.966	-1.82871e-05\\
1.968	-1.81521e-05\\
1.97	-1.94657e-05\\
1.972	-2.17093e-05\\
1.974	-2.38696e-05\\
1.976	-2.59431e-05\\
1.978	-2.79269e-05\\
1.98	-2.98179e-05\\
1.982	-3.16136e-05\\
1.984	-3.33115e-05\\
1.986	-3.49097e-05\\
1.988	-3.64061e-05\\
1.99	-3.77991e-05\\
1.992	-3.90874e-05\\
1.994	-4.02699e-05\\
1.996	-4.13456e-05\\
1.998	-4.2314e-05\\
2	-4.31746e-05\\
};
\addplot [color=mycolor10, line width=1.0pt, forget plot]
  table[row sep=crcr]{%
-0.024	0\\
-0.022	0\\
-0.02	0\\
-0.018	0\\
-0.016	0\\
-0.014	0\\
-0.012	0\\
-0.01	0\\
-0.008	0\\
-0.006	0\\
-0.004	0\\
-0.002	0\\
0	0\\
0.002	-2.46543e-08\\
0.004	-3.10727e-07\\
0.006	-1.40427e-06\\
0.008	-4.11486e-06\\
0.01	-9.4757e-06\\
0.012	-1.87127e-05\\
0.014	-3.32183e-05\\
0.016	-5.45275e-05\\
0.018	-8.42926e-05\\
0.02	-0.000124255\\
0.022	-0.000176214\\
0.024	-0.00024199\\
0.026	-0.000323382\\
0.028	-0.000422123\\
0.03	-0.000539831\\
0.032	-0.000677956\\
0.034	-0.00083773\\
0.036	-0.00102011\\
0.038	-0.00122575\\
0.04	-0.00317548\\
0.042	-0.00949191\\
0.044	-0.0100737\\
0.046	-0.0106017\\
0.048	-0.0110662\\
0.05	-0.0114583\\
0.052	-0.0117702\\
0.054	-0.0119955\\
0.056	-0.0121289\\
0.058	-0.0121671\\
0.06	-0.012108\\
0.062	-0.0119514\\
0.064	-0.0116986\\
0.066	-0.0113527\\
0.068	-0.0109182\\
0.07	-0.0104011\\
0.072	-0.00980892\\
0.074	-0.00915012\\
0.076	-0.00843433\\
0.078	-0.0138918\\
0.08	-0.0207498\\
0.082	-0.021656\\
0.084	-0.0225429\\
0.086	-0.0234064\\
0.088	-0.0242428\\
0.09	-0.025048\\
0.092	-0.0258183\\
0.094	-0.0265499\\
0.096	-0.0272393\\
0.098	-0.0278829\\
0.1	-0.0284773\\
0.102	-0.0290193\\
0.104	-0.0295058\\
0.106	-0.029934\\
0.108	-0.030301\\
0.11	-0.0306045\\
0.112	-0.030842\\
0.114	-0.0310115\\
0.116	-0.0297724\\
0.118	-0.0282287\\
0.12	-0.028298\\
0.122	-0.0288443\\
0.124	-0.0293383\\
0.126	-0.0297787\\
0.128	-0.0301646\\
0.13	-0.0304954\\
0.132	-0.0307708\\
0.134	-0.0309908\\
0.136	-0.0311554\\
0.138	-0.0312651\\
0.14	-0.0313205\\
0.142	-0.0313225\\
0.144	-0.0312719\\
0.146	-0.03117\\
0.148	-0.0310179\\
0.15	-0.030817\\
0.152	-0.0305688\\
0.154	-0.0302748\\
0.156	-0.0299365\\
0.158	-0.0295556\\
0.16	-0.0291338\\
0.162	-0.0286727\\
0.164	-0.0281739\\
0.166	-0.0276392\\
0.168	-0.0270702\\
0.17	-0.0264686\\
0.172	-0.025836\\
0.174	-0.0251738\\
0.176	-0.0244838\\
0.178	-0.0237675\\
0.18	-0.0230262\\
0.182	-0.0222615\\
0.184	-0.0214748\\
0.186	-0.0206674\\
0.188	-0.0198407\\
0.19	-0.0189961\\
0.192	-0.0181347\\
0.194	-0.0172578\\
0.196	-0.0163667\\
0.198	-0.0154626\\
0.2	-0.0145466\\
0.202	-0.01362\\
0.204	-0.0126837\\
0.206	-0.011739\\
0.208	-0.010787\\
0.21	-0.00982873\\
0.212	-0.00886528\\
0.214	-0.00789771\\
0.216	-0.00692709\\
0.218	-0.00595444\\
0.22	-0.00498078\\
0.222	-0.00400714\\
0.224	-0.0030345\\
0.226	-0.00206385\\
0.228	-0.00109616\\
0.23	-0.000132377\\
0.232	0.00082655\\
0.234	0.0017797\\
0.236	0.00201371\\
0.238	0.00198112\\
0.24	0.00198051\\
0.242	0.00201034\\
0.244	0.00206891\\
0.246	0.00215435\\
0.248	0.00122692\\
0.25	-8.71109e-06\\
0.252	-0.00111616\\
0.254	-0.00209782\\
0.256	-0.00295626\\
0.258	-0.00369428\\
0.26	-0.00431484\\
0.262	-0.00482109\\
0.264	-0.00521636\\
0.266	-0.0055041\\
0.268	-0.00568793\\
0.27	-0.00558881\\
0.272	-0.00535407\\
0.274	-0.00503723\\
0.276	-0.00464131\\
0.278	-0.00416949\\
0.28	-0.00362516\\
0.282	-0.00288047\\
0.284	-0.00188101\\
0.286	-0.000833948\\
0.288	0.000256479\\
0.29	0.000601488\\
0.292	0.00070288\\
0.294	0.000802149\\
0.296	0.000520254\\
0.298	-0.00060004\\
0.3	-0.00167437\\
0.302	-0.00270277\\
0.304	-0.00368531\\
0.306	-0.00462205\\
0.308	-0.00551311\\
0.31	-0.00635861\\
0.312	-0.00715873\\
0.314	-0.00791364\\
0.316	-0.00862356\\
0.318	-0.00928873\\
0.32	-0.00990944\\
0.322	-0.010486\\
0.324	-0.0110187\\
0.326	-0.011508\\
0.328	-0.0119542\\
0.33	-0.0123578\\
0.332	-0.0127193\\
0.334	-0.0130391\\
0.336	-0.0133178\\
0.338	-0.0135561\\
0.34	-0.0137544\\
0.342	-0.0139134\\
0.344	-0.0140338\\
0.346	-0.0141163\\
0.348	-0.0141405\\
0.35	-0.0139201\\
0.352	-0.013678\\
0.354	-0.0140128\\
0.356	-0.0145128\\
0.358	-0.014965\\
0.36	-0.0153688\\
0.362	-0.0157237\\
0.364	-0.0160295\\
0.366	-0.016286\\
0.368	-0.016493\\
0.37	-0.0166507\\
0.372	-0.0167593\\
0.374	-0.016819\\
0.376	-0.0168301\\
0.378	-0.0167933\\
0.38	-0.0167091\\
0.382	-0.0165783\\
0.384	-0.0164015\\
0.386	-0.0161799\\
0.388	-0.0159143\\
0.39	-0.0156058\\
0.392	-0.0152557\\
0.394	-0.0148652\\
0.396	-0.0144356\\
0.398	-0.0139685\\
0.4	-0.0134652\\
0.402	-0.0129274\\
0.404	-0.0123566\\
0.406	-0.0117546\\
0.408	-0.0111232\\
0.41	-0.010464\\
0.412	-0.00977902\\
0.414	-0.00907004\\
0.416	-0.00833899\\
0.418	-0.00758783\\
0.42	-0.00681853\\
0.422	-0.00603308\\
0.424	-0.00523349\\
0.426	-0.00442178\\
0.428	-0.00359998\\
0.43	-0.00277013\\
0.432	-0.00193425\\
0.434	-0.000973646\\
0.436	4.47046e-06\\
0.438	0.000982646\\
0.44	0.00195856\\
0.442	0.00254464\\
0.444	0.00240066\\
0.446	0.00193277\\
0.448	0.00123925\\
0.45	0.000560989\\
0.452	-0.00010135\\
0.454	-0.000747133\\
0.456	-0.00137575\\
0.458	-0.00198663\\
0.46	-0.00257923\\
0.462	-0.00315302\\
0.464	-0.00370753\\
0.466	-0.00424231\\
0.468	-0.00509758\\
0.47	-0.00593001\\
0.472	-0.00672892\\
0.474	-0.00749313\\
0.476	-0.00822156\\
0.478	-0.00891322\\
0.48	-0.0095672\\
0.482	-0.0101827\\
0.484	-0.010759\\
0.486	-0.0112955\\
0.488	-0.0117916\\
0.49	-0.012247\\
0.492	-0.0126356\\
0.494	-0.0129391\\
0.496	-0.0132047\\
0.498	-0.0134324\\
0.5	-0.0136221\\
0.502	-0.013774\\
0.504	-0.0138881\\
0.506	-0.0139647\\
0.508	-0.0140042\\
0.51	-0.0140068\\
0.512	-0.0139732\\
0.514	-0.0139038\\
0.516	-0.0137993\\
0.518	-0.0136604\\
0.52	-0.0134877\\
0.522	-0.0132822\\
0.524	-0.0130448\\
0.526	-0.0127762\\
0.528	-0.0124777\\
0.53	-0.0121501\\
0.532	-0.0117947\\
0.534	-0.0114125\\
0.536	-0.0110048\\
0.538	-0.0105728\\
0.54	-0.0101177\\
0.542	-0.00964095\\
0.544	-0.0091438\\
0.546	-0.00862764\\
0.548	-0.00809386\\
0.55	-0.00754387\\
0.552	-0.00697911\\
0.554	-0.00640102\\
0.556	-0.00609571\\
0.558	-0.00579025\\
0.56	-0.00547798\\
0.562	-0.00515951\\
0.564	-0.00483543\\
0.566	-0.00450635\\
0.568	-0.00417286\\
0.57	-0.00383557\\
0.572	-0.00349508\\
0.574	-0.00315198\\
0.576	-0.00280688\\
0.578	-0.00246036\\
0.58	-0.00211302\\
0.582	-0.00176543\\
0.584	-0.00141817\\
0.586	-0.00107182\\
0.588	-0.000726918\\
0.59	-0.000644753\\
0.592	-0.00132403\\
0.594	-0.00198874\\
0.596	-0.00263772\\
0.598	-0.00326984\\
0.6	-0.00388403\\
0.602	-0.00447927\\
0.604	-0.00505459\\
0.606	-0.00558803\\
0.608	-0.00609075\\
0.61	-0.0065714\\
0.612	-0.00702928\\
0.614	-0.00746376\\
0.616	-0.00787425\\
0.618	-0.00826024\\
0.62	-0.00862127\\
0.622	-0.00895693\\
0.624	-0.00926688\\
0.626	-0.00955083\\
0.628	-0.00980856\\
0.63	-0.0100399\\
0.632	-0.0102447\\
0.634	-0.010423\\
0.636	-0.0105747\\
0.638	-0.0106922\\
0.64	-0.0107602\\
0.642	-0.0108039\\
0.644	-0.0108234\\
0.646	-0.010819\\
0.648	-0.010791\\
0.65	-0.0107398\\
0.652	-0.0106657\\
0.654	-0.0105692\\
0.656	-0.0104507\\
0.658	-0.0103109\\
0.66	-0.0101502\\
0.662	-0.00996924\\
0.664	-0.0097687\\
0.666	-0.00954922\\
0.668	-0.0093115\\
0.67	-0.00905627\\
0.672	-0.00878429\\
0.674	-0.00849635\\
0.676	-0.00819326\\
0.678	-0.00787584\\
0.68	-0.00754496\\
0.682	-0.00720148\\
0.684	-0.00684629\\
0.686	-0.00648029\\
0.688	-0.00610438\\
0.69	-0.00571951\\
0.692	-0.00532658\\
0.694	-0.00507499\\
0.696	-0.00497191\\
0.698	-0.00486141\\
0.7	-0.00474378\\
0.702	-0.00461934\\
0.704	-0.00448841\\
0.706	-0.00435132\\
0.708	-0.00420841\\
0.71	-0.00406001\\
0.712	-0.00390647\\
0.714	-0.00374812\\
0.716	-0.00358534\\
0.718	-0.00341845\\
0.72	-0.00324783\\
0.722	-0.00307381\\
0.724	-0.00289677\\
0.726	-0.00271705\\
0.728	-0.00253502\\
0.73	-0.00235101\\
0.732	-0.0021654\\
0.734	-0.00207297\\
0.736	-0.00245016\\
0.738	-0.00281588\\
0.74	-0.00316956\\
0.742	-0.00351062\\
0.744	-0.00383854\\
0.746	-0.00415283\\
0.748	-0.00445303\\
0.75	-0.00473872\\
0.752	-0.00500952\\
0.754	-0.00526508\\
0.756	-0.00550508\\
0.758	-0.00572926\\
0.76	-0.00593738\\
0.762	-0.00612923\\
0.764	-0.00630467\\
0.766	-0.00646356\\
0.768	-0.00660581\\
0.77	-0.00673138\\
0.772	-0.00684025\\
0.774	-0.00693243\\
0.776	-0.00700799\\
0.778	-0.00706701\\
0.78	-0.00710962\\
0.782	-0.00713597\\
0.784	-0.00714625\\
0.786	-0.00714069\\
0.788	-0.00711953\\
0.79	-0.00708305\\
0.792	-0.00703156\\
0.794	-0.0069654\\
0.796	-0.00688493\\
0.798	-0.00679054\\
0.8	-0.00668263\\
0.802	-0.00656165\\
0.804	-0.00642805\\
0.806	-0.0062823\\
0.808	-0.00612489\\
0.81	-0.00595635\\
0.812	-0.0057772\\
0.814	-0.00558798\\
0.816	-0.00538925\\
0.818	-0.00518158\\
0.82	-0.00496555\\
0.822	-0.00474061\\
0.824	-0.00450435\\
0.826	-0.0042618\\
0.828	-0.00401355\\
0.83	-0.00376023\\
0.832	-0.00350243\\
0.834	-0.00324077\\
0.836	-0.00305334\\
0.838	-0.00301695\\
0.84	-0.00297553\\
0.842	-0.00292921\\
0.844	-0.00287817\\
0.846	-0.00282259\\
0.848	-0.00276263\\
0.85	-0.00269848\\
0.852	-0.00263032\\
0.854	-0.00255835\\
0.856	-0.00248276\\
0.858	-0.00240375\\
0.86	-0.00232152\\
0.862	-0.00223627\\
0.864	-0.00214821\\
0.866	-0.00205754\\
0.868	-0.00196448\\
0.87	-0.00186923\\
0.872	-0.00177201\\
0.874	-0.00167303\\
0.876	-0.0016464\\
0.878	-0.00185926\\
0.88	-0.0020646\\
0.882	-0.0022621\\
0.884	-0.00245147\\
0.886	-0.00263244\\
0.888	-0.00280475\\
0.89	-0.00296817\\
0.892	-0.00312251\\
0.894	-0.00326757\\
0.896	-0.00340321\\
0.898	-0.00352927\\
0.9	-0.00364564\\
0.902	-0.00375224\\
0.904	-0.00384898\\
0.906	-0.00393583\\
0.908	-0.00401275\\
0.91	-0.00407975\\
0.912	-0.00413683\\
0.914	-0.00418403\\
0.916	-0.00422141\\
0.918	-0.00424906\\
0.92	-0.00426706\\
0.922	-0.00427553\\
0.924	-0.0042746\\
0.926	-0.00426444\\
0.928	-0.00424521\\
0.93	-0.00421709\\
0.932	-0.00418029\\
0.934	-0.00413503\\
0.936	-0.00407695\\
0.938	-0.00401092\\
0.94	-0.00393737\\
0.942	-0.00385657\\
0.944	-0.00376881\\
0.946	-0.00367437\\
0.948	-0.00357357\\
0.95	-0.00346672\\
0.952	-0.00335413\\
0.954	-0.00323613\\
0.956	-0.00311306\\
0.958	-0.00298526\\
0.96	-0.00285308\\
0.962	-0.00271685\\
0.964	-0.00257693\\
0.966	-0.00243368\\
0.968	-0.00228745\\
0.97	-0.0021386\\
0.972	-0.00198747\\
0.974	-0.00183444\\
0.976	-0.00167985\\
0.978	-0.00160999\\
0.98	-0.00160336\\
0.982	-0.00159373\\
0.984	-0.00158118\\
0.986	-0.00156578\\
0.988	-0.00154761\\
0.99	-0.00152675\\
0.992	-0.0015033\\
0.994	-0.00147736\\
0.996	-0.001449\\
0.998	-0.00141834\\
1	-0.00138547\\
1.002	-0.00135049\\
1.004	-0.00131352\\
1.006	-0.00127467\\
1.008	-0.00123403\\
1.01	-0.00119173\\
1.012	-0.00114787\\
1.014	-0.00110257\\
1.016	-0.00110533\\
1.018	-0.00122499\\
1.02	-0.00134004\\
1.022	-0.0014503\\
1.024	-0.00155562\\
1.026	-0.00165485\\
1.028	-0.00174854\\
1.03	-0.00183691\\
1.032	-0.00191987\\
1.034	-0.00199732\\
1.036	-0.00206918\\
1.038	-0.00213541\\
1.04	-0.00219594\\
1.042	-0.00225074\\
1.044	-0.0022998\\
1.046	-0.00234309\\
1.048	-0.00238063\\
1.05	-0.00241243\\
1.052	-0.0024385\\
1.054	-0.0024589\\
1.056	-0.00247367\\
1.058	-0.00248287\\
1.06	-0.00248656\\
1.062	-0.00248485\\
1.064	-0.0024778\\
1.066	-0.00246554\\
1.068	-0.00244816\\
1.07	-0.0024258\\
1.072	-0.00239857\\
1.074	-0.00236661\\
1.076	-0.00233008\\
1.078	-0.00228912\\
1.08	-0.00224389\\
1.082	-0.00219456\\
1.084	-0.0021413\\
1.086	-0.00208429\\
1.088	-0.00202372\\
1.09	-0.00195976\\
1.092	-0.00189261\\
1.094	-0.00182247\\
1.096	-0.00174953\\
1.098	-0.001674\\
1.1	-0.00159608\\
1.102	-0.00151597\\
1.104	-0.00143389\\
1.106	-0.00135004\\
1.108	-0.00126463\\
1.11	-0.00117786\\
1.112	-0.00108995\\
1.114	-0.00100109\\
1.116	-0.000911509\\
1.118	-0.000821394\\
1.12	-0.000796041\\
1.122	-0.000795254\\
1.124	-0.000792863\\
1.126	-0.000788907\\
1.128	-0.000783424\\
1.13	-0.000776458\\
1.132	-0.000768052\\
1.134	-0.000758254\\
1.136	-0.00074711\\
1.138	-0.000734672\\
1.14	-0.000720991\\
1.142	-0.000706119\\
1.144	-0.000690111\\
1.146	-0.000673022\\
1.148	-0.000654909\\
1.15	-0.00063583\\
1.152	-0.000615842\\
1.154	-0.000638122\\
1.156	-0.000705023\\
1.158	-0.000769157\\
1.16	-0.000830558\\
1.162	-0.000889139\\
1.164	-0.000944821\\
1.166	-0.000997533\\
1.168	-0.00104721\\
1.17	-0.0010938\\
1.172	-0.00113724\\
1.174	-0.0011775\\
1.176	-0.00121455\\
1.178	-0.00124836\\
1.18	-0.0012789\\
1.182	-0.00130617\\
1.184	-0.00133016\\
1.186	-0.00135088\\
1.188	-0.00136833\\
1.19	-0.00138253\\
1.192	-0.0013935\\
1.194	-0.00140128\\
1.196	-0.0014059\\
1.198	-0.0014074\\
1.2	-0.00140583\\
1.202	-0.00140124\\
1.204	-0.0013937\\
1.206	-0.00138327\\
1.208	-0.00137001\\
1.21	-0.00135401\\
1.212	-0.00133535\\
1.214	-0.0013141\\
1.216	-0.00129036\\
1.218	-0.00126423\\
1.22	-0.00123578\\
1.222	-0.00120514\\
1.224	-0.00117239\\
1.226	-0.00113765\\
1.228	-0.00110102\\
1.23	-0.00106261\\
1.232	-0.00102254\\
1.234	-0.000980914\\
1.236	-0.000937855\\
1.238	-0.000893478\\
1.24	-0.000847901\\
1.242	-0.000801241\\
1.244	-0.000753618\\
1.246	-0.00070501\\
1.248	-0.000655602\\
1.25	-0.000605599\\
1.252	-0.000555119\\
1.254	-0.000504277\\
1.256	-0.000453191\\
1.258	-0.000401973\\
1.26	-0.000368006\\
1.262	-0.000368781\\
1.264	-0.000368762\\
1.266	-0.000367968\\
1.268	-0.000366414\\
1.27	-0.000364121\\
1.272	-0.000361108\\
1.274	-0.000357396\\
1.276	-0.000353008\\
1.278	-0.000347966\\
1.28	-0.000342296\\
1.282	-0.00033602\\
1.284	-0.000329166\\
1.286	-0.00032176\\
1.288	-0.000313828\\
1.29	-0.000335894\\
1.292	-0.000374503\\
1.294	-0.00041164\\
1.296	-0.000447249\\
1.298	-0.00048128\\
1.3	-0.000513671\\
1.302	-0.000544334\\
1.304	-0.000573288\\
1.306	-0.0006005\\
1.308	-0.00062594\\
1.31	-0.000649584\\
1.312	-0.00067141\\
1.314	-0.000691402\\
1.316	-0.000709548\\
1.318	-0.000725838\\
1.32	-0.000740271\\
1.322	-0.000752844\\
1.324	-0.000763564\\
1.326	-0.000772437\\
1.328	-0.000779475\\
1.33	-0.000784695\\
1.332	-0.000788116\\
1.334	-0.000789761\\
1.336	-0.000789656\\
1.338	-0.00078783\\
1.34	-0.000784318\\
1.342	-0.000779155\\
1.344	-0.00077238\\
1.346	-0.000764034\\
1.348	-0.000754164\\
1.35	-0.000742815\\
1.352	-0.000730038\\
1.354	-0.000715884\\
1.356	-0.000700408\\
1.358	-0.000683664\\
1.36	-0.000665712\\
1.362	-0.00064661\\
1.364	-0.000626419\\
1.366	-0.000605202\\
1.368	-0.000583021\\
1.37	-0.000559942\\
1.372	-0.00053603\\
1.374	-0.00051135\\
1.376	-0.00048597\\
1.378	-0.000459956\\
1.38	-0.000433376\\
1.382	-0.000406297\\
1.384	-0.000378786\\
1.386	-0.00035091\\
1.388	-0.000322736\\
1.39	-0.000294331\\
1.392	-0.00026576\\
1.394	-0.000237089\\
1.396	-0.00020838\\
1.398	-0.000179698\\
1.4	-0.00017018\\
1.402	-0.00017036\\
1.404	-0.00017015\\
1.406	-0.000169561\\
1.408	-0.000168602\\
1.41	-0.000167282\\
1.412	-0.000165612\\
1.414	-0.000163604\\
1.416	-0.000161268\\
1.418	-0.000158617\\
1.42	-0.000155663\\
1.422	-0.00015242\\
1.424	-0.000159693\\
1.426	-0.000182268\\
1.428	-0.000204058\\
1.43	-0.00022503\\
1.432	-0.000245152\\
1.434	-0.000264395\\
1.436	-0.00028273\\
1.438	-0.000300135\\
1.44	-0.000316586\\
1.442	-0.000332064\\
1.444	-0.000346551\\
1.446	-0.000360034\\
1.448	-0.000372499\\
1.45	-0.000383936\\
1.452	-0.000394338\\
1.454	-0.0004037\\
1.456	-0.000412018\\
1.458	-0.000419293\\
1.46	-0.000425526\\
1.462	-0.000430721\\
1.464	-0.000434885\\
1.466	-0.000438025\\
1.468	-0.000440151\\
1.47	-0.000441278\\
1.472	-0.000441418\\
1.474	-0.000440588\\
1.476	-0.000438806\\
1.478	-0.000436093\\
1.48	-0.00043247\\
1.482	-0.000427959\\
1.484	-0.000422587\\
1.486	-0.000416379\\
1.488	-0.000409362\\
1.49	-0.000401567\\
1.492	-0.000393022\\
1.494	-0.000383759\\
1.496	-0.000373811\\
1.498	-0.00036321\\
1.5	-0.000351992\\
1.502	-0.00034019\\
1.504	-0.000327841\\
1.506	-0.00031498\\
1.508	-0.000301644\\
1.51	-0.000287871\\
1.512	-0.000273698\\
1.514	-0.000259163\\
1.516	-0.000244303\\
1.518	-0.000229157\\
1.52	-0.000213762\\
1.522	-0.000198156\\
1.524	-0.000182377\\
1.526	-0.000166463\\
1.528	-0.00015045\\
1.53	-0.000134375\\
1.532	-0.000118274\\
1.534	-0.000102184\\
1.536	-8.61386e-05\\
1.538	-8.12398e-05\\
1.54	-8.09986e-05\\
1.542	-8.05692e-05\\
1.544	-7.99564e-05\\
1.546	-7.91653e-05\\
1.548	-7.82014e-05\\
1.55	-7.70702e-05\\
1.552	-7.57777e-05\\
1.554	-7.43303e-05\\
1.556	-7.27343e-05\\
1.558	-7.64331e-05\\
1.56	-8.95393e-05\\
1.562	-0.000102227\\
1.564	-0.000114476\\
1.566	-0.000126266\\
1.568	-0.000137579\\
1.57	-0.000148399\\
1.572	-0.00015871\\
1.574	-0.000168499\\
1.576	-0.000177752\\
1.578	-0.000186458\\
1.58	-0.000194608\\
1.582	-0.000202192\\
1.584	-0.000209205\\
1.586	-0.00021564\\
1.588	-0.000221492\\
1.59	-0.00022676\\
1.592	-0.00023144\\
1.594	-0.000235533\\
1.596	-0.000239039\\
1.598	-0.000241961\\
1.6	-0.000244303\\
1.602	-0.000246068\\
1.604	-0.000247262\\
1.606	-0.000247893\\
1.608	-0.000247969\\
1.61	-0.000247499\\
1.612	-0.000246492\\
1.614	-0.000244961\\
1.616	-0.000242917\\
1.618	-0.000240373\\
1.62	-0.000237344\\
1.622	-0.000233844\\
1.624	-0.000229889\\
1.626	-0.000225494\\
1.628	-0.000220678\\
1.63	-0.000215457\\
1.632	-0.000209851\\
1.634	-0.000203877\\
1.636	-0.000197554\\
1.638	-0.000190904\\
1.64	-0.000183945\\
1.642	-0.000176698\\
1.644	-0.000169184\\
1.646	-0.000161423\\
1.648	-0.000153438\\
1.65	-0.000145248\\
1.652	-0.000136876\\
1.654	-0.000128343\\
1.656	-0.000119671\\
1.658	-0.00011088\\
1.66	-0.000101991\\
1.662	-9.30271e-05\\
1.664	-8.40077e-05\\
1.666	-7.49539e-05\\
1.668	-6.58861e-05\\
1.67	-5.68244e-05\\
1.672	-4.77886e-05\\
1.674	-4.24319e-05\\
1.676	-4.20574e-05\\
1.678	-4.1587e-05\\
1.68	-4.10233e-05\\
1.682	-4.03694e-05\\
1.684	-3.96285e-05\\
1.686	-3.88037e-05\\
1.688	-3.78986e-05\\
1.69	-3.69167e-05\\
1.692	-3.65325e-05\\
1.694	-4.41239e-05\\
1.696	-5.14916e-05\\
1.698	-5.86232e-05\\
1.7	-6.5507e-05\\
1.702	-7.21322e-05\\
1.704	-7.84886e-05\\
1.706	-8.45666e-05\\
1.708	-9.03578e-05\\
1.71	-9.58542e-05\\
1.712	-0.000101049\\
1.714	-0.000105935\\
1.716	-0.000110507\\
1.718	-0.000114761\\
1.72	-0.000118693\\
1.722	-0.000122298\\
1.724	-0.000125576\\
1.726	-0.000128523\\
1.728	-0.00013114\\
1.73	-0.000133425\\
1.732	-0.00013538\\
1.734	-0.000137006\\
1.736	-0.000138304\\
1.738	-0.000139278\\
1.74	-0.00013993\\
1.742	-0.000140264\\
1.744	-0.000140285\\
1.746	-0.000139999\\
1.748	-0.00013941\\
1.75	-0.000138525\\
1.752	-0.000137352\\
1.754	-0.000135896\\
1.756	-0.000134167\\
1.758	-0.000132173\\
1.76	-0.000129921\\
1.762	-0.000127423\\
1.764	-0.000124686\\
1.766	-0.000121721\\
1.768	-0.000118539\\
1.77	-0.000115149\\
1.772	-0.000111563\\
1.774	-0.000107791\\
1.776	-0.000103846\\
1.778	-9.97382e-05\\
1.78	-9.54794e-05\\
1.782	-9.10817e-05\\
1.784	-8.65567e-05\\
1.786	-8.19167e-05\\
1.788	-7.71735e-05\\
1.79	-7.23392e-05\\
1.792	-6.74259e-05\\
1.794	-6.24457e-05\\
1.796	-5.74104e-05\\
1.798	-5.2332e-05\\
1.8	-4.72224e-05\\
1.802	-4.20931e-05\\
1.804	-3.69557e-05\\
1.806	-3.18215e-05\\
1.808	-2.67016e-05\\
1.81	-2.43202e-05\\
1.812	-2.39372e-05\\
1.814	-2.35032e-05\\
1.816	-2.302e-05\\
1.818	-2.24895e-05\\
1.82	-2.19137e-05\\
1.822	-2.12946e-05\\
1.824	-2.06344e-05\\
1.826	-1.99353e-05\\
1.828	-2.12225e-05\\
1.83	-2.55263e-05\\
1.832	-2.97039e-05\\
1.834	-3.37481e-05\\
1.836	-3.76524e-05\\
1.838	-4.14105e-05\\
1.84	-4.50166e-05\\
1.842	-4.84654e-05\\
1.844	-5.17519e-05\\
1.846	-5.48716e-05\\
1.848	-5.78205e-05\\
1.85	-6.05949e-05\\
1.852	-6.31916e-05\\
1.854	-6.56079e-05\\
1.856	-6.78414e-05\\
1.858	-6.98903e-05\\
1.86	-7.17532e-05\\
1.862	-7.34289e-05\\
1.864	-7.4917e-05\\
1.866	-7.62172e-05\\
1.868	-7.73297e-05\\
1.87	-7.82552e-05\\
1.872	-7.89948e-05\\
1.874	-7.95498e-05\\
1.876	-7.9922e-05\\
1.878	-8.01137e-05\\
1.88	-8.01272e-05\\
1.882	-7.99656e-05\\
1.884	-7.9632e-05\\
1.886	-7.91298e-05\\
1.888	-7.84631e-05\\
1.89	-7.76358e-05\\
1.892	-7.66523e-05\\
1.894	-7.55175e-05\\
1.896	-7.42361e-05\\
1.898	-7.28134e-05\\
1.9	-7.12548e-05\\
1.902	-6.95658e-05\\
1.904	-6.77522e-05\\
1.906	-6.582e-05\\
1.908	-6.37754e-05\\
1.91	-6.16246e-05\\
1.912	-5.9374e-05\\
1.914	-5.70301e-05\\
1.916	-5.45996e-05\\
1.918	-5.2089e-05\\
1.92	-4.95053e-05\\
1.922	-4.68552e-05\\
1.924	-4.41455e-05\\
1.926	-4.13832e-05\\
1.928	-3.8575e-05\\
1.93	-3.57278e-05\\
1.932	-3.28485e-05\\
1.934	-2.99438e-05\\
1.936	-2.70205e-05\\
1.938	-2.40852e-05\\
1.94	-2.11446e-05\\
1.942	-1.82051e-05\\
1.944	-1.58838e-05\\
1.946	-1.63954e-05\\
1.948	-1.68496e-05\\
1.95	-1.72454e-05\\
1.952	-1.75822e-05\\
1.954	-1.78596e-05\\
1.956	-1.80773e-05\\
1.958	-1.82352e-05\\
1.96	-1.83335e-05\\
1.962	-1.83726e-05\\
1.964	-1.83529e-05\\
1.966	-1.82751e-05\\
1.968	-1.81403e-05\\
1.97	-1.94541e-05\\
1.972	-2.16978e-05\\
1.974	-2.38583e-05\\
1.976	-2.5932e-05\\
1.978	-2.79159e-05\\
1.98	-2.9807e-05\\
1.982	-3.16029e-05\\
1.984	-3.3301e-05\\
1.986	-3.48993e-05\\
1.988	-3.63958e-05\\
1.99	-3.7789e-05\\
1.992	-3.90775e-05\\
1.994	-4.02601e-05\\
1.996	-4.1336e-05\\
1.998	-4.23045e-05\\
2	-4.31652e-05\\
};
\addplot [color=mycolor11, line width=1.0pt, forget plot]
  table[row sep=crcr]{%
-0.024	0\\
-0.022	0\\
-0.02	0\\
-0.018	0\\
-0.016	0\\
-0.014	0\\
-0.012	0\\
-0.01	0\\
-0.008	0\\
-0.006	0\\
-0.004	0\\
-0.002	0\\
0	0\\
0.002	-0.000271864\\
0.004	-0.00101003\\
0.006	-0.00209995\\
0.008	-0.00343154\\
0.01	-0.00490177\\
0.012	-0.0064168\\
0.014	-0.0078937\\
0.016	-0.0092617\\
0.018	-0.010463\\
0.02	-0.0114529\\
0.022	-0.0122002\\
0.024	-0.0126861\\
0.026	-0.0129041\\
0.028	-0.0128585\\
0.03	-0.0125632\\
0.032	-0.0120405\\
0.034	-0.0113192\\
0.036	-0.0104335\\
0.038	-0.00942091\\
0.04	-0.0100415\\
0.042	-0.01427\\
0.044	-0.0152754\\
0.046	-0.0161847\\
0.048	-0.0169792\\
0.05	-0.017642\\
0.052	-0.0181591\\
0.054	-0.0185191\\
0.056	-0.0187137\\
0.058	-0.0187379\\
0.06	-0.01859\\
0.062	-0.0182715\\
0.064	-0.0177871\\
0.066	-0.0171446\\
0.068	-0.0163544\\
0.07	-0.0154298\\
0.072	-0.014386\\
0.074	-0.0132401\\
0.076	-0.0120107\\
0.078	-0.0169373\\
0.08	-0.0237315\\
0.082	-0.024676\\
0.084	-0.0255908\\
0.086	-0.0264728\\
0.088	-0.0273193\\
0.09	-0.0281273\\
0.092	-0.0288939\\
0.094	-0.0296166\\
0.096	-0.0302924\\
0.098	-0.0309189\\
0.1	-0.0314933\\
0.102	-0.032013\\
0.104	-0.0324755\\
0.106	-0.0328784\\
0.108	-0.0332192\\
0.11	-0.0334957\\
0.112	-0.0337057\\
0.114	-0.033847\\
0.116	-0.0329671\\
0.118	-0.0314324\\
0.12	-0.0315033\\
0.122	-0.0320445\\
0.124	-0.0325269\\
0.126	-0.0329498\\
0.128	-0.0333128\\
0.13	-0.0336157\\
0.132	-0.0338586\\
0.134	-0.0340419\\
0.136	-0.0341663\\
0.138	-0.0342326\\
0.14	-0.0342417\\
0.142	-0.0341948\\
0.144	-0.0340931\\
0.146	-0.0339381\\
0.148	-0.0337313\\
0.15	-0.0334743\\
0.152	-0.0331686\\
0.154	-0.032816\\
0.156	-0.0324182\\
0.158	-0.0319768\\
0.16	-0.0314938\\
0.162	-0.0309706\\
0.164	-0.0304092\\
0.166	-0.0298111\\
0.168	-0.0291781\\
0.17	-0.0285118\\
0.172	-0.0278139\\
0.174	-0.0270859\\
0.176	-0.0263294\\
0.178	-0.0255459\\
0.18	-0.024737\\
0.182	-0.023904\\
0.184	-0.0230485\\
0.186	-0.0221719\\
0.188	-0.0212755\\
0.19	-0.0203608\\
0.192	-0.019429\\
0.194	-0.0184816\\
0.196	-0.0175198\\
0.198	-0.016545\\
0.2	-0.0155584\\
0.202	-0.0145613\\
0.204	-0.0135551\\
0.206	-0.0125408\\
0.208	-0.0115199\\
0.21	-0.0104935\\
0.212	-0.00946278\\
0.214	-0.00842899\\
0.216	-0.00739331\\
0.218	-0.0063569\\
0.22	-0.0053209\\
0.222	-0.00428644\\
0.224	-0.00325461\\
0.226	-0.00222649\\
0.228	-0.00120315\\
0.23	-0.000185592\\
0.232	0.000825173\\
0.234	0.00182818\\
0.236	0.00211002\\
0.238	0.00212321\\
0.24	0.00216634\\
0.242	0.00223787\\
0.244	0.0023361\\
0.246	0.00245918\\
0.248	0.00156743\\
0.25	0.00036555\\
0.252	-0.000710027\\
0.254	-0.00166162\\
0.256	-0.00249174\\
0.258	-0.0032031\\
0.26	-0.00379861\\
0.262	-0.00428132\\
0.264	-0.00465447\\
0.266	-0.00492144\\
0.268	-0.00508577\\
0.27	-0.00515111\\
0.272	-0.00497859\\
0.274	-0.00461384\\
0.276	-0.00417176\\
0.278	-0.00365582\\
0.28	-0.00306964\\
0.282	-0.00241693\\
0.284	-0.00170156\\
0.286	-0.000845075\\
0.288	0.000256479\\
0.29	0.000612501\\
0.292	0.000724792\\
0.294	0.000834838\\
0.296	0.000563601\\
0.298	-0.000546161\\
0.3	-0.00161009\\
0.302	-0.00262822\\
0.304	-0.0036006\\
0.306	-0.00452733\\
0.308	-0.00540852\\
0.31	-0.00624428\\
0.312	-0.00703479\\
0.314	-0.00778024\\
0.316	-0.00848083\\
0.318	-0.00913683\\
0.32	-0.00974849\\
0.322	-0.0103161\\
0.324	-0.0108401\\
0.326	-0.0113208\\
0.328	-0.0117585\\
0.33	-0.0121538\\
0.332	-0.0125071\\
0.334	-0.0128189\\
0.336	-0.0130898\\
0.338	-0.0133203\\
0.34	-0.013511\\
0.342	-0.0136626\\
0.344	-0.0137757\\
0.346	-0.0138511\\
0.348	-0.0137281\\
0.35	-0.0135124\\
0.352	-0.0132742\\
0.354	-0.0136121\\
0.356	-0.0141147\\
0.358	-0.0145687\\
0.36	-0.0149736\\
0.362	-0.015329\\
0.364	-0.0156347\\
0.366	-0.0158904\\
0.368	-0.0160963\\
0.37	-0.0162523\\
0.372	-0.0163587\\
0.374	-0.0164157\\
0.376	-0.0164238\\
0.378	-0.0163835\\
0.38	-0.0162955\\
0.382	-0.0161605\\
0.384	-0.0159793\\
0.386	-0.015753\\
0.388	-0.0154824\\
0.39	-0.0151687\\
0.392	-0.0148132\\
0.394	-0.0144172\\
0.396	-0.013982\\
0.398	-0.013509\\
0.4	-0.0129998\\
0.402	-0.012456\\
0.404	-0.0118792\\
0.406	-0.0112712\\
0.408	-0.0106336\\
0.41	-0.00996839\\
0.412	-0.00927732\\
0.414	-0.00856232\\
0.416	-0.00782532\\
0.418	-0.00706828\\
0.42	-0.0062932\\
0.422	-0.00550207\\
0.424	-0.00469693\\
0.426	-0.00387981\\
0.428	-0.00305276\\
0.43	-0.00221781\\
0.432	-0.00137702\\
0.434	-0.000532407\\
0.436	0.000436152\\
0.438	0.00141489\\
0.44	0.00239125\\
0.442	0.00297767\\
0.444	0.00283392\\
0.446	0.00236616\\
0.448	0.00167267\\
0.45	0.000994336\\
0.452	0.000331826\\
0.454	-0.000314221\\
0.456	-0.000943198\\
0.458	-0.00155453\\
0.46	-0.00214766\\
0.462	-0.00272207\\
0.464	-0.00327729\\
0.466	-0.00381285\\
0.468	-0.00466899\\
0.47	-0.00550236\\
0.472	-0.00630228\\
0.474	-0.00706758\\
0.476	-0.00779716\\
0.478	-0.00849005\\
0.48	-0.00914532\\
0.482	-0.00976218\\
0.484	-0.0103399\\
0.486	-0.0108778\\
0.488	-0.0113521\\
0.49	-0.0117374\\
0.492	-0.0120851\\
0.494	-0.0123951\\
0.496	-0.0126671\\
0.498	-0.012901\\
0.5	-0.0130969\\
0.502	-0.0132548\\
0.504	-0.0133749\\
0.506	-0.0134574\\
0.508	-0.0135026\\
0.51	-0.0135109\\
0.512	-0.0134828\\
0.514	-0.0134188\\
0.516	-0.0133197\\
0.518	-0.0131859\\
0.52	-0.0130184\\
0.522	-0.0128179\\
0.524	-0.0125853\\
0.526	-0.0123216\\
0.528	-0.0120278\\
0.53	-0.0117049\\
0.532	-0.011354\\
0.534	-0.0109763\\
0.536	-0.0105729\\
0.538	-0.0101452\\
0.54	-0.00969434\\
0.542	-0.00922169\\
0.544	-0.00872858\\
0.546	-0.00821639\\
0.548	-0.0076865\\
0.55	-0.00714034\\
0.552	-0.00657933\\
0.554	-0.00600493\\
0.556	-0.00570326\\
0.558	-0.00540136\\
0.56	-0.0050926\\
0.562	-0.00477758\\
0.564	-0.00445691\\
0.566	-0.00413118\\
0.568	-0.003801\\
0.57	-0.00346697\\
0.572	-0.0031297\\
0.574	-0.00278979\\
0.576	-0.00244783\\
0.578	-0.00210443\\
0.58	-0.00176016\\
0.582	-0.00141562\\
0.584	-0.00107139\\
0.586	-0.000728026\\
0.588	-0.000386101\\
0.59	-0.000306888\\
0.592	-0.000989093\\
0.594	-0.00165672\\
0.596	-0.0023086\\
0.598	-0.00294361\\
0.6	-0.00356067\\
0.602	-0.00415877\\
0.604	-0.00473695\\
0.606	-0.00529429\\
0.608	-0.00581342\\
0.61	-0.00629651\\
0.612	-0.00675681\\
0.614	-0.0071937\\
0.616	-0.00760659\\
0.618	-0.00799498\\
0.62	-0.00835838\\
0.622	-0.00869641\\
0.624	-0.00900872\\
0.626	-0.00929502\\
0.628	-0.00955508\\
0.63	-0.00978873\\
0.632	-0.00999586\\
0.634	-0.0101764\\
0.636	-0.0103304\\
0.638	-0.0104579\\
0.64	-0.0105589\\
0.642	-0.0106268\\
0.644	-0.0106492\\
0.646	-0.0106476\\
0.648	-0.0106223\\
0.65	-0.0105736\\
0.652	-0.010502\\
0.654	-0.0104079\\
0.656	-0.0102918\\
0.658	-0.0101541\\
0.66	-0.00999555\\
0.662	-0.00981665\\
0.664	-0.00961805\\
0.666	-0.00940043\\
0.668	-0.00916449\\
0.67	-0.00891096\\
0.672	-0.00864061\\
0.674	-0.00835422\\
0.676	-0.00805262\\
0.678	-0.00773662\\
0.68	-0.00740709\\
0.682	-0.0070649\\
0.684	-0.00671094\\
0.686	-0.00634611\\
0.688	-0.00597133\\
0.69	-0.00558752\\
0.692	-0.00519562\\
0.694	-0.00494501\\
0.696	-0.00484286\\
0.698	-0.00473325\\
0.7	-0.00461647\\
0.702	-0.00449285\\
0.704	-0.00436271\\
0.706	-0.00422638\\
0.708	-0.0040842\\
0.71	-0.0039365\\
0.712	-0.00378364\\
0.714	-0.00362595\\
0.716	-0.00346381\\
0.718	-0.00329754\\
0.72	-0.00312753\\
0.722	-0.00295411\\
0.724	-0.00277766\\
0.726	-0.00259852\\
0.728	-0.00241327\\
0.73	-0.00221551\\
0.732	-0.00201616\\
0.734	-0.00192387\\
0.736	-0.00230265\\
0.738	-0.00266995\\
0.74	-0.0030252\\
0.742	-0.00336781\\
0.744	-0.00369728\\
0.746	-0.00401309\\
0.748	-0.00431481\\
0.75	-0.004602\\
0.752	-0.00487428\\
0.754	-0.00513131\\
0.756	-0.00537277\\
0.758	-0.0055984\\
0.76	-0.00580795\\
0.762	-0.00600123\\
0.764	-0.00617807\\
0.766	-0.00633835\\
0.768	-0.00648198\\
0.77	-0.00660892\\
0.772	-0.00671914\\
0.774	-0.00681266\\
0.776	-0.00688955\\
0.778	-0.00694989\\
0.78	-0.0069938\\
0.782	-0.00702144\\
0.784	-0.007033\\
0.786	-0.0070287\\
0.788	-0.00700879\\
0.79	-0.00697355\\
0.792	-0.00692328\\
0.794	-0.00685834\\
0.796	-0.00677907\\
0.798	-0.00668586\\
0.8	-0.00657913\\
0.802	-0.00645931\\
0.804	-0.00632686\\
0.806	-0.00618225\\
0.808	-0.00602597\\
0.81	-0.00585855\\
0.812	-0.0056805\\
0.814	-0.00549238\\
0.816	-0.00529473\\
0.818	-0.00508813\\
0.82	-0.00487315\\
0.822	-0.00465039\\
0.824	-0.00442044\\
0.826	-0.0041839\\
0.828	-0.00393667\\
0.83	-0.00368422\\
0.832	-0.00342728\\
0.834	-0.00316647\\
0.836	-0.00297986\\
0.838	-0.00294429\\
0.84	-0.00290366\\
0.842	-0.00285812\\
0.844	-0.00280785\\
0.846	-0.00275302\\
0.848	-0.0026938\\
0.85	-0.00263038\\
0.852	-0.00256293\\
0.854	-0.00249167\\
0.856	-0.00241677\\
0.858	-0.00233843\\
0.86	-0.00225687\\
0.862	-0.00217227\\
0.864	-0.00208486\\
0.866	-0.00199482\\
0.868	-0.00190239\\
0.87	-0.00180776\\
0.872	-0.00171114\\
0.874	-0.00161275\\
0.876	-0.00158671\\
0.878	-0.00180015\\
0.88	-0.00200606\\
0.882	-0.00220413\\
0.884	-0.00239406\\
0.886	-0.00257557\\
0.888	-0.00274842\\
0.89	-0.00291239\\
0.892	-0.00306725\\
0.894	-0.00321284\\
0.896	-0.00334899\\
0.898	-0.00347557\\
0.9	-0.00359245\\
0.902	-0.00369955\\
0.904	-0.00379679\\
0.906	-0.00388413\\
0.908	-0.00396155\\
0.91	-0.00402903\\
0.912	-0.00408659\\
0.914	-0.00413427\\
0.916	-0.00417213\\
0.918	-0.00420024\\
0.92	-0.00421871\\
0.922	-0.00422765\\
0.924	-0.00422719\\
0.926	-0.00421749\\
0.928	-0.00419871\\
0.93	-0.00417105\\
0.932	-0.0041347\\
0.934	-0.00408989\\
0.936	-0.00403685\\
0.938	-0.00397383\\
0.94	-0.00390307\\
0.942	-0.00382444\\
0.944	-0.00373712\\
0.946	-0.00364312\\
0.948	-0.00354274\\
0.95	-0.00343629\\
0.952	-0.0033241\\
0.954	-0.00320649\\
0.956	-0.00308381\\
0.958	-0.00295638\\
0.96	-0.00282455\\
0.962	-0.00268867\\
0.964	-0.0025491\\
0.966	-0.00240619\\
0.968	-0.00226029\\
0.97	-0.00211175\\
0.972	-0.00196095\\
0.974	-0.00180822\\
0.976	-0.00165393\\
0.978	-0.00158436\\
0.98	-0.00157802\\
0.982	-0.00156867\\
0.984	-0.0015564\\
0.986	-0.00154127\\
0.988	-0.00152336\\
0.99	-0.00150277\\
0.992	-0.00147958\\
0.994	-0.00145388\\
0.996	-0.00142578\\
0.998	-0.00139536\\
1	-0.00136273\\
1.002	-0.00132799\\
1.004	-0.00129126\\
1.006	-0.00125263\\
1.008	-0.00121222\\
1.01	-0.00117014\\
1.012	-0.00112651\\
1.014	-0.00108048\\
1.016	-0.00108226\\
1.018	-0.00120095\\
1.02	-0.00131501\\
1.022	-0.00142427\\
1.024	-0.0015286\\
1.026	-0.00162784\\
1.028	-0.00172187\\
1.03	-0.00181058\\
1.032	-0.00189386\\
1.034	-0.00197164\\
1.036	-0.00204382\\
1.038	-0.00211037\\
1.04	-0.00217121\\
1.042	-0.00222633\\
1.044	-0.00227569\\
1.046	-0.00231929\\
1.048	-0.00235713\\
1.05	-0.00238922\\
1.052	-0.00241559\\
1.054	-0.00243627\\
1.056	-0.00245132\\
1.058	-0.0024608\\
1.06	-0.00246478\\
1.062	-0.00246334\\
1.064	-0.00245657\\
1.066	-0.00244457\\
1.068	-0.00242746\\
1.07	-0.00240535\\
1.072	-0.00237838\\
1.074	-0.00234668\\
1.076	-0.0023104\\
1.078	-0.00226969\\
1.08	-0.00222471\\
1.082	-0.00217562\\
1.084	-0.00212261\\
1.086	-0.00206583\\
1.088	-0.00200549\\
1.09	-0.00194176\\
1.092	-0.00187484\\
1.094	-0.00180493\\
1.096	-0.00173221\\
1.098	-0.0016569\\
1.1	-0.0015792\\
1.102	-0.00149931\\
1.104	-0.00141743\\
1.106	-0.00133379\\
1.108	-0.00124858\\
1.11	-0.00116202\\
1.112	-0.00107431\\
1.114	-0.000985654\\
1.116	-0.000896265\\
1.118	-0.000806344\\
1.12	-0.000781182\\
1.122	-0.000780584\\
1.124	-0.000778379\\
1.126	-0.000774607\\
1.128	-0.000769306\\
1.13	-0.000762519\\
1.132	-0.00075429\\
1.134	-0.000744666\\
1.136	-0.000733695\\
1.138	-0.000721427\\
1.14	-0.000707914\\
1.142	-0.000693208\\
1.144	-0.000677364\\
1.146	-0.000660437\\
1.148	-0.000642484\\
1.15	-0.000623562\\
1.152	-0.00060373\\
1.154	-0.000626164\\
1.156	-0.000693355\\
1.158	-0.000757888\\
1.16	-0.000819669\\
1.162	-0.000878612\\
1.164	-0.000934637\\
1.166	-0.000987673\\
1.168	-0.00103766\\
1.17	-0.00108453\\
1.172	-0.00112825\\
1.174	-0.00116877\\
1.176	-0.00120603\\
1.178	-0.00123994\\
1.18	-0.00127058\\
1.182	-0.00129794\\
1.184	-0.00132202\\
1.186	-0.00134283\\
1.188	-0.00136037\\
1.19	-0.00137466\\
1.192	-0.00138572\\
1.194	-0.00139358\\
1.196	-0.00139828\\
1.198	-0.00139986\\
1.2	-0.00139837\\
1.202	-0.00139386\\
1.204	-0.00138639\\
1.206	-0.00137603\\
1.208	-0.00136286\\
1.21	-0.00134693\\
1.212	-0.00132834\\
1.214	-0.00130717\\
1.216	-0.0012835\\
1.218	-0.00125743\\
1.22	-0.00122906\\
1.222	-0.00119848\\
1.224	-0.0011658\\
1.226	-0.00113113\\
1.228	-0.00109456\\
1.23	-0.00105622\\
1.232	-0.00101616\\
1.234	-0.000974516\\
1.236	-0.000931432\\
1.238	-0.000887028\\
1.24	-0.000841422\\
1.242	-0.000794731\\
1.244	-0.000747077\\
1.246	-0.000698577\\
1.248	-0.000649351\\
1.25	-0.000599518\\
1.252	-0.000549195\\
1.254	-0.0004985\\
1.256	-0.000447549\\
1.258	-0.000396457\\
1.26	-0.000362604\\
1.262	-0.000363484\\
1.264	-0.000363561\\
1.266	-0.000362852\\
1.268	-0.000361376\\
1.27	-0.000359151\\
1.272	-0.000356199\\
1.274	-0.000352541\\
1.276	-0.000348199\\
1.278	-0.000343197\\
1.28	-0.000337559\\
1.282	-0.00033131\\
1.284	-0.000324478\\
1.286	-0.000317087\\
1.288	-0.000309166\\
1.29	-0.000331239\\
1.292	-0.00036985\\
1.294	-0.000406985\\
1.296	-0.00044259\\
1.298	-0.000476613\\
1.3	-0.000509007\\
1.302	-0.000539731\\
1.304	-0.000568745\\
1.306	-0.000596016\\
1.308	-0.000621515\\
1.31	-0.000645217\\
1.312	-0.0006671\\
1.314	-0.000687149\\
1.316	-0.00070535\\
1.318	-0.000721696\\
1.32	-0.000736182\\
1.322	-0.00074881\\
1.324	-0.000759582\\
1.326	-0.000768507\\
1.328	-0.000775598\\
1.33	-0.000780869\\
1.332	-0.00078434\\
1.334	-0.000786034\\
1.336	-0.000785978\\
1.338	-0.000784201\\
1.34	-0.000780737\\
1.342	-0.000775621\\
1.344	-0.000768892\\
1.346	-0.000760593\\
1.348	-0.000750768\\
1.35	-0.000739464\\
1.352	-0.000726731\\
1.354	-0.000712621\\
1.356	-0.000697188\\
1.358	-0.000680487\\
1.36	-0.000662577\\
1.362	-0.000643516\\
1.364	-0.000623366\\
1.366	-0.000602189\\
1.368	-0.000580049\\
1.37	-0.000557009\\
1.372	-0.000533136\\
1.374	-0.000508494\\
1.376	-0.000483152\\
1.378	-0.000457176\\
1.38	-0.000430632\\
1.382	-0.000403589\\
1.384	-0.000376114\\
1.386	-0.000348274\\
1.388	-0.000320135\\
1.39	-0.000291765\\
1.392	-0.000263228\\
1.394	-0.00023459\\
1.396	-0.000205914\\
1.398	-0.000177265\\
1.4	-0.00016778\\
1.402	-0.000167991\\
1.404	-0.000167813\\
1.406	-0.000167255\\
1.408	-0.000166327\\
1.41	-0.000165037\\
1.412	-0.000163397\\
1.414	-0.000161418\\
1.416	-0.000159111\\
1.418	-0.000156489\\
1.42	-0.000153563\\
1.422	-0.000150348\\
1.424	-0.000157649\\
1.426	-0.000180251\\
1.428	-0.000202068\\
1.43	-0.000223066\\
1.432	-0.000243214\\
1.434	-0.000262483\\
1.436	-0.000280844\\
1.438	-0.000298273\\
1.44	-0.000314749\\
1.442	-0.000330251\\
1.444	-0.000344763\\
1.446	-0.000358269\\
1.448	-0.000370758\\
1.45	-0.000382218\\
1.452	-0.000392643\\
1.454	-0.000402027\\
1.456	-0.000410368\\
1.458	-0.000417665\\
1.46	-0.00042392\\
1.462	-0.000429136\\
1.464	-0.000433321\\
1.466	-0.000436481\\
1.468	-0.000438629\\
1.47	-0.000439775\\
1.472	-0.000439935\\
1.474	-0.000439125\\
1.476	-0.000437363\\
1.478	-0.000434669\\
1.48	-0.000431065\\
1.482	-0.000426573\\
1.484	-0.000421219\\
1.486	-0.000415029\\
1.488	-0.000408031\\
1.49	-0.000400253\\
1.492	-0.000391725\\
1.494	-0.00038248\\
1.496	-0.000372549\\
1.498	-0.000361965\\
1.5	-0.000350763\\
1.502	-0.000338977\\
1.504	-0.000326644\\
1.506	-0.000313799\\
1.508	-0.00030048\\
1.51	-0.000286722\\
1.512	-0.000272564\\
1.514	-0.000258044\\
1.516	-0.000243199\\
1.518	-0.000228067\\
1.52	-0.000212687\\
1.522	-0.000197096\\
1.524	-0.000181331\\
1.526	-0.00016543\\
1.528	-0.000149431\\
1.53	-0.00013337\\
1.532	-0.000117283\\
1.534	-0.000101206\\
1.536	-8.51735e-05\\
1.538	-8.02876e-05\\
1.54	-8.00591e-05\\
1.542	-7.96422e-05\\
1.544	-7.90418e-05\\
1.546	-7.8263e-05\\
1.548	-7.7311e-05\\
1.55	-7.61918e-05\\
1.552	-7.49111e-05\\
1.554	-7.34753e-05\\
1.556	-7.18907e-05\\
1.558	-7.56007e-05\\
1.56	-8.87181e-05\\
1.562	-0.000101417\\
1.564	-0.000113676\\
1.566	-0.000125477\\
1.568	-0.000136801\\
1.57	-0.000147631\\
1.572	-0.000157953\\
1.574	-0.000167751\\
1.576	-0.000177014\\
1.578	-0.000185731\\
1.58	-0.00019389\\
1.582	-0.000201485\\
1.584	-0.000208507\\
1.586	-0.000214951\\
1.588	-0.000220813\\
1.59	-0.000226089\\
1.592	-0.000230778\\
1.594	-0.00023488\\
1.596	-0.000238395\\
1.598	-0.000241326\\
1.6	-0.000243676\\
1.602	-0.00024545\\
1.604	-0.000246652\\
1.606	-0.000247292\\
1.608	-0.000247376\\
1.61	-0.000246913\\
1.612	-0.000245915\\
1.614	-0.000244391\\
1.616	-0.000242355\\
1.618	-0.000239819\\
1.62	-0.000236797\\
1.622	-0.000233304\\
1.624	-0.000229356\\
1.626	-0.000224969\\
1.628	-0.00022016\\
1.63	-0.000214946\\
1.632	-0.000209346\\
1.634	-0.000203379\\
1.636	-0.000197064\\
1.638	-0.00019042\\
1.64	-0.000183467\\
1.642	-0.000176227\\
1.644	-0.000168719\\
1.646	-0.000160965\\
1.648	-0.000152986\\
1.65	-0.000144802\\
1.652	-0.000136437\\
1.654	-0.000127909\\
1.656	-0.000119243\\
1.658	-0.000110457\\
1.66	-0.000101575\\
1.662	-9.26162e-05\\
1.664	-8.36024e-05\\
1.666	-7.45541e-05\\
1.668	-6.54918e-05\\
1.67	-5.64354e-05\\
1.672	-4.74048e-05\\
1.674	-4.20534e-05\\
1.676	-4.1684e-05\\
1.678	-4.12186e-05\\
1.68	-4.066e-05\\
1.682	-4.0011e-05\\
1.684	-3.92749e-05\\
1.686	-3.8455e-05\\
1.688	-3.75546e-05\\
1.69	-3.65774e-05\\
1.692	-3.61978e-05\\
1.694	-4.37938e-05\\
1.696	-5.11659e-05\\
1.698	-5.8302e-05\\
1.7	-6.51902e-05\\
1.702	-7.18197e-05\\
1.704	-7.81803e-05\\
1.706	-8.42626e-05\\
1.708	-9.00579e-05\\
1.71	-9.55583e-05\\
1.712	-0.000100757\\
1.714	-0.000105647\\
1.716	-0.000110223\\
1.718	-0.000114481\\
1.72	-0.000118416\\
1.722	-0.000122026\\
1.724	-0.000125307\\
1.726	-0.000128258\\
1.728	-0.000130878\\
1.73	-0.000133168\\
1.732	-0.000135126\\
1.734	-0.000136755\\
1.736	-0.000138057\\
1.738	-0.000139034\\
1.74	-0.000139689\\
1.742	-0.000140026\\
1.744	-0.000140051\\
1.746	-0.000139768\\
1.748	-0.000139182\\
1.75	-0.0001383\\
1.752	-0.00013713\\
1.754	-0.000135678\\
1.756	-0.000133951\\
1.758	-0.00013196\\
1.76	-0.000129712\\
1.762	-0.000127216\\
1.764	-0.000124482\\
1.766	-0.00012152\\
1.768	-0.00011834\\
1.77	-0.000114953\\
1.772	-0.00011137\\
1.774	-0.000107601\\
1.776	-0.000103658\\
1.778	-9.95527e-05\\
1.78	-9.52965e-05\\
1.782	-9.09012e-05\\
1.784	-8.63788e-05\\
1.786	-8.17411e-05\\
1.788	-7.70003e-05\\
1.79	-7.21684e-05\\
1.792	-6.72575e-05\\
1.794	-6.22795e-05\\
1.796	-5.72465e-05\\
1.798	-5.21704e-05\\
1.8	-4.7063e-05\\
1.802	-4.19358e-05\\
1.804	-3.68006e-05\\
1.806	-3.16685e-05\\
1.808	-2.65508e-05\\
1.81	-2.41714e-05\\
1.812	-2.37905e-05\\
1.814	-2.33585e-05\\
1.816	-2.28772e-05\\
1.818	-2.23487e-05\\
1.82	-2.17748e-05\\
1.822	-2.11576e-05\\
1.824	-2.04993e-05\\
1.826	-1.9802e-05\\
1.828	-2.1091e-05\\
1.83	-2.53967e-05\\
1.832	-2.9576e-05\\
1.834	-3.3622e-05\\
1.836	-3.7528e-05\\
1.838	-4.12878e-05\\
1.84	-4.48956e-05\\
1.842	-4.8346e-05\\
1.844	-5.16342e-05\\
1.846	-5.47555e-05\\
1.848	-5.7706e-05\\
1.85	-6.04819e-05\\
1.852	-6.30802e-05\\
1.854	-6.5498e-05\\
1.856	-6.7733e-05\\
1.858	-6.97834e-05\\
1.86	-7.16478e-05\\
1.862	-7.33249e-05\\
1.864	-7.48144e-05\\
1.866	-7.6116e-05\\
1.868	-7.72299e-05\\
1.87	-7.81568e-05\\
1.872	-7.88977e-05\\
1.874	-7.94541e-05\\
1.876	-7.98276e-05\\
1.878	-8.00205e-05\\
1.88	-8.00354e-05\\
1.882	-7.9875e-05\\
1.884	-7.95426e-05\\
1.886	-7.90417e-05\\
1.888	-7.83761e-05\\
1.89	-7.755e-05\\
1.892	-7.65678e-05\\
1.894	-7.54341e-05\\
1.896	-7.41539e-05\\
1.898	-7.27323e-05\\
1.9	-7.11748e-05\\
1.902	-6.94869e-05\\
1.904	-6.76744e-05\\
1.906	-6.57433e-05\\
1.908	-6.36997e-05\\
1.91	-6.15499e-05\\
1.912	-5.93004e-05\\
1.914	-5.69575e-05\\
1.916	-5.45279e-05\\
1.918	-5.20184e-05\\
1.92	-4.94357e-05\\
1.922	-4.67865e-05\\
1.924	-4.40778e-05\\
1.926	-4.13163e-05\\
1.928	-3.85091e-05\\
1.93	-3.56628e-05\\
1.932	-3.27844e-05\\
1.934	-2.98806e-05\\
1.936	-2.69581e-05\\
1.938	-2.40237e-05\\
1.94	-2.10839e-05\\
1.942	-1.81452e-05\\
1.944	-1.58248e-05\\
1.946	-1.63372e-05\\
1.948	-1.67922e-05\\
1.95	-1.71888e-05\\
1.952	-1.75264e-05\\
1.954	-1.78046e-05\\
1.956	-1.8023e-05\\
1.958	-1.81817e-05\\
1.96	-1.82807e-05\\
1.962	-1.83205e-05\\
1.964	-1.83015e-05\\
1.966	-1.82245e-05\\
1.968	-1.80903e-05\\
1.97	-1.94048e-05\\
1.972	-2.16493e-05\\
1.974	-2.38104e-05\\
1.976	-2.58847e-05\\
1.978	-2.78693e-05\\
1.98	-2.97611e-05\\
1.982	-3.15576e-05\\
1.984	-3.32563e-05\\
1.986	-3.48552e-05\\
1.988	-3.63524e-05\\
1.99	-3.77462e-05\\
1.992	-3.90352e-05\\
1.994	-4.02184e-05\\
1.996	-4.12949e-05\\
1.998	-4.22639e-05\\
2	-4.31252e-05\\
};
\end{axis}

\begin{axis}[%
width={\figscale*1.528in},
height={\figscale*0.46in},
at={(\figscale*2.841in,\figscale*0.56in)},
scale only axis,
separate axis lines,
every outer x axis line/.append style={black},
every x tick label/.append style={font=\figticfont},
every x tick/.append style={black},
xmin=-0.025,
xmax=2,
xlabel={Time since step in s},
every outer y axis line/.append style={black},
every y tick label/.append style={font=\figticfont},
every y tick/.append style={black},
ymin=-0.0359538,
ymax=0.00159155,
ytick={-0.03,0},
axis background/.style={fill=white},
xmajorgrids,
ymajorgrids,
grid style={white!55!black, opacity=0.3},
ylabel style={rotate=-90, at={(axis description cs:-0.145,0.5)}, anchor=east}
]
\addplot [color=mycolor1, line width=1.4pt, forget plot]
  table[row sep=crcr]{%
-0.024	0\\
-0.022	0\\
-0.02	0\\
-0.018	0\\
-0.016	0\\
-0.014	0\\
-0.012	0\\
-0.01	0\\
-0.008	0\\
-0.006	0\\
-0.004	0\\
-0.002	0\\
0	0\\
0.002	0\\
0.004	0\\
0.006	0\\
0.008	0\\
0.01	0\\
0.012	0\\
0.014	0\\
0.016	0\\
0.018	0\\
0.02	0\\
0.022	0\\
0.024	0\\
0.026	0\\
0.028	0\\
0.03	0\\
0.032	0\\
0.034	0\\
0.036	0\\
0.038	0\\
0.04	0\\
0.042	0\\
0.044	0\\
0.046	0\\
0.048	0\\
0.05	0\\
0.052	0\\
0.054	0\\
0.056	0\\
0.058	0\\
0.06	0\\
0.062	0\\
0.064	0\\
0.066	0\\
0.068	0\\
0.07	0\\
0.072	0\\
0.074	0\\
0.076	0\\
0.078	0\\
0.08	0\\
0.082	0\\
0.084	0\\
0.086	0\\
0.088	0\\
0.09	0\\
0.092	0\\
0.094	0\\
0.096	0\\
0.098	0\\
0.1	0\\
0.102	0\\
0.104	0\\
0.106	0\\
0.108	0\\
0.11	0\\
0.112	0\\
0.114	0\\
0.116	0\\
0.118	0\\
0.12	0\\
0.122	0\\
0.124	0\\
0.126	0\\
0.128	0\\
0.13	0\\
0.132	0\\
0.134	0\\
0.136	0\\
0.138	0\\
0.14	0\\
0.142	0\\
0.144	0\\
0.146	0\\
0.148	0\\
0.15	0\\
0.152	0\\
0.154	0\\
0.156	0\\
0.158	0\\
0.16	0\\
0.162	0\\
0.164	0\\
0.166	0\\
0.168	0\\
0.17	0\\
0.172	0\\
0.174	0\\
0.176	0\\
0.178	0\\
0.18	0\\
0.182	0\\
0.184	0\\
0.186	0\\
0.188	0\\
0.19	0\\
0.192	0\\
0.194	0\\
0.196	0\\
0.198	0\\
0.2	0\\
0.202	0\\
0.204	0\\
0.206	0\\
0.208	0\\
0.21	0\\
0.212	0\\
0.214	0\\
0.216	0\\
0.218	0\\
0.22	0\\
0.222	0\\
0.224	0\\
0.226	0\\
0.228	0\\
0.23	0\\
0.232	0\\
0.234	0\\
0.236	0\\
0.238	0\\
0.24	0\\
0.242	0\\
0.244	0\\
0.246	0\\
0.248	0\\
0.25	0\\
0.252	0\\
0.254	0\\
0.256	0\\
0.258	0\\
0.26	0\\
0.262	0\\
0.264	0\\
0.266	0\\
0.268	0\\
0.27	0\\
0.272	0\\
0.274	0\\
0.276	0\\
0.278	0\\
0.28	0\\
0.282	0\\
0.284	0\\
0.286	0\\
0.288	0\\
0.29	0\\
0.292	0\\
0.294	0\\
0.296	0\\
0.298	0\\
0.3	0\\
0.302	0\\
0.304	0\\
0.306	0\\
0.308	0\\
0.31	0\\
0.312	0\\
0.314	0\\
0.316	0\\
0.318	0\\
0.32	0\\
0.322	0\\
0.324	0\\
0.326	0\\
0.328	0\\
0.33	0\\
0.332	0\\
0.334	0\\
0.336	0\\
0.338	0\\
0.34	0\\
0.342	0\\
0.344	0\\
0.346	0\\
0.348	0\\
0.35	0\\
0.352	0\\
0.354	0\\
0.356	0\\
0.358	0\\
0.36	0\\
0.362	0\\
0.364	0\\
0.366	0\\
0.368	0\\
0.37	0\\
0.372	0\\
0.374	0\\
0.376	0\\
0.378	0\\
0.38	0\\
0.382	0\\
0.384	0\\
0.386	0\\
0.388	0\\
0.39	0\\
0.392	0\\
0.394	0\\
0.396	0\\
0.398	0\\
0.4	0\\
0.402	0\\
0.404	0\\
0.406	0\\
0.408	0\\
0.41	0\\
0.412	0\\
0.414	0\\
0.416	0\\
0.418	0\\
0.42	0\\
0.422	0\\
0.424	0\\
0.426	0\\
0.428	0\\
0.43	0\\
0.432	0\\
0.434	0\\
0.436	0\\
0.438	0\\
0.44	0\\
0.442	0\\
0.444	0\\
0.446	0\\
0.448	0\\
0.45	0\\
0.452	0\\
0.454	0\\
0.456	0\\
0.458	0\\
0.46	0\\
0.462	0\\
0.464	0\\
0.466	0\\
0.468	0\\
0.47	0\\
0.472	0\\
0.474	0\\
0.476	0\\
0.478	0\\
0.48	0\\
0.482	0\\
0.484	0\\
0.486	0\\
0.488	0\\
0.49	0\\
0.492	0\\
0.494	0\\
0.496	0\\
0.498	0\\
0.5	0\\
0.502	0\\
0.504	0\\
0.506	0\\
0.508	0\\
0.51	0\\
0.512	0\\
0.514	0\\
0.516	0\\
0.518	0\\
0.52	0\\
0.522	0\\
0.524	0\\
0.526	0\\
0.528	0\\
0.53	0\\
0.532	0\\
0.534	0\\
0.536	0\\
0.538	0\\
0.54	0\\
0.542	0\\
0.544	0\\
0.546	0\\
0.548	0\\
0.55	0\\
0.552	0\\
0.554	0\\
0.556	0\\
0.558	0\\
0.56	0\\
0.562	0\\
0.564	0\\
0.566	0\\
0.568	0\\
0.57	0\\
0.572	0\\
0.574	0\\
0.576	0\\
0.578	0\\
0.58	0\\
0.582	0\\
0.584	0\\
0.586	0\\
0.588	0\\
0.59	0\\
0.592	0\\
0.594	0\\
0.596	0\\
0.598	0\\
0.6	0\\
0.602	0\\
0.604	0\\
0.606	0\\
0.608	0\\
0.61	0\\
0.612	0\\
0.614	0\\
0.616	0\\
0.618	0\\
0.62	0\\
0.622	0\\
0.624	0\\
0.626	0\\
0.628	0\\
0.63	0\\
0.632	0\\
0.634	0\\
0.636	0\\
0.638	0\\
0.64	0\\
0.642	0\\
0.644	0\\
0.646	0\\
0.648	0\\
0.65	0\\
0.652	0\\
0.654	0\\
0.656	0\\
0.658	0\\
0.66	0\\
0.662	0\\
0.664	0\\
0.666	0\\
0.668	0\\
0.67	0\\
0.672	0\\
0.674	0\\
0.676	0\\
0.678	0\\
0.68	0\\
0.682	0\\
0.684	0\\
0.686	0\\
0.688	0\\
0.69	0\\
0.692	0\\
0.694	0\\
0.696	0\\
0.698	0\\
0.7	0\\
0.702	0\\
0.704	0\\
0.706	0\\
0.708	0\\
0.71	0\\
0.712	0\\
0.714	0\\
0.716	0\\
0.718	0\\
0.72	0\\
0.722	0\\
0.724	0\\
0.726	0\\
0.728	0\\
0.73	0\\
0.732	0\\
0.734	0\\
0.736	0\\
0.738	0\\
0.74	0\\
0.742	0\\
0.744	0\\
0.746	0\\
0.748	0\\
0.75	0\\
0.752	0\\
0.754	0\\
0.756	0\\
0.758	0\\
0.76	0\\
0.762	0\\
0.764	0\\
0.766	0\\
0.768	0\\
0.77	0\\
0.772	0\\
0.774	0\\
0.776	0\\
0.778	0\\
0.78	0\\
0.782	0\\
0.784	0\\
0.786	0\\
0.788	0\\
0.79	0\\
0.792	0\\
0.794	0\\
0.796	0\\
0.798	0\\
0.8	0\\
0.802	0\\
0.804	0\\
0.806	0\\
0.808	0\\
0.81	0\\
0.812	0\\
0.814	0\\
0.816	0\\
0.818	0\\
0.82	0\\
0.822	0\\
0.824	0\\
0.826	0\\
0.828	0\\
0.83	0\\
0.832	0\\
0.834	0\\
0.836	0\\
0.838	0\\
0.84	0\\
0.842	0\\
0.844	0\\
0.846	0\\
0.848	0\\
0.85	0\\
0.852	0\\
0.854	0\\
0.856	0\\
0.858	0\\
0.86	0\\
0.862	0\\
0.864	0\\
0.866	0\\
0.868	0\\
0.87	0\\
0.872	0\\
0.874	0\\
0.876	0\\
0.878	0\\
0.88	0\\
0.882	0\\
0.884	0\\
0.886	0\\
0.888	0\\
0.89	0\\
0.892	0\\
0.894	0\\
0.896	0\\
0.898	0\\
0.9	0\\
0.902	0\\
0.904	0\\
0.906	0\\
0.908	0\\
0.91	0\\
0.912	0\\
0.914	0\\
0.916	0\\
0.918	0\\
0.92	0\\
0.922	0\\
0.924	0\\
0.926	0\\
0.928	0\\
0.93	0\\
0.932	0\\
0.934	0\\
0.936	0\\
0.938	0\\
0.94	0\\
0.942	0\\
0.944	0\\
0.946	0\\
0.948	0\\
0.95	0\\
0.952	0\\
0.954	0\\
0.956	0\\
0.958	0\\
0.96	0\\
0.962	0\\
0.964	0\\
0.966	0\\
0.968	0\\
0.97	0\\
0.972	0\\
0.974	0\\
0.976	0\\
0.978	0\\
0.98	0\\
0.982	0\\
0.984	0\\
0.986	0\\
0.988	0\\
0.99	0\\
0.992	0\\
0.994	0\\
0.996	0\\
0.998	0\\
1	0\\
1.002	0\\
1.004	0\\
1.006	0\\
1.008	0\\
1.01	0\\
1.012	0\\
1.014	0\\
1.016	0\\
1.018	0\\
1.02	0\\
1.022	0\\
1.024	0\\
1.026	0\\
1.028	0\\
1.03	0\\
1.032	0\\
1.034	0\\
1.036	0\\
1.038	0\\
1.04	0\\
1.042	0\\
1.044	0\\
1.046	0\\
1.048	0\\
1.05	0\\
1.052	0\\
1.054	0\\
1.056	0\\
1.058	0\\
1.06	0\\
1.062	0\\
1.064	0\\
1.066	0\\
1.068	0\\
1.07	0\\
1.072	0\\
1.074	0\\
1.076	0\\
1.078	0\\
1.08	0\\
1.082	0\\
1.084	0\\
1.086	0\\
1.088	0\\
1.09	0\\
1.092	0\\
1.094	0\\
1.096	0\\
1.098	0\\
1.1	0\\
1.102	0\\
1.104	0\\
1.106	0\\
1.108	0\\
1.11	0\\
1.112	0\\
1.114	0\\
1.116	0\\
1.118	0\\
1.12	0\\
1.122	0\\
1.124	0\\
1.126	0\\
1.128	0\\
1.13	0\\
1.132	0\\
1.134	0\\
1.136	0\\
1.138	0\\
1.14	0\\
1.142	0\\
1.144	0\\
1.146	0\\
1.148	0\\
1.15	0\\
1.152	0\\
1.154	0\\
1.156	0\\
1.158	0\\
1.16	0\\
1.162	0\\
1.164	0\\
1.166	0\\
1.168	0\\
1.17	0\\
1.172	0\\
1.174	0\\
1.176	0\\
1.178	0\\
1.18	0\\
1.182	0\\
1.184	0\\
1.186	0\\
1.188	0\\
1.19	0\\
1.192	0\\
1.194	0\\
1.196	0\\
1.198	0\\
1.2	0\\
1.202	0\\
1.204	0\\
1.206	0\\
1.208	0\\
1.21	0\\
1.212	0\\
1.214	0\\
1.216	0\\
1.218	0\\
1.22	0\\
1.222	0\\
1.224	0\\
1.226	0\\
1.228	0\\
1.23	0\\
1.232	0\\
1.234	0\\
1.236	0\\
1.238	0\\
1.24	0\\
1.242	0\\
1.244	0\\
1.246	0\\
1.248	0\\
1.25	0\\
1.252	0\\
1.254	0\\
1.256	0\\
1.258	0\\
1.26	0\\
1.262	0\\
1.264	0\\
1.266	0\\
1.268	0\\
1.27	0\\
1.272	0\\
1.274	0\\
1.276	0\\
1.278	0\\
1.28	0\\
1.282	0\\
1.284	0\\
1.286	0\\
1.288	0\\
1.29	0\\
1.292	0\\
1.294	0\\
1.296	0\\
1.298	0\\
1.3	0\\
1.302	0\\
1.304	0\\
1.306	0\\
1.308	0\\
1.31	0\\
1.312	0\\
1.314	0\\
1.316	0\\
1.318	0\\
1.32	0\\
1.322	0\\
1.324	0\\
1.326	0\\
1.328	0\\
1.33	0\\
1.332	0\\
1.334	0\\
1.336	0\\
1.338	0\\
1.34	0\\
1.342	0\\
1.344	0\\
1.346	0\\
1.348	0\\
1.35	0\\
1.352	0\\
1.354	0\\
1.356	0\\
1.358	0\\
1.36	0\\
1.362	0\\
1.364	0\\
1.366	0\\
1.368	0\\
1.37	0\\
1.372	0\\
1.374	0\\
1.376	0\\
1.378	0\\
1.38	0\\
1.382	0\\
1.384	0\\
1.386	0\\
1.388	0\\
1.39	0\\
1.392	0\\
1.394	0\\
1.396	0\\
1.398	0\\
1.4	0\\
1.402	0\\
1.404	0\\
1.406	0\\
1.408	0\\
1.41	0\\
1.412	0\\
1.414	0\\
1.416	0\\
1.418	0\\
1.42	0\\
1.422	0\\
1.424	0\\
1.426	0\\
1.428	0\\
1.43	0\\
1.432	0\\
1.434	0\\
1.436	0\\
1.438	0\\
1.44	0\\
1.442	0\\
1.444	0\\
1.446	0\\
1.448	0\\
1.45	0\\
1.452	0\\
1.454	0\\
1.456	0\\
1.458	0\\
1.46	0\\
1.462	0\\
1.464	0\\
1.466	0\\
1.468	0\\
1.47	0\\
1.472	0\\
1.474	0\\
1.476	0\\
1.478	0\\
1.48	0\\
1.482	0\\
1.484	0\\
1.486	0\\
1.488	0\\
1.49	0\\
1.492	0\\
1.494	0\\
1.496	0\\
1.498	0\\
1.5	0\\
1.502	0\\
1.504	0\\
1.506	0\\
1.508	0\\
1.51	0\\
1.512	0\\
1.514	0\\
1.516	0\\
1.518	0\\
1.52	0\\
1.522	0\\
1.524	0\\
1.526	0\\
1.528	0\\
1.53	0\\
1.532	0\\
1.534	0\\
1.536	0\\
1.538	0\\
1.54	0\\
1.542	0\\
1.544	0\\
1.546	0\\
1.548	0\\
1.55	0\\
1.552	0\\
1.554	0\\
1.556	0\\
1.558	0\\
1.56	0\\
1.562	0\\
1.564	0\\
1.566	0\\
1.568	0\\
1.57	0\\
1.572	0\\
1.574	0\\
1.576	0\\
1.578	0\\
1.58	0\\
1.582	0\\
1.584	0\\
1.586	0\\
1.588	0\\
1.59	0\\
1.592	0\\
1.594	0\\
1.596	0\\
1.598	0\\
1.6	0\\
1.602	0\\
1.604	0\\
1.606	0\\
1.608	0\\
1.61	0\\
1.612	0\\
1.614	0\\
1.616	0\\
1.618	0\\
1.62	0\\
1.622	0\\
1.624	0\\
1.626	0\\
1.628	0\\
1.63	0\\
1.632	0\\
1.634	0\\
1.636	0\\
1.638	0\\
1.64	0\\
1.642	0\\
1.644	0\\
1.646	0\\
1.648	0\\
1.65	0\\
1.652	0\\
1.654	0\\
1.656	0\\
1.658	0\\
1.66	0\\
1.662	0\\
1.664	0\\
1.666	0\\
1.668	0\\
1.67	0\\
1.672	0\\
1.674	0\\
1.676	0\\
1.678	0\\
1.68	0\\
1.682	0\\
1.684	0\\
1.686	0\\
1.688	0\\
1.69	0\\
1.692	0\\
1.694	0\\
1.696	0\\
1.698	0\\
1.7	0\\
1.702	0\\
1.704	0\\
1.706	0\\
1.708	0\\
1.71	0\\
1.712	0\\
1.714	0\\
1.716	0\\
1.718	0\\
1.72	0\\
1.722	0\\
1.724	0\\
1.726	0\\
1.728	0\\
1.73	0\\
1.732	0\\
1.734	0\\
1.736	0\\
1.738	0\\
1.74	0\\
1.742	0\\
1.744	0\\
1.746	0\\
1.748	0\\
1.75	0\\
1.752	0\\
1.754	0\\
1.756	0\\
1.758	0\\
1.76	0\\
1.762	0\\
1.764	0\\
1.766	0\\
1.768	0\\
1.77	0\\
1.772	0\\
1.774	0\\
1.776	0\\
1.778	0\\
1.78	0\\
1.782	0\\
1.784	0\\
1.786	0\\
1.788	0\\
1.79	0\\
1.792	0\\
1.794	0\\
1.796	0\\
1.798	0\\
1.8	0\\
1.802	0\\
1.804	0\\
1.806	0\\
1.808	0\\
1.81	0\\
1.812	0\\
1.814	0\\
1.816	0\\
1.818	0\\
1.82	0\\
1.822	0\\
1.824	0\\
1.826	0\\
1.828	0\\
1.83	0\\
1.832	0\\
1.834	0\\
1.836	0\\
1.838	0\\
1.84	0\\
1.842	0\\
1.844	0\\
1.846	0\\
1.848	0\\
1.85	0\\
1.852	0\\
1.854	0\\
1.856	0\\
1.858	0\\
1.86	0\\
1.862	0\\
1.864	0\\
1.866	0\\
1.868	0\\
1.87	0\\
1.872	0\\
1.874	0\\
1.876	0\\
1.878	0\\
1.88	0\\
1.882	0\\
1.884	0\\
1.886	0\\
1.888	0\\
1.89	0\\
1.892	0\\
1.894	0\\
1.896	0\\
1.898	0\\
1.9	0\\
1.902	0\\
1.904	0\\
1.906	0\\
1.908	0\\
1.91	0\\
1.912	0\\
1.914	0\\
1.916	0\\
1.918	0\\
1.92	0\\
1.922	0\\
1.924	0\\
1.926	0\\
1.928	0\\
1.93	0\\
1.932	0\\
1.934	0\\
1.936	0\\
1.938	0\\
1.94	0\\
1.942	0\\
1.944	0\\
1.946	0\\
1.948	0\\
1.95	0\\
1.952	0\\
1.954	0\\
1.956	0\\
1.958	0\\
1.96	0\\
1.962	0\\
1.964	0\\
1.966	0\\
1.968	0\\
1.97	0\\
1.972	0\\
1.974	0\\
1.976	0\\
1.978	0\\
1.98	0\\
1.982	0\\
1.984	0\\
1.986	0\\
1.988	0\\
1.99	0\\
1.992	0\\
1.994	0\\
1.996	0\\
1.998	0\\
2	0\\
};
\addplot [color=mycolor2, line width=1.0pt, forget plot]
  table[row sep=crcr]{%
-0.024	0\\
-0.022	0\\
-0.02	0\\
-0.018	0\\
-0.016	0\\
-0.014	0\\
-0.012	0\\
-0.01	0\\
-0.008	0\\
-0.006	0\\
-0.004	0\\
-0.002	0\\
0	0\\
0.002	1.18692e-10\\
0.004	-5.30024e-10\\
0.006	-5.69843e-08\\
0.008	-1.76852e-07\\
0.01	-4.17683e-07\\
0.012	-8.31334e-07\\
0.014	-1.47122e-06\\
0.016	-2.38958e-06\\
0.018	-2.49501e-06\\
0.02	-3.66426e-06\\
0.022	-4.62084e-06\\
0.024	-0.000968861\\
0.026	-0.00329765\\
0.028	-0.00519444\\
0.03	-0.00730522\\
0.032	-0.00908246\\
0.034	-0.00985143\\
0.036	-0.0105968\\
0.038	-0.0113117\\
0.04	-0.0119896\\
0.042	-0.0126245\\
0.044	-0.0132107\\
0.046	-0.0137427\\
0.048	-0.0142158\\
0.05	-0.0146255\\
0.052	-0.0149678\\
0.054	-0.0152391\\
0.056	-0.0154364\\
0.058	-0.0155572\\
0.06	-0.0155994\\
0.062	-0.0155613\\
0.064	-0.0154419\\
0.066	-0.0152405\\
0.068	-0.014957\\
0.07	-0.0145916\\
0.072	-0.0141452\\
0.074	-0.0136189\\
0.076	-0.0130144\\
0.078	-0.0123338\\
0.08	-0.0115794\\
0.082	-0.0107543\\
0.084	-0.00986147\\
0.086	-0.0089046\\
0.088	-0.00788753\\
0.09	-0.00681444\\
0.092	-0.00568977\\
0.094	-0.00451823\\
0.096	-0.00330474\\
0.098	-0.00205443\\
0.1	-0.000772589\\
0.102	0.000535335\\
0.104	0.00163667\\
0.106	0.000908885\\
0.108	0.000210661\\
0.11	-0.000456487\\
0.112	-0.00109122\\
0.114	-0.00118987\\
0.116	-0.00111826\\
0.118	-0.00103995\\
0.12	-0.000955031\\
0.122	-0.00128832\\
0.124	-0.00229293\\
0.126	-0.00305017\\
0.128	-0.00354901\\
0.13	-0.00404598\\
0.132	-0.00453936\\
0.134	-0.00502743\\
0.136	-0.00550849\\
0.138	-0.00598083\\
0.14	-0.00644278\\
0.142	-0.00739914\\
0.144	-0.00962672\\
0.146	-0.0117647\\
0.148	-0.0138074\\
0.15	-0.0157497\\
0.152	-0.0175869\\
0.154	-0.0193147\\
0.156	-0.0209293\\
0.158	-0.0224273\\
0.16	-0.0238058\\
0.162	-0.0250624\\
0.164	-0.0257011\\
0.166	-0.0253683\\
0.168	-0.0249495\\
0.17	-0.0244463\\
0.172	-0.0238605\\
0.174	-0.0231942\\
0.176	-0.0224498\\
0.178	-0.0216301\\
0.18	-0.020738\\
0.182	-0.0197769\\
0.184	-0.0187501\\
0.186	-0.0176615\\
0.188	-0.016515\\
0.19	-0.0153148\\
0.192	-0.0140652\\
0.194	-0.0127708\\
0.196	-0.0114362\\
0.198	-0.0100661\\
0.2	-0.00866564\\
0.202	-0.00723967\\
0.204	-0.00579331\\
0.206	-0.00433167\\
0.208	-0.00285993\\
0.21	-0.00144583\\
0.212	-0.0016804\\
0.214	-0.00190027\\
0.216	-0.00210472\\
0.218	-0.00229311\\
0.22	-0.00246487\\
0.222	-0.00261952\\
0.224	-0.00275666\\
0.226	-0.0028566\\
0.228	-0.00271963\\
0.23	-0.00257279\\
0.232	-0.00241644\\
0.234	-0.00225099\\
0.236	-0.00207687\\
0.238	-0.00185821\\
0.24	-0.00253329\\
0.242	-0.00397527\\
0.244	-0.00537917\\
0.246	-0.00674039\\
0.248	-0.00805455\\
0.25	-0.00931752\\
0.252	-0.0105254\\
0.254	-0.0116746\\
0.256	-0.0127618\\
0.258	-0.0137839\\
0.26	-0.0147381\\
0.262	-0.015622\\
0.264	-0.0164334\\
0.266	-0.0171704\\
0.268	-0.0178314\\
0.27	-0.0184152\\
0.272	-0.0189208\\
0.274	-0.0193475\\
0.276	-0.019695\\
0.278	-0.0199632\\
0.28	-0.0201523\\
0.282	-0.0202629\\
0.284	-0.0202957\\
0.286	-0.019768\\
0.288	-0.0189684\\
0.29	-0.0181118\\
0.292	-0.0172015\\
0.294	-0.0162412\\
0.296	-0.0152344\\
0.298	-0.0141848\\
0.3	-0.0130965\\
0.302	-0.0119735\\
0.304	-0.0108197\\
0.306	-0.00963951\\
0.308	-0.008437\\
0.31	-0.00721648\\
0.312	-0.00598225\\
0.314	-0.00473862\\
0.316	-0.00348989\\
0.318	-0.00331771\\
0.32	-0.00340356\\
0.322	-0.00347565\\
0.324	-0.00353373\\
0.326	-0.00357756\\
0.328	-0.003607\\
0.33	-0.00362196\\
0.332	-0.00362238\\
0.334	-0.00360829\\
0.336	-0.00357975\\
0.338	-0.0035369\\
0.34	-0.00347993\\
0.342	-0.00340906\\
0.344	-0.00324243\\
0.346	-0.00304949\\
0.348	-0.00407807\\
0.35	-0.00508737\\
0.352	-0.0059553\\
0.354	-0.00678799\\
0.356	-0.00758284\\
0.358	-0.00833741\\
0.36	-0.00904946\\
0.362	-0.00971692\\
0.364	-0.0103379\\
0.366	-0.0109109\\
0.368	-0.0114342\\
0.37	-0.0119067\\
0.372	-0.0123274\\
0.374	-0.0126953\\
0.376	-0.0130099\\
0.378	-0.0132707\\
0.38	-0.0134775\\
0.382	-0.0136303\\
0.384	-0.0137292\\
0.386	-0.0137556\\
0.388	-0.0136878\\
0.39	-0.0135682\\
0.392	-0.013398\\
0.394	-0.0131785\\
0.396	-0.0129113\\
0.398	-0.012598\\
0.4	-0.0122405\\
0.402	-0.0118408\\
0.404	-0.011401\\
0.406	-0.0109232\\
0.408	-0.0104099\\
0.41	-0.00972129\\
0.412	-0.00883836\\
0.414	-0.00793557\\
0.416	-0.0075888\\
0.418	-0.00725763\\
0.42	-0.00691405\\
0.422	-0.00655904\\
0.424	-0.00619355\\
0.426	-0.00581857\\
0.428	-0.0054351\\
0.43	-0.00504414\\
0.432	-0.00464671\\
0.434	-0.00424382\\
0.436	-0.00407501\\
0.438	-0.00406039\\
0.44	-0.00403213\\
0.442	-0.00399034\\
0.444	-0.00393519\\
0.446	-0.00386688\\
0.448	-0.00378565\\
0.45	-0.0036918\\
0.452	-0.00358567\\
0.454	-0.00380527\\
0.456	-0.00449925\\
0.458	-0.00516532\\
0.46	-0.00580145\\
0.462	-0.00640574\\
0.464	-0.00691721\\
0.466	-0.00736838\\
0.468	-0.00778635\\
0.47	-0.00816999\\
0.472	-0.00851828\\
0.474	-0.00883037\\
0.476	-0.00910555\\
0.478	-0.00934324\\
0.48	-0.00954303\\
0.482	-0.00970465\\
0.484	-0.00982795\\
0.486	-0.00991295\\
0.488	-0.00995981\\
0.49	-0.00994136\\
0.492	-0.00986016\\
0.494	-0.00974285\\
0.496	-0.00959032\\
0.498	-0.00940356\\
0.5	-0.00918369\\
0.502	-0.00893195\\
0.504	-0.00864966\\
0.506	-0.00833825\\
0.508	-0.00799924\\
0.51	-0.00763423\\
0.512	-0.00724489\\
0.514	-0.00683298\\
0.516	-0.00640029\\
0.518	-0.0059487\\
0.52	-0.00548011\\
0.522	-0.00554268\\
0.524	-0.00563945\\
0.526	-0.0057203\\
0.528	-0.00578525\\
0.53	-0.00583433\\
0.532	-0.00586764\\
0.534	-0.00587821\\
0.536	-0.00572809\\
0.538	-0.00556648\\
0.54	-0.00539393\\
0.542	-0.00521101\\
0.544	-0.00501829\\
0.546	-0.00481638\\
0.548	-0.0046059\\
0.55	-0.00438746\\
0.552	-0.00416171\\
0.554	-0.00392928\\
0.556	-0.00377036\\
0.558	-0.00369224\\
0.56	-0.00360268\\
0.562	-0.00381027\\
0.564	-0.00423147\\
0.566	-0.00463081\\
0.568	-0.00500713\\
0.57	-0.00535933\\
0.572	-0.00568646\\
0.574	-0.00598765\\
0.576	-0.00626215\\
0.578	-0.0065093\\
0.58	-0.00672858\\
0.582	-0.00691957\\
0.584	-0.00708195\\
0.586	-0.00719604\\
0.588	-0.00728005\\
0.59	-0.00733576\\
0.592	-0.00736333\\
0.594	-0.00736304\\
0.596	-0.00733524\\
0.598	-0.0072671\\
0.6	-0.00714724\\
0.602	-0.00700264\\
0.604	-0.00683413\\
0.606	-0.00664261\\
0.608	-0.00642904\\
0.61	-0.00619447\\
0.612	-0.00594\\
0.614	-0.00566679\\
0.616	-0.00537607\\
0.618	-0.00506909\\
0.62	-0.00474718\\
0.622	-0.00441168\\
0.624	-0.00406397\\
0.626	-0.00370547\\
0.628	-0.0033376\\
0.63	-0.00296182\\
0.632	-0.00284444\\
0.634	-0.00297853\\
0.636	-0.00310309\\
0.638	-0.00321794\\
0.64	-0.00332291\\
0.642	-0.00341789\\
0.644	-0.00350276\\
0.646	-0.00357746\\
0.648	-0.00364194\\
0.65	-0.00369617\\
0.652	-0.00374017\\
0.654	-0.00377396\\
0.656	-0.00379761\\
0.658	-0.0038112\\
0.66	-0.00374034\\
0.662	-0.00365782\\
0.664	-0.00356764\\
0.666	-0.00347016\\
0.668	-0.00336571\\
0.67	-0.00340494\\
0.672	-0.00366143\\
0.674	-0.00390091\\
0.676	-0.00412274\\
0.678	-0.00432631\\
0.68	-0.00451113\\
0.682	-0.00467676\\
0.684	-0.00482283\\
0.686	-0.00494907\\
0.688	-0.00505527\\
0.69	-0.0051413\\
0.692	-0.00520712\\
0.694	-0.00525274\\
0.696	-0.00527828\\
0.698	-0.00528388\\
0.7	-0.00526981\\
0.702	-0.00523637\\
0.704	-0.00518394\\
0.706	-0.00511297\\
0.708	-0.00502111\\
0.71	-0.0049043\\
0.712	-0.00477091\\
0.714	-0.00462164\\
0.716	-0.00445723\\
0.718	-0.0042744\\
0.72	-0.00406487\\
0.722	-0.00384324\\
0.724	-0.00361045\\
0.726	-0.00336749\\
0.728	-0.00311534\\
0.73	-0.00285503\\
0.732	-0.00258759\\
0.734	-0.00231407\\
0.736	-0.00203553\\
0.738	-0.00175302\\
0.74	-0.00146762\\
0.742	-0.00137747\\
0.744	-0.00148636\\
0.746	-0.00162871\\
0.748	-0.00176418\\
0.75	-0.00189241\\
0.752	-0.00201303\\
0.754	-0.00212574\\
0.756	-0.00223025\\
0.758	-0.00232629\\
0.76	-0.00241364\\
0.762	-0.00249209\\
0.764	-0.00256149\\
0.766	-0.0026217\\
0.768	-0.00267261\\
0.77	-0.00271416\\
0.772	-0.00274629\\
0.774	-0.00276901\\
0.776	-0.00278234\\
0.778	-0.00278633\\
0.78	-0.00298109\\
0.782	-0.00313789\\
0.784	-0.00326668\\
0.786	-0.00338179\\
0.788	-0.00348297\\
0.79	-0.00357004\\
0.792	-0.00364284\\
0.794	-0.00370129\\
0.796	-0.00374536\\
0.798	-0.00377508\\
0.8	-0.00379051\\
0.802	-0.00379178\\
0.804	-0.00377906\\
0.806	-0.00375259\\
0.808	-0.00371263\\
0.81	-0.0036595\\
0.812	-0.00359356\\
0.814	-0.00351523\\
0.816	-0.00342494\\
0.818	-0.00332318\\
0.82	-0.00321047\\
0.822	-0.00308736\\
0.824	-0.00295444\\
0.826	-0.00281232\\
0.828	-0.00266164\\
0.83	-0.00250305\\
0.832	-0.0023346\\
0.834	-0.00215716\\
0.836	-0.00197405\\
0.838	-0.00178602\\
0.84	-0.0015938\\
0.842	-0.00139816\\
0.844	-0.00119298\\
0.846	-0.000985976\\
0.848	-0.000879727\\
0.85	-0.000940083\\
0.852	-0.00105584\\
0.854	-0.00116746\\
0.856	-0.00127423\\
0.858	-0.00137579\\
0.86	-0.00147186\\
0.862	-0.00156219\\
0.864	-0.00164652\\
0.866	-0.00172465\\
0.868	-0.00179639\\
0.87	-0.00186155\\
0.872	-0.00192\\
0.874	-0.00197161\\
0.876	-0.00201628\\
0.878	-0.00205393\\
0.88	-0.00208452\\
0.882	-0.002108\\
0.884	-0.00212437\\
0.886	-0.00213366\\
0.888	-0.00221085\\
0.89	-0.00231612\\
0.892	-0.00241123\\
0.894	-0.00249598\\
0.896	-0.00257022\\
0.898	-0.00263382\\
0.9	-0.0026867\\
0.902	-0.00272727\\
0.904	-0.00273777\\
0.906	-0.0027381\\
0.908	-0.00272838\\
0.91	-0.00270877\\
0.912	-0.00267948\\
0.914	-0.00264073\\
0.916	-0.00259277\\
0.918	-0.0025359\\
0.92	-0.00247044\\
0.922	-0.00239673\\
0.924	-0.00231515\\
0.926	-0.00222608\\
0.928	-0.00212996\\
0.93	-0.00202722\\
0.932	-0.00191832\\
0.934	-0.00180373\\
0.936	-0.00168395\\
0.938	-0.00155947\\
0.94	-0.00143081\\
0.942	-0.00129849\\
0.944	-0.00116304\\
0.946	-0.001025\\
0.948	-0.000884902\\
0.95	-0.000743285\\
0.952	-0.000650113\\
0.954	-0.00064688\\
0.956	-0.000639145\\
0.958	-0.000629286\\
0.96	-0.000685112\\
0.962	-0.000776822\\
0.964	-0.000865348\\
0.966	-0.000950142\\
0.968	-0.00103097\\
0.97	-0.00110761\\
0.972	-0.00117987\\
0.974	-0.00124599\\
0.976	-0.001306\\
0.978	-0.00136118\\
0.98	-0.00141141\\
0.982	-0.00145657\\
0.984	-0.00149658\\
0.986	-0.00153135\\
0.988	-0.00156083\\
0.99	-0.00158497\\
0.992	-0.00160375\\
0.994	-0.00161821\\
0.996	-0.00168159\\
0.998	-0.00173778\\
1	-0.00178669\\
1.002	-0.00182649\\
1.004	-0.00185705\\
1.006	-0.00188029\\
1.008	-0.00189622\\
1.01	-0.00190489\\
1.012	-0.00190636\\
1.014	-0.00190072\\
1.016	-0.00188809\\
1.018	-0.00186862\\
1.02	-0.00184247\\
1.022	-0.00180983\\
1.024	-0.0017709\\
1.026	-0.00172591\\
1.028	-0.00167512\\
1.03	-0.00161878\\
1.032	-0.00154956\\
1.034	-0.00147522\\
1.036	-0.00139639\\
1.038	-0.00131341\\
1.04	-0.00122663\\
1.042	-0.00113642\\
1.044	-0.00104314\\
1.046	-0.000947171\\
1.048	-0.000848899\\
1.05	-0.000748709\\
1.052	-0.000646988\\
1.054	-0.000581249\\
1.056	-0.000560729\\
1.058	-0.000539333\\
1.06	-0.000517133\\
1.062	-0.000494202\\
1.064	-0.000472575\\
1.066	-0.000469591\\
1.068	-0.000465166\\
1.07	-0.000465833\\
1.072	-0.000536217\\
1.074	-0.000604205\\
1.076	-0.000669484\\
1.078	-0.000730986\\
1.08	-0.000789035\\
1.082	-0.000843906\\
1.084	-0.000895458\\
1.086	-0.000943563\\
1.088	-0.000988104\\
1.09	-0.00102898\\
1.092	-0.00106609\\
1.094	-0.00109937\\
1.096	-0.00112903\\
1.098	-0.00117474\\
1.1	-0.0012155\\
1.102	-0.00125122\\
1.104	-0.00128185\\
1.106	-0.00130732\\
1.108	-0.00132761\\
1.11	-0.00134271\\
1.112	-0.00135264\\
1.114	-0.00135741\\
1.116	-0.00135707\\
1.118	-0.0013517\\
1.12	-0.00134136\\
1.122	-0.00132618\\
1.124	-0.00130625\\
1.126	-0.00128171\\
1.128	-0.00125271\\
1.13	-0.00121942\\
1.132	-0.001182\\
1.134	-0.00114065\\
1.136	-0.00109558\\
1.138	-0.00104698\\
1.14	-0.000995082\\
1.142	-0.000940124\\
1.144	-0.000882341\\
1.146	-0.000821983\\
1.148	-0.000759302\\
1.15	-0.000694559\\
1.152	-0.000628018\\
1.154	-0.00055595\\
1.156	-0.000481661\\
1.158	-0.00040649\\
1.16	-0.000363209\\
1.162	-0.000362737\\
1.164	-0.000361515\\
1.166	-0.000359563\\
1.168	-0.000356901\\
1.17	-0.000353554\\
1.172	-0.000349542\\
1.174	-0.00034489\\
1.176	-0.000339623\\
1.178	-0.000333765\\
1.18	-0.000327341\\
1.182	-0.000381065\\
1.184	-0.000443592\\
1.186	-0.000503705\\
1.188	-0.000561212\\
1.19	-0.000615932\\
1.192	-0.000667697\\
1.194	-0.000716351\\
1.196	-0.000761753\\
1.198	-0.000803774\\
1.2	-0.000841198\\
1.202	-0.00087364\\
1.204	-0.000902483\\
1.206	-0.000927664\\
1.208	-0.000949132\\
1.21	-0.000966852\\
1.212	-0.000980802\\
1.214	-0.000990976\\
1.216	-0.000997379\\
1.218	-0.00100003\\
1.22	-0.00099897\\
1.222	-0.000994238\\
1.224	-0.000985896\\
1.226	-0.000974017\\
1.228	-0.000958686\\
1.23	-0.000939998\\
1.232	-0.00091806\\
1.234	-0.000892991\\
1.236	-0.000864919\\
1.238	-0.00083398\\
1.24	-0.000800321\\
1.242	-0.000764096\\
1.244	-0.000725468\\
1.246	-0.000684606\\
1.248	-0.000641684\\
1.25	-0.000596884\\
1.252	-0.000556862\\
1.254	-0.000520926\\
1.256	-0.000483757\\
1.258	-0.000445488\\
1.26	-0.000406254\\
1.262	-0.000366191\\
1.264	-0.000325437\\
1.266	-0.00028413\\
1.268	-0.000241548\\
1.27	-0.000198442\\
1.272	-0.000185956\\
1.274	-0.000189887\\
1.276	-0.000193139\\
1.278	-0.000195713\\
1.28	-0.000197611\\
1.282	-0.000203398\\
1.284	-0.000252015\\
1.286	-0.000299151\\
1.288	-0.00034465\\
1.29	-0.000388364\\
1.292	-0.000430152\\
1.294	-0.000469884\\
1.296	-0.000507437\\
1.298	-0.000542698\\
1.3	-0.000575565\\
1.302	-0.000605943\\
1.304	-0.000633751\\
1.306	-0.0006583\\
1.308	-0.000679672\\
1.31	-0.000698332\\
1.312	-0.000714244\\
1.314	-0.000727378\\
1.316	-0.000737719\\
1.318	-0.000745258\\
1.32	-0.000750001\\
1.322	-0.00075196\\
1.324	-0.000751159\\
1.326	-0.000747631\\
1.328	-0.00074142\\
1.33	-0.000732577\\
1.332	-0.000721163\\
1.334	-0.000707249\\
1.336	-0.000690912\\
1.338	-0.000672238\\
1.34	-0.000651321\\
1.342	-0.000628262\\
1.344	-0.000603167\\
1.346	-0.00057615\\
1.348	-0.000558105\\
1.35	-0.000542205\\
1.352	-0.000524765\\
1.354	-0.000505861\\
1.356	-0.00048557\\
1.358	-0.000463974\\
1.36	-0.000441159\\
1.362	-0.000417212\\
1.364	-0.000392086\\
1.366	-0.000365301\\
1.368	-0.000337677\\
1.37	-0.000309313\\
1.372	-0.000280309\\
1.374	-0.000250768\\
1.376	-0.000220791\\
1.378	-0.000190481\\
1.38	-0.000159939\\
1.382	-0.000129267\\
1.384	-0.000109818\\
1.386	-0.000147987\\
1.388	-0.000185183\\
1.39	-0.000221282\\
1.392	-0.000256168\\
1.394	-0.000289728\\
1.396	-0.000321856\\
1.398	-0.000352452\\
1.4	-0.000381424\\
1.402	-0.000408685\\
1.404	-0.000433287\\
1.406	-0.00045579\\
1.408	-0.000476237\\
1.41	-0.000494747\\
1.412	-0.000511274\\
1.414	-0.000525781\\
1.416	-0.000538237\\
1.418	-0.000548619\\
1.42	-0.000556912\\
1.422	-0.000563109\\
1.424	-0.00056721\\
1.426	-0.000569224\\
1.428	-0.000569165\\
1.43	-0.000567056\\
1.432	-0.000562927\\
1.434	-0.000556815\\
1.436	-0.000548763\\
1.438	-0.000538821\\
1.44	-0.000527046\\
1.442	-0.000513499\\
1.444	-0.000498249\\
1.446	-0.00048137\\
1.448	-0.000462939\\
1.45	-0.000443039\\
1.452	-0.000424923\\
1.454	-0.000417986\\
1.456	-0.000409791\\
1.458	-0.000400381\\
1.46	-0.000389801\\
1.462	-0.0003781\\
1.464	-0.000365104\\
1.466	-0.000350728\\
1.468	-0.000335408\\
1.47	-0.000319204\\
1.472	-0.000302181\\
1.474	-0.000284404\\
1.476	-0.000265941\\
1.478	-0.00024686\\
1.48	-0.000227231\\
1.482	-0.000207125\\
1.484	-0.000186612\\
1.486	-0.000165765\\
1.488	-0.000144655\\
1.49	-0.000123354\\
1.492	-0.000127784\\
1.494	-0.000154832\\
1.496	-0.000181017\\
1.498	-0.000206256\\
1.5	-0.000230467\\
1.502	-0.000253576\\
1.504	-0.000275512\\
1.506	-0.00029621\\
1.508	-0.000315608\\
1.51	-0.000333652\\
1.512	-0.000350293\\
1.514	-0.000365486\\
1.516	-0.000379193\\
1.518	-0.000391383\\
1.52	-0.000402027\\
1.522	-0.000410984\\
1.524	-0.000417976\\
1.526	-0.000423396\\
1.528	-0.000427245\\
1.53	-0.000429523\\
1.532	-0.000430241\\
1.534	-0.000429413\\
1.536	-0.000427059\\
1.538	-0.000423203\\
1.54	-0.000417877\\
1.542	-0.000411114\\
1.544	-0.000402954\\
1.546	-0.000393443\\
1.548	-0.000382629\\
1.55	-0.000370564\\
1.552	-0.000357198\\
1.554	-0.000342642\\
1.556	-0.000327016\\
1.558	-0.000310388\\
1.56	-0.000301384\\
1.562	-0.000297679\\
1.564	-0.000293055\\
1.566	-0.000287541\\
1.568	-0.000281167\\
1.57	-0.000273964\\
1.572	-0.000265959\\
1.574	-0.000257047\\
1.576	-0.000247374\\
1.578	-0.00023703\\
1.58	-0.000226057\\
1.582	-0.000214498\\
1.584	-0.000202401\\
1.586	-0.000189809\\
1.588	-0.000176772\\
1.59	-0.000163338\\
1.592	-0.000149555\\
1.594	-0.000135472\\
1.596	-0.00012114\\
1.598	-0.000106607\\
1.6	-0.000125111\\
1.602	-0.000144201\\
1.604	-0.000162559\\
1.606	-0.000180129\\
1.608	-0.000196856\\
1.61	-0.00021269\\
1.612	-0.000227584\\
1.614	-0.000241496\\
1.616	-0.000254386\\
1.618	-0.00026622\\
1.62	-0.000276968\\
1.622	-0.000286603\\
1.624	-0.000295103\\
1.626	-0.000302452\\
1.628	-0.000308635\\
1.63	-0.000313643\\
1.632	-0.000317473\\
1.634	-0.000320123\\
1.636	-0.000321597\\
1.638	-0.000321904\\
1.64	-0.000321055\\
1.642	-0.000319067\\
1.644	-0.00031596\\
1.646	-0.000311757\\
1.648	-0.000306388\\
1.65	-0.000299759\\
1.652	-0.000292133\\
1.654	-0.00028355\\
1.656	-0.000274052\\
1.658	-0.000263681\\
1.66	-0.000252485\\
1.662	-0.000240514\\
1.664	-0.000227817\\
1.666	-0.00021445\\
1.668	-0.000209399\\
1.67	-0.000207678\\
1.672	-0.000205305\\
1.674	-0.000202296\\
1.676	-0.000198671\\
1.678	-0.000194448\\
1.68	-0.00018965\\
1.682	-0.0001843\\
1.684	-0.000178423\\
1.686	-0.000172047\\
1.688	-0.000165197\\
1.69	-0.000157903\\
1.692	-0.000150195\\
1.694	-0.000142103\\
1.696	-0.000133631\\
1.698	-0.000124795\\
1.7	-0.000115674\\
1.702	-0.000106303\\
1.704	-9.67158e-05\\
1.706	-9.86472e-05\\
1.708	-0.000112603\\
1.71	-0.000126\\
1.712	-0.000138796\\
1.714	-0.000150953\\
1.716	-0.000162434\\
1.718	-0.000173205\\
1.72	-0.000183237\\
1.722	-0.000192502\\
1.724	-0.000200976\\
1.726	-0.000208636\\
1.728	-0.000215466\\
1.73	-0.00022145\\
1.732	-0.000226576\\
1.734	-0.000230836\\
1.736	-0.000234224\\
1.738	-0.000236739\\
1.74	-0.00023838\\
1.742	-0.000239152\\
1.744	-0.000239063\\
1.746	-0.000238121\\
1.748	-0.000236341\\
1.75	-0.000233738\\
1.752	-0.00023033\\
1.754	-0.00022614\\
1.756	-0.00022119\\
1.758	-0.000215463\\
1.76	-0.000208999\\
1.762	-0.000201865\\
1.764	-0.000194042\\
1.766	-0.000185603\\
1.768	-0.0001766\\
1.77	-0.000167069\\
1.772	-0.000157053\\
1.774	-0.000146592\\
1.776	-0.000143853\\
1.778	-0.000143299\\
1.78	-0.000142248\\
1.782	-0.000140749\\
1.784	-0.000138813\\
1.786	-0.000136453\\
1.788	-0.000133681\\
1.79	-0.000130514\\
1.792	-0.000126966\\
1.794	-0.000123055\\
1.796	-0.000118798\\
1.798	-0.000114213\\
1.8	-0.000109321\\
1.802	-0.000104142\\
1.804	-9.86953e-05\\
1.806	-9.30039e-05\\
1.808	-8.70891e-05\\
1.81	-8.09734e-05\\
1.812	-7.62041e-05\\
1.814	-8.64524e-05\\
1.816	-9.62765e-05\\
1.818	-0.000105646\\
1.82	-0.000114533\\
1.822	-0.000122895\\
1.824	-0.000130723\\
1.826	-0.000137996\\
1.828	-0.000144696\\
1.83	-0.000150804\\
1.832	-0.000156307\\
1.834	-0.000161191\\
1.836	-0.000165446\\
1.838	-0.000169064\\
1.84	-0.00017204\\
1.842	-0.000174369\\
1.844	-0.000176052\\
1.846	-0.000177089\\
1.848	-0.000177484\\
1.85	-0.000177244\\
1.852	-0.000176374\\
1.854	-0.000174877\\
1.856	-0.000172706\\
1.858	-0.000169931\\
1.86	-0.000166584\\
1.862	-0.000162684\\
1.864	-0.00015825\\
1.866	-0.000153304\\
1.868	-0.00014787\\
1.87	-0.000141973\\
1.872	-0.000135638\\
1.874	-0.000128894\\
1.876	-0.000121769\\
1.878	-0.000114293\\
1.88	-0.000106496\\
1.882	-9.84096e-05\\
1.884	-9.70625e-05\\
1.886	-9.71316e-05\\
1.888	-9.68852e-05\\
1.89	-9.63288e-05\\
1.892	-9.54686e-05\\
1.894	-9.43115e-05\\
1.896	-9.28657e-05\\
1.898	-9.11399e-05\\
1.9	-8.91438e-05\\
1.902	-8.68875e-05\\
1.904	-8.43822e-05\\
1.906	-8.16395e-05\\
1.908	-7.86717e-05\\
1.91	-7.54915e-05\\
1.912	-7.20579e-05\\
1.914	-6.84343e-05\\
1.916	-6.46403e-05\\
1.918	-6.06906e-05\\
1.92	-6.5678e-05\\
1.922	-7.28632e-05\\
1.924	-7.97079e-05\\
1.926	-8.61918e-05\\
1.928	-9.22957e-05\\
1.93	-9.80026e-05\\
1.932	-0.000103296\\
1.934	-0.000108163\\
1.936	-0.000112591\\
1.938	-0.000116568\\
1.94	-0.000120087\\
1.942	-0.000123132\\
1.944	-0.000125702\\
1.946	-0.000127797\\
1.948	-0.000129415\\
1.95	-0.000130555\\
1.952	-0.000131219\\
1.954	-0.000131409\\
1.956	-0.000131078\\
1.958	-0.000130281\\
1.96	-0.000129031\\
1.962	-0.000127338\\
1.964	-0.000125213\\
1.966	-0.000122668\\
1.968	-0.000119717\\
1.97	-0.000116376\\
1.972	-0.000112661\\
1.974	-0.000108589\\
1.976	-0.000104178\\
1.978	-9.94488e-05\\
1.98	-9.44209e-05\\
1.982	-8.91158e-05\\
1.984	-8.35555e-05\\
1.986	-7.77624e-05\\
1.988	-7.17599e-05\\
1.99	-6.55716e-05\\
1.992	-6.48005e-05\\
1.994	-6.51697e-05\\
1.996	-6.53238e-05\\
1.998	-6.52653e-05\\
2	-6.49976e-05\\
};
\addplot [color=mycolor3, line width=1.0pt, forget plot]
  table[row sep=crcr]{%
-0.024	0\\
-0.022	0\\
-0.02	0\\
-0.018	0\\
-0.016	0\\
-0.014	0\\
-0.012	0\\
-0.01	0\\
-0.008	0\\
-0.006	0\\
-0.004	0\\
-0.002	0\\
0	0\\
0.002	3.62726e-09\\
0.004	1.65765e-08\\
0.006	1.79777e-08\\
0.008	-3.99277e-08\\
0.01	-2.2146e-07\\
0.012	-6.01995e-07\\
0.014	-1.26193e-06\\
0.016	-2.28196e-06\\
0.018	-1.69294e-06\\
0.02	-2.84817e-06\\
0.022	-3.35355e-06\\
0.024	-0.000967628\\
0.026	-0.00329502\\
0.028	-0.00519172\\
0.03	-0.00730253\\
0.032	-0.00960979\\
0.034	-0.0119182\\
0.036	-0.0128467\\
0.038	-0.0137442\\
0.04	-0.014603\\
0.042	-0.015416\\
0.044	-0.0161765\\
0.046	-0.0168782\\
0.048	-0.0175153\\
0.05	-0.0180825\\
0.052	-0.0185751\\
0.054	-0.0189887\\
0.056	-0.0193195\\
0.058	-0.0195644\\
0.06	-0.0197205\\
0.062	-0.0197859\\
0.064	-0.0197587\\
0.066	-0.019638\\
0.068	-0.0194231\\
0.07	-0.019114\\
0.072	-0.0187112\\
0.074	-0.0182155\\
0.076	-0.0176285\\
0.078	-0.0169519\\
0.08	-0.0161882\\
0.082	-0.0153401\\
0.084	-0.0144108\\
0.086	-0.0134039\\
0.088	-0.0123232\\
0.09	-0.0111732\\
0.092	-0.00995826\\
0.094	-0.0086834\\
0.096	-0.00735374\\
0.098	-0.00597466\\
0.1	-0.00455176\\
0.102	-0.00309084\\
0.104	-0.00182492\\
0.106	-0.00237695\\
0.108	-0.0028887\\
0.11	-0.0033591\\
0.112	-0.00378733\\
0.114	-0.00417274\\
0.116	-0.00396378\\
0.118	-0.00367811\\
0.12	-0.00337727\\
0.122	-0.00348658\\
0.124	-0.00425971\\
0.126	-0.00510289\\
0.128	-0.00592597\\
0.13	-0.00655439\\
0.132	-0.00692062\\
0.134	-0.00727332\\
0.136	-0.00761089\\
0.138	-0.00793175\\
0.14	-0.00823439\\
0.142	-0.00902383\\
0.144	-0.0110771\\
0.146	-0.0130336\\
0.148	-0.0148605\\
0.15	-0.0164066\\
0.152	-0.0178374\\
0.154	-0.0191499\\
0.156	-0.0203415\\
0.158	-0.0214101\\
0.16	-0.022354\\
0.162	-0.0231721\\
0.164	-0.0238635\\
0.166	-0.024428\\
0.168	-0.0248659\\
0.17	-0.0251777\\
0.172	-0.0253645\\
0.174	-0.0254278\\
0.176	-0.0253696\\
0.178	-0.0251922\\
0.18	-0.0248984\\
0.182	-0.0244912\\
0.184	-0.0239743\\
0.186	-0.0233514\\
0.188	-0.0224007\\
0.19	-0.0212779\\
0.192	-0.0200919\\
0.194	-0.0188472\\
0.196	-0.0175482\\
0.198	-0.0161997\\
0.2	-0.0148065\\
0.202	-0.0133738\\
0.204	-0.0119067\\
0.206	-0.0104103\\
0.208	-0.00889005\\
0.21	-0.0074138\\
0.212	-0.00757283\\
0.214	-0.00770404\\
0.216	-0.007807\\
0.218	-0.00788139\\
0.22	-0.00792704\\
0.222	-0.00794385\\
0.224	-0.00793186\\
0.226	-0.00789123\\
0.228	-0.0078222\\
0.23	-0.00772514\\
0.232	-0.00760053\\
0.234	-0.00727782\\
0.236	-0.00693588\\
0.238	-0.00657651\\
0.24	-0.00735801\\
0.242	-0.00888953\\
0.244	-0.0103659\\
0.246	-0.0117825\\
0.248	-0.0131349\\
0.25	-0.0144189\\
0.252	-0.0155533\\
0.254	-0.0165446\\
0.256	-0.017463\\
0.258	-0.0183058\\
0.26	-0.0190709\\
0.262	-0.0197563\\
0.264	-0.0203604\\
0.266	-0.020882\\
0.268	-0.0213203\\
0.27	-0.0216744\\
0.272	-0.0219443\\
0.274	-0.0221299\\
0.276	-0.0222317\\
0.278	-0.0222502\\
0.28	-0.0221864\\
0.282	-0.0220416\\
0.284	-0.0218174\\
0.286	-0.0215156\\
0.288	-0.0209335\\
0.29	-0.0202557\\
0.292	-0.0195155\\
0.294	-0.0187159\\
0.296	-0.0178605\\
0.298	-0.0169526\\
0.3	-0.0159959\\
0.302	-0.0149943\\
0.304	-0.0139517\\
0.306	-0.012872\\
0.308	-0.0117593\\
0.31	-0.010618\\
0.312	-0.00945216\\
0.314	-0.00826613\\
0.316	-0.00706421\\
0.318	-0.00692811\\
0.32	-0.00703937\\
0.322	-0.00712631\\
0.324	-0.00718881\\
0.326	-0.0072268\\
0.328	-0.00724032\\
0.33	-0.00722948\\
0.332	-0.00719449\\
0.334	-0.00713562\\
0.336	-0.00705324\\
0.338	-0.00694777\\
0.34	-0.00681972\\
0.342	-0.00666969\\
0.344	-0.00649833\\
0.346	-0.00630635\\
0.348	-0.00729291\\
0.35	-0.00825815\\
0.352	-0.00917921\\
0.354	-0.0100533\\
0.356	-0.0108779\\
0.358	-0.0116507\\
0.36	-0.0123696\\
0.362	-0.0130327\\
0.364	-0.0136382\\
0.366	-0.014185\\
0.368	-0.0146716\\
0.37	-0.0150972\\
0.372	-0.0154432\\
0.374	-0.0156951\\
0.376	-0.0158866\\
0.378	-0.0160177\\
0.38	-0.0160888\\
0.382	-0.0161002\\
0.384	-0.0160529\\
0.386	-0.0159475\\
0.388	-0.0157854\\
0.39	-0.0155677\\
0.392	-0.015296\\
0.394	-0.0149719\\
0.396	-0.0145974\\
0.398	-0.0141742\\
0.4	-0.0137047\\
0.402	-0.0131911\\
0.404	-0.0126359\\
0.406	-0.0120414\\
0.408	-0.0114105\\
0.41	-0.0107459\\
0.412	-0.0100503\\
0.414	-0.00932674\\
0.416	-0.00911889\\
0.418	-0.00887189\\
0.42	-0.00860658\\
0.422	-0.0083238\\
0.424	-0.00802445\\
0.426	-0.00770944\\
0.428	-0.00737973\\
0.43	-0.0070363\\
0.432	-0.00668014\\
0.434	-0.00631228\\
0.436	-0.00617229\\
0.438	-0.00618033\\
0.44	-0.00616861\\
0.442	-0.00613733\\
0.444	-0.00608674\\
0.446	-0.00601715\\
0.448	-0.00592892\\
0.45	-0.00582249\\
0.452	-0.00569834\\
0.454	-0.00589463\\
0.456	-0.00656019\\
0.458	-0.00719292\\
0.46	-0.00779095\\
0.462	-0.00835261\\
0.464	-0.00887278\\
0.466	-0.00934108\\
0.468	-0.00976951\\
0.47	-0.010157\\
0.472	-0.0105028\\
0.474	-0.0108061\\
0.476	-0.0110664\\
0.478	-0.0112832\\
0.48	-0.0114565\\
0.482	-0.0115861\\
0.484	-0.0116721\\
0.486	-0.0117148\\
0.488	-0.0117146\\
0.49	-0.011672\\
0.492	-0.01158\\
0.494	-0.0114382\\
0.496	-0.0112572\\
0.498	-0.0110382\\
0.5	-0.0107824\\
0.502	-0.0104913\\
0.504	-0.0101662\\
0.506	-0.0098089\\
0.508	-0.00942095\\
0.51	-0.00900416\\
0.512	-0.00856039\\
0.514	-0.00809157\\
0.516	-0.00759969\\
0.518	-0.00708682\\
0.52	-0.00655504\\
0.522	-0.00655272\\
0.524	-0.00658309\\
0.526	-0.00659624\\
0.528	-0.00659236\\
0.53	-0.00657171\\
0.532	-0.00653456\\
0.534	-0.00648126\\
0.536	-0.00641215\\
0.538	-0.00632765\\
0.54	-0.00620409\\
0.542	-0.00606459\\
0.544	-0.00591213\\
0.546	-0.00574726\\
0.548	-0.00557058\\
0.55	-0.00538267\\
0.552	-0.00518418\\
0.554	-0.00497574\\
0.556	-0.00483753\\
0.558	-0.00477688\\
0.56	-0.00470156\\
0.562	-0.00492018\\
0.564	-0.00534926\\
0.566	-0.00575337\\
0.568	-0.00613139\\
0.57	-0.00648232\\
0.572	-0.00680524\\
0.574	-0.00709937\\
0.576	-0.00736403\\
0.578	-0.00759866\\
0.58	-0.00780283\\
0.582	-0.0079762\\
0.584	-0.00811406\\
0.586	-0.00822004\\
0.588	-0.00829514\\
0.59	-0.0083395\\
0.592	-0.00835336\\
0.594	-0.00833709\\
0.596	-0.00829112\\
0.598	-0.00821603\\
0.6	-0.00811245\\
0.602	-0.00798114\\
0.604	-0.00782294\\
0.606	-0.00763877\\
0.608	-0.00742963\\
0.61	-0.00719661\\
0.612	-0.00694085\\
0.614	-0.006658\\
0.616	-0.00635484\\
0.618	-0.006033\\
0.62	-0.0056939\\
0.622	-0.00533896\\
0.624	-0.00496966\\
0.626	-0.00458752\\
0.628	-0.00419406\\
0.63	-0.00379084\\
0.632	-0.00364428\\
0.634	-0.00374756\\
0.636	-0.00383981\\
0.638	-0.00392094\\
0.64	-0.00399092\\
0.642	-0.00404974\\
0.644	-0.00409741\\
0.646	-0.004134\\
0.648	-0.00415956\\
0.65	-0.00417421\\
0.652	-0.00417806\\
0.654	-0.00417129\\
0.656	-0.00415405\\
0.658	-0.00412657\\
0.66	-0.00408905\\
0.662	-0.00404176\\
0.664	-0.00397773\\
0.666	-0.00390323\\
0.668	-0.00382015\\
0.67	-0.00387914\\
0.672	-0.00415374\\
0.674	-0.00440969\\
0.676	-0.00464631\\
0.678	-0.00486302\\
0.68	-0.0050593\\
0.682	-0.00523473\\
0.684	-0.00538896\\
0.686	-0.00552171\\
0.688	-0.0056328\\
0.69	-0.00572212\\
0.692	-0.00578965\\
0.694	-0.00583543\\
0.696	-0.0058596\\
0.698	-0.00586236\\
0.7	-0.005844\\
0.702	-0.00580486\\
0.704	-0.005745\\
0.706	-0.00566383\\
0.708	-0.00556342\\
0.71	-0.00544439\\
0.712	-0.00530742\\
0.714	-0.00515324\\
0.716	-0.00498264\\
0.718	-0.00479646\\
0.72	-0.00459557\\
0.722	-0.00438089\\
0.724	-0.00415338\\
0.726	-0.00391404\\
0.728	-0.00366389\\
0.73	-0.00340396\\
0.732	-0.00313533\\
0.734	-0.00285909\\
0.736	-0.00257475\\
0.738	-0.00228464\\
0.74	-0.0019903\\
0.742	-0.00188993\\
0.744	-0.00198738\\
0.746	-0.00211713\\
0.748	-0.00223889\\
0.75	-0.00235236\\
0.752	-0.00245725\\
0.754	-0.00255331\\
0.756	-0.00264032\\
0.758	-0.00271808\\
0.76	-0.00278644\\
0.762	-0.00284525\\
0.764	-0.00289444\\
0.766	-0.00293393\\
0.768	-0.00296368\\
0.77	-0.00298371\\
0.772	-0.00299403\\
0.774	-0.0029947\\
0.776	-0.00298582\\
0.778	-0.0029675\\
0.78	-0.00313994\\
0.782	-0.00330103\\
0.784	-0.00344789\\
0.786	-0.00358019\\
0.788	-0.00369522\\
0.79	-0.00379485\\
0.792	-0.00387938\\
0.794	-0.00394869\\
0.796	-0.00400276\\
0.798	-0.00404159\\
0.8	-0.00406525\\
0.802	-0.00407385\\
0.804	-0.00406758\\
0.806	-0.00404666\\
0.808	-0.00401136\\
0.81	-0.003962\\
0.812	-0.00389896\\
0.814	-0.00382265\\
0.816	-0.00373353\\
0.818	-0.00363209\\
0.82	-0.00351888\\
0.822	-0.00339445\\
0.824	-0.00325941\\
0.826	-0.00311441\\
0.828	-0.00295961\\
0.83	-0.00279608\\
0.832	-0.00262466\\
0.834	-0.00244607\\
0.836	-0.00226106\\
0.838	-0.00207038\\
0.84	-0.00187482\\
0.842	-0.00167513\\
0.844	-0.00147211\\
0.846	-0.00126653\\
0.848	-0.00116087\\
0.85	-0.00122099\\
0.852	-0.00133572\\
0.854	-0.00144553\\
0.856	-0.00154973\\
0.858	-0.00164784\\
0.86	-0.00173925\\
0.862	-0.00182425\\
0.864	-0.00190265\\
0.866	-0.00197426\\
0.868	-0.00203892\\
0.87	-0.00209649\\
0.872	-0.00214686\\
0.874	-0.00218994\\
0.876	-0.00222566\\
0.878	-0.00225398\\
0.88	-0.00227489\\
0.882	-0.00228839\\
0.884	-0.00229452\\
0.886	-0.00229332\\
0.888	-0.00235982\\
0.89	-0.00245424\\
0.892	-0.00253837\\
0.894	-0.00261205\\
0.896	-0.00267515\\
0.898	-0.00272759\\
0.9	-0.00276933\\
0.902	-0.00280034\\
0.904	-0.00282067\\
0.906	-0.0028304\\
0.908	-0.00282961\\
0.91	-0.00281847\\
0.912	-0.00279654\\
0.914	-0.00276434\\
0.916	-0.00272246\\
0.918	-0.0026712\\
0.92	-0.00261085\\
0.922	-0.00254177\\
0.924	-0.00246432\\
0.926	-0.0023789\\
0.928	-0.00228592\\
0.93	-0.00218582\\
0.932	-0.00207908\\
0.934	-0.00196616\\
0.936	-0.00184756\\
0.938	-0.00172379\\
0.94	-0.00159537\\
0.942	-0.00146283\\
0.944	-0.00132671\\
0.946	-0.00118756\\
0.948	-0.00104593\\
0.95	-0.000902357\\
0.952	-0.000806656\\
0.954	-0.000800492\\
0.956	-0.000791908\\
0.958	-0.000780961\\
0.96	-0.0008353\\
0.962	-0.000925138\\
0.964	-0.00101142\\
0.966	-0.00109361\\
0.968	-0.00117148\\
0.97	-0.00124485\\
0.972	-0.00131352\\
0.974	-0.00137731\\
0.976	-0.00143609\\
0.978	-0.00148972\\
0.98	-0.00153807\\
0.982	-0.00158098\\
0.984	-0.0016183\\
0.986	-0.00165011\\
0.988	-0.00167637\\
0.99	-0.00169706\\
0.992	-0.00171218\\
0.994	-0.00172277\\
0.996	-0.00178209\\
0.998	-0.00183406\\
1	-0.00187859\\
1.002	-0.0019156\\
1.004	-0.00194507\\
1.006	-0.00196698\\
1.008	-0.00198136\\
1.01	-0.00198824\\
1.012	-0.0019877\\
1.014	-0.00197984\\
1.016	-0.00196478\\
1.018	-0.00194268\\
1.02	-0.0019137\\
1.022	-0.00187804\\
1.024	-0.00183592\\
1.026	-0.00178758\\
1.028	-0.00173327\\
1.03	-0.00167328\\
1.032	-0.00160789\\
1.034	-0.00153742\\
1.036	-0.00146199\\
1.038	-0.00138203\\
1.04	-0.00129802\\
1.042	-0.00121031\\
1.044	-0.00111927\\
1.046	-0.00102527\\
1.048	-0.000928708\\
1.05	-0.00082996\\
1.052	-0.000729421\\
1.054	-0.000664602\\
1.056	-0.000644746\\
1.058	-0.000623761\\
1.06	-0.000601723\\
1.062	-0.000578711\\
1.064	-0.000556766\\
1.066	-0.000553232\\
1.068	-0.000548033\\
1.07	-0.000547709\\
1.072	-0.000616892\\
1.074	-0.000683478\\
1.076	-0.000747107\\
1.078	-0.000807657\\
1.08	-0.00086498\\
1.082	-0.000918928\\
1.084	-0.000969369\\
1.086	-0.00101618\\
1.088	-0.00105925\\
1.09	-0.00109849\\
1.092	-0.00113381\\
1.094	-0.00116515\\
1.096	-0.00119354\\
1.098	-0.00124068\\
1.1	-0.00128265\\
1.102	-0.00131936\\
1.104	-0.00135074\\
1.106	-0.00137672\\
1.108	-0.00139725\\
1.11	-0.00141237\\
1.112	-0.0014221\\
1.114	-0.00142645\\
1.116	-0.0014255\\
1.118	-0.00141929\\
1.12	-0.00140793\\
1.122	-0.00139151\\
1.124	-0.00137017\\
1.126	-0.00134403\\
1.128	-0.00131326\\
1.13	-0.00127802\\
1.132	-0.0012385\\
1.134	-0.0011949\\
1.136	-0.00114743\\
1.138	-0.0010963\\
1.14	-0.00104176\\
1.142	-0.000984043\\
1.144	-0.0009234\\
1.146	-0.000860091\\
1.148	-0.00079438\\
1.15	-0.000726538\\
1.152	-0.000656841\\
1.154	-0.000585566\\
1.156	-0.000512995\\
1.158	-0.00043941\\
1.16	-0.000397531\\
1.162	-0.000398278\\
1.164	-0.000398145\\
1.166	-0.00039715\\
1.168	-0.000395317\\
1.17	-0.000392667\\
1.172	-0.000389226\\
1.174	-0.000385019\\
1.176	-0.000380071\\
1.178	-0.00037441\\
1.18	-0.000368065\\
1.182	-0.00042175\\
1.184	-0.000484124\\
1.186	-0.000543975\\
1.188	-0.000601112\\
1.19	-0.00065536\\
1.192	-0.000706553\\
1.194	-0.000754542\\
1.196	-0.000799187\\
1.198	-0.000840367\\
1.2	-0.000877954\\
1.202	-0.000911865\\
1.204	-0.000942026\\
1.206	-0.00096837\\
1.208	-0.000990847\\
1.21	-0.00100942\\
1.212	-0.00102407\\
1.214	-0.00103478\\
1.216	-0.00104157\\
1.218	-0.00104445\\
1.22	-0.00104347\\
1.222	-0.00103867\\
1.224	-0.00103011\\
1.226	-0.00101787\\
1.228	-0.00100204\\
1.23	-0.000982709\\
1.232	-0.000959981\\
1.234	-0.000933997\\
1.236	-0.00090489\\
1.238	-0.000872803\\
1.24	-0.000837889\\
1.242	-0.000800307\\
1.244	-0.000760227\\
1.246	-0.000717825\\
1.248	-0.000673283\\
1.25	-0.000626788\\
1.252	-0.000585004\\
1.254	-0.000547247\\
1.256	-0.000508204\\
1.258	-0.000468018\\
1.26	-0.000426829\\
1.262	-0.000384782\\
1.264	-0.000342022\\
1.266	-0.000298694\\
1.268	-0.000254943\\
1.27	-0.000210915\\
1.272	-0.000197627\\
1.274	-0.000202417\\
1.276	-0.00020647\\
1.278	-0.000209788\\
1.28	-0.000212372\\
1.282	-0.000218785\\
1.284	-0.000267955\\
1.286	-0.000315568\\
1.288	-0.000361488\\
1.29	-0.000405565\\
1.292	-0.000447661\\
1.294	-0.000487644\\
1.296	-0.000525395\\
1.298	-0.0005608\\
1.3	-0.000593758\\
1.302	-0.000624176\\
1.304	-0.000651974\\
1.306	-0.00067708\\
1.308	-0.000699433\\
1.31	-0.000718984\\
1.312	-0.000735693\\
1.314	-0.000749533\\
1.316	-0.000760485\\
1.318	-0.000768543\\
1.32	-0.000773711\\
1.322	-0.000776004\\
1.324	-0.000775439\\
1.326	-0.000772056\\
1.328	-0.000765901\\
1.33	-0.000757027\\
1.332	-0.000745499\\
1.334	-0.000731387\\
1.336	-0.000714774\\
1.338	-0.000695747\\
1.34	-0.000674403\\
1.342	-0.000650847\\
1.344	-0.000625188\\
1.346	-0.000597544\\
1.348	-0.000578813\\
1.35	-0.00056217\\
1.352	-0.000543937\\
1.354	-0.000524192\\
1.356	-0.000503009\\
1.358	-0.000480483\\
1.36	-0.000456703\\
1.362	-0.000431761\\
1.364	-0.000405753\\
1.366	-0.000378775\\
1.368	-0.000350925\\
1.37	-0.000322306\\
1.372	-0.000293018\\
1.374	-0.000263165\\
1.376	-0.000232849\\
1.378	-0.000202173\\
1.38	-0.000171242\\
1.382	-0.000140157\\
1.384	-0.000120274\\
1.386	-0.000157988\\
1.388	-0.00019471\\
1.39	-0.000230318\\
1.392	-0.000264697\\
1.394	-0.000297735\\
1.396	-0.000329329\\
1.398	-0.000359379\\
1.4	-0.000387796\\
1.402	-0.000414494\\
1.404	-0.000439396\\
1.406	-0.000462432\\
1.408	-0.000483537\\
1.41	-0.000502653\\
1.412	-0.000519739\\
1.414	-0.000534758\\
1.416	-0.000547677\\
1.418	-0.000558475\\
1.42	-0.000567135\\
1.422	-0.000573651\\
1.424	-0.000578023\\
1.426	-0.000580259\\
1.428	-0.000580375\\
1.43	-0.000578395\\
1.432	-0.000574349\\
1.434	-0.000568275\\
1.436	-0.000560218\\
1.438	-0.000550229\\
1.44	-0.000538365\\
1.442	-0.000524691\\
1.444	-0.000509276\\
1.446	-0.000492195\\
1.448	-0.000473527\\
1.45	-0.000453358\\
1.452	-0.000434944\\
1.454	-0.000427679\\
1.456	-0.000419131\\
1.458	-0.000409345\\
1.46	-0.000398367\\
1.462	-0.000386249\\
1.464	-0.000373045\\
1.466	-0.000358811\\
1.468	-0.000343608\\
1.47	-0.000327497\\
1.472	-0.000310541\\
1.474	-0.000292807\\
1.476	-0.000274362\\
1.478	-0.000255275\\
1.48	-0.000235614\\
1.482	-0.000215452\\
1.484	-0.000194861\\
1.486	-0.000173914\\
1.488	-0.000152682\\
1.49	-0.000131239\\
1.492	-0.000135506\\
1.494	-0.000162372\\
1.496	-0.000188357\\
1.498	-0.000213377\\
1.5	-0.000237354\\
1.502	-0.000260213\\
1.504	-0.000281884\\
1.506	-0.000302303\\
1.508	-0.00032141\\
1.51	-0.000339151\\
1.512	-0.000355478\\
1.514	-0.000370348\\
1.516	-0.000383724\\
1.518	-0.000395575\\
1.52	-0.000405875\\
1.522	-0.000414604\\
1.524	-0.000421749\\
1.526	-0.000427302\\
1.528	-0.000431261\\
1.53	-0.000433629\\
1.532	-0.000434415\\
1.534	-0.000433632\\
1.536	-0.000431303\\
1.538	-0.000427453\\
1.54	-0.000422115\\
1.542	-0.000415322\\
1.544	-0.000407117\\
1.546	-0.000397544\\
1.548	-0.000386652\\
1.55	-0.000374497\\
1.552	-0.000361134\\
1.554	-0.000346627\\
1.556	-0.000331039\\
1.558	-0.000314439\\
1.56	-0.000305452\\
1.562	-0.000301753\\
1.564	-0.000297127\\
1.566	-0.0002916\\
1.568	-0.000285202\\
1.57	-0.000277967\\
1.572	-0.000269928\\
1.574	-0.000261122\\
1.576	-0.000251589\\
1.578	-0.000241369\\
1.58	-0.000230504\\
1.582	-0.000219039\\
1.584	-0.000207019\\
1.586	-0.00019449\\
1.588	-0.0001815\\
1.59	-0.000168097\\
1.592	-0.00015433\\
1.594	-0.000140249\\
1.596	-0.000125904\\
1.598	-0.000111345\\
1.6	-0.000129808\\
1.602	-0.000148843\\
1.604	-0.000167134\\
1.606	-0.000184623\\
1.608	-0.000201258\\
1.61	-0.000216988\\
1.612	-0.000231767\\
1.614	-0.000245553\\
1.616	-0.000258308\\
1.618	-0.000269998\\
1.62	-0.000280594\\
1.622	-0.000290069\\
1.624	-0.000298402\\
1.626	-0.000305576\\
1.628	-0.00031158\\
1.63	-0.000316403\\
1.632	-0.000320044\\
1.634	-0.000322501\\
1.636	-0.000323779\\
1.638	-0.000323888\\
1.64	-0.000322839\\
1.642	-0.00032065\\
1.644	-0.000317341\\
1.646	-0.000312937\\
1.648	-0.000307465\\
1.65	-0.000300958\\
1.652	-0.00029345\\
1.654	-0.000284979\\
1.656	-0.000275586\\
1.658	-0.000265315\\
1.66	-0.000254212\\
1.662	-0.000242328\\
1.664	-0.000229712\\
1.666	-0.000216419\\
1.668	-0.000211435\\
1.67	-0.000209775\\
1.672	-0.000207457\\
1.674	-0.000204495\\
1.676	-0.00020091\\
1.678	-0.000196721\\
1.68	-0.00019195\\
1.682	-0.000186621\\
1.684	-0.000180758\\
1.686	-0.000174388\\
1.688	-0.000167539\\
1.69	-0.000160239\\
1.692	-0.000152519\\
1.694	-0.000144409\\
1.696	-0.000135941\\
1.698	-0.000127147\\
1.7	-0.000118061\\
1.702	-0.000108716\\
1.704	-9.9146e-05\\
1.706	-0.000101087\\
1.708	-0.000115045\\
1.71	-0.000128436\\
1.712	-0.000141219\\
1.714	-0.000153355\\
1.716	-0.000164808\\
1.718	-0.000175545\\
1.72	-0.000185536\\
1.722	-0.000194754\\
1.724	-0.000203174\\
1.726	-0.000210777\\
1.728	-0.000217542\\
1.73	-0.000223457\\
1.732	-0.00022851\\
1.734	-0.000232693\\
1.736	-0.000236\\
1.738	-0.000238429\\
1.74	-0.000239983\\
1.742	-0.000240665\\
1.744	-0.000240482\\
1.746	-0.000239446\\
1.748	-0.000237569\\
1.75	-0.000234869\\
1.752	-0.000231362\\
1.754	-0.000227073\\
1.756	-0.000222023\\
1.758	-0.000216241\\
1.76	-0.000209755\\
1.762	-0.000202596\\
1.764	-0.000194797\\
1.766	-0.000186393\\
1.768	-0.000177421\\
1.77	-0.000167919\\
1.772	-0.000157928\\
1.774	-0.000147487\\
1.776	-0.000144766\\
1.778	-0.000144231\\
1.78	-0.000143235\\
1.782	-0.000141787\\
1.784	-0.000139898\\
1.786	-0.00013758\\
1.788	-0.000134848\\
1.79	-0.000131715\\
1.792	-0.000128197\\
1.794	-0.000124312\\
1.796	-0.000120077\\
1.798	-0.00011551\\
1.8	-0.000110632\\
1.802	-0.000105461\\
1.804	-0.00010002\\
1.806	-9.43303e-05\\
1.808	-8.84129e-05\\
1.81	-8.22907e-05\\
1.812	-7.75111e-05\\
1.814	-8.77455e-05\\
1.816	-9.7552e-05\\
1.818	-0.000106901\\
1.82	-0.000115763\\
1.822	-0.000124114\\
1.824	-0.000131929\\
1.826	-0.000139187\\
1.828	-0.000145868\\
1.83	-0.000151954\\
1.832	-0.000157432\\
1.834	-0.000162288\\
1.836	-0.000166513\\
1.838	-0.000170098\\
1.84	-0.000173038\\
1.842	-0.00017533\\
1.844	-0.000176973\\
1.846	-0.000177968\\
1.848	-0.000178321\\
1.85	-0.000178036\\
1.852	-0.000177121\\
1.854	-0.000175588\\
1.856	-0.000173449\\
1.858	-0.000170718\\
1.86	-0.000167411\\
1.862	-0.000163548\\
1.864	-0.000159148\\
1.866	-0.000154232\\
1.868	-0.000148825\\
1.87	-0.00014295\\
1.872	-0.000136635\\
1.874	-0.000129906\\
1.876	-0.000122793\\
1.878	-0.000115324\\
1.88	-0.000107531\\
1.882	-9.94442e-05\\
1.884	-9.80929e-05\\
1.886	-9.81539e-05\\
1.888	-9.78958e-05\\
1.89	-9.73238e-05\\
1.892	-9.64443e-05\\
1.894	-9.52646e-05\\
1.896	-9.37926e-05\\
1.898	-9.20373e-05\\
1.9	-9.00084e-05\\
1.902	-8.77165e-05\\
1.904	-8.51726e-05\\
1.906	-8.23886e-05\\
1.908	-7.93769e-05\\
1.91	-7.61506e-05\\
1.912	-7.27232e-05\\
1.914	-6.91087e-05\\
1.916	-6.53215e-05\\
1.918	-6.13764e-05\\
1.92	-6.63661e-05\\
1.922	-7.35516e-05\\
1.924	-8.03944e-05\\
1.926	-8.68742e-05\\
1.928	-9.29723e-05\\
1.93	-9.86712e-05\\
1.932	-0.000103955\\
1.934	-0.000108811\\
1.936	-0.000113225\\
1.938	-0.000117187\\
1.94	-0.000120689\\
1.942	-0.000123723\\
1.944	-0.000126284\\
1.946	-0.000128368\\
1.948	-0.000129973\\
1.95	-0.000131099\\
1.952	-0.000131747\\
1.954	-0.000131921\\
1.956	-0.000131625\\
1.958	-0.000130866\\
1.96	-0.000129652\\
1.962	-0.000127992\\
1.964	-0.000125897\\
1.966	-0.000123379\\
1.968	-0.000120453\\
1.97	-0.000117134\\
1.972	-0.000113437\\
1.974	-0.00010938\\
1.976	-0.000104981\\
1.978	-0.000100261\\
1.98	-9.52388e-05\\
1.982	-8.99364e-05\\
1.984	-8.43757e-05\\
1.986	-7.85793e-05\\
1.988	-7.25705e-05\\
1.99	-6.63729e-05\\
1.992	-6.55897e-05\\
1.994	-6.5944e-05\\
1.996	-6.60805e-05\\
1.998	-6.6002e-05\\
2	-6.57117e-05\\
};
\addplot [color=mycolor4, line width=1.0pt, forget plot]
  table[row sep=crcr]{%
-0.024	0\\
-0.022	0\\
-0.02	0\\
-0.018	0\\
-0.016	0\\
-0.014	0\\
-0.012	0\\
-0.01	0\\
-0.008	0\\
-0.006	0\\
-0.004	0\\
-0.002	0\\
0	0\\
0.002	7.99727e-09\\
0.004	4.55661e-08\\
0.006	1.25842e-07\\
0.008	1.92322e-07\\
0.01	1.8907e-07\\
0.012	3.72659e-08\\
0.014	-3.51783e-07\\
0.016	-1.07113e-06\\
0.018	1.30383e-07\\
0.02	-5.96655e-07\\
0.022	-4.43691e-07\\
0.024	-0.000964249\\
0.026	-0.00329052\\
0.028	-0.00518669\\
0.03	-0.00729705\\
0.032	-0.00960397\\
0.034	-0.0120852\\
0.036	-0.0136363\\
0.038	-0.0145965\\
0.04	-0.0155176\\
0.042	-0.0163922\\
0.044	-0.0172132\\
0.046	-0.0179743\\
0.048	-0.0186694\\
0.05	-0.0192929\\
0.052	-0.0198399\\
0.054	-0.0203058\\
0.056	-0.0206867\\
0.058	-0.0209792\\
0.06	-0.0211803\\
0.062	-0.0212878\\
0.064	-0.0212997\\
0.066	-0.0212149\\
0.068	-0.0210325\\
0.07	-0.0207524\\
0.072	-0.0203748\\
0.074	-0.0199004\\
0.076	-0.0193306\\
0.078	-0.0186671\\
0.08	-0.017912\\
0.082	-0.017068\\
0.084	-0.0161382\\
0.086	-0.0151259\\
0.088	-0.014035\\
0.09	-0.0128696\\
0.092	-0.0116343\\
0.094	-0.0103339\\
0.096	-0.00897339\\
0.098	-0.0075582\\
0.1	-0.00609389\\
0.102	-0.00458624\\
0.104	-0.0032683\\
0.106	-0.00376307\\
0.108	-0.00421236\\
0.11	-0.00461522\\
0.112	-0.00497092\\
0.114	-0.00527894\\
0.116	-0.0051929\\
0.118	-0.00483033\\
0.12	-0.00444801\\
0.122	-0.00447146\\
0.124	-0.00515453\\
0.126	-0.0059037\\
0.128	-0.00662906\\
0.13	-0.00732823\\
0.132	-0.0078681\\
0.134	-0.00811435\\
0.136	-0.00834189\\
0.138	-0.00854949\\
0.14	-0.00873594\\
0.142	-0.00940659\\
0.144	-0.0113388\\
0.146	-0.0131724\\
0.148	-0.0149025\\
0.15	-0.0165245\\
0.152	-0.0180345\\
0.154	-0.019429\\
0.156	-0.0207049\\
0.158	-0.0218599\\
0.16	-0.0228918\\
0.162	-0.023799\\
0.164	-0.0245806\\
0.166	-0.0252359\\
0.168	-0.025755\\
0.17	-0.0260775\\
0.172	-0.0262723\\
0.174	-0.0263409\\
0.176	-0.0262852\\
0.178	-0.0261076\\
0.18	-0.0258109\\
0.182	-0.0253983\\
0.184	-0.0248732\\
0.186	-0.0242398\\
0.188	-0.0235022\\
0.19	-0.0226649\\
0.192	-0.021733\\
0.194	-0.0207116\\
0.196	-0.019606\\
0.198	-0.018422\\
0.2	-0.0171654\\
0.202	-0.0158114\\
0.204	-0.0143256\\
0.206	-0.0128044\\
0.208	-0.0112532\\
0.21	-0.00974019\\
0.212	-0.00985678\\
0.214	-0.00994007\\
0.216	-0.00998987\\
0.218	-0.0100061\\
0.22	-0.00998888\\
0.222	-0.00993835\\
0.224	-0.00985484\\
0.226	-0.00973881\\
0.228	-0.0095908\\
0.23	-0.0094115\\
0.232	-0.00920169\\
0.234	-0.00896227\\
0.236	-0.00862165\\
0.238	-0.00824612\\
0.24	-0.0089608\\
0.242	-0.0104088\\
0.244	-0.0117997\\
0.246	-0.0131289\\
0.248	-0.0143925\\
0.25	-0.0155865\\
0.252	-0.0167075\\
0.254	-0.0177523\\
0.256	-0.018718\\
0.258	-0.0196023\\
0.26	-0.0204028\\
0.262	-0.021101\\
0.264	-0.0216133\\
0.266	-0.0220424\\
0.268	-0.0223876\\
0.27	-0.0226485\\
0.272	-0.0228251\\
0.274	-0.0229177\\
0.276	-0.0229269\\
0.278	-0.0228534\\
0.28	-0.0226986\\
0.282	-0.0224638\\
0.284	-0.0221508\\
0.286	-0.0217616\\
0.288	-0.0212985\\
0.29	-0.0207639\\
0.292	-0.0201607\\
0.294	-0.0194918\\
0.296	-0.0187603\\
0.298	-0.0179407\\
0.3	-0.0170159\\
0.302	-0.0160425\\
0.304	-0.0150243\\
0.306	-0.0139653\\
0.308	-0.0128698\\
0.31	-0.011742\\
0.312	-0.0105862\\
0.314	-0.00940678\\
0.316	-0.00820808\\
0.318	-0.00807194\\
0.32	-0.00817998\\
0.322	-0.00826063\\
0.324	-0.00831388\\
0.326	-0.00833976\\
0.328	-0.00833845\\
0.33	-0.00831017\\
0.332	-0.00825525\\
0.334	-0.00817409\\
0.336	-0.00806718\\
0.338	-0.00793509\\
0.34	-0.00777847\\
0.342	-0.00759804\\
0.344	-0.00739457\\
0.346	-0.00716894\\
0.348	-0.00815332\\
0.35	-0.00913187\\
0.352	-0.0100628\\
0.354	-0.0109436\\
0.356	-0.011756\\
0.358	-0.0124968\\
0.36	-0.0131823\\
0.362	-0.0138107\\
0.364	-0.0143805\\
0.366	-0.0148904\\
0.368	-0.0153394\\
0.37	-0.0157266\\
0.372	-0.0160514\\
0.374	-0.0163135\\
0.376	-0.0165127\\
0.378	-0.0166492\\
0.38	-0.0167233\\
0.382	-0.0167356\\
0.384	-0.0166868\\
0.386	-0.0165779\\
0.388	-0.0164101\\
0.39	-0.0161848\\
0.392	-0.0159005\\
0.394	-0.0155364\\
0.396	-0.0151215\\
0.398	-0.0146578\\
0.4	-0.0141476\\
0.402	-0.0135932\\
0.404	-0.0129973\\
0.406	-0.0123624\\
0.408	-0.0116913\\
0.41	-0.0109868\\
0.412	-0.0102519\\
0.414	-0.00948956\\
0.416	-0.0092754\\
0.418	-0.00906886\\
0.42	-0.00884176\\
0.422	-0.00859493\\
0.424	-0.00832924\\
0.426	-0.0080456\\
0.428	-0.00774495\\
0.43	-0.00742826\\
0.432	-0.00709654\\
0.434	-0.00675083\\
0.436	-0.00663072\\
0.438	-0.00665639\\
0.44	-0.00665101\\
0.442	-0.00661941\\
0.444	-0.00656717\\
0.446	-0.00649464\\
0.448	-0.00640223\\
0.45	-0.00629041\\
0.452	-0.00615973\\
0.454	-0.00634838\\
0.456	-0.00700527\\
0.458	-0.00762832\\
0.46	-0.00821575\\
0.462	-0.00876591\\
0.464	-0.00927733\\
0.466	-0.0097487\\
0.468	-0.0101788\\
0.47	-0.0105668\\
0.472	-0.0109116\\
0.474	-0.0112128\\
0.476	-0.0114697\\
0.478	-0.0116773\\
0.48	-0.0118377\\
0.482	-0.0119538\\
0.484	-0.0120257\\
0.486	-0.0120538\\
0.488	-0.0120384\\
0.49	-0.0119803\\
0.492	-0.0118802\\
0.494	-0.011739\\
0.496	-0.0115577\\
0.498	-0.0113374\\
0.5	-0.0110796\\
0.502	-0.0107856\\
0.504	-0.0104569\\
0.506	-0.0100951\\
0.508	-0.009702\\
0.51	-0.00927939\\
0.512	-0.00882915\\
0.514	-0.00835325\\
0.516	-0.00785374\\
0.518	-0.00733271\\
0.52	-0.00679229\\
0.522	-0.00677845\\
0.524	-0.00679296\\
0.526	-0.00679021\\
0.528	-0.00677042\\
0.53	-0.00673391\\
0.532	-0.00668098\\
0.534	-0.00661201\\
0.536	-0.00652739\\
0.538	-0.00642757\\
0.54	-0.00631301\\
0.542	-0.0061842\\
0.544	-0.00604169\\
0.546	-0.00588602\\
0.548	-0.00571778\\
0.55	-0.00553758\\
0.552	-0.00534603\\
0.554	-0.0051438\\
0.556	-0.00501107\\
0.558	-0.00495517\\
0.56	-0.00488389\\
0.562	-0.00510379\\
0.564	-0.00553275\\
0.566	-0.00593623\\
0.568	-0.00631312\\
0.57	-0.00666244\\
0.572	-0.0069833\\
0.574	-0.00727492\\
0.576	-0.00753667\\
0.578	-0.00776799\\
0.58	-0.00796847\\
0.582	-0.00813781\\
0.584	-0.00827581\\
0.586	-0.0083824\\
0.588	-0.00845762\\
0.59	-0.00850162\\
0.592	-0.00851466\\
0.594	-0.00849711\\
0.596	-0.00844944\\
0.598	-0.00837224\\
0.6	-0.00826503\\
0.602	-0.00812927\\
0.604	-0.00796635\\
0.606	-0.00777721\\
0.608	-0.00756289\\
0.61	-0.00732447\\
0.612	-0.00706315\\
0.614	-0.00678016\\
0.616	-0.00647679\\
0.618	-0.00615441\\
0.62	-0.00581443\\
0.622	-0.00545829\\
0.624	-0.00508748\\
0.626	-0.00470353\\
0.628	-0.00430799\\
0.63	-0.00390243\\
0.632	-0.00375326\\
0.634	-0.00385371\\
0.636	-0.0039429\\
0.638	-0.00402078\\
0.64	-0.00408731\\
0.642	-0.00414252\\
0.644	-0.00418643\\
0.646	-0.0042191\\
0.648	-0.00424063\\
0.65	-0.00425114\\
0.652	-0.00424969\\
0.654	-0.00423727\\
0.656	-0.00421441\\
0.658	-0.00418131\\
0.66	-0.00413823\\
0.662	-0.00408542\\
0.664	-0.00402316\\
0.666	-0.00395177\\
0.668	-0.00387156\\
0.67	-0.00393316\\
0.672	-0.00421012\\
0.674	-0.00446818\\
0.676	-0.00470666\\
0.678	-0.00492499\\
0.68	-0.00512266\\
0.682	-0.00529923\\
0.684	-0.00545437\\
0.686	-0.00558781\\
0.688	-0.00569936\\
0.69	-0.0057887\\
0.692	-0.00585575\\
0.694	-0.00590088\\
0.696	-0.00592423\\
0.698	-0.00592602\\
0.7	-0.00590653\\
0.702	-0.00586613\\
0.704	-0.00580524\\
0.706	-0.00572436\\
0.708	-0.00562406\\
0.71	-0.00550497\\
0.712	-0.00536775\\
0.714	-0.00521317\\
0.716	-0.005042\\
0.718	-0.0048551\\
0.72	-0.00465334\\
0.722	-0.00443748\\
0.724	-0.00420849\\
0.726	-0.00396757\\
0.728	-0.00371573\\
0.73	-0.00345404\\
0.732	-0.00318356\\
0.734	-0.0029054\\
0.736	-0.00262065\\
0.738	-0.00233043\\
0.74	-0.00203586\\
0.742	-0.00193514\\
0.744	-0.00203212\\
0.746	-0.00216128\\
0.748	-0.00228236\\
0.75	-0.00239504\\
0.752	-0.00249906\\
0.754	-0.00259415\\
0.756	-0.0026801\\
0.758	-0.00275673\\
0.76	-0.00282389\\
0.762	-0.00288144\\
0.764	-0.00292929\\
0.766	-0.0029674\\
0.768	-0.00299573\\
0.77	-0.00301428\\
0.772	-0.00302308\\
0.774	-0.00302221\\
0.776	-0.00301176\\
0.778	-0.00299186\\
0.78	-0.00316258\\
0.782	-0.00332177\\
0.784	-0.00346673\\
0.786	-0.00359715\\
0.788	-0.00371274\\
0.79	-0.0038133\\
0.792	-0.00389867\\
0.794	-0.00396874\\
0.796	-0.00402349\\
0.798	-0.0040629\\
0.8	-0.00408707\\
0.802	-0.0040961\\
0.804	-0.00409018\\
0.806	-0.00406952\\
0.808	-0.00403442\\
0.81	-0.00398518\\
0.812	-0.00392219\\
0.814	-0.00384586\\
0.816	-0.00375656\\
0.818	-0.00365487\\
0.82	-0.00354135\\
0.822	-0.00341657\\
0.824	-0.00328114\\
0.826	-0.00313568\\
0.828	-0.00298087\\
0.83	-0.00281739\\
0.832	-0.00264596\\
0.834	-0.0024673\\
0.836	-0.00228216\\
0.838	-0.00209131\\
0.84	-0.00189551\\
0.842	-0.00169554\\
0.844	-0.00149215\\
0.846	-0.00128612\\
0.848	-0.00117997\\
0.85	-0.00123956\\
0.852	-0.00135373\\
0.854	-0.00146294\\
0.856	-0.00156652\\
0.858	-0.00166414\\
0.86	-0.00175554\\
0.862	-0.0018405\\
0.864	-0.00191881\\
0.866	-0.00199029\\
0.868	-0.00205477\\
0.87	-0.00211213\\
0.872	-0.00216224\\
0.874	-0.00220504\\
0.876	-0.00224044\\
0.878	-0.00226841\\
0.88	-0.00228894\\
0.882	-0.00230204\\
0.884	-0.00230773\\
0.886	-0.00230608\\
0.888	-0.00237211\\
0.89	-0.00246603\\
0.892	-0.00254965\\
0.894	-0.00262281\\
0.896	-0.00268538\\
0.898	-0.00273727\\
0.9	-0.00277844\\
0.902	-0.0028089\\
0.904	-0.00282866\\
0.906	-0.00283781\\
0.908	-0.00283645\\
0.91	-0.00282467\\
0.912	-0.00280273\\
0.914	-0.00277083\\
0.916	-0.00272922\\
0.918	-0.0026782\\
0.92	-0.00261806\\
0.922	-0.00254916\\
0.924	-0.00247187\\
0.926	-0.00238658\\
0.928	-0.0022937\\
0.93	-0.00219369\\
0.932	-0.00208699\\
0.934	-0.0019741\\
0.936	-0.0018555\\
0.938	-0.00173171\\
0.94	-0.00160324\\
0.942	-0.00147063\\
0.944	-0.00133442\\
0.946	-0.00119515\\
0.948	-0.00105338\\
0.95	-0.000909662\\
0.952	-0.000813971\\
0.954	-0.000807818\\
0.956	-0.000799225\\
0.958	-0.000788248\\
0.96	-0.000842538\\
0.962	-0.000932306\\
0.964	-0.0010185\\
0.966	-0.00110058\\
0.968	-0.00117832\\
0.97	-0.00125153\\
0.972	-0.00132003\\
0.974	-0.00138366\\
0.976	-0.00144225\\
0.978	-0.00149567\\
0.98	-0.00154382\\
0.982	-0.00158658\\
0.984	-0.00162389\\
0.986	-0.00165568\\
0.988	-0.00168191\\
0.99	-0.00170256\\
0.992	-0.00171761\\
0.994	-0.00172812\\
0.996	-0.00178735\\
0.998	-0.00183923\\
1	-0.00188364\\
1.002	-0.00192054\\
1.004	-0.00194987\\
1.006	-0.00197165\\
1.008	-0.00198587\\
1.01	-0.0019926\\
1.012	-0.00199189\\
1.014	-0.00198387\\
1.016	-0.00196863\\
1.018	-0.00194635\\
1.02	-0.00191719\\
1.022	-0.00188134\\
1.024	-0.00183903\\
1.026	-0.0017905\\
1.028	-0.001736\\
1.03	-0.00167581\\
1.032	-0.00161022\\
1.034	-0.00153956\\
1.036	-0.00146413\\
1.038	-0.00138427\\
1.04	-0.00130034\\
1.042	-0.00121271\\
1.044	-0.00112173\\
1.046	-0.0010278\\
1.048	-0.000931277\\
1.05	-0.000832567\\
1.052	-0.000732056\\
1.054	-0.000667258\\
1.056	-0.000647413\\
1.058	-0.000626432\\
1.06	-0.000604391\\
1.062	-0.000581368\\
1.064	-0.000559404\\
1.066	-0.000555846\\
1.068	-0.000550615\\
1.07	-0.000550253\\
1.072	-0.000619392\\
1.074	-0.000685929\\
1.076	-0.000749558\\
1.078	-0.00081011\\
1.08	-0.000867428\\
1.082	-0.000921365\\
1.084	-0.000971787\\
1.086	-0.00101857\\
1.088	-0.00106161\\
1.09	-0.00110081\\
1.092	-0.00113609\\
1.094	-0.00116738\\
1.096	-0.00119572\\
1.098	-0.00124279\\
1.1	-0.0012847\\
1.102	-0.00132134\\
1.104	-0.00135266\\
1.106	-0.0013786\\
1.108	-0.00139912\\
1.11	-0.00141424\\
1.112	-0.00142395\\
1.114	-0.00142829\\
1.116	-0.00142731\\
1.118	-0.00142108\\
1.12	-0.00140969\\
1.122	-0.00139324\\
1.124	-0.00137186\\
1.126	-0.00134568\\
1.128	-0.00131487\\
1.13	-0.00127958\\
1.132	-0.00124002\\
1.134	-0.00119636\\
1.136	-0.00114884\\
1.138	-0.00109766\\
1.14	-0.00104305\\
1.142	-0.000985277\\
1.144	-0.000924573\\
1.146	-0.000861201\\
1.148	-0.000795427\\
1.15	-0.000727521\\
1.152	-0.000657759\\
1.154	-0.000586419\\
1.156	-0.000513783\\
1.158	-0.000440133\\
1.16	-0.000398239\\
1.162	-0.000399019\\
1.164	-0.000398913\\
1.166	-0.000397943\\
1.168	-0.000396131\\
1.17	-0.0003935\\
1.172	-0.000390075\\
1.174	-0.00038588\\
1.176	-0.000380942\\
1.178	-0.000375288\\
1.18	-0.000368947\\
1.182	-0.000422633\\
1.184	-0.000485006\\
1.186	-0.000544853\\
1.188	-0.000601985\\
1.19	-0.000656225\\
1.192	-0.000707408\\
1.194	-0.000755384\\
1.196	-0.000800015\\
1.198	-0.000841178\\
1.2	-0.000878764\\
1.202	-0.000912677\\
1.204	-0.000942836\\
1.206	-0.000969176\\
1.208	-0.000991646\\
1.21	-0.00101021\\
1.212	-0.00102485\\
1.214	-0.00103555\\
1.216	-0.00104232\\
1.218	-0.00104519\\
1.22	-0.00104419\\
1.222	-0.00103937\\
1.224	-0.00103079\\
1.226	-0.00101853\\
1.228	-0.00100267\\
1.23	-0.000983326\\
1.232	-0.000960599\\
1.234	-0.000934613\\
1.236	-0.000905502\\
1.238	-0.00087341\\
1.24	-0.000838489\\
1.242	-0.000800899\\
1.244	-0.00076081\\
1.246	-0.000718397\\
1.248	-0.000673842\\
1.25	-0.000627334\\
1.252	-0.000585536\\
1.254	-0.000547763\\
1.256	-0.000508705\\
1.258	-0.000468501\\
1.26	-0.000427295\\
1.262	-0.000385229\\
1.264	-0.00034245\\
1.266	-0.000299102\\
1.268	-0.000255331\\
1.27	-0.000211283\\
1.272	-0.000197973\\
1.274	-0.000202742\\
1.276	-0.000206775\\
1.278	-0.000210071\\
1.28	-0.000212633\\
1.282	-0.000219025\\
1.284	-0.000268188\\
1.286	-0.000315811\\
1.288	-0.000361739\\
1.29	-0.000405825\\
1.292	-0.000447928\\
1.294	-0.000487918\\
1.296	-0.000525673\\
1.298	-0.000561082\\
1.3	-0.000594043\\
1.302	-0.000624464\\
1.304	-0.000652263\\
1.306	-0.000677369\\
1.308	-0.000699722\\
1.31	-0.000719272\\
1.312	-0.000735979\\
1.314	-0.000749816\\
1.316	-0.000760765\\
1.318	-0.00076882\\
1.32	-0.000773983\\
1.322	-0.00077627\\
1.324	-0.000775704\\
1.326	-0.000772322\\
1.328	-0.000766166\\
1.33	-0.000757291\\
1.332	-0.00074576\\
1.334	-0.000731646\\
1.336	-0.000715029\\
1.338	-0.000695998\\
1.34	-0.00067465\\
1.342	-0.000651088\\
1.344	-0.000625424\\
1.346	-0.000597774\\
1.348	-0.000579036\\
1.35	-0.000562387\\
1.352	-0.000544146\\
1.354	-0.000524394\\
1.356	-0.000503211\\
1.358	-0.000480684\\
1.36	-0.000456903\\
1.362	-0.00043196\\
1.364	-0.000405949\\
1.366	-0.000378969\\
1.368	-0.000351116\\
1.37	-0.000322494\\
1.372	-0.000293202\\
1.374	-0.000263344\\
1.376	-0.000233023\\
1.378	-0.000202343\\
1.38	-0.000171406\\
1.382	-0.000140316\\
1.384	-0.000120427\\
1.386	-0.000158135\\
1.388	-0.000194851\\
1.39	-0.000230453\\
1.392	-0.000264825\\
1.394	-0.000297856\\
1.396	-0.000329443\\
1.398	-0.000359487\\
1.4	-0.000387896\\
1.402	-0.000414587\\
1.404	-0.000439482\\
1.406	-0.000462511\\
1.408	-0.000483613\\
1.41	-0.000502731\\
1.412	-0.000519821\\
1.414	-0.000534842\\
1.416	-0.000547764\\
1.418	-0.000558564\\
1.42	-0.000567226\\
1.422	-0.000573743\\
1.424	-0.000578116\\
1.426	-0.000580353\\
1.428	-0.00058047\\
1.43	-0.00057849\\
1.432	-0.000574444\\
1.434	-0.000568369\\
1.436	-0.000560312\\
1.438	-0.000550321\\
1.44	-0.000538457\\
1.442	-0.000524781\\
1.444	-0.000509365\\
1.446	-0.000492282\\
1.448	-0.000473614\\
1.45	-0.000453445\\
1.452	-0.00043503\\
1.454	-0.000427765\\
1.456	-0.000419217\\
1.458	-0.000409429\\
1.46	-0.000398451\\
1.462	-0.000386331\\
1.464	-0.000373126\\
1.466	-0.00035889\\
1.468	-0.000343685\\
1.47	-0.000327572\\
1.472	-0.000310614\\
1.474	-0.000292878\\
1.476	-0.00027443\\
1.478	-0.000255341\\
1.48	-0.00023568\\
1.482	-0.000215518\\
1.484	-0.000194927\\
1.486	-0.000173979\\
1.488	-0.000152746\\
1.49	-0.000131302\\
1.492	-0.000135568\\
1.494	-0.000162433\\
1.496	-0.000188417\\
1.498	-0.000213436\\
1.5	-0.000237411\\
1.502	-0.000260268\\
1.504	-0.000281938\\
1.506	-0.000302354\\
1.508	-0.000321459\\
1.51	-0.000339199\\
1.512	-0.000355524\\
1.514	-0.000370392\\
1.516	-0.000383766\\
1.518	-0.000395614\\
1.52	-0.000405912\\
1.522	-0.000414639\\
1.524	-0.000421782\\
1.526	-0.000427333\\
1.528	-0.000431289\\
1.53	-0.000433655\\
1.532	-0.000434439\\
1.534	-0.000433657\\
1.536	-0.000431329\\
1.538	-0.000427481\\
1.54	-0.000422143\\
1.542	-0.000415351\\
1.544	-0.000407146\\
1.546	-0.000397574\\
1.548	-0.000386683\\
1.55	-0.000374527\\
1.552	-0.000361165\\
1.554	-0.000346658\\
1.556	-0.00033107\\
1.558	-0.00031447\\
1.56	-0.000305483\\
1.562	-0.000301784\\
1.564	-0.000297157\\
1.566	-0.000291629\\
1.568	-0.000285231\\
1.57	-0.000277995\\
1.572	-0.000269956\\
1.574	-0.00026115\\
1.576	-0.000251617\\
1.578	-0.000241397\\
1.58	-0.000230532\\
1.582	-0.000219067\\
1.584	-0.000207046\\
1.586	-0.000194517\\
1.588	-0.000181526\\
1.59	-0.000168123\\
1.592	-0.000154355\\
1.594	-0.000140274\\
1.596	-0.000125928\\
1.598	-0.000111368\\
1.6	-0.00012983\\
1.602	-0.000148865\\
1.604	-0.000167155\\
1.606	-0.000184644\\
1.608	-0.000201279\\
1.61	-0.000217009\\
1.612	-0.000231788\\
1.614	-0.000245573\\
1.616	-0.000258328\\
1.618	-0.000270018\\
1.62	-0.000280613\\
1.622	-0.000290088\\
1.624	-0.00029842\\
1.626	-0.000305594\\
1.628	-0.000311597\\
1.63	-0.00031642\\
1.632	-0.00032006\\
1.634	-0.000322516\\
1.636	-0.000323794\\
1.638	-0.000323902\\
1.64	-0.000322853\\
1.642	-0.000320663\\
1.644	-0.000317353\\
1.646	-0.000312948\\
1.648	-0.000307476\\
1.65	-0.000300968\\
1.652	-0.000293459\\
1.654	-0.000284988\\
1.656	-0.000275594\\
1.658	-0.000265323\\
1.66	-0.000254221\\
1.662	-0.000242337\\
1.664	-0.000229721\\
1.666	-0.000216428\\
1.668	-0.000211445\\
1.67	-0.000209785\\
1.672	-0.000207466\\
1.674	-0.000204505\\
1.676	-0.00020092\\
1.678	-0.000196731\\
1.68	-0.00019196\\
1.682	-0.000186631\\
1.684	-0.000180768\\
1.686	-0.000174398\\
1.688	-0.000167549\\
1.69	-0.000160249\\
1.692	-0.000152528\\
1.694	-0.000144418\\
1.696	-0.00013595\\
1.698	-0.000127156\\
1.7	-0.00011807\\
1.702	-0.000108725\\
1.704	-9.91551e-05\\
1.706	-0.000101096\\
1.708	-0.000115054\\
1.71	-0.000128445\\
1.712	-0.000141227\\
1.714	-0.000153363\\
1.716	-0.000164816\\
1.718	-0.000175553\\
1.72	-0.000185544\\
1.722	-0.000194761\\
1.724	-0.000203182\\
1.726	-0.000210784\\
1.728	-0.000217549\\
1.73	-0.000223464\\
1.732	-0.000228517\\
1.734	-0.0002327\\
1.736	-0.000236006\\
1.738	-0.000238436\\
1.74	-0.000239989\\
1.742	-0.000240671\\
1.744	-0.000240488\\
1.746	-0.000239452\\
1.748	-0.000237575\\
1.75	-0.000234874\\
1.752	-0.000231368\\
1.754	-0.000227078\\
1.756	-0.000222029\\
1.758	-0.000216246\\
1.76	-0.00020976\\
1.762	-0.0002026\\
1.764	-0.000194801\\
1.766	-0.000186397\\
1.768	-0.000177425\\
1.77	-0.000167923\\
1.772	-0.000157931\\
1.774	-0.00014749\\
1.776	-0.000144768\\
1.778	-0.000144234\\
1.78	-0.000143238\\
1.782	-0.00014179\\
1.784	-0.000139901\\
1.786	-0.000137583\\
1.788	-0.000134851\\
1.79	-0.000131718\\
1.792	-0.0001282\\
1.794	-0.000124315\\
1.796	-0.00012008\\
1.798	-0.000115513\\
1.8	-0.000110635\\
1.802	-0.000105465\\
1.804	-0.000100024\\
1.806	-9.43336e-05\\
1.808	-8.84161e-05\\
1.81	-8.22939e-05\\
1.812	-7.75143e-05\\
1.814	-8.77486e-05\\
1.816	-9.75551e-05\\
1.818	-0.000106904\\
1.82	-0.000115766\\
1.822	-0.000124117\\
1.824	-0.000131932\\
1.826	-0.00013919\\
1.828	-0.000145871\\
1.83	-0.000151957\\
1.832	-0.000157435\\
1.834	-0.000162291\\
1.836	-0.000166515\\
1.838	-0.0001701\\
1.84	-0.00017304\\
1.842	-0.000175332\\
1.844	-0.000176975\\
1.846	-0.000177971\\
1.848	-0.000178323\\
1.85	-0.000178038\\
1.852	-0.000177123\\
1.854	-0.00017559\\
1.856	-0.000173451\\
1.858	-0.00017072\\
1.86	-0.000167413\\
1.862	-0.00016355\\
1.864	-0.00015915\\
1.866	-0.000154234\\
1.868	-0.000148827\\
1.87	-0.000142952\\
1.872	-0.000136637\\
1.874	-0.000129908\\
1.876	-0.000122795\\
1.878	-0.000115326\\
1.88	-0.000107532\\
1.882	-9.94459e-05\\
1.884	-9.80944e-05\\
1.886	-9.81554e-05\\
1.888	-9.78972e-05\\
1.89	-9.73252e-05\\
1.892	-9.64456e-05\\
1.894	-9.52658e-05\\
1.896	-9.37937e-05\\
1.898	-9.20383e-05\\
1.9	-9.00094e-05\\
1.902	-8.77174e-05\\
1.904	-8.51734e-05\\
1.906	-8.23894e-05\\
1.908	-7.93778e-05\\
1.91	-7.61515e-05\\
1.912	-7.27242e-05\\
1.914	-6.91097e-05\\
1.916	-6.53225e-05\\
1.918	-6.13774e-05\\
1.92	-6.63672e-05\\
1.922	-7.35526e-05\\
1.924	-8.03954e-05\\
1.926	-8.68753e-05\\
1.928	-9.29733e-05\\
1.93	-9.86722e-05\\
1.932	-0.000103956\\
1.934	-0.000108812\\
1.936	-0.000113226\\
1.938	-0.000117188\\
1.94	-0.00012069\\
1.942	-0.000123724\\
1.944	-0.000126285\\
1.946	-0.000128369\\
1.948	-0.000129974\\
1.95	-0.0001311\\
1.952	-0.000131748\\
1.954	-0.000131922\\
1.956	-0.000131626\\
1.958	-0.000130867\\
1.96	-0.000129653\\
1.962	-0.000127993\\
1.964	-0.000125898\\
1.966	-0.00012338\\
1.968	-0.000120454\\
1.97	-0.000117135\\
1.972	-0.000113437\\
1.974	-0.00010938\\
1.976	-0.000104982\\
1.978	-0.000100262\\
1.98	-9.52395e-05\\
1.982	-8.99371e-05\\
1.984	-8.43764e-05\\
1.986	-7.858e-05\\
1.988	-7.25712e-05\\
1.99	-6.63736e-05\\
1.992	-6.55904e-05\\
1.994	-6.59447e-05\\
1.996	-6.60812e-05\\
1.998	-6.60026e-05\\
2	-6.57123e-05\\
};
\addplot [color=mycolor5, line width=1.0pt, forget plot]
  table[row sep=crcr]{%
-0.024	0\\
-0.022	0\\
-0.02	0\\
-0.018	0\\
-0.016	0\\
-0.014	0\\
-0.012	0\\
-0.01	0\\
-0.008	0\\
-0.006	0\\
-0.004	0\\
-0.002	0\\
0	0\\
0.002	9.94461e-09\\
0.004	5.45797e-08\\
0.006	1.72568e-07\\
0.008	2.80949e-07\\
0.01	3.21808e-07\\
0.012	2.01488e-07\\
0.014	-1.87637e-07\\
0.016	-9.60776e-07\\
0.018	3.88015e-07\\
0.02	-4.66587e-07\\
0.022	-2.70978e-07\\
0.024	-0.000964371\\
0.026	-0.00329026\\
0.028	-0.00518686\\
0.03	-0.00729784\\
0.032	-0.0096056\\
0.034	-0.012088\\
0.036	-0.0139619\\
0.038	-0.0149523\\
0.04	-0.0159044\\
0.042	-0.0168106\\
0.044	-0.0176641\\
0.046	-0.0184582\\
0.048	-0.0191868\\
0.05	-0.0198442\\
0.052	-0.0204255\\
0.054	-0.0209259\\
0.056	-0.0213413\\
0.058	-0.0216683\\
0.06	-0.0219036\\
0.062	-0.0220449\\
0.064	-0.0220901\\
0.066	-0.0220376\\
0.068	-0.0218867\\
0.07	-0.0216367\\
0.072	-0.0212878\\
0.074	-0.0208406\\
0.076	-0.020296\\
0.078	-0.0196555\\
0.08	-0.0189212\\
0.082	-0.0180955\\
0.084	-0.0171811\\
0.086	-0.0161814\\
0.088	-0.0151\\
0.09	-0.0139409\\
0.092	-0.0127084\\
0.094	-0.0114073\\
0.096	-0.0100425\\
0.098	-0.00861924\\
0.1	-0.00714315\\
0.102	-0.00561991\\
0.104	-0.00428259\\
0.106	-0.00475419\\
0.108	-0.00517657\\
0.11	-0.00554884\\
0.112	-0.00587037\\
0.114	-0.00614076\\
0.116	-0.00611391\\
0.118	-0.00571583\\
0.12	-0.00529502\\
0.122	-0.00527715\\
0.124	-0.00591625\\
0.126	-0.00661896\\
0.128	-0.00729559\\
0.13	-0.00794394\\
0.132	-0.00856192\\
0.134	-0.0088457\\
0.136	-0.0090204\\
0.138	-0.00917349\\
0.14	-0.00930399\\
0.142	-0.00991746\\
0.144	-0.0117915\\
0.146	-0.0135661\\
0.148	-0.0152365\\
0.15	-0.0167984\\
0.152	-0.0182482\\
0.154	-0.0195825\\
0.156	-0.0207985\\
0.158	-0.0218938\\
0.16	-0.0228665\\
0.162	-0.0237154\\
0.164	-0.0244395\\
0.166	-0.0250384\\
0.168	-0.025512\\
0.17	-0.0258611\\
0.172	-0.0260864\\
0.174	-0.0261893\\
0.176	-0.0261718\\
0.178	-0.0260361\\
0.18	-0.0257847\\
0.182	-0.0254207\\
0.184	-0.0249476\\
0.186	-0.024369\\
0.188	-0.0236891\\
0.19	-0.0228311\\
0.192	-0.0218399\\
0.194	-0.0207607\\
0.196	-0.019599\\
0.198	-0.0183607\\
0.2	-0.0170518\\
0.202	-0.0156785\\
0.204	-0.0142471\\
0.206	-0.0127639\\
0.208	-0.0112357\\
0.21	-0.00973144\\
0.212	-0.00984371\\
0.214	-0.00991063\\
0.216	-0.009933\\
0.218	-0.00991176\\
0.22	-0.00984793\\
0.222	-0.00974263\\
0.224	-0.00959711\\
0.226	-0.00941266\\
0.228	-0.0091907\\
0.23	-0.0089327\\
0.232	-0.00864023\\
0.234	-0.00831491\\
0.236	-0.00795845\\
0.238	-0.00757258\\
0.24	-0.00831641\\
0.242	-0.00979986\\
0.244	-0.0112185\\
0.246	-0.0125685\\
0.248	-0.0138462\\
0.25	-0.0150481\\
0.252	-0.0161713\\
0.254	-0.0172129\\
0.256	-0.0181707\\
0.258	-0.0190425\\
0.26	-0.0198266\\
0.262	-0.0205214\\
0.264	-0.021126\\
0.266	-0.0216395\\
0.268	-0.0220616\\
0.27	-0.022392\\
0.272	-0.0226311\\
0.274	-0.0227793\\
0.276	-0.0228375\\
0.278	-0.0228068\\
0.28	-0.0226887\\
0.282	-0.022485\\
0.284	-0.0221975\\
0.286	-0.0218287\\
0.288	-0.021381\\
0.29	-0.0208572\\
0.292	-0.0202604\\
0.294	-0.0195937\\
0.296	-0.0188607\\
0.298	-0.0180052\\
0.3	-0.0170946\\
0.302	-0.0161342\\
0.304	-0.0151282\\
0.306	-0.0140804\\
0.308	-0.0129952\\
0.31	-0.0118767\\
0.312	-0.0107292\\
0.314	-0.00955721\\
0.316	-0.00836501\\
0.318	-0.00823445\\
0.32	-0.00834717\\
0.322	-0.00843162\\
0.324	-0.0084878\\
0.326	-0.00851578\\
0.328	-0.00851574\\
0.33	-0.00848795\\
0.332	-0.00843275\\
0.334	-0.00835059\\
0.336	-0.00824198\\
0.338	-0.00810752\\
0.34	-0.00794789\\
0.342	-0.00776385\\
0.344	-0.00755621\\
0.346	-0.00732588\\
0.348	-0.00830507\\
0.35	-0.00927797\\
0.352	-0.0102029\\
0.354	-0.0110772\\
0.356	-0.0118984\\
0.358	-0.0126385\\
0.36	-0.013321\\
0.362	-0.0139458\\
0.364	-0.0145117\\
0.366	-0.0150173\\
0.368	-0.0154615\\
0.37	-0.0158436\\
0.372	-0.0161631\\
0.374	-0.0164196\\
0.376	-0.016613\\
0.378	-0.0167435\\
0.38	-0.0168113\\
0.382	-0.0168172\\
0.384	-0.0167619\\
0.386	-0.0166465\\
0.388	-0.016472\\
0.39	-0.01624\\
0.392	-0.0159521\\
0.394	-0.0156099\\
0.396	-0.0152156\\
0.398	-0.0147647\\
0.4	-0.0142533\\
0.402	-0.0136973\\
0.404	-0.0130993\\
0.406	-0.0124619\\
0.408	-0.011788\\
0.41	-0.0110799\\
0.412	-0.0103387\\
0.414	-0.00957002\\
0.416	-0.00934946\\
0.418	-0.00913648\\
0.42	-0.00890293\\
0.422	-0.00864967\\
0.424	-0.00837757\\
0.426	-0.00808756\\
0.428	-0.0077806\\
0.43	-0.00745768\\
0.432	-0.00711982\\
0.434	-0.00676808\\
0.436	-0.00664205\\
0.438	-0.00666193\\
0.44	-0.00665999\\
0.442	-0.00663646\\
0.444	-0.00659165\\
0.446	-0.00652591\\
0.448	-0.00643969\\
0.45	-0.00633346\\
0.452	-0.00620778\\
0.454	-0.00640089\\
0.456	-0.00706167\\
0.458	-0.0076881\\
0.46	-0.00827841\\
0.462	-0.00883097\\
0.464	-0.00934433\\
0.466	-0.00981718\\
0.468	-0.0102484\\
0.47	-0.010637\\
0.472	-0.0109822\\
0.474	-0.0112833\\
0.476	-0.0115398\\
0.478	-0.0117515\\
0.48	-0.0119182\\
0.482	-0.012036\\
0.484	-0.0121066\\
0.486	-0.0121329\\
0.488	-0.0121157\\
0.49	-0.0120556\\
0.492	-0.0119534\\
0.494	-0.0118099\\
0.496	-0.0116264\\
0.498	-0.0114038\\
0.5	-0.0111436\\
0.502	-0.0108472\\
0.504	-0.010516\\
0.506	-0.0101517\\
0.508	-0.00975603\\
0.51	-0.00933085\\
0.512	-0.00887803\\
0.514	-0.00839955\\
0.516	-0.00789626\\
0.518	-0.00737111\\
0.52	-0.00682653\\
0.522	-0.00681092\\
0.524	-0.00682759\\
0.526	-0.00682667\\
0.528	-0.00680842\\
0.53	-0.00677315\\
0.532	-0.0067212\\
0.534	-0.00665294\\
0.536	-0.00656879\\
0.538	-0.0064692\\
0.54	-0.00635464\\
0.542	-0.00622417\\
0.544	-0.00607894\\
0.546	-0.00592057\\
0.548	-0.00574964\\
0.55	-0.00556678\\
0.552	-0.00537261\\
0.554	-0.0051678\\
0.556	-0.00503253\\
0.558	-0.00497415\\
0.56	-0.00490044\\
0.562	-0.00512003\\
0.564	-0.00554946\\
0.566	-0.00595334\\
0.568	-0.00633054\\
0.57	-0.00668009\\
0.572	-0.00700111\\
0.574	-0.00729282\\
0.576	-0.00755458\\
0.578	-0.00778586\\
0.58	-0.00798624\\
0.582	-0.00815541\\
0.584	-0.00829319\\
0.586	-0.00839951\\
0.588	-0.00847441\\
0.59	-0.00851805\\
0.592	-0.00853068\\
0.594	-0.00851268\\
0.596	-0.00846454\\
0.598	-0.00838682\\
0.6	-0.00827966\\
0.602	-0.00814384\\
0.604	-0.0079808\\
0.606	-0.00779151\\
0.608	-0.00757699\\
0.61	-0.00733834\\
0.612	-0.00707675\\
0.614	-0.00679344\\
0.616	-0.00648974\\
0.618	-0.00616699\\
0.62	-0.0058266\\
0.622	-0.00547004\\
0.624	-0.00509879\\
0.626	-0.00471438\\
0.628	-0.00431836\\
0.63	-0.00391229\\
0.632	-0.00376262\\
0.634	-0.00386255\\
0.636	-0.00395121\\
0.638	-0.00402855\\
0.64	-0.00409454\\
0.642	-0.0041492\\
0.644	-0.00419256\\
0.646	-0.00422469\\
0.648	-0.00424568\\
0.65	-0.00425565\\
0.652	-0.00425474\\
0.654	-0.00424313\\
0.656	-0.00422101\\
0.658	-0.00418859\\
0.66	-0.00414613\\
0.662	-0.00409389\\
0.664	-0.00403147\\
0.666	-0.00395973\\
0.668	-0.00387916\\
0.67	-0.00394039\\
0.672	-0.00421698\\
0.674	-0.00447466\\
0.676	-0.00471276\\
0.678	-0.0049307\\
0.68	-0.00512798\\
0.682	-0.00530416\\
0.684	-0.00545891\\
0.686	-0.00559196\\
0.688	-0.00570314\\
0.69	-0.00579233\\
0.692	-0.00585953\\
0.694	-0.00590479\\
0.696	-0.00592824\\
0.698	-0.00593012\\
0.7	-0.00591069\\
0.702	-0.00587033\\
0.704	-0.00580946\\
0.706	-0.00572859\\
0.708	-0.00562827\\
0.71	-0.00550914\\
0.712	-0.00537188\\
0.714	-0.00521723\\
0.716	-0.00504599\\
0.718	-0.00485899\\
0.72	-0.00465714\\
0.722	-0.00444134\\
0.724	-0.00421258\\
0.726	-0.00397185\\
0.728	-0.00372018\\
0.73	-0.00345863\\
0.732	-0.00318827\\
0.734	-0.0029102\\
0.736	-0.00262549\\
0.738	-0.00233517\\
0.74	-0.00204048\\
0.742	-0.00193964\\
0.744	-0.0020365\\
0.746	-0.00216553\\
0.748	-0.00228647\\
0.75	-0.00239901\\
0.752	-0.00250288\\
0.754	-0.00259782\\
0.756	-0.00268362\\
0.758	-0.00276009\\
0.76	-0.00282709\\
0.762	-0.00288448\\
0.764	-0.00293217\\
0.766	-0.00297012\\
0.768	-0.00299828\\
0.77	-0.00301667\\
0.772	-0.00302532\\
0.774	-0.00302429\\
0.776	-0.00301368\\
0.778	-0.00299362\\
0.78	-0.0031643\\
0.782	-0.00332362\\
0.784	-0.00346864\\
0.786	-0.00359901\\
0.788	-0.00371455\\
0.79	-0.00381506\\
0.792	-0.00390037\\
0.794	-0.00397039\\
0.796	-0.00402507\\
0.798	-0.00406443\\
0.8	-0.00408853\\
0.802	-0.0040975\\
0.804	-0.00409152\\
0.806	-0.0040708\\
0.808	-0.00403562\\
0.81	-0.00398632\\
0.812	-0.00392327\\
0.814	-0.00384687\\
0.816	-0.0037576\\
0.818	-0.00365595\\
0.82	-0.00354246\\
0.822	-0.0034177\\
0.824	-0.00328228\\
0.826	-0.00313683\\
0.828	-0.00298203\\
0.83	-0.00281856\\
0.832	-0.00264713\\
0.834	-0.00246847\\
0.836	-0.00228332\\
0.838	-0.00209245\\
0.84	-0.00189661\\
0.842	-0.00169657\\
0.844	-0.00149315\\
0.846	-0.00128712\\
0.848	-0.00118097\\
0.85	-0.00124057\\
0.852	-0.00135473\\
0.854	-0.00146394\\
0.856	-0.00156751\\
0.858	-0.00166512\\
0.86	-0.00175651\\
0.862	-0.00184145\\
0.864	-0.00191974\\
0.866	-0.0019912\\
0.868	-0.00205566\\
0.87	-0.00211299\\
0.872	-0.00216308\\
0.874	-0.00220585\\
0.876	-0.00224122\\
0.878	-0.00226916\\
0.88	-0.00228966\\
0.882	-0.00230273\\
0.884	-0.00230839\\
0.886	-0.0023067\\
0.888	-0.0023727\\
0.89	-0.00246659\\
0.892	-0.00255018\\
0.894	-0.0026233\\
0.896	-0.00268583\\
0.898	-0.00273769\\
0.9	-0.00277883\\
0.902	-0.00280925\\
0.904	-0.00282898\\
0.906	-0.0028381\\
0.908	-0.00283671\\
0.91	-0.00282496\\
0.912	-0.00280304\\
0.914	-0.00277116\\
0.916	-0.00272958\\
0.918	-0.00267854\\
0.92	-0.0026184\\
0.922	-0.0025495\\
0.924	-0.00247219\\
0.926	-0.00238689\\
0.928	-0.00229401\\
0.93	-0.00219398\\
0.932	-0.00208728\\
0.934	-0.00197438\\
0.936	-0.00185577\\
0.938	-0.00173197\\
0.94	-0.0016035\\
0.942	-0.00147089\\
0.944	-0.00133468\\
0.946	-0.00119542\\
0.948	-0.00105365\\
0.95	-0.000909925\\
0.952	-0.000814228\\
0.954	-0.000808068\\
0.956	-0.000799467\\
0.958	-0.000788482\\
0.96	-0.000842764\\
0.962	-0.000932523\\
0.964	-0.00101871\\
0.966	-0.00110078\\
0.968	-0.00117853\\
0.97	-0.00125175\\
0.972	-0.00132025\\
0.974	-0.00138388\\
0.976	-0.00144247\\
0.978	-0.0014959\\
0.98	-0.00154404\\
0.982	-0.00158681\\
0.984	-0.00162411\\
0.986	-0.0016559\\
0.988	-0.00168213\\
0.99	-0.00170277\\
0.992	-0.00171781\\
0.994	-0.00172832\\
0.996	-0.00178755\\
0.998	-0.00183942\\
1	-0.00188383\\
1.002	-0.00192071\\
1.004	-0.00195004\\
1.006	-0.00197181\\
1.008	-0.00198603\\
1.01	-0.00199274\\
1.012	-0.00199203\\
1.014	-0.001984\\
1.016	-0.00196876\\
1.018	-0.00194647\\
1.02	-0.0019173\\
1.022	-0.00188144\\
1.024	-0.00183912\\
1.026	-0.00179058\\
1.028	-0.00173607\\
1.03	-0.00167587\\
1.032	-0.00161028\\
1.034	-0.0015396\\
1.036	-0.00146416\\
1.038	-0.00138431\\
1.04	-0.00130038\\
1.042	-0.00121274\\
1.044	-0.00112177\\
1.046	-0.00102783\\
1.048	-0.000931317\\
1.05	-0.000832609\\
1.052	-0.000732099\\
1.054	-0.000667302\\
1.056	-0.000647459\\
1.058	-0.000626479\\
1.06	-0.000604438\\
1.062	-0.000581416\\
1.064	-0.000559452\\
1.066	-0.000555893\\
1.068	-0.000550662\\
1.07	-0.000550301\\
1.072	-0.00061944\\
1.074	-0.000685976\\
1.076	-0.000749605\\
1.078	-0.000810159\\
1.08	-0.000867477\\
1.082	-0.000921415\\
1.084	-0.000971836\\
1.086	-0.00101862\\
1.088	-0.00106166\\
1.09	-0.00110086\\
1.092	-0.00113614\\
1.094	-0.00116743\\
1.096	-0.00119577\\
1.098	-0.00124285\\
1.1	-0.00128475\\
1.102	-0.0013214\\
1.104	-0.00135271\\
1.106	-0.00137865\\
1.108	-0.00139918\\
1.11	-0.00141429\\
1.112	-0.001424\\
1.114	-0.00142834\\
1.116	-0.00142736\\
1.118	-0.00142113\\
1.12	-0.00140974\\
1.122	-0.00139329\\
1.124	-0.0013719\\
1.126	-0.00134572\\
1.128	-0.00131491\\
1.13	-0.00127962\\
1.132	-0.00124005\\
1.134	-0.0011964\\
1.136	-0.00114887\\
1.138	-0.00109769\\
1.14	-0.00104308\\
1.142	-0.000985303\\
1.144	-0.000924597\\
1.146	-0.000861223\\
1.148	-0.000795447\\
1.15	-0.000727539\\
1.152	-0.000657775\\
1.154	-0.000586434\\
1.156	-0.000513796\\
1.158	-0.000440144\\
1.16	-0.000398248\\
1.162	-0.000399028\\
1.164	-0.000398923\\
1.166	-0.000397953\\
1.168	-0.000396142\\
1.17	-0.000393511\\
1.172	-0.000390086\\
1.174	-0.000385891\\
1.176	-0.000380954\\
1.178	-0.0003753\\
1.18	-0.000368959\\
1.182	-0.000422645\\
1.184	-0.000485019\\
1.186	-0.000544866\\
1.188	-0.000601997\\
1.19	-0.000656237\\
1.192	-0.00070742\\
1.194	-0.000755396\\
1.196	-0.000800027\\
1.198	-0.00084119\\
1.2	-0.000878775\\
1.202	-0.000912688\\
1.204	-0.000942846\\
1.206	-0.000969186\\
1.208	-0.000991656\\
1.21	-0.00101022\\
1.212	-0.00102486\\
1.214	-0.00103556\\
1.216	-0.00104233\\
1.218	-0.0010452\\
1.22	-0.0010442\\
1.222	-0.00103938\\
1.224	-0.0010308\\
1.226	-0.00101854\\
1.228	-0.00100269\\
1.23	-0.000983338\\
1.232	-0.00096061\\
1.234	-0.000934624\\
1.236	-0.000905513\\
1.238	-0.000873421\\
1.24	-0.0008385\\
1.242	-0.000800909\\
1.244	-0.00076082\\
1.246	-0.000718406\\
1.248	-0.000673851\\
1.25	-0.000627343\\
1.252	-0.000585544\\
1.254	-0.000547771\\
1.256	-0.000508712\\
1.258	-0.000468508\\
1.26	-0.000427301\\
1.262	-0.000385235\\
1.264	-0.000342455\\
1.266	-0.000299107\\
1.268	-0.000255336\\
1.27	-0.000211287\\
1.272	-0.000197977\\
1.274	-0.000202745\\
1.276	-0.000206777\\
1.278	-0.000210073\\
1.28	-0.000212635\\
1.282	-0.000219026\\
1.284	-0.000268189\\
1.286	-0.000315813\\
1.288	-0.000361741\\
1.29	-0.000405827\\
1.292	-0.00044793\\
1.294	-0.00048792\\
1.296	-0.000525675\\
1.298	-0.000561085\\
1.3	-0.000594046\\
1.302	-0.000624466\\
1.304	-0.000652265\\
1.306	-0.000677372\\
1.308	-0.000699725\\
1.31	-0.000719274\\
1.312	-0.000735982\\
1.314	-0.000749818\\
1.316	-0.000760767\\
1.318	-0.000768822\\
1.32	-0.000773985\\
1.322	-0.000776272\\
1.324	-0.000775706\\
1.326	-0.000772323\\
1.328	-0.000766168\\
1.33	-0.000757293\\
1.332	-0.000745762\\
1.334	-0.000731648\\
1.336	-0.000715031\\
1.338	-0.000696\\
1.34	-0.000674652\\
1.342	-0.00065109\\
1.344	-0.000625426\\
1.346	-0.000597776\\
1.348	-0.000579038\\
1.35	-0.000562389\\
1.352	-0.000544149\\
1.354	-0.000524396\\
1.356	-0.000503213\\
1.358	-0.000480687\\
1.36	-0.000456905\\
1.362	-0.000431962\\
1.364	-0.000405952\\
1.366	-0.000378971\\
1.368	-0.000351118\\
1.37	-0.000322496\\
1.372	-0.000293204\\
1.374	-0.000263346\\
1.376	-0.000233025\\
1.378	-0.000202345\\
1.38	-0.000171408\\
1.382	-0.000140317\\
1.384	-0.000120429\\
1.386	-0.000158136\\
1.388	-0.000194852\\
1.39	-0.000230454\\
1.392	-0.000264826\\
1.394	-0.000297857\\
1.396	-0.000329444\\
1.398	-0.000359487\\
1.4	-0.000387897\\
1.402	-0.000414588\\
1.404	-0.000439482\\
1.406	-0.000462512\\
1.408	-0.000483613\\
1.41	-0.000502732\\
1.412	-0.000519822\\
1.414	-0.000534843\\
1.416	-0.000547765\\
1.418	-0.000558564\\
1.42	-0.000567227\\
1.422	-0.000573744\\
1.424	-0.000578117\\
1.426	-0.000580353\\
1.428	-0.00058047\\
1.43	-0.00057849\\
1.432	-0.000574444\\
1.434	-0.00056837\\
1.436	-0.000560312\\
1.438	-0.000550322\\
1.44	-0.000538457\\
1.442	-0.000524782\\
1.444	-0.000509365\\
1.446	-0.000492283\\
1.448	-0.000473614\\
1.45	-0.000453445\\
1.452	-0.000435031\\
1.454	-0.000427766\\
1.456	-0.000419217\\
1.458	-0.00040943\\
1.46	-0.000398451\\
1.462	-0.000386332\\
1.464	-0.000373126\\
1.466	-0.000358891\\
1.468	-0.000343686\\
1.47	-0.000327572\\
1.472	-0.000310614\\
1.474	-0.000292878\\
1.476	-0.000274431\\
1.478	-0.000255342\\
1.48	-0.00023568\\
1.482	-0.000215518\\
1.484	-0.000194927\\
1.486	-0.000173979\\
1.488	-0.000152747\\
1.49	-0.000131302\\
1.492	-0.000135568\\
1.494	-0.000162433\\
1.496	-0.000188417\\
1.498	-0.000213436\\
1.5	-0.000237411\\
1.502	-0.000260269\\
1.504	-0.000281938\\
1.506	-0.000302355\\
1.508	-0.00032146\\
1.51	-0.000339199\\
1.512	-0.000355524\\
1.514	-0.000370392\\
1.516	-0.000383766\\
1.518	-0.000395614\\
1.52	-0.000405912\\
1.522	-0.000414639\\
1.524	-0.000421782\\
1.526	-0.000427333\\
1.528	-0.000431289\\
1.53	-0.000433655\\
1.532	-0.000434439\\
1.534	-0.000433658\\
1.536	-0.00043133\\
1.538	-0.000427481\\
1.54	-0.000422143\\
1.542	-0.000415351\\
1.544	-0.000407147\\
1.546	-0.000397574\\
1.548	-0.000386683\\
1.55	-0.000374527\\
1.552	-0.000361165\\
1.554	-0.000346658\\
1.556	-0.00033107\\
1.558	-0.00031447\\
1.56	-0.000305483\\
1.562	-0.000301784\\
1.564	-0.000297157\\
1.566	-0.000291629\\
1.568	-0.000285231\\
1.57	-0.000277995\\
1.572	-0.000269956\\
1.574	-0.00026115\\
1.576	-0.000251617\\
1.578	-0.000241397\\
1.58	-0.000230532\\
1.582	-0.000219067\\
1.584	-0.000207046\\
1.586	-0.000194517\\
1.588	-0.000181526\\
1.59	-0.000168123\\
1.592	-0.000154355\\
1.594	-0.000140274\\
1.596	-0.000125928\\
1.598	-0.000111368\\
1.6	-0.00012983\\
1.602	-0.000148865\\
1.604	-0.000167155\\
1.606	-0.000184645\\
1.608	-0.000201279\\
1.61	-0.000217009\\
1.612	-0.000231788\\
1.614	-0.000245574\\
1.616	-0.000258328\\
1.618	-0.000270018\\
1.62	-0.000280613\\
1.622	-0.000290088\\
1.624	-0.00029842\\
1.626	-0.000305594\\
1.628	-0.000311597\\
1.63	-0.00031642\\
1.632	-0.00032006\\
1.634	-0.000322517\\
1.636	-0.000323794\\
1.638	-0.000323902\\
1.64	-0.000322853\\
1.642	-0.000320663\\
1.644	-0.000317353\\
1.646	-0.000312948\\
1.648	-0.000307476\\
1.65	-0.000300968\\
1.652	-0.000293459\\
1.654	-0.000284988\\
1.656	-0.000275594\\
1.658	-0.000265323\\
1.66	-0.000254221\\
1.662	-0.000242337\\
1.664	-0.000229721\\
1.666	-0.000216428\\
1.668	-0.000211445\\
1.67	-0.000209785\\
1.672	-0.000207466\\
1.674	-0.000204505\\
1.676	-0.00020092\\
1.678	-0.000196731\\
1.68	-0.00019196\\
1.682	-0.000186631\\
1.684	-0.000180768\\
1.686	-0.000174398\\
1.688	-0.000167549\\
1.69	-0.000160249\\
1.692	-0.000152528\\
1.694	-0.000144418\\
1.696	-0.00013595\\
1.698	-0.000127156\\
1.7	-0.00011807\\
1.702	-0.000108725\\
1.704	-9.91551e-05\\
1.706	-0.000101096\\
1.708	-0.000115054\\
1.71	-0.000128445\\
1.712	-0.000141227\\
1.714	-0.000153363\\
1.716	-0.000164816\\
1.718	-0.000175553\\
1.72	-0.000185544\\
1.722	-0.000194761\\
1.724	-0.000203182\\
1.726	-0.000210784\\
1.728	-0.000217549\\
1.73	-0.000223464\\
1.732	-0.000228517\\
1.734	-0.0002327\\
1.736	-0.000236006\\
1.738	-0.000238436\\
1.74	-0.000239989\\
1.742	-0.000240671\\
1.744	-0.000240488\\
1.746	-0.000239452\\
1.748	-0.000237575\\
1.75	-0.000234874\\
1.752	-0.000231368\\
1.754	-0.000227078\\
1.756	-0.000222029\\
1.758	-0.000216246\\
1.76	-0.00020976\\
1.762	-0.0002026\\
1.764	-0.000194801\\
1.766	-0.000186397\\
1.768	-0.000177425\\
1.77	-0.000167923\\
1.772	-0.000157931\\
1.774	-0.00014749\\
1.776	-0.000144768\\
1.778	-0.000144234\\
1.78	-0.000143238\\
1.782	-0.00014179\\
1.784	-0.000139901\\
1.786	-0.000137583\\
1.788	-0.000134851\\
1.79	-0.000131718\\
1.792	-0.000128201\\
1.794	-0.000124315\\
1.796	-0.00012008\\
1.798	-0.000115513\\
1.8	-0.000110635\\
1.802	-0.000105465\\
1.804	-0.000100024\\
1.806	-9.43336e-05\\
1.808	-8.84161e-05\\
1.81	-8.22939e-05\\
1.812	-7.75143e-05\\
1.814	-8.77486e-05\\
1.816	-9.75551e-05\\
1.818	-0.000106904\\
1.82	-0.000115766\\
1.822	-0.000124117\\
1.824	-0.000131932\\
1.826	-0.00013919\\
1.828	-0.000145871\\
1.83	-0.000151957\\
1.832	-0.000157435\\
1.834	-0.000162291\\
1.836	-0.000166515\\
1.838	-0.0001701\\
1.84	-0.00017304\\
1.842	-0.000175332\\
1.844	-0.000176975\\
1.846	-0.000177971\\
1.848	-0.000178323\\
1.85	-0.000178038\\
1.852	-0.000177123\\
1.854	-0.00017559\\
1.856	-0.000173451\\
1.858	-0.00017072\\
1.86	-0.000167413\\
1.862	-0.00016355\\
1.864	-0.00015915\\
1.866	-0.000154234\\
1.868	-0.000148827\\
1.87	-0.000142952\\
1.872	-0.000136637\\
1.874	-0.000129908\\
1.876	-0.000122795\\
1.878	-0.000115326\\
1.88	-0.000107532\\
1.882	-9.94459e-05\\
1.884	-9.80944e-05\\
1.886	-9.81554e-05\\
1.888	-9.78972e-05\\
1.89	-9.73252e-05\\
1.892	-9.64456e-05\\
1.894	-9.52658e-05\\
1.896	-9.37937e-05\\
1.898	-9.20383e-05\\
1.9	-9.00094e-05\\
1.902	-8.77174e-05\\
1.904	-8.51735e-05\\
1.906	-8.23894e-05\\
1.908	-7.93778e-05\\
1.91	-7.61515e-05\\
1.912	-7.27242e-05\\
1.914	-6.91097e-05\\
1.916	-6.53225e-05\\
1.918	-6.13774e-05\\
1.92	-6.63672e-05\\
1.922	-7.35526e-05\\
1.924	-8.03954e-05\\
1.926	-8.68753e-05\\
1.928	-9.29733e-05\\
1.93	-9.86722e-05\\
1.932	-0.000103956\\
1.934	-0.000108812\\
1.936	-0.000113226\\
1.938	-0.000117188\\
1.94	-0.00012069\\
1.942	-0.000123724\\
1.944	-0.000126285\\
1.946	-0.000128369\\
1.948	-0.000129974\\
1.95	-0.0001311\\
1.952	-0.000131748\\
1.954	-0.000131922\\
1.956	-0.000131626\\
1.958	-0.000130867\\
1.96	-0.000129653\\
1.962	-0.000127993\\
1.964	-0.000125898\\
1.966	-0.00012338\\
1.968	-0.000120454\\
1.97	-0.000117135\\
1.972	-0.000113437\\
1.974	-0.00010938\\
1.976	-0.000104982\\
1.978	-0.000100262\\
1.98	-9.52395e-05\\
1.982	-8.99371e-05\\
1.984	-8.43764e-05\\
1.986	-7.858e-05\\
1.988	-7.25712e-05\\
1.99	-6.63736e-05\\
1.992	-6.55904e-05\\
1.994	-6.59447e-05\\
1.996	-6.60812e-05\\
1.998	-6.60026e-05\\
2	-6.57123e-05\\
};
\addplot [color=mycolor6, line width=1.0pt, forget plot]
  table[row sep=crcr]{%
-0.024	0\\
-0.022	0\\
-0.02	0\\
-0.018	0\\
-0.016	0\\
-0.014	0\\
-0.012	0\\
-0.01	0\\
-0.008	0\\
-0.006	0\\
-0.004	0\\
-0.002	0\\
0	0\\
0.002	-2.8883e-08\\
0.004	-3.88405e-07\\
0.006	-2.07896e-06\\
0.008	-6.03621e-06\\
0.01	-1.35816e-05\\
0.012	-2.59806e-05\\
0.014	-4.44021e-05\\
0.016	-6.98382e-05\\
0.018	-6.04909e-05\\
0.02	-8.64363e-05\\
0.022	-0.000106878\\
0.024	-0.00110519\\
0.026	-0.00344324\\
0.028	-0.00538033\\
0.03	-0.00753745\\
0.032	-0.00989684\\
0.034	-0.012436\\
0.036	-0.0135021\\
0.038	-0.0144808\\
0.04	-0.0154271\\
0.042	-0.0163339\\
0.044	-0.0171944\\
0.046	-0.0180022\\
0.048	-0.0187514\\
0.05	-0.0194364\\
0.052	-0.0200518\\
0.054	-0.0205929\\
0.056	-0.0210554\\
0.058	-0.0214351\\
0.06	-0.0217287\\
0.062	-0.0219331\\
0.064	-0.0220456\\
0.066	-0.0220641\\
0.068	-0.0219871\\
0.07	-0.0218133\\
0.072	-0.0215421\\
0.074	-0.0211734\\
0.076	-0.0207073\\
0.078	-0.0201447\\
0.08	-0.0194869\\
0.082	-0.0187356\\
0.084	-0.0178931\\
0.086	-0.0169619\\
0.088	-0.0159452\\
0.09	-0.0148465\\
0.092	-0.0136697\\
0.094	-0.0124191\\
0.096	-0.0110995\\
0.098	-0.00971569\\
0.1	-0.00827313\\
0.102	-0.00677741\\
0.104	-0.00546148\\
0.106	-0.00594836\\
0.108	-0.0063799\\
0.11	-0.00675532\\
0.112	-0.00707411\\
0.114	-0.00733604\\
0.116	-0.00722986\\
0.118	-0.00682076\\
0.12	-0.00638466\\
0.122	-0.00634747\\
0.124	-0.00696347\\
0.126	-0.0076396\\
0.128	-0.00828643\\
0.13	-0.00890206\\
0.132	-0.00948471\\
0.134	-0.00980893\\
0.136	-0.00997598\\
0.138	-0.0101085\\
0.14	-0.0102162\\
0.142	-0.0108049\\
0.144	-0.0126524\\
0.146	-0.0143988\\
0.148	-0.0160396\\
0.15	-0.0175706\\
0.152	-0.0189883\\
0.154	-0.0202895\\
0.156	-0.0214715\\
0.158	-0.0225322\\
0.16	-0.0234697\\
0.162	-0.0242829\\
0.164	-0.024971\\
0.166	-0.0255336\\
0.168	-0.0259709\\
0.17	-0.0262835\\
0.172	-0.0264725\\
0.174	-0.0265392\\
0.176	-0.0264857\\
0.178	-0.0263143\\
0.18	-0.0260276\\
0.182	-0.0256287\\
0.184	-0.0251212\\
0.186	-0.0244479\\
0.188	-0.0236066\\
0.19	-0.0226685\\
0.192	-0.0216386\\
0.194	-0.0205225\\
0.196	-0.0193257\\
0.198	-0.0180539\\
0.2	-0.0167133\\
0.202	-0.0153099\\
0.204	-0.01385\\
0.206	-0.0123401\\
0.208	-0.0107866\\
0.21	-0.00925864\\
0.212	-0.00934865\\
0.214	-0.00939472\\
0.216	-0.0093976\\
0.218	-0.00935816\\
0.22	-0.00927736\\
0.222	-0.00915629\\
0.224	-0.0089961\\
0.226	-0.00879805\\
0.228	-0.00856348\\
0.23	-0.00829383\\
0.232	-0.0079906\\
0.234	-0.00765537\\
0.236	-0.00728978\\
0.238	-0.00689555\\
0.24	-0.00763171\\
0.242	-0.00910816\\
0.244	-0.0105205\\
0.246	-0.0118647\\
0.248	-0.0131371\\
0.25	-0.0143344\\
0.252	-0.0154534\\
0.254	-0.0164914\\
0.256	-0.0174459\\
0.258	-0.018315\\
0.26	-0.0190967\\
0.262	-0.0197896\\
0.264	-0.0203927\\
0.266	-0.0209051\\
0.268	-0.0213264\\
0.27	-0.0216565\\
0.272	-0.0218956\\
0.274	-0.0220442\\
0.276	-0.0221031\\
0.278	-0.0220736\\
0.28	-0.0219569\\
0.282	-0.0217549\\
0.284	-0.0214695\\
0.286	-0.0211031\\
0.288	-0.0206582\\
0.29	-0.0201375\\
0.292	-0.0195441\\
0.294	-0.0188812\\
0.296	-0.0181522\\
0.298	-0.0173607\\
0.3	-0.0165106\\
0.302	-0.0156057\\
0.304	-0.0146503\\
0.306	-0.0136484\\
0.308	-0.0126045\\
0.31	-0.0115228\\
0.312	-0.0104081\\
0.314	-0.00926465\\
0.316	-0.00809719\\
0.318	-0.0079877\\
0.32	-0.00811798\\
0.322	-0.00821669\\
0.324	-0.008284\\
0.326	-0.00832017\\
0.328	-0.00832556\\
0.33	-0.00830062\\
0.332	-0.00824588\\
0.334	-0.00816196\\
0.336	-0.00804954\\
0.338	-0.00790941\\
0.34	-0.0077424\\
0.342	-0.00754944\\
0.344	-0.00733151\\
0.346	-0.00708965\\
0.348	-0.00805624\\
0.35	-0.00901562\\
0.352	-0.00992622\\
0.354	-0.0107855\\
0.356	-0.0115912\\
0.358	-0.0123413\\
0.36	-0.0130338\\
0.362	-0.0136671\\
0.364	-0.0142398\\
0.366	-0.0147508\\
0.368	-0.0151992\\
0.37	-0.0155842\\
0.372	-0.0159053\\
0.374	-0.0161624\\
0.376	-0.0163554\\
0.378	-0.0164846\\
0.38	-0.0165504\\
0.382	-0.0165534\\
0.384	-0.0164947\\
0.386	-0.0163751\\
0.388	-0.0161961\\
0.39	-0.015959\\
0.392	-0.0156657\\
0.394	-0.0153178\\
0.396	-0.0149175\\
0.398	-0.0144669\\
0.4	-0.0139683\\
0.402	-0.0134242\\
0.404	-0.0128371\\
0.406	-0.0122098\\
0.408	-0.0115451\\
0.41	-0.0108459\\
0.412	-0.0101152\\
0.414	-0.00935604\\
0.416	-0.00914412\\
0.418	-0.00893891\\
0.42	-0.00871232\\
0.422	-0.00846524\\
0.424	-0.00819859\\
0.426	-0.00791333\\
0.428	-0.00761047\\
0.43	-0.00729104\\
0.432	-0.00695609\\
0.434	-0.00660672\\
0.436	-0.00648257\\
0.438	-0.00650386\\
0.44	-0.00650291\\
0.442	-0.00647997\\
0.444	-0.0064354\\
0.446	-0.00636958\\
0.448	-0.00628298\\
0.45	-0.00617611\\
0.452	-0.00604956\\
0.454	-0.00624159\\
0.456	-0.00690113\\
0.458	-0.00752616\\
0.46	-0.00811494\\
0.462	-0.00866587\\
0.464	-0.00917753\\
0.466	-0.00964863\\
0.468	-0.0100781\\
0.47	-0.0104648\\
0.472	-0.0108082\\
0.474	-0.0111075\\
0.476	-0.0113623\\
0.478	-0.0115723\\
0.48	-0.0117374\\
0.482	-0.0118575\\
0.484	-0.0119329\\
0.486	-0.011964\\
0.488	-0.0119511\\
0.49	-0.011895\\
0.492	-0.0117964\\
0.494	-0.0116563\\
0.496	-0.0114758\\
0.498	-0.0112559\\
0.5	-0.0109982\\
0.502	-0.0107039\\
0.504	-0.0103747\\
0.506	-0.0100121\\
0.508	-0.00961806\\
0.51	-0.00919424\\
0.512	-0.00874263\\
0.514	-0.00826521\\
0.516	-0.00776404\\
0.518	-0.00724123\\
0.52	-0.00669894\\
0.522	-0.00668559\\
0.524	-0.00670449\\
0.526	-0.00670576\\
0.528	-0.00668967\\
0.53	-0.00665652\\
0.532	-0.00660665\\
0.534	-0.00654044\\
0.536	-0.00645831\\
0.538	-0.0063607\\
0.54	-0.00624811\\
0.542	-0.00612103\\
0.544	-0.00598002\\
0.546	-0.00582564\\
0.548	-0.0056585\\
0.55	-0.00547921\\
0.552	-0.00528841\\
0.554	-0.00508676\\
0.556	-0.00495447\\
0.558	-0.00489888\\
0.56	-0.00482777\\
0.562	-0.00504981\\
0.564	-0.00548152\\
0.566	-0.00588751\\
0.568	-0.00626668\\
0.57	-0.00661806\\
0.572	-0.00694076\\
0.574	-0.00723403\\
0.576	-0.00749723\\
0.578	-0.00772982\\
0.58	-0.00793141\\
0.582	-0.00810168\\
0.584	-0.00824047\\
0.586	-0.00834771\\
0.588	-0.00842345\\
0.59	-0.00846784\\
0.592	-0.00848117\\
0.594	-0.00846379\\
0.596	-0.00841621\\
0.598	-0.00833899\\
0.6	-0.00823283\\
0.602	-0.0080985\\
0.604	-0.00793685\\
0.606	-0.00774885\\
0.608	-0.00753553\\
0.61	-0.007298\\
0.612	-0.00703743\\
0.614	-0.00675508\\
0.616	-0.00645225\\
0.618	-0.00613031\\
0.62	-0.00579067\\
0.622	-0.00543478\\
0.624	-0.00506415\\
0.626	-0.00468031\\
0.628	-0.0042848\\
0.63	-0.0038792\\
0.632	-0.00372995\\
0.634	-0.00383027\\
0.636	-0.00391927\\
0.638	-0.00399692\\
0.64	-0.0040632\\
0.642	-0.00411812\\
0.644	-0.00416172\\
0.646	-0.00419406\\
0.648	-0.00421524\\
0.65	-0.00422539\\
0.652	-0.00422464\\
0.654	-0.00421318\\
0.656	-0.0041912\\
0.658	-0.00415892\\
0.66	-0.00411659\\
0.662	-0.00406447\\
0.664	-0.00400285\\
0.666	-0.00393203\\
0.668	-0.00385234\\
0.67	-0.00391441\\
0.672	-0.00419179\\
0.674	-0.00445023\\
0.676	-0.00468906\\
0.678	-0.00490769\\
0.68	-0.00510562\\
0.682	-0.00528242\\
0.684	-0.00543775\\
0.686	-0.00557136\\
0.688	-0.00568306\\
0.69	-0.00577275\\
0.692	-0.00584042\\
0.694	-0.00588612\\
0.696	-0.00591\\
0.698	-0.00591227\\
0.7	-0.00589322\\
0.702	-0.00585322\\
0.704	-0.00579269\\
0.706	-0.00571213\\
0.708	-0.00561212\\
0.71	-0.00549328\\
0.712	-0.00535629\\
0.714	-0.0052019\\
0.716	-0.0050309\\
0.718	-0.00484414\\
0.72	-0.0046425\\
0.722	-0.00442692\\
0.724	-0.00419837\\
0.726	-0.00395783\\
0.728	-0.00370635\\
0.73	-0.00344498\\
0.732	-0.00317479\\
0.734	-0.00289689\\
0.736	-0.00261237\\
0.738	-0.00232235\\
0.74	-0.00202796\\
0.742	-0.00192739\\
0.744	-0.00202451\\
0.746	-0.00215378\\
0.748	-0.00227495\\
0.75	-0.00238771\\
0.752	-0.00249178\\
0.754	-0.00258691\\
0.756	-0.00267289\\
0.758	-0.00274954\\
0.76	-0.0028167\\
0.762	-0.00287424\\
0.764	-0.00292208\\
0.766	-0.00296016\\
0.768	-0.00298846\\
0.77	-0.00300697\\
0.772	-0.00301574\\
0.774	-0.00301483\\
0.776	-0.00300433\\
0.778	-0.00298438\\
0.78	-0.00315516\\
0.782	-0.00331458\\
0.784	-0.00345977\\
0.786	-0.00359039\\
0.788	-0.00370618\\
0.79	-0.00380692\\
0.792	-0.00389246\\
0.794	-0.0039627\\
0.796	-0.00401759\\
0.798	-0.00405714\\
0.8	-0.00408144\\
0.802	-0.00409059\\
0.804	-0.00408478\\
0.806	-0.00406422\\
0.808	-0.00402921\\
0.81	-0.00398006\\
0.812	-0.00391715\\
0.814	-0.0038409\\
0.816	-0.00375176\\
0.818	-0.00365024\\
0.82	-0.00353687\\
0.822	-0.00341222\\
0.824	-0.00327691\\
0.826	-0.00313158\\
0.828	-0.00297687\\
0.83	-0.0028135\\
0.832	-0.00264216\\
0.834	-0.00246359\\
0.836	-0.00227853\\
0.838	-0.00208775\\
0.84	-0.00189202\\
0.842	-0.00169211\\
0.844	-0.00148881\\
0.846	-0.0012829\\
0.848	-0.00117686\\
0.85	-0.00123657\\
0.852	-0.00135083\\
0.854	-0.00146014\\
0.856	-0.0015638\\
0.858	-0.0016615\\
0.86	-0.00175297\\
0.862	-0.001838\\
0.864	-0.00191636\\
0.866	-0.00198789\\
0.868	-0.00205242\\
0.87	-0.00210982\\
0.872	-0.00215998\\
0.874	-0.0022028\\
0.876	-0.00223824\\
0.878	-0.00226624\\
0.88	-0.00228679\\
0.882	-0.00229991\\
0.884	-0.00230562\\
0.886	-0.00230398\\
0.888	-0.00237002\\
0.89	-0.00246395\\
0.892	-0.00254758\\
0.894	-0.00262074\\
0.896	-0.00268332\\
0.898	-0.00273521\\
0.9	-0.00277639\\
0.902	-0.00280684\\
0.904	-0.00282661\\
0.906	-0.00283576\\
0.908	-0.0028344\\
0.91	-0.00282269\\
0.912	-0.0028008\\
0.914	-0.00276895\\
0.916	-0.0027274\\
0.918	-0.00267642\\
0.92	-0.00261633\\
0.922	-0.00254748\\
0.924	-0.00247022\\
0.926	-0.00238497\\
0.928	-0.00229213\\
0.93	-0.00219215\\
0.932	-0.00208549\\
0.934	-0.00197263\\
0.936	-0.00185406\\
0.938	-0.0017303\\
0.94	-0.00160186\\
0.942	-0.00146929\\
0.944	-0.00133311\\
0.946	-0.00119388\\
0.948	-0.00105215\\
0.95	-0.000908461\\
0.952	-0.000812802\\
0.954	-0.000806679\\
0.956	-0.000798113\\
0.958	-0.000787162\\
0.96	-0.000841477\\
0.962	-0.000931268\\
0.964	-0.00101748\\
0.966	-0.00109959\\
0.968	-0.00117736\\
0.97	-0.00125061\\
0.972	-0.00131914\\
0.974	-0.00138279\\
0.976	-0.00144141\\
0.978	-0.00149486\\
0.98	-0.00154303\\
0.982	-0.00158582\\
0.984	-0.00162315\\
0.986	-0.00165496\\
0.988	-0.0016812\\
0.99	-0.00170186\\
0.992	-0.00171693\\
0.994	-0.00172745\\
0.996	-0.0017867\\
0.998	-0.00183858\\
1	-0.00188301\\
1.002	-0.00191991\\
1.004	-0.00194926\\
1.006	-0.00197104\\
1.008	-0.00198527\\
1.01	-0.001992\\
1.012	-0.0019913\\
1.014	-0.00198328\\
1.016	-0.00196805\\
1.018	-0.00194577\\
1.02	-0.00191661\\
1.022	-0.00188077\\
1.024	-0.00183847\\
1.026	-0.00178993\\
1.028	-0.00173544\\
1.03	-0.00167525\\
1.032	-0.00160967\\
1.034	-0.001539\\
1.036	-0.00146358\\
1.038	-0.00138374\\
1.04	-0.00129982\\
1.042	-0.0012122\\
1.044	-0.00112124\\
1.046	-0.00102731\\
1.048	-0.000930809\\
1.05	-0.000832112\\
1.052	-0.000731613\\
1.054	-0.000666826\\
1.056	-0.000646993\\
1.058	-0.000626022\\
1.06	-0.000603991\\
1.062	-0.000580977\\
1.064	-0.000559022\\
1.066	-0.000555472\\
1.068	-0.000550249\\
1.07	-0.000549895\\
1.072	-0.000619042\\
1.074	-0.000685585\\
1.076	-0.000749222\\
1.078	-0.000809783\\
1.08	-0.000867108\\
1.082	-0.000921052\\
1.084	-0.000971481\\
1.086	-0.00101827\\
1.088	-0.00106132\\
1.09	-0.00110053\\
1.092	-0.00113582\\
1.094	-0.00116711\\
1.096	-0.00119546\\
1.098	-0.00124255\\
1.1	-0.00128446\\
1.102	-0.00132111\\
1.104	-0.00135243\\
1.106	-0.00137837\\
1.108	-0.00139891\\
1.11	-0.00141403\\
1.112	-0.00142374\\
1.114	-0.00142809\\
1.116	-0.00142712\\
1.118	-0.00142089\\
1.12	-0.0014095\\
1.122	-0.00139306\\
1.124	-0.00137168\\
1.126	-0.0013455\\
1.128	-0.00131469\\
1.13	-0.00127941\\
1.132	-0.00123984\\
1.134	-0.00119619\\
1.136	-0.00114867\\
1.138	-0.00109749\\
1.14	-0.00104289\\
1.142	-0.000985115\\
1.144	-0.000924412\\
1.146	-0.000861042\\
1.148	-0.000795269\\
1.15	-0.000727365\\
1.152	-0.000657605\\
1.154	-0.000586266\\
1.156	-0.000513632\\
1.158	-0.000439983\\
1.16	-0.00039809\\
1.162	-0.000398873\\
1.164	-0.000398772\\
1.166	-0.000397806\\
1.168	-0.000395997\\
1.17	-0.00039337\\
1.172	-0.000389948\\
1.174	-0.000385756\\
1.176	-0.000380821\\
1.178	-0.00037517\\
1.18	-0.000368832\\
1.182	-0.000422521\\
1.184	-0.000484897\\
1.186	-0.000544746\\
1.188	-0.00060188\\
1.19	-0.000656122\\
1.192	-0.000707308\\
1.194	-0.000755286\\
1.196	-0.000799919\\
1.198	-0.000841084\\
1.2	-0.000878671\\
1.202	-0.000912586\\
1.204	-0.000942747\\
1.206	-0.000969089\\
1.208	-0.000991561\\
1.21	-0.00101013\\
1.212	-0.00102477\\
1.214	-0.00103547\\
1.216	-0.00104225\\
1.218	-0.00104512\\
1.22	-0.00104412\\
1.222	-0.0010393\\
1.224	-0.00103072\\
1.226	-0.00101846\\
1.228	-0.00100261\\
1.23	-0.000983264\\
1.232	-0.000960538\\
1.234	-0.000934553\\
1.236	-0.000905444\\
1.238	-0.000873353\\
1.24	-0.000838433\\
1.242	-0.000800845\\
1.244	-0.000760756\\
1.246	-0.000718345\\
1.248	-0.000673791\\
1.25	-0.000627284\\
1.252	-0.000585487\\
1.254	-0.000547715\\
1.256	-0.000508657\\
1.258	-0.000468454\\
1.26	-0.000427248\\
1.262	-0.000385183\\
1.264	-0.000342405\\
1.266	-0.000299057\\
1.268	-0.000255287\\
1.27	-0.00021124\\
1.272	-0.000197931\\
1.274	-0.0002027\\
1.276	-0.000206733\\
1.278	-0.00021003\\
1.28	-0.000212593\\
1.282	-0.000218985\\
1.284	-0.000268149\\
1.286	-0.000315773\\
1.288	-0.000361702\\
1.29	-0.000405789\\
1.292	-0.000447892\\
1.294	-0.000487883\\
1.296	-0.000525639\\
1.298	-0.000561049\\
1.3	-0.000594011\\
1.302	-0.000624433\\
1.304	-0.000652232\\
1.306	-0.000677339\\
1.308	-0.000699693\\
1.31	-0.000719243\\
1.312	-0.000735951\\
1.314	-0.000749789\\
1.316	-0.000760738\\
1.318	-0.000768793\\
1.32	-0.000773957\\
1.322	-0.000776244\\
1.324	-0.000775679\\
1.326	-0.000772297\\
1.328	-0.000766142\\
1.33	-0.000757268\\
1.332	-0.000745738\\
1.334	-0.000731624\\
1.336	-0.000715008\\
1.338	-0.000695977\\
1.34	-0.000674629\\
1.342	-0.000651068\\
1.344	-0.000625404\\
1.346	-0.000597755\\
1.348	-0.000579018\\
1.35	-0.000562369\\
1.352	-0.000544129\\
1.354	-0.000524376\\
1.356	-0.000503194\\
1.358	-0.000480668\\
1.36	-0.000456887\\
1.362	-0.000431944\\
1.364	-0.000405934\\
1.366	-0.000378954\\
1.368	-0.000351102\\
1.37	-0.000322479\\
1.372	-0.000293188\\
1.374	-0.00026333\\
1.376	-0.00023301\\
1.378	-0.000202329\\
1.38	-0.000171393\\
1.382	-0.000140303\\
1.384	-0.000120414\\
1.386	-0.000158122\\
1.388	-0.000194838\\
1.39	-0.00023044\\
1.392	-0.000264813\\
1.394	-0.000297845\\
1.396	-0.000329431\\
1.398	-0.000359475\\
1.4	-0.000387885\\
1.402	-0.000414576\\
1.404	-0.000439471\\
1.406	-0.000462501\\
1.408	-0.000483602\\
1.41	-0.000502721\\
1.412	-0.000519811\\
1.414	-0.000534833\\
1.416	-0.000547755\\
1.418	-0.000558555\\
1.42	-0.000567217\\
1.422	-0.000573734\\
1.424	-0.000578107\\
1.426	-0.000580344\\
1.428	-0.000580461\\
1.43	-0.000578482\\
1.432	-0.000574436\\
1.434	-0.000568362\\
1.436	-0.000560304\\
1.438	-0.000550314\\
1.44	-0.000538449\\
1.442	-0.000524774\\
1.444	-0.000509358\\
1.446	-0.000492275\\
1.448	-0.000473607\\
1.45	-0.000453439\\
1.452	-0.000435024\\
1.454	-0.000427759\\
1.456	-0.000419211\\
1.458	-0.000409423\\
1.46	-0.000398445\\
1.462	-0.000386326\\
1.464	-0.00037312\\
1.466	-0.000358885\\
1.468	-0.00034368\\
1.47	-0.000327567\\
1.472	-0.000310609\\
1.474	-0.000292873\\
1.476	-0.000274426\\
1.478	-0.000255336\\
1.48	-0.000235675\\
1.482	-0.000215513\\
1.484	-0.000194922\\
1.486	-0.000173975\\
1.488	-0.000152742\\
1.49	-0.000131298\\
1.492	-0.000135564\\
1.494	-0.000162429\\
1.496	-0.000188413\\
1.498	-0.000213432\\
1.5	-0.000237407\\
1.502	-0.000260265\\
1.504	-0.000281934\\
1.506	-0.000302351\\
1.508	-0.000321456\\
1.51	-0.000339195\\
1.512	-0.00035552\\
1.514	-0.000370389\\
1.516	-0.000383763\\
1.518	-0.000395611\\
1.52	-0.000405909\\
1.522	-0.000414636\\
1.524	-0.000421779\\
1.526	-0.00042733\\
1.528	-0.000431286\\
1.53	-0.000433652\\
1.532	-0.000434437\\
1.534	-0.000433655\\
1.536	-0.000431327\\
1.538	-0.000427478\\
1.54	-0.000422141\\
1.542	-0.000415349\\
1.544	-0.000407144\\
1.546	-0.000397571\\
1.548	-0.00038668\\
1.55	-0.000374525\\
1.552	-0.000361163\\
1.554	-0.000346656\\
1.556	-0.000331068\\
1.558	-0.000314467\\
1.56	-0.000305481\\
1.562	-0.000301782\\
1.564	-0.000297155\\
1.566	-0.000291627\\
1.568	-0.000285229\\
1.57	-0.000277993\\
1.572	-0.000269954\\
1.574	-0.000261149\\
1.576	-0.000251615\\
1.578	-0.000241395\\
1.58	-0.00023053\\
1.582	-0.000219065\\
1.584	-0.000207045\\
1.586	-0.000194515\\
1.588	-0.000181525\\
1.59	-0.000168121\\
1.592	-0.000154354\\
1.594	-0.000140272\\
1.596	-0.000125926\\
1.598	-0.000111366\\
1.6	-0.000129829\\
1.602	-0.000148864\\
1.604	-0.000167154\\
1.606	-0.000184643\\
1.608	-0.000201278\\
1.61	-0.000217007\\
1.612	-0.000231786\\
1.614	-0.000245572\\
1.616	-0.000258327\\
1.618	-0.000270017\\
1.62	-0.000280612\\
1.622	-0.000290087\\
1.624	-0.000298419\\
1.626	-0.000305593\\
1.628	-0.000311596\\
1.63	-0.000316419\\
1.632	-0.000320059\\
1.634	-0.000322516\\
1.636	-0.000323793\\
1.638	-0.000323901\\
1.64	-0.000322852\\
1.642	-0.000320662\\
1.644	-0.000317352\\
1.646	-0.000312948\\
1.648	-0.000307475\\
1.65	-0.000300967\\
1.652	-0.000293458\\
1.654	-0.000284987\\
1.656	-0.000275594\\
1.658	-0.000265323\\
1.66	-0.00025422\\
1.662	-0.000242336\\
1.664	-0.000229721\\
1.666	-0.000216427\\
1.668	-0.000211444\\
1.67	-0.000209784\\
1.672	-0.000207466\\
1.674	-0.000204505\\
1.676	-0.00020092\\
1.678	-0.000196731\\
1.68	-0.00019196\\
1.682	-0.00018663\\
1.684	-0.000180768\\
1.686	-0.000174398\\
1.688	-0.000167548\\
1.69	-0.000160248\\
1.692	-0.000152528\\
1.694	-0.000144418\\
1.696	-0.00013595\\
1.698	-0.000127156\\
1.7	-0.00011807\\
1.702	-0.000108725\\
1.704	-9.91547e-05\\
1.706	-0.000101096\\
1.708	-0.000115053\\
1.71	-0.000128444\\
1.712	-0.000141227\\
1.714	-0.000153363\\
1.716	-0.000164816\\
1.718	-0.000175552\\
1.72	-0.000185543\\
1.722	-0.000194761\\
1.724	-0.000203181\\
1.726	-0.000210783\\
1.728	-0.000217549\\
1.73	-0.000223464\\
1.732	-0.000228517\\
1.734	-0.000232699\\
1.736	-0.000236006\\
1.738	-0.000238436\\
1.74	-0.000239989\\
1.742	-0.000240671\\
1.744	-0.000240488\\
1.746	-0.000239452\\
1.748	-0.000237575\\
1.75	-0.000234874\\
1.752	-0.000231368\\
1.754	-0.000227078\\
1.756	-0.000222028\\
1.758	-0.000216246\\
1.76	-0.000209759\\
1.762	-0.0002026\\
1.764	-0.000194801\\
1.766	-0.000186397\\
1.768	-0.000177424\\
1.77	-0.000167923\\
1.772	-0.000157931\\
1.774	-0.00014749\\
1.776	-0.000144768\\
1.778	-0.000144234\\
1.78	-0.000143237\\
1.782	-0.000141789\\
1.784	-0.000139901\\
1.786	-0.000137583\\
1.788	-0.000134851\\
1.79	-0.000131718\\
1.792	-0.0001282\\
1.794	-0.000124315\\
1.796	-0.00012008\\
1.798	-0.000115513\\
1.8	-0.000110635\\
1.802	-0.000105464\\
1.804	-0.000100024\\
1.806	-9.43334e-05\\
1.808	-8.8416e-05\\
1.81	-8.22937e-05\\
1.812	-7.75142e-05\\
1.814	-8.77485e-05\\
1.816	-9.7555e-05\\
1.818	-0.000106904\\
1.82	-0.000115766\\
1.822	-0.000124117\\
1.824	-0.000131932\\
1.826	-0.00013919\\
1.828	-0.00014587\\
1.83	-0.000151957\\
1.832	-0.000157435\\
1.834	-0.000162291\\
1.836	-0.000166515\\
1.838	-0.0001701\\
1.84	-0.00017304\\
1.842	-0.000175332\\
1.844	-0.000176975\\
1.846	-0.000177971\\
1.848	-0.000178323\\
1.85	-0.000178038\\
1.852	-0.000177123\\
1.854	-0.00017559\\
1.856	-0.000173451\\
1.858	-0.00017072\\
1.86	-0.000167413\\
1.862	-0.00016355\\
1.864	-0.00015915\\
1.866	-0.000154234\\
1.868	-0.000148827\\
1.87	-0.000142952\\
1.872	-0.000136637\\
1.874	-0.000129908\\
1.876	-0.000122794\\
1.878	-0.000115326\\
1.88	-0.000107532\\
1.882	-9.94458e-05\\
1.884	-9.80944e-05\\
1.886	-9.81554e-05\\
1.888	-9.78972e-05\\
1.89	-9.73251e-05\\
1.892	-9.64456e-05\\
1.894	-9.52657e-05\\
1.896	-9.37936e-05\\
1.898	-9.20383e-05\\
1.9	-9.00094e-05\\
1.902	-8.77173e-05\\
1.904	-8.51734e-05\\
1.906	-8.23894e-05\\
1.908	-7.93777e-05\\
1.91	-7.61515e-05\\
1.912	-7.27241e-05\\
1.914	-6.91097e-05\\
1.916	-6.53225e-05\\
1.918	-6.13773e-05\\
1.92	-6.63671e-05\\
1.922	-7.35526e-05\\
1.924	-8.03954e-05\\
1.926	-8.68752e-05\\
1.928	-9.29733e-05\\
1.93	-9.86722e-05\\
1.932	-0.000103956\\
1.934	-0.000108812\\
1.936	-0.000113226\\
1.938	-0.000117188\\
1.94	-0.00012069\\
1.942	-0.000123724\\
1.944	-0.000126285\\
1.946	-0.000128369\\
1.948	-0.000129974\\
1.95	-0.0001311\\
1.952	-0.000131748\\
1.954	-0.000131922\\
1.956	-0.000131626\\
1.958	-0.000130867\\
1.96	-0.000129653\\
1.962	-0.000127993\\
1.964	-0.000125898\\
1.966	-0.00012338\\
1.968	-0.000120454\\
1.97	-0.000117134\\
1.972	-0.000113437\\
1.974	-0.00010938\\
1.976	-0.000104982\\
1.978	-0.000100262\\
1.98	-9.52395e-05\\
1.982	-8.99371e-05\\
1.984	-8.43764e-05\\
1.986	-7.858e-05\\
1.988	-7.25712e-05\\
1.99	-6.63735e-05\\
1.992	-6.55904e-05\\
1.994	-6.59447e-05\\
1.996	-6.60811e-05\\
1.998	-6.60025e-05\\
2	-6.57122e-05\\
};
\addplot [color=mycolor7, line width=1.0pt, forget plot]
  table[row sep=crcr]{%
-0.024	0\\
-0.022	0\\
-0.02	0\\
-0.018	0\\
-0.016	0\\
-0.014	0\\
-0.012	0\\
-0.01	0\\
-0.008	0\\
-0.006	0\\
-0.004	0\\
-0.002	0\\
0	0\\
0.002	-8.9764e-08\\
0.004	-9.74028e-07\\
0.006	-5.1302e-06\\
0.008	-1.42162e-05\\
0.01	-3.09798e-05\\
0.012	-5.7834e-05\\
0.014	-9.68775e-05\\
0.016	-0.000149738\\
0.018	-0.000182472\\
0.02	-0.000253796\\
0.022	-0.000308821\\
0.024	-0.00136372\\
0.026	-0.00421295\\
0.028	-0.00728118\\
0.03	-0.0100042\\
0.032	-0.0110189\\
0.034	-0.01203\\
0.036	-0.0130299\\
0.038	-0.0140111\\
0.04	-0.0149661\\
0.042	-0.0158879\\
0.044	-0.0167696\\
0.046	-0.0176046\\
0.048	-0.0183865\\
0.05	-0.0191092\\
0.052	-0.0197671\\
0.054	-0.0203547\\
0.056	-0.020867\\
0.058	-0.0212993\\
0.06	-0.0216476\\
0.062	-0.021908\\
0.064	-0.0220772\\
0.066	-0.0221526\\
0.068	-0.0221317\\
0.07	-0.0220128\\
0.072	-0.0217948\\
0.074	-0.0214768\\
0.076	-0.0210589\\
0.078	-0.0205413\\
0.08	-0.019925\\
0.082	-0.0192115\\
0.084	-0.0184027\\
0.086	-0.0175013\\
0.088	-0.0165103\\
0.09	-0.0154331\\
0.092	-0.0142736\\
0.094	-0.0130363\\
0.096	-0.011726\\
0.098	-0.0103479\\
0.1	-0.00890731\\
0.102	-0.00741023\\
0.104	-0.0060898\\
0.106	-0.00656926\\
0.108	-0.00699074\\
0.11	-0.0073537\\
0.112	-0.00765788\\
0.114	-0.00790329\\
0.116	-0.00773885\\
0.118	-0.00732228\\
0.12	-0.00687744\\
0.122	-0.00683035\\
0.124	-0.00743542\\
0.126	-0.00809968\\
0.128	-0.0087338\\
0.13	-0.00933596\\
0.132	-0.00990445\\
0.134	-0.0102604\\
0.136	-0.010414\\
0.138	-0.0105416\\
0.14	-0.0106437\\
0.142	-0.011226\\
0.144	-0.0130663\\
0.146	-0.0148047\\
0.148	-0.0164368\\
0.15	-0.0179584\\
0.152	-0.019366\\
0.154	-0.0206563\\
0.156	-0.0218269\\
0.158	-0.0228754\\
0.16	-0.0238001\\
0.162	-0.0245999\\
0.164	-0.025274\\
0.166	-0.0258222\\
0.168	-0.0262446\\
0.17	-0.0265418\\
0.172	-0.0267151\\
0.174	-0.0267658\\
0.176	-0.026696\\
0.178	-0.0265081\\
0.18	-0.0262048\\
0.182	-0.0257892\\
0.184	-0.0252068\\
0.186	-0.0244707\\
0.188	-0.0236327\\
0.19	-0.0226977\\
0.192	-0.0216707\\
0.194	-0.0205571\\
0.196	-0.0193626\\
0.198	-0.0180929\\
0.2	-0.0167541\\
0.202	-0.0153523\\
0.204	-0.0138937\\
0.206	-0.0123849\\
0.208	-0.0108322\\
0.21	-0.00930484\\
0.212	-0.00939525\\
0.214	-0.00944152\\
0.216	-0.00944442\\
0.218	-0.00940482\\
0.22	-0.00932372\\
0.222	-0.00920219\\
0.224	-0.0090414\\
0.226	-0.00884264\\
0.228	-0.00860726\\
0.23	-0.00833669\\
0.232	-0.00803245\\
0.234	-0.00769614\\
0.236	-0.0073294\\
0.238	-0.00693396\\
0.24	-0.00766886\\
0.242	-0.009144\\
0.244	-0.010555\\
0.246	-0.0118978\\
0.248	-0.0131688\\
0.25	-0.0143646\\
0.252	-0.0154822\\
0.254	-0.0165187\\
0.256	-0.0174717\\
0.258	-0.0183392\\
0.26	-0.0191194\\
0.262	-0.0198108\\
0.264	-0.0204123\\
0.266	-0.0209231\\
0.268	-0.0213428\\
0.27	-0.0216713\\
0.272	-0.0219088\\
0.274	-0.0220558\\
0.276	-0.0221131\\
0.278	-0.0220819\\
0.28	-0.0219636\\
0.282	-0.02176\\
0.284	-0.0214731\\
0.286	-0.0211051\\
0.288	-0.0206586\\
0.29	-0.0201364\\
0.292	-0.0195415\\
0.294	-0.018877\\
0.296	-0.0181466\\
0.298	-0.0173537\\
0.3	-0.0165022\\
0.302	-0.015596\\
0.304	-0.0146392\\
0.306	-0.0136361\\
0.308	-0.012591\\
0.31	-0.0115082\\
0.312	-0.0103923\\
0.314	-0.00924786\\
0.316	-0.00807942\\
0.318	-0.007969\\
0.32	-0.00809842\\
0.322	-0.00819632\\
0.324	-0.00826287\\
0.326	-0.00829835\\
0.328	-0.00830311\\
0.33	-0.0082776\\
0.332	-0.00822234\\
0.334	-0.00813796\\
0.336	-0.00802514\\
0.338	-0.00788466\\
0.34	-0.00771735\\
0.342	-0.00752415\\
0.344	-0.00730602\\
0.346	-0.00706402\\
0.348	-0.00803051\\
0.35	-0.00898983\\
0.352	-0.00990042\\
0.354	-0.0107598\\
0.356	-0.0115655\\
0.358	-0.0123157\\
0.36	-0.0130083\\
0.362	-0.0136418\\
0.364	-0.0142147\\
0.366	-0.0147259\\
0.368	-0.0151745\\
0.37	-0.0155598\\
0.372	-0.0158812\\
0.374	-0.0161386\\
0.376	-0.016332\\
0.378	-0.0164615\\
0.38	-0.0165277\\
0.382	-0.0165311\\
0.384	-0.0164727\\
0.386	-0.0163535\\
0.388	-0.0161749\\
0.39	-0.0159383\\
0.392	-0.0156453\\
0.394	-0.0152979\\
0.396	-0.014898\\
0.398	-0.0144478\\
0.4	-0.0139497\\
0.402	-0.013406\\
0.404	-0.0128193\\
0.406	-0.0121925\\
0.408	-0.0115282\\
0.41	-0.0108294\\
0.412	-0.0100991\\
0.414	-0.00934035\\
0.416	-0.00912883\\
0.418	-0.00892402\\
0.42	-0.00869782\\
0.422	-0.00845111\\
0.424	-0.00818483\\
0.426	-0.00789994\\
0.428	-0.00759743\\
0.43	-0.00727833\\
0.432	-0.00694372\\
0.434	-0.00659467\\
0.436	-0.00647083\\
0.438	-0.00649243\\
0.44	-0.00649177\\
0.442	-0.00646912\\
0.444	-0.00642482\\
0.446	-0.00635926\\
0.448	-0.00627291\\
0.45	-0.00616629\\
0.452	-0.00603997\\
0.454	-0.00623222\\
0.456	-0.00689198\\
0.458	-0.00751721\\
0.46	-0.00810619\\
0.462	-0.00865732\\
0.464	-0.00916915\\
0.466	-0.00964042\\
0.468	-0.01007\\
0.47	-0.010457\\
0.472	-0.0108005\\
0.474	-0.0110999\\
0.476	-0.0113549\\
0.478	-0.011565\\
0.48	-0.0117302\\
0.482	-0.0118504\\
0.484	-0.011926\\
0.486	-0.0119571\\
0.488	-0.0119443\\
0.49	-0.0118883\\
0.492	-0.0117898\\
0.494	-0.0116498\\
0.496	-0.0114693\\
0.498	-0.0112496\\
0.5	-0.0109919\\
0.502	-0.0106977\\
0.504	-0.0103686\\
0.506	-0.0100061\\
0.508	-0.00961205\\
0.51	-0.0091883\\
0.512	-0.00873674\\
0.514	-0.00825938\\
0.516	-0.00775825\\
0.518	-0.00723549\\
0.52	-0.00669326\\
0.522	-0.00667996\\
0.524	-0.00669889\\
0.526	-0.00670021\\
0.528	-0.00668416\\
0.53	-0.00665105\\
0.532	-0.00660122\\
0.534	-0.00653505\\
0.536	-0.00645296\\
0.538	-0.00635539\\
0.54	-0.00624283\\
0.542	-0.00611579\\
0.544	-0.00597482\\
0.546	-0.00582048\\
0.548	-0.00565337\\
0.55	-0.00547411\\
0.552	-0.00528335\\
0.554	-0.00508174\\
0.556	-0.00494948\\
0.558	-0.00489393\\
0.56	-0.00482287\\
0.562	-0.00504494\\
0.564	-0.00547669\\
0.566	-0.00588272\\
0.568	-0.00626193\\
0.57	-0.00661334\\
0.572	-0.00693608\\
0.574	-0.0072294\\
0.576	-0.00749263\\
0.578	-0.00772527\\
0.58	-0.0079269\\
0.582	-0.00809722\\
0.584	-0.00823606\\
0.586	-0.00834334\\
0.588	-0.00841912\\
0.59	-0.00846356\\
0.592	-0.00847693\\
0.594	-0.0084596\\
0.596	-0.00841206\\
0.598	-0.0083349\\
0.6	-0.00822878\\
0.602	-0.0080945\\
0.604	-0.0079329\\
0.606	-0.00774495\\
0.608	-0.00753168\\
0.61	-0.0072942\\
0.612	-0.00703368\\
0.614	-0.00675138\\
0.616	-0.0064486\\
0.618	-0.00612671\\
0.62	-0.00578712\\
0.622	-0.00543128\\
0.624	-0.0050607\\
0.626	-0.00467691\\
0.628	-0.00428145\\
0.63	-0.00387591\\
0.632	-0.00372671\\
0.634	-0.00382707\\
0.636	-0.00391613\\
0.638	-0.00399383\\
0.64	-0.00406016\\
0.642	-0.00411513\\
0.644	-0.00415878\\
0.646	-0.00419117\\
0.648	-0.0042124\\
0.65	-0.00422259\\
0.652	-0.00422189\\
0.654	-0.00421048\\
0.656	-0.00418854\\
0.658	-0.00415631\\
0.66	-0.00411403\\
0.662	-0.00406195\\
0.664	-0.00400038\\
0.666	-0.0039296\\
0.668	-0.00384996\\
0.67	-0.00391207\\
0.672	-0.0041895\\
0.674	-0.00444798\\
0.676	-0.00468685\\
0.678	-0.00490552\\
0.68	-0.00510349\\
0.682	-0.00528033\\
0.684	-0.0054357\\
0.686	-0.00556935\\
0.688	-0.00568108\\
0.69	-0.00577081\\
0.692	-0.00583852\\
0.694	-0.00588426\\
0.696	-0.00590817\\
0.698	-0.00591048\\
0.7	-0.00589146\\
0.702	-0.00585149\\
0.704	-0.00579099\\
0.706	-0.00571047\\
0.708	-0.00561049\\
0.71	-0.00549168\\
0.712	-0.00535472\\
0.714	-0.00520036\\
0.716	-0.0050294\\
0.718	-0.00484266\\
0.72	-0.00464105\\
0.722	-0.0044255\\
0.724	-0.00419697\\
0.726	-0.00395646\\
0.728	-0.00370501\\
0.73	-0.00344366\\
0.732	-0.0031735\\
0.734	-0.00289562\\
0.736	-0.00261112\\
0.738	-0.00232113\\
0.74	-0.00202676\\
0.742	-0.00192622\\
0.744	-0.00202336\\
0.746	-0.00215265\\
0.748	-0.00227384\\
0.75	-0.00238662\\
0.752	-0.00249071\\
0.754	-0.00258586\\
0.756	-0.00267186\\
0.758	-0.00274853\\
0.76	-0.0028157\\
0.762	-0.00287327\\
0.764	-0.00292113\\
0.766	-0.00295923\\
0.768	-0.00298754\\
0.77	-0.00300607\\
0.772	-0.00301485\\
0.774	-0.00301396\\
0.776	-0.00300347\\
0.778	-0.00298354\\
0.78	-0.00315433\\
0.782	-0.00331377\\
0.784	-0.00345897\\
0.786	-0.00358961\\
0.788	-0.00370541\\
0.79	-0.00380617\\
0.792	-0.00389172\\
0.794	-0.00396197\\
0.796	-0.00401687\\
0.798	-0.00405644\\
0.8	-0.00408075\\
0.802	-0.00408991\\
0.804	-0.00408411\\
0.806	-0.00406357\\
0.808	-0.00402857\\
0.81	-0.00397943\\
0.812	-0.00391653\\
0.814	-0.00384029\\
0.816	-0.00375116\\
0.818	-0.00364965\\
0.82	-0.00353629\\
0.822	-0.00341165\\
0.824	-0.00327635\\
0.826	-0.00313103\\
0.828	-0.00297633\\
0.83	-0.00281297\\
0.832	-0.00264164\\
0.834	-0.00246308\\
0.836	-0.00227803\\
0.838	-0.00208726\\
0.84	-0.00189153\\
0.842	-0.00169163\\
0.844	-0.00148834\\
0.846	-0.00128244\\
0.848	-0.00117641\\
0.85	-0.00123612\\
0.852	-0.00135039\\
0.854	-0.00145971\\
0.856	-0.00156338\\
0.858	-0.00166108\\
0.86	-0.00175256\\
0.862	-0.00183759\\
0.864	-0.00191597\\
0.866	-0.0019875\\
0.868	-0.00205204\\
0.87	-0.00210945\\
0.872	-0.00215961\\
0.874	-0.00220244\\
0.876	-0.00223788\\
0.878	-0.00226589\\
0.88	-0.00228645\\
0.882	-0.00229957\\
0.884	-0.00230529\\
0.886	-0.00230365\\
0.888	-0.0023697\\
0.89	-0.00246364\\
0.892	-0.00254727\\
0.894	-0.00262044\\
0.896	-0.00268302\\
0.898	-0.00273492\\
0.9	-0.0027761\\
0.902	-0.00280656\\
0.904	-0.00282633\\
0.906	-0.00283549\\
0.908	-0.00283414\\
0.91	-0.00282243\\
0.912	-0.00280054\\
0.914	-0.0027687\\
0.916	-0.00272715\\
0.918	-0.00267618\\
0.92	-0.00261609\\
0.922	-0.00254724\\
0.924	-0.00246999\\
0.926	-0.00238474\\
0.928	-0.00229191\\
0.93	-0.00219193\\
0.932	-0.00208527\\
0.934	-0.00197242\\
0.936	-0.00185385\\
0.938	-0.00173009\\
0.94	-0.00160166\\
0.942	-0.00146909\\
0.944	-0.00133292\\
0.946	-0.00119369\\
0.948	-0.00105196\\
0.95	-0.00090828\\
0.952	-0.000812624\\
0.954	-0.000806504\\
0.956	-0.000797942\\
0.958	-0.000786994\\
0.96	-0.000841312\\
0.962	-0.000931107\\
0.964	-0.00101732\\
0.966	-0.00109943\\
0.968	-0.00117721\\
0.97	-0.00125046\\
0.972	-0.001319\\
0.974	-0.00138265\\
0.976	-0.00144127\\
0.978	-0.00149472\\
0.98	-0.00154289\\
0.982	-0.00158568\\
0.984	-0.00162301\\
0.986	-0.00165483\\
0.988	-0.00168108\\
0.99	-0.00170174\\
0.992	-0.0017168\\
0.994	-0.00172733\\
0.996	-0.00178658\\
0.998	-0.00183847\\
1	-0.00188289\\
1.002	-0.0019198\\
1.004	-0.00194915\\
1.006	-0.00197093\\
1.008	-0.00198516\\
1.01	-0.0019919\\
1.012	-0.0019912\\
1.014	-0.00198318\\
1.016	-0.00196796\\
1.018	-0.00194568\\
1.02	-0.00191652\\
1.022	-0.00188068\\
1.024	-0.00183838\\
1.026	-0.00178985\\
1.028	-0.00173535\\
1.03	-0.00167517\\
1.032	-0.00160959\\
1.034	-0.00153892\\
1.036	-0.0014635\\
1.038	-0.00138366\\
1.04	-0.00129975\\
1.042	-0.00121213\\
1.044	-0.00112117\\
1.046	-0.00102724\\
1.048	-0.000930739\\
1.05	-0.000832043\\
1.052	-0.000731546\\
1.054	-0.00066676\\
1.056	-0.000646928\\
1.058	-0.000625959\\
1.06	-0.000603929\\
1.062	-0.000580916\\
1.064	-0.000558963\\
1.066	-0.000555414\\
1.068	-0.000550192\\
1.07	-0.000549839\\
1.072	-0.000618987\\
1.074	-0.000685531\\
1.076	-0.000749169\\
1.078	-0.000809731\\
1.08	-0.000867057\\
1.082	-0.000921002\\
1.084	-0.000971432\\
1.086	-0.00101823\\
1.088	-0.00106127\\
1.09	-0.00110048\\
1.092	-0.00113577\\
1.094	-0.00116707\\
1.096	-0.00119542\\
1.098	-0.0012425\\
1.1	-0.00128442\\
1.102	-0.00132107\\
1.104	-0.00135239\\
1.106	-0.00137833\\
1.108	-0.00139887\\
1.11	-0.00141399\\
1.112	-0.00142371\\
1.114	-0.00142805\\
1.116	-0.00142708\\
1.118	-0.00142085\\
1.12	-0.00140947\\
1.122	-0.00139302\\
1.124	-0.00137164\\
1.126	-0.00134547\\
1.128	-0.00131466\\
1.13	-0.00127938\\
1.132	-0.00123981\\
1.134	-0.00119616\\
1.136	-0.00114864\\
1.138	-0.00109746\\
1.14	-0.00104286\\
1.142	-0.000985087\\
1.144	-0.000924385\\
1.146	-0.000861015\\
1.148	-0.000795243\\
1.15	-0.000727339\\
1.152	-0.000657579\\
1.154	-0.000586242\\
1.156	-0.000513607\\
1.158	-0.000439959\\
1.16	-0.000398067\\
1.162	-0.000398851\\
1.164	-0.000398749\\
1.166	-0.000397784\\
1.168	-0.000395976\\
1.17	-0.000393349\\
1.172	-0.000389927\\
1.174	-0.000385736\\
1.176	-0.000380801\\
1.178	-0.000375151\\
1.18	-0.000368813\\
1.182	-0.000422502\\
1.184	-0.000484878\\
1.186	-0.000544728\\
1.188	-0.000601863\\
1.19	-0.000656105\\
1.192	-0.000707291\\
1.194	-0.000755269\\
1.196	-0.000799903\\
1.198	-0.000841068\\
1.2	-0.000878655\\
1.202	-0.00091257\\
1.204	-0.000942731\\
1.206	-0.000969074\\
1.208	-0.000991547\\
1.21	-0.00101011\\
1.212	-0.00102475\\
1.214	-0.00103546\\
1.216	-0.00104223\\
1.218	-0.0010451\\
1.22	-0.0010441\\
1.222	-0.00103928\\
1.224	-0.00103071\\
1.226	-0.00101845\\
1.228	-0.0010026\\
1.23	-0.000983252\\
1.232	-0.000960526\\
1.234	-0.000934542\\
1.236	-0.000905433\\
1.238	-0.000873342\\
1.24	-0.000838422\\
1.242	-0.000800834\\
1.244	-0.000760746\\
1.246	-0.000718334\\
1.248	-0.000673781\\
1.25	-0.000627274\\
1.252	-0.000585477\\
1.254	-0.000547705\\
1.256	-0.000508648\\
1.258	-0.000468445\\
1.26	-0.000427239\\
1.262	-0.000385175\\
1.264	-0.000342396\\
1.266	-0.000299049\\
1.268	-0.000255279\\
1.27	-0.000211232\\
1.272	-0.000197923\\
1.274	-0.000202692\\
1.276	-0.000206725\\
1.278	-0.000210022\\
1.28	-0.000212585\\
1.282	-0.000218978\\
1.284	-0.000268142\\
1.286	-0.000315766\\
1.288	-0.000361696\\
1.29	-0.000405782\\
1.292	-0.000447886\\
1.294	-0.000487877\\
1.296	-0.000525633\\
1.298	-0.000561043\\
1.3	-0.000594005\\
1.302	-0.000624427\\
1.304	-0.000652227\\
1.306	-0.000677334\\
1.308	-0.000699687\\
1.31	-0.000719238\\
1.312	-0.000735946\\
1.314	-0.000749783\\
1.316	-0.000760733\\
1.318	-0.000768788\\
1.32	-0.000773952\\
1.322	-0.000776239\\
1.324	-0.000775675\\
1.326	-0.000772293\\
1.328	-0.000766137\\
1.33	-0.000757263\\
1.332	-0.000745733\\
1.334	-0.00073162\\
1.336	-0.000715003\\
1.338	-0.000695973\\
1.34	-0.000674625\\
1.342	-0.000651064\\
1.344	-0.000625401\\
1.346	-0.000597751\\
1.348	-0.000579014\\
1.35	-0.000562365\\
1.352	-0.000544125\\
1.354	-0.000524373\\
1.356	-0.000503191\\
1.358	-0.000480665\\
1.36	-0.000456884\\
1.362	-0.000431941\\
1.364	-0.000405931\\
1.366	-0.00037895\\
1.368	-0.000351099\\
1.37	-0.000322476\\
1.372	-0.000293185\\
1.374	-0.000263327\\
1.376	-0.000233007\\
1.378	-0.000202327\\
1.38	-0.00017139\\
1.382	-0.0001403\\
1.384	-0.000120412\\
1.386	-0.00015812\\
1.388	-0.000194836\\
1.39	-0.000230438\\
1.392	-0.00026481\\
1.394	-0.000297842\\
1.396	-0.000329429\\
1.398	-0.000359473\\
1.4	-0.000387883\\
1.402	-0.000414574\\
1.404	-0.000439469\\
1.406	-0.000462499\\
1.408	-0.0004836\\
1.41	-0.000502719\\
1.412	-0.000519809\\
1.414	-0.000534831\\
1.416	-0.000547753\\
1.418	-0.000558553\\
1.42	-0.000567215\\
1.422	-0.000573733\\
1.424	-0.000578106\\
1.426	-0.000580343\\
1.428	-0.00058046\\
1.43	-0.00057848\\
1.432	-0.000574434\\
1.434	-0.00056836\\
1.436	-0.000560302\\
1.438	-0.000550312\\
1.44	-0.000538448\\
1.442	-0.000524773\\
1.444	-0.000509356\\
1.446	-0.000492274\\
1.448	-0.000473606\\
1.45	-0.000453437\\
1.452	-0.000435022\\
1.454	-0.000427758\\
1.456	-0.00041921\\
1.458	-0.000409422\\
1.46	-0.000398444\\
1.462	-0.000386324\\
1.464	-0.000373119\\
1.466	-0.000358884\\
1.468	-0.000343679\\
1.47	-0.000327566\\
1.472	-0.000310608\\
1.474	-0.000292872\\
1.476	-0.000274425\\
1.478	-0.000255335\\
1.48	-0.000235674\\
1.482	-0.000215512\\
1.484	-0.000194921\\
1.486	-0.000173974\\
1.488	-0.000152741\\
1.49	-0.000131297\\
1.492	-0.000135563\\
1.494	-0.000162428\\
1.496	-0.000188412\\
1.498	-0.000213431\\
1.5	-0.000237407\\
1.502	-0.000260264\\
1.504	-0.000281933\\
1.506	-0.00030235\\
1.508	-0.000321455\\
1.51	-0.000339195\\
1.512	-0.00035552\\
1.514	-0.000370388\\
1.516	-0.000383762\\
1.518	-0.00039561\\
1.52	-0.000405908\\
1.522	-0.000414635\\
1.524	-0.000421778\\
1.526	-0.000427329\\
1.528	-0.000431286\\
1.53	-0.000433651\\
1.532	-0.000434436\\
1.534	-0.000433654\\
1.536	-0.000431326\\
1.538	-0.000427478\\
1.54	-0.00042214\\
1.542	-0.000415348\\
1.544	-0.000407143\\
1.546	-0.000397571\\
1.548	-0.00038668\\
1.55	-0.000374524\\
1.552	-0.000361162\\
1.554	-0.000346655\\
1.556	-0.000331067\\
1.558	-0.000314467\\
1.56	-0.00030548\\
1.562	-0.000301781\\
1.564	-0.000297154\\
1.566	-0.000291627\\
1.568	-0.000285229\\
1.57	-0.000277993\\
1.572	-0.000269954\\
1.574	-0.000261148\\
1.576	-0.000251615\\
1.578	-0.000241395\\
1.58	-0.00023053\\
1.582	-0.000219065\\
1.584	-0.000207044\\
1.586	-0.000194515\\
1.588	-0.000181524\\
1.59	-0.000168121\\
1.592	-0.000154353\\
1.594	-0.000140272\\
1.596	-0.000125926\\
1.598	-0.000111366\\
1.6	-0.000129828\\
1.602	-0.000148863\\
1.604	-0.000167154\\
1.606	-0.000184643\\
1.608	-0.000201277\\
1.61	-0.000217007\\
1.612	-0.000231786\\
1.614	-0.000245572\\
1.616	-0.000258327\\
1.618	-0.000270017\\
1.62	-0.000280612\\
1.622	-0.000290086\\
1.624	-0.000298419\\
1.626	-0.000305593\\
1.628	-0.000311596\\
1.63	-0.000316419\\
1.632	-0.000320059\\
1.634	-0.000322515\\
1.636	-0.000323793\\
1.638	-0.000323901\\
1.64	-0.000322852\\
1.642	-0.000320662\\
1.644	-0.000317352\\
1.646	-0.000312947\\
1.648	-0.000307475\\
1.65	-0.000300967\\
1.652	-0.000293458\\
1.654	-0.000284987\\
1.656	-0.000275594\\
1.658	-0.000265322\\
1.66	-0.00025422\\
1.662	-0.000242336\\
1.664	-0.00022972\\
1.666	-0.000216427\\
1.668	-0.000211444\\
1.67	-0.000209784\\
1.672	-0.000207466\\
1.674	-0.000204505\\
1.676	-0.000200919\\
1.678	-0.000196731\\
1.68	-0.00019196\\
1.682	-0.00018663\\
1.684	-0.000180767\\
1.686	-0.000174398\\
1.688	-0.000167548\\
1.69	-0.000160248\\
1.692	-0.000152528\\
1.694	-0.000144418\\
1.696	-0.000135949\\
1.698	-0.000127156\\
1.7	-0.000118069\\
1.702	-0.000108724\\
1.704	-9.91546e-05\\
1.706	-0.000101096\\
1.708	-0.000115053\\
1.71	-0.000128444\\
1.712	-0.000141227\\
1.714	-0.000153363\\
1.716	-0.000164815\\
1.718	-0.000175552\\
1.72	-0.000185543\\
1.722	-0.000194761\\
1.724	-0.000203181\\
1.726	-0.000210783\\
1.728	-0.000217549\\
1.73	-0.000223464\\
1.732	-0.000228517\\
1.734	-0.000232699\\
1.736	-0.000236006\\
1.738	-0.000238436\\
1.74	-0.000239989\\
1.742	-0.000240671\\
1.744	-0.000240488\\
1.746	-0.000239452\\
1.748	-0.000237575\\
1.75	-0.000234874\\
1.752	-0.000231368\\
1.754	-0.000227078\\
1.756	-0.000222028\\
1.758	-0.000216246\\
1.76	-0.000209759\\
1.762	-0.0002026\\
1.764	-0.000194801\\
1.766	-0.000186397\\
1.768	-0.000177424\\
1.77	-0.000167922\\
1.772	-0.000157931\\
1.774	-0.00014749\\
1.776	-0.000144768\\
1.778	-0.000144234\\
1.78	-0.000143237\\
1.782	-0.000141789\\
1.784	-0.000139901\\
1.786	-0.000137583\\
1.788	-0.000134851\\
1.79	-0.000131718\\
1.792	-0.0001282\\
1.794	-0.000124315\\
1.796	-0.00012008\\
1.798	-0.000115513\\
1.8	-0.000110635\\
1.802	-0.000105464\\
1.804	-0.000100024\\
1.806	-9.43334e-05\\
1.808	-8.8416e-05\\
1.81	-8.22937e-05\\
1.812	-7.75141e-05\\
1.814	-8.77485e-05\\
1.816	-9.75549e-05\\
1.818	-0.000106904\\
1.82	-0.000115766\\
1.822	-0.000124117\\
1.824	-0.000131932\\
1.826	-0.00013919\\
1.828	-0.00014587\\
1.83	-0.000151957\\
1.832	-0.000157435\\
1.834	-0.000162291\\
1.836	-0.000166515\\
1.838	-0.0001701\\
1.84	-0.00017304\\
1.842	-0.000175332\\
1.844	-0.000176975\\
1.846	-0.000177971\\
1.848	-0.000178323\\
1.85	-0.000178038\\
1.852	-0.000177123\\
1.854	-0.00017559\\
1.856	-0.000173451\\
1.858	-0.00017072\\
1.86	-0.000167413\\
1.862	-0.00016355\\
1.864	-0.00015915\\
1.866	-0.000154234\\
1.868	-0.000148827\\
1.87	-0.000142952\\
1.872	-0.000136637\\
1.874	-0.000129908\\
1.876	-0.000122794\\
1.878	-0.000115326\\
1.88	-0.000107532\\
1.882	-9.94458e-05\\
1.884	-9.80944e-05\\
1.886	-9.81553e-05\\
1.888	-9.78971e-05\\
1.89	-9.73251e-05\\
1.892	-9.64456e-05\\
1.894	-9.52657e-05\\
1.896	-9.37936e-05\\
1.898	-9.20383e-05\\
1.9	-9.00093e-05\\
1.902	-8.77173e-05\\
1.904	-8.51734e-05\\
1.906	-8.23894e-05\\
1.908	-7.93777e-05\\
1.91	-7.61515e-05\\
1.912	-7.27241e-05\\
1.914	-6.91096e-05\\
1.916	-6.53225e-05\\
1.918	-6.13773e-05\\
1.92	-6.63671e-05\\
1.922	-7.35526e-05\\
1.924	-8.03954e-05\\
1.926	-8.68752e-05\\
1.928	-9.29733e-05\\
1.93	-9.86722e-05\\
1.932	-0.000103956\\
1.934	-0.000108812\\
1.936	-0.000113226\\
1.938	-0.000117188\\
1.94	-0.00012069\\
1.942	-0.000123724\\
1.944	-0.000126285\\
1.946	-0.000128369\\
1.948	-0.000129974\\
1.95	-0.0001311\\
1.952	-0.000131748\\
1.954	-0.000131922\\
1.956	-0.000131626\\
1.958	-0.000130867\\
1.96	-0.000129653\\
1.962	-0.000127993\\
1.964	-0.000125898\\
1.966	-0.00012338\\
1.968	-0.000120454\\
1.97	-0.000117134\\
1.972	-0.000113437\\
1.974	-0.00010938\\
1.976	-0.000104982\\
1.978	-0.000100262\\
1.98	-9.52395e-05\\
1.982	-8.99371e-05\\
1.984	-8.43764e-05\\
1.986	-7.858e-05\\
1.988	-7.25711e-05\\
1.99	-6.63735e-05\\
1.992	-6.55904e-05\\
1.994	-6.59447e-05\\
1.996	-6.60811e-05\\
1.998	-6.60025e-05\\
2	-6.57122e-05\\
};
\addplot [color=mycolor8, line width=1.0pt, forget plot]
  table[row sep=crcr]{%
-0.024	0\\
-0.022	0\\
-0.02	0\\
-0.018	0\\
-0.016	0\\
-0.014	0\\
-0.012	0\\
-0.01	0\\
-0.008	0\\
-0.006	0\\
-0.004	0\\
-0.002	0\\
0	0\\
0.002	-1.62551e-07\\
0.004	-1.74859e-06\\
0.006	-9.93482e-06\\
0.008	-2.73142e-05\\
0.01	-5.89081e-05\\
0.012	-0.000108651\\
0.014	-0.00017958\\
0.016	-0.000273579\\
0.018	-0.000365224\\
0.02	-0.000498998\\
0.022	-0.000974127\\
0.024	-0.00272713\\
0.026	-0.00621601\\
0.028	-0.00894339\\
0.03	-0.00996574\\
0.032	-0.0109953\\
0.034	-0.012024\\
0.036	-0.013044\\
0.038	-0.0140474\\
0.04	-0.0150266\\
0.042	-0.0159741\\
0.044	-0.0168827\\
0.046	-0.0177454\\
0.048	-0.0185554\\
0.05	-0.0193063\\
0.052	-0.019992\\
0.054	-0.0206068\\
0.056	-0.0211454\\
0.058	-0.0216028\\
0.06	-0.0219746\\
0.062	-0.0222569\\
0.064	-0.0224461\\
0.066	-0.0225394\\
0.068	-0.0225344\\
0.07	-0.0224292\\
0.072	-0.0222226\\
0.074	-0.0219138\\
0.076	-0.0215028\\
0.078	-0.0209899\\
0.08	-0.0203763\\
0.082	-0.0196634\\
0.084	-0.0188533\\
0.086	-0.0179487\\
0.088	-0.0169527\\
0.09	-0.015869\\
0.092	-0.0147015\\
0.094	-0.0134549\\
0.096	-0.012134\\
0.098	-0.010744\\
0.1	-0.00929073\\
0.102	-0.00777995\\
0.104	-0.00644502\\
0.106	-0.00690925\\
0.108	-0.00731486\\
0.11	-0.00766138\\
0.112	-0.00794863\\
0.114	-0.0081767\\
0.116	-0.00804455\\
0.118	-0.00761446\\
0.12	-0.00715563\\
0.122	-0.00709412\\
0.124	-0.00768438\\
0.126	-0.00833349\\
0.128	-0.00895216\\
0.13	-0.00953861\\
0.132	-0.0100912\\
0.134	-0.0105116\\
0.136	-0.0106519\\
0.138	-0.0107671\\
0.14	-0.0108565\\
0.142	-0.0114258\\
0.144	-0.0132528\\
0.146	-0.0149778\\
0.148	-0.0165962\\
0.15	-0.018104\\
0.152	-0.0194977\\
0.154	-0.0207741\\
0.156	-0.0219306\\
0.158	-0.0229651\\
0.16	-0.0238759\\
0.162	-0.0246619\\
0.164	-0.0253222\\
0.166	-0.0258568\\
0.168	-0.0262657\\
0.17	-0.0265498\\
0.172	-0.02671\\
0.174	-0.026748\\
0.176	-0.0266658\\
0.178	-0.0264657\\
0.18	-0.0261505\\
0.182	-0.0257235\\
0.184	-0.0251881\\
0.186	-0.0244994\\
0.188	-0.0236504\\
0.19	-0.0227047\\
0.192	-0.0216673\\
0.194	-0.0205436\\
0.196	-0.0193392\\
0.198	-0.01806\\
0.2	-0.0167119\\
0.202	-0.015301\\
0.204	-0.0138337\\
0.206	-0.0123164\\
0.208	-0.0107555\\
0.21	-0.0092201\\
0.212	-0.00930276\\
0.214	-0.0093415\\
0.216	-0.0093371\\
0.218	-0.00929041\\
0.22	-0.00920242\\
0.222	-0.0090742\\
0.224	-0.00890692\\
0.226	-0.00870186\\
0.228	-0.00846035\\
0.23	-0.00818384\\
0.232	-0.00787384\\
0.234	-0.00753193\\
0.236	-0.00715976\\
0.238	-0.00675907\\
0.24	-0.00748888\\
0.242	-0.0089591\\
0.244	-0.0103653\\
0.246	-0.0117036\\
0.248	-0.0129702\\
0.25	-0.0141617\\
0.252	-0.0152752\\
0.254	-0.0163078\\
0.256	-0.0172571\\
0.258	-0.0181211\\
0.26	-0.0188979\\
0.262	-0.0195861\\
0.264	-0.0201846\\
0.266	-0.0206926\\
0.268	-0.0211097\\
0.27	-0.0214357\\
0.272	-0.0216709\\
0.274	-0.0218158\\
0.276	-0.0218713\\
0.278	-0.0218383\\
0.28	-0.0217185\\
0.282	-0.0215136\\
0.284	-0.0212254\\
0.286	-0.0208564\\
0.288	-0.0204091\\
0.29	-0.0198862\\
0.292	-0.0192907\\
0.294	-0.0186259\\
0.296	-0.0178952\\
0.298	-0.0171023\\
0.3	-0.0162508\\
0.302	-0.0153448\\
0.304	-0.0143884\\
0.306	-0.0133857\\
0.308	-0.0123411\\
0.31	-0.011259\\
0.312	-0.0101439\\
0.314	-0.00900029\\
0.316	-0.00783282\\
0.318	-0.00772345\\
0.32	-0.00785401\\
0.322	-0.00795313\\
0.324	-0.00802097\\
0.326	-0.00805781\\
0.328	-0.00806399\\
0.33	-0.00803997\\
0.332	-0.00798625\\
0.334	-0.00790346\\
0.336	-0.00779228\\
0.338	-0.00765348\\
0.34	-0.0074879\\
0.342	-0.00729646\\
0.344	-0.00708013\\
0.346	-0.00683995\\
0.348	-0.0078083\\
0.35	-0.0087695\\
0.352	-0.00968199\\
0.354	-0.0105432\\
0.356	-0.011351\\
0.358	-0.0121031\\
0.36	-0.0127976\\
0.362	-0.0134331\\
0.364	-0.0140081\\
0.366	-0.0145213\\
0.368	-0.0149719\\
0.37	-0.0153592\\
0.372	-0.0156827\\
0.374	-0.0159421\\
0.376	-0.0161375\\
0.378	-0.016269\\
0.38	-0.0163372\\
0.382	-0.0163427\\
0.384	-0.0162863\\
0.386	-0.0161691\\
0.388	-0.0159925\\
0.39	-0.0157579\\
0.392	-0.015467\\
0.394	-0.0151215\\
0.396	-0.0147236\\
0.398	-0.0142754\\
0.4	-0.0137792\\
0.402	-0.0132375\\
0.404	-0.0126528\\
0.406	-0.0120279\\
0.408	-0.0113656\\
0.41	-0.0106687\\
0.412	-0.00994033\\
0.414	-0.00918346\\
0.416	-0.00897383\\
0.418	-0.00877089\\
0.42	-0.00854655\\
0.422	-0.0083017\\
0.424	-0.00803725\\
0.426	-0.00775418\\
0.428	-0.00745348\\
0.43	-0.00713618\\
0.432	-0.00680335\\
0.434	-0.00645607\\
0.436	-0.00633398\\
0.438	-0.00635731\\
0.44	-0.00635837\\
0.442	-0.00633743\\
0.444	-0.00629482\\
0.446	-0.00623094\\
0.448	-0.00614625\\
0.45	-0.00604127\\
0.452	-0.00591658\\
0.454	-0.00611044\\
0.456	-0.00677179\\
0.458	-0.00739861\\
0.46	-0.00798914\\
0.462	-0.00854182\\
0.464	-0.00905518\\
0.466	-0.00952796\\
0.468	-0.00995904\\
0.47	-0.0103475\\
0.472	-0.0106924\\
0.474	-0.0109934\\
0.476	-0.0112497\\
0.478	-0.0114613\\
0.48	-0.0116278\\
0.482	-0.0117495\\
0.484	-0.0118264\\
0.486	-0.0118588\\
0.488	-0.0118474\\
0.49	-0.0117927\\
0.492	-0.0116955\\
0.494	-0.0115568\\
0.496	-0.0113776\\
0.498	-0.0111591\\
0.5	-0.0109026\\
0.502	-0.0106097\\
0.504	-0.0102817\\
0.506	-0.00992043\\
0.508	-0.00952758\\
0.51	-0.00910499\\
0.512	-0.00865458\\
0.514	-0.00817835\\
0.516	-0.00767834\\
0.518	-0.00715669\\
0.52	-0.00661555\\
0.522	-0.00660333\\
0.524	-0.00662333\\
0.526	-0.0066257\\
0.528	-0.00661069\\
0.53	-0.00657861\\
0.532	-0.00652979\\
0.534	-0.00646463\\
0.536	-0.00638353\\
0.538	-0.00628694\\
0.54	-0.00617534\\
0.542	-0.00604926\\
0.544	-0.00590923\\
0.546	-0.00575582\\
0.548	-0.00558963\\
0.55	-0.00541129\\
0.552	-0.00522142\\
0.554	-0.0050207\\
0.556	-0.00488932\\
0.558	-0.00483463\\
0.56	-0.00476442\\
0.562	-0.00498734\\
0.564	-0.00541993\\
0.566	-0.00582678\\
0.568	-0.0062068\\
0.57	-0.00655902\\
0.572	-0.00688256\\
0.574	-0.00717665\\
0.576	-0.00744066\\
0.578	-0.00767407\\
0.58	-0.00787645\\
0.582	-0.00804752\\
0.584	-0.00818709\\
0.586	-0.0082951\\
0.588	-0.0083716\\
0.59	-0.00841674\\
0.592	-0.00843081\\
0.594	-0.00841418\\
0.596	-0.00836732\\
0.598	-0.00829082\\
0.6	-0.00818537\\
0.602	-0.00805174\\
0.604	-0.00789079\\
0.606	-0.00770348\\
0.608	-0.00749084\\
0.61	-0.00725398\\
0.612	-0.00699408\\
0.614	-0.00671238\\
0.616	-0.0064102\\
0.618	-0.00608889\\
0.62	-0.00574988\\
0.622	-0.00539462\\
0.624	-0.00502461\\
0.626	-0.00464137\\
0.628	-0.00424646\\
0.63	-0.00384146\\
0.632	-0.0036928\\
0.634	-0.00379369\\
0.636	-0.00388326\\
0.638	-0.00396148\\
0.64	-0.00402831\\
0.642	-0.00408378\\
0.644	-0.00412792\\
0.646	-0.00416079\\
0.648	-0.0041825\\
0.65	-0.00419317\\
0.652	-0.00419293\\
0.654	-0.00418198\\
0.656	-0.00416049\\
0.658	-0.00412871\\
0.66	-0.00408686\\
0.662	-0.00403522\\
0.664	-0.00397406\\
0.666	-0.00390371\\
0.668	-0.00382448\\
0.67	-0.003887\\
0.672	-0.00416483\\
0.674	-0.00442371\\
0.676	-0.00466296\\
0.678	-0.00488202\\
0.68	-0.00508036\\
0.682	-0.00525758\\
0.684	-0.00541332\\
0.686	-0.00554732\\
0.688	-0.00565942\\
0.69	-0.0057495\\
0.692	-0.00581755\\
0.694	-0.00586363\\
0.696	-0.00588788\\
0.698	-0.00589051\\
0.7	-0.00587182\\
0.702	-0.00583217\\
0.704	-0.00577199\\
0.706	-0.00569178\\
0.708	-0.0055921\\
0.71	-0.00547359\\
0.712	-0.00533693\\
0.714	-0.00518286\\
0.716	-0.00501218\\
0.718	-0.00482573\\
0.72	-0.0046244\\
0.722	-0.00440912\\
0.724	-0.00418086\\
0.726	-0.00394062\\
0.728	-0.00368943\\
0.73	-0.00342834\\
0.732	-0.00315843\\
0.734	-0.0028808\\
0.736	-0.00259655\\
0.738	-0.0023068\\
0.74	-0.00201266\\
0.742	-0.00191236\\
0.744	-0.00200972\\
0.746	-0.00213925\\
0.748	-0.00226066\\
0.75	-0.00237366\\
0.752	-0.00247796\\
0.754	-0.00257333\\
0.756	-0.00265953\\
0.758	-0.0027364\\
0.76	-0.00280378\\
0.762	-0.00286154\\
0.764	-0.0029096\\
0.766	-0.00294789\\
0.768	-0.00297639\\
0.77	-0.00299511\\
0.772	-0.00300408\\
0.774	-0.00300336\\
0.776	-0.00299306\\
0.778	-0.00297329\\
0.78	-0.00314426\\
0.782	-0.00330387\\
0.784	-0.00344923\\
0.786	-0.00358004\\
0.788	-0.003696\\
0.79	-0.00379691\\
0.792	-0.00388262\\
0.794	-0.00395302\\
0.796	-0.00400807\\
0.798	-0.00404779\\
0.8	-0.00407224\\
0.802	-0.00408155\\
0.804	-0.00407589\\
0.806	-0.00405549\\
0.808	-0.00402062\\
0.81	-0.00397162\\
0.812	-0.00390885\\
0.814	-0.00383274\\
0.816	-0.00374374\\
0.818	-0.00364235\\
0.82	-0.00352911\\
0.822	-0.0034046\\
0.824	-0.00326942\\
0.826	-0.00312421\\
0.828	-0.00296963\\
0.83	-0.00280638\\
0.832	-0.00263516\\
0.834	-0.00245671\\
0.836	-0.00227177\\
0.838	-0.00208111\\
0.84	-0.00188548\\
0.842	-0.00168569\\
0.844	-0.00148249\\
0.846	-0.00127669\\
0.848	-0.00117076\\
0.85	-0.00123057\\
0.852	-0.00134494\\
0.854	-0.00145434\\
0.856	-0.00155811\\
0.858	-0.0016559\\
0.86	-0.00174747\\
0.862	-0.00183259\\
0.864	-0.00191105\\
0.866	-0.00198267\\
0.868	-0.00204729\\
0.87	-0.00210477\\
0.872	-0.00215502\\
0.874	-0.00219793\\
0.876	-0.00223344\\
0.878	-0.00226152\\
0.88	-0.00228216\\
0.882	-0.00229535\\
0.884	-0.00230114\\
0.886	-0.00229958\\
0.888	-0.0023657\\
0.89	-0.0024597\\
0.892	-0.00254341\\
0.894	-0.00261664\\
0.896	-0.00267928\\
0.898	-0.00273125\\
0.9	-0.0027725\\
0.902	-0.00280302\\
0.904	-0.00282285\\
0.906	-0.00283206\\
0.908	-0.00283077\\
0.91	-0.00281912\\
0.912	-0.00279729\\
0.914	-0.0027655\\
0.916	-0.00272401\\
0.918	-0.00267309\\
0.92	-0.00261306\\
0.922	-0.00254426\\
0.924	-0.00246706\\
0.926	-0.00238186\\
0.928	-0.00228908\\
0.93	-0.00218915\\
0.932	-0.00208254\\
0.934	-0.00196973\\
0.936	-0.00185122\\
0.938	-0.0017275\\
0.94	-0.00159912\\
0.942	-0.00146659\\
0.944	-0.00133046\\
0.946	-0.00119128\\
0.948	-0.00104959\\
0.95	-0.000905947\\
0.952	-0.000810332\\
0.954	-0.000804252\\
0.956	-0.000795729\\
0.958	-0.00078482\\
0.96	-0.000839175\\
0.962	-0.000929007\\
0.964	-0.00101526\\
0.966	-0.00109741\\
0.968	-0.00117522\\
0.97	-0.0012485\\
0.972	-0.00131707\\
0.974	-0.00138076\\
0.976	-0.00143941\\
0.978	-0.0014929\\
0.98	-0.0015411\\
0.982	-0.00158392\\
0.984	-0.00162129\\
0.986	-0.00165313\\
0.988	-0.00167941\\
0.99	-0.0017001\\
0.992	-0.00171519\\
0.994	-0.00172575\\
0.996	-0.00178503\\
0.998	-0.00183694\\
1	-0.00188139\\
1.002	-0.00191833\\
1.004	-0.0019477\\
1.006	-0.00196951\\
1.008	-0.00198377\\
1.01	-0.00199052\\
1.012	-0.00198985\\
1.014	-0.00198186\\
1.016	-0.00196666\\
1.018	-0.0019444\\
1.02	-0.00191527\\
1.022	-0.00187945\\
1.024	-0.00183716\\
1.026	-0.00178866\\
1.028	-0.00173418\\
1.03	-0.00167402\\
1.032	-0.00160846\\
1.034	-0.00153781\\
1.036	-0.00146241\\
1.038	-0.00138259\\
1.04	-0.0012987\\
1.042	-0.00121109\\
1.044	-0.00112015\\
1.046	-0.00102625\\
1.048	-0.000929761\\
1.05	-0.000831082\\
1.052	-0.000730602\\
1.054	-0.000665833\\
1.056	-0.000646017\\
1.058	-0.000625064\\
1.06	-0.00060305\\
1.062	-0.000580053\\
1.064	-0.000558115\\
1.066	-0.000554581\\
1.068	-0.000549374\\
1.07	-0.000549036\\
1.072	-0.000618198\\
1.074	-0.000684756\\
1.076	-0.000748408\\
1.078	-0.000808983\\
1.08	-0.000866323\\
1.082	-0.000920281\\
1.084	-0.000970724\\
1.086	-0.00101753\\
1.088	-0.00106059\\
1.09	-0.00109981\\
1.092	-0.00113511\\
1.094	-0.00116642\\
1.096	-0.00119478\\
1.098	-0.00124188\\
1.1	-0.0012838\\
1.102	-0.00132047\\
1.104	-0.0013518\\
1.106	-0.00137775\\
1.108	-0.0013983\\
1.11	-0.00141343\\
1.112	-0.00142316\\
1.114	-0.00142751\\
1.116	-0.00142655\\
1.118	-0.00142033\\
1.12	-0.00140895\\
1.122	-0.00139252\\
1.124	-0.00137115\\
1.126	-0.00134498\\
1.128	-0.00131418\\
1.13	-0.00127891\\
1.132	-0.00123935\\
1.134	-0.00119571\\
1.136	-0.00114819\\
1.138	-0.00109702\\
1.14	-0.00104243\\
1.142	-0.000984667\\
1.144	-0.000923972\\
1.146	-0.00086061\\
1.148	-0.000794845\\
1.15	-0.000726949\\
1.152	-0.000657196\\
1.154	-0.000585865\\
1.156	-0.000513237\\
1.158	-0.000439595\\
1.16	-0.00039771\\
1.162	-0.0003985\\
1.164	-0.000398405\\
1.166	-0.000397446\\
1.168	-0.000395644\\
1.17	-0.000393023\\
1.172	-0.000389607\\
1.174	-0.000385421\\
1.176	-0.000380492\\
1.178	-0.000374847\\
1.18	-0.000368515\\
1.182	-0.00042221\\
1.184	-0.000484591\\
1.186	-0.000544446\\
1.188	-0.000601585\\
1.19	-0.000655833\\
1.192	-0.000707023\\
1.194	-0.000755007\\
1.196	-0.000799645\\
1.198	-0.000840815\\
1.2	-0.000878407\\
1.202	-0.000912326\\
1.204	-0.000942492\\
1.206	-0.000968839\\
1.208	-0.000991316\\
1.21	-0.00100989\\
1.212	-0.00102453\\
1.214	-0.00103524\\
1.216	-0.00104202\\
1.218	-0.00104489\\
1.22	-0.0010439\\
1.222	-0.00103908\\
1.224	-0.00103051\\
1.226	-0.00101825\\
1.228	-0.0010024\\
1.23	-0.000983063\\
1.232	-0.000960341\\
1.234	-0.00093436\\
1.236	-0.000905254\\
1.238	-0.000873166\\
1.24	-0.00083825\\
1.242	-0.000800664\\
1.244	-0.00076058\\
1.246	-0.000718171\\
1.248	-0.00067362\\
1.25	-0.000627116\\
1.252	-0.000585322\\
1.254	-0.000547553\\
1.256	-0.000508498\\
1.258	-0.000468298\\
1.26	-0.000427095\\
1.262	-0.000385033\\
1.264	-0.000342257\\
1.266	-0.000298913\\
1.268	-0.000255145\\
1.27	-0.0002111\\
1.272	-0.000197794\\
1.274	-0.000202566\\
1.276	-0.000206601\\
1.278	-0.0002099\\
1.28	-0.000212466\\
1.282	-0.00021886\\
1.284	-0.000268026\\
1.286	-0.000315653\\
1.288	-0.000361584\\
1.29	-0.000405673\\
1.292	-0.000447778\\
1.294	-0.000487771\\
1.296	-0.00052553\\
1.298	-0.000560942\\
1.3	-0.000593905\\
1.302	-0.000624329\\
1.304	-0.000652131\\
1.306	-0.000677239\\
1.308	-0.000699595\\
1.31	-0.000719147\\
1.312	-0.000735856\\
1.314	-0.000749696\\
1.316	-0.000760647\\
1.318	-0.000768703\\
1.32	-0.000773869\\
1.322	-0.000776158\\
1.324	-0.000775595\\
1.326	-0.000772214\\
1.328	-0.00076606\\
1.33	-0.000757187\\
1.332	-0.000745659\\
1.334	-0.000731547\\
1.336	-0.000714932\\
1.338	-0.000695903\\
1.34	-0.000674556\\
1.342	-0.000650997\\
1.344	-0.000625334\\
1.346	-0.000597686\\
1.348	-0.00057895\\
1.35	-0.000562302\\
1.352	-0.000544063\\
1.354	-0.000524312\\
1.356	-0.000503131\\
1.358	-0.000480606\\
1.36	-0.000456827\\
1.362	-0.000431885\\
1.364	-0.000405876\\
1.366	-0.000378896\\
1.368	-0.000351045\\
1.37	-0.000322424\\
1.372	-0.000293133\\
1.374	-0.000263277\\
1.376	-0.000232957\\
1.378	-0.000202278\\
1.38	-0.000171342\\
1.382	-0.000140253\\
1.384	-0.000120366\\
1.386	-0.000158075\\
1.388	-0.000194791\\
1.39	-0.000230395\\
1.392	-0.000264768\\
1.394	-0.0002978\\
1.396	-0.000329388\\
1.398	-0.000359433\\
1.4	-0.000387843\\
1.402	-0.000414535\\
1.404	-0.000439431\\
1.406	-0.000462461\\
1.408	-0.000483563\\
1.41	-0.000502683\\
1.412	-0.000519774\\
1.414	-0.000534796\\
1.416	-0.000547719\\
1.418	-0.000558519\\
1.42	-0.000567182\\
1.422	-0.0005737\\
1.424	-0.000578074\\
1.426	-0.000580312\\
1.428	-0.000580429\\
1.43	-0.00057845\\
1.432	-0.000574405\\
1.434	-0.000568331\\
1.436	-0.000560274\\
1.438	-0.000550284\\
1.44	-0.00053842\\
1.442	-0.000524746\\
1.444	-0.00050933\\
1.446	-0.000492248\\
1.448	-0.00047358\\
1.45	-0.000453412\\
1.452	-0.000434998\\
1.454	-0.000427734\\
1.456	-0.000419186\\
1.458	-0.000409399\\
1.46	-0.000398421\\
1.462	-0.000386302\\
1.464	-0.000373097\\
1.466	-0.000358862\\
1.468	-0.000343658\\
1.47	-0.000327545\\
1.472	-0.000310587\\
1.474	-0.000292852\\
1.476	-0.000274405\\
1.478	-0.000255316\\
1.48	-0.000235655\\
1.482	-0.000215494\\
1.484	-0.000194903\\
1.486	-0.000173956\\
1.488	-0.000152724\\
1.49	-0.00013128\\
1.492	-0.000135546\\
1.494	-0.000162411\\
1.496	-0.000188396\\
1.498	-0.000213415\\
1.5	-0.000237391\\
1.502	-0.000260248\\
1.504	-0.000281918\\
1.506	-0.000302335\\
1.508	-0.000321441\\
1.51	-0.00033918\\
1.512	-0.000355506\\
1.514	-0.000370374\\
1.516	-0.000383748\\
1.518	-0.000395597\\
1.52	-0.000405895\\
1.522	-0.000414622\\
1.524	-0.000421766\\
1.526	-0.000427317\\
1.528	-0.000431273\\
1.53	-0.00043364\\
1.532	-0.000434424\\
1.534	-0.000433643\\
1.536	-0.000431315\\
1.538	-0.000427467\\
1.54	-0.000422129\\
1.542	-0.000415338\\
1.544	-0.000407133\\
1.546	-0.000397561\\
1.548	-0.00038667\\
1.55	-0.000374515\\
1.552	-0.000361153\\
1.554	-0.000346646\\
1.556	-0.000331058\\
1.558	-0.000314458\\
1.56	-0.000305471\\
1.562	-0.000301772\\
1.564	-0.000297146\\
1.566	-0.000291618\\
1.568	-0.000285221\\
1.57	-0.000277985\\
1.572	-0.000269946\\
1.574	-0.00026114\\
1.576	-0.000251607\\
1.578	-0.000241387\\
1.58	-0.000230523\\
1.582	-0.000219057\\
1.584	-0.000207037\\
1.586	-0.000194508\\
1.588	-0.000181517\\
1.59	-0.000168114\\
1.592	-0.000154347\\
1.594	-0.000140265\\
1.596	-0.00012592\\
1.598	-0.00011136\\
1.6	-0.000129822\\
1.602	-0.000148857\\
1.604	-0.000167148\\
1.606	-0.000184637\\
1.608	-0.000201272\\
1.61	-0.000217002\\
1.612	-0.000231781\\
1.614	-0.000245567\\
1.616	-0.000258322\\
1.618	-0.000270012\\
1.62	-0.000280607\\
1.622	-0.000290081\\
1.624	-0.000298414\\
1.626	-0.000305588\\
1.628	-0.000311591\\
1.63	-0.000316414\\
1.632	-0.000320054\\
1.634	-0.000322511\\
1.636	-0.000323789\\
1.638	-0.000323897\\
1.64	-0.000322847\\
1.642	-0.000320658\\
1.644	-0.000317348\\
1.646	-0.000312943\\
1.648	-0.000307471\\
1.65	-0.000300963\\
1.652	-0.000293455\\
1.654	-0.000284983\\
1.656	-0.00027559\\
1.658	-0.000265319\\
1.66	-0.000254217\\
1.662	-0.000242332\\
1.664	-0.000229717\\
1.666	-0.000216424\\
1.668	-0.000211441\\
1.67	-0.000209781\\
1.672	-0.000207462\\
1.674	-0.000204501\\
1.676	-0.000200916\\
1.678	-0.000196728\\
1.68	-0.000191957\\
1.682	-0.000186627\\
1.684	-0.000180765\\
1.686	-0.000174395\\
1.688	-0.000167546\\
1.69	-0.000160246\\
1.692	-0.000152525\\
1.694	-0.000144415\\
1.696	-0.000135947\\
1.698	-0.000127153\\
1.7	-0.000118067\\
1.702	-0.000108722\\
1.704	-9.91522e-05\\
1.706	-0.000101093\\
1.708	-0.000115051\\
1.71	-0.000128442\\
1.712	-0.000141225\\
1.714	-0.00015336\\
1.716	-0.000164813\\
1.718	-0.00017555\\
1.72	-0.000185541\\
1.722	-0.000194759\\
1.724	-0.000203179\\
1.726	-0.000210781\\
1.728	-0.000217547\\
1.73	-0.000223462\\
1.732	-0.000228515\\
1.734	-0.000232697\\
1.736	-0.000236004\\
1.738	-0.000238434\\
1.74	-0.000239987\\
1.742	-0.000240669\\
1.744	-0.000240487\\
1.746	-0.00023945\\
1.748	-0.000237573\\
1.75	-0.000234873\\
1.752	-0.000231366\\
1.754	-0.000227076\\
1.756	-0.000222027\\
1.758	-0.000216244\\
1.76	-0.000209758\\
1.762	-0.000202599\\
1.764	-0.000194799\\
1.766	-0.000186395\\
1.768	-0.000177423\\
1.77	-0.000167921\\
1.772	-0.000157929\\
1.774	-0.000147489\\
1.776	-0.000144767\\
1.778	-0.000144232\\
1.78	-0.000143236\\
1.782	-0.000141788\\
1.784	-0.000139899\\
1.786	-0.000137582\\
1.788	-0.00013485\\
1.79	-0.000131717\\
1.792	-0.000128199\\
1.794	-0.000124314\\
1.796	-0.000120079\\
1.798	-0.000115512\\
1.8	-0.000110634\\
1.802	-0.000105464\\
1.804	-0.000100023\\
1.806	-9.43325e-05\\
1.808	-8.84151e-05\\
1.81	-8.22928e-05\\
1.812	-7.75133e-05\\
1.814	-8.77476e-05\\
1.816	-9.75541e-05\\
1.818	-0.000106903\\
1.82	-0.000115765\\
1.822	-0.000124116\\
1.824	-0.000131931\\
1.826	-0.000139189\\
1.828	-0.00014587\\
1.83	-0.000151956\\
1.832	-0.000157434\\
1.834	-0.00016229\\
1.836	-0.000166515\\
1.838	-0.0001701\\
1.84	-0.00017304\\
1.842	-0.000175331\\
1.844	-0.000176974\\
1.846	-0.00017797\\
1.848	-0.000178322\\
1.85	-0.000178037\\
1.852	-0.000177123\\
1.854	-0.00017559\\
1.856	-0.00017345\\
1.858	-0.000170719\\
1.86	-0.000167413\\
1.862	-0.000163549\\
1.864	-0.000159149\\
1.866	-0.000154234\\
1.868	-0.000148826\\
1.87	-0.000142952\\
1.872	-0.000136636\\
1.874	-0.000129908\\
1.876	-0.000122794\\
1.878	-0.000115325\\
1.88	-0.000107532\\
1.882	-9.94453e-05\\
1.884	-9.80939e-05\\
1.886	-9.81549e-05\\
1.888	-9.78967e-05\\
1.89	-9.73247e-05\\
1.892	-9.64452e-05\\
1.894	-9.52653e-05\\
1.896	-9.37932e-05\\
1.898	-9.20379e-05\\
1.9	-9.0009e-05\\
1.902	-8.7717e-05\\
1.904	-8.5173e-05\\
1.906	-8.2389e-05\\
1.908	-7.93774e-05\\
1.91	-7.61511e-05\\
1.912	-7.27238e-05\\
1.914	-6.91093e-05\\
1.916	-6.53221e-05\\
1.918	-6.1377e-05\\
1.92	-6.63668e-05\\
1.922	-7.35523e-05\\
1.924	-8.03951e-05\\
1.926	-8.68749e-05\\
1.928	-9.2973e-05\\
1.93	-9.86719e-05\\
1.932	-0.000103956\\
1.934	-0.000108811\\
1.936	-0.000113226\\
1.938	-0.000117188\\
1.94	-0.00012069\\
1.942	-0.000123724\\
1.944	-0.000126285\\
1.946	-0.000128369\\
1.948	-0.000129974\\
1.95	-0.0001311\\
1.952	-0.000131748\\
1.954	-0.000131922\\
1.956	-0.000131626\\
1.958	-0.000130867\\
1.96	-0.000129652\\
1.962	-0.000127992\\
1.964	-0.000125897\\
1.966	-0.00012338\\
1.968	-0.000120454\\
1.97	-0.000117134\\
1.972	-0.000113437\\
1.974	-0.00010938\\
1.976	-0.000104982\\
1.978	-0.000100261\\
1.98	-9.52393e-05\\
1.982	-8.99369e-05\\
1.984	-8.43762e-05\\
1.986	-7.85798e-05\\
1.988	-7.2571e-05\\
1.99	-6.63734e-05\\
1.992	-6.55902e-05\\
1.994	-6.59445e-05\\
1.996	-6.6081e-05\\
1.998	-6.60024e-05\\
2	-6.57121e-05\\
};
\addplot [color=mycolor9, line width=1.0pt, forget plot]
  table[row sep=crcr]{%
-0.024	0\\
-0.022	0\\
-0.02	0\\
-0.018	0\\
-0.016	0\\
-0.014	0\\
-0.012	0\\
-0.01	0\\
-0.008	0\\
-0.006	0\\
-0.004	0\\
-0.002	0\\
0	0\\
0.002	-3.1616e-07\\
0.004	-3.21585e-06\\
0.006	-1.7405e-05\\
0.008	-4.68068e-05\\
0.01	-9.90896e-05\\
0.012	-0.000179703\\
0.014	-0.00029228\\
0.016	-0.00043832\\
0.018	-0.000629209\\
0.02	-0.00181929\\
0.022	-0.00326677\\
0.024	-0.00514553\\
0.026	-0.00772145\\
0.028	-0.00873144\\
0.03	-0.00976122\\
0.032	-0.0108024\\
0.034	-0.0118465\\
0.036	-0.0128853\\
0.038	-0.0139105\\
0.04	-0.0149139\\
0.042	-0.0158877\\
0.044	-0.016824\\
0.046	-0.0177153\\
0.048	-0.0185546\\
0.05	-0.0193349\\
0.052	-0.0200498\\
0.054	-0.0206931\\
0.056	-0.0212593\\
0.058	-0.0217433\\
0.06	-0.0221403\\
0.062	-0.0224463\\
0.064	-0.0226577\\
0.066	-0.0227714\\
0.068	-0.0227851\\
0.07	-0.0226968\\
0.072	-0.0225054\\
0.074	-0.0222101\\
0.076	-0.0218109\\
0.078	-0.0213082\\
0.08	-0.0207031\\
0.082	-0.0199973\\
0.084	-0.0191929\\
0.086	-0.0182926\\
0.088	-0.0172997\\
0.09	-0.0162177\\
0.092	-0.0150508\\
0.094	-0.0138035\\
0.096	-0.0124809\\
0.098	-0.0110883\\
0.1	-0.00963123\\
0.102	-0.00811574\\
0.104	-0.00677517\\
0.106	-0.00723286\\
0.108	-0.00763106\\
0.11	-0.00796935\\
0.112	-0.00824758\\
0.114	-0.00846587\\
0.116	-0.00832595\\
0.118	-0.00789055\\
0.12	-0.00742571\\
0.122	-0.00735752\\
0.124	-0.00794046\\
0.126	-0.00858164\\
0.128	-0.00919181\\
0.13	-0.00976923\\
0.132	-0.0103123\\
0.134	-0.0107616\\
0.136	-0.0108979\\
0.138	-0.0110086\\
0.14	-0.0110929\\
0.142	-0.0116568\\
0.144	-0.013478\\
0.146	-0.0151967\\
0.148	-0.0168084\\
0.15	-0.0183093\\
0.152	-0.0196956\\
0.154	-0.0209644\\
0.156	-0.0221131\\
0.158	-0.0231394\\
0.16	-0.0240419\\
0.162	-0.0248194\\
0.164	-0.025471\\
0.166	-0.0259968\\
0.168	-0.0263968\\
0.17	-0.0266718\\
0.172	-0.0268229\\
0.174	-0.0268518\\
0.176	-0.0267604\\
0.178	-0.0265511\\
0.18	-0.0262268\\
0.182	-0.0257906\\
0.184	-0.0252461\\
0.186	-0.0245152\\
0.188	-0.0236664\\
0.19	-0.0227207\\
0.192	-0.0216833\\
0.194	-0.0205596\\
0.196	-0.0193551\\
0.198	-0.0180757\\
0.2	-0.0167274\\
0.202	-0.0153163\\
0.204	-0.0138486\\
0.206	-0.0123309\\
0.208	-0.0107696\\
0.21	-0.00923384\\
0.212	-0.00931603\\
0.214	-0.00935426\\
0.216	-0.00934932\\
0.218	-0.00930207\\
0.22	-0.00921349\\
0.222	-0.00908466\\
0.224	-0.00891675\\
0.226	-0.00871103\\
0.228	-0.00846886\\
0.23	-0.00819166\\
0.232	-0.00788095\\
0.234	-0.00753833\\
0.236	-0.00716544\\
0.238	-0.00676401\\
0.24	-0.00749308\\
0.242	-0.00896255\\
0.244	-0.010368\\
0.246	-0.0117055\\
0.248	-0.0129713\\
0.25	-0.0141621\\
0.252	-0.0152748\\
0.254	-0.0163066\\
0.256	-0.0172552\\
0.258	-0.0181184\\
0.26	-0.0188944\\
0.262	-0.0195818\\
0.264	-0.0201796\\
0.266	-0.0206868\\
0.268	-0.0211031\\
0.27	-0.0214284\\
0.272	-0.0216629\\
0.274	-0.0218071\\
0.276	-0.0218617\\
0.278	-0.0218281\\
0.28	-0.0217076\\
0.282	-0.0215019\\
0.284	-0.0212131\\
0.286	-0.0208435\\
0.288	-0.0203955\\
0.29	-0.0198719\\
0.292	-0.0192759\\
0.294	-0.0186105\\
0.296	-0.0178792\\
0.298	-0.0170857\\
0.3	-0.0162337\\
0.302	-0.0153272\\
0.304	-0.0143702\\
0.306	-0.0133671\\
0.308	-0.012322\\
0.31	-0.0112395\\
0.312	-0.010124\\
0.314	-0.00897998\\
0.316	-0.00781214\\
0.318	-0.00770243\\
0.32	-0.00783267\\
0.322	-0.00793149\\
0.324	-0.00799907\\
0.326	-0.00803566\\
0.328	-0.00804162\\
0.33	-0.00801738\\
0.332	-0.00796349\\
0.334	-0.00788054\\
0.336	-0.00776922\\
0.338	-0.00763031\\
0.34	-0.00746464\\
0.342	-0.00727312\\
0.344	-0.00705674\\
0.346	-0.00681653\\
0.348	-0.00778486\\
0.35	-0.00874607\\
0.352	-0.00965858\\
0.354	-0.0105199\\
0.356	-0.0113276\\
0.358	-0.0120798\\
0.36	-0.0127745\\
0.362	-0.01341\\
0.364	-0.0139851\\
0.366	-0.0144985\\
0.368	-0.0149492\\
0.37	-0.0153366\\
0.372	-0.0156603\\
0.374	-0.0159198\\
0.376	-0.0161154\\
0.378	-0.0162471\\
0.38	-0.0163155\\
0.382	-0.0163212\\
0.384	-0.016265\\
0.386	-0.016148\\
0.388	-0.0159716\\
0.39	-0.0157372\\
0.392	-0.0154465\\
0.394	-0.0151013\\
0.396	-0.0147037\\
0.398	-0.0142557\\
0.4	-0.0137597\\
0.402	-0.0132182\\
0.404	-0.0126338\\
0.406	-0.0120091\\
0.408	-0.011347\\
0.41	-0.0106504\\
0.412	-0.00992221\\
0.414	-0.00916558\\
0.416	-0.00895618\\
0.418	-0.00875347\\
0.42	-0.00852936\\
0.422	-0.00828474\\
0.424	-0.00802052\\
0.426	-0.00773767\\
0.428	-0.00743719\\
0.43	-0.00712012\\
0.432	-0.0067875\\
0.434	-0.00644044\\
0.436	-0.00631857\\
0.438	-0.00634211\\
0.44	-0.00634338\\
0.442	-0.00632264\\
0.444	-0.00628023\\
0.446	-0.00621655\\
0.448	-0.00613205\\
0.45	-0.00602727\\
0.452	-0.00590277\\
0.454	-0.00609682\\
0.456	-0.00675836\\
0.458	-0.00738535\\
0.46	-0.00797607\\
0.462	-0.00852892\\
0.464	-0.00904246\\
0.466	-0.00951542\\
0.468	-0.00994667\\
0.47	-0.0103353\\
0.472	-0.0106804\\
0.474	-0.0109815\\
0.476	-0.011238\\
0.478	-0.0114497\\
0.48	-0.0116164\\
0.482	-0.0117382\\
0.484	-0.0118152\\
0.486	-0.0118479\\
0.488	-0.0118366\\
0.49	-0.011782\\
0.492	-0.011685\\
0.494	-0.0115464\\
0.496	-0.0113673\\
0.498	-0.011149\\
0.5	-0.0108926\\
0.502	-0.0105998\\
0.504	-0.010272\\
0.506	-0.00991081\\
0.508	-0.00951809\\
0.51	-0.00909562\\
0.512	-0.00864533\\
0.514	-0.00816921\\
0.516	-0.00766933\\
0.518	-0.00714779\\
0.52	-0.00660676\\
0.522	-0.00659465\\
0.524	-0.00661476\\
0.526	-0.00661724\\
0.528	-0.00660234\\
0.53	-0.00657036\\
0.532	-0.00652165\\
0.534	-0.00645659\\
0.536	-0.00637559\\
0.538	-0.0062791\\
0.54	-0.0061676\\
0.542	-0.00604162\\
0.544	-0.00590168\\
0.546	-0.00574837\\
0.548	-0.00558227\\
0.55	-0.00540402\\
0.552	-0.00521425\\
0.554	-0.00501361\\
0.556	-0.00488232\\
0.558	-0.00482772\\
0.56	-0.0047576\\
0.562	-0.0049806\\
0.564	-0.00541327\\
0.566	-0.0058202\\
0.568	-0.00620031\\
0.57	-0.00655261\\
0.572	-0.00687623\\
0.574	-0.0071704\\
0.576	-0.00743449\\
0.578	-0.00766797\\
0.58	-0.00787043\\
0.582	-0.00804158\\
0.584	-0.00818122\\
0.586	-0.00828931\\
0.588	-0.00836588\\
0.59	-0.0084111\\
0.592	-0.00842524\\
0.594	-0.00840867\\
0.596	-0.00836189\\
0.598	-0.00828546\\
0.6	-0.00818008\\
0.602	-0.00804652\\
0.604	-0.00788564\\
0.606	-0.00769839\\
0.608	-0.00748582\\
0.61	-0.00724902\\
0.612	-0.00698918\\
0.614	-0.00670755\\
0.616	-0.00640543\\
0.618	-0.00608419\\
0.62	-0.00574524\\
0.622	-0.00539004\\
0.624	-0.00502009\\
0.626	-0.00463691\\
0.628	-0.00424206\\
0.63	-0.00383712\\
0.632	-0.00368851\\
0.634	-0.00378946\\
0.636	-0.00387909\\
0.638	-0.00395737\\
0.64	-0.00402426\\
0.642	-0.00407978\\
0.644	-0.00412397\\
0.646	-0.0041569\\
0.648	-0.00417866\\
0.65	-0.00418938\\
0.652	-0.0041892\\
0.654	-0.00417829\\
0.656	-0.00415686\\
0.658	-0.00412512\\
0.66	-0.00408333\\
0.662	-0.00403173\\
0.664	-0.00397063\\
0.666	-0.00390033\\
0.668	-0.00382114\\
0.67	-0.00388371\\
0.672	-0.00416158\\
0.674	-0.00442051\\
0.676	-0.00465981\\
0.678	-0.00487891\\
0.68	-0.0050773\\
0.682	-0.00525456\\
0.684	-0.00541034\\
0.686	-0.00554439\\
0.688	-0.00565652\\
0.69	-0.00574665\\
0.692	-0.00581474\\
0.694	-0.00586086\\
0.696	-0.00588515\\
0.698	-0.00588783\\
0.7	-0.00586918\\
0.702	-0.00582956\\
0.704	-0.00576942\\
0.706	-0.00568925\\
0.708	-0.00558961\\
0.71	-0.00547113\\
0.712	-0.00533451\\
0.714	-0.00518048\\
0.716	-0.00500983\\
0.718	-0.00482342\\
0.72	-0.00462212\\
0.722	-0.00440688\\
0.724	-0.00417865\\
0.726	-0.00393844\\
0.728	-0.00368728\\
0.73	-0.00342622\\
0.732	-0.00315635\\
0.734	-0.00287875\\
0.736	-0.00259453\\
0.738	-0.00230481\\
0.74	-0.0020107\\
0.742	-0.00191043\\
0.744	-0.00200782\\
0.746	-0.00213737\\
0.748	-0.00225882\\
0.75	-0.00237184\\
0.752	-0.00247617\\
0.754	-0.00257157\\
0.756	-0.0026578\\
0.758	-0.0027347\\
0.76	-0.0028021\\
0.762	-0.00285989\\
0.764	-0.00290797\\
0.766	-0.00294629\\
0.768	-0.00297482\\
0.77	-0.00299356\\
0.772	-0.00300255\\
0.774	-0.00300186\\
0.776	-0.00299158\\
0.778	-0.00297183\\
0.78	-0.00314282\\
0.782	-0.00330245\\
0.784	-0.00344784\\
0.786	-0.00357867\\
0.788	-0.00369465\\
0.79	-0.00379559\\
0.792	-0.00388132\\
0.794	-0.00395174\\
0.796	-0.00400681\\
0.798	-0.00404655\\
0.8	-0.00407102\\
0.802	-0.00408035\\
0.804	-0.00407471\\
0.806	-0.00405432\\
0.808	-0.00401948\\
0.81	-0.00397049\\
0.812	-0.00390774\\
0.814	-0.00383165\\
0.816	-0.00374266\\
0.818	-0.00364129\\
0.82	-0.00352807\\
0.822	-0.00340358\\
0.824	-0.00326841\\
0.826	-0.00312322\\
0.828	-0.00296866\\
0.83	-0.00280542\\
0.832	-0.00263422\\
0.834	-0.00245578\\
0.836	-0.00227086\\
0.838	-0.00208021\\
0.84	-0.0018846\\
0.842	-0.00168481\\
0.844	-0.00148164\\
0.846	-0.00127585\\
0.848	-0.00116993\\
0.85	-0.00122975\\
0.852	-0.00134413\\
0.854	-0.00145355\\
0.856	-0.00155733\\
0.858	-0.00165513\\
0.86	-0.00174672\\
0.862	-0.00183184\\
0.864	-0.00191032\\
0.866	-0.00198195\\
0.868	-0.00204658\\
0.87	-0.00210408\\
0.872	-0.00215433\\
0.874	-0.00219725\\
0.876	-0.00223278\\
0.878	-0.00226087\\
0.88	-0.00228152\\
0.882	-0.00229472\\
0.884	-0.00230052\\
0.886	-0.00229897\\
0.888	-0.0023651\\
0.89	-0.00245911\\
0.892	-0.00254283\\
0.894	-0.00261607\\
0.896	-0.00267872\\
0.898	-0.0027307\\
0.9	-0.00277195\\
0.902	-0.00280248\\
0.904	-0.00282232\\
0.906	-0.00283154\\
0.908	-0.00283026\\
0.91	-0.00281862\\
0.912	-0.0027968\\
0.914	-0.00276502\\
0.916	-0.00272353\\
0.918	-0.00267262\\
0.92	-0.0026126\\
0.922	-0.00254381\\
0.924	-0.00246662\\
0.926	-0.00238142\\
0.928	-0.00228865\\
0.93	-0.00218872\\
0.932	-0.00208212\\
0.934	-0.00196932\\
0.936	-0.00185081\\
0.938	-0.0017271\\
0.94	-0.00159872\\
0.942	-0.0014662\\
0.944	-0.00133008\\
0.946	-0.0011909\\
0.948	-0.00104922\\
0.95	-0.000905584\\
0.952	-0.000809975\\
0.954	-0.000803901\\
0.956	-0.000795383\\
0.958	-0.00078448\\
0.96	-0.000838841\\
0.962	-0.000928678\\
0.964	-0.00101494\\
0.966	-0.00109709\\
0.968	-0.00117491\\
0.97	-0.0012482\\
0.972	-0.00131677\\
0.974	-0.00138046\\
0.976	-0.00143912\\
0.978	-0.00149261\\
0.98	-0.00154082\\
0.982	-0.00158364\\
0.984	-0.00162101\\
0.986	-0.00165286\\
0.988	-0.00167914\\
0.99	-0.00169984\\
0.992	-0.00171494\\
0.994	-0.0017255\\
0.996	-0.00178478\\
0.998	-0.00183669\\
1	-0.00188115\\
1.002	-0.00191809\\
1.004	-0.00194747\\
1.006	-0.00196928\\
1.008	-0.00198354\\
1.01	-0.0019903\\
1.012	-0.00198964\\
1.014	-0.00198164\\
1.016	-0.00196644\\
1.018	-0.00194419\\
1.02	-0.00191506\\
1.022	-0.00187925\\
1.024	-0.00183697\\
1.026	-0.00178846\\
1.028	-0.00173399\\
1.03	-0.00167383\\
1.032	-0.00160827\\
1.034	-0.00153763\\
1.036	-0.00146223\\
1.038	-0.00138241\\
1.04	-0.00129852\\
1.042	-0.00121093\\
1.044	-0.00111999\\
1.046	-0.00102608\\
1.048	-0.0009296\\
1.05	-0.000830924\\
1.052	-0.000730446\\
1.054	-0.00066568\\
1.056	-0.000645867\\
1.058	-0.000624916\\
1.06	-0.000602905\\
1.062	-0.000579911\\
1.064	-0.000557975\\
1.066	-0.000554443\\
1.068	-0.000549238\\
1.07	-0.000548902\\
1.072	-0.000618067\\
1.074	-0.000684627\\
1.076	-0.000748281\\
1.078	-0.000808858\\
1.08	-0.0008662\\
1.082	-0.000920161\\
1.084	-0.000970605\\
1.086	-0.00101741\\
1.088	-0.00106048\\
1.09	-0.0010997\\
1.092	-0.001135\\
1.094	-0.00116631\\
1.096	-0.00119467\\
1.098	-0.00124177\\
1.1	-0.0012837\\
1.102	-0.00132036\\
1.104	-0.0013517\\
1.106	-0.00137765\\
1.108	-0.0013982\\
1.11	-0.00141333\\
1.112	-0.00142306\\
1.114	-0.00142742\\
1.116	-0.00142646\\
1.118	-0.00142024\\
1.12	-0.00140887\\
1.122	-0.00139243\\
1.124	-0.00137106\\
1.126	-0.0013449\\
1.128	-0.0013141\\
1.13	-0.00127883\\
1.132	-0.00123927\\
1.134	-0.00119563\\
1.136	-0.00114812\\
1.138	-0.00109695\\
1.14	-0.00104236\\
1.142	-0.000984595\\
1.144	-0.000923902\\
1.146	-0.000860541\\
1.148	-0.000794777\\
1.15	-0.000726881\\
1.152	-0.000657129\\
1.154	-0.0005858\\
1.156	-0.000513173\\
1.158	-0.000439533\\
1.16	-0.000397648\\
1.162	-0.000398439\\
1.164	-0.000398345\\
1.166	-0.000397387\\
1.168	-0.000395586\\
1.17	-0.000392966\\
1.172	-0.000389551\\
1.174	-0.000385366\\
1.176	-0.000380438\\
1.178	-0.000374794\\
1.18	-0.000368463\\
1.182	-0.000422158\\
1.184	-0.000484541\\
1.186	-0.000544397\\
1.188	-0.000601537\\
1.19	-0.000655785\\
1.192	-0.000706976\\
1.194	-0.00075496\\
1.196	-0.0007996\\
1.198	-0.00084077\\
1.2	-0.000878363\\
1.202	-0.000912283\\
1.204	-0.00094245\\
1.206	-0.000968797\\
1.208	-0.000991275\\
1.21	-0.00100985\\
1.212	-0.00102449\\
1.214	-0.0010352\\
1.216	-0.00104198\\
1.218	-0.00104485\\
1.22	-0.00104386\\
1.222	-0.00103904\\
1.224	-0.00103047\\
1.226	-0.00101822\\
1.228	-0.00100237\\
1.23	-0.000983029\\
1.232	-0.000960307\\
1.234	-0.000934327\\
1.236	-0.000905222\\
1.238	-0.000873135\\
1.24	-0.000838219\\
1.242	-0.000800634\\
1.244	-0.00076055\\
1.246	-0.000718141\\
1.248	-0.000673592\\
1.25	-0.000627088\\
1.252	-0.000585294\\
1.254	-0.000547526\\
1.256	-0.000508471\\
1.258	-0.000468272\\
1.26	-0.000427069\\
1.262	-0.000385008\\
1.264	-0.000342232\\
1.266	-0.000298888\\
1.268	-0.000255121\\
1.27	-0.000211076\\
1.272	-0.00019777\\
1.274	-0.000202543\\
1.276	-0.000206578\\
1.278	-0.000209878\\
1.28	-0.000212444\\
1.282	-0.000218839\\
1.284	-0.000268005\\
1.286	-0.000315632\\
1.288	-0.000361564\\
1.29	-0.000405653\\
1.292	-0.000447759\\
1.294	-0.000487752\\
1.296	-0.000525511\\
1.298	-0.000560923\\
1.3	-0.000593887\\
1.302	-0.000624311\\
1.304	-0.000652113\\
1.306	-0.000677222\\
1.308	-0.000699577\\
1.31	-0.00071913\\
1.312	-0.00073584\\
1.314	-0.000749679\\
1.316	-0.000760631\\
1.318	-0.000768688\\
1.32	-0.000773854\\
1.322	-0.000776143\\
1.324	-0.00077558\\
1.326	-0.000772199\\
1.328	-0.000766046\\
1.33	-0.000757173\\
1.332	-0.000745645\\
1.334	-0.000731533\\
1.336	-0.000714918\\
1.338	-0.000695889\\
1.34	-0.000674543\\
1.342	-0.000650984\\
1.344	-0.000625321\\
1.346	-0.000597674\\
1.348	-0.000578938\\
1.35	-0.00056229\\
1.352	-0.000544052\\
1.354	-0.000524301\\
1.356	-0.00050312\\
1.358	-0.000480595\\
1.36	-0.000456816\\
1.362	-0.000431874\\
1.364	-0.000405865\\
1.366	-0.000378886\\
1.368	-0.000351035\\
1.37	-0.000322414\\
1.372	-0.000293124\\
1.374	-0.000263267\\
1.376	-0.000232948\\
1.378	-0.000202269\\
1.38	-0.000171333\\
1.382	-0.000140244\\
1.384	-0.000120357\\
1.386	-0.000158066\\
1.388	-0.000194783\\
1.39	-0.000230386\\
1.392	-0.000264759\\
1.394	-0.000297792\\
1.396	-0.00032938\\
1.398	-0.000359425\\
1.4	-0.000387835\\
1.402	-0.000414527\\
1.404	-0.000439423\\
1.406	-0.000462454\\
1.408	-0.000483556\\
1.41	-0.000502676\\
1.412	-0.000519767\\
1.414	-0.000534789\\
1.416	-0.000547712\\
1.418	-0.000558513\\
1.42	-0.000567176\\
1.422	-0.000573694\\
1.424	-0.000578068\\
1.426	-0.000580305\\
1.428	-0.000580423\\
1.43	-0.000578444\\
1.432	-0.000574399\\
1.434	-0.000568325\\
1.436	-0.000560268\\
1.438	-0.000550279\\
1.44	-0.000538415\\
1.442	-0.00052474\\
1.444	-0.000509325\\
1.446	-0.000492243\\
1.448	-0.000473575\\
1.45	-0.000453407\\
1.452	-0.000434993\\
1.454	-0.000427729\\
1.456	-0.000419181\\
1.458	-0.000409394\\
1.46	-0.000398416\\
1.462	-0.000386298\\
1.464	-0.000373093\\
1.466	-0.000358858\\
1.468	-0.000343654\\
1.47	-0.000327541\\
1.472	-0.000310583\\
1.474	-0.000292848\\
1.476	-0.000274401\\
1.478	-0.000255312\\
1.48	-0.000235652\\
1.482	-0.00021549\\
1.484	-0.0001949\\
1.486	-0.000173952\\
1.488	-0.00015272\\
1.49	-0.000131276\\
1.492	-0.000135543\\
1.494	-0.000162408\\
1.496	-0.000188392\\
1.498	-0.000213412\\
1.5	-0.000237388\\
1.502	-0.000260245\\
1.504	-0.000281915\\
1.506	-0.000302332\\
1.508	-0.000321438\\
1.51	-0.000339177\\
1.512	-0.000355503\\
1.514	-0.000370371\\
1.516	-0.000383746\\
1.518	-0.000395594\\
1.52	-0.000405892\\
1.522	-0.00041462\\
1.524	-0.000421763\\
1.526	-0.000427314\\
1.528	-0.000431271\\
1.53	-0.000433637\\
1.532	-0.000434422\\
1.534	-0.00043364\\
1.536	-0.000431313\\
1.538	-0.000427465\\
1.54	-0.000422127\\
1.542	-0.000415335\\
1.544	-0.000407131\\
1.546	-0.000397558\\
1.548	-0.000386668\\
1.55	-0.000374513\\
1.552	-0.000361151\\
1.554	-0.000346644\\
1.556	-0.000331056\\
1.558	-0.000314456\\
1.56	-0.00030547\\
1.562	-0.000301771\\
1.564	-0.000297144\\
1.566	-0.000291617\\
1.568	-0.000285219\\
1.57	-0.000277983\\
1.572	-0.000269944\\
1.574	-0.000261139\\
1.576	-0.000251606\\
1.578	-0.000241386\\
1.58	-0.000230521\\
1.582	-0.000219056\\
1.584	-0.000207036\\
1.586	-0.000194506\\
1.588	-0.000181516\\
1.59	-0.000168113\\
1.592	-0.000154345\\
1.594	-0.000140264\\
1.596	-0.000125918\\
1.598	-0.000111358\\
1.6	-0.000129821\\
1.602	-0.000148856\\
1.604	-0.000167147\\
1.606	-0.000184636\\
1.608	-0.00020127\\
1.61	-0.000217\\
1.612	-0.000231779\\
1.614	-0.000245566\\
1.616	-0.000258321\\
1.618	-0.000270011\\
1.62	-0.000280606\\
1.622	-0.00029008\\
1.624	-0.000298413\\
1.626	-0.000305587\\
1.628	-0.00031159\\
1.63	-0.000316413\\
1.632	-0.000320053\\
1.634	-0.00032251\\
1.636	-0.000323788\\
1.638	-0.000323896\\
1.64	-0.000322846\\
1.642	-0.000320657\\
1.644	-0.000317347\\
1.646	-0.000312942\\
1.648	-0.00030747\\
1.65	-0.000300962\\
1.652	-0.000293454\\
1.654	-0.000284982\\
1.656	-0.000275589\\
1.658	-0.000265318\\
1.66	-0.000254216\\
1.662	-0.000242332\\
1.664	-0.000229716\\
1.666	-0.000216423\\
1.668	-0.00021144\\
1.67	-0.00020978\\
1.672	-0.000207462\\
1.674	-0.000204501\\
1.676	-0.000200916\\
1.678	-0.000196727\\
1.68	-0.000191956\\
1.682	-0.000186627\\
1.684	-0.000180764\\
1.686	-0.000174394\\
1.688	-0.000167545\\
1.69	-0.000160245\\
1.692	-0.000152525\\
1.694	-0.000144414\\
1.696	-0.000135946\\
1.698	-0.000127153\\
1.7	-0.000118067\\
1.702	-0.000108722\\
1.704	-9.91517e-05\\
1.706	-0.000101093\\
1.708	-0.000115051\\
1.71	-0.000128441\\
1.712	-0.000141224\\
1.714	-0.00015336\\
1.716	-0.000164813\\
1.718	-0.00017555\\
1.72	-0.000185541\\
1.722	-0.000194758\\
1.724	-0.000203179\\
1.726	-0.000210781\\
1.728	-0.000217547\\
1.73	-0.000223462\\
1.732	-0.000228515\\
1.734	-0.000232697\\
1.736	-0.000236004\\
1.738	-0.000238433\\
1.74	-0.000239987\\
1.742	-0.000240669\\
1.744	-0.000240486\\
1.746	-0.00023945\\
1.748	-0.000237573\\
1.75	-0.000234872\\
1.752	-0.000231366\\
1.754	-0.000227076\\
1.756	-0.000222026\\
1.758	-0.000216244\\
1.76	-0.000209758\\
1.762	-0.000202598\\
1.764	-0.000194799\\
1.766	-0.000186395\\
1.768	-0.000177423\\
1.77	-0.000167921\\
1.772	-0.000157929\\
1.774	-0.000147489\\
1.776	-0.000144767\\
1.778	-0.000144232\\
1.78	-0.000143236\\
1.782	-0.000141788\\
1.784	-0.000139899\\
1.786	-0.000137582\\
1.788	-0.000134849\\
1.79	-0.000131716\\
1.792	-0.000128199\\
1.794	-0.000124314\\
1.796	-0.000120079\\
1.798	-0.000115512\\
1.8	-0.000110633\\
1.802	-0.000105463\\
1.804	-0.000100022\\
1.806	-9.43323e-05\\
1.808	-8.84149e-05\\
1.81	-8.22927e-05\\
1.812	-7.75131e-05\\
1.814	-8.77474e-05\\
1.816	-9.75539e-05\\
1.818	-0.000106903\\
1.82	-0.000115765\\
1.822	-0.000124116\\
1.824	-0.000131931\\
1.826	-0.000139189\\
1.828	-0.000145869\\
1.83	-0.000151956\\
1.832	-0.000157434\\
1.834	-0.00016229\\
1.836	-0.000166514\\
1.838	-0.000170099\\
1.84	-0.000173039\\
1.842	-0.000175331\\
1.844	-0.000176974\\
1.846	-0.00017797\\
1.848	-0.000178322\\
1.85	-0.000178037\\
1.852	-0.000177123\\
1.854	-0.00017559\\
1.856	-0.00017345\\
1.858	-0.000170719\\
1.86	-0.000167413\\
1.862	-0.000163549\\
1.864	-0.000159149\\
1.866	-0.000154233\\
1.868	-0.000148826\\
1.87	-0.000142952\\
1.872	-0.000136636\\
1.874	-0.000129907\\
1.876	-0.000122794\\
1.878	-0.000115325\\
1.88	-0.000107532\\
1.882	-9.94452e-05\\
1.884	-9.80938e-05\\
1.886	-9.81548e-05\\
1.888	-9.78966e-05\\
1.89	-9.73246e-05\\
1.892	-9.64451e-05\\
1.894	-9.52652e-05\\
1.896	-9.37931e-05\\
1.898	-9.20378e-05\\
1.9	-9.00089e-05\\
1.902	-8.77169e-05\\
1.904	-8.51729e-05\\
1.906	-8.23889e-05\\
1.908	-7.93773e-05\\
1.91	-7.6151e-05\\
1.912	-7.27237e-05\\
1.914	-6.91092e-05\\
1.916	-6.53221e-05\\
1.918	-6.13769e-05\\
1.92	-6.63668e-05\\
1.922	-7.35522e-05\\
1.924	-8.0395e-05\\
1.926	-8.68749e-05\\
1.928	-9.29729e-05\\
1.93	-9.86718e-05\\
1.932	-0.000103956\\
1.934	-0.000108811\\
1.936	-0.000113226\\
1.938	-0.000117188\\
1.94	-0.00012069\\
1.942	-0.000123724\\
1.944	-0.000126285\\
1.946	-0.000128369\\
1.948	-0.000129974\\
1.95	-0.0001311\\
1.952	-0.000131748\\
1.954	-0.000131922\\
1.956	-0.000131626\\
1.958	-0.000130867\\
1.96	-0.000129652\\
1.962	-0.000127992\\
1.964	-0.000125897\\
1.966	-0.00012338\\
1.968	-0.000120454\\
1.97	-0.000117134\\
1.972	-0.000113437\\
1.974	-0.00010938\\
1.976	-0.000104982\\
1.978	-0.000100261\\
1.98	-9.52393e-05\\
1.982	-8.99369e-05\\
1.984	-8.43762e-05\\
1.986	-7.85798e-05\\
1.988	-7.25709e-05\\
1.99	-6.63733e-05\\
1.992	-6.55902e-05\\
1.994	-6.59445e-05\\
1.996	-6.6081e-05\\
1.998	-6.60024e-05\\
2	-6.57121e-05\\
};
\addplot [color=mycolor10, line width=1.0pt, forget plot]
  table[row sep=crcr]{%
-0.024	0\\
-0.022	0\\
-0.02	0\\
-0.018	0\\
-0.016	0\\
-0.014	0\\
-0.012	0\\
-0.01	0\\
-0.008	0\\
-0.006	0\\
-0.004	0\\
-0.002	0\\
0	0\\
0.002	-5.08171e-07\\
0.004	-5.08654e-06\\
0.006	-0.00109264\\
0.008	-0.0016561\\
0.01	-0.00219135\\
0.012	-0.00265435\\
0.014	-0.00301899\\
0.016	-0.00327403\\
0.018	-0.00343211\\
0.02	-0.00391241\\
0.022	-0.00435342\\
0.024	-0.00550761\\
0.026	-0.00763047\\
0.028	-0.0086522\\
0.03	-0.00969627\\
0.032	-0.0107538\\
0.034	-0.011816\\
0.036	-0.012874\\
0.038	-0.0139192\\
0.04	-0.0149431\\
0.042	-0.0159373\\
0.044	-0.016894\\
0.046	-0.0178052\\
0.048	-0.0186637\\
0.05	-0.0194624\\
0.052	-0.0201949\\
0.054	-0.0208548\\
0.056	-0.0214365\\
0.058	-0.0219349\\
0.06	-0.0223453\\
0.062	-0.0226636\\
0.064	-0.0228862\\
0.066	-0.0230101\\
0.068	-0.0230329\\
0.07	-0.0229527\\
0.072	-0.0227683\\
0.074	-0.022479\\
0.076	-0.0220848\\
0.078	-0.0215862\\
0.08	-0.0209841\\
0.082	-0.0202804\\
0.084	-0.0194771\\
0.086	-0.0185769\\
0.088	-0.0175831\\
0.09	-0.0164992\\
0.092	-0.0153296\\
0.094	-0.0140787\\
0.096	-0.0127516\\
0.098	-0.0113535\\
0.1	-0.00989024\\
0.102	-0.00836772\\
0.104	-0.00701936\\
0.106	-0.00746855\\
0.108	-0.0078576\\
0.11	-0.00818612\\
0.112	-0.00845402\\
0.114	-0.00866149\\
0.116	-0.00853665\\
0.118	-0.00809331\\
0.12	-0.00762008\\
0.122	-0.00754308\\
0.124	-0.00811685\\
0.126	-0.00874853\\
0.128	-0.00934892\\
0.13	-0.0099163\\
0.132	-0.0104491\\
0.134	-0.010941\\
0.136	-0.0110706\\
0.138	-0.0111744\\
0.14	-0.0112515\\
0.142	-0.0118079\\
0.144	-0.0136214\\
0.146	-0.0153322\\
0.148	-0.016936\\
0.15	-0.0184286\\
0.152	-0.0198066\\
0.154	-0.021067\\
0.156	-0.0222072\\
0.158	-0.023225\\
0.16	-0.0241189\\
0.162	-0.0248877\\
0.164	-0.0255308\\
0.166	-0.0260479\\
0.168	-0.0264394\\
0.17	-0.0267059\\
0.172	-0.0268487\\
0.174	-0.0268692\\
0.176	-0.0267696\\
0.178	-0.0265522\\
0.18	-0.0262199\\
0.182	-0.025776\\
0.184	-0.0252239\\
0.186	-0.0245525\\
0.188	-0.0237007\\
0.19	-0.0227521\\
0.192	-0.0217118\\
0.194	-0.0205852\\
0.196	-0.0193779\\
0.198	-0.0180957\\
0.2	-0.0167446\\
0.202	-0.0153308\\
0.204	-0.0138605\\
0.206	-0.0123402\\
0.208	-0.0107764\\
0.21	-0.00923811\\
0.212	-0.00931787\\
0.214	-0.00935376\\
0.216	-0.00934653\\
0.218	-0.00929706\\
0.22	-0.00920633\\
0.222	-0.00907541\\
0.224	-0.00890548\\
0.226	-0.00869782\\
0.228	-0.00845375\\
0.23	-0.00817474\\
0.232	-0.00786228\\
0.234	-0.00751797\\
0.236	-0.00714345\\
0.238	-0.00674045\\
0.24	-0.00746801\\
0.242	-0.00893603\\
0.244	-0.0103401\\
0.246	-0.0116762\\
0.248	-0.0129408\\
0.25	-0.0141303\\
0.252	-0.0152419\\
0.254	-0.0162726\\
0.256	-0.01722\\
0.258	-0.0180822\\
0.26	-0.0188572\\
0.262	-0.0195437\\
0.264	-0.0201405\\
0.266	-0.0206469\\
0.268	-0.0210624\\
0.27	-0.0213869\\
0.272	-0.0216206\\
0.274	-0.0217641\\
0.276	-0.0218181\\
0.278	-0.0217839\\
0.28	-0.0216628\\
0.282	-0.0214565\\
0.284	-0.0211672\\
0.286	-0.0207971\\
0.288	-0.0203486\\
0.29	-0.0198247\\
0.292	-0.0192282\\
0.294	-0.0185625\\
0.296	-0.0178309\\
0.298	-0.0170371\\
0.3	-0.0161848\\
0.302	-0.0152781\\
0.304	-0.014321\\
0.306	-0.0133177\\
0.308	-0.0122725\\
0.31	-0.0111898\\
0.312	-0.0100743\\
0.314	-0.00893024\\
0.316	-0.0077624\\
0.318	-0.0076527\\
0.32	-0.00778299\\
0.322	-0.00788188\\
0.324	-0.00794955\\
0.326	-0.00798626\\
0.328	-0.00799237\\
0.33	-0.00796831\\
0.332	-0.00791461\\
0.334	-0.00783187\\
0.336	-0.00772079\\
0.338	-0.00758213\\
0.34	-0.00741674\\
0.342	-0.00722552\\
0.344	-0.00700945\\
0.346	-0.00676957\\
0.348	-0.00773824\\
0.35	-0.00869981\\
0.352	-0.00961269\\
0.354	-0.0104744\\
0.356	-0.0112825\\
0.358	-0.0120351\\
0.36	-0.0127302\\
0.362	-0.0133662\\
0.364	-0.0139417\\
0.366	-0.0144555\\
0.368	-0.0149067\\
0.37	-0.0152946\\
0.372	-0.0156187\\
0.374	-0.0158788\\
0.376	-0.0160748\\
0.378	-0.016207\\
0.38	-0.0162759\\
0.382	-0.016282\\
0.384	-0.0162263\\
0.386	-0.0161099\\
0.388	-0.0159339\\
0.39	-0.0157\\
0.392	-0.0154098\\
0.394	-0.0150651\\
0.396	-0.0146679\\
0.398	-0.0142204\\
0.4	-0.0137249\\
0.402	-0.0131839\\
0.404	-0.0126\\
0.406	-0.0119757\\
0.408	-0.0113141\\
0.41	-0.0106179\\
0.412	-0.00989024\\
0.414	-0.00913406\\
0.416	-0.00892511\\
0.418	-0.00872285\\
0.42	-0.00849918\\
0.422	-0.008255\\
0.424	-0.00799121\\
0.426	-0.00770879\\
0.428	-0.00740873\\
0.43	-0.00709207\\
0.432	-0.00675986\\
0.434	-0.0064132\\
0.436	-0.00629173\\
0.438	-0.00631566\\
0.44	-0.00631732\\
0.442	-0.00629697\\
0.444	-0.00625493\\
0.446	-0.00619162\\
0.448	-0.0061075\\
0.45	-0.00600307\\
0.452	-0.00587893\\
0.454	-0.00607334\\
0.456	-0.00673522\\
0.458	-0.00736256\\
0.46	-0.00795361\\
0.462	-0.00850679\\
0.464	-0.00902066\\
0.466	-0.00949393\\
0.468	-0.0099255\\
0.47	-0.0103144\\
0.472	-0.0106599\\
0.474	-0.0109612\\
0.476	-0.0112181\\
0.478	-0.01143\\
0.48	-0.0115971\\
0.482	-0.0117191\\
0.484	-0.0117965\\
0.486	-0.0118294\\
0.488	-0.0118183\\
0.49	-0.011764\\
0.492	-0.0116673\\
0.494	-0.0115289\\
0.496	-0.0113501\\
0.498	-0.011132\\
0.5	-0.0108759\\
0.502	-0.0105833\\
0.504	-0.0102558\\
0.506	-0.00989485\\
0.508	-0.00950236\\
0.51	-0.00908012\\
0.512	-0.00863006\\
0.514	-0.00815417\\
0.516	-0.00765451\\
0.518	-0.00713318\\
0.52	-0.00659237\\
0.522	-0.00658047\\
0.524	-0.00660079\\
0.526	-0.00660347\\
0.528	-0.00658877\\
0.53	-0.00655699\\
0.532	-0.00650848\\
0.534	-0.00644361\\
0.536	-0.0063628\\
0.538	-0.0062665\\
0.54	-0.00615519\\
0.542	-0.00602938\\
0.544	-0.00588963\\
0.546	-0.00573649\\
0.548	-0.00557057\\
0.55	-0.00539248\\
0.552	-0.00520288\\
0.554	-0.00500241\\
0.556	-0.00487129\\
0.558	-0.00481685\\
0.56	-0.00474688\\
0.562	-0.00497004\\
0.564	-0.00540286\\
0.566	-0.00580995\\
0.568	-0.00619021\\
0.57	-0.00654265\\
0.572	-0.00686642\\
0.574	-0.00716074\\
0.576	-0.00742497\\
0.578	-0.00765859\\
0.58	-0.00786119\\
0.582	-0.00803246\\
0.584	-0.00817224\\
0.586	-0.00828046\\
0.588	-0.00835716\\
0.59	-0.00840251\\
0.592	-0.00841677\\
0.594	-0.00840033\\
0.596	-0.00835367\\
0.598	-0.00827737\\
0.6	-0.0081721\\
0.602	-0.00803866\\
0.604	-0.00787789\\
0.606	-0.00769077\\
0.608	-0.0074783\\
0.61	-0.00724162\\
0.612	-0.00698189\\
0.614	-0.00670036\\
0.616	-0.00639835\\
0.618	-0.00607721\\
0.62	-0.00573837\\
0.622	-0.00538327\\
0.624	-0.00501342\\
0.626	-0.00463034\\
0.628	-0.00423559\\
0.63	-0.00383075\\
0.632	-0.00368224\\
0.634	-0.00378328\\
0.636	-0.00387301\\
0.638	-0.00395137\\
0.64	-0.00401835\\
0.642	-0.00407396\\
0.644	-0.00411824\\
0.646	-0.00415126\\
0.648	-0.0041731\\
0.65	-0.0041839\\
0.652	-0.00418381\\
0.654	-0.00417298\\
0.656	-0.00415163\\
0.658	-0.00411997\\
0.66	-0.00407825\\
0.662	-0.00402674\\
0.664	-0.00396571\\
0.666	-0.00389548\\
0.668	-0.00381637\\
0.67	-0.00387901\\
0.672	-0.00415696\\
0.674	-0.00441595\\
0.676	-0.00465533\\
0.678	-0.0048745\\
0.68	-0.00507295\\
0.682	-0.00525028\\
0.684	-0.00540613\\
0.686	-0.00554024\\
0.688	-0.00565244\\
0.69	-0.00574263\\
0.692	-0.00581078\\
0.694	-0.00585696\\
0.696	-0.00588131\\
0.698	-0.00588405\\
0.7	-0.00586546\\
0.702	-0.0058259\\
0.704	-0.00576582\\
0.706	-0.0056857\\
0.708	-0.00558612\\
0.71	-0.0054677\\
0.712	-0.00533113\\
0.714	-0.00517715\\
0.716	-0.00500656\\
0.718	-0.00482019\\
0.72	-0.00461895\\
0.722	-0.00440375\\
0.724	-0.00417557\\
0.726	-0.00393541\\
0.728	-0.0036843\\
0.73	-0.00342329\\
0.732	-0.00315346\\
0.734	-0.00287591\\
0.736	-0.00259174\\
0.738	-0.00230206\\
0.74	-0.002008\\
0.742	-0.00190777\\
0.744	-0.00200521\\
0.746	-0.0021348\\
0.748	-0.00225628\\
0.75	-0.00236935\\
0.752	-0.00247372\\
0.754	-0.00256915\\
0.756	-0.00265543\\
0.758	-0.00273236\\
0.76	-0.0027998\\
0.762	-0.00285763\\
0.764	-0.00290575\\
0.766	-0.0029441\\
0.768	-0.00297266\\
0.77	-0.00299144\\
0.772	-0.00300046\\
0.774	-0.00299981\\
0.776	-0.00298956\\
0.778	-0.00296985\\
0.78	-0.00314087\\
0.782	-0.00330053\\
0.784	-0.00344595\\
0.786	-0.00357681\\
0.788	-0.00369283\\
0.79	-0.00379379\\
0.792	-0.00387955\\
0.794	-0.00395\\
0.796	-0.0040051\\
0.798	-0.00404487\\
0.8	-0.00406936\\
0.802	-0.00407872\\
0.804	-0.0040731\\
0.806	-0.00405275\\
0.808	-0.00401792\\
0.81	-0.00396897\\
0.812	-0.00390624\\
0.814	-0.00383017\\
0.816	-0.00374121\\
0.818	-0.00363986\\
0.82	-0.00352667\\
0.822	-0.0034022\\
0.824	-0.00326705\\
0.826	-0.00312188\\
0.828	-0.00296734\\
0.83	-0.00280413\\
0.832	-0.00263295\\
0.834	-0.00245453\\
0.836	-0.00226963\\
0.838	-0.002079\\
0.84	-0.00188341\\
0.842	-0.00168364\\
0.844	-0.00148049\\
0.846	-0.00127472\\
0.848	-0.00116882\\
0.85	-0.00122866\\
0.852	-0.00134306\\
0.854	-0.00145249\\
0.856	-0.00155629\\
0.858	-0.00165411\\
0.86	-0.00174571\\
0.862	-0.00183085\\
0.864	-0.00190934\\
0.866	-0.00198099\\
0.868	-0.00204564\\
0.87	-0.00210315\\
0.872	-0.00215342\\
0.874	-0.00219636\\
0.876	-0.0022319\\
0.878	-0.00226001\\
0.88	-0.00228066\\
0.882	-0.00229388\\
0.884	-0.0022997\\
0.886	-0.00229816\\
0.888	-0.0023643\\
0.89	-0.00245833\\
0.892	-0.00254205\\
0.894	-0.00261531\\
0.896	-0.00267797\\
0.898	-0.00272996\\
0.9	-0.00277123\\
0.902	-0.00280177\\
0.904	-0.00282162\\
0.906	-0.00283086\\
0.908	-0.00282958\\
0.91	-0.00281795\\
0.912	-0.00279614\\
0.914	-0.00276437\\
0.916	-0.0027229\\
0.918	-0.002672\\
0.92	-0.00261198\\
0.922	-0.0025432\\
0.924	-0.00246602\\
0.926	-0.00238084\\
0.928	-0.00228807\\
0.93	-0.00218816\\
0.932	-0.00208157\\
0.934	-0.00196877\\
0.936	-0.00185027\\
0.938	-0.00172658\\
0.94	-0.0015982\\
0.942	-0.00146569\\
0.944	-0.00132958\\
0.946	-0.00119041\\
0.948	-0.00104873\\
0.95	-0.000905107\\
0.952	-0.000809505\\
0.954	-0.000803439\\
0.956	-0.000794929\\
0.958	-0.000784033\\
0.96	-0.000838402\\
0.962	-0.000928246\\
0.964	-0.00101451\\
0.966	-0.00109667\\
0.968	-0.0011745\\
0.97	-0.00124779\\
0.972	-0.00131637\\
0.974	-0.00138007\\
0.976	-0.00143874\\
0.978	-0.00149223\\
0.98	-0.00154045\\
0.982	-0.00158328\\
0.984	-0.00162065\\
0.986	-0.0016525\\
0.988	-0.00167879\\
0.99	-0.00169949\\
0.992	-0.0017146\\
0.994	-0.00172517\\
0.996	-0.00178445\\
0.998	-0.00183637\\
1	-0.00188084\\
1.002	-0.00191778\\
1.004	-0.00194716\\
1.006	-0.00196898\\
1.008	-0.00198325\\
1.01	-0.00199001\\
1.012	-0.00198935\\
1.014	-0.00198136\\
1.016	-0.00196617\\
1.018	-0.00194392\\
1.02	-0.0019148\\
1.022	-0.00187899\\
1.024	-0.00183671\\
1.026	-0.00178821\\
1.028	-0.00173374\\
1.03	-0.00167359\\
1.032	-0.00160803\\
1.034	-0.0015374\\
1.036	-0.001462\\
1.038	-0.00138219\\
1.04	-0.0012983\\
1.042	-0.00121071\\
1.044	-0.00111977\\
1.046	-0.00102587\\
1.048	-0.00092939\\
1.05	-0.000830718\\
1.052	-0.000730244\\
1.054	-0.000665481\\
1.056	-0.000645671\\
1.058	-0.000624724\\
1.06	-0.000602716\\
1.062	-0.000579725\\
1.064	-0.000557792\\
1.066	-0.000554263\\
1.068	-0.000549062\\
1.07	-0.000548729\\
1.072	-0.000617896\\
1.074	-0.00068446\\
1.076	-0.000748116\\
1.078	-0.000808696\\
1.08	-0.000866041\\
1.082	-0.000920004\\
1.084	-0.000970451\\
1.086	-0.00101726\\
1.088	-0.00106033\\
1.09	-0.00109955\\
1.092	-0.00113486\\
1.094	-0.00116617\\
1.096	-0.00119454\\
1.098	-0.00124164\\
1.1	-0.00128357\\
1.102	-0.00132023\\
1.104	-0.00135157\\
1.106	-0.00137753\\
1.108	-0.00139808\\
1.11	-0.00141321\\
1.112	-0.00142294\\
1.114	-0.0014273\\
1.116	-0.00142634\\
1.118	-0.00142013\\
1.12	-0.00140875\\
1.122	-0.00139232\\
1.124	-0.00137096\\
1.126	-0.00134479\\
1.128	-0.00131399\\
1.13	-0.00127873\\
1.132	-0.00123917\\
1.134	-0.00119553\\
1.136	-0.00114802\\
1.138	-0.00109685\\
1.14	-0.00104227\\
1.142	-0.000984501\\
1.144	-0.00092381\\
1.146	-0.00086045\\
1.148	-0.000794688\\
1.15	-0.000726794\\
1.152	-0.000657044\\
1.154	-0.000585715\\
1.156	-0.00051309\\
1.158	-0.000439451\\
1.16	-0.000397568\\
1.162	-0.00039836\\
1.164	-0.000398268\\
1.166	-0.000397311\\
1.168	-0.000395511\\
1.17	-0.000392893\\
1.172	-0.000389479\\
1.174	-0.000385295\\
1.176	-0.000380369\\
1.178	-0.000374726\\
1.18	-0.000368395\\
1.182	-0.000422092\\
1.184	-0.000484476\\
1.186	-0.000544333\\
1.188	-0.000601474\\
1.19	-0.000655723\\
1.192	-0.000706916\\
1.194	-0.000754901\\
1.196	-0.000799541\\
1.198	-0.000840712\\
1.2	-0.000878306\\
1.202	-0.000912227\\
1.204	-0.000942395\\
1.206	-0.000968743\\
1.208	-0.000991222\\
1.21	-0.00100979\\
1.212	-0.00102444\\
1.214	-0.00103515\\
1.216	-0.00104193\\
1.218	-0.00104481\\
1.22	-0.00104381\\
1.222	-0.001039\\
1.224	-0.00103043\\
1.226	-0.00101817\\
1.228	-0.00100233\\
1.23	-0.000982986\\
1.232	-0.000960264\\
1.234	-0.000934285\\
1.236	-0.00090518\\
1.238	-0.000873094\\
1.24	-0.000838178\\
1.242	-0.000800594\\
1.244	-0.000760511\\
1.246	-0.000718103\\
1.248	-0.000673554\\
1.25	-0.000627051\\
1.252	-0.000585258\\
1.254	-0.00054749\\
1.256	-0.000508436\\
1.258	-0.000468237\\
1.26	-0.000427035\\
1.262	-0.000384974\\
1.264	-0.0003422\\
1.266	-0.000298856\\
1.268	-0.00025509\\
1.27	-0.000211045\\
1.272	-0.00019774\\
1.274	-0.000202513\\
1.276	-0.000206549\\
1.278	-0.000209849\\
1.28	-0.000212415\\
1.282	-0.000218811\\
1.284	-0.000267977\\
1.286	-0.000315605\\
1.288	-0.000361537\\
1.29	-0.000405626\\
1.292	-0.000447733\\
1.294	-0.000487727\\
1.296	-0.000525486\\
1.298	-0.000560899\\
1.3	-0.000593863\\
1.302	-0.000624287\\
1.304	-0.00065209\\
1.306	-0.000677199\\
1.308	-0.000699555\\
1.31	-0.000719108\\
1.312	-0.000735818\\
1.314	-0.000749658\\
1.316	-0.00076061\\
1.318	-0.000768667\\
1.32	-0.000773833\\
1.322	-0.000776123\\
1.324	-0.00077556\\
1.326	-0.00077218\\
1.328	-0.000766027\\
1.33	-0.000757155\\
1.332	-0.000745627\\
1.334	-0.000731515\\
1.336	-0.000714901\\
1.338	-0.000695872\\
1.34	-0.000674526\\
1.342	-0.000650967\\
1.344	-0.000625305\\
1.346	-0.000597658\\
1.348	-0.000578922\\
1.35	-0.000562275\\
1.352	-0.000544037\\
1.354	-0.000524286\\
1.356	-0.000503105\\
1.358	-0.000480581\\
1.36	-0.000456801\\
1.362	-0.00043186\\
1.364	-0.000405851\\
1.366	-0.000378872\\
1.368	-0.000351022\\
1.37	-0.000322401\\
1.372	-0.000293111\\
1.374	-0.000263255\\
1.376	-0.000232935\\
1.378	-0.000202257\\
1.38	-0.000171321\\
1.382	-0.000140233\\
1.384	-0.000120345\\
1.386	-0.000158055\\
1.388	-0.000194772\\
1.39	-0.000230375\\
1.392	-0.000264749\\
1.394	-0.000297782\\
1.396	-0.00032937\\
1.398	-0.000359415\\
1.4	-0.000387825\\
1.402	-0.000414517\\
1.404	-0.000439414\\
1.406	-0.000462444\\
1.408	-0.000483547\\
1.41	-0.000502667\\
1.412	-0.000519758\\
1.414	-0.00053478\\
1.416	-0.000547703\\
1.418	-0.000558504\\
1.42	-0.000567167\\
1.422	-0.000573686\\
1.424	-0.00057806\\
1.426	-0.000580298\\
1.428	-0.000580415\\
1.43	-0.000578436\\
1.432	-0.000574391\\
1.434	-0.000568318\\
1.436	-0.000560261\\
1.438	-0.000550272\\
1.44	-0.000538408\\
1.442	-0.000524733\\
1.444	-0.000509318\\
1.446	-0.000492236\\
1.448	-0.000473569\\
1.45	-0.000453401\\
1.452	-0.000434987\\
1.454	-0.000427723\\
1.456	-0.000419175\\
1.458	-0.000409388\\
1.46	-0.00039841\\
1.462	-0.000386292\\
1.464	-0.000373087\\
1.466	-0.000358853\\
1.468	-0.000343648\\
1.47	-0.000327535\\
1.472	-0.000310578\\
1.474	-0.000292842\\
1.476	-0.000274396\\
1.478	-0.000255307\\
1.48	-0.000235647\\
1.482	-0.000215485\\
1.484	-0.000194895\\
1.486	-0.000173947\\
1.488	-0.000152716\\
1.49	-0.000131272\\
1.492	-0.000135538\\
1.494	-0.000162404\\
1.496	-0.000188388\\
1.498	-0.000213408\\
1.5	-0.000237384\\
1.502	-0.000260241\\
1.504	-0.000281911\\
1.506	-0.000302328\\
1.508	-0.000321434\\
1.51	-0.000339174\\
1.512	-0.000355499\\
1.514	-0.000370368\\
1.516	-0.000383742\\
1.518	-0.000395591\\
1.52	-0.000405889\\
1.522	-0.000414616\\
1.524	-0.00042176\\
1.526	-0.000427311\\
1.528	-0.000431268\\
1.53	-0.000433634\\
1.532	-0.000434419\\
1.534	-0.000433637\\
1.536	-0.00043131\\
1.538	-0.000427462\\
1.54	-0.000422124\\
1.542	-0.000415333\\
1.544	-0.000407128\\
1.546	-0.000397556\\
1.548	-0.000386665\\
1.55	-0.00037451\\
1.552	-0.000361148\\
1.554	-0.000346641\\
1.556	-0.000331053\\
1.558	-0.000314453\\
1.56	-0.000305467\\
1.562	-0.000301768\\
1.564	-0.000297141\\
1.566	-0.000291614\\
1.568	-0.000285217\\
1.57	-0.000277981\\
1.572	-0.000269942\\
1.574	-0.000261136\\
1.576	-0.000251603\\
1.578	-0.000241383\\
1.58	-0.000230519\\
1.582	-0.000219054\\
1.584	-0.000207034\\
1.586	-0.000194504\\
1.588	-0.000181514\\
1.59	-0.000168111\\
1.592	-0.000154344\\
1.594	-0.000140262\\
1.596	-0.000125916\\
1.598	-0.000111357\\
1.6	-0.000129819\\
1.602	-0.000148854\\
1.604	-0.000167145\\
1.606	-0.000184634\\
1.608	-0.000201269\\
1.61	-0.000216999\\
1.612	-0.000231778\\
1.614	-0.000245564\\
1.616	-0.000258319\\
1.618	-0.000270009\\
1.62	-0.000280604\\
1.622	-0.000290079\\
1.624	-0.000298412\\
1.626	-0.000305586\\
1.628	-0.000311589\\
1.63	-0.000316412\\
1.632	-0.000320052\\
1.634	-0.000322509\\
1.636	-0.000323787\\
1.638	-0.000323895\\
1.64	-0.000322845\\
1.642	-0.000320656\\
1.644	-0.000317346\\
1.646	-0.000312941\\
1.648	-0.000307469\\
1.65	-0.000300961\\
1.652	-0.000293453\\
1.654	-0.000284981\\
1.656	-0.000275588\\
1.658	-0.000265317\\
1.66	-0.000254215\\
1.662	-0.000242331\\
1.664	-0.000229715\\
1.666	-0.000216422\\
1.668	-0.000211439\\
1.67	-0.000209779\\
1.672	-0.000207461\\
1.674	-0.0002045\\
1.676	-0.000200915\\
1.678	-0.000196726\\
1.68	-0.000191955\\
1.682	-0.000186626\\
1.684	-0.000180763\\
1.686	-0.000174393\\
1.688	-0.000167544\\
1.69	-0.000160244\\
1.692	-0.000152524\\
1.694	-0.000144414\\
1.696	-0.000135946\\
1.698	-0.000127152\\
1.7	-0.000118066\\
1.702	-0.000108721\\
1.704	-9.9151e-05\\
1.706	-0.000101092\\
1.708	-0.00011505\\
1.71	-0.000128441\\
1.712	-0.000141224\\
1.714	-0.000153359\\
1.716	-0.000164812\\
1.718	-0.000175549\\
1.72	-0.00018554\\
1.722	-0.000194758\\
1.724	-0.000203178\\
1.726	-0.00021078\\
1.728	-0.000217546\\
1.73	-0.000223461\\
1.732	-0.000228514\\
1.734	-0.000232696\\
1.736	-0.000236003\\
1.738	-0.000238433\\
1.74	-0.000239986\\
1.742	-0.000240668\\
1.744	-0.000240486\\
1.746	-0.000239449\\
1.748	-0.000237573\\
1.75	-0.000234872\\
1.752	-0.000231365\\
1.754	-0.000227076\\
1.756	-0.000222026\\
1.758	-0.000216244\\
1.76	-0.000209757\\
1.762	-0.000202598\\
1.764	-0.000194799\\
1.766	-0.000186395\\
1.768	-0.000177422\\
1.77	-0.000167921\\
1.772	-0.000157929\\
1.774	-0.000147488\\
1.776	-0.000144766\\
1.778	-0.000144232\\
1.78	-0.000143236\\
1.782	-0.000141788\\
1.784	-0.000139899\\
1.786	-0.000137581\\
1.788	-0.000134849\\
1.79	-0.000131716\\
1.792	-0.000128199\\
1.794	-0.000124313\\
1.796	-0.000120078\\
1.798	-0.000115512\\
1.8	-0.000110633\\
1.802	-0.000105463\\
1.804	-0.000100022\\
1.806	-9.4332e-05\\
1.808	-8.84146e-05\\
1.81	-8.22924e-05\\
1.812	-7.75128e-05\\
1.814	-8.77472e-05\\
1.816	-9.75537e-05\\
1.818	-0.000106902\\
1.82	-0.000115765\\
1.822	-0.000124116\\
1.824	-0.000131931\\
1.826	-0.000139189\\
1.828	-0.000145869\\
1.83	-0.000151956\\
1.832	-0.000157434\\
1.834	-0.00016229\\
1.836	-0.000166514\\
1.838	-0.000170099\\
1.84	-0.000173039\\
1.842	-0.000175331\\
1.844	-0.000176974\\
1.846	-0.00017797\\
1.848	-0.000178322\\
1.85	-0.000178037\\
1.852	-0.000177122\\
1.854	-0.000175589\\
1.856	-0.00017345\\
1.858	-0.000170719\\
1.86	-0.000167413\\
1.862	-0.000163549\\
1.864	-0.000159149\\
1.866	-0.000154233\\
1.868	-0.000148826\\
1.87	-0.000142951\\
1.872	-0.000136636\\
1.874	-0.000129907\\
1.876	-0.000122794\\
1.878	-0.000115325\\
1.88	-0.000107532\\
1.882	-9.94451e-05\\
1.884	-9.80937e-05\\
1.886	-9.81547e-05\\
1.888	-9.78965e-05\\
1.89	-9.73245e-05\\
1.892	-9.64449e-05\\
1.894	-9.52651e-05\\
1.896	-9.3793e-05\\
1.898	-9.20377e-05\\
1.9	-9.00088e-05\\
1.902	-8.77168e-05\\
1.904	-8.51728e-05\\
1.906	-8.23888e-05\\
1.908	-7.93772e-05\\
1.91	-7.61509e-05\\
1.912	-7.27236e-05\\
1.914	-6.91091e-05\\
1.916	-6.5322e-05\\
1.918	-6.13768e-05\\
1.92	-6.63667e-05\\
1.922	-7.35521e-05\\
1.924	-8.03949e-05\\
1.926	-8.68748e-05\\
1.928	-9.29728e-05\\
1.93	-9.86717e-05\\
1.932	-0.000103956\\
1.934	-0.000108811\\
1.936	-0.000113225\\
1.938	-0.000117188\\
1.94	-0.00012069\\
1.942	-0.000123724\\
1.944	-0.000126285\\
1.946	-0.000128369\\
1.948	-0.000129974\\
1.95	-0.000131099\\
1.952	-0.000131748\\
1.954	-0.000131922\\
1.956	-0.000131626\\
1.958	-0.000130867\\
1.96	-0.000129652\\
1.962	-0.000127992\\
1.964	-0.000125897\\
1.966	-0.00012338\\
1.968	-0.000120454\\
1.97	-0.000117134\\
1.972	-0.000113437\\
1.974	-0.00010938\\
1.976	-0.000104982\\
1.978	-0.000100261\\
1.98	-9.52392e-05\\
1.982	-8.99369e-05\\
1.984	-8.43761e-05\\
1.986	-7.85797e-05\\
1.988	-7.25709e-05\\
1.99	-6.63733e-05\\
1.992	-6.55901e-05\\
1.994	-6.59444e-05\\
1.996	-6.60809e-05\\
1.998	-6.60023e-05\\
2	-6.5712e-05\\
};
\addplot [color=mycolor11, line width=1.0pt, forget plot]
  table[row sep=crcr]{%
-0.024	0\\
-0.022	0\\
-0.02	0\\
-0.018	0\\
-0.016	0\\
-0.014	0\\
-0.012	0\\
-0.01	0\\
-0.008	0\\
-0.006	0\\
-0.004	0\\
-0.002	0\\
0	0\\
0.002	-0.000172477\\
0.004	-0.000515473\\
0.006	-0.00106806\\
0.008	-0.00163349\\
0.01	-0.00218144\\
0.012	-0.00266741\\
0.014	-0.00306338\\
0.016	-0.00335549\\
0.018	-0.00355353\\
0.02	-0.00405204\\
0.022	-0.00450585\\
0.024	-0.00569698\\
0.026	-0.00752273\\
0.028	-0.00855137\\
0.03	-0.00960365\\
0.032	-0.0106705\\
0.034	-0.0117429\\
0.036	-0.0128118\\
0.038	-0.0138685\\
0.04	-0.0149043\\
0.042	-0.0159109\\
0.044	-0.01688\\
0.046	-0.0178039\\
0.048	-0.0186752\\
0.05	-0.0194867\\
0.052	-0.0202317\\
0.054	-0.0209042\\
0.056	-0.0214983\\
0.058	-0.0220088\\
0.06	-0.022431\\
0.062	-0.0227607\\
0.064	-0.0229943\\
0.066	-0.0231288\\
0.068	-0.0231617\\
0.07	-0.0230912\\
0.072	-0.0229159\\
0.074	-0.0226351\\
0.076	-0.0222488\\
0.078	-0.0217573\\
0.08	-0.0211619\\
0.082	-0.0204641\\
0.084	-0.019666\\
0.086	-0.0187703\\
0.088	-0.0177803\\
0.09	-0.0166997\\
0.092	-0.0155325\\
0.094	-0.0142834\\
0.096	-0.0129573\\
0.098	-0.0115596\\
0.1	-0.0100961\\
0.102	-0.00857263\\
0.104	-0.00722273\\
0.106	-0.00766979\\
0.108	-0.00805611\\
0.11	-0.00838134\\
0.112	-0.00864541\\
0.114	-0.00884853\\
0.116	-0.00871699\\
0.118	-0.00827106\\
0.12	-0.00779479\\
0.122	-0.00771434\\
0.124	-0.00828426\\
0.126	-0.0089117\\
0.128	-0.0095075\\
0.13	-0.01007\\
0.132	-0.0105975\\
0.134	-0.0110391\\
0.136	-0.0112338\\
0.138	-0.0113347\\
0.14	-0.0114087\\
0.142	-0.0119617\\
0.144	-0.0137715\\
0.146	-0.0154784\\
0.148	-0.0170779\\
0.15	-0.0185661\\
0.152	-0.0199395\\
0.154	-0.021195\\
0.156	-0.0223301\\
0.158	-0.0233427\\
0.16	-0.0242312\\
0.162	-0.0249945\\
0.164	-0.0256319\\
0.166	-0.0261433\\
0.168	-0.0265289\\
0.17	-0.0267895\\
0.172	-0.0269263\\
0.174	-0.0269408\\
0.176	-0.0268351\\
0.178	-0.0266117\\
0.18	-0.0262734\\
0.182	-0.0258234\\
0.184	-0.0252652\\
0.186	-0.0245596\\
0.188	-0.0237062\\
0.19	-0.022756\\
0.192	-0.0217141\\
0.194	-0.0205859\\
0.196	-0.019377\\
0.198	-0.0180932\\
0.2	-0.0167406\\
0.202	-0.0153252\\
0.204	-0.0138534\\
0.206	-0.0123317\\
0.208	-0.0107664\\
0.21	-0.00922668\\
0.212	-0.00930506\\
0.214	-0.00933959\\
0.216	-0.00933104\\
0.218	-0.00928027\\
0.22	-0.00918828\\
0.222	-0.00905614\\
0.224	-0.00888502\\
0.226	-0.00867618\\
0.228	-0.00843099\\
0.23	-0.00815088\\
0.232	-0.00783735\\
0.234	-0.007492\\
0.236	-0.00711649\\
0.238	-0.00671252\\
0.24	-0.00743914\\
0.242	-0.00890625\\
0.244	-0.0103094\\
0.246	-0.0116447\\
0.248	-0.0129085\\
0.25	-0.0140972\\
0.252	-0.015208\\
0.254	-0.016238\\
0.256	-0.0171847\\
0.258	-0.0180462\\
0.26	-0.0188206\\
0.262	-0.0195064\\
0.264	-0.0201027\\
0.266	-0.0206085\\
0.268	-0.0210235\\
0.27	-0.0213474\\
0.272	-0.0215807\\
0.274	-0.0217237\\
0.276	-0.0217772\\
0.278	-0.0217426\\
0.28	-0.0216211\\
0.282	-0.0214145\\
0.284	-0.0211248\\
0.286	-0.0207543\\
0.288	-0.0203056\\
0.29	-0.0197814\\
0.292	-0.0191846\\
0.294	-0.0185187\\
0.296	-0.0177869\\
0.298	-0.0169929\\
0.3	-0.0161405\\
0.302	-0.0152336\\
0.304	-0.0142763\\
0.306	-0.0132729\\
0.308	-0.0122276\\
0.31	-0.0111449\\
0.312	-0.0100293\\
0.314	-0.00888524\\
0.316	-0.00771738\\
0.318	-0.0076077\\
0.32	-0.00773801\\
0.322	-0.00783695\\
0.324	-0.00790468\\
0.326	-0.00794147\\
0.328	-0.00794768\\
0.33	-0.00792373\\
0.332	-0.00787015\\
0.334	-0.00778756\\
0.336	-0.00767664\\
0.338	-0.00753815\\
0.34	-0.00737295\\
0.342	-0.00718192\\
0.344	-0.00696607\\
0.346	-0.00672641\\
0.348	-0.00769532\\
0.35	-0.00865714\\
0.352	-0.00957028\\
0.354	-0.0104322\\
0.356	-0.0112407\\
0.358	-0.0119935\\
0.36	-0.0126889\\
0.362	-0.0133252\\
0.364	-0.013901\\
0.366	-0.0144152\\
0.368	-0.0148667\\
0.37	-0.0152549\\
0.372	-0.0155793\\
0.374	-0.0158398\\
0.376	-0.0160361\\
0.378	-0.0161687\\
0.38	-0.0162379\\
0.382	-0.0162444\\
0.384	-0.0161891\\
0.386	-0.016073\\
0.388	-0.0158974\\
0.39	-0.0156639\\
0.392	-0.015374\\
0.394	-0.0150297\\
0.396	-0.0146329\\
0.398	-0.0141858\\
0.4	-0.0136907\\
0.402	-0.01315\\
0.404	-0.0125664\\
0.406	-0.0119426\\
0.408	-0.0112813\\
0.41	-0.0105855\\
0.412	-0.00985819\\
0.414	-0.00910238\\
0.416	-0.0088938\\
0.418	-0.0086919\\
0.42	-0.0084686\\
0.422	-0.00822477\\
0.424	-0.00796134\\
0.426	-0.00767927\\
0.428	-0.00737957\\
0.43	-0.00706326\\
0.432	-0.00673141\\
0.434	-0.00638509\\
0.436	-0.00626396\\
0.438	-0.00628824\\
0.44	-0.00629024\\
0.442	-0.00627022\\
0.444	-0.00622852\\
0.446	-0.00616554\\
0.448	-0.00608174\\
0.45	-0.00597765\\
0.452	-0.00585383\\
0.454	-0.00604855\\
0.456	-0.00671075\\
0.458	-0.0073384\\
0.46	-0.00792976\\
0.462	-0.00848325\\
0.464	-0.00899742\\
0.466	-0.009471\\
0.468	-0.00990287\\
0.47	-0.0102921\\
0.472	-0.0106378\\
0.474	-0.0109395\\
0.476	-0.0111966\\
0.478	-0.0114089\\
0.48	-0.0115762\\
0.482	-0.0116985\\
0.484	-0.0117761\\
0.486	-0.0118093\\
0.488	-0.0117985\\
0.49	-0.0117445\\
0.492	-0.011648\\
0.494	-0.0115099\\
0.496	-0.0113314\\
0.498	-0.0111135\\
0.5	-0.0108577\\
0.502	-0.0105653\\
0.504	-0.010238\\
0.506	-0.00987733\\
0.508	-0.00948507\\
0.51	-0.00906307\\
0.512	-0.00861325\\
0.514	-0.00813759\\
0.516	-0.00763815\\
0.518	-0.00711706\\
0.52	-0.00657647\\
0.522	-0.00656479\\
0.524	-0.00658533\\
0.526	-0.00658822\\
0.528	-0.00657374\\
0.53	-0.00654217\\
0.532	-0.00649386\\
0.534	-0.0064292\\
0.536	-0.00634859\\
0.538	-0.00625248\\
0.54	-0.00614137\\
0.542	-0.00601576\\
0.544	-0.0058762\\
0.546	-0.00572325\\
0.548	-0.00555752\\
0.55	-0.00537962\\
0.552	-0.00519019\\
0.554	-0.00498991\\
0.556	-0.00485896\\
0.558	-0.00480469\\
0.56	-0.0047349\\
0.562	-0.00495823\\
0.564	-0.00539122\\
0.566	-0.00579847\\
0.568	-0.0061789\\
0.57	-0.0065315\\
0.572	-0.00685542\\
0.574	-0.0071499\\
0.576	-0.00741429\\
0.578	-0.00764806\\
0.58	-0.00785081\\
0.582	-0.00802224\\
0.584	-0.00816216\\
0.586	-0.00827052\\
0.588	-0.00834737\\
0.59	-0.00839286\\
0.592	-0.00840726\\
0.594	-0.00839096\\
0.596	-0.00834443\\
0.598	-0.00826826\\
0.6	-0.00816313\\
0.602	-0.00802982\\
0.604	-0.00786918\\
0.606	-0.00768218\\
0.608	-0.00746984\\
0.61	-0.00723328\\
0.612	-0.00697367\\
0.614	-0.00669226\\
0.616	-0.00639037\\
0.618	-0.00606935\\
0.62	-0.00573062\\
0.622	-0.00537564\\
0.624	-0.0050059\\
0.626	-0.00462293\\
0.628	-0.00422829\\
0.63	-0.00382355\\
0.632	-0.00367515\\
0.634	-0.00377629\\
0.636	-0.00386612\\
0.638	-0.00394459\\
0.64	-0.00401167\\
0.642	-0.00406738\\
0.644	-0.00411175\\
0.646	-0.00414486\\
0.648	-0.00416681\\
0.65	-0.0041777\\
0.652	-0.00417769\\
0.654	-0.00416696\\
0.656	-0.0041457\\
0.658	-0.00411413\\
0.66	-0.0040725\\
0.662	-0.00402107\\
0.664	-0.00396012\\
0.666	-0.00388998\\
0.668	-0.00381095\\
0.67	-0.00387367\\
0.672	-0.0041517\\
0.674	-0.00441077\\
0.676	-0.00465022\\
0.678	-0.00486946\\
0.68	-0.005068\\
0.682	-0.0052454\\
0.684	-0.00540132\\
0.686	-0.00553551\\
0.688	-0.00564778\\
0.69	-0.00573803\\
0.692	-0.00580625\\
0.694	-0.00585251\\
0.696	-0.00587692\\
0.698	-0.00587973\\
0.7	-0.0058612\\
0.702	-0.00582171\\
0.704	-0.00576169\\
0.706	-0.00568163\\
0.708	-0.00558211\\
0.71	-0.00546375\\
0.712	-0.00532724\\
0.714	-0.00517332\\
0.716	-0.00500279\\
0.718	-0.00481648\\
0.72	-0.00461529\\
0.722	-0.00440015\\
0.724	-0.00417203\\
0.726	-0.00393192\\
0.728	-0.00368086\\
0.73	-0.00341991\\
0.732	-0.00315013\\
0.734	-0.00287263\\
0.736	-0.00258851\\
0.738	-0.00229888\\
0.74	-0.00200487\\
0.742	-0.00190468\\
0.744	-0.00200217\\
0.746	-0.00213181\\
0.748	-0.00225334\\
0.75	-0.00236645\\
0.752	-0.00247087\\
0.754	-0.00256634\\
0.756	-0.00265266\\
0.758	-0.00272964\\
0.76	-0.00279712\\
0.762	-0.00285499\\
0.764	-0.00290315\\
0.766	-0.00294154\\
0.768	-0.00297015\\
0.77	-0.00298896\\
0.772	-0.00299803\\
0.774	-0.00299741\\
0.776	-0.0029872\\
0.778	-0.00296752\\
0.78	-0.00313858\\
0.782	-0.00329828\\
0.784	-0.00344374\\
0.786	-0.00357463\\
0.788	-0.00369068\\
0.79	-0.00379168\\
0.792	-0.00387747\\
0.794	-0.00394795\\
0.796	-0.00400309\\
0.798	-0.00404288\\
0.8	-0.00406741\\
0.802	-0.0040768\\
0.804	-0.00407121\\
0.806	-0.00405089\\
0.808	-0.00401609\\
0.81	-0.00396716\\
0.812	-0.00390447\\
0.814	-0.00382843\\
0.816	-0.00373949\\
0.818	-0.00363817\\
0.82	-0.003525\\
0.822	-0.00340056\\
0.824	-0.00326544\\
0.826	-0.0031203\\
0.828	-0.00296578\\
0.83	-0.00280259\\
0.832	-0.00263144\\
0.834	-0.00245304\\
0.836	-0.00226816\\
0.838	-0.00207756\\
0.84	-0.00188199\\
0.842	-0.00168225\\
0.844	-0.00147912\\
0.846	-0.00127337\\
0.848	-0.00116749\\
0.85	-0.00122735\\
0.852	-0.00134177\\
0.854	-0.00145123\\
0.856	-0.00155504\\
0.858	-0.00165289\\
0.86	-0.0017445\\
0.862	-0.00182967\\
0.864	-0.00190818\\
0.866	-0.00197984\\
0.868	-0.00204451\\
0.87	-0.00210204\\
0.872	-0.00215233\\
0.874	-0.00219528\\
0.876	-0.00223084\\
0.878	-0.00225897\\
0.88	-0.00227964\\
0.882	-0.00229288\\
0.884	-0.00229871\\
0.886	-0.00229718\\
0.888	-0.00236334\\
0.89	-0.00245738\\
0.892	-0.00254112\\
0.894	-0.0026144\\
0.896	-0.00267708\\
0.898	-0.00272908\\
0.9	-0.00277036\\
0.902	-0.00280091\\
0.904	-0.00282078\\
0.906	-0.00283003\\
0.908	-0.00282877\\
0.91	-0.00281715\\
0.912	-0.00279535\\
0.914	-0.0027636\\
0.916	-0.00272213\\
0.918	-0.00267125\\
0.92	-0.00261125\\
0.922	-0.00254248\\
0.924	-0.00246531\\
0.926	-0.00238014\\
0.928	-0.00228738\\
0.93	-0.00218748\\
0.932	-0.0020809\\
0.934	-0.00196811\\
0.936	-0.00184962\\
0.938	-0.00172594\\
0.94	-0.00159758\\
0.942	-0.00146507\\
0.944	-0.00132897\\
0.946	-0.00118981\\
0.948	-0.00104815\\
0.95	-0.000904529\\
0.952	-0.000808937\\
0.954	-0.00080288\\
0.956	-0.000794379\\
0.958	-0.000783492\\
0.96	-0.000837869\\
0.962	-0.000927723\\
0.964	-0.001014\\
0.966	-0.00109616\\
0.968	-0.001174\\
0.97	-0.0012473\\
0.972	-0.00131589\\
0.974	-0.0013796\\
0.976	-0.00143827\\
0.978	-0.00149177\\
0.98	-0.00153999\\
0.982	-0.00158283\\
0.984	-0.00162021\\
0.986	-0.00165207\\
0.988	-0.00167837\\
0.99	-0.00169908\\
0.992	-0.00171419\\
0.994	-0.00172477\\
0.996	-0.00178406\\
0.998	-0.00183599\\
1	-0.00188046\\
1.002	-0.0019174\\
1.004	-0.00194679\\
1.006	-0.00196862\\
1.008	-0.00198289\\
1.01	-0.00198966\\
1.012	-0.00198901\\
1.014	-0.00198102\\
1.016	-0.00196583\\
1.018	-0.00194359\\
1.02	-0.00191447\\
1.022	-0.00187867\\
1.024	-0.0018364\\
1.026	-0.0017879\\
1.028	-0.00173344\\
1.03	-0.00167329\\
1.032	-0.00160774\\
1.034	-0.00153711\\
1.036	-0.00146172\\
1.038	-0.00138191\\
1.04	-0.00129803\\
1.042	-0.00121043\\
1.044	-0.0011195\\
1.046	-0.00102561\\
1.048	-0.000929133\\
1.05	-0.000830465\\
1.052	-0.000729995\\
1.054	-0.000665236\\
1.056	-0.000645431\\
1.058	-0.000624487\\
1.06	-0.000602483\\
1.062	-0.000579496\\
1.064	-0.000557567\\
1.066	-0.000554042\\
1.068	-0.000548844\\
1.07	-0.000548515\\
1.072	-0.000617685\\
1.074	-0.000684253\\
1.076	-0.000747913\\
1.078	-0.000808496\\
1.08	-0.000865844\\
1.082	-0.00091981\\
1.084	-0.000970261\\
1.086	-0.00101707\\
1.088	-0.00106014\\
1.09	-0.00109937\\
1.092	-0.00113468\\
1.094	-0.00116599\\
1.096	-0.00119436\\
1.098	-0.00124147\\
1.1	-0.0012834\\
1.102	-0.00132007\\
1.104	-0.00135141\\
1.106	-0.00137737\\
1.108	-0.00139792\\
1.11	-0.00141306\\
1.112	-0.00142279\\
1.114	-0.00142715\\
1.116	-0.0014262\\
1.118	-0.00141999\\
1.12	-0.00140861\\
1.122	-0.00139218\\
1.124	-0.00137082\\
1.126	-0.00134466\\
1.128	-0.00131386\\
1.13	-0.0012786\\
1.132	-0.00123905\\
1.134	-0.00119541\\
1.136	-0.0011479\\
1.138	-0.00109673\\
1.14	-0.00104215\\
1.142	-0.000984384\\
1.144	-0.000923695\\
1.146	-0.000860337\\
1.148	-0.000794577\\
1.15	-0.000726685\\
1.152	-0.000656936\\
1.154	-0.00058561\\
1.156	-0.000512986\\
1.158	-0.000439349\\
1.16	-0.000397468\\
1.162	-0.000398262\\
1.164	-0.000398171\\
1.166	-0.000397215\\
1.168	-0.000395418\\
1.17	-0.0003928\\
1.172	-0.000389388\\
1.174	-0.000385206\\
1.176	-0.000380281\\
1.178	-0.00037464\\
1.18	-0.00036831\\
1.182	-0.000422009\\
1.184	-0.000484394\\
1.186	-0.000544252\\
1.188	-0.000601395\\
1.19	-0.000655645\\
1.192	-0.000706839\\
1.194	-0.000754825\\
1.196	-0.000799467\\
1.198	-0.00084064\\
1.2	-0.000878235\\
1.202	-0.000912157\\
1.204	-0.000942326\\
1.206	-0.000968675\\
1.208	-0.000991155\\
1.21	-0.00100973\\
1.212	-0.00102437\\
1.214	-0.00103508\\
1.216	-0.00104187\\
1.218	-0.00104474\\
1.22	-0.00104375\\
1.222	-0.00103894\\
1.224	-0.00103037\\
1.226	-0.00101812\\
1.228	-0.00100227\\
1.23	-0.00098293\\
1.232	-0.00096021\\
1.234	-0.000934231\\
1.236	-0.000905127\\
1.238	-0.000873042\\
1.24	-0.000838127\\
1.242	-0.000800544\\
1.244	-0.000760462\\
1.246	-0.000718055\\
1.248	-0.000673506\\
1.25	-0.000627004\\
1.252	-0.000585212\\
1.254	-0.000547445\\
1.256	-0.000508392\\
1.258	-0.000468194\\
1.26	-0.000426992\\
1.262	-0.000384932\\
1.264	-0.000342158\\
1.266	-0.000298815\\
1.268	-0.000255049\\
1.27	-0.000211006\\
1.272	-0.000197701\\
1.274	-0.000202475\\
1.276	-0.000206512\\
1.278	-0.000209812\\
1.28	-0.000212379\\
1.282	-0.000218775\\
1.284	-0.000267942\\
1.286	-0.000315571\\
1.288	-0.000361503\\
1.29	-0.000405593\\
1.292	-0.0004477\\
1.294	-0.000487695\\
1.296	-0.000525454\\
1.298	-0.000560867\\
1.3	-0.000593832\\
1.302	-0.000624257\\
1.304	-0.00065206\\
1.306	-0.00067717\\
1.308	-0.000699527\\
1.31	-0.00071908\\
1.312	-0.000735791\\
1.314	-0.000749631\\
1.316	-0.000760583\\
1.318	-0.000768641\\
1.32	-0.000773808\\
1.322	-0.000776098\\
1.324	-0.000775536\\
1.326	-0.000772156\\
1.328	-0.000766003\\
1.33	-0.000757131\\
1.332	-0.000745604\\
1.334	-0.000731492\\
1.336	-0.000714878\\
1.338	-0.00069585\\
1.34	-0.000674505\\
1.342	-0.000650946\\
1.344	-0.000625284\\
1.346	-0.000597637\\
1.348	-0.000578902\\
1.35	-0.000562255\\
1.352	-0.000544017\\
1.354	-0.000524267\\
1.356	-0.000503086\\
1.358	-0.000480562\\
1.36	-0.000456783\\
1.362	-0.000431842\\
1.364	-0.000405834\\
1.366	-0.000378855\\
1.368	-0.000351005\\
1.37	-0.000322384\\
1.372	-0.000293094\\
1.374	-0.000263238\\
1.376	-0.000232919\\
1.378	-0.000202241\\
1.38	-0.000171306\\
1.382	-0.000140218\\
1.384	-0.000120331\\
1.386	-0.00015804\\
1.388	-0.000194757\\
1.39	-0.000230361\\
1.392	-0.000264735\\
1.394	-0.000297768\\
1.396	-0.000329356\\
1.398	-0.000359401\\
1.4	-0.000387812\\
1.402	-0.000414505\\
1.404	-0.000439401\\
1.406	-0.000462432\\
1.408	-0.000483535\\
1.41	-0.000502655\\
1.412	-0.000519746\\
1.414	-0.000534769\\
1.416	-0.000547692\\
1.418	-0.000558493\\
1.42	-0.000567156\\
1.422	-0.000573675\\
1.424	-0.000578049\\
1.426	-0.000580287\\
1.428	-0.000580405\\
1.43	-0.000578426\\
1.432	-0.000574382\\
1.434	-0.000568308\\
1.436	-0.000560252\\
1.438	-0.000550263\\
1.44	-0.000538399\\
1.442	-0.000524724\\
1.444	-0.000509309\\
1.446	-0.000492227\\
1.448	-0.00047356\\
1.45	-0.000453392\\
1.452	-0.000434978\\
1.454	-0.000427715\\
1.456	-0.000419167\\
1.458	-0.000409381\\
1.46	-0.000398403\\
1.462	-0.000386284\\
1.464	-0.000373079\\
1.466	-0.000358845\\
1.468	-0.000343641\\
1.47	-0.000327528\\
1.472	-0.000310571\\
1.474	-0.000292836\\
1.476	-0.000274389\\
1.478	-0.000255301\\
1.48	-0.00023564\\
1.482	-0.000215479\\
1.484	-0.000194889\\
1.486	-0.000173941\\
1.488	-0.00015271\\
1.49	-0.000131266\\
1.492	-0.000135532\\
1.494	-0.000162398\\
1.496	-0.000188382\\
1.498	-0.000213402\\
1.5	-0.000237378\\
1.502	-0.000260236\\
1.504	-0.000281906\\
1.506	-0.000302323\\
1.508	-0.000321429\\
1.51	-0.000339169\\
1.512	-0.000355494\\
1.514	-0.000370363\\
1.516	-0.000383737\\
1.518	-0.000395586\\
1.52	-0.000405884\\
1.522	-0.000414612\\
1.524	-0.000421755\\
1.526	-0.000427306\\
1.528	-0.000431263\\
1.53	-0.00043363\\
1.532	-0.000434415\\
1.534	-0.000433633\\
1.536	-0.000431306\\
1.538	-0.000427458\\
1.54	-0.00042212\\
1.542	-0.000415329\\
1.544	-0.000407124\\
1.546	-0.000397552\\
1.548	-0.000386661\\
1.55	-0.000374506\\
1.552	-0.000361145\\
1.554	-0.000346638\\
1.556	-0.00033105\\
1.558	-0.00031445\\
1.56	-0.000305464\\
1.562	-0.000301765\\
1.564	-0.000297138\\
1.566	-0.000291611\\
1.568	-0.000285214\\
1.57	-0.000277978\\
1.572	-0.000269939\\
1.574	-0.000261134\\
1.576	-0.000251601\\
1.578	-0.000241381\\
1.58	-0.000230516\\
1.582	-0.000219051\\
1.584	-0.000207031\\
1.586	-0.000194502\\
1.588	-0.000181512\\
1.59	-0.000168108\\
1.592	-0.000154341\\
1.594	-0.00014026\\
1.596	-0.000125914\\
1.598	-0.000111354\\
1.6	-0.000129817\\
1.602	-0.000148852\\
1.604	-0.000167143\\
1.606	-0.000184632\\
1.608	-0.000201267\\
1.61	-0.000216997\\
1.612	-0.000231776\\
1.614	-0.000245562\\
1.616	-0.000258317\\
1.618	-0.000270007\\
1.62	-0.000280602\\
1.622	-0.000290077\\
1.624	-0.00029841\\
1.626	-0.000305584\\
1.628	-0.000311587\\
1.63	-0.00031641\\
1.632	-0.00032005\\
1.634	-0.000322507\\
1.636	-0.000323785\\
1.638	-0.000323893\\
1.64	-0.000322844\\
1.642	-0.000320654\\
1.644	-0.000317345\\
1.646	-0.00031294\\
1.648	-0.000307468\\
1.65	-0.00030096\\
1.652	-0.000293451\\
1.654	-0.00028498\\
1.656	-0.000275587\\
1.658	-0.000265316\\
1.66	-0.000254214\\
1.662	-0.000242329\\
1.664	-0.000229714\\
1.666	-0.000216421\\
1.668	-0.000211438\\
1.67	-0.000209778\\
1.672	-0.00020746\\
1.674	-0.000204499\\
1.676	-0.000200914\\
1.678	-0.000196725\\
1.68	-0.000191954\\
1.682	-0.000186625\\
1.684	-0.000180762\\
1.686	-0.000174392\\
1.688	-0.000167543\\
1.69	-0.000160243\\
1.692	-0.000152523\\
1.694	-0.000144413\\
1.696	-0.000135945\\
1.698	-0.000127151\\
1.7	-0.000118065\\
1.702	-0.00010872\\
1.704	-9.91501e-05\\
1.706	-0.000101091\\
1.708	-0.000115049\\
1.71	-0.00012844\\
1.712	-0.000141223\\
1.714	-0.000153358\\
1.716	-0.000164811\\
1.718	-0.000175548\\
1.72	-0.000185539\\
1.722	-0.000194757\\
1.724	-0.000203177\\
1.726	-0.000210779\\
1.728	-0.000217545\\
1.73	-0.00022346\\
1.732	-0.000228513\\
1.734	-0.000232696\\
1.736	-0.000236003\\
1.738	-0.000238432\\
1.74	-0.000239986\\
1.742	-0.000240668\\
1.744	-0.000240485\\
1.746	-0.000239449\\
1.748	-0.000237572\\
1.75	-0.000234871\\
1.752	-0.000231365\\
1.754	-0.000227075\\
1.756	-0.000222025\\
1.758	-0.000216243\\
1.76	-0.000209757\\
1.762	-0.000202597\\
1.764	-0.000194798\\
1.766	-0.000186394\\
1.768	-0.000177422\\
1.77	-0.00016792\\
1.772	-0.000157928\\
1.774	-0.000147488\\
1.776	-0.000144766\\
1.778	-0.000144231\\
1.78	-0.000143235\\
1.782	-0.000141787\\
1.784	-0.000139898\\
1.786	-0.000137581\\
1.788	-0.000134849\\
1.79	-0.000131716\\
1.792	-0.000128198\\
1.794	-0.000124313\\
1.796	-0.000120078\\
1.798	-0.000115511\\
1.8	-0.000110633\\
1.802	-0.000105463\\
1.804	-0.000100022\\
1.806	-9.43316e-05\\
1.808	-8.84142e-05\\
1.81	-8.2292e-05\\
1.812	-7.75125e-05\\
1.814	-8.77468e-05\\
1.816	-9.75533e-05\\
1.818	-0.000106902\\
1.82	-0.000115765\\
1.822	-0.000124115\\
1.824	-0.000131931\\
1.826	-0.000139188\\
1.828	-0.000145869\\
1.83	-0.000151956\\
1.832	-0.000157433\\
1.834	-0.000162289\\
1.836	-0.000166514\\
1.838	-0.000170099\\
1.84	-0.000173039\\
1.842	-0.000175331\\
1.844	-0.000176974\\
1.846	-0.00017797\\
1.848	-0.000178322\\
1.85	-0.000178037\\
1.852	-0.000177122\\
1.854	-0.000175589\\
1.856	-0.00017345\\
1.858	-0.000170719\\
1.86	-0.000167412\\
1.862	-0.000163549\\
1.864	-0.000159149\\
1.866	-0.000154233\\
1.868	-0.000148826\\
1.87	-0.000142951\\
1.872	-0.000136636\\
1.874	-0.000129907\\
1.876	-0.000122794\\
1.878	-0.000115325\\
1.88	-0.000107531\\
1.882	-9.94449e-05\\
1.884	-9.80935e-05\\
1.886	-9.81545e-05\\
1.888	-9.78963e-05\\
1.89	-9.73243e-05\\
1.892	-9.64448e-05\\
1.894	-9.52649e-05\\
1.896	-9.37929e-05\\
1.898	-9.20375e-05\\
1.9	-9.00086e-05\\
1.902	-8.77166e-05\\
1.904	-8.51727e-05\\
1.906	-8.23887e-05\\
1.908	-7.9377e-05\\
1.91	-7.61508e-05\\
1.912	-7.27234e-05\\
1.914	-6.9109e-05\\
1.916	-6.53218e-05\\
1.918	-6.13767e-05\\
1.92	-6.63665e-05\\
1.922	-7.3552e-05\\
1.924	-8.03948e-05\\
1.926	-8.68747e-05\\
1.928	-9.29727e-05\\
1.93	-9.86716e-05\\
1.932	-0.000103956\\
1.934	-0.000108811\\
1.936	-0.000113225\\
1.938	-0.000117188\\
1.94	-0.00012069\\
1.942	-0.000123724\\
1.944	-0.000126285\\
1.946	-0.000128368\\
1.948	-0.000129973\\
1.95	-0.000131099\\
1.952	-0.000131748\\
1.954	-0.000131922\\
1.956	-0.000131626\\
1.958	-0.000130867\\
1.96	-0.000129652\\
1.962	-0.000127992\\
1.964	-0.000125897\\
1.966	-0.00012338\\
1.968	-0.000120454\\
1.97	-0.000117134\\
1.972	-0.000113437\\
1.974	-0.00010938\\
1.976	-0.000104982\\
1.978	-0.000100261\\
1.98	-9.52391e-05\\
1.982	-8.99368e-05\\
1.984	-8.43761e-05\\
1.986	-7.85797e-05\\
1.988	-7.25708e-05\\
1.99	-6.63732e-05\\
1.992	-6.559e-05\\
1.994	-6.59443e-05\\
1.996	-6.60808e-05\\
1.998	-6.60022e-05\\
2	-6.57119e-05\\
};
\end{axis}

\begin{axis}[%
width={\figscale*1.528in},
height={\figscale*0.46in},
at={(\figscale*4.752in,\figscale*0.56in)},
scale only axis,
separate axis lines,
every outer x axis line/.append style={black},
every x tick label/.append style={font=\figticfont},
every x tick/.append style={black},
xmin=-0.025,
xmax=2,
xlabel={Time since step in s},
every outer y axis line/.append style={black},
every y tick label/.append style={font=\figticfont},
every y tick/.append style={black},
ymin=-0.0359538,
ymax=0.00159155,
ytick={-0.03,0},
axis background/.style={fill=white},
xmajorgrids,
ymajorgrids,
grid style={white!55!black, opacity=0.3},
ylabel style={rotate=-90, at={(axis description cs:-0.145,0.5)}, anchor=east}
]
\addplot [color=mycolor1, line width=1.4pt, forget plot]
  table[row sep=crcr]{%
-0.024	0\\
-0.022	0\\
-0.02	0\\
-0.018	0\\
-0.016	0\\
-0.014	0\\
-0.012	0\\
-0.01	0\\
-0.008	0\\
-0.006	0\\
-0.004	0\\
-0.002	0\\
0	0\\
0.002	0\\
0.004	0\\
0.006	0\\
0.008	0\\
0.01	0\\
0.012	0\\
0.014	0\\
0.016	0\\
0.018	0\\
0.02	0\\
0.022	0\\
0.024	0\\
0.026	0\\
0.028	0\\
0.03	0\\
0.032	0\\
0.034	0\\
0.036	0\\
0.038	0\\
0.04	0\\
0.042	0\\
0.044	0\\
0.046	0\\
0.048	0\\
0.05	0\\
0.052	0\\
0.054	0\\
0.056	0\\
0.058	0\\
0.06	0\\
0.062	0\\
0.064	0\\
0.066	0\\
0.068	0\\
0.07	0\\
0.072	0\\
0.074	0\\
0.076	0\\
0.078	0\\
0.08	0\\
0.082	0\\
0.084	0\\
0.086	0\\
0.088	0\\
0.09	0\\
0.092	0\\
0.094	0\\
0.096	0\\
0.098	0\\
0.1	0\\
0.102	0\\
0.104	0\\
0.106	0\\
0.108	0\\
0.11	0\\
0.112	0\\
0.114	0\\
0.116	0\\
0.118	0\\
0.12	0\\
0.122	0\\
0.124	0\\
0.126	0\\
0.128	0\\
0.13	0\\
0.132	0\\
0.134	0\\
0.136	0\\
0.138	0\\
0.14	0\\
0.142	0\\
0.144	0\\
0.146	0\\
0.148	0\\
0.15	0\\
0.152	0\\
0.154	0\\
0.156	0\\
0.158	0\\
0.16	0\\
0.162	0\\
0.164	0\\
0.166	0\\
0.168	0\\
0.17	0\\
0.172	0\\
0.174	0\\
0.176	0\\
0.178	0\\
0.18	0\\
0.182	0\\
0.184	0\\
0.186	0\\
0.188	0\\
0.19	0\\
0.192	0\\
0.194	0\\
0.196	0\\
0.198	0\\
0.2	0\\
0.202	0\\
0.204	0\\
0.206	0\\
0.208	0\\
0.21	0\\
0.212	0\\
0.214	0\\
0.216	0\\
0.218	0\\
0.22	0\\
0.222	0\\
0.224	0\\
0.226	0\\
0.228	0\\
0.23	0\\
0.232	0\\
0.234	0\\
0.236	0\\
0.238	0\\
0.24	0\\
0.242	0\\
0.244	0\\
0.246	0\\
0.248	0\\
0.25	0\\
0.252	0\\
0.254	0\\
0.256	0\\
0.258	0\\
0.26	0\\
0.262	0\\
0.264	0\\
0.266	0\\
0.268	0\\
0.27	0\\
0.272	0\\
0.274	0\\
0.276	0\\
0.278	0\\
0.28	0\\
0.282	0\\
0.284	0\\
0.286	0\\
0.288	0\\
0.29	0\\
0.292	0\\
0.294	0\\
0.296	0\\
0.298	0\\
0.3	0\\
0.302	0\\
0.304	0\\
0.306	0\\
0.308	0\\
0.31	0\\
0.312	0\\
0.314	0\\
0.316	0\\
0.318	0\\
0.32	0\\
0.322	0\\
0.324	0\\
0.326	0\\
0.328	0\\
0.33	0\\
0.332	0\\
0.334	0\\
0.336	0\\
0.338	0\\
0.34	0\\
0.342	0\\
0.344	0\\
0.346	0\\
0.348	0\\
0.35	0\\
0.352	0\\
0.354	0\\
0.356	0\\
0.358	0\\
0.36	0\\
0.362	0\\
0.364	0\\
0.366	0\\
0.368	0\\
0.37	0\\
0.372	0\\
0.374	0\\
0.376	0\\
0.378	0\\
0.38	0\\
0.382	0\\
0.384	0\\
0.386	0\\
0.388	0\\
0.39	0\\
0.392	0\\
0.394	0\\
0.396	0\\
0.398	0\\
0.4	0\\
0.402	0\\
0.404	0\\
0.406	0\\
0.408	0\\
0.41	0\\
0.412	0\\
0.414	0\\
0.416	0\\
0.418	0\\
0.42	0\\
0.422	0\\
0.424	0\\
0.426	0\\
0.428	0\\
0.43	0\\
0.432	0\\
0.434	0\\
0.436	0\\
0.438	0\\
0.44	0\\
0.442	0\\
0.444	0\\
0.446	0\\
0.448	0\\
0.45	0\\
0.452	0\\
0.454	0\\
0.456	0\\
0.458	0\\
0.46	0\\
0.462	0\\
0.464	0\\
0.466	0\\
0.468	0\\
0.47	0\\
0.472	0\\
0.474	0\\
0.476	0\\
0.478	0\\
0.48	0\\
0.482	0\\
0.484	0\\
0.486	0\\
0.488	0\\
0.49	0\\
0.492	0\\
0.494	0\\
0.496	0\\
0.498	0\\
0.5	0\\
0.502	0\\
0.504	0\\
0.506	0\\
0.508	0\\
0.51	0\\
0.512	0\\
0.514	0\\
0.516	0\\
0.518	0\\
0.52	0\\
0.522	0\\
0.524	0\\
0.526	0\\
0.528	0\\
0.53	0\\
0.532	0\\
0.534	0\\
0.536	0\\
0.538	0\\
0.54	0\\
0.542	0\\
0.544	0\\
0.546	0\\
0.548	0\\
0.55	0\\
0.552	0\\
0.554	0\\
0.556	0\\
0.558	0\\
0.56	0\\
0.562	0\\
0.564	0\\
0.566	0\\
0.568	0\\
0.57	0\\
0.572	0\\
0.574	0\\
0.576	0\\
0.578	0\\
0.58	0\\
0.582	0\\
0.584	0\\
0.586	0\\
0.588	0\\
0.59	0\\
0.592	0\\
0.594	0\\
0.596	0\\
0.598	0\\
0.6	0\\
0.602	0\\
0.604	0\\
0.606	0\\
0.608	0\\
0.61	0\\
0.612	0\\
0.614	0\\
0.616	0\\
0.618	0\\
0.62	0\\
0.622	0\\
0.624	0\\
0.626	0\\
0.628	0\\
0.63	0\\
0.632	0\\
0.634	0\\
0.636	0\\
0.638	0\\
0.64	0\\
0.642	0\\
0.644	0\\
0.646	0\\
0.648	0\\
0.65	0\\
0.652	0\\
0.654	0\\
0.656	0\\
0.658	0\\
0.66	0\\
0.662	0\\
0.664	0\\
0.666	0\\
0.668	0\\
0.67	0\\
0.672	0\\
0.674	0\\
0.676	0\\
0.678	0\\
0.68	0\\
0.682	0\\
0.684	0\\
0.686	0\\
0.688	0\\
0.69	0\\
0.692	0\\
0.694	0\\
0.696	0\\
0.698	0\\
0.7	0\\
0.702	0\\
0.704	0\\
0.706	0\\
0.708	0\\
0.71	0\\
0.712	0\\
0.714	0\\
0.716	0\\
0.718	0\\
0.72	0\\
0.722	0\\
0.724	0\\
0.726	0\\
0.728	0\\
0.73	0\\
0.732	0\\
0.734	0\\
0.736	0\\
0.738	0\\
0.74	0\\
0.742	0\\
0.744	0\\
0.746	0\\
0.748	0\\
0.75	0\\
0.752	0\\
0.754	0\\
0.756	0\\
0.758	0\\
0.76	0\\
0.762	0\\
0.764	0\\
0.766	0\\
0.768	0\\
0.77	0\\
0.772	0\\
0.774	0\\
0.776	0\\
0.778	0\\
0.78	0\\
0.782	0\\
0.784	0\\
0.786	0\\
0.788	0\\
0.79	0\\
0.792	0\\
0.794	0\\
0.796	0\\
0.798	0\\
0.8	0\\
0.802	0\\
0.804	0\\
0.806	0\\
0.808	0\\
0.81	0\\
0.812	0\\
0.814	0\\
0.816	0\\
0.818	0\\
0.82	0\\
0.822	0\\
0.824	0\\
0.826	0\\
0.828	0\\
0.83	0\\
0.832	0\\
0.834	0\\
0.836	0\\
0.838	0\\
0.84	0\\
0.842	0\\
0.844	0\\
0.846	0\\
0.848	0\\
0.85	0\\
0.852	0\\
0.854	0\\
0.856	0\\
0.858	0\\
0.86	0\\
0.862	0\\
0.864	0\\
0.866	0\\
0.868	0\\
0.87	0\\
0.872	0\\
0.874	0\\
0.876	0\\
0.878	0\\
0.88	0\\
0.882	0\\
0.884	0\\
0.886	0\\
0.888	0\\
0.89	0\\
0.892	0\\
0.894	0\\
0.896	0\\
0.898	0\\
0.9	0\\
0.902	0\\
0.904	0\\
0.906	0\\
0.908	0\\
0.91	0\\
0.912	0\\
0.914	0\\
0.916	0\\
0.918	0\\
0.92	0\\
0.922	0\\
0.924	0\\
0.926	0\\
0.928	0\\
0.93	0\\
0.932	0\\
0.934	0\\
0.936	0\\
0.938	0\\
0.94	0\\
0.942	0\\
0.944	0\\
0.946	0\\
0.948	0\\
0.95	0\\
0.952	0\\
0.954	0\\
0.956	0\\
0.958	0\\
0.96	0\\
0.962	0\\
0.964	0\\
0.966	0\\
0.968	0\\
0.97	0\\
0.972	0\\
0.974	0\\
0.976	0\\
0.978	0\\
0.98	0\\
0.982	0\\
0.984	0\\
0.986	0\\
0.988	0\\
0.99	0\\
0.992	0\\
0.994	0\\
0.996	0\\
0.998	0\\
1	0\\
1.002	0\\
1.004	0\\
1.006	0\\
1.008	0\\
1.01	0\\
1.012	0\\
1.014	0\\
1.016	0\\
1.018	0\\
1.02	0\\
1.022	0\\
1.024	0\\
1.026	0\\
1.028	0\\
1.03	0\\
1.032	0\\
1.034	0\\
1.036	0\\
1.038	0\\
1.04	0\\
1.042	0\\
1.044	0\\
1.046	0\\
1.048	0\\
1.05	0\\
1.052	0\\
1.054	0\\
1.056	0\\
1.058	0\\
1.06	0\\
1.062	0\\
1.064	0\\
1.066	0\\
1.068	0\\
1.07	0\\
1.072	0\\
1.074	0\\
1.076	0\\
1.078	0\\
1.08	0\\
1.082	0\\
1.084	0\\
1.086	0\\
1.088	0\\
1.09	0\\
1.092	0\\
1.094	0\\
1.096	0\\
1.098	0\\
1.1	0\\
1.102	0\\
1.104	0\\
1.106	0\\
1.108	0\\
1.11	0\\
1.112	0\\
1.114	0\\
1.116	0\\
1.118	0\\
1.12	0\\
1.122	0\\
1.124	0\\
1.126	0\\
1.128	0\\
1.13	0\\
1.132	0\\
1.134	0\\
1.136	0\\
1.138	0\\
1.14	0\\
1.142	0\\
1.144	0\\
1.146	0\\
1.148	0\\
1.15	0\\
1.152	0\\
1.154	0\\
1.156	0\\
1.158	0\\
1.16	0\\
1.162	0\\
1.164	0\\
1.166	0\\
1.168	0\\
1.17	0\\
1.172	0\\
1.174	0\\
1.176	0\\
1.178	0\\
1.18	0\\
1.182	0\\
1.184	0\\
1.186	0\\
1.188	0\\
1.19	0\\
1.192	0\\
1.194	0\\
1.196	0\\
1.198	0\\
1.2	0\\
1.202	0\\
1.204	0\\
1.206	0\\
1.208	0\\
1.21	0\\
1.212	0\\
1.214	0\\
1.216	0\\
1.218	0\\
1.22	0\\
1.222	0\\
1.224	0\\
1.226	0\\
1.228	0\\
1.23	0\\
1.232	0\\
1.234	0\\
1.236	0\\
1.238	0\\
1.24	0\\
1.242	0\\
1.244	0\\
1.246	0\\
1.248	0\\
1.25	0\\
1.252	0\\
1.254	0\\
1.256	0\\
1.258	0\\
1.26	0\\
1.262	0\\
1.264	0\\
1.266	0\\
1.268	0\\
1.27	0\\
1.272	0\\
1.274	0\\
1.276	0\\
1.278	0\\
1.28	0\\
1.282	0\\
1.284	0\\
1.286	0\\
1.288	0\\
1.29	0\\
1.292	0\\
1.294	0\\
1.296	0\\
1.298	0\\
1.3	0\\
1.302	0\\
1.304	0\\
1.306	0\\
1.308	0\\
1.31	0\\
1.312	0\\
1.314	0\\
1.316	0\\
1.318	0\\
1.32	0\\
1.322	0\\
1.324	0\\
1.326	0\\
1.328	0\\
1.33	0\\
1.332	0\\
1.334	0\\
1.336	0\\
1.338	0\\
1.34	0\\
1.342	0\\
1.344	0\\
1.346	0\\
1.348	0\\
1.35	0\\
1.352	0\\
1.354	0\\
1.356	0\\
1.358	0\\
1.36	0\\
1.362	0\\
1.364	0\\
1.366	0\\
1.368	0\\
1.37	0\\
1.372	0\\
1.374	0\\
1.376	0\\
1.378	0\\
1.38	0\\
1.382	0\\
1.384	0\\
1.386	0\\
1.388	0\\
1.39	0\\
1.392	0\\
1.394	0\\
1.396	0\\
1.398	0\\
1.4	0\\
1.402	0\\
1.404	0\\
1.406	0\\
1.408	0\\
1.41	0\\
1.412	0\\
1.414	0\\
1.416	0\\
1.418	0\\
1.42	0\\
1.422	0\\
1.424	0\\
1.426	0\\
1.428	0\\
1.43	0\\
1.432	0\\
1.434	0\\
1.436	0\\
1.438	0\\
1.44	0\\
1.442	0\\
1.444	0\\
1.446	0\\
1.448	0\\
1.45	0\\
1.452	0\\
1.454	0\\
1.456	0\\
1.458	0\\
1.46	0\\
1.462	0\\
1.464	0\\
1.466	0\\
1.468	0\\
1.47	0\\
1.472	0\\
1.474	0\\
1.476	0\\
1.478	0\\
1.48	0\\
1.482	0\\
1.484	0\\
1.486	0\\
1.488	0\\
1.49	0\\
1.492	0\\
1.494	0\\
1.496	0\\
1.498	0\\
1.5	0\\
1.502	0\\
1.504	0\\
1.506	0\\
1.508	0\\
1.51	0\\
1.512	0\\
1.514	0\\
1.516	0\\
1.518	0\\
1.52	0\\
1.522	0\\
1.524	0\\
1.526	0\\
1.528	0\\
1.53	0\\
1.532	0\\
1.534	0\\
1.536	0\\
1.538	0\\
1.54	0\\
1.542	0\\
1.544	0\\
1.546	0\\
1.548	0\\
1.55	0\\
1.552	0\\
1.554	0\\
1.556	0\\
1.558	0\\
1.56	0\\
1.562	0\\
1.564	0\\
1.566	0\\
1.568	0\\
1.57	0\\
1.572	0\\
1.574	0\\
1.576	0\\
1.578	0\\
1.58	0\\
1.582	0\\
1.584	0\\
1.586	0\\
1.588	0\\
1.59	0\\
1.592	0\\
1.594	0\\
1.596	0\\
1.598	0\\
1.6	0\\
1.602	0\\
1.604	0\\
1.606	0\\
1.608	0\\
1.61	0\\
1.612	0\\
1.614	0\\
1.616	0\\
1.618	0\\
1.62	0\\
1.622	0\\
1.624	0\\
1.626	0\\
1.628	0\\
1.63	0\\
1.632	0\\
1.634	0\\
1.636	0\\
1.638	0\\
1.64	0\\
1.642	0\\
1.644	0\\
1.646	0\\
1.648	0\\
1.65	0\\
1.652	0\\
1.654	0\\
1.656	0\\
1.658	0\\
1.66	0\\
1.662	0\\
1.664	0\\
1.666	0\\
1.668	0\\
1.67	0\\
1.672	0\\
1.674	0\\
1.676	0\\
1.678	0\\
1.68	0\\
1.682	0\\
1.684	0\\
1.686	0\\
1.688	0\\
1.69	0\\
1.692	0\\
1.694	0\\
1.696	0\\
1.698	0\\
1.7	0\\
1.702	0\\
1.704	0\\
1.706	0\\
1.708	0\\
1.71	0\\
1.712	0\\
1.714	0\\
1.716	0\\
1.718	0\\
1.72	0\\
1.722	0\\
1.724	0\\
1.726	0\\
1.728	0\\
1.73	0\\
1.732	0\\
1.734	0\\
1.736	0\\
1.738	0\\
1.74	0\\
1.742	0\\
1.744	0\\
1.746	0\\
1.748	0\\
1.75	0\\
1.752	0\\
1.754	0\\
1.756	0\\
1.758	0\\
1.76	0\\
1.762	0\\
1.764	0\\
1.766	0\\
1.768	0\\
1.77	0\\
1.772	0\\
1.774	0\\
1.776	0\\
1.778	0\\
1.78	0\\
1.782	0\\
1.784	0\\
1.786	0\\
1.788	0\\
1.79	0\\
1.792	0\\
1.794	0\\
1.796	0\\
1.798	0\\
1.8	0\\
1.802	0\\
1.804	0\\
1.806	0\\
1.808	0\\
1.81	0\\
1.812	0\\
1.814	0\\
1.816	0\\
1.818	0\\
1.82	0\\
1.822	0\\
1.824	0\\
1.826	0\\
1.828	0\\
1.83	0\\
1.832	0\\
1.834	0\\
1.836	0\\
1.838	0\\
1.84	0\\
1.842	0\\
1.844	0\\
1.846	0\\
1.848	0\\
1.85	0\\
1.852	0\\
1.854	0\\
1.856	0\\
1.858	0\\
1.86	0\\
1.862	0\\
1.864	0\\
1.866	0\\
1.868	0\\
1.87	0\\
1.872	0\\
1.874	0\\
1.876	0\\
1.878	0\\
1.88	0\\
1.882	0\\
1.884	0\\
1.886	0\\
1.888	0\\
1.89	0\\
1.892	0\\
1.894	0\\
1.896	0\\
1.898	0\\
1.9	0\\
1.902	0\\
1.904	0\\
1.906	0\\
1.908	0\\
1.91	0\\
1.912	0\\
1.914	0\\
1.916	0\\
1.918	0\\
1.92	0\\
1.922	0\\
1.924	0\\
1.926	0\\
1.928	0\\
1.93	0\\
1.932	0\\
1.934	0\\
1.936	0\\
1.938	0\\
1.94	0\\
1.942	0\\
1.944	0\\
1.946	0\\
1.948	0\\
1.95	0\\
1.952	0\\
1.954	0\\
1.956	0\\
1.958	0\\
1.96	0\\
1.962	0\\
1.964	0\\
1.966	0\\
1.968	0\\
1.97	0\\
1.972	0\\
1.974	0\\
1.976	0\\
1.978	0\\
1.98	0\\
1.982	0\\
1.984	0\\
1.986	0\\
1.988	0\\
1.99	0\\
1.992	0\\
1.994	0\\
1.996	0\\
1.998	0\\
2	0\\
};
\addplot [color=mycolor2, line width=1.0pt, forget plot]
  table[row sep=crcr]{%
-0.024	0\\
-0.022	0\\
-0.02	0\\
-0.018	0\\
-0.016	0\\
-0.014	0\\
-0.012	0\\
-0.01	0\\
-0.008	0\\
-0.006	0\\
-0.004	0\\
-0.002	0\\
0	0\\
0.002	-1.73174e-09\\
0.004	-2.6482e-08\\
0.006	-1.25426e-07\\
0.008	-3.70293e-07\\
0.01	-8.44495e-07\\
0.012	-1.63617e-06\\
0.014	-2.83284e-06\\
0.016	-4.51712e-06\\
0.018	-4.65848e-06\\
0.02	-6.69424e-06\\
0.022	-9.24411e-06\\
0.024	-4.5812e-06\\
0.026	-6.171e-06\\
0.028	-8.12105e-06\\
0.03	-1.04743e-05\\
0.032	-1.32739e-05\\
0.034	-1.65627e-05\\
0.036	-2.03828e-05\\
0.038	-2.47751e-05\\
0.04	-2.97786e-05\\
0.042	-3.54302e-05\\
0.044	-4.1764e-05\\
0.046	-4.8811e-05\\
0.048	-5.65986e-05\\
0.05	-6.51499e-05\\
0.052	-7.44839e-05\\
0.054	-8.46146e-05\\
0.056	-9.55513e-05\\
0.058	-0.000107297\\
0.06	-0.000119851\\
0.062	-0.000133205\\
0.064	-7.92865e-05\\
0.066	-8.84203e-05\\
0.068	-9.81825e-05\\
0.07	-0.000108576\\
0.072	-0.000119602\\
0.074	-0.000131256\\
0.076	-0.000143533\\
0.078	-0.000156421\\
0.08	-0.000169907\\
0.082	-0.000183975\\
0.084	-0.000565156\\
0.086	-0.000489512\\
0.088	-0.000410982\\
0.09	-0.000329905\\
0.092	-0.000249645\\
0.094	-0.000260515\\
0.096	-0.000277151\\
0.098	-0.000452149\\
0.1	-0.0023632\\
0.102	-0.00424243\\
0.104	-0.00608183\\
0.106	-0.00787374\\
0.108	-0.00961091\\
0.11	-0.0112865\\
0.112	-0.0122766\\
0.114	-0.0105751\\
0.116	-0.00883824\\
0.118	-0.00707119\\
0.12	-0.00527928\\
0.122	-0.00346799\\
0.124	-0.00164298\\
0.126	-0.00034124\\
0.128	-0.000358811\\
0.13	-0.000288584\\
0.132	-0.000302645\\
0.134	-0.000316894\\
0.136	-0.0003313\\
0.138	-0.000345833\\
0.14	-0.000360462\\
0.142	-0.000375153\\
0.144	-0.000389871\\
0.146	-0.000404982\\
0.148	-0.000427123\\
0.15	-0.000442731\\
0.152	-0.000458238\\
0.154	-0.0004736\\
0.156	-0.000488772\\
0.158	-0.000503705\\
0.16	-0.000518353\\
0.162	-0.000532667\\
0.164	-0.0005466\\
0.166	-0.000541268\\
0.168	-0.000520194\\
0.17	-0.000498134\\
0.172	-0.000475163\\
0.174	-0.000451362\\
0.176	-0.000426815\\
0.178	-0.000401608\\
0.18	-0.000375829\\
0.182	-0.000349569\\
0.184	-0.000322918\\
0.186	-0.00029597\\
0.188	-0.000268818\\
0.19	-0.000241555\\
0.192	-0.000214274\\
0.194	-0.000187066\\
0.196	-0.000160024\\
0.198	-0.000133236\\
0.2	-0.000106791\\
0.202	-8.0774e-05\\
0.204	8.18852e-05\\
0.206	0.000127206\\
0.208	0.000171494\\
0.21	0.000214581\\
0.212	0.000256301\\
0.214	0.000296497\\
0.216	0.000335015\\
0.218	0.000371709\\
0.22	0.000406442\\
0.222	0.000439084\\
0.224	0.000469514\\
0.226	0.000497622\\
0.228	0.000523307\\
0.23	0.000546478\\
0.232	0.000567057\\
0.234	0.000584977\\
0.236	0.000600182\\
0.238	0.000612628\\
0.24	0.000622285\\
0.242	0.000629134\\
0.244	0.000633168\\
0.246	0.000634394\\
0.248	0.000632831\\
0.25	0.000628511\\
0.252	0.000621476\\
0.254	0.000344914\\
0.256	0.000364441\\
0.258	0.00038262\\
0.26	0.000399401\\
0.262	0.000414735\\
0.264	0.000428581\\
0.266	0.000440901\\
0.268	0.000451662\\
0.27	0.000460836\\
0.272	0.000468401\\
0.274	0.000474339\\
0.276	0.000478639\\
0.278	0.000481294\\
0.28	0.000482303\\
0.282	0.00048167\\
0.284	0.000479405\\
0.286	0.000475523\\
0.288	0.000470044\\
0.29	0.000462995\\
0.292	0.000454405\\
0.294	8.72421e-05\\
0.296	-0.000525347\\
0.298	-0.000394655\\
0.3	-0.000262526\\
0.302	-0.000400492\\
0.304	-0.000688298\\
0.306	-0.000955806\\
0.308	-0.0012024\\
0.31	-0.00142761\\
0.312	-0.00163113\\
0.314	-0.00181278\\
0.316	-0.00197255\\
0.318	-0.00211055\\
0.32	-0.00222703\\
0.322	-0.00232236\\
0.324	-0.00239704\\
0.326	-0.0024517\\
0.328	-0.00248703\\
0.33	-0.00250388\\
0.332	-0.00250314\\
0.334	-0.00248581\\
0.336	-0.00241984\\
0.338	-0.00127777\\
0.34	-0.00127884\\
0.342	-0.00128553\\
0.344	-0.0012975\\
0.346	-0.00131439\\
0.348	-0.00133575\\
0.35	-0.00136111\\
0.352	-0.000787078\\
0.354	-1.6155e-05\\
0.356	0.00036418\\
0.358	0.000349015\\
0.36	0.000332968\\
0.362	0.000316091\\
0.364	0.000298443\\
0.366	0.00028008\\
0.368	0.000261065\\
0.37	0.000241458\\
0.372	0.000221326\\
0.374	0.000200732\\
0.376	0.000179745\\
0.378	0.000158432\\
0.38	0.000136861\\
0.382	0.000115102\\
0.384	9.32241e-05\\
0.386	7.12978e-05\\
0.388	4.93926e-05\\
0.39	8.70629e-05\\
0.392	7.40692e-05\\
0.394	6.1056e-05\\
0.396	4.80611e-05\\
0.398	3.5122e-05\\
0.4	2.22761e-05\\
0.402	9.56003e-06\\
0.404	-2.99015e-06\\
0.406	-0.000509111\\
0.408	-0.000451252\\
0.41	-0.000324948\\
0.412	-0.000198143\\
0.414	-7.12634e-05\\
0.416	5.52655e-05\\
0.418	0.000181022\\
0.42	0.000288147\\
0.422	0.000303053\\
0.424	0.000283814\\
0.426	0.000263493\\
0.428	0.00024216\\
0.43	0.000219883\\
0.432	0.000196736\\
0.434	7.14362e-05\\
0.436	5.74121e-05\\
0.438	4.34487e-05\\
0.44	2.9579e-05\\
0.442	1.58359e-05\\
0.444	2.25147e-06\\
0.446	-1.11427e-05\\
0.448	-2.43158e-05\\
0.45	-3.72376e-05\\
0.452	-4.98788e-05\\
0.454	-6.22109e-05\\
0.456	-7.42066e-05\\
0.458	-8.58393e-05\\
0.46	-9.70836e-05\\
0.462	-0.000107915\\
0.464	-0.000118311\\
0.466	-0.000128249\\
0.468	-0.000137709\\
0.47	-0.000134239\\
0.472	0.000414863\\
0.474	0.000429591\\
0.476	0.00044308\\
0.478	0.00045529\\
0.48	0.000466187\\
0.482	0.000475739\\
0.484	0.000483921\\
0.486	0.000490712\\
0.488	0.000496096\\
0.49	0.000500062\\
0.492	0.000502604\\
0.494	0.00050372\\
0.496	0.000503415\\
0.498	0.000501695\\
0.5	0.000498576\\
0.502	0.000494074\\
0.504	0.000488213\\
0.506	0.00048102\\
0.508	0.000212565\\
0.51	-0.00058217\\
0.512	-0.000592159\\
0.514	-0.000480313\\
0.516	-0.000367128\\
0.518	-0.000252983\\
0.52	-0.00013826\\
0.522	-2.33394e-05\\
0.524	9.13984e-05\\
0.526	0.000205576\\
0.528	0.000318819\\
0.53	0.000414279\\
0.532	0.000437518\\
0.534	0.000421022\\
0.536	0.000403534\\
0.538	0.000385114\\
0.54	0.000365829\\
0.542	0.000345743\\
0.544	0.000324928\\
0.546	0.000303454\\
0.548	0.000281394\\
0.55	0.000258823\\
0.552	0.000235818\\
0.554	0.000212455\\
0.556	0.000188811\\
0.558	0.000369966\\
0.56	9.5078e-05\\
0.562	-0.000255165\\
0.564	-0.000602643\\
0.566	-0.000946333\\
0.568	-0.00128523\\
0.57	-0.00161835\\
0.572	-0.00194472\\
0.574	-0.00226341\\
0.576	-0.00257351\\
0.578	-0.00287413\\
0.58	-0.00297172\\
0.582	-0.00267746\\
0.584	-0.00237951\\
0.586	-0.0020787\\
0.588	-0.00177588\\
0.59	-0.00147188\\
0.592	-0.00116753\\
0.594	-0.000863642\\
0.596	-0.000561025\\
0.598	-0.000260466\\
0.6	3.72623e-05\\
0.602	0.000192953\\
0.604	0.000176681\\
0.606	0.00016025\\
0.608	0.000143713\\
0.61	0.000127122\\
0.612	0.000110528\\
0.614	9.39822e-05\\
0.616	-0.000156818\\
0.618	-0.000868285\\
0.62	-0.000774444\\
0.622	-0.000678908\\
0.624	-0.000581999\\
0.626	-0.000484041\\
0.628	-0.000385362\\
0.63	-0.000286289\\
0.632	-0.000187148\\
0.634	-8.8263e-05\\
0.636	1.00435e-05\\
0.638	0.000107454\\
0.64	0.000190452\\
0.642	0.000217175\\
0.644	0.000206034\\
0.646	0.000191524\\
0.648	0.000176234\\
0.65	0.000160221\\
0.652	0.000143543\\
0.654	0.000126261\\
0.656	0.000108434\\
0.658	9.01267e-05\\
0.66	7.14026e-05\\
0.662	-1.99724e-05\\
0.664	-2.59989e-05\\
0.666	-3.17367e-05\\
0.668	-3.71744e-05\\
0.67	-4.23017e-05\\
0.672	-4.71092e-05\\
0.674	-5.15885e-05\\
0.676	-5.57322e-05\\
0.678	-8.1045e-05\\
0.68	-0.000194007\\
0.682	8.26721e-05\\
0.684	0.000285507\\
0.686	0.00029655\\
0.688	0.000306731\\
0.69	0.000316022\\
0.692	0.000324395\\
0.694	0.000331829\\
0.696	0.000338302\\
0.698	0.0003438\\
0.7	0.00034831\\
0.702	0.000351822\\
0.704	0.000354331\\
0.706	0.000355836\\
0.708	0.000356338\\
0.71	0.000355842\\
0.712	0.000354357\\
0.714	0.000351894\\
0.716	0.000348469\\
0.718	0.0003441\\
0.72	0.000338809\\
0.722	0.000277004\\
0.724	-0.000245644\\
0.726	-0.000678953\\
0.728	-0.000602709\\
0.73	-0.000524824\\
0.732	-0.000445564\\
0.734	-0.0003652\\
0.736	-0.000284001\\
0.738	-0.000202239\\
0.74	-0.000120187\\
0.742	-3.81142e-05\\
0.744	4.371e-05\\
0.746	0.00012502\\
0.748	0.000205553\\
0.75	0.000276703\\
0.752	0.000298451\\
0.754	0.000297312\\
0.756	0.0002848\\
0.758	0.000271652\\
0.76	0.000257915\\
0.762	0.000243637\\
0.764	0.000228869\\
0.766	0.000213662\\
0.768	0.000198067\\
0.77	0.000182138\\
0.772	0.000165929\\
0.774	0.000149493\\
0.776	0.000132886\\
0.778	0.000116162\\
0.78	0.00027232\\
0.782	0.000261452\\
0.784	0.000250033\\
0.786	0.000238102\\
0.788	0.0002257\\
0.79	0.000212868\\
0.792	-8.64292e-06\\
0.794	-0.000252304\\
0.796	-0.00049324\\
0.798	-0.000730754\\
0.8	-0.000964162\\
0.802	-0.0011928\\
0.804	-0.00141602\\
0.806	-0.00163321\\
0.808	-0.00184375\\
0.81	-0.00204706\\
0.812	-0.00205253\\
0.814	-0.00194339\\
0.816	-0.00183022\\
0.818	-0.00171338\\
0.82	-0.00159325\\
0.822	-0.00147018\\
0.824	-0.00134454\\
0.826	-0.00121671\\
0.828	-0.00108706\\
0.83	-0.000955959\\
0.832	-0.000823776\\
0.834	-0.000690883\\
0.836	-0.000557645\\
0.838	-0.000424425\\
0.84	-0.000545272\\
0.842	-0.000480116\\
0.844	-0.000414062\\
0.846	-0.000347332\\
0.848	-0.000280146\\
0.85	-0.000212726\\
0.852	-0.000145294\\
0.854	-7.80677e-05\\
0.856	-1.12637e-05\\
0.858	5.49041e-05\\
0.86	0.000116783\\
0.862	0.000135021\\
0.864	0.000142262\\
0.866	0.000132566\\
0.868	0.000122359\\
0.87	0.000111677\\
0.872	0.000100562\\
0.874	8.90536e-05\\
0.876	7.71932e-05\\
0.878	6.50233e-05\\
0.88	5.25869e-05\\
0.882	3.99277e-05\\
0.884	2.70897e-05\\
0.886	1.41169e-05\\
0.888	1.05376e-06\\
0.89	3.0302e-05\\
0.892	2.6804e-05\\
0.894	2.33741e-05\\
0.896	2.00198e-05\\
0.898	1.67481e-05\\
0.9	7.80718e-05\\
0.902	0.00017247\\
0.904	0.000198893\\
0.906	0.000204143\\
0.908	0.000208791\\
0.91	0.000212824\\
0.912	0.000216235\\
0.914	0.000219016\\
0.916	0.000221163\\
0.918	0.000222672\\
0.92	0.000223545\\
0.922	0.000223782\\
0.924	0.000223388\\
0.926	0.00022237\\
0.928	0.000220736\\
0.93	0.000218497\\
0.932	0.000161495\\
0.934	-9.38129e-05\\
0.936	-0.000345926\\
0.938	-0.000593983\\
0.94	-0.000837148\\
0.942	-0.00096006\\
0.944	-0.000814727\\
0.946	-0.000667401\\
0.948	-0.000518583\\
0.95	-0.000368777\\
0.952	-0.000315812\\
0.954	-0.000263131\\
0.956	-0.000209897\\
0.958	-0.000156289\\
0.96	-0.000102485\\
0.962	-4.86622e-05\\
0.964	5.00278e-06\\
0.966	5.83364e-05\\
0.968	0.000111167\\
0.97	0.000163053\\
0.972	0.000177367\\
0.974	0.000188842\\
0.976	0.000180706\\
0.978	0.000172167\\
0.98	0.000163255\\
0.982	0.000154003\\
0.984	0.000144443\\
0.986	0.000134608\\
0.988	0.000124532\\
0.99	0.000114251\\
0.992	0.000103798\\
0.994	9.32094e-05\\
0.996	8.25204e-05\\
0.998	7.17663e-05\\
1	6.09824e-05\\
1.002	9.85731e-05\\
1.004	0.000173341\\
1.006	0.000166725\\
1.008	0.000159769\\
1.01	0.000152429\\
1.012	0.000144788\\
1.014	0.000136874\\
1.016	0.000128712\\
1.018	0.00012033\\
1.02	0.000111755\\
1.022	4.95681e-05\\
1.024	-0.000105271\\
1.026	-0.000258363\\
1.028	-0.000409265\\
1.03	-0.000557546\\
1.032	-0.000702784\\
1.034	-0.000844572\\
1.036	-0.000982517\\
1.038	-0.00111624\\
1.04	-0.00122233\\
1.042	-0.00122399\\
1.044	-0.00119266\\
1.046	-0.00115841\\
1.048	-0.00112139\\
1.05	-0.00108172\\
1.052	-0.00103954\\
1.054	-0.000994998\\
1.056	-0.000948236\\
1.058	-0.000825176\\
1.06	-0.000615126\\
1.062	-0.000408651\\
1.064	-0.00027772\\
1.066	-0.000235203\\
1.068	-0.000192383\\
1.07	-0.000149399\\
1.072	-0.000106393\\
1.074	-6.35052e-05\\
1.076	-2.08729e-05\\
1.078	2.13675e-05\\
1.08	6.30826e-05\\
1.082	7.63483e-05\\
1.084	8.75361e-05\\
1.086	8.4179e-05\\
1.088	7.76524e-05\\
1.09	7.08347e-05\\
1.092	6.37514e-05\\
1.094	5.64287e-05\\
1.096	4.8893e-05\\
1.098	4.11716e-05\\
1.1	3.32919e-05\\
1.102	2.52818e-05\\
1.104	1.7169e-05\\
1.106	8.98181e-06\\
1.108	7.48139e-07\\
1.11	-7.50401e-06\\
1.112	0.000109224\\
1.114	0.000113421\\
1.116	0.000117277\\
1.118	0.000120782\\
1.12	0.000123925\\
1.122	0.000126699\\
1.124	0.000129098\\
1.126	0.000131117\\
1.128	0.000132751\\
1.13	0.000133998\\
1.132	0.000134857\\
1.134	0.000135329\\
1.136	0.000135414\\
1.138	0.000135116\\
1.14	0.000134439\\
1.142	8.22779e-05\\
1.144	2.05401e-05\\
1.146	-4.30327e-05\\
1.148	-0.000186196\\
1.15	-0.000298774\\
1.152	-0.000337886\\
1.154	-0.000373877\\
1.156	-0.000406646\\
1.158	-0.00043611\\
1.16	-0.0004622\\
1.162	-0.000484865\\
1.164	-0.000497471\\
1.166	-0.000397421\\
1.168	-0.00029684\\
1.17	-0.000242874\\
1.172	-0.000210557\\
1.174	-0.000177767\\
1.176	-0.000144616\\
1.178	-0.000111215\\
1.18	-7.76767e-05\\
1.182	-4.41105e-05\\
1.184	-1.06265e-05\\
1.186	2.26665e-05\\
1.188	5.56617e-05\\
1.19	8.8254e-05\\
1.192	9.97268e-05\\
1.194	0.000108285\\
1.196	0.000108925\\
1.198	0.000103529\\
1.2	9.791e-05\\
1.202	9.20888e-05\\
1.204	8.60858e-05\\
1.206	7.9922e-05\\
1.208	7.3619e-05\\
1.21	6.71983e-05\\
1.212	6.06819e-05\\
1.214	5.40917e-05\\
1.216	4.74498e-05\\
1.218	4.07782e-05\\
1.22	3.40989e-05\\
1.222	2.74335e-05\\
1.224	2.08037e-05\\
1.226	7.95205e-05\\
1.228	0.000102669\\
1.23	9.88338e-05\\
1.232	9.48023e-05\\
1.234	9.06523e-05\\
1.236	8.62467e-05\\
1.238	8.16788e-05\\
1.24	7.69642e-05\\
1.242	7.21189e-05\\
1.244	6.71592e-05\\
1.246	6.21015e-05\\
1.248	5.69623e-05\\
1.25	5.17583e-05\\
1.252	4.46819e-05\\
1.254	-4.82354e-05\\
1.256	-0.000140083\\
1.258	-0.000230598\\
1.26	-0.00031952\\
1.262	-0.000406599\\
1.264	-0.000491591\\
1.266	-0.000574262\\
1.268	-0.000654387\\
1.27	-0.000673436\\
1.272	-0.000671784\\
1.274	-0.000668289\\
1.276	-0.000601405\\
1.278	-0.000519528\\
1.28	-0.000438447\\
1.282	-0.000358406\\
1.284	-0.000279644\\
1.286	-0.000202388\\
1.288	-0.000126858\\
1.29	-9.9625e-05\\
1.292	-7.32735e-05\\
1.294	-4.69782e-05\\
1.296	-2.08238e-05\\
1.298	5.10619e-06\\
1.3	3.07297e-05\\
1.302	4.08404e-05\\
1.304	4.76756e-05\\
1.306	5.10661e-05\\
1.308	4.69876e-05\\
1.31	4.27393e-05\\
1.312	3.83369e-05\\
1.314	3.37969e-05\\
1.316	2.91356e-05\\
1.318	2.43699e-05\\
1.32	1.95168e-05\\
1.322	1.45934e-05\\
1.324	9.61692e-06\\
1.326	4.60449e-06\\
1.328	-4.26689e-07\\
1.33	-5.45953e-06\\
1.332	5.37116e-05\\
1.334	7.30736e-05\\
1.336	7.46845e-05\\
1.338	7.60717e-05\\
1.34	7.72323e-05\\
1.342	7.81643e-05\\
1.344	7.88663e-05\\
1.346	7.9338e-05\\
1.348	6.04758e-05\\
1.35	2.80268e-05\\
1.352	-3.99929e-06\\
1.354	-3.54902e-05\\
1.356	-6.63369e-05\\
1.358	-9.64342e-05\\
1.36	-0.000125681\\
1.362	-0.000161817\\
1.364	-0.000241571\\
1.366	-0.000319593\\
1.368	-0.000395626\\
1.37	-0.000459836\\
1.372	-0.000518399\\
1.374	-0.000575009\\
1.376	-0.000529507\\
1.378	-0.000470683\\
1.38	-0.000410802\\
1.382	-0.000350071\\
1.384	-0.000288697\\
1.386	-0.000226888\\
1.388	-0.000174123\\
1.39	-0.000155051\\
1.392	-0.000135613\\
1.394	-0.000115875\\
1.396	-9.59046e-05\\
1.398	-7.57686e-05\\
1.4	-5.55347e-05\\
1.402	-3.52697e-05\\
1.404	-1.50401e-05\\
1.406	5.0884e-06\\
1.408	2.50511e-05\\
1.41	4.47843e-05\\
1.412	5.31745e-05\\
1.414	5.83127e-05\\
1.416	6.28595e-05\\
1.418	5.95179e-05\\
1.42	5.6049e-05\\
1.422	5.24654e-05\\
1.424	4.87796e-05\\
1.426	4.50047e-05\\
1.428	4.11536e-05\\
1.43	3.72398e-05\\
1.432	3.32764e-05\\
1.434	2.92768e-05\\
1.436	2.52544e-05\\
1.438	2.12225e-05\\
1.44	1.71944e-05\\
1.442	1.3183e-05\\
1.444	9.20143e-06\\
1.446	5.26224e-06\\
1.448	2.31656e-06\\
1.45	5.8967e-05\\
1.452	5.7006e-05\\
1.454	5.49225e-05\\
1.456	5.27242e-05\\
1.458	5.04201e-05\\
1.46	4.8082e-05\\
1.462	4.55921e-05\\
1.464	4.30206e-05\\
1.466	4.03762e-05\\
1.468	3.7668e-05\\
1.47	3.49049e-05\\
1.472	3.20961e-05\\
1.474	2.92508e-05\\
1.476	2.63782e-05\\
1.478	2.34874e-05\\
1.48	-9.98916e-06\\
1.482	-4.84217e-05\\
1.484	-8.55195e-05\\
1.486	-0.000121183\\
1.488	-0.000155319\\
1.49	-0.000187841\\
1.492	-0.000229454\\
1.494	-0.000264676\\
1.496	-0.000236897\\
1.498	-0.000209394\\
1.5	-0.000182256\\
1.502	-0.000155565\\
1.504	-0.000129403\\
1.506	-0.000103847\\
1.508	-7.96693e-05\\
1.51	-6.3994e-05\\
1.512	-4.82895e-05\\
1.514	-3.26069e-05\\
1.516	-1.69969e-05\\
1.518	-1.50949e-06\\
1.52	1.38063e-05\\
1.522	2.08848e-05\\
1.524	2.49203e-05\\
1.526	2.88066e-05\\
1.528	2.75731e-05\\
1.53	2.49578e-05\\
1.532	2.22565e-05\\
1.534	1.94792e-05\\
1.536	1.66358e-05\\
1.538	1.37366e-05\\
1.54	1.07919e-05\\
1.542	7.81201e-06\\
1.544	4.80728e-06\\
1.546	1.78804e-06\\
1.548	-1.2354e-06\\
1.55	-4.25285e-06\\
1.552	-7.25422e-06\\
1.554	-3.12889e-06\\
1.556	2.20879e-05\\
1.558	5.16217e-06\\
1.56	-1.1515e-05\\
1.562	-2.78854e-05\\
1.564	-4.38927e-05\\
1.566	-5.94825e-05\\
1.568	-7.46027e-05\\
1.57	-7.65321e-05\\
1.572	-7.02118e-05\\
1.574	-6.37387e-05\\
1.576	-5.719e-05\\
1.578	-7.17842e-05\\
1.58	-9.69582e-05\\
1.582	-0.000121671\\
1.584	-0.000145843\\
1.586	-0.000169396\\
1.588	-0.000192258\\
1.59	-0.000214359\\
1.592	-0.000235634\\
1.594	-0.000256019\\
1.596	-0.000275459\\
1.598	-0.000269235\\
1.6	-0.000233023\\
1.602	-0.00019637\\
1.604	-0.000159402\\
1.606	-0.000122244\\
1.608	-0.000107256\\
1.61	-9.60657e-05\\
1.612	-8.46499e-05\\
1.614	-7.3048e-05\\
1.616	-6.12996e-05\\
1.618	-4.94445e-05\\
1.62	-3.75224e-05\\
1.622	-2.55729e-05\\
1.624	-1.36354e-05\\
1.626	-1.74852e-06\\
1.628	1.00493e-05\\
1.63	2.17203e-05\\
1.632	2.73941e-05\\
1.634	3.03894e-05\\
1.636	3.32756e-05\\
1.638	3.33982e-05\\
1.64	3.12821e-05\\
1.642	2.91029e-05\\
1.644	2.68683e-05\\
1.646	2.45863e-05\\
1.648	2.22647e-05\\
1.65	1.99115e-05\\
1.652	1.75346e-05\\
1.654	1.51421e-05\\
1.656	1.2742e-05\\
1.658	1.03421e-05\\
1.66	7.9503e-06\\
1.662	5.5744e-06\\
1.664	3.22199e-06\\
1.666	9.0056e-07\\
1.668	-1.3826e-06\\
1.67	-2.77684e-06\\
1.672	4.13263e-06\\
1.674	3.80082e-06\\
1.676	1.09505e-05\\
1.678	1.80041e-05\\
1.68	2.47332e-05\\
1.682	2.68323e-05\\
1.684	2.55843e-05\\
1.686	2.43333e-05\\
1.688	2.29976e-05\\
1.69	2.16237e-05\\
1.692	2.02164e-05\\
1.694	1.50114e-05\\
1.696	5.31776e-07\\
1.698	-1.34439e-05\\
1.7	-2.68768e-05\\
1.702	-3.97306e-05\\
1.704	-5.19715e-05\\
1.706	-6.3569e-05\\
1.708	-7.44948e-05\\
1.71	-8.47241e-05\\
1.712	-9.42346e-05\\
1.714	-0.000103007\\
1.716	-0.000101747\\
1.718	-9.33897e-05\\
1.72	-8.48269e-05\\
1.722	-7.60891e-05\\
1.724	-6.72062e-05\\
1.726	-5.82085e-05\\
1.728	-4.91264e-05\\
1.73	-3.99901e-05\\
1.732	-3.08299e-05\\
1.734	-2.16758e-05\\
1.736	-1.25574e-05\\
1.738	-3.50389e-06\\
1.74	5.45589e-06\\
1.742	1.00348e-05\\
1.744	1.23481e-05\\
1.746	1.45819e-05\\
1.748	1.55709e-05\\
1.75	1.39714e-05\\
1.752	1.23255e-05\\
1.754	1.06392e-05\\
1.756	8.9186e-06\\
1.758	7.16984e-06\\
1.76	5.39905e-06\\
1.762	3.61243e-06\\
1.764	1.81614e-06\\
1.766	1.63176e-08\\
1.768	-1.78094e-06\\
1.77	-3.56959e-06\\
1.772	-5.3437e-06\\
1.774	-7.09743e-06\\
1.776	-2.14136e-05\\
1.778	-4.23257e-05\\
1.78	-5.86441e-05\\
1.782	-6.56575e-05\\
1.784	-7.23152e-05\\
1.786	-7.13899e-05\\
1.788	-6.14271e-05\\
1.79	-5.14069e-05\\
1.792	-4.41373e-05\\
1.794	-5.32077e-05\\
1.796	-6.20715e-05\\
1.798	-7.06665e-05\\
1.8	-7.89647e-05\\
1.802	-8.69716e-05\\
1.804	-9.46621e-05\\
1.806	-0.000102012\\
1.808	-0.000109\\
1.81	-0.000115606\\
1.812	-0.000121702\\
1.814	-0.000126174\\
1.816	-0.000130055\\
1.818	-0.000133342\\
1.82	-0.00012502\\
1.822	-0.0001037\\
1.824	-8.22417e-05\\
1.826	-7.06051e-05\\
1.828	-6.43259e-05\\
1.83	-5.78885e-05\\
1.832	-5.13155e-05\\
1.834	-4.46298e-05\\
1.836	-3.78544e-05\\
1.838	-3.10123e-05\\
1.84	-2.41264e-05\\
1.842	-1.72199e-05\\
1.844	-1.03154e-05\\
1.846	-3.43542e-06\\
1.848	3.3977e-06\\
1.85	1.01621e-05\\
1.852	1.36979e-05\\
1.854	1.539e-05\\
1.856	1.70246e-05\\
1.858	1.8347e-05\\
1.86	1.70563e-05\\
1.862	1.57321e-05\\
1.864	1.4379e-05\\
1.866	1.30018e-05\\
1.868	1.16052e-05\\
1.87	1.01939e-05\\
1.872	8.77272e-06\\
1.874	7.34644e-06\\
1.876	5.91975e-06\\
1.878	4.49735e-06\\
1.88	3.08384e-06\\
1.882	1.68377e-06\\
1.884	3.01614e-07\\
1.886	-1.05827e-06\\
1.888	-2.39164e-06\\
1.89	-3.69437e-06\\
1.892	-4.9625e-06\\
1.894	-7.1043e-06\\
1.896	1.87805e-06\\
1.898	1.09236e-05\\
1.9	1.51293e-05\\
1.902	1.37987e-05\\
1.904	8.38708e-06\\
1.906	3.16938e-06\\
1.908	-1.83837e-06\\
1.91	-6.62103e-06\\
1.912	-1.11169e-05\\
1.914	-1.5362e-05\\
1.916	-1.93452e-05\\
1.918	-2.30561e-05\\
1.92	-2.6486e-05\\
1.922	-2.96269e-05\\
1.924	-3.24725e-05\\
1.926	-3.50175e-05\\
1.928	-3.7258e-05\\
1.93	-3.91911e-05\\
1.932	-3.98736e-05\\
1.934	-3.81818e-05\\
1.936	-4.67247e-05\\
1.938	-5.46437e-05\\
1.94	-4.97564e-05\\
1.942	-4.47652e-05\\
1.944	-3.96873e-05\\
1.946	-3.45402e-05\\
1.948	-2.93411e-05\\
1.95	-2.41076e-05\\
1.952	-1.88569e-05\\
1.954	-1.36065e-05\\
1.956	-8.37323e-06\\
1.958	-3.17403e-06\\
1.96	1.97455e-06\\
1.962	4.72901e-06\\
1.964	6.01794e-06\\
1.966	7.26531e-06\\
1.968	8.46758e-06\\
1.97	7.76694e-06\\
1.972	6.76426e-06\\
1.974	5.74076e-06\\
1.976	4.70007e-06\\
1.978	2.82786e-06\\
1.98	-2.2387e-06\\
1.982	-7.39601e-06\\
1.984	-1.26249e-05\\
1.986	-1.7906e-05\\
1.988	-2.32198e-05\\
1.99	-2.85466e-05\\
1.992	-3.38668e-05\\
1.994	-3.91609e-05\\
1.996	-4.44094e-05\\
1.998	-4.95932e-05\\
2	-5.46933e-05\\
};
\addplot [color=mycolor3, line width=1.0pt, forget plot]
  table[row sep=crcr]{%
-0.024	0\\
-0.022	0\\
-0.02	0\\
-0.018	0\\
-0.016	0\\
-0.014	0\\
-0.012	0\\
-0.01	0\\
-0.008	0\\
-0.006	0\\
-0.004	0\\
-0.002	0\\
0	0\\
0.002	-2.73542e-09\\
0.004	-4.17916e-08\\
0.006	-1.98834e-07\\
0.008	-5.90026e-07\\
0.01	-1.35265e-06\\
0.012	-2.63442e-06\\
0.014	-4.5852e-06\\
0.016	-7.35026e-06\\
0.018	-7.6714e-06\\
0.02	-1.10851e-05\\
0.022	-1.53927e-05\\
0.024	-6.49926e-06\\
0.026	-8.81252e-06\\
0.028	-1.16733e-05\\
0.03	-1.51536e-05\\
0.032	-1.9327e-05\\
0.034	-2.42682e-05\\
0.036	-3.00519e-05\\
0.038	-3.67521e-05\\
0.04	-4.44414e-05\\
0.042	-5.31897e-05\\
0.044	-6.3064e-05\\
0.046	-7.41267e-05\\
0.048	-8.64355e-05\\
0.05	-0.000100042\\
0.052	-0.000114992\\
0.054	-0.000131321\\
0.056	-0.000149061\\
0.058	-0.000168232\\
0.06	-0.000188846\\
0.062	-0.000210904\\
0.064	-0.000144109\\
0.066	-0.000161214\\
0.068	-0.000179562\\
0.07	-0.000199167\\
0.072	-0.000220037\\
0.074	-0.000242173\\
0.076	-0.000265569\\
0.078	-0.000290212\\
0.08	-0.000316083\\
0.082	-0.000343154\\
0.084	-0.000947786\\
0.086	-0.00159925\\
0.088	-0.00226246\\
0.09	-0.00293446\\
0.092	-0.00361525\\
0.094	-0.00439172\\
0.096	-0.00517503\\
0.098	-0.00611429\\
0.1	-0.00878371\\
0.102	-0.0114119\\
0.104	-0.0139874\\
0.106	-0.0164991\\
0.108	-0.0157186\\
0.11	-0.0141138\\
0.112	-0.0124642\\
0.114	-0.0107746\\
0.116	-0.00905002\\
0.118	-0.00729555\\
0.12	-0.0055165\\
0.122	-0.00371836\\
0.124	-0.00190673\\
0.126	-0.000618596\\
0.128	-0.000649957\\
0.13	-0.000565347\\
0.132	-0.00059253\\
0.134	-0.00062003\\
0.136	-0.000647788\\
0.138	-0.000675742\\
0.14	-0.000703828\\
0.142	-0.000731978\\
0.144	-0.000760122\\
0.146	-0.000766235\\
0.148	-0.000794925\\
0.15	-0.000823467\\
0.152	-0.00085178\\
0.154	-0.000879781\\
0.156	-0.000907383\\
0.158	-0.000934499\\
0.16	-0.000961041\\
0.162	-0.00098692\\
0.164	-0.00101205\\
0.166	-0.00100688\\
0.168	-0.000970892\\
0.17	-0.000932982\\
0.172	-0.000893285\\
0.174	-0.000851942\\
0.176	-0.000809098\\
0.178	-0.000764903\\
0.18	-0.000719513\\
0.182	-0.000673087\\
0.184	-0.000625787\\
0.186	-0.000577776\\
0.188	-0.000529219\\
0.19	-0.000480283\\
0.192	-0.000431134\\
0.194	-0.000381935\\
0.196	-0.000332851\\
0.198	-0.000284043\\
0.2	-0.000235669\\
0.202	-0.000187883\\
0.204	-8.61518e-05\\
0.206	-1.52029e-05\\
0.208	5.47603e-05\\
0.21	0.000123483\\
0.212	0.000190714\\
0.214	0.000256209\\
0.216	0.000319731\\
0.218	0.000381052\\
0.22	0.000439952\\
0.222	0.000496222\\
0.224	0.000549664\\
0.226	0.000600093\\
0.228	0.000647339\\
0.23	0.000691241\\
0.232	0.000731658\\
0.234	0.000768462\\
0.236	0.00080154\\
0.238	0.000830798\\
0.24	0.000856157\\
0.242	0.000877556\\
0.244	0.00089495\\
0.246	0.000908314\\
0.248	0.000917638\\
0.25	0.000922933\\
0.252	0.000924223\\
0.254	0.000843011\\
0.256	0.000883247\\
0.258	0.000920507\\
0.26	0.000954676\\
0.262	0.000985654\\
0.264	0.00101335\\
0.266	0.00103767\\
0.268	0.00105856\\
0.27	0.00107595\\
0.272	0.0010898\\
0.274	0.00110007\\
0.276	0.00110673\\
0.278	0.00110977\\
0.28	0.0011092\\
0.282	0.00110501\\
0.284	0.00109725\\
0.286	0.00108594\\
0.288	0.00107112\\
0.29	0.00105287\\
0.292	0.00103124\\
0.294	0.000649264\\
0.296	-0.000325374\\
0.298	-0.0012981\\
0.3	-0.00226543\\
0.302	-0.00349499\\
0.304	-0.00486263\\
0.306	-0.0061944\\
0.308	-0.00748591\\
0.31	-0.00873301\\
0.312	-0.00993179\\
0.314	-0.0110786\\
0.316	-0.0114174\\
0.318	-0.0106989\\
0.32	-0.00993979\\
0.322	-0.00914329\\
0.324	-0.00831295\\
0.326	-0.00745241\\
0.328	-0.00656543\\
0.33	-0.00565586\\
0.332	-0.00472762\\
0.334	-0.00336898\\
0.336	-0.00198598\\
0.338	-0.000859242\\
0.34	-0.000876758\\
0.342	-0.000900939\\
0.344	-0.000931396\\
0.346	-0.000967686\\
0.348	-0.00100932\\
0.35	-0.00105575\\
0.352	-0.000369052\\
0.354	0.00038723\\
0.356	0.000752033\\
0.358	0.000720496\\
0.36	0.000687289\\
0.362	0.000652522\\
0.364	0.000616307\\
0.366	0.000578762\\
0.368	0.000540009\\
0.37	0.000500173\\
0.372	0.000459382\\
0.374	0.000417765\\
0.376	0.000375455\\
0.378	0.000332587\\
0.38	0.000289296\\
0.382	0.000245719\\
0.384	0.000201992\\
0.386	0.000158253\\
0.388	0.000114638\\
0.39	-0.000139563\\
0.392	-0.00106028\\
0.394	-0.00124286\\
0.396	-0.00113914\\
0.398	-0.00109046\\
0.4	-0.00109563\\
0.402	-0.00115325\\
0.404	-0.00126173\\
0.406	-0.00191303\\
0.408	-0.00281458\\
0.41	-0.00309159\\
0.412	-0.00298165\\
0.414	-0.00281667\\
0.416	-0.00259793\\
0.418	-0.0023269\\
0.42	-0.00202267\\
0.422	-0.00176022\\
0.424	-0.00148305\\
0.426	-0.00115983\\
0.428	-0.000792375\\
0.43	-0.000382639\\
0.432	6.72839e-05\\
0.434	7.59275e-05\\
0.436	-6.00153e-05\\
0.438	-0.000195301\\
0.44	-0.000329654\\
0.442	-0.000462798\\
0.444	-0.000594461\\
0.446	-0.000724376\\
0.448	-0.000852278\\
0.45	-0.000977906\\
0.452	-0.00110101\\
0.454	-0.00122133\\
0.456	-0.00133863\\
0.458	-0.00145267\\
0.46	-0.00156323\\
0.462	-0.00167008\\
0.464	-0.001773\\
0.466	-0.00187181\\
0.468	-0.000874951\\
0.47	0.000400978\\
0.472	0.000969793\\
0.474	0.00100269\\
0.476	0.00103275\\
0.478	0.00105989\\
0.48	0.00108402\\
0.482	0.00110507\\
0.484	0.001123\\
0.486	0.00113775\\
0.488	0.00114929\\
0.49	0.00115759\\
0.492	0.00116264\\
0.494	0.00116444\\
0.496	0.00116299\\
0.498	0.00115832\\
0.5	0.00115045\\
0.502	0.00113942\\
0.504	0.00112529\\
0.506	0.00110812\\
0.508	0.000828009\\
0.51	1.99982e-05\\
0.512	-0.00078276\\
0.514	-0.00157751\\
0.516	-0.00236154\\
0.518	-0.0031322\\
0.52	-0.00388688\\
0.522	-0.00462307\\
0.524	-0.00533832\\
0.526	-0.00603028\\
0.528	-0.00669668\\
0.53	-0.00649671\\
0.532	-0.00599842\\
0.534	-0.00552281\\
0.536	-0.0050329\\
0.538	-0.00453044\\
0.54	-0.00401719\\
0.542	-0.00349492\\
0.544	-0.00296543\\
0.546	-0.00243052\\
0.548	-0.00189199\\
0.55	-0.00135164\\
0.552	-0.00081124\\
0.554	-0.000272566\\
0.556	0.000262639\\
0.558	0.000618099\\
0.56	0.000307262\\
0.562	-7.92286e-05\\
0.564	-0.000463141\\
0.566	-0.000843341\\
0.568	-0.00121871\\
0.57	-0.00158816\\
0.572	-0.00195061\\
0.574	-0.00230501\\
0.576	-0.00265034\\
0.578	-0.00291917\\
0.58	-0.00264139\\
0.582	-0.00235978\\
0.584	-0.00207514\\
0.586	-0.00178826\\
0.588	-0.00149996\\
0.59	-0.00121101\\
0.592	-0.000922204\\
0.594	-0.000634298\\
0.596	-0.000348051\\
0.598	-6.41995e-05\\
0.6	0.000216535\\
0.602	0.000354998\\
0.604	0.000321316\\
0.606	0.000287348\\
0.608	0.000253197\\
0.61	0.000218969\\
0.612	0.000215381\\
0.614	0.000303672\\
0.616	0.000134878\\
0.618	-0.00054059\\
0.62	-0.00120985\\
0.622	-0.00187064\\
0.624	-0.00252071\\
0.626	-0.0031579\\
0.628	-0.00378007\\
0.63	-0.00433169\\
0.632	-0.00449174\\
0.634	-0.00460849\\
0.636	-0.00468231\\
0.638	-0.0047137\\
0.64	-0.00471649\\
0.642	-0.004733\\
0.644	-0.00466275\\
0.646	-0.00428526\\
0.648	-0.00389795\\
0.65	-0.00350217\\
0.652	-0.00309928\\
0.654	-0.00269064\\
0.656	-0.00227762\\
0.658	-0.0018616\\
0.66	-0.00144395\\
0.662	-0.00116168\\
0.664	-0.0010217\\
0.666	-0.000893796\\
0.668	-0.000778044\\
0.67	-0.000674454\\
0.672	-0.000582994\\
0.674	-0.000503589\\
0.676	-0.000436122\\
0.678	-0.000246768\\
0.68	4.75276e-05\\
0.682	0.00034181\\
0.684	0.000561347\\
0.686	0.00058814\\
0.688	0.000613072\\
0.69	0.000636073\\
0.692	0.000657075\\
0.694	0.000676021\\
0.696	0.00069286\\
0.698	0.000707549\\
0.7	0.000720053\\
0.702	0.000730345\\
0.704	0.000738404\\
0.706	0.000744219\\
0.708	0.000747787\\
0.71	0.000749111\\
0.712	0.000748201\\
0.714	0.000745079\\
0.716	0.000739769\\
0.718	0.000732306\\
0.72	0.000722731\\
0.722	0.000655475\\
0.724	0.000126236\\
0.726	-0.000423019\\
0.728	-0.000966035\\
0.73	-0.00150099\\
0.732	-0.00202611\\
0.734	-0.00253967\\
0.736	-0.00303999\\
0.738	-0.00352545\\
0.74	-0.00399451\\
0.742	-0.00444569\\
0.744	-0.00485047\\
0.746	-0.00462588\\
0.748	-0.00438743\\
0.75	-0.00404861\\
0.752	-0.00374748\\
0.754	-0.00345852\\
0.756	-0.00317125\\
0.758	-0.002876\\
0.76	-0.00257377\\
0.762	-0.00226561\\
0.764	-0.00195256\\
0.766	-0.00163568\\
0.768	-0.00131602\\
0.77	-0.000994628\\
0.772	-0.000672558\\
0.774	-0.000350851\\
0.776	-3.05379e-05\\
0.778	0.000287365\\
0.78	0.000491739\\
0.782	0.000454688\\
0.784	0.000416702\\
0.786	0.000377906\\
0.788	0.000338421\\
0.79	0.000298375\\
0.792	4.96015e-05\\
0.794	-0.000221288\\
0.796	-0.000489335\\
0.798	-0.00075376\\
0.8	-0.0010138\\
0.802	-0.00126871\\
0.804	-0.00151777\\
0.806	-0.00176028\\
0.808	-0.00199555\\
0.81	-0.00199441\\
0.812	-0.00189875\\
0.814	-0.00179908\\
0.816	-0.00169572\\
0.818	-0.00158901\\
0.82	-0.00147927\\
0.822	-0.00136684\\
0.824	-0.00125204\\
0.826	-0.00113521\\
0.828	-0.0010167\\
0.83	-0.000896839\\
0.832	-0.000775958\\
0.834	-0.000654394\\
0.836	-0.000532478\\
0.838	-0.000410536\\
0.84	-0.000682119\\
0.842	-0.000976667\\
0.844	-0.00125788\\
0.846	-0.00152508\\
0.848	-0.00177761\\
0.85	-0.00201491\\
0.852	-0.00223644\\
0.854	-0.00244176\\
0.856	-0.00262896\\
0.858	-0.00264589\\
0.86	-0.00264513\\
0.862	-0.00266617\\
0.864	-0.00267659\\
0.866	-0.00268263\\
0.868	-0.00256879\\
0.87	-0.00233668\\
0.872	-0.00209955\\
0.874	-0.00185823\\
0.876	-0.00161351\\
0.878	-0.00136621\\
0.88	-0.00111714\\
0.882	-0.000867119\\
0.884	-0.000616949\\
0.886	-0.000367432\\
0.888	-0.00011936\\
0.89	-4.93738e-05\\
0.892	-0.000117288\\
0.894	-0.000113145\\
0.896	7.8448e-06\\
0.898	0.000123748\\
0.9	0.000234177\\
0.902	0.00033878\\
0.904	0.000374795\\
0.906	0.000388995\\
0.908	0.000401928\\
0.91	0.000413561\\
0.912	0.000423866\\
0.914	0.00043282\\
0.916	0.000440405\\
0.918	0.000446606\\
0.92	0.000451416\\
0.922	0.000454831\\
0.924	0.000456852\\
0.926	0.000457484\\
0.928	0.000456739\\
0.93	0.000454631\\
0.932	0.000397012\\
0.934	0.000140346\\
0.936	-0.000113852\\
0.938	-0.000364708\\
0.94	-0.000611369\\
0.942	-0.000853006\\
0.944	-0.00108882\\
0.946	-0.00131803\\
0.948	-0.00153992\\
0.95	-0.00175376\\
0.952	-0.00205622\\
0.954	-0.0023496\\
0.956	-0.00263199\\
0.958	-0.00264199\\
0.96	-0.00256241\\
0.962	-0.00247536\\
0.964	-0.00238122\\
0.966	-0.0022804\\
0.968	-0.0021733\\
0.97	-0.00206063\\
0.972	-0.00197929\\
0.974	-0.00189485\\
0.976	-0.00182439\\
0.978	-0.00174902\\
0.98	-0.00166905\\
0.982	-0.00158188\\
0.984	-0.00139849\\
0.986	-0.0012121\\
0.988	-0.00102333\\
0.99	-0.00083281\\
0.992	-0.000641161\\
0.994	-0.000449006\\
0.996	-0.000256962\\
0.998	-6.5643e-05\\
1	0.000124349\\
1.002	0.000312419\\
1.004	0.00037141\\
1.006	0.000348555\\
1.008	0.000324656\\
1.01	0.00030008\\
1.012	0.000274966\\
1.014	0.000249397\\
1.016	0.000223455\\
1.018	0.000197223\\
1.02	0.000170784\\
1.022	9.07726e-05\\
1.024	-8.17963e-05\\
1.026	-0.000252469\\
1.028	-0.000420752\\
1.03	-0.000586159\\
1.032	-0.000748222\\
1.034	-0.000906482\\
1.036	-0.0010605\\
1.038	-0.00115799\\
1.04	-0.00113884\\
1.042	-0.00111678\\
1.044	-0.0010919\\
1.046	-0.0010643\\
1.048	-0.00103408\\
1.05	-0.00100137\\
1.052	-0.000966263\\
1.054	-0.000928889\\
1.056	-0.000889371\\
1.058	-0.000847838\\
1.06	-0.000804421\\
1.062	-0.000759257\\
1.064	-0.000783871\\
1.066	-0.000890714\\
1.068	-0.000990675\\
1.07	-0.00108353\\
1.072	-0.00116908\\
1.074	-0.00124716\\
1.076	-0.00131763\\
1.078	-0.00138036\\
1.08	-0.00143528\\
1.082	-0.00150732\\
1.084	-0.00145163\\
1.086	-0.00140633\\
1.088	-0.00136023\\
1.09	-0.00131069\\
1.092	-0.0012579\\
1.094	-0.00120206\\
1.096	-0.0011434\\
1.098	-0.00103996\\
1.1	-0.000893086\\
1.102	-0.00070318\\
1.104	-0.000512039\\
1.106	-0.000320333\\
1.108	-0.000128721\\
1.11	6.21471e-05\\
1.112	0.000204382\\
1.114	0.000215047\\
1.116	0.000225031\\
1.118	0.000234306\\
1.12	0.000242844\\
1.122	0.000250624\\
1.124	0.000257626\\
1.126	0.00026383\\
1.128	0.000269224\\
1.13	0.000273795\\
1.132	0.000277535\\
1.134	0.000280438\\
1.136	0.000282501\\
1.138	0.000283723\\
1.14	0.000284109\\
1.142	0.000232552\\
1.144	0.000170963\\
1.146	0.000107085\\
1.148	-3.68305e-05\\
1.15	-0.000178975\\
1.152	-0.000318872\\
1.154	-0.000456055\\
1.156	-0.000590075\\
1.158	-0.000720497\\
1.16	-0.000846904\\
1.162	-0.000968897\\
1.164	-0.0010861\\
1.166	-0.00119814\\
1.168	-0.00129257\\
1.17	-0.0013381\\
1.172	-0.00138046\\
1.174	-0.00135706\\
1.176	-0.00132943\\
1.178	-0.00129769\\
1.18	-0.00126202\\
1.182	-0.00122257\\
1.184	-0.00117953\\
1.186	-0.00113307\\
1.188	-0.00108341\\
1.19	-0.00103073\\
1.192	-0.000995862\\
1.194	-0.000960721\\
1.196	-0.000930431\\
1.198	-0.000903245\\
1.2	-0.000873489\\
1.202	-0.000841291\\
1.204	-0.000806784\\
1.206	-0.000770106\\
1.208	-0.000731401\\
1.21	-0.000666189\\
1.212	-0.000556391\\
1.214	-0.000445765\\
1.216	-0.000334672\\
1.218	-0.000223471\\
1.22	-0.000112519\\
1.222	-2.16785e-06\\
1.224	0.000107237\\
1.226	0.000215355\\
1.228	0.000229122\\
1.23	0.00021563\\
1.232	0.000201691\\
1.234	0.000187183\\
1.236	0.000172246\\
1.238	0.000157008\\
1.24	0.000141519\\
1.242	0.000125829\\
1.244	0.000109988\\
1.246	9.40463e-05\\
1.248	7.80545e-05\\
1.25	6.20624e-05\\
1.252	4.42951e-05\\
1.254	-5.91838e-05\\
1.256	-0.000161433\\
1.258	-0.000262156\\
1.26	-0.000361066\\
1.262	-0.000457883\\
1.264	-0.000552335\\
1.266	-0.00060688\\
1.268	-0.000612778\\
1.27	-0.000616916\\
1.272	-0.000619298\\
1.274	-0.000619936\\
1.276	-0.000618846\\
1.278	-0.000616045\\
1.28	-0.00061156\\
1.282	-0.000605417\\
1.284	-0.000597649\\
1.286	-0.000588291\\
1.288	-0.000577385\\
1.29	-0.000611334\\
1.292	-0.00064257\\
1.294	-0.000670129\\
1.296	-0.000693971\\
1.298	-0.00071407\\
1.3	-0.000730412\\
1.302	-0.000758124\\
1.304	-0.000735347\\
1.306	-0.000713601\\
1.308	-0.000697169\\
1.31	-0.000678845\\
1.312	-0.00065871\\
1.314	-0.000595907\\
1.316	-0.000520475\\
1.318	-0.000443709\\
1.32	-0.000365885\\
1.322	-0.00028728\\
1.324	-0.000208168\\
1.326	-0.000128822\\
1.328	-4.95097e-05\\
1.33	2.95064e-05\\
1.332	0.000107968\\
1.334	0.000130962\\
1.336	0.000135996\\
1.338	0.000140585\\
1.34	0.00014472\\
1.342	0.00014839\\
1.344	0.000151587\\
1.346	0.000154305\\
1.348	0.000137435\\
1.35	0.000106721\\
1.352	7.61682e-05\\
1.354	4.58872e-05\\
1.356	1.59855e-05\\
1.358	-1.34324e-05\\
1.36	-4.22648e-05\\
1.362	-7.82498e-05\\
1.364	-0.000158116\\
1.366	-0.000236507\\
1.368	-0.000313165\\
1.37	-0.00038784\\
1.372	-0.000460292\\
1.374	-0.00053029\\
1.376	-0.000597615\\
1.378	-0.000662061\\
1.38	-0.000697103\\
1.382	-0.000702983\\
1.384	-0.000704967\\
1.386	-0.0007031\\
1.388	-0.000704941\\
1.39	-0.000698566\\
1.392	-0.000689931\\
1.394	-0.000679094\\
1.396	-0.000666117\\
1.398	-0.000651072\\
1.4	-0.000634032\\
1.402	-0.00061508\\
1.404	-0.000594301\\
1.406	-0.000571786\\
1.408	-0.00054763\\
1.41	-0.000521935\\
1.412	-0.000505854\\
1.414	-0.000491343\\
1.416	-0.000475792\\
1.418	-0.000466554\\
1.42	-0.000455926\\
1.422	-0.000443962\\
1.424	-0.000430715\\
1.426	-0.000416243\\
1.428	-0.000400608\\
1.43	-0.000383874\\
1.432	-0.000366107\\
1.434	-0.000347377\\
1.436	-0.000290103\\
1.438	-0.000227033\\
1.44	-0.000163815\\
1.442	-0.000100655\\
1.444	-3.77524e-05\\
1.446	2.4693e-05\\
1.448	8.64864e-05\\
1.45	0.000137836\\
1.452	0.000130385\\
1.454	0.000122639\\
1.456	0.000114627\\
1.458	0.000106375\\
1.46	9.78288e-05\\
1.462	8.90402e-05\\
1.464	8.00995e-05\\
1.466	7.10361e-05\\
1.468	6.18797e-05\\
1.47	5.26596e-05\\
1.472	4.34053e-05\\
1.474	3.4146e-05\\
1.476	2.49109e-05\\
1.478	1.57285e-05\\
1.48	-2.39497e-05\\
1.482	-6.84752e-05\\
1.484	-0.000111539\\
1.486	-0.000153025\\
1.488	-0.000192821\\
1.49	-0.000230826\\
1.492	-0.000277729\\
1.494	-0.00029995\\
1.496	-0.00031083\\
1.498	-0.000320725\\
1.5	-0.000329616\\
1.502	-0.000337488\\
1.504	-0.000344328\\
1.506	-0.000350124\\
1.508	-0.000355567\\
1.51	-0.000367708\\
1.512	-0.000377948\\
1.514	-0.000386277\\
1.516	-0.000392695\\
1.518	-0.000397204\\
1.52	-0.000393213\\
1.522	-0.000381991\\
1.524	-0.0003726\\
1.526	-0.000362179\\
1.528	-0.000355734\\
1.53	-0.000331657\\
1.532	-0.000304643\\
1.534	-0.000276735\\
1.536	-0.000248039\\
1.538	-0.000218663\\
1.54	-0.000188713\\
1.542	-0.000158298\\
1.544	-0.000127524\\
1.546	-9.64984e-05\\
1.548	-6.53255e-05\\
1.55	-3.41096e-05\\
1.552	-2.9531e-06\\
1.554	2.80437e-05\\
1.556	5.51723e-05\\
1.558	4.00386e-05\\
1.56	2.50295e-05\\
1.562	1.01992e-05\\
1.564	-4.39912e-06\\
1.566	-1.87142e-05\\
1.568	-3.26962e-05\\
1.57	-3.82236e-05\\
1.572	-3.67019e-05\\
1.574	-3.51595e-05\\
1.576	-3.36012e-05\\
1.578	-5.32292e-05\\
1.58	-8.34646e-05\\
1.582	-0.000113251\\
1.584	-0.000142491\\
1.586	-0.000171094\\
1.588	-0.00019897\\
1.59	-0.000226034\\
1.592	-0.000252206\\
1.594	-0.000277409\\
1.596	-0.000301572\\
1.598	-0.000324628\\
1.6	-0.000346514\\
1.602	-0.000367175\\
1.604	-0.000377856\\
1.606	-0.000355646\\
1.608	-0.00034959\\
1.61	-0.000346437\\
1.612	-0.000342148\\
1.614	-0.000336751\\
1.616	-0.000330278\\
1.618	-0.000322763\\
1.62	-0.000314244\\
1.622	-0.000304762\\
1.624	-0.000294359\\
1.626	-0.000283082\\
1.628	-0.000270977\\
1.63	-0.000258096\\
1.632	-0.000250324\\
1.634	-0.00024436\\
1.636	-0.000237655\\
1.638	-0.000232886\\
1.64	-0.000229554\\
1.642	-0.00022551\\
1.644	-0.000220775\\
1.646	-0.000215374\\
1.648	-0.000209331\\
1.65	-0.000202675\\
1.652	-0.000195433\\
1.654	-0.000187636\\
1.656	-0.000179316\\
1.658	-0.000170504\\
1.66	-0.000135905\\
1.662	-0.000100347\\
1.664	-6.47934e-05\\
1.666	-2.93563e-05\\
1.668	5.85162e-06\\
1.67	4.07198e-05\\
1.672	5.28669e-05\\
1.674	5.73668e-05\\
1.676	6.14534e-05\\
1.678	6.51193e-05\\
1.68	6.83591e-05\\
1.682	6.68786e-05\\
1.684	6.19726e-05\\
1.686	5.69117e-05\\
1.688	5.17128e-05\\
1.69	4.64348e-05\\
1.692	4.10947e-05\\
1.694	3.19406e-05\\
1.696	1.35077e-05\\
1.698	-4.41323e-06\\
1.7	-2.17713e-05\\
1.702	-3.85184e-05\\
1.704	-5.46093e-05\\
1.706	-7.00017e-05\\
1.708	-8.46568e-05\\
1.71	-9.85384e-05\\
1.712	-0.000111614\\
1.714	-0.000123855\\
1.716	-0.000140458\\
1.718	-0.000171054\\
1.72	-0.000200467\\
1.722	-0.000207882\\
1.724	-0.000213995\\
1.726	-0.000219084\\
1.728	-0.000223146\\
1.73	-0.000226179\\
1.732	-0.000228184\\
1.734	-0.000225592\\
1.736	-0.000217449\\
1.738	-0.000208718\\
1.74	-0.000199435\\
1.742	-0.000193895\\
1.744	-0.000189995\\
1.746	-0.000179399\\
1.748	-0.000168764\\
1.75	-0.000160115\\
1.752	-0.00015095\\
1.754	-0.000141308\\
1.756	-0.000131225\\
1.758	-0.000120742\\
1.76	-0.000109896\\
1.762	-9.87282e-05\\
1.764	-8.72782e-05\\
1.766	-7.55862e-05\\
1.768	-6.36925e-05\\
1.77	-5.16373e-05\\
1.772	-3.94607e-05\\
1.774	-2.72025e-05\\
1.776	-2.74904e-05\\
1.778	-3.44025e-05\\
1.78	-4.10641e-05\\
1.782	-4.74537e-05\\
1.784	-4.14464e-05\\
1.786	-3.38583e-05\\
1.788	-2.6281e-05\\
1.79	-1.87383e-05\\
1.792	-1.40295e-05\\
1.794	-2.57352e-05\\
1.796	-3.73002e-05\\
1.798	-4.86853e-05\\
1.8	-5.98528e-05\\
1.802	-7.07667e-05\\
1.804	-8.13921e-05\\
1.806	-9.1696e-05\\
1.808	-0.000101647\\
1.81	-0.000111216\\
1.812	-0.000120375\\
1.814	-0.000129098\\
1.816	-0.000137361\\
1.818	-0.000145144\\
1.82	-0.000152425\\
1.822	-0.000159189\\
1.824	-0.000165419\\
1.826	-0.000178756\\
1.828	-0.000177738\\
1.83	-0.00017613\\
1.832	-0.000173946\\
1.834	-0.0001712\\
1.836	-0.000167907\\
1.838	-0.000164087\\
1.84	-0.000159756\\
1.842	-0.000154938\\
1.844	-0.000149652\\
1.846	-0.000143923\\
1.848	-0.000137774\\
1.85	-0.000131231\\
1.852	-0.00012746\\
1.854	-0.000125081\\
1.856	-0.000122317\\
1.858	-0.00011943\\
1.86	-0.000118732\\
1.862	-0.000117656\\
1.864	-0.000116208\\
1.866	-0.000114399\\
1.868	-0.000112238\\
1.87	-0.000109737\\
1.872	-0.000106907\\
1.874	-0.000103761\\
1.876	-0.000100314\\
1.878	-9.65789e-05\\
1.88	-9.25722e-05\\
1.882	-8.48414e-05\\
1.884	-6.49762e-05\\
1.886	-4.50927e-05\\
1.888	-2.5255e-05\\
1.89	-5.5259e-06\\
1.892	1.40326e-05\\
1.894	2.91843e-05\\
1.896	3.6564e-05\\
1.898	4.39214e-05\\
1.9	4.63593e-05\\
1.902	4.31877e-05\\
1.904	3.5868e-05\\
1.906	2.86817e-05\\
1.908	2.16516e-05\\
1.91	1.47993e-05\\
1.912	8.14524e-06\\
1.914	1.70909e-06\\
1.916	-4.4908e-06\\
1.918	-1.04372e-05\\
1.92	-1.61143e-05\\
1.922	-2.15071e-05\\
1.924	-2.66025e-05\\
1.926	-3.13882e-05\\
1.928	-3.58534e-05\\
1.93	-3.99889e-05\\
1.932	-4.37865e-05\\
1.934	-4.83183e-05\\
1.936	-6.30653e-05\\
1.938	-7.72986e-05\\
1.94	-9.0979e-05\\
1.942	-0.00010407\\
1.944	-0.000116536\\
1.946	-0.000127226\\
1.948	-0.000124443\\
1.95	-0.000121299\\
1.952	-0.000117807\\
1.954	-0.000113982\\
1.956	-0.000109843\\
1.958	-0.000105405\\
1.96	-0.000100687\\
1.962	-9.80356e-05\\
1.964	-9.56288e-05\\
1.966	-9.18945e-05\\
1.968	-8.78533e-05\\
1.97	-8.53767e-05\\
1.972	-8.28782e-05\\
1.974	-8.0092e-05\\
1.976	-7.70304e-05\\
1.978	-7.45243e-05\\
1.98	-7.49538e-05\\
1.982	-7.5233e-05\\
1.984	-7.5361e-05\\
1.986	-7.53371e-05\\
1.988	-7.51608e-05\\
1.99	-7.48317e-05\\
1.992	-7.43499e-05\\
1.994	-7.37154e-05\\
1.996	-7.29288e-05\\
1.998	-7.19907e-05\\
2	-7.03277e-05\\
};
\addplot [color=mycolor4, line width=1.0pt, forget plot]
  table[row sep=crcr]{%
-0.024	0\\
-0.022	0\\
-0.02	0\\
-0.018	0\\
-0.016	0\\
-0.014	0\\
-0.012	0\\
-0.01	0\\
-0.008	0\\
-0.006	0\\
-0.004	0\\
-0.002	0\\
0	0\\
0.002	-2.93741e-09\\
0.004	-4.49626e-08\\
0.006	-2.14278e-07\\
0.008	-6.37046e-07\\
0.01	-1.46344e-06\\
0.012	-2.85657e-06\\
0.014	-4.98378e-06\\
0.016	-8.00977e-06\\
0.018	-8.70268e-06\\
0.02	-1.26244e-05\\
0.022	-1.76016e-05\\
0.024	-7.23053e-06\\
0.026	-9.83571e-06\\
0.028	-1.30714e-05\\
0.03	-1.7025e-05\\
0.032	-2.17874e-05\\
0.034	-2.74514e-05\\
0.036	-3.4112e-05\\
0.038	-4.18645e-05\\
0.04	-5.08042e-05\\
0.042	-6.1025e-05\\
0.044	-7.26189e-05\\
0.046	-8.56746e-05\\
0.048	-0.000100276\\
0.05	-0.000116504\\
0.052	-0.00013443\\
0.054	-0.00015412\\
0.056	-0.000175632\\
0.058	-0.000199016\\
0.06	-0.000224311\\
0.062	-0.000251545\\
0.064	-0.000171883\\
0.066	-0.000192986\\
0.068	-0.000215743\\
0.07	-0.00024019\\
0.072	-0.000266359\\
0.074	-0.000294272\\
0.076	-0.000323944\\
0.078	-0.000355384\\
0.08	-0.000388591\\
0.082	-0.000423555\\
0.084	-0.00103666\\
0.086	-0.00169717\\
0.088	-0.00237004\\
0.09	-0.00305232\\
0.092	-0.00374401\\
0.094	-0.004532\\
0.096	-0.00532748\\
0.098	-0.00627955\\
0.1	-0.0089624\\
0.102	-0.0116047\\
0.104	-0.0141949\\
0.106	-0.0167219\\
0.108	-0.0191751\\
0.11	-0.0215443\\
0.112	-0.0196757\\
0.114	-0.0173928\\
0.116	-0.0150511\\
0.118	-0.0126578\\
0.12	-0.0102204\\
0.122	-0.00774661\\
0.124	-0.00524434\\
0.126	-0.00325287\\
0.128	-0.00257056\\
0.13	-0.00190422\\
0.132	-0.00179963\\
0.134	-0.00170802\\
0.136	-0.00162911\\
0.138	-0.00156249\\
0.14	-0.00150762\\
0.142	-0.00146388\\
0.144	-0.0014305\\
0.146	-0.00140704\\
0.148	-0.00146273\\
0.15	-0.00151831\\
0.152	-0.00157361\\
0.154	-0.00162848\\
0.156	-0.00168275\\
0.158	-0.00173623\\
0.16	-0.00178876\\
0.162	-0.00184014\\
0.164	-0.0018902\\
0.166	-0.00181529\\
0.168	-0.0017828\\
0.17	-0.00174692\\
0.172	-0.00170776\\
0.174	-0.00166542\\
0.176	-0.00162003\\
0.178	-0.00157172\\
0.18	-0.00152062\\
0.182	-0.00146688\\
0.184	-0.00141066\\
0.186	-0.00135213\\
0.188	-0.00129146\\
0.19	-0.00122882\\
0.192	-0.00116439\\
0.194	-0.00109837\\
0.196	-0.00103093\\
0.198	-0.000962284\\
0.2	-0.000892617\\
0.202	-0.000822126\\
0.204	-0.000842182\\
0.206	-0.000747091\\
0.208	-0.000651177\\
0.21	-0.000554741\\
0.212	-0.000458084\\
0.214	-0.000361511\\
0.216	-0.000265325\\
0.218	-0.000169827\\
0.22	-7.53156e-05\\
0.222	1.7916e-05\\
0.224	0.000109579\\
0.226	0.000199392\\
0.228	0.000287081\\
0.23	0.000372382\\
0.232	0.00045504\\
0.234	0.000534811\\
0.236	0.000611466\\
0.238	0.000684785\\
0.24	0.000754565\\
0.242	0.000820617\\
0.244	0.000882766\\
0.246	0.000940854\\
0.248	0.000994741\\
0.25	0.0010443\\
0.252	0.00108943\\
0.254	0.000831094\\
0.256	0.000918292\\
0.258	0.00100265\\
0.26	0.00108394\\
0.262	0.00116193\\
0.264	0.0012364\\
0.266	0.00130714\\
0.268	0.00137396\\
0.27	0.00143667\\
0.272	0.00149509\\
0.274	0.00154907\\
0.276	0.00159844\\
0.278	0.00164309\\
0.28	0.00168288\\
0.282	0.00171771\\
0.284	0.00174748\\
0.286	0.00177213\\
0.288	0.0017916\\
0.29	0.00180582\\
0.292	0.00181479\\
0.294	0.00146142\\
0.296	0.000513312\\
0.298	-0.000435043\\
0.3	-0.00138023\\
0.302	-0.00258995\\
0.304	-0.00394013\\
0.306	-0.00525686\\
0.308	-0.00653582\\
0.31	-0.00777288\\
0.312	-0.00896418\\
0.314	-0.0101061\\
0.316	-0.0111952\\
0.318	-0.0122284\\
0.32	-0.0132029\\
0.322	-0.014116\\
0.324	-0.0149655\\
0.326	-0.0157492\\
0.328	-0.015932\\
0.33	-0.0143704\\
0.332	-0.0127744\\
0.334	-0.0111499\\
0.336	-0.00950296\\
0.338	-0.00809057\\
0.34	-0.00780161\\
0.342	-0.00749952\\
0.344	-0.00718495\\
0.346	-0.00685858\\
0.348	-0.00652107\\
0.35	-0.00617308\\
0.352	-0.00581526\\
0.354	-0.00544827\\
0.356	-0.00545302\\
0.358	-0.00583317\\
0.36	-0.0061929\\
0.362	-0.00653105\\
0.364	-0.00684652\\
0.366	-0.00713829\\
0.368	-0.00740544\\
0.37	-0.0076471\\
0.372	-0.00786254\\
0.374	-0.00805108\\
0.376	-0.00772813\\
0.378	-0.00686218\\
0.38	-0.00602918\\
0.382	-0.00523106\\
0.384	-0.00446935\\
0.386	-0.00361318\\
0.388	-0.0023445\\
0.39	-0.00135185\\
0.392	-0.00109386\\
0.394	-0.000893989\\
0.396	-0.000751647\\
0.398	-0.000666006\\
0.4	-0.000636008\\
0.402	-0.000660372\\
0.404	-0.000737603\\
0.406	-0.00135977\\
0.408	-0.0022344\\
0.41	-0.00309877\\
0.412	-0.00394992\\
0.414	-0.00478493\\
0.416	-0.00560097\\
0.418	-0.00639527\\
0.42	-0.00718261\\
0.422	-0.00757667\\
0.424	-0.00755034\\
0.426	-0.00744353\\
0.428	-0.007258\\
0.43	-0.00699573\\
0.432	-0.00665895\\
0.434	-0.00672932\\
0.436	-0.00691044\\
0.438	-0.00705747\\
0.44	-0.00717057\\
0.442	-0.00725001\\
0.444	-0.00729616\\
0.446	-0.00730942\\
0.448	-0.00729032\\
0.45	-0.00723941\\
0.452	-0.00715736\\
0.454	-0.00704486\\
0.456	-0.00690271\\
0.458	-0.00673173\\
0.46	-0.00653283\\
0.462	-0.00521819\\
0.464	-0.00381946\\
0.466	-0.00244263\\
0.468	-0.00109202\\
0.47	0.000228214\\
0.472	0.000842309\\
0.474	0.000921311\\
0.476	0.000998146\\
0.478	0.00107257\\
0.48	0.00114435\\
0.482	0.00121327\\
0.484	0.00127909\\
0.486	0.00134163\\
0.488	0.00140068\\
0.49	0.00145607\\
0.492	0.00150762\\
0.494	0.00155518\\
0.496	0.0015986\\
0.498	0.00163777\\
0.5	0.00167256\\
0.502	0.00170287\\
0.504	0.00172862\\
0.506	0.00174974\\
0.508	0.00150621\\
0.51	0.000732947\\
0.512	-3.70127e-05\\
0.514	-0.000801019\\
0.516	-0.00155646\\
0.518	-0.00230078\\
0.52	-0.00303145\\
0.522	-0.00374605\\
0.524	-0.00444219\\
0.526	-0.00511757\\
0.528	-0.00577\\
0.53	-0.00641384\\
0.532	-0.00710116\\
0.534	-0.0077972\\
0.536	-0.00846109\\
0.538	-0.00909076\\
0.54	-0.00968427\\
0.542	-0.0102398\\
0.544	-0.0102895\\
0.546	-0.00982881\\
0.548	-0.00929869\\
0.55	-0.0086759\\
0.552	-0.00803282\\
0.554	-0.00737167\\
0.556	-0.00669468\\
0.558	-0.0061787\\
0.56	-0.00631086\\
0.562	-0.00650117\\
0.564	-0.00667208\\
0.566	-0.00682317\\
0.568	-0.00688467\\
0.57	-0.00689182\\
0.572	-0.00686373\\
0.574	-0.00680098\\
0.576	-0.00670427\\
0.578	-0.00657434\\
0.58	-0.00641205\\
0.582	-0.00621833\\
0.584	-0.00599418\\
0.586	-0.00568422\\
0.588	-0.0049428\\
0.59	-0.0042186\\
0.592	-0.00351353\\
0.594	-0.00282939\\
0.596	-0.00216787\\
0.598	-0.00153053\\
0.6	-0.000918822\\
0.602	-0.000472503\\
0.604	-0.000221956\\
0.606	4.27854e-06\\
0.608	0.00020601\\
0.61	0.000383155\\
0.612	0.000535732\\
0.614	0.000663858\\
0.616	0.000533395\\
0.618	-0.000105362\\
0.62	-0.000739646\\
0.622	-0.00136729\\
0.624	-0.00198615\\
0.626	-0.00259414\\
0.628	-0.00318925\\
0.63	-0.00376948\\
0.632	-0.00407486\\
0.634	-0.00417771\\
0.636	-0.00423838\\
0.638	-0.00425738\\
0.64	-0.00424857\\
0.642	-0.00425431\\
0.644	-0.00425676\\
0.646	-0.00422216\\
0.648	-0.00414877\\
0.65	-0.00403748\\
0.652	-0.0038893\\
0.654	-0.00370535\\
0.656	-0.00348684\\
0.658	-0.00323508\\
0.66	-0.00295146\\
0.662	-0.00277311\\
0.664	-0.00270714\\
0.666	-0.00262365\\
0.668	-0.00252308\\
0.67	-0.00240591\\
0.672	-0.00227265\\
0.674	-0.00212385\\
0.676	-0.00142723\\
0.678	-0.000741954\\
0.68	-0.000415913\\
0.682	-8.84213e-05\\
0.684	0.000165674\\
0.686	0.000228252\\
0.688	0.000290073\\
0.69	0.00035094\\
0.692	0.000410656\\
0.694	0.000469032\\
0.696	0.000525884\\
0.698	0.000581036\\
0.7	0.000634315\\
0.702	0.000685561\\
0.704	0.000734617\\
0.706	0.000781338\\
0.708	0.000825585\\
0.71	0.00086723\\
0.712	0.000906155\\
0.714	0.000942251\\
0.716	0.000975419\\
0.718	0.00100557\\
0.72	0.00103263\\
0.722	0.00100091\\
0.724	0.000505996\\
0.726	-1.02436e-05\\
0.728	-0.000521659\\
0.73	-0.00102652\\
0.732	-0.00152315\\
0.734	-0.0020099\\
0.736	-0.00248518\\
0.738	-0.00294743\\
0.74	-0.00339519\\
0.742	-0.00382701\\
0.744	-0.00424157\\
0.746	-0.00463756\\
0.748	-0.00501379\\
0.75	-0.00537747\\
0.752	-0.00574074\\
0.754	-0.00569509\\
0.756	-0.00563168\\
0.758	-0.00554044\\
0.76	-0.00542206\\
0.762	-0.00527732\\
0.764	-0.00510707\\
0.766	-0.00491223\\
0.768	-0.00469378\\
0.77	-0.00445278\\
0.772	-0.00419034\\
0.774	-0.00390762\\
0.776	-0.00360584\\
0.778	-0.00328626\\
0.78	-0.00306031\\
0.782	-0.00305623\\
0.784	-0.00303397\\
0.786	-0.00299386\\
0.788	-0.00293631\\
0.79	-0.00286176\\
0.792	-0.00297899\\
0.794	-0.00310205\\
0.796	-0.00319481\\
0.798	-0.00305717\\
0.8	-0.00292205\\
0.802	-0.00278984\\
0.804	-0.00266088\\
0.806	-0.00253553\\
0.808	-0.00241407\\
0.81	-0.00229679\\
0.812	-0.00218392\\
0.814	-0.00207569\\
0.816	-0.00197227\\
0.818	-0.00187383\\
0.82	-0.00178048\\
0.822	-0.00169232\\
0.824	-0.00160942\\
0.826	-0.00153181\\
0.828	-0.00145952\\
0.83	-0.00139251\\
0.832	-0.00133075\\
0.834	-0.00127418\\
0.836	-0.0012227\\
0.838	-0.00117621\\
0.84	-0.00152779\\
0.842	-0.00186756\\
0.844	-0.00202741\\
0.846	-0.00216552\\
0.848	-0.0022818\\
0.85	-0.0023762\\
0.852	-0.00244877\\
0.854	-0.00249962\\
0.856	-0.00252893\\
0.858	-0.00253695\\
0.86	-0.00252743\\
0.862	-0.00253991\\
0.864	-0.00254197\\
0.866	-0.00253988\\
0.868	-0.00251757\\
0.87	-0.00247545\\
0.872	-0.00241396\\
0.874	-0.00233363\\
0.876	-0.00223502\\
0.878	-0.00211875\\
0.88	-0.00198549\\
0.882	-0.00183596\\
0.884	-0.00167091\\
0.886	-0.00149116\\
0.888	-0.00110365\\
0.89	-0.00082156\\
0.892	-0.00067461\\
0.894	-0.000530539\\
0.896	-0.000389873\\
0.898	-0.000253109\\
0.9	-0.000120711\\
0.902	6.88948e-06\\
0.904	6.68469e-05\\
0.906	0.000105851\\
0.908	0.000144361\\
0.91	0.000182256\\
0.912	0.000219419\\
0.914	0.000255733\\
0.916	0.000291088\\
0.918	0.000325378\\
0.92	0.000358501\\
0.922	0.00039036\\
0.924	0.000420864\\
0.926	0.000449925\\
0.928	0.000477464\\
0.93	0.000503406\\
0.932	0.000473512\\
0.934	0.000244162\\
0.936	1.67851e-05\\
0.938	-0.000207826\\
0.94	-0.000428899\\
0.942	-0.000645681\\
0.944	-0.000857446\\
0.946	-0.00106349\\
0.948	-0.00126315\\
0.95	-0.00145577\\
0.952	-0.00173808\\
0.954	-0.00201244\\
0.956	-0.00227698\\
0.958	-0.00253086\\
0.96	-0.00277333\\
0.962	-0.00300366\\
0.964	-0.00306721\\
0.966	-0.00306342\\
0.968	-0.00304406\\
0.97	-0.00300971\\
0.972	-0.00299714\\
0.974	-0.00297185\\
0.976	-0.00295086\\
0.978	-0.00291527\\
0.98	-0.00286541\\
0.982	-0.00280164\\
0.984	-0.00272435\\
0.986	-0.00263401\\
0.988	-0.00253108\\
0.99	-0.00241608\\
0.992	-0.00228956\\
0.994	-0.0021521\\
0.996	-0.0020043\\
0.998	-0.00184679\\
1	-0.00168023\\
1.002	-0.00150531\\
1.004	-0.00144929\\
1.006	-0.00145536\\
1.008	-0.00134988\\
1.01	-0.0012456\\
1.012	-0.00114278\\
1.014	-0.00104171\\
1.016	-0.000942685\\
1.018	-0.000845992\\
1.02	-0.000751885\\
1.022	-0.000714058\\
1.024	-0.000771806\\
1.026	-0.000830967\\
1.028	-0.00089132\\
1.03	-0.000952635\\
1.032	-0.00101468\\
1.034	-0.0010772\\
1.036	-0.00113997\\
1.038	-0.00120272\\
1.04	-0.0012652\\
1.042	-0.00132717\\
1.044	-0.00138836\\
1.046	-0.00144851\\
1.048	-0.00150739\\
1.05	-0.00156472\\
1.052	-0.00162028\\
1.054	-0.00162358\\
1.056	-0.00156891\\
1.058	-0.00150681\\
1.06	-0.00143757\\
1.062	-0.00136149\\
1.064	-0.00135028\\
1.066	-0.0014166\\
1.068	-0.00147154\\
1.07	-0.00151508\\
1.072	-0.00154728\\
1.074	-0.00156819\\
1.076	-0.00157791\\
1.078	-0.0015766\\
1.08	-0.00156441\\
1.082	-0.00156936\\
1.084	-0.00156519\\
1.086	-0.0015645\\
1.088	-0.00155611\\
1.09	-0.00153734\\
1.092	-0.00150843\\
1.094	-0.00146962\\
1.096	-0.00142119\\
1.098	-0.00136345\\
1.1	-0.00119206\\
1.102	-0.000997364\\
1.104	-0.000800468\\
1.106	-0.000602065\\
1.108	-0.00040285\\
1.11	-0.000203506\\
1.112	-5.19606e-05\\
1.114	-3.119e-05\\
1.116	-1.03477e-05\\
1.118	1.04937e-05\\
1.12	3.12623e-05\\
1.122	5.18873e-05\\
1.124	7.22991e-05\\
1.126	9.24298e-05\\
1.128	0.000112213\\
1.13	0.000131585\\
1.132	0.000150482\\
1.134	0.000168846\\
1.136	0.000186619\\
1.138	0.000203746\\
1.14	0.000220175\\
1.142	0.000184747\\
1.144	0.000139316\\
1.146	9.15738e-05\\
1.148	-3.62835e-05\\
1.15	-0.000162499\\
1.152	-0.000286646\\
1.154	-0.000408311\\
1.156	-0.000527091\\
1.158	-0.000642601\\
1.16	-0.000754469\\
1.162	-0.000862339\\
1.164	-0.000965874\\
1.166	-0.00106476\\
1.168	-0.00115869\\
1.17	-0.0012942\\
1.172	-0.00144574\\
1.174	-0.00150746\\
1.176	-0.00154969\\
1.178	-0.00158372\\
1.18	-0.00160955\\
1.182	-0.0016272\\
1.184	-0.00163672\\
1.186	-0.0016382\\
1.188	-0.00163174\\
1.19	-0.00161746\\
1.192	-0.00161613\\
1.194	-0.00160963\\
1.196	-0.00160306\\
1.198	-0.00159465\\
1.2	-0.00157876\\
1.202	-0.00155552\\
1.204	-0.00152512\\
1.206	-0.00148602\\
1.208	-0.0014369\\
1.21	-0.00138131\\
1.212	-0.0013195\\
1.214	-0.00125178\\
1.216	-0.00117305\\
1.218	-0.00103806\\
1.22	-0.000903003\\
1.222	-0.000768321\\
1.224	-0.000634439\\
1.226	-0.000501773\\
1.228	-0.000463463\\
1.23	-0.000452485\\
1.232	-0.000442097\\
1.234	-0.000432483\\
1.236	-0.000423564\\
1.238	-0.000415269\\
1.24	-0.0004076\\
1.242	-0.000400556\\
1.244	-0.000394132\\
1.246	-0.000388317\\
1.248	-0.000383099\\
1.25	-0.000378461\\
1.252	-0.000376206\\
1.254	-0.000460296\\
1.256	-0.000543811\\
1.258	-0.000626474\\
1.26	-0.000708011\\
1.262	-0.000788152\\
1.264	-0.000866634\\
1.266	-0.000908116\\
1.268	-0.000929505\\
1.27	-0.000946239\\
1.272	-0.000958335\\
1.274	-0.000965823\\
1.276	-0.000968748\\
1.278	-0.000967165\\
1.28	-0.000961143\\
1.282	-0.000950764\\
1.284	-0.000936117\\
1.286	-0.000917308\\
1.288	-0.00089445\\
1.29	-0.000914029\\
1.292	-0.000928563\\
1.294	-0.000937182\\
1.296	-0.000939943\\
1.298	-0.000936925\\
1.3	-0.000928223\\
1.302	-0.000929075\\
1.304	-0.000927296\\
1.306	-0.000923125\\
1.308	-0.000920674\\
1.31	-0.000906646\\
1.312	-0.000834745\\
1.314	-0.000760379\\
1.316	-0.000683834\\
1.318	-0.000605398\\
1.32	-0.000525363\\
1.322	-0.000444018\\
1.324	-0.000361654\\
1.326	-0.00027856\\
1.328	-0.000195022\\
1.33	-0.000111323\\
1.332	-2.77439e-05\\
1.334	7.8169e-07\\
1.336	1.17353e-05\\
1.338	2.26102e-05\\
1.34	3.33694e-05\\
1.342	4.39768e-05\\
1.344	5.43971e-05\\
1.346	6.45963e-05\\
1.348	5.54376e-05\\
1.35	3.26351e-05\\
1.352	1.01663e-05\\
1.354	-1.18882e-05\\
1.356	-3.34501e-05\\
1.358	-5.44443e-05\\
1.36	-7.47985e-05\\
1.362	-0.00010228\\
1.364	-0.000173646\\
1.366	-0.00024357\\
1.368	-0.000311821\\
1.37	-0.000378176\\
1.372	-0.000442422\\
1.374	-0.000504356\\
1.376	-0.000563784\\
1.378	-0.000620524\\
1.38	-0.000674405\\
1.382	-0.000715818\\
1.384	-0.000718417\\
1.386	-0.000717101\\
1.388	-0.0007212\\
1.39	-0.00075523\\
1.392	-0.000785008\\
1.394	-0.000810494\\
1.396	-0.000831662\\
1.398	-0.000848501\\
1.4	-0.000861013\\
1.402	-0.000869215\\
1.404	-0.000873139\\
1.406	-0.00087283\\
1.408	-0.000868344\\
1.41	-0.000859754\\
1.412	-0.000858194\\
1.414	-0.000855609\\
1.416	-0.000849387\\
1.418	-0.000846886\\
1.42	-0.000840421\\
1.422	-0.000830065\\
1.424	-0.000815906\\
1.426	-0.000798042\\
1.428	-0.00077658\\
1.43	-0.000737004\\
1.432	-0.000682822\\
1.434	-0.000627884\\
1.436	-0.000572374\\
1.438	-0.000516475\\
1.44	-0.000460367\\
1.442	-0.000404229\\
1.444	-0.000348238\\
1.446	-0.000292565\\
1.448	-0.000237381\\
1.45	-0.00019245\\
1.452	-0.000206103\\
1.454	-0.000219807\\
1.456	-0.000233508\\
1.458	-0.000247152\\
1.46	-0.00026077\\
1.462	-0.000274283\\
1.464	-0.000287577\\
1.466	-0.000300599\\
1.468	-0.000313296\\
1.47	-0.000325616\\
1.472	-0.00033751\\
1.474	-0.000348928\\
1.476	-0.000350694\\
1.478	-0.000335795\\
1.48	-0.000349709\\
1.482	-0.000366893\\
1.484	-0.000381148\\
1.486	-0.000392469\\
1.488	-0.000400863\\
1.49	-0.000406344\\
1.492	-0.000419723\\
1.494	-0.000439442\\
1.496	-0.000456836\\
1.498	-0.000471886\\
1.5	-0.00048458\\
1.502	-0.000494914\\
1.504	-0.00050289\\
1.506	-0.000508516\\
1.508	-0.000512505\\
1.51	-0.000521935\\
1.512	-0.000528236\\
1.514	-0.000531433\\
1.516	-0.00053156\\
1.518	-0.000528661\\
1.52	-0.00052279\\
1.522	-0.000522027\\
1.524	-0.000495859\\
1.526	-0.000466315\\
1.528	-0.000440278\\
1.53	-0.000414123\\
1.532	-0.000386674\\
1.534	-0.000358041\\
1.536	-0.000328335\\
1.538	-0.000297672\\
1.54	-0.000266166\\
1.542	-0.000233934\\
1.544	-0.000201091\\
1.546	-0.000167755\\
1.548	-0.000134041\\
1.55	-0.000100064\\
1.552	-6.59392e-05\\
1.554	-3.17793e-05\\
1.556	-1.30582e-06\\
1.558	-1.29263e-05\\
1.56	-2.42679e-05\\
1.562	-3.52906e-05\\
1.564	-4.59559e-05\\
1.566	-5.62273e-05\\
1.568	-6.60702e-05\\
1.57	-7.54519e-05\\
1.572	-8.43418e-05\\
1.574	-9.27118e-05\\
1.576	-0.000100536\\
1.578	-0.000128987\\
1.58	-0.00016746\\
1.582	-0.000204872\\
1.584	-0.000241103\\
1.586	-0.000276039\\
1.588	-0.000309571\\
1.59	-0.000341599\\
1.592	-0.000364015\\
1.594	-0.000371333\\
1.596	-0.000376529\\
1.598	-0.000379617\\
1.6	-0.000380613\\
1.602	-0.000379545\\
1.604	-0.000376444\\
1.606	-0.00037135\\
1.608	-0.000386542\\
1.61	-0.000403577\\
1.612	-0.000418366\\
1.614	-0.00043089\\
1.616	-0.000441138\\
1.618	-0.000449108\\
1.62	-0.000454804\\
1.622	-0.000458237\\
1.624	-0.000459428\\
1.626	-0.000458402\\
1.628	-0.000455193\\
1.63	-0.000449841\\
1.632	-0.000448226\\
1.634	-0.000447047\\
1.636	-0.000443758\\
1.638	-0.000441046\\
1.64	-0.000438424\\
1.642	-0.00043376\\
1.644	-0.000418746\\
1.646	-0.000399437\\
1.648	-0.000379429\\
1.65	-0.000358794\\
1.652	-0.000337602\\
1.654	-0.000315928\\
1.656	-0.000293844\\
1.658	-0.000271424\\
1.66	-0.00024874\\
1.662	-0.000225867\\
1.664	-0.000202876\\
1.666	-0.000179838\\
1.668	-0.000156825\\
1.67	-0.000133906\\
1.672	-0.000133422\\
1.674	-0.000140261\\
1.676	-0.000147148\\
1.678	-0.000154057\\
1.68	-0.000160956\\
1.682	-0.000172106\\
1.684	-0.000186179\\
1.686	-0.000192217\\
1.688	-0.000194154\\
1.69	-0.000194953\\
1.692	-0.00019463\\
1.694	-0.000196972\\
1.696	-0.000207482\\
1.698	-0.00021641\\
1.7	-0.00022375\\
1.702	-0.000229501\\
1.704	-0.000233667\\
1.706	-0.000236256\\
1.708	-0.000237285\\
1.71	-0.000236772\\
1.712	-0.000234742\\
1.714	-0.000231225\\
1.716	-0.000231478\\
1.718	-0.000245193\\
1.72	-0.000257256\\
1.722	-0.000267651\\
1.724	-0.000276368\\
1.726	-0.000283402\\
1.728	-0.000288754\\
1.73	-0.000292429\\
1.732	-0.000294441\\
1.734	-0.000294806\\
1.736	-0.000285772\\
1.738	-0.000271313\\
1.74	-0.000255994\\
1.742	-0.000244138\\
1.744	-0.000233667\\
1.746	-0.000222435\\
1.748	-0.000211648\\
1.75	-0.000202694\\
1.752	-0.000193074\\
1.754	-0.000182831\\
1.756	-0.000172005\\
1.758	-0.00016064\\
1.76	-0.00014878\\
1.762	-0.000136469\\
1.764	-0.000123753\\
1.766	-0.000110678\\
1.768	-9.72908e-05\\
1.77	-8.36372e-05\\
1.772	-6.97644e-05\\
1.774	-5.5719e-05\\
1.776	-5.41361e-05\\
1.778	-5.9101e-05\\
1.78	-6.37466e-05\\
1.782	-6.80591e-05\\
1.784	-7.20259e-05\\
1.786	-7.5636e-05\\
1.788	-7.88796e-05\\
1.79	-8.17485e-05\\
1.792	-8.70114e-05\\
1.794	-0.00010822\\
1.796	-0.000128791\\
1.798	-0.000148661\\
1.8	-0.000167767\\
1.802	-0.000180533\\
1.804	-0.00018782\\
1.806	-0.000193988\\
1.808	-0.000199034\\
1.81	-0.000202954\\
1.812	-0.000205752\\
1.814	-0.000207433\\
1.816	-0.000208008\\
1.818	-0.00020749\\
1.82	-0.000205895\\
1.822	-0.000203245\\
1.824	-0.000199563\\
1.826	-0.000204761\\
1.828	-0.000214334\\
1.83	-0.000222729\\
1.832	-0.000229933\\
1.834	-0.000235939\\
1.836	-0.000240743\\
1.838	-0.000244345\\
1.84	-0.00024675\\
1.842	-0.000247966\\
1.844	-0.000248004\\
1.846	-0.000246881\\
1.848	-0.000244616\\
1.85	-0.000241232\\
1.852	-0.000239893\\
1.854	-0.000239224\\
1.856	-0.000237451\\
1.858	-0.000230001\\
1.86	-0.000224588\\
1.862	-0.000218685\\
1.864	-0.000212315\\
1.866	-0.000205504\\
1.868	-0.000198277\\
1.87	-0.000190661\\
1.872	-0.000182684\\
1.874	-0.000174374\\
1.876	-0.000165759\\
1.878	-0.000156869\\
1.88	-0.000147735\\
1.882	-0.000138385\\
1.884	-0.000128851\\
1.886	-0.000119163\\
1.888	-0.000109352\\
1.89	-9.94471e-05\\
1.892	-8.94798e-05\\
1.894	-8.15001e-05\\
1.896	-7.8265e-05\\
1.898	-7.43903e-05\\
1.9	-7.4772e-05\\
1.902	-8.01016e-05\\
1.904	-8.89206e-05\\
1.906	-9.69526e-05\\
1.908	-0.000104182\\
1.91	-0.000110596\\
1.912	-0.000116186\\
1.914	-0.000120944\\
1.916	-0.000124866\\
1.918	-0.000127951\\
1.92	-0.000130202\\
1.922	-0.000131622\\
1.924	-0.000132219\\
1.926	-0.000132003\\
1.928	-0.000130986\\
1.93	-0.000129184\\
1.932	-0.000126615\\
1.934	-0.000124376\\
1.936	-0.000131977\\
1.938	-0.000138716\\
1.94	-0.000144585\\
1.942	-0.000149576\\
1.944	-0.000153685\\
1.946	-0.00015691\\
1.948	-0.000157241\\
1.95	-0.000153266\\
1.952	-0.000148765\\
1.954	-0.000143763\\
1.956	-0.000138281\\
1.958	-0.000132345\\
1.96	-0.000125981\\
1.962	-0.000121542\\
1.964	-0.000118111\\
1.966	-0.000114277\\
1.968	-0.000110059\\
1.97	-0.000107329\\
1.972	-0.000104502\\
1.974	-0.000101316\\
1.976	-9.7783e-05\\
1.978	-9.47383e-05\\
1.98	-9.4564e-05\\
1.982	-9.41772e-05\\
1.984	-9.358e-05\\
1.986	-9.27749e-05\\
1.988	-9.17647e-05\\
1.99	-9.05525e-05\\
1.992	-8.91419e-05\\
1.994	-8.75369e-05\\
1.996	-8.57415e-05\\
1.998	-8.37604e-05\\
2	-8.15986e-05\\
};
\addplot [color=mycolor5, line width=1.0pt, forget plot]
  table[row sep=crcr]{%
-0.024	0\\
-0.022	0\\
-0.02	0\\
-0.018	0\\
-0.016	0\\
-0.014	0\\
-0.012	0\\
-0.01	0\\
-0.008	0\\
-0.006	0\\
-0.004	0\\
-0.002	0\\
0	0\\
0.002	-3.98949e-09\\
0.004	-6.1449e-08\\
0.006	-2.94594e-07\\
0.008	-8.8029e-07\\
0.01	-2.03152e-06\\
0.012	-3.98234e-06\\
0.014	-6.97574e-06\\
0.016	-1.12539e-05\\
0.018	-1.20274e-05\\
0.02	-1.75004e-05\\
0.022	-2.44696e-05\\
0.024	-9.02481e-06\\
0.026	-1.23298e-05\\
0.028	-1.64543e-05\\
0.03	-2.1517e-05\\
0.032	-2.76421e-05\\
0.034	-3.49577e-05\\
0.036	-4.35951e-05\\
0.038	-5.36873e-05\\
0.04	-6.53683e-05\\
0.042	-7.87707e-05\\
0.044	-9.40254e-05\\
0.046	-0.000111259\\
0.048	-0.000130595\\
0.05	-0.000152146\\
0.052	-0.000176022\\
0.054	-0.000202318\\
0.056	-0.000231124\\
0.058	-0.000262512\\
0.06	-0.000296546\\
0.062	-0.000333271\\
0.064	-0.00587162\\
0.066	-0.0131861\\
0.068	-0.0139208\\
0.07	-0.0141888\\
0.072	-0.0144069\\
0.074	-0.0145727\\
0.076	-0.0146841\\
0.078	-0.0147393\\
0.08	-0.0147369\\
0.082	-0.0146755\\
0.084	-0.0151307\\
0.086	-0.0155711\\
0.088	-0.0159616\\
0.09	-0.016299\\
0.092	-0.0165839\\
0.094	-0.0169038\\
0.096	-0.0171709\\
0.098	-0.0175354\\
0.1	-0.0195733\\
0.102	-0.0215149\\
0.104	-0.023351\\
0.106	-0.0250728\\
0.108	-0.0266724\\
0.11	-0.0281424\\
0.112	-0.0266778\\
0.114	-0.0246641\\
0.116	-0.0225566\\
0.118	-0.020361\\
0.12	-0.0180833\\
0.122	-0.0157301\\
0.124	-0.0133082\\
0.126	-0.011258\\
0.128	-0.010454\\
0.13	-0.00963899\\
0.132	-0.0093585\\
0.134	-0.00906378\\
0.136	-0.00875465\\
0.138	-0.00843092\\
0.14	-0.00809241\\
0.142	-0.00773892\\
0.144	-0.0073703\\
0.146	-0.00698677\\
0.148	-0.00665838\\
0.15	-0.00630681\\
0.152	-0.00593291\\
0.154	-0.00553768\\
0.156	-0.00512218\\
0.158	-0.00468759\\
0.16	-0.00423516\\
0.162	-0.00376623\\
0.164	-0.00328224\\
0.166	-0.00241569\\
0.168	-0.00236603\\
0.17	-0.00231163\\
0.172	-0.00225265\\
0.174	-0.00218925\\
0.176	-0.00212164\\
0.178	-0.00205001\\
0.18	-0.00197458\\
0.182	-0.00189558\\
0.184	-0.00181326\\
0.186	-0.00172786\\
0.188	-0.00163966\\
0.19	-0.00154891\\
0.192	-0.0014559\\
0.194	-0.00136091\\
0.196	-0.00126423\\
0.198	-0.00116615\\
0.2	-0.00106695\\
0.202	-0.000966948\\
0.204	-0.00124734\\
0.206	-0.00111644\\
0.208	-0.000983697\\
0.21	-0.000849513\\
0.212	-0.000714311\\
0.214	-0.000578513\\
0.216	-0.000442549\\
0.218	-0.000306848\\
0.22	-0.00017184\\
0.222	-3.79508e-05\\
0.224	9.43946e-05\\
0.226	0.000224778\\
0.228	0.00035279\\
0.23	0.000478027\\
0.232	0.0006001\\
0.234	0.000718628\\
0.236	0.000833248\\
0.238	0.000943611\\
0.24	0.00104938\\
0.242	0.00115025\\
0.244	0.00124592\\
0.246	0.00133611\\
0.248	0.00142058\\
0.25	0.00149909\\
0.252	0.00157144\\
0.254	0.000118663\\
0.256	-0.00364163\\
0.258	-0.00735789\\
0.26	-0.0110118\\
0.262	-0.0139086\\
0.264	-0.0155225\\
0.266	-0.0161536\\
0.268	-0.0167232\\
0.27	-0.0172293\\
0.272	-0.0176704\\
0.274	-0.018045\\
0.276	-0.0183521\\
0.278	-0.0185907\\
0.28	-0.0187603\\
0.282	-0.0188606\\
0.284	-0.0188915\\
0.286	-0.0188532\\
0.288	-0.0187461\\
0.29	-0.018571\\
0.292	-0.0183287\\
0.294	-0.0183776\\
0.296	-0.0189515\\
0.298	-0.0194577\\
0.3	-0.0198946\\
0.302	-0.0205321\\
0.304	-0.0212486\\
0.306	-0.0218727\\
0.308	-0.022403\\
0.31	-0.0228383\\
0.312	-0.0231779\\
0.314	-0.0234216\\
0.316	-0.0235695\\
0.318	-0.023622\\
0.32	-0.02358\\
0.322	-0.0234449\\
0.324	-0.0232183\\
0.326	-0.0229021\\
0.328	-0.0224987\\
0.33	-0.0213775\\
0.332	-0.0201026\\
0.334	-0.0187675\\
0.336	-0.0173777\\
0.338	-0.0161896\\
0.34	-0.0160919\\
0.342	-0.0156524\\
0.344	-0.0150486\\
0.346	-0.0144096\\
0.348	-0.0137374\\
0.35	-0.013034\\
0.352	-0.0123015\\
0.354	-0.0115418\\
0.356	-0.0111375\\
0.358	-0.0110936\\
0.36	-0.011016\\
0.362	-0.010905\\
0.364	-0.0107613\\
0.366	-0.0105853\\
0.368	-0.0103778\\
0.37	-0.0101396\\
0.372	-0.0098716\\
0.374	-0.00902934\\
0.376	-0.00813216\\
0.378	-0.00726359\\
0.38	-0.00642601\\
0.382	-0.0056214\\
0.384	-0.00485132\\
0.386	-0.00411695\\
0.388	-0.0034191\\
0.39	-0.00296904\\
0.392	-0.00322294\\
0.394	-0.00350177\\
0.396	-0.00380247\\
0.398	-0.00412171\\
0.4	-0.004456\\
0.402	-0.00480168\\
0.404	-0.00515498\\
0.406	-0.00600581\\
0.408	-0.0068627\\
0.41	-0.0072519\\
0.412	-0.00755171\\
0.414	-0.00776266\\
0.416	-0.00788558\\
0.418	-0.00792155\\
0.42	-0.00788937\\
0.422	-0.00786387\\
0.424	-0.00778843\\
0.426	-0.00763175\\
0.428	-0.00739572\\
0.43	-0.00708249\\
0.432	-0.00669444\\
0.434	-0.0067134\\
0.436	-0.00684314\\
0.438	-0.00693894\\
0.44	-0.00700114\\
0.442	-0.00703018\\
0.444	-0.00702653\\
0.446	-0.00699077\\
0.448	-0.00692356\\
0.45	-0.00682559\\
0.452	-0.00669766\\
0.454	-0.00654059\\
0.456	-0.0063553\\
0.458	-0.00614273\\
0.46	-0.00590389\\
0.462	-0.00563984\\
0.464	-0.00535168\\
0.466	-0.00504055\\
0.468	-0.00470762\\
0.47	-0.00435411\\
0.472	-0.00465308\\
0.474	-0.005431\\
0.476	-0.00615269\\
0.478	-0.0068164\\
0.48	-0.00742058\\
0.482	-0.00796396\\
0.484	-0.00844549\\
0.486	-0.00886434\\
0.488	-0.00921996\\
0.49	-0.009512\\
0.492	-0.00974038\\
0.494	-0.00990523\\
0.496	-0.0100069\\
0.498	-0.010046\\
0.5	-0.0100234\\
0.502	-0.00994007\\
0.504	-0.00979724\\
0.506	-0.00959637\\
0.508	-0.00959905\\
0.51	-0.0100721\\
0.512	-0.0104846\\
0.514	-0.0108358\\
0.516	-0.0111256\\
0.518	-0.0113539\\
0.52	-0.0115208\\
0.522	-0.0116265\\
0.524	-0.0116717\\
0.526	-0.0116569\\
0.528	-0.0115832\\
0.53	-0.0114681\\
0.532	-0.0113668\\
0.534	-0.0112481\\
0.536	-0.0110743\\
0.538	-0.0108469\\
0.54	-0.0105673\\
0.542	-0.0102372\\
0.544	-0.00985845\\
0.546	-0.00943289\\
0.548	-0.00896261\\
0.55	-0.00844978\\
0.552	-0.00789669\\
0.554	-0.00730573\\
0.556	-0.00667937\\
0.558	-0.00619473\\
0.56	-0.00633931\\
0.562	-0.00652352\\
0.564	-0.00667028\\
0.566	-0.00677971\\
0.568	-0.006852\\
0.57	-0.00688749\\
0.572	-0.00688659\\
0.574	-0.00684981\\
0.576	-0.00677776\\
0.578	-0.00667116\\
0.58	-0.00653078\\
0.582	-0.00635752\\
0.584	-0.00615233\\
0.586	-0.00591626\\
0.588	-0.00565041\\
0.59	-0.00535599\\
0.592	-0.00503423\\
0.594	-0.00468646\\
0.596	-0.00431405\\
0.598	-0.00391844\\
0.6	-0.00350108\\
0.602	-0.00320196\\
0.604	-0.00305183\\
0.606	-0.00287977\\
0.608	-0.00268665\\
0.61	-0.0024734\\
0.612	-0.00224101\\
0.614	-0.00199048\\
0.616	-0.00195724\\
0.618	-0.00239239\\
0.62	-0.00278469\\
0.622	-0.0031336\\
0.624	-0.00343875\\
0.626	-0.0036999\\
0.628	-0.00391695\\
0.63	-0.00408996\\
0.632	-0.00421912\\
0.634	-0.00430477\\
0.636	-0.00434738\\
0.638	-0.00434754\\
0.64	-0.0043192\\
0.642	-0.00430477\\
0.644	-0.00428651\\
0.646	-0.00423074\\
0.648	-0.00413579\\
0.65	-0.00400263\\
0.652	-0.00383235\\
0.654	-0.00362616\\
0.656	-0.00338534\\
0.658	-0.00311127\\
0.66	-0.00280544\\
0.662	-0.00260503\\
0.664	-0.00251724\\
0.666	-0.00241224\\
0.668	-0.00229053\\
0.67	-0.00215266\\
0.672	-0.00199921\\
0.674	-0.00183079\\
0.676	-0.00164805\\
0.678	-0.00147315\\
0.68	-0.00163286\\
0.682	-0.00176518\\
0.684	-0.00194377\\
0.686	-0.00228565\\
0.688	-0.00259911\\
0.69	-0.00288349\\
0.692	-0.00313827\\
0.694	-0.00336301\\
0.696	-0.00355742\\
0.698	-0.00372131\\
0.7	-0.00385459\\
0.702	-0.00395731\\
0.704	-0.0040296\\
0.706	-0.00407172\\
0.708	-0.00408402\\
0.71	-0.00406696\\
0.712	-0.00402111\\
0.714	-0.00394711\\
0.716	-0.00384571\\
0.718	-0.00371776\\
0.72	-0.00356416\\
0.722	-0.00344153\\
0.724	-0.00375531\\
0.726	-0.00406472\\
0.728	-0.0043448\\
0.73	-0.00459506\\
0.732	-0.00481512\\
0.734	-0.00500471\\
0.736	-0.00516365\\
0.738	-0.00529187\\
0.74	-0.00538941\\
0.742	-0.00545638\\
0.744	-0.00549302\\
0.746	-0.00549964\\
0.748	-0.00547666\\
0.75	-0.00543293\\
0.752	-0.00540881\\
0.754	-0.00537822\\
0.756	-0.00533016\\
0.758	-0.00525451\\
0.76	-0.0051519\\
0.762	-0.00502307\\
0.764	-0.00486879\\
0.766	-0.00468995\\
0.768	-0.00448747\\
0.77	-0.00426237\\
0.772	-0.00401569\\
0.774	-0.00374856\\
0.776	-0.00346215\\
0.778	-0.00315767\\
0.78	-0.00294651\\
0.782	-0.00295687\\
0.784	-0.00294865\\
0.786	-0.00292215\\
0.788	-0.00287773\\
0.79	-0.00281582\\
0.792	-0.00294515\\
0.794	-0.00307974\\
0.796	-0.00319533\\
0.798	-0.00329184\\
0.8	-0.00336928\\
0.802	-0.00342768\\
0.804	-0.00346716\\
0.806	-0.00348786\\
0.808	-0.00349001\\
0.81	-0.00347385\\
0.812	-0.00343969\\
0.814	-0.0033879\\
0.816	-0.00331886\\
0.818	-0.00323305\\
0.82	-0.00313093\\
0.822	-0.00301305\\
0.824	-0.00287997\\
0.826	-0.0027323\\
0.828	-0.00257069\\
0.83	-0.00239581\\
0.832	-0.00220837\\
0.834	-0.00200909\\
0.836	-0.00179875\\
0.838	-0.00157811\\
0.84	-0.00174121\\
0.842	-0.001918\\
0.844	-0.00207269\\
0.846	-0.00220515\\
0.848	-0.0023153\\
0.85	-0.00240312\\
0.852	-0.00246871\\
0.854	-0.00251219\\
0.856	-0.00253378\\
0.858	-0.00253377\\
0.86	-0.00251594\\
0.862	-0.00251986\\
0.864	-0.00251315\\
0.866	-0.00250211\\
0.868	-0.00247071\\
0.87	-0.00241939\\
0.872	-0.00234863\\
0.874	-0.00225899\\
0.876	-0.00215107\\
0.878	-0.00202552\\
0.88	-0.00188304\\
0.882	-0.00172438\\
0.884	-0.00155033\\
0.886	-0.00136173\\
0.888	-0.00115944\\
0.89	-0.00112023\\
0.892	-0.00120408\\
0.894	-0.00127812\\
0.896	-0.00134233\\
0.898	-0.0013967\\
0.9	-0.00144125\\
0.902	-0.00147603\\
0.904	-0.00156355\\
0.906	-0.00165685\\
0.908	-0.00173525\\
0.91	-0.00179871\\
0.912	-0.00184725\\
0.914	-0.00188095\\
0.916	-0.00189991\\
0.918	-0.00190432\\
0.92	-0.00189439\\
0.922	-0.00187041\\
0.924	-0.00183268\\
0.926	-0.00178157\\
0.928	-0.00171749\\
0.93	-0.00164089\\
0.932	-0.00160641\\
0.934	-0.00175815\\
0.936	-0.00189516\\
0.938	-0.00201721\\
0.94	-0.0021241\\
0.942	-0.00221567\\
0.944	-0.00229184\\
0.946	-0.00235258\\
0.948	-0.00239791\\
0.95	-0.0024279\\
0.952	-0.00254001\\
0.954	-0.00263735\\
0.956	-0.00271881\\
0.958	-0.00278434\\
0.96	-0.00283397\\
0.962	-0.00286774\\
0.964	-0.00288578\\
0.966	-0.00288824\\
0.968	-0.00287535\\
0.97	-0.00284762\\
0.972	-0.00284183\\
0.974	-0.00282344\\
0.976	-0.00280946\\
0.978	-0.00278096\\
0.98	-0.00273824\\
0.982	-0.00268165\\
0.984	-0.00261156\\
0.986	-0.0025284\\
0.988	-0.00243262\\
0.99	-0.00232472\\
0.992	-0.00220522\\
0.994	-0.00207468\\
0.996	-0.00193369\\
0.998	-0.00178286\\
1	-0.00162283\\
1.002	-0.00145426\\
1.004	-0.0014044\\
1.006	-0.00142615\\
1.008	-0.00143887\\
1.01	-0.00144237\\
1.012	-0.00143674\\
1.014	-0.00142213\\
1.016	-0.00139874\\
1.018	-0.00136677\\
1.02	-0.00132646\\
1.022	-0.00133152\\
1.024	-0.00142129\\
1.026	-0.00150171\\
1.028	-0.00157267\\
1.03	-0.00163411\\
1.032	-0.001686\\
1.034	-0.00172832\\
1.036	-0.00176109\\
1.038	-0.00178435\\
1.04	-0.00179818\\
1.042	-0.00180268\\
1.044	-0.00179796\\
1.046	-0.00178417\\
1.048	-0.00176149\\
1.05	-0.00173012\\
1.052	-0.00169026\\
1.054	-0.00164217\\
1.056	-0.00158611\\
1.058	-0.00152235\\
1.06	-0.00145119\\
1.062	-0.00137296\\
1.064	-0.00135937\\
1.066	-0.00142308\\
1.068	-0.00147521\\
1.07	-0.00151577\\
1.072	-0.00154479\\
1.074	-0.00156236\\
1.076	-0.00156861\\
1.078	-0.00156367\\
1.08	-0.00154776\\
1.082	-0.00154887\\
1.084	-0.00154077\\
1.086	-0.00153608\\
1.088	-0.00152363\\
1.09	-0.00150077\\
1.092	-0.00146773\\
1.094	-0.00142478\\
1.096	-0.00137221\\
1.098	-0.00131035\\
1.1	-0.00123955\\
1.102	-0.00116016\\
1.104	-0.0010726\\
1.106	-0.000977274\\
1.108	-0.000874624\\
1.11	-0.000765104\\
1.112	-0.000696438\\
1.114	-0.00075143\\
1.116	-0.000799084\\
1.118	-0.000839355\\
1.12	-0.000872223\\
1.122	-0.000897695\\
1.124	-0.000915805\\
1.126	-0.000926611\\
1.128	-0.000930198\\
1.13	-0.000926675\\
1.132	-0.000916176\\
1.134	-0.000898857\\
1.136	-0.000874895\\
1.138	-0.000844492\\
1.14	-0.000807866\\
1.142	-0.000816368\\
1.144	-0.000828355\\
1.146	-0.000836367\\
1.148	-0.000918459\\
1.15	-0.000993142\\
1.152	-0.00106028\\
1.154	-0.00111975\\
1.156	-0.00117147\\
1.158	-0.00121538\\
1.16	-0.00125145\\
1.162	-0.00127968\\
1.164	-0.00130008\\
1.166	-0.00131271\\
1.168	-0.00131763\\
1.17	-0.00136176\\
1.172	-0.00141991\\
1.174	-0.00146977\\
1.176	-0.00151131\\
1.178	-0.0015445\\
1.18	-0.00156935\\
1.182	-0.00158589\\
1.184	-0.00159419\\
1.186	-0.00159433\\
1.188	-0.00158643\\
1.19	-0.00157063\\
1.192	-0.00156771\\
1.194	-0.00155953\\
1.196	-0.00155123\\
1.198	-0.00154106\\
1.2	-0.00152335\\
1.202	-0.00149828\\
1.204	-0.00146603\\
1.206	-0.0014268\\
1.208	-0.00138082\\
1.21	-0.00132835\\
1.212	-0.00126963\\
1.214	-0.00120496\\
1.216	-0.00113462\\
1.218	-0.00105893\\
1.22	-0.000978214\\
1.222	-0.000892809\\
1.224	-0.000803064\\
1.226	-0.000709337\\
1.228	-0.000704726\\
1.23	-0.000722189\\
1.232	-0.000734982\\
1.234	-0.000743305\\
1.236	-0.000747115\\
1.238	-0.000746393\\
1.24	-0.000741212\\
1.242	-0.00073166\\
1.244	-0.000717834\\
1.246	-0.000699845\\
1.248	-0.000677815\\
1.25	-0.000651876\\
1.252	-0.000623996\\
1.254	-0.000678312\\
1.256	-0.000728092\\
1.258	-0.000773261\\
1.26	-0.000813753\\
1.262	-0.000849519\\
1.264	-0.000880525\\
1.266	-0.000906749\\
1.268	-0.000928185\\
1.27	-0.000944839\\
1.272	-0.000956733\\
1.274	-0.0009639\\
1.276	-0.000966388\\
1.278	-0.000964256\\
1.28	-0.000957579\\
1.282	-0.00094644\\
1.284	-0.000930938\\
1.286	-0.00091118\\
1.288	-0.000887285\\
1.29	-0.000905746\\
1.292	-0.000919085\\
1.294	-0.000926439\\
1.296	-0.000927872\\
1.298	-0.000923466\\
1.3	-0.000913325\\
1.302	-0.000912692\\
1.304	-0.000909388\\
1.306	-0.000903659\\
1.308	-0.000899623\\
1.31	-0.000890093\\
1.312	-0.000875182\\
1.314	-0.000855019\\
1.316	-0.000829749\\
1.318	-0.00079953\\
1.32	-0.000764535\\
1.322	-0.000724952\\
1.324	-0.00068098\\
1.326	-0.00063283\\
1.328	-0.000580723\\
1.33	-0.000524892\\
1.332	-0.000465579\\
1.334	-0.000457693\\
1.336	-0.000463741\\
1.338	-0.000466231\\
1.34	-0.000465216\\
1.342	-0.000460543\\
1.344	-0.000452375\\
1.346	-0.00044093\\
1.348	-0.000445409\\
1.35	-0.000460171\\
1.352	-0.000471325\\
1.354	-0.000478886\\
1.356	-0.000482881\\
1.358	-0.000483351\\
1.36	-0.000480347\\
1.362	-0.000481768\\
1.364	-0.000524513\\
1.366	-0.000563401\\
1.368	-0.000598355\\
1.37	-0.000629314\\
1.372	-0.000656227\\
1.374	-0.000679062\\
1.376	-0.000697798\\
1.378	-0.000712428\\
1.38	-0.000722962\\
1.382	-0.000729421\\
1.384	-0.000731839\\
1.386	-0.000730267\\
1.388	-0.000734035\\
1.39	-0.000767662\\
1.392	-0.000796968\\
1.394	-0.000821918\\
1.396	-0.000842487\\
1.398	-0.000858668\\
1.4	-0.000870468\\
1.402	-0.000877907\\
1.404	-0.000881019\\
1.406	-0.000879855\\
1.408	-0.000874475\\
1.41	-0.000864954\\
1.412	-0.000862433\\
1.414	-0.000858858\\
1.416	-0.000851623\\
1.418	-0.000848091\\
1.42	-0.000840579\\
1.422	-0.000829165\\
1.424	-0.000813941\\
1.426	-0.000795008\\
1.428	-0.000772479\\
1.43	-0.000746477\\
1.432	-0.000717134\\
1.434	-0.000684592\\
1.436	-0.000649002\\
1.438	-0.000610522\\
1.44	-0.000569319\\
1.442	-0.000525566\\
1.444	-0.000479442\\
1.446	-0.000431132\\
1.448	-0.000380825\\
1.45	-0.000338315\\
1.452	-0.000351973\\
1.454	-0.000363309\\
1.456	-0.000372324\\
1.458	-0.000379023\\
1.46	-0.000383505\\
1.462	-0.000385768\\
1.464	-0.000385777\\
1.466	-0.000383566\\
1.468	-0.000379175\\
1.47	-0.000372652\\
1.472	-0.000364048\\
1.474	-0.000353423\\
1.476	-0.000340838\\
1.478	-0.000326365\\
1.48	-0.000340652\\
1.482	-0.000358158\\
1.484	-0.000372685\\
1.486	-0.000384228\\
1.488	-0.000392793\\
1.49	-0.000398396\\
1.492	-0.000411849\\
1.494	-0.000431595\\
1.496	-0.00044897\\
1.498	-0.000463956\\
1.5	-0.000476543\\
1.502	-0.000486727\\
1.504	-0.000494512\\
1.506	-0.000499909\\
1.508	-0.000503633\\
1.51	-0.000512762\\
1.512	-0.00051873\\
1.514	-0.000521562\\
1.516	-0.000521296\\
1.518	-0.000517977\\
1.52	-0.000511662\\
1.522	-0.000510432\\
1.524	-0.000509121\\
1.526	-0.000504871\\
1.528	-0.000502698\\
1.53	-0.000498921\\
1.532	-0.000492309\\
1.534	-0.000482924\\
1.536	-0.000470838\\
1.538	-0.000456131\\
1.54	-0.00043889\\
1.542	-0.000419209\\
1.544	-0.000397189\\
1.546	-0.000372938\\
1.548	-0.000346569\\
1.55	-0.0003182\\
1.552	-0.000287956\\
1.554	-0.000255965\\
1.556	-0.00022597\\
1.558	-0.000236406\\
1.56	-0.000244932\\
1.562	-0.000251547\\
1.564	-0.000256256\\
1.566	-0.000259069\\
1.568	-0.000260004\\
1.57	-0.000259087\\
1.572	-0.000256348\\
1.574	-0.000251823\\
1.576	-0.000245556\\
1.578	-0.000258791\\
1.58	-0.000280998\\
1.582	-0.000301171\\
1.584	-0.000319272\\
1.586	-0.000335267\\
1.588	-0.000349133\\
1.59	-0.000360852\\
1.592	-0.000370414\\
1.594	-0.000377818\\
1.596	-0.000383068\\
1.598	-0.000386177\\
1.6	-0.000387164\\
1.602	-0.000386056\\
1.604	-0.000382886\\
1.606	-0.000377694\\
1.608	-0.00039276\\
1.61	-0.000409643\\
1.612	-0.000424254\\
1.614	-0.000436576\\
1.616	-0.0004466\\
1.618	-0.000454323\\
1.62	-0.000459752\\
1.622	-0.0004629\\
1.624	-0.000463787\\
1.626	-0.000462442\\
1.628	-0.000458899\\
1.63	-0.0004532\\
1.632	-0.000451227\\
1.634	-0.000449679\\
1.636	-0.000446014\\
1.638	-0.000442893\\
1.64	-0.000439851\\
1.642	-0.000434765\\
1.644	-0.000427679\\
1.646	-0.000418643\\
1.648	-0.000407714\\
1.65	-0.000394954\\
1.652	-0.00038043\\
1.654	-0.000364214\\
1.656	-0.000346383\\
1.658	-0.000327017\\
1.66	-0.000306202\\
1.662	-0.000284026\\
1.664	-0.000260581\\
1.666	-0.000235963\\
1.668	-0.00021027\\
1.67	-0.000183601\\
1.672	-0.000178331\\
1.674	-0.000179386\\
1.676	-0.000179533\\
1.678	-0.000178785\\
1.68	-0.000177159\\
1.682	-0.000178963\\
1.684	-0.00018292\\
1.686	-0.000185781\\
1.688	-0.000187554\\
1.69	-0.00018821\\
1.692	-0.000187763\\
1.694	-0.000190002\\
1.696	-0.000200428\\
1.698	-0.000209293\\
1.7	-0.000216591\\
1.702	-0.000222318\\
1.704	-0.00022648\\
1.706	-0.000229085\\
1.708	-0.000230148\\
1.71	-0.000229688\\
1.712	-0.000227729\\
1.714	-0.0002243\\
1.716	-0.000224657\\
1.718	-0.000238493\\
1.72	-0.000250693\\
1.722	-0.000261239\\
1.724	-0.000270121\\
1.726	-0.000277333\\
1.728	-0.000282875\\
1.73	-0.000286752\\
1.732	-0.000288976\\
1.734	-0.000289563\\
1.736	-0.000288533\\
1.738	-0.000285914\\
1.74	-0.000281736\\
1.742	-0.000280295\\
1.744	-0.000279489\\
1.746	-0.00027715\\
1.748	-0.000274467\\
1.75	-0.000272812\\
1.752	-0.000269675\\
1.754	-0.000265089\\
1.756	-0.000259088\\
1.758	-0.000251715\\
1.76	-0.000243011\\
1.762	-0.000233026\\
1.764	-0.000221811\\
1.766	-0.000209421\\
1.768	-0.000195913\\
1.77	-0.000181349\\
1.772	-0.000165793\\
1.774	-0.00014931\\
1.776	-0.000144557\\
1.778	-0.000145645\\
1.78	-0.000145731\\
1.782	-0.000144831\\
1.784	-0.000142963\\
1.786	-0.000140146\\
1.788	-0.000136406\\
1.79	-0.00013177\\
1.792	-0.000129042\\
1.794	-0.000141812\\
1.796	-0.000153537\\
1.798	-0.00016419\\
1.8	-0.000173752\\
1.802	-0.000182205\\
1.804	-0.000189537\\
1.806	-0.000195737\\
1.808	-0.000200802\\
1.81	-0.000204729\\
1.812	-0.000207521\\
1.814	-0.000209184\\
1.816	-0.000209729\\
1.818	-0.000209169\\
1.82	-0.000207522\\
1.822	-0.000204808\\
1.824	-0.000201051\\
1.826	-0.000206165\\
1.828	-0.000215644\\
1.83	-0.000223936\\
1.832	-0.000231028\\
1.834	-0.000236914\\
1.836	-0.00024159\\
1.838	-0.000245058\\
1.84	-0.000247322\\
1.842	-0.000248391\\
1.844	-0.000248277\\
1.846	-0.000246998\\
1.848	-0.000244573\\
1.85	-0.000241025\\
1.852	-0.000239521\\
1.854	-0.000238684\\
1.856	-0.000236741\\
1.858	-0.000233964\\
1.86	-0.000232674\\
1.862	-0.000230316\\
1.864	-0.00022691\\
1.866	-0.000222484\\
1.868	-0.000217066\\
1.87	-0.000210686\\
1.872	-0.000203379\\
1.874	-0.000195183\\
1.876	-0.000186135\\
1.878	-0.000176278\\
1.88	-0.000165656\\
1.882	-0.000154313\\
1.884	-0.000142298\\
1.886	-0.00012966\\
1.888	-0.000116449\\
1.89	-0.000102716\\
1.892	-8.85152e-05\\
1.894	-7.80751e-05\\
1.896	-7.47561e-05\\
1.898	-7.08057e-05\\
1.9	-7.112e-05\\
1.902	-7.63907e-05\\
1.904	-8.51593e-05\\
1.906	-9.31495e-05\\
1.908	-0.000100346\\
1.91	-0.000106736\\
1.912	-0.000112309\\
1.914	-0.00011706\\
1.916	-0.000120983\\
1.918	-0.000124078\\
1.92	-0.000126347\\
1.922	-0.000127793\\
1.924	-0.000128424\\
1.926	-0.000128249\\
1.928	-0.000127282\\
1.93	-0.000125536\\
1.932	-0.000123029\\
1.934	-0.00012086\\
1.936	-0.000128536\\
1.938	-0.000135358\\
1.94	-0.000141314\\
1.942	-0.000146397\\
1.944	-0.000150602\\
1.946	-0.000153929\\
1.948	-0.000156378\\
1.95	-0.000157954\\
1.952	-0.000158664\\
1.954	-0.000158517\\
1.956	-0.000157526\\
1.958	-0.000155706\\
1.96	-0.000153076\\
1.962	-0.000151982\\
1.964	-0.000151502\\
1.966	-0.000150223\\
1.968	-0.000148159\\
1.97	-0.000147183\\
1.972	-0.000145711\\
1.974	-0.000143481\\
1.976	-0.000140512\\
1.978	-0.000137642\\
1.98	-0.00013726\\
1.982	-0.000136291\\
1.984	-0.000134744\\
1.986	-0.000132633\\
1.988	-0.000129971\\
1.99	-0.000126774\\
1.992	-0.000123058\\
1.994	-0.000118841\\
1.996	-0.000114142\\
1.998	-0.000108982\\
2	-0.000103382\\
};
\addplot [color=mycolor6, line width=1.0pt, forget plot]
  table[row sep=crcr]{%
-0.024	0\\
-0.022	0\\
-0.02	0\\
-0.018	0\\
-0.016	0\\
-0.014	0\\
-0.012	0\\
-0.01	0\\
-0.008	0\\
-0.006	0\\
-0.004	0\\
-0.002	0\\
0	0\\
0.002	-4.6808e-09\\
0.004	-7.23594e-08\\
0.006	-3.47822e-07\\
0.008	-1.04157e-06\\
0.01	-2.40836e-06\\
0.012	-4.72972e-06\\
0.014	-8.29978e-06\\
0.016	-1.34139e-05\\
0.018	-1.51926e-05\\
0.02	-2.21422e-05\\
0.022	-3.10099e-05\\
0.024	-2.68702e-05\\
0.026	-3.64379e-05\\
0.028	-4.8236e-05\\
0.03	-6.25351e-05\\
0.032	-7.96038e-05\\
0.034	-9.97054e-05\\
0.036	-0.000123093\\
0.038	-0.000150008\\
0.04	-0.000180671\\
0.042	-0.000215284\\
0.044	-0.000254024\\
0.046	-0.000297039\\
0.048	-0.000344444\\
0.05	-0.000396322\\
0.052	-0.000452717\\
0.054	-0.000513635\\
0.056	-0.000579038\\
0.058	-0.000648849\\
0.06	-0.000722945\\
0.062	-0.000801159\\
0.064	-0.00638218\\
0.066	-0.0137269\\
0.068	-0.0140563\\
0.07	-0.0143394\\
0.072	-0.0145736\\
0.074	-0.0147566\\
0.076	-0.0148863\\
0.078	-0.0149608\\
0.08	-0.0149788\\
0.082	-0.0149389\\
0.084	-0.0154166\\
0.086	-0.0158805\\
0.088	-0.0162954\\
0.09	-0.0166582\\
0.092	-0.0169692\\
0.094	-0.017316\\
0.096	-0.0176106\\
0.098	-0.0180033\\
0.1	-0.0200698\\
0.102	-0.0220405\\
0.104	-0.0239059\\
0.106	-0.0256571\\
0.108	-0.0272861\\
0.11	-0.0287855\\
0.112	-0.0275677\\
0.114	-0.0255944\\
0.116	-0.0235267\\
0.118	-0.0213701\\
0.12	-0.0191304\\
0.122	-0.0166625\\
0.124	-0.0141154\\
0.126	-0.0120516\\
0.128	-0.0112688\\
0.13	-0.0104735\\
0.132	-0.0102113\\
0.134	-0.00993309\\
0.136	-0.00963869\\
0.138	-0.00932776\\
0.14	-0.00899997\\
0.142	-0.00865503\\
0.144	-0.00829266\\
0.146	-0.00791302\\
0.148	-0.00758605\\
0.15	-0.00723335\\
0.152	-0.00685572\\
0.154	-0.00645409\\
0.156	-0.00602948\\
0.158	-0.00558305\\
0.16	-0.00511602\\
0.162	-0.00462973\\
0.164	-0.00412563\\
0.166	-0.0014101\\
0.168	-0.00131104\\
0.17	-0.00121433\\
0.172	-0.00112051\\
0.174	-0.00103009\\
0.176	-0.000943535\\
0.178	-0.000861274\\
0.18	-0.000783691\\
0.182	-0.000711124\\
0.184	-0.00064386\\
0.186	-0.000582136\\
0.188	-0.000526137\\
0.19	-0.000475992\\
0.192	-0.000431776\\
0.194	-0.000393509\\
0.196	-0.000361156\\
0.198	-0.000334629\\
0.2	-0.000313785\\
0.202	-0.00029843\\
0.204	-0.00195012\\
0.206	-0.00615007\\
0.208	-0.00969156\\
0.21	-0.0100579\\
0.212	-0.0103631\\
0.214	-0.0106053\\
0.216	-0.0107829\\
0.218	-0.0108948\\
0.22	-0.0109403\\
0.222	-0.0109189\\
0.224	-0.010831\\
0.226	-0.0106771\\
0.228	-0.0104581\\
0.23	-0.0101756\\
0.232	-0.00983124\\
0.234	-0.00942738\\
0.236	-0.00896657\\
0.238	-0.00845178\\
0.24	-0.00788629\\
0.242	-0.00727371\\
0.244	-0.00661794\\
0.246	-0.00592314\\
0.248	-0.0051937\\
0.25	-0.00443421\\
0.252	-0.00364944\\
0.254	-0.00436307\\
0.256	-0.00731765\\
0.258	-0.00929032\\
0.26	-0.0111648\\
0.262	-0.0129357\\
0.264	-0.0145983\\
0.266	-0.0153956\\
0.268	-0.0160007\\
0.27	-0.0165453\\
0.272	-0.0170275\\
0.274	-0.0174457\\
0.276	-0.0177986\\
0.278	-0.0180852\\
0.28	-0.0183045\\
0.282	-0.018456\\
0.284	-0.0185395\\
0.286	-0.0185548\\
0.288	-0.0185021\\
0.29	-0.018382\\
0.292	-0.018195\\
0.294	-0.0182992\\
0.296	-0.0189283\\
0.298	-0.0194891\\
0.3	-0.0199799\\
0.302	-0.0206704\\
0.304	-0.0214387\\
0.306	-0.0221132\\
0.308	-0.0226922\\
0.31	-0.0231744\\
0.312	-0.0235588\\
0.314	-0.0238451\\
0.316	-0.0240331\\
0.318	-0.0241233\\
0.32	-0.0241162\\
0.322	-0.0240132\\
0.324	-0.0238156\\
0.326	-0.0235255\\
0.328	-0.0230918\\
0.33	-0.0219208\\
0.332	-0.0206815\\
0.334	-0.0193792\\
0.336	-0.0180193\\
0.338	-0.0168581\\
0.34	-0.0167842\\
0.342	-0.0166284\\
0.344	-0.0160359\\
0.346	-0.0154038\\
0.348	-0.0147341\\
0.35	-0.014029\\
0.352	-0.0132905\\
0.354	-0.0125209\\
0.356	-0.0121025\\
0.358	-0.0120409\\
0.36	-0.0119418\\
0.362	-0.0118058\\
0.364	-0.0116336\\
0.366	-0.0114261\\
0.368	-0.0111841\\
0.37	-0.0109086\\
0.372	-0.0106007\\
0.374	-0.0102616\\
0.376	-0.00989262\\
0.378	-0.00949508\\
0.38	-0.00907041\\
0.382	-0.00862014\\
0.384	-0.00814584\\
0.386	-0.00764914\\
0.388	-0.00705828\\
0.39	-0.00645212\\
0.392	-0.00648501\\
0.394	-0.00648109\\
0.396	-0.00644079\\
0.398	-0.00636461\\
0.4	-0.00625314\\
0.402	-0.00610702\\
0.404	-0.00592699\\
0.406	-0.00620763\\
0.408	-0.00665924\\
0.41	-0.00702213\\
0.412	-0.00729658\\
0.414	-0.00748319\\
0.416	-0.0075828\\
0.418	-0.00759655\\
0.42	-0.00754327\\
0.422	-0.0074978\\
0.424	-0.00740356\\
0.426	-0.00722923\\
0.428	-0.00697674\\
0.43	-0.00664824\\
0.432	-0.0062461\\
0.434	-0.00625215\\
0.436	-0.00637015\\
0.438	-0.00645537\\
0.44	-0.00650813\\
0.442	-0.00652885\\
0.444	-0.00651798\\
0.446	-0.00647609\\
0.448	-0.00640378\\
0.45	-0.00630174\\
0.452	-0.00617071\\
0.454	-0.00601151\\
0.456	-0.00582499\\
0.458	-0.00561208\\
0.46	-0.00537374\\
0.462	-0.005111\\
0.464	-0.00482491\\
0.466	-0.00451657\\
0.468	-0.00418712\\
0.47	-0.00383773\\
0.472	-0.00414142\\
0.474	-0.00492463\\
0.476	-0.00565213\\
0.478	-0.00632212\\
0.48	-0.00693304\\
0.482	-0.00748355\\
0.484	-0.00797257\\
0.486	-0.00839926\\
0.488	-0.00876299\\
0.49	-0.0090634\\
0.492	-0.00930037\\
0.494	-0.009474\\
0.496	-0.00958461\\
0.498	-0.00963279\\
0.5	-0.0096193\\
0.502	-0.00954514\\
0.504	-0.00941154\\
0.506	-0.00921989\\
0.508	-0.00923176\\
0.51	-0.009714\\
0.512	-0.0101355\\
0.514	-0.0104957\\
0.516	-0.0107944\\
0.518	-0.0110313\\
0.52	-0.0112067\\
0.522	-0.0113208\\
0.524	-0.0113741\\
0.526	-0.0113673\\
0.528	-0.0113013\\
0.53	-0.0111935\\
0.532	-0.0110995\\
0.534	-0.0109876\\
0.536	-0.0108204\\
0.538	-0.0105993\\
0.54	-0.0103257\\
0.542	-0.0100012\\
0.544	-0.00962777\\
0.546	-0.00920718\\
0.548	-0.00874152\\
0.55	-0.00823296\\
0.552	-0.00768378\\
0.554	-0.00709637\\
0.556	-0.00647319\\
0.558	-0.00599137\\
0.56	-0.00613841\\
0.562	-0.00632471\\
0.564	-0.0064732\\
0.566	-0.00658399\\
0.568	-0.00665729\\
0.57	-0.00669343\\
0.572	-0.00669282\\
0.574	-0.00665599\\
0.576	-0.00658356\\
0.578	-0.00647624\\
0.58	-0.00633483\\
0.582	-0.00616022\\
0.584	-0.00595338\\
0.586	-0.00571536\\
0.588	-0.00544729\\
0.59	-0.00515036\\
0.592	-0.00482586\\
0.594	-0.0044751\\
0.596	-0.00409947\\
0.598	-0.00370042\\
0.6	-0.00327944\\
0.602	-0.00297651\\
0.604	-0.0028224\\
0.606	-0.00264621\\
0.608	-0.00244883\\
0.61	-0.00223122\\
0.612	-0.00199435\\
0.614	-0.00173928\\
0.616	-0.00170143\\
0.618	-0.00213193\\
0.62	-0.00251954\\
0.622	-0.00286377\\
0.624	-0.00316424\\
0.626	-0.00342073\\
0.628	-0.00363317\\
0.63	-0.00380163\\
0.632	-0.00392632\\
0.634	-0.00400759\\
0.636	-0.00404592\\
0.638	-0.00404194\\
0.64	-0.00400958\\
0.642	-0.00399129\\
0.644	-0.00396934\\
0.646	-0.00391005\\
0.648	-0.00381177\\
0.65	-0.00367548\\
0.652	-0.00350228\\
0.654	-0.00329339\\
0.656	-0.00305009\\
0.658	-0.00277379\\
0.66	-0.00246595\\
0.662	-0.0022638\\
0.664	-0.0021745\\
0.666	-0.00206825\\
0.668	-0.00194556\\
0.67	-0.00180696\\
0.672	-0.00165304\\
0.674	-0.00148442\\
0.676	-0.00130172\\
0.678	-0.00112713\\
0.68	-0.0012874\\
0.682	-0.00142053\\
0.684	-0.00160017\\
0.686	-0.00194335\\
0.688	-0.00225834\\
0.69	-0.00254447\\
0.692	-0.00280122\\
0.694	-0.00302815\\
0.696	-0.00322494\\
0.698	-0.00339141\\
0.7	-0.00352745\\
0.702	-0.00363309\\
0.704	-0.00370846\\
0.706	-0.00375381\\
0.708	-0.00376947\\
0.71	-0.00375591\\
0.712	-0.00371365\\
0.714	-0.00364334\\
0.716	-0.00354573\\
0.718	-0.00342162\\
0.72	-0.00327193\\
0.722	-0.00315326\\
0.724	-0.00347102\\
0.726	-0.00378444\\
0.728	-0.00406853\\
0.73	-0.0043228\\
0.732	-0.00454684\\
0.734	-0.00474039\\
0.736	-0.00490325\\
0.738	-0.00503534\\
0.74	-0.00513668\\
0.742	-0.00520738\\
0.744	-0.00524767\\
0.746	-0.00525784\\
0.748	-0.00523831\\
0.75	-0.00519793\\
0.752	-0.00517703\\
0.754	-0.00514955\\
0.756	-0.00510446\\
0.758	-0.00503164\\
0.76	-0.00493173\\
0.762	-0.00480544\\
0.764	-0.00465356\\
0.766	-0.00447697\\
0.768	-0.00427659\\
0.77	-0.00405342\\
0.772	-0.00380852\\
0.774	-0.00354301\\
0.776	-0.00325806\\
0.778	-0.00295488\\
0.78	-0.00274487\\
0.782	-0.00275622\\
0.784	-0.00274884\\
0.786	-0.00272304\\
0.788	-0.00267917\\
0.79	-0.00261766\\
0.792	-0.00274726\\
0.794	-0.002882\\
0.796	-0.00299759\\
0.798	-0.003094\\
0.8	-0.00317121\\
0.802	-0.00322928\\
0.804	-0.00326832\\
0.806	-0.0032885\\
0.808	-0.00329002\\
0.81	-0.00327316\\
0.812	-0.00323824\\
0.814	-0.0031856\\
0.816	-0.00311568\\
0.818	-0.00302891\\
0.82	-0.00292581\\
0.822	-0.0028069\\
0.824	-0.00267277\\
0.826	-0.00252404\\
0.828	-0.00236134\\
0.83	-0.00218538\\
0.832	-0.00199685\\
0.834	-0.00179649\\
0.836	-0.00158509\\
0.838	-0.00136341\\
0.84	-0.00152551\\
0.842	-0.00170132\\
0.844	-0.00185508\\
0.846	-0.00198665\\
0.848	-0.00209596\\
0.85	-0.00218301\\
0.852	-0.00224788\\
0.854	-0.00229071\\
0.856	-0.00231172\\
0.858	-0.00231119\\
0.86	-0.00229293\\
0.862	-0.00229649\\
0.864	-0.0022895\\
0.866	-0.00227827\\
0.868	-0.00224676\\
0.87	-0.00219541\\
0.872	-0.00212472\\
0.874	-0.00203523\\
0.876	-0.00192754\\
0.878	-0.00180231\\
0.88	-0.00166024\\
0.882	-0.00150208\\
0.884	-0.00132861\\
0.886	-0.00114067\\
0.888	-0.000939131\\
0.89	-0.000900744\\
0.892	-0.000985494\\
0.894	-0.00106051\\
0.896	-0.00112577\\
0.898	-0.00118126\\
0.9	-0.001227\\
0.902	-0.00126303\\
0.904	-0.00135186\\
0.906	-0.00144653\\
0.908	-0.00152634\\
0.91	-0.00159127\\
0.912	-0.00164133\\
0.914	-0.00167657\\
0.916	-0.00169713\\
0.918	-0.00170316\\
0.92	-0.00169489\\
0.922	-0.00167259\\
0.924	-0.00163656\\
0.926	-0.00158717\\
0.928	-0.00152482\\
0.93	-0.00144996\\
0.932	-0.00141723\\
0.934	-0.00157072\\
0.936	-0.00170949\\
0.938	-0.00183328\\
0.94	-0.0019419\\
0.942	-0.0020352\\
0.944	-0.00211308\\
0.946	-0.00217552\\
0.948	-0.00222252\\
0.95	-0.00225417\\
0.952	-0.0023679\\
0.954	-0.00246684\\
0.956	-0.00254987\\
0.958	-0.00261694\\
0.96	-0.00266807\\
0.962	-0.00270331\\
0.964	-0.00272278\\
0.966	-0.00272664\\
0.968	-0.0027151\\
0.97	-0.00268868\\
0.972	-0.00268417\\
0.974	-0.00266702\\
0.976	-0.00265423\\
0.978	-0.00262688\\
0.98	-0.00258527\\
0.982	-0.00252974\\
0.984	-0.00246068\\
0.986	-0.0023785\\
0.988	-0.00228367\\
0.99	-0.00217667\\
0.992	-0.00205803\\
0.994	-0.00192832\\
0.996	-0.00178812\\
0.998	-0.00163803\\
1	-0.00147872\\
1.002	-0.00131082\\
1.004	-0.00126161\\
1.006	-0.00128397\\
1.008	-0.00129727\\
1.01	-0.00130133\\
1.012	-0.00129623\\
1.014	-0.00128212\\
1.016	-0.0012592\\
1.018	-0.00122769\\
1.02	-0.00118781\\
1.022	-0.00119328\\
1.024	-0.00128344\\
1.026	-0.00136423\\
1.028	-0.00143556\\
1.03	-0.00149735\\
1.032	-0.00154957\\
1.034	-0.00159221\\
1.036	-0.00162529\\
1.038	-0.00164886\\
1.04	-0.00166299\\
1.042	-0.00166778\\
1.044	-0.00166335\\
1.046	-0.00164985\\
1.048	-0.00162746\\
1.05	-0.00159637\\
1.052	-0.00155681\\
1.054	-0.00150901\\
1.056	-0.00145324\\
1.058	-0.00138977\\
1.06	-0.00131891\\
1.062	-0.00124099\\
1.064	-0.00122771\\
1.066	-0.00129175\\
1.068	-0.0013442\\
1.07	-0.00138509\\
1.072	-0.00141446\\
1.074	-0.00143238\\
1.076	-0.00143899\\
1.078	-0.00143443\\
1.08	-0.00141889\\
1.082	-0.0014204\\
1.084	-0.0014127\\
1.086	-0.00140842\\
1.088	-0.0013964\\
1.09	-0.00137397\\
1.092	-0.00134138\\
1.094	-0.00129888\\
1.096	-0.00124678\\
1.098	-0.0011854\\
1.1	-0.00111508\\
1.102	-0.00103619\\
1.104	-0.000949141\\
1.106	-0.000854335\\
1.108	-0.000752213\\
1.11	-0.000643231\\
1.112	-0.000575113\\
1.114	-0.000630661\\
1.116	-0.000678881\\
1.118	-0.000719725\\
1.12	-0.000753175\\
1.122	-0.000779236\\
1.124	-0.000797942\\
1.126	-0.000809351\\
1.128	-0.000813549\\
1.13	-0.000810643\\
1.132	-0.000800766\\
1.134	-0.000784075\\
1.136	-0.000760747\\
1.138	-0.000730981\\
1.14	-0.000694997\\
1.142	-0.000704146\\
1.144	-0.000716782\\
1.146	-0.000725447\\
1.148	-0.000808195\\
1.15	-0.000883536\\
1.152	-0.000951329\\
1.154	-0.00101146\\
1.156	-0.00106385\\
1.158	-0.00110842\\
1.16	-0.00114516\\
1.162	-0.00117405\\
1.164	-0.00119512\\
1.166	-0.00120841\\
1.168	-0.00121399\\
1.17	-0.00125878\\
1.172	-0.00131759\\
1.174	-0.00136811\\
1.176	-0.0014103\\
1.178	-0.00144415\\
1.18	-0.00146964\\
1.182	-0.00148683\\
1.184	-0.00149577\\
1.186	-0.00149655\\
1.188	-0.00148929\\
1.19	-0.00147412\\
1.192	-0.00147181\\
1.194	-0.00146426\\
1.196	-0.00145657\\
1.198	-0.001447\\
1.2	-0.0014299\\
1.202	-0.00140542\\
1.204	-0.00137376\\
1.206	-0.00133512\\
1.208	-0.00128972\\
1.21	-0.00123781\\
1.212	-0.00117966\\
1.214	-0.00111554\\
1.216	-0.00104575\\
1.218	-0.000970599\\
1.22	-0.000890416\\
1.222	-0.000805536\\
1.224	-0.000716308\\
1.226	-0.00062309\\
1.228	-0.000618981\\
1.23	-0.000636937\\
1.232	-0.000650215\\
1.234	-0.000659016\\
1.236	-0.000663295\\
1.238	-0.000663034\\
1.24	-0.000658308\\
1.242	-0.000649202\\
1.244	-0.000635814\\
1.246	-0.000618256\\
1.248	-0.000596649\\
1.25	-0.000571125\\
1.252	-0.000543653\\
1.254	-0.000598371\\
1.256	-0.000648546\\
1.258	-0.000694102\\
1.26	-0.000734975\\
1.262	-0.000771116\\
1.264	-0.000802491\\
1.266	-0.000829078\\
1.268	-0.000850871\\
1.27	-0.000867877\\
1.272	-0.000880117\\
1.274	-0.000887626\\
1.276	-0.000890451\\
1.278	-0.000888653\\
1.28	-0.000882304\\
1.282	-0.00087149\\
1.284	-0.000856309\\
1.286	-0.00083687\\
1.288	-0.000813291\\
1.29	-0.000832065\\
1.292	-0.000845715\\
1.294	-0.000853377\\
1.296	-0.000855116\\
1.298	-0.000851016\\
1.3	-0.000841179\\
1.302	-0.000840849\\
1.304	-0.000837848\\
1.306	-0.000832422\\
1.308	-0.000828688\\
1.31	-0.00081946\\
1.312	-0.000804851\\
1.314	-0.000784991\\
1.316	-0.000760024\\
1.318	-0.000730109\\
1.32	-0.000695421\\
1.322	-0.000656145\\
1.324	-0.000612481\\
1.326	-0.00056464\\
1.328	-0.000512844\\
1.33	-0.000457327\\
1.332	-0.000398328\\
1.334	-0.000390759\\
1.336	-0.000397126\\
1.338	-0.000399937\\
1.34	-0.000399245\\
1.342	-0.000395113\\
1.344	-0.000387614\\
1.346	-0.000376836\\
1.348	-0.000381977\\
1.35	-0.000397219\\
1.352	-0.000408688\\
1.354	-0.000416565\\
1.356	-0.000420879\\
1.358	-0.000421668\\
1.36	-0.000418985\\
1.362	-0.00042073\\
1.364	-0.000463799\\
1.366	-0.000503013\\
1.368	-0.000538295\\
1.37	-0.000569582\\
1.372	-0.000596825\\
1.374	-0.00061999\\
1.376	-0.000639057\\
1.378	-0.00065402\\
1.38	-0.000664886\\
1.382	-0.000671678\\
1.384	-0.00067443\\
1.386	-0.000673191\\
1.388	-0.000677292\\
1.39	-0.000711252\\
1.392	-0.000740892\\
1.394	-0.000766174\\
1.396	-0.000787075\\
1.398	-0.000803587\\
1.4	-0.000815717\\
1.402	-0.000823484\\
1.404	-0.000826925\\
1.406	-0.000826087\\
1.408	-0.000821032\\
1.41	-0.000811835\\
1.412	-0.000809635\\
1.414	-0.00080638\\
1.416	-0.000799463\\
1.418	-0.000796246\\
1.42	-0.000789047\\
1.422	-0.000777945\\
1.424	-0.000763029\\
1.426	-0.000744403\\
1.428	-0.000722177\\
1.43	-0.000696476\\
1.432	-0.000667431\\
1.434	-0.000635184\\
1.436	-0.000599887\\
1.438	-0.000561697\\
1.44	-0.000520781\\
1.442	-0.000477312\\
1.444	-0.000431469\\
1.446	-0.000383437\\
1.448	-0.000333405\\
1.45	-0.000291167\\
1.452	-0.000305094\\
1.454	-0.000316696\\
1.456	-0.000325974\\
1.458	-0.000332934\\
1.46	-0.000337673\\
1.462	-0.00034019\\
1.464	-0.00034045\\
1.466	-0.000338488\\
1.468	-0.000334344\\
1.47	-0.000328064\\
1.472	-0.000319702\\
1.474	-0.000309314\\
1.476	-0.000296965\\
1.478	-0.000282725\\
1.48	-0.000297243\\
1.482	-0.000314978\\
1.484	-0.000329731\\
1.486	-0.000341498\\
1.488	-0.000350285\\
1.49	-0.000356109\\
1.492	-0.00036978\\
1.494	-0.000389742\\
1.496	-0.000407331\\
1.498	-0.00042253\\
1.5	-0.000435328\\
1.502	-0.000445721\\
1.504	-0.000453715\\
1.506	-0.000459318\\
1.508	-0.000463247\\
1.51	-0.00047258\\
1.512	-0.000478751\\
1.514	-0.000481785\\
1.516	-0.000481719\\
1.518	-0.0004786\\
1.52	-0.000472483\\
1.522	-0.000471451\\
1.524	-0.000470336\\
1.526	-0.000466282\\
1.528	-0.000464305\\
1.53	-0.000460722\\
1.532	-0.000454303\\
1.534	-0.000445111\\
1.536	-0.000433218\\
1.538	-0.000418703\\
1.54	-0.000401654\\
1.542	-0.000382164\\
1.544	-0.000360334\\
1.546	-0.000336273\\
1.548	-0.000310093\\
1.55	-0.000281913\\
1.552	-0.000251858\\
1.554	-0.000220055\\
1.556	-0.000190248\\
1.558	-0.000200872\\
1.56	-0.000209585\\
1.562	-0.000216387\\
1.564	-0.000221282\\
1.566	-0.000224281\\
1.568	-0.000225402\\
1.57	-0.00022467\\
1.572	-0.000222116\\
1.574	-0.000217776\\
1.576	-0.000211693\\
1.578	-0.000225112\\
1.58	-0.000247502\\
1.582	-0.000267858\\
1.584	-0.000286141\\
1.586	-0.000302319\\
1.588	-0.000316366\\
1.59	-0.000328266\\
1.592	-0.000338008\\
1.594	-0.000345592\\
1.596	-0.000351021\\
1.598	-0.000354309\\
1.6	-0.000355474\\
1.602	-0.000354544\\
1.604	-0.000351551\\
1.606	-0.000346536\\
1.608	-0.000361778\\
1.61	-0.000378836\\
1.612	-0.000393622\\
1.614	-0.000406117\\
1.616	-0.000416314\\
1.618	-0.00042421\\
1.62	-0.00042981\\
1.622	-0.000433129\\
1.624	-0.000434186\\
1.626	-0.00043301\\
1.628	-0.000429635\\
1.63	-0.000424104\\
1.632	-0.000422297\\
1.634	-0.000420916\\
1.636	-0.000417416\\
1.638	-0.000414484\\
1.64	-0.000411637\\
1.642	-0.000406742\\
1.644	-0.000399847\\
1.646	-0.000391\\
1.648	-0.000380259\\
1.65	-0.000367685\\
1.652	-0.000353346\\
1.654	-0.000337313\\
1.656	-0.000319663\\
1.658	-0.000300478\\
1.66	-0.000279842\\
1.662	-0.000257844\\
1.664	-0.000234577\\
1.666	-0.000210134\\
1.668	-0.000184615\\
1.67	-0.000158119\\
1.672	-0.000153022\\
1.674	-0.000154247\\
1.676	-0.000154564\\
1.678	-0.000153985\\
1.68	-0.000152526\\
1.682	-0.000154496\\
1.684	-0.000158619\\
1.686	-0.000161644\\
1.688	-0.00016358\\
1.69	-0.000164399\\
1.692	-0.000164113\\
1.694	-0.000166512\\
1.696	-0.000177098\\
1.698	-0.000186122\\
1.7	-0.000193576\\
1.702	-0.000199461\\
1.704	-0.000203778\\
1.706	-0.000206538\\
1.708	-0.000207755\\
1.71	-0.000207448\\
1.712	-0.000205641\\
1.714	-0.000202363\\
1.716	-0.000202871\\
1.718	-0.000216857\\
1.72	-0.000229206\\
1.722	-0.0002399\\
1.724	-0.000248929\\
1.726	-0.000256288\\
1.728	-0.000261976\\
1.73	-0.000265998\\
1.732	-0.000268366\\
1.734	-0.000269096\\
1.736	-0.000268209\\
1.738	-0.000265732\\
1.74	-0.000261695\\
1.742	-0.000260395\\
1.744	-0.000259728\\
1.746	-0.000257528\\
1.748	-0.000254983\\
1.75	-0.000253465\\
1.752	-0.000250464\\
1.754	-0.000246014\\
1.756	-0.000240148\\
1.758	-0.000232909\\
1.76	-0.000224339\\
1.762	-0.000214486\\
1.764	-0.000203403\\
1.766	-0.000191143\\
1.768	-0.000177766\\
1.77	-0.000163331\\
1.772	-0.000147903\\
1.774	-0.000131548\\
1.776	-0.000126923\\
1.778	-0.000128136\\
1.78	-0.000128348\\
1.782	-0.000127572\\
1.784	-0.000125827\\
1.786	-0.000123133\\
1.788	-0.000119515\\
1.79	-0.000115\\
1.792	-0.000112392\\
1.794	-0.000125282\\
1.796	-0.000137124\\
1.798	-0.000147895\\
1.8	-0.000157574\\
1.802	-0.000166143\\
1.804	-0.000173589\\
1.806	-0.000179904\\
1.808	-0.000185082\\
1.81	-0.000189121\\
1.812	-0.000192024\\
1.814	-0.000193798\\
1.816	-0.000194453\\
1.818	-0.000194002\\
1.82	-0.000192462\\
1.822	-0.000189855\\
1.824	-0.000186204\\
1.826	-0.000191423\\
1.828	-0.000201007\\
1.83	-0.000209402\\
1.832	-0.000216597\\
1.834	-0.000222584\\
1.836	-0.000227362\\
1.838	-0.000230929\\
1.84	-0.000233293\\
1.842	-0.00023446\\
1.844	-0.000234444\\
1.846	-0.000233262\\
1.848	-0.000230932\\
1.85	-0.00022748\\
1.852	-0.00022607\\
1.854	-0.000225326\\
1.856	-0.000223476\\
1.858	-0.000220792\\
1.86	-0.000219593\\
1.862	-0.000217325\\
1.864	-0.00021401\\
1.866	-0.000209673\\
1.868	-0.000204342\\
1.87	-0.00019805\\
1.872	-0.000190831\\
1.874	-0.00018272\\
1.876	-0.000173758\\
1.878	-0.000163986\\
1.88	-0.000153448\\
1.882	-0.000142189\\
1.884	-0.000130257\\
1.886	-0.000117701\\
1.888	-0.000104572\\
1.89	-9.09203e-05\\
1.892	-7.68001e-05\\
1.894	-6.644e-05\\
1.896	-6.32004e-05\\
1.898	-5.93288e-05\\
1.9	-5.97215e-05\\
1.902	-6.50699e-05\\
1.904	-7.39157e-05\\
1.906	-8.19826e-05\\
1.908	-8.92551e-05\\
1.91	-9.57206e-05\\
1.912	-0.000101369\\
1.914	-0.000106195\\
1.916	-0.000110192\\
1.918	-0.000113361\\
1.92	-0.000115703\\
1.922	-0.000117222\\
1.924	-0.000117925\\
1.926	-0.000117822\\
1.928	-0.000116926\\
1.93	-0.000115251\\
1.932	-0.000112815\\
1.934	-0.000110715\\
1.936	-0.000118462\\
1.938	-0.000125352\\
1.94	-0.000131377\\
1.942	-0.000136529\\
1.944	-0.000140802\\
1.946	-0.000144197\\
1.948	-0.000146713\\
1.95	-0.000148356\\
1.952	-0.000149131\\
1.954	-0.000149051\\
1.956	-0.000148125\\
1.958	-0.000146371\\
1.96	-0.000143805\\
1.962	-0.000142775\\
1.964	-0.00014236\\
1.966	-0.000141144\\
1.968	-0.000139143\\
1.97	-0.00013823\\
1.972	-0.000136821\\
1.974	-0.000134653\\
1.976	-0.000131746\\
1.978	-0.000128937\\
1.98	-0.000128616\\
1.982	-0.000127707\\
1.984	-0.000126221\\
1.986	-0.000124169\\
1.988	-0.000121567\\
1.99	-0.000118429\\
1.992	-0.000114771\\
1.994	-0.000110613\\
1.996	-0.000105972\\
1.998	-0.000100869\\
2	-9.53267e-05\\
};
\addplot [color=mycolor7, line width=1.0pt, forget plot]
  table[row sep=crcr]{%
-0.024	0\\
-0.022	0\\
-0.02	0\\
-0.018	0\\
-0.016	0\\
-0.014	0\\
-0.012	0\\
-0.01	0\\
-0.008	0\\
-0.006	0\\
-0.004	0\\
-0.002	0\\
0	0\\
0.002	-1.43576e-08\\
0.004	-1.997e-07\\
0.006	-9.27297e-07\\
0.008	-2.72768e-06\\
0.01	-6.23441e-06\\
0.012	-1.21365e-05\\
0.014	-2.11387e-05\\
0.016	-3.3928e-05\\
0.018	-4.28913e-05\\
0.02	-6.1852e-05\\
0.022	-8.56835e-05\\
0.024	-0.00165765\\
0.026	-0.00838178\\
0.028	-0.00925358\\
0.03	-0.0100812\\
0.032	-0.01085\\
0.034	-0.0115464\\
0.036	-0.0121575\\
0.038	-0.0126717\\
0.04	-0.0130788\\
0.042	-0.0133699\\
0.044	-0.0135377\\
0.046	-0.0135768\\
0.048	-0.0134831\\
0.05	-0.0132546\\
0.052	-0.0128909\\
0.054	-0.0123935\\
0.056	-0.0117655\\
0.058	-0.0110119\\
0.06	-0.010139\\
0.062	-0.00915494\\
0.064	-0.0135118\\
0.066	-0.0139274\\
0.068	-0.0143027\\
0.07	-0.0146345\\
0.072	-0.01492\\
0.074	-0.0151566\\
0.076	-0.0153418\\
0.078	-0.0154735\\
0.08	-0.0155499\\
0.082	-0.0155692\\
0.084	-0.0161066\\
0.086	-0.0166303\\
0.088	-0.0171043\\
0.09	-0.0175255\\
0.092	-0.0178934\\
0.094	-0.0182953\\
0.096	-0.0186427\\
0.098	-0.0190855\\
0.1	-0.0211988\\
0.102	-0.0232128\\
0.104	-0.0251174\\
0.106	-0.0269036\\
0.108	-0.0285629\\
0.11	-0.0300874\\
0.112	-0.0286627\\
0.114	-0.0267282\\
0.116	-0.024696\\
0.118	-0.0225712\\
0.12	-0.0203595\\
0.122	-0.0179976\\
0.124	-0.0154666\\
0.126	-0.0134138\\
0.128	-0.0126372\\
0.13	-0.011843\\
0.132	-0.0115767\\
0.134	-0.0112892\\
0.136	-0.0109804\\
0.138	-0.01065\\
0.14	-0.0102976\\
0.142	-0.00992315\\
0.144	-0.00952646\\
0.146	-0.00910783\\
0.148	-0.00873742\\
0.15	-0.00833703\\
0.152	-0.00790772\\
0.154	-0.00745066\\
0.156	-0.00696717\\
0.158	-0.0064587\\
0.16	-0.0059268\\
0.162	-0.00537315\\
0.164	-0.00479953\\
0.166	-0.00537696\\
0.168	-0.00845325\\
0.17	-0.0112403\\
0.172	-0.0137318\\
0.174	-0.0150929\\
0.176	-0.0151117\\
0.178	-0.0150115\\
0.18	-0.014795\\
0.182	-0.0144658\\
0.184	-0.0140285\\
0.186	-0.0134885\\
0.188	-0.0128521\\
0.19	-0.0121262\\
0.192	-0.0113186\\
0.194	-0.0104375\\
0.196	-0.00949153\\
0.198	-0.00848994\\
0.2	-0.00744214\\
0.202	-0.00615507\\
0.204	-0.00549349\\
0.206	-0.00733967\\
0.208	-0.00904854\\
0.21	-0.0106107\\
0.212	-0.0117439\\
0.214	-0.0121315\\
0.216	-0.0124415\\
0.218	-0.0126717\\
0.22	-0.0128206\\
0.222	-0.0128872\\
0.224	-0.0128709\\
0.226	-0.012772\\
0.228	-0.0125911\\
0.23	-0.0123295\\
0.232	-0.0119891\\
0.234	-0.0115721\\
0.236	-0.0110815\\
0.238	-0.0105206\\
0.24	-0.00989327\\
0.242	-0.00920379\\
0.244	-0.00845685\\
0.246	-0.00765749\\
0.248	-0.00681113\\
0.25	-0.00592343\\
0.252	-0.00500035\\
0.254	-0.00496917\\
0.256	-0.00704142\\
0.258	-0.00902036\\
0.26	-0.0109004\\
0.262	-0.0126762\\
0.264	-0.0143432\\
0.266	-0.0154967\\
0.268	-0.0161254\\
0.27	-0.0166918\\
0.272	-0.0171938\\
0.274	-0.0176301\\
0.276	-0.0179992\\
0.278	-0.0183\\
0.28	-0.0185319\\
0.282	-0.0186943\\
0.284	-0.0187868\\
0.286	-0.0188096\\
0.288	-0.0187629\\
0.29	-0.0186472\\
0.292	-0.0184632\\
0.294	-0.0185692\\
0.296	-0.0191986\\
0.298	-0.0197587\\
0.3	-0.0202478\\
0.302	-0.0209354\\
0.304	-0.0217\\
0.306	-0.0223699\\
0.308	-0.0229436\\
0.31	-0.0234197\\
0.312	-0.0237975\\
0.314	-0.0240766\\
0.316	-0.024257\\
0.318	-0.0243389\\
0.32	-0.0243234\\
0.322	-0.0242115\\
0.324	-0.0240049\\
0.326	-0.0237054\\
0.328	-0.0232724\\
0.33	-0.0221058\\
0.332	-0.0208704\\
0.334	-0.0195713\\
0.336	-0.0182141\\
0.338	-0.0170551\\
0.34	-0.0169829\\
0.342	-0.0168606\\
0.344	-0.0162804\\
0.346	-0.0156498\\
0.348	-0.0149806\\
0.35	-0.0142751\\
0.352	-0.0135354\\
0.354	-0.0127636\\
0.356	-0.0123423\\
0.358	-0.0122768\\
0.36	-0.0121731\\
0.362	-0.0120317\\
0.364	-0.0118534\\
0.366	-0.0116389\\
0.368	-0.0113892\\
0.37	-0.0111053\\
0.372	-0.0107884\\
0.374	-0.0104396\\
0.376	-0.0100603\\
0.378	-0.00965192\\
0.38	-0.00921582\\
0.382	-0.0087536\\
0.384	-0.00826686\\
0.386	-0.00775728\\
0.388	-0.00709516\\
0.39	-0.00647929\\
0.392	-0.00650222\\
0.394	-0.00648815\\
0.396	-0.00643752\\
0.398	-0.00635089\\
0.4	-0.00622887\\
0.402	-0.00607215\\
0.404	-0.00588151\\
0.406	-0.00615158\\
0.408	-0.00659267\\
0.41	-0.00694514\\
0.412	-0.00720933\\
0.414	-0.00738585\\
0.416	-0.00747559\\
0.418	-0.00747971\\
0.42	-0.00741709\\
0.422	-0.00736258\\
0.424	-0.00725965\\
0.426	-0.00707701\\
0.428	-0.0068166\\
0.43	-0.0064806\\
0.432	-0.0060714\\
0.434	-0.00607085\\
0.436	-0.00618272\\
0.438	-0.0062623\\
0.44	-0.00630994\\
0.442	-0.00632603\\
0.444	-0.00631106\\
0.446	-0.00626559\\
0.448	-0.00619023\\
0.45	-0.00608565\\
0.452	-0.00595261\\
0.454	-0.00579191\\
0.456	-0.0056044\\
0.458	-0.00539098\\
0.46	-0.00515263\\
0.462	-0.00489033\\
0.464	-0.00460513\\
0.466	-0.00429812\\
0.468	-0.00397041\\
0.47	-0.00362315\\
0.472	-0.00392933\\
0.474	-0.00471537\\
0.476	-0.00544602\\
0.478	-0.00611944\\
0.48	-0.00673406\\
0.482	-0.0072885\\
0.484	-0.00778165\\
0.486	-0.00821264\\
0.488	-0.00858082\\
0.49	-0.0088858\\
0.492	-0.00912743\\
0.494	-0.00930576\\
0.496	-0.00942113\\
0.498	-0.00947405\\
0.5	-0.00946529\\
0.502	-0.00939581\\
0.504	-0.00926681\\
0.506	-0.00907968\\
0.508	-0.00909594\\
0.51	-0.00958244\\
0.512	-0.010008\\
0.514	-0.0103722\\
0.516	-0.0106746\\
0.518	-0.0109151\\
0.52	-0.0110938\\
0.522	-0.0112109\\
0.524	-0.0112671\\
0.526	-0.0112629\\
0.528	-0.0111992\\
0.53	-0.0110936\\
0.532	-0.0110013\\
0.534	-0.010891\\
0.536	-0.0107252\\
0.538	-0.0105051\\
0.54	-0.0102323\\
0.542	-0.00990847\\
0.544	-0.00953532\\
0.546	-0.00911482\\
0.548	-0.00864903\\
0.55	-0.00814012\\
0.552	-0.00759039\\
0.554	-0.00700223\\
0.556	-0.00637812\\
0.558	-0.0058952\\
0.56	-0.00604098\\
0.562	-0.00622588\\
0.564	-0.00637284\\
0.566	-0.00648198\\
0.568	-0.00655353\\
0.57	-0.00658783\\
0.572	-0.00658531\\
0.574	-0.00654651\\
0.576	-0.00647206\\
0.578	-0.00636269\\
0.58	-0.0062192\\
0.582	-0.0060425\\
0.584	-0.00583357\\
0.586	-0.00559348\\
0.588	-0.00532336\\
0.59	-0.00502443\\
0.592	-0.00469795\\
0.594	-0.00434527\\
0.596	-0.00396778\\
0.598	-0.00356693\\
0.6	-0.00314422\\
0.602	-0.00283964\\
0.604	-0.00268396\\
0.606	-0.00250629\\
0.608	-0.00230751\\
0.61	-0.00208859\\
0.612	-0.00185052\\
0.614	-0.00159433\\
0.616	-0.00155546\\
0.618	-0.00198503\\
0.62	-0.00237181\\
0.622	-0.0027153\\
0.624	-0.00301513\\
0.626	-0.00327107\\
0.628	-0.00348304\\
0.63	-0.00365112\\
0.632	-0.00377552\\
0.634	-0.00385657\\
0.636	-0.00389477\\
0.638	-0.00389072\\
0.64	-0.00385838\\
0.642	-0.00384017\\
0.644	-0.00381836\\
0.646	-0.00375927\\
0.648	-0.00366124\\
0.65	-0.00352527\\
0.652	-0.00335244\\
0.654	-0.00314396\\
0.656	-0.00290112\\
0.658	-0.00262531\\
0.66	-0.00231801\\
0.662	-0.00211643\\
0.664	-0.00202775\\
0.666	-0.00192213\\
0.668	-0.0018001\\
0.67	-0.0016622\\
0.672	-0.001509\\
0.674	-0.00134111\\
0.676	-0.00115918\\
0.678	-0.000985374\\
0.68	-0.00114645\\
0.682	-0.0012804\\
0.684	-0.00146087\\
0.686	-0.0018049\\
0.688	-0.00212075\\
0.69	-0.00240777\\
0.692	-0.00266541\\
0.694	-0.00289325\\
0.696	-0.00309096\\
0.698	-0.00325835\\
0.7	-0.00339533\\
0.702	-0.00350192\\
0.704	-0.00357826\\
0.706	-0.00362457\\
0.708	-0.00364121\\
0.71	-0.00362863\\
0.712	-0.00358736\\
0.714	-0.00351804\\
0.716	-0.00342143\\
0.718	-0.00329832\\
0.72	-0.00314963\\
0.722	-0.00303196\\
0.724	-0.00335072\\
0.726	-0.00366515\\
0.728	-0.00395023\\
0.73	-0.00420549\\
0.732	-0.00443052\\
0.734	-0.00462505\\
0.736	-0.00478887\\
0.738	-0.00492192\\
0.74	-0.0050242\\
0.742	-0.00509583\\
0.744	-0.00513702\\
0.746	-0.00514809\\
0.748	-0.00512943\\
0.75	-0.0050899\\
0.752	-0.00506983\\
0.754	-0.00504315\\
0.756	-0.00499884\\
0.758	-0.00492677\\
0.76	-0.00482758\\
0.762	-0.00470198\\
0.764	-0.00455077\\
0.766	-0.0043748\\
0.768	-0.00417501\\
0.77	-0.0039524\\
0.772	-0.00370803\\
0.774	-0.00344302\\
0.776	-0.00315853\\
0.778	-0.00285578\\
0.78	-0.00264615\\
0.782	-0.00265786\\
0.784	-0.0026508\\
0.786	-0.00262528\\
0.788	-0.00258167\\
0.79	-0.00252039\\
0.792	-0.00265018\\
0.794	-0.00278508\\
0.796	-0.00290081\\
0.798	-0.00299732\\
0.8	-0.00307461\\
0.802	-0.00313274\\
0.804	-0.00317181\\
0.806	-0.003192\\
0.808	-0.00319352\\
0.81	-0.00317664\\
0.812	-0.00314168\\
0.814	-0.00308899\\
0.816	-0.00301901\\
0.818	-0.00293217\\
0.82	-0.00282899\\
0.822	-0.00271\\
0.824	-0.00257579\\
0.826	-0.00242697\\
0.828	-0.00226419\\
0.83	-0.00208815\\
0.832	-0.00189954\\
0.834	-0.00169912\\
0.836	-0.00148766\\
0.838	-0.00126593\\
0.84	-0.00142799\\
0.842	-0.00160378\\
0.844	-0.00175753\\
0.846	-0.00188911\\
0.848	-0.00199845\\
0.85	-0.00208554\\
0.852	-0.00215047\\
0.854	-0.00219338\\
0.856	-0.00221449\\
0.858	-0.00221409\\
0.86	-0.00219598\\
0.862	-0.0021997\\
0.864	-0.00219291\\
0.866	-0.00218189\\
0.868	-0.00215061\\
0.87	-0.00209952\\
0.872	-0.00202911\\
0.874	-0.00193992\\
0.876	-0.00183255\\
0.878	-0.00170767\\
0.88	-0.00156596\\
0.882	-0.00140818\\
0.884	-0.00123511\\
0.886	-0.00104759\\
0.888	-0.000846479\\
0.89	-0.000808542\\
0.892	-0.000893755\\
0.894	-0.000969252\\
0.896	-0.001035\\
0.898	-0.00109099\\
0.9	-0.00113725\\
0.902	-0.0011738\\
0.904	-0.00126316\\
0.906	-0.00135836\\
0.908	-0.00143873\\
0.91	-0.0015042\\
0.912	-0.00155481\\
0.914	-0.00159062\\
0.916	-0.00161173\\
0.918	-0.00161832\\
0.92	-0.00161061\\
0.922	-0.00158887\\
0.924	-0.0015534\\
0.926	-0.00150457\\
0.928	-0.00144278\\
0.93	-0.00136847\\
0.932	-0.00133629\\
0.934	-0.00149032\\
0.936	-0.00162963\\
0.938	-0.00175395\\
0.94	-0.0018631\\
0.942	-0.00195691\\
0.944	-0.00203531\\
0.946	-0.00209825\\
0.948	-0.00214576\\
0.95	-0.00217789\\
0.952	-0.00229211\\
0.954	-0.00239152\\
0.956	-0.00247502\\
0.958	-0.00254256\\
0.96	-0.00259414\\
0.962	-0.00262983\\
0.964	-0.00264973\\
0.966	-0.00265402\\
0.968	-0.0026429\\
0.97	-0.0026169\\
0.972	-0.0026128\\
0.974	-0.00259605\\
0.976	-0.00258365\\
0.978	-0.00255669\\
0.98	-0.00251546\\
0.982	-0.00246031\\
0.984	-0.00239162\\
0.986	-0.0023098\\
0.988	-0.00221533\\
0.99	-0.00210868\\
0.992	-0.00199039\\
0.994	-0.00186102\\
0.996	-0.00172115\\
0.998	-0.0015714\\
1	-0.0014124\\
1.002	-0.00124483\\
1.004	-0.00119594\\
1.006	-0.00121861\\
1.008	-0.00123222\\
1.01	-0.00123658\\
1.012	-0.00123178\\
1.014	-0.00121797\\
1.016	-0.00119534\\
1.018	-0.00116412\\
1.02	-0.00112452\\
1.022	-0.00113027\\
1.024	-0.00122072\\
1.026	-0.00130178\\
1.028	-0.00137337\\
1.03	-0.00143543\\
1.032	-0.00148792\\
1.034	-0.00153082\\
1.036	-0.00156416\\
1.038	-0.00158799\\
1.04	-0.00160237\\
1.042	-0.00160741\\
1.044	-0.00160323\\
1.046	-0.00158998\\
1.048	-0.00156783\\
1.05	-0.00153698\\
1.052	-0.00149765\\
1.054	-0.00145008\\
1.056	-0.00139454\\
1.058	-0.00133131\\
1.06	-0.00126068\\
1.062	-0.00118298\\
1.064	-0.00116993\\
1.066	-0.00123419\\
1.068	-0.00128687\\
1.07	-0.00132798\\
1.072	-0.00135756\\
1.074	-0.00137571\\
1.076	-0.00138253\\
1.078	-0.00137819\\
1.08	-0.00136287\\
1.082	-0.00136459\\
1.084	-0.00135711\\
1.086	-0.00135305\\
1.088	-0.00134124\\
1.09	-0.00131903\\
1.092	-0.00128665\\
1.094	-0.00124437\\
1.096	-0.00119249\\
1.098	-0.00113133\\
1.1	-0.00106123\\
1.102	-0.000982562\\
1.104	-0.00089573\\
1.106	-0.000801146\\
1.108	-0.000699249\\
1.11	-0.000590492\\
1.112	-0.000522601\\
1.114	-0.000578377\\
1.116	-0.000626827\\
1.118	-0.000667902\\
1.12	-0.000701585\\
1.122	-0.000727881\\
1.124	-0.000746823\\
1.126	-0.00075847\\
1.128	-0.000762907\\
1.13	-0.000760241\\
1.132	-0.000750607\\
1.134	-0.000734159\\
1.136	-0.000711076\\
1.138	-0.000681556\\
1.14	-0.00064582\\
1.142	-0.000655217\\
1.144	-0.000668103\\
1.146	-0.000677018\\
1.148	-0.000760017\\
1.15	-0.00083561\\
1.152	-0.000903656\\
1.154	-0.000964041\\
1.156	-0.00101668\\
1.158	-0.00106151\\
1.16	-0.0010985\\
1.162	-0.00112764\\
1.164	-0.00114896\\
1.166	-0.0011625\\
1.168	-0.00116834\\
1.17	-0.00121338\\
1.172	-0.00127244\\
1.174	-0.00132321\\
1.176	-0.00136565\\
1.178	-0.00139974\\
1.18	-0.00142548\\
1.182	-0.00144291\\
1.184	-0.0014521\\
1.186	-0.00145312\\
1.188	-0.0014461\\
1.19	-0.00143116\\
1.192	-0.00142909\\
1.194	-0.00142177\\
1.196	-0.00141432\\
1.198	-0.00140498\\
1.2	-0.0013881\\
1.202	-0.00136385\\
1.204	-0.00133241\\
1.206	-0.00129399\\
1.208	-0.00124881\\
1.21	-0.00119712\\
1.212	-0.00113918\\
1.214	-0.00107527\\
1.216	-0.00100569\\
1.218	-0.000930753\\
1.22	-0.000850777\\
1.222	-0.000766101\\
1.224	-0.000677074\\
1.226	-0.000584055\\
1.228	-0.000580144\\
1.23	-0.000598295\\
1.232	-0.000611767\\
1.234	-0.00062076\\
1.236	-0.000625228\\
1.238	-0.000625156\\
1.24	-0.000620616\\
1.242	-0.000611694\\
1.244	-0.00059849\\
1.246	-0.000581113\\
1.248	-0.000559686\\
1.25	-0.000534342\\
1.252	-0.000507047\\
1.254	-0.00056194\\
1.256	-0.00061229\\
1.258	-0.00065802\\
1.26	-0.000699065\\
1.262	-0.000735378\\
1.264	-0.000766923\\
1.266	-0.000793679\\
1.268	-0.00081564\\
1.27	-0.000832814\\
1.272	-0.00084522\\
1.274	-0.000852895\\
1.276	-0.000855884\\
1.278	-0.000854249\\
1.28	-0.000848064\\
1.282	-0.000837412\\
1.284	-0.000822393\\
1.286	-0.000803114\\
1.288	-0.000779695\\
1.29	-0.000798628\\
1.292	-0.000812437\\
1.294	-0.000820257\\
1.296	-0.000822153\\
1.298	-0.00081821\\
1.3	-0.000808528\\
1.302	-0.000808354\\
1.304	-0.000805507\\
1.306	-0.000800234\\
1.308	-0.000796653\\
1.31	-0.000787578\\
1.312	-0.000773121\\
1.314	-0.000753412\\
1.316	-0.000728595\\
1.318	-0.00069883\\
1.32	-0.000664291\\
1.322	-0.000625163\\
1.324	-0.000581647\\
1.326	-0.000533953\\
1.328	-0.000482304\\
1.33	-0.000426751\\
1.332	-0.000367492\\
1.334	-0.000359668\\
1.336	-0.000365786\\
1.338	-0.000368354\\
1.34	-0.000367423\\
1.342	-0.000363057\\
1.344	-0.000355331\\
1.346	-0.000344329\\
1.348	-0.000349253\\
1.35	-0.000364461\\
1.352	-0.000376063\\
1.354	-0.000384073\\
1.356	-0.000388519\\
1.358	-0.00038944\\
1.36	-0.000386889\\
1.362	-0.000388764\\
1.364	-0.000431965\\
1.366	-0.000471309\\
1.368	-0.000506722\\
1.37	-0.000538139\\
1.372	-0.000565512\\
1.374	-0.000588806\\
1.376	-0.000608003\\
1.378	-0.000623095\\
1.38	-0.00063409\\
1.382	-0.00064101\\
1.384	-0.00064389\\
1.386	-0.000642779\\
1.388	-0.000647008\\
1.39	-0.000681096\\
1.392	-0.000710863\\
1.394	-0.000736272\\
1.396	-0.000757301\\
1.398	-0.00077394\\
1.4	-0.000786196\\
1.402	-0.00079409\\
1.404	-0.000797657\\
1.406	-0.000796945\\
1.408	-0.000792016\\
1.41	-0.000782944\\
1.412	-0.00078087\\
1.414	-0.00077774\\
1.416	-0.000770948\\
1.418	-0.000767855\\
1.42	-0.000760781\\
1.422	-0.000749802\\
1.424	-0.000735011\\
1.426	-0.000716508\\
1.428	-0.000694405\\
1.43	-0.000668827\\
1.432	-0.000639904\\
1.434	-0.00060778\\
1.436	-0.000572604\\
1.438	-0.000534536\\
1.44	-0.000493741\\
1.442	-0.000450392\\
1.444	-0.000404669\\
1.446	-0.000356756\\
1.448	-0.000306844\\
1.45	-0.000264725\\
1.452	-0.00027877\\
1.454	-0.00029049\\
1.456	-0.000299885\\
1.458	-0.000306962\\
1.46	-0.000311817\\
1.462	-0.00031445\\
1.464	-0.000314826\\
1.466	-0.000312978\\
1.468	-0.000308948\\
1.47	-0.000302781\\
1.472	-0.000294532\\
1.474	-0.000284256\\
1.476	-0.000272019\\
1.478	-0.00025789\\
1.48	-0.000272519\\
1.482	-0.000290364\\
1.484	-0.000305226\\
1.486	-0.000317102\\
1.488	-0.000325997\\
1.49	-0.000331928\\
1.492	-0.000345707\\
1.494	-0.000365775\\
1.496	-0.00038347\\
1.498	-0.000398774\\
1.5	-0.000411677\\
1.502	-0.000422175\\
1.504	-0.000430272\\
1.506	-0.000435978\\
1.508	-0.000440009\\
1.51	-0.000449444\\
1.512	-0.000455716\\
1.514	-0.000458851\\
1.516	-0.000458886\\
1.518	-0.000455867\\
1.52	-0.000449849\\
1.522	-0.000448916\\
1.524	-0.0004479\\
1.526	-0.000443944\\
1.528	-0.000442064\\
1.53	-0.000438578\\
1.532	-0.000432256\\
1.534	-0.00042316\\
1.536	-0.000411363\\
1.538	-0.000396943\\
1.54	-0.000379989\\
1.542	-0.000360593\\
1.544	-0.000338858\\
1.546	-0.00031489\\
1.548	-0.000288804\\
1.55	-0.000260717\\
1.552	-0.000230755\\
1.554	-0.000199044\\
1.556	-0.000169329\\
1.558	-0.000180044\\
1.56	-0.000188849\\
1.562	-0.000195741\\
1.564	-0.000200727\\
1.566	-0.000203816\\
1.568	-0.000205027\\
1.57	-0.000204385\\
1.572	-0.00020192\\
1.574	-0.000197668\\
1.576	-0.000191674\\
1.578	-0.000205181\\
1.58	-0.000227659\\
1.582	-0.000248103\\
1.584	-0.000266474\\
1.586	-0.000282738\\
1.588	-0.000296872\\
1.59	-0.000308858\\
1.592	-0.000318687\\
1.594	-0.000326356\\
1.596	-0.000331871\\
1.598	-0.000335244\\
1.6	-0.000336495\\
1.602	-0.000335649\\
1.604	-0.00033274\\
1.606	-0.000327809\\
1.608	-0.000343135\\
1.61	-0.000360276\\
1.612	-0.000375145\\
1.614	-0.000387724\\
1.616	-0.000398003\\
1.618	-0.000405981\\
1.62	-0.000411664\\
1.622	-0.000415064\\
1.624	-0.000416203\\
1.626	-0.000415108\\
1.628	-0.000411815\\
1.63	-0.000406364\\
1.632	-0.000404638\\
1.634	-0.000403336\\
1.636	-0.000399916\\
1.638	-0.000397065\\
1.64	-0.000394296\\
1.642	-0.000389481\\
1.644	-0.000382664\\
1.646	-0.000373896\\
1.648	-0.000363233\\
1.65	-0.000350737\\
1.652	-0.000336476\\
1.654	-0.00032052\\
1.656	-0.000302948\\
1.658	-0.00028384\\
1.66	-0.00026328\\
1.662	-0.000241359\\
1.664	-0.000218167\\
1.666	-0.000193801\\
1.668	-0.000168357\\
1.67	-0.000141937\\
1.672	-0.000136914\\
1.674	-0.000138214\\
1.676	-0.000138605\\
1.678	-0.0001381\\
1.68	-0.000136716\\
1.682	-0.00013876\\
1.684	-0.000142955\\
1.686	-0.000146053\\
1.688	-0.000148063\\
1.69	-0.000148953\\
1.692	-0.00014874\\
1.694	-0.000151211\\
1.696	-0.000161869\\
1.698	-0.000170964\\
1.7	-0.000178489\\
1.702	-0.000184444\\
1.704	-0.000188832\\
1.706	-0.000191663\\
1.708	-0.000192949\\
1.71	-0.000192712\\
1.712	-0.000190974\\
1.714	-0.000187765\\
1.716	-0.000188342\\
1.718	-0.000202396\\
1.72	-0.000214812\\
1.722	-0.000225574\\
1.724	-0.000234671\\
1.726	-0.000242097\\
1.728	-0.000247851\\
1.73	-0.00025194\\
1.732	-0.000254374\\
1.734	-0.00025517\\
1.736	-0.000254348\\
1.738	-0.000251936\\
1.74	-0.000247964\\
1.742	-0.000246729\\
1.744	-0.000246126\\
1.746	-0.00024399\\
1.748	-0.000241509\\
1.75	-0.000240054\\
1.752	-0.000237116\\
1.754	-0.000232728\\
1.756	-0.000226925\\
1.758	-0.000219748\\
1.76	-0.000211239\\
1.762	-0.000201448\\
1.764	-0.000190426\\
1.766	-0.000178228\\
1.768	-0.000164911\\
1.77	-0.000150537\\
1.772	-0.000135169\\
1.774	-0.000118873\\
1.776	-0.000114307\\
1.778	-0.000115579\\
1.78	-0.00011585\\
1.782	-0.000115133\\
1.784	-0.000113446\\
1.786	-0.000110811\\
1.788	-0.00010725\\
1.79	-0.000102793\\
1.792	-0.000100242\\
1.794	-0.000113189\\
1.796	-0.000125088\\
1.798	-0.000135916\\
1.8	-0.000145651\\
1.802	-0.000154276\\
1.804	-0.000161778\\
1.806	-0.000168148\\
1.808	-0.000173381\\
1.81	-0.000177475\\
1.812	-0.000180433\\
1.814	-0.000182262\\
1.816	-0.000182971\\
1.818	-0.000182574\\
1.82	-0.000181088\\
1.822	-0.000178535\\
1.824	-0.000174938\\
1.826	-0.00018021\\
1.828	-0.000189847\\
1.83	-0.000198294\\
1.832	-0.000205542\\
1.834	-0.000211582\\
1.836	-0.000216411\\
1.838	-0.000220031\\
1.84	-0.000222446\\
1.842	-0.000223665\\
1.844	-0.000223701\\
1.846	-0.000222569\\
1.848	-0.000220291\\
1.85	-0.000216889\\
1.852	-0.00021553\\
1.854	-0.000214837\\
1.856	-0.000213037\\
1.858	-0.000210402\\
1.86	-0.000209253\\
1.862	-0.000207035\\
1.864	-0.000203769\\
1.866	-0.000199481\\
1.868	-0.0001942\\
1.87	-0.000187956\\
1.872	-0.000180785\\
1.874	-0.000172723\\
1.876	-0.00016381\\
1.878	-0.000154086\\
1.88	-0.000143595\\
1.882	-0.000132384\\
1.884	-0.0001205\\
1.886	-0.000107991\\
1.888	-9.49082e-05\\
1.89	-8.13036e-05\\
1.892	-6.723e-05\\
1.894	-5.69163e-05\\
1.896	-5.37229e-05\\
1.898	-4.98974e-05\\
1.9	-5.03359e-05\\
1.902	-5.57299e-05\\
1.904	-6.46212e-05\\
1.906	-7.27333e-05\\
1.908	-8.00509e-05\\
1.91	-8.65612e-05\\
1.912	-9.22547e-05\\
1.914	-9.71244e-05\\
1.916	-0.000101166\\
1.918	-0.000104379\\
1.92	-0.000106765\\
1.922	-0.000108327\\
1.924	-0.000109074\\
1.926	-0.000109014\\
1.928	-0.000108161\\
1.93	-0.000106529\\
1.932	-0.000104136\\
1.934	-0.000102079\\
1.936	-0.000109868\\
1.938	-0.0001168\\
1.94	-0.000122867\\
1.942	-0.00012806\\
1.944	-0.000132376\\
1.946	-0.000135811\\
1.948	-0.000138369\\
1.95	-0.000140052\\
1.952	-0.000140869\\
1.954	-0.000140828\\
1.956	-0.000139944\\
1.958	-0.000138229\\
1.96	-0.000135704\\
1.962	-0.000134714\\
1.964	-0.000134338\\
1.966	-0.000133161\\
1.968	-0.0001312\\
1.97	-0.000130326\\
1.972	-0.000128955\\
1.974	-0.000126826\\
1.976	-0.000123958\\
1.978	-0.000121188\\
1.98	-0.000120904\\
1.982	-0.000120033\\
1.984	-0.000118585\\
1.986	-0.000116571\\
1.988	-0.000114006\\
1.99	-0.000110905\\
1.992	-0.000107284\\
1.994	-0.000103163\\
1.996	-9.85587e-05\\
1.998	-9.34927e-05\\
2	-8.79863e-05\\
};
\addplot [color=mycolor8, line width=1.0pt, forget plot]
  table[row sep=crcr]{%
-0.024	0\\
-0.022	0\\
-0.02	0\\
-0.018	0\\
-0.016	0\\
-0.014	0\\
-0.012	0\\
-0.01	0\\
-0.008	0\\
-0.006	0\\
-0.004	0\\
-0.002	0\\
0	0\\
0.002	-2.14629e-08\\
0.004	-2.77985e-07\\
0.006	-1.25511e-06\\
0.008	-3.6368e-06\\
0.01	-8.2357e-06\\
0.012	-1.59347e-05\\
0.014	-2.76383e-05\\
0.016	-4.42321e-05\\
0.018	-5.81629e-05\\
0.02	-8.37077e-05\\
0.022	-0.000115792\\
0.024	-0.00169783\\
0.026	-0.00836309\\
0.028	-0.00923591\\
0.03	-0.0100657\\
0.032	-0.0108382\\
0.034	-0.0115397\\
0.036	-0.0121578\\
0.038	-0.0126809\\
0.04	-0.0130991\\
0.042	-0.0134035\\
0.044	-0.0135871\\
0.046	-0.0136443\\
0.048	-0.0135713\\
0.05	-0.0133661\\
0.052	-0.0130283\\
0.054	-0.0125593\\
0.056	-0.0119621\\
0.058	-0.0112417\\
0.06	-0.0104042\\
0.062	-0.00945749\\
0.064	-0.0135889\\
0.066	-0.0140265\\
0.068	-0.0144255\\
0.07	-0.0147828\\
0.072	-0.0150952\\
0.074	-0.0153603\\
0.076	-0.0155753\\
0.078	-0.015738\\
0.08	-0.0158463\\
0.082	-0.0158985\\
0.084	-0.0164694\\
0.086	-0.0170271\\
0.088	-0.0175355\\
0.09	-0.017991\\
0.092	-0.0183933\\
0.094	-0.0188291\\
0.096	-0.01921\\
0.098	-0.0196856\\
0.1	-0.0218309\\
0.102	-0.0238757\\
0.104	-0.0258099\\
0.106	-0.0276243\\
0.108	-0.0293101\\
0.11	-0.0308596\\
0.112	-0.0297619\\
0.114	-0.0278512\\
0.116	-0.0258398\\
0.118	-0.023733\\
0.12	-0.0215362\\
0.122	-0.0192554\\
0.124	-0.0168973\\
0.126	-0.0148607\\
0.128	-0.0140833\\
0.13	-0.0132839\\
0.132	-0.0130077\\
0.134	-0.012706\\
0.136	-0.0123785\\
0.138	-0.012025\\
0.14	-0.0116455\\
0.142	-0.0112397\\
0.144	-0.0108079\\
0.146	-0.0103504\\
0.148	-0.00993761\\
0.15	-0.00949158\\
0.152	-0.00901358\\
0.154	-0.00850505\\
0.156	-0.00796755\\
0.158	-0.0074028\\
0.16	-0.00681262\\
0.162	-0.00619897\\
0.164	-0.0055639\\
0.166	-0.00607874\\
0.168	-0.00909153\\
0.17	-0.0118145\\
0.172	-0.0142415\\
0.174	-0.0149325\\
0.176	-0.0149378\\
0.178	-0.0148271\\
0.18	-0.014603\\
0.182	-0.0142689\\
0.184	-0.0138294\\
0.186	-0.0132898\\
0.188	-0.0126562\\
0.19	-0.0119353\\
0.192	-0.0111348\\
0.194	-0.0102626\\
0.196	-0.0093274\\
0.198	-0.008338\\
0.2	-0.00730371\\
0.202	-0.00623398\\
0.204	-0.0057355\\
0.206	-0.00756481\\
0.208	-0.00925736\\
0.21	-0.0108038\\
0.212	-0.0114706\\
0.214	-0.0118733\\
0.216	-0.0121998\\
0.218	-0.0124476\\
0.22	-0.0126152\\
0.222	-0.0127012\\
0.224	-0.0127052\\
0.226	-0.012627\\
0.228	-0.0124672\\
0.23	-0.012227\\
0.232	-0.0119079\\
0.234	-0.0115122\\
0.236	-0.0110426\\
0.238	-0.0105024\\
0.24	-0.00989517\\
0.242	-0.00922519\\
0.244	-0.00849698\\
0.246	-0.00771549\\
0.248	-0.00688602\\
0.25	-0.00601415\\
0.252	-0.00510575\\
0.254	-0.00499263\\
0.256	-0.00706571\\
0.258	-0.00904605\\
0.26	-0.010928\\
0.262	-0.0127062\\
0.264	-0.0143759\\
0.266	-0.0152131\\
0.268	-0.0158516\\
0.27	-0.0164284\\
0.272	-0.0169415\\
0.274	-0.0173892\\
0.276	-0.0177703\\
0.278	-0.0180837\\
0.28	-0.0183283\\
0.282	-0.0185037\\
0.284	-0.0186096\\
0.286	-0.0186458\\
0.288	-0.0186126\\
0.29	-0.0185104\\
0.292	-0.01834\\
0.294	-0.0184594\\
0.296	-0.0191021\\
0.298	-0.0196753\\
0.3	-0.0201772\\
0.302	-0.0208775\\
0.304	-0.0216542\\
0.306	-0.022336\\
0.308	-0.0229211\\
0.31	-0.0234082\\
0.312	-0.0237965\\
0.314	-0.0240855\\
0.316	-0.0242752\\
0.318	-0.024366\\
0.32	-0.0243586\\
0.322	-0.0242543\\
0.324	-0.0240545\\
0.326	-0.0237613\\
0.328	-0.0232437\\
0.33	-0.0220851\\
0.332	-0.0208574\\
0.334	-0.0195658\\
0.336	-0.0182156\\
0.338	-0.0170632\\
0.34	-0.0169972\\
0.342	-0.0168807\\
0.344	-0.0163699\\
0.346	-0.0157469\\
0.348	-0.0150845\\
0.35	-0.0143848\\
0.352	-0.01365\\
0.354	-0.0128823\\
0.356	-0.0124642\\
0.358	-0.012401\\
0.36	-0.0122988\\
0.362	-0.0121581\\
0.364	-0.0119796\\
0.366	-0.0117642\\
0.368	-0.0115129\\
0.37	-0.0112266\\
0.372	-0.0109065\\
0.374	-0.0105539\\
0.376	-0.0101701\\
0.378	-0.00975657\\
0.38	-0.00931476\\
0.382	-0.00884626\\
0.384	-0.00835271\\
0.386	-0.00783583\\
0.388	-0.00720813\\
0.39	-0.00658752\\
0.392	-0.00660573\\
0.394	-0.00658696\\
0.396	-0.00653169\\
0.398	-0.00644047\\
0.4	-0.00631391\\
0.402	-0.00615274\\
0.404	-0.00595773\\
0.406	-0.00622351\\
0.408	-0.00666042\\
0.41	-0.00700882\\
0.412	-0.00726905\\
0.414	-0.00744173\\
0.416	-0.00752776\\
0.418	-0.00752831\\
0.42	-0.00746226\\
0.422	-0.00740447\\
0.424	-0.0072984\\
0.426	-0.00711276\\
0.428	-0.00684952\\
0.43	-0.00651084\\
0.432	-0.00609912\\
0.434	-0.00609621\\
0.436	-0.00620588\\
0.438	-0.00628342\\
0.44	-0.00632916\\
0.442	-0.00634353\\
0.444	-0.00632698\\
0.446	-0.00628008\\
0.448	-0.00620344\\
0.45	-0.00609773\\
0.452	-0.0059637\\
0.454	-0.00580213\\
0.456	-0.00561389\\
0.458	-0.00539987\\
0.46	-0.00516102\\
0.462	-0.00489835\\
0.464	-0.00461288\\
0.466	-0.0043057\\
0.468	-0.00397791\\
0.47	-0.00363065\\
0.472	-0.00393692\\
0.474	-0.00472311\\
0.476	-0.00545398\\
0.478	-0.00612768\\
0.48	-0.00674261\\
0.482	-0.00729742\\
0.484	-0.00779097\\
0.486	-0.00822239\\
0.488	-0.00859102\\
0.49	-0.00889646\\
0.492	-0.00913855\\
0.494	-0.00931736\\
0.496	-0.00943318\\
0.498	-0.00948656\\
0.5	-0.00947823\\
0.502	-0.00940917\\
0.504	-0.00928055\\
0.506	-0.00909377\\
0.508	-0.00911035\\
0.51	-0.00959712\\
0.512	-0.0100229\\
0.514	-0.0103873\\
0.516	-0.0106898\\
0.518	-0.0109304\\
0.52	-0.0111091\\
0.522	-0.0112262\\
0.524	-0.0112822\\
0.526	-0.0112779\\
0.528	-0.011214\\
0.53	-0.011108\\
0.532	-0.0110154\\
0.534	-0.0109047\\
0.536	-0.0107384\\
0.538	-0.0105178\\
0.54	-0.0102444\\
0.542	-0.00991981\\
0.544	-0.00954592\\
0.546	-0.00912461\\
0.548	-0.00865796\\
0.55	-0.00814813\\
0.552	-0.00759743\\
0.554	-0.00700824\\
0.556	-0.00638306\\
0.558	-0.00589902\\
0.56	-0.00604364\\
0.562	-0.00622734\\
0.564	-0.00637305\\
0.566	-0.00648091\\
0.568	-0.00655116\\
0.57	-0.00658412\\
0.572	-0.00658025\\
0.574	-0.00654007\\
0.576	-0.00646423\\
0.578	-0.00635345\\
0.58	-0.00620854\\
0.582	-0.00603042\\
0.584	-0.00582008\\
0.586	-0.00557856\\
0.588	-0.00530703\\
0.59	-0.00500669\\
0.592	-0.00467882\\
0.594	-0.00432476\\
0.596	-0.00394591\\
0.598	-0.00354372\\
0.6	-0.0031197\\
0.602	-0.00281384\\
0.604	-0.00265691\\
0.606	-0.00247801\\
0.608	-0.00227806\\
0.61	-0.00205799\\
0.612	-0.00181881\\
0.614	-0.00156156\\
0.616	-0.00152166\\
0.618	-0.00195025\\
0.62	-0.00233611\\
0.622	-0.00267871\\
0.624	-0.0029777\\
0.626	-0.00323286\\
0.628	-0.0034441\\
0.63	-0.00361149\\
0.632	-0.00373525\\
0.634	-0.00381572\\
0.636	-0.00385338\\
0.638	-0.00384885\\
0.64	-0.00381607\\
0.642	-0.00379746\\
0.644	-0.00377531\\
0.646	-0.00371593\\
0.648	-0.00361766\\
0.65	-0.00348148\\
0.652	-0.00330849\\
0.654	-0.00309989\\
0.656	-0.00285698\\
0.658	-0.00258113\\
0.66	-0.00227383\\
0.662	-0.00207227\\
0.664	-0.00198365\\
0.666	-0.00187814\\
0.668	-0.00175623\\
0.67	-0.00161848\\
0.672	-0.00146546\\
0.674	-0.00129777\\
0.676	-0.00111606\\
0.678	-0.00094249\\
0.68	-0.00110382\\
0.682	-0.00123804\\
0.684	-0.0014188\\
0.686	-0.00176312\\
0.688	-0.00207928\\
0.69	-0.0023666\\
0.692	-0.00262456\\
0.694	-0.00285271\\
0.696	-0.00305074\\
0.698	-0.00321845\\
0.7	-0.00335574\\
0.702	-0.00346264\\
0.704	-0.00353928\\
0.706	-0.0035859\\
0.708	-0.00360283\\
0.71	-0.00359053\\
0.712	-0.00354954\\
0.714	-0.00348049\\
0.716	-0.00338412\\
0.718	-0.00326125\\
0.72	-0.00311279\\
0.722	-0.00299533\\
0.724	-0.00331429\\
0.726	-0.00362889\\
0.728	-0.00391415\\
0.73	-0.00416955\\
0.732	-0.00439472\\
0.734	-0.00458936\\
0.736	-0.00475329\\
0.738	-0.00488642\\
0.74	-0.00498877\\
0.742	-0.00506045\\
0.744	-0.00510168\\
0.746	-0.00511277\\
0.748	-0.00509412\\
0.75	-0.00505457\\
0.752	-0.00503448\\
0.754	-0.00500776\\
0.756	-0.0049634\\
0.758	-0.00489127\\
0.76	-0.004792\\
0.762	-0.00466631\\
0.764	-0.004515\\
0.766	-0.00433892\\
0.768	-0.00413901\\
0.77	-0.00391627\\
0.772	-0.00367176\\
0.774	-0.0034066\\
0.776	-0.00312195\\
0.778	-0.00281904\\
0.78	-0.00260924\\
0.782	-0.00262078\\
0.784	-0.00261355\\
0.786	-0.00258785\\
0.788	-0.00254405\\
0.79	-0.00248258\\
0.792	-0.00261218\\
0.794	-0.00274689\\
0.796	-0.00286243\\
0.798	-0.00295875\\
0.8	-0.00303585\\
0.802	-0.00309379\\
0.804	-0.00313267\\
0.806	-0.00315268\\
0.808	-0.00315402\\
0.81	-0.00313695\\
0.812	-0.00310181\\
0.814	-0.00304896\\
0.816	-0.0029788\\
0.818	-0.0028918\\
0.82	-0.00278846\\
0.822	-0.00266932\\
0.824	-0.00253496\\
0.826	-0.002386\\
0.828	-0.00222309\\
0.83	-0.00204691\\
0.832	-0.00185818\\
0.834	-0.00165765\\
0.836	-0.00144607\\
0.838	-0.00122425\\
0.84	-0.00138621\\
0.842	-0.00156192\\
0.844	-0.00171559\\
0.846	-0.0018471\\
0.848	-0.00195637\\
0.85	-0.00204341\\
0.852	-0.00210829\\
0.854	-0.00215117\\
0.856	-0.00217225\\
0.858	-0.00217182\\
0.86	-0.00215369\\
0.862	-0.00215741\\
0.864	-0.00215061\\
0.866	-0.0021396\\
0.868	-0.00210834\\
0.87	-0.00205727\\
0.872	-0.00198688\\
0.874	-0.00189773\\
0.876	-0.00179041\\
0.878	-0.00166557\\
0.88	-0.00152392\\
0.882	-0.0013662\\
0.884	-0.0011932\\
0.886	-0.00100575\\
0.888	-0.000804725\\
0.89	-0.000766875\\
0.892	-0.000852181\\
0.894	-0.000927775\\
0.896	-0.000993627\\
0.898	-0.00104973\\
0.9	-0.00109609\\
0.902	-0.00113276\\
0.904	-0.00122224\\
0.906	-0.00131757\\
0.908	-0.00139806\\
0.91	-0.00146367\\
0.912	-0.00151442\\
0.914	-0.00155036\\
0.916	-0.00157161\\
0.918	-0.00157835\\
0.92	-0.00157079\\
0.922	-0.00154919\\
0.924	-0.00151387\\
0.926	-0.00146519\\
0.928	-0.00140355\\
0.93	-0.00132939\\
0.932	-0.00129737\\
0.934	-0.00145155\\
0.936	-0.00159102\\
0.938	-0.0017155\\
0.94	-0.0018248\\
0.942	-0.00191877\\
0.944	-0.00199732\\
0.946	-0.00206042\\
0.948	-0.00210808\\
0.95	-0.00214037\\
0.952	-0.00225474\\
0.954	-0.0023543\\
0.956	-0.00243795\\
0.958	-0.00250564\\
0.96	-0.00255736\\
0.962	-0.0025932\\
0.964	-0.00261325\\
0.966	-0.00261768\\
0.968	-0.0026067\\
0.97	-0.00258084\\
0.972	-0.00257688\\
0.974	-0.00256026\\
0.976	-0.00254799\\
0.978	-0.00252116\\
0.98	-0.00248006\\
0.982	-0.00242504\\
0.984	-0.00235647\\
0.986	-0.00227478\\
0.988	-0.00218042\\
0.99	-0.00207388\\
0.992	-0.00195571\\
0.994	-0.00182645\\
0.996	-0.00168669\\
0.998	-0.00153705\\
1	-0.00137816\\
1.002	-0.00121069\\
1.004	-0.00116189\\
1.006	-0.00118467\\
1.008	-0.00119838\\
1.01	-0.00120283\\
1.012	-0.00119812\\
1.014	-0.0011844\\
1.016	-0.00116187\\
1.018	-0.00113073\\
1.02	-0.00109122\\
1.022	-0.00109706\\
1.024	-0.00118759\\
1.026	-0.00126873\\
1.028	-0.00134041\\
1.03	-0.00140255\\
1.032	-0.00145511\\
1.034	-0.00149809\\
1.036	-0.00153151\\
1.038	-0.00155542\\
1.04	-0.00156987\\
1.042	-0.00157499\\
1.044	-0.00157088\\
1.046	-0.0015577\\
1.048	-0.00153563\\
1.05	-0.00150485\\
1.052	-0.0014656\\
1.054	-0.0014181\\
1.056	-0.00136263\\
1.058	-0.00129947\\
1.06	-0.00122892\\
1.062	-0.00115129\\
1.064	-0.00113831\\
1.066	-0.00120265\\
1.068	-0.0012554\\
1.07	-0.00129658\\
1.072	-0.00132624\\
1.074	-0.00134446\\
1.076	-0.00135136\\
1.078	-0.00134709\\
1.08	-0.00133185\\
1.082	-0.00133364\\
1.084	-0.00132624\\
1.086	-0.00132226\\
1.088	-0.00131053\\
1.09	-0.0012884\\
1.092	-0.0012561\\
1.094	-0.0012139\\
1.096	-0.0011621\\
1.098	-0.00110102\\
1.1	-0.001031\\
1.102	-0.00095242\\
1.104	-0.000865672\\
1.106	-0.000771173\\
1.108	-0.000669361\\
1.11	-0.00056069\\
1.112	-0.000492886\\
1.114	-0.000548749\\
1.116	-0.000597287\\
1.118	-0.000638451\\
1.12	-0.000672223\\
1.122	-0.000698609\\
1.124	-0.000717641\\
1.126	-0.00072938\\
1.128	-0.000733908\\
1.13	-0.000731335\\
1.132	-0.000721794\\
1.134	-0.000705439\\
1.136	-0.00068245\\
1.138	-0.000653025\\
1.14	-0.000617383\\
1.142	-0.000626875\\
1.144	-0.000639857\\
1.146	-0.000648868\\
1.148	-0.000731962\\
1.15	-0.000807652\\
1.152	-0.000875794\\
1.154	-0.000936277\\
1.156	-0.000989011\\
1.158	-0.00103394\\
1.16	-0.00107103\\
1.162	-0.00110027\\
1.164	-0.00112169\\
1.166	-0.00113533\\
1.168	-0.00114126\\
1.17	-0.0011864\\
1.172	-0.00124555\\
1.174	-0.00129643\\
1.176	-0.00133896\\
1.178	-0.00137315\\
1.18	-0.00139899\\
1.182	-0.00141652\\
1.184	-0.0014258\\
1.186	-0.00142692\\
1.188	-0.00142\\
1.19	-0.00140516\\
1.192	-0.00140319\\
1.194	-0.00139597\\
1.196	-0.00138861\\
1.198	-0.00137937\\
1.2	-0.00136259\\
1.202	-0.00133843\\
1.204	-0.00130709\\
1.206	-0.00126876\\
1.208	-0.00122368\\
1.21	-0.00117208\\
1.212	-0.00111424\\
1.214	-0.00105042\\
1.216	-0.000980939\\
1.218	-0.000906092\\
1.22	-0.000826209\\
1.222	-0.000741626\\
1.224	-0.000652692\\
1.226	-0.000559766\\
1.228	-0.000555946\\
1.23	-0.00057419\\
1.232	-0.000587752\\
1.234	-0.000596836\\
1.236	-0.000601395\\
1.238	-0.000601412\\
1.24	-0.000596962\\
1.242	-0.000588129\\
1.244	-0.000575013\\
1.246	-0.000557725\\
1.248	-0.000536386\\
1.25	-0.000511128\\
1.252	-0.00048392\\
1.254	-0.0005389\\
1.256	-0.000589336\\
1.258	-0.000635151\\
1.26	-0.000676282\\
1.262	-0.000712679\\
1.264	-0.000744308\\
1.266	-0.000771148\\
1.268	-0.000793193\\
1.27	-0.000810449\\
1.272	-0.000822939\\
1.274	-0.000830695\\
1.276	-0.000833766\\
1.278	-0.000832213\\
1.28	-0.000826108\\
1.282	-0.000815537\\
1.284	-0.000800597\\
1.286	-0.000781398\\
1.288	-0.000758059\\
1.29	-0.00077707\\
1.292	-0.000790957\\
1.294	-0.000798856\\
1.296	-0.00080083\\
1.298	-0.000796964\\
1.3	-0.000787359\\
1.302	-0.000787262\\
1.304	-0.000784491\\
1.306	-0.000779294\\
1.308	-0.000775789\\
1.31	-0.000766789\\
1.312	-0.000752014\\
1.314	-0.000731899\\
1.316	-0.000706684\\
1.318	-0.000676527\\
1.32	-0.000641602\\
1.322	-0.000602096\\
1.324	-0.000558207\\
1.326	-0.000510147\\
1.328	-0.000458139\\
1.33	-0.000402414\\
1.332	-0.000343213\\
1.334	-0.000335447\\
1.336	-0.000341623\\
1.338	-0.000344249\\
1.34	-0.000343377\\
1.342	-0.00033907\\
1.344	-0.000331402\\
1.346	-0.000320459\\
1.348	-0.000325441\\
1.35	-0.000340709\\
1.352	-0.000352369\\
1.354	-0.000360439\\
1.356	-0.000364944\\
1.358	-0.000365924\\
1.36	-0.000363433\\
1.362	-0.000365368\\
1.364	-0.000408628\\
1.366	-0.000448033\\
1.368	-0.000483505\\
1.37	-0.000514982\\
1.372	-0.000542416\\
1.374	-0.000565771\\
1.376	-0.000585028\\
1.378	-0.000600181\\
1.38	-0.000611237\\
1.382	-0.000618218\\
1.384	-0.00062116\\
1.386	-0.00062011\\
1.388	-0.000624401\\
1.39	-0.000658551\\
1.392	-0.00068838\\
1.394	-0.000713851\\
1.396	-0.000734941\\
1.398	-0.000751643\\
1.4	-0.000763962\\
1.402	-0.000771919\\
1.404	-0.000775548\\
1.406	-0.000774899\\
1.408	-0.000770032\\
1.41	-0.000761024\\
1.412	-0.000759013\\
1.414	-0.000755947\\
1.416	-0.000749218\\
1.418	-0.000746189\\
1.42	-0.000739178\\
1.422	-0.000728264\\
1.424	-0.000713536\\
1.426	-0.000695097\\
1.428	-0.000673059\\
1.43	-0.000647544\\
1.432	-0.000618686\\
1.434	-0.000586626\\
1.436	-0.000551514\\
1.438	-0.00051351\\
1.44	-0.00047278\\
1.442	-0.000429495\\
1.444	-0.000383837\\
1.446	-0.000335989\\
1.448	-0.00028614\\
1.45	-0.000244086\\
1.452	-0.000258196\\
1.454	-0.00026998\\
1.456	-0.00027944\\
1.458	-0.000286581\\
1.46	-0.000291501\\
1.462	-0.000294198\\
1.464	-0.000294638\\
1.466	-0.000292854\\
1.468	-0.000288888\\
1.47	-0.000282786\\
1.472	-0.0002746\\
1.474	-0.000264389\\
1.476	-0.000252216\\
1.478	-0.000238151\\
1.48	-0.000252844\\
1.482	-0.000270752\\
1.484	-0.000285678\\
1.486	-0.000297617\\
1.488	-0.000306576\\
1.49	-0.000312571\\
1.492	-0.000326412\\
1.494	-0.000346544\\
1.496	-0.000364301\\
1.498	-0.000379668\\
1.5	-0.000392633\\
1.502	-0.000403194\\
1.504	-0.000411353\\
1.506	-0.000417122\\
1.508	-0.000421215\\
1.51	-0.000430713\\
1.512	-0.000437046\\
1.514	-0.000440243\\
1.516	-0.000440339\\
1.518	-0.000437381\\
1.52	-0.000431425\\
1.522	-0.000430553\\
1.524	-0.000429598\\
1.526	-0.000425702\\
1.528	-0.000423883\\
1.53	-0.000420458\\
1.532	-0.000414196\\
1.534	-0.000405161\\
1.536	-0.000393423\\
1.538	-0.000379064\\
1.54	-0.000362169\\
1.542	-0.000342833\\
1.544	-0.000321157\\
1.546	-0.000297249\\
1.548	-0.000271221\\
1.55	-0.000243194\\
1.552	-0.00021329\\
1.554	-0.000181638\\
1.556	-0.000151982\\
1.558	-0.000162755\\
1.56	-0.000171618\\
1.562	-0.000178569\\
1.564	-0.000183613\\
1.566	-0.00018676\\
1.568	-0.000188029\\
1.57	-0.000187445\\
1.572	-0.000185037\\
1.574	-0.000180843\\
1.576	-0.000174906\\
1.578	-0.000188471\\
1.58	-0.000211006\\
1.582	-0.000231507\\
1.584	-0.000249934\\
1.586	-0.000266256\\
1.588	-0.000280447\\
1.59	-0.000292489\\
1.592	-0.000302375\\
1.594	-0.000310101\\
1.596	-0.000315672\\
1.598	-0.000319101\\
1.6	-0.000320408\\
1.602	-0.000319619\\
1.604	-0.000316766\\
1.606	-0.000311891\\
1.608	-0.000327273\\
1.61	-0.00034447\\
1.612	-0.000359394\\
1.614	-0.000372029\\
1.616	-0.000382364\\
1.618	-0.000390397\\
1.62	-0.000396135\\
1.622	-0.000399591\\
1.624	-0.000400785\\
1.626	-0.000399746\\
1.628	-0.000396507\\
1.63	-0.000391111\\
1.632	-0.00038944\\
1.634	-0.000388194\\
1.636	-0.000384829\\
1.638	-0.000382032\\
1.64	-0.000379318\\
1.642	-0.000374557\\
1.644	-0.000367795\\
1.646	-0.000359081\\
1.648	-0.000348473\\
1.65	-0.000336031\\
1.652	-0.000321824\\
1.654	-0.000305922\\
1.656	-0.000288404\\
1.658	-0.00026935\\
1.66	-0.000248844\\
1.662	-0.000226976\\
1.664	-0.000203838\\
1.666	-0.000179525\\
1.668	-0.000154135\\
1.67	-0.000127768\\
1.672	-0.000122799\\
1.674	-0.000124152\\
1.676	-0.000124596\\
1.678	-0.000124144\\
1.68	-0.000122812\\
1.682	-0.000124908\\
1.684	-0.000129156\\
1.686	-0.000132307\\
1.688	-0.000134369\\
1.69	-0.000135312\\
1.692	-0.000135151\\
1.694	-0.000137673\\
1.696	-0.000148383\\
1.698	-0.000157529\\
1.7	-0.000165106\\
1.702	-0.000171113\\
1.704	-0.000175552\\
1.706	-0.000178434\\
1.708	-0.000179771\\
1.71	-0.000179585\\
1.712	-0.000177897\\
1.714	-0.000174739\\
1.716	-0.000175366\\
1.718	-0.000189471\\
1.72	-0.000201937\\
1.722	-0.000212749\\
1.724	-0.000221896\\
1.726	-0.000229371\\
1.728	-0.000235175\\
1.73	-0.000239314\\
1.732	-0.000241797\\
1.734	-0.000242642\\
1.736	-0.000241869\\
1.738	-0.000239506\\
1.74	-0.000235583\\
1.742	-0.000234395\\
1.744	-0.000233842\\
1.746	-0.000231754\\
1.748	-0.00022932\\
1.75	-0.000227913\\
1.752	-0.000225023\\
1.754	-0.000220683\\
1.756	-0.000214927\\
1.758	-0.000207797\\
1.76	-0.000199336\\
1.762	-0.000189592\\
1.764	-0.000178617\\
1.766	-0.000166465\\
1.768	-0.000153194\\
1.77	-0.000138866\\
1.772	-0.000123545\\
1.774	-0.000107295\\
1.776	-0.000102775\\
1.778	-0.000104093\\
1.78	-0.000104409\\
1.782	-0.000103738\\
1.784	-0.000102096\\
1.786	-9.95061e-05\\
1.788	-9.59909e-05\\
1.79	-9.15778e-05\\
1.792	-8.90721e-05\\
1.794	-0.000102063\\
1.796	-0.000114007\\
1.798	-0.000124879\\
1.8	-0.000134658\\
1.802	-0.000143327\\
1.804	-0.000150873\\
1.806	-0.000157287\\
1.808	-0.000162564\\
1.81	-0.000166702\\
1.812	-0.000169703\\
1.814	-0.000171575\\
1.816	-0.000172327\\
1.818	-0.000171973\\
1.82	-0.00017053\\
1.822	-0.000168019\\
1.824	-0.000164465\\
1.826	-0.00016978\\
1.828	-0.000179459\\
1.83	-0.000187949\\
1.832	-0.000195238\\
1.834	-0.00020132\\
1.836	-0.000206191\\
1.838	-0.000209853\\
1.84	-0.00021231\\
1.842	-0.00021357\\
1.844	-0.000213647\\
1.846	-0.000212557\\
1.848	-0.00021032\\
1.85	-0.000206959\\
1.852	-0.00020564\\
1.854	-0.000204988\\
1.856	-0.000203229\\
1.858	-0.000200634\\
1.86	-0.000199526\\
1.862	-0.000197348\\
1.864	-0.000194122\\
1.866	-0.000189874\\
1.868	-0.000184632\\
1.87	-0.000178429\\
1.872	-0.000171298\\
1.874	-0.000163275\\
1.876	-0.000154401\\
1.878	-0.000144716\\
1.88	-0.000134265\\
1.882	-0.000123093\\
1.884	-0.000111247\\
1.886	-9.87771e-05\\
1.888	-8.57332e-05\\
1.89	-7.21672e-05\\
1.892	-5.8132e-05\\
1.894	-4.78567e-05\\
1.896	-4.47015e-05\\
1.898	-4.09141e-05\\
1.9	-4.13905e-05\\
1.902	-4.68224e-05\\
1.904	-5.57514e-05\\
1.906	-6.39011e-05\\
1.908	-7.12561e-05\\
1.91	-7.78038e-05\\
1.912	-8.35345e-05\\
1.914	-8.84413e-05\\
1.916	-9.25201e-05\\
1.918	-9.57698e-05\\
1.92	-9.81921e-05\\
1.922	-9.97912e-05\\
1.924	-0.000100574\\
1.926	-0.000100551\\
1.928	-9.97343e-05\\
1.93	-9.81383e-05\\
1.932	-9.57808e-05\\
1.934	-9.37599e-05\\
1.936	-0.000101584\\
1.938	-0.000108553\\
1.94	-0.000114655\\
1.942	-0.000119883\\
1.944	-0.000124234\\
1.946	-0.000127704\\
1.948	-0.000130297\\
1.95	-0.000132015\\
1.952	-0.000132866\\
1.954	-0.000132861\\
1.956	-0.00013201\\
1.958	-0.00013033\\
1.96	-0.000127839\\
1.962	-0.000126883\\
1.964	-0.000126541\\
1.966	-0.000125398\\
1.968	-0.00012347\\
1.97	-0.00012263\\
1.972	-0.000121293\\
1.974	-0.000119197\\
1.976	-0.000116362\\
1.978	-0.000113625\\
1.98	-0.000113374\\
1.982	-0.000112536\\
1.984	-0.00011112\\
1.986	-0.000109139\\
1.988	-0.000106606\\
1.99	-0.000103537\\
1.992	-9.99492e-05\\
1.994	-9.58595e-05\\
1.996	-9.12874e-05\\
1.998	-8.62533e-05\\
2	-8.07785e-05\\
};
\addplot [color=mycolor9, line width=1.0pt, forget plot]
  table[row sep=crcr]{%
-0.024	0\\
-0.022	0\\
-0.02	0\\
-0.018	0\\
-0.016	0\\
-0.014	0\\
-0.012	0\\
-0.01	0\\
-0.008	0\\
-0.006	0\\
-0.004	0\\
-0.002	0\\
0	0\\
0.002	-6.27632e-08\\
0.004	-7.13565e-07\\
0.006	-3.01388e-06\\
0.008	-8.35023e-06\\
0.01	-1.82635e-05\\
0.012	-3.43144e-05\\
0.014	-5.79801e-05\\
0.016	-9.05761e-05\\
0.018	-0.000121777\\
0.02	-0.000171167\\
0.022	-0.000231404\\
0.024	-0.00184571\\
0.026	-0.00834713\\
0.028	-0.0092303\\
0.03	-0.0100737\\
0.032	-0.0108631\\
0.034	-0.011585\\
0.036	-0.012227\\
0.038	-0.0127774\\
0.04	-0.0132261\\
0.042	-0.0135643\\
0.044	-0.0137845\\
0.046	-0.0138809\\
0.048	-0.0138494\\
0.05	-0.0136875\\
0.052	-0.0133943\\
0.054	-0.0129709\\
0.056	-0.0124199\\
0.058	-0.0117455\\
0.06	-0.0109536\\
0.062	-0.0100513\\
0.064	-0.0138067\\
0.066	-0.0142743\\
0.068	-0.0147038\\
0.07	-0.0150917\\
0.072	-0.0154349\\
0.074	-0.0157306\\
0.076	-0.0159759\\
0.078	-0.0161684\\
0.08	-0.016306\\
0.082	-0.0163866\\
0.084	-0.016985\\
0.086	-0.0175691\\
0.088	-0.0181028\\
0.09	-0.0185825\\
0.092	-0.0190076\\
0.094	-0.0194649\\
0.096	-0.0198658\\
0.098	-0.0203599\\
0.1	-0.0225222\\
0.102	-0.0245825\\
0.104	-0.0265307\\
0.106	-0.0283574\\
0.108	-0.030054\\
0.11	-0.0316126\\
0.112	-0.0307275\\
0.114	-0.02883\\
0.116	-0.02683\\
0.118	-0.0247326\\
0.12	-0.0225433\\
0.122	-0.0202682\\
0.124	-0.0179138\\
0.126	-0.0160181\\
0.128	-0.0152351\\
0.13	-0.0144244\\
0.132	-0.0141345\\
0.134	-0.0138166\\
0.136	-0.0134705\\
0.138	-0.0130959\\
0.14	-0.012693\\
0.142	-0.0122617\\
0.144	-0.0118021\\
0.146	-0.0113148\\
0.148	-0.0108703\\
0.15	-0.0103906\\
0.152	-0.00987726\\
0.154	-0.00933174\\
0.156	-0.00875577\\
0.158	-0.00815117\\
0.16	-0.00751993\\
0.162	-0.00686413\\
0.164	-0.00618598\\
0.166	-0.00665693\\
0.168	-0.0096252\\
0.17	-0.0123031\\
0.172	-0.0146848\\
0.174	-0.0149052\\
0.176	-0.0149091\\
0.178	-0.0147981\\
0.18	-0.0145746\\
0.182	-0.0142422\\
0.184	-0.0138051\\
0.186	-0.0132686\\
0.188	-0.0126388\\
0.19	-0.0119225\\
0.192	-0.0111269\\
0.194	-0.0102602\\
0.196	-0.00933078\\
0.198	-0.00834747\\
0.2	-0.00731947\\
0.202	-0.00625618\\
0.204	-0.00602498\\
0.206	-0.00784502\\
0.208	-0.00952882\\
0.21	-0.0108646\\
0.212	-0.0113522\\
0.214	-0.0117671\\
0.216	-0.0121063\\
0.218	-0.0123674\\
0.22	-0.0125486\\
0.222	-0.0126485\\
0.224	-0.0126664\\
0.226	-0.0126021\\
0.228	-0.0124562\\
0.23	-0.0122295\\
0.232	-0.0119237\\
0.234	-0.0115408\\
0.236	-0.0110835\\
0.238	-0.010555\\
0.24	-0.00995885\\
0.242	-0.00929915\\
0.244	-0.0085804\\
0.246	-0.0078075\\
0.248	-0.00698568\\
0.25	-0.00612051\\
0.252	-0.0052178\\
0.254	-0.00520012\\
0.256	-0.00727629\\
0.258	-0.00925991\\
0.26	-0.0111453\\
0.262	-0.012927\\
0.264	-0.0143294\\
0.266	-0.0150353\\
0.268	-0.0156819\\
0.27	-0.0162671\\
0.272	-0.016789\\
0.274	-0.0172458\\
0.276	-0.0176361\\
0.278	-0.0179588\\
0.28	-0.0182129\\
0.282	-0.0183979\\
0.284	-0.0185132\\
0.286	-0.018559\\
0.288	-0.0185352\\
0.29	-0.0184423\\
0.292	-0.0182811\\
0.294	-0.0184094\\
0.296	-0.0190609\\
0.298	-0.0196426\\
0.3	-0.0201526\\
0.302	-0.0208607\\
0.304	-0.021645\\
0.306	-0.0223338\\
0.308	-0.0229256\\
0.31	-0.023419\\
0.312	-0.0238131\\
0.314	-0.0241074\\
0.316	-0.024302\\
0.318	-0.0243971\\
0.32	-0.0243936\\
0.322	-0.0242926\\
0.324	-0.0240957\\
0.326	-0.0238047\\
0.328	-0.0232298\\
0.33	-0.0220759\\
0.332	-0.0208525\\
0.334	-0.019565\\
0.336	-0.0182187\\
0.338	-0.0170698\\
0.34	-0.017007\\
0.342	-0.0168934\\
0.344	-0.0164235\\
0.346	-0.0158065\\
0.348	-0.0151495\\
0.35	-0.0144545\\
0.352	-0.0137237\\
0.354	-0.0129593\\
0.356	-0.0125438\\
0.358	-0.0124826\\
0.36	-0.0123815\\
0.362	-0.0122414\\
0.364	-0.0120628\\
0.366	-0.0118466\\
0.368	-0.0115938\\
0.37	-0.0113055\\
0.372	-0.0109829\\
0.374	-0.0106271\\
0.376	-0.0102397\\
0.378	-0.00982189\\
0.38	-0.00937537\\
0.382	-0.0089017\\
0.384	-0.00840254\\
0.386	-0.00787963\\
0.388	-0.00733477\\
0.39	-0.00684644\\
0.392	-0.00686139\\
0.394	-0.00683937\\
0.396	-0.00678085\\
0.398	-0.00668638\\
0.4	-0.00655661\\
0.402	-0.00639225\\
0.404	-0.00619408\\
0.406	-0.00645674\\
0.408	-0.00689058\\
0.41	-0.00723595\\
0.412	-0.00749321\\
0.414	-0.00766298\\
0.416	-0.00774616\\
0.418	-0.00774393\\
0.42	-0.00767517\\
0.422	-0.00761474\\
0.424	-0.00750611\\
0.426	-0.00731799\\
0.428	-0.00705234\\
0.43	-0.00671135\\
0.432	-0.00629739\\
0.434	-0.00629233\\
0.436	-0.00639993\\
0.438	-0.00647549\\
0.44	-0.00651933\\
0.442	-0.00653188\\
0.444	-0.00651361\\
0.446	-0.00646506\\
0.448	-0.00638684\\
0.45	-0.00627964\\
0.452	-0.00614419\\
0.454	-0.00598129\\
0.456	-0.00579177\\
0.458	-0.00557655\\
0.46	-0.00533657\\
0.462	-0.00507282\\
0.464	-0.00478634\\
0.466	-0.0044782\\
0.468	-0.00414951\\
0.47	-0.00380139\\
0.472	-0.00410685\\
0.474	-0.00489227\\
0.476	-0.0056224\\
0.478	-0.00629541\\
0.48	-0.00690967\\
0.482	-0.00746384\\
0.484	-0.00795677\\
0.486	-0.00838757\\
0.488	-0.00875561\\
0.49	-0.00906048\\
0.492	-0.00930199\\
0.494	-0.00948023\\
0.496	-0.00959548\\
0.498	-0.00964828\\
0.5	-0.00963937\\
0.502	-0.00956972\\
0.504	-0.00944049\\
0.506	-0.00925308\\
0.508	-0.00926902\\
0.51	-0.00975513\\
0.512	-0.0101803\\
0.514	-0.0105439\\
0.516	-0.0108457\\
0.518	-0.0110855\\
0.52	-0.0112634\\
0.522	-0.0113797\\
0.524	-0.0114348\\
0.526	-0.0114296\\
0.528	-0.0113647\\
0.53	-0.0112578\\
0.532	-0.0111642\\
0.534	-0.0110524\\
0.536	-0.010885\\
0.538	-0.0106633\\
0.54	-0.0103887\\
0.542	-0.0100629\\
0.544	-0.00968782\\
0.546	-0.00926525\\
0.548	-0.0087973\\
0.55	-0.00828614\\
0.552	-0.00773407\\
0.554	-0.00714349\\
0.556	-0.00651689\\
0.558	-0.0060314\\
0.56	-0.00617455\\
0.562	-0.00635675\\
0.564	-0.00650095\\
0.566	-0.00660728\\
0.568	-0.00667597\\
0.57	-0.00670737\\
0.572	-0.00670192\\
0.574	-0.00666015\\
0.576	-0.00658271\\
0.578	-0.00647033\\
0.58	-0.00632382\\
0.582	-0.00614409\\
0.584	-0.00593213\\
0.586	-0.00568902\\
0.588	-0.00541589\\
0.59	-0.00511395\\
0.592	-0.0047845\\
0.594	-0.00442887\\
0.596	-0.00404846\\
0.598	-0.00364474\\
0.6	-0.00321919\\
0.602	-0.00291182\\
0.604	-0.00275341\\
0.606	-0.00257305\\
0.608	-0.00237165\\
0.61	-0.00215017\\
0.612	-0.0019096\\
0.614	-0.00165099\\
0.616	-0.00160976\\
0.618	-0.00203705\\
0.62	-0.00242163\\
0.622	-0.00276299\\
0.624	-0.00306077\\
0.626	-0.00331474\\
0.628	-0.00352483\\
0.63	-0.00369111\\
0.632	-0.00381378\\
0.634	-0.00389319\\
0.636	-0.00392983\\
0.638	-0.00392431\\
0.64	-0.00389057\\
0.642	-0.00387104\\
0.644	-0.003848\\
0.646	-0.00378775\\
0.648	-0.00368864\\
0.65	-0.00355165\\
0.652	-0.00337788\\
0.654	-0.00316852\\
0.656	-0.00292487\\
0.658	-0.00264832\\
0.66	-0.00234034\\
0.662	-0.00213813\\
0.664	-0.00204887\\
0.666	-0.00194274\\
0.668	-0.00182024\\
0.67	-0.00168191\\
0.672	-0.00152833\\
0.674	-0.0013601\\
0.676	-0.00117785\\
0.678	-0.00100377\\
0.68	-0.00116459\\
0.682	-0.00129832\\
0.684	-0.0014786\\
0.686	-0.00182245\\
0.688	-0.00213813\\
0.69	-0.002425\\
0.692	-0.0026825\\
0.694	-0.00291021\\
0.696	-0.00310779\\
0.698	-0.00327506\\
0.7	-0.00341192\\
0.702	-0.00351838\\
0.704	-0.00359459\\
0.706	-0.00364076\\
0.708	-0.00365726\\
0.71	-0.00364452\\
0.712	-0.00360308\\
0.714	-0.00353359\\
0.716	-0.00343677\\
0.718	-0.00331346\\
0.72	-0.00316454\\
0.722	-0.00304662\\
0.724	-0.00336512\\
0.726	-0.00367926\\
0.728	-0.00396404\\
0.73	-0.00421897\\
0.732	-0.00444365\\
0.734	-0.00463781\\
0.736	-0.00480125\\
0.738	-0.00493388\\
0.74	-0.00503573\\
0.742	-0.0051069\\
0.744	-0.00514762\\
0.746	-0.0051582\\
0.748	-0.00513903\\
0.75	-0.00509896\\
0.752	-0.00507835\\
0.754	-0.00505109\\
0.756	-0.0050062\\
0.758	-0.00493353\\
0.76	-0.00483372\\
0.762	-0.00470749\\
0.764	-0.00455563\\
0.766	-0.004379\\
0.768	-0.00417855\\
0.77	-0.00395526\\
0.772	-0.0037102\\
0.774	-0.00344449\\
0.776	-0.0031593\\
0.778	-0.00285583\\
0.78	-0.00264549\\
0.782	-0.00265648\\
0.784	-0.0026487\\
0.786	-0.00262246\\
0.788	-0.00257812\\
0.79	-0.0025161\\
0.792	-0.00264517\\
0.794	-0.00277934\\
0.796	-0.00289435\\
0.798	-0.00299014\\
0.8	-0.00306672\\
0.802	-0.00312414\\
0.804	-0.00316251\\
0.806	-0.003182\\
0.808	-0.00318283\\
0.81	-0.00316527\\
0.812	-0.00312963\\
0.814	-0.00307628\\
0.816	-0.00300564\\
0.818	-0.00291816\\
0.82	-0.00281435\\
0.822	-0.00269474\\
0.824	-0.00255991\\
0.826	-0.00241049\\
0.828	-0.00224713\\
0.83	-0.00207051\\
0.832	-0.00188135\\
0.834	-0.00168038\\
0.836	-0.00146838\\
0.838	-0.00124614\\
0.84	-0.00140769\\
0.842	-0.00158298\\
0.844	-0.00173626\\
0.846	-0.00186738\\
0.848	-0.00197626\\
0.85	-0.00206292\\
0.852	-0.00212743\\
0.854	-0.00216993\\
0.856	-0.00219065\\
0.858	-0.00218988\\
0.86	-0.00217139\\
0.862	-0.00217477\\
0.864	-0.00216764\\
0.866	-0.0021563\\
0.868	-0.00212471\\
0.87	-0.00207333\\
0.872	-0.00200263\\
0.874	-0.00191317\\
0.876	-0.00180555\\
0.878	-0.00168042\\
0.88	-0.00153848\\
0.882	-0.00138048\\
0.884	-0.00120721\\
0.886	-0.00101949\\
0.888	-0.000818196\\
0.89	-0.000780085\\
0.892	-0.000865136\\
0.894	-0.00094048\\
0.896	-0.00100609\\
0.898	-0.00106194\\
0.9	-0.00110807\\
0.902	-0.00114451\\
0.904	-0.00123376\\
0.906	-0.00132887\\
0.908	-0.00140914\\
0.91	-0.00147454\\
0.912	-0.00152507\\
0.914	-0.0015608\\
0.916	-0.00158185\\
0.918	-0.00158839\\
0.92	-0.00158063\\
0.922	-0.00155884\\
0.924	-0.00152333\\
0.926	-0.00147446\\
0.928	-0.00141264\\
0.93	-0.0013383\\
0.932	-0.0013061\\
0.934	-0.0014601\\
0.936	-0.00159939\\
0.938	-0.0017237\\
0.94	-0.00183283\\
0.942	-0.00192664\\
0.944	-0.00200502\\
0.946	-0.00206796\\
0.948	-0.00211545\\
0.95	-0.00214758\\
0.952	-0.0022618\\
0.954	-0.0023612\\
0.956	-0.0024447\\
0.958	-0.00251223\\
0.96	-0.0025638\\
0.962	-0.00259949\\
0.964	-0.00261939\\
0.966	-0.00262367\\
0.968	-0.00261254\\
0.97	-0.00258653\\
0.972	-0.00258242\\
0.974	-0.00256565\\
0.976	-0.00255325\\
0.978	-0.00252627\\
0.98	-0.00248503\\
0.982	-0.00242986\\
0.984	-0.00236115\\
0.986	-0.00227932\\
0.988	-0.00218482\\
0.99	-0.00207814\\
0.992	-0.00195983\\
0.994	-0.00183043\\
0.996	-0.00169053\\
0.998	-0.00154075\\
1	-0.00138173\\
1.002	-0.00121412\\
1.004	-0.00116519\\
1.006	-0.00118783\\
1.008	-0.0012014\\
1.01	-0.00120572\\
1.012	-0.00120088\\
1.014	-0.00118703\\
1.016	-0.00116436\\
1.018	-0.00113309\\
1.02	-0.00109345\\
1.022	-0.00109916\\
1.024	-0.00118955\\
1.026	-0.00127057\\
1.028	-0.00134212\\
1.03	-0.00140413\\
1.032	-0.00145656\\
1.034	-0.00149942\\
1.036	-0.00153271\\
1.038	-0.00155649\\
1.04	-0.00157083\\
1.042	-0.00157582\\
1.044	-0.00157159\\
1.046	-0.00155829\\
1.048	-0.0015361\\
1.05	-0.0015052\\
1.052	-0.00146583\\
1.054	-0.00141822\\
1.056	-0.00136264\\
1.058	-0.00129936\\
1.06	-0.00122869\\
1.062	-0.00115095\\
1.064	-0.00113787\\
1.066	-0.00120209\\
1.068	-0.00125473\\
1.07	-0.00129581\\
1.072	-0.00132536\\
1.074	-0.00134348\\
1.076	-0.00135027\\
1.078	-0.00134591\\
1.08	-0.00133056\\
1.082	-0.00133226\\
1.084	-0.00132476\\
1.086	-0.00132068\\
1.088	-0.00130885\\
1.09	-0.00128663\\
1.092	-0.00125424\\
1.094	-0.00121195\\
1.096	-0.00116006\\
1.098	-0.00109889\\
1.1	-0.00102879\\
1.102	-0.000950121\\
1.104	-0.000863289\\
1.106	-0.000768707\\
1.108	-0.000666814\\
1.11	-0.000558064\\
1.112	-0.000490181\\
1.114	-0.000545967\\
1.116	-0.000594429\\
1.118	-0.000635519\\
1.12	-0.000669218\\
1.122	-0.000695532\\
1.124	-0.000714494\\
1.126	-0.000726163\\
1.128	-0.000730623\\
1.13	-0.000727984\\
1.132	-0.000718376\\
1.134	-0.000701958\\
1.136	-0.000678905\\
1.138	-0.000649418\\
1.14	-0.000613715\\
1.142	-0.000623148\\
1.144	-0.000636071\\
1.146	-0.000645025\\
1.148	-0.000728063\\
1.15	-0.000803697\\
1.152	-0.000871786\\
1.154	-0.000932215\\
1.156	-0.000984898\\
1.158	-0.00102977\\
1.16	-0.00106681\\
1.162	-0.00109601\\
1.164	-0.00111737\\
1.166	-0.00113097\\
1.168	-0.00113685\\
1.17	-0.00118195\\
1.172	-0.00124106\\
1.174	-0.00129189\\
1.176	-0.00133438\\
1.178	-0.00136853\\
1.18	-0.00139433\\
1.182	-0.00141182\\
1.184	-0.00142106\\
1.186	-0.00142214\\
1.188	-0.00141518\\
1.19	-0.0014003\\
1.192	-0.0013983\\
1.194	-0.00139104\\
1.196	-0.00138365\\
1.198	-0.00137437\\
1.2	-0.00135756\\
1.202	-0.00133337\\
1.204	-0.00130199\\
1.206	-0.00126364\\
1.208	-0.00121852\\
1.21	-0.0011669\\
1.212	-0.00110902\\
1.214	-0.00104518\\
1.216	-0.000975664\\
1.218	-0.000900789\\
1.22	-0.000820878\\
1.222	-0.000736268\\
1.224	-0.000647308\\
1.226	-0.000554355\\
1.228	-0.00055051\\
1.23	-0.000568728\\
1.232	-0.000582267\\
1.234	-0.000591326\\
1.236	-0.000595861\\
1.238	-0.000595855\\
1.24	-0.000591382\\
1.242	-0.000582526\\
1.244	-0.000569389\\
1.246	-0.000552078\\
1.248	-0.000530718\\
1.25	-0.000505439\\
1.252	-0.000478211\\
1.254	-0.00053317\\
1.256	-0.000583586\\
1.258	-0.000629382\\
1.26	-0.000670493\\
1.262	-0.000706871\\
1.264	-0.000738481\\
1.266	-0.000765303\\
1.268	-0.00078733\\
1.27	-0.000804569\\
1.272	-0.00081704\\
1.274	-0.000824779\\
1.276	-0.000827834\\
1.278	-0.000826264\\
1.28	-0.000820142\\
1.282	-0.000809556\\
1.284	-0.0007946\\
1.286	-0.000775386\\
1.288	-0.000752031\\
1.29	-0.000771028\\
1.292	-0.0007849\\
1.294	-0.000792784\\
1.296	-0.000794744\\
1.298	-0.000790584\\
1.3	-0.000780418\\
1.302	-0.000779766\\
1.304	-0.00077645\\
1.306	-0.000770714\\
1.308	-0.000766679\\
1.31	-0.000757156\\
1.312	-0.00074226\\
1.314	-0.000722118\\
1.316	-0.000696876\\
1.318	-0.000666693\\
1.32	-0.000631743\\
1.322	-0.000592211\\
1.324	-0.000548298\\
1.326	-0.000500215\\
1.328	-0.000448183\\
1.33	-0.000392435\\
1.332	-0.000333212\\
1.334	-0.000325425\\
1.336	-0.00033158\\
1.338	-0.000334185\\
1.34	-0.000333293\\
1.342	-0.000328966\\
1.344	-0.000321279\\
1.346	-0.000310318\\
1.348	-0.000315282\\
1.35	-0.000330533\\
1.352	-0.000342177\\
1.354	-0.00035023\\
1.356	-0.000354719\\
1.358	-0.000355685\\
1.36	-0.000353179\\
1.362	-0.0003551\\
1.364	-0.000398347\\
1.366	-0.000437739\\
1.368	-0.000473199\\
1.37	-0.000504664\\
1.372	-0.000532086\\
1.374	-0.000555431\\
1.376	-0.000574678\\
1.378	-0.000589821\\
1.38	-0.000600868\\
1.382	-0.000607841\\
1.384	-0.000610774\\
1.386	-0.000609717\\
1.388	-0.000614001\\
1.39	-0.000648144\\
1.392	-0.000677967\\
1.394	-0.000703433\\
1.396	-0.000724518\\
1.398	-0.000741215\\
1.4	-0.00075353\\
1.402	-0.000761484\\
1.404	-0.00076511\\
1.406	-0.000764458\\
1.408	-0.00075959\\
1.41	-0.00075058\\
1.412	-0.000748568\\
1.414	-0.000745501\\
1.416	-0.000738772\\
1.418	-0.000735744\\
1.42	-0.000728733\\
1.422	-0.00071782\\
1.424	-0.000703094\\
1.426	-0.000684657\\
1.428	-0.000662621\\
1.43	-0.000637109\\
1.432	-0.000608254\\
1.434	-0.000576198\\
1.436	-0.000541091\\
1.438	-0.000503091\\
1.44	-0.000462365\\
1.442	-0.000419086\\
1.444	-0.000373433\\
1.446	-0.000325591\\
1.448	-0.000275749\\
1.45	-0.000233702\\
1.452	-0.000247818\\
1.454	-0.00025961\\
1.456	-0.000269077\\
1.458	-0.000276226\\
1.46	-0.000281154\\
1.462	-0.00028386\\
1.464	-0.000284309\\
1.466	-0.000282534\\
1.468	-0.000278578\\
1.47	-0.000272485\\
1.472	-0.00026431\\
1.474	-0.000254109\\
1.476	-0.000241946\\
1.478	-0.000227891\\
1.48	-0.000242595\\
1.482	-0.000260515\\
1.484	-0.000275452\\
1.486	-0.000287403\\
1.488	-0.000296373\\
1.49	-0.00030238\\
1.492	-0.000316233\\
1.494	-0.000336377\\
1.496	-0.000354148\\
1.498	-0.000369527\\
1.5	-0.000382506\\
1.502	-0.000393079\\
1.504	-0.000401252\\
1.506	-0.000407034\\
1.508	-0.000411141\\
1.51	-0.000420652\\
1.512	-0.000427\\
1.514	-0.000430211\\
1.516	-0.000430321\\
1.518	-0.000427378\\
1.52	-0.000421436\\
1.522	-0.000420579\\
1.524	-0.000419639\\
1.526	-0.000415758\\
1.528	-0.000413954\\
1.53	-0.000410544\\
1.532	-0.000404298\\
1.534	-0.000395278\\
1.536	-0.000383556\\
1.538	-0.000369212\\
1.54	-0.000352333\\
1.542	-0.000333013\\
1.544	-0.000311353\\
1.546	-0.000287462\\
1.548	-0.000261451\\
1.55	-0.000233439\\
1.552	-0.000203552\\
1.554	-0.000171917\\
1.556	-0.000142278\\
1.558	-0.000153068\\
1.56	-0.000161948\\
1.562	-0.000168916\\
1.564	-0.000173977\\
1.566	-0.000177142\\
1.568	-0.000178428\\
1.57	-0.000177861\\
1.572	-0.000175472\\
1.574	-0.000171296\\
1.576	-0.000165377\\
1.578	-0.000178959\\
1.58	-0.000201513\\
1.582	-0.000222032\\
1.584	-0.000240478\\
1.586	-0.000256818\\
1.588	-0.000271027\\
1.59	-0.000283089\\
1.592	-0.000292994\\
1.594	-0.000300738\\
1.596	-0.000306329\\
1.598	-0.000309777\\
1.6	-0.000311103\\
1.602	-0.000310333\\
1.604	-0.000307501\\
1.606	-0.000302645\\
1.608	-0.000318047\\
1.61	-0.000335264\\
1.612	-0.000350208\\
1.614	-0.000362863\\
1.616	-0.000373218\\
1.618	-0.000381272\\
1.62	-0.00038703\\
1.622	-0.000390507\\
1.624	-0.000391722\\
1.626	-0.000390703\\
1.628	-0.000387485\\
1.63	-0.00038211\\
1.632	-0.00038046\\
1.634	-0.000379235\\
1.636	-0.000375891\\
1.638	-0.000373115\\
1.64	-0.000370423\\
1.642	-0.000365684\\
1.644	-0.000358943\\
1.646	-0.000350251\\
1.648	-0.000339664\\
1.65	-0.000327244\\
1.652	-0.000313059\\
1.654	-0.00029718\\
1.656	-0.000279684\\
1.658	-0.000260651\\
1.66	-0.000240168\\
1.662	-0.000218322\\
1.664	-0.000195207\\
1.666	-0.000170916\\
1.668	-0.000145549\\
1.67	-0.000119204\\
1.672	-0.000114257\\
1.674	-0.000115633\\
1.676	-0.0001161\\
1.678	-0.000115671\\
1.68	-0.000114362\\
1.682	-0.000116481\\
1.684	-0.000120752\\
1.686	-0.000123926\\
1.688	-0.00012601\\
1.69	-0.000126976\\
1.692	-0.000126838\\
1.694	-0.000129384\\
1.696	-0.000140117\\
1.698	-0.000149286\\
1.7	-0.000156887\\
1.702	-0.000162916\\
1.704	-0.000167379\\
1.706	-0.000170284\\
1.708	-0.000171645\\
1.71	-0.000171481\\
1.712	-0.000169818\\
1.714	-0.000166683\\
1.716	-0.000167333\\
1.718	-0.000181461\\
1.72	-0.000193951\\
1.722	-0.000204786\\
1.724	-0.000213956\\
1.726	-0.000221455\\
1.728	-0.000227282\\
1.73	-0.000231444\\
1.732	-0.000233951\\
1.734	-0.000234819\\
1.736	-0.00023407\\
1.738	-0.00023173\\
1.74	-0.00022783\\
1.742	-0.000226666\\
1.744	-0.000226135\\
1.746	-0.000224071\\
1.748	-0.00022166\\
1.75	-0.000220277\\
1.752	-0.00021741\\
1.754	-0.000213093\\
1.756	-0.000207361\\
1.758	-0.000200254\\
1.76	-0.000191816\\
1.762	-0.000182095\\
1.764	-0.000171143\\
1.766	-0.000159014\\
1.768	-0.000145767\\
1.77	-0.000131462\\
1.772	-0.000116164\\
1.774	-9.99373e-05\\
1.776	-9.54402e-05\\
1.778	-9.67815e-05\\
1.78	-9.71207e-05\\
1.782	-9.64722e-05\\
1.784	-9.48538e-05\\
1.786	-9.22864e-05\\
1.788	-8.87941e-05\\
1.79	-8.4404e-05\\
1.792	-8.19211e-05\\
1.794	-9.49351e-05\\
1.796	-0.000106902\\
1.798	-0.000117797\\
1.8	-0.000127599\\
1.802	-0.000136291\\
1.804	-0.000143859\\
1.806	-0.000150296\\
1.808	-0.000155595\\
1.81	-0.000159756\\
1.812	-0.00016278\\
1.814	-0.000164675\\
1.816	-0.000165449\\
1.818	-0.000165118\\
1.82	-0.000163698\\
1.822	-0.00016121\\
1.824	-0.000157678\\
1.826	-0.000163015\\
1.828	-0.000172717\\
1.83	-0.000181229\\
1.832	-0.000188541\\
1.834	-0.000194646\\
1.836	-0.000199539\\
1.838	-0.000203223\\
1.84	-0.000205702\\
1.842	-0.000206985\\
1.844	-0.000207084\\
1.846	-0.000206017\\
1.848	-0.000203802\\
1.85	-0.000200463\\
1.852	-0.000199167\\
1.854	-0.000198537\\
1.856	-0.0001968\\
1.858	-0.000194228\\
1.86	-0.000193141\\
1.862	-0.000190985\\
1.864	-0.000187782\\
1.866	-0.000183556\\
1.868	-0.000178336\\
1.87	-0.000172155\\
1.872	-0.000165045\\
1.874	-0.000157045\\
1.876	-0.000148193\\
1.878	-0.00013853\\
1.88	-0.0001281\\
1.882	-0.00011695\\
1.884	-0.000105126\\
1.886	-9.26781e-05\\
1.888	-7.9656e-05\\
1.89	-6.61117e-05\\
1.892	-5.20982e-05\\
1.894	-4.18445e-05\\
1.896	-3.8711e-05\\
1.898	-3.49451e-05\\
1.9	-3.54431e-05\\
1.902	-4.08965e-05\\
1.904	-4.98469e-05\\
1.906	-5.80181e-05\\
1.908	-6.53945e-05\\
1.91	-7.19635e-05\\
1.912	-7.77155e-05\\
1.914	-8.26436e-05\\
1.916	-8.67437e-05\\
1.918	-9.00146e-05\\
1.92	-9.2458e-05\\
1.922	-9.40783e-05\\
1.924	-9.48825e-05\\
1.926	-9.48804e-05\\
1.928	-9.40844e-05\\
1.93	-9.25094e-05\\
1.932	-9.01727e-05\\
1.934	-8.81727e-05\\
1.936	-9.60177e-05\\
1.938	-0.000103007\\
1.94	-0.00010913\\
1.942	-0.000114379\\
1.944	-0.00011875\\
1.946	-0.000122241\\
1.948	-0.000124854\\
1.95	-0.000126593\\
1.952	-0.000127465\\
1.954	-0.000127479\\
1.956	-0.000126649\\
1.958	-0.00012499\\
1.96	-0.000122518\\
1.962	-0.000121583\\
1.964	-0.000121261\\
1.966	-0.000120138\\
1.968	-0.00011823\\
1.97	-0.00011741\\
1.972	-0.000116092\\
1.974	-0.000114017\\
1.976	-0.000111201\\
1.978	-0.000108484\\
1.98	-0.000108253\\
1.982	-0.000107434\\
1.984	-0.000106038\\
1.986	-0.000104076\\
1.988	-0.000101563\\
1.99	-9.8514e-05\\
1.992	-9.49452e-05\\
1.994	-9.08748e-05\\
1.996	-8.6322e-05\\
1.998	-8.13071e-05\\
2	-7.58515e-05\\
};
\addplot [color=mycolor10, line width=1.0pt, forget plot]
  table[row sep=crcr]{%
-0.024	0\\
-0.022	0\\
-0.02	0\\
-0.018	0\\
-0.016	0\\
-0.014	0\\
-0.012	0\\
-0.01	0\\
-0.008	0\\
-0.006	0\\
-0.004	0\\
-0.002	0\\
0	0\\
0.002	-4.01004e-07\\
0.004	-4.71276e-06\\
0.006	-1.97146e-05\\
0.008	-5.31182e-05\\
0.01	-0.000111761\\
0.012	-0.000200425\\
0.014	-0.000321208\\
0.016	-0.000473364\\
0.018	-0.00203606\\
0.02	-0.00433736\\
0.022	-0.00424938\\
0.024	-0.00520985\\
0.026	-0.00847091\\
0.028	-0.00940308\\
0.03	-0.0103002\\
0.032	-0.0111471\\
0.034	-0.0119292\\
0.036	-0.0126327\\
0.038	-0.0132451\\
0.04	-0.013755\\
0.042	-0.0141526\\
0.044	-0.0144295\\
0.046	-0.0145791\\
0.048	-0.0145966\\
0.05	-0.0144789\\
0.052	-0.0142249\\
0.054	-0.013835\\
0.056	-0.0133117\\
0.058	-0.0126593\\
0.06	-0.0118836\\
0.062	-0.0109919\\
0.064	-0.014406\\
0.066	-0.0148961\\
0.068	-0.0153462\\
0.07	-0.015753\\
0.072	-0.0161134\\
0.074	-0.0164244\\
0.076	-0.0166835\\
0.078	-0.0168882\\
0.08	-0.0170363\\
0.082	-0.0171259\\
0.084	-0.0177319\\
0.086	-0.0183221\\
0.088	-0.0188605\\
0.09	-0.0193434\\
0.092	-0.0197703\\
0.094	-0.0202281\\
0.096	-0.0206281\\
0.098	-0.0211199\\
0.1	-0.0232786\\
0.102	-0.025334\\
0.104	-0.027276\\
0.106	-0.0290952\\
0.108	-0.0307831\\
0.11	-0.0323317\\
0.112	-0.0315973\\
0.114	-0.0297078\\
0.116	-0.0277141\\
0.118	-0.0256212\\
0.12	-0.0234348\\
0.122	-0.0211609\\
0.124	-0.0188059\\
0.126	-0.0168715\\
0.128	-0.0160886\\
0.13	-0.0152792\\
0.132	-0.0149888\\
0.134	-0.0146687\\
0.136	-0.0143186\\
0.138	-0.0139384\\
0.14	-0.0135281\\
0.142	-0.0130878\\
0.144	-0.0126177\\
0.146	-0.0121182\\
0.148	-0.0116601\\
0.15	-0.0111653\\
0.152	-0.0106355\\
0.154	-0.0100722\\
0.156	-0.00947711\\
0.158	-0.00885227\\
0.16	-0.00819966\\
0.162	-0.00752147\\
0.164	-0.00681996\\
0.166	-0.0072667\\
0.168	-0.01021\\
0.17	-0.0128622\\
0.172	-0.0148768\\
0.174	-0.0150016\\
0.176	-0.0150096\\
0.178	-0.0149023\\
0.18	-0.0146823\\
0.182	-0.0143528\\
0.184	-0.0139183\\
0.186	-0.0133839\\
0.188	-0.0127558\\
0.19	-0.0120405\\
0.192	-0.0112457\\
0.194	-0.0103791\\
0.196	-0.00944938\\
0.198	-0.00846528\\
0.2	-0.00743603\\
0.202	-0.00637102\\
0.204	-0.00640949\\
0.206	-0.00822246\\
0.208	-0.00989943\\
0.21	-0.0108449\\
0.212	-0.0113431\\
0.214	-0.0117687\\
0.216	-0.0121187\\
0.218	-0.0123905\\
0.22	-0.0125822\\
0.222	-0.0126924\\
0.224	-0.0127203\\
0.226	-0.0126657\\
0.228	-0.0125288\\
0.23	-0.0123108\\
0.232	-0.012013\\
0.234	-0.0116376\\
0.236	-0.011187\\
0.238	-0.0106645\\
0.24	-0.0100735\\
0.242	-0.00941826\\
0.244	-0.00870312\\
0.246	-0.00793299\\
0.248	-0.00711311\\
0.25	-0.00624901\\
0.252	-0.00534652\\
0.254	-0.00549471\\
0.256	-0.00757139\\
0.258	-0.00955561\\
0.26	-0.0114416\\
0.262	-0.0132242\\
0.264	-0.0142686\\
0.266	-0.0149822\\
0.268	-0.0156364\\
0.27	-0.0162291\\
0.272	-0.0167582\\
0.274	-0.0172221\\
0.276	-0.0176193\\
0.278	-0.0179487\\
0.28	-0.0182092\\
0.282	-0.0184002\\
0.284	-0.0185213\\
0.286	-0.0185725\\
0.288	-0.0185537\\
0.29	-0.0184656\\
0.292	-0.0183086\\
0.294	-0.0184409\\
0.296	-0.0190959\\
0.298	-0.0196806\\
0.3	-0.0201933\\
0.302	-0.0209036\\
0.304	-0.0216896\\
0.306	-0.0223798\\
0.308	-0.0229724\\
0.31	-0.0234662\\
0.312	-0.0238603\\
0.314	-0.0241542\\
0.316	-0.0243478\\
0.318	-0.0244417\\
0.32	-0.0244365\\
0.322	-0.0243333\\
0.324	-0.0241339\\
0.326	-0.02384\\
0.328	-0.023266\\
0.33	-0.0221133\\
0.332	-0.0208908\\
0.334	-0.0196039\\
0.336	-0.0182577\\
0.338	-0.0171087\\
0.34	-0.0170454\\
0.342	-0.016931\\
0.344	-0.0164787\\
0.346	-0.0158653\\
0.348	-0.0152113\\
0.35	-0.0145188\\
0.352	-0.0137898\\
0.354	-0.0130266\\
0.356	-0.0126118\\
0.358	-0.0125507\\
0.36	-0.0124492\\
0.362	-0.0123081\\
0.364	-0.012128\\
0.366	-0.0119098\\
0.368	-0.0116545\\
0.37	-0.0113632\\
0.372	-0.0110371\\
0.374	-0.0106775\\
0.376	-0.0102857\\
0.378	-0.00986322\\
0.38	-0.00941162\\
0.382	-0.00893251\\
0.384	-0.0084276\\
0.386	-0.00789863\\
0.388	-0.00734744\\
0.39	-0.00698673\\
0.392	-0.00715119\\
0.394	-0.00712685\\
0.396	-0.00706599\\
0.398	-0.00696917\\
0.4	-0.00683703\\
0.402	-0.00667028\\
0.404	-0.00646973\\
0.406	-0.00673\\
0.408	-0.00716145\\
0.41	-0.00750445\\
0.412	-0.00775934\\
0.414	-0.00792676\\
0.416	-0.0080076\\
0.418	-0.00800306\\
0.42	-0.00793201\\
0.422	-0.00786932\\
0.424	-0.00775846\\
0.426	-0.00756815\\
0.428	-0.00730034\\
0.43	-0.00695722\\
0.432	-0.00654118\\
0.434	-0.00653408\\
0.436	-0.00663969\\
0.438	-0.00671329\\
0.44	-0.00675523\\
0.442	-0.00676592\\
0.444	-0.00674583\\
0.446	-0.00669552\\
0.448	-0.00661558\\
0.45	-0.00650671\\
0.452	-0.00636963\\
0.454	-0.00620515\\
0.456	-0.0060141\\
0.458	-0.00579739\\
0.46	-0.00555596\\
0.462	-0.00529081\\
0.464	-0.00500297\\
0.466	-0.0046935\\
0.468	-0.00436352\\
0.47	-0.00401415\\
0.472	-0.00431839\\
0.474	-0.00510263\\
0.476	-0.0058316\\
0.478	-0.00650347\\
0.48	-0.00711663\\
0.482	-0.00766971\\
0.484	-0.00816158\\
0.486	-0.00859134\\
0.488	-0.00895835\\
0.49	-0.0092622\\
0.492	-0.00950271\\
0.494	-0.00967995\\
0.496	-0.00979421\\
0.498	-0.00984602\\
0.5	-0.00983612\\
0.502	-0.00976548\\
0.504	-0.00963526\\
0.506	-0.00944686\\
0.508	-0.0094618\\
0.51	-0.00994689\\
0.512	-0.010371\\
0.514	-0.0107336\\
0.516	-0.0110343\\
0.518	-0.0112731\\
0.52	-0.0114499\\
0.522	-0.0115651\\
0.524	-0.0116192\\
0.526	-0.0116128\\
0.528	-0.0115468\\
0.53	-0.0114387\\
0.532	-0.0113439\\
0.534	-0.0112309\\
0.536	-0.0110623\\
0.538	-0.0108393\\
0.54	-0.0105635\\
0.542	-0.0102364\\
0.544	-0.00986\\
0.546	-0.0094361\\
0.548	-0.0089668\\
0.55	-0.00845428\\
0.552	-0.00790083\\
0.554	-0.00730885\\
0.556	-0.00668084\\
0.558	-0.00619392\\
0.56	-0.00633562\\
0.562	-0.00651637\\
0.564	-0.0066591\\
0.566	-0.00676395\\
0.568	-0.00683116\\
0.57	-0.00686106\\
0.572	-0.00685411\\
0.574	-0.00681083\\
0.576	-0.00673188\\
0.578	-0.00661798\\
0.58	-0.00646996\\
0.582	-0.00628872\\
0.584	-0.00607525\\
0.586	-0.00583062\\
0.588	-0.00555598\\
0.59	-0.00525255\\
0.592	-0.00492161\\
0.594	-0.00456449\\
0.596	-0.00418261\\
0.598	-0.00377741\\
0.6	-0.00335041\\
0.602	-0.0030416\\
0.604	-0.00288176\\
0.606	-0.00269999\\
0.608	-0.0024972\\
0.61	-0.00227433\\
0.612	-0.0020324\\
0.614	-0.00177244\\
0.616	-0.00172988\\
0.618	-0.00215586\\
0.62	-0.00253915\\
0.622	-0.00287924\\
0.624	-0.00317576\\
0.626	-0.0034285\\
0.628	-0.00363738\\
0.63	-0.00380246\\
0.632	-0.00392396\\
0.634	-0.00400223\\
0.636	-0.00403774\\
0.638	-0.0040311\\
0.64	-0.00399628\\
0.642	-0.00397568\\
0.644	-0.00395158\\
0.646	-0.0038903\\
0.648	-0.00379017\\
0.65	-0.00365219\\
0.652	-0.00347744\\
0.654	-0.00326713\\
0.656	-0.00302254\\
0.658	-0.00274506\\
0.66	-0.00243617\\
0.662	-0.00223307\\
0.664	-0.00214293\\
0.666	-0.00203593\\
0.668	-0.00191258\\
0.67	-0.00177341\\
0.672	-0.001619\\
0.674	-0.00144995\\
0.676	-0.0012669\\
0.678	-0.00109203\\
0.68	-0.00125207\\
0.682	-0.00138502\\
0.684	-0.00156453\\
0.686	-0.00190762\\
0.688	-0.00222256\\
0.69	-0.00250868\\
0.692	-0.00276544\\
0.694	-0.00299241\\
0.696	-0.00318927\\
0.698	-0.00335582\\
0.7	-0.00349196\\
0.702	-0.00359771\\
0.704	-0.0036732\\
0.706	-0.00371868\\
0.708	-0.00373447\\
0.71	-0.00372103\\
0.712	-0.00367889\\
0.714	-0.0036087\\
0.716	-0.0035112\\
0.718	-0.00338719\\
0.72	-0.00323758\\
0.722	-0.00311897\\
0.724	-0.00343679\\
0.726	-0.00375024\\
0.728	-0.00403433\\
0.73	-0.00428858\\
0.732	-0.00451259\\
0.734	-0.00470606\\
0.736	-0.00486882\\
0.738	-0.00500078\\
0.74	-0.00510195\\
0.742	-0.00517245\\
0.744	-0.0052125\\
0.746	-0.0052224\\
0.748	-0.00520256\\
0.75	-0.00516182\\
0.752	-0.00514054\\
0.754	-0.00511262\\
0.756	-0.00506707\\
0.758	-0.00499374\\
0.76	-0.00489327\\
0.762	-0.00476639\\
0.764	-0.00461387\\
0.766	-0.0044366\\
0.768	-0.00423549\\
0.77	-0.00401156\\
0.772	-0.00376586\\
0.774	-0.00349951\\
0.776	-0.00321368\\
0.778	-0.00290959\\
0.78	-0.00269862\\
0.782	-0.00270899\\
0.784	-0.00270059\\
0.786	-0.00267373\\
0.788	-0.00262878\\
0.79	-0.00256615\\
0.792	-0.00269461\\
0.794	-0.00282819\\
0.796	-0.0029426\\
0.798	-0.0030378\\
0.8	-0.00311379\\
0.802	-0.00317063\\
0.804	-0.00320842\\
0.806	-0.00322734\\
0.808	-0.0032276\\
0.81	-0.00320948\\
0.812	-0.00317328\\
0.814	-0.00311938\\
0.816	-0.00304819\\
0.818	-0.00296017\\
0.82	-0.00285582\\
0.822	-0.00273567\\
0.824	-0.00260032\\
0.826	-0.00245038\\
0.828	-0.0022865\\
0.83	-0.00210936\\
0.832	-0.00191969\\
0.834	-0.00171822\\
0.836	-0.00150573\\
0.838	-0.00128299\\
0.84	-0.00144406\\
0.842	-0.00161887\\
0.844	-0.00177167\\
0.846	-0.00190231\\
0.848	-0.00201073\\
0.85	-0.00209693\\
0.852	-0.00216098\\
0.854	-0.00220304\\
0.856	-0.00222331\\
0.858	-0.00222209\\
0.86	-0.00220318\\
0.862	-0.00220612\\
0.864	-0.00219857\\
0.866	-0.0021868\\
0.868	-0.00215481\\
0.87	-0.00210301\\
0.872	-0.00203191\\
0.874	-0.00194205\\
0.876	-0.00183404\\
0.878	-0.00170851\\
0.88	-0.00156619\\
0.882	-0.00140781\\
0.884	-0.00123416\\
0.886	-0.00104607\\
0.888	-0.000844411\\
0.89	-0.000805939\\
0.892	-0.000890632\\
0.894	-0.000965624\\
0.896	-0.00103088\\
0.898	-0.0010864\\
0.9	-0.00113219\\
0.902	-0.00116829\\
0.904	-0.00125721\\
0.906	-0.00135199\\
0.908	-0.00143194\\
0.91	-0.00149702\\
0.912	-0.00154724\\
0.914	-0.00158266\\
0.916	-0.00160341\\
0.918	-0.00160964\\
0.92	-0.00160158\\
0.922	-0.0015795\\
0.924	-0.0015437\\
0.926	-0.00149455\\
0.928	-0.00143244\\
0.93	-0.00135782\\
0.932	-0.00132534\\
0.934	-0.00147907\\
0.936	-0.00161809\\
0.938	-0.00174213\\
0.94	-0.001851\\
0.942	-0.00194454\\
0.944	-0.00202267\\
0.946	-0.00208535\\
0.948	-0.00213259\\
0.95	-0.00216447\\
0.952	-0.00227844\\
0.954	-0.0023776\\
0.956	-0.00246085\\
0.958	-0.00252814\\
0.96	-0.00257948\\
0.962	-0.00261493\\
0.964	-0.0026346\\
0.966	-0.00263865\\
0.968	-0.00262729\\
0.97	-0.00260106\\
0.972	-0.00259672\\
0.974	-0.00257974\\
0.976	-0.00256711\\
0.978	-0.00253992\\
0.98	-0.00249846\\
0.982	-0.00244308\\
0.984	-0.00237415\\
0.986	-0.0022921\\
0.988	-0.00219739\\
0.99	-0.00209052\\
0.992	-0.00197199\\
0.994	-0.00184239\\
0.996	-0.00170229\\
0.998	-0.00155231\\
1	-0.00139308\\
1.002	-0.00122528\\
1.004	-0.00117615\\
1.006	-0.00119859\\
1.008	-0.00121197\\
1.01	-0.0012161\\
1.012	-0.00121107\\
1.014	-0.00119703\\
1.016	-0.00117417\\
1.018	-0.00114271\\
1.02	-0.00110289\\
1.022	-0.00110841\\
1.024	-0.00119863\\
1.026	-0.00127946\\
1.028	-0.00135083\\
1.03	-0.00141266\\
1.032	-0.00146492\\
1.034	-0.0015076\\
1.036	-0.00154072\\
1.038	-0.00156433\\
1.04	-0.00157849\\
1.042	-0.00158332\\
1.044	-0.00157892\\
1.046	-0.00156545\\
1.048	-0.00154309\\
1.05	-0.00151204\\
1.052	-0.0014725\\
1.054	-0.00142473\\
1.056	-0.00136899\\
1.058	-0.00130556\\
1.06	-0.00123473\\
1.062	-0.00115684\\
1.064	-0.0011436\\
1.066	-0.00120767\\
1.068	-0.00126017\\
1.07	-0.00130109\\
1.072	-0.0013305\\
1.074	-0.00134847\\
1.076	-0.00135512\\
1.078	-0.00135061\\
1.08	-0.00133513\\
1.082	-0.00133669\\
1.084	-0.00132904\\
1.086	-0.00132483\\
1.088	-0.00131287\\
1.09	-0.00129052\\
1.092	-0.001258\\
1.094	-0.00121558\\
1.096	-0.00116356\\
1.098	-0.00110226\\
1.1	-0.00103203\\
1.102	-0.000953244\\
1.104	-0.00086629\\
1.106	-0.000771589\\
1.108	-0.000669577\\
1.11	-0.000560709\\
1.112	-0.000492711\\
1.114	-0.000548383\\
1.116	-0.000596733\\
1.118	-0.000637712\\
1.12	-0.000671301\\
1.122	-0.000697507\\
1.124	-0.000716362\\
1.126	-0.000727926\\
1.128	-0.000732283\\
1.13	-0.000729541\\
1.132	-0.000719833\\
1.134	-0.000703315\\
1.136	-0.000680165\\
1.138	-0.000650581\\
1.14	-0.000614783\\
1.142	-0.000624122\\
1.144	-0.000636953\\
1.146	-0.000645815\\
1.148	-0.000728764\\
1.15	-0.00080431\\
1.152	-0.000872311\\
1.154	-0.000932654\\
1.156	-0.000985252\\
1.158	-0.00103005\\
1.16	-0.001067\\
1.162	-0.00109611\\
1.164	-0.0011174\\
1.166	-0.00113092\\
1.168	-0.00113673\\
1.17	-0.00118174\\
1.172	-0.00124078\\
1.174	-0.00129153\\
1.176	-0.00133396\\
1.178	-0.00136803\\
1.18	-0.00139376\\
1.182	-0.00141118\\
1.184	-0.00142035\\
1.186	-0.00142137\\
1.188	-0.00141434\\
1.19	-0.0013994\\
1.192	-0.00139733\\
1.194	-0.00139\\
1.196	-0.00138255\\
1.198	-0.00137321\\
1.2	-0.00135633\\
1.202	-0.00133209\\
1.204	-0.00130065\\
1.206	-0.00126224\\
1.208	-0.00121706\\
1.21	-0.00116538\\
1.212	-0.00110745\\
1.214	-0.00104355\\
1.216	-0.000973983\\
1.218	-0.000899055\\
1.22	-0.000819091\\
1.222	-0.000734429\\
1.224	-0.000645417\\
1.226	-0.000552414\\
1.228	-0.000548518\\
1.23	-0.000566687\\
1.232	-0.000580177\\
1.234	-0.000589188\\
1.236	-0.000593676\\
1.238	-0.000593623\\
1.24	-0.000589103\\
1.242	-0.000580203\\
1.244	-0.00056702\\
1.246	-0.000549665\\
1.248	-0.000528261\\
1.25	-0.00050294\\
1.252	-0.000475668\\
1.254	-0.000530586\\
1.256	-0.000580961\\
1.258	-0.000626715\\
1.26	-0.000667786\\
1.262	-0.000704124\\
1.264	-0.000735695\\
1.266	-0.000762478\\
1.268	-0.000784467\\
1.27	-0.000801668\\
1.272	-0.000814102\\
1.274	-0.000821804\\
1.276	-0.000824822\\
1.278	-0.000823217\\
1.28	-0.00081706\\
1.282	-0.000806438\\
1.284	-0.000791449\\
1.286	-0.0007722\\
1.288	-0.000748716\\
1.29	-0.00076708\\
1.292	-0.000780328\\
1.294	-0.000787597\\
1.296	-0.00078895\\
1.298	-0.000784472\\
1.3	-0.000774265\\
1.302	-0.000773572\\
1.304	-0.000770216\\
1.306	-0.000764442\\
1.308	-0.000760368\\
1.31	-0.000750807\\
1.312	-0.000735873\\
1.314	-0.000715695\\
1.316	-0.000690417\\
1.318	-0.000660199\\
1.32	-0.000625213\\
1.322	-0.000585648\\
1.324	-0.000541701\\
1.326	-0.000493584\\
1.328	-0.00044152\\
1.33	-0.00038574\\
1.332	-0.000326486\\
1.334	-0.000318668\\
1.336	-0.000324793\\
1.338	-0.000327369\\
1.34	-0.000326448\\
1.342	-0.000322093\\
1.344	-0.000314379\\
1.346	-0.00030339\\
1.348	-0.000308328\\
1.35	-0.000323553\\
1.352	-0.000335171\\
1.354	-0.0003432\\
1.356	-0.000347665\\
1.358	-0.000348607\\
1.36	-0.000346078\\
1.362	-0.000347976\\
1.364	-0.000391201\\
1.366	-0.000430572\\
1.368	-0.000466011\\
1.37	-0.000497456\\
1.372	-0.000524858\\
1.374	-0.000548184\\
1.376	-0.000567412\\
1.378	-0.000582537\\
1.38	-0.000593566\\
1.382	-0.000600522\\
1.384	-0.000603439\\
1.386	-0.000602366\\
1.388	-0.000606634\\
1.39	-0.000640763\\
1.392	-0.000670571\\
1.394	-0.000696023\\
1.396	-0.000717095\\
1.398	-0.000733779\\
1.4	-0.000746082\\
1.402	-0.000754023\\
1.404	-0.000757638\\
1.406	-0.000756975\\
1.408	-0.000752097\\
1.41	-0.000743077\\
1.412	-0.000741055\\
1.414	-0.000737979\\
1.416	-0.000731241\\
1.418	-0.000728205\\
1.42	-0.000721187\\
1.422	-0.000710266\\
1.424	-0.000695533\\
1.426	-0.000677089\\
1.428	-0.000655048\\
1.43	-0.00062953\\
1.432	-0.00060067\\
1.434	-0.000568609\\
1.436	-0.000533497\\
1.438	-0.000495493\\
1.44	-0.000454764\\
1.442	-0.000411482\\
1.444	-0.000365826\\
1.446	-0.000317981\\
1.448	-0.000268137\\
1.45	-0.000226088\\
1.452	-0.000240203\\
1.454	-0.000251993\\
1.456	-0.000261459\\
1.458	-0.000268608\\
1.46	-0.000273536\\
1.462	-0.000276241\\
1.464	-0.00027669\\
1.466	-0.000274917\\
1.468	-0.000270961\\
1.47	-0.00026487\\
1.472	-0.000256696\\
1.474	-0.000246497\\
1.476	-0.000234336\\
1.478	-0.000220284\\
1.48	-0.00023499\\
1.482	-0.000252913\\
1.484	-0.000267853\\
1.486	-0.000279807\\
1.488	-0.000288781\\
1.49	-0.000294792\\
1.492	-0.000308649\\
1.494	-0.000328797\\
1.496	-0.000346572\\
1.498	-0.000361957\\
1.5	-0.00037494\\
1.502	-0.000385518\\
1.504	-0.000393697\\
1.506	-0.000399484\\
1.508	-0.000403597\\
1.51	-0.000413114\\
1.512	-0.000419468\\
1.514	-0.000422685\\
1.516	-0.000422802\\
1.518	-0.000419865\\
1.52	-0.00041393\\
1.522	-0.00041308\\
1.524	-0.000412147\\
1.526	-0.000408273\\
1.528	-0.000406477\\
1.53	-0.000403075\\
1.532	-0.000396836\\
1.534	-0.000387824\\
1.536	-0.00037611\\
1.538	-0.000361774\\
1.54	-0.000344904\\
1.542	-0.000325593\\
1.544	-0.000303942\\
1.546	-0.000280059\\
1.548	-0.000254057\\
1.55	-0.000226055\\
1.552	-0.000196177\\
1.554	-0.000164551\\
1.556	-0.000134921\\
1.558	-0.000145721\\
1.56	-0.000154611\\
1.562	-0.00016159\\
1.564	-0.00016666\\
1.566	-0.000169835\\
1.568	-0.000171133\\
1.57	-0.000170576\\
1.572	-0.000168197\\
1.574	-0.000164032\\
1.576	-0.000158124\\
1.578	-0.000171718\\
1.58	-0.000194283\\
1.582	-0.000214813\\
1.584	-0.000233271\\
1.586	-0.000249622\\
1.588	-0.000263843\\
1.59	-0.000275917\\
1.592	-0.000285833\\
1.594	-0.00029359\\
1.596	-0.000299193\\
1.598	-0.000302654\\
1.6	-0.000303992\\
1.602	-0.000303235\\
1.604	-0.000300415\\
1.606	-0.000295572\\
1.608	-0.000310987\\
1.61	-0.000328217\\
1.612	-0.000343175\\
1.614	-0.000355843\\
1.616	-0.000366211\\
1.618	-0.000374279\\
1.62	-0.000380051\\
1.622	-0.000383541\\
1.624	-0.00038477\\
1.626	-0.000383765\\
1.628	-0.000380561\\
1.63	-0.000375201\\
1.632	-0.000373565\\
1.634	-0.000372354\\
1.636	-0.000369025\\
1.638	-0.000366264\\
1.64	-0.000363586\\
1.642	-0.000358862\\
1.644	-0.000352136\\
1.646	-0.000343459\\
1.648	-0.000332887\\
1.65	-0.000320483\\
1.652	-0.000306313\\
1.654	-0.000290449\\
1.656	-0.000272968\\
1.658	-0.000253951\\
1.66	-0.000233484\\
1.662	-0.000211654\\
1.664	-0.000188554\\
1.666	-0.000164279\\
1.668	-0.000138927\\
1.67	-0.000112599\\
1.672	-0.000107668\\
1.674	-0.00010906\\
1.676	-0.000109543\\
1.678	-0.00010913\\
1.68	-0.000107838\\
1.682	-0.000109973\\
1.684	-0.000114261\\
1.686	-0.000117451\\
1.688	-0.000119552\\
1.69	-0.000120534\\
1.692	-0.000120413\\
1.694	-0.000122975\\
1.696	-0.000133724\\
1.698	-0.000142911\\
1.7	-0.000150528\\
1.702	-0.000156575\\
1.704	-0.000161054\\
1.706	-0.000163976\\
1.708	-0.000165354\\
1.71	-0.000165207\\
1.712	-0.00016356\\
1.714	-0.000160442\\
1.716	-0.00016111\\
1.718	-0.000175255\\
1.72	-0.000187762\\
1.722	-0.000198614\\
1.724	-0.000207801\\
1.726	-0.000215317\\
1.728	-0.000221162\\
1.73	-0.000225341\\
1.732	-0.000227865\\
1.734	-0.00022875\\
1.736	-0.000228018\\
1.738	-0.000225695\\
1.74	-0.000221813\\
1.742	-0.000220666\\
1.744	-0.000220153\\
1.746	-0.000218105\\
1.748	-0.000215712\\
1.75	-0.000214346\\
1.752	-0.000211497\\
1.754	-0.000207197\\
1.756	-0.000201482\\
1.758	-0.000194392\\
1.76	-0.000185971\\
1.762	-0.000176268\\
1.764	-0.000165333\\
1.766	-0.000153222\\
1.768	-0.000139992\\
1.77	-0.000125704\\
1.772	-0.000110423\\
1.774	-9.4214e-05\\
1.776	-8.97342e-05\\
1.778	-9.10928e-05\\
1.78	-9.14493e-05\\
1.782	-9.08181e-05\\
1.784	-8.9217e-05\\
1.786	-8.66669e-05\\
1.788	-8.31919e-05\\
1.79	-7.8819e-05\\
1.792	-7.63535e-05\\
1.794	-8.93847e-05\\
1.796	-0.000101369\\
1.798	-0.000112281\\
1.8	-0.0001221\\
1.802	-0.000130809\\
1.804	-0.000138395\\
1.806	-0.000144849\\
1.808	-0.000150166\\
1.81	-0.000154343\\
1.812	-0.000157385\\
1.814	-0.000159297\\
1.816	-0.000160089\\
1.818	-0.000159774\\
1.82	-0.000158372\\
1.822	-0.000155901\\
1.824	-0.000152386\\
1.826	-0.00015774\\
1.828	-0.000167459\\
1.83	-0.000175989\\
1.832	-0.000183318\\
1.834	-0.000189439\\
1.836	-0.00019435\\
1.838	-0.000198051\\
1.84	-0.000200547\\
1.842	-0.000201847\\
1.844	-0.000201964\\
1.846	-0.000200913\\
1.848	-0.000198715\\
1.85	-0.000195394\\
1.852	-0.000194114\\
1.854	-0.000193501\\
1.856	-0.000191781\\
1.858	-0.000189226\\
1.86	-0.000188157\\
1.862	-0.000186018\\
1.864	-0.000182831\\
1.866	-0.000178622\\
1.868	-0.000173419\\
1.87	-0.000167255\\
1.872	-0.000160162\\
1.874	-0.000152179\\
1.876	-0.000143343\\
1.878	-0.000133697\\
1.88	-0.000123285\\
1.882	-0.000112151\\
1.884	-0.000100344\\
1.886	-8.79125e-05\\
1.888	-7.49071e-05\\
1.89	-6.13795e-05\\
1.892	-4.73827e-05\\
1.894	-3.71457e-05\\
1.896	-3.40288e-05\\
1.898	-3.02795e-05\\
1.9	-3.07941e-05\\
1.902	-3.6264e-05\\
1.904	-4.5231e-05\\
1.906	-5.34187e-05\\
1.908	-6.08116e-05\\
1.91	-6.73971e-05\\
1.912	-7.31656e-05\\
1.914	-7.811e-05\\
1.916	-8.22265e-05\\
1.918	-8.55138e-05\\
1.92	-8.79735e-05\\
1.922	-8.961e-05\\
1.924	-9.04305e-05\\
1.926	-9.04447e-05\\
1.928	-8.96649e-05\\
1.93	-8.81061e-05\\
1.932	-8.57855e-05\\
1.934	-8.38016e-05\\
1.936	-9.16627e-05\\
1.938	-9.86679e-05\\
1.94	-0.000104807\\
1.942	-0.000110072\\
1.944	-0.000114459\\
1.946	-0.000117966\\
1.948	-0.000120595\\
1.95	-0.00012235\\
1.952	-0.000123237\\
1.954	-0.000123268\\
1.956	-0.000122453\\
1.958	-0.000120809\\
1.96	-0.000118354\\
1.962	-0.000117434\\
1.964	-0.000117127\\
1.966	-0.00011602\\
1.968	-0.000114128\\
1.97	-0.000113323\\
1.972	-0.000112021\\
1.974	-0.000109961\\
1.976	-0.000107161\\
1.978	-0.000104458\\
1.98	-0.000104243\\
1.982	-0.00010344\\
1.984	-0.000102058\\
1.986	-0.000100112\\
1.988	-9.76141e-05\\
1.99	-9.45799e-05\\
1.992	-9.10262e-05\\
1.994	-8.69708e-05\\
1.996	-8.2433e-05\\
1.998	-7.74331e-05\\
2	-7.19924e-05\\
};
\addplot [color=mycolor11, line width=1.0pt, forget plot]
  table[row sep=crcr]{%
-0.024	0\\
-0.022	0\\
-0.02	0\\
-0.018	0\\
-0.016	0\\
-0.014	0\\
-0.012	0\\
-0.01	0\\
-0.008	0\\
-0.006	0\\
-0.004	0\\
-0.002	0\\
0	0\\
0.002	-0.000283196\\
0.004	-0.000978526\\
0.006	-0.00188336\\
0.008	-0.00283483\\
0.01	-0.00370974\\
0.012	-0.00442248\\
0.014	-0.00492154\\
0.016	-0.00518487\\
0.018	-0.00549102\\
0.02	-0.00542819\\
0.022	-0.00516914\\
0.024	-0.00628811\\
0.026	-0.00871419\\
0.028	-0.00970693\\
0.03	-0.0106671\\
0.032	-0.0115779\\
0.034	-0.0124233\\
0.036	-0.0131882\\
0.038	-0.013859\\
0.04	-0.0144231\\
0.042	-0.0148697\\
0.044	-0.01519\\
0.046	-0.0153767\\
0.048	-0.0154248\\
0.05	-0.0153309\\
0.052	-0.0150939\\
0.054	-0.0147146\\
0.056	-0.0141957\\
0.058	-0.0135417\\
0.06	-0.0127589\\
0.062	-0.0118553\\
0.064	-0.0150795\\
0.066	-0.0155751\\
0.068	-0.0160291\\
0.07	-0.0164384\\
0.072	-0.0168\\
0.074	-0.0171112\\
0.076	-0.0173694\\
0.078	-0.0175723\\
0.08	-0.0177178\\
0.082	-0.0178041\\
0.084	-0.018406\\
0.086	-0.0189914\\
0.088	-0.0195243\\
0.09	-0.020001\\
0.092	-0.020421\\
0.094	-0.0208712\\
0.096	-0.0212628\\
0.098	-0.0217456\\
0.1	-0.0238946\\
0.102	-0.0259394\\
0.104	-0.0278702\\
0.106	-0.0296775\\
0.108	-0.0313528\\
0.11	-0.0328882\\
0.112	-0.0323977\\
0.114	-0.0305144\\
0.116	-0.0285254\\
0.118	-0.0264359\\
0.12	-0.0242513\\
0.122	-0.0219777\\
0.124	-0.0196216\\
0.126	-0.0177211\\
0.128	-0.0170738\\
0.13	-0.0162916\\
0.132	-0.0159793\\
0.134	-0.0156352\\
0.136	-0.0152591\\
0.138	-0.014851\\
0.14	-0.014411\\
0.142	-0.0139393\\
0.144	-0.0134361\\
0.146	-0.0129022\\
0.148	-0.0124082\\
0.15	-0.0118763\\
0.152	-0.0113082\\
0.154	-0.0107057\\
0.156	-0.0100705\\
0.158	-0.00940486\\
0.16	-0.00871082\\
0.162	-0.00799074\\
0.164	-0.00724701\\
0.166	-0.00765131\\
0.168	-0.0105521\\
0.17	-0.0131619\\
0.172	-0.0149777\\
0.174	-0.0151003\\
0.176	-0.0151058\\
0.178	-0.0149958\\
0.18	-0.0147728\\
0.182	-0.01444\\
0.184	-0.014002\\
0.186	-0.013464\\
0.188	-0.0128319\\
0.19	-0.0121126\\
0.192	-0.0113135\\
0.194	-0.0104425\\
0.196	-0.00950828\\
0.198	-0.00851958\\
0.2	-0.00748566\\
0.202	-0.00641594\\
0.204	-0.00660514\\
0.206	-0.00840924\\
0.208	-0.0100777\\
0.21	-0.0107908\\
0.212	-0.0112993\\
0.214	-0.0117351\\
0.216	-0.0120952\\
0.218	-0.0123768\\
0.22	-0.0125781\\
0.222	-0.0126975\\
0.224	-0.0127341\\
0.226	-0.0126878\\
0.228	-0.0125587\\
0.23	-0.0123478\\
0.232	-0.0120565\\
0.234	-0.0116869\\
0.236	-0.0112414\\
0.238	-0.0107232\\
0.24	-0.0101358\\
0.242	-0.00948335\\
0.244	-0.0087702\\
0.246	-0.00800127\\
0.248	-0.00718178\\
0.25	-0.00631729\\
0.252	-0.00541363\\
0.254	-0.00558127\\
0.256	-0.00765726\\
0.258	-0.00964095\\
0.26	-0.0115266\\
0.262	-0.0133088\\
0.264	-0.0142067\\
0.266	-0.0149256\\
0.268	-0.015585\\
0.27	-0.0161827\\
0.272	-0.0167167\\
0.274	-0.0171853\\
0.276	-0.0175869\\
0.278	-0.0179205\\
0.28	-0.0181849\\
0.282	-0.0183796\\
0.284	-0.0185042\\
0.286	-0.0185585\\
0.288	-0.0185426\\
0.29	-0.0184569\\
0.292	-0.0183022\\
0.294	-0.0184363\\
0.296	-0.0190929\\
0.298	-0.0196789\\
0.3	-0.0201926\\
0.302	-0.0209035\\
0.304	-0.0216898\\
0.306	-0.02238\\
0.308	-0.0229723\\
0.31	-0.0234654\\
0.312	-0.0238585\\
0.314	-0.0241512\\
0.316	-0.0243434\\
0.318	-0.0244354\\
0.32	-0.0244281\\
0.322	-0.0243227\\
0.324	-0.0241206\\
0.326	-0.023824\\
0.328	-0.0232671\\
0.33	-0.0221134\\
0.332	-0.0208898\\
0.334	-0.0196014\\
0.336	-0.0182536\\
0.338	-0.0171027\\
0.34	-0.0170374\\
0.342	-0.0169209\\
0.344	-0.0164709\\
0.346	-0.0158622\\
0.348	-0.0152125\\
0.35	-0.0145238\\
0.352	-0.0137983\\
0.354	-0.0130381\\
0.356	-0.0126258\\
0.358	-0.0125669\\
0.36	-0.0124671\\
0.362	-0.0123273\\
0.364	-0.0121481\\
0.366	-0.0119303\\
0.368	-0.0116751\\
0.37	-0.0113835\\
0.372	-0.0110567\\
0.374	-0.010696\\
0.376	-0.0103028\\
0.378	-0.00987858\\
0.38	-0.00942489\\
0.382	-0.00894337\\
0.384	-0.00843574\\
0.386	-0.00790377\\
0.388	-0.00734931\\
0.39	-0.00698507\\
0.392	-0.00725678\\
0.394	-0.00723112\\
0.396	-0.00716893\\
0.398	-0.00707078\\
0.4	-0.00693731\\
0.402	-0.00676925\\
0.404	-0.00656737\\
0.406	-0.00682634\\
0.408	-0.00725649\\
0.41	-0.0075982\\
0.412	-0.00785182\\
0.414	-0.00801798\\
0.416	-0.00809758\\
0.418	-0.00809182\\
0.42	-0.00801957\\
0.422	-0.0079557\\
0.424	-0.00784368\\
0.426	-0.00765224\\
0.428	-0.00738333\\
0.43	-0.00703913\\
0.432	-0.00662203\\
0.434	-0.0066139\\
0.436	-0.00671851\\
0.438	-0.00679114\\
0.44	-0.00683214\\
0.442	-0.00684191\\
0.444	-0.00682093\\
0.446	-0.00676975\\
0.448	-0.00668899\\
0.45	-0.00657931\\
0.452	-0.00644146\\
0.454	-0.00627622\\
0.456	-0.00608444\\
0.458	-0.00586703\\
0.46	-0.00562493\\
0.462	-0.00535912\\
0.464	-0.00507064\\
0.466	-0.00476056\\
0.468	-0.00442999\\
0.47	-0.00408005\\
0.472	-0.00438373\\
0.474	-0.00516743\\
0.476	-0.00589588\\
0.478	-0.00656724\\
0.48	-0.0071799\\
0.482	-0.00773249\\
0.484	-0.00822388\\
0.486	-0.00865317\\
0.488	-0.00901972\\
0.49	-0.0093231\\
0.492	-0.00956316\\
0.494	-0.00973994\\
0.496	-0.00985375\\
0.498	-0.00990511\\
0.5	-0.00989476\\
0.502	-0.00982366\\
0.504	-0.00969299\\
0.506	-0.00950413\\
0.508	-0.0095186\\
0.51	-0.0100032\\
0.512	-0.0104269\\
0.514	-0.010789\\
0.516	-0.0110892\\
0.518	-0.0113274\\
0.52	-0.0115037\\
0.522	-0.0116184\\
0.524	-0.0116719\\
0.526	-0.011665\\
0.528	-0.0115984\\
0.53	-0.0114898\\
0.532	-0.0113944\\
0.534	-0.0112809\\
0.536	-0.0111116\\
0.538	-0.010888\\
0.54	-0.0106115\\
0.542	-0.0102839\\
0.544	-0.00990679\\
0.546	-0.00948224\\
0.548	-0.00901227\\
0.55	-0.00849907\\
0.552	-0.00794494\\
0.554	-0.00735227\\
0.556	-0.00672355\\
0.558	-0.00623593\\
0.56	-0.00637692\\
0.562	-0.00655695\\
0.564	-0.00669896\\
0.566	-0.00680308\\
0.568	-0.00686956\\
0.57	-0.00689873\\
0.572	-0.00689105\\
0.574	-0.00684704\\
0.576	-0.00676736\\
0.578	-0.00665273\\
0.58	-0.00650398\\
0.582	-0.00632201\\
0.584	-0.00610783\\
0.586	-0.00586249\\
0.588	-0.00558714\\
0.59	-0.005283\\
0.592	-0.00495136\\
0.594	-0.00459356\\
0.596	-0.004211\\
0.598	-0.00380514\\
0.6	-0.00337747\\
0.602	-0.00306802\\
0.604	-0.00290754\\
0.606	-0.00272514\\
0.608	-0.00252173\\
0.61	-0.00229826\\
0.612	-0.00205574\\
0.614	-0.0017952\\
0.616	-0.00175207\\
0.618	-0.0021775\\
0.62	-0.00256025\\
0.622	-0.00289982\\
0.624	-0.00319583\\
0.626	-0.00344807\\
0.628	-0.00365647\\
0.63	-0.00382108\\
0.632	-0.00394213\\
0.634	-0.00401995\\
0.636	-0.00405503\\
0.638	-0.00404799\\
0.64	-0.00401276\\
0.642	-0.00399177\\
0.644	-0.00396729\\
0.646	-0.00390565\\
0.648	-0.00380517\\
0.65	-0.00366685\\
0.652	-0.00349177\\
0.654	-0.00328114\\
0.656	-0.00303624\\
0.658	-0.00275846\\
0.66	-0.00244928\\
0.662	-0.0022459\\
0.664	-0.00215549\\
0.666	-0.00204823\\
0.668	-0.00192462\\
0.67	-0.0017852\\
0.672	-0.00163055\\
0.674	-0.00146126\\
0.676	-0.00127798\\
0.678	-0.00110288\\
0.68	-0.0012627\\
0.682	-0.00139544\\
0.684	-0.00157474\\
0.686	-0.00191762\\
0.688	-0.00223235\\
0.69	-0.00251827\\
0.692	-0.00277483\\
0.694	-0.00300161\\
0.696	-0.00319827\\
0.698	-0.00336462\\
0.7	-0.00350057\\
0.702	-0.00360613\\
0.704	-0.00368143\\
0.706	-0.00372671\\
0.708	-0.00374232\\
0.71	-0.00372868\\
0.712	-0.00368636\\
0.714	-0.00361598\\
0.716	-0.00351828\\
0.718	-0.00339408\\
0.72	-0.00324428\\
0.722	-0.00312548\\
0.724	-0.0034431\\
0.726	-0.00375636\\
0.728	-0.00404026\\
0.73	-0.00429431\\
0.732	-0.00451812\\
0.734	-0.0047114\\
0.736	-0.00487396\\
0.738	-0.00500571\\
0.74	-0.00510669\\
0.742	-0.00517699\\
0.744	-0.00521683\\
0.746	-0.00522653\\
0.748	-0.00520649\\
0.75	-0.00516555\\
0.752	-0.00514407\\
0.754	-0.00511595\\
0.756	-0.00507019\\
0.758	-0.00499665\\
0.76	-0.00489599\\
0.762	-0.0047689\\
0.764	-0.00461618\\
0.766	-0.0044387\\
0.768	-0.0042374\\
0.77	-0.00401326\\
0.772	-0.00376736\\
0.774	-0.00350081\\
0.776	-0.00321478\\
0.778	-0.00291049\\
0.78	-0.00269932\\
0.782	-0.0027095\\
0.784	-0.0027009\\
0.786	-0.00267385\\
0.788	-0.0026287\\
0.79	-0.00256589\\
0.792	-0.00269416\\
0.794	-0.00282755\\
0.796	-0.00294178\\
0.798	-0.00303679\\
0.8	-0.0031126\\
0.802	-0.00316926\\
0.804	-0.00320687\\
0.806	-0.00322562\\
0.808	-0.00322571\\
0.81	-0.00320741\\
0.812	-0.00317104\\
0.814	-0.00311698\\
0.816	-0.00304562\\
0.818	-0.00295743\\
0.82	-0.00285292\\
0.822	-0.00273262\\
0.824	-0.00259711\\
0.826	-0.00244701\\
0.828	-0.00228298\\
0.83	-0.0021057\\
0.832	-0.00191588\\
0.834	-0.00171426\\
0.836	-0.00150163\\
0.838	-0.00127875\\
0.84	-0.00143968\\
0.842	-0.00161436\\
0.844	-0.00176702\\
0.846	-0.00189754\\
0.848	-0.00200583\\
0.85	-0.0020919\\
0.852	-0.00215583\\
0.854	-0.00219777\\
0.856	-0.00221793\\
0.858	-0.00221659\\
0.86	-0.00219756\\
0.862	-0.0022004\\
0.864	-0.00219274\\
0.866	-0.00218086\\
0.868	-0.00214876\\
0.87	-0.00209687\\
0.872	-0.00202567\\
0.874	-0.00193571\\
0.876	-0.0018276\\
0.878	-0.00170199\\
0.88	-0.00155957\\
0.882	-0.00140111\\
0.884	-0.00122737\\
0.886	-0.0010392\\
0.888	-0.000837459\\
0.89	-0.000798908\\
0.892	-0.000883525\\
0.894	-0.000958442\\
0.896	-0.00102363\\
0.898	-0.00107907\\
0.9	-0.00112479\\
0.902	-0.00116083\\
0.904	-0.00124969\\
0.906	-0.00134441\\
0.908	-0.0014243\\
0.91	-0.00148931\\
0.912	-0.00153948\\
0.914	-0.00157484\\
0.916	-0.00159554\\
0.918	-0.00160172\\
0.92	-0.00159361\\
0.922	-0.00157148\\
0.924	-0.00153563\\
0.926	-0.00148643\\
0.928	-0.00142428\\
0.93	-0.00134962\\
0.932	-0.00131709\\
0.934	-0.00147079\\
0.936	-0.00160977\\
0.938	-0.00173377\\
0.94	-0.0018426\\
0.942	-0.00193611\\
0.944	-0.00201421\\
0.946	-0.00207685\\
0.948	-0.00212406\\
0.95	-0.00215591\\
0.952	-0.00226984\\
0.954	-0.00236898\\
0.956	-0.0024522\\
0.958	-0.00251946\\
0.96	-0.00257077\\
0.962	-0.0026062\\
0.964	-0.00262584\\
0.966	-0.00262986\\
0.968	-0.00261848\\
0.97	-0.00259223\\
0.972	-0.00258787\\
0.974	-0.00257086\\
0.976	-0.00255821\\
0.978	-0.00253099\\
0.98	-0.00248951\\
0.982	-0.00243411\\
0.984	-0.00236516\\
0.986	-0.0022831\\
0.988	-0.00218837\\
0.99	-0.00208147\\
0.992	-0.00196293\\
0.994	-0.00183331\\
0.996	-0.00169319\\
0.998	-0.00154319\\
1	-0.00138394\\
1.002	-0.00121612\\
1.004	-0.00116698\\
1.006	-0.0011894\\
1.008	-0.00120276\\
1.01	-0.00120688\\
1.012	-0.00120183\\
1.014	-0.00118777\\
1.016	-0.0011649\\
1.018	-0.00113342\\
1.02	-0.00109359\\
1.022	-0.00109909\\
1.024	-0.00118929\\
1.026	-0.00127011\\
1.028	-0.00134146\\
1.03	-0.00140328\\
1.032	-0.00145553\\
1.034	-0.00149819\\
1.036	-0.0015313\\
1.038	-0.00155489\\
1.04	-0.00156904\\
1.042	-0.00157385\\
1.044	-0.00156944\\
1.046	-0.00155596\\
1.048	-0.00153359\\
1.05	-0.00150252\\
1.052	-0.00146298\\
1.054	-0.00141519\\
1.056	-0.00135944\\
1.058	-0.001296\\
1.06	-0.00122516\\
1.062	-0.00114726\\
1.064	-0.00113401\\
1.066	-0.00119807\\
1.068	-0.00125056\\
1.07	-0.00129147\\
1.072	-0.00132087\\
1.074	-0.00133883\\
1.076	-0.00134548\\
1.078	-0.00134096\\
1.08	-0.00132547\\
1.082	-0.00132702\\
1.084	-0.00131937\\
1.086	-0.00131515\\
1.088	-0.00130319\\
1.09	-0.00128083\\
1.092	-0.0012483\\
1.094	-0.00120588\\
1.096	-0.00115386\\
1.098	-0.00109256\\
1.1	-0.00102232\\
1.102	-0.000943531\\
1.104	-0.000856574\\
1.106	-0.000761871\\
1.108	-0.000659857\\
1.11	-0.000550989\\
1.112	-0.000482989\\
1.114	-0.000538662\\
1.116	-0.000587011\\
1.118	-0.000627991\\
1.12	-0.000661581\\
1.122	-0.000687788\\
1.124	-0.000706645\\
1.126	-0.000718212\\
1.128	-0.000722571\\
1.13	-0.000719832\\
1.132	-0.000710128\\
1.134	-0.000693614\\
1.136	-0.000670467\\
1.138	-0.000640888\\
1.14	-0.000605096\\
1.142	-0.00061444\\
1.144	-0.000627277\\
1.146	-0.000636146\\
1.148	-0.000719101\\
1.15	-0.000794654\\
1.152	-0.000862663\\
1.154	-0.000923014\\
1.156	-0.00097562\\
1.158	-0.00102042\\
1.16	-0.00105739\\
1.162	-0.00108651\\
1.164	-0.00110781\\
1.166	-0.00112133\\
1.168	-0.00112715\\
1.17	-0.00117218\\
1.172	-0.00123123\\
1.174	-0.00128199\\
1.176	-0.00132443\\
1.178	-0.00135851\\
1.18	-0.00138425\\
1.182	-0.00140168\\
1.184	-0.00141087\\
1.186	-0.0014119\\
1.188	-0.00140488\\
1.19	-0.00138995\\
1.192	-0.0013879\\
1.194	-0.00138059\\
1.196	-0.00137315\\
1.198	-0.00136382\\
1.2	-0.00134696\\
1.202	-0.00132273\\
1.204	-0.00129131\\
1.206	-0.00125291\\
1.208	-0.00120775\\
1.21	-0.00115608\\
1.212	-0.00109817\\
1.214	-0.00103429\\
1.216	-0.000964735\\
1.218	-0.000889822\\
1.22	-0.000809875\\
1.222	-0.000725229\\
1.224	-0.000636234\\
1.226	-0.000543247\\
1.228	-0.000539369\\
1.23	-0.000557555\\
1.232	-0.000571061\\
1.234	-0.000580089\\
1.236	-0.000584595\\
1.238	-0.000584559\\
1.24	-0.000580057\\
1.242	-0.000571174\\
1.244	-0.000558009\\
1.246	-0.000540672\\
1.248	-0.000519285\\
1.25	-0.000493982\\
1.252	-0.000466728\\
1.254	-0.000521664\\
1.256	-0.000572056\\
1.258	-0.000617829\\
1.26	-0.000658917\\
1.262	-0.000695274\\
1.264	-0.000726863\\
1.266	-0.000753664\\
1.268	-0.00077567\\
1.27	-0.000792889\\
1.272	-0.000805342\\
1.274	-0.000813062\\
1.276	-0.000816098\\
1.278	-0.00081451\\
1.28	-0.000808372\\
1.282	-0.000797763\\
1.284	-0.000782141\\
1.286	-0.000762269\\
1.288	-0.000738267\\
1.29	-0.000756625\\
1.292	-0.000769868\\
1.294	-0.000777131\\
1.296	-0.00077848\\
1.298	-0.000773997\\
1.3	-0.000763785\\
1.302	-0.000763089\\
1.304	-0.000759729\\
1.306	-0.000753951\\
1.308	-0.000749874\\
1.31	-0.000740311\\
1.312	-0.000725374\\
1.314	-0.000705194\\
1.316	-0.000679914\\
1.318	-0.000649694\\
1.32	-0.000614707\\
1.322	-0.00057514\\
1.324	-0.000531193\\
1.326	-0.000483076\\
1.328	-0.000431011\\
1.33	-0.000375231\\
1.332	-0.000315978\\
1.334	-0.00030816\\
1.336	-0.000314286\\
1.338	-0.000316863\\
1.34	-0.000315944\\
1.342	-0.00031159\\
1.344	-0.000303878\\
1.346	-0.000292892\\
1.348	-0.000297832\\
1.35	-0.000313059\\
1.352	-0.000324681\\
1.354	-0.000332713\\
1.356	-0.000337182\\
1.358	-0.000338127\\
1.36	-0.000335602\\
1.362	-0.000337506\\
1.364	-0.000380735\\
1.366	-0.000420111\\
1.368	-0.000455555\\
1.37	-0.000487006\\
1.372	-0.000514414\\
1.374	-0.000537746\\
1.376	-0.00055698\\
1.378	-0.000572112\\
1.38	-0.000583148\\
1.382	-0.000590111\\
1.384	-0.000593036\\
1.386	-0.00059197\\
1.388	-0.000596247\\
1.39	-0.000630383\\
1.392	-0.0006602\\
1.394	-0.000685661\\
1.396	-0.000706742\\
1.398	-0.000723435\\
1.4	-0.000735747\\
1.402	-0.000743698\\
1.404	-0.000747323\\
1.406	-0.000746671\\
1.408	-0.000741802\\
1.41	-0.000732793\\
1.412	-0.000730782\\
1.414	-0.000727718\\
1.416	-0.000720991\\
1.418	-0.000717967\\
1.42	-0.00071096\\
1.422	-0.000700052\\
1.424	-0.000685331\\
1.426	-0.0006669\\
1.428	-0.000644871\\
1.43	-0.000619367\\
1.432	-0.00059052\\
1.434	-0.000558472\\
1.436	-0.000523374\\
1.438	-0.000485384\\
1.44	-0.000444668\\
1.442	-0.0004014\\
1.444	-0.000355759\\
1.446	-0.000307928\\
1.448	-0.000258099\\
1.45	-0.000216065\\
1.452	-0.000230195\\
1.454	-0.000242\\
1.456	-0.000251482\\
1.458	-0.000258646\\
1.46	-0.000263589\\
1.462	-0.000266311\\
1.464	-0.000266776\\
1.466	-0.000265018\\
1.468	-0.000261079\\
1.47	-0.000255004\\
1.472	-0.000246846\\
1.474	-0.000236664\\
1.476	-0.00022452\\
1.478	-0.000210485\\
1.48	-0.000225208\\
1.482	-0.000243148\\
1.484	-0.000258105\\
1.486	-0.000270077\\
1.488	-0.000279068\\
1.49	-0.000285096\\
1.492	-0.000298971\\
1.494	-0.000319137\\
1.496	-0.00033693\\
1.498	-0.000352332\\
1.5	-0.000365333\\
1.502	-0.00037593\\
1.504	-0.000384127\\
1.506	-0.000389933\\
1.508	-0.000394064\\
1.51	-0.000403599\\
1.512	-0.000409971\\
1.514	-0.000413207\\
1.516	-0.000413343\\
1.518	-0.000410425\\
1.52	-0.000404509\\
1.522	-0.000403677\\
1.524	-0.000402763\\
1.526	-0.000398909\\
1.528	-0.000397131\\
1.53	-0.000393748\\
1.532	-0.000387529\\
1.534	-0.000378536\\
1.536	-0.000366842\\
1.538	-0.000352526\\
1.54	-0.000335674\\
1.542	-0.000316383\\
1.544	-0.000294751\\
1.546	-0.000270888\\
1.548	-0.000244906\\
1.55	-0.000216923\\
1.552	-0.000187065\\
1.554	-0.00015546\\
1.556	-0.000125849\\
1.558	-0.00013667\\
1.56	-0.00014558\\
1.562	-0.000152578\\
1.564	-0.000157669\\
1.566	-0.000160864\\
1.568	-0.000162181\\
1.57	-0.000161645\\
1.572	-0.000159286\\
1.574	-0.000155142\\
1.576	-0.000149254\\
1.578	-0.000162868\\
1.58	-0.000185454\\
1.582	-0.000206005\\
1.584	-0.000224483\\
1.586	-0.000240855\\
1.588	-0.000255097\\
1.59	-0.000267191\\
1.592	-0.000277128\\
1.594	-0.000284906\\
1.596	-0.00029053\\
1.598	-0.000294012\\
1.6	-0.000295371\\
1.602	-0.000294635\\
1.604	-0.000291836\\
1.606	-0.000287014\\
1.608	-0.00030245\\
1.61	-0.000319701\\
1.612	-0.00033468\\
1.614	-0.000347369\\
1.616	-0.000357759\\
1.618	-0.000365848\\
1.62	-0.000371642\\
1.622	-0.000375153\\
1.624	-0.000376403\\
1.626	-0.00037542\\
1.628	-0.000372238\\
1.63	-0.000366899\\
1.632	-0.000365285\\
1.634	-0.000364096\\
1.636	-0.000360788\\
1.638	-0.000358048\\
1.64	-0.000355392\\
1.642	-0.00035069\\
1.644	-0.000343986\\
1.646	-0.000335331\\
1.648	-0.000324781\\
1.65	-0.000312398\\
1.652	-0.00029825\\
1.654	-0.000282408\\
1.656	-0.000264949\\
1.658	-0.000245954\\
1.66	-0.000225509\\
1.662	-0.000203701\\
1.664	-0.000180623\\
1.666	-0.00015637\\
1.668	-0.00013104\\
1.67	-0.000104734\\
1.672	-9.98251e-05\\
1.674	-0.000101239\\
1.676	-0.000101744\\
1.678	-0.000101353\\
1.68	-0.000100083\\
1.682	-0.00010224\\
1.684	-0.00010655\\
1.686	-0.000109762\\
1.688	-0.000111885\\
1.69	-0.00011289\\
1.692	-0.000112791\\
1.694	-0.000115375\\
1.696	-0.000126147\\
1.698	-0.000135355\\
1.7	-0.000142995\\
1.702	-0.000149063\\
1.704	-0.000153565\\
1.706	-0.000156508\\
1.708	-0.000157908\\
1.71	-0.000157784\\
1.712	-0.000156159\\
1.714	-0.000153063\\
1.716	-0.000153753\\
1.718	-0.00016792\\
1.72	-0.000180449\\
1.722	-0.000191323\\
1.724	-0.000200532\\
1.726	-0.00020807\\
1.728	-0.000213937\\
1.73	-0.000218138\\
1.732	-0.000220684\\
1.734	-0.000221591\\
1.736	-0.000220881\\
1.738	-0.00021858\\
1.74	-0.000214719\\
1.742	-0.000213594\\
1.744	-0.000213103\\
1.746	-0.000211077\\
1.748	-0.000208706\\
1.75	-0.000207362\\
1.752	-0.000204534\\
1.754	-0.000200256\\
1.756	-0.000194563\\
1.758	-0.000187495\\
1.76	-0.000179096\\
1.762	-0.000169414\\
1.764	-0.000158501\\
1.766	-0.000146411\\
1.768	-0.000133202\\
1.77	-0.000118936\\
1.772	-0.000103677\\
1.774	-8.74891e-05\\
1.776	-8.30307e-05\\
1.778	-8.44107e-05\\
1.78	-8.47886e-05\\
1.782	-8.41788e-05\\
1.784	-8.2599e-05\\
1.786	-8.00702e-05\\
1.788	-7.66165e-05\\
1.79	-7.22648e-05\\
1.792	-6.98205e-05\\
1.794	-8.2873e-05\\
1.796	-9.48783e-05\\
1.798	-0.000105811\\
1.8	-0.000115652\\
1.802	-0.000124382\\
1.804	-0.000131989\\
1.806	-0.000138464\\
1.808	-0.000143801\\
1.81	-0.000148\\
1.812	-0.000151063\\
1.814	-0.000152995\\
1.816	-0.000153808\\
1.818	-0.000153515\\
1.82	-0.000152132\\
1.822	-0.000149682\\
1.824	-0.000146188\\
1.826	-0.000151564\\
1.828	-0.000161303\\
1.83	-0.000169853\\
1.832	-0.000177203\\
1.834	-0.000183345\\
1.836	-0.000188276\\
1.838	-0.000191998\\
1.84	-0.000194514\\
1.842	-0.000195835\\
1.844	-0.000195972\\
1.846	-0.000194941\\
1.848	-0.000192764\\
1.85	-0.000189463\\
1.852	-0.000188204\\
1.854	-0.000187611\\
1.856	-0.000185911\\
1.858	-0.000183376\\
1.86	-0.000182327\\
1.862	-0.000180208\\
1.864	-0.000177041\\
1.866	-0.000172852\\
1.868	-0.00016767\\
1.87	-0.000161525\\
1.872	-0.000154453\\
1.874	-0.000146489\\
1.876	-0.000137673\\
1.878	-0.000128047\\
1.88	-0.000117654\\
1.882	-0.000106541\\
1.884	-9.47533e-05\\
1.886	-8.23415e-05\\
1.888	-6.93557e-05\\
1.89	-5.58477e-05\\
1.892	-4.18705e-05\\
1.894	-3.1653e-05\\
1.896	-2.85555e-05\\
1.898	-2.48257e-05\\
1.9	-2.53597e-05\\
1.902	-3.08491e-05\\
1.904	-3.98354e-05\\
1.906	-4.80423e-05\\
1.908	-5.54545e-05\\
1.91	-6.20592e-05\\
1.912	-6.78468e-05\\
1.914	-7.28104e-05\\
1.916	-7.6946e-05\\
1.918	-8.02523e-05\\
1.92	-8.2731e-05\\
1.922	-8.43865e-05\\
1.924	-8.52259e-05\\
1.926	-8.5259e-05\\
1.928	-8.4498e-05\\
1.93	-8.2958e-05\\
1.932	-8.06562e-05\\
1.934	-7.8691e-05\\
1.936	-8.65707e-05\\
1.938	-9.35946e-05\\
1.94	-9.9752e-05\\
1.942	-0.000105036\\
1.944	-0.000109441\\
1.946	-0.000112967\\
1.948	-0.000115614\\
1.95	-0.000117387\\
1.952	-0.000118293\\
1.954	-0.000118341\\
1.956	-0.000117545\\
1.958	-0.00011592\\
1.96	-0.000113482\\
1.962	-0.00011258\\
1.964	-0.000112292\\
1.966	-0.000111203\\
1.968	-0.000109328\\
1.97	-0.000108541\\
1.972	-0.000107257\\
1.974	-0.000105215\\
1.976	-0.000102432\\
1.978	-9.97478e-05\\
1.98	-9.95503e-05\\
1.982	-9.87644e-05\\
1.984	-9.74007e-05\\
1.986	-9.54718e-05\\
1.988	-9.29913e-05\\
1.99	-8.99746e-05\\
1.992	-8.64382e-05\\
1.994	-8.24002e-05\\
1.996	-7.78797e-05\\
1.998	-7.28969e-05\\
2	-6.74735e-05\\
};
\end{axis}

\begin{axis}[%
width={\figscale*0.642in},
height={\figscale*0.17in},
at={(\figscale*1.756in,\figscale*1.895in)},
scale only axis,
separate axis lines,
every outer x axis line/.append style={black},
every x tick label/.append style={font=\figticfont, yshift=\figinsettickshift},
every x tick/.append style={black},
xmin=0,
xmax=0.5,
xtick={  0, 0.4},
every outer y axis line/.append style={black},
every y tick label/.append style={font=\figticfont},
every y tick/.append style={black},
ymin=-0.0437676,
ymax=-0.0413803,
ytick={-0.0435, -0.0415},
axis background/.style={fill=white},
ylabel style={rotate=-90, at={(axis description cs:-0.145,0.5)}, anchor=east}
]
\addplot [color=mycolor1, line width=1.4pt, forget plot]
  table[row sep=crcr]{%
0	-0.017507\\
0.002	-0.017507\\
0.004	-0.017507\\
0.006	-0.017507\\
0.008	-0.017507\\
0.01	-0.017507\\
0.012	-0.017507\\
0.014	-0.017507\\
0.016	-0.017507\\
0.018	-0.017507\\
0.02	-0.017507\\
0.022	-0.017507\\
0.024	-0.017507\\
0.026	-0.017507\\
0.028	-0.017507\\
0.03	-0.017507\\
0.032	-0.0177659\\
0.034	-0.0183895\\
0.036	-0.0189935\\
0.038	-0.0195812\\
0.04	-0.0201555\\
0.042	-0.0207188\\
0.044	-0.0212731\\
0.046	-0.0218202\\
0.048	-0.0223611\\
0.05	-0.022897\\
0.052	-0.0234283\\
0.054	-0.0239556\\
0.056	-0.024479\\
0.058	-0.0249985\\
0.06	-0.0255139\\
0.062	-0.0260251\\
0.064	-0.0265317\\
0.066	-0.0270332\\
0.068	-0.0275293\\
0.07	-0.0280194\\
0.072	-0.0285032\\
0.074	-0.0289802\\
0.076	-0.0294501\\
0.078	-0.0299125\\
0.08	-0.0303672\\
0.082	-0.030814\\
0.084	-0.0312528\\
0.086	-0.0316834\\
0.088	-0.0321058\\
0.09	-0.0325201\\
0.092	-0.0329264\\
0.094	-0.0333248\\
0.096	-0.0337155\\
0.098	-0.0340986\\
0.1	-0.0344744\\
0.102	-0.0348431\\
0.104	-0.0352048\\
0.106	-0.0355598\\
0.108	-0.0359082\\
0.11	-0.0362501\\
0.112	-0.0365857\\
0.114	-0.036915\\
0.116	-0.037238\\
0.118	-0.0375547\\
0.12	-0.037865\\
0.122	-0.0381689\\
0.124	-0.038466\\
0.126	-0.0387562\\
0.128	-0.0390394\\
0.13	-0.0393151\\
0.132	-0.0395831\\
0.134	-0.0398431\\
0.136	-0.0400948\\
0.138	-0.0403377\\
0.14	-0.0405716\\
0.142	-0.0407961\\
0.144	-0.041011\\
0.146	-0.0412158\\
0.148	-0.0414105\\
0.15	-0.0415946\\
0.152	-0.0417681\\
0.154	-0.0419308\\
0.156	-0.0420826\\
0.158	-0.0422233\\
0.16	-0.0423532\\
0.162	-0.042472\\
0.164	-0.0425801\\
0.166	-0.0426775\\
0.168	-0.0427643\\
0.17	-0.042841\\
0.172	-0.0429077\\
0.174	-0.0429648\\
0.176	-0.0430126\\
0.178	-0.0430516\\
0.18	-0.0430821\\
0.182	-0.0431047\\
0.184	-0.0431198\\
0.186	-0.043128\\
0.188	-0.0431296\\
0.19	-0.0431252\\
0.192	-0.0431154\\
0.194	-0.0431006\\
0.196	-0.0430813\\
0.198	-0.0430581\\
0.2	-0.0430313\\
0.202	-0.0430016\\
0.204	-0.0429692\\
0.206	-0.0429348\\
0.208	-0.0428985\\
0.21	-0.042861\\
0.212	-0.0428224\\
0.214	-0.0427831\\
0.216	-0.0427435\\
0.218	-0.0427039\\
0.22	-0.0426644\\
0.222	-0.0426254\\
0.224	-0.042587\\
0.226	-0.0425495\\
0.228	-0.042513\\
0.23	-0.0424777\\
0.232	-0.0424437\\
0.234	-0.0424111\\
0.236	-0.04238\\
0.238	-0.0423505\\
0.24	-0.0423227\\
0.242	-0.0422966\\
0.244	-0.0422722\\
0.246	-0.0422497\\
0.248	-0.042229\\
0.25	-0.0422101\\
0.252	-0.0421931\\
0.254	-0.0421779\\
0.256	-0.0421646\\
0.258	-0.042153\\
0.26	-0.0421433\\
0.262	-0.0421354\\
0.264	-0.0421292\\
0.266	-0.0421248\\
0.268	-0.042122\\
0.27	-0.0421209\\
0.272	-0.0421214\\
0.274	-0.0421235\\
0.276	-0.0421271\\
0.278	-0.0421323\\
0.28	-0.0421388\\
0.282	-0.0421467\\
0.284	-0.042156\\
0.286	-0.0421666\\
0.288	-0.0421784\\
0.29	-0.0421913\\
0.292	-0.0422054\\
0.294	-0.0422206\\
0.296	-0.0422367\\
0.298	-0.0422538\\
0.3	-0.0422718\\
0.302	-0.0422906\\
0.304	-0.0423102\\
0.306	-0.0423305\\
0.308	-0.0423514\\
0.31	-0.042373\\
0.312	-0.0423951\\
0.314	-0.0424177\\
0.316	-0.0424407\\
0.318	-0.0424642\\
0.32	-0.0424879\\
0.322	-0.042512\\
0.324	-0.0425363\\
0.326	-0.0425608\\
0.328	-0.0425855\\
0.33	-0.0426102\\
0.332	-0.0426351\\
0.334	-0.04266\\
0.336	-0.042685\\
0.338	-0.0427099\\
0.34	-0.0427348\\
0.342	-0.0427596\\
0.344	-0.0427843\\
0.346	-0.0428089\\
0.348	-0.0428333\\
0.35	-0.0428576\\
0.352	-0.0428817\\
0.354	-0.0429057\\
0.356	-0.0429294\\
0.358	-0.0429529\\
0.36	-0.0429762\\
0.362	-0.0429992\\
0.364	-0.043022\\
0.366	-0.0430446\\
0.368	-0.0430668\\
0.37	-0.0430888\\
0.372	-0.0431106\\
0.374	-0.043132\\
0.376	-0.0431532\\
0.378	-0.043174\\
0.38	-0.0431946\\
0.382	-0.0432149\\
0.384	-0.0432348\\
0.386	-0.0432544\\
0.388	-0.0432737\\
0.39	-0.0432927\\
0.392	-0.0433114\\
0.394	-0.0433297\\
0.396	-0.0433476\\
0.398	-0.0433652\\
0.4	-0.0433825\\
0.402	-0.0433994\\
0.404	-0.0434159\\
0.406	-0.0434321\\
0.408	-0.0434478\\
0.41	-0.0434632\\
0.412	-0.0434782\\
0.414	-0.0434927\\
0.416	-0.0435069\\
0.418	-0.0435207\\
0.42	-0.043534\\
0.422	-0.0435469\\
0.424	-0.0435594\\
0.426	-0.0435715\\
0.428	-0.0435831\\
0.43	-0.0435942\\
0.432	-0.043605\\
0.434	-0.0436153\\
0.436	-0.0436251\\
0.438	-0.0436345\\
0.44	-0.0436434\\
0.442	-0.0436519\\
0.444	-0.04366\\
0.446	-0.0436676\\
0.448	-0.0436748\\
0.45	-0.0436815\\
0.452	-0.0436878\\
0.454	-0.0436937\\
0.456	-0.0436992\\
0.458	-0.0437043\\
0.46	-0.0437089\\
0.462	-0.0437132\\
0.464	-0.0437171\\
0.466	-0.0437206\\
0.468	-0.0437238\\
0.47	-0.0437266\\
0.472	-0.043729\\
0.474	-0.0437312\\
0.476	-0.043733\\
0.478	-0.0437346\\
0.48	-0.0437358\\
0.482	-0.0437368\\
0.484	-0.0437375\\
0.486	-0.043738\\
0.488	-0.0437383\\
0.49	-0.0437384\\
0.492	-0.0437382\\
0.494	-0.0437379\\
0.496	-0.0437375\\
0.498	-0.0437369\\
0.5	-0.0437361\\
};
\addplot [color=mycolor2, line width=1.0pt, forget plot]
  table[row sep=crcr]{%
0	-0.017507\\
0.002	-0.017507\\
0.004	-0.017507\\
0.006	-0.017507\\
0.008	-0.017507\\
0.01	-0.017507\\
0.012	-0.017507\\
0.014	-0.017507\\
0.016	-0.017507\\
0.018	-0.017507\\
0.02	-0.017507\\
0.022	-0.017507\\
0.024	-0.017507\\
0.026	-0.017507\\
0.028	-0.017507\\
0.03	-0.017507\\
0.032	-0.0177661\\
0.034	-0.0183898\\
0.036	-0.0189939\\
0.038	-0.0195817\\
0.04	-0.0201561\\
0.042	-0.0207196\\
0.044	-0.0212741\\
0.046	-0.0218214\\
0.048	-0.0223626\\
0.05	-0.0228988\\
0.052	-0.0234305\\
0.054	-0.0239582\\
0.056	-0.024482\\
0.058	-0.0250021\\
0.06	-0.0255182\\
0.062	-0.02603\\
0.064	-0.0265374\\
0.066	-0.0270397\\
0.068	-0.0275367\\
0.07	-0.0280279\\
0.072	-0.0285128\\
0.074	-0.0289911\\
0.076	-0.0294623\\
0.078	-0.0299262\\
0.08	-0.0303824\\
0.082	-0.0308308\\
0.084	-0.0312713\\
0.086	-0.0317038\\
0.088	-0.0321282\\
0.09	-0.0325446\\
0.092	-0.0329531\\
0.094	-0.0333538\\
0.096	-0.0337469\\
0.098	-0.0341325\\
0.1	-0.0345109\\
0.102	-0.0348823\\
0.104	-0.0352468\\
0.106	-0.0356047\\
0.108	-0.035956\\
0.11	-0.036301\\
0.112	-0.0366396\\
0.114	-0.036972\\
0.116	-0.0372982\\
0.118	-0.0376182\\
0.12	-0.0379317\\
0.122	-0.0382388\\
0.124	-0.0385392\\
0.126	-0.0388328\\
0.128	-0.0391192\\
0.13	-0.0393981\\
0.132	-0.0396694\\
0.134	-0.0399326\\
0.136	-0.0401874\\
0.138	-0.0404334\\
0.14	-0.0406703\\
0.142	-0.0408977\\
0.144	-0.0411154\\
0.146	-0.041323\\
0.148	-0.0415203\\
0.15	-0.0417069\\
0.152	-0.0418828\\
0.154	-0.0420477\\
0.156	-0.0422016\\
0.158	-0.0423443\\
0.16	-0.0424758\\
0.162	-0.0425963\\
0.164	-0.0427058\\
0.166	-0.0428043\\
0.168	-0.0428923\\
0.17	-0.0429697\\
0.172	-0.043037\\
0.174	-0.0430945\\
0.176	-0.0431426\\
0.178	-0.0431815\\
0.18	-0.0432118\\
0.182	-0.043234\\
0.184	-0.0432484\\
0.186	-0.0432556\\
0.188	-0.0432561\\
0.19	-0.0432504\\
0.192	-0.043239\\
0.194	-0.0432224\\
0.196	-0.0432011\\
0.198	-0.0431757\\
0.2	-0.0431465\\
0.202	-0.0431141\\
0.204	-0.043079\\
0.206	-0.0430415\\
0.208	-0.043002\\
0.21	-0.0429611\\
0.212	-0.0429189\\
0.214	-0.042876\\
0.216	-0.0428325\\
0.218	-0.0427888\\
0.22	-0.0427452\\
0.222	-0.0427019\\
0.224	-0.0426592\\
0.226	-0.0426172\\
0.228	-0.0425761\\
0.23	-0.0425361\\
0.232	-0.0424973\\
0.234	-0.0424599\\
0.236	-0.042424\\
0.238	-0.0423896\\
0.24	-0.0423569\\
0.242	-0.0423258\\
0.244	-0.0422966\\
0.246	-0.0422691\\
0.248	-0.0422435\\
0.25	-0.0422198\\
0.252	-0.0421979\\
0.254	-0.0421779\\
0.256	-0.0421598\\
0.258	-0.0421436\\
0.26	-0.0421293\\
0.262	-0.0421169\\
0.264	-0.0421063\\
0.266	-0.0420975\\
0.268	-0.0420906\\
0.27	-0.0420854\\
0.272	-0.0420819\\
0.274	-0.0420802\\
0.276	-0.0420801\\
0.278	-0.0420817\\
0.28	-0.0420848\\
0.282	-0.0420894\\
0.284	-0.0420956\\
0.286	-0.0421032\\
0.288	-0.0421121\\
0.29	-0.0421224\\
0.292	-0.0421339\\
0.294	-0.0421467\\
0.296	-0.0421606\\
0.298	-0.0421757\\
0.3	-0.0421917\\
0.302	-0.0422088\\
0.304	-0.0422267\\
0.306	-0.0422456\\
0.308	-0.0422652\\
0.31	-0.0422855\\
0.312	-0.0423066\\
0.314	-0.0423283\\
0.316	-0.0423505\\
0.318	-0.0423732\\
0.32	-0.0423964\\
0.322	-0.0424201\\
0.324	-0.042444\\
0.326	-0.0424683\\
0.328	-0.0424928\\
0.33	-0.0425176\\
0.332	-0.0425425\\
0.334	-0.0425675\\
0.336	-0.0425927\\
0.338	-0.0426178\\
0.34	-0.0426431\\
0.342	-0.0426683\\
0.344	-0.0426934\\
0.346	-0.0427185\\
0.348	-0.0427435\\
0.35	-0.0427684\\
0.352	-0.0427932\\
0.354	-0.0428178\\
0.356	-0.0428422\\
0.358	-0.0428664\\
0.36	-0.0428905\\
0.362	-0.0429143\\
0.364	-0.0429379\\
0.366	-0.0429612\\
0.368	-0.0429843\\
0.37	-0.0430071\\
0.372	-0.0430297\\
0.374	-0.0430519\\
0.376	-0.0430739\\
0.378	-0.0430956\\
0.38	-0.0431171\\
0.382	-0.0431382\\
0.384	-0.043159\\
0.386	-0.0431795\\
0.388	-0.0431996\\
0.39	-0.0432195\\
0.392	-0.043239\\
0.394	-0.0432582\\
0.396	-0.0432771\\
0.398	-0.0432956\\
0.4	-0.0433138\\
0.402	-0.0433316\\
0.404	-0.0433491\\
0.406	-0.0433662\\
0.408	-0.043383\\
0.41	-0.0433994\\
0.412	-0.0434154\\
0.414	-0.043431\\
0.416	-0.0434463\\
0.418	-0.0434612\\
0.42	-0.0434757\\
0.422	-0.0434898\\
0.424	-0.0435035\\
0.426	-0.0435169\\
0.428	-0.0435298\\
0.43	-0.0435424\\
0.432	-0.0435545\\
0.434	-0.0435663\\
0.436	-0.0435776\\
0.438	-0.0435886\\
0.44	-0.0435991\\
0.442	-0.0436093\\
0.444	-0.0436191\\
0.446	-0.0436284\\
0.448	-0.0436374\\
0.45	-0.043646\\
0.452	-0.0436543\\
0.454	-0.0436621\\
0.456	-0.0436696\\
0.458	-0.0436767\\
0.46	-0.0436835\\
0.462	-0.0436899\\
0.464	-0.0436959\\
0.466	-0.0437017\\
0.468	-0.0437071\\
0.47	-0.0437121\\
0.472	-0.0437169\\
0.474	-0.0437214\\
0.476	-0.0437255\\
0.478	-0.0437294\\
0.48	-0.0437331\\
0.482	-0.0437364\\
0.484	-0.0437395\\
0.486	-0.0437424\\
0.488	-0.043745\\
0.49	-0.0437474\\
0.492	-0.0437496\\
0.494	-0.0437515\\
0.496	-0.0437533\\
0.498	-0.0437549\\
0.5	-0.0437563\\
};
\addplot [color=mycolor3, line width=1.0pt, forget plot]
  table[row sep=crcr]{%
0	-0.017507\\
0.002	-0.017507\\
0.004	-0.017507\\
0.006	-0.017507\\
0.008	-0.017507\\
0.01	-0.017507\\
0.012	-0.017507\\
0.014	-0.017507\\
0.016	-0.017507\\
0.018	-0.017507\\
0.02	-0.017507\\
0.022	-0.017507\\
0.024	-0.017507\\
0.026	-0.017507\\
0.028	-0.017507\\
0.03	-0.017507\\
0.032	-0.0177654\\
0.034	-0.0183889\\
0.036	-0.0189928\\
0.038	-0.0195803\\
0.04	-0.0201545\\
0.042	-0.0207176\\
0.044	-0.0212719\\
0.046	-0.0218187\\
0.048	-0.0223595\\
0.05	-0.0228952\\
0.052	-0.0234263\\
0.054	-0.0239534\\
0.056	-0.0244765\\
0.058	-0.0249958\\
0.06	-0.025511\\
0.062	-0.026022\\
0.064	-0.0265283\\
0.066	-0.0270296\\
0.068	-0.0275254\\
0.07	-0.0280153\\
0.072	-0.0284988\\
0.074	-0.0289756\\
0.076	-0.0294452\\
0.078	-0.0299074\\
0.08	-0.0303618\\
0.082	-0.0308083\\
0.084	-0.0312467\\
0.086	-0.031677\\
0.088	-0.0320991\\
0.09	-0.0325131\\
0.092	-0.0329191\\
0.094	-0.0333171\\
0.096	-0.0337074\\
0.098	-0.0340902\\
0.1	-0.0344655\\
0.102	-0.0348338\\
0.104	-0.035195\\
0.106	-0.0355495\\
0.108	-0.0358973\\
0.11	-0.0362387\\
0.112	-0.0365737\\
0.114	-0.0369024\\
0.116	-0.0372247\\
0.118	-0.0375407\\
0.12	-0.0378503\\
0.122	-0.0381533\\
0.124	-0.0384496\\
0.126	-0.038739\\
0.128	-0.0390212\\
0.13	-0.039296\\
0.132	-0.0395631\\
0.134	-0.0398221\\
0.136	-0.0400727\\
0.138	-0.0403146\\
0.14	-0.0405475\\
0.142	-0.0407709\\
0.144	-0.0409847\\
0.146	-0.0411885\\
0.148	-0.041382\\
0.15	-0.0415651\\
0.152	-0.0417375\\
0.154	-0.0418991\\
0.156	-0.0420499\\
0.158	-0.0421896\\
0.16	-0.0423185\\
0.162	-0.0424364\\
0.164	-0.0425436\\
0.166	-0.0426402\\
0.168	-0.0427263\\
0.17	-0.0428022\\
0.172	-0.0428683\\
0.174	-0.0429249\\
0.176	-0.0429723\\
0.178	-0.0430109\\
0.18	-0.0430412\\
0.182	-0.0430637\\
0.184	-0.0430788\\
0.186	-0.0430871\\
0.188	-0.0430889\\
0.19	-0.043085\\
0.192	-0.0430757\\
0.194	-0.0430615\\
0.196	-0.0430431\\
0.198	-0.0430209\\
0.2	-0.0429953\\
0.202	-0.0429669\\
0.204	-0.0429361\\
0.206	-0.0429033\\
0.208	-0.0428689\\
0.21	-0.0428333\\
0.212	-0.042797\\
0.214	-0.0427601\\
0.216	-0.0427231\\
0.218	-0.0426863\\
0.22	-0.0426498\\
0.222	-0.0426139\\
0.224	-0.0425789\\
0.226	-0.0425449\\
0.228	-0.0425121\\
0.23	-0.0424806\\
0.232	-0.0424507\\
0.234	-0.0424223\\
0.236	-0.0423956\\
0.238	-0.0423707\\
0.24	-0.0423476\\
0.242	-0.0423263\\
0.244	-0.042307\\
0.246	-0.0422896\\
0.248	-0.0422742\\
0.25	-0.0422607\\
0.252	-0.0422492\\
0.254	-0.0422396\\
0.256	-0.0422319\\
0.258	-0.0422261\\
0.26	-0.0422222\\
0.262	-0.0422201\\
0.264	-0.0422199\\
0.266	-0.0422213\\
0.268	-0.0422245\\
0.27	-0.0422293\\
0.272	-0.0422358\\
0.274	-0.0422438\\
0.276	-0.0422532\\
0.278	-0.0422641\\
0.28	-0.0422764\\
0.282	-0.0422899\\
0.284	-0.0423047\\
0.286	-0.0423206\\
0.288	-0.0423377\\
0.29	-0.0423557\\
0.292	-0.0423747\\
0.294	-0.0423946\\
0.296	-0.0424153\\
0.298	-0.0424367\\
0.3	-0.0424588\\
0.302	-0.0424815\\
0.304	-0.0425047\\
0.306	-0.0425284\\
0.308	-0.0425524\\
0.31	-0.0425767\\
0.312	-0.0426013\\
0.314	-0.0426261\\
0.316	-0.042651\\
0.318	-0.042676\\
0.32	-0.0427009\\
0.322	-0.0427258\\
0.324	-0.0427507\\
0.326	-0.0427753\\
0.328	-0.0427998\\
0.33	-0.042824\\
0.332	-0.042848\\
0.334	-0.0428716\\
0.336	-0.0428949\\
0.338	-0.0429178\\
0.34	-0.0429403\\
0.342	-0.0429624\\
0.344	-0.042984\\
0.346	-0.0430052\\
0.348	-0.0430259\\
0.35	-0.0430461\\
0.352	-0.0430658\\
0.354	-0.0430849\\
0.356	-0.0431036\\
0.358	-0.0431217\\
0.36	-0.0431393\\
0.362	-0.0431565\\
0.364	-0.043173\\
0.366	-0.0431891\\
0.368	-0.0432047\\
0.37	-0.0432198\\
0.372	-0.0432344\\
0.374	-0.0432486\\
0.376	-0.0432622\\
0.378	-0.0432755\\
0.38	-0.0432883\\
0.382	-0.0433006\\
0.384	-0.0433126\\
0.386	-0.0433241\\
0.388	-0.0433353\\
0.39	-0.043346\\
0.392	-0.0433565\\
0.394	-0.0433665\\
0.396	-0.0433763\\
0.398	-0.0433857\\
0.4	-0.0433948\\
0.402	-0.0434036\\
0.404	-0.0434121\\
0.406	-0.0434203\\
0.408	-0.0434282\\
0.41	-0.0434359\\
0.412	-0.0434434\\
0.414	-0.0434506\\
0.416	-0.0434576\\
0.418	-0.0434643\\
0.42	-0.0434709\\
0.422	-0.0434772\\
0.424	-0.0434834\\
0.426	-0.0434893\\
0.428	-0.0434951\\
0.43	-0.0435006\\
0.432	-0.043506\\
0.434	-0.0435113\\
0.436	-0.0435163\\
0.438	-0.0435212\\
0.44	-0.043526\\
0.442	-0.0435306\\
0.444	-0.0435351\\
0.446	-0.0435394\\
0.448	-0.0435436\\
0.45	-0.0435476\\
0.452	-0.0435515\\
0.454	-0.0435553\\
0.456	-0.043559\\
0.458	-0.0435626\\
0.46	-0.043566\\
0.462	-0.0435694\\
0.464	-0.0435726\\
0.466	-0.0435757\\
0.468	-0.0435788\\
0.47	-0.0435817\\
0.472	-0.0435846\\
0.474	-0.0435873\\
0.476	-0.04359\\
0.478	-0.0435926\\
0.48	-0.0435952\\
0.482	-0.0435976\\
0.484	-0.0436\\
0.486	-0.0436023\\
0.488	-0.0436046\\
0.49	-0.0436068\\
0.492	-0.0436089\\
0.494	-0.043611\\
0.496	-0.043613\\
0.498	-0.0436149\\
0.5	-0.0436169\\
};
\addplot [color=mycolor4, line width=1.0pt, forget plot]
  table[row sep=crcr]{%
0	-0.017507\\
0.002	-0.017507\\
0.004	-0.017507\\
0.006	-0.017507\\
0.008	-0.017507\\
0.01	-0.017507\\
0.012	-0.017507\\
0.014	-0.017507\\
0.016	-0.017507\\
0.018	-0.017507\\
0.02	-0.017507\\
0.022	-0.017507\\
0.024	-0.017507\\
0.026	-0.017507\\
0.028	-0.017507\\
0.03	-0.017507\\
0.032	-0.0177653\\
0.034	-0.0183889\\
0.036	-0.0189927\\
0.038	-0.0195803\\
0.04	-0.0201545\\
0.042	-0.0207177\\
0.044	-0.0212719\\
0.046	-0.0218189\\
0.048	-0.0223597\\
0.05	-0.0228955\\
0.052	-0.0234268\\
0.054	-0.023954\\
0.056	-0.0244774\\
0.058	-0.0249968\\
0.06	-0.0255123\\
0.062	-0.0260236\\
0.064	-0.0265302\\
0.066	-0.0270319\\
0.068	-0.0275281\\
0.07	-0.0280185\\
0.072	-0.0285025\\
0.074	-0.0289799\\
0.076	-0.0294501\\
0.078	-0.0299129\\
0.08	-0.0303681\\
0.082	-0.0308154\\
0.084	-0.0312547\\
0.086	-0.0316858\\
0.088	-0.0321089\\
0.09	-0.0325239\\
0.092	-0.0329309\\
0.094	-0.03333\\
0.096	-0.0337214\\
0.098	-0.0341053\\
0.1	-0.0344818\\
0.102	-0.0348513\\
0.104	-0.0352137\\
0.106	-0.0355694\\
0.108	-0.0359185\\
0.11	-0.0362611\\
0.112	-0.0365973\\
0.114	-0.0369271\\
0.116	-0.0372507\\
0.118	-0.0375678\\
0.12	-0.0378785\\
0.122	-0.0381825\\
0.124	-0.0384798\\
0.126	-0.0387701\\
0.128	-0.0390532\\
0.13	-0.0393288\\
0.132	-0.0395965\\
0.134	-0.0398561\\
0.136	-0.0401072\\
0.138	-0.0403495\\
0.14	-0.0405826\\
0.142	-0.0408062\\
0.144	-0.04102\\
0.146	-0.0412237\\
0.148	-0.041417\\
0.15	-0.0415996\\
0.152	-0.0417715\\
0.154	-0.0419325\\
0.156	-0.0420824\\
0.158	-0.0422213\\
0.16	-0.042349\\
0.162	-0.0424657\\
0.164	-0.0425716\\
0.166	-0.0426666\\
0.168	-0.0427511\\
0.17	-0.0428253\\
0.172	-0.0428895\\
0.174	-0.0429441\\
0.176	-0.0429894\\
0.178	-0.0430259\\
0.18	-0.0430539\\
0.182	-0.0430741\\
0.184	-0.0430868\\
0.186	-0.0430925\\
0.188	-0.0430919\\
0.19	-0.0430853\\
0.192	-0.0430734\\
0.194	-0.0430567\\
0.196	-0.0430356\\
0.198	-0.0430107\\
0.2	-0.0429826\\
0.202	-0.0429516\\
0.204	-0.0429182\\
0.206	-0.0428829\\
0.208	-0.0428462\\
0.21	-0.0428083\\
0.212	-0.0427697\\
0.214	-0.0427307\\
0.216	-0.0426917\\
0.218	-0.042653\\
0.22	-0.0426147\\
0.222	-0.0425773\\
0.224	-0.0425408\\
0.226	-0.0425055\\
0.228	-0.0424716\\
0.23	-0.0424393\\
0.232	-0.0424086\\
0.234	-0.0423797\\
0.236	-0.0423526\\
0.238	-0.0423276\\
0.24	-0.0423045\\
0.242	-0.0422836\\
0.244	-0.0422648\\
0.246	-0.0422481\\
0.248	-0.0422336\\
0.25	-0.0422212\\
0.252	-0.042211\\
0.254	-0.0422029\\
0.256	-0.042197\\
0.258	-0.0421931\\
0.26	-0.0421913\\
0.262	-0.0421915\\
0.264	-0.0421937\\
0.266	-0.0421978\\
0.268	-0.0422037\\
0.27	-0.0422115\\
0.272	-0.0422209\\
0.274	-0.0422321\\
0.276	-0.0422448\\
0.278	-0.042259\\
0.28	-0.0422747\\
0.282	-0.0422917\\
0.284	-0.04231\\
0.286	-0.0423295\\
0.288	-0.0423501\\
0.29	-0.0423717\\
0.292	-0.0423943\\
0.294	-0.0424177\\
0.296	-0.0424419\\
0.298	-0.0424667\\
0.3	-0.0424922\\
0.302	-0.0425181\\
0.304	-0.0425444\\
0.306	-0.0425711\\
0.308	-0.042598\\
0.31	-0.0426251\\
0.312	-0.0426523\\
0.314	-0.0426795\\
0.316	-0.0427066\\
0.318	-0.0427336\\
0.32	-0.0427604\\
0.322	-0.042787\\
0.324	-0.0428132\\
0.326	-0.0428391\\
0.328	-0.0428645\\
0.33	-0.0428894\\
0.332	-0.0429139\\
0.334	-0.0429378\\
0.336	-0.042961\\
0.338	-0.0429837\\
0.34	-0.0430057\\
0.342	-0.0430271\\
0.344	-0.0430477\\
0.346	-0.0430676\\
0.348	-0.0430869\\
0.35	-0.0431054\\
0.352	-0.0431232\\
0.354	-0.0431402\\
0.356	-0.0431566\\
0.358	-0.0431722\\
0.36	-0.0431871\\
0.362	-0.0432013\\
0.364	-0.0432148\\
0.366	-0.0432277\\
0.368	-0.04324\\
0.37	-0.0432516\\
0.372	-0.0432626\\
0.374	-0.043273\\
0.376	-0.0432829\\
0.378	-0.0432922\\
0.38	-0.0433011\\
0.382	-0.0433095\\
0.384	-0.0433174\\
0.386	-0.0433249\\
0.388	-0.0433321\\
0.39	-0.0433388\\
0.392	-0.0433453\\
0.394	-0.0433514\\
0.396	-0.0433572\\
0.398	-0.0433628\\
0.4	-0.0433681\\
0.402	-0.0433733\\
0.404	-0.0433782\\
0.406	-0.043383\\
0.408	-0.0433877\\
0.41	-0.0433922\\
0.412	-0.0433966\\
0.414	-0.043401\\
0.416	-0.0434052\\
0.418	-0.0434095\\
0.42	-0.0434136\\
0.422	-0.0434178\\
0.424	-0.043422\\
0.426	-0.0434261\\
0.428	-0.0434303\\
0.43	-0.0434344\\
0.432	-0.0434386\\
0.434	-0.0434429\\
0.436	-0.0434471\\
0.438	-0.0434515\\
0.44	-0.0434558\\
0.442	-0.0434602\\
0.444	-0.0434647\\
0.446	-0.0434692\\
0.448	-0.0434738\\
0.45	-0.0434784\\
0.452	-0.043483\\
0.454	-0.0434878\\
0.456	-0.0434925\\
0.458	-0.0434973\\
0.46	-0.0435022\\
0.462	-0.043507\\
0.464	-0.043512\\
0.466	-0.0435169\\
0.468	-0.0435218\\
0.47	-0.0435268\\
0.472	-0.0435318\\
0.474	-0.0435368\\
0.476	-0.0435418\\
0.478	-0.0435467\\
0.48	-0.0435517\\
0.482	-0.0435566\\
0.484	-0.0435615\\
0.486	-0.0435664\\
0.488	-0.0435712\\
0.49	-0.043576\\
0.492	-0.0435807\\
0.494	-0.0435854\\
0.496	-0.04359\\
0.498	-0.0435946\\
0.5	-0.043599\\
};
\addplot [color=mycolor5, line width=1.0pt, forget plot]
  table[row sep=crcr]{%
0	-0.017507\\
0.002	-0.017507\\
0.004	-0.017507\\
0.006	-0.017507\\
0.008	-0.017507\\
0.01	-0.017507\\
0.012	-0.017507\\
0.014	-0.017507\\
0.016	-0.017507\\
0.018	-0.017507\\
0.02	-0.017507\\
0.022	-0.017507\\
0.024	-0.017507\\
0.026	-0.017507\\
0.028	-0.017507\\
0.03	-0.017507\\
0.032	-0.0177638\\
0.034	-0.0183869\\
0.036	-0.0189904\\
0.038	-0.0195775\\
0.04	-0.0201511\\
0.042	-0.0207138\\
0.044	-0.0212674\\
0.046	-0.0218136\\
0.048	-0.0223537\\
0.05	-0.0228886\\
0.052	-0.0234189\\
0.054	-0.0239452\\
0.056	-0.0244675\\
0.058	-0.0249858\\
0.06	-0.0255001\\
0.062	-0.0260101\\
0.064	-0.0265154\\
0.066	-0.0270157\\
0.068	-0.0275105\\
0.07	-0.0279994\\
0.072	-0.0284819\\
0.074	-0.0289577\\
0.076	-0.0294264\\
0.078	-0.0298877\\
0.08	-0.0303412\\
0.082	-0.0307869\\
0.084	-0.0312246\\
0.086	-0.0316542\\
0.088	-0.0320757\\
0.09	-0.0324892\\
0.092	-0.0328947\\
0.094	-0.0332924\\
0.096	-0.0336824\\
0.098	-0.034065\\
0.1	-0.0344403\\
0.102	-0.0348086\\
0.104	-0.03517\\
0.106	-0.0355248\\
0.108	-0.0358731\\
0.11	-0.0362149\\
0.112	-0.0365506\\
0.114	-0.03688\\
0.116	-0.0372032\\
0.118	-0.0375201\\
0.12	-0.0378308\\
0.122	-0.038135\\
0.124	-0.0384325\\
0.126	-0.0387233\\
0.128	-0.039007\\
0.13	-0.0392833\\
0.132	-0.039552\\
0.134	-0.0398127\\
0.136	-0.0400651\\
0.138	-0.0403088\\
0.14	-0.0405435\\
0.142	-0.0407689\\
0.144	-0.0409846\\
0.146	-0.0411903\\
0.148	-0.0413859\\
0.15	-0.041571\\
0.152	-0.0417454\\
0.154	-0.0419091\\
0.156	-0.0420618\\
0.158	-0.0422036\\
0.16	-0.0423345\\
0.162	-0.0424545\\
0.164	-0.0425636\\
0.166	-0.0426621\\
0.168	-0.0427501\\
0.17	-0.0428279\\
0.172	-0.0428959\\
0.174	-0.0429542\\
0.176	-0.0430033\\
0.178	-0.0430436\\
0.18	-0.0430756\\
0.182	-0.0430996\\
0.184	-0.0431162\\
0.186	-0.0431259\\
0.188	-0.0431292\\
0.19	-0.0431265\\
0.192	-0.0431184\\
0.194	-0.0431055\\
0.196	-0.0430882\\
0.198	-0.0430669\\
0.2	-0.0430423\\
0.202	-0.0430148\\
0.204	-0.0429848\\
0.206	-0.0429527\\
0.208	-0.042919\\
0.21	-0.0428841\\
0.212	-0.0428483\\
0.214	-0.042812\\
0.216	-0.0427754\\
0.218	-0.0427389\\
0.22	-0.0427027\\
0.222	-0.0426671\\
0.224	-0.0426322\\
0.226	-0.0425984\\
0.228	-0.0425657\\
0.23	-0.0425343\\
0.232	-0.0425044\\
0.234	-0.042476\\
0.236	-0.0424492\\
0.238	-0.0424242\\
0.24	-0.042401\\
0.242	-0.0423796\\
0.244	-0.0423601\\
0.246	-0.0423425\\
0.248	-0.0423269\\
0.25	-0.0423132\\
0.252	-0.0423014\\
0.254	-0.0422916\\
0.256	-0.0422837\\
0.258	-0.0422776\\
0.26	-0.0422735\\
0.262	-0.0422712\\
0.264	-0.0422706\\
0.266	-0.0422718\\
0.268	-0.0422748\\
0.27	-0.0422793\\
0.272	-0.0422855\\
0.274	-0.0422932\\
0.276	-0.0423024\\
0.278	-0.0423131\\
0.28	-0.042325\\
0.282	-0.0423383\\
0.284	-0.0423528\\
0.286	-0.0423685\\
0.288	-0.0423852\\
0.29	-0.0424029\\
0.292	-0.0424216\\
0.294	-0.0424411\\
0.296	-0.0424614\\
0.298	-0.0424824\\
0.3	-0.0425041\\
0.302	-0.0425263\\
0.304	-0.0425489\\
0.306	-0.042572\\
0.308	-0.0425954\\
0.31	-0.0426191\\
0.312	-0.042643\\
0.314	-0.042667\\
0.316	-0.0426911\\
0.318	-0.0427151\\
0.32	-0.0427391\\
0.322	-0.042763\\
0.324	-0.0427867\\
0.326	-0.0428102\\
0.328	-0.0428334\\
0.33	-0.0428563\\
0.332	-0.0428788\\
0.334	-0.042901\\
0.336	-0.0429227\\
0.338	-0.042944\\
0.34	-0.0429647\\
0.342	-0.042985\\
0.344	-0.0430047\\
0.346	-0.0430239\\
0.348	-0.0430426\\
0.35	-0.0430607\\
0.352	-0.0430782\\
0.354	-0.0430951\\
0.356	-0.0431115\\
0.358	-0.0431273\\
0.36	-0.0431425\\
0.362	-0.0431572\\
0.364	-0.0431713\\
0.366	-0.0431848\\
0.368	-0.0431979\\
0.37	-0.0432104\\
0.372	-0.0432224\\
0.374	-0.0432339\\
0.376	-0.043245\\
0.378	-0.0432556\\
0.38	-0.0432658\\
0.382	-0.0432756\\
0.384	-0.043285\\
0.386	-0.0432941\\
0.388	-0.0433027\\
0.39	-0.0433111\\
0.392	-0.0433191\\
0.394	-0.0433269\\
0.396	-0.0433344\\
0.398	-0.0433416\\
0.4	-0.0433486\\
0.402	-0.0433554\\
0.404	-0.0433619\\
0.406	-0.0433683\\
0.408	-0.0433746\\
0.41	-0.0433806\\
0.412	-0.0433866\\
0.414	-0.0433924\\
0.416	-0.0433981\\
0.418	-0.0434036\\
0.42	-0.0434091\\
0.422	-0.0434146\\
0.424	-0.0434199\\
0.426	-0.0434252\\
0.428	-0.0434304\\
0.43	-0.0434355\\
0.432	-0.0434406\\
0.434	-0.0434457\\
0.436	-0.0434507\\
0.438	-0.0434557\\
0.44	-0.0434607\\
0.442	-0.0434656\\
0.444	-0.0434705\\
0.446	-0.0434754\\
0.448	-0.0434803\\
0.45	-0.0434851\\
0.452	-0.0434899\\
0.454	-0.0434947\\
0.456	-0.0434995\\
0.458	-0.0435043\\
0.46	-0.043509\\
0.462	-0.0435137\\
0.464	-0.0435184\\
0.466	-0.043523\\
0.468	-0.0435276\\
0.47	-0.0435322\\
0.472	-0.0435367\\
0.474	-0.0435412\\
0.476	-0.0435457\\
0.478	-0.0435501\\
0.48	-0.0435544\\
0.482	-0.0435587\\
0.484	-0.043563\\
0.486	-0.0435672\\
0.488	-0.0435713\\
0.49	-0.0435754\\
0.492	-0.0435795\\
0.494	-0.0435834\\
0.496	-0.0435873\\
0.498	-0.0435911\\
0.5	-0.0435949\\
};
\addplot [color=mycolor6, line width=1.0pt, forget plot]
  table[row sep=crcr]{%
0	-0.017507\\
0.002	-0.017507\\
0.004	-0.017507\\
0.006	-0.017507\\
0.008	-0.017507\\
0.01	-0.017507\\
0.012	-0.017507\\
0.014	-0.017507\\
0.016	-0.017507\\
0.018	-0.017507\\
0.02	-0.017507\\
0.022	-0.017507\\
0.024	-0.017507\\
0.026	-0.017507\\
0.028	-0.017507\\
0.03	-0.017507\\
0.032	-0.0177615\\
0.034	-0.0183842\\
0.036	-0.018987\\
0.038	-0.0195735\\
0.04	-0.0201463\\
0.042	-0.0207081\\
0.044	-0.0212607\\
0.046	-0.0218059\\
0.048	-0.0223447\\
0.05	-0.0228783\\
0.052	-0.0234072\\
0.054	-0.0239319\\
0.056	-0.0244525\\
0.058	-0.024969\\
0.06	-0.0254814\\
0.062	-0.0259893\\
0.064	-0.0264924\\
0.066	-0.0269904\\
0.068	-0.0274828\\
0.07	-0.0279691\\
0.072	-0.028449\\
0.074	-0.0289221\\
0.076	-0.029388\\
0.078	-0.0298464\\
0.08	-0.030297\\
0.082	-0.0307397\\
0.084	-0.0311744\\
0.086	-0.0316009\\
0.088	-0.0320194\\
0.09	-0.0324298\\
0.092	-0.0328323\\
0.094	-0.033227\\
0.096	-0.0336141\\
0.098	-0.0339938\\
0.1	-0.0343664\\
0.102	-0.034732\\
0.104	-0.0350909\\
0.106	-0.0354433\\
0.108	-0.0357893\\
0.11	-0.0361291\\
0.112	-0.0364628\\
0.114	-0.0367905\\
0.116	-0.0371122\\
0.118	-0.0374279\\
0.12	-0.0377375\\
0.122	-0.038041\\
0.124	-0.038338\\
0.126	-0.0386285\\
0.128	-0.0389123\\
0.13	-0.0391889\\
0.132	-0.0394582\\
0.134	-0.0397197\\
0.136	-0.0399733\\
0.138	-0.0402185\\
0.14	-0.0404549\\
0.142	-0.0406824\\
0.144	-0.0409004\\
0.146	-0.0411088\\
0.148	-0.0413073\\
0.15	-0.0414956\\
0.152	-0.0416735\\
0.154	-0.0418409\\
0.156	-0.0419976\\
0.158	-0.0421436\\
0.16	-0.0422789\\
0.162	-0.0424035\\
0.164	-0.0425175\\
0.166	-0.042621\\
0.168	-0.0427143\\
0.17	-0.0427974\\
0.172	-0.0428709\\
0.174	-0.0429348\\
0.176	-0.0429896\\
0.178	-0.0430358\\
0.18	-0.0430736\\
0.182	-0.0431035\\
0.184	-0.0431261\\
0.186	-0.0431417\\
0.188	-0.0431509\\
0.19	-0.0431542\\
0.192	-0.043152\\
0.194	-0.0431448\\
0.196	-0.0431332\\
0.198	-0.0431176\\
0.2	-0.0430985\\
0.202	-0.0430763\\
0.204	-0.0430515\\
0.206	-0.0430244\\
0.208	-0.0429956\\
0.21	-0.0429653\\
0.212	-0.042934\\
0.214	-0.0429018\\
0.216	-0.0428693\\
0.218	-0.0428365\\
0.22	-0.0428038\\
0.222	-0.0427714\\
0.224	-0.0427396\\
0.226	-0.0427084\\
0.228	-0.0426781\\
0.23	-0.0426489\\
0.232	-0.0426207\\
0.234	-0.0425939\\
0.236	-0.0425684\\
0.238	-0.0425443\\
0.24	-0.0425218\\
0.242	-0.0425008\\
0.244	-0.0424814\\
0.246	-0.0424636\\
0.248	-0.0424475\\
0.25	-0.042433\\
0.252	-0.0424202\\
0.254	-0.0424091\\
0.256	-0.0423997\\
0.258	-0.0423919\\
0.26	-0.0423857\\
0.262	-0.0423811\\
0.264	-0.0423781\\
0.266	-0.0423767\\
0.268	-0.0423768\\
0.27	-0.0423783\\
0.272	-0.0423813\\
0.274	-0.0423857\\
0.276	-0.0423914\\
0.278	-0.0423984\\
0.28	-0.0424066\\
0.282	-0.0424161\\
0.284	-0.0424267\\
0.286	-0.0424383\\
0.288	-0.042451\\
0.29	-0.0424646\\
0.292	-0.0424791\\
0.294	-0.0424945\\
0.296	-0.0425107\\
0.298	-0.0425275\\
0.3	-0.042545\\
0.302	-0.0425631\\
0.304	-0.0425817\\
0.306	-0.0426008\\
0.308	-0.0426203\\
0.31	-0.0426401\\
0.312	-0.0426602\\
0.314	-0.0426805\\
0.316	-0.0427009\\
0.318	-0.0427215\\
0.32	-0.0427421\\
0.322	-0.0427628\\
0.324	-0.0427834\\
0.326	-0.0428039\\
0.328	-0.0428243\\
0.33	-0.0428445\\
0.332	-0.0428645\\
0.334	-0.0428843\\
0.336	-0.0429038\\
0.338	-0.042923\\
0.34	-0.0429418\\
0.342	-0.0429604\\
0.344	-0.0429785\\
0.346	-0.0429963\\
0.348	-0.0430137\\
0.35	-0.0430307\\
0.352	-0.0430472\\
0.354	-0.0430634\\
0.356	-0.0430791\\
0.358	-0.0430944\\
0.36	-0.0431093\\
0.362	-0.0431237\\
0.364	-0.0431377\\
0.366	-0.0431514\\
0.368	-0.0431646\\
0.37	-0.0431774\\
0.372	-0.0431898\\
0.374	-0.0432019\\
0.376	-0.0432136\\
0.378	-0.0432249\\
0.38	-0.0432359\\
0.382	-0.0432465\\
0.384	-0.0432569\\
0.386	-0.043267\\
0.388	-0.0432767\\
0.39	-0.0432862\\
0.392	-0.0432954\\
0.394	-0.0433044\\
0.396	-0.0433132\\
0.398	-0.0433217\\
0.4	-0.04333\\
0.402	-0.0433381\\
0.404	-0.043346\\
0.406	-0.0433537\\
0.408	-0.0433613\\
0.41	-0.0433687\\
0.412	-0.0433759\\
0.414	-0.043383\\
0.416	-0.04339\\
0.418	-0.0433968\\
0.42	-0.0434035\\
0.422	-0.0434101\\
0.424	-0.0434166\\
0.426	-0.043423\\
0.428	-0.0434292\\
0.43	-0.0434354\\
0.432	-0.0434415\\
0.434	-0.0434474\\
0.436	-0.0434533\\
0.438	-0.0434591\\
0.44	-0.0434648\\
0.442	-0.0434704\\
0.444	-0.043476\\
0.446	-0.0434814\\
0.448	-0.0434868\\
0.45	-0.043492\\
0.452	-0.0434972\\
0.454	-0.0435023\\
0.456	-0.0435074\\
0.458	-0.0435123\\
0.46	-0.0435172\\
0.462	-0.043522\\
0.464	-0.0435267\\
0.466	-0.0435313\\
0.468	-0.0435359\\
0.47	-0.0435403\\
0.472	-0.0435447\\
0.474	-0.043549\\
0.476	-0.0435532\\
0.478	-0.0435573\\
0.48	-0.0435614\\
0.482	-0.0435654\\
0.484	-0.0435692\\
0.486	-0.043573\\
0.488	-0.0435768\\
0.49	-0.0435804\\
0.492	-0.0435839\\
0.494	-0.0435874\\
0.496	-0.0435908\\
0.498	-0.0435941\\
0.5	-0.0435973\\
};
\addplot [color=mycolor7, line width=1.0pt, forget plot]
  table[row sep=crcr]{%
0	-0.017507\\
0.002	-0.017507\\
0.004	-0.017507\\
0.006	-0.017507\\
0.008	-0.017507\\
0.01	-0.017507\\
0.012	-0.017507\\
0.014	-0.017507\\
0.016	-0.017507\\
0.018	-0.017507\\
0.02	-0.017507\\
0.022	-0.017507\\
0.024	-0.017507\\
0.026	-0.017507\\
0.028	-0.017507\\
0.03	-0.017507\\
0.032	-0.0177605\\
0.034	-0.0183829\\
0.036	-0.0189855\\
0.038	-0.0195717\\
0.04	-0.0201442\\
0.042	-0.0207055\\
0.044	-0.0212577\\
0.046	-0.0218024\\
0.048	-0.0223407\\
0.05	-0.0228737\\
0.052	-0.0234019\\
0.054	-0.0239258\\
0.056	-0.0244457\\
0.058	-0.0249614\\
0.06	-0.0254728\\
0.062	-0.0259798\\
0.064	-0.0264819\\
0.066	-0.0269787\\
0.068	-0.02747\\
0.07	-0.0279551\\
0.072	-0.0284338\\
0.074	-0.0289055\\
0.076	-0.0293701\\
0.078	-0.029827\\
0.08	-0.0302762\\
0.082	-0.0307174\\
0.084	-0.0311506\\
0.086	-0.0315756\\
0.088	-0.0319925\\
0.09	-0.0324013\\
0.092	-0.0328022\\
0.094	-0.0331953\\
0.096	-0.0335808\\
0.098	-0.033959\\
0.1	-0.03433\\
0.102	-0.034694\\
0.104	-0.0350514\\
0.106	-0.0354023\\
0.108	-0.0357469\\
0.11	-0.0360853\\
0.112	-0.0364177\\
0.114	-0.0367441\\
0.116	-0.0370647\\
0.118	-0.0373793\\
0.12	-0.0376879\\
0.122	-0.0379904\\
0.124	-0.0382866\\
0.126	-0.0385764\\
0.128	-0.0388596\\
0.13	-0.0391357\\
0.132	-0.0394046\\
0.134	-0.039666\\
0.136	-0.0399195\\
0.138	-0.0401647\\
0.14	-0.0404013\\
0.142	-0.0406291\\
0.144	-0.0408476\\
0.146	-0.0410566\\
0.148	-0.0412558\\
0.15	-0.041445\\
0.152	-0.0416239\\
0.154	-0.0417925\\
0.156	-0.0419506\\
0.158	-0.0420981\\
0.16	-0.042235\\
0.162	-0.0423614\\
0.164	-0.0424772\\
0.166	-0.0425828\\
0.168	-0.0426782\\
0.17	-0.0427636\\
0.172	-0.0428394\\
0.174	-0.0429059\\
0.176	-0.0429633\\
0.178	-0.0430122\\
0.18	-0.0430528\\
0.182	-0.0430856\\
0.184	-0.0431111\\
0.186	-0.0431297\\
0.188	-0.043142\\
0.19	-0.0431483\\
0.192	-0.0431492\\
0.194	-0.0431453\\
0.196	-0.0431368\\
0.198	-0.0431244\\
0.2	-0.0431085\\
0.202	-0.0430895\\
0.204	-0.0430678\\
0.206	-0.043044\\
0.208	-0.0430182\\
0.21	-0.042991\\
0.212	-0.0429626\\
0.214	-0.0429334\\
0.216	-0.0429036\\
0.218	-0.0428736\\
0.22	-0.0428436\\
0.222	-0.0428137\\
0.224	-0.0427843\\
0.226	-0.0427555\\
0.228	-0.0427274\\
0.23	-0.0427002\\
0.232	-0.0426741\\
0.234	-0.042649\\
0.236	-0.0426252\\
0.238	-0.0426026\\
0.24	-0.0425815\\
0.242	-0.0425617\\
0.244	-0.0425433\\
0.246	-0.0425265\\
0.248	-0.0425111\\
0.25	-0.0424973\\
0.252	-0.0424849\\
0.254	-0.0424741\\
0.256	-0.0424648\\
0.258	-0.042457\\
0.26	-0.0424507\\
0.262	-0.0424458\\
0.264	-0.0424424\\
0.266	-0.0424404\\
0.268	-0.0424398\\
0.27	-0.0424405\\
0.272	-0.0424425\\
0.274	-0.0424458\\
0.276	-0.0424503\\
0.278	-0.0424561\\
0.28	-0.0424629\\
0.282	-0.0424709\\
0.284	-0.0424799\\
0.286	-0.0424899\\
0.288	-0.0425009\\
0.29	-0.0425127\\
0.292	-0.0425254\\
0.294	-0.0425389\\
0.296	-0.0425531\\
0.298	-0.042568\\
0.3	-0.0425835\\
0.302	-0.0425995\\
0.304	-0.0426161\\
0.306	-0.0426331\\
0.308	-0.0426505\\
0.31	-0.0426683\\
0.312	-0.0426863\\
0.314	-0.0427046\\
0.316	-0.0427231\\
0.318	-0.0427417\\
0.32	-0.0427604\\
0.322	-0.0427791\\
0.324	-0.0427978\\
0.326	-0.0428165\\
0.328	-0.0428351\\
0.33	-0.0428536\\
0.332	-0.0428719\\
0.334	-0.04289\\
0.336	-0.042908\\
0.338	-0.0429257\\
0.34	-0.0429431\\
0.342	-0.0429603\\
0.344	-0.0429771\\
0.346	-0.0429937\\
0.348	-0.0430099\\
0.35	-0.0430258\\
0.352	-0.0430414\\
0.354	-0.0430566\\
0.356	-0.0430714\\
0.358	-0.0430859\\
0.36	-0.0431001\\
0.362	-0.0431139\\
0.364	-0.0431273\\
0.366	-0.0431404\\
0.368	-0.0431532\\
0.37	-0.0431656\\
0.372	-0.0431777\\
0.374	-0.0431895\\
0.376	-0.043201\\
0.378	-0.0432122\\
0.38	-0.0432231\\
0.382	-0.0432337\\
0.384	-0.0432441\\
0.386	-0.0432542\\
0.388	-0.043264\\
0.39	-0.0432736\\
0.392	-0.0432831\\
0.394	-0.0432922\\
0.396	-0.0433012\\
0.398	-0.04331\\
0.4	-0.0433186\\
0.402	-0.0433271\\
0.404	-0.0433353\\
0.406	-0.0433434\\
0.408	-0.0433514\\
0.41	-0.0433592\\
0.412	-0.0433668\\
0.414	-0.0433743\\
0.416	-0.0433817\\
0.418	-0.043389\\
0.42	-0.0433961\\
0.422	-0.0434032\\
0.424	-0.0434101\\
0.426	-0.0434169\\
0.428	-0.0434236\\
0.43	-0.0434301\\
0.432	-0.0434366\\
0.434	-0.043443\\
0.436	-0.0434492\\
0.438	-0.0434554\\
0.44	-0.0434615\\
0.442	-0.0434674\\
0.444	-0.0434733\\
0.446	-0.0434791\\
0.448	-0.0434847\\
0.45	-0.0434903\\
0.452	-0.0434957\\
0.454	-0.0435011\\
0.456	-0.0435064\\
0.458	-0.0435115\\
0.46	-0.0435166\\
0.462	-0.0435216\\
0.464	-0.0435264\\
0.466	-0.0435312\\
0.468	-0.0435359\\
0.47	-0.0435405\\
0.472	-0.0435449\\
0.474	-0.0435493\\
0.476	-0.0435536\\
0.478	-0.0435578\\
0.48	-0.0435618\\
0.482	-0.0435658\\
0.484	-0.0435697\\
0.486	-0.0435735\\
0.488	-0.0435772\\
0.49	-0.0435808\\
0.492	-0.0435843\\
0.494	-0.0435877\\
0.496	-0.043591\\
0.498	-0.0435943\\
0.5	-0.0435974\\
};
\addplot [color=mycolor8, line width=1.0pt, forget plot]
  table[row sep=crcr]{%
0	-0.017507\\
0.002	-0.017507\\
0.004	-0.017507\\
0.006	-0.017507\\
0.008	-0.017507\\
0.01	-0.017507\\
0.012	-0.017507\\
0.014	-0.017507\\
0.016	-0.017507\\
0.018	-0.017507\\
0.02	-0.017507\\
0.022	-0.017507\\
0.024	-0.017507\\
0.026	-0.017507\\
0.028	-0.017507\\
0.03	-0.017507\\
0.032	-0.0177587\\
0.034	-0.0183807\\
0.036	-0.0189828\\
0.038	-0.0195684\\
0.04	-0.0201402\\
0.042	-0.0207009\\
0.044	-0.0212522\\
0.046	-0.021796\\
0.048	-0.0223333\\
0.05	-0.0228651\\
0.052	-0.0233922\\
0.054	-0.0239148\\
0.056	-0.0244331\\
0.058	-0.0249473\\
0.06	-0.025457\\
0.062	-0.0259621\\
0.064	-0.0264623\\
0.066	-0.0269572\\
0.068	-0.0274462\\
0.07	-0.0279291\\
0.072	-0.0284054\\
0.074	-0.0288746\\
0.076	-0.0293366\\
0.078	-0.0297909\\
0.08	-0.0302373\\
0.082	-0.0306756\\
0.084	-0.0311058\\
0.086	-0.0315278\\
0.088	-0.0319417\\
0.09	-0.0323474\\
0.092	-0.0327452\\
0.094	-0.0331352\\
0.096	-0.0335176\\
0.098	-0.0338925\\
0.1	-0.0342604\\
0.102	-0.0346214\\
0.104	-0.0349757\\
0.106	-0.0353236\\
0.108	-0.0356652\\
0.11	-0.0360008\\
0.112	-0.0363304\\
0.114	-0.0366543\\
0.116	-0.0369723\\
0.118	-0.0372845\\
0.12	-0.0375908\\
0.122	-0.0378913\\
0.124	-0.0381856\\
0.126	-0.0384736\\
0.128	-0.0387552\\
0.13	-0.03903\\
0.132	-0.0392978\\
0.134	-0.0395582\\
0.136	-0.0398109\\
0.138	-0.0400556\\
0.14	-0.040292\\
0.142	-0.0405197\\
0.144	-0.0407385\\
0.146	-0.0409479\\
0.148	-0.0411478\\
0.15	-0.0413379\\
0.152	-0.0415181\\
0.154	-0.0416881\\
0.156	-0.0418478\\
0.158	-0.0419972\\
0.16	-0.0421363\\
0.162	-0.042265\\
0.164	-0.0423835\\
0.166	-0.0424919\\
0.168	-0.0425903\\
0.17	-0.042679\\
0.172	-0.0427582\\
0.174	-0.0428283\\
0.176	-0.0428895\\
0.178	-0.0429423\\
0.18	-0.042987\\
0.182	-0.0430241\\
0.184	-0.043054\\
0.186	-0.0430772\\
0.188	-0.0430941\\
0.19	-0.0431052\\
0.192	-0.043111\\
0.194	-0.0431119\\
0.196	-0.0431085\\
0.198	-0.0431011\\
0.2	-0.0430902\\
0.202	-0.0430763\\
0.204	-0.0430598\\
0.206	-0.043041\\
0.208	-0.0430203\\
0.21	-0.0429981\\
0.212	-0.0429748\\
0.214	-0.0429505\\
0.216	-0.0429257\\
0.218	-0.0429005\\
0.22	-0.0428752\\
0.222	-0.0428501\\
0.224	-0.0428252\\
0.226	-0.0428008\\
0.228	-0.042777\\
0.23	-0.042754\\
0.232	-0.0427319\\
0.234	-0.0427108\\
0.236	-0.0426907\\
0.238	-0.0426717\\
0.24	-0.042654\\
0.242	-0.0426374\\
0.244	-0.0426222\\
0.246	-0.0426082\\
0.248	-0.0425956\\
0.25	-0.0425843\\
0.252	-0.0425743\\
0.254	-0.0425656\\
0.256	-0.0425583\\
0.258	-0.0425523\\
0.26	-0.0425476\\
0.262	-0.0425441\\
0.264	-0.0425419\\
0.266	-0.042541\\
0.268	-0.0425412\\
0.27	-0.0425426\\
0.272	-0.0425451\\
0.274	-0.0425488\\
0.276	-0.0425535\\
0.278	-0.0425592\\
0.28	-0.0425659\\
0.282	-0.0425736\\
0.284	-0.0425822\\
0.286	-0.0425916\\
0.288	-0.0426018\\
0.29	-0.0426128\\
0.292	-0.0426246\\
0.294	-0.0426369\\
0.296	-0.04265\\
0.298	-0.0426635\\
0.3	-0.0426776\\
0.302	-0.0426922\\
0.304	-0.0427072\\
0.306	-0.0427226\\
0.308	-0.0427383\\
0.31	-0.0427543\\
0.312	-0.0427705\\
0.314	-0.0427869\\
0.316	-0.0428034\\
0.318	-0.04282\\
0.32	-0.0428367\\
0.322	-0.0428534\\
0.324	-0.0428701\\
0.326	-0.0428867\\
0.328	-0.0429033\\
0.33	-0.0429197\\
0.332	-0.042936\\
0.334	-0.0429521\\
0.336	-0.042968\\
0.338	-0.0429836\\
0.34	-0.0429991\\
0.342	-0.0430143\\
0.344	-0.0430292\\
0.346	-0.0430438\\
0.348	-0.0430582\\
0.35	-0.0430722\\
0.352	-0.043086\\
0.354	-0.0430994\\
0.356	-0.0431125\\
0.358	-0.0431253\\
0.36	-0.0431379\\
0.362	-0.0431501\\
0.364	-0.043162\\
0.366	-0.0431736\\
0.368	-0.0431849\\
0.37	-0.043196\\
0.372	-0.0432067\\
0.374	-0.0432172\\
0.376	-0.0432275\\
0.378	-0.0432375\\
0.38	-0.0432473\\
0.382	-0.0432568\\
0.384	-0.0432662\\
0.386	-0.0432753\\
0.388	-0.0432842\\
0.39	-0.0432929\\
0.392	-0.0433015\\
0.394	-0.0433099\\
0.396	-0.0433181\\
0.398	-0.0433262\\
0.4	-0.0433341\\
0.402	-0.0433419\\
0.404	-0.0433496\\
0.406	-0.0433571\\
0.408	-0.0433645\\
0.41	-0.0433718\\
0.412	-0.043379\\
0.414	-0.043386\\
0.416	-0.043393\\
0.418	-0.0433999\\
0.42	-0.0434066\\
0.422	-0.0434133\\
0.424	-0.0434198\\
0.426	-0.0434263\\
0.428	-0.0434327\\
0.43	-0.043439\\
0.432	-0.0434451\\
0.434	-0.0434512\\
0.436	-0.0434572\\
0.438	-0.0434631\\
0.44	-0.0434689\\
0.442	-0.0434747\\
0.444	-0.0434803\\
0.446	-0.0434858\\
0.448	-0.0434912\\
0.45	-0.0434965\\
0.452	-0.0435018\\
0.454	-0.0435069\\
0.456	-0.0435119\\
0.458	-0.0435169\\
0.46	-0.0435217\\
0.462	-0.0435264\\
0.464	-0.0435311\\
0.466	-0.0435356\\
0.468	-0.04354\\
0.47	-0.0435443\\
0.472	-0.0435486\\
0.474	-0.0435527\\
0.476	-0.0435567\\
0.478	-0.0435606\\
0.48	-0.0435645\\
0.482	-0.0435682\\
0.484	-0.0435718\\
0.486	-0.0435753\\
0.488	-0.0435788\\
0.49	-0.0435821\\
0.492	-0.0435853\\
0.494	-0.0435885\\
0.496	-0.0435915\\
0.498	-0.0435945\\
0.5	-0.0435973\\
};
\addplot [color=mycolor9, line width=1.0pt, forget plot]
  table[row sep=crcr]{%
0	-0.017507\\
0.002	-0.017507\\
0.004	-0.017507\\
0.006	-0.017507\\
0.008	-0.017507\\
0.01	-0.017507\\
0.012	-0.017507\\
0.014	-0.017507\\
0.016	-0.017507\\
0.018	-0.017507\\
0.02	-0.017507\\
0.022	-0.017507\\
0.024	-0.017507\\
0.026	-0.017507\\
0.028	-0.017507\\
0.03	-0.017507\\
0.032	-0.0178038\\
0.034	-0.018435\\
0.036	-0.0190475\\
0.038	-0.0196446\\
0.04	-0.0202293\\
0.042	-0.0208039\\
0.044	-0.0213704\\
0.046	-0.0219303\\
0.048	-0.022485\\
0.05	-0.0230353\\
0.052	-0.0235816\\
0.054	-0.0241243\\
0.056	-0.0246635\\
0.058	-0.025199\\
0.06	-0.0257306\\
0.062	-0.0262577\\
0.064	-0.02678\\
0.066	-0.0272967\\
0.068	-0.0278074\\
0.07	-0.0283113\\
0.072	-0.0288078\\
0.074	-0.0292964\\
0.076	-0.0297765\\
0.078	-0.0302475\\
0.08	-0.0307092\\
0.082	-0.0311611\\
0.084	-0.0316029\\
0.086	-0.0320345\\
0.088	-0.0324558\\
0.09	-0.0328667\\
0.092	-0.0332672\\
0.094	-0.0336574\\
0.096	-0.0340376\\
0.098	-0.0344079\\
0.1	-0.0347685\\
0.102	-0.0351197\\
0.104	-0.0354618\\
0.106	-0.035795\\
0.108	-0.0361198\\
0.11	-0.0364362\\
0.112	-0.0367447\\
0.114	-0.0370455\\
0.116	-0.0373387\\
0.118	-0.0376245\\
0.12	-0.0379031\\
0.122	-0.0381746\\
0.124	-0.038439\\
0.126	-0.0386964\\
0.128	-0.0389468\\
0.13	-0.0391901\\
0.132	-0.0394263\\
0.134	-0.0396553\\
0.136	-0.0398768\\
0.138	-0.0400909\\
0.14	-0.0402974\\
0.142	-0.040496\\
0.144	-0.0406867\\
0.146	-0.0408694\\
0.148	-0.0410438\\
0.15	-0.0412099\\
0.152	-0.0413676\\
0.154	-0.0415168\\
0.156	-0.0416574\\
0.158	-0.0417895\\
0.16	-0.041913\\
0.162	-0.0420281\\
0.164	-0.0421347\\
0.166	-0.042233\\
0.168	-0.0423232\\
0.17	-0.0424055\\
0.172	-0.04248\\
0.174	-0.042547\\
0.176	-0.0426068\\
0.178	-0.0426596\\
0.18	-0.0427059\\
0.182	-0.0427458\\
0.184	-0.0427798\\
0.186	-0.0428083\\
0.188	-0.0428315\\
0.19	-0.0428499\\
0.192	-0.0428638\\
0.194	-0.0428736\\
0.196	-0.0428797\\
0.198	-0.0428825\\
0.2	-0.0428822\\
0.202	-0.0428793\\
0.204	-0.0428739\\
0.206	-0.0428666\\
0.208	-0.0428575\\
0.21	-0.0428469\\
0.212	-0.0428352\\
0.214	-0.0428224\\
0.216	-0.0428089\\
0.218	-0.0427949\\
0.22	-0.0427805\\
0.222	-0.042766\\
0.224	-0.0427514\\
0.226	-0.042737\\
0.228	-0.0427227\\
0.23	-0.0427088\\
0.232	-0.0426954\\
0.234	-0.0426824\\
0.236	-0.0426701\\
0.238	-0.0426584\\
0.24	-0.0426473\\
0.242	-0.042637\\
0.244	-0.0426275\\
0.246	-0.0426187\\
0.248	-0.0426108\\
0.25	-0.0426036\\
0.252	-0.0425973\\
0.254	-0.0425918\\
0.256	-0.0425872\\
0.258	-0.0425833\\
0.26	-0.0425803\\
0.262	-0.0425781\\
0.264	-0.0425767\\
0.266	-0.0425761\\
0.268	-0.0425762\\
0.27	-0.0425771\\
0.272	-0.0425788\\
0.274	-0.0425812\\
0.276	-0.0425843\\
0.278	-0.042588\\
0.28	-0.0425925\\
0.282	-0.0425976\\
0.284	-0.0426033\\
0.286	-0.0426096\\
0.288	-0.0426165\\
0.29	-0.0426239\\
0.292	-0.0426319\\
0.294	-0.0426404\\
0.296	-0.0426494\\
0.298	-0.0426588\\
0.3	-0.0426686\\
0.302	-0.0426789\\
0.304	-0.0426895\\
0.306	-0.0427005\\
0.308	-0.0427117\\
0.31	-0.0427233\\
0.312	-0.0427352\\
0.314	-0.0427472\\
0.316	-0.0427595\\
0.318	-0.042772\\
0.32	-0.0427847\\
0.322	-0.0427975\\
0.324	-0.0428104\\
0.326	-0.0428234\\
0.328	-0.0428365\\
0.33	-0.0428496\\
0.332	-0.0428627\\
0.334	-0.0428759\\
0.336	-0.0428891\\
0.338	-0.0429023\\
0.34	-0.0429154\\
0.342	-0.0429285\\
0.344	-0.0429415\\
0.346	-0.0429544\\
0.348	-0.0429673\\
0.35	-0.0429801\\
0.352	-0.0429927\\
0.354	-0.0430053\\
0.356	-0.0430177\\
0.358	-0.0430301\\
0.36	-0.0430423\\
0.362	-0.0430543\\
0.364	-0.0430663\\
0.366	-0.0430781\\
0.368	-0.0430897\\
0.37	-0.0431012\\
0.372	-0.0431126\\
0.374	-0.0431238\\
0.376	-0.0431349\\
0.378	-0.0431458\\
0.38	-0.0431566\\
0.382	-0.0431673\\
0.384	-0.0431778\\
0.386	-0.0431882\\
0.388	-0.0431984\\
0.39	-0.0432085\\
0.392	-0.0432184\\
0.394	-0.0432282\\
0.396	-0.0432379\\
0.398	-0.0432474\\
0.4	-0.0432568\\
0.402	-0.0432661\\
0.404	-0.0432752\\
0.406	-0.0432842\\
0.408	-0.043293\\
0.41	-0.0433018\\
0.412	-0.0433103\\
0.414	-0.0433188\\
0.416	-0.0433271\\
0.418	-0.0433353\\
0.42	-0.0433434\\
0.422	-0.0433513\\
0.424	-0.0433591\\
0.426	-0.0433668\\
0.428	-0.0433743\\
0.43	-0.0433817\\
0.432	-0.0433889\\
0.434	-0.0433961\\
0.436	-0.0434031\\
0.438	-0.0434099\\
0.44	-0.0434167\\
0.442	-0.0434233\\
0.444	-0.0434297\\
0.446	-0.0434361\\
0.448	-0.0434423\\
0.45	-0.0434484\\
0.452	-0.0434543\\
0.454	-0.0434601\\
0.456	-0.0434658\\
0.458	-0.0434713\\
0.46	-0.0434767\\
0.462	-0.043482\\
0.464	-0.0434872\\
0.466	-0.0434922\\
0.468	-0.0434971\\
0.47	-0.0435019\\
0.472	-0.0435066\\
0.474	-0.0435111\\
0.476	-0.0435156\\
0.478	-0.0435199\\
0.48	-0.0435241\\
0.482	-0.0435282\\
0.484	-0.0435321\\
0.486	-0.043536\\
0.488	-0.0435398\\
0.49	-0.0435434\\
0.492	-0.043547\\
0.494	-0.0435504\\
0.496	-0.0435538\\
0.498	-0.043557\\
0.5	-0.0435602\\
};
\addplot [color=mycolor10, line width=1.0pt, forget plot]
  table[row sep=crcr]{%
0	-0.017507\\
0.002	-0.017507\\
0.004	-0.017507\\
0.006	-0.017507\\
0.008	-0.017507\\
0.01	-0.017507\\
0.012	-0.017507\\
0.014	-0.017507\\
0.016	-0.017507\\
0.018	-0.017507\\
0.02	-0.017507\\
0.022	-0.017507\\
0.024	-0.017507\\
0.026	-0.017507\\
0.028	-0.017507\\
0.03	-0.017507\\
0.032	-0.0177992\\
0.034	-0.0184296\\
0.036	-0.0190411\\
0.038	-0.0196371\\
0.04	-0.0202206\\
0.042	-0.020794\\
0.044	-0.0213592\\
0.046	-0.0219178\\
0.048	-0.022471\\
0.05	-0.0230196\\
0.052	-0.0235644\\
0.054	-0.0241054\\
0.056	-0.0246429\\
0.058	-0.0251767\\
0.06	-0.0257065\\
0.062	-0.0262318\\
0.064	-0.0267523\\
0.066	-0.0272672\\
0.068	-0.0277761\\
0.07	-0.0282783\\
0.072	-0.0287731\\
0.074	-0.02926\\
0.076	-0.0297385\\
0.078	-0.030208\\
0.08	-0.0306682\\
0.082	-0.0311186\\
0.084	-0.0315592\\
0.086	-0.0319896\\
0.088	-0.0324097\\
0.09	-0.0328196\\
0.092	-0.0332192\\
0.094	-0.0336086\\
0.096	-0.0339881\\
0.098	-0.0343579\\
0.1	-0.034718\\
0.102	-0.0350689\\
0.104	-0.0354108\\
0.106	-0.035744\\
0.108	-0.0360688\\
0.11	-0.0363855\\
0.112	-0.0366943\\
0.114	-0.0369955\\
0.116	-0.0372892\\
0.118	-0.0375757\\
0.12	-0.0378551\\
0.122	-0.0381274\\
0.124	-0.0383928\\
0.126	-0.0386513\\
0.128	-0.0389029\\
0.13	-0.0391474\\
0.132	-0.0393849\\
0.134	-0.0396152\\
0.136	-0.0398382\\
0.138	-0.0400538\\
0.14	-0.0402618\\
0.142	-0.040462\\
0.144	-0.0406544\\
0.146	-0.0408387\\
0.148	-0.0410148\\
0.15	-0.0411825\\
0.152	-0.0413419\\
0.154	-0.0414927\\
0.156	-0.0416351\\
0.158	-0.0417688\\
0.16	-0.041894\\
0.162	-0.0420107\\
0.164	-0.042119\\
0.166	-0.0422189\\
0.168	-0.0423107\\
0.17	-0.0423945\\
0.172	-0.0424706\\
0.174	-0.0425391\\
0.176	-0.0426004\\
0.178	-0.0426547\\
0.18	-0.0427024\\
0.182	-0.0427438\\
0.184	-0.0427792\\
0.186	-0.042809\\
0.188	-0.0428336\\
0.19	-0.0428534\\
0.192	-0.0428686\\
0.194	-0.0428798\\
0.196	-0.0428872\\
0.198	-0.0428912\\
0.2	-0.0428922\\
0.202	-0.0428905\\
0.204	-0.0428865\\
0.206	-0.0428804\\
0.208	-0.0428725\\
0.21	-0.0428631\\
0.212	-0.0428525\\
0.214	-0.042841\\
0.216	-0.0428286\\
0.218	-0.0428157\\
0.22	-0.0428024\\
0.222	-0.0427889\\
0.224	-0.0427754\\
0.226	-0.0427619\\
0.228	-0.0427487\\
0.23	-0.0427357\\
0.232	-0.0427231\\
0.234	-0.042711\\
0.236	-0.0426995\\
0.238	-0.0426885\\
0.24	-0.0426782\\
0.242	-0.0426685\\
0.244	-0.0426596\\
0.246	-0.0426514\\
0.248	-0.042644\\
0.25	-0.0426373\\
0.252	-0.0426313\\
0.254	-0.0426262\\
0.256	-0.0426219\\
0.258	-0.0426183\\
0.26	-0.0426155\\
0.262	-0.0426134\\
0.264	-0.0426121\\
0.266	-0.0426116\\
0.268	-0.0426118\\
0.27	-0.0426127\\
0.272	-0.0426143\\
0.274	-0.0426166\\
0.276	-0.0426196\\
0.278	-0.0426233\\
0.28	-0.0426276\\
0.282	-0.0426325\\
0.284	-0.042638\\
0.286	-0.0426442\\
0.288	-0.0426508\\
0.29	-0.0426581\\
0.292	-0.0426658\\
0.294	-0.042674\\
0.296	-0.0426827\\
0.298	-0.0426919\\
0.3	-0.0427014\\
0.302	-0.0427114\\
0.304	-0.0427217\\
0.306	-0.0427324\\
0.308	-0.0427434\\
0.31	-0.0427547\\
0.312	-0.0427662\\
0.314	-0.042778\\
0.316	-0.04279\\
0.318	-0.0428022\\
0.32	-0.0428145\\
0.322	-0.042827\\
0.324	-0.0428395\\
0.326	-0.0428522\\
0.328	-0.0428649\\
0.33	-0.0428777\\
0.332	-0.0428905\\
0.334	-0.0429034\\
0.336	-0.0429162\\
0.338	-0.0429289\\
0.34	-0.0429417\\
0.342	-0.0429544\\
0.344	-0.042967\\
0.346	-0.0429795\\
0.348	-0.042992\\
0.35	-0.0430043\\
0.352	-0.0430166\\
0.354	-0.0430287\\
0.356	-0.0430407\\
0.358	-0.0430526\\
0.36	-0.0430643\\
0.362	-0.0430759\\
0.364	-0.0430874\\
0.366	-0.0430987\\
0.368	-0.0431099\\
0.37	-0.0431209\\
0.372	-0.0431318\\
0.374	-0.0431425\\
0.376	-0.0431531\\
0.378	-0.0431636\\
0.38	-0.0431739\\
0.382	-0.043184\\
0.384	-0.043194\\
0.386	-0.0432039\\
0.388	-0.0432137\\
0.39	-0.0432233\\
0.392	-0.0432327\\
0.394	-0.043242\\
0.396	-0.0432512\\
0.398	-0.0432603\\
0.4	-0.0432692\\
0.402	-0.0432779\\
0.404	-0.0432866\\
0.406	-0.0432951\\
0.408	-0.0433035\\
0.41	-0.0433118\\
0.412	-0.0433199\\
0.414	-0.0433279\\
0.416	-0.0433357\\
0.418	-0.0433435\\
0.42	-0.0433511\\
0.422	-0.0433586\\
0.424	-0.0433659\\
0.426	-0.0433732\\
0.428	-0.0433803\\
0.43	-0.0433872\\
0.432	-0.0433941\\
0.434	-0.0434008\\
0.436	-0.0434074\\
0.438	-0.0434138\\
0.44	-0.0434201\\
0.442	-0.0434263\\
0.444	-0.0434324\\
0.446	-0.0434384\\
0.448	-0.0434442\\
0.45	-0.0434499\\
0.452	-0.0434555\\
0.454	-0.0434609\\
0.456	-0.0434663\\
0.458	-0.0434715\\
0.46	-0.0434766\\
0.462	-0.0434815\\
0.464	-0.0434864\\
0.466	-0.0434911\\
0.468	-0.0434958\\
0.47	-0.0435003\\
0.472	-0.0435047\\
0.474	-0.043509\\
0.476	-0.0435131\\
0.478	-0.0435172\\
0.48	-0.0435212\\
0.482	-0.0435251\\
0.484	-0.0435288\\
0.486	-0.0435325\\
0.488	-0.0435361\\
0.49	-0.0435396\\
0.492	-0.043543\\
0.494	-0.0435463\\
0.496	-0.0435495\\
0.498	-0.0435527\\
0.5	-0.0435557\\
};
\addplot [color=mycolor11, line width=1.0pt, forget plot]
  table[row sep=crcr]{%
0	-0.017507\\
0.002	-0.017507\\
0.004	-0.017507\\
0.006	-0.017507\\
0.008	-0.017507\\
0.01	-0.017507\\
0.012	-0.017507\\
0.014	-0.017507\\
0.016	-0.017507\\
0.018	-0.017507\\
0.02	-0.017507\\
0.022	-0.017507\\
0.024	-0.017507\\
0.026	-0.017507\\
0.028	-0.017507\\
0.03	-0.017507\\
0.032	-0.0180743\\
0.034	-0.0187326\\
0.036	-0.0193709\\
0.038	-0.0199919\\
0.04	-0.0205983\\
0.042	-0.0211921\\
0.044	-0.021775\\
0.046	-0.0223484\\
0.048	-0.0229133\\
0.05	-0.0234705\\
0.052	-0.0240205\\
0.054	-0.0245636\\
0.056	-0.0251001\\
0.058	-0.0256298\\
0.06	-0.0261526\\
0.062	-0.0266685\\
0.064	-0.027177\\
0.066	-0.027678\\
0.068	-0.028171\\
0.07	-0.0286558\\
0.072	-0.0291321\\
0.074	-0.0295995\\
0.076	-0.0300579\\
0.078	-0.030507\\
0.08	-0.0309467\\
0.082	-0.0313768\\
0.084	-0.0317974\\
0.086	-0.0322083\\
0.088	-0.0326098\\
0.09	-0.0330018\\
0.092	-0.0333844\\
0.094	-0.0337579\\
0.096	-0.0341224\\
0.098	-0.0344782\\
0.1	-0.0348253\\
0.102	-0.0351642\\
0.104	-0.0354949\\
0.106	-0.0358176\\
0.108	-0.0361327\\
0.11	-0.0364402\\
0.112	-0.0367404\\
0.114	-0.0370333\\
0.116	-0.0373191\\
0.118	-0.0375978\\
0.12	-0.0378695\\
0.122	-0.0381342\\
0.124	-0.0383918\\
0.126	-0.0386425\\
0.128	-0.0388861\\
0.13	-0.0391225\\
0.132	-0.0393516\\
0.134	-0.0395735\\
0.136	-0.0397878\\
0.138	-0.0399946\\
0.14	-0.0401937\\
0.142	-0.040385\\
0.144	-0.0405685\\
0.146	-0.0407439\\
0.148	-0.0409114\\
0.15	-0.0410707\\
0.152	-0.0412219\\
0.154	-0.041365\\
0.156	-0.0415\\
0.158	-0.041627\\
0.16	-0.0417459\\
0.162	-0.0418569\\
0.164	-0.0419602\\
0.166	-0.0420559\\
0.168	-0.0421441\\
0.17	-0.0422251\\
0.172	-0.0422991\\
0.174	-0.0423663\\
0.176	-0.042427\\
0.178	-0.0424815\\
0.18	-0.04253\\
0.182	-0.0425729\\
0.184	-0.0426104\\
0.186	-0.042643\\
0.188	-0.0426708\\
0.19	-0.0426943\\
0.192	-0.0427137\\
0.194	-0.0427294\\
0.196	-0.0427417\\
0.198	-0.0427508\\
0.2	-0.0427571\\
0.202	-0.0427609\\
0.204	-0.0427624\\
0.206	-0.0427618\\
0.208	-0.0427596\\
0.21	-0.0427558\\
0.212	-0.0427506\\
0.214	-0.0427444\\
0.216	-0.0427373\\
0.218	-0.0427295\\
0.22	-0.0427211\\
0.222	-0.0427123\\
0.224	-0.0427033\\
0.226	-0.0426941\\
0.228	-0.0426849\\
0.23	-0.0426757\\
0.232	-0.0426667\\
0.234	-0.042658\\
0.236	-0.0426496\\
0.238	-0.0426416\\
0.24	-0.042634\\
0.242	-0.0426269\\
0.244	-0.0426203\\
0.246	-0.0426143\\
0.248	-0.0426089\\
0.25	-0.0426041\\
0.252	-0.0425999\\
0.254	-0.0425964\\
0.256	-0.0425935\\
0.258	-0.0425912\\
0.26	-0.0425896\\
0.262	-0.0425887\\
0.264	-0.0425885\\
0.266	-0.0425888\\
0.268	-0.0425899\\
0.27	-0.0425915\\
0.272	-0.0425938\\
0.274	-0.0425967\\
0.276	-0.0426002\\
0.278	-0.0426043\\
0.28	-0.0426089\\
0.282	-0.0426141\\
0.284	-0.0426198\\
0.286	-0.042626\\
0.288	-0.0426327\\
0.29	-0.0426399\\
0.292	-0.0426475\\
0.294	-0.0426555\\
0.296	-0.042664\\
0.298	-0.0426728\\
0.3	-0.042682\\
0.302	-0.0426915\\
0.304	-0.0427013\\
0.306	-0.0427114\\
0.308	-0.0427218\\
0.31	-0.0427324\\
0.312	-0.0427432\\
0.314	-0.0427543\\
0.316	-0.0427655\\
0.318	-0.0427769\\
0.32	-0.0427884\\
0.322	-0.0428\\
0.324	-0.0428118\\
0.326	-0.0428236\\
0.328	-0.0428355\\
0.33	-0.0428474\\
0.332	-0.0428594\\
0.334	-0.0428714\\
0.336	-0.0428834\\
0.338	-0.0428954\\
0.34	-0.0429074\\
0.342	-0.0429193\\
0.344	-0.0429313\\
0.346	-0.0429431\\
0.348	-0.0429549\\
0.35	-0.0429667\\
0.352	-0.0429783\\
0.354	-0.0429899\\
0.356	-0.0430014\\
0.358	-0.0430129\\
0.36	-0.0430242\\
0.362	-0.0430354\\
0.364	-0.0430465\\
0.366	-0.0430575\\
0.368	-0.0430685\\
0.37	-0.0430793\\
0.372	-0.04309\\
0.374	-0.0431005\\
0.376	-0.043111\\
0.378	-0.0431214\\
0.38	-0.0431316\\
0.382	-0.0431417\\
0.384	-0.0431517\\
0.386	-0.0431616\\
0.388	-0.0431713\\
0.39	-0.043181\\
0.392	-0.0431905\\
0.394	-0.0431999\\
0.396	-0.0432092\\
0.398	-0.0432184\\
0.4	-0.0432274\\
0.402	-0.0432363\\
0.404	-0.0432451\\
0.406	-0.0432538\\
0.408	-0.0432623\\
0.41	-0.0432708\\
0.412	-0.0432791\\
0.414	-0.0432873\\
0.416	-0.0432954\\
0.418	-0.0433033\\
0.42	-0.0433111\\
0.422	-0.0433188\\
0.424	-0.0433264\\
0.426	-0.0433338\\
0.428	-0.0433412\\
0.43	-0.0433484\\
0.432	-0.0433554\\
0.434	-0.0433624\\
0.436	-0.0433692\\
0.438	-0.0433759\\
0.44	-0.0433825\\
0.442	-0.0433889\\
0.444	-0.0433952\\
0.446	-0.0434014\\
0.448	-0.0434075\\
0.45	-0.0434135\\
0.452	-0.0434193\\
0.454	-0.043425\\
0.456	-0.0434306\\
0.458	-0.043436\\
0.46	-0.0434414\\
0.462	-0.0434466\\
0.464	-0.0434517\\
0.466	-0.0434567\\
0.468	-0.0434616\\
0.47	-0.0434664\\
0.472	-0.043471\\
0.474	-0.0434756\\
0.476	-0.04348\\
0.478	-0.0434844\\
0.48	-0.0434886\\
0.482	-0.0434927\\
0.484	-0.0434968\\
0.486	-0.0435007\\
0.488	-0.0435045\\
0.49	-0.0435083\\
0.492	-0.0435119\\
0.494	-0.0435155\\
0.496	-0.043519\\
0.498	-0.0435224\\
0.5	-0.0435257\\
};
\end{axis}

\begin{axis}[%
width={\figscale*0.642in},
height={\figscale*0.17in},
at={(\figscale*3.666in,\figscale*1.895in)},
scale only axis,
separate axis lines,
every outer x axis line/.append style={black},
every x tick label/.append style={font=\figticfont, yshift=\figinsettickshift},
every x tick/.append style={black},
xmin=0,
xmax=0.5,
xtick={  0, 0.4},
every outer y axis line/.append style={black},
every y tick label/.append style={font=\figticfont},
every y tick/.append style={black},
ymin=-0.0206901,
ymax=-0.017507,
ytick={ -0.02, -0.018},
axis background/.style={fill=white},
ylabel style={rotate=-90, at={(axis description cs:-0.145,0.5)}, anchor=east}
]
\addplot [color=mycolor1, line width=1.4pt, forget plot]
  table[row sep=crcr]{%
0	0\\
0.002	-0.00183136\\
0.004	-0.00337796\\
0.006	-0.00470392\\
0.008	-0.00585756\\
0.01	-0.00687514\\
0.012	-0.00778386\\
0.014	-0.00860411\\
0.016	-0.00935119\\
0.018	-0.0100367\\
0.02	-0.0106695\\
0.022	-0.0112563\\
0.024	-0.0118026\\
0.026	-0.0123126\\
0.028	-0.0127899\\
0.03	-0.0132372\\
0.032	-0.013657\\
0.034	-0.0140516\\
0.036	-0.0144226\\
0.038	-0.0147718\\
0.04	-0.0151007\\
0.042	-0.0154107\\
0.044	-0.0157029\\
0.046	-0.0159786\\
0.048	-0.0162389\\
0.05	-0.0164849\\
0.052	-0.0167174\\
0.054	-0.0169375\\
0.056	-0.017146\\
0.058	-0.0173437\\
0.06	-0.0175314\\
0.062	-0.0177098\\
0.064	-0.0178796\\
0.066	-0.0180412\\
0.068	-0.0181953\\
0.07	-0.0183424\\
0.072	-0.0184828\\
0.074	-0.0186169\\
0.076	-0.0187451\\
0.078	-0.0188675\\
0.08	-0.0189844\\
0.082	-0.0190959\\
0.084	-0.0192021\\
0.086	-0.0193032\\
0.088	-0.0193992\\
0.09	-0.0194901\\
0.092	-0.019576\\
0.094	-0.0196567\\
0.096	-0.0197323\\
0.098	-0.0198028\\
0.1	-0.0198681\\
0.102	-0.0199282\\
0.104	-0.0199831\\
0.106	-0.0200328\\
0.108	-0.0200772\\
0.11	-0.0201164\\
0.112	-0.0201504\\
0.114	-0.0201792\\
0.116	-0.0202028\\
0.118	-0.0202214\\
0.12	-0.0202349\\
0.122	-0.0202436\\
0.124	-0.0202474\\
0.126	-0.0202466\\
0.128	-0.0202412\\
0.13	-0.0202315\\
0.132	-0.0202175\\
0.134	-0.0201996\\
0.136	-0.0201777\\
0.138	-0.0201522\\
0.14	-0.0201233\\
0.142	-0.020091\\
0.144	-0.0200558\\
0.146	-0.0200177\\
0.148	-0.0199771\\
0.15	-0.019934\\
0.152	-0.0198888\\
0.154	-0.0198417\\
0.156	-0.0197929\\
0.158	-0.0197425\\
0.16	-0.0196909\\
0.162	-0.0196382\\
0.164	-0.0195846\\
0.166	-0.0195304\\
0.168	-0.0194757\\
0.17	-0.0194207\\
0.172	-0.0193656\\
0.174	-0.0193106\\
0.176	-0.0192559\\
0.178	-0.0192016\\
0.18	-0.0191479\\
0.182	-0.0190949\\
0.184	-0.0190428\\
0.186	-0.0189917\\
0.188	-0.0189418\\
0.19	-0.0188931\\
0.192	-0.0188458\\
0.194	-0.0188\\
0.196	-0.0187558\\
0.198	-0.0187132\\
0.2	-0.0186724\\
0.202	-0.0186335\\
0.204	-0.0185964\\
0.206	-0.0185613\\
0.208	-0.0185282\\
0.21	-0.0184972\\
0.212	-0.0184682\\
0.214	-0.0184413\\
0.216	-0.0184166\\
0.218	-0.018394\\
0.22	-0.0183735\\
0.222	-0.0183552\\
0.224	-0.018339\\
0.226	-0.0183249\\
0.228	-0.0183129\\
0.23	-0.018303\\
0.232	-0.0182951\\
0.234	-0.0182891\\
0.236	-0.0182851\\
0.238	-0.018283\\
0.24	-0.0182826\\
0.242	-0.0182841\\
0.244	-0.0182873\\
0.246	-0.0182921\\
0.248	-0.0182984\\
0.25	-0.0183063\\
0.252	-0.0183156\\
0.254	-0.0183262\\
0.256	-0.0183381\\
0.258	-0.0183512\\
0.26	-0.0183654\\
0.262	-0.0183806\\
0.264	-0.0183967\\
0.266	-0.0184137\\
0.268	-0.0184315\\
0.27	-0.01845\\
0.272	-0.0184691\\
0.274	-0.0184887\\
0.276	-0.0185088\\
0.278	-0.0185293\\
0.28	-0.01855\\
0.282	-0.018571\\
0.284	-0.0185922\\
0.286	-0.0186134\\
0.288	-0.0186346\\
0.29	-0.0186558\\
0.292	-0.0186769\\
0.294	-0.0186978\\
0.296	-0.0187184\\
0.298	-0.0187388\\
0.3	-0.0187588\\
0.302	-0.0187784\\
0.304	-0.0187975\\
0.306	-0.0188162\\
0.308	-0.0188344\\
0.31	-0.018852\\
0.312	-0.018869\\
0.314	-0.0188853\\
0.316	-0.018901\\
0.318	-0.0189161\\
0.32	-0.0189304\\
0.322	-0.018944\\
0.324	-0.0189569\\
0.326	-0.0189691\\
0.328	-0.0189806\\
0.33	-0.0189913\\
0.332	-0.0190012\\
0.334	-0.0190104\\
0.336	-0.0190189\\
0.338	-0.0190267\\
0.34	-0.0190338\\
0.342	-0.0190402\\
0.344	-0.0190459\\
0.346	-0.0190509\\
0.348	-0.0190553\\
0.35	-0.0190592\\
0.352	-0.0190624\\
0.354	-0.019065\\
0.356	-0.0190672\\
0.358	-0.0190688\\
0.36	-0.01907\\
0.362	-0.0190707\\
0.364	-0.019071\\
0.366	-0.0190709\\
0.368	-0.0190705\\
0.37	-0.0190697\\
0.372	-0.0190687\\
0.374	-0.0190674\\
0.376	-0.0190659\\
0.378	-0.0190641\\
0.38	-0.0190622\\
0.382	-0.0190602\\
0.384	-0.019058\\
0.386	-0.0190558\\
0.388	-0.0190534\\
0.39	-0.019051\\
0.392	-0.0190486\\
0.394	-0.0190462\\
0.396	-0.0190437\\
0.398	-0.0190413\\
0.4	-0.0190389\\
0.402	-0.0190366\\
0.404	-0.0190343\\
0.406	-0.0190321\\
0.408	-0.0190299\\
0.41	-0.0190279\\
0.412	-0.0190259\\
0.414	-0.019024\\
0.416	-0.0190222\\
0.418	-0.0190204\\
0.42	-0.0190188\\
0.422	-0.0190172\\
0.424	-0.0190157\\
0.426	-0.0190143\\
0.428	-0.019013\\
0.43	-0.0190117\\
0.432	-0.0190105\\
0.434	-0.0190093\\
0.436	-0.0190082\\
0.438	-0.0190071\\
0.44	-0.0190061\\
0.442	-0.019005\\
0.444	-0.019004\\
0.446	-0.0190029\\
0.448	-0.0190019\\
0.45	-0.0190008\\
0.452	-0.0189997\\
0.454	-0.0189985\\
0.456	-0.0189973\\
0.458	-0.018996\\
0.46	-0.0189947\\
0.462	-0.0189933\\
0.464	-0.0189918\\
0.466	-0.0189903\\
0.468	-0.0189886\\
0.47	-0.0189868\\
0.472	-0.018985\\
0.474	-0.018983\\
0.476	-0.0189809\\
0.478	-0.0189787\\
0.48	-0.0189763\\
0.482	-0.0189739\\
0.484	-0.0189713\\
0.486	-0.0189686\\
0.488	-0.0189658\\
0.49	-0.0189628\\
0.492	-0.0189598\\
0.494	-0.0189566\\
0.496	-0.0189533\\
0.498	-0.0189499\\
0.5	-0.0189464\\
};
\addplot [color=mycolor2, line width=1.0pt, forget plot]
  table[row sep=crcr]{%
0	0\\
0.002	-0.00183136\\
0.004	-0.00337796\\
0.006	-0.00470391\\
0.008	-0.00585754\\
0.01	-0.00687511\\
0.012	-0.00778382\\
0.014	-0.00860405\\
0.016	-0.00935112\\
0.018	-0.0100366\\
0.02	-0.0106694\\
0.022	-0.0112563\\
0.024	-0.0118026\\
0.026	-0.0123128\\
0.028	-0.0127901\\
0.03	-0.0132376\\
0.032	-0.0136577\\
0.034	-0.0140526\\
0.036	-0.014424\\
0.038	-0.0147737\\
0.04	-0.0151031\\
0.042	-0.0154138\\
0.044	-0.0157068\\
0.046	-0.0159835\\
0.048	-0.0162448\\
0.05	-0.016492\\
0.052	-0.0167259\\
0.054	-0.0169474\\
0.056	-0.0171576\\
0.058	-0.0173571\\
0.06	-0.0175467\\
0.062	-0.0177273\\
0.064	-0.0178993\\
0.066	-0.0180633\\
0.068	-0.0182201\\
0.07	-0.0183699\\
0.072	-0.0185132\\
0.074	-0.0186503\\
0.076	-0.0187816\\
0.078	-0.0189073\\
0.08	-0.0190276\\
0.082	-0.0191425\\
0.084	-0.0192524\\
0.086	-0.0193571\\
0.088	-0.0194568\\
0.09	-0.0195514\\
0.092	-0.019641\\
0.094	-0.0197254\\
0.096	-0.0198048\\
0.098	-0.0198789\\
0.1	-0.0199479\\
0.102	-0.0200115\\
0.104	-0.0200699\\
0.106	-0.0201228\\
0.108	-0.0201704\\
0.11	-0.0202126\\
0.112	-0.0202493\\
0.114	-0.0202807\\
0.116	-0.0203066\\
0.118	-0.0203273\\
0.12	-0.0203426\\
0.122	-0.0203527\\
0.124	-0.0203578\\
0.126	-0.0203578\\
0.128	-0.0203529\\
0.13	-0.0203433\\
0.132	-0.020329\\
0.134	-0.0203103\\
0.136	-0.0202874\\
0.138	-0.0202604\\
0.14	-0.0202295\\
0.142	-0.0201948\\
0.144	-0.0201567\\
0.146	-0.0201154\\
0.148	-0.020071\\
0.15	-0.0200237\\
0.152	-0.0199738\\
0.154	-0.0199216\\
0.156	-0.0198672\\
0.158	-0.0198108\\
0.16	-0.0197527\\
0.162	-0.0196932\\
0.164	-0.0196324\\
0.166	-0.0195705\\
0.168	-0.0195078\\
0.17	-0.0194444\\
0.172	-0.0193806\\
0.174	-0.0193166\\
0.176	-0.0192526\\
0.178	-0.0191887\\
0.18	-0.0191252\\
0.182	-0.0190623\\
0.184	-0.019\\
0.186	-0.0189386\\
0.188	-0.0188783\\
0.19	-0.0188191\\
0.192	-0.0187613\\
0.194	-0.018705\\
0.196	-0.0186503\\
0.198	-0.0185973\\
0.2	-0.0185462\\
0.202	-0.0184971\\
0.204	-0.01845\\
0.206	-0.0184051\\
0.208	-0.0183625\\
0.21	-0.0183222\\
0.212	-0.0182844\\
0.214	-0.018249\\
0.216	-0.0182161\\
0.218	-0.0181858\\
0.22	-0.0181581\\
0.222	-0.018133\\
0.224	-0.0181106\\
0.226	-0.0180908\\
0.228	-0.0180737\\
0.23	-0.0180593\\
0.232	-0.0180475\\
0.234	-0.0180383\\
0.236	-0.0180318\\
0.238	-0.0180278\\
0.24	-0.0180263\\
0.242	-0.0180273\\
0.244	-0.0180307\\
0.246	-0.0180366\\
0.248	-0.0180447\\
0.25	-0.018055\\
0.252	-0.0180675\\
0.254	-0.0180821\\
0.256	-0.0180987\\
0.258	-0.0181172\\
0.26	-0.0181375\\
0.262	-0.0181595\\
0.264	-0.0181832\\
0.266	-0.0182084\\
0.268	-0.0182351\\
0.27	-0.018263\\
0.272	-0.0182922\\
0.274	-0.0183225\\
0.276	-0.0183539\\
0.278	-0.0183861\\
0.28	-0.0184191\\
0.282	-0.0184528\\
0.284	-0.018487\\
0.286	-0.0185218\\
0.288	-0.0185568\\
0.29	-0.0185922\\
0.292	-0.0186277\\
0.294	-0.0186632\\
0.296	-0.0186987\\
0.298	-0.018734\\
0.3	-0.018769\\
0.302	-0.0188037\\
0.304	-0.018838\\
0.306	-0.0188717\\
0.308	-0.0189048\\
0.31	-0.0189372\\
0.312	-0.0189688\\
0.314	-0.0189996\\
0.316	-0.0190294\\
0.318	-0.0190583\\
0.32	-0.0190861\\
0.322	-0.0191128\\
0.324	-0.0191383\\
0.326	-0.0191626\\
0.328	-0.0191857\\
0.33	-0.0192074\\
0.332	-0.0192279\\
0.334	-0.019247\\
0.336	-0.0192647\\
0.338	-0.019281\\
0.34	-0.0192959\\
0.342	-0.0193094\\
0.344	-0.0193214\\
0.346	-0.019332\\
0.348	-0.0193412\\
0.35	-0.019349\\
0.352	-0.0193554\\
0.354	-0.0193604\\
0.356	-0.0193641\\
0.358	-0.0193664\\
0.36	-0.0193674\\
0.362	-0.0193671\\
0.364	-0.0193656\\
0.366	-0.0193628\\
0.368	-0.0193589\\
0.37	-0.0193538\\
0.372	-0.0193477\\
0.374	-0.0193405\\
0.376	-0.0193323\\
0.378	-0.0193232\\
0.38	-0.0193132\\
0.382	-0.0193023\\
0.384	-0.0192907\\
0.386	-0.0192783\\
0.388	-0.0192652\\
0.39	-0.0192515\\
0.392	-0.0192372\\
0.394	-0.0192224\\
0.396	-0.0192071\\
0.398	-0.0191915\\
0.4	-0.0191754\\
0.402	-0.0191591\\
0.404	-0.0191425\\
0.406	-0.0191257\\
0.408	-0.0191087\\
0.41	-0.0190917\\
0.412	-0.0190746\\
0.414	-0.0190575\\
0.416	-0.0190404\\
0.418	-0.0190234\\
0.42	-0.0190066\\
0.422	-0.0189899\\
0.424	-0.0189734\\
0.426	-0.0189572\\
0.428	-0.0189413\\
0.43	-0.0189257\\
0.432	-0.0189104\\
0.434	-0.0188956\\
0.436	-0.0188811\\
0.438	-0.0188671\\
0.44	-0.0188536\\
0.442	-0.0188406\\
0.444	-0.0188282\\
0.446	-0.0188162\\
0.448	-0.0188049\\
0.45	-0.0187941\\
0.452	-0.0187839\\
0.454	-0.0187743\\
0.456	-0.0187654\\
0.458	-0.0187571\\
0.46	-0.0187494\\
0.462	-0.0187424\\
0.464	-0.0187361\\
0.466	-0.0187304\\
0.468	-0.0187254\\
0.47	-0.018721\\
0.472	-0.0187173\\
0.474	-0.0187142\\
0.476	-0.0187119\\
0.478	-0.0187101\\
0.48	-0.018709\\
0.482	-0.0187085\\
0.484	-0.0187086\\
0.486	-0.0187094\\
0.488	-0.0187107\\
0.49	-0.0187126\\
0.492	-0.018715\\
0.494	-0.018718\\
0.496	-0.0187215\\
0.498	-0.0187255\\
0.5	-0.01873\\
};
\addplot [color=mycolor3, line width=1.0pt, forget plot]
  table[row sep=crcr]{%
0	0\\
0.002	-0.00183136\\
0.004	-0.00337792\\
0.006	-0.00470378\\
0.008	-0.00585724\\
0.01	-0.00687453\\
0.012	-0.00778283\\
0.014	-0.00860249\\
0.016	-0.00934882\\
0.018	-0.0100334\\
0.02	-0.010665\\
0.022	-0.0112505\\
0.024	-0.0117952\\
0.026	-0.0123034\\
0.028	-0.0127786\\
0.03	-0.0132236\\
0.032	-0.0136409\\
0.034	-0.0140326\\
0.036	-0.0144005\\
0.038	-0.0147464\\
0.04	-0.0150717\\
0.042	-0.0153778\\
0.044	-0.0156659\\
0.046	-0.0159373\\
0.048	-0.0161931\\
0.05	-0.0164343\\
0.052	-0.0166619\\
0.054	-0.0168769\\
0.056	-0.0170801\\
0.058	-0.0172724\\
0.06	-0.0174546\\
0.062	-0.0176274\\
0.064	-0.0177914\\
0.066	-0.0179474\\
0.068	-0.0180958\\
0.07	-0.0182372\\
0.072	-0.0183719\\
0.074	-0.0185005\\
0.076	-0.0186232\\
0.078	-0.0187403\\
0.08	-0.018852\\
0.082	-0.0189585\\
0.084	-0.0190601\\
0.086	-0.0191567\\
0.088	-0.0192485\\
0.09	-0.0193355\\
0.092	-0.0194178\\
0.094	-0.0194953\\
0.096	-0.0195681\\
0.098	-0.0196362\\
0.1	-0.0196995\\
0.102	-0.0197582\\
0.104	-0.019812\\
0.106	-0.0198611\\
0.108	-0.0199055\\
0.11	-0.0199451\\
0.112	-0.01998\\
0.114	-0.0200103\\
0.116	-0.0200359\\
0.118	-0.020057\\
0.12	-0.0200736\\
0.122	-0.0200858\\
0.124	-0.0200938\\
0.126	-0.0200976\\
0.128	-0.0200974\\
0.13	-0.0200933\\
0.132	-0.0200855\\
0.134	-0.0200741\\
0.136	-0.0200592\\
0.138	-0.0200412\\
0.14	-0.0200201\\
0.142	-0.0199961\\
0.144	-0.0199694\\
0.146	-0.0199403\\
0.148	-0.0199088\\
0.15	-0.0198753\\
0.152	-0.0198398\\
0.154	-0.0198026\\
0.156	-0.0197638\\
0.158	-0.0197237\\
0.16	-0.0196823\\
0.162	-0.01964\\
0.164	-0.0195969\\
0.166	-0.0195531\\
0.168	-0.0195087\\
0.17	-0.0194641\\
0.172	-0.0194192\\
0.174	-0.0193743\\
0.176	-0.0193295\\
0.178	-0.0192849\\
0.18	-0.0192406\\
0.182	-0.0191969\\
0.184	-0.0191537\\
0.186	-0.0191112\\
0.188	-0.0190696\\
0.19	-0.0190288\\
0.192	-0.0189891\\
0.194	-0.0189504\\
0.196	-0.0189129\\
0.198	-0.0188767\\
0.2	-0.0188418\\
0.202	-0.0188082\\
0.204	-0.0187761\\
0.206	-0.0187454\\
0.208	-0.0187163\\
0.21	-0.0186887\\
0.212	-0.0186627\\
0.214	-0.0186384\\
0.216	-0.0186157\\
0.218	-0.0185946\\
0.22	-0.0185752\\
0.222	-0.0185575\\
0.224	-0.0185414\\
0.226	-0.018527\\
0.228	-0.0185143\\
0.23	-0.0185032\\
0.232	-0.0184937\\
0.234	-0.0184858\\
0.236	-0.0184794\\
0.238	-0.0184746\\
0.24	-0.0184713\\
0.242	-0.0184694\\
0.244	-0.0184689\\
0.246	-0.0184698\\
0.248	-0.018472\\
0.25	-0.0184755\\
0.252	-0.0184801\\
0.254	-0.018486\\
0.256	-0.0184929\\
0.258	-0.0185009\\
0.26	-0.0185098\\
0.262	-0.0185197\\
0.264	-0.0185305\\
0.266	-0.018542\\
0.268	-0.0185543\\
0.27	-0.0185673\\
0.272	-0.0185809\\
0.274	-0.018595\\
0.276	-0.0186096\\
0.278	-0.0186247\\
0.28	-0.0186402\\
0.282	-0.0186559\\
0.284	-0.0186719\\
0.286	-0.0186881\\
0.288	-0.0187044\\
0.29	-0.0187209\\
0.292	-0.0187373\\
0.294	-0.0187537\\
0.296	-0.01877\\
0.298	-0.0187862\\
0.3	-0.0188022\\
0.302	-0.018818\\
0.304	-0.0188335\\
0.306	-0.0188488\\
0.308	-0.0188636\\
0.31	-0.0188781\\
0.312	-0.0188921\\
0.314	-0.0189057\\
0.316	-0.0189188\\
0.318	-0.0189314\\
0.32	-0.0189435\\
0.322	-0.018955\\
0.324	-0.0189659\\
0.326	-0.0189762\\
0.328	-0.0189859\\
0.33	-0.018995\\
0.332	-0.0190035\\
0.334	-0.0190114\\
0.336	-0.0190186\\
0.338	-0.0190252\\
0.34	-0.0190311\\
0.342	-0.0190364\\
0.344	-0.0190411\\
0.346	-0.0190452\\
0.348	-0.0190487\\
0.35	-0.0190516\\
0.352	-0.0190539\\
0.354	-0.0190557\\
0.356	-0.019057\\
0.358	-0.0190577\\
0.36	-0.0190579\\
0.362	-0.0190576\\
0.364	-0.0190569\\
0.366	-0.0190558\\
0.368	-0.0190543\\
0.37	-0.0190524\\
0.372	-0.0190501\\
0.374	-0.0190475\\
0.376	-0.0190446\\
0.378	-0.0190414\\
0.38	-0.0190379\\
0.382	-0.0190342\\
0.384	-0.0190303\\
0.386	-0.0190263\\
0.388	-0.019022\\
0.39	-0.0190176\\
0.392	-0.0190131\\
0.394	-0.0190086\\
0.396	-0.0190039\\
0.398	-0.0189992\\
0.4	-0.0189944\\
0.402	-0.0189896\\
0.404	-0.0189848\\
0.406	-0.01898\\
0.408	-0.0189753\\
0.41	-0.0189706\\
0.412	-0.0189659\\
0.414	-0.0189613\\
0.416	-0.0189568\\
0.418	-0.0189523\\
0.42	-0.018948\\
0.422	-0.0189437\\
0.424	-0.0189396\\
0.426	-0.0189356\\
0.428	-0.0189317\\
0.43	-0.0189279\\
0.432	-0.0189243\\
0.434	-0.0189208\\
0.436	-0.0189174\\
0.438	-0.0189142\\
0.44	-0.0189111\\
0.442	-0.0189082\\
0.444	-0.0189055\\
0.446	-0.0189029\\
0.448	-0.0189004\\
0.45	-0.0188981\\
0.452	-0.018896\\
0.454	-0.018894\\
0.456	-0.0188922\\
0.458	-0.0188905\\
0.46	-0.018889\\
0.462	-0.0188876\\
0.464	-0.0188864\\
0.466	-0.0188853\\
0.468	-0.0188844\\
0.47	-0.0188836\\
0.472	-0.018883\\
0.474	-0.0188825\\
0.476	-0.0188822\\
0.478	-0.018882\\
0.48	-0.0188819\\
0.482	-0.018882\\
0.484	-0.0188821\\
0.486	-0.0188825\\
0.488	-0.0188829\\
0.49	-0.0188834\\
0.492	-0.0188841\\
0.494	-0.0188849\\
0.496	-0.0188858\\
0.498	-0.0188867\\
0.5	-0.0188878\\
};
\addplot [color=mycolor4, line width=1.0pt, forget plot]
  table[row sep=crcr]{%
0	0\\
0.002	-0.00183135\\
0.004	-0.00337788\\
0.006	-0.00470365\\
0.008	-0.00585694\\
0.01	-0.00687396\\
0.012	-0.00778188\\
0.014	-0.00860104\\
0.016	-0.00934673\\
0.018	-0.0100305\\
0.02	-0.0106612\\
0.022	-0.0112456\\
0.024	-0.0117891\\
0.026	-0.012296\\
0.028	-0.0127697\\
0.03	-0.013213\\
0.032	-0.0136286\\
0.034	-0.0140185\\
0.036	-0.0143846\\
0.038	-0.0147285\\
0.04	-0.0150518\\
0.042	-0.0153558\\
0.044	-0.015642\\
0.046	-0.0159113\\
0.048	-0.0161651\\
0.05	-0.0164044\\
0.052	-0.0166301\\
0.054	-0.0168433\\
0.056	-0.0170448\\
0.058	-0.0172355\\
0.06	-0.0174162\\
0.062	-0.0175876\\
0.064	-0.0177504\\
0.066	-0.0179053\\
0.068	-0.0180528\\
0.07	-0.0181933\\
0.072	-0.0183275\\
0.074	-0.0184555\\
0.076	-0.0185779\\
0.078	-0.0186947\\
0.08	-0.0188064\\
0.082	-0.018913\\
0.084	-0.0190147\\
0.086	-0.0191115\\
0.088	-0.0192037\\
0.09	-0.0192911\\
0.092	-0.0193738\\
0.094	-0.0194519\\
0.096	-0.0195252\\
0.098	-0.0195939\\
0.1	-0.0196578\\
0.102	-0.019717\\
0.104	-0.0197714\\
0.106	-0.019821\\
0.108	-0.0198659\\
0.11	-0.019906\\
0.112	-0.0199413\\
0.114	-0.0199719\\
0.116	-0.0199978\\
0.118	-0.0200191\\
0.12	-0.0200359\\
0.122	-0.0200482\\
0.124	-0.0200562\\
0.126	-0.02006\\
0.128	-0.0200598\\
0.13	-0.0200557\\
0.132	-0.0200478\\
0.134	-0.0200363\\
0.136	-0.0200215\\
0.138	-0.0200034\\
0.14	-0.0199824\\
0.142	-0.0199585\\
0.144	-0.0199321\\
0.146	-0.0199033\\
0.148	-0.0198723\\
0.15	-0.0198393\\
0.152	-0.0198046\\
0.154	-0.0197684\\
0.156	-0.0197308\\
0.158	-0.019692\\
0.16	-0.0196523\\
0.162	-0.0196118\\
0.164	-0.0195708\\
0.166	-0.0195294\\
0.168	-0.0194877\\
0.17	-0.0194459\\
0.172	-0.0194043\\
0.174	-0.0193629\\
0.176	-0.0193218\\
0.178	-0.0192812\\
0.18	-0.0192413\\
0.182	-0.0192021\\
0.184	-0.0191637\\
0.186	-0.0191262\\
0.188	-0.0190898\\
0.19	-0.0190544\\
0.192	-0.0190203\\
0.194	-0.0189873\\
0.196	-0.0189556\\
0.198	-0.0189253\\
0.2	-0.0188963\\
0.202	-0.0188687\\
0.204	-0.0188425\\
0.206	-0.0188178\\
0.208	-0.0187945\\
0.21	-0.0187727\\
0.212	-0.0187523\\
0.214	-0.0187334\\
0.216	-0.0187159\\
0.218	-0.0186998\\
0.22	-0.0186851\\
0.222	-0.0186717\\
0.224	-0.0186598\\
0.226	-0.0186491\\
0.228	-0.0186396\\
0.23	-0.0186314\\
0.232	-0.0186244\\
0.234	-0.0186186\\
0.236	-0.0186138\\
0.238	-0.0186101\\
0.24	-0.0186074\\
0.242	-0.0186056\\
0.244	-0.0186048\\
0.246	-0.0186048\\
0.248	-0.0186056\\
0.25	-0.0186072\\
0.252	-0.0186095\\
0.254	-0.0186125\\
0.256	-0.0186162\\
0.258	-0.0186204\\
0.26	-0.0186252\\
0.262	-0.0186305\\
0.264	-0.0186363\\
0.266	-0.0186425\\
0.268	-0.0186491\\
0.27	-0.018656\\
0.272	-0.0186633\\
0.274	-0.0186709\\
0.276	-0.0186788\\
0.278	-0.0186869\\
0.28	-0.0186952\\
0.282	-0.0187037\\
0.284	-0.0187124\\
0.286	-0.0187212\\
0.288	-0.0187301\\
0.29	-0.018739\\
0.292	-0.0187481\\
0.294	-0.0187572\\
0.296	-0.0187663\\
0.298	-0.0187754\\
0.3	-0.0187845\\
0.302	-0.0187936\\
0.304	-0.0188026\\
0.306	-0.0188115\\
0.308	-0.0188203\\
0.31	-0.0188291\\
0.312	-0.0188376\\
0.314	-0.0188461\\
0.316	-0.0188544\\
0.318	-0.0188625\\
0.32	-0.0188704\\
0.322	-0.0188782\\
0.324	-0.0188857\\
0.326	-0.018893\\
0.328	-0.0189\\
0.33	-0.0189068\\
0.332	-0.0189134\\
0.334	-0.0189197\\
0.336	-0.0189257\\
0.338	-0.0189314\\
0.34	-0.0189369\\
0.342	-0.018942\\
0.344	-0.0189469\\
0.346	-0.0189515\\
0.348	-0.0189557\\
0.35	-0.0189597\\
0.352	-0.0189634\\
0.354	-0.0189668\\
0.356	-0.0189699\\
0.358	-0.0189727\\
0.36	-0.0189752\\
0.362	-0.0189775\\
0.364	-0.0189794\\
0.366	-0.0189811\\
0.368	-0.0189826\\
0.37	-0.0189837\\
0.372	-0.0189847\\
0.374	-0.0189854\\
0.376	-0.0189859\\
0.378	-0.0189861\\
0.38	-0.0189862\\
0.382	-0.0189861\\
0.384	-0.0189857\\
0.386	-0.0189852\\
0.388	-0.0189846\\
0.39	-0.0189838\\
0.392	-0.0189828\\
0.394	-0.0189818\\
0.396	-0.0189806\\
0.398	-0.0189792\\
0.4	-0.0189778\\
0.402	-0.0189764\\
0.404	-0.0189748\\
0.406	-0.0189732\\
0.408	-0.0189715\\
0.41	-0.0189697\\
0.412	-0.018968\\
0.414	-0.0189661\\
0.416	-0.0189643\\
0.418	-0.0189625\\
0.42	-0.0189606\\
0.422	-0.0189587\\
0.424	-0.0189569\\
0.426	-0.018955\\
0.428	-0.0189532\\
0.43	-0.0189514\\
0.432	-0.0189496\\
0.434	-0.0189479\\
0.436	-0.0189462\\
0.438	-0.0189445\\
0.44	-0.0189428\\
0.442	-0.0189412\\
0.444	-0.0189397\\
0.446	-0.0189382\\
0.448	-0.0189367\\
0.45	-0.0189353\\
0.452	-0.018934\\
0.454	-0.0189327\\
0.456	-0.0189314\\
0.458	-0.0189303\\
0.46	-0.0189291\\
0.462	-0.0189281\\
0.464	-0.018927\\
0.466	-0.0189261\\
0.468	-0.0189252\\
0.47	-0.0189243\\
0.472	-0.0189235\\
0.474	-0.0189228\\
0.476	-0.0189221\\
0.478	-0.0189215\\
0.48	-0.0189209\\
0.482	-0.0189204\\
0.484	-0.0189199\\
0.486	-0.0189194\\
0.488	-0.018919\\
0.49	-0.0189187\\
0.492	-0.0189184\\
0.494	-0.0189182\\
0.496	-0.018918\\
0.498	-0.0189178\\
0.5	-0.0189177\\
};
\addplot [color=mycolor5, line width=1.0pt, forget plot]
  table[row sep=crcr]{%
0	0\\
0.002	-0.00183135\\
0.004	-0.00337785\\
0.006	-0.00470358\\
0.008	-0.00585677\\
0.01	-0.00687365\\
0.012	-0.00778136\\
0.014	-0.00860025\\
0.016	-0.0093456\\
0.018	-0.010029\\
0.02	-0.0106592\\
0.022	-0.0112431\\
0.024	-0.011786\\
0.026	-0.0122921\\
0.028	-0.0127651\\
0.03	-0.0132077\\
0.032	-0.0136224\\
0.034	-0.0140115\\
0.036	-0.0143766\\
0.038	-0.0147196\\
0.04	-0.015042\\
0.042	-0.0153451\\
0.044	-0.0156302\\
0.046	-0.0158986\\
0.048	-0.0161513\\
0.05	-0.0163896\\
0.052	-0.0166142\\
0.054	-0.0168263\\
0.056	-0.0170267\\
0.058	-0.0172162\\
0.06	-0.0173957\\
0.062	-0.0175659\\
0.064	-0.0177274\\
0.066	-0.0178808\\
0.068	-0.0180268\\
0.07	-0.0181658\\
0.072	-0.0182983\\
0.074	-0.0184245\\
0.076	-0.018545\\
0.078	-0.0186598\\
0.08	-0.0187693\\
0.082	-0.0188736\\
0.084	-0.0189729\\
0.086	-0.0190672\\
0.088	-0.0191566\\
0.09	-0.0192413\\
0.092	-0.0193211\\
0.094	-0.0193962\\
0.096	-0.0194665\\
0.098	-0.019532\\
0.1	-0.0195928\\
0.102	-0.0196488\\
0.104	-0.0197\\
0.106	-0.0197465\\
0.108	-0.0197883\\
0.11	-0.0198253\\
0.112	-0.0198578\\
0.114	-0.0198857\\
0.116	-0.0199091\\
0.118	-0.0199281\\
0.12	-0.0199428\\
0.122	-0.0199534\\
0.124	-0.0199599\\
0.126	-0.0199627\\
0.128	-0.0199617\\
0.13	-0.0199572\\
0.132	-0.0199493\\
0.134	-0.0199384\\
0.136	-0.0199244\\
0.138	-0.0199078\\
0.14	-0.0198886\\
0.142	-0.0198671\\
0.144	-0.0198434\\
0.146	-0.0198178\\
0.148	-0.0197905\\
0.15	-0.0197617\\
0.152	-0.0197316\\
0.154	-0.0197003\\
0.156	-0.0196681\\
0.158	-0.019635\\
0.16	-0.0196014\\
0.162	-0.0195673\\
0.164	-0.0195329\\
0.166	-0.0194983\\
0.168	-0.0194637\\
0.17	-0.0194292\\
0.172	-0.0193948\\
0.174	-0.0193608\\
0.176	-0.0193272\\
0.178	-0.0192941\\
0.18	-0.0192615\\
0.182	-0.0192296\\
0.184	-0.0191984\\
0.186	-0.0191679\\
0.188	-0.0191383\\
0.19	-0.0191096\\
0.192	-0.0190818\\
0.194	-0.0190549\\
0.196	-0.0190289\\
0.198	-0.019004\\
0.2	-0.0189801\\
0.202	-0.0189572\\
0.204	-0.0189354\\
0.206	-0.0189146\\
0.208	-0.0188949\\
0.21	-0.0188762\\
0.212	-0.0188586\\
0.214	-0.018842\\
0.216	-0.0188265\\
0.218	-0.018812\\
0.22	-0.0187985\\
0.222	-0.018786\\
0.224	-0.0187744\\
0.226	-0.0187639\\
0.228	-0.0187542\\
0.23	-0.0187455\\
0.232	-0.0187377\\
0.234	-0.0187308\\
0.236	-0.0187247\\
0.238	-0.0187194\\
0.24	-0.0187149\\
0.242	-0.0187111\\
0.244	-0.0187081\\
0.246	-0.0187058\\
0.248	-0.0187042\\
0.25	-0.0187032\\
0.252	-0.0187028\\
0.254	-0.018703\\
0.256	-0.0187037\\
0.258	-0.018705\\
0.26	-0.0187068\\
0.262	-0.0187091\\
0.264	-0.0187118\\
0.266	-0.0187149\\
0.268	-0.0187184\\
0.27	-0.0187222\\
0.272	-0.0187264\\
0.274	-0.0187309\\
0.276	-0.0187356\\
0.278	-0.0187407\\
0.28	-0.0187459\\
0.282	-0.0187514\\
0.284	-0.018757\\
0.286	-0.0187628\\
0.288	-0.0187687\\
0.29	-0.0187748\\
0.292	-0.0187809\\
0.294	-0.0187871\\
0.296	-0.0187934\\
0.298	-0.0187997\\
0.3	-0.018806\\
0.302	-0.0188124\\
0.304	-0.0188187\\
0.306	-0.0188249\\
0.308	-0.0188311\\
0.31	-0.0188373\\
0.312	-0.0188434\\
0.314	-0.0188493\\
0.316	-0.0188552\\
0.318	-0.018861\\
0.32	-0.0188666\\
0.322	-0.0188721\\
0.324	-0.0188774\\
0.326	-0.0188826\\
0.328	-0.0188877\\
0.33	-0.0188925\\
0.332	-0.0188972\\
0.334	-0.0189017\\
0.336	-0.018906\\
0.338	-0.0189102\\
0.34	-0.0189141\\
0.342	-0.0189179\\
0.344	-0.0189214\\
0.346	-0.0189248\\
0.348	-0.018928\\
0.35	-0.018931\\
0.352	-0.0189338\\
0.354	-0.0189364\\
0.356	-0.0189388\\
0.358	-0.0189411\\
0.36	-0.0189431\\
0.362	-0.018945\\
0.364	-0.0189468\\
0.366	-0.0189483\\
0.368	-0.0189497\\
0.37	-0.018951\\
0.372	-0.0189521\\
0.374	-0.018953\\
0.376	-0.0189538\\
0.378	-0.0189545\\
0.38	-0.0189551\\
0.382	-0.0189555\\
0.384	-0.0189558\\
0.386	-0.018956\\
0.388	-0.0189561\\
0.39	-0.0189561\\
0.392	-0.0189561\\
0.394	-0.0189559\\
0.396	-0.0189557\\
0.398	-0.0189553\\
0.4	-0.0189549\\
0.402	-0.0189545\\
0.404	-0.018954\\
0.406	-0.0189534\\
0.408	-0.0189528\\
0.41	-0.0189522\\
0.412	-0.0189515\\
0.414	-0.0189507\\
0.416	-0.01895\\
0.418	-0.0189492\\
0.42	-0.0189484\\
0.422	-0.0189476\\
0.424	-0.0189468\\
0.426	-0.0189459\\
0.428	-0.0189451\\
0.43	-0.0189442\\
0.432	-0.0189434\\
0.434	-0.0189425\\
0.436	-0.0189417\\
0.438	-0.0189408\\
0.44	-0.01894\\
0.442	-0.0189392\\
0.444	-0.0189384\\
0.446	-0.0189376\\
0.448	-0.0189368\\
0.45	-0.018936\\
0.452	-0.0189353\\
0.454	-0.0189346\\
0.456	-0.0189339\\
0.458	-0.0189332\\
0.46	-0.0189326\\
0.462	-0.018932\\
0.464	-0.0189314\\
0.466	-0.0189308\\
0.468	-0.0189302\\
0.47	-0.0189297\\
0.472	-0.0189292\\
0.474	-0.0189288\\
0.476	-0.0189283\\
0.478	-0.0189279\\
0.48	-0.0189275\\
0.482	-0.0189271\\
0.484	-0.0189268\\
0.486	-0.0189265\\
0.488	-0.0189262\\
0.49	-0.0189259\\
0.492	-0.0189257\\
0.494	-0.0189254\\
0.496	-0.0189252\\
0.498	-0.0189251\\
0.5	-0.0189249\\
};
\addplot [color=mycolor6, line width=1.0pt, forget plot]
  table[row sep=crcr]{%
0	0\\
0.002	-0.00183137\\
0.004	-0.00337798\\
0.006	-0.00470399\\
0.008	-0.00585772\\
0.01	-0.00687546\\
0.012	-0.00778441\\
0.014	-0.00860498\\
0.016	-0.00935252\\
0.018	-0.0100386\\
0.02	-0.0106722\\
0.022	-0.0112601\\
0.024	-0.0118076\\
0.026	-0.0123191\\
0.028	-0.0127981\\
0.03	-0.0132474\\
0.032	-0.0136694\\
0.034	-0.0140662\\
0.036	-0.0144397\\
0.038	-0.0147914\\
0.04	-0.0151229\\
0.042	-0.0154352\\
0.044	-0.0157297\\
0.046	-0.0160074\\
0.048	-0.0162694\\
0.05	-0.0165164\\
0.052	-0.0167495\\
0.054	-0.0169695\\
0.056	-0.017177\\
0.058	-0.0173729\\
0.06	-0.0175578\\
0.062	-0.0177322\\
0.064	-0.0178969\\
0.066	-0.0180523\\
0.068	-0.0181989\\
0.07	-0.0183372\\
0.072	-0.0184675\\
0.074	-0.0185903\\
0.076	-0.0187058\\
0.078	-0.0188144\\
0.08	-0.0189163\\
0.082	-0.0190117\\
0.084	-0.0191009\\
0.086	-0.0191841\\
0.088	-0.0192613\\
0.09	-0.0193329\\
0.092	-0.0193988\\
0.094	-0.0194593\\
0.096	-0.0195146\\
0.098	-0.0195646\\
0.1	-0.0196097\\
0.102	-0.0196499\\
0.104	-0.0196854\\
0.106	-0.0197163\\
0.108	-0.0197428\\
0.11	-0.0197651\\
0.112	-0.0197833\\
0.114	-0.0197976\\
0.116	-0.0198081\\
0.118	-0.0198152\\
0.12	-0.0198189\\
0.122	-0.0198195\\
0.124	-0.0198171\\
0.126	-0.0198119\\
0.128	-0.0198042\\
0.13	-0.0197941\\
0.132	-0.0197819\\
0.134	-0.0197676\\
0.136	-0.0197515\\
0.138	-0.0197338\\
0.14	-0.0197146\\
0.142	-0.0196941\\
0.144	-0.0196725\\
0.146	-0.0196498\\
0.148	-0.0196263\\
0.15	-0.019602\\
0.152	-0.0195771\\
0.154	-0.0195516\\
0.156	-0.0195258\\
0.158	-0.0194997\\
0.16	-0.0194733\\
0.162	-0.0194468\\
0.164	-0.0194203\\
0.166	-0.0193937\\
0.168	-0.0193673\\
0.17	-0.0193409\\
0.172	-0.0193148\\
0.174	-0.0192889\\
0.176	-0.0192633\\
0.178	-0.019238\\
0.18	-0.019213\\
0.182	-0.0191885\\
0.184	-0.0191644\\
0.186	-0.0191408\\
0.188	-0.0191177\\
0.19	-0.0190951\\
0.192	-0.0190731\\
0.194	-0.0190517\\
0.196	-0.0190309\\
0.198	-0.0190107\\
0.2	-0.0189912\\
0.202	-0.0189723\\
0.204	-0.0189542\\
0.206	-0.0189368\\
0.208	-0.0189201\\
0.21	-0.0189042\\
0.212	-0.018889\\
0.214	-0.0188746\\
0.216	-0.0188609\\
0.218	-0.0188481\\
0.22	-0.018836\\
0.222	-0.0188247\\
0.224	-0.0188142\\
0.226	-0.0188044\\
0.228	-0.0187954\\
0.23	-0.0187872\\
0.232	-0.0187797\\
0.234	-0.0187729\\
0.236	-0.0187669\\
0.238	-0.0187616\\
0.24	-0.0187569\\
0.242	-0.0187529\\
0.244	-0.0187495\\
0.246	-0.0187467\\
0.248	-0.0187446\\
0.25	-0.0187429\\
0.252	-0.0187419\\
0.254	-0.0187413\\
0.256	-0.0187412\\
0.258	-0.0187416\\
0.26	-0.0187424\\
0.262	-0.0187436\\
0.264	-0.0187452\\
0.266	-0.0187472\\
0.268	-0.0187495\\
0.27	-0.0187521\\
0.272	-0.0187549\\
0.274	-0.018758\\
0.276	-0.0187614\\
0.278	-0.018765\\
0.28	-0.0187688\\
0.282	-0.0187727\\
0.284	-0.0187768\\
0.286	-0.018781\\
0.288	-0.0187854\\
0.29	-0.0187898\\
0.292	-0.0187943\\
0.294	-0.0187989\\
0.296	-0.0188036\\
0.298	-0.0188082\\
0.3	-0.0188129\\
0.302	-0.0188177\\
0.304	-0.0188224\\
0.306	-0.0188271\\
0.308	-0.0188317\\
0.31	-0.0188364\\
0.312	-0.018841\\
0.314	-0.0188455\\
0.316	-0.01885\\
0.318	-0.0188544\\
0.32	-0.0188587\\
0.322	-0.018863\\
0.324	-0.0188671\\
0.326	-0.0188712\\
0.328	-0.0188751\\
0.33	-0.018879\\
0.332	-0.0188827\\
0.334	-0.0188863\\
0.336	-0.0188898\\
0.338	-0.0188932\\
0.34	-0.0188965\\
0.342	-0.0188996\\
0.344	-0.0189026\\
0.346	-0.0189055\\
0.348	-0.0189082\\
0.35	-0.0189108\\
0.352	-0.0189133\\
0.354	-0.0189157\\
0.356	-0.0189179\\
0.358	-0.01892\\
0.36	-0.018922\\
0.362	-0.0189238\\
0.364	-0.0189256\\
0.366	-0.0189272\\
0.368	-0.0189287\\
0.37	-0.0189301\\
0.372	-0.0189314\\
0.374	-0.0189326\\
0.376	-0.0189337\\
0.378	-0.0189347\\
0.38	-0.0189355\\
0.382	-0.0189363\\
0.384	-0.0189371\\
0.386	-0.0189377\\
0.388	-0.0189382\\
0.39	-0.0189387\\
0.392	-0.0189391\\
0.394	-0.0189394\\
0.396	-0.0189397\\
0.398	-0.0189399\\
0.4	-0.01894\\
0.402	-0.0189401\\
0.404	-0.0189402\\
0.406	-0.0189401\\
0.408	-0.0189401\\
0.41	-0.01894\\
0.412	-0.0189399\\
0.414	-0.0189397\\
0.416	-0.0189395\\
0.418	-0.0189392\\
0.42	-0.018939\\
0.422	-0.0189387\\
0.424	-0.0189384\\
0.426	-0.0189381\\
0.428	-0.0189377\\
0.43	-0.0189374\\
0.432	-0.018937\\
0.434	-0.0189366\\
0.436	-0.0189362\\
0.438	-0.0189359\\
0.44	-0.0189355\\
0.442	-0.0189351\\
0.444	-0.0189347\\
0.446	-0.0189342\\
0.448	-0.0189338\\
0.45	-0.0189334\\
0.452	-0.018933\\
0.454	-0.0189326\\
0.456	-0.0189323\\
0.458	-0.0189319\\
0.46	-0.0189315\\
0.462	-0.0189311\\
0.464	-0.0189308\\
0.466	-0.0189304\\
0.468	-0.0189301\\
0.47	-0.0189298\\
0.472	-0.0189294\\
0.474	-0.0189291\\
0.476	-0.0189288\\
0.478	-0.0189285\\
0.48	-0.0189283\\
0.482	-0.018928\\
0.484	-0.0189278\\
0.486	-0.0189275\\
0.488	-0.0189273\\
0.49	-0.0189271\\
0.492	-0.0189269\\
0.494	-0.0189267\\
0.496	-0.0189265\\
0.498	-0.0189263\\
0.5	-0.0189262\\
};
\addplot [color=mycolor7, line width=1.0pt, forget plot]
  table[row sep=crcr]{%
0	0\\
0.002	-0.00183141\\
0.004	-0.00337829\\
0.006	-0.00470498\\
0.008	-0.00585993\\
0.01	-0.00687955\\
0.012	-0.00779111\\
0.014	-0.00861509\\
0.016	-0.00936686\\
0.018	-0.0100581\\
0.02	-0.0106975\\
0.022	-0.0112921\\
0.024	-0.0118471\\
0.026	-0.0123667\\
0.028	-0.0128543\\
0.03	-0.0133127\\
0.032	-0.0137441\\
0.034	-0.0141506\\
0.036	-0.0145337\\
0.038	-0.0148949\\
0.04	-0.0152355\\
0.042	-0.0155565\\
0.044	-0.0158591\\
0.046	-0.0161441\\
0.048	-0.0164124\\
0.05	-0.0166647\\
0.052	-0.0169019\\
0.054	-0.0171246\\
0.056	-0.0173336\\
0.058	-0.0175296\\
0.06	-0.0177131\\
0.062	-0.0178847\\
0.064	-0.0180452\\
0.066	-0.018195\\
0.068	-0.0183348\\
0.07	-0.018465\\
0.072	-0.0185861\\
0.074	-0.0186987\\
0.076	-0.0188031\\
0.078	-0.0188999\\
0.08	-0.0189895\\
0.082	-0.0190721\\
0.084	-0.0191482\\
0.086	-0.0192182\\
0.088	-0.0192823\\
0.09	-0.0193409\\
0.092	-0.0193941\\
0.094	-0.0194424\\
0.096	-0.019486\\
0.098	-0.019525\\
0.1	-0.0195597\\
0.102	-0.0195903\\
0.104	-0.0196171\\
0.106	-0.0196402\\
0.108	-0.0196597\\
0.11	-0.019676\\
0.112	-0.019689\\
0.114	-0.019699\\
0.116	-0.0197062\\
0.118	-0.0197107\\
0.12	-0.0197125\\
0.122	-0.019712\\
0.124	-0.0197091\\
0.126	-0.0197041\\
0.128	-0.019697\\
0.13	-0.0196881\\
0.132	-0.0196773\\
0.134	-0.0196649\\
0.136	-0.0196509\\
0.138	-0.0196355\\
0.14	-0.0196188\\
0.142	-0.0196009\\
0.144	-0.0195819\\
0.146	-0.019562\\
0.148	-0.0195412\\
0.15	-0.0195196\\
0.152	-0.0194974\\
0.154	-0.0194747\\
0.156	-0.0194514\\
0.158	-0.0194279\\
0.16	-0.019404\\
0.162	-0.01938\\
0.164	-0.0193558\\
0.166	-0.0193317\\
0.168	-0.0193076\\
0.17	-0.0192836\\
0.172	-0.0192597\\
0.174	-0.0192362\\
0.176	-0.0192129\\
0.178	-0.01919\\
0.18	-0.0191675\\
0.182	-0.0191454\\
0.184	-0.0191238\\
0.186	-0.0191027\\
0.188	-0.0190822\\
0.19	-0.0190622\\
0.192	-0.0190428\\
0.194	-0.0190241\\
0.196	-0.019006\\
0.198	-0.0189886\\
0.2	-0.0189718\\
0.202	-0.0189557\\
0.204	-0.0189403\\
0.206	-0.0189256\\
0.208	-0.0189116\\
0.21	-0.0188983\\
0.212	-0.0188856\\
0.214	-0.0188737\\
0.216	-0.0188624\\
0.218	-0.0188519\\
0.22	-0.018842\\
0.222	-0.0188327\\
0.224	-0.0188241\\
0.226	-0.0188162\\
0.228	-0.0188089\\
0.23	-0.0188022\\
0.232	-0.0187961\\
0.234	-0.0187906\\
0.236	-0.0187857\\
0.238	-0.0187813\\
0.24	-0.0187774\\
0.242	-0.0187741\\
0.244	-0.0187713\\
0.246	-0.018769\\
0.248	-0.0187671\\
0.25	-0.0187657\\
0.252	-0.0187647\\
0.254	-0.0187642\\
0.256	-0.018764\\
0.258	-0.0187641\\
0.26	-0.0187647\\
0.262	-0.0187655\\
0.264	-0.0187667\\
0.266	-0.0187682\\
0.268	-0.0187699\\
0.27	-0.0187719\\
0.272	-0.0187741\\
0.274	-0.0187765\\
0.276	-0.0187791\\
0.278	-0.018782\\
0.28	-0.0187849\\
0.282	-0.0187881\\
0.284	-0.0187913\\
0.286	-0.0187947\\
0.288	-0.0187981\\
0.29	-0.0188017\\
0.292	-0.0188053\\
0.294	-0.018809\\
0.296	-0.0188128\\
0.298	-0.0188165\\
0.3	-0.0188203\\
0.302	-0.0188242\\
0.304	-0.018828\\
0.306	-0.0188318\\
0.308	-0.0188356\\
0.31	-0.0188394\\
0.312	-0.0188431\\
0.314	-0.0188468\\
0.316	-0.0188504\\
0.318	-0.018854\\
0.32	-0.0188576\\
0.322	-0.0188611\\
0.324	-0.0188645\\
0.326	-0.0188678\\
0.328	-0.0188711\\
0.33	-0.0188742\\
0.332	-0.0188773\\
0.334	-0.0188803\\
0.336	-0.0188832\\
0.338	-0.0188861\\
0.34	-0.0188888\\
0.342	-0.0188914\\
0.344	-0.018894\\
0.346	-0.0188964\\
0.348	-0.0188988\\
0.35	-0.018901\\
0.352	-0.0189032\\
0.354	-0.0189052\\
0.356	-0.0189072\\
0.358	-0.0189091\\
0.36	-0.0189109\\
0.362	-0.0189126\\
0.364	-0.0189142\\
0.366	-0.0189157\\
0.368	-0.0189171\\
0.37	-0.0189185\\
0.372	-0.0189197\\
0.374	-0.0189209\\
0.376	-0.018922\\
0.378	-0.0189231\\
0.38	-0.0189241\\
0.382	-0.0189249\\
0.384	-0.0189258\\
0.386	-0.0189265\\
0.388	-0.0189272\\
0.39	-0.0189279\\
0.392	-0.0189285\\
0.394	-0.018929\\
0.396	-0.0189295\\
0.398	-0.0189299\\
0.4	-0.0189303\\
0.402	-0.0189306\\
0.404	-0.0189309\\
0.406	-0.0189312\\
0.408	-0.0189314\\
0.41	-0.0189316\\
0.412	-0.0189318\\
0.414	-0.0189319\\
0.416	-0.018932\\
0.418	-0.018932\\
0.42	-0.0189321\\
0.422	-0.0189321\\
0.424	-0.0189321\\
0.426	-0.0189321\\
0.428	-0.018932\\
0.43	-0.018932\\
0.432	-0.0189319\\
0.434	-0.0189318\\
0.436	-0.0189317\\
0.438	-0.0189316\\
0.44	-0.0189315\\
0.442	-0.0189313\\
0.444	-0.0189312\\
0.446	-0.0189311\\
0.448	-0.0189309\\
0.45	-0.0189307\\
0.452	-0.0189306\\
0.454	-0.0189304\\
0.456	-0.0189303\\
0.458	-0.0189301\\
0.46	-0.0189299\\
0.462	-0.0189298\\
0.464	-0.0189296\\
0.466	-0.0189294\\
0.468	-0.0189293\\
0.47	-0.0189291\\
0.472	-0.018929\\
0.474	-0.0189288\\
0.476	-0.0189286\\
0.478	-0.0189285\\
0.48	-0.0189284\\
0.482	-0.0189282\\
0.484	-0.0189281\\
0.486	-0.0189279\\
0.488	-0.0189278\\
0.49	-0.0189277\\
0.492	-0.0189276\\
0.494	-0.0189275\\
0.496	-0.0189273\\
0.498	-0.0189272\\
0.5	-0.0189271\\
};
\addplot [color=mycolor8, line width=1.0pt, forget plot]
  table[row sep=crcr]{%
0	0\\
0.002	-0.00183144\\
0.004	-0.00337855\\
0.006	-0.00470578\\
0.008	-0.00586167\\
0.01	-0.00688263\\
0.012	-0.00779596\\
0.014	-0.0086221\\
0.016	-0.00937638\\
0.018	-0.0100703\\
0.02	-0.0107128\\
0.022	-0.0113104\\
0.024	-0.0118685\\
0.026	-0.0123911\\
0.028	-0.0128815\\
0.03	-0.0133425\\
0.032	-0.0137761\\
0.034	-0.0141843\\
0.036	-0.0145687\\
0.038	-0.0149306\\
0.04	-0.0152713\\
0.042	-0.0155919\\
0.044	-0.0158933\\
0.046	-0.0161765\\
0.048	-0.0164423\\
0.05	-0.0166917\\
0.052	-0.0169254\\
0.054	-0.0171441\\
0.056	-0.0173487\\
0.058	-0.0175399\\
0.06	-0.0177183\\
0.062	-0.0178848\\
0.064	-0.0180399\\
0.066	-0.0181844\\
0.068	-0.0183188\\
0.07	-0.0184438\\
0.072	-0.0185599\\
0.074	-0.0186677\\
0.076	-0.0187677\\
0.078	-0.0188604\\
0.08	-0.0189461\\
0.082	-0.0190254\\
0.084	-0.0190986\\
0.086	-0.0191661\\
0.088	-0.0192282\\
0.09	-0.0192851\\
0.092	-0.0193372\\
0.094	-0.0193846\\
0.096	-0.0194277\\
0.098	-0.0194667\\
0.1	-0.0195017\\
0.102	-0.0195329\\
0.104	-0.0195604\\
0.106	-0.0195846\\
0.108	-0.0196054\\
0.11	-0.019623\\
0.112	-0.0196377\\
0.114	-0.0196494\\
0.116	-0.0196584\\
0.118	-0.0196647\\
0.12	-0.0196685\\
0.122	-0.01967\\
0.124	-0.0196692\\
0.126	-0.0196662\\
0.128	-0.0196612\\
0.13	-0.0196544\\
0.132	-0.0196458\\
0.134	-0.0196355\\
0.136	-0.0196237\\
0.138	-0.0196104\\
0.14	-0.0195959\\
0.142	-0.0195802\\
0.144	-0.0195634\\
0.146	-0.0195457\\
0.148	-0.0195271\\
0.15	-0.0195077\\
0.152	-0.0194878\\
0.154	-0.0194672\\
0.156	-0.0194462\\
0.158	-0.0194249\\
0.16	-0.0194032\\
0.162	-0.0193814\\
0.164	-0.0193594\\
0.166	-0.0193374\\
0.168	-0.0193154\\
0.17	-0.0192934\\
0.172	-0.0192716\\
0.174	-0.01925\\
0.176	-0.0192286\\
0.178	-0.0192075\\
0.18	-0.0191867\\
0.182	-0.0191663\\
0.184	-0.0191462\\
0.186	-0.0191266\\
0.188	-0.0191075\\
0.19	-0.0190888\\
0.192	-0.0190707\\
0.194	-0.0190531\\
0.196	-0.019036\\
0.198	-0.0190195\\
0.2	-0.0190035\\
0.202	-0.0189882\\
0.204	-0.0189734\\
0.206	-0.0189592\\
0.208	-0.0189457\\
0.21	-0.0189327\\
0.212	-0.0189204\\
0.214	-0.0189086\\
0.216	-0.0188975\\
0.218	-0.018887\\
0.22	-0.018877\\
0.222	-0.0188677\\
0.224	-0.018859\\
0.226	-0.0188508\\
0.228	-0.0188432\\
0.23	-0.0188362\\
0.232	-0.0188297\\
0.234	-0.0188238\\
0.236	-0.0188184\\
0.238	-0.0188135\\
0.24	-0.0188091\\
0.242	-0.0188051\\
0.244	-0.0188017\\
0.246	-0.0187986\\
0.248	-0.0187961\\
0.25	-0.0187939\\
0.252	-0.0187921\\
0.254	-0.0187907\\
0.256	-0.0187897\\
0.258	-0.018789\\
0.26	-0.0187886\\
0.262	-0.0187886\\
0.264	-0.0187888\\
0.266	-0.0187893\\
0.268	-0.0187901\\
0.27	-0.0187912\\
0.272	-0.0187924\\
0.274	-0.0187939\\
0.276	-0.0187956\\
0.278	-0.0187974\\
0.28	-0.0187994\\
0.282	-0.0188016\\
0.284	-0.0188039\\
0.286	-0.0188063\\
0.288	-0.0188089\\
0.29	-0.0188115\\
0.292	-0.0188143\\
0.294	-0.0188171\\
0.296	-0.0188199\\
0.298	-0.0188229\\
0.3	-0.0188258\\
0.302	-0.0188289\\
0.304	-0.0188319\\
0.306	-0.0188349\\
0.308	-0.018838\\
0.31	-0.0188411\\
0.312	-0.0188441\\
0.314	-0.0188471\\
0.316	-0.0188501\\
0.318	-0.0188531\\
0.32	-0.0188561\\
0.322	-0.018859\\
0.324	-0.0188619\\
0.326	-0.0188647\\
0.328	-0.0188675\\
0.33	-0.0188702\\
0.332	-0.0188728\\
0.334	-0.0188754\\
0.336	-0.018878\\
0.338	-0.0188805\\
0.34	-0.0188829\\
0.342	-0.0188852\\
0.344	-0.0188874\\
0.346	-0.0188896\\
0.348	-0.0188918\\
0.35	-0.0188938\\
0.352	-0.0188958\\
0.354	-0.0188977\\
0.356	-0.0188995\\
0.358	-0.0189012\\
0.36	-0.0189029\\
0.362	-0.0189045\\
0.364	-0.0189061\\
0.366	-0.0189075\\
0.368	-0.0189089\\
0.37	-0.0189103\\
0.372	-0.0189115\\
0.374	-0.0189127\\
0.376	-0.0189139\\
0.378	-0.018915\\
0.38	-0.018916\\
0.382	-0.0189169\\
0.384	-0.0189179\\
0.386	-0.0189187\\
0.388	-0.0189195\\
0.39	-0.0189203\\
0.392	-0.018921\\
0.394	-0.0189216\\
0.396	-0.0189222\\
0.398	-0.0189228\\
0.4	-0.0189233\\
0.402	-0.0189238\\
0.404	-0.0189243\\
0.406	-0.0189247\\
0.408	-0.0189251\\
0.41	-0.0189254\\
0.412	-0.0189258\\
0.414	-0.018926\\
0.416	-0.0189263\\
0.418	-0.0189266\\
0.42	-0.0189268\\
0.422	-0.018927\\
0.424	-0.0189271\\
0.426	-0.0189273\\
0.428	-0.0189274\\
0.43	-0.0189275\\
0.432	-0.0189276\\
0.434	-0.0189277\\
0.436	-0.0189278\\
0.438	-0.0189278\\
0.44	-0.0189279\\
0.442	-0.0189279\\
0.444	-0.0189279\\
0.446	-0.0189279\\
0.448	-0.0189279\\
0.45	-0.0189279\\
0.452	-0.0189279\\
0.454	-0.0189279\\
0.456	-0.0189278\\
0.458	-0.0189278\\
0.46	-0.0189278\\
0.462	-0.0189277\\
0.464	-0.0189277\\
0.466	-0.0189276\\
0.468	-0.0189276\\
0.47	-0.0189275\\
0.472	-0.0189274\\
0.474	-0.0189274\\
0.476	-0.0189273\\
0.478	-0.0189273\\
0.48	-0.0189272\\
0.482	-0.0189271\\
0.484	-0.0189271\\
0.486	-0.018927\\
0.488	-0.018927\\
0.49	-0.0189269\\
0.492	-0.0189268\\
0.494	-0.0189268\\
0.496	-0.0189267\\
0.498	-0.0189267\\
0.5	-0.0189266\\
};
\addplot [color=mycolor9, line width=1.0pt, forget plot]
  table[row sep=crcr]{%
0	0\\
0.002	-0.00183154\\
0.004	-0.00337927\\
0.006	-0.00470796\\
0.008	-0.00586631\\
0.01	-0.00689082\\
0.012	-0.00780871\\
0.014	-0.0086403\\
0.016	-0.00940078\\
0.018	-0.0101015\\
0.02	-0.010751\\
0.022	-0.0113558\\
0.024	-0.0119209\\
0.026	-0.0124502\\
0.028	-0.0129468\\
0.03	-0.0134133\\
0.032	-0.0138515\\
0.034	-0.0142634\\
0.036	-0.0146505\\
0.038	-0.015014\\
0.04	-0.0153552\\
0.042	-0.0156752\\
0.044	-0.015975\\
0.046	-0.0162557\\
0.048	-0.0165182\\
0.05	-0.0167634\\
0.052	-0.0169923\\
0.054	-0.0172058\\
0.056	-0.0174046\\
0.058	-0.0175898\\
0.06	-0.0177621\\
0.062	-0.0179224\\
0.064	-0.0180713\\
0.066	-0.0182098\\
0.068	-0.0183384\\
0.07	-0.0184578\\
0.072	-0.0185687\\
0.074	-0.0186716\\
0.076	-0.0187671\\
0.078	-0.0188556\\
0.08	-0.0189376\\
0.082	-0.0190135\\
0.084	-0.0190837\\
0.086	-0.0191484\\
0.088	-0.019208\\
0.09	-0.0192628\\
0.092	-0.0193129\\
0.094	-0.0193586\\
0.096	-0.0194001\\
0.098	-0.0194376\\
0.1	-0.0194712\\
0.102	-0.0195011\\
0.104	-0.0195275\\
0.106	-0.0195505\\
0.108	-0.0195702\\
0.11	-0.0195868\\
0.112	-0.0196005\\
0.114	-0.0196113\\
0.116	-0.0196193\\
0.118	-0.0196248\\
0.12	-0.0196279\\
0.122	-0.0196286\\
0.124	-0.0196272\\
0.126	-0.0196238\\
0.128	-0.0196184\\
0.13	-0.0196112\\
0.132	-0.0196024\\
0.134	-0.0195921\\
0.136	-0.0195803\\
0.138	-0.0195673\\
0.14	-0.0195531\\
0.142	-0.0195378\\
0.144	-0.0195216\\
0.146	-0.0195045\\
0.148	-0.0194868\\
0.15	-0.0194683\\
0.152	-0.0194493\\
0.154	-0.0194299\\
0.156	-0.0194101\\
0.158	-0.01939\\
0.16	-0.0193696\\
0.162	-0.0193491\\
0.164	-0.0193286\\
0.166	-0.019308\\
0.168	-0.0192874\\
0.17	-0.019267\\
0.172	-0.0192466\\
0.174	-0.0192265\\
0.176	-0.0192066\\
0.178	-0.019187\\
0.18	-0.0191677\\
0.182	-0.0191488\\
0.184	-0.0191302\\
0.186	-0.0191121\\
0.188	-0.0190943\\
0.19	-0.0190771\\
0.192	-0.0190603\\
0.194	-0.019044\\
0.196	-0.0190283\\
0.198	-0.019013\\
0.2	-0.0189983\\
0.202	-0.0189842\\
0.204	-0.0189706\\
0.206	-0.0189575\\
0.208	-0.0189451\\
0.21	-0.0189331\\
0.212	-0.0189218\\
0.214	-0.018911\\
0.216	-0.0189008\\
0.218	-0.0188912\\
0.22	-0.0188821\\
0.222	-0.0188735\\
0.224	-0.0188655\\
0.226	-0.018858\\
0.228	-0.0188511\\
0.23	-0.0188446\\
0.232	-0.0188387\\
0.234	-0.0188332\\
0.236	-0.0188282\\
0.238	-0.0188237\\
0.24	-0.0188196\\
0.242	-0.0188159\\
0.244	-0.0188127\\
0.246	-0.0188099\\
0.248	-0.0188074\\
0.25	-0.0188053\\
0.252	-0.0188036\\
0.254	-0.0188022\\
0.256	-0.0188012\\
0.258	-0.0188004\\
0.26	-0.0188\\
0.262	-0.0187998\\
0.264	-0.0187999\\
0.266	-0.0188003\\
0.268	-0.0188008\\
0.27	-0.0188016\\
0.272	-0.0188026\\
0.274	-0.0188038\\
0.276	-0.0188052\\
0.278	-0.0188067\\
0.28	-0.0188084\\
0.282	-0.0188102\\
0.284	-0.0188122\\
0.286	-0.0188142\\
0.288	-0.0188164\\
0.29	-0.0188187\\
0.292	-0.018821\\
0.294	-0.0188234\\
0.296	-0.0188259\\
0.298	-0.0188284\\
0.3	-0.018831\\
0.302	-0.0188336\\
0.304	-0.0188362\\
0.306	-0.0188388\\
0.308	-0.0188415\\
0.31	-0.0188441\\
0.312	-0.0188468\\
0.314	-0.0188494\\
0.316	-0.018852\\
0.318	-0.0188546\\
0.32	-0.0188572\\
0.322	-0.0188598\\
0.324	-0.0188623\\
0.326	-0.0188648\\
0.328	-0.0188672\\
0.33	-0.0188696\\
0.332	-0.018872\\
0.334	-0.0188743\\
0.336	-0.0188765\\
0.338	-0.0188787\\
0.34	-0.0188808\\
0.342	-0.0188829\\
0.344	-0.0188849\\
0.346	-0.0188869\\
0.348	-0.0188888\\
0.35	-0.0188906\\
0.352	-0.0188924\\
0.354	-0.0188941\\
0.356	-0.0188958\\
0.358	-0.0188974\\
0.36	-0.0188989\\
0.362	-0.0189004\\
0.364	-0.0189018\\
0.366	-0.0189032\\
0.368	-0.0189045\\
0.37	-0.0189058\\
0.372	-0.018907\\
0.374	-0.0189081\\
0.376	-0.0189092\\
0.378	-0.0189103\\
0.38	-0.0189113\\
0.382	-0.0189122\\
0.384	-0.0189131\\
0.386	-0.018914\\
0.388	-0.0189148\\
0.39	-0.0189155\\
0.392	-0.0189163\\
0.394	-0.018917\\
0.396	-0.0189176\\
0.398	-0.0189182\\
0.4	-0.0189188\\
0.402	-0.0189193\\
0.404	-0.0189199\\
0.406	-0.0189203\\
0.408	-0.0189208\\
0.41	-0.0189212\\
0.412	-0.0189216\\
0.414	-0.018922\\
0.416	-0.0189223\\
0.418	-0.0189227\\
0.42	-0.018923\\
0.422	-0.0189233\\
0.424	-0.0189235\\
0.426	-0.0189238\\
0.428	-0.018924\\
0.43	-0.0189242\\
0.432	-0.0189244\\
0.434	-0.0189246\\
0.436	-0.0189247\\
0.438	-0.0189249\\
0.44	-0.018925\\
0.442	-0.0189251\\
0.444	-0.0189253\\
0.446	-0.0189254\\
0.448	-0.0189255\\
0.45	-0.0189256\\
0.452	-0.0189256\\
0.454	-0.0189257\\
0.456	-0.0189258\\
0.458	-0.0189258\\
0.46	-0.0189259\\
0.462	-0.0189259\\
0.464	-0.0189259\\
0.466	-0.018926\\
0.468	-0.018926\\
0.47	-0.018926\\
0.472	-0.018926\\
0.474	-0.0189261\\
0.476	-0.0189261\\
0.478	-0.0189261\\
0.48	-0.0189261\\
0.482	-0.0189261\\
0.484	-0.0189261\\
0.486	-0.0189261\\
0.488	-0.0189261\\
0.49	-0.0189261\\
0.492	-0.0189261\\
0.494	-0.0189261\\
0.496	-0.0189261\\
0.498	-0.0189261\\
0.5	-0.0189261\\
};
\addplot [color=mycolor10, line width=1.0pt, forget plot]
  table[row sep=crcr]{%
0	0\\
0.002	-0.00183165\\
0.004	-0.00338005\\
0.006	-0.00471028\\
0.008	-0.00587111\\
0.01	-0.00689895\\
0.012	-0.00782084\\
0.014	-0.00865687\\
0.016	-0.00942191\\
0.018	-0.010127\\
0.02	-0.0107805\\
0.022	-0.0113887\\
0.024	-0.0119563\\
0.026	-0.0124871\\
0.028	-0.0129842\\
0.03	-0.01345\\
0.032	-0.0138867\\
0.034	-0.014296\\
0.036	-0.0146796\\
0.038	-0.0150389\\
0.04	-0.0153752\\
0.042	-0.01569\\
0.044	-0.0159842\\
0.046	-0.0162592\\
0.048	-0.0165161\\
0.05	-0.0167557\\
0.052	-0.0169793\\
0.054	-0.0171878\\
0.056	-0.0173821\\
0.058	-0.0175631\\
0.06	-0.0177318\\
0.062	-0.0178888\\
0.064	-0.0180351\\
0.066	-0.0181712\\
0.068	-0.0182979\\
0.07	-0.0184159\\
0.072	-0.0185256\\
0.074	-0.0186277\\
0.076	-0.0187225\\
0.078	-0.0188106\\
0.08	-0.0188924\\
0.082	-0.0189682\\
0.084	-0.0190383\\
0.086	-0.0191031\\
0.088	-0.0191628\\
0.09	-0.0192177\\
0.092	-0.019268\\
0.094	-0.019314\\
0.096	-0.0193557\\
0.098	-0.0193935\\
0.1	-0.0194275\\
0.102	-0.0194578\\
0.104	-0.0194847\\
0.106	-0.0195082\\
0.108	-0.0195285\\
0.11	-0.0195458\\
0.112	-0.0195602\\
0.114	-0.0195719\\
0.116	-0.0195809\\
0.118	-0.0195875\\
0.12	-0.0195917\\
0.122	-0.0195937\\
0.124	-0.0195936\\
0.126	-0.0195915\\
0.128	-0.0195876\\
0.13	-0.019582\\
0.132	-0.0195748\\
0.134	-0.0195661\\
0.136	-0.0195561\\
0.138	-0.0195448\\
0.14	-0.0195324\\
0.142	-0.0195189\\
0.144	-0.0195045\\
0.146	-0.0194893\\
0.148	-0.0194733\\
0.15	-0.0194567\\
0.152	-0.0194395\\
0.154	-0.0194218\\
0.156	-0.0194037\\
0.158	-0.0193854\\
0.16	-0.0193667\\
0.162	-0.0193479\\
0.164	-0.019329\\
0.166	-0.01931\\
0.168	-0.019291\\
0.17	-0.0192721\\
0.172	-0.0192533\\
0.174	-0.0192346\\
0.176	-0.0192161\\
0.178	-0.0191978\\
0.18	-0.0191798\\
0.182	-0.0191621\\
0.184	-0.0191448\\
0.186	-0.0191278\\
0.188	-0.0191111\\
0.19	-0.0190949\\
0.192	-0.0190791\\
0.194	-0.0190638\\
0.196	-0.0190489\\
0.198	-0.0190345\\
0.2	-0.0190205\\
0.202	-0.0190071\\
0.204	-0.0189941\\
0.206	-0.0189817\\
0.208	-0.0189697\\
0.21	-0.0189583\\
0.212	-0.0189474\\
0.214	-0.018937\\
0.216	-0.0189271\\
0.218	-0.0189177\\
0.22	-0.0189088\\
0.222	-0.0189004\\
0.224	-0.0188925\\
0.226	-0.0188851\\
0.228	-0.0188781\\
0.23	-0.0188716\\
0.232	-0.0188656\\
0.234	-0.01886\\
0.236	-0.0188549\\
0.238	-0.0188501\\
0.24	-0.0188458\\
0.242	-0.0188419\\
0.244	-0.0188384\\
0.246	-0.0188352\\
0.248	-0.0188324\\
0.25	-0.0188299\\
0.252	-0.0188278\\
0.254	-0.018826\\
0.256	-0.0188244\\
0.258	-0.0188232\\
0.26	-0.0188223\\
0.262	-0.0188216\\
0.264	-0.0188211\\
0.266	-0.0188209\\
0.268	-0.0188209\\
0.27	-0.0188211\\
0.272	-0.0188216\\
0.274	-0.0188222\\
0.276	-0.0188229\\
0.278	-0.0188239\\
0.28	-0.0188249\\
0.282	-0.0188261\\
0.284	-0.0188275\\
0.286	-0.0188289\\
0.288	-0.0188305\\
0.29	-0.0188321\\
0.292	-0.0188339\\
0.294	-0.0188357\\
0.296	-0.0188376\\
0.298	-0.0188395\\
0.3	-0.0188415\\
0.302	-0.0188435\\
0.304	-0.0188456\\
0.306	-0.0188476\\
0.308	-0.0188497\\
0.31	-0.0188519\\
0.312	-0.018854\\
0.314	-0.0188561\\
0.316	-0.0188583\\
0.318	-0.0188604\\
0.32	-0.0188625\\
0.322	-0.0188646\\
0.324	-0.0188667\\
0.326	-0.0188687\\
0.328	-0.0188707\\
0.33	-0.0188727\\
0.332	-0.0188747\\
0.334	-0.0188766\\
0.336	-0.0188785\\
0.338	-0.0188803\\
0.34	-0.0188822\\
0.342	-0.0188839\\
0.344	-0.0188856\\
0.346	-0.0188873\\
0.348	-0.018889\\
0.35	-0.0188905\\
0.352	-0.0188921\\
0.354	-0.0188936\\
0.356	-0.018895\\
0.358	-0.0188964\\
0.36	-0.0188978\\
0.362	-0.0188991\\
0.364	-0.0189003\\
0.366	-0.0189016\\
0.368	-0.0189027\\
0.37	-0.0189039\\
0.372	-0.0189049\\
0.374	-0.018906\\
0.376	-0.018907\\
0.378	-0.0189079\\
0.38	-0.0189089\\
0.382	-0.0189097\\
0.384	-0.0189106\\
0.386	-0.0189114\\
0.388	-0.0189122\\
0.39	-0.0189129\\
0.392	-0.0189136\\
0.394	-0.0189143\\
0.396	-0.0189149\\
0.398	-0.0189155\\
0.4	-0.0189161\\
0.402	-0.0189166\\
0.404	-0.0189172\\
0.406	-0.0189177\\
0.408	-0.0189181\\
0.41	-0.0189186\\
0.412	-0.018919\\
0.414	-0.0189194\\
0.416	-0.0189198\\
0.418	-0.0189202\\
0.42	-0.0189205\\
0.422	-0.0189209\\
0.424	-0.0189212\\
0.426	-0.0189215\\
0.428	-0.0189218\\
0.43	-0.018922\\
0.432	-0.0189223\\
0.434	-0.0189225\\
0.436	-0.0189227\\
0.438	-0.018923\\
0.44	-0.0189232\\
0.442	-0.0189233\\
0.444	-0.0189235\\
0.446	-0.0189237\\
0.448	-0.0189238\\
0.45	-0.018924\\
0.452	-0.0189241\\
0.454	-0.0189243\\
0.456	-0.0189244\\
0.458	-0.0189245\\
0.46	-0.0189246\\
0.462	-0.0189247\\
0.464	-0.0189248\\
0.466	-0.0189249\\
0.468	-0.018925\\
0.47	-0.0189251\\
0.472	-0.0189251\\
0.474	-0.0189252\\
0.476	-0.0189253\\
0.478	-0.0189253\\
0.48	-0.0189254\\
0.482	-0.0189254\\
0.484	-0.0189255\\
0.486	-0.0189255\\
0.488	-0.0189256\\
0.49	-0.0189256\\
0.492	-0.0189256\\
0.494	-0.0189257\\
0.496	-0.0189257\\
0.498	-0.0189257\\
0.5	-0.0189258\\
};
\addplot [color=mycolor11, line width=1.0pt, forget plot]
  table[row sep=crcr]{%
0	0\\
0.002	-0.00183223\\
0.004	-0.0033837\\
0.006	-0.00472\\
0.008	-0.0058893\\
0.01	-0.00692694\\
0.012	-0.00785889\\
0.014	-0.00870434\\
0.016	-0.0094775\\
0.018	-0.010189\\
0.02	-0.0108471\\
0.022	-0.0114578\\
0.024	-0.0120262\\
0.026	-0.0125563\\
0.028	-0.0130512\\
0.03	-0.0135137\\
0.032	-0.0139463\\
0.034	-0.0143508\\
0.036	-0.0147292\\
0.038	-0.0150831\\
0.04	-0.0154141\\
0.042	-0.0157235\\
0.044	-0.0160126\\
0.046	-0.0162828\\
0.048	-0.016535\\
0.05	-0.0167705\\
0.052	-0.0169902\\
0.054	-0.0171952\\
0.056	-0.0173864\\
0.058	-0.0175646\\
0.06	-0.0177307\\
0.062	-0.0178855\\
0.064	-0.0180297\\
0.066	-0.0181639\\
0.068	-0.0182889\\
0.07	-0.0184053\\
0.072	-0.0185135\\
0.074	-0.0186141\\
0.076	-0.0187076\\
0.078	-0.0187943\\
0.08	-0.0188748\\
0.082	-0.0189493\\
0.084	-0.0190182\\
0.086	-0.0190817\\
0.088	-0.0191403\\
0.09	-0.0191941\\
0.092	-0.0192433\\
0.094	-0.0192882\\
0.096	-0.0193291\\
0.098	-0.019366\\
0.1	-0.0193991\\
0.102	-0.0194287\\
0.104	-0.0194549\\
0.106	-0.0194779\\
0.108	-0.0194978\\
0.11	-0.0195147\\
0.112	-0.0195288\\
0.114	-0.0195403\\
0.116	-0.0195493\\
0.118	-0.0195558\\
0.12	-0.0195601\\
0.122	-0.0195623\\
0.124	-0.0195624\\
0.126	-0.0195607\\
0.128	-0.0195572\\
0.13	-0.019552\\
0.132	-0.0195454\\
0.134	-0.0195373\\
0.136	-0.0195279\\
0.138	-0.0195174\\
0.14	-0.0195057\\
0.142	-0.019493\\
0.144	-0.0194795\\
0.146	-0.0194652\\
0.148	-0.0194501\\
0.15	-0.0194345\\
0.152	-0.0194183\\
0.154	-0.0194017\\
0.156	-0.0193846\\
0.158	-0.0193673\\
0.16	-0.0193498\\
0.162	-0.019332\\
0.164	-0.0193142\\
0.166	-0.0192963\\
0.168	-0.0192783\\
0.17	-0.0192605\\
0.172	-0.0192427\\
0.174	-0.0192251\\
0.176	-0.0192076\\
0.178	-0.0191904\\
0.18	-0.0191734\\
0.182	-0.0191566\\
0.184	-0.0191402\\
0.186	-0.0191242\\
0.188	-0.0191084\\
0.19	-0.0190931\\
0.192	-0.0190781\\
0.194	-0.0190636\\
0.196	-0.0190495\\
0.198	-0.0190358\\
0.2	-0.0190226\\
0.202	-0.0190099\\
0.204	-0.0189976\\
0.206	-0.0189858\\
0.208	-0.0189744\\
0.21	-0.0189635\\
0.212	-0.0189532\\
0.214	-0.0189432\\
0.216	-0.0189338\\
0.218	-0.0189248\\
0.22	-0.0189163\\
0.222	-0.0189083\\
0.224	-0.0189007\\
0.226	-0.0188936\\
0.228	-0.0188869\\
0.23	-0.0188807\\
0.232	-0.0188748\\
0.234	-0.0188694\\
0.236	-0.0188644\\
0.238	-0.0188598\\
0.24	-0.0188555\\
0.242	-0.0188517\\
0.244	-0.0188482\\
0.246	-0.018845\\
0.248	-0.0188422\\
0.25	-0.0188397\\
0.252	-0.0188375\\
0.254	-0.0188356\\
0.256	-0.018834\\
0.258	-0.0188326\\
0.26	-0.0188315\\
0.262	-0.0188307\\
0.264	-0.0188301\\
0.266	-0.0188297\\
0.268	-0.0188295\\
0.27	-0.0188295\\
0.272	-0.0188297\\
0.274	-0.0188301\\
0.276	-0.0188306\\
0.278	-0.0188313\\
0.28	-0.0188321\\
0.282	-0.0188331\\
0.284	-0.0188341\\
0.286	-0.0188353\\
0.288	-0.0188366\\
0.29	-0.018838\\
0.292	-0.0188394\\
0.294	-0.018841\\
0.296	-0.0188426\\
0.298	-0.0188442\\
0.3	-0.0188459\\
0.302	-0.0188477\\
0.304	-0.0188494\\
0.306	-0.0188513\\
0.308	-0.0188531\\
0.31	-0.0188549\\
0.312	-0.0188568\\
0.314	-0.0188587\\
0.316	-0.0188606\\
0.318	-0.0188624\\
0.32	-0.0188643\\
0.322	-0.0188661\\
0.324	-0.018868\\
0.326	-0.0188698\\
0.328	-0.0188716\\
0.33	-0.0188734\\
0.332	-0.0188752\\
0.334	-0.0188769\\
0.336	-0.0188786\\
0.338	-0.0188802\\
0.34	-0.0188819\\
0.342	-0.0188835\\
0.344	-0.018885\\
0.346	-0.0188865\\
0.348	-0.018888\\
0.35	-0.0188895\\
0.352	-0.0188909\\
0.354	-0.0188922\\
0.356	-0.0188936\\
0.358	-0.0188949\\
0.36	-0.0188961\\
0.362	-0.0188973\\
0.364	-0.0188985\\
0.366	-0.0188996\\
0.368	-0.0189007\\
0.37	-0.0189018\\
0.372	-0.0189028\\
0.374	-0.0189038\\
0.376	-0.0189047\\
0.378	-0.0189056\\
0.38	-0.0189065\\
0.382	-0.0189073\\
0.384	-0.0189081\\
0.386	-0.0189089\\
0.388	-0.0189097\\
0.39	-0.0189104\\
0.392	-0.0189111\\
0.394	-0.0189117\\
0.396	-0.0189124\\
0.398	-0.018913\\
0.4	-0.0189136\\
0.402	-0.0189141\\
0.404	-0.0189146\\
0.406	-0.0189152\\
0.408	-0.0189157\\
0.41	-0.0189161\\
0.412	-0.0189166\\
0.414	-0.018917\\
0.416	-0.0189174\\
0.418	-0.0189178\\
0.42	-0.0189182\\
0.422	-0.0189186\\
0.424	-0.0189189\\
0.426	-0.0189192\\
0.428	-0.0189195\\
0.43	-0.0189198\\
0.432	-0.0189201\\
0.434	-0.0189204\\
0.436	-0.0189207\\
0.438	-0.0189209\\
0.44	-0.0189212\\
0.442	-0.0189214\\
0.444	-0.0189216\\
0.446	-0.0189218\\
0.448	-0.018922\\
0.45	-0.0189222\\
0.452	-0.0189224\\
0.454	-0.0189226\\
0.456	-0.0189227\\
0.458	-0.0189229\\
0.46	-0.018923\\
0.462	-0.0189232\\
0.464	-0.0189233\\
0.466	-0.0189235\\
0.468	-0.0189236\\
0.47	-0.0189237\\
0.472	-0.0189238\\
0.474	-0.0189239\\
0.476	-0.018924\\
0.478	-0.0189241\\
0.48	-0.0189242\\
0.482	-0.0189243\\
0.484	-0.0189244\\
0.486	-0.0189245\\
0.488	-0.0189245\\
0.49	-0.0189246\\
0.492	-0.0189247\\
0.494	-0.0189248\\
0.496	-0.0189248\\
0.498	-0.0189249\\
0.5	-0.0189249\\
};
\end{axis}

\begin{axis}[%
width={\figscale*0.642in},
height={\figscale*0.17in},
at={(\figscale*5.577in,\figscale*1.895in)},
scale only axis,
separate axis lines,
every outer x axis line/.append style={black},
every x tick label/.append style={font=\figticfont, yshift=\figinsettickshift},
every x tick/.append style={black},
xmin=0,
xmax=0.5,
xtick={  0, 0.4},
every outer y axis line/.append style={black},
every y tick label/.append style={font=\figticfont},
every y tick/.append style={black},
ymin=-0.023,
ymax=-0.0215,
ytick={-0.023,-0.0215},
axis background/.style={fill=white},
ylabel style={rotate=-90, at={(axis description cs:-0.145,0.5)}, anchor=east}
]
\addplot [color=mycolor1, line width=1.4pt, forget plot]
  table[row sep=crcr]{%
0	-0.0065\\
0.002	-0.0065\\
0.004	-0.0065\\
0.006	-0.0065\\
0.008	-0.0065\\
0.01	-0.00699977\\
0.012	-0.00791183\\
0.014	-0.00872528\\
0.016	-0.00945819\\
0.018	-0.0101252\\
0.02	-0.0107383\\
0.022	-0.0113069\\
0.024	-0.0118387\\
0.026	-0.0123396\\
0.028	-0.0128143\\
0.03	-0.0132664\\
0.032	-0.0136986\\
0.034	-0.0141129\\
0.036	-0.0145108\\
0.038	-0.0148935\\
0.04	-0.0152615\\
0.042	-0.0156154\\
0.044	-0.0159555\\
0.046	-0.016282\\
0.048	-0.016595\\
0.05	-0.0168948\\
0.052	-0.0171813\\
0.054	-0.0174547\\
0.056	-0.0177152\\
0.058	-0.017963\\
0.06	-0.0181983\\
0.062	-0.0184213\\
0.064	-0.0186325\\
0.066	-0.0188322\\
0.068	-0.0190207\\
0.07	-0.0191986\\
0.072	-0.0193662\\
0.074	-0.019524\\
0.076	-0.0196726\\
0.078	-0.0198123\\
0.08	-0.0199438\\
0.082	-0.0200675\\
0.084	-0.0201838\\
0.086	-0.0202933\\
0.088	-0.0203963\\
0.09	-0.0204935\\
0.092	-0.0205851\\
0.094	-0.0206715\\
0.096	-0.0207533\\
0.098	-0.0208307\\
0.1	-0.020904\\
0.102	-0.0209737\\
0.104	-0.02104\\
0.106	-0.0211031\\
0.108	-0.0211633\\
0.11	-0.0212208\\
0.112	-0.0212759\\
0.114	-0.0213286\\
0.116	-0.0213791\\
0.118	-0.0214276\\
0.12	-0.0214742\\
0.122	-0.0215189\\
0.124	-0.0215618\\
0.126	-0.0216029\\
0.128	-0.0216422\\
0.13	-0.0216798\\
0.132	-0.0217157\\
0.134	-0.0217498\\
0.136	-0.0217822\\
0.138	-0.0218127\\
0.14	-0.0218414\\
0.142	-0.0218682\\
0.144	-0.021893\\
0.146	-0.0219159\\
0.148	-0.0219368\\
0.15	-0.0219556\\
0.152	-0.0219723\\
0.154	-0.021987\\
0.156	-0.0219995\\
0.158	-0.02201\\
0.16	-0.0220184\\
0.162	-0.0220247\\
0.164	-0.022029\\
0.166	-0.0220314\\
0.168	-0.0220318\\
0.17	-0.0220304\\
0.172	-0.0220273\\
0.174	-0.0220225\\
0.176	-0.0220162\\
0.178	-0.0220086\\
0.18	-0.0219997\\
0.182	-0.0219896\\
0.184	-0.0219786\\
0.186	-0.0219668\\
0.188	-0.0219544\\
0.19	-0.0219415\\
0.192	-0.0219283\\
0.194	-0.0219149\\
0.196	-0.0219017\\
0.198	-0.0218886\\
0.2	-0.0218759\\
0.202	-0.0218638\\
0.204	-0.0218524\\
0.206	-0.0218418\\
0.208	-0.0218323\\
0.21	-0.0218239\\
0.212	-0.0218168\\
0.214	-0.0218111\\
0.216	-0.0218069\\
0.218	-0.0218044\\
0.22	-0.0218035\\
0.222	-0.0218044\\
0.224	-0.0218072\\
0.226	-0.0218119\\
0.228	-0.0218185\\
0.23	-0.0218271\\
0.232	-0.0218377\\
0.234	-0.0218503\\
0.236	-0.021865\\
0.238	-0.0218816\\
0.24	-0.0219002\\
0.242	-0.0219208\\
0.244	-0.0219432\\
0.246	-0.0219676\\
0.248	-0.0219937\\
0.25	-0.0220216\\
0.252	-0.0220512\\
0.254	-0.0220823\\
0.256	-0.022115\\
0.258	-0.0221492\\
0.26	-0.0221847\\
0.262	-0.0222214\\
0.264	-0.0222593\\
0.266	-0.0222983\\
0.268	-0.0223382\\
0.27	-0.022379\\
0.272	-0.0224206\\
0.274	-0.0224628\\
0.276	-0.0225056\\
0.278	-0.0225489\\
0.28	-0.0225926\\
0.282	-0.0226365\\
0.284	-0.0226806\\
0.286	-0.0227248\\
0.288	-0.022769\\
0.29	-0.0228132\\
0.292	-0.0228572\\
0.294	-0.0229009\\
0.296	-0.0229443\\
0.298	-0.0229873\\
0.3	-0.0230299\\
0.302	-0.0230719\\
0.304	-0.0231134\\
0.306	-0.0231542\\
0.308	-0.0231943\\
0.31	-0.0232337\\
0.312	-0.0232723\\
0.314	-0.02331\\
0.316	-0.0233469\\
0.318	-0.0233828\\
0.32	-0.0234179\\
0.322	-0.0234519\\
0.324	-0.023485\\
0.326	-0.023517\\
0.328	-0.023548\\
0.33	-0.023578\\
0.332	-0.0236069\\
0.334	-0.0236347\\
0.336	-0.0236615\\
0.338	-0.0236871\\
0.34	-0.0237117\\
0.342	-0.0237352\\
0.344	-0.0237577\\
0.346	-0.023779\\
0.348	-0.0237993\\
0.35	-0.0238186\\
0.352	-0.0238368\\
0.354	-0.023854\\
0.356	-0.0238701\\
0.358	-0.0238853\\
0.36	-0.0238996\\
0.362	-0.0239129\\
0.364	-0.0239252\\
0.366	-0.0239367\\
0.368	-0.0239474\\
0.37	-0.0239572\\
0.372	-0.0239662\\
0.374	-0.0239744\\
0.376	-0.0239819\\
0.378	-0.0239887\\
0.38	-0.0239948\\
0.382	-0.0240004\\
0.384	-0.0240053\\
0.386	-0.0240097\\
0.388	-0.0240136\\
0.39	-0.024017\\
0.392	-0.02402\\
0.394	-0.0240226\\
0.396	-0.0240248\\
0.398	-0.0240268\\
0.4	-0.0240285\\
0.402	-0.02403\\
0.404	-0.0240313\\
0.406	-0.0240325\\
0.408	-0.0240336\\
0.41	-0.0240347\\
0.412	-0.0240357\\
0.414	-0.0240368\\
0.416	-0.0240379\\
0.418	-0.0240391\\
0.42	-0.0240405\\
0.422	-0.024042\\
0.424	-0.0240437\\
0.426	-0.0240457\\
0.428	-0.0240479\\
0.43	-0.0240504\\
0.432	-0.0240532\\
0.434	-0.0240563\\
0.436	-0.0240598\\
0.438	-0.0240637\\
0.44	-0.0240679\\
0.442	-0.0240725\\
0.444	-0.0240775\\
0.446	-0.024083\\
0.448	-0.0240889\\
0.45	-0.0240952\\
0.452	-0.0241019\\
0.454	-0.024109\\
0.456	-0.0241166\\
0.458	-0.0241246\\
0.46	-0.024133\\
0.462	-0.0241418\\
0.464	-0.024151\\
0.466	-0.0241606\\
0.468	-0.0241705\\
0.47	-0.0241807\\
0.472	-0.0241913\\
0.474	-0.0242022\\
0.476	-0.0242133\\
0.478	-0.0242247\\
0.48	-0.0242363\\
0.482	-0.0242481\\
0.484	-0.0242601\\
0.486	-0.0242722\\
0.488	-0.0242844\\
0.49	-0.0242966\\
0.492	-0.0243089\\
0.494	-0.0243213\\
0.496	-0.0243335\\
0.498	-0.0243458\\
0.5	-0.0243579\\
};
\addplot [color=mycolor2, line width=1.0pt, forget plot]
  table[row sep=crcr]{%
0	-0.0065\\
0.002	-0.0065\\
0.004	-0.0065\\
0.006	-0.0065\\
0.008	-0.0065\\
0.01	-0.00699977\\
0.012	-0.00791183\\
0.014	-0.00872528\\
0.016	-0.00945819\\
0.018	-0.0101252\\
0.02	-0.0107383\\
0.022	-0.0113069\\
0.024	-0.0118387\\
0.026	-0.0123396\\
0.028	-0.0128143\\
0.03	-0.0132664\\
0.032	-0.0136986\\
0.034	-0.0141129\\
0.036	-0.0145109\\
0.038	-0.0148936\\
0.04	-0.0152616\\
0.042	-0.0156155\\
0.044	-0.0159557\\
0.046	-0.0162822\\
0.048	-0.0165954\\
0.05	-0.0168951\\
0.052	-0.0171817\\
0.054	-0.0174553\\
0.056	-0.0177158\\
0.058	-0.0179637\\
0.06	-0.0181991\\
0.062	-0.0184223\\
0.064	-0.0186336\\
0.066	-0.0188334\\
0.068	-0.0190222\\
0.07	-0.0192002\\
0.072	-0.019368\\
0.074	-0.0195261\\
0.076	-0.0196749\\
0.078	-0.0198149\\
0.08	-0.0199467\\
0.082	-0.0200707\\
0.084	-0.0201873\\
0.086	-0.0202971\\
0.088	-0.0204006\\
0.09	-0.0204981\\
0.092	-0.0205902\\
0.094	-0.0206771\\
0.096	-0.0207593\\
0.098	-0.0208372\\
0.1	-0.0209112\\
0.102	-0.0209814\\
0.104	-0.0210482\\
0.106	-0.021112\\
0.108	-0.0211728\\
0.11	-0.021231\\
0.112	-0.0212867\\
0.114	-0.0213402\\
0.116	-0.0213915\\
0.118	-0.0214407\\
0.12	-0.0214881\\
0.122	-0.0215336\\
0.124	-0.0215773\\
0.126	-0.0216193\\
0.128	-0.0216595\\
0.13	-0.021698\\
0.132	-0.0217348\\
0.134	-0.0217698\\
0.136	-0.0218031\\
0.138	-0.0218346\\
0.14	-0.0218642\\
0.142	-0.0218919\\
0.144	-0.0219177\\
0.146	-0.0219416\\
0.148	-0.0219634\\
0.15	-0.0219832\\
0.152	-0.0220009\\
0.154	-0.0220165\\
0.156	-0.02203\\
0.158	-0.0220414\\
0.16	-0.0220507\\
0.162	-0.0220579\\
0.164	-0.0220631\\
0.166	-0.0220663\\
0.168	-0.0220676\\
0.17	-0.022067\\
0.172	-0.0220647\\
0.174	-0.0220607\\
0.176	-0.0220552\\
0.178	-0.0220482\\
0.18	-0.0220399\\
0.182	-0.0220305\\
0.184	-0.0220201\\
0.186	-0.0220089\\
0.188	-0.0219969\\
0.19	-0.0219845\\
0.192	-0.0219717\\
0.194	-0.0219587\\
0.196	-0.0219458\\
0.198	-0.021933\\
0.2	-0.0219205\\
0.202	-0.0219086\\
0.204	-0.0218973\\
0.206	-0.0218868\\
0.208	-0.0218773\\
0.21	-0.0218688\\
0.212	-0.0218617\\
0.214	-0.0218558\\
0.216	-0.0218514\\
0.218	-0.0218486\\
0.22	-0.0218474\\
0.222	-0.021848\\
0.224	-0.0218504\\
0.226	-0.0218546\\
0.228	-0.0218607\\
0.23	-0.0218688\\
0.232	-0.0218788\\
0.234	-0.0218908\\
0.236	-0.0219047\\
0.238	-0.0219206\\
0.24	-0.0219385\\
0.242	-0.0219582\\
0.244	-0.0219799\\
0.246	-0.0220034\\
0.248	-0.0220286\\
0.25	-0.0220556\\
0.252	-0.0220842\\
0.254	-0.0221144\\
0.256	-0.0221461\\
0.258	-0.0221793\\
0.26	-0.0222137\\
0.262	-0.0222494\\
0.264	-0.0222863\\
0.266	-0.0223242\\
0.268	-0.0223631\\
0.27	-0.0224029\\
0.272	-0.0224434\\
0.274	-0.0224846\\
0.276	-0.0225264\\
0.278	-0.0225686\\
0.28	-0.0226112\\
0.282	-0.0226542\\
0.284	-0.0226973\\
0.286	-0.0227405\\
0.288	-0.0227838\\
0.29	-0.022827\\
0.292	-0.0228701\\
0.294	-0.0229129\\
0.296	-0.0229555\\
0.298	-0.0229977\\
0.3	-0.0230394\\
0.302	-0.0230807\\
0.304	-0.0231214\\
0.306	-0.0231615\\
0.308	-0.023201\\
0.31	-0.0232397\\
0.312	-0.0232777\\
0.314	-0.0233149\\
0.316	-0.0233513\\
0.318	-0.0233867\\
0.32	-0.0234213\\
0.322	-0.023455\\
0.324	-0.0234877\\
0.326	-0.0235194\\
0.328	-0.0235501\\
0.33	-0.0235798\\
0.332	-0.0236084\\
0.334	-0.0236361\\
0.336	-0.0236627\\
0.338	-0.0236882\\
0.34	-0.0237127\\
0.342	-0.0237361\\
0.344	-0.0237585\\
0.346	-0.0237799\\
0.348	-0.0238002\\
0.35	-0.0238194\\
0.352	-0.0238377\\
0.354	-0.023855\\
0.356	-0.0238712\\
0.358	-0.0238865\\
0.36	-0.0239009\\
0.362	-0.0239143\\
0.364	-0.0239268\\
0.366	-0.0239385\\
0.368	-0.0239493\\
0.37	-0.0239592\\
0.372	-0.0239684\\
0.374	-0.0239768\\
0.376	-0.0239844\\
0.378	-0.0239914\\
0.38	-0.0239977\\
0.382	-0.0240034\\
0.384	-0.0240084\\
0.386	-0.0240129\\
0.388	-0.0240169\\
0.39	-0.0240204\\
0.392	-0.0240235\\
0.394	-0.0240262\\
0.396	-0.0240285\\
0.398	-0.0240305\\
0.4	-0.0240323\\
0.402	-0.0240338\\
0.404	-0.024035\\
0.406	-0.0240362\\
0.408	-0.0240372\\
0.41	-0.0240382\\
0.412	-0.0240391\\
0.414	-0.02404\\
0.416	-0.024041\\
0.418	-0.024042\\
0.42	-0.0240432\\
0.422	-0.0240444\\
0.424	-0.0240459\\
0.426	-0.0240476\\
0.428	-0.0240495\\
0.43	-0.0240516\\
0.432	-0.0240541\\
0.434	-0.0240568\\
0.436	-0.0240599\\
0.438	-0.0240633\\
0.44	-0.0240671\\
0.442	-0.0240713\\
0.444	-0.0240759\\
0.446	-0.0240809\\
0.448	-0.0240863\\
0.45	-0.0240921\\
0.452	-0.0240983\\
0.454	-0.0241049\\
0.456	-0.0241119\\
0.458	-0.0241194\\
0.46	-0.0241273\\
0.462	-0.0241355\\
0.464	-0.0241442\\
0.466	-0.0241532\\
0.468	-0.0241626\\
0.47	-0.0241723\\
0.472	-0.0241824\\
0.474	-0.0241927\\
0.476	-0.0242033\\
0.478	-0.0242142\\
0.48	-0.0242253\\
0.482	-0.0242366\\
0.484	-0.0242481\\
0.486	-0.0242598\\
0.488	-0.0242715\\
0.49	-0.0242834\\
0.492	-0.0242953\\
0.494	-0.0243072\\
0.496	-0.0243192\\
0.498	-0.0243311\\
0.5	-0.0243429\\
};
\addplot [color=mycolor3, line width=1.0pt, forget plot]
  table[row sep=crcr]{%
0	-0.0065\\
0.002	-0.0065\\
0.004	-0.0065\\
0.006	-0.0065\\
0.008	-0.0065\\
0.01	-0.00699977\\
0.012	-0.00791183\\
0.014	-0.00872529\\
0.016	-0.0094582\\
0.018	-0.0101253\\
0.02	-0.0107383\\
0.022	-0.011307\\
0.024	-0.0118388\\
0.026	-0.0123397\\
0.028	-0.0128144\\
0.03	-0.0132666\\
0.032	-0.0136988\\
0.034	-0.0141133\\
0.036	-0.0145113\\
0.038	-0.0148941\\
0.04	-0.0152623\\
0.042	-0.0156164\\
0.044	-0.0159567\\
0.046	-0.0162835\\
0.048	-0.0165968\\
0.05	-0.0168969\\
0.052	-0.0171838\\
0.054	-0.0174577\\
0.056	-0.0177187\\
0.058	-0.017967\\
0.06	-0.0182029\\
0.062	-0.0184267\\
0.064	-0.0186386\\
0.066	-0.0188391\\
0.068	-0.0190285\\
0.07	-0.0192073\\
0.072	-0.019376\\
0.074	-0.019535\\
0.076	-0.0196847\\
0.078	-0.0198258\\
0.08	-0.0199586\\
0.082	-0.0200838\\
0.084	-0.0202016\\
0.086	-0.0203127\\
0.088	-0.0204175\\
0.09	-0.0205165\\
0.092	-0.02061\\
0.094	-0.0206985\\
0.096	-0.0207823\\
0.098	-0.0208618\\
0.1	-0.0209374\\
0.102	-0.0210094\\
0.104	-0.0210781\\
0.106	-0.0211437\\
0.108	-0.0212064\\
0.11	-0.0212665\\
0.112	-0.0213243\\
0.114	-0.0213797\\
0.116	-0.021433\\
0.118	-0.0214844\\
0.12	-0.0215338\\
0.122	-0.0215814\\
0.124	-0.0216272\\
0.126	-0.0216713\\
0.128	-0.0217136\\
0.13	-0.0217543\\
0.132	-0.0217932\\
0.134	-0.0218303\\
0.136	-0.0218656\\
0.138	-0.0218991\\
0.14	-0.0219308\\
0.142	-0.0219605\\
0.144	-0.0219882\\
0.146	-0.022014\\
0.148	-0.0220376\\
0.15	-0.0220592\\
0.152	-0.0220786\\
0.154	-0.0220959\\
0.156	-0.022111\\
0.158	-0.0221239\\
0.16	-0.0221346\\
0.162	-0.0221432\\
0.164	-0.0221496\\
0.166	-0.022154\\
0.168	-0.0221563\\
0.17	-0.0221567\\
0.172	-0.0221552\\
0.174	-0.0221519\\
0.176	-0.022147\\
0.178	-0.0221405\\
0.18	-0.0221327\\
0.182	-0.0221235\\
0.184	-0.0221133\\
0.186	-0.0221021\\
0.188	-0.02209\\
0.19	-0.0220774\\
0.192	-0.0220642\\
0.194	-0.0220508\\
0.196	-0.0220373\\
0.198	-0.0220238\\
0.2	-0.0220105\\
0.202	-0.0219976\\
0.204	-0.0219852\\
0.206	-0.0219736\\
0.208	-0.0219628\\
0.21	-0.0219529\\
0.212	-0.0219442\\
0.214	-0.0219368\\
0.216	-0.0219307\\
0.218	-0.021926\\
0.22	-0.0219229\\
0.222	-0.0219215\\
0.224	-0.0219218\\
0.226	-0.0219238\\
0.228	-0.0219277\\
0.23	-0.0219335\\
0.232	-0.0219411\\
0.234	-0.0219507\\
0.236	-0.0219621\\
0.238	-0.0219755\\
0.24	-0.0219908\\
0.242	-0.022008\\
0.244	-0.022027\\
0.246	-0.0220478\\
0.248	-0.0220704\\
0.25	-0.0220947\\
0.252	-0.0221206\\
0.254	-0.0221482\\
0.256	-0.0221772\\
0.258	-0.0222077\\
0.26	-0.0222395\\
0.262	-0.0222726\\
0.264	-0.0223069\\
0.266	-0.0223423\\
0.268	-0.0223787\\
0.27	-0.022416\\
0.272	-0.0224541\\
0.274	-0.022493\\
0.276	-0.0225325\\
0.278	-0.0225725\\
0.28	-0.022613\\
0.282	-0.0226539\\
0.284	-0.022695\\
0.286	-0.0227364\\
0.288	-0.0227778\\
0.29	-0.0228193\\
0.292	-0.0228608\\
0.294	-0.0229021\\
0.296	-0.0229432\\
0.298	-0.0229841\\
0.3	-0.0230246\\
0.302	-0.0230647\\
0.304	-0.0231044\\
0.306	-0.0231435\\
0.308	-0.0231821\\
0.31	-0.0232201\\
0.312	-0.0232574\\
0.314	-0.023294\\
0.316	-0.0233299\\
0.318	-0.023365\\
0.32	-0.0233993\\
0.322	-0.0234327\\
0.324	-0.0234653\\
0.326	-0.023497\\
0.328	-0.0235277\\
0.33	-0.0235575\\
0.332	-0.0235864\\
0.334	-0.0236143\\
0.336	-0.0236412\\
0.338	-0.0236671\\
0.34	-0.023692\\
0.342	-0.0237159\\
0.344	-0.0237388\\
0.346	-0.0237607\\
0.348	-0.0237816\\
0.35	-0.0238015\\
0.352	-0.0238204\\
0.354	-0.0238383\\
0.356	-0.0238552\\
0.358	-0.0238712\\
0.36	-0.0238863\\
0.362	-0.0239004\\
0.364	-0.0239136\\
0.366	-0.0239259\\
0.368	-0.0239373\\
0.37	-0.0239479\\
0.372	-0.0239576\\
0.374	-0.0239666\\
0.376	-0.0239748\\
0.378	-0.0239822\\
0.38	-0.023989\\
0.382	-0.023995\\
0.384	-0.0240005\\
0.386	-0.0240053\\
0.388	-0.0240095\\
0.39	-0.0240132\\
0.392	-0.0240163\\
0.394	-0.0240191\\
0.396	-0.0240213\\
0.398	-0.0240232\\
0.4	-0.0240248\\
0.402	-0.024026\\
0.404	-0.024027\\
0.406	-0.0240277\\
0.408	-0.0240282\\
0.41	-0.0240286\\
0.412	-0.0240289\\
0.414	-0.0240291\\
0.416	-0.0240292\\
0.418	-0.0240294\\
0.42	-0.0240296\\
0.422	-0.0240298\\
0.424	-0.0240302\\
0.426	-0.0240307\\
0.428	-0.0240313\\
0.43	-0.0240322\\
0.432	-0.0240333\\
0.434	-0.0240346\\
0.436	-0.0240362\\
0.438	-0.0240381\\
0.44	-0.0240403\\
0.442	-0.0240428\\
0.444	-0.0240457\\
0.446	-0.024049\\
0.448	-0.0240526\\
0.45	-0.0240566\\
0.452	-0.024061\\
0.454	-0.0240658\\
0.456	-0.0240711\\
0.458	-0.0240767\\
0.46	-0.0240827\\
0.462	-0.0240891\\
0.464	-0.0240959\\
0.466	-0.0241031\\
0.468	-0.0241107\\
0.47	-0.0241186\\
0.472	-0.0241269\\
0.474	-0.0241355\\
0.476	-0.0241444\\
0.478	-0.0241536\\
0.48	-0.0241631\\
0.482	-0.0241729\\
0.484	-0.0241828\\
0.486	-0.024193\\
0.488	-0.0242034\\
0.49	-0.0242139\\
0.492	-0.0242246\\
0.494	-0.0242354\\
0.496	-0.0242462\\
0.498	-0.0242571\\
0.5	-0.024268\\
};
\addplot [color=mycolor4, line width=1.0pt, forget plot]
  table[row sep=crcr]{%
0	-0.0065\\
0.002	-0.0065\\
0.004	-0.0065\\
0.006	-0.0065\\
0.008	-0.0065\\
0.01	-0.00699977\\
0.012	-0.00791183\\
0.014	-0.00872529\\
0.016	-0.00945821\\
0.018	-0.0101253\\
0.02	-0.0107384\\
0.022	-0.011307\\
0.024	-0.0118388\\
0.026	-0.0123398\\
0.028	-0.0128146\\
0.03	-0.0132668\\
0.032	-0.0136991\\
0.034	-0.0141136\\
0.036	-0.0145118\\
0.038	-0.0148946\\
0.04	-0.0152629\\
0.042	-0.0156172\\
0.044	-0.0159577\\
0.046	-0.0162847\\
0.048	-0.0165984\\
0.05	-0.0168987\\
0.052	-0.017186\\
0.054	-0.0174603\\
0.056	-0.0177217\\
0.058	-0.0179705\\
0.06	-0.018207\\
0.062	-0.0184313\\
0.064	-0.0186439\\
0.066	-0.0188452\\
0.068	-0.0190354\\
0.07	-0.0192151\\
0.072	-0.0193848\\
0.074	-0.0195448\\
0.076	-0.0196957\\
0.078	-0.019838\\
0.08	-0.0199722\\
0.082	-0.0200987\\
0.084	-0.0202181\\
0.086	-0.0203308\\
0.088	-0.0204373\\
0.09	-0.0205381\\
0.092	-0.0206335\\
0.094	-0.020724\\
0.096	-0.02081\\
0.098	-0.0208918\\
0.1	-0.0209697\\
0.102	-0.0210441\\
0.104	-0.0211153\\
0.106	-0.0211836\\
0.108	-0.0212491\\
0.11	-0.0213121\\
0.112	-0.0213727\\
0.114	-0.0214312\\
0.116	-0.0214877\\
0.118	-0.0215423\\
0.12	-0.021595\\
0.122	-0.021646\\
0.124	-0.0216952\\
0.126	-0.0217428\\
0.128	-0.0217888\\
0.13	-0.021833\\
0.132	-0.0218756\\
0.134	-0.0219165\\
0.136	-0.0219556\\
0.138	-0.0219929\\
0.14	-0.0220284\\
0.142	-0.0220619\\
0.144	-0.0220935\\
0.146	-0.022123\\
0.148	-0.0221505\\
0.15	-0.0221758\\
0.152	-0.022199\\
0.154	-0.0222199\\
0.156	-0.0222386\\
0.158	-0.0222551\\
0.16	-0.0222693\\
0.162	-0.0222813\\
0.164	-0.022291\\
0.166	-0.0222986\\
0.168	-0.022304\\
0.17	-0.0223073\\
0.172	-0.0223086\\
0.174	-0.022308\\
0.176	-0.0223056\\
0.178	-0.0223015\\
0.18	-0.0222958\\
0.182	-0.0222886\\
0.184	-0.0222801\\
0.186	-0.0222705\\
0.188	-0.0222598\\
0.19	-0.0222483\\
0.192	-0.0222361\\
0.194	-0.0222234\\
0.196	-0.0222103\\
0.198	-0.022197\\
0.2	-0.0221837\\
0.202	-0.0221705\\
0.204	-0.0221576\\
0.206	-0.0221451\\
0.208	-0.0221332\\
0.21	-0.0221221\\
0.212	-0.0221118\\
0.214	-0.0221025\\
0.216	-0.0220943\\
0.218	-0.0220873\\
0.22	-0.0220816\\
0.222	-0.0220773\\
0.224	-0.0220744\\
0.226	-0.0220731\\
0.228	-0.0220734\\
0.23	-0.0220753\\
0.232	-0.0220788\\
0.234	-0.022084\\
0.236	-0.022091\\
0.238	-0.0220996\\
0.24	-0.02211\\
0.242	-0.022122\\
0.244	-0.0221357\\
0.246	-0.0221511\\
0.248	-0.022168\\
0.25	-0.0221866\\
0.252	-0.0222066\\
0.254	-0.0222281\\
0.256	-0.022251\\
0.258	-0.0222752\\
0.26	-0.0223007\\
0.262	-0.0223274\\
0.264	-0.0223553\\
0.266	-0.0223842\\
0.268	-0.022414\\
0.27	-0.0224448\\
0.272	-0.0224764\\
0.274	-0.0225087\\
0.276	-0.0225418\\
0.278	-0.0225754\\
0.28	-0.0226095\\
0.282	-0.022644\\
0.284	-0.0226789\\
0.286	-0.0227142\\
0.288	-0.0227496\\
0.29	-0.0227852\\
0.292	-0.0228208\\
0.294	-0.0228565\\
0.296	-0.0228922\\
0.298	-0.0229278\\
0.3	-0.0229632\\
0.302	-0.0229984\\
0.304	-0.0230333\\
0.306	-0.023068\\
0.308	-0.0231023\\
0.31	-0.0231362\\
0.312	-0.0231697\\
0.314	-0.0232027\\
0.316	-0.0232353\\
0.318	-0.0232672\\
0.32	-0.0232987\\
0.322	-0.0233295\\
0.324	-0.0233597\\
0.326	-0.0233893\\
0.328	-0.0234182\\
0.33	-0.0234465\\
0.332	-0.023474\\
0.334	-0.0235008\\
0.336	-0.0235269\\
0.338	-0.0235523\\
0.34	-0.0235769\\
0.342	-0.0236008\\
0.344	-0.0236239\\
0.346	-0.0236462\\
0.348	-0.0236678\\
0.35	-0.0236885\\
0.352	-0.0237085\\
0.354	-0.0237278\\
0.356	-0.0237462\\
0.358	-0.0237639\\
0.36	-0.0237808\\
0.362	-0.0237969\\
0.364	-0.0238123\\
0.366	-0.023827\\
0.368	-0.0238409\\
0.37	-0.0238541\\
0.372	-0.0238666\\
0.374	-0.0238783\\
0.376	-0.0238894\\
0.378	-0.0238999\\
0.38	-0.0239097\\
0.382	-0.0239189\\
0.384	-0.0239274\\
0.386	-0.0239354\\
0.388	-0.0239428\\
0.39	-0.0239497\\
0.392	-0.0239561\\
0.394	-0.023962\\
0.396	-0.0239674\\
0.398	-0.0239724\\
0.4	-0.023977\\
0.402	-0.0239812\\
0.404	-0.023985\\
0.406	-0.0239886\\
0.408	-0.0239918\\
0.41	-0.0239948\\
0.412	-0.0239976\\
0.414	-0.0240002\\
0.416	-0.0240026\\
0.418	-0.0240048\\
0.42	-0.0240069\\
0.422	-0.024009\\
0.424	-0.024011\\
0.426	-0.024013\\
0.428	-0.024015\\
0.43	-0.024017\\
0.432	-0.024019\\
0.434	-0.0240211\\
0.436	-0.0240234\\
0.438	-0.0240257\\
0.44	-0.0240282\\
0.442	-0.0240308\\
0.444	-0.0240336\\
0.446	-0.0240366\\
0.448	-0.0240398\\
0.45	-0.0240432\\
0.452	-0.0240468\\
0.454	-0.0240506\\
0.456	-0.0240547\\
0.458	-0.0240591\\
0.46	-0.0240637\\
0.462	-0.0240685\\
0.464	-0.0240737\\
0.466	-0.024079\\
0.468	-0.0240846\\
0.47	-0.0240905\\
0.472	-0.0240966\\
0.474	-0.0241029\\
0.476	-0.0241095\\
0.478	-0.0241163\\
0.48	-0.0241233\\
0.482	-0.0241305\\
0.484	-0.0241378\\
0.486	-0.0241454\\
0.488	-0.0241531\\
0.49	-0.0241609\\
0.492	-0.0241689\\
0.494	-0.0241769\\
0.496	-0.0241851\\
0.498	-0.0241933\\
0.5	-0.0242015\\
};
\addplot [color=mycolor5, line width=1.0pt, forget plot]
  table[row sep=crcr]{%
0	-0.0065\\
0.002	-0.0065\\
0.004	-0.0065\\
0.006	-0.0065\\
0.008	-0.0065\\
0.01	-0.00699977\\
0.012	-0.00791183\\
0.014	-0.00872528\\
0.016	-0.0094582\\
0.018	-0.0101253\\
0.02	-0.0107383\\
0.022	-0.011307\\
0.024	-0.0118388\\
0.026	-0.0123398\\
0.028	-0.0128146\\
0.03	-0.0132668\\
0.032	-0.0136992\\
0.034	-0.0141137\\
0.036	-0.0145119\\
0.038	-0.0148949\\
0.04	-0.0152633\\
0.042	-0.0156177\\
0.044	-0.0159584\\
0.046	-0.0162856\\
0.048	-0.0165994\\
0.05	-0.0169001\\
0.052	-0.0171876\\
0.054	-0.0174622\\
0.056	-0.0177241\\
0.058	-0.0179733\\
0.06	-0.0182103\\
0.062	-0.0184352\\
0.064	-0.0186484\\
0.066	-0.0188504\\
0.068	-0.0190414\\
0.07	-0.0192219\\
0.072	-0.0193925\\
0.074	-0.0195535\\
0.076	-0.0197055\\
0.078	-0.019849\\
0.08	-0.0199844\\
0.082	-0.0201123\\
0.084	-0.0202331\\
0.086	-0.0203474\\
0.088	-0.0204555\\
0.09	-0.020558\\
0.092	-0.0206552\\
0.094	-0.0207476\\
0.096	-0.0208356\\
0.098	-0.0209195\\
0.1	-0.0209996\\
0.102	-0.0210763\\
0.104	-0.0211498\\
0.106	-0.0212205\\
0.108	-0.0212885\\
0.11	-0.0213541\\
0.112	-0.0214174\\
0.114	-0.0214787\\
0.116	-0.0215379\\
0.118	-0.0215954\\
0.12	-0.021651\\
0.122	-0.0217049\\
0.124	-0.0217572\\
0.126	-0.0218078\\
0.128	-0.0218568\\
0.13	-0.0219041\\
0.132	-0.0219497\\
0.134	-0.0219936\\
0.136	-0.0220358\\
0.138	-0.0220761\\
0.14	-0.0221146\\
0.142	-0.0221511\\
0.144	-0.0221856\\
0.146	-0.022218\\
0.148	-0.0222483\\
0.15	-0.0222764\\
0.152	-0.0223022\\
0.154	-0.0223257\\
0.156	-0.0223469\\
0.158	-0.0223657\\
0.16	-0.0223821\\
0.162	-0.0223962\\
0.164	-0.0224079\\
0.166	-0.0224173\\
0.168	-0.0224244\\
0.17	-0.0224292\\
0.172	-0.0224319\\
0.174	-0.0224325\\
0.176	-0.022431\\
0.178	-0.0224277\\
0.18	-0.0224226\\
0.182	-0.0224159\\
0.184	-0.0224077\\
0.186	-0.0223981\\
0.188	-0.0223872\\
0.19	-0.0223753\\
0.192	-0.0223625\\
0.194	-0.022349\\
0.196	-0.0223349\\
0.198	-0.0223204\\
0.2	-0.0223056\\
0.202	-0.0222908\\
0.204	-0.0222761\\
0.206	-0.0222616\\
0.208	-0.0222474\\
0.21	-0.0222338\\
0.212	-0.0222209\\
0.214	-0.0222088\\
0.216	-0.0221976\\
0.218	-0.0221874\\
0.22	-0.0221784\\
0.222	-0.0221706\\
0.224	-0.0221641\\
0.226	-0.022159\\
0.228	-0.0221554\\
0.23	-0.0221532\\
0.232	-0.0221526\\
0.234	-0.0221536\\
0.236	-0.0221563\\
0.238	-0.0221605\\
0.24	-0.0221664\\
0.242	-0.0221739\\
0.244	-0.022183\\
0.246	-0.0221938\\
0.248	-0.0222061\\
0.25	-0.02222\\
0.252	-0.0222354\\
0.254	-0.0222522\\
0.256	-0.0222705\\
0.258	-0.0222902\\
0.26	-0.0223112\\
0.262	-0.0223335\\
0.264	-0.022357\\
0.266	-0.0223816\\
0.268	-0.0224072\\
0.27	-0.0224339\\
0.272	-0.0224615\\
0.274	-0.02249\\
0.276	-0.0225193\\
0.278	-0.0225493\\
0.28	-0.02258\\
0.282	-0.0226112\\
0.284	-0.022643\\
0.286	-0.0226752\\
0.288	-0.0227079\\
0.29	-0.0227408\\
0.292	-0.022774\\
0.294	-0.0228074\\
0.296	-0.022841\\
0.298	-0.0228747\\
0.3	-0.0229084\\
0.302	-0.022942\\
0.304	-0.0229757\\
0.306	-0.0230091\\
0.308	-0.0230425\\
0.31	-0.0230756\\
0.312	-0.0231085\\
0.314	-0.0231411\\
0.316	-0.0231733\\
0.318	-0.0232053\\
0.32	-0.0232368\\
0.322	-0.0232678\\
0.324	-0.0232984\\
0.326	-0.0233286\\
0.328	-0.0233582\\
0.33	-0.0233872\\
0.332	-0.0234157\\
0.334	-0.0234436\\
0.336	-0.0234708\\
0.338	-0.0234975\\
0.34	-0.0235235\\
0.342	-0.0235488\\
0.344	-0.0235734\\
0.346	-0.0235973\\
0.348	-0.0236205\\
0.35	-0.0236429\\
0.352	-0.0236646\\
0.354	-0.0236856\\
0.356	-0.0237058\\
0.358	-0.0237253\\
0.36	-0.023744\\
0.362	-0.0237619\\
0.364	-0.023779\\
0.366	-0.0237954\\
0.368	-0.023811\\
0.37	-0.0238258\\
0.372	-0.0238399\\
0.374	-0.0238532\\
0.376	-0.0238658\\
0.378	-0.0238776\\
0.38	-0.0238887\\
0.382	-0.0238991\\
0.384	-0.0239089\\
0.386	-0.0239179\\
0.388	-0.0239263\\
0.39	-0.023934\\
0.392	-0.0239411\\
0.394	-0.0239477\\
0.396	-0.0239536\\
0.398	-0.023959\\
0.4	-0.0239639\\
0.402	-0.0239684\\
0.404	-0.0239723\\
0.406	-0.0239758\\
0.408	-0.0239789\\
0.41	-0.0239817\\
0.412	-0.0239841\\
0.414	-0.0239862\\
0.416	-0.023988\\
0.418	-0.0239896\\
0.42	-0.023991\\
0.422	-0.0239922\\
0.424	-0.0239933\\
0.426	-0.0239942\\
0.428	-0.0239951\\
0.43	-0.0239959\\
0.432	-0.0239967\\
0.434	-0.0239975\\
0.436	-0.0239984\\
0.438	-0.0239993\\
0.44	-0.0240003\\
0.442	-0.0240015\\
0.444	-0.0240028\\
0.446	-0.0240043\\
0.448	-0.0240059\\
0.45	-0.0240078\\
0.452	-0.0240098\\
0.454	-0.0240122\\
0.456	-0.0240148\\
0.458	-0.0240176\\
0.46	-0.0240208\\
0.462	-0.0240242\\
0.464	-0.0240279\\
0.466	-0.0240319\\
0.468	-0.0240363\\
0.47	-0.0240409\\
0.472	-0.0240459\\
0.474	-0.0240512\\
0.476	-0.0240567\\
0.478	-0.0240626\\
0.48	-0.0240688\\
0.482	-0.0240752\\
0.484	-0.024082\\
0.486	-0.024089\\
0.488	-0.0240962\\
0.49	-0.0241037\\
0.492	-0.0241114\\
0.494	-0.0241194\\
0.496	-0.0241275\\
0.498	-0.0241358\\
0.5	-0.0241443\\
};
\addplot [color=mycolor6, line width=1.0pt, forget plot]
  table[row sep=crcr]{%
0	-0.0065\\
0.002	-0.0065\\
0.004	-0.0065\\
0.006	-0.0065\\
0.008	-0.0065\\
0.01	-0.00699976\\
0.012	-0.00791182\\
0.014	-0.00872527\\
0.016	-0.00945819\\
0.018	-0.0101253\\
0.02	-0.0107383\\
0.022	-0.011307\\
0.024	-0.0118388\\
0.026	-0.0123398\\
0.028	-0.0128146\\
0.03	-0.0132669\\
0.032	-0.0136993\\
0.034	-0.014114\\
0.036	-0.0145123\\
0.038	-0.0148953\\
0.04	-0.0152639\\
0.042	-0.0156184\\
0.044	-0.0159593\\
0.046	-0.0162867\\
0.048	-0.0166009\\
0.05	-0.0169018\\
0.052	-0.0171898\\
0.054	-0.0174648\\
0.056	-0.0177271\\
0.058	-0.0179769\\
0.06	-0.0182144\\
0.062	-0.01844\\
0.064	-0.0186539\\
0.066	-0.0188566\\
0.068	-0.0190484\\
0.07	-0.0192299\\
0.072	-0.0194014\\
0.074	-0.0195634\\
0.076	-0.0197165\\
0.078	-0.0198611\\
0.08	-0.0199977\\
0.082	-0.0201268\\
0.084	-0.020249\\
0.086	-0.0203646\\
0.088	-0.0204741\\
0.09	-0.020578\\
0.092	-0.0206767\\
0.094	-0.0207705\\
0.096	-0.02086\\
0.098	-0.0209453\\
0.1	-0.0210269\\
0.102	-0.0211051\\
0.104	-0.0211801\\
0.106	-0.0212522\\
0.108	-0.0213216\\
0.11	-0.0213885\\
0.112	-0.0214531\\
0.114	-0.0215155\\
0.116	-0.0215759\\
0.118	-0.0216343\\
0.12	-0.0216909\\
0.122	-0.0217456\\
0.124	-0.0217986\\
0.126	-0.0218498\\
0.128	-0.0218993\\
0.13	-0.0219469\\
0.132	-0.0219928\\
0.134	-0.0220367\\
0.136	-0.0220788\\
0.138	-0.0221189\\
0.14	-0.022157\\
0.142	-0.0221929\\
0.144	-0.0222267\\
0.146	-0.0222583\\
0.148	-0.0222876\\
0.15	-0.0223145\\
0.152	-0.0223391\\
0.154	-0.0223612\\
0.156	-0.0223808\\
0.158	-0.022398\\
0.16	-0.0224127\\
0.162	-0.0224249\\
0.164	-0.0224347\\
0.166	-0.0224421\\
0.168	-0.0224471\\
0.17	-0.0224498\\
0.172	-0.0224503\\
0.174	-0.0224487\\
0.176	-0.0224451\\
0.178	-0.0224396\\
0.18	-0.0224324\\
0.182	-0.0224235\\
0.184	-0.0224132\\
0.186	-0.0224016\\
0.188	-0.0223889\\
0.19	-0.0223752\\
0.192	-0.0223607\\
0.194	-0.0223457\\
0.196	-0.0223302\\
0.198	-0.0223144\\
0.2	-0.0222986\\
0.202	-0.0222829\\
0.204	-0.0222675\\
0.206	-0.0222525\\
0.208	-0.0222381\\
0.21	-0.0222245\\
0.212	-0.0222117\\
0.214	-0.0222\\
0.216	-0.0221894\\
0.218	-0.0221801\\
0.22	-0.0221722\\
0.222	-0.0221657\\
0.224	-0.0221607\\
0.226	-0.0221574\\
0.228	-0.0221557\\
0.23	-0.0221558\\
0.232	-0.0221575\\
0.234	-0.0221611\\
0.236	-0.0221664\\
0.238	-0.0221736\\
0.24	-0.0221825\\
0.242	-0.0221932\\
0.244	-0.0222056\\
0.246	-0.0222197\\
0.248	-0.0222354\\
0.25	-0.0222528\\
0.252	-0.0222717\\
0.254	-0.0222922\\
0.256	-0.022314\\
0.258	-0.0223372\\
0.26	-0.0223616\\
0.262	-0.0223873\\
0.264	-0.0224141\\
0.266	-0.0224418\\
0.268	-0.0224706\\
0.27	-0.0225001\\
0.272	-0.0225305\\
0.274	-0.0225615\\
0.276	-0.022593\\
0.278	-0.0226251\\
0.28	-0.0226576\\
0.282	-0.0226904\\
0.284	-0.0227235\\
0.286	-0.0227568\\
0.288	-0.0227901\\
0.29	-0.0228235\\
0.292	-0.0228568\\
0.294	-0.02289\\
0.296	-0.022923\\
0.298	-0.0229558\\
0.3	-0.0229883\\
0.302	-0.0230205\\
0.304	-0.0230523\\
0.306	-0.0230836\\
0.308	-0.0231145\\
0.31	-0.0231449\\
0.312	-0.0231747\\
0.314	-0.023204\\
0.316	-0.0232327\\
0.318	-0.0232607\\
0.32	-0.0232881\\
0.322	-0.0233149\\
0.324	-0.023341\\
0.326	-0.0233664\\
0.328	-0.0233911\\
0.33	-0.0234152\\
0.332	-0.0234385\\
0.334	-0.0234612\\
0.336	-0.0234831\\
0.338	-0.0235044\\
0.34	-0.023525\\
0.342	-0.0235449\\
0.344	-0.0235641\\
0.346	-0.0235827\\
0.348	-0.0236006\\
0.35	-0.0236179\\
0.352	-0.0236346\\
0.354	-0.0236506\\
0.356	-0.023666\\
0.358	-0.0236809\\
0.36	-0.0236952\\
0.362	-0.0237089\\
0.364	-0.023722\\
0.366	-0.0237347\\
0.368	-0.0237468\\
0.37	-0.0237585\\
0.372	-0.0237696\\
0.374	-0.0237803\\
0.376	-0.0237906\\
0.378	-0.0238005\\
0.38	-0.0238099\\
0.382	-0.023819\\
0.384	-0.0238276\\
0.386	-0.023836\\
0.388	-0.023844\\
0.39	-0.0238516\\
0.392	-0.023859\\
0.394	-0.0238661\\
0.396	-0.0238729\\
0.398	-0.0238795\\
0.4	-0.0238859\\
0.402	-0.023892\\
0.404	-0.0238979\\
0.406	-0.0239037\\
0.408	-0.0239093\\
0.41	-0.0239147\\
0.412	-0.02392\\
0.414	-0.0239252\\
0.416	-0.0239303\\
0.418	-0.0239353\\
0.42	-0.0239402\\
0.422	-0.0239451\\
0.424	-0.0239499\\
0.426	-0.0239547\\
0.428	-0.0239594\\
0.43	-0.0239641\\
0.432	-0.0239689\\
0.434	-0.0239736\\
0.436	-0.0239784\\
0.438	-0.0239832\\
0.44	-0.023988\\
0.442	-0.0239929\\
0.444	-0.0239978\\
0.446	-0.0240028\\
0.448	-0.0240079\\
0.45	-0.0240131\\
0.452	-0.0240183\\
0.454	-0.0240236\\
0.456	-0.024029\\
0.458	-0.0240345\\
0.46	-0.02404\\
0.462	-0.0240457\\
0.464	-0.0240515\\
0.466	-0.0240574\\
0.468	-0.0240633\\
0.47	-0.0240694\\
0.472	-0.0240755\\
0.474	-0.0240818\\
0.476	-0.0240881\\
0.478	-0.0240946\\
0.48	-0.0241011\\
0.482	-0.0241077\\
0.484	-0.0241144\\
0.486	-0.0241211\\
0.488	-0.0241279\\
0.49	-0.0241348\\
0.492	-0.0241417\\
0.494	-0.0241487\\
0.496	-0.0241557\\
0.498	-0.0241628\\
0.5	-0.0241698\\
};
\addplot [color=mycolor7, line width=1.0pt, forget plot]
  table[row sep=crcr]{%
0	-0.0065\\
0.002	-0.0065\\
0.004	-0.0065\\
0.006	-0.0065\\
0.008	-0.0065\\
0.01	-0.00700007\\
0.012	-0.00791235\\
0.014	-0.00872612\\
0.016	-0.00945946\\
0.018	-0.0101271\\
0.02	-0.0107409\\
0.022	-0.0113104\\
0.024	-0.0118432\\
0.026	-0.0123455\\
0.028	-0.0128217\\
0.03	-0.0132757\\
0.032	-0.01371\\
0.034	-0.0141269\\
0.036	-0.0145276\\
0.038	-0.0149134\\
0.04	-0.015285\\
0.042	-0.0156428\\
0.044	-0.0159872\\
0.046	-0.0163185\\
0.048	-0.0166366\\
0.05	-0.0169419\\
0.052	-0.0172343\\
0.054	-0.017514\\
0.056	-0.0177812\\
0.058	-0.018036\\
0.06	-0.0182785\\
0.062	-0.0185092\\
0.064	-0.0187283\\
0.066	-0.0189361\\
0.068	-0.0191329\\
0.07	-0.0193192\\
0.072	-0.0194954\\
0.074	-0.0196619\\
0.076	-0.0198191\\
0.078	-0.0199676\\
0.08	-0.0201076\\
0.082	-0.0202398\\
0.084	-0.0203644\\
0.086	-0.020482\\
0.088	-0.020593\\
0.09	-0.0206977\\
0.092	-0.0207966\\
0.094	-0.0208901\\
0.096	-0.0209784\\
0.098	-0.021062\\
0.1	-0.0211412\\
0.102	-0.0212162\\
0.104	-0.0212873\\
0.106	-0.0213549\\
0.108	-0.0214191\\
0.11	-0.0214802\\
0.112	-0.0215383\\
0.114	-0.0215937\\
0.116	-0.0216466\\
0.118	-0.021697\\
0.12	-0.0217451\\
0.122	-0.021791\\
0.124	-0.0218348\\
0.126	-0.0218766\\
0.128	-0.0219164\\
0.13	-0.0219544\\
0.132	-0.0219904\\
0.134	-0.0220247\\
0.136	-0.0220571\\
0.138	-0.0220877\\
0.14	-0.0221165\\
0.142	-0.0221435\\
0.144	-0.0221687\\
0.146	-0.022192\\
0.148	-0.0222136\\
0.15	-0.0222334\\
0.152	-0.0222514\\
0.154	-0.0222676\\
0.156	-0.0222821\\
0.158	-0.0222947\\
0.16	-0.0223057\\
0.162	-0.0223149\\
0.164	-0.0223225\\
0.166	-0.0223285\\
0.168	-0.0223329\\
0.17	-0.0223358\\
0.172	-0.0223373\\
0.174	-0.0223375\\
0.176	-0.0223365\\
0.178	-0.0223343\\
0.18	-0.022331\\
0.182	-0.0223268\\
0.184	-0.0223218\\
0.186	-0.0223161\\
0.188	-0.0223098\\
0.19	-0.022303\\
0.192	-0.0222959\\
0.194	-0.0222885\\
0.196	-0.022281\\
0.198	-0.0222736\\
0.2	-0.0222663\\
0.202	-0.0222592\\
0.204	-0.0222525\\
0.206	-0.0222463\\
0.208	-0.0222406\\
0.21	-0.0222356\\
0.212	-0.0222313\\
0.214	-0.0222279\\
0.216	-0.0222254\\
0.218	-0.0222238\\
0.22	-0.0222233\\
0.222	-0.0222239\\
0.224	-0.0222256\\
0.226	-0.0222285\\
0.228	-0.0222326\\
0.23	-0.0222379\\
0.232	-0.0222445\\
0.234	-0.0222524\\
0.236	-0.0222614\\
0.238	-0.0222718\\
0.24	-0.0222833\\
0.242	-0.0222962\\
0.244	-0.0223102\\
0.246	-0.0223253\\
0.248	-0.0223417\\
0.25	-0.0223591\\
0.252	-0.0223776\\
0.254	-0.0223972\\
0.256	-0.0224177\\
0.258	-0.0224391\\
0.26	-0.0224614\\
0.262	-0.0224846\\
0.264	-0.0225085\\
0.266	-0.0225331\\
0.268	-0.0225584\\
0.27	-0.0225843\\
0.272	-0.0226107\\
0.274	-0.0226377\\
0.276	-0.022665\\
0.278	-0.0226927\\
0.28	-0.0227207\\
0.282	-0.022749\\
0.284	-0.0227775\\
0.286	-0.0228061\\
0.288	-0.0228348\\
0.29	-0.0228636\\
0.292	-0.0228924\\
0.294	-0.0229211\\
0.296	-0.0229498\\
0.298	-0.0229783\\
0.3	-0.0230066\\
0.302	-0.0230347\\
0.304	-0.0230626\\
0.306	-0.0230902\\
0.308	-0.0231176\\
0.31	-0.0231445\\
0.312	-0.0231711\\
0.314	-0.0231973\\
0.316	-0.0232232\\
0.318	-0.0232485\\
0.32	-0.0232735\\
0.322	-0.0232979\\
0.324	-0.0233219\\
0.326	-0.0233454\\
0.328	-0.0233683\\
0.33	-0.0233908\\
0.332	-0.0234127\\
0.334	-0.0234341\\
0.336	-0.023455\\
0.338	-0.0234753\\
0.34	-0.0234951\\
0.342	-0.0235143\\
0.344	-0.023533\\
0.346	-0.0235512\\
0.348	-0.0235688\\
0.35	-0.0235859\\
0.352	-0.0236024\\
0.354	-0.0236184\\
0.356	-0.0236339\\
0.358	-0.0236489\\
0.36	-0.0236634\\
0.362	-0.0236774\\
0.364	-0.0236909\\
0.366	-0.0237039\\
0.368	-0.0237165\\
0.37	-0.0237286\\
0.372	-0.0237402\\
0.374	-0.0237514\\
0.376	-0.0237623\\
0.378	-0.0237727\\
0.38	-0.0237827\\
0.382	-0.0237923\\
0.384	-0.0238016\\
0.386	-0.0238106\\
0.388	-0.0238192\\
0.39	-0.0238275\\
0.392	-0.0238355\\
0.394	-0.0238432\\
0.396	-0.0238506\\
0.398	-0.0238578\\
0.4	-0.0238648\\
0.402	-0.0238716\\
0.404	-0.0238781\\
0.406	-0.0238845\\
0.408	-0.0238907\\
0.41	-0.0238967\\
0.412	-0.0239026\\
0.414	-0.0239083\\
0.416	-0.023914\\
0.418	-0.0239196\\
0.42	-0.0239251\\
0.422	-0.0239305\\
0.424	-0.0239359\\
0.426	-0.0239412\\
0.428	-0.0239465\\
0.43	-0.0239518\\
0.432	-0.023957\\
0.434	-0.0239623\\
0.436	-0.0239676\\
0.438	-0.023973\\
0.44	-0.0239783\\
0.442	-0.0239837\\
0.444	-0.0239892\\
0.446	-0.0239947\\
0.448	-0.0240003\\
0.45	-0.0240059\\
0.452	-0.0240117\\
0.454	-0.0240174\\
0.456	-0.0240233\\
0.458	-0.0240293\\
0.46	-0.0240353\\
0.462	-0.0240414\\
0.464	-0.0240476\\
0.466	-0.0240539\\
0.468	-0.0240603\\
0.47	-0.0240667\\
0.472	-0.0240732\\
0.474	-0.0240798\\
0.476	-0.0240864\\
0.478	-0.0240932\\
0.48	-0.0240999\\
0.482	-0.0241068\\
0.484	-0.0241137\\
0.486	-0.0241206\\
0.488	-0.0241275\\
0.49	-0.0241345\\
0.492	-0.0241415\\
0.494	-0.0241486\\
0.496	-0.0241556\\
0.498	-0.0241626\\
0.5	-0.0241697\\
};
\addplot [color=mycolor8, line width=1.0pt, forget plot]
  table[row sep=crcr]{%
0	-0.0065\\
0.002	-0.0065\\
0.004	-0.0065\\
0.006	-0.0065\\
0.008	-0.0065\\
0.01	-0.00700049\\
0.012	-0.00791307\\
0.014	-0.00872724\\
0.016	-0.00946111\\
0.018	-0.0101294\\
0.02	-0.010744\\
0.022	-0.0113146\\
0.024	-0.0118486\\
0.026	-0.0123523\\
0.028	-0.0128302\\
0.03	-0.0132859\\
0.032	-0.0137224\\
0.034	-0.0141415\\
0.036	-0.0145448\\
0.038	-0.0149333\\
0.04	-0.0153078\\
0.042	-0.0156688\\
0.044	-0.0160166\\
0.046	-0.0163514\\
0.048	-0.0166732\\
0.05	-0.0169823\\
0.052	-0.0172786\\
0.054	-0.0175623\\
0.056	-0.0178334\\
0.058	-0.0180922\\
0.06	-0.0183387\\
0.062	-0.0185733\\
0.064	-0.0187962\\
0.066	-0.0190076\\
0.068	-0.0192079\\
0.07	-0.0193975\\
0.072	-0.0195767\\
0.074	-0.019746\\
0.076	-0.0199057\\
0.078	-0.0200563\\
0.08	-0.0201982\\
0.082	-0.0203318\\
0.084	-0.0204575\\
0.086	-0.0205758\\
0.088	-0.0206871\\
0.09	-0.0207918\\
0.092	-0.0208903\\
0.094	-0.0209829\\
0.096	-0.02107\\
0.098	-0.021152\\
0.1	-0.0212293\\
0.102	-0.0213021\\
0.104	-0.0213707\\
0.106	-0.0214354\\
0.108	-0.0214966\\
0.11	-0.0215544\\
0.112	-0.0216091\\
0.114	-0.0216608\\
0.116	-0.0217099\\
0.118	-0.0217564\\
0.12	-0.0218006\\
0.122	-0.0218425\\
0.124	-0.0218824\\
0.126	-0.0219202\\
0.128	-0.0219561\\
0.13	-0.0219902\\
0.132	-0.0220225\\
0.134	-0.0220531\\
0.136	-0.0220821\\
0.138	-0.0221094\\
0.14	-0.0221351\\
0.142	-0.0221591\\
0.144	-0.0221816\\
0.146	-0.0222025\\
0.148	-0.0222218\\
0.15	-0.0222396\\
0.152	-0.0222558\\
0.154	-0.0222704\\
0.156	-0.0222835\\
0.158	-0.022295\\
0.16	-0.0223051\\
0.162	-0.0223136\\
0.164	-0.0223207\\
0.166	-0.0223264\\
0.168	-0.0223307\\
0.17	-0.0223337\\
0.172	-0.0223355\\
0.174	-0.0223361\\
0.176	-0.0223356\\
0.178	-0.0223341\\
0.18	-0.0223316\\
0.182	-0.0223283\\
0.184	-0.0223243\\
0.186	-0.0223197\\
0.188	-0.0223145\\
0.19	-0.0223089\\
0.192	-0.022303\\
0.194	-0.022297\\
0.196	-0.0222908\\
0.198	-0.0222847\\
0.2	-0.0222787\\
0.202	-0.022273\\
0.204	-0.0222676\\
0.206	-0.0222626\\
0.208	-0.0222582\\
0.21	-0.0222544\\
0.212	-0.0222514\\
0.214	-0.0222491\\
0.216	-0.0222476\\
0.218	-0.0222471\\
0.22	-0.0222476\\
0.222	-0.0222491\\
0.224	-0.0222517\\
0.226	-0.0222553\\
0.228	-0.0222601\\
0.23	-0.0222661\\
0.232	-0.0222732\\
0.234	-0.0222815\\
0.236	-0.022291\\
0.238	-0.0223016\\
0.24	-0.0223134\\
0.242	-0.0223264\\
0.244	-0.0223404\\
0.246	-0.0223556\\
0.248	-0.0223718\\
0.25	-0.0223891\\
0.252	-0.0224073\\
0.254	-0.0224265\\
0.256	-0.0224466\\
0.258	-0.0224675\\
0.26	-0.0224893\\
0.262	-0.0225118\\
0.264	-0.0225349\\
0.266	-0.0225588\\
0.268	-0.0225832\\
0.27	-0.0226082\\
0.272	-0.0226336\\
0.274	-0.0226595\\
0.276	-0.0226858\\
0.278	-0.0227124\\
0.28	-0.0227392\\
0.282	-0.0227663\\
0.284	-0.0227935\\
0.286	-0.0228209\\
0.288	-0.0228483\\
0.29	-0.0228757\\
0.292	-0.0229032\\
0.294	-0.0229305\\
0.296	-0.0229578\\
0.298	-0.022985\\
0.3	-0.0230119\\
0.302	-0.0230387\\
0.304	-0.0230652\\
0.306	-0.0230915\\
0.308	-0.0231174\\
0.31	-0.0231431\\
0.312	-0.0231684\\
0.314	-0.0231933\\
0.316	-0.0232179\\
0.318	-0.0232421\\
0.32	-0.0232658\\
0.322	-0.0232891\\
0.324	-0.023312\\
0.326	-0.0233344\\
0.328	-0.0233563\\
0.33	-0.0233777\\
0.332	-0.0233987\\
0.334	-0.0234192\\
0.336	-0.0234392\\
0.338	-0.0234587\\
0.34	-0.0234777\\
0.342	-0.0234962\\
0.344	-0.0235142\\
0.346	-0.0235317\\
0.348	-0.0235488\\
0.35	-0.0235653\\
0.352	-0.0235813\\
0.354	-0.0235969\\
0.356	-0.023612\\
0.358	-0.0236266\\
0.36	-0.0236408\\
0.362	-0.0236545\\
0.364	-0.0236678\\
0.366	-0.0236806\\
0.368	-0.023693\\
0.37	-0.023705\\
0.372	-0.0237166\\
0.374	-0.0237278\\
0.376	-0.0237386\\
0.378	-0.0237491\\
0.38	-0.0237592\\
0.382	-0.0237689\\
0.384	-0.0237783\\
0.386	-0.0237874\\
0.388	-0.0237963\\
0.39	-0.0238048\\
0.392	-0.023813\\
0.394	-0.023821\\
0.396	-0.0238288\\
0.398	-0.0238363\\
0.4	-0.0238436\\
0.402	-0.0238507\\
0.404	-0.0238576\\
0.406	-0.0238643\\
0.408	-0.0238709\\
0.41	-0.0238774\\
0.412	-0.0238837\\
0.414	-0.0238899\\
0.416	-0.023896\\
0.418	-0.023902\\
0.42	-0.0239079\\
0.422	-0.0239138\\
0.424	-0.0239196\\
0.426	-0.0239254\\
0.428	-0.0239311\\
0.43	-0.0239368\\
0.432	-0.0239425\\
0.434	-0.0239483\\
0.436	-0.023954\\
0.438	-0.0239597\\
0.44	-0.0239655\\
0.442	-0.0239713\\
0.444	-0.0239771\\
0.446	-0.023983\\
0.448	-0.0239889\\
0.45	-0.0239949\\
0.452	-0.0240009\\
0.454	-0.024007\\
0.456	-0.0240132\\
0.458	-0.0240194\\
0.46	-0.0240256\\
0.462	-0.024032\\
0.464	-0.0240384\\
0.466	-0.0240448\\
0.468	-0.0240514\\
0.47	-0.0240579\\
0.472	-0.0240646\\
0.474	-0.0240713\\
0.476	-0.024078\\
0.478	-0.0240848\\
0.48	-0.0240916\\
0.482	-0.0240984\\
0.484	-0.0241053\\
0.486	-0.0241122\\
0.488	-0.0241191\\
0.49	-0.0241261\\
0.492	-0.024133\\
0.494	-0.02414\\
0.496	-0.0241469\\
0.498	-0.0241538\\
0.5	-0.0241607\\
};
\addplot [color=mycolor9, line width=1.0pt, forget plot]
  table[row sep=crcr]{%
0	-0.0065\\
0.002	-0.0065\\
0.004	-0.0065\\
0.006	-0.0065\\
0.008	-0.0065\\
0.01	-0.00700279\\
0.012	-0.00791681\\
0.014	-0.00873287\\
0.016	-0.00946905\\
0.018	-0.0101401\\
0.02	-0.0107579\\
0.022	-0.0113321\\
0.024	-0.0118702\\
0.026	-0.0123782\\
0.028	-0.0128608\\
0.03	-0.0133216\\
0.032	-0.0137632\\
0.034	-0.0141876\\
0.036	-0.0145963\\
0.038	-0.0149903\\
0.04	-0.0153702\\
0.042	-0.0157364\\
0.044	-0.0160893\\
0.046	-0.0164289\\
0.048	-0.0167552\\
0.05	-0.0170685\\
0.052	-0.0173686\\
0.054	-0.0176556\\
0.056	-0.0179296\\
0.058	-0.0181908\\
0.06	-0.0184392\\
0.062	-0.0186752\\
0.064	-0.018899\\
0.066	-0.0191109\\
0.068	-0.0193112\\
0.07	-0.0195004\\
0.072	-0.0196788\\
0.074	-0.0198468\\
0.076	-0.020005\\
0.078	-0.0201538\\
0.08	-0.0202937\\
0.082	-0.020425\\
0.084	-0.0205484\\
0.086	-0.0206641\\
0.088	-0.0207728\\
0.09	-0.0208748\\
0.092	-0.0209706\\
0.094	-0.0210606\\
0.096	-0.0211451\\
0.098	-0.0212245\\
0.1	-0.0212992\\
0.102	-0.0213696\\
0.104	-0.0214359\\
0.106	-0.0214984\\
0.108	-0.0215574\\
0.11	-0.0216131\\
0.112	-0.0216658\\
0.114	-0.0217157\\
0.116	-0.021763\\
0.118	-0.0218078\\
0.12	-0.0218503\\
0.122	-0.0218906\\
0.124	-0.0219288\\
0.126	-0.0219651\\
0.128	-0.0219995\\
0.13	-0.0220321\\
0.132	-0.022063\\
0.134	-0.0220921\\
0.136	-0.0221196\\
0.138	-0.0221455\\
0.14	-0.0221697\\
0.142	-0.0221923\\
0.144	-0.0222134\\
0.146	-0.0222328\\
0.148	-0.0222508\\
0.15	-0.0222671\\
0.152	-0.0222819\\
0.154	-0.0222952\\
0.156	-0.022307\\
0.158	-0.0223173\\
0.16	-0.0223262\\
0.162	-0.0223336\\
0.164	-0.0223397\\
0.166	-0.0223444\\
0.168	-0.0223478\\
0.17	-0.02235\\
0.172	-0.022351\\
0.174	-0.022351\\
0.176	-0.0223499\\
0.178	-0.022348\\
0.18	-0.0223451\\
0.182	-0.0223416\\
0.184	-0.0223374\\
0.186	-0.0223326\\
0.188	-0.0223275\\
0.19	-0.0223219\\
0.192	-0.0223162\\
0.194	-0.0223103\\
0.196	-0.0223044\\
0.198	-0.0222986\\
0.2	-0.0222929\\
0.202	-0.0222876\\
0.204	-0.0222826\\
0.206	-0.0222781\\
0.208	-0.0222741\\
0.21	-0.0222708\\
0.212	-0.0222681\\
0.214	-0.0222663\\
0.216	-0.0222653\\
0.218	-0.0222652\\
0.22	-0.022266\\
0.222	-0.0222679\\
0.224	-0.0222707\\
0.226	-0.0222747\\
0.228	-0.0222797\\
0.23	-0.0222858\\
0.232	-0.0222931\\
0.234	-0.0223014\\
0.236	-0.0223109\\
0.238	-0.0223215\\
0.24	-0.0223332\\
0.242	-0.022346\\
0.244	-0.0223598\\
0.246	-0.0223747\\
0.248	-0.0223905\\
0.25	-0.0224074\\
0.252	-0.0224251\\
0.254	-0.0224438\\
0.256	-0.0224633\\
0.258	-0.0224836\\
0.26	-0.0225046\\
0.262	-0.0225263\\
0.264	-0.0225487\\
0.266	-0.0225718\\
0.268	-0.0225953\\
0.27	-0.0226194\\
0.272	-0.0226439\\
0.274	-0.0226687\\
0.276	-0.022694\\
0.278	-0.0227195\\
0.28	-0.0227453\\
0.282	-0.0227712\\
0.284	-0.0227974\\
0.286	-0.0228236\\
0.288	-0.0228499\\
0.29	-0.0228762\\
0.292	-0.0229024\\
0.294	-0.0229287\\
0.296	-0.0229548\\
0.298	-0.0229808\\
0.3	-0.0230066\\
0.302	-0.0230322\\
0.304	-0.0230576\\
0.306	-0.0230827\\
0.308	-0.0231076\\
0.31	-0.0231322\\
0.312	-0.0231564\\
0.314	-0.0231803\\
0.316	-0.0232039\\
0.318	-0.023227\\
0.32	-0.0232498\\
0.322	-0.0232722\\
0.324	-0.0232941\\
0.326	-0.0233156\\
0.328	-0.0233367\\
0.33	-0.0233574\\
0.332	-0.0233776\\
0.334	-0.0233973\\
0.336	-0.0234166\\
0.338	-0.0234355\\
0.34	-0.0234538\\
0.342	-0.0234718\\
0.344	-0.0234892\\
0.346	-0.0235062\\
0.348	-0.0235228\\
0.35	-0.0235389\\
0.352	-0.0235545\\
0.354	-0.0235697\\
0.356	-0.0235845\\
0.358	-0.0235989\\
0.36	-0.0236128\\
0.362	-0.0236263\\
0.364	-0.0236394\\
0.366	-0.0236521\\
0.368	-0.0236644\\
0.37	-0.0236763\\
0.372	-0.0236879\\
0.374	-0.023699\\
0.376	-0.0237099\\
0.378	-0.0237204\\
0.38	-0.0237306\\
0.382	-0.0237404\\
0.384	-0.02375\\
0.386	-0.0237592\\
0.388	-0.0237682\\
0.39	-0.0237769\\
0.392	-0.0237854\\
0.394	-0.0237936\\
0.396	-0.0238016\\
0.398	-0.0238094\\
0.4	-0.023817\\
0.402	-0.0238244\\
0.404	-0.0238316\\
0.406	-0.0238387\\
0.408	-0.0238456\\
0.41	-0.0238524\\
0.412	-0.023859\\
0.414	-0.0238656\\
0.416	-0.023872\\
0.418	-0.0238784\\
0.42	-0.0238847\\
0.422	-0.0238909\\
0.424	-0.0238971\\
0.426	-0.0239032\\
0.428	-0.0239093\\
0.43	-0.0239154\\
0.432	-0.0239215\\
0.434	-0.0239275\\
0.436	-0.0239336\\
0.438	-0.0239396\\
0.44	-0.0239457\\
0.442	-0.0239518\\
0.444	-0.023958\\
0.446	-0.0239641\\
0.448	-0.0239704\\
0.45	-0.0239766\\
0.452	-0.0239829\\
0.454	-0.0239892\\
0.456	-0.0239956\\
0.458	-0.024002\\
0.46	-0.0240085\\
0.462	-0.0240151\\
0.464	-0.0240216\\
0.466	-0.0240283\\
0.468	-0.0240349\\
0.47	-0.0240417\\
0.472	-0.0240484\\
0.474	-0.0240552\\
0.476	-0.0240621\\
0.478	-0.024069\\
0.48	-0.0240759\\
0.482	-0.0240828\\
0.484	-0.0240897\\
0.486	-0.0240967\\
0.488	-0.0241036\\
0.49	-0.0241106\\
0.492	-0.0241176\\
0.494	-0.0241245\\
0.496	-0.0241315\\
0.498	-0.0241384\\
0.5	-0.0241453\\
};
\addplot [color=mycolor10, line width=1.0pt, forget plot]
  table[row sep=crcr]{%
0	-0.0065\\
0.002	-0.0065\\
0.004	-0.0065\\
0.006	-0.0065\\
0.008	-0.0065\\
0.01	-0.00701266\\
0.012	-0.00793228\\
0.014	-0.00875511\\
0.016	-0.00949912\\
0.018	-0.0101788\\
0.02	-0.010806\\
0.022	-0.0113899\\
0.024	-0.0119379\\
0.026	-0.0124559\\
0.028	-0.0129482\\
0.03	-0.0134184\\
0.032	-0.0138689\\
0.034	-0.0143016\\
0.036	-0.0147179\\
0.038	-0.0151185\\
0.04	-0.0155043\\
0.042	-0.0158755\\
0.044	-0.0162324\\
0.046	-0.0165751\\
0.048	-0.0169037\\
0.05	-0.0172183\\
0.052	-0.017519\\
0.054	-0.0178058\\
0.056	-0.0180789\\
0.058	-0.0183386\\
0.06	-0.018585\\
0.062	-0.0188184\\
0.064	-0.0190392\\
0.066	-0.0192477\\
0.068	-0.0194444\\
0.07	-0.0196297\\
0.072	-0.019804\\
0.074	-0.0199679\\
0.076	-0.0201218\\
0.078	-0.0202662\\
0.08	-0.0204017\\
0.082	-0.0205288\\
0.084	-0.0206479\\
0.086	-0.0207595\\
0.088	-0.0208641\\
0.09	-0.0209621\\
0.092	-0.0210541\\
0.094	-0.0211403\\
0.096	-0.0212212\\
0.098	-0.0212972\\
0.1	-0.0213686\\
0.102	-0.0214358\\
0.104	-0.0214991\\
0.106	-0.0215587\\
0.108	-0.0216149\\
0.11	-0.021668\\
0.112	-0.0217181\\
0.114	-0.0217656\\
0.116	-0.0218105\\
0.118	-0.021853\\
0.12	-0.0218933\\
0.122	-0.0219316\\
0.124	-0.0219678\\
0.126	-0.0220021\\
0.128	-0.0220346\\
0.13	-0.0220654\\
0.132	-0.0220945\\
0.134	-0.0221219\\
0.136	-0.0221477\\
0.138	-0.0221719\\
0.14	-0.0221945\\
0.142	-0.0222156\\
0.144	-0.0222351\\
0.146	-0.0222532\\
0.148	-0.0222697\\
0.15	-0.0222847\\
0.152	-0.0222982\\
0.154	-0.0223103\\
0.156	-0.0223209\\
0.158	-0.0223301\\
0.16	-0.022338\\
0.162	-0.0223444\\
0.164	-0.0223496\\
0.166	-0.0223534\\
0.168	-0.0223561\\
0.17	-0.0223576\\
0.172	-0.0223581\\
0.174	-0.0223575\\
0.176	-0.022356\\
0.178	-0.0223537\\
0.18	-0.0223505\\
0.182	-0.0223467\\
0.184	-0.0223424\\
0.186	-0.0223375\\
0.188	-0.0223323\\
0.19	-0.0223268\\
0.192	-0.0223211\\
0.194	-0.0223153\\
0.196	-0.0223096\\
0.198	-0.022304\\
0.2	-0.0222985\\
0.202	-0.0222934\\
0.204	-0.0222887\\
0.206	-0.0222845\\
0.208	-0.0222808\\
0.21	-0.0222778\\
0.212	-0.0222755\\
0.214	-0.0222739\\
0.216	-0.0222732\\
0.218	-0.0222734\\
0.22	-0.0222745\\
0.222	-0.0222766\\
0.224	-0.0222796\\
0.226	-0.0222837\\
0.228	-0.0222889\\
0.23	-0.0222951\\
0.232	-0.0223024\\
0.234	-0.0223108\\
0.236	-0.0223202\\
0.238	-0.0223307\\
0.24	-0.0223423\\
0.242	-0.0223549\\
0.244	-0.0223685\\
0.246	-0.0223831\\
0.248	-0.0223987\\
0.25	-0.0224152\\
0.252	-0.0224325\\
0.254	-0.0224507\\
0.256	-0.0224697\\
0.258	-0.0224895\\
0.26	-0.02251\\
0.262	-0.0225311\\
0.264	-0.0225529\\
0.266	-0.0225752\\
0.268	-0.022598\\
0.27	-0.0226213\\
0.272	-0.0226451\\
0.274	-0.0226692\\
0.276	-0.0226936\\
0.278	-0.0227183\\
0.28	-0.0227432\\
0.282	-0.0227683\\
0.284	-0.0227935\\
0.286	-0.0228188\\
0.288	-0.0228442\\
0.29	-0.0228696\\
0.292	-0.022895\\
0.294	-0.0229203\\
0.296	-0.0229455\\
0.298	-0.0229706\\
0.3	-0.0229955\\
0.302	-0.0230202\\
0.304	-0.0230447\\
0.306	-0.023069\\
0.308	-0.023093\\
0.31	-0.0231167\\
0.312	-0.0231401\\
0.314	-0.0231632\\
0.316	-0.023186\\
0.318	-0.0232083\\
0.32	-0.0232303\\
0.322	-0.023252\\
0.324	-0.0232732\\
0.326	-0.0232941\\
0.328	-0.0233145\\
0.33	-0.0233345\\
0.332	-0.0233541\\
0.334	-0.0233732\\
0.336	-0.023392\\
0.338	-0.0234103\\
0.34	-0.0234281\\
0.342	-0.0234456\\
0.344	-0.0234626\\
0.346	-0.0234792\\
0.348	-0.0234953\\
0.35	-0.023511\\
0.352	-0.0235263\\
0.354	-0.0235412\\
0.356	-0.0235557\\
0.358	-0.0235698\\
0.36	-0.0235835\\
0.362	-0.0235968\\
0.364	-0.0236097\\
0.366	-0.0236222\\
0.368	-0.0236344\\
0.37	-0.0236462\\
0.372	-0.0236576\\
0.374	-0.0236688\\
0.376	-0.0236796\\
0.378	-0.02369\\
0.38	-0.0237002\\
0.382	-0.0237101\\
0.384	-0.0237197\\
0.386	-0.023729\\
0.388	-0.0237381\\
0.39	-0.0237469\\
0.392	-0.0237555\\
0.394	-0.0237638\\
0.396	-0.023772\\
0.398	-0.0237799\\
0.4	-0.0237877\\
0.402	-0.0237953\\
0.404	-0.0238027\\
0.406	-0.02381\\
0.408	-0.0238171\\
0.41	-0.0238241\\
0.412	-0.023831\\
0.414	-0.0238378\\
0.416	-0.0238445\\
0.418	-0.0238511\\
0.42	-0.0238577\\
0.422	-0.0238641\\
0.424	-0.0238706\\
0.426	-0.023877\\
0.428	-0.0238834\\
0.43	-0.0238897\\
0.432	-0.023896\\
0.434	-0.0239024\\
0.436	-0.0239087\\
0.438	-0.023915\\
0.44	-0.0239214\\
0.442	-0.0239277\\
0.444	-0.0239341\\
0.446	-0.0239405\\
0.448	-0.023947\\
0.45	-0.0239535\\
0.452	-0.02396\\
0.454	-0.0239666\\
0.456	-0.0239732\\
0.458	-0.0239798\\
0.46	-0.0239865\\
0.462	-0.0239932\\
0.464	-0.024\\
0.466	-0.0240068\\
0.468	-0.0240137\\
0.47	-0.0240205\\
0.472	-0.0240275\\
0.474	-0.0240344\\
0.476	-0.0240414\\
0.478	-0.0240484\\
0.48	-0.0240554\\
0.482	-0.0240625\\
0.484	-0.0240695\\
0.486	-0.0240766\\
0.488	-0.0240836\\
0.49	-0.0240907\\
0.492	-0.0240977\\
0.494	-0.0241047\\
0.496	-0.0241117\\
0.498	-0.0241187\\
0.5	-0.0241257\\
};
\addplot [color=mycolor11, line width=1.0pt, forget plot]
  table[row sep=crcr]{%
0	-0.0065\\
0.002	-0.0065\\
0.004	-0.0065\\
0.006	-0.0065\\
0.008	-0.0065\\
0.01	-0.00702436\\
0.012	-0.0079498\\
0.014	-0.0087792\\
0.016	-0.00953022\\
0.018	-0.0102171\\
0.02	-0.0108512\\
0.022	-0.0114418\\
0.024	-0.0119959\\
0.026	-0.0125192\\
0.028	-0.0130161\\
0.03	-0.0134899\\
0.032	-0.0139432\\
0.034	-0.0143777\\
0.036	-0.014795\\
0.038	-0.0151959\\
0.04	-0.0155811\\
0.042	-0.0159512\\
0.044	-0.0163064\\
0.046	-0.016647\\
0.048	-0.0169731\\
0.05	-0.017285\\
0.052	-0.0175828\\
0.054	-0.0178667\\
0.056	-0.0181369\\
0.058	-0.0183936\\
0.06	-0.0186372\\
0.062	-0.0188679\\
0.064	-0.019086\\
0.066	-0.0192921\\
0.068	-0.0194865\\
0.07	-0.0196697\\
0.072	-0.0198421\\
0.074	-0.0200042\\
0.076	-0.0201565\\
0.078	-0.0202995\\
0.08	-0.0204336\\
0.082	-0.0205595\\
0.084	-0.0206774\\
0.086	-0.020788\\
0.088	-0.0208917\\
0.09	-0.0209888\\
0.092	-0.0210799\\
0.094	-0.0211654\\
0.096	-0.0212456\\
0.098	-0.0213208\\
0.1	-0.0213916\\
0.102	-0.0214581\\
0.104	-0.0215207\\
0.106	-0.0215797\\
0.108	-0.0216354\\
0.11	-0.0216879\\
0.112	-0.0217375\\
0.114	-0.0217844\\
0.116	-0.0218288\\
0.118	-0.0218709\\
0.12	-0.0219108\\
0.122	-0.0219485\\
0.124	-0.0219843\\
0.126	-0.0220183\\
0.128	-0.0220504\\
0.13	-0.0220808\\
0.132	-0.0221096\\
0.134	-0.0221367\\
0.136	-0.0221622\\
0.138	-0.0221861\\
0.14	-0.0222086\\
0.142	-0.0222295\\
0.144	-0.0222488\\
0.146	-0.0222667\\
0.148	-0.0222832\\
0.15	-0.0222981\\
0.152	-0.0223116\\
0.154	-0.0223237\\
0.156	-0.0223343\\
0.158	-0.0223436\\
0.16	-0.0223515\\
0.162	-0.0223581\\
0.164	-0.0223634\\
0.166	-0.0223675\\
0.168	-0.0223703\\
0.17	-0.0223721\\
0.172	-0.0223728\\
0.174	-0.0223726\\
0.176	-0.0223714\\
0.178	-0.0223694\\
0.18	-0.0223667\\
0.182	-0.0223633\\
0.184	-0.0223593\\
0.186	-0.0223549\\
0.188	-0.0223501\\
0.19	-0.022345\\
0.192	-0.0223398\\
0.194	-0.0223345\\
0.196	-0.0223291\\
0.198	-0.022324\\
0.2	-0.022319\\
0.202	-0.0223143\\
0.204	-0.0223099\\
0.206	-0.0223061\\
0.208	-0.0223028\\
0.21	-0.0223001\\
0.212	-0.022298\\
0.214	-0.0222968\\
0.216	-0.0222963\\
0.218	-0.0222966\\
0.22	-0.0222979\\
0.222	-0.0223001\\
0.224	-0.0223033\\
0.226	-0.0223074\\
0.228	-0.0223125\\
0.23	-0.0223187\\
0.232	-0.0223259\\
0.234	-0.0223342\\
0.236	-0.0223434\\
0.238	-0.0223537\\
0.24	-0.022365\\
0.242	-0.0223773\\
0.244	-0.0223906\\
0.246	-0.0224048\\
0.248	-0.0224199\\
0.25	-0.0224359\\
0.252	-0.0224527\\
0.254	-0.0224704\\
0.256	-0.0224888\\
0.258	-0.0225079\\
0.26	-0.0225278\\
0.262	-0.0225482\\
0.264	-0.0225693\\
0.266	-0.0225908\\
0.268	-0.0226129\\
0.27	-0.0226355\\
0.272	-0.0226584\\
0.274	-0.0226817\\
0.276	-0.0227053\\
0.278	-0.0227291\\
0.28	-0.0227532\\
0.282	-0.0227774\\
0.284	-0.0228018\\
0.286	-0.0228262\\
0.288	-0.0228507\\
0.29	-0.0228753\\
0.292	-0.0228998\\
0.294	-0.0229242\\
0.296	-0.0229486\\
0.298	-0.0229728\\
0.3	-0.0229969\\
0.302	-0.0230207\\
0.304	-0.0230444\\
0.306	-0.0230679\\
0.308	-0.0230911\\
0.31	-0.023114\\
0.312	-0.0231367\\
0.314	-0.023159\\
0.316	-0.023181\\
0.318	-0.0232027\\
0.32	-0.0232241\\
0.322	-0.023245\\
0.324	-0.0232656\\
0.326	-0.0232858\\
0.328	-0.0233057\\
0.33	-0.0233251\\
0.332	-0.0233441\\
0.334	-0.0233628\\
0.336	-0.023381\\
0.338	-0.0233988\\
0.34	-0.0234162\\
0.342	-0.0234332\\
0.344	-0.0234498\\
0.346	-0.023466\\
0.348	-0.0234818\\
0.35	-0.0234972\\
0.352	-0.0235122\\
0.354	-0.0235268\\
0.356	-0.023541\\
0.358	-0.0235548\\
0.36	-0.0235683\\
0.362	-0.0235814\\
0.364	-0.0235941\\
0.366	-0.0236065\\
0.368	-0.0236185\\
0.37	-0.0236302\\
0.372	-0.0236415\\
0.374	-0.0236526\\
0.376	-0.0236633\\
0.378	-0.0236737\\
0.38	-0.0236839\\
0.382	-0.0236937\\
0.384	-0.0237033\\
0.386	-0.0237127\\
0.388	-0.0237218\\
0.39	-0.0237306\\
0.392	-0.0237392\\
0.394	-0.0237477\\
0.396	-0.0237559\\
0.398	-0.0237639\\
0.4	-0.0237718\\
0.402	-0.0237794\\
0.404	-0.023787\\
0.406	-0.0237944\\
0.408	-0.0238016\\
0.41	-0.0238088\\
0.412	-0.0238158\\
0.414	-0.0238227\\
0.416	-0.0238296\\
0.418	-0.0238363\\
0.42	-0.023843\\
0.422	-0.0238496\\
0.424	-0.0238562\\
0.426	-0.0238628\\
0.428	-0.0238693\\
0.43	-0.0238758\\
0.432	-0.0238823\\
0.434	-0.0238887\\
0.436	-0.0238952\\
0.438	-0.0239017\\
0.44	-0.0239082\\
0.442	-0.0239147\\
0.444	-0.0239212\\
0.446	-0.0239277\\
0.448	-0.0239343\\
0.45	-0.0239409\\
0.452	-0.0239475\\
0.454	-0.0239542\\
0.456	-0.0239609\\
0.458	-0.0239676\\
0.46	-0.0239744\\
0.462	-0.0239812\\
0.464	-0.0239881\\
0.466	-0.023995\\
0.468	-0.0240019\\
0.47	-0.0240088\\
0.472	-0.0240158\\
0.474	-0.0240228\\
0.476	-0.0240298\\
0.478	-0.0240369\\
0.48	-0.0240439\\
0.482	-0.024051\\
0.484	-0.0240581\\
0.486	-0.0240651\\
0.488	-0.0240722\\
0.49	-0.0240793\\
0.492	-0.0240863\\
0.494	-0.0240934\\
0.496	-0.0241004\\
0.498	-0.0241073\\
0.5	-0.0241143\\
};
\end{axis}

\begin{axis}[%
width={\figscale*5.349in},
height={\figscale*0.097in},
at={(\figscale*0.931in,\figscale*0.13in)},
scale only axis,
axis on top,
separate axis lines,
every outer x axis line/.append style={black},
every x tick label/.append style={font=\figticfont},
every x tick/.append style={black},
xmin=0,
xmax=100,
xtick={0,20,40,60,80,100},
every outer y axis line/.append style={black},
every y tick label/.append style={font=\figticfont},
every y tick/.append style={black},
y dir=reverse,
ymin=0,
ymax=1,
ytick={\empty},
axis background/.style={fill=mycolor12},
ylabel style={rotate=-90, at={(axis description cs:-0.145,0.5)}, anchor=east}
]
\addplot [forget plot] graphics [xmin=-0.196078, xmax=100.196, ymin=-0.5, ymax=1.5] {fig7_coi_envelope_3cases-1.png};

\addplot[area legend, draw=none, fill=mycolor1, forget plot]
table[row sep=crcr] {%
x	y\\
0	0\\
1.86951	0\\
1.86951	1\\
0	1\\
}--cycle;
\end{axis}
\end{tikzpicture}%
  \caption{Fig.~7 --- COI envelope over the activation sweep (full text width).}
\end{figure*}

\end{document}
